\documentclass[11pt,a4paper]{article}
\usepackage[margin=25mm,headheight=15pt]{geometry}
\usepackage[T1]{fontenc}
\usepackage[utf8]{inputenc}
\usepackage{lmodern}
\usepackage{amsmath,amssymb,amsthm,amscd,mathtools,mathrsfs}
\usepackage{microtype,booktabs,longtable,array,xcolor,graphicx,enumitem,needspace,fancyhdr}
\usepackage[round,authoryear]{natbib}
\usepackage[unicode,colorlinks=true,linkcolor=blue!40!black,citecolor=blue!40!black,urlcolor=blue!40!black]{hyperref}
\usepackage{xurl}
\usepackage{listings}
\usepackage{tabularx}
\lstset{basicstyle=\ttfamily\small,breaklines=true,columns=fullflexible,keepspaces=true}
\usepackage{tikz}
\usepackage{tikz-cd}
\usetikzlibrary{arrows.meta,calc,patterns,positioning}
\usepackage{aliascnt}
\usepackage[nameinlink,noabbrev]{cleveref}
\allowdisplaybreaks[2]
\numberwithin{equation}{section}
\makeatletter
\newcommand{\SsixEquationAliases}[1]{%
  \let\SsixEquationOriginalLabel\ltx@label
  \def\ltx@label##1{%
    \SsixEquationOriginalLabel{##1}%
    \begingroup\let\label\cref@label
    #1%
    \endgroup}}
\makeatother
\newcommand{\SsixRepositoryCommit}{a02f2eacf5a8c3db7cf8cdfcdba50c0002ee91c7}
\newcommand{\SsixPublishRecords}{https://github.com/KokunoYumeto/s6-complex-structure-workbench/blob/main/RECORDS.md\#paper-ii}
\setcounter{MaxMatrixCols}{32}
\setlength{\emergencystretch}{3em}
\setlength{\parskip}{2pt}
\newtheorem{theorem}{Theorem}[section]
\newaliascnt{proposition}{theorem}
\newtheorem{proposition}[proposition]{Proposition}
\aliascntresetthe{proposition}
\crefname{proposition}{proposition}{propositions}
\newaliascnt{lemma}{theorem}
\newtheorem{lemma}[lemma]{Lemma}
\aliascntresetthe{lemma}
\crefname{lemma}{lemma}{lemmas}
\newaliascnt{corollary}{theorem}
\newtheorem{corollary}[corollary]{Corollary}
\aliascntresetthe{corollary}
\crefname{corollary}{corollary}{corollaries}
\newaliascnt{claim}{theorem}
\newtheorem{claim}[claim]{Claim}
\aliascntresetthe{claim}
\crefname{claim}{claim}{claims}
\newaliascnt{statement}{theorem}
\newtheorem{statement}[statement]{Statement}
\aliascntresetthe{statement}
\crefname{statement}{statement}{statements}
\theoremstyle{definition}
\newaliascnt{definition}{theorem}
\newtheorem{definition}[definition]{Definition}
\aliascntresetthe{definition}
\crefname{definition}{definition}{definitions}
\newaliascnt{remark}{theorem}
\newtheorem{remark}[remark]{Remark}
\aliascntresetthe{remark}
\crefname{remark}{remark}{remarks}
\newaliascnt{example}{theorem}
\newtheorem{example}[example]{Example}
\aliascntresetthe{example}
\crefname{example}{example}{examples}
\newaliascnt{problem}{theorem}
\newtheorem{problem}[problem]{Open problem}
\aliascntresetthe{problem}
\crefname{problem}{open problem}{open problems}
\newaliascnt{question}{theorem}
\newtheorem{question}[question]{Question}
\aliascntresetthe{question}
\crefname{question}{question}{questions}
\newaliascnt{auditrecord}{theorem}
\newtheorem{auditrecord}[auditrecord]{Verification record}
\aliascntresetthe{auditrecord}
\crefname{auditrecord}{verification record}{verification records}
\newaliascnt{test}{theorem}
\newtheorem{test}[test]{Verification}
\aliascntresetthe{test}
\crefname{test}{verification}{verifications}
\newaliascnt{conjecture}{theorem}
\newtheorem{conjecture}[conjecture]{Conjecture}
\aliascntresetthe{conjecture}
\crefname{conjecture}{conjecture}{conjectures}
\newcommand{\C}{\mathbb C}\newcommand{\R}{\mathbb R}\newcommand{\Z}{\mathbb Z}\newcommand{\Q}{\mathbb Q}\newcommand{\N}{\mathbb N}\newcommand{\F}{\mathbb F}
\newcommand{\CC}{\mathbb C}\newcommand{\RR}{\mathbb R}\newcommand{\ZZ}{\mathbb Z}\newcommand{\QQ}{\mathbb Q}\newcommand{\NN}{\mathbb N}
\newcommand{\PP}{\mathbb P}\newcommand{\Pp}{\mathbb P}\newcommand{\Pone}{\mathbb P^1}\newcommand{\HH}{\mathbb H}
\newcommand{\RP}{\mathbb{RP}}\newcommand{\CP}{\mathbb{CP}}
\newcommand{\OO}{\mathcal O}\newcommand{\cO}{\mathcal O}\newcommand{\cM}{\mathcal M}\newcommand{\cF}{\mathcal F}
\newcommand{\h}{\mathfrak h}\newcommand{\hz}{\mathfrak h_z}\newcommand{\Acon}{\mathfrak A}\newcommand{\wt}[1]{\widehat{#1}}
\newcommand{\red}{\mathrm{red}}\newcommand{\tor}{\mathrm{tor}}\newcommand{\dP}{\mathrm{dP}}
\newcommand{\gh}{\widehat\gamma}\newcommand{\uh}{\widehat u}\newcommand{\wh}{\widehat w}\newcommand{\dhat}{\widehat\delta}\renewcommand{\dh}{\widehat\delta}
\newcommand{\dd}{\mathrm d}\newcommand{\ii}{\mathrm i}\newcommand{\eps}{\varepsilon}\newcommand{\partt}{\mathrm{part}}\newcommand{\Bcirc}{B^\circ}\newcommand{\ddbar}{\partial\bar\partial}
\DeclareMathOperator{\Pic}{Pic}\DeclareMathOperator{\Alb}{Alb}\DeclareMathOperator{\Aut}{Aut}\DeclareMathOperator{\NS}{NS}
\DeclareMathOperator{\Hom}{Hom}\DeclareMathOperator{\End}{End}\DeclareMathOperator{\Ext}{Ext}\DeclareMathOperator{\Tor}{Tor}
\DeclareMathOperator{\Spec}{Spec}\DeclareMathOperator{\Proj}{Proj}\DeclareMathOperator{\Sym}{Sym}\DeclareMathOperator{\Diff}{Diff}
\DeclareMathOperator{\GL}{GL}\DeclareMathOperator{\SL}{SL}\DeclareMathOperator{\PSL}{PSL}\DeclareMathOperator{\Sp}{Sp}
\DeclareMathOperator{\im}{im}\DeclareMathOperator{\coker}{coker}\DeclareMathOperator{\rank}{rank}\DeclareMathOperator{\diag}{diag}
\DeclareMathOperator{\ord}{ord}\DeclareMathOperator{\vol}{vol}\DeclareMathOperator{\pr}{pr}\DeclareMathOperator{\id}{id}
\DeclareMathOperator{\divisor}{div}\DeclareMathOperator{\Tr}{Tr}\DeclareMathOperator{\Nm}{Nm}\DeclareMathOperator{\Sha}{Sha}
\DeclareMathOperator{\Imn}{Im}\DeclareMathOperator{\Ren}{Re}\DeclareMathOperator{\trdeg}{tr.deg}
\newcommand{\Uses}[1]{\par\smallskip\noindent\textit{Uses:} #1.\par\smallskip}
\newcommand{\Verified}{}\newcommand{\AuthorSelfCheck}{}\newcommand{\NeedsCheck}{}\newcommand{\Open}{\textbf{Open.}\ }

\newcommand{\RungZero}[1]{\par\noindent\textbf{Result and mechanism.} #1\par\smallskip}
\newcommand{\RungOne}{\par\Needspace{6\baselineskip}\noindent\textbf{Proof at the first rung.}\par\nobreak\smallskip}
\newcommand{\CurrentChecker}{GPT-6 Astra (OpenAI, ChatGPT Pro mode, effort Pro), 18 November 2026}
\newaliascnt{cited}{theorem}
\newtheorem{cited}[cited]{Cited result}
\aliascntresetthe{cited}
\crefname{cited}{cited result}{cited results}
\newaliascnt{refuted}{theorem}
\newtheorem{refuted}[refuted]{Correction}
\aliascntresetthe{refuted}
\crefname{refuted}{correction}{corrections}
\newaliascnt{computed}{theorem}
\newtheorem{computed}[computed]{Computation}
\aliascntresetthe{computed}
\crefname{computed}{computation}{computations}
\newaliascnt{openq}{theorem}
\newtheorem{openq}[openq]{Open question}
\aliascntresetthe{openq}
\crefname{openq}{open question}{open questions}

\providecommand{\rk}{\operatorname{rk}}
\providecommand{\Ann}{\operatorname{Ann}}
\providecommand{\Log}{\operatorname{Log}}
\providecommand{\vir}{\mathrm{vir}}

\usepackage{etoolbox}
\hypersetup{bookmarksnumbered=true,bookmarksdepth=2}
\setcounter{tocdepth}{2}
\makeatletter\renewcommand{\@pnumwidth}{2.5em}\renewcommand{\@tocrmarg}{3.5em}\renewcommand*{\l@subsection}{\@dottedtocline{2}{1.5em}{3.5em}}\makeatother
\newtheorem{maintheorem}{Theorem}
\renewcommand{\themaintheorem}{\Alph{maintheorem}}
\crefname{maintheorem}{theorem}{theorems}
\newcommand{\StatementRecord}[4]{{\upshape\footnotesize[tag~\texttt{#3}]}\enspace\ignorespaces}
\newcommand{\ExternalPaperRef}[3]{#1#2 (Paper~#3)}
\newcommand{\CompendiumRef}[2]{\ifmmode\text{\upshape S6\_09, Paper~#1, label~\nolinkurl{#2}}\else S6\_09, Paper~#1, label~\nolinkurl{#2}\fi}
\newcommand{\SeriesRefText}[1]{\ifcsname SRef@#1\endcsname\csname SRef@#1\endcsname\else\textbf{UNRESOLVED}\fi}
\newcommand{\SeriesRef}[1]{\ifmmode\text{\SeriesRefText{#1}}\else\SeriesRefText{#1}\fi}


\newcommand{\CurrentPaper}{II}
\expandafter\def\csname SRef@wb:pshift:eq:periods\endcsname{\href{S6_SERIES_GRIND_I.pdf\#nameddest=tag-03CN}{I.{PS1} (p.~349)} }
\expandafter\def\csname SRef@wb:pshift:eq:rank-one-map\endcsname{\href{S6_SERIES_GRIND_I.pdf\#nameddest=tag-03CO}{I.{PS2} (p.~349)} }
\expandafter\def\csname SRef@wb:pshift:eq:finite-data\endcsname{\href{S6_SERIES_GRIND_I.pdf\#nameddest=tag-03C8}{I.{PS3} (p.~349)} }
\expandafter\def\csname SRef@wb:pshift:eq:covariance\endcsname{\href{S6_SERIES_GRIND_I.pdf\#nameddest=tag-03C2}{I.{PS4} (p.~349)} }
\expandafter\def\csname SRef@wb:pshift:eq:integer-shear\endcsname{\href{S6_SERIES_GRIND_I.pdf\#nameddest=tag-03CF}{I.{PS6} (p.~349)} }
\expandafter\def\csname SRef@wb:pshift:eq:centralizer\endcsname{\href{S6_SERIES_GRIND_I.pdf\#nameddest=tag-03BZ}{I.{PS7} (p.~349)} }
\expandafter\def\csname SRef@wb:pshift:eq:period-morphism\endcsname{\href{S6_SERIES_GRIND_I.pdf\#nameddest=tag-03CM}{I.{PS8} (p.~350)} }
\expandafter\def\csname SRef@wb:pshift:eq:norms\endcsname{\href{S6_SERIES_GRIND_I.pdf\#nameddest=tag-03CK}{I.{PS9} (p.~351)} }
\expandafter\def\csname SRef@wb:pshift:eq:finite-equation\endcsname{\href{S6_SERIES_GRIND_I.pdf\#nameddest=tag-03C9}{I.{PS10} (p.~351)} }
\expandafter\def\csname SRef@wb:pshift:eq:finite-norm-obstruction\endcsname{\href{S6_SERIES_GRIND_I.pdf\#nameddest=tag-03CB}{I.{PS11} (p.~351)} }
\expandafter\def\csname SRef@wb:pshift:eq:finite-class\endcsname{\href{S6_SERIES_GRIND_I.pdf\#nameddest=tag-03C7}{I.{PS12} (p.~351)} }
\expandafter\def\csname SRef@wb:pshift:eq:local-corrections\endcsname{\href{S6_SERIES_GRIND_I.pdf\#nameddest=tag-03CH}{I.{PS13} (p.~352)} }
\expandafter\def\csname SRef@wb:pshift:eq:local-matrix-identity\endcsname{\href{S6_SERIES_GRIND_I.pdf\#nameddest=tag-03CI}{I.{PS14} (p.~352)} }
\expandafter\def\csname SRef@wb:pshift:eq:finite-middle-sections\endcsname{\href{S6_SERIES_GRIND_I.pdf\#nameddest=tag-03CA}{I.{PS15} (p.~352)} }
\expandafter\def\csname SRef@wb:pshift:eq:local-affine-cocycle\endcsname{\href{S6_SERIES_GRIND_I.pdf\#nameddest=tag-03CG}{I.{PS16} (p.~352)} }
\expandafter\def\csname SRef@wb:pshift:eq:actual-overlap-difference\endcsname{\href{S6_SERIES_GRIND_I.pdf\#nameddest=tag-03BX}{I.{PS17} (p.~352)} }
\expandafter\def\csname SRef@wb:pshift:eq:cusp-periods\endcsname{\href{S6_SERIES_GRIND_I.pdf\#nameddest=tag-03C5}{I.{PS18} (p.~352)} }
\expandafter\def\csname SRef@wb:pshift:eq:cusp-actions-identical\endcsname{\href{S6_SERIES_GRIND_I.pdf\#nameddest=tag-03C3}{I.{PS19} (p.~352)} }
\expandafter\def\csname SRef@wb:pshift:eq:cocycle-law\endcsname{\href{S6_SERIES_GRIND_I.pdf\#nameddest=tag-03C0}{I.{PS20} (p.~352)} }
\expandafter\def\csname SRef@wb:pshift:eq:cusp-cocycle\endcsname{\href{S6_SERIES_GRIND_I.pdf\#nameddest=tag-03C4}{I.{PS21} (p.~352)} }
\expandafter\def\csname SRef@wb:pshift:eq:cusp-translation\endcsname{\href{S6_SERIES_GRIND_I.pdf\#nameddest=tag-03C6}{I.{PS22} (p.~353)} }
\expandafter\def\csname SRef@wb:pshift:eq:affine-bundle-action\endcsname{\href{S6_SERIES_GRIND_I.pdf\#nameddest=tag-03BY}{I.{PS23} (p.~353)} }
\expandafter\def\csname SRef@wb:pshift:eq:vector-sheaf\endcsname{\href{S6_SERIES_GRIND_I.pdf\#nameddest=tag-03CP}{I.{PS24} (p.~353)} }
\expandafter\def\csname SRef@wb:pshift:eq:no-analytic-obstruction\endcsname{\href{S6_SERIES_GRIND_I.pdf\#nameddest=tag-03CJ}{I.{PS25} (p.~354)} }
\expandafter\def\csname SRef@wb:pshift:eq:global-correction\endcsname{\href{S6_SERIES_GRIND_I.pdf\#nameddest=tag-03CE}{I.{PS26} (p.~354)} }
\expandafter\def\csname SRef@wb:pshift:eq:full-isomorphism\endcsname{\href{S6_SERIES_GRIND_I.pdf\#nameddest=tag-03CD}{I.{PS27} (p.~354)} }
\expandafter\def\csname SRef@wb:pshift:eq:finite-overlap-identity\endcsname{\href{S6_SERIES_GRIND_I.pdf\#nameddest=tag-03CC}{I.{PS28} (p.~354)} }
\expandafter\def\csname SRef@wb:pshift:eq:coherent-shift-action\endcsname{\href{S6_SERIES_GRIND_I.pdf\#nameddest=tag-03C1}{I.{PS28a} (p.~355)} }
\expandafter\def\csname SRef@wb:pshift:eq:parameter-disc\endcsname{\href{S6_SERIES_GRIND_I.pdf\#nameddest=tag-03CL}{I.{PS29} (p.~355)} }
\expandafter\def\csname SRef@mapsI:quotient-eq-2\endcsname{\href{S6_SERIES_GRIND_I.pdf\#nameddest=tag-02OZ}{I.22 (p.~296)} }
\expandafter\def\csname SRef@mapsI:quotient-eq-3\endcsname{\href{S6_SERIES_GRIND_I.pdf\#nameddest=tag-02P0}{I.{{$\langle 1\rangle $1.}} (p.~297)} }
\expandafter\def\csname SRef@mapsI:quotient-eq-4\endcsname{\href{S6_SERIES_GRIND_I.pdf\#nameddest=tag-02P1}{I.{{$\langle 1\rangle $2.}} (p.~297)} }
\expandafter\def\csname SRef@mapsI:quotient-eq-5\endcsname{\href{S6_SERIES_GRIND_I.pdf\#nameddest=tag-02P2}{I.22 (p.~297)} }
\expandafter\def\csname SRef@mapsI:quotient-eq-7\endcsname{\href{S6_SERIES_GRIND_I.pdf\#nameddest=tag-02P3}{I.{{$\langle 1\rangle $1.}} (p.~298)} }
\expandafter\def\csname SRef@app:rederive:fam-src-D.1.2\endcsname{\href{S6_SERIES_GRIND_I.pdf\#nameddest=tag-000F}{I.{{$\langle 1\rangle $1.}} (p.~77)} }
\expandafter\def\csname SRef@app:rederive:fam-src-D.1.6\endcsname{\href{S6_SERIES_GRIND_I.pdf\#nameddest=tag-000G}{I.{{$\langle 1\rangle $1.}} (p.~80)} }
\expandafter\def\csname SRef@app:rederive:fam-src-D.2.3\endcsname{\href{S6_SERIES_GRIND_I.pdf\#nameddest=tag-000H}{I.{{$\langle 1\rangle $1.}} (p.~83)} }
\expandafter\def\csname SRef@mapsI:real-eq-1\endcsname{\href{S6_SERIES_GRIND_I.pdf\#nameddest=tag-02P4}{I.12 (p.~89)} }
\expandafter\def\csname SRef@mapsI:real-eq-2\endcsname{\href{S6_SERIES_GRIND_I.pdf\#nameddest=tag-02P4}{I.12 (p.~89)} }
\expandafter\def\csname SRef@mapsI:real-eq-4\endcsname{\href{S6_SERIES_GRIND_I.pdf\#nameddest=tag-02P5}{I.{{$\langle 1\rangle $1.}} (p.~89)} }
\expandafter\def\csname SRef@mapsI:real-eq-5\endcsname{\href{S6_SERIES_GRIND_I.pdf\#nameddest=tag-02P6}{I.{{$\langle 1\rangle $4.}} (p.~90)} }
\expandafter\def\csname SRef@mapsI:real-eq-6\endcsname{\href{S6_SERIES_GRIND_I.pdf\#nameddest=tag-02P6}{I.{{$\langle 1\rangle $4.}} (p.~90)} }
\expandafter\def\csname SRef@mapsI:real-eq-7\endcsname{\href{S6_SERIES_GRIND_I.pdf\#nameddest=tag-02P7}{I.{{$\langle 1\rangle $8.}} (p.~91)} }
\expandafter\def\csname SRef@mapsI:real-eq-8\endcsname{\href{S6_SERIES_GRIND_I.pdf\#nameddest=tag-02P8}{I.{{$\langle 1\rangle $9.}} (p.~91)} }
\expandafter\def\csname SRef@ari:open:norm\endcsname{\href{S6_SERIES_GRIND_I.pdf\#nameddest=tag-003T}{I.F (p.~406)} }
\expandafter\def\csname SRef@push:src:C.3.5\endcsname{\href{S6_SERIES_GRIND_I.pdf\#nameddest=tag-0322}{I.12.5 (p.~93)} }
\expandafter\def\csname SRef@push:src:B.4.3\endcsname{\href{S6_SERIES_GRIND_I.pdf\#nameddest=tag-030Z}{I.15.16 (p.~144)} }
\expandafter\def\csname SRef@push:src:D.4.6\endcsname{\href{S6_SERIES_GRIND_I.pdf\#nameddest=tag-032O}{I.16.15 (p.~241)} }
\expandafter\def\csname SRef@inv:A1.0\endcsname{\href{S6_SERIES_GRIND_I.pdf\#nameddest=tag-01RQ}{I.18.1 (p.~254)} }
\expandafter\def\csname SRef@wb:lem:cusp-structure-sheaf-audit\endcsname{\href{S6_SERIES_GRIND_I.pdf\#nameddest=tag-01RQ}{I.18.1 (p.~254)} }
\expandafter\def\csname SRef@wb:thm:structure-direct-images-audit\endcsname{\href{S6_SERIES_GRIND_I.pdf\#nameddest=tag-01RQ}{I.18.1 (p.~254)} }
\expandafter\def\csname SRef@inv:C3.6\endcsname{\href{S6_SERIES_GRIND_I.pdf\#nameddest=tag-01SJ}{I.19.10 (p.~273)} }
\expandafter\def\csname SRef@push:src:E.3.1\endcsname{\href{S6_SERIES_PUBLISH_I.pdf\#nameddest=tag-033I}{I.12.1 (p.~77)} }
\expandafter\def\csname SRef@ari:C.5.1\endcsname{\href{S6_SERIES_GRIND_I.pdf\#nameddest=tag-000U}{I.22.9 (p.~304)} }
\expandafter\def\csname SRef@ari:C.6.R4\endcsname{\href{S6_SERIES_GRIND_I.pdf\#nameddest=tag-0013}{I.23.9 (p.~315)} }
\expandafter\def\csname SRef@ari:C.4.1\endcsname{\href{S6_SERIES_GRIND_I.pdf\#nameddest=tag-000P}{I.5.4 (p.~38)} }
\expandafter\def\csname SRef@ari:C.7.7\endcsname{\href{S6_SERIES_GRIND_I.pdf\#nameddest=tag-001G}{I.E.6 (p.~388)} }
\expandafter\def\csname SRef@ari:C.7.8\endcsname{\href{S6_SERIES_GRIND_I.pdf\#nameddest=tag-001H}{I.E.7 (p.~390)} }
\expandafter\def\csname SRef@push:src:A.4.5\endcsname{\href{S6_SERIES_GRIND_I.pdf\#nameddest=tag-02ZV}{I.14.14 (p.~120)} }
\expandafter\def\csname SRef@push:A45\endcsname{\href{S6_SERIES_GRIND_I.pdf\#nameddest=tag-02ZV}{I.14.14 (p.~120)} }
\expandafter\def\csname SRef@push:src:B.3.5\endcsname{\href{S6_SERIES_GRIND_I.pdf\#nameddest=tag-0309}{I.15.10 (p.~138)} }
\expandafter\def\csname SRef@push:B35\endcsname{\href{S6_SERIES_GRIND_I.pdf\#nameddest=tag-0309}{I.15.10 (p.~138)} }
\expandafter\def\csname SRef@push:src:B.4.1\endcsname{\href{S6_SERIES_GRIND_I.pdf\#nameddest=tag-030D}{I.15.14 (p.~143)} }
\expandafter\def\csname SRef@push:B41\endcsname{\href{S6_SERIES_GRIND_I.pdf\#nameddest=tag-030D}{I.15.14 (p.~143)} }
\expandafter\def\csname SRef@push:src:B.2.2\endcsname{\href{S6_SERIES_GRIND_I.pdf\#nameddest=tag-0303}{I.15.4 (p.~132)} }
\expandafter\def\csname SRef@push:B22\endcsname{\href{S6_SERIES_GRIND_I.pdf\#nameddest=tag-0303}{I.15.4 (p.~132)} }
\expandafter\def\csname SRef@push:src:B.5.1\endcsname{\href{S6_SERIES_GRIND_I.pdf\#nameddest=tag-031A}{I.15.47 (p.~227)} }
\expandafter\def\csname SRef@push:B51\endcsname{\href{S6_SERIES_GRIND_I.pdf\#nameddest=tag-031A}{I.15.47 (p.~227)} }
\expandafter\def\csname SRef@push:src:B.2.3\endcsname{\href{S6_SERIES_GRIND_I.pdf\#nameddest=tag-0304}{I.15.5 (p.~132)} }
\expandafter\def\csname SRef@push:B23\endcsname{\href{S6_SERIES_GRIND_I.pdf\#nameddest=tag-0304}{I.15.5 (p.~132)} }
\expandafter\def\csname SRef@push:src:D.3.2\endcsname{\href{S6_SERIES_GRIND_I.pdf\#nameddest=tag-032I}{I.16.9 (p.~235)} }
\expandafter\def\csname SRef@push:src:E.1.3\endcsname{\href{S6_SERIES_GRIND_I.pdf\#nameddest=tag-032T}{I.9.3 (p.~67)} }
\expandafter\def\csname SRef@fam:eq:cd-beta\endcsname{\href{S6_SERIES_GRIND_I.pdf\#nameddest=tag-01P5}{I.1 (p.~29)} }
\expandafter\def\csname SRef@ari:eq:ca-GHbounds\endcsname{\href{S6_SERIES_GRIND_I.pdf\#nameddest=tag-0024}{I.10 (p.~44)} }
\expandafter\def\csname SRef@eq:P5L\endcsname{\href{S6_SERIES_GRIND_I.pdf\#nameddest=tag-01OG}{I.100 (p.~367)} }
\expandafter\def\csname SRef@eq:P5floor\endcsname{\href{S6_SERIES_GRIND_I.pdf\#nameddest=tag-01OI}{I.101 (p.~367)} }
\expandafter\def\csname SRef@eq:P5normalization\endcsname{\href{S6_SERIES_GRIND_I.pdf\#nameddest=tag-01OJ}{I.102 (p.~368)} }
\expandafter\def\csname SRef@eq:P5exponential\endcsname{\href{S6_SERIES_GRIND_I.pdf\#nameddest=tag-01OH}{I.103 (p.~371)} }
\expandafter\def\csname SRef@eq:P5H2\endcsname{\href{S6_SERIES_GRIND_I.pdf\#nameddest=tag-01OF}{I.104 (p.~372)} }
\expandafter\def\csname SRef@ari:eq:cp-sigma\endcsname{\href{S6_SERIES_GRIND_I.pdf\#nameddest=tag-003K}{I.119 (p.~383)} }
\expandafter\def\csname SRef@ari:eq:cp-klein\endcsname{\href{S6_SERIES_GRIND_I.pdf\#nameddest=tag-003F}{I.120 (p.~383)} }
\expandafter\def\csname SRef@ari:eq:cp-periods\endcsname{\href{S6_SERIES_GRIND_I.pdf\#nameddest=tag-003G}{I.121 (p.~383)} }
\expandafter\def\csname SRef@ari:eq:cp-v\endcsname{\href{S6_SERIES_GRIND_I.pdf\#nameddest=tag-003L}{I.122 (p.~384)} }
\expandafter\def\csname SRef@ari:eq:cp-chi\endcsname{\href{S6_SERIES_GRIND_I.pdf\#nameddest=tag-003E}{I.123 (p.~384)} }
\expandafter\def\csname SRef@ari:eq:ca-derivatives\endcsname{\href{S6_SERIES_GRIND_I.pdf\#nameddest=tag-002B}{I.11 (p.~44)} }
\expandafter\def\csname SRef@ari:eq:cp-MTdim\endcsname{\href{S6_SERIES_GRIND_I.pdf\#nameddest=tag-003D}{I.124 (p.~385)} }
\expandafter\def\csname SRef@ari:eq:cp-positive-centre\endcsname{\href{S6_SERIES_GRIND_I.pdf\#nameddest=tag-003H}{I.125 (p.~386)} }
\expandafter\def\csname SRef@ari:eq:cp-real-linear\endcsname{\href{S6_SERIES_GRIND_I.pdf\#nameddest=tag-003I}{I.126 (p.~387)} }
\expandafter\def\csname SRef@ari:eq:cp-relation\endcsname{\href{S6_SERIES_GRIND_I.pdf\#nameddest=tag-003J}{I.127 (p.~387)} }
\expandafter\def\csname SRef@ari:eq:ct-CMmult\endcsname{\href{S6_SERIES_GRIND_I.pdf\#nameddest=tag-003Q}{I.128 (p.~388)} }
\expandafter\def\csname SRef@ari:eq:ct-Gaussian\endcsname{\href{S6_SERIES_GRIND_I.pdf\#nameddest=tag-003R}{I.129 (p.~389)} }
\expandafter\def\csname SRef@ari:eq:cd-slope\endcsname{\href{S6_SERIES_GRIND_I.pdf\#nameddest=tag-0033}{I.130 (p.~392)} }
\expandafter\def\csname SRef@ari:eq:cd-compatibility\endcsname{\href{S6_SERIES_GRIND_I.pdf\#nameddest=tag-0032}{I.131 (p.~392)} }
\expandafter\def\csname SRef@ari:eq:ce-ribet-r\endcsname{\href{S6_SERIES_GRIND_I.pdf\#nameddest=tag-0036}{I.132 (p.~395)} }
\expandafter\def\csname SRef@ari:eq:ce-zeta\endcsname{\href{S6_SERIES_GRIND_I.pdf\#nameddest=tag-0037}{I.133 (p.~395)} }
\expandafter\def\csname SRef@ari:eq:ca-delta-h\endcsname{\href{S6_SERIES_GRIND_I.pdf\#nameddest=tag-0028}{I.12 (p.~45)} }
\expandafter\def\csname SRef@ari:eq:ce-gap\endcsname{\href{S6_SERIES_GRIND_I.pdf\#nameddest=tag-0035}{I.134 (p.~396)} }
\expandafter\def\csname SRef@ari:eq:ce-explicit-witness\endcsname{\href{S6_SERIES_GRIND_I.pdf\#nameddest=tag-0034}{I.135 (p.~396)} }
\expandafter\def\csname SRef@ari:eq:cc-fixedfield\endcsname{\href{S6_SERIES_GRIND_I.pdf\#nameddest=tag-002Q}{I.136 (p.~399)} }
\expandafter\def\csname SRef@ari:eq:cc-rankcriterion\endcsname{\href{S6_SERIES_GRIND_I.pdf\#nameddest=tag-002X}{I.137 (p.~400)} }
\expandafter\def\csname SRef@ari:eq:division-ratio-model\endcsname{\href{S6_SERIES_GRIND_I.pdf\#nameddest=tag-003S}{I.138 (p.~401)} }
\expandafter\def\csname SRef@ari:eq:ca-ladder\endcsname{\href{S6_SERIES_GRIND_I.pdf\#nameddest=tag-002D}{I.13 (p.~47)} }
\expandafter\def\csname SRef@fam:eq:F-g\endcsname{\href{S6_SERIES_GRIND_I.pdf\#nameddest=tag-01OT}{I.14 (p.~69)} }
\expandafter\def\csname SRef@fam:eq:F-Gamma\endcsname{\href{S6_SERIES_GRIND_I.pdf\#nameddest=tag-01OQ}{I.15 (p.~69)} }
\expandafter\def\csname SRef@fam:eq:F-law\endcsname{\href{S6_SERIES_GRIND_I.pdf\#nameddest=tag-01OV}{I.16 (p.~69)} }
\expandafter\def\csname SRef@fam:eq:F-SLpresentation\endcsname{\href{S6_SERIES_GRIND_I.pdf\#nameddest=tag-01OS}{I.17 (p.~69)} }
\expandafter\def\csname SRef@fam:eq:F-H\endcsname{\href{S6_SERIES_GRIND_I.pdf\#nameddest=tag-01OR}{I.18 (p.~70)} }
\expandafter\def\csname SRef@fam:eq:F-qD\endcsname{\href{S6_SERIES_GRIND_I.pdf\#nameddest=tag-01OW}{I.19 (p.~73)} }
\expandafter\def\csname SRef@fam:eq:cd-wp\endcsname{\href{S6_SERIES_GRIND_I.pdf\#nameddest=tag-01P8}{I.2 (p.~30)} }
\expandafter\def\csname SRef@fam:eq:F-tauaction\endcsname{\href{S6_SERIES_GRIND_I.pdf\#nameddest=tag-01OX}{I.20 (p.~74)} }
\expandafter\def\csname SRef@fam:eq:F-j\endcsname{\href{S6_SERIES_GRIND_I.pdf\#nameddest=tag-01OU}{I.21 (p.~74)} }
\expandafter\def\csname SRef@fam:eq:cd-NS\endcsname{\href{S6_SERIES_GRIND_I.pdf\#nameddest=tag-01P4}{I.22 (p.~79)} }
\expandafter\def\csname SRef@fam:eq:cd-theta-sections\endcsname{\href{S6_SERIES_GRIND_I.pdf\#nameddest=tag-01P7}{I.23 (p.~84)} }
\expandafter\def\csname SRef@mapsI:closing:eq:C-explicit-closing\endcsname{\href{S6_SERIES_GRIND_I.pdf\#nameddest=tag-02OP}{I.24 (p.~93)} }
\expandafter\def\csname SRef@mapsI:closing:eq:C-twelve-scalar\endcsname{\href{S6_SERIES_GRIND_I.pdf\#nameddest=tag-02OQ}{I.25 (p.~93)} }
\expandafter\def\csname SRef@eq:d-period-normal\endcsname{\href{S6_SERIES_GRIND_I.pdf\#nameddest=tag-01ON}{I.26 (p.~230)} }
\expandafter\def\csname SRef@eq:d-limit-group\endcsname{\href{S6_SERIES_GRIND_I.pdf\#nameddest=tag-01OM}{I.27 (p.~233)} }
\expandafter\def\csname SRef@eq:d-dual-exact\endcsname{\href{S6_SERIES_GRIND_I.pdf\#nameddest=tag-01OL}{I.28 (p.~234)} }
\expandafter\def\csname SRef@eq:P2H\endcsname{\href{S6_SERIES_GRIND_I.pdf\#nameddest=tag-01O8}{I.29 (p.~243)} }
\expandafter\def\csname SRef@fam:eq:cd-norm\endcsname{\href{S6_SERIES_GRIND_I.pdf\#nameddest=tag-01P6}{I.3 (p.~31)} }
\expandafter\def\csname SRef@eq:P2zg\endcsname{\href{S6_SERIES_GRIND_I.pdf\#nameddest=tag-01OC}{I.30 (p.~243)} }
\expandafter\def\csname SRef@eq:P2F\endcsname{\href{S6_SERIES_GRIND_I.pdf\#nameddest=tag-01O5}{I.31 (p.~244)} }
\expandafter\def\csname SRef@eq:P2FQ\endcsname{\href{S6_SERIES_GRIND_I.pdf\#nameddest=tag-01O6}{I.32 (p.~244)} }
\expandafter\def\csname SRef@eq:P2FR\endcsname{\href{S6_SERIES_GRIND_I.pdf\#nameddest=tag-01O7}{I.33 (p.~244)} }
\expandafter\def\csname SRef@eq:P2chain\endcsname{\href{S6_SERIES_GRIND_I.pdf\#nameddest=tag-01O9}{I.34 (p.~245)} }
\expandafter\def\csname SRef@eq:P2split\endcsname{\href{S6_SERIES_GRIND_I.pdf\#nameddest=tag-01OB}{I.35 (p.~247)} }
\expandafter\def\csname SRef@eq:P2ratio\endcsname{\href{S6_SERIES_GRIND_I.pdf\#nameddest=tag-01OA}{I.36 (p.~248)} }
\expandafter\def\csname SRef@inv:eq:A-E-line\endcsname{\href{S6_SERIES_GRIND_I.pdf\#nameddest=tag-01SZ}{I.37 (p.~254)} }
\expandafter\def\csname SRef@inv:eq:A-H-images\endcsname{\href{S6_SERIES_GRIND_I.pdf\#nameddest=tag-01T0}{I.38 (p.~255)} }
\expandafter\def\csname SRef@inv:eq:A-N-images\endcsname{\href{S6_SERIES_GRIND_I.pdf\#nameddest=tag-01T1}{I.39 (p.~255)} }
\expandafter\def\csname SRef@ari:eq:cs-isogeny\endcsname{\href{S6_SERIES_GRIND_I.pdf\#nameddest=tag-003O}{I.4 (p.~36)} }
\expandafter\def\csname SRef@inv:eq:A-absolute-table\endcsname{\href{S6_SERIES_GRIND_I.pdf\#nameddest=tag-01T2}{I.40 (p.~256)} }
\expandafter\def\csname SRef@inv:eq:A-s-branch\endcsname{\href{S6_SERIES_GRIND_I.pdf\#nameddest=tag-01T5}{I.41 (p.~257)} }
\expandafter\def\csname SRef@inv:eq:A-modular-K\endcsname{\href{S6_SERIES_GRIND_I.pdf\#nameddest=tag-01T3}{I.42 (p.~257)} }
\expandafter\def\csname SRef@inv:eq:A-mu-value\endcsname{\href{S6_SERIES_GRIND_I.pdf\#nameddest=tag-01T4}{I.43 (p.~258)} }
\expandafter\def\csname SRef@inv:eq:A-DI-branch\endcsname{\href{S6_SERIES_GRIND_I.pdf\#nameddest=tag-01SY}{I.44 (p.~260)} }
\expandafter\def\csname SRef@inv:eq:C-hodge-input\endcsname{\href{S6_SERIES_GRIND_I.pdf\#nameddest=tag-01TF}{I.45 (p.~266)} }
\expandafter\def\csname SRef@inv:eq:C-BC-A\endcsname{\href{S6_SERIES_GRIND_I.pdf\#nameddest=tag-01T6}{I.46 (p.~267)} }
\expandafter\def\csname SRef@inv:eq:C-arrows\endcsname{\href{S6_SERIES_GRIND_I.pdf\#nameddest=tag-01TA}{I.47 (p.~268)} }
\expandafter\def\csname SRef@inv:eq:C-Pic-map\endcsname{\href{S6_SERIES_GRIND_I.pdf\#nameddest=tag-01T8}{I.48 (p.~269)} }
\expandafter\def\csname SRef@inv:eq:C-G-Green\endcsname{\href{S6_SERIES_GRIND_I.pdf\#nameddest=tag-01T7}{I.49 (p.~269)} }
\expandafter\def\csname SRef@ari:eq:cs-division\endcsname{\href{S6_SERIES_GRIND_I.pdf\#nameddest=tag-003N}{I.5 (p.~39)} }
\expandafter\def\csname SRef@inv:eq:C-factor12\endcsname{\href{S6_SERIES_GRIND_I.pdf\#nameddest=tag-01TE}{I.50 (p.~270)} }
\expandafter\def\csname SRef@inv:eq:C-collar-sum\endcsname{\href{S6_SERIES_GRIND_I.pdf\#nameddest=tag-01TB}{I.51 (p.~271)} }
\expandafter\def\csname SRef@inv:eq:C-delta\endcsname{\href{S6_SERIES_GRIND_I.pdf\#nameddest=tag-01TC}{I.52 (p.~273)} }
\expandafter\def\csname SRef@inv:eq:C-Ricci\endcsname{\href{S6_SERIES_GRIND_I.pdf\#nameddest=tag-01T9}{I.53 (p.~274)} }
\expandafter\def\csname SRef@inv:eq:C-eta-negative\endcsname{\href{S6_SERIES_GRIND_I.pdf\#nameddest=tag-01TD}{I.54 (p.~274)} }
\expandafter\def\csname SRef@ari:eq:cs-Li\endcsname{\href{S6_SERIES_GRIND_I.pdf\#nameddest=tag-003M}{I.55 (p.~301)} }
\expandafter\def\csname SRef@ari:eq:ch-global\endcsname{\href{S6_SERIES_GRIND_I.pdf\#nameddest=tag-0039}{I.56 (p.~304)} }
\expandafter\def\csname SRef@ari:eq:ch-local-mult\endcsname{\href{S6_SERIES_GRIND_I.pdf\#nameddest=tag-003B}{I.57 (p.~304)} }
\expandafter\def\csname SRef@ari:eq:ch-local-exact\endcsname{\href{S6_SERIES_GRIND_I.pdf\#nameddest=tag-003A}{I.58 (p.~305)} }
\expandafter\def\csname SRef@ari:eq:ch-bound\endcsname{\href{S6_SERIES_GRIND_I.pdf\#nameddest=tag-0038}{I.59 (p.~308)} }
\expandafter\def\csname SRef@ari:eq:cs-torsionvalue\endcsname{\href{S6_SERIES_GRIND_I.pdf\#nameddest=tag-003P}{I.6 (p.~39)} }
\expandafter\def\csname SRef@ari:eq:ch-wp-bracket\endcsname{\href{S6_SERIES_GRIND_I.pdf\#nameddest=tag-003C}{I.60 (p.~309)} }
\expandafter\def\csname SRef@ari:eq:Cr-first\endcsname{\href{S6_SERIES_GRIND_I.pdf\#nameddest=tag-0020}{I.61 (p.~312)} }
\expandafter\def\csname SRef@ari:eq:Cr-second\endcsname{\href{S6_SERIES_GRIND_I.pdf\#nameddest=tag-0022}{I.62 (p.~312)} }
\expandafter\def\csname SRef@ari:eq:Cr-tails\endcsname{\href{S6_SERIES_GRIND_I.pdf\#nameddest=tag-0023}{I.63 (p.~312)} }
\expandafter\def\csname SRef@ari:eq:Cr-rational-factor\endcsname{\href{S6_SERIES_GRIND_I.pdf\#nameddest=tag-0021}{I.64 (p.~314)} }
\expandafter\def\csname SRef@ari:eq:Cgz-central-integrality\endcsname{\href{S6_SERIES_GRIND_I.pdf\#nameddest=tag-001X}{I.65 (p.~318)} }
\expandafter\def\csname SRef@ari:eq:Cgz-product\endcsname{\href{S6_SERIES_GRIND_I.pdf\#nameddest=tag-001Z}{I.66 (p.~319)} }
\expandafter\def\csname SRef@ari:eq:Cgz-denominator\endcsname{\href{S6_SERIES_GRIND_I.pdf\#nameddest=tag-001Y}{I.67 (p.~319)} }
\expandafter\def\csname SRef@ari:eq:ca-mahler\endcsname{\href{S6_SERIES_GRIND_I.pdf\#nameddest=tag-002E}{I.68 (p.~325)} }
\expandafter\def\csname SRef@ari:eq:ca-wpmean\endcsname{\href{S6_SERIES_GRIND_I.pdf\#nameddest=tag-002H}{I.69 (p.~325)} }
\expandafter\def\csname SRef@ari:eq:ca-sup\endcsname{\href{S6_SERIES_GRIND_I.pdf\#nameddest=tag-002G}{I.7 (p.~43)} }
\expandafter\def\csname SRef@ari:eq:ca-nu\endcsname{\href{S6_SERIES_GRIND_I.pdf\#nameddest=tag-002F}{I.70 (p.~325)} }
\expandafter\def\csname SRef@ari:eq:ca-current\endcsname{\href{S6_SERIES_GRIND_I.pdf\#nameddest=tag-0026}{I.71 (p.~327)} }
\expandafter\def\csname SRef@ari:eq:ca-density0\endcsname{\href{S6_SERIES_GRIND_I.pdf\#nameddest=tag-0029}{I.72 (p.~328)} }
\expandafter\def\csname SRef@ari:eq:ca-density1\endcsname{\href{S6_SERIES_GRIND_I.pdf\#nameddest=tag-002A}{I.73 (p.~328)} }
\expandafter\def\csname SRef@ari:eq:ca-cuspexp\endcsname{\href{S6_SERIES_GRIND_I.pdf\#nameddest=tag-0027}{I.74 (p.~328)} }
\expandafter\def\csname SRef@ari:eq:cc-EZ\endcsname{\href{S6_SERIES_GRIND_I.pdf\#nameddest=tag-002I}{I.75 (p.~330)} }
\expandafter\def\csname SRef@ari:eq:cc-radial\endcsname{\href{S6_SERIES_GRIND_I.pdf\#nameddest=tag-002W}{I.76 (p.~331)} }
\expandafter\def\csname SRef@ari:eq:cc-FR\endcsname{\href{S6_SERIES_GRIND_I.pdf\#nameddest=tag-002J}{I.77 (p.~331)} }
\expandafter\def\csname SRef@ari:eq:cc-cuspbound\endcsname{\href{S6_SERIES_GRIND_I.pdf\#nameddest=tag-002O}{I.78 (p.~332)} }
\expandafter\def\csname SRef@ari:eq:cc-ode\endcsname{\href{S6_SERIES_GRIND_I.pdf\#nameddest=tag-002T}{I.79 (p.~333)} }
\expandafter\def\csname SRef@ari:eq:ca-exp\endcsname{\href{S6_SERIES_GRIND_I.pdf\#nameddest=tag-002C}{I.8 (p.~43)} }
\expandafter\def\csname SRef@ari:eq:cc-h\endcsname{\href{S6_SERIES_GRIND_I.pdf\#nameddest=tag-002R}{I.80 (p.~333)} }
\expandafter\def\csname SRef@ari:eq:cc-Fseries\endcsname{\href{S6_SERIES_GRIND_I.pdf\#nameddest=tag-002L}{I.81 (p.~333)} }
\expandafter\def\csname SRef@ari:eq:cc-Ffactored\endcsname{\href{S6_SERIES_GRIND_I.pdf\#nameddest=tag-002K}{I.82 (p.~334)} }
\expandafter\def\csname SRef@ari:eq:cc-norms\endcsname{\href{S6_SERIES_GRIND_I.pdf\#nameddest=tag-002S}{I.83 (p.~334)} }
\expandafter\def\csname SRef@ari:eq:cc-recurrence\endcsname{\href{S6_SERIES_GRIND_I.pdf\#nameddest=tag-002Y}{I.84 (p.~334)} }
\expandafter\def\csname SRef@ari:eq:cc-power\endcsname{\href{S6_SERIES_GRIND_I.pdf\#nameddest=tag-002V}{I.85 (p.~334)} }
\expandafter\def\csname SRef@ari:eq:cc-constants\endcsname{\href{S6_SERIES_GRIND_I.pdf\#nameddest=tag-002N}{I.86 (p.~334)} }
\expandafter\def\csname SRef@ari:eq:cc-outer\endcsname{\href{S6_SERIES_GRIND_I.pdf\#nameddest=tag-002U}{I.87 (p.~335)} }
\expandafter\def\csname SRef@ari:eq:cc-tail60\endcsname{\href{S6_SERIES_GRIND_I.pdf\#nameddest=tag-0031}{I.88 (p.~335)} }
\expandafter\def\csname SRef@ari:eq:cc-tail100\endcsname{\href{S6_SERIES_GRIND_I.pdf\#nameddest=tag-0030}{I.89 (p.~335)} }
\expandafter\def\csname SRef@ari:eq:ca-H\endcsname{\href{S6_SERIES_GRIND_I.pdf\#nameddest=tag-0025}{I.9 (p.~44)} }
\expandafter\def\csname SRef@ari:eq:cc-finite-table\endcsname{\href{S6_SERIES_GRIND_I.pdf\#nameddest=tag-002P}{I.90 (p.~336)} }
\expandafter\def\csname SRef@ari:eq:cc-scalar\endcsname{\href{S6_SERIES_GRIND_I.pdf\#nameddest=tag-002Z}{I.91 (p.~336)} }
\expandafter\def\csname SRef@ari:eq:cc-P\endcsname{\href{S6_SERIES_GRIND_I.pdf\#nameddest=tag-002M}{I.92 (p.~336)} }
\expandafter\def\csname SRef@fam:eq:ab-root-degree\endcsname{\href{S6_SERIES_GRIND_I.pdf\#nameddest=tag-01P3}{I.93 (p.~356)} }
\expandafter\def\csname SRef@fam:eq:ab-cusp-resolution\endcsname{\href{S6_SERIES_GRIND_I.pdf\#nameddest=tag-01OY}{I.94 (p.~357)} }
\expandafter\def\csname SRef@fam:eq:ab-log-tangent\endcsname{\href{S6_SERIES_GRIND_I.pdf\#nameddest=tag-01P0}{I.95 (p.~358)} }
\expandafter\def\csname SRef@fam:eq:ab-relative-dualizing\endcsname{\href{S6_SERIES_GRIND_I.pdf\#nameddest=tag-01P2}{I.96 (p.~358)} }
\expandafter\def\csname SRef@fam:eq:ab-direct-images\endcsname{\href{S6_SERIES_GRIND_I.pdf\#nameddest=tag-01OZ}{I.97 (p.~358)} }
\expandafter\def\csname SRef@fam:eq:ab-normal-quotient\endcsname{\href{S6_SERIES_GRIND_I.pdf\#nameddest=tag-01P1}{I.98 (p.~359)} }
\expandafter\def\csname SRef@eq:P5tangent\endcsname{\href{S6_SERIES_GRIND_I.pdf\#nameddest=tag-01OK}{I.99 (p.~366)} }
\expandafter\def\csname SRef@fig:I:orbifold\endcsname{\href{S6_SERIES_GRIND_I.pdf\#nameddest=tag-01RL}{I.2 (p.~94)} }
\expandafter\def\csname SRef@fig:I:parameter-disc\endcsname{\href{S6_SERIES_GRIND_I.pdf\#nameddest=tag-01RM}{I.3 (p.~94)} }
\expandafter\def\csname SRef@fig:I:monodromy-domain\endcsname{\href{S6_SERIES_GRIND_I.pdf\#nameddest=tag-01RK}{I.4 (p.~95)} }
\expandafter\def\csname SRef@push:src:A.3.1\endcsname{\href{S6_SERIES_GRIND_I.pdf\#nameddest=tag-02ZO}{I.14.7 (p.~113)} }
\expandafter\def\csname SRef@push:A31\endcsname{\href{S6_SERIES_GRIND_I.pdf\#nameddest=tag-02ZO}{I.14.7 (p.~113)} }
\expandafter\def\csname SRef@push:src:D.1.2\endcsname{\href{S6_SERIES_GRIND_I.pdf\#nameddest=tag-032B}{I.16.2 (p.~228)} }
\expandafter\def\csname SRef@push:src:D.2.1\endcsname{\href{S6_SERIES_GRIND_I.pdf\#nameddest=tag-032D}{I.16.4 (p.~230)} }
\expandafter\def\csname SRef@I:weierstrass-comparison\endcsname{\href{S6_SERIES_GRIND_I.pdf\#nameddest=tag-0002}{I.3.1 (p.~24)} }
\expandafter\def\csname SRef@wb:pshift:lem:centralizer\endcsname{\href{S6_SERIES_GRIND_I.pdf\#nameddest=tag-03CQ}{I.A.4 (p.~349)} }
\expandafter\def\csname SRef@fam:src:F.1.2\endcsname{\href{S6_SERIES_GRIND_I.pdf\#nameddest=tag-01RA}{I.10.2 (p.~70)} }
\expandafter\def\csname SRef@fam:st:F.1.2\endcsname{\href{S6_SERIES_GRIND_I.pdf\#nameddest=tag-01RA}{I.10.2 (p.~70)} }
\expandafter\def\csname SRef@fam:src:F.2.1\endcsname{\href{S6_SERIES_GRIND_I.pdf\#nameddest=tag-01RB}{I.10.3 (p.~71)} }
\expandafter\def\csname SRef@fam:st:F.2.1\endcsname{\href{S6_SERIES_GRIND_I.pdf\#nameddest=tag-01RB}{I.10.3 (p.~71)} }
\expandafter\def\csname SRef@fam:src:F.2.3\endcsname{\href{S6_SERIES_GRIND_I.pdf\#nameddest=tag-01RD}{I.10.5 (p.~73)} }
\expandafter\def\csname SRef@fam:st:F.2.3\endcsname{\href{S6_SERIES_GRIND_I.pdf\#nameddest=tag-01RD}{I.10.5 (p.~73)} }
\expandafter\def\csname SRef@fam:src:F.3.2\endcsname{\href{S6_SERIES_GRIND_I.pdf\#nameddest=tag-01RF}{I.10.7 (p.~75)} }
\expandafter\def\csname SRef@fam:st:F.3.2\endcsname{\href{S6_SERIES_GRIND_I.pdf\#nameddest=tag-01RF}{I.10.7 (p.~75)} }
\expandafter\def\csname SRef@fam:src:F.3.3\endcsname{\href{S6_SERIES_GRIND_I.pdf\#nameddest=tag-01RG}{I.10.8 (p.~75)} }
\expandafter\def\csname SRef@fam:st:F.3.3\endcsname{\href{S6_SERIES_GRIND_I.pdf\#nameddest=tag-01RG}{I.10.8 (p.~75)} }
\expandafter\def\csname SRef@fam:src:D.1.1\endcsname{\href{S6_SERIES_GRIND_I.pdf\#nameddest=tag-01QR}{I.11.1 (p.~76)} }
\expandafter\def\csname SRef@fam:st:D.1.1\endcsname{\href{S6_SERIES_GRIND_I.pdf\#nameddest=tag-01QR}{I.11.1 (p.~76)} }
\expandafter\def\csname SRef@fam:src:C.3.1\endcsname{\href{S6_SERIES_GRIND_I.pdf\#nameddest=tag-01QM}{I.11.11 (p.~84)} }
\expandafter\def\csname SRef@fam:st:C.3.1\endcsname{\href{S6_SERIES_GRIND_I.pdf\#nameddest=tag-01QM}{I.11.11 (p.~84)} }
\expandafter\def\csname SRef@fam:src:C.3.2\endcsname{\href{S6_SERIES_GRIND_I.pdf\#nameddest=tag-01QN}{I.11.12 (p.~85)} }
\expandafter\def\csname SRef@fam:st:C.3.2\endcsname{\href{S6_SERIES_GRIND_I.pdf\#nameddest=tag-01QN}{I.11.12 (p.~85)} }
\expandafter\def\csname SRef@fam:src:C.3.4\endcsname{\href{S6_SERIES_GRIND_I.pdf\#nameddest=tag-01QP}{I.11.13 (p.~86)} }
\expandafter\def\csname SRef@fam:st:C.3.4\endcsname{\href{S6_SERIES_GRIND_I.pdf\#nameddest=tag-01QP}{I.11.13 (p.~86)} }
\expandafter\def\csname SRef@fam:src:D.1.2\endcsname{\href{S6_SERIES_GRIND_I.pdf\#nameddest=tag-01QS}{I.11.2 (p.~77)} }
\expandafter\def\csname SRef@fam:st:D.1.2\endcsname{\href{S6_SERIES_GRIND_I.pdf\#nameddest=tag-01QS}{I.11.2 (p.~77)} }
\expandafter\def\csname SRef@fam:src:D.1.4\endcsname{\href{S6_SERIES_GRIND_I.pdf\#nameddest=tag-01QU}{I.11.4 (p.~79)} }
\expandafter\def\csname SRef@fam:st:D.1.4\endcsname{\href{S6_SERIES_GRIND_I.pdf\#nameddest=tag-01QU}{I.11.4 (p.~79)} }
\expandafter\def\csname SRef@fam:src:D.1.5\endcsname{\href{S6_SERIES_GRIND_I.pdf\#nameddest=tag-01QV}{I.11.5 (p.~79)} }
\expandafter\def\csname SRef@fam:st:D.1.5\endcsname{\href{S6_SERIES_GRIND_I.pdf\#nameddest=tag-01QV}{I.11.5 (p.~79)} }
\expandafter\def\csname SRef@fam:src:D.1.7\endcsname{\href{S6_SERIES_GRIND_I.pdf\#nameddest=tag-01QX}{I.11.7 (p.~81)} }
\expandafter\def\csname SRef@fam:st:D.1.7\endcsname{\href{S6_SERIES_GRIND_I.pdf\#nameddest=tag-01QX}{I.11.7 (p.~81)} }
\expandafter\def\csname SRef@fam:src:D.2.1\endcsname{\href{S6_SERIES_GRIND_I.pdf\#nameddest=tag-01QY}{I.11.8 (p.~82)} }
\expandafter\def\csname SRef@fam:st:D.2.1\endcsname{\href{S6_SERIES_GRIND_I.pdf\#nameddest=tag-01QY}{I.11.8 (p.~82)} }
\expandafter\def\csname SRef@fam:src:D.2.2\endcsname{\href{S6_SERIES_GRIND_I.pdf\#nameddest=tag-01QZ}{I.11.9 (p.~83)} }
\expandafter\def\csname SRef@fam:st:D.2.2\endcsname{\href{S6_SERIES_GRIND_I.pdf\#nameddest=tag-01QZ}{I.11.9 (p.~83)} }
\expandafter\def\csname SRef@mapsI:src:A.3.1\endcsname{\href{S6_SERIES_GRIND_I.pdf\#nameddest=tag-02PE}{I.12.1 (p.~88)} }
\expandafter\def\csname SRef@mapsI:src:A.3.4\endcsname{\href{S6_SERIES_GRIND_I.pdf\#nameddest=tag-02PH}{I.12.4 (p.~93)} }
\expandafter\def\csname SRef@mapsI:real-eq-9\endcsname{\href{S6_SERIES_GRIND_I.pdf\#nameddest=tag-02PH}{I.12.4 (p.~93)} }
\expandafter\def\csname SRef@I:monodromy-fundamental-set\endcsname{\href{S6_SERIES_GRIND_I.pdf\#nameddest=tag-0001}{I.12.6 (p.~94)} }
\expandafter\def\csname SRef@fam:src:E.1.1\endcsname{\href{S6_SERIES_GRIND_I.pdf\#nameddest=tag-01R1}{I.13.1 (p.~96)} }
\expandafter\def\csname SRef@fam:st:E.1.1\endcsname{\href{S6_SERIES_GRIND_I.pdf\#nameddest=tag-01R1}{I.13.1 (p.~96)} }
\expandafter\def\csname SRef@fam:src:E.1.2\endcsname{\href{S6_SERIES_GRIND_I.pdf\#nameddest=tag-01R2}{I.13.2 (p.~97)} }
\expandafter\def\csname SRef@fam:st:E.1.2\endcsname{\href{S6_SERIES_GRIND_I.pdf\#nameddest=tag-01R2}{I.13.2 (p.~97)} }
\expandafter\def\csname SRef@fam:src:E.1.3\endcsname{\href{S6_SERIES_GRIND_I.pdf\#nameddest=tag-01R3}{I.13.3 (p.~98)} }
\expandafter\def\csname SRef@fam:st:E.1.3\endcsname{\href{S6_SERIES_GRIND_I.pdf\#nameddest=tag-01R3}{I.13.3 (p.~98)} }
\expandafter\def\csname SRef@fam:src:E.1.4\endcsname{\href{S6_SERIES_GRIND_I.pdf\#nameddest=tag-01R4}{I.13.4 (p.~99)} }
\expandafter\def\csname SRef@fam:st:E.1.4\endcsname{\href{S6_SERIES_GRIND_I.pdf\#nameddest=tag-01R4}{I.13.4 (p.~99)} }
\expandafter\def\csname SRef@fam:src:E.2.1\endcsname{\href{S6_SERIES_GRIND_I.pdf\#nameddest=tag-01R6}{I.13.5 (p.~101)} }
\expandafter\def\csname SRef@fam:st:E.2.1\endcsname{\href{S6_SERIES_GRIND_I.pdf\#nameddest=tag-01R6}{I.13.5 (p.~101)} }
\expandafter\def\csname SRef@push:src:A.1.1\endcsname{\href{S6_SERIES_GRIND_I.pdf\#nameddest=tag-02ZH}{I.14.1 (p.~103)} }
\expandafter\def\csname SRef@push:A11\endcsname{\href{S6_SERIES_GRIND_I.pdf\#nameddest=tag-02ZH}{I.14.1 (p.~103)} }
\expandafter\def\csname SRef@push:src:A.4.1\endcsname{\href{S6_SERIES_GRIND_I.pdf\#nameddest=tag-02ZR}{I.14.10 (p.~116)} }
\expandafter\def\csname SRef@push:A41\endcsname{\href{S6_SERIES_GRIND_I.pdf\#nameddest=tag-02ZR}{I.14.10 (p.~116)} }
\expandafter\def\csname SRef@push:src:A.4.2\endcsname{\href{S6_SERIES_GRIND_I.pdf\#nameddest=tag-02ZS}{I.14.11 (p.~118)} }
\expandafter\def\csname SRef@push:A42\endcsname{\href{S6_SERIES_GRIND_I.pdf\#nameddest=tag-02ZS}{I.14.11 (p.~118)} }
\expandafter\def\csname SRef@push:src:A.4.3\endcsname{\href{S6_SERIES_GRIND_I.pdf\#nameddest=tag-02ZT}{I.14.12 (p.~119)} }
\expandafter\def\csname SRef@push:A43\endcsname{\href{S6_SERIES_GRIND_I.pdf\#nameddest=tag-02ZT}{I.14.12 (p.~119)} }
\expandafter\def\csname SRef@push:src:A.4.4\endcsname{\href{S6_SERIES_GRIND_I.pdf\#nameddest=tag-02ZU}{I.14.13 (p.~120)} }
\expandafter\def\csname SRef@push:A44\endcsname{\href{S6_SERIES_GRIND_I.pdf\#nameddest=tag-02ZU}{I.14.13 (p.~120)} }
\expandafter\def\csname SRef@push:src:A.4.6\endcsname{\href{S6_SERIES_GRIND_I.pdf\#nameddest=tag-02ZW}{I.14.15 (p.~121)} }
\expandafter\def\csname SRef@push:A46\endcsname{\href{S6_SERIES_GRIND_I.pdf\#nameddest=tag-02ZW}{I.14.15 (p.~121)} }
\expandafter\def\csname SRef@push:src:A.4.9\endcsname{\href{S6_SERIES_GRIND_I.pdf\#nameddest=tag-02ZZ}{I.14.18 (p.~126)} }
\expandafter\def\csname SRef@push:A49\endcsname{\href{S6_SERIES_GRIND_I.pdf\#nameddest=tag-02ZZ}{I.14.18 (p.~126)} }
\expandafter\def\csname SRef@push:src:A.2.1\endcsname{\href{S6_SERIES_GRIND_I.pdf\#nameddest=tag-02ZJ}{I.14.3 (p.~105)} }
\expandafter\def\csname SRef@push:A21\endcsname{\href{S6_SERIES_GRIND_I.pdf\#nameddest=tag-02ZJ}{I.14.3 (p.~105)} }
\expandafter\def\csname SRef@push:src:A.2.5\endcsname{\href{S6_SERIES_GRIND_I.pdf\#nameddest=tag-02ZN}{I.14.6 (p.~111)} }
\expandafter\def\csname SRef@push:A25\endcsname{\href{S6_SERIES_GRIND_I.pdf\#nameddest=tag-02ZN}{I.14.6 (p.~111)} }
\expandafter\def\csname SRef@push:src:A.3.3\endcsname{\href{S6_SERIES_GRIND_I.pdf\#nameddest=tag-02ZQ}{I.14.9 (p.~115)} }
\expandafter\def\csname SRef@push:A33\endcsname{\href{S6_SERIES_GRIND_I.pdf\#nameddest=tag-02ZQ}{I.14.9 (p.~115)} }
\expandafter\def\csname SRef@push:src:B.1.1\endcsname{\href{S6_SERIES_GRIND_I.pdf\#nameddest=tag-0300}{I.15.1 (p.~128)} }
\expandafter\def\csname SRef@push:B11\endcsname{\href{S6_SERIES_GRIND_I.pdf\#nameddest=tag-0300}{I.15.1 (p.~128)} }
\expandafter\def\csname SRef@push:src:B.3.6\endcsname{\href{S6_SERIES_GRIND_I.pdf\#nameddest=tag-030A}{I.15.11 (p.~138)} }
\expandafter\def\csname SRef@push:B36\endcsname{\href{S6_SERIES_GRIND_I.pdf\#nameddest=tag-030A}{I.15.11 (p.~138)} }
\expandafter\def\csname SRef@push:src:B.3.7\endcsname{\href{S6_SERIES_GRIND_I.pdf\#nameddest=tag-030B}{I.15.12 (p.~140)} }
\expandafter\def\csname SRef@push:B37\endcsname{\href{S6_SERIES_GRIND_I.pdf\#nameddest=tag-030B}{I.15.12 (p.~140)} }
\expandafter\def\csname SRef@push:src:B.3.8\endcsname{\href{S6_SERIES_GRIND_I.pdf\#nameddest=tag-030C}{I.15.13 (p.~142)} }
\expandafter\def\csname SRef@push:B38\endcsname{\href{S6_SERIES_GRIND_I.pdf\#nameddest=tag-030C}{I.15.13 (p.~142)} }
\expandafter\def\csname SRef@push:src:B.4.4\endcsname{\href{S6_SERIES_GRIND_I.pdf\#nameddest=tag-0314}{I.15.17 (p.~145)} }
\expandafter\def\csname SRef@push:src:B.4.6\endcsname{\href{S6_SERIES_GRIND_I.pdf\#nameddest=tag-0316}{I.15.19 (p.~148)} }
\expandafter\def\csname SRef@push:B46\endcsname{\href{S6_SERIES_GRIND_I.pdf\#nameddest=tag-0316}{I.15.19 (p.~148)} }
\expandafter\def\csname SRef@push:src:B.4.7\endcsname{\href{S6_SERIES_GRIND_I.pdf\#nameddest=tag-0317}{I.15.20 (p.~151)} }
\expandafter\def\csname SRef@push:src:B.4.8\endcsname{\href{S6_SERIES_GRIND_I.pdf\#nameddest=tag-0318}{I.15.21 (p.~155)} }
\expandafter\def\csname SRef@push:src:B.4.9\endcsname{\href{S6_SERIES_GRIND_I.pdf\#nameddest=tag-0319}{I.15.22 (p.~158)} }
\expandafter\def\csname SRef@push:src:B.4.10\endcsname{\href{S6_SERIES_GRIND_I.pdf\#nameddest=tag-030E}{I.15.23 (p.~160)} }
\expandafter\def\csname SRef@push:src:B.4.11\endcsname{\href{S6_SERIES_GRIND_I.pdf\#nameddest=tag-030F}{I.15.24 (p.~164)} }
\expandafter\def\csname SRef@push:src:B.4.12\endcsname{\href{S6_SERIES_GRIND_I.pdf\#nameddest=tag-030G}{I.15.25 (p.~167)} }
\expandafter\def\csname SRef@push:src:B.4.13\endcsname{\href{S6_SERIES_GRIND_I.pdf\#nameddest=tag-030H}{I.15.26 (p.~169)} }
\expandafter\def\csname SRef@push:src:B.4.14\endcsname{\href{S6_SERIES_GRIND_I.pdf\#nameddest=tag-030I}{I.15.27 (p.~171)} }
\expandafter\def\csname SRef@push:src:B.4.15\endcsname{\href{S6_SERIES_GRIND_I.pdf\#nameddest=tag-030J}{I.15.28 (p.~174)} }
\expandafter\def\csname SRef@push:B415\endcsname{\href{S6_SERIES_GRIND_I.pdf\#nameddest=tag-030J}{I.15.28 (p.~174)} }
\expandafter\def\csname SRef@push:src:B.4.16\endcsname{\href{S6_SERIES_GRIND_I.pdf\#nameddest=tag-030K}{I.15.29 (p.~177)} }
\expandafter\def\csname SRef@push:B416\endcsname{\href{S6_SERIES_GRIND_I.pdf\#nameddest=tag-030K}{I.15.29 (p.~177)} }
\expandafter\def\csname SRef@push:src:B.2.1\endcsname{\href{S6_SERIES_GRIND_I.pdf\#nameddest=tag-0302}{I.15.3 (p.~131)} }
\expandafter\def\csname SRef@push:B21\endcsname{\href{S6_SERIES_GRIND_I.pdf\#nameddest=tag-0302}{I.15.3 (p.~131)} }
\expandafter\def\csname SRef@push:src:B.4.17\endcsname{\href{S6_SERIES_GRIND_I.pdf\#nameddest=tag-030L}{I.15.30 (p.~178)} }
\expandafter\def\csname SRef@push:src:B.4.18\endcsname{\href{S6_SERIES_GRIND_I.pdf\#nameddest=tag-030M}{I.15.31 (p.~180)} }
\expandafter\def\csname SRef@push:B418\endcsname{\href{S6_SERIES_GRIND_I.pdf\#nameddest=tag-030M}{I.15.31 (p.~180)} }
\expandafter\def\csname SRef@push:src:B.4.19\endcsname{\href{S6_SERIES_GRIND_I.pdf\#nameddest=tag-030N}{I.15.32 (p.~182)} }
\expandafter\def\csname SRef@push:B419\endcsname{\href{S6_SERIES_GRIND_I.pdf\#nameddest=tag-030N}{I.15.32 (p.~182)} }
\expandafter\def\csname SRef@push:src:B.4.20\endcsname{\href{S6_SERIES_GRIND_I.pdf\#nameddest=tag-030P}{I.15.33 (p.~185)} }
\expandafter\def\csname SRef@push:B420\endcsname{\href{S6_SERIES_GRIND_I.pdf\#nameddest=tag-030P}{I.15.33 (p.~185)} }
\expandafter\def\csname SRef@push:src:B.4.21\endcsname{\href{S6_SERIES_GRIND_I.pdf\#nameddest=tag-030Q}{I.15.34 (p.~187)} }
\expandafter\def\csname SRef@push:src:B.4.22\endcsname{\href{S6_SERIES_GRIND_I.pdf\#nameddest=tag-030R}{I.15.35 (p.~191)} }
\expandafter\def\csname SRef@push:B422\endcsname{\href{S6_SERIES_GRIND_I.pdf\#nameddest=tag-030R}{I.15.35 (p.~191)} }
\expandafter\def\csname SRef@push:src:B.4.23\endcsname{\href{S6_SERIES_GRIND_I.pdf\#nameddest=tag-030S}{I.15.36 (p.~193)} }
\expandafter\def\csname SRef@push:src:B.4.24\endcsname{\href{S6_SERIES_GRIND_I.pdf\#nameddest=tag-030T}{I.15.37 (p.~195)} }
\expandafter\def\csname SRef@push:B424\endcsname{\href{S6_SERIES_GRIND_I.pdf\#nameddest=tag-030T}{I.15.37 (p.~195)} }
\expandafter\def\csname SRef@push:src:B.4.25\endcsname{\href{S6_SERIES_GRIND_I.pdf\#nameddest=tag-030U}{I.15.38 (p.~198)} }
\expandafter\def\csname SRef@push:B425\endcsname{\href{S6_SERIES_GRIND_I.pdf\#nameddest=tag-030U}{I.15.38 (p.~198)} }
\expandafter\def\csname SRef@push:src:B.4.26\endcsname{\href{S6_SERIES_GRIND_I.pdf\#nameddest=tag-030V}{I.15.39 (p.~200)} }
\expandafter\def\csname SRef@push:B426\endcsname{\href{S6_SERIES_GRIND_I.pdf\#nameddest=tag-030V}{I.15.39 (p.~200)} }
\expandafter\def\csname SRef@push:src:B.4.27\endcsname{\href{S6_SERIES_GRIND_I.pdf\#nameddest=tag-030W}{I.15.40 (p.~202)} }
\expandafter\def\csname SRef@push:B427\endcsname{\href{S6_SERIES_GRIND_I.pdf\#nameddest=tag-030W}{I.15.40 (p.~202)} }
\expandafter\def\csname SRef@push:src:B.4.28\endcsname{\href{S6_SERIES_GRIND_I.pdf\#nameddest=tag-030X}{I.15.41 (p.~204)} }
\expandafter\def\csname SRef@push:src:B.4.29\endcsname{\href{S6_SERIES_GRIND_I.pdf\#nameddest=tag-030Y}{I.15.42 (p.~208)} }
\expandafter\def\csname SRef@push:B429\endcsname{\href{S6_SERIES_GRIND_I.pdf\#nameddest=tag-030Y}{I.15.42 (p.~208)} }
\expandafter\def\csname SRef@push:src:B.4.30\endcsname{\href{S6_SERIES_GRIND_I.pdf\#nameddest=tag-0310}{I.15.43 (p.~210)} }
\expandafter\def\csname SRef@push:B430\endcsname{\href{S6_SERIES_GRIND_I.pdf\#nameddest=tag-0310}{I.15.43 (p.~210)} }
\expandafter\def\csname SRef@push:src:B.4.31\endcsname{\href{S6_SERIES_GRIND_I.pdf\#nameddest=tag-0311}{I.15.44 (p.~216)} }
\expandafter\def\csname SRef@push:src:B.4.32\endcsname{\href{S6_SERIES_GRIND_I.pdf\#nameddest=tag-0312}{I.15.45 (p.~219)} }
\expandafter\def\csname SRef@push:B432\endcsname{\href{S6_SERIES_GRIND_I.pdf\#nameddest=tag-0312}{I.15.45 (p.~219)} }
\expandafter\def\csname SRef@push:src:B.3.2\endcsname{\href{S6_SERIES_GRIND_I.pdf\#nameddest=tag-0306}{I.15.7 (p.~134)} }
\expandafter\def\csname SRef@push:B32\endcsname{\href{S6_SERIES_GRIND_I.pdf\#nameddest=tag-0306}{I.15.7 (p.~134)} }
\expandafter\def\csname SRef@push:src:B.3.3\endcsname{\href{S6_SERIES_GRIND_I.pdf\#nameddest=tag-0307}{I.15.8 (p.~134)} }
\expandafter\def\csname SRef@push:B33\endcsname{\href{S6_SERIES_GRIND_I.pdf\#nameddest=tag-0307}{I.15.8 (p.~134)} }
\expandafter\def\csname SRef@push:src:D.1.1\endcsname{\href{S6_SERIES_GRIND_I.pdf\#nameddest=tag-032A}{I.16.1 (p.~227)} }
\expandafter\def\csname SRef@push:src:D.4.1\endcsname{\href{S6_SERIES_GRIND_I.pdf\#nameddest=tag-032J}{I.16.10 (p.~235)} }
\expandafter\def\csname SRef@push:src:D.4.2\endcsname{\href{S6_SERIES_GRIND_I.pdf\#nameddest=tag-032K}{I.16.11 (p.~236)} }
\expandafter\def\csname SRef@push:src:D.4.3\endcsname{\href{S6_SERIES_GRIND_I.pdf\#nameddest=tag-032L}{I.16.12 (p.~237)} }
\expandafter\def\csname SRef@push:src:D.4.4\endcsname{\href{S6_SERIES_GRIND_I.pdf\#nameddest=tag-032M}{I.16.13 (p.~238)} }
\expandafter\def\csname SRef@push:src:D.4.5\endcsname{\href{S6_SERIES_GRIND_I.pdf\#nameddest=tag-032N}{I.16.14 (p.~238)} }
\expandafter\def\csname SRef@push:src:D.4.7\endcsname{\href{S6_SERIES_GRIND_I.pdf\#nameddest=tag-032P}{I.16.16 (p.~241)} }
\expandafter\def\csname SRef@push:src:D.4.8\endcsname{\href{S6_SERIES_GRIND_I.pdf\#nameddest=tag-032Q}{I.16.17 (p.~242)} }
\expandafter\def\csname SRef@push:src:D.2.2\endcsname{\href{S6_SERIES_GRIND_I.pdf\#nameddest=tag-032E}{I.16.5 (p.~231)} }
\expandafter\def\csname SRef@push:src:D.2.4\endcsname{\href{S6_SERIES_GRIND_I.pdf\#nameddest=tag-032G}{I.16.7 (p.~232)} }
\expandafter\def\csname SRef@inv:A3.2\endcsname{\href{S6_SERIES_GRIND_I.pdf\#nameddest=tag-01RZ}{I.18.10 (p.~264)} }
\expandafter\def\csname SRef@wb:lem:bcpr-eta-extension\endcsname{\href{S6_SERIES_GRIND_I.pdf\#nameddest=tag-01RZ}{I.18.10 (p.~264)} }
\expandafter\def\csname SRef@inv:A4.1\endcsname{\href{S6_SERIES_GRIND_I.pdf\#nameddest=tag-01S0}{I.18.11 (p.~264)} }
\expandafter\def\csname SRef@inv:A4.2\endcsname{\href{S6_SERIES_GRIND_I.pdf\#nameddest=tag-01S1}{I.18.12 (p.~265)} }
\expandafter\def\csname SRef@inv:A1.1\endcsname{\href{S6_SERIES_GRIND_I.pdf\#nameddest=tag-01RR}{I.18.2 (p.~254)} }
\expandafter\def\csname SRef@inv:A1.2\endcsname{\href{S6_SERIES_GRIND_I.pdf\#nameddest=tag-01RS}{I.18.3 (p.~255)} }
\expandafter\def\csname SRef@wb:prop:log-relative-direct-images\endcsname{\href{S6_SERIES_GRIND_I.pdf\#nameddest=tag-01RS}{I.18.3 (p.~255)} }
\expandafter\def\csname SRef@inv:A1.3\endcsname{\href{S6_SERIES_GRIND_I.pdf\#nameddest=tag-01RT}{I.18.4 (p.~256)} }
\expandafter\def\csname SRef@wb:prop:absolute-cotangent-direct-images\endcsname{\href{S6_SERIES_GRIND_I.pdf\#nameddest=tag-01RT}{I.18.4 (p.~256)} }
\expandafter\def\csname SRef@inv:A2.1\endcsname{\href{S6_SERIES_GRIND_I.pdf\#nameddest=tag-01RU}{I.18.5 (p.~257)} }
\expandafter\def\csname SRef@wb:prop:log-ks-divisor\endcsname{\href{S6_SERIES_GRIND_I.pdf\#nameddest=tag-01RU}{I.18.5 (p.~257)} }
\expandafter\def\csname SRef@wb:prop:period-derivative-determinant\endcsname{\href{S6_SERIES_GRIND_I.pdf\#nameddest=tag-01RU}{I.18.5 (p.~257)} }
\expandafter\def\csname SRef@inv:A2.2\endcsname{\href{S6_SERIES_GRIND_I.pdf\#nameddest=tag-01RV}{I.18.6 (p.~258)} }
\expandafter\def\csname SRef@inv:A2.3\endcsname{\href{S6_SERIES_GRIND_I.pdf\#nameddest=tag-01RW}{I.18.7 (p.~260)} }
\expandafter\def\csname SRef@inv:A2.4\endcsname{\href{S6_SERIES_GRIND_I.pdf\#nameddest=tag-01RX}{I.18.8 (p.~262)} }
\expandafter\def\csname SRef@wb:thm:h12-scalar\endcsname{\href{S6_SERIES_GRIND_I.pdf\#nameddest=tag-01RX}{I.18.8 (p.~262)} }
\expandafter\def\csname SRef@inv:A3.1\endcsname{\href{S6_SERIES_GRIND_I.pdf\#nameddest=tag-01RY}{I.18.9 (p.~263)} }
\expandafter\def\csname SRef@wb:cor:exact-hodge-diamond\endcsname{\href{S6_SERIES_GRIND_I.pdf\#nameddest=tag-01RY}{I.18.9 (p.~263)} }
\expandafter\def\csname SRef@wb:thm:parametric-hodge-diamond\endcsname{\href{S6_SERIES_GRIND_I.pdf\#nameddest=tag-01RY}{I.18.9 (p.~263)} }
\expandafter\def\csname SRef@wb:prop:workbench-holomorphic-form-vanishing\endcsname{\href{S6_SERIES_GRIND_I.pdf\#nameddest=tag-01RY}{I.18.9 (p.~263)} }
\expandafter\def\csname SRef@inv:C1.1\endcsname{\href{S6_SERIES_GRIND_I.pdf\#nameddest=tag-01SA}{I.19.1 (p.~266)} }
\expandafter\def\csname SRef@wb:thm:frolicher-e2-candidate\endcsname{\href{S6_SERIES_GRIND_I.pdf\#nameddest=tag-01SA}{I.19.1 (p.~266)} }
\expandafter\def\csname SRef@inv:C3.7\endcsname{\href{S6_SERIES_GRIND_I.pdf\#nameddest=tag-01SK}{I.19.11 (p.~274)} }
\expandafter\def\csname SRef@inv:C1.2\endcsname{\href{S6_SERIES_GRIND_I.pdf\#nameddest=tag-01SB}{I.19.2 (p.~267)} }
\expandafter\def\csname SRef@inv:C2.1\endcsname{\href{S6_SERIES_GRIND_I.pdf\#nameddest=tag-01SC}{I.19.3 (p.~267)} }
\expandafter\def\csname SRef@wb:thm:exact-bott-chern-aeppli\endcsname{\href{S6_SERIES_GRIND_I.pdf\#nameddest=tag-01SC}{I.19.3 (p.~267)} }
\expandafter\def\csname SRef@inv:C2.2\endcsname{\href{S6_SERIES_GRIND_I.pdf\#nameddest=tag-01SD}{I.19.4 (p.~268)} }
\expandafter\def\csname SRef@inv:C3.1\endcsname{\href{S6_SERIES_GRIND_I.pdf\#nameddest=tag-01SE}{I.19.5 (p.~269)} }
\expandafter\def\csname SRef@wb:cor:picard-line-audit\endcsname{\href{S6_SERIES_GRIND_I.pdf\#nameddest=tag-01SE}{I.19.5 (p.~269)} }
\expandafter\def\csname SRef@inv:C3.2\endcsname{\href{S6_SERIES_GRIND_I.pdf\#nameddest=tag-01SF}{I.19.6 (p.~270)} }
\expandafter\def\csname SRef@wb:lem:bcpr-degree-one-input\endcsname{\href{S6_SERIES_GRIND_I.pdf\#nameddest=tag-01SF}{I.19.6 (p.~270)} }
\expandafter\def\csname SRef@inv:C3.3\endcsname{\href{S6_SERIES_GRIND_I.pdf\#nameddest=tag-01SG}{I.19.7 (p.~272)} }
\expandafter\def\csname SRef@wb:thm:bcpr-square-zero-algebra\endcsname{\href{S6_SERIES_GRIND_I.pdf\#nameddest=tag-01SG}{I.19.7 (p.~272)} }
\expandafter\def\csname SRef@wb:prop:bcpr-generators-are-bases\endcsname{\href{S6_SERIES_GRIND_I.pdf\#nameddest=tag-01SG}{I.19.7 (p.~272)} }
\expandafter\def\csname SRef@inv:C3.4\endcsname{\href{S6_SERIES_GRIND_I.pdf\#nameddest=tag-01SH}{I.19.8 (p.~273)} }
\expandafter\def\csname SRef@inv:C3.5\endcsname{\href{S6_SERIES_GRIND_I.pdf\#nameddest=tag-01SI}{I.19.9 (p.~273)} }
\expandafter\def\csname SRef@inv:D4.1\endcsname{\href{S6_SERIES_GRIND_I.pdf\#nameddest=tag-01SO}{I.20.1 (p.~275)} }
\expandafter\def\csname SRef@inv:D4.2\endcsname{\href{S6_SERIES_GRIND_I.pdf\#nameddest=tag-01SP}{I.20.2 (p.~276)} }
\expandafter\def\csname SRef@inv:D4.3\endcsname{\href{S6_SERIES_GRIND_I.pdf\#nameddest=tag-01SQ}{I.20.3 (p.~276)} }
\expandafter\def\csname SRef@inv:D4.4\endcsname{\href{S6_SERIES_GRIND_I.pdf\#nameddest=tag-01SR}{I.20.4 (p.~277)} }
\expandafter\def\csname SRef@inv:D4.5\endcsname{\href{S6_SERIES_GRIND_I.pdf\#nameddest=tag-01SS}{I.20.5 (p.~280)} }
\expandafter\def\csname SRef@push:src:E.2.2\endcsname{\href{S6_SERIES_GRIND_I.pdf\#nameddest=tag-0335}{I.20.7 (p.~284)} }
\expandafter\def\csname SRef@push:src:E.2.3\endcsname{\href{S6_SERIES_GRIND_I.pdf\#nameddest=tag-033B}{I.20.8 (p.~284)} }
\expandafter\def\csname SRef@inv:B1.0\endcsname{\href{S6_SERIES_GRIND_I.pdf\#nameddest=tag-01S2}{I.21.1 (p.~286)} }
\expandafter\def\csname SRef@inv:B1.1\endcsname{\href{S6_SERIES_GRIND_I.pdf\#nameddest=tag-01S3}{I.21.2 (p.~287)} }
\expandafter\def\csname SRef@inv:B1.2\endcsname{\href{S6_SERIES_GRIND_I.pdf\#nameddest=tag-01S4}{I.21.3 (p.~288)} }
\expandafter\def\csname SRef@inv:B2.1\endcsname{\href{S6_SERIES_GRIND_I.pdf\#nameddest=tag-01S5}{I.21.4 (p.~288)} }
\expandafter\def\csname SRef@inv:B2.2\endcsname{\href{S6_SERIES_GRIND_I.pdf\#nameddest=tag-01S6}{I.21.5 (p.~289)} }
\expandafter\def\csname SRef@inv:B3.1\endcsname{\href{S6_SERIES_PUBLISH_I.pdf\#nameddest=tag-01S7}{I.23.7 (p.~259)} }
\expandafter\def\csname SRef@inv:B3.2\endcsname{\href{S6_SERIES_GRIND_I.pdf\#nameddest=tag-01S8}{I.21.7 (p.~292)} }
\expandafter\def\csname SRef@inv:B3.3\endcsname{\href{S6_SERIES_GRIND_I.pdf\#nameddest=tag-01S9}{I.21.8 (p.~292)} }
\expandafter\def\csname SRef@mapsI:src:A.1.1\endcsname{\href{S6_SERIES_GRIND_I.pdf\#nameddest=tag-02PA}{I.22.1 (p.~295)} }
\expandafter\def\csname SRef@mapsI:quotient-eq-1\endcsname{\href{S6_SERIES_GRIND_I.pdf\#nameddest=tag-02PA}{I.22.1 (p.~295)} }
\expandafter\def\csname SRef@mapsI:src:A.1.2\endcsname{\href{S6_SERIES_PUBLISH_I.pdf\#nameddest=tag-02PB}{I.16.4 (p.~114)} }
\expandafter\def\csname SRef@mapsI:src:A.1.3\endcsname{\href{S6_SERIES_GRIND_I.pdf\#nameddest=tag-02PC}{I.22.3 (p.~297)} }
\expandafter\def\csname SRef@mapsI:quotient-eq-6\endcsname{\href{S6_SERIES_GRIND_I.pdf\#nameddest=tag-02PC}{I.22.3 (p.~298)} }
\expandafter\def\csname SRef@ari:C.2.1\endcsname{\href{S6_SERIES_GRIND_I.pdf\#nameddest=tag-000L}{I.22.5 (p.~299)} }
\expandafter\def\csname SRef@ari:C.2.2\endcsname{\href{S6_SERIES_GRIND_I.pdf\#nameddest=tag-000M}{I.22.6 (p.~300)} }
\expandafter\def\csname SRef@ari:C.3.1\endcsname{\href{S6_SERIES_GRIND_I.pdf\#nameddest=tag-000N}{I.22.7 (p.~301)} }
\expandafter\def\csname SRef@ari:C.3.2\endcsname{\href{S6_SERIES_GRIND_I.pdf\#nameddest=tag-000O}{I.22.8 (p.~302)} }
\expandafter\def\csname SRef@ari:C.6.1\endcsname{\href{S6_SERIES_GRIND_I.pdf\#nameddest=tag-000V}{I.23.1 (p.~304)} }
\expandafter\def\csname SRef@ari:C.6.R5\endcsname{\href{S6_SERIES_GRIND_I.pdf\#nameddest=tag-0014}{I.23.10 (p.~316)} }
\expandafter\def\csname SRef@ari:C.6.T1\endcsname{\href{S6_SERIES_GRIND_I.pdf\#nameddest=tag-0015}{I.23.11 (p.~322)} }
\expandafter\def\csname SRef@ari:C.6.T2\endcsname{\href{S6_SERIES_GRIND_I.pdf\#nameddest=tag-0016}{I.23.12 (p.~323)} }
\expandafter\def\csname SRef@ari:C.8.6\endcsname{\href{S6_SERIES_GRIND_I.pdf\#nameddest=tag-001T}{I.23.13 (p.~325)} }
\expandafter\def\csname SRef@ari:C.8.7\endcsname{\href{S6_SERIES_GRIND_I.pdf\#nameddest=tag-001U}{I.23.14 (p.~326)} }
\expandafter\def\csname SRef@ari:C.8.8\endcsname{\href{S6_SERIES_GRIND_I.pdf\#nameddest=tag-001V}{I.23.15 (p.~326)} }
\expandafter\def\csname SRef@ari:C.8.9\endcsname{\href{S6_SERIES_GRIND_I.pdf\#nameddest=tag-001W}{I.23.16 (p.~327)} }
\expandafter\def\csname SRef@ari:C.8.10\endcsname{\href{S6_SERIES_GRIND_I.pdf\#nameddest=tag-001K}{I.23.17 (p.~329)} }
\expandafter\def\csname SRef@ari:C.8.11\endcsname{\href{S6_SERIES_GRIND_I.pdf\#nameddest=tag-001L}{I.23.18 (p.~330)} }
\expandafter\def\csname SRef@ari:C.6.2\endcsname{\href{S6_SERIES_GRIND_I.pdf\#nameddest=tag-000W}{I.23.2 (p.~305)} }
\expandafter\def\csname SRef@ari:C.6.3\endcsname{\href{S6_SERIES_GRIND_I.pdf\#nameddest=tag-000X}{I.23.3 (p.~305)} }
\expandafter\def\csname SRef@ari:C.6.4\endcsname{\href{S6_SERIES_GRIND_I.pdf\#nameddest=tag-000Y}{I.23.4 (p.~308)} }
\expandafter\def\csname SRef@ari:C.6.5\endcsname{\href{S6_SERIES_GRIND_I.pdf\#nameddest=tag-000Z}{I.23.5 (p.~309)} }
\expandafter\def\csname SRef@ari:C.6.R1\endcsname{\href{S6_SERIES_GRIND_I.pdf\#nameddest=tag-0010}{I.23.6 (p.~310)} }
\expandafter\def\csname SRef@ari:C.6.R2\endcsname{\href{S6_SERIES_GRIND_I.pdf\#nameddest=tag-0011}{I.23.7 (p.~311)} }
\expandafter\def\csname SRef@ari:C.6.R3\endcsname{\href{S6_SERIES_GRIND_I.pdf\#nameddest=tag-0012}{I.23.8 (p.~314)} }
\expandafter\def\csname SRef@ari:C.8.12\endcsname{\href{S6_SERIES_GRIND_I.pdf\#nameddest=tag-001M}{I.24.1 (p.~331)} }
\expandafter\def\csname SRef@ari:C.8.13\endcsname{\href{S6_SERIES_GRIND_I.pdf\#nameddest=tag-001N}{I.24.2 (p.~332)} }
\expandafter\def\csname SRef@ari:C.8.14\endcsname{\href{S6_SERIES_GRIND_I.pdf\#nameddest=tag-001O}{I.24.3 (p.~332)} }
\expandafter\def\csname SRef@ari:C.8.15\endcsname{\href{S6_SERIES_GRIND_I.pdf\#nameddest=tag-001P}{I.24.4 (p.~337)} }
\expandafter\def\csname SRef@fam:src:B.1.1\endcsname{\href{S6_SERIES_GRIND_I.pdf\#nameddest=tag-01Q7}{I.3.2 (p.~24)} }
\expandafter\def\csname SRef@fam:st:B.1.1\endcsname{\href{S6_SERIES_GRIND_I.pdf\#nameddest=tag-01Q7}{I.3.2 (p.~24)} }
\expandafter\def\csname SRef@fam:src:B.1.2\endcsname{\href{S6_SERIES_GRIND_I.pdf\#nameddest=tag-01Q8}{I.3.3 (p.~26)} }
\expandafter\def\csname SRef@fam:st:B.1.2\endcsname{\href{S6_SERIES_GRIND_I.pdf\#nameddest=tag-01Q8}{I.3.3 (p.~26)} }
\expandafter\def\csname SRef@fam:src:B.1.3\endcsname{\href{S6_SERIES_GRIND_I.pdf\#nameddest=tag-01Q9}{I.3.4 (p.~26)} }
\expandafter\def\csname SRef@fam:st:B.1.3\endcsname{\href{S6_SERIES_GRIND_I.pdf\#nameddest=tag-01Q9}{I.3.4 (p.~26)} }
\expandafter\def\csname SRef@fam:src:B.1.4\endcsname{\href{S6_SERIES_GRIND_I.pdf\#nameddest=tag-01QA}{I.3.5 (p.~28)} }
\expandafter\def\csname SRef@fam:st:B.1.4\endcsname{\href{S6_SERIES_GRIND_I.pdf\#nameddest=tag-01QA}{I.3.5 (p.~28)} }
\expandafter\def\csname SRef@fam:src:C.1.1\endcsname{\href{S6_SERIES_PUBLISH_I.pdf\#nameddest=tag-01QH}{I.3.2 (p.~25)} }
\expandafter\def\csname SRef@I:def:finite-coordinates\endcsname{\href{S6_SERIES_PUBLISH_I.pdf\#nameddest=tag-PI07}{I.12.8 (p.~82)} }
\expandafter\def\csname SRef@fam:st:C.1.1\endcsname{\href{S6_SERIES_GRIND_I.pdf\#nameddest=tag-01QH}{I.4.1 (p.~29)} }
\expandafter\def\csname SRef@fam:src:C.1.2\endcsname{\href{S6_SERIES_GRIND_I.pdf\#nameddest=tag-01QI}{I.4.2 (p.~30)} }
\expandafter\def\csname SRef@fam:st:C.1.2\endcsname{\href{S6_SERIES_GRIND_I.pdf\#nameddest=tag-01QI}{I.4.2 (p.~30)} }
\expandafter\def\csname SRef@fam:src:C.2.1\endcsname{\href{S6_SERIES_GRIND_I.pdf\#nameddest=tag-01QJ}{I.4.3 (p.~31)} }
\expandafter\def\csname SRef@fam:st:C.2.1\endcsname{\href{S6_SERIES_GRIND_I.pdf\#nameddest=tag-01QJ}{I.4.3 (p.~31)} }
\expandafter\def\csname SRef@fam:src:C.2.2\endcsname{\href{S6_SERIES_GRIND_I.pdf\#nameddest=tag-01QK}{I.4.4 (p.~32)} }
\expandafter\def\csname SRef@fam:st:C.2.2\endcsname{\href{S6_SERIES_GRIND_I.pdf\#nameddest=tag-01QK}{I.4.4 (p.~32)} }
\expandafter\def\csname SRef@ari:C.1.1\endcsname{\href{S6_SERIES_GRIND_I.pdf\#nameddest=tag-000I}{I.5.1 (p.~36)} }
\expandafter\def\csname SRef@ari:C.1.2\endcsname{\href{S6_SERIES_GRIND_I.pdf\#nameddest=tag-000J}{I.5.2 (p.~37)} }
\expandafter\def\csname SRef@ari:C.1.3\endcsname{\href{S6_SERIES_GRIND_I.pdf\#nameddest=tag-000K}{I.5.3 (p.~37)} }
\expandafter\def\csname SRef@ari:C.4.2\endcsname{\href{S6_SERIES_GRIND_I.pdf\#nameddest=tag-000Q}{I.5.5 (p.~39)} }
\expandafter\def\csname SRef@ari:C.4.3\endcsname{\href{S6_SERIES_GRIND_I.pdf\#nameddest=tag-000R}{I.5.6 (p.~39)} }
\expandafter\def\csname SRef@ari:C.4.4\endcsname{\href{S6_SERIES_GRIND_I.pdf\#nameddest=tag-000S}{I.5.7 (p.~40)} }
\expandafter\def\csname SRef@ari:C.4.5\endcsname{\href{S6_SERIES_GRIND_I.pdf\#nameddest=tag-000T}{I.5.8 (p.~42)} }
\expandafter\def\csname SRef@ari:C.8.1\endcsname{\href{S6_SERIES_GRIND_I.pdf\#nameddest=tag-001J}{I.6.1 (p.~42)} }
\expandafter\def\csname SRef@ari:C.8.2\endcsname{\href{S6_SERIES_GRIND_I.pdf\#nameddest=tag-001Q}{I.6.2 (p.~44)} }
\expandafter\def\csname SRef@ari:C.8.3\endcsname{\href{S6_SERIES_GRIND_I.pdf\#nameddest=tag-001R}{I.6.3 (p.~46)} }
\expandafter\def\csname SRef@ari:C.8.4\endcsname{\href{S6_SERIES_GRIND_I.pdf\#nameddest=tag-001S}{I.6.4 (p.~47)} }
\expandafter\def\csname SRef@fam:src:A.1.1\endcsname{\href{S6_SERIES_PUBLISH_I.pdf\#nameddest=tag-01PZ}{I.7.2 (p.~42)} }
\expandafter\def\csname SRef@fam:st:A.1.1\endcsname{\href{S6_SERIES_GRIND_I.pdf\#nameddest=tag-01PZ}{I.7.1 (p.~49)} }
\expandafter\def\csname SRef@fam:src:A.1.2\endcsname{\href{S6_SERIES_PUBLISH_I.pdf\#nameddest=tag-01Q0}{I.7.4 (p.~43)} }
\expandafter\def\csname SRef@fam:st:A.1.2\endcsname{\href{S6_SERIES_GRIND_I.pdf\#nameddest=tag-01Q0}{I.7.2 (p.~49)} }
\expandafter\def\csname SRef@fam:src:A.2.1\endcsname{\href{S6_SERIES_GRIND_I.pdf\#nameddest=tag-01Q1}{I.7.3 (p.~50)} }
\expandafter\def\csname SRef@fam:st:A.2.1\endcsname{\href{S6_SERIES_GRIND_I.pdf\#nameddest=tag-01Q1}{I.7.3 (p.~50)} }
\expandafter\def\csname SRef@fam:src:A.3.1\endcsname{\href{S6_SERIES_PUBLISH_I.pdf\#nameddest=tag-01Q2}{I.7.5 (p.~44)} }
\expandafter\def\csname SRef@pub:I:eq:PD7\endcsname{\href{S6_SERIES_PUBLISH_I.pdf\#nameddest=equation.7.8}{I.(7.8) (p.~45)} }
\expandafter\def\csname SRef@pub:I:eq:PD8\endcsname{\href{S6_SERIES_PUBLISH_I.pdf\#nameddest=equation.7.9}{I.(7.9) (p.~45)} }
\expandafter\def\csname SRef@pub:I:eq:PD11\endcsname{\href{S6_SERIES_PUBLISH_I.pdf\#nameddest=equation.7.12}{I.(7.12) (p.~45)} }
\expandafter\def\csname SRef@fam:st:A.3.1\endcsname{\href{S6_SERIES_GRIND_I.pdf\#nameddest=tag-01Q2}{I.7.4 (p.~51)} }
\expandafter\def\csname SRef@fam:src:A.4.1\endcsname{\href{S6_SERIES_GRIND_I.pdf\#nameddest=tag-01Q3}{I.7.5 (p.~53)} }
\expandafter\def\csname SRef@fam:st:A.4.1\endcsname{\href{S6_SERIES_GRIND_I.pdf\#nameddest=tag-01Q3}{I.7.5 (p.~53)} }
\expandafter\def\csname SRef@fam:src:A.4.2\endcsname{\href{S6_SERIES_GRIND_I.pdf\#nameddest=tag-01Q4}{I.7.6 (p.~54)} }
\expandafter\def\csname SRef@fam:st:A.4.2\endcsname{\href{S6_SERIES_GRIND_I.pdf\#nameddest=tag-01Q4}{I.7.6 (p.~54)} }
\expandafter\def\csname SRef@fam:src:A.4.3\endcsname{\href{S6_SERIES_GRIND_I.pdf\#nameddest=tag-01Q5}{I.7.7 (p.~55)} }
\expandafter\def\csname SRef@fam:st:A.4.3\endcsname{\href{S6_SERIES_GRIND_I.pdf\#nameddest=tag-01Q5}{I.7.7 (p.~55)} }
\expandafter\def\csname SRef@fam:src:B.2.1\endcsname{\href{S6_SERIES_PUBLISH_I.pdf\#nameddest=tag-01QB}{I.8.1 (p.~52)} }
\expandafter\def\csname SRef@fam:st:B.2.1\endcsname{\href{S6_SERIES_GRIND_I.pdf\#nameddest=tag-01QB}{I.8.1 (p.~57)} }
\expandafter\def\csname SRef@fam:src:B.2.2\endcsname{\href{S6_SERIES_GRIND_I.pdf\#nameddest=tag-01QC}{I.8.2 (p.~59)} }
\expandafter\def\csname SRef@fam:st:B.2.2\endcsname{\href{S6_SERIES_GRIND_I.pdf\#nameddest=tag-01QC}{I.8.2 (p.~59)} }
\expandafter\def\csname SRef@fam:src:B.2.3\endcsname{\href{S6_SERIES_GRIND_I.pdf\#nameddest=tag-01QD}{I.8.3 (p.~60)} }
\expandafter\def\csname SRef@fam:st:B.2.3\endcsname{\href{S6_SERIES_GRIND_I.pdf\#nameddest=tag-01QD}{I.8.3 (p.~60)} }
\expandafter\def\csname SRef@fam:src:B.3.1\endcsname{\href{S6_SERIES_GRIND_I.pdf\#nameddest=tag-01QE}{I.8.4 (p.~60)} }
\expandafter\def\csname SRef@fam:st:B.3.1\endcsname{\href{S6_SERIES_GRIND_I.pdf\#nameddest=tag-01QE}{I.8.4 (p.~60)} }
\expandafter\def\csname SRef@fam:src:B.4.1\endcsname{\href{S6_SERIES_GRIND_I.pdf\#nameddest=tag-01QF}{I.8.5 (p.~62)} }
\expandafter\def\csname SRef@fam:st:B.4.1\endcsname{\href{S6_SERIES_GRIND_I.pdf\#nameddest=tag-01QF}{I.8.5 (p.~62)} }
\expandafter\def\csname SRef@fam:src:B.4.2\endcsname{\href{S6_SERIES_GRIND_I.pdf\#nameddest=tag-01QG}{I.8.6 (p.~62)} }
\expandafter\def\csname SRef@fam:st:B.4.2\endcsname{\href{S6_SERIES_GRIND_I.pdf\#nameddest=tag-01QG}{I.8.6 (p.~62)} }
\expandafter\def\csname SRef@push:src:E.1.1\endcsname{\href{S6_SERIES_GRIND_I.pdf\#nameddest=tag-032R}{I.9.1 (p.~64)} }
\expandafter\def\csname SRef@inv:D1.1\endcsname{\href{S6_SERIES_GRIND_I.pdf\#nameddest=tag-01SM}{I.B.1 (p.~360)} }
\expandafter\def\csname SRef@inv:D1.2\endcsname{\href{S6_SERIES_GRIND_I.pdf\#nameddest=tag-01SN}{I.B.2 (p.~361)} }
\expandafter\def\csname SRef@ari:C.7.1\endcsname{\href{S6_SERIES_GRIND_I.pdf\#nameddest=tag-0017}{I.E.1 (p.~383)} }
\expandafter\def\csname SRef@ari:C.7.11\endcsname{\href{S6_SERIES_GRIND_I.pdf\#nameddest=tag-0019}{I.E.10 (p.~398)} }
\expandafter\def\csname SRef@ari:C.7.12\endcsname{\href{S6_SERIES_GRIND_I.pdf\#nameddest=tag-001A}{I.E.11 (p.~400)} }
\expandafter\def\csname SRef@ari:C.7.2\endcsname{\href{S6_SERIES_GRIND_I.pdf\#nameddest=tag-001B}{I.E.2 (p.~384)} }
\expandafter\def\csname SRef@ari:C.7.3\endcsname{\href{S6_SERIES_GRIND_I.pdf\#nameddest=tag-001C}{I.E.3 (p.~385)} }
\expandafter\def\csname SRef@ari:C.7.4\endcsname{\href{S6_SERIES_GRIND_I.pdf\#nameddest=tag-001D}{I.E.4 (p.~386)} }
\expandafter\def\csname SRef@ari:C.7.6\endcsname{\href{S6_SERIES_GRIND_I.pdf\#nameddest=tag-001F}{I.E.5 (p.~387)} }
\expandafter\def\csname SRef@ari:C.7.9\endcsname{\href{S6_SERIES_GRIND_I.pdf\#nameddest=tag-001I}{I.E.8 (p.~391)} }
\expandafter\def\csname SRef@ari:C.7.10\endcsname{\href{S6_SERIES_GRIND_I.pdf\#nameddest=tag-0018}{I.E.9 (p.~394)} }
\expandafter\def\csname SRef@inv:C3.8\endcsname{\href{S6_SERIES_GRIND_I.pdf\#nameddest=tag-01SL}{I.D.1 (p.~374)} }
\expandafter\def\csname SRef@fam:src:C.3.3\endcsname{\href{S6_SERIES_GRIND_I.pdf\#nameddest=tag-01QO}{I.11.3 (p.~86)} }
\expandafter\def\csname SRef@fam:st:C.3.3\endcsname{\href{S6_SERIES_GRIND_I.pdf\#nameddest=tag-01QO}{I.11.3 (p.~86)} }
\expandafter\def\csname SRef@fam:src:E.1.5\endcsname{\href{S6_SERIES_GRIND_I.pdf\#nameddest=tag-01R5}{I.13.1 (p.~100)} }
\expandafter\def\csname SRef@fam:st:E.1.5\endcsname{\href{S6_SERIES_GRIND_I.pdf\#nameddest=tag-01R5}{I.13.1 (p.~100)} }
\expandafter\def\csname SRef@wb:rem:bcpr-literal-gauge\endcsname{\href{S6_SERIES_GRIND_I.pdf\#nameddest=tag-03CT}{I.D.3 (p.~382)} }
\expandafter\def\csname SRef@wb:rem:bcpr-global-gauge-boundary\endcsname{\href{S6_SERIES_GRIND_I.pdf\#nameddest=tag-03CS}{I.D.3 (p.~382)} }
\expandafter\def\csname SRef@push:src:A.2.3\endcsname{\href{S6_SERIES_GRIND_I.pdf\#nameddest=tag-02ZL}{I.14.2 (p.~107)} }
\expandafter\def\csname SRef@push:A23\endcsname{\href{S6_SERIES_GRIND_I.pdf\#nameddest=tag-02ZL}{I.14.2 (p.~107)} }
\expandafter\def\csname SRef@fam:section:F\endcsname{\href{S6_SERIES_GRIND_I.pdf\#nameddest=tag-01PV}{I.10 (p.~68)} }
\expandafter\def\csname SRef@fam:section:D\endcsname{\href{S6_SERIES_GRIND_I.pdf\#nameddest=tag-01PP}{I.11 (p.~76)} }
\expandafter\def\csname SRef@mapsI:real-forms\endcsname{\href{S6_SERIES_GRIND_I.pdf\#nameddest=tag-02P9}{I.12 (p.~88)} }
\expandafter\def\csname SRef@fam:section:E\endcsname{\href{S6_SERIES_GRIND_I.pdf\#nameddest=tag-01PS}{I.13 (p.~96)} }
\expandafter\def\csname SRef@push:section:A\endcsname{\href{S6_SERIES_GRIND_I.pdf\#nameddest=tag-02YZ}{I.14 (p.~103)} }
\expandafter\def\csname SRef@push:section:B\endcsname{\href{S6_SERIES_GRIND_I.pdf\#nameddest=tag-02Z4}{I.15 (p.~128)} }
\expandafter\def\csname SRef@push:section:D\endcsname{\href{S6_SERIES_GRIND_I.pdf\#nameddest=tag-02Z8}{I.16 (p.~227)} }
\expandafter\def\csname SRef@inv:sec:partA\endcsname{\href{S6_SERIES_GRIND_I.pdf\#nameddest=tag-01TH}{I.18 (p.~254)} }
\expandafter\def\csname SRef@inv:sec:partC\endcsname{\href{S6_SERIES_GRIND_I.pdf\#nameddest=tag-01TI}{I.19 (p.~266)} }
\expandafter\def\csname SRef@def:common-conventions\endcsname{\href{S6_SERIES_GRIND_I.pdf\#nameddest=tag-01O1}{I.2 (p.~21)} }
\expandafter\def\csname SRef@mapsI:circle-quotient\endcsname{\href{S6_SERIES_GRIND_I.pdf\#nameddest=tag-02OO}{I.22 (p.~295)} }
\expandafter\def\csname SRef@inv:notation-dictionary\endcsname{\href{S6_SERIES_GRIND_I.pdf\#nameddest=tag-01TG}{I.3 (p.~21)} }
\expandafter\def\csname SRef@fam:section:C\endcsname{\href{S6_SERIES_GRIND_I.pdf\#nameddest=tag-01PK}{I.4 (p.~29)} }
\expandafter\def\csname SRef@fam:section:A\endcsname{\href{S6_SERIES_GRIND_I.pdf\#nameddest=tag-01P9}{I.7 (p.~48)} }
\expandafter\def\csname SRef@fam:section:B\endcsname{\href{S6_SERIES_GRIND_I.pdf\#nameddest=tag-01PF}{I.8 (p.~57)} }
\expandafter\def\csname SRef@push:section:E\endcsname{\href{S6_SERIES_GRIND_I.pdf\#nameddest=tag-02ZD}{I.9 (p.~64)} }
\expandafter\def\csname SRef@fam:section:F.1\endcsname{\href{S6_SERIES_GRIND_I.pdf\#nameddest=tag-01PW}{I.10.1 (p.~68)} }
\expandafter\def\csname SRef@fam:section:F.2\endcsname{\href{S6_SERIES_GRIND_I.pdf\#nameddest=tag-01PX}{I.10.2 (p.~71)} }
\expandafter\def\csname SRef@fam:section:F.3\endcsname{\href{S6_SERIES_GRIND_I.pdf\#nameddest=tag-01PY}{I.10.3 (p.~74)} }
\expandafter\def\csname SRef@fam:section:D.1\endcsname{\href{S6_SERIES_GRIND_I.pdf\#nameddest=tag-01PQ}{I.11.1 (p.~76)} }
\expandafter\def\csname SRef@fam:section:D.2\endcsname{\href{S6_SERIES_GRIND_I.pdf\#nameddest=tag-01PR}{I.11.2 (p.~82)} }
\expandafter\def\csname SRef@fam:section:C.3\endcsname{\href{S6_SERIES_GRIND_I.pdf\#nameddest=tag-01PN}{I.11.3 (p.~84)} }
\expandafter\def\csname SRef@fam:section:C.4\endcsname{\href{S6_SERIES_GRIND_I.pdf\#nameddest=tag-01PO}{I.11.4 (p.~86)} }
\expandafter\def\csname SRef@fam:section:E.1\endcsname{\href{S6_SERIES_GRIND_I.pdf\#nameddest=tag-01PT}{I.13.1 (p.~96)} }
\expandafter\def\csname SRef@fam:section:E.2\endcsname{\href{S6_SERIES_GRIND_I.pdf\#nameddest=tag-01PU}{I.13.2 (p.~101)} }
\expandafter\def\csname SRef@push:section:A.1\endcsname{\href{S6_SERIES_GRIND_I.pdf\#nameddest=tag-02Z0}{I.14.1 (p.~103)} }
\expandafter\def\csname SRef@push:section:A.2\endcsname{\href{S6_SERIES_GRIND_I.pdf\#nameddest=tag-02Z1}{I.14.2 (p.~105)} }
\expandafter\def\csname SRef@push:section:A.3\endcsname{\href{S6_SERIES_GRIND_I.pdf\#nameddest=tag-02Z2}{I.14.3 (p.~113)} }
\expandafter\def\csname SRef@push:section:A.4\endcsname{\href{S6_SERIES_GRIND_I.pdf\#nameddest=tag-02Z3}{I.14.4 (p.~116)} }
\expandafter\def\csname SRef@push:section:B.1\endcsname{\href{S6_SERIES_GRIND_I.pdf\#nameddest=tag-02Z5}{I.15.1 (p.~128)} }
\expandafter\def\csname SRef@push:section:B.2\endcsname{\href{S6_SERIES_GRIND_I.pdf\#nameddest=tag-02Z6}{I.15.2 (p.~131)} }
\expandafter\def\csname SRef@push:section:B.3\endcsname{\href{S6_SERIES_GRIND_I.pdf\#nameddest=tag-02Z7}{I.15.3 (p.~133)} }
\expandafter\def\csname SRef@push:section:D.1\endcsname{\href{S6_SERIES_GRIND_I.pdf\#nameddest=tag-02Z9}{I.16.1 (p.~227)} }
\expandafter\def\csname SRef@push:section:D.2\endcsname{\href{S6_SERIES_GRIND_I.pdf\#nameddest=tag-02ZA}{I.16.2 (p.~230)} }
\expandafter\def\csname SRef@push:section:D.3\endcsname{\href{S6_SERIES_GRIND_I.pdf\#nameddest=tag-02ZB}{I.16.3 (p.~233)} }
\expandafter\def\csname SRef@push:section:D.4\endcsname{\href{S6_SERIES_GRIND_I.pdf\#nameddest=tag-02ZC}{I.16.4 (p.~235)} }
\expandafter\def\csname SRef@push:section:E.2\endcsname{\href{S6_SERIES_GRIND_I.pdf\#nameddest=tag-02ZF}{I.20.2 (p.~282)} }
\expandafter\def\csname SRef@push:section:E.3\endcsname{\href{S6_SERIES_GRIND_I.pdf\#nameddest=tag-02ZG}{I.20.3 (p.~285)} }
\expandafter\def\csname SRef@sec:e4\endcsname{\href{S6_SERIES_GRIND_I.pdf\#nameddest=tag-034N}{I.24.1 (p.~338)} }
\expandafter\def\csname SRef@push:section:E.4\endcsname{\href{S6_SERIES_GRIND_I.pdf\#nameddest=tag-034N}{I.24.1 (p.~338)} }
\expandafter\def\csname SRef@fam:section:B.1\endcsname{\href{S6_SERIES_GRIND_I.pdf\#nameddest=tag-01PG}{I.3.1 (p.~24)} }
\expandafter\def\csname SRef@fam:section:C.1\endcsname{\href{S6_SERIES_GRIND_I.pdf\#nameddest=tag-01PL}{I.4.1 (p.~29)} }
\expandafter\def\csname SRef@fam:section:C.2\endcsname{\href{S6_SERIES_GRIND_I.pdf\#nameddest=tag-01PM}{I.4.2 (p.~31)} }
\expandafter\def\csname SRef@fam:section:A.1\endcsname{\href{S6_SERIES_GRIND_I.pdf\#nameddest=tag-01PA}{I.7.1 (p.~49)} }
\expandafter\def\csname SRef@fam:section:A.2\endcsname{\href{S6_SERIES_GRIND_I.pdf\#nameddest=tag-01PB}{I.7.2 (p.~50)} }
\expandafter\def\csname SRef@fam:section:A.3\endcsname{\href{S6_SERIES_GRIND_I.pdf\#nameddest=tag-01PC}{I.7.3 (p.~51)} }
\expandafter\def\csname SRef@fam:section:A.4\endcsname{\href{S6_SERIES_GRIND_I.pdf\#nameddest=tag-01PD}{I.7.4 (p.~53)} }
\expandafter\def\csname SRef@fam:section:A.5\endcsname{\href{S6_SERIES_GRIND_I.pdf\#nameddest=tag-01PE}{I.7.5 (p.~56)} }
\expandafter\def\csname SRef@fam:section:B.2\endcsname{\href{S6_SERIES_GRIND_I.pdf\#nameddest=tag-01PH}{I.8.1 (p.~57)} }
\expandafter\def\csname SRef@fam:section:B.3\endcsname{\href{S6_SERIES_GRIND_I.pdf\#nameddest=tag-01PI}{I.8.2 (p.~60)} }
\expandafter\def\csname SRef@fam:section:B.4\endcsname{\href{S6_SERIES_GRIND_I.pdf\#nameddest=tag-01PJ}{I.8.3 (p.~62)} }
\expandafter\def\csname SRef@push:section:E.1\endcsname{\href{S6_SERIES_GRIND_I.pdf\#nameddest=tag-02ZE}{I.9.1 (p.~64)} }
\expandafter\def\csname SRef@wb:sec:canonical-ring-recovery\endcsname{\href{S6_SERIES_GRIND_I.pdf\#nameddest=tag-03CU}{I.A.1 (p.~342)} }
\expandafter\def\csname SRef@wb:sec:retained-period-deformation\endcsname{\href{S6_SERIES_GRIND_I.pdf\#nameddest=tag-03CV}{I.A.2 (p.~344)} }
\expandafter\def\csname SRef@wb:pshift:sec:classification\endcsname{\href{S6_SERIES_GRIND_I.pdf\#nameddest=tag-03CR}{I.A.3 (p.~349)} }
\expandafter\def\csname SRef@app:rederive:fam-src-A.5.1\endcsname{\href{S6_SERIES_GRIND_I.pdf\#nameddest=tag-000C}{I.A.4 (p.~355)} }
\expandafter\def\csname SRef@app:rederive:fam-src-A.5.2\endcsname{\href{S6_SERIES_GRIND_I.pdf\#nameddest=tag-000D}{I.A.5 (p.~356)} }
\expandafter\def\csname SRef@app:rederive:fam-src-B.4.1\endcsname{\href{S6_SERIES_GRIND_I.pdf\#nameddest=tag-000E}{I.A.6 (p.~359)} }
\expandafter\def\csname SRef@fam:src:F.1.1\endcsname{\href{S6_SERIES_GRIND_I.pdf\#nameddest=tag-01R9}{I.10.1 (p.~68)} }
\expandafter\def\csname SRef@fam:st:F.1.1\endcsname{\href{S6_SERIES_GRIND_I.pdf\#nameddest=tag-01R9}{I.10.1 (p.~68)} }
\expandafter\def\csname SRef@cdp:st:D.1.14\endcsname{\href{S6_SERIES_GRIND_I.pdf\#nameddest=tag-01R9}{I.10.1 (p.~68)} }
\expandafter\def\csname SRef@fam:src:F.2.2\endcsname{\href{S6_SERIES_GRIND_I.pdf\#nameddest=tag-01RC}{I.10.4 (p.~72)} }
\expandafter\def\csname SRef@fam:st:F.2.2\endcsname{\href{S6_SERIES_GRIND_I.pdf\#nameddest=tag-01RC}{I.10.4 (p.~72)} }
\expandafter\def\csname SRef@fam:src:F.3.1\endcsname{\href{S6_SERIES_GRIND_I.pdf\#nameddest=tag-01RE}{I.10.6 (p.~74)} }
\expandafter\def\csname SRef@fam:st:F.3.1\endcsname{\href{S6_SERIES_GRIND_I.pdf\#nameddest=tag-01RE}{I.10.6 (p.~74)} }
\expandafter\def\csname SRef@fam:src:D.2.3\endcsname{\href{S6_SERIES_GRIND_I.pdf\#nameddest=tag-01R0}{I.11.10 (p.~83)} }
\expandafter\def\csname SRef@fam:st:D.2.3\endcsname{\href{S6_SERIES_GRIND_I.pdf\#nameddest=tag-01R0}{I.11.10 (p.~83)} }
\expandafter\def\csname SRef@fam:src:C.4.1\endcsname{\href{S6_SERIES_GRIND_I.pdf\#nameddest=tag-01QQ}{I.11.14 (p.~86)} }
\expandafter\def\csname SRef@fam:st:C.4.1\endcsname{\href{S6_SERIES_GRIND_I.pdf\#nameddest=tag-01QQ}{I.11.14 (p.~86)} }
\expandafter\def\csname SRef@fam:src:D.1.3\endcsname{\href{S6_SERIES_GRIND_I.pdf\#nameddest=tag-01QT}{I.11.3 (p.~78)} }
\expandafter\def\csname SRef@fam:st:D.1.3\endcsname{\href{S6_SERIES_GRIND_I.pdf\#nameddest=tag-01QT}{I.11.3 (p.~78)} }
\expandafter\def\csname SRef@fam:src:D.1.6\endcsname{\href{S6_SERIES_GRIND_I.pdf\#nameddest=tag-01QW}{I.11.6 (p.~80)} }
\expandafter\def\csname SRef@fam:st:D.1.6\endcsname{\href{S6_SERIES_GRIND_I.pdf\#nameddest=tag-01QW}{I.11.6 (p.~80)} }
\expandafter\def\csname SRef@mapsI:src:A.3.2\endcsname{\href{S6_SERIES_GRIND_I.pdf\#nameddest=tag-02PF}{I.12.2 (p.~89)} }
\expandafter\def\csname SRef@mapsI:real-eq-3\endcsname{\href{S6_SERIES_GRIND_I.pdf\#nameddest=tag-02PF}{I.12.2 (p.~89)} }
\expandafter\def\csname SRef@mapsI:src:A.3.3\endcsname{\href{S6_SERIES_GRIND_I.pdf\#nameddest=tag-02PG}{I.12.3 (p.~92)} }
\expandafter\def\csname SRef@push:src:A.4.7\endcsname{\href{S6_SERIES_GRIND_I.pdf\#nameddest=tag-02ZX}{I.14.16 (p.~122)} }
\expandafter\def\csname SRef@push:A47\endcsname{\href{S6_SERIES_GRIND_I.pdf\#nameddest=tag-02ZX}{I.14.16 (p.~122)} }
\expandafter\def\csname SRef@push:src:A.4.8\endcsname{\href{S6_SERIES_GRIND_I.pdf\#nameddest=tag-02ZY}{I.14.17 (p.~123)} }
\expandafter\def\csname SRef@push:A48\endcsname{\href{S6_SERIES_GRIND_I.pdf\#nameddest=tag-02ZY}{I.14.17 (p.~123)} }
\expandafter\def\csname SRef@push:src:A.1.2\endcsname{\href{S6_SERIES_GRIND_I.pdf\#nameddest=tag-02ZI}{I.14.2 (p.~103)} }
\expandafter\def\csname SRef@push:A12\endcsname{\href{S6_SERIES_GRIND_I.pdf\#nameddest=tag-02ZI}{I.14.2 (p.~103)} }
\expandafter\def\csname SRef@push:src:A.2.2\endcsname{\href{S6_SERIES_GRIND_I.pdf\#nameddest=tag-02ZK}{I.14.4 (p.~106)} }
\expandafter\def\csname SRef@push:A22\endcsname{\href{S6_SERIES_GRIND_I.pdf\#nameddest=tag-02ZK}{I.14.4 (p.~106)} }
\expandafter\def\csname SRef@push:src:A.2.4\endcsname{\href{S6_SERIES_GRIND_I.pdf\#nameddest=tag-02ZM}{I.14.5 (p.~108)} }
\expandafter\def\csname SRef@push:A24\endcsname{\href{S6_SERIES_GRIND_I.pdf\#nameddest=tag-02ZM}{I.14.5 (p.~108)} }
\expandafter\def\csname SRef@push:src:A.3.2\endcsname{\href{S6_SERIES_GRIND_I.pdf\#nameddest=tag-02ZP}{I.14.8 (p.~114)} }
\expandafter\def\csname SRef@push:A32\endcsname{\href{S6_SERIES_GRIND_I.pdf\#nameddest=tag-02ZP}{I.14.8 (p.~114)} }
\expandafter\def\csname SRef@push:src:B.4.2\endcsname{\href{S6_SERIES_GRIND_I.pdf\#nameddest=tag-030O}{I.15.15 (p.~143)} }
\expandafter\def\csname SRef@push:B42\endcsname{\href{S6_SERIES_GRIND_I.pdf\#nameddest=tag-030O}{I.15.15 (p.~143)} }
\expandafter\def\csname SRef@push:src:B.4.5\endcsname{\href{S6_SERIES_GRIND_I.pdf\#nameddest=tag-0315}{I.15.18 (p.~147)} }
\expandafter\def\csname SRef@push:src:B.1.2\endcsname{\href{S6_SERIES_GRIND_I.pdf\#nameddest=tag-0301}{I.15.2 (p.~129)} }
\expandafter\def\csname SRef@push:B12\endcsname{\href{S6_SERIES_GRIND_I.pdf\#nameddest=tag-0301}{I.15.2 (p.~129)} }
\expandafter\def\csname SRef@push:src:B.4.33\endcsname{\href{S6_SERIES_GRIND_I.pdf\#nameddest=tag-0313}{I.15.46 (p.~222)} }
\expandafter\def\csname SRef@push:B433\endcsname{\href{S6_SERIES_GRIND_I.pdf\#nameddest=tag-0313}{I.15.46 (p.~222)} }
\expandafter\def\csname SRef@push:src:B.3.1\endcsname{\href{S6_SERIES_GRIND_I.pdf\#nameddest=tag-0305}{I.15.6 (p.~133)} }
\expandafter\def\csname SRef@push:B31\endcsname{\href{S6_SERIES_GRIND_I.pdf\#nameddest=tag-0305}{I.15.6 (p.~133)} }
\expandafter\def\csname SRef@push:src:B.3.4\endcsname{\href{S6_SERIES_GRIND_I.pdf\#nameddest=tag-0308}{I.15.9 (p.~135)} }
\expandafter\def\csname SRef@push:B34\endcsname{\href{S6_SERIES_GRIND_I.pdf\#nameddest=tag-0308}{I.15.9 (p.~135)} }
\expandafter\def\csname SRef@push:src:D.1.3\endcsname{\href{S6_SERIES_GRIND_I.pdf\#nameddest=tag-032C}{I.16.3 (p.~229)} }
\expandafter\def\csname SRef@push:src:D.2.3\endcsname{\href{S6_SERIES_GRIND_I.pdf\#nameddest=tag-032F}{I.16.6 (p.~231)} }
\expandafter\def\csname SRef@push:src:D.3.1\endcsname{\href{S6_SERIES_GRIND_I.pdf\#nameddest=tag-032H}{I.16.8 (p.~233)} }
\expandafter\def\csname SRef@night:P2.1\endcsname{\href{S6_SERIES_GRIND_I.pdf\#nameddest=tag-02R9}{I.17.1 (p.~243)} }
\expandafter\def\csname SRef@stmt:P2.1\endcsname{\href{S6_SERIES_GRIND_I.pdf\#nameddest=tag-02R9}{I.17.1 (p.~243)} }
\expandafter\def\csname SRef@night:P2.10\endcsname{\href{S6_SERIES_GRIND_I.pdf\#nameddest=tag-02RA}{I.17.10 (p.~251)} }
\expandafter\def\csname SRef@stmt:P2.10\endcsname{\href{S6_SERIES_GRIND_I.pdf\#nameddest=tag-02RA}{I.17.10 (p.~251)} }
\expandafter\def\csname SRef@night:P2.11\endcsname{\href{S6_SERIES_GRIND_I.pdf\#nameddest=tag-02RB}{I.17.11 (p.~253)} }
\expandafter\def\csname SRef@stmt:P2.11\endcsname{\href{S6_SERIES_GRIND_I.pdf\#nameddest=tag-02RB}{I.17.11 (p.~253)} }
\expandafter\def\csname SRef@night:P2.2\endcsname{\href{S6_SERIES_GRIND_I.pdf\#nameddest=tag-02RC}{I.17.2 (p.~244)} }
\expandafter\def\csname SRef@stmt:P2.2\endcsname{\href{S6_SERIES_GRIND_I.pdf\#nameddest=tag-02RC}{I.17.2 (p.~244)} }
\expandafter\def\csname SRef@night:P2.3\endcsname{\href{S6_SERIES_GRIND_I.pdf\#nameddest=tag-02RD}{I.17.3 (p.~245)} }
\expandafter\def\csname SRef@stmt:P2.3\endcsname{\href{S6_SERIES_GRIND_I.pdf\#nameddest=tag-02RD}{I.17.3 (p.~245)} }
\expandafter\def\csname SRef@night:P2.4\endcsname{\href{S6_SERIES_GRIND_I.pdf\#nameddest=tag-02RE}{I.17.4 (p.~246)} }
\expandafter\def\csname SRef@stmt:P2.4\endcsname{\href{S6_SERIES_GRIND_I.pdf\#nameddest=tag-02RE}{I.17.4 (p.~246)} }
\expandafter\def\csname SRef@night:P2.5\endcsname{\href{S6_SERIES_GRIND_I.pdf\#nameddest=tag-02RF}{I.17.5 (p.~246)} }
\expandafter\def\csname SRef@stmt:P2.5\endcsname{\href{S6_SERIES_GRIND_I.pdf\#nameddest=tag-02RF}{I.17.5 (p.~246)} }
\expandafter\def\csname SRef@night:P2.6\endcsname{\href{S6_SERIES_GRIND_I.pdf\#nameddest=tag-02RG}{I.17.6 (p.~247)} }
\expandafter\def\csname SRef@stmt:P2.6\endcsname{\href{S6_SERIES_GRIND_I.pdf\#nameddest=tag-02RG}{I.17.6 (p.~247)} }
\expandafter\def\csname SRef@night:P2.7\endcsname{\href{S6_SERIES_GRIND_I.pdf\#nameddest=tag-02RH}{I.17.7 (p.~248)} }
\expandafter\def\csname SRef@stmt:P2.7\endcsname{\href{S6_SERIES_GRIND_I.pdf\#nameddest=tag-02RH}{I.17.7 (p.~248)} }
\expandafter\def\csname SRef@night:P2.8\endcsname{\href{S6_SERIES_GRIND_I.pdf\#nameddest=tag-02RI}{I.17.8 (p.~249)} }
\expandafter\def\csname SRef@stmt:P2.8\endcsname{\href{S6_SERIES_GRIND_I.pdf\#nameddest=tag-02RI}{I.17.8 (p.~249)} }
\expandafter\def\csname SRef@night:P2.9\endcsname{\href{S6_SERIES_GRIND_I.pdf\#nameddest=tag-02RJ}{I.17.9 (p.~249)} }
\expandafter\def\csname SRef@stmt:P2.9\endcsname{\href{S6_SERIES_GRIND_I.pdf\#nameddest=tag-02RJ}{I.17.9 (p.~249)} }
\expandafter\def\csname SRef@push:src:E.2.1\endcsname{\href{S6_SERIES_GRIND_I.pdf\#nameddest=tag-032U}{I.20.6 (p.~282)} }
\expandafter\def\csname SRef@mapsI:src:A.2.1\endcsname{\href{S6_SERIES_GRIND_I.pdf\#nameddest=tag-02PD}{I.22.4 (p.~298)} }
\expandafter\def\csname SRef@push:src:E.4.1\endcsname{\href{S6_SERIES_GRIND_I.pdf\#nameddest=tag-033L}{I.24.5 (p.~338)} }
\expandafter\def\csname SRef@push:src:E.4.2\endcsname{\href{S6_SERIES_GRIND_I.pdf\#nameddest=tag-033M}{I.24.6 (p.~341)} }
\expandafter\def\csname SRef@fam:src:C.2.3\endcsname{\href{S6_SERIES_GRIND_I.pdf\#nameddest=tag-01QL}{I.4.5 (p.~34)} }
\expandafter\def\csname SRef@ari:C.8.5\endcsname{\href{S6_SERIES_GRIND_I.pdf\#nameddest=tag-01QL}{I.4.5 (p.~34)} }
\expandafter\def\csname SRef@fam:st:C.2.3\endcsname{\href{S6_SERIES_GRIND_I.pdf\#nameddest=tag-01QL}{I.4.5 (p.~34)} }
\expandafter\def\csname SRef@push:src:E.4.3\endcsname{\href{S6_SERIES_GRIND_I.pdf\#nameddest=tag-01QL}{I.4.5 (p.~34)} }
\expandafter\def\csname SRef@fam:src:A.5.1\endcsname{\href{S6_SERIES_PUBLISH_I.pdf\#nameddest=tag-01Q6}{I.7.8 (p.~51)} }
\expandafter\def\csname SRef@wb:pshift:thm:classification\endcsname{\href{S6_SERIES_GRIND_I.pdf\#nameddest=tag-01Q6}{I.7.8 (p.~56)} }
\expandafter\def\csname SRef@fam:st:A.5.1\endcsname{\href{S6_SERIES_GRIND_I.pdf\#nameddest=tag-01Q6}{I.7.8 (p.~56)} }
\expandafter\def\csname SRef@fam:st:A.5.2\endcsname{\href{S6_SERIES_GRIND_I.pdf\#nameddest=tag-01RH}{I.7.9 (p.~56)} }
\expandafter\def\csname SRef@fam:src:A.5.2\endcsname{\href{S6_SERIES_GRIND_I.pdf\#nameddest=tag-01RH}{I.7.9 (p.~56)} }
\expandafter\def\csname SRef@inv:D1.3\endcsname{\href{S6_SERIES_GRIND_I.pdf\#nameddest=tag-01RH}{I.7.9 (p.~56)} }
\expandafter\def\csname SRef@inv:D2.1\endcsname{\href{S6_SERIES_GRIND_I.pdf\#nameddest=tag-01RH}{I.7.9 (p.~56)} }
\expandafter\def\csname SRef@inv:D3.1\endcsname{\href{S6_SERIES_GRIND_I.pdf\#nameddest=tag-01RH}{I.7.9 (p.~56)} }
\expandafter\def\csname SRef@cdp:st:E9\endcsname{\href{S6_SERIES_GRIND_I.pdf\#nameddest=tag-01RH}{I.7.9 (p.~56)} }
\expandafter\def\csname SRef@push:src:E.1.2\endcsname{\href{S6_SERIES_GRIND_I.pdf\#nameddest=tag-032S}{I.9.2 (p.~65)} }
\expandafter\def\csname SRef@wb:thm:canonical-ring-recovery\endcsname{\href{S6_SERIES_GRIND_I.pdf\#nameddest=tag-03DM}{I.A.2 (p.~342)} }
\expandafter\def\csname SRef@wb:thm:retained-period-deformation\endcsname{\href{S6_SERIES_GRIND_I.pdf\#nameddest=tag-03DP}{I.A.3 (p.~344)} }
\expandafter\def\csname SRef@night:P5.1\endcsname{\href{S6_SERIES_GRIND_I.pdf\#nameddest=tag-02YM}{I.C.1 (p.~366)} }
\expandafter\def\csname SRef@stmt:P5.1\endcsname{\href{S6_SERIES_GRIND_I.pdf\#nameddest=tag-02YM}{I.C.1 (p.~366)} }
\expandafter\def\csname SRef@night:P5.2\endcsname{\href{S6_SERIES_GRIND_I.pdf\#nameddest=tag-02YN}{I.C.2 (p.~367)} }
\expandafter\def\csname SRef@stmt:P5.2\endcsname{\href{S6_SERIES_GRIND_I.pdf\#nameddest=tag-02YN}{I.C.2 (p.~367)} }
\expandafter\def\csname SRef@night:P5.3\endcsname{\href{S6_SERIES_GRIND_I.pdf\#nameddest=tag-02YO}{I.C.3 (p.~369)} }
\expandafter\def\csname SRef@stmt:P5.3\endcsname{\href{S6_SERIES_GRIND_I.pdf\#nameddest=tag-02YO}{I.C.3 (p.~369)} }
\expandafter\def\csname SRef@night:P5.4\endcsname{\href{S6_SERIES_GRIND_I.pdf\#nameddest=tag-02YP}{I.C.4 (p.~369)} }
\expandafter\def\csname SRef@stmt:P5.4\endcsname{\href{S6_SERIES_GRIND_I.pdf\#nameddest=tag-02YP}{I.C.4 (p.~369)} }
\expandafter\def\csname SRef@night:P5.5\endcsname{\href{S6_SERIES_GRIND_I.pdf\#nameddest=tag-02YQ}{I.C.5 (p.~370)} }
\expandafter\def\csname SRef@stmt:P5.5\endcsname{\href{S6_SERIES_GRIND_I.pdf\#nameddest=tag-02YQ}{I.C.5 (p.~370)} }
\expandafter\def\csname SRef@night:P5.6\endcsname{\href{S6_SERIES_GRIND_I.pdf\#nameddest=tag-02YR}{I.C.6 (p.~372)} }
\expandafter\def\csname SRef@stmt:P5.6\endcsname{\href{S6_SERIES_GRIND_I.pdf\#nameddest=tag-02YR}{I.C.6 (p.~372)} }
\expandafter\def\csname SRef@cdp:eq:A.RC.2\endcsname{\hyperlink{tag-006V}{II.\ref*{cdp:eq:A.RC.2}}}
\expandafter\def\csname SRef@cdp:eq:A.RC.3\endcsname{\hyperlink{tag-006W}{II.\ref*{cdp:eq:A.RC.3}}}
\expandafter\def\csname SRef@cdp:eq:A.RC.4\endcsname{\hyperlink{tag-006X}{II.\ref*{cdp:eq:A.RC.4}}}
\expandafter\def\csname SRef@cdp:eq:A.RC.5\endcsname{\hyperlink{tag-006Y}{II.\ref*{cdp:eq:A.RC.5}}}
\expandafter\def\csname SRef@cdp:eq:A.RC.6\endcsname{\hyperlink{tag-006Z}{II.\ref*{cdp:eq:A.RC.6}}}
\expandafter\def\csname SRef@cdp:eq:A.RC.7\endcsname{\hyperlink{tag-0070}{II.\ref*{cdp:eq:A.RC.7}}}
\expandafter\def\csname SRef@cdp:eq:A.Pi.1\endcsname{\hyperlink{tag-006O}{II.\ref*{cdp:eq:A.Pi.1}}}
\expandafter\def\csname SRef@cdp:eq:A.Pi.2\endcsname{\hyperlink{tag-006P}{II.\ref*{cdp:eq:A.Pi.2}}}
\expandafter\def\csname SRef@cdp:eq:A.Pi.3\endcsname{\hyperlink{tag-006Q}{II.\ref*{cdp:eq:A.Pi.3}}}
\expandafter\def\csname SRef@cdp:eq:A.Pi.4\endcsname{\hyperlink{tag-006R}{II.\ref*{cdp:eq:A.Pi.4}}}
\expandafter\def\csname SRef@cdp:eq:A.CD.1\endcsname{\hyperlink{tag-0058}{II.\ref*{cdp:eq:A.CD.1}}}
\expandafter\def\csname SRef@cdp:eq:A.CD.2\endcsname{\hyperlink{tag-0059}{II.\ref*{cdp:eq:A.CD.2}}}
\expandafter\def\csname SRef@cdp:eq:A.CD.3\endcsname{\hyperlink{tag-005A}{II.\ref*{cdp:eq:A.CD.3}}}
\expandafter\def\csname SRef@cdp:eq:A.CD.4\endcsname{\hyperlink{tag-005B}{II.\ref*{cdp:eq:A.CD.4}}}
\expandafter\def\csname SRef@cdp:eq:A.Q1.1\endcsname{\hyperlink{tag-006S}{II.\ref*{cdp:eq:A.Q1.1}}}
\expandafter\def\csname SRef@cdp:eq:A.Q1.2\endcsname{\hyperlink{tag-006T}{II.\ref*{cdp:eq:A.Q1.2}}}
\expandafter\def\csname SRef@cdp:eq:A.J.1\endcsname{\hyperlink{tag-005C}{II.\ref*{cdp:eq:A.J.1}}}
\expandafter\def\csname SRef@cdp:eq:A.J.2\endcsname{\hyperlink{tag-005N}{II.\ref*{cdp:eq:A.J.2}}}
\expandafter\def\csname SRef@cdp:eq:A.J.3\endcsname{\hyperlink{tag-005Y}{II.\ref*{cdp:eq:A.J.3}}}
\expandafter\def\csname SRef@cdp:eq:A.J.4\endcsname{\hyperlink{tag-006C}{II.\ref*{cdp:eq:A.J.4}}}
\expandafter\def\csname SRef@cdp:eq:A.J.5\endcsname{\hyperlink{tag-006E}{II.\ref*{cdp:eq:A.J.5}}}
\expandafter\def\csname SRef@cdp:eq:A.J.6\endcsname{\hyperlink{tag-006F}{II.\ref*{cdp:eq:A.J.6}}}
\expandafter\def\csname SRef@cdp:eq:A.J.7\endcsname{\hyperlink{tag-006G}{II.\ref*{cdp:eq:A.J.7}}}
\expandafter\def\csname SRef@cdp:eq:A.J.8\endcsname{\hyperlink{tag-006H}{II.\ref*{cdp:eq:A.J.8}}}
\expandafter\def\csname SRef@cdp:eq:A.J.9\endcsname{\hyperlink{tag-006I}{II.\ref*{cdp:eq:A.J.9}}}
\expandafter\def\csname SRef@cdp:eq:A.J.10\endcsname{\hyperlink{tag-005D}{II.\ref*{cdp:eq:A.J.10}}}
\expandafter\def\csname SRef@cdp:eq:A.J.11\endcsname{\hyperlink{tag-005E}{II.\ref*{cdp:eq:A.J.11}}}
\expandafter\def\csname SRef@cdp:eq:A.J.12\endcsname{\hyperlink{tag-005F}{II.\ref*{cdp:eq:A.J.12}}}
\expandafter\def\csname SRef@cdp:eq:A.J.13\endcsname{\hyperlink{tag-005G}{II.\ref*{cdp:eq:A.J.13}}}
\expandafter\def\csname SRef@cdp:eq:A.J.14\endcsname{\hyperlink{tag-005H}{II.\ref*{cdp:eq:A.J.14}}}
\expandafter\def\csname SRef@cdp:eq:A.J.15\endcsname{\hyperlink{tag-005I}{II.\ref*{cdp:eq:A.J.15}}}
\expandafter\def\csname SRef@cdp:eq:A.J.16\endcsname{\hyperlink{tag-005J}{II.\ref*{cdp:eq:A.J.16}}}
\expandafter\def\csname SRef@cdp:eq:A.J.17\endcsname{\hyperlink{tag-005K}{II.\ref*{cdp:eq:A.J.17}}}
\expandafter\def\csname SRef@cdp:eq:A.J.18\endcsname{\hyperlink{tag-005L}{II.\ref*{cdp:eq:A.J.18}}}
\expandafter\def\csname SRef@cdp:eq:A.J.19\endcsname{\hyperlink{tag-005M}{II.\ref*{cdp:eq:A.J.19}}}
\expandafter\def\csname SRef@cdp:eq:A.J.20\endcsname{\hyperlink{tag-005O}{II.\ref*{cdp:eq:A.J.20}}}
\expandafter\def\csname SRef@cdp:eq:A.J.30\endcsname{\hyperlink{tag-005Z}{II.\ref*{cdp:eq:A.J.30}}}
\expandafter\def\csname SRef@cdp:eq:A.J.31\endcsname{\hyperlink{tag-0060}{II.\ref*{cdp:eq:A.J.31}}}
\expandafter\def\csname SRef@cdp:eq:A.J.32\endcsname{\hyperlink{tag-0061}{II.\ref*{cdp:eq:A.J.32}}}
\expandafter\def\csname SRef@cdp:eq:A.J.21\endcsname{\hyperlink{tag-005P}{II.\ref*{cdp:eq:A.J.21}}}
\expandafter\def\csname SRef@cdp:eq:A.J.22\endcsname{\hyperlink{tag-005Q}{II.\ref*{cdp:eq:A.J.22}}}
\expandafter\def\csname SRef@cdp:eq:A.J.23\endcsname{\hyperlink{tag-005R}{II.\ref*{cdp:eq:A.J.23}}}
\expandafter\def\csname SRef@cdp:eq:A.J.24\endcsname{\hyperlink{tag-005S}{II.\ref*{cdp:eq:A.J.24}}}
\expandafter\def\csname SRef@cdp:eq:A.J.25\endcsname{\hyperlink{tag-005T}{II.\ref*{cdp:eq:A.J.25}}}
\expandafter\def\csname SRef@cdp:eq:A.J.26\endcsname{\hyperlink{tag-005U}{II.\ref*{cdp:eq:A.J.26}}}
\expandafter\def\csname SRef@cdp:eq:A.J.27\endcsname{\hyperlink{tag-005V}{II.\ref*{cdp:eq:A.J.27}}}
\expandafter\def\csname SRef@cdp:eq:A.J.28\endcsname{\hyperlink{tag-005W}{II.\ref*{cdp:eq:A.J.28}}}
\expandafter\def\csname SRef@cdp:eq:A.J.29\endcsname{\hyperlink{tag-005X}{II.\ref*{cdp:eq:A.J.29}}}
\expandafter\def\csname SRef@cdp:eq:A.J.33\endcsname{\hyperlink{tag-0062}{II.\ref*{cdp:eq:A.J.33}}}
\expandafter\def\csname SRef@cdp:eq:A.J.34\endcsname{\hyperlink{tag-0065}{II.\ref*{cdp:eq:A.J.34}}}
\expandafter\def\csname SRef@cdp:eq:A.J.33a\endcsname{\hyperlink{tag-0063}{II.\ref*{cdp:eq:A.J.33a}}}
\expandafter\def\csname SRef@cdp:eq:A.J.33b\endcsname{\hyperlink{tag-0064}{II.\ref*{cdp:eq:A.J.33b}}}
\expandafter\def\csname SRef@cdp:eq:A.J.35\endcsname{\hyperlink{tag-0066}{II.\ref*{cdp:eq:A.J.35}}}
\expandafter\def\csname SRef@cdp:eq:A.J.36\endcsname{\hyperlink{tag-0067}{II.\ref*{cdp:eq:A.J.36}}}
\expandafter\def\csname SRef@cdp:eq:A.J.37\endcsname{\hyperlink{tag-0068}{II.\ref*{cdp:eq:A.J.37}}}
\expandafter\def\csname SRef@cdp:eq:A.J.38\endcsname{\hyperlink{tag-0069}{II.\ref*{cdp:eq:A.J.38}}}
\expandafter\def\csname SRef@cdp:eq:A.J.38a\endcsname{\hyperlink{tag-006A}{II.\ref*{cdp:eq:A.J.38a}}}
\expandafter\def\csname SRef@cdp:eq:A.J.39\endcsname{\hyperlink{tag-006B}{II.\ref*{cdp:eq:A.J.39}}}
\expandafter\def\csname SRef@cdp:eq:A.J.40\endcsname{\hyperlink{tag-006D}{II.\ref*{cdp:eq:A.J.40}}}
\expandafter\def\csname SRef@cdp:eq:A.T.1\endcsname{\hyperlink{tag-007M}{II.\ref*{cdp:eq:A.T.1}}}
\expandafter\def\csname SRef@cdp:eq:A.T.2\endcsname{\hyperlink{tag-007X}{II.\ref*{cdp:eq:A.T.2}}}
\expandafter\def\csname SRef@cdp:eq:A.T.3\endcsname{\hyperlink{tag-0082}{II.\ref*{cdp:eq:A.T.3}}}
\expandafter\def\csname SRef@cdp:eq:A.T.4\endcsname{\hyperlink{tag-0083}{II.\ref*{cdp:eq:A.T.4}}}
\expandafter\def\csname SRef@cdp:eq:A.T.5\endcsname{\hyperlink{tag-0084}{II.\ref*{cdp:eq:A.T.5}}}
\expandafter\def\csname SRef@cdp:eq:A.T.6\endcsname{\hyperlink{tag-0085}{II.\ref*{cdp:eq:A.T.6}}}
\expandafter\def\csname SRef@cdp:eq:A.T.7\endcsname{\hyperlink{tag-0086}{II.\ref*{cdp:eq:A.T.7}}}
\expandafter\def\csname SRef@cdp:eq:A.T.8\endcsname{\hyperlink{tag-0087}{II.\ref*{cdp:eq:A.T.8}}}
\expandafter\def\csname SRef@cdp:eq:A.T.9\endcsname{\hyperlink{tag-0088}{II.\ref*{cdp:eq:A.T.9}}}
\expandafter\def\csname SRef@cdp:eq:A.T.10\endcsname{\hyperlink{tag-007N}{II.\ref*{cdp:eq:A.T.10}}}
\expandafter\def\csname SRef@cdp:eq:A.T.11\endcsname{\hyperlink{tag-007O}{II.\ref*{cdp:eq:A.T.11}}}
\expandafter\def\csname SRef@cdp:eq:A.T.12\endcsname{\hyperlink{tag-007P}{II.\ref*{cdp:eq:A.T.12}}}
\expandafter\def\csname SRef@cdp:eq:A.T.13\endcsname{\hyperlink{tag-007Q}{II.\ref*{cdp:eq:A.T.13}}}
\expandafter\def\csname SRef@cdp:eq:A.T.14\endcsname{\hyperlink{tag-007R}{II.\ref*{cdp:eq:A.T.14}}}
\expandafter\def\csname SRef@cdp:eq:A.T.15\endcsname{\hyperlink{tag-007S}{II.\ref*{cdp:eq:A.T.15}}}
\expandafter\def\csname SRef@cdp:eq:A.T.16\endcsname{\hyperlink{tag-007T}{II.\ref*{cdp:eq:A.T.16}}}
\expandafter\def\csname SRef@cdp:eq:A.T.17\endcsname{\hyperlink{tag-007U}{II.\ref*{cdp:eq:A.T.17}}}
\expandafter\def\csname SRef@cdp:eq:A.T.18\endcsname{\hyperlink{tag-007V}{II.\ref*{cdp:eq:A.T.18}}}
\expandafter\def\csname SRef@cdp:eq:A.T.19\endcsname{\hyperlink{tag-007W}{II.\ref*{cdp:eq:A.T.19}}}
\expandafter\def\csname SRef@cdp:eq:A.T.20\endcsname{\hyperlink{tag-007Y}{II.\ref*{cdp:eq:A.T.20}}}
\expandafter\def\csname SRef@cdp:eq:A.T.21\endcsname{\hyperlink{tag-007Z}{II.\ref*{cdp:eq:A.T.21}}}
\expandafter\def\csname SRef@cdp:eq:A.T.22\endcsname{\hyperlink{tag-0080}{II.\ref*{cdp:eq:A.T.22}}}
\expandafter\def\csname SRef@cdp:eq:A.T.23\endcsname{\hyperlink{tag-0081}{II.\ref*{cdp:eq:A.T.23}}}
\expandafter\def\csname SRef@cdp:eq:A.U.1\endcsname{\hyperlink{tag-0089}{II.\ref*{cdp:eq:A.U.1}}}
\expandafter\def\csname SRef@cdp:eq:A.U.2\endcsname{\hyperlink{tag-008A}{II.\ref*{cdp:eq:A.U.2}}}
\expandafter\def\csname SRef@cdp:eq:A.U.3\endcsname{\hyperlink{tag-008B}{II.\ref*{cdp:eq:A.U.3}}}
\expandafter\def\csname SRef@cdp:eq:A.U.4\endcsname{\hyperlink{tag-008C}{II.\ref*{cdp:eq:A.U.4}}}
\expandafter\def\csname SRef@cdp:eq:A.U.5\endcsname{\hyperlink{tag-008D}{II.\ref*{cdp:eq:A.U.5}}}
\expandafter\def\csname SRef@cdp:eq:A.U.6\endcsname{\hyperlink{tag-008E}{II.\ref*{cdp:eq:A.U.6}}}
\expandafter\def\csname SRef@cdp:eq:A.U.7\endcsname{\hyperlink{tag-008F}{II.\ref*{cdp:eq:A.U.7}}}
\expandafter\def\csname SRef@cdp:eq:A.RT.1\endcsname{\hyperlink{tag-0071}{II.\ref*{cdp:eq:A.RT.1}}}
\expandafter\def\csname SRef@cdp:eq:A.RT.2\endcsname{\hyperlink{tag-007C}{II.\ref*{cdp:eq:A.RT.2}}}
\expandafter\def\csname SRef@cdp:eq:A.RT.3\endcsname{\hyperlink{tag-007F}{II.\ref*{cdp:eq:A.RT.3}}}
\expandafter\def\csname SRef@cdp:eq:A.RT.4\endcsname{\hyperlink{tag-007G}{II.\ref*{cdp:eq:A.RT.4}}}
\expandafter\def\csname SRef@cdp:eq:A.RT.5\endcsname{\hyperlink{tag-007H}{II.\ref*{cdp:eq:A.RT.5}}}
\expandafter\def\csname SRef@cdp:eq:A.RT.6\endcsname{\hyperlink{tag-007I}{II.\ref*{cdp:eq:A.RT.6}}}
\expandafter\def\csname SRef@cdp:eq:A.RT.7\endcsname{\hyperlink{tag-007J}{II.\ref*{cdp:eq:A.RT.7}}}
\expandafter\def\csname SRef@cdp:eq:A.RT.8\endcsname{\hyperlink{tag-007K}{II.\ref*{cdp:eq:A.RT.8}}}
\expandafter\def\csname SRef@cdp:eq:A.RT.9\endcsname{\hyperlink{tag-007L}{II.\ref*{cdp:eq:A.RT.9}}}
\expandafter\def\csname SRef@cdp:eq:A.RT.10\endcsname{\hyperlink{tag-0072}{II.\ref*{cdp:eq:A.RT.10}}}
\expandafter\def\csname SRef@cdp:eq:A.RT.11\endcsname{\hyperlink{tag-0073}{II.\ref*{cdp:eq:A.RT.11}}}
\expandafter\def\csname SRef@cdp:eq:A.RT.12\endcsname{\hyperlink{tag-0074}{II.\ref*{cdp:eq:A.RT.12}}}
\expandafter\def\csname SRef@cdp:eq:A.RT.13\endcsname{\hyperlink{tag-0075}{II.\ref*{cdp:eq:A.RT.13}}}
\expandafter\def\csname SRef@cdp:eq:A.RT.14\endcsname{\hyperlink{tag-0076}{II.\ref*{cdp:eq:A.RT.14}}}
\expandafter\def\csname SRef@cdp:eq:A.RT.15\endcsname{\hyperlink{tag-0077}{II.\ref*{cdp:eq:A.RT.15}}}
\expandafter\def\csname SRef@cdp:eq:A.RT.16\endcsname{\hyperlink{tag-0078}{II.\ref*{cdp:eq:A.RT.16}}}
\expandafter\def\csname SRef@cdp:eq:A.RT.17\endcsname{\hyperlink{tag-0079}{II.\ref*{cdp:eq:A.RT.17}}}
\expandafter\def\csname SRef@cdp:eq:A.RT.18\endcsname{\hyperlink{tag-007A}{II.\ref*{cdp:eq:A.RT.18}}}
\expandafter\def\csname SRef@cdp:eq:A.RT.19\endcsname{\hyperlink{tag-007B}{II.\ref*{cdp:eq:A.RT.19}}}
\expandafter\def\csname SRef@cdp:eq:A.RT.20\endcsname{\hyperlink{tag-007D}{II.\ref*{cdp:eq:A.RT.20}}}
\expandafter\def\csname SRef@cdp:eq:A.RT.21\endcsname{\hyperlink{tag-007E}{II.\ref*{cdp:eq:A.RT.21}}}
\expandafter\def\csname SRef@cdp:eq:A.BP.1\endcsname{\hyperlink{tag-004W}{II.\ref*{cdp:eq:A.BP.1}}}
\expandafter\def\csname SRef@cdp:eq:A.BP.4\endcsname{\hyperlink{tag-0052}{II.\ref*{cdp:eq:A.BP.4}}}
\expandafter\def\csname SRef@cdp:eq:A.BP.2\endcsname{\hyperlink{tag-0050}{II.\ref*{cdp:eq:A.BP.2}}}
\expandafter\def\csname SRef@cdp:eq:A.BP.3\endcsname{\hyperlink{tag-0051}{II.\ref*{cdp:eq:A.BP.3}}}
\expandafter\def\csname SRef@cdp:eq:A.BP.5\endcsname{\hyperlink{tag-0053}{II.\ref*{cdp:eq:A.BP.5}}}
\expandafter\def\csname SRef@cdp:eq:A.BP.6\endcsname{\hyperlink{tag-0054}{II.\ref*{cdp:eq:A.BP.6}}}
\expandafter\def\csname SRef@cdp:eq:A.BP.7\endcsname{\hyperlink{tag-0055}{II.\ref*{cdp:eq:A.BP.7}}}
\expandafter\def\csname SRef@cdp:eq:A.BP.8\endcsname{\hyperlink{tag-0056}{II.\ref*{cdp:eq:A.BP.8}}}
\expandafter\def\csname SRef@cdp:eq:A.BP.9\endcsname{\hyperlink{tag-0057}{II.\ref*{cdp:eq:A.BP.9}}}
\expandafter\def\csname SRef@cdp:eq:A.BP.10\endcsname{\hyperlink{tag-004X}{II.\ref*{cdp:eq:A.BP.10}}}
\expandafter\def\csname SRef@cdp:eq:A.BP.11\endcsname{\hyperlink{tag-004Y}{II.\ref*{cdp:eq:A.BP.11}}}
\expandafter\def\csname SRef@cdp:eq:A.BP.12\endcsname{\hyperlink{tag-004Z}{II.\ref*{cdp:eq:A.BP.12}}}
\expandafter\def\csname SRef@cdp:eq:C.NC.4\endcsname{\hyperlink{tag-00S3}{II.\ref*{cdp:eq:C.NC.4}}}
\expandafter\def\csname SRef@cdp:eq:C.NC.5\endcsname{\hyperlink{tag-00S4}{II.\ref*{cdp:eq:C.NC.5}}}
\expandafter\def\csname SRef@cdp:eq:C.NC.6\endcsname{\hyperlink{tag-00S5}{II.\ref*{cdp:eq:C.NC.6}}}
\expandafter\def\csname SRef@cdp:eq:C.NC.7\endcsname{\hyperlink{tag-00S6}{II.\ref*{cdp:eq:C.NC.7}}}
\expandafter\def\csname SRef@cdp:eq:C.NC.8\endcsname{\hyperlink{tag-00S7}{II.\ref*{cdp:eq:C.NC.8}}}
\expandafter\def\csname SRef@cdp:eq:C.NC.9\endcsname{\hyperlink{tag-00S8}{II.\ref*{cdp:eq:C.NC.9}}}
\expandafter\def\csname SRef@cdp:eq:C.NC.10\endcsname{\hyperlink{tag-00S1}{II.\ref*{cdp:eq:C.NC.10}}}
\expandafter\def\csname SRef@cdp:eq:C.NC.11\endcsname{\hyperlink{tag-00S2}{II.\ref*{cdp:eq:C.NC.11}}}
\expandafter\def\csname SRef@cdp:eq:C.J.1\endcsname{\hyperlink{tag-00OV}{II.\ref*{cdp:eq:C.J.1}}}
\expandafter\def\csname SRef@cdp:eq:C.J.2\endcsname{\hyperlink{tag-00OW}{II.\ref*{cdp:eq:C.J.2}}}
\expandafter\def\csname SRef@cdp:eq:C.J.3\endcsname{\hyperlink{tag-00OX}{II.\ref*{cdp:eq:C.J.3}}}
\expandafter\def\csname SRef@cdp:eq:D1\endcsname{\hyperlink{tag-019T}{II.\ref*{cdp:eq:D1}}}
\expandafter\def\csname SRef@cdp:eq:D5\endcsname{\hyperlink{tag-019X}{II.\ref*{cdp:eq:D5}}}
\expandafter\def\csname SRef@cdp:eq:D4\endcsname{\hyperlink{tag-019W}{II.\ref*{cdp:eq:D4}}}
\expandafter\def\csname SRef@cdp:eq:D2\endcsname{\hyperlink{tag-019U}{II.\ref*{cdp:eq:D2}}}
\expandafter\def\csname SRef@cdp:eq:D3\endcsname{\hyperlink{tag-019V}{II.\ref*{cdp:eq:D3}}}
\expandafter\def\csname SRef@cdp:eq:C.SK.1\endcsname{\hyperlink{tag-011H}{II.\ref*{cdp:eq:C.SK.1}}}
\expandafter\def\csname SRef@cdp:eq:C.SK.2\endcsname{\hyperlink{tag-011I}{II.\ref*{cdp:eq:C.SK.2}}}
\expandafter\def\csname SRef@a.-one-dimensional-first-jets\endcsname{\hyperlink{tag-0005}{II.\ref*{a.-one-dimensional-first-jets}}}
\expandafter\def\csname SRef@three-smooth-branches-the-complete-excess-two-classification\endcsname{\hyperlink{tag-0005}{II.\ref*{three-smooth-branches-the-complete-excess-two-classification}}}
\expandafter\def\csname SRef@cdp:eq:C.SK.3\endcsname{\hyperlink{tag-011J}{II.\ref*{cdp:eq:C.SK.3}}}
\expandafter\def\csname SRef@b.-two-dimensional-first-jets\endcsname{\hyperlink{tag-003V}{II.\ref*{b.-two-dimensional-first-jets}}}
\expandafter\def\csname SRef@multiplicity-four-one-cusp-and-two-smooth-branches\endcsname{\hyperlink{tag-02QH}{II.\ref*{multiplicity-four-one-cusp-and-two-smooth-branches}}}
\expandafter\def\csname SRef@a.-a-nonzero-coefficient-of-the-cusp-cubic\endcsname{\hyperlink{tag-0004}{II.\ref*{a.-a-nonzero-coefficient-of-the-cusp-cubic}}}
\expandafter\def\csname SRef@b.-a-zero-coefficient-of-the-cusp-cubic\endcsname{\hyperlink{tag-003U}{II.\ref*{b.-a-zero-coefficient-of-the-cusp-cubic}}}
\expandafter\def\csname SRef@cdp:eq:C.SK.4\endcsname{\hyperlink{tag-011K}{II.\ref*{cdp:eq:C.SK.4}}}
\expandafter\def\csname SRef@every-nonordinary-triple-exceeds-the-budget\endcsname{\hyperlink{tag-01OO}{II.\ref*{every-nonordinary-triple-exceeds-the-budget}}}
\expandafter\def\csname SRef@cdp:eq:C.SK.5\endcsname{\hyperlink{tag-011L}{II.\ref*{cdp:eq:C.SK.5}}}
\expandafter\def\csname SRef@cdp:eq:C.15\endcsname{\hyperlink{tag-008T}{II.\ref*{cdp:eq:C.15}}}
\expandafter\def\csname SRef@cdp:eq:C.16\endcsname{\hyperlink{tag-008U}{II.\ref*{cdp:eq:C.16}}}
\expandafter\def\csname SRef@cdp:eq:C.17\endcsname{\hyperlink{tag-008V}{II.\ref*{cdp:eq:C.17}}}
\expandafter\def\csname SRef@cdp:eq:C.18\endcsname{\hyperlink{tag-008W}{II.\ref*{cdp:eq:C.18}}}
\expandafter\def\csname SRef@cdp:eq:C.SK.6\endcsname{\hyperlink{tag-011M}{II.\ref*{cdp:eq:C.SK.6}}}
\expandafter\def\csname SRef@cdp:eq:C.19\endcsname{\hyperlink{tag-008X}{II.\ref*{cdp:eq:C.19}}}
\expandafter\def\csname SRef@cdp:eq:C.20\endcsname{\hyperlink{tag-008Z}{II.\ref*{cdp:eq:C.20}}}
\expandafter\def\csname SRef@cdp:eq:C.21\endcsname{\hyperlink{tag-0090}{II.\ref*{cdp:eq:C.21}}}
\expandafter\def\csname SRef@cdp:eq:C.22\endcsname{\hyperlink{tag-0091}{II.\ref*{cdp:eq:C.22}}}
\expandafter\def\csname SRef@cdp:eq:C.22b\endcsname{\hyperlink{tag-0093}{II.\ref*{cdp:eq:C.22b}}}
\expandafter\def\csname SRef@cdp:eq:C.22a\endcsname{\hyperlink{tag-0092}{II.\ref*{cdp:eq:C.22a}}}
\expandafter\def\csname SRef@cdp:eq:C.VC.1\endcsname{\hyperlink{tag-015G}{II.\ref*{cdp:eq:C.VC.1}}}
\expandafter\def\csname SRef@cdp:eq:C.VC.2\endcsname{\hyperlink{tag-015H}{II.\ref*{cdp:eq:C.VC.2}}}
\expandafter\def\csname SRef@cdp:eq:C.VC.3\endcsname{\hyperlink{tag-015I}{II.\ref*{cdp:eq:C.VC.3}}}
\expandafter\def\csname SRef@cdp:eq:C.VC.4\endcsname{\hyperlink{tag-015J}{II.\ref*{cdp:eq:C.VC.4}}}
\expandafter\def\csname SRef@cdp:eq:C.PK.1\endcsname{\hyperlink{tag-00V8}{II.\ref*{cdp:eq:C.PK.1}}}
\expandafter\def\csname SRef@cdp:eq:C.TEST.1\endcsname{\hyperlink{tag-012O}{II.\ref*{cdp:eq:C.TEST.1}}}
\expandafter\def\csname SRef@cdp:eq:C.TEST.2\endcsname{\hyperlink{tag-012P}{II.\ref*{cdp:eq:C.TEST.2}}}
\expandafter\def\csname SRef@cdp:eq:C.3\endcsname{\hyperlink{tag-0094}{II.\ref*{cdp:eq:C.3}}}
\expandafter\def\csname SRef@cdp:eq:C.4\endcsname{\hyperlink{tag-0095}{II.\ref*{cdp:eq:C.4}}}
\expandafter\def\csname SRef@cdp:eq:C.5\endcsname{\hyperlink{tag-0096}{II.\ref*{cdp:eq:C.5}}}
\expandafter\def\csname SRef@cdp:eq:C.6\endcsname{\hyperlink{tag-0097}{II.\ref*{cdp:eq:C.6}}}
\expandafter\def\csname SRef@cdp:eq:C.7\endcsname{\hyperlink{tag-0098}{II.\ref*{cdp:eq:C.7}}}
\expandafter\def\csname SRef@cdp:eq:C.8\endcsname{\hyperlink{tag-0099}{II.\ref*{cdp:eq:C.8}}}
\expandafter\def\csname SRef@cdp:eq:C.PK.2\endcsname{\hyperlink{tag-00V9}{II.\ref*{cdp:eq:C.PK.2}}}
\expandafter\def\csname SRef@cdp:eq:C.9\endcsname{\hyperlink{tag-009A}{II.\ref*{cdp:eq:C.9}}}
\expandafter\def\csname SRef@cdp:eq:C.10\endcsname{\hyperlink{tag-008O}{II.\ref*{cdp:eq:C.10}}}
\expandafter\def\csname SRef@cdp:eq:C.11\endcsname{\hyperlink{tag-008P}{II.\ref*{cdp:eq:C.11}}}
\expandafter\def\csname SRef@cdp:eq:C.12\endcsname{\hyperlink{tag-008Q}{II.\ref*{cdp:eq:C.12}}}
\expandafter\def\csname SRef@cdp:eq:C.13\endcsname{\hyperlink{tag-008R}{II.\ref*{cdp:eq:C.13}}}
\expandafter\def\csname SRef@cdp:eq:C.14\endcsname{\hyperlink{tag-008S}{II.\ref*{cdp:eq:C.14}}}
\expandafter\def\csname SRef@cdp:eq:D.18\endcsname{\hyperlink{tag-0195}{II.\ref*{cdp:eq:D.18}}}
\expandafter\def\csname SRef@cdp:eq:D.19\endcsname{\hyperlink{tag-0196}{II.\ref*{cdp:eq:D.19}}}
\expandafter\def\csname SRef@cdp:eq:D.S.1\endcsname{\hyperlink{tag-019O}{II.\ref*{cdp:eq:D.S.1}}}
\expandafter\def\csname SRef@cdp:eq:C.PK.3\endcsname{\hyperlink{tag-00VA}{II.\ref*{cdp:eq:C.PK.3}}}
\expandafter\def\csname SRef@cdp:eq:D.S.2\endcsname{\hyperlink{tag-019P}{II.\ref*{cdp:eq:D.S.2}}}
\expandafter\def\csname SRef@cdp:eq:D.S.3\endcsname{\hyperlink{tag-019Q}{II.\ref*{cdp:eq:D.S.3}}}
\expandafter\def\csname SRef@cdp:eq:D.S.4\endcsname{\hyperlink{tag-019R}{II.\ref*{cdp:eq:D.S.4}}}
\expandafter\def\csname SRef@cdp:eq:D.S.5\endcsname{\hyperlink{tag-019S}{II.\ref*{cdp:eq:D.S.5}}}
\expandafter\def\csname SRef@cdp:eq:D.20\endcsname{\hyperlink{tag-0197}{II.\ref*{cdp:eq:D.20}}}
\expandafter\def\csname SRef@cdp:eq:D.21\endcsname{\hyperlink{tag-0198}{II.\ref*{cdp:eq:D.21}}}
\expandafter\def\csname SRef@cdp:eq:D.22\endcsname{\hyperlink{tag-0199}{II.\ref*{cdp:eq:D.22}}}
\expandafter\def\csname SRef@cdp:eq:D.23\endcsname{\hyperlink{tag-019A}{II.\ref*{cdp:eq:D.23}}}
\expandafter\def\csname SRef@cdp:eq:D.24\endcsname{\hyperlink{tag-019B}{II.\ref*{cdp:eq:D.24}}}
\expandafter\def\csname SRef@cdp:eq:C.PK.4\endcsname{\hyperlink{tag-00VB}{II.\ref*{cdp:eq:C.PK.4}}}
\expandafter\def\csname SRef@cdp:eq:D.25\endcsname{\hyperlink{tag-019C}{II.\ref*{cdp:eq:D.25}}}
\expandafter\def\csname SRef@cdp:eq:D.26\endcsname{\hyperlink{tag-019D}{II.\ref*{cdp:eq:D.26}}}
\expandafter\def\csname SRef@cdp:eq:D.27\endcsname{\hyperlink{tag-019E}{II.\ref*{cdp:eq:D.27}}}
\expandafter\def\csname SRef@cdp:eq:D.28\endcsname{\hyperlink{tag-019F}{II.\ref*{cdp:eq:D.28}}}
\expandafter\def\csname SRef@cdp:eq:D.29\endcsname{\hyperlink{tag-019G}{II.\ref*{cdp:eq:D.29}}}
\expandafter\def\csname SRef@cdp:eq:D.30\endcsname{\hyperlink{tag-019H}{II.\ref*{cdp:eq:D.30}}}
\expandafter\def\csname SRef@cdp:eq:D.31\endcsname{\hyperlink{tag-019I}{II.\ref*{cdp:eq:D.31}}}
\expandafter\def\csname SRef@cdp:eq:D.32\endcsname{\hyperlink{tag-019J}{II.\ref*{cdp:eq:D.32}}}
\expandafter\def\csname SRef@cdp:eq:D.33\endcsname{\hyperlink{tag-019K}{II.\ref*{cdp:eq:D.33}}}
\expandafter\def\csname SRef@cdp:eq:C.PK.5\endcsname{\hyperlink{tag-00VC}{II.\ref*{cdp:eq:C.PK.5}}}
\expandafter\def\csname SRef@cdp:eq:D.34\endcsname{\hyperlink{tag-019L}{II.\ref*{cdp:eq:D.34}}}
\expandafter\def\csname SRef@cdp:eq:D.35\endcsname{\hyperlink{tag-019M}{II.\ref*{cdp:eq:D.35}}}
\expandafter\def\csname SRef@cdp:eq:D.36\endcsname{\hyperlink{tag-019N}{II.\ref*{cdp:eq:D.36}}}
\expandafter\def\csname SRef@cdp:eq:E.16.1\endcsname{\hyperlink{tag-01A2}{II.\ref*{cdp:eq:E.16.1}}}
\expandafter\def\csname SRef@cdp:eq:E.16.2\endcsname{\hyperlink{tag-01A3}{II.\ref*{cdp:eq:E.16.2}}}
\expandafter\def\csname SRef@cdp:eq:E.16.3\endcsname{\hyperlink{tag-01A4}{II.\ref*{cdp:eq:E.16.3}}}
\expandafter\def\csname SRef@cdp:eq:C.PK.6\endcsname{\hyperlink{tag-00VD}{II.\ref*{cdp:eq:C.PK.6}}}
\expandafter\def\csname SRef@cdp:eq:C.PK.7\endcsname{\hyperlink{tag-00VE}{II.\ref*{cdp:eq:C.PK.7}}}
\expandafter\def\csname SRef@cdp:eq:C.PK.8\endcsname{\hyperlink{tag-00VF}{II.\ref*{cdp:eq:C.PK.8}}}
\expandafter\def\csname SRef@cdp:eq:C.IN.1\endcsname{\hyperlink{tag-00NK}{II.\ref*{cdp:eq:C.IN.1}}}
\expandafter\def\csname SRef@invII:eq:E.2\endcsname{\hyperlink{tag-01TR}{II.\ref*{invII:eq:E.2}}}
\expandafter\def\csname SRef@cdp:eq:C.IN.2\endcsname{\hyperlink{tag-00NL}{II.\ref*{cdp:eq:C.IN.2}}}
\expandafter\def\csname SRef@invII:eq:E.3\endcsname{\hyperlink{tag-01TS}{II.\ref*{invII:eq:E.3}}}
\expandafter\def\csname SRef@invII:eq:E.4a\endcsname{\hyperlink{tag-01TT}{II.\ref*{invII:eq:E.4a}}}
\expandafter\def\csname SRef@invII:eq:E.5\endcsname{\hyperlink{tag-01TU}{II.\ref*{invII:eq:E.5}}}
\expandafter\def\csname SRef@invII:eq:E.6\endcsname{\hyperlink{tag-01TV}{II.\ref*{invII:eq:E.6}}}
\expandafter\def\csname SRef@inv:wb:eq:9.19\endcsname{\hyperlink{tag-01TJ}{II.\ref*{inv:wb:eq:9.19}}}
\expandafter\def\csname SRef@inv:wb:eq:9.20\endcsname{\hyperlink{tag-01TK}{II.\ref*{inv:wb:eq:9.20}}}
\expandafter\def\csname SRef@inv:wb:eq:9.21\endcsname{\hyperlink{tag-01TL}{II.\ref*{inv:wb:eq:9.21}}}
\expandafter\def\csname SRef@inv:wb:eq:9.22\endcsname{\hyperlink{tag-01TM}{II.\ref*{inv:wb:eq:9.22}}}
\expandafter\def\csname SRef@inv:wb:eq:9.23\endcsname{\hyperlink{tag-01TN}{II.\ref*{inv:wb:eq:9.23}}}
\expandafter\def\csname SRef@inv:wb:eq:9.24\endcsname{\hyperlink{tag-01TO}{II.\ref*{inv:wb:eq:9.24}}}
\expandafter\def\csname SRef@cdp:eq:C.IN.3\endcsname{\hyperlink{tag-00NM}{II.\ref*{cdp:eq:C.IN.3}}}
\expandafter\def\csname SRef@inv:wb:eq:9.25\endcsname{\hyperlink{tag-01TP}{II.\ref*{inv:wb:eq:9.25}}}
\expandafter\def\csname SRef@wbII:eq:CN:2.2\endcsname{\hyperlink{tag-03DW}{II.\ref*{wbII:eq:CN:2.2}}}
\expandafter\def\csname SRef@wbII:eq:W:5.3\endcsname{\hyperlink{tag-03HS}{II.\ref*{wbII:eq:W:5.3}}}
\expandafter\def\csname SRef@wbII:eq:W:5.4\endcsname{\hyperlink{tag-03HX}{II.\ref*{wbII:eq:W:5.4}}}
\expandafter\def\csname SRef@wbII:eq:W:5.5\endcsname{\hyperlink{tag-03HY}{II.\ref*{wbII:eq:W:5.5}}}
\expandafter\def\csname SRef@wbII:eq:W:5.6\endcsname{\hyperlink{tag-03HZ}{II.\ref*{wbII:eq:W:5.6}}}
\expandafter\def\csname SRef@cdp:eq:C.IN.4\endcsname{\hyperlink{tag-00NN}{II.\ref*{cdp:eq:C.IN.4}}}
\expandafter\def\csname SRef@wbII:eq:W:5.7\endcsname{\hyperlink{tag-03I0}{II.\ref*{wbII:eq:W:5.7}}}
\expandafter\def\csname SRef@wbII:eq:W:5.8\endcsname{\hyperlink{tag-03I1}{II.\ref*{wbII:eq:W:5.8}}}
\expandafter\def\csname SRef@wbII:eq:W:5.9\endcsname{\hyperlink{tag-03I2}{II.\ref*{wbII:eq:W:5.9}}}
\expandafter\def\csname SRef@wbII:eq:W:5.10\endcsname{\hyperlink{tag-03EA}{II.\ref*{wbII:eq:W:5.10}}}
\expandafter\def\csname SRef@wbII:eq:CN:4.3\endcsname{\hyperlink{tag-03DX}{II.\ref*{wbII:eq:CN:4.3}}}
\expandafter\def\csname SRef@wbII:eq:CN:4.4\endcsname{\hyperlink{tag-03DY}{II.\ref*{wbII:eq:CN:4.4}}}
\expandafter\def\csname SRef@wbII:eq:CN:4.5a\endcsname{\hyperlink{tag-03E0}{II.\ref*{wbII:eq:CN:4.5a}}}
\expandafter\def\csname SRef@cdp:eq:C.IN.5\endcsname{\hyperlink{tag-00NO}{II.\ref*{cdp:eq:C.IN.5}}}
\expandafter\def\csname SRef@wbII:eq:CN:4.5\endcsname{\hyperlink{tag-03DZ}{II.\ref*{wbII:eq:CN:4.5}}}
\expandafter\def\csname SRef@wbII:eq:CN:4.6\endcsname{\hyperlink{tag-03E1}{II.\ref*{wbII:eq:CN:4.6}}}
\expandafter\def\csname SRef@wbII:eq:CN:4.7\endcsname{\hyperlink{tag-03E2}{II.\ref*{wbII:eq:CN:4.7}}}
\expandafter\def\csname SRef@wbII:eq:CN:4.8\endcsname{\hyperlink{tag-03E3}{II.\ref*{wbII:eq:CN:4.8}}}
\expandafter\def\csname SRef@wbII:eq:CN:4.9\endcsname{\hyperlink{tag-03E4}{II.\ref*{wbII:eq:CN:4.9}}}
\expandafter\def\csname SRef@wbII:eq:W:5.11\endcsname{\hyperlink{tag-03EB}{II.\ref*{wbII:eq:W:5.11}}}
\expandafter\def\csname SRef@wbII:eq:W:5.15a\endcsname{\hyperlink{tag-03EC}{II.\ref*{wbII:eq:W:5.15a}}}
\expandafter\def\csname SRef@wbII:eq:W:5.16\endcsname{\hyperlink{tag-03ED}{II.\ref*{wbII:eq:W:5.16}}}
\expandafter\def\csname SRef@wbII:eq:W:5.17\endcsname{\hyperlink{tag-03EE}{II.\ref*{wbII:eq:W:5.17}}}
\expandafter\def\csname SRef@cdp:eq:C.IN.6\endcsname{\hyperlink{tag-00NP}{II.\ref*{cdp:eq:C.IN.6}}}
\expandafter\def\csname SRef@wbII:eq:W:5.17-V\endcsname{\hyperlink{tag-03FU}{II.\ref*{wbII:eq:W:5.17-V}}}
\expandafter\def\csname SRef@wbII:eq:W:5.17-T\endcsname{\hyperlink{tag-03FS}{II.\ref*{wbII:eq:W:5.17-T}}}
\expandafter\def\csname SRef@wbII:eq:W:5.17-U\endcsname{\hyperlink{tag-03FT}{II.\ref*{wbII:eq:W:5.17-U}}}
\expandafter\def\csname SRef@wbII:eq:W:5.17a\endcsname{\hyperlink{tag-03FV}{II.\ref*{wbII:eq:W:5.17a}}}
\expandafter\def\csname SRef@wbII:eq:W:5.17b\endcsname{\hyperlink{tag-03GL}{II.\ref*{wbII:eq:W:5.17b}}}
\expandafter\def\csname SRef@wbII:eq:W:5.17c\endcsname{\hyperlink{tag-03GM}{II.\ref*{wbII:eq:W:5.17c}}}
\expandafter\def\csname SRef@wbII:eq:W:5.17d\endcsname{\hyperlink{tag-03GN}{II.\ref*{wbII:eq:W:5.17d}}}
\expandafter\def\csname SRef@wbII:eq:W:5.17e\endcsname{\hyperlink{tag-03GO}{II.\ref*{wbII:eq:W:5.17e}}}
\expandafter\def\csname SRef@wbII:eq:W:5.17g\endcsname{\hyperlink{tag-03GQ}{II.\ref*{wbII:eq:W:5.17g}}}
\expandafter\def\csname SRef@wbII:eq:W:5.17h\endcsname{\hyperlink{tag-03GR}{II.\ref*{wbII:eq:W:5.17h}}}
\expandafter\def\csname SRef@cdp:eq:C.PK.9\endcsname{\hyperlink{tag-00VG}{II.\ref*{cdp:eq:C.PK.9}}}
\expandafter\def\csname SRef@wbII:eq:W:5.17i\endcsname{\hyperlink{tag-03GS}{II.\ref*{wbII:eq:W:5.17i}}}
\expandafter\def\csname SRef@wbII:eq:W:5.17j'\endcsname{\hyperlink{tag-03GU}{II.\ref*{wbII:eq:W:5.17j'}}}
\expandafter\def\csname SRef@wbII:eq:W:5.17j\endcsname{\hyperlink{tag-03GT}{II.\ref*{wbII:eq:W:5.17j}}}
\expandafter\def\csname SRef@wbII:eq:W:5.17k\endcsname{\hyperlink{tag-03GV}{II.\ref*{wbII:eq:W:5.17k}}}
\expandafter\def\csname SRef@wbII:eq:W:5.17l\endcsname{\hyperlink{tag-03GW}{II.\ref*{wbII:eq:W:5.17l}}}
\expandafter\def\csname SRef@wbII:eq:W:5.17l'\endcsname{\hyperlink{tag-03GX}{II.\ref*{wbII:eq:W:5.17l'}}}
\expandafter\def\csname SRef@wbII:eq:W:5.17l''\endcsname{\hyperlink{tag-03GY}{II.\ref*{wbII:eq:W:5.17l''}}}
\expandafter\def\csname SRef@wbII:eq:W:5.17l'''\endcsname{\hyperlink{tag-03GZ}{II.\ref*{wbII:eq:W:5.17l'''}}}
\expandafter\def\csname SRef@wbII:eq:W:5.17m\endcsname{\hyperlink{tag-03H0}{II.\ref*{wbII:eq:W:5.17m}}}
\expandafter\def\csname SRef@wbII:eq:W:5.17m'\endcsname{\hyperlink{tag-03H1}{II.\ref*{wbII:eq:W:5.17m'}}}
\expandafter\def\csname SRef@cdp:eq:C.H.1\endcsname{\hyperlink{tag-00L8}{II.\ref*{cdp:eq:C.H.1}}}
\expandafter\def\csname SRef@wbII:eq:W:5.17m''\endcsname{\hyperlink{tag-03H2}{II.\ref*{wbII:eq:W:5.17m''}}}
\expandafter\def\csname SRef@wbII:eq:W:5.17n\endcsname{\hyperlink{tag-03H3}{II.\ref*{wbII:eq:W:5.17n}}}
\expandafter\def\csname SRef@wbII:eq:W:5.17o\endcsname{\hyperlink{tag-03H4}{II.\ref*{wbII:eq:W:5.17o}}}
\expandafter\def\csname SRef@wbII:eq:W:5.17p\endcsname{\hyperlink{tag-03H5}{II.\ref*{wbII:eq:W:5.17p}}}
\expandafter\def\csname SRef@wbII:eq:W:5.17p0\endcsname{\hyperlink{tag-03H7}{II.\ref*{wbII:eq:W:5.17p0}}}
\expandafter\def\csname SRef@wbII:eq:W:5.17p'\endcsname{\hyperlink{tag-03H6}{II.\ref*{wbII:eq:W:5.17p'}}}
\expandafter\def\csname SRef@wbII:eq:W:5.17q\endcsname{\hyperlink{tag-03H8}{II.\ref*{wbII:eq:W:5.17q}}}
\expandafter\def\csname SRef@wbII:eq:W:5.17r\endcsname{\hyperlink{tag-03H9}{II.\ref*{wbII:eq:W:5.17r}}}
\expandafter\def\csname SRef@wbII:eq:W:5.17s\endcsname{\hyperlink{tag-03HA}{II.\ref*{wbII:eq:W:5.17s}}}
\expandafter\def\csname SRef@wbII:eq:W:5.17t\endcsname{\hyperlink{tag-03HB}{II.\ref*{wbII:eq:W:5.17t}}}
\expandafter\def\csname SRef@cdp:eq:A.1\endcsname{\hyperlink{tag-003Y}{II.\ref*{cdp:eq:A.1}}}
\expandafter\def\csname SRef@wbII:eq:W:5.17u\endcsname{\hyperlink{tag-03HC}{II.\ref*{wbII:eq:W:5.17u}}}
\expandafter\def\csname SRef@wbII:eq:W:5.17v\endcsname{\hyperlink{tag-03HD}{II.\ref*{wbII:eq:W:5.17v}}}
\expandafter\def\csname SRef@wbII:eq:W:5.17w\endcsname{\hyperlink{tag-03HE}{II.\ref*{wbII:eq:W:5.17w}}}
\expandafter\def\csname SRef@wbII:eq:W:5.17x\endcsname{\hyperlink{tag-03HF}{II.\ref*{wbII:eq:W:5.17x}}}
\expandafter\def\csname SRef@wbII:eq:W:5.17y\endcsname{\hyperlink{tag-03HG}{II.\ref*{wbII:eq:W:5.17y}}}
\expandafter\def\csname SRef@wbII:eq:W:5.17z\endcsname{\hyperlink{tag-03HH}{II.\ref*{wbII:eq:W:5.17z}}}
\expandafter\def\csname SRef@wbII:eq:W:5.17aa\endcsname{\hyperlink{tag-03FW}{II.\ref*{wbII:eq:W:5.17aa}}}
\expandafter\def\csname SRef@wbII:eq:W:5.17ab\endcsname{\hyperlink{tag-03FX}{II.\ref*{wbII:eq:W:5.17ab}}}
\expandafter\def\csname SRef@wbII:eq:W:5.17ac\endcsname{\hyperlink{tag-03FY}{II.\ref*{wbII:eq:W:5.17ac}}}
\expandafter\def\csname SRef@wbII:eq:W:5.17ad\endcsname{\hyperlink{tag-03FZ}{II.\ref*{wbII:eq:W:5.17ad}}}
\expandafter\def\csname SRef@cdp:eq:A.2\endcsname{\hyperlink{tag-004B}{II.\ref*{cdp:eq:A.2}}}
\expandafter\def\csname SRef@wbII:eq:W:5.17ae\endcsname{\hyperlink{tag-03G0}{II.\ref*{wbII:eq:W:5.17ae}}}
\expandafter\def\csname SRef@wbII:eq:W:5.17af\endcsname{\hyperlink{tag-03G1}{II.\ref*{wbII:eq:W:5.17af}}}
\expandafter\def\csname SRef@wbII:eq:W:5.17ag\endcsname{\hyperlink{tag-03G2}{II.\ref*{wbII:eq:W:5.17ag}}}
\expandafter\def\csname SRef@wbII:eq:W:5.17ah\endcsname{\hyperlink{tag-03G3}{II.\ref*{wbII:eq:W:5.17ah}}}
\expandafter\def\csname SRef@wbII:eq:W:5.17ai\endcsname{\hyperlink{tag-03G4}{II.\ref*{wbII:eq:W:5.17ai}}}
\expandafter\def\csname SRef@wbII:eq:W:5.17aj\endcsname{\hyperlink{tag-03G5}{II.\ref*{wbII:eq:W:5.17aj}}}
\expandafter\def\csname SRef@wbII:eq:W:5.17ak\endcsname{\hyperlink{tag-03G6}{II.\ref*{wbII:eq:W:5.17ak}}}
\expandafter\def\csname SRef@wbII:eq:W:5.17al\endcsname{\hyperlink{tag-03G7}{II.\ref*{wbII:eq:W:5.17al}}}
\expandafter\def\csname SRef@cdp:eq:A.3\endcsname{\hyperlink{tag-004M}{II.\ref*{cdp:eq:A.3}}}
\expandafter\def\csname SRef@wbII:eq:W:5.17am\endcsname{\hyperlink{tag-03G8}{II.\ref*{wbII:eq:W:5.17am}}}
\expandafter\def\csname SRef@wbII:eq:W:5.17an\endcsname{\hyperlink{tag-03G9}{II.\ref*{wbII:eq:W:5.17an}}}
\expandafter\def\csname SRef@wbII:eq:W:5.17ao\endcsname{\hyperlink{tag-03GA}{II.\ref*{wbII:eq:W:5.17ao}}}
\expandafter\def\csname SRef@wbII:eq:W:5.17ap\endcsname{\hyperlink{tag-03GB}{II.\ref*{wbII:eq:W:5.17ap}}}
\expandafter\def\csname SRef@wbII:eq:W:5.17aq\endcsname{\hyperlink{tag-03GC}{II.\ref*{wbII:eq:W:5.17aq}}}
\expandafter\def\csname SRef@wbII:eq:W:5.17ar\endcsname{\hyperlink{tag-03GD}{II.\ref*{wbII:eq:W:5.17ar}}}
\expandafter\def\csname SRef@wbII:eq:W:5.17as\endcsname{\hyperlink{tag-03GE}{II.\ref*{wbII:eq:W:5.17as}}}
\expandafter\def\csname SRef@wbII:eq:W:5.17at\endcsname{\hyperlink{tag-03GF}{II.\ref*{wbII:eq:W:5.17at}}}
\expandafter\def\csname SRef@wbII:eq:W:5.17au\endcsname{\hyperlink{tag-03GG}{II.\ref*{wbII:eq:W:5.17au}}}
\expandafter\def\csname SRef@wbII:eq:W:5.17av\endcsname{\hyperlink{tag-03GH}{II.\ref*{wbII:eq:W:5.17av}}}
\expandafter\def\csname SRef@cdp:eq:A.4\endcsname{\hyperlink{tag-004O}{II.\ref*{cdp:eq:A.4}}}
\expandafter\def\csname SRef@wbII:eq:W:5.17aw\endcsname{\hyperlink{tag-03GI}{II.\ref*{wbII:eq:W:5.17aw}}}
\expandafter\def\csname SRef@wbII:eq:W:5.17ax\endcsname{\hyperlink{tag-03GJ}{II.\ref*{wbII:eq:W:5.17ax}}}
\expandafter\def\csname SRef@wbII:eq:W:5.17ay\endcsname{\hyperlink{tag-03GK}{II.\ref*{wbII:eq:W:5.17ay}}}
\expandafter\def\csname SRef@wbII:eq:W:5.17-A1\endcsname{\hyperlink{tag-03EF}{II.\ref*{wbII:eq:W:5.17-A1}}}
\expandafter\def\csname SRef@wbII:eq:W:5.17-A2\endcsname{\hyperlink{tag-03EY}{II.\ref*{wbII:eq:W:5.17-A2}}}
\expandafter\def\csname SRef@wbII:eq:W:5.17-A3\endcsname{\hyperlink{tag-03F9}{II.\ref*{wbII:eq:W:5.17-A3}}}
\expandafter\def\csname SRef@wbII:eq:W:5.17-A4\endcsname{\hyperlink{tag-03FM}{II.\ref*{wbII:eq:W:5.17-A4}}}
\expandafter\def\csname SRef@wbII:eq:W:5.17-A5\endcsname{\hyperlink{tag-03FN}{II.\ref*{wbII:eq:W:5.17-A5}}}
\expandafter\def\csname SRef@wbII:eq:W:5.17-A6\endcsname{\hyperlink{tag-03FO}{II.\ref*{wbII:eq:W:5.17-A6}}}
\expandafter\def\csname SRef@wbII:eq:W:5.17-A7\endcsname{\hyperlink{tag-03FP}{II.\ref*{wbII:eq:W:5.17-A7}}}
\expandafter\def\csname SRef@wbII:eq:W:5.17-A8\endcsname{\hyperlink{tag-03FQ}{II.\ref*{wbII:eq:W:5.17-A8}}}
\expandafter\def\csname SRef@wbII:eq:W:5.17-A9\endcsname{\hyperlink{tag-03FR}{II.\ref*{wbII:eq:W:5.17-A9}}}
\expandafter\def\csname SRef@wbII:eq:W:5.17-A10\endcsname{\hyperlink{tag-03EG}{II.\ref*{wbII:eq:W:5.17-A10}}}
\expandafter\def\csname SRef@wbII:eq:W:5.17-A11\endcsname{\hyperlink{tag-03EH}{II.\ref*{wbII:eq:W:5.17-A11}}}
\expandafter\def\csname SRef@wbII:eq:W:5.17-A12\endcsname{\hyperlink{tag-03EI}{II.\ref*{wbII:eq:W:5.17-A12}}}
\expandafter\def\csname SRef@wbII:eq:W:5.17-A13\endcsname{\hyperlink{tag-03EJ}{II.\ref*{wbII:eq:W:5.17-A13}}}
\expandafter\def\csname SRef@wbII:eq:W:5.17-A14\endcsname{\hyperlink{tag-03EK}{II.\ref*{wbII:eq:W:5.17-A14}}}
\expandafter\def\csname SRef@wbII:eq:W:5.17-A14a\endcsname{\hyperlink{tag-03EL}{II.\ref*{wbII:eq:W:5.17-A14a}}}
\expandafter\def\csname SRef@cdp:eq:A.5\endcsname{\hyperlink{tag-004P}{II.\ref*{cdp:eq:A.5}}}
\expandafter\def\csname SRef@wbII:eq:W:5.17-A14c\endcsname{\hyperlink{tag-03EM}{II.\ref*{wbII:eq:W:5.17-A14c}}}
\expandafter\def\csname SRef@wbII:eq:W:5.17-A14d\endcsname{\hyperlink{tag-03EN}{II.\ref*{wbII:eq:W:5.17-A14d}}}
\expandafter\def\csname SRef@wbII:eq:W:5.17-A14e\endcsname{\hyperlink{tag-03EO}{II.\ref*{wbII:eq:W:5.17-A14e}}}
\expandafter\def\csname SRef@wbII:eq:W:5.17-A14f\endcsname{\hyperlink{tag-03EP}{II.\ref*{wbII:eq:W:5.17-A14f}}}
\expandafter\def\csname SRef@wbII:eq:W:5.17-A15\endcsname{\hyperlink{tag-03EQ}{II.\ref*{wbII:eq:W:5.17-A15}}}
\expandafter\def\csname SRef@wbII:eq:W:5.17-A16\endcsname{\hyperlink{tag-03ER}{II.\ref*{wbII:eq:W:5.17-A16}}}
\expandafter\def\csname SRef@wbII:eq:W:5.17-A17\endcsname{\hyperlink{tag-03ES}{II.\ref*{wbII:eq:W:5.17-A17}}}
\expandafter\def\csname SRef@wbII:eq:W:5.17-A18\endcsname{\hyperlink{tag-03ET}{II.\ref*{wbII:eq:W:5.17-A18}}}
\expandafter\def\csname SRef@wbII:eq:W:5.17-A19\endcsname{\hyperlink{tag-03EU}{II.\ref*{wbII:eq:W:5.17-A19}}}
\expandafter\def\csname SRef@wbII:eq:W:5.17-A19a\endcsname{\hyperlink{tag-03EV}{II.\ref*{wbII:eq:W:5.17-A19a}}}
\expandafter\def\csname SRef@cdp:eq:A.6\endcsname{\hyperlink{tag-004Q}{II.\ref*{cdp:eq:A.6}}}
\expandafter\def\csname SRef@wbII:eq:W:5.17-A19b\endcsname{\hyperlink{tag-03EW}{II.\ref*{wbII:eq:W:5.17-A19b}}}
\expandafter\def\csname SRef@wbII:eq:W:5.17-A19c\endcsname{\hyperlink{tag-03EX}{II.\ref*{wbII:eq:W:5.17-A19c}}}
\expandafter\def\csname SRef@wbII:eq:W:5.17-A20\endcsname{\hyperlink{tag-03EZ}{II.\ref*{wbII:eq:W:5.17-A20}}}
\expandafter\def\csname SRef@wbII:eq:W:5.17-A21\endcsname{\hyperlink{tag-03F0}{II.\ref*{wbII:eq:W:5.17-A21}}}
\expandafter\def\csname SRef@wbII:eq:W:5.17-A22\endcsname{\hyperlink{tag-03F1}{II.\ref*{wbII:eq:W:5.17-A22}}}
\expandafter\def\csname SRef@wbII:eq:W:5.17-A23\endcsname{\hyperlink{tag-03F2}{II.\ref*{wbII:eq:W:5.17-A23}}}
\expandafter\def\csname SRef@wbII:eq:W:5.17-A24\endcsname{\hyperlink{tag-03F3}{II.\ref*{wbII:eq:W:5.17-A24}}}
\expandafter\def\csname SRef@wbII:eq:W:5.17-A25\endcsname{\hyperlink{tag-03F4}{II.\ref*{wbII:eq:W:5.17-A25}}}
\expandafter\def\csname SRef@wbII:eq:W:5.17-A26\endcsname{\hyperlink{tag-03F5}{II.\ref*{wbII:eq:W:5.17-A26}}}
\expandafter\def\csname SRef@wbII:eq:W:5.17-A27\endcsname{\hyperlink{tag-03F6}{II.\ref*{wbII:eq:W:5.17-A27}}}
\expandafter\def\csname SRef@cdp:eq:A.7\endcsname{\hyperlink{tag-004R}{II.\ref*{cdp:eq:A.7}}}
\expandafter\def\csname SRef@wbII:eq:W:5.17-A28\endcsname{\hyperlink{tag-03F7}{II.\ref*{wbII:eq:W:5.17-A28}}}
\expandafter\def\csname SRef@wbII:eq:W:5.17-A29\endcsname{\hyperlink{tag-03F8}{II.\ref*{wbII:eq:W:5.17-A29}}}
\expandafter\def\csname SRef@wbII:eq:W:5.17-A30\endcsname{\hyperlink{tag-03FA}{II.\ref*{wbII:eq:W:5.17-A30}}}
\expandafter\def\csname SRef@wbII:eq:W:5.17-A31\endcsname{\hyperlink{tag-03FB}{II.\ref*{wbII:eq:W:5.17-A31}}}
\expandafter\def\csname SRef@wbII:eq:W:5.17-A32\endcsname{\hyperlink{tag-03FC}{II.\ref*{wbII:eq:W:5.17-A32}}}
\expandafter\def\csname SRef@wbII:eq:W:5.17-A33\endcsname{\hyperlink{tag-03FD}{II.\ref*{wbII:eq:W:5.17-A33}}}
\expandafter\def\csname SRef@wbII:eq:W:5.17-A34\endcsname{\hyperlink{tag-03FE}{II.\ref*{wbII:eq:W:5.17-A34}}}
\expandafter\def\csname SRef@wbII:eq:W:5.17-A34a\endcsname{\hyperlink{tag-03FF}{II.\ref*{wbII:eq:W:5.17-A34a}}}
\expandafter\def\csname SRef@wbII:eq:W:5.17-A34b\endcsname{\hyperlink{tag-03FG}{II.\ref*{wbII:eq:W:5.17-A34b}}}
\expandafter\def\csname SRef@wbII:eq:W:5.17-A34c\endcsname{\hyperlink{tag-03FH}{II.\ref*{wbII:eq:W:5.17-A34c}}}
\expandafter\def\csname SRef@cdp:eq:A.8\endcsname{\hyperlink{tag-004S}{II.\ref*{cdp:eq:A.8}}}
\expandafter\def\csname SRef@wbII:eq:W:5.17-A35\endcsname{\hyperlink{tag-03FI}{II.\ref*{wbII:eq:W:5.17-A35}}}
\expandafter\def\csname SRef@wbII:eq:W:5.17-A36\endcsname{\hyperlink{tag-03FJ}{II.\ref*{wbII:eq:W:5.17-A36}}}
\expandafter\def\csname SRef@wbII:eq:W:5.17-A37\endcsname{\hyperlink{tag-03FK}{II.\ref*{wbII:eq:W:5.17-A37}}}
\expandafter\def\csname SRef@wbII:eq:W:5.17-A38\endcsname{\hyperlink{tag-03FL}{II.\ref*{wbII:eq:W:5.17-A38}}}
\expandafter\def\csname SRef@wbII:eq:W:5.17f\endcsname{\hyperlink{tag-03GP}{II.\ref*{wbII:eq:W:5.17f}}}
\expandafter\def\csname SRef@wbII:eq:W:5.18\endcsname{\hyperlink{tag-03HI}{II.\ref*{wbII:eq:W:5.18}}}
\expandafter\def\csname SRef@wbII:eq:W:5.19\endcsname{\hyperlink{tag-03HJ}{II.\ref*{wbII:eq:W:5.19}}}
\expandafter\def\csname SRef@wbII:eq:W:5.20\endcsname{\hyperlink{tag-03HK}{II.\ref*{wbII:eq:W:5.20}}}
\expandafter\def\csname SRef@wbII:eq:CN:5.2\endcsname{\hyperlink{tag-03E5}{II.\ref*{wbII:eq:CN:5.2}}}
\expandafter\def\csname SRef@cdp:eq:A.9\endcsname{\hyperlink{tag-004T}{II.\ref*{cdp:eq:A.9}}}
\expandafter\def\csname SRef@wbII:eq:W:5.23\endcsname{\hyperlink{tag-03HL}{II.\ref*{wbII:eq:W:5.23}}}
\expandafter\def\csname SRef@wbII:eq:W:5.24\endcsname{\hyperlink{tag-03HM}{II.\ref*{wbII:eq:W:5.24}}}
\expandafter\def\csname SRef@wbII:eq:W:5.25\endcsname{\hyperlink{tag-03HN}{II.\ref*{wbII:eq:W:5.25}}}
\expandafter\def\csname SRef@wbII:eq:W:5.26\endcsname{\hyperlink{tag-03HO}{II.\ref*{wbII:eq:W:5.26}}}
\expandafter\def\csname SRef@wbII:eq:W:5.27\endcsname{\hyperlink{tag-03HP}{II.\ref*{wbII:eq:W:5.27}}}
\expandafter\def\csname SRef@wbII:eq:W:5.28\endcsname{\hyperlink{tag-03HQ}{II.\ref*{wbII:eq:W:5.28}}}
\expandafter\def\csname SRef@wbII:eq:W:5.29\endcsname{\hyperlink{tag-03HR}{II.\ref*{wbII:eq:W:5.29}}}
\expandafter\def\csname SRef@wbII:eq:W:5.30\endcsname{\hyperlink{tag-03HT}{II.\ref*{wbII:eq:W:5.30}}}
\expandafter\def\csname SRef@wbII:eq:W:5.31\endcsname{\hyperlink{tag-03HU}{II.\ref*{wbII:eq:W:5.31}}}
\expandafter\def\csname SRef@wbII:eq:W:5.32\endcsname{\hyperlink{tag-03HV}{II.\ref*{wbII:eq:W:5.32}}}
\expandafter\def\csname SRef@cdp:eq:A.9a\endcsname{\hyperlink{tag-004U}{II.\ref*{cdp:eq:A.9a}}}
\expandafter\def\csname SRef@wbII:eq:W:5.33\endcsname{\hyperlink{tag-03HW}{II.\ref*{wbII:eq:W:5.33}}}
\expandafter\def\csname SRef@wbII:eq:CN:6.1\endcsname{\hyperlink{tag-03E6}{II.\ref*{wbII:eq:CN:6.1}}}
\expandafter\def\csname SRef@wbII:eq:CN:6.2\endcsname{\hyperlink{tag-03E7}{II.\ref*{wbII:eq:CN:6.2}}}
\expandafter\def\csname SRef@wbII:eq:CN:6.3\endcsname{\hyperlink{tag-03E8}{II.\ref*{wbII:eq:CN:6.3}}}
\expandafter\def\csname SRef@wbII:eq:CN:6.4\endcsname{\hyperlink{tag-03E9}{II.\ref*{wbII:eq:CN:6.4}}}
\expandafter\def\csname SRef@cdp:eq:A.9b\endcsname{\hyperlink{tag-004V}{II.\ref*{cdp:eq:A.9b}}}
\expandafter\def\csname SRef@cdp:eq:A.NL.1\endcsname{\hyperlink{tag-006J}{II.\ref*{cdp:eq:A.NL.1}}}
\expandafter\def\csname SRef@invII:eq:E.7\endcsname{\hyperlink{tag-01TW}{II.\ref*{invII:eq:E.7}}}
\expandafter\def\csname SRef@cdp:eq:A.NL.2\endcsname{\hyperlink{tag-006K}{II.\ref*{cdp:eq:A.NL.2}}}
\expandafter\def\csname SRef@cdp:eq:A.NL.3\endcsname{\hyperlink{tag-006L}{II.\ref*{cdp:eq:A.NL.3}}}
\expandafter\def\csname SRef@cdp:eq:A.NL.4\endcsname{\hyperlink{tag-006M}{II.\ref*{cdp:eq:A.NL.4}}}
\expandafter\def\csname SRef@cdp:eq:A.NL.5\endcsname{\hyperlink{tag-006N}{II.\ref*{cdp:eq:A.NL.5}}}
\expandafter\def\csname SRef@cdp:eq:B.1\endcsname{\hyperlink{tag-008G}{II.\ref*{cdp:eq:B.1}}}
\expandafter\def\csname SRef@cdp:eq:B.2\endcsname{\hyperlink{tag-008H}{II.\ref*{cdp:eq:B.2}}}
\expandafter\def\csname SRef@cdp:eq:B.3\endcsname{\hyperlink{tag-008I}{II.\ref*{cdp:eq:B.3}}}
\expandafter\def\csname SRef@cdp:eq:B.4\endcsname{\hyperlink{tag-008J}{II.\ref*{cdp:eq:B.4}}}
\expandafter\def\csname SRef@cdp:eq:B.5\endcsname{\hyperlink{tag-008K}{II.\ref*{cdp:eq:B.5}}}
\expandafter\def\csname SRef@cdp:eq:B.6\endcsname{\hyperlink{tag-008L}{II.\ref*{cdp:eq:B.6}}}
\expandafter\def\csname SRef@cdp:eq:B.7\endcsname{\hyperlink{tag-008M}{II.\ref*{cdp:eq:B.7}}}
\expandafter\def\csname SRef@cdp:eq:A.10\endcsname{\hyperlink{tag-003Z}{II.\ref*{cdp:eq:A.10}}}
\expandafter\def\csname SRef@cdp:eq:A.11\endcsname{\hyperlink{tag-0040}{II.\ref*{cdp:eq:A.11}}}
\expandafter\def\csname SRef@cdp:eq:A.12\endcsname{\hyperlink{tag-0041}{II.\ref*{cdp:eq:A.12}}}
\expandafter\def\csname SRef@cdp:eq:A.13\endcsname{\hyperlink{tag-0042}{II.\ref*{cdp:eq:A.13}}}
\expandafter\def\csname SRef@cdp:eq:A.14\endcsname{\hyperlink{tag-0043}{II.\ref*{cdp:eq:A.14}}}
\expandafter\def\csname SRef@cdp:eq:A.15a\endcsname{\hyperlink{tag-0045}{II.\ref*{cdp:eq:A.15a}}}
\expandafter\def\csname SRef@cdp:eq:A.15b\endcsname{\hyperlink{tag-0046}{II.\ref*{cdp:eq:A.15b}}}
\expandafter\def\csname SRef@cdp:eq:A.15\endcsname{\hyperlink{tag-0044}{II.\ref*{cdp:eq:A.15}}}
\expandafter\def\csname SRef@cdp:eq:A.16\endcsname{\hyperlink{tag-0047}{II.\ref*{cdp:eq:A.16}}}
\expandafter\def\csname SRef@cdp:eq:A.17\endcsname{\hyperlink{tag-0048}{II.\ref*{cdp:eq:A.17}}}
\expandafter\def\csname SRef@cdp:eq:A.18\endcsname{\hyperlink{tag-0049}{II.\ref*{cdp:eq:A.18}}}
\expandafter\def\csname SRef@cdp:eq:A.19\endcsname{\hyperlink{tag-004A}{II.\ref*{cdp:eq:A.19}}}
\expandafter\def\csname SRef@cdp:eq:A.20\endcsname{\hyperlink{tag-004C}{II.\ref*{cdp:eq:A.20}}}
\expandafter\def\csname SRef@cdp:eq:A.21\endcsname{\hyperlink{tag-004D}{II.\ref*{cdp:eq:A.21}}}
\expandafter\def\csname SRef@cdp:eq:A.22\endcsname{\hyperlink{tag-004E}{II.\ref*{cdp:eq:A.22}}}
\expandafter\def\csname SRef@cdp:eq:A.23\endcsname{\hyperlink{tag-004F}{II.\ref*{cdp:eq:A.23}}}
\expandafter\def\csname SRef@cdp:eq:A.24\endcsname{\hyperlink{tag-004G}{II.\ref*{cdp:eq:A.24}}}
\expandafter\def\csname SRef@cdp:eq:A.25\endcsname{\hyperlink{tag-004H}{II.\ref*{cdp:eq:A.25}}}
\expandafter\def\csname SRef@cdp:eq:A.26\endcsname{\hyperlink{tag-004I}{II.\ref*{cdp:eq:A.26}}}
\expandafter\def\csname SRef@cdp:eq:A.27\endcsname{\hyperlink{tag-004J}{II.\ref*{cdp:eq:A.27}}}
\expandafter\def\csname SRef@cdp:eq:A.28\endcsname{\hyperlink{tag-004K}{II.\ref*{cdp:eq:A.28}}}
\expandafter\def\csname SRef@cdp:eq:A.29\endcsname{\hyperlink{tag-004L}{II.\ref*{cdp:eq:A.29}}}
\expandafter\def\csname SRef@cdp:eq:A.30\endcsname{\hyperlink{tag-004N}{II.\ref*{cdp:eq:A.30}}}
\expandafter\def\csname SRef@cdp:eq:A.RC.1\endcsname{\hyperlink{tag-006U}{II.\ref*{cdp:eq:A.RC.1}}}
\expandafter\def\csname SRef@push:src:E.2.4\endcsname{\hyperlink{tag-033C}{II.\ref*{push:src:E.2.4}}}
\expandafter\def\csname SRef@push:src:E.2.24\endcsname{\hyperlink{tag-033A}{II.\ref*{push:src:E.2.24}}}
\expandafter\def\csname SRef@wbII:cor:surface-top-base-change\endcsname{\hyperlink{tag-03DV}{II.\ref*{wbII:cor:surface-top-base-change}}}
\expandafter\def\csname SRef@wbII:cor:candidate-pic0-route\endcsname{\hyperlink{tag-03DS}{II.\ref*{wbII:cor:candidate-pic0-route}}}
\expandafter\def\csname SRef@wbII:cor:candidate-conductor-actual\endcsname{\hyperlink{tag-03DR}{II.\ref*{wbII:cor:candidate-conductor-actual}}}
\expandafter\def\csname SRef@wbII:cor:corrected-cdp-relative\endcsname{\hyperlink{tag-03DT}{II.\ref*{wbII:cor:corrected-cdp-relative}}}
\expandafter\def\csname SRef@wbII:cor:relative\endcsname{\hyperlink{tag-03DT}{II.\ref*{wbII:cor:relative}}}
\expandafter\def\csname SRef@wbII:cor:normal-crossing-kernel\endcsname{\hyperlink{tag-03DU}{II.\ref*{wbII:cor:normal-crossing-kernel}}}
\expandafter\def\csname SRef@wbII:test:cdp-comparison\endcsname{\hyperlink{tag-03IV}{II.\ref*{wbII:test:cdp-comparison}}}
\expandafter\def\csname SRef@cdp:A:filtration1\endcsname{\hyperlink{tag-003W}{II.\ref*{cdp:A:filtration1}}}
\expandafter\def\csname SRef@cdp:eq:E13-residue\endcsname{\hyperlink{tag-01A6}{II.\ref*{cdp:eq:E13-residue}}}
\expandafter\def\csname SRef@cdp:eq:E14-elliptic-family\endcsname{\hyperlink{tag-01A8}{II.\ref*{cdp:eq:E14-elliptic-family}}}
\expandafter\def\csname SRef@cdp:eq:E14-canonical\endcsname{\hyperlink{tag-01A7}{II.\ref*{cdp:eq:E14-canonical}}}
\expandafter\def\csname SRef@cdp:eq:E14-log-coho\endcsname{\hyperlink{tag-01AA}{II.\ref*{cdp:eq:E14-log-coho}}}
\expandafter\def\csname SRef@cdp:eq:E14-jet-map\endcsname{\hyperlink{tag-01A9}{II.\ref*{cdp:eq:E14-jet-map}}}
\expandafter\def\csname SRef@cdp:eq:E15-positive\endcsname{\hyperlink{tag-01AB}{II.\ref*{cdp:eq:E15-positive}}}
\expandafter\def\csname SRef@cdp:eq:E15-products\endcsname{\hyperlink{tag-01AC}{II.\ref*{cdp:eq:E15-products}}}
\expandafter\def\csname SRef@wbII:eq:analytic-toric-open-directions\endcsname{\hyperlink{tag-03I5}{II.\ref*{wbII:eq:analytic-toric-open-directions}}}
\expandafter\def\csname SRef@wbII:eq:analytic-residual-descent-factor\endcsname{\hyperlink{tag-03I4}{II.\ref*{wbII:eq:analytic-residual-descent-factor}}}
\expandafter\def\csname SRef@wbII:eq:analytic-factor-horizontal-winding\endcsname{\hyperlink{tag-03I3}{II.\ref*{wbII:eq:analytic-factor-horizontal-winding}}}
\expandafter\def\csname SRef@cdp:A:filtration2\endcsname{\hyperlink{tag-003X}{II.\ref*{cdp:A:filtration2}}}
\expandafter\def\csname SRef@eq:P4chi-resolution\endcsname{\hyperlink{tag-01OD}{II.\ref*{eq:P4chi-resolution}}}
\expandafter\def\csname SRef@eq:P4exceptional-genus\endcsname{\hyperlink{tag-01OE}{II.\ref*{eq:P4exceptional-genus}}}
\expandafter\def\csname SRef@cdp:eq:E-Leray-map\endcsname{\hyperlink{tag-019Y}{II.\ref*{cdp:eq:E-Leray-map}}}
\expandafter\def\csname SRef@cdp:eq:E-surj-embedding\endcsname{\hyperlink{tag-01A0}{II.\ref*{cdp:eq:E-surj-embedding}}}
\expandafter\def\csname SRef@cdp:eq:E-log-period\endcsname{\hyperlink{tag-019Z}{II.\ref*{cdp:eq:E-log-period}}}
\expandafter\def\csname SRef@cdp:eq:E-theta-general\endcsname{\hyperlink{tag-01A1}{II.\ref*{cdp:eq:E-theta-general}}}
\expandafter\def\csname SRef@cdp:eq:E13-connection\endcsname{\hyperlink{tag-01A5}{II.\ref*{cdp:eq:E13-connection}}}
\expandafter\def\csname SRef@wbII:lem:torsion-identification\endcsname{\hyperlink{tag-03IF}{II.\ref*{wbII:lem:torsion-identification}}}
\expandafter\def\csname SRef@wbII:lem:picard-crossed-defect\endcsname{\hyperlink{tag-03IC}{II.\ref*{wbII:lem:picard-crossed-defect}}}
\expandafter\def\csname SRef@wbII:lem:picard-finite-norm\endcsname{\hyperlink{tag-03IE}{II.\ref*{wbII:lem:picard-finite-norm}}}
\expandafter\def\csname SRef@wbII:lem:picard-cusp-defect\endcsname{\hyperlink{tag-03ID}{II.\ref*{wbII:lem:picard-cusp-defect}}}
\expandafter\def\csname SRef@wbII:lem:analytic-support-specialization\endcsname{\hyperlink{tag-03I6}{II.\ref*{wbII:lem:analytic-support-specialization}}}
\expandafter\def\csname SRef@wbII:lem:generic-endpoints\endcsname{\hyperlink{tag-03IA}{II.\ref*{wbII:lem:generic-endpoints}}}
\expandafter\def\csname SRef@wbII:lem:cdp-generic-endpoints\endcsname{\hyperlink{tag-03I7}{II.\ref*{wbII:lem:cdp-generic-endpoints}}}
\expandafter\def\csname SRef@wbII:lem:gluing-bundle\endcsname{\hyperlink{tag-03IB}{II.\ref*{wbII:lem:gluing-bundle}}}
\expandafter\def\csname SRef@wbII:lem:extending-gluing\endcsname{\hyperlink{tag-03I9}{II.\ref*{wbII:lem:extending-gluing}}}
\expandafter\def\csname SRef@wbII:lem:curve-flatness\endcsname{\hyperlink{tag-03I8}{II.\ref*{wbII:lem:curve-flatness}}}
\expandafter\def\csname SRef@inv:E2.3\endcsname{\hyperlink{tag-01SW}{II.\ref*{inv:E2.3}}}
\expandafter\def\csname SRef@cdp:st:D.2.1\endcsname{\hyperlink{tag-01G8}{II.\ref*{cdp:st:D.2.1}}}
\expandafter\def\csname SRef@cdp:st:D.2.3\endcsname{\hyperlink{tag-01GA}{II.\ref*{cdp:st:D.2.3}}}
\expandafter\def\csname SRef@cdp:st:D.2.4\endcsname{\hyperlink{tag-01GB}{II.\ref*{cdp:st:D.2.4}}}
\expandafter\def\csname SRef@cdp:st:D.2.5\endcsname{\hyperlink{tag-01GC}{II.\ref*{cdp:st:D.2.5}}}
\expandafter\def\csname SRef@cdp:st:D.2.6\endcsname{\hyperlink{tag-01GD}{II.\ref*{cdp:st:D.2.6}}}
\expandafter\def\csname SRef@cdp:st:D.2.7\endcsname{\hyperlink{tag-01GE}{II.\ref*{cdp:st:D.2.7}}}
\expandafter\def\csname SRef@cdp:st:D.2.8\endcsname{\hyperlink{tag-01GF}{II.\ref*{cdp:st:D.2.8}}}
\expandafter\def\csname SRef@cdp:st:D.2.2\endcsname{\hyperlink{tag-01G9}{II.\ref*{cdp:st:D.2.2}}}
\expandafter\def\csname SRef@cdp:st:E5\endcsname{\hyperlink{tag-01GU}{II.\ref*{cdp:st:E5}}}
\expandafter\def\csname SRef@cdp:st:E15\endcsname{\hyperlink{tag-01GM}{II.\ref*{cdp:st:E15}}}
\expandafter\def\csname SRef@cdp:st:E16\endcsname{\hyperlink{tag-01GN}{II.\ref*{cdp:st:E16}}}
\expandafter\def\csname SRef@cdp:st:E17\endcsname{\hyperlink{tag-01GO}{II.\ref*{cdp:st:E17}}}
\expandafter\def\csname SRef@cdp:st:E18\endcsname{\hyperlink{tag-01GP}{II.\ref*{cdp:st:E18}}}
\expandafter\def\csname SRef@cdp:st:E6\endcsname{\hyperlink{tag-01GV}{II.\ref*{cdp:st:E6}}}
\expandafter\def\csname SRef@cdp:st:E7\endcsname{\hyperlink{tag-01GW}{II.\ref*{cdp:st:E7}}}
\expandafter\def\csname SRef@cdp:st:E8\endcsname{\hyperlink{tag-01GX}{II.\ref*{cdp:st:E8}}}
\expandafter\def\csname SRef@cdp:st:E10\endcsname{\hyperlink{tag-01GH}{II.\ref*{cdp:st:E10}}}
\expandafter\def\csname SRef@cdp:st:E11\endcsname{\hyperlink{tag-01GI}{II.\ref*{cdp:st:E11}}}
\expandafter\def\csname SRef@cdp:st:E12\endcsname{\hyperlink{tag-01GJ}{II.\ref*{cdp:st:E12}}}
\expandafter\def\csname SRef@cdp:st:E13\endcsname{\hyperlink{tag-01GK}{II.\ref*{cdp:st:E13}}}
\expandafter\def\csname SRef@cdp:st:E14\endcsname{\hyperlink{tag-01GL}{II.\ref*{cdp:st:E14}}}
\expandafter\def\csname SRef@push:src:E.2.13\endcsname{\hyperlink{tag-032Y}{II.\ref*{push:src:E.2.13}}}
\expandafter\def\csname SRef@push:src:E.2.14\endcsname{\hyperlink{tag-032Z}{II.\ref*{push:src:E.2.14}}}
\expandafter\def\csname SRef@push:src:E.2.15\endcsname{\hyperlink{tag-0330}{II.\ref*{push:src:E.2.15}}}
\expandafter\def\csname SRef@push:src:E.2.16\endcsname{\hyperlink{tag-0331}{II.\ref*{push:src:E.2.16}}}
\expandafter\def\csname SRef@push:src:E.2.17\endcsname{\hyperlink{tag-0332}{II.\ref*{push:src:E.2.17}}}
\expandafter\def\csname SRef@push:src:E.2.18\endcsname{\hyperlink{tag-0333}{II.\ref*{push:src:E.2.18}}}
\expandafter\def\csname SRef@push:src:E.2.19\endcsname{\hyperlink{tag-0334}{II.\ref*{push:src:E.2.19}}}
\expandafter\def\csname SRef@push:src:E.2.20\endcsname{\hyperlink{tag-0336}{II.\ref*{push:src:E.2.20}}}
\expandafter\def\csname SRef@push:src:E.2.21\endcsname{\hyperlink{tag-0337}{II.\ref*{push:src:E.2.21}}}
\expandafter\def\csname SRef@push:src:E.2.22\endcsname{\hyperlink{tag-0338}{II.\ref*{push:src:E.2.22}}}
\expandafter\def\csname SRef@push:src:E.2.5\endcsname{\hyperlink{tag-033D}{II.\ref*{push:src:E.2.5}}}
\expandafter\def\csname SRef@push:src:E.5.1\endcsname{\hyperlink{tag-033N}{II.\ref*{push:src:E.5.1}}}
\expandafter\def\csname SRef@push:src:E.5.2\endcsname{\hyperlink{tag-033O}{II.\ref*{push:src:E.5.2}}}
\expandafter\def\csname SRef@push:src:E.2.6\endcsname{\hyperlink{tag-033E}{II.\ref*{push:src:E.2.6}}}
\expandafter\def\csname SRef@push:src:E.2.7\endcsname{\hyperlink{tag-033F}{II.\ref*{push:src:E.2.7}}}
\expandafter\def\csname SRef@push:src:E.2.8\endcsname{\hyperlink{tag-033G}{II.\ref*{push:src:E.2.8}}}
\expandafter\def\csname SRef@push:src:E.2.9\endcsname{\hyperlink{tag-033H}{II.\ref*{push:src:E.2.9}}}
\expandafter\def\csname SRef@push:src:E.2.10\endcsname{\hyperlink{tag-032V}{II.\ref*{push:src:E.2.10}}}
\expandafter\def\csname SRef@push:src:E.2.11\endcsname{\hyperlink{tag-032W}{II.\ref*{push:src:E.2.11}}}
\expandafter\def\csname SRef@push:src:E.2.12\endcsname{\hyperlink{tag-032X}{II.\ref*{push:src:E.2.12}}}
\expandafter\def\csname SRef@inv:E2.2\endcsname{\hyperlink{tag-01SV}{II.\ref*{inv:E2.2}}}
\expandafter\def\csname SRef@inv:E2.4\endcsname{\hyperlink{tag-01SX}{II.\ref*{inv:E2.4}}}
\expandafter\def\csname SRef@wbII:prop:pic0-pullback\endcsname{\hyperlink{tag-03IO}{II.\ref*{wbII:prop:pic0-pullback}}}
\expandafter\def\csname SRef@wbII:prop:all-pic-restriction\endcsname{\hyperlink{tag-03IG}{II.\ref*{wbII:prop:all-pic-restriction}}}
\expandafter\def\csname SRef@wbII:prop:leray-edge-restriction\endcsname{\hyperlink{tag-03IM}{II.\ref*{wbII:prop:leray-edge-restriction}}}
\expandafter\def\csname SRef@wbII:prop:direct-boundary-picard\endcsname{\hyperlink{tag-03IL}{II.\ref*{wbII:prop:direct-boundary-picard}}}
\expandafter\def\csname SRef@wbII:prop:local-cusp-picard-obstruction\endcsname{\hyperlink{tag-03IN}{II.\ref*{wbII:prop:local-cusp-picard-obstruction}}}
\expandafter\def\csname SRef@wbII:prop:analytic-ah-boundary\endcsname{\hyperlink{tag-03IH}{II.\ref*{wbII:prop:analytic-ah-boundary}}}
\expandafter\def\csname SRef@wbII:prop:relative-local\endcsname{\hyperlink{tag-03IP}{II.\ref*{wbII:prop:relative-local}}}
\expandafter\def\csname SRef@wbII:prop:cartier-duality-absolute\endcsname{\hyperlink{tag-03IJ}{II.\ref*{wbII:prop:cartier-duality-absolute}}}
\expandafter\def\csname SRef@wbII:prop:analytic-base-change\endcsname{\hyperlink{tag-03II}{II.\ref*{wbII:prop:analytic-base-change}}}
\expandafter\def\csname SRef@cdp:st:C.1.1\endcsname{\hyperlink{tag-01BN}{II.\ref*{cdp:st:C.1.1}}}
\expandafter\def\csname SRef@cdp:st:C.1.3\endcsname{\hyperlink{tag-01BP}{II.\ref*{cdp:st:C.1.3}}}
\expandafter\def\csname SRef@cdp:sec:E\endcsname{\hyperlink{tag-01BP}{II.\ref*{cdp:sec:E}}}
\expandafter\def\csname SRef@cdp:st:E1\endcsname{\hyperlink{tag-01GG}{II.\ref*{cdp:st:E1}}}
\expandafter\def\csname SRef@cdp:st:D.1.2\endcsname{\hyperlink{tag-01G0}{II.\ref*{cdp:st:D.1.2}}}
\expandafter\def\csname SRef@cdp:st:D.1.4\endcsname{\hyperlink{tag-01G2}{II.\ref*{cdp:st:D.1.4}}}
\expandafter\def\csname SRef@cdp:st:D.1.13\endcsname{\hyperlink{tag-01FY}{II.\ref*{cdp:st:D.1.13}}}
\expandafter\def\csname SRef@cdp:st:D.1.15\endcsname{\hyperlink{tag-01FZ}{II.\ref*{cdp:st:D.1.15}}}
\expandafter\def\csname SRef@cdp:st:D.1.5\endcsname{\hyperlink{tag-01G3}{II.\ref*{cdp:st:D.1.5}}}
\expandafter\def\csname SRef@cdp:st:D.1.6\endcsname{\hyperlink{tag-01G4}{II.\ref*{cdp:st:D.1.6}}}
\expandafter\def\csname SRef@cdp:st:D.1.7\endcsname{\hyperlink{tag-01G5}{II.\ref*{cdp:st:D.1.7}}}
\expandafter\def\csname SRef@cdp:st:D.1.8\endcsname{\hyperlink{tag-01G6}{II.\ref*{cdp:st:D.1.8}}}
\expandafter\def\csname SRef@cdp:st:D.1.9\endcsname{\hyperlink{tag-01G7}{II.\ref*{cdp:st:D.1.9}}}
\expandafter\def\csname SRef@cdp:st:D.1.10\endcsname{\hyperlink{tag-01FV}{II.\ref*{cdp:st:D.1.10}}}
\expandafter\def\csname SRef@cdp:st:D.1.11\endcsname{\hyperlink{tag-01FW}{II.\ref*{cdp:st:D.1.11}}}
\expandafter\def\csname SRef@cdp:st:D.1.12\endcsname{\hyperlink{tag-01FX}{II.\ref*{cdp:st:D.1.12}}}
\expandafter\def\csname SRef@cdp:st:C.3.2\endcsname{\hyperlink{tag-01F8}{II.\ref*{cdp:st:C.3.2}}}
\expandafter\def\csname SRef@cdp:st:C.3.11\endcsname{\hyperlink{tag-01F6}{II.\ref*{cdp:st:C.3.11}}}
\expandafter\def\csname SRef@cdp:st:C.3.5\endcsname{\hyperlink{tag-01FB}{II.\ref*{cdp:st:C.3.5}}}
\expandafter\def\csname SRef@cdp:st:C.3.6\endcsname{\hyperlink{tag-01FC}{II.\ref*{cdp:st:C.3.6}}}
\expandafter\def\csname SRef@cdp:st:C.3.3\endcsname{\hyperlink{tag-01F9}{II.\ref*{cdp:st:C.3.3}}}
\expandafter\def\csname SRef@cdp:st:C.3.4\endcsname{\hyperlink{tag-01FA}{II.\ref*{cdp:st:C.3.4}}}
\expandafter\def\csname SRef@cdp:st:C.3.7\endcsname{\hyperlink{tag-01FD}{II.\ref*{cdp:st:C.3.7}}}
\expandafter\def\csname SRef@cdp:st:C.3.8\endcsname{\hyperlink{tag-01FE}{II.\ref*{cdp:st:C.3.8}}}
\expandafter\def\csname SRef@cdp:st:C.3.9\endcsname{\hyperlink{tag-01FF}{II.\ref*{cdp:st:C.3.9}}}
\expandafter\def\csname SRef@cdp:st:C.3.10\endcsname{\hyperlink{tag-01F5}{II.\ref*{cdp:st:C.3.10}}}
\expandafter\def\csname SRef@cdp:st:C.3.12\endcsname{\hyperlink{tag-01F7}{II.\ref*{cdp:st:C.3.12}}}
\expandafter\def\csname SRef@cdp:st:C.3.1\endcsname{\hyperlink{tag-01F4}{II.\ref*{cdp:st:C.3.1}}}
\expandafter\def\csname SRef@cdp:st:A.1.1\endcsname{\hyperlink{tag-01AD}{II.\ref*{cdp:st:A.1.1}}}
\expandafter\def\csname SRef@wbII:lem:conormal\endcsname{\hyperlink{tag-01AD}{II.\ref*{wbII:lem:conormal}}}
\expandafter\def\csname SRef@wbII:prop:criterion\endcsname{\hyperlink{tag-01AD}{II.\ref*{wbII:prop:criterion}}}
\expandafter\def\csname SRef@wbII:prop:filtration\endcsname{\hyperlink{tag-01AD}{II.\ref*{wbII:prop:filtration}}}
\expandafter\def\csname SRef@wbII:lem:conormal-injective\endcsname{\hyperlink{tag-01AD}{II.\ref*{wbII:lem:conormal-injective}}}
\expandafter\def\csname SRef@wbII:prop:missing-cotangent-filtration\endcsname{\hyperlink{tag-01AD}{II.\ref*{wbII:prop:missing-cotangent-filtration}}}
\expandafter\def\csname SRef@wbII:prop:missing-global-section-criterion\endcsname{\hyperlink{tag-01AD}{II.\ref*{wbII:prop:missing-global-section-criterion}}}
\expandafter\def\csname SRef@cdp:st:B3\endcsname{\hyperlink{tag-01BM}{II.\ref*{cdp:st:B3}}}
\expandafter\def\csname SRef@cdp:st:A.3.1\endcsname{\hyperlink{tag-01AL}{II.\ref*{cdp:st:A.3.1}}}
\expandafter\def\csname SRef@cdp:st:A.3.2\endcsname{\hyperlink{tag-01AM}{II.\ref*{cdp:st:A.3.2}}}
\expandafter\def\csname SRef@cdp:st:A.3.3\endcsname{\hyperlink{tag-01AN}{II.\ref*{cdp:st:A.3.3}}}
\expandafter\def\csname SRef@cdp:st:A.3.4\endcsname{\hyperlink{tag-01AO}{II.\ref*{cdp:st:A.3.4}}}
\expandafter\def\csname SRef@cdp:st:A.3.5\endcsname{\hyperlink{tag-01AP}{II.\ref*{cdp:st:A.3.5}}}
\expandafter\def\csname SRef@cdp:st:A.4.1\endcsname{\hyperlink{tag-01AQ}{II.\ref*{cdp:st:A.4.1}}}
\expandafter\def\csname SRef@cdp:st:A.4.2\endcsname{\hyperlink{tag-01AR}{II.\ref*{cdp:st:A.4.2}}}
\expandafter\def\csname SRef@cdp:st:A.4.3\endcsname{\hyperlink{tag-01AS}{II.\ref*{cdp:st:A.4.3}}}
\expandafter\def\csname SRef@cdp:st:A.4.4\endcsname{\hyperlink{tag-01AT}{II.\ref*{cdp:st:A.4.4}}}
\expandafter\def\csname SRef@cdp:st:A.2.1\endcsname{\hyperlink{tag-01AF}{II.\ref*{cdp:st:A.2.1}}}
\expandafter\def\csname SRef@cdp:st:A.4.5\endcsname{\hyperlink{tag-01AU}{II.\ref*{cdp:st:A.4.5}}}
\expandafter\def\csname SRef@cdp:st:A.4.6\endcsname{\hyperlink{tag-01AV}{II.\ref*{cdp:st:A.4.6}}}
\expandafter\def\csname SRef@cdp:st:A.4.7\endcsname{\hyperlink{tag-01AW}{II.\ref*{cdp:st:A.4.7}}}
\expandafter\def\csname SRef@cdp:st:A.4.8\endcsname{\hyperlink{tag-01AX}{II.\ref*{cdp:st:A.4.8}}}
\expandafter\def\csname SRef@cdp:st:A.4.9\endcsname{\hyperlink{tag-01AY}{II.\ref*{cdp:st:A.4.9}}}
\expandafter\def\csname SRef@cdp:st:A.2.2\endcsname{\hyperlink{tag-01AG}{II.\ref*{cdp:st:A.2.2}}}
\expandafter\def\csname SRef@wbII:lem:conductor\endcsname{\hyperlink{tag-01AG}{II.\ref*{wbII:lem:conductor}}}
\expandafter\def\csname SRef@wbII:lem:conductor-square\endcsname{\hyperlink{tag-01AG}{II.\ref*{wbII:lem:conductor-square}}}
\expandafter\def\csname SRef@cdp:st:A.1.2\endcsname{\hyperlink{tag-01AE}{II.\ref*{cdp:st:A.1.2}}}
\expandafter\def\csname SRef@cdp:st:A.2.4\endcsname{\hyperlink{tag-01AI}{II.\ref*{cdp:st:A.2.4}}}
\expandafter\def\csname SRef@cdp:st:A.2.5\endcsname{\hyperlink{tag-01AJ}{II.\ref*{cdp:st:A.2.5}}}
\expandafter\def\csname SRef@cdp:st:A.2.6\endcsname{\hyperlink{tag-01AK}{II.\ref*{cdp:st:A.2.6}}}
\expandafter\def\csname SRef@cdp:st:B1\endcsname{\hyperlink{tag-01BK}{II.\ref*{cdp:st:B1}}}
\expandafter\def\csname SRef@wbII:thm:nodal-test\endcsname{\hyperlink{tag-01BK}{II.\ref*{wbII:thm:nodal-test}}}
\expandafter\def\csname SRef@wbII:prop:family-geometry\endcsname{\hyperlink{tag-01BK}{II.\ref*{wbII:prop:family-geometry}}}
\expandafter\def\csname SRef@wbII:prop:nodal-counterexample\endcsname{\hyperlink{tag-01BK}{II.\ref*{wbII:prop:nodal-counterexample}}}
\expandafter\def\csname SRef@cdp:st:B2\endcsname{\hyperlink{tag-01BL}{II.\ref*{cdp:st:B2}}}
\expandafter\def\csname SRef@cdp:st:A.5.1\endcsname{\hyperlink{tag-01AZ}{II.\ref*{cdp:st:A.5.1}}}
\expandafter\def\csname SRef@cdp:st:A.5.5\endcsname{\hyperlink{tag-01BF}{II.\ref*{cdp:st:A.5.5}}}
\expandafter\def\csname SRef@cdp:st:A.5.7\endcsname{\hyperlink{tag-01BH}{II.\ref*{cdp:st:A.5.7}}}
\expandafter\def\csname SRef@cdp:st:A.5.8\endcsname{\hyperlink{tag-01BI}{II.\ref*{cdp:st:A.5.8}}}
\expandafter\def\csname SRef@cdp:st:A.5.14\endcsname{\hyperlink{tag-01B4}{II.\ref*{cdp:st:A.5.14}}}
\expandafter\def\csname SRef@cdp:st:A.5.10\endcsname{\hyperlink{tag-01B0}{II.\ref*{cdp:st:A.5.10}}}
\expandafter\def\csname SRef@cdp:st:A.5.20\endcsname{\hyperlink{tag-01BB}{II.\ref*{cdp:st:A.5.20}}}
\expandafter\def\csname SRef@cdp:st:A.5.15\endcsname{\hyperlink{tag-01B5}{II.\ref*{cdp:st:A.5.15}}}
\expandafter\def\csname SRef@cdp:st:A.5.17\endcsname{\hyperlink{tag-01B7}{II.\ref*{cdp:st:A.5.17}}}
\expandafter\def\csname SRef@cdp:st:A.5.18\endcsname{\hyperlink{tag-01B8}{II.\ref*{cdp:st:A.5.18}}}
\expandafter\def\csname SRef@cdp:st:A.5.19\endcsname{\hyperlink{tag-01B9}{II.\ref*{cdp:st:A.5.19}}}
\expandafter\def\csname SRef@cdp:st:A.5.2\endcsname{\hyperlink{tag-01BA}{II.\ref*{cdp:st:A.5.2}}}
\expandafter\def\csname SRef@cdp:st:A.5.16\endcsname{\hyperlink{tag-01B6}{II.\ref*{cdp:st:A.5.16}}}
\expandafter\def\csname SRef@def:II:minimal-boundary-data\endcsname{\hyperlink{tag-P209}{II.\ref*{def:II:minimal-boundary-data}}}
\expandafter\def\csname SRef@prop:II:minimal-boundary-complements\endcsname{\hyperlink{tag-P210}{II.\ref*{prop:II:minimal-boundary-complements}}}
\expandafter\def\csname SRef@cdp:st:A.5.21\endcsname{\hyperlink{tag-01BC}{II.\ref*{cdp:st:A.5.21}}}
\expandafter\def\csname SRef@cdp:st:A.5.3\endcsname{\hyperlink{tag-01BD}{II.\ref*{cdp:st:A.5.3}}}
\expandafter\def\csname SRef@cdp:st:A.5.4\endcsname{\hyperlink{tag-01BE}{II.\ref*{cdp:st:A.5.4}}}
\expandafter\def\csname SRef@cdp:st:A.5.6\endcsname{\hyperlink{tag-01BG}{II.\ref*{cdp:st:A.5.6}}}
\expandafter\def\csname SRef@cdp:st:A.5.9\endcsname{\hyperlink{tag-01BJ}{II.\ref*{cdp:st:A.5.9}}}
\expandafter\def\csname SRef@cdp:st:A.5.11\endcsname{\hyperlink{tag-01B1}{II.\ref*{cdp:st:A.5.11}}}
\expandafter\def\csname SRef@cdp:st:A.5.12\endcsname{\hyperlink{tag-01B2}{II.\ref*{cdp:st:A.5.12}}}
\expandafter\def\csname SRef@cdp:st:A.5.13\endcsname{\hyperlink{tag-01B3}{II.\ref*{cdp:st:A.5.13}}}
\expandafter\def\csname SRef@cdp:st:C.2.NC.2\endcsname{\hyperlink{tag-01C9}{II.\ref*{cdp:st:C.2.NC.2}}}
\expandafter\def\csname SRef@cdp:st:C.2.NC.3\endcsname{\hyperlink{tag-01CA}{II.\ref*{cdp:st:C.2.NC.3}}}
\expandafter\def\csname SRef@cdp:st:C.2.NC.4\endcsname{\hyperlink{tag-01CB}{II.\ref*{cdp:st:C.2.NC.4}}}
\expandafter\def\csname SRef@cdp:st:C.2.NC.5\endcsname{\hyperlink{tag-01CC}{II.\ref*{cdp:st:C.2.NC.5}}}
\expandafter\def\csname SRef@cdp:st:C2d\endcsname{\hyperlink{tag-01FJ}{II.\ref*{cdp:st:C2d}}}
\expandafter\def\csname SRef@night:P4.1\endcsname{\hyperlink{tag-02WG}{II.\ref*{night:P4.1}}}
\expandafter\def\csname SRef@stmt:P4.1\endcsname{\hyperlink{tag-02WG}{II.\ref*{stmt:P4.1}}}
\expandafter\def\csname SRef@night:P4.10\endcsname{\hyperlink{tag-02WH}{II.\ref*{night:P4.10}}}
\expandafter\def\csname SRef@stmt:P4.10\endcsname{\hyperlink{tag-02WH}{II.\ref*{stmt:P4.10}}}
\expandafter\def\csname SRef@night:P4.11\endcsname{\hyperlink{tag-02WI}{II.\ref*{night:P4.11}}}
\expandafter\def\csname SRef@stmt:P4.11\endcsname{\hyperlink{tag-02WI}{II.\ref*{stmt:P4.11}}}
\expandafter\def\csname SRef@night:P4.12\endcsname{\hyperlink{tag-02WJ}{II.\ref*{night:P4.12}}}
\expandafter\def\csname SRef@stmt:P4.12\endcsname{\hyperlink{tag-02WJ}{II.\ref*{stmt:P4.12}}}
\expandafter\def\csname SRef@night:P4.13\endcsname{\hyperlink{tag-02WK}{II.\ref*{night:P4.13}}}
\expandafter\def\csname SRef@stmt:P4.13\endcsname{\hyperlink{tag-02WK}{II.\ref*{stmt:P4.13}}}
\expandafter\def\csname SRef@night:P4.14\endcsname{\hyperlink{tag-02WL}{II.\ref*{night:P4.14}}}
\expandafter\def\csname SRef@stmt:P4.14\endcsname{\hyperlink{tag-02WL}{II.\ref*{stmt:P4.14}}}
\expandafter\def\csname SRef@night:P4.15\endcsname{\hyperlink{tag-02WM}{II.\ref*{night:P4.15}}}
\expandafter\def\csname SRef@stmt:P4.15\endcsname{\hyperlink{tag-02WM}{II.\ref*{stmt:P4.15}}}
\expandafter\def\csname SRef@night:P4.16\endcsname{\hyperlink{tag-02WN}{II.\ref*{night:P4.16}}}
\expandafter\def\csname SRef@stmt:P4.16\endcsname{\hyperlink{tag-02WN}{II.\ref*{stmt:P4.16}}}
\expandafter\def\csname SRef@night:P4.17\endcsname{\hyperlink{tag-02WO}{II.\ref*{night:P4.17}}}
\expandafter\def\csname SRef@stmt:P4.17\endcsname{\hyperlink{tag-02WO}{II.\ref*{stmt:P4.17}}}
\expandafter\def\csname SRef@night:P4.18\endcsname{\hyperlink{tag-02WP}{II.\ref*{night:P4.18}}}
\expandafter\def\csname SRef@stmt:P4.18\endcsname{\hyperlink{tag-02WP}{II.\ref*{stmt:P4.18}}}
\expandafter\def\csname SRef@night:P4.19\endcsname{\hyperlink{tag-02WQ}{II.\ref*{night:P4.19}}}
\expandafter\def\csname SRef@stmt:P4.19\endcsname{\hyperlink{tag-02WQ}{II.\ref*{stmt:P4.19}}}
\expandafter\def\csname SRef@night:P4.2\endcsname{\hyperlink{tag-02WR}{II.\ref*{night:P4.2}}}
\expandafter\def\csname SRef@stmt:P4.2\endcsname{\hyperlink{tag-02WR}{II.\ref*{stmt:P4.2}}}
\expandafter\def\csname SRef@night:P4.20\endcsname{\hyperlink{tag-02WS}{II.\ref*{night:P4.20}}}
\expandafter\def\csname SRef@stmt:P4.20\endcsname{\hyperlink{tag-02WS}{II.\ref*{stmt:P4.20}}}
\expandafter\def\csname SRef@night:P4.21\endcsname{\hyperlink{tag-02WT}{II.\ref*{night:P4.21}}}
\expandafter\def\csname SRef@stmt:P4.21\endcsname{\hyperlink{tag-02WT}{II.\ref*{stmt:P4.21}}}
\expandafter\def\csname SRef@night:P4.22\endcsname{\hyperlink{tag-02WU}{II.\ref*{night:P4.22}}}
\expandafter\def\csname SRef@stmt:P4.22\endcsname{\hyperlink{tag-02WU}{II.\ref*{stmt:P4.22}}}
\expandafter\def\csname SRef@night:P4.23\endcsname{\hyperlink{tag-02WV}{II.\ref*{night:P4.23}}}
\expandafter\def\csname SRef@stmt:P4.23\endcsname{\hyperlink{tag-02WV}{II.\ref*{stmt:P4.23}}}
\expandafter\def\csname SRef@night:P4.24\endcsname{\hyperlink{tag-02WW}{II.\ref*{night:P4.24}}}
\expandafter\def\csname SRef@stmt:P4.24\endcsname{\hyperlink{tag-02WW}{II.\ref*{stmt:P4.24}}}
\expandafter\def\csname SRef@night:P4.25\endcsname{\hyperlink{tag-02WX}{II.\ref*{night:P4.25}}}
\expandafter\def\csname SRef@stmt:P4.25\endcsname{\hyperlink{tag-02WX}{II.\ref*{stmt:P4.25}}}
\expandafter\def\csname SRef@night:P4.26\endcsname{\hyperlink{tag-02WY}{II.\ref*{night:P4.26}}}
\expandafter\def\csname SRef@stmt:P4.26\endcsname{\hyperlink{tag-02WY}{II.\ref*{stmt:P4.26}}}
\expandafter\def\csname SRef@night:P4.27\endcsname{\hyperlink{tag-02WZ}{II.\ref*{night:P4.27}}}
\expandafter\def\csname SRef@stmt:P4.27\endcsname{\hyperlink{tag-02WZ}{II.\ref*{stmt:P4.27}}}
\expandafter\def\csname SRef@night:P4.28\endcsname{\hyperlink{tag-02X0}{II.\ref*{night:P4.28}}}
\expandafter\def\csname SRef@stmt:P4.28\endcsname{\hyperlink{tag-02X0}{II.\ref*{stmt:P4.28}}}
\expandafter\def\csname SRef@night:P4.29\endcsname{\hyperlink{tag-02X1}{II.\ref*{night:P4.29}}}
\expandafter\def\csname SRef@stmt:P4.29\endcsname{\hyperlink{tag-02X1}{II.\ref*{stmt:P4.29}}}
\expandafter\def\csname SRef@night:P4.3\endcsname{\hyperlink{tag-02X2}{II.\ref*{night:P4.3}}}
\expandafter\def\csname SRef@stmt:P4.3\endcsname{\hyperlink{tag-02X2}{II.\ref*{stmt:P4.3}}}
\expandafter\def\csname SRef@night:P4.30\endcsname{\hyperlink{tag-02X3}{II.\ref*{night:P4.30}}}
\expandafter\def\csname SRef@stmt:P4.30\endcsname{\hyperlink{tag-02X3}{II.\ref*{stmt:P4.30}}}
\expandafter\def\csname SRef@night:P4.31\endcsname{\hyperlink{tag-02X4}{II.\ref*{night:P4.31}}}
\expandafter\def\csname SRef@stmt:P4.31\endcsname{\hyperlink{tag-02X4}{II.\ref*{stmt:P4.31}}}
\expandafter\def\csname SRef@night:P4.32\endcsname{\hyperlink{tag-02X5}{II.\ref*{night:P4.32}}}
\expandafter\def\csname SRef@stmt:P4.32\endcsname{\hyperlink{tag-02X5}{II.\ref*{stmt:P4.32}}}
\expandafter\def\csname SRef@night:P4.33\endcsname{\hyperlink{tag-02X6}{II.\ref*{night:P4.33}}}
\expandafter\def\csname SRef@stmt:P4.33\endcsname{\hyperlink{tag-02X6}{II.\ref*{stmt:P4.33}}}
\expandafter\def\csname SRef@night:P4.34\endcsname{\hyperlink{tag-02X7}{II.\ref*{night:P4.34}}}
\expandafter\def\csname SRef@stmt:P4.34\endcsname{\hyperlink{tag-02X7}{II.\ref*{stmt:P4.34}}}
\expandafter\def\csname SRef@night:P4.35\endcsname{\hyperlink{tag-02X8}{II.\ref*{night:P4.35}}}
\expandafter\def\csname SRef@stmt:P4.35\endcsname{\hyperlink{tag-02X8}{II.\ref*{stmt:P4.35}}}
\expandafter\def\csname SRef@night:P4.36\endcsname{\hyperlink{tag-02X9}{II.\ref*{night:P4.36}}}
\expandafter\def\csname SRef@stmt:P4.36\endcsname{\hyperlink{tag-02X9}{II.\ref*{stmt:P4.36}}}
\expandafter\def\csname SRef@night:P4.37\endcsname{\hyperlink{tag-02XA}{II.\ref*{night:P4.37}}}
\expandafter\def\csname SRef@stmt:P4.37\endcsname{\hyperlink{tag-02XA}{II.\ref*{stmt:P4.37}}}
\expandafter\def\csname SRef@night:P4.38\endcsname{\hyperlink{tag-02XB}{II.\ref*{night:P4.38}}}
\expandafter\def\csname SRef@stmt:P4.38\endcsname{\hyperlink{tag-02XB}{II.\ref*{stmt:P4.38}}}
\expandafter\def\csname SRef@night:P4.39\endcsname{\hyperlink{tag-02XC}{II.\ref*{night:P4.39}}}
\expandafter\def\csname SRef@stmt:P4.39\endcsname{\hyperlink{tag-02XC}{II.\ref*{stmt:P4.39}}}
\expandafter\def\csname SRef@night:P4.4\endcsname{\hyperlink{tag-02XD}{II.\ref*{night:P4.4}}}
\expandafter\def\csname SRef@stmt:P4.4\endcsname{\hyperlink{tag-02XD}{II.\ref*{stmt:P4.4}}}
\expandafter\def\csname SRef@night:P4.40\endcsname{\hyperlink{tag-02XE}{II.\ref*{night:P4.40}}}
\expandafter\def\csname SRef@stmt:P4.40\endcsname{\hyperlink{tag-02XE}{II.\ref*{stmt:P4.40}}}
\expandafter\def\csname SRef@night:P4.41\endcsname{\hyperlink{tag-02XF}{II.\ref*{night:P4.41}}}
\expandafter\def\csname SRef@stmt:P4.41\endcsname{\hyperlink{tag-02XF}{II.\ref*{stmt:P4.41}}}
\expandafter\def\csname SRef@night:P4.42\endcsname{\hyperlink{tag-02XG}{II.\ref*{night:P4.42}}}
\expandafter\def\csname SRef@stmt:P4.42\endcsname{\hyperlink{tag-02XG}{II.\ref*{stmt:P4.42}}}
\expandafter\def\csname SRef@night:P4.43\endcsname{\hyperlink{tag-02XH}{II.\ref*{night:P4.43}}}
\expandafter\def\csname SRef@stmt:P4.43\endcsname{\hyperlink{tag-02XH}{II.\ref*{stmt:P4.43}}}
\expandafter\def\csname SRef@night:P4.44\endcsname{\hyperlink{tag-02XI}{II.\ref*{night:P4.44}}}
\expandafter\def\csname SRef@stmt:P4.44\endcsname{\hyperlink{tag-02XI}{II.\ref*{stmt:P4.44}}}
\expandafter\def\csname SRef@night:P4.45\endcsname{\hyperlink{tag-02XJ}{II.\ref*{night:P4.45}}}
\expandafter\def\csname SRef@stmt:P4.45\endcsname{\hyperlink{tag-02XJ}{II.\ref*{stmt:P4.45}}}
\expandafter\def\csname SRef@night:P4.46\endcsname{\hyperlink{tag-02XK}{II.\ref*{night:P4.46}}}
\expandafter\def\csname SRef@stmt:P4.46\endcsname{\hyperlink{tag-02XK}{II.\ref*{stmt:P4.46}}}
\expandafter\def\csname SRef@night:P4.47\endcsname{\hyperlink{tag-02XL}{II.\ref*{night:P4.47}}}
\expandafter\def\csname SRef@stmt:P4.47\endcsname{\hyperlink{tag-02XL}{II.\ref*{stmt:P4.47}}}
\expandafter\def\csname SRef@night:P4.48\endcsname{\hyperlink{tag-02XM}{II.\ref*{night:P4.48}}}
\expandafter\def\csname SRef@stmt:P4.48\endcsname{\hyperlink{tag-02XM}{II.\ref*{stmt:P4.48}}}
\expandafter\def\csname SRef@night:P4.49\endcsname{\hyperlink{tag-02XN}{II.\ref*{night:P4.49}}}
\expandafter\def\csname SRef@stmt:P4.49\endcsname{\hyperlink{tag-02XN}{II.\ref*{stmt:P4.49}}}
\expandafter\def\csname SRef@night:P4.5\endcsname{\hyperlink{tag-02XO}{II.\ref*{night:P4.5}}}
\expandafter\def\csname SRef@stmt:P4.5\endcsname{\hyperlink{tag-02XO}{II.\ref*{stmt:P4.5}}}
\expandafter\def\csname SRef@night:P4.50\endcsname{\hyperlink{tag-02XP}{II.\ref*{night:P4.50}}}
\expandafter\def\csname SRef@stmt:P4.50\endcsname{\hyperlink{tag-02XP}{II.\ref*{stmt:P4.50}}}
\expandafter\def\csname SRef@night:P4.51\endcsname{\hyperlink{tag-02XQ}{II.\ref*{night:P4.51}}}
\expandafter\def\csname SRef@stmt:P4.51\endcsname{\hyperlink{tag-02XQ}{II.\ref*{stmt:P4.51}}}
\expandafter\def\csname SRef@night:P4.52\endcsname{\hyperlink{tag-02XR}{II.\ref*{night:P4.52}}}
\expandafter\def\csname SRef@stmt:P4.52\endcsname{\hyperlink{tag-02XR}{II.\ref*{stmt:P4.52}}}
\expandafter\def\csname SRef@night:P4.53\endcsname{\hyperlink{tag-02XS}{II.\ref*{night:P4.53}}}
\expandafter\def\csname SRef@stmt:P4.53\endcsname{\hyperlink{tag-02XS}{II.\ref*{stmt:P4.53}}}
\expandafter\def\csname SRef@night:P4.54\endcsname{\hyperlink{tag-02XT}{II.\ref*{night:P4.54}}}
\expandafter\def\csname SRef@stmt:P4.54\endcsname{\hyperlink{tag-02XT}{II.\ref*{stmt:P4.54}}}
\expandafter\def\csname SRef@night:P4.55\endcsname{\hyperlink{tag-02XU}{II.\ref*{night:P4.55}}}
\expandafter\def\csname SRef@stmt:P4.55\endcsname{\hyperlink{tag-02XU}{II.\ref*{stmt:P4.55}}}
\expandafter\def\csname SRef@night:P4.56\endcsname{\hyperlink{tag-02XV}{II.\ref*{night:P4.56}}}
\expandafter\def\csname SRef@stmt:P4.56\endcsname{\hyperlink{tag-02XV}{II.\ref*{stmt:P4.56}}}
\expandafter\def\csname SRef@night:P4.57\endcsname{\hyperlink{tag-02XW}{II.\ref*{night:P4.57}}}
\expandafter\def\csname SRef@stmt:P4.57\endcsname{\hyperlink{tag-02XW}{II.\ref*{stmt:P4.57}}}
\expandafter\def\csname SRef@night:P4.58\endcsname{\hyperlink{tag-02XX}{II.\ref*{night:P4.58}}}
\expandafter\def\csname SRef@stmt:P4.58\endcsname{\hyperlink{tag-02XX}{II.\ref*{stmt:P4.58}}}
\expandafter\def\csname SRef@night:P4R.1\endcsname{\hyperlink{tag-02Y2}{II.\ref*{night:P4R.1}}}
\expandafter\def\csname SRef@stmt:P4R.1\endcsname{\hyperlink{tag-02Y2}{II.\ref*{stmt:P4R.1}}}
\expandafter\def\csname SRef@night:P4.6\endcsname{\hyperlink{tag-02XY}{II.\ref*{night:P4.6}}}
\expandafter\def\csname SRef@stmt:P4.6\endcsname{\hyperlink{tag-02XY}{II.\ref*{stmt:P4.6}}}
\expandafter\def\csname SRef@night:P4R.2\endcsname{\hyperlink{tag-02YD}{II.\ref*{night:P4R.2}}}
\expandafter\def\csname SRef@stmt:P4R.2\endcsname{\hyperlink{tag-02YD}{II.\ref*{stmt:P4R.2}}}
\expandafter\def\csname SRef@night:P4R.3\endcsname{\hyperlink{tag-02YF}{II.\ref*{night:P4R.3}}}
\expandafter\def\csname SRef@stmt:P4R.3\endcsname{\hyperlink{tag-02YF}{II.\ref*{stmt:P4R.3}}}
\expandafter\def\csname SRef@night:P4R.4\endcsname{\hyperlink{tag-02YG}{II.\ref*{night:P4R.4}}}
\expandafter\def\csname SRef@stmt:P4R.4\endcsname{\hyperlink{tag-02YG}{II.\ref*{stmt:P4R.4}}}
\expandafter\def\csname SRef@night:P4R.5\endcsname{\hyperlink{tag-02YH}{II.\ref*{night:P4R.5}}}
\expandafter\def\csname SRef@stmt:P4R.5\endcsname{\hyperlink{tag-02YH}{II.\ref*{stmt:P4R.5}}}
\expandafter\def\csname SRef@night:P4R.6\endcsname{\hyperlink{tag-02YI}{II.\ref*{night:P4R.6}}}
\expandafter\def\csname SRef@stmt:P4R.6\endcsname{\hyperlink{tag-02YI}{II.\ref*{stmt:P4R.6}}}
\expandafter\def\csname SRef@night:P4R.7\endcsname{\hyperlink{tag-02YJ}{II.\ref*{night:P4R.7}}}
\expandafter\def\csname SRef@stmt:P4R.7\endcsname{\hyperlink{tag-02YJ}{II.\ref*{stmt:P4R.7}}}
\expandafter\def\csname SRef@night:P4R.8\endcsname{\hyperlink{tag-02YK}{II.\ref*{night:P4R.8}}}
\expandafter\def\csname SRef@stmt:P4R.8\endcsname{\hyperlink{tag-02YK}{II.\ref*{stmt:P4R.8}}}
\expandafter\def\csname SRef@night:P4R.9\endcsname{\hyperlink{tag-02YL}{II.\ref*{night:P4R.9}}}
\expandafter\def\csname SRef@stmt:P4R.9\endcsname{\hyperlink{tag-02YL}{II.\ref*{stmt:P4R.9}}}
\expandafter\def\csname SRef@night:P4R.10\endcsname{\hyperlink{tag-02Y3}{II.\ref*{night:P4R.10}}}
\expandafter\def\csname SRef@stmt:P4R.10\endcsname{\hyperlink{tag-02Y3}{II.\ref*{stmt:P4R.10}}}
\expandafter\def\csname SRef@night:P4R.11\endcsname{\hyperlink{tag-02Y4}{II.\ref*{night:P4R.11}}}
\expandafter\def\csname SRef@stmt:P4R.11\endcsname{\hyperlink{tag-02Y4}{II.\ref*{stmt:P4R.11}}}
\expandafter\def\csname SRef@night:P4.7\endcsname{\hyperlink{tag-02XZ}{II.\ref*{night:P4.7}}}
\expandafter\def\csname SRef@stmt:P4.7\endcsname{\hyperlink{tag-02XZ}{II.\ref*{stmt:P4.7}}}
\expandafter\def\csname SRef@night:P4R.12\endcsname{\hyperlink{tag-02Y5}{II.\ref*{night:P4R.12}}}
\expandafter\def\csname SRef@stmt:P4R.12\endcsname{\hyperlink{tag-02Y5}{II.\ref*{stmt:P4R.12}}}
\expandafter\def\csname SRef@night:P4R.13\endcsname{\hyperlink{tag-02Y6}{II.\ref*{night:P4R.13}}}
\expandafter\def\csname SRef@stmt:P4R.13\endcsname{\hyperlink{tag-02Y6}{II.\ref*{stmt:P4R.13}}}
\expandafter\def\csname SRef@night:P4R.14\endcsname{\hyperlink{tag-02Y7}{II.\ref*{night:P4R.14}}}
\expandafter\def\csname SRef@stmt:P4R.14\endcsname{\hyperlink{tag-02Y7}{II.\ref*{stmt:P4R.14}}}
\expandafter\def\csname SRef@night:P4R.15\endcsname{\hyperlink{tag-02Y8}{II.\ref*{night:P4R.15}}}
\expandafter\def\csname SRef@stmt:P4R.15\endcsname{\hyperlink{tag-02Y8}{II.\ref*{stmt:P4R.15}}}
\expandafter\def\csname SRef@night:P4R.16\endcsname{\hyperlink{tag-02Y9}{II.\ref*{night:P4R.16}}}
\expandafter\def\csname SRef@stmt:P4R.16\endcsname{\hyperlink{tag-02Y9}{II.\ref*{stmt:P4R.16}}}
\expandafter\def\csname SRef@night:P4R.17\endcsname{\hyperlink{tag-02YA}{II.\ref*{night:P4R.17}}}
\expandafter\def\csname SRef@stmt:P4R.17\endcsname{\hyperlink{tag-02YA}{II.\ref*{stmt:P4R.17}}}
\expandafter\def\csname SRef@bounds-without-an-ordinary-square-assumption\endcsname{\hyperlink{tag-02YA}{II.\ref*{bounds-without-an-ordinary-square-assumption}}}
\expandafter\def\csname SRef@night:P4R.18\endcsname{\hyperlink{tag-02YB}{II.\ref*{night:P4R.18}}}
\expandafter\def\csname SRef@stmt:P4R.18\endcsname{\hyperlink{tag-02YB}{II.\ref*{stmt:P4R.18}}}
\expandafter\def\csname SRef@night:P4R.19\endcsname{\hyperlink{tag-02YC}{II.\ref*{night:P4R.19}}}
\expandafter\def\csname SRef@stmt:P4R.19\endcsname{\hyperlink{tag-02YC}{II.\ref*{stmt:P4R.19}}}
\expandafter\def\csname SRef@night:P4R.20\endcsname{\hyperlink{tag-02YE}{II.\ref*{night:P4R.20}}}
\expandafter\def\csname SRef@stmt:P4R.20\endcsname{\hyperlink{tag-02YE}{II.\ref*{stmt:P4R.20}}}
\expandafter\def\csname SRef@night:P4.8\endcsname{\hyperlink{tag-02Y0}{II.\ref*{night:P4.8}}}
\expandafter\def\csname SRef@stmt:P4.8\endcsname{\hyperlink{tag-02Y0}{II.\ref*{stmt:P4.8}}}
\expandafter\def\csname SRef@night:P4.9\endcsname{\hyperlink{tag-02Y1}{II.\ref*{night:P4.9}}}
\expandafter\def\csname SRef@stmt:P4.9\endcsname{\hyperlink{tag-02Y1}{II.\ref*{stmt:P4.9}}}
\expandafter\def\csname SRef@cdp:st:C4a\endcsname{\hyperlink{tag-01FK}{II.\ref*{cdp:st:C4a}}}
\expandafter\def\csname SRef@cdp:st:C5c\endcsname{\hyperlink{tag-01FT}{II.\ref*{cdp:st:C5c}}}
\expandafter\def\csname SRef@cdp:st:C4b\endcsname{\hyperlink{tag-01FL}{II.\ref*{cdp:st:C4b}}}
\expandafter\def\csname SRef@cdp:st:C4c\endcsname{\hyperlink{tag-01FM}{II.\ref*{cdp:st:C4c}}}
\expandafter\def\csname SRef@cdp:st:C4d\endcsname{\hyperlink{tag-01FN}{II.\ref*{cdp:st:C4d}}}
\expandafter\def\csname SRef@cdp:st:C4e\endcsname{\hyperlink{tag-01FO}{II.\ref*{cdp:st:C4e}}}
\expandafter\def\csname SRef@cdp:st:C4f\endcsname{\hyperlink{tag-01FP}{II.\ref*{cdp:st:C4f}}}
\expandafter\def\csname SRef@cdp:st:C4g\endcsname{\hyperlink{tag-01FQ}{II.\ref*{cdp:st:C4g}}}
\expandafter\def\csname SRef@cdp:st:C5a\endcsname{\hyperlink{tag-01FR}{II.\ref*{cdp:st:C5a}}}
\expandafter\def\csname SRef@cdp:st:C5b\endcsname{\hyperlink{tag-01FS}{II.\ref*{cdp:st:C5b}}}
\expandafter\def\csname SRef@night:P3K.36\endcsname{\hyperlink{tag-02UH}{II.\ref*{night:P3K.36}}}
\expandafter\def\csname SRef@stmt:P3K.36\endcsname{\hyperlink{tag-02UH}{II.\ref*{stmt:P3K.36}}}
\expandafter\def\csname SRef@inv:wb:uniform-nontorsion-provenance\endcsname{\hyperlink{tag-01TQ}{II.\ref*{inv:wb:uniform-nontorsion-provenance}}}
\expandafter\def\csname SRef@wbII:rem:scope\endcsname{\hyperlink{tag-03IS}{II.\ref*{wbII:rem:scope}}}
\expandafter\def\csname SRef@wbII:rem:local-cusp-picard-obstruction\endcsname{\hyperlink{tag-03IQ}{II.\ref*{wbII:rem:local-cusp-picard-obstruction}}}
\expandafter\def\csname SRef@wbII:rem:relative-local-scope\endcsname{\hyperlink{tag-03IR}{II.\ref*{wbII:rem:relative-local-scope}}}
\expandafter\def\csname SRef@sec:II:conventions\endcsname{\hyperlink{tag-034E}{II.\ref*{sec:II:conventions}}}
\expandafter\def\csname SRef@sec:II:13\endcsname{\hyperlink{tag-0349}{II.\ref*{sec:II:13}}}
\expandafter\def\csname SRef@sec:II:14\endcsname{\hyperlink{tag-034A}{II.\ref*{sec:II:14}}}
\expandafter\def\csname SRef@sec:II:push-stability\endcsname{\hyperlink{tag-034G}{II.\ref*{sec:II:push-stability}}}
\expandafter\def\csname SRef@sec:II:combs\endcsname{\hyperlink{tag-034D}{II.\ref*{sec:II:combs}}}
\expandafter\def\csname SRef@sec:II:uniform-nontorsion\endcsname{\hyperlink{tag-034H}{II.\ref*{sec:II:uniform-nontorsion}}}
\expandafter\def\csname SRef@app:II:programs\endcsname{\hyperlink{tag-000B}{II.\ref*{app:II:programs}}}
\expandafter\def\csname SRef@sec:II:01\endcsname{\hyperlink{tag-033X}{II.\ref*{sec:II:01}}}
\expandafter\def\csname SRef@sec:II:02\endcsname{\hyperlink{tag-033Y}{II.\ref*{sec:II:02}}}
\expandafter\def\csname SRef@sec:II:03\endcsname{\hyperlink{tag-033Z}{II.\ref*{sec:II:03}}}
\expandafter\def\csname SRef@sec:II:04\endcsname{\hyperlink{tag-0340}{II.\ref*{sec:II:04}}}
\expandafter\def\csname SRef@sec:II:05\endcsname{\hyperlink{tag-0341}{II.\ref*{sec:II:05}}}
\expandafter\def\csname SRef@sec:II:08\endcsname{\hyperlink{tag-0344}{II.\ref*{sec:II:08}}}
\expandafter\def\csname SRef@sec:II:night-sign\endcsname{\hyperlink{tag-034F}{II.\ref*{sec:II:night-sign}}}
\expandafter\def\csname SRef@sec:II:12\endcsname{\hyperlink{tag-0348}{II.\ref*{sec:II:12}}}
\expandafter\def\csname SRef@wbII:subsec:local-cusp-picard-obstruction\endcsname{\hyperlink{tag-03IU}{II.\ref*{wbII:subsec:local-cusp-picard-obstruction}}}
\expandafter\def\csname SRef@wbII:subsec:analytic-all-picard\endcsname{\hyperlink{tag-03IT}{II.\ref*{wbII:subsec:analytic-all-picard}}}
\expandafter\def\csname SRef@push:src:E.2.23\endcsname{\hyperlink{tag-0339}{II.\ref*{push:src:E.2.23}}}
\expandafter\def\csname SRef@push:src:E.5.3\endcsname{\hyperlink{tag-033P}{II.\ref*{push:src:E.5.3}}}
\expandafter\def\csname SRef@wb:thm:uniform-nontorsion-audit\endcsname{\hyperlink{tag-03DQ}{II.\ref*{wb:thm:uniform-nontorsion-audit}}}
\expandafter\def\csname SRef@wbII:prop:conductor\endcsname{\hyperlink{tag-03IK}{II.\ref*{wbII:prop:conductor}}}
\expandafter\def\csname SRef@wbII:prop:conductor-differential\endcsname{\hyperlink{tag-03IK}{II.\ref*{wbII:prop:conductor-differential}}}
\expandafter\def\csname SRef@wbII:thm:analytic-all-picard\endcsname{\hyperlink{tag-03IW}{II.\ref*{wbII:thm:analytic-all-picard}}}
\expandafter\def\csname SRef@wbII:thm:candidate-conductor\endcsname{\hyperlink{tag-03IX}{II.\ref*{wbII:thm:candidate-conductor}}}
\expandafter\def\csname SRef@wbII:thm:corrected-cdp-global\endcsname{\hyperlink{tag-03J0}{II.\ref*{wbII:thm:corrected-cdp-global}}}
\expandafter\def\csname SRef@wbII:thm:corrected-global\endcsname{\hyperlink{tag-03J0}{II.\ref*{wbII:thm:corrected-global}}}
\expandafter\def\csname SRef@wbII:thm:corrected-cdp-fibrewise\endcsname{\hyperlink{tag-03IZ}{II.\ref*{wbII:thm:corrected-cdp-fibrewise}}}
\expandafter\def\csname SRef@wbII:thm:cartier-duality\endcsname{\hyperlink{tag-03IY}{II.\ref*{wbII:thm:cartier-duality}}}
\expandafter\def\csname SRef@ii:LR:covering\endcsname{\hyperlink{tag-01RO}{II.\ref*{ii:LR:covering}}}
\expandafter\def\csname SRef@maps:eq:D.O9\endcsname{\href{S6_SERIES_GRIND_III.pdf\#nameddest=tag-020U}{III.{D.O9} (p.~85)} }
\expandafter\def\csname SRef@maps:eq:D.FI1\endcsname{\href{S6_SERIES_GRIND_III.pdf\#nameddest=tag-01YT}{III.{D.FI1} (p.~87)} }
\expandafter\def\csname SRef@maps:eq:D.FI2\endcsname{\href{S6_SERIES_GRIND_III.pdf\#nameddest=tag-01YU}{III.{D.FI2} (p.~88)} }
\expandafter\def\csname SRef@maps:eq:D.FI3\endcsname{\href{S6_SERIES_GRIND_III.pdf\#nameddest=tag-01YV}{III.{D.FI3} (p.~88)} }
\expandafter\def\csname SRef@maps:eq:D.FI4\endcsname{\href{S6_SERIES_GRIND_III.pdf\#nameddest=tag-01YW}{III.{D.FI4} (p.~90)} }
\expandafter\def\csname SRef@maps:eq:D.FI5\endcsname{\href{S6_SERIES_GRIND_III.pdf\#nameddest=tag-01YX}{III.{D.FI5} (p.~91)} }
\expandafter\def\csname SRef@maps:eq:D.FG1\endcsname{\href{S6_SERIES_GRIND_III.pdf\#nameddest=tag-01YQ}{III.{D.FG1} (p.~92)} }
\expandafter\def\csname SRef@maps:eq:D.FG2\endcsname{\href{S6_SERIES_GRIND_III.pdf\#nameddest=tag-01YR}{III.{D.FG2} (p.~92)} }
\expandafter\def\csname SRef@maps:eq:D.FG3\endcsname{\href{S6_SERIES_GRIND_III.pdf\#nameddest=tag-01YS}{III.{D.FG3} (p.~92)} }
\expandafter\def\csname SRef@maps:eq:D.IS1\endcsname{\href{S6_SERIES_GRIND_III.pdf\#nameddest=tag-0207}{III.{D.IS1} (p.~93)} }
\expandafter\def\csname SRef@maps:eq:D.IS2\endcsname{\href{S6_SERIES_GRIND_III.pdf\#nameddest=tag-0208}{III.{D.IS2} (p.~93)} }
\expandafter\def\csname SRef@maps:eq:D.SC1\endcsname{\href{S6_SERIES_GRIND_III.pdf\#nameddest=tag-023Z}{III.{D.SC1} (p.~94)} }
\expandafter\def\csname SRef@maps:eq:D.SC2\endcsname{\href{S6_SERIES_GRIND_III.pdf\#nameddest=tag-0240}{III.{D.SC2} (p.~95)} }
\expandafter\def\csname SRef@maps:eq:D.SC3\endcsname{\href{S6_SERIES_GRIND_III.pdf\#nameddest=tag-0241}{III.{D.SC3} (p.~95)} }
\expandafter\def\csname SRef@maps:eq:D.SC4\endcsname{\href{S6_SERIES_GRIND_III.pdf\#nameddest=tag-0242}{III.{D.SC4} (p.~95)} }
\expandafter\def\csname SRef@maps:eq:D.SC5\endcsname{\href{S6_SERIES_GRIND_III.pdf\#nameddest=tag-0243}{III.{D.SC5} (p.~95)} }
\expandafter\def\csname SRef@maps:eq:D.SC6\endcsname{\href{S6_SERIES_GRIND_III.pdf\#nameddest=tag-0244}{III.{D.SC6} (p.~95)} }
\expandafter\def\csname SRef@maps:eq:D.SC7\endcsname{\href{S6_SERIES_GRIND_III.pdf\#nameddest=tag-0245}{III.{D.SC7} (p.~97)} }
\expandafter\def\csname SRef@maps:eq:D.SC8\endcsname{\href{S6_SERIES_GRIND_III.pdf\#nameddest=tag-0246}{III.{D.SC8} (p.~97)} }
\expandafter\def\csname SRef@maps:eq:D.HY1\endcsname{\href{S6_SERIES_GRIND_III.pdf\#nameddest=tag-01ZN}{III.{D.HY1} (p.~98)} }
\expandafter\def\csname SRef@maps:eq:D.HY2\endcsname{\href{S6_SERIES_GRIND_III.pdf\#nameddest=tag-01ZO}{III.{D.HY2} (p.~98)} }
\expandafter\def\csname SRef@maps:eq:D.HY3\endcsname{\href{S6_SERIES_GRIND_III.pdf\#nameddest=tag-01ZP}{III.{D.HY3} (p.~99)} }
\expandafter\def\csname SRef@maps:eq:D.HY4\endcsname{\href{S6_SERIES_GRIND_III.pdf\#nameddest=tag-01ZQ}{III.{D.HY4} (p.~99)} }
\expandafter\def\csname SRef@maps:eq:D.HY5\endcsname{\href{S6_SERIES_GRIND_III.pdf\#nameddest=tag-01ZR}{III.{D.HY5} (p.~99)} }
\expandafter\def\csname SRef@maps:eq:D.HY6\endcsname{\href{S6_SERIES_GRIND_III.pdf\#nameddest=tag-01ZS}{III.{D.HY6} (p.~100)} }
\expandafter\def\csname SRef@maps:eq:D.HY7\endcsname{\href{S6_SERIES_GRIND_III.pdf\#nameddest=tag-01ZT}{III.{D.HY7} (p.~100)} }
\expandafter\def\csname SRef@maps:eq:D.ST1\endcsname{\href{S6_SERIES_GRIND_III.pdf\#nameddest=tag-024V}{III.{D.ST1} (p.~102)} }
\expandafter\def\csname SRef@maps:eq:D.ST2\endcsname{\href{S6_SERIES_GRIND_III.pdf\#nameddest=tag-0250}{III.{D.ST2} (p.~103)} }
\expandafter\def\csname SRef@maps:eq:D.ST3\endcsname{\href{S6_SERIES_GRIND_III.pdf\#nameddest=tag-0251}{III.{D.ST3} (p.~103)} }
\expandafter\def\csname SRef@maps:eq:D.ST4\endcsname{\href{S6_SERIES_GRIND_III.pdf\#nameddest=tag-0252}{III.{D.ST4} (p.~103)} }
\expandafter\def\csname SRef@maps:eq:D.ST5\endcsname{\href{S6_SERIES_GRIND_III.pdf\#nameddest=tag-0253}{III.{D.ST5} (p.~103)} }
\expandafter\def\csname SRef@maps:eq:D.ST6\endcsname{\href{S6_SERIES_GRIND_III.pdf\#nameddest=tag-0254}{III.{D.ST6} (p.~103)} }
\expandafter\def\csname SRef@maps:eq:D.ST7\endcsname{\href{S6_SERIES_GRIND_III.pdf\#nameddest=tag-0255}{III.{D.ST7} (p.~104)} }
\expandafter\def\csname SRef@maps:eq:D.ST8\endcsname{\href{S6_SERIES_GRIND_III.pdf\#nameddest=tag-0256}{III.{D.ST8} (p.~104)} }
\expandafter\def\csname SRef@maps:eq:D.ST9\endcsname{\href{S6_SERIES_GRIND_III.pdf\#nameddest=tag-0257}{III.{D.ST9} (p.~104)} }
\expandafter\def\csname SRef@maps:eq:D.ST10\endcsname{\href{S6_SERIES_GRIND_III.pdf\#nameddest=tag-024W}{III.{D.ST10} (p.~105)} }
\expandafter\def\csname SRef@maps:eq:D.ST11\endcsname{\href{S6_SERIES_GRIND_III.pdf\#nameddest=tag-024X}{III.{D.ST11} (p.~105)} }
\expandafter\def\csname SRef@maps:eq:D.ST12\endcsname{\href{S6_SERIES_GRIND_III.pdf\#nameddest=tag-024Y}{III.{D.ST12} (p.~105)} }
\expandafter\def\csname SRef@maps:eq:D.ST13\endcsname{\href{S6_SERIES_GRIND_III.pdf\#nameddest=tag-024Z}{III.{D.ST13} (p.~106)} }
\expandafter\def\csname SRef@maps:eq:D.UF1\endcsname{\href{S6_SERIES_GRIND_III.pdf\#nameddest=tag-025Z}{III.{D.UF1} (p.~107)} }
\expandafter\def\csname SRef@maps:eq:D.UF2\endcsname{\href{S6_SERIES_GRIND_III.pdf\#nameddest=tag-0261}{III.{D.UF2} (p.~107)} }
\expandafter\def\csname SRef@maps:eq:D.UF3\endcsname{\href{S6_SERIES_GRIND_III.pdf\#nameddest=tag-0262}{III.{D.UF3} (p.~107)} }
\expandafter\def\csname SRef@maps:eq:D.UF4\endcsname{\href{S6_SERIES_GRIND_III.pdf\#nameddest=tag-0263}{III.{D.UF4} (p.~107)} }
\expandafter\def\csname SRef@maps:eq:D.UF5\endcsname{\href{S6_SERIES_GRIND_III.pdf\#nameddest=tag-0264}{III.{D.UF5} (p.~108)} }
\expandafter\def\csname SRef@maps:eq:D.UF6\endcsname{\href{S6_SERIES_GRIND_III.pdf\#nameddest=tag-0265}{III.{D.UF6} (p.~108)} }
\expandafter\def\csname SRef@maps:eq:D.UF7\endcsname{\href{S6_SERIES_GRIND_III.pdf\#nameddest=tag-0266}{III.{D.UF7} (p.~108)} }
\expandafter\def\csname SRef@maps:eq:D.UF8\endcsname{\href{S6_SERIES_GRIND_III.pdf\#nameddest=tag-0267}{III.{D.UF8} (p.~109)} }
\expandafter\def\csname SRef@maps:eq:D.UF9\endcsname{\href{S6_SERIES_GRIND_III.pdf\#nameddest=tag-0268}{III.{D.UF9} (p.~109)} }
\expandafter\def\csname SRef@maps:eq:D.UF10\endcsname{\href{S6_SERIES_GRIND_III.pdf\#nameddest=tag-0260}{III.{D.UF10} (p.~109)} }
\expandafter\def\csname SRef@maps:eq:D.AF1\endcsname{\href{S6_SERIES_GRIND_III.pdf\#nameddest=tag-01XI}{III.{D.AF1} (p.~110)} }
\expandafter\def\csname SRef@maps:eq:D.AF2\endcsname{\href{S6_SERIES_GRIND_III.pdf\#nameddest=tag-01XJ}{III.{D.AF2} (p.~110)} }
\expandafter\def\csname SRef@maps:eq:D.AF3\endcsname{\href{S6_SERIES_GRIND_III.pdf\#nameddest=tag-01XK}{III.{D.AF3} (p.~111)} }
\expandafter\def\csname SRef@maps:eq:D.AF4\endcsname{\href{S6_SERIES_GRIND_III.pdf\#nameddest=tag-01XL}{III.{D.AF4} (p.~111)} }
\expandafter\def\csname SRef@maps:eq:D.AF5\endcsname{\href{S6_SERIES_GRIND_III.pdf\#nameddest=tag-01XM}{III.{D.AF5} (p.~111)} }
\expandafter\def\csname SRef@maps:eq:D.AF6\endcsname{\href{S6_SERIES_GRIND_III.pdf\#nameddest=tag-01XN}{III.{D.AF6} (p.~111)} }
\expandafter\def\csname SRef@maps:eq:D.AF7\endcsname{\href{S6_SERIES_GRIND_III.pdf\#nameddest=tag-01XO}{III.{D.AF7} (p.~111)} }
\expandafter\def\csname SRef@maps:eq:D.AF8\endcsname{\href{S6_SERIES_GRIND_III.pdf\#nameddest=tag-01XP}{III.{D.AF8} (p.~111)} }
\expandafter\def\csname SRef@maps:eq:D.AF9\endcsname{\href{S6_SERIES_GRIND_III.pdf\#nameddest=tag-01XQ}{III.{D.AF9} (p.~112)} }
\expandafter\def\csname SRef@maps:eq:D.OD1\endcsname{\href{S6_SERIES_GRIND_III.pdf\#nameddest=tag-0212}{III.{D.OD1} (p.~113)} }
\expandafter\def\csname SRef@maps:eq:D.OD2\endcsname{\href{S6_SERIES_GRIND_III.pdf\#nameddest=tag-0213}{III.{D.OD2} (p.~113)} }
\expandafter\def\csname SRef@maps:eq:D.OD3\endcsname{\href{S6_SERIES_GRIND_III.pdf\#nameddest=tag-0214}{III.{D.OD3} (p.~113)} }
\expandafter\def\csname SRef@maps:eq:D.OD4\endcsname{\href{S6_SERIES_GRIND_III.pdf\#nameddest=tag-0215}{III.{D.OD4} (p.~114)} }
\expandafter\def\csname SRef@maps:eq:D.OD5\endcsname{\href{S6_SERIES_GRIND_III.pdf\#nameddest=tag-0216}{III.{D.OD5} (p.~114)} }
\expandafter\def\csname SRef@maps:eq:D.OC1\endcsname{\href{S6_SERIES_GRIND_III.pdf\#nameddest=tag-020V}{III.{D.OC1} (p.~116)} }
\expandafter\def\csname SRef@maps:eq:D.OC2\endcsname{\href{S6_SERIES_GRIND_III.pdf\#nameddest=tag-020W}{III.{D.OC2} (p.~116)} }
\expandafter\def\csname SRef@maps:eq:D.OC3\endcsname{\href{S6_SERIES_GRIND_III.pdf\#nameddest=tag-020X}{III.{D.OC3} (p.~116)} }
\expandafter\def\csname SRef@maps:eq:D.OC4\endcsname{\href{S6_SERIES_GRIND_III.pdf\#nameddest=tag-020Y}{III.{D.OC4} (p.~117)} }
\expandafter\def\csname SRef@maps:eq:D.OC5\endcsname{\href{S6_SERIES_GRIND_III.pdf\#nameddest=tag-020Z}{III.{D.OC5} (p.~117)} }
\expandafter\def\csname SRef@maps:eq:D.OC6\endcsname{\href{S6_SERIES_GRIND_III.pdf\#nameddest=tag-0210}{III.{D.OC6} (p.~117)} }
\expandafter\def\csname SRef@maps:eq:D.OC7\endcsname{\href{S6_SERIES_GRIND_III.pdf\#nameddest=tag-0211}{III.{D.OC7} (p.~118)} }
\expandafter\def\csname SRef@maps:eq:D.OM1\endcsname{\href{S6_SERIES_GRIND_III.pdf\#nameddest=tag-0222}{III.{D.OM1} (p.~119)} }
\expandafter\def\csname SRef@maps:eq:D.OM2\endcsname{\href{S6_SERIES_GRIND_III.pdf\#nameddest=tag-0224}{III.{D.OM2} (p.~119)} }
\expandafter\def\csname SRef@maps:eq:D.OM3\endcsname{\href{S6_SERIES_GRIND_III.pdf\#nameddest=tag-0225}{III.{D.OM3} (p.~120)} }
\expandafter\def\csname SRef@maps:eq:D.OM4\endcsname{\href{S6_SERIES_GRIND_III.pdf\#nameddest=tag-0226}{III.{D.OM4} (p.~120)} }
\expandafter\def\csname SRef@maps:eq:D.OM5\endcsname{\href{S6_SERIES_GRIND_III.pdf\#nameddest=tag-0227}{III.{D.OM5} (p.~120)} }
\expandafter\def\csname SRef@maps:eq:D.OM6\endcsname{\href{S6_SERIES_GRIND_III.pdf\#nameddest=tag-0228}{III.{D.OM6} (p.~120)} }
\expandafter\def\csname SRef@maps:eq:D.OM7\endcsname{\href{S6_SERIES_GRIND_III.pdf\#nameddest=tag-0229}{III.{D.OM7} (p.~120)} }
\expandafter\def\csname SRef@maps:eq:D.OM8\endcsname{\href{S6_SERIES_GRIND_III.pdf\#nameddest=tag-022A}{III.{D.OM8} (p.~120)} }
\expandafter\def\csname SRef@maps:eq:D.OM9\endcsname{\href{S6_SERIES_GRIND_III.pdf\#nameddest=tag-022B}{III.{D.OM9} (p.~121)} }
\expandafter\def\csname SRef@maps:eq:B.1.1\endcsname{\href{S6_SERIES_GRIND_III.pdf\#nameddest=tag-01UV}{III.{B.1.1} (p.~14)} }
\expandafter\def\csname SRef@maps:eq:D.OM10\endcsname{\href{S6_SERIES_GRIND_III.pdf\#nameddest=tag-0223}{III.{D.OM10} (p.~121)} }
\expandafter\def\csname SRef@maps:eq:D.WE1\endcsname{\href{S6_SERIES_GRIND_III.pdf\#nameddest=tag-026O}{III.{D.WE1} (p.~122)} }
\expandafter\def\csname SRef@maps:eq:D.WE2\endcsname{\href{S6_SERIES_GRIND_III.pdf\#nameddest=tag-026U}{III.{D.WE2} (p.~122)} }
\expandafter\def\csname SRef@maps:eq:D.WE3\endcsname{\href{S6_SERIES_GRIND_III.pdf\#nameddest=tag-026V}{III.{D.WE3} (p.~122)} }
\expandafter\def\csname SRef@maps:eq:D.WE4\endcsname{\href{S6_SERIES_GRIND_III.pdf\#nameddest=tag-026W}{III.{D.WE4} (p.~122)} }
\expandafter\def\csname SRef@maps:eq:D.WE5\endcsname{\href{S6_SERIES_GRIND_III.pdf\#nameddest=tag-026X}{III.{D.WE5} (p.~122)} }
\expandafter\def\csname SRef@maps:eq:D.WE6\endcsname{\href{S6_SERIES_GRIND_III.pdf\#nameddest=tag-026Y}{III.{D.WE6} (p.~122)} }
\expandafter\def\csname SRef@maps:eq:D.WE7\endcsname{\href{S6_SERIES_GRIND_III.pdf\#nameddest=tag-026Z}{III.{D.WE7} (p.~123)} }
\expandafter\def\csname SRef@maps:eq:D.WE8\endcsname{\href{S6_SERIES_GRIND_III.pdf\#nameddest=tag-0270}{III.{D.WE8} (p.~123)} }
\expandafter\def\csname SRef@maps:eq:D.WE9\endcsname{\href{S6_SERIES_GRIND_III.pdf\#nameddest=tag-0271}{III.{D.WE9} (p.~123)} }
\expandafter\def\csname SRef@maps:eq:B.1.2\endcsname{\href{S6_SERIES_GRIND_III.pdf\#nameddest=tag-01UW}{III.{B.1.2} (p.~14)} }
\expandafter\def\csname SRef@maps:eq:D.WE10\endcsname{\href{S6_SERIES_GRIND_III.pdf\#nameddest=tag-026P}{III.{D.WE10} (p.~123)} }
\expandafter\def\csname SRef@maps:eq:D.WE11\endcsname{\href{S6_SERIES_GRIND_III.pdf\#nameddest=tag-026Q}{III.{D.WE11} (p.~124)} }
\expandafter\def\csname SRef@maps:eq:D.WE12\endcsname{\href{S6_SERIES_GRIND_III.pdf\#nameddest=tag-026R}{III.{D.WE12} (p.~125)} }
\expandafter\def\csname SRef@maps:eq:D.WE13\endcsname{\href{S6_SERIES_GRIND_III.pdf\#nameddest=tag-026S}{III.{D.WE13} (p.~125)} }
\expandafter\def\csname SRef@maps:eq:D.WE14\endcsname{\href{S6_SERIES_GRIND_III.pdf\#nameddest=tag-026T}{III.{D.WE14} (p.~126)} }
\expandafter\def\csname SRef@maps:eq:D.C21\endcsname{\href{S6_SERIES_GRIND_III.pdf\#nameddest=tag-01Y7}{III.{D.C21} (p.~126)} }
\expandafter\def\csname SRef@maps:eq:D.C22\endcsname{\href{S6_SERIES_GRIND_III.pdf\#nameddest=tag-01YF}{III.{D.C22} (p.~126)} }
\expandafter\def\csname SRef@maps:eq:D.C23\endcsname{\href{S6_SERIES_GRIND_III.pdf\#nameddest=tag-01YG}{III.{D.C23} (p.~127)} }
\expandafter\def\csname SRef@maps:eq:D.C24\endcsname{\href{S6_SERIES_GRIND_III.pdf\#nameddest=tag-01YH}{III.{D.C24} (p.~127)} }
\expandafter\def\csname SRef@maps:eq:D.C25\endcsname{\href{S6_SERIES_GRIND_III.pdf\#nameddest=tag-01YI}{III.{D.C25} (p.~127)} }
\expandafter\def\csname SRef@maps:eq:B.1.3\endcsname{\href{S6_SERIES_GRIND_III.pdf\#nameddest=tag-01UX}{III.{B.1.3} (p.~14)} }
\expandafter\def\csname SRef@maps:eq:D.C26\endcsname{\href{S6_SERIES_GRIND_III.pdf\#nameddest=tag-01YJ}{III.{D.C26} (p.~127)} }
\expandafter\def\csname SRef@maps:eq:D.C27\endcsname{\href{S6_SERIES_GRIND_III.pdf\#nameddest=tag-01YK}{III.{D.C27} (p.~127)} }
\expandafter\def\csname SRef@maps:eq:D.C28\endcsname{\href{S6_SERIES_GRIND_III.pdf\#nameddest=tag-01YL}{III.{D.C28} (p.~127)} }
\expandafter\def\csname SRef@maps:eq:D.C29\endcsname{\href{S6_SERIES_GRIND_III.pdf\#nameddest=tag-01YM}{III.{D.C29} (p.~128)} }
\expandafter\def\csname SRef@maps:eq:D.C210\endcsname{\href{S6_SERIES_GRIND_III.pdf\#nameddest=tag-01Y8}{III.{D.C210} (p.~128)} }
\expandafter\def\csname SRef@maps:eq:D.C211\endcsname{\href{S6_SERIES_GRIND_III.pdf\#nameddest=tag-01Y9}{III.{D.C211} (p.~128)} }
\expandafter\def\csname SRef@maps:eq:D.C212\endcsname{\href{S6_SERIES_GRIND_III.pdf\#nameddest=tag-01YA}{III.{D.C212} (p.~129)} }
\expandafter\def\csname SRef@maps:eq:D.C213\endcsname{\href{S6_SERIES_GRIND_III.pdf\#nameddest=tag-01YB}{III.{D.C213} (p.~129)} }
\expandafter\def\csname SRef@maps:eq:D.C214\endcsname{\href{S6_SERIES_GRIND_III.pdf\#nameddest=tag-01YC}{III.{D.C214} (p.~129)} }
\expandafter\def\csname SRef@maps:eq:D.C215\endcsname{\href{S6_SERIES_GRIND_III.pdf\#nameddest=tag-01YD}{III.{D.C215} (p.~129)} }
\expandafter\def\csname SRef@maps:eq:B.2.1\endcsname{\href{S6_SERIES_GRIND_III.pdf\#nameddest=tag-01UY}{III.{B.2.1} (p.~15)} }
\expandafter\def\csname SRef@maps:eq:D.C216\endcsname{\href{S6_SERIES_GRIND_III.pdf\#nameddest=tag-01YE}{III.{D.C216} (p.~130)} }
\expandafter\def\csname SRef@maps:eq:D.OP1\endcsname{\href{S6_SERIES_GRIND_III.pdf\#nameddest=tag-022C}{III.{D.OP1} (p.~131)} }
\expandafter\def\csname SRef@maps:eq:D.OP2\endcsname{\href{S6_SERIES_GRIND_III.pdf\#nameddest=tag-022F}{III.{D.OP2} (p.~131)} }
\expandafter\def\csname SRef@maps:eq:D.OP3\endcsname{\href{S6_SERIES_GRIND_III.pdf\#nameddest=tag-022G}{III.{D.OP3} (p.~131)} }
\expandafter\def\csname SRef@maps:eq:D.OP4\endcsname{\href{S6_SERIES_GRIND_III.pdf\#nameddest=tag-022H}{III.{D.OP4} (p.~131)} }
\expandafter\def\csname SRef@maps:eq:D.OP5\endcsname{\href{S6_SERIES_GRIND_III.pdf\#nameddest=tag-022I}{III.{D.OP5} (p.~132)} }
\expandafter\def\csname SRef@maps:eq:D.OP6\endcsname{\href{S6_SERIES_GRIND_III.pdf\#nameddest=tag-022K}{III.{D.OP6} (p.~132)} }
\expandafter\def\csname SRef@maps:eq:D.OP5a\endcsname{\href{S6_SERIES_GRIND_III.pdf\#nameddest=tag-022J}{III.{D.OP5a} (p.~132)} }
\expandafter\def\csname SRef@maps:eq:D.OP7\endcsname{\href{S6_SERIES_GRIND_III.pdf\#nameddest=tag-022L}{III.{D.OP7} (p.~133)} }
\expandafter\def\csname SRef@maps:eq:D.OP8\endcsname{\href{S6_SERIES_GRIND_III.pdf\#nameddest=tag-022M}{III.{D.OP8} (p.~133)} }
\expandafter\def\csname SRef@maps:eq:B.3.1\endcsname{\href{S6_SERIES_GRIND_III.pdf\#nameddest=tag-01UZ}{III.{B.3.1} (p.~16)} }
\expandafter\def\csname SRef@maps:eq:D.OP9\endcsname{\href{S6_SERIES_GRIND_III.pdf\#nameddest=tag-022N}{III.{D.OP9} (p.~133)} }
\expandafter\def\csname SRef@maps:eq:D.OP10\endcsname{\href{S6_SERIES_GRIND_III.pdf\#nameddest=tag-022D}{III.{D.OP10} (p.~133)} }
\expandafter\def\csname SRef@maps:eq:D.OP11\endcsname{\href{S6_SERIES_GRIND_III.pdf\#nameddest=tag-022E}{III.{D.OP11} (p.~133)} }
\expandafter\def\csname SRef@maps:eq:D.WP1\endcsname{\href{S6_SERIES_GRIND_III.pdf\#nameddest=tag-027O}{III.{D.WP1} (p.~134)} }
\expandafter\def\csname SRef@maps:eq:D.WP2\endcsname{\href{S6_SERIES_GRIND_III.pdf\#nameddest=tag-027S}{III.{D.WP2} (p.~134)} }
\expandafter\def\csname SRef@maps:eq:D.WP3\endcsname{\href{S6_SERIES_GRIND_III.pdf\#nameddest=tag-027T}{III.{D.WP3} (p.~134)} }
\expandafter\def\csname SRef@maps:eq:D.WP4\endcsname{\href{S6_SERIES_GRIND_III.pdf\#nameddest=tag-027U}{III.{D.WP4} (p.~134)} }
\expandafter\def\csname SRef@maps:eq:D.WP5\endcsname{\href{S6_SERIES_GRIND_III.pdf\#nameddest=tag-027V}{III.{D.WP5} (p.~135)} }
\expandafter\def\csname SRef@maps:eq:D.WP6\endcsname{\href{S6_SERIES_GRIND_III.pdf\#nameddest=tag-027W}{III.{D.WP6} (p.~135)} }
\expandafter\def\csname SRef@maps:eq:D.WP7\endcsname{\href{S6_SERIES_GRIND_III.pdf\#nameddest=tag-027X}{III.{D.WP7} (p.~135)} }
\expandafter\def\csname SRef@maps:eq:B.3.2\endcsname{\href{S6_SERIES_GRIND_III.pdf\#nameddest=tag-01V0}{III.{B.3.2} (p.~16)} }
\expandafter\def\csname SRef@maps:eq:D.WP8\endcsname{\href{S6_SERIES_GRIND_III.pdf\#nameddest=tag-027Y}{III.{D.WP8} (p.~135)} }
\expandafter\def\csname SRef@maps:eq:D.WP9\endcsname{\href{S6_SERIES_GRIND_III.pdf\#nameddest=tag-027Z}{III.{D.WP9} (p.~135)} }
\expandafter\def\csname SRef@maps:eq:D.WP10\endcsname{\href{S6_SERIES_GRIND_III.pdf\#nameddest=tag-027P}{III.{D.WP10} (p.~136)} }
\expandafter\def\csname SRef@maps:eq:D.WP11\endcsname{\href{S6_SERIES_GRIND_III.pdf\#nameddest=tag-027Q}{III.{D.WP11} (p.~137)} }
\expandafter\def\csname SRef@maps:eq:D.WP12\endcsname{\href{S6_SERIES_GRIND_III.pdf\#nameddest=tag-027R}{III.{D.WP12} (p.~138)} }
\expandafter\def\csname SRef@maps:eq:D.MP1\endcsname{\href{S6_SERIES_GRIND_III.pdf\#nameddest=tag-0209}{III.{D.MP1} (p.~138)} }
\expandafter\def\csname SRef@maps:eq:D.MP2\endcsname{\href{S6_SERIES_GRIND_III.pdf\#nameddest=tag-020B}{III.{D.MP2} (p.~139)} }
\expandafter\def\csname SRef@maps:eq:D.MP3\endcsname{\href{S6_SERIES_GRIND_III.pdf\#nameddest=tag-020C}{III.{D.MP3} (p.~139)} }
\expandafter\def\csname SRef@maps:eq:D.MP4\endcsname{\href{S6_SERIES_GRIND_III.pdf\#nameddest=tag-020D}{III.{D.MP4} (p.~139)} }
\expandafter\def\csname SRef@maps:eq:D.MP5\endcsname{\href{S6_SERIES_GRIND_III.pdf\#nameddest=tag-020E}{III.{D.MP5} (p.~140)} }
\expandafter\def\csname SRef@maps:eq:B.3.3\endcsname{\href{S6_SERIES_GRIND_III.pdf\#nameddest=tag-01V1}{III.{B.3.3} (p.~16)} }
\expandafter\def\csname SRef@maps:eq:D.MP6\endcsname{\href{S6_SERIES_GRIND_III.pdf\#nameddest=tag-020F}{III.{D.MP6} (p.~140)} }
\expandafter\def\csname SRef@maps:eq:D.MP7\endcsname{\href{S6_SERIES_GRIND_III.pdf\#nameddest=tag-020G}{III.{D.MP7} (p.~140)} }
\expandafter\def\csname SRef@maps:eq:D.MP8\endcsname{\href{S6_SERIES_GRIND_III.pdf\#nameddest=tag-020H}{III.{D.MP8} (p.~140)} }
\expandafter\def\csname SRef@maps:eq:D.MP9\endcsname{\href{S6_SERIES_GRIND_III.pdf\#nameddest=tag-020I}{III.{D.MP9} (p.~141)} }
\expandafter\def\csname SRef@maps:eq:D.MP10\endcsname{\href{S6_SERIES_GRIND_III.pdf\#nameddest=tag-020A}{III.{D.MP10} (p.~141)} }
\expandafter\def\csname SRef@maps:eq:D.T31\endcsname{\href{S6_SERIES_GRIND_III.pdf\#nameddest=tag-025C}{III.{D.T31} (p.~142)} }
\expandafter\def\csname SRef@maps:eq:D.T32\endcsname{\href{S6_SERIES_GRIND_III.pdf\#nameddest=tag-025D}{III.{D.T32} (p.~142)} }
\expandafter\def\csname SRef@maps:eq:D.T33\endcsname{\href{S6_SERIES_GRIND_III.pdf\#nameddest=tag-025E}{III.{D.T33} (p.~142)} }
\expandafter\def\csname SRef@maps:eq:D.T34\endcsname{\href{S6_SERIES_GRIND_III.pdf\#nameddest=tag-025F}{III.{D.T34} (p.~143)} }
\expandafter\def\csname SRef@maps:eq:D.T35\endcsname{\href{S6_SERIES_GRIND_III.pdf\#nameddest=tag-025G}{III.{D.T35} (p.~143)} }
\expandafter\def\csname SRef@maps:eq:B.3.4\endcsname{\href{S6_SERIES_GRIND_III.pdf\#nameddest=tag-01V2}{III.{B.3.4} (p.~16)} }
\expandafter\def\csname SRef@maps:eq:D.T36\endcsname{\href{S6_SERIES_GRIND_III.pdf\#nameddest=tag-025H}{III.{D.T36} (p.~143)} }
\expandafter\def\csname SRef@maps:eq:D.T37\endcsname{\href{S6_SERIES_GRIND_III.pdf\#nameddest=tag-025I}{III.{D.T37} (p.~144)} }
\expandafter\def\csname SRef@maps:eq:D.T38\endcsname{\href{S6_SERIES_GRIND_III.pdf\#nameddest=tag-025J}{III.{D.T38} (p.~144)} }
\expandafter\def\csname SRef@maps:eq:D.T39\endcsname{\href{S6_SERIES_GRIND_III.pdf\#nameddest=tag-025K}{III.{D.T39} (p.~144)} }
\expandafter\def\csname SRef@maps:eq:D.OW1\endcsname{\href{S6_SERIES_GRIND_III.pdf\#nameddest=tag-022W}{III.{D.OW1} (p.~145)} }
\expandafter\def\csname SRef@maps:eq:D.OW2\endcsname{\href{S6_SERIES_GRIND_III.pdf\#nameddest=tag-022X}{III.{D.OW2} (p.~145)} }
\expandafter\def\csname SRef@maps:eq:D.OW3\endcsname{\href{S6_SERIES_GRIND_III.pdf\#nameddest=tag-022Y}{III.{D.OW3} (p.~145)} }
\expandafter\def\csname SRef@maps:eq:D.OW4\endcsname{\href{S6_SERIES_GRIND_III.pdf\#nameddest=tag-022Z}{III.{D.OW4} (p.~145)} }
\expandafter\def\csname SRef@maps:eq:D.OW5\endcsname{\href{S6_SERIES_GRIND_III.pdf\#nameddest=tag-0230}{III.{D.OW5} (p.~145)} }
\expandafter\def\csname SRef@maps:eq:D.OW6\endcsname{\href{S6_SERIES_GRIND_III.pdf\#nameddest=tag-0231}{III.{D.OW6} (p.~146)} }
\expandafter\def\csname SRef@maps:eq:B.3.5\endcsname{\href{S6_SERIES_GRIND_III.pdf\#nameddest=tag-01V3}{III.{B.3.5} (p.~18)} }
\expandafter\def\csname SRef@maps:eq:D.OW7\endcsname{\href{S6_SERIES_GRIND_III.pdf\#nameddest=tag-0232}{III.{D.OW7} (p.~147)} }
\expandafter\def\csname SRef@maps:eq:D.BC1\endcsname{\href{S6_SERIES_GRIND_III.pdf\#nameddest=tag-01XR}{III.{D.BC1} (p.~148)} }
\expandafter\def\csname SRef@maps:eq:D.BC2\endcsname{\href{S6_SERIES_GRIND_III.pdf\#nameddest=tag-01XY}{III.{D.BC2} (p.~148)} }
\expandafter\def\csname SRef@maps:eq:D.BC3\endcsname{\href{S6_SERIES_GRIND_III.pdf\#nameddest=tag-01XZ}{III.{D.BC3} (p.~148)} }
\expandafter\def\csname SRef@maps:eq:D.BC4\endcsname{\href{S6_SERIES_GRIND_III.pdf\#nameddest=tag-01Y0}{III.{D.BC4} (p.~148)} }
\expandafter\def\csname SRef@maps:eq:D.BC5\endcsname{\href{S6_SERIES_GRIND_III.pdf\#nameddest=tag-01Y1}{III.{D.BC5} (p.~148)} }
\expandafter\def\csname SRef@maps:eq:D.BC6\endcsname{\href{S6_SERIES_GRIND_III.pdf\#nameddest=tag-01Y2}{III.{D.BC6} (p.~149)} }
\expandafter\def\csname SRef@maps:eq:D.BC7\endcsname{\href{S6_SERIES_GRIND_III.pdf\#nameddest=tag-01Y3}{III.{D.BC7} (p.~149)} }
\expandafter\def\csname SRef@maps:eq:D.BC8\endcsname{\href{S6_SERIES_GRIND_III.pdf\#nameddest=tag-01Y4}{III.{D.BC8} (p.~149)} }
\expandafter\def\csname SRef@maps:eq:D.BC9\endcsname{\href{S6_SERIES_GRIND_III.pdf\#nameddest=tag-01Y5}{III.{D.BC9} (p.~149)} }
\expandafter\def\csname SRef@maps:eq:B.4.1\endcsname{\href{S6_SERIES_GRIND_III.pdf\#nameddest=tag-01V4}{III.{B.4.1} (p.~19)} }
\expandafter\def\csname SRef@maps:eq:D.BC10\endcsname{\href{S6_SERIES_GRIND_III.pdf\#nameddest=tag-01XS}{III.{D.BC10} (p.~149)} }
\expandafter\def\csname SRef@maps:eq:D.BC11\endcsname{\href{S6_SERIES_GRIND_III.pdf\#nameddest=tag-01XT}{III.{D.BC11} (p.~149)} }
\expandafter\def\csname SRef@maps:eq:D.BC12\endcsname{\href{S6_SERIES_GRIND_III.pdf\#nameddest=tag-01XU}{III.{D.BC12} (p.~149)} }
\expandafter\def\csname SRef@maps:eq:D.BC13\endcsname{\href{S6_SERIES_GRIND_III.pdf\#nameddest=tag-01XV}{III.{D.BC13} (p.~150)} }
\expandafter\def\csname SRef@maps:eq:D.BC14\endcsname{\href{S6_SERIES_GRIND_III.pdf\#nameddest=tag-01XW}{III.{D.BC14} (p.~150)} }
\expandafter\def\csname SRef@maps:eq:D.BC15\endcsname{\href{S6_SERIES_GRIND_III.pdf\#nameddest=tag-01XX}{III.{D.BC15} (p.~150)} }
\expandafter\def\csname SRef@maps:eq:D.MT1\endcsname{\href{S6_SERIES_GRIND_III.pdf\#nameddest=tag-020J}{III.{D.MT1} (p.~151)} }
\expandafter\def\csname SRef@maps:eq:D.MT2\endcsname{\href{S6_SERIES_GRIND_III.pdf\#nameddest=tag-020K}{III.{D.MT2} (p.~152)} }
\expandafter\def\csname SRef@maps:eq:D.MT3\endcsname{\href{S6_SERIES_GRIND_III.pdf\#nameddest=tag-020L}{III.{D.MT3} (p.~153)} }
\expandafter\def\csname SRef@maps:eq:D.GC1\endcsname{\href{S6_SERIES_GRIND_III.pdf\#nameddest=tag-01YY}{III.{D.GC1} (p.~153)} }
\expandafter\def\csname SRef@maps:eq:B.4.2\endcsname{\href{S6_SERIES_GRIND_III.pdf\#nameddest=tag-01V5}{III.{B.4.2} (p.~19)} }
\expandafter\def\csname SRef@maps:eq:D.GC2\endcsname{\href{S6_SERIES_GRIND_III.pdf\#nameddest=tag-01Z4}{III.{D.GC2} (p.~153)} }
\expandafter\def\csname SRef@maps:eq:D.GC3\endcsname{\href{S6_SERIES_GRIND_III.pdf\#nameddest=tag-01Z5}{III.{D.GC3} (p.~154)} }
\expandafter\def\csname SRef@maps:eq:D.GC4\endcsname{\href{S6_SERIES_GRIND_III.pdf\#nameddest=tag-01Z6}{III.{D.GC4} (p.~154)} }
\expandafter\def\csname SRef@maps:eq:D.GC5\endcsname{\href{S6_SERIES_GRIND_III.pdf\#nameddest=tag-01Z7}{III.{D.GC5} (p.~154)} }
\expandafter\def\csname SRef@maps:eq:D.GC6\endcsname{\href{S6_SERIES_GRIND_III.pdf\#nameddest=tag-01Z8}{III.{D.GC6} (p.~154)} }
\expandafter\def\csname SRef@maps:eq:D.GC7\endcsname{\href{S6_SERIES_GRIND_III.pdf\#nameddest=tag-01Z9}{III.{D.GC7} (p.~154)} }
\expandafter\def\csname SRef@maps:eq:D.GC8\endcsname{\href{S6_SERIES_GRIND_III.pdf\#nameddest=tag-01ZA}{III.{D.GC8} (p.~155)} }
\expandafter\def\csname SRef@maps:eq:D.GC9\endcsname{\href{S6_SERIES_GRIND_III.pdf\#nameddest=tag-01ZB}{III.{D.GC9} (p.~155)} }
\expandafter\def\csname SRef@maps:eq:D.GC10\endcsname{\href{S6_SERIES_GRIND_III.pdf\#nameddest=tag-01YZ}{III.{D.GC10} (p.~155)} }
\expandafter\def\csname SRef@maps:eq:D.GC11\endcsname{\href{S6_SERIES_GRIND_III.pdf\#nameddest=tag-01Z0}{III.{D.GC11} (p.~155)} }
\expandafter\def\csname SRef@maps:eq:B.4.3\endcsname{\href{S6_SERIES_GRIND_III.pdf\#nameddest=tag-01V6}{III.{B.4.3} (p.~20)} }
\expandafter\def\csname SRef@maps:eq:D.GC12\endcsname{\href{S6_SERIES_GRIND_III.pdf\#nameddest=tag-01Z1}{III.{D.GC12} (p.~155)} }
\expandafter\def\csname SRef@maps:eq:D.GC13\endcsname{\href{S6_SERIES_GRIND_III.pdf\#nameddest=tag-01Z2}{III.{D.GC13} (p.~155)} }
\expandafter\def\csname SRef@maps:eq:D.GC14\endcsname{\href{S6_SERIES_GRIND_III.pdf\#nameddest=tag-01Z3}{III.{D.GC14} (p.~156)} }
\expandafter\def\csname SRef@maps:eq:D44.1\endcsname{\href{S6_SERIES_GRIND_III.pdf\#nameddest=tag-0282}{III.{D44.1} (p.~157)} }
\expandafter\def\csname SRef@maps:eq:D44.0a\endcsname{\href{S6_SERIES_GRIND_III.pdf\#nameddest=tag-0280}{III.{D44.0a} (p.~157)} }
\expandafter\def\csname SRef@maps:eq:D44.0b\endcsname{\href{S6_SERIES_GRIND_III.pdf\#nameddest=tag-0281}{III.{D44.0b} (p.~157)} }
\expandafter\def\csname SRef@maps:eq:D44.2\endcsname{\href{S6_SERIES_GRIND_III.pdf\#nameddest=tag-0288}{III.{D44.2} (p.~158)} }
\expandafter\def\csname SRef@maps:eq:D44.3\endcsname{\href{S6_SERIES_GRIND_III.pdf\#nameddest=tag-0289}{III.{D44.3} (p.~158)} }
\expandafter\def\csname SRef@maps:eq:D44.4\endcsname{\href{S6_SERIES_GRIND_III.pdf\#nameddest=tag-028A}{III.{D44.4} (p.~159)} }
\expandafter\def\csname SRef@maps:eq:D44.5\endcsname{\href{S6_SERIES_GRIND_III.pdf\#nameddest=tag-028B}{III.{D44.5} (p.~159)} }
\expandafter\def\csname SRef@maps:eq:B.5.1\endcsname{\href{S6_SERIES_GRIND_III.pdf\#nameddest=tag-01V7}{III.{B.5.1} (p.~20)} }
\expandafter\def\csname SRef@maps:eq:D44.6\endcsname{\href{S6_SERIES_GRIND_III.pdf\#nameddest=tag-028C}{III.{D44.6} (p.~159)} }
\expandafter\def\csname SRef@maps:eq:D44.7\endcsname{\href{S6_SERIES_GRIND_III.pdf\#nameddest=tag-028D}{III.{D44.7} (p.~159)} }
\expandafter\def\csname SRef@maps:eq:D44.8\endcsname{\href{S6_SERIES_GRIND_III.pdf\#nameddest=tag-028E}{III.{D44.8} (p.~160)} }
\expandafter\def\csname SRef@maps:eq:D44.9\endcsname{\href{S6_SERIES_GRIND_III.pdf\#nameddest=tag-028F}{III.{D44.9} (p.~160)} }
\expandafter\def\csname SRef@maps:eq:D44.10\endcsname{\href{S6_SERIES_GRIND_III.pdf\#nameddest=tag-0283}{III.{D44.10} (p.~160)} }
\expandafter\def\csname SRef@maps:eq:D44.11\endcsname{\href{S6_SERIES_GRIND_III.pdf\#nameddest=tag-0284}{III.{D44.11} (p.~160)} }
\expandafter\def\csname SRef@maps:eq:D44.12\endcsname{\href{S6_SERIES_GRIND_III.pdf\#nameddest=tag-0285}{III.{D44.12} (p.~160)} }
\expandafter\def\csname SRef@maps:eq:D44.13\endcsname{\href{S6_SERIES_GRIND_III.pdf\#nameddest=tag-0286}{III.{D44.13} (p.~161)} }
\expandafter\def\csname SRef@maps:eq:D44.14\endcsname{\href{S6_SERIES_GRIND_III.pdf\#nameddest=tag-0287}{III.{D44.14} (p.~161)} }
\expandafter\def\csname SRef@maps:eq:D45.1\endcsname{\href{S6_SERIES_GRIND_III.pdf\#nameddest=tag-028G}{III.{D45.1} (p.~162)} }
\expandafter\def\csname SRef@maps:eq:B.5.2\endcsname{\href{S6_SERIES_GRIND_III.pdf\#nameddest=tag-01VH}{III.{B.5.2} (p.~21)} }
\expandafter\def\csname SRef@maps:eq:D45.2\endcsname{\href{S6_SERIES_GRIND_III.pdf\#nameddest=tag-028R}{III.{D45.2} (p.~162)} }
\expandafter\def\csname SRef@maps:eq:D45.3\endcsname{\href{S6_SERIES_GRIND_III.pdf\#nameddest=tag-028V}{III.{D45.3} (p.~162)} }
\expandafter\def\csname SRef@maps:eq:D45.4\endcsname{\href{S6_SERIES_GRIND_III.pdf\#nameddest=tag-028W}{III.{D45.4} (p.~162)} }
\expandafter\def\csname SRef@maps:eq:D45.5\endcsname{\href{S6_SERIES_GRIND_III.pdf\#nameddest=tag-028X}{III.{D45.5} (p.~162)} }
\expandafter\def\csname SRef@maps:eq:D45.6\endcsname{\href{S6_SERIES_GRIND_III.pdf\#nameddest=tag-028Y}{III.{D45.6} (p.~163)} }
\expandafter\def\csname SRef@maps:eq:D45.7\endcsname{\href{S6_SERIES_GRIND_III.pdf\#nameddest=tag-028Z}{III.{D45.7} (p.~163)} }
\expandafter\def\csname SRef@maps:eq:D45.8\endcsname{\href{S6_SERIES_GRIND_III.pdf\#nameddest=tag-0290}{III.{D45.8} (p.~163)} }
\expandafter\def\csname SRef@maps:eq:D45.9\endcsname{\href{S6_SERIES_GRIND_III.pdf\#nameddest=tag-0291}{III.{D45.9} (p.~163)} }
\expandafter\def\csname SRef@maps:eq:D45.10\endcsname{\href{S6_SERIES_GRIND_III.pdf\#nameddest=tag-028H}{III.{D45.10} (p.~164)} }
\expandafter\def\csname SRef@maps:eq:D45.11\endcsname{\href{S6_SERIES_GRIND_III.pdf\#nameddest=tag-028I}{III.{D45.11} (p.~164)} }
\expandafter\def\csname SRef@maps:eq:B.5.3\endcsname{\href{S6_SERIES_GRIND_III.pdf\#nameddest=tag-01VI}{III.{B.5.3} (p.~21)} }
\expandafter\def\csname SRef@maps:eq:D45.12\endcsname{\href{S6_SERIES_GRIND_III.pdf\#nameddest=tag-028J}{III.{D45.12} (p.~164)} }
\expandafter\def\csname SRef@maps:eq:D45.13\endcsname{\href{S6_SERIES_GRIND_III.pdf\#nameddest=tag-028K}{III.{D45.13} (p.~164)} }
\expandafter\def\csname SRef@maps:eq:D45.14\endcsname{\href{S6_SERIES_GRIND_III.pdf\#nameddest=tag-028L}{III.{D45.14} (p.~164)} }
\expandafter\def\csname SRef@maps:eq:D45.15\endcsname{\href{S6_SERIES_GRIND_III.pdf\#nameddest=tag-028M}{III.{D45.15} (p.~164)} }
\expandafter\def\csname SRef@maps:eq:D45.16\endcsname{\href{S6_SERIES_GRIND_III.pdf\#nameddest=tag-028N}{III.{D45.16} (p.~165)} }
\expandafter\def\csname SRef@maps:eq:D45.17\endcsname{\href{S6_SERIES_GRIND_III.pdf\#nameddest=tag-028O}{III.{D45.17} (p.~165)} }
\expandafter\def\csname SRef@maps:eq:D45.18\endcsname{\href{S6_SERIES_GRIND_III.pdf\#nameddest=tag-028P}{III.{D45.18} (p.~165)} }
\expandafter\def\csname SRef@maps:eq:D45.19\endcsname{\href{S6_SERIES_GRIND_III.pdf\#nameddest=tag-028Q}{III.{D45.19} (p.~165)} }
\expandafter\def\csname SRef@maps:eq:D45.20\endcsname{\href{S6_SERIES_GRIND_III.pdf\#nameddest=tag-028S}{III.{D45.20} (p.~165)} }
\expandafter\def\csname SRef@maps:eq:D45.21\endcsname{\href{S6_SERIES_GRIND_III.pdf\#nameddest=tag-028T}{III.{D45.21} (p.~166)} }
\expandafter\def\csname SRef@maps:eq:B.5.4\endcsname{\href{S6_SERIES_GRIND_III.pdf\#nameddest=tag-01VJ}{III.{B.5.4} (p.~21)} }
\expandafter\def\csname SRef@maps:eq:D45.22\endcsname{\href{S6_SERIES_GRIND_III.pdf\#nameddest=tag-028U}{III.{D45.22} (p.~166)} }
\expandafter\def\csname SRef@maps:eq:D.GI1\endcsname{\href{S6_SERIES_GRIND_III.pdf\#nameddest=tag-01ZC}{III.{D.GI1} (p.~167)} }
\expandafter\def\csname SRef@maps:eq:D.GI2\endcsname{\href{S6_SERIES_GRIND_III.pdf\#nameddest=tag-01ZF}{III.{D.GI2} (p.~167)} }
\expandafter\def\csname SRef@maps:eq:D.GI3\endcsname{\href{S6_SERIES_GRIND_III.pdf\#nameddest=tag-01ZG}{III.{D.GI3} (p.~167)} }
\expandafter\def\csname SRef@maps:eq:D.GI4\endcsname{\href{S6_SERIES_GRIND_III.pdf\#nameddest=tag-01ZH}{III.{D.GI4} (p.~167)} }
\expandafter\def\csname SRef@maps:eq:D.GI5\endcsname{\href{S6_SERIES_GRIND_III.pdf\#nameddest=tag-01ZI}{III.{D.GI5} (p.~168)} }
\expandafter\def\csname SRef@maps:eq:D.GI6\endcsname{\href{S6_SERIES_GRIND_III.pdf\#nameddest=tag-01ZJ}{III.{D.GI6} (p.~168)} }
\expandafter\def\csname SRef@maps:eq:D.GI7\endcsname{\href{S6_SERIES_GRIND_III.pdf\#nameddest=tag-01ZK}{III.{D.GI7} (p.~168)} }
\expandafter\def\csname SRef@maps:eq:D.GI8\endcsname{\href{S6_SERIES_GRIND_III.pdf\#nameddest=tag-01ZL}{III.{D.GI8} (p.~168)} }
\expandafter\def\csname SRef@maps:eq:D.GI9\endcsname{\href{S6_SERIES_GRIND_III.pdf\#nameddest=tag-01ZM}{III.{D.GI9} (p.~170)} }
\expandafter\def\csname SRef@maps:eq:B.5.5\endcsname{\href{S6_SERIES_GRIND_III.pdf\#nameddest=tag-01VK}{III.{B.5.5} (p.~22)} }
\expandafter\def\csname SRef@maps:eq:D.GI10\endcsname{\href{S6_SERIES_GRIND_III.pdf\#nameddest=tag-01ZD}{III.{D.GI10} (p.~171)} }
\expandafter\def\csname SRef@maps:eq:D.GI11\endcsname{\href{S6_SERIES_GRIND_III.pdf\#nameddest=tag-01ZE}{III.{D.GI11} (p.~171)} }
\expandafter\def\csname SRef@maps:eq:D47.1\endcsname{\href{S6_SERIES_GRIND_III.pdf\#nameddest=tag-0292}{III.{D47.1} (p.~172)} }
\expandafter\def\csname SRef@maps:eq:D47.2\endcsname{\href{S6_SERIES_GRIND_III.pdf\#nameddest=tag-0297}{III.{D47.2} (p.~172)} }
\expandafter\def\csname SRef@maps:eq:D47.3\endcsname{\href{S6_SERIES_GRIND_III.pdf\#nameddest=tag-0298}{III.{D47.3} (p.~172)} }
\expandafter\def\csname SRef@maps:eq:D47.4\endcsname{\href{S6_SERIES_GRIND_III.pdf\#nameddest=tag-0299}{III.{D47.4} (p.~172)} }
\expandafter\def\csname SRef@maps:eq:D47.5\endcsname{\href{S6_SERIES_GRIND_III.pdf\#nameddest=tag-029A}{III.{D47.5} (p.~173)} }
\expandafter\def\csname SRef@maps:eq:D47.6\endcsname{\href{S6_SERIES_GRIND_III.pdf\#nameddest=tag-029B}{III.{D47.6} (p.~173)} }
\expandafter\def\csname SRef@maps:eq:D47.7\endcsname{\href{S6_SERIES_GRIND_III.pdf\#nameddest=tag-029C}{III.{D47.7} (p.~173)} }
\expandafter\def\csname SRef@maps:eq:D47.8\endcsname{\href{S6_SERIES_GRIND_III.pdf\#nameddest=tag-029D}{III.{D47.8} (p.~173)} }
\expandafter\def\csname SRef@maps:eq:B.5.6\endcsname{\href{S6_SERIES_GRIND_III.pdf\#nameddest=tag-01VL}{III.{B.5.6} (p.~22)} }
\expandafter\def\csname SRef@maps:eq:D47.9\endcsname{\href{S6_SERIES_GRIND_III.pdf\#nameddest=tag-029E}{III.{D47.9} (p.~173)} }
\expandafter\def\csname SRef@maps:eq:D47.10\endcsname{\href{S6_SERIES_GRIND_III.pdf\#nameddest=tag-0293}{III.{D47.10} (p.~174)} }
\expandafter\def\csname SRef@maps:eq:D47.11\endcsname{\href{S6_SERIES_GRIND_III.pdf\#nameddest=tag-0294}{III.{D47.11} (p.~174)} }
\expandafter\def\csname SRef@maps:eq:D47.12\endcsname{\href{S6_SERIES_GRIND_III.pdf\#nameddest=tag-0295}{III.{D47.12} (p.~174)} }
\expandafter\def\csname SRef@maps:eq:D47.13\endcsname{\href{S6_SERIES_GRIND_III.pdf\#nameddest=tag-0296}{III.{D47.13} (p.~174)} }
\expandafter\def\csname SRef@maps:eq:D.WI1\endcsname{\href{S6_SERIES_GRIND_III.pdf\#nameddest=tag-0272}{III.{D.WI1} (p.~175)} }
\expandafter\def\csname SRef@maps:eq:D.WI2\endcsname{\href{S6_SERIES_GRIND_III.pdf\#nameddest=tag-027D}{III.{D.WI2} (p.~175)} }
\expandafter\def\csname SRef@maps:eq:D.WI3\endcsname{\href{S6_SERIES_GRIND_III.pdf\#nameddest=tag-027H}{III.{D.WI3} (p.~175)} }
\expandafter\def\csname SRef@maps:eq:D.WI14\endcsname{\href{S6_SERIES_GRIND_III.pdf\#nameddest=tag-0277}{III.{D.WI14} (p.~175)} }
\expandafter\def\csname SRef@maps:eq:D.WI15\endcsname{\href{S6_SERIES_GRIND_III.pdf\#nameddest=tag-0278}{III.{D.WI15} (p.~175)} }
\expandafter\def\csname SRef@maps:eq:B.5.7\endcsname{\href{S6_SERIES_GRIND_III.pdf\#nameddest=tag-01VM}{III.{B.5.7} (p.~22)} }
\expandafter\def\csname SRef@maps:eq:D.WI16\endcsname{\href{S6_SERIES_GRIND_III.pdf\#nameddest=tag-0279}{III.{D.WI16} (p.~175)} }
\expandafter\def\csname SRef@maps:eq:D.WI4\endcsname{\href{S6_SERIES_GRIND_III.pdf\#nameddest=tag-027I}{III.{D.WI4} (p.~176)} }
\expandafter\def\csname SRef@maps:eq:D.WI5\endcsname{\href{S6_SERIES_GRIND_III.pdf\#nameddest=tag-027J}{III.{D.WI5} (p.~176)} }
\expandafter\def\csname SRef@maps:eq:D.WI6\endcsname{\href{S6_SERIES_GRIND_III.pdf\#nameddest=tag-027K}{III.{D.WI6} (p.~177)} }
\expandafter\def\csname SRef@maps:eq:D.WI7\endcsname{\href{S6_SERIES_GRIND_III.pdf\#nameddest=tag-027L}{III.{D.WI7} (p.~177)} }
\expandafter\def\csname SRef@maps:eq:D.WI8\endcsname{\href{S6_SERIES_GRIND_III.pdf\#nameddest=tag-027M}{III.{D.WI8} (p.~177)} }
\expandafter\def\csname SRef@maps:eq:D.WI9\endcsname{\href{S6_SERIES_GRIND_III.pdf\#nameddest=tag-027N}{III.{D.WI9} (p.~177)} }
\expandafter\def\csname SRef@maps:eq:D.WI10\endcsname{\href{S6_SERIES_GRIND_III.pdf\#nameddest=tag-0273}{III.{D.WI10} (p.~177)} }
\expandafter\def\csname SRef@maps:eq:D.WI11\endcsname{\href{S6_SERIES_GRIND_III.pdf\#nameddest=tag-0274}{III.{D.WI11} (p.~177)} }
\expandafter\def\csname SRef@maps:eq:D.WI12\endcsname{\href{S6_SERIES_GRIND_III.pdf\#nameddest=tag-0275}{III.{D.WI12} (p.~177)} }
\expandafter\def\csname SRef@maps:eq:B.5.8\endcsname{\href{S6_SERIES_GRIND_III.pdf\#nameddest=tag-01VN}{III.{B.5.8} (p.~23)} }
\expandafter\def\csname SRef@maps:eq:D.WI13\endcsname{\href{S6_SERIES_GRIND_III.pdf\#nameddest=tag-0276}{III.{D.WI13} (p.~178)} }
\expandafter\def\csname SRef@maps:eq:D.WI17\endcsname{\href{S6_SERIES_GRIND_III.pdf\#nameddest=tag-027A}{III.{D.WI17} (p.~178)} }
\expandafter\def\csname SRef@maps:eq:D.WI18\endcsname{\href{S6_SERIES_GRIND_III.pdf\#nameddest=tag-027B}{III.{D.WI18} (p.~178)} }
\expandafter\def\csname SRef@maps:eq:D.WI19\endcsname{\href{S6_SERIES_GRIND_III.pdf\#nameddest=tag-027C}{III.{D.WI19} (p.~178)} }
\expandafter\def\csname SRef@maps:eq:D.WI20\endcsname{\href{S6_SERIES_GRIND_III.pdf\#nameddest=tag-027E}{III.{D.WI20} (p.~179)} }
\expandafter\def\csname SRef@maps:eq:D.WI21\endcsname{\href{S6_SERIES_GRIND_III.pdf\#nameddest=tag-027F}{III.{D.WI21} (p.~179)} }
\expandafter\def\csname SRef@maps:eq:D.TP1\endcsname{\href{S6_SERIES_GRIND_III.pdf\#nameddest=tag-025R}{III.{D.TP1} (p.~179)} }
\expandafter\def\csname SRef@maps:eq:D.TP2\endcsname{\href{S6_SERIES_GRIND_III.pdf\#nameddest=tag-025S}{III.{D.TP2} (p.~180)} }
\expandafter\def\csname SRef@maps:eq:D.TP3\endcsname{\href{S6_SERIES_GRIND_III.pdf\#nameddest=tag-025T}{III.{D.TP3} (p.~180)} }
\expandafter\def\csname SRef@maps:eq:B.5.9\endcsname{\href{S6_SERIES_GRIND_III.pdf\#nameddest=tag-01VO}{III.{B.5.9} (p.~24)} }
\expandafter\def\csname SRef@maps:eq:D.TP4\endcsname{\href{S6_SERIES_GRIND_III.pdf\#nameddest=tag-025U}{III.{D.TP4} (p.~180)} }
\expandafter\def\csname SRef@maps:eq:D.TP5\endcsname{\href{S6_SERIES_GRIND_III.pdf\#nameddest=tag-025V}{III.{D.TP5} (p.~180)} }
\expandafter\def\csname SRef@maps:eq:D.TP6\endcsname{\href{S6_SERIES_GRIND_III.pdf\#nameddest=tag-025W}{III.{D.TP6} (p.~180)} }
\expandafter\def\csname SRef@maps:eq:D.TP7\endcsname{\href{S6_SERIES_GRIND_III.pdf\#nameddest=tag-025X}{III.{D.TP7} (p.~180)} }
\expandafter\def\csname SRef@maps:eq:D.TP8\endcsname{\href{S6_SERIES_GRIND_III.pdf\#nameddest=tag-025Y}{III.{D.TP8} (p.~180)} }
\expandafter\def\csname SRef@maps:eq:D.WI22\endcsname{\href{S6_SERIES_GRIND_III.pdf\#nameddest=tag-027G}{III.{D.WI22} (p.~181)} }
\expandafter\def\csname SRef@maps:eq:D.SL1\endcsname{\href{S6_SERIES_GRIND_III.pdf\#nameddest=tag-0247}{III.{D.SL1} (p.~182)} }
\expandafter\def\csname SRef@maps:eq:D.SL2\endcsname{\href{S6_SERIES_GRIND_III.pdf\#nameddest=tag-024I}{III.{D.SL2} (p.~182)} }
\expandafter\def\csname SRef@maps:eq:D.SL3\endcsname{\href{S6_SERIES_GRIND_III.pdf\#nameddest=tag-024O}{III.{D.SL3} (p.~182)} }
\expandafter\def\csname SRef@maps:eq:D.SL4\endcsname{\href{S6_SERIES_GRIND_III.pdf\#nameddest=tag-024P}{III.{D.SL4} (p.~182)} }
\expandafter\def\csname SRef@maps:eq:B.5.10\endcsname{\href{S6_SERIES_GRIND_III.pdf\#nameddest=tag-01V8}{III.{B.5.10} (p.~24)} }
\expandafter\def\csname SRef@maps:eq:D.SL5\endcsname{\href{S6_SERIES_GRIND_III.pdf\#nameddest=tag-024Q}{III.{D.SL5} (p.~182)} }
\expandafter\def\csname SRef@maps:eq:D.SL6\endcsname{\href{S6_SERIES_GRIND_III.pdf\#nameddest=tag-024R}{III.{D.SL6} (p.~183)} }
\expandafter\def\csname SRef@maps:eq:D.SL7\endcsname{\href{S6_SERIES_GRIND_III.pdf\#nameddest=tag-024S}{III.{D.SL7} (p.~183)} }
\expandafter\def\csname SRef@maps:eq:D.SL8\endcsname{\href{S6_SERIES_GRIND_III.pdf\#nameddest=tag-024T}{III.{D.SL8} (p.~183)} }
\expandafter\def\csname SRef@maps:eq:D.SL9\endcsname{\href{S6_SERIES_GRIND_III.pdf\#nameddest=tag-024U}{III.{D.SL9} (p.~183)} }
\expandafter\def\csname SRef@maps:eq:D.SL10\endcsname{\href{S6_SERIES_GRIND_III.pdf\#nameddest=tag-0248}{III.{D.SL10} (p.~183)} }
\expandafter\def\csname SRef@maps:eq:D.SL11\endcsname{\href{S6_SERIES_GRIND_III.pdf\#nameddest=tag-0249}{III.{D.SL11} (p.~183)} }
\expandafter\def\csname SRef@maps:eq:D.SL12\endcsname{\href{S6_SERIES_GRIND_III.pdf\#nameddest=tag-024A}{III.{D.SL12} (p.~184)} }
\expandafter\def\csname SRef@maps:eq:D.SL13\endcsname{\href{S6_SERIES_GRIND_III.pdf\#nameddest=tag-024B}{III.{D.SL13} (p.~184)} }
\expandafter\def\csname SRef@maps:eq:D.SL14\endcsname{\href{S6_SERIES_GRIND_III.pdf\#nameddest=tag-024C}{III.{D.SL14} (p.~184)} }
\expandafter\def\csname SRef@maps:eq:B.5.11\endcsname{\href{S6_SERIES_GRIND_III.pdf\#nameddest=tag-01V9}{III.{B.5.11} (p.~24)} }
\expandafter\def\csname SRef@maps:eq:D.SL15\endcsname{\href{S6_SERIES_GRIND_III.pdf\#nameddest=tag-024D}{III.{D.SL15} (p.~185)} }
\expandafter\def\csname SRef@maps:eq:D.SL16\endcsname{\href{S6_SERIES_GRIND_III.pdf\#nameddest=tag-024E}{III.{D.SL16} (p.~185)} }
\expandafter\def\csname SRef@maps:eq:D.SL17\endcsname{\href{S6_SERIES_GRIND_III.pdf\#nameddest=tag-024F}{III.{D.SL17} (p.~185)} }
\expandafter\def\csname SRef@maps:eq:D.SL18\endcsname{\href{S6_SERIES_GRIND_III.pdf\#nameddest=tag-024G}{III.{D.SL18} (p.~185)} }
\expandafter\def\csname SRef@maps:eq:D.SL19\endcsname{\href{S6_SERIES_GRIND_III.pdf\#nameddest=tag-024H}{III.{D.SL19} (p.~185)} }
\expandafter\def\csname SRef@maps:eq:D.SL20\endcsname{\href{S6_SERIES_GRIND_III.pdf\#nameddest=tag-024J}{III.{D.SL20} (p.~186)} }
\expandafter\def\csname SRef@maps:eq:D.SL21\endcsname{\href{S6_SERIES_GRIND_III.pdf\#nameddest=tag-024K}{III.{D.SL21} (p.~186)} }
\expandafter\def\csname SRef@maps:eq:D.SL22\endcsname{\href{S6_SERIES_GRIND_III.pdf\#nameddest=tag-024L}{III.{D.SL22} (p.~186)} }
\expandafter\def\csname SRef@maps:eq:D.SL24\endcsname{\href{S6_SERIES_GRIND_III.pdf\#nameddest=tag-024N}{III.{D.SL24} (p.~186)} }
\expandafter\def\csname SRef@maps:eq:D.SL23\endcsname{\href{S6_SERIES_GRIND_III.pdf\#nameddest=tag-024M}{III.{D.SL23} (p.~187)} }
\expandafter\def\csname SRef@maps:eq:B.5.12\endcsname{\href{S6_SERIES_GRIND_III.pdf\#nameddest=tag-01VA}{III.{B.5.12} (p.~24)} }
\expandafter\def\csname SRef@maps:eq:D.VC1\endcsname{\href{S6_SERIES_GRIND_III.pdf\#nameddest=tag-026E}{III.{D.VC1} (p.~187)} }
\expandafter\def\csname SRef@maps:eq:D.VC7\endcsname{\href{S6_SERIES_GRIND_III.pdf\#nameddest=tag-026L}{III.{D.VC7} (p.~188)} }
\expandafter\def\csname SRef@maps:eq:D.VC2\endcsname{\href{S6_SERIES_GRIND_III.pdf\#nameddest=tag-026G}{III.{D.VC2} (p.~188)} }
\expandafter\def\csname SRef@maps:eq:D.VC3\endcsname{\href{S6_SERIES_GRIND_III.pdf\#nameddest=tag-026H}{III.{D.VC3} (p.~188)} }
\expandafter\def\csname SRef@maps:eq:D.VC4\endcsname{\href{S6_SERIES_GRIND_III.pdf\#nameddest=tag-026I}{III.{D.VC4} (p.~188)} }
\expandafter\def\csname SRef@maps:eq:D.VC5\endcsname{\href{S6_SERIES_GRIND_III.pdf\#nameddest=tag-026J}{III.{D.VC5} (p.~188)} }
\expandafter\def\csname SRef@maps:eq:D.VC8\endcsname{\href{S6_SERIES_GRIND_III.pdf\#nameddest=tag-026M}{III.{D.VC8} (p.~189)} }
\expandafter\def\csname SRef@maps:eq:D.VC6\endcsname{\href{S6_SERIES_GRIND_III.pdf\#nameddest=tag-026K}{III.{D.VC6} (p.~189)} }
\expandafter\def\csname SRef@maps:eq:D.VC9\endcsname{\href{S6_SERIES_GRIND_III.pdf\#nameddest=tag-026N}{III.{D.VC9} (p.~190)} }
\expandafter\def\csname SRef@maps:eq:D.VC10\endcsname{\href{S6_SERIES_GRIND_III.pdf\#nameddest=tag-026F}{III.{D.VC10} (p.~190)} }
\expandafter\def\csname SRef@maps:eq:B.5.13\endcsname{\href{S6_SERIES_GRIND_III.pdf\#nameddest=tag-01VB}{III.{B.5.13} (p.~25)} }
\expandafter\def\csname SRef@maps:eq:D.IC1\endcsname{\href{S6_SERIES_GRIND_III.pdf\#nameddest=tag-01ZY}{III.{D.IC1} (p.~191)} }
\expandafter\def\csname SRef@maps:eq:D.IC2\endcsname{\href{S6_SERIES_GRIND_III.pdf\#nameddest=tag-01ZZ}{III.{D.IC2} (p.~191)} }
\expandafter\def\csname SRef@maps:eq:D.IC3\endcsname{\href{S6_SERIES_GRIND_III.pdf\#nameddest=tag-0200}{III.{D.IC3} (p.~191)} }
\expandafter\def\csname SRef@maps:eq:D.IC4\endcsname{\href{S6_SERIES_GRIND_III.pdf\#nameddest=tag-0201}{III.{D.IC4} (p.~191)} }
\expandafter\def\csname SRef@maps:eq:D.IC5\endcsname{\href{S6_SERIES_GRIND_III.pdf\#nameddest=tag-0202}{III.{D.IC5} (p.~191)} }
\expandafter\def\csname SRef@maps:eq:D.IC6\endcsname{\href{S6_SERIES_GRIND_III.pdf\#nameddest=tag-0203}{III.{D.IC6} (p.~191)} }
\expandafter\def\csname SRef@maps:eq:D.IC7\endcsname{\href{S6_SERIES_GRIND_III.pdf\#nameddest=tag-0204}{III.{D.IC7} (p.~191)} }
\expandafter\def\csname SRef@maps:eq:D.IC8\endcsname{\href{S6_SERIES_GRIND_III.pdf\#nameddest=tag-0205}{III.{D.IC8} (p.~192)} }
\expandafter\def\csname SRef@maps:eq:D.IC9\endcsname{\href{S6_SERIES_GRIND_III.pdf\#nameddest=tag-0206}{III.{D.IC9} (p.~192)} }
\expandafter\def\csname SRef@maps:eq:D.CI1\endcsname{\href{S6_SERIES_GRIND_III.pdf\#nameddest=tag-01YN}{III.{D.CI1} (p.~195)} }
\expandafter\def\csname SRef@maps:eq:B.5.14\endcsname{\href{S6_SERIES_GRIND_III.pdf\#nameddest=tag-01VC}{III.{B.5.14} (p.~25)} }
\expandafter\def\csname SRef@maps:eq:D.OS1\endcsname{\href{S6_SERIES_GRIND_III.pdf\#nameddest=tag-022O}{III.{D.OS1} (p.~196)} }
\expandafter\def\csname SRef@maps:eq:D.OS2\endcsname{\href{S6_SERIES_GRIND_III.pdf\#nameddest=tag-022P}{III.{D.OS2} (p.~196)} }
\expandafter\def\csname SRef@maps:eq:D.OS3\endcsname{\href{S6_SERIES_GRIND_III.pdf\#nameddest=tag-022Q}{III.{D.OS3} (p.~196)} }
\expandafter\def\csname SRef@maps:eq:D.OS4\endcsname{\href{S6_SERIES_GRIND_III.pdf\#nameddest=tag-022R}{III.{D.OS4} (p.~196)} }
\expandafter\def\csname SRef@maps:eq:D.PC1\endcsname{\href{S6_SERIES_GRIND_III.pdf\#nameddest=tag-0233}{III.{D.PC1} (p.~196)} }
\expandafter\def\csname SRef@maps:eq:D.PC2\endcsname{\href{S6_SERIES_GRIND_III.pdf\#nameddest=tag-0239}{III.{D.PC2} (p.~197)} }
\expandafter\def\csname SRef@maps:eq:D.PC3\endcsname{\href{S6_SERIES_GRIND_III.pdf\#nameddest=tag-023A}{III.{D.PC3} (p.~197)} }
\expandafter\def\csname SRef@maps:eq:D.PC4\endcsname{\href{S6_SERIES_GRIND_III.pdf\#nameddest=tag-023B}{III.{D.PC4} (p.~197)} }
\expandafter\def\csname SRef@maps:eq:D.PC5\endcsname{\href{S6_SERIES_GRIND_III.pdf\#nameddest=tag-023C}{III.{D.PC5} (p.~197)} }
\expandafter\def\csname SRef@maps:eq:D.PC6\endcsname{\href{S6_SERIES_GRIND_III.pdf\#nameddest=tag-023D}{III.{D.PC6} (p.~197)} }
\expandafter\def\csname SRef@maps:eq:B.5.15\endcsname{\href{S6_SERIES_GRIND_III.pdf\#nameddest=tag-01VD}{III.{B.5.15} (p.~25)} }
\expandafter\def\csname SRef@maps:eq:D.PC7\endcsname{\href{S6_SERIES_GRIND_III.pdf\#nameddest=tag-023E}{III.{D.PC7} (p.~197)} }
\expandafter\def\csname SRef@maps:eq:D.PC8\endcsname{\href{S6_SERIES_GRIND_III.pdf\#nameddest=tag-023F}{III.{D.PC8} (p.~197)} }
\expandafter\def\csname SRef@maps:eq:D.PC9\endcsname{\href{S6_SERIES_GRIND_III.pdf\#nameddest=tag-023G}{III.{D.PC9} (p.~198)} }
\expandafter\def\csname SRef@maps:eq:D.PC10\endcsname{\href{S6_SERIES_GRIND_III.pdf\#nameddest=tag-0234}{III.{D.PC10} (p.~198)} }
\expandafter\def\csname SRef@maps:eq:D.PC11\endcsname{\href{S6_SERIES_GRIND_III.pdf\#nameddest=tag-0235}{III.{D.PC11} (p.~198)} }
\expandafter\def\csname SRef@maps:eq:D.PC12\endcsname{\href{S6_SERIES_GRIND_III.pdf\#nameddest=tag-0236}{III.{D.PC12} (p.~198)} }
\expandafter\def\csname SRef@maps:eq:D.PC13\endcsname{\href{S6_SERIES_GRIND_III.pdf\#nameddest=tag-0237}{III.{D.PC13} (p.~198)} }
\expandafter\def\csname SRef@maps:eq:D.PC14\endcsname{\href{S6_SERIES_GRIND_III.pdf\#nameddest=tag-0238}{III.{D.PC14} (p.~198)} }
\expandafter\def\csname SRef@maps:eq:D.OS5\endcsname{\href{S6_SERIES_GRIND_III.pdf\#nameddest=tag-022S}{III.{D.OS5} (p.~199)} }
\expandafter\def\csname SRef@maps:eq:B.5.16\endcsname{\href{S6_SERIES_GRIND_III.pdf\#nameddest=tag-01VE}{III.{B.5.16} (p.~25)} }
\expandafter\def\csname SRef@maps:eq:D.OS6\endcsname{\href{S6_SERIES_GRIND_III.pdf\#nameddest=tag-022T}{III.{D.OS6} (p.~199)} }
\expandafter\def\csname SRef@maps:eq:D.OS7\endcsname{\href{S6_SERIES_GRIND_III.pdf\#nameddest=tag-022U}{III.{D.OS7} (p.~199)} }
\expandafter\def\csname SRef@maps:eq:D.OS8\endcsname{\href{S6_SERIES_GRIND_III.pdf\#nameddest=tag-022V}{III.{D.OS8} (p.~200)} }
\expandafter\def\csname SRef@maps:eq:D.OF1\endcsname{\href{S6_SERIES_GRIND_III.pdf\#nameddest=tag-0217}{III.{D.OF1} (p.~200)} }
\expandafter\def\csname SRef@maps:eq:D.OF2\endcsname{\href{S6_SERIES_GRIND_III.pdf\#nameddest=tag-021I}{III.{D.OF2} (p.~200)} }
\expandafter\def\csname SRef@maps:eq:D.OF3\endcsname{\href{S6_SERIES_GRIND_III.pdf\#nameddest=tag-021J}{III.{D.OF3} (p.~201)} }
\expandafter\def\csname SRef@maps:eq:D.OF4\endcsname{\href{S6_SERIES_GRIND_III.pdf\#nameddest=tag-021K}{III.{D.OF4} (p.~201)} }
\expandafter\def\csname SRef@maps:eq:D.OF5\endcsname{\href{S6_SERIES_GRIND_III.pdf\#nameddest=tag-021L}{III.{D.OF5} (p.~201)} }
\expandafter\def\csname SRef@maps:eq:B.5.17\endcsname{\href{S6_SERIES_GRIND_III.pdf\#nameddest=tag-01VF}{III.{B.5.17} (p.~25)} }
\expandafter\def\csname SRef@maps:eq:D.OF6\endcsname{\href{S6_SERIES_GRIND_III.pdf\#nameddest=tag-021M}{III.{D.OF6} (p.~201)} }
\expandafter\def\csname SRef@maps:eq:D.OF7\endcsname{\href{S6_SERIES_GRIND_III.pdf\#nameddest=tag-021N}{III.{D.OF7} (p.~201)} }
\expandafter\def\csname SRef@maps:eq:D.OF8\endcsname{\href{S6_SERIES_GRIND_III.pdf\#nameddest=tag-021O}{III.{D.OF8} (p.~201)} }
\expandafter\def\csname SRef@maps:eq:D.OF9\endcsname{\href{S6_SERIES_GRIND_III.pdf\#nameddest=tag-021P}{III.{D.OF9} (p.~202)} }
\expandafter\def\csname SRef@maps:eq:D.OF10\endcsname{\href{S6_SERIES_GRIND_III.pdf\#nameddest=tag-0218}{III.{D.OF10} (p.~202)} }
\expandafter\def\csname SRef@maps:eq:D.OF11\endcsname{\href{S6_SERIES_GRIND_III.pdf\#nameddest=tag-0219}{III.{D.OF11} (p.~202)} }
\expandafter\def\csname SRef@maps:eq:D.OF12\endcsname{\href{S6_SERIES_GRIND_III.pdf\#nameddest=tag-021A}{III.{D.OF12} (p.~202)} }
\expandafter\def\csname SRef@maps:eq:D.OF13\endcsname{\href{S6_SERIES_GRIND_III.pdf\#nameddest=tag-021B}{III.{D.OF13} (p.~202)} }
\expandafter\def\csname SRef@maps:eq:D.OF14\endcsname{\href{S6_SERIES_GRIND_III.pdf\#nameddest=tag-021C}{III.{D.OF14} (p.~202)} }
\expandafter\def\csname SRef@maps:eq:B.5.18\endcsname{\href{S6_SERIES_GRIND_III.pdf\#nameddest=tag-01VG}{III.{B.5.18} (p.~26)} }
\expandafter\def\csname SRef@maps:eq:D.OF15\endcsname{\href{S6_SERIES_GRIND_III.pdf\#nameddest=tag-021D}{III.{D.OF15} (p.~203)} }
\expandafter\def\csname SRef@maps:eq:D.OF16\endcsname{\href{S6_SERIES_GRIND_III.pdf\#nameddest=tag-021E}{III.{D.OF16} (p.~203)} }
\expandafter\def\csname SRef@maps:eq:D.OF17\endcsname{\href{S6_SERIES_GRIND_III.pdf\#nameddest=tag-021F}{III.{D.OF17} (p.~203)} }
\expandafter\def\csname SRef@maps:eq:D.OF18\endcsname{\href{S6_SERIES_GRIND_III.pdf\#nameddest=tag-021G}{III.{D.OF18} (p.~203)} }
\expandafter\def\csname SRef@maps:eq:D.OL1\endcsname{\href{S6_SERIES_GRIND_III.pdf\#nameddest=tag-021Q}{III.{D.OL1} (p.~204)} }
\expandafter\def\csname SRef@maps:eq:D.OL2\endcsname{\href{S6_SERIES_GRIND_III.pdf\#nameddest=tag-021U}{III.{D.OL2} (p.~204)} }
\expandafter\def\csname SRef@maps:eq:D.OL3\endcsname{\href{S6_SERIES_GRIND_III.pdf\#nameddest=tag-021V}{III.{D.OL3} (p.~204)} }
\expandafter\def\csname SRef@maps:eq:D.OL4\endcsname{\href{S6_SERIES_GRIND_III.pdf\#nameddest=tag-021W}{III.{D.OL4} (p.~204)} }
\expandafter\def\csname SRef@maps:eq:D.OL5\endcsname{\href{S6_SERIES_GRIND_III.pdf\#nameddest=tag-021X}{III.{D.OL5} (p.~204)} }
\expandafter\def\csname SRef@maps:eq:C.1.1\endcsname{\href{S6_SERIES_GRIND_III.pdf\#nameddest=tag-01XE}{III.{C.1.1} (p.~27)} }
\expandafter\def\csname SRef@maps:eq:D.OL6\endcsname{\href{S6_SERIES_GRIND_III.pdf\#nameddest=tag-021Y}{III.{D.OL6} (p.~204)} }
\expandafter\def\csname SRef@maps:eq:D.OL7\endcsname{\href{S6_SERIES_GRIND_III.pdf\#nameddest=tag-021Z}{III.{D.OL7} (p.~204)} }
\expandafter\def\csname SRef@maps:eq:D.OL8\endcsname{\href{S6_SERIES_GRIND_III.pdf\#nameddest=tag-0220}{III.{D.OL8} (p.~205)} }
\expandafter\def\csname SRef@maps:eq:D.OL9\endcsname{\href{S6_SERIES_GRIND_III.pdf\#nameddest=tag-0221}{III.{D.OL9} (p.~205)} }
\expandafter\def\csname SRef@maps:eq:D.OL10\endcsname{\href{S6_SERIES_GRIND_III.pdf\#nameddest=tag-021R}{III.{D.OL10} (p.~205)} }
\expandafter\def\csname SRef@maps:eq:D.OL11\endcsname{\href{S6_SERIES_GRIND_III.pdf\#nameddest=tag-021S}{III.{D.OL11} (p.~205)} }
\expandafter\def\csname SRef@maps:eq:D.OL12\endcsname{\href{S6_SERIES_GRIND_III.pdf\#nameddest=tag-021T}{III.{D.OL12} (p.~205)} }
\expandafter\def\csname SRef@maps:eq:D.OF19\endcsname{\href{S6_SERIES_GRIND_III.pdf\#nameddest=tag-021H}{III.{D.OF19} (p.~206)} }
\expandafter\def\csname SRef@maps:eq:C.1.2\endcsname{\href{S6_SERIES_GRIND_III.pdf\#nameddest=tag-01XF}{III.{C.1.2} (p.~28)} }
\expandafter\def\csname SRef@maps:eq:C.1.3\endcsname{\href{S6_SERIES_GRIND_III.pdf\#nameddest=tag-01XG}{III.{C.1.3} (p.~30)} }
\expandafter\def\csname SRef@maps:eq:E.1\endcsname{\href{S6_SERIES_GRIND_V.pdf\#nameddest=tag-029F}{V.{E.1} (p.~14)} }
\expandafter\def\csname SRef@maps:eq:E.2\endcsname{\href{S6_SERIES_GRIND_V.pdf\#nameddest=tag-029O}{V.{E.2} (p.~15)} }
\expandafter\def\csname SRef@maps:eq:C.1.4\endcsname{\href{S6_SERIES_GRIND_III.pdf\#nameddest=tag-01XH}{III.{C.1.4} (p.~30)} }
\expandafter\def\csname SRef@maps:eq:E.3\endcsname{\href{S6_SERIES_GRIND_V.pdf\#nameddest=tag-029Q}{V.{E.3} (p.~15)} }
\expandafter\def\csname SRef@maps:eq:E.4\endcsname{\href{S6_SERIES_GRIND_V.pdf\#nameddest=tag-029R}{V.{E.4} (p.~15)} }
\expandafter\def\csname SRef@maps:eq:E.5\endcsname{\href{S6_SERIES_GRIND_V.pdf\#nameddest=tag-029S}{V.{E.5} (p.~15)} }
\expandafter\def\csname SRef@maps:eq:E.6\endcsname{\href{S6_SERIES_GRIND_V.pdf\#nameddest=tag-029T}{V.{E.6} (p.~16)} }
\expandafter\def\csname SRef@maps:eq:E.7\endcsname{\href{S6_SERIES_GRIND_V.pdf\#nameddest=tag-029U}{V.{E.7} (p.~16)} }
\expandafter\def\csname SRef@maps:eq:E.8\endcsname{\href{S6_SERIES_GRIND_V.pdf\#nameddest=tag-029V}{V.{E.8} (p.~18)} }
\expandafter\def\csname SRef@maps:eq:E.9\endcsname{\href{S6_SERIES_GRIND_V.pdf\#nameddest=tag-029W}{V.{E.9} (p.~18)} }
\expandafter\def\csname SRef@maps:eq:E.10\endcsname{\href{S6_SERIES_GRIND_V.pdf\#nameddest=tag-029G}{V.{E.10} (p.~18)} }
\expandafter\def\csname SRef@maps:eq:E.11\endcsname{\href{S6_SERIES_GRIND_V.pdf\#nameddest=tag-029H}{V.{E.11} (p.~18)} }
\expandafter\def\csname SRef@maps:eq:E.I1\endcsname{\href{S6_SERIES_GRIND_V.pdf\#nameddest=tag-02B5}{V.{E.I1} (p.~19)} }
\expandafter\def\csname SRef@maps:eq:D.C1\endcsname{\href{S6_SERIES_GRIND_III.pdf\#nameddest=tag-01Y6}{III.{D.C1} (p.~59)} }
\expandafter\def\csname SRef@maps:eq:E.I2\endcsname{\href{S6_SERIES_GRIND_V.pdf\#nameddest=tag-02B6}{V.{E.I2} (p.~20)} }
\expandafter\def\csname SRef@maps:eq:E.I3\endcsname{\href{S6_SERIES_GRIND_V.pdf\#nameddest=tag-02B7}{V.{E.I3} (p.~20)} }
\expandafter\def\csname SRef@maps:eq:E.I4\endcsname{\href{S6_SERIES_GRIND_V.pdf\#nameddest=tag-02B8}{V.{E.I4} (p.~20)} }
\expandafter\def\csname SRef@maps:eq:E.I5\endcsname{\href{S6_SERIES_GRIND_V.pdf\#nameddest=tag-02B9}{V.{E.I5} (p.~21)} }
\expandafter\def\csname SRef@maps:eq:E.12\endcsname{\href{S6_SERIES_GRIND_V.pdf\#nameddest=tag-029I}{V.{E.12} (p.~21)} }
\expandafter\def\csname SRef@maps:eq:E.13\endcsname{\href{S6_SERIES_GRIND_V.pdf\#nameddest=tag-029J}{V.{E.13} (p.~21)} }
\expandafter\def\csname SRef@maps:eq:E.14\endcsname{\href{S6_SERIES_GRIND_V.pdf\#nameddest=tag-029K}{V.{E.14} (p.~22)} }
\expandafter\def\csname SRef@maps:eq:E.16\endcsname{\href{S6_SERIES_GRIND_V.pdf\#nameddest=tag-029M}{V.{E.16} (p.~23)} }
\expandafter\def\csname SRef@maps:eq:E.17\endcsname{\href{S6_SERIES_GRIND_V.pdf\#nameddest=tag-029N}{V.{E.17} (p.~23)} }
\expandafter\def\csname SRef@maps:eq:E.20\endcsname{\href{S6_SERIES_GRIND_V.pdf\#nameddest=tag-029P}{V.{E.20} (p.~25)} }
\expandafter\def\csname SRef@maps:eq:D.R1\endcsname{\href{S6_SERIES_GRIND_III.pdf\#nameddest=tag-023H}{III.{D.R1} (p.~64)} }
\expandafter\def\csname SRef@maps:eq:E.MV1\endcsname{\href{S6_SERIES_GRIND_V.pdf\#nameddest=tag-02BA}{V.{E.MV1} (p.~26)} }
\expandafter\def\csname SRef@maps:eq:E.MV2\endcsname{\href{S6_SERIES_GRIND_V.pdf\#nameddest=tag-02BB}{V.{E.MV2} (p.~27)} }
\expandafter\def\csname SRef@maps:eq:E.MV3\endcsname{\href{S6_SERIES_GRIND_V.pdf\#nameddest=tag-02BC}{V.{E.MV3} (p.~27)} }
\expandafter\def\csname SRef@maps:eq:E.MV4\endcsname{\href{S6_SERIES_GRIND_V.pdf\#nameddest=tag-02BD}{V.{E.MV4} (p.~27)} }
\expandafter\def\csname SRef@maps:eq:E.MV5\endcsname{\href{S6_SERIES_GRIND_V.pdf\#nameddest=tag-02BE}{V.{E.MV5} (p.~27)} }
\expandafter\def\csname SRef@maps:eq:E.MV6\endcsname{\href{S6_SERIES_GRIND_V.pdf\#nameddest=tag-02BF}{V.{E.MV6} (p.~28)} }
\expandafter\def\csname SRef@maps:eq:E.MV7\endcsname{\href{S6_SERIES_GRIND_V.pdf\#nameddest=tag-02BG}{V.{E.MV7} (p.~28)} }
\expandafter\def\csname SRef@maps:eq:E.MV8\endcsname{\href{S6_SERIES_GRIND_V.pdf\#nameddest=tag-02BH}{V.{E.MV8} (p.~28)} }
\expandafter\def\csname SRef@maps:eq:E.MV9\endcsname{\href{S6_SERIES_GRIND_V.pdf\#nameddest=tag-02BI}{V.{E.MV9} (p.~29)} }
\expandafter\def\csname SRef@maps:eq:E.B1\endcsname{\href{S6_SERIES_GRIND_V.pdf\#nameddest=tag-029X}{V.{E.B1} (p.~30)} }
\expandafter\def\csname SRef@maps:eq:D.R2\endcsname{\href{S6_SERIES_GRIND_III.pdf\#nameddest=tag-023I}{III.{D.R2} (p.~64)} }
\expandafter\def\csname SRef@maps:eq:E.B2\endcsname{\href{S6_SERIES_GRIND_V.pdf\#nameddest=tag-029Y}{V.{E.B2} (p.~30)} }
\expandafter\def\csname SRef@maps:eq:E.B3\endcsname{\href{S6_SERIES_GRIND_V.pdf\#nameddest=tag-029Z}{V.{E.B3} (p.~31)} }
\expandafter\def\csname SRef@maps:eq:E.B4\endcsname{\href{S6_SERIES_GRIND_V.pdf\#nameddest=tag-02A0}{V.{E.B4} (p.~31)} }
\expandafter\def\csname SRef@maps:eq:E.B5\endcsname{\href{S6_SERIES_GRIND_V.pdf\#nameddest=tag-02A1}{V.{E.B5} (p.~31)} }
\expandafter\def\csname SRef@maps:eq:E.B6\endcsname{\href{S6_SERIES_GRIND_V.pdf\#nameddest=tag-02A2}{V.{E.B6} (p.~32)} }
\expandafter\def\csname SRef@maps:eq:E.B7\endcsname{\href{S6_SERIES_GRIND_V.pdf\#nameddest=tag-02A3}{V.{E.B7} (p.~32)} }
\expandafter\def\csname SRef@maps:eq:E.S1\endcsname{\href{S6_SERIES_GRIND_V.pdf\#nameddest=tag-02CC}{V.{E.S1} (p.~33)} }
\expandafter\def\csname SRef@maps:eq:E.S2\endcsname{\href{S6_SERIES_GRIND_V.pdf\#nameddest=tag-02CN}{V.{E.S2} (p.~33)} }
\expandafter\def\csname SRef@maps:eq:E.S3\endcsname{\href{S6_SERIES_GRIND_V.pdf\#nameddest=tag-02CY}{V.{E.S3} (p.~34)} }
\expandafter\def\csname SRef@maps:eq:E.S4\endcsname{\href{S6_SERIES_GRIND_V.pdf\#nameddest=tag-02D9}{V.{E.S4} (p.~34)} }
\expandafter\def\csname SRef@maps:eq:D.R3\endcsname{\href{S6_SERIES_GRIND_III.pdf\#nameddest=tag-023J}{III.{D.R3} (p.~64)} }
\expandafter\def\csname SRef@maps:eq:E.S5\endcsname{\href{S6_SERIES_GRIND_V.pdf\#nameddest=tag-02DK}{V.{E.S5} (p.~34)} }
\expandafter\def\csname SRef@maps:eq:E.S12\endcsname{\href{S6_SERIES_GRIND_V.pdf\#nameddest=tag-02CF}{V.{E.S12} (p.~35)} }
\expandafter\def\csname SRef@maps:eq:E.S6\endcsname{\href{S6_SERIES_GRIND_V.pdf\#nameddest=tag-02DR}{V.{E.S6} (p.~35)} }
\expandafter\def\csname SRef@maps:eq:E.S7\endcsname{\href{S6_SERIES_GRIND_V.pdf\#nameddest=tag-02DS}{V.{E.S7} (p.~35)} }
\expandafter\def\csname SRef@maps:eq:E.S8\endcsname{\href{S6_SERIES_GRIND_V.pdf\#nameddest=tag-02DT}{V.{E.S8} (p.~36)} }
\expandafter\def\csname SRef@maps:eq:E.S9\endcsname{\href{S6_SERIES_GRIND_V.pdf\#nameddest=tag-02DU}{V.{E.S9} (p.~36)} }
\expandafter\def\csname SRef@maps:eq:E.S10\endcsname{\href{S6_SERIES_GRIND_V.pdf\#nameddest=tag-02CD}{V.{E.S10} (p.~36)} }
\expandafter\def\csname SRef@maps:eq:E.S11\endcsname{\href{S6_SERIES_GRIND_V.pdf\#nameddest=tag-02CE}{V.{E.S11} (p.~36)} }
\expandafter\def\csname SRef@maps:eq:E.S13\endcsname{\href{S6_SERIES_GRIND_V.pdf\#nameddest=tag-02CG}{V.{E.S13} (p.~37)} }
\expandafter\def\csname SRef@maps:eq:E.S14\endcsname{\href{S6_SERIES_GRIND_V.pdf\#nameddest=tag-02CH}{V.{E.S14} (p.~37)} }
\expandafter\def\csname SRef@maps:eq:D.R4\endcsname{\href{S6_SERIES_GRIND_III.pdf\#nameddest=tag-023K}{III.{D.R4} (p.~64)} }
\expandafter\def\csname SRef@maps:eq:E.S15\endcsname{\href{S6_SERIES_GRIND_V.pdf\#nameddest=tag-02CI}{V.{E.S15} (p.~37)} }
\expandafter\def\csname SRef@maps:eq:E.S16\endcsname{\href{S6_SERIES_GRIND_V.pdf\#nameddest=tag-02CJ}{V.{E.S16} (p.~37)} }
\expandafter\def\csname SRef@maps:eq:E.S20\endcsname{\href{S6_SERIES_GRIND_V.pdf\#nameddest=tag-02CO}{V.{E.S20} (p.~37)} }
\expandafter\def\csname SRef@maps:eq:E.S17\endcsname{\href{S6_SERIES_GRIND_V.pdf\#nameddest=tag-02CK}{V.{E.S17} (p.~37)} }
\expandafter\def\csname SRef@maps:eq:E.S18\endcsname{\href{S6_SERIES_GRIND_V.pdf\#nameddest=tag-02CL}{V.{E.S18} (p.~37)} }
\expandafter\def\csname SRef@maps:eq:E.S19\endcsname{\href{S6_SERIES_GRIND_V.pdf\#nameddest=tag-02CM}{V.{E.S19} (p.~38)} }
\expandafter\def\csname SRef@maps:eq:E.S21\endcsname{\href{S6_SERIES_GRIND_V.pdf\#nameddest=tag-02CP}{V.{E.S21} (p.~39)} }
\expandafter\def\csname SRef@maps:eq:E.S22\endcsname{\href{S6_SERIES_GRIND_V.pdf\#nameddest=tag-02CQ}{V.{E.S22} (p.~40)} }
\expandafter\def\csname SRef@maps:eq:E.S23\endcsname{\href{S6_SERIES_GRIND_V.pdf\#nameddest=tag-02CR}{V.{E.S23} (p.~40)} }
\expandafter\def\csname SRef@maps:eq:E.S24\endcsname{\href{S6_SERIES_GRIND_V.pdf\#nameddest=tag-02CS}{V.{E.S24} (p.~40)} }
\expandafter\def\csname SRef@maps:eq:D.S1\endcsname{\href{S6_SERIES_GRIND_III.pdf\#nameddest=tag-023L}{III.{D.S1} (p.~65)} }
\expandafter\def\csname SRef@maps:eq:E.S25\endcsname{\href{S6_SERIES_GRIND_V.pdf\#nameddest=tag-02CT}{V.{E.S25} (p.~40)} }
\expandafter\def\csname SRef@maps:eq:E.S26\endcsname{\href{S6_SERIES_GRIND_V.pdf\#nameddest=tag-02CU}{V.{E.S26} (p.~40)} }
\expandafter\def\csname SRef@maps:eq:E.S27\endcsname{\href{S6_SERIES_GRIND_V.pdf\#nameddest=tag-02CV}{V.{E.S27} (p.~41)} }
\expandafter\def\csname SRef@maps:eq:E.S28\endcsname{\href{S6_SERIES_GRIND_V.pdf\#nameddest=tag-02CW}{V.{E.S28} (p.~41)} }
\expandafter\def\csname SRef@maps:eq:E.S29\endcsname{\href{S6_SERIES_GRIND_V.pdf\#nameddest=tag-02CX}{V.{E.S29} (p.~42)} }
\expandafter\def\csname SRef@maps:eq:E.S30\endcsname{\href{S6_SERIES_GRIND_V.pdf\#nameddest=tag-02CZ}{V.{E.S30} (p.~42)} }
\expandafter\def\csname SRef@maps:eq:E.S31\endcsname{\href{S6_SERIES_GRIND_V.pdf\#nameddest=tag-02D0}{V.{E.S31} (p.~42)} }
\expandafter\def\csname SRef@maps:eq:E.S32\endcsname{\href{S6_SERIES_GRIND_V.pdf\#nameddest=tag-02D1}{V.{E.S32} (p.~42)} }
\expandafter\def\csname SRef@maps:eq:E.S33\endcsname{\href{S6_SERIES_GRIND_V.pdf\#nameddest=tag-02D2}{V.{E.S33} (p.~42)} }
\expandafter\def\csname SRef@maps:eq:E.S34\endcsname{\href{S6_SERIES_GRIND_V.pdf\#nameddest=tag-02D3}{V.{E.S34} (p.~43)} }
\expandafter\def\csname SRef@maps:eq:D.S2\endcsname{\href{S6_SERIES_GRIND_III.pdf\#nameddest=tag-023R}{III.{D.S2} (p.~65)} }
\expandafter\def\csname SRef@maps:eq:E.S35\endcsname{\href{S6_SERIES_GRIND_V.pdf\#nameddest=tag-02D4}{V.{E.S35} (p.~43)} }
\expandafter\def\csname SRef@maps:eq:E.S36\endcsname{\href{S6_SERIES_GRIND_V.pdf\#nameddest=tag-02D5}{V.{E.S36} (p.~43)} }
\expandafter\def\csname SRef@maps:eq:E.S37\endcsname{\href{S6_SERIES_GRIND_V.pdf\#nameddest=tag-02D6}{V.{E.S37} (p.~43)} }
\expandafter\def\csname SRef@maps:eq:E.S38\endcsname{\href{S6_SERIES_GRIND_V.pdf\#nameddest=tag-02D7}{V.{E.S38} (p.~43)} }
\expandafter\def\csname SRef@maps:eq:E.S39\endcsname{\href{S6_SERIES_GRIND_V.pdf\#nameddest=tag-02D8}{V.{E.S39} (p.~44)} }
\expandafter\def\csname SRef@maps:eq:E.S40\endcsname{\href{S6_SERIES_GRIND_V.pdf\#nameddest=tag-02DA}{V.{E.S40} (p.~44)} }
\expandafter\def\csname SRef@maps:eq:E.S41\endcsname{\href{S6_SERIES_GRIND_V.pdf\#nameddest=tag-02DB}{V.{E.S41} (p.~44)} }
\expandafter\def\csname SRef@maps:eq:E.S42\endcsname{\href{S6_SERIES_GRIND_V.pdf\#nameddest=tag-02DC}{V.{E.S42} (p.~44)} }
\expandafter\def\csname SRef@maps:eq:E.S43\endcsname{\href{S6_SERIES_GRIND_V.pdf\#nameddest=tag-02DD}{V.{E.S43} (p.~45)} }
\expandafter\def\csname SRef@maps:eq:E.S44\endcsname{\href{S6_SERIES_GRIND_V.pdf\#nameddest=tag-02DE}{V.{E.S44} (p.~45)} }
\expandafter\def\csname SRef@maps:eq:D.S3\endcsname{\href{S6_SERIES_GRIND_III.pdf\#nameddest=tag-023S}{III.{D.S3} (p.~65)} }
\expandafter\def\csname SRef@maps:eq:E.S45\endcsname{\href{S6_SERIES_GRIND_V.pdf\#nameddest=tag-02DF}{V.{E.S45} (p.~45)} }
\expandafter\def\csname SRef@maps:eq:E.S46\endcsname{\href{S6_SERIES_GRIND_V.pdf\#nameddest=tag-02DG}{V.{E.S46} (p.~45)} }
\expandafter\def\csname SRef@maps:eq:E.S47\endcsname{\href{S6_SERIES_GRIND_V.pdf\#nameddest=tag-02DH}{V.{E.S47} (p.~45)} }
\expandafter\def\csname SRef@maps:eq:E.S48\endcsname{\href{S6_SERIES_GRIND_V.pdf\#nameddest=tag-02DI}{V.{E.S48} (p.~46)} }
\expandafter\def\csname SRef@maps:eq:E.S49\endcsname{\href{S6_SERIES_GRIND_V.pdf\#nameddest=tag-02DJ}{V.{E.S49} (p.~46)} }
\expandafter\def\csname SRef@maps:eq:E.S50\endcsname{\href{S6_SERIES_GRIND_V.pdf\#nameddest=tag-02DL}{V.{E.S50} (p.~46)} }
\expandafter\def\csname SRef@maps:eq:E.S51\endcsname{\href{S6_SERIES_GRIND_V.pdf\#nameddest=tag-02DM}{V.{E.S51} (p.~46)} }
\expandafter\def\csname SRef@maps:eq:E.S52\endcsname{\href{S6_SERIES_GRIND_V.pdf\#nameddest=tag-02DN}{V.{E.S52} (p.~46)} }
\expandafter\def\csname SRef@maps:eq:E.S53\endcsname{\href{S6_SERIES_GRIND_V.pdf\#nameddest=tag-02DO}{V.{E.S53} (p.~46)} }
\expandafter\def\csname SRef@maps:eq:E.S54\endcsname{\href{S6_SERIES_GRIND_V.pdf\#nameddest=tag-02DP}{V.{E.S54} (p.~46)} }
\expandafter\def\csname SRef@maps:eq:D.S4\endcsname{\href{S6_SERIES_GRIND_III.pdf\#nameddest=tag-023T}{III.{D.S4} (p.~65)} }
\expandafter\def\csname SRef@maps:eq:E.S55\endcsname{\href{S6_SERIES_GRIND_V.pdf\#nameddest=tag-02DQ}{V.{E.S55} (p.~47)} }
\expandafter\def\csname SRef@maps:eq:E.C1\endcsname{\href{S6_SERIES_GRIND_V.pdf\#nameddest=tag-02A4}{V.{E.C1} (p.~48)} }
\expandafter\def\csname SRef@maps:eq:E.C2\endcsname{\href{S6_SERIES_GRIND_V.pdf\#nameddest=tag-02AG}{V.{E.C2} (p.~48)} }
\expandafter\def\csname SRef@maps:eq:E.C3\endcsname{\href{S6_SERIES_GRIND_V.pdf\#nameddest=tag-02AJ}{V.{E.C3} (p.~48)} }
\expandafter\def\csname SRef@maps:eq:E.C4\endcsname{\href{S6_SERIES_GRIND_V.pdf\#nameddest=tag-02AK}{V.{E.C4} (p.~48)} }
\expandafter\def\csname SRef@maps:eq:E.C19a\endcsname{\href{S6_SERIES_GRIND_V.pdf\#nameddest=tag-02AE}{V.{E.C19a} (p.~48)} }
\expandafter\def\csname SRef@maps:eq:E.C19b\endcsname{\href{S6_SERIES_GRIND_V.pdf\#nameddest=tag-02AF}{V.{E.C19b} (p.~48)} }
\expandafter\def\csname SRef@maps:eq:E.C5\endcsname{\href{S6_SERIES_GRIND_V.pdf\#nameddest=tag-02AL}{V.{E.C5} (p.~49)} }
\expandafter\def\csname SRef@maps:eq:E.C6\endcsname{\href{S6_SERIES_GRIND_V.pdf\#nameddest=tag-02AM}{V.{E.C6} (p.~49)} }
\expandafter\def\csname SRef@maps:eq:E.C7\endcsname{\href{S6_SERIES_GRIND_V.pdf\#nameddest=tag-02AN}{V.{E.C7} (p.~49)} }
\expandafter\def\csname SRef@maps:eq:D.S5\endcsname{\href{S6_SERIES_GRIND_III.pdf\#nameddest=tag-023U}{III.{D.S5} (p.~66)} }
\expandafter\def\csname SRef@maps:eq:E.C8\endcsname{\href{S6_SERIES_GRIND_V.pdf\#nameddest=tag-02AO}{V.{E.C8} (p.~49)} }
\expandafter\def\csname SRef@maps:eq:E.C9\endcsname{\href{S6_SERIES_GRIND_V.pdf\#nameddest=tag-02AP}{V.{E.C9} (p.~49)} }
\expandafter\def\csname SRef@maps:eq:E.C10\endcsname{\href{S6_SERIES_GRIND_V.pdf\#nameddest=tag-02A5}{V.{E.C10} (p.~49)} }
\expandafter\def\csname SRef@maps:eq:E.C11\endcsname{\href{S6_SERIES_GRIND_V.pdf\#nameddest=tag-02A6}{V.{E.C11} (p.~49)} }
\expandafter\def\csname SRef@maps:eq:E.C12\endcsname{\href{S6_SERIES_GRIND_V.pdf\#nameddest=tag-02A7}{V.{E.C12} (p.~50)} }
\expandafter\def\csname SRef@maps:eq:E.C20\endcsname{\href{S6_SERIES_GRIND_V.pdf\#nameddest=tag-02AH}{V.{E.C20} (p.~50)} }
\expandafter\def\csname SRef@maps:eq:E.C13\endcsname{\href{S6_SERIES_GRIND_V.pdf\#nameddest=tag-02A8}{V.{E.C13} (p.~51)} }
\expandafter\def\csname SRef@maps:eq:E.C14\endcsname{\href{S6_SERIES_GRIND_V.pdf\#nameddest=tag-02A9}{V.{E.C14} (p.~51)} }
\expandafter\def\csname SRef@maps:eq:E.C15\endcsname{\href{S6_SERIES_GRIND_V.pdf\#nameddest=tag-02AA}{V.{E.C15} (p.~51)} }
\expandafter\def\csname SRef@maps:eq:E.C16\endcsname{\href{S6_SERIES_GRIND_V.pdf\#nameddest=tag-02AB}{V.{E.C16} (p.~51)} }
\expandafter\def\csname SRef@maps:eq:D.S6\endcsname{\href{S6_SERIES_GRIND_III.pdf\#nameddest=tag-023V}{III.{D.S6} (p.~66)} }
\expandafter\def\csname SRef@maps:eq:E.C17\endcsname{\href{S6_SERIES_GRIND_V.pdf\#nameddest=tag-02AC}{V.{E.C17} (p.~51)} }
\expandafter\def\csname SRef@maps:eq:E.C18\endcsname{\href{S6_SERIES_GRIND_V.pdf\#nameddest=tag-02AD}{V.{E.C18} (p.~52)} }
\expandafter\def\csname SRef@maps:eq:E.C21\endcsname{\href{S6_SERIES_GRIND_V.pdf\#nameddest=tag-02AI}{V.{E.C21} (p.~52)} }
\expandafter\def\csname SRef@maps:eq:E.P1\endcsname{\href{S6_SERIES_GRIND_V.pdf\#nameddest=tag-02BJ}{V.{E.P1} (p.~52)} }
\expandafter\def\csname SRef@maps:eq:E.P2\endcsname{\href{S6_SERIES_GRIND_V.pdf\#nameddest=tag-02BS}{V.{E.P2} (p.~53)} }
\expandafter\def\csname SRef@maps:eq:E.P3\endcsname{\href{S6_SERIES_GRIND_V.pdf\#nameddest=tag-02BT}{V.{E.P3} (p.~53)} }
\expandafter\def\csname SRef@maps:eq:E.P4\endcsname{\href{S6_SERIES_GRIND_V.pdf\#nameddest=tag-02BU}{V.{E.P4} (p.~53)} }
\expandafter\def\csname SRef@maps:eq:E.P5\endcsname{\href{S6_SERIES_GRIND_V.pdf\#nameddest=tag-02BV}{V.{E.P5} (p.~53)} }
\expandafter\def\csname SRef@maps:eq:E.P6\endcsname{\href{S6_SERIES_GRIND_V.pdf\#nameddest=tag-02BW}{V.{E.P6} (p.~53)} }
\expandafter\def\csname SRef@maps:eq:E.P7\endcsname{\href{S6_SERIES_GRIND_V.pdf\#nameddest=tag-02BX}{V.{E.P7} (p.~53)} }
\expandafter\def\csname SRef@maps:eq:D.S7\endcsname{\href{S6_SERIES_GRIND_III.pdf\#nameddest=tag-023W}{III.{D.S7} (p.~67)} }
\expandafter\def\csname SRef@maps:eq:E.P8\endcsname{\href{S6_SERIES_GRIND_V.pdf\#nameddest=tag-02BY}{V.{E.P8} (p.~53)} }
\expandafter\def\csname SRef@maps:eq:E.P9\endcsname{\href{S6_SERIES_GRIND_V.pdf\#nameddest=tag-02BZ}{V.{E.P9} (p.~53)} }
\expandafter\def\csname SRef@maps:eq:E.P10\endcsname{\href{S6_SERIES_GRIND_V.pdf\#nameddest=tag-02BK}{V.{E.P10} (p.~54)} }
\expandafter\def\csname SRef@maps:eq:E.P11\endcsname{\href{S6_SERIES_GRIND_V.pdf\#nameddest=tag-02BL}{V.{E.P11} (p.~54)} }
\expandafter\def\csname SRef@maps:eq:E.P12\endcsname{\href{S6_SERIES_GRIND_V.pdf\#nameddest=tag-02BM}{V.{E.P12} (p.~54)} }
\expandafter\def\csname SRef@maps:eq:E.P13\endcsname{\href{S6_SERIES_GRIND_V.pdf\#nameddest=tag-02BN}{V.{E.P13} (p.~54)} }
\expandafter\def\csname SRef@maps:eq:E.P14\endcsname{\href{S6_SERIES_GRIND_V.pdf\#nameddest=tag-02BO}{V.{E.P14} (p.~55)} }
\expandafter\def\csname SRef@maps:eq:E.P15\endcsname{\href{S6_SERIES_GRIND_V.pdf\#nameddest=tag-02BP}{V.{E.P15} (p.~55)} }
\expandafter\def\csname SRef@maps:eq:E.P16\endcsname{\href{S6_SERIES_GRIND_V.pdf\#nameddest=tag-02BQ}{V.{E.P16} (p.~55)} }
\expandafter\def\csname SRef@maps:eq:E.P17\endcsname{\href{S6_SERIES_GRIND_V.pdf\#nameddest=tag-02BR}{V.{E.P17} (p.~56)} }
\expandafter\def\csname SRef@maps:eq:D.S8\endcsname{\href{S6_SERIES_GRIND_III.pdf\#nameddest=tag-023X}{III.{D.S8} (p.~67)} }
\expandafter\def\csname SRef@maps:eq:E.Q1\endcsname{\href{S6_SERIES_GRIND_V.pdf\#nameddest=tag-02C0}{V.{E.Q1} (p.~57)} }
\expandafter\def\csname SRef@maps:eq:E.Q2\endcsname{\href{S6_SERIES_GRIND_V.pdf\#nameddest=tag-02C4}{V.{E.Q2} (p.~57)} }
\expandafter\def\csname SRef@maps:eq:E.Q3\endcsname{\href{S6_SERIES_GRIND_V.pdf\#nameddest=tag-02C5}{V.{E.Q3} (p.~57)} }
\expandafter\def\csname SRef@maps:eq:E.Q4\endcsname{\href{S6_SERIES_GRIND_V.pdf\#nameddest=tag-02C6}{V.{E.Q4} (p.~57)} }
\expandafter\def\csname SRef@maps:eq:E.Q5\endcsname{\href{S6_SERIES_GRIND_V.pdf\#nameddest=tag-02C7}{V.{E.Q5} (p.~58)} }
\expandafter\def\csname SRef@maps:eq:E.Q6\endcsname{\href{S6_SERIES_GRIND_V.pdf\#nameddest=tag-02C8}{V.{E.Q6} (p.~58)} }
\expandafter\def\csname SRef@maps:eq:E.Q7\endcsname{\href{S6_SERIES_GRIND_V.pdf\#nameddest=tag-02C9}{V.{E.Q7} (p.~58)} }
\expandafter\def\csname SRef@maps:eq:E.Q8\endcsname{\href{S6_SERIES_GRIND_V.pdf\#nameddest=tag-02CA}{V.{E.Q8} (p.~58)} }
\expandafter\def\csname SRef@maps:eq:E.Q9\endcsname{\href{S6_SERIES_GRIND_V.pdf\#nameddest=tag-02CB}{V.{E.Q9} (p.~59)} }
\expandafter\def\csname SRef@maps:eq:E.Q10\endcsname{\href{S6_SERIES_GRIND_V.pdf\#nameddest=tag-02C1}{V.{E.Q10} (p.~59)} }
\expandafter\def\csname SRef@maps:eq:D.S9\endcsname{\href{S6_SERIES_GRIND_III.pdf\#nameddest=tag-023Y}{III.{D.S9} (p.~68)} }
\expandafter\def\csname SRef@maps:eq:E.Q11\endcsname{\href{S6_SERIES_GRIND_V.pdf\#nameddest=tag-02C2}{V.{E.Q11} (p.~59)} }
\expandafter\def\csname SRef@maps:eq:E.Q12\endcsname{\href{S6_SERIES_GRIND_V.pdf\#nameddest=tag-02C3}{V.{E.Q12} (p.~60)} }
\expandafter\def\csname SRef@maps:eq:E.T1\endcsname{\href{S6_SERIES_GRIND_V.pdf\#nameddest=tag-02DV}{V.{E.T1} (p.~60)} }
\expandafter\def\csname SRef@maps:eq:E.T2\endcsname{\href{S6_SERIES_GRIND_V.pdf\#nameddest=tag-02DW}{V.{E.T2} (p.~60)} }
\expandafter\def\csname SRef@maps:eq:E.T3\endcsname{\href{S6_SERIES_GRIND_V.pdf\#nameddest=tag-02DX}{V.{E.T3} (p.~61)} }
\expandafter\def\csname SRef@maps:eq:E.T4\endcsname{\href{S6_SERIES_GRIND_V.pdf\#nameddest=tag-02DY}{V.{E.T4} (p.~61)} }
\expandafter\def\csname SRef@maps:eq:E.T5\endcsname{\href{S6_SERIES_GRIND_V.pdf\#nameddest=tag-02DZ}{V.{E.T5} (p.~61)} }
\expandafter\def\csname SRef@maps:eq:E.T6\endcsname{\href{S6_SERIES_GRIND_V.pdf\#nameddest=tag-02E0}{V.{E.T6} (p.~62)} }
\expandafter\def\csname SRef@maps:eq:E.T7\endcsname{\href{S6_SERIES_GRIND_V.pdf\#nameddest=tag-02E1}{V.{E.T7} (p.~62)} }
\expandafter\def\csname SRef@maps:eq:E.T9\endcsname{\href{S6_SERIES_GRIND_V.pdf\#nameddest=tag-02E3}{V.{E.T9} (p.~62)} }
\expandafter\def\csname SRef@maps:eq:D.S10\endcsname{\href{S6_SERIES_GRIND_III.pdf\#nameddest=tag-023M}{III.{D.S10} (p.~68)} }
\expandafter\def\csname SRef@maps:eq:E.T8\endcsname{\href{S6_SERIES_GRIND_V.pdf\#nameddest=tag-02E2}{V.{E.T8} (p.~62)} }
\expandafter\def\csname SRef@maps:eq:E.U1\endcsname{\href{S6_SERIES_GRIND_V.pdf\#nameddest=tag-02E4}{V.{E.U1} (p.~63)} }
\expandafter\def\csname SRef@maps:eq:E.U2\endcsname{\href{S6_SERIES_GRIND_V.pdf\#nameddest=tag-02E7}{V.{E.U2} (p.~63)} }
\expandafter\def\csname SRef@maps:eq:E.U3\endcsname{\href{S6_SERIES_GRIND_V.pdf\#nameddest=tag-02E8}{V.{E.U3} (p.~63)} }
\expandafter\def\csname SRef@maps:eq:E.U4\endcsname{\href{S6_SERIES_GRIND_V.pdf\#nameddest=tag-02E9}{V.{E.U4} (p.~63)} }
\expandafter\def\csname SRef@maps:eq:E.U5\endcsname{\href{S6_SERIES_GRIND_V.pdf\#nameddest=tag-02EA}{V.{E.U5} (p.~63)} }
\expandafter\def\csname SRef@maps:eq:E.U6\endcsname{\href{S6_SERIES_GRIND_V.pdf\#nameddest=tag-02EB}{V.{E.U6} (p.~64)} }
\expandafter\def\csname SRef@maps:eq:E.U7\endcsname{\href{S6_SERIES_GRIND_V.pdf\#nameddest=tag-02EC}{V.{E.U7} (p.~64)} }
\expandafter\def\csname SRef@maps:eq:E.U8\endcsname{\href{S6_SERIES_GRIND_V.pdf\#nameddest=tag-02ED}{V.{E.U8} (p.~64)} }
\expandafter\def\csname SRef@maps:eq:E.U9\endcsname{\href{S6_SERIES_GRIND_V.pdf\#nameddest=tag-02EE}{V.{E.U9} (p.~64)} }
\expandafter\def\csname SRef@maps:eq:D.S11\endcsname{\href{S6_SERIES_GRIND_III.pdf\#nameddest=tag-023N}{III.{D.S11} (p.~69)} }
\expandafter\def\csname SRef@maps:eq:E.U10\endcsname{\href{S6_SERIES_GRIND_V.pdf\#nameddest=tag-02E5}{V.{E.U10} (p.~64)} }
\expandafter\def\csname SRef@maps:eq:E.U11\endcsname{\href{S6_SERIES_GRIND_V.pdf\#nameddest=tag-02E6}{V.{E.U11} (p.~64)} }
\expandafter\def\csname SRef@maps:eq:E.H0\endcsname{\href{S6_SERIES_GRIND_V.pdf\#nameddest=tag-02AQ}{V.{E.H0} (p.~65)} }
\expandafter\def\csname SRef@maps:eq:E.H1\endcsname{\href{S6_SERIES_GRIND_V.pdf\#nameddest=tag-02AR}{V.{E.H1} (p.~65)} }
\expandafter\def\csname SRef@maps:eq:E.H2\endcsname{\href{S6_SERIES_GRIND_V.pdf\#nameddest=tag-02AW}{V.{E.H2} (p.~66)} }
\expandafter\def\csname SRef@maps:eq:E.H3\endcsname{\href{S6_SERIES_GRIND_V.pdf\#nameddest=tag-02AX}{V.{E.H3} (p.~66)} }
\expandafter\def\csname SRef@maps:eq:E.H4\endcsname{\href{S6_SERIES_GRIND_V.pdf\#nameddest=tag-02AY}{V.{E.H4} (p.~66)} }
\expandafter\def\csname SRef@maps:eq:E.H5\endcsname{\href{S6_SERIES_GRIND_V.pdf\#nameddest=tag-02AZ}{V.{E.H5} (p.~66)} }
\expandafter\def\csname SRef@maps:eq:E.H6\endcsname{\href{S6_SERIES_GRIND_V.pdf\#nameddest=tag-02B0}{V.{E.H6} (p.~66)} }
\expandafter\def\csname SRef@maps:eq:E.H7\endcsname{\href{S6_SERIES_GRIND_V.pdf\#nameddest=tag-02B1}{V.{E.H7} (p.~66)} }
\expandafter\def\csname SRef@maps:eq:D.S12\endcsname{\href{S6_SERIES_GRIND_III.pdf\#nameddest=tag-023O}{III.{D.S12} (p.~69)} }
\expandafter\def\csname SRef@maps:eq:E.H8\endcsname{\href{S6_SERIES_GRIND_V.pdf\#nameddest=tag-02B2}{V.{E.H8} (p.~67)} }
\expandafter\def\csname SRef@maps:eq:E.H9\endcsname{\href{S6_SERIES_GRIND_V.pdf\#nameddest=tag-02B3}{V.{E.H9} (p.~67)} }
\expandafter\def\csname SRef@maps:eq:E.H9a\endcsname{\href{S6_SERIES_GRIND_V.pdf\#nameddest=tag-02B4}{V.{E.H9a} (p.~67)} }
\expandafter\def\csname SRef@maps:eq:E.H10\endcsname{\href{S6_SERIES_GRIND_V.pdf\#nameddest=tag-02AS}{V.{E.H10} (p.~68)} }
\expandafter\def\csname SRef@maps:eq:E.H11\endcsname{\href{S6_SERIES_GRIND_V.pdf\#nameddest=tag-02AT}{V.{E.H11} (p.~68)} }
\expandafter\def\csname SRef@maps:eq:E.H12\endcsname{\href{S6_SERIES_GRIND_V.pdf\#nameddest=tag-02AU}{V.{E.H12} (p.~68)} }
\expandafter\def\csname SRef@maps:eq:E.H13\endcsname{\href{S6_SERIES_GRIND_V.pdf\#nameddest=tag-02AV}{V.{E.H13} (p.~69)} }
\expandafter\def\csname SRef@maps:eq:D.S13\endcsname{\href{S6_SERIES_GRIND_III.pdf\#nameddest=tag-023P}{III.{D.S13} (p.~70)} }
\expandafter\def\csname SRef@maps:eq:D.S14\endcsname{\href{S6_SERIES_GRIND_III.pdf\#nameddest=tag-023Q}{III.{D.S14} (p.~70)} }
\expandafter\def\csname SRef@maps:eq:D.I1\endcsname{\href{S6_SERIES_GRIND_III.pdf\#nameddest=tag-01ZU}{III.{D.I1} (p.~71)} }
\expandafter\def\csname SRef@maps:eq:D.I2\endcsname{\href{S6_SERIES_GRIND_III.pdf\#nameddest=tag-01ZV}{III.{D.I2} (p.~71)} }
\expandafter\def\csname SRef@maps:eq:E.15\endcsname{\href{S6_SERIES_GRIND_III.pdf\#nameddest=tag-029L}{III.{E.15} (p.~317)} }
\expandafter\def\csname SRef@maps:eq:D.I3\endcsname{\href{S6_SERIES_GRIND_III.pdf\#nameddest=tag-01ZW}{III.{D.I3} (p.~71)} }
\expandafter\def\csname SRef@maps:eq:D.I4\endcsname{\href{S6_SERIES_GRIND_III.pdf\#nameddest=tag-01ZX}{III.{D.I4} (p.~72)} }
\expandafter\def\csname SRef@maps:eq:D.T1\endcsname{\href{S6_SERIES_GRIND_III.pdf\#nameddest=tag-0258}{III.{D.T1} (p.~73)} }
\expandafter\def\csname SRef@maps:eq:D.T2\endcsname{\href{S6_SERIES_GRIND_III.pdf\#nameddest=tag-025A}{III.{D.T2} (p.~74)} }
\expandafter\def\csname SRef@maps:eq:D.T3\endcsname{\href{S6_SERIES_GRIND_III.pdf\#nameddest=tag-025B}{III.{D.T3} (p.~74)} }
\expandafter\def\csname SRef@maps:eq:D.T4\endcsname{\href{S6_SERIES_GRIND_III.pdf\#nameddest=tag-025L}{III.{D.T4} (p.~74)} }
\expandafter\def\csname SRef@maps:eq:D.T5\endcsname{\href{S6_SERIES_GRIND_III.pdf\#nameddest=tag-025M}{III.{D.T5} (p.~75)} }
\expandafter\def\csname SRef@maps:eq:D.T6\endcsname{\href{S6_SERIES_GRIND_III.pdf\#nameddest=tag-025N}{III.{D.T6} (p.~75)} }
\expandafter\def\csname SRef@maps:eq:D.T7\endcsname{\href{S6_SERIES_GRIND_III.pdf\#nameddest=tag-025O}{III.{D.T7} (p.~75)} }
\expandafter\def\csname SRef@maps:eq:D.T8\endcsname{\href{S6_SERIES_GRIND_III.pdf\#nameddest=tag-025P}{III.{D.T8} (p.~75)} }
\expandafter\def\csname SRef@maps:eq:D.T9\endcsname{\href{S6_SERIES_GRIND_III.pdf\#nameddest=tag-025Q}{III.{D.T9} (p.~76)} }
\expandafter\def\csname SRef@maps:eq:D.T10\endcsname{\href{S6_SERIES_GRIND_III.pdf\#nameddest=tag-0259}{III.{D.T10} (p.~76)} }
\expandafter\def\csname SRef@maps:eq:D.V1\endcsname{\href{S6_SERIES_GRIND_III.pdf\#nameddest=tag-0269}{III.{D.V1} (p.~78)} }
\expandafter\def\csname SRef@maps:eq:D.V2\endcsname{\href{S6_SERIES_GRIND_III.pdf\#nameddest=tag-026A}{III.{D.V2} (p.~78)} }
\expandafter\def\csname SRef@maps:eq:D.V3\endcsname{\href{S6_SERIES_GRIND_III.pdf\#nameddest=tag-026B}{III.{D.V3} (p.~78)} }
\expandafter\def\csname SRef@maps:eq:D.V4\endcsname{\href{S6_SERIES_GRIND_III.pdf\#nameddest=tag-026C}{III.{D.V4} (p.~79)} }
\expandafter\def\csname SRef@maps:eq:D.V5\endcsname{\href{S6_SERIES_GRIND_III.pdf\#nameddest=tag-026D}{III.{D.V5} (p.~79)} }
\expandafter\def\csname SRef@maps:eq:D.F1\endcsname{\href{S6_SERIES_GRIND_III.pdf\#nameddest=tag-01YO}{III.{D.F1} (p.~81)} }
\expandafter\def\csname SRef@maps:eq:D.F2\endcsname{\href{S6_SERIES_GRIND_III.pdf\#nameddest=tag-01YP}{III.{D.F2} (p.~81)} }
\expandafter\def\csname SRef@maps:eq:D.O1\endcsname{\href{S6_SERIES_GRIND_III.pdf\#nameddest=tag-020M}{III.{D.O1} (p.~82)} }
\expandafter\def\csname SRef@maps:eq:D.O2\endcsname{\href{S6_SERIES_GRIND_III.pdf\#nameddest=tag-020N}{III.{D.O2} (p.~82)} }
\expandafter\def\csname SRef@maps:eq:D.O3\endcsname{\href{S6_SERIES_GRIND_III.pdf\#nameddest=tag-020O}{III.{D.O3} (p.~82)} }
\expandafter\def\csname SRef@maps:eq:D.O4\endcsname{\href{S6_SERIES_GRIND_III.pdf\#nameddest=tag-020P}{III.{D.O4} (p.~83)} }
\expandafter\def\csname SRef@maps:eq:D.O5\endcsname{\href{S6_SERIES_GRIND_III.pdf\#nameddest=tag-020Q}{III.{D.O5} (p.~83)} }
\expandafter\def\csname SRef@maps:eq:D.O6\endcsname{\href{S6_SERIES_GRIND_III.pdf\#nameddest=tag-020R}{III.{D.O6} (p.~83)} }
\expandafter\def\csname SRef@maps:eq:D.O7\endcsname{\href{S6_SERIES_GRIND_III.pdf\#nameddest=tag-020S}{III.{D.O7} (p.~84)} }
\expandafter\def\csname SRef@maps:eq:D.O8\endcsname{\href{S6_SERIES_GRIND_III.pdf\#nameddest=tag-020T}{III.{D.O8} (p.~84)} }
\expandafter\def\csname SRef@maps:src:D.2.2\endcsname{\href{S6_SERIES_GRIND_III.pdf\#nameddest=tag-02KU}{III.5.12 (p.~42)} }
\expandafter\def\csname SRef@maps:src:D.2.17\endcsname{\href{S6_SERIES_GRIND_III.pdf\#nameddest=tag-02KR}{III.5.26 (p.~77)} }
\expandafter\def\csname SRef@maps:src:D.2.28\endcsname{\href{S6_SERIES_GRIND_III.pdf\#nameddest=tag-02L3}{III.5.37 (p.~106)} }
\expandafter\def\csname SRef@wb:cusp-cone-S12-sphere-map-null\endcsname{\href{S6_SERIES_GRIND_III.pdf\#nameddest=tag-035Z}{III.7.14 (p.~270)} }
\expandafter\def\csname SRef@wb:cusp-cone-S24-sphere-map-null\endcsname{\href{S6_SERIES_GRIND_III.pdf\#nameddest=tag-036E}{III.7.15 (p.~270)} }
\expandafter\def\csname SRef@wb:cusp-cone-literal-R4-open-cells\endcsname{\href{S6_SERIES_GRIND_III.pdf\#nameddest=tag-036V}{III.7.16 (p.~271)} }
\expandafter\def\csname SRef@wb:cusp-cone-sharp-global-R4-scope\endcsname{\href{S6_SERIES_GRIND_III.pdf\#nameddest=tag-037J}{III.7.20 (p.~278)} }
\expandafter\def\csname SRef@wb:cusp-s12-all-higher-smashes-null\endcsname{\href{S6_SERIES_GRIND_III.pdf\#nameddest=tag-038T}{III.7.22 (p.~280)} }
\expandafter\def\csname SRef@wb:cusp_s12_carrier_homotopy:L567\endcsname{\href{S6_SERIES_GRIND_III.pdf\#nameddest=tag-03A6}{III.7.24 (p.~286)} }
\expandafter\def\csname SRef@wb:sphere-metric-unit-volume\endcsname{\href{S6_SERIES_GRIND_III.pdf\#nameddest=tag-03DH}{III.7.28 (p.~294)} }
\expandafter\def\csname SRef@wb:h24:ret:q2-detected-quotient\endcsname{\href{S6_SERIES_GRIND_III.pdf\#nameddest=tag-03BE}{III.7.33 (p.~306)} }
\expandafter\def\csname SRef@wb:cusp-cone-target-slice-embeddings\endcsname{\href{S6_SERIES_GRIND_III.pdf\#nameddest=tag-037N}{III.7.9 (p.~261)} }
\expandafter\def\csname SRef@wb:cusp-cover-target-action\endcsname{\href{S6_SERIES_GRIND_III.pdf\#nameddest=tag-0388}{III.1 (p.~238)} }
\expandafter\def\csname SRef@wb:cusp-cover-normalization-diagram\endcsname{\href{S6_SERIES_GRIND_III.pdf\#nameddest=tag-0383}{III.10 (p.~239)} }
\expandafter\def\csname SRef@wb:cusp-cone-products-not-spheres\endcsname{\href{S6_SERIES_GRIND_III.pdf\#nameddest=tag-0378}{III.100 (p.~263)} }
\expandafter\def\csname SRef@wb:cusp-cone-join-dimensions\endcsname{\href{S6_SERIES_GRIND_III.pdf\#nameddest=tag-036T}{III.101 (p.~263)} }
\expandafter\def\csname SRef@wb:cusp-cone-kernel-contained\endcsname{\href{S6_SERIES_GRIND_III.pdf\#nameddest=tag-036U}{III.102 (p.~263)} }
\expandafter\def\csname SRef@wb:cusp-cone-finite-index-orbit\endcsname{\href{S6_SERIES_GRIND_III.pdf\#nameddest=tag-036Q}{III.103 (p.~263)} }
\expandafter\def\csname SRef@wb:cusp-cone-fibre-H4-nonzero\endcsname{\href{S6_SERIES_GRIND_III.pdf\#nameddest=tag-036P}{III.104 (p.~264)} }
\expandafter\def\csname SRef@wb:cusp-cone-sphere-to-sphere-maps\endcsname{\href{S6_SERIES_GRIND_III.pdf\#nameddest=tag-037L}{III.105 (p.~266)} }
\expandafter\def\csname SRef@wb:cusp-cone-S12-S4-radial-map\endcsname{\href{S6_SERIES_GRIND_III.pdf\#nameddest=tag-035R}{III.106 (p.~267)} }
\expandafter\def\csname SRef@wb:cusp-cone-S24-S4-radial-map\endcsname{\href{S6_SERIES_GRIND_III.pdf\#nameddest=tag-0361}{III.107 (p.~267)} }
\expandafter\def\csname SRef@wb:cusp-cone-S12-exact-collapsed-locus\endcsname{\href{S6_SERIES_GRIND_III.pdf\#nameddest=tag-035T}{III.108 (p.~267)} }
\expandafter\def\csname SRef@wb:cusp-cone-S24-exact-collapsed-locus\endcsname{\href{S6_SERIES_GRIND_III.pdf\#nameddest=tag-0362}{III.109 (p.~267)} }
\expandafter\def\csname SRef@wb:cusp-cover-monodromy-conjugacy\endcsname{\href{S6_SERIES_GRIND_III.pdf\#nameddest=tag-0381}{III.11 (p.~239)} }
\expandafter\def\csname SRef@wb:cusp-cone-moving-S4-R4-limits\endcsname{\href{S6_SERIES_GRIND_III.pdf\#nameddest=tag-036Y}{III.110 (p.~267)} }
\expandafter\def\csname SRef@wb:cusp-cone-transported-skeleton-collapse\endcsname{\href{S6_SERIES_GRIND_III.pdf\#nameddest=tag-037O}{III.111 (p.~268)} }
\expandafter\def\csname SRef@wb:cusp-cone-distance-to-skeleton\endcsname{\href{S6_SERIES_GRIND_III.pdf\#nameddest=tag-036M}{III.112 (p.~268)} }
\expandafter\def\csname SRef@wb:cusp-cone-S12-onto-exact-ball\endcsname{\href{S6_SERIES_GRIND_III.pdf\#nameddest=tag-035W}{III.113 (p.~269)} }
\expandafter\def\csname SRef@wb:cusp-cone-S24-onto-exact-ball\endcsname{\href{S6_SERIES_GRIND_III.pdf\#nameddest=tag-0369}{III.114 (p.~270)} }
\expandafter\def\csname SRef@wb:cusp-cone-S24-null-maps\endcsname{\href{S6_SERIES_GRIND_III.pdf\#nameddest=tag-0367}{III.115 (p.~270)} }
\expandafter\def\csname SRef@wb:cusp-cone-S24-null-special-cases\endcsname{\href{S6_SERIES_GRIND_III.pdf\#nameddest=tag-0368}{III.116 (p.~270)} }
\expandafter\def\csname SRef@wb:cusp-cone-S24-smash-square-factorization\endcsname{\href{S6_SERIES_GRIND_III.pdf\#nameddest=tag-036C}{III.117 (p.~271)} }
\expandafter\def\csname SRef@wb:cusp-cone-S24-sixfold-suspension\endcsname{\href{S6_SERIES_GRIND_III.pdf\#nameddest=tag-036B}{III.118 (p.~271)} }
\expandafter\def\csname SRef@wb:cusp-cone-fourth-stable-stem\endcsname{\href{S6_SERIES_GRIND_III.pdf\#nameddest=tag-036R}{III.119 (p.~271)} }
\expandafter\def\csname SRef@wb:cusp-cover-maximal-chart-map\endcsname{\href{S6_SERIES_GRIND_III.pdf\#nameddest=tag-0380}{III.12 (p.~241)} }
\expandafter\def\csname SRef@wb:cusp-cone-U12-homeomorphism\endcsname{\href{S6_SERIES_GRIND_III.pdf\#nameddest=tag-036G}{III.121 (p.~271)} }
\expandafter\def\csname SRef@wb:cusp-cone-U24-homeomorphism\endcsname{\href{S6_SERIES_GRIND_III.pdf\#nameddest=tag-036H}{III.122 (p.~271)} }
\expandafter\def\csname SRef@wb:cusp-cone-D12-embedding\endcsname{\href{S6_SERIES_GRIND_III.pdf\#nameddest=tag-035G}{III.123 (p.~273)} }
\expandafter\def\csname SRef@wb:cusp-cone-D24-embedding\endcsname{\href{S6_SERIES_GRIND_III.pdf\#nameddest=tag-035I}{III.124 (p.~273)} }
\expandafter\def\csname SRef@wb:cusp-cone-D12-slice\endcsname{\href{S6_SERIES_GRIND_III.pdf\#nameddest=tag-035H}{III.125 (p.~273)} }
\expandafter\def\csname SRef@wb:cusp-cone-D24-slice\endcsname{\href{S6_SERIES_GRIND_III.pdf\#nameddest=tag-035J}{III.126 (p.~273)} }
\expandafter\def\csname SRef@wb:cusp-cone-carrier-local-homology\endcsname{\href{S6_SERIES_GRIND_III.pdf\#nameddest=tag-036K}{III.127 (p.~273)} }
\expandafter\def\csname SRef@wb:cusp-cone-A-smash-wedge\endcsname{\href{S6_SERIES_GRIND_III.pdf\#nameddest=tag-0357}{III.128 (p.~274)} }
\expandafter\def\csname SRef@wb:cusp-cone-C-smash-wedge\endcsname{\href{S6_SERIES_GRIND_III.pdf\#nameddest=tag-035F}{III.129 (p.~274)} }
\expandafter\def\csname SRef@wb:cusp-cone-E12-fibres\endcsname{\href{S6_SERIES_GRIND_III.pdf\#nameddest=tag-035K}{III.130 (p.~276)} }
\expandafter\def\csname SRef@wb:cusp-cone-E24-fibres\endcsname{\href{S6_SERIES_GRIND_III.pdf\#nameddest=tag-035L}{III.131 (p.~276)} }
\expandafter\def\csname SRef@wb:cusp-s12-smash-factorization\endcsname{\href{S6_SERIES_GRIND_III.pdf\#nameddest=tag-0390}{III.132 (p.~280)} }
\expandafter\def\csname SRef@wb:cusp-s12-toda-smash-formula\endcsname{\href{S6_SERIES_GRIND_III.pdf\#nameddest=tag-0392}{III.133 (p.~280)} }
\expandafter\def\csname SRef@wb:cusp-s12-three-stem-groups\endcsname{\href{S6_SERIES_GRIND_III.pdf\#nameddest=tag-0391}{III.134 (p.~280)} }
\expandafter\def\csname SRef@wb:cusp-s12-sixfold-suspension-class\endcsname{\href{S6_SERIES_GRIND_III.pdf\#nameddest=tag-038Z}{III.135 (p.~281)} }
\expandafter\def\csname SRef@wb:cusp-s12-regular-inverse-image\endcsname{\href{S6_SERIES_GRIND_III.pdf\#nameddest=tag-038Y}{III.136 (p.~281)} }
\expandafter\def\csname SRef@wb:cusp-null:eq:periods\endcsname{\href{S6_SERIES_GRIND_III.pdf\#nameddest=tag-038N}{III.137 (p.~282)} }
\expandafter\def\csname SRef@wb:cusp-null:eq:action\endcsname{\href{S6_SERIES_GRIND_III.pdf\#nameddest=tag-038E}{III.138 (p.~282)} }
\expandafter\def\csname SRef@wb:cusp-null:eq:g\endcsname{\href{S6_SERIES_GRIND_III.pdf\#nameddest=tag-038J}{III.139 (p.~282)} }
\expandafter\def\csname SRef@wb:cusp-null:eq:chart\endcsname{\href{S6_SERIES_GRIND_III.pdf\#nameddest=tag-038H}{III.140 (p.~282)} }
\expandafter\def\csname SRef@wb:cusp-null:eq:moduli\endcsname{\href{S6_SERIES_GRIND_III.pdf\#nameddest=tag-038L}{III.141 (p.~282)} }
\expandafter\def\csname SRef@wb:cusp-null:eq:invertibility\endcsname{\href{S6_SERIES_GRIND_III.pdf\#nameddest=tag-038K}{III.142 (p.~283)} }
\expandafter\def\csname SRef@wb:cusp-null:eq:straightening\endcsname{\href{S6_SERIES_GRIND_III.pdf\#nameddest=tag-038Q}{III.144 (p.~283)} }
\expandafter\def\csname SRef@wb:cusp-null:eq:torus\endcsname{\href{S6_SERIES_GRIND_III.pdf\#nameddest=tag-038R}{III.145 (p.~283)} }
\expandafter\def\csname SRef@wb:cusp-null:eq:spatial-bound\endcsname{\href{S6_SERIES_GRIND_III.pdf\#nameddest=tag-038P}{III.146 (p.~283)} }
\expandafter\def\csname SRef@wb:cusp-null:eq:nearby-moduli\endcsname{\href{S6_SERIES_GRIND_III.pdf\#nameddest=tag-038M}{III.147 (p.~284)} }
\expandafter\def\csname SRef@wb:cusp-null:eq:J\endcsname{\href{S6_SERIES_GRIND_III.pdf\#nameddest=tag-038B}{III.148 (p.~284)} }
\expandafter\def\csname SRef@wb:cusp-null:eq:actual-frame\endcsname{\href{S6_SERIES_GRIND_III.pdf\#nameddest=tag-038F}{III.149 (p.~284)} }
\expandafter\def\csname SRef@wb:cusp-cover-equivariance-computation\endcsname{\href{S6_SERIES_GRIND_III.pdf\#nameddest=tag-037W}{III.15 (p.~241)} }
\expandafter\def\csname SRef@wb:cusp-null:eq:J0\endcsname{\href{S6_SERIES_GRIND_III.pdf\#nameddest=tag-038C}{III.150 (p.~284)} }
\expandafter\def\csname SRef@wb:cusp-null:eq:P\endcsname{\href{S6_SERIES_GRIND_III.pdf\#nameddest=tag-038D}{III.151 (p.~284)} }
\expandafter\def\csname SRef@wb:cusp-null:eq:q\endcsname{\href{S6_SERIES_GRIND_III.pdf\#nameddest=tag-038O}{III.152 (p.~285)} }
\expandafter\def\csname SRef@wb:cusp-null:eq:bordism\endcsname{\href{S6_SERIES_GRIND_III.pdf\#nameddest=tag-038G}{III.153 (p.~285)} }
\expandafter\def\csname SRef@wb:cusp-null:eq:extended-frame\endcsname{\href{S6_SERIES_GRIND_III.pdf\#nameddest=tag-038I}{III.154 (p.~285)} }
\expandafter\def\csname SRef@wb:sphere-metric-normalized-Q-limit\endcsname{\href{S6_SERIES_GRIND_III.pdf\#nameddest=tag-03DA}{III.155 (p.~287)} }
\expandafter\def\csname SRef@wb:sphere-metric-Hausdorff-collapse\endcsname{\href{S6_SERIES_GRIND_III.pdf\#nameddest=tag-03D2}{III.158 (p.~288)} }
\expandafter\def\csname SRef@wb:sphere-metric-A-fixed-metric\endcsname{\href{S6_SERIES_GRIND_III.pdf\#nameddest=tag-03CW}{III.159 (p.~288)} }
\expandafter\def\csname SRef@wb:cusp-cover-basechange-local-ring\endcsname{\href{S6_SERIES_GRIND_III.pdf\#nameddest=tag-037S}{III.16 (p.~242)} }
\expandafter\def\csname SRef@wb:sphere-metric-BC-fixed-metric\endcsname{\href{S6_SERIES_GRIND_III.pdf\#nameddest=tag-03CY}{III.160 (p.~288)} }
\expandafter\def\csname SRef@wb:sphere-metric-A-slice-exact\endcsname{\href{S6_SERIES_GRIND_III.pdf\#nameddest=tag-03CX}{III.161 (p.~289)} }
\expandafter\def\csname SRef@wb:sphere-metric-BC-slice-exact\endcsname{\href{S6_SERIES_GRIND_III.pdf\#nameddest=tag-03CZ}{III.162 (p.~289)} }
\expandafter\def\csname SRef@wb:sphere-metric-slice-fixed-blowups\endcsname{\href{S6_SERIES_GRIND_III.pdf\#nameddest=tag-03DE}{III.164 (p.~289)} }
\expandafter\def\csname SRef@wb:sphere-metric-ambient-fixed-blowups\endcsname{\href{S6_SERIES_GRIND_III.pdf\#nameddest=tag-03D3}{III.166 (p.~289)} }
\expandafter\def\csname SRef@wb:sphere-metric-scaling-trichotomy-result\endcsname{\href{S6_SERIES_GRIND_III.pdf\#nameddest=tag-03DD}{III.167 (p.~291)} }
\expandafter\def\csname SRef@wb:sphere-metric-vertex-distance\endcsname{\href{S6_SERIES_GRIND_III.pdf\#nameddest=tag-03DI}{III.168 (p.~292)} }
\expandafter\def\csname SRef@wb:sphere-metric-intrinsic-ambient-agreement\endcsname{\href{S6_SERIES_GRIND_III.pdf\#nameddest=tag-03D7}{III.169 (p.~292)} }
\expandafter\def\csname SRef@wb:cusp-cover-invariant-ring\endcsname{\href{S6_SERIES_GRIND_III.pdf\#nameddest=tag-037Z}{III.17 (p.~242)} }
\expandafter\def\csname SRef@wb:sphere-metric-intrinsic-ambient-ball-comparison\endcsname{\href{S6_SERIES_GRIND_III.pdf\#nameddest=tag-03D8}{III.170 (p.~292)} }
\expandafter\def\csname SRef@wb:sphere-metric-arbitrary-unrolling-ambient-bubbles\endcsname{\href{S6_SERIES_GRIND_III.pdf\#nameddest=tag-03D4}{III.172 (p.~292)} }
\expandafter\def\csname SRef@wb:sphere-metric-outward-shortcut-exclusion\endcsname{\href{S6_SERIES_GRIND_III.pdf\#nameddest=tag-03DB}{III.173 (p.~293)} }
\expandafter\def\csname SRef@wb:sphere-metric-C-smash-local-homology\endcsname{\href{S6_SERIES_GRIND_III.pdf\#nameddest=tag-03D1}{III.174 (p.~295)} }
\expandafter\def\csname SRef@wb:sphere-metric-C-boundary-quotient-homeomorphism\endcsname{\href{S6_SERIES_GRIND_III.pdf\#nameddest=tag-03D0}{III.175 (p.~296)} }
\expandafter\def\csname SRef@wb:sphere-metric-boundary-link-local-shift\endcsname{\href{S6_SERIES_GRIND_III.pdf\#nameddest=tag-03D5}{III.176 (p.~296)} }
\expandafter\def\csname SRef@wb:h24:ret:eq:original-periods\endcsname{\href{S6_SERIES_GRIND_III.pdf\#nameddest=tag-03B4}{III.177 (p.~298)} }
\expandafter\def\csname SRef@wb:h24:ret:eq:A\endcsname{\href{S6_SERIES_GRIND_III.pdf\#nameddest=tag-03AF}{III.178 (p.~298)} }
\expandafter\def\csname SRef@wb:h24:ret:eq:deck\endcsname{\href{S6_SERIES_GRIND_III.pdf\#nameddest=tag-03AS}{III.179 (p.~298)} }
\expandafter\def\csname SRef@wb:h24:ret:eq:logx\endcsname{\href{S6_SERIES_GRIND_III.pdf\#nameddest=tag-03B0}{III.180 (p.~298)} }
\expandafter\def\csname SRef@wb:h24:ret:eq:invertibility\endcsname{\href{S6_SERIES_GRIND_III.pdf\#nameddest=tag-03AZ}{III.181 (p.~298)} }
\expandafter\def\csname SRef@wb:h24:ret:eq:native\endcsname{\href{S6_SERIES_GRIND_III.pdf\#nameddest=tag-03B1}{III.182 (p.~298)} }
\expandafter\def\csname SRef@wb:h24:ret:eq:two-skeleton\endcsname{\href{S6_SERIES_GRIND_III.pdf\#nameddest=tag-03B9}{III.183 (p.~298)} }
\expandafter\def\csname SRef@wb:h24:ret:eq:A-target\endcsname{\href{S6_SERIES_GRIND_III.pdf\#nameddest=tag-03AH}{III.184 (p.~299)} }
\expandafter\def\csname SRef@wb:h24:ret:eq:q2\endcsname{\href{S6_SERIES_GRIND_III.pdf\#nameddest=tag-03B5}{III.185 (p.~299)} }
\expandafter\def\csname SRef@wb:h24:ret:eq:H\endcsname{\href{S6_SERIES_GRIND_III.pdf\#nameddest=tag-03AN}{III.186 (p.~299)} }
\expandafter\def\csname SRef@wb:h24:ret:eq:four-cell\endcsname{\href{S6_SERIES_GRIND_III.pdf\#nameddest=tag-03AW}{III.187 (p.~299)} }
\expandafter\def\csname SRef@wb:h24:ret:eq:chosen-value\endcsname{\href{S6_SERIES_GRIND_III.pdf\#nameddest=tag-03AR}{III.188 (p.~299)} }
\expandafter\def\csname SRef@wb:h24:ret:eq:z-chart\endcsname{\href{S6_SERIES_GRIND_III.pdf\#nameddest=tag-03BB}{III.189 (p.~299)} }
\expandafter\def\csname SRef@wb:cusp-cover-triple-fibre-length\endcsname{\href{S6_SERIES_GRIND_III.pdf\#nameddest=tag-038A}{III.19 (p.~242)} }
\expandafter\def\csname SRef@wb:h24:ret:eq:units\endcsname{\href{S6_SERIES_GRIND_III.pdf\#nameddest=tag-03BA}{III.190 (p.~299)} }
\expandafter\def\csname SRef@wb:h24:ret:eq:F\endcsname{\href{S6_SERIES_GRIND_III.pdf\#nameddest=tag-03AL}{III.191 (p.~299)} }
\expandafter\def\csname SRef@wb:h24:ret:eq:Finv\endcsname{\href{S6_SERIES_GRIND_III.pdf\#nameddest=tag-03AM}{III.192 (p.~299)} }
\expandafter\def\csname SRef@wb:h24:ret:eq:M\endcsname{\href{S6_SERIES_GRIND_III.pdf\#nameddest=tag-03AP}{III.193 (p.~300)} }
\expandafter\def\csname SRef@wb:h24:ret:eq:exact-u\endcsname{\href{S6_SERIES_GRIND_III.pdf\#nameddest=tag-03AV}{III.194 (p.~300)} }
\expandafter\def\csname SRef@wb:h24:ret:eq:full-derivative\endcsname{\href{S6_SERIES_GRIND_III.pdf\#nameddest=tag-03AX}{III.195 (p.~300)} }
\expandafter\def\csname SRef@wb:h24:ret:eq:derivative-on-M\endcsname{\href{S6_SERIES_GRIND_III.pdf\#nameddest=tag-03AT}{III.196 (p.~300)} }
\expandafter\def\csname SRef@wb:h24:ret:eq:LQ\endcsname{\href{S6_SERIES_GRIND_III.pdf\#nameddest=tag-03AO}{III.197 (p.~301)} }
\expandafter\def\csname SRef@wb:h24:ret:eq:tangents\endcsname{\href{S6_SERIES_GRIND_III.pdf\#nameddest=tag-03B8}{III.198 (p.~301)} }
\expandafter\def\csname SRef@wb:h24:ret:eq:normal-first\endcsname{\href{S6_SERIES_GRIND_III.pdf\#nameddest=tag-03B2}{III.199 (p.~301)} }
\expandafter\def\csname SRef@wb:cusp-cover-base-change\endcsname{\href{S6_SERIES_GRIND_III.pdf\#nameddest=tag-037Q}{III.2 (p.~238)} }
\expandafter\def\csname SRef@wb:cusp-cover-double-fibre-length\endcsname{\href{S6_SERIES_GRIND_III.pdf\#nameddest=tag-037V}{III.20 (p.~242)} }
\expandafter\def\csname SRef@wb:h24:ret:eq:normal-u\endcsname{\href{S6_SERIES_GRIND_III.pdf\#nameddest=tag-03B3}{III.200 (p.~301)} }
\expandafter\def\csname SRef@wb:h24:ret:eq:differential-table\endcsname{\href{S6_SERIES_GRIND_III.pdf\#nameddest=tag-03AU}{III.201 (p.~301)} }
\expandafter\def\csname SRef@wb:h24:ret:eq:A-homotopy\endcsname{\href{S6_SERIES_GRIND_III.pdf\#nameddest=tag-03AG}{III.202 (p.~301)} }
\expandafter\def\csname SRef@wb:h24:ret:eq:sign\endcsname{\href{S6_SERIES_GRIND_III.pdf\#nameddest=tag-03B7}{III.203 (p.~302)} }
\expandafter\def\csname SRef@wb:h24:ret:eq:D\endcsname{\href{S6_SERIES_GRIND_III.pdf\#nameddest=tag-03AJ}{III.204 (p.~302)} }
\expandafter\def\csname SRef@wb:h24:ret:eq:full-frame\endcsname{\href{S6_SERIES_GRIND_III.pdf\#nameddest=tag-03AY}{III.205 (p.~302)} }
\expandafter\def\csname SRef@wb:h24:ret:eq:D-homotopy\endcsname{\href{S6_SERIES_GRIND_III.pdf\#nameddest=tag-03AK}{III.206 (p.~302)} }
\expandafter\def\csname SRef@wb:h24:ret:eq:quadratic\endcsname{\href{S6_SERIES_GRIND_III.pdf\#nameddest=tag-03B6}{III.207 (p.~303)} }
\expandafter\def\csname SRef@wb:h24:ret:eq:Arf\endcsname{\href{S6_SERIES_GRIND_III.pdf\#nameddest=tag-03AI}{III.208 (p.~303)} }
\expandafter\def\csname SRef@wb:h24:ret:eq:actual-Arf\endcsname{\href{S6_SERIES_GRIND_III.pdf\#nameddest=tag-03AQ}{III.209 (p.~304)} }
\expandafter\def\csname SRef@wb:cusp-cover-additive-fibre-map\endcsname{\href{S6_SERIES_GRIND_III.pdf\#nameddest=tag-037P}{III.21 (p.~243)} }
\expandafter\def\csname SRef@mapsD:repair:cocycle-relation\endcsname{\href{S6_SERIES_GRIND_III.pdf\#nameddest=tag-02OF}{III.210 (p.~310)} }
\expandafter\def\csname SRef@mapsD:repair:canonical-matrices\endcsname{\href{S6_SERIES_GRIND_III.pdf\#nameddest=tag-02OE}{III.211 (p.~312)} }
\expandafter\def\csname SRef@mapsD:repair:free-peripherals\endcsname{\href{S6_SERIES_GRIND_III.pdf\#nameddest=tag-02OH}{III.212 (p.~312)} }
\expandafter\def\csname SRef@mapsD:repair:50-weierstrass\endcsname{\href{S6_SERIES_GRIND_III.pdf\#nameddest=tag-02OD}{III.213 (p.~312)} }
\expandafter\def\csname SRef@mapsD:repair:height-formula\endcsname{\href{S6_SERIES_GRIND_III.pdf\#nameddest=tag-02OI}{III.214 (p.~315)} }
\expandafter\def\csname SRef@wb:cusp-spec-radial-notation\endcsname{\href{S6_SERIES_GRIND_III.pdf\#nameddest=tag-03A3}{III.22 (p.~244)} }
\expandafter\def\csname SRef@wb:cusp-spec-holomorphic-errors\endcsname{\href{S6_SERIES_GRIND_III.pdf\#nameddest=tag-039V}{III.23 (p.~244)} }
\expandafter\def\csname SRef@wb:cusp-spec-full-inverse-transpose\endcsname{\href{S6_SERIES_GRIND_III.pdf\#nameddest=tag-039O}{III.24 (p.~244)} }
\expandafter\def\csname SRef@wb:cusp-spec-global-circle-map\endcsname{\href{S6_SERIES_GRIND_III.pdf\#nameddest=tag-039R}{III.25 (p.~246)} }
\expandafter\def\csname SRef@wb:cusp-spec-global-circle-quotient\endcsname{\href{S6_SERIES_GRIND_III.pdf\#nameddest=tag-039S}{III.26 (p.~246)} }
\expandafter\def\csname SRef@wb:cusp-spec-global-circle-coordinate\endcsname{\href{S6_SERIES_GRIND_III.pdf\#nameddest=tag-039Q}{III.27 (p.~246)} }
\expandafter\def\csname SRef@wb:cusp-spec-global-kernel-family\endcsname{\href{S6_SERIES_GRIND_III.pdf\#nameddest=tag-039T}{III.28 (p.~246)} }
\expandafter\def\csname SRef@wb:cusp-spec-radial-circle-character\endcsname{\href{S6_SERIES_GRIND_III.pdf\#nameddest=tag-03A2}{III.29 (p.~247)} }
\expandafter\def\csname SRef@wb:cusp-cover-source-factor\endcsname{\href{S6_SERIES_GRIND_III.pdf\#nameddest=tag-0387}{III.3 (p.~238)} }
\expandafter\def\csname SRef@wb:cusp-spec-global-circle-circumference\endcsname{\href{S6_SERIES_GRIND_III.pdf\#nameddest=tag-039P}{III.30 (p.~247)} }
\expandafter\def\csname SRef@wb:cusp-spec-fourfold-fibre-product\endcsname{\href{S6_SERIES_GRIND_III.pdf\#nameddest=tag-039N}{III.31 (p.~247)} }
\expandafter\def\csname SRef@wb:cusp-spec-fourfold-circle-map\endcsname{\href{S6_SERIES_GRIND_III.pdf\#nameddest=tag-039M}{III.32 (p.~247)} }
\expandafter\def\csname SRef@wb:cusp-spec-four-circle-fibre\endcsname{\href{S6_SERIES_GRIND_III.pdf\#nameddest=tag-039J}{III.33 (p.~247)} }
\expandafter\def\csname SRef@wb:cusp-spec-four-circle-limit\endcsname{\href{S6_SERIES_GRIND_III.pdf\#nameddest=tag-039K}{III.34 (p.~247)} }
\expandafter\def\csname SRef@wb:cusp-spec-doubled-local-quotient\endcsname{\href{S6_SERIES_GRIND_III.pdf\#nameddest=tag-0398}{III.35 (p.~247)} }
\expandafter\def\csname SRef@wb:cusp-spec-finite-affine-action\endcsname{\href{S6_SERIES_GRIND_III.pdf\#nameddest=tag-039A}{III.36 (p.~248)} }
\expandafter\def\csname SRef@wb:cusp-spec-finite-character-shift\endcsname{\href{S6_SERIES_GRIND_III.pdf\#nameddest=tag-039B}{III.37 (p.~248)} }
\expandafter\def\csname SRef@wb:cusp-spec-finite-Seifert-target\endcsname{\href{S6_SERIES_GRIND_III.pdf\#nameddest=tag-0399}{III.38 (p.~248)} }
\expandafter\def\csname SRef@wb:cusp-spec-finite-fourfold-normalization\endcsname{\href{S6_SERIES_GRIND_III.pdf\#nameddest=tag-039D}{III.39 (p.~248)} }
\expandafter\def\csname SRef@wb:cusp-cover-source-action\endcsname{\href{S6_SERIES_GRIND_III.pdf\#nameddest=tag-0386}{III.4 (p.~238)} }
\expandafter\def\csname SRef@wb:cusp-spec-finite-fourfold-target\endcsname{\href{S6_SERIES_GRIND_III.pdf\#nameddest=tag-039E}{III.40 (p.~248)} }
\expandafter\def\csname SRef@wb:cusp-spec-finite-normalization-map\endcsname{\href{S6_SERIES_GRIND_III.pdf\#nameddest=tag-039F}{III.41 (p.~248)} }
\expandafter\def\csname SRef@wb:cusp-spec-product-not-sphere\endcsname{\href{S6_SERIES_GRIND_III.pdf\#nameddest=tag-03A1}{III.42 (p.~248)} }
\expandafter\def\csname SRef@wb:cusp-spec-four-product-not-sphere\endcsname{\href{S6_SERIES_GRIND_III.pdf\#nameddest=tag-039L}{III.43 (p.~248)} }
\expandafter\def\csname SRef@wb:cusp-spec-smash-spheres\endcsname{\href{S6_SERIES_GRIND_III.pdf\#nameddest=tag-03A4}{III.44 (p.~249)} }
\expandafter\def\csname SRef@wb:cusp-spec-target-smash-map\endcsname{\href{S6_SERIES_GRIND_III.pdf\#nameddest=tag-03A5}{III.45 (p.~249)} }
\expandafter\def\csname SRef@wb:cusp-spec-cusp-kernel\endcsname{\href{S6_SERIES_GRIND_III.pdf\#nameddest=tag-0394}{III.46 (p.~249)} }
\expandafter\def\csname SRef@wb:cusp-spec-invariant-flag-calculation\endcsname{\href{S6_SERIES_GRIND_III.pdf\#nameddest=tag-039X}{III.47 (p.~249)} }
\expandafter\def\csname SRef@wb:cusp-spec-invariant-character-lattice\endcsname{\href{S6_SERIES_GRIND_III.pdf\#nameddest=tag-039W}{III.48 (p.~249)} }
\expandafter\def\csname SRef@wb:cusp-spec-circle-dual-vector\endcsname{\href{S6_SERIES_GRIND_III.pdf\#nameddest=tag-0393}{III.49 (p.~250)} }
\expandafter\def\csname SRef@wb:cusp-cover-toric-power-map\endcsname{\href{S6_SERIES_GRIND_III.pdf\#nameddest=tag-0389}{III.5 (p.~239)} }
\expandafter\def\csname SRef@wb:cusp-spec-no-finite-cover-matrix\endcsname{\href{S6_SERIES_GRIND_III.pdf\#nameddest=tag-039Z}{III.50 (p.~251)} }
\expandafter\def\csname SRef@wb:cusp-spec-no-finite-cover-orbit\endcsname{\href{S6_SERIES_GRIND_III.pdf\#nameddest=tag-03A0}{III.51 (p.~251)} }
\expandafter\def\csname SRef@wb:cusp-spec-finite-target-conjugacy\endcsname{\href{S6_SERIES_GRIND_III.pdf\#nameddest=tag-039I}{III.52 (p.~251)} }
\expandafter\def\csname SRef@wb:cusp-spec-finite-target-cocycle\endcsname{\href{S6_SERIES_GRIND_III.pdf\#nameddest=tag-039H}{III.53 (p.~251)} }
\expandafter\def\csname SRef@wb:cusp-spec-finite-overlap-identity\endcsname{\href{S6_SERIES_GRIND_III.pdf\#nameddest=tag-039G}{III.54 (p.~251)} }
\expandafter\def\csname SRef@wb:cusp-spec-finite-fibre-product-branches\endcsname{\href{S6_SERIES_GRIND_III.pdf\#nameddest=tag-039C}{III.55 (p.~252)} }
\expandafter\def\csname SRef@wb:cusp-spec-cusp-nonextension-position\endcsname{\href{S6_SERIES_GRIND_III.pdf\#nameddest=tag-0397}{III.56 (p.~252)} }
\expandafter\def\csname SRef@wb:cusp-spec-cusp-nonextension-character\endcsname{\href{S6_SERIES_GRIND_III.pdf\#nameddest=tag-0395}{III.57 (p.~252)} }
\expandafter\def\csname SRef@wb:cusp-spec-cusp-nonextension-chart\endcsname{\href{S6_SERIES_GRIND_III.pdf\#nameddest=tag-0396}{III.58 (p.~253)} }
\expandafter\def\csname SRef@wb:cusp-cone-disc\endcsname{\href{S6_SERIES_GRIND_III.pdf\#nameddest=tag-036L}{III.59 (p.~255)} }
\expandafter\def\csname SRef@wb:cusp-cover-completed-map\endcsname{\href{S6_SERIES_GRIND_III.pdf\#nameddest=tag-037T}{III.6 (p.~239)} }
\expandafter\def\csname SRef@wb:cusp-cone-quotient-lattice\endcsname{\href{S6_SERIES_GRIND_III.pdf\#nameddest=tag-037F}{III.60 (p.~255)} }
\expandafter\def\csname SRef@wb:cusp-cone-native-family-map\endcsname{\href{S6_SERIES_GRIND_III.pdf\#nameddest=tag-0371}{III.61 (p.~255)} }
\expandafter\def\csname SRef@wb:cusp-cone-radial-isometry\endcsname{\href{S6_SERIES_GRIND_III.pdf\#nameddest=tag-037G}{III.62 (p.~255)} }
\expandafter\def\csname SRef@wb:cusp-cone-native-restriction\endcsname{\href{S6_SERIES_GRIND_III.pdf\#nameddest=tag-0372}{III.63 (p.~255)} }
\expandafter\def\csname SRef@wb:cusp-cone-native-systole\endcsname{\href{S6_SERIES_GRIND_III.pdf\#nameddest=tag-0373}{III.64 (p.~255)} }
\expandafter\def\csname SRef@wb:cusp-cone-monodromy-image\endcsname{\href{S6_SERIES_GRIND_III.pdf\#nameddest=tag-036X}{III.65 (p.~255)} }
\expandafter\def\csname SRef@wb:cusp-cone-A-definition\endcsname{\href{S6_SERIES_GRIND_III.pdf\#nameddest=tag-0354}{III.66 (p.~256)} }
\expandafter\def\csname SRef@wb:cusp-cone-A-suspension\endcsname{\href{S6_SERIES_GRIND_III.pdf\#nameddest=tag-0358}{III.67 (p.~256)} }
\expandafter\def\csname SRef@wb:cusp-cone-q2\endcsname{\href{S6_SERIES_GRIND_III.pdf\#nameddest=tag-037D}{III.68 (p.~256)} }
\expandafter\def\csname SRef@wb:cusp-cone-q2-formula\endcsname{\href{S6_SERIES_GRIND_III.pdf\#nameddest=tag-037E}{III.69 (p.~256)} }
\expandafter\def\csname SRef@wb:cusp-cover-direct-ramification\endcsname{\href{S6_SERIES_GRIND_III.pdf\#nameddest=tag-037U}{III.7 (p.~239)} }
\expandafter\def\csname SRef@wb:cusp-cone-S12-map\endcsname{\href{S6_SERIES_GRIND_III.pdf\#nameddest=tag-035V}{III.70 (p.~256)} }
\expandafter\def\csname SRef@wb:cusp-cone-S12-source\endcsname{\href{S6_SERIES_GRIND_III.pdf\#nameddest=tag-035Y}{III.71 (p.~256)} }
\expandafter\def\csname SRef@wb:cusp-cone-S12-fibre-map\endcsname{\href{S6_SERIES_GRIND_III.pdf\#nameddest=tag-035U}{III.72 (p.~256)} }
\expandafter\def\csname SRef@wb:cusp-cone-S12-R4-limit\endcsname{\href{S6_SERIES_GRIND_III.pdf\#nameddest=tag-035Q}{III.73 (p.~256)} }
\expandafter\def\csname SRef@wb:cusp-cone-A-homology\endcsname{\href{S6_SERIES_GRIND_III.pdf\#nameddest=tag-0355}{III.74 (p.~256)} }
\expandafter\def\csname SRef@wb:cusp-cone-A-smash-homology\endcsname{\href{S6_SERIES_GRIND_III.pdf\#nameddest=tag-0356}{III.75 (p.~257)} }
\expandafter\def\csname SRef@wb:cusp-cone-B-definition\endcsname{\href{S6_SERIES_GRIND_III.pdf\#nameddest=tag-0359}{III.76 (p.~258)} }
\expandafter\def\csname SRef@wb:cusp-cone-q1-local\endcsname{\href{S6_SERIES_GRIND_III.pdf\#nameddest=tag-037B}{III.77 (p.~258)} }
\expandafter\def\csname SRef@wb:cusp-cone-q1-local-formula\endcsname{\href{S6_SERIES_GRIND_III.pdf\#nameddest=tag-037C}{III.78 (p.~258)} }
\expandafter\def\csname SRef@wb:cusp-cone-S24-local-fibre-map\endcsname{\href{S6_SERIES_GRIND_III.pdf\#nameddest=tag-0364}{III.79 (p.~258)} }
\expandafter\def\csname SRef@wb:cusp-cover-base-changed-target\endcsname{\href{S6_SERIES_GRIND_III.pdf\#nameddest=tag-037R}{III.8 (p.~239)} }
\expandafter\def\csname SRef@wb:cusp-cone-B-homology\endcsname{\href{S6_SERIES_GRIND_III.pdf\#nameddest=tag-035A}{III.80 (p.~258)} }
\expandafter\def\csname SRef@wb:cusp-cone-B-smash-homology\endcsname{\href{S6_SERIES_GRIND_III.pdf\#nameddest=tag-035B}{III.81 (p.~258)} }
\expandafter\def\csname SRef@wb:cusp-cone-pair-of-pants\endcsname{\href{S6_SERIES_GRIND_III.pdf\#nameddest=tag-0377}{III.82 (p.~259)} }
\expandafter\def\csname SRef@wb:cusp-cone-M-definition\endcsname{\href{S6_SERIES_GRIND_III.pdf\#nameddest=tag-035P}{III.83 (p.~259)} }
\expandafter\def\csname SRef@wb:cusp-cone-M-boundary\endcsname{\href{S6_SERIES_GRIND_III.pdf\#nameddest=tag-035N}{III.84 (p.~259)} }
\expandafter\def\csname SRef@wb:cusp-cone-C-definition\endcsname{\href{S6_SERIES_GRIND_III.pdf\#nameddest=tag-035C}{III.85 (p.~259)} }
\expandafter\def\csname SRef@wb:cusp-cone-q1\endcsname{\href{S6_SERIES_GRIND_III.pdf\#nameddest=tag-0379}{III.86 (p.~259)} }
\expandafter\def\csname SRef@wb:cusp-cone-q1-cusp-formula\endcsname{\href{S6_SERIES_GRIND_III.pdf\#nameddest=tag-037A}{III.87 (p.~259)} }
\expandafter\def\csname SRef@wb:cusp-cone-S24-map\endcsname{\href{S6_SERIES_GRIND_III.pdf\#nameddest=tag-0366}{III.88 (p.~260)} }
\expandafter\def\csname SRef@wb:cusp-cone-S24-source\endcsname{\href{S6_SERIES_GRIND_III.pdf\#nameddest=tag-036D}{III.89 (p.~260)} }
\expandafter\def\csname SRef@wb:cusp-cover-relative-factorization\endcsname{\href{S6_SERIES_GRIND_III.pdf\#nameddest=tag-0384}{III.9 (p.~239)} }
\expandafter\def\csname SRef@wb:cusp-cone-S24-fibre-map\endcsname{\href{S6_SERIES_GRIND_III.pdf\#nameddest=tag-0363}{III.90 (p.~260)} }
\expandafter\def\csname SRef@wb:cusp-cone-S24-R4-limit\endcsname{\href{S6_SERIES_GRIND_III.pdf\#nameddest=tag-0360}{III.91 (p.~260)} }
\expandafter\def\csname SRef@wb:cusp-cone-C-homology\endcsname{\href{S6_SERIES_GRIND_III.pdf\#nameddest=tag-035D}{III.92 (p.~260)} }
\expandafter\def\csname SRef@wb:cusp-cone-C-smash-homology\endcsname{\href{S6_SERIES_GRIND_III.pdf\#nameddest=tag-035E}{III.93 (p.~260)} }
\expandafter\def\csname SRef@wb:cusp-cone-Seifert-meridians\endcsname{\href{S6_SERIES_GRIND_III.pdf\#nameddest=tag-036F}{III.94 (p.~261)} }
\expandafter\def\csname SRef@wb:cusp-cone-M-H1\endcsname{\href{S6_SERIES_GRIND_III.pdf\#nameddest=tag-035M}{III.95 (p.~261)} }
\expandafter\def\csname SRef@wb:cusp-cone-M-boundary-map\endcsname{\href{S6_SERIES_GRIND_III.pdf\#nameddest=tag-035O}{III.96 (p.~261)} }
\expandafter\def\csname SRef@wb:cusp-cone-ambient-collapse\endcsname{\href{S6_SERIES_GRIND_III.pdf\#nameddest=tag-036I}{III.97 (p.~262)} }
\expandafter\def\csname SRef@wb:cusp-cone-intrinsic-expansion\endcsname{\href{S6_SERIES_GRIND_III.pdf\#nameddest=tag-036S}{III.98 (p.~262)} }
\expandafter\def\csname SRef@wb:cusp-cone-rank-bound\endcsname{\href{S6_SERIES_GRIND_III.pdf\#nameddest=tag-037H}{III.99 (p.~262)} }
\expandafter\def\csname SRef@maps:src:B.1.1\endcsname{\href{S6_SERIES_GRIND_III.pdf\#nameddest=tag-02J4}{III.3.2 (p.~13)} }
\expandafter\def\csname SRef@maps:src:C.1.2\endcsname{\href{S6_SERIES_GRIND_III.pdf\#nameddest=tag-02K3}{III.4.2 (p.~28)} }
\expandafter\def\csname SRef@maps:src:D.1.10\endcsname{\href{S6_SERIES_GRIND_III.pdf\#nameddest=tag-02K8}{III.5.10 (p.~40)} }
\expandafter\def\csname SRef@wb:h24:ret:chart-lemma\endcsname{\href{S6_SERIES_GRIND_III.pdf\#nameddest=tag-03AE}{III.7.30 (p.~300)} }
\expandafter\def\csname SRef@wb:h24:ret:orientation-lemma\endcsname{\href{S6_SERIES_GRIND_III.pdf\#nameddest=tag-03BD}{III.7.31 (p.~301)} }
\expandafter\def\csname SRef@mapsIII:carrier-three-complex-wedge\endcsname{\href{S6_SERIES_GRIND_III.pdf\#nameddest=tag-02QC}{III.7.34 (p.~306)} }
\expandafter\def\csname SRef@mapsIII:carrier-extension-lemma\endcsname{\href{S6_SERIES_GRIND_III.pdf\#nameddest=tag-02Q8}{III.7.36 (p.~308)} }
\expandafter\def\csname SRef@wb:cusp-cone-native-descent\endcsname{\href{S6_SERIES_GRIND_III.pdf\#nameddest=tag-0370}{III.7.5 (p.~255)} }
\expandafter\def\csname SRef@maps:src:D.2.4\endcsname{\href{S6_SERIES_GRIND_III.pdf\#nameddest=tag-02LG}{III.10.1 (p.~319)} }
\expandafter\def\csname SRef@maps:src:A.2.2\endcsname{\href{S6_SERIES_GRIND_III.pdf\#nameddest=tag-02J3}{III.3.1 (p.~12)} }
\expandafter\def\csname SRef@maps:src:C.1.4\endcsname{\href{S6_SERIES_GRIND_III.pdf\#nameddest=tag-02K5}{III.4.4 (p.~30)} }
\expandafter\def\csname SRef@maps:src:D.2.3\endcsname{\href{S6_SERIES_GRIND_III.pdf\#nameddest=tag-02L5}{III.5.13 (p.~43)} }
\expandafter\def\csname SRef@maps:src:D.2.19\endcsname{\href{S6_SERIES_GRIND_III.pdf\#nameddest=tag-02KT}{III.5.28 (p.~80)} }
\expandafter\def\csname SRef@maps:src:D.1.5\endcsname{\href{S6_SERIES_GRIND_III.pdf\#nameddest=tag-02KE}{III.5.5 (p.~36)} }
\expandafter\def\csname SRef@maps:src:E.1.1\endcsname{\href{S6_SERIES_GRIND_V.pdf\#nameddest=tag-02M4}{V.2.1 (p.~14)} }
\expandafter\def\csname SRef@maps:src:E.2.7\endcsname{\href{S6_SERIES_GRIND_V.pdf\#nameddest=tag-02MR}{V.2.11 (p.~29)} }
\expandafter\def\csname SRef@maps:src:E.2.8\endcsname{\href{S6_SERIES_GRIND_V.pdf\#nameddest=tag-02MS}{V.2.12 (p.~30)} }
\expandafter\def\csname SRef@maps:src:E.2.10\endcsname{\href{S6_SERIES_GRIND_V.pdf\#nameddest=tag-02MA}{V.2.14 (p.~33)} }
\expandafter\def\csname SRef@maps:src:E.2.18\endcsname{\href{S6_SERIES_GRIND_V.pdf\#nameddest=tag-02MI}{V.2.22 (p.~56)} }
\expandafter\def\csname SRef@maps:src:E.2.19\endcsname{\href{S6_SERIES_GRIND_V.pdf\#nameddest=tag-02MJ}{V.2.23 (p.~60)} }
\expandafter\def\csname SRef@maps:src:E.2.20\endcsname{\href{S6_SERIES_GRIND_V.pdf\#nameddest=tag-02ML}{V.2.24 (p.~63)} }
\expandafter\def\csname SRef@maps:src:E.2.4\endcsname{\href{S6_SERIES_GRIND_V.pdf\#nameddest=tag-02MO}{V.2.8 (p.~24)} }
\expandafter\def\csname SRef@maps:src:E.2.5\endcsname{\href{S6_SERIES_GRIND_V.pdf\#nameddest=tag-02MP}{V.2.9 (p.~25)} }
\expandafter\def\csname SRef@wb:cusp-cone-metric-scope\endcsname{\href{S6_SERIES_GRIND_III.pdf\#nameddest=tag-036W}{III.7.10 (p.~262)} }
\expandafter\def\csname SRef@wb:cusp-cone-explicit-carriers\endcsname{\href{S6_SERIES_GRIND_III.pdf\#nameddest=tag-036O}{III.7.17 (p.~273)} }
\expandafter\def\csname SRef@mapsIII:carrier-real-periods\endcsname{\href{S6_SERIES_GRIND_III.pdf\#nameddest=tag-02QB}{III.7.2 (p.~244)} }
\expandafter\def\csname SRef@wb:sphere-metric-tangent-scope\endcsname{\href{S6_SERIES_GRIND_III.pdf\#nameddest=tag-03DF}{III.7.29 (p.~294)} }
\expandafter\def\csname SRef@wb:cusp-spec-native-quotient\endcsname{\href{S6_SERIES_GRIND_III.pdf\#nameddest=tag-039Y}{III.7.3 (p.~245)} }
\expandafter\def\csname SRef@wb:h24:ret:framing-proposition\endcsname{\href{S6_SERIES_GRIND_III.pdf\#nameddest=tag-03BC}{III.7.32 (p.~302)} }
\expandafter\def\csname SRef@mapsaudit:finite-degree-three\endcsname{\href{S6_SERIES_GRIND_III.pdf\#nameddest=tag-02QE}{III.7.35 (p.~307)} }
\expandafter\def\csname SRef@mapsIII:angle-forgotten-one-copy-null\endcsname{\href{S6_SERIES_GRIND_III.pdf\#nameddest=tag-02Q7}{III.7.37 (p.~308)} }
\expandafter\def\csname SRef@wb:cusp-cone-S24-local-route\endcsname{\href{S6_SERIES_GRIND_III.pdf\#nameddest=tag-0365}{III.7.7 (p.~258)} }
\expandafter\def\csname SRef@maps:src:D.3.4\endcsname{\href{S6_SERIES_GRIND_III.pdf\#nameddest=tag-02M3}{III.9.2 (p.~317)} }
\expandafter\def\csname SRef@qa:III:layout:12\endcsname{\href{S6_SERIES_GRIND_III.pdf\#nameddest=tag-033S}{III.3 (p.~13)} }
\expandafter\def\csname SRef@wb:cusp-cover-naive-action-failure\endcsname{\href{S6_SERIES_GRIND_III.pdf\#nameddest=tag-0382}{III.7.1 (p.~243)} }
\expandafter\def\csname SRef@wb:cusp-cover-global-scope\endcsname{\href{S6_SERIES_GRIND_III.pdf\#nameddest=tag-037Y}{III.7.1 (p.~243)} }
\expandafter\def\csname SRef@qa:III:layout:11\endcsname{\href{S6_SERIES_GRIND_III.pdf\#nameddest=tag-033R}{III.{{$\langle 1\rangle $1.}} (p.~271)} }
\expandafter\def\csname SRef@wb:cusp-cone-moving-S4-null-homotopy\endcsname{\href{S6_SERIES_GRIND_III.pdf\#nameddest=tag-036Z}{III.7.3.1 (p.~272)} }
\expandafter\def\csname SRef@wb:cusp-cone-sphere-map-homotopy-status\endcsname{\href{S6_SERIES_GRIND_III.pdf\#nameddest=tag-037K}{III.7.3.1 (p.~273)} }
\expandafter\def\csname SRef@wb:cusp-s12-one-copy-invariant\endcsname{\href{S6_SERIES_GRIND_III.pdf\#nameddest=tag-038X}{III.7.4 (p.~281)} }
\expandafter\def\csname SRef@wb:cusp-s12-nullity-scope\endcsname{\href{S6_SERIES_GRIND_III.pdf\#nameddest=tag-038W}{III.7.4 (p.~281)} }
\expandafter\def\csname SRef@qa:III:layout:1\endcsname{\href{S6_SERIES_GRIND_III.pdf\#nameddest=tag-033Q}{III.{{$\langle 1\rangle $5.}} (p.~283)} }
\expandafter\def\csname SRef@qa:III:layout:2\endcsname{\href{S6_SERIES_GRIND_III.pdf\#nameddest=tag-033T}{III.{{$\langle 1\rangle $6.}} (p.~316)} }
\expandafter\def\csname SRef@qa:III:layout:3\endcsname{\href{S6_SERIES_GRIND_III.pdf\#nameddest=tag-033U}{III.{{$\langle 1\rangle $6.}} (p.~316)} }
\expandafter\def\csname SRef@qa:III:layout:4\endcsname{\href{S6_SERIES_GRIND_III.pdf\#nameddest=tag-033V}{III.8 (p.~316)} }
\expandafter\def\csname SRef@sec:III:conventions\endcsname{\href{S6_SERIES_GRIND_III.pdf\#nameddest=tag-034J}{III.2 (p.~12)} }
\expandafter\def\csname SRef@sec:III:maps\endcsname{\href{S6_SERIES_GRIND_III.pdf\#nameddest=tag-034K}{III.3 (p.~12)} }
\expandafter\def\csname SRef@sec:III:os49-night\endcsname{\href{S6_SERIES_GRIND_III.pdf\#nameddest=tag-034L}{III.6 (p.~215)} }
\expandafter\def\csname SRef@sec:III:carriers\endcsname{\href{S6_SERIES_GRIND_III.pdf\#nameddest=tag-034I}{III.7 (p.~238)} }
\expandafter\def\csname SRef@mapsD:repair:modular-collisions\endcsname{\href{S6_SERIES_GRIND_III.pdf\#nameddest=tag-02OJ}{III.8 (p.~310)} }
\expandafter\def\csname SRef@wb:cusp-cover-section\endcsname{\href{S6_SERIES_GRIND_III.pdf\#nameddest=tag-0385}{III.7.1 (p.~238)} }
\expandafter\def\csname SRef@mapsIII:carrier-radial-prerequisites\endcsname{\href{S6_SERIES_GRIND_III.pdf\#nameddest=tag-02QA}{III.7.2 (p.~244)} }
\expandafter\def\csname SRef@wb:cusp-cone-section\endcsname{\href{S6_SERIES_GRIND_III.pdf\#nameddest=tag-037I}{III.7.3 (p.~255)} }
\expandafter\def\csname SRef@wb:cusp-s12-carrier-homotopy\endcsname{\href{S6_SERIES_GRIND_III.pdf\#nameddest=tag-038U}{III.7.4 (p.~279)} }
\expandafter\def\csname SRef@wb:sphere-metric-limit-note\endcsname{\href{S6_SERIES_GRIND_III.pdf\#nameddest=tag-03D9}{III.7.6 (p.~287)} }
\expandafter\def\csname SRef@wb:h24:ret:section\endcsname{\href{S6_SERIES_GRIND_III.pdf\#nameddest=tag-03BF}{III.7.7 (p.~298)} }
\expandafter\def\csname SRef@mapsIII:carrier-extension-section\endcsname{\href{S6_SERIES_GRIND_III.pdf\#nameddest=tag-02Q9}{III.7.8 (p.~308)} }
\expandafter\def\csname SRef@the-source-section-and-component-divisor\endcsname{\href{S6_SERIES_GRIND_III.pdf\#nameddest=tag-0351}{III.6.0.1 (p.~215)} }
\expandafter\def\csname SRef@the-page-including-the-finite-cap-line\endcsname{\href{S6_SERIES_GRIND_III.pdf\#nameddest=tag-034Z}{III.6.0.10 (p.~231)} }
\expandafter\def\csname SRef@the-actual-graph-line-and-return\endcsname{\href{S6_SERIES_GRIND_III.pdf\#nameddest=tag-034Q}{III.6.0.11 (p.~232)} }
\expandafter\def\csname SRef@all-stabilizers-on-the-actual-root-chart\endcsname{\href{S6_SERIES_GRIND_III.pdf\#nameddest=tag-0008}{III.6.0.12 (p.~233)} }
\expandafter\def\csname SRef@the-full-fundamental-group-presentation\endcsname{\href{S6_SERIES_GRIND_III.pdf\#nameddest=tag-034U}{III.6.0.13 (p.~234)} }
\expandafter\def\csname SRef@the-fixed-exterior-and-the-actual-boundary-circle\endcsname{\href{S6_SERIES_GRIND_III.pdf\#nameddest=tag-034T}{III.6.0.14 (p.~235)} }
\expandafter\def\csname SRef@an-integral-basis-of-the-boundary-two-cohomology\endcsname{\href{S6_SERIES_GRIND_III.pdf\#nameddest=tag-000A}{III.6.0.15 (p.~236)} }
\expandafter\def\csname SRef@determine-the-exterior-restriction-without-guessing-a-splitting\endcsname{\href{S6_SERIES_GRIND_III.pdf\#nameddest=tag-01O2}{III.6.0.16 (p.~236)} }
\expandafter\def\csname SRef@the-two-smith-invariants-and-every-homology-group\endcsname{\href{S6_SERIES_GRIND_III.pdf\#nameddest=tag-0353}{III.6.0.17 (p.~237)} }
\expandafter\def\csname SRef@actual-local-fillings-and-their-gluing\endcsname{\href{S6_SERIES_GRIND_III.pdf\#nameddest=tag-0007}{III.6.0.2 (p.~216)} }
\expandafter\def\csname SRef@marked-monodromy-and-its-completeness\endcsname{\href{S6_SERIES_GRIND_III.pdf\#nameddest=tag-02QG}{III.6.0.3 (p.~217)} }
\expandafter\def\csname SRef@integral-topology\endcsname{\href{S6_SERIES_GRIND_III.pdf\#nameddest=tag-01RP}{III.6.0.4 (p.~218)} }
\expandafter\def\csname SRef@the-exact-line-subfan\endcsname{\href{S6_SERIES_GRIND_III.pdf\#nameddest=tag-034S}{III.6.0.5 (p.~225)} }
\expandafter\def\csname SRef@a-page-carrying-the-actual-source-line\endcsname{\href{S6_SERIES_GRIND_III.pdf\#nameddest=tag-0003}{III.6.0.6 (p.~227)} }
\expandafter\def\csname SRef@the-full-graph-to-page-comparison-by-regular-bands\endcsname{\href{S6_SERIES_GRIND_III.pdf\#nameddest=tag-034V}{III.6.0.7 (p.~228)} }
\expandafter\def\csname SRef@characteristic-calculation-and-the-scope-of-the-result\endcsname{\href{S6_SERIES_GRIND_III.pdf\#nameddest=tag-01GY}{III.6.0.8 (p.~229)} }
\expandafter\def\csname SRef@the-actual-closing-map-and-its-origins\endcsname{\href{S6_SERIES_GRIND_III.pdf\#nameddest=tag-034P}{III.6.0.9 (p.~231)} }
\expandafter\def\csname SRef@wb:cusp-cone-optimal-strengthening\endcsname{\href{S6_SERIES_GRIND_III.pdf\#nameddest=tag-0376}{III.7.3.1 (p.~264)} }
\expandafter\def\csname SRef@maps:src:B.5.2\endcsname{\href{S6_SERIES_GRIND_III.pdf\#nameddest=tag-02JC}{III.3.10 (p.~23)} }
\expandafter\def\csname SRef@maps:src:B.1.2\endcsname{\href{S6_SERIES_GRIND_III.pdf\#nameddest=tag-02J5}{III.3.3 (p.~14)} }
\expandafter\def\csname SRef@maps:src:B.2.1\endcsname{\href{S6_SERIES_GRIND_III.pdf\#nameddest=tag-02J6}{III.3.4 (p.~15)} }
\expandafter\def\csname SRef@maps:src:B.3.1\endcsname{\href{S6_SERIES_GRIND_III.pdf\#nameddest=tag-02J7}{III.3.5 (p.~16)} }
\expandafter\def\csname SRef@maps:src:B.3.2\endcsname{\href{S6_SERIES_GRIND_III.pdf\#nameddest=tag-02J8}{III.3.6 (p.~17)} }
\expandafter\def\csname SRef@maps:src:B.3.3\endcsname{\href{S6_SERIES_GRIND_III.pdf\#nameddest=tag-02J9}{III.3.7 (p.~18)} }
\expandafter\def\csname SRef@maps:src:B.4.1\endcsname{\href{S6_SERIES_GRIND_III.pdf\#nameddest=tag-02JA}{III.3.8 (p.~18)} }
\expandafter\def\csname SRef@maps:src:B.5.1\endcsname{\href{S6_SERIES_GRIND_III.pdf\#nameddest=tag-02JB}{III.3.9 (p.~20)} }
\expandafter\def\csname SRef@maps:src:C.1.1\endcsname{\href{S6_SERIES_GRIND_III.pdf\#nameddest=tag-02K2}{III.4.1 (p.~27)} }
\expandafter\def\csname SRef@maps:src:C.1.3\endcsname{\href{S6_SERIES_GRIND_III.pdf\#nameddest=tag-02K4}{III.4.3 (p.~29)} }
\expandafter\def\csname SRef@maps:src:C.1.5\endcsname{\href{S6_SERIES_GRIND_III.pdf\#nameddest=tag-02K6}{III.4.5 (p.~30)} }
\expandafter\def\csname SRef@maps:src:D.1.1\endcsname{\href{S6_SERIES_GRIND_III.pdf\#nameddest=tag-02K7}{III.5.1 (p.~32)} }
\expandafter\def\csname SRef@maps:src:D.1.12\endcsname{\href{S6_SERIES_GRIND_III.pdf\#nameddest=tag-02KA}{III.5.11 (p.~41)} }
\expandafter\def\csname SRef@maps:src:D.2.5\endcsname{\href{S6_SERIES_GRIND_III.pdf\#nameddest=tag-02LR}{III.5.14 (p.~44)} }
\expandafter\def\csname SRef@maps:src:D.2.6\endcsname{\href{S6_SERIES_GRIND_III.pdf\#nameddest=tag-02LW}{III.5.15 (p.~46)} }
\expandafter\def\csname SRef@maps:src:D.2.7\endcsname{\href{S6_SERIES_GRIND_III.pdf\#nameddest=tag-02LX}{III.5.16 (p.~48)} }
\expandafter\def\csname SRef@maps:src:D.2.8\endcsname{\href{S6_SERIES_GRIND_III.pdf\#nameddest=tag-02LY}{III.5.17 (p.~50)} }
\expandafter\def\csname SRef@maps:src:D.2.9\endcsname{\href{S6_SERIES_GRIND_III.pdf\#nameddest=tag-02LZ}{III.5.18 (p.~53)} }
\expandafter\def\csname SRef@maps:src:D.2.10\endcsname{\href{S6_SERIES_GRIND_III.pdf\#nameddest=tag-02KK}{III.5.19 (p.~55)} }
\expandafter\def\csname SRef@maps:src:D.1.2\endcsname{\href{S6_SERIES_GRIND_III.pdf\#nameddest=tag-02KB}{III.5.2 (p.~32)} }
\expandafter\def\csname SRef@maps:src:D.2.11\endcsname{\href{S6_SERIES_GRIND_III.pdf\#nameddest=tag-02KL}{III.5.20 (p.~59)} }
\expandafter\def\csname SRef@maps:src:D.2.12\endcsname{\href{S6_SERIES_GRIND_III.pdf\#nameddest=tag-02KM}{III.5.21 (p.~62)} }
\expandafter\def\csname SRef@maps:src:D.2.13\endcsname{\href{S6_SERIES_GRIND_III.pdf\#nameddest=tag-02KN}{III.5.22 (p.~65)} }
\expandafter\def\csname SRef@maps:src:D.2.14\endcsname{\href{S6_SERIES_GRIND_III.pdf\#nameddest=tag-02KO}{III.5.23 (p.~68)} }
\expandafter\def\csname SRef@maps:src:D.2.15\endcsname{\href{S6_SERIES_GRIND_III.pdf\#nameddest=tag-02KP}{III.5.24 (p.~70)} }
\expandafter\def\csname SRef@maps:src:D.2.16\endcsname{\href{S6_SERIES_GRIND_III.pdf\#nameddest=tag-02KQ}{III.5.25 (p.~73)} }
\expandafter\def\csname SRef@maps:src:D.2.18\endcsname{\href{S6_SERIES_GRIND_III.pdf\#nameddest=tag-02KS}{III.5.27 (p.~78)} }
\expandafter\def\csname SRef@maps:src:D.2.20\endcsname{\href{S6_SERIES_GRIND_III.pdf\#nameddest=tag-02KV}{III.5.29 (p.~82)} }
\expandafter\def\csname SRef@maps:src:D.1.3\endcsname{\href{S6_SERIES_GRIND_III.pdf\#nameddest=tag-02KC}{III.5.3 (p.~33)} }
\expandafter\def\csname SRef@maps:src:D.2.21\endcsname{\href{S6_SERIES_GRIND_III.pdf\#nameddest=tag-02KW}{III.5.30 (p.~86)} }
\expandafter\def\csname SRef@maps:src:D.2.22\endcsname{\href{S6_SERIES_GRIND_III.pdf\#nameddest=tag-02KX}{III.5.31 (p.~87)} }
\expandafter\def\csname SRef@maps:src:D.2.23\endcsname{\href{S6_SERIES_GRIND_III.pdf\#nameddest=tag-02KY}{III.5.32 (p.~91)} }
\expandafter\def\csname SRef@maps:src:D.2.24\endcsname{\href{S6_SERIES_GRIND_III.pdf\#nameddest=tag-02KZ}{III.5.33 (p.~93)} }
\expandafter\def\csname SRef@maps:src:D.2.25\endcsname{\href{S6_SERIES_GRIND_III.pdf\#nameddest=tag-02L0}{III.5.34 (p.~94)} }
\expandafter\def\csname SRef@maps:src:D.2.26\endcsname{\href{S6_SERIES_GRIND_III.pdf\#nameddest=tag-02L1}{III.5.35 (p.~98)} }
\expandafter\def\csname SRef@maps:src:D.2.27\endcsname{\href{S6_SERIES_GRIND_III.pdf\#nameddest=tag-02L2}{III.5.36 (p.~102)} }
\expandafter\def\csname SRef@maps:src:D.2.29\endcsname{\href{S6_SERIES_GRIND_III.pdf\#nameddest=tag-02L4}{III.5.38 (p.~107)} }
\expandafter\def\csname SRef@maps:src:D.2.30\endcsname{\href{S6_SERIES_GRIND_III.pdf\#nameddest=tag-02L6}{III.5.39 (p.~110)} }
\expandafter\def\csname SRef@maps:src:D.1.4\endcsname{\href{S6_SERIES_GRIND_III.pdf\#nameddest=tag-02KD}{III.5.4 (p.~34)} }
\expandafter\def\csname SRef@maps:src:D.2.31\endcsname{\href{S6_SERIES_GRIND_III.pdf\#nameddest=tag-02L7}{III.5.40 (p.~113)} }
\expandafter\def\csname SRef@maps:src:D.2.32\endcsname{\href{S6_SERIES_GRIND_III.pdf\#nameddest=tag-02L8}{III.5.41 (p.~116)} }
\expandafter\def\csname SRef@maps:src:D.2.33\endcsname{\href{S6_SERIES_GRIND_III.pdf\#nameddest=tag-02L9}{III.5.42 (p.~119)} }
\expandafter\def\csname SRef@maps:src:D.2.34\endcsname{\href{S6_SERIES_GRIND_III.pdf\#nameddest=tag-02LA}{III.5.43 (p.~121)} }
\expandafter\def\csname SRef@maps:src:D.2.35\endcsname{\href{S6_SERIES_GRIND_III.pdf\#nameddest=tag-02LB}{III.5.44 (p.~126)} }
\expandafter\def\csname SRef@maps:src:D.2.36\endcsname{\href{S6_SERIES_GRIND_III.pdf\#nameddest=tag-02LC}{III.5.45 (p.~130)} }
\expandafter\def\csname SRef@maps:src:D.2.37\endcsname{\href{S6_SERIES_GRIND_III.pdf\#nameddest=tag-02LD}{III.5.46 (p.~134)} }
\expandafter\def\csname SRef@maps:src:D.2.38\endcsname{\href{S6_SERIES_GRIND_III.pdf\#nameddest=tag-02LE}{III.5.47 (p.~138)} }
\expandafter\def\csname SRef@maps:src:D.2.39\endcsname{\href{S6_SERIES_GRIND_III.pdf\#nameddest=tag-02LF}{III.5.48 (p.~141)} }
\expandafter\def\csname SRef@maps:src:D.2.40\endcsname{\href{S6_SERIES_GRIND_III.pdf\#nameddest=tag-02LH}{III.5.49 (p.~145)} }
\expandafter\def\csname SRef@maps:src:D.2.41\endcsname{\href{S6_SERIES_GRIND_III.pdf\#nameddest=tag-02LI}{III.5.50 (p.~148)} }
\expandafter\def\csname SRef@maps:src:D.2.42\endcsname{\href{S6_SERIES_GRIND_III.pdf\#nameddest=tag-02LJ}{III.5.51 (p.~151)} }
\expandafter\def\csname SRef@maps:src:D.2.43\endcsname{\href{S6_SERIES_GRIND_III.pdf\#nameddest=tag-02LK}{III.5.52 (p.~153)} }
\expandafter\def\csname SRef@maps:src:D.2.44\endcsname{\href{S6_SERIES_GRIND_III.pdf\#nameddest=tag-02LL}{III.5.53 (p.~156)} }
\expandafter\def\csname SRef@maps:src:D.2.45\endcsname{\href{S6_SERIES_GRIND_III.pdf\#nameddest=tag-02LM}{III.5.54 (p.~161)} }
\expandafter\def\csname SRef@maps:src:D.2.46\endcsname{\href{S6_SERIES_GRIND_III.pdf\#nameddest=tag-02LN}{III.5.55 (p.~167)} }
\expandafter\def\csname SRef@maps:src:D.2.47\endcsname{\href{S6_SERIES_GRIND_III.pdf\#nameddest=tag-02LO}{III.5.56 (p.~172)} }
\expandafter\def\csname SRef@maps:src:D.2.48\endcsname{\href{S6_SERIES_GRIND_III.pdf\#nameddest=tag-02LP}{III.5.57 (p.~175)} }
\expandafter\def\csname SRef@maps:src:D.2.49\endcsname{\href{S6_SERIES_GRIND_III.pdf\#nameddest=tag-02LQ}{III.5.58 (p.~182)} }
\expandafter\def\csname SRef@maps:src:D.2.50\endcsname{\href{S6_SERIES_GRIND_III.pdf\#nameddest=tag-02LS}{III.5.59 (p.~187)} }
\expandafter\def\csname SRef@maps:src:D.1.6\endcsname{\href{S6_SERIES_GRIND_III.pdf\#nameddest=tag-02KF}{III.5.6 (p.~36)} }
\expandafter\def\csname SRef@maps:src:D.2.51\endcsname{\href{S6_SERIES_GRIND_III.pdf\#nameddest=tag-02LT}{III.5.60 (p.~191)} }
\expandafter\def\csname SRef@maps:src:D.2.52\endcsname{\href{S6_SERIES_GRIND_III.pdf\#nameddest=tag-02LU}{III.5.61 (p.~194)} }
\expandafter\def\csname SRef@maps:src:D.2.53\endcsname{\href{S6_SERIES_GRIND_III.pdf\#nameddest=tag-02LV}{III.5.62 (p.~196)} }
\expandafter\def\csname SRef@maps:src:D.3.1\endcsname{\href{S6_SERIES_GRIND_III.pdf\#nameddest=tag-02M0}{III.5.63 (p.~211)} }
\expandafter\def\csname SRef@maps:src:D.3.3\endcsname{\href{S6_SERIES_GRIND_III.pdf\#nameddest=tag-02M2}{III.5.64 (p.~213)} }
\expandafter\def\csname SRef@maps:src:D.1.11\endcsname{\href{S6_SERIES_GRIND_III.pdf\#nameddest=tag-02K9}{III.5.65 (p.~214)} }
\expandafter\def\csname SRef@maps:src:D.1.7\endcsname{\href{S6_SERIES_GRIND_III.pdf\#nameddest=tag-02KG}{III.5.7 (p.~37)} }
\expandafter\def\csname SRef@maps:src:D.1.8\endcsname{\href{S6_SERIES_GRIND_III.pdf\#nameddest=tag-02KH}{III.5.8 (p.~38)} }
\expandafter\def\csname SRef@maps:src:D.1.9\endcsname{\href{S6_SERIES_GRIND_III.pdf\#nameddest=tag-02KI}{III.5.9 (p.~39)} }
\expandafter\def\csname SRef@night:P3F.26\endcsname{\href{S6_SERIES_GRIND_III.pdf\#nameddest=tag-02S2}{III.6.1 (p.~215)} }
\expandafter\def\csname SRef@stmt:P3F.26\endcsname{\href{S6_SERIES_GRIND_III.pdf\#nameddest=tag-02S2}{III.6.1 (p.~215)} }
\expandafter\def\csname SRef@night:P3F.27\endcsname{\href{S6_SERIES_GRIND_III.pdf\#nameddest=tag-02S3}{III.6.2 (p.~220)} }
\expandafter\def\csname SRef@stmt:P3F.27\endcsname{\href{S6_SERIES_GRIND_III.pdf\#nameddest=tag-02S3}{III.6.2 (p.~220)} }
\expandafter\def\csname SRef@night:P3F.28\endcsname{\href{S6_SERIES_GRIND_III.pdf\#nameddest=tag-02S4}{III.6.3 (p.~221)} }
\expandafter\def\csname SRef@stmt:P3F.28\endcsname{\href{S6_SERIES_GRIND_III.pdf\#nameddest=tag-02S4}{III.6.3 (p.~221)} }
\expandafter\def\csname SRef@night:P3F.29\endcsname{\href{S6_SERIES_GRIND_III.pdf\#nameddest=tag-02S5}{III.6.4 (p.~223)} }
\expandafter\def\csname SRef@stmt:P3F.29\endcsname{\href{S6_SERIES_GRIND_III.pdf\#nameddest=tag-02S5}{III.6.4 (p.~223)} }
\expandafter\def\csname SRef@night:P3F.31\endcsname{\href{S6_SERIES_GRIND_III.pdf\#nameddest=tag-02S8}{III.6.5 (p.~225)} }
\expandafter\def\csname SRef@stmt:P3F.31\endcsname{\href{S6_SERIES_GRIND_III.pdf\#nameddest=tag-02S8}{III.6.5 (p.~225)} }
\expandafter\def\csname SRef@night:P3F.32\endcsname{\href{S6_SERIES_GRIND_III.pdf\#nameddest=tag-02S9}{III.6.6 (p.~230)} }
\expandafter\def\csname SRef@stmt:P3F.32\endcsname{\href{S6_SERIES_GRIND_III.pdf\#nameddest=tag-02S9}{III.6.6 (p.~230)} }
\expandafter\def\csname SRef@night:P3F.33\endcsname{\href{S6_SERIES_GRIND_III.pdf\#nameddest=tag-02SA}{III.6.7 (p.~233)} }
\expandafter\def\csname SRef@stmt:P3F.33\endcsname{\href{S6_SERIES_GRIND_III.pdf\#nameddest=tag-02SA}{III.6.7 (p.~233)} }
\expandafter\def\csname SRef@night:P3F.34\endcsname{\href{S6_SERIES_GRIND_III.pdf\#nameddest=tag-02SB}{III.6.8 (p.~234)} }
\expandafter\def\csname SRef@stmt:P3F.34\endcsname{\href{S6_SERIES_GRIND_III.pdf\#nameddest=tag-02SB}{III.6.8 (p.~234)} }
\expandafter\def\csname SRef@maps:src:E.2.6\endcsname{\href{S6_SERIES_GRIND_V.pdf\#nameddest=tag-02MQ}{V.2.10 (p.~26)} }
\expandafter\def\csname SRef@maps:src:E.2.9\endcsname{\href{S6_SERIES_GRIND_V.pdf\#nameddest=tag-02MT}{V.2.13 (p.~31)} }
\expandafter\def\csname SRef@maps:src:E.2.11\endcsname{\href{S6_SERIES_GRIND_V.pdf\#nameddest=tag-02MB}{V.2.15 (p.~35)} }
\expandafter\def\csname SRef@maps:src:E.2.12\endcsname{\href{S6_SERIES_GRIND_V.pdf\#nameddest=tag-02MC}{V.2.16 (p.~37)} }
\expandafter\def\csname SRef@maps:src:E.2.13\endcsname{\href{S6_SERIES_GRIND_V.pdf\#nameddest=tag-02MD}{V.2.17 (p.~40)} }
\expandafter\def\csname SRef@maps:src:E.2.14\endcsname{\href{S6_SERIES_GRIND_V.pdf\#nameddest=tag-02ME}{V.2.18 (p.~42)} }
\expandafter\def\csname SRef@maps:src:E.2.15\endcsname{\href{S6_SERIES_GRIND_V.pdf\#nameddest=tag-02MF}{V.2.19 (p.~45)} }
\expandafter\def\csname SRef@maps:src:E.1.2\endcsname{\href{S6_SERIES_GRIND_V.pdf\#nameddest=tag-02M5}{V.2.2 (p.~14)} }
\expandafter\def\csname SRef@maps:src:E.2.16\endcsname{\href{S6_SERIES_GRIND_V.pdf\#nameddest=tag-02MG}{V.2.20 (p.~48)} }
\expandafter\def\csname SRef@maps:src:E.2.17\endcsname{\href{S6_SERIES_GRIND_V.pdf\#nameddest=tag-02MH}{V.2.21 (p.~52)} }
\expandafter\def\csname SRef@maps:src:E.2.21\endcsname{\href{S6_SERIES_GRIND_V.pdf\#nameddest=tag-02MM}{V.2.25 (p.~65)} }
\expandafter\def\csname SRef@maps:src:E.1.3\endcsname{\href{S6_SERIES_GRIND_V.pdf\#nameddest=tag-02M6}{V.2.3 (p.~16)} }
\expandafter\def\csname SRef@maps:src:E.1.4\endcsname{\href{S6_SERIES_GRIND_V.pdf\#nameddest=tag-02M7}{V.2.4 (p.~18)} }
\expandafter\def\csname SRef@maps:src:E.1.5\endcsname{\href{S6_SERIES_GRIND_V.pdf\#nameddest=tag-02M8}{V.2.5 (p.~19)} }
\expandafter\def\csname SRef@maps:src:E.2.1\endcsname{\href{S6_SERIES_GRIND_V.pdf\#nameddest=tag-02M9}{V.2.6 (p.~21)} }
\expandafter\def\csname SRef@maps:src:E.2.3\endcsname{\href{S6_SERIES_GRIND_V.pdf\#nameddest=tag-02MN}{V.2.7 (p.~23)} }
\expandafter\def\csname SRef@wb:cusp-cover-extension-theorem\endcsname{\href{S6_SERIES_GRIND_III.pdf\#nameddest=tag-037X}{III.7.1 (p.~238)} }
\expandafter\def\csname SRef@wb:cusp-cone-obstructions\endcsname{\href{S6_SERIES_GRIND_III.pdf\#nameddest=tag-0374}{III.7.11 (p.~262)} }
\expandafter\def\csname SRef@wb:cusp-cone-optimal-exponents\endcsname{\href{S6_SERIES_GRIND_III.pdf\#nameddest=tag-0375}{III.7.12 (p.~265)} }
\expandafter\def\csname SRef@wb:cusp-cone-angle-forgotten-sphere-carriers\endcsname{\href{S6_SERIES_GRIND_III.pdf\#nameddest=tag-036J}{III.7.13 (p.~266)} }
\expandafter\def\csname SRef@wb:cusp-cone-target-homotopy-types\endcsname{\href{S6_SERIES_GRIND_III.pdf\#nameddest=tag-037M}{III.7.18 (p.~274)} }
\expandafter\def\csname SRef@wb:cusp-cone-exact-decompactification-space\endcsname{\href{S6_SERIES_GRIND_III.pdf\#nameddest=tag-036N}{III.7.19 (p.~276)} }
\expandafter\def\csname SRef@wb:cusp-s12-full-carrier-null\endcsname{\href{S6_SERIES_GRIND_III.pdf\#nameddest=tag-038V}{III.7.21 (p.~280)} }
\expandafter\def\csname SRef@wb:cusp-null:thm\endcsname{\href{S6_SERIES_GRIND_III.pdf\#nameddest=tag-038S}{III.7.23 (p.~282)} }
\expandafter\def\csname SRef@wb:sphere-metric-uniform-collapse\endcsname{\href{S6_SERIES_GRIND_III.pdf\#nameddest=tag-03DG}{III.7.25 (p.~288)} }
\expandafter\def\csname SRef@wb:sphere-metric-fixed-realization\endcsname{\href{S6_SERIES_GRIND_III.pdf\#nameddest=tag-03D6}{III.7.26 (p.~288)} }
\expandafter\def\csname SRef@wb:sphere-metric-scaling-trichotomy\endcsname{\href{S6_SERIES_GRIND_III.pdf\#nameddest=tag-03DC}{III.7.27 (p.~291)} }
\expandafter\def\csname SRef@wb:cusp-spec-globalization-obstruction\endcsname{\href{S6_SERIES_GRIND_III.pdf\#nameddest=tag-039U}{III.7.4 (p.~246)} }
\expandafter\def\csname SRef@wb:cusp-cone-S12-route\endcsname{\href{S6_SERIES_GRIND_III.pdf\#nameddest=tag-035X}{III.7.6 (p.~256)} }
\expandafter\def\csname SRef@wb:cusp-cone-S24-route\endcsname{\href{S6_SERIES_GRIND_III.pdf\#nameddest=tag-036A}{III.7.8 (p.~259)} }
\expandafter\def\csname SRef@mapsD:repair:collision-realization\endcsname{\href{S6_SERIES_GRIND_III.pdf\#nameddest=tag-02OG}{III.8.1 (p.~310)} }
\expandafter\def\csname SRef@maps:src:D.3.2\endcsname{\href{S6_SERIES_GRIND_III.pdf\#nameddest=tag-02M1}{III.9.1 (p.~316)} }
\expandafter\def\csname SRef@maps:src:E.2.2\endcsname{\href{S6_SERIES_GRIND_III.pdf\#nameddest=tag-02MK}{III.9.3 (p.~317)} }
\expandafter\def\csname SRef@maps:eq:B.6.2\endcsname{\href{S6_SERIES_GRIND_IV.pdf\#nameddest=tag-01VZ}{IV.{B.6.2} (p.~68)} }
\expandafter\def\csname SRef@maps:eq:B.6.3\endcsname{\href{S6_SERIES_GRIND_IV.pdf\#nameddest=tag-01WA}{IV.{B.6.3} (p.~68)} }
\expandafter\def\csname SRef@maps:eq:B.6.4\endcsname{\href{S6_SERIES_GRIND_IV.pdf\#nameddest=tag-01WL}{IV.{B.6.4} (p.~69)} }
\expandafter\def\csname SRef@maps:eq:B.6.5\endcsname{\href{S6_SERIES_GRIND_IV.pdf\#nameddest=tag-01WR}{IV.{B.6.5} (p.~69)} }
\expandafter\def\csname SRef@maps:eq:B.6.6\endcsname{\href{S6_SERIES_GRIND_IV.pdf\#nameddest=tag-01WS}{IV.{B.6.6} (p.~69)} }
\expandafter\def\csname SRef@maps:eq:B.6.7\endcsname{\href{S6_SERIES_GRIND_IV.pdf\#nameddest=tag-01WT}{IV.{B.6.7} (p.~70)} }
\expandafter\def\csname SRef@maps:eq:B.6.8\endcsname{\href{S6_SERIES_GRIND_IV.pdf\#nameddest=tag-01WU}{IV.{B.6.8} (p.~70)} }
\expandafter\def\csname SRef@maps:eq:B.6.9\endcsname{\href{S6_SERIES_GRIND_IV.pdf\#nameddest=tag-01WV}{IV.{B.6.9} (p.~70)} }
\expandafter\def\csname SRef@maps:eq:B.6.10\endcsname{\href{S6_SERIES_GRIND_IV.pdf\#nameddest=tag-01VQ}{IV.{B.6.10} (p.~70)} }
\expandafter\def\csname SRef@maps:eq:B.6.11\endcsname{\href{S6_SERIES_GRIND_IV.pdf\#nameddest=tag-01VR}{IV.{B.6.11} (p.~72)} }
\expandafter\def\csname SRef@maps:eq:B.6.12\endcsname{\href{S6_SERIES_GRIND_IV.pdf\#nameddest=tag-01VS}{IV.{B.6.12} (p.~72)} }
\expandafter\def\csname SRef@maps:eq:B.6.13\endcsname{\href{S6_SERIES_GRIND_IV.pdf\#nameddest=tag-01VT}{IV.{B.6.13} (p.~73)} }
\expandafter\def\csname SRef@maps:eq:B.6.14\endcsname{\href{S6_SERIES_GRIND_IV.pdf\#nameddest=tag-01VU}{IV.{B.6.14} (p.~73)} }
\expandafter\def\csname SRef@maps:eq:B.6.15\endcsname{\href{S6_SERIES_GRIND_IV.pdf\#nameddest=tag-01VV}{IV.{B.6.15} (p.~75)} }
\expandafter\def\csname SRef@maps:eq:B.6.16\endcsname{\href{S6_SERIES_GRIND_IV.pdf\#nameddest=tag-01VW}{IV.{B.6.16} (p.~75)} }
\expandafter\def\csname SRef@maps:eq:B.6.17\endcsname{\href{S6_SERIES_GRIND_IV.pdf\#nameddest=tag-01VX}{IV.{B.6.17} (p.~76)} }
\expandafter\def\csname SRef@maps:eq:B.C1\endcsname{\href{S6_SERIES_GRIND_IV.pdf\#nameddest=tag-01WW}{IV.{B.C1} (p.~77)} }
\expandafter\def\csname SRef@maps:eq:B.C2\endcsname{\href{S6_SERIES_GRIND_IV.pdf\#nameddest=tag-01WX}{IV.{B.C2} (p.~77)} }
\expandafter\def\csname SRef@maps:eq:B.C3\endcsname{\href{S6_SERIES_GRIND_IV.pdf\#nameddest=tag-01WY}{IV.{B.C3} (p.~77)} }
\expandafter\def\csname SRef@maps:eq:B.C4\endcsname{\href{S6_SERIES_GRIND_IV.pdf\#nameddest=tag-01WZ}{IV.{B.C4} (p.~78)} }
\expandafter\def\csname SRef@maps:eq:B.P1\endcsname{\href{S6_SERIES_GRIND_IV.pdf\#nameddest=tag-01X0}{IV.{B.P1} (p.~78)} }
\expandafter\def\csname SRef@maps:eq:B.P2\endcsname{\href{S6_SERIES_GRIND_IV.pdf\#nameddest=tag-01X1}{IV.{B.P2} (p.~79)} }
\expandafter\def\csname SRef@maps:eq:B.P3\endcsname{\href{S6_SERIES_GRIND_IV.pdf\#nameddest=tag-01X2}{IV.{B.P3} (p.~79)} }
\expandafter\def\csname SRef@maps:eq:B.P4\endcsname{\href{S6_SERIES_GRIND_IV.pdf\#nameddest=tag-01X3}{IV.{B.P4} (p.~79)} }
\expandafter\def\csname SRef@maps:eq:B.6.19\endcsname{\href{S6_SERIES_GRIND_IV.pdf\#nameddest=tag-01VY}{IV.{B.6.19} (p.~80)} }
\expandafter\def\csname SRef@maps:eq:B.6.20\endcsname{\href{S6_SERIES_GRIND_IV.pdf\#nameddest=tag-01W0}{IV.{B.6.20} (p.~81)} }
\expandafter\def\csname SRef@maps:eq:B.6.21\endcsname{\href{S6_SERIES_GRIND_IV.pdf\#nameddest=tag-01W1}{IV.{B.6.21} (p.~83)} }
\expandafter\def\csname SRef@maps:eq:B.6.22\endcsname{\href{S6_SERIES_GRIND_IV.pdf\#nameddest=tag-01W2}{IV.{B.6.22} (p.~83)} }
\expandafter\def\csname SRef@maps:eq:B.6.23\endcsname{\href{S6_SERIES_GRIND_IV.pdf\#nameddest=tag-01W3}{IV.{B.6.23} (p.~83)} }
\expandafter\def\csname SRef@maps:eq:B.6.24\endcsname{\href{S6_SERIES_GRIND_IV.pdf\#nameddest=tag-01W4}{IV.{B.6.24} (p.~84)} }
\expandafter\def\csname SRef@maps:eq:B.6.25\endcsname{\href{S6_SERIES_GRIND_IV.pdf\#nameddest=tag-01W5}{IV.{B.6.25} (p.~84)} }
\expandafter\def\csname SRef@maps:eq:B.T1\endcsname{\href{S6_SERIES_GRIND_IV.pdf\#nameddest=tag-01XD}{IV.{B.T1} (p.~87)} }
\expandafter\def\csname SRef@maps:eq:B.6.26\endcsname{\href{S6_SERIES_GRIND_IV.pdf\#nameddest=tag-01W6}{IV.{B.6.26} (p.~87)} }
\expandafter\def\csname SRef@maps:eq:B.6.27\endcsname{\href{S6_SERIES_GRIND_IV.pdf\#nameddest=tag-01W7}{IV.{B.6.27} (p.~88)} }
\expandafter\def\csname SRef@maps:eq:B.6.28\endcsname{\href{S6_SERIES_GRIND_IV.pdf\#nameddest=tag-01W8}{IV.{B.6.28} (p.~89)} }
\expandafter\def\csname SRef@maps:eq:B.6.29\endcsname{\href{S6_SERIES_GRIND_IV.pdf\#nameddest=tag-01W9}{IV.{B.6.29} (p.~89)} }
\expandafter\def\csname SRef@maps:eq:B.6.30\endcsname{\href{S6_SERIES_GRIND_IV.pdf\#nameddest=tag-01WB}{IV.{B.6.30} (p.~89)} }
\expandafter\def\csname SRef@maps:eq:B.6.31\endcsname{\href{S6_SERIES_GRIND_IV.pdf\#nameddest=tag-01WC}{IV.{B.6.31} (p.~90)} }
\expandafter\def\csname SRef@maps:eq:B.Q1\endcsname{\href{S6_SERIES_GRIND_IV.pdf\#nameddest=tag-01X4}{IV.{B.Q1} (p.~91)} }
\expandafter\def\csname SRef@maps:eq:B.Q2\endcsname{\href{S6_SERIES_GRIND_IV.pdf\#nameddest=tag-01X5}{IV.{B.Q2} (p.~91)} }
\expandafter\def\csname SRef@maps:eq:B.Q3\endcsname{\href{S6_SERIES_GRIND_IV.pdf\#nameddest=tag-01X6}{IV.{B.Q3} (p.~91)} }
\expandafter\def\csname SRef@maps:eq:B.Q4\endcsname{\href{S6_SERIES_GRIND_IV.pdf\#nameddest=tag-01X7}{IV.{B.Q4} (p.~92)} }
\expandafter\def\csname SRef@maps:eq:B.Q5\endcsname{\href{S6_SERIES_GRIND_IV.pdf\#nameddest=tag-01X8}{IV.{B.Q5} (p.~92)} }
\expandafter\def\csname SRef@maps:eq:B.Q6\endcsname{\href{S6_SERIES_GRIND_IV.pdf\#nameddest=tag-01X9}{IV.{B.Q6} (p.~92)} }
\expandafter\def\csname SRef@maps:eq:B.Q7\endcsname{\href{S6_SERIES_GRIND_IV.pdf\#nameddest=tag-01XA}{IV.{B.Q7} (p.~92)} }
\expandafter\def\csname SRef@maps:eq:B.Q8\endcsname{\href{S6_SERIES_GRIND_IV.pdf\#nameddest=tag-01XB}{IV.{B.Q8} (p.~92)} }
\expandafter\def\csname SRef@maps:eq:B.Q9\endcsname{\href{S6_SERIES_GRIND_IV.pdf\#nameddest=tag-01XC}{IV.{B.Q9} (p.~93)} }
\expandafter\def\csname SRef@maps:eq:B.6.32\endcsname{\href{S6_SERIES_GRIND_IV.pdf\#nameddest=tag-01WD}{IV.{B.6.32} (p.~93)} }
\expandafter\def\csname SRef@maps:eq:B.6.33\endcsname{\href{S6_SERIES_GRIND_IV.pdf\#nameddest=tag-01WE}{IV.{B.6.33} (p.~94)} }
\expandafter\def\csname SRef@maps:eq:B.6.34\endcsname{\href{S6_SERIES_GRIND_IV.pdf\#nameddest=tag-01WF}{IV.{B.6.34} (p.~94)} }
\expandafter\def\csname SRef@famIII:eq:E9\endcsname{\href{S6_SERIES_GRIND_IV.pdf\#nameddest=tag-01RJ}{IV.{E9} (p.~20)} }
\expandafter\def\csname SRef@maps:eq:B.6.35\endcsname{\href{S6_SERIES_GRIND_IV.pdf\#nameddest=tag-01WG}{IV.{B.6.35} (p.~95)} }
\expandafter\def\csname SRef@maps:eq:B.6.36\endcsname{\href{S6_SERIES_GRIND_IV.pdf\#nameddest=tag-01WH}{IV.{B.6.36} (p.~95)} }
\expandafter\def\csname SRef@maps:eq:B.6.37\endcsname{\href{S6_SERIES_GRIND_IV.pdf\#nameddest=tag-01WI}{IV.{B.6.37} (p.~95)} }
\expandafter\def\csname SRef@maps:eq:B.6.38\endcsname{\href{S6_SERIES_GRIND_IV.pdf\#nameddest=tag-01WJ}{IV.{B.6.38} (p.~99)} }
\expandafter\def\csname SRef@maps:eq:B.6.39\endcsname{\href{S6_SERIES_GRIND_IV.pdf\#nameddest=tag-01WK}{IV.{B.6.39} (p.~100)} }
\expandafter\def\csname SRef@famIII:eq:E10\endcsname{\href{S6_SERIES_GRIND_IV.pdf\#nameddest=tag-01RI}{IV.{E10} (p.~20)} }
\expandafter\def\csname SRef@maps:eq:B.6.40\endcsname{\href{S6_SERIES_GRIND_IV.pdf\#nameddest=tag-01WM}{IV.{B.6.40} (p.~100)} }
\expandafter\def\csname SRef@maps:eq:B.6.41\endcsname{\href{S6_SERIES_GRIND_IV.pdf\#nameddest=tag-01WN}{IV.{B.6.41} (p.~100)} }
\expandafter\def\csname SRef@maps:eq:B.6.42\endcsname{\href{S6_SERIES_GRIND_IV.pdf\#nameddest=tag-01WO}{IV.{B.6.42} (p.~100)} }
\expandafter\def\csname SRef@maps:eq:B.6.43\endcsname{\href{S6_SERIES_GRIND_IV.pdf\#nameddest=tag-01WP}{IV.{B.6.43} (p.~103)} }
\expandafter\def\csname SRef@maps:eq:B.6.44\endcsname{\href{S6_SERIES_GRIND_IV.pdf\#nameddest=tag-01WQ}{IV.{B.6.44} (p.~103)} }
\expandafter\def\csname SRef@maps:eq:B.6.1\endcsname{\href{S6_SERIES_GRIND_IV.pdf\#nameddest=tag-01VP}{IV.{B.6.1} (p.~68)} }
\expandafter\def\csname SRef@push:src:C.3.15\endcsname{\href{S6_SERIES_GRIND_IV.pdf\#nameddest=tag-031M}{IV.2.16 (p.~40)} }
\expandafter\def\csname SRef@push:src:C.3.27\endcsname{\href{S6_SERIES_GRIND_IV.pdf\#nameddest=tag-031Z}{IV.2.28 (p.~63)} }
\expandafter\def\csname SRef@push:src:C.5.1\endcsname{\href{S6_SERIES_GRIND_IV.pdf\#nameddest=tag-0329}{IV.2.31 (p.~65)} }
\expandafter\def\csname SRef@maps:src:B.6.8\endcsname{\href{S6_SERIES_GRIND_IV.pdf\#nameddest=tag-02K0}{IV.2.39 (p.~76)} }
\expandafter\def\csname SRef@maps:src:B.6.25\endcsname{\href{S6_SERIES_GRIND_IV.pdf\#nameddest=tag-02JU}{IV.2.55 (p.~104)} }
\expandafter\def\csname SRef@push:eq:C12-u-flow\endcsname{\href{S6_SERIES_GRIND_IV.pdf\#nameddest=tag-02YV}{IV.1 (p.~33)} }
\expandafter\def\csname SRef@push:eq:C12-limiting-identification\endcsname{\href{S6_SERIES_GRIND_IV.pdf\#nameddest=tag-02YU}{IV.2 (p.~34)} }
\expandafter\def\csname SRef@push:eq:C13-real-overlap\endcsname{\href{S6_SERIES_GRIND_IV.pdf\#nameddest=tag-02YX}{IV.3 (p.~35)} }
\expandafter\def\csname SRef@push:eq:C13-weighted-map\endcsname{\href{S6_SERIES_GRIND_IV.pdf\#nameddest=tag-02YY}{IV.4 (p.~35)} }
\expandafter\def\csname SRef@push:eq:C13-cusp-gauge\endcsname{\href{S6_SERIES_GRIND_IV.pdf\#nameddest=tag-02YW}{IV.5 (p.~37)} }
\expandafter\def\csname SRef@push:src:C.3.9\endcsname{\href{S6_SERIES_GRIND_IV.pdf\#nameddest=tag-0326}{IV.2.10 (p.~29)} }
\expandafter\def\csname SRef@push:src:C.3.22\endcsname{\href{S6_SERIES_GRIND_IV.pdf\#nameddest=tag-031U}{IV.2.23 (p.~54)} }
\expandafter\def\csname SRef@push:src:C.3.8\endcsname{\href{S6_SERIES_GRIND_IV.pdf\#nameddest=tag-0325}{IV.2.9 (p.~28)} }
\expandafter\def\csname SRef@iv:rank_two_specialization\endcsname{\href{S6_SERIES_GRIND_IV.pdf\#nameddest=tag-01UU}{IV.3.20 (p.~131)} }
\expandafter\def\csname SRef@fam:src:E.2.2\endcsname{\href{S6_SERIES_GRIND_IV.pdf\#nameddest=tag-01R7}{IV.2.1 (p.~20)} }
\expandafter\def\csname SRef@push:src:C.3.10\endcsname{\href{S6_SERIES_GRIND_IV.pdf\#nameddest=tag-031H}{IV.2.11 (p.~30)} }
\expandafter\def\csname SRef@push:src:C.3.11\endcsname{\href{S6_SERIES_GRIND_IV.pdf\#nameddest=tag-031I}{IV.2.12 (p.~32)} }
\expandafter\def\csname SRef@push:src:C.3.12\endcsname{\href{S6_SERIES_GRIND_IV.pdf\#nameddest=tag-031J}{IV.2.13 (p.~32)} }
\expandafter\def\csname SRef@push:src:C.3.13\endcsname{\href{S6_SERIES_GRIND_IV.pdf\#nameddest=tag-031K}{IV.2.14 (p.~34)} }
\expandafter\def\csname SRef@push:src:C.3.14\endcsname{\href{S6_SERIES_GRIND_IV.pdf\#nameddest=tag-031L}{IV.2.15 (p.~37)} }
\expandafter\def\csname SRef@push:src:C.3.16\endcsname{\href{S6_SERIES_GRIND_IV.pdf\#nameddest=tag-031N}{IV.2.17 (p.~40)} }
\expandafter\def\csname SRef@push:src:C.3.17\endcsname{\href{S6_SERIES_GRIND_IV.pdf\#nameddest=tag-031O}{IV.2.18 (p.~41)} }
\expandafter\def\csname SRef@push:src:C.3.18\endcsname{\href{S6_SERIES_GRIND_IV.pdf\#nameddest=tag-031P}{IV.2.19 (p.~44)} }
\expandafter\def\csname SRef@push:src:C.3.19\endcsname{\href{S6_SERIES_GRIND_IV.pdf\#nameddest=tag-031Q}{IV.2.20 (p.~47)} }
\expandafter\def\csname SRef@push:src:C.3.20\endcsname{\href{S6_SERIES_GRIND_IV.pdf\#nameddest=tag-031S}{IV.2.21 (p.~49)} }
\expandafter\def\csname SRef@push:src:C.3.21\endcsname{\href{S6_SERIES_GRIND_IV.pdf\#nameddest=tag-031T}{IV.2.22 (p.~51)} }
\expandafter\def\csname SRef@push:src:C.3.23\endcsname{\href{S6_SERIES_GRIND_IV.pdf\#nameddest=tag-031V}{IV.2.24 (p.~55)} }
\expandafter\def\csname SRef@push:src:C.2.2\endcsname{\href{S6_SERIES_GRIND_IV.pdf\#nameddest=tag-031E}{IV.2.3 (p.~22)} }
\expandafter\def\csname SRef@push:src:C.3.25\endcsname{\href{S6_SERIES_GRIND_IV.pdf\#nameddest=tag-031X}{IV.2.26 (p.~58)} }
\expandafter\def\csname SRef@push:src:C.3.26\endcsname{\href{S6_SERIES_GRIND_IV.pdf\#nameddest=tag-031Y}{IV.2.27 (p.~61)} }
\expandafter\def\csname SRef@push:src:C.4.1\endcsname{\href{S6_SERIES_GRIND_IV.pdf\#nameddest=tag-0327}{IV.2.29 (p.~64)} }
\expandafter\def\csname SRef@push:src:C.4.2\endcsname{\href{S6_SERIES_GRIND_IV.pdf\#nameddest=tag-0328}{IV.2.30 (p.~65)} }
\expandafter\def\csname SRef@push:src:C.2.3\endcsname{\href{S6_SERIES_GRIND_IV.pdf\#nameddest=tag-031F}{IV.2.4 (p.~23)} }
\expandafter\def\csname SRef@maps:src:B.6.1\endcsname{\href{S6_SERIES_GRIND_IV.pdf\#nameddest=tag-02JD}{IV.2.32 (p.~66)} }
\expandafter\def\csname SRef@maps:src:B.6.3\endcsname{\href{S6_SERIES_GRIND_IV.pdf\#nameddest=tag-02JV}{IV.2.34 (p.~67)} }
\expandafter\def\csname SRef@maps:src:B.6.4\endcsname{\href{S6_SERIES_GRIND_IV.pdf\#nameddest=tag-02JW}{IV.2.35 (p.~69)} }
\expandafter\def\csname SRef@maps:src:B.6.5\endcsname{\href{S6_SERIES_GRIND_IV.pdf\#nameddest=tag-02JX}{IV.2.36 (p.~71)} }
\expandafter\def\csname SRef@maps:src:B.6.6\endcsname{\href{S6_SERIES_GRIND_IV.pdf\#nameddest=tag-02JY}{IV.2.37 (p.~73)} }
\expandafter\def\csname SRef@maps:src:B.6.7\endcsname{\href{S6_SERIES_GRIND_IV.pdf\#nameddest=tag-02JZ}{IV.2.38 (p.~75)} }
\expandafter\def\csname SRef@maps:src:B.6.10\endcsname{\href{S6_SERIES_GRIND_IV.pdf\#nameddest=tag-02JE}{IV.2.40 (p.~77)} }
\expandafter\def\csname SRef@maps:src:B.6.11\endcsname{\href{S6_SERIES_GRIND_IV.pdf\#nameddest=tag-02JF}{IV.2.41 (p.~78)} }
\expandafter\def\csname SRef@push:src:C.3.2\endcsname{\href{S6_SERIES_GRIND_IV.pdf\#nameddest=tag-031R}{IV.2.5 (p.~25)} }
\expandafter\def\csname SRef@maps:src:B.6.12\endcsname{\href{S6_SERIES_GRIND_IV.pdf\#nameddest=tag-02JG}{IV.2.42 (p.~80)} }
\expandafter\def\csname SRef@maps:src:B.6.13\endcsname{\href{S6_SERIES_GRIND_IV.pdf\#nameddest=tag-02JH}{IV.2.43 (p.~83)} }
\expandafter\def\csname SRef@maps:src:B.6.14\endcsname{\href{S6_SERIES_GRIND_IV.pdf\#nameddest=tag-02JI}{IV.2.44 (p.~86)} }
\expandafter\def\csname SRef@maps:src:B.6.15\endcsname{\href{S6_SERIES_GRIND_IV.pdf\#nameddest=tag-02JJ}{IV.2.45 (p.~87)} }
\expandafter\def\csname SRef@maps:src:B.6.16\endcsname{\href{S6_SERIES_GRIND_IV.pdf\#nameddest=tag-02JK}{IV.2.46 (p.~88)} }
\expandafter\def\csname SRef@maps:src:B.6.17\endcsname{\href{S6_SERIES_GRIND_IV.pdf\#nameddest=tag-02JL}{IV.2.47 (p.~91)} }
\expandafter\def\csname SRef@maps:src:B.6.18\endcsname{\href{S6_SERIES_GRIND_IV.pdf\#nameddest=tag-02JM}{IV.2.48 (p.~93)} }
\expandafter\def\csname SRef@maps:src:B.6.19\endcsname{\href{S6_SERIES_GRIND_IV.pdf\#nameddest=tag-02JN}{IV.2.49 (p.~94)} }
\expandafter\def\csname SRef@push:src:C.3.4\endcsname{\href{S6_SERIES_GRIND_IV.pdf\#nameddest=tag-0321}{IV.2.6 (p.~25)} }
\expandafter\def\csname SRef@maps:src:B.6.22\endcsname{\href{S6_SERIES_GRIND_IV.pdf\#nameddest=tag-02JR}{IV.2.52 (p.~99)} }
\expandafter\def\csname SRef@push:src:C.3.6\endcsname{\href{S6_SERIES_GRIND_IV.pdf\#nameddest=tag-0323}{IV.2.7 (p.~25)} }
\expandafter\def\csname SRef@push:src:C.3.7\endcsname{\href{S6_SERIES_GRIND_IV.pdf\#nameddest=tag-0324}{IV.2.8 (p.~26)} }
\expandafter\def\csname SRef@iv:correction:P3F.25\endcsname{\href{S6_SERIES_GRIND_IV.pdf\#nameddest=tag-01TX}{IV.8.1 (p.~933)} }
\expandafter\def\csname SRef@iv:correction:P3K.29\endcsname{\href{S6_SERIES_GRIND_IV.pdf\#nameddest=tag-01U6}{IV.8.10 (p.~937)} }
\expandafter\def\csname SRef@iv:correction:P3K.30\endcsname{\href{S6_SERIES_GRIND_IV.pdf\#nameddest=tag-01U7}{IV.8.11 (p.~938)} }
\expandafter\def\csname SRef@iv:correction:P3K.31\endcsname{\href{S6_SERIES_GRIND_IV.pdf\#nameddest=tag-01U8}{IV.8.12 (p.~938)} }
\expandafter\def\csname SRef@iv:correction:P3K.32\endcsname{\href{S6_SERIES_GRIND_IV.pdf\#nameddest=tag-01U9}{IV.8.13 (p.~939)} }
\expandafter\def\csname SRef@iv:correction:P3K.33\endcsname{\href{S6_SERIES_GRIND_IV.pdf\#nameddest=tag-01UA}{IV.8.14 (p.~939)} }
\expandafter\def\csname SRef@iv:correction:P3K.34\endcsname{\href{S6_SERIES_GRIND_IV.pdf\#nameddest=tag-01UB}{IV.8.15 (p.~940)} }
\expandafter\def\csname SRef@iv:correction:P3F.53\endcsname{\href{S6_SERIES_GRIND_IV.pdf\#nameddest=tag-01TY}{IV.8.16 (p.~940)} }
\expandafter\def\csname SRef@iv:correction:P3K.35\endcsname{\href{S6_SERIES_GRIND_IV.pdf\#nameddest=tag-01UC}{IV.8.17 (p.~940)} }
\expandafter\def\csname SRef@iv:correction:P3K.36\endcsname{\href{S6_SERIES_GRIND_IV.pdf\#nameddest=tag-01UD}{IV.8.18 (p.~941)} }
\expandafter\def\csname SRef@iv:correction:P3K.37\endcsname{\href{S6_SERIES_GRIND_IV.pdf\#nameddest=tag-01UE}{IV.8.19 (p.~941)} }
\expandafter\def\csname SRef@iv:correction:P3K.6\endcsname{\href{S6_SERIES_GRIND_IV.pdf\#nameddest=tag-01US}{IV.8.2 (p.~934)} }
\expandafter\def\csname SRef@iv:correction:P3K.40\endcsname{\href{S6_SERIES_GRIND_IV.pdf\#nameddest=tag-01UF}{IV.8.20 (p.~942)} }
\expandafter\def\csname SRef@iv:correction:P3K.41\endcsname{\href{S6_SERIES_GRIND_IV.pdf\#nameddest=tag-01UG}{IV.8.21 (p.~942)} }
\expandafter\def\csname SRef@iv:correction:P3K.43\endcsname{\href{S6_SERIES_GRIND_IV.pdf\#nameddest=tag-01UH}{IV.8.22 (p.~943)} }
\expandafter\def\csname SRef@iv:correction:P3K.44\endcsname{\href{S6_SERIES_GRIND_IV.pdf\#nameddest=tag-01UI}{IV.8.23 (p.~943)} }
\expandafter\def\csname SRef@iv:correction:P3K.45\endcsname{\href{S6_SERIES_GRIND_IV.pdf\#nameddest=tag-01UJ}{IV.8.24 (p.~944)} }
\expandafter\def\csname SRef@iv:correction:P3K.46\endcsname{\href{S6_SERIES_GRIND_IV.pdf\#nameddest=tag-01UK}{IV.8.25 (p.~944)} }
\expandafter\def\csname SRef@iv:correction:P3K.50\endcsname{\href{S6_SERIES_GRIND_IV.pdf\#nameddest=tag-01UL}{IV.8.26 (p.~944)} }
\expandafter\def\csname SRef@iv:correction:P3K.52\endcsname{\href{S6_SERIES_GRIND_IV.pdf\#nameddest=tag-01UM}{IV.8.27 (p.~944)} }
\expandafter\def\csname SRef@iv:correction:P3K.54\endcsname{\href{S6_SERIES_GRIND_IV.pdf\#nameddest=tag-01UN}{IV.8.28 (p.~945)} }
\expandafter\def\csname SRef@iv:correction:P3K.56\endcsname{\href{S6_SERIES_GRIND_IV.pdf\#nameddest=tag-01UO}{IV.8.29 (p.~945)} }
\expandafter\def\csname SRef@iv:correction:P3K.10\endcsname{\href{S6_SERIES_GRIND_IV.pdf\#nameddest=tag-01TZ}{IV.8.3 (p.~934)} }
\expandafter\def\csname SRef@iv:correction:P3K.57\endcsname{\href{S6_SERIES_GRIND_IV.pdf\#nameddest=tag-01UP}{IV.8.30 (p.~945)} }
\expandafter\def\csname SRef@iv:correction:P3K.58\endcsname{\href{S6_SERIES_GRIND_IV.pdf\#nameddest=tag-01UQ}{IV.8.31 (p.~946)} }
\expandafter\def\csname SRef@iv:correction:P3K.59\endcsname{\href{S6_SERIES_GRIND_IV.pdf\#nameddest=tag-01UR}{IV.8.32 (p.~946)} }
\expandafter\def\csname SRef@iv:correction:P3K.60\endcsname{\href{S6_SERIES_GRIND_IV.pdf\#nameddest=tag-01UT}{IV.8.33 (p.~947)} }
\expandafter\def\csname SRef@iv:correction:P3K.17\endcsname{\href{S6_SERIES_GRIND_IV.pdf\#nameddest=tag-01U0}{IV.8.4 (p.~935)} }
\expandafter\def\csname SRef@iv:correction:P3K.22\endcsname{\href{S6_SERIES_GRIND_IV.pdf\#nameddest=tag-01U1}{IV.8.5 (p.~935)} }
\expandafter\def\csname SRef@iv:correction:P3K.23\endcsname{\href{S6_SERIES_GRIND_IV.pdf\#nameddest=tag-01U2}{IV.8.6 (p.~936)} }
\expandafter\def\csname SRef@iv:correction:P3K.24\endcsname{\href{S6_SERIES_GRIND_IV.pdf\#nameddest=tag-01U3}{IV.8.7 (p.~936)} }
\expandafter\def\csname SRef@iv:correction:P3K.25\endcsname{\href{S6_SERIES_GRIND_IV.pdf\#nameddest=tag-01U4}{IV.8.8 (p.~937)} }
\expandafter\def\csname SRef@iv:correction:P3K.27\endcsname{\href{S6_SERIES_GRIND_IV.pdf\#nameddest=tag-01U5}{IV.8.9 (p.~937)} }
\expandafter\def\csname SRef@fam:src:E.2.3\endcsname{\href{S6_SERIES_GRIND_IV.pdf\#nameddest=tag-01R8}{IV.2.2 (p.~21)} }
\expandafter\def\csname SRef@construction-and-proof\endcsname{\href{S6_SERIES_GRIND_IV.pdf\#nameddest=tag-01GZ}{IV.3.0.1 (p.~127)} }
\expandafter\def\csname SRef@the-marked-double-and-its-bordism\endcsname{\href{S6_SERIES_GRIND_IV.pdf\#nameddest=tag-034Y}{IV.{{$\langle 1\rangle $1.}} (p.~144)} }
\expandafter\def\csname SRef@removal-of-the-cubic-coordinate-by-an-extending-map\endcsname{\href{S6_SERIES_GRIND_IV.pdf\#nameddest=tag-033W}{IV.{{$\langle 1\rangle $2.}} (p.~144)} }
\expandafter\def\csname SRef@the-surgery-and-its-bounding-eight-manifold\endcsname{\href{S6_SERIES_GRIND_IV.pdf\#nameddest=tag-0352}{IV.{{$\langle 1\rangle $3.}} (p.~144)} }
\expandafter\def\csname SRef@the-invariant-and-the-resulting-homomorphism\endcsname{\href{S6_SERIES_GRIND_IV.pdf\#nameddest=tag-034W}{IV.{{$\langle 1\rangle $4.}} (p.~145)} }
\expandafter\def\csname SRef@proof\endcsname{\href{S6_SERIES_GRIND_IV.pdf\#nameddest=tag-02YT}{IV.3.0.2 (p.~129)} }
\expandafter\def\csname SRef@an-explicit-rational-elliptic-surface-and-primitive-section\endcsname{\href{S6_SERIES_GRIND_IV.pdf\#nameddest=tag-0009}{IV.3.0.3 (p.~135)} }
\expandafter\def\csname SRef@the-linearization-and-every-local-filling\endcsname{\href{S6_SERIES_GRIND_IV.pdf\#nameddest=tag-034X}{IV.3.0.4 (p.~136)} }
\expandafter\def\csname SRef@actual-integral-monodromy-not-an-abstract-substitute-tuple\endcsname{\href{S6_SERIES_GRIND_IV.pdf\#nameddest=tag-0006}{IV.3.0.5 (p.~138)} }
\expandafter\def\csname SRef@fundamental-group-and-every-integral-cohomology-differential\endcsname{\href{S6_SERIES_GRIND_IV.pdf\#nameddest=tag-01RN}{IV.3.0.6 (p.~139)} }
\expandafter\def\csname SRef@the-punctured-holomorphic-family-and-its-closing-map\endcsname{\href{S6_SERIES_GRIND_IV.pdf\#nameddest=tag-0350}{IV.3.0.7 (p.~142)} }
\expandafter\def\csname SRef@the-actual-page-and-its-section-line\endcsname{\href{S6_SERIES_GRIND_IV.pdf\#nameddest=tag-034R}{IV.3.0.8 (p.~142)} }
\expandafter\def\csname SRef@executed-verification\endcsname{\href{S6_SERIES_GRIND_IV.pdf\#nameddest=tag-01OP}{IV.3.0.9 (p.~143)} }
\expandafter\def\csname SRef@night:P3K.75\endcsname{\href{S6_SERIES_GRIND_IV.pdf\#nameddest=tag-02VN}{IV.11.1 (p.~1392)} }
\expandafter\def\csname SRef@night:P3K.117\endcsname{\href{S6_SERIES_GRIND_IV.pdf\#nameddest=tag-02TO}{IV.10.7 (p.~1213)} }
\expandafter\def\csname SRef@night:P3K.118\endcsname{\href{S6_SERIES_GRIND_IV.pdf\#nameddest=tag-02TP}{IV.10.8 (p.~1250)} }
\expandafter\def\csname SRef@night:P3K.119\endcsname{\href{S6_SERIES_GRIND_IV.pdf\#nameddest=tag-02TQ}{IV.10.9 (p.~1278)} }
\expandafter\def\csname SRef@night:P3K.109\endcsname{\href{S6_SERIES_GRIND_IV.pdf\#nameddest=tag-02TF}{IV.10.1 (p.~1116)} }
\expandafter\def\csname SRef@night:P3K.110\endcsname{\href{S6_SERIES_GRIND_IV.pdf\#nameddest=tag-02TH}{IV.10.2 (p.~1127)} }
\expandafter\def\csname SRef@night:P3K.111\endcsname{\href{S6_SERIES_GRIND_IV.pdf\#nameddest=tag-02TI}{IV.10.10 (p.~1305)} }
\expandafter\def\csname SRef@night:P3K.112\endcsname{\href{S6_SERIES_GRIND_IV.pdf\#nameddest=tag-02TJ}{IV.10.11 (p.~1369)} }
\expandafter\def\csname SRef@night:P3K.113\endcsname{\href{S6_SERIES_GRIND_IV.pdf\#nameddest=tag-02TK}{IV.10.3 (p.~1135)} }
\expandafter\def\csname SRef@night:P3K.114\endcsname{\href{S6_SERIES_GRIND_IV.pdf\#nameddest=tag-02TL}{IV.10.4 (p.~1147)} }
\expandafter\def\csname SRef@night:P3K.115\endcsname{\href{S6_SERIES_GRIND_IV.pdf\#nameddest=tag-02TM}{IV.10.5 (p.~1177)} }
\expandafter\def\csname SRef@night:P3K.116\endcsname{\href{S6_SERIES_GRIND_IV.pdf\#nameddest=tag-02TN}{IV.10.6 (p.~1193)} }
\expandafter\def\csname SRef@push:src:C.3.24\endcsname{\href{S6_SERIES_GRIND_IV.pdf\#nameddest=tag-031W}{IV.2.25 (p.~57)} }
\expandafter\def\csname SRef@maps:src:B.6.2\endcsname{\href{S6_SERIES_GRIND_IV.pdf\#nameddest=tag-02JO}{IV.2.33 (p.~67)} }
\expandafter\def\csname SRef@maps:src:B.6.20\endcsname{\href{S6_SERIES_GRIND_IV.pdf\#nameddest=tag-02JP}{IV.2.50 (p.~96)} }
\expandafter\def\csname SRef@maps:src:B.6.21\endcsname{\href{S6_SERIES_GRIND_IV.pdf\#nameddest=tag-02JQ}{IV.2.51 (p.~99)} }
\expandafter\def\csname SRef@maps:src:B.6.23\endcsname{\href{S6_SERIES_GRIND_IV.pdf\#nameddest=tag-02JS}{IV.2.53 (p.~101)} }
\expandafter\def\csname SRef@maps:src:B.6.24\endcsname{\href{S6_SERIES_GRIND_IV.pdf\#nameddest=tag-02JT}{IV.2.54 (p.~102)} }
\expandafter\def\csname SRef@night:P3F.1\endcsname{\href{S6_SERIES_GRIND_IV.pdf\#nameddest=tag-02RK}{IV.3.1 (p.~104)} }
\expandafter\def\csname SRef@night:P3F.10\endcsname{\href{S6_SERIES_GRIND_IV.pdf\#nameddest=tag-02RL}{IV.3.10 (p.~117)} }
\expandafter\def\csname SRef@night:P3F.11\endcsname{\href{S6_SERIES_GRIND_IV.pdf\#nameddest=tag-02RM}{IV.3.11 (p.~119)} }
\expandafter\def\csname SRef@night:P3F.12\endcsname{\href{S6_SERIES_GRIND_IV.pdf\#nameddest=tag-02RN}{IV.3.12 (p.~121)} }
\expandafter\def\csname SRef@night:P3F.13\endcsname{\href{S6_SERIES_GRIND_IV.pdf\#nameddest=tag-02RO}{IV.3.13 (p.~123)} }
\expandafter\def\csname SRef@night:P3F.14\endcsname{\href{S6_SERIES_GRIND_IV.pdf\#nameddest=tag-02RP}{IV.3.14 (p.~124)} }
\expandafter\def\csname SRef@night:P3F.15\endcsname{\href{S6_SERIES_GRIND_IV.pdf\#nameddest=tag-02RQ}{IV.3.15 (p.~125)} }
\expandafter\def\csname SRef@night:P3F.16\endcsname{\href{S6_SERIES_GRIND_IV.pdf\#nameddest=tag-02RR}{IV.3.16 (p.~126)} }
\expandafter\def\csname SRef@night:P3F.17\endcsname{\href{S6_SERIES_GRIND_IV.pdf\#nameddest=tag-02RS}{IV.3.17 (p.~128)} }
\expandafter\def\csname SRef@night:P3F.18\endcsname{\href{S6_SERIES_GRIND_IV.pdf\#nameddest=tag-02RT}{IV.3.18 (p.~129)} }
\expandafter\def\csname SRef@night:P3F.19\endcsname{\href{S6_SERIES_GRIND_IV.pdf\#nameddest=tag-02RU}{IV.3.19 (p.~130)} }
\expandafter\def\csname SRef@night:P3F.2\endcsname{\href{S6_SERIES_GRIND_IV.pdf\#nameddest=tag-02RV}{IV.3.2 (p.~107)} }
\expandafter\def\csname SRef@night:P3F.20\endcsname{\href{S6_SERIES_GRIND_IV.pdf\#nameddest=tag-02RW}{IV.3.21 (p.~134)} }
\expandafter\def\csname SRef@night:P3F.21\endcsname{\href{S6_SERIES_GRIND_IV.pdf\#nameddest=tag-02RX}{IV.3.22 (p.~134)} }
\expandafter\def\csname SRef@night:P3F.22\endcsname{\href{S6_SERIES_GRIND_IV.pdf\#nameddest=tag-02RY}{IV.3.23 (p.~144)} }
\expandafter\def\csname SRef@night:P3F.23\endcsname{\href{S6_SERIES_GRIND_IV.pdf\#nameddest=tag-02RZ}{IV.3.24 (p.~145)} }
\expandafter\def\csname SRef@night:P3F.24\endcsname{\href{S6_SERIES_GRIND_IV.pdf\#nameddest=tag-02S0}{IV.3.25 (p.~147)} }
\expandafter\def\csname SRef@night:P3F.30\endcsname{\href{S6_SERIES_GRIND_IV.pdf\#nameddest=tag-02S7}{IV.3.26 (p.~148)} }
\expandafter\def\csname SRef@night:P3F.35\endcsname{\href{S6_SERIES_GRIND_IV.pdf\#nameddest=tag-02SC}{IV.3.28 (p.~152)} }
\expandafter\def\csname SRef@night:P3F.36\endcsname{\href{S6_SERIES_GRIND_IV.pdf\#nameddest=tag-02SD}{IV.3.29 (p.~154)} }
\expandafter\def\csname SRef@night:P3F.37\endcsname{\href{S6_SERIES_GRIND_IV.pdf\#nameddest=tag-02SE}{IV.3.30 (p.~159)} }
\expandafter\def\csname SRef@night:P3F.3\endcsname{\href{S6_SERIES_GRIND_IV.pdf\#nameddest=tag-02S6}{IV.3.3 (p.~108)} }
\expandafter\def\csname SRef@night:P3F.38\endcsname{\href{S6_SERIES_GRIND_IV.pdf\#nameddest=tag-02SF}{IV.3.31 (p.~164)} }
\expandafter\def\csname SRef@night:P3F.39\endcsname{\href{S6_SERIES_GRIND_IV.pdf\#nameddest=tag-02SG}{IV.3.32 (p.~169)} }
\expandafter\def\csname SRef@night:P3F.40\endcsname{\href{S6_SERIES_GRIND_IV.pdf\#nameddest=tag-02SI}{IV.3.33 (p.~172)} }
\expandafter\def\csname SRef@night:P3F.41\endcsname{\href{S6_SERIES_GRIND_IV.pdf\#nameddest=tag-02SJ}{IV.3.34 (p.~175)} }
\expandafter\def\csname SRef@night:P3F.42\endcsname{\href{S6_SERIES_GRIND_IV.pdf\#nameddest=tag-02SK}{IV.3.35 (p.~180)} }
\expandafter\def\csname SRef@night:P3F.43\endcsname{\href{S6_SERIES_GRIND_IV.pdf\#nameddest=tag-02SL}{IV.3.36 (p.~186)} }
\expandafter\def\csname SRef@night:P3F.44\endcsname{\href{S6_SERIES_GRIND_IV.pdf\#nameddest=tag-02SM}{IV.3.37 (p.~188)} }
\expandafter\def\csname SRef@night:P3F.45\endcsname{\href{S6_SERIES_GRIND_IV.pdf\#nameddest=tag-02SN}{IV.3.38 (p.~191)} }
\expandafter\def\csname SRef@night:P3F.46\endcsname{\href{S6_SERIES_GRIND_IV.pdf\#nameddest=tag-02SO}{IV.3.39 (p.~195)} }
\expandafter\def\csname SRef@night:P3F.47\endcsname{\href{S6_SERIES_GRIND_IV.pdf\#nameddest=tag-02SP}{IV.3.40 (p.~202)} }
\expandafter\def\csname SRef@night:P3F.4\endcsname{\href{S6_SERIES_GRIND_IV.pdf\#nameddest=tag-02SH}{IV.3.4 (p.~109)} }
\expandafter\def\csname SRef@night:P3F.48\endcsname{\href{S6_SERIES_GRIND_IV.pdf\#nameddest=tag-02SQ}{IV.3.41 (p.~207)} }
\expandafter\def\csname SRef@night:P3F.49\endcsname{\href{S6_SERIES_GRIND_IV.pdf\#nameddest=tag-02SR}{IV.3.42 (p.~210)} }
\expandafter\def\csname SRef@night:P3F.50\endcsname{\href{S6_SERIES_GRIND_IV.pdf\#nameddest=tag-02ST}{IV.3.43 (p.~212)} }
\expandafter\def\csname SRef@night:P3F.51\endcsname{\href{S6_SERIES_GRIND_IV.pdf\#nameddest=tag-02SU}{IV.3.44 (p.~216)} }
\expandafter\def\csname SRef@night:P3F.52\endcsname{\href{S6_SERIES_GRIND_IV.pdf\#nameddest=tag-02SV}{IV.3.45 (p.~222)} }
\expandafter\def\csname SRef@night:P3F.54\endcsname{\href{S6_SERIES_GRIND_IV.pdf\#nameddest=tag-02SX}{IV.3.46 (p.~225)} }
\expandafter\def\csname SRef@night:P3F.55\endcsname{\href{S6_SERIES_GRIND_IV.pdf\#nameddest=tag-02SY}{IV.3.47 (p.~229)} }
\expandafter\def\csname SRef@night:P3F.56\endcsname{\href{S6_SERIES_GRIND_IV.pdf\#nameddest=tag-02SZ}{IV.3.48 (p.~234)} }
\expandafter\def\csname SRef@night:P3F.5\endcsname{\href{S6_SERIES_GRIND_IV.pdf\#nameddest=tag-02SS}{IV.3.5 (p.~110)} }
\expandafter\def\csname SRef@night:P3F.6\endcsname{\href{S6_SERIES_GRIND_IV.pdf\#nameddest=tag-02T0}{IV.3.6 (p.~111)} }
\expandafter\def\csname SRef@night:P3F.7\endcsname{\href{S6_SERIES_GRIND_IV.pdf\#nameddest=tag-02T1}{IV.3.7 (p.~112)} }
\expandafter\def\csname SRef@night:P3F.8\endcsname{\href{S6_SERIES_GRIND_IV.pdf\#nameddest=tag-02T2}{IV.3.8 (p.~114)} }
\expandafter\def\csname SRef@night:P3F.9\endcsname{\href{S6_SERIES_GRIND_IV.pdf\#nameddest=tag-02T3}{IV.3.9 (p.~115)} }
\expandafter\def\csname SRef@night:P3R.1\endcsname{\href{S6_SERIES_GRIND_IV.pdf\#nameddest=tag-02WE}{IV.4.1 (p.~243)} }
\expandafter\def\csname SRef@night:P3K.9\endcsname{\href{S6_SERIES_GRIND_IV.pdf\#nameddest=tag-02W3}{IV.4.10 (p.~254)} }
\expandafter\def\csname SRef@night:P3K.11\endcsname{\href{S6_SERIES_GRIND_IV.pdf\#nameddest=tag-02TG}{IV.4.11 (p.~257)} }
\expandafter\def\csname SRef@night:P3K.12\endcsname{\href{S6_SERIES_GRIND_IV.pdf\#nameddest=tag-02TR}{IV.4.12 (p.~260)} }
\expandafter\def\csname SRef@night:P3K.13\endcsname{\href{S6_SERIES_GRIND_IV.pdf\#nameddest=tag-02TS}{IV.4.13 (p.~261)} }
\expandafter\def\csname SRef@night:P3K.14\endcsname{\href{S6_SERIES_GRIND_IV.pdf\#nameddest=tag-02TT}{IV.4.14 (p.~263)} }
\expandafter\def\csname SRef@night:P3K.15\endcsname{\href{S6_SERIES_GRIND_IV.pdf\#nameddest=tag-02TU}{IV.4.15 (p.~264)} }
\expandafter\def\csname SRef@night:P3K.16\endcsname{\href{S6_SERIES_GRIND_IV.pdf\#nameddest=tag-02TV}{IV.4.16 (p.~268)} }
\expandafter\def\csname SRef@night:P3K.18\endcsname{\href{S6_SERIES_GRIND_IV.pdf\#nameddest=tag-02TX}{IV.4.17 (p.~270)} }
\expandafter\def\csname SRef@night:P3K.19\endcsname{\href{S6_SERIES_GRIND_IV.pdf\#nameddest=tag-02TY}{IV.4.18 (p.~271)} }
\expandafter\def\csname SRef@night:P3K.20\endcsname{\href{S6_SERIES_GRIND_IV.pdf\#nameddest=tag-02U0}{IV.4.19 (p.~273)} }
\expandafter\def\csname SRef@night:P3R.2\endcsname{\href{S6_SERIES_GRIND_IV.pdf\#nameddest=tag-02WF}{IV.4.2 (p.~243)} }
\expandafter\def\csname SRef@night:P3K.21\endcsname{\href{S6_SERIES_GRIND_IV.pdf\#nameddest=tag-02U1}{IV.4.20 (p.~276)} }
\expandafter\def\csname SRef@night:P3K.26\endcsname{\href{S6_SERIES_GRIND_IV.pdf\#nameddest=tag-02U6}{IV.4.21 (p.~278)} }
\expandafter\def\csname SRef@night:P3K.28\endcsname{\href{S6_SERIES_GRIND_IV.pdf\#nameddest=tag-02U8}{IV.4.22 (p.~283)} }
\expandafter\def\csname SRef@night:P3K.1\endcsname{\href{S6_SERIES_GRIND_IV.pdf\#nameddest=tag-02T4}{IV.4.3 (p.~244)} }
\expandafter\def\csname SRef@night:P3K.2\endcsname{\href{S6_SERIES_GRIND_IV.pdf\#nameddest=tag-02TZ}{IV.4.4 (p.~245)} }
\expandafter\def\csname SRef@night:P3K.3\endcsname{\href{S6_SERIES_GRIND_IV.pdf\#nameddest=tag-02UA}{IV.4.5 (p.~246)} }
\expandafter\def\csname SRef@night:P3K.4\endcsname{\href{S6_SERIES_GRIND_IV.pdf\#nameddest=tag-02UK}{IV.4.6 (p.~247)} }
\expandafter\def\csname SRef@night:P3K.5\endcsname{\href{S6_SERIES_GRIND_IV.pdf\#nameddest=tag-02UV}{IV.4.7 (p.~248)} }
\expandafter\def\csname SRef@night:P3K.7\endcsname{\href{S6_SERIES_GRIND_IV.pdf\#nameddest=tag-02VH}{IV.4.8 (p.~249)} }
\expandafter\def\csname SRef@night:P3K.8\endcsname{\href{S6_SERIES_GRIND_IV.pdf\#nameddest=tag-02VS}{IV.4.9 (p.~251)} }
\expandafter\def\csname SRef@night:P3K.38\endcsname{\href{S6_SERIES_GRIND_IV.pdf\#nameddest=tag-02UJ}{IV.5.1 (p.~291)} }
\expandafter\def\csname SRef@night:P3K.62\endcsname{\href{S6_SERIES_GRIND_IV.pdf\#nameddest=tag-02V9}{IV.5.10 (p.~333)} }
\expandafter\def\csname SRef@night:P3K.63\endcsname{\href{S6_SERIES_GRIND_IV.pdf\#nameddest=tag-02VA}{IV.5.11 (p.~340)} }
\expandafter\def\csname SRef@night:P3K.64\endcsname{\href{S6_SERIES_GRIND_IV.pdf\#nameddest=tag-02VB}{IV.5.12 (p.~348)} }
\expandafter\def\csname SRef@night:P3K.65\endcsname{\href{S6_SERIES_GRIND_IV.pdf\#nameddest=tag-02VC}{IV.5.13 (p.~352)} }
\expandafter\def\csname SRef@night:P3K.66\endcsname{\href{S6_SERIES_GRIND_IV.pdf\#nameddest=tag-02VD}{IV.5.14 (p.~362)} }
\expandafter\def\csname SRef@night:P3K.67\endcsname{\href{S6_SERIES_GRIND_IV.pdf\#nameddest=tag-02VE}{IV.5.15 (p.~369)} }
\expandafter\def\csname SRef@night:P3K.42\endcsname{\href{S6_SERIES_GRIND_IV.pdf\#nameddest=tag-02UN}{IV.5.2 (p.~298)} }
\expandafter\def\csname SRef@night:P3K.47\endcsname{\href{S6_SERIES_GRIND_IV.pdf\#nameddest=tag-02US}{IV.5.3 (p.~304)} }
\expandafter\def\csname SRef@night:P3K.48\endcsname{\href{S6_SERIES_GRIND_IV.pdf\#nameddest=tag-02UT}{IV.5.4 (p.~307)} }
\expandafter\def\csname SRef@night:P3K.49\endcsname{\href{S6_SERIES_GRIND_IV.pdf\#nameddest=tag-02UU}{IV.5.5 (p.~311)} }
\expandafter\def\csname SRef@night:P3K.51\endcsname{\href{S6_SERIES_GRIND_IV.pdf\#nameddest=tag-02UX}{IV.5.6 (p.~318)} }
\expandafter\def\csname SRef@night:P3K.53\endcsname{\href{S6_SERIES_GRIND_IV.pdf\#nameddest=tag-02UZ}{IV.5.7 (p.~321)} }
\expandafter\def\csname SRef@night:P3K.55\endcsname{\href{S6_SERIES_GRIND_IV.pdf\#nameddest=tag-02V1}{IV.5.8 (p.~325)} }
\expandafter\def\csname SRef@night:P3K.61\endcsname{\href{S6_SERIES_GRIND_IV.pdf\#nameddest=tag-02V8}{IV.5.9 (p.~330)} }
\expandafter\def\csname SRef@night:P3K.68\endcsname{\href{S6_SERIES_GRIND_IV.pdf\#nameddest=tag-02VF}{IV.6.1 (p.~374)} }
\expandafter\def\csname SRef@night:P3K.69\endcsname{\href{S6_SERIES_GRIND_IV.pdf\#nameddest=tag-02VG}{IV.6.2 (p.~380)} }
\expandafter\def\csname SRef@night:P3K.70\endcsname{\href{S6_SERIES_GRIND_IV.pdf\#nameddest=tag-02VI}{IV.6.3 (p.~384)} }
\expandafter\def\csname SRef@night:P3K.71\endcsname{\href{S6_SERIES_GRIND_IV.pdf\#nameddest=tag-02VJ}{IV.6.4 (p.~389)} }
\expandafter\def\csname SRef@night:P3K.72\endcsname{\href{S6_SERIES_GRIND_IV.pdf\#nameddest=tag-02VK}{IV.6.5 (p.~393)} }
\expandafter\def\csname SRef@night:P3K.73\endcsname{\href{S6_SERIES_GRIND_IV.pdf\#nameddest=tag-02VL}{IV.6.6 (p.~400)} }
\expandafter\def\csname SRef@night:P3K.74\endcsname{\href{S6_SERIES_GRIND_IV.pdf\#nameddest=tag-02VM}{IV.6.7 (p.~405)} }
\expandafter\def\csname SRef@night:P3K.76\endcsname{\href{S6_SERIES_GRIND_IV.pdf\#nameddest=tag-02VO}{IV.7.1 (p.~407)} }
\expandafter\def\csname SRef@night:P3K.85\endcsname{\href{S6_SERIES_GRIND_IV.pdf\#nameddest=tag-02VY}{IV.7.10 (p.~494)} }
\expandafter\def\csname SRef@night:P3K.86\endcsname{\href{S6_SERIES_GRIND_IV.pdf\#nameddest=tag-02VZ}{IV.7.11 (p.~499)} }
\expandafter\def\csname SRef@night:P3K.87\endcsname{\href{S6_SERIES_GRIND_IV.pdf\#nameddest=tag-02W0}{IV.7.12 (p.~505)} }
\expandafter\def\csname SRef@night:P3K.88\endcsname{\href{S6_SERIES_GRIND_IV.pdf\#nameddest=tag-02W1}{IV.7.13 (p.~509)} }
\expandafter\def\csname SRef@night:P3K.89\endcsname{\href{S6_SERIES_GRIND_IV.pdf\#nameddest=tag-02W2}{IV.7.14 (p.~513)} }
\expandafter\def\csname SRef@night:P3K.90\endcsname{\href{S6_SERIES_GRIND_IV.pdf\#nameddest=tag-02W4}{IV.7.15 (p.~521)} }
\expandafter\def\csname SRef@night:P3K.91\endcsname{\href{S6_SERIES_GRIND_IV.pdf\#nameddest=tag-02W5}{IV.7.16 (p.~525)} }
\expandafter\def\csname SRef@night:P3K.92\endcsname{\href{S6_SERIES_GRIND_IV.pdf\#nameddest=tag-02W6}{IV.7.17 (p.~531)} }
\expandafter\def\csname SRef@night:P3K.93\endcsname{\href{S6_SERIES_GRIND_IV.pdf\#nameddest=tag-02W7}{IV.7.18 (p.~538)} }
\expandafter\def\csname SRef@night:P3K.94\endcsname{\href{S6_SERIES_GRIND_IV.pdf\#nameddest=tag-02W8}{IV.7.19 (p.~543)} }
\expandafter\def\csname SRef@night:P3K.77\endcsname{\href{S6_SERIES_GRIND_IV.pdf\#nameddest=tag-02VP}{IV.7.2 (p.~414)} }
\expandafter\def\csname SRef@night:P3K.95\endcsname{\href{S6_SERIES_GRIND_IV.pdf\#nameddest=tag-02W9}{IV.7.20 (p.~551)} }
\expandafter\def\csname SRef@night:P3K.96\endcsname{\href{S6_SERIES_GRIND_IV.pdf\#nameddest=tag-02WA}{IV.7.21 (p.~562)} }
\expandafter\def\csname SRef@night:P3K.97\endcsname{\href{S6_SERIES_GRIND_IV.pdf\#nameddest=tag-02WB}{IV.7.22 (p.~569)} }
\expandafter\def\csname SRef@night:P3K.98\endcsname{\href{S6_SERIES_GRIND_IV.pdf\#nameddest=tag-02WC}{IV.7.23 (p.~578)} }
\expandafter\def\csname SRef@night:P3K.99\endcsname{\href{S6_SERIES_GRIND_IV.pdf\#nameddest=tag-02WD}{IV.7.24 (p.~587)} }
\expandafter\def\csname SRef@night:P3K.100\endcsname{\href{S6_SERIES_GRIND_IV.pdf\#nameddest=tag-02T6}{IV.7.25 (p.~662)} }
\expandafter\def\csname SRef@night:P3K.101\endcsname{\href{S6_SERIES_GRIND_IV.pdf\#nameddest=tag-02T7}{IV.7.26 (p.~669)} }
\expandafter\def\csname SRef@night:P3K.102\endcsname{\href{S6_SERIES_GRIND_IV.pdf\#nameddest=tag-02T8}{IV.7.27 (p.~758)} }
\expandafter\def\csname SRef@night:P3K.103\endcsname{\href{S6_SERIES_GRIND_IV.pdf\#nameddest=tag-02T9}{IV.7.28 (p.~768)} }
\expandafter\def\csname SRef@night:P3K.104\endcsname{\href{S6_SERIES_GRIND_IV.pdf\#nameddest=tag-02TA}{IV.7.29 (p.~774)} }
\expandafter\def\csname SRef@night:P3K.78\endcsname{\href{S6_SERIES_GRIND_IV.pdf\#nameddest=tag-02VQ}{IV.7.3 (p.~420)} }
\expandafter\def\csname SRef@night:P3K.105\endcsname{\href{S6_SERIES_GRIND_IV.pdf\#nameddest=tag-02TB}{IV.7.30 (p.~785)} }
\expandafter\def\csname SRef@night:P3K.106\endcsname{\href{S6_SERIES_GRIND_IV.pdf\#nameddest=tag-02TC}{IV.7.31 (p.~798)} }
\expandafter\def\csname SRef@night:P3K.108\endcsname{\href{S6_SERIES_GRIND_IV.pdf\#nameddest=tag-02TE}{IV.7.32 (p.~861)} }
\expandafter\def\csname SRef@iv:i3_all_reduced_caps\endcsname{\href{S6_SERIES_GRIND_IV.pdf\#nameddest=tag-02TE}{IV.7.32 (p.~861)} }
\expandafter\def\csname SRef@night:P3K.107\endcsname{\href{S6_SERIES_GRIND_IV.pdf\#nameddest=tag-02TD}{IV.7.33 (p.~867)} }
\expandafter\def\csname SRef@night:P3K.79\endcsname{\href{S6_SERIES_GRIND_IV.pdf\#nameddest=tag-02VR}{IV.7.4 (p.~426)} }
\expandafter\def\csname SRef@night:P3K.80\endcsname{\href{S6_SERIES_GRIND_IV.pdf\#nameddest=tag-02VT}{IV.7.5 (p.~431)} }
\expandafter\def\csname SRef@night:P3K.81\endcsname{\href{S6_SERIES_GRIND_IV.pdf\#nameddest=tag-02VU}{IV.7.6 (p.~467)} }
\expandafter\def\csname SRef@night:P3K.82\endcsname{\href{S6_SERIES_GRIND_IV.pdf\#nameddest=tag-02VV}{IV.7.7 (p.~475)} }
\expandafter\def\csname SRef@night:P3K.83\endcsname{\href{S6_SERIES_GRIND_IV.pdf\#nameddest=tag-02VW}{IV.7.8 (p.~478)} }
\expandafter\def\csname SRef@night:P3K.84\endcsname{\href{S6_SERIES_GRIND_IV.pdf\#nameddest=tag-02VX}{IV.7.9 (p.~490)} }
\expandafter\def\csname SRef@night:P3F.25\endcsname{\href{S6_SERIES_GRIND_IV.pdf\#nameddest=tag-02S1}{IV.9.1 (p.~947)} }
\expandafter\def\csname SRef@night:P3K.29\endcsname{\href{S6_SERIES_GRIND_IV.pdf\#nameddest=tag-02U9}{IV.9.10 (p.~988)} }
\expandafter\def\csname SRef@night:P3K.30\endcsname{\href{S6_SERIES_GRIND_IV.pdf\#nameddest=tag-02UB}{IV.9.11 (p.~992)} }
\expandafter\def\csname SRef@night:P3K.31\endcsname{\href{S6_SERIES_GRIND_IV.pdf\#nameddest=tag-02UC}{IV.9.12 (p.~997)} }
\expandafter\def\csname SRef@night:P3K.32\endcsname{\href{S6_SERIES_GRIND_IV.pdf\#nameddest=tag-02UD}{IV.9.13 (p.~1000)} }
\expandafter\def\csname SRef@night:P3K.33\endcsname{\href{S6_SERIES_GRIND_IV.pdf\#nameddest=tag-02UE}{IV.9.14 (p.~1010)} }
\expandafter\def\csname SRef@night:P3K.34\endcsname{\href{S6_SERIES_GRIND_IV.pdf\#nameddest=tag-02UF}{IV.9.15 (p.~1015)} }
\expandafter\def\csname SRef@night:P3F.53\endcsname{\href{S6_SERIES_GRIND_IV.pdf\#nameddest=tag-02SW}{IV.9.16 (p.~1020)} }
\expandafter\def\csname SRef@night:P3K.35\endcsname{\href{S6_SERIES_GRIND_IV.pdf\#nameddest=tag-02UG}{IV.9.17 (p.~1025)} }
\expandafter\def\csname SRef@night:P3K.37\endcsname{\href{S6_SERIES_GRIND_IV.pdf\#nameddest=tag-02UI}{IV.9.18 (p.~1032)} }
\expandafter\def\csname SRef@night:P3K.40\endcsname{\href{S6_SERIES_GRIND_IV.pdf\#nameddest=tag-02UL}{IV.9.19 (p.~1036)} }
\expandafter\def\csname SRef@night:P3K.6\endcsname{\href{S6_SERIES_GRIND_IV.pdf\#nameddest=tag-02V6}{IV.9.2 (p.~950)} }
\expandafter\def\csname SRef@night:P3K.41\endcsname{\href{S6_SERIES_GRIND_IV.pdf\#nameddest=tag-02UM}{IV.9.20 (p.~1040)} }
\expandafter\def\csname SRef@night:P3K.43\endcsname{\href{S6_SERIES_GRIND_IV.pdf\#nameddest=tag-02UO}{IV.9.21 (p.~1047)} }
\expandafter\def\csname SRef@night:P3K.44\endcsname{\href{S6_SERIES_GRIND_IV.pdf\#nameddest=tag-02UP}{IV.9.22 (p.~1050)} }
\expandafter\def\csname SRef@night:P3K.45\endcsname{\href{S6_SERIES_GRIND_IV.pdf\#nameddest=tag-02UQ}{IV.9.23 (p.~1055)} }
\expandafter\def\csname SRef@night:P3K.46\endcsname{\href{S6_SERIES_GRIND_IV.pdf\#nameddest=tag-02UR}{IV.9.24 (p.~1058)} }
\expandafter\def\csname SRef@night:P3K.50\endcsname{\href{S6_SERIES_GRIND_IV.pdf\#nameddest=tag-02UW}{IV.9.25 (p.~1061)} }
\expandafter\def\csname SRef@night:P3K.52\endcsname{\href{S6_SERIES_GRIND_IV.pdf\#nameddest=tag-02UY}{IV.9.26 (p.~1064)} }
\expandafter\def\csname SRef@night:P3K.54\endcsname{\href{S6_SERIES_GRIND_IV.pdf\#nameddest=tag-02V0}{IV.9.27 (p.~1067)} }
\expandafter\def\csname SRef@night:P3K.56\endcsname{\href{S6_SERIES_GRIND_IV.pdf\#nameddest=tag-02V2}{IV.9.28 (p.~1069)} }
\expandafter\def\csname SRef@night:P3K.57\endcsname{\href{S6_SERIES_GRIND_IV.pdf\#nameddest=tag-02V3}{IV.9.29 (p.~1072)} }
\expandafter\def\csname SRef@night:P3K.10\endcsname{\href{S6_SERIES_GRIND_IV.pdf\#nameddest=tag-02T5}{IV.9.3 (p.~952)} }
\expandafter\def\csname SRef@night:P3K.58\endcsname{\href{S6_SERIES_GRIND_IV.pdf\#nameddest=tag-02V4}{IV.9.30 (p.~1095)} }
\expandafter\def\csname SRef@night:P3K.59\endcsname{\href{S6_SERIES_GRIND_IV.pdf\#nameddest=tag-02V5}{IV.9.31 (p.~1104)} }
\expandafter\def\csname SRef@night:P3K.60\endcsname{\href{S6_SERIES_GRIND_IV.pdf\#nameddest=tag-02V7}{IV.9.32 (p.~1109)} }
\expandafter\def\csname SRef@night:P3K.17\endcsname{\href{S6_SERIES_GRIND_IV.pdf\#nameddest=tag-02TW}{IV.9.4 (p.~961)} }
\expandafter\def\csname SRef@night:P3K.22\endcsname{\href{S6_SERIES_GRIND_IV.pdf\#nameddest=tag-02U2}{IV.9.5 (p.~965)} }
\expandafter\def\csname SRef@night:P3K.23\endcsname{\href{S6_SERIES_GRIND_IV.pdf\#nameddest=tag-02U3}{IV.9.6 (p.~971)} }
\expandafter\def\csname SRef@night:P3K.24\endcsname{\href{S6_SERIES_GRIND_IV.pdf\#nameddest=tag-02U4}{IV.9.7 (p.~975)} }
\expandafter\def\csname SRef@night:P3K.25\endcsname{\href{S6_SERIES_GRIND_IV.pdf\#nameddest=tag-02U5}{IV.9.8 (p.~980)} }
\expandafter\def\csname SRef@night:P3K.27\endcsname{\href{S6_SERIES_GRIND_IV.pdf\#nameddest=tag-02U7}{IV.9.9 (p.~983)} }
\expandafter\def\csname SRef@maps:eq:F.6\endcsname{\href{S6_SERIES_GRIND_V.pdf\#nameddest=tag-02FA}{V.{F.6} (p.~85)} }
\expandafter\def\csname SRef@maps:eq:F.L2\endcsname{\href{S6_SERIES_GRIND_V.pdf\#nameddest=tag-02GV}{V.{F.L2} (p.~174)} }
\expandafter\def\csname SRef@maps:eq:F.L3\endcsname{\href{S6_SERIES_GRIND_V.pdf\#nameddest=tag-02GW}{V.{F.L3} (p.~174)} }
\expandafter\def\csname SRef@maps:eq:F.L4\endcsname{\href{S6_SERIES_GRIND_V.pdf\#nameddest=tag-02GX}{V.{F.L4} (p.~175)} }
\expandafter\def\csname SRef@maps:eq:F.L5\endcsname{\href{S6_SERIES_GRIND_V.pdf\#nameddest=tag-02GY}{V.{F.L5} (p.~175)} }
\expandafter\def\csname SRef@maps:eq:F.L6\endcsname{\href{S6_SERIES_GRIND_V.pdf\#nameddest=tag-02GZ}{V.{F.L6} (p.~175)} }
\expandafter\def\csname SRef@maps:eq:F.L7\endcsname{\href{S6_SERIES_GRIND_V.pdf\#nameddest=tag-02H0}{V.{F.L7} (p.~175)} }
\expandafter\def\csname SRef@maps:eq:F.L8\endcsname{\href{S6_SERIES_GRIND_V.pdf\#nameddest=tag-02H1}{V.{F.L8} (p.~175)} }
\expandafter\def\csname SRef@maps:eq:F.L9\endcsname{\href{S6_SERIES_GRIND_V.pdf\#nameddest=tag-02H2}{V.{F.L9} (p.~176)} }
\expandafter\def\csname SRef@maps:eq:F.L10\endcsname{\href{S6_SERIES_GRIND_V.pdf\#nameddest=tag-02GU}{V.{F.L10} (p.~176)} }
\expandafter\def\csname SRef@maps:eq:F.T1\endcsname{\href{S6_SERIES_GRIND_V.pdf\#nameddest=tag-02I7}{V.{F.T1} (p.~177)} }
\expandafter\def\csname SRef@maps:eq:F.T2\endcsname{\href{S6_SERIES_GRIND_V.pdf\#nameddest=tag-02IC}{V.{F.T2} (p.~177)} }
\expandafter\def\csname SRef@maps:eq:F.T3\endcsname{\href{S6_SERIES_GRIND_V.pdf\#nameddest=tag-02ID}{V.{F.T3} (p.~177)} }
\expandafter\def\csname SRef@maps:eq:F.T4\endcsname{\href{S6_SERIES_GRIND_V.pdf\#nameddest=tag-02IE}{V.{F.T4} (p.~178)} }
\expandafter\def\csname SRef@maps:eq:F.T5\endcsname{\href{S6_SERIES_GRIND_V.pdf\#nameddest=tag-02IF}{V.{F.T5} (p.~178)} }
\expandafter\def\csname SRef@maps:eq:F.T6\endcsname{\href{S6_SERIES_GRIND_V.pdf\#nameddest=tag-02IG}{V.{F.T6} (p.~178)} }
\expandafter\def\csname SRef@maps:eq:F.T7\endcsname{\href{S6_SERIES_GRIND_V.pdf\#nameddest=tag-02IH}{V.{F.T7} (p.~178)} }
\expandafter\def\csname SRef@maps:eq:F.T8\endcsname{\href{S6_SERIES_GRIND_V.pdf\#nameddest=tag-02II}{V.{F.T8} (p.~178)} }
\expandafter\def\csname SRef@maps:eq:F.T9\endcsname{\href{S6_SERIES_GRIND_V.pdf\#nameddest=tag-02IJ}{V.{F.T9} (p.~179)} }
\expandafter\def\csname SRef@maps:eq:F.T10\endcsname{\href{S6_SERIES_GRIND_V.pdf\#nameddest=tag-02I8}{V.{F.T10} (p.~179)} }
\expandafter\def\csname SRef@maps:eq:F.S1\endcsname{\href{S6_SERIES_GRIND_V.pdf\#nameddest=tag-02HX}{V.{F.S1} (p.~180)} }
\expandafter\def\csname SRef@maps:eq:F.S2\endcsname{\href{S6_SERIES_GRIND_V.pdf\#nameddest=tag-02HY}{V.{F.S2} (p.~180)} }
\expandafter\def\csname SRef@maps:eq:F.S3\endcsname{\href{S6_SERIES_GRIND_V.pdf\#nameddest=tag-02HZ}{V.{F.S3} (p.~180)} }
\expandafter\def\csname SRef@maps:eq:F.S4\endcsname{\href{S6_SERIES_GRIND_V.pdf\#nameddest=tag-02I0}{V.{F.S4} (p.~180)} }
\expandafter\def\csname SRef@maps:eq:F.S5\endcsname{\href{S6_SERIES_GRIND_V.pdf\#nameddest=tag-02I1}{V.{F.S5} (p.~180)} }
\expandafter\def\csname SRef@maps:eq:F.S6\endcsname{\href{S6_SERIES_GRIND_V.pdf\#nameddest=tag-02I2}{V.{F.S6} (p.~181)} }
\expandafter\def\csname SRef@maps:eq:F.S7\endcsname{\href{S6_SERIES_GRIND_V.pdf\#nameddest=tag-02I3}{V.{F.S7} (p.~181)} }
\expandafter\def\csname SRef@maps:eq:F.S8\endcsname{\href{S6_SERIES_GRIND_V.pdf\#nameddest=tag-02I4}{V.{F.S8} (p.~181)} }
\expandafter\def\csname SRef@maps:eq:F.S9\endcsname{\href{S6_SERIES_GRIND_V.pdf\#nameddest=tag-02I5}{V.{F.S9} (p.~181)} }
\expandafter\def\csname SRef@maps:eq:F.R10\endcsname{\href{S6_SERIES_GRIND_V.pdf\#nameddest=tag-02HK}{V.{F.R10} (p.~182)} }
\expandafter\def\csname SRef@maps:eq:F.R11\endcsname{\href{S6_SERIES_GRIND_V.pdf\#nameddest=tag-02HL}{V.{F.R11} (p.~182)} }
\expandafter\def\csname SRef@maps:eq:F.R12\endcsname{\href{S6_SERIES_GRIND_V.pdf\#nameddest=tag-02HM}{V.{F.R12} (p.~182)} }
\expandafter\def\csname SRef@maps:eq:F.R13\endcsname{\href{S6_SERIES_GRIND_V.pdf\#nameddest=tag-02HN}{V.{F.R13} (p.~182)} }
\expandafter\def\csname SRef@maps:eq:F.R14\endcsname{\href{S6_SERIES_GRIND_V.pdf\#nameddest=tag-02HO}{V.{F.R14} (p.~182)} }
\expandafter\def\csname SRef@maps:eq:F.X1\endcsname{\href{S6_SERIES_GRIND_V.pdf\#nameddest=tag-02IV}{V.{F.X1} (p.~183)} }
\expandafter\def\csname SRef@maps:eq:F.X2\endcsname{\href{S6_SERIES_GRIND_V.pdf\#nameddest=tag-02IW}{V.{F.X2} (p.~183)} }
\expandafter\def\csname SRef@maps:eq:F.X3\endcsname{\href{S6_SERIES_GRIND_V.pdf\#nameddest=tag-02IX}{V.{F.X3} (p.~184)} }
\expandafter\def\csname SRef@maps:eq:F.X4\endcsname{\href{S6_SERIES_GRIND_V.pdf\#nameddest=tag-02IY}{V.{F.X4} (p.~184)} }
\expandafter\def\csname SRef@maps:eq:F.X5\endcsname{\href{S6_SERIES_GRIND_V.pdf\#nameddest=tag-02IZ}{V.{F.X5} (p.~184)} }
\expandafter\def\csname SRef@maps:eq:F.T11\endcsname{\href{S6_SERIES_GRIND_V.pdf\#nameddest=tag-02I9}{V.{F.T11} (p.~185)} }
\expandafter\def\csname SRef@maps:eq:F.T12\endcsname{\href{S6_SERIES_GRIND_V.pdf\#nameddest=tag-02IA}{V.{F.T12} (p.~185)} }
\expandafter\def\csname SRef@maps:eq:F.T13\endcsname{\href{S6_SERIES_GRIND_V.pdf\#nameddest=tag-02IB}{V.{F.T13} (p.~185)} }
\expandafter\def\csname SRef@maps:eq:F.HC1\endcsname{\href{S6_SERIES_GRIND_V.pdf\#nameddest=tag-02G9}{V.{F.HC1} (p.~186)} }
\expandafter\def\csname SRef@maps:eq:F.HC2\endcsname{\href{S6_SERIES_GRIND_V.pdf\#nameddest=tag-02GE}{V.{F.HC2} (p.~186)} }
\expandafter\def\csname SRef@maps:eq:F.HC3\endcsname{\href{S6_SERIES_GRIND_V.pdf\#nameddest=tag-02GF}{V.{F.HC3} (p.~186)} }
\expandafter\def\csname SRef@maps:eq:F.HC4\endcsname{\href{S6_SERIES_GRIND_V.pdf\#nameddest=tag-02GG}{V.{F.HC4} (p.~187)} }
\expandafter\def\csname SRef@maps:eq:F.HC5\endcsname{\href{S6_SERIES_GRIND_V.pdf\#nameddest=tag-02GH}{V.{F.HC5} (p.~187)} }
\expandafter\def\csname SRef@maps:eq:F.HC6\endcsname{\href{S6_SERIES_GRIND_V.pdf\#nameddest=tag-02GI}{V.{F.HC6} (p.~187)} }
\expandafter\def\csname SRef@maps:eq:F.HC7\endcsname{\href{S6_SERIES_GRIND_V.pdf\#nameddest=tag-02GJ}{V.{F.HC7} (p.~188)} }
\expandafter\def\csname SRef@maps:eq:F.HC8\endcsname{\href{S6_SERIES_GRIND_V.pdf\#nameddest=tag-02GK}{V.{F.HC8} (p.~188)} }
\expandafter\def\csname SRef@maps:eq:F.HC9\endcsname{\href{S6_SERIES_GRIND_V.pdf\#nameddest=tag-02GL}{V.{F.HC9} (p.~189)} }
\expandafter\def\csname SRef@maps:eq:F.HC10\endcsname{\href{S6_SERIES_GRIND_V.pdf\#nameddest=tag-02GA}{V.{F.HC10} (p.~189)} }
\expandafter\def\csname SRef@maps:eq:F.HC11\endcsname{\href{S6_SERIES_GRIND_V.pdf\#nameddest=tag-02GB}{V.{F.HC11} (p.~189)} }
\expandafter\def\csname SRef@maps:eq:F.HC12\endcsname{\href{S6_SERIES_GRIND_V.pdf\#nameddest=tag-02GC}{V.{F.HC12} (p.~189)} }
\expandafter\def\csname SRef@maps:eq:F.HC13\endcsname{\href{S6_SERIES_GRIND_V.pdf\#nameddest=tag-02GD}{V.{F.HC13} (p.~190)} }
\expandafter\def\csname SRef@maps:eq:F.G1\endcsname{\href{S6_SERIES_GRIND_V.pdf\#nameddest=tag-02FO}{V.{F.G1} (p.~191)} }
\expandafter\def\csname SRef@maps:eq:F.G2\endcsname{\href{S6_SERIES_GRIND_V.pdf\#nameddest=tag-02FR}{V.{F.G2} (p.~191)} }
\expandafter\def\csname SRef@maps:eq:F.G3\endcsname{\href{S6_SERIES_GRIND_V.pdf\#nameddest=tag-02FS}{V.{F.G3} (p.~191)} }
\expandafter\def\csname SRef@maps:eq:F.G4\endcsname{\href{S6_SERIES_GRIND_V.pdf\#nameddest=tag-02FT}{V.{F.G4} (p.~191)} }
\expandafter\def\csname SRef@maps:eq:F.G5\endcsname{\href{S6_SERIES_GRIND_V.pdf\#nameddest=tag-02FU}{V.{F.G5} (p.~191)} }
\expandafter\def\csname SRef@maps:eq:F.G6\endcsname{\href{S6_SERIES_GRIND_V.pdf\#nameddest=tag-02FV}{V.{F.G6} (p.~191)} }
\expandafter\def\csname SRef@maps:eq:F.G7\endcsname{\href{S6_SERIES_GRIND_V.pdf\#nameddest=tag-02FW}{V.{F.G7} (p.~192)} }
\expandafter\def\csname SRef@maps:eq:F.G8\endcsname{\href{S6_SERIES_GRIND_V.pdf\#nameddest=tag-02FX}{V.{F.G8} (p.~192)} }
\expandafter\def\csname SRef@maps:eq:F.G9\endcsname{\href{S6_SERIES_GRIND_V.pdf\#nameddest=tag-02FY}{V.{F.G9} (p.~192)} }
\expandafter\def\csname SRef@maps:eq:F.G10\endcsname{\href{S6_SERIES_GRIND_V.pdf\#nameddest=tag-02FP}{V.{F.G10} (p.~192)} }
\expandafter\def\csname SRef@maps:eq:F.G11\endcsname{\href{S6_SERIES_GRIND_V.pdf\#nameddest=tag-02FQ}{V.{F.G11} (p.~194)} }
\expandafter\def\csname SRef@maps:eq:F.LT1\endcsname{\href{S6_SERIES_GRIND_V.pdf\#nameddest=tag-02H3}{V.{F.LT1} (p.~195)} }
\expandafter\def\csname SRef@maps:eq:F.LT2\endcsname{\href{S6_SERIES_GRIND_V.pdf\#nameddest=tag-02H4}{V.{F.LT2} (p.~195)} }
\expandafter\def\csname SRef@maps:eq:F.LT3\endcsname{\href{S6_SERIES_GRIND_V.pdf\#nameddest=tag-02H5}{V.{F.LT3} (p.~195)} }
\expandafter\def\csname SRef@maps:eq:F.LT4\endcsname{\href{S6_SERIES_GRIND_V.pdf\#nameddest=tag-02H6}{V.{F.LT4} (p.~195)} }
\expandafter\def\csname SRef@maps:eq:F.LT5\endcsname{\href{S6_SERIES_GRIND_V.pdf\#nameddest=tag-02H7}{V.{F.LT5} (p.~195)} }
\expandafter\def\csname SRef@maps:eq:F.LT6\endcsname{\href{S6_SERIES_GRIND_V.pdf\#nameddest=tag-02H8}{V.{F.LT6} (p.~196)} }
\expandafter\def\csname SRef@maps:eq:F.LT7\endcsname{\href{S6_SERIES_GRIND_V.pdf\#nameddest=tag-02H9}{V.{F.LT7} (p.~196)} }
\expandafter\def\csname SRef@maps:eq:F.LT8\endcsname{\href{S6_SERIES_GRIND_V.pdf\#nameddest=tag-02HA}{V.{F.LT8} (p.~196)} }
\expandafter\def\csname SRef@maps:eq:F.LT9\endcsname{\href{S6_SERIES_GRIND_V.pdf\#nameddest=tag-02HB}{V.{F.LT9} (p.~197)} }
\expandafter\def\csname SRef@maps:eq:F.CA1\endcsname{\href{S6_SERIES_GRIND_V.pdf\#nameddest=tag-02FE}{V.{F.CA1} (p.~197)} }
\expandafter\def\csname SRef@maps:eq:F.CA2\endcsname{\href{S6_SERIES_GRIND_V.pdf\#nameddest=tag-02FF}{V.{F.CA2} (p.~198)} }
\expandafter\def\csname SRef@maps:eq:F.CA3\endcsname{\href{S6_SERIES_GRIND_V.pdf\#nameddest=tag-02FG}{V.{F.CA3} (p.~198)} }
\expandafter\def\csname SRef@maps:eq:F.CA4\endcsname{\href{S6_SERIES_GRIND_V.pdf\#nameddest=tag-02FH}{V.{F.CA4} (p.~198)} }
\expandafter\def\csname SRef@maps:eq:F.CA5\endcsname{\href{S6_SERIES_GRIND_V.pdf\#nameddest=tag-02FI}{V.{F.CA5} (p.~198)} }
\expandafter\def\csname SRef@maps:eq:F.CK1\endcsname{\href{S6_SERIES_GRIND_V.pdf\#nameddest=tag-02FJ}{V.{F.CK1} (p.~200)} }
\expandafter\def\csname SRef@maps:eq:F.CK2\endcsname{\href{S6_SERIES_GRIND_V.pdf\#nameddest=tag-02FK}{V.{F.CK2} (p.~201)} }
\expandafter\def\csname SRef@maps:eq:F.GC1\endcsname{\href{S6_SERIES_GRIND_V.pdf\#nameddest=tag-02FZ}{V.{F.GC1} (p.~202)} }
\expandafter\def\csname SRef@maps:eq:F.GC2\endcsname{\href{S6_SERIES_GRIND_V.pdf\#nameddest=tag-02G0}{V.{F.GC2} (p.~203)} }
\expandafter\def\csname SRef@maps:eq:F.GC3\endcsname{\href{S6_SERIES_GRIND_V.pdf\#nameddest=tag-02G1}{V.{F.GC3} (p.~203)} }
\expandafter\def\csname SRef@maps:eq:F.GC4\endcsname{\href{S6_SERIES_GRIND_V.pdf\#nameddest=tag-02G2}{V.{F.GC4} (p.~203)} }
\expandafter\def\csname SRef@maps:eq:F.GC5\endcsname{\href{S6_SERIES_GRIND_V.pdf\#nameddest=tag-02G3}{V.{F.GC5} (p.~204)} }
\expandafter\def\csname SRef@maps:eq:F.GC5a\endcsname{\href{S6_SERIES_GRIND_V.pdf\#nameddest=tag-02G4}{V.{F.GC5a} (p.~204)} }
\expandafter\def\csname SRef@maps:eq:F.GC6\endcsname{\href{S6_SERIES_GRIND_V.pdf\#nameddest=tag-02G5}{V.{F.GC6} (p.~204)} }
\expandafter\def\csname SRef@maps:eq:F.GC7\endcsname{\href{S6_SERIES_GRIND_V.pdf\#nameddest=tag-02G6}{V.{F.GC7} (p.~205)} }
\expandafter\def\csname SRef@maps:eq:F.GC8\endcsname{\href{S6_SERIES_GRIND_V.pdf\#nameddest=tag-02G7}{V.{F.GC8} (p.~205)} }
\expandafter\def\csname SRef@maps:eq:F.GC9\endcsname{\href{S6_SERIES_GRIND_V.pdf\#nameddest=tag-02G8}{V.{F.GC9} (p.~205)} }
\expandafter\def\csname SRef@maps:eq:F.IC1\endcsname{\href{S6_SERIES_GRIND_V.pdf\#nameddest=tag-02GM}{V.{F.IC1} (p.~207)} }
\expandafter\def\csname SRef@maps:eq:F.IC2\endcsname{\href{S6_SERIES_GRIND_V.pdf\#nameddest=tag-02GN}{V.{F.IC2} (p.~207)} }
\expandafter\def\csname SRef@maps:eq:F.IC3\endcsname{\href{S6_SERIES_GRIND_V.pdf\#nameddest=tag-02GO}{V.{F.IC3} (p.~207)} }
\expandafter\def\csname SRef@maps:eq:F.CN1\endcsname{\href{S6_SERIES_GRIND_V.pdf\#nameddest=tag-02FL}{V.{F.CN1} (p.~209)} }
\expandafter\def\csname SRef@maps:eq:F.CN2\endcsname{\href{S6_SERIES_GRIND_V.pdf\#nameddest=tag-02FM}{V.{F.CN2} (p.~209)} }
\expandafter\def\csname SRef@maps:eq:F.CN3\endcsname{\href{S6_SERIES_GRIND_V.pdf\#nameddest=tag-02FN}{V.{F.CN3} (p.~210)} }
\expandafter\def\csname SRef@maps:eq:F.IV1\endcsname{\href{S6_SERIES_GRIND_V.pdf\#nameddest=tag-02GP}{V.{F.IV1} (p.~210)} }
\expandafter\def\csname SRef@maps:eq:F.IV2\endcsname{\href{S6_SERIES_GRIND_V.pdf\#nameddest=tag-02GQ}{V.{F.IV2} (p.~210)} }
\expandafter\def\csname SRef@maps:eq:F.IV3\endcsname{\href{S6_SERIES_GRIND_V.pdf\#nameddest=tag-02GR}{V.{F.IV3} (p.~211)} }
\expandafter\def\csname SRef@maps:eq:F.IV4\endcsname{\href{S6_SERIES_GRIND_V.pdf\#nameddest=tag-02GS}{V.{F.IV4} (p.~211)} }
\expandafter\def\csname SRef@maps:eq:F.SY1\endcsname{\href{S6_SERIES_GRIND_V.pdf\#nameddest=tag-02I6}{V.{F.SY1} (p.~212)} }
\expandafter\def\csname SRef@maps:eq:F.O1\endcsname{\href{S6_SERIES_GRIND_V.pdf\#nameddest=tag-02HC}{V.{F.O1} (p.~216)} }
\expandafter\def\csname SRef@maps:eq:F.O2\endcsname{\href{S6_SERIES_GRIND_V.pdf\#nameddest=tag-02HD}{V.{F.O2} (p.~216)} }
\expandafter\def\csname SRef@maps:eq:F.O3\endcsname{\href{S6_SERIES_GRIND_V.pdf\#nameddest=tag-02HE}{V.{F.O3} (p.~216)} }
\expandafter\def\csname SRef@maps:eq:F.O4\endcsname{\href{S6_SERIES_GRIND_V.pdf\#nameddest=tag-02HF}{V.{F.O4} (p.~217)} }
\expandafter\def\csname SRef@maps:eq:F.O5\endcsname{\href{S6_SERIES_GRIND_V.pdf\#nameddest=tag-02HG}{V.{F.O5} (p.~218)} }
\expandafter\def\csname SRef@maps:eq:F.O6\endcsname{\href{S6_SERIES_GRIND_V.pdf\#nameddest=tag-02HH}{V.{F.O6} (p.~219)} }
\expandafter\def\csname SRef@maps:eq:F.O7\endcsname{\href{S6_SERIES_GRIND_V.pdf\#nameddest=tag-02HI}{V.{F.O7} (p.~220)} }
\expandafter\def\csname SRef@maps:eq:F.7\endcsname{\href{S6_SERIES_GRIND_V.pdf\#nameddest=tag-02FB}{V.{F.7} (p.~104)} }
\expandafter\def\csname SRef@maps:eq:F.8\endcsname{\href{S6_SERIES_GRIND_V.pdf\#nameddest=tag-02FC}{V.{F.8} (p.~104)} }
\expandafter\def\csname SRef@maps:eq:F.9\endcsname{\href{S6_SERIES_GRIND_V.pdf\#nameddest=tag-02FD}{V.{F.9} (p.~105)} }
\expandafter\def\csname SRef@maps:eq:F.10\endcsname{\href{S6_SERIES_GRIND_V.pdf\#nameddest=tag-02EG}{V.{F.10} (p.~105)} }
\expandafter\def\csname SRef@maps:eq:F.11\endcsname{\href{S6_SERIES_GRIND_V.pdf\#nameddest=tag-02EH}{V.{F.11} (p.~106)} }
\expandafter\def\csname SRef@maps:eq:F.12\endcsname{\href{S6_SERIES_GRIND_V.pdf\#nameddest=tag-02EI}{V.{F.12} (p.~106)} }
\expandafter\def\csname SRef@maps:eq:F.13\endcsname{\href{S6_SERIES_GRIND_V.pdf\#nameddest=tag-02EJ}{V.{F.13} (p.~110)} }
\expandafter\def\csname SRef@maps:eq:F.14\endcsname{\href{S6_SERIES_GRIND_V.pdf\#nameddest=tag-02EK}{V.{F.14} (p.~110)} }
\expandafter\def\csname SRef@maps:eq:F.15\endcsname{\href{S6_SERIES_GRIND_V.pdf\#nameddest=tag-02EL}{V.{F.15} (p.~110)} }
\expandafter\def\csname SRef@maps:eq:F.16\endcsname{\href{S6_SERIES_GRIND_V.pdf\#nameddest=tag-02EM}{V.{F.16} (p.~110)} }
\expandafter\def\csname SRef@maps:eq:F.17\endcsname{\href{S6_SERIES_GRIND_V.pdf\#nameddest=tag-02EN}{V.{F.17} (p.~110)} }
\expandafter\def\csname SRef@maps:eq:F.18\endcsname{\href{S6_SERIES_GRIND_V.pdf\#nameddest=tag-02EO}{V.{F.18} (p.~110)} }
\expandafter\def\csname SRef@maps:eq:F.19\endcsname{\href{S6_SERIES_GRIND_V.pdf\#nameddest=tag-02EP}{V.{F.19} (p.~111)} }
\expandafter\def\csname SRef@maps:eq:F.20\endcsname{\href{S6_SERIES_GRIND_V.pdf\#nameddest=tag-02ER}{V.{F.20} (p.~111)} }
\expandafter\def\csname SRef@maps:eq:F.21\endcsname{\href{S6_SERIES_GRIND_V.pdf\#nameddest=tag-02ES}{V.{F.21} (p.~111)} }
\expandafter\def\csname SRef@maps:eq:F.22\endcsname{\href{S6_SERIES_GRIND_V.pdf\#nameddest=tag-02ET}{V.{F.22} (p.~111)} }
\expandafter\def\csname SRef@maps:eq:F.23\endcsname{\href{S6_SERIES_GRIND_V.pdf\#nameddest=tag-02EU}{V.{F.23} (p.~111)} }
\expandafter\def\csname SRef@maps:eq:F.24\endcsname{\href{S6_SERIES_GRIND_V.pdf\#nameddest=tag-02EV}{V.{F.24} (p.~115)} }
\expandafter\def\csname SRef@maps:eq:F.25\endcsname{\href{S6_SERIES_GRIND_V.pdf\#nameddest=tag-02EW}{V.{F.25} (p.~117)} }
\expandafter\def\csname SRef@maps:eq:F.3\endcsname{\href{S6_SERIES_GRIND_V.pdf\#nameddest=tag-02EZ}{V.{F.3} (p.~73)} }
\expandafter\def\csname SRef@maps:eq:F.27\endcsname{\href{S6_SERIES_GRIND_V.pdf\#nameddest=tag-02EX}{V.{F.27} (p.~147)} }
\expandafter\def\csname SRef@maps:eq:F.28\endcsname{\href{S6_SERIES_GRIND_V.pdf\#nameddest=tag-02EY}{V.{F.28} (p.~147)} }
\expandafter\def\csname SRef@maps:eq:F.4\endcsname{\href{S6_SERIES_GRIND_V.pdf\#nameddest=tag-02F7}{V.{F.4} (p.~75)} }
\expandafter\def\csname SRef@maps:eq:F.30\endcsname{\href{S6_SERIES_GRIND_V.pdf\#nameddest=tag-02F0}{V.{F.30} (p.~150)} }
\expandafter\def\csname SRef@maps:eq:F.31\endcsname{\href{S6_SERIES_GRIND_V.pdf\#nameddest=tag-02F1}{V.{F.31} (p.~150)} }
\expandafter\def\csname SRef@maps:eq:F.32\endcsname{\href{S6_SERIES_GRIND_V.pdf\#nameddest=tag-02F2}{V.{F.32} (p.~151)} }
\expandafter\def\csname SRef@maps:eq:F.33\endcsname{\href{S6_SERIES_GRIND_V.pdf\#nameddest=tag-02F3}{V.{F.33} (p.~152)} }
\expandafter\def\csname SRef@maps:eq:F.34\endcsname{\href{S6_SERIES_GRIND_V.pdf\#nameddest=tag-02F4}{V.{F.34} (p.~152)} }
\expandafter\def\csname SRef@maps:eq:F.36\endcsname{\href{S6_SERIES_GRIND_V.pdf\#nameddest=tag-02F5}{V.{F.36} (p.~155)} }
\expandafter\def\csname SRef@maps:eq:F.39\endcsname{\href{S6_SERIES_GRIND_V.pdf\#nameddest=tag-02F6}{V.{F.39} (p.~161)} }
\expandafter\def\csname SRef@maps:eq:F.40\endcsname{\href{S6_SERIES_GRIND_V.pdf\#nameddest=tag-02F8}{V.{F.40} (p.~162)} }
\expandafter\def\csname SRef@maps:eq:F.U1\endcsname{\href{S6_SERIES_GRIND_V.pdf\#nameddest=tag-02IK}{V.{F.U1} (p.~163)} }
\expandafter\def\csname SRef@maps:eq:F.5\endcsname{\href{S6_SERIES_GRIND_V.pdf\#nameddest=tag-02F9}{V.{F.5} (p.~76)} }
\expandafter\def\csname SRef@maps:eq:F.U2\endcsname{\href{S6_SERIES_GRIND_V.pdf\#nameddest=tag-02IO}{V.{F.U2} (p.~163)} }
\expandafter\def\csname SRef@maps:eq:F.U3\endcsname{\href{S6_SERIES_GRIND_V.pdf\#nameddest=tag-02IP}{V.{F.U3} (p.~163)} }
\expandafter\def\csname SRef@maps:eq:F.U4\endcsname{\href{S6_SERIES_GRIND_V.pdf\#nameddest=tag-02IQ}{V.{F.U4} (p.~164)} }
\expandafter\def\csname SRef@maps:eq:F.U5\endcsname{\href{S6_SERIES_GRIND_V.pdf\#nameddest=tag-02IR}{V.{F.U5} (p.~164)} }
\expandafter\def\csname SRef@maps:eq:F.U6\endcsname{\href{S6_SERIES_GRIND_V.pdf\#nameddest=tag-02IS}{V.{F.U6} (p.~164)} }
\expandafter\def\csname SRef@maps:eq:F.U7\endcsname{\href{S6_SERIES_GRIND_V.pdf\#nameddest=tag-02IT}{V.{F.U7} (p.~166)} }
\expandafter\def\csname SRef@maps:eq:F.U9\endcsname{\href{S6_SERIES_GRIND_V.pdf\#nameddest=tag-02IU}{V.{F.U9} (p.~166)} }
\expandafter\def\csname SRef@maps:eq:F.U10\endcsname{\href{S6_SERIES_GRIND_V.pdf\#nameddest=tag-02IL}{V.{F.U10} (p.~166)} }
\expandafter\def\csname SRef@maps:eq:F.U11\endcsname{\href{S6_SERIES_GRIND_V.pdf\#nameddest=tag-02IM}{V.{F.U11} (p.~167)} }
\expandafter\def\csname SRef@maps:eq:F.U12\endcsname{\href{S6_SERIES_GRIND_V.pdf\#nameddest=tag-02IN}{V.{F.U12} (p.~169)} }
\expandafter\def\csname SRef@maps:eq:FT.1\endcsname{\href{S6_SERIES_GRIND_V.pdf\#nameddest=tag-02J0}{V.{FT.1} (p.~81)} }
\expandafter\def\csname SRef@maps:eq:F.R1\endcsname{\href{S6_SERIES_GRIND_V.pdf\#nameddest=tag-02HJ}{V.{F.R1} (p.~170)} }
\expandafter\def\csname SRef@maps:eq:F.R2\endcsname{\href{S6_SERIES_GRIND_V.pdf\#nameddest=tag-02HP}{V.{F.R2} (p.~170)} }
\expandafter\def\csname SRef@maps:eq:F.R3\endcsname{\href{S6_SERIES_GRIND_V.pdf\#nameddest=tag-02HQ}{V.{F.R3} (p.~171)} }
\expandafter\def\csname SRef@maps:eq:F.R4\endcsname{\href{S6_SERIES_GRIND_V.pdf\#nameddest=tag-02HR}{V.{F.R4} (p.~171)} }
\expandafter\def\csname SRef@maps:eq:F.R5\endcsname{\href{S6_SERIES_GRIND_V.pdf\#nameddest=tag-02HS}{V.{F.R5} (p.~171)} }
\expandafter\def\csname SRef@maps:eq:F.R6\endcsname{\href{S6_SERIES_GRIND_V.pdf\#nameddest=tag-02HT}{V.{F.R6} (p.~171)} }
\expandafter\def\csname SRef@maps:eq:F.R7\endcsname{\href{S6_SERIES_GRIND_V.pdf\#nameddest=tag-02HU}{V.{F.R7} (p.~172)} }
\expandafter\def\csname SRef@maps:eq:F.R8\endcsname{\href{S6_SERIES_GRIND_V.pdf\#nameddest=tag-02HV}{V.{F.R8} (p.~172)} }
\expandafter\def\csname SRef@maps:eq:F.R9\endcsname{\href{S6_SERIES_GRIND_V.pdf\#nameddest=tag-02HW}{V.{F.R9} (p.~172)} }
\expandafter\def\csname SRef@maps:eq:F.L1\endcsname{\href{S6_SERIES_GRIND_V.pdf\#nameddest=tag-02GT}{V.{F.L1} (p.~174)} }
\expandafter\def\csname SRef@maps:mapsF:repair:cyclic-inclusion\endcsname{\href{S6_SERIES_GRIND_V.pdf\#nameddest=tag-02J1}{V.{{$\langle 1\rangle $6.}} (p.~199)} }
\expandafter\def\csname SRef@maps:mapsF:repair:log-cylinder-scope\endcsname{\href{S6_SERIES_GRIND_V.pdf\#nameddest=tag-02J2}{V.{{$\langle 1\rangle $5.}} (p.~201)} }
\expandafter\def\csname SRef@sec:fourfold\endcsname{\href{S6_SERIES_GRIND_V.pdf\#nameddest=tag-034O}{V.5 (p.~80)} }
\expandafter\def\csname SRef@maps:src:F.2.21\endcsname{\href{S6_SERIES_GRIND_V.pdf\#nameddest=tag-02NB}{V.11.15 (p.~173)} }
\expandafter\def\csname SRef@maps:src:F.2.30\endcsname{\href{S6_SERIES_GRIND_V.pdf\#nameddest=tag-02NL}{V.11.24 (p.~190)} }
\expandafter\def\csname SRef@maps:eq:F.1\endcsname{\href{S6_SERIES_GRIND_V.pdf\#nameddest=tag-02EF}{V.1 (p.~72)} }
\expandafter\def\csname SRef@maps:eq:F.2\endcsname{\href{S6_SERIES_GRIND_V.pdf\#nameddest=tag-02EQ}{V.2 (p.~72)} }
\expandafter\def\csname SRef@eq:P1Adegree\endcsname{\href{S6_SERIES_GRIND_V.pdf\#nameddest=tag-01O3}{V.3 (p.~134)} }
\expandafter\def\csname SRef@eq:P1Alogcotangent\endcsname{\href{S6_SERIES_GRIND_V.pdf\#nameddest=tag-01O4}{V.4 (p.~135)} }
\expandafter\def\csname SRef@maps:src:F.3.10\endcsname{\href{S6_SERIES_GRIND_V.pdf\#nameddest=tag-02O2}{V.12.5 (p.~218)} }
\expandafter\def\csname SRef@mapsaudit:order-four-torsion-transgression\endcsname{\href{S6_SERIES_GRIND_V.pdf\#nameddest=tag-02QF}{V.5.4 (p.~81)} }
\expandafter\def\csname SRef@maps:src:F.2.17\endcsname{\href{S6_SERIES_GRIND_V.pdf\#nameddest=tag-02N6}{V.11.11 (p.~166)} }
\expandafter\def\csname SRef@maps:src:F.2.8\endcsname{\href{S6_SERIES_GRIND_V.pdf\#nameddest=tag-02NZ}{V.11.2 (p.~147)} }
\expandafter\def\csname SRef@maps:src:F.2.9\endcsname{\href{S6_SERIES_GRIND_V.pdf\#nameddest=tag-02O0}{V.11.3 (p.~148)} }
\expandafter\def\csname SRef@maps:src:F.2.11\endcsname{\href{S6_SERIES_GRIND_V.pdf\#nameddest=tag-02N0}{V.11.5 (p.~151)} }
\expandafter\def\csname SRef@maps:src:F.2.15\endcsname{\href{S6_SERIES_GRIND_V.pdf\#nameddest=tag-02N4}{V.11.9 (p.~162)} }
\expandafter\def\csname SRef@maps:src:F.3.3\endcsname{\href{S6_SERIES_GRIND_V.pdf\#nameddest=tag-02O5}{V.12.1 (p.~215)} }
\expandafter\def\csname SRef@maps:src:F.1.4\endcsname{\href{S6_SERIES_GRIND_V.pdf\#nameddest=tag-02MX}{V.4.2 (p.~72)} }
\expandafter\def\csname SRef@maps:src:F.2.1\endcsname{\href{S6_SERIES_GRIND_V.pdf\#nameddest=tag-02MY}{V.4.3 (p.~73)} }
\expandafter\def\csname SRef@maps:src:F.2.2\endcsname{\href{S6_SERIES_GRIND_V.pdf\#nameddest=tag-02N9}{V.4.4 (p.~73)} }
\expandafter\def\csname SRef@maps:src:F.2.3\endcsname{\href{S6_SERIES_GRIND_V.pdf\#nameddest=tag-02NK}{V.4.5 (p.~74)} }
\expandafter\def\csname SRef@maps:src:F.3.2\endcsname{\href{S6_SERIES_GRIND_V.pdf\#nameddest=tag-02O4}{V.4.7 (p.~75)} }
\expandafter\def\csname SRef@maps:src:F.2.4\endcsname{\href{S6_SERIES_GRIND_V.pdf\#nameddest=tag-02NV}{V.4.8 (p.~75)} }
\expandafter\def\csname SRef@sec:V-known-inputs\endcsname{\href{S6_SERIES_GRIND_V.pdf\#nameddest=tag-034M}{V.3 (p.~72)} }
\expandafter\def\csname SRef@maps:src:F.2.7\endcsname{\href{S6_SERIES_GRIND_V.pdf\#nameddest=tag-02NY}{V.11.1 (p.~146)} }
\expandafter\def\csname SRef@maps:src:F.2.16\endcsname{\href{S6_SERIES_GRIND_V.pdf\#nameddest=tag-02N5}{V.11.10 (p.~163)} }
\expandafter\def\csname SRef@maps:src:F.2.18\endcsname{\href{S6_SERIES_GRIND_V.pdf\#nameddest=tag-02N7}{V.11.12 (p.~167)} }
\expandafter\def\csname SRef@maps:src:F.2.19\endcsname{\href{S6_SERIES_GRIND_V.pdf\#nameddest=tag-02N8}{V.11.13 (p.~169)} }
\expandafter\def\csname SRef@maps:src:F.2.20\endcsname{\href{S6_SERIES_GRIND_V.pdf\#nameddest=tag-02NA}{V.11.14 (p.~170)} }
\expandafter\def\csname SRef@maps:src:F.2.22\endcsname{\href{S6_SERIES_GRIND_V.pdf\#nameddest=tag-02NC}{V.11.16 (p.~174)} }
\expandafter\def\csname SRef@maps:src:F.2.23\endcsname{\href{S6_SERIES_GRIND_V.pdf\#nameddest=tag-02ND}{V.11.17 (p.~176)} }
\expandafter\def\csname SRef@maps:src:F.2.24\endcsname{\href{S6_SERIES_GRIND_V.pdf\#nameddest=tag-02NE}{V.11.18 (p.~177)} }
\expandafter\def\csname SRef@maps:src:F.2.25\endcsname{\href{S6_SERIES_GRIND_V.pdf\#nameddest=tag-02NF}{V.11.19 (p.~179)} }
\expandafter\def\csname SRef@maps:src:F.2.26\endcsname{\href{S6_SERIES_GRIND_V.pdf\#nameddest=tag-02NG}{V.11.20 (p.~181)} }
\expandafter\def\csname SRef@maps:src:F.2.27\endcsname{\href{S6_SERIES_GRIND_V.pdf\#nameddest=tag-02NH}{V.11.21 (p.~183)} }
\expandafter\def\csname SRef@maps:src:F.2.28\endcsname{\href{S6_SERIES_GRIND_V.pdf\#nameddest=tag-02NI}{V.11.22 (p.~185)} }
\expandafter\def\csname SRef@maps:src:F.2.29\endcsname{\href{S6_SERIES_GRIND_V.pdf\#nameddest=tag-02NJ}{V.11.23 (p.~186)} }
\expandafter\def\csname SRef@maps:src:F.2.31\endcsname{\href{S6_SERIES_GRIND_V.pdf\#nameddest=tag-02NM}{V.11.25 (p.~191)} }
\expandafter\def\csname SRef@maps:src:F.2.32\endcsname{\href{S6_SERIES_GRIND_V.pdf\#nameddest=tag-02NN}{V.11.26 (p.~195)} }
\expandafter\def\csname SRef@maps:src:F.2.33\endcsname{\href{S6_SERIES_GRIND_V.pdf\#nameddest=tag-02NO}{V.11.27 (p.~197)} }
\expandafter\def\csname SRef@maps:src:F.2.34\endcsname{\href{S6_SERIES_GRIND_V.pdf\#nameddest=tag-02NP}{V.11.28 (p.~200)} }
\expandafter\def\csname SRef@maps:src:F.2.35\endcsname{\href{S6_SERIES_GRIND_V.pdf\#nameddest=tag-02NQ}{V.11.29 (p.~202)} }
\expandafter\def\csname SRef@maps:src:F.2.36\endcsname{\href{S6_SERIES_GRIND_V.pdf\#nameddest=tag-02NR}{V.11.30 (p.~206)} }
\expandafter\def\csname SRef@maps:src:F.2.37\endcsname{\href{S6_SERIES_GRIND_V.pdf\#nameddest=tag-02NS}{V.11.31 (p.~208)} }
\expandafter\def\csname SRef@maps:src:F.2.38\endcsname{\href{S6_SERIES_GRIND_V.pdf\#nameddest=tag-02NT}{V.11.32 (p.~210)} }
\expandafter\def\csname SRef@maps:src:F.2.39\endcsname{\href{S6_SERIES_GRIND_V.pdf\#nameddest=tag-02NU}{V.11.33 (p.~212)} }
\expandafter\def\csname SRef@night:P1B.1\endcsname{\href{S6_SERIES_GRIND_V.pdf\#nameddest=tag-02R7}{V.11.34 (p.~213)} }
\expandafter\def\csname SRef@night:P1B.2\endcsname{\href{S6_SERIES_GRIND_V.pdf\#nameddest=tag-02R8}{V.11.35 (p.~214)} }
\expandafter\def\csname SRef@maps:src:F.2.10\endcsname{\href{S6_SERIES_GRIND_V.pdf\#nameddest=tag-02MZ}{V.11.4 (p.~150)} }
\expandafter\def\csname SRef@maps:src:F.2.12\endcsname{\href{S6_SERIES_GRIND_V.pdf\#nameddest=tag-02N1}{V.11.6 (p.~152)} }
\expandafter\def\csname SRef@maps:src:F.2.13\endcsname{\href{S6_SERIES_GRIND_V.pdf\#nameddest=tag-02N2}{V.11.7 (p.~155)} }
\expandafter\def\csname SRef@maps:src:F.2.14\endcsname{\href{S6_SERIES_GRIND_V.pdf\#nameddest=tag-02N3}{V.11.8 (p.~161)} }
\expandafter\def\csname SRef@maps:src:F.3.4\endcsname{\href{S6_SERIES_GRIND_V.pdf\#nameddest=tag-02O6}{V.12.2 (p.~215)} }
\expandafter\def\csname SRef@maps:src:F.3.8\endcsname{\href{S6_SERIES_GRIND_V.pdf\#nameddest=tag-02OA}{V.12.3 (p.~216)} }
\expandafter\def\csname SRef@maps:src:F.3.9\endcsname{\href{S6_SERIES_GRIND_V.pdf\#nameddest=tag-02OB}{V.12.4 (p.~217)} }
\expandafter\def\csname SRef@maps:src:F.3.11\endcsname{\href{S6_SERIES_GRIND_V.pdf\#nameddest=tag-02O3}{V.12.6 (p.~219)} }
\expandafter\def\csname SRef@maps:src:F.1.1\endcsname{\href{S6_SERIES_GRIND_V.pdf\#nameddest=tag-02MU}{V.4.1 (p.~72)} }
\expandafter\def\csname SRef@maps:src:F.3.1\endcsname{\href{S6_SERIES_GRIND_V.pdf\#nameddest=tag-02O1}{V.4.6 (p.~74)} }
\expandafter\def\csname SRef@maps:src:F.2.5\endcsname{\href{S6_SERIES_GRIND_V.pdf\#nameddest=tag-02NW}{V.5.1 (p.~76)} }
\expandafter\def\csname SRef@night:P1.1\endcsname{\href{S6_SERIES_GRIND_V.pdf\#nameddest=tag-02QI}{V.5.2 (p.~80)} }
\expandafter\def\csname SRef@night:P1.3\endcsname{\href{S6_SERIES_GRIND_V.pdf\#nameddest=tag-02QU}{V.5.3 (p.~80)} }
\expandafter\def\csname SRef@maps:src:F.2.6\endcsname{\href{S6_SERIES_GRIND_V.pdf\#nameddest=tag-02NX}{V.5.5 (p.~82)} }
\expandafter\def\csname SRef@night:P1.10\endcsname{\href{S6_SERIES_GRIND_V.pdf\#nameddest=tag-02QJ}{V.5.6 (p.~86)} }
\expandafter\def\csname SRef@night:P1.11\endcsname{\href{S6_SERIES_GRIND_V.pdf\#nameddest=tag-02QK}{V.5.7 (p.~87)} }
\expandafter\def\csname SRef@night:P1.12\endcsname{\href{S6_SERIES_GRIND_V.pdf\#nameddest=tag-02QL}{V.5.8 (p.~91)} }
\expandafter\def\csname SRef@night:P1.19\endcsname{\href{S6_SERIES_GRIND_V.pdf\#nameddest=tag-02QS}{V.5.9 (p.~94)} }
\expandafter\def\csname SRef@night:P1.4\endcsname{\href{S6_SERIES_GRIND_V.pdf\#nameddest=tag-02QV}{V.6.1 (p.~103)} }
\expandafter\def\csname SRef@night:P1.5\endcsname{\href{S6_SERIES_GRIND_V.pdf\#nameddest=tag-02QW}{V.6.2 (p.~104)} }
\expandafter\def\csname SRef@maps:src:F.3.5\endcsname{\href{S6_SERIES_GRIND_V.pdf\#nameddest=tag-02O7}{V.6.3 (p.~104)} }
\expandafter\def\csname SRef@night:P1.2\endcsname{\href{S6_SERIES_GRIND_V.pdf\#nameddest=tag-02QT}{V.6.4 (p.~106)} }
\expandafter\def\csname SRef@night:P1.7\endcsname{\href{S6_SERIES_GRIND_V.pdf\#nameddest=tag-02QY}{V.6.5 (p.~108)} }
\expandafter\def\csname SRef@maps:src:F.3.6\endcsname{\href{S6_SERIES_GRIND_V.pdf\#nameddest=tag-02O8}{V.6.6 (p.~109)} }
\expandafter\def\csname SRef@night:P1.8\endcsname{\href{S6_SERIES_GRIND_V.pdf\#nameddest=tag-02QZ}{V.6.7 (p.~112)} }
\expandafter\def\csname SRef@night:P1.9\endcsname{\href{S6_SERIES_GRIND_V.pdf\#nameddest=tag-02R0}{V.6.8 (p.~113)} }
\expandafter\def\csname SRef@maps:src:F.3.7\endcsname{\href{S6_SERIES_GRIND_V.pdf\#nameddest=tag-02O9}{V.6.9 (p.~114)} }
\expandafter\def\csname SRef@night:P1.13\endcsname{\href{S6_SERIES_GRIND_V.pdf\#nameddest=tag-02QM}{V.7.1 (p.~116)} }
\expandafter\def\csname SRef@maps:src:F.4.1\endcsname{\href{S6_SERIES_GRIND_V.pdf\#nameddest=tag-02OC}{V.7.2 (p.~117)} }
\expandafter\def\csname SRef@night:P1.15\endcsname{\href{S6_SERIES_GRIND_V.pdf\#nameddest=tag-02QO}{V.8.1 (p.~119)} }
\expandafter\def\csname SRef@night:P1.16\endcsname{\href{S6_SERIES_GRIND_V.pdf\#nameddest=tag-02QP}{V.8.2 (p.~121)} }
\expandafter\def\csname SRef@night:P1.18\endcsname{\href{S6_SERIES_GRIND_V.pdf\#nameddest=tag-02QR}{V.8.3 (p.~124)} }
\expandafter\def\csname SRef@night:P1A.1\endcsname{\href{S6_SERIES_GRIND_V.pdf\#nameddest=tag-02R1}{V.9.1 (p.~134)} }
\expandafter\def\csname SRef@night:P1A.2\endcsname{\href{S6_SERIES_GRIND_V.pdf\#nameddest=tag-02R2}{V.9.2 (p.~135)} }
\expandafter\def\csname SRef@night:P1A.3\endcsname{\href{S6_SERIES_GRIND_V.pdf\#nameddest=tag-02R3}{V.9.3 (p.~137)} }
\expandafter\def\csname SRef@night:P1A.4\endcsname{\href{S6_SERIES_GRIND_V.pdf\#nameddest=tag-02R4}{V.9.4 (p.~139)} }
\expandafter\def\csname SRef@night:P1.17\endcsname{\href{S6_SERIES_GRIND_V.pdf\#nameddest=tag-02QQ}{V.9.5 (p.~141)} }
\expandafter\def\csname SRef@night:P1A.5\endcsname{\href{S6_SERIES_GRIND_V.pdf\#nameddest=tag-02R5}{V.10.1 (p.~143)} }
\expandafter\def\csname SRef@night:P1A.6\endcsname{\href{S6_SERIES_GRIND_V.pdf\#nameddest=tag-02R6}{V.10.2 (p.~146)} }
\expandafter\def\csname SRef@cdp:eq:C.UN.1\endcsname{\hyperlink{tag-0152}{II.\ref*{cdp:eq:C.UN.1}}}
\expandafter\def\csname SRef@cdp:eq:C.FL.1\endcsname{\hyperlink{tag-00KR}{II.\ref*{cdp:eq:C.FL.1}}}
\expandafter\def\csname SRef@cdp:eq:C.QT.11\endcsname{\hyperlink{tag-00WF}{II.\ref*{cdp:eq:C.QT.11}}}
\expandafter\def\csname SRef@cdp:eq:C.QT.12\endcsname{\hyperlink{tag-00WG}{II.\ref*{cdp:eq:C.QT.12}}}
\expandafter\def\csname SRef@cdp:eq:C.QT.13\endcsname{\hyperlink{tag-00WH}{II.\ref*{cdp:eq:C.QT.13}}}
\expandafter\def\csname SRef@cdp:eq:C.WT.1\endcsname{\hyperlink{tag-017O}{II.\ref*{cdp:eq:C.WT.1}}}
\expandafter\def\csname SRef@cdp:eq:C.WT.2\endcsname{\hyperlink{tag-017Z}{II.\ref*{cdp:eq:C.WT.2}}}
\expandafter\def\csname SRef@cdp:eq:C.WT.3\endcsname{\hyperlink{tag-018B}{II.\ref*{cdp:eq:C.WT.3}}}
\expandafter\def\csname SRef@cdp:eq:C.WT.4\endcsname{\hyperlink{tag-018M}{II.\ref*{cdp:eq:C.WT.4}}}
\expandafter\def\csname SRef@cdp:eq:C.WT.5\endcsname{\hyperlink{tag-018R}{II.\ref*{cdp:eq:C.WT.5}}}
\expandafter\def\csname SRef@cdp:eq:C.WT.6\endcsname{\hyperlink{tag-018S}{II.\ref*{cdp:eq:C.WT.6}}}
\expandafter\def\csname SRef@cdp:eq:C.WT.7\endcsname{\hyperlink{tag-018T}{II.\ref*{cdp:eq:C.WT.7}}}
\expandafter\def\csname SRef@cdp:eq:C.FL.2\endcsname{\hyperlink{tag-00KS}{II.\ref*{cdp:eq:C.FL.2}}}
\expandafter\def\csname SRef@cdp:eq:C.WT.8\endcsname{\hyperlink{tag-018U}{II.\ref*{cdp:eq:C.WT.8}}}
\expandafter\def\csname SRef@cdp:eq:C.WT.9\endcsname{\hyperlink{tag-018V}{II.\ref*{cdp:eq:C.WT.9}}}
\expandafter\def\csname SRef@cdp:eq:C.WT.10\endcsname{\hyperlink{tag-017P}{II.\ref*{cdp:eq:C.WT.10}}}
\expandafter\def\csname SRef@cdp:eq:C.WT.11\endcsname{\hyperlink{tag-017Q}{II.\ref*{cdp:eq:C.WT.11}}}
\expandafter\def\csname SRef@cdp:eq:C.WT.12\endcsname{\hyperlink{tag-017R}{II.\ref*{cdp:eq:C.WT.12}}}
\expandafter\def\csname SRef@cdp:eq:C.WT.13\endcsname{\hyperlink{tag-017S}{II.\ref*{cdp:eq:C.WT.13}}}
\expandafter\def\csname SRef@cdp:eq:C.WT.14\endcsname{\hyperlink{tag-017T}{II.\ref*{cdp:eq:C.WT.14}}}
\expandafter\def\csname SRef@cdp:eq:C.WT.15\endcsname{\hyperlink{tag-017U}{II.\ref*{cdp:eq:C.WT.15}}}
\expandafter\def\csname SRef@cdp:eq:C.WT.16\endcsname{\hyperlink{tag-017V}{II.\ref*{cdp:eq:C.WT.16}}}
\expandafter\def\csname SRef@cdp:eq:C.WT.17\endcsname{\hyperlink{tag-017W}{II.\ref*{cdp:eq:C.WT.17}}}
\expandafter\def\csname SRef@cdp:eq:C.FL.3\endcsname{\hyperlink{tag-00KT}{II.\ref*{cdp:eq:C.FL.3}}}
\expandafter\def\csname SRef@cdp:eq:C.WT.18\endcsname{\hyperlink{tag-017X}{II.\ref*{cdp:eq:C.WT.18}}}
\expandafter\def\csname SRef@cdp:eq:C.WT.19\endcsname{\hyperlink{tag-017Y}{II.\ref*{cdp:eq:C.WT.19}}}
\expandafter\def\csname SRef@cdp:eq:C.WT.20\endcsname{\hyperlink{tag-0180}{II.\ref*{cdp:eq:C.WT.20}}}
\expandafter\def\csname SRef@cdp:eq:C.WT.21\endcsname{\hyperlink{tag-0181}{II.\ref*{cdp:eq:C.WT.21}}}
\expandafter\def\csname SRef@cdp:eq:C.WT.22\endcsname{\hyperlink{tag-0182}{II.\ref*{cdp:eq:C.WT.22}}}
\expandafter\def\csname SRef@cdp:eq:C.WT.23\endcsname{\hyperlink{tag-0183}{II.\ref*{cdp:eq:C.WT.23}}}
\expandafter\def\csname SRef@cdp:eq:C.WT.23a\endcsname{\hyperlink{tag-0184}{II.\ref*{cdp:eq:C.WT.23a}}}
\expandafter\def\csname SRef@cdp:eq:C.WT.24\endcsname{\hyperlink{tag-0185}{II.\ref*{cdp:eq:C.WT.24}}}
\expandafter\def\csname SRef@cdp:eq:C.WT.25\endcsname{\hyperlink{tag-0186}{II.\ref*{cdp:eq:C.WT.25}}}
\expandafter\def\csname SRef@cdp:eq:C.WT.26\endcsname{\hyperlink{tag-0187}{II.\ref*{cdp:eq:C.WT.26}}}
\expandafter\def\csname SRef@cdp:eq:C.FL.4\endcsname{\hyperlink{tag-00KU}{II.\ref*{cdp:eq:C.FL.4}}}
\expandafter\def\csname SRef@cdp:eq:C.WT.27\endcsname{\hyperlink{tag-0188}{II.\ref*{cdp:eq:C.WT.27}}}
\expandafter\def\csname SRef@cdp:eq:C.WT.28\endcsname{\hyperlink{tag-0189}{II.\ref*{cdp:eq:C.WT.28}}}
\expandafter\def\csname SRef@cdp:eq:C.WT.29\endcsname{\hyperlink{tag-018A}{II.\ref*{cdp:eq:C.WT.29}}}
\expandafter\def\csname SRef@cdp:eq:C.WT.30\endcsname{\hyperlink{tag-018C}{II.\ref*{cdp:eq:C.WT.30}}}
\expandafter\def\csname SRef@cdp:eq:C.WT.31\endcsname{\hyperlink{tag-018D}{II.\ref*{cdp:eq:C.WT.31}}}
\expandafter\def\csname SRef@cdp:eq:C.WT.32\endcsname{\hyperlink{tag-018E}{II.\ref*{cdp:eq:C.WT.32}}}
\expandafter\def\csname SRef@cdp:eq:C.WT.33\endcsname{\hyperlink{tag-018F}{II.\ref*{cdp:eq:C.WT.33}}}
\expandafter\def\csname SRef@cdp:eq:C.WT.34\endcsname{\hyperlink{tag-018G}{II.\ref*{cdp:eq:C.WT.34}}}
\expandafter\def\csname SRef@cdp:eq:C.WT.35\endcsname{\hyperlink{tag-018H}{II.\ref*{cdp:eq:C.WT.35}}}
\expandafter\def\csname SRef@cdp:eq:C.WT.36\endcsname{\hyperlink{tag-018I}{II.\ref*{cdp:eq:C.WT.36}}}
\expandafter\def\csname SRef@cdp:eq:C.FL.5\endcsname{\hyperlink{tag-00KV}{II.\ref*{cdp:eq:C.FL.5}}}
\expandafter\def\csname SRef@cdp:eq:C.WT.37\endcsname{\hyperlink{tag-018J}{II.\ref*{cdp:eq:C.WT.37}}}
\expandafter\def\csname SRef@cdp:eq:C.WT.38\endcsname{\hyperlink{tag-018K}{II.\ref*{cdp:eq:C.WT.38}}}
\expandafter\def\csname SRef@cdp:eq:C.WT.39\endcsname{\hyperlink{tag-018L}{II.\ref*{cdp:eq:C.WT.39}}}
\expandafter\def\csname SRef@cdp:eq:C.WT.40\endcsname{\hyperlink{tag-018N}{II.\ref*{cdp:eq:C.WT.40}}}
\expandafter\def\csname SRef@cdp:eq:C.WT.41\endcsname{\hyperlink{tag-018O}{II.\ref*{cdp:eq:C.WT.41}}}
\expandafter\def\csname SRef@cdp:eq:C.WT.42\endcsname{\hyperlink{tag-018P}{II.\ref*{cdp:eq:C.WT.42}}}
\expandafter\def\csname SRef@cdp:eq:C.WT.43\endcsname{\hyperlink{tag-018Q}{II.\ref*{cdp:eq:C.WT.43}}}
\expandafter\def\csname SRef@cdp:eq:C.NM.1\endcsname{\hyperlink{tag-00SR}{II.\ref*{cdp:eq:C.NM.1}}}
\expandafter\def\csname SRef@cdp:eq:C.NM.2\endcsname{\hyperlink{tag-00ST}{II.\ref*{cdp:eq:C.NM.2}}}
\expandafter\def\csname SRef@cdp:eq:C.NM.3\endcsname{\hyperlink{tag-00SU}{II.\ref*{cdp:eq:C.NM.3}}}
\expandafter\def\csname SRef@cdp:eq:C.PG.1\endcsname{\hyperlink{tag-00UZ}{II.\ref*{cdp:eq:C.PG.1}}}
\expandafter\def\csname SRef@cdp:eq:C.NM.4\endcsname{\hyperlink{tag-00SV}{II.\ref*{cdp:eq:C.NM.4}}}
\expandafter\def\csname SRef@cdp:eq:C.NM.5\endcsname{\hyperlink{tag-00SW}{II.\ref*{cdp:eq:C.NM.5}}}
\expandafter\def\csname SRef@cdp:eq:C.NM.6\endcsname{\hyperlink{tag-00SX}{II.\ref*{cdp:eq:C.NM.6}}}
\expandafter\def\csname SRef@cdp:eq:C.NM.7\endcsname{\hyperlink{tag-00SY}{II.\ref*{cdp:eq:C.NM.7}}}
\expandafter\def\csname SRef@cdp:eq:C.NM.8\endcsname{\hyperlink{tag-00SZ}{II.\ref*{cdp:eq:C.NM.8}}}
\expandafter\def\csname SRef@cdp:eq:C.NM.9\endcsname{\hyperlink{tag-00T0}{II.\ref*{cdp:eq:C.NM.9}}}
\expandafter\def\csname SRef@cdp:eq:C.NM.10\endcsname{\hyperlink{tag-00SS}{II.\ref*{cdp:eq:C.NM.10}}}
\expandafter\def\csname SRef@cdp:eq:C.NP.1\endcsname{\hyperlink{tag-00T1}{II.\ref*{cdp:eq:C.NP.1}}}
\expandafter\def\csname SRef@cdp:eq:C.NP.2\endcsname{\hyperlink{tag-00T2}{II.\ref*{cdp:eq:C.NP.2}}}
\expandafter\def\csname SRef@cdp:eq:C.NP.3\endcsname{\hyperlink{tag-00T3}{II.\ref*{cdp:eq:C.NP.3}}}
\expandafter\def\csname SRef@cdp:eq:C.PG.2\endcsname{\hyperlink{tag-00V0}{II.\ref*{cdp:eq:C.PG.2}}}
\expandafter\def\csname SRef@cdp:eq:C.NP.4\endcsname{\hyperlink{tag-00T4}{II.\ref*{cdp:eq:C.NP.4}}}
\expandafter\def\csname SRef@cdp:eq:C.NP.5\endcsname{\hyperlink{tag-00T5}{II.\ref*{cdp:eq:C.NP.5}}}
\expandafter\def\csname SRef@cdp:eq:C.NP.6\endcsname{\hyperlink{tag-00T6}{II.\ref*{cdp:eq:C.NP.6}}}
\expandafter\def\csname SRef@cdp:eq:C.NP.7\endcsname{\hyperlink{tag-00T7}{II.\ref*{cdp:eq:C.NP.7}}}
\expandafter\def\csname SRef@cdp:eq:C.NP.8\endcsname{\hyperlink{tag-00T8}{II.\ref*{cdp:eq:C.NP.8}}}
\expandafter\def\csname SRef@cdp:eq:C.NP.9\endcsname{\hyperlink{tag-00T9}{II.\ref*{cdp:eq:C.NP.9}}}
\expandafter\def\csname SRef@cdp:eq:C.TG.1\endcsname{\hyperlink{tag-012Q}{II.\ref*{cdp:eq:C.TG.1}}}
\expandafter\def\csname SRef@cdp:eq:C.TG.2\endcsname{\hyperlink{tag-012X}{II.\ref*{cdp:eq:C.TG.2}}}
\expandafter\def\csname SRef@cdp:eq:C.TG.3\endcsname{\hyperlink{tag-012Y}{II.\ref*{cdp:eq:C.TG.3}}}
\expandafter\def\csname SRef@cdp:eq:C.TG.4\endcsname{\hyperlink{tag-012Z}{II.\ref*{cdp:eq:C.TG.4}}}
\expandafter\def\csname SRef@cdp:eq:C.PG.3\endcsname{\hyperlink{tag-00V1}{II.\ref*{cdp:eq:C.PG.3}}}
\expandafter\def\csname SRef@cdp:eq:C.TG.5\endcsname{\hyperlink{tag-0130}{II.\ref*{cdp:eq:C.TG.5}}}
\expandafter\def\csname SRef@cdp:eq:C.TG.6\endcsname{\hyperlink{tag-0131}{II.\ref*{cdp:eq:C.TG.6}}}
\expandafter\def\csname SRef@cdp:eq:C.TG.7\endcsname{\hyperlink{tag-0132}{II.\ref*{cdp:eq:C.TG.7}}}
\expandafter\def\csname SRef@cdp:eq:C.TG.8\endcsname{\hyperlink{tag-0133}{II.\ref*{cdp:eq:C.TG.8}}}
\expandafter\def\csname SRef@cdp:eq:C.TG.9\endcsname{\hyperlink{tag-0134}{II.\ref*{cdp:eq:C.TG.9}}}
\expandafter\def\csname SRef@cdp:eq:C.TG.10\endcsname{\hyperlink{tag-012R}{II.\ref*{cdp:eq:C.TG.10}}}
\expandafter\def\csname SRef@cdp:eq:C.TG.11\endcsname{\hyperlink{tag-012S}{II.\ref*{cdp:eq:C.TG.11}}}
\expandafter\def\csname SRef@cdp:eq:C.TG.12\endcsname{\hyperlink{tag-012T}{II.\ref*{cdp:eq:C.TG.12}}}
\expandafter\def\csname SRef@cdp:eq:C.TG.13\endcsname{\hyperlink{tag-012U}{II.\ref*{cdp:eq:C.TG.13}}}
\expandafter\def\csname SRef@cdp:eq:C.TG.14\endcsname{\hyperlink{tag-012V}{II.\ref*{cdp:eq:C.TG.14}}}
\expandafter\def\csname SRef@cdp:eq:C.PG.4\endcsname{\hyperlink{tag-00V2}{II.\ref*{cdp:eq:C.PG.4}}}
\expandafter\def\csname SRef@cdp:eq:C.TG.15\endcsname{\hyperlink{tag-012W}{II.\ref*{cdp:eq:C.TG.15}}}
\expandafter\def\csname SRef@cdp:eq:C.TX.1\endcsname{\hyperlink{tag-013Y}{II.\ref*{cdp:eq:C.TX.1}}}
\expandafter\def\csname SRef@cdp:eq:C.TX.2\endcsname{\hyperlink{tag-0144}{II.\ref*{cdp:eq:C.TX.2}}}
\expandafter\def\csname SRef@cdp:eq:C.TX.3\endcsname{\hyperlink{tag-0145}{II.\ref*{cdp:eq:C.TX.3}}}
\expandafter\def\csname SRef@cdp:eq:C.TX.4\endcsname{\hyperlink{tag-0146}{II.\ref*{cdp:eq:C.TX.4}}}
\expandafter\def\csname SRef@cdp:eq:C.TX.5\endcsname{\hyperlink{tag-0147}{II.\ref*{cdp:eq:C.TX.5}}}
\expandafter\def\csname SRef@cdp:eq:C.TX.6\endcsname{\hyperlink{tag-0148}{II.\ref*{cdp:eq:C.TX.6}}}
\expandafter\def\csname SRef@cdp:eq:C.TX.7\endcsname{\hyperlink{tag-0149}{II.\ref*{cdp:eq:C.TX.7}}}
\expandafter\def\csname SRef@cdp:eq:C.TX.8\endcsname{\hyperlink{tag-014A}{II.\ref*{cdp:eq:C.TX.8}}}
\expandafter\def\csname SRef@cdp:eq:C.TX.9\endcsname{\hyperlink{tag-014B}{II.\ref*{cdp:eq:C.TX.9}}}
\expandafter\def\csname SRef@cdp:eq:C.PG.5\endcsname{\hyperlink{tag-00V3}{II.\ref*{cdp:eq:C.PG.5}}}
\expandafter\def\csname SRef@cdp:eq:C.TX.10\endcsname{\hyperlink{tag-013Z}{II.\ref*{cdp:eq:C.TX.10}}}
\expandafter\def\csname SRef@cdp:eq:C.TX.11\endcsname{\hyperlink{tag-0140}{II.\ref*{cdp:eq:C.TX.11}}}
\expandafter\def\csname SRef@cdp:eq:C.TX.12\endcsname{\hyperlink{tag-0141}{II.\ref*{cdp:eq:C.TX.12}}}
\expandafter\def\csname SRef@cdp:eq:C.TX.13\endcsname{\hyperlink{tag-0142}{II.\ref*{cdp:eq:C.TX.13}}}
\expandafter\def\csname SRef@cdp:eq:C.TX.14\endcsname{\hyperlink{tag-0143}{II.\ref*{cdp:eq:C.TX.14}}}
\expandafter\def\csname SRef@cdp:eq:C.WF.1\endcsname{\hyperlink{tag-0165}{II.\ref*{cdp:eq:C.WF.1}}}
\expandafter\def\csname SRef@cdp:eq:C.WF.2\endcsname{\hyperlink{tag-016G}{II.\ref*{cdp:eq:C.WF.2}}}
\expandafter\def\csname SRef@cdp:eq:C.WF.3\endcsname{\hyperlink{tag-016Q}{II.\ref*{cdp:eq:C.WF.3}}}
\expandafter\def\csname SRef@cdp:eq:C.WF.4\endcsname{\hyperlink{tag-016R}{II.\ref*{cdp:eq:C.WF.4}}}
\expandafter\def\csname SRef@cdp:eq:C.WF.5\endcsname{\hyperlink{tag-016S}{II.\ref*{cdp:eq:C.WF.5}}}
\expandafter\def\csname SRef@cdp:eq:C.UN.2\endcsname{\hyperlink{tag-0153}{II.\ref*{cdp:eq:C.UN.2}}}
\expandafter\def\csname SRef@cdp:eq:C.PG.6\endcsname{\hyperlink{tag-00V4}{II.\ref*{cdp:eq:C.PG.6}}}
\expandafter\def\csname SRef@cdp:eq:C.WF.6\endcsname{\hyperlink{tag-016T}{II.\ref*{cdp:eq:C.WF.6}}}
\expandafter\def\csname SRef@cdp:eq:C.WF.7\endcsname{\hyperlink{tag-016U}{II.\ref*{cdp:eq:C.WF.7}}}
\expandafter\def\csname SRef@cdp:eq:C.WF.8\endcsname{\hyperlink{tag-016V}{II.\ref*{cdp:eq:C.WF.8}}}
\expandafter\def\csname SRef@cdp:eq:C.WF.9\endcsname{\hyperlink{tag-016W}{II.\ref*{cdp:eq:C.WF.9}}}
\expandafter\def\csname SRef@cdp:eq:C.WF.10\endcsname{\hyperlink{tag-0166}{II.\ref*{cdp:eq:C.WF.10}}}
\expandafter\def\csname SRef@cdp:eq:C.WF.11\endcsname{\hyperlink{tag-0167}{II.\ref*{cdp:eq:C.WF.11}}}
\expandafter\def\csname SRef@cdp:eq:C.WF.12\endcsname{\hyperlink{tag-0168}{II.\ref*{cdp:eq:C.WF.12}}}
\expandafter\def\csname SRef@cdp:eq:C.WF.13\endcsname{\hyperlink{tag-0169}{II.\ref*{cdp:eq:C.WF.13}}}
\expandafter\def\csname SRef@cdp:eq:C.WF.14\endcsname{\hyperlink{tag-016A}{II.\ref*{cdp:eq:C.WF.14}}}
\expandafter\def\csname SRef@cdp:eq:C.WF.15\endcsname{\hyperlink{tag-016B}{II.\ref*{cdp:eq:C.WF.15}}}
\expandafter\def\csname SRef@cdp:eq:C.PG.7\endcsname{\hyperlink{tag-00V5}{II.\ref*{cdp:eq:C.PG.7}}}
\expandafter\def\csname SRef@cdp:eq:C.WF.16\endcsname{\hyperlink{tag-016C}{II.\ref*{cdp:eq:C.WF.16}}}
\expandafter\def\csname SRef@cdp:eq:C.WF.17\endcsname{\hyperlink{tag-016D}{II.\ref*{cdp:eq:C.WF.17}}}
\expandafter\def\csname SRef@cdp:eq:C.WF.18\endcsname{\hyperlink{tag-016E}{II.\ref*{cdp:eq:C.WF.18}}}
\expandafter\def\csname SRef@cdp:eq:C.WF.19\endcsname{\hyperlink{tag-016F}{II.\ref*{cdp:eq:C.WF.19}}}
\expandafter\def\csname SRef@cdp:eq:C.WF.20\endcsname{\hyperlink{tag-016H}{II.\ref*{cdp:eq:C.WF.20}}}
\expandafter\def\csname SRef@cdp:eq:C.WF.21\endcsname{\hyperlink{tag-016I}{II.\ref*{cdp:eq:C.WF.21}}}
\expandafter\def\csname SRef@cdp:eq:C.WF.22\endcsname{\hyperlink{tag-016J}{II.\ref*{cdp:eq:C.WF.22}}}
\expandafter\def\csname SRef@cdp:eq:C.WF.23\endcsname{\hyperlink{tag-016K}{II.\ref*{cdp:eq:C.WF.23}}}
\expandafter\def\csname SRef@cdp:eq:C.WF.24\endcsname{\hyperlink{tag-016L}{II.\ref*{cdp:eq:C.WF.24}}}
\expandafter\def\csname SRef@cdp:eq:C.WF.25\endcsname{\hyperlink{tag-016M}{II.\ref*{cdp:eq:C.WF.25}}}
\expandafter\def\csname SRef@cdp:eq:C.PG.8\endcsname{\hyperlink{tag-00V6}{II.\ref*{cdp:eq:C.PG.8}}}
\expandafter\def\csname SRef@cdp:eq:C.WF.26\endcsname{\hyperlink{tag-016N}{II.\ref*{cdp:eq:C.WF.26}}}
\expandafter\def\csname SRef@cdp:eq:C.WF.27\endcsname{\hyperlink{tag-016O}{II.\ref*{cdp:eq:C.WF.27}}}
\expandafter\def\csname SRef@cdp:eq:C.WF.28\endcsname{\hyperlink{tag-016P}{II.\ref*{cdp:eq:C.WF.28}}}
\expandafter\def\csname SRef@cdp:eq:C.RC.1\endcsname{\hyperlink{tag-00XN}{II.\ref*{cdp:eq:C.RC.1}}}
\expandafter\def\csname SRef@cdp:eq:C.RC.2\endcsname{\hyperlink{tag-00XR}{II.\ref*{cdp:eq:C.RC.2}}}
\expandafter\def\csname SRef@cdp:eq:C.RC.3\endcsname{\hyperlink{tag-00XS}{II.\ref*{cdp:eq:C.RC.3}}}
\expandafter\def\csname SRef@cdp:eq:C.RC.4\endcsname{\hyperlink{tag-00XT}{II.\ref*{cdp:eq:C.RC.4}}}
\expandafter\def\csname SRef@cdp:eq:C.RC.5\endcsname{\hyperlink{tag-00XU}{II.\ref*{cdp:eq:C.RC.5}}}
\expandafter\def\csname SRef@cdp:eq:C.RC.6\endcsname{\hyperlink{tag-00XV}{II.\ref*{cdp:eq:C.RC.6}}}
\expandafter\def\csname SRef@cdp:eq:C.RC.7\endcsname{\hyperlink{tag-00XW}{II.\ref*{cdp:eq:C.RC.7}}}
\expandafter\def\csname SRef@cdp:eq:C.PG.9\endcsname{\hyperlink{tag-00V7}{II.\ref*{cdp:eq:C.PG.9}}}
\expandafter\def\csname SRef@cdp:eq:C.RC.8\endcsname{\hyperlink{tag-00XX}{II.\ref*{cdp:eq:C.RC.8}}}
\expandafter\def\csname SRef@cdp:eq:C.RC.9\endcsname{\hyperlink{tag-00XY}{II.\ref*{cdp:eq:C.RC.9}}}
\expandafter\def\csname SRef@cdp:eq:C.RC.10\endcsname{\hyperlink{tag-00XO}{II.\ref*{cdp:eq:C.RC.10}}}
\expandafter\def\csname SRef@cdp:eq:C.RC.11\endcsname{\hyperlink{tag-00XP}{II.\ref*{cdp:eq:C.RC.11}}}
\expandafter\def\csname SRef@cdp:eq:C.RC.12\endcsname{\hyperlink{tag-00XQ}{II.\ref*{cdp:eq:C.RC.12}}}
\expandafter\def\csname SRef@cdp:eq:C.KG.1\endcsname{\hyperlink{tag-00P8}{II.\ref*{cdp:eq:C.KG.1}}}
\expandafter\def\csname SRef@cdp:eq:C.KG.2\endcsname{\hyperlink{tag-00PA}{II.\ref*{cdp:eq:C.KG.2}}}
\expandafter\def\csname SRef@cdp:eq:C.KG.3\endcsname{\hyperlink{tag-00PB}{II.\ref*{cdp:eq:C.KG.3}}}
\expandafter\def\csname SRef@cdp:eq:C.KG.4\endcsname{\hyperlink{tag-00PC}{II.\ref*{cdp:eq:C.KG.4}}}
\expandafter\def\csname SRef@cdp:eq:C.KG.5\endcsname{\hyperlink{tag-00PD}{II.\ref*{cdp:eq:C.KG.5}}}
\expandafter\def\csname SRef@cdp:eq:C.RG.1\endcsname{\hyperlink{tag-00Z3}{II.\ref*{cdp:eq:C.RG.1}}}
\expandafter\def\csname SRef@cdp:eq:C.KG.6\endcsname{\hyperlink{tag-00PE}{II.\ref*{cdp:eq:C.KG.6}}}
\expandafter\def\csname SRef@cdp:eq:C.KG.7\endcsname{\hyperlink{tag-00PF}{II.\ref*{cdp:eq:C.KG.7}}}
\expandafter\def\csname SRef@cdp:eq:C.KG.8\endcsname{\hyperlink{tag-00PG}{II.\ref*{cdp:eq:C.KG.8}}}
\expandafter\def\csname SRef@cdp:eq:C.KG.9\endcsname{\hyperlink{tag-00PH}{II.\ref*{cdp:eq:C.KG.9}}}
\expandafter\def\csname SRef@cdp:eq:C.KG.10\endcsname{\hyperlink{tag-00P9}{II.\ref*{cdp:eq:C.KG.10}}}
\expandafter\def\csname SRef@cdp:eq:C.HM.1\endcsname{\hyperlink{tag-00MG}{II.\ref*{cdp:eq:C.HM.1}}}
\expandafter\def\csname SRef@cdp:eq:C.HM.2\endcsname{\hyperlink{tag-00MI}{II.\ref*{cdp:eq:C.HM.2}}}
\expandafter\def\csname SRef@cdp:eq:C.HM.3\endcsname{\hyperlink{tag-00MJ}{II.\ref*{cdp:eq:C.HM.3}}}
\expandafter\def\csname SRef@cdp:eq:C.HM.4\endcsname{\hyperlink{tag-00MK}{II.\ref*{cdp:eq:C.HM.4}}}
\expandafter\def\csname SRef@cdp:eq:C.HM.5\endcsname{\hyperlink{tag-00ML}{II.\ref*{cdp:eq:C.HM.5}}}
\expandafter\def\csname SRef@cdp:eq:C.RG.2\endcsname{\hyperlink{tag-00Z8}{II.\ref*{cdp:eq:C.RG.2}}}
\expandafter\def\csname SRef@cdp:eq:C.HM.6\endcsname{\hyperlink{tag-00MM}{II.\ref*{cdp:eq:C.HM.6}}}
\expandafter\def\csname SRef@cdp:eq:C.HM.7\endcsname{\hyperlink{tag-00MN}{II.\ref*{cdp:eq:C.HM.7}}}
\expandafter\def\csname SRef@cdp:eq:C.HM.8\endcsname{\hyperlink{tag-00MO}{II.\ref*{cdp:eq:C.HM.8}}}
\expandafter\def\csname SRef@cdp:eq:C.HM.9\endcsname{\hyperlink{tag-00MP}{II.\ref*{cdp:eq:C.HM.9}}}
\expandafter\def\csname SRef@cdp:eq:C.HM.10\endcsname{\hyperlink{tag-00MH}{II.\ref*{cdp:eq:C.HM.10}}}
\expandafter\def\csname SRef@cdp:eq:C.BC.1\endcsname{\hyperlink{tag-00AZ}{II.\ref*{cdp:eq:C.BC.1}}}
\expandafter\def\csname SRef@cdp:eq:C.BC.2\endcsname{\hyperlink{tag-00B4}{II.\ref*{cdp:eq:C.BC.2}}}
\expandafter\def\csname SRef@cdp:eq:C.BC.3\endcsname{\hyperlink{tag-00B5}{II.\ref*{cdp:eq:C.BC.3}}}
\expandafter\def\csname SRef@cdp:eq:C.BC.4\endcsname{\hyperlink{tag-00B6}{II.\ref*{cdp:eq:C.BC.4}}}
\expandafter\def\csname SRef@cdp:eq:C.BC.5\endcsname{\hyperlink{tag-00B7}{II.\ref*{cdp:eq:C.BC.5}}}
\expandafter\def\csname SRef@cdp:eq:C.RG.3\endcsname{\hyperlink{tag-00Z9}{II.\ref*{cdp:eq:C.RG.3}}}
\expandafter\def\csname SRef@cdp:eq:C.BC.6\endcsname{\hyperlink{tag-00B8}{II.\ref*{cdp:eq:C.BC.6}}}
\expandafter\def\csname SRef@cdp:eq:C.BC.7\endcsname{\hyperlink{tag-00B9}{II.\ref*{cdp:eq:C.BC.7}}}
\expandafter\def\csname SRef@cdp:eq:C.BC.8\endcsname{\hyperlink{tag-00BA}{II.\ref*{cdp:eq:C.BC.8}}}
\expandafter\def\csname SRef@cdp:eq:C.BC.9\endcsname{\hyperlink{tag-00BB}{II.\ref*{cdp:eq:C.BC.9}}}
\expandafter\def\csname SRef@cdp:eq:C.BC.10\endcsname{\hyperlink{tag-00B0}{II.\ref*{cdp:eq:C.BC.10}}}
\expandafter\def\csname SRef@cdp:eq:C.BC.11\endcsname{\hyperlink{tag-00B1}{II.\ref*{cdp:eq:C.BC.11}}}
\expandafter\def\csname SRef@cdp:eq:C.BC.12\endcsname{\hyperlink{tag-00B2}{II.\ref*{cdp:eq:C.BC.12}}}
\expandafter\def\csname SRef@cdp:eq:C.BC.13\endcsname{\hyperlink{tag-00B3}{II.\ref*{cdp:eq:C.BC.13}}}
\expandafter\def\csname SRef@cdp:eq:C.EP.1\endcsname{\hyperlink{tag-00I1}{II.\ref*{cdp:eq:C.EP.1}}}
\expandafter\def\csname SRef@cdp:eq:C.EP.2\endcsname{\hyperlink{tag-00I4}{II.\ref*{cdp:eq:C.EP.2}}}
\expandafter\def\csname SRef@cdp:eq:C.RG.4\endcsname{\hyperlink{tag-00ZA}{II.\ref*{cdp:eq:C.RG.4}}}
\expandafter\def\csname SRef@cdp:eq:C.EP.3\endcsname{\hyperlink{tag-00I5}{II.\ref*{cdp:eq:C.EP.3}}}
\expandafter\def\csname SRef@cdp:eq:C.EP.4\endcsname{\hyperlink{tag-00I6}{II.\ref*{cdp:eq:C.EP.4}}}
\expandafter\def\csname SRef@cdp:eq:C.EP.5\endcsname{\hyperlink{tag-00I7}{II.\ref*{cdp:eq:C.EP.5}}}
\expandafter\def\csname SRef@cdp:eq:C.EP.6\endcsname{\hyperlink{tag-00I8}{II.\ref*{cdp:eq:C.EP.6}}}
\expandafter\def\csname SRef@cdp:eq:C.EP.7\endcsname{\hyperlink{tag-00I9}{II.\ref*{cdp:eq:C.EP.7}}}
\expandafter\def\csname SRef@cdp:eq:C.EP.8\endcsname{\hyperlink{tag-00IA}{II.\ref*{cdp:eq:C.EP.8}}}
\expandafter\def\csname SRef@cdp:eq:C.EP.9\endcsname{\hyperlink{tag-00IB}{II.\ref*{cdp:eq:C.EP.9}}}
\expandafter\def\csname SRef@cdp:eq:C.EP.10\endcsname{\hyperlink{tag-00I2}{II.\ref*{cdp:eq:C.EP.10}}}
\expandafter\def\csname SRef@cdp:eq:C.EP.11\endcsname{\hyperlink{tag-00I3}{II.\ref*{cdp:eq:C.EP.11}}}
\expandafter\def\csname SRef@cdp:eq:C.AP.1\endcsname{\hyperlink{tag-009M}{II.\ref*{cdp:eq:C.AP.1}}}
\expandafter\def\csname SRef@cdp:eq:C.RG.5\endcsname{\hyperlink{tag-00ZB}{II.\ref*{cdp:eq:C.RG.5}}}
\expandafter\def\csname SRef@cdp:eq:C.AP.2\endcsname{\hyperlink{tag-009S}{II.\ref*{cdp:eq:C.AP.2}}}
\expandafter\def\csname SRef@cdp:eq:C.AP.3\endcsname{\hyperlink{tag-009T}{II.\ref*{cdp:eq:C.AP.3}}}
\expandafter\def\csname SRef@cdp:eq:C.AP.4\endcsname{\hyperlink{tag-009U}{II.\ref*{cdp:eq:C.AP.4}}}
\expandafter\def\csname SRef@cdp:eq:C.AP.5\endcsname{\hyperlink{tag-009V}{II.\ref*{cdp:eq:C.AP.5}}}
\expandafter\def\csname SRef@cdp:eq:C.AP.6\endcsname{\hyperlink{tag-009W}{II.\ref*{cdp:eq:C.AP.6}}}
\expandafter\def\csname SRef@cdp:eq:C.AP.7\endcsname{\hyperlink{tag-009X}{II.\ref*{cdp:eq:C.AP.7}}}
\expandafter\def\csname SRef@cdp:eq:C.AP.8\endcsname{\hyperlink{tag-009Y}{II.\ref*{cdp:eq:C.AP.8}}}
\expandafter\def\csname SRef@cdp:eq:C.AP.9\endcsname{\hyperlink{tag-009Z}{II.\ref*{cdp:eq:C.AP.9}}}
\expandafter\def\csname SRef@cdp:eq:C.AP.10\endcsname{\hyperlink{tag-009N}{II.\ref*{cdp:eq:C.AP.10}}}
\expandafter\def\csname SRef@cdp:eq:C.AP.11\endcsname{\hyperlink{tag-009O}{II.\ref*{cdp:eq:C.AP.11}}}
\expandafter\def\csname SRef@cdp:eq:C.RG.6\endcsname{\hyperlink{tag-00ZC}{II.\ref*{cdp:eq:C.RG.6}}}
\expandafter\def\csname SRef@cdp:eq:C.AP.12\endcsname{\hyperlink{tag-009P}{II.\ref*{cdp:eq:C.AP.12}}}
\expandafter\def\csname SRef@cdp:eq:C.AP.13\endcsname{\hyperlink{tag-009Q}{II.\ref*{cdp:eq:C.AP.13}}}
\expandafter\def\csname SRef@cdp:eq:C.AP.14\endcsname{\hyperlink{tag-009R}{II.\ref*{cdp:eq:C.AP.14}}}
\expandafter\def\csname SRef@cdp:eq:C.ML.1\endcsname{\hyperlink{tag-00QR}{II.\ref*{cdp:eq:C.ML.1}}}
\expandafter\def\csname SRef@cdp:eq:C.ML.2\endcsname{\hyperlink{tag-00QT}{II.\ref*{cdp:eq:C.ML.2}}}
\expandafter\def\csname SRef@cdp:eq:C.ML.3\endcsname{\hyperlink{tag-00QU}{II.\ref*{cdp:eq:C.ML.3}}}
\expandafter\def\csname SRef@cdp:eq:C.ML.4\endcsname{\hyperlink{tag-00QV}{II.\ref*{cdp:eq:C.ML.4}}}
\expandafter\def\csname SRef@cdp:eq:C.ML.5\endcsname{\hyperlink{tag-00QW}{II.\ref*{cdp:eq:C.ML.5}}}
\expandafter\def\csname SRef@cdp:eq:C.ML.6\endcsname{\hyperlink{tag-00QX}{II.\ref*{cdp:eq:C.ML.6}}}
\expandafter\def\csname SRef@cdp:eq:C.ML.7\endcsname{\hyperlink{tag-00QY}{II.\ref*{cdp:eq:C.ML.7}}}
\expandafter\def\csname SRef@cdp:eq:C.UN.3\endcsname{\hyperlink{tag-0154}{II.\ref*{cdp:eq:C.UN.3}}}
\expandafter\def\csname SRef@cdp:eq:C.RG.7\endcsname{\hyperlink{tag-00ZD}{II.\ref*{cdp:eq:C.RG.7}}}
\expandafter\def\csname SRef@cdp:eq:C.ML.8\endcsname{\hyperlink{tag-00QZ}{II.\ref*{cdp:eq:C.ML.8}}}
\expandafter\def\csname SRef@cdp:eq:C.ML.9\endcsname{\hyperlink{tag-00R0}{II.\ref*{cdp:eq:C.ML.9}}}
\expandafter\def\csname SRef@cdp:eq:C.ML.10\endcsname{\hyperlink{tag-00QS}{II.\ref*{cdp:eq:C.ML.10}}}
\expandafter\def\csname SRef@cdp:eq:C.TI.1\endcsname{\hyperlink{tag-0135}{II.\ref*{cdp:eq:C.TI.1}}}
\expandafter\def\csname SRef@cdp:eq:C.TI.2\endcsname{\hyperlink{tag-013B}{II.\ref*{cdp:eq:C.TI.2}}}
\expandafter\def\csname SRef@cdp:eq:C.TI.3\endcsname{\hyperlink{tag-013C}{II.\ref*{cdp:eq:C.TI.3}}}
\expandafter\def\csname SRef@cdp:eq:C.TI.4\endcsname{\hyperlink{tag-013D}{II.\ref*{cdp:eq:C.TI.4}}}
\expandafter\def\csname SRef@cdp:eq:C.TI.5\endcsname{\hyperlink{tag-013E}{II.\ref*{cdp:eq:C.TI.5}}}
\expandafter\def\csname SRef@cdp:eq:C.TI.6\endcsname{\hyperlink{tag-013F}{II.\ref*{cdp:eq:C.TI.6}}}
\expandafter\def\csname SRef@cdp:eq:C.TI.7\endcsname{\hyperlink{tag-013G}{II.\ref*{cdp:eq:C.TI.7}}}
\expandafter\def\csname SRef@cdp:eq:C.RG.8\endcsname{\hyperlink{tag-00ZE}{II.\ref*{cdp:eq:C.RG.8}}}
\expandafter\def\csname SRef@cdp:eq:C.TI.8\endcsname{\hyperlink{tag-013H}{II.\ref*{cdp:eq:C.TI.8}}}
\expandafter\def\csname SRef@cdp:eq:C.TI.9\endcsname{\hyperlink{tag-013I}{II.\ref*{cdp:eq:C.TI.9}}}
\expandafter\def\csname SRef@cdp:eq:C.TI.10\endcsname{\hyperlink{tag-0136}{II.\ref*{cdp:eq:C.TI.10}}}
\expandafter\def\csname SRef@cdp:eq:C.TI.11\endcsname{\hyperlink{tag-0137}{II.\ref*{cdp:eq:C.TI.11}}}
\expandafter\def\csname SRef@cdp:eq:C.TI.12\endcsname{\hyperlink{tag-0138}{II.\ref*{cdp:eq:C.TI.12}}}
\expandafter\def\csname SRef@cdp:eq:C.TI.13\endcsname{\hyperlink{tag-0139}{II.\ref*{cdp:eq:C.TI.13}}}
\expandafter\def\csname SRef@cdp:eq:C.TI.14\endcsname{\hyperlink{tag-013A}{II.\ref*{cdp:eq:C.TI.14}}}
\expandafter\def\csname SRef@cdp:eq:C.TK.1\endcsname{\hyperlink{tag-013J}{II.\ref*{cdp:eq:C.TK.1}}}
\expandafter\def\csname SRef@cdp:eq:C.TK.2\endcsname{\hyperlink{tag-013K}{II.\ref*{cdp:eq:C.TK.2}}}
\expandafter\def\csname SRef@cdp:eq:C.TK.3\endcsname{\hyperlink{tag-013L}{II.\ref*{cdp:eq:C.TK.3}}}
\expandafter\def\csname SRef@cdp:eq:C.RG.9\endcsname{\hyperlink{tag-00ZF}{II.\ref*{cdp:eq:C.RG.9}}}
\expandafter\def\csname SRef@cdp:eq:C.TK.4\endcsname{\hyperlink{tag-013M}{II.\ref*{cdp:eq:C.TK.4}}}
\expandafter\def\csname SRef@cdp:eq:C.TK.5\endcsname{\hyperlink{tag-013N}{II.\ref*{cdp:eq:C.TK.5}}}
\expandafter\def\csname SRef@cdp:eq:C.TK.6\endcsname{\hyperlink{tag-013O}{II.\ref*{cdp:eq:C.TK.6}}}
\expandafter\def\csname SRef@cdp:eq:C.TK.7\endcsname{\hyperlink{tag-013P}{II.\ref*{cdp:eq:C.TK.7}}}
\expandafter\def\csname SRef@cdp:eq:C.RK.1\endcsname{\hyperlink{tag-00ZQ}{II.\ref*{cdp:eq:C.RK.1}}}
\expandafter\def\csname SRef@cdp:eq:C.RK.2\endcsname{\hyperlink{tag-0101}{II.\ref*{cdp:eq:C.RK.2}}}
\expandafter\def\csname SRef@cdp:eq:C.RK.3\endcsname{\hyperlink{tag-0105}{II.\ref*{cdp:eq:C.RK.3}}}
\expandafter\def\csname SRef@cdp:eq:C.RK.4\endcsname{\hyperlink{tag-0106}{II.\ref*{cdp:eq:C.RK.4}}}
\expandafter\def\csname SRef@cdp:eq:C.RK.5\endcsname{\hyperlink{tag-0107}{II.\ref*{cdp:eq:C.RK.5}}}
\expandafter\def\csname SRef@cdp:eq:C.RK.6\endcsname{\hyperlink{tag-0108}{II.\ref*{cdp:eq:C.RK.6}}}
\expandafter\def\csname SRef@cdp:eq:C.RG.10\endcsname{\hyperlink{tag-00Z4}{II.\ref*{cdp:eq:C.RG.10}}}
\expandafter\def\csname SRef@cdp:eq:C.RK.7\endcsname{\hyperlink{tag-0109}{II.\ref*{cdp:eq:C.RK.7}}}
\expandafter\def\csname SRef@cdp:eq:C.RK.8\endcsname{\hyperlink{tag-010A}{II.\ref*{cdp:eq:C.RK.8}}}
\expandafter\def\csname SRef@cdp:eq:C.RK.18\endcsname{\hyperlink{tag-00ZZ}{II.\ref*{cdp:eq:C.RK.18}}}
\expandafter\def\csname SRef@cdp:eq:C.RK.19\endcsname{\hyperlink{tag-0100}{II.\ref*{cdp:eq:C.RK.19}}}
\expandafter\def\csname SRef@cdp:eq:C.RK.20\endcsname{\hyperlink{tag-0102}{II.\ref*{cdp:eq:C.RK.20}}}
\expandafter\def\csname SRef@cdp:eq:C.RK.9\endcsname{\hyperlink{tag-010B}{II.\ref*{cdp:eq:C.RK.9}}}
\expandafter\def\csname SRef@cdp:eq:C.RK.10\endcsname{\hyperlink{tag-00ZR}{II.\ref*{cdp:eq:C.RK.10}}}
\expandafter\def\csname SRef@cdp:eq:C.RK.11\endcsname{\hyperlink{tag-00ZS}{II.\ref*{cdp:eq:C.RK.11}}}
\expandafter\def\csname SRef@cdp:eq:C.RK.12\endcsname{\hyperlink{tag-00ZT}{II.\ref*{cdp:eq:C.RK.12}}}
\expandafter\def\csname SRef@cdp:eq:C.RK.13\endcsname{\hyperlink{tag-00ZU}{II.\ref*{cdp:eq:C.RK.13}}}
\expandafter\def\csname SRef@cdp:eq:C.RG.11\endcsname{\hyperlink{tag-00Z5}{II.\ref*{cdp:eq:C.RG.11}}}
\expandafter\def\csname SRef@cdp:eq:C.RK.14\endcsname{\hyperlink{tag-00ZV}{II.\ref*{cdp:eq:C.RK.14}}}
\expandafter\def\csname SRef@cdp:eq:C.RK.15\endcsname{\hyperlink{tag-00ZW}{II.\ref*{cdp:eq:C.RK.15}}}
\expandafter\def\csname SRef@cdp:eq:C.RK.16\endcsname{\hyperlink{tag-00ZX}{II.\ref*{cdp:eq:C.RK.16}}}
\expandafter\def\csname SRef@cdp:eq:C.RK.17\endcsname{\hyperlink{tag-00ZY}{II.\ref*{cdp:eq:C.RK.17}}}
\expandafter\def\csname SRef@cdp:eq:C.RK.21\endcsname{\hyperlink{tag-0103}{II.\ref*{cdp:eq:C.RK.21}}}
\expandafter\def\csname SRef@cdp:eq:C.RK.22\endcsname{\hyperlink{tag-0104}{II.\ref*{cdp:eq:C.RK.22}}}
\expandafter\def\csname SRef@cdp:eq:C.NT.1\endcsname{\hyperlink{tag-00TU}{II.\ref*{cdp:eq:C.NT.1}}}
\expandafter\def\csname SRef@cdp:eq:C.NT.2\endcsname{\hyperlink{tag-00TV}{II.\ref*{cdp:eq:C.NT.2}}}
\expandafter\def\csname SRef@cdp:eq:C.NT.3\endcsname{\hyperlink{tag-00TW}{II.\ref*{cdp:eq:C.NT.3}}}
\expandafter\def\csname SRef@cdp:eq:C.NT.4\endcsname{\hyperlink{tag-00TX}{II.\ref*{cdp:eq:C.NT.4}}}
\expandafter\def\csname SRef@cdp:eq:C.RG.12\endcsname{\hyperlink{tag-00Z6}{II.\ref*{cdp:eq:C.RG.12}}}
\expandafter\def\csname SRef@cdp:eq:C.NT.5\endcsname{\hyperlink{tag-00TY}{II.\ref*{cdp:eq:C.NT.5}}}
\expandafter\def\csname SRef@cdp:eq:C.NT.6\endcsname{\hyperlink{tag-00TZ}{II.\ref*{cdp:eq:C.NT.6}}}
\expandafter\def\csname SRef@cdp:eq:C.NT.7\endcsname{\hyperlink{tag-00U0}{II.\ref*{cdp:eq:C.NT.7}}}
\expandafter\def\csname SRef@cdp:eq:C.NT.8\endcsname{\hyperlink{tag-00U1}{II.\ref*{cdp:eq:C.NT.8}}}
\expandafter\def\csname SRef@cdp:eq:C.GP.1\endcsname{\hyperlink{tag-00KY}{II.\ref*{cdp:eq:C.GP.1}}}
\expandafter\def\csname SRef@cdp:eq:C.GP.2\endcsname{\hyperlink{tag-00L0}{II.\ref*{cdp:eq:C.GP.2}}}
\expandafter\def\csname SRef@cdp:eq:C.GP.3\endcsname{\hyperlink{tag-00L1}{II.\ref*{cdp:eq:C.GP.3}}}
\expandafter\def\csname SRef@cdp:eq:C.GP.4\endcsname{\hyperlink{tag-00L2}{II.\ref*{cdp:eq:C.GP.4}}}
\expandafter\def\csname SRef@cdp:eq:C.GP.5\endcsname{\hyperlink{tag-00L3}{II.\ref*{cdp:eq:C.GP.5}}}
\expandafter\def\csname SRef@cdp:eq:C.GP.6\endcsname{\hyperlink{tag-00L4}{II.\ref*{cdp:eq:C.GP.6}}}
\expandafter\def\csname SRef@cdp:eq:C.RG.13\endcsname{\hyperlink{tag-00Z7}{II.\ref*{cdp:eq:C.RG.13}}}
\expandafter\def\csname SRef@cdp:eq:C.GP.7\endcsname{\hyperlink{tag-00L5}{II.\ref*{cdp:eq:C.GP.7}}}
\expandafter\def\csname SRef@cdp:eq:C.GP.8\endcsname{\hyperlink{tag-00L6}{II.\ref*{cdp:eq:C.GP.8}}}
\expandafter\def\csname SRef@cdp:eq:C.GP.9\endcsname{\hyperlink{tag-00L7}{II.\ref*{cdp:eq:C.GP.9}}}
\expandafter\def\csname SRef@cdp:eq:C.GP.10\endcsname{\hyperlink{tag-00KZ}{II.\ref*{cdp:eq:C.GP.10}}}
\expandafter\def\csname SRef@cdp:eq:C.AR.1\endcsname{\hyperlink{tag-00A0}{II.\ref*{cdp:eq:C.AR.1}}}
\expandafter\def\csname SRef@cdp:eq:C.AR.2\endcsname{\hyperlink{tag-00A8}{II.\ref*{cdp:eq:C.AR.2}}}
\expandafter\def\csname SRef@cdp:eq:C.AR.3\endcsname{\hyperlink{tag-00A9}{II.\ref*{cdp:eq:C.AR.3}}}
\expandafter\def\csname SRef@cdp:eq:C.AR.4\endcsname{\hyperlink{tag-00AA}{II.\ref*{cdp:eq:C.AR.4}}}
\expandafter\def\csname SRef@cdp:eq:C.AR.5\endcsname{\hyperlink{tag-00AB}{II.\ref*{cdp:eq:C.AR.5}}}
\expandafter\def\csname SRef@cdp:eq:C.CS.1\endcsname{\hyperlink{tag-00DH}{II.\ref*{cdp:eq:C.CS.1}}}
\expandafter\def\csname SRef@cdp:eq:C.AR.6\endcsname{\hyperlink{tag-00AC}{II.\ref*{cdp:eq:C.AR.6}}}
\expandafter\def\csname SRef@cdp:eq:C.AR.7\endcsname{\hyperlink{tag-00AD}{II.\ref*{cdp:eq:C.AR.7}}}
\expandafter\def\csname SRef@cdp:eq:C.AR.8\endcsname{\hyperlink{tag-00AE}{II.\ref*{cdp:eq:C.AR.8}}}
\expandafter\def\csname SRef@cdp:eq:C.AR.9\endcsname{\hyperlink{tag-00AF}{II.\ref*{cdp:eq:C.AR.9}}}
\expandafter\def\csname SRef@cdp:eq:C.AR.10\endcsname{\hyperlink{tag-00A1}{II.\ref*{cdp:eq:C.AR.10}}}
\expandafter\def\csname SRef@cdp:eq:C.AR.11\endcsname{\hyperlink{tag-00A2}{II.\ref*{cdp:eq:C.AR.11}}}
\expandafter\def\csname SRef@cdp:eq:C.AR.12\endcsname{\hyperlink{tag-00A3}{II.\ref*{cdp:eq:C.AR.12}}}
\expandafter\def\csname SRef@cdp:eq:C.AR.13\endcsname{\hyperlink{tag-00A4}{II.\ref*{cdp:eq:C.AR.13}}}
\expandafter\def\csname SRef@cdp:eq:C.AR.14\endcsname{\hyperlink{tag-00A5}{II.\ref*{cdp:eq:C.AR.14}}}
\expandafter\def\csname SRef@cdp:eq:C.AR.15\endcsname{\hyperlink{tag-00A6}{II.\ref*{cdp:eq:C.AR.15}}}
\expandafter\def\csname SRef@cdp:eq:C.CS.2\endcsname{\hyperlink{tag-00DI}{II.\ref*{cdp:eq:C.CS.2}}}
\expandafter\def\csname SRef@cdp:eq:C.AR.16\endcsname{\hyperlink{tag-00A7}{II.\ref*{cdp:eq:C.AR.16}}}
\expandafter\def\csname SRef@cdp:eq:C.G.1\endcsname{\hyperlink{tag-00KW}{II.\ref*{cdp:eq:C.G.1}}}
\expandafter\def\csname SRef@cdp:eq:C.G.2\endcsname{\hyperlink{tag-00KX}{II.\ref*{cdp:eq:C.G.2}}}
\expandafter\def\csname SRef@cdp:eq:C.CS.3\endcsname{\hyperlink{tag-00DJ}{II.\ref*{cdp:eq:C.CS.3}}}
\expandafter\def\csname SRef@cdp:eq:C.UN.4\endcsname{\hyperlink{tag-0155}{II.\ref*{cdp:eq:C.UN.4}}}
\expandafter\def\csname SRef@cdp:eq:C.CS.4\endcsname{\hyperlink{tag-00DK}{II.\ref*{cdp:eq:C.CS.4}}}
\expandafter\def\csname SRef@cdp:eq:C.CS.5\endcsname{\hyperlink{tag-00DL}{II.\ref*{cdp:eq:C.CS.5}}}
\expandafter\def\csname SRef@cdp:eq:C.CS.6\endcsname{\hyperlink{tag-00DM}{II.\ref*{cdp:eq:C.CS.6}}}
\expandafter\def\csname SRef@cdp:eq:C.EQ.1\endcsname{\hyperlink{tag-00IC}{II.\ref*{cdp:eq:C.EQ.1}}}
\expandafter\def\csname SRef@cdp:eq:C.EQ.2\endcsname{\hyperlink{tag-00IF}{II.\ref*{cdp:eq:C.EQ.2}}}
\expandafter\def\csname SRef@cdp:eq:C.EQ.3\endcsname{\hyperlink{tag-00IG}{II.\ref*{cdp:eq:C.EQ.3}}}
\expandafter\def\csname SRef@cdp:eq:C.EQ.4\endcsname{\hyperlink{tag-00IH}{II.\ref*{cdp:eq:C.EQ.4}}}
\expandafter\def\csname SRef@cdp:eq:C.EQ.5\endcsname{\hyperlink{tag-00II}{II.\ref*{cdp:eq:C.EQ.5}}}
\expandafter\def\csname SRef@cdp:eq:C.EQ.6\endcsname{\hyperlink{tag-00IJ}{II.\ref*{cdp:eq:C.EQ.6}}}
\expandafter\def\csname SRef@cdp:eq:C.EQ.7\endcsname{\hyperlink{tag-00IK}{II.\ref*{cdp:eq:C.EQ.7}}}
\expandafter\def\csname SRef@cdp:eq:C.UN.5\endcsname{\hyperlink{tag-0156}{II.\ref*{cdp:eq:C.UN.5}}}
\expandafter\def\csname SRef@cdp:eq:C.EQ.8\endcsname{\hyperlink{tag-00IL}{II.\ref*{cdp:eq:C.EQ.8}}}
\expandafter\def\csname SRef@cdp:eq:C.EQ.9\endcsname{\hyperlink{tag-00IM}{II.\ref*{cdp:eq:C.EQ.9}}}
\expandafter\def\csname SRef@cdp:eq:C.EQ.10\endcsname{\hyperlink{tag-00ID}{II.\ref*{cdp:eq:C.EQ.10}}}
\expandafter\def\csname SRef@cdp:eq:C.EQ.11\endcsname{\hyperlink{tag-00IE}{II.\ref*{cdp:eq:C.EQ.11}}}
\expandafter\def\csname SRef@cdp:eq:C.CL.1\endcsname{\hyperlink{tag-00BW}{II.\ref*{cdp:eq:C.CL.1}}}
\expandafter\def\csname SRef@cdp:eq:C.CL.2\endcsname{\hyperlink{tag-00C1}{II.\ref*{cdp:eq:C.CL.2}}}
\expandafter\def\csname SRef@cdp:eq:C.CL.3\endcsname{\hyperlink{tag-00C2}{II.\ref*{cdp:eq:C.CL.3}}}
\expandafter\def\csname SRef@cdp:eq:C.CL.4\endcsname{\hyperlink{tag-00C3}{II.\ref*{cdp:eq:C.CL.4}}}
\expandafter\def\csname SRef@cdp:eq:C.CL.5\endcsname{\hyperlink{tag-00C4}{II.\ref*{cdp:eq:C.CL.5}}}
\expandafter\def\csname SRef@cdp:eq:C.CL.6\endcsname{\hyperlink{tag-00C5}{II.\ref*{cdp:eq:C.CL.6}}}
\expandafter\def\csname SRef@cdp:eq:C.UN.6\endcsname{\hyperlink{tag-0157}{II.\ref*{cdp:eq:C.UN.6}}}
\expandafter\def\csname SRef@cdp:eq:C.CL.7\endcsname{\hyperlink{tag-00C6}{II.\ref*{cdp:eq:C.CL.7}}}
\expandafter\def\csname SRef@cdp:eq:C.CL.8\endcsname{\hyperlink{tag-00C7}{II.\ref*{cdp:eq:C.CL.8}}}
\expandafter\def\csname SRef@cdp:eq:C.CL.9\endcsname{\hyperlink{tag-00C8}{II.\ref*{cdp:eq:C.CL.9}}}
\expandafter\def\csname SRef@cdp:eq:C.CL.10\endcsname{\hyperlink{tag-00BX}{II.\ref*{cdp:eq:C.CL.10}}}
\expandafter\def\csname SRef@cdp:eq:C.CL.11\endcsname{\hyperlink{tag-00BY}{II.\ref*{cdp:eq:C.CL.11}}}
\expandafter\def\csname SRef@cdp:eq:C.CL.12\endcsname{\hyperlink{tag-00BZ}{II.\ref*{cdp:eq:C.CL.12}}}
\expandafter\def\csname SRef@cdp:eq:C.CL.13\endcsname{\hyperlink{tag-00C0}{II.\ref*{cdp:eq:C.CL.13}}}
\expandafter\def\csname SRef@cdp:eq:C.HB.1\endcsname{\hyperlink{tag-00LF}{II.\ref*{cdp:eq:C.HB.1}}}
\expandafter\def\csname SRef@cdp:eq:C.HB.2\endcsname{\hyperlink{tag-00LG}{II.\ref*{cdp:eq:C.HB.2}}}
\expandafter\def\csname SRef@cdp:eq:C.HB.3\endcsname{\hyperlink{tag-00LH}{II.\ref*{cdp:eq:C.HB.3}}}
\expandafter\def\csname SRef@cdp:eq:C.UC.1\endcsname{\hyperlink{tag-014C}{II.\ref*{cdp:eq:C.UC.1}}}
\expandafter\def\csname SRef@cdp:eq:C.HB.4\endcsname{\hyperlink{tag-00LI}{II.\ref*{cdp:eq:C.HB.4}}}
\expandafter\def\csname SRef@cdp:eq:C.HB.5\endcsname{\hyperlink{tag-00LJ}{II.\ref*{cdp:eq:C.HB.5}}}
\expandafter\def\csname SRef@cdp:eq:C.HB.6\endcsname{\hyperlink{tag-00LK}{II.\ref*{cdp:eq:C.HB.6}}}
\expandafter\def\csname SRef@cdp:eq:C.HB.7\endcsname{\hyperlink{tag-00LL}{II.\ref*{cdp:eq:C.HB.7}}}
\expandafter\def\csname SRef@cdp:eq:C.HB.8\endcsname{\hyperlink{tag-00LM}{II.\ref*{cdp:eq:C.HB.8}}}
\expandafter\def\csname SRef@cdp:eq:C.HB.9\endcsname{\hyperlink{tag-00LN}{II.\ref*{cdp:eq:C.HB.9}}}
\expandafter\def\csname SRef@cdp:eq:C.ZC.1\endcsname{\hyperlink{tag-018W}{II.\ref*{cdp:eq:C.ZC.1}}}
\expandafter\def\csname SRef@cdp:eq:C.ZC.2\endcsname{\hyperlink{tag-018X}{II.\ref*{cdp:eq:C.ZC.2}}}
\expandafter\def\csname SRef@cdp:eq:C.ZC.3\endcsname{\hyperlink{tag-018Y}{II.\ref*{cdp:eq:C.ZC.3}}}
\expandafter\def\csname SRef@cdp:eq:C.ZC.4\endcsname{\hyperlink{tag-018Z}{II.\ref*{cdp:eq:C.ZC.4}}}
\expandafter\def\csname SRef@cdp:eq:C.UC.2\endcsname{\hyperlink{tag-014E}{II.\ref*{cdp:eq:C.UC.2}}}
\expandafter\def\csname SRef@cdp:eq:C.ZC.5\endcsname{\hyperlink{tag-0190}{II.\ref*{cdp:eq:C.ZC.5}}}
\expandafter\def\csname SRef@cdp:eq:C.ZC.6\endcsname{\hyperlink{tag-0191}{II.\ref*{cdp:eq:C.ZC.6}}}
\expandafter\def\csname SRef@cdp:eq:C.ZC.7\endcsname{\hyperlink{tag-0192}{II.\ref*{cdp:eq:C.ZC.7}}}
\expandafter\def\csname SRef@cdp:eq:C.ZC.8\endcsname{\hyperlink{tag-0193}{II.\ref*{cdp:eq:C.ZC.8}}}
\expandafter\def\csname SRef@cdp:eq:C.ZC.9\endcsname{\hyperlink{tag-0194}{II.\ref*{cdp:eq:C.ZC.9}}}
\expandafter\def\csname SRef@cdp:eq:C.FG.1\endcsname{\hyperlink{tag-00K8}{II.\ref*{cdp:eq:C.FG.1}}}
\expandafter\def\csname SRef@cdp:eq:C.FG.2\endcsname{\hyperlink{tag-00KJ}{II.\ref*{cdp:eq:C.FG.2}}}
\expandafter\def\csname SRef@cdp:eq:C.FG.3\endcsname{\hyperlink{tag-00KK}{II.\ref*{cdp:eq:C.FG.3}}}
\expandafter\def\csname SRef@cdp:eq:C.FG.4\endcsname{\hyperlink{tag-00KL}{II.\ref*{cdp:eq:C.FG.4}}}
\expandafter\def\csname SRef@cdp:eq:C.FG.5\endcsname{\hyperlink{tag-00KM}{II.\ref*{cdp:eq:C.FG.5}}}
\expandafter\def\csname SRef@cdp:eq:C.UC.3\endcsname{\hyperlink{tag-014F}{II.\ref*{cdp:eq:C.UC.3}}}
\expandafter\def\csname SRef@cdp:eq:C.FG.6\endcsname{\hyperlink{tag-00KN}{II.\ref*{cdp:eq:C.FG.6}}}
\expandafter\def\csname SRef@cdp:eq:C.FG.7\endcsname{\hyperlink{tag-00KO}{II.\ref*{cdp:eq:C.FG.7}}}
\expandafter\def\csname SRef@cdp:eq:C.FG.8\endcsname{\hyperlink{tag-00KP}{II.\ref*{cdp:eq:C.FG.8}}}
\expandafter\def\csname SRef@cdp:eq:C.FG.9\endcsname{\hyperlink{tag-00KQ}{II.\ref*{cdp:eq:C.FG.9}}}
\expandafter\def\csname SRef@cdp:eq:C.FG.10\endcsname{\hyperlink{tag-00K9}{II.\ref*{cdp:eq:C.FG.10}}}
\expandafter\def\csname SRef@cdp:eq:C.FG.11\endcsname{\hyperlink{tag-00KA}{II.\ref*{cdp:eq:C.FG.11}}}
\expandafter\def\csname SRef@cdp:eq:C.FG.12\endcsname{\hyperlink{tag-00KB}{II.\ref*{cdp:eq:C.FG.12}}}
\expandafter\def\csname SRef@cdp:eq:C.FG.13\endcsname{\hyperlink{tag-00KC}{II.\ref*{cdp:eq:C.FG.13}}}
\expandafter\def\csname SRef@cdp:eq:C.FG.14\endcsname{\hyperlink{tag-00KD}{II.\ref*{cdp:eq:C.FG.14}}}
\expandafter\def\csname SRef@cdp:eq:C.FG.15\endcsname{\hyperlink{tag-00KE}{II.\ref*{cdp:eq:C.FG.15}}}
\expandafter\def\csname SRef@cdp:eq:C.UC.4\endcsname{\hyperlink{tag-014G}{II.\ref*{cdp:eq:C.UC.4}}}
\expandafter\def\csname SRef@cdp:eq:C.FG.16\endcsname{\hyperlink{tag-00KF}{II.\ref*{cdp:eq:C.FG.16}}}
\expandafter\def\csname SRef@cdp:eq:C.FG.17\endcsname{\hyperlink{tag-00KG}{II.\ref*{cdp:eq:C.FG.17}}}
\expandafter\def\csname SRef@cdp:eq:C.FG.18\endcsname{\hyperlink{tag-00KH}{II.\ref*{cdp:eq:C.FG.18}}}
\expandafter\def\csname SRef@cdp:eq:C.FG.19\endcsname{\hyperlink{tag-00KI}{II.\ref*{cdp:eq:C.FG.19}}}
\expandafter\def\csname SRef@cdp:eq:C.PF.1\endcsname{\hyperlink{tag-00UN}{II.\ref*{cdp:eq:C.PF.1}}}
\expandafter\def\csname SRef@cdp:eq:C.PF.2\endcsname{\hyperlink{tag-00UR}{II.\ref*{cdp:eq:C.PF.2}}}
\expandafter\def\csname SRef@cdp:eq:C.PF.3\endcsname{\hyperlink{tag-00US}{II.\ref*{cdp:eq:C.PF.3}}}
\expandafter\def\csname SRef@cdp:eq:C.PF.4\endcsname{\hyperlink{tag-00UT}{II.\ref*{cdp:eq:C.PF.4}}}
\expandafter\def\csname SRef@cdp:eq:C.PF.5\endcsname{\hyperlink{tag-00UU}{II.\ref*{cdp:eq:C.PF.5}}}
\expandafter\def\csname SRef@cdp:eq:C.PF.6\endcsname{\hyperlink{tag-00UV}{II.\ref*{cdp:eq:C.PF.6}}}
\expandafter\def\csname SRef@cdp:eq:C.UC.5\endcsname{\hyperlink{tag-014H}{II.\ref*{cdp:eq:C.UC.5}}}
\expandafter\def\csname SRef@cdp:eq:C.PF.7\endcsname{\hyperlink{tag-00UW}{II.\ref*{cdp:eq:C.PF.7}}}
\expandafter\def\csname SRef@cdp:eq:C.PF.8\endcsname{\hyperlink{tag-00UX}{II.\ref*{cdp:eq:C.PF.8}}}
\expandafter\def\csname SRef@cdp:eq:C.PF.9\endcsname{\hyperlink{tag-00UY}{II.\ref*{cdp:eq:C.PF.9}}}
\expandafter\def\csname SRef@cdp:eq:C.PF.10\endcsname{\hyperlink{tag-00UO}{II.\ref*{cdp:eq:C.PF.10}}}
\expandafter\def\csname SRef@cdp:eq:C.PF.11\endcsname{\hyperlink{tag-00UP}{II.\ref*{cdp:eq:C.PF.11}}}
\expandafter\def\csname SRef@cdp:eq:C.PF.12\endcsname{\hyperlink{tag-00UQ}{II.\ref*{cdp:eq:C.PF.12}}}
\expandafter\def\csname SRef@cdp:eq:C.UNN.1\endcsname{\hyperlink{tag-015B}{II.\ref*{cdp:eq:C.UNN.1}}}
\expandafter\def\csname SRef@cdp:eq:C.UNN.2\endcsname{\hyperlink{tag-015C}{II.\ref*{cdp:eq:C.UNN.2}}}
\expandafter\def\csname SRef@cdp:eq:C.CR.1\endcsname{\hyperlink{tag-00CV}{II.\ref*{cdp:eq:C.CR.1}}}
\expandafter\def\csname SRef@cdp:eq:C.CR.2\endcsname{\hyperlink{tag-00D6}{II.\ref*{cdp:eq:C.CR.2}}}
\expandafter\def\csname SRef@cdp:eq:C.UC.6\endcsname{\hyperlink{tag-014I}{II.\ref*{cdp:eq:C.UC.6}}}
\expandafter\def\csname SRef@cdp:eq:C.CR.3\endcsname{\hyperlink{tag-00DA}{II.\ref*{cdp:eq:C.CR.3}}}
\expandafter\def\csname SRef@cdp:eq:C.CR.4\endcsname{\hyperlink{tag-00DB}{II.\ref*{cdp:eq:C.CR.4}}}
\expandafter\def\csname SRef@cdp:eq:C.CR.5\endcsname{\hyperlink{tag-00DC}{II.\ref*{cdp:eq:C.CR.5}}}
\expandafter\def\csname SRef@cdp:eq:C.CR.6\endcsname{\hyperlink{tag-00DD}{II.\ref*{cdp:eq:C.CR.6}}}
\expandafter\def\csname SRef@cdp:eq:C.CR.7\endcsname{\hyperlink{tag-00DE}{II.\ref*{cdp:eq:C.CR.7}}}
\expandafter\def\csname SRef@cdp:eq:C.CR.8\endcsname{\hyperlink{tag-00DF}{II.\ref*{cdp:eq:C.CR.8}}}
\expandafter\def\csname SRef@cdp:eq:C.CR.9\endcsname{\hyperlink{tag-00DG}{II.\ref*{cdp:eq:C.CR.9}}}
\expandafter\def\csname SRef@cdp:eq:C.CR.10\endcsname{\hyperlink{tag-00CW}{II.\ref*{cdp:eq:C.CR.10}}}
\expandafter\def\csname SRef@cdp:eq:C.CR.11\endcsname{\hyperlink{tag-00CX}{II.\ref*{cdp:eq:C.CR.11}}}
\expandafter\def\csname SRef@cdp:eq:C.CR.12\endcsname{\hyperlink{tag-00CY}{II.\ref*{cdp:eq:C.CR.12}}}
\expandafter\def\csname SRef@cdp:eq:C.UC.7\endcsname{\hyperlink{tag-014J}{II.\ref*{cdp:eq:C.UC.7}}}
\expandafter\def\csname SRef@cdp:eq:C.CR.13\endcsname{\hyperlink{tag-00CZ}{II.\ref*{cdp:eq:C.CR.13}}}
\expandafter\def\csname SRef@cdp:eq:C.CR.14\endcsname{\hyperlink{tag-00D0}{II.\ref*{cdp:eq:C.CR.14}}}
\expandafter\def\csname SRef@cdp:eq:C.CR.15\endcsname{\hyperlink{tag-00D1}{II.\ref*{cdp:eq:C.CR.15}}}
\expandafter\def\csname SRef@cdp:eq:C.CR.16\endcsname{\hyperlink{tag-00D2}{II.\ref*{cdp:eq:C.CR.16}}}
\expandafter\def\csname SRef@cdp:eq:C.CR.17\endcsname{\hyperlink{tag-00D3}{II.\ref*{cdp:eq:C.CR.17}}}
\expandafter\def\csname SRef@cdp:eq:C.CR.18\endcsname{\hyperlink{tag-00D4}{II.\ref*{cdp:eq:C.CR.18}}}
\expandafter\def\csname SRef@cdp:eq:C.CR.19\endcsname{\hyperlink{tag-00D5}{II.\ref*{cdp:eq:C.CR.19}}}
\expandafter\def\csname SRef@cdp:eq:C.CR.20\endcsname{\hyperlink{tag-00D7}{II.\ref*{cdp:eq:C.CR.20}}}
\expandafter\def\csname SRef@cdp:eq:C.CR.21\endcsname{\hyperlink{tag-00D8}{II.\ref*{cdp:eq:C.CR.21}}}
\expandafter\def\csname SRef@cdp:eq:C.CR.22\endcsname{\hyperlink{tag-00D9}{II.\ref*{cdp:eq:C.CR.22}}}
\expandafter\def\csname SRef@cdp:eq:C.UC.8\endcsname{\hyperlink{tag-014K}{II.\ref*{cdp:eq:C.UC.8}}}
\expandafter\def\csname SRef@cdp:eq:C.UNNODE.1\endcsname{\hyperlink{tag-015D}{II.\ref*{cdp:eq:C.UNNODE.1}}}
\expandafter\def\csname SRef@cdp:eq:C.UNNODE.2\endcsname{\hyperlink{tag-015E}{II.\ref*{cdp:eq:C.UNNODE.2}}}
\expandafter\def\csname SRef@cdp:eq:C.UNNODE.3\endcsname{\hyperlink{tag-015F}{II.\ref*{cdp:eq:C.UNNODE.3}}}
\expandafter\def\csname SRef@cdp:eq:C.UNCR.1\endcsname{\hyperlink{tag-0158}{II.\ref*{cdp:eq:C.UNCR.1}}}
\expandafter\def\csname SRef@cdp:eq:C.UNCR.2\endcsname{\hyperlink{tag-0159}{II.\ref*{cdp:eq:C.UNCR.2}}}
\expandafter\def\csname SRef@cdp:eq:C.UNCR.3\endcsname{\hyperlink{tag-015A}{II.\ref*{cdp:eq:C.UNCR.3}}}
\expandafter\def\csname SRef@cdp:eq:C.IR.1\endcsname{\hyperlink{tag-00OH}{II.\ref*{cdp:eq:C.IR.1}}}
\expandafter\def\csname SRef@cdp:eq:C.IR.2\endcsname{\hyperlink{tag-00OI}{II.\ref*{cdp:eq:C.IR.2}}}
\expandafter\def\csname SRef@cdp:eq:C.IR.4\endcsname{\hyperlink{tag-00OJ}{II.\ref*{cdp:eq:C.IR.4}}}
\expandafter\def\csname SRef@cdp:eq:C.IR.5\endcsname{\hyperlink{tag-00OK}{II.\ref*{cdp:eq:C.IR.5}}}
\expandafter\def\csname SRef@cdp:eq:C.UC.9\endcsname{\hyperlink{tag-014L}{II.\ref*{cdp:eq:C.UC.9}}}
\expandafter\def\csname SRef@cdp:eq:C.IR.6\endcsname{\hyperlink{tag-00OL}{II.\ref*{cdp:eq:C.IR.6}}}
\expandafter\def\csname SRef@cdp:eq:C.IR.9\endcsname{\hyperlink{tag-00OM}{II.\ref*{cdp:eq:C.IR.9}}}
\expandafter\def\csname SRef@cdp:eq:C.VII.1\endcsname{\hyperlink{tag-015K}{II.\ref*{cdp:eq:C.VII.1}}}
\expandafter\def\csname SRef@cdp:eq:C.VII.2\endcsname{\hyperlink{tag-015O}{II.\ref*{cdp:eq:C.VII.2}}}
\expandafter\def\csname SRef@cdp:eq:C.VII.3\endcsname{\hyperlink{tag-015P}{II.\ref*{cdp:eq:C.VII.3}}}
\expandafter\def\csname SRef@cdp:eq:C.VII.4\endcsname{\hyperlink{tag-015Q}{II.\ref*{cdp:eq:C.VII.4}}}
\expandafter\def\csname SRef@cdp:eq:C.VII.5\endcsname{\hyperlink{tag-015R}{II.\ref*{cdp:eq:C.VII.5}}}
\expandafter\def\csname SRef@cdp:eq:C.VII.6\endcsname{\hyperlink{tag-015S}{II.\ref*{cdp:eq:C.VII.6}}}
\expandafter\def\csname SRef@cdp:eq:C.VII.7\endcsname{\hyperlink{tag-015T}{II.\ref*{cdp:eq:C.VII.7}}}
\expandafter\def\csname SRef@cdp:eq:C.VII.8\endcsname{\hyperlink{tag-015U}{II.\ref*{cdp:eq:C.VII.8}}}
\expandafter\def\csname SRef@cdp:eq:C.UC.10\endcsname{\hyperlink{tag-014D}{II.\ref*{cdp:eq:C.UC.10}}}
\expandafter\def\csname SRef@cdp:eq:C.VII.9\endcsname{\hyperlink{tag-015V}{II.\ref*{cdp:eq:C.VII.9}}}
\expandafter\def\csname SRef@cdp:eq:C.VII.10\endcsname{\hyperlink{tag-015L}{II.\ref*{cdp:eq:C.VII.10}}}
\expandafter\def\csname SRef@cdp:eq:C.VII.11\endcsname{\hyperlink{tag-015M}{II.\ref*{cdp:eq:C.VII.11}}}
\expandafter\def\csname SRef@cdp:eq:C.VII.12\endcsname{\hyperlink{tag-015N}{II.\ref*{cdp:eq:C.VII.12}}}
\expandafter\def\csname SRef@cdp:eq:C.RA.1\endcsname{\hyperlink{tag-00WQ}{II.\ref*{cdp:eq:C.RA.1}}}
\expandafter\def\csname SRef@cdp:eq:C.RA.2\endcsname{\hyperlink{tag-00X1}{II.\ref*{cdp:eq:C.RA.2}}}
\expandafter\def\csname SRef@cdp:eq:C.RA.3\endcsname{\hyperlink{tag-00XC}{II.\ref*{cdp:eq:C.RA.3}}}
\expandafter\def\csname SRef@cdp:eq:C.RA.4\endcsname{\hyperlink{tag-00XH}{II.\ref*{cdp:eq:C.RA.4}}}
\expandafter\def\csname SRef@cdp:eq:C.RA.5\endcsname{\hyperlink{tag-00XI}{II.\ref*{cdp:eq:C.RA.5}}}
\expandafter\def\csname SRef@cdp:eq:C.RA.6\endcsname{\hyperlink{tag-00XJ}{II.\ref*{cdp:eq:C.RA.6}}}
\expandafter\def\csname SRef@cdp:eq:C.VR.1\endcsname{\hyperlink{tag-015W}{II.\ref*{cdp:eq:C.VR.1}}}
\expandafter\def\csname SRef@cdp:eq:C.RA.7\endcsname{\hyperlink{tag-00XK}{II.\ref*{cdp:eq:C.RA.7}}}
\expandafter\def\csname SRef@cdp:eq:C.RA.8\endcsname{\hyperlink{tag-00XL}{II.\ref*{cdp:eq:C.RA.8}}}
\expandafter\def\csname SRef@cdp:eq:C.RA.9\endcsname{\hyperlink{tag-00XM}{II.\ref*{cdp:eq:C.RA.9}}}
\expandafter\def\csname SRef@cdp:eq:C.RA.10\endcsname{\hyperlink{tag-00WR}{II.\ref*{cdp:eq:C.RA.10}}}
\expandafter\def\csname SRef@cdp:eq:C.RA.11\endcsname{\hyperlink{tag-00WS}{II.\ref*{cdp:eq:C.RA.11}}}
\expandafter\def\csname SRef@cdp:eq:C.RA.12\endcsname{\hyperlink{tag-00WT}{II.\ref*{cdp:eq:C.RA.12}}}
\expandafter\def\csname SRef@cdp:eq:C.RA.13\endcsname{\hyperlink{tag-00WU}{II.\ref*{cdp:eq:C.RA.13}}}
\expandafter\def\csname SRef@cdp:eq:C.RA.14\endcsname{\hyperlink{tag-00WV}{II.\ref*{cdp:eq:C.RA.14}}}
\expandafter\def\csname SRef@cdp:eq:C.RA.15\endcsname{\hyperlink{tag-00WW}{II.\ref*{cdp:eq:C.RA.15}}}
\expandafter\def\csname SRef@cdp:eq:C.RA.16\endcsname{\hyperlink{tag-00WX}{II.\ref*{cdp:eq:C.RA.16}}}
\expandafter\def\csname SRef@cdp:eq:C.VR.2\endcsname{\hyperlink{tag-015X}{II.\ref*{cdp:eq:C.VR.2}}}
\expandafter\def\csname SRef@cdp:eq:C.RA.17\endcsname{\hyperlink{tag-00WY}{II.\ref*{cdp:eq:C.RA.17}}}
\expandafter\def\csname SRef@cdp:eq:C.RA.18\endcsname{\hyperlink{tag-00WZ}{II.\ref*{cdp:eq:C.RA.18}}}
\expandafter\def\csname SRef@cdp:eq:C.RA.19\endcsname{\hyperlink{tag-00X0}{II.\ref*{cdp:eq:C.RA.19}}}
\expandafter\def\csname SRef@cdp:eq:C.RA.20\endcsname{\hyperlink{tag-00X2}{II.\ref*{cdp:eq:C.RA.20}}}
\expandafter\def\csname SRef@cdp:eq:C.RA.21\endcsname{\hyperlink{tag-00X3}{II.\ref*{cdp:eq:C.RA.21}}}
\expandafter\def\csname SRef@cdp:eq:C.RA.22\endcsname{\hyperlink{tag-00X4}{II.\ref*{cdp:eq:C.RA.22}}}
\expandafter\def\csname SRef@cdp:eq:C.RA.23\endcsname{\hyperlink{tag-00X5}{II.\ref*{cdp:eq:C.RA.23}}}
\expandafter\def\csname SRef@cdp:eq:C.RA.24\endcsname{\hyperlink{tag-00X6}{II.\ref*{cdp:eq:C.RA.24}}}
\expandafter\def\csname SRef@cdp:eq:C.RA.25\endcsname{\hyperlink{tag-00X7}{II.\ref*{cdp:eq:C.RA.25}}}
\expandafter\def\csname SRef@cdp:eq:C.RA.26\endcsname{\hyperlink{tag-00X8}{II.\ref*{cdp:eq:C.RA.26}}}
\expandafter\def\csname SRef@cdp:eq:C.VR.3\endcsname{\hyperlink{tag-015Y}{II.\ref*{cdp:eq:C.VR.3}}}
\expandafter\def\csname SRef@cdp:eq:C.RA.27\endcsname{\hyperlink{tag-00X9}{II.\ref*{cdp:eq:C.RA.27}}}
\expandafter\def\csname SRef@cdp:eq:C.RA.28\endcsname{\hyperlink{tag-00XA}{II.\ref*{cdp:eq:C.RA.28}}}
\expandafter\def\csname SRef@cdp:eq:C.RA.29\endcsname{\hyperlink{tag-00XB}{II.\ref*{cdp:eq:C.RA.29}}}
\expandafter\def\csname SRef@cdp:eq:C.RA.30\endcsname{\hyperlink{tag-00XD}{II.\ref*{cdp:eq:C.RA.30}}}
\expandafter\def\csname SRef@cdp:eq:C.RA.31\endcsname{\hyperlink{tag-00XE}{II.\ref*{cdp:eq:C.RA.31}}}
\expandafter\def\csname SRef@cdp:eq:C.RA.31a\endcsname{\hyperlink{tag-00XF}{II.\ref*{cdp:eq:C.RA.31a}}}
\expandafter\def\csname SRef@cdp:eq:C.RA.32\endcsname{\hyperlink{tag-00XG}{II.\ref*{cdp:eq:C.RA.32}}}
\expandafter\def\csname SRef@cdp:eq:C.RD.1\endcsname{\hyperlink{tag-00XZ}{II.\ref*{cdp:eq:C.RD.1}}}
\expandafter\def\csname SRef@cdp:eq:C.RD.2\endcsname{\hyperlink{tag-00Y9}{II.\ref*{cdp:eq:C.RD.2}}}
\expandafter\def\csname SRef@cdp:eq:C.RD.3\endcsname{\hyperlink{tag-00YK}{II.\ref*{cdp:eq:C.RD.3}}}
\expandafter\def\csname SRef@cdp:eq:C.VR.4\endcsname{\hyperlink{tag-015Z}{II.\ref*{cdp:eq:C.VR.4}}}
\expandafter\def\csname SRef@cdp:eq:C.RD.4\endcsname{\hyperlink{tag-00YT}{II.\ref*{cdp:eq:C.RD.4}}}
\expandafter\def\csname SRef@cdp:eq:C.RD.11\endcsname{\hyperlink{tag-00Y0}{II.\ref*{cdp:eq:C.RD.11}}}
\expandafter\def\csname SRef@cdp:eq:C.RD.12\endcsname{\hyperlink{tag-00Y1}{II.\ref*{cdp:eq:C.RD.12}}}
\expandafter\def\csname SRef@cdp:eq:C.RD.13\endcsname{\hyperlink{tag-00Y2}{II.\ref*{cdp:eq:C.RD.13}}}
\expandafter\def\csname SRef@cdp:eq:C.RD.14\endcsname{\hyperlink{tag-00Y3}{II.\ref*{cdp:eq:C.RD.14}}}
\expandafter\def\csname SRef@cdp:eq:C.RD.15\endcsname{\hyperlink{tag-00Y4}{II.\ref*{cdp:eq:C.RD.15}}}
\expandafter\def\csname SRef@cdp:eq:C.RD.16\endcsname{\hyperlink{tag-00Y5}{II.\ref*{cdp:eq:C.RD.16}}}
\expandafter\def\csname SRef@cdp:eq:C.RD.17\endcsname{\hyperlink{tag-00Y6}{II.\ref*{cdp:eq:C.RD.17}}}
\expandafter\def\csname SRef@cdp:eq:C.RD.18\endcsname{\hyperlink{tag-00Y7}{II.\ref*{cdp:eq:C.RD.18}}}
\expandafter\def\csname SRef@cdp:eq:C.VR.5\endcsname{\hyperlink{tag-0160}{II.\ref*{cdp:eq:C.VR.5}}}
\expandafter\def\csname SRef@cdp:eq:C.RD.19\endcsname{\hyperlink{tag-00Y8}{II.\ref*{cdp:eq:C.RD.19}}}
\expandafter\def\csname SRef@cdp:eq:C.RD.20\endcsname{\hyperlink{tag-00YA}{II.\ref*{cdp:eq:C.RD.20}}}
\expandafter\def\csname SRef@cdp:eq:C.RD.21\endcsname{\hyperlink{tag-00YB}{II.\ref*{cdp:eq:C.RD.21}}}
\expandafter\def\csname SRef@cdp:eq:C.RD.22\endcsname{\hyperlink{tag-00YC}{II.\ref*{cdp:eq:C.RD.22}}}
\expandafter\def\csname SRef@cdp:eq:C.RD.23\endcsname{\hyperlink{tag-00YD}{II.\ref*{cdp:eq:C.RD.23}}}
\expandafter\def\csname SRef@cdp:eq:C.RD.24\endcsname{\hyperlink{tag-00YE}{II.\ref*{cdp:eq:C.RD.24}}}
\expandafter\def\csname SRef@cdp:eq:C.RD.25\endcsname{\hyperlink{tag-00YF}{II.\ref*{cdp:eq:C.RD.25}}}
\expandafter\def\csname SRef@cdp:eq:C.RD.26\endcsname{\hyperlink{tag-00YG}{II.\ref*{cdp:eq:C.RD.26}}}
\expandafter\def\csname SRef@cdp:eq:C.RD.27\endcsname{\hyperlink{tag-00YH}{II.\ref*{cdp:eq:C.RD.27}}}
\expandafter\def\csname SRef@cdp:eq:C.VR.6\endcsname{\hyperlink{tag-0161}{II.\ref*{cdp:eq:C.VR.6}}}
\expandafter\def\csname SRef@cdp:eq:C.RD.28\endcsname{\hyperlink{tag-00YI}{II.\ref*{cdp:eq:C.RD.28}}}
\expandafter\def\csname SRef@cdp:eq:C.RD.29\endcsname{\hyperlink{tag-00YJ}{II.\ref*{cdp:eq:C.RD.29}}}
\expandafter\def\csname SRef@cdp:eq:C.RD.30\endcsname{\hyperlink{tag-00YL}{II.\ref*{cdp:eq:C.RD.30}}}
\expandafter\def\csname SRef@cdp:eq:C.RD.31\endcsname{\hyperlink{tag-00YM}{II.\ref*{cdp:eq:C.RD.31}}}
\expandafter\def\csname SRef@cdp:eq:C.RD.32\endcsname{\hyperlink{tag-00YN}{II.\ref*{cdp:eq:C.RD.32}}}
\expandafter\def\csname SRef@cdp:eq:C.RD.33\endcsname{\hyperlink{tag-00YO}{II.\ref*{cdp:eq:C.RD.33}}}
\expandafter\def\csname SRef@cdp:eq:C.RD.34\endcsname{\hyperlink{tag-00YP}{II.\ref*{cdp:eq:C.RD.34}}}
\expandafter\def\csname SRef@cdp:eq:C.RD.35\endcsname{\hyperlink{tag-00YQ}{II.\ref*{cdp:eq:C.RD.35}}}
\expandafter\def\csname SRef@cdp:eq:C.VR.7\endcsname{\hyperlink{tag-0162}{II.\ref*{cdp:eq:C.VR.7}}}
\expandafter\def\csname SRef@cdp:eq:C.RD.36\endcsname{\hyperlink{tag-00YR}{II.\ref*{cdp:eq:C.RD.36}}}
\expandafter\def\csname SRef@cdp:eq:C.RD.37\endcsname{\hyperlink{tag-00YS}{II.\ref*{cdp:eq:C.RD.37}}}
\expandafter\def\csname SRef@cdp:eq:C.D4.1\endcsname{\hyperlink{tag-00E1}{II.\ref*{cdp:eq:C.D4.1}}}
\expandafter\def\csname SRef@cdp:eq:C.D4.2\endcsname{\hyperlink{tag-00EC}{II.\ref*{cdp:eq:C.D4.2}}}
\expandafter\def\csname SRef@cdp:eq:C.D4.3\endcsname{\hyperlink{tag-00EE}{II.\ref*{cdp:eq:C.D4.3}}}
\expandafter\def\csname SRef@cdp:eq:C.D4.4\endcsname{\hyperlink{tag-00EF}{II.\ref*{cdp:eq:C.D4.4}}}
\expandafter\def\csname SRef@cdp:eq:C.D4.5\endcsname{\hyperlink{tag-00EG}{II.\ref*{cdp:eq:C.D4.5}}}
\expandafter\def\csname SRef@cdp:eq:C.VR.8\endcsname{\hyperlink{tag-0163}{II.\ref*{cdp:eq:C.VR.8}}}
\expandafter\def\csname SRef@cdp:eq:C.D4.6\endcsname{\hyperlink{tag-00EH}{II.\ref*{cdp:eq:C.D4.6}}}
\expandafter\def\csname SRef@cdp:eq:C.D4.7\endcsname{\hyperlink{tag-00EI}{II.\ref*{cdp:eq:C.D4.7}}}
\expandafter\def\csname SRef@cdp:eq:C.D4.8\endcsname{\hyperlink{tag-00EJ}{II.\ref*{cdp:eq:C.D4.8}}}
\expandafter\def\csname SRef@cdp:eq:C.D4.9\endcsname{\hyperlink{tag-00EK}{II.\ref*{cdp:eq:C.D4.9}}}
\expandafter\def\csname SRef@cdp:eq:C.D4.10\endcsname{\hyperlink{tag-00E2}{II.\ref*{cdp:eq:C.D4.10}}}
\expandafter\def\csname SRef@cdp:eq:C.D4.11\endcsname{\hyperlink{tag-00E3}{II.\ref*{cdp:eq:C.D4.11}}}
\expandafter\def\csname SRef@cdp:eq:C.D4.12\endcsname{\hyperlink{tag-00E4}{II.\ref*{cdp:eq:C.D4.12}}}
\expandafter\def\csname SRef@cdp:eq:C.DD.1\endcsname{\hyperlink{tag-00EL}{II.\ref*{cdp:eq:C.DD.1}}}
\expandafter\def\csname SRef@cdp:eq:C.DD.2\endcsname{\hyperlink{tag-00EW}{II.\ref*{cdp:eq:C.DD.2}}}
\expandafter\def\csname SRef@cdp:eq:C.VR.9\endcsname{\hyperlink{tag-0164}{II.\ref*{cdp:eq:C.VR.9}}}
\expandafter\def\csname SRef@cdp:eq:C.DD.3\endcsname{\hyperlink{tag-00EZ}{II.\ref*{cdp:eq:C.DD.3}}}
\expandafter\def\csname SRef@cdp:eq:C.DD.4\endcsname{\hyperlink{tag-00F0}{II.\ref*{cdp:eq:C.DD.4}}}
\expandafter\def\csname SRef@cdp:eq:C.DD.5\endcsname{\hyperlink{tag-00F1}{II.\ref*{cdp:eq:C.DD.5}}}
\expandafter\def\csname SRef@cdp:eq:C.DD.6\endcsname{\hyperlink{tag-00F2}{II.\ref*{cdp:eq:C.DD.6}}}
\expandafter\def\csname SRef@cdp:eq:C.DD.7\endcsname{\hyperlink{tag-00F3}{II.\ref*{cdp:eq:C.DD.7}}}
\expandafter\def\csname SRef@cdp:eq:C.DD.8\endcsname{\hyperlink{tag-00F4}{II.\ref*{cdp:eq:C.DD.8}}}
\expandafter\def\csname SRef@cdp:eq:C.DD.9\endcsname{\hyperlink{tag-00F5}{II.\ref*{cdp:eq:C.DD.9}}}
\expandafter\def\csname SRef@cdp:eq:C.DD.10\endcsname{\hyperlink{tag-00EM}{II.\ref*{cdp:eq:C.DD.10}}}
\expandafter\def\csname SRef@cdp:eq:C.EC.1\endcsname{\hyperlink{tag-00FO}{II.\ref*{cdp:eq:C.EC.1}}}
\expandafter\def\csname SRef@cdp:eq:C.DD.11\endcsname{\hyperlink{tag-00EN}{II.\ref*{cdp:eq:C.DD.11}}}
\expandafter\def\csname SRef@cdp:eq:C.DD.12\endcsname{\hyperlink{tag-00EO}{II.\ref*{cdp:eq:C.DD.12}}}
\expandafter\def\csname SRef@cdp:eq:C.DD.13\endcsname{\hyperlink{tag-00EP}{II.\ref*{cdp:eq:C.DD.13}}}
\expandafter\def\csname SRef@cdp:eq:C.DD.14\endcsname{\hyperlink{tag-00EQ}{II.\ref*{cdp:eq:C.DD.14}}}
\expandafter\def\csname SRef@cdp:eq:C.DD.15\endcsname{\hyperlink{tag-00ER}{II.\ref*{cdp:eq:C.DD.15}}}
\expandafter\def\csname SRef@cdp:eq:C.DD.16\endcsname{\hyperlink{tag-00ES}{II.\ref*{cdp:eq:C.DD.16}}}
\expandafter\def\csname SRef@cdp:eq:C.EC.2\endcsname{\hyperlink{tag-00FP}{II.\ref*{cdp:eq:C.EC.2}}}
\expandafter\def\csname SRef@cdp:eq:C.DD.17\endcsname{\hyperlink{tag-00ET}{II.\ref*{cdp:eq:C.DD.17}}}
\expandafter\def\csname SRef@cdp:eq:C.DD.18\endcsname{\hyperlink{tag-00EU}{II.\ref*{cdp:eq:C.DD.18}}}
\expandafter\def\csname SRef@cdp:eq:C.DD.19\endcsname{\hyperlink{tag-00EV}{II.\ref*{cdp:eq:C.DD.19}}}
\expandafter\def\csname SRef@cdp:eq:C.DD.20\endcsname{\hyperlink{tag-00EX}{II.\ref*{cdp:eq:C.DD.20}}}
\expandafter\def\csname SRef@cdp:eq:C.DD.21\endcsname{\hyperlink{tag-00EY}{II.\ref*{cdp:eq:C.DD.21}}}
\expandafter\def\csname SRef@cdp:eq:C.D4.13\endcsname{\hyperlink{tag-00E5}{II.\ref*{cdp:eq:C.D4.13}}}
\expandafter\def\csname SRef@cdp:eq:C.D4.14\endcsname{\hyperlink{tag-00E6}{II.\ref*{cdp:eq:C.D4.14}}}
\expandafter\def\csname SRef@cdp:eq:C.D4.15\endcsname{\hyperlink{tag-00E7}{II.\ref*{cdp:eq:C.D4.15}}}
\expandafter\def\csname SRef@cdp:eq:C.SC.1\endcsname{\hyperlink{tag-010L}{II.\ref*{cdp:eq:C.SC.1}}}
\expandafter\def\csname SRef@cdp:eq:C.D4.16\endcsname{\hyperlink{tag-00E8}{II.\ref*{cdp:eq:C.D4.16}}}
\expandafter\def\csname SRef@cdp:eq:C.D4.17\endcsname{\hyperlink{tag-00E9}{II.\ref*{cdp:eq:C.D4.17}}}
\expandafter\def\csname SRef@cdp:eq:C.D4.18\endcsname{\hyperlink{tag-00EA}{II.\ref*{cdp:eq:C.D4.18}}}
\expandafter\def\csname SRef@cdp:eq:C.D4.19\endcsname{\hyperlink{tag-00EB}{II.\ref*{cdp:eq:C.D4.19}}}
\expandafter\def\csname SRef@cdp:eq:C.D4.20\endcsname{\hyperlink{tag-00ED}{II.\ref*{cdp:eq:C.D4.20}}}
\expandafter\def\csname SRef@cdp:eq:C.UI.1\endcsname{\hyperlink{tag-014M}{II.\ref*{cdp:eq:C.UI.1}}}
\expandafter\def\csname SRef@cdp:eq:C.UI.2\endcsname{\hyperlink{tag-014U}{II.\ref*{cdp:eq:C.UI.2}}}
\expandafter\def\csname SRef@cdp:eq:C.UI.3\endcsname{\hyperlink{tag-014V}{II.\ref*{cdp:eq:C.UI.3}}}
\expandafter\def\csname SRef@cdp:eq:C.UI.4\endcsname{\hyperlink{tag-014W}{II.\ref*{cdp:eq:C.UI.4}}}
\expandafter\def\csname SRef@cdp:eq:C.SC.2\endcsname{\hyperlink{tag-010M}{II.\ref*{cdp:eq:C.SC.2}}}
\expandafter\def\csname SRef@cdp:eq:C.UI.5\endcsname{\hyperlink{tag-014X}{II.\ref*{cdp:eq:C.UI.5}}}
\expandafter\def\csname SRef@cdp:eq:C.UI.6\endcsname{\hyperlink{tag-014Y}{II.\ref*{cdp:eq:C.UI.6}}}
\expandafter\def\csname SRef@cdp:eq:C.UI.7\endcsname{\hyperlink{tag-014Z}{II.\ref*{cdp:eq:C.UI.7}}}
\expandafter\def\csname SRef@cdp:eq:C.UI.8\endcsname{\hyperlink{tag-0150}{II.\ref*{cdp:eq:C.UI.8}}}
\expandafter\def\csname SRef@cdp:eq:C.UI.9\endcsname{\hyperlink{tag-0151}{II.\ref*{cdp:eq:C.UI.9}}}
\expandafter\def\csname SRef@cdp:eq:C.UI.10\endcsname{\hyperlink{tag-014N}{II.\ref*{cdp:eq:C.UI.10}}}
\expandafter\def\csname SRef@cdp:eq:C.UI.11\endcsname{\hyperlink{tag-014O}{II.\ref*{cdp:eq:C.UI.11}}}
\expandafter\def\csname SRef@cdp:eq:C.UI.12\endcsname{\hyperlink{tag-014P}{II.\ref*{cdp:eq:C.UI.12}}}
\expandafter\def\csname SRef@cdp:eq:C.UI.13\endcsname{\hyperlink{tag-014Q}{II.\ref*{cdp:eq:C.UI.13}}}
\expandafter\def\csname SRef@cdp:eq:C.UI.14\endcsname{\hyperlink{tag-014R}{II.\ref*{cdp:eq:C.UI.14}}}
\expandafter\def\csname SRef@cdp:eq:C.SC.3\endcsname{\hyperlink{tag-010N}{II.\ref*{cdp:eq:C.SC.3}}}
\expandafter\def\csname SRef@cdp:eq:C.UI.15\endcsname{\hyperlink{tag-014S}{II.\ref*{cdp:eq:C.UI.15}}}
\expandafter\def\csname SRef@cdp:eq:C.UI.16\endcsname{\hyperlink{tag-014T}{II.\ref*{cdp:eq:C.UI.16}}}
\expandafter\def\csname SRef@cdp:eq:C.IZ.1\endcsname{\hyperlink{tag-00ON}{II.\ref*{cdp:eq:C.IZ.1}}}
\expandafter\def\csname SRef@cdp:eq:C.IZ.2\endcsname{\hyperlink{tag-00OO}{II.\ref*{cdp:eq:C.IZ.2}}}
\expandafter\def\csname SRef@cdp:eq:C.IZ.3\endcsname{\hyperlink{tag-00OP}{II.\ref*{cdp:eq:C.IZ.3}}}
\expandafter\def\csname SRef@cdp:eq:C.IZ.4\endcsname{\hyperlink{tag-00OQ}{II.\ref*{cdp:eq:C.IZ.4}}}
\expandafter\def\csname SRef@cdp:eq:C.IZ.5\endcsname{\hyperlink{tag-00OR}{II.\ref*{cdp:eq:C.IZ.5}}}
\expandafter\def\csname SRef@cdp:eq:C.IZ.6\endcsname{\hyperlink{tag-00OS}{II.\ref*{cdp:eq:C.IZ.6}}}
\expandafter\def\csname SRef@cdp:eq:C.IZ.7\endcsname{\hyperlink{tag-00OT}{II.\ref*{cdp:eq:C.IZ.7}}}
\expandafter\def\csname SRef@cdp:eq:C.IZ.8\endcsname{\hyperlink{tag-00OU}{II.\ref*{cdp:eq:C.IZ.8}}}
\expandafter\def\csname SRef@cdp:eq:C.SC.4\endcsname{\hyperlink{tag-010O}{II.\ref*{cdp:eq:C.SC.4}}}
\expandafter\def\csname SRef@cdp:eq:C.INT.1\endcsname{\hyperlink{tag-00NQ}{II.\ref*{cdp:eq:C.INT.1}}}
\expandafter\def\csname SRef@cdp:eq:C.INT.2\endcsname{\hyperlink{tag-00O1}{II.\ref*{cdp:eq:C.INT.2}}}
\expandafter\def\csname SRef@cdp:eq:C.INT.3\endcsname{\hyperlink{tag-00OA}{II.\ref*{cdp:eq:C.INT.3}}}
\expandafter\def\csname SRef@cdp:eq:C.INT.4\endcsname{\hyperlink{tag-00OB}{II.\ref*{cdp:eq:C.INT.4}}}
\expandafter\def\csname SRef@cdp:eq:C.INT.5\endcsname{\hyperlink{tag-00OC}{II.\ref*{cdp:eq:C.INT.5}}}
\expandafter\def\csname SRef@cdp:eq:C.INT.6\endcsname{\hyperlink{tag-00OD}{II.\ref*{cdp:eq:C.INT.6}}}
\expandafter\def\csname SRef@cdp:eq:C.INT.7\endcsname{\hyperlink{tag-00OE}{II.\ref*{cdp:eq:C.INT.7}}}
\expandafter\def\csname SRef@cdp:eq:C.INT.8\endcsname{\hyperlink{tag-00OF}{II.\ref*{cdp:eq:C.INT.8}}}
\expandafter\def\csname SRef@cdp:eq:C.INT.9\endcsname{\hyperlink{tag-00OG}{II.\ref*{cdp:eq:C.INT.9}}}
\expandafter\def\csname SRef@cdp:eq:C.INT.10\endcsname{\hyperlink{tag-00NR}{II.\ref*{cdp:eq:C.INT.10}}}
\expandafter\def\csname SRef@cdp:eq:C.SC.5\endcsname{\hyperlink{tag-010P}{II.\ref*{cdp:eq:C.SC.5}}}
\expandafter\def\csname SRef@cdp:eq:C.INT.11\endcsname{\hyperlink{tag-00NS}{II.\ref*{cdp:eq:C.INT.11}}}
\expandafter\def\csname SRef@cdp:eq:C.INT.12\endcsname{\hyperlink{tag-00NT}{II.\ref*{cdp:eq:C.INT.12}}}
\expandafter\def\csname SRef@cdp:eq:C.INT.13\endcsname{\hyperlink{tag-00NU}{II.\ref*{cdp:eq:C.INT.13}}}
\expandafter\def\csname SRef@cdp:eq:C.INT.14\endcsname{\hyperlink{tag-00NV}{II.\ref*{cdp:eq:C.INT.14}}}
\expandafter\def\csname SRef@cdp:eq:C.INT.15\endcsname{\hyperlink{tag-00NW}{II.\ref*{cdp:eq:C.INT.15}}}
\expandafter\def\csname SRef@cdp:eq:C.INT.16\endcsname{\hyperlink{tag-00NX}{II.\ref*{cdp:eq:C.INT.16}}}
\expandafter\def\csname SRef@cdp:eq:C.INT.17\endcsname{\hyperlink{tag-00NY}{II.\ref*{cdp:eq:C.INT.17}}}
\expandafter\def\csname SRef@cdp:eq:C.INT.18\endcsname{\hyperlink{tag-00NZ}{II.\ref*{cdp:eq:C.INT.18}}}
\expandafter\def\csname SRef@cdp:eq:C.INT.19\endcsname{\hyperlink{tag-00O0}{II.\ref*{cdp:eq:C.INT.19}}}
\expandafter\def\csname SRef@cdp:eq:C.INT.20\endcsname{\hyperlink{tag-00O2}{II.\ref*{cdp:eq:C.INT.20}}}
\expandafter\def\csname SRef@cdp:eq:C.SC.6\endcsname{\hyperlink{tag-010Q}{II.\ref*{cdp:eq:C.SC.6}}}
\expandafter\def\csname SRef@cdp:eq:C.INT.21\endcsname{\hyperlink{tag-00O3}{II.\ref*{cdp:eq:C.INT.21}}}
\expandafter\def\csname SRef@cdp:eq:C.INT.22\endcsname{\hyperlink{tag-00O4}{II.\ref*{cdp:eq:C.INT.22}}}
\expandafter\def\csname SRef@cdp:eq:C.INT.23\endcsname{\hyperlink{tag-00O5}{II.\ref*{cdp:eq:C.INT.23}}}
\expandafter\def\csname SRef@cdp:eq:C.INT.24\endcsname{\hyperlink{tag-00O6}{II.\ref*{cdp:eq:C.INT.24}}}
\expandafter\def\csname SRef@cdp:eq:C.INT.25\endcsname{\hyperlink{tag-00O7}{II.\ref*{cdp:eq:C.INT.25}}}
\expandafter\def\csname SRef@cdp:eq:C.INT.26\endcsname{\hyperlink{tag-00O8}{II.\ref*{cdp:eq:C.INT.26}}}
\expandafter\def\csname SRef@cdp:eq:C.INT.27\endcsname{\hyperlink{tag-00O9}{II.\ref*{cdp:eq:C.INT.27}}}
\expandafter\def\csname SRef@cdp:eq:C.MH.1\endcsname{\hyperlink{tag-00QA}{II.\ref*{cdp:eq:C.MH.1}}}
\expandafter\def\csname SRef@cdp:eq:C.MH.2\endcsname{\hyperlink{tag-00QJ}{II.\ref*{cdp:eq:C.MH.2}}}
\expandafter\def\csname SRef@cdp:eq:C.MH.3\endcsname{\hyperlink{tag-00QK}{II.\ref*{cdp:eq:C.MH.3}}}
\expandafter\def\csname SRef@cdp:eq:C.CF.1\endcsname{\hyperlink{tag-00BN}{II.\ref*{cdp:eq:C.CF.1}}}
\expandafter\def\csname SRef@cdp:eq:C.MH.4\endcsname{\hyperlink{tag-00QL}{II.\ref*{cdp:eq:C.MH.4}}}
\expandafter\def\csname SRef@cdp:eq:C.MH.5\endcsname{\hyperlink{tag-00QM}{II.\ref*{cdp:eq:C.MH.5}}}
\expandafter\def\csname SRef@cdp:eq:C.MH.6\endcsname{\hyperlink{tag-00QN}{II.\ref*{cdp:eq:C.MH.6}}}
\expandafter\def\csname SRef@cdp:eq:C.MH.7\endcsname{\hyperlink{tag-00QO}{II.\ref*{cdp:eq:C.MH.7}}}
\expandafter\def\csname SRef@cdp:eq:C.MH.8\endcsname{\hyperlink{tag-00QP}{II.\ref*{cdp:eq:C.MH.8}}}
\expandafter\def\csname SRef@cdp:eq:C.MH.9\endcsname{\hyperlink{tag-00QQ}{II.\ref*{cdp:eq:C.MH.9}}}
\expandafter\def\csname SRef@cdp:eq:C.MH.10\endcsname{\hyperlink{tag-00QB}{II.\ref*{cdp:eq:C.MH.10}}}
\expandafter\def\csname SRef@cdp:eq:C.MH.11\endcsname{\hyperlink{tag-00QC}{II.\ref*{cdp:eq:C.MH.11}}}
\expandafter\def\csname SRef@cdp:eq:C.MH.12\endcsname{\hyperlink{tag-00QD}{II.\ref*{cdp:eq:C.MH.12}}}
\expandafter\def\csname SRef@cdp:eq:C.MH.13\endcsname{\hyperlink{tag-00QE}{II.\ref*{cdp:eq:C.MH.13}}}
\expandafter\def\csname SRef@cdp:eq:C.CF.2\endcsname{\hyperlink{tag-00BO}{II.\ref*{cdp:eq:C.CF.2}}}
\expandafter\def\csname SRef@cdp:eq:C.MH.14\endcsname{\hyperlink{tag-00QF}{II.\ref*{cdp:eq:C.MH.14}}}
\expandafter\def\csname SRef@cdp:eq:C.MH.15\endcsname{\hyperlink{tag-00QG}{II.\ref*{cdp:eq:C.MH.15}}}
\expandafter\def\csname SRef@cdp:eq:C.MH.16\endcsname{\hyperlink{tag-00QH}{II.\ref*{cdp:eq:C.MH.16}}}
\expandafter\def\csname SRef@cdp:eq:C.MH.17\endcsname{\hyperlink{tag-00QI}{II.\ref*{cdp:eq:C.MH.17}}}
\expandafter\def\csname SRef@cdp:eq:C.1\endcsname{\hyperlink{tag-008N}{II.\ref*{cdp:eq:C.1}}}
\expandafter\def\csname SRef@cdp:eq:C.2\endcsname{\hyperlink{tag-008Y}{II.\ref*{cdp:eq:C.2}}}
\expandafter\def\csname SRef@cdp:eq:C.MS.1\endcsname{\hyperlink{tag-00RH}{II.\ref*{cdp:eq:C.MS.1}}}
\expandafter\def\csname SRef@cdp:eq:C.MS.2\endcsname{\hyperlink{tag-00RJ}{II.\ref*{cdp:eq:C.MS.2}}}
\expandafter\def\csname SRef@cdp:eq:C.MS.3\endcsname{\hyperlink{tag-00RK}{II.\ref*{cdp:eq:C.MS.3}}}
\expandafter\def\csname SRef@cdp:eq:C.MS.10\endcsname{\hyperlink{tag-00RI}{II.\ref*{cdp:eq:C.MS.10}}}
\expandafter\def\csname SRef@cdp:eq:C.CF.3\endcsname{\hyperlink{tag-00BP}{II.\ref*{cdp:eq:C.CF.3}}}
\expandafter\def\csname SRef@cdp:eq:C.MS.4\endcsname{\hyperlink{tag-00RL}{II.\ref*{cdp:eq:C.MS.4}}}
\expandafter\def\csname SRef@cdp:eq:C.MS.5\endcsname{\hyperlink{tag-00RM}{II.\ref*{cdp:eq:C.MS.5}}}
\expandafter\def\csname SRef@cdp:eq:C.MS.6\endcsname{\hyperlink{tag-00RN}{II.\ref*{cdp:eq:C.MS.6}}}
\expandafter\def\csname SRef@cdp:eq:C.MS.7\endcsname{\hyperlink{tag-00RO}{II.\ref*{cdp:eq:C.MS.7}}}
\expandafter\def\csname SRef@cdp:eq:C.MS.8\endcsname{\hyperlink{tag-00RP}{II.\ref*{cdp:eq:C.MS.8}}}
\expandafter\def\csname SRef@cdp:eq:C.MS.9\endcsname{\hyperlink{tag-00RQ}{II.\ref*{cdp:eq:C.MS.9}}}
\expandafter\def\csname SRef@cdp:eq:C.K3.1\endcsname{\hyperlink{tag-00OY}{II.\ref*{cdp:eq:C.K3.1}}}
\expandafter\def\csname SRef@cdp:eq:C.K3.2\endcsname{\hyperlink{tag-00P0}{II.\ref*{cdp:eq:C.K3.2}}}
\expandafter\def\csname SRef@cdp:eq:C.K3.3\endcsname{\hyperlink{tag-00P1}{II.\ref*{cdp:eq:C.K3.3}}}
\expandafter\def\csname SRef@cdp:eq:C.K3.4\endcsname{\hyperlink{tag-00P2}{II.\ref*{cdp:eq:C.K3.4}}}
\expandafter\def\csname SRef@cdp:eq:C.CF.4\endcsname{\hyperlink{tag-00BQ}{II.\ref*{cdp:eq:C.CF.4}}}
\expandafter\def\csname SRef@cdp:eq:C.K3.5\endcsname{\hyperlink{tag-00P3}{II.\ref*{cdp:eq:C.K3.5}}}
\expandafter\def\csname SRef@cdp:eq:C.K3.6\endcsname{\hyperlink{tag-00P4}{II.\ref*{cdp:eq:C.K3.6}}}
\expandafter\def\csname SRef@cdp:eq:C.K3.7\endcsname{\hyperlink{tag-00P5}{II.\ref*{cdp:eq:C.K3.7}}}
\expandafter\def\csname SRef@cdp:eq:C.K3.8\endcsname{\hyperlink{tag-00P6}{II.\ref*{cdp:eq:C.K3.8}}}
\expandafter\def\csname SRef@cdp:eq:C.K3.9\endcsname{\hyperlink{tag-00P7}{II.\ref*{cdp:eq:C.K3.9}}}
\expandafter\def\csname SRef@cdp:eq:C.K3.10\endcsname{\hyperlink{tag-00OZ}{II.\ref*{cdp:eq:C.K3.10}}}
\expandafter\def\csname SRef@cdp:eq:C.KX.1\endcsname{\hyperlink{tag-00PW}{II.\ref*{cdp:eq:C.KX.1}}}
\expandafter\def\csname SRef@cdp:eq:C.KX.2\endcsname{\hyperlink{tag-00Q2}{II.\ref*{cdp:eq:C.KX.2}}}
\expandafter\def\csname SRef@cdp:eq:C.KX.3\endcsname{\hyperlink{tag-00Q3}{II.\ref*{cdp:eq:C.KX.3}}}
\expandafter\def\csname SRef@cdp:eq:C.KX.4\endcsname{\hyperlink{tag-00Q4}{II.\ref*{cdp:eq:C.KX.4}}}
\expandafter\def\csname SRef@cdp:eq:C.CF.5\endcsname{\hyperlink{tag-00BR}{II.\ref*{cdp:eq:C.CF.5}}}
\expandafter\def\csname SRef@cdp:eq:C.KX.5\endcsname{\hyperlink{tag-00Q5}{II.\ref*{cdp:eq:C.KX.5}}}
\expandafter\def\csname SRef@cdp:eq:C.KX.6\endcsname{\hyperlink{tag-00Q6}{II.\ref*{cdp:eq:C.KX.6}}}
\expandafter\def\csname SRef@cdp:eq:C.KX.7\endcsname{\hyperlink{tag-00Q7}{II.\ref*{cdp:eq:C.KX.7}}}
\expandafter\def\csname SRef@cdp:eq:C.KX.8\endcsname{\hyperlink{tag-00Q8}{II.\ref*{cdp:eq:C.KX.8}}}
\expandafter\def\csname SRef@cdp:eq:C.KX.9\endcsname{\hyperlink{tag-00Q9}{II.\ref*{cdp:eq:C.KX.9}}}
\expandafter\def\csname SRef@cdp:eq:C.KX.10\endcsname{\hyperlink{tag-00PX}{II.\ref*{cdp:eq:C.KX.10}}}
\expandafter\def\csname SRef@cdp:eq:C.KX.11\endcsname{\hyperlink{tag-00PY}{II.\ref*{cdp:eq:C.KX.11}}}
\expandafter\def\csname SRef@cdp:eq:C.KX.12\endcsname{\hyperlink{tag-00PZ}{II.\ref*{cdp:eq:C.KX.12}}}
\expandafter\def\csname SRef@cdp:eq:C.KX.13\endcsname{\hyperlink{tag-00Q0}{II.\ref*{cdp:eq:C.KX.13}}}
\expandafter\def\csname SRef@cdp:eq:C.KX.14\endcsname{\hyperlink{tag-00Q1}{II.\ref*{cdp:eq:C.KX.14}}}
\expandafter\def\csname SRef@cdp:eq:C.CF.6\endcsname{\hyperlink{tag-00BS}{II.\ref*{cdp:eq:C.CF.6}}}
\expandafter\def\csname SRef@cdp:eq:C.EL.1\endcsname{\hyperlink{tag-00FY}{II.\ref*{cdp:eq:C.EL.1}}}
\expandafter\def\csname SRef@cdp:eq:C.EL.2\endcsname{\hyperlink{tag-00G1}{II.\ref*{cdp:eq:C.EL.2}}}
\expandafter\def\csname SRef@cdp:eq:C.EL.3\endcsname{\hyperlink{tag-00G2}{II.\ref*{cdp:eq:C.EL.3}}}
\expandafter\def\csname SRef@cdp:eq:C.EL.4\endcsname{\hyperlink{tag-00G3}{II.\ref*{cdp:eq:C.EL.4}}}
\expandafter\def\csname SRef@cdp:eq:C.EL.5\endcsname{\hyperlink{tag-00G4}{II.\ref*{cdp:eq:C.EL.5}}}
\expandafter\def\csname SRef@cdp:eq:C.EL.6\endcsname{\hyperlink{tag-00G5}{II.\ref*{cdp:eq:C.EL.6}}}
\expandafter\def\csname SRef@cdp:eq:C.EL.7\endcsname{\hyperlink{tag-00G6}{II.\ref*{cdp:eq:C.EL.7}}}
\expandafter\def\csname SRef@cdp:eq:C.EL.8\endcsname{\hyperlink{tag-00G7}{II.\ref*{cdp:eq:C.EL.8}}}
\expandafter\def\csname SRef@cdp:eq:C.EL.9\endcsname{\hyperlink{tag-00G8}{II.\ref*{cdp:eq:C.EL.9}}}
\expandafter\def\csname SRef@cdp:eq:C.EL.10\endcsname{\hyperlink{tag-00FZ}{II.\ref*{cdp:eq:C.EL.10}}}
\expandafter\def\csname SRef@cdp:eq:C.CF.7\endcsname{\hyperlink{tag-00BT}{II.\ref*{cdp:eq:C.CF.7}}}
\expandafter\def\csname SRef@cdp:eq:C.EL.11\endcsname{\hyperlink{tag-00G0}{II.\ref*{cdp:eq:C.EL.11}}}
\expandafter\def\csname SRef@cdp:eq:C.CZ.1\endcsname{\hyperlink{tag-00DN}{II.\ref*{cdp:eq:C.CZ.1}}}
\expandafter\def\csname SRef@cdp:eq:C.CZ.2\endcsname{\hyperlink{tag-00DT}{II.\ref*{cdp:eq:C.CZ.2}}}
\expandafter\def\csname SRef@cdp:eq:C.CZ.3\endcsname{\hyperlink{tag-00DU}{II.\ref*{cdp:eq:C.CZ.3}}}
\expandafter\def\csname SRef@cdp:eq:C.CZ.4\endcsname{\hyperlink{tag-00DV}{II.\ref*{cdp:eq:C.CZ.4}}}
\expandafter\def\csname SRef@cdp:eq:C.CZ.5\endcsname{\hyperlink{tag-00DW}{II.\ref*{cdp:eq:C.CZ.5}}}
\expandafter\def\csname SRef@cdp:eq:C.CZ.6\endcsname{\hyperlink{tag-00DX}{II.\ref*{cdp:eq:C.CZ.6}}}
\expandafter\def\csname SRef@cdp:eq:C.CZ.7\endcsname{\hyperlink{tag-00DY}{II.\ref*{cdp:eq:C.CZ.7}}}
\expandafter\def\csname SRef@cdp:eq:C.CZ.8\endcsname{\hyperlink{tag-00DZ}{II.\ref*{cdp:eq:C.CZ.8}}}
\expandafter\def\csname SRef@cdp:eq:C.CZ.9\endcsname{\hyperlink{tag-00E0}{II.\ref*{cdp:eq:C.CZ.9}}}
\expandafter\def\csname SRef@cdp:eq:C.CF.8\endcsname{\hyperlink{tag-00BU}{II.\ref*{cdp:eq:C.CF.8}}}
\expandafter\def\csname SRef@cdp:eq:C.CZ.10\endcsname{\hyperlink{tag-00DO}{II.\ref*{cdp:eq:C.CZ.10}}}
\expandafter\def\csname SRef@cdp:eq:C.CZ.11\endcsname{\hyperlink{tag-00DP}{II.\ref*{cdp:eq:C.CZ.11}}}
\expandafter\def\csname SRef@cdp:eq:C.CZ.12\endcsname{\hyperlink{tag-00DQ}{II.\ref*{cdp:eq:C.CZ.12}}}
\expandafter\def\csname SRef@cdp:eq:C.CZ.13\endcsname{\hyperlink{tag-00DR}{II.\ref*{cdp:eq:C.CZ.13}}}
\expandafter\def\csname SRef@cdp:eq:C.CZ.14\endcsname{\hyperlink{tag-00DS}{II.\ref*{cdp:eq:C.CZ.14}}}
\expandafter\def\csname SRef@cdp:eq:C.MR.1\endcsname{\hyperlink{tag-00R1}{II.\ref*{cdp:eq:C.MR.1}}}
\expandafter\def\csname SRef@cdp:eq:C.MR.2\endcsname{\hyperlink{tag-00R9}{II.\ref*{cdp:eq:C.MR.2}}}
\expandafter\def\csname SRef@cdp:eq:C.MR.3\endcsname{\hyperlink{tag-00RA}{II.\ref*{cdp:eq:C.MR.3}}}
\expandafter\def\csname SRef@cdp:eq:C.MR.4\endcsname{\hyperlink{tag-00RB}{II.\ref*{cdp:eq:C.MR.4}}}
\expandafter\def\csname SRef@cdp:eq:C.MR.5\endcsname{\hyperlink{tag-00RC}{II.\ref*{cdp:eq:C.MR.5}}}
\expandafter\def\csname SRef@cdp:eq:C.CF.9\endcsname{\hyperlink{tag-00BV}{II.\ref*{cdp:eq:C.CF.9}}}
\expandafter\def\csname SRef@cdp:eq:C.MR.6\endcsname{\hyperlink{tag-00RD}{II.\ref*{cdp:eq:C.MR.6}}}
\expandafter\def\csname SRef@cdp:eq:C.MR.7\endcsname{\hyperlink{tag-00RE}{II.\ref*{cdp:eq:C.MR.7}}}
\expandafter\def\csname SRef@cdp:eq:C.MR.8\endcsname{\hyperlink{tag-00RF}{II.\ref*{cdp:eq:C.MR.8}}}
\expandafter\def\csname SRef@cdp:eq:C.MR.9\endcsname{\hyperlink{tag-00RG}{II.\ref*{cdp:eq:C.MR.9}}}
\expandafter\def\csname SRef@cdp:eq:C.MR.10\endcsname{\hyperlink{tag-00R2}{II.\ref*{cdp:eq:C.MR.10}}}
\expandafter\def\csname SRef@cdp:eq:C.MR.11\endcsname{\hyperlink{tag-00R3}{II.\ref*{cdp:eq:C.MR.11}}}
\expandafter\def\csname SRef@cdp:eq:C.MR.12\endcsname{\hyperlink{tag-00R4}{II.\ref*{cdp:eq:C.MR.12}}}
\expandafter\def\csname SRef@cdp:eq:C.MR.13\endcsname{\hyperlink{tag-00R5}{II.\ref*{cdp:eq:C.MR.13}}}
\expandafter\def\csname SRef@cdp:eq:C.MR.14\endcsname{\hyperlink{tag-00R6}{II.\ref*{cdp:eq:C.MR.14}}}
\expandafter\def\csname SRef@cdp:eq:C.MR.15\endcsname{\hyperlink{tag-00R7}{II.\ref*{cdp:eq:C.MR.15}}}
\expandafter\def\csname SRef@cdp:eq:C.CA.1\endcsname{\hyperlink{tag-00BC}{II.\ref*{cdp:eq:C.CA.1}}}
\expandafter\def\csname SRef@cdp:eq:C.MR.16\endcsname{\hyperlink{tag-00R8}{II.\ref*{cdp:eq:C.MR.16}}}
\expandafter\def\csname SRef@cdp:eq:C.RH.1\endcsname{\hyperlink{tag-00ZG}{II.\ref*{cdp:eq:C.RH.1}}}
\expandafter\def\csname SRef@cdp:eq:C.RH.2\endcsname{\hyperlink{tag-00ZI}{II.\ref*{cdp:eq:C.RH.2}}}
\expandafter\def\csname SRef@cdp:eq:C.RH.3\endcsname{\hyperlink{tag-00ZJ}{II.\ref*{cdp:eq:C.RH.3}}}
\expandafter\def\csname SRef@cdp:eq:C.RH.4\endcsname{\hyperlink{tag-00ZK}{II.\ref*{cdp:eq:C.RH.4}}}
\expandafter\def\csname SRef@cdp:eq:C.RH.5\endcsname{\hyperlink{tag-00ZL}{II.\ref*{cdp:eq:C.RH.5}}}
\expandafter\def\csname SRef@cdp:eq:C.RH.6\endcsname{\hyperlink{tag-00ZM}{II.\ref*{cdp:eq:C.RH.6}}}
\expandafter\def\csname SRef@cdp:eq:C.RH.7\endcsname{\hyperlink{tag-00ZN}{II.\ref*{cdp:eq:C.RH.7}}}
\expandafter\def\csname SRef@cdp:eq:C.RH.8\endcsname{\hyperlink{tag-00ZO}{II.\ref*{cdp:eq:C.RH.8}}}
\expandafter\def\csname SRef@cdp:eq:C.RH.9\endcsname{\hyperlink{tag-00ZP}{II.\ref*{cdp:eq:C.RH.9}}}
\expandafter\def\csname SRef@cdp:eq:C.CA.2\endcsname{\hyperlink{tag-00BF}{II.\ref*{cdp:eq:C.CA.2}}}
\expandafter\def\csname SRef@cdp:eq:C.RH.10\endcsname{\hyperlink{tag-00ZH}{II.\ref*{cdp:eq:C.RH.10}}}
\expandafter\def\csname SRef@cdp:eq:C.TC.1\endcsname{\hyperlink{tag-012C}{II.\ref*{cdp:eq:C.TC.1}}}
\expandafter\def\csname SRef@cdp:eq:C.TC.2\endcsname{\hyperlink{tag-012G}{II.\ref*{cdp:eq:C.TC.2}}}
\expandafter\def\csname SRef@cdp:eq:C.TC.3\endcsname{\hyperlink{tag-012H}{II.\ref*{cdp:eq:C.TC.3}}}
\expandafter\def\csname SRef@cdp:eq:C.TC.4\endcsname{\hyperlink{tag-012I}{II.\ref*{cdp:eq:C.TC.4}}}
\expandafter\def\csname SRef@cdp:eq:C.TC.5\endcsname{\hyperlink{tag-012J}{II.\ref*{cdp:eq:C.TC.5}}}
\expandafter\def\csname SRef@cdp:eq:C.TC.6\endcsname{\hyperlink{tag-012K}{II.\ref*{cdp:eq:C.TC.6}}}
\expandafter\def\csname SRef@cdp:eq:C.TC.7\endcsname{\hyperlink{tag-012L}{II.\ref*{cdp:eq:C.TC.7}}}
\expandafter\def\csname SRef@cdp:eq:C.TC.8\endcsname{\hyperlink{tag-012M}{II.\ref*{cdp:eq:C.TC.8}}}
\expandafter\def\csname SRef@cdp:eq:C.TC.9\endcsname{\hyperlink{tag-012N}{II.\ref*{cdp:eq:C.TC.9}}}
\expandafter\def\csname SRef@cdp:eq:C.CA.3\endcsname{\hyperlink{tag-00BG}{II.\ref*{cdp:eq:C.CA.3}}}
\expandafter\def\csname SRef@cdp:eq:C.TC.10\endcsname{\hyperlink{tag-012D}{II.\ref*{cdp:eq:C.TC.10}}}
\expandafter\def\csname SRef@cdp:eq:C.TC.11\endcsname{\hyperlink{tag-012E}{II.\ref*{cdp:eq:C.TC.11}}}
\expandafter\def\csname SRef@cdp:eq:C.TC.12\endcsname{\hyperlink{tag-012F}{II.\ref*{cdp:eq:C.TC.12}}}
\expandafter\def\csname SRef@cdp:eq:C.ER.1\endcsname{\hyperlink{tag-00IN}{II.\ref*{cdp:eq:C.ER.1}}}
\expandafter\def\csname SRef@cdp:eq:C.ER.2\endcsname{\hyperlink{tag-00IX}{II.\ref*{cdp:eq:C.ER.2}}}
\expandafter\def\csname SRef@cdp:eq:C.ER.3\endcsname{\hyperlink{tag-00IY}{II.\ref*{cdp:eq:C.ER.3}}}
\expandafter\def\csname SRef@cdp:eq:C.ER.4\endcsname{\hyperlink{tag-00IZ}{II.\ref*{cdp:eq:C.ER.4}}}
\expandafter\def\csname SRef@cdp:eq:C.ER.5\endcsname{\hyperlink{tag-00J0}{II.\ref*{cdp:eq:C.ER.5}}}
\expandafter\def\csname SRef@cdp:eq:C.ER.6\endcsname{\hyperlink{tag-00J1}{II.\ref*{cdp:eq:C.ER.6}}}
\expandafter\def\csname SRef@cdp:eq:C.ER.7\endcsname{\hyperlink{tag-00J2}{II.\ref*{cdp:eq:C.ER.7}}}
\expandafter\def\csname SRef@cdp:eq:C.CA.4\endcsname{\hyperlink{tag-00BH}{II.\ref*{cdp:eq:C.CA.4}}}
\expandafter\def\csname SRef@cdp:eq:C.ER.8\endcsname{\hyperlink{tag-00J3}{II.\ref*{cdp:eq:C.ER.8}}}
\expandafter\def\csname SRef@cdp:eq:C.ER.9\endcsname{\hyperlink{tag-00J4}{II.\ref*{cdp:eq:C.ER.9}}}
\expandafter\def\csname SRef@cdp:eq:C.ER.10\endcsname{\hyperlink{tag-00IO}{II.\ref*{cdp:eq:C.ER.10}}}
\expandafter\def\csname SRef@cdp:eq:C.ER.11\endcsname{\hyperlink{tag-00IP}{II.\ref*{cdp:eq:C.ER.11}}}
\expandafter\def\csname SRef@cdp:eq:C.ER.12\endcsname{\hyperlink{tag-00IQ}{II.\ref*{cdp:eq:C.ER.12}}}
\expandafter\def\csname SRef@cdp:eq:C.ER.13\endcsname{\hyperlink{tag-00IR}{II.\ref*{cdp:eq:C.ER.13}}}
\expandafter\def\csname SRef@cdp:eq:C.ER.14\endcsname{\hyperlink{tag-00IS}{II.\ref*{cdp:eq:C.ER.14}}}
\expandafter\def\csname SRef@cdp:eq:C.ER.15\endcsname{\hyperlink{tag-00IT}{II.\ref*{cdp:eq:C.ER.15}}}
\expandafter\def\csname SRef@cdp:eq:C.ER.16\endcsname{\hyperlink{tag-00IU}{II.\ref*{cdp:eq:C.ER.16}}}
\expandafter\def\csname SRef@cdp:eq:C.CA.5\endcsname{\hyperlink{tag-00BI}{II.\ref*{cdp:eq:C.CA.5}}}
\expandafter\def\csname SRef@cdp:eq:C.ER.17\endcsname{\hyperlink{tag-00IV}{II.\ref*{cdp:eq:C.ER.17}}}
\expandafter\def\csname SRef@cdp:eq:C.ER.18\endcsname{\hyperlink{tag-00IW}{II.\ref*{cdp:eq:C.ER.18}}}
\expandafter\def\csname SRef@cdp:eq:C.CP.1\endcsname{\hyperlink{tag-00C9}{II.\ref*{cdp:eq:C.CP.1}}}
\expandafter\def\csname SRef@cdp:eq:C.CP.2\endcsname{\hyperlink{tag-00CB}{II.\ref*{cdp:eq:C.CP.2}}}
\expandafter\def\csname SRef@cdp:eq:C.CP.3\endcsname{\hyperlink{tag-00CC}{II.\ref*{cdp:eq:C.CP.3}}}
\expandafter\def\csname SRef@cdp:eq:C.CP.4\endcsname{\hyperlink{tag-00CD}{II.\ref*{cdp:eq:C.CP.4}}}
\expandafter\def\csname SRef@cdp:eq:C.CP.5\endcsname{\hyperlink{tag-00CE}{II.\ref*{cdp:eq:C.CP.5}}}
\expandafter\def\csname SRef@cdp:eq:C.CP.6\endcsname{\hyperlink{tag-00CF}{II.\ref*{cdp:eq:C.CP.6}}}
\expandafter\def\csname SRef@cdp:eq:C.CP.7\endcsname{\hyperlink{tag-00CG}{II.\ref*{cdp:eq:C.CP.7}}}
\expandafter\def\csname SRef@cdp:eq:C.CA.6\endcsname{\hyperlink{tag-00BJ}{II.\ref*{cdp:eq:C.CA.6}}}
\expandafter\def\csname SRef@cdp:eq:C.CP.8\endcsname{\hyperlink{tag-00CH}{II.\ref*{cdp:eq:C.CP.8}}}
\expandafter\def\csname SRef@cdp:eq:C.CP.9\endcsname{\hyperlink{tag-00CI}{II.\ref*{cdp:eq:C.CP.9}}}
\expandafter\def\csname SRef@cdp:eq:C.CP.10\endcsname{\hyperlink{tag-00CA}{II.\ref*{cdp:eq:C.CP.10}}}
\expandafter\def\csname SRef@cdp:eq:C.KM.1\endcsname{\hyperlink{tag-00PI}{II.\ref*{cdp:eq:C.KM.1}}}
\expandafter\def\csname SRef@cdp:eq:C.KM.2\endcsname{\hyperlink{tag-00PL}{II.\ref*{cdp:eq:C.KM.2}}}
\expandafter\def\csname SRef@cdp:eq:C.KM.S1\endcsname{\hyperlink{tag-00PT}{II.\ref*{cdp:eq:C.KM.S1}}}
\expandafter\def\csname SRef@cdp:eq:C.KM.3\endcsname{\hyperlink{tag-00PM}{II.\ref*{cdp:eq:C.KM.3}}}
\expandafter\def\csname SRef@cdp:eq:C.KM.4\endcsname{\hyperlink{tag-00PN}{II.\ref*{cdp:eq:C.KM.4}}}
\expandafter\def\csname SRef@cdp:eq:C.KM.5\endcsname{\hyperlink{tag-00PO}{II.\ref*{cdp:eq:C.KM.5}}}
\expandafter\def\csname SRef@cdp:eq:C.KM.6\endcsname{\hyperlink{tag-00PP}{II.\ref*{cdp:eq:C.KM.6}}}
\expandafter\def\csname SRef@cdp:eq:C.CA.7\endcsname{\hyperlink{tag-00BK}{II.\ref*{cdp:eq:C.CA.7}}}
\expandafter\def\csname SRef@cdp:eq:C.KM.7\endcsname{\hyperlink{tag-00PQ}{II.\ref*{cdp:eq:C.KM.7}}}
\expandafter\def\csname SRef@cdp:eq:C.KM.8\endcsname{\hyperlink{tag-00PR}{II.\ref*{cdp:eq:C.KM.8}}}
\expandafter\def\csname SRef@cdp:eq:C.KM.9\endcsname{\hyperlink{tag-00PS}{II.\ref*{cdp:eq:C.KM.9}}}
\expandafter\def\csname SRef@cdp:eq:C.KM.10\endcsname{\hyperlink{tag-00PJ}{II.\ref*{cdp:eq:C.KM.10}}}
\expandafter\def\csname SRef@cdp:eq:C.KM.S2\endcsname{\hyperlink{tag-00PU}{II.\ref*{cdp:eq:C.KM.S2}}}
\expandafter\def\csname SRef@cdp:eq:C.KM.S3\endcsname{\hyperlink{tag-00PV}{II.\ref*{cdp:eq:C.KM.S3}}}
\expandafter\def\csname SRef@cdp:eq:C.KM.11\endcsname{\hyperlink{tag-00PK}{II.\ref*{cdp:eq:C.KM.11}}}
\expandafter\def\csname SRef@cdp:eq:C.PE.1\endcsname{\hyperlink{tag-00U2}{II.\ref*{cdp:eq:C.PE.1}}}
\expandafter\def\csname SRef@cdp:eq:C.PE.2\endcsname{\hyperlink{tag-00UD}{II.\ref*{cdp:eq:C.PE.2}}}
\expandafter\def\csname SRef@cdp:eq:C.PE.3\endcsname{\hyperlink{tag-00UG}{II.\ref*{cdp:eq:C.PE.3}}}
\expandafter\def\csname SRef@cdp:eq:C.CA.8\endcsname{\hyperlink{tag-00BL}{II.\ref*{cdp:eq:C.CA.8}}}
\expandafter\def\csname SRef@cdp:eq:C.PE.4\endcsname{\hyperlink{tag-00UH}{II.\ref*{cdp:eq:C.PE.4}}}
\expandafter\def\csname SRef@cdp:eq:C.PE.5\endcsname{\hyperlink{tag-00UI}{II.\ref*{cdp:eq:C.PE.5}}}
\expandafter\def\csname SRef@cdp:eq:C.PE.6\endcsname{\hyperlink{tag-00UJ}{II.\ref*{cdp:eq:C.PE.6}}}
\expandafter\def\csname SRef@cdp:eq:C.PE.7\endcsname{\hyperlink{tag-00UK}{II.\ref*{cdp:eq:C.PE.7}}}
\expandafter\def\csname SRef@cdp:eq:C.PE.8\endcsname{\hyperlink{tag-00UL}{II.\ref*{cdp:eq:C.PE.8}}}
\expandafter\def\csname SRef@cdp:eq:C.PE.9\endcsname{\hyperlink{tag-00UM}{II.\ref*{cdp:eq:C.PE.9}}}
\expandafter\def\csname SRef@cdp:eq:C.PE.10\endcsname{\hyperlink{tag-00U3}{II.\ref*{cdp:eq:C.PE.10}}}
\expandafter\def\csname SRef@cdp:eq:C.PE.11\endcsname{\hyperlink{tag-00U4}{II.\ref*{cdp:eq:C.PE.11}}}
\expandafter\def\csname SRef@cdp:eq:C.PE.12\endcsname{\hyperlink{tag-00U5}{II.\ref*{cdp:eq:C.PE.12}}}
\expandafter\def\csname SRef@cdp:eq:C.PE.13\endcsname{\hyperlink{tag-00U6}{II.\ref*{cdp:eq:C.PE.13}}}
\expandafter\def\csname SRef@cdp:eq:C.CA.9\endcsname{\hyperlink{tag-00BM}{II.\ref*{cdp:eq:C.CA.9}}}
\expandafter\def\csname SRef@cdp:eq:C.PE.14\endcsname{\hyperlink{tag-00U7}{II.\ref*{cdp:eq:C.PE.14}}}
\expandafter\def\csname SRef@cdp:eq:C.PE.15\endcsname{\hyperlink{tag-00U8}{II.\ref*{cdp:eq:C.PE.15}}}
\expandafter\def\csname SRef@cdp:eq:C.PE.16\endcsname{\hyperlink{tag-00U9}{II.\ref*{cdp:eq:C.PE.16}}}
\expandafter\def\csname SRef@cdp:eq:C.PE.17\endcsname{\hyperlink{tag-00UA}{II.\ref*{cdp:eq:C.PE.17}}}
\expandafter\def\csname SRef@cdp:eq:C.PE.18\endcsname{\hyperlink{tag-00UB}{II.\ref*{cdp:eq:C.PE.18}}}
\expandafter\def\csname SRef@cdp:eq:C.PE.19\endcsname{\hyperlink{tag-00UC}{II.\ref*{cdp:eq:C.PE.19}}}
\expandafter\def\csname SRef@cdp:eq:C.PE.20\endcsname{\hyperlink{tag-00UE}{II.\ref*{cdp:eq:C.PE.20}}}
\expandafter\def\csname SRef@cdp:eq:C.PE.21\endcsname{\hyperlink{tag-00UF}{II.\ref*{cdp:eq:C.PE.21}}}
\expandafter\def\csname SRef@cdp:eq:C.SKD.1\endcsname{\hyperlink{tag-011N}{II.\ref*{cdp:eq:C.SKD.1}}}
\expandafter\def\csname SRef@cdp:eq:C.SKD.2\endcsname{\hyperlink{tag-011U}{II.\ref*{cdp:eq:C.SKD.2}}}
\expandafter\def\csname SRef@cdp:eq:C.CA.10\endcsname{\hyperlink{tag-00BD}{II.\ref*{cdp:eq:C.CA.10}}}
\expandafter\def\csname SRef@cdp:eq:C.SKD.3\endcsname{\hyperlink{tag-011V}{II.\ref*{cdp:eq:C.SKD.3}}}
\expandafter\def\csname SRef@cdp:eq:C.SKD.4\endcsname{\hyperlink{tag-011W}{II.\ref*{cdp:eq:C.SKD.4}}}
\expandafter\def\csname SRef@cdp:eq:C.SKD.5\endcsname{\hyperlink{tag-011X}{II.\ref*{cdp:eq:C.SKD.5}}}
\expandafter\def\csname SRef@cdp:eq:C.SKD.6\endcsname{\hyperlink{tag-011Y}{II.\ref*{cdp:eq:C.SKD.6}}}
\expandafter\def\csname SRef@cdp:eq:C.SKD.7\endcsname{\hyperlink{tag-011Z}{II.\ref*{cdp:eq:C.SKD.7}}}
\expandafter\def\csname SRef@cdp:eq:C.SKD.8\endcsname{\hyperlink{tag-0120}{II.\ref*{cdp:eq:C.SKD.8}}}
\expandafter\def\csname SRef@cdp:eq:C.SKD.9\endcsname{\hyperlink{tag-0121}{II.\ref*{cdp:eq:C.SKD.9}}}
\expandafter\def\csname SRef@cdp:eq:C.SKD.10\endcsname{\hyperlink{tag-011O}{II.\ref*{cdp:eq:C.SKD.10}}}
\expandafter\def\csname SRef@cdp:eq:C.SKD.11\endcsname{\hyperlink{tag-011P}{II.\ref*{cdp:eq:C.SKD.11}}}
\expandafter\def\csname SRef@cdp:eq:C.SKD.14\endcsname{\hyperlink{tag-011S}{II.\ref*{cdp:eq:C.SKD.14}}}
\expandafter\def\csname SRef@cdp:eq:C.CA.11\endcsname{\hyperlink{tag-00BE}{II.\ref*{cdp:eq:C.CA.11}}}
\expandafter\def\csname SRef@cdp:eq:C.SKD.15\endcsname{\hyperlink{tag-011T}{II.\ref*{cdp:eq:C.SKD.15}}}
\expandafter\def\csname SRef@cdp:eq:C.SKD.12\endcsname{\hyperlink{tag-011Q}{II.\ref*{cdp:eq:C.SKD.12}}}
\expandafter\def\csname SRef@cdp:eq:C.SKD.13\endcsname{\hyperlink{tag-011R}{II.\ref*{cdp:eq:C.SKD.13}}}
\expandafter\def\csname SRef@cdp:eq:C.EN.1\endcsname{\hyperlink{tag-00G9}{II.\ref*{cdp:eq:C.EN.1}}}
\expandafter\def\csname SRef@cdp:eq:C.EN.18\endcsname{\hyperlink{tag-00GI}{II.\ref*{cdp:eq:C.EN.18}}}
\expandafter\def\csname SRef@cdp:eq:C.EN.2\endcsname{\hyperlink{tag-00GJ}{II.\ref*{cdp:eq:C.EN.2}}}
\expandafter\def\csname SRef@cdp:eq:C.EN.3\endcsname{\hyperlink{tag-00GK}{II.\ref*{cdp:eq:C.EN.3}}}
\expandafter\def\csname SRef@cdp:eq:C.EN.4\endcsname{\hyperlink{tag-00GL}{II.\ref*{cdp:eq:C.EN.4}}}
\expandafter\def\csname SRef@cdp:eq:C.EN.5\endcsname{\hyperlink{tag-00GM}{II.\ref*{cdp:eq:C.EN.5}}}
\expandafter\def\csname SRef@cdp:eq:C.EN.6\endcsname{\hyperlink{tag-00GN}{II.\ref*{cdp:eq:C.EN.6}}}
\expandafter\def\csname SRef@cdp:eq:C.HQ.1\endcsname{\hyperlink{tag-00MQ}{II.\ref*{cdp:eq:C.HQ.1}}}
\expandafter\def\csname SRef@cdp:eq:C.EN.7\endcsname{\hyperlink{tag-00GO}{II.\ref*{cdp:eq:C.EN.7}}}
\expandafter\def\csname SRef@cdp:eq:C.EN.8\endcsname{\hyperlink{tag-00GP}{II.\ref*{cdp:eq:C.EN.8}}}
\expandafter\def\csname SRef@cdp:eq:C.EN.9\endcsname{\hyperlink{tag-00GQ}{II.\ref*{cdp:eq:C.EN.9}}}
\expandafter\def\csname SRef@cdp:eq:C.EN.10\endcsname{\hyperlink{tag-00GA}{II.\ref*{cdp:eq:C.EN.10}}}
\expandafter\def\csname SRef@cdp:eq:C.EN.11\endcsname{\hyperlink{tag-00GB}{II.\ref*{cdp:eq:C.EN.11}}}
\expandafter\def\csname SRef@cdp:eq:C.EN.12\endcsname{\hyperlink{tag-00GC}{II.\ref*{cdp:eq:C.EN.12}}}
\expandafter\def\csname SRef@cdp:eq:C.EN.13\endcsname{\hyperlink{tag-00GD}{II.\ref*{cdp:eq:C.EN.13}}}
\expandafter\def\csname SRef@cdp:eq:C.EN.14\endcsname{\hyperlink{tag-00GE}{II.\ref*{cdp:eq:C.EN.14}}}
\expandafter\def\csname SRef@cdp:eq:C.EN.15\endcsname{\hyperlink{tag-00GF}{II.\ref*{cdp:eq:C.EN.15}}}
\expandafter\def\csname SRef@cdp:eq:C.EN.16\endcsname{\hyperlink{tag-00GG}{II.\ref*{cdp:eq:C.EN.16}}}
\expandafter\def\csname SRef@cdp:eq:C.HQ.2\endcsname{\hyperlink{tag-00MS}{II.\ref*{cdp:eq:C.HQ.2}}}
\expandafter\def\csname SRef@cdp:eq:C.EN.17\endcsname{\hyperlink{tag-00GH}{II.\ref*{cdp:eq:C.EN.17}}}
\expandafter\def\csname SRef@cdp:eq:C.EN.S1\endcsname{\hyperlink{tag-00H7}{II.\ref*{cdp:eq:C.EN.S1}}}
\expandafter\def\csname SRef@cdp:eq:C.EN.S2\endcsname{\hyperlink{tag-00H8}{II.\ref*{cdp:eq:C.EN.S2}}}
\expandafter\def\csname SRef@cdp:eq:C.EN.S3\endcsname{\hyperlink{tag-00H9}{II.\ref*{cdp:eq:C.EN.S3}}}
\expandafter\def\csname SRef@cdp:eq:C.EN.S4\endcsname{\hyperlink{tag-00HA}{II.\ref*{cdp:eq:C.EN.S4}}}
\expandafter\def\csname SRef@cdp:eq:C.EN.S5\endcsname{\hyperlink{tag-00HB}{II.\ref*{cdp:eq:C.EN.S5}}}
\expandafter\def\csname SRef@cdp:eq:C.EN.S6\endcsname{\hyperlink{tag-00HC}{II.\ref*{cdp:eq:C.EN.S6}}}
\expandafter\def\csname SRef@cdp:eq:C.EN.S7\endcsname{\hyperlink{tag-00HD}{II.\ref*{cdp:eq:C.EN.S7}}}
\expandafter\def\csname SRef@cdp:eq:C.EN.R1\endcsname{\hyperlink{tag-00GV}{II.\ref*{cdp:eq:C.EN.R1}}}
\expandafter\def\csname SRef@cdp:eq:C.EN.R2\endcsname{\hyperlink{tag-00GZ}{II.\ref*{cdp:eq:C.EN.R2}}}
\expandafter\def\csname SRef@cdp:eq:C.HQ.3\endcsname{\hyperlink{tag-00MT}{II.\ref*{cdp:eq:C.HQ.3}}}
\expandafter\def\csname SRef@cdp:eq:C.EN.R3\endcsname{\hyperlink{tag-00H0}{II.\ref*{cdp:eq:C.EN.R3}}}
\expandafter\def\csname SRef@cdp:eq:C.EN.R4\endcsname{\hyperlink{tag-00H1}{II.\ref*{cdp:eq:C.EN.R4}}}
\expandafter\def\csname SRef@cdp:eq:C.EN.R5\endcsname{\hyperlink{tag-00H2}{II.\ref*{cdp:eq:C.EN.R5}}}
\expandafter\def\csname SRef@cdp:eq:C.EN.R6\endcsname{\hyperlink{tag-00H3}{II.\ref*{cdp:eq:C.EN.R6}}}
\expandafter\def\csname SRef@cdp:eq:C.EN.R7\endcsname{\hyperlink{tag-00H4}{II.\ref*{cdp:eq:C.EN.R7}}}
\expandafter\def\csname SRef@cdp:eq:C.EN.R8\endcsname{\hyperlink{tag-00H5}{II.\ref*{cdp:eq:C.EN.R8}}}
\expandafter\def\csname SRef@cdp:eq:C.EN.R9\endcsname{\hyperlink{tag-00H6}{II.\ref*{cdp:eq:C.EN.R9}}}
\expandafter\def\csname SRef@cdp:eq:C.EN.R10\endcsname{\hyperlink{tag-00GW}{II.\ref*{cdp:eq:C.EN.R10}}}
\expandafter\def\csname SRef@cdp:eq:C.EN.R11\endcsname{\hyperlink{tag-00GX}{II.\ref*{cdp:eq:C.EN.R11}}}
\expandafter\def\csname SRef@cdp:eq:C.EN.R12\endcsname{\hyperlink{tag-00GY}{II.\ref*{cdp:eq:C.EN.R12}}}
\expandafter\def\csname SRef@cdp:eq:C.HQ.4\endcsname{\hyperlink{tag-00MU}{II.\ref*{cdp:eq:C.HQ.4}}}
\expandafter\def\csname SRef@cdp:eq:C.EN.G1\endcsname{\hyperlink{tag-00GR}{II.\ref*{cdp:eq:C.EN.G1}}}
\expandafter\def\csname SRef@cdp:eq:C.EN.G2\endcsname{\hyperlink{tag-00GS}{II.\ref*{cdp:eq:C.EN.G2}}}
\expandafter\def\csname SRef@cdp:eq:C.EN.G3\endcsname{\hyperlink{tag-00GT}{II.\ref*{cdp:eq:C.EN.G3}}}
\expandafter\def\csname SRef@cdp:eq:C.EN.G4\endcsname{\hyperlink{tag-00GU}{II.\ref*{cdp:eq:C.EN.G4}}}
\expandafter\def\csname SRef@cdp:eq:C.ET.1\endcsname{\hyperlink{tag-00J5}{II.\ref*{cdp:eq:C.ET.1}}}
\expandafter\def\csname SRef@cdp:eq:C.ET.2\endcsname{\hyperlink{tag-00J9}{II.\ref*{cdp:eq:C.ET.2}}}
\expandafter\def\csname SRef@cdp:eq:C.ET.3\endcsname{\hyperlink{tag-00JA}{II.\ref*{cdp:eq:C.ET.3}}}
\expandafter\def\csname SRef@cdp:eq:C.ET.4\endcsname{\hyperlink{tag-00JB}{II.\ref*{cdp:eq:C.ET.4}}}
\expandafter\def\csname SRef@cdp:eq:C.ET.5\endcsname{\hyperlink{tag-00JC}{II.\ref*{cdp:eq:C.ET.5}}}
\expandafter\def\csname SRef@cdp:eq:C.ET.6\endcsname{\hyperlink{tag-00JD}{II.\ref*{cdp:eq:C.ET.6}}}
\expandafter\def\csname SRef@cdp:eq:C.HQ.5\endcsname{\hyperlink{tag-00MV}{II.\ref*{cdp:eq:C.HQ.5}}}
\expandafter\def\csname SRef@cdp:eq:C.ET.7\endcsname{\hyperlink{tag-00JE}{II.\ref*{cdp:eq:C.ET.7}}}
\expandafter\def\csname SRef@cdp:eq:C.ET.8\endcsname{\hyperlink{tag-00JF}{II.\ref*{cdp:eq:C.ET.8}}}
\expandafter\def\csname SRef@cdp:eq:C.ET.9\endcsname{\hyperlink{tag-00JG}{II.\ref*{cdp:eq:C.ET.9}}}
\expandafter\def\csname SRef@cdp:eq:C.ET.10\endcsname{\hyperlink{tag-00J6}{II.\ref*{cdp:eq:C.ET.10}}}
\expandafter\def\csname SRef@cdp:eq:C.ET.11\endcsname{\hyperlink{tag-00J7}{II.\ref*{cdp:eq:C.ET.11}}}
\expandafter\def\csname SRef@cdp:eq:C.ET.12\endcsname{\hyperlink{tag-00J8}{II.\ref*{cdp:eq:C.ET.12}}}
\expandafter\def\csname SRef@cdp:eq:C.SF.1\endcsname{\hyperlink{tag-010R}{II.\ref*{cdp:eq:C.SF.1}}}
\expandafter\def\csname SRef@cdp:eq:C.SF.2\endcsname{\hyperlink{tag-010W}{II.\ref*{cdp:eq:C.SF.2}}}
\expandafter\def\csname SRef@cdp:eq:C.SF.3\endcsname{\hyperlink{tag-010X}{II.\ref*{cdp:eq:C.SF.3}}}
\expandafter\def\csname SRef@cdp:eq:C.SF.4\endcsname{\hyperlink{tag-010Y}{II.\ref*{cdp:eq:C.SF.4}}}
\expandafter\def\csname SRef@cdp:eq:C.HQ.6\endcsname{\hyperlink{tag-00MW}{II.\ref*{cdp:eq:C.HQ.6}}}
\expandafter\def\csname SRef@cdp:eq:C.SF.5\endcsname{\hyperlink{tag-010Z}{II.\ref*{cdp:eq:C.SF.5}}}
\expandafter\def\csname SRef@cdp:eq:C.SF.6\endcsname{\hyperlink{tag-0110}{II.\ref*{cdp:eq:C.SF.6}}}
\expandafter\def\csname SRef@cdp:eq:C.SF.7\endcsname{\hyperlink{tag-0111}{II.\ref*{cdp:eq:C.SF.7}}}
\expandafter\def\csname SRef@cdp:eq:C.SF.8\endcsname{\hyperlink{tag-0112}{II.\ref*{cdp:eq:C.SF.8}}}
\expandafter\def\csname SRef@cdp:eq:C.SF.9\endcsname{\hyperlink{tag-0113}{II.\ref*{cdp:eq:C.SF.9}}}
\expandafter\def\csname SRef@cdp:eq:C.SF.10\endcsname{\hyperlink{tag-010S}{II.\ref*{cdp:eq:C.SF.10}}}
\expandafter\def\csname SRef@cdp:eq:C.SF.11\endcsname{\hyperlink{tag-010T}{II.\ref*{cdp:eq:C.SF.11}}}
\expandafter\def\csname SRef@cdp:eq:C.SF.12\endcsname{\hyperlink{tag-010U}{II.\ref*{cdp:eq:C.SF.12}}}
\expandafter\def\csname SRef@cdp:eq:C.SF.13\endcsname{\hyperlink{tag-010V}{II.\ref*{cdp:eq:C.SF.13}}}
\expandafter\def\csname SRef@cdp:eq:C.EO.1\endcsname{\hyperlink{tag-00HE}{II.\ref*{cdp:eq:C.EO.1}}}
\expandafter\def\csname SRef@cdp:eq:C.HQ.7\endcsname{\hyperlink{tag-00MX}{II.\ref*{cdp:eq:C.HQ.7}}}
\expandafter\def\csname SRef@cdp:eq:C.EO.2\endcsname{\hyperlink{tag-00HM}{II.\ref*{cdp:eq:C.EO.2}}}
\expandafter\def\csname SRef@cdp:eq:C.EO.3\endcsname{\hyperlink{tag-00HN}{II.\ref*{cdp:eq:C.EO.3}}}
\expandafter\def\csname SRef@cdp:eq:C.EO.4\endcsname{\hyperlink{tag-00HO}{II.\ref*{cdp:eq:C.EO.4}}}
\expandafter\def\csname SRef@cdp:eq:C.EO.5\endcsname{\hyperlink{tag-00HP}{II.\ref*{cdp:eq:C.EO.5}}}
\expandafter\def\csname SRef@cdp:eq:C.EO.6\endcsname{\hyperlink{tag-00HQ}{II.\ref*{cdp:eq:C.EO.6}}}
\expandafter\def\csname SRef@cdp:eq:C.EO.H1\endcsname{\hyperlink{tag-00HU}{II.\ref*{cdp:eq:C.EO.H1}}}
\expandafter\def\csname SRef@cdp:eq:C.EO.H2\endcsname{\hyperlink{tag-00HV}{II.\ref*{cdp:eq:C.EO.H2}}}
\expandafter\def\csname SRef@cdp:eq:C.EO.H3\endcsname{\hyperlink{tag-00HW}{II.\ref*{cdp:eq:C.EO.H3}}}
\expandafter\def\csname SRef@cdp:eq:C.EO.H4\endcsname{\hyperlink{tag-00HX}{II.\ref*{cdp:eq:C.EO.H4}}}
\expandafter\def\csname SRef@cdp:eq:C.EO.H5\endcsname{\hyperlink{tag-00HY}{II.\ref*{cdp:eq:C.EO.H5}}}
\expandafter\def\csname SRef@cdp:eq:C.HQ.8\endcsname{\hyperlink{tag-00MY}{II.\ref*{cdp:eq:C.HQ.8}}}
\expandafter\def\csname SRef@cdp:eq:C.EO.H6\endcsname{\hyperlink{tag-00HZ}{II.\ref*{cdp:eq:C.EO.H6}}}
\expandafter\def\csname SRef@cdp:eq:C.EO.H7\endcsname{\hyperlink{tag-00I0}{II.\ref*{cdp:eq:C.EO.H7}}}
\expandafter\def\csname SRef@cdp:eq:C.EO.7\endcsname{\hyperlink{tag-00HR}{II.\ref*{cdp:eq:C.EO.7}}}
\expandafter\def\csname SRef@cdp:eq:C.EO.8\endcsname{\hyperlink{tag-00HS}{II.\ref*{cdp:eq:C.EO.8}}}
\expandafter\def\csname SRef@cdp:eq:C.EO.9\endcsname{\hyperlink{tag-00HT}{II.\ref*{cdp:eq:C.EO.9}}}
\expandafter\def\csname SRef@cdp:eq:C.EO.10\endcsname{\hyperlink{tag-00HF}{II.\ref*{cdp:eq:C.EO.10}}}
\expandafter\def\csname SRef@cdp:eq:C.EO.11\endcsname{\hyperlink{tag-00HG}{II.\ref*{cdp:eq:C.EO.11}}}
\expandafter\def\csname SRef@cdp:eq:C.EO.12\endcsname{\hyperlink{tag-00HH}{II.\ref*{cdp:eq:C.EO.12}}}
\expandafter\def\csname SRef@cdp:eq:C.EO.13\endcsname{\hyperlink{tag-00HI}{II.\ref*{cdp:eq:C.EO.13}}}
\expandafter\def\csname SRef@cdp:eq:C.EO.14\endcsname{\hyperlink{tag-00HJ}{II.\ref*{cdp:eq:C.EO.14}}}
\expandafter\def\csname SRef@cdp:eq:C.HQ.9\endcsname{\hyperlink{tag-00MZ}{II.\ref*{cdp:eq:C.HQ.9}}}
\expandafter\def\csname SRef@cdp:eq:C.EO.15\endcsname{\hyperlink{tag-00HK}{II.\ref*{cdp:eq:C.EO.15}}}
\expandafter\def\csname SRef@cdp:eq:C.EO.16\endcsname{\hyperlink{tag-00HL}{II.\ref*{cdp:eq:C.EO.16}}}
\expandafter\def\csname SRef@cdp:eq:C.EH.1\endcsname{\hyperlink{tag-00FQ}{II.\ref*{cdp:eq:C.EH.1}}}
\expandafter\def\csname SRef@cdp:eq:C.EH.2\endcsname{\hyperlink{tag-00FR}{II.\ref*{cdp:eq:C.EH.2}}}
\expandafter\def\csname SRef@cdp:eq:C.EH.3\endcsname{\hyperlink{tag-00FS}{II.\ref*{cdp:eq:C.EH.3}}}
\expandafter\def\csname SRef@cdp:eq:C.EH.4\endcsname{\hyperlink{tag-00FT}{II.\ref*{cdp:eq:C.EH.4}}}
\expandafter\def\csname SRef@cdp:eq:C.EH.5\endcsname{\hyperlink{tag-00FU}{II.\ref*{cdp:eq:C.EH.5}}}
\expandafter\def\csname SRef@cdp:eq:C.EH.6\endcsname{\hyperlink{tag-00FV}{II.\ref*{cdp:eq:C.EH.6}}}
\expandafter\def\csname SRef@cdp:eq:C.EH.7\endcsname{\hyperlink{tag-00FW}{II.\ref*{cdp:eq:C.EH.7}}}
\expandafter\def\csname SRef@cdp:eq:C.EH.8\endcsname{\hyperlink{tag-00FX}{II.\ref*{cdp:eq:C.EH.8}}}
\expandafter\def\csname SRef@cdp:eq:C.HQ.10\endcsname{\hyperlink{tag-00MR}{II.\ref*{cdp:eq:C.HQ.10}}}
\expandafter\def\csname SRef@cdp:eq:C.DH.1\endcsname{\hyperlink{tag-00F6}{II.\ref*{cdp:eq:C.DH.1}}}
\expandafter\def\csname SRef@cdp:eq:C.DH.2\endcsname{\hyperlink{tag-00FG}{II.\ref*{cdp:eq:C.DH.2}}}
\expandafter\def\csname SRef@cdp:eq:C.DH.3\endcsname{\hyperlink{tag-00FH}{II.\ref*{cdp:eq:C.DH.3}}}
\expandafter\def\csname SRef@cdp:eq:C.DH.4\endcsname{\hyperlink{tag-00FI}{II.\ref*{cdp:eq:C.DH.4}}}
\expandafter\def\csname SRef@cdp:eq:C.DH.5\endcsname{\hyperlink{tag-00FJ}{II.\ref*{cdp:eq:C.DH.5}}}
\expandafter\def\csname SRef@cdp:eq:C.DH.6\endcsname{\hyperlink{tag-00FK}{II.\ref*{cdp:eq:C.DH.6}}}
\expandafter\def\csname SRef@cdp:eq:C.DH.7\endcsname{\hyperlink{tag-00FL}{II.\ref*{cdp:eq:C.DH.7}}}
\expandafter\def\csname SRef@cdp:eq:C.DH.8\endcsname{\hyperlink{tag-00FM}{II.\ref*{cdp:eq:C.DH.8}}}
\expandafter\def\csname SRef@cdp:eq:C.DH.9\endcsname{\hyperlink{tag-00FN}{II.\ref*{cdp:eq:C.DH.9}}}
\expandafter\def\csname SRef@cdp:eq:C.DH.10\endcsname{\hyperlink{tag-00F7}{II.\ref*{cdp:eq:C.DH.10}}}
\expandafter\def\csname SRef@cdp:eq:C.FD.1\endcsname{\hyperlink{tag-00K2}{II.\ref*{cdp:eq:C.FD.1}}}
\expandafter\def\csname SRef@cdp:eq:C.DH.11\endcsname{\hyperlink{tag-00F8}{II.\ref*{cdp:eq:C.DH.11}}}
\expandafter\def\csname SRef@cdp:eq:C.DH.12\endcsname{\hyperlink{tag-00F9}{II.\ref*{cdp:eq:C.DH.12}}}
\expandafter\def\csname SRef@cdp:eq:C.DH.13\endcsname{\hyperlink{tag-00FA}{II.\ref*{cdp:eq:C.DH.13}}}
\expandafter\def\csname SRef@cdp:eq:C.DH.14\endcsname{\hyperlink{tag-00FB}{II.\ref*{cdp:eq:C.DH.14}}}
\expandafter\def\csname SRef@cdp:eq:C.DH.14a\endcsname{\hyperlink{tag-00FC}{II.\ref*{cdp:eq:C.DH.14a}}}
\expandafter\def\csname SRef@cdp:eq:C.DH.14b\endcsname{\hyperlink{tag-00FD}{II.\ref*{cdp:eq:C.DH.14b}}}
\expandafter\def\csname SRef@cdp:eq:C.DH.15\endcsname{\hyperlink{tag-00FE}{II.\ref*{cdp:eq:C.DH.15}}}
\expandafter\def\csname SRef@cdp:eq:C.DH.16\endcsname{\hyperlink{tag-00FF}{II.\ref*{cdp:eq:C.DH.16}}}
\expandafter\def\csname SRef@cdp:eq:C.NH.1\endcsname{\hyperlink{tag-00S9}{II.\ref*{cdp:eq:C.NH.1}}}
\expandafter\def\csname SRef@cdp:eq:C.NH.2\endcsname{\hyperlink{tag-00SJ}{II.\ref*{cdp:eq:C.NH.2}}}
\expandafter\def\csname SRef@cdp:eq:C.FD.2\endcsname{\hyperlink{tag-00K3}{II.\ref*{cdp:eq:C.FD.2}}}
\expandafter\def\csname SRef@cdp:eq:C.NH.3\endcsname{\hyperlink{tag-00SK}{II.\ref*{cdp:eq:C.NH.3}}}
\expandafter\def\csname SRef@cdp:eq:C.NH.4\endcsname{\hyperlink{tag-00SL}{II.\ref*{cdp:eq:C.NH.4}}}
\expandafter\def\csname SRef@cdp:eq:C.NH.5\endcsname{\hyperlink{tag-00SM}{II.\ref*{cdp:eq:C.NH.5}}}
\expandafter\def\csname SRef@cdp:eq:C.NH.6\endcsname{\hyperlink{tag-00SN}{II.\ref*{cdp:eq:C.NH.6}}}
\expandafter\def\csname SRef@cdp:eq:C.NH.7\endcsname{\hyperlink{tag-00SO}{II.\ref*{cdp:eq:C.NH.7}}}
\expandafter\def\csname SRef@cdp:eq:C.NH.8\endcsname{\hyperlink{tag-00SP}{II.\ref*{cdp:eq:C.NH.8}}}
\expandafter\def\csname SRef@cdp:eq:C.NH.9\endcsname{\hyperlink{tag-00SQ}{II.\ref*{cdp:eq:C.NH.9}}}
\expandafter\def\csname SRef@cdp:eq:C.NH.10\endcsname{\hyperlink{tag-00SA}{II.\ref*{cdp:eq:C.NH.10}}}
\expandafter\def\csname SRef@cdp:eq:C.NH.11\endcsname{\hyperlink{tag-00SB}{II.\ref*{cdp:eq:C.NH.11}}}
\expandafter\def\csname SRef@cdp:eq:C.NH.12\endcsname{\hyperlink{tag-00SC}{II.\ref*{cdp:eq:C.NH.12}}}
\expandafter\def\csname SRef@cdp:eq:C.FD.3\endcsname{\hyperlink{tag-00K4}{II.\ref*{cdp:eq:C.FD.3}}}
\expandafter\def\csname SRef@cdp:eq:C.NH.13\endcsname{\hyperlink{tag-00SD}{II.\ref*{cdp:eq:C.NH.13}}}
\expandafter\def\csname SRef@cdp:eq:C.NH.14\endcsname{\hyperlink{tag-00SE}{II.\ref*{cdp:eq:C.NH.14}}}
\expandafter\def\csname SRef@cdp:eq:C.NH.15\endcsname{\hyperlink{tag-00SF}{II.\ref*{cdp:eq:C.NH.15}}}
\expandafter\def\csname SRef@cdp:eq:C.NH.16\endcsname{\hyperlink{tag-00SG}{II.\ref*{cdp:eq:C.NH.16}}}
\expandafter\def\csname SRef@cdp:eq:C.NH.17\endcsname{\hyperlink{tag-00SH}{II.\ref*{cdp:eq:C.NH.17}}}
\expandafter\def\csname SRef@cdp:eq:C.NH.18\endcsname{\hyperlink{tag-00SI}{II.\ref*{cdp:eq:C.NH.18}}}
\expandafter\def\csname SRef@cdp:eq:C.SH.1\endcsname{\hyperlink{tag-0114}{II.\ref*{cdp:eq:C.SH.1}}}
\expandafter\def\csname SRef@cdp:eq:C.SH.2\endcsname{\hyperlink{tag-0119}{II.\ref*{cdp:eq:C.SH.2}}}
\expandafter\def\csname SRef@cdp:eq:C.SH.3\endcsname{\hyperlink{tag-011A}{II.\ref*{cdp:eq:C.SH.3}}}
\expandafter\def\csname SRef@cdp:eq:C.SH.4\endcsname{\hyperlink{tag-011B}{II.\ref*{cdp:eq:C.SH.4}}}
\expandafter\def\csname SRef@cdp:eq:C.FD.4\endcsname{\hyperlink{tag-00K5}{II.\ref*{cdp:eq:C.FD.4}}}
\expandafter\def\csname SRef@cdp:eq:C.SH.5\endcsname{\hyperlink{tag-011C}{II.\ref*{cdp:eq:C.SH.5}}}
\expandafter\def\csname SRef@cdp:eq:C.SH.6\endcsname{\hyperlink{tag-011D}{II.\ref*{cdp:eq:C.SH.6}}}
\expandafter\def\csname SRef@cdp:eq:C.SH.7\endcsname{\hyperlink{tag-011E}{II.\ref*{cdp:eq:C.SH.7}}}
\expandafter\def\csname SRef@cdp:eq:C.SH.8\endcsname{\hyperlink{tag-011F}{II.\ref*{cdp:eq:C.SH.8}}}
\expandafter\def\csname SRef@cdp:eq:C.SH.9\endcsname{\hyperlink{tag-011G}{II.\ref*{cdp:eq:C.SH.9}}}
\expandafter\def\csname SRef@cdp:eq:C.SH.10\endcsname{\hyperlink{tag-0115}{II.\ref*{cdp:eq:C.SH.10}}}
\expandafter\def\csname SRef@cdp:eq:C.SH.11\endcsname{\hyperlink{tag-0116}{II.\ref*{cdp:eq:C.SH.11}}}
\expandafter\def\csname SRef@cdp:eq:C.SH.12\endcsname{\hyperlink{tag-0117}{II.\ref*{cdp:eq:C.SH.12}}}
\expandafter\def\csname SRef@cdp:eq:C.SH.13\endcsname{\hyperlink{tag-0118}{II.\ref*{cdp:eq:C.SH.13}}}
\expandafter\def\csname SRef@cdp:eq:C.RU.1\endcsname{\hyperlink{tag-010C}{II.\ref*{cdp:eq:C.RU.1}}}
\expandafter\def\csname SRef@cdp:eq:C.FD.5\endcsname{\hyperlink{tag-00K6}{II.\ref*{cdp:eq:C.FD.5}}}
\expandafter\def\csname SRef@cdp:eq:C.RU.2\endcsname{\hyperlink{tag-010D}{II.\ref*{cdp:eq:C.RU.2}}}
\expandafter\def\csname SRef@cdp:eq:C.RU.3\endcsname{\hyperlink{tag-010E}{II.\ref*{cdp:eq:C.RU.3}}}
\expandafter\def\csname SRef@cdp:eq:C.RU.4\endcsname{\hyperlink{tag-010F}{II.\ref*{cdp:eq:C.RU.4}}}
\expandafter\def\csname SRef@cdp:eq:C.RU.5\endcsname{\hyperlink{tag-010G}{II.\ref*{cdp:eq:C.RU.5}}}
\expandafter\def\csname SRef@cdp:eq:C.RU.6\endcsname{\hyperlink{tag-010H}{II.\ref*{cdp:eq:C.RU.6}}}
\expandafter\def\csname SRef@cdp:eq:C.RU.7\endcsname{\hyperlink{tag-010I}{II.\ref*{cdp:eq:C.RU.7}}}
\expandafter\def\csname SRef@cdp:eq:C.RU.8\endcsname{\hyperlink{tag-010J}{II.\ref*{cdp:eq:C.RU.8}}}
\expandafter\def\csname SRef@cdp:eq:C.RU.9\endcsname{\hyperlink{tag-010K}{II.\ref*{cdp:eq:C.RU.9}}}
\expandafter\def\csname SRef@cdp:eq:C.HA.1\endcsname{\hyperlink{tag-00L9}{II.\ref*{cdp:eq:C.HA.1}}}
\expandafter\def\csname SRef@cdp:eq:C.HA.2\endcsname{\hyperlink{tag-00LA}{II.\ref*{cdp:eq:C.HA.2}}}
\expandafter\def\csname SRef@cdp:eq:C.FD.6\endcsname{\hyperlink{tag-00K7}{II.\ref*{cdp:eq:C.FD.6}}}
\expandafter\def\csname SRef@cdp:eq:C.HA.3\endcsname{\hyperlink{tag-00LB}{II.\ref*{cdp:eq:C.HA.3}}}
\expandafter\def\csname SRef@cdp:eq:C.HA.4\endcsname{\hyperlink{tag-00LC}{II.\ref*{cdp:eq:C.HA.4}}}
\expandafter\def\csname SRef@cdp:eq:C.HA.5\endcsname{\hyperlink{tag-00LD}{II.\ref*{cdp:eq:C.HA.5}}}
\expandafter\def\csname SRef@cdp:eq:C.HA.6\endcsname{\hyperlink{tag-00LE}{II.\ref*{cdp:eq:C.HA.6}}}
\expandafter\def\csname SRef@cdp:eq:C.HC.1\endcsname{\hyperlink{tag-00LO}{II.\ref*{cdp:eq:C.HC.1}}}
\expandafter\def\csname SRef@cdp:eq:C.HC.2\endcsname{\hyperlink{tag-00LR}{II.\ref*{cdp:eq:C.HC.2}}}
\expandafter\def\csname SRef@cdp:eq:C.HC.3\endcsname{\hyperlink{tag-00LS}{II.\ref*{cdp:eq:C.HC.3}}}
\expandafter\def\csname SRef@cdp:eq:C.HC.4\endcsname{\hyperlink{tag-00LT}{II.\ref*{cdp:eq:C.HC.4}}}
\expandafter\def\csname SRef@cdp:eq:C.HC.5\endcsname{\hyperlink{tag-00LU}{II.\ref*{cdp:eq:C.HC.5}}}
\expandafter\def\csname SRef@cdp:eq:C.HC.6\endcsname{\hyperlink{tag-00LV}{II.\ref*{cdp:eq:C.HC.6}}}
\expandafter\def\csname SRef@cdp:eq:C.RF.5\endcsname{\hyperlink{tag-00YY}{II.\ref*{cdp:eq:C.RF.5}}}
\expandafter\def\csname SRef@cdp:eq:C.HC.7\endcsname{\hyperlink{tag-00LW}{II.\ref*{cdp:eq:C.HC.7}}}
\expandafter\def\csname SRef@cdp:eq:C.HC.8\endcsname{\hyperlink{tag-00LX}{II.\ref*{cdp:eq:C.HC.8}}}
\expandafter\def\csname SRef@cdp:eq:C.HC.9\endcsname{\hyperlink{tag-00LY}{II.\ref*{cdp:eq:C.HC.9}}}
\expandafter\def\csname SRef@cdp:eq:C.HC.10\endcsname{\hyperlink{tag-00LP}{II.\ref*{cdp:eq:C.HC.10}}}
\expandafter\def\csname SRef@cdp:eq:C.HC.11\endcsname{\hyperlink{tag-00LQ}{II.\ref*{cdp:eq:C.HC.11}}}
\expandafter\def\csname SRef@cdp:eq:C.HG.1\endcsname{\hyperlink{tag-00LZ}{II.\ref*{cdp:eq:C.HG.1}}}
\expandafter\def\csname SRef@cdp:eq:C.HG.2\endcsname{\hyperlink{tag-00M8}{II.\ref*{cdp:eq:C.HG.2}}}
\expandafter\def\csname SRef@cdp:eq:C.HG.3\endcsname{\hyperlink{tag-00M9}{II.\ref*{cdp:eq:C.HG.3}}}
\expandafter\def\csname SRef@cdp:eq:C.HG.4\endcsname{\hyperlink{tag-00MA}{II.\ref*{cdp:eq:C.HG.4}}}
\expandafter\def\csname SRef@cdp:eq:C.HG.5\endcsname{\hyperlink{tag-00MB}{II.\ref*{cdp:eq:C.HG.5}}}
\expandafter\def\csname SRef@cdp:eq:C.RF.1\endcsname{\hyperlink{tag-00YU}{II.\ref*{cdp:eq:C.RF.1}}}
\expandafter\def\csname SRef@cdp:eq:C.HG.6\endcsname{\hyperlink{tag-00MC}{II.\ref*{cdp:eq:C.HG.6}}}
\expandafter\def\csname SRef@cdp:eq:C.HG.7\endcsname{\hyperlink{tag-00MD}{II.\ref*{cdp:eq:C.HG.7}}}
\expandafter\def\csname SRef@cdp:eq:C.HG.8\endcsname{\hyperlink{tag-00ME}{II.\ref*{cdp:eq:C.HG.8}}}
\expandafter\def\csname SRef@cdp:eq:C.HG.9\endcsname{\hyperlink{tag-00MF}{II.\ref*{cdp:eq:C.HG.9}}}
\expandafter\def\csname SRef@cdp:eq:C.HG.10\endcsname{\hyperlink{tag-00M0}{II.\ref*{cdp:eq:C.HG.10}}}
\expandafter\def\csname SRef@cdp:eq:C.HG.11\endcsname{\hyperlink{tag-00M1}{II.\ref*{cdp:eq:C.HG.11}}}
\expandafter\def\csname SRef@cdp:eq:C.HG.12\endcsname{\hyperlink{tag-00M2}{II.\ref*{cdp:eq:C.HG.12}}}
\expandafter\def\csname SRef@cdp:eq:C.HG.13\endcsname{\hyperlink{tag-00M3}{II.\ref*{cdp:eq:C.HG.13}}}
\expandafter\def\csname SRef@cdp:eq:C.HG.14\endcsname{\hyperlink{tag-00M4}{II.\ref*{cdp:eq:C.HG.14}}}
\expandafter\def\csname SRef@cdp:eq:C.HG.15\endcsname{\hyperlink{tag-00M5}{II.\ref*{cdp:eq:C.HG.15}}}
\expandafter\def\csname SRef@cdp:eq:C.RF.2\endcsname{\hyperlink{tag-00YV}{II.\ref*{cdp:eq:C.RF.2}}}
\expandafter\def\csname SRef@cdp:eq:C.HG.16\endcsname{\hyperlink{tag-00M6}{II.\ref*{cdp:eq:C.HG.16}}}
\expandafter\def\csname SRef@cdp:eq:C.HG.17\endcsname{\hyperlink{tag-00M7}{II.\ref*{cdp:eq:C.HG.17}}}
\expandafter\def\csname SRef@cdp:eq:C.NS.1\endcsname{\hyperlink{tag-00TA}{II.\ref*{cdp:eq:C.NS.1}}}
\expandafter\def\csname SRef@cdp:eq:C.NS.2\endcsname{\hyperlink{tag-00TL}{II.\ref*{cdp:eq:C.NS.2}}}
\expandafter\def\csname SRef@cdp:eq:C.NS.3\endcsname{\hyperlink{tag-00TN}{II.\ref*{cdp:eq:C.NS.3}}}
\expandafter\def\csname SRef@cdp:eq:C.NS.4\endcsname{\hyperlink{tag-00TO}{II.\ref*{cdp:eq:C.NS.4}}}
\expandafter\def\csname SRef@cdp:eq:C.NS.5\endcsname{\hyperlink{tag-00TP}{II.\ref*{cdp:eq:C.NS.5}}}
\expandafter\def\csname SRef@cdp:eq:C.NS.6\endcsname{\hyperlink{tag-00TQ}{II.\ref*{cdp:eq:C.NS.6}}}
\expandafter\def\csname SRef@cdp:eq:C.NS.7\endcsname{\hyperlink{tag-00TR}{II.\ref*{cdp:eq:C.NS.7}}}
\expandafter\def\csname SRef@cdp:eq:C.NS.8\endcsname{\hyperlink{tag-00TS}{II.\ref*{cdp:eq:C.NS.8}}}
\expandafter\def\csname SRef@cdp:eq:C.RF.3\endcsname{\hyperlink{tag-00YW}{II.\ref*{cdp:eq:C.RF.3}}}
\expandafter\def\csname SRef@cdp:eq:C.NS.9\endcsname{\hyperlink{tag-00TT}{II.\ref*{cdp:eq:C.NS.9}}}
\expandafter\def\csname SRef@cdp:eq:C.NS.10\endcsname{\hyperlink{tag-00TB}{II.\ref*{cdp:eq:C.NS.10}}}
\expandafter\def\csname SRef@cdp:eq:C.NS.11\endcsname{\hyperlink{tag-00TC}{II.\ref*{cdp:eq:C.NS.11}}}
\expandafter\def\csname SRef@cdp:eq:C.NS.12\endcsname{\hyperlink{tag-00TD}{II.\ref*{cdp:eq:C.NS.12}}}
\expandafter\def\csname SRef@cdp:eq:C.NS.13\endcsname{\hyperlink{tag-00TE}{II.\ref*{cdp:eq:C.NS.13}}}
\expandafter\def\csname SRef@cdp:eq:C.NS.14\endcsname{\hyperlink{tag-00TF}{II.\ref*{cdp:eq:C.NS.14}}}
\expandafter\def\csname SRef@cdp:eq:C.NS.15\endcsname{\hyperlink{tag-00TG}{II.\ref*{cdp:eq:C.NS.15}}}
\expandafter\def\csname SRef@cdp:eq:C.NS.16\endcsname{\hyperlink{tag-00TH}{II.\ref*{cdp:eq:C.NS.16}}}
\expandafter\def\csname SRef@cdp:eq:C.NS.17\endcsname{\hyperlink{tag-00TI}{II.\ref*{cdp:eq:C.NS.17}}}
\expandafter\def\csname SRef@cdp:eq:C.NS.18\endcsname{\hyperlink{tag-00TJ}{II.\ref*{cdp:eq:C.NS.18}}}
\expandafter\def\csname SRef@cdp:eq:C.RF.4\endcsname{\hyperlink{tag-00YX}{II.\ref*{cdp:eq:C.RF.4}}}
\expandafter\def\csname SRef@cdp:eq:C.NS.19\endcsname{\hyperlink{tag-00TK}{II.\ref*{cdp:eq:C.NS.19}}}
\expandafter\def\csname SRef@cdp:eq:C.NS.20\endcsname{\hyperlink{tag-00TM}{II.\ref*{cdp:eq:C.NS.20}}}
\expandafter\def\csname SRef@cdp:eq:C.SN.1\endcsname{\hyperlink{tag-0122}{II.\ref*{cdp:eq:C.SN.1}}}
\expandafter\def\csname SRef@cdp:eq:C.SN.2\endcsname{\hyperlink{tag-0124}{II.\ref*{cdp:eq:C.SN.2}}}
\expandafter\def\csname SRef@cdp:eq:C.SN.3\endcsname{\hyperlink{tag-0125}{II.\ref*{cdp:eq:C.SN.3}}}
\expandafter\def\csname SRef@cdp:eq:C.SN.4\endcsname{\hyperlink{tag-0126}{II.\ref*{cdp:eq:C.SN.4}}}
\expandafter\def\csname SRef@cdp:eq:C.SN.5\endcsname{\hyperlink{tag-0127}{II.\ref*{cdp:eq:C.SN.5}}}
\expandafter\def\csname SRef@cdp:eq:C.SN.6\endcsname{\hyperlink{tag-0128}{II.\ref*{cdp:eq:C.SN.6}}}
\expandafter\def\csname SRef@cdp:eq:C.SN.7\endcsname{\hyperlink{tag-0129}{II.\ref*{cdp:eq:C.SN.7}}}
\expandafter\def\csname SRef@cdp:eq:C.SN.8\endcsname{\hyperlink{tag-012A}{II.\ref*{cdp:eq:C.SN.8}}}
\expandafter\def\csname SRef@cdp:eq:C.RF.9\endcsname{\hyperlink{tag-00Z2}{II.\ref*{cdp:eq:C.RF.9}}}
\expandafter\def\csname SRef@cdp:eq:C.SN.9\endcsname{\hyperlink{tag-012B}{II.\ref*{cdp:eq:C.SN.9}}}
\expandafter\def\csname SRef@cdp:eq:C.SN.10\endcsname{\hyperlink{tag-0123}{II.\ref*{cdp:eq:C.SN.10}}}
\expandafter\def\csname SRef@cdp:eq:C.MX.1\endcsname{\hyperlink{tag-00RR}{II.\ref*{cdp:eq:C.MX.1}}}
\expandafter\def\csname SRef@cdp:eq:C.MX.2\endcsname{\hyperlink{tag-00RT}{II.\ref*{cdp:eq:C.MX.2}}}
\expandafter\def\csname SRef@cdp:eq:C.MX.3\endcsname{\hyperlink{tag-00RU}{II.\ref*{cdp:eq:C.MX.3}}}
\expandafter\def\csname SRef@cdp:eq:C.MX.4\endcsname{\hyperlink{tag-00RV}{II.\ref*{cdp:eq:C.MX.4}}}
\expandafter\def\csname SRef@cdp:eq:C.MX.5\endcsname{\hyperlink{tag-00RW}{II.\ref*{cdp:eq:C.MX.5}}}
\expandafter\def\csname SRef@cdp:eq:C.MX.6\endcsname{\hyperlink{tag-00RX}{II.\ref*{cdp:eq:C.MX.6}}}
\expandafter\def\csname SRef@cdp:eq:C.MX.7\endcsname{\hyperlink{tag-00RY}{II.\ref*{cdp:eq:C.MX.7}}}
\expandafter\def\csname SRef@cdp:eq:C.MX.8\endcsname{\hyperlink{tag-00RZ}{II.\ref*{cdp:eq:C.MX.8}}}
\expandafter\def\csname SRef@cdp:eq:C.RF.6\endcsname{\hyperlink{tag-00YZ}{II.\ref*{cdp:eq:C.RF.6}}}
\expandafter\def\csname SRef@cdp:eq:C.MX.9\endcsname{\hyperlink{tag-00S0}{II.\ref*{cdp:eq:C.MX.9}}}
\expandafter\def\csname SRef@cdp:eq:C.MX.10\endcsname{\hyperlink{tag-00RS}{II.\ref*{cdp:eq:C.MX.10}}}
\expandafter\def\csname SRef@cdp:eq:C.WH.0\endcsname{\hyperlink{tag-016X}{II.\ref*{cdp:eq:C.WH.0}}}
\expandafter\def\csname SRef@cdp:eq:C.WH.1\endcsname{\hyperlink{tag-016Y}{II.\ref*{cdp:eq:C.WH.1}}}
\expandafter\def\csname SRef@cdp:eq:C.WH.2\endcsname{\hyperlink{tag-0177}{II.\ref*{cdp:eq:C.WH.2}}}
\expandafter\def\csname SRef@cdp:eq:C.WH.3\endcsname{\hyperlink{tag-0178}{II.\ref*{cdp:eq:C.WH.3}}}
\expandafter\def\csname SRef@cdp:eq:C.WH.4\endcsname{\hyperlink{tag-0179}{II.\ref*{cdp:eq:C.WH.4}}}
\expandafter\def\csname SRef@cdp:eq:C.WH.5\endcsname{\hyperlink{tag-017A}{II.\ref*{cdp:eq:C.WH.5}}}
\expandafter\def\csname SRef@cdp:eq:C.WH.6\endcsname{\hyperlink{tag-017B}{II.\ref*{cdp:eq:C.WH.6}}}
\expandafter\def\csname SRef@cdp:eq:C.WH.7\endcsname{\hyperlink{tag-017C}{II.\ref*{cdp:eq:C.WH.7}}}
\expandafter\def\csname SRef@cdp:eq:C.RF.7\endcsname{\hyperlink{tag-00Z0}{II.\ref*{cdp:eq:C.RF.7}}}
\expandafter\def\csname SRef@cdp:eq:C.WH.8\endcsname{\hyperlink{tag-017D}{II.\ref*{cdp:eq:C.WH.8}}}
\expandafter\def\csname SRef@cdp:eq:C.WH.9\endcsname{\hyperlink{tag-017E}{II.\ref*{cdp:eq:C.WH.9}}}
\expandafter\def\csname SRef@cdp:eq:C.WH.10\endcsname{\hyperlink{tag-016Z}{II.\ref*{cdp:eq:C.WH.10}}}
\expandafter\def\csname SRef@cdp:eq:C.WH.11\endcsname{\hyperlink{tag-0170}{II.\ref*{cdp:eq:C.WH.11}}}
\expandafter\def\csname SRef@cdp:eq:C.WH.12\endcsname{\hyperlink{tag-0171}{II.\ref*{cdp:eq:C.WH.12}}}
\expandafter\def\csname SRef@cdp:eq:C.WH.13\endcsname{\hyperlink{tag-0172}{II.\ref*{cdp:eq:C.WH.13}}}
\expandafter\def\csname SRef@cdp:eq:C.WH.14\endcsname{\hyperlink{tag-0173}{II.\ref*{cdp:eq:C.WH.14}}}
\expandafter\def\csname SRef@cdp:eq:C.WH.15\endcsname{\hyperlink{tag-0174}{II.\ref*{cdp:eq:C.WH.15}}}
\expandafter\def\csname SRef@cdp:eq:C.WH.16\endcsname{\hyperlink{tag-0175}{II.\ref*{cdp:eq:C.WH.16}}}
\expandafter\def\csname SRef@cdp:eq:C.WH.17\endcsname{\hyperlink{tag-0176}{II.\ref*{cdp:eq:C.WH.17}}}
\expandafter\def\csname SRef@cdp:eq:C.RF.8\endcsname{\hyperlink{tag-00Z1}{II.\ref*{cdp:eq:C.RF.8}}}
\expandafter\def\csname SRef@cdp:eq:C.AS.1\endcsname{\hyperlink{tag-00AG}{II.\ref*{cdp:eq:C.AS.1}}}
\expandafter\def\csname SRef@cdp:eq:C.AS.2\endcsname{\hyperlink{tag-00AR}{II.\ref*{cdp:eq:C.AS.2}}}
\expandafter\def\csname SRef@cdp:eq:C.AS.3\endcsname{\hyperlink{tag-00AS}{II.\ref*{cdp:eq:C.AS.3}}}
\expandafter\def\csname SRef@cdp:eq:C.AS.4\endcsname{\hyperlink{tag-00AT}{II.\ref*{cdp:eq:C.AS.4}}}
\expandafter\def\csname SRef@cdp:eq:C.AS.5\endcsname{\hyperlink{tag-00AU}{II.\ref*{cdp:eq:C.AS.5}}}
\expandafter\def\csname SRef@cdp:eq:C.AS.6\endcsname{\hyperlink{tag-00AV}{II.\ref*{cdp:eq:C.AS.6}}}
\expandafter\def\csname SRef@cdp:eq:C.AS.7\endcsname{\hyperlink{tag-00AW}{II.\ref*{cdp:eq:C.AS.7}}}
\expandafter\def\csname SRef@cdp:eq:C.AS.8\endcsname{\hyperlink{tag-00AX}{II.\ref*{cdp:eq:C.AS.8}}}
\expandafter\def\csname SRef@cdp:eq:C.AS.9\endcsname{\hyperlink{tag-00AY}{II.\ref*{cdp:eq:C.AS.9}}}
\expandafter\def\csname SRef@cdp:eq:C.AS.10\endcsname{\hyperlink{tag-00AH}{II.\ref*{cdp:eq:C.AS.10}}}
\expandafter\def\csname SRef@cdp:eq:C.CQ.1\endcsname{\hyperlink{tag-00CJ}{II.\ref*{cdp:eq:C.CQ.1}}}
\expandafter\def\csname SRef@cdp:eq:C.AS.11\endcsname{\hyperlink{tag-00AI}{II.\ref*{cdp:eq:C.AS.11}}}
\expandafter\def\csname SRef@cdp:eq:C.AS.12\endcsname{\hyperlink{tag-00AJ}{II.\ref*{cdp:eq:C.AS.12}}}
\expandafter\def\csname SRef@cdp:eq:C.AS.13\endcsname{\hyperlink{tag-00AK}{II.\ref*{cdp:eq:C.AS.13}}}
\expandafter\def\csname SRef@cdp:eq:C.AS.14\endcsname{\hyperlink{tag-00AL}{II.\ref*{cdp:eq:C.AS.14}}}
\expandafter\def\csname SRef@cdp:eq:C.AS.15\endcsname{\hyperlink{tag-00AM}{II.\ref*{cdp:eq:C.AS.15}}}
\expandafter\def\csname SRef@cdp:eq:C.AS.16\endcsname{\hyperlink{tag-00AN}{II.\ref*{cdp:eq:C.AS.16}}}
\expandafter\def\csname SRef@cdp:eq:C.AS.17\endcsname{\hyperlink{tag-00AO}{II.\ref*{cdp:eq:C.AS.17}}}
\expandafter\def\csname SRef@cdp:eq:C.AS.18\endcsname{\hyperlink{tag-00AP}{II.\ref*{cdp:eq:C.AS.18}}}
\expandafter\def\csname SRef@cdp:eq:C.AS.19\endcsname{\hyperlink{tag-00AQ}{II.\ref*{cdp:eq:C.AS.19}}}
\expandafter\def\csname SRef@cdp:eq:C.HR.1\endcsname{\hyperlink{tag-00N0}{II.\ref*{cdp:eq:C.HR.1}}}
\expandafter\def\csname SRef@cdp:eq:C.CQ.2\endcsname{\hyperlink{tag-00CN}{II.\ref*{cdp:eq:C.CQ.2}}}
\expandafter\def\csname SRef@cdp:eq:C.HR.2\endcsname{\hyperlink{tag-00N1}{II.\ref*{cdp:eq:C.HR.2}}}
\expandafter\def\csname SRef@cdp:eq:C.HR.3\endcsname{\hyperlink{tag-00N2}{II.\ref*{cdp:eq:C.HR.3}}}
\expandafter\def\csname SRef@cdp:eq:C.HR.4\endcsname{\hyperlink{tag-00N3}{II.\ref*{cdp:eq:C.HR.4}}}
\expandafter\def\csname SRef@cdp:eq:C.HR.5\endcsname{\hyperlink{tag-00N4}{II.\ref*{cdp:eq:C.HR.5}}}
\expandafter\def\csname SRef@cdp:eq:C.HR.6\endcsname{\hyperlink{tag-00N5}{II.\ref*{cdp:eq:C.HR.6}}}
\expandafter\def\csname SRef@cdp:eq:C.HR.7\endcsname{\hyperlink{tag-00N6}{II.\ref*{cdp:eq:C.HR.7}}}
\expandafter\def\csname SRef@cdp:eq:C.HR.8\endcsname{\hyperlink{tag-00N7}{II.\ref*{cdp:eq:C.HR.8}}}
\expandafter\def\csname SRef@cdp:eq:C.TL.1\endcsname{\hyperlink{tag-013Q}{II.\ref*{cdp:eq:C.TL.1}}}
\expandafter\def\csname SRef@cdp:eq:C.TL.2\endcsname{\hyperlink{tag-013R}{II.\ref*{cdp:eq:C.TL.2}}}
\expandafter\def\csname SRef@cdp:eq:C.TL.3\endcsname{\hyperlink{tag-013S}{II.\ref*{cdp:eq:C.TL.3}}}
\expandafter\def\csname SRef@cdp:eq:C.CQ.3\endcsname{\hyperlink{tag-00CO}{II.\ref*{cdp:eq:C.CQ.3}}}
\expandafter\def\csname SRef@cdp:eq:C.TL.4\endcsname{\hyperlink{tag-013T}{II.\ref*{cdp:eq:C.TL.4}}}
\expandafter\def\csname SRef@cdp:eq:C.TL.5\endcsname{\hyperlink{tag-013U}{II.\ref*{cdp:eq:C.TL.5}}}
\expandafter\def\csname SRef@cdp:eq:C.TL.6\endcsname{\hyperlink{tag-013V}{II.\ref*{cdp:eq:C.TL.6}}}
\expandafter\def\csname SRef@cdp:eq:C.TL.7\endcsname{\hyperlink{tag-013W}{II.\ref*{cdp:eq:C.TL.7}}}
\expandafter\def\csname SRef@cdp:eq:C.TL.8\endcsname{\hyperlink{tag-013X}{II.\ref*{cdp:eq:C.TL.8}}}
\expandafter\def\csname SRef@cdp:eq:C.WM.1\endcsname{\hyperlink{tag-017F}{II.\ref*{cdp:eq:C.WM.1}}}
\expandafter\def\csname SRef@cdp:eq:C.WM.2\endcsname{\hyperlink{tag-017G}{II.\ref*{cdp:eq:C.WM.2}}}
\expandafter\def\csname SRef@cdp:eq:C.WM.3\endcsname{\hyperlink{tag-017H}{II.\ref*{cdp:eq:C.WM.3}}}
\expandafter\def\csname SRef@cdp:eq:C.WM.4\endcsname{\hyperlink{tag-017I}{II.\ref*{cdp:eq:C.WM.4}}}
\expandafter\def\csname SRef@cdp:eq:C.WM.5\endcsname{\hyperlink{tag-017J}{II.\ref*{cdp:eq:C.WM.5}}}
\expandafter\def\csname SRef@cdp:eq:C.CQ.4\endcsname{\hyperlink{tag-00CP}{II.\ref*{cdp:eq:C.CQ.4}}}
\expandafter\def\csname SRef@cdp:eq:C.WM.6\endcsname{\hyperlink{tag-017K}{II.\ref*{cdp:eq:C.WM.6}}}
\expandafter\def\csname SRef@cdp:eq:C.WM.7\endcsname{\hyperlink{tag-017L}{II.\ref*{cdp:eq:C.WM.7}}}
\expandafter\def\csname SRef@cdp:eq:C.WM.8\endcsname{\hyperlink{tag-017M}{II.\ref*{cdp:eq:C.WM.8}}}
\expandafter\def\csname SRef@cdp:eq:C.WM.9\endcsname{\hyperlink{tag-017N}{II.\ref*{cdp:eq:C.WM.9}}}
\expandafter\def\csname SRef@cdp:eq:C.AD.1\endcsname{\hyperlink{tag-009B}{II.\ref*{cdp:eq:C.AD.1}}}
\expandafter\def\csname SRef@cdp:eq:C.AD.2\endcsname{\hyperlink{tag-009E}{II.\ref*{cdp:eq:C.AD.2}}}
\expandafter\def\csname SRef@cdp:eq:C.AD.3\endcsname{\hyperlink{tag-009F}{II.\ref*{cdp:eq:C.AD.3}}}
\expandafter\def\csname SRef@cdp:eq:C.AD.4\endcsname{\hyperlink{tag-009G}{II.\ref*{cdp:eq:C.AD.4}}}
\expandafter\def\csname SRef@cdp:eq:C.AD.5\endcsname{\hyperlink{tag-009H}{II.\ref*{cdp:eq:C.AD.5}}}
\expandafter\def\csname SRef@cdp:eq:C.AD.6\endcsname{\hyperlink{tag-009I}{II.\ref*{cdp:eq:C.AD.6}}}
\expandafter\def\csname SRef@cdp:eq:C.CQ.5\endcsname{\hyperlink{tag-00CQ}{II.\ref*{cdp:eq:C.CQ.5}}}
\expandafter\def\csname SRef@cdp:eq:C.AD.7\endcsname{\hyperlink{tag-009J}{II.\ref*{cdp:eq:C.AD.7}}}
\expandafter\def\csname SRef@cdp:eq:C.AD.8\endcsname{\hyperlink{tag-009K}{II.\ref*{cdp:eq:C.AD.8}}}
\expandafter\def\csname SRef@cdp:eq:C.AD.9\endcsname{\hyperlink{tag-009L}{II.\ref*{cdp:eq:C.AD.9}}}
\expandafter\def\csname SRef@cdp:eq:C.AD.10\endcsname{\hyperlink{tag-009C}{II.\ref*{cdp:eq:C.AD.10}}}
\expandafter\def\csname SRef@cdp:eq:C.AD.11\endcsname{\hyperlink{tag-009D}{II.\ref*{cdp:eq:C.AD.11}}}
\expandafter\def\csname SRef@cdp:eq:C.HT.1\endcsname{\hyperlink{tag-00N8}{II.\ref*{cdp:eq:C.HT.1}}}
\expandafter\def\csname SRef@cdp:eq:C.HT.2\endcsname{\hyperlink{tag-00NC}{II.\ref*{cdp:eq:C.HT.2}}}
\expandafter\def\csname SRef@cdp:eq:C.HT.3\endcsname{\hyperlink{tag-00ND}{II.\ref*{cdp:eq:C.HT.3}}}
\expandafter\def\csname SRef@cdp:eq:C.HT.4\endcsname{\hyperlink{tag-00NE}{II.\ref*{cdp:eq:C.HT.4}}}
\expandafter\def\csname SRef@cdp:eq:C.HT.5\endcsname{\hyperlink{tag-00NF}{II.\ref*{cdp:eq:C.HT.5}}}
\expandafter\def\csname SRef@cdp:eq:C.CQ.6\endcsname{\hyperlink{tag-00CR}{II.\ref*{cdp:eq:C.CQ.6}}}
\expandafter\def\csname SRef@cdp:eq:C.HT.6\endcsname{\hyperlink{tag-00NG}{II.\ref*{cdp:eq:C.HT.6}}}
\expandafter\def\csname SRef@cdp:eq:C.HT.7\endcsname{\hyperlink{tag-00NH}{II.\ref*{cdp:eq:C.HT.7}}}
\expandafter\def\csname SRef@cdp:eq:C.HT.8\endcsname{\hyperlink{tag-00NI}{II.\ref*{cdp:eq:C.HT.8}}}
\expandafter\def\csname SRef@cdp:eq:C.HT.9\endcsname{\hyperlink{tag-00NJ}{II.\ref*{cdp:eq:C.HT.9}}}
\expandafter\def\csname SRef@cdp:eq:C.HT.10\endcsname{\hyperlink{tag-00N9}{II.\ref*{cdp:eq:C.HT.10}}}
\expandafter\def\csname SRef@cdp:eq:C.HT.11\endcsname{\hyperlink{tag-00NA}{II.\ref*{cdp:eq:C.HT.11}}}
\expandafter\def\csname SRef@cdp:eq:C.HT.12\endcsname{\hyperlink{tag-00NB}{II.\ref*{cdp:eq:C.HT.12}}}
\expandafter\def\csname SRef@cdp:eq:C.EW.1\endcsname{\hyperlink{tag-00JH}{II.\ref*{cdp:eq:C.EW.1}}}
\expandafter\def\csname SRef@cdp:eq:C.EW.2\endcsname{\hyperlink{tag-00JS}{II.\ref*{cdp:eq:C.EW.2}}}
\expandafter\def\csname SRef@cdp:eq:C.EW.3\endcsname{\hyperlink{tag-00JV}{II.\ref*{cdp:eq:C.EW.3}}}
\expandafter\def\csname SRef@cdp:eq:C.CQ.7\endcsname{\hyperlink{tag-00CS}{II.\ref*{cdp:eq:C.CQ.7}}}
\expandafter\def\csname SRef@cdp:eq:C.EW.4\endcsname{\hyperlink{tag-00JW}{II.\ref*{cdp:eq:C.EW.4}}}
\expandafter\def\csname SRef@cdp:eq:C.EW.5\endcsname{\hyperlink{tag-00JX}{II.\ref*{cdp:eq:C.EW.5}}}
\expandafter\def\csname SRef@cdp:eq:C.EW.6\endcsname{\hyperlink{tag-00JY}{II.\ref*{cdp:eq:C.EW.6}}}
\expandafter\def\csname SRef@cdp:eq:C.EW.7\endcsname{\hyperlink{tag-00JZ}{II.\ref*{cdp:eq:C.EW.7}}}
\expandafter\def\csname SRef@cdp:eq:C.EW.8\endcsname{\hyperlink{tag-00K0}{II.\ref*{cdp:eq:C.EW.8}}}
\expandafter\def\csname SRef@cdp:eq:C.EW.9\endcsname{\hyperlink{tag-00K1}{II.\ref*{cdp:eq:C.EW.9}}}
\expandafter\def\csname SRef@cdp:eq:C.EW.10\endcsname{\hyperlink{tag-00JI}{II.\ref*{cdp:eq:C.EW.10}}}
\expandafter\def\csname SRef@cdp:eq:C.EW.11\endcsname{\hyperlink{tag-00JJ}{II.\ref*{cdp:eq:C.EW.11}}}
\expandafter\def\csname SRef@cdp:eq:C.EW.12\endcsname{\hyperlink{tag-00JK}{II.\ref*{cdp:eq:C.EW.12}}}
\expandafter\def\csname SRef@cdp:eq:C.EW.13\endcsname{\hyperlink{tag-00JL}{II.\ref*{cdp:eq:C.EW.13}}}
\expandafter\def\csname SRef@cdp:eq:C.CQ.8\endcsname{\hyperlink{tag-00CT}{II.\ref*{cdp:eq:C.CQ.8}}}
\expandafter\def\csname SRef@cdp:eq:C.EW.14\endcsname{\hyperlink{tag-00JM}{II.\ref*{cdp:eq:C.EW.14}}}
\expandafter\def\csname SRef@cdp:eq:C.EW.15\endcsname{\hyperlink{tag-00JN}{II.\ref*{cdp:eq:C.EW.15}}}
\expandafter\def\csname SRef@cdp:eq:C.EW.16\endcsname{\hyperlink{tag-00JO}{II.\ref*{cdp:eq:C.EW.16}}}
\expandafter\def\csname SRef@cdp:eq:C.EW.17\endcsname{\hyperlink{tag-00JP}{II.\ref*{cdp:eq:C.EW.17}}}
\expandafter\def\csname SRef@cdp:eq:C.EW.18\endcsname{\hyperlink{tag-00JQ}{II.\ref*{cdp:eq:C.EW.18}}}
\expandafter\def\csname SRef@cdp:eq:C.EW.19\endcsname{\hyperlink{tag-00JR}{II.\ref*{cdp:eq:C.EW.19}}}
\expandafter\def\csname SRef@cdp:eq:C.EW.20\endcsname{\hyperlink{tag-00JT}{II.\ref*{cdp:eq:C.EW.20}}}
\expandafter\def\csname SRef@cdp:eq:C.EW.21\endcsname{\hyperlink{tag-00JU}{II.\ref*{cdp:eq:C.EW.21}}}
\expandafter\def\csname SRef@cdp:eq:C.PR.1\endcsname{\hyperlink{tag-00VH}{II.\ref*{cdp:eq:C.PR.1}}}
\expandafter\def\csname SRef@cdp:eq:C.PR.2\endcsname{\hyperlink{tag-00VR}{II.\ref*{cdp:eq:C.PR.2}}}
\expandafter\def\csname SRef@cdp:eq:C.CQ.9\endcsname{\hyperlink{tag-00CU}{II.\ref*{cdp:eq:C.CQ.9}}}
\expandafter\def\csname SRef@cdp:eq:C.PR.3\endcsname{\hyperlink{tag-00VS}{II.\ref*{cdp:eq:C.PR.3}}}
\expandafter\def\csname SRef@cdp:eq:C.PR.4\endcsname{\hyperlink{tag-00VT}{II.\ref*{cdp:eq:C.PR.4}}}
\expandafter\def\csname SRef@cdp:eq:C.PR.5\endcsname{\hyperlink{tag-00VU}{II.\ref*{cdp:eq:C.PR.5}}}
\expandafter\def\csname SRef@cdp:eq:C.PR.6\endcsname{\hyperlink{tag-00VV}{II.\ref*{cdp:eq:C.PR.6}}}
\expandafter\def\csname SRef@cdp:eq:C.PR.7\endcsname{\hyperlink{tag-00VW}{II.\ref*{cdp:eq:C.PR.7}}}
\expandafter\def\csname SRef@cdp:eq:C.PR.8\endcsname{\hyperlink{tag-00VX}{II.\ref*{cdp:eq:C.PR.8}}}
\expandafter\def\csname SRef@cdp:eq:C.PR.9\endcsname{\hyperlink{tag-00VY}{II.\ref*{cdp:eq:C.PR.9}}}
\expandafter\def\csname SRef@cdp:eq:C.PR.10\endcsname{\hyperlink{tag-00VI}{II.\ref*{cdp:eq:C.PR.10}}}
\expandafter\def\csname SRef@cdp:eq:C.PR.11\endcsname{\hyperlink{tag-00VJ}{II.\ref*{cdp:eq:C.PR.11}}}
\expandafter\def\csname SRef@cdp:eq:C.PR.12\endcsname{\hyperlink{tag-00VK}{II.\ref*{cdp:eq:C.PR.12}}}
\expandafter\def\csname SRef@cdp:eq:C.CQ.10\endcsname{\hyperlink{tag-00CK}{II.\ref*{cdp:eq:C.CQ.10}}}
\expandafter\def\csname SRef@cdp:eq:C.PR.13\endcsname{\hyperlink{tag-00VL}{II.\ref*{cdp:eq:C.PR.13}}}
\expandafter\def\csname SRef@cdp:eq:C.PR.14\endcsname{\hyperlink{tag-00VM}{II.\ref*{cdp:eq:C.PR.14}}}
\expandafter\def\csname SRef@cdp:eq:C.PR.15\endcsname{\hyperlink{tag-00VN}{II.\ref*{cdp:eq:C.PR.15}}}
\expandafter\def\csname SRef@cdp:eq:C.PR.16\endcsname{\hyperlink{tag-00VO}{II.\ref*{cdp:eq:C.PR.16}}}
\expandafter\def\csname SRef@cdp:eq:C.PR.17\endcsname{\hyperlink{tag-00VP}{II.\ref*{cdp:eq:C.PR.17}}}
\expandafter\def\csname SRef@cdp:eq:C.PR.18\endcsname{\hyperlink{tag-00VQ}{II.\ref*{cdp:eq:C.PR.18}}}
\expandafter\def\csname SRef@cdp:eq:C.PT.1\endcsname{\hyperlink{tag-00VZ}{II.\ref*{cdp:eq:C.PT.1}}}
\expandafter\def\csname SRef@cdp:eq:C.PT.2\endcsname{\hyperlink{tag-00W5}{II.\ref*{cdp:eq:C.PT.2}}}
\expandafter\def\csname SRef@cdp:eq:C.PT.3\endcsname{\hyperlink{tag-00W6}{II.\ref*{cdp:eq:C.PT.3}}}
\expandafter\def\csname SRef@cdp:eq:C.PT.4\endcsname{\hyperlink{tag-00W7}{II.\ref*{cdp:eq:C.PT.4}}}
\expandafter\def\csname SRef@cdp:eq:C.CQ.11\endcsname{\hyperlink{tag-00CL}{II.\ref*{cdp:eq:C.CQ.11}}}
\expandafter\def\csname SRef@cdp:eq:C.PT.5\endcsname{\hyperlink{tag-00W8}{II.\ref*{cdp:eq:C.PT.5}}}
\expandafter\def\csname SRef@cdp:eq:C.PT.6\endcsname{\hyperlink{tag-00W9}{II.\ref*{cdp:eq:C.PT.6}}}
\expandafter\def\csname SRef@cdp:eq:C.PT.7\endcsname{\hyperlink{tag-00WA}{II.\ref*{cdp:eq:C.PT.7}}}
\expandafter\def\csname SRef@cdp:eq:C.PT.8\endcsname{\hyperlink{tag-00WB}{II.\ref*{cdp:eq:C.PT.8}}}
\expandafter\def\csname SRef@cdp:eq:C.PT.9\endcsname{\hyperlink{tag-00WC}{II.\ref*{cdp:eq:C.PT.9}}}
\expandafter\def\csname SRef@cdp:eq:C.PT.10\endcsname{\hyperlink{tag-00W0}{II.\ref*{cdp:eq:C.PT.10}}}
\expandafter\def\csname SRef@cdp:eq:C.PT.11\endcsname{\hyperlink{tag-00W1}{II.\ref*{cdp:eq:C.PT.11}}}
\expandafter\def\csname SRef@cdp:eq:C.PT.12\endcsname{\hyperlink{tag-00W2}{II.\ref*{cdp:eq:C.PT.12}}}
\expandafter\def\csname SRef@cdp:eq:C.PT.13\endcsname{\hyperlink{tag-00W3}{II.\ref*{cdp:eq:C.PT.13}}}
\expandafter\def\csname SRef@cdp:eq:C.PT.14\endcsname{\hyperlink{tag-00W4}{II.\ref*{cdp:eq:C.PT.14}}}
\expandafter\def\csname SRef@cdp:eq:C.CQ.12\endcsname{\hyperlink{tag-00CM}{II.\ref*{cdp:eq:C.CQ.12}}}
\expandafter\def\csname SRef@cdp:eq:C.QT.1\endcsname{\hyperlink{tag-00WD}{II.\ref*{cdp:eq:C.QT.1}}}
\expandafter\def\csname SRef@cdp:eq:C.QT.2\endcsname{\hyperlink{tag-00WI}{II.\ref*{cdp:eq:C.QT.2}}}
\expandafter\def\csname SRef@cdp:eq:C.QT.3\endcsname{\hyperlink{tag-00WJ}{II.\ref*{cdp:eq:C.QT.3}}}
\expandafter\def\csname SRef@cdp:eq:C.QT.4\endcsname{\hyperlink{tag-00WK}{II.\ref*{cdp:eq:C.QT.4}}}
\expandafter\def\csname SRef@cdp:eq:C.QT.5\endcsname{\hyperlink{tag-00WL}{II.\ref*{cdp:eq:C.QT.5}}}
\expandafter\def\csname SRef@cdp:eq:C.QT.6\endcsname{\hyperlink{tag-00WM}{II.\ref*{cdp:eq:C.QT.6}}}
\expandafter\def\csname SRef@cdp:eq:C.QT.7\endcsname{\hyperlink{tag-00WN}{II.\ref*{cdp:eq:C.QT.7}}}
\expandafter\def\csname SRef@cdp:eq:C.QT.8\endcsname{\hyperlink{tag-00WO}{II.\ref*{cdp:eq:C.QT.8}}}
\expandafter\def\csname SRef@cdp:eq:C.QT.9\endcsname{\hyperlink{tag-00WP}{II.\ref*{cdp:eq:C.QT.9}}}
\expandafter\def\csname SRef@cdp:eq:C.QT.10\endcsname{\hyperlink{tag-00WE}{II.\ref*{cdp:eq:C.QT.10}}}
\expandafter\def\csname SRef@cdp:st:C.2.UN.1\endcsname{\hyperlink{tag-01CM}{II.\ref*{cdp:st:C.2.UN.1}}}
\expandafter\def\csname SRef@cdp:st:C.2.UN.20\endcsname{\hyperlink{tag-01CY}{II.\ref*{cdp:st:C.2.UN.20}}}
\expandafter\def\csname SRef@cdp:st:C.2.UN.22\endcsname{\hyperlink{tag-01D0}{II.\ref*{cdp:st:C.2.UN.22}}}
\expandafter\def\csname SRef@cdp:st:C.2.UN.23\endcsname{\hyperlink{tag-01D1}{II.\ref*{cdp:st:C.2.UN.23}}}
\expandafter\def\csname SRef@cdp:st:C.2.UN.26\endcsname{\hyperlink{tag-01D4}{II.\ref*{cdp:st:C.2.UN.26}}}
\expandafter\def\csname SRef@cdp:st:C.2.UN.28\endcsname{\hyperlink{tag-01D6}{II.\ref*{cdp:st:C.2.UN.28}}}
\expandafter\def\csname SRef@cdp:st:C.2.UN.29\endcsname{\hyperlink{tag-01D7}{II.\ref*{cdp:st:C.2.UN.29}}}
\expandafter\def\csname SRef@cdp:st:C.2.UN.31\endcsname{\hyperlink{tag-01DA}{II.\ref*{cdp:st:C.2.UN.31}}}
\expandafter\def\csname SRef@cdp:st:C.2.UN.33\endcsname{\hyperlink{tag-01DC}{II.\ref*{cdp:st:C.2.UN.33}}}
\expandafter\def\csname SRef@cdp:st:C.2.UN.36\endcsname{\hyperlink{tag-01DF}{II.\ref*{cdp:st:C.2.UN.36}}}
\expandafter\def\csname SRef@cdp:st:C.2.UN.38\endcsname{\hyperlink{tag-01DH}{II.\ref*{cdp:st:C.2.UN.38}}}
\expandafter\def\csname SRef@cdp:st:C.2.UN.2\endcsname{\hyperlink{tag-01CX}{II.\ref*{cdp:st:C.2.UN.2}}}
\expandafter\def\csname SRef@cdp:st:C.2.UN.39\endcsname{\hyperlink{tag-01DI}{II.\ref*{cdp:st:C.2.UN.39}}}
\expandafter\def\csname SRef@cdp:st:C.2.UN.44\endcsname{\hyperlink{tag-01DO}{II.\ref*{cdp:st:C.2.UN.44}}}
\expandafter\def\csname SRef@cdp:st:C.2.UN.47\endcsname{\hyperlink{tag-01DR}{II.\ref*{cdp:st:C.2.UN.47}}}
\expandafter\def\csname SRef@cdp:st:C.2.UN.3\endcsname{\hyperlink{tag-01D8}{II.\ref*{cdp:st:C.2.UN.3}}}
\expandafter\def\csname SRef@cdp:st:C.2.UN.4\endcsname{\hyperlink{tag-01DJ}{II.\ref*{cdp:st:C.2.UN.4}}}
\expandafter\def\csname SRef@cdp:st:C.2.UN.14\endcsname{\hyperlink{tag-01CR}{II.\ref*{cdp:st:C.2.UN.14}}}
\expandafter\def\csname SRef@cdp:st:C.2.UN.15\endcsname{\hyperlink{tag-01CS}{II.\ref*{cdp:st:C.2.UN.15}}}
\expandafter\def\csname SRef@cdp:st:C.2.UN.16\endcsname{\hyperlink{tag-01CT}{II.\ref*{cdp:st:C.2.UN.16}}}
\expandafter\def\csname SRef@cdp:st:C.2.UN.17\endcsname{\hyperlink{tag-01CU}{II.\ref*{cdp:st:C.2.UN.17}}}
\expandafter\def\csname SRef@cdp:st:C.2.UN.19\endcsname{\hyperlink{tag-01CW}{II.\ref*{cdp:st:C.2.UN.19}}}
\expandafter\def\csname SRef@cdp:st:C.2.UN.5\endcsname{\hyperlink{tag-01DU}{II.\ref*{cdp:st:C.2.UN.5}}}
\expandafter\def\csname SRef@cdp:st:C.2.IRR.1\endcsname{\hyperlink{tag-01C3}{II.\ref*{cdp:st:C.2.IRR.1}}}
\expandafter\def\csname SRef@cdp:st:C.2.VII.1\endcsname{\hyperlink{tag-01F0}{II.\ref*{cdp:st:C.2.VII.1}}}
\expandafter\def\csname SRef@cdp:st:C.2.VII.2\endcsname{\hyperlink{tag-01F1}{II.\ref*{cdp:st:C.2.VII.2}}}
\expandafter\def\csname SRef@cdp:st:C.2.VII.3\endcsname{\hyperlink{tag-01F2}{II.\ref*{cdp:st:C.2.VII.3}}}
\expandafter\def\csname SRef@cdp:st:C.2.VII.4\endcsname{\hyperlink{tag-01F3}{II.\ref*{cdp:st:C.2.VII.4}}}
\expandafter\def\csname SRef@cdp:st:C.2.RAT.1\endcsname{\hyperlink{tag-01CD}{II.\ref*{cdp:st:C.2.RAT.1}}}
\expandafter\def\csname SRef@cdp:st:C.2.RAT.2\endcsname{\hyperlink{tag-01CE}{II.\ref*{cdp:st:C.2.RAT.2}}}
\expandafter\def\csname SRef@cdp:st:C.2.RAT.3\endcsname{\hyperlink{tag-01CF}{II.\ref*{cdp:st:C.2.RAT.3}}}
\expandafter\def\csname SRef@cdp:st:C.2.RAT.4\endcsname{\hyperlink{tag-01CG}{II.\ref*{cdp:st:C.2.RAT.4}}}
\expandafter\def\csname SRef@cdp:st:C.2.RAT.5\endcsname{\hyperlink{tag-01CH}{II.\ref*{cdp:st:C.2.RAT.5}}}
\expandafter\def\csname SRef@cdp:st:C.2.CR.1\endcsname{\hyperlink{tag-01BQ}{II.\ref*{cdp:st:C.2.CR.1}}}
\expandafter\def\csname SRef@cdp:st:C.2.RAT.7\endcsname{\hyperlink{tag-01CJ}{II.\ref*{cdp:st:C.2.RAT.7}}}
\expandafter\def\csname SRef@cdp:st:C.2.RAT.8\endcsname{\hyperlink{tag-01CK}{II.\ref*{cdp:st:C.2.RAT.8}}}
\expandafter\def\csname SRef@cdp:st:C.2.RAT.9\endcsname{\hyperlink{tag-01CL}{II.\ref*{cdp:st:C.2.RAT.9}}}
\expandafter\def\csname SRef@cdp:st:C.2.D4.1\endcsname{\hyperlink{tag-01BW}{II.\ref*{cdp:st:C.2.D4.1}}}
\expandafter\def\csname SRef@cdp:st:C.2.D4.2\endcsname{\hyperlink{tag-01BX}{II.\ref*{cdp:st:C.2.D4.2}}}
\expandafter\def\csname SRef@cdp:st:C.2.D4.3\endcsname{\hyperlink{tag-01BY}{II.\ref*{cdp:st:C.2.D4.3}}}
\expandafter\def\csname SRef@cdp:st:C.2.UN.7\endcsname{\hyperlink{tag-01EG}{II.\ref*{cdp:st:C.2.UN.7}}}
\expandafter\def\csname SRef@cdp:st:C.2.UN.12\endcsname{\hyperlink{tag-01CP}{II.\ref*{cdp:st:C.2.UN.12}}}
\expandafter\def\csname SRef@cdp:st:C.2.INT.1\endcsname{\hyperlink{tag-01C0}{II.\ref*{cdp:st:C.2.INT.1}}}
\expandafter\def\csname SRef@cdp:st:C.2.INT.2\endcsname{\hyperlink{tag-01C1}{II.\ref*{cdp:st:C.2.INT.2}}}
\expandafter\def\csname SRef@cdp:st:C.2.CR.2\endcsname{\hyperlink{tag-01BR}{II.\ref*{cdp:st:C.2.CR.2}}}
\expandafter\def\csname SRef@cdp:st:C.2.INT.3\endcsname{\hyperlink{tag-01C2}{II.\ref*{cdp:st:C.2.INT.3}}}
\expandafter\def\csname SRef@cdp:st:C.2.MH.1\endcsname{\hyperlink{tag-01C4}{II.\ref*{cdp:st:C.2.MH.1}}}
\expandafter\def\csname SRef@cdp:st:C.2.MH.2\endcsname{\hyperlink{tag-01C5}{II.\ref*{cdp:st:C.2.MH.2}}}
\expandafter\def\csname SRef@cdp:st:C.2.MH.3\endcsname{\hyperlink{tag-01C6}{II.\ref*{cdp:st:C.2.MH.3}}}
\expandafter\def\csname SRef@cdp:st:C.2.MH.4\endcsname{\hyperlink{tag-01C7}{II.\ref*{cdp:st:C.2.MH.4}}}
\expandafter\def\csname SRef@cdp:st:C.2.CR.3\endcsname{\hyperlink{tag-01BS}{II.\ref*{cdp:st:C.2.CR.3}}}
\expandafter\def\csname SRef@cdp:st:C.2.CR.4\endcsname{\hyperlink{tag-01BT}{II.\ref*{cdp:st:C.2.CR.4}}}
\expandafter\def\csname SRef@cdp:st:C.2.CR.5\endcsname{\hyperlink{tag-01BU}{II.\ref*{cdp:st:C.2.CR.5}}}
\expandafter\def\csname SRef@cdp:st:C.2.CR.6\endcsname{\hyperlink{tag-01BV}{II.\ref*{cdp:st:C.2.CR.6}}}
\expandafter\def\csname SRef@cdp:st:C.2.UN.9\endcsname{\hyperlink{tag-01EZ}{II.\ref*{cdp:st:C.2.UN.9}}}
\expandafter\def\csname SRef@cdp:st:C.2.UN.10\endcsname{\hyperlink{tag-01CN}{II.\ref*{cdp:st:C.2.UN.10}}}
\expandafter\def\csname SRef@cdp:st:C2a\endcsname{\hyperlink{tag-01FG}{II.\ref*{cdp:st:C2a}}}
\expandafter\def\csname SRef@cdp:st:C.2.UN.67\endcsname{\hyperlink{tag-01ED}{II.\ref*{cdp:st:C.2.UN.67}}}
\expandafter\def\csname SRef@cdp:st:C.2.UN.68\endcsname{\hyperlink{tag-01EE}{II.\ref*{cdp:st:C.2.UN.68}}}
\expandafter\def\csname SRef@cdp:st:C2b\endcsname{\hyperlink{tag-01FH}{II.\ref*{cdp:st:C2b}}}
\expandafter\def\csname SRef@cdp:st:C2c\endcsname{\hyperlink{tag-01FI}{II.\ref*{cdp:st:C2c}}}
\expandafter\def\csname SRef@cdp:st:C.2.UN.11\endcsname{\hyperlink{tag-01CO}{II.\ref*{cdp:st:C.2.UN.11}}}
\expandafter\def\csname SRef@cdp:st:C.2.UN.24\endcsname{\hyperlink{tag-01D2}{II.\ref*{cdp:st:C.2.UN.24}}}
\expandafter\def\csname SRef@cdp:st:C.2.UN.32\endcsname{\hyperlink{tag-01DB}{II.\ref*{cdp:st:C.2.UN.32}}}
\expandafter\def\csname SRef@cdp:st:C.2.UN.64\endcsname{\hyperlink{tag-01EA}{II.\ref*{cdp:st:C.2.UN.64}}}
\expandafter\def\csname SRef@cdp:st:C.2.UN.65\endcsname{\hyperlink{tag-01EB}{II.\ref*{cdp:st:C.2.UN.65}}}
\expandafter\def\csname SRef@cdp:st:C.2.UN.66\endcsname{\hyperlink{tag-01EC}{II.\ref*{cdp:st:C.2.UN.66}}}
\expandafter\def\csname SRef@cdp:st:C.2.UN.6\endcsname{\hyperlink{tag-01E5}{II.\ref*{cdp:st:C.2.UN.6}}}
\expandafter\def\csname SRef@cdp:st:C.2.UN.35\endcsname{\hyperlink{tag-01DE}{II.\ref*{cdp:st:C.2.UN.35}}}
\expandafter\def\csname SRef@cdp:st:C.2.UN.37\endcsname{\hyperlink{tag-01DG}{II.\ref*{cdp:st:C.2.UN.37}}}
\expandafter\def\csname SRef@cdp:st:C.2.UN.41\endcsname{\hyperlink{tag-01DL}{II.\ref*{cdp:st:C.2.UN.41}}}
\expandafter\def\csname SRef@cdp:st:C.2.UN.42\endcsname{\hyperlink{tag-01DM}{II.\ref*{cdp:st:C.2.UN.42}}}
\expandafter\def\csname SRef@cdp:st:C.2.UN.43\endcsname{\hyperlink{tag-01DN}{II.\ref*{cdp:st:C.2.UN.43}}}
\expandafter\def\csname SRef@cdp:st:C.2.UN.45\endcsname{\hyperlink{tag-01DP}{II.\ref*{cdp:st:C.2.UN.45}}}
\expandafter\def\csname SRef@cdp:st:C.2.UN.8\endcsname{\hyperlink{tag-01ER}{II.\ref*{cdp:st:C.2.UN.8}}}
\expandafter\def\csname SRef@cdp:st:C.2.UN.13\endcsname{\hyperlink{tag-01CQ}{II.\ref*{cdp:st:C.2.UN.13}}}
\expandafter\def\csname SRef@cdp:st:C.2.UN.18\endcsname{\hyperlink{tag-01CV}{II.\ref*{cdp:st:C.2.UN.18}}}
\expandafter\def\csname SRef@cdp:st:C.2.UN.21\endcsname{\hyperlink{tag-01CZ}{II.\ref*{cdp:st:C.2.UN.21}}}
\expandafter\def\csname SRef@cdp:st:C.2.UN.25\endcsname{\hyperlink{tag-01D3}{II.\ref*{cdp:st:C.2.UN.25}}}
\expandafter\def\csname SRef@cdp:st:C.2.UN.27\endcsname{\hyperlink{tag-01D5}{II.\ref*{cdp:st:C.2.UN.27}}}
\expandafter\def\csname SRef@cdp:st:C.2.UN.30\endcsname{\hyperlink{tag-01D9}{II.\ref*{cdp:st:C.2.UN.30}}}
\expandafter\def\csname SRef@cdp:st:C.2.UN.34\endcsname{\hyperlink{tag-01DD}{II.\ref*{cdp:st:C.2.UN.34}}}
\expandafter\def\csname SRef@cdp:st:C.2.UN.40\endcsname{\hyperlink{tag-01DK}{II.\ref*{cdp:st:C.2.UN.40}}}
\expandafter\def\csname SRef@cdp:st:C.2.UN.55\endcsname{\hyperlink{tag-01E0}{II.\ref*{cdp:st:C.2.UN.55}}}
\expandafter\def\csname SRef@cdp:st:C.2.UN.56\endcsname{\hyperlink{tag-01E1}{II.\ref*{cdp:st:C.2.UN.56}}}
\expandafter\def\csname SRef@cdp:st:C.2.UN.57\endcsname{\hyperlink{tag-01E2}{II.\ref*{cdp:st:C.2.UN.57}}}
\expandafter\def\csname SRef@cdp:st:C.2.UN.58\endcsname{\hyperlink{tag-01E3}{II.\ref*{cdp:st:C.2.UN.58}}}
\expandafter\def\csname SRef@cdp:st:C.2.UN.59\endcsname{\hyperlink{tag-01E4}{II.\ref*{cdp:st:C.2.UN.59}}}
\expandafter\def\csname SRef@cdp:st:C.2.UN.60\endcsname{\hyperlink{tag-01E6}{II.\ref*{cdp:st:C.2.UN.60}}}
\expandafter\def\csname SRef@cdp:st:C.2.UN.61\endcsname{\hyperlink{tag-01E7}{II.\ref*{cdp:st:C.2.UN.61}}}
\expandafter\def\csname SRef@cdp:st:C.2.UN.62\endcsname{\hyperlink{tag-01E8}{II.\ref*{cdp:st:C.2.UN.62}}}
\expandafter\def\csname SRef@cdp:st:C.2.UN.63\endcsname{\hyperlink{tag-01E9}{II.\ref*{cdp:st:C.2.UN.63}}}
\expandafter\def\csname SRef@cdp:st:C.2.UN.46\endcsname{\hyperlink{tag-01DQ}{II.\ref*{cdp:st:C.2.UN.46}}}
\expandafter\def\csname SRef@cdp:st:C.2.UN.48\endcsname{\hyperlink{tag-01DS}{II.\ref*{cdp:st:C.2.UN.48}}}
\expandafter\def\csname SRef@cdp:st:C.2.UN.49\endcsname{\hyperlink{tag-01DT}{II.\ref*{cdp:st:C.2.UN.49}}}
\expandafter\def\csname SRef@cdp:st:C.2.UN.50\endcsname{\hyperlink{tag-01DV}{II.\ref*{cdp:st:C.2.UN.50}}}
\expandafter\def\csname SRef@cdp:st:C.2.UN.51\endcsname{\hyperlink{tag-01DW}{II.\ref*{cdp:st:C.2.UN.51}}}
\expandafter\def\csname SRef@cdp:st:C.2.UN.52\endcsname{\hyperlink{tag-01DX}{II.\ref*{cdp:st:C.2.UN.52}}}
\expandafter\def\csname SRef@cdp:st:C.2.UN.53\endcsname{\hyperlink{tag-01DY}{II.\ref*{cdp:st:C.2.UN.53}}}
\expandafter\def\csname SRef@cdp:st:C.2.UN.54\endcsname{\hyperlink{tag-01DZ}{II.\ref*{cdp:st:C.2.UN.54}}}
\expandafter\def\csname SRef@cdp:st:C.2.UN.69\endcsname{\hyperlink{tag-01EF}{II.\ref*{cdp:st:C.2.UN.69}}}
\expandafter\def\csname SRef@cdp:st:C.2.UN.80\endcsname{\hyperlink{tag-01ES}{II.\ref*{cdp:st:C.2.UN.80}}}
\expandafter\def\csname SRef@cdp:st:C.2.UN.81\endcsname{\hyperlink{tag-01ET}{II.\ref*{cdp:st:C.2.UN.81}}}
\expandafter\def\csname SRef@cdp:st:C.2.UN.82\endcsname{\hyperlink{tag-01EU}{II.\ref*{cdp:st:C.2.UN.82}}}
\expandafter\def\csname SRef@cdp:st:C.2.UN.83\endcsname{\hyperlink{tag-01EV}{II.\ref*{cdp:st:C.2.UN.83}}}
\expandafter\def\csname SRef@cdp:st:C.2.UN.84\endcsname{\hyperlink{tag-01EW}{II.\ref*{cdp:st:C.2.UN.84}}}
\expandafter\def\csname SRef@cdp:st:C.2.UN.78\endcsname{\hyperlink{tag-01EP}{II.\ref*{cdp:st:C.2.UN.78}}}
\expandafter\def\csname SRef@cdp:st:C.2.UN.72\endcsname{\hyperlink{tag-01EJ}{II.\ref*{cdp:st:C.2.UN.72}}}
\expandafter\def\csname SRef@cdp:st:C.2.UN.85\endcsname{\hyperlink{tag-01EX}{II.\ref*{cdp:st:C.2.UN.85}}}
\expandafter\def\csname SRef@cdp:st:C.2.UN.86\endcsname{\hyperlink{tag-01EY}{II.\ref*{cdp:st:C.2.UN.86}}}
\expandafter\def\csname SRef@cdp:st:C.2.UN.70\endcsname{\hyperlink{tag-01EH}{II.\ref*{cdp:st:C.2.UN.70}}}
\expandafter\def\csname SRef@cdp:st:C.2.UN.71\endcsname{\hyperlink{tag-01EI}{II.\ref*{cdp:st:C.2.UN.71}}}
\expandafter\def\csname SRef@cdp:st:C.2.UN.73\endcsname{\hyperlink{tag-01EK}{II.\ref*{cdp:st:C.2.UN.73}}}
\expandafter\def\csname SRef@cdp:st:C.2.UN.76\endcsname{\hyperlink{tag-01EN}{II.\ref*{cdp:st:C.2.UN.76}}}
\expandafter\def\csname SRef@cdp:st:C.2.UN.74\endcsname{\hyperlink{tag-01EL}{II.\ref*{cdp:st:C.2.UN.74}}}
\expandafter\def\csname SRef@cdp:st:C.2.UN.75\endcsname{\hyperlink{tag-01EM}{II.\ref*{cdp:st:C.2.UN.75}}}
\expandafter\def\csname SRef@cdp:st:C.2.UN.77\endcsname{\hyperlink{tag-01EO}{II.\ref*{cdp:st:C.2.UN.77}}}
\expandafter\def\csname SRef@cdp:st:C.2.UN.79\endcsname{\hyperlink{tag-01EQ}{II.\ref*{cdp:st:C.2.UN.79}}}
\expandafter\def\csname SRef@sec:II:06\endcsname{\hyperlink{tag-0342}{II.\ref*{sec:II:06}}}
\expandafter\def\csname SRef@sec:II:07\endcsname{\hyperlink{tag-0343}{II.\ref*{sec:II:07}}}
\expandafter\def\csname SRef@sec:II:all:08\endcsname{\hyperlink{tag-034B}{II.\ref*{sec:II:all:08}}}
\expandafter\def\csname SRef@sec:II:09\endcsname{\hyperlink{tag-0345}{II.\ref*{sec:II:09}}}
\expandafter\def\csname SRef@sec:II:10\endcsname{\hyperlink{tag-0346}{II.\ref*{sec:II:10}}}
\expandafter\def\csname SRef@sec:II:11\endcsname{\hyperlink{tag-0347}{II.\ref*{sec:II:11}}}
\expandafter\def\csname SRef@sec:II:all:12\endcsname{\hyperlink{tag-034C}{II.\ref*{sec:II:all:12}}}
\expandafter\def\csname SRef@cdp:st:C.2.G.1\endcsname{\hyperlink{tag-01BZ}{II.\ref*{cdp:st:C.2.G.1}}}
\expandafter\def\csname SRef@inv:E1.1\endcsname{\hyperlink{tag-01ST}{II.\ref*{inv:E1.1}}}
\expandafter\def\csname SRef@inv:E2.1\endcsname{\hyperlink{tag-01SU}{II.\ref*{inv:E2.1}}}
\expandafter\def\csname SRef@push:src:C.1.1\endcsname{\href{S6_SERIES_GRIND_III.pdf\#nameddest=tag-031B}{III.? (p.~?)} }
\expandafter\def\csname SRef@push:src:C.1.2\endcsname{\href{S6_SERIES_GRIND_III.pdf\#nameddest=tag-031C}{III.? (p.~?)} }
\expandafter\def\csname SRef@push:src:C.2.1\endcsname{\href{S6_SERIES_GRIND_III.pdf\#nameddest=tag-031D}{III.? (p.~?)} }
\expandafter\def\csname SRef@push:src:C.3.1\endcsname{\href{S6_SERIES_GRIND_III.pdf\#nameddest=tag-031G}{III.? (p.~?)} }
\expandafter\def\csname SRef@push:src:C.3.3\endcsname{\href{S6_SERIES_GRIND_III.pdf\#nameddest=tag-0320}{III.? (p.~?)} }
\expandafter\def\csname SRef@push:src:E.3.2\endcsname{\href{S6_SERIES_GRIND_I.pdf\#nameddest=tag-033J}{I.? (p.~?)} }
\expandafter\def\csname SRef@push:src:E.3.3\endcsname{\href{S6_SERIES_GRIND_I.pdf\#nameddest=tag-033K}{I.? (p.~?)} }
\expandafter\def\csname SRef@cdp:st:A.2.3\endcsname{\hyperlink{tag-01AH}{II.\ref*{cdp:st:A.2.3}}}
\expandafter\def\csname SRef@cdp:st:C.1.2\endcsname{\hyperlink{tag-01BO}{II.\ref*{cdp:st:C.1.2}}}
\expandafter\def\csname SRef@cdp:st:C.2.NC.1\endcsname{\hyperlink{tag-01C8}{II.\ref*{cdp:st:C.2.NC.1}}}
\expandafter\def\csname SRef@cdp:st:C.2.RAT.6\endcsname{\hyperlink{tag-01CI}{II.\ref*{cdp:st:C.2.RAT.6}}}
\expandafter\def\csname SRef@cdp:st:D.1.1\endcsname{\hyperlink{tag-01FU}{II.\ref*{cdp:st:D.1.1}}}
\expandafter\def\csname SRef@cdp:st:D.1.3\endcsname{\hyperlink{tag-01G1}{II.\ref*{cdp:st:D.1.3}}}
\expandafter\def\csname SRef@cdp:st:E2\endcsname{\hyperlink{tag-01GR}{II.\ref*{cdp:st:E2}}}
\expandafter\def\csname SRef@cdp:st:E3\endcsname{\hyperlink{tag-01GS}{II.\ref*{cdp:st:E3}}}
\expandafter\def\csname SRef@cdp:st:E4\endcsname{\hyperlink{tag-01GT}{II.\ref*{cdp:st:E4}}}
\expandafter\def\csname SRef@cdp:st:E19\endcsname{\hyperlink{tag-01GQ}{II.\ref*{cdp:st:E19}}}
\expandafter\def\csname SRef@maps:src:B.6.9\endcsname{\href{S6_SERIES_GRIND_IV.pdf\#nameddest=tag-02K1}{IV.? (p.~?)} }
\expandafter\def\csname SRef@maps:src:D.2.1\endcsname{\href{S6_SERIES_GRIND_III.pdf\#nameddest=tag-02KJ}{III.? (p.~?)} }
\expandafter\def\csname SRef@maps:src:F.1.2\endcsname{\href{S6_SERIES_GRIND_V.pdf\#nameddest=tag-02MV}{V.? (p.~?)} }
\expandafter\def\csname SRef@maps:src:F.1.3\endcsname{\href{S6_SERIES_GRIND_V.pdf\#nameddest=tag-02MW}{V.? (p.~?)} }
\expandafter\def\csname SRef@ari:C.7.5\endcsname{\href{S6_SERIES_GRIND_I.pdf\#nameddest=tag-001E}{I.? (p.~?)} }
\expandafter\def\csname SRef@wb:lem:triangle-generator-orientations\endcsname{\href{S6_SERIES_GRIND_I.pdf\#nameddest=tag-03BM}{I.? (p.~?)} }
\expandafter\def\csname SRef@wb:prop:period-input-audit\endcsname{\href{S6_SERIES_GRIND_I.pdf\#nameddest=tag-03BS}{I.? (p.~?)} }
\expandafter\def\csname SRef@wb:prop:period-exact-admissible-constants\endcsname{\href{S6_SERIES_GRIND_I.pdf\#nameddest=tag-03BR}{I.? (p.~?)} }
\expandafter\def\csname SRef@wb:lem:toric-chart-control\endcsname{\href{S6_SERIES_GRIND_I.pdf\#nameddest=tag-03BK}{I.? (p.~?)} }
\expandafter\def\csname SRef@wb:thm:candidate-cusp-geometry\endcsname{\href{S6_SERIES_GRIND_I.pdf\#nameddest=tag-03DK}{I.? (p.~?)} }
\expandafter\def\csname SRef@wb:prop:candidate-cusp-overlap\endcsname{\href{S6_SERIES_GRIND_I.pdf\#nameddest=tag-03BN}{I.? (p.~?)} }
\expandafter\def\csname SRef@wb:lem:middle-period-proper\endcsname{\href{S6_SERIES_GRIND_I.pdf\#nameddest=tag-03BI}{I.? (p.~?)} }
\expandafter\def\csname SRef@wb:prop:candidate-finite-fillings\endcsname{\href{S6_SERIES_GRIND_I.pdf\#nameddest=tag-03BO}{I.? (p.~?)} }
\expandafter\def\csname SRef@wb:lem:finite-filling-stalk-map\endcsname{\href{S6_SERIES_GRIND_I.pdf\#nameddest=tag-03BG}{I.? (p.~?)} }
\expandafter\def\csname SRef@wb:thm:candidate-global-geometry\endcsname{\href{S6_SERIES_GRIND_I.pdf\#nameddest=tag-03DL}{I.? (p.~?)} }
\expandafter\def\csname SRef@wb:prop:monodromy-certificate\endcsname{\href{S6_SERIES_GRIND_I.pdf\#nameddest=tag-03BQ}{I.? (p.~?)} }
\expandafter\def\csname SRef@wb:thm:cusp-lattices\endcsname{\href{S6_SERIES_GRIND_I.pdf\#nameddest=tag-03DO}{I.? (p.~?)} }
\expandafter\def\csname SRef@wb:thm:coinvariants\endcsname{\href{S6_SERIES_GRIND_I.pdf\#nameddest=tag-03DN}{I.? (p.~?)} }
\expandafter\def\csname SRef@wb:prop:finite-indices\endcsname{\href{S6_SERIES_GRIND_I.pdf\#nameddest=tag-03BP}{I.? (p.~?)} }
\expandafter\def\csname SRef@wb:prop:terminal-group\endcsname{\href{S6_SERIES_GRIND_I.pdf\#nameddest=tag-03BV}{I.? (p.~?)} }
\expandafter\def\csname SRef@wb:prop:terminal-degree-two\endcsname{\href{S6_SERIES_GRIND_I.pdf\#nameddest=tag-03BU}{I.? (p.~?)} }
\expandafter\def\csname SRef@wb:prop:terminal-all-degree\endcsname{\href{S6_SERIES_GRIND_I.pdf\#nameddest=tag-03BT}{I.? (p.~?)} }
\expandafter\def\csname SRef@wb:ffcert:prop:integral-summands\endcsname{\href{S6_SERIES_GRIND_I.pdf\#nameddest=tag-03A8}{I.? (p.~?)} }
\expandafter\def\csname SRef@wb:ffcert:thm:product-covers\endcsname{\href{S6_SERIES_GRIND_I.pdf\#nameddest=tag-03AD}{I.? (p.~?)} }
\expandafter\def\csname SRef@wb:ffcert:prop:quotient-charts\endcsname{\href{S6_SERIES_GRIND_I.pdf\#nameddest=tag-03A9}{I.? (p.~?)} }
\expandafter\def\csname SRef@wb:ffcert:thm:normalization\endcsname{\href{S6_SERIES_GRIND_I.pdf\#nameddest=tag-03AB}{I.? (p.~?)} }
\expandafter\def\csname SRef@wb:ffcert:thm:peripheral-kernel\endcsname{\href{S6_SERIES_GRIND_I.pdf\#nameddest=tag-03AC}{I.? (p.~?)} }
\expandafter\def\csname SRef@wb:ffcert:cor:covering-h1\endcsname{\href{S6_SERIES_GRIND_I.pdf\#nameddest=tag-03A7}{I.? (p.~?)} }
\expandafter\def\csname SRef@wb:ffcert:thm:normal-trivialization\endcsname{\href{S6_SERIES_GRIND_I.pdf\#nameddest=tag-03AA}{I.? (p.~?)} }
\expandafter\def\csname SRef@wb:lem:torus-algebraic-reduction-audit\endcsname{\href{S6_SERIES_GRIND_I.pdf\#nameddest=tag-03BL}{I.? (p.~?)} }
\expandafter\def\csname SRef@wb:prop:very-general-torus-audit\endcsname{\href{S6_SERIES_GRIND_I.pdf\#nameddest=tag-03BW}{I.? (p.~?)} }
\expandafter\def\csname SRef@wb:lem:meromorphic-descent-audit\endcsname{\href{S6_SERIES_GRIND_I.pdf\#nameddest=tag-03BH}{I.? (p.~?)} }
\expandafter\def\csname SRef@wb:thm:algebraic-dimension-audit\endcsname{\href{S6_SERIES_GRIND_I.pdf\#nameddest=tag-03DJ}{I.? (p.~?)} }
\expandafter\def\csname SRef@wb:lem:semidirect-abelianization\endcsname{\href{S6_SERIES_GRIND_I.pdf\#nameddest=tag-03BJ}{I.? (p.~?)} }
\expandafter\def\csname SRef@mapsI:top:thm:p0-integral-specialization\endcsname{\href{S6_SERIES_GRIND_I.pdf\#nameddest=tag-02Q5}{I.? (p.~?)} }
\expandafter\def\csname SRef@mapsI:top:prop:leray-e2\endcsname{\href{S6_SERIES_GRIND_I.pdf\#nameddest=tag-02PU}{I.? (p.~?)} }
\expandafter\def\csname SRef@mapsI:top:lem:first-transgression\endcsname{\href{S6_SERIES_GRIND_I.pdf\#nameddest=tag-02PN}{I.? (p.~?)} }
\expandafter\def\csname SRef@mapsI:top:lem:higher-transgressions\endcsname{\href{S6_SERIES_GRIND_I.pdf\#nameddest=tag-02PO}{I.? (p.~?)} }
\expandafter\def\csname SRef@mapsI:top:thm:h2-vanishing\endcsname{\href{S6_SERIES_GRIND_I.pdf\#nameddest=tag-02Q3}{I.? (p.~?)} }
\expandafter\def\csname SRef@mapsI:top:prop:geometric-pi1\endcsname{\href{S6_SERIES_GRIND_I.pdf\#nameddest=tag-02PT}{I.? (p.~?)} }
\expandafter\def\csname SRef@mapsI:top:cor:candidate-sphere\endcsname{\href{S6_SERIES_GRIND_I.pdf\#nameddest=tag-02PK}{I.? (p.~?)} }
\expandafter\def\csname SRef@mapsI:top:prop:topological-spine\endcsname{\href{S6_SERIES_GRIND_I.pdf\#nameddest=tag-02PV}{I.? (p.~?)} }
\expandafter\def\csname SRef@mapsI:top:lem:collapse-straightening\endcsname{\href{S6_SERIES_GRIND_I.pdf\#nameddest=tag-02PM}{I.? (p.~?)} }
\expandafter\def\csname SRef@mapsI:top:lem:positive-collapse\endcsname{\href{S6_SERIES_GRIND_I.pdf\#nameddest=tag-02PQ}{I.? (p.~?)} }
\expandafter\def\csname SRef@mapsI:top:lem:phase-spread-collapse\endcsname{\href{S6_SERIES_GRIND_I.pdf\#nameddest=tag-02PP}{I.? (p.~?)} }
\expandafter\def\csname SRef@mapsI:top:prop:verified-cellular-collapse\endcsname{\href{S6_SERIES_GRIND_I.pdf\#nameddest=tag-02PW}{I.? (p.~?)} }
\expandafter\def\csname SRef@mapsI:top:prop:verified-collapse-homology\endcsname{\href{S6_SERIES_GRIND_I.pdf\#nameddest=tag-02PX}{I.? (p.~?)} }
\expandafter\def\csname SRef@mapsI:top:thm:collapse-equals-specialization\endcsname{\href{S6_SERIES_GRIND_I.pdf\#nameddest=tag-02Q2}{I.? (p.~?)} }
\expandafter\def\csname SRef@mapsI:top:remark:L600\endcsname{\href{S6_SERIES_GRIND_I.pdf\#nameddest=tag-02Q0}{I.? (p.~?)} }
\expandafter\def\csname SRef@mapsI:top:record:vankampen\endcsname{\href{S6_SERIES_GRIND_I.pdf\#nameddest=tag-02PY}{I.? (p.~?)} }
\expandafter\def\csname SRef@mapsI:top:prop:actual-finite-band-maps\endcsname{\href{S6_SERIES_GRIND_I.pdf\#nameddest=tag-02PR}{I.? (p.~?)} }
\expandafter\def\csname SRef@mapsI:top:prop:cusp-mv-blocks\endcsname{\href{S6_SERIES_GRIND_I.pdf\#nameddest=tag-02PS}{I.? (p.~?)} }
\expandafter\def\csname SRef@mapsI:top:remark:L920\endcsname{\href{S6_SERIES_GRIND_I.pdf\#nameddest=tag-02Q1}{I.? (p.~?)} }
\expandafter\def\csname SRef@mapsI:top:thm:mv-leray-comparison\endcsname{\href{S6_SERIES_GRIND_I.pdf\#nameddest=tag-02Q4}{I.? (p.~?)} }
\expandafter\def\csname SRef@mapsI:top:remark:L1113\endcsname{\href{S6_SERIES_GRIND_I.pdf\#nameddest=tag-02PZ}{I.? (p.~?)} }
\expandafter\def\csname SRef@mapsI:top:lem:cohomological-closure\endcsname{\href{S6_SERIES_GRIND_I.pdf\#nameddest=tag-02PL}{I.? (p.~?)} }
\expandafter\def\csname SRef@mapsI:top:thm:sphere-recognition\endcsname{\href{S6_SERIES_GRIND_I.pdf\#nameddest=tag-02Q6}{I.? (p.~?)} }
\expandafter\def\csname SRef@mapsI:top:cor:abstract-s6\endcsname{\href{S6_SERIES_GRIND_I.pdf\#nameddest=tag-02PI}{I.? (p.~?)} }
\expandafter\def\csname SRef@mapsI:top:cor:candidate-mv-sphere\endcsname{\href{S6_SERIES_GRIND_I.pdf\#nameddest=tag-02PJ}{I.? (p.~?)} }
\expandafter\def\csname SRef@wb:cusp-cone-S12-almost-complex-obstruction\endcsname{\href{S6_SERIES_GRIND_III.pdf\#nameddest=tag-035S}{III.? (p.~?)} }
\expandafter\def\csname SRef@courseIII:almost-complex-structures-and-integrability:section-1\endcsname{\href{S6_SERIES_GRIND_III.pdf\#nameddest=tag-01I4}{III.? (p.~?)} }
\expandafter\def\csname SRef@courseIII:almost-complex-structures-and-integrability:section-2\endcsname{\href{S6_SERIES_GRIND_III.pdf\#nameddest=tag-01I7}{III.? (p.~?)} }
\expandafter\def\csname SRef@courseIII:almost-complex-structures-and-integrability:section-3\endcsname{\href{S6_SERIES_GRIND_III.pdf\#nameddest=tag-01I8}{III.? (p.~?)} }
\expandafter\def\csname SRef@courseIII:almost-complex-structures-and-integrability:section-4\endcsname{\href{S6_SERIES_GRIND_III.pdf\#nameddest=tag-01I9}{III.? (p.~?)} }
\expandafter\def\csname SRef@courseIII:almost-complex-structures-and-integrability:section-5\endcsname{\href{S6_SERIES_GRIND_III.pdf\#nameddest=tag-01IA}{III.? (p.~?)} }
\expandafter\def\csname SRef@courseIII:almost-complex-structures-and-integrability:section-6\endcsname{\href{S6_SERIES_GRIND_III.pdf\#nameddest=tag-01IB}{III.? (p.~?)} }
\expandafter\def\csname SRef@courseIII:almost-complex-structures-and-integrability:section-7\endcsname{\href{S6_SERIES_GRIND_III.pdf\#nameddest=tag-01IC}{III.? (p.~?)} }
\expandafter\def\csname SRef@courseIII:almost-complex-structures-and-integrability:section-8\endcsname{\href{S6_SERIES_GRIND_III.pdf\#nameddest=tag-01ID}{III.? (p.~?)} }
\expandafter\def\csname SRef@courseIII:almost-complex-structures-and-integrability:section-9\endcsname{\href{S6_SERIES_GRIND_III.pdf\#nameddest=tag-01IE}{III.? (p.~?)} }
\expandafter\def\csname SRef@courseIII:almost-complex-structures-and-integrability:section-10\endcsname{\href{S6_SERIES_GRIND_III.pdf\#nameddest=tag-01I5}{III.? (p.~?)} }
\expandafter\def\csname SRef@courseIII:almost-complex-structures-and-integrability:section-11\endcsname{\href{S6_SERIES_GRIND_III.pdf\#nameddest=tag-01I6}{III.? (p.~?)} }
\expandafter\def\csname SRef@mapsI:course:CG-S6-07-HOMOTOPY\endcsname{\href{S6_SERIES_GRIND_I.pdf\#nameddest=tag-02OR}{I.? (p.~?)} }
\expandafter\def\csname SRef@courseIII:integral-homology-and-sphere-recognition:section\endcsname{\href{S6_SERIES_GRIND_III.pdf\#nameddest=tag-01K9}{III.? (p.~?)} }
\expandafter\def\csname SRef@courseIII:sphere-twists-and-the-original-smooth-six-sphere-comparison:twist-actual-tube\endcsname{\href{S6_SERIES_GRIND_III.pdf\#nameddest=tag-01MN}{III.? (p.~?)} }
\expandafter\def\csname SRef@courseIII:sphere-twists-and-the-original-smooth-six-sphere-comparison:twist-signed-action\endcsname{\href{S6_SERIES_GRIND_III.pdf\#nameddest=tag-01MT}{III.? (p.~?)} }
\expandafter\def\csname SRef@courseIII:sphere-twists-and-the-original-smooth-six-sphere-comparison:twist-integral-lagrangians\endcsname{\href{S6_SERIES_GRIND_III.pdf\#nameddest=tag-01MP}{III.? (p.~?)} }
\expandafter\def\csname SRef@courseIII:sphere-twists-and-the-original-smooth-six-sphere-comparison:twist-original-middle-level\endcsname{\href{S6_SERIES_GRIND_III.pdf\#nameddest=tag-01MR}{III.? (p.~?)} }
\expandafter\def\csname SRef@courseIII:sphere-twists-and-the-original-smooth-six-sphere-comparison:twist-original-gluing\endcsname{\href{S6_SERIES_GRIND_III.pdf\#nameddest=tag-01MQ}{III.? (p.~?)} }
\expandafter\def\csname SRef@courseIII:sphere-twists-and-the-original-smooth-six-sphere-comparison:twist-actual-cancellation\endcsname{\href{S6_SERIES_GRIND_III.pdf\#nameddest=tag-01MM}{III.? (p.~?)} }
\expandafter\def\csname SRef@courseIII:sphere-twists-and-the-original-smooth-six-sphere-comparison:twist-original-smooth-map\endcsname{\href{S6_SERIES_GRIND_III.pdf\#nameddest=tag-01MS}{III.? (p.~?)} }
\expandafter\def\csname SRef@courseIII:sphere-twists-and-the-original-smooth-six-sphere-comparison:twist-general-stable-framing\endcsname{\href{S6_SERIES_GRIND_III.pdf\#nameddest=tag-01MO}{III.? (p.~?)} }
\expandafter\def\csname SRef@courseIII:sphere-twists-and-the-original-smooth-six-sphere-comparison:twist-worked-exercises\endcsname{\href{S6_SERIES_GRIND_III.pdf\#nameddest=tag-01MU}{III.? (p.~?)} }
\expandafter\def\csname SRef@mapsaudit:almost-complex-homotopy\endcsname{\href{S6_SERIES_GRIND_III.pdf\#nameddest=tag-02QD}{III.? (p.~?)} }
\expandafter\def\csname SRef@mapsI:hodge:line\endcsname{\href{S6_SERIES_GRIND_I.pdf\#nameddest=tag-02OV}{I.? (p.~?)} }
\expandafter\def\csname SRef@mapsI:chain:lemma:1.1\endcsname{\href{S6_SERIES_GRIND_I.pdf\#nameddest=tag-02OL}{I.? (p.~?)} }
\expandafter\def\csname SRef@mapsI:chain:lemma:2.1\endcsname{\href{S6_SERIES_GRIND_I.pdf\#nameddest=tag-02OM}{I.? (p.~?)} }
\expandafter\def\csname SRef@mapsI:chain:corollary:2.2\endcsname{\href{S6_SERIES_GRIND_I.pdf\#nameddest=tag-02OK}{I.? (p.~?)} }
\expandafter\def\csname SRef@mapsI:chain:lemma:4.1\endcsname{\href{S6_SERIES_GRIND_I.pdf\#nameddest=tag-02ON}{I.? (p.~?)} }
\expandafter\def\csname SRef@mapsI:cw:hurewicz-theorem-result\endcsname{\href{S6_SERIES_GRIND_I.pdf\#nameddest=tag-02OU}{I.? (p.~?)} }
\expandafter\def\csname SRef@mapsI:cw:exercise:5.1\endcsname{\href{S6_SERIES_GRIND_I.pdf\#nameddest=tag-02OS}{I.? (p.~?)} }
\expandafter\def\csname SRef@mapsI:cw:exercise:5.2\endcsname{\href{S6_SERIES_GRIND_I.pdf\#nameddest=tag-02OT}{I.? (p.~?)} }
\expandafter\def\csname SRef@mapsI:pf:theorem:5.1\endcsname{\href{S6_SERIES_GRIND_I.pdf\#nameddest=tag-02OY}{I.? (p.~?)} }
\expandafter\def\csname SRef@mapsI:pf:exercise:6.1\endcsname{\href{S6_SERIES_GRIND_I.pdf\#nameddest=tag-02OW}{I.? (p.~?)} }
\expandafter\def\csname SRef@mapsI:pf:exercise:6.2\endcsname{\href{S6_SERIES_GRIND_I.pdf\#nameddest=tag-02OX}{I.? (p.~?)} }
\expandafter\def\csname SRef@courseIII:thom-euler-for-lesson11:CG-S6-CHAIN-LEMMA-5.1\endcsname{\href{S6_SERIES_GRIND_III.pdf\#nameddest=tag-01N3}{III.? (p.~?)} }
\expandafter\def\csname SRef@courseIII:thom-euler-for-lesson11:CG-S6-CHAIN-THEOREM-5.2\endcsname{\href{S6_SERIES_GRIND_III.pdf\#nameddest=tag-01N7}{III.? (p.~?)} }
\expandafter\def\csname SRef@courseIII:thom-euler-for-lesson11:CG-S6-CHAIN-PROPOSITION-6.1\endcsname{\href{S6_SERIES_GRIND_III.pdf\#nameddest=tag-01N4}{III.? (p.~?)} }
\expandafter\def\csname SRef@courseIII:thom-euler-for-lesson11:CG-S6-CHAIN-THEOREM-7.1\endcsname{\href{S6_SERIES_GRIND_III.pdf\#nameddest=tag-01N8}{III.? (p.~?)} }
\expandafter\def\csname SRef@courseIII:thom-euler-for-lesson11:CG-S6-CHAIN-PROPOSITION-8.1\endcsname{\href{S6_SERIES_GRIND_III.pdf\#nameddest=tag-01N5}{III.? (p.~?)} }
\expandafter\def\csname SRef@courseIII:thom-euler-for-lesson11:CG-S6-CHAIN-PROPOSITION-8.2\endcsname{\href{S6_SERIES_GRIND_III.pdf\#nameddest=tag-01N6}{III.? (p.~?)} }
\expandafter\def\csname SRef@courseIII:projective-splitting-for-lesson11:CG-S6-PROJECTIVE-1\endcsname{\href{S6_SERIES_GRIND_III.pdf\#nameddest=tag-01M7}{III.? (p.~?)} }
\expandafter\def\csname SRef@courseIII:projective-splitting-for-lesson11:CG-S6-PROJECTIVE-2\endcsname{\href{S6_SERIES_GRIND_III.pdf\#nameddest=tag-01M8}{III.? (p.~?)} }
\expandafter\def\csname SRef@courseIII:projective-splitting-for-lesson11:CG-S6-PROJECTIVE-3\endcsname{\href{S6_SERIES_GRIND_III.pdf\#nameddest=tag-01M9}{III.? (p.~?)} }
\expandafter\def\csname SRef@courseIII:projective-splitting-for-lesson11:CG-S6-PROJECTIVE-4\endcsname{\href{S6_SERIES_GRIND_III.pdf\#nameddest=tag-01MA}{III.? (p.~?)} }
\expandafter\def\csname SRef@courseIII:projective-splitting-for-lesson11:CG-S6-PROJECTIVE-5\endcsname{\href{S6_SERIES_GRIND_III.pdf\#nameddest=tag-01MB}{III.? (p.~?)} }
\expandafter\def\csname SRef@courseIII:integral-chern-classes-for-lesson11:CG-S6-CHERN-1\endcsname{\href{S6_SERIES_GRIND_III.pdf\#nameddest=tag-01K6}{III.? (p.~?)} }
\expandafter\def\csname SRef@courseIII:integral-chern-classes-for-lesson11:CG-S6-CHERN-2\endcsname{\href{S6_SERIES_GRIND_III.pdf\#nameddest=tag-01K7}{III.? (p.~?)} }
\expandafter\def\csname SRef@courseIII:integral-chern-classes-for-lesson11:CG-S6-CHERN-3\endcsname{\href{S6_SERIES_GRIND_III.pdf\#nameddest=tag-01K8}{III.? (p.~?)} }
\expandafter\def\csname SRef@courseIII:framed6:1\endcsname{\href{S6_SERIES_GRIND_III.pdf\#nameddest=tag-01J8}{III.? (p.~?)} }
\expandafter\def\csname SRef@courseIII:framed6:2\endcsname{\href{S6_SERIES_GRIND_III.pdf\#nameddest=tag-01JB}{III.? (p.~?)} }
\expandafter\def\csname SRef@courseIII:framed6:3\endcsname{\href{S6_SERIES_GRIND_III.pdf\#nameddest=tag-01JC}{III.? (p.~?)} }
\expandafter\def\csname SRef@courseIII:framed6:4\endcsname{\href{S6_SERIES_GRIND_III.pdf\#nameddest=tag-01JD}{III.? (p.~?)} }
\expandafter\def\csname SRef@courseIII:framed6:5\endcsname{\href{S6_SERIES_GRIND_III.pdf\#nameddest=tag-01JE}{III.? (p.~?)} }
\expandafter\def\csname SRef@courseIII:framed6:6\endcsname{\href{S6_SERIES_GRIND_III.pdf\#nameddest=tag-01JF}{III.? (p.~?)} }
\expandafter\def\csname SRef@courseIII:framed6:7\endcsname{\href{S6_SERIES_GRIND_III.pdf\#nameddest=tag-01JG}{III.? (p.~?)} }
\expandafter\def\csname SRef@courseIII:framed6:8\endcsname{\href{S6_SERIES_GRIND_III.pdf\#nameddest=tag-01JH}{III.? (p.~?)} }
\expandafter\def\csname SRef@courseIII:framed6:9\endcsname{\href{S6_SERIES_GRIND_III.pdf\#nameddest=tag-01JI}{III.? (p.~?)} }
\expandafter\def\csname SRef@courseIII:framed6:10.1\endcsname{\href{S6_SERIES_GRIND_III.pdf\#nameddest=tag-01J9}{III.? (p.~?)} }
\expandafter\def\csname SRef@courseIII:framed6:10.2\endcsname{\href{S6_SERIES_GRIND_III.pdf\#nameddest=tag-01JA}{III.? (p.~?)} }
\expandafter\def\csname SRef@courseIII:whitney-receiving:1\endcsname{\href{S6_SERIES_GRIND_III.pdf\#nameddest=tag-01NG}{III.? (p.~?)} }
\expandafter\def\csname SRef@courseIII:whitney-receiving:2\endcsname{\href{S6_SERIES_GRIND_III.pdf\#nameddest=tag-01NH}{III.? (p.~?)} }
\expandafter\def\csname SRef@courseIII:whitney-receiving:3\endcsname{\href{S6_SERIES_GRIND_III.pdf\#nameddest=tag-01NI}{III.? (p.~?)} }
\expandafter\def\csname SRef@courseIII:whitney-receiving:4\endcsname{\href{S6_SERIES_GRIND_III.pdf\#nameddest=tag-01NJ}{III.? (p.~?)} }
\expandafter\def\csname SRef@courseIII:whitney-receiving:5.1\endcsname{\href{S6_SERIES_GRIND_III.pdf\#nameddest=tag-01NK}{III.? (p.~?)} }
\expandafter\def\csname SRef@courseIII:whitney-receiving:5.2\endcsname{\href{S6_SERIES_GRIND_III.pdf\#nameddest=tag-01NL}{III.? (p.~?)} }
\expandafter\def\csname SRef@courseIII:higher-slides:1\endcsname{\href{S6_SERIES_GRIND_III.pdf\#nameddest=tag-01JY}{III.? (p.~?)} }
\expandafter\def\csname SRef@courseIII:higher-slides:2\endcsname{\href{S6_SERIES_GRIND_III.pdf\#nameddest=tag-01JZ}{III.? (p.~?)} }
\expandafter\def\csname SRef@courseIII:higher-slides:3\endcsname{\href{S6_SERIES_GRIND_III.pdf\#nameddest=tag-01K0}{III.? (p.~?)} }
\expandafter\def\csname SRef@courseIII:higher-slides:4\endcsname{\href{S6_SERIES_GRIND_III.pdf\#nameddest=tag-01K1}{III.? (p.~?)} }
\expandafter\def\csname SRef@courseIII:higher-slides:5\endcsname{\href{S6_SERIES_GRIND_III.pdf\#nameddest=tag-01K2}{III.? (p.~?)} }
\expandafter\def\csname SRef@courseIII:higher-slides:6\endcsname{\href{S6_SERIES_GRIND_III.pdf\#nameddest=tag-01K3}{III.? (p.~?)} }
\expandafter\def\csname SRef@courseIII:higher-slides:8.1\endcsname{\href{S6_SERIES_GRIND_III.pdf\#nameddest=tag-01K4}{III.? (p.~?)} }
\expandafter\def\csname SRef@courseIII:higher-slides:8.2\endcsname{\href{S6_SERIES_GRIND_III.pdf\#nameddest=tag-01K5}{III.? (p.~?)} }
\expandafter\def\csname SRef@courseIII:middle-cancel:1\endcsname{\href{S6_SERIES_GRIND_III.pdf\#nameddest=tag-01KT}{III.? (p.~?)} }
\expandafter\def\csname SRef@courseIII:middle-cancel:2\endcsname{\href{S6_SERIES_GRIND_III.pdf\#nameddest=tag-01KU}{III.? (p.~?)} }
\expandafter\def\csname SRef@courseIII:middle-cancel:3\endcsname{\href{S6_SERIES_GRIND_III.pdf\#nameddest=tag-01KV}{III.? (p.~?)} }
\expandafter\def\csname SRef@courseIII:middle-cancel:4\endcsname{\href{S6_SERIES_GRIND_III.pdf\#nameddest=tag-01KW}{III.? (p.~?)} }
\expandafter\def\csname SRef@courseIII:middle-cancel:5\endcsname{\href{S6_SERIES_GRIND_III.pdf\#nameddest=tag-01KX}{III.? (p.~?)} }
\expandafter\def\csname SRef@courseIII:middle-cancel:6\endcsname{\href{S6_SERIES_GRIND_III.pdf\#nameddest=tag-01KY}{III.? (p.~?)} }
\expandafter\def\csname SRef@courseIII:middle-cancel:7.1\endcsname{\href{S6_SERIES_GRIND_III.pdf\#nameddest=tag-01KZ}{III.? (p.~?)} }
\expandafter\def\csname SRef@courseIII:middle-cancel:7.2\endcsname{\href{S6_SERIES_GRIND_III.pdf\#nameddest=tag-01L0}{III.? (p.~?)} }
\expandafter\def\csname SRef@courseIII:morse-cancel:1\endcsname{\href{S6_SERIES_GRIND_III.pdf\#nameddest=tag-01L5}{III.? (p.~?)} }
\expandafter\def\csname SRef@courseIII:morse-cancel:2.1\endcsname{\href{S6_SERIES_GRIND_III.pdf\#nameddest=tag-01L6}{III.? (p.~?)} }
\expandafter\def\csname SRef@courseIII:morse-cancel:6\endcsname{\href{S6_SERIES_GRIND_III.pdf\#nameddest=tag-01L7}{III.? (p.~?)} }
\expandafter\def\csname SRef@courseIII:morse-cancel:7.1\endcsname{\href{S6_SERIES_GRIND_III.pdf\#nameddest=tag-01L8}{III.? (p.~?)} }
\expandafter\def\csname SRef@courseIII:morse-cancel:7.2\endcsname{\href{S6_SERIES_GRIND_III.pdf\#nameddest=tag-01L9}{III.? (p.~?)} }
\expandafter\def\csname SRef@courseIII:spin6:1\endcsname{\href{S6_SERIES_GRIND_III.pdf\#nameddest=tag-01MV}{III.? (p.~?)} }
\expandafter\def\csname SRef@courseIII:spin6:2\endcsname{\href{S6_SERIES_GRIND_III.pdf\#nameddest=tag-01MW}{III.? (p.~?)} }
\expandafter\def\csname SRef@courseIII:spin6:3\endcsname{\href{S6_SERIES_GRIND_III.pdf\#nameddest=tag-01MX}{III.? (p.~?)} }
\expandafter\def\csname SRef@courseIII:spin6:4\endcsname{\href{S6_SERIES_GRIND_III.pdf\#nameddest=tag-01MY}{III.? (p.~?)} }
\expandafter\def\csname SRef@courseIII:spin6:5\endcsname{\href{S6_SERIES_GRIND_III.pdf\#nameddest=tag-01MZ}{III.? (p.~?)} }
\expandafter\def\csname SRef@courseIII:spin6:6\endcsname{\href{S6_SERIES_GRIND_III.pdf\#nameddest=tag-01N0}{III.? (p.~?)} }
\expandafter\def\csname SRef@courseIII:spin6:7.1\endcsname{\href{S6_SERIES_GRIND_III.pdf\#nameddest=tag-01N1}{III.? (p.~?)} }
\expandafter\def\csname SRef@courseIII:spin6:7.2\endcsname{\href{S6_SERIES_GRIND_III.pdf\#nameddest=tag-01N2}{III.? (p.~?)} }
\expandafter\def\csname SRef@courseIII:integral-stable:1\endcsname{\href{S6_SERIES_GRIND_III.pdf\#nameddest=tag-01KA}{III.? (p.~?)} }
\expandafter\def\csname SRef@courseIII:integral-stable:2\endcsname{\href{S6_SERIES_GRIND_III.pdf\#nameddest=tag-01KB}{III.? (p.~?)} }
\expandafter\def\csname SRef@courseIII:integral-stable:3\endcsname{\href{S6_SERIES_GRIND_III.pdf\#nameddest=tag-01KC}{III.? (p.~?)} }
\expandafter\def\csname SRef@courseIII:integral-stable:4\endcsname{\href{S6_SERIES_GRIND_III.pdf\#nameddest=tag-01KD}{III.? (p.~?)} }
\expandafter\def\csname SRef@courseIII:integral-stable:5\endcsname{\href{S6_SERIES_GRIND_III.pdf\#nameddest=tag-01KE}{III.? (p.~?)} }
\expandafter\def\csname SRef@courseIII:integral-stable:6\endcsname{\href{S6_SERIES_GRIND_III.pdf\#nameddest=tag-01KF}{III.? (p.~?)} }
\expandafter\def\csname SRef@courseIII:integral-stable:7\endcsname{\href{S6_SERIES_GRIND_III.pdf\#nameddest=tag-01KG}{III.? (p.~?)} }
\expandafter\def\csname SRef@courseIII:integral-stable:8.1\endcsname{\href{S6_SERIES_GRIND_III.pdf\#nameddest=tag-01KH}{III.? (p.~?)} }
\expandafter\def\csname SRef@courseIII:integral-stable:8.2\endcsname{\href{S6_SERIES_GRIND_III.pdf\#nameddest=tag-01KI}{III.? (p.~?)} }
\expandafter\def\csname SRef@courseIII:pontryagin-thom:1\endcsname{\href{S6_SERIES_GRIND_III.pdf\#nameddest=tag-01LQ}{III.? (p.~?)} }
\expandafter\def\csname SRef@courseIII:pontryagin-thom:2\endcsname{\href{S6_SERIES_GRIND_III.pdf\#nameddest=tag-01LR}{III.? (p.~?)} }
\expandafter\def\csname SRef@courseIII:pontryagin-thom:3\endcsname{\href{S6_SERIES_GRIND_III.pdf\#nameddest=tag-01LS}{III.? (p.~?)} }
\expandafter\def\csname SRef@courseIII:pontryagin-thom:4\endcsname{\href{S6_SERIES_GRIND_III.pdf\#nameddest=tag-01LT}{III.? (p.~?)} }
\expandafter\def\csname SRef@courseIII:pontryagin-thom:5\endcsname{\href{S6_SERIES_GRIND_III.pdf\#nameddest=tag-01LU}{III.? (p.~?)} }
\expandafter\def\csname SRef@courseIII:pontryagin-thom:6\endcsname{\href{S6_SERIES_GRIND_III.pdf\#nameddest=tag-01LV}{III.? (p.~?)} }
\expandafter\def\csname SRef@courseIII:pontryagin-thom:7\endcsname{\href{S6_SERIES_GRIND_III.pdf\#nameddest=tag-01LW}{III.? (p.~?)} }
\expandafter\def\csname SRef@courseIII:pontryagin-thom:8\endcsname{\href{S6_SERIES_GRIND_III.pdf\#nameddest=tag-01LX}{III.? (p.~?)} }
\expandafter\def\csname SRef@courseIII:pontryagin-thom:9.1\endcsname{\href{S6_SERIES_GRIND_III.pdf\#nameddest=tag-01LY}{III.? (p.~?)} }
\expandafter\def\csname SRef@courseIII:pontryagin-thom:9.2\endcsname{\href{S6_SERIES_GRIND_III.pdf\#nameddest=tag-01LZ}{III.? (p.~?)} }
\expandafter\def\csname SRef@courseIII:relative-morse:1\endcsname{\href{S6_SERIES_GRIND_III.pdf\#nameddest=tag-01MG}{III.? (p.~?)} }
\expandafter\def\csname SRef@courseIII:relative-morse:2\endcsname{\href{S6_SERIES_GRIND_III.pdf\#nameddest=tag-01MH}{III.? (p.~?)} }
\expandafter\def\csname SRef@courseIII:relative-morse:3\endcsname{\href{S6_SERIES_GRIND_III.pdf\#nameddest=tag-01MI}{III.? (p.~?)} }
\expandafter\def\csname SRef@courseIII:relative-morse:4\endcsname{\href{S6_SERIES_GRIND_III.pdf\#nameddest=tag-01MJ}{III.? (p.~?)} }
\expandafter\def\csname SRef@courseIII:relative-morse:5.1\endcsname{\href{S6_SERIES_GRIND_III.pdf\#nameddest=tag-01MK}{III.? (p.~?)} }
\expandafter\def\csname SRef@courseIII:relative-morse:5.2\endcsname{\href{S6_SERIES_GRIND_III.pdf\#nameddest=tag-01ML}{III.? (p.~?)} }
\expandafter\def\csname SRef@courseIII:handle-arrangement:1\endcsname{\href{S6_SERIES_GRIND_III.pdf\#nameddest=tag-01JJ}{III.? (p.~?)} }
\expandafter\def\csname SRef@courseIII:handle-arrangement:2\endcsname{\href{S6_SERIES_GRIND_III.pdf\#nameddest=tag-01JK}{III.? (p.~?)} }
\expandafter\def\csname SRef@courseIII:handle-arrangement:3\endcsname{\href{S6_SERIES_GRIND_III.pdf\#nameddest=tag-01JL}{III.? (p.~?)} }
\expandafter\def\csname SRef@courseIII:handle-arrangement:4\endcsname{\href{S6_SERIES_GRIND_III.pdf\#nameddest=tag-01JM}{III.? (p.~?)} }
\expandafter\def\csname SRef@courseIII:handle-arrangement:5.1\endcsname{\href{S6_SERIES_GRIND_III.pdf\#nameddest=tag-01JN}{III.? (p.~?)} }
\expandafter\def\csname SRef@courseIII:handle-arrangement:5.2\endcsname{\href{S6_SERIES_GRIND_III.pdf\#nameddest=tag-01JO}{III.? (p.~?)} }
\expandafter\def\csname SRef@courseIII:zero-top:1\endcsname{\href{S6_SERIES_GRIND_III.pdf\#nameddest=tag-01NU}{III.? (p.~?)} }
\expandafter\def\csname SRef@courseIII:zero-top:2\endcsname{\href{S6_SERIES_GRIND_III.pdf\#nameddest=tag-01NV}{III.? (p.~?)} }
\expandafter\def\csname SRef@courseIII:zero-top:3\endcsname{\href{S6_SERIES_GRIND_III.pdf\#nameddest=tag-01NW}{III.? (p.~?)} }
\expandafter\def\csname SRef@courseIII:zero-top:4\endcsname{\href{S6_SERIES_GRIND_III.pdf\#nameddest=tag-01NX}{III.? (p.~?)} }
\expandafter\def\csname SRef@courseIII:zero-top:5\endcsname{\href{S6_SERIES_GRIND_III.pdf\#nameddest=tag-01NY}{III.? (p.~?)} }
\expandafter\def\csname SRef@courseIII:zero-top:6.1\endcsname{\href{S6_SERIES_GRIND_III.pdf\#nameddest=tag-01NZ}{III.? (p.~?)} }
\expandafter\def\csname SRef@courseIII:zero-top:6.2\endcsname{\href{S6_SERIES_GRIND_III.pdf\#nameddest=tag-01O0}{III.? (p.~?)} }
\expandafter\def\csname SRef@courseIII:trajectory-lowering:1\endcsname{\href{S6_SERIES_GRIND_III.pdf\#nameddest=tag-01N9}{III.? (p.~?)} }
\expandafter\def\csname SRef@courseIII:trajectory-lowering:2\endcsname{\href{S6_SERIES_GRIND_III.pdf\#nameddest=tag-01NA}{III.? (p.~?)} }
\expandafter\def\csname SRef@courseIII:trajectory-lowering:3\endcsname{\href{S6_SERIES_GRIND_III.pdf\#nameddest=tag-01NB}{III.? (p.~?)} }
\expandafter\def\csname SRef@courseIII:trajectory-lowering:4\endcsname{\href{S6_SERIES_GRIND_III.pdf\#nameddest=tag-01NC}{III.? (p.~?)} }
\expandafter\def\csname SRef@courseIII:trajectory-lowering:5\endcsname{\href{S6_SERIES_GRIND_III.pdf\#nameddest=tag-01ND}{III.? (p.~?)} }
\expandafter\def\csname SRef@courseIII:trajectory-lowering:6.1\endcsname{\href{S6_SERIES_GRIND_III.pdf\#nameddest=tag-01NE}{III.? (p.~?)} }
\expandafter\def\csname SRef@courseIII:trajectory-lowering:6.2\endcsname{\href{S6_SERIES_GRIND_III.pdf\#nameddest=tag-01NF}{III.? (p.~?)} }
\expandafter\def\csname SRef@courseIII:belt-complement:1\endcsname{\href{S6_SERIES_GRIND_III.pdf\#nameddest=tag-01IF}{III.? (p.~?)} }
\expandafter\def\csname SRef@courseIII:belt-complement:2\endcsname{\href{S6_SERIES_GRIND_III.pdf\#nameddest=tag-01IG}{III.? (p.~?)} }
\expandafter\def\csname SRef@courseIII:belt-complement:3\endcsname{\href{S6_SERIES_GRIND_III.pdf\#nameddest=tag-01IH}{III.? (p.~?)} }
\expandafter\def\csname SRef@courseIII:belt-complement:4\endcsname{\href{S6_SERIES_GRIND_III.pdf\#nameddest=tag-01II}{III.? (p.~?)} }
\expandafter\def\csname SRef@courseIII:belt-complement:5\endcsname{\href{S6_SERIES_GRIND_III.pdf\#nameddest=tag-01IJ}{III.? (p.~?)} }
\expandafter\def\csname SRef@courseIII:belt-complement:6\endcsname{\href{S6_SERIES_GRIND_III.pdf\#nameddest=tag-01IK}{III.? (p.~?)} }
\expandafter\def\csname SRef@courseIII:belt-complement:7\endcsname{\href{S6_SERIES_GRIND_III.pdf\#nameddest=tag-01IL}{III.? (p.~?)} }
\expandafter\def\csname SRef@courseIII:whitney-support:1\endcsname{\href{S6_SERIES_GRIND_III.pdf\#nameddest=tag-01NM}{III.? (p.~?)} }
\expandafter\def\csname SRef@courseIII:whitney-support:2\endcsname{\href{S6_SERIES_GRIND_III.pdf\#nameddest=tag-01NN}{III.? (p.~?)} }
\expandafter\def\csname SRef@courseIII:whitney-support:3\endcsname{\href{S6_SERIES_GRIND_III.pdf\#nameddest=tag-01NO}{III.? (p.~?)} }
\expandafter\def\csname SRef@courseIII:whitney-support:4\endcsname{\href{S6_SERIES_GRIND_III.pdf\#nameddest=tag-01NP}{III.? (p.~?)} }
\expandafter\def\csname SRef@courseIII:whitney-support:5\endcsname{\href{S6_SERIES_GRIND_III.pdf\#nameddest=tag-01NQ}{III.? (p.~?)} }
\expandafter\def\csname SRef@courseIII:whitney-support:6\endcsname{\href{S6_SERIES_GRIND_III.pdf\#nameddest=tag-01NR}{III.? (p.~?)} }
\expandafter\def\csname SRef@courseIII:whitney-support:7.1\endcsname{\href{S6_SERIES_GRIND_III.pdf\#nameddest=tag-01NS}{III.? (p.~?)} }
\expandafter\def\csname SRef@courseIII:whitney-support:7.2\endcsname{\href{S6_SERIES_GRIND_III.pdf\#nameddest=tag-01NT}{III.? (p.~?)} }
\expandafter\def\csname SRef@courseIII:handle-chains:1\endcsname{\href{S6_SERIES_GRIND_III.pdf\#nameddest=tag-01JP}{III.? (p.~?)} }
\expandafter\def\csname SRef@courseIII:handle-chains:2\endcsname{\href{S6_SERIES_GRIND_III.pdf\#nameddest=tag-01JQ}{III.? (p.~?)} }
\expandafter\def\csname SRef@courseIII:handle-chains:3\endcsname{\href{S6_SERIES_GRIND_III.pdf\#nameddest=tag-01JR}{III.? (p.~?)} }
\expandafter\def\csname SRef@courseIII:handle-chains:4\endcsname{\href{S6_SERIES_GRIND_III.pdf\#nameddest=tag-01JS}{III.? (p.~?)} }
\expandafter\def\csname SRef@courseIII:handle-chains:5\endcsname{\href{S6_SERIES_GRIND_III.pdf\#nameddest=tag-01JT}{III.? (p.~?)} }
\expandafter\def\csname SRef@courseIII:handle-chains:6.1\endcsname{\href{S6_SERIES_GRIND_III.pdf\#nameddest=tag-01JU}{III.? (p.~?)} }
\expandafter\def\csname SRef@courseIII:handle-chains:6.2\endcsname{\href{S6_SERIES_GRIND_III.pdf\#nameddest=tag-01JV}{III.? (p.~?)} }
\expandafter\def\csname SRef@courseIII:handle-chains:8.1\endcsname{\href{S6_SERIES_GRIND_III.pdf\#nameddest=tag-01JW}{III.? (p.~?)} }
\expandafter\def\csname SRef@courseIII:handle-chains:8.2\endcsname{\href{S6_SERIES_GRIND_III.pdf\#nameddest=tag-01JX}{III.? (p.~?)} }
\expandafter\def\csname SRef@courseIII:pair-creation:1\endcsname{\href{S6_SERIES_GRIND_III.pdf\#nameddest=tag-01LI}{III.? (p.~?)} }
\expandafter\def\csname SRef@courseIII:pair-creation:2\endcsname{\href{S6_SERIES_GRIND_III.pdf\#nameddest=tag-01LJ}{III.? (p.~?)} }
\expandafter\def\csname SRef@courseIII:pair-creation:3\endcsname{\href{S6_SERIES_GRIND_III.pdf\#nameddest=tag-01LK}{III.? (p.~?)} }
\expandafter\def\csname SRef@courseIII:pair-creation:4\endcsname{\href{S6_SERIES_GRIND_III.pdf\#nameddest=tag-01LL}{III.? (p.~?)} }
\expandafter\def\csname SRef@courseIII:pair-creation:5\endcsname{\href{S6_SERIES_GRIND_III.pdf\#nameddest=tag-01LM}{III.? (p.~?)} }
\expandafter\def\csname SRef@courseIII:pair-creation:6\endcsname{\href{S6_SERIES_GRIND_III.pdf\#nameddest=tag-01LN}{III.? (p.~?)} }
\expandafter\def\csname SRef@courseIII:pair-creation:7.1\endcsname{\href{S6_SERIES_GRIND_III.pdf\#nameddest=tag-01LO}{III.? (p.~?)} }
\expandafter\def\csname SRef@courseIII:pair-creation:7.2\endcsname{\href{S6_SERIES_GRIND_III.pdf\#nameddest=tag-01LP}{III.? (p.~?)} }
\expandafter\def\csname SRef@courseIII:low-index-slides:1\endcsname{\href{S6_SERIES_GRIND_III.pdf\#nameddest=tag-01KJ}{III.? (p.~?)} }
\expandafter\def\csname SRef@courseIII:low-index-slides:2\endcsname{\href{S6_SERIES_GRIND_III.pdf\#nameddest=tag-01KK}{III.? (p.~?)} }
\expandafter\def\csname SRef@courseIII:low-index-slides:3\endcsname{\href{S6_SERIES_GRIND_III.pdf\#nameddest=tag-01KL}{III.? (p.~?)} }
\expandafter\def\csname SRef@courseIII:low-index-slides:4\endcsname{\href{S6_SERIES_GRIND_III.pdf\#nameddest=tag-01KM}{III.? (p.~?)} }
\expandafter\def\csname SRef@courseIII:low-index-slides:5\endcsname{\href{S6_SERIES_GRIND_III.pdf\#nameddest=tag-01KN}{III.? (p.~?)} }
\expandafter\def\csname SRef@courseIII:low-index-slides:6\endcsname{\href{S6_SERIES_GRIND_III.pdf\#nameddest=tag-01KO}{III.? (p.~?)} }
\expandafter\def\csname SRef@courseIII:low-index-slides:7\endcsname{\href{S6_SERIES_GRIND_III.pdf\#nameddest=tag-01KP}{III.? (p.~?)} }
\expandafter\def\csname SRef@courseIII:low-index-slides:8\endcsname{\href{S6_SERIES_GRIND_III.pdf\#nameddest=tag-01KQ}{III.? (p.~?)} }
\expandafter\def\csname SRef@courseIII:low-index-slides:9.1\endcsname{\href{S6_SERIES_GRIND_III.pdf\#nameddest=tag-01KR}{III.? (p.~?)} }
\expandafter\def\csname SRef@courseIII:low-index-slides:9.2\endcsname{\href{S6_SERIES_GRIND_III.pdf\#nameddest=tag-01KS}{III.? (p.~?)} }
\expandafter\def\csname SRef@courseIII:adams-arithmetic:1\endcsname{\href{S6_SERIES_GRIND_III.pdf\#nameddest=tag-01HJ}{III.? (p.~?)} }
\expandafter\def\csname SRef@courseIII:adams-arithmetic:2\endcsname{\href{S6_SERIES_GRIND_III.pdf\#nameddest=tag-01HK}{III.? (p.~?)} }
\expandafter\def\csname SRef@courseIII:adams-arithmetic:3\endcsname{\href{S6_SERIES_GRIND_III.pdf\#nameddest=tag-01HL}{III.? (p.~?)} }
\expandafter\def\csname SRef@courseIII:adams-arithmetic:4\endcsname{\href{S6_SERIES_GRIND_III.pdf\#nameddest=tag-01HM}{III.? (p.~?)} }
\expandafter\def\csname SRef@courseIII:adams-arithmetic:5\endcsname{\href{S6_SERIES_GRIND_III.pdf\#nameddest=tag-01HN}{III.? (p.~?)} }
\expandafter\def\csname SRef@courseIII:adams-arithmetic:6\endcsname{\href{S6_SERIES_GRIND_III.pdf\#nameddest=tag-01HO}{III.? (p.~?)} }
\expandafter\def\csname SRef@courseIII:adams-arithmetic:7\endcsname{\href{S6_SERIES_GRIND_III.pdf\#nameddest=tag-01HP}{III.? (p.~?)} }
\expandafter\def\csname SRef@courseIII:adams-arithmetic:8\endcsname{\href{S6_SERIES_GRIND_III.pdf\#nameddest=tag-01HQ}{III.? (p.~?)} }
\expandafter\def\csname SRef@courseIII:adams-arithmetic:9\endcsname{\href{S6_SERIES_GRIND_III.pdf\#nameddest=tag-01HR}{III.? (p.~?)} }
\expandafter\def\csname SRef@courseIII:adams-tower:1\endcsname{\href{S6_SERIES_GRIND_III.pdf\#nameddest=tag-01HS}{III.? (p.~?)} }
\expandafter\def\csname SRef@courseIII:adams-tower:2\endcsname{\href{S6_SERIES_GRIND_III.pdf\#nameddest=tag-01HV}{III.? (p.~?)} }
\expandafter\def\csname SRef@courseIII:adams-tower:3\endcsname{\href{S6_SERIES_GRIND_III.pdf\#nameddest=tag-01HW}{III.? (p.~?)} }
\expandafter\def\csname SRef@courseIII:adams-tower:4\endcsname{\href{S6_SERIES_GRIND_III.pdf\#nameddest=tag-01HX}{III.? (p.~?)} }
\expandafter\def\csname SRef@courseIII:adams-tower:5\endcsname{\href{S6_SERIES_GRIND_III.pdf\#nameddest=tag-01HY}{III.? (p.~?)} }
\expandafter\def\csname SRef@courseIII:adams-tower:6\endcsname{\href{S6_SERIES_GRIND_III.pdf\#nameddest=tag-01HZ}{III.? (p.~?)} }
\expandafter\def\csname SRef@courseIII:adams-tower:7.1\endcsname{\href{S6_SERIES_GRIND_III.pdf\#nameddest=tag-01I0}{III.? (p.~?)} }
\expandafter\def\csname SRef@courseIII:adams-tower:7.2\endcsname{\href{S6_SERIES_GRIND_III.pdf\#nameddest=tag-01I1}{III.? (p.~?)} }
\expandafter\def\csname SRef@courseIII:adams-tower:8\endcsname{\href{S6_SERIES_GRIND_III.pdf\#nameddest=tag-01I2}{III.? (p.~?)} }
\expandafter\def\csname SRef@courseIII:adams-tower:9\endcsname{\href{S6_SERIES_GRIND_III.pdf\#nameddest=tag-01I3}{III.? (p.~?)} }
\expandafter\def\csname SRef@courseIII:adams-tower:10\endcsname{\href{S6_SERIES_GRIND_III.pdf\#nameddest=tag-01HT}{III.? (p.~?)} }
\expandafter\def\csname SRef@courseIII:adams-tower:11\endcsname{\href{S6_SERIES_GRIND_III.pdf\#nameddest=tag-01HU}{III.? (p.~?)} }
\expandafter\def\csname SRef@courseIII:formal-char2:1\endcsname{\href{S6_SERIES_GRIND_III.pdf\#nameddest=tag-01J2}{III.? (p.~?)} }
\expandafter\def\csname SRef@courseIII:formal-char2:2\endcsname{\href{S6_SERIES_GRIND_III.pdf\#nameddest=tag-01J3}{III.? (p.~?)} }
\expandafter\def\csname SRef@courseIII:formal-char2:3\endcsname{\href{S6_SERIES_GRIND_III.pdf\#nameddest=tag-01J4}{III.? (p.~?)} }
\expandafter\def\csname SRef@courseIII:formal-char2:4\endcsname{\href{S6_SERIES_GRIND_III.pdf\#nameddest=tag-01J5}{III.? (p.~?)} }
\expandafter\def\csname SRef@courseIII:formal-char2:4.1\endcsname{\href{S6_SERIES_GRIND_III.pdf\#nameddest=tag-01J6}{III.? (p.~?)} }
\expandafter\def\csname SRef@courseIII:formal-char2:4.2\endcsname{\href{S6_SERIES_GRIND_III.pdf\#nameddest=tag-01J7}{III.? (p.~?)} }
\expandafter\def\csname SRef@courseIII:projective-coop:1\endcsname{\href{S6_SERIES_GRIND_III.pdf\#nameddest=tag-01M0}{III.? (p.~?)} }
\expandafter\def\csname SRef@courseIII:projective-coop:2\endcsname{\href{S6_SERIES_GRIND_III.pdf\#nameddest=tag-01M1}{III.? (p.~?)} }
\expandafter\def\csname SRef@courseIII:projective-coop:3\endcsname{\href{S6_SERIES_GRIND_III.pdf\#nameddest=tag-01M2}{III.? (p.~?)} }
\expandafter\def\csname SRef@courseIII:projective-coop:4\endcsname{\href{S6_SERIES_GRIND_III.pdf\#nameddest=tag-01M3}{III.? (p.~?)} }
\expandafter\def\csname SRef@courseIII:projective-coop:5\endcsname{\href{S6_SERIES_GRIND_III.pdf\#nameddest=tag-01M4}{III.? (p.~?)} }
\expandafter\def\csname SRef@courseIII:projective-coop:6.1\endcsname{\href{S6_SERIES_GRIND_III.pdf\#nameddest=tag-01M5}{III.? (p.~?)} }
\expandafter\def\csname SRef@courseIII:projective-coop:6.2\endcsname{\href{S6_SERIES_GRIND_III.pdf\#nameddest=tag-01M6}{III.? (p.~?)} }
\expandafter\def\csname SRef@courseIII:actual-degree8:1\endcsname{\href{S6_SERIES_GRIND_III.pdf\#nameddest=tag-01H0}{III.? (p.~?)} }
\expandafter\def\csname SRef@courseIII:actual-degree8:2\endcsname{\href{S6_SERIES_GRIND_III.pdf\#nameddest=tag-01H1}{III.? (p.~?)} }
\expandafter\def\csname SRef@courseIII:actual-degree8:3\endcsname{\href{S6_SERIES_GRIND_III.pdf\#nameddest=tag-01H2}{III.? (p.~?)} }
\expandafter\def\csname SRef@courseIII:actual-degree8:4\endcsname{\href{S6_SERIES_GRIND_III.pdf\#nameddest=tag-01H3}{III.? (p.~?)} }
\expandafter\def\csname SRef@courseIII:actual-degree8:5\endcsname{\href{S6_SERIES_GRIND_III.pdf\#nameddest=tag-01H4}{III.? (p.~?)} }
\expandafter\def\csname SRef@courseIII:actual-degree8:6\endcsname{\href{S6_SERIES_GRIND_III.pdf\#nameddest=tag-01H5}{III.? (p.~?)} }
\expandafter\def\csname SRef@courseIII:actual-degree8:7\endcsname{\href{S6_SERIES_GRIND_III.pdf\#nameddest=tag-01H6}{III.? (p.~?)} }
\expandafter\def\csname SRef@courseIII:actual-degree8:8\endcsname{\href{S6_SERIES_GRIND_III.pdf\#nameddest=tag-01H7}{III.? (p.~?)} }
\expandafter\def\csname SRef@courseIII:actual-degree8:9.1\endcsname{\href{S6_SERIES_GRIND_III.pdf\#nameddest=tag-01H8}{III.? (p.~?)} }
\expandafter\def\csname SRef@courseIII:actual-degree8:9.2\endcsname{\href{S6_SERIES_GRIND_III.pdf\#nameddest=tag-01H9}{III.? (p.~?)} }
\expandafter\def\csname SRef@courseIII:pt-relative-embedding\endcsname{\href{S6_SERIES_GRIND_III.pdf\#nameddest=tag-01ME}{III.? (p.~?)} }
\expandafter\def\csname SRef@courseIII:pt-relative-normal-frame\endcsname{\href{S6_SERIES_GRIND_III.pdf\#nameddest=tag-01MF}{III.? (p.~?)} }
\expandafter\def\csname SRef@courseIII:pt-normal-inverse-comparison\endcsname{\href{S6_SERIES_GRIND_III.pdf\#nameddest=tag-01MD}{III.? (p.~?)} }
\expandafter\def\csname SRef@courseIII:pt-abstract-framed-comparison\endcsname{\href{S6_SERIES_GRIND_III.pdf\#nameddest=tag-01MC}{III.? (p.~?)} }
\expandafter\def\csname SRef@courseIII:actual-resolution:1\endcsname{\href{S6_SERIES_GRIND_III.pdf\#nameddest=tag-01HA}{III.? (p.~?)} }
\expandafter\def\csname SRef@courseIII:actual-resolution:2\endcsname{\href{S6_SERIES_GRIND_III.pdf\#nameddest=tag-01HB}{III.? (p.~?)} }
\expandafter\def\csname SRef@courseIII:actual-resolution:3\endcsname{\href{S6_SERIES_GRIND_III.pdf\#nameddest=tag-01HC}{III.? (p.~?)} }
\expandafter\def\csname SRef@courseIII:actual-resolution:4\endcsname{\href{S6_SERIES_GRIND_III.pdf\#nameddest=tag-01HD}{III.? (p.~?)} }
\expandafter\def\csname SRef@courseIII:actual-resolution:5\endcsname{\href{S6_SERIES_GRIND_III.pdf\#nameddest=tag-01HE}{III.? (p.~?)} }
\expandafter\def\csname SRef@courseIII:actual-resolution:6\endcsname{\href{S6_SERIES_GRIND_III.pdf\#nameddest=tag-01HF}{III.? (p.~?)} }
\expandafter\def\csname SRef@courseIII:actual-resolution:7\endcsname{\href{S6_SERIES_GRIND_III.pdf\#nameddest=tag-01HG}{III.? (p.~?)} }
\expandafter\def\csname SRef@courseIII:actual-resolution:8.1\endcsname{\href{S6_SERIES_GRIND_III.pdf\#nameddest=tag-01HH}{III.? (p.~?)} }
\expandafter\def\csname SRef@courseIII:actual-resolution:8.2\endcsname{\href{S6_SERIES_GRIND_III.pdf\#nameddest=tag-01HI}{III.? (p.~?)} }
\expandafter\def\csname SRef@courseIII:milnor-cobar:1.1\endcsname{\href{S6_SERIES_GRIND_III.pdf\#nameddest=tag-01L1}{III.? (p.~?)} }
\expandafter\def\csname SRef@courseIII:milnor-cobar:6.1\endcsname{\href{S6_SERIES_GRIND_III.pdf\#nameddest=tag-01L2}{III.? (p.~?)} }
\expandafter\def\csname SRef@courseIII:milnor-cobar:7.1\endcsname{\href{S6_SERIES_GRIND_III.pdf\#nameddest=tag-01L3}{III.? (p.~?)} }
\expandafter\def\csname SRef@courseIII:milnor-cobar:7.2\endcsname{\href{S6_SERIES_GRIND_III.pdf\#nameddest=tag-01L4}{III.? (p.~?)} }
\expandafter\def\csname SRef@courseIII:odd-coop:1\endcsname{\href{S6_SERIES_GRIND_III.pdf\#nameddest=tag-01LA}{III.? (p.~?)} }
\expandafter\def\csname SRef@courseIII:odd-coop:2\endcsname{\href{S6_SERIES_GRIND_III.pdf\#nameddest=tag-01LB}{III.? (p.~?)} }
\expandafter\def\csname SRef@courseIII:odd-coop:3\endcsname{\href{S6_SERIES_GRIND_III.pdf\#nameddest=tag-01LC}{III.? (p.~?)} }
\expandafter\def\csname SRef@courseIII:odd-coop:4\endcsname{\href{S6_SERIES_GRIND_III.pdf\#nameddest=tag-01LD}{III.? (p.~?)} }
\expandafter\def\csname SRef@courseIII:odd-coop:5\endcsname{\href{S6_SERIES_GRIND_III.pdf\#nameddest=tag-01LE}{III.? (p.~?)} }
\expandafter\def\csname SRef@courseIII:odd-coop:6\endcsname{\href{S6_SERIES_GRIND_III.pdf\#nameddest=tag-01LF}{III.? (p.~?)} }
\expandafter\def\csname SRef@courseIII:odd-coop:7.1\endcsname{\href{S6_SERIES_GRIND_III.pdf\#nameddest=tag-01LG}{III.? (p.~?)} }
\expandafter\def\csname SRef@courseIII:odd-coop:7.2\endcsname{\href{S6_SERIES_GRIND_III.pdf\#nameddest=tag-01LH}{III.? (p.~?)} }
\expandafter\def\csname SRef@courseIII:coefficient-rings:1\endcsname{\href{S6_SERIES_GRIND_III.pdf\#nameddest=tag-01IM}{III.? (p.~?)} }
\expandafter\def\csname SRef@courseIII:coefficient-rings:2\endcsname{\href{S6_SERIES_GRIND_III.pdf\#nameddest=tag-01IN}{III.? (p.~?)} }
\expandafter\def\csname SRef@courseIII:coefficient-rings:3\endcsname{\href{S6_SERIES_GRIND_III.pdf\#nameddest=tag-01IO}{III.? (p.~?)} }
\expandafter\def\csname SRef@courseIII:coefficient-rings:4\endcsname{\href{S6_SERIES_GRIND_III.pdf\#nameddest=tag-01IP}{III.? (p.~?)} }
\expandafter\def\csname SRef@courseIII:coefficient-rings:5\endcsname{\href{S6_SERIES_GRIND_III.pdf\#nameddest=tag-01IQ}{III.? (p.~?)} }
\expandafter\def\csname SRef@courseIII:coefficient-rings:6\endcsname{\href{S6_SERIES_GRIND_III.pdf\#nameddest=tag-01IR}{III.? (p.~?)} }
\expandafter\def\csname SRef@courseIII:coefficient-rings:7.1\endcsname{\href{S6_SERIES_GRIND_III.pdf\#nameddest=tag-01IS}{III.? (p.~?)} }
\expandafter\def\csname SRef@courseIII:coefficient-rings:7.2\endcsname{\href{S6_SERIES_GRIND_III.pdf\#nameddest=tag-01IT}{III.? (p.~?)} }
\expandafter\def\csname SRef@courseIII:duality:2.1\endcsname{\href{S6_SERIES_GRIND_III.pdf\#nameddest=tag-01IU}{III.? (p.~?)} }
\expandafter\def\csname SRef@courseIII:duality:2.2\endcsname{\href{S6_SERIES_GRIND_III.pdf\#nameddest=tag-01IV}{III.? (p.~?)} }
\expandafter\def\csname SRef@courseIII:duality:2.3\endcsname{\href{S6_SERIES_GRIND_III.pdf\#nameddest=tag-01IW}{III.? (p.~?)} }
\expandafter\def\csname SRef@courseIII:duality:3.1\endcsname{\href{S6_SERIES_GRIND_III.pdf\#nameddest=tag-01IX}{III.? (p.~?)} }
\expandafter\def\csname SRef@courseIII:duality:3.2\endcsname{\href{S6_SERIES_GRIND_III.pdf\#nameddest=tag-01IY}{III.? (p.~?)} }
\expandafter\def\csname SRef@courseIII:duality:3.3\endcsname{\href{S6_SERIES_GRIND_III.pdf\#nameddest=tag-01IZ}{III.? (p.~?)} }
\expandafter\def\csname SRef@courseIII:duality:4.1\endcsname{\href{S6_SERIES_GRIND_III.pdf\#nameddest=tag-01J0}{III.? (p.~?)} }
\expandafter\def\csname SRef@courseIII:duality:4.2\endcsname{\href{S6_SERIES_GRIND_III.pdf\#nameddest=tag-01J1}{III.? (p.~?)} }
\expandafter\def\csname SRef@night:P1.6\endcsname{\href{S6_SERIES_GRIND_V.pdf\#nameddest=tag-02QX}{V.? (p.~?)} }
\expandafter\def\csname SRef@night:P1.14\endcsname{\href{S6_SERIES_GRIND_V.pdf\#nameddest=tag-02QN}{V.? (p.~?)} }
\expandafter\def\csname SRef@night:P5.7\endcsname{\href{S6_SERIES_GRIND_I.pdf\#nameddest=tag-02YS}{I.A.1 (p.~342)} }
\expandafter\def\csname SRef@ari:A.1.1\endcsname{\href{S6_SERIES_GRIND_I.pdf\#nameddest=tag-03J1}{I.? (p.~?)} }
\expandafter\def\csname SRef@ari:A.1.2\endcsname{\href{S6_SERIES_GRIND_I.pdf\#nameddest=tag-03J2}{I.? (p.~?)} }
\expandafter\def\csname SRef@ari:A.1.3\endcsname{\href{S6_SERIES_GRIND_I.pdf\#nameddest=tag-03J3}{I.? (p.~?)} }
\expandafter\def\csname SRef@ari:A.1.4\endcsname{\href{S6_SERIES_GRIND_I.pdf\#nameddest=tag-03J4}{I.? (p.~?)} }
\expandafter\def\csname SRef@ari:A.2.1\endcsname{\href{S6_SERIES_GRIND_I.pdf\#nameddest=tag-03J5}{I.? (p.~?)} }
\expandafter\def\csname SRef@ari:A.2.2\endcsname{\href{S6_SERIES_GRIND_I.pdf\#nameddest=tag-03J6}{I.? (p.~?)} }
\expandafter\def\csname SRef@ari:A.2.3\endcsname{\href{S6_SERIES_GRIND_I.pdf\#nameddest=tag-03J7}{I.? (p.~?)} }
\expandafter\def\csname SRef@ari:A.2.4\endcsname{\href{S6_SERIES_GRIND_I.pdf\#nameddest=tag-03J8}{I.? (p.~?)} }
\expandafter\def\csname SRef@ari:A.2.5\endcsname{\href{S6_SERIES_GRIND_I.pdf\#nameddest=tag-03J9}{I.? (p.~?)} }
\expandafter\def\csname SRef@ari:A.3.1\endcsname{\href{S6_SERIES_GRIND_I.pdf\#nameddest=tag-03JA}{I.? (p.~?)} }
\expandafter\def\csname SRef@ari:A.3.2\endcsname{\href{S6_SERIES_GRIND_I.pdf\#nameddest=tag-03JB}{I.? (p.~?)} }
\expandafter\def\csname SRef@ari:A.4.1\endcsname{\href{S6_SERIES_GRIND_I.pdf\#nameddest=tag-03JC}{I.? (p.~?)} }
\expandafter\def\csname SRef@ari:A.4.2\endcsname{\href{S6_SERIES_GRIND_I.pdf\#nameddest=tag-03JD}{I.? (p.~?)} }
\expandafter\def\csname SRef@ari:A.4.3\endcsname{\href{S6_SERIES_GRIND_I.pdf\#nameddest=tag-03JE}{I.? (p.~?)} }
\expandafter\def\csname SRef@ari:A.4.4\endcsname{\href{S6_SERIES_GRIND_I.pdf\#nameddest=tag-03JF}{I.? (p.~?)} }
\expandafter\def\csname SRef@ari:A.4.5\endcsname{\href{S6_SERIES_GRIND_I.pdf\#nameddest=tag-03JG}{I.? (p.~?)} }
\expandafter\def\csname SRef@ari:A.4.6\endcsname{\href{S6_SERIES_GRIND_I.pdf\#nameddest=tag-03JH}{I.? (p.~?)} }
\expandafter\def\csname SRef@ari:A.4.7\endcsname{\href{S6_SERIES_GRIND_I.pdf\#nameddest=tag-03JI}{I.? (p.~?)} }
\expandafter\def\csname SRef@ari:A.4.8\endcsname{\href{S6_SERIES_GRIND_I.pdf\#nameddest=tag-03JJ}{I.? (p.~?)} }
\expandafter\def\csname SRef@ari:A.5.1\endcsname{\href{S6_SERIES_GRIND_I.pdf\#nameddest=tag-03JK}{I.? (p.~?)} }
\expandafter\def\csname SRef@ari:A.5.2\endcsname{\href{S6_SERIES_GRIND_I.pdf\#nameddest=tag-03JL}{I.? (p.~?)} }
\expandafter\def\csname SRef@ari:A.5.3\endcsname{\href{S6_SERIES_GRIND_I.pdf\#nameddest=tag-03JM}{I.? (p.~?)} }
\expandafter\def\csname SRef@ari:A.5.4\endcsname{\href{S6_SERIES_GRIND_I.pdf\#nameddest=tag-03JN}{I.? (p.~?)} }
\expandafter\def\csname SRef@ari:A.5.5\endcsname{\href{S6_SERIES_GRIND_I.pdf\#nameddest=tag-03JO}{I.? (p.~?)} }
\expandafter\def\csname SRef@ari:A.5.6\endcsname{\href{S6_SERIES_GRIND_I.pdf\#nameddest=tag-03JP}{I.? (p.~?)} }
\expandafter\def\csname SRef@ari:A.5.8\endcsname{\href{S6_SERIES_GRIND_I.pdf\#nameddest=tag-03JR}{I.? (p.~?)} }
\expandafter\def\csname SRef@ari:A.5.7\endcsname{\href{S6_SERIES_GRIND_I.pdf\#nameddest=tag-03JQ}{I.? (p.~?)} }
\expandafter\def\csname SRef@series:I:main-theorems\endcsname{\href{S6_SERIES_GRIND_I.pdf\#nameddest=tag-03JS}{I.1.2 (p.~6)} }
\expandafter\def\csname SRef@series:I:A\endcsname{\href{S6_SERIES_GRIND_I.pdf\#nameddest=tag-03JT}{I.A (p.~6)} }
\expandafter\def\csname SRef@series:I:B\endcsname{\href{S6_SERIES_GRIND_I.pdf\#nameddest=tag-03JU}{I.B (p.~8)} }
\expandafter\def\csname SRef@series:I:C\endcsname{\href{S6_SERIES_GRIND_I.pdf\#nameddest=tag-03JV}{I.C (p.~9)} }
\expandafter\def\csname SRef@series:I:D\endcsname{\href{S6_SERIES_GRIND_I.pdf\#nameddest=tag-03JW}{I.D (p.~10)} }
\expandafter\def\csname SRef@series:I:E\endcsname{\href{S6_SERIES_GRIND_I.pdf\#nameddest=tag-03JX}{I.E (p.~11)} }
\expandafter\def\csname SRef@series:I:F\endcsname{\href{S6_SERIES_GRIND_I.pdf\#nameddest=tag-03JY}{I.F (p.~12)} }
\expandafter\def\csname SRef@series:I:G\endcsname{\href{S6_SERIES_GRIND_I.pdf\#nameddest=tag-03JZ}{I.G (p.~12)} }
\expandafter\def\csname SRef@series:I:outlines\endcsname{\href{S6_SERIES_GRIND_I.pdf\#nameddest=tag-03K0}{I.1.3 (p.~15)} }
\expandafter\def\csname SRef@series:II:main-theorems\endcsname{\hyperlink{tag-03K1}{II.\ref*{series:II:main-theorems}}}
\expandafter\def\csname SRef@series:II:A\endcsname{\hyperlink{tag-03K2}{II.\ref*{series:II:A}}}
\expandafter\def\csname SRef@series:II:B\endcsname{\hyperlink{tag-03K3}{II.\ref*{series:II:B}}}
\expandafter\def\csname SRef@series:II:C\endcsname{\hyperlink{tag-03K4}{II.\ref*{series:II:C}}}
\expandafter\def\csname SRef@series:II:D\endcsname{\hyperlink{tag-03K5}{II.\ref*{series:II:D}}}
\expandafter\def\csname SRef@series:II:E\endcsname{\hyperlink{tag-03K6}{II.\ref*{series:II:E}}}
\expandafter\def\csname SRef@series:II:F\endcsname{\hyperlink{tag-03K7}{II.\ref*{series:II:F}}}
\expandafter\def\csname SRef@series:II:G\endcsname{\hyperlink{tag-03K8}{II.\ref*{series:II:G}}}
\expandafter\def\csname SRef@series:II:outline\endcsname{\hyperlink{tag-03K9}{II.\ref*{series:II:outline}}}
\expandafter\def\csname SRef@series:III:main-theorems\endcsname{\href{S6_SERIES_GRIND_III.pdf\#nameddest=tag-03KA}{III.1.2 (p.~3)} }
\expandafter\def\csname SRef@series:III:A\endcsname{\href{S6_SERIES_GRIND_III.pdf\#nameddest=tag-03KB}{III.A (p.~3)} }
\expandafter\def\csname SRef@series:III:B\endcsname{\href{S6_SERIES_GRIND_III.pdf\#nameddest=tag-03KC}{III.B (p.~3)} }
\expandafter\def\csname SRef@series:III:C\endcsname{\href{S6_SERIES_GRIND_III.pdf\#nameddest=tag-03KD}{III.C (p.~4)} }
\expandafter\def\csname SRef@series:III:D\endcsname{\href{S6_SERIES_GRIND_III.pdf\#nameddest=tag-03KE}{III.D (p.~4)} }
\expandafter\def\csname SRef@series:III:E\endcsname{\href{S6_SERIES_GRIND_III.pdf\#nameddest=tag-03KF}{III.E (p.~4)} }
\expandafter\def\csname SRef@series:III:F\endcsname{\href{S6_SERIES_GRIND_III.pdf\#nameddest=tag-03KG}{III.F (p.~5)} }
\expandafter\def\csname SRef@series:III:G\endcsname{\href{S6_SERIES_GRIND_III.pdf\#nameddest=tag-03KH}{III.G (p.~5)} }
\expandafter\def\csname SRef@series:IV:main-theorems\endcsname{\href{S6_SERIES_GRIND_IV.pdf\#nameddest=tag-03KI}{IV.1.2 (p.~3)} }
\expandafter\def\csname SRef@series:IV:A\endcsname{\href{S6_SERIES_GRIND_IV.pdf\#nameddest=tag-03KJ}{IV.A (p.~3)} }
\expandafter\def\csname SRef@series:IV:B\endcsname{\href{S6_SERIES_GRIND_IV.pdf\#nameddest=tag-03KK}{IV.B (p.~4)} }
\expandafter\def\csname SRef@series:IV:C\endcsname{\href{S6_SERIES_GRIND_IV.pdf\#nameddest=tag-03KL}{IV.C (p.~5)} }
\expandafter\def\csname SRef@series:IV:D\endcsname{\href{S6_SERIES_GRIND_IV.pdf\#nameddest=tag-03KM}{IV.D (p.~6)} }
\expandafter\def\csname SRef@series:IV:E\endcsname{\href{S6_SERIES_GRIND_IV.pdf\#nameddest=tag-03KN}{IV.E (p.~7)} }
\expandafter\def\csname SRef@series:IV:F\endcsname{\href{S6_SERIES_GRIND_IV.pdf\#nameddest=tag-03KO}{IV.F (p.~8)} }
\expandafter\def\csname SRef@series:IV:G\endcsname{\href{S6_SERIES_GRIND_IV.pdf\#nameddest=tag-03KP}{IV.G (p.~8)} }
\expandafter\def\csname SRef@series:V:main-theorems\endcsname{\href{S6_SERIES_GRIND_V.pdf\#nameddest=tag-03KQ}{V.1.2 (p.~3)} }
\expandafter\def\csname SRef@series:V:A\endcsname{\href{S6_SERIES_GRIND_V.pdf\#nameddest=tag-03KR}{V.A (p.~3)} }
\expandafter\def\csname SRef@series:V:B\endcsname{\href{S6_SERIES_GRIND_V.pdf\#nameddest=tag-03KS}{V.B (p.~3)} }
\expandafter\def\csname SRef@series:V:C\endcsname{\href{S6_SERIES_GRIND_V.pdf\#nameddest=tag-03KT}{V.C (p.~4)} }
\expandafter\def\csname SRef@series:V:D\endcsname{\href{S6_SERIES_GRIND_V.pdf\#nameddest=tag-03KU}{V.D (p.~4)} }
\expandafter\def\csname SRef@series:V:E\endcsname{\href{S6_SERIES_GRIND_V.pdf\#nameddest=tag-03KV}{V.E (p.~5)} }
\expandafter\def\csname SRef@series:V:F\endcsname{\href{S6_SERIES_GRIND_V.pdf\#nameddest=tag-03KW}{V.F (p.~5)} }
\expandafter\def\csname SRef@layout:V:fresh-gauges\endcsname{\href{S6_SERIES_GRIND_V.pdf\#nameddest=tag-03KX}{V.{{$\langle 2\rangle $1.}} (p.~126)} }
\expandafter\def\csname SRef@layout:V:fresh-minimum\endcsname{\href{S6_SERIES_GRIND_V.pdf\#nameddest=tag-03KY}{V.{{$\langle 1\rangle $4.}} (p.~130)} }
\expandafter\def\csname SRef@layout:V:fresh-bounds\endcsname{\href{S6_SERIES_GRIND_V.pdf\#nameddest=tag-03KZ}{V.{1.18.4} (p.~130)} }
\expandafter\def\csname SRef@layout:V:fresh-marking\endcsname{\href{S6_SERIES_GRIND_V.pdf\#nameddest=tag-03L0}{V.{{$\langle 1\rangle $6.}} (p.~132)} }
\expandafter\def\csname SRef@layout:V:fresh-splitting\endcsname{\href{S6_SERIES_GRIND_V.pdf\#nameddest=tag-03L1}{V.{{$\langle 1\rangle $7.}} (p.~133)} }
\expandafter\def\csname SRef@iv:source_transports\endcsname{\href{S6_SERIES_GRIND_IV.pdf\#nameddest=tag-03L2}{IV.3.27 (p.~150)} }
\expandafter\def\csname SRef@inv:section:coefficient-atlas\endcsname{\href{S6_SERIES_GRIND_I.pdf\#nameddest=tag-03L3}{I.D (p.~374)} }
\expandafter\def\csname SRef@inv:eq:atlas-theta-factor\endcsname{\href{S6_SERIES_GRIND_I.pdf\#nameddest=tag-03L4}{I.105 (p.~375)} }
\expandafter\def\csname SRef@inv:eq:atlas-semiflat\endcsname{\href{S6_SERIES_GRIND_I.pdf\#nameddest=tag-03L5}{I.106 (p.~375)} }
\expandafter\def\csname SRef@inv:eq:atlas-affine-lift\endcsname{\href{S6_SERIES_GRIND_I.pdf\#nameddest=tag-03L6}{I.107 (p.~376)} }
\expandafter\def\csname SRef@inv:eq:atlas-cusp-factor\endcsname{\href{S6_SERIES_GRIND_I.pdf\#nameddest=tag-03L7}{I.108 (p.~377)} }
\expandafter\def\csname SRef@inv:eq:atlas-cusp-unit\endcsname{\href{S6_SERIES_GRIND_I.pdf\#nameddest=tag-03L8}{I.109 (p.~377)} }
\expandafter\def\csname SRef@inv:eq:atlas-connection-inverse\endcsname{\href{S6_SERIES_GRIND_I.pdf\#nameddest=tag-03L9}{I.110 (p.~377)} }
\expandafter\def\csname SRef@inv:eq:atlas-cusp-connection\endcsname{\href{S6_SERIES_GRIND_I.pdf\#nameddest=tag-03LA}{I.111 (p.~377)} }
\expandafter\def\csname SRef@inv:eq:atlas-finite-connection\endcsname{\href{S6_SERIES_GRIND_I.pdf\#nameddest=tag-03LB}{I.112 (p.~378)} }
\expandafter\def\csname SRef@inv:eq:atlas-averaged-connection\endcsname{\href{S6_SERIES_GRIND_I.pdf\#nameddest=tag-03LC}{I.113 (p.~378)} }
\expandafter\def\csname SRef@inv:eq:atlas-cocycle\endcsname{\href{S6_SERIES_GRIND_I.pdf\#nameddest=tag-03LD}{I.114 (p.~378)} }
\expandafter\def\csname SRef@inv:eq:atlas-weights\endcsname{\href{S6_SERIES_GRIND_I.pdf\#nameddest=tag-03LE}{I.115 (p.~379)} }
\expandafter\def\csname SRef@inv:eq:atlas-weight-derivative\endcsname{\href{S6_SERIES_GRIND_I.pdf\#nameddest=tag-03LF}{I.116 (p.~379)} }
\expandafter\def\csname SRef@inv:eq:atlas-final-forms\endcsname{\href{S6_SERIES_GRIND_I.pdf\#nameddest=tag-03LG}{I.117 (p.~379)} }
\expandafter\def\csname SRef@inv:eq:atlas-final-coefficients\endcsname{\href{S6_SERIES_GRIND_I.pdf\#nameddest=tag-03LH}{I.118 (p.~380)} }
\expandafter\def\csname SRef@inv:C3.8:operator-comparison\endcsname{\href{S6_SERIES_GRIND_I.pdf\#nameddest=tag-03LI}{I.D.2 (p.~381)} }
\expandafter\def\csname SRef@iv:open:K111:original_auxiliary_arrays\endcsname{\href{S6_SERIES_GRIND_IV.pdf\#nameddest=tag-03LJ}{IV.12.1 (p.~1396)} }
\expandafter\def\csname SRef@maps:proof:D.OS.refined\endcsname{\href{S6_SERIES_GRIND_III.pdf\#nameddest=tag-03LK}{III.5.2 (p.~207)} }
\expandafter\def\csname SRef@maps:eq:D.OR1\endcsname{\href{S6_SERIES_GRIND_III.pdf\#nameddest=tag-03LL}{III.{D.OR1} (p.~207)} }
\expandafter\def\csname SRef@maps:eq:D.OR2\endcsname{\href{S6_SERIES_GRIND_III.pdf\#nameddest=tag-03LM}{III.{D.OR2} (p.~207)} }
\expandafter\def\csname SRef@maps:eq:D.OR3\endcsname{\href{S6_SERIES_GRIND_III.pdf\#nameddest=tag-03LN}{III.{D.OR3} (p.~207)} }
\expandafter\def\csname SRef@maps:eq:D.OR4\endcsname{\href{S6_SERIES_GRIND_III.pdf\#nameddest=tag-03LO}{III.{D.OR4} (p.~207)} }
\expandafter\def\csname SRef@maps:eq:D.OR5\endcsname{\href{S6_SERIES_GRIND_III.pdf\#nameddest=tag-03LP}{III.{D.OR5} (p.~208)} }
\expandafter\def\csname SRef@maps:eq:D.OR6\endcsname{\href{S6_SERIES_GRIND_III.pdf\#nameddest=tag-03LQ}{III.{D.OR6} (p.~208)} }
\expandafter\def\csname SRef@maps:eq:D.OR7\endcsname{\href{S6_SERIES_GRIND_III.pdf\#nameddest=tag-03LR}{III.{D.OR7} (p.~208)} }
\expandafter\def\csname SRef@maps:eq:D.OR8\endcsname{\href{S6_SERIES_GRIND_III.pdf\#nameddest=tag-03LS}{III.{D.OR8} (p.~208)} }
\expandafter\def\csname SRef@maps:eq:D.OR9\endcsname{\href{S6_SERIES_GRIND_III.pdf\#nameddest=tag-03LT}{III.{D.OR9} (p.~208)} }
\expandafter\def\csname SRef@maps:eq:D.OR10\endcsname{\href{S6_SERIES_GRIND_III.pdf\#nameddest=tag-03LU}{III.{D.OR10} (p.~209)} }
\expandafter\def\csname SRef@maps:eq:D.OR11\endcsname{\href{S6_SERIES_GRIND_III.pdf\#nameddest=tag-03LV}{III.{D.OR11} (p.~209)} }
\expandafter\def\csname SRef@maps:eq:D.OR12\endcsname{\href{S6_SERIES_GRIND_III.pdf\#nameddest=tag-03LW}{III.{D.OR12} (p.~210)} }
\expandafter\def\csname SRef@maps:eq:D.OR13\endcsname{\href{S6_SERIES_GRIND_III.pdf\#nameddest=tag-03LX}{III.{D.OR13} (p.~210)} }
\expandafter\def\csname SRef@maps:eq:D.OR14\endcsname{\href{S6_SERIES_GRIND_III.pdf\#nameddest=tag-03LY}{III.{D.OR14} (p.~211)} }
\expandafter\def\csname SRef@inv:E2.5\endcsname{\hyperlink{tag-03LZ}{II.\ref*{inv:E2.5}}}
\expandafter\def\csname SRef@invII:eq:E12-node\endcsname{\hyperlink{tag-03M0}{II.\ref*{invII:eq:E12-node}}}
\expandafter\def\csname SRef@invII:eq:E12-jets\endcsname{\hyperlink{tag-03M1}{II.\ref*{invII:eq:E12-jets}}}
\expandafter\def\csname SRef@invII:eq:E12-open\endcsname{\hyperlink{tag-03M2}{II.\ref*{invII:eq:E12-open}}}
\expandafter\def\csname SBack@cdp:st:A.1.1\endcsname{\hyperlink{step-II-A-1}{proves step 1 of Theorem A}}
\expandafter\def\csname SBack@cdp:st:B1\endcsname{\hyperlink{step-II-A-2}{proves step 2 of Theorem A}}
\expandafter\def\csname SBack@cdp:st:A.4.7\endcsname{\hyperlink{step-II-A-3}{proves step 3 of Theorem A}}
\expandafter\def\csname SBack@cdp:st:A.4.9\endcsname{\hyperlink{step-II-A-3}{proves step 3 of Theorem A}}
\expandafter\def\csname SBack@cdp:st:C4f\endcsname{\hyperlink{step-II-A-4}{proves step 4 of Theorem A}}
\expandafter\def\csname SBack@cdp:st:E12\endcsname{\hyperlink{step-II-A-5}{proves step 5 of Theorem A}}
\expandafter\def\csname SBack@cdp:st:E17\endcsname{\hyperlink{step-II-A-5}{proves step 5 of Theorem A}}
\expandafter\def\csname SBack@cdp:st:D.1.2\endcsname{\hyperlink{step-II-B-1}{proves step 1 of Theorem B}}
\expandafter\def\csname SBack@cdp:st:C.3.1\endcsname{\hyperlink{step-II-B-2}{proves step 2 of Theorem B}}
\expandafter\def\csname SBack@cdp:st:C.3.7\endcsname{\hyperlink{step-II-B-3}{proves step 3 of Theorem B}}
\expandafter\def\csname SBack@cdp:st:C.3.8\endcsname{\hyperlink{step-II-B-3}{proves step 3 of Theorem B}}
\expandafter\def\csname SBack@cdp:st:C.3.9\endcsname{\hyperlink{step-II-B-3}{proves step 3 of Theorem B}}
\expandafter\def\csname SBack@cdp:st:C.3.10\endcsname{\hyperlink{step-II-B-3}{proves step 3 of Theorem B}}
\expandafter\def\csname SBack@cdp:st:C.3.12\endcsname{\hyperlink{step-II-B-3}{proves step 3 of Theorem B}}
\expandafter\def\csname SBack@cdp:st:C.3.6\endcsname{\hyperlink{step-II-B-4}{proves step 4 of Theorem B}}
\expandafter\def\csname SBack@cdp:st:D.2.7\endcsname{\hyperlink{step-II-B-4}{proves step 4 of Theorem B}; \hyperlink{step-II-C-6}{proves step 6 of Theorem C}}
\expandafter\def\csname SBack@night:P4.58\endcsname{\hyperlink{step-II-C-1}{proves step 1 of Theorem C}}
\expandafter\def\csname SBack@cdp:st:C.2.MH.2\endcsname{\hyperlink{step-II-C-1}{proves step 1 of Theorem C}}
\expandafter\def\csname SBack@cdp:st:C.2.INT.1\endcsname{\hyperlink{step-II-C-2}{proves step 2 of Theorem C}}
\expandafter\def\csname SBack@cdp:st:C.2.INT.2\endcsname{\hyperlink{step-II-C-2}{proves step 2 of Theorem C}}
\expandafter\def\csname SBack@cdp:st:C.2.INT.3\endcsname{\hyperlink{step-II-C-3}{proves step 3 of Theorem C}}
\expandafter\def\csname SBack@cdp:st:C.2.G.1\endcsname{\hyperlink{step-II-C-4}{proves step 4 of Theorem C}}
\expandafter\def\csname SBack@cdp:st:D.2.3\endcsname{\hyperlink{step-II-C-5}{proves step 5 of Theorem C}}
\expandafter\def\csname SBack@cdp:st:D.2.5\endcsname{\hyperlink{step-II-C-5}{proves step 5 of Theorem C}}
\expandafter\def\csname SBack@cdp:st:C.2.UN.1\endcsname{\hyperlink{step-II-D-1}{proves step 1 of Theorem D}; \hyperlink{step-II-D-2}{proves step 2 of Theorem D}}
\expandafter\def\csname SBack@cdp:st:C.2.UN.4\endcsname{\hyperlink{step-II-D-3}{proves step 3 of Theorem D}}
\expandafter\def\csname SBack@cdp:st:C.2.UN.14\endcsname{\hyperlink{step-II-D-4}{proves step 4 of Theorem D}}
\expandafter\def\csname SBack@cdp:st:C.2.UN.2\endcsname{\hyperlink{step-II-D-5}{proves step 5 of Theorem D}}
\expandafter\def\csname SBack@cdp:st:C5b\endcsname{\hyperlink{step-II-E-1}{proves step 1 of Theorem E}; \hyperlink{step-II-E-2}{proves step 2 of Theorem E}}
\expandafter\def\csname SBack@cdp:st:C5c\endcsname{\hyperlink{step-II-E-3}{proves step 3 of Theorem E}; \hyperlink{step-II-E-4}{proves step 4 of Theorem E}}
\expandafter\def\csname SBack@cdp:st:E8\endcsname{\hyperlink{step-II-E-5}{proves step 5 of Theorem E}}
\expandafter\def\csname SBack@cdp:st:E18\endcsname{\hyperlink{step-II-E-5}{proves step 5 of Theorem E}}
\expandafter\def\csname SBack@cdp:st:E5\endcsname{\hyperlink{step-II-F-1}{proves step 1 of Theorem F}}
\expandafter\def\csname SBack@cdp:st:E6\endcsname{\hyperlink{step-II-F-2}{proves step 2 of Theorem F}}
\expandafter\def\csname SBack@cdp:st:E13\endcsname{\hyperlink{step-II-F-3}{proves step 3 of Theorem F}}
\expandafter\def\csname SBack@cdp:st:E14\endcsname{\hyperlink{step-II-F-3}{proves step 3 of Theorem F}}
\expandafter\def\csname SBack@cdp:st:E15\endcsname{\hyperlink{step-II-F-4}{proves step 4 of Theorem F}}
\expandafter\def\csname SBack@push:src:E.2.23\endcsname{\hyperlink{step-II-F-5}{proves step 5 of Theorem F}}
\expandafter\def\csname SBack@push:src:E.2.24\endcsname{\hyperlink{step-II-F-6}{proves step 6 of Theorem F}}
\expandafter\def\csname SBack@inv:E2.5\endcsname{\hyperlink{step-II-F-7}{proves step 7 of Theorem F}}
\expandafter\def\csname SBack@inv:E2.2\endcsname{\hyperlink{step-II-F-7}{proves step 7 of Theorem F}}
\expandafter\def\csname SBack@ii:LR:covering\endcsname{\hyperlink{step-II-G-1}{proves step 1 of Theorem G}; \hyperlink{step-II-G-2}{proves step 2 of Theorem G}; \hyperlink{step-II-G-3}{proves step 3 of Theorem G}; \hyperlink{step-II-G-4}{proves step 4 of Theorem G}}

\begin{document}
\title{The $(3,4,\infty)$ family of complex six-spheres\\[.5em]II. Torus-fibred threefolds and the sign theorem}
\author{AI-produced mathematical text}
\date{23 November 2026}
\maketitle
\begin{abstract}
We study holomorphic fibrations of compact complex threefolds with
vanishing first and second Betti numbers. For a torus general fibre,
every degeneration has nonnegative Euler number, with arbitrary
multiplicities and support singularities, and the total Euler number
is zero or two. We describe the equality supports and the equality
forced by identity monodromy, through integral component counts,
local parity, conductor maps and compact constructions. A local
nodal family gives a counterexample to a vanishing implication that
omits the cotangent-torsion connecting map. Under the integral
Campana--Demailly--Peternell hypotheses we prove the exact
normalization-Picard criterion for generic tangent-cohomology
vanishing outside degree one. We compute tangent twists and deformation spaces,
prove persistence of the fibration in the stated families, and
determine the two-turn covering parameter of the Liu--Rollenske
family. We construct a compact torus fibration with three doubled
nodal supports and elliptic ruled normalizations. A flat-kernel and
conic two-jet calculation gives a nonempty open family of
twelve-tooth smoothings with kernel
$\mathcal O(-8)\oplus\mathcal O(-9)^2$ and Hodge numbers
$h^{1,1}=49$, $h^{1,2}=48$. We state the remaining splitting
problem for general and special comb members.
\end{abstract}
\clearpage
\tableofcontents
\clearpage
\section{Results and their proofs in outline}\label{sec:results}
\subsection{The series and this paper}
This series studies the $(3,4,\infty)$ family of torus-fibred complex
structures on $S^6$ \citep[Theorem~8.1, pp.~63--64]{AC26}
\citep[Theorem~4.1, p.~12]{Engel2026}, its periods, maps and smooth
monodromy, and the higher-dimensional complex manifolds built from it.
\begin{center}\small
\begin{tabularx}{\textwidth}{@{}lXl@{}}\toprule
Paper & Main result & Uses\\\midrule
I & Exact periods, moduli, fibre strata and the two ends & ---\\
II & Torus-fibred threefolds with $b_1=b_2=0$: the Campana--Demailly--Peternell argument and the sign theorem for all degenerations & I, IV\\
III & Maps from the six-spheres and the manifolds they give & I, II, IV, V\\
IV & Smooth monodromy in $\Theta_7$ & I, II, III, V\\
V & Complex fourfolds on $3\mathbb{HP}^2\#(-\mathbb{HP}^2)$ and a resolved fibre product & I, III, IV\\
\bottomrule\end{tabularx}\end{center}
\subsection{The main theorems}
\hypertarget{tag-03K1}{}\label{series:II:main-theorems}
\begin{definition}[Integral Campana--Demailly--Peternell hypotheses]
\hypertarget{tag-P201}{}\label{def:II:integral-cdp}
\StatementRecord{def:II:integral-cdp}{proved}{P201}{}
We write $F_b=f^{-1}(b)$ for a full Cartier fibre, including its
multiplicities, and $S_b=(F_b)_{\rm red}$ for its reduced support.
The integral CDP hypotheses are that $Y$ is a connected
smooth compact complex threefold, $H^1(Y,\Z)=H^2(Y,\Z)=0$,
$a(Y)=1$, and $f:Y\to\PP^1$ is its holomorphic algebraic reduction
with connected fibres.
\end{definition}
\hypertarget{main-II-A}{}
\begin{maintheorem}[The cotangent term and the exact generic-vanishing criterion]
\hypertarget{tag-03K2}{}\label{series:II:A}
\StatementRecord{series:II:A}{proved}{03K2}{\SeriesRef{cdp:st:A.4.9}, \SeriesRef{cdp:st:E12}, \SeriesRef{cdp:st:E17}, \SeriesRef{cdp:st:A.1.1}, \SeriesRef{cdp:st:B1}, \SeriesRef{cdp:st:C4f}}
Let $S$ be a reduced hypersurface in a smooth complex threefold $Y$,
let $A$ be a holomorphic line bundle on $S$, and put
\[
 N^*=\OO_Y(-S)|_S,\quad V=\Omega_Y^1|_S,\quad
 T=\operatorname{Tors}\Omega_S^1,\quad Q=\Omega_S^1/T,
 \quad K=\ker(V\to Q).
\]
The conormal map and its inverse image of torsion give
\[
 0\to N^*\to K\to T\to0,
 \qquad 0\to K\to V\to Q\to0.
\]
Writing $\delta_A:H^0(T\otimes A)\to H^1(N^*\otimes A)$
for the first connecting map, the two vanishings
$H^0(N^*\otimes A)=H^0(Q\otimes A)=0$ identify
$H^0(V\otimes A)$ canonically with $\ker\delta_A$.
In the family
\[
 \mathcal C=\{y^2z=x^2(x+z)+tz^3\}\subset\PP^2\times\Delta,
 \qquad Y=\mathcal C\times\PP^1,
 \qquad 0<\operatorname{radius}(\Delta)<4/27,
\]
the disc bound excludes the second singular value $t=-4/27$ of the same cubic pencil. The central support and every nontrivial normalization-gluing line
$A_\lambda$, $\lambda\in\C^*\setminus\{1\}$, have both displayed
endpoint vanishings and $h^0(V\otimes A_\lambda)=1$.

For a map satisfying the integral CDP hypotheses, a general
$L\in\Pic(Y)$ has $H^q(Y,T_Y\otimes L)=0$ for $q\ne1$
exactly when the normalization pullback
\[
 \Pic(Y)\longrightarrow\Pic(T_S),\qquad
 T_S\xrightarrow{\eta_S}S,
\]
is nonzero for every irreducible nonnormal critical support $S$.
A zero pullback gives $H^2(Y,T_Y\otimes L)\ne0$ for every $L$.
This criterion permits arbitrary full-fibre multiplicities and
arbitrary support singularities. In its vanishing case,
$\chi(Y,T_Y\otimes L)\le0$ and $c_3(Y)\le0$.

The following direct-image and divisor-twisting identities also hold. For an elliptic curve $E$ with origin, its normalized
Poincar\'e line $P$ on $E\times E$, the projection
$g:E\times E\times\PP^1\to E$ to the second factor, and
$V=P\boxtimes\OO_{\PP^1}(-2)$, one has
$g_*V=R^1g_*V=0$ and $R^2g_*V=\C_0$, of length one.
For an effective Cartier divisor $D$ on a compact complex space
$Z$ and locally free $V$, the decreasing spaces
$H^0(Z,V(-kD))\subset H^0(Z,V)$ stabilize at the sections whose
support is disjoint from $|D|$. They eventually vanish when $Z$
has no embedded associated components and $D$ meets every
irreducible component.

\end{maintheorem}

\textit{Mechanism.} The conormal filtration includes the torsion
connecting map, and normalization Picard pullback decides whether
its conductor section survives every twist.
\textit{Relation to literature.} The normalization criterion extends
\cite[Lemma~10.3, Theorem~10.5 and \S10.4, pp.~87--90]{AC26}
and corrects \cite[Theorem~2.2, p.~680]{CDP2020}; the two auxiliary
identities correct \cite[Proposition~1.1, p.~78; Lemma~1.5, pp.~79--80]{CDP98}.


\hypertarget{main-II-B}{}
\begin{maintheorem}[Component count and the smooth general fibre]
\hypertarget{tag-03K3}{}\label{series:II:B}
\StatementRecord{series:II:B}{proved}{03K3}{\SeriesRef{cdp:st:D.1.2}, \SeriesRef{cdp:st:C.3.1}, \SeriesRef{cdp:st:C.3.11}}
Let $f:Y\to C$ be a proper holomorphic surjection with connected
fibres from a connected smooth compact complex threefold to a
smooth compact curve. Delete a nonempty finite set $\Delta$
containing all critical values, put $s=|\Delta|$, $g=g(C)$,
let $r$ count the irreducible components over $\Delta$, and set
\[
 t'=\dim_\Q H_1(F,\Q)_{\pi_1(C\setminus\Delta)}
\]
for a smooth fibre $F$. With $b_2(Y)=0$,
\[
 r=2g+s-1+t'-b_1(Y).
\]
Under $H^1(Y,\Z)=H^2(Y,\Z)=0$, the base is $\PP^1$,
$r=s-1+t'$, and the integral coinvariant quotient
$H_1(F,\Z)_{\pi_1(C\setminus\Delta)}$ is free abelian.
For a holomorphic algebraic reduction with $a(Y)=1$ and
$b_1(Y)=b_2(Y)=0$, every smooth fibre satisfies
\[
 c_1(K_F)=0\text{ in }H^2(F,\Q),\qquad
 c_2(F)=0,\qquad \chi(F,\OO_F)=0.
\]
Its very general fibre is a complex two-torus or a primary or
secondary Hopf surface of algebraic dimension zero. Each of these
three types occurs in the constructions stated in the body.

\end{maintheorem}

\textit{Mechanism.} The exact homology sequence of the complement
identifies the monodromy coinvariants. Relative Albanese and Picard
comparisons, surface classification and period calculations then
determine the possible general fibres.
\textit{Relation to literature.} The rational component formula is
\cite[\S10.5, p.~91]{AC26}; its integral form and the smooth-fibre
classification are \cref{cdp:st:D.1.2,cdp:st:C.3.1,cdp:st:C.3.11}.


\hypertarget{main-II-C}{}
\begin{maintheorem}[The sign for every degeneration and its equality criterion]
\hypertarget{tag-03K4}{}\label{series:II:C}
\StatementRecord{series:II:C}{proved}{03K4}{\SeriesRef{cdp:st:D.2.7}, \SeriesRef{cdp:st:C.2.G.1}, \SeriesRef{night:P4.6}, \SeriesRef{cdp:st:D.2.5}}
Let $Y$ be a connected smooth compact complex threefold with
$b_1(Y)=b_2(Y)=0$, and let $f:Y\to B$ be a proper surjective
holomorphic map with connected fibres to a smooth compact curve.
Suppose the smooth general fibre is a complex two-torus. For
arbitrary positive full-fibre multiplicities and arbitrary
reduced-support singularities,
\[
 B\simeq\PP^1,\qquad
 c_3(Y)=e(Y)\in\{0,2\},\qquad b_3(Y)=2-c_3(Y).
\]
Every reduced support has nonnegative Euler number. A singular
irreducible support with rational normalization resolution has
Euler number two; every other irreducible support has Euler
number zero. A normal-crossing support has Euler number equal
to its number of triple points.

The equality $c_3(Y)=0$ holds exactly when every support has its
corresponding description below. A smooth irreducible support
has a finite unramified torus cover. A singular irreducible support
has smooth minimal elliptic ruled or smooth minimal Hopf
normalization. A reducible support has normalization a disjoint
union of smooth minimal Hopf or non-K\"ahler elliptic surfaces;
the reduced conductor image and inverse image are disjoint unions
of smooth elliptic curves, and their finite branch maps are
unramified. For $Y_b=f^*(b)=\sum_i m_iS_i$ and $S=(Y_b)_{\rm red}$,
the normalization is $\nu:\coprod_iT_i\to S$ and its map
to the full fibre is $j\circ\nu$, where $j:S\hookrightarrow Y_b$
is the reduction inclusion. The scheme conductor is defined by
$\operatorname{Ann}_{\OO_S}(\nu_*\OO_{\coprod_iT_i}/\OO_S)$.

For identity monodromy on $H_1(F,\Q)$ one always has
$c_3(Y)=0$ and $b_3(Y)=2$; the compact construction in
\cref{cdp:st:D.2.7} attains this value with
$H_1(Y,\Z)=H_2(Y,\Z)=0$.

\end{maintheorem}

\textit{Mechanism.} The normalization and conductor yield a sum of
nonnegative local Euler contributions, while identity monodromy gives
the opposite inequality through bounded logarithmic periods.
\textit{Relation to literature.} The sign and equality criterion in
\cref{cdp:st:C.2.G.1} extend the standard-family Euler value of
\cite[Theorem~7.22, p.~60]{AC26} and
\cite[Proposition~7.1, p.~17]{Engel2026}.


\hypertarget{main-II-D}{}
\begin{maintheorem}[Local parity and prescribed Euler numbers]
\hypertarget{tag-03K5}{}\label{series:II:D}
\StatementRecord{series:II:D}{proved}{03K5}{\SeriesRef{cdp:st:C.2.UN.1}, \SeriesRef{cdp:st:C.2.UN.2}, \SeriesRef{cdp:st:C.2.UN.4}, \SeriesRef{cdp:st:C.2.UN.14}}
Let $f:Y\to B$ be a proper holomorphic map from a smooth,
possibly noncompact, complex threefold to a smooth complex curve,
with connected fibres and smooth general fibre a complex two-torus.
Every reduced fibre support has even Euler number, with arbitrary
singularities and multiplicities.
For every fixed $b\in B$ and every $k\in\Z$, a finite composition
of blowups in smooth compact centres over $b$ gives
\[
 \begin{gathered}
 \pi:\widetilde Y\to Y,\qquad e((\widetilde Y_b)_{\rm red})=2k,\\
 \pi:\widetilde Y|_{B\setminus\{b\}}
       \xrightarrow{\ \simeq\ }Y|_{B\setminus\{b\}}.
 \end{gathered}
\]
The strict transforms retain their original full-fibre
multiplicities. The operation preserves the finite multiset of
bimeromorphism classes of non-uniruled normalization resolutions,
through the induced normalization and blowdown maps. Every such
filling has zero-Euler refinements with arbitrarily many components.
Every even integer also occurs in a smooth projective total space.

\end{maintheorem}

\textit{Mechanism.} Mapping-torus duality and the symmetric wedge
pairing force evenness, and exceptional-divisor replacement gives Euler
changes $+2$ at a point and $-4$ along a plane quartic after a point
blowup.
\textit{Relation to literature.} The parity and refinement assertions
are \cref{cdp:st:C.2.UN.1,cdp:st:C.2.UN.4,cdp:st:C.2.UN.14},
proved by mapping-torus duality and blowup formulas.


\hypertarget{main-II-E}{}
\begin{maintheorem}[Twisted tangent cohomology of the standard family]
\hypertarget{tag-03K6}{}\label{series:II:E}
\StatementRecord{series:II:E}{proved}{03K6}{\SeriesRef{cdp:st:C5b}, \SeriesRef{cdp:st:C5c}, \SeriesRef{cdp:st:E8}, \SeriesRef{cdp:st:E18}}
Let $X=X_c$ be a member of the standard $(3,4,\infty)$ family,
with the fixed source, physical line $M$, affine attachment data,
period normalization and fixed reference cusp fan of Paper~I.
Write $M_z$ for the ambient Picard line with parameter $z\in\C$
in that normalization. For $n\in\Z$, define
\[
 \begin{gathered}
 0\le r_1(n)<3,\ r_1(n)\equiv n\pmod3,\qquad
 0\le r_2(n)<4,\ r_2(n)\equiv3n\pmod4,\\
 d(n)=\frac{n-4r_1(n)-3r_2(n)}{12},\quad
 u(n)=\max\{d(n)+1,0\},\quad v(n)=\max\{-d(n)-1,0\},\\
 \epsilon_n=0\ (3\mid n),\qquad \epsilon_n=1\ (3\nmid n).
 \end{gathered}
\]
Then
\[
\begin{aligned}
 h^0(T_X\otimes M_n)&=u(n)+u(n-1),\\
 h^1(T_X\otimes M_n)&=v(n)+v(n-1)+u(n)+u(n-1)+u(n+2),\\
 h^2(T_X\otimes M_n)&=v(n)+v(n-1)+v(n+2)+u(n+6)+\epsilon_n,\\
 h^3(T_X\otimes M_n)&=v(n+6).
\end{aligned}
\]
The map to $H^2(\mathrm dP_6,K_{\mathrm dP_6})=\C$ has rank one
exactly for $n\in\{-6,-3,-2\}\cup\Z_{\ge0}$, and rank zero
otherwise. At $n=-1$ its source has dimension one.

For $z\notin\Z$, let $z_1,z_2$ be the marked fixed lifts of
$p_1=0,p_2=1$ on the $(3,4,\infty)$ uniformizing half-plane, with
$\tau(z_1)=e^{\pi i/3}$ and $\tau(z_2)=i$ and with the cubic and
quartic root coordinates of \cref{def:II:marked-conventions}.
Use the marked period functions $\tau,\mu,\beta_c$ and let
$\mathcal R_z$ be the regular base points where integers
$(n,p,q,r)$ satisfy
\[
 q\tau+r\mu=p,\qquad z=n-6q\mu-r\beta_c.
\]
Put $\rho=\tau(z_1)$,
$\Lambda_1=\Z+(3\beta_c(z_1)+4-2\rho)\Z$ and
$\Lambda_2=\Z+(2\beta_c(z_2)+3-3i)\Z$.
Define $\ell_1(z)=1+\epsilon_n$ at
$z=n+k(3\beta_c(z_1)+4-2\rho)$, $k\ne0$, and zero elsewhere.
Define $\ell_2(z)=1$ at
$z=n+k(2\beta_c(z_2)+3-3i)$, $k\ne0$,
$n\equiv1,2\pmod4$, and zero elsewhere. The set $\mathcal R_z$
is finite. With $J_z=2|\mathcal R_z|+\ell_1(z)+\ell_2(z)$,
\[
 h^\bullet(X,T_X\otimes M_z)=(0,J_z,1+J_z,0),
\]
and its map to the same normalization target is surjective.

For the symmetric period matrix
$\Omega=\left(\begin{smallmatrix}\tau&\mu\\\mu&\beta/6\end{smallmatrix}\right)$,
set $\beta=\beta_c=\beta^{\rm part}+c$. With Jacobi theta periods
$1,\tau$ and
$\beta/6=(i/\pi)\log(\vartheta_1(\mu,\tau)/\eta(\tau))+\mathrm{constant}$,
put $h=\partial_\mu\log\vartheta_1$. The exact identity is
\[
 \det\frac{d\Omega}{d\tau}
 =-\left(\mu'-\frac{i}{2\pi}h\right)^2
   -\frac{\wp(\mu,\tau)}{4\pi^2}.
\]
On the original branch $\wp(\mu,\tau)=0$ it equals
$-g_2^2/(36\pi^2g_3)$. The rank on the smooth root base is one
at the order-three point and two at every other finite point;
the determinant has order two there, and the logarithmic cusp
tensor has rank two.

\end{maintheorem}

\textit{Mechanism.} Root-curve direct images and exact local
resonance modules determine the normalization map together with
its source dimension.
\textit{Relation to literature.} The formulas refine
\cite[Theorem~10.5 and Corollary~10.6, pp.~88--89]{AC26}; the
untwisted tuple is \cite[Theorem~6.12, pp.~81--82]{DGKXv1},
arXiv version~1.


\hypertarget{main-II-F}{}
\begin{maintheorem}[Fibrations, deformation spaces and comb Hodge numbers]
\hypertarget{tag-03K7}{}\label{series:II:F}
\StatementRecord{series:II:F}{proved}{03K7}{\SeriesRef{inv:E2.4}, \SeriesRef{cdp:st:E6}, \SeriesRef{cdp:st:E13}, \SeriesRef{cdp:st:E14}, \SeriesRef{cdp:st:E15}, \SeriesRef{push:src:E.2.23}, \SeriesRef{push:src:E.2.24}, \SeriesRef{inv:E2.2}, \SeriesRef{inv:E2.5}}
Let $g:Y\to\PP^1$ be proper and surjective with connected fibres
from a smooth connected compact complex threefold, and let its
smooth general fibre be a complex two-torus. Rank two of its
Kodaira--Spencer tensor at one regular value ensures that every
proper deformation $\mathcal Y\to(S,0)$, including one over a
nonreduced complex analytic germ, admits after shrinking $S$ a
map $G:\mathcal Y\to\PP^1\times S$ extending $g$. Every $G_s$
is surjective with connected fibres and smooth general fibre a
two-torus. Under $b_2(Y)=0$ its sufficiently small deformations
have algebraic dimension one.

For the constructed Viaclovsky two-parameter family, with
$\mu$ close to zero and nonzero sufficiently small scale
$\lambda$, and for the Liu--Rollenske family at sufficiently
small nonzero scale, the tangent tuple is $(1,2,2,0)$ and the
Kuranishi germ is a smooth reduced surface. The displayed
Viaclovsky family is locally miniversal; the second
Liu--Rollenske direction splits its cuspidal fibre into two
rank-one nodal fibres. For every positive-comb manifold $X_m$
in the construction of \cite[Theorem~1.1, p.~2]{Viaclovsky2026Combs}, in particular
for its constructed $m\ge12$ range,
\[
 h^\bullet(X_m,T_{X_m})=(2m+1,8m+1,6m+1,0),
\]
and the Kuranishi germ is smooth and reduced of dimension $8m+1$.
Its local period-section family retains the affine attachment
classes, three critical values, multiplicities $3,4$, and the
hexagonal fibre type.
For the original marked section on $\mathbb P(3,4)$, there is a
nonempty open family obtained by embedded smoothings of twelve
admissible deformed twistor conics at distinct ordinary points,
fixing the three marked period planes, whose cup kernel is
\[
 K_{12}=\mathcal O_{\mathbb P^1}(-8)
       \oplus\mathcal O_{\mathbb P^1}(-9)^2,
 \qquad \operatorname{coker}\kappa_1=0.
\]
Here $K_{12}$ and $\kappa_1$ are the sheaf and cup map of
\eqref{invII:eq:E.6}. These members have
$h^{1,1}=49$ and $h^{1,2}=48$.

In any proper smooth holomorphic family over a connected reduced
complex analytic base containing a constructed sphere member with
coprime ordered multiplicities $(3k,4l)$, every fibre is a
constructed member with that same pair. Two such members are
deformation equivalent through finite chains precisely when their
ordered pairs agree. On the fixed smooth sphere, the structures
with one such pair form an open and closed subset of the
$C^\infty$ space of integrable complex structures, and have the
same separation in the quotients by $\operatorname{Diff}_0(S^6)$
and by orientation-preserving diffeomorphisms.

\end{maintheorem}

\textit{Mechanism.} Fibrewise cup products remove map obstructions,
period sections realize tangent directions, propagation preserves the
ordered multiplicities, and the cyclic two-jet quotient of the flat
comb kernel, conic frame maps and tooth collapse give the stated
balanced cohomology.
\textit{Relation to literature.} The comb germ answers
\cite[Section~10(1), p.~33]{Viaclovsky2026Combs}; the Hodge values
address its Section~10(2) for the stated open twelve-tooth family.
The ordered-pair deformation classes answer
\cite[Section~1.1(1), p.~2]{Viaclovsky2026Log}, and the vector-field
value corrects \cite[Corollary~1.3, preprint p.~2]{HKP2000}.


\hypertarget{main-II-G}{}
\begin{maintheorem}[The Liu--Rollenske covering parameter]
\hypertarget{tag-03K8}{}\label{series:II:G}
\StatementRecord{series:II:G}{proved}{03K8}{\SeriesRef{ii:LR:covering}}
Fix the exact scalar $u$ in the Liu--Rollenske action $B_u=uB_0$,
the good fourth root at the $III^*$ source fibre, its fixed
root-line character, the logarithmic vector $a=e_4/4$, and the
fixed periodic cusp fan with its source and line attachment
maps. The construction gives a holomorphic family on the universal
logarithmic cover of a sufficiently small punctured $u$-disc.
One positive $u$-turn reaches a threefold having a different
biholomorphism class. This family descends through
$u=w^2$, and this degree two is minimal. Every constant root-line character
in this fixed construction has the same covering degree.

\end{maintheorem}

\textit{Mechanism.} One scalar turn changes the multiple fibre's
Albanese period by a half-integer, and the two-turn source-line lift
has the required closing norm on every fixed cap.
\textit{Relation to literature.} The minimal parameter covering in
\cref{ii:LR:covering} is determined for the scalar action and good-root
data of \cite[equation~(2.3), p.~4; Lemma~2.7 and Corollary~2.8,
pp.~6--7; Proposition~3.6 and Definition~3.7, pp.~9--10]{LiuRollenske2026},
arXiv version~2.


\subsection{The proofs in outline}
\hypertarget{tag-03K9}{}\label{series:II:outline}

\hypertarget{outline-II-A}{}
\subsubsection*{Theorem A}
\begin{enumerate}[label=\arabic*.,leftmargin=*]
\item\hypertarget{step-II-A-1}{} The conormal map is injective at
each generic smooth point of the reduced hypersurface. Its inverse
image of cotangent torsion gives the two exact sequences, and their
cohomology identifies the ambient section space with the stated
connecting-map kernel
(\cref{cdp:st:A.1.1}).
\item\hypertarget{step-II-A-2}{} The node has normalization
$u\mapsto(u^2-1,u(u^2-1))$ and identifies $u=1,-1$. Its gluing-line
complex is $\C\xrightarrow{1-\lambda}\C$, while the inverse
image of torsion is the pushforward of the trivial normalization
line. Taking the product with $\PP^1$ gives the asserted
counterexample dimensions
(\cref{cdp:st:B1}).
\item\hypertarget{step-II-A-3}{} Transgression gives unitary
infinite-order twists on every fibre support. The explicit
fibre-group calculation and neighbourhood comparison make
degree-one local-system cohomology vanish. The canonical
injection then kills the twisted torsion-free forms
(\cref{cdp:st:A.4.7},
\cref{cdp:st:A.4.9}).
\item\hypertarget{step-II-A-4}{} Full Cartier filtrations reduce
generic ambient vanishing to irreducible nonnormal supports.
Normalization Picard pullback determines persistence of their
conductor section, yielding exactly the displayed global
criterion (\cref{cdp:st:C4f}).
\item\hypertarget{step-II-A-5}{} The Poincar\'e transform and
cohomology of $\OO_{\PP^1}(-2)$ produce the length-one higher
direct image. Stabilization of the finite-dimensional section
spaces and the associated-component calculation give the exact
divisor-twisting statement
(\cref{cdp:st:E12},
\cref{cdp:st:E17}).
\end{enumerate}

\hypertarget{outline-II-B}{}
\subsubsection*{Theorem B}
\begin{enumerate}[label=\arabic*.,leftmargin=*]
\item\hypertarget{step-II-B-1}{} Lifts of the free generators of
$\pi_1(C\setminus\Delta)$ split the punctured bundle group.
Its abelianization is $\Z^{2g+s-1}\oplus H_1(F,\Z)_{\pi_1}$;
the compact-support sequence then gives the component count and
integral freeness (\cref{cdp:st:D.1.2}).
\item\hypertarget{step-II-B-2}{} Fibre adjunction and the tangent
sequence give the displayed Chern identities; Noether's formula
gives the coherent Euler number. Surface classification supplies
the finite necessary list (\cref{cdp:st:C.3.1}).
\item\hypertarget{step-II-B-3}{} Relative Albanese maps and the
ambient second-cohomology vanishing exclude bielliptic fibres.
Picard quotients and period comparisons exclude the Kodaira
cases. For Inoue fibres, the form $d\log\operatorname{Im}w$
has zero meridian periods and gives a nonzero class in $H^1(Y,\R)$.
Anticanonical relative linear systems make a very general Hopf
fibre have algebraic dimension zero
(\cref{cdp:st:C.3.7},
\cref{cdp:st:C.3.8},
\cref{cdp:st:C.3.9},
\cref{cdp:st:C.3.10},
\cref{cdp:st:C.3.12}).
\item\hypertarget{step-II-B-4}{} The quotient and filling maps
produce the primary and secondary Hopf types and a fixed
abelian general fibre (\cref{cdp:st:C.3.6},
\cref{cdp:st:D.2.7}).
\end{enumerate}

\hypertarget{outline-II-C}{}
\subsubsection*{Theorem C}
\begin{enumerate}[label=\arabic*.,leftmargin=*]
\item\hypertarget{step-II-C-1}{} Connected-fibre Leray makes the
base rational. The finite-character comparison for the same full
irreducible fibre and its normalization gives the values zero and
two; the normalization analysis gives their stated geometric forms.
For a normal-crossing support $S$, the divisor calculation applied
to $(X,F)=(Y,S)$ gives the triple-point count
(\cref{cdp:st:C.2.MH.2,cdp:st:C.2.G.1,night:P4.6}).
\item\hypertarget{step-II-C-2}{} An adjacent component pulls back
to a nonzero effective divisor with rational Chern class zero.
The surface and Picard comparisons leave the Hopf, non-K\"ahler
elliptic and Enoki possibilities and then exclude the Enoki
interfaces through the explicit local vanishing sphere
(\cref{cdp:st:C.2.INT.1},
\cref{cdp:st:C.2.INT.2}).
\item\hypertarget{step-II-C-3}{} Coherent conductor duality, curve
normalization and exact blowup counts give nonnegative summands
for every reducible support allowed by the preceding component classification. Equality forces the stated
minimal normalizations and unramified elliptic conductor maps
(\cref{cdp:st:C.2.INT.3}).
\item\hypertarget{step-II-C-4}{} Euler additivity gives
$e(Y)=\sum_b e(S_b)$. The nondegenerate skew pairing on $H^3$
and $e(Y)=2-b_3(Y)$ give the two values and the global equality
criterion (\cref{cdp:st:C.2.G.1}).
\item\hypertarget{step-II-C-5}{} For identity monodromy, a
persistent positive-index obstruction has a fixed irreducible
support. The finite \'etale root cover, twisted cotangent
comparison and Cartier-layer filtration exclude this obstruction
on a smooth support. For a singular support, Gauduchon degree
and the nonzero conductor make its normalization have geometric
genus zero; vertical resolution preserves that property on all
components. The logarithmic
period obstruction contradicts identity monodromy, giving
$c_3(Y)\le0$ and hence equality
(\cref{cdp:st:D.2.3},
\cref{cdp:st:D.2.5}).
\item\hypertarget{step-II-C-6}{} The three-piece
holomorphic filling and its integral Mayer--Vietoris maps give
$H_1=H_2=0$ and Euler number zero
(\cref{cdp:st:D.2.7}).
\end{enumerate}

\hypertarget{outline-II-D}{}
\subsubsection*{Theorem D}
\begin{enumerate}[label=\arabic*.,leftmargin=*]
\item\hypertarget{step-II-D-1}{} A full fibre tube $U$ retracts
onto the support, and its boundary $M$ is the nearby torus mapping
torus. Pair duality gives
$e(U)\equiv b_0(M)+b_1(M)+b_2(M)$ modulo two
(\cref{cdp:st:C.2.UN.1}).
\item\hypertarget{step-II-D-2}{} Wang identifies that parity with
$\dim\ker(\Lambda^2T-I)$. The wedge pairing and Cayley transform
make this dimension even for every orientation-preserving torus
monodromy (\cref{cdp:st:C.2.UN.1}).
\item\hypertarget{step-II-D-3}{} Replacing a smooth codimension-$r$
centre by its exceptional projective bundle changes support Euler
number by $(r-1)e(C)$. A point adds two; a plane quartic after a
point blowup gives net change minus two. Iteration gives every
prescribed value and cofinal zero-Euler models
(\cref{cdp:st:C.2.UN.4}).
\item\hypertarget{step-II-D-4}{} Strict transforms are
bimeromorphic on normalization resolutions and preserve their
base-function valuations. Every new exceptional surface is
explicitly ruled or rational, proving the multiset assertion
(\cref{cdp:st:C.2.UN.14}).
\item\hypertarget{step-II-D-5}{} The genus-two curve in the
projective product family and disjoint point blowups realize all
values $2(m-k)$ in compact projective threefolds
(\cref{cdp:st:C.2.UN.2}).
\end{enumerate}

\hypertarget{outline-II-E}{}
\subsubsection*{Theorem E}
\begin{enumerate}[label=\arabic*.,leftmargin=*]
\item\hypertarget{step-II-E-1}{} The stabilizer characters on the
root curve give the integral degrees $d(n)$, with direct-image
summands $d(n)$, $d(n-1)$, $d(n+2)$ and $d(n+6)$ and the
specified finite root contribution
(\cref{cdp:st:C5b}).
\item\hypertarget{step-II-E-2}{} Cohomology of each line on
$\PP^1$ gives the four displayed dimensions. The cusp
map has the stated rank locus, including its zero value at
$n=-1$ (\cref{cdp:st:C5b}).
\item\hypertarget{step-II-E-3}{} The exact translation-character
complex is acyclic off resonance. At a regular resonance its
local modules have two degree-one generators and a length-two
degree-two generator; finite-root characters give
$\ell_1,\ell_2$ (\cref{cdp:st:C5c}).
\item\hypertarget{step-II-E-4}{} The countable isolated resonance
loci and the Fourier contractions give a regular point with zero
fibre cohomology. The coherent direct-image support is therefore
a proper analytic subset of the compact base curve, hence finite.
The cusp adds one reduced degree-two skyscraper and its nonzero
map to the normalization target. Leray gives $(0,J_z,1+J_z,0)$
(\cref{cdp:st:C5c}).
\item\hypertarget{step-II-E-5}{} The identity-induced degree-six
isogeny gives the symmetric period matrix. The theta heat equation
and the eta derivative give the exact determinant with its
Weierstrass term; the original branch equation then gives its stated
rank locus (\cref{cdp:st:E8},
\cref{cdp:st:E18}).
\end{enumerate}

\hypertarget{outline-II-F}{}
\subsubsection*{Theorem F}
\begin{enumerate}[label=\arabic*.,leftmargin=*]
\item\hypertarget{step-II-F-1}{} At a smooth fibre the boundary
map sends $a\eta_1+b\eta_2$ to
$(bu_1-au_2)\otimes\eta_1\wedge\eta_2$. Rank two makes it
injective and kills the infinitesimal map obstruction over
$\PP^1$ (\cref{cdp:st:E5}).
\item\hypertarget{step-II-F-2}{} The same cup-product comparison
at every Artin order, followed by the relative graph construction,
gives the extending map over the analytic germ, with its
connected torus fibres and stated algebraic dimension
(\cref{cdp:st:E6}).
\item\hypertarget{step-II-F-3}{} Root-curve tangent cohomology
and the two-parameter period sections give the two smooth
surface germs and their tuples
(\cref{cdp:st:E13},
\cref{cdp:st:E14}).
\item\hypertarget{step-II-F-4}{} Admissible comb period sections
realize every tangent direction while retaining affine attachments.
Their tangent isomorphism makes the Kuranishi ideal zero and gives
the displayed comb dimensions
(\cref{cdp:st:E15}).
\item\hypertarget{step-II-F-5}{} Disc propagation and finite
curve projections make the fixed ordered-pair locus open and
closed in a connected reduced analytic base. This gives
separation through arbitrary finite deformation chains
(\cref{push:src:E.2.23}).
\item\hypertarget{step-II-F-6}{} The Kuranishi retraction transports
that separation to the $C^\infty$ space of integrable structures
and to the indicated diffeomorphism quotients
(\cref{push:src:E.2.24}).
\item\hypertarget{step-II-F-7}{} The flat cup kernel in the comb smoothing family
imposes a cyclic two-jet condition at each central node. The explicit
conic frame map realizes a dense open set of these quotients;
twelve teeth attain the balanced cohomology. The exact collapse
and semicontinuity maps give $K_{12}=\mathcal O(-8)\oplus\mathcal O(-9)^2$
on the stated open family, and the Hodge formula gives
$(h^{1,1},h^{1,2})=(49,48)$
(\cref{inv:E2.5},
\cref{inv:E2.2}).
\end{enumerate}

\hypertarget{outline-II-G}{}
\subsubsection*{Theorem G}
\begin{enumerate}[label=\arabic*.,leftmargin=*]
\item\hypertarget{step-II-G-1}{} The exact scalar is $\alpha=cu$
with fixed $c\ne0$. Setting $\tau=\log(cu)/(2\pi i)$ gives
the central lattice generated by $(2,0),(1+i,0),(1,1/2),(0,\tau)$
and affine fourth-root action
$(z,\eta)\mapsto(-iz,\eta+\tau/4+\kappa)$, where
$4\kappa\in\Z$ (\cref{ii:LR:covering}).
\item\hypertarget{step-II-G-2}{} Projection to $\eta$ is the
Albanese map with period lattice
$\tfrac12\Z+(\tau/4+\kappa)\Z$. The inverse maps
$[\eta]\mapsto[2\eta]$ and $[v]\mapsto[v/2]$ give its modulus
$\omega=\tau/2+2\kappa$, and one $u$-turn gives $\omega+1/2$.
Its Tate parameter $Q=e^{2\pi i\omega}$ therefore changes to $-Q$.
For $\operatorname{Im}\omega>1$, a relation
$\omega+1/2=(A\omega+B)/(C\omega+D)$ with
$\left(\begin{smallmatrix}A&B\\C&D\end{smallmatrix}\right)\in\operatorname{SL}_2(\Z)$
would force $|C\omega+D|=1$, hence $C=0$ and an integral period change,
which excludes this half-integer shift
(\cref{ii:LR:covering}).
\item\hypertarget{step-II-G-3}{} Ambient $H^2(Y,\R)=0$ makes
every effective divisor vertical. Meromorphic functions therefore
descend through the torus fibration; any endpoint biholomorphism
preserves its unique multiplicity-four fibre. The previous step
excludes every one-turn endpoint
(\cref{ii:LR:covering}).
\item\hypertarget{step-II-G-4}{} The source-line lift $C$ satisfies
$C^2=tB_0$. Put $U=C^{-1}$ and $u=r^4$. The two-turn map
$\widetilde D_r=-r^{-2}U$ has cap expression
$-r^{-2}\zeta^2U$ and extends through every cusp chart.
Its exact norm is
$\widetilde D_{-r}\widetilde D_r=t^{-1}B_u^{-1}$.
On the $B_u$-quotient this gives the four-turn closing multiplier
$t^{-1}$, with cap multiplier $\theta/t$ for $\theta=\zeta^4$
and the same cusp period relabelling. The resulting free involution
on the family closed after four turns has quotient parameter
$w=r^2$ and $u=w^2$ (\cref{ii:LR:covering}).
\end{enumerate}
\subsection{Known and new}
\begingroup\small
\begin{longtable}{@{}p{.10\textwidth}p{.85\textwidth}@{}}
\toprule Theorem & Status and source\\\midrule\endhead
A & New normalization criterion extending \cite[\S10.4, pp.~89--90]{AC26};
corrects \cite[Theorem~2.2, p.~680]{CDP2020} and
\cite[Proposition~1.1, p.~78; Lemma~1.5, pp.~79--80]{CDP98}.\\
B & New integral refinement; rational component formula:
\cite[\S10.5, p.~91]{AC26}; fibre classification: \cref{cdp:st:C.3.1,cdp:st:C.3.11}.\\
C & New unrestricted sign and equality theorem; standard-family Euler value:
\cite[Theorem~7.22, p.~60]{AC26} and \cite[Proposition~7.1, p.~17]{Engel2026}.\\
D & New local parity and refinements of each fixed punctured family:
\cref{cdp:st:C.2.UN.1,cdp:st:C.2.UN.4,cdp:st:C.2.UN.14}.\\
E & New twisted formulas and normalization-map ranks; untwisted values:
\cite[Theorem~6.12, pp.~81--82]{DGKXv1}, arXiv version~1.\\
F & New answers to \cite[\S10(1), p.~33]{Viaclovsky2026Combs} and
\cite[\S1.1(1), p.~2]{Viaclovsky2026Log}, with Hodge values on the stated open twelve-tooth locus;
corrects \cite[Corollary~1.3, preprint p.~2]{HKP2000}.\\
G & New minimal covering for the fixed scalar and root data of
\cite[equation~(2.3), p.~4; Lemma~2.7 and Corollary~2.8, pp.~6--7;
Proposition~3.6 and Definition~3.7, pp.~9--10]{LiuRollenske2026}, arXiv version~2.\\
\bottomrule\end{longtable}\endgroup
\subsection{Methodology}
This text is AI-produced; no part of it was written by a human.
GPT-5.6 Sol (OpenAI), in Codex at reasoning effort ultra, contributed
to the working edition of 24 August--9 September 2026. The pinned
repository's authorship record also credits GPT-6 Astra (OpenAI),
at reasoning effort ultra, from 5 September.
GPT-6 Astra (OpenAI), in ChatGPT Pro mode, with the effort setting
not separately exposed, carried out the research of 12--16 November,
the independent re-derivation of 16--18 November,
and the writing of 16--20 November 2026, followed by the publication
revision and further checks of 20--23 November 2026.

Claude Fable 5.2 (Anthropic), at effort high and at effort max,
and Claude Opus 5.5 (Anthropic), at effort max, undertook readings,
intakes, checks and recomputations on 12--20 November 2026.
Within the series, Claude Opus 5.5 also rewrote Paper~V's abstract
and \S1 on 20 November and read Paper~I in full on 21--22 November.
On 22 November it corrected two accents and one family index in
Paper~I and updated its records links. Their objections and
calculations were compared with the statements, proofs and cited
sources, and the resulting corrections were incorporated into the
text. The records specify the passages and calculations examined
by each model.

The records associate 3,391 entries with statements of this paper:
1,145 examinations of proof steps, 36 independent
re-derivations, 166 computations or formal certificates,
1,776 readings of proofs and applications, 263 uses of
background results without a fault found, and 5 refutation
records associated with their corrected statements.
These entries belong to 2,804 recorded groups; one argument or
calculation can be associated with several statements and kinds of
examination. The counts describe those recorded associations, and
each outcome specifies its scope and prior proof exposure.

The records distinguish derivations completed before the target proof
was opened from comparisons made after reading it. Computational
comparisons specify the prior and newly used methods, the displayed
calculations they cover, and any unavailable historical program.
The compact example with three doubled nodal supports supplies the
two conductor-scope refutations through one construction. It was
developed after one participating model had read the affected proofs;
its subsequent comparisons and computations add no independent
re-derivation to these totals. The local curve-ring calculation and
the reviews of its applications form a separate recorded group.
No human has checked the mathematics.

The check records and programs are available on the
\href{\SsixPublishRecords}{records page} for this paper; the repository
commit is \texttt{\SsixRepositoryCommit}.
The text, programs and check records are dedicated to the public domain under
\href{https://creativecommons.org/publicdomain/zero/1.0/}{CC0 1.0 Universal}.

\subsection{How the paper is organized}
\Cref{part:II:cotangent-fibre-geometry} develops the index,
Picard and direct-image calculations, the smooth-fibre classification,
and the cotangent and conductor arguments. It ends with the full
normalization criterion of Theorem~A in
\cref{sec:II:global-normalization-criterion}.
\Cref{part:II:irreducible-euler} treats irreducible supports:
normalization topology, canonical boundaries, local transverse
singularities and the equality cases for their Euler numbers.
It includes the compact construction with three doubled nodal
supports in \cref{ii:three-doubled-nodal-fibres}.

\Cref{part:II:periods-tangent-deformations} computes period ranks,
Picard resonance, tangent cohomology and deformation spaces for the
marked families. It contains the conductor nonvanishing arguments,
the open twelve-tooth splitting calculation in \cref{inv:E2.5},
and the Liu--Rollenske parameter covering.
\Cref{part:II:component-geometry} studies component covers,
meromorphic fields, integral component lattices and horizontal
conductors. Its elliptic-interface arguments give the reducible-fibre
Euler equality criterion in \cref{cdp:st:C.2.INT.3}.
\Cref{part:II:compact-occurrence} constructs compact torus fibrations
with the stated component types and classifies marked Hopf, ruled
and elliptic supports with their full multiplicities.

\Cref{part:II:canonical-reduction} develops the component contractions,
canonical-order reduction and finite rooted generation, ending with
the Euler sign and equality criterion for arbitrary torus fibres in
\cref{sec:II:all:12}.
\Cref{part:II:remaining-splitting} states the remaining comb bundle
splitting problem in \cref{inv:E2.3}. Its degree bounds, cup-duality
reduction and twelve-tooth case are proved in
\cref{prop:II:comb-kernel-bounds,prop:II:comb-cup-duality,inv:E2.5}
of \cref{part:II:periods-tangent-deformations}.
The figure identifies five premise uses in the proof of Theorem~A.

\begin{figure}[htbp]\centering
\begin{tikzpicture}[>=stealth,every node/.style={font=\small,align=center}]
\node[draw,rounded corners,text width=4.1cm,minimum height=1cm] (SB) at (0,0)
 {Smooth general fibres\\\hyperref[cdp:st:C.3.11]{\ref*{cdp:st:C.3.11}}};
\node[draw,rounded corners,text width=4.1cm,minimum height=1cm] (TV) at (6,0)
 {Twisted cotangent\\vanishing\\\hyperref[cdp:st:A.4.9]{\ref*{cdp:st:A.4.9}}};
\node[draw,rounded corners,text width=4.6cm,minimum height=1cm] (NT) at (3,-2)
 {Persistent normalization\\pullback\\\hyperref[lem:II:persistent-normalization-pullback]{\ref*{lem:II:persistent-normalization-pullback}}};
\node[draw,rounded corners,text width=4.1cm,minimum height=1cm] (GV) at (3,-4)
 {Generic-vanishing criterion\\\hyperref[cdp:st:C4f]{\ref*{cdp:st:C4f}}};
\draw[->,semithick] (SB) -- (NT);
\draw[->,semithick] (TV) -- (NT);
\draw[->,semithick] (NT) -- (GV);
\draw[->,semithick] (SB.west) -- ++(-.5,0) |- (GV.west);
\draw[->,semithick] (TV.east) -- ++(.5,0) |- (GV.east);
\end{tikzpicture}
\caption{The smooth-fibre and normalization premises of Theorem~A.}
\par\medskip
\begingroup\small
\begin{tabularx}{\textwidth}{@{}p{.25\textwidth}X@{}}
\toprule Arrow & Supplying and receiving arguments\\\midrule
\ref{cdp:st:C.3.11}$\to$\ref{lem:II:persistent-normalization-pullback} &
The torus and Hopf alternatives give the smooth-fibre cohomology
used in the conductor bound.\\
\ref{cdp:st:A.4.9}$\to$\ref{lem:II:persistent-normalization-pullback} &
The full-fibre filtration and twisted quotient vanishing turn
persistent ambient sections into inverse-Jacobian sections.\\
\ref{lem:II:persistent-normalization-pullback}$\to$\ref{cdp:st:C4f} &
Persistent sections on a fixed full fibre force its normalization
Picard pullback to vanish.\\
\ref{cdp:st:C.3.11}$\to$\ref{cdp:st:C4f} &
The two smooth-fibre types select the torus and Hopf direct-image
arguments in the criterion.\\
\ref{cdp:st:A.4.9}$\to$\ref{cdp:st:C4f} &
On the Hopf branch, the transgression vanishings make the second
direct image torsion. The full-fibre obstruction and analytic
continuation give nonvanishing for every twist.\\
\bottomrule\end{tabularx}\endgroup\end{figure}

\paragraph{Further results.}
\begingroup\small
\begin{itemize}[label={},leftmargin=0pt,itemsep=1pt,parsep=0pt,topsep=3pt]
\item \cref{cdp:st:C.1.1}\quad Tangent index and Picard components.
\item \cref{cdp:st:C.1.3}\quad A lower bound for twisted tangent cohomology.
\item \cref{cdp:st:E1}\quad Vertical divisors and the field of meromorphic functions.
\item \cref{cdp:st:D.1.5}\quad Two projective torus fibrations.
\item \cref{cdp:st:D.1.6}\quad Components in Wang's torus fibration.
\item \cref{cdp:st:D.1.7}\quad Real cohomology and the structure sheaf.
\item \cref{cdp:st:D.1.8}\quad Injectivity of integral generization.
\item \cref{cdp:st:D.1.9}\quad Line bundles of infinite order on every fibre.
\item \cref{cdp:st:D.1.10}\quad Invariant classes in the restriction image.
\item \cref{cdp:st:D.1.11}\quad Generic line bundles on every fibre.
\item \cref{cdp:st:D.1.12}\quad Cup products and cohomological flatness.
\item \cref{cdp:st:D.1.4}\quad Restriction maps and splitting degrees.
\item \cref{cdp:st:D.1.13}\quad Higher direct images of Wang's threefold.
\item \cref{cdp:st:D.1.15}\quad Constant integral monodromy and coherent splitting.
\item \cref{cdp:st:C.3.2}\quad A secondary Kodaira surface with torsion-free first homology.
\item \cref{cdp:st:C.3.3}\quad A secondary Hopf surface with torsion-free first homology.
\item \cref{cdp:st:C.3.4}\quad Obstruction to K\"ahler resolutions.
\item \cref{cdp:st:C.3.7}\quad Exclusion of bielliptic fibres.
\item \cref{cdp:st:C.3.8}\quad Exclusion of elliptic Hopf general fibres.
\item \cref{cdp:st:C.3.9}\quad Exclusion of secondary Kodaira general fibres.
\item \cref{cdp:st:C.3.10}\quad Exclusion of primary Kodaira general fibres.
\item \cref{cdp:st:C.3.12}\quad Exclusion of Inoue general fibres.
\item \cref{cdp:st:C.3.11}\quad Complete smooth-fibre classification, occurrence and full-fibre cohomology.
\item \cref{cdp:st:C.3.5}\quad Euler numbers of fibre supports.
\item \cref{cdp:st:C.3.6}\quad Bundles with primary and secondary Hopf fibres.
\item \cref{def:II:conormal-data}\quad Conormal and torsion notation.
\item \cref{wbII:lem:torsion-identification}\quad Identification of the torsion subsheaf.
\item \cref{cdp:st:A.2.1}\quad The Jacobian fractional ideal for every reduced hypersurface.
\item \cref{cdp:st:A.2.2}\quad Normalization, conductor, and the excess kernel.
\item \cref{cdp:st:A.1.2}\quad The canonical normal-crossings formula.
\item \cref{cdp:st:A.2.4}\quad Pinch points and cuspidal edges.
\item \cref{cdp:st:A.2.5}\quad Vertical residue poles at a tacnode--cusp intersection.
\item \cref{cdp:st:A.2.6}\quad A cusp along a non-complete-intersection curve.
\item \cref{cdp:st:B1}\quad A complete nodal-family computation, including tensor order, explicit extending twists and the endpoint maps.
\item \cref{cdp:st:B2}\quad A projective example with twelve length-one contributions.
\item \cref{cdp:st:B3}\quad Topology of the Weierstrass product.
\item \cref{cdp:st:A.3.1}\quad All thickenings and their lifting maps.
\item \cref{cdp:st:A.3.2}\quad Sections on a multiple abelian fibre.
\item \cref{cdp:st:A.3.3}\quad Finite Picard translations for multiple fibres.
\item \cref{cdp:st:A.3.4}\quad Ambient cotangent vanishing on smooth multiple fibres.
\item \cref{cdp:st:A.3.5}\quad Every reducible full fibre has generic ambient vanishing.
\item \cref{cdp:st:A.4.1}\quad Cotangent groups on the rational normalization fibre.
\item \cref{cdp:st:A.4.2}\quad Constant twists on an elliptic ruled surface.
\item \cref{cdp:st:A.4.3}\quad Twisted differentials on a ruled normal-crossings surface.
\item \cref{cdp:st:A.4.4}\quad Canonical triviality on an irreducible reduced torus fibre.
\item \cref{cdp:st:A.4.5}\quad Twisted cotangent groups of irreducible normal-crossings torus fibres.
\item \cref{cdp:st:A.4.6}\quad Euler numbers of irreducible normal-crossings torus fibres.
\item \cref{cdp:st:A.4.7}\quad Unitary-flat differentials on every compact reduced complex surface.
\item \cref{cdp:st:A.4.8}\quad Cotangent vanishing for the transgression family.
\item \cref{cdp:st:A.5.1}\quad Conormal sections with trivial normalization pullback.
\item \cref{cdp:st:A.5.2}\quad Normal Gorenstein surfaces with vanishing first Chern class and Euler characteristic.
\item \cref{cdp:st:A.5.3}\quad The effective conductor divisor on a minimal resolution.
\item \cref{cdp:st:A.5.4}\quad The remaining normalization class for a persistent ambient obstruction.
\item \cref{cdp:st:A.5.6}\quad Uniform negative-degree vanishing for a fixed vector bundle.
\item \cref{cdp:st:A.5.9}\quad The ruled obstruction with a reduced conductor image.
\item \cref{cdp:st:A.5.11}\quad The genus of a horizontal conductor image from a local torus group.
\item \cref{cdp:st:A.5.12}\quad The arithmetic genus of conductor images from coherent duality.
\item \cref{cdp:st:A.5.13}\quad The conductor bound when the general fibre has $q=1$ and $p_g=0$.
\item \cref{cdp:st:A.5.5}\quad Conductor-cleared jets on a complete-intersection critical curve.
\item \cref{cdp:st:A.5.7}\quad A pure double conductor forces a smooth cusp image.
\item \cref{cdp:st:A.5.8}\quad Rational singularities for a vertically resolved ruled normalization.
\item \cref{cdp:st:A.5.14}\quad The genus bound and delta-one curve germs.
\item \cref{cdp:st:A.5.10}\quad Eliminating vertical conductor on a smooth ruled cusp normalization.
\item \cref{cdp:st:A.5.20}\quad A noncontracted horizontal cusp has zero Euler number.
\item \cref{cdp:st:A.5.15}\quad Binomial germs and spherical intersection classes.
\item \cref{cdp:st:A.5.17}\quad Isolated singularities are excluded and suspension germs are transverse products.
\item \cref{cdp:st:A.5.18}\quad Every horizontal tacnode image is elliptic.
\item \cref{cdp:st:A.5.19}\quad Ordinary-triple images have arithmetic genus at most one.
\item \cref{cdp:st:A.5.16}\quad Rational normalization bounds with the exact constant-sheaf map.
\item \cref{def:II:minimal-boundary-data}\quad A specified contraction sequence defines the boundary multiplicities and the count of multiplicity-one centres.
\item \cref{prop:II:minimal-boundary-complements}\quad Minimal rational boundary complements and the conductor-image bound for their negative cases.
\item \cref{cdp:st:A.5.21}\quad A homogeneous binary-product grid is excluded.
\item \cref{cdp:st:C2a}\quad Euler localization for torus fibrations.
\item \cref{cdp:st:C.2.NC.2}\quad The normalized double curves and the coherent Euler formula.
\item \cref{cdp:st:C.2.NC.3}\quad Uniform multiple normal-crossings fibres.
\item \cref{cdp:st:C.2.NC.4}\quad Reducible normal-crossings fibres with unequal multiplicities.
\item \cref{cdp:st:C.2.NC.5}\quad Global positivity and the all-twist normalization obstruction in the normal-crossings class.
\item \cref{cdp:st:C2d}\quad The normal jets along cusp, tacnode and triple curves.
\item \cref{night:P4.1}\quad An elliptic bundle degeneration.
\item \cref{cdp:st:D.2.7}\quad A simply connected compact torus fibration and its integral homology.
\item \cref{night:P4.3}\quad Integral homology of the local degeneration.
\item \cref{night:P4.6}\quad Euler characteristic of a normal crossing divisor.
\item \cref{night:P4.7}\quad The Chern number of a normal crossing fibration.
\item \cref{night:P4.12}\quad Irreducible fibres of torus fibrations.
\item \cref{night:P4.8}\quad Normal crossing torus fibres and their normalizations.
\item \cref{night:P4.13}\quad Monodromy and the smooth locus of an irreducible fibre.
\item \cref{night:P4.14}\quad Genus-one curves with nonreduced conductor.
\item \cref{night:P4.15}\quad Finite rank drop of a normalization.
\item \cref{night:P4R.1}\quad Fundamental group of an irreducible fibre.
\item \cref{night:P4R.6}\quad Conductor and exceptional pole divisor.
\item \cref{night:P4.33}\quad Normalization cohomology and the conductor Euler characteristic.
\item \cref{night:P4R.2}\quad Euler number for a bijective normalization.
\item \cref{night:P4.16}\quad Spherical intersection lattices of $x^a-zy^n=1$ and the torus-fibre obstruction.
\item \cref{night:P4.19}\quad A ramified conductor with vanishing nearby homology.
\item \cref{night:P4.20}\quad Ordinary double strata and normal crossings.
\item \cref{night:P4.28}\quad Local rigidity of a transverse cusp.
\item \cref{night:P4R.3}\quad Conductor topology over a higher-genus ruling.
\item \cref{night:P4R.4}\quad Euler number for a rational incidence tree.
\item \cref{night:P4.21}\quad Positive Euler number for rational normalization.
\item \cref{night:P4.22}\quad Euler number of an elliptic-ruled normalization.
\item \cref{night:P4.25}\quad Semi-log-canonical fibre supports.
\item \cref{night:P4R.7}\quad Nonnegative Euler number for an irrational ruled normalization resolution.
\item \cref{night:P4.26}\quad Euler numbers of irreducible and normal crossing fibres.
\item \cref{night:P4.29}\quad Zero-Euler cuspidal ruled supports.
\item \cref{lem:II:rational-conormal-defects}\quad Conormal defects of compact curves.
\item \cref{lem:II:triple-conductor-symbolic-square}\quad Primary ordinary-triple conductor schemes.
\item \cref{night:P4.30}\quad Symbolic double of a unibranch elliptic curve.
\item \cref{night:P4.31}\quad Zero-Euler ruled supports with ordinary triple curves.
\item \cref{night:P4R.10}\quad Euler characteristic of a tacnodal conductor.
\item \cref{night:P4R.11}\quad Tacnodal supports of Euler number zero.
\item \cref{night:P4.32}\quad The zero-Euler fibre supports.
\item \cref{night:P4R.12}\quad Components outside the pole divisor of a ruled fibre.
\item \cref{night:P4R.13}\quad Elliptic components of the conductor.
\item \cref{ii:three-doubled-nodal-fibres}\quad A compact fibration with three doubled nodal supports.
\item \cref{lem:II:conductor-two-branch-cycle}\quad Two image branches force a conductor cycle.
\item \cref{night:P4.35}\quad Acyclic normalization conductors.
\item \cref{night:P4.36}\quad Smooth minimal Hopf normalization, unramified elliptic conductor and zero Euler number.
\item \cref{night:P4.37}\quad Every class VII normalization is a smooth minimal Hopf surface.
\item \cref{night:P4.38}\quad Local exclusions for ordinary double curves.
\item \cref{night:P4.39}\quad Pole coefficients over an ordinary triple curve.
\item \cref{night:P4R.14}\quad Local models along a smooth tacnodal critical curve.
\item \cref{night:P4R.15}\quad Product germs from a nonzero Hessian.
\item \cref{night:P4.51}\quad Vanishing local intersection forms from a nonzero first Betti number.
\item \cref{ii:excess-one-curve-germs}\quad The reduced curve germs of excess one.
\item \cref{night:P4.40}\quad Exclusion of ordinary double curves at positive Euler number.
\item \cref{night:P4.43}\quad Euler bound for uniform cuspidal ruled supports.
\item \cref{night:P4.41}\quad Hessian defect at a nodal critical curve.
\item \cref{night:P4.42}\quad Two identifications in the tacnodal ruled case.
\item \cref{night:P4.44}\quad The contracted section and three-branch critical point.
\item \cref{night:P4.45}\quad Planar critical curves of ordinary triple type.
\item \cref{night:P4.46}\quad Euler bound for uniform ordinary-triple ruled supports.
\item \cref{night:P4.47}\quad Rational normalization for a nondegenerate quartic square.
\item \cref{night:P4.48}\quad Conductor defects at three-branch critical points.
\item \cref{night:P4.49}\quad Rational normalization at a three-axis critical point.
\item \cref{night:P4R.16}\quad The defect of a tacnode with a transverse line.
\item \cref{night:P4R.17}\quad Defect bounds at nonordinary three-branch points.
\item \cref{night:P4.50}\quad Euler bound for uniform tacnodal ruled supports.
\item \cref{night:P4R.18}\quad Transverse type along a noncontracted horizontal section.
\item \cref{night:P4.52}\quad Multiplicity of contracted horizontal pole sections.
\item \cref{night:P4.53}\quad Normality after the multiplicity-three blowup.
\item \cref{night:P4.54}\quad Noncontraction of the horizontal pole section.
\item \cref{night:P4R.19}\quad Cubic normal tensors along a $D_4$ branch.
\item \cref{night:P4.55}\quad The horizontal conductor cover and its rational third sheet.
\item \cref{night:P4.56}\quad At least three branches along the second conductor component.
\item \cref{night:P4.57}\quad Exclusion of two distinct conductor junctions.
\item \cref{night:P4R.20}\quad Exclusion of a three-branch conductor junction.
\item \cref{night:P4.10}\quad A projective torus fibration with Chern number twenty-four.
\item \cref{night:P4.11}\quad A primary Hopf bundle on a product of spheres.
\item \cref{cdp:st:C4c}\quad A finite cohomology test and a fixed fibre.
\item \cref{cdp:st:C4a}\quad Generic tangent cohomology vanishing.
\item \cref{cdp:st:C4b}\quad Twisted tangent cohomology on rational homology six-spheres.
\item \cref{cdp:st:C4d}\quad A finite transgression criterion for tangent cohomology.
\item \cref{cdp:st:C4e}\quad The unrestricted normalization target under the integral CDP hypotheses.
\item \cref{cdp:st:C4g}\quad Hexagonal fibres and tangent cohomology in three families.
\item \cref{cdp:st:C5a}\quad The normalization quotient and its cohomology map.
\item \cref{cdp:st:D.2.1}\quad Components under trivial torus monodromy.
\item \cref{cdp:st:D.2.3}\quad Logarithmic period obstruction.
\item \cref{cdp:st:D.2.6}\quad Extension of the marked torus period.
\item \cref{cdp:st:D.2.8}\quad Coherent and integral cohomology of the trivial-monodromy example.
\item \cref{cdp:st:D.2.2}\quad Euler equality under trivial torus monodromy, including normal-crossing supports.
\item \cref{cdp:st:E5}\quad The infinitesimal map and the rank-two boundary.
\item \cref{cdp:st:E7}\quad Deformations of the constructed torus fibrations.
\item \cref{cdp:st:E10}\quad Deformations over a fixed genus-two curve.
\item \cref{cdp:st:E11}\quad All coprime $(3k,4l)$ multiplicities.
\item \cref{cdp:st:E16}\quad Rank divisors from the root Kodaira--Spencer morphism.
\item \cref{push:src:E.2.4}\quad Arbitrary connected deformation chains.
\item \cref{push:src:E.2.5}\quad A fractional canonical pencil at a smooth limit.
\item \cref{push:src:E.2.6}\quad Positive canonical slope preserves the fibration.
\item \cref{push:src:E.2.7}\quad An extended fibration cannot reach the zero period.
\item \cref{push:src:E.2.8}\quad Horizontal divisors force non-K\"ahler fibre normalizations.
\item \cref{push:src:E.2.9}\quad No horizontal divisor when all fibres are irreducible.
\item \cref{push:src:E.2.10}\quad Persistence for either canonical slope.
\item \cref{push:src:E.2.11}\quad Purity with a smooth torus Jacobian.
\item \cref{push:src:E.2.12}\quad No additional central critical values.
\item \cref{push:src:E.2.13}\quad Persistence of an isolated finite multiple fibre.
\item \cref{push:src:E.2.14}\quad Extension of the circle field and exclusion of finite-root collapse.
\item \cref{push:src:E.2.15}\quad An isotrivial volume obstruction.
\item \cref{push:src:E.2.16}\quad The central cusp is reduced and irreducible.
\item \cref{push:src:E.2.17}\quad Persistence of the normal-crossing cusp.
\item \cref{push:src:E.2.18}\quad The normalized cusp and its marked hexagon.
\item \cref{push:src:E.2.19}\quad Extension of the marked cusp torus action.
\item \cref{push:src:E.2.20}\quad Identification of the analytic cusp.
\item \cref{push:src:E.2.21}\quad Global identification of a smooth completion.
\item \cref{push:src:E.2.22}\quad Propagation from a known open subset of a disc.
\item \cref{push:src:E.5.1}\quad Singular fibres of the two-parameter threefold.
\item \cref{push:src:E.5.2}\quad Monodromy, centralizer, and integral lattices.
\item \cref{def:II:two-parameter-period-data}\quad Two-parameter period data.
\item \cref{push:src:E.5.3}\quad Biholomorphism classes in the two-parameter family.
\item \cref{wb:thm:uniform-nontorsion-audit}\quad Uniform non-torsion.
\item \cref{wbII:prop:conductor}\quad Conductor-supported ambient differential.
\item \cref{wbII:cor:normal-crossing-kernel}\quad Nonzero cotangent connecting kernel.
\item \cref{wbII:thm:cartier-duality}\quad Duality and trace pairings for compact Cartier divisors and full scheme-theoretic fibres.
\item \cref{wbII:lem:curve-flatness}\quad Flatness and Cartier fibres over a curve.
\item \cref{wbII:prop:analytic-base-change}\quad Analytic base change in degree two.
\item \cref{wbII:cor:surface-top-base-change}\quad Top-degree base change for surface fibres.
\item \cref{wbII:prop:pic0-pullback}\quad Normalization pullback and dualization, including the topologically trivial case.
\item \cref{wbII:prop:analytic-ah-boundary}\quad Boundary cocycles, finite-cover norms and cusp Chern-class specialization in the standard family.
\item \cref{wbII:prop:local-cusp-picard-obstruction}\quad An explicit nonzero local cusp restriction.
\item \cref{wbII:lem:analytic-support-specialization}\quad The toric support function detects central restriction.
\item \cref{wbII:thm:analytic-all-picard}\quad Vanishing of the global edge section, fibre Chern class and normalized cusp restriction.
\item \cref{wbII:thm:candidate-conductor}\quad Conductor nonvanishing for the standard family.
\item \cref{wbII:cor:candidate-conductor-actual}\quad Tangent-cohomology nonvanishing on the standard fibre for every ambient line bundle.
\item \cref{wbII:lem:generic-endpoints}\quad Generic endpoints in an analytic parameter space.
\item \cref{wbII:lem:cdp-generic-endpoints}\quad Generic endpoints on \texorpdfstring{$\operatorname{Pic}^0$}{Pic0}.
\item \cref{wbII:thm:corrected-cdp-global}\quad Global vanishing from one degree-two vanishing.
\item \cref{wbII:cor:corrected-cdp-relative}\quad Relative vanishing criterion.
\item \cref{wbII:thm:corrected-cdp-fibrewise}\quad Reduced-fibre vanishing criterion.
\item \cref{wbII:test:cdp-comparison}\quad Positive Euler number and cohomology nonvanishing.
\item \cref{wbII:prop:relative-local}\quad Crossing and multiple-fibre modules.
\item \cref{cdp:st:C.2.UN.3}\quad The same punctured torus family has non-bimeromorphic zero-Euler fillings.
\item \cref{cdp:st:C.2.UN.15}\quad No component of a torus degeneration has general-type normalization.
\item \cref{cdp:st:C.2.UN.16}\quad Commuting semicomplete fields on a component cover.
\item \cref{cdp:st:C.2.UN.17}\quad Canonical support forces the numerically anticanonical class VII case.
\item \cref{cdp:st:C.2.UN.19}\quad Every class VII normalization component has Hopf minimal model.
\item \cref{cdp:st:C.2.UN.20}\quad Descent of the rank-one component fibration.
\item \cref{cdp:st:C.2.UN.22}\quad A component cover with a rank-two translation frame has coherent Euler characteristic zero.
\item \cref{cdp:st:C.2.UN.23}\quad A cyclic torus quotient with nonnegative Kodaira dimension is a torus, bielliptic, or K3.
\item \cref{cdp:st:C.2.UN.26}\quad Every compact elliptic surface satisfies the intrinsic rank-one field condition.
\item \cref{cdp:st:C.2.UN.28}\quad The sum of the normalization geometric genera in one torus fibre is at most one.
\item \cref{cdp:st:C.2.UN.29}\quad Rank-one elliptic component covers have only smooth multiple fibres.
\item \cref{cdp:st:C.2.UN.31}\quad Cyclic covers of components with nonnegative Kodaira dimension.
\item \cref{cdp:st:C.2.UN.33}\quad A properly elliptic normalization component has one of eleven positive-coherent-Euler fibre configurations.
\item \cref{cdp:st:C.2.UN.36}\quad Compact component cores give a multiplicity-weighted invariant lattice in one nearby torus.
\item \cref{cdp:st:C.2.UN.38}\quad Every curve of a smooth minimal algebraic-dimension-zero Hopf normalization lies in the fibre boundary.
\item \cref{cdp:st:C.2.UN.39}\quad Curve genera, triple points and the finite monodromy constant decide zero Euler number.
\item \cref{cdp:st:C.2.UN.44}\quad The fundamental group and the component-incidence constraint of every torus fibre.
\item \cref{cdp:st:C.2.UN.47}\quad Positive geometric genus controls the whole support group.
\item \cref{cdp:st:C.2.UN.5}\quad Every normal irreducible support is a smooth torus quotient.
\item \cref{cdp:st:C.2.CR.1}\quad Topology and coefficients in the contracted-section case.
\item \cref{lem:II:contracted-section-duality}\quad Duality and ruling genus for a contracted section.
\item \cref{cdp:st:C.2.CR.3}\quad The genus-two node alternative is excluded.
\item \cref{cdp:st:C.2.CR.6}\quad The elliptic contraction has a planar conductor and only double transverse branches at its image.
\item \cref{cdp:st:C.2.UN.9}\quad A smooth full critical curve with generic node is a product node.
\item \cref{cdp:st:C.2.UN.10}\quad Contracted irrational sections and their Euler number.
\item \cref{cdp:st:C.2.IRR.1}\quad The reduced horizontal case is a minimal elliptic ruled double-curve fibre.
\item \cref{cdp:st:C.2.VII.1}\quad Anticanonical supports on class VII surfaces.
\item \cref{cdp:st:C.2.VII.2}\quad Conductor Euler bounds.
\item \cref{cdp:st:C.2.VII.3}\quad Euler number for class VII normalization.
\item \cref{cdp:st:C.2.VII.4}\quad Zero Euler number for class VII normalization.
\item \cref{prop:II:rational-torus-fibre-euler}\quad Positive Euler number for every irreducible rational-normalized torus fibre.
\item \cref{cdp:st:C.2.RAT.1}\quad Euler and attachment bounds for rational boundary classes.
\item \cref{cdp:st:C.2.RAT.3}\quad Vertical-tree attachment loss and the conductor Euler bound.
\item \cref{cdp:st:C.2.RAT.2}\quad Attachment bounds for rational tacnode covers.
\item \cref{lem:II:normal-cubic-compact-core}\quad Compact normal-cubic cycles and their transport into a nearby torus.
\item \cref{lem:II:cyclic-cubic-boundary}\quad Cyclic-cover character sheaves, boundary relations and compact intersection pairings.
\item \cref{cdp:st:C.2.D4.1}\quad Connected triple covers over a rational image.
\item \cref{cdp:st:C.2.D4.2}\quad Double covers over a rational ordinary-triple image.
\item \cref{prop:II:cubic-root-attachments}\quad Common cubic zeros, normalization branches and vertical attachments.
\item \cref{cdp:st:C.2.UN.7}\quad The irreducible local comparison and its rational and ruled consequences.
\item \cref{cdp:st:C.2.UN.12}\quad The unrestricted irreducible zero-Euler criterion.
\item \cref{cdp:st:C.2.INT.1}\quad Elliptic images and Picard restriction.
\item \cref{cdp:st:C.2.INT.2}\quad Exclusion of Enoki components.
\item \cref{cdp:st:C.2.INT.3}\quad Euler sign and equality for reducible fibres.
\item \cref{cdp:st:C.2.MH.1}\quad Purity of support cohomology.
\item \cref{cdp:st:C.2.MH.2}\quad Euler number from a ruled normalization.
\item \cref{cdp:st:C.2.MH.4}\quad Every irreducible class VII normalization contributes zero.
\item \cref{cdp:st:C2b}\quad Torus fibres of Euler numbers $-2$ and zero.
\item \cref{cdp:st:C2c}\quad Euler numbers in semistable toric rank-two models.
\item \cref{cdp:st:C.2.UN.11}\quad All rank-four monodromy stratum signatures.
\item \cref{cdp:st:C.2.UN.32}\quad K3 components, cyclic quotients and elliptic pencils.
\item \cref{cdp:st:C.2.UN.64}\quad Contractions and normal extensions of terminal ruled pairs.
\item \cref{cdp:st:C.2.UN.65}\quad Trivial local canonical class, triple points and the complete zero case.
\item \cref{cdp:st:C.2.UN.66}\quad Leaf contractions, elliptic cores and Euler number.
\item \cref{cdp:st:C.2.UN.67}\quad A torsion-normal elliptic boundary cannot lie on a rational component.
\item \cref{cdp:st:C.2.UN.68}\quad Rational toric component pairs and the exact dual torus.
\item \cref{cdp:st:C.2.UN.6}\quad Zero-Euler product fillings with essential Kodaira-dimension-one components.
\item \cref{cdp:st:C.2.UN.8}\quad A simply connected irreducible rational-normalized fibre.
\item \cref{cdp:st:C.2.UN.13}\quad A projective zero-Euler filling with an essential Kummer component.
\item \cref{cdp:st:C.2.UN.18}\quad Opposite logarithmic data produce a projective essential component.
\item \cref{cdp:st:C.2.UN.21}\quad A free half-period quotient realizes secondary Kodaira normalization.
\item \cref{cdp:st:C.2.UN.27}\quad Essential Enriques and projective elliptic components in zero-Euler fillings of one punctured torus family.
\item \cref{cdp:st:C.2.UN.30}\quad The fixed Kummer punctured family has zero-Euler fillings with arbitrarily large canonical growth on a projective elliptic component.
\item \cref{cdp:st:C.2.UN.34}\quad Every fixed-modulus elliptic quasi-bundle over a rational or elliptic curve occurs unchanged.
\item \cref{cdp:st:C.2.UN.35}\quad Every prescribed cyclic smooth-multiple elliptic quotient occurs in a zero-Euler torus filling.
\item \cref{cdp:st:C.2.UN.37}\quad Occurrence of every Hopf surface of algebraic dimension one.
\item \cref{cdp:st:C.2.UN.41}\quad Every diagonal primary Hopf surface occurs unchanged in a compact torus filling.
\item \cref{cdp:st:C.2.UN.42}\quad Every nonlinear primary Hopf surface occurs unchanged over a rational base.
\item \cref{cdp:st:C.2.UN.43}\quad Every secondary Hopf surface of algebraic dimension zero occurs unchanged.
\item \cref{cdp:st:C.2.UN.45}\quad Ruling bases of normalization components.
\item \cref{cdp:st:C.2.UN.40}\quad Hopf-axis adjacency with arbitrary multiplicities and point blowups.
\item \cref{cdp:st:C.2.UN.48}\quad The full compatible cycle is realized over a compact rational base.
\item \cref{cdp:st:C.2.UN.56}\quad The weighted Hopf and ruled root criterion, with exact inertia and compact elliptic-suspension realizations.
\item \cref{cdp:st:C.2.UN.53}\quad Reduced primary and secondary Hopf stars: the full marked criterion and compact realization.
\item \cref{cdp:st:C.2.UN.54}\quad Every reduced all-Hopf support has one cycle, with an exact algebraic-dimension-zero classification.
\item \cref{cdp:st:C.2.UN.57}\quad The exact marked Hopf and ruled cycle criterion and its compact realizations.
\item \cref{cdp:st:C.2.UN.58}\quad The complete canonical and period reduction for all-Hopf trees.
\item \cref{cdp:st:C.2.UN.59}\quad All multiplicities and all diagonal Hopf and ruled vertices.
\item \cref{cdp:st:C.2.UN.60}\quad Exact pruning and compact reinsertion classify every reduced all-Hopf support.
\item \cref{cdp:st:C.2.UN.61}\quad The exact rooted elliptic-tree criterion and compact occurrence.
\item \cref{cdp:st:C.2.UN.62}\quad Unicyclic supports with elliptic quasi-bundle torsors.
\item \cref{cdp:st:C.2.UN.63}\quad The weighted-root criterion for elliptic and Hopf trees.
\item \cref{cdp:st:C.2.UN.69}\quad Maximum components without a minimality assumption.
\item \cref{cdp:st:C.2.UN.83}\quad Torsion-normal elliptic boundaries have explicit elliptic, nonlinear Hopf or one-section ruled models.
\item \cref{cdp:st:C.2.UN.71}\quad All positive canonical-flat central log spaces over a disc.
\item \cref{cdp:st:C.2.UN.73}\quad The prescribed toric fibre and its conjugate over the projective line.
\item \cref{cdp:st:C.2.UN.76}\quad Exact weighted toric descent and compact occurrence.
\item \cref{cdp:st:C.2.UN.74}\quad A section and ruling fibres give the exact transverse curve blowup.
\item \cref{cdp:st:C.2.UN.75}\quad Canonical-flat cores in the Hopf-minimum case.
\item \cref{cdp:st:C.2.UN.77}\quad Hopf minimal models at maximal canonical order.
\item \cref{cdp:st:C.2.UN.79}\quad The \((-1,-2)\) curve with canonical orders \((0,0,1)\) occurs in a compact zero-Euler torus fibre.
\item \cref{cdp:st:C.2.UN.80}\quad Exceptional divisors in the component boundary.
\item \cref{cdp:st:C.2.UN.81}\quad Elliptic parents at marked Hopf interfaces.
\item \cref{cdp:st:C.2.UN.82}\quad Every positive maximal Hopf after minimalization is a torsion-normal leaf.
\item \cref{cdp:st:C.2.UN.84}\quad The highest Hopf order can be made the global maximum.
\item \cref{cdp:st:C.2.UN.78}\quad Root base change and birational descent from canonical-flat cores.
\item \cref{cdp:st:C.2.UN.72}\quad All no-triple fibres are ordinary vertical blowups of the classified elliptic cores.
\item \cref{cdp:st:C.2.UN.85}\quad Every reduced SNC torus degeneration has a reversible finite reduction to a classified core.
\item \cref{cdp:st:C.2.UN.86}\quad Complete rooted generation and the zero-Euler criterion for arbitrary fibres.
\item \cref{inv:E2.3}\quad The comb bundle splitting problem.
\end{itemize}\endgroup

\clearpage
\section{Common conventions}
\hypertarget{tag-034E}{}\label{sec:II:conventions}
\begin{definition}[Marked periods, fibres and orientations]
\hypertarget{tag-P202}{}\label{def:II:marked-conventions}
\StatementRecord{def:II:marked-conventions}{proved}{P202}{\SeriesRef{fam:src:C.1.1}, \SeriesRef{I:def:finite-coordinates}, \SeriesRef{fam:src:A.5.1}, \SeriesRef{fam:src:B.2.1}}
All spaces and coherent sheaves are complex analytic unless a polynomial
ring is specified. We use $B=\PP^1_t$, $(p_1,p_2,p_0)=(0,1,\infty)$,
the ordered lattice
$\Lambda=\Z\langle\widehat\gamma,\widehat u,\widehat w,\widehat\delta\rangle$,
and the marked periods, writing $\beta_c=\beta^{\rm part}+c$,
\[
\Pi_c=\begin{pmatrix}6\mu&\tau&1&0\\\beta^{\rm part}+c&\mu&0&1\end{pmatrix},
\qquad w=e^{-\pi i c},\qquad (\beta^{\rm part}+\tau)(p_0)=0.
\]
Let $\pi:\mathfrak H\to B\setminus\{p_0\}$ be the marked
$(3,4,\infty)$ orbifold uniformization, with marked deck generators
$g_1,g_2$ of orders three and four. Let $z_j\in\mathfrak H$
be the fixed point of $g_j$ above $p_j$. Thus
$\pi(z_j)=p_j$, $g_jz_j=z_j$,
$\rho=\tau(z_1)=e^{\pi i/3}$ and $\tau(z_2)=i$.
The marked values are
$\mu(z_1)=(2-\rho)/3$ and $\mu(z_2)=(1-i)/2$
(\SeriesRef{fam:src:C.1.1}). The local roots are
$r^3=t\circ\pi$ at $z_1$ and $s^4=(t-1)\circ\pi$ at $z_2$.
We take $r>0$ on $0<t<1$, where $\tau$ follows the unit-circle
arc from $\rho$ to $i$. At $z_2$ we use the branch with
$\operatorname{Im}\mu_s(0)>0$ and
\[
 \tau=i+\frac{i}{\pi\sqrt{E_4(i)}}s^2+O(s^4),
\]
where $E_4$ is the normalized Eisenstein series of weight four;
on the marked arc $t>1$ one has $\tau=iy$, $y>1$
(\SeriesRef{I:def:finite-coordinates}).
Base meridians are counterclockwise. The analytic root decks are clockwise,
$s_j(g_jz)=e^{-2\pi i/m_j}s_j(z)$, with $g_1g_2g_0=1$;
the resulting homology monodromies are $A_1,A_2,M_0$.
The primitive invariant vectors are
$\varepsilon=(1,2,-4,0)^t$ and
$\varepsilon'=(1,3,-3,0)^t$.
For an admissible finite affine datum, with $3\nmid l_1$ and
$2\nmid l_2$, we write
$v_1=l_1\varepsilon+k_1\widehat\delta$ and
$v_2=l_2\varepsilon'+k_2\widehat\delta$, using the coefficients
$l_1,l_2$ of the peripheral presentation.
Define
\[
 D_{\rm part}=\operatorname{Im}\beta^{\rm part}
       -\frac{6(\operatorname{Im}\mu)^2}{\operatorname{Im}\tau},
 \qquad M=\max_{B\setminus\{p_0\}}D_{\rm part}
       =-\frac3\pi\log6.
\]
The admissible parameter inequality is $\operatorname{Im}c<-M$;
the maximum of $D_{\rm part}$ is attained only at $t=1$, as proved in \SeriesRef{fam:src:C.2.3}.
The standard multiplicities are $m_1=3,m_2=4$,
$f^*p_j=m_jS_j$, and $W=f^{-1}(p_0)$ is reduced, with normalization
$\eta:\mathrm dP_6\to W$.
Let $\gamma:\Lambda\to\Z$ be the first coordinate and
$\Lambda_{\rm tor}=\Z\langle\widehat w,\widehat\delta\rangle$.
For the cusp collar section $\sigma$, define
$l_0=\gamma(c(\sigma))$, where
$c(\sigma)=[\lambda_\sigma]\in\Lambda/\Lambda_{\rm tor}$ is determined
by a holomorphic lift $\widetilde\sigma$ on the cusp cover:
\[
 t^{-1}=e^{2\pi i s_0},\qquad
 \widetilde\sigma(s_0+1)-\widetilde\sigma(s_0)
     =\Pi_c(s_0+1)\lambda_\sigma.
\]
Changing the lift adds $(I-M_0^{-1})\lambda$ to $\lambda_\sigma$;
this lies in $\Lambda_{\rm tor}$, on which $\gamma$ vanishes.
This is the winding convention of
\cite[Definition~6.6 and Proposition~6.7, pp.~39--40]{AC26}.
Since $\gamma(v_j)=l_j$, the peripheral presentation in
\cite[Theorem~7.17, p.~57]{AC26}, using the clockwise analytic
generators, gives
\[
 p=12l_0-4l_1-3l_2,\qquad \pi_1(X)\simeq\Z/|p|\Z.
\]
Those peripheral generators lift the inverses of our
counterclockwise base meridians. The standard fixed data are
those of \SeriesRef{fam:src:A.5.1}; the arbitrary affine vectors
and full cusp collar section are those of
\SeriesRef{fam:src:B.2.1}.
An unqualified fibre is its full Cartier scheme; its reduction is always
specified explicitly. The coefficient group is displayed in each homology or cohomology group.
Complex manifolds carry their complex orientations.
\end{definition}


\Needspace{13\baselineskip}
\part{Cotangent obstructions and fibre geometry}\label{part:II:cotangent-fibre-geometry}

\section{Index, Picard components and endpoint cohomology}\hypertarget{tag-033X}{}\label{sec:II:01}


\begin{proposition}[Tangent index and Picard components]\hypertarget{tag-01BN}{}\label{cdp:st:C.1.1}
\StatementRecord{cdp:st:C.1.1}{proved}{01BN}{the definitions and constructions stated in the proof}

Let $Y$ be a connected smooth compact complex threefold with
$b_1(Y)=b_2(Y)=0$. Every holomorphic line bundle $L$ on $Y$ satisfies
\[
 \chi(Y,T_Y\otimes L)=\frac12\int_Yc_3(T_Y)
 =1-\frac{b_3(Y)}2.
\]
The orientation in the integral is the complex orientation. The group
$\Pic(Y)$ has finitely many connected components, each a torsor for the
complex vector group $H^1(Y,\OO_Y)$, and
$\pi_0\Pic(Y)\simeq H^2(Y;\Z)$ through the first Chern class.
\end{proposition}


\begin{proof}

\begin{enumerate}[label=\arabic*.,leftmargin=*]
\item The universal coefficient theorem gives
$H^1(Y;\Z)=\Hom(H_1(Y;\Z),\Z)=0$ from $b_1(Y)=0$.
The group $H^2(Y;\Z)$ is finite because $Y$ has a finite CW model and
$b_2(Y)=0$. In the exponential sequence, normalized by
$\exp(2\pi i\,\cdot)$, the relevant maps are
\[
0=H^1(Y;\Z)\longrightarrow H^1(Y,\OO_Y)
 \xrightarrow{\exp} \Pic(Y)
 \xrightarrow{c_1} H^2(Y;\Z).
\]
They are group homomorphisms; the middle map is holomorphic on the
Picard group. Consequently its image $\Pic^0(Y)$ is isomorphic to the
complex vector space $H^1(Y,\OO_Y)$. The next map in the exponential
sequence is $H^2(Y;\Z)\to H^2(Y,\OO_Y)$; it is zero because its
source is finite and its target is a complex vector space. Thus
$c_1$ is surjective, and the component group is exactly $H^2(Y;\Z)$.

\item Poincar\'e duality gives $H^4(Y;\Q)=0$. Thus
$c_1(T_Y)$, $c_2(T_Y)$ and $c_1(L)$ vanish in rational cohomology.
The degree-six component of
$\operatorname{ch}(T_Y)\exp(c_1(L))\operatorname{td}(T_Y)$ is then
$c_3(T_Y)/2$. The Hirzebruch--Riemann--Roch formula for the Dolbeault
index gives the first equality. Chern--Gauss--Bonnet identifies
$\int_Yc_3(T_Y)$ with $e(Y)$. Since $Y$ is a connected closed oriented
six-manifold, its Betti numbers satisfy
$b_0=b_6=1$, $b_1=b_5=0$, and $b_2=b_4=0$. Hence
$e(Y)=2-b_3(Y)$. The pairing on $H^3(Y;\Q)$ is nondegenerate and
alternating, so $b_3(Y)$ is even, consistently with the integral index.


\end{enumerate}
\end{proof}


\begin{proposition}[A lower bound for twisted tangent cohomology]\hypertarget{tag-01BP}{}\label{cdp:st:C.1.3}
\StatementRecord{cdp:st:C.1.3}{proved}{01BP}{\SeriesRef{cdp:st:C.1.1}}

Let $Y$ be a smooth connected compact complex threefold with
$b_1(Y)=b_2(Y)=0$, carrying a nonzero effective Cartier divisor $D$.
For every $L\in\Pic(Y)$,
\[
 h^2(Y,T_Y\otimes L)\geq
 \max\left\{0,1-\frac{b_3(Y)}2\right\}.
\]
In particular, $b_3(Y)=0$ gives $H^2(Y,T_Y\otimes L)\ne0$ for
every $L$. A threefold in this class possessing a twist with vanishing
second tangent cohomology has $b_3(Y)\geq2$.
\end{proposition}


\begin{proof}
\begin{enumerate}[label=\arabic*.,leftmargin=*]
\item 
Fix a component $P$ of $\Pic(Y)$ and a representative $L_P$.
Apply \cite[Corollary~1.3, p.~79]{CDP98} to
$T_Y\otimes L_P$ and $T_Y^*\otimes K_Y\otimes L_P^{-1}$.
Tensoring by $L_P$ identifies $\Pic^0(Y)$ with $P$, and
Serre duality identifies the latter degree-zero vanishing with
the degree-three vanishing for the inverse Picard parameter.
The two dense open sets intersect in a dense open subset of $P$.
On this subset, \cref{cdp:st:C.1.1} gives
\[
 h^2(Y,T_Y\otimes L)-h^1(Y,T_Y\otimes L)
 =1-b_3(Y)/2.
\]
The asserted lower bound holds there. By \cref{cdp:st:C.1.1}, the component $P$ is a torsor for
$H^1(Y,\OO_Y)$. Exponentiating additive \v Cech representatives
of this vector space and tensoring by $L_P$ gives a holomorphic
family of line bundles on $Y\times P$. Proper cohomology
semicontinuity makes the locus satisfying the lower bound on
$h^2$ closed analytic in $P$. It contains the dense open subset
above and therefore all of $P$. 

\end{enumerate}
\end{proof}


\label{cdp:sec:E}

\begin{proposition}[Vertical divisors and the field of meromorphic functions]\hypertarget{tag-01GG}{}\label{cdp:st:E1}
\StatementRecord{cdp:st:E1}{proved}{01GG}{the definitions and constructions stated in the proof}

Let $Y$ be a connected compact complex manifold of dimension $n\geq2$
with $H^2(Y,\mathbb R)=0$. Let $g:Y\to C$ be a surjective holomorphic
map to a connected compact Riemann surface, with connected fibres.
Suppose that its smooth fibres over a nonempty open subset of $C$ are
K\"ahler. Every effective divisor of $Y$ is supported on finitely many
fibres, and pullback induces an isomorphism of fields
\[
 g^*: \mathcal M(C)\xrightarrow{\ \sim\ }\mathcal M(Y).
\]
In particular $a(Y)=1$. 
\end{proposition}


\begin{proof}

\begin{enumerate}[label=\arabic*.,leftmargin=*]
\item For a prime divisor $D$, the image $g(D)$ is analytic by properness. It
is a point or all of $C$. In the latter case choose a regular value $c$
in the stated open subset for which $F=g^{-1}(c)$ is not contained in
$D$. The canonical section of $\mathcal O_Y(D)$ restricts to a nonzero
section with nonempty effective zero divisor $D_F$ on $F$. For a
K\"ahler form $\omega$ on $F$, the divisor class formula gives
\[
 0=\int_F c_1(\mathcal O_Y(D)|_F)\wedge\omega^{n-2}
   =\int_{D_F}\omega^{n-2}>0.
\]
The first equality follows from the restriction homomorphism
$H^2(Y,\mathbb R)\to H^2(F,\mathbb R)$; the last follows from positivity
on every component with its positive multiplicity. Thus $g(D)$ is a
point. A divisor has finitely many prime components, proving the first
assertion.

\item For $h\in\mathcal M(Y)$ remove from $C$ the finitely many values
supporting its polar divisor. On the inverse image of the complement,
$h$ is holomorphic: a meromorphic function on a manifold having no
codimension-one pole extends across its codimension-two indeterminacy
set by the local normality of the structure sheaf. Its restriction to
each compact connected fibre is constant. The identity
$g_*\mathcal O_Y=\mathcal O_C$, given by Stein factorization and
the normality of $C$, therefore gives a holomorphic descended function
there. To check meromorphic extension at the removed values, take the
closure of the graph of $h$ in $Y\times\mathbb P^1$ and apply the proper
map $g\times1$. Its image is a one-dimensional analytic subset of
$C\times\mathbb P^1$, generically the graph of the descended function.
Its component dominating $C$ has generic degree one, and hence is the
graph of a meromorphic function on $C$. Pullback equals $h$ on a dense
open subset and therefore in the meromorphic function field. Pullback
is injective because $g$ is surjective. The field isomorphism preserves transcendence degree.

\end{enumerate}
\end{proof}


\section{Direct images and divisor vanishing}\hypertarget{tag-P270}{}\label{sec:II:direct-image-comparisons}

\begin{cited}[The non-holomorphic reduction theorem]\hypertarget{tag-01GR}{}\label{cdp:st:E2}
\StatementRecord{cdp:st:E2}{proved}{01GR}{\cite[Theorem~3.1, arXiv v1, p.~9]{LRS2019}}
Let $Y$ be a smooth compact connected complex threefold with
$b_2(Y)=0$, and let $g:Y\dashrightarrow C$ be a nonconstant
meromorphic map to a compact nonsingular curve which does
not extend holomorphically. For every holomorphic vector
bundle $V$ on $Y$, a dense open subset of each component of
$\operatorname{Pic}(Y)$ satisfies
\[
H^q(Y,V\otimes L)=0\qquad(0\leq q\leq3).
\]
Consequently $\chi(Y,V\otimes L)=0$ for every
$L\in\operatorname{Pic}(Y)$, and $\int_Yc_3(Y)=0$.
\end{cited}

\begin{proof}
The vanishing on $\operatorname{Pic}^0(Y)$, the Euler-index
conclusion for every line in that component, and the third-Chern
conclusion are \cite[Theorem~3.1, arXiv v1, p.~9]{LRS2019}.
For a component represented by $L_*$, tensoring defines the
holomorphic isomorphism
\[
\operatorname{Pic}^0(Y)\longrightarrow
L_*\operatorname{Pic}^0(Y),\qquad Q\longmapsto L_*\otimes Q,
\]
with inverse $L\mapsto L_*^{-1}\otimes L$.
Apply the cited theorem to the holomorphic vector bundle
$V\otimes L_*$. The tensor associativity isomorphism
$(V\otimes L_*)\otimes Q\simeq V\otimes(L_*\otimes Q)$
identifies every cohomology group and every Euler index.
The component isomorphism carries the theorem's dense open
set to the claimed dense open set and proves the assertion
for every line in that component.
\end{proof}

\begin{proposition}[A length-one torsion direct image]\hypertarget{tag-01GJ}{}\label{cdp:st:E12}
\StatementRecord{cdp:st:E12}{proved}{01GJ}{the definitions and constructions stated in the proof}

Let $A$ be an elliptic curve with origin $0$, let
$P=\mathcal O_{A\times A}(\Delta-A\times\{0\}-\{0\}\times A)$,
and let $Y=A\times A\times\mathbb P^1$, with $g:Y\to A$
projection to the second factor. For
$V=P\boxtimes\mathcal O_{\mathbb P^1}(-2)$ one has
\[
 g_*V=R^1g_*V=0,\qquad R^2g_*V=\mathbb C_0,
 \qquad h^\bullet(Y,V)=(0,0,1,0).
\]
The sheaf in degree two has length one. For $t\ne0$ the fibre
cohomology is zero in every degree; for $t=0$ it is $(0,1,1)$.
The torsion in $R^2g_*V$ contradicts the local-freeness assertion of
\citet[Lemma~1.5, pp.~79--80]{CDP98}.
\end{proposition}


\begin{proof}

\begin{enumerate}[label=\arabic*.,leftmargin=*]
\item Each nontrivial degree-zero line on an elliptic curve has no
sections: a nonzero section would have an effective divisor of
degree zero and hence trivialize it. Riemann--Roch and Serre
duality then give zero first cohomology. The trivial line has
cohomology $(1,1)$. Put $p:A\times A\to A$ and $M=P\otimes\operatorname{pr}_1^*
\mathcal O_A(0)$. Its fibres have degree one, so
$R^1p_*M=0$ and $p_*M$ is a line bundle. The exact sequence for
the section $\{0\}\times A$ gives the two-term complex
\[
 Rp_*P=\bigl[p_*M\longrightarrow\mathcal O_A\bigr].
\]
The arrow evaluates the unique section of the moving degree-one
line at the origin. Its zero is exactly $t=0$ and is simple:
in a local coordinate the moving zero divisor is the diagonal,
with equation $x-t$. Thus the arrow is the inclusion
$\mathcal O_A(-0)\to\mathcal O_A$, and
$Rp_*P=\mathbb C_0[-1]$. Tensoring with
$R\Gamma(\mathbb P^1,\mathcal O(-2))=\mathbb C[-1]$
gives $Rg_*V=\mathbb C_0[-2]$. Fibre K\"unneth gives
the stated fibre cohomology.

\item For any proper flat family over a nonsingular curve and a locally
free sheaf $V$, multiplication by a local parameter gives the
base-change exact sequence
\[
0\to (R^qg_*V)_c/t(R^qg_*V)_c
 \to H^q(Y_c,V|_{Y_c})
 \to (R^{q+1}g_*V)_c[t]\to0.
\]
It follows directly from the long direct-image sequence of
$0\to V\xrightarrow{t}V\to V|_{Y_c}\to0$. Here $[t]$ denotes
the kernel of multiplication by $t$. Local constancy of all
fibre cohomology dimensions implies local freeness and base
change by Grauert's theorem; generic vanishing alone permits
the torsion explicitly present above.

\end{enumerate}
\end{proof}

\begin{proposition}[Divisor twisting and embedded associated points]\hypertarget{tag-01GO}{}\label{cdp:st:E17}
\StatementRecord{cdp:st:E17}{proved}{01GO}{the definitions and constructions stated in the proof}

For an effective Cartier divisor $D$ on a compact complex
space $Z$ and a locally free sheaf $\mathcal V$, the spaces
$H^0(Z,\mathcal V(-kD))$, viewed inside $H^0(Z,\mathcal V)$,
stabilize for $k\gg0$ at
\[
 \{s\in H^0(Z,\mathcal V):\operatorname{supp}(s)\cap|D|
       =\varnothing\}.
\]
In particular they eventually vanish when $Z$ has no embedded
associated components and $D$ meets every irreducible component,
as in \cite[Lemma~3.5, p.~10]{LRS2019}.
For the graph of a nonconstant meromorphic map from a smooth connected threefold
$Y$ to $\mathbb P^1$, every full fibre is isomorphic to a Cartier divisor in $Y$
and is Cohen--Macaulay with no embedded associated components.
\end{proposition}


\begin{proof}

\begin{enumerate}[label=\arabic*.,leftmargin=*]
\item The descending spaces in the statement are
finite dimensional, hence stabilize. A section in their
intersection belongs to $\bigcap_k I_D^k\mathcal V_z=0$ at
every $z\in|D|$, by Krull intersection in the analytic local
ring. Its coherent support is therefore disjoint from $D$.
Conversely a section with support disjoint from $D$ vanishes
near $D$ and is a section of $\mathcal V(-kD)$ for every $k$.
When there are no
embedded associated components, every nonzero submodule of a
locally free sheaf has support containing an irreducible
component of $Z$. Such a support cannot avoid $D$ under the
stated intersection hypothesis.

\item Locally write the meromorphic map as $[a:b]$, where $a,b$ have
no common prime factor. Its graph $Z$ has equation
$av-bu=0$ in $Y\times\mathbb P^1$. At $c=[u_c:v_c]$, the maps
\[
 \{av_c-bu_c=0\}\longrightarrow Z_c,\qquad y\longmapsto(y,c),
\]
and the restriction of the projection to $Y$ are inverse holomorphic
maps of analytic spaces: both local rings are
$\mathcal O_Y/(av_c-bu_c)$. The local equation is a nonzero
element of the regular local ring of $Y$, hence a nonzerodivisor.
The quotient is Cohen--Macaulay and Gorenstein and has no embedded
associated components. This identifies the full analytic fibre,
including its nilpotents and multiplicities.

\end{enumerate}
\end{proof}

\begin{example}[An embedded point surviving every divisor twist]\hypertarget{tag-P267}{}\label{ex:II:embedded-point-divisor-twist}
\StatementRecord{ex:II:embedded-point-divisor-twist}{proved}{P267}{\cite[Proposition~1.1, p.~78]{CDP98}}
There is a compact connected one-dimensional complex space $Z$
with one embedded associated point and an effective Cartier
divisor $D$ meeting its irreducible component such that
\[
h^0(Z,\mathcal O_Z(-kD))=\max(1-k,0)+1,\qquad
h^1(Z,\mathcal O_Z(-kD))=\max(k-1,0)
\]
for every integer $k$, and every higher cohomology group is zero.
\end{example}

\begin{proof}
Let $C=\mathbb P^1$, let $p,q$ be distinct points, and form the
square-zero complex space
\[
 |Z|=C,\qquad\mathcal O_Z=\mathcal O_C\oplus\mathbb C_p,
 \qquad(a,u)(b,v)=(ab,av+bu).
\]
Near $p$ its analytic local ring is
$\mathbb C\{z,e\}/(ze,e^2)$, and away from $p$ it is the
ordinary smooth curve. These charts glue through the quotient
by $e$, so this specifies a compact connected complex space.
It has one irreducible component, of positive dimension, and an
embedded associated point at $p$. The divisor $D=q$ is effective
Cartier: its local equation is a coordinate at $q$ and a unit
near $p$. For every integer $k$ there is the exact, split
description as an $\mathcal O_C$-module
\[
 \mathcal O_Z(-kD)=\mathcal O_C(-kq)\oplus\mathbb C_p.
\]
It gives every cohomology group:
\[
 h^0=\max(1-k,0)+1,\qquad h^1=\max(k-1,0),\qquad
 h^j=0\quad(j\geq2).
\]
In particular the nonzero section of the square-zero summand
survives for every $k\geq1$, contrary to
\citet[Proposition~1.1, p.~78]{CDP98}.
\end{proof}

\begin{lemma}[Higher images of the exceptional divisor of a pencil]\hypertarget{tag-P265}{}\label{lem:II:pencil-exceptional-thickenings}
\StatementRecord{lem:II:pencil-exceptional-thickenings}{proved}{P265}{}
Let $Y$ be a smooth compact complex threefold, and resolve a
meromorphic pencil without a fixed divisorial component by a
composition of blowups with smooth centres
$\sigma:\widehat Y\to Y$. Write $h:\widehat Y\to\mathbb P^1$
for the resolved map, $P$ for the line of the pencil, and $E$
for the effective divisor of its pulled-back base ideal, so that
\[
h^*\mathcal O_{\mathbb P^1}(1)
 =\sigma^*P\otimes\mathcal O_{\widehat Y}(-E).
\]
For every positive integer $t$ one has
\[
R^1\sigma_*\mathcal O_{tE}=R^2\sigma_*\mathcal O_{tE}=0.
\]
For every holomorphic vector bundle $W$ on $Y$,
$H^2(tE,\sigma^*W|_{tE})=0$.
\end{lemma}

\begin{proof}
A smooth blowup has structure-sheaf derived image equal to
the original structure sheaf; composition gives
$R\sigma_*\mathcal O_{\widehat Y}=\mathcal O_Y$.
Pull back the evaluation sequence for
$\mathcal O_{\mathbb P^1}(t)$ and tensor by $\sigma^*P^{-t}$.
It gives
\[
\begin{split}
0\longrightarrow&\bigl(\mathcal O_{\widehat Y}(E)\otimes
              \sigma^*P^{-(t+1)}\bigr)^{\oplus t}
 \longrightarrow(\sigma^*P^{-t})^{\oplus(t+1)}\\
 &\longrightarrow\mathcal O_{\widehat Y}(-tE)
 \longrightarrow0.
\end{split}
\]
Pullback preserves exactness because the evaluation sequence
has locally free quotient. The middle term has zero
$R^2\sigma_*$ by the projection formula. Every coherent sheaf
has zero $R^3\sigma_*$ since the fibres have dimension at most
two. The resulting direct-image sequence gives
$R^2\sigma_*\mathcal O_{\widehat Y}(-tE)=0$. Applying it to
\[
0\to\mathcal O_{\widehat Y}(-tE)\to
\mathcal O_{\widehat Y}\to\mathcal O_{tE}\to0
\]
gives the asserted degree-one and degree-two vanishings.
The support of $\sigma_*\mathcal O_{tE}$ is contained in the
base locus of the pencil, which has dimension at most one
because the pencil has no fixed divisorial component. The
projection formula identifies
$R\sigma_*(\mathcal O_{tE}\otimes\sigma^*W)$ with
$\sigma_*\mathcal O_{tE}\otimes W$ in degree zero. Leray and
the dimension of that support give the asserted second-cohomology
vanishing.
\end{proof}

\begin{example}[An elliptic exceptional higher image]\hypertarget{tag-P266}{}\label{ex:II:exceptional-elliptic-direct-image}
\StatementRecord{ex:II:exceptional-elliptic-direct-image}{proved}{P266}{\cite[Example~3.4(2), p.~10]{LRS2019}}
Blow up a point $p\in\mathbb P^3$, then blow up a smooth
plane cubic $C$ in the exceptional $\mathbb P^2$. For the
second exceptional divisor $S$ and the composite map
$\sigma$ to $\mathbb P^3$, one has
\[
h^\bullet(S,\mathcal O_S)=(1,1,0),\qquad
R^1\sigma_*\mathcal O_S=\mathbb C_p.
\]
\end{example}

\begin{proof}
This construction is \cite[Example~3.4(2), p.~10]{LRS2019}.
The exceptional divisor is the projectivization of the rank-two
normal bundle of $C$ in the first blowup. Its projection
$r:S\to C$ is a $\mathbb P^1$-bundle, with
$r_*\mathcal O_S=\mathcal O_C$ and higher direct images zero.
The restriction of $\sigma$ to $S$ is the composite of $r$,
the map $C\to\{p\}$ and the closed inclusion
$\{p\}\hookrightarrow\mathbb P^3$. Since $C$ has genus one,
Leray gives the displayed cohomology and the length-one
skyscraper in degree one.
\end{proof}

\begin{remark}[Adjunction on a full Cartier fibre]\hypertarget{tag-P272}{}\label{rem:II:full-fibre-adjunction}
\StatementRecord{rem:II:full-fibre-adjunction}{proved}{P272}{\cite[p.~14, line following (3.5)]{LRS2019}}
For an effective Cartier divisor $F$ in a smooth threefold $Y$,
including a full nonreduced fibre, its dualizing line is
\[
\omega_F=(K_Y\otimes\mathcal O_Y(F))|_F.
\]
This corrects the sign of $F$ in the formula following
equation~(3.5) on p.~14 of \cite{LRS2019}.
\end{remark}

\begin{proof}
Apply $\mathcal Hom_Y(-,K_Y)$ to
$0\to\mathcal O_Y(-F)\to\mathcal O_Y\to\mathcal O_F\to0$.
Its degree-one cokernel is
$(K_Y\otimes\mathcal O_Y(F))|_F$.
Cartier duality identifies this cokernel
$\mathcal Ext^1_Y(\mathcal O_F,K_Y)$ with $\omega_F$.
The calculation uses the full nonzerodivisor equation of $F$.
\end{proof}



\section{Integral specialization, coherent restrictions and component counts}\hypertarget{tag-033Y}{}\label{sec:II:02}

\subsection{Component counts and projective examples}\label{sec:II:component-counts}

\begin{proposition}[Components of exceptional fibres]\hypertarget{tag-01G0}{}\label{cdp:st:D.1.2}
\StatementRecord{cdp:st:D.1.2}{proved}{01G0}{\SeriesRef{cdp:eq:D1}}

Let $f:Y\to C$ be a proper holomorphic surjection from a connected
smooth compact complex threefold to a compact Riemann surface, with
connected fibres. Let $\Delta$ be a nonempty finite set containing
the critical values, $s=|\Delta|$, $g=g(C)$,
$D=f^{-1}(\Delta)_{\rm red}$, and let $r$ be the number of
irreducible components of $D$. Write $F$ for a smooth fibre and
\[
t'=\dim_{\Q}H_1(F;\Q)_{\pi_1(C\setminus\Delta)}
   =\dim_{\Q}H^1(F;\Q)^{\pi_1(C\setminus\Delta)}.
\]
For $b_2(Y)=0$ the exact component count is
\[
\boxed{\quad r=2g+s-1+t'-b_1(Y).\quad}
\]
Under the additional integral hypotheses
$H^1(Y;\Z)=H^2(Y;\Z)=0$, the base is $\PP^1$ and
\[
r=s-1+t',\qquad
H_1(F;\Z)_{\pi_1(C\setminus\Delta)}\text{ is free abelian}.
\]
The freeness assertion concerns the displayed coinvariant quotient of
$H_1(F;\Z)$.
\end{proposition}





\begin{proof}
Put $U=Y\setminus D$ and $C^\circ=C\setminus\Delta$.
The rational component formula and the isomorphism
\begin{equation}
H_1(U;\Z)\simeq
\Z^{2g+s-1}\oplus H_1(F;\Z)_{\pi_1(C^\circ)}
\hypertarget{tag-019T}{}\label{cdp:eq:D1}
\end{equation}
are proved in \cite[Section~10.5, pp.~90--91, including
Lemma~10.7]{AC26}. The splitting in \eqref{cdp:eq:D1} comes from
lifting free generators of $\pi_1(C^\circ)$. The dual of its
coinvariant summand is the invariant space in the definition of $t'$:
a functional descends precisely when it vanishes on every
$\rho(g)v-v$.

For the integral assertion, the universal coefficient theorem first
makes $H_1(Y;\Z)$ finite from $H^1(Y;\Z)=0$, and then gives
$H_1(Y;\Z)=0$ from
$\operatorname{Ext}^1(H_1(Y;\Z),\Z)\subset H^2(Y;\Z)=0$.
It also makes $H_2(Y;\Z)$ finite. The compact-support sequence is
\[
H^4(Y;\Z)\longrightarrow H^4(D;\Z)
\longrightarrow H_c^5(U;\Z)
\longrightarrow H^5(Y;\Z)\longrightarrow H^5(D;\Z)=0.
\]
The complex orientation identifies $H^4(D;\Z)$ with $\Z^r$:
the singular and pairwise intersection loci have real dimension at
most two, so top cohomology has one generator for each oriented
irreducible component. Poincar\'e duality identifies $H^4(Y;\Z)$
with the finite group $H_2(Y;\Z)$, $H_c^5(U;\Z)$ with
$H_1(U;\Z)$, and $H^5(Y;\Z)$ with $H_1(Y;\Z)=0$.
The map from the finite group to $\Z^r$ is zero. Hence the
connecting map identifies $H_1(U;\Z)$ with $\Z^r$, and
\eqref{cdp:eq:D1} makes the coinvariant summand free abelian.
Finally, the Leray edge injection
$H^1(C;\Q)\hookrightarrow H^1(Y;\Q)=0$ gives $g=0$.
Substitution in the rational formula gives $r=s-1+t'$.
\end{proof}

\begin{proposition}[Two projective torus fibrations]\hypertarget{tag-01G3}{}\label{cdp:st:D.1.5}
\StatementRecord{cdp:st:D.1.5}{proved}{01G3}{\citet[Section~10.5, p.~91]{AC26}}
Let $E_0=F_0=E'=\C/(\Z+i\Z)$ and $a=1/2\in E_0$.
The involution
\[
(e,z)\longmapsto(e+a,-z)
\]
acts freely on $E_0\times F_0$. Put
$S=(E_0\times F_0)/\langle(a,-1)\rangle$,
$Y=S\times E'$, and $C=E_0/\langle a\rangle$.
The induced map $Y\to C$, denoted by $f$, is a smooth torus fibration with fibre
$F_0\times E'$ and monodromy
$\operatorname{diag}(-I_2,I_2)$ around the nontrivial deck character.
Its integral homology coinvariants are
$\Z^2\oplus(\Z/2\Z)^2$. After removing $s\geq1$ fibres,
\[
H_1(U;\Z)=\Z^{s+3}\oplus(\Z/2\Z)^2.
\]
The sum of the first Betti numbers of the punctured base and the
fibre is $s+5$.

Let $\eta$ be the nontrivial order-two degree-zero line bundle on
$C$ associated to the double cover $E_0\to C$. The coherent direct
images are
\[
R^1f_*\OO_Y=\eta\oplus\OO_C,\qquad R^2f_*\OO_Y=\eta.
\]
The cohomology dimensions are $h^q(Y,\OO_Y)=(1,2,1,0)$.
Restriction from $H^1(Y,\OO_Y)$ to the two-dimensional
$H^1(F_0\times E',\OO)$ has image precisely the $E'$ direction
and has rank one. These homomorphisms are induced by the finite
\'etale quotient and projection, and $b_1(Y)=4$.

For an example over $\PP^1$, let
$h:R\to\PP^1$ be the relatively minimal elliptic surface with
Weierstrass equation
\[
y^2=x^3+(t^4+1)x+t^6.
\]
Its fundamental line bundle is $\OO(1)$ and its discriminant has
twelve simple zeroes. Moreover, $R^1h_*\OO_R=\OO(-1)$.
For $Y'=R\times E'\to\PP^1$, write $f$ for the projection.
The product K\"unneth maps give
\[
R^1f_*\OO_{Y'}=\OO\oplus\OO(-1),\qquad
R^2f_*\OO_{Y'}=\OO(-1),
\]
and restriction in degree one again has rank one. This projective
threefold has $b_1=2$, $b_2=11$ and $a=3$.
\end{proposition}
\begin{proof}
The first construction specializes the bielliptic example of
\citet[Section~10.5, p.~91]{AC26} to the displayed involution.
Translation by $a$ has no fixed point, so the quotient is smooth.
The map $[e]\mapsto([e,0],0)$ is a section of the fibration. Over
the punctured base its fundamental group is therefore the semidirect
product of the fibre lattice and the free base group of rank $s+1$.
The map $(k,q)\mapsto([q],[k])$ identifies its abelianization with
the direct sum of the base abelianization and the fibre coinvariants:
the relations are $\rho(q)k-k$, which here are twice the two
$F_0$ generators. The same section over the complete elliptic base
gives the stated value of the first Betti number.

On the covering product, the three antiholomorphic one-form
directions have signs $(1,-1,1)$. Descent of the fibre cohomology
gives the two direct images, and taking invariants in the exterior
algebra gives the stated cohomology dimensions. The base direction
restricts to zero on a fibre, while the invariant fibre direction is
exactly that of $E'$. This proves the description of the restriction
map.

For the Weierstrass surface, the coordinates $u=1/t$,
$x'=u^2x$, $y'=u^3y$ give
\[
y'^2=x'^3+(1+u^4)x'+1.
\]
The relative residue differential transforms by
$dx'/(2y')=u^{-1}dx/(2y)$, so the fundamental line is
$\OO_{\PP^1}(1)$. Its discriminant is
\[
-16P(t),\qquad P(t)=31t^{12}+12t^8+12t^4+4=Q(t^4),
\qquad Q(z)=31z^3+12z^2+12z+4.
\]
The identity
\[
\left(\frac5{36}-\frac{31z}{54}\right)Q(z)
+\left(\frac1{27}-\frac{7z}{324}+\frac{31z^2}{162}\right)Q'(z)=1
\]
shows that the roots of $Q$ are simple; its constant term is
nonzero. Thus $P$ has twelve simple zeroes. At infinity the
transformed discriminant is $-16\cdot31$, so the fibre there is
smooth. The twelve singular fibres are nodal cubics. Relative
duality gives the stated first direct image. Leray gives
$h^1(R,\OO_R)=h^2(R,\OO_R)=0$; the twelve nodal fibres give
$e(R)=12$, and hence $b_1(R)=0$ and $b_2(R)=10$.
K\"unneth now gives both direct images and the Betti numbers of
the product. Its degree-one coherent cohomology is the $E'$
summand, giving the restriction map. The product is projective
of dimension three, so its algebraic dimension is three.
\end{proof}

\begin{proposition}[Components in Wang's torus fibration]\hypertarget{tag-01G4}{}\label{cdp:st:D.1.6}
\StatementRecord{cdp:st:D.1.6}{proved}{01G4}{\SeriesRef{cdp:st:D.1.13}; \cite{Wang2026}}
Wang's threefold $X_2(\lambda)$ has integral homology
$H_0=H_6=\Z$, $H_2=H_3=\Z/2\Z$, and the other reduced
homology groups zero. Its integral cohomology satisfies
$H^1=H^2=0$. Its fibration has four exceptional values, over which
the reduced fibres are irreducible: the multiplicity-three
bielliptic fibre, two rank-one fillings, and the rank-two filling.
In the rank-one filling over the $I_2$ value, the primitive period
exponent is $(2,1)$ and the fibre has one component.

In its cohomology basis the four monodromies are
\[
\begin{aligned}
B_*&=\begin{pmatrix}1&0&-6&-2\\0&-1&-1&1\\0&1&0&-1\\0&0&0&1\end{pmatrix},&\qquad
B_2&=\begin{pmatrix}1&0&6&3\\0&1&2&1\\0&0&1&0\\0&0&0&1\end{pmatrix},\\
B_+&=\begin{pmatrix}1&0&0&0\\0&2&1&0\\0&-1&0&0\\0&0&0&1\end{pmatrix},&\qquad
B_-&=\begin{pmatrix}1&0&0&-1\\0&1&1&0\\0&0&1&0\\0&0&0&1\end{pmatrix}.
\end{aligned}
\]
Their product in the displayed order is the identity, and their
common invariant cohomology group has rank one. On homology the
matrices are $B_i^{-t}$; the Smith diagonal of their difference
matrix is $(1,1,1,0)$. Thus $t'=1$ and the component count is
$r=4-1+1=4$. The expression $s-1+b_1(F)$ equals seven.

\end{proposition}
\begin{proof}
The construction and its integral homology are
\citet[Theorems~1.1 and~3.1, pp.~2 and~8]{Wang2026}.
The rank-one, rank-two and logarithmic fillings are
\citet[Propositions~2.7, 2.8 and~2.10, pp.~4, 6 and~7]{Wang2026}.
The primitive period vector is given in
\citet[Lemma~2.6, p.~4]{Wang2026}: the monomial change
$(z,v)\mapsto(x,y)=(v,v^2/z)$, with inverse
$(x,y)\mapsto(x^2/y,x)$, sends the vector to $(1,0)$.
The integral universal coefficient sequence gives the asserted
cohomology vanishing, since
$\operatorname{Ext}(H_1,\Z)=\operatorname{Hom}(H_2,\Z)=0$.

The matrices are \citet[Proposition~4.4, equation~(20), p.~11]{Wang2026}.
Direct multiplication gives their ordered product. For a column
$(a,b,c,d)^t$, invariance under $B_-$ gives $c=d=0$, and invariance
under $B_+$ then gives $b=0$. Every matrix fixes the first basis
vector. Preservation of evaluation under transport gives
$B_i^tA_i=I$, so $A_i=B_i^{-t}$ on the homology lattice.
Put $D_i=A_i-I$. Their columns satisfy
\[
D_-e_1=e_4,\qquad D_-e_2=-e_3,\qquad
D_+e_2=-e_2-e_3.
\]
Every image column has first coordinate zero, and these three
columns generate its full kernel. The coinvariant quotient is
therefore identified with $\Z$ by first coordinate, with inverse
$n\mapsto[ne_1]$. This gives the Smith diagonal and the invariant
rank. The component formula of \cref{cdp:st:D.1.2} gives $r=4$.
\end{proof}

\subsection{Cohomology and line-bundle restrictions}\label{sec:II:coherent-restrictions}

\begin{proposition}[Real cohomology and the structure sheaf]\hypertarget{tag-01G5}{}\label{cdp:st:D.1.7}
\StatementRecord{cdp:st:D.1.7}{proved}{01G5}{the definitions and constructions stated in the proof}

For a connected compact complex space $S$, including a nonreduced
one, the coefficient map
\[
H^1(S;\mathbb R)\longrightarrow H^1(S,\OO_S)
\]
is injective as a map of real vector spaces.
\end{proposition}


\begin{proof}
First take $S$ reduced. Represent a real cohomology class by a
locally constant real \v{C}ech cocycle $a_{ij}$. If its image is zero,
after refining the open cover there are holomorphic functions $u_i$
with $a_{ij}=u_i-u_j$. The imaginary parts of the $u_i$ agree on
overlaps and give a global continuous pluriharmonic function.
Pull it back to a compact resolution of each irreducible component.
The maximum principle makes each such pullback constant; intersection
of components and connectedness of $S$ make the constants equal.
Thus, after subtracting one imaginary constant from all $u_i$,
each $u_i$ is real-valued and holomorphic, hence locally constant.
They form a real coboundary for $a_{ij}$. For nonreduced $S$, a
class zero in $H^1(S,\OO_S)$ also becomes zero after applying
$\OO_S\to\OO_{S_{\rm red}}$. The underlying topological space
and its real cohomology do not change under reduction, so the
reduced assertion proves the result.

\end{proof}

\begin{proposition}[Injectivity of integral generization]\hypertarget{tag-01G6}{}\label{cdp:st:D.1.8}
\StatementRecord{cdp:st:D.1.8}{proved}{01G6}{\SeriesRef{cdp:st:D.1.2}}

Let $f:Y\to C$ be as in \ref{cdp:st:D.1.2}, and let $j:C^\circ\hookrightarrow C$
be the smooth-fibre locus. The natural generization homomorphism
\[
R^1f_*\Z_Y\longrightarrow
j_*\bigl((R^1f_*\Z_Y)|_{C^\circ}\bigr)
\]
is injective. Every nonzero global section of $R^1f_*\Z_Y$ is
nonzero in $H^1(f^{-1}(c);\Z)$ at every $c\in C$.
\end{proposition}


\begin{proof}

\begin{enumerate}[leftmargin=*]
\item It suffices to check a critical value $c$. Represent a stalk class
by $a\in H^1(V;\Z)$ for $V=f^{-1}(\Delta_c)$ and a sufficiently
small disc $\Delta_c$. Put $V^*=f^{-1}(\Delta_c\setminus\{c\})$.
Assume that the restriction of $a$ to a smooth fibre is zero.
The homotopy sequence for $V^*\to\Delta_c^*$ shows that
$a|_{V^*}$ is the pullback of an integer multiple $k$ of the
positive meridian class of the punctured disc.

\item Choose a smooth point of one irreducible component of the central
divisor, away from all its other components. There the map has
the form $t=z^m u$ with $m>0$ and $u$ a holomorphic unit.
A small positive normal circle has base winding $m$ and bounds
its normal disc in $V$. Evaluating $a$ on that circle gives
$km=0$, whence $k=0$.
The inclusion $V^*\hookrightarrow V$ induces a surjection of
fundamental groups: a loop can be perturbed off a complex divisor,
including its higher-codimension singular strata. Therefore
$a|_{V^*}=0$ implies $a=0$.
Taking the direct limit over the discs proves the asserted stalk
injection; proper topological base change identifies its source
with $H^1(f^{-1}(c);\Z)$.

\item If a global section vanished at one point, it would vanish on a
neighbourhood of that point. Its restriction to the connected
smooth locus is a section of a local system, so it would vanish
everywhere on that locus. The just-proved injection then makes
it zero everywhere, as required.

\end{enumerate}
\end{proof}

\begin{proposition}[Restriction maps and splitting degrees]\hypertarget{tag-01G2}{}\label{cdp:st:D.1.4}
\StatementRecord{cdp:st:D.1.4}{proved}{01G2}{\SeriesRef{cdp:st:D.1.2}}

For any $f:Y\to C$ as in \ref{cdp:st:D.1.2} put $E=R^1f_*\OO_Y$.
The coherent Leray sequence gives
\begin{equation}
0\longrightarrow H^1(C,\OO_C)\xrightarrow{f^*}H^1(Y,\OO_Y)
\xrightarrow{e}H^0(C,E)\longrightarrow0.               \hypertarget{tag-019U}{}\label{cdp:eq:D2}
\end{equation}
At a smooth fibre where cohomological base change holds, the
restriction homomorphism is exactly the composite
\begin{equation}
H^1(Y,\OO_Y)\xrightarrow{e}H^0(C,E)
\xrightarrow{\operatorname{ev}_c}E\otimes\C(c)
\xrightarrow{\simeq}H^1(F_c,\OO_{F_c}).               \hypertarget{tag-019V}{}\label{cdp:eq:D3}
\end{equation}
By \cite[Theorem~2.2, pp.~4--5]{Viaclovsky2026Combs}, $E$ is locally free. For $C=\PP^1$, write
$E=\bigoplus_{j=1}^q\OO(a_j)$.
Then the rank of \eqref{cdp:eq:D3} at a general point is exactly
$\#\{j:a_j\geq0\}$. For a locally free sheaf $E$ of rank two, the nonnegative splitting
degrees in \cite[Proposition~5.5, p.~688]{CDP2020} are therefore
equivalent to surjectivity of the general-fibre restriction map.
\end{proposition}


\begin{proof}
Stein factorization gives $f_*\OO_Y=\OO_C$. A curve has no coherent
cohomology in degree two, so the low-degree Leray sequence is \eqref{cdp:eq:D2}.
The edge map is induced by restriction to inverse images of open
sets in $C$. Restricting these same representatives to $F_c$ gives
\eqref{cdp:eq:D3}, which proves the asserted identification of the maps.
Local freeness is \cite[Theorem~2.2, pp.~4--5]{Viaclovsky2026Combs}. For a line bundle
$\OO(a)$ on $\PP^1$, its evaluation has rank one at every point
for $a\geq0$, and its space of global sections is zero for $a<0$.
Applying this separately to the direct summands gives the formula.

\end{proof}

\Needspace{5\baselineskip}
\begin{proposition}[Cup products and cohomological flatness]\hypertarget{tag-01FX}{}\label{cdp:st:D.1.12}
\StatementRecord{cdp:st:D.1.12}{proved}{01FX}{\cite[\S3, equations (9)--(12), pp.~376--377]{Borcea1984}}

For a proper connected-fibre holomorphic map $f:Y\to C$ from a
compact complex threefold with $b_2(Y)=0$ and general fibre a
two-torus or a primary Kodaira surface, the cup-product homomorphism
is an isomorphism
\[
\bigwedge^2R^1f_*\OO_Y\xrightarrow{\simeq}R^2f_*\OO_Y.
\]
Every full fibre has structure-sheaf cohomology dimensions
$(h^0,h^1,h^2)=(1,2,1)$, and these two direct images satisfy
base change at every fibre. These are the torus and primary Kodaira cases of
\cite[Proposition~5.1, pp.~687--688]{CDP2020}.
\end{proposition}





\begin{proof}
The cup-product index argument for torus-fibred six-spheres is due to
\citet[Theorem~2.3, pp.~5--6]{Viaclovsky2026Combs}. We apply it under
the displayed hypotheses and verify the cup product for primary Kodaira fibres.

\begin{enumerate}[leftmargin=*]
\item Write $E=R^1f_*\OO_Y$, which is locally free of rank two by \cite[Theorem~2.2, pp.~4--5]{Viaclovsky2026Combs},
and $G=R^2f_*\OO_Y$, which has rank one. Put
$d=\deg E$, $g=\deg(G/G_{\rm tor})$, and
$\ell=\operatorname{length}(G_{\rm tor})$.
Since $c_1(Y)=0$ rationally, Hirzebruch--Riemann--Roch gives
$\chi(Y,\OO_Y)=\int_Yc_1c_2/24=0$.
Coherent Leray and Riemann--Roch on $C$ give, with the rank terms
cancelled,
\[
0=\chi(\OO_C)-\chi(E)+\chi(G)=g+\ell-d.
\]
The cup product $\det E\to G$ is an isomorphism on a general
fibre. For a two-torus this is its exterior-algebra description.
For a primary Kodaira surface, use invariant $(1,0)$-forms
$\alpha,\eta$ with $d\alpha=0$ and
$d\eta=c\alpha\wedge\overline\alpha$, $c\ne0$.
Their conjugates are $\overline\partial$-closed. The class
$[\overline\alpha\wedge\overline\eta]$ is nonzero because its
Serre pairing with the holomorphic two-form $\alpha\wedge\eta$
is the nonzero integral of a volume form. This also proves that
$[\overline\alpha],[\overline\eta]$ are independent; since
$h^1(\OO)=2$ and $h^2(\OO)=1$, it proves the required cup-product
isomorphism. The invariant forms and their differential are
the forms in \cite[\S3, equations (9)--(12), pp.~376--377]{Borcea1984}.
The composite of $\det E\to G$ with the torsion-free quotient is a
nonzero homomorphism of line bundles, so $g\geq d$. The displayed
identity and $\ell\geq0$ force $g=d$, $\ell=0$. A nonzero map
of line bundles of equal degree has no zero divisor, proving the
claimed global isomorphism.

\item Top-degree base change gives
$G\otimes\C(c)\simeq H^2(Y_c,\OO_{Y_c})$, of dimension one.
The multiplication-parameter exact sequence then identifies
$E\otimes\C(c)$ with $H^1(Y_c,\OO_{Y_c})$, since $G$ has no
parameter-torsion. \cite[Theorem~2.2, pp.~4--5]{Viaclovsky2026Combs} gives $h^0=1$
on every full fibre.

\end{enumerate}
\end{proof}

\begin{proposition}[Line bundles of infinite order on every fibre]\hypertarget{tag-01G7}{}\label{cdp:st:D.1.9}
\StatementRecord{cdp:st:D.1.9}{proved}{01G7}{\SeriesRef{cdp:eq:D5}, \SeriesRef{cdp:st:D.1.8}, \SeriesRef{cdp:st:D.1.7}}

Let $Y$ be a connected compact complex manifold with
$b_1(Y)=b_2(Y)=0$, and let $f:Y\to\PP^1$ be a proper
holomorphic surjection with connected fibres. There is a class
$\beta\in H^1(Y,\OO_Y)$ such that, for every irrational real number
$r$, the line bundle
\[
L_r=\exp(2\pi i\,r\beta)
\]
has infinite order on every scheme-theoretic fibre of $f$.
Every connected component of $\Pic(Y)$ likewise contains a real
one-parameter family with this property at every irrational real
parameter.

\end{proposition}


\begin{proof}

\begin{enumerate}[leftmargin=*]
\item The low-degree integral Leray sequence and connectedness of the
fibres give an injection with finite cokernel
\begin{equation}
d_2:H^0(\PP^1,R^1f_*\Z_Y)
 \hookrightarrow H^2(\PP^1;\Z)=\Z.                   \hypertarget{tag-019X}{}\label{cdp:eq:D5}
\end{equation}
Indeed $H^1(Y;\Z)=0$ and $H^2(Y;\Z)$ is finite.
Thus the image in \eqref{cdp:eq:D5} is $d\Z$ for an integer $d>0$; for the
stronger integral hypothesis $H^2(Y;\Z)=0$ one has $d=1$.
Choose $\gamma$ mapped to $d$. The local meridian argument in
the proof of Proposition~\ref{cdp:st:D.1.8} applies in every complex
dimension: the local equation is $t=z^m u$ with $u$ a unit and $m>0$,
and evaluation on a positive normal circle gives $km=0$, hence $k=0$.
Thus each $\gamma_c\in H^1(f^{-1}(c);\Z)$ is nonzero.
The coefficient morphism $\Z_Y\to\OO_Y$ sends $\gamma$ to a
global section of $R^1f_*\OO_Y$. Coherent Leray and
$H^1(\PP^1,\OO)=H^2(\PP^1,\OO)=0$ identify this space with
$H^1(Y,\OO_Y)$; let $\beta$ be its inverse image.

\item For every full fibre $Y_c$, naturality of the edge maps and of
restriction gives
\[
\beta|_{Y_c}=\iota_c(\gamma_c)\in H^1(Y_c,\OO_{Y_c}), 
\]
where $\iota_c$ is the coefficient map. This is the image under the natural coherent base-change
homomorphism. If $L_r|_{Y_c}$ had finite order, the exponential
sequence on the full fibre would give an integer $n>0$ and
$\lambda\in H^1(Y_c;\Z)$ with
\[
nr\,\iota_c(\gamma_c)=\iota_c(\lambda).
\]
Proposition~\ref{cdp:st:D.1.7} implies $nr\gamma_c=\lambda$ in real cohomology.
The group $H^1(Y_c;\Z)$ is a finitely generated free abelian group,
and $\gamma_c$ is a nonzero element of it. A nonzero integral
coordinate of $\gamma_c$ now implies that $r$ is rational, a
contradiction. This works for every $c$ and proves the simultaneous
assertion in $\Pic^0(Y)$.

\item For a component represented by a line bundle $A$, choose $N>0$
with $A^N\in\Pic^0(Y)$ and write $A^N=\exp(2\pi i b)$.
Then
\[
A_r=A\otimes\exp\bigl(2\pi i(r\beta-b)/N\bigr)
\]
belongs to that component and satisfies $A_r^N=L_r$. A group
element has finite order exactly when any fixed positive power of
it has finite order. Therefore each $A_r|_{Y_c}$ is non-torsion
for every irrational real $r$, proving the componentwise statement.

\end{enumerate}
\end{proof}

\begin{proposition}[Invariant classes in the restriction image]\hypertarget{tag-01FV}{}\label{cdp:st:D.1.10}
\StatementRecord{cdp:st:D.1.10}{proved}{01FV}{\SeriesRef{cdp:st:D.1.2}, \SeriesRef{cdp:st:D.1.7}}

Let $f:Y\to\PP^1$ be a holomorphic algebraic reduction of a smooth compact
complex threefold with $b_1(Y)=b_2(Y)=0$ and $a(Y)=1$.
At a general smooth fibre, the image of the restriction
$H^1(Y,\OO_Y)\to H^1(F,\OO_F)$ contains the complex span of
the coefficient homomorphism
\begin{equation}
\iota_F:H^1(F;\Z)^{\pi_1(C^\circ)}
 \longrightarrow H^1(F,\OO_F),                         \hypertarget{tag-019W}{}\label{cdp:eq:D4}
\end{equation}
which is induced by $\Z_F\hookrightarrow\OO_F$.
The final spanning step in \cite[Theorem~3.1, proof, final paragraph, p.~86]{CDP98}
reaches the whole target exactly when the indicated invariant lattice
spans $H^1(F,\OO_F)$ over $\C$.
\end{proposition}


\begin{proof}

\begin{enumerate}[leftmargin=*]
\item Let $D_1,\ldots,D_r$ be the prime components above $\Delta$, and
choose $N>0$ annihilating the finite group $H^2(Y;\Z)$.
The logarithms, in the exponential identification
$\Pic^0(Y)=H^1(Y,\OO_Y)$, of
$\OO_Y(ND_1),\ldots,\OO_Y(ND_r)$ generate a free abelian
subgroup of rank $r-s+1=t'$.
Consider the homomorphism
$\Z^r\to\Pic^0(Y)$, $(n_i)\mapsto\OO_Y(N\sum n_iD_i)$.
A relation means that $N\sum n_iD_i$ is the divisor of a
meromorphic function. Since $f$ is the algebraic reduction, that
function is $f^*h$ for $h\in\mathcal M(\PP^1)$.
Its divisor is a sum of full fibres with total base degree zero.
Such relations have rank $s-1$, and those given by divisors of
rational functions supported on $\Delta$ already have this rank.
\Cref{cdp:st:D.1.2} gives the last equality.

\item Over $U=f^{-1}(C^\circ)$ each of these divisor line bundles has
its canonical trivialization. Exactness of the exponential sequence
therefore places its logarithm in the coefficient image of
$H^1(U;\Z)$. Put $E=R^1f_*\OO_Y$. Coherent Leray identifies restriction
$H^1(Y,\OO_Y)\to H^1(U,\OO_U)$ with restriction of global
sections of $E$; this is injective because \cite[Theorem~2.2, pp.~4--5]{Viaclovsky2026Combs} makes $E$ locally free.
Thus the rank-$t'$ group of logarithms remains injective on $U$.

\item Topological Leray on the punctured curve gives
\[
0\to H^1(C^\circ;\Z)\to H^1(U;\Z)
\to H^1(F;\Z)^{\pi_1(C^\circ)}\to0.
\]
The first group's coefficient image vanishes because
$H^1(C^\circ,\OO)=0$. Conversely a class whose coefficient image
in $H^1(U,\OO)$ vanishes restricts to zero in $H^1(F,\OO_F)$;
Proposition~\ref{cdp:st:D.1.7} makes its integral restriction to $F$ zero. Thus the
coefficient image of $H^1(U;\Z)$ identifies exactly with the
invariant lattice in \eqref{cdp:eq:D4}. The logarithm group maps to a subgroup
of the same rank $t'$, hence of finite index, in that lattice.
Its complex span equals the span of the entire lattice, and the
restriction image is a complex vector subspace. This proves the
assertion and the statement about the last spanning step.

\end{enumerate}
\end{proof}

\begin{proposition}[Generic line bundles on every fibre]\hypertarget{tag-01FW}{}\label{cdp:st:D.1.11}
\StatementRecord{cdp:st:D.1.11}{proved}{01FW}{\SeriesRef{cdp:st:D.1.9}, \SeriesRef{cdp:st:D.1.7}, \SeriesRef{cdp:st:D.1.8}}

Let $f:Y\to\PP^1$ be a proper connected-fibre holomorphic map from
a connected compact complex manifold with $b_1(Y)=b_2(Y)=0$.
In every component of $\Pic(Y)$, the line
bundles which are non-torsion on every full fibre form a dense
$G_\delta$ set of full Lebesgue measure in its complex vector-space
coordinates. If evaluation of global sections of $R^1f_*\OO_Y$
has rank one, then $\Pic^0(Y)=\C$ and the exceptional set is
countable. For two-torus fibres this alternative holds whenever
the splitting has a negative summand.
This gives the general-line-bundle conclusion of
\citet[Proposition~5.6, p.~689]{CDP2020} under the displayed hypotheses.
\end{proposition}


\begin{proof}

\begin{enumerate}[leftmargin=*]
\item Let $E=R^1f_*\OO_Y$, a vector bundle by \cite[Theorem~2.2, pp.~4--5]{Viaclovsky2026Combs}, and use
$\gamma,\beta$ from \ref{cdp:st:D.1.9}. The section of $E$ corresponding to
$\beta$ has no zero. In fact its value maps naturally to the
nonzero class $\iota_c(\gamma_c)$ in the cohomology of every full
fibre; Proposition~\ref{cdp:st:D.1.7} and the dimension-independent local
meridian argument in the proof of Proposition~\ref{cdp:st:D.1.8}
prove nonvanishing. Write
$E=\bigoplus_j\OO(a_j)$ and let $k$ count the nonnegative degrees.
Evaluation has constant rank $k$ onto their direct-summand subbundle
$E_+$. Nonvanishing of $\beta$ gives $k\geq1$. When $k=1$, its
sole nonnegative degree is zero, since a positive-degree line bundle
has no nowhere-zero section. In that case $H^0(E)=\C\beta$.

\item Suppose first $k\geq2$, and put $V=H^1(Y,\OO_Y)=H^0(E)$ and
$N=\dim_{\C}V$. On a simply connected disc in the smooth-fibre
locus where degree-one base change holds, trivialize the integral
local system. For a flat integral
class $\lambda$, the equation
\[
z|_{F_c}=\iota_c(\lambda),\qquad (z,c)\in V\times\Delta
\]
defines an analytic set. The parameter values where
$\iota_c(\lambda)\in(E_+)_c$ form an analytic subset of the disc;
over it the fibres are affine spaces of dimension $N-k$.
The incidence set thus has complex dimension at most
$N-k+1\leq N-1$. Its projection to $V$ has Lebesgue
measure zero. It is also meagre: stratify and exhaust that set
by countably many compact sets; the images are closed sets with
empty interior, by their real dimension at most $2N-2$.
Take countably many discs, integral classes and positive integer
denominators to include every torsion condition.

\item There are finitely many excluded fibres, including the exceptional
values of degree-one base change. At each, the natural map
$E\otimes\C(c)\to H^1(Y_c,\OO_{Y_c})$ is injective. This follows
directly from the long exact sequence for multiplication by a local
base parameter: its cokernel is the parameter-torsion in
$R^2f_*\OO_Y$. The restriction $V\to H^1(Y_c,\OO_{Y_c})$
therefore has rank $k\geq2$. The torsion conditions are contained in
countably many affine subspaces of complex codimension at least
two. Combining these with the preceding projections gives a
meagre null set. The complement is dense and $G_\delta$, as claimed.

\item It remains to prove the case $k=1$.
Now $V=\C\beta$. On a simply connected disc in the same open locus,
the projection of a flat integral class $\lambda$ to
$E/(\OO\beta)$ is holomorphic. Unless it vanishes identically,
its zero set is discrete. Such zeroes give only countably many
parameters $z$ for which $\iota_c(\lambda)=z\beta_c$.
We prove that an identically vanishing projection gives a constant
parameter, including after analytic continuation.

\item Pass to the universal cover of the smooth locus, on which
$H^1(F;\Z)$ is marked. The identity theorem continues the relation
$\iota_c(\lambda)\in\C\iota_c(\gamma)$ everywhere on this cover.
If $\lambda$ and $\gamma$ are real-dependent, their ratio is a
constant rational number. Otherwise let
$W_{\Q}=\Q\gamma+\Q\lambda$. The coefficient map
$H^1(F;\R)\to H^1(F,\OO_F)$ is real-linearly injective by \ref{cdp:st:D.1.7}.
The images of the two real-independent classes span the entire
real two-plane underlying $\C\beta_c$, so $W_{\R}$ is precisely
its inverse image. Analytic continuation around a base loop
sends $\lambda$ to a class whose coefficient image again belongs
to this line. Hence it preserves $W_{\Q}$ and fixes $\gamma$.

\item Choose an integral basis $(\gamma_0,\lambda_0)$ of
$W_{\Q}\cap H^1(F;\Z)$ with $\gamma$ a positive integral multiple
of $\gamma_0$. The induced monodromy on this lattice fixes
$\gamma_0$ and preserves the orientation of $W_{\R}$: the complex
structure induced on that plane by the coefficient map varies
continuously along the loop, so its complex orientation cannot change.
It therefore sends $\lambda_0$ to $\lambda_0+k\gamma_0$ for an
integer $k$. Put
\[
\tau(c)=\frac{\iota_c(\lambda_0)}{\iota_c(\gamma_0)}.
\]
This holomorphic function on the universal cover never takes a
real value, by real injectivity of the coefficient map, and its
imaginary part has constant sign. Monodromy changes $\tau$ by
integers. Accordingly either $\exp(2\pi i\tau)$ or
$\exp(-2\pi i\tau)$ descends to a holomorphic function of absolute
value less than one on $\PP^1$ minus finitely many points. The
removable singularity theorem extends it to $\PP^1$, where it is
constant. Its logarithmic derivative gives $d\tau=0$. Thus the
original ratio is constant as asserted.

\item There are countably many integral classes on a countable disc
cover. The isolated-zero alternatives and the constant alternatives
just proved give a countable exceptional set on smooth fibres.
Each exceptional full fibre contributes another countable set,
because $z\mapsto z\beta_c$ is injective and its intersection with
the rational integral coefficient image is countable. This proves
the assertion in $\Pic^0(Y)$.

\item Finally for any component $A\Pic^0(Y)$, the map taking the $N$th
power, for an integer $N$ killing its component class, is the affine
isomorphism $z\mapsto b+Nz$ of the complex parameter space used
in \ref{cdp:st:D.1.9}. It preserves the torsion condition in every fibre, and
preserves countability, meagreness and measure-zero sets. This
proves all the componentwise assertions.

\end{enumerate}
\end{proof}

\begin{remark}[Dense torsion restrictions]
\hypertarget{tag-P211}{}\label{rem:II:dense-torsion-restrictions}
\StatementRecord{rem:II:dense-torsion-restrictions}{proved}{P211}{}
In the globally generated two-torus case of \cref{cdp:st:D.1.11},
the line bundles in $\Pic^0(Y)$ which are torsion on any fixed
smooth fibre form a dense subset. Indeed, \cref{cdp:st:D.1.4,cdp:st:D.1.12}
identify restriction with a surjective complex-linear map onto
$H^1(F,\OO_F)$. The coefficient map identifies $H^1(F;\R)$
with the underlying real vector space of $H^1(F,\OO_F)$, so the
image of $H^1(F;\Q)$ is dense there. Its inverse image under
restriction is dense because a surjective linear map between
finite-dimensional vector spaces is open. The exponential sequence
identifies that inverse image precisely with the torsion restrictions.
Consequently the line bundles which are non-torsion on every fibre
contain no nonempty open subset of $\Pic^0(Y)$ in this case.
\end{remark}

\Needspace{8\baselineskip}
\subsection{Coherent splitting in explicit torus families}\label{sec:II:coherent-splitting}

\begin{proposition}[Higher direct images of Wang's threefold]\hypertarget{tag-01FY}{}\label{cdp:st:D.1.13}
\StatementRecord{cdp:st:D.1.13}{proved}{01FY}{\SeriesRef{cdp:st:D.1.6}, \SeriesRef{cdp:st:D.1.12}, \SeriesRef{cdp:st:D.1.9}, \SeriesRef{cdp:st:D.1.4}, \SeriesRef{cdp:st:E1}}


For Wang's threefold $Y=X_2(\lambda)$ in \ref{cdp:st:D.1.6} one has
\[
K_Y\simeq f^*\OO_{\PP^1}(-1),\qquad
R^1f_*\OO_Y\simeq\OO\oplus\OO(-1),\qquad
R^2f_*\OO_Y\simeq\OO(-1).
\]
Consequently $h^q(Y,\OO_Y)=(1,1,0,0)$, and restriction to every
smooth two-torus fibre has rank one. This threefold satisfies
$H^1(Y;\Z)=H^2(Y;\Z)=0$ and $a(Y)=1$. Its rank-one
restriction map and direct-image splitting contradict the conclusions of
\cite[Corollary~5.2 and Proposition~5.5, pp.~688--689]{CDP2020}.
For the threefold $Y_0$ of \SeriesRef{maps:src:B.5.2}, one has
$H_2(Y_0;\Z)=0$, whereas $H_2(Y;\Z)=\Z/2\Z$; hence a
biholomorphism $Y_0\to Y$ would induce an impossible isomorphism of
these two homology groups.
\end{proposition}


\begin{proof}

\begin{enumerate}[label=\arabic*.,leftmargin=*]
\item We compute the canonical line through the local fillings.
The rational elliptic surface $\pi:S\to\PP^1$ in this construction
has Weierstrass equation
\[
y^2=x^3-3(1+t)x+(2+3t+2t^2)
\]
and fundamental line bundle $\OO(1)$; hence
$K_S=\pi^*\OO(-1)$. On the complement of the zero section in
$M=\OO_S(P-O)$, the determinant of the relative cotangent sequence
identifies its canonical line with the pullback of $K_S$.
In a local multiplicative coordinate $v$, this identification sends
a local two-form $\omega_S$ to $\omega_S\wedge d\log v$.
Changing $v$ by a nonzero holomorphic function on $S$ adds a
one-form from $S$ to $d\log v$, whose wedge with $\omega_S$ is zero.
This verifies the coordinate independence of the identification.
The contracting generator covers translation by $-6P$ on the
elliptic fibres. Translation preserves the relative elliptic
differential, and the extra base terms disappear after wedging
with the base differential. Its line-bundle lift changes $d\log v$
by a one-form pulled back from $S$. It therefore preserves this canonical
identification. The latter descends to
$K_Y=f^*\OO(-1)$ on the open quotient.

\item At either rank-one filling, let $q$ be the base coordinate and
$(x,y)$ the two multiplicative coordinates of the construction.
At $t=0$ the explicit change is $x=v$, $y=v^2/z$; its exponent
matrix is unimodular. At $\alpha_+$ no change is required.
Thus the preceding three-form is a nonzero holomorphic unit times
\[
dq\wedge d\log x\wedge d\log y.
\]
The nondegenerating $y$ direction becomes an elliptic curve after
quotienting by its contracting period. Its logarithmic differential
is a nonvanishing relative one-form; changes by a base differential
do not affect the displayed wedge. At a double curve the equation
is $uv=q$, and
$dq\wedge d\log u=-du\wedge dv$. Therefore the three-form extends
without a zero or pole over the full rank-one filling.

\item At $\alpha_-$ the rank-two filling uses a unimodular period lattice
and the standard toric charts $q=z_0z_1z_2$. A pair of independent
fibre characters, together with $q$, is a unimodular system of
characters in each smooth chart. The same form is, up to a nonzero
constant and a holomorphic unit,
\[
q\,d\log z_0\wedge d\log z_1\wedge d\log z_2
=dz_0\wedge dz_1\wedge dz_2.
\]
It extends without zero or pole, including at the triple points.

\item It remains to check the order-three logarithmic filling at infinity.
Put $s=1/t$, $x_s=s^2x$ and $y_s=s^3y$. Its Weierstrass
coefficients have orders three and four in $s$, giving type $IV^*$.
On the cubic base change $s=r^3$, put
$x_s=r^4X$, $y_s=r^6Y$. The good-reduction equation is
\[
Y^2=X^3-3r(r^3+1)X+(2+3r^3+2r^6),
\]
whose central elliptic curve is smooth. The canonical form satisfies
\[
ds\wedge\frac{dx_s}{y_s}
=3\,dr\wedge\frac{dX}{Y}.
\]
It has no zero or pole. For $\zeta=e^{2\pi i/3}$ the root deck
action has weights $r\mapsto\zeta r$, $X\mapsto\zeta^2X$,
$Y\mapsto Y$; the weights on $dr$ and $dX/Y$ multiply to one.
The additional invariant translation used in the logarithmic
transform preserves the holomorphic two-form on the torus cover.
Hence the three-form is invariant and descends through the free
cyclic action as a nonvanishing three-form. This proves that the
canonical identification from the open quotient extends everywhere,
giving $K_Y=f^*\OO(-1)$.

\item The map is flat and Gorenstein, since its source is smooth and its
base is a smooth curve. Relative Serre duality gives
\[
(R^2f_*\OO_Y)^*=f_*\omega_{Y/\PP^1}
=f_*f^*\OO(1)=\OO(1).
\]
Proposition~\ref{cdp:st:D.1.12} gives local freeness and
$\det(R^1f_*\OO_Y)=R^2f_*\OO_Y=\OO(-1)$.
The nowhere-zero section from \ref{cdp:st:D.1.9} embeds $\OO$ as a subbundle
of $R^1f_*\OO_Y$ with quotient $\OO(-1)$. The extension splits
because $H^1(\PP^1,\OO(1))=0$. Coherent Leray now computes all
four displayed structure-sheaf cohomology groups, and \ref{cdp:st:D.1.4}
identifies the degree-one restriction with rank-one evaluation.

\item The homology computation in \ref{cdp:st:D.1.6} and the universal coefficient
theorem give $H^1(Y;\Z)=H^2(Y;\Z)=0$. More explicitly,
$H_1(Y;\Z)=0$ gives
$H^2(Y;\Z)\simeq\operatorname{Hom}(H_2(Y;\Z),\Z)=0$,
and the universal-coefficient injection in degree three is
$\operatorname{Ext}^1(H_2(Y;\Z),\Z)\simeq\Z/2\Z
\hookrightarrow H^3(Y;\Z)$.
The general fibre is K\"ahler and $b_2(Y)=0$, so \ref{cdp:st:E1} proves that
this map is the algebraic reduction and $a(Y)=1$.
The threefold of \SeriesRef{maps:src:B.5.2} has zero second
integral homology, whereas $H_2(Y;\Z)=\Z/2\Z$. A biholomorphism
and its inverse induce mutually inverse homomorphisms on $H_2$;
the displayed groups therefore exclude such a biholomorphism.

\end{enumerate}
\end{proof}

\begin{proposition}[Constant integral monodromy and coherent splitting]\hypertarget{tag-01FZ}{}\label{cdp:st:D.1.15}
\StatementRecord{cdp:st:D.1.15}{proved}{01FZ}{the definitions and constructions stated in the proof}


Two smooth proper families of two-tori over $\PP^1$ have the same
constant marked $H^1(-;\Z)$ local system but have
$R^1f_*\OO=\OO^2$ and $R^1f_*\OO=\OO(1)^2$, respectively.
Thus the integral representation alone does not determine the
coherent splitting for general compact torus fibrations. Both
examples have $b_1=4$.
\end{proposition}


\begin{proof}

\begin{enumerate}[leftmargin=*]
\item Let $V=\OO_{\PP^1}(1)^{\oplus2}$. On the affine coordinate $t$,
four global sections of $V$ are the columns of
\[
P(t)=\begin{pmatrix}1&i&t&it\\-t&it&1&-i\end{pmatrix}.
\]
They are a real basis in every fibre. Indeed, writing
$z_1=a+ib$, $z_2=c+id$, their real linear combination is
$(z_1+t z_2,\overline z_2-t\overline z_1)$.
Its vanishing gives $(1+|t|^2)\overline z_2=0$ and then $z_1=0$.
At infinity, the $\OO(1)$ frame changes $P(t)$ to $t^{-1}P(t)$,
whose leading columns again form a real basis. The lattice
generated by these sections therefore acts freely and properly by
holomorphic fibre translations on $V$. Its quotient $Z$ is a
compact smooth holomorphic torus bundle over $\PP^1$.
The displayed real basis identifies it differentiably with
$\PP^1\times(\R^4/\Z^4)$ and marks its constant integral local
system.

\item The fibrewise period map includes the bundle of holomorphic
one-forms by $P(t)^t$. The Dolbeault coefficient map gives the
exact sequence
\[
0\longrightarrow\OO(-1)^2
 \xrightarrow{P(t)^t}\OO^4
 \longrightarrow R^1f_*\OO_Z\longrightarrow0.
\]
The last map can be represented explicitly by
\[
B(t)=\begin{pmatrix}-it&-t&i&1\\-i&1&-it&t\end{pmatrix}.
\]
The identity $B(t)P(t)^t=0$ holds polynomially. Its two rows have
rank two at every finite $t$, as is seen from their two-by-two
minors, and their leading coefficients have rank two at infinity.
Thus this is the exact complex of vector bundles
\[
0\longrightarrow\OO(-1)^2\longrightarrow\OO^4
 \xrightarrow{B}\OO(1)^2\longrightarrow0,
\]
proving $R^1f_*\OO_Z\simeq\OO(1)^2$. More explicitly, with
\[
J=\begin{pmatrix}0&1&0&0\\-1&0&0&0\\0&0&0&1\\0&0&-1&0\end{pmatrix},
\qquad S=\begin{pmatrix}0&1\\1&0\end{pmatrix},
\]
one has $B=SPJ$ and $PJP^t=0$, so $BP^t=0$.

\item For the constant family $Z_0=Z_{t=0}\times\PP^1$, the same
marked fibre lattice gives the identical constant integral local
system, whereas coherent K\"unneth gives
$R^1f_*\OO_{Z_0}=H^1(Z_{t=0},\OO)\otimes\OO=\OO^2$.
The different splittings are therefore compared by an explicit
identity of the integral representations and explicit coherent
quotient maps.

\end{enumerate}
\end{proof}

\section{Smooth fibres and Hopf bundles}\hypertarget{tag-033Z}{}\label{sec:II:03}

\subsection{Necessary properties and Euler numbers}\label{sec:II:fibre-properties}

\begin{proposition}[Smooth fibres of an algebraic reduction]\hypertarget{tag-01F4}{}\label{cdp:st:C.3.1}
\StatementRecord{cdp:st:C.3.1}{proved}{01F4}{\cite[Chapter~1, Theorem~3.1, p.~10]{Peters1990}; \cite[p.~2]{Teleman2008}; \cite[arXiv version~3, Theorem~1, p.~1]{BBK2019}}

Let $f:Y\to C$ be a holomorphic algebraic reduction of a connected
smooth compact complex threefold with $b_1(Y)=b_2(Y)=0$ and $a(Y)=1$.
The base is $C\simeq\PP^1$. Every smooth fibre $F$ has
\[
 c_1(K_F)=0\text{ in }H^2(F;\Q),\qquad
 c_2(F)=0\text{ in }H^4(F;\Z),\qquad
 \chi(F,\OO_F)=0,
\]
and contains no exceptional curve of the first kind. A very general
smooth fibre has $\kappa(F)\leq0$. The very general fibre's type belongs to the following
list: a complex two-torus, a bielliptic surface, a primary or secondary
Kodaira surface, a primary or secondary Hopf surface, or an Inoue
surface. Every listed type has Euler characteristic zero.
\end{proposition}


\begin{proof}

\begin{enumerate}[leftmargin=*]
\item Connected fibres give $f_*\Q_Y=\Q_C$ and the degree-one edge injection
$H^1(C;\Q)\hookrightarrow H^1(Y;\Q)$. Therefore the genus of $C$ is
zero. The holomorphic tangent sequence on a smooth fibre, with the
last map induced by the differential of $f$, is
\[
0\longrightarrow T_F\longrightarrow T_Y|_F
 \xrightarrow{df} \OO_F\otimes T_{C,f(F)}\longrightarrow0.
\]
The normal factor is a trivial holomorphic line bundle. Thus
$K_F=K_Y|_F$ and $c_2(F)=c_2(Y)|_F$. Rational $H^2(Y)$ and $H^4(Y)$
vanish, and $H^4(F;\Z)=\Z$ has no torsion. These facts give the two
Chern-class assertions. Noether's formula gives $\chi(\OO_F)=0$.
An exceptional curve $E\simeq\PP^1$ with $E^2=-1$ would have
$K_F\cdot E=-1$ by adjunction, contradicting the rational vanishing of
$c_1(K_F)$.

\item Fix $m>0$ and put $G_m=f_*(K_Y^m)$. This is a coherent torsion-free
sheaf on the smooth curve $C$: a section killed by a nonzero base
function vanishes on a dense open subset of $Y$ and hence everywhere.
It is consequently a vector bundle. On the nonempty open set of
cohomological base change its fibre is $H^0(F,K_F^m)$. A rank at least
two would give a meromorphic map
\[
 Y\dashrightarrow \PP_C(G_m^*)
\]
whose restriction to a general fibre is defined by at least two
independent sections of a line bundle and has positive-dimensional
image. Its image has dimension at least two. The target is projective:
a vector bundle on a compact Riemann surface is algebraic, and its
projectivization is a projective variety. The image is therefore
projective and its rational functions pull back to a subfield of
$\mathcal M(Y)$ of transcendence degree at least two. This contradicts
$a(Y)=1$. Thus $\rk G_m\leq1$ for every $m$. Removing the countable
union of the exceptional base-change sets gives the claimed very
general-fibre assertion.

\item The Enriques--Kodaira classification is given in
\cite[Chapter~1, Theorem~3.1, p.~10]{Peters1990}. The surface is minimal by the
adjunction argument. Rational and ruled surfaces have a rational
curve on which the canonical class has nonzero degree, so they are
excluded. K3 and Enriques surfaces have Euler numbers $24$ and $12$,
respectively. Properly elliptic and general-type surfaces have
positive Kodaira dimension. The remaining K\"ahler possibilities are
tori and bielliptic surfaces; the remaining Kodaira-dimension-zero
non-K\"ahler possibilities are the two types of Kodaira surfaces.
For class VII, $b_1(F)=1$, and $e(F)=0$ gives $b_2(F)=0$. The
classification of this case gives Hopf and Inoue surfaces; its
precise scope, including finite covers, is stated in
\cite[p.~2]{Teleman2008} and \cite[arXiv version~3, Theorem~1, p.~1]{BBK2019}.

\end{enumerate}
\end{proof}

\begin{proposition}[Euler numbers of fibre supports]\hypertarget{tag-01FB}{}\label{cdp:st:C.3.5}
\StatementRecord{cdp:st:C.3.5}{proved}{01FB}{\SeriesRef{cdp:st:C.3.1}}

For every threefold and fibration in \ref{cdp:st:C.3.1}, the Euler-number identity is
\[
\int_Y c_3(Y)=e(Y)=\sum_{c\in\Delta}e\bigl((Y_c)_{\rm red}\bigr),
\]
where $\Delta$ contains the finitely many critical values.
\end{proposition}


\begin{proof}
All the listed smooth surfaces have Euler number zero by \ref{cdp:st:C.3.1}.
Additivity of compactly supported Euler characteristic for the
stratification by the smooth-fibre locus and the finitely many
exceptional fibres gives
$e(Y)=e_c(C\setminus\Delta)e(F)+\sum_{c\in\Delta}e(Y_c)$.
The first term vanishes. A scheme-theoretic fibre and its reduction
have the same underlying topological space, so their topological
Euler numbers agree. Chern--Gauss--Bonnet gives the first equality.

\end{proof}

\begin{proposition}[Obstruction to K\"ahler resolutions]\hypertarget{tag-01FA}{}\label{cdp:st:C.3.4}
\StatementRecord{cdp:st:C.3.4}{proved}{01FA}{the definitions and constructions stated in the proof}

Let $Y$ be a smooth compact complex threefold with $b_2(Y)=0$, and let
$f:Y\to C$ be proper, holomorphic, surjective, with connected fibres.
Every irreducible component of a reducible reduced fibre has no
compact K\"ahler resolution. In particular, a fibre all of whose
components admit compact K\"ahler resolutions is irreducible.
\end{proposition}


\begin{proof}
\begin{enumerate}[label=\arabic*.,leftmargin=*]
\item 
Let $D$ be a component of a reducible fibre. The incidence graph of
the components is connected, so some other component $D'$ meets $D$.
Both are Cartier divisors on the smooth threefold. Let
$\mu:\widehat D\to D\hookrightarrow Y$ be a resolution by a compact
smooth surface. Pullback of the canonical section of $\OO_Y(D')$
is a nonzero holomorphic section of $\mu^*\OO_Y(D')$. Its zero
divisor $E$ is effective and nonzero: the intersection $D\cap D'$
is a nonempty divisor on the irreducible surface $D$, and pullback
to the resolution preserves its nonzero vanishing order along a
component dominating that intersection. On the other hand
$c_1(\mu^*\OO_Y(D'))=0$ in $H^2(\widehat D;\R)$.
A K\"ahler form $\omega$ on $\widehat D$ would give
\[
0=\int_{\widehat D}c_1(\OO_{\widehat D}(E))\wedge\omega
 =\int_E\omega>0,
\]
where $E$ carries its positive divisor multiplicities.
This contradiction proves the assertion for every component.\qedhere

\end{enumerate}
\end{proof}

\subsection{Secondary surfaces and integral cohomology}\label{sec:II:secondary-surfaces}

\begin{proposition}[A secondary Kodaira surface with torsion-free first homology]\hypertarget{tag-01F8}{}\label{cdp:st:C.3.2}
\StatementRecord{cdp:st:C.3.2}{proved}{01F8}{the definitions and constructions stated in the proof}


There is a secondary Kodaira surface $Q$ with
\[
H_1(Q;\Z)=\Z,\qquad H^2(Q;\Z)=0,
\]
whose canonical bundle has holomorphic order six.
\end{proposition}


\begin{proof}

\begin{enumerate}[label=\arabic*.,leftmargin=*]
\item Let $H$ be the real Heisenberg group in coordinates $(x,y,z)$ with law
\[
(x,y,z)(x',y',z')
=\left(x+x',y+y',z+z'+\frac{xy'-yx'}2\right).
\]
Let $\Gamma\subset H$ be generated by
$a=(1,0,0)$, $b=(0,1,0)$ and $c=(0,0,1)$. Thus
\[
\Gamma=\{(x,y,z):x,y\in\Z,\ z-xy/2\in\Z\},
\qquad [a,b]=c.
\]
It is a discrete cocompact subgroup. On $H\times\mathbb R_t$ write
\[
\tau=e^{\pi i/3},\quad
\eta=dz-\frac{x\,dy-y\,dx}{2},\quad
\alpha=dx+\tau\,dy,\quad
\beta=\eta-\frac{dx+dy}{2}+i\,dt.
\]
These two complex one-forms and their conjugates are linearly
independent over $\C$. Declaring $\alpha,\beta$ to have type $(1,0)$
defines a left-invariant almost complex structure. It is integrable:
$d\alpha=0$ and $d\beta=-dx\wedge dy$ has type $(1,1)$, so the
ideal generated by the $(1,0)$ forms has no $(0,2)$ obstruction.
The structure descends to
$P=\Gamma\backslash H\times(\mathbb R/\Z)$.
The nonzero form $\alpha\wedge\beta$ is $\bar\partial$-closed and
trivializes $K_P$. The projection induced by $x+\tau y$ is a
holomorphic elliptic bundle over $\C/(\Z+\tau\Z)$, with nonzero
central Chern class $[a,b]=c$. Its first Betti number is three, so
$P$ is a primary Kodaira surface.

\item Define the Lie-group automorphism
\[
\sigma(x,y,z)=(-y,x+y,z-y/2).
\]
It preserves $\Gamma$: the difference between the new central
coordinate and half the product of the new horizontal coordinates
is
\[
(z-xy/2)+xy+y(y-1)/2.
\]
This is integral on $\Gamma$. The inverse preserves the lattice too,
and direct substitution gives $\sigma^6=1$. Moreover
$\sigma^*\alpha=\tau\alpha$ and $\sigma^*\beta=\beta$.
Consequently
\[
g:P\longrightarrow P,\qquad
g([x,y,z],t)=([\sigma(x,y,z)],t+1/6)
\]
is a holomorphic automorphism of order six. Every nonidentity power
is fixed-point-free because of its translation of the $t$-circle.
The quotient $Q=P/\langle g\rangle$ is therefore a compact complex
surface and $P\to Q$ is finite \emph{\'etale} of degree six.

\item The contractible manifold $H\times\mathbb R$ is its universal cover.
The deck group is $\Gamma\rtimes_\sigma\Z$, where the infinite
generator is the lift $(\sigma,t\mapsto t+1/6)$; its sixth power is
the original integral translation in $t$. Since
$\Gamma^{\rm ab}=\Z^2$ and the induced matrix is
\[
A=\begin{pmatrix}0&-1\\1&1\end{pmatrix},\qquad
\det(A-I)=1,
\]
abelianization gives $H_1(Q;\Z)=\Z$.
The finite cover and $e(P)=0$ give $e(Q)=0$. Poincar\'e duality then
gives $b_2(Q)=0$, and the universal coefficient theorem gives
$H^2(Q;\Z)=0$ because $H_1(Q;\Z)$ has no torsion.
Finally $g^*(\alpha\wedge\beta)=\tau\alpha\wedge\beta$.
Every holomorphic section of $K_P^m$ is a constant multiple of
$(\alpha\wedge\beta)^m$, by compactness. Such a section descends
precisely when $6$ divides $m$. Thus $K_Q$ has holomorphic order six.
The surface has $b_1=1$ and a primary Kodaira surface as a finite
\'etale cover, so it is secondary Kodaira.

\end{enumerate}
\end{proof}

\begin{proposition}[A secondary Hopf surface with torsion-free first homology]\hypertarget{tag-01F9}{}\label{cdp:st:C.3.3}
\StatementRecord{cdp:st:C.3.3}{proved}{01F9}{the definitions and constructions stated in the proof}


There is a secondary Hopf surface $Q_H$ with
$H_1(Q_H;\Z)=\Z$ and $H^2(Q_H;\Z)=0$.
\end{proposition}


\begin{proof}

\begin{enumerate}[leftmargin=*]
\item Write a quaternion as $(q_0,q_1,q_2,q_3)$ and put
$\varphi=(1+\sqrt5)/2$. Let $G$ be the following set of 120 unit
quaternions: the eight signed coordinate units; the sixteen points
$(\pm1,\pm1,\pm1,\pm1)/2$; and the 96 points obtained by even
coordinate permutations of
$(0,\pm1,\pm\varphi,\pm\varphi^{-1})/2$, with independent signs.
Quaternion multiplication makes this the binary icosahedral group.
Representing each coordinate as $a+b\sqrt5$ with $a,b\in\Q$,
Hamilton multiplication and $g^{-1}=\overline g$ compute both the
products and the commutators $gh\overline g\,\overline h$ exactly.
Enumerating the ordered pairs of the displayed elements gives
products in the same set and a commutator image equal to that set.
Thus every element is a commutator and $G^{\rm ab}=0$.
The program is listed in \cref{tab:II:programs}.

\item Use the faithful holomorphic representation
\[
(a,b,c,d)\longmapsto
\begin{pmatrix}a+ib&c+id\\-c+id&a-ib\end{pmatrix}\in SU(2)
\]
and let the generator of $\Z$ act on $\C^2\setminus\{0\}$ by
$v\mapsto2v$. The action of $\Z\times G$ is free. An equality
$2^ngv=v$ first gives $n=0$ by norms, and a unit quaternion whose
matrix has eigenvalue one is the identity. Proper discontinuity
follows from the radial scaling and finiteness of $G$.
Polar coordinates identify the quotient smoothly with
\[
Q_H=(\C^2\setminus\{0\})/(2^{\Z}\times G)
\simeq (\mathbb R/(\log2)\Z)\times(S^3/G).
\]
It is compact, and its universal cover is $\C^2\setminus\{0\}$.
The fundamental group is $\Z\times G$, whose abelianization is
$\Z$. The quotient by $2^{\Z}$ alone is a primary Hopf surface
covering $Q_H$ with degree 120. The fundamental group of $Q_H$ is
noncyclic, so $Q_H$ is secondary Hopf. The product description gives
$e(Q_H)=0$, and the same duality and universal-coefficient argument
as before gives $H^2(Q_H;\Z)=0$.

\end{enumerate}
\end{proof}

\begin{remark}[First homology of secondary surfaces]
\hypertarget{tag-P212}{}\label{rem:II:secondary-surface-homology}
\StatementRecord{rem:II:secondary-surface-homology}{proved}{P212}{}
The finite \'etale quotients in \cref{cdp:st:C.3.2,cdp:st:C.3.3}
contradict the two assertions in \cite[Lemma~4.3, proof, p.~685]{CDP2020}
that secondary Kodaira and secondary Hopf surfaces always have
nonzero torsion in their first integral homology. Their first homology
groups are computed by abelianizing the displayed deck groups.
\end{remark}

\subsection{Bielliptic and elliptic Hopf fibres}\label{sec:II:elliptic-exclusions}

\begin{proposition}[Exclusion of bielliptic fibres]\hypertarget{tag-01FD}{}\label{cdp:st:C.3.7}
\StatementRecord{cdp:st:C.3.7}{proved}{01FD}{\cite[Proposition~2, pp.~605--608]{Kawai1969}; \cite[Theorem~11, p.~297]{Buchdahl1999}}


A connected smooth compact complex threefold $Y$ with $b_1(Y)=b_2(Y)=0$ has no
proper connected-fibre holomorphic map to a compact curve whose
general fibre is bielliptic.
\end{proposition}


\begin{proof}

Suppose that $f:Y\to C$ is a map with the stated properties.
Since $Y$ is reduced, $f$ factors as $f=\iota\circ f_{\mathrm{red}}$
through the inclusion $\iota:C_{\mathrm{red}}\hookrightarrow C$ of the
reduced image curve. Let $\nu:C^\nu\to C_{\mathrm{red}}$ be its
normalization. The reduced closure of the regular-locus lift graph in
$Y\times_{C_{\mathrm{red}}}C^\nu$ is finite and birational over $Y$.
Since $Y$ is normal, its projection to $Y$ is an isomorphism. The other
projection gives a holomorphic map $f^\nu:Y\to C^\nu$ with
$f_{\mathrm{red}}=\nu\circ f^\nu$. This map is surjective because its
proper image contains the regular open set of $C^\nu$. For each
$c\in C_{\mathrm{red}}$, the finitely many closed sets
$(f^\nu)^{-1}(c')$, for $c'\in\nu^{-1}(c)$, partition $f^{-1}(c)$;
each is also open in that fibre. Connectedness of $f^{-1}(c)$ forces
a single nonempty piece, so $f^\nu$ has connected fibres. Over the
regular locus $\nu$ is an isomorphism, and the general bielliptic fibre
is unchanged. Write $f:Y\to C$ for the lifted map $f^\nu:Y\to C^\nu$ in the
remainder of the proof.

\begin{enumerate}[label=\arabic*.,leftmargin=*]
\item A bielliptic surface has
$p_g=0$ and $h^0(\Omega^1)=1$. The relative Albanese construction of
\cite[Proposition~2, pp.~605--608]{Kawai1969} gives a compact complex surface
$B$, a holomorphic map $h:B\to C$, and a meromorphic map
$\Phi:Y\dashrightarrow B$ over $C$. On the complement $C^\circ$
of finitely many values, $B^\circ\to C^\circ$ is a smooth family
of elliptic curves and
\[
 \Phi|_{F_c}:F_c\longrightarrow B_c
\]
is the Albanese morphism of the bielliptic surface, with connected
elliptic fibres. The
construction forms the relative period-lattice quotient, glues its
local Albanese maps by fibre translations, and compactifies at each
missing value by a normalized finite base change and the projective
space associated with the direct images of a divisor which is the
inverse image of a local Albanese section. The sum of the finite orbit of that divisor is invariant under
the covering group. Resolve the normalized base change first if needed, so the
divisor defines a line bundle. The orbit sum restricts on every
general fibre to the pullback of a positive-degree divisor on its
Albanese elliptic curve; sufficiently large powers give the same
Albanese image. Its naturally linearized line bundle and direct
images carry the finite covering-group action, and their quotient
glues the construction to $B^\circ$.

\item Resolve $B$ over $C\setminus C^\circ$ and resolve the induced
meromorphic map by a sequence of blowups with smooth centres
supported over those same values. This produces a diagram
\[
 \widetilde Y\xrightarrow{\ g\ }\widetilde B,
 \qquad \pi:\widetilde Y\longrightarrow Y,
 \qquad (h\circ\rho)\circ g=f\circ\pi,
\]
where $\rho:\widetilde B\to B$ is a resolution and both spaces
are smooth and compact. The map $g$ has connected fibres:
its general fibres are the connected Albanese fibres just described,
and its Stein factorization has finite degree one over the normal
surface $\widetilde B$. Hence
$H^1(\widetilde B;\Q)\hookrightarrow H^1(\widetilde Y;\Q)$.
Blowing up smooth centres does not change first cohomology, so
$b_1(\widetilde B)=0$. The compact-surface K\"ahler criterion
\cite[Theorem~11, p.~297]{Buchdahl1999} therefore gives a K\"ahler form
$\omega_B$ on $\widetilde B$.

\item The blowup cohomology formula gives
\[
 H^2(\widetilde Y;\R)
 =\pi^*H^2(Y;\R)\oplus\bigoplus_j\R[E_j]
 =\bigoplus_j\R[E_j],
\]
where the $E_j$ are the exceptional prime-divisor classes (with
the corresponding total transforms for successive blowups).
Every $E_j$ is supported above one of the removed values in $C$.
Consequently its restriction to a general smooth fibre $F_c$,
identified with its unchanged inverse image in $\widetilde Y$,
is zero. In particular
\[
 [g^*\omega_B]|_{F_c}=0\quad\hbox{in }H^2(F_c;\R).
\]
Choose a K\"ahler form $\omega_F$ on the bielliptic surface $F_c$.
The Albanese morphism $g|_{F_c}$ is a submersion to the elliptic
curve $\widetilde B_c$. Thus its pullback of $\omega_B$ is a
nonzero semipositive $(1,1)$-form and
\[
 \int_{F_c}g^*\omega_B\wedge\omega_F>0.
\]
The vanishing cohomology class makes the same integral zero by
Stokes' theorem, a contradiction.

\end{enumerate}
\end{proof}

\begin{proposition}[Exclusion of elliptic Hopf general fibres]\hypertarget{tag-01FE}{}\label{cdp:st:C.3.8}
\StatementRecord{cdp:st:C.3.8}{proved}{01FE}{\SeriesRef{cdp:st:C.3.1}; \cite[Proposition~3.1 and its proof, pp.~6--7]{SS2002}}


The very general Hopf fibre of the algebraic reduction in \ref{cdp:st:C.3.1}
has algebraic dimension zero. In particular neither primary nor
secondary Hopf surfaces of algebraic dimension one can be its
general fibre.
\end{proposition}


\begin{proof}

\begin{enumerate}[label=\arabic*.,leftmargin=*]
\item A primary Hopf surface $H$ with $a(H)=1$ satisfies
$\kappa(H,K_H^{-1})=1$. The contraction
normal forms are diagonal, or
\[
A(z,w)=(\alpha z+\lambda w^n,\beta w),\qquad
\lambda\ne0,\quad \alpha=\beta^n.
\]
The invariant bivectors $zw\,\partial_z\wedge\partial_w$ in the
diagonal case and $w^{n+1}\,\partial_z\wedge\partial_w$ in the
second case give a nonzero effective anticanonical divisor.
The invariance is an equality under the derivative of the
displayed contraction: its determinant is $\alpha\beta$, and
the displayed coefficient transforms by the same scalar.
The normal forms and elliptic-fibration description are given in \cite[Proposition~3.1 and its proof, pp.~6--7]{SS2002}.

\item The algebraic reduction of this primary elliptic Hopf surface is a
holomorphic map to $\PP^1$, with smooth reductions of all fibres.
Every curve is vertical: its intersection with a general fibre
is zero since $H^2(H;\R)=0$, whereas a horizontal curve would
have positive intersection equal to its degree over the base.
Thus the nonzero anticanonical divisor is a positive sum of
reduced fibres. If $m_iD_i$ is the full fibre supported on
$D_i$, each $\OO(m_iD_i)$ is the pullback of $\OO_{\PP^1}(1)$.
A common multiple of the $m_i$ therefore makes a positive power
of $K_H^{-1}$ the pullback of a positive-degree line bundle on
$\PP^1$. Its Iitaka dimension is one.

\item For a secondary Hopf surface take a normal finite primary cover
$p:H'\to H$. The cover is unramified, so
$K_{H'}^{-1}=p^*K_H^{-1}$, and $a(H')=a(H)=1$ by invariance of
transcendence degree under a finite map. The anticanonical
section ring of $H$ is the invariant subring of that of $H'$.
The latter is integral over the invariant subring: a homogeneous
section $s$ satisfies the monic orbit polynomial
$\prod_{g\in\operatorname{Gal}(H'/H)}(T-gs)$. Their fraction
fields have the same transcendence degree, proving
$\kappa(H,K_H^{-1})=1$ also in this case.

\item For each positive integer $m$, the locus
\[
Z_m=\{c\in C^\circ:
 h^0(Y_c,K_Y^{-m}|_{Y_c})\ge2\}
\]
is a closed analytic subset of the smooth base $C^\circ$.
If the very general Hopf fibre had algebraic dimension one,
the preceding calculation would make the union of these
countably many loci contain the complement of a countable
subset of $C^\circ$. A proper closed analytic subset of a
curve is discrete and countable, so some $Z_m$ is the whole
connected smooth base. Coherent base change then gives
$\operatorname{rank}f_*K_Y^{-m}\ge2$.
The relative linear system defines a meromorphic map of $Y$
to $\PP(f_*K_Y^{-m})$ over $C$ with positive-dimensional image
on a general fibre. Its image has dimension at least two.
The projective bundle over the compact curve is projective,
so pulling back its meromorphic functions gives $a(Y)\ge2$,
contrary to the hypothesis. In fact the same argument shows
that, when the smooth fibres are Hopf surfaces, the values with
$a(Y_c)=1$ form at most a countable set.

\end{enumerate}
\end{proof}

\subsection{Kodaira fibres}\label{sec:II:kodaira-exclusions}

\begin{proposition}[Exclusion of secondary Kodaira general fibres]\hypertarget{tag-01FF}{}\label{cdp:st:C.3.9}
\StatementRecord{cdp:st:C.3.9}{proved}{01FF}{\SeriesRef{cdp:st:D.1.9}, \SeriesRef{cdp:st:D.1.7}, \SeriesRef{cdp:eq:C.SK.1}, \SeriesRef{cdp:eq:C.SK.2}, \SeriesRef{cdp:eq:C.SK.3}, \SeriesRef{cdp:eq:C.SK.4}, \SeriesRef{cdp:eq:C.SK.5}; \cite[Theorem~4 and Remark~1, pp.~377, 380]{Borcea1984}; \cite[\S1, pp.~2--3]{SS2002}}


A connected smooth compact complex threefold $Y$ with $b_1(Y)=b_2(Y)=0$ and
$a(Y)=1$ cannot have a holomorphic algebraic reduction
$f:Y\to\PP^1$ whose general fibre is a secondary Kodaira
surface.
\end{proposition}


\begin{proof}

\begin{enumerate}[label=\arabic*.,leftmargin=*]
\item Suppose that the smooth general fibres $S_c$ are secondary
Kodaira surfaces. They have $h^1(\OO_{S_c})=b_1(S_c)=1$.
The nonzero integral transgression section of \ref{cdp:st:D.1.9} has
coefficient image $\beta\in H^1(Y,\OO_Y)$ and gives a
nowhere-zero section of $\mathscr E=R^1f_*\OO_Y$.
This coherent sheaf is a line bundle by \cite[Theorem~2.2, pp.~4--5]{Viaclovsky2026Combs}. Thus
$\mathscr E=\OO_{\PP^1}$ and
$H^1(Y,\OO_Y)=\C\beta$. For every $z\in\C$ set
\begin{equation}
 M_z=\exp(2\pi\mathrm i z\beta)\in\Pic^0(Y).
\hypertarget{tag-011H}{}\label{cdp:eq:C.SK.1}
\end{equation}
Its restriction satisfies $M_{z+1}|_{S_c}\simeq M_z|_{S_c}$,
because $\beta|_{S_c}$ is the coefficient image of an
integral class. It is nonzero as a tangent vector of
$\Pic^0(S_c)$ by \ref{cdp:st:D.1.7}. The exponential sequence therefore
identifies $\Pic^0(S_c)$ with $\C^*$ and makes \eqref{cdp:eq:C.SK.1}, on
each smooth fibre, a nonconstant homomorphism from
$\C/\Z$ to this group.

\item Let $m>1$ be the order of $K_{S_c}$, and let
\[
 p_c:P_c\longrightarrow S_c
\]
be its connected cyclic canonical cover, with deck group
$G=\mu_m$. It is finite \emph{\'etale}. Its canonical bundle
is trivial, and it is non-K\"ahler, because a finite
\emph{\'etale} quotient of a K\"ahler surface is K\"ahler
by averaging a K\"ahler form. The classification of
surfaces with trivial canonical bundle consequently makes
$P_c$ a primary Kodaira surface. Write
$\pi_c:P_c\to E_c$ for its Albanese elliptic bundle, and
$d>0$ for its topological degree. These covers exist in
holomorphic families locally on the smooth base: the
evaluation of a local frame of $f_*K_{Y/\PP^1}^{m}$
trivializes $K_{Y/\PP^1}^{m}$, since it is nowhere zero on
each fibre, and taking its cyclic root gives the required
finite \emph{\'etale} covering family. The integers $m,d$
are locally constant and hence constant on the connected
smooth base.

\item The primary-surface invariant is
\begin{equation}
 Q_c=\Pic^0(P_c)/\pi_c^*\Pic^0(E_c)\simeq\C^*.
\hypertarget{tag-011I}{}\label{cdp:eq:C.SK.2}
\end{equation}
The Picard pullback in the denominator is injective:
$\pi_{c*}\OO_{P_c}=\OO_{E_c}$ and the projection formula
show that a line with trivial pullback is trivial. For
any line $L_c$ of degree $d$ on $E_c$, the pullback
$\pi_c^*L_c$ belongs to $\Pic^0(P_c)$, and its class in
$Q_c$ is independent of $L_c$. Choose the group
isomorphism $q_c:Q_c\to\C^*$ that places this
positive-degree class in the punctured unit disc, and
denote its value by $\alpha(c)$.
This fixes the possible inversion ambiguity of the
group coordinate. The descriptions in
\cite[Theorem~4 and Remark~1, pp.~377, 380]{Borcea1984} and
\cite[\S1, pp.~2--3]{SS2002} identify $\alpha$ as the
intrinsic holomorphic punctured-disc coordinate of the
moduli of primary Kodaira surfaces. More explicitly,
$P_c$ is the associated elliptic bundle of a line of
degree $d$ with fibre $\C^*/\langle\alpha(c)\rangle$;
quotienting by $\pi_c^*\Pic^0(E_c)$ removes the choice
of that line within its degree. In a holomorphic
family the local integral marking and the coefficient
map give this quotient a holomorphic $\C^*$ coordinate,
and the chosen positive-degree orientation is locally
constant. Borcea's complete period family gives these local holomorphic
coordinates. Thus $\alpha(c)$ is holomorphic.

\item The locally constructed canonical covers have the same
isomorphism class on overlaps, so this intrinsic
function glues on the smooth base. It has
$0<|\alpha(c)|<1$. It extends across the finitely many
omitted values by the removable singularity theorem
and is constant on $\PP^1$. Write its nonzero constant
value as $\alpha$.

\item Consider the exact homomorphism of connected Picard
groups
\begin{equation}
 \Pic^0(S_c)\xrightarrow{\ p_c^*\ }\Pic^0(P_c)
          \longrightarrow Q_c\xrightarrow{\ q_c\ }\C^*.
\hypertarget{tag-011J}{}\label{cdp:eq:C.SK.3}
\end{equation}
Its differential is nonzero. Indeed finite-cover trace
identifies $H^1(S_c,\OO_{S_c})$ with the $G$-invariant
line in $H^1(P_c,\OO_{P_c})$. The other line,
$\pi_c^*H^1(E_c,\OO_{E_c})$, has nontrivial $G$ action:
otherwise the unique holomorphic one-form of $P_c$
would descend to $S_c$, whereas a secondary Kodaira
surface has no holomorphic one-form. Hence the invariant
line has zero intersection with the denominator's
tangent line in \eqref{cdp:eq:C.SK.2}.
Consequently \eqref{cdp:eq:C.SK.3} is a nonzero isogeny of $\C^*$
groups. Composing it with \eqref{cdp:eq:C.SK.1} gives
\begin{equation}
 q_c\bigl([p_c^*(M_z|_{S_c})]\bigr)
       =\exp(2\pi\mathrm i k z)
       \qquad(k\in\Z\setminus\{0\}).
\hypertarget{tag-011K}{}\label{cdp:eq:C.SK.4}
\end{equation}
Here integrality follows from the period $z\mapsto z+1$;
every holomorphic group homomorphism $\C^*\to\C^*$ is
an integral power. The integer $k$ is locally constant
in the holomorphic families just described, and is
therefore a single integer on the connected smooth base.

\item Let $\rho_c:S_c\to\PP^1$ be its elliptic algebraic
reduction. The induced finite map $r_c:E_c\to\PP^1$
satisfies $\rho_c p_c=r_c\pi_c$. Write $h=\deg r_c$;
this positive integer is constant. Choose one positive
integer $N$ killing the torsion group $H^2(S_c;\Z)$
and satisfying $d\mid hN$. This choice is uniform
because the smooth fibres are differentiably locally
trivial. Then
$A_c=\rho_c^*\OO_{\PP^1}(N)$ belongs to $\Pic^0(S_c)$,
and \eqref{cdp:eq:C.SK.2} gives
\begin{equation}
 q_c([p_c^*A_c])=\alpha^{hN/d}.
\hypertarget{tag-011L}{}\label{cdp:eq:C.SK.5}
\end{equation}
This identity follows by taking degrees on the Albanese elliptic
curve.

\item Choose $z_0\in\C$ with
$\exp(2\pi\mathrm i k z_0)=\alpha^{hN/d}$ and put
$M=M_{z_0}$. By \eqref{cdp:eq:C.SK.4}--\eqref{cdp:eq:C.SK.5}, the difference of
$p_c^*(M|_{S_c})$ and $p_c^*A_c$ is
$\pi_c^*R_c$ for some $R_c\in\Pic^0(E_c)$.
Both of the first two lines come from $S_c$ and are
$G$-invariant. Injectivity of $\pi_c^*$ implies that
$R_c$ is fixed by the induced nontrivial linear action
on $\Pic^0(E_c)$. An elliptic-curve linear automorphism
of finite order other than one has order
$r\in\{2,3,4,6\}$. If it fixes a point $R_c$, then
\[
 rR_c=(1+\zeta+\cdots+\zeta^{r-1})R_c=0
\]
in the elliptic group. Thus $R_c^{\otimes12}$ is
trivial. The norm of the finite cover satisfies
$\operatorname{Nm}_{p_c}(p_c^*L)=L^{\otimes m}$.
Applying it to the difference gives the line-bundle identity
\begin{equation}
 M^{\otimes12m}|_{S_c}
       \simeq \rho_c^*\OO_{\PP^1}(12mN)
\hypertarget{tag-011M}{}\label{cdp:eq:C.SK.6}
\end{equation}
for every smooth general fibre.

\item The right side has at least two linearly independent
sections, by pullback from $\PP^1$. Coherent generic
base change gives
$\operatorname{rank}f_*M^{\otimes12m}\ge2$.
Its relative linear system has a positive-dimensional
image on a general fibre and therefore a surface
image in the projective bundle over $\PP^1$.
Pullback of meromorphic functions from this image
gives $a(Y)\ge2$, the desired contradiction.

\end{enumerate}
\end{proof}

\begin{proposition}[Exclusion of primary Kodaira general fibres]\hypertarget{tag-01F5}{}\label{cdp:st:C.3.10}
\StatementRecord{cdp:st:C.3.10}{proved}{01F5}{\SeriesRef{cdp:st:C.3.9}, \SeriesRef{cdp:st:D.1.12}, \SeriesRef{cdp:st:D.1.9}, \SeriesRef{cdp:st:D.1.7}, \SeriesRef{cdp:eq:C.PK.4}, \SeriesRef{cdp:eq:C.PK.5}, \SeriesRef{cdp:eq:C.PK.1}, \SeriesRef{cdp:eq:C.PK.2}, \SeriesRef{cdp:eq:C.PK.7}; \cite[pp.~373--374]{Borcea1984}; \cite[\S1, pp.~2--3]{SS2002}; \cite[Remark~1, p.~377]{Borcea1984}; \cite[\S3, equations (8)--(12), pp.~376--377]{Borcea1984}; \cite[Theorem~2.18 and \S\S2.19--2.21]{Steenbrink1976}}


A smooth compact complex threefold with $b_1(Y)=b_2(Y)=0$ and
$a(Y)=1$ cannot have a holomorphic algebraic reduction
$f:Y\to\PP^1$ whose general fibre is a primary Kodaira
surface.
\end{proposition}


\begin{proof}

\begin{enumerate}[leftmargin=*]
\item Let $B=\PP^1\setminus\Sigma$ be the smooth base and write
$P_c=f^{-1}(c)$ and $\pi_c:P_c\to E_c$ for the canonical
elliptic algebraic reduction of a smooth fibre. The degree
$d>0$ is constant. All the smooth fibres have the fixed
primary Kodaira differentiable type; every complex structure
on that type is again a primary Kodaira structure
\cite[pp.~373--374]{Borcea1984}.
Use the intrinsic quotient and its orientation from \eqref{cdp:eq:C.SK.2}:
\[
 Q_c=\Pic^0(P_c)/\pi_c^*\Pic^0(E_c),\qquad
 q_c:Q_c\xrightarrow{\simeq}\C^*.
\]
The class of the pullback of a degree-$d$ line has coordinate
$\alpha(c)\in\Delta^*$. The holomorphic moduli comparison
in the proof of \ref{cdp:st:C.3.9} applies directly here, so this bounded
function extends to $\PP^1$ and has a single nonzero
constant value $\alpha$.

\item Put $\mathscr E=R^1f_*\OO_Y$ and
$\mathscr G=R^2f_*\OO_Y$. The primary Kodaira case of
\ref{cdp:st:D.1.12} proves, on the original threefold and every full fibre,
\begin{equation}
 \operatorname{rank}\mathscr E=2,\qquad
 \operatorname{rank}\mathscr G=1,\qquad
 \bigwedge^2\mathscr E\xrightarrow{\simeq}\mathscr G.
\hypertarget{tag-00V8}{}\label{cdp:eq:C.PK.1}
\end{equation}
Both sheaves are locally free and satisfy base change.
Let $\gamma$ be the integral transgression section of \ref{cdp:st:D.1.9},
and let $\beta\in H^0(\PP^1,\mathscr E)=H^1(Y,\OO_Y)$
be its coefficient image. This section is nowhere zero.
For $u\in H^1(P_c,\OO_{P_c})$ define the complex-linear
functional $\lambda_c$ by the exact identity
\begin{equation}
 q_c\bigl([\exp(2\pi\mathrm i z u)]\bigr)
     =\exp(2\pi\mathrm i z\lambda_c(u))
     \quad(z\in\C).
\hypertarget{tag-00V9}{}\label{cdp:eq:C.PK.2}
\end{equation}
Its kernel is the line
$A_c=\pi_c^*H^1(E_c,\OO_{E_c})$.
The function $k(c)=\lambda_c(\beta_c)$ takes integral
values, because $\beta_c$ is the coefficient image of
an integral class. It is holomorphic, by the relative
Picard quotient in \eqref{cdp:eq:C.SK.2}, hence is a constant integer $k$.

\item When $k\ne0$, choose $n>0$ with $nd\ge2$ and choose
$z\in\C$ with $\exp(2\pi\mathrm i kz)=\alpha^n$.
The ambient line $M=\exp(2\pi\mathrm i z\beta)$ then
has, on every smooth fibre, the same $Q_c$ class as the
pullback of a degree-$nd$ line on $E_c$. Its difference
from that pullback belongs to $\pi_c^*\Pic^0(E_c)$.
Thus
$M|_{P_c}=\pi_c^*L_c$ for a line $L_c$ of degree $nd$,
and $h^0(P_c,M|_{P_c})=nd\ge2$. The relative linear
system gives $a(Y)\ge2$, as at the end of \ref{cdp:st:C.3.9}.
\item For $k=0$, $\beta_c\in A_c$. The coefficient map
$H^1(E_c;\R)\to H^1(E_c,\OO_{E_c})$ is an isomorphism
of real vector spaces. Its pullback, together with the real
injectivity of \ref{cdp:st:D.1.7} on $P_c$, proves the exact implication
\begin{equation}
 \beta_c\in A_c
 \quad\Longrightarrow\quad
 \gamma_c\in\pi_c^*H^1(E_c;\R).
\hypertarget{tag-00VA}{}\label{cdp:eq:C.PK.3}
\end{equation}
The lattice $\pi_c^*H^1(E_c;\Z)$ is saturated in
$H^1(P_c;\Z)$: the elliptic bundle induces a surjection
on fundamental groups, so the cokernel of the induced
injection on integral $H^1$ is torsion free. Thus $\gamma$
lies in this rank-two integral local system. Divide it by
its fixed positive divisibility to obtain a primitive fixed
section $\gamma_0$. On a simply connected open set choose
a complementary integral section $\delta$ with oriented
intersection number one. The ratio of their nonzero
coefficient images is a holomorphic elliptic period
\[
 \tau=\frac{\iota(\delta)}{\iota(\gamma_0)},
 \qquad \operatorname{Im}\tau>0,
\]
after choosing the sign convention of the ordered pair.
Monodromy fixes $\gamma_0$ and replaces $\delta$ by
$\delta+N\gamma_0$, $N\in\Z$. Therefore
$\exp(2\pi\mathrm i\tau)$ is a single-valued bounded
holomorphic function on $B$. It extends to $\PP^1$ and
is constant. Its logarithmic derivative gives $d\tau=0$,
and a nonzero integral translation cannot fix this constant
period. Consequently the Albanese integral $H^1$ local
system has identity monodromy, and $j(E_c)$ is constant.

\item Both moduli $(j,\alpha)$ and the degree $d$ are now
constant. Borcea's classification, in the precise form
\cite[\S1, pp.~2--3]{SS2002}, makes all smooth fibres biholomorphic
to one primary Kodaira surface $P$. Borcea's complete
holomorphic marked period family \cite[Remark~1, p.~377]{Borcea1984}
expresses our family, near any smooth base point, as the
pullback by a holomorphic map into his period domain
$\mathcal D\subset\PP(H^2(P;\C))$. By his Theorem~3 \cite[Theorem~3 and the ensuing quotient description, pp.~378--379]{Borcea1984},
constant isomorphism class means that the image of this
map lies in one orbit of the integral cohomology action.
That orbit is countable, since the action factors through
$\operatorname{GL}(H^2(P;\Z)/\operatorname{torsion})$.
A connected subset of a countable subset of a complex
manifold is one point: a real coordinate separating two
points would have a nontrivial interval as its connected
image. Therefore the period parameter is constant and
its pullback is a product. 
\item Write the fundamental group as
\begin{equation}
 \Gamma=\langle x,y,w,u\mid
       w,u\text{ central},\ [x,y]=w^d\rangle.
\hypertarget{tag-00VB}{}\label{cdp:eq:C.PK.4}
\end{equation}
Its centre $\langle w,u\rangle$ is the fundamental group
of the elliptic fibre, and its quotient is the fundamental
group of the elliptic base. A biholomorphism inducing the
identity on the base $H^1$ induces the identity on this
quotient. The relation in \eqref{cdp:eq:C.PK.4} then forces it to fix
$w$. On an elliptic fibre its linear part is a holomorphic
group automorphism. Such an automorphism acts by a complex
scalar on the universal cover; fixing the nonzero period
$w$ forces that scalar to be one. It therefore fixes both
$w$ and $u$. The remaining possible actions on $x,y$ are
exactly multiplication by central elements.

\item In the real
nilmanifold presentation choose invariant one-forms
$a,b,c,e$ with
\[
 da=db=de=0,\qquad dc=d\,a\wedge b.
\]
The four classes $[a\wedge c],[b\wedge c],
[a\wedge e],[b\wedge e]$ form a basis of
$H^2(P;\R)$, as in \cite[\S3, equations (8)--(12), pp.~376--377]{Borcea1984}.
An automorphism of \eqref{cdp:eq:C.PK.4} with the stated action extends
to the real nilpotent group: in logarithmic coordinates it
fixes the two central generators and adds central vectors
to the two base generators. Its pullback fixes $a,b$ and
replaces each of $c,e$ by itself plus a real linear
combination of $a,b$. Each of the four displayed classes
therefore changes only by a multiple of the exact form
$a\wedge b$. Since the universal cover of $P$ is
contractible, a map of $P$ and this group-model map with
the same induced fundamental-group map are homotopic.
The original biholomorphism therefore acts as the identity on $H^2$.
In particular every monodromy of our isotrivial family,
whose action on the base $H^1$ was just proved to be the
identity, acts identically on $H^2(P;\C)$.

\item Resolve the union of the exceptional-fibre supports to
normal crossings by successive blowups with smooth centres contained
in those fibres, writing
$\rho:\widetilde Y\to Y$ and $h=f\rho$. For each such
blowup, $\rho_*\OO_{\widetilde Y}=\OO_Y$ and
$R^j\rho_*\OO_{\widetilde Y}=0$ for $j>0$.
One local verification uses the formal neighbourhoods of
the exceptional projective bundle: their successive
conormal layers restrict to nonnegative twists
$\OO_{\PP^{r-1}}(l)$, $l\ge0$, which have no higher
cohomology. The analytic theorem on formal functions
gives the asserted vanishing; normality gives the
degree-zero equality. The Leray edge is consequently the
canonical isomorphism
\begin{equation}
 R^2h_*\OO_{\widetilde Y}\xrightarrow{\simeq}
 R^2f_*\OO_Y=\mathscr G.
\hypertarget{tag-00VC}{}\label{cdp:eq:C.PK.5}
\end{equation}
The same identities hold for a composite by derived
pushforward.

\item Let $D=(h^{-1}(\Sigma))_{\mathrm{red}}$ and use the
relative logarithmic de Rham complex
\[
 \mathscr K^\bullet=
 \Omega_{\widetilde Y}^\bullet(\log D)\big/
 \bigl(h^*\Omega_{\PP^1}^1(\log\Sigma)
       \wedge\Omega_{\widetilde Y}^{\bullet-1}(\log D)\bigr).
\]
Projection onto degree zero is a morphism of complexes
$\mathscr K^\bullet\to\OO_{\widetilde Y}[0]$:
it is the identity in degree zero and zero in every
positive degree. The analytic comparison of
\cite[\S2.1, p.~234; Theorem~2.18, pp.~237--238;
\S2.19 and Proposition~2.20, pp.~238--239; Theorem~2.21, p.~239]{Steenbrink1976}
applies on each base disc: the total space is smooth, the map is
proper and smooth off the central value, and the full central fibre
has normal-crossings support with its positive multiplicities.
It gives the logarithmic de Rham bundle
$\mathscr H^2=\mathbb R^2h_*\mathscr K^\bullet$.
Identity monodromy makes every residue zero, by the
residue interval $[0,1)$ and the exponential comparison.
The global flat frame on $B$ therefore extends across
every point of $\Sigma$, proving
$\mathscr H^2\simeq H^2(P;\C)\otimes\OO_{\PP^1}$.
Taking degree-two direct hypercohomology of the displayed
projection and then \eqref{cdp:eq:C.PK.5} gives a holomorphic map
\begin{equation}
 p:H^2(P;\C)\otimes\OO_{\PP^1}
          \longrightarrow\mathscr G.
\hypertarget{tag-00VD}{}\label{cdp:eq:C.PK.6}
\end{equation}
Over $B$ it is precisely the coefficient projection
$H^2(P_c;\C)\to H^2(P_c,\OO_{P_c})$.

\item In an isotrivial chart this latter map is constant. It is
surjective: the class of
$\overline\alpha\wedge\overline\eta$ used in the proof
of \ref{cdp:st:D.1.12} is represented by a closed invariant two-form
and has nonzero Dolbeault class. The charts also show
that its kernel is a flat subspace. Choose a constant flat
class $v_0\in H^2(P;\C)$ whose projection is nonzero,
and put $s=p(v_0)\in H^0(\PP^1,\mathscr G)$.
This section is nowhere zero on $B$.

\item By \eqref{cdp:eq:C.PK.1} and the nowhere-zero section $\beta$, the
following is an exact sequence of vector bundles on the
entire projective line:
\begin{equation}
 0\longrightarrow\OO_{\PP^1}\beta
  \longrightarrow\mathscr E
  \xrightarrow{\ \beta\wedge-\ }\mathscr G
  \longrightarrow0.
\hypertarget{tag-00VE}{}\label{cdp:eq:C.PK.7}
\end{equation}
Since $H^1(\PP^1,\OO)=0$, choose
$v\in H^0(\PP^1,\mathscr E)=H^1(Y,\OO_Y)$ with
$\beta\wedge v=s$.
In each isotrivial chart,
$\beta$ is the coefficient image of a fixed integral
class, $v_0$ is a fixed de Rham class, and $s=p(v_0)$
is a fixed nonzero class in $H^2(P,\OO_P)$.
The intrinsically oriented quotient $q_c$ in \eqref{cdp:eq:C.PK.2} is also fixed
by the chart: every automorphism preserves its
positive-degree class $\alpha\in\Delta^*$, so it cannot
induce inversion of $\C^*$.
Both linear maps $u\mapsto\beta\wedge u$ and
$u\mapsto\lambda(u)s$ have kernel $A$ and are nonzero.
Thus they differ by a constant $\kappa\in\C^*$ on
that chart. The constants agree on overlaps, because
all the maps and sections used in their definition are
globally defined. The connectedness of $B$ gives
one constant $\kappa$, and \eqref{cdp:eq:C.PK.7} yields
\begin{equation}
 \lambda_c(v_c)=\kappa^{-1}\ne0
 \quad\hbox{for every }c\in B.
\hypertarget{tag-00VF}{}\label{cdp:eq:C.PK.8}
\end{equation}

\item Choose $n>0$ with $nd\ge2$ and $z\in\C$ with
$\exp(2\pi\mathrm i z\kappa^{-1})=\alpha^n$.
The single ambient line $M=\exp(2\pi\mathrm i z v)$
has this quotient coordinate on every smooth fibre.
Exactly as in the $k\ne0$ case,
$M|_{P_c}=\pi_c^*L_c$ with $\deg L_c=nd\ge2$.
Its direct image has rank at least two and its relative
linear system gives a surface image over $\PP^1$,
contradicting $a(Y)=1$.

\end{enumerate}
\end{proof}

\subsection{Inoue fibres}\label{sec:II:inoue-exclusion}

\begin{proposition}[Exclusion of Inoue general fibres]\hypertarget{tag-01F7}{}\label{cdp:st:C.3.12}
\StatementRecord{cdp:st:C.3.12}{proved}{01F7}{\SeriesRef{cdp:st:C.3.1}, \SeriesRef{cdp:eq:C.IN.2}, \SeriesRef{cdp:eq:C.IN.3}, \SeriesRef{cdp:eq:C.IN.6}; \cite[pp.~274, 278, 280]{Inoue1974}; \cite[arXiv version~4, \S\S2.1--2.3, pp.~6--11]{KhaledTeleman2026}; \cite[Theorem~3.2, p.~16]{BMS2012}; \cite[\S3.3, equation~(1), pp.~17--18]{BMS2012}}


Under the hypotheses of \ref{cdp:st:C.3.1}, the very general smooth fibre is
not an Inoue surface. More precisely, a proper holomorphic map from
a compact connected complex manifold to a compact Riemann surface,
with arbitrary fibre multiplicities and singularities over a finite
set and Inoue surfaces over its complement, has a class in
$H^1(Y;\mathbb R)$ restricting nontrivially to each Inoue fibre.
Thus these families cannot have $b_1(Y)=0$.
\end{proposition}


\begin{proof}

\begin{enumerate}[leftmargin=*]
\item Write $B^\circ$ for the complement of the
finite set and $Y^\circ=f^{-1}(B^\circ)$. Put
$\mathbb H=\{w\in\mathbb C:\operatorname{Im}w>0\}$.
Each Inoue surface is a quotient of $\mathbb H\times\mathbb C$.
In the three standard presentations the first coordinates of all
deck transformations have the form
\begin{equation}
 w\longmapsto a w+b,\qquad a>0,\quad b\in\mathbb R.
\hypertarget{tag-00NK}{}\label{cdp:eq:C.IN.1}
\end{equation}
They include a nonzero real translation and a transformation
$w\mapsto\alpha w$, with $\alpha>1$. Inoue's original
Propositions 2 and 5 establish infinitesimal rigidity for types
$S_M$ and $S^{-}$; his Proposition 4 gives a locally complete
one-parameter family for type $S^{+}$
\cite[pp.~274, 278, 280]{Inoue1974}.
In that family the parameter changes only the second coordinate:
local coordinate transitions have the form
\begin{equation}
\begin{aligned}
 w_i&=\alpha^{m_{ij}}w_j+a_{ij},\\
 z_i&=z_j+b_{ij}w_j+c_{ij}+t\,m_{ij},
\end{aligned}
\hypertarget{tag-00NL}{}\label{cdp:eq:C.IN.2}
\end{equation}
where $m_{ij}\in\mathbb Z$ and the other coefficients are
independent of $t$. In particular the real affine transformations
in the first coordinate are constant throughout this complete
family. The standard deck presentations are also recorded in
\cite[arXiv version~4, \S\S2.1--2.3, pp.~6--11]{KhaledTeleman2026}.

\item After shrinking a simply connected base disc $V\subset B^\circ$,
local completeness identifies $f^{-1}(V)$ with the pullback of one
of these families. Its universal cover over $V$ is therefore
$\mathbb H\times\mathbb C\times V$, with holomorphic covering
map to $f^{-1}(V)$. In the $S^{+}$ case the parameter $t$ in
\eqref{cdp:eq:C.IN.2} is a holomorphic function on $V$; the other two cases
are products. The smooth real one-form
\begin{equation}
 \theta_V=d\log\operatorname{Im}w
\hypertarget{tag-00NM}{}\label{cdp:eq:C.IN.3}
\end{equation}
is invariant under every deck transformation, where $d$ is the
exterior derivative on the \emph{whole} covering space. It thus
descends to a closed real one-form on $f^{-1}(V)$.

\item Over a simply connected smaller overlap, a transition
between the two family descriptions lifts holomorphically to the
universal covers. On each fibre its first coordinate is independent
of $z$: composing with a Cayley transform gives a bounded entire
function of $z$, which is constant. Applying the same argument to
the inverse proves that the resulting first coordinate is an
automorphism $h$ of $\mathbb H$. This automorphism conjugates the
first-coordinate deck groups. Their common fixed-point set on
$\partial\mathbb H=\mathbb R\cup\{\infty\}$ is exactly
$\{\infty\}$: the nonzero translation has this as its sole fixed
point, and all their elements fix it. Consequently $h(\infty)=
\infty$ and $h(w)=a w+b$ with $a>0$ and $b\in\mathbb R$.
In a holomorphic family of these transitions, $a$ and $b$ are
holomorphic functions of the base parameter; this follows, for
example, by taking the $w$ derivative and then evaluating at a
fixed $w$. Their real values force them to be locally constant.
Therefore
\begin{equation}
 d\log\operatorname{Im}(a w+b)
 =d\log(a\operatorname{Im}w)
 =d\log\operatorname{Im}w
\hypertarget{tag-00NN}{}\label{cdp:eq:C.IN.4}
\end{equation}
as forms on the total family. Thus the $\theta_V$ glue to a
closed form $\theta$ on $Y^\circ$.
The same argument includes a finite \'etale quotient of an Inoue
surface. Locally over a simply connected base disc its finite
cover extends to the family: lift a differentiable marking by
Ehresmann's theorem and equip the resulting covering space with
the pulled-back complex structure. Passing to the Galois closure
when necessary gives a family of finite covers of Inoue surfaces;
these are again Inoue surfaces, as is also seen from their
finite-index lattice presentations. Every deck transformation of
the quotient normalizes a finite-index Inoue deck group on
$\mathbb H\times\mathbb C$. Its first-coordinate group still
contains a nonzero translation and has common boundary fixed
point precisely $\infty$. The preceding conjugation argument
therefore makes each quotient deck transformation real affine in
$w$, and its coefficients are constant over the base disc.
Consequently \eqref{cdp:eq:C.IN.3} descends through the finite quotient as
well. Its pullback to the Inoue cover has the nonzero period
$\log\alpha$, so the descended fibre class is nonzero.
For a fibre loop represented on its universal cover by the deck
transformation $w\mapsto\alpha w$, its integral is
$\log\alpha\ne0$. Hence $[\theta]$ restricts nontrivially to
the fibre.

\item Put
$D=Y\setminus Y^\circ$. Let $\iota:\Delta\to Y$ be a small
holomorphic transverse disc at a smooth point of a reduced
irreducible component of $D$, avoiding all other components.
The pullback of the preceding transverse coordinate atlas on
$\Delta^*$ has constant real affine transitions. On the universal
cover
\[
 \mathbb H_c=\{\operatorname{Im}\tau>c\}
 \longrightarrow\Delta^*,\qquad
 \tau\longmapsto\exp(2\pi\mathrm i\tau),
\]
analytic continuation consequently gives a holomorphic developing
map $W:\mathbb H_c\to\mathbb H$ with
$W(\tau+1)=aW(\tau)+b$ for some $a>0$, $b\in\mathbb R$.
This follows by trivializing the pulled-back constant transition
cocycle on the simply connected covering space.
The period of $\iota^*\theta$ on the positive puncture loop is
$\log a$.

\item Use the Poincar\'e metrics
$|dw|/\operatorname{Im}w$ and
$|d\tau|/(\operatorname{Im}\tau-c)$.
The function $\log\operatorname{Im}w$ has differential of norm
one for the first metric. Schwarz--Pick and the horizontal segment
from $\tau$ to $\tau+1$ therefore give
\begin{equation}
\begin{aligned}
 |\log a|
 &\le d_{\mathbb H}\bigl(W(\tau),W(\tau+1)\bigr)\\
 &\le d_{\mathbb H_c}(\tau,\tau+1)
 \le\frac{1}{\operatorname{Im}\tau-c}.
\end{aligned}
\hypertarget{tag-00NO}{}\label{cdp:eq:C.IN.5}
\end{equation}
Letting $\operatorname{Im}\tau\to+\infty$ proves $a=1$.
Every meridian period of $\theta$ is thus zero. This calculation
applies to the reduced divisor and is unaffected by the
scheme-theoretic multiplicities of its components.

\item The map $H_1(Y^\circ;\mathbb R)\to H_1(Y;\mathbb R)$ is
surjective: loops can be moved away from a complex divisor.
Its kernel is generated by the meridians of the irreducible
components of $D$. Indeed, a bounding real two-chain can be put
in general position, transverse to the smooth divisor locus and
disjoint from its singular locus, which has real codimension at
least four. Removing small discs about its intersection points
leaves a chain in $Y^\circ$ whose additional boundary circles
are meridians. The period homomorphism of $\theta$ therefore factors uniquely
through $H_1(Y;\mathbb R)$ and defines a class
\begin{equation}
 \widetilde\theta\in H^1(Y;\mathbb R),\qquad
 \widetilde\theta|_{Y^\circ}=[\theta].
\hypertarget{tag-00NP}{}\label{cdp:eq:C.IN.6}
\end{equation}
Its restriction to an Inoue fibre is nonzero by \eqref{cdp:eq:C.IN.3}.
Consequently $b_1(Y)>0$.

\item Suppose a very general smooth fibre under the hypotheses of
\ref{cdp:st:C.3.1} is an Inoue surface. Ehresmann's
theorem identifies the diffeomorphism type, and in particular
the fundamental group, of all smooth fibres. The first two
Betti numbers are then $(1,0)$. These surfaces are minimal by
\ref{cdp:st:C.3.1} and non-K\"ahler. The compact-surface classification
gives class VII surfaces or non-K\"ahler elliptic surfaces
\cite[Theorem~3.2, p.~16]{BMS2012}.
The Inoue group is solvable and is not virtually nilpotent.
For the latter assertion, its normal lattice is $\mathbb Z^3$
in type $S_M$, and a three-dimensional integral Heisenberg group
in types $S^{+}$ and $S^{-}$. Conjugation by a generator of the
cyclic quotient has an eigenvalue $\alpha>1$ on the rational
abelianization of this lattice. Every finite-index subgroup
contains a positive power of that generator and a finite-index
sublattice, so the eigenvalue becomes $\alpha^k>1$. In a
nilpotent group repeated commutators make every such conjugation
unipotent on this rational abelianization. 
\item A Hopf group is virtually cyclic, and a Kodaira group is virtually
nilpotent; neither can equal this group. For a properly elliptic
surface of Euler number zero the canonical bundle formula makes
the base orbifold hyperbolic: its orbifold canonical degree is
positive exactly in Kodaira dimension one. The elliptic
fibration induces a surjection of its fundamental group onto the
orbifold group. That group has a finite-index closed surface
group of genus at least two, and hence is not solvable. This
also cannot be a quotient of an Inoue group. The orbifold exact
sequence is the one used in \cite[\S3.3, equation~(1), pp.~17--18]{BMS2012}.
Thus every smooth fibre is an Inoue surface. Applying \eqref{cdp:eq:C.IN.6}
contradicts $b_1(Y)=0$.

\end{enumerate}
\end{proof}

\subsection{Compact torus and Hopf fibrations}\label{sec:II:fibre-classification}

\begin{definition}[The square elliptic curve and its lattice]
\hypertarget{tag-P241}{}\label{def:II:identity-example-lattice}
\StatementRecord{def:II:identity-example-lattice}{proved}{P241}{}
Put
\[
 E=\C/(\Z+i\Z),\qquad q=e^{-2\pi},\qquad
 A=E^2,\qquad \Lambda=H_1(A;\Z)=\bigoplus_{j=1}^4\Z e_j,
\]
where $(e_1,e_2,e_3,e_4)=((1,0),(i,0),(0,1),(0,i))$.
We identify the additive and multiplicative presentations of $E$ by
the holomorphic group isomorphism
\begin{equation}
 \C/(\Z+i\Z)\longrightarrow\C^*/q^\Z,
 \qquad z\longmapsto e^{2\pi i z}.
\hypertarget{tag-0199}{}\label{cdp:eq:D.22}
\end{equation}
The lattice generators $1$ and $i$ multiply the exponential by $1$
and $q$, respectively. Homology in these torus constructions has
coefficients in $\Z$. A base meridian is counterclockwise in its
holomorphic coordinate; at infinity that coordinate is $1/t$.
\end{definition}


\begin{proposition}[An elliptic bundle degeneration]\hypertarget{tag-02WG}{}\label{night:P4.1}\label{stmt:P4.1}
\StatementRecord{night:P4.1}{proved}{02WG}{\SeriesRef{def:II:identity-example-lattice}}
Take labelled copies $C=B=E$ and identify $C\times B$ with $A$
by $(z,w)\mapsto(z,w)$. Put
$V=\operatorname{Bl}_{(0,0)}(B\times\Delta)$, and let $b:V\to\Delta$
be the composition of the blowdown with the projection. Let $R$ be
the exceptional curve and put $L=\OO_V(R)$, with $L^\times$ denoting
the complement of its zero section. The holomorphic quotient
\begin{equation}
 U=L^\times/q^\Z\xrightarrow{\ p\ }V\xrightarrow{\ b\ }\Delta
\hypertarget{tag-019A}{}\label{cdp:eq:D.23}
\end{equation}
is a smooth Hausdorff threefold, proper over $\Delta$.
Its noncentral fibres are $C\times B$. Its central fibre is the reduced
union of a primary Kodaira surface and the scalar Hopf surface
$(\C^2\setminus\{0\})/q^\Z$, meeting transversally in one elliptic curve.
\end{proposition}
\begin{proof}

\begin{enumerate}[label=\arabic*.,leftmargin=*]
\item We use coordinates $(z,t)$ at the blown-up point. The blowup has charts
\[
 (z,t)=(u,uv),\qquad (z,t)=(rs,s),\qquad r=v^{-1},\quad s=uv.
\]
The frames $u^{-1}$ and $s^{-1}$ for $L$ give fibre coordinates related by
$\lambda_B=v\lambda_A$. Each nonzero fibre coordinate, modulo multiplication
by $q$, gives a product chart with fibre $E$. The transitions are holomorphic
translations on $E$. The scalar action is free and properly discontinuous:
the logarithm of the norm of a nonzero fibre coordinate changes by the
nonzero constant $\log q$, and compact subsets meet only finitely many
translates. Thus the quotient is Hausdorff and smooth. These product charts
have compact fibre $E$, so $U\to V$ is proper. The blowup and $B\times\Delta
\to\Delta$ are proper, proving properness of their composite.

\item The canonical section of $\OO(R)$ has coefficients $u$ and $s$. Over
$t\ne0$ it never vanishes. The ratios $\lambda_A/u=\lambda_B/s$ therefore
give the holomorphic product trivialization
\begin{equation}
 U|_{\Delta^*}\simeq C\times B\times\Delta^*.
\hypertarget{tag-019B}{}\label{cdp:eq:D.24}
\end{equation}

\item The central divisor in $V$ is $B'\cup R$, with $R^2=-1$ and $R\cdot B'=1$.
Thus $L|_R=\OO_{\Pp^1}(-1)$ and $L|_{B'}$ has degree one. Over $B'$ the
quotient is a principal elliptic bundle with primitive nonzero circle Chern
class. Its canonical bundle is trivial: the base differential wedged with
$d\lambda/\lambda$ is invariant under the line-bundle transition, because
the additional term is pulled back from the one-dimensional base. The
circle bundle over $B'$ has primitive Euler class one: in its
fundamental-group presentation the commutator of the two base generators
is the fibre generator. Abelianization therefore gives $H_1=\Z^2$.
The logarithm of the fibre norm modulo $\log q$ gives the further
circle factor, so the surface has $b_1=3$. The classification
\cite[Chapter~VI, Theorem~1.1 and Table~10, pp.~243--244]{BHPV2004E}
identifies this surface, whose canonical bundle is trivial, as a primary
Kodaira surface.

\item The tautological evaluation morphism
$\OO_{\Pp^1}(-1)^\times\to\C^2\setminus\{0\}$ is a biholomorphism: a
nonzero vector uniquely determines its line. It intertwines both scalar
actions, and hence identifies the other component with the stated Hopf
surface. Over $B'\cap R$ the components meet in $E$. In local product
coordinates the map is $(u,v,c)\mapsto uv$; this proves transversality,
multiplicity one, and the absence of other critical points.

\end{enumerate}
\end{proof}

\begin{proposition}[Integral homology of the local degeneration]\hypertarget{tag-02X2}{}\label{night:P4.3}\label{stmt:P4.3}
\StatementRecord{night:P4.3}{proved}{02X2}{\SeriesRef{night:P4.1}}
For the degeneration $U$ of Proposition~\ref{night:P4.1}, order the
regular fibre lattice as $(a,c,b,d)$, where $(a,c)$ are the cycles
$(1,i)$ in $C$ and $(b,d)$ are the cycles $(1,i)$ in $B$.
Use the product trivialization \eqref{cdp:eq:D.24}.
The space $U$ deformation retracts onto its central fibre.
The regular-fibre inclusions induce the quotient maps
\begin{equation}
 H_1(C\times B)=\Lambda\longrightarrow\Lambda/\Z a=H_1(U),\qquad
 H_2(C\times B)=\bigwedge^2\Lambda\longrightarrow
 \bigwedge^2\Lambda/\Z(a\wedge c)=H_2(U).
\hypertarget{tag-019C}{}\label{cdp:eq:D.25}
\end{equation}
Both target groups are free. For the positive base meridian $\mu$,
the boundary inclusions satisfy
\begin{equation}
 \mu\longmapsto0\quad\text{in }H_1(U),\qquad
 x\wedge\mu\longmapsto0\quad\text{in }H_2(U)
 \quad(x\in\Lambda).
\hypertarget{tag-019D}{}\label{cdp:eq:D.26}
\end{equation}
\end{proposition}
\begin{proof}

\begin{enumerate}[label=\arabic*.,leftmargin=*]
\item A Hermitian metric gives a diffeomorphism
$L^\times/q^\Z\cong S(L)\times S^1_c$, with factors the circle bundle and
the logarithm of the norm modulo $\log q$. Orient $S^1_c$ by the
lattice period $i$, whose norm-logarithm increment is $\log q$.
We first specify a homotopy model for $V$. Blowing up a four-ball replaces it by the disc bundle of
$\OO(-1)$ along its Hopf boundary. The complement of the ball in
$B\times\Delta$ has homotopy type $B\vee S^3$; the boundary sphere is
the $S^3$ summand, since its map to $B$ is nullhomotopic. The resulting
homotopy pushout is $B\vee S^2$. Equivalently, on simply connected
covers the replacement adds one sphere at each lift of the centre and
no other homology; the comparison is an isomorphism on $\pi_1$ and on
the homology of those covers, and hence a homotopy equivalence. We may
choose the $B$ summand to be a regular $B_t$ and the sphere to be $R$,
joining their base points by a path.

Let $B'$ be the strict transform of $B\times\{0\}$.
The inclusion $B'\cup R\hookrightarrow V$ induces an isomorphism
on fundamental groups: the point blowup preserves $\pi_1$, and
the blowdown maps $B'$ isomorphically onto $B$.
On the simply connected covers, the lift of $B'$ is contractible
and each exceptional sphere maps to the corresponding sphere
added by the blowup. The same homology comparison therefore
makes this inclusion a homotopy equivalence. A triangulation
compatible with the normal-crossings divisor $B'\cup R$ makes
the inclusion a cofibration, so it admits a strong deformation
retraction. The relative homotopy lifting property of the
principal elliptic bundle, with the identity lift on the central
fibre, gives the asserted deformation retraction of $U$.

\item The circle Chern class is zero on $B_t$ and $-1$ on $R$. Therefore
$S(L)$ is represented by the homotopy pushout
\[
 T^3_{(a,b,d)}\longleftarrow S^1_a\longrightarrow S^3.
\]
The second map is a Hopf fibre and is nullhomotopic. Thus this pushout
has the homotopy type $(T^3\cup_aD^2)\vee S^3$. We multiply by $S^1_c$.
The resulting integral cellular complex starts with the exterior-cell
complex of $T^4$, ordered by $(a,c,b,d)$, and adds cells $r$ in degree
two, $rc,h$ in degree three, and $hc$ in degree four. Its only nonzero
differentials are
\[
 \partial r=a,\qquad \partial(rc)=a\wedge c.
\]
We choose the orientation of $r$ to give the first equality and use the
product orientation for $rc$. The pairs $(r,a)$ and $(rc,a\wedge c)$
are integral contractible summands: the degree-one homotopy sends
$a$ to $r$ and $a\wedge c$ to $rc$, and is zero on the remaining
basis cells. Deleting these pairs leaves the zero differential and
free homology of ranks $(1,3,5,5,2)$. The inclusion of the torus
subcomplex kills exactly $a$ in degree one and $a\wedge c$ in
degree two, giving the asserted quotient maps.

\item We represent a base point, and both base generators $b,d$, away from
$0\in B$. Their products with the disc miss $R$, so the canonical section
trivializes $L$ over the resulting chains. The meridian and the products
with $b,d$ consequently bound. The products with $a,c$ bound over one
such base point. This proves all the remaining assertions on the maps.

\end{enumerate}
\end{proof}

\begin{proposition}[A simply connected compact torus fibration]
\hypertarget{tag-01GE}{}\label{cdp:st:D.2.7}
\hypertarget{tag-02WR}{}\label{night:P4.2}\label{stmt:P4.2}
\hypertarget{tag-02XD}{}\label{night:P4.4}\label{stmt:P4.4}
\hypertarget{tag-02XO}{}\label{night:P4.5}\label{stmt:P4.5}
\StatementRecord{cdp:st:D.2.7}{proved}{01GE}{\SeriesRef{night:P4.1}, \SeriesRef{night:P4.3}, \SeriesRef{cdp:st:E1}}
Let $E$, $A$ and $\Lambda$ be as in
Definition~\ref{def:II:identity-example-lattice}. Define a principal
$A$-bundle $Y_0\to\PP^1$ by the clutching translation
\begin{equation}
 a_{\rm out}=a_{\rm in}+(\ell(t),0),\qquad
 \ell(t)=\frac{\log t}{2\pi i}\pmod{\Z+i\Z}
\hypertarget{tag-019F}{}\label{cdp:eq:D.28}
\end{equation}
between a coordinate disc at zero and the exterior chart.
On disjoint discs about $0,1,\infty$, replace the product models
by the degeneration of Proposition~\ref{night:P4.1}, using the
canonical-section trivialization \eqref{cdp:eq:D.24} and the
ordered product markings, with $C_i=B_i=E$,
\begin{equation}
 \Psi_i:C_i\times B_i\longrightarrow A,\qquad
 \Psi_1(z,w)=(w,z),\quad
 \Psi_2(z,w)=(iz,iz+w),\quad
 \Psi_3(z,w)=(z,iz+w).
\hypertarget{tag-019E}{}\label{cdp:eq:D.27}
\end{equation}
These transitions define a smooth compact connected complex
threefold $f:Y\to\PP^1$ with connected fibres and
\begin{equation}
 H_1(Y;\Z)=H_2(Y;\Z)=0,\qquad
 a(Y)=1,\qquad c_3(Y)=0,\qquad b_3(Y)=2.
\hypertarget{tag-0198}{}\label{cdp:eq:D.21}
\end{equation}
Moreover $\pi_1(Y)=1$, and all integral homology groups are free,
of ranks $(1,0,0,2,0,0,1)$ in degrees zero through six.
Pullback gives the function-field isomorphism
$f^*:\mathcal M(\PP^1)\xrightarrow{\simeq}\mathcal M(Y)$.
Every smooth fibre is the fixed abelian surface $A$, and the
marked integral monodromy on $H_1(A;\Z)$ is the identity.
There are exactly three singular fibres. Each is reduced with
simple normal crossings and consists of a primary Kodaira surface
and a primary Hopf surface meeting along one smooth elliptic curve.
\end{proposition}
\begin{proof}
\begin{enumerate}[label=\arabic*.,leftmargin=*]
\item The map $\ell(t)$ in \eqref{cdp:eq:D.28} is single-valued and
holomorphic on an annulus because changing the logarithm adds an
integer. Thus it defines a holomorphic transition for a principal
$A$-bundle $Y_0$ on $\PP^1$. Its product charts make $Y_0$ smooth
and Hausdorff, and their compact fibre makes it proper over the
compact base. Hence $Y_0$ is compact.

Multiplication by $i$ preserves $\Z+i\Z$. The three inverse maps
in \eqref{cdp:eq:D.27} are
$(x,y)\mapsto(y,x)$, $(x,y)\mapsto(-ix,y-x)$ and
$(x,y)\mapsto(x,y-ix)$, respectively. Their complex matrices have
determinants $-1,i,1$, which are units in $\Z[i]$; their real
lattice determinants are one. Glue the three local models using
these maps, the canonical-section product identification, and
the local trivializations of $Y_0$. Different base points separate
in the base, and points over one base point lie in a Hausdorff
local model. The gluing is therefore Hausdorff. Properness holds
over each base chart and hence over the base; compactness of
the base implies that $Y$ is compact.

The two central components have connected elliptic fibres over
connected curves and meet each other. All remaining fibres equal
$A$, so all fibres and $Y$ are connected. The equations $uv=t$
from Proposition~\ref{night:P4.1} give precisely the three stated
singular fibres, with multiplicity one and simple normal crossings.
All changes of the marked smooth fibre are translations, so they
induce the identity on $H_1(A;\Z)$.

\item For the integral maps, write $(a_i,c_i,b_i,d_i)$ for the
images under $\Psi_i$ of the ordered local basis $(a,c,b,d)$ in
Proposition~\ref{night:P4.3}. The lattice matrices, with these
columns in the common basis $(e_1,e_2,e_3,e_4)$, are
\[
S_1=\begin{pmatrix}0&0&1&0\\0&0&0&1\\1&0&0&0\\0&1&0&0\end{pmatrix},\quad
S_2=\begin{pmatrix}0&-1&0&0\\1&0&0&0\\0&-1&1&0\\1&0&0&1\end{pmatrix},\quad
S_3=\begin{pmatrix}1&0&0&0\\0&1&0&0\\0&-1&1&0\\1&0&0&1\end{pmatrix}.
\]
Put $Q_i=S_i^{-1}$. Deleting the first row of $Q_i$ gives $R_{1,i}$.
The row indices of $\bigwedge^2Q_i$ are $(ac,ab,ad,cb,cd,bd)$ and its
column indices are $(12,13,14,23,24,34)$; deleting its first row gives
$R_{2,i}$. These are the integer matrices of the local quotient maps
\eqref{cdp:eq:D.25}, with target bases $(c_i,b_i,d_i)$ and
$(a_ib_i,a_id_i,c_ib_i,c_id_i,b_id_i)$.

Let $P$ be the complement of three small open base discs and
$W=f^{-1}(P)=A\times P$. Let $s_1,s_2$ generate $H_1(P)$ and
put $s_3=-s_1-s_2$. The overlaps are $N_i=A\times S^1$.
Positive continuation of $\ell$ around zero adds $e_1$. With
\[
 g_1=e_1,\quad g_2=g_3=0,\qquad
 v_1=\binom10,\quad v_2=\binom01,\quad v_3=\binom{-1}{-1},
\]
the overlap maps into $W$, in the basis $(s_1,s_2)$ for its base
factor, are
\begin{equation}
 x\longmapsto x\quad(x\in\Lambda),\qquad
 s_i^{\rm loc}\longmapsto g_i+v_i.
\hypertarget{tag-019G}{}\label{cdp:eq:D.29}
\end{equation}
The induced maps on products of these cycles are
\begin{equation}
 x\wedge y\longmapsto x\wedge y,\qquad
 x\wedge s_i^{\rm loc}\longmapsto x\wedge g_i+x\otimes v_i.
\hypertarget{tag-019H}{}\label{cdp:eq:D.30}
\end{equation}
They include the degree-two contribution of the clutching
translation, with fibre cycles preceding the base cycle.

\item Use the ordered basis $(e_1,e_2,e_3,e_4,s_1,s_2)$ for $H_1(W)$
and append the three bases $(c_i,b_i,d_i)$ for the local groups.
On the source concatenate the three bases
$(e_1,e_2,e_3,e_4,s_i^{\rm loc})$. The $15\times15$
Mayer--Vietoris matrix $F_1$ has, in its $i$th five-column block,
upper block
\[
 \begin{pmatrix}I_4&g_i\\0_{2\times4}&v_i\end{pmatrix}
\]
and local block $(-R_{1,i},0)$ in the $i$th summand; its other local
blocks vanish. The sign is the difference of the two inclusions.

For $H_2(W)$ use
\[
 (12,13,14,23,24,34,e_1s_1,\ldots,e_4s_1,e_1s_2,\ldots,e_4s_2)
\]
and append the three local five-element bases. The source is the
concatenation of
$(12,13,14,23,24,34,e_1s_i^{\rm loc},\ldots,e_4s_i^{\rm loc})$.
The $29\times30$ matrix $F_2$ has top six rows
$[I_6,(e_j\wedge g_i)_{j=1}^4]$ in its $i$th block; its next eight
rows have the coefficients of $e_j\otimes v_i$ in the displayed
order, and its $i$th local block is $[-R_{2,i},0]$. All other blocks
vanish, by \eqref{cdp:eq:D.26}. Finally
\[
 F_0=\begin{pmatrix}1&1&1\\-1&0&0\\0&-1&0\\0&0&-1\end{pmatrix}
\]
has zero kernel. Thus the relevant exact sequence is
\[
 \Z^{30}\xrightarrow{F_2}\Z^{29}\longrightarrow H_2(Y)
 \longrightarrow\Z^{15}\xrightarrow{F_1}\Z^{15}\longrightarrow H_1(Y)
 \longrightarrow\Z^3\xrightarrow{F_0}\Z^4.
\]

\item Eliminate the local $H_1$ groups and two independent base
meridians. The remaining quotient is
\begin{equation}
 H_1(Y;\Z)=\Lambda/\langle e_1,a_1,a_2,a_3\rangle=0.
\hypertarget{tag-019I}{}\label{cdp:eq:D.31}
\end{equation}
Indeed its relation matrix is
\[
 R_{H_1}=\begin{pmatrix}1&0&0&1\\0&0&1&0\\0&1&0&0\\0&0&1&1\end{pmatrix},
 \qquad\det R_{H_1}=-1.
\]
Since $F_1$ has equal source and target rank fifteen, its integral
surjectivity makes it an isomorphism. Its kernel contributes
nothing to $H_2(Y)$.

In degree two, the pure fibre columns eliminate the fifteen local
generators, leaving the three relations $\eta_i=a_i\wedge c_i$.
The second meridian columns kill $\Lambda\otimes s_2$, the third
then kill $\Lambda\otimes s_1$, and the first leave
$e_1\wedge\Lambda$. Hence
\begin{equation}
 H_2(Y;\Z)=
 \frac{\bigwedge^2\Lambda}
 {e_1\wedge\Lambda+\Z\eta_1+\Z\eta_2+\Z\eta_3}=0.
\hypertarget{tag-019J}{}\label{cdp:eq:D.32}
\end{equation}
The six relations in the order
$e_{12},e_{13},e_{14},\eta_1,\eta_2,\eta_3$ have matrix
\[
R_{H_2}=\begin{pmatrix}
1&0&0&0&1&1\\0&1&0&0&0&-1\\0&0&1&0&1&0\\
0&0&0&0&-1&0\\0&0&0&0&0&-1\\0&0&0&1&1&1
\end{pmatrix},\qquad \det R_{H_2}=1.
\]
Thus $F_2$ is surjective over $\Z$, including all integral torsion.
In the ordered bases above, $\det F_1=-1$, and the Smith normal
forms of $F_1$ and $F_2$ have fifteen and twenty-nine unit invariant
factors, respectively. The relation matrix $R_{H_2}$ also has only
unit invariant factors. The programs are listed in
\cref{tab:II:programs}.

\item The local pushout in Proposition~\ref{night:P4.3} gives
$\pi_1(U_i)=\Lambda/\Z a_i$, and the marked meridian is trivial.
For van Kampen, choose based positive meridians with
$s_1s_2s_3=1$ in $\pi_1(P)$. Write
$\pi_1(W)=\Lambda\times F(s_1,s_2)$, and denote the image of
$x\in\Lambda$ by $t_x$, so $t_xt_y=t_{x+y}$ and each $t_x$
commutes with the two base generators. The boundary maps impose
\[
 t_{a_1}=t_{a_2}=t_{a_3}=1,\qquad
 t_{e_1}s_1=1,\quad s_2=1,\quad (s_1s_2)^{-1}=1.
\]
Thus $s_1=s_2=t_{e_1}=1$. The same unimodular relation matrix
$R_{H_1}$ kills the fibre lattice, proving $\pi_1(Y)=1$.

\item The regular fibre has Euler number zero. Both components of
each special fibre are elliptic bundles over compact curves, and
their intersection is elliptic, so additivity gives zero for each
special support. Stratifying the base and applying multiplicativity
to the smooth bundle gives $e(Y)=0$. The Euler class of the real
bundle underlying a complex rank-three tangent bundle is its top
Chern class with the complex orientation. Hence $c_3(Y)=0$.
Duality and \eqref{cdp:eq:D.31}--\eqref{cdp:eq:D.32} give
$b_4=b_5=0$ and $e(Y)=2-b_3$, so $b_3=2$.
Integral duality and the universal coefficient theorem identify
the torsion of $H_3$ with that of $H_2$, which is zero. They also
give $H_4=H_5=0$, while connectedness and the complex orientation
give $H_0=H_6=\Z$. This proves all the stated integral homology
groups.

\item The universal coefficient theorem gives
$H^1(Y;\Z)=H^2(Y;\Z)=0$. The smooth fibre $A=E^2$ is K\"ahler,
and the constructed $f$ is a proper surjection with connected
fibres from a smooth compact connected threefold. Apply
Proposition~\ref{cdp:st:E1} to this same $f$: every effective
divisor is vertical, and pullback identifies the function fields
through $f^*:\mathcal M(\PP^1)\xrightarrow{\simeq}\mathcal M(Y)$.
Consequently $a(Y)=1$.
\end{enumerate}
\end{proof}


\begin{remark}[Integral coordinates for the quotient maps]
\hypertarget{tag-P214}{}\label{rem:II:torus-quotient-coordinates}
\StatementRecord{rem:II:torus-quotient-coordinates}{proved}{P214}{\SeriesRef{cdp:st:D.2.7}}
For a vector $v\in\Z^n$ with a coordinate equal to $\pm1$, let
$p$ be its first such coordinate and define
\[
 Q_v(x)=\bigl(x_j-(v_j/v_p)x_p\bigr)_{j\ne p},
\]
with the remaining coordinates in their original order. This is
the coordinate quotient by $\Z v$: its kernel is $\Z v$, and
inserting zero in coordinate $p$ gives a right inverse.
For the vectors $a_i$ and $\eta_i=a_i\wedge c_i$ in
Proposition~\ref{cdp:st:D.2.7}, put
$\widehat R_{1,i}=Q_{a_i}$ and $\widehat R_{2,i}=Q_{\eta_i}$,
using the lexicographic exterior-square basis. The exact changes
to the quotient coordinates $R_{1,i},R_{2,i}$ used in its proof are
\[
 R_{1,i}=M_{1,i}\widehat R_{1,i},\qquad
 R_{2,i}=M_{2,i}\widehat R_{2,i},
\]
where
\[
 M_{1,1}=\begin{pmatrix}0&0&1\\1&0&0\\0&1&0\end{pmatrix},\quad
 M_{1,2}=\begin{pmatrix}-1&0&0\\-1&1&0\\0&0&1\end{pmatrix},\quad
 M_{1,3}=\begin{pmatrix}1&0&0\\1&1&0\\0&0&1\end{pmatrix}
\]
have determinants $1,-1,1$, respectively, and for $y\in\Z^5$,
\[
\begin{aligned}
 M_{2,1}y&=(-y_2,-y_4,-y_3,-y_5,y_1),\\
 M_{2,2}y&=(y_3,y_4,-y_1,-y_2,-y_2+y_3+y_5),\\
 M_{2,3}y&=(y_1,y_2,y_3,y_4,y_1+y_4+y_5).
\end{aligned}
\]
Their determinants are $-1,1,1$, respectively. These identities
follow by applying the displayed maps to the basis vectors, using
$R_{1,i}$ from $S_i^{-1}$ and $R_{2,i}$ from
$\bigwedge^2S_i^{-1}$. Thus all six changes are integral
isomorphisms.

Let $\widehat F_1,\widehat F_2$ use these $Q_v$ target coordinates,
with the same source and exterior target bases as $F_1,F_2$.
The signed inclusion maps then give the identities
\[
\begin{aligned}
 F_1&=\operatorname{diag}(I_6,M_{1,1},M_{1,2},M_{1,3})\widehat F_1,\\
 F_2&=\operatorname{diag}(I_{14},M_{2,1},M_{2,2},M_{2,3})\widehat F_2.
\end{aligned}
\]
In particular $\det\widehat F_1=+1$, while $\det F_1=-1$,
and the two coordinate choices have the same Smith invariant
factors in each degree.
\end{remark}

\begin{proposition}[Bundles with primary and secondary Hopf fibres]\hypertarget{tag-01FC}{}\label{cdp:st:C.3.6}
\StatementRecord{cdp:st:C.3.6}{proved}{01FC}{\SeriesRef{cdp:eq:C.H.1}}


There are compact complex threefolds with $b_1=b_2=0$ and $a=1$
whose holomorphic algebraic reduction is a smooth bundle over
$\PP^1$ with primary Hopf fibre. There are also such threefolds
with secondary Hopf fibre. In these examples every fibre has Euler
number zero and $c_3=0$.
\end{proposition}


\begin{proof}

\begin{enumerate}[label=\arabic*.,leftmargin=*]
\item Put $a_1=1$, $a_2=\sqrt2$ and
$W=(\C^2\setminus\{0\})_z\times(\C^2\setminus\{0\})_w$.
The additive complex Lie group $\C$ acts holomorphically by
\begin{equation}
t\cdot(z,w_1,w_2)
 =\bigl(e^t z,e^{ia_1t}w_1,e^{ia_2t}w_2\bigr).
\hypertarget{tag-00L8}{}\label{cdp:eq:C.H.1}
\end{equation}
This action is free. Indeed $e^t z=z$ implies $t=2\pi in$ with
$n\in\Z$; a nonzero $w_j$ then gives $e^{-2\pi a_jn}=1$ and
therefore $n=0$. It is proper as well. If a point and its translate
range in fixed compact subsets of $W$, the $z$-norm bounds
$\operatorname{Re}t$. Writing $s=\operatorname{Im}t$, the identity
\[
\|t\cdot w\|^2
 =e^{-2a_1s}|w_1|^2+e^{-2a_2s}|w_2|^2
\]
and the positive constants $a_1,a_2$ bound $s$ both above and
below, since both $w$-norms are bounded away from zero. The inverse
image of a compact subset under the action graph is consequently
closed and bounded in the group variable and compact in the other
variable.

\item Every orbit meets the real submanifold
$\{\|z\|=\|w\|=1\}=S^3\times S^3$ exactly once.
The first equation fixes $\operatorname{Re}t=-\log\|z\|$.
The displayed squared $w$-norm is strictly decreasing in $s$,
with limits $+\infty$ and $0$, so the second equation has one
solution. Its derivative is strictly negative, and the implicit
function theorem makes the resulting cross-section smooth.
Thus the complex quotient $Y=W/\C$ is a smooth compact complex
threefold diffeomorphic to $S^3\times S^3$. In particular
$b_1(Y)=b_2(Y)=0$ and $e(Y)=0$.

\item The map
\[
f:Y\longrightarrow\PP^1,\qquad [z,w]\longmapsto[z_1:z_2]
\]
is well-defined and holomorphic. Choose a holomorphic nonzero
frame of the tautological line on a small base disc. Scaling the
$z$-coordinate to this frame leaves exactly the transformations
$t=2\pi in$, which act on $w$ by powers of
\[
Q=\operatorname{diag}(q_1,q_2),\qquad
q_j=e^{-2\pi a_j}.
\]
This identifies $f$ over that disc with the product having fibre
$H=(\C^2\setminus\{0\})/\langle Q\rangle$.
The quotient is a primary Hopf surface, and these local products
prove that $f$ is a proper holomorphic submersion with connected
fibres.

\item A meromorphic function on $H$ lifts
to a $Q$-invariant meromorphic function on $\C^2\setminus\{0\}$.
It extends meromorphically across the origin: the closure of its
polar divisor is analytic by the codimension-two extension theorem
for divisors, is locally principal on the smooth ball, and
multiplying by a local defining function reduces extension to
Hartogs' theorem for a holomorphic function. Write the germ as
$g/h$ with relatively prime convergent power series. The identity
$(g\circ Q)h=g(h\circ Q)$ gives
\[
g\circ Q=ug,\qquad h\circ Q=uh
\]
for one holomorphic unit $u$: divisibility follows from relative
primality, and preservation of the vanishing order makes the
quotient a unit. Write $u/u(0)=e^\ell$ with $\ell(0)=0$ and set
\[
v(w)=\exp\left(\sum_{n\ge0}\ell(Q^nw)\right).
\]
The sum converges normally on a sufficiently small ball, because
$Q$ is a strict contraction. It gives
$v(Qw)/v(w)=u(0)/u(w)$. Consequently both $vg$ and $vh$ are
eigenfunctions of pullback by $Q$ with the same eigenvalue $u(0)$.
The monomial $w_1^rw_2^s$ has eigenvalue
$e^{-2\pi(r+s\sqrt2)}$. These numbers are distinct for distinct
$(r,s)\in\Z_{\ge0}^2$. Comparing power-series coefficients shows
that each nonzero eigenfunction is a single monomial and that the
two eigenfunctions are proportional. The meromorphic function is
therefore constant, first near the origin and then everywhere by
analytic continuation. Hence $a(H)=0$.

\item Every meromorphic function on $Y$ is constant on each general
fibre of $f$. On a product chart, restricting to a local section
gives the descended meromorphic function on the base; equality on
the fibres and the identity theorem identify its pullback with
the original function. These local functions agree on overlaps
and extend at the finitely many possible poles on the base.
Thus $\mathcal M(Y)=f^*\mathcal M(\PP^1)$, proving $a(Y)=1$
and proving that $f$ is the holomorphic algebraic reduction.

\item For any integer $m>1$, let $\mu_m$ act on $W$ by
$(z,w)\mapsto(z,\zeta w)$. It commutes with \eqref{cdp:eq:C.H.1} and induces
a free action on $Y$. A fixed orbit would give $t=2\pi in$ from
the $z$-coordinate; equality of norms in a nonzero $w_j$ then
gives $n=0$ and $\zeta=1$. The quotient
$Y_m=Y/\mu_m$ is a smooth compact complex threefold. Transfer
injects $H^j(Y_m;\Q)$ into $H^j(Y;\Q)$, so $b_1(Y_m)=b_2(Y_m)=0$.
The simply connected cover $Y\simeq S^3\times S^3$ gives
$\pi_1(Y_m)=\mu_m$, and the universal coefficient theorem gives
$H^2(Y_m;\Z)\simeq\Z/m\Z$: its free part is zero by the
rational calculation, and its torsion is
$\operatorname{Ext}(H_1(Y_m;\Z),\Z)$. The primary
example has $H^1(Y;\Z)=H^2(Y;\Z)=0$.
The induced map to $\PP^1$ has fibre
\[
H_m=(\C^2\setminus\{0\})/
       \bigl(\langle Q\rangle\times\mu_m\bigr).
\]
Its universal cover is $\C^2\setminus\{0\}$ and its fundamental
group is $\Z\times\mu_m$, which is noncyclic. It is therefore a
secondary Hopf surface, with the displayed primary cover of degree
$m$. Any meromorphic function on $H_m$ pulls back to one on $H$,
so $a(H_m)=0$. The same local descent proves
$\mathcal M(Y_m)=f_m^*\mathcal M(\PP^1)$ and $a(Y_m)=1$.
Transfer of Euler characteristic under the finite covering,
or the smooth-fibration Euler formula, gives $e(Y_m)=0$.

\end{enumerate}
\end{proof}

\subsection{Complete classification and cohomology}\label{sec:II:complete-fibres}

\begin{proposition}[Complete classification and full-fibre cohomology]\hypertarget{tag-01F6}{}\label{cdp:st:C.3.11}
\StatementRecord{cdp:st:C.3.11}{proved}{01F6}{\SeriesRef{cdp:st:C.3.1}, \SeriesRef{cdp:st:C.3.12}, \SeriesRef{cdp:st:C.3.6}, \SeriesRef{cdp:st:D.2.7}, \SeriesRef{cdp:st:D.1.9}, \SeriesRef{cdp:st:C.3.7}, \SeriesRef{cdp:st:C.3.9}, \SeriesRef{cdp:st:C.3.10}, \SeriesRef{cdp:st:C.3.8}, \SeriesRef{cdp:st:D.1.12}}


Under the hypotheses of \ref{cdp:st:C.3.1}, a very general smooth fibre
is a two-torus or a primary or secondary Hopf surface of
algebraic dimension zero. Both Hopf types occur, and the torus case
occurs with a fixed abelian surface as general fibre.
For the entire smooth locus there are exactly the following
two cohomological alternatives:
\begin{equation}
\begin{array}{c|c|c}
\text{smooth fibres}&
 (h^0(\OO),h^1(\OO),h^2(\OO))&
 \text{coherent direct images}\\ \hline
\text{two-tori}&(1,2,1)&
 R^2f_*\OO_Y\simeq\det R^1f_*\OO_Y\\
b_1=1,\ b_2=0&(1,1,0)&
 R^1f_*\OO_Y\simeq\OO_{\PP^1},\
 R^2f_*\OO_Y=0 .
\end{array}
\hypertarget{tag-00VG}{}\label{cdp:eq:C.PK.9}
\end{equation}
The indicated cohomology dimensions hold on every
scheme-theoretic full fibre in the respective branch,
and both direct images satisfy base change there.
In the second branch
$H^q(Y,\OO_Y)=(\C,\C,0,0)$ and
$\Pic^0(Y)=\C\beta$ with the exponential convention
of \ref{cdp:st:D.1.9}.
\end{proposition}


\begin{proof}

\begin{enumerate}[leftmargin=*]
\item \ref{cdp:st:C.3.7}, \ref{cdp:st:C.3.9}, \ref{cdp:st:C.3.10} and \ref{cdp:st:C.3.12} remove the bielliptic,
Kodaira and Inoue types from \ref{cdp:st:C.3.1}; \ref{cdp:st:C.3.8} gives the Hopf algebraic
dimension. The constructions in \cref{cdp:st:C.3.6,cdp:st:D.2.7}
give the two Hopf types and the fixed abelian-surface occurrence.
All smooth fibres are diffeomorphic over the
connected smooth base. In the torus branch they have
$b_1=4$, $b_2=6$, and are K\"ahler by the compact-surface
criterion. Their rational first Chern class and second
Chern number vanish by \ref{cdp:st:C.3.1}. The surface classification
then makes them two-tori: an effective nonzero canonical
divisor has strictly positive degree against a K\"ahler
form, whereas this rational first Chern class has zero
degree. The minimal K\"ahler surfaces of Kodaira dimension
zero with these Betti numbers are exactly the two-tori.
\ref{cdp:st:D.1.12} supplies this branch's full-fibre assertions.

\item In the Hopf branch every smooth fibre has
$b_1=1$ and $b_2=0$. A nonzero holomorphic two-form
is closed and has strictly positive integral against
its conjugate, so it defines a nonzero de Rham
cohomology class. Thus $b_2=0$ forces $p_g=0$.
Since $\chi(\OO)=0$ by \ref{cdp:st:C.3.1}, one has $q=1$.
By \cite[Theorem~2.2, pp.~4--5]{Viaclovsky2026Combs}, $\mathscr E=R^1f_*\OO_Y$ is a line bundle,
and its nowhere-zero transgression section $\beta$
from \ref{cdp:st:D.1.9} trivializes it. The coherent sheaf
$\mathscr G=R^2f_*\OO_Y$ has rank zero, hence finite
length. Hirzebruch--Riemann--Roch on $Y$ and coherent
Leray give the exact equality
\[
 0=\chi(Y,\OO_Y)
   =\chi(\OO_{\PP^1})-\chi(\mathscr E)+\chi(\mathscr G)
   =1-1+\operatorname{length}(\mathscr G).
\]
Therefore $\mathscr G=0$. Top-degree base change and
the parameter-multiplication exact sequence give
$h^2(\OO_{Y_c})=0$ and $h^1(\OO_{Y_c})=1$ on each
full fibre; \cite[Theorem~2.2, pp.~4--5]{Viaclovsky2026Combs} gives $h^0(\OO_{Y_c})=1$.
Leray now gives the stated global structure-sheaf
cohomology. Finally $H^1(Y;\Z)=0$ and the exponential
sequence identify $\Pic^0(Y)$ with
$H^1(Y,\OO_Y)=\C\beta$.

\end{enumerate}
\end{proof}

\section{Cotangent torsion and the Campana--Demailly--Peternell endpoint argument}\hypertarget{tag-0340}{}\label{sec:II:04}

\subsection{Conormal filtration and the Jacobian fractional ideal}\label{subsec:II:cotangent-filtration}

\begin{definition}[Conormal and torsion notation]
\hypertarget{tag-P207}{}\label{def:II:conormal-data}
\StatementRecord{def:II:conormal-data}{proved}{P207}{}
Throughout this section $i:S\hookrightarrow Y$ is the inclusion of a reduced effective Cartier divisor into a complex manifold of dimension $n+1$, $j:S_{\rm reg}\hookrightarrow S$ is the inclusion of its regular locus, and $A$ is an invertible analytic $\OO_S$-module. We write
\[
 N^*=\OO_Y(-S)|_S,\qquad V=\Omega_Y^1|_S,\qquad
 T=\ker(\Omega_S^1\longrightarrow j_*\Omega^1_{S_{\rm reg}}),\qquad Q=\Omega_S^1/T.
\]
All maps between sheaves below are $\OO_S$-linear maps of coherent analytic sheaves, except for the displayed restriction into $j_*\Omega^1_{S_{\rm reg}}$. The latter is also $\OO_S$-linear; its kernel and image have the coherent descriptions proved below.
Put $K=\ker(V\to Q)$; we also write $\widetilde\Omega_S^1=Q$ for the quotient of $\Omega_S^1$ by $T$.
\end{definition}

\begin{proposition}[The conormal filtration and its exact section criterion]\hypertarget{tag-01AD}{}\label{cdp:st:A.1.1}
\StatementRecord{cdp:st:A.1.1}{proved}{01AD}{the definitions and constructions stated in the proof}
\label{wbII:prop:criterion}\label{wbII:prop:filtration}\label{wbII:lem:conormal}\label{wbII:lem:conormal-injective}\label{wbII:prop:missing-global-section-criterion}\label{wbII:prop:missing-cotangent-filtration}
Use the notation of \cref{def:II:conormal-data}. The conormal map $d:N^*\to V$, defined by $[F]\mapsto dF|_S$ for a local defining equation $F$, is injective. There are canonical exact sequences
\begin{align}
0&\longrightarrow N^*\xrightarrow{d}K\longrightarrow T\longrightarrow0,
\hypertarget{tag-003W}{}\label{cdp:A:filtration1}\\
0&\longrightarrow K\longrightarrow V\longrightarrow Q\longrightarrow0.
\hypertarget{tag-003X}{}\label{cdp:A:filtration2}
\end{align}
Let $\partial_A:H^0(S,Q\otimes A)\to H^1(S,K\otimes A)$ and $\delta_A:H^0(S,T\otimes A)\to H^1(S,N^*\otimes A)$ be their connecting maps. Then
\begin{equation}
 H^0(S,V\otimes A)=0
 \quad\Longleftrightarrow\quad
 H^0(S,K\otimes A)=0\ \hbox{and}\ \partial_A\hbox{ is injective}.
\hypertarget{tag-003Y}{}\label{cdp:eq:A.1}
\end{equation}
Under the explicit hypotheses $H^0(S,N^*\otimes A)=0$ and $H^0(S,Q\otimes A)=0$, $H^0(S,V\otimes A)$ is canonically $\ker\delta_A$.
\end{proposition}


\begin{proof}
At a point of $S$, let $R=\mathcal O_{Y,s}$ and choose a reduced local defining equation $g$.
The generator $[g]$ of the conormal line identifies its stalk map with
\begin{equation}
 R/(g)\longrightarrow (\Omega^1_R/g\Omega^1_R),
 \qquad \bar h\longmapsto h\,dg.
 \hypertarget{tag-03DW}{}\label{wbII:eq:CN:2.2}
\end{equation}
The quotient $R/(g)$ is reduced, and its total quotient ring is the product
of the fraction fields of its irreducible components. At the generic point
of every component, $S$ is smooth and $dg$ is nonzero. Thus a class
$\bar h$ in the kernel of \eqref{wbII:eq:CN:2.2} vanishes in each of those
fields. The reduced ring $R/(g)$ injects into their product, so
$\bar h=0$. This proves injectivity in the usual right-exact conormal sequence. Taking the inverse image of $T\subset\Omega_S^1$ gives \eqref{cdp:A:filtration1} and \eqref{cdp:A:filtration2}. Tensoring by $A$ is exact because $A$ is locally free. The cohomology sequence of \eqref{cdp:A:filtration2} identifies the image of $H^0(VA)$ in $H^0(QA)$ with $\ker\partial_A$ and its kernel with $H^0(KA)$, proving \eqref{cdp:eq:A.1}. When $H^0(QA)=0$, it identifies $H^0(VA)$ with $H^0(KA)$. The first sequence then identifies this group with $\ker\delta_A$ when $H^0(N^*A)=0$.
\end{proof}

\begin{proposition}[The Jacobian fractional ideal for every reduced hypersurface]\hypertarget{tag-01AF}{}\label{cdp:st:A.2.1}
\StatementRecord{cdp:st:A.2.1}{proved}{01AF}{\SeriesRef{cdp:A:filtration1}}

Let $\mathcal J\subset\OO_S$ be the Jacobian ideal, generated locally by $F_{x_0},\ldots,F_{x_n}$ modulo $F$. In the total quotient sheaf $\mathcal M_S$ define
\[
 \mathcal J^{-1}=\{q\in\mathcal M_S:q\mathcal J\subset\OO_S\}.
\]
There are canonical identifications
\begin{equation}
 K\simeq N^*\otimes\mathcal J^{-1},\qquad
 T\simeq N^*\otimes(\mathcal J^{-1}/\OO_S),
\hypertarget{tag-004B}{}\label{cdp:eq:A.2}
\end{equation}
under which the first map in \eqref{cdp:A:filtration1} is induced by $\OO_S\hookrightarrow\mathcal J^{-1}$. Thus \eqref{cdp:eq:A.1}, with this explicitly defined fractional ideal, is an exact necessary and sufficient criterion for every reduced hypersurface.
\end{proposition}


\begin{proof}

\begin{enumerate}[label=\arabic*.,leftmargin=*]
\item The ideal $J$ contains a nonzerodivisor: it is nonzero at every minimal prime, and prime avoidance supplies an element outside all minimal primes. Every $R$-linear map $\varphi:J\to R$ is multiplication by a unique element of the total quotient ring. Indeed, for a nonzerodivisor $a\in J$, the identity $a\varphi(b)=b\varphi(a)$ for $b\in J$ gives the multiplier $q=\varphi(a)/a$. Hence $J^{-1}=\Hom_R(J,R)$ is coherent.

\item At the generic points, an element $\alpha\in K$ is uniquely $q\,dF$ for a tuple of meromorphic functions $q$. The coefficients of $\alpha$ belong to $R$ exactly when $qF_{x_i}\in R$ for every $i$, which is the condition $q\in J^{-1}$. Conversely such a vector restricts to zero in the cotangent space of every regular component, hence belongs to $K$. The map $[F]\otimes q\mapsto q\,dF|_S$ proves the first identification. Its quotient by $R\,dF$ gives the second. Replacing $F$ by $uF$, with $u$ a holomorphic unit, replaces $dF|_S$ by $u\,dF|_S$ and replaces the conormal generator by $u[F]$. Coordinate changes transform the derivative vector by an invertible Jacobian matrix. These are exactly the transitions in \eqref{cdp:eq:A.2}, so the identifications are canonical.

\end{enumerate}
\end{proof}

\begin{proposition}[Normalization, conductor, and the excess kernel]\hypertarget{tag-01AG}{}\label{cdp:st:A.2.2}
\StatementRecord{cdp:st:A.2.2}{proved}{01AG}{the definitions and constructions stated in the proof}
\label{wbII:lem:conductor}\label{wbII:lem:conductor-square}
Let $\eta:\widetilde S\to S$ be the finite analytic normalization, $\mathcal B=\eta_*\OO_{\widetilde S}$, and $\mathfrak c=\Ann_{\OO_S}(\mathcal B/\OO_S)$ its conductor. There are canonical inclusions
\begin{equation}
 \mathcal J\subset\mathfrak c,\qquad
 \mathcal B\subset\mathcal J^{-1},\qquad
 \eta_*\eta^*N^*\hookrightarrow K.
\hypertarget{tag-004M}{}\label{cdp:eq:A.3}
\end{equation}
In a local ring $R$ with normalization $B$, put $D=\Hom_B(JB,B)$, viewed as a fractional ideal. Then
\begin{equation}
 J^{-1}=\ker\left[D\longrightarrow
 \Hom_R\bigl(J,(B/\mathfrak c)/(R/\mathfrak c)\bigr)\right],
\quad q\longmapsto\bigl(b\mapsto qb\bmod R\bigr).
\hypertarget{tag-004O}{}\label{cdp:eq:A.4}
\end{equation}
For any locally free $\OO_S$-module $E$, the conductor gives a second exact section criterion:
\begin{multline}
H^0(S,E)=
\ker\left[H^0(\widetilde S,\eta^*E)\oplus H^0(D_c,E|_{D_c})\right.\\
\left.\longrightarrow H^0(\widetilde D_c,\eta^*E|_{\widetilde D_c})\right],
\qquad (a,b)\longmapsto a|_{\widetilde D_c}-\eta^*b,
\hypertarget{tag-004P}{}\label{cdp:eq:A.5}
\end{multline}
where $D_c$ and $\widetilde D_c$ are defined by $\mathfrak c$ in $S$ and $\widetilde S$. In particular, the ambient group vanishes exactly when the map in \eqref{cdp:eq:A.5}, for $E=V\otimes A$, is injective.
\end{proposition}


\begin{proof}

\begin{enumerate}[label=\arabic*.,leftmargin=*]
\item Apply \cite[Proposition~3.4, pp.~9--10, and Corollary~3.5, p.~10, arXiv v7]{GS2014} with their reduced hypersurface $D=S$ and local equation $h=F$. Their residue module is
$\mathcal R_D=\operatorname{Hom}_R(J,R)=J^{-1}$, where a meromorphic multiplier $q$ acts by $b\mapsto qb$; their conductor is $\mathcal C_D=\operatorname{Ann}_R(B/R)=\mathfrak c$.
Their inclusions give $J\subset\mathfrak c$ and $B\subset J^{-1}$.
Tensoring the latter inclusion by $N^*$ and using \eqref{cdp:eq:A.2} and the finite projection formula gives the last inclusion in \eqref{cdp:eq:A.3}.

\item The module $D$ consists exactly of the meromorphic $q$ satisfying $qJ\subset B$. Requiring that these products descend to $R$ is precisely the vanishing of their classes in $B/R=(B/\mathfrak c)/(R/\mathfrak c)$, proving \eqref{cdp:eq:A.4}. Finally the ring map
\[
 R\longrightarrow B\oplus R/\mathfrak c,
 \qquad r\longmapsto(r,r\bmod\mathfrak c)
\]
has image the kernel of $(b,\bar r)\mapsto b-\bar r\in B/\mathfrak c$. Surjectivity onto $B/\mathfrak c$ comes from the first summand. Tensoring this exact sequence by the free module corresponding to $E$ and taking global sections gives \eqref{cdp:eq:A.5}. All conductor restrictions are scheme-theoretic, so this argument includes nonreduced conductors.

\end{enumerate}
\end{proof}

\begin{cited}[The normalization criterion]\hypertarget{tag-01AH}{}\label{cdp:st:A.2.3}
\StatementRecord{cdp:st:A.2.3}{proved}{01AH}{\SeriesRef{cdp:st:A.2.1}, \SeriesRef{cdp:st:A.2.2}; \cite[Theorems~1.1 and~1.3, p.~3; proof of Theorem~1.3, p.~12; Proposition~3.4 and Corollary~3.5, pp.~9--10, arXiv v7]{GS2014}}
For a reduced hypersurface $S$ with the notation of
\cref{cdp:st:A.2.1,cdp:st:A.2.2}, the canonical inclusion
$\eta_*\eta^*N^*\hookrightarrow K$ is an isomorphism exactly when
$S$ has normal crossings outside an analytic subset of codimension
at least two in $S$.
\end{cited}
\begin{proof}[Application of the cited residue criterion]
Apply \cite[Proposition~3.4 and Corollary~3.5, pp.~9--10,
arXiv v7]{GS2014} with their hypersurface $D=S$, local equation
$h=F$, local ring $R=\mathcal O_S$, and normalization ring
$B=\eta_*\mathcal O_{\widetilde S}$. The multiplier map
$q\mapsto(b\mapsto qb)$ identifies their residue module
$\mathcal R_D$ with $J^{-1}=\operatorname{Hom}_R(J,R)$ as
submodules of the same total quotient ring. Their weakly
holomorphic functions form $B$, and their inclusion
$B\subset\mathcal R_D$ is the inclusion of \cref{cdp:st:A.2.2}.
Theorem~1.1, implication $(B)\Rightarrow(C)$, and Theorem~1.3
in \cite[p.~3; proof of Theorem~1.3, p.~12]{GS2014} give
$J^{-1}=B$ exactly for the asserted normal-crossing locus.
Tensor this equality by the invertible sheaf $N^*$.
The finite projection formula and the map
$[F]\otimes q\mapsto q\,dF|_S$ of \cref{cdp:st:A.2.1}
identify its two sides with $\eta_*\eta^*N^*$ and $K$, and
identify their inclusion with the canonical map in the statement.
\end{proof}

\begin{proposition}[The canonical normal-crossings formula]\hypertarget{tag-01AE}{}\label{cdp:st:A.1.2}
\StatementRecord{cdp:st:A.1.2}{proved}{01AE}{\SeriesRef{cdp:st:A.2.1}}

For a reduced normal-crossings divisor, the map in \eqref{cdp:eq:A.3} is an isomorphism. In a chart with $F=x_1\cdots x_r$, it is explicitly
\begin{equation}
 \bigoplus_{i=1}^r R/(x_i)\longrightarrow K,
 \qquad (a_i)_i\longmapsto
 \sum_{i=1}^r a_i\prod_{j\ne i}x_j\,dx_i.
\hypertarget{tag-004Q}{}\label{cdp:eq:A.6}
\end{equation}
The tensor factor $N^*$ is trivialized here by $[F]$. At a double crossing $T\simeq R/(x,y)$; at a triple crossing,
\begin{equation}
 T\simeq R^3/\bigl(R(1,1,1)+R(x,0,0)+R(0,y,0)+R(0,0,z)\bigr).
\hypertarget{tag-004R}{}\label{cdp:eq:A.7}
\end{equation}
\end{proposition}


\begin{proof}

\begin{enumerate}[label=\arabic*.,leftmargin=*]
\item On the normalization branch $x_i=0$, an ambient form maps to zero in the cotangent bundle of that branch precisely when all coefficients except that of $dx_i$ vanish. Thus a coefficient of $dx_i$ must vanish on every branch $x_j=0$ for $j\ne i$. In the convergent power-series ring it is divisible by $\prod_{j\ne i}x_j$. Its remaining factor is defined modulo $x_i$, which gives exactly \eqref{cdp:eq:A.6}. The coefficients of different $dx_i$ show injectivity. Coefficients of the coordinates not among $x_1,\ldots,x_r$ vanish on every branch and hence vanish in $R$. This proves surjectivity in every dimension.

\item For $r=1$ the map is $a\mapsto a\,dx$. For $r=2$ it is $(a,b)\mapsto ay\,dx+bx\,dy$, and quotienting by $dF=y\,dx+x\,dy$ leaves the generator $y\,dx=-x\,dy$, annihilated by $(x,y)$. For $r=3$ the branch generators are $yz\,dx,xz\,dy,xy\,dz$. Their sum is $dF$ and their respective annihilators are $(x),(y),(z)$, proving \eqref{cdp:eq:A.7}. The equation and coordinate transition maps are those of \cref{cdp:st:A.2.1}. They preserve the canonical fractional-ideal map whose restriction is \eqref{cdp:eq:A.6}; permutation of the branches permutes the same direct-sum factors and terms. Thus the local formulas glue to that canonical morphism.

\end{enumerate}
\end{proof}

\subsection{Residue poles and singular critical curves}\label{subsec:II:cotangent-residues}

\begin{proposition}[Pinch points and cuspidal edges]\hypertarget{tag-01AI}{}\label{cdp:st:A.2.4}
\StatementRecord{cdp:st:A.2.4}{proved}{01AI}{\SeriesRef{cdp:st:A.2.1}, \SeriesRef{cdp:st:A.1.2}; \cite[Proposition 3.4, pp.~9--10, and Corollary 3.5, p.~10, author version]{GS2014}; \cite[Theorem 1, p.~1]{Simonic2016Hartogs}}

At the pinch point $F=x^2-y^2z$, with normalization $x=uv,y=u,z=v^2$, write $B=\C\{u,v\}$. Then
\begin{equation}
 K=B\,dF,\qquad T\simeq B/R\simeq R/(x,y),
\quad \tau=2yz\,dx-2xz\,dy-xy\,dz
\hypertarget{tag-004S}{}\label{cdp:eq:A.8}
\end{equation}
is a torsion generator. At the cuspidal edge $F=x^2-y^3$, with normalization $x=t^3,y=t^2$ and unchanged coordinate $z$, write $B=\C\{t,z\}$. Then
\begin{equation}
 K=t^{-1}B\,dF,\qquad
 T\simeq t^{-1}B/R\simeq R/(x,y^2),
 \quad\tau=2y\,dx-3x\,dy.
\hypertarget{tag-004T}{}\label{cdp:eq:A.9}
\end{equation}
The conductor is $(x,y)$ for the pinch point and $(x,y)$ for the cusp; in the cusp normalization it is $t^2B$. An ordinary surface, meaning a surface whose singularities are double normal crossings, triple normal crossings, and pinch points, satisfies canonically $K=\eta_*\eta^*N^*$ and $T=N^*\otimes(\eta_*\OO_{\widetilde S}/\OO_S)$. For the cuspidal edge the additional quotient $J^{-1}/B$ is generated by $t^{-1}$ over $\C\{z\}$.
\end{proposition}


\begin{proof}

\begin{enumerate}[label=\arabic*.,leftmargin=*]
\item For the pinch point,
\[
 R=\C\{u,v^2,uv\}=\C\{v^2\}+uB,
 \quad J=(uv,uv^2,u^2)=u(u,v)B.
\]
The equality of the two descriptions of $J$ follows by separating even and odd powers of $v$. An element $q\in J^{-1}$ satisfies $(qu)(u,v)\subset B$. Since $B$ is a regular local ring, $\Hom_B((u,v),B)=B$: a meromorphic function whose products with $u$ and $v$ are holomorphic has no pole along any divisor, and its possible isolated singularity is removable by the holomorphic extension theorem \cite[arXiv version~1, Theorem~1, p.~1]{Simonic2016Hartogs}. Hence $qu=h\in B$. The requirements $hv,hv^2\in R$, restricted to $u=0$, say respectively that $h(0,v)v$ and $h(0,v)v^2$ are even power series. The first requires $h(0,v)$ to be odd and the second to be even. Therefore $h(0,v)=0$, so $h\in uB$ and $q\in B$. The reverse inclusion follows either directly or from \eqref{cdp:eq:A.3}. Also $B/R$ is generated by $v$ and its annihilator is $(u,uv)=(y,x)$. Multiplying $dF$ by $v=x/y$ gives the displayed $\tau$.

\item For the cusp, $R=\C\{t^2,t^3,z\}$ and $J=(t^3,t^4)$. A meromorphic $q$ with $qt^3,qt^4\in R$ has no poles away from $t=0$; writing its Laurent expansion in $t$ shows that the lowest allowed exponent is $-1$. Indeed an exponent $-3$ produces a forbidden $t$ term after multiplication by $t^4$, an exponent $-2$ produces one after multiplication by $t^3$, and smaller exponents produce a pole. Conversely every $t^{-1}B$ has both products in $R$. Thus $J^{-1}=t^{-1}B$. Modulo $R$, the quotient has the $\C\{z\}$-basis $t^{-1},t$; multiplication by $y$ sends the first basis vector to the second, while $x$ and $y^2$ kill both. This proves the cyclic presentation and \eqref{cdp:eq:A.9}. The conductor calculations follow directly from the displayed descriptions of $R\subset B$. Formula \eqref{cdp:eq:A.2} and the three local computations prove the ordinary-surface assertion by gluing.

\end{enumerate}
\end{proof}

\begin{proposition}[Vertical residue poles at a tacnode--cusp intersection]\hypertarget{tag-01AJ}{}\label{cdp:st:A.2.5}
\StatementRecord{cdp:st:A.2.5}{proved}{01AJ}{the definitions and constructions stated in the proof}

For $F=x^2-y^4z^3$, normalization is $x=u^2v^3$, $y=u$, $z=v^2$. Put $A_0=\C\{u,v^2\}$ and $B=\C\{u,v\}$. Then
\begin{equation}
 \begin{aligned}
 R&=A_0\oplus u^2v^3A_0,&\mathfrak c&=u^2v^2B,\\
 J^{-1}&=A_0\oplus(uv)^{-1}A_0,&
 J^{-1}/B&\simeq A_0/(uv^2).
 \end{aligned}
 \hypertarget{tag-004U}{}\label{cdp:eq:A.9a}
\end{equation}
In a trivialization of $N^*$ the torsion module is
\begin{equation}
 T\simeq R/(x,y^3z^2),\qquad
 \tau=2yz\,dx-4xz\,dy-3xy\,dz.
 \hypertarget{tag-004V}{}\label{cdp:eq:A.9b}
\end{equation}
In particular $v^{-1}\in J^{-1}\setminus B$. For the ruling with base coordinate $v$, this is a residue constant on its generic fibres with a pole on the vertical divisor $v=0$, although the horizontal conductor $u=0$ has coefficient two.
\end{proposition}


\begin{proof}

\begin{enumerate}[label=\arabic*.,leftmargin=*]
\item The element $v=x/(y^2z)$ belongs to the fraction field and satisfies $v^2=z$. Adjoining it gives the regular ring $B$, which is therefore the finite normalization. Separating even and odd powers of $v$ gives the displayed description of $R$. An even monomial $u^av^b$ belongs to the conductor exactly when $a\ge2$ and $b\ge2$; an odd monomial already in $R$ has $a\ge2$, $b\ge3$ and is then in the conductor. Thus $\mathfrak c=u^2v^2B$.

\item The normalized Jacobian generators, with irrelevant nonzero constants removed, are $u^2v^3,u^3v^6,u^4v^4$. If $qJ\subset R$, the first product shows that $q$ has no other polar divisor, with possible orders at most two along $u=0$ and three along $v=0$. A term of order $v^{-3}$ would produce a forbidden $v^1$ term in $qu^4v^4$; a term of order $v^{-2}$ would produce one in $qu^2v^3$. Both coefficients must vanish, so $q$ has order at least $-1$ in $v$. A coefficient $a(v)u^{-2}$ would require both $a(v)v^3$ and $a(v)v^6$ to be even Laurent series, from the $u^0$ and $u^1$ coefficients of the first two products. The first requirement makes $a$ odd and the second makes it even, so $a=0$. Thus $q=t/(uv)$ with $t\in B$.

\item Write $t=a(u,v^2)+v b(u,v^2)$. The first Jacobian product equals $uv^2t$ and belongs to $R$ exactly when $b$ is divisible by $u$. The other two products automatically belong to $R$ under this condition. This gives $J^{-1}=A_0\oplus(uv)^{-1}A_0$. The odd part of $B$ is $(uv)^{-1}(uv^2)A_0$, giving $J^{-1}/B$ in \eqref{cdp:eq:A.9a}. The odd part of $R$ is $(uv)^{-1}u^3v^4A_0$, so $J^{-1}/R$ is generated by $(uv)^{-1}$ with annihilator $(x,y^3z^2)$. Multiplication of $dF$ by $(uv)^{-1}$ gives \eqref{cdp:eq:A.9b}.

\item The particular vertical pole has the completely regular gradient products
\[
 v^{-1}dF=2y^2z\,dx-4xyz\,dy-3xy^2\,dz.
\]
Pulling this expression to $B$ gives zero. It is a nonzero ambient-kernel vector with multiplier $v^{-1}$, which has a pole on the normalization.  

\end{enumerate}
\end{proof}

\begin{proposition}[A cusp along a non-complete-intersection curve]\hypertarget{tag-01AK}{}\label{cdp:st:A.2.6}
\StatementRecord{cdp:st:A.2.6}{proved}{01AK}{\SeriesRef{cdp:st:A.5.5}}

Put
\begin{equation}
 \begin{aligned}
 I&=(y^2-xz,\ yz^2-x^4,\ z^3-x^3y)\subset\C\{x,y,z\},\\
 F&=4x^7-15x^3yz^2+10x^2y^3z-3xy^5+4z^5.
 \end{aligned}                                                   \hypertarget{tag-006J}{}\label{cdp:eq:A.NL.1}
\end{equation}
The reduced critical locus of the reduced hypersurface $F=0$ is the
irreducible curve $D$ parametrized by $(x,y,z)=(t^5,t^6,t^7)$.
Its transverse type at every point with $t\ne0$ is the ordinary cusp
$A_2$. The curve is not a local complete intersection at the origin,
and
\begin{equation}
             F\in I^{(2)}\setminus I^2.                         \hypertarget{tag-006K}{}\label{cdp:eq:A.NL.2}
\end{equation}
Thus the ordinary-square membership used to define the conormal
quadratic jet in \ref{cdp:st:A.5.5} does not extend to an arbitrary reduced
critical curve, even with generic transverse type $A_2$.

\end{proposition}


\begin{proof}

\begin{enumerate}[label=\arabic*.,leftmargin=*]
\item Write the three displayed generators of $I$ as $g_1,g_2,g_3$.
The quotient is a finite $\C\{x\}$-module generated by
$1,y,z,yz,z^2$: their products by $y$ and $z$ reduce to this list
using $y^2=xz$, $yz^2=x^4$, and $z^3=x^3y$.
For example these identities also give $y^5=x^6$ and $z^5=x^7$,
which make the analytic quotient finite over $\C\{x\}$.
The parametrization sends the five generators to
$1,t^6,t^7,t^{13},t^{14}$. Their exponents have distinct residues
modulo five, so they are linearly independent over $\C\{t^5\}$.
Consequently this ring homomorphism is injective and has
image $\C\{t^5,t^6,t^7\}$; it proves the stated normalization and
the primeness of $I$. The three generators are minimal. Indeed their
initial terms have degrees two, three, three; after the degree-two
coefficient vanishes, the cubic terms $yz^2$ and $z^3$ are independent
modulo the span of $xg_1,yg_1,zg_1$. A height-two ideal with three
minimal generators is not a complete intersection.

\item The following polynomial identities determine the critical locus:
\begin{equation}
 \begin{aligned}
 F_y&=-15xg_1^2,\\
 F_x&=(17xyz-3y^3)g_1-28x^2g_2,\\
 F_z&=10x^2yg_1+20zg_3.
 \end{aligned}                                                   \hypertarget{tag-006L}{}\label{cdp:eq:A.NL.3}
\end{equation}
At $x=0$, the equations $F_x=F_z=0$ give $y=z=0$.
At $x\ne0$, $F_y=0$ gives $g_1=0$, and then $F_x=0$
gives $g_2=0$. The value $z=0$ is impossible there, so $F_z=0$
also gives $g_3=0$. Conversely all three partials vanish on $D$.
Weighted Euler's identity, with weights $(5,6,7)$ and degree $35$,
puts $F$ in its Jacobian ideal. The critical locus is therefore
exactly $D$ as a reduced set. A repeated factor of $F$ would put
a two-dimensional component in this critical locus, which proves
that $F$ is reduced.

\item The Hessian along the normalization has the exact factorization
\begin{equation}
 \operatorname{Hess}(F)|_D
 =2t^{21}
 \begin{pmatrix}7t^2\\0\\-5\end{pmatrix}
 \begin{pmatrix}7t^2&0&-5\end{pmatrix},
 \qquad \frac{1}{6}F_{yyy}|_D=-20t^{17}.                       \hypertarget{tag-006M}{}\label{cdp:eq:A.NL.4}
\end{equation}
For $t\ne0$, the plane $x=t^5$ is transverse to $D$.
The restricted Hessian has rank one and its radical is the
$y$ direction; its cubic coefficient in that direction is nonzero.
The holomorphic splitting lemma and a local cube root of a unit
give the transverse cusp equation. At the generic point of $D$,
its conormal is locally free and injects into ambient differentials.
Since $F$ and $dF$ vanish there, $F\in I^2\C\{x,y,z\}_I$,
which is precisely $F\in I^{(2)}$.

\item To decide ordinary-square membership, use the positive weights
$(5,6,7)$. The three generators have weights $12,20,21$.
Among products of two generators, only $g_1^2$ can contribute
to weight $35$: it needs a coefficient of weight $11$, whose
space is spanned by $xy$. Products $g_1g_2$ and $g_1g_3$ would
need weights $3$ and $2$, which do not occur; the other products
already have greater weight. Thus the weight-$35$ part of $I^2$
is exactly $\C\,xyg_1^2$. Its $x^7$ coefficient is zero, whereas
that of $F$ is four. Taking the weight-$35$ component of a
convergent series identity is legitimate because all weights are
positive. This proves $F\notin I^2$ in the analytic local ring.

\item The normalization conductor
is $t^{10}\C\{t\}$: the semigroup gaps are $1,2,3,4,8,9$.
The $B$-module generated by the differential image, where $B=\C\{t\}$, is $t^4B\,dt$.
The smallest valuation of a two-by-two minor of the displayed
generators' Jacobian matrix is $19$. Thus the normalized conormal
determinant correction is $19-4=15$, while the conductor exponent
is $10$; the complete-intersection adjunction comparison cannot
identify these two quantities here.

\item The symbolic-square correction is visible directly on the free
normalized conormal. Set $a=dg_1|_D$ and
$b=dg_2|_D+2t^8a$. Then
\begin{equation}
 \begin{aligned}
 a&=(-t^7,2t^6,-t^5),& b&=(-6t^{15},5t^{14},0),\\
 dg_3|_D&=-3t^9a+tb,&
 q&=t^{-5}(5t^8a-2b)^2.
 \end{aligned}                                                   \hypertarget{tag-006N}{}\label{cdp:eq:A.NL.5}
\end{equation}
The first two rows generate the entire normalized conormal image,
and $q$ is the quadratic term, namely half the Hessian.
If $A,B$ are the dual basis, its radical frame $2A+5t^8B$
is represented in the ambient normal quotient by $t^{-6}\partial_y$.
Its intrinsic cubic is therefore $-20t^{-1}$. These are
quadratic and cubic poles of orders five and one. The corresponding
combined invariant $q^3k^2$ has pole order $17$.


\end{enumerate}
\end{proof}

\subsection{Nodal families and the endpoint groups}\label{subsec:II:cotangent-nodal}

\begin{proposition}[A complete nodal-family computation]\hypertarget{tag-01BK}{}\label{cdp:st:B1}
\StatementRecord{cdp:st:B1}{proved}{01BK}{the definitions and constructions stated in the proof}
\hypertarget{tag-03IB}{}\label{wbII:lem:gluing-bundle}
\StatementRecord{wbII:lem:gluing-bundle}{proved}{03IB}{}
\hypertarget{tag-03I9}{}\label{wbII:lem:extending-gluing}
\StatementRecord{wbII:lem:extending-gluing}{proved}{03I9}{}
\hypertarget{tag-03IS}{}\label{wbII:rem:scope}
\StatementRecord{wbII:rem:scope}{proved}{03IS}{}
\label{wbII:thm:nodal-test}\label{wbII:prop:family-geometry}\label{wbII:prop:nodal-counterexample}
Let $0<\varepsilon<4/27$, let $\Delta=\{|t|<\varepsilon\}$, and define
\begin{equation}
\mathcal C=
\{([x:y:z],t)\in\mathbb P^2\times\Delta:
y^2z=x^2(x+z)+tz^3\}.
\hypertarget{tag-03I1}{}\label{wbII:eq:W:5.8}
\end{equation}
Put $X=\mathcal C\times\PP^1$.

The projection $X\to\Delta$ is proper, its total space is smooth, and its central fibre is $S=C_0\times\PP^1$, where $C_0$ is an irreducible nodal cubic. For each $\lambda\in\C^*$ let $A_\lambda$ be the line bundle whose pullback to the normalization $\nu:\PP^1_u\to C_0$ is trivial and whose sections in this trivialization satisfy $h(1)=\lambda h(-1)$. Put $A=\operatorname{pr}_1^*A_\lambda$ and $d_\lambda=1$ for $\lambda=1$, and $d_\lambda=0$ otherwise. In the order $q=0,1,2$, the coherent cohomology dimensions are
\begin{equation}
\begin{array}{c|c}
\text{sheaf on }S & (h^0,h^1,h^2)\\\hline
 N^*_{S/X}\otimes A &(d_\lambda,d_\lambda,0)\\
 K_S\otimes A &(1,0,0)\\
 T_S\otimes A &(1,0,0)\\
 \widetilde\Omega^1_S\otimes A &(0,1+d_\lambda,d_\lambda)\\
 \Omega^1_S\otimes A &(1,1+d_\lambda,d_\lambda)\\
 \Omega^1_X|_S\otimes A &(1,1+d_\lambda,d_\lambda).
\end{array}                                                    \hypertarget{tag-008G}{}\label{cdp:eq:B.1}
\hypertarget{tag-03DZ}{}\hypertarget{tag-03E1}{}\SsixEquationAliases{\label{wbII:eq:CN:4.5}\label{wbII:eq:CN:4.6}}
\end{equation}
Here $T_S$ denotes the torsion of the cotangent sheaf, and $K_S$ its inverse image in the ambient cotangent sheaf. In particular, $H^0(S,N^*_{S/X}\otimes A)$ and $H^0(S,\widetilde\Omega^1_S\otimes A)$ both vanish exactly for $\lambda\ne1$, and the ambient space of sections has dimension one for every $\lambda\in\C^*$. After shrinking $\Delta$, every $A_\lambda$ extends to a line bundle on $\mathcal C$ of the form $\OO_{\mathcal C}(P-Q)$ for holomorphic sections $P,Q$.
The tensor order of $A_\lambda$ equals the multiplicative order of
$\lambda$; in particular the bundles for non-roots of unity have infinite
order. The explicit sections \eqref{wbII:eq:CN:4.5a} give the value
$\lambda=1/3$ with the stated gluing convention.
\end{proposition}



\begin{proof}

\begin{enumerate}[label=\arabic*.,leftmargin=*]
\item In the affine chart $z=1$, the defining equation is $F=y^2-x^3-x^2-t=0$. Its partial derivative in $t$ is $-1$, so the total space is smooth there. Its only point with $z=0$ is $[0:1:0]$. In the chart $y=1$ the equation is $z-x^3-x^2z-tz^3=0$, whose $z$ derivative is one at this point. Fibre singularities in the affine chart satisfy $y=0$ and $x(3x+2)=0$; they occur at $t=0,-4/27$. Our disc contains only the first. Near the node choose the holomorphic square root
$h(x)=\sqrt{1+x}$ with $h(0)=1$ and put
$v=y-xh(x)$, $w=y+xh(x)$. Then $vw=y^2-x^2-x^3=t$.
The Jacobian of $(x,y)\mapsto(v,w)$ at the origin is
$\left(\begin{smallmatrix}-1&1\\1&1\end{smallmatrix}\right)$,
of determinant $-2$, so these are local coordinates. Properness follows from the closed inclusion into $\PP^2\times\Delta$.

\item The normalization is
\begin{equation}
 \nu:\PP^1\longrightarrow C_0,\qquad
 [u:v]\longmapsto
 [(u^2-v^2)v:\,u(u^2-v^2):\,v^3]
 \hypertarget{tag-03DX}{}\label{wbII:eq:CN:4.3}
\end{equation}
On the chart $v=1$ this is $u\mapsto(x,y)=(u^2-1,u(u^2-1))$.
At $v=0$ it maps to the smooth point at infinity. On the affine part
with $x\ne0$ its inverse parameter is $u=y/x$; the equation of $C_0$
then gives $x=u^2-1$. The only affine point with $x=0$ is the node,
whose preimages are $u=1,-1$. These formulas show surjectivity and
generic injectivity. The nonconstant map from the projective line is
finite, hence is the normalization. Its product with the identity of
$\PP^1$ is the normalization $\eta$ of $S$. The descent sequence, an exact sequence of coherent $\OO_{C_0}$-modules, is
\begin{equation}
0\longrightarrow A_\lambda\longrightarrow\nu_*\OO_{\PP^1}
 \xrightarrow{\ h\mapsto h(1)-\lambda h(-1)\ }\C_{\rm node}
 \longrightarrow0.                                      \hypertarget{tag-008H}{}\label{cdp:eq:B.2}
\end{equation}
Its cohomology is the two-term complex $\C\xrightarrow{1-\lambda}\C$. Thus $h^0(A_\lambda)=h^1(A_\lambda)=d_\lambda$.
Pullback along $\operatorname{pr}_1$ and restriction to
$C_0\times\{r_0\}$ for a fixed $r_0\in\PP^1$ are inverse on global
sections: in a local frame for $A_\lambda$, holomorphic dependence on
the compact second factor is constant. Multiplication of the normalized
frames multiplies their gluing scalars. Consequently
\begin{equation}
\begin{gathered}
 \nu^*A_\lambda\simeq\cO_{\PP^1},\qquad
 H^0(S,A)\simeq H^0(C_0,A_\lambda)\simeq\C^{d_\lambda},\\
 A_\lambda^{\otimes m}\simeq A_{\lambda^m}\qquad(m\ge1).
\end{gathered}
 \hypertarget{tag-03DY}{}\label{wbII:eq:CN:4.4}
\hypertarget{tag-03I2}{}\hypertarget{tag-03E2}{}\SsixEquationAliases{\label{wbII:eq:W:5.9}\label{wbII:eq:CN:4.7}}
\end{equation}
The sequence \eqref{cdp:eq:B.2} shows that $A_\mu$ is trivial exactly
when $\mu=1$, so the tensor identity proves the assertion about order.

\item Write $V=\Omega^1_{\mathcal C}|_{C_0}$. The normalization map induces a surjective $\OO_{C_0}$-linear morphism $\Omega^1_{C_0}\to\nu_*\Omega^1_{\PP^1}$; at a node $vw=0$, it sends $dv$ to $(dv,0)$ and $dw$ to $(0,dw)$. Its kernel is the length-one module generated by $v\,dw=-w\,dv$. Consequently
\[
 T_{C_0}=\C_{\rm node},\qquad
 \widetilde\Omega^1_{C_0}=\nu_*\OO_{\PP^1}(-2).
\]
The kernel of $V\to\nu_*\Omega^1_{\PP^1}$ is $\nu_*\OO_{\PP^1}$. Locally its generators are $w\,dv$ and $v\,dw$ on the two branches; the identification is the normal-crossing conormal identification of \cref{cdp:st:A.1.2} with the globally fixed conormal $dt$. Tensoring by $A_\lambda$ preserves exactness, and finite pushforward and the projection formula identify the two resulting normalization terms with $\nu_*\OO$ and $\nu_*\OO(-2)$. Hence
\[
 (h^0,h^1)(C_0,V\otimes A_\lambda)=(1,1),\qquad
 (h^0,h^1)(C_0,\Omega^1_{C_0}\otimes A_\lambda)=(1,1).
\]
The product morphisms give the direct decompositions
\[
 \Omega^1_X|_S\otimes A
 =(V\otimes A_\lambda)\boxtimes\OO
 \ \oplus\ A_\lambda\boxtimes\OO(-2),
\]
and the analogous decomposition for $\Omega^1_S$. The torsion-free quotient is obtained by replacing the first factor by $\nu_*\OO(-2)$. The torsion is $\C_{\rm node}\boxtimes\OO$, its ambient inverse image is $\nu_*\OO\boxtimes\OO$, and $N^*_{S/X}=\OO_S$. The normalization pullback on the torsion-free quotient is the
identity on $\nu_*\Omega^1_{\PP^1}\boxtimes\OO$ and is induced by
$\OO_{C_0}\hookrightarrow\nu_*\OO_{\PP^1}$ on the other summand.
This gives the injective sheaf map
\begin{equation}
 \widetilde\Omega^1_S\hookrightarrow\eta_*\Omega^1_{\PP^1\times\PP^1}.
 \hypertarget{tag-03E3}{}\label{wbII:eq:CN:4.8}
\end{equation}
Tensoring by $A$ and using $\eta^*A\simeq\OO$ gives
\begin{equation}
 H^0(S,\widetilde\Omega^1_S\otimes A)
 \hookrightarrow
 H^0(\PP^1\times\PP^1,\Omega^1_{\PP^1\times\PP^1})=0.
 \hypertarget{tag-03E4}{}\label{wbII:eq:CN:4.9}
\hypertarget{tag-03EA}{}\SsixEquationAliases{\label{wbII:eq:W:5.10}}
\end{equation}
Here the two product summands of $\Omega^1_{\PP^1\times\PP^1}$ have no
global sections because $H^0(\PP^1,\OO(-2))=0$.
The coherent K\"unneth isomorphism and $H^0(\PP^1,\OO)=\C$, $H^1(\PP^1,\OO(-2))=\C$ now give every entry of \eqref{cdp:eq:B.1}. All other cohomology of these two line bundles on $\PP^1$ vanishes.

\item We give the extension explicitly. Fix $Q(0)$ at $u=0$, and take $P(0)$ at $u=a=(1-\lambda)/(1+\lambda)$; for $\lambda=-1$ take $a=\infty$. For $\lambda\ne0$ both points are in the smooth locus. With $b=0$, the rational function $r(u)=(u-a)/(u-b)$ has divisor $P(0)-Q(0)$, and a local section in its pullback trivialization descends precisely with scalar
\[
 \frac{r(1)}{r(-1)}
 =\frac{(1-a)(1+b)}{(1-b)(1+a)}=\lambda.
\]
For $a=\infty$ the same expression is its limit, namely $-1$. The projection $\mathcal C\to\Delta$ is a submersion at these points, so the holomorphic implicit-function theorem extends them to sections on a smaller disc. The point at infinity itself is a constant section. For $\lambda=1$ take $P=Q$ identically. Sections through the smooth fibre locus are effective Cartier divisors in the smooth surface $\mathcal C$; their difference therefore defines the claimed extension. Pullback along $X\to\mathcal C$ extends $A$. This proves the claim also for all roots of unity, including the identity. With the gluing convention of \eqref{cdp:eq:B.2}, the only exceptional scalar for the conormal vanishing is $1$.
One concrete extending choice is
\begin{equation}
 O(t)=[0:1:0],\qquad
 P(t)=[3:\sqrt{36+t}:1]
\hypertarget{tag-03E0}{}\label{wbII:eq:CN:4.5a}
\end{equation}
where the square root has value $6$ at $t=0$. The sections are disjoint
and lie in the smooth locus on a smaller disc, so
$B=\OO_{\mathcal C}(P-O)$ is a line bundle. On the normalization,
$O(0)$ is $u=\infty$ and $P(0)$ is $u=2$.
The function $r(u)=u-2$ has divisor $[2]-[\infty]$;
the frame $1/r$ of its associated line bundle gives descent scalar
$r(1)/r(-1)=1/3$ in \eqref{cdp:eq:B.2}. Thus
$\widetilde A=\operatorname{pr}_1^*B$ restricts to $A$ for
$\lambda=1/3$. Put $\widetilde L=\omega_X\otimes\widetilde A^{-1}$.
Then $(\widetilde L^\vee\otimes\omega_X)|_S\simeq A$, so the same
example realizes the Serre-dual twist in the endpoint inference of
\cite[p.~692]{CDP2020}. The projection $X\to\Delta$ is surjective;
its total space is noncompact, since otherwise its image $\Delta$
would be compact. The cited Theorem~2.2 concerns compact threefolds.


\end{enumerate}
\end{proof}

\begin{proposition}[A projective example with twelve length-one contributions]\hypertarget{tag-01BL}{}\label{cdp:st:B2}
\StatementRecord{cdp:st:B2}{proved}{01BL}{\SeriesRef{cdp:eq:B.6}, \SeriesRef{cdp:eq:B.5}, \SeriesRef{cdp:eq:B.7}, \SeriesRef{cdp:st:B1}}

Let $h:E\to\PP^1_t$ be the Weierstrass surface defined on the finite chart by
\begin{equation}
 y^2z=x^3+(t^4+1)xz^2+t^6z^3,                          \hypertarget{tag-008I}{}\label{cdp:eq:B.3}
\end{equation}
and at $s=1/t$ by $x_\infty=s^2x$, $y_\infty=s^3y$, $z_\infty=z$. Let $O=[0:1:0]$ and $P=(0,t^3,1)$ be its two sections and put
\[
 M_E=\OO_E(P-O),\quad Y=E\times\PP^1,\quad
 g=h\circ\operatorname{pr}_E,\quad
 L=K_Y\otimes\operatorname{pr}_E^*M_E^{-1},\quad
 \mathcal E=T_Y\otimes L.
\]
The surface $E$ is smooth and projective. Its singular fibres are twelve irreducible nodal cubics, at the distinct roots $t_i$ of
\[
 D(t)=31t^{12}+12t^8+12t^4+4.
\]
For every such fibre $S_i=g^{-1}(t_i)$, with $A=(K_Y\otimes L^{-1})|_{S_i}$, we have
\[
H^0(S_i,N^*_{S_i/Y}\otimes A)=0,
\qquad H^0(S_i,\widetilde\Omega^1_{S_i}\otimes A)=0,
\]
\[
 (h^0,h^1,h^2)(S_i,\mathcal E|_{S_i})=(0,1,1).
\]
As coherent sheaves on the base, with noncanonical identifications of their one-dimensional fibres,
\begin{equation}
 R^qg_*\mathcal E=0\ (q\ne2),\qquad
 R^2g_*\mathcal E=\bigoplus_{i=1}^{12}\C_{t_i}.        \hypertarget{tag-008J}{}\label{cdp:eq:B.4}
\end{equation}
In particular $(h^0,h^1,h^2,h^3)(Y,\mathcal E)=(0,0,12,0)$.
\end{proposition}



\begin{proof}

\begin{enumerate}[label=\arabic*.,leftmargin=*]
\item The transition rule realizes \eqref{cdp:eq:B.3} as a cubic in a projective plane bundle over $\PP^1$, with Weierstrass coefficients in $\OO(4)$ and $\OO(6)$; hence its total space is projective. Its discriminant is $-16D(t)$. Exact polynomial division gives
\[
 \gcd(D,D')=1,\qquad \gcd(D,t^4+1)=1.
\]
The leading coefficient $31$ is its nonzero value in the chart at infinity. Each discriminant root therefore has nonzero $c_4=-48(t^4+1)$ and is a nodal cubic. A simple discriminant zero is transverse to the nodal discriminant and makes the surface smooth at that node. Alternatively the Gr\"obner basis of the affine total-space singularity ideal $(F,F_x,F_y,F_t)$ is $(1)$. At the zero section smoothness follows from $F_z=1$ in the chart $y=1$. At infinity the coefficients are $a=1+s^4,b=1$, with discriminant $-16\cdot31$ at $s=0$.

\item The sections $P,O$ are disjoint and meet the smooth locus of every fibre. For $P$, a singular point would require both $a(t)=0$ and $t=0$, which is impossible. Each fibre is an integral arithmetic-genus-one curve with trivial dualizing sheaf. The line bundle $M_E$ has degree zero on every fibre and is nontrivial there. On a smooth fibre, a trivialization would give a meromorphic function with one simple pole and one simple zero, hence a degree-one map to $\PP^1$, contradicting genus one. On a nodal fibre the normalization calculation in \eqref{cdp:eq:B.2} says that two distinct smooth points determine a nontrivial degree-zero line bundle. Thus $H^q(E_t,M_E^{\pm1}|_{E_t})=0$ for both $q=0,1$: a nonzero section of a degree-zero line bundle on an integral curve is a trivialization, and Serre duality gives the assertion in degree one. Grauert's theorem, applied to the proper flat map $h$, yields
\begin{equation}
 Rh_*M_E=Rh_*M_E^{-1}=0.                               \hypertarget{tag-008K}{}\label{cdp:eq:B.5}
\end{equation}

\item Let $Z\subset E$ be the reduced set of the twelve critical points. The relative dualizing sheaf is $\omega_{E/\PP^1}=h^*\OO(1)$. Indeed the residue frame $dx/(2y)$ has this transition under $x_\infty=s^2x,y_\infty=s^3y$, and extends as a dualizing frame at a node. It follows that $K_E=h^*\OO(-1)$. There is an exact sequence of coherent sheaves on $E$
\begin{equation}
0\longrightarrow h^*\Omega^1_{\PP^1}
 \xrightarrow{dh}\Omega^1_E
 \longrightarrow\mathcal I_Z\otimes\omega_{E/\PP^1}
 \longrightarrow0.                                   \hypertarget{tag-008L}{}\label{cdp:eq:B.6}
\end{equation}
Away from $Z$ this is the ordinary relative cotangent sequence. At a node choose coordinates $vw=t$. The map from $\Omega^1_E$ to the relative dualizing line sends $dv\mapsto v\omega$, $dw\mapsto-w\omega$; its image is $(v,w)\omega$ and its kernel is generated by $w\,dv+v\,dw=dt$. This proves \eqref{cdp:eq:B.6} on every stalk, with the critical scheme reduced.

\item Tensor \eqref{cdp:eq:B.6} by $M_E^{-1}$. By \eqref{cdp:eq:B.5}, projection formula, and the sequence for the ideal of $Z$, its derived direct image has only degree one, namely
\begin{equation}
 R^1h_*(\Omega^1_E\otimes M_E^{-1})
 =h_*(\omega_{E/\PP^1}\otimes M_E^{-1}|_Z),
 \qquad h_*(\Omega^1_E\otimes M_E^{-1})=0.             \hypertarget{tag-008M}{}\label{cdp:eq:B.7}
\end{equation}
The right side is the sum of twelve reduced skyscrapers. Each summand in \eqref{cdp:eq:B.7} has length one.

\item Contraction on the smooth surface gives the canonical isomorphism $T_E\otimes K_E\simeq\Omega^1_E$. Consequently the product tangent decomposition identifies
\[
\mathcal E=(\Omega^1_E\otimes M_E^{-1})\boxtimes\OO(-2)
 \ \oplus\ (K_E\otimes M_E^{-1})\boxtimes\OO.
\]
The second summand has zero derived direct image by \eqref{cdp:eq:B.5}. The first has only the tensor product of degree-one terms: \eqref{cdp:eq:B.7} and $H^1(\PP^1,\OO(-2))=\C$. The relative coherent K\"unneth comparison gives \eqref{cdp:eq:B.4}. The Leray spectral sequence over a curve then gives the four stated global dimensions.

\item On a singular fibre $S_i=E_{t_i}\times\PP^1$, the restriction of $M_E$ is a nontrivial $A_\lambda$. \Cref{cdp:st:B1} gives the two vanishings and ambient dimensions $(1,1,0)$. Adjunction for the Cartier fibre gives $\omega_{S_i}=K_Y|_{S_i}$, because its normal line is trivial. Serre duality for this compact Gorenstein surface identifies, for $q=0,1,2$,
\[
 H^q(S_i,\mathcal E|_{S_i})^\vee
 =H^{2-q}(S_i,\Omega_Y^1|_{S_i}\otimes A).
\]
The fibre value $h^1=1$ with $R^1g_*\mathcal E=0$ is the degree-one Tor term obtained by tensoring a reduced skyscraper in degree two with the residue field.

\end{enumerate}
\end{proof}

\begin{proposition}[Topology of the Weierstrass product]\hypertarget{tag-01BM}{}\label{cdp:st:B3}
\StatementRecord{cdp:st:B3}{proved}{01BM}{\SeriesRef{cdp:st:B2}; \cite{CDP2020}}

The threefold $Y$ of \cref{cdp:st:B2} has $H^1(Y,\Z)=0$, $H^2(Y,\Z)\simeq\Z^{11}$, algebraic dimension $3$, and $e(Y)=24$. Its map $g$ is holomorphic, proper, surjective, and has connected fibres. For the line $A$ and the nodal fibre $S_i$ in \cref{cdp:st:B2}, the two conormal-sequence endpoint groups vanish and $h^0(S_i,\Omega_Y^1|_{S_i}\otimes A)=1$. These are the groups used in the fibrewise inference of \cite[Section~7, p.~692]{CDP2020}. The line $M_E$ restricts to a nontorsion line on each nodal cubic. Theorem~2.2 of that paper assumes $H^2(Y,\Z)=0$ and algebraic dimension one.
\end{proposition}


\begin{proof}
\begin{enumerate}[label=\arabic*.,leftmargin=*]
\item 
The residue calculation above gives $K_E=h^*\OO(-1)$. Relative duality gives $R^1h_*\OO_E=\OO(-1)$, since every fibre is an integral genus-one curve, and $h_*\OO_E=\OO$. Thus $q(E)=p_g(E)=0$ and $H^0(E,K_E^{\otimes m})=0$ for every positive $m$. The invariants $q$ and $P_2$ are unchanged under point blowups. Castelnuovo's criterion \cite[Theorem~2.26, p.~178]{Hassett2009Rational}, applied to a smooth minimal model of $E$, therefore makes $E$ rational. A rational smooth projective surface is obtained from a minimal rational surface by point blowups; its first integral homology vanishes and its second integral cohomology is free. Euler additivity gives $e(E)=12$, since the twelve nodal fibres each have Euler characteristic one and the smooth elliptic fibre has Euler characteristic zero. Poincar\'e duality gives $b_2(E)=10$. The integral K\"unneth isomorphism for $Y=E\times\PP^1$ gives the asserted groups and Euler number. Projectivity and dimension give $a(Y)=3$. Connectedness and surjectivity follow directly from the integral cubic fibres of $h$ and the second product factor.

\item At a nodal fibre write its cubic as $y^2=(x-r)^2(x+2r)$, where $r\ne0$. Its normalization is $x=v^2-2r$, $y=v(v^2-3r)$, and the two preimages of the node are $v=\pm\sqrt{3r}$. At $P$ one has $v_P^2=2r$, while $O$ is the point at infinity. The rational function with divisor $P-O$ on the normalization gives the gluing multiplier $\lambda=(v_P-\sqrt{3r})/(v_P+\sqrt{3r})$, up to inversion. Writing $q=v_P/\sqrt{3r}$ gives $q^2=2/3$ and $\lambda^2+10\lambda+1=0$. Both roots are real and neither is $1$ or $-1$, so $\lambda$ has infinite order. Thus $M_E$ restricts to a nontorsion line on every nodal fibre. The two endpoint vanishings and the one-dimensional ambient cotangent group follow from \cref{cdp:st:B2} and its Serre-duality identification.

\end{enumerate}
\end{proof}

\subsection{Unitary twists on reduced fibre supports}\label{subsec:II:cotangent-unitary}

\begin{proposition}[Constant twists on an elliptic ruled surface]\hypertarget{tag-01AR}{}\label{cdp:st:A.4.2}
\StatementRecord{cdp:st:A.4.2}{proved}{01AR}{the definitions and constructions stated in the proof}

For an elliptic curve $E$, the minimal ruled surface $S_0=E\times\PP^1$ and the constant family $M_t=\OO_{S_0}$ satisfy $H^0(S_0,\Omega^1_{S_0}\otimes M_t)\ne0$ for every parameter. The ruling has genus one. These data give a counterexample to the genus-at-least-two inference in \cite[Proposition~7.3 and Step~1, p.~693]{CDP2020}, where the line-bundle family is permitted to be constant.
\end{proposition}


\begin{proof}
Take $S_0=E\times\PP^1$ for an elliptic curve $E$ and take $M_t=\OO_{S_0}$ for every parameter. Every rational curve on $S_0$ maps constantly to $E$ by Riemann--Hurwitz, so it is a ruling fibre and has self-intersection zero. Thus $S_0$ has no exceptional curve of self-intersection $-1$ and is minimal. The product cotangent decomposition gives
\[
 h^q(S_0,\Omega^1_{S_0})=(1,2,1),\qquad q=0,1,2.
\]
Indeed the $E$ cotangent summand contributes $(1,1,0)$, while the $\PP^1$ cotangent summand contributes $(0,1,1)$ by the K\"unneth formula and $h^1(\PP^1,\OO(-2))=1$. A nonzero invariant form on $E$ provides the asserted sections for every parameter, whereas the ruling has genus one.

\end{proof}

\begin{proposition}[Unitary-flat differentials on every compact reduced complex surface]\hypertarget{tag-01AW}{}\label{cdp:st:A.4.7}
\StatementRecord{cdp:st:A.4.7}{proved}{01AW}{the definitions and constructions stated in the proof}

Let $S$ be a compact reduced complex surface, let $\mathbb A$ be a rank-one unitary complex local system on its underlying topological space, and let $A=\mathbb A\otimes_{\C}\OO_S$ be its associated holomorphic line bundle. There is a canonical complex-linear injection
\begin{equation}
 H^0(S,\widetilde\Omega_S^1\otimes A)
       \longrightarrow H^1(S,\mathbb A).
 \hypertarget{tag-004C}{}\label{cdp:eq:A.20}
\end{equation}
The map is constructed from continuous local primitives whose pullbacks to the normalization are holomorphic.
\end{proposition}


\begin{proof}

\begin{enumerate}[label=\arabic*.,leftmargin=*]
\item Let $\eta:\widetilde S\to S$ be normalization and choose a resolution $\pi:Z\to\widetilde S$ such that the exceptional curves have normal crossings; the components of $Z$ are compact smooth surfaces. Pullback of Kähler differentials along $h=\eta\circ\pi:Z\to S$ kills the torsion of $\Omega_S^1$, because the resulting form vanishes on the dense inverse image of $S_{\rm reg}$ and the target is locally free. Thus a section $\alpha$ on the left of \eqref{cdp:eq:A.20} gives a holomorphic one-form $\alpha_Z$ with values in the unitary-flat line $h^*A$. This map on sections is injective for the same dense-open reason.

\item Write $\nabla$ for the flat unitary connection. Holomorphicity gives $\nabla\alpha_Z$ type $(2,0)$. The Hermitian pairing identifies the tensor product of the flat line with its conjugate with the trivial local system. Stokes' theorem gives, on each connected component of $Z$,
\[
 \int_Z \nabla\alpha_Z\wedge\overline{\nabla\alpha_Z}
 =\int_Z d\bigl(\alpha_Z\wedge\overline{\nabla\alpha_Z}\bigr)=0.
\]
After multiplication by the fixed dimensional sign, the integrand is the nonnegative squared norm of a $(2,0)$-form times a positive volume form. Hence $\nabla\alpha_Z=0$.

\item Fix a point $x\in S$. Trivialize $\mathbb A$ on a sufficiently small connected neighbourhood $V$ of $x$; complex analytic spaces are locally contractible, so such a flat trivialization exists. The section $\alpha$ has a local lift to $\Omega_S^1\otimes A$, since the latter surjects onto the torsion-free quotient. Its pullback to each reduced fibre of $h$ over $x$ is zero: the morphism of that fibre to $S$ is constant, so the induced morphism on Kähler differentials is zero.

\item For each point $\widetilde x\in\eta^{-1}(x)$, a sufficiently small regular neighbourhood of the exceptional fibre $\pi^{-1}(\widetilde x)$ in the good resolution retracts onto that fibre. This neighbourhood can be constructed by plumbing the normal disc bundles of the exceptional curves at their transverse intersection points; it is a neighbourhood of the entire fibre. A sufficiently small inverse image of a neighbourhood of $\widetilde x$ lies inside it by properness. In the smooth case one instead uses a coordinate ball. Every period of $\alpha_Z$ on the regular neighbourhood is zero, because the form is closed and its restriction to the retracting fibre is zero. Integration from a fixed point therefore produces a single-valued holomorphic primitive: a smooth primitive has differential of type $(1,0)$, so its $\bar\partial$ is zero. Normality and $\pi_*\OO_Z=\OO_{\widetilde S}$ descend that primitive to a holomorphic function $b$ on a neighbourhood of $\widetilde x$. Choose its constant so that $b(\widetilde x)=0$. Perform this independently at every point over $x$ and shrink $V$ so that all its inverse image is covered by these neighbourhoods.

\item These primitives have equal values on every fibre of $\eta$ over $V$. To see this, form the finite analytic fibre product
$W=\widetilde S\times_S\widetilde S$ with projections $p_1,p_2$. On the normalization of every positive-dimensional irreducible component of $W_{\rm red}$, the differential of $p_1^*b-p_2^*b$ is zero on a dense open subset: both differentials are pullbacks of the same local lift of $\alpha$ from $S$. The isolated singularities of $\widetilde S$ can be omitted for this comparison; holomorphicity then extends the equality. Thus the difference is constant on each such component. The same conclusion is automatic on a zero-dimensional component. Shrink $V$ once more so that every component of the finite space $W|_V$ meets the fibre over $x$; components disjoint from that fibre can be removed by properness. Our normalization $b(\widetilde x)=0$ makes all these constants zero. Hence $p_1^*b=p_2^*b$ as continuous functions, proving equality on the fibres of $\eta$.

\item A finite proper surjection is a topological quotient map. The functions $b$ therefore descend to continuous functions on $V$, holomorphic after pullback to normalization, whose differential on the regular locus is $\alpha$. On overlaps, any two such primitives differ by a locally constant function: their difference is constant on each local normalized branch, and continuity identifies the constants whenever branches meet. With the flat transition functions of $\mathbb A$, these differences define a \v Cech class in $H^1(S,\mathbb A)$. Changing the chosen primitives changes it by a locally constant coboundary. This defines the canonical linear map \eqref{cdp:eq:A.20}.

\item If the resulting class is zero, adjust the local primitives by locally constant sections so that they glue. Their pullback is then a global holomorphic section of the unitary-flat line on each compact component of $Z$. Its squared norm is plurisubharmonic, so the maximum principle makes the section parallel; its covariant differential is zero. Consequently $\alpha_Z=0$, and injectivity of pullback gives $\alpha=0$ for every element of the kernel of \eqref{cdp:eq:A.20}.

\end{enumerate}
\end{proof}

\begin{proposition}[Cotangent vanishing for the transgression family]\hypertarget{tag-01AX}{}\label{cdp:st:A.4.8}
\StatementRecord{cdp:st:A.4.8}{proved}{01AX}{\SeriesRef{cdp:st:D.1.9}, \SeriesRef{cdp:st:D.1.8}, \SeriesRef{cdp:eq:A.22}, \SeriesRef{cdp:st:A.4.7}}

Let $g:Y\to\PP^1$ be a proper surjective holomorphic map with connected fibres from a smooth compact complex threefold with $b_1(Y)=b_2(Y)=0$, and suppose its general fibre is a complex two-torus. Let $L_r=\exp(2\pi\mathrm i r\beta)\in\Pic^0(Y)$ be the transgression family constructed in \ref{cdp:st:D.1.9}. For every real irrational $r$ and every reduced fibre support $S=(Y_c)_{\rm red}$,
\begin{equation}
 H^1(S,\mathbb A_{r,S})=0,
 \qquad
 H^0(S,\widetilde\Omega_S^1\otimes L_r|_S)=0,
 \hypertarget{tag-004D}{}\label{cdp:eq:A.21}
\end{equation}
where $\mathbb A_{r,S}$ is the unitary local system whose associated holomorphic line is $L_r|_S$. The second group is the torsion-free cotangent group in \cite[Equation~(11), p.~692]{CDP2020}.
\end{proposition}


\begin{proof}

\begin{enumerate}[label=\arabic*.,leftmargin=*]
\item Take the transgression section of \ref{cdp:st:D.1.9}. On a sufficiently small disc $\Delta$ about $c$, put $U=g^{-1}(\Delta)$ and $U^*=U\setminus S$. The integral section $\gamma$ of $R^1g_*\Z$ has a lift $\gamma_U\in H^1(U,\Z)$ by Leray over the contractible disc. Its image in $H^1(U,\OO_U)$ equals $\beta|_U$, by the coherent Leray sequence. Thus the restriction of the holomorphic line $L_r$ to $U$ is the line associated to the flat character
\begin{equation}
 \pi_1(U)\longrightarrow S^1,\qquad
 a\longmapsto\exp\bigl(2\pi\mathrm i r\langle\gamma_U,a\rangle\bigr).
 \hypertarget{tag-004E}{}\label{cdp:eq:A.22}
\end{equation}
Denote this local system by $\mathbb A_U$. \ref{cdp:st:D.1.8}--\ref{cdp:st:D.1.9} prove that $\gamma_U$ restricts nontrivially to every nearby smooth fibre. Irrationality of $r$ therefore makes \eqref{cdp:eq:A.22} nontrivial on such a fibre $F$.

\item A nontrivial rank-one local system on the real torus underlying $F$ has zero cohomology in every degree. Explicitly, use the standard cellular cochain complex of the product of four circles with coefficients in the character. It is the Koszul complex for the four scalars $\chi(e_j)-1$. At least one scalar is nonzero; contraction in that factor, divided by this scalar, is a contracting homotopy. Therefore $H^q(F,\mathbb A_U|_F)=0$ for every $q$. The map $U^*\to\Delta^*$ is a differentiable fibre bundle. Its local-system Leray spectral sequence has zero fibre-cohomology sheaves, hence
$H^1(U^*,\mathbb A_U|_{U^*})=0$.

\item Inclusion induces a surjection $\pi_1(U^*)\to\pi_1(U)$: a loop in the smooth real six-manifold $U$ can be perturbed off the divisor, whose real codimension is at least two, relative to a base point in $U^*$. The induced map on first cohomology with these local coefficients is injective. For an explicit verification, identify first cohomology with crossed homomorphisms modulo coboundaries. If a crossed homomorphism $a:G\to M$ becomes $a(q(h))=(\chi(q(h))-1)v$ after pullback by a surjection $q:H\to G$, surjectivity gives $a(g)=(\chi(g)-1)v$ for every $g\in G$. Thus its original class was already zero. Applied to the displayed fundamental-group map, this proves $H^1(U,\mathbb A_U)=0$.

\item The argument applies after every sufficiently small shrinking of $\Delta$, with the restrictions of the same local system. Taking the direct limit and using proper base change identifies the resulting stalk with $H^1(S,\mathbb A_U|_S)$; scheme-theoretic nilpotents do not affect this topological local system. The second vanishing in \eqref{cdp:eq:A.21} follows from \ref{cdp:st:A.4.7} applied to the restriction of \eqref{cdp:eq:A.22}. All singular values are treated by the same global parameter $r$, so the assertion is simultaneous.

\end{enumerate}
\end{proof}

\begin{proposition}[Twisted cotangent vanishing on reduced fibre supports]\hypertarget{tag-01AY}{}\label{cdp:st:A.4.9}
\StatementRecord{cdp:st:A.4.9}{proved}{01AY}{\SeriesRef{cdp:st:D.1.9}, \SeriesRef{cdp:st:C.3.1}, \SeriesRef{cdp:st:A.4.8}, \SeriesRef{cdp:st:D.1.8}, \SeriesRef{cdp:st:A.4.7}, \SeriesRef{cdp:st:A.4.2}, \SeriesRef{cdp:eq:A.24}; \cite[\S1 and the central-extension diagram in the proof of Lemma~1.1, pp.~2--3]{SS2002}; \cite[\S2.1.1 and Theorem~2.2, p.~3]{Albanese2020}; \cite[\S2.1.2 and Proposition~2.3, pp.~3--6]{Albanese2020}}

Let $Y$ be a connected smooth compact complex threefold with $H^1(Y,\Z)=H^2(Y,\Z)=0$, algebraic dimension one and holomorphic algebraic reduction $g:Y\to C$. There is a dense open subset of $\Pic(Y)$, whose complement is a proper complex analytic subset, such that
\begin{equation}
 H^0\bigl((Y_c)_{\rm red},
       \widetilde\Omega^1_{(Y_c)_{\rm red}}\otimes L|_{(Y_c)_{\rm red}}\bigr)=0
 \hypertarget{tag-004F}{}\label{cdp:eq:A.23}
\end{equation}
for every singular reduced fibre support and every $L$ in that subset. In fact the transgression twists $L_r$ of \ref{cdp:st:D.1.9} satisfy \eqref{cdp:eq:A.23} on every reduced fibre support, including smooth and reducible supports, for every real irrational $r$.
\end{proposition}


\begin{proof}

\begin{enumerate}[label=\arabic*.,leftmargin=*]
\item Choose one very general smooth fibre $F_0$. The classification \ref{cdp:st:C.3.1} places this fibre among complex tori, bielliptic surfaces, primary or secondary Kodaira surfaces, and primary or secondary Hopf surfaces or Inoue surfaces. We first prove that for each surface $F$ in this list and each nonzero $\gamma_F\in H^1(F,\Z)$ the local system
\[
 \mathbb A_{r,F}:\quad a\longmapsto
       \exp(2\pi\mathrm i r\langle\gamma_F,a\rangle)
 \quad(r\in\mathbb R\setminus\Q)
\]
satisfies $H^0(F,\mathbb A_{r,F})=H^1(F,\mathbb A_{r,F})=0$. Its character has infinite image, so the assertion in degree zero is immediate.

\item For a torus the Koszul argument in \ref{cdp:st:A.4.8} applies. A bielliptic surface has a finite étale torus cover. Pullback of the integral class is nonzero: transfer composed with pullback is multiplication by the degree. The character stays nontrivial after pullback, and transfer with local coefficients makes pullback on cohomology injective over $\C$. Thus degree-one vanishing descends from the torus cover.

\item A primary Kodaira surface is a principal elliptic bundle over an elliptic curve, as described in \cite[\S1 and the central-extension diagram in the proof of Lemma~1.1, pp.~2--3]{SS2002}. Its fundamental group has a central extension
\begin{equation}
 1\longrightarrow\Z^2\longrightarrow\Gamma\longrightarrow\Z^2
       \longrightarrow1.
 \hypertarget{tag-004G}{}\label{cdp:eq:A.24}
\end{equation}
For any nontrivial one-dimensional character $\chi$ of $\Gamma$, if its restriction to the central subgroup is nontrivial, the latter's cohomology with these coefficients is zero by the two-variable Koszul contraction. The Hochschild--Serre spectral sequence \cite[Chapter~VII, \S6, Theorem~6.3 and its cohomology formulation, p.~171]{Brown1982} then gives $H^1(\Gamma,\C_\chi)=0$. If its central restriction is trivial, $\chi$ descends to a nontrivial character of the quotient $\Z^2$. Centrality makes its action on the cohomology of the central subgroup equal to that scalar character, and the same Koszul contraction kills every term on the second page. Again $H^1(\Gamma,\C_\chi)=0$. First cohomology of a connected space with local coefficients is first cohomology of its fundamental group, so this proves the required assertion. A secondary Kodaira surface has a finite primary cover; the same transfer argument as above applies. Hopf surfaces have finite primary covers, and a primary Hopf surface is diffeomorphic to $S^1\times S^3$; the cover statement and topology are recorded in \cite[\S2.1.1 and Theorem~2.2, p.~3]{Albanese2020}. The nontrivial circle character has an invertible cellular differential, giving the required vanishing and its finite-cover descent.

\item For an Inoue surface, we compute the same two local-system groups from its mapping-torus presentation. The explicit constructions in \cite[\S2.1.2 and Proposition~2.3, pp.~3--6]{Albanese2020} identify every Inoue surface with a mapping torus $N_\varphi$ of a compact three-manifold $N$. In the $S_M$ cases $N$ is a three-torus and the eigenvalues of $\varphi^*$ on $H^1(N,\C)$ are $\alpha,\beta,\bar\beta$, with $\alpha>1$ and $|\beta|=\alpha^{-1/2}<1$, after reversing the circle orientation when necessary. In the $S^+$ and $S^-$ cases $N$ is the nontrivial circle bundle over a two-torus; the nonzero Euler class makes $H^1(N,\C)$ the pullback of $H^1(T^2,\C)$ by its Gysin sequence. The monodromy on it has eigenvalues $\alpha,\alpha^{-1}$ or $\alpha,-\alpha^{-1}$, respectively. These assertions follow as well directly from the conjugation relations displayed in that construction: on the abelian rank-three kernel the matrix is $M$, and on the quotient of the Heisenberg kernel by its centre the matrix is $N\in\mathrm{GL}_2(\Z)$ with precisely those eigenvalues.

\item None of these degree-one eigenvalues has modulus one. In particular there are no invariant degree-one classes on $N$. The ordinary mapping-torus sequence therefore identifies $H^1(F,\Z)$ with the copy of $\Z$ from the circle, since the invariant subgroup of the torsion-free group $H^1(N,\Z)$ is zero. Thus $\mathbb A_{r,F}$ is pulled back from the circle, with monodromy $\lambda\in S^1$ of infinite order. The local-coefficient Leray sequence for $N\to N_\varphi\to S^1$ gives
\[
\begin{split}
0\longrightarrow
 \coker(\lambda-1:H^0(N,\C)\to H^0(N,\C))
 \longrightarrow H^1(F,\mathbb A_{r,F})\\
 \longrightarrow
 \ker(\lambda\varphi^*-1:H^1(N,\C)\to H^1(N,\C))
 \longrightarrow0.
\end{split}
\]
The first term is zero because $\lambda\ne1$, and the last term is zero by the eigenvalue moduli. Using the reciprocal convention for the circle local system replaces $\lambda$ by $\lambda^{-1}$ and gives the same conclusion. 

\item Let $F_1$ be any smooth fibre. The smooth locus in the connected base curve is the complement of finitely many points and is connected. A path from the value of $F_0$ to that of $F_1$, together with proper smooth local triviality, gives a diffeomorphism $d:F_0\to F_1$. For any nonzero $\gamma_1\in H^1(F_1,\Z)$, the class $d^*\gamma_1$ is nonzero and the character formula gives the local-system isomorphism
\[
 d^*\mathbb A_{r,F_1,\gamma_1}
       \simeq\mathbb A_{r,F_0,d^*\gamma_1}.
\]
The induced cohomology isomorphism identifies degrees zero and one on $F_1$ with the groups just proved zero on $F_0$. This proves both vanishings on every smooth fibre. 

\item Now repeat the neighbourhood argument of \ref{cdp:st:A.4.8} with these two fibre-cohomology vanishings. \ref{cdp:st:D.1.8}--\ref{cdp:st:D.1.9} supply the same nonzero integral class on every nearby smooth fibre and the unitary local system on every fibre neighbourhood. The only terms that can contribute to first cohomology of the punctured neighbourhood are its degree-zero and degree-one fibre cohomology, just proved zero. The fundamental-group surjection and proper base change therefore give $H^1((Y_c)_{\rm red},\mathbb A_{r,c})=0$ for every $c$. Applying \ref{cdp:st:A.4.7} proves \eqref{cdp:eq:A.23} for the single global twist $L_r$, simultaneously on all reduced supports.

\item Finally, the exponential sequence and the integral vanishings identify $\Pic(Y)$ with the complex vector space $H^1(Y,\OO_Y)$. Choose holomorphic Cech cocycles representing a basis; exponentiating their linear combinations gives a holomorphic family of line bundles on $Y\times H^1(Y,\OO_Y)$ whose fibre is the indicated line bundle. There are finitely many singular fibre supports. For each of them, tensor its fixed coherent torsion-free cotangent sheaf by this family and apply proper semicontinuity: the locus where its space of sections is nonzero is a closed complex analytic subset of the parameter space. Each is proper because it excludes $L_r$. Their finite union is proper and has empty interior, giving the asserted dense open set. 

\end{enumerate}
\end{proof}

\subsection{Multiple and reducible fibres}\label{subsec:II:cotangent-multiple}

\begin{proposition}[All thickenings and their lifting maps]\hypertarget{tag-01AL}{}\label{cdp:st:A.3.1}
\StatementRecord{cdp:st:A.3.1}{proved}{01AL}{\cite[Tags 0B3P and 0C4T, Divisors chapter p.~26; Tag 01DZ, Cohomology chapter pp.~2--3]{StacksProject}}

Let $S$ be any effective Cartier divisor in a complex manifold, let $E$ be any locally free sheaf on that manifold, and let $E_k=E|_{kS}$. For each $k\ge2$ there is a canonical exact sequence
\begin{equation}
0\longrightarrow (E|_S)\otimes(N^*)^{k-1}
 \longrightarrow E_k\longrightarrow E_{k-1}\longrightarrow0.
\hypertarget{tag-003Z}{}\label{cdp:eq:A.10}
\end{equation}
Write $G_{k-1}=(E|_S)\otimes(N^*)^{k-1}$ and let
$\epsilon_k:H^0((k-1)S,E_{k-1})\to H^1(S,G_{k-1})$
be the connecting map. Then
\begin{equation}
0\longrightarrow H^0(S,G_{k-1})\longrightarrow H^0(kS,E_k)
 \longrightarrow\ker\epsilon_k\longrightarrow0.
\hypertarget{tag-0040}{}\label{cdp:eq:A.11}
\end{equation}
In particular $H^0(kS,E_k)=0$ exactly when $H^0(S,G_{k-1})=0$ and $\epsilon_k$ is injective. Starting with $H^0(S,E|_S)$, equations \eqref{cdp:eq:A.10}--\eqref{cdp:eq:A.11} compute the group on $mS$ through all its extension maps. The vanishing of every $H^0(S,G_j)$, $0\le j<m$, is a sufficient criterion. These statements apply to $E=\Omega_Y^1\otimes B$ for every line bundle $B$, and include smooth, normal-crossings, and all other reduced supports.
\end{proposition}


\begin{proof}
With $I=\OO_Y(-S)$, the sequence is the tensor product of $E$ with
$0\to I^{k-1}/I^k\to\OO_Y/I^k\to\OO_Y/I^{k-1}\to0$.
The Cartier property identifies $I^{k-1}/I^k$ with $(N^*)^{k-1}$ on $S$; local freeness makes the tensor product exact. Taking global sections proves \eqref{cdp:eq:A.11} and its equivalence. Induction on $k$ proves the sufficient criterion.
\end{proof}

\begin{proposition}[Sections on a multiple abelian fibre]\hypertarget{tag-01AM}{}\label{cdp:st:A.3.2}
\StatementRecord{cdp:st:A.3.2}{proved}{01AM}{the definitions and constructions stated in the proof}

There is a compact smooth complex threefold $g:Y\to\PP^1$ with connected general fibre an abelian surface and a smooth fibre support $S$ of multiplicity $m$, for every $m\ge3$, and a line bundle $B$ on $Y$, such that
\begin{equation}
H^0(S,(\Omega_Y^1\otimes B)|_S)=0,
\qquad h^0(mS,(\Omega_Y^1\otimes B)|_{mS})=3.
\hypertarget{tag-0041}{}\label{cdp:eq:A.12}
\end{equation}
Its conormal has exact order $m$.
\end{proposition}


\begin{proof}

\begin{enumerate}[label=\arabic*.,leftmargin=*]
\item Take an abelian surface $F$ and a translation $a:F\to F$ of exact order $m$. Let $\zeta$ be a primitive $m$th root of unity. The action generated by $(p,z)\mapsto(a(p),\zeta z)$ on $F\times\PP^1$ is free, so the quotient $Y$ is a smooth compact threefold and the quotient map is finite \'etale. The invariant map $t=z^m$ gives $g:Y\to\PP^1$. Away from $0$ and $\infty$ its fibre is isomorphic to $F$: the $m$ lifts of a base point are permuted transitively and freely. At $0$ the fibre is $mS$, where $S=F/\langle a\rangle$. Set $B=\OO_Y(-S)$ and $E=\Omega_Y^1\otimes B$.

\item The connected cyclic cover $F\to S$ identifies $N^*=B|_S$ with the character line of weight one. Its order is exactly $m$: otherwise a nontrivial proper character power would have a nonzero invariant trivialization on the connected compact cover, whose holomorphic functions are constant. Pullback to $F$ splits the ambient cotangent bundle into two invariant tangent forms and the normal form $dz$ of weight one. Consequently
\[
 (\Omega_Y^1\otimes B)|_S\simeq (N^*)^{\oplus2}\oplus(N^*)^2.
\]
The nontrivial flat lines on the abelian surface $S$ have no holomorphic sections. Their pullbacks to $F$ are trivial; a section would be a constant function transforming by a nontrivial character, which is zero. Since $m\ge3$, both characters in the preceding splitting are nontrivial, proving the first equality.

\item The inverse image of $mS$ is the ringed space with structure sheaf $\OO_F\otimes\C[z]/(z^m)$. Write $e$ for the local generator of the pulled-back line $B$; it has weight one. All global sections on this inverse image are linear combinations of $z^k\alpha_1e,z^k\alpha_2e,z^kdz\,e$ for $0\le k<m$, where $\alpha_1,\alpha_2$ are a basis of invariant one-forms on $F$. The first two have weight $k+1$ and the third weight $k+2$. The invariant monomials are exactly
$z^{m-1}\alpha_1e,z^{m-1}\alpha_2e,z^{m-2}dz\,e$.
Taking invariants is exact over $\C$, by averaging over the finite group. These three monomials form a basis. Translations act identically on $H^q(F,\OO_F)$, whose dimensions are $\binom{2}{q}$. The same calculation gives all fibre groups: $h^q(S,E|_S)=(0,0,0)$ and $h^q(mS,E|_{mS})=(3,6,3)$ for $q=0,1,2$. On the compact threefold the two tangent summands pull back to $\OO_F\boxtimes\OO_{\PP^1}(-1)$ and have no cohomology. The normal summand is $\OO_F\boxtimes\Omega^1_{\PP^1}(-\{0\})$. Its first cohomology on $\PP^1$ has the \v Cech basis $dz/z,dz$, of weights zero and one, respectively. Indeed sections near zero are multiples of $z\,dz$, and sections near infinity are multiples of $z^{-2}dz$, leaving exactly those two Laurent classes. Its invariant part is one-dimensional. K\"unneth and finite-cover descent therefore give $h^q(Y,E)=(0,1,2,1)$ for $q=0,1,2,3$.


For this example the complete dimensions on $kS$, $1\le k\le m$, are $0$ for $k\le m-2$, $1$ for $k=m-1$, and $3$ for $k=m$; the same character calculation proves this. More generally the character count with $B=\OO_Y(-bS)$ is
\begin{equation}
\begin{aligned}
h^0(kS,(\Omega_Y^1\otimes B)|_{kS})
&=2\#\{0\le j<k:j+b\equiv0\pmod m\}\\
&\quad+\#\{0\le j<k:j+b+1\equiv0\pmod m\}.
\end{aligned}
\hypertarget{tag-0042}{}\label{cdp:eq:A.13}
\end{equation}

\end{enumerate}
\end{proof}

\Needspace{5\baselineskip}
\begin{proposition}[Finite Picard translations for multiple fibres]\hypertarget{tag-01AN}{}\label{cdp:st:A.3.3}
\StatementRecord{cdp:st:A.3.3}{proved}{01AN}{\SeriesRef{cdp:st:A.3.1}}

For a fixed multiple fibre $mS$, suppose that the set of line bundles $B$ for which $H^0(S,(\Omega_Y^1\otimes B)|_S)=0$ contains a dense open subset $U$ of the relevant Picard component. The intersection
\begin{equation}
\bigcap_{j=0}^{m-1}\{B:B\otimes\OO_Y(-jS)\in U\}
\hypertarget{tag-0043}{}\label{cdp:eq:A.14}
\end{equation}
is a dense open subset on which the full-fibre ambient group vanishes, whenever the displayed translations preserve that component. The same assertion for a finite union of components follows by using the appropriate dense open subset in each translated component.
\end{proposition}


\begin{proof}
Tensoring by the fixed line bundle $\OO_Y(-jS)$ is a biholomorphic translation between Picard components. Its inverse image of a dense open set is dense and open. The finite intersection retains these properties. Its defining condition is exactly the vanishing of every graded group in \ref{cdp:st:A.3.1}, which proves the result.
\end{proof}

\begin{remark}[Thickening and the conormal criterion]\hypertarget{tag-P221}{}\label{rem:II:thickening-endpoint}
\StatementRecord{rem:II:thickening-endpoint}{proved}{P221}{}
On a smooth support the conormal filtration in \cref{cdp:st:A.1.1} has $K=N^*$ and $Q=\Omega_S^1$. Thus the vanishing of the two end groups gives ambient vanishing, and the finite translations in \cref{cdp:st:A.3.3} give the thickening step of \cite[Remark~7.2, p.~693]{CDP2020}. On a singular support the same section sequence contains $H^0(S,K\otimes A)$, with $K=N^*\otimes\mathcal J^{-1}$ by \cref{cdp:st:A.2.1}; this is the additional kernel group required in the application to \cite[Remark~7.4, p.~695]{CDP2020}. In \cref{cdp:st:A.3.2}, the successive invariant characters give zero sections on $S$ and three on $mS$; the finite-translation criterion requires vanishing at every displayed translated twist.
\end{remark}

\begin{proposition}[Ambient cotangent vanishing on smooth multiple fibres]\hypertarget{tag-01AO}{}\label{cdp:st:A.3.4}
\StatementRecord{cdp:st:A.3.4}{proved}{01AO}{\SeriesRef{cdp:st:D.1.9}, \SeriesRef{cdp:st:A.4.9}, \SeriesRef{cdp:st:A.4.7}, \SeriesRef{cdp:st:A.3.1}}

Under the hypotheses of \cite[Proposition~7.3, p.~693]{CDP2020}, let $mS$, with $m\ge2$, be a full fibre whose reduced support $S$ is smooth. For every real irrational transgression parameter $r$ of \ref{cdp:st:D.1.9},
\begin{equation}
 H^0\bigl(mS,(\Omega_Y^1\otimes L_r)|_{mS}\bigr)=0.
 \hypertarget{tag-0045}{}\label{cdp:eq:A.15a}
\end{equation}
The same simultaneous vanishing on all smooth multiple fibres $(m\ge2)$ holds on the complement of a proper complex analytic subset of $\Pic(Y)$. 
\end{proposition}


\begin{proof}

\begin{enumerate}[label=\arabic*.,leftmargin=*]
\item Choose a disc $\Delta$ about the fibre value and write $U=g^{-1}(\Delta)$. Form the base change $t=s^m$ and normalize it, obtaining $p:U'\to U$ and $g':U'\to\Delta'$. This map $p$ is finite étale of degree $m$. Indeed, near a point of the smooth support write $g=u z^m$ with $u$ a holomorphic unit. The normalization is given by
\[
 w^m=u,\qquad s=wz.
\]
The derivative $m w^{m-1}$ is invertible, so these equations give an unramified cover of the original smooth chart. They cover the normalized base change, since they coincide away from $z=0$ and are normal. The local constructions glue by uniqueness of normalization. The divisor $S'=(g')^{-1}(0)$ is smooth and reduced, $g'$ is a submersion near it, and $p^{-1}(mS)=mS'$ as Cartier divisors.

\item Pull back the unitary local system \eqref{cdp:eq:A.22} on $U$ to $U'$; its associated holomorphic line is $A'=p^*L_r$. Every nearby smooth fibre of $g'$ maps isomorphically onto the corresponding nearby fibre of $g$. Its character is therefore the irrational exponential of a nonzero integral class. The smooth-fibre calculation in \ref{cdp:st:A.4.9} gives zero cohomology in degrees zero and one for that local system. Proper smooth local triviality at $S'$ transports the same cohomology statement to $S'$. If a component presentation of the cover is used, the argument applies on each component and the groups add.

\item A holomorphic section of a unitary-flat line on a compact smooth complex manifold is parallel: its squared norm is plurisubharmonic and hence constant by the maximum principle. Thus $H^0(S',A')=0$ follows from the zero degree-zero local-system cohomology. \ref{cdp:st:A.4.7} gives $H^0(S',\Omega_{S'}^1\otimes A')=0$ from degree-one vanishing. The conormal of the reduced full fibre $S'$ is trivial, with global generator $dg'$. Its conormal sequence, tensored by $A'$, now gives
$H^0(S',(\Omega_{U'}^1\otimes A')|_{S'})=0$.

\item All the graded factors for $mS'$ in \ref{cdp:st:A.3.1} equal this same restricted bundle, because the conormal is trivial. Induction in \eqref{cdp:eq:A.10} proves $H^0(mS',(\Omega_{U'}^1\otimes A')|_{mS'})=0$. Since $p$ is étale, $p^*\Omega_U^1=\Omega_{U'}^1$, and finite faithfully flat pullback injects sections on $mS$ into sections on $mS'$. One can verify that injection locally by tensoring the inclusion of a module into its faithfully flat scalar extension, or by the trace divided by $m$. This proves \eqref{cdp:eq:A.15a}.

\item There are finitely many non-submersion values: the critical locus is a compact analytic set with finitely many irreducible components, and $dg=0$ makes $g$ constant on the smooth locus, hence on the closure, of each component. Fix the root cover above each smooth multiple fibre and write its smooth central surface as $S'_c$. The exponential sequence identifies $\Pic(Y)$ with the complex vector space $H^1(Y,\OO_Y)$. Exponentiating linear combinations of a basis of holomorphic Cech cocycles gives the Picard family. On $S'_c$ consider the holomorphic family of bundles
\[
 E_B=(\Omega^1_{U'_c}\otimes p_c^*B)|_{S'_c}.
\]
We give the analytic nonvanishing locus by an explicit finite kernel map. Near a parameter $B_0$, smooth gauges obtained from a partition of unity and holomorphic parameter-dependent logarithms of the transition ratios identify the underlying smooth bundles. Their Dolbeault operators have the form $D_b=D_0+A_b$, where $A_b$ has order zero, depends holomorphically on $b$, vanishes at $B_0$, and satisfies
\[
 D_0A_b+A_bD_0+A_b^2=0.
\]
Choose Hermitian metrics on $S'_c$ and the central bundle. Let $i$ include the finite-dimensional harmonic spaces, let $\pi$ be harmonic projection, let $G$ be the Green operator for the Dolbeault Laplacian, and put $h=D_0^*G$. The classical Green and Hodge identities give
\[
 D_0h+hD_0=1-i\pi,\qquad \pi i=1,
 \qquad \pi h=0,\qquad hi=0;
\]
see \cite[Part~I, Theorem~3.15, author-PDF body p.~15, and Theorems~7.1--7.2, body p.~32]{DemaillyHodgeAuthor2002}. These results apply to compact Hermitian manifolds. On fixed sufficiently high Sobolev spaces, shrink the parameter neighbourhood so that the Neumann inverses $T_b^q=(1+hA_b)^{-1}$ exist in the degrees and adjacent regularities used below. Define the finite holomorphic matrix
\[
 M_b=\pi^1 A_b T_b^0 i^0:
 \mathcal H^0(E_{B_0})\longrightarrow\mathcal H^1(E_{B_0}).
\]
The two inverse kernel maps are $v\mapsto T_b^0i^0v$ and $u\mapsto\pi^0u$. Indeed, $D_bu=0$ implies
\[
 u-i\pi u=hD_0u=-hA_bu,
 \qquad u=T_b^0i^0\pi^0u.
\]
Conversely, for $u=T_b^0i^0v$, the contraction and integrability identities give
\[
 \begin{aligned}
 D_bu&=i\pi A_bu+hD_0A_bu
       =i\pi A_bu-hA_bD_bu,\\
 D_bu&=T_b^1i^1M_bv,\qquad
 \pi^0T_b^0i^0v=v.
 \end{aligned}
\]
Thus $\ker M_b$ is isomorphic by these maps to the section space of $E_B$; elliptic regularity makes every resulting Sobolev solution smooth. The nonzero-section locus is the zero set of the maximal minors expressing failure of full column rank of $M_b$, with the zero-domain and insufficient-row cases interpreted directly. It is therefore a closed complex analytic subset $Z_c$ of the parameter space.  Step~3 shows that $Z_c$ excludes every real irrational $L_r$. Outside $Z_c$, all the graded factors on $mS'_c$ vanish, and Step~4 gives the full-fibre vanishing after descent. The finite union of the $Z_c$ is proper because the same $L_r$ lies outside every member. 

\end{enumerate}
\end{proof}

\begin{proposition}[Every reducible full fibre has generic ambient vanishing]\hypertarget{tag-01AP}{}\label{cdp:st:A.3.5}
\StatementRecord{cdp:st:A.3.5}{proved}{01AP}{the definitions and constructions stated in the proof}

Let $Y$ be a smooth compact complex threefold with $H^2(Y,\Q)=0$, and let $F=\sum_i m_iS_i$ be an effective Cartier divisor with connected reducible support. Every connected component of $\Pic(Y)$ contains a dense open subset, with proper closed analytic complement, on which
\begin{equation}
 H^0(F,(\Omega_Y^1\otimes L)|_F)=0.
 \hypertarget{tag-0046}{}\label{cdp:eq:A.15b}
\end{equation}
The multiplicities may be unequal. For a proper connected fibration over a compact curve, the vanishing is simultaneous on all reducible full fibres after intersecting finitely many such open sets. Consequently a full fibre with ambient sections for every Picard twist has irreducible support.
\end{proposition}


\begin{proof}

\begin{enumerate}[label=\arabic*.,leftmargin=*]
\item Every component $S_i$ meets a different component $S_j$, because the support is connected and reducible. The canonical section of $\OO_Y(S_j)$ restricts to a nonzero section on $S_i$. It is a nonzerodivisor in every local ring of $S_i$: otherwise one local branch of $S_i$ would be contained in $S_j$, and analytic continuation on the global irreducible component would give $S_i=S_j$. Its zero scheme is therefore a nonzero effective Cartier divisor $D$ on $S_i$.

\item For any fixed locally free sheaf $E$ on the reduced irreducible compact surface $S_i$, the injections
\[
 H^0(S_i,E(-nD))\hookrightarrow H^0(S_i,E)
\]
form a decreasing sequence of finite-dimensional subspaces. Their intersection is zero. At a point of $D$, a section in the intersection has its germ in every power of the Cartier ideal; the Krull intersection theorem for the Noetherian analytic local ring makes this germ zero. The identity theorem on the reduced irreducible surface, applied in local trivializations of $E$, then makes the global section zero. A descending sequence of finite-dimensional subspaces stabilizes, so $H^0(S_i,E(-nD))=0$ for all sufficiently large $n$.

\item Fix a Picard component and a representative $L_0$ in it. The integral class $c_1(\OO_Y(S_j))$ is torsion because $H^2(Y,\Q)=0$; choose a positive integer $q$ killing it. Then $\OO_Y(-qnS_j)$ belongs to $\Pic^0(Y)$ for every $n$. Apply the preceding argument to $E=(\Omega_Y^1\otimes L_0)|_{S_i}$. It supplies, in the chosen Picard component, a twist $L_0\otimes\OO_Y(-qnS_j)$ with zero ambient group on $S_i$. Local holomorphic Picard families and proper semicontinuity show that the nonvanishing locus on that component is a proper closed analytic subset. Its complement is dense and open. The same argument works in every Picard component.

\item Choose an ordered list $A_1,\ldots,A_N$ of the prime divisors $S_i$, with $S_i$ repeated $m_i$ times, and put $F_0=0$, $F_k=A_1+\cdots+A_k$. The exact Cartier sequences, tensored with the locally free ambient bundle, are
\[
\begin{split}
0\longrightarrow
 (\Omega_Y^1\otimes L\otimes\OO_Y(-F_{k-1}))|_{A_k}
 \longrightarrow (\Omega_Y^1\otimes L)|_{F_k}\\
 \longrightarrow(\Omega_Y^1\otimes L)|_{F_{k-1}}
 \longrightarrow0.
\end{split}
\]
This follows locally from multiplication by an equation of $F_{k-1}$, and remains exact when a component occurs repeatedly. For each $k$ take the dense open set just proved on $A_k$ and translate it by the fixed global line $\OO_Y(-F_{k-1})$, using the appropriate Picard component. Their finite intersection is dense open. Every displayed first term has no sections there, so induction from $F_0=0$ proves \eqref{cdp:eq:A.15b} on the full scheme $F_N=F$.

\item A proper holomorphic map from a smooth compact manifold to a smooth curve has finitely many critical values, since its proper analytic critical-value set does not contain a general point. Its smooth connected fibres are irreducible. Thus there are finitely many reducible fibres, and a further finite intersection gives the asserted simultaneous vanishing.

\end{enumerate}
\end{proof}

\subsection{Irreducible normal-crossings fibres}\label{subsec:II:cotangent-normal-crossings}

\begin{proposition}[Cotangent groups on the rational normalization fibre]\hypertarget{tag-01AQ}{}\label{cdp:st:A.4.1}
\StatementRecord{cdp:st:A.4.1}{proved}{01AQ}{\SeriesRef{cdp:st:A.1.2}, \SeriesRef{cdp:st:A.1.1}}

For the full reduced normal-crossings fibre $W$ in the standard $(3,4,\infty)$ six-sphere member $X$, with normalization $\eta:\operatorname{dP}_6\to W$, and every $L\in\Pic(X)$, put $A=(K_X\otimes L^{-1})|_W$. Then
\begin{equation}
H^0(W,Q\otimes A)=0,
\qquad H^0(W,\Omega_X^1|_W\otimes A)\simeq\C.
\hypertarget{tag-0044}{}\label{cdp:eq:A.15}
\end{equation}
The same two conclusions hold with $A=B|_W$ for every $B\in\Pic(X)$.
\end{proposition}





\begin{proof}
For every $B\in\Pic(X)$, the triviality of $\eta^*(B|_W)$ is \cite[Theorem~10.5, proof Step~1, p.~88]{AC26}. The normalization differential and its injection on the torsion-free quotient are the maps of \cite[\S10.4, pp.~89--90]{AC26}:
\[
 Q\otimes B|_W\hookrightarrow
 \eta_*\bigl(\Omega^1_{\operatorname{dP}_6}\otimes\eta^*(B|_W)\bigr).
\]
The same source gives $H^0(\operatorname{dP}_6,\Omega^1)=0$, so the projection formula proves the quotient vanishing. This applies to $B=K_X\otimes L^{-1}$ and to arbitrary $B$.

The full reduced fibre $W$ has trivial conormal. By \cref{cdp:st:A.1.2},
$K\otimes A=\eta_*\eta^*A=\eta_*\OO_{\operatorname{dP}_6}$.
The normalization is connected and compact, so this sheaf has section space $\C$. The exact sequence \eqref{cdp:A:filtration2}, together with the quotient vanishing just proved, identifies that section space with $H^0(W,\Omega_X^1|_W\otimes A)$.
\end{proof}

\begin{proposition}[Twisted differentials on a ruled normal-crossings surface]\hypertarget{tag-01AS}{}\label{cdp:st:A.4.3}
\StatementRecord{cdp:st:A.4.3}{proved}{01AS}{the definitions and constructions stated in the proof}

Let $E$ be an elliptic curve, let $P\in\Pic(E)$, and let $T=\PP_E(\OO_E\oplus P)$ with its zero and infinity sections $D_0,D_\infty$. Glue $D_0$ to $D_\infty$ by a translation of $E$ to form the compact normal-crossings surface $S$, and denote its normalization by $\eta:T\to S$. Every line bundle $A$ on $S$ whose pullback $\eta^*A$ is topologically trivial satisfies
\begin{equation}
H^0(S,\widetilde\Omega_S^1\otimes A)=
\begin{cases}\C,&A\simeq\OO_S,\\0,&A\not\simeq\OO_S.\end{cases}
\hypertarget{tag-0047}{}\label{cdp:eq:A.16}
\end{equation}
\end{proposition}


\begin{proof}

\begin{enumerate}[label=\arabic*.,leftmargin=*]
\item For the ruling $p:T\to E$, every topologically trivial line bundle on $T$ is $p^*P_A$ for a unique $P_A\in\Pic^0(E)$. Indeed its degree on every ruling fibre is zero, so the adjunction map from the pullback of its rank-one direct image is an isomorphism, checked on each fibre by $H^0(\PP^1,\OO)=\C$ and $H^1(\PP^1,\OO)=0$. The line on $E$ has degree zero because its first Chern class pulls back to zero and $p$ has a section.

\item The relative cotangent sequence and restriction to ruling fibres give
\[
 H^0(T,\Omega_T^1\otimes p^*P_A)
 =p^*H^0(E,\Omega_E^1\otimes P_A).
\]
The relative quotient has no sections because its restriction to every fibre is $\OO_{\PP^1}(-2)$. Since $\Omega_E^1\simeq\OO_E$, the displayed group vanishes for nontrivial $P_A$. The injection into normalization differentials gives the vanishing in \eqref{cdp:eq:A.16}.

\item For $P_A=\OO_E$, choose a trivialization of $\eta^*A$. The descent datum for $A$ is a nonzero constant $\lambda$ identifying the two restrictions to $D_0,D_\infty$, because these are trivial lines on the compact connected curve $E$. Multiplying the chosen trivialization by a constant changes neither $\lambda$ nor its meaning. The local calculation for $xy=0$ identifies the image of $\widetilde\Omega_S^1$ in $\eta_*\Omega_T^1$ with pairs of forms whose tangential restrictions to the double curve agree. The unique normalization form is a scalar multiple of $p^*\alpha$, where $\alpha$ is a nonzero invariant form on $E$. Translation pulls $\alpha$ back to itself. Twisted descent is therefore the equation $\alpha=\lambda\alpha$, so a nonzero section exists exactly when $\lambda=1$. That equality also says precisely that the trivialization descends to $A\simeq\OO_S$.

\end{enumerate}
\end{proof}

\begin{proposition}[Canonical triviality on an irreducible reduced torus fibre]\hypertarget{tag-01AT}{}\label{cdp:st:A.4.4}
\StatementRecord{cdp:st:A.4.4}{proved}{01AT}{\SeriesRef{cdp:st:A.5.2}}

Let $g:Y\to C$ be a proper holomorphic map from a connected smooth compact threefold onto a smooth compact curve, with connected fibres and general fibre a complex two-torus. Every irreducible reduced full fibre $S=g^{-1}(c)$ has $\omega_S\simeq\OO_S$. When $H^2(Y,\Q)=0$, it therefore satisfies $h^q(S,\OO_S)=(1,2,1)$.
\end{proposition}


\begin{proof}

\begin{enumerate}[label=\arabic*.,leftmargin=*]
\item The coherent sheaf $g_*\omega_Y$ is torsion-free: a nonzero function on an open subset of the base cannot annihilate a holomorphic section of a line bundle on the dominating smooth total space. It has generic rank one because the smooth torus fibre has one holomorphic top form; after restricting to a base open set where Grauert base change holds, this identifies its rank. On the smooth curve it is consequently a line bundle $M$. The evaluation morphism $g^*M\to\omega_Y$ is generically nonzero, so its zero scheme is an effective Cartier divisor $D$ and $\omega_Y\simeq g^*M\otimes\OO_Y(D)$. No component of $D$ dominates the base: on a general smooth torus fibre its defining section is a nonzero holomorphic top form and is nowhere zero. Thus $D$ is vertical.

\item All vertical components meeting the specified irreducible full fibre are that fibre itself. Consequently $\OO_Y(D)|_S$ is a power of $\OO_Y(S)|_S$, and the latter is trivial because a coordinate on $C$ defines the full reduced fibre. The pullback $g^*M|_S$ is also trivial. Adjunction gives $\omega_S=(\omega_Y\otimes\OO_Y(S))|_S\simeq\OO_S$. The fibre is reduced and connected, so $h^0(\OO_S)=1$; it is a Cartier divisor, hence Cohen--Macaulay and Gorenstein, so Serre duality gives $h^2(\OO_S)=1$. Under $H^2(Y,\Q)=0$, one has $c_1(\OO_Y(S))=0$ rationally. The divisor sequence and Hirzebruch--Riemann--Roch therefore give $\chi(\OO_S)=\chi(\OO_Y)-\chi(\OO_Y(-S))=0$, yielding $h^1=2$.

\end{enumerate}
\end{proof}

\begin{proposition}[Twisted cotangent groups of irreducible normal-crossings torus fibres]\hypertarget{tag-01AU}{}\label{cdp:st:A.4.5}
\StatementRecord{cdp:st:A.4.5}{proved}{01AU}{\SeriesRef{cdp:st:A.4.4}, \SeriesRef{cdp:st:A.4.1}, \SeriesRef{cdp:st:A.4.3}, \SeriesRef{cdp:st:A.1.2}; \cite[Definition~2.1 and the following paragraph, p.~4]{SS2002}}

Let $g:Y\to C$ have the hypotheses of \ref{cdp:st:A.4.4} and let $S$ be a singular irreducible reduced full fibre with normal crossings. For every nontrivial unitary-flat line bundle $A$ on $S$,
\begin{equation}
H^0(S,\widetilde\Omega_S^1\otimes A)=0,
\qquad
H^0(S,\Omega_Y^1|_S\otimes A)
 \simeq H^0(\widetilde S,\eta^*A).
\hypertarget{tag-0048}{}\label{cdp:eq:A.17}
\end{equation}
The latter group has dimension one when $\eta^*A$ is holomorphically trivial, and zero otherwise. The normalization is exactly one of the following: a diagonal Hopf surface with two disjoint elliptic curves identified; a ruled surface over an elliptic curve with two disjoint sections identified by translation; or a rational surface with an anticanonical hexagon as in Schr\"oer--Siebert Proposition~2.6 \cite[Proposition~2.6, p.~6]{SS2002}.
\end{proposition}


\begin{proof}

\begin{enumerate}[label=\arabic*.,leftmargin=*]
\item \ref{cdp:st:A.4.4} gives $\omega_S\simeq\OO_S$. Dualizing the locally free conormal resolution of $\Omega_S^1$ yields
\[
 \mathcal T^1_S=\mathcal Ext^1_{\OO_S}(\Omega_S^1,\OO_S)
 \simeq\OO_S/\mathcal J=\OO_{S_{\rm sing,red}}.
\]
Here the conormal is trivial because $S$ is the full reduced fibre; the local identities $\mathcal J=(x,y)$ at a double point and $\mathcal J=(xy,xz,yz)$ at a triple point show that this is the reduced singular locus. Hence $S$ is d-semistable in the exact sense of \cite[Definition~2.1 and the following paragraph, p.~4]{SS2002}. Its hypotheses are irreducibility, normal crossings, d-semistability and trivial dualizing sheaf. Propositions~2.4, 2.5 and 2.6 there \cite[Propositions~2.4--2.6, pp.~5--6; Proposition~3.1 and its proof, pp.~6--7]{SS2002} give the three normalization cases. In the ruled case their Proposition~4.1 \cite[Proposition~4.1, p.~13]{SS2002} says precisely that trivial $\omega_S$ forces the section identification to be a translation. The comparison with their notation is $X_{\mathrm{SS}}=S$ and $S_{\mathrm{SS}}=\widetilde S$.

\item In the rational case $\eta^*A$ is trivial: its first Chern class is torsion because the line is unitary-flat, whereas the integral second cohomology of a smooth rational surface is torsion-free and its Picard group injects into that group by the exponential sequence. To compare its one-forms with those on $\PP^2$, take a common resolution by point blowups. At each blowup, restriction away from the exceptional curve and Hartogs extension of the coefficients over the blown-up point give the inverse of pullback on global one-forms. The resulting isomorphism identifies this group with $H^0(\PP^2,\Omega^1_{\PP^2})=0$. The injection $\widetilde\Omega_S^1\to\eta_*\Omega^1_{\widetilde S}$ therefore proves the first vanishing.

\item In the elliptic ruled case the proof of \ref{cdp:st:A.4.3} applies to a decomposable ruled surface of any degree: the two disjoint sections express it as $\PP_E(\OO_E\oplus P)$, with $P\in\Pic(E)$. The integral second cohomology of this ruled surface is torsion-free, by the projective-bundle formula. A unitary-flat line has torsion first Chern class, so its pullback is topologically trivial and comes from $\Pic^0(E)$. The cotangent calculation of \ref{cdp:st:A.4.3} applies to every $P\in\Pic(E)$ and gives the vanishing for every nontrivial $A$.

\item In the Hopf case, pull back to a finite primary diagonal Hopf cover. Its universal covering is $\C^2\setminus\{0\}$ with contraction $(z_1,z_2)\mapsto(a_1z_1,a_2z_2)$, where $|a_i|<1$. A holomorphic one-form valued in a unitary-flat line lifts to $f_1(z)dz_1+f_2(z)dz_2$, with each $f_i$ holomorphic on the punctured space and hence on $\C^2$ by Hartogs' theorem. For a contraction character $\rho$ of modulus one, equivariance says that the coefficient of $z_1^{k_1}z_2^{k_2}dz_i$ can be nonzero only when
\[
 a_1^{k_1}a_2^{k_2}a_i=\rho
\]
(or its inverse, according to the convention for the associated line). The modulus on the left is strictly less than one and the character has modulus one. Thus every coefficient is zero. Finite \'etale pullback is injective on sections, so the twisted one-form group on the Hopf normalization also vanishes, and again the normalization injection proves the first assertion on $S$.

\item Proposition~\ref{cdp:st:A.1.2} and the first vanishing identify the ambient group with $H^0(\widetilde S,\eta^*A)$. A section of a unitary-flat line on a compact connected complex manifold is either zero or parallel and nowhere zero: its squared norm is locally a sum of squares of holomorphic functions, hence plurisubharmonic, and the maximum principle makes it constant. It therefore trivializes the line, and in the trivial case the section space is $\C$.

\end{enumerate}
\end{proof}

\begin{proposition}[Euler numbers of irreducible normal-crossings torus fibres]\hypertarget{tag-01AV}{}\label{cdp:st:A.4.6}
\StatementRecord{cdp:st:A.4.6}{proved}{01AV}{\SeriesRef{cdp:st:A.4.5}}

For $S$ in \ref{cdp:st:A.4.5},
\begin{equation}
e(S)=\begin{cases}
2,&\widetilde S\text{ is rational},\\
0,&\widetilde S\text{ is elliptic ruled or Hopf}.
\end{cases}
\hypertarget{tag-0049}{}\label{cdp:eq:A.18}
\end{equation}
Thus an arbitrary compact torus fibration whose fibres are irreducible reduced normal-crossings surfaces or multiples of smooth surfaces of Euler number zero has
\begin{equation}
c_3(Y)=e(Y)=2r,
\hypertarget{tag-004A}{}\label{cdp:eq:A.19}
\end{equation}
where $r$ is the number of rational-normalization fibres. For $b_1(Y)=b_2(Y)=0$, one has $r\in\{0,1\}$ and $b_3(Y)=2-2r$.
\end{proposition}


\begin{proof}

\begin{enumerate}[label=\arabic*.,leftmargin=*]
\item Write $D\subset\widetilde S$ for the conductor support and $N\subset S$ for the singular curve. Normalization is a homeomorphism off these sets, so additivity of the topological Euler characteristic gives $e(S)=e(\widetilde S)-e(D)+e(N)$. In the Hopf case $e(\widetilde S)=0$, $D$ is a disjoint union of two elliptic curves, and $N$ is one elliptic curve; all three Euler numbers vanish. The ruled case has the same conductor curves and $e(\widetilde S)=e(\PP^1)e(E)=0$.

\item In the rational case the normalization is a Hirzebruch surface blown up twice, so its Euler number is $4+2=6$. The conductor is a cycle of six rational curves, with six nodes, so $e(D)=12-6=6$. Its image consists of three rational curves with exactly two triple points; counting the three preimages of each triple point once gives $e(N)=6-2(3-1)=2$. Hence $e(S)=2$.

\item Away from finitely many critical values the proper holomorphic map is a differentiable fibre bundle, and a two-torus has Euler number zero. Additivity for a stratification by this open base and the exceptional points therefore gives $e(Y)=\sum_c e((Y_c)_{\rm red})$. Nilpotents do not change the underlying topological space. The assumed fibre classes and \eqref{cdp:eq:A.18} give \eqref{cdp:eq:A.19}. The equality with $c_3$ is the complex-oriented Gauss--Bonnet theorem. Poincar\'e duality for the connected oriented six-manifold gives $e(Y)=2-2b_1+2b_2-b_3$. Under the two Betti vanishings this is $2-b_3$, so $b_3=2-2r\ge0$ and $r$ is zero or one.

\end{enumerate}
\end{proof}

\section{Euler identities, normal-crossing fibres and conductor jets}\hypertarget{tag-0344}{}\label{sec:II:08}

\subsection{Topology of blowups}\label{subsec:II:blowup-topology}

\begin{lemma}[Topology of a blowup]\hypertarget{tag-P234}{}\label{lem:II:blowup-topology}
\StatementRecord{lem:II:blowup-topology}{proved}{P234}{}
Let $Y$ be a connected smooth complex manifold and $Z\subset Y$ a smooth compact centre of complex codimension $r\ge2$. Write $\pi_Z:Y'\to Y$ for its blowup, $j:E_Z\hookrightarrow Y'$ for the exceptional inclusion, $p:E_Z=\PP(N_{Z/Y})\to Z$ for the projection, and $\xi=c_1(\OO_{E_Z}(1))$. The map $\pi_Z$ induces an isomorphism on fundamental groups and the integral cohomology isomorphism \eqref{cdp:eq:C.VR.9}. For every compact analytic subset $S\subset Y$ containing $Z$, its full inverse image $S'=\pi_Z^{-1}(S)$ satisfies \eqref{cdp:eq:C.VR.4} and the rational cohomology isomorphism \eqref{cdp:eq:C.KM.S2}. These conclusions also hold on a compact smooth fibre tube with the centre in its interior and the blowup unchanged on its boundary. For a smooth divisor centre, the blowup is the identity and the formulas hold with $r=1$ and an empty extra cohomology sum.
\end{lemma}
\begin{proof}
For a divisor centre its ideal is invertible; the relative Proj of its Rees algebra is identified with $Y$ by each local generator, and these identifications glue to the blowdown isomorphism. Now take $r\ge2$. The restrictions of the blowup give
\begin{equation}
 S'\setminus\pi_Z^{-1}(Z)\simeq S\setminus Z,
 \qquad \pi_Z^{-1}(Z)=E_Z=\PP(N_{Z/Y})\xrightarrow{p}Z.
 \hypertarget{tag-015Y}{}\label{cdp:eq:C.VR.3}
 \hypertarget{tag-00FP}{}\SsixEquationAliases{\label{cdp:eq:C.EC.2}}
\end{equation}
Compatible finite triangulations and Euler additivity with compact supports on the complements give
\begin{equation}
 \begin{aligned}
 e(S')-e(S)&=e(\PP(N_{Z/Y}))-e(Z)\\
           &=(r-1)e(Z).
 \end{aligned}
 \hypertarget{tag-015Z}{}\label{cdp:eq:C.VR.4}
 \hypertarget{tag-0193}{}\SsixEquationAliases{\label{cdp:eq:C.ZC.8}}
\end{equation}
Here the projective bundle has fibre $\PP^{r-1}$ of Euler number $r$; the formula uses the full inverse image even where $Z$ meets several components of $S$.

The integral map is
\begin{equation}
 \begin{split}
 H^q(Y,\Z)\oplus\bigoplus_{i=1}^{r-1}H^{q-2i}(Z,\Z)
   &\longrightarrow H^q(Y',\Z),\\
 (\alpha,\beta_1,\ldots,\beta_{r-1})
   &\longmapsto \pi_Z^*\alpha+
       \sum_{i=1}^{r-1}j_*(\xi^{i-1}p^*\beta_i).
 \end{split}
 \hypertarget{tag-0164}{}\label{cdp:eq:C.VR.9}
\end{equation}
Blowdown induces a proper continuous bijection $Y'/E_Z\to Y/Z$ and hence a homeomorphism. Tubular neighbourhoods make both pairs cofibrations, so this identifies their relative cohomology. The projective-bundle theorem decomposes $H^q(E_Z,\Z)$ into $\xi^i p^*H^{q-2i}(Z,\Z)$ for $0\le i<r$. The self-intersection map is $j^*j_*\eta=-\xi\eta$, so the extra terms in \eqref{cdp:eq:C.VR.9} map isomorphically, with sign minus, onto the terms with $i>0$. Comparing the relative exact sequences proves surjectivity, with the remaining $i=0$ term supplied by pullback from $Y$. The proper degree-one trace makes $\pi_Z^*$ injective: under Poincar\'e duality with Borel--Moore homology it is the identity $\pi_{Z*}(\pi_Z^*\alpha\mathbin\cap[Y'])=\alpha\mathbin\cap[Y]$. Restriction to $E_Z$ proves independence of the extra summands. This establishes the integral isomorphism, including its torsion groups. On a compact tube the same argument uses relative fundamental classes for its unchanged boundary.

For the full inverse image of a compact analytic $S\supset Z$, let $\iota:S'\hookrightarrow Y'$ and $j_S:E_Z\hookrightarrow S'$ be the inclusions, and put $\pi_S=\pi_Z|_{S'}$. The support cohomology map is
\begin{equation}
 \begin{aligned}
 H^q(S,\Q)\oplus\bigoplus_{i=1}^{r-1}H^{q-2i}(Z,\Q)
   &\xrightarrow{\ \simeq\ }H^q(S',\Q),\\
 (\alpha,\beta_1,\ldots,\beta_{r-1})
   &\longmapsto\pi_S^*\alpha+
     \sum_{i=1}^{r-1}\iota^*j_*(\xi^{i-1}p^*\beta_i).
 \end{aligned}
 \hypertarget{tag-00PU}{}\label{cdp:eq:C.KM.S2}
\end{equation}
The proper map induces a continuous bijection $S'/E_Z\to S/Z$. These are compact Hausdorff spaces, so it is a homeomorphism. Compatible triangulations make $(S,Z)$ and $(S',E_Z)$ cofibrations. Their relative cohomology is consequently identified by pullback. Comparing the two relative exact sequences gives
\[
 \begin{aligned}
 \cdots\longrightarrow H^q(S,\Q)
 &\xrightarrow{(\pi_S^*,i_Z^*)}H^q(S',\Q)\oplus H^q(Z,\Q)\\
 &\xrightarrow{\ j_S^*-p^*\ }H^q(E_Z,\Q)
 \longrightarrow H^{q+1}(S,\Q)\longrightarrow\cdots,
 \end{aligned}
\]
where $i_Z:Z\hookrightarrow S$. The $i=0$ projective-bundle summand is the image of the second factor under $-p^*$. For every $i>0$, the restriction of $\iota^*j_*(\xi^{i-1}p^*\beta_i)$ to $E_Z$ is $-\xi^i p^*\beta_i$. Thus the middle map is surjective in every degree, and the displayed sequence becomes short exact in every degree.

Pullback $\pi_S^*$ is injective: if $\pi_S^*\alpha=0$, restriction to $E_Z$ gives $p^*i_Z^*\alpha=0$. The projective-bundle decomposition makes $p^*$ injective, so $i_Z^*\alpha=0$ as well; injectivity of $(\pi_S^*,i_Z^*)$ then gives $\alpha=0$. For an arbitrary class on $S'$, subtract the displayed Gysin terms to remove all its positive powers of $\xi$ on $E_Z$. The remaining restriction has the form $p^*\gamma$, and exactness makes that remaining class a pullback from $S$. Restriction to $E_Z$ also shows that a sum in \eqref{cdp:eq:C.KM.S2} can vanish only when each $\beta_i=0$; injectivity of $\pi_S^*$ then gives $\alpha=0$. This proves the stated isomorphism for the full inverse image, including centres meeting several components of $S$.

For the fundamental group, choose a tubular neighbourhood of each connected component of $Z$. Its boundary is the sphere bundle $S(N_{Z/Y})$, with simply connected fibre $S^{2r-1}$. The old neighbourhood retracts onto $Z$; its blowup retracts onto $\PP(N_{Z/Y})$, whose fibre $\PP^{r-1}$ is also simply connected. Projection identifies all three fundamental groups with $\pi_1(Z)$, and the two boundary inclusions induce those identifications. Gluing the same exterior by van Kampen therefore proves that $\pi_Z$ induces an isomorphism. The argument applies component by component and also inside a tube.
\end{proof}

\subsection{The specialization fundamental group}\label{subsec:II:specialization-fundamental-group}

\begin{lemma}[Compact analytic fibre tubes]\hypertarget{tag-P233}{}\label{lem:II:compact-analytic-tube}
\StatementRecord{lem:II:compact-analytic-tube}{proved}{P233}{\cite[p.~85, paragraph following Definition~1]{BobadillaKollar2012}; \cite[Proposition and Lemma, p.~154]{Ehresmann1952Connections}}
Let $f:X\to\Delta$ be a proper holomorphic map from a smooth complex threefold to a disc, smooth over $\Delta\setminus\{0\}$. Put $S=(f^{-1}(0))_{\rm red}$. For every sufficiently small radius $R>0$, the closed tube $U_R=f^{-1}(\overline\Delta_R)$ is a compact oriented smooth six-manifold with boundary, and the inclusion $S\hookrightarrow U_R$ is a homotopy equivalence. Its boundary is the mapping torus of the monodromy on a smooth fibre. For $0<\rho<R$, the radial collar over $\rho\le |t|\le R$ gives a deformation retraction $U_R\to U_\rho$ fixed on $U_\rho$.
\end{lemma}
\begin{proof}
Choose $R$ with $\overline\Delta_R\subset\Delta$. Properness gives compactness, and smoothness over $|t|=R$ makes that inverse image a smooth boundary. The complex orientation orients the tube and the smooth fibres. Ehresmann transport \cite[Proposition and Lemma, p.~154]{Ehresmann1952Connections} on the punctured disc identifies the boundary with the mapping torus. On the closed annulus $\rho\le|t|\le R$, lift the inward radial vector field using a smooth horizontal distribution. Properness makes its transport complete across the annulus. It identifies the annular inverse image with $f^{-1}(|t|=\rho)\times[\rho,R]$. Clamping the second coordinate to $\rho$, while fixing $U_\rho$, gives a homotopy $R_s:U_R\to U_R$ with $R_0=\operatorname{id}$ and $R_1(U_R)=U_\rho$.

Apply the neighbourhood-retraction statement of \cite[p.~85, paragraph following Definition~1]{BobadillaKollar2012} to the proper restriction over the open disc $\Delta_R$. It gives an open neighbourhood $0\in V\subset\Delta_R$ and a deformation retraction $H_W$ of $W=f^{-1}(V)$ onto $S$, with final map $r:W\to S$. Choose $\rho>0$ with $\overline\Delta_\rho\subset V$. Then $U_\rho\subset W\subset U_R$, and the composite
\[
 q=r\circ R_1:U_R\longrightarrow S
\]
is a homotopy inverse to the inclusion. Indeed, the concatenated homotopy is $R_{2s}$ for $0\le s\le\tfrac12$ and $H_W(R_1(-),2s-1)$ for $\tfrac12\le s\le1$. Its two pieces agree at the joining time; their images lie in $U_R$, and they fix $S$. Its endpoint is the inclusion composed with $q$, while $q|_S=\operatorname{id}$. The full fibre and its reduction have the same underlying space, so this comparison includes all multiplicities.
\end{proof}

\begin{proposition}[Fundamental group of an irreducible fibre]\hypertarget{tag-02Y2}{}\label{night:P4R.1}\label{stmt:P4R.1}
\StatementRecord{night:P4R.1}{proved}{02Y2}{the definitions and constructions stated in the proof}
Let $f:X\to\Delta$ be a proper holomorphic degeneration of complex
two-tori with smooth total threefold and central fibre $mS$, where
$S$ is irreducible and reduced. Choose a positive small-loop
monodromy $A$ on $\Lambda=H_1(T^4;\mathbb Z)$ and a lift $t$ of
the positive base generator, with $tvt^{-1}=A(v)$ for
$v\in\Lambda$. A positive meridian of $S$ has the form
$\mu=\lambda t^m$ for some $\lambda\in\Lambda$. The inclusion
of $S$ in a sufficiently small full tube induces an isomorphism
on fundamental groups, and the regular-fibre inclusion gives
\[
 \pi_1(S)\simeq
 (\Lambda\rtimes_A\langle t\rangle)/\langle\!\langle\mu\rangle\!\rangle,
 \qquad
 \pi_1(S)/\operatorname{im}\Lambda\simeq\mathbb Z/m,
 \qquad
 b_1(S)=\dim_{\mathbb Q}\ker(A-I).
\]
In particular $\operatorname{im}\Lambda$ is an abelian normal
subgroup with finite cyclic quotient, so $\pi_1(S)$ is virtually
abelian.
\end{proposition}
\begin{proof}
\begin{enumerate}[label=\arabic*.,leftmargin=*]
\item Let $U$ be the small closed tube of \cref{lem:II:compact-analytic-tube}, and put $U^*=U\setminus S$. The inclusion $S\hookrightarrow U$ and its homotopy inverse identify their fundamental groups. The proper smooth fibration $U^*\to\overline\Delta^*$ has fibre $T^4$. Its homotopy sequence, with $\pi_2(\overline\Delta^*)=0$, gives
\[
 \pi_1(U^*)=\Lambda\rtimes_A\langle t\rangle.
\]
The splitting chooses $t$ over the positive generator of the infinite cyclic base group; every element is uniquely $\lambda t^n$.

\item The inclusion $U^*\to U$ is surjective on fundamental groups. A loop can be perturbed off the real-codimension-two smooth stratum of $S$ and its higher-codimension singular strata. A null-homotopy disc avoids the singular strata and meets $S_{\rm reg}$ transversely in finitely many points. Deleting small discs around those intersections expresses the kernel as the normal closure of smooth-support meridians. The smooth locus of the irreducible $S$ is connected. A path in it and a trivialization of its complex normal line give an annulus of positive meridians, so all of them are conjugate.

At a smooth point of $S$, write $f=vz^m$ with $v$ a holomorphic unit. On a small coordinate neighbourhood, $z'=z\exp((\log v)/m)$ is a holomorphic normal coordinate and $f=(z')^m$. The positive circle $z'=a\exp(2\pi i\theta)$, $0\le\theta\le1$, maps to $a^m\exp(2\pi im\theta)$. Its image in the base group is $m$, so its class is $\mu=\lambda t^m$. The kernel calculation and the inclusion-induced identification $\pi_1(S)=\pi_1(U)$ give the stated normal-closure presentation, including the regular-fibre map.

\item The image of $\Lambda$ is abelian and normal. Dividing
by it leaves exactly the presentation $\langle t\mid t^m=1\rangle$,
which gives the stated cyclic quotient. Rational abelianization
of the full presentation is
\[
 \frac{(\Lambda/(A-I)\Lambda)\otimes\mathbb Q\ \oplus\ \mathbb Qt}
      {\mathbb Q(\overline\lambda+mt)}.
\]
Put $V=(\Lambda/(A-I)\Lambda)\otimes\mathbb Q$. The map
$(v,c)\mapsto v-c\overline\lambda/m$ induces an isomorphism of
this quotient with $V$, whose inverse is $v\mapsto[(v,0)]$.
Rank-nullity for $A-I$ identifies the dimension of its
cokernel with the dimension of its kernel. This proves the
formula for $b_1(S)$.
\end{enumerate}
\end{proof}

\subsection{Euler localization}\label{subsec:II:euler-localization}

\begin{proposition}[Euler localization for torus fibrations]\hypertarget{tag-01FG}{}\label{cdp:st:C2a}
\StatementRecord{cdp:st:C2a}{proved}{01FG}{the definitions and constructions stated in the proof}

Let $g:Y\to C$ be a proper surjective holomorphic map from a connected compact complex manifold to a compact Riemann surface, with connected general fibre $F$. Let $\Sigma\subset C$ contain every critical value. Then $\Sigma$ can be chosen finite, and
\begin{equation}
 e(Y)=e(C\setminus\Sigma)e(F)+\sum_{c\in\Sigma}e((Y_c)_{\rm red}). \hypertarget{tag-008N}{}\label{cdp:eq:C.1}
\end{equation}
For a complex threefold with general fibre a complex two-torus this becomes
$\int_Yc_3(T_Y)=e(Y)=\sum_{c\in\Sigma}e((Y_c)_{\rm red})$.
\end{proposition}


\begin{proof}
The critical locus is a closed analytic subset of $Y$. Properness and the proper mapping theorem make its image analytic in $C$; generic smoothness over characteristic zero makes this image proper. Thus it is finite. Over $C\setminus\Sigma$, the underlying differentiable map is a proper submersion. The finite map in the Stein factorization of $g$ has generic degree one because $F$ is connected; normality of $C$ makes it an isomorphism, so the fibres of $g$ are connected. The compact-fibre proposition and horizontal-lift lemma in \cite[p.~154]{Ehresmann1952Connections} therefore give a smooth locally trivial fibre bundle on this punctured base. Euler characteristic with compact supports is additive for the finite stratification by this open part and the fibres; triangulating the base and trivializing over its cells shows that the open part contributes $e_c(C\setminus\Sigma)e(F)$. The real dimensions of a complex curve and its punctured strata are even, so $e_c(C\setminus\Sigma)=e(C\setminus\Sigma)$. Nilpotents change no underlying topological space. This proves \eqref{cdp:eq:C.1}. A two-torus has the cohomology of the real four-torus and Euler number $(1-1)^4=0$. Finally the top Chern class of the complex tangent bundle is the Euler class for the complex orientation, giving the last identity.


\end{proof}

\subsection{Normal-crossing fibre identities}\label{subsec:II:normal-crossing-identities}


\begin{cited}[The local-to-global Euler formula with multiplicities and self-intersections]\hypertarget{tag-01C8}{}\label{cdp:st:C.2.NC.1}
\StatementRecord{cdp:st:C.2.NC.1}{proved}{01C8}{\cite[\S2.1--2.2, pp.~237--241, and \S2.4--3, pp.~242--245]{ACampo1975}}
Let $g:U\to\Delta$ be a proper holomorphic map from a smooth complex threefold to a disc, smooth over $\Delta^*$, whose central fibre is $\sum_i m_iS_i$ and whose reduced support $S$ has normal crossings. Let $\nu:\coprod_i T_i\to S$ be normalization and let $D_i\subset T_i$ be its reduced conductor divisor. For a sufficiently small nonzero $t$,
\begin{equation}
 e(U_t)=\sum_i m_i e(T_i\setminus D_i).
 \label{cdp:eq:C.NC.1}
\end{equation}
\end{cited}

\begin{proof}
The proper analytic boundary fibration is the construction and lemma of \cite[\S2.1--2.2, pp.~237--241]{ACampo1975}. In a chart where
$g=z_1^{a_1}\cdots z_r^{a_r}$, write $z_j=\rho_j u_j$ with $u_j\in S^1$. Its boundary fibre over an argument $\theta$ maps to $S$, with fibre over that $r$-branch stratum
\begin{equation}
 \{(u_1,\ldots,u_r)\in(S^1)^r:
                 u_1^{a_1}\cdots u_r^{a_r}=e^{i\theta}\}.
 \hypertarget{tag-014Z}{}\label{cdp:eq:C.UI.7}
\end{equation}
An integral change of angular basis identifies this fibre with
$\gcd(a_1,\ldots,a_r)$ copies of $(S^1)^{r-1}$.
For self-intersections, changing a local branch equation multiplies its modulus by a positive function and translates its argument; a branch permutation permutes the corresponding factors. These transition maps preserve the lifted equation of $g$ and glue the same boundary projection. Thus the local product comparison in the cited lemma applies to these charts as well.

The boundary fibre is homeomorphic to $U_t$. The constructible cohomology of its projection to $S$ is the proper comparison of \cite[\S2.4--3, pp.~242--245]{ACampo1975}. Euler integration on the strata gives zero from $r\ge2$, since their angular tori have positive dimension, and gives $m_i$ from a one-branch point of $S_i$. Normalization identifies that one-branch stratum with $T_i\setminus D_i$, proving the displayed formula.
\end{proof}

\begin{proposition}[The normalized double curves and the coherent Euler formula]\hypertarget{tag-01C9}{}\label{cdp:st:C.2.NC.2}
\StatementRecord{cdp:st:C.2.NC.2}{proved}{01C9}{\SeriesRef{cdp:eq:C.NC.6}, \SeriesRef{cdp:eq:C.NC.7}}

Let $S$ be a connected compact normal-crossings divisor in a smooth compact complex threefold $Y$ with $H^2(Y,\Q)=0$. Write $\nu:\coprod_iT_i\to S$ for normalization, $D_i\subset T_i$ for its conductor, $C_\alpha$ for the connected components of the normalization of the double locus, and $\tau$ for the number of triple points. Let $R_\alpha$ be the reduced divisor on $C_\alpha$ of preimages of triple points, counted separately for the three double-curve branches at each triple point. Then
\begin{equation}
 e(C_\alpha)=\deg R_\alpha,\qquad
 \sum_\alpha e(C_\alpha)=3\tau,
 \qquad
 \sum_i\chi(T_i,\OO_{T_i})=\frac{\tau}{2}.
 \hypertarget{tag-00S3}{}\label{cdp:eq:C.NC.4}
\end{equation}
In particular every normalized double curve has genus at most one, and $\tau$ is even. The remaining topological identities are
\begin{equation}
\begin{aligned}
 \sum_i e(D_i)&=2\sum_\alpha e(C_\alpha)-3\tau,\\
 e(S)&=\sum_i e(T_i)-\sum_\alpha e(C_\alpha)+\tau.
\end{aligned}
 \hypertarget{tag-00S4}{}\label{cdp:eq:C.NC.5}
\end{equation}
\end{proposition}


\begin{proof}

\begin{enumerate}[label=\arabic*.,leftmargin=*]
\item Adjunction and the local residue comparison give a canonical holomorphic line-bundle isomorphism
\begin{equation}
 \nu^*\omega_S|_{T_i}\simeq\omega_{T_i}(D_i),
 \qquad \omega_S=(\omega_Y\otimes\OO_Y(S))|_S.
 \hypertarget{tag-00S5}{}\label{cdp:eq:C.NC.6}
\end{equation}
For example, the residue of $dx\wedge dy\wedge dz/(xy)$ on the branch $x=0$ is $\pm dy\wedge dz/y$, and at $xyz=0$ it is $\pm dy\wedge dz/(yz)$. These expressions agree under changes of local equations by the usual determinant and unit factors of the displayed line bundles. Since both line bundles on the right in the expression for $\omega_S$ come from $Y$, its first Chern class and that of $\omega_{T_i}(D_i)$ vanish rationally.

\item Normalize every irreducible curve in every $D_i$. The resulting disjoint union $\widehat D$ maps to $\coprod_\alpha C_\alpha$ by an étale map of total degree two over each $C_\alpha$. Away from a triple point this follows from the two branches of $xy=0$. At a triple point, each of the three normalized double branches has exactly two preimages, and their maps in the remaining coordinate are local isomorphisms. The covering over a given $C_\alpha$ may be connected.

\item Adjunction for a nodal divisor on a smooth surface, pulled back to $\widehat D$, identifies the restriction of $\omega_{T_i}(D_i)$ with the canonical line of that normalized conductor curve twisted by all its preimages of nodes of $D_i$. Such nodes are exactly the preimages of triple points of $S$. This is the elementary local dualizing-form identity $dw/w$ on each branch of a node. The resulting degree is zero by \eqref{cdp:eq:C.NC.6}. Since the covering of $C_\alpha$ is étale and its marked divisor is the pullback of $R_\alpha$, it follows that
$\deg\omega_{C_\alpha}+\deg R_\alpha=0$.
Thus $e(C_\alpha)=\deg R_\alpha$. There are exactly three marks over every triple point, even when two or three belong to one $C_\alpha$.

\item For the coherent calculation, the alternating branch-resolution sequence is
\begin{equation}
 0\longrightarrow\OO_S\longrightarrow\nu_*\OO_{\coprod T_i}
 \longrightarrow\bigoplus_\alpha(j_\alpha)_*\varepsilon_\alpha
 \longrightarrow\bigoplus_{p\in S^{[3]}}\C_p
 \longrightarrow0.
 \hypertarget{tag-00S6}{}\label{cdp:eq:C.NC.7}
\end{equation}
Here $j_\alpha:C_\alpha\to S$ is the finite normalization map of that double curve, and $\varepsilon_\alpha$ is the sign line of its two-sheeted conductor covering. Locally order the branches and send $(f_i)$ to $(f_j-f_i)_{i<j}$; the next map is the alternating sum of the three values at a triple point. The sequence is exact at $xy=0$ because a pair of functions glues precisely when its restrictions agree, and at $xyz=0$ because pairwise differences satisfy exactly the single alternating relation. Equivalently, compare coefficients of convergent power series on the three coordinate planes; equality on their coordinate axes is exactly the condition for these series to come from one series modulo $xyz$. Changing a local branch ordering multiplies the appropriate difference by $-1$, producing $\varepsilon_\alpha$. At a triple point this sign convention only chooses a basis of the one-dimensional last term. Hence the sequence is globally well-defined with the indicated local systems, including self-intersections.

\item The sign line satisfies $\varepsilon_\alpha^{\otimes2}\simeq\OO_{C_\alpha}$, so it has degree zero and $\chi(C_\alpha,\varepsilon_\alpha)=1-g(C_\alpha)=e(C_\alpha)/2$. The divisor sequence on $Y$ and Hirzebruch--Riemann--Roch give $\chi(S,\OO_S)=0$, since $c_1(\OO_Y(S))=0$ rationally. Taking Euler characteristics in \eqref{cdp:eq:C.NC.7} now gives
\[
 0=\sum_i\chi(T_i,\OO_{T_i})-\frac12\sum_\alpha e(C_\alpha)+\tau.
\]

\item $\widehat D$ has Euler number $2\sum_\alpha e(C_\alpha)$. Each triple point produces three nodes in $\coprod D_i$, and identifying the two preimages of each node subtracts one. The double locus in $S$ has normalization $\coprod C_\alpha$ and three preimages over each triple point, so its Euler number is $\sum_\alpha e(C_\alpha)-2\tau$. Normalization is an isomorphism off the conductor. Additivity of Euler characteristic therefore gives the second identity in \eqref{cdp:eq:C.NC.5}.

\end{enumerate}
\end{proof}






\begin{proposition}[Uniform multiple normal-crossings fibres]\hypertarget{tag-01CA}{}\label{cdp:st:C.2.NC.3}
\StatementRecord{cdp:st:C.2.NC.3}{proved}{01CA}{the definitions and constructions stated in the proof}

Let $g:Y\to C$ be a proper holomorphic map from a connected smooth compact complex threefold with $H^2(Y,\Q)=0$ to a smooth compact curve, with connected fibres and smooth general fibre of Euler characteristic zero. A fibre $mS$ with reduced normal-crossings support has
\begin{equation}
 e(S)=\tau(S)\ge0.
 \hypertarget{tag-00S7}{}\label{cdp:eq:C.NC.8}
\end{equation}
The normalized base change $t=s^m$ is finite étale of degree $m$ over a neighbourhood of $S$ and has reduced normal-crossings central fibre $S'$. The induced cover of reduced fibres satisfies
$e(S')=m e(S)$ and $\tau(S')=m\tau(S)$.
\end{proposition}


\begin{proof}

\begin{enumerate}[label=\arabic*.,leftmargin=*]
\item The local formula \cref{cdp:st:C.2.NC.1} gives
\[
 0=e(Y_t)=m\sum_i e(T_i\setminus D_i).
\]
By \eqref{cdp:eq:C.NC.5},
\[
 e(S)-\sum_i e(T_i\setminus D_i)
 =\sum_\alpha e(C_\alpha)-2\tau=\tau.
\]

\item Locally write $t=u(z_1\cdots z_r)^m$ with $u$ nowhere zero. The normalization of $s^m=u(z_1\cdots z_r)^m$ is given by
$w^m=u$ and $s=wz_1\cdots z_r$.
The derivative of $w^m-u$ with respect to $w$ is nonzero because $w$ is a unit. Thus its map to the original chart is étale of degree $m$, and the lifted central fibre has equation $z_1\cdots z_r=0$. These local normalizations glue by the universal property of normalization; their degree and étaleness are local properties. Restriction is therefore an étale degree-$m$ covering $S'\to S$, and each triple point has exactly $m$ preimages. A lifted triangulation proves the Euler-number comparison. If the cover has several connected components, the degrees and Euler numbers add over them.

\end{enumerate}
\end{proof}






\begin{proposition}[Reducible normal-crossings fibres with unequal multiplicities]\hypertarget{tag-01CB}{}\label{cdp:st:C.2.NC.4}
\StatementRecord{cdp:st:C.2.NC.4}{proved}{01CB}{\SeriesRef{cdp:st:C.2.NC.3}, \SeriesRef{cdp:st:C.2.NC.2}; \cite[Theorem~2.2, p.~843]{Peters2021}; \cite[Theorem~3.2, p.~16]{BMS2012}}

Under the hypotheses on $g:Y\to C$ in \ref{cdp:st:C.2.NC.3}, let a fibre be $\sum_i m_iS_i$ with connected reducible normal-crossings support $S$. Its multiplicities may be unequal. Then
\begin{equation}
 \tau(S)=0,\qquad e(S)=0.
 \hypertarget{tag-00S8}{}\label{cdp:eq:C.NC.9}
\end{equation}
Every normalized conductor is a disjoint union of smooth elliptic curves, and every normalization component has Euler number zero. More generally, any singular normal-crossings fibre support under these hypotheses has either no triple point or exactly two triple points; in the latter case it is irreducible and its normalization is rational.
\end{proposition}


\begin{proof}

\begin{enumerate}[label=\arabic*.,leftmargin=*]
\item First consider a normalization component $T_i$ whose conductor $D_i$ is nonempty. If $T_i$ is Kähler and has nonnegative Kodaira dimension, choose a nonzero section of $\omega_{T_i}^{\otimes n}$ and denote its effective zero divisor by $E$. The effective divisor $E+nD_i$ has rational Chern class zero by \eqref{cdp:eq:C.NC.6}. Its integral against a Kähler form is positive because $D_i$ is nonempty, whereas its cohomology class gives integral zero. This is impossible. By \cite[Theorem~2.2, p.~843]{Peters2021}, a Kähler $T_i$ is consequently rational or a blown-up ruled surface. For a ruling $T_i\to B$ and a general ruling fibre $F$, one has $K_{T_i}\cdot F=-2$, hence $D_i\cdot F=2$. Some conductor curve therefore dominates $B$. Its normalization has genus at most one by \ref{cdp:st:C.2.NC.2}, so Riemann--Hurwitz forces $g(B)\le1$. The Kähler possibilities have
$\chi(\OO_{T_i})=1$ in the rational case and $0$ in the elliptic ruled case; their Euler numbers are nonnegative.

\item For a non-Kähler $T_i$, the Enriques--Kodaira table \cite[Theorem~2.2, p.~843]{Peters2021} leaves class VII, Kodaira and properly elliptic minimal models. Class VII has $c_1^2=-b_2$ and $c_2=b_2$, so Noether's formula gives $\chi(\OO)=0$. Kodaira surfaces have both Chern numbers zero. A minimal non-Kähler elliptic surface has topological Euler number zero by Kodaira's theorem in the precise form \cite[Theorem~3.2, p.~16]{BMS2012}; its canonical square is zero by the same classification table, so Noether's formula again gives $\chi(\OO)=0$. Blowing up preserves the holomorphic Euler characteristic and adds one to the topological Euler characteristic. Thus every non-Kähler $T_i$ has $\chi(\OO_{T_i})=0$ and $e(T_i)\ge0$.

\item A rational normalization component cannot meet a different global component of $S$. Indeed, for $j\ne i$, pull back the canonical section of $\OO_Y(S_j)$ to $T_i$. It is a nonzero section whose zeros are precisely the inverse image of $S_i\cap S_j$. Its line bundle has zero rational first Chern class. For a smooth rational surface $T_i$, $H^1(T_i,\OO)=0$ and $H^2(T_i,\Z)$ is torsion-free, by the corresponding facts for $\PP^2$ or a Hirzebruch surface and the blowup formula. The exponential sequence therefore makes that pulled-back line bundle trivial. A nonzero section of a trivial line bundle on a connected compact surface is a nonzero constant and has no zeros. Hence $S_i\cap S_j$ is empty.

\item The intersection graph of the global components of a connected reducible fibre is connected. Each component has nonempty conductor, and the preceding paragraph excludes rational normalizations. Thus every $\chi(\OO_{T_i})$ is zero, and \eqref{cdp:eq:C.NC.4} gives $\tau=0$. Each $C_\alpha$ is now elliptic and has no marked point. The conductor covering is étale, so the normalized conductor components are elliptic; there are no nodes of $D_i$, since these would be preimages of triple points. Consequently $e(D_i)=0$. Apply \cref{cdp:st:C.2.NC.1} to the restriction of $g$ over a disc about this fibre, retaining its multiplicities $m_i$. It gives $e(Y_t)=\sum_i m_i e(T_i\setminus D_i)$. Since $e(D_i)=0$, this is
\[
 0=e(Y_t)=\sum_i m_i e(T_i).
\]
All summands are nonnegative and every $m_i$ is positive, so each $e(T_i)=0$. Formula \eqref{cdp:eq:C.NC.5} gives $e(S)=0$.

\item For a singular irreducible support the conductor is nonempty and the same classification argument gives $\chi(\OO_{T_1})\in\{0,1\}$, with value one exactly in the rational case. Formula \eqref{cdp:eq:C.NC.4} then gives $\tau=0$ or $2$, respectively. Thus any triple point lies on an irreducible support with rational normalization.

\end{enumerate}
\end{proof}






\begin{proposition}[Global positivity and the all-twist normalization obstruction in the normal-crossings class]\hypertarget{tag-01CC}{}\label{cdp:st:C.2.NC.5}
\StatementRecord{cdp:st:C.2.NC.5}{proved}{01CC}{\SeriesRef{cdp:st:C.2.NC.3}, \SeriesRef{cdp:st:C.3.5}, \SeriesRef{cdp:st:C.3.1}, \SeriesRef{cdp:st:C.2.NC.4}, \SeriesRef{cdp:st:A.1.2}}

Let $g:Y\to C$ have the hypotheses of \ref{cdp:st:C.2.NC.3}, and assume every fibre reduction has normal crossings; smooth reductions are included. Then
\begin{equation}
 c_3(Y)=e(Y)=\sum_{c\in C}\tau((Y_c)_{\rm red})=2r\ge0,
 \hypertarget{tag-00S1}{}\label{cdp:eq:C.NC.10}
\end{equation}
where $r$ is the number of fibre supports with rational normalization. Every such support is irreducible, nonnormal, and has exactly two triple points. Equality $c_3(Y)=0$ holds exactly when every fibre support is free of triple points. If also $b_1(Y)=0$, then $r\in\{0,1\}$ and $b_3(Y)=2-2r$.

In the positive case there is an irreducible nonnormal reduced support $S$ with rational normalization $\eta:T\to S$ such that, for every $L\in\Pic(Y)$,
\begin{equation}
 H^2\!\left(T,K_T\otimes\eta^*((K_Y^{-1}\otimes L)|_S)\right)\simeq\C,
 \qquad
 H^0\!\left(S,\Omega_Y^1|_S\otimes(K_Y\otimes L^{-1})|_S\right)\simeq\C.
 \hypertarget{tag-00S2}{}\label{cdp:eq:C.NC.11}
\end{equation}
The scheme-theoretic fibre may be a multiple of $S$. By \ref{cdp:st:C.3.5}, all these conclusions apply to every fibration under the CDP hypotheses whose fibre reductions have normal crossings, for every smooth-fibre type in \ref{cdp:st:C.3.1}.
\end{proposition}


\begin{proof}

\begin{enumerate}[label=\arabic*.,leftmargin=*]
\item Apply Proposition~\ref{cdp:st:C2a} to this same map
$g:Y\to C$, with $F$ its smooth general fibre. The hypothesis
$e(F)=0$ makes \eqref{cdp:eq:C.1} the equality
$e(Y)=\sum_c e((Y_c)_{\rm red})$; its complex-oriented
Gauss--Bonnet comparison gives $c_3(Y)=e(Y)$.
An irreducible support has a single multiplicity, so
\ref{cdp:st:C.2.NC.3} applies; a reducible support is covered by
\ref{cdp:st:C.2.NC.4}. A smooth support is connected and irreducible,
so the uniform-multiplicity formula gives Euler number zero.
Thus every summand equals its triple-point count, and
\ref{cdp:st:C.2.NC.4} identifies the positive summands and their
value two. Poincar\'e duality gives
$e(Y)=2-2b_1(Y)+2b_2(Y)-b_3(Y)$. Under the additional hypothesis
$b_1(Y)=0$, this gives $b_3(Y)=2-2r\ge0$, hence $r\in\{0,1\}$.
Under the CDP hypotheses, \ref{cdp:st:C.3.5} gives Euler number
zero for every smooth fibre.

\item For a positive support, every line bundle on $Y$ pulls back trivially to its smooth rational normalization: its rational first Chern class is zero, and the exponential sequence identifies $\Pic(T)$ with the torsion-free group $H^2(T,\Z)$ by the first Chern class, as in \ref{cdp:st:C.2.NC.4}. In particular $\eta^*((K_Y^{-1}\otimes L)|_S)$ is trivial. Serre duality on the connected smooth compact surface $T$ identifies the first group in \eqref{cdp:eq:C.NC.11} with $H^0(T,\OO_T)^\vee\simeq\C$.

\item For the ambient group set $A=(K_Y\otimes L^{-1})|_S$ and $N^*=\OO_Y(-S)|_S$. Both $\eta^*A$ and $\eta^*N^*$ are trivial, even when the full fibre is $mS$. The differential pullback has kernel $\operatorname{Tor}(\Omega_S^1)$ and induces $\widetilde\Omega_S^1\hookrightarrow\eta_*\Omega_T^1$. Tensoring by $A$ and applying the projection formula gives
$H^0(S,\widetilde\Omega_S^1\otimes A)\hookrightarrow H^0(T,\Omega_T^1)=0$.
Put $K=\ker(\Omega_Y^1|_S\to\widetilde\Omega_S^1)$. The exact sequence in \ref{cdp:st:A.1.1} and the isomorphism $K\simeq\eta_*\eta^*N^*$ of \ref{cdp:st:A.1.2}, followed by the projection formula, give
\[
 H^0(S,\Omega_Y^1|_S\otimes A)
 =H^0(T,\eta^*(N^*\otimes A))=\C.
\]


\end{enumerate}
\end{proof}






\subsection{Normal jets along conductor curves}\label{subsec:II:normal-conductor-jets}

\begin{proposition}[The normal jets along cusp, tacnode and triple curves]\hypertarget{tag-01FJ}{}\label{cdp:st:C2d}
\StatementRecord{cdp:st:C2d}{proved}{01FJ}{\SeriesRef{cdp:eq:C.J.2}, \SeriesRef{cdp:eq:C.J.3}, \SeriesRef{cdp:st:A.2.4}}

Let $Y$ be a smooth complex threefold with $b_2(Y)=0$, let $S$ be a reduced Cartier divisor, and let $D\subset\operatorname{Sing}S$ be a smooth compact connected curve. Suppose the general plane section transverse to $D$ has one of the equations
\begin{equation}
 x^2-y^3=0,\qquad x^2-y^4=0,\qquad xy(x-y)=0.
 \hypertarget{tag-00OV}{}\label{cdp:eq:C.J.1}
\end{equation}
Then $g(D)\le1$. For $g(D)=1$, the indicated transverse analytic type is constant along all of $D$. The inverse image of $D$ in the normalization is respectively an isomorphic copy, a finite \'etale cover of degree two, or a finite \'etale cover of degree three. Each connected component $C$ of that cover has genus one and normal line of degree zero in the normalization, which is smooth near $C$. At these three transverse types, the inverse Jacobian ideal pulls back into the line allowing at most a simple pole on each normalization branch.
\end{proposition}


\begin{proof}

\begin{enumerate}[label=\arabic*.,leftmargin=*]
\item Put $N=N_{D/Y}$ and $B=\OO_Y(S)|_D$. Since every ambient line bundle has zero rational first Chern class, $\deg B=\deg K_Y|_D=0$. Adjunction for the smooth curve gives
\begin{equation}
 \det N=K_D\otimes K_Y^{-1}|_D,\qquad
 \deg N=2g(D)-2.                                      \hypertarget{tag-00OW}{}\label{cdp:eq:C.J.2}
\end{equation}
A local defining section of $S$ belongs to $I_D^2$. Its quadratic normal jet is a well-defined section of $\operatorname{Sym}^2N^*\otimes B$. In the first two cases its generic rank is one. Its saturated radical $R\subset N$ is a line subbundle on the smooth curve, and $Q=N/R$ is a line bundle. The quadratic jet gives a nonzero section of $B\otimes Q^{-2}$, so
$\deg Q\le0$.

\item In the cusp case, the cubic jet on the radical is an intrinsic nonzero section of $B\otimes R^{-3}$. To check its independence, choose coordinates $(x,y,z)$ with $D=(x=y=0)$ and $R$ the $y$ direction. A change of lift of a radical vector changes the cubic by a Hessian term with one radical argument, which is zero. Changing the local retraction to $D$ differentiates the identity that the quadratic jet vanishes on $R$, again giving zero. Thus the local cubics glue across all points, including possible zeros of the quadratic coefficient. Consequently $\deg R\le0$. Equation \eqref{cdp:eq:C.J.2} proves $g(D)\le1$.

\item In the tacnode case the radical cubic vanishes identically. Write the relevant normal terms as
\[
 a x^2+bxy^2+cy^4.
\]
Completing the square has invariant numerator
\begin{equation}
 \Delta=ac-b^2/4\in
 H^0(D,B^2\otimes Q^{-2}\otimes R^{-4}).               \hypertarget{tag-00OX}{}\label{cdp:eq:C.J.3}
\end{equation}
This numerator is holomorphic even at a zero of $a$. Under $x\mapsto x+uy^2$, the coefficients become $b+2au$ and $c+bu+au^2$, and \eqref{cdp:eq:C.J.3} is unchanged. Linear changes preserving the radical give precisely the displayed line-bundle weights. Changes of the retraction to $D$ add terms containing the radical cubic or a radical Hessian argument, so also leave this quartic numerator unchanged. These checks give its global definition. It is nonzero by the generic tacnode assumption. Hence
\[
 \deg Q\le0,\qquad \deg Q+2\deg R\le0,
 \qquad \deg N\le\tfrac12\deg Q\le0.
\]
This proves the same genus bound.

\item For the ordinary triple case the quadratic jet is identically zero. The binary cubic normal jet is a section of $\operatorname{Sym}^3N^*\otimes B$. Its discriminant is a nonzero section of
$B^4\otimes(\det N^*)^6$. The discriminant has degree four in the cubic coefficients and weight six under a linear change of the two normal coordinates; these weights also follow by writing it as the square of the product of the three pairwise determinants of the linear factors. Its nonnegative divisor degree gives $\deg N\le0$ and the genus bound again.

\item When $g(D)=1$, all the inequalities are equalities. In the cusp case both the quadratic and radical cubic sections have no zeros. In the tacnode case $a$ and $\Delta$ have no zeros. In the triple case the cubic discriminant has no zeros. The holomorphic splitting lemma with parameter $z$ completes the square in the first two cases. The remaining one-variable term is respectively a unit times $y^3$ or $y^4$; taking a local root of the unit gives the corresponding equation in \eqref{cdp:eq:C.J.1}. In the triple case its three distinct tangent lines split after an \'etale ordering cover; the holomorphic implicit function theorem gives three smooth branches meeting transversely in each normal plane. Straightening two branches to the axes and rescaling the third branch's unit gives the third equation. Thus these are constant transverse types throughout $D$. In particular no other singular curve meets $D$.

\item In the cusp normalization, $x=u^3,y=u^2$, up to the choice of which coordinate is the radical. The leading terms $y=u^2$ and $x=u^3$ identify
$P^2$ with $R$ and $P^3$ with $Q$, where
$P=N_{C/\widetilde S}$. Adapted coordinate changes preserve these
leading terms through the corresponding line transitions. Thus
$\deg P=0$. For the tacnode, the two branches have $x=\pm y^2$ and normal line the pullback of $R$, again of degree zero. Their permutation cover is \'etale because $\Delta$ is nowhere zero. For the ordinary triple, order the three tangent lines by a finite \'etale cover. Three pairwise transverse line subbundles in a rank-two bundle make it a sum $P\oplus P$: the third is a graph of an isomorphism between the first two. Since the determinant has degree zero, $\deg P=0$. Each normalization branch has this tangent-root normal line. \'{E}tale descent gives the claimed degrees and covers before ordering.

\item Compute the inverse Jacobian ideal in the parameter-preserving
normal forms, with coefficients holomorphic in the curve parameter $z$.
A biholomorphic coordinate change transforms the gradient by its
invertible Jacobian matrix; multiplication of the defining equation
by a unit multiplies the gradient by that unit modulo the equation.
These maps transport the following fractional ideals to $S$.

For the cusp set $B=\C\{u,z\}$ and
$A=\C\{u^2,u^3,z\}=\C\{z\}+u^2B$.
The Jacobian ideal has images $u^3,u^4$. For a multiplier $h$, put
$a=u^3h$. The conditions $a,ua\in A$ kill respectively the linear
and constant coefficients of $a$. Hence
$\mathcal J^{-1}=u^{-1}B$, agreeing with
Proposition~\ref{cdp:st:A.2.4} under the displayed normalization map.

For the tacnode the two maps are $x=\pm u^2,y=u$, and the image ring is
\[
 A=\{(f,g)\in B^2:f-g\in u^2B\}.
\]
Indeed, the inverse expression is
$(f+g)/2+x(f-g)/(2u^2)$. The Jacobian generators have images
$(u^2,-u^2)$ and $(u^3,u^3)$, up to nonzero scalar factors.
For $a=u^2h_1$ and $b=u^2h_2$ the membership conditions are
$a+b\in u^2B$ and $a-b\in uB$. Their constant coefficients vanish,
so each $h_i$ has at most a simple pole. Their residues are opposite:
\[
 \mathcal J^{-1}=B^2+\C\{z\}(u^{-1},-u^{-1}).
\]

For the ordinary triple point use the branches
$(0,u),(u,0),(u,u)$. The image ring consists of triples with a common
constant and with third linear coefficient equal to the sum of the
first two. All coefficients of degree at least two are independent:
for $k\ge2$ the polynomials
$y^k-xy^{k-1}$, $x^k-x^{k-1}y$, and $x^{k-1}y$ have respectively
the three individual order-$k$ branch values.
The Jacobian generators have images
$(-u^2,0,u^2)$ and $(0,u^2,-u^2)$.
For $a_i=u^2h_i$, the common-constant condition on both products forces
all three constant coefficients to vanish. The linear relations
give residues $(a,a,-a)$. Thus
\[
 \mathcal J^{-1}=B^3+\C\{z\}(u^{-1},u^{-1},-u^{-1}).
\]
For the tacnode and triple point, regularity of the products on each
normalization branch excludes poles of order greater than two before
these coefficient conditions are imposed. The three displayed modules are contained in the
invertible sheaf allowing a simple pole on the reduced inverse-image
curve. This proves the claimed bound at every point.

\end{enumerate}
\end{proof}







\section{Normalization, Picard twists and ruled surfaces}\hypertarget{tag-0341}{}\label{sec:II:05}

\subsection{Duality and normal surface geometry}\label{subsec:II:normalization-duality}

\begin{lemma}[Conductor duality and its section maps]\hypertarget{tag-P222}{}\label{lem:II:conductor-duality}
\StatementRecord{lem:II:conductor-duality}{proved}{P222}{\cite[Proposition~2.3, p.~703, and Theorem~2.6, pp.~704--705]{Reid1994Nonnormal}}
Let $\eta:T\to S$ be the finite normalization of a reduced Gorenstein complex surface. Put $\mathcal B=\eta_*\OO_T$, $\mathcal Q=\mathcal B/\OO_S$, $\mathfrak c=\operatorname{Ann}(\mathcal Q)$, and $\Gamma=\operatorname{Spec}_S(\OO_S/\mathfrak c)$. The conductor scheme $\Gamma$ is Cohen--Macaulay of pure dimension one wherever nonempty, and its dualizing sheaf gives the canonical normalization sequence and comparison
\begin{equation}
 0\longrightarrow\OO_S\longrightarrow\mathcal B
 \longrightarrow
 \mathcal Ext^1_S(\OO_\Gamma,\OO_S)\longrightarrow0,
 \qquad
 \mathcal Q\simeq\omega_\Gamma\otimes\omega_S^{-1}|_\Gamma.
                                                               \hypertarget{tag-006W}{}\label{cdp:eq:A.RC.3}
\end{equation}
For a local ring $R=\OO_{S,p}$, its semilocal normalization $\overline R$, $A=R/\mathfrak c$ and $Q=\overline R/R$, the local comparisons are
\begin{equation}
 \operatorname{Ext}_R^1(Q,R)=A.                    \hypertarget{tag-005E}{}\label{cdp:eq:A.J.11}
\end{equation}
\begin{equation}
 Q\simeq\operatorname{Ext}_R^1(A,R)=\omega_A,
 \qquad \operatorname{End}_A(Q)=A.                \hypertarget{tag-005F}{}\label{cdp:eq:A.J.12}
\end{equation}
The endomorphism equality is the multiplication map. On $T$, with $\mathfrak c_T$ the conductor ideal, finite duality gives
\begin{equation}
 \omega_T=\mathfrak c_T\otimes\eta^*\omega_S
       \ \subset\ \eta^*\omega_S,
\hypertarget{tag-007H}{}\label{cdp:eq:A.RT.5}
\end{equation}
The ideal $\mathfrak c_T$ is reflexive. At a point
$q\in T$ with $(\mathfrak c_T)_q=\OO_{T,q}$, the
normalization is an isomorphism over a neighbourhood
of $\eta(q)$, including every local branch of $S$ there.
For compact $S$, these maps induce
\begin{equation}
 H^0(S,\mathcal Q)^\vee\simeq H^1(\Gamma,\omega_S|_\Gamma).
 \label{eq:II:conductor-section-duality}
\end{equation}
\end{lemma}
\begin{proof}
The finite-duality and subadjunction maps are those of
\cite[Proposition~2.3 and its proof, p.~703; Theorem~2.6 and its proof, pp.~704--705]{Reid1994Nonnormal}.
Their local form applies to the analytic rings here as follows. The normalization $\overline R$ is maximal Cohen--Macaulay over the two-dimensional Gorenstein ring $R$. A homomorphism $\overline R\to R$ is multiplication in the total quotient ring, and its multipliers are exactly $\mathfrak c$. Applying $\operatorname{Hom}_R(-,R)$ to $0\to R\to\overline R\to Q\to0$ gives \eqref{cdp:eq:A.J.11}; the positive Ext groups of $\overline R$ vanish. The depth lemma makes the nonzero $Q$ Cohen--Macaulay of dimension one. Gorenstein biduality in that dimension is $M\mapsto\operatorname{Ext}_R^1(M,R)$, so it makes $A$ Cohen--Macaulay and gives \eqref{cdp:eq:A.J.12}. Duality identifies its endomorphism ring with $\operatorname{End}_A(A)=A$, through multiplication. Applying the same duality to $0\to\mathfrak c\to R\to A\to0$ identifies the first arrow of \eqref{cdp:eq:A.RC.3} with the original inclusion $R\hookrightarrow\overline R$. Tensoring by the invertible dualizing sheaf supplies exactly the displayed global twists and \eqref{cdp:eq:A.RT.5}.

The canonical sheaf on the normal surface $T$ is reflexive,
so \eqref{cdp:eq:A.RT.5} gives the asserted reflexivity.
For the local isomorphism, put $R=\OO_{S,\eta(q)}$, let
$B$ be its finite semilocal normalization, and write
$\mathfrak q\subset B$ for the maximal ideal of $q$.
The equality $\mathfrak c B_{\mathfrak q}=B_{\mathfrak q}$
gives $c\in\mathfrak c\setminus\mathfrak q$.
Since $c\in R$ and $\mathfrak q\cap R=\mathfrak m_R$,
the element $c$ is a unit of $R$. The defining inclusion
$cB\subset R$ then gives $B=R$. The finite coherent
quotient $\eta_*\OO_T/\OO_S$ vanishes at $\eta(q)$,
and hence on a neighbourhood of that point.

Finally $\mathcal Ext^1_S(\mathcal Q,\omega_S)=\omega_S|_\Gamma$, and the other sheaf Ext groups vanish. The local-to-global Ext sequence gives
$\operatorname{Ext}^2_S(\mathcal Q,\omega_S)=H^1(\Gamma,\omega_S|_\Gamma)$.
Serre duality on compact $S$ identifies its left side with $H^0(S,\mathcal Q)^\vee$, proving \eqref{eq:II:conductor-section-duality}.
\end{proof}

\begin{lemma}[Conductor ideals of hypersurface branches]\hypertarget{tag-P251}{}\label{lem:II:hypersurface-branch-conductor}
\StatementRecord{lem:II:hypersurface-branch-conductor}{proved}{P251}{}
Let $R$ be the local analytic ring of a smooth complex manifold and let $f_1,\ldots,f_r$ be pairwise nonassociate irreducible elements. Put
\[
 A=R/(\textstyle\prod_i f_i),\qquad A_i=R/(f_i),\qquad
 g_i=(\textstyle\prod_{j\ne i}f_j)|_{A_i},
\]
and let $B_i$ normalize $A_i$. Write $\mathfrak c_i=\{b\in B_i:bB_i\subset A_i\}$ for its intrinsic conductor. Under the normalization map $A\hookrightarrow B=\prod_i B_i$, the conductor is
\[
 \mathfrak c_{B/A}=\prod_i g_i\mathfrak c_i.
\]
The identification of its $i$th summand with $\mathfrak c_i$ is multiplication by the nonzero element $g_i$.
\end{lemma}
\begin{proof}
The ring $R$ is factorial. Thus $\bigcap_{j\ne i}(f_j)=(\prod_{j\ne i}f_j)$, and the elements of $A$ supported only on the $i$th branch have $i$th components exactly $g_iA_i$. An element $b=(b_i)\in B$ lies in the conductor precisely when multiplication by every element supported on one branch stays in $A$, namely when $b_iB_i\subset g_iA_i$ for every $i$. Taking the multiplier $1$ gives $b_i=g_i a_i$ with $a_i\in A_i$. Since $g_i$ is a nonzerodivisor in the domain $B_i$, division by $g_i$ identifies the preceding condition with $a_iB_i\subset A_i$, or $a_i\in\mathfrak c_i$. This proves the formula and both directions of its multiplication-map identification.
\end{proof}

\begin{lemma}[Conductor exponents of plane branches]\hypertarget{tag-P282}{}\label{lem:II:plane-branch-conductor}
\StatementRecord{lem:II:plane-branch-conductor}{proved}{P282}{\cite[Tag~0A7M, Lemmas~47.16.1, 47.16.4 and~47.16.7]{StacksProject}}
Let $h=u\prod_{j=1}^r f_j$ be a reduced plane-curve equation, with $u$ a unit and $f_j$ the equations of its distinct irreducible branches. On the normalization of branch $C_i$, let $c_i$ be the exponent of the conductor of the whole curve and let $c_i^{\rm int}$ be the conductor exponent of that branch alone. Then
\[
 c_i=c_i^{\rm int}+\sum_{j\ne i}I(C_i,C_j),\qquad
 c_i^{\rm int}=2\delta(C_i),\qquad
 \sum_i c_i=2\delta(h).
\]
In particular, $c_i=1$ implies that $C_i$ is smooth and meets precisely one other branch, smoothly and transversely. If $c_i=2$, the other branches, counted with their positive intersection multiplicities, number at most two.
\end{lemma}
\begin{proof}
Put $A=\mathbb C\{x,y\}/(h)$, let $B=\prod_i\mathbb C\{t_i\}$ be its finite normalization, and put $Q=B/A$. The one-dimensional hypersurface ring $A$ is Gorenstein, with normalized dualizing complex $A[1]$. A nonzerodivisor of $A$ acts injectively on $B$, so $B$ is a maximal Cohen--Macaulay $A$-module. Lemma~47.16.7 of \cite[Tag~0A7M]{StacksProject}, applied with this dualizing complex, gives $\operatorname{Ext}^1_A(B,A)=0$.

Inside the total quotient ring of $A$, every $A$-linear map $B\to A$ is multiplication by its value at $1$. Thus
\[
 I=\operatorname{Hom}_A(B,A)=\{a\in A:aB\subset A\}
\]
is the conductor. Since $Q$ has finite length and $A$ has depth one, $\operatorname{Hom}_A(Q,A)=0$. Applying $\operatorname{Hom}_A(-,A)$ to $0\to A\to B\to Q\to0$ gives
\[
 0\longrightarrow I\longrightarrow A\longrightarrow
 \operatorname{Ext}^1_A(Q,A)\longrightarrow0.
\]
For the normalized complex $A[1]$, Lemma~47.16.4 of \cite[Tag~0A7M]{StacksProject} identifies $\operatorname{Ext}^1_A(Q,A)$ with the Matlis dual of $Q$. This exact duality preserves length: it sends the unique simple module $\mathbb C$ to $\mathbb C$ and reverses a composition series. Hence
\[
 \operatorname{length}(A/I)=\operatorname{length}(B/A)=\delta(h),
 \qquad \operatorname{length}(B/I)=2\delta(h).
\]
The conductor is the product of the ideals $(t_i^{c_i})$ in the normalization, so the second length is $\sum_i c_i$. Applying the same argument to the individual branch ring gives $c_i^{\rm int}=2\delta(C_i)$. Apply \cref{lem:II:hypersurface-branch-conductor} in $R=\mathbb C\{x,y\}$ to these same factors $f_i$. Its multiplication map by $g_i=\prod_{j\ne i}f_j|_{C_i}$ adds $\operatorname{ord}_{t_i}g_i=\sum_{j\ne i}I(C_i,C_j)$ to the intrinsic conductor exponent, proving the branch formula.

Choose a linear coordinate $x$ nonconstant on every branch and a complementary coordinate $y$ generically. Hypersurface adjunction gives the residue generators $dx/h_y$ for $\omega_A$ and $dx/(f_i)_y$ for $\omega_{C_i}$. The finite duality of \cite[Tag~0A7M, Lemma~47.16.1]{StacksProject}, applied to each local finite map $A\to\mathbb C\{t_i\}$ and then taking their product, identifies
\[
 \omega_B=\operatorname{Hom}_A(B,\omega_A)=I\omega_A.
\]
These are submodules of the meromorphic differentials on the normalized branches: the birational normalization identifies them at the generic points, where the finite-duality trace is the identity. On $C_i$ one has
\[
 h_y|_{C_i}=u|_{C_i}\,(f_i)_y|_{C_i}\prod_{j\ne i}f_j|_{C_i}.
\]
Thus the residue generators realize the same multiplication by $g_i$ in the conductor formula, with the unit $u|_{C_i}$ invertible. Intersection multiplicities of distinct plane branches are positive integers. The cases $c_i=1,2$ follow; intersection multiplicity one makes both branches smooth and transverse.
\end{proof}



\begin{proposition}[Conormal sections with trivial normalization pullback]\hypertarget{tag-01AZ}{}\label{cdp:st:A.5.1}
\StatementRecord{cdp:st:A.5.1}{proved}{01AZ}{\SeriesRef{cdp:st:A.1.1}, \SeriesRef{cdp:st:A.4.1}}

Let $S$ be a connected reduced full fibre of a proper map from a complex manifold to a smooth curve. For every line bundle $A$ on $S$ with $\eta^*A\simeq\OO_{\widetilde S}$, the group $H^0(S,V\otimes A)$ contains the canonical image of $H^0(\widetilde S,\eta^*(N^*\otimes A))$. A local base coordinate and a choice of trivialization of $\eta^*A$ identify this subspace with $H^0(\widetilde S,\OO_{\widetilde S})$. In particular it is nonzero. This applies to the subcase of \cite[Proposition~7.3, pp.~693--694]{CDP2020} in which all $\eta^*(L_t|_S)$ are trivial. It is compatible with the vanishing of $H^0(S,Q\otimes A)$ through the maps in Proposition~\ref{cdp:st:A.1.1}, as \ref{cdp:st:A.4.1} demonstrates.
\end{proposition}


\begin{proof}
\begin{enumerate}[label=\arabic*.,leftmargin=*]
\item Proposition~\ref{cdp:st:A.2.2} gives the inclusion in
\eqref{cdp:eq:A.3}. Tensoring it by the invertible sheaf $A$ and applying
the finite projection formula gives
\[
 \eta_*\eta^*(N^*\otimes A)\hookrightarrow K\otimes A
                         \hookrightarrow V\otimes A.
\]
Locally, for a defining equation $F$ of $S$, the first map sends
$b\otimes[F]$ to $b\,dF|_S$. The inclusion of the Jacobian ideal in the
conductor makes its coefficients holomorphic on $S$. Under $F'=uF$,
both the conormal frame and $dF|_S$ are multiplied by $u$, so these maps
glue. Taking global sections gives the claimed canonical injection.

\item A coordinate $t$ about the image point on the base defines the full
fibre by $F=t\circ f$. Since this full fibre is reduced, $[F]$ is a frame
of $N^*$. Its pullback, together with the chosen trivialization of
$\eta^*A$, identifies the source with
$H^0(\widetilde S,\OO_{\widetilde S})$. The nonempty compact normalization
has the nonzero constant function $1$ on each connected component.
This proves nonvanishing and applies to each $A=L_t|_S$ in the stated
published subcase.

\item The image lies in $H^0(S,K\otimes A)$, the kernel of
$H^0(S,V\otimes A)\to H^0(S,Q\otimes A)$ by
\eqref{cdp:A:filtration2}. Thus these nonzero sections map to zero in
the quotient group. Proposition~\ref{cdp:st:A.4.1} gives the stated
example with vanishing quotient group.
\end{enumerate}
\end{proof}

\begin{remark}[The top dualizing-sheaf group]\hypertarget{tag-P223}{}\label{rem:II:normal-duality-endpoint}
\StatementRecord{rem:II:normal-duality-endpoint}{proved}{P223}{}
For a connected compact normal complex surface $T$, Serre duality gives
$H^2(T,\omega_T)^\vee\simeq H^0(T,\OO_T)=\C$.
This is the group printed as zero in the proof of
\cite[Proposition~7.3, Step~1, p.~693]{CDP2020}.
Replacing $\omega_T$ there by $\OO_T$ recovers the vanishing stated at the start of that step. The genus-one ruled case and its twisted cotangent sections are computed in \cref{cdp:st:A.4.2}; the unitary local-system proof of the proposition is \cref{cdp:st:A.4.9}.
\end{remark}

\begin{lemma}[Exceptional intersection matrices]
\hypertarget{tag-P243}{}\label{lem:II:exceptional-intersection}
\StatementRecord{lem:II:exceptional-intersection}{proved}{P243}{\cite[p.~6]{Mumford1961}}
Let $\pi:T\to N$ be a resolution of a normal complex surface and
let $E_1,\ldots,E_r$ be the exceptional curves over a singular
point $p$. Their intersection matrix $M=(E_i\cdot E_j)$ is
negative definite, and every entry of $-M^{-1}$ is strictly
positive. Consequently the equation $Md=-b$, with $b_i\ge0$,
gives $d_i\ge0$; when $b\ne0$, every $d_i$ is positive.
\end{lemma}
\begin{proof}
The fibre $\pi^{-1}(p)$ is connected. A decomposition into two
disjoint compact parts would, by properness, give disjoint
neighbourhoods whose union contains $\pi^{-1}(U)$ for a sufficiently
small neighbourhood $U$ of $p$. The function with values $0$
and $1$ on those neighbourhoods descends to a bounded holomorphic
function on $U\setminus\{p\}$. Normality extends it across $p$,
whereas approach to the two exceptional parts gives different
limiting values.

Choose a nonzero local holomorphic function vanishing at $p$.
Its pullback divisor is $H+\sum_i n_iE_i$, where $n_i>0$ and
$H$ is effective without exceptional components. Properness
makes $H$ meet the exceptional fibre. The principal divisor
has degree zero on each $E_i$, so
\[
 Mn=-h,\qquad h_i=H\cdot E_i\ge0,\qquad h\ne0.
\]
For $v_i=n_iw_i$, expansion gives
\[
 v^tMv
 =-\sum_{i<j}M_{ij}n_in_j(w_i-w_j)^2
  -\sum_i n_ih_iw_i^2.
\]
This is the weighted intersection identity of
\cite[p.~6]{Mumford1961}: the local principal divisor supplies
the vector $n$, and $v=\operatorname{diag}(n_i)w$ is the change
of variables. The off-diagonal entries of $M$ are nonnegative.
Connectedness of the exceptional fibre and $h\ne0$ make the
right-hand side strictly negative for $v\ne0$.

Put $Q=-M$ and choose $c$ larger than its largest eigenvalue.
Then $A=I-Q/c$ is entrywise nonnegative and has spectral radius
less than one. Thus
\[
 -M^{-1}=Q^{-1}=c^{-1}\sum_{k\ge0}A^k.
\]
For each pair of vertices, a path in the connected intersection
graph gives a positive entry in one power of $A$. Every entry
of the sum is therefore positive. Multiplication by $b$ proves
the final assertions.
\end{proof}

\begin{proposition}[Rational singularities for a vertically resolved ruled normalization]\hypertarget{tag-01BI}{}\label{cdp:st:A.5.8}
\StatementRecord{cdp:st:A.5.8}{proved}{01BI}{\SeriesRef{cdp:eq:A.J.19}}

Let $\pi:Z\to T$ be a resolution of a compact connected
normal complex surface, where $Z$ is an irrational ruled
surface with ruling $p:Z\to B$. Suppose that every
exceptional curve of $\pi$ is vertical for $p$.
Then there is a unique holomorphic map
$\widetilde p:T\to B$ satisfying
$p=\widetilde p\circ\pi$, and
\begin{equation}
 R\pi_*\mathcal O_Z=\mathcal O_T,
 \qquad R^1\widetilde p_*\mathcal O_T=0.           \hypertarget{tag-005K}{}\label{cdp:eq:A.J.17}
\end{equation}
In particular all singularities of $T$ are rational. For
$g=g(B)$, its complete structure-sheaf cohomology dimensions
are
\begin{equation}
 (h^0(T,\mathcal O_T),h^1(T,\mathcal O_T),
                         h^2(T,\mathcal O_T))=(1,g,0).
                                                        \hypertarget{tag-005L}{}\label{cdp:eq:A.J.18}
\end{equation}
\end{proposition}


\begin{proof}

\begin{enumerate}[label=\arabic*.,leftmargin=*]
\item The exceptional fibre of the proper birational map $\pi$
over a point is connected, because $T$ is normal and
$\pi_*\mathcal O_Z=\mathcal O_T$. Its curve components
are vertical by assumption, so their intersection
connectivity makes $p$ constant on that whole fibre.
It is also constant on the one-point fibres. The equality $\pi_*\OO_Z=\OO_T$ gives holomorphic descent as follows.

\item Fix $t\in T$ and let $U\subset B$ be a coordinate disc
containing $p(\pi^{-1}(t))$. Properness makes
$\pi(Z\setminus p^{-1}(U))$ closed and disjoint from
$t$. There is therefore an open neighbourhood $W$ of $t$
with $\pi^{-1}(W)\subset p^{-1}(U)$. The coordinate
function of $p$ on $\pi^{-1}(W)$ is a section of
$\pi_*\mathcal O_Z(W)=\mathcal O_T(W)$, so descends
to a holomorphic function on $W$. Its values lie in the
coordinate image of $U$; the inverse coordinate gives a
holomorphic map $W\to U$. These local maps agree after
pullback to the surjective map $\pi$, hence agree on
overlaps. They construct the required unique
$\widetilde p$.

\item Contracting vertical ruling $(-1)$-curves factors $p$ as
\[
 Z\xrightarrow{b} Z_0\xrightarrow{p_0}B,
\]
where $Z_0\to B$ is a $\mathbb P^1$-bundle and $b$
is a finite sequence of point blowups. Every such blowup
has direct image $\mathcal O$ and no higher direct image
of $\mathcal O$: the successive layers on its
infinitesimal exceptional curves are
$\mathcal O_{\mathbb P^1}(k)$ for $k\geq0$, whose
degree-one cohomology is zero, and formal functions gives
the assertion. The bundle $p_0$ has the same direct-image
property by the cohomology of $\mathbb P^1$ and local
triviality. Composition consequently gives
\begin{equation}
 p_*\mathcal O_Z=\mathcal O_B,
 \qquad R^j p_*\mathcal O_Z=0\quad(j>0).          \hypertarget{tag-005M}{}\label{cdp:eq:A.J.19}
\end{equation}

\item The sheaf $\mathcal T=R^1\pi_*\mathcal O_Z$ has
finite support at the singular points of $T$. The
low-degree Leray sequence for $p=\widetilde p\circ\pi$
contains the arrows
\begin{equation}
 0\longrightarrow R^1\widetilde p_*\mathcal O_T
 \longrightarrow R^1p_*\mathcal O_Z
 \longrightarrow\widetilde p_*\mathcal T
 \longrightarrow R^2\widetilde p_*\mathcal O_T.    \hypertarget{tag-005O}{}\label{cdp:eq:A.J.20}
\end{equation}
The last term is zero: the nonconstant proper map
$\widetilde p$ has fibres of complex dimension one.
The middle term is zero by \eqref{cdp:eq:A.J.19}. Thus
$R^1\widetilde p_*\mathcal O_T=0$ and
$\widetilde p_*\mathcal T=0$. Pushforward is faithful
on a sheaf with finite support, since its stalk sections
at finitely many points cannot disappear under a proper
map. Hence $\mathcal T=0$. Higher $R^j\pi_*$ vanish
for $j>1$ by the dimension of its fibres. This proves
\eqref{cdp:eq:A.J.17} and the defining rational-singularity property.

\item Finally \eqref{cdp:eq:A.J.17} identifies every
$H^q(T,\mathcal O_T)$ with $H^q(Z,\mathcal O_Z)$.
Equation \eqref{cdp:eq:A.J.19} identifies the latter with
$H^q(B,\mathcal O_B)$. A compact connected smooth curve
has dimensions $(1,g)$ in degrees zero and one and no
higher cohomology, proving \eqref{cdp:eq:A.J.18}.

\end{enumerate}
\end{proof}

\begin{proposition}[Normal Gorenstein surfaces with vanishing first Chern class and Euler characteristic]\hypertarget{tag-01BA}{}\label{cdp:st:A.5.2}
\StatementRecord{cdp:st:A.5.2}{proved}{01BA}{\cite[Theorem~2.2, p.~843]{Peters2021}}

A compact connected normal Gorenstein complex surface $S$ with $c_1(\omega_S)=0$ in $H^2(S,\Q)$ and $\chi(S,\OO_S)=0$ is smooth. Consequently every singular irreducible reduced divisor in a compact complex threefold $Y$ with $H^2(Y,\Q)=0$ is nonnormal. The normal-surface smoothness step in \cite[proof of Proposition~4.11, p.~687]{CDP2020} follows under these numerical hypotheses.
\end{proposition}


\begin{proof}

\begin{enumerate}[label=\arabic*.,leftmargin=*]
\item Let $\pi:Z\to S$ be the minimal resolution and set $\ell=h^0(S,R^1\pi_*\OO_Z)$. Normality gives $\pi_*\OO_Z=\OO_S$, and the fibres have dimension at most one, so Leray gives $\chi(Z,\OO_Z)=-\ell$. We first prove $\ell=0$. The classification table in \cite[Theorem~2.2, p.~843]{Peters2021}, together with Noether's formula and invariance of $\chi(\OO)$ under blowing up, shows that a smooth compact surface of negative holomorphic Euler characteristic is a blown-up ruled surface over a smooth curve $B$ of genus $g\ge2$. Indeed the sums $c_1^2+c_2$ in every other row of that table are nonnegative. Thus, when $\ell>0$, there is a ruling $p:Z\to B$ and $\chi(Z,\OO_Z)=1-g$, giving $\ell=g-1$.

\item An exceptional irreducible curve $E$ of $\pi$ dominating $B$ would have $h^1(E,\OO_E)\ge g$: normalize $E$, apply Riemann--Hurwitz to its finite map to $B$, and use the normalization exact sequence on $E$. The exact sequence $0\to\mathcal I_E\to\OO_Z\to\OO_E\to0$, pushed forward by $\pi$, gives a surjection $R^1\pi_*\OO_Z\to R^1\pi_*\OO_E$ because $R^2\pi_*\mathcal I_E=0$. It follows that $\ell\ge h^1(E,\OO_E)\ge g$, a contradiction. Hence every exceptional curve is vertical for $p$. Apply \cref{cdp:st:A.5.8} to this ruling and the resolution $\pi:Z\to S$. The verticality just proved gives its hypotheses, so $R^1\pi_*\OO_Z=0$, contradicting $\ell>0$. Thus every singularity of $S$ is rational.

\item The canonical-cycle conclusion is \cite[Theorem~4.20, p.~99, arXiv v1]{Reid1997Chapters}. Its local analytic canonical-bundle comparison here is as follows. Proper duality applied to $R\pi_*\OO_Z=\OO_S$ identifies $\pi_*\omega_Z$ with the invertible dualizing sheaf $\omega_S$. A local generator of $\omega_S$ therefore pulls back to a holomorphic two-form on $Z$, agreeing with the usual pullback off the exceptional set. Its zero divisor is effective and exceptional. These divisors glue to $D\ge0$ with $K_Z=\pi^*K_S+D$. Every exceptional irreducible curve has arithmetic genus zero by the same surjection to its first cohomology, hence is a smooth rational curve. Minimality of the resolution gives $E^2\le-2$, and adjunction gives $D\cdot E=K_Z\cdot E=-2-E^2\ge0$. The exceptional intersection matrix is negative definite. Thus a nonzero $D$ would satisfy both $D^2<0$ and $D^2=\sum_E d_E(D\cdot E)\ge0$. Therefore $D=0$.

\item Consequently $c_1(K_Z)=0$ in rational cohomology. It follows by adjunction that $Z$ contains no $(-1)$-curve, so $Z$ is a minimal compact surface; Noether's formula gives $c_2(Z)=0$. We exclude rational curves on each surviving row of the classification. A ruled surface is excluded by $K_Z\cdot F=-2$ on a ruling fibre. General-type surfaces have $K_Z^2>0$ and K3 and Enriques surfaces have positive $c_2$, so none occurs. A torus has no rational curve because a map from $\PP^1$ lifts to its vector-space universal cover and is constant; the same holds for a finite etale quotient of a torus. A primary Kodaira surface is an elliptic bundle over an elliptic curve, so a rational curve would map constantly to the base and then constantly to a fibre; finite etale quotients have the same property. A class VII surface here has $b_2=c_2=0$, and the Hopf--Inoue classification gives no rational curves. Equivalently one can use their standard universal covers: $\C^2\setminus\{0\}$ for a primary Hopf cover and $\mathbb H\times\C$ for an Inoue surface. A map from $\PP^1$ lifts and has constant coordinate functions.

\item It remains to check the properly elliptic row. For a minimal elliptic surface, the Kodaira fibre list expresses $c_2$ as a sum of nonnegative fibre Euler numbers, with zero contributions exactly at multiples of smooth elliptic curves. Since $c_2(Z)=0$, all fibre reductions are smooth elliptic curves. The canonical bundle formula on a Kähler elliptic surface with $\chi(\OO_Z)=0$ has rational first Chern class
\[
 c_1(K_Z)=\left(2g(B)-2+\sum_i(1-m_i^{-1})\right)[F].
\]
In Kodaira dimension one the coefficient is positive, and the Kähler class has positive integral over $F$, contradicting $c_1(K_Z)=0$. In the non-Kähler case, a horizontal rational curve would be a multisection $C$. For $D=C+nF$ with $n$ sufficiently large, $D^2>0$ and $(K_Z-mD)\cdot F<0$ for $m>0$. Riemann--Roch therefore gives quadratic growth of $h^0(Z,\OO_Z(mD))$, making $Z$ Moishezon and hence projective as a smooth compact surface. This contradicts the non-Kähler case. Vertical rational curves have already been excluded by the fibre list. Thus $Z$ has no rational curves in any case. The exceptional curves of $\pi$ were rational, so there are none and $S$ is smooth.

\item Finally take an irreducible reduced divisor $S\subset Y$ with $H^2(Y,\Q)=0$. It is a Cartier divisor in the smooth manifold $Y$, hence Gorenstein, and adjunction gives $\omega_S=(K_Y\otimes\OO_Y(S))|_S$, with rational first Chern class zero. The divisor sequence and Hirzebruch--Riemann--Roch give
$\chi(\OO_S)=\chi(\OO_Y)-\chi(\OO_Y(-S))=0$, since $c_1(\OO_Y(S))=0$ rationally. A normal $S$ is therefore smooth by the result just proved.

\end{enumerate}
\end{proof}

\begin{proposition}[The effective conductor divisor on a minimal resolution]\hypertarget{tag-01BD}{}\label{cdp:st:A.5.3}
\StatementRecord{cdp:st:A.5.3}{proved}{01BD}{\cref{lem:II:conductor-duality,lem:II:exceptional-intersection}; \cite[Lemma~1.1 and Corollary~1.4, p.~5]{Furushima1993}}

Let $S$ be a compact reduced hypersurface in a smooth threefold, let $\eta:\coprod_i T_i\to S$ be its normalization, and let $\pi_i:Z_i\to T_i$ be the minimal resolutions. Put $h_i=\eta\pi_i$. There are canonical effective integral divisors $D_i$ on $Z_i$ with
\begin{equation}
 \omega_{Z_i}(D_i)\simeq h_i^*\omega_S.
 \hypertarget{tag-004H}{}\label{cdp:eq:A.25}
\end{equation}
The nonexceptional part of $D_i$ is exactly the divisorial conductor, with its scheme-theoretic multiplicities. On each connected exceptional configuration, all coefficients are strictly positive when the configuration meets the strict conductor or contains a curve of positive geometric genus. A zero coefficient forces that configuration to consist of smooth rational $(-2)$-curves disjoint from the strict conductor. When $S\subset Y$ and $H^2(Y,\Q)=0$, every $K_{Z_i}+D_i$ has rational first Chern class zero. A ruled $Z_i\to B$ with $g(B)\ge2$ then has
\begin{equation}
 D_i=2C+V,
 \hypertarget{tag-004I}{}\label{cdp:eq:A.26}
\end{equation}
where $C$ is a section and $V$ is vertical. In the case where $C$ is $\pi_i$-exceptional, $\eta^*(L|_S)|_{T_i}$ is trivial for every $L\in\Pic(Y)$.
\end{proposition}


\begin{proof}

\begin{enumerate}[label=\arabic*.,leftmargin=*]
\item Apply the finite-duality map of \cref{lem:II:conductor-duality} to the normalization component $T_i$. It identifies $\omega_{T_i}$ with $\mathfrak c_i\otimes\eta^*\omega_S$, where $\mathfrak c_i$ is the conductor ideal on $T_i$. The multiplier identification there makes $\mathfrak c_i$ maximal Cohen--Macaulay. On a normal surface a rank-one maximal Cohen--Macaulay module is reflexive, so the displayed equality is determined by its codimension-one orders. Those orders define an effective conductor divisor $D_i^0$ on $T_i$.

\item Off the exceptional set the canonical identification pulls back to $\omega_{Z_i}(D_i^{\rm str})\simeq h_i^*\omega_S$, where $D_i^{\rm str}$ denotes strict transform. The meromorphic extension of this line-bundle identification gives integral exceptional coefficients $a_j$ with
$K_{Z_i}+D_i^{\rm str}+\sum_j a_jE_j=h_i^*K_S$.
We verify their signs. For each exceptional irreducible curve $E_j$, minimality excludes a smooth rational $(-1)$-curve. Adjunction therefore gives
$K_{Z_i}\cdot E_j=2p_a(E_j)-2-E_j^2\ge0$;
also $D_i^{\rm str}\cdot E_j\ge0$. The exceptional intersection matrix $M=(E_j\cdot E_k)$ is negative definite with nonnegative off-diagonal entries, and
\[
 M a=-v,\qquad v_j=(K_{Z_i}+D_i^{\rm str})\cdot E_j\ge0.
\]
By \cref{lem:II:exceptional-intersection}, every entry of $-M^{-1}$ is strictly positive on a connected exceptional block. Hence $a=-M^{-1}v$ has nonnegative entries, and all entries are positive when $v\ne0$. A meeting with $D_i^{\rm str}$ makes some $v_j>0$. A component of positive geometric genus also does so: $p_a(E_j)\ge1$ and $E_j^2<0$ give $K_{Z_i}\cdot E_j>0$. Conversely, a zero coefficient forces $v=0$ on its block. Then every $K_{Z_i}\cdot E_j=0$ and $D_i^{\rm str}\cdot E_j=0$; adjunction and minimality make every $E_j$ a smooth rational $(-2)$-curve. Thus $D_i=D_i^{\rm str}+\sum_j a_jE_j$ proves \eqref{cdp:eq:A.25} and all assertions about exceptional coefficients. Furushima's Lemma~1.1 treats nonnormal irreducible reduced projective Gorenstein surfaces; his Corollary~1.4 is stated for such a surface with trivial dualizing sheaf, $h^1(\OO)=0$, and a very ample polarization \cite[Lemma~1.1, Definition~1.2 and Corollary~1.4, p.~5]{Furushima1993}. The preceding finite-duality and exceptional-matrix maps apply to the analytic Gorenstein surface of the statement.

\item For a divisor in $Y$, adjunction gives $\omega_S=(K_Y\otimes\OO_Y(S))|_S$, proving the numerical assertion. On a general ruling fibre $f$, $K_{Z_i}\cdot f=-2$, so $D_i\cdot f=2$. Every horizontal irreducible curve has normalization of genus at least $g(B)\ge2$. Such a curve $C_0$ cannot occur in $D_i$ with coefficient one: otherwise
\[
 0=(K_{Z_i}+D_i)\cdot C_0
   =2p_a(C_0)-2+(D_i-C_0)\cdot C_0>0.
\]
Every horizontal coefficient is consequently at least two. The degree-two intersection with $f$ forces exactly one horizontal component, coefficient two and degree one over $B$, giving \eqref{cdp:eq:A.26}. The degree-one finite morphism $C\to B$ is an isomorphism because $B$ is normal; thus $C$ is a smooth section.

\item Finally, integral $H^2$ of a smooth ruled surface and of its blowups is torsion-free, and its topologically trivial line bundles are precisely the pullbacks of $\Pic^0(B)$. These facts follow from the projective-bundle and blowup cohomology formulas and from $p_*\OO_{Z_i}=\OO_B$, $R^1p_*\OO_{Z_i}=0$. Restriction to the section $C$ is inverse to this Picard pullback. The line $h_i^*L$ has torsion, hence zero, integral first Chern class. When $C$ is exceptional, its restriction to $C$ is trivial because $h_i$ maps $C$ to a point. Therefore $h_i^*L$ is trivial on $Z_i$. Normality gives $(\pi_i)_*\OO_{Z_i}=\OO_{T_i}$, and the projection formula descends this triviality to $T_i$. This proves the last assertion.

\end{enumerate}
\end{proof}

\subsection{Persistent Picard twists}\label{subsec:II:normalization-picard}

\begin{proposition}[The remaining normalization class for a persistent ambient obstruction]\hypertarget{tag-01BE}{}\label{cdp:st:A.5.4}
\StatementRecord{cdp:st:A.5.4}{proved}{01BE}{\SeriesRef{cdp:st:A.5.3}, \SeriesRef{cdp:st:A.4.9}, \SeriesRef{cdp:st:A.1.1}; \cite[Lemma~2.3(1)--(2), pp.~5--6, arXiv v3]{Teleman2006}}

Let $Y$ satisfy the CDP hypotheses, and let $S$ be a connected nonnormal reduced fibre support for which
\begin{equation}
 H^0(S,\Omega_Y^1|_S\otimes L|_S)\ne0
 \quad\hbox{for every }L\in\Pic(Y).
 \hypertarget{tag-004J}{}\label{cdp:eq:A.27}
\end{equation}
There is a normalization component $T_i$ whose minimal resolution $Z_i$ has Kodaira dimension $-\infty$. Every line from $\Pic(Y)$ pulls back to a unitary-flat line on that $Z_i$. A rational or class VII $Z_i$ has trivial Picard pullback from $Y$, and then
\begin{equation}
 H^2\!\left(T_i,\omega_{T_i}\otimes
             \eta^*((K_Y^{-1}\otimes L)|_S)|_{T_i}\right)\simeq\C
 \quad\hbox{for every }L\in\Pic(Y).
 \hypertarget{tag-004K}{}\label{cdp:eq:A.28}
\end{equation}
The same conclusion holds for the exceptional-section case of \ref{cdp:st:A.5.3}. Thus the only normalization type left by this argument, when \eqref{cdp:eq:A.28} has not been obtained, is an irrational ruled surface with nonconstant Picard pullback; in base genus at least two its divisor \eqref{cdp:eq:A.26} has its horizontal section in the strict conductor.
\end{proposition}


\begin{proof}

\begin{enumerate}[label=\arabic*.,leftmargin=*]
\item \ref{cdp:st:A.4.9} gives $H^0(S,\widetilde\Omega_S^1\otimes L|_S)=0$ on a dense open subset of $\Pic(Y)$. By \ref{cdp:st:A.1.1} and \eqref{cdp:eq:A.27}, $H^0(S,K\otimes L|_S)$ is nonzero there. Its nonvanishing locus is closed analytic in the holomorphic Picard family, hence is all of $\Pic(Y)$.

\item Pull $K$ to each $Z_i$, quotient the resulting sheaf by its torsion, and take its double dual. The result is a fixed holomorphic line bundle $M_i$, since its generic rank on each component is one and $Z_i$ is a smooth surface. The induced maps give an injection
\begin{equation}
 H^0(S,K\otimes L|_S)\hookrightarrow
       \bigoplus_i H^0(Z_i,M_i\otimes h_i^*L).
 \hypertarget{tag-004L}{}\label{cdp:eq:A.29}
\end{equation}
Indeed a section in the kernel vanishes generically on every component, and $K\subset\Omega_Y^1|_S$ is torsion-free over the reduced space. Each summand's nonvanishing locus is closed analytic in the connected complex vector space $\Pic(Y)=H^1(Y,\OO_Y)$. A finite union of proper analytic subsets cannot cover that space. Consequently one component $i$ has
$H^0(Z_i,M_i\otimes h_i^*L)\ne0$ for every $L$.

\item Fix a Gauduchon metric on $Z_i$, with fundamental form $\omega$ satisfying $\partial\bar\partial\omega=0$. For Chern curvature $F_G$, use $\deg_\omega G=(2\pi)^{-1}\int_{Z_i}\mathrm iF_G\wedge\omega$; it is independent of the chosen Hermitian metric and additive under tensor product. A nonzero section has an effective zero divisor, so its line has nonnegative degree. For any $G\in\Pic(Y)$ the preceding nonvanishing applied to $G^{\otimes n}$ therefore gives
\[
 \deg_\omega(M_i)+n\deg_\omega(h_i^*G)\ge0
 \quad\hbox{for every }n\in\Z.
\]
It follows that $\deg_\omega(h_i^*G)=0$ for every $G$ and every Gauduchon metric.

\item Let $F_G$ be the Chern curvature of a Hermitian metric on $h_i^*G$. Since $H^2(Y,\Z)=0$, the real closed form $\mathrm iF_G$ is exact. Its pairing with $\omega$ is zero by the degree calculation. The second exact sequence in \cite[Lemma~2.3(1), p.~5, arXiv v3]{Teleman2006} therefore gives a real function $u$ with $\mathrm iF_G=\mathrm i\partial\bar\partial u$. With the Chern-curvature convention $F=-\partial\bar\partial\log h$, replacing the line metric $h$ by $e^u h$ changes its curvature to $F_G-\partial\bar\partial u=0$. Its Chern connection is consequently flat and unitary.

\item The divisor $D_i$ of \ref{cdp:st:A.5.3} is nonzero. For an irreducible nonnormal support this follows from its nonempty conductor; for a connected reducible support every component meets another component, again giving a nonempty conductor. If $Z_i$ had nonnegative Kodaira dimension, some positive pluricanonical power would have a nonzero section, making $\deg_\omega(K_{Z_i})\ge0$. But $D_i$ is a nonzero effective divisor and therefore has strictly positive degree. Equation \eqref{cdp:eq:A.25}, with $\omega_S=(K_Y\otimes\OO_Y(S))|_S$, would give
\[
 0=\deg_\omega h_i^*(K_Y\otimes\OO_Y(S))
   =\deg_\omega(K_{Z_i})+\deg_\omega(D_i)>0,
\]
a contradiction. Thus $\kappa(Z_i)=-\infty$, whose classification consists of rational, irrational ruled, and class VII surfaces.

\item For a rational $Z_i$, topologically trivial lines are trivial. For a class VII $Z_i$, $h^1(\OO_{Z_i})=1$ and $b_1(Z_i)=1$. With our exponential convention $u\mapsto\exp(2\pi\mathrm i u)$, the unitary-flat subgroup of $\Pic^0(Z_i)$ has real tangent space equal to the injective image of $H^1(Z_i,\mathbb R)$ in $H^1(Z_i,\OO_{Z_i})$. The injective map is \cite[Lemma~2.3(2), pp.~5--6, arXiv v3]{Teleman2006}, after multiplication by the fixed scalar $\mathrm i$ to match the convention $u\mapsto\exp(2\pi\mathrm i u)$. That real tangent space has dimension one. The holomorphic Picard pullback has complex-linear differential and, by the proved flatness, takes values in this real one-dimensional subspace. Its differential is zero, so the homomorphism is constant and equals zero at the trivial line. Every $h_i^*L$ is therefore trivial. Projection formula and $(\pi_i)_*\OO_{Z_i}=\OO_{T_i}$ again descend triviality to $T_i$. The exceptional-section case has the same conclusion by \ref{cdp:st:A.5.3}.

\item In each of these cases, the line multiplying $\omega_{T_i}$ in \eqref{cdp:eq:A.28} is trivial, since it too comes from $Y$. A normal complex surface is Cohen--Macaulay, so its dualizing sheaf and Serre duality give $H^2(T_i,\omega_{T_i})=H^0(T_i,\OO_{T_i})^\vee=\C$. This proves \eqref{cdp:eq:A.28} and the stated exact reduction of the remaining normalization problem.

\end{enumerate}
\end{proof}

\begin{proposition}[Uniform negative-degree vanishing for a fixed vector bundle]\hypertarget{tag-01BG}{}\label{cdp:st:A.5.6}
\StatementRecord{cdp:st:A.5.6}{proved}{01BG}{\SeriesRef{cdp:st:A.5.4}; \cite[Lemma~2.3]{Teleman2006}}

Let $Z$ be a compact smooth complex surface, $\omega$ a Gauduchon metric, and $E$ a fixed holomorphic vector bundle. There is a real constant $b(E,\omega)$ such that
\begin{equation}
 \deg_\omega P<b(E,\omega)
 \quad\Longrightarrow\quad H^0(Z,E\otimes P)=0
 \hypertarget{tag-004N}{}\label{cdp:eq:A.30}
\end{equation}
for every holomorphic line bundle $P$. Thus persistence of sections of $E\otimes h^*L$ for every $L\in\Pic(Y)$ forces $\deg_\omega h^*L=0$ for every $L$, for any holomorphic map $h:Z\to Y$. Persistence on one Picard component already forces degree zero on the pullback of $\Pic^0(Y)$. Under $H^2(Y,\Q)=0$, it forces degree zero and unitary flatness for every ambient Picard pullback. On a class VII surface $Z$ with $H^2(Y,\Z)=0$, the entire holomorphic Picard pullback is the zero map.
\end{proposition}


\begin{proof}

\begin{enumerate}[label=\arabic*.,leftmargin=*]
\item Choose a Hermitian metric on $E$ and let $C_E$ bound the largest eigenvalue of the Hermitian curvature endomorphism $\mathrm i\Lambda_\omega F_E$ on the compact surface. Write $V_\omega=\int_Z\omega^2/2$ and normalize degree by
$\deg_\omega P=(2\pi)^{-1}\int_Z\mathrm iF_P\wedge\omega$.
The scalar Gauduchon elliptic equation changes a Hermitian line metric so that its contracted curvature is the constant
\[
 \mathrm i\Lambda_\omega F_P
       =c_P:=2\pi\deg_\omega(P)/V_\omega.
\]
The proof of \cite[Lemma~2.3(1), p.~5, arXiv v3]{Teleman2006} identifies the image of $u\mapsto\mathrm i\Lambda_\omega\partial\bar\partial u$ with the functions of integral zero for a Gauduchon metric. Solve
\[
 \mathrm i\Lambda_\omega\partial\bar\partial u
       =c_P-\mathrm i\Lambda_\omega F_P.
\]
For the curvature convention $F=-\partial\bar\partial\log h$, replacing $h$ by $e^{-u}h$ adds $\partial\bar\partial u$ to the curvature, giving the stated constant.

\item For a holomorphic section $s$ of $E\otimes P$, the Chern-connection identity is
\[
 \mathrm i\Lambda_\omega\partial\bar\partial|s|^2
 =|\nabla^{1,0}s|^2
  -\langle\mathrm i\Lambda_\omega F_{E\otimes P}s,s\rangle.
\]
Differentiating the Hermitian norm in a holomorphic frame gives this pointwise Chern-connection identity. At a positive maximum of $|s|^2$, the left side is nonpositive, whereas the right side is at least
$-(C_E+c_P)|s|^2$.
If $c_P<-C_E$, this is strictly positive, a contradiction. We may take $b(E,\omega)=-C_EV_\omega/(2\pi)$, proving \eqref{cdp:eq:A.30}.

\item Apply this bound to $P=h^*(L^{\otimes n})$ for every positive and negative integer $n$. Nonzero sections for all $n$ are incompatible with a nonzero degree, proving the persistence assertion. For persistence on a single component, absorb one representative line into the fixed bundle $E$ and apply the same argument to all powers of every line in $\Pic^0(Y)$. If $H^2(Y,\Q)=0$, a positive power of every ambient line lies in $\Pic^0(Y)$, so additivity of degree gives degree zero for every pullback. Its real first Chern class is zero, even if its integral class is torsion. The Chern curvature is therefore exact. Apply the kernel-to-potential map of \cite[Lemma~2.3(1), p.~5, arXiv v3]{Teleman2006} and the explicit metric change in \cref{cdp:st:A.5.4} to obtain a flat unitary Chern connection. On a class VII surface the differential of the holomorphic Picard pullback then lies in a real one-dimensional unitary tangent space. Complex linearity makes that differential zero. The exponential sequence makes $\Pic(Y)$ a connected quotient of $H^1(Y,\OO_Y)$ when $H^2(Y,\Z)=0$, so under that integral hypothesis the entire map is identically zero.

\end{enumerate}
\end{proof}

\subsection{Ruled normalizations and conductor genus}\label{subsec:II:normalization-genus}

\begin{lemma}[Canonical boundaries under point blowups]\hypertarget{tag-P232}{}\label{lem:II:canonical-boundary-blowup}
\StatementRecord{lem:II:canonical-boundary-blowup}{proved}{P232}{}
Let $\sigma:Z_1\to Z_2$ be the blowup of a point on a smooth compact complex surface, with exceptional curve $E$. Let $D_1$ be an effective integral divisor with $c_1(K_{Z_1}+D_1)=0$ in rational cohomology. Put $D_2=\sigma_*D_1$, let $a$ be the coefficient of $E$ in $D_1$, and let $m$ be the multiplicity of $D_2$ at the centre. Writing canonical line bundles additively, their Jacobian comparison is
\begin{equation}
 K_{Z_1}+D_1=\sigma^*(K_{Z_2}+D_2)+(1+a-m)E.
                                                               \hypertarget{tag-006Y}{}\label{cdp:eq:A.RC.5}
\hypertarget{tag-0080}{}\hypertarget{tag-0073}{}\hypertarget{tag-0093}{}\hypertarget{tag-00DC}{}\hypertarget{tag-00OJ}{}\SsixEquationAliases{\label{cdp:eq:A.T.22}\label{cdp:eq:A.RT.11}\label{cdp:eq:C.22b}\label{cdp:eq:C.CR.5}\label{cdp:eq:C.IR.4}}
\end{equation}
One has $a=m-1$, $m\ge1$, and $K_{Z_2}+D_2$ also has zero rational Chern class. The support of $D_1$ is connected exactly when the support of $D_2$ is connected. The complement Euler numbers satisfy
\[
 e(Z_1\setminus(D_1)_{\rm red})-e(Z_2\setminus(D_2)_{\rm red})
 =\begin{cases}1,&m=1,\\0,&m\ge2.\end{cases}
\]
\end{lemma}
\begin{proof}
In a blowup chart $(u,v)\mapsto(u,uv)$, pullback of the local two-form multiplies it by $u$. These Jacobian maps glue to $K_{Z_1}\simeq\sigma^*K_{Z_2}\otimes\OO_{Z_1}(E)$. Pullback of a local equation of $D_2$ has order $m$ on $E$, while $D_1$ has order $a$, proving the line-bundle identity \eqref{cdp:eq:A.RC.5}. Intersecting with $E$, whose square is $-1$ and whose intersection with a pulled-back divisor is zero, gives $1+a-m=0$ in the numerically trivial case. The formula then identifies the rational class upstairs with the pullback of the class downstairs; injectivity of that pullback gives its vanishing downstairs.

Since $a\ge0$, the centre belongs to $D_2$. For $m=1$ it is a smooth point of one component of coefficient one, $a=0$, and the strict boundary maps isomorphically to the old boundary. For $m\ge2$, $a>0$ and $(D_1)_{\rm red}=\sigma^{-1}((D_2)_{\rm red})$; its proper map to the old boundary has connected fibres, with $\PP^1$ over the centre. These descriptions prove the connectedness assertion. For $m=1$ the surface Euler number increases by one and the boundary Euler number is unchanged. For $m\ge2$, $\sigma$ is an isomorphism on the two complements. This proves the Euler formula.
\end{proof}

\begin{lemma}[Topology of a normalized compact curve]
\hypertarget{tag-P245}{}\label{lem:II:curve-normalization-graph}
\StatementRecord{lem:II:curve-normalization-graph}{proved}{P245}{}
Let $C$ be a connected compact reduced complex curve, with normalization
$\coprod_{i=1}^c C_i^{\rm norm}\to C$. Write $g_i$ for the genus of
$C_i^{\rm norm}$. Form the incidence graph with one vertex for each
normalization component, one vertex for each singular point, and one
edge for each normalization preimage of a singular point. Let $b$ be
its cycle rank. Base points and paths along a spanning tree give
\[
 \pi_1(C)\simeq
 \Bigl(*_{i=1}^c\pi_1(C_i^{\rm norm})\Bigr)*F_b,
 \qquad e(C)=c+1-2\sum_i g_i-b.
\]
For a virtually abelian $\pi_1(C)$, the possibilities are all $g_i=0$
with $b\le1$, or exactly one $g_i=1$, all other $g_i=0$, and $b=0$.
Consequently $e(C)\ge0$, with equality precisely when $c=1$, $g_1=1$
and $b=0$. The group and Euler formulas apply separately to every
connected component of a disconnected compact reduced curve.
\end{lemma}
\begin{proof}
Join the finitely many points in each singular normalization fibre
by a new star. Collapsing each star identifies exactly that fibre
and gives the topological curve $C$. Compatible triangulations make
these contractible stars subcomplexes, so their collapse is a
homotopy equivalence by \cite[Proposition~0.17, p.~15]{Hatcher2002}.
A spanning tree of the incidence graph joins the normalized surfaces;
the remaining $b$ edges add $b$ circles. Van Kampen gives the displayed
free product with the indicated component maps. If $r_p$ is the number
of normalization preimages of a singular point, then
\[
 b=\sum_p r_p-c-\#\operatorname{Sing}(C)+1,
 \qquad e(C)=\sum_i(2-2g_i)-\sum_p(r_p-1),
\]
which gives the Euler formula. A unibranch singular point has $r_p=1$
and contributes zero to the subtraction.

A surface factor of genus at least two maps onto $F_2$: in the
presentation $\prod_j[a_j,b_j]=1$, send every $b_j$ and all but two
$a_j$ to the identity. Two nontrivial surface factors, two free
generators, or one elliptic factor and a free generator contain
$F_2$ by the free-product normal form. A finite-index subgroup of
index $d$ in $F_2$ is represented by a connected cover of the
two-circle bouquet with $d$ vertices and $2d$ edges, so has rank
$d+1\ge2$. Thus $F_2$ has no finite-index abelian subgroup. Every
subgroup and quotient of a virtually abelian group is virtually
abelian, as follows by intersecting with, or mapping, a finite-index
abelian subgroup. This proves the stated alternatives. Substitution
in the Euler formula proves its nonnegativity and equality assertion.
\end{proof}

\begin{lemma}[Fundamental groups under a canonical-boundary blowup]
\hypertarget{tag-P246}{}\label{lem:II:boundary-blowup-groups}
\StatementRecord{lem:II:boundary-blowup-groups}{proved}{P246}{\SeriesRef{lem:II:canonical-boundary-blowup}}
Use the blowup $\sigma:Z_1\to Z_2$ and divisors $D_1,D_2$ of
\cref{lem:II:canonical-boundary-blowup}. For each connected component
$C_2$ of $(D_2)_{\rm red}$, let $C_1$ be its inverse image in
$(D_1)_{\rm red}$. The maps $\sigma|_{C_1}:C_1\to C_2$ and
$\sigma:Z_1\to Z_2$ induce isomorphisms on fundamental groups.
With base points chosen away from the centre and its exceptional
curve, these isomorphisms commute with the homomorphisms induced
by $C_i\hookrightarrow Z_i$.
\end{lemma}
\begin{proof}
For a component avoiding the centre the boundary map is an
isomorphism. At multiplicity one, \cref{lem:II:canonical-boundary-blowup}
identifies the strict boundary isomorphically with the old boundary.
At multiplicity at least two its inverse image contains the exceptional
$E\simeq\mathbb P^1$, and $C_1/E\to C_2$ is a continuous bijection
from a compact space to a Hausdorff space. It is a homeomorphism.
A compatible triangulation makes $E\subset C_1$ a subcomplex;
van Kampen identifies the quotient group with
$\pi_1(C_1)/\langle\!\langle\pi_1(E)\rangle\!\rangle=\pi_1(C_1)$.
For the surfaces, apply \cref{lem:II:blowup-topology} to the point centre, of complex codimension two. Its inclusion and blowdown maps give the stated surface isomorphism. The original inclusions and
$\sigma$ commute as continuous maps, proving the last assertion.
\end{proof}

\begin{lemma}[A square-zero elliptic section and its conductor]\hypertarget{tag-P226}{}\label{lem:II:elliptic-ruled-boundary}
\StatementRecord{lem:II:elliptic-ruled-boundary}{proved}{P226}{\cref{lem:II:exceptional-intersection}}
Let $\mu:Z\to T$ be the minimal resolution of a normal compact complex surface, with an irrational ruling $p:Z\to B$ over an elliptic curve. Suppose that every $\mu$-exceptional curve is vertical. Let $C$ be a section with $C^2=0$, and let $V$ be an effective vertical divisor disjoint from $C$, with $c_1(K_Z+2C+V)=0$ in rational cohomology. Then $V=0$, $Z$ is a minimal ruled bundle, and $\mu$ is an isomorphism.
\end{lemma}
\begin{proof}
Contract the vertical ruling $(-1)$-curves to a bundle $\rho:Z\to Z_0$ over $B$, and put $C_0=\rho(C)$, $V_0=\rho_*V$. The image $C_0$ is a section. Its self-intersection is $C_0^2=C^2+n=n\ge0$, where $n$ counts the blowup centres on the successive section transforms. Pushing the canonical-boundary class forward and using adjunction on the elliptic section gives
\[
 0=(K_{Z_0}+2C_0+V_0)\cdot C_0
   =n+V_0\cdot C_0.
\]
The divisor $V_0$ is an effective sum of fibres, so $n=0$ and $V_0=0$. Thus $C=\rho^*C_0$ and $V$ is $\rho$-exceptional. In
$K_Z=\rho^*K_{Z_0}+E_\rho$, the divisor $E_\rho$ is effective and positive on every exceptional prime. The class of $E_\rho+V$ is numerically zero. The total transforms of the point-blowup exceptional curves form an orthogonal basis with square $-1$, so their intersection matrix is negative definite. Therefore $E_\rho+V=0$. Effectivity gives $E_\rho=V=0$, and $\rho$ is an isomorphism.

Every vertical irreducible curve of this bundle is a fibre and has square zero. By \cref{lem:II:exceptional-intersection}, an exceptional curve of $\mu$ belongs to a negative-definite block, so no such fibre is exceptional. The proper map $\mu$ has finite fibres and is bimeromorphic. It is therefore finite, and normality of $T$ makes it an isomorphism.
\end{proof}

\begin{proposition}[The ruled obstruction with a reduced conductor image]\hypertarget{tag-01BJ}{}\label{cdp:st:A.5.9}
\StatementRecord{cdp:st:A.5.9}{proved}{01BJ}{\SeriesRef{cdp:st:A.5.3}, \SeriesRef{cdp:st:A.1.1}, \SeriesRef{cdp:st:A.5.8}, \SeriesRef{cdp:st:A.5.2}, \SeriesRef{cdp:eq:A.RC.3}, \SeriesRef{cdp:eq:A.RC.4}, \SeriesRef{cdp:st:C2d}, \SeriesRef{cdp:eq:A.RC.5}, \SeriesRef{cdp:st:A.2.4}}

Let $S\subset Y$ be an irreducible nonnormal reduced Cartier
divisor in a compact smooth complex threefold with
$H^2(Y,\Z)=0$ and $\chi(\OO_S)=0$.
Let $\eta:T\to S$ be its normalization and
$\pi:Z\to T$ its minimal resolution, and put $h=\eta\pi$.
Suppose that $Z$ is irrational ruled, with ruling $p:Z\to B$,
that every exceptional curve of $\pi$ is vertical, and that
the effective divisor of \ref{cdp:st:A.5.3} has the form
\begin{equation}
 K_Z+D_Z=h^*\omega_S,\qquad D_Z=2C+V,                 \hypertarget{tag-006U}{}\label{cdp:eq:A.RC.1}
\end{equation}
where $C$ is a section in the strict conductor and $V$ is
vertical. Suppose also that the conductor image scheme
$\Gamma=\operatorname{Spec}_S(\OO_S/\mathfrak c)$ is reduced,
and that $H^0(S,K\otimes L|_S)\ne0$ for every $L\in\Pic(Y)$,
with $K$ as in \ref{cdp:st:A.1.1}.
Then every ambient line bundle pulls back trivially to $T$.
In particular
\begin{equation}
 H^2\bigl(T,\omega_T\otimes
       \eta^*((K_Y^{-1}\otimes L)|_S)\bigr)\simeq\C
       \quad\hbox{for every }L\in\Pic(Y).            \hypertarget{tag-006V}{}\label{cdp:eq:A.RC.2}
\end{equation}
The curves $B$ and $\Gamma$ are smooth elliptic curves,
$Z=T$ is minimal ruled, and $D_Z=2C$ with $C^2=0$.

\end{proposition}


\begin{proof}

\begin{enumerate}[label=\arabic*.,leftmargin=*]
\item Write $\mathcal B=\eta_*\OO_T$ and $\mathcal Q=\mathcal B/\OO_S$. Apply \cref{lem:II:conductor-duality} to this normalization. Equation~\eqref{cdp:eq:A.RC.3} identifies $\mathcal Q$ with $\omega_\Gamma\otimes\omega_S^{-1}|_\Gamma$ through the original inclusion $\OO_S\hookrightarrow\mathcal B$.
Adjunction makes $\omega_S$ the restriction of
$K_Y\otimes\OO_Y(S)$. Its degree is zero on every compact
curve, because $H^2(Y,\Z)=0$. Curve Riemann--Roch and Serre
duality therefore give
$\chi(\mathcal Q)=\chi(\omega_\Gamma)=-\chi(\OO_\Gamma)$.
On the other hand, \ref{cdp:st:A.5.8} gives rational singularities of $T$
and $\chi(\OO_T)=1-g(B)$. The normalization sequence and
$\chi(\OO_S)=0$ now give the exact numerical identity
\begin{equation}
                   \chi(\OO_\Gamma)=g(B)-1.          \hypertarget{tag-006X}{}\label{cdp:eq:A.RC.4}
\end{equation}

\item Contract ruling $(-1)$-curves from $Z$ to a minimal ruled bundle $Z_0\to B$. Every pushed-forward boundary is effective. Apply \cref{lem:II:canonical-boundary-blowup} at each contraction, with $D_1$ and $D_2$ the boundaries before and after contraction. It gives $a=m-1$, numerical triviality downstairs and equivalence of connectedness of their supports. On $Z_0$ the boundary is $2C_0+V_0$, with $C_0$ a section and $V_0$ an effective sum of fibres. Its support is connected, so the support of $D_Z$ is connected.

\item Every positive exceptional component of $D_Z$ maps into
$\Gamma$. Indeed away from $\Gamma$ the surface $T$ equals
the Gorenstein surface $S$ and has rational singularities by
\ref{cdp:st:A.5.8}. The local rational-duality and exceptional-matrix comparison in the proof of \cref{cdp:st:A.5.2}, applied to this Gorenstein germ and its minimal resolution, gives $\omega_Z=\pi^*\omega_T$ there. Thus \eqref{cdp:eq:A.RC.1} has no positive exceptional
coefficient there. The strict conductor maps onto the whole
conductor support. It follows that
$h(\operatorname{Supp}D_Z)=\operatorname{Supp}\Gamma$,
which proves connectedness.

\item Let $D\subset\Gamma$ be the image of $C$. Localize \eqref{cdp:eq:A.RC.3}
at the generic point of $D$. Because $\Gamma$ is reduced,
$\OO_\Gamma$ there is the function field $k(D)$, and its
dualizing module has dimension one over that field. The
elementary conductor quotient sequence is
\[
 0\longrightarrow\OO_\Gamma\longrightarrow
          \mathcal B/\mathfrak c
   \longrightarrow\mathcal Q\longrightarrow0.
\]
It follows that $\mathcal B/\mathfrak c$ has dimension two
over $k(D)$. The conductor coefficient two on $C$ contributes
$2\deg(C\to D)$ to this dimension. Therefore this degree is
one and there is no other normalization branch over $D$.
The generic transverse ring has conductor $t^2\C\{t\}$ and
is exactly $\C+t^2\C\{t\}=\C\{t^2,t^3\}$, the cusp ring.
Thus $D^\nu\simeq C\simeq B$. This deduction also rules out
any additional vertical branch over the same image curve.

\item A connected compact reduced curve has only constant global
holomorphic functions. From \eqref{cdp:eq:A.RC.4} we therefore obtain
$h^1(\Gamma,\OO_\Gamma)=2-g(B)$.
Restriction to the reduced component $D$ gives a surjection
on $H^1$, because its kernel is a coherent sheaf on a curve
and has no $H^2$. Normalization of $D$ gives
$h^1(D,\OO_D)\ge g(B)$. Hence
\begin{equation}
             2-g(B)\ge h^1(D,\OO_D)\ge g(B).         \hypertarget{tag-006Z}{}\label{cdp:eq:A.RC.6}
\end{equation}
As $g(B)\ge1$, all quantities equal one. The normalization
quotient of $D$ has length zero, so $D$ itself is smooth
elliptic. \Cref{cdp:st:C2d} applies to this smooth curve
in $Y$ and makes the transversal cusp constant along it.
No other singular curve meets $D$. Connectedness of $\Gamma$
therefore gives $\Gamma=D$. Its normalization inverse image
is the single smooth curve $C$, the normalization is smooth
near it, and \ref{cdp:st:C2d} gives $C^2=0$.

\item Adjunction in \eqref{cdp:eq:A.RC.1} gives $V\cdot C=0$. Since $V$ is effective and vertical, it is disjoint from the section $C$. Apply \cref{lem:II:elliptic-ruled-boundary} to the minimal resolution $\pi:Z\to T$, its ruling, $C$ and $V$. The preceding step gives the elliptic base and $C^2=0$, the statement assumes vertical exceptional curves, and \eqref{cdp:eq:A.RC.1} gives the zero canonical-boundary class. Hence $V=0$, $Z$ is a minimal ruled bundle and $\pi$ is an isomorphism.

\item The constant-cusp computation \ref{cdp:st:A.2.4} now gives the injection
$K\hookrightarrow\eta_*(h^*\OO_Y(-S)\otimes\OO_Z(C))$.
Every ambient line is topologically trivial on $Z$ and is
the pullback of a degree-zero line on the elliptic curve $B$.
The section exact sequence, together with $R^1p_*\OO_Z=0$,
gives
\begin{equation}
 0\longrightarrow\OO_B\longrightarrow p_*\OO_Z(C)
 \longrightarrow N_{C/Z}\longrightarrow0,
 \qquad\deg N_{C/Z}=0.                              \hypertarget{tag-0070}{}\label{cdp:eq:A.RC.7}
\end{equation}
After a degree-zero twist $P$, both end terms have zero
sections unless $P=\OO_B$ or $P=N_{C/Z}^{-1}$.
The image of $\Pic(Y)$ in $\Pic^0(B)$ is connected:
$H^2(Y,\Z)=0$ makes $\Pic(Y)$ the connected image of
$H^1(Y,\OO_Y)$ under the exponential map.
All-twist persistence of $K$, and translation by the ambient
line $\OO_Y(-S)$, put this entire image in those at most two
points. It is therefore the identity. This proves triviality
of every ambient pullback. Serre duality on $T$ proves
\eqref{cdp:eq:A.RC.2}.

\end{enumerate}
\end{proof}

\begin{proposition}[The genus of a horizontal conductor image from a local torus group]\hypertarget{tag-01B1}{}\label{cdp:st:A.5.11}
\StatementRecord{cdp:st:A.5.11}{proved}{01B1}{\SeriesRef{cdp:st:A.5.8}, \SeriesRef{cdp:eq:A.Pi.1}, \SeriesRef{cdp:eq:A.Pi.2}, \SeriesRef{cdp:eq:A.Pi.4}; \cite[\S3, Theorem~2, p.~464]{Lojasiewicz1964}}

Let $f:Y\to C_0$ be a proper holomorphic map from a smooth
compact complex threefold to a smooth curve, with connected
fibres and complex two-torus general fibre. Let $S$ be an
irreducible reduced fibre support, $\eta:T\to S$ its
normalization, and $\pi:Z\to T$ a resolution which is
irrational ruled, with ruling $p:Z\to B$.
Suppose every exceptional curve of $\pi$ is vertical, and
the strict conductor has just one horizontal irreducible
component $C$, which is a section of $p$.
Let $D=\eta\pi(C)$ and let $\Sigma=D^\nu$ be its normalization.
Then $g(\Sigma)\le1$.
In particular, when $C\to D$ has degree one, $B$ is elliptic.
This applies to the horizontal-cusp case of $D_Z=2C+V$;
it allows arbitrary multiplicity of the full fibre and
arbitrary singularities or nonreduced conductor image on
the finitely many vertical curves.
\end{proposition}


\begin{proof}

\begin{enumerate}[label=\arabic*.,leftmargin=*]
\item There is nothing to prove when $\Sigma$ is rational, so
assume $g(\Sigma)\ge1$.
\Cref{cdp:st:A.5.8} gives a holomorphic ruling
$\widetilde p:T\to B$ with $p=\widetilde p\pi$.
Identifying $C$ with $B$ by its section map, normalization
of its nonconstant map to $D$ gives a finite surjective
holomorphic map $\phi:B\to\Sigma$.
Put $\psi=\phi\widetilde p:T\to\Sigma$.

\item We describe precisely the finite identifications needed to
make this map descend to $S$. Every vertical conductor
component of $T$ lies over one of finitely many points of
$B$, so $\psi$ is constant on it. Its image under $\eta$
has rational normalization: a vertical component is the
image of a rational curve on $Z$, and a nonconstant map
from that curve to its image normalization is finite.
Such an image cannot equal $D$, because $g(\Sigma)\ge1$.
Consequently these vertical images meet $D$ at finitely
many points. Away from those points and from the finitely
many singular points of $D$, every normalization preimage
of a point of $D$ lies on $C$, and its $\psi$ value is the
unique corresponding point of $\Sigma$.
At a point outside $D$ with more than one normalization
preimage, all those preimages lie on vertical conductor
components. There are no other horizontal conductor
components by hypothesis. Thus there is a finite set
$A\subset\Sigma$ such that, for every $s\in S$, either
$\psi(\eta^{-1}(s))$ is a singleton or it is contained
in $A$. One can take $A$ to contain the constant values
on all vertical conductor components and all points of
$\Sigma$ above the singularities of $D$ and its
intersections with the vertical image curves.

\item Choose an embedded finite tree $J$ in the underlying
oriented surface $\Sigma$ containing $A$, and let
$q:\Sigma\to\Sigma/J$ collapse that tree to one point.
Such a tree is obtained by joining the finite set by
arcs and retaining a spanning tree after subdividing
their finitely many crossings. Triangulate the surface
with $J$ a subcomplex. The quotient map $q$ is a homotopy equivalence by \cite[Propositions~0.16--0.17, pp.~15--16]{Hatcher2002}, since $J$ is a contractible CW subcomplex. In particular
\begin{equation}
              q_*:\pi_1(\Sigma)\xrightarrow{\sim}
                         \pi_1(\Sigma/J).          \hypertarget{tag-006O}{}\label{cdp:eq:A.Pi.1}
\end{equation}
The map $q\psi$ is constant on every fibre of the finite
proper surjection $\eta$. A proper continuous surjection
between these compact Hausdorff spaces is a quotient
map, so there is a unique continuous map
\begin{equation}
 a:S\longrightarrow\Sigma/J,
 \qquad a\eta=q\psi.                              \hypertarget{tag-006P}{}\label{cdp:eq:A.Pi.2}
\end{equation}
Let $\delta:\Sigma\to D\subset S$ be the curve normalization
map. On its dense smooth locus, lifting through $C$ gives
$a\delta=q$. Continuity gives this equality everywhere.
Equation \eqref{cdp:eq:A.Pi.1} therefore makes $a_*$ surjective on
fundamental groups.

\item We next extend the map \eqref{cdp:eq:A.Pi.2} to a neighbourhood
of $S$ in $Y$. The simultaneous analytic triangulation
theorem of \cite[\S3, Theorem~2, p.~464]{Lojasiewicz1964}
triangulates the pair $(Y,S)$, since a complex analytic
subset is real semianalytic. Compactness makes the
subcomplex representing $S$ finite. After barycentric
subdivision it is full: vertices lying in it span a
face in that same subcomplex. Write $\lambda_v$ for
barycentric coordinates and let $L$ be its vertex set.
On the open neighbourhood where $\sum_{v\in L}\lambda_v>0$,
the formula
\begin{equation}
 r(x)=\frac{\sum_{v\in L}\lambda_v(x)v}
                  {\sum_{v\in L}\lambda_v(x)}       \hypertarget{tag-006Q}{}\label{cdp:eq:A.Pi.3}
\end{equation}
defines a continuous retraction to $S$ in triangulation
coordinates. Its value lies in the face spanned by the
$L$ vertices of the same simplex, so these formulas
agree across faces. This supplies an open
neighbourhood $N$ and a map $r:N\to S$ fixed on $S$.

\item Let $c=f(S)$. Properness allows a small disc $\Delta$
about $c$ with $U=f^{-1}(\Delta)\subset N$ and no other
critical value. The map $(ar)|_U:U\to\Sigma/J$ is
surjective on fundamental groups: its restriction to
$\delta(\Sigma)\subset S\subset U$ is $q$, whose
fundamental-group map is an isomorphism.
Set $U^*=U\setminus S=f^{-1}(\Delta^*)$.
The inclusion induces a surjection
$\pi_1(U^*)\twoheadrightarrow\pi_1(U)$.
Indeed a loop in the smooth six-manifold $U$ can be
perturbed relative to its base point to be transverse
to an analytic stratification of $S$; every stratum
has real codimension at least two, so the resulting
one-dimensional loop misses every stratum.

\item Over $\Delta^*$, the proper submersion is a smooth
torus bundle. The fundamental-group exact sequence
and $\pi_2(\Delta^*)=0$ give
\[
 1\longrightarrow\Z^4\longrightarrow\pi_1(U^*)
       \longrightarrow\Z\longrightarrow1.
\]
A lift of the generator of $\Z$ splits this group
extension, hence
\begin{equation}
 \pi_1(U^*)\simeq\Z^4\rtimes_T\Z
 \twoheadrightarrow\pi_1(U)
 \twoheadrightarrow\pi_1(\Sigma/J)
 \simeq\pi_1(\Sigma).                             \hypertarget{tag-006R}{}\label{cdp:eq:A.Pi.4}
\end{equation}
The first group is metabelian: its commutator subgroup
is contained in the abelian normal subgroup $\Z^4$.
Every quotient in \eqref{cdp:eq:A.Pi.4} is therefore metabelian.
A closed orientable surface of genus at least two
has a fundamental group surjecting onto the free
group on two generators: in its usual surface
presentation, retain two $a$ generators and send all
$b$ generators and the other $a$ generators to the
identity. That free group surjects onto $S_4$ by
sending its generators to $(1234)$ and $(12)$;
their conjugates supply the adjacent transpositions.
The commutator subgroup of $S_4$ is the nonabelian
group $A_4$, so this quotient is not metabelian.
This contradicts \eqref{cdp:eq:A.Pi.4}. Thus $g(\Sigma)=1$ in the
positive-genus case. When $C\to D$ has degree one,
$\Sigma\simeq C\simeq B$, which proves the final
assertion.

\end{enumerate}
\end{proof}

\begin{proposition}[The arithmetic genus of conductor images from coherent duality]\hypertarget{tag-01B2}{}\label{cdp:st:A.5.12}
\StatementRecord{cdp:st:A.5.12}{proved}{01B2}{\SeriesRef{cdp:eq:A.CD.2}, \SeriesRef{cdp:eq:A.CD.3}, \SeriesRef{cdp:eq:A.CD.1}}

Let $f:Y\to C_0$ be a proper connected-fibre map from a
smooth compact complex threefold with $H^2(Y,\Z)=0$ to
a smooth curve, and let its smooth general fibre have
trivial canonical bundle. This includes complex
two-tori and primary Kodaira surfaces.
Let $mS$ be a full fibre with irreducible
reduced nonnormal support $S$. Write $\eta:T\to S$ for
normalization and $\pi:Z\to T$ for a resolution.
If the pullback homomorphism
$\Pic(Y)\to\Pic(Z)$ is nonconstant, then every reduced
irreducible component $D$ of the conductor image scheme
has arithmetic genus at most one.
In particular every such component whose normalization
has positive genus is a smooth elliptic curve.
Here $\Gamma$ denotes the full conductor image scheme, and its components are arbitrary reduced irreducible curves.
\end{proposition}


\begin{proof}

\begin{enumerate}[label=\arabic*.,leftmargin=*]
\item The canonical line on the reduced support is determined by evaluation. The sheaf $f_*K_Y$ is torsion-free of
generic rank one on the smooth curve and hence is a
line bundle. Choose a frame near $c=f(S)$ and evaluate
it in $K_Y$. A coordinate $t$ on the base disc supplies
the conormal frame $dt$, and adjunction identifies
$K_Y|_F$ with $K_F$ after dividing by this frame.
Under that line-bundle isomorphism, the evaluated
section restricts on a general smooth fibre to a
nonzero canonical two-form, which is nowhere zero.
Its zero divisor is therefore
vertical. Shrinking the base disc to remove the other
vertical zero divisors, irreducibility of the support
gives an effective divisor $kS$, with $k\ge0$, and a
line-bundle comparison on $U=f^{-1}(\Delta)$:
\begin{equation}
 K_Y|_U\simeq\OO_U(kS),\qquad
 \omega_S\simeq\OO_Y((k+1)S)|_S,\qquad
 \omega_S^{\otimes m}\simeq\OO_S.                 \hypertarget{tag-0058}{}\label{cdp:eq:A.CD.1}
\end{equation}
The last isomorphism uses the local base equation whose
zero divisor is exactly the full fibre $mS$.


\item Hirzebruch--Riemann--Roch in
$0\to\OO_Y(-S)\to\OO_Y\to\OO_S\to0$ gives
$\chi(\OO_S)=0$: the rational Chern character of
$\OO_Y(-S)$ equals that of $\OO_Y$, since its first
Chern class is zero. A nonzero section of the torsion
line $\omega_S$ trivializes it. Indeed its $m$th power
is a nonzero holomorphic function on the compact
connected reduced space and hence is a nonzero constant.
Thus, with $e=h^0(S,\omega_S)$, we have
\begin{equation}
 e\in\{0,1\},\qquad
 e=1\Longleftrightarrow\omega_S\simeq\OO_S,
 \qquad h^1(S,\OO_S)=1+e.                         \hypertarget{tag-0059}{}\label{cdp:eq:A.CD.2}
\end{equation}
Here $h^0(\OO_S)=1$ and Serre duality identifies
$h^2(\OO_S)$ with $e$.

\item Let $\mathcal Q=\eta_*\OO_T/\OO_S$ and let
$\Gamma\subset S$ be the conductor image scheme.
Both $S$ and $T$ are connected and compact, so their
global holomorphic functions are the same constants.
The normalization sequence consequently gives an
injection and the indicated kernel:
\begin{equation}
 H^0(S,\mathcal Q)\hookrightarrow H^1(S,\OO_S),
 \qquad
 \operatorname{im}H^0(S,\mathcal Q)
 =\ker\bigl(H^1(S,\OO_S)\to H^1(T,\OO_T)\bigr).
                                                               \hypertarget{tag-005A}{}\label{cdp:eq:A.CD.3}
\end{equation}
Because $H^2(Y,\Z)=0$, $\Pic(Y)$ is the connected
exponential image of $H^1(Y,\OO_Y)$. A nonconstant
holomorphic group homomorphism out of it has nonzero
differential. The differential of the stated Picard
pullback is the composite
\[
 H^1(Y,\OO_Y)\longrightarrow H^1(S,\OO_S)
 \longrightarrow H^1(T,\OO_T)
 \longrightarrow H^1(Z,\OO_Z).
\]
Thus its middle arrow has positive rank. Equations
\eqref{cdp:eq:A.CD.2}--\eqref{cdp:eq:A.CD.3} imply $h^0(\mathcal Q)\le e$.

\item Apply \cref{lem:II:conductor-duality} to the full conductor scheme $\Gamma$. Its section-duality map \eqref{eq:II:conductor-section-duality} gives
\begin{equation}
 h^1(\Gamma,\omega_S|_\Gamma)
       =h^0(S,\mathcal Q)\le e.                   \hypertarget{tag-005B}{}\label{cdp:eq:A.CD.4}
\end{equation}
For any reduced irreducible component $D$ of $\Gamma$,
restriction of the line gives an exact quotient
$\omega_S|_\Gamma\to\omega_S|_D$. Its kernel is a
coherent sheaf on a one-dimensional compact space and
has no $H^2$, so the corresponding map on $H^1$ is
surjective. Hence
$h^1(D,\omega_S|_D)\le e$.

\item For $e=1$, \eqref{cdp:eq:A.CD.2} makes this precisely
$p_a(D)=h^1(D,\OO_D)\le1$.
For $e=0$, the same inequality makes
$h^1(D,\omega_S|_D)=0$.
The restriction $\omega_S|_D$ is torsion by \eqref{cdp:eq:A.CD.1}
and has degree zero. Curve Riemann--Roch now gives
\[
 0\le h^0(D,\omega_S|_D)
       =\chi(D,\omega_S|_D)=1-p_a(D),
\]
again proving $p_a(D)\le1$.
Finally, for a reduced irreducible curve the
normalization exact sequence gives
$p_a(D)=g(D^\nu)+\operatorname{length}
(\nu_*\OO_{D^\nu}/\OO_D)$.
Positive normalization genus therefore forces genus
one and zero normalization quotient. A normal complex
curve is smooth, which proves the final assertion.

\end{enumerate}
\end{proof}

\begin{proposition}[The conductor bound when the general fibre has $q=1$ and $p_g=0$]\hypertarget{tag-01B3}{}\label{cdp:st:A.5.13}
\StatementRecord{cdp:st:A.5.13}{proved}{01B3}{\SeriesRef{cdp:st:D.1.9}, \SeriesRef{cdp:st:A.5.12}}

Let $f:Y\to\PP^1$ be a proper connected-fibre map from
a smooth compact complex threefold with
$H^1(Y,\Z)=H^2(Y,\Z)=0$. Suppose its smooth general
fibre has $h^1(\OO)=1$ and $h^2(\OO)=0$.
Then
\begin{equation}
 R^1f_*\OO_Y=\OO_{\PP^1},\qquad R^2f_*\OO_Y=0,
 \qquad\Pic(Y)=\C\beta,                            \hypertarget{tag-006S}{}\label{cdp:eq:A.Q1.1}
\end{equation}
where the last identification uses the exponential
coordinates of \ref{cdp:st:D.1.9} and its nonzero transgression
class $\beta$. Every full fibre has structure-sheaf
cohomology dimensions $(1,1,0)$.
For an irreducible reduced support $S$, every
nonconstant ambient Picard pullback to a normalization
resolution forces every reduced conductor component
to have arithmetic genus at most one. In particular
each positive-genus component is smooth elliptic.
\end{proposition}


\begin{proof}

\begin{enumerate}[label=\arabic*.,leftmargin=*]
\item By \cite[Theorem~2.2, pp.~4--5]{Viaclovsky2026Combs}, $E=R^1f_*\OO_Y$ is a line bundle. The section
induced by \ref{cdp:st:D.1.9}'s integral transgression class is
nowhere zero: its image under the natural base-change
map at each point is the nonzero coefficient class on
that full fibre. A zero of the section in $E$ would
have zero image under that map. Thus $E\simeq\OO$.
Coherent Leray now gives $H^1(Y,\OO_Y)=\C\beta$.
The exponential sequence and the two stated integral
cohomology vanishings give the Picard identification
in \eqref{cdp:eq:A.Q1.1}.

\item The coherent sheaf $G=R^2f_*\OO_Y$ has generic rank
zero, hence is a torsion sheaf. Riemann--Roch on $Y$
gives $\chi(\OO_Y)=\int_Yc_1c_2/24=0$, since
$c_1(Y)$ is zero in rational cohomology. The Euler
characteristic of coherent Leray is therefore
\[
 0=\chi(\OO_{\PP^1})-\chi(E)+\chi(G)
      =\operatorname{length}G.
\]
Thus $G=0$. Top-degree base change gives
$H^2(F,\OO_F)=0$ for every full Cartier fibre $F$.
The multiplication-by-parameter sequence then gives
$H^1(F,\OO_F)=E\otimes\C(c)$, of dimension one;
\cite[Theorem~2.2, pp.~4--5]{Viaclovsky2026Combs} gives $h^0(F,\OO_F)=1$.

\item For a full fibre $F=mS$ with irreducible reduced
support, the quotient $\OO_F\to\OO_S$ induces a
surjection on $H^2$, because its kernel has support
of dimension two and has no $H^3$.
Hence $h^2(S,\OO_S)=0$. Riemann--Roch in the Cartier
sequence on $Y$ gives $\chi(\OO_S)=0$ as in \ref{cdp:st:A.5.12},
so $h^1(S,\OO_S)=1$.

\item Put $\mathcal Q=\eta_*\OO_T/\OO_S$ for the
normalization $T$ and write $\Gamma$ for the conductor
image scheme. Nonconstant Picard pullback has nonzero
differential factoring through
$H^1(S,\OO_S)\to H^1(T,\OO_T)$.
The source has dimension one, so this map has zero
kernel. The normalization sequence \eqref{cdp:eq:A.CD.3} gives
$H^0(\mathcal Q)=0$. Conductor duality and Serre
duality give
\begin{equation}
 H^1(\Gamma,\omega_S|_\Gamma)=0,
 \qquad H^1(D,\omega_S|_D)=0                      \hypertarget{tag-006T}{}\label{cdp:eq:A.Q1.2}
\end{equation}
for each reduced irreducible component $D$, where
the second map is the restriction surjection in the
proof of \ref{cdp:st:A.5.12}. Adjunction makes $\omega_S$ ambient,
so $\deg(\omega_S|_D)=0$ by $H^2(Y,\Z)=0$.
Curve Riemann--Roch yields
$0\le h^0(D,\omega_S|_D)=1-p_a(D)$.
The normalization genus and its nonnegative delta
length then give the asserted smooth elliptic
conclusion whenever that genus is positive.

\end{enumerate}
\end{proof}

\Needspace{12\baselineskip}
\section{Conductor jets and binomial torus exclusions}\label{sec:II:conductor-binomial-exclusions}

\subsection{Conductor jets and the genus bound}\label{subsec:II:normalization-jets}

\begin{proposition}[Conductor-cleared jets on a complete-intersection critical curve]\hypertarget{tag-01BF}{}\label{cdp:st:A.5.5}
\StatementRecord{cdp:st:A.5.5}{proved}{01BF}{\SeriesRef{cdp:eq:A.J.4}, \SeriesRef{cdp:eq:A.J.5}, \SeriesRef{cdp:eq:A.J.2}, \SeriesRef{cdp:eq:A.J.3}, \SeriesRef{cdp:eq:A.J.7}, \SeriesRef{cdp:eq:A.J.9}, \SeriesRef{cdp:eq:A.J.8}, \SeriesRef{cdp:eq:A.J.10}}

Let $Y$ be a smooth compact complex threefold with
$H^2(Y,\mathbb Q)=0$, let $S\subset Y$ be a reduced divisor,
and let $D\subset\operatorname{Sing}(S)$ be a reduced
irreducible compact curve which is a local complete intersection
in $Y$. Assume that the generic plane-curve transversal to $D$
in $S$ is an ordinary cusp, an ordinary tacnode, or an ordinary
triple point.
If the normalization $C\to D$ has positive genus, then $D$ is
a smooth elliptic curve. The transversal type is the indicated
cusp, tacnode, or triple point at every point of $D$.

In the cusp case, the intrinsic cubic on the radical of the Hessian
extends to a meromorphic section on $C$ whose pole divisor is bounded
by twice the conductor. The tacnode fourth-order invariant extends
with poles bounded by four times the conductor.
\end{proposition}


\begin{proof}

\begin{enumerate}[label=\arabic*.,leftmargin=*]
\item Let $\nu:C\to D$ be the normalization and $j:C\to Y$
its composite with the inclusion,
$I=\mathcal I_{D/Y}$, and
\[
 N=\bigl(\nu^*(I/I^2)\bigr)^*,\qquad
 M=j^*\mathcal O_Y(S).
\]
These are a rank-two vector bundle and a line bundle on $C$.
Let $\Xi$ be the conductor divisor on $C$, so that
$\nu_*\mathcal O_C(-\Xi)$ is the conductor ideal in
$\mathcal O_D$, and put $c=\deg\Xi\geq0$.
Adjunction for the complete-intersection curve and finite
duality give exact line-bundle comparisons
\begin{equation}
 \det N=\nu^*\omega_D\otimes j^*K_Y^{-1},\qquad
 \omega_C=\nu^*\omega_D\otimes\mathcal O_C(-\Xi).
                                                        \hypertarget{tag-005C}{}\label{cdp:eq:A.J.1}
\end{equation}
For the second comparison one can work directly in the total
quotient ring: $\operatorname{Hom}_{\mathcal O_D}
(\nu_*\mathcal O_C,\mathcal O_D)$ consists precisely of
elements multiplying the normalization into $\mathcal O_D$,
which is its conductor. Since $D$ is Gorenstein, tensoring
this identity by its invertible dualizing sheaf is finite
duality. Both $M$ and $j^*K_Y$ have degree zero, by the
stated rational cohomology hypothesis. Consequently
\begin{equation}
 \deg\det N=2g(C)-2+c.                              \hypertarget{tag-005N}{}\label{cdp:eq:A.J.2}
\end{equation}
This is the normal bundle of the complete-intersection ideal,
with its specified normalization pullback.

\item For an ordinary cusp, choose a local defining section
$f$ of $S$, using a local frame of $\mathcal O_Y(S)$.
Since $D\subset\operatorname{Sing}(S)$, its conormal class
maps to zero under
$I/I^2\to\Omega_Y^1|_D$. This homomorphism is injective:
its source is locally free on the reduced curve and the map
is injective at every generic point. Therefore $f\in I^2$.
The associated-graded identity for a regular sequence gives
$I^2/I^3=\operatorname{Sym}^2(I/I^2)$. The class of $f$
is thus a quadratic form
\[
 q\in H^0\bigl(C,\operatorname{Sym}^2N^*\otimes M\bigr).
\]
It has rank one generically. Let $R\subset N$ be its
saturated radical and $Q=N/R$. Both are line bundles.
The quadratic form annihilates $R$ everywhere by holomorphic
continuation, and gives a nonzero section of $M\otimes Q^{-2}$.
In particular
\begin{equation}
 \deg Q\leq0.                                      \hypertarget{tag-005Y}{}\label{cdp:eq:A.J.3}
\end{equation}

\item On $D_{\rm reg}$, the third-order term on the radical is
intrinsic. Indeed, after writing the quadratic form as
$a\ell^2$, choose holomorphic lifts $\widehat a$ and $\widehat\ell$
and subtract $\widehat a\,\widehat\ell^2$ from $f$.
Restrict the resulting cubic class to $R^3$. Changing these two lifts
adds cubic terms divisible by $\ell$, which vanish on $R^3$.
Changing normal generators has the same effect. Its nonzero
value characterizes the ordinary-cusp transversal among the
rank-one quadratic cases. Denote this cubic by $k$.

\item At a point of $D$, choose two generators $X,Y$ of $I$. After a constant invertible linear
change of these generators, the normalized linear factor of
$q$ has the form $X+hY$ on every normalization branch over
that point. Such a choice exists: the saturated quotient
$N\to Q$ is surjective at each of the finitely many points,
so a constant vector can be chosen outside their finitely
many kernels. Thus $h$ is holomorphic on the normalization,
and
\[
 q=a(X+hY)^2,
\]
where $a=q_{XX}$ is the pullback of a coefficient
in $\mathcal O_D$. Choose a local conductor generator $d$
on the normalization. Both $d$ and $dh$ belong to
$\mathcal O_D$, since the conductor is an ideal in the
normalization contained in that ring. Choose lifts
$\widehat d,\widehat{dh},\widehat a$ to $\mathcal O_Y$,
and set
\begin{equation}
 G=\widehat d^{\,2} f-
    \widehat a\bigl(\widehat d X+\widehat{dh}Y\bigr)^2.
                                                        \hypertarget{tag-006C}{}\label{cdp:eq:A.J.4}
\end{equation}
The quadratic class of $G$ is zero after normalization.
It is already zero on $D$: the map from the locally free
sheaf $\operatorname{Sym}^2(I/I^2)$ to its normalization
pullback is injective because $D$ is reduced. Hence
$G\in I^3$.

\item Its class in $I^3/I^4=\operatorname{Sym}^3(I/I^2)$,
pulled back to $C$ and evaluated on $R^3$, is holomorphic.
Dividing that class by $d^2$ gives $k$ on the dense regular
locus. To check this equality, there $d$ is a unit and
\eqref{cdp:eq:A.J.4} is exactly the subtraction of the square used in the
definition of the radical cubic. Different choices of lifts
add cubic terms containing its linear factor, and therefore
have zero restriction to $R^3$. Thus these local expressions
are the same meromorphic section wherever they overlap.
Since $d$ has divisor $\Xi$, the resulting bound is
\begin{equation}
 0\ne k\in H^0\bigl(C,M\otimes R^{-3}
                               \otimes\mathcal O_C(2\Xi)\bigr).
                                                        \hypertarget{tag-006E}{}\label{cdp:eq:A.J.5}
\end{equation}

\item Taking degrees in \eqref{cdp:eq:A.J.5} gives $\deg R\leq2c/3$.
Together with \eqref{cdp:eq:A.J.2}--\eqref{cdp:eq:A.J.3}, this yields
\begin{equation}
 2g(C)-2+c=\deg R+\deg Q\leq\frac{2c}{3},
 \qquad 2g(C)-2+\frac c3\leq0.                     \hypertarget{tag-006F}{}\label{cdp:eq:A.J.6}
\end{equation}
For $g(C)\geq1$, this forces $g(C)=1$ and $c=0$.
Conductor zero means $\mathcal O_D=\nu_*\mathcal O_C$;
a normal complex curve is smooth. Hence $D=C$ is elliptic.
Now \eqref{cdp:eq:A.J.2}--\eqref{cdp:eq:A.J.5} give $\deg R=\deg Q=0$.
The Hessian coefficient and radical cubic are nonzero
sections of degree-zero line bundles on this elliptic curve,
so neither has a zero. Locally the holomorphic Morse lemma
in the nondegenerate normal direction, with the curve
coordinate as parameter, writes $f=y^2+g(x,z)$.
The conditions just proved give $g=x^3u(x,z)$ with $u$ a
holomorphic unit. Replacing $x$ by a local cube root of $u$
times $x$ gives the local equation $y^2+x^3=0$.

\item In the tacnode case, put $P=\widehat d X+\widehat{dh}Y$ and
$G=\widehat d^{\,2}f-\widehat aP^2\in I^3$.
The generic radical cubic is now zero, so its continuation
$[G]_3|_{R^3}$ is zero everywhere on the normalization.
Since $\ell=X+hY$ is monic in $X$, polynomial division gives
a holomorphic quadratic $C_2$ on the normalization
such that
\[
 [G]_3=\ell C_2.
\]
Each coefficient of $dC_2$ is in the conductor and hence in
$\mathcal O_D$. Lift those coefficients and let $H$ be the
resulting homogeneous quadratic in $X,Y$. Set
\[
 K=\widehat d^{\,2}G-PH.
\]
Its cubic class pulls back to
$d^2\ell C_2-d\ell\cdot dC_2=0$.
The same locally free graded-piece injection used above gives
$K\in I^4$. There is an exact identity in the local ring of
the ambient smooth threefold:
\begin{equation}
 4\widehat a\widehat d^{\,6}f
 -(2\widehat a\widehat d^{\,2}P+H)^2
 =4\widehat a\widehat d^{\,2}K-H^2.               \hypertarget{tag-006G}{}\label{cdp:eq:A.J.7}
\end{equation}
On the dense smooth locus, complete the square in the
nondegenerate normal direction. If the relevant terms there
are $a\ell^2+b\ell r^2+er^4$, the fourth-order invariant is
$\Delta=ae-b^2/4$. A substitution
$\ell\mapsto\ell+ur^2$ leaves this expression unchanged;
terms divisible by $\ell^2$ and higher changes in the radical
coordinate do not alter it. Thus it is the intrinsic section
of $M^2\otimes Q^{-2}\otimes R^{-4}$ on that dense locus.
It is nonzero precisely for the ordinary-tacnode transversal
among the rank-one Hessian cases with zero radical cubic.

\item Divide \eqref{cdp:eq:A.J.7} by $4\widehat d^{\,6}$ and take the
fourth-order class on $R^4$. Since
$[H]_2=dC_2$, this expresses exactly that same invariant as
\begin{equation}
 \Delta=
 \frac{a[K]_4|_{R^4}-(C_2|_{R^2})^2/4}{d^4}.
                                                        \hypertarget{tag-006H}{}\label{cdp:eq:A.J.8}
\end{equation}
Its numerator is holomorphic on the normalization.
Different choices agree on the dense regular locus and hence
give one meromorphic section. Therefore
\begin{equation}
 0\ne\Delta\in H^0\bigl(C,M^2\otimes Q^{-2}
                    \otimes R^{-4}\otimes\mathcal O_C(4\Xi)\bigr).
                                                        \hypertarget{tag-006I}{}\label{cdp:eq:A.J.9}
\end{equation}
Write $q_0=\deg Q\leq0$. Using
$\deg R=2g(C)-2+c-q_0$, the degree of the line bundle
in \eqref{cdp:eq:A.J.9} is
\[
 -2q_0-4\bigl(2g(C)-2+c-q_0\bigr)+4c
       =-8g(C)+8+2q_0.
\]
Its nonzero section and $g(C)\geq1$ force
$g(C)=1$ and $q_0=0$.

\item The section $q$ of $M\otimes Q^{-2}$ is now nowhere zero.
In the local presentation above, $a=q_{XX}$ is therefore a
unit on every normalization branch. Since it belongs to
$\mathcal O_D$, it is a unit in that local ring as well.
The coefficient $h=q_{XY}/a$ consequently belongs to
$\mathcal O_D$; here $q_{XY}$ denotes the coefficient in
the symmetric matrix, so the cross term is $2q_{XY}XY$.
We can thus repeat the entire construction with $d=1$.
Polynomial division by the monic $X+hY$ then takes place
over $\mathcal O_D$ itself; the zero radical cubic descends
by the same injectivity argument. Formula \eqref{cdp:eq:A.J.8} now has
no poles. The resulting nonzero section of
$M^2\otimes Q^{-2}\otimes R^{-4}$ has degree $-4c$,
which forces $c=0$.

\item Thus again $D=C$ is smooth elliptic. Both the nonzero
quadratic coefficient and $\Delta$ are sections of
degree-zero line bundles and are nowhere zero. The parametric
Morse lemma first gives $f=y^2+g(x,z)$, and the vanished
radical cubic and nonzero fourth invariant give
$g=x^4u(x,z)$ with $u$ a unit. Taking a local fourth root
produces the equation $y^2+x^4=0$ at every point of $D$.

\item For an ordinary triple transversal, the quadratic class
vanishes generically, hence vanishes in the locally free
sheaf $I^2/I^3$. Thus $f\in I^3$, and its initial cubic
is a global section
$b\in H^0(C,\operatorname{Sym}^3N^*\otimes M)$.
The discriminant of a binary cubic is homogeneous of degree
four in its coefficients and transforms by the sixth power
of the determinant of a change of variables. Its value is
therefore a section
\begin{equation}
 0\ne\operatorname{disc}(b)\in
 H^0\bigl(C,M^4\otimes(\det N)^{-6}\bigr).
                                                        \hypertarget{tag-005D}{}\label{cdp:eq:A.J.10}
\end{equation}
It is nonzero because the three generic tangent lines are
distinct. Equations \eqref{cdp:eq:A.J.2} and \eqref{cdp:eq:A.J.10} give
$2g(C)-2+c\leq0$, again forcing $g(C)=1$ and $c=0$
under the genus assumption. The discriminant has degree
zero and is nowhere zero. The three simple tangent roots
extend holomorphically on each small curve disc. The implicit
function theorem gives three smooth local branches of $S$
with those pairwise distinct normal tangents. Choose the
first two branches as the coordinate axes; the third is
$y=xu(x,z)$ with $u$ a unit, and replacing $y$ by $y/u(x,z)$
straightens it to $y=x$. This gives the local equation
$xy(x-y)=0$ up to a unit at every point of $D$.


\end{enumerate}
\end{proof}

\begin{example}[A conductor pole on a singular critical curve]\hypertarget{tag-P224}{}\label{ex:II:conductor-pole}
\StatementRecord{ex:II:conductor-pole}{proved}{P224}{}
Put $h=z^2-x^3$, $D=(h,w)=0$ in $\mathbb C^3$, and
\[
 f=xw^2+2zwh+x^2h^2+h^3.
\]
Its radical cubic on the normalization of $D$ has a pole of order two, and multiplication by the square of the conductor generator gives the holomorphic jet in \eqref{cdp:eq:A.J.4}.
\end{example}
\begin{proof}
On the normalization $x=s^2,z=s^3$, the quadratic class is $s^2(w+sh)^2$. Completing the square on $D_{\rm reg}$ gives
\[
 f=x\left(w+\frac zx h\right)^2+
                   \left(1-\frac1x\right)h^3.
\]
Thus the radical cubic has a pole of order two at $s=0$. The conductor generator is $d=x=s^2$, and
\[
 x^2f-x(xw+zh)^2=x(x-1)h^3\in I_D^3.
\]
This is the conductor-cleared cubic of \eqref{cdp:eq:A.J.4} for the specified map and generator.
\end{proof}

\begin{proposition}[The genus bound and delta-one curve germs]\hypertarget{tag-01B4}{}\label{cdp:st:A.5.14}
\StatementRecord{cdp:st:A.5.14}{proved}{01B4}{\SeriesRef{cdp:eq:A.J.31}, \SeriesRef{cdp:st:A.5.5}}

Let $f:Y\to C_0$ be a proper connected-fibre map from a
smooth compact complex threefold with $H^2(Y,\Q)=0$ to
a smooth curve. Suppose its smooth general fibre has
trivial canonical bundle. Let $mS$ be a full fibre with
irreducible reduced nonnormal support, and let $D$ be
a reduced irreducible component of its conductor image.
Then
\begin{equation}
                         p_a(D)\le2.                \hypertarget{tag-005Z}{}\label{cdp:eq:A.J.30}
\end{equation}
If the normalization of $D$ has positive genus, then
$D$ is a local complete intersection in $Y$; every
singular germ of $D$ is an ordinary node or an ordinary
cusp, and the sum of its delta invariants is at most
one. In particular, when the generic transverse type
along such a positive-genus $D$ is $A_2$, $A_3$, or an
ordinary triple point, $D$ is smooth elliptic and the
transverse analytic type is constant along all of $D$.
\end{proposition}


\begin{proof}

\begin{enumerate}[label=\arabic*.,leftmargin=*]
\item Restrict $f$ to a disc about the image of $mS$ and trivialize the
canonical line of the disc. The sheaf $f_*K_Y$ is torsion-free of rank
one, hence locally free there. The evaluation of a generator is
nonzero on the smooth nearby fibres, where their canonical lines are
trivial. Its zero divisor is therefore $aS$ for an integer $a\ge0$.
Thus $K_Y\simeq\OO_Y(aS)$ on this neighbourhood. The full-fibre equation
trivializes $\OO_Y(mS)$, and adjunction gives
$\omega_S\simeq\OO_Y((a+1)S)|_S$, a torsion line.
The Cartier sequence and Riemann--Roch on $Y$ give $\chi(\OO_S)=0$,
since $c_1(\OO_Y(S))=0$ in rational cohomology. A nonzero section of a
torsion line on the compact irreducible surface has a nonzero constant
tensor power; it therefore trivializes the line. Consequently
$e=h^0(S,\omega_S)\in\{0,1\}$, and Serre duality gives
$h^1(\OO_S)=1+e$. For $e=1$, the line $\omega_S$ is trivial.

\item Write $\eta:T\to S$ for normalization,
$\mathcal Q=\eta_*\OO_T/\OO_S$, and $\Gamma$ for
the full conductor image scheme. Equality of the
constant global functions on $S$ and $T$ gives the
injection $H^0(\mathcal Q)\hookrightarrow
H^1(\OO_S)$. The section-duality isomorphism \eqref{eq:II:conductor-section-duality}, for this $\eta$ and $\Gamma$, identifies $H^0(S,\mathcal Q)^\vee$ with $H^1(\Gamma,\omega_S|_\Gamma)$. Taking dimensions gives
\begin{equation}
 h^1(\Gamma,\omega_S|_\Gamma)
       =h^0(S,\mathcal Q)\le1+e.                  \hypertarget{tag-0060}{}\label{cdp:eq:A.J.31}
\end{equation}
The restriction quotient
$\omega_S|_\Gamma\twoheadrightarrow\omega_S|_D$
has a kernel supported in dimension at most one, so
its $H^1$ map is surjective. For $e=1$ it follows that
$p_a(D)=h^1(\OO_D)\le2$. For $e=0$, put
$L=\omega_S|_D$. This torsion line has degree zero,
and \eqref{cdp:eq:A.J.31} gives $h^1(D,L)\le1$. Curve
Riemann--Roch gives
\[
 p_a(D)=1+h^1(D,L)-h^0(D,L)\le2,
\]
proving \eqref{cdp:eq:A.J.30} in both cases.

\item Let $\nu:\widehat D\to D$ be normalization. The
normalization sequence gives the identity
\begin{equation}
 p_a(D)=g(\widehat D)+
        \sum_{x\in D}\dim_{\C}
          (\nu_*\OO_{\widehat D}/\OO_D)_x.       \hypertarget{tag-0061}{}\label{cdp:eq:A.J.32}
\end{equation}
If $g(\widehat D)>0$, every local delta is consequently
at most one and their sum is at most one. We give the
analytic local rings at a point with delta one.
Let $A=\OO_{D,x}$ and let
$B=\bigoplus_{i=1}^r\C\{t_i\}$ be its finite
normalization. Here $\dim_{\C}(B/A)=1$. The residue map
of the local ring $A$ has the same value on every
normalization branch. The resulting quotient of
$\C^r$ by its diagonal has dimension $r-1$, so
$r\le2$.

\item For $r=2$, the subring
\[
 A_0=\{(b_1,b_2)\in\C\{t_1\}\oplus\C\{t_2\}:
                                  b_1(0)=b_2(0)\}
\]
contains $A$ and has colength one in $B$. Hence
$A=A_0$. The explicit homomorphism
$\C\{u,v\}/(uv)\to A_0$ sends
$u$ to $(t_1,0)$, $v$ to $(0,t_2)$, and a constant
to the corresponding equal pair. Expanding a
convergent series shows that it is an isomorphism.
This is the ordinary node germ.

\item For $r=1$, an element of $A$ with valuation one would
be a local coordinate on the normalization disc by
the holomorphic inverse function theorem. Convergent
composition in the analytic algebra $A$ would then
give $\C\{t\}\subset A$, contrary to its colength
one. Therefore every nonconstant element of $A$ has
valuation at least two. We have
$A\subset\C+t^2\C\{t\}$; both rings have
colength one in $\C\{t\}$, so equality holds.
The map
\[
 \C\{u,v\}/(v^2-u^3)\longrightarrow A,
                 \qquad (u,v)\longmapsto(t^2,t^3)
\]
is an isomorphism: every series with no linear term
is uniquely the sum of an even series in $t^2$ and
$t^3$ times a series in $t^2$. This is the ordinary
cusp germ. A delta-zero germ is its smooth
normalization disc.

\item We finally check the embedding into the threefold $Y$. These analytic rings have embedding
dimension at most two. In local coordinates of the
smooth threefold, the tangent space of $D$ therefore
has codimension at least one. There is an element
$z$ of its ideal with $dz(x)\ne0$. The inverse
function theorem makes $z$ an ambient coordinate,
and the ideal of $D$ is the inverse image of a
height-one ideal in the regular analytic ring
$\C\{u,v\}=\OO_{Y,x}/(z)$. This ideal is principal:
a pure reduced curve in a smooth surface is a Cartier
divisor, or equivalently its reduced height-one ideal
is a product of the prime factors in this factorial
ring. Thus $I_{D,x}=(z,h)$, a regular sequence.
This proves the asserted lci embedding, including
at the node and cusp. Proposition~\ref{cdp:st:A.5.5} applies
to the three specified generic transverse types
and proves the final smooth elliptic and constant-type
conclusion.

\end{enumerate}
\end{proof}

\subsection{Cusp images on ruled normalizations}\label{subsec:II:normalization-cusps}

\begin{proposition}[A pure double conductor forces a smooth cusp image]\hypertarget{tag-01BH}{}\label{cdp:st:A.5.7}
\StatementRecord{cdp:st:A.5.7}{proved}{01BH}{\SeriesRef{cdp:eq:A.J.13}, \SeriesRef{cdp:eq:A.J.14}, \SeriesRef{cdp:eq:A.J.12}}

Let $S$ be a reduced Gorenstein complex surface and
$\pi:\overline S\to S$ its finite normalization.
Suppose that, over a neighbourhood of a point $p\in S$,
$\overline S$ is smooth and its conductor ideal is exactly
$\mathcal O_{\overline S}(-2C)$ for a smooth reduced
divisor $C$; its local branches over $p$ are disjoint.
Suppose that the generic transversal of $S$ along the
reduced image $D=\pi(C)$ is an ordinary cusp.
Then $D$ is smooth at $p$, $C\to D$ is an isomorphism
near the preimage of $p$, and the germ of $S$ at $p$ is
analytically a product of an ordinary cusp and a smooth disc.

\end{proposition}


\begin{proof}
The graph of a rational derivation on a square-zero conductor extension is the comparison of \cite[\S3.9, pp.~716--717; \S4.4--4.8, pp.~720--724]{Reid1994Nonnormal}. The following local maps identify that graph and its canonical module in the analytic rings of the statement.


\begin{enumerate}[label=\arabic*.,leftmargin=*]
\item Write $R=\mathcal O_{S,p}$, let $\overline R$ be its finite semilocal normalization, and put
\[
 A=R/\mathfrak c,\qquad \overline A=\overline R/\mathfrak c,
 \qquad Q=\overline R/R=\overline A/A.
\]
The normalization is a product of regular two-dimensional local analytic rings. The local maps in \cref{lem:II:conductor-duality} give $Q\simeq\omega_A$ and $\operatorname{End}_A(Q)=A$ by \eqref{cdp:eq:A.J.11}--\eqref{cdp:eq:A.J.12}.

\item The ordinary-cusp hypothesis says that $A$ is generically
reduced. Its nilradical is therefore supported at finitely
many points; a one-dimensional Cohen--Macaulay ring has
no nonzero submodule with that support. Hence $A$ is reduced.
Choose local coordinates $(s_i,z_i)$ on the normalization
branches with $C=(z_i=0)$. Put
\[
 B=\prod_i\mathbb C\{s_i\},\qquad
 K=\prod_i\operatorname{Frac}\bigl(\mathbb C\{s_i\}\bigr).
\]
The specified conductor gives the ring comparison
\begin{equation}
 \overline A=B[\epsilon]/(\epsilon^2),             \hypertarget{tag-005G}{}\label{cdp:eq:A.J.13}
\end{equation}
component by component. The projection
$\overline A\to B$ is injective on $A$, since its kernel
is nilpotent and $A$ is reduced. It identifies $B$ with
the finite normalization of the curve ring $A$.

\item Under \eqref{cdp:eq:A.J.13} the inclusion of $A$ takes the form
\[
 a\longmapsto a+\delta(a)\epsilon,
\]
where $\delta:A\to B$ is a complex-linear analytic
derivation. It extends to the total quotient ring, and on
each normalization branch has the form
$\delta(a)=u_i\,da/ds_i$ for a meromorphic germ $u_i$.
Let $u=(u_i)$ and write derivatives componentwise.
The map
\begin{equation}
 \Phi:B[\epsilon]/(\epsilon^2)\longrightarrow K,
 \qquad b+c\epsilon\longmapsto c-u b'
                                                        \hypertarget{tag-005H}{}\label{cdp:eq:A.J.14}
\end{equation}
is $A$-linear for this specified action: multiplication by
$a+u a'\epsilon$ changes its value to
$ac+u a'b-u(ab)'=a(c-u b')$. Its kernel at every generic
point is exactly the graph of the extended derivation on
$K$. It therefore induces a generic isomorphism from $Q$
to its image in $K$. The module $Q\simeq\omega_A$ is
torsion-free on the reduced curve, so this induced map is
injective.

\item Differentiation of convergent power-series germs onto
convergent power-series germs is surjective, by termwise
integration. Hence the image in \eqref{cdp:eq:A.J.14} is exactly
\begin{equation}
 I=B+uB\subset K.                                \hypertarget{tag-005I}{}\label{cdp:eq:A.J.15}
\end{equation}
It is an invertible fractional ideal on the semilocal
normalization $B$. Every $A$-linear endomorphism of this
rank-one module is multiplication by an element of $K$;
preserving $I$ is equivalent, branch by branch, to belonging
to $B$. Thus $\operatorname{End}_A(I)=B$ as subrings of
$K$. Combining this with the $A$-linear isomorphism
$Q\xrightarrow{\Phi}I$ and \eqref{cdp:eq:A.J.12} gives $A=B$.
Because $A$ is local, there is a single branch over $p$,
and its ring is regular. This proves that $D$ is smooth
and that $C\to D$ is an isomorphism there.

\item Finally $A=B=\mathbb C\{s\}$ makes $u=\delta(s)$
holomorphic. The analytic coordinate change
$s'=s+u(s)z$ has invertible differential at $z=0$ and
identifies the image of $A$ modulo $z^2$ with the functions
of $s'$ alone. Since $R$ is the inverse image of $A$ under
$\overline R\to\overline A$, it follows that
\begin{equation}
 R=\mathbb C\{s'\}+z^2\mathbb C\{s',z\}
   =\mathbb C\{s',z^2,z^3\}
   \simeq\mathbb C\{s',x,y\}/(y^2-x^3).
                                                        \hypertarget{tag-005J}{}\label{cdp:eq:A.J.16}
\end{equation}
This is the asserted product germ.\qedhere

\end{enumerate}
\end{proof}

\begin{proposition}[Eliminating vertical conductor on a smooth ruled cusp normalization]\hypertarget{tag-01B0}{}\label{cdp:st:A.5.10}
\StatementRecord{cdp:st:A.5.10}{proved}{01B0}{\SeriesRef{cdp:eq:A.J.26}, \SeriesRef{cdp:eq:A.J.23}, \SeriesRef{cdp:eq:A.J.21}, \SeriesRef{cdp:eq:A.J.27}, \SeriesRef{cdp:eq:A.J.22}, \SeriesRef{cdp:st:A.5.7}}

Let $Y$ be a smooth compact complex threefold with
$H^2(Y,\mathbb Q)=0$, let $S\subset Y$ be a reduced
irreducible divisor, and suppose its normalization
$h:Z\to S\subset Y$ is a smooth irrational
ruled surface $p:Z\to B$.
Suppose its conductor divisor is $D_Z=2C+V$, where
$C$ is a ruling section and $V$ is vertical, and
the generic transversal along $h(C)$ is an ordinary cusp.
Then $B$ is elliptic, $V=0$, $C^2=0$, and $Z$ is
already minimal.
The image $h(C)$ is smooth elliptic and $S$ is locally
an ordinary-cusp product along it.  The conductor-image scheme along the vertical
components may initially be nonreduced.
\end{proposition}


\begin{proof}

\begin{enumerate}[label=\arabic*.,leftmargin=*]
\item Let $g=g(B)\geq1$, let $F$ be the full ruling fibre
class, and write $V=\sum_i c_i\Gamma_i$ for its
irreducible vertical components. Put
$v=V\cdot C$, $n=C^2$ and $k=2g-2$.
Conductor duality and Cartier adjunction give the exact
line-bundle comparison
\begin{equation}
 K_Z\otimes\mathcal O_Z(2C+V)
       =h^*(K_Y\otimes\mathcal O_Y(S)).
                                                        \hypertarget{tag-005P}{}\label{cdp:eq:A.J.21}
\end{equation}
Its right side has rational first Chern class zero.
Intersecting with $C$ and using adjunction on the section
gives
\begin{equation}
 n=2-2g-v=-k-v.                                  \hypertarget{tag-005Q}{}\label{cdp:eq:A.J.22}
\end{equation}
Let $D_p$ be the divisorial common zero of
$dp:T_Z\to p^*T_B$. Its remaining zero scheme is
zero-dimensional, say $\Xi$, so the image is
$I_\Xi\otimes p^*T_B(-D_p)$.
The kernel $T_p$ is an invertible sheaf, by the depth
lemma. Taking determinants gives the comparison
\[
 T_p=K_Z^{-1}\otimes p^*K_B\otimes\mathcal O_Z(D_p).
\]
The section makes $dp$ surjective at every point of $C$.
Consequently $D_p$ and $\Xi$ are disjoint from $C$,
and $T_p|_C=N_{C/Z}$.

\item Put $\mathcal F=h^*T_Y$. Poincar\'e duality on the compact
oriented six-manifold $Y$ also gives
$H^4(Y,\mathbb Q)=0$. Therefore
\begin{equation}
 c_1(\mathcal F)=0,\qquad c_2(\mathcal F)=0
        \quad\hbox{in rational cohomology}.       \hypertarget{tag-005R}{}\label{cdp:eq:A.J.23}
\end{equation}
The differential is generically injective. Let
$L\subset\mathcal F$ be the saturation of the image
of $T_p\xrightarrow{dh}\mathcal F$.
It is invertible: a saturated rank-one subsheaf of a
locally free sheaf on a smooth surface is reflexive.
The map $T_p\to L$ has no vertical divisorial zero.
Indeed $h$ is finite, so no ruling fibre is contracted;
in characteristic zero its derivative along that fibre
is nonzero generically.

\item At the generic point of $C$, an ordinary cusp has
normalization differential equivalent to that of
$(u,s)\mapsto(u^2,u^3,s)$.
The ruling tangent either lies in its kernel along $C$
or does not. The vanishing order in the first case is
exactly one, and in the second case it is zero.
Write $\epsilon=1$ or $\epsilon=0$ for these respective
alternatives. There are no other horizontal conductor
components, and the normalization is an immersion away
from the conductor. Hence the line comparison is
\begin{equation}
 L=T_p\otimes\mathcal O_Z(\epsilon C).
                                                        \hypertarget{tag-005S}{}\label{cdp:eq:A.J.24}
\end{equation}
This checks both positions of the derivative kernel
relative to the ruling tangent.

\item The quotient $E=\mathcal F/L$ is torsion-free of rank
two. The differential induces a generically nonzero map
\[
 T_Z/T_p=I_\Xi\otimes p^*T_B(-D_p)
                    \longrightarrow E\longrightarrow E^{**}.
\]
It extends uniquely to a map from $p^*T_B(-D_p)$ into
$E^{**}$, since both are locally free and the omitted
set has codimension two; locally this is Hartogs
extension of its coefficient functions.
Let $M$ be its saturated image line in the rank-two
bundle $E^{**}$. At the generic point of $C$, the
two-by-two differential minors have common vanishing
order one. Removing the order $\epsilon$ already
assigned to $L$ leaves order $1-\epsilon$ for this map.
At the generic point of $\Gamma_i$, take a plane transversal
to its image and a normalization parameter $u$ for the
corresponding plane-curve branch. If $m_i$ is its
multiplicity, its derivative modulo the nonzero tangent
along $\Gamma_i$ vanishes to order $r_i=m_i-1$.
These are all codimension-one zeros. Put
$R=\sum_i r_i\Gamma_i$ and $r=R\cdot C$. We obtain
\begin{equation}
 M=p^*T_B(-D_p)\otimes
           \mathcal O_Z\bigl((1-\epsilon)C+R\bigr).
                                                        \hypertarget{tag-005T}{}\label{cdp:eq:A.J.25}
\end{equation}

\item The plane-curve conductor gives the inequalities
\begin{equation}
 2R\leq V,\qquad 2r\leq v.                      \hypertarget{tag-005U}{}\label{cdp:eq:A.J.26}
\end{equation}
To spell out their local length content, on a single
plane branch of multiplicity $m_i$ the classes
$u,u^2,\ldots,u^{m_i-1}$ are linearly independent in
the normalization quotient. Its delta invariant is thus
at least $m_i-1$. The branch is Gorenstein. Finite
duality identifies the length of its ring modulo its
conductor with the length of its normalization quotient;
the intrinsic conductor exponent is consequently twice
that delta invariant. Other branches add their
nonnegative intersection multiplicities to this exponent.
For the last comparison, writing the plane equation
as $\prod_j f_j$ multiplies the residue denominator on
branch $i$ by $\prod_{j\ne i}f_j$, whose valuation is
exactly the sum of those intersection multiplicities.
This proves \eqref{cdp:eq:A.J.26} for every branch, including cases
with nonreduced conductor-image scheme.

\item Let $\ell_1$ be the length of $E^{**}/E$. Saturation
of $M$ gives a zero-dimensional subscheme $\Xi_M$ and a
line bundle $N$ with
\[
 0\longrightarrow M\longrightarrow E^{**}
       \longrightarrow I_{\Xi_M}\otimes N\longrightarrow0.
\]
Write $\ell_2=\operatorname{length}\Xi_M$.
First Chern classes give $c_1(N)=-c_1(L)-c_1(M)$.
The second Chern classes in these typed quotient
sequences and \eqref{cdp:eq:A.J.23} give the intersection identity
\begin{equation}
 L^2+L\cdot M+M^2=\ell_1+\ell_2\geq0.            \hypertarget{tag-005V}{}\label{cdp:eq:A.J.27}
\end{equation}
Indeed $c_2(E)=L^2$,
$c_2(E^{**})=L^2-\ell_1$, and
$c_2(E^{**})=-L\cdot M-M^2+\ell_2$.

\item Set $W=V+D_p$ and $B_v=R-D_p$, both vertical
divisors, the first effective. The exact comparisons
above and \eqref{cdp:eq:A.J.21} give the numerical classes
\[
 L\equiv(2+\epsilon)C+W+kF,\qquad
 M\equiv(1-\epsilon)C+B_v-kF.
\]
The intersection form on vertical divisors is negative
semidefinite. One can check this directly by contracting
the ruling blowups: every vertical class is a multiple
of the pullback fibre class plus exceptional classes,
on which the form is negative definite. Hence
\[
 q_v=W^2+W\cdot B_v+B_v^2
      =(W+B_v/2)^2+3B_v^2/4\leq0.
\]
Its vanishing implies $W^2=B_v^2=0$; these classes
are then numerical multiples of the fibre class.

\item For $\epsilon=1$, using $D_p\cdot C=0$ in the
left side of \eqref{cdp:eq:A.J.27} gives
\[
 \ell_1+\ell_2=-6k-3v+3r+q_v
                  \leq-6k-3v/2.
\]
Nonnegativity, with $k\geq0$ and $v\geq0$,
forces $g=1$, $v=0$, $r=0$ and $q_v=0$.

\item For $\epsilon=0$, the same intersection calculation reads
\begin{equation}
 0\leq\ell_1+\ell_2
       =-6k-2v+4r+q_v\leq-6k.                   \hypertarget{tag-005W}{}\label{cdp:eq:A.J.28}
\end{equation}
It forces $g=1$, $\ell_1=\ell_2=0$, $2r=v$,
and $q_v=0$.
Thus $E=E^{**}$ is locally free, and
$L\to\mathcal F$ is a subbundle inclusion.
Along $C$, the composite
$T_C\to\mathcal F|_C\to E|_C$ vanishes generically,
because $dh$ has rank one there and its nonzero ruling
tangent image is $L|_C$. It therefore vanishes
everywhere. The exact subbundle sequence gives an
nonzero holomorphic line-bundle map
\begin{equation}
 T_C\longrightarrow L|_C=T_p|_C=N_{C/Z}.
                                                        \hypertarget{tag-005X}{}\label{cdp:eq:A.J.29}
\end{equation}
It is nonzero since the finite map $C\to h(C)$ has
nonzero derivative generically. Its Hom line has degree
$n-(2-2g)=-v$ by \eqref{cdp:eq:A.J.22}. A nonzero section on a
compact curve has nonnegative degree, so $v=0$ in this
alternative too.

\item In both alternatives, $v=r=0$ and $q_v=0$.
Thus $W$ is numerically a fibre multiple and
$W\cdot C=V\cdot C+D_p\cdot C=0$, so it is
numerically zero. An effective nonzero divisor on the
projective ruled surface has positive intersection
with an ample class. It follows that $W=0$ and hence
$V=D_p=0$. Now $B_v=R$ is effective, has square zero
and has zero intersection with $C$, so the same
argument gives $R=0$. Equation \eqref{cdp:eq:A.J.22} gives $C^2=0$.
Finally $K_Z+2C$ is numerically trivial. A ruling
$(-1)$-curve $E_0$ would satisfy
$0=(K_Z+2C)\cdot E_0=-1+2C\cdot E_0$, impossible
because $C\cdot E_0$ is an integer. There are therefore
no ruling blowups and $Z$ is minimal.

\item \Cref{cdp:st:A.5.7} now applies everywhere along
the conductor: the normalization is smooth and its
conductor is exactly $2C$. It makes the image smooth
and the surface locally the cusp product. The image is
elliptic because $C\to h(C)$ is an isomorphism.

\end{enumerate}
\end{proof}

\begin{proposition}[A noncontracted horizontal cusp has zero Euler number]\hypertarget{tag-01BB}{}\label{cdp:st:A.5.20}
\StatementRecord{cdp:st:A.5.20}{proved}{01BB}{\SeriesRef{cdp:st:A.5.14}, \SeriesRef{cdp:st:A.5.5}, \SeriesRef{cdp:st:C2d}, \cref{lem:II:elliptic-ruled-boundary}}

Let $f:Y\to C_0$ be a proper holomorphic map from a smooth
compact complex threefold with $b_2(Y)=0$, with connected
fibres and smooth general fibre a complex two-torus.
Let $mS$ be a full fibre with irreducible reduced support,
$\eta:T\to S$ its normalization and $\pi:Z\to T$ an
irrational ruled resolution with ruling $p:Z\to B$.
Suppose every exceptional curve of $\pi$ is vertical
and the resolution conductor is $2C+V$, where $C$
is a noncontracted section and $V$ is vertical.
If the generic transversal at $D=\eta\pi(C)$ is an
ordinary cusp, then $D$ and $C$ are smooth elliptic
curves and $C\to D$ is an isomorphism. The cusp type
is constant along $D$. The normalization $T$ is a
smooth minimal ruled surface over this elliptic curve,
its conductor is twice a section, and $e(S)=0$.
\end{proposition}


\begin{proof}
\begin{enumerate}[label=\arabic*.,leftmargin=*]
\item The section identifies $C$ with the positive-genus curve $B$.
The proper nonconstant map $C\to D$ is finite. At a general point of
$D$, the cusp normalization
\[
 (u,z)\longmapsto(u^3,u^2,z)
\]
has one branch and degree one on its singular curve. Thus $C\to D$ is
the normalization of $D$.
Proposition~\ref{cdp:st:A.5.14}, applied to the same $Y,S,D$, makes
$D$ smooth elliptic with constant cusp transversal. It follows that
$C\to D$ is an isomorphism and $B$ is elliptic.
Let $C_T\subset T$ be the inverse-image curve.
Proposition~\ref{cdp:st:C2d} gives a smooth neighbourhood of $C_T$,
the isomorphism $C_T\to D$, and
$\deg N_{C_T/T}=0$. Its leading normal terms identify the square and
cube of this normal line with the radical and Hessian quotient lines.

\item Factor the given resolution as
\[
 Z\xrightarrow{\sigma}Z'\xrightarrow{\mu}T,
 \qquad \pi=\mu\sigma,
\]
where $\mu$ is the minimal resolution and $\sigma$ is a sequence of
point blowups. Every $\sigma$-exceptional curve is $\pi$-exceptional
and is vertical for $p$. The map $p$ therefore descends to
$p':Z'\to B$. Explicitly, the proper image of $(\sigma,p)$ in
$Z'\times B$ projects finitely and bimeromorphically to the normal
surface $Z'$ with one point in every fibre. This projection is an
isomorphism and defines $p'$, with $p=p'\sigma$.
The curve $C'=\sigma(C)$ is a section and $\sigma|_C$ is an isomorphism.
Since $\mu$ is an isomorphism near $C_T$, one has
$(C')^2=\deg N_{C_T/T}=0$.

The strict transform of every $\mu$-exceptional curve is
$\pi$-exceptional, so those curves are vertical for $p'$.
Proposition~\ref{cdp:st:A.5.3} supplies the canonical effective divisor
on this minimal resolution. Its nonexceptional components have the
same conductor valuations as on $Z$. Consequently it has the form
$2C'+V'$, with $V'$ effective and vertical, and
\[
 \omega_{Z'}(2C'+V')\simeq(\eta\mu)^*\omega_S,
 \qquad c_1(K_{Z'}+2C'+V')=0\quad\hbox{in }H^2(Z',\Q).
\]
The constant cusp charts and the isomorphism of $\mu$ there show that
$V'$ is disjoint from $C'$.

\item Apply \cref{lem:II:elliptic-ruled-boundary} to $\mu:Z'\to T$, $p':Z'\to B$, $C'$ and $V'$. The previous step gives every hypothesis: $B$ is elliptic, $(C')^2=0$, $V'$ is effective, vertical and disjoint from $C'$, the canonical-boundary class is zero, and the exceptional curves of the minimal resolution are vertical. Thus $V'=0$ and $T=Z'$ is a smooth minimal ruled bundle.

\item Finite duality on the smooth normalization identifies the
conductor ideal with the invertible ideal of its divisor, now
$2C_T$. Thus $\eta$ is an isomorphism off $C_T$, and its restriction
$C_T\to D$ is an isomorphism. The normalization is consequently a
proper continuous bijection and a homeomorphism of the underlying
compact spaces. Equivalently, additivity of Euler characteristic
gives
\[
 e(S)=e(T)-e(C_T)+e(D)=0.
\]
Here $e(T)=e(\mathbf P^1)e(B)=2\cdot0$, and both curves have genus one.
The full scheme fibre $mS$ has the same underlying space as $S$,
so its topological Euler characteristic is also zero.
\end{enumerate}
\end{proof}

\subsection{Binomial Milnor fibres}\label{subsec:II:normalization-binomial}

\begin{lemma}[Matching spheres from a curve and a monomial]\hypertarget{tag-P297}{}\label{lem:II:monomial-curve-matching-sphere}
\hypertarget{tag-P247}{}\label{lem:II:conic-bundle-matching-sphere}
\StatementRecord{lem:II:monomial-curve-matching-sphere}{proved}{P297}{}
\StatementRecord{lem:II:conic-bundle-matching-sphere}{proved}{P247}{}
Let $B\subset\C^2_{x,v}$ be a connected smooth complex curve such that $v|_B$ has two distinct simple zeros. Let $b',c'$ be relatively prime positive integers, and put
\[
 X_{B;b',c'}=\{(x,y,z)\in\C^3:(x,y^{b'}z^{c'})\in B\}.
\]
Choose an embedded arc $\gamma$ between two such zeros, avoiding every other zero. The smooth locus of $X_{B;b',c'}$ contains the compact spherical cycle \eqref{cdp:eq:A.J.33a}, whose class has square $-2b'c'$. For $b'=c'=1$, write $X_B=X_{B;1,1}$. The cycle is then a smoothly embedded sphere, its parametrization is \eqref{cdp:eq:A.BP.3}, and its square is $-2$.
\end{lemma}
\begin{proof}
\begin{enumerate}[label=\arabic*.,leftmargin=*]
\item Over $v\ne0$, both $y,z$ are nonzero and the derivative of $y^{b'}z^{c'}$ is nonzero, so the projection to $B$ is locally a smooth one-dimensional fibration. At a simple zero of $v|_B$, write $x=\chi(v)$. Then
\[
 (y,z)\longmapsto\bigl(\chi(y^{b'}z^{c'}),y,z\bigr)
\]
is a holomorphic chart on the surface, with inverse projection to $(y,z)$. Thus it is smooth near the inverse image of the chosen arc.

\item Make the arc radial in the local $v$ coordinate at its endpoints $P,Q$. Choose smooth complex functions $A(t),B_0(t)$ on its interior such that
\[
 A(t)^{b'}B_0(t)^{c'}=v(\gamma(t)),\qquad
 |A(t)|=|v(\gamma(t))|^{1/(2b')},\quad
 |B_0(t)|=|v(\gamma(t))|^{1/(2c')}.
\]
They exist by choosing a continuous argument along the arc. The continuous map
\begin{equation}
 (t,\theta)\longmapsto
 \bigl(x(\gamma(t)),A(t)e^{\mathrm i c'\theta},
                 B_0(t)e^{-\mathrm i b'\theta}\bigr)
 \hypertarget{tag-0063}{}\label{cdp:eq:A.J.33a}
\end{equation}
from $[0,1]\times S^1$ collapses its two boundary circles to $P,Q$, giving an oriented spherical cycle $\sigma_\gamma:S^2\to X_{B;b',c'}$. Coprimality makes each interior circle embedded. The projection to $B$ recovers the arc, so the whole map is injective away from the collapsed ends.

\item Take a nearby arc $\gamma'$ with the same endpoints and disjoint interior, by an isotopy in a thin disc neighbourhood of $\gamma$ avoiding all other zeros. Choose the arguments in \eqref{cdp:eq:A.J.33a} continuously throughout the isotopy. The coordinates tend uniformly to zero at its endpoints, so this is a continuous homotopy of maps from $S^2$. The two cycles are homologous and meet only at $P,Q$.

In the holomorphic chart $(y,z)$ at either endpoint, arrange the arcs radial with distinct arguments in $v$. The links of the two cycles in a small three-sphere are parallel curves of slope $(c',-b')$ on a Clifford torus. Their oriented linking number is $-b'c'$: the torus framing of a $(p,q)$ curve has linking number $pq$ with its parallel copy, obtained by pushing the curves to opposite solid tori in the Heegaard splitting of $S^3$ and counting meridian discs. Replacing the two cones by transverse relative chains with those boundary links identifies their local intersection with this linking number. At the second endpoint both cycle orientations reverse relative to the outgoing radial coordinate, leaving the intersection sign unchanged. Consequently
\begin{equation}
 [\sigma_\gamma]^2
   =[\sigma_\gamma]\cdot[\sigma_{\gamma'}]
   =-2b'c'\ne0.
 \hypertarget{tag-0064}{}\label{cdp:eq:A.J.33b}
\end{equation}
The cycles and their homotopy all have compact support.

\item For $b'=c'=1$, the map over $v\ne0$ is explicitly
\[
 (b,y)\longmapsto(x(b),y,v(b)/y),\qquad
 (x,y,z)\longmapsto((x,yz),y)
\]
in the two directions. The sphere is
\begin{equation}
 (t,\theta)\longmapsto
   \bigl(x(\gamma(t)),A(t)e^{\mathrm i\theta},
                        B_0(t)e^{-\mathrm i\theta}\bigr).
 \hypertarget{tag-0051}{}\label{cdp:eq:A.BP.3}
\end{equation}
At an endpoint with $v=e^{\mathrm i\phi}|y|^2$, its disk is
\[
 y\longmapsto
 \bigl(\chi(e^{\mathrm i\phi}|y|^2),y,
                         e^{\mathrm i\phi}\overline y\bigr),
\]
with inverse disk coordinate given by projection to $y$. Together with the interior circle coordinate $y/A(t)$, these maps prove that the sphere is smoothly embedded. For a parallel arc, rotating $z$ puts the two oriented disks in the form $z=\overline y$ and $z=\eta\overline y$, with $\eta=a+\mathrm i b\ne1$. Their tangent-frame determinant in $(\Re y,\Im y,\Re z,\Im z)$ is
\[
                         -((a-1)^2+b^2)<0.
\]
Thus each endpoint intersection is $-1$, agreeing with the framing calculation above, and the square is $-2$.
\end{enumerate}
\end{proof}

\begin{proposition}[Binomial germs and spherical intersection classes]\hypertarget{tag-01B5}{}\label{cdp:st:A.5.15}
\StatementRecord{cdp:st:A.5.15}{proved}{01B5}{\cref{lem:II:monomial-curve-matching-sphere}}
\hypertarget{tag-02WP}{}\label{night:P4.18}\label{stmt:P4.18}

Let $f:Y\to C_0$ be a proper holomorphic map from a
smooth complex threefold to a smooth curve, with smooth
nearby fibres complex two-tori. There is no point of
a fibre at which, in holomorphic ambient coordinates
$(x,y,z)$ and a coordinate on the base, the map has
the equation
\begin{equation}
 f(x,y,z)=u(x,y,z)\bigl(x^a-y^bz^c\bigr)^m,
 \quad a\ge2,\quad b,c,m\ge1,                                    \hypertarget{tag-0062}{}\label{cdp:eq:A.J.33}
\end{equation}
where $u$ is a holomorphic unit. This includes the pinch
point, every germ $x^2-y^{2k}z^{2r+1}$ with
$k\ge1$, $r\ge0$, and the reducible germ
$x^2-y^2z^2$.
For every $a\ge2$ and $b,c\ge1$, the affine nearby fibre
of $x^a-y^bz^c$ contains a spherical
class of nonzero square. In particular the exclusion applies
to $x^2-z^3y^4$, whose smooth normalization has a divisorial
differential rank drop.
\end{proposition}


\begin{proof}

\begin{enumerate}[label=\arabic*.,leftmargin=*]
\item Put $h=x^a-y^bz^c$. This polynomial is reduced:
a repeated factor would divide both $h$ and $a x^{a-1}$,
hence divide $x$, which does not divide $h$.
We compare the map with this reduced equation. This polynomial is weighted
homogeneous of degree $a(b+c)$ for weights
$(b+c,a,a)$. On a small simply connected coordinate
neighbourhood choose a holomorphic logarithm of $u$.
Define the holomorphic coordinate change
\begin{equation}
 \Phi(x,y,z)=\left(
 x u^{1/(ma)},\;
 y u^{1/(m(b+c))},\;
 z u^{1/(m(b+c))}\right).
                                                        \hypertarget{tag-0065}{}\label{cdp:eq:A.J.34}
\end{equation}
Its derivative at the origin is an invertible diagonal
matrix. The holomorphic inverse function theorem and
weighted homogeneity give
$h\circ\Phi=u^{1/m}h$ and therefore
$f=(h\circ\Phi)^m$ exactly. To specify the inverse, put
$D=a(b+c)$, choose $v=u^{1/(mD)}$, and solve
\[
 \tau^{mD}=u(\tau^{-(b+c)}X,\tau^{-a}Y,\tau^{-a}Z),
 \qquad \tau(0)=v(0).
\]
The derivative in $\tau$ at the origin is
$mD\,v(0)^{mD-1}\ne0$. The inverse is
$\Psi(X,Y,Z)=(\tau^{-(b+c)}X,\tau^{-a}Y,\tau^{-a}Z)$:
uniqueness of the chosen root gives $v\circ\Psi=\tau$
and $\tau\circ\Phi=v$, hence both inverse identities.
For a sufficiently small
nonzero regular value $t$, a chosen Milnor fibre of
$h$ at one of the $m$ roots of $t$ is consequently a
compact oriented submanifold with boundary of the
smooth global fibre $Y_t$. Denote this Milnor fibre
by $M$.

\item We construct a spherical cycle. Put $d=\gcd(b,c)$,
$b'=b/d$ and $c'=c/d$, and consider the smooth affine
curve
\[
 B=\{(x,v)\in\C^2:x^a-v^d=1\}.
\]
The affine surface $H=\{h=1\}$ maps holomorphically
to $B$ by $v=y^{b'}z^{c'}$. The curve $B$ is
connected: projection to the $v$-line has $a$ sheets,
and a small loop about any root of $v^d+1$ cyclically
permutes them. It is smooth because its two partial
derivatives cannot vanish on $x^a-v^d=1$.
The function $v$ has $a$ distinct simple zeros
$P_j=(\zeta_a^j,0)$ on $B$.

\item Apply \cref{lem:II:monomial-curve-matching-sphere} to the curve $B$ just constructed and the relatively prime integers $b'=b/d$, $c'=c/d$. The identity on $(x,y,z)$ identifies its surface $X_{B;b',c'}$ with $H$, because $(y^{b'}z^{c'})^d=y^bz^c$. The lemma gives the exact cycle \eqref{cdp:eq:A.J.33a}, its parallel-arc homotopy, and the square \eqref{cdp:eq:A.J.33b}. Their supports are compact.
The weighted homothety
$(x,y,z)\mapsto(s^{b+c}x,s^a y,s^a z)$ maps
$H$ biholomorphically onto $\{h=s^{a(b+c)}\}$.
For sufficiently small positive $s$, it puts those
compact supports into any prescribed Milnor
neighbourhood of the origin. Thus \eqref{cdp:eq:A.J.33b} is a
spherical class of nonzero square in a local Milnor fibre $M$.

\item Let $j:M\hookrightarrow Y_t$ be the embedding obtained from \eqref{cdp:eq:A.J.34}. A complex two-torus
has contractible universal cover $\C^2$, so
$\pi_2(Y_t)=0$. In particular the spherical class
constructed in \eqref{cdp:eq:A.J.33a} maps to zero in
$H_2(Y_t,\Z)$. Intersection numbers of compact
two-cycles in the interior of $M$ are preserved by
this oriented codimension-zero embedding: the same
small transverse perturbations compute them in
both manifolds. Its square in $Y_t$ is therefore
both zero and $-2b'c'$, contradicting \eqref{cdp:eq:A.J.33b}.
This proves the exclusion,
including every multiplicity in \eqref{cdp:eq:A.J.33}.

\item For the mixed example the normalization is the finite
birational map
\[
 (u,v)\longmapsto(u^3v^2,v,u^2).
\]
Finiteness follows from $u^2=z$, and $u=x/(zy^2)$ gives
birationality. The image ring is
$R=\mathbb C\{u^2,v\}\oplus u^3v^2\mathbb C\{u^2,v\}$:
expanding a series in the generators gives this expression,
and every term on the right is in $R$. Equivalently, the
odd monomials have $u$-exponent at least three and
$v$-exponent at least two.
Since the normalization is generated over the image ring by $1,u$,
its conductor consists of the series $h$ for which both $h$ and $uh$
lie in that ring. The even--odd monomial description therefore gives
$u^2v^2\mathbb C\{u,v\}$, with divisor
$2\{u=0\}+2\{v=0\}$. Its differential minors are
$3u^2v^2,-4u^4v,-2u$, and generate the ideal $(u)$. Thus the rank-drop locus is exactly
the divisor $\{u=0\}$, which maps by $v\mapsto(0,v,0)$
to a line. 

\end{enumerate}
\end{proof}

\begin{proposition}[Two monomial-separated branches with arbitrary multiplicities]\hypertarget{tag-P296}{}\label{prop:II:weighted-two-branch-sphere}
\StatementRecord{prop:II:weighted-two-branch-sphere}{proved}{P296}{\cref{lem:II:monomial-curve-matching-sphere}}
Let $a,b,k,l$ be positive integers and put
\[
 G(X,Y,Z)=(X-Y^kZ^l)^a(X+Y^kZ^l)^b,\qquad e=\gcd(k,l).
\]
The affine fibre $G=1$ contains a compact spherical class of square $-2(k/e)(l/e)$. Weighted homothety puts that class, together with its intersection comparison, into arbitrarily small nonzero fibres near the origin. Consequently a map from a smooth complex threefold to a smooth curve with nearby smooth fibres complex two-tori cannot have local full-fibre equation $uG$, where $u$ is a holomorphic unit. This conclusion retains the original multiplicities $a,b$.
\end{proposition}
\begin{proof}
\begin{enumerate}[label=\arabic*.,leftmargin=*]
\item Put $d=\gcd(a,b)$, $a=da_0$, $b=db_0$ and $N=a_0+b_0\ge2$. On the part of $G=1$ where
\[
 (X-Y^kZ^l)^{a_0}(X+Y^kZ^l)^{b_0}=1,
\]
both factors are nonzero. The maps
\[
 X-Y^kZ^l=s^{b_0},\qquad X+Y^kZ^l=s^{-a_0}
\]
identify this part with
\[
 \mathcal X=\{(s,Y,Z)\in\C^*\times\C^2:
              Y^kZ^l=\psi(s)\},\qquad
 \psi(s)=\frac{s^{-a_0}-s^{b_0}}2.
\]
The map from $\mathcal X$ to the original fibre is
\[
 \Theta(s,Y,Z)=\left(\frac{s^{b_0}+s^{-a_0}}2,Y,Z\right).
\]
For integers $r,t$ with $b_0r-a_0t=1$, its holomorphic inverse is
\[
 (X,Y,Z)\longmapsto
 \bigl((X-Y^kZ^l)^r(X+Y^kZ^l)^t,Y,Z\bigr).
\]
The relation between the two factors and the B\'ezout identity verify both inverse identities, including points with $Y=Z=0$.

\item Write $k=ek'$, $l=el'$ with $\gcd(k',l')=1$, and form the curve
\[
 B'=\{(s,v)\in\C^*\times\C:v^e=\psi(s)\}.
\]
The zeros of $\psi$ are the $N$ roots of $s^N=1$, and at each of them
\[
                    \psi'(s)=-\frac N2s^{-a_0-1}\ne0.
\]
The curve is smooth: away from $v=0$ its derivative in $v$ is nonzero, and at $v=0$ its derivative in $s$ is nonzero. It is connected. Over the punctured $s$-line with the zeros removed it is the cyclic $e$-sheeted covering obtained by taking an $e$th root of $\psi$; a loop around any simple zero cyclically permutes all $e$ sheets. This monodromy is transitive, and adjoining the points over the zeros preserves connectedness. The assertion is immediate also for $e=1$.

At each zero, $v$ is a local coordinate on $B'$, so $v|_{B'}$ has exactly $N$ distinct simple zeros. The identity map on $(s,Y,Z)$ identifies $\mathcal X$ with $X_{B';k',l'}$ of \cref{lem:II:monomial-curve-matching-sphere}, since
\[
                 (Y^{k'}Z^{l'})^e=Y^kZ^l=\psi(s).
\]
Apply that lemma to any arc between two of these zeros. Its cycle has square $-2k'l'$. The biholomorphism $\Theta$ and its displayed inverse carry the cycle, its parallel-arc comparison and both endpoints into the original full $G$-fibre. Thus their compact intersection number is unchanged.

\item The polynomial $G$ has weights $(k+l,1,1)$ and degree $D=(k+l)(a+b)$. The homothety
\[
 \delta_\lambda(X,Y,Z)
       =(\lambda^{k+l}X,\lambda Y,\lambda Z)
\]
maps $G=1$ biholomorphically to $G=\lambda^D$, with inverse $\delta_{1/\lambda}$, and shrinks all the compact supports to the origin as $\lambda\to0$.

Choose a local logarithm of $u$ and set
\[
 \Phi(X,Y,Z)
   =\bigl(u^{1/(a+b)}X,u^{1/D}Y,u^{1/D}Z\bigr).
\]
Its derivative at the origin is invertible and weighted homogeneity gives $G\circ\Phi=uG$. For an explicit inverse, choose $v=u^{1/D}$ and solve
\[
 \tau^D=u(\tau^{-(k+l)}X',\tau^{-1}Y',\tau^{-1}Z'),
 \qquad \tau(0)=v(0).
\]
Its derivative in $\tau$ at the origin is $D v(0)^{D-1}\ne0$. The inverse is
\[
 \Psi(X',Y',Z')
       =(\tau^{-(k+l)}X',\tau^{-1}Y',\tau^{-1}Z'),
\]
because uniqueness of the chosen root gives $v\circ\Psi=\tau$ and $\tau\circ\Phi=v$. The map $\Psi\circ\delta_\lambda\circ\Theta$ therefore embeds a neighbourhood of the constructed compact cycles in the local full fibre $uG=\lambda^D$, with the original exponents $a,b$.

Choose $\lambda$ sufficiently small so that the corresponding nearby global fibre is a smooth two-torus. The inclusion preserves the compact intersection number, while the composite sphere map lifts to the contractible universal cover of that torus and has zero homology class. Its square would then be both zero and $-2k'l'$, a contradiction.
\end{enumerate}
\end{proof}

\begin{proposition}[The homotopy and intersection form of an irreducible binomial Milnor fibre]\hypertarget{tag-P225}{}\label{prop:II:binomial-milnor-rank}
\StatementRecord{prop:II:binomial-milnor-rank}{proved}{P225}{\cite[Theorem~L1 and Example~2.4(a), pp.~2--3]{Nemethi1991Join}; \cite[Theorem~4.1 and Lemma~4.6]{MPW2005Hirzebruch}}
Let $a\ge2$, $b,c\ge1$, $\gcd(a,b,c)=1$, and $d=\gcd(b,c)$. For the Milnor fibre $M$ of $x^a-y^bz^c$, the homotopy type is \eqref{cdp:eq:A.J.35}, and its rational intersection morphism has rank $a-1$. Every class in $H_2(M,\Z)$ is a sum of spherical classes. For $a=2$, $b=1$, $c=\ell\ge1$, its boundary is $L(2\ell,1)$ and the rank-one integral intersection form has absolute determinant $2\ell$.
\end{proposition}
\begin{proof}
\begin{enumerate}[label=\arabic*.,leftmargin=*]
\item When $\gcd(a,b,c)=1$, the homotopy type and
boundary invariants give the full rational intersection rank.
The join theorem
\cite[Theorem~L1 and Example~2.4(a), pp.~2--3]{Nemethi1991Join}
applies to arbitrary analytic germs, including
non-isolated ones. Set $d=\gcd(b,c)$. The Milnor
fibre of $y^bz^c$ is a disjoint union of $d$ annuli:
its equations in a Milnor polydisc are
$y^{b/d}z^{c/d}=\zeta\epsilon^{1/d}$, and each
component is parametrized by
$y=A w^{c/d}$, $z=B w^{-b/d}$. For integers $r,s$ with $(c/d)r-(b/d)s=1$, its inverse is
$w=(y/A)^r(z/B)^s$. In a polydisc of radius $R$, the bounds are
$( |B|/R)^{d/b}\le |w|\le(R/|A|)^{d/c}$;
they define a nonempty annulus for sufficiently small
$\epsilon$. Thus the parametrization and its inverse
identify each component with that annulus.
The Milnor fibre of $x^a$ consists of $a$ points.
The join theorem identifies $M$ with the homotopy pushout
$(\operatorname{Cone}G\times H)\cup_{G\times H}
(G\times\operatorname{Cone}H)$, where $G$ is the set of
$a$ points and $H$ is the union of the $d$ circles.
This pushout is homeomorphic to the union $J$ of $a$ cones
on $H$ with their common base identified. For cones with
base parameter zero and apex parameter one, the homeomorphism
sends $([g,t],h)$ to $[g,h,(1-t)/2]$ and
$(g,[h,t])$ to $[g,h,(1+t)/2]$.
Its inverse on $[g,h,r]$ uses $([g,1-2r],h)$ for
$r\le1/2$ and $(g,[h,2r-1])$ for $r\ge1/2$.
The quotient map $J\to J/\operatorname{Cone}_1H$ collapses a contractible CW subcomplex and is a homotopy equivalence by \cite[Propositions~0.16--0.17, pp.~15--16]{Hatcher2002}. Its quotient is a wedge of $a-1$ unreduced suspensions of $H$. Each suspension has two vertices, $d$ meridian edges and $d$ two-cells attached along a meridian followed by its reverse. Collapse one meridian by the same quotient theorem. The remaining $d-1$ edges are loops, and the backtracking attaching maps are homotopic to constants; \cite[Proposition~0.18, p.~16]{Hatcher2002} identifies the resulting attachment with $d$ spheres. This proves
\begin{equation}
 M\simeq
 \bigvee^{(a-1)d}S^2\ \vee\
 \bigvee^{(a-1)(d-1)}S^1.                          \hypertarget{tag-0066}{}\label{cdp:eq:A.J.35}
\end{equation}
In particular, every class in $H_2(M,\Z)$ is a sum
of images of spherical fundamental classes.

\item For the same irreducible germ, the boundary calculation
of \cite[Theorem~4.1, p.~6, and Lemma~4.6, p.~11]{MPW2005Hirzebruch}
gives a Seifert three-manifold with base genus
\begin{equation}
 g_\partial=(a-1)(d-1),\qquad
 e_0=\frac{a}{(b/d)(c/d)}>0.                     \hypertarget{tag-0067}{}\label{cdp:eq:A.J.36}
\end{equation}
Here $e_0$ is its rational Seifert Euler number in
the orientation convention of that paper. Its variables
$(z,x,y;m,k,\ell)$ correspond to
$(x,y,z;a,b,c)$, exchanging $y,z$ when needed to arrange
$b\le c$; the inverse is the reverse coordinate permutation.
An exceptional-fibre order equal to one contributes zero
to the Euler correction. The formula therefore includes
$b/d=1$ or $c/d=1$. Its
nonzero value gives
$b_1(\partial M)=2g_\partial$. To make this last
deduction explicit, the rational abelianized Seifert
presentation has a regular-fibre generator $v$,
exceptional-fibre generators $q_i$, the relations
$\alpha_iq_i+\beta_iv=0$ and
$\sum_iq_i+ev=0$, and $2g_\partial$ free base
generators. Eliminating the $q_i$ gives $e_0v=0$.
Thus all fibre generators vanish over $\Q$, leaving
exactly the $2g_\partial$ base generators.

\item Let $I_M:H_2(M,\Q)\to H_2(M,\partial M;\Q)$ be
the intersection morphism. Poincar\'e--Lefschetz duality
and the homology sequence of the pair give
\begin{equation}
 \operatorname{rank} I_M
   =b_2(M)+b_1(M)-b_1(\partial M)=a-1>0.          \hypertarget{tag-0068}{}\label{cdp:eq:A.J.37}
\end{equation}
Indeed $H_3(M)=0$ by \eqref{cdp:eq:A.J.35}, so
$H_3(M,\partial M)\simeq H^1(M)$ injects into
$H_2(\partial M)$. Its cokernel is the kernel of
$I_M$, which has dimension
$b_1(\partial M)-b_1(M)$. This proves the displayed
rank formula. For $a=2$, $b=1$, $c=\ell\ge2$, the
boundary is $L(2\ell,1)$ by
\cite[Theorem~4.2(2), p.~7]{MPW2005Hirzebruch};
for $\ell=1$ the same description follows from
\cite[Theorem~4.2(1), p.~7]{MPW2005Hirzebruch}.
The same
pair sequence over $\Z$ makes the rank-one
intersection form have absolute determinant
$2\ell$. The special case $\ell=1$ is the ordinary
quadratic singularity and has determinant two.

\end{enumerate}
\end{proof}

\subsection{Isolated singularities and suspension germs}\label{subsec:II:normalization-suspensions}

\begin{proposition}[Local vanishing spheres and transverse suspension products]\hypertarget{tag-01B7}{}\label{cdp:st:A.5.17}
\StatementRecord{cdp:st:A.5.17}{proved}{01B7}{}

Let $H:(\C^3,0)\to(\C,0)$ be a holomorphic germ with
an isolated critical point at the origin. Every sufficiently
small local Milnor fibre of $H$ contains a compact embedded
two-sphere of self-intersection $-2$.

Let $f:Y\to C_0$ have the hypotheses of \ref{cdp:st:A.5.15}.
No local equation $f=u h^m$, with $u$ a holomorphic
unit and $m\ge1$, can have a reduced equation $h$
with an isolated critical point at the origin.
Suppose its local equation at a point of a fibre is
\begin{equation}
             f=(x^a-g(y,z))^m,
             \qquad a\ge2,\quad m\ge1,             \hypertarget{tag-0069}{}\label{cdp:eq:A.J.38}
\end{equation}
where $g\ne0$ and $g(0,0)=0$. Then the reduced
plane germ $\{g=0\}_{\rm red}$ is smooth. If
$b$ is its multiplicity in $\operatorname{div}(g)$,
there are local holomorphic coordinates in which
$f=(x^a-y^b)^m$.
Consequently, at every point where the full fibre
is a uniform multiple of its reduced support and
the reduced equation has a nonzero Hessian, its
critical surface germ is a constant
transverse $A_{b-1}$ product. A nonordinary change
of transverse type at a point with nonzero Hessian
is excluded.
\end{proposition}


\begin{proof}

\begin{enumerate}[label=\arabic*.,leftmargin=*]
\item We first use the germ $H$ in the local assertion.
Choose a small closed ball $K\subset\C^3$ such that
$H^{-1}(0)$ is smooth along its boundary intersection
and transverse to $\partial K$.
Such a radius is a regular value of the norm on the
punctured smooth zero fibre. The boundary argument is
\cite[Lemma~6.1, pp.~55--56]{Milnor1968}; choosing a
regular radius gives the same construction for these
holomorphic germs.
The restriction $H|_{\partial K}$ is a submersion
at its zero set. Compactness therefore supplies
$\eta>0$ for which it is a submersion wherever
$|H|<2\eta$ on the boundary. For the linear
deformation
\begin{equation}
 H_\lambda(w)=H(w)-\lambda_1w_1
                         -\lambda_2w_2-\lambda_3w_3,
                                                        \hypertarget{tag-006A}{}\label{cdp:eq:A.J.38a}
\end{equation}
this boundary submersion persists for all parameters
in a sufficiently small polydisc $P$ and all values
of modulus less than $\eta$. The gradient map of
$H$ has an isolated zero and consequently is a finite
holomorphic map germ of positive local degree
\cite[Appendix~B, pp.~111--114]{Milnor1968}.
Here the boundary degree of the holomorphic gradient is
preserved by small constant perturbations; a regular
value counts its inverse images with positive complex
Jacobian sign. The positivity argument there uses
$G-\alpha w$ with $\alpha$ outside the eigenvalues of
$DG(0)$, followed by a small regular value in its local
inverse-function image, and applies to holomorphic $G$.
Choose a small regular value $\lambda$ of the gradient map. Its
nonempty finite inverse image consists of critical
points of $H_\lambda$ with invertible Hessian. All
these points and their critical values tend to the
origin as $\lambda\to0$.

\item Remove the critical-value discriminant from
$P\times\Delta_\eta$. The relative critical locus
over this domain misses the boundary, so its
projection is proper and its image is complex analytic
\cite[Satz~22$'$, pp.~355--356, and Satz~23, p.~356]{Remmert1957Mappings}.
The map on the relative critical set is
$(\lambda,w)\mapsto(\lambda,H_\lambda(w))$;
boundary regularity makes inverse images of compact
subsets compact in the interior. Regular values show
that the image is proper. Its complement is path
connected: in a coordinate ball the complex line through
two points outside the analytic subset meets it discretely,
and small arcs detour around the finitely many intersections
of a compact segment. A chain of such balls joins any two
points in the complement. The projection of the corresponding
smooth level sets in $K$ is a proper submersion
with boundary, and hence an Ehresmann bundle.
At a nondegenerate critical point, complete squares
holomorphically. In a direction with nonzero second
derivative, solve $\partial F/\partial w_1=0$ as
$w_1=\alpha(w')$ and write
$F=F(\alpha(w'),w')+(w_1-\alpha(w'))^2A$,
where $A$ is a holomorphic unit. The coordinate
$X=(w_1-\alpha(w'))\sqrt A$ has inverse given by the
unique small solution of this equation; its derivative
in $w_1$ is nonzero. The remaining Hessian is the
invertible Schur complement, so iteration gives
$H_\lambda-H_\lambda(p)=X^2+Y^2+Z^2$.
A local Milnor fibre at a sufficiently small regular
value embeds in the corresponding level set in $K$.
The bundle transports it, by an oriented embedding,
into $K\cap H^{-1}(\epsilon_0)$ for a nonzero small
$\epsilon_0$. Varying $\epsilon_0$ in the punctured
Milnor disc transports the same sphere to every
sufficiently small nonzero local level.

\item For a positive real local Milnor value, the real
locus $X,Y,Z\in\mathbb R$ is a two-sphere. It is
Lagrangian for the standard K\"ahler form; multiplication
by $\mathrm i$ identifies its tangent bundle with
its normal bundle with the opposite oriented Euler
class. Its self-intersection is therefore
$-\chi(S^2)=-2$. Equivalently this is the case
$a=2$, $b=c=1$ of the explicit cycle in \ref{cdp:st:A.5.15},
under $U=Y+\mathrm iZ$, $V=-Y+\mathrm iZ$, which gives
$X^2+Y^2+Z^2=X^2-UV$. Its inverse is
$Y=(U-V)/2$, $Z=(U+V)/(2\mathrm i)$.
The oriented transport preserves its self-intersection,
proving the local assertion.

\item For the torus assertion, replace $h$
by $H=u^{1/m}h$, taking a holomorphic root on a small
coordinate neighbourhood. Then $f=H^m$ exactly.
This replacement preserves an isolated critical point
as a germ. Indeed every local irreducible component
of the critical locus of $H$ has constant $H$-value,
necessarily zero when it contains the origin. On
$H=0$ one has $dH=u^{1/m}dh$. Thus its critical
locus has the same zero-dimensional germ as that
of $h$.

The local sphere just constructed lies in
$K\cap H^{-1}(\epsilon_0)$. The identity inclusion of
this coordinate neighbourhood maps it into the same
global fibre $Y_{\epsilon_0^m}$, because $f=H^m$.
Its square is still $-2$. The torus has $\pi_2=0$,
so the sphere is null-homotopic and has square zero,
a contradiction.

\item Factor the analytic plane function as
\[
                   g=u\prod_{i=1}^r h_i^{b_i},
\]
where the $h_i$ are distinct reduced irreducible
germs and $u$ is a unit. Deform each $h_i$
holomorphically to $h_{i,s}$ in a finite-dimensional
parameter polydisc $P$, keeping each exponent $b_i$ fixed:
\begin{equation}
 g_s=u\prod_i h_{i,s}^{b_i},\qquad
 H_s=x^a-g_s.                                    \hypertarget{tag-006B}{}\label{cdp:eq:A.J.39}
\end{equation}
We use only sufficiently small parameters.

\item Choose a small closed ball $K\subset\C^2_{y,z}$
such that each $h_i=0$ meets $\partial K$
transversely in its smooth part, and their boundary
links are disjoint. Choose its representative so
that every critical value of $g$ on $K$ is zero.
This is possible because the critical locus of a
holomorphic function has finitely many local
irreducible components, the function is constant
on each of them, and each component of the germ
contains the origin. The boundary smoothness,
transversality and disjointness persist for
$h_{i,s}$ by compactness and the implicit function
theorem. Near any of those links, the other factors
are units. Write $g_s=U_{i,s}h_{i,s}^{b_i}$ there.
On the boundary its differential factors as
$h_{i,s}^{b_i-1}(b_iU_{i,s}\,dh_{i,s}+h_{i,s}\,dU_{i,s})$.
The bracket has real rank two on the deformed link,
since $h_{i,s}|_{\partial K}$ is a submersion there,
and remains of rank two in a uniform neighbourhood.
Consequently $g_s|_{\partial K}$ is a submersion where
$0<|g_s|<\eta_0$, for a uniform sufficiently small
$\eta_0>0$. Away from these link neighbourhoods
$|g_s|$ has a uniform positive lower bound on the
boundary. This proves the asserted uniform
transversality, including at parameters for which
the germs inside $K$ change their singularities.

\item Take a small disc $D_x$ of radius $\rho$ with
$\rho^a<\eta_0/2$, then shrink $P$ so that every
critical value of $g_s$ on $K$ has modulus less than
$\rho^a/4$. This shrinking follows by compactness:
a sequence of critical points with parameters
tending to zero would otherwise limit to a
nonzero critical value of $g$ in the chosen small
representative. The latter has no such value.
Fix $0<\eta<\rho^a/4$ and put
$W=\overline D_x\times K$, a compact manifold
with corners contained in the original coordinate
chart in $Y$.

\item Every $H_s^{-1}(\epsilon)$ with
$0<|\epsilon|<\eta$ is transverse to this boundary.
On the face $\partial D_x\times K$, the equation
makes $|g_s|>3\rho^a/4$, so $dg_s\ne0$; its plane
directions give the submersion. On the face
$\overline D_x\times\partial K$, the $x$
derivative gives it when $x\ne0$. At $x=0$ we
have $g_s=-\epsilon\ne0$, and the preceding
boundary calculation gives it. At the corner
$|x|=\rho$, that same calculation for $g_s$
supplies both real target directions. The corner
stratum is therefore also transverse. We use
horizontal transport on this fixed domain
\cite[Proposition and Lemma, p.~154]{Ehresmann1952Connections}.
At a face or corner, adjoining its boundary defining
functions to the real and imaginary parts of $H_s$
gives a submersion. Complete these to local coordinates
by the implicit function theorem; the inverse coordinate
map lifts target tangent vectors while fixing every
boundary defining function. A partition of unity patches
these lifts, and compactness integrates them along the
whole path. The reverse path gives the inverse transport.
The resulting fibre diffeomorphisms preserve orientation
by continuity from the identity.

\item Inside $P\times\{0<|\epsilon|<\eta\}$ remove the
critical-value discriminant $\mathcal D$ of the
maps $H_s$ on $W$. This is a proper complex
analytic subset. The relative critical locus has
no point on the boundary at such a value, so its
projection $(s,p)\mapsto(s,H_s(p))$ is proper over
every compact subset of this parameter domain. The
proper mapping argument above therefore applies. For fixed $s$ its critical
values are discrete, since a holomorphic function
is constant on each connected critical component.
The complement of a proper complex analytic
subset of a connected complex manifold is path
connected. The transversality just proved and
the implicit function theorem give a proper
smooth submersion, including its boundary,
\begin{equation}
 \{(s,p):p\in W,\ (s,H_s(p))\in
                    (P\times\Delta_\eta^*)\setminus\mathcal D\}
 \longrightarrow
       (P\times\Delta_\eta^*)\setminus\mathcal D,
 \quad (s,p)\longmapsto(s,H_s(p)).                 \hypertarget{tag-006D}{}\label{cdp:eq:A.J.40}
\end{equation}
Ehresmann's theorem identifies all its compact
fibres by orientation-preserving diffeomorphisms
of manifolds with corners.
In particular a sufficiently small local Milnor
fibre of an interior singularity of some $H_s$,
at a value near zero, embeds into a smooth fibre
in \eqref{cdp:eq:A.J.40}, hence into
$W\cap H_0^{-1}(\epsilon_0)$ for a sufficiently
small nonzero $\epsilon_0$. This last space is
a submanifold of the global torus
$Y_{\epsilon_0^m}$ by \eqref{cdp:eq:A.J.38}.
Choose the value to avoid every other
critical value, and choose its local Milnor ball
in the interior of $W$. The exponents along the
boundary links remain $b_i$ throughout
\eqref{cdp:eq:A.J.40}.

\item Suppose first that $r\ge2$. Choose two distinct
factors $h_i,h_j$. Their map
$(h_i,h_j):(\C^2,0)\to(\C^2,0)$ is finite,
since their common zero germ is zero dimensional.
Its positive local degree is their intersection
multiplicity. A sufficiently small regular value
$(s_i,s_j)$ consequently has nonempty finite
inverse image. At each inverse-image point,
$h_i-s_i$ and $h_j-s_j$ are transverse coordinates.
Choose the remaining small shifts to avoid those
finitely many points. The deformation \eqref{cdp:eq:A.J.39}
then has an interior point at which
\[
            H_s=x^a-U v^{b_i}w^{b_j},
\]
with $U$ a unit independent of $x$. Replacing
$v$ by $vU^{1/b_i}$ makes this exactly
$x^a-v^{b_i}w^{b_j}$ by an invertible holomorphic
coordinate change: with $w$ fixed, its inverse solves
$vU(v,w)^{1/b_i}=v'$ uniquely near zero, since the
derivative is the nonzero chosen root of $U(0,0)$.
\ref{cdp:st:A.5.15} constructs a spherical
class of nonzero square in its local Milnor fibre.
The embedding \eqref{cdp:eq:A.J.40} puts that class in the
smooth torus fibre of $f=H_0^m$.
Preservation of intersection
numbers and $\pi_2$-vanishing give the same
contradiction as in \ref{cdp:st:A.5.15}. Thus $r=1$.

\item Suppose the reduced factor $h_0=h_1$ is singular.
Write $g=u h_0^b$ and put $h=u^{1/b}h_0$, using a local
root of the unit. Then $g=h^b$, with the same reduced
zero germ. Since $h$ is reduced, its gradient
has an isolated zero at the origin: a critical
curve would have constant value zero and would
be a repeated component of $h$. The finite
gradient map has a positive local degree.
Choose a small regular value
$(\lambda_1,\lambda_2)$ of this map. Then
\[
        h_s=h-\lambda_1y-\lambda_2z-\lambda_0
\]
has nondegenerate critical points near zero;
regularity of the gradient value says exactly
that their Hessians are invertible. Use
$g_s=h_s^b=u(u^{-1/b}h_s)^b$ in \eqref{cdp:eq:A.J.39};
thus $h_{0,s}=u^{-1/b}h_s$ is a deformation of $h_0$
with the original unit and multiplicity fixed. Choose
$\lambda_0$ so that their critical values are
nonzero. At one such point $p$ put
$v=h_s(p)\ne0$. The critical values tend to zero
with the parameters; choose them small enough that
$|v|^b<\eta/2$. The Hessian of $h_s^b$ is
$b v^{b-1}$ times the Hessian of $h_s$, and
is invertible. For $a=2$, the coordinate $x=\mathrm iX$
and target negation give the sign convention $X^2+h_s^b$.
Writing $Q=\operatorname{Hess}h_s(p)$, its critical value,
Hessian and determinant are exactly
\[
 v^b,\qquad
 \begin{pmatrix}2&0\\0&b v^{b-1}Q\end{pmatrix},\qquad
              2b^2v^{2b-2}\det Q\ne0.
\]
This identity holds for every positive integer $b$. In particular
for $X^2+(YZ-\varepsilon)^b$, the origin has critical value
$(-\varepsilon)^b$ and plane Hessian
$\left(\begin{smallmatrix}0&1\\1&0\end{smallmatrix}\right)$,
so the full Hessian determinant is $-2b^2\varepsilon^{2b-2}\ne0$.
The holomorphic Morse lemma on
the two plane variables identifies the germ of
$H_s-H_s(p)$ with
\[
                         x^a-Y^2-Z^2.
\]
Put $U=Y+\mathrm iZ$ and $V=Y-\mathrm iZ$; then
$Y^2+Z^2=UV$, with inverse
$Y=(U+V)/2$, $Z=(U-V)/(2\mathrm i)$. Thus the germ is
$x^a-UV$, and \ref{cdp:st:A.5.15} supplies its spherical
class of nonzero square.
Its Milnor fibre embeds by \eqref{cdp:eq:A.J.40}, producing
the same contradiction. Hence $h$ is smooth.
Take $y=h=u^{1/b}h_0$ and a complementary coordinate
$z$. At zero its differential is $u(0)^{1/b}dh_0$,
so the coordinate map is invertible; its inverse is
the unique solution of these two coordinate equations.
Substitution gives $g=y^b$ and the asserted transverse
product.

\item For the Hessian consequence, take a local $m$th
root of the map, absorbing its unit. If $H=v h$
with $v$ a unit, then at $h=dh=0$ one has
$\operatorname{Hess}H=v(0)\operatorname{Hess}h$.
A nonzero symmetric Hessian has a direction of
nonzero quadratic value. The first completing-square
coordinate map above splits off that direction and
gives $f=(x^2-g(y,z))^m$ exactly. If $g=0$,
the equation would not be reduced, so $g\ne0$.
At a critical point $dg(0)=0$. The result just
proved gives $g=y^b$ with $b\ge2$ in suitable
holomorphic coordinates, which is the constant
transverse $A_{b-1}$ product.

\end{enumerate}
\end{proof}

\section{Plane-curve conductors and rational boundary complements}\label{sec:II:plane-conductor-boundaries}

\subsection{Binary products and cubic jets}\label{subsec:II:normalization-local}



\begin{proposition}[A homogeneous binary-product grid is excluded]\hypertarget{tag-01BC}{}\label{cdp:st:A.5.21}
\StatementRecord{cdp:st:A.5.21}{proved}{01BC}{}

Let $f:Y\to C_0$ be a proper holomorphic map from a smooth
complex threefold with smooth nearby fibres complex
two-tori. Let $P(X,V)$ be a homogeneous polynomial of
degree $d\ge2$, a product of distinct complex linear
forms, and suppose $P(1,0)\ne0$.
There is no point at which the fibre map has
the local equation
\begin{equation}
             f=u(x,y,z)\,P(x,yz)^m,
             \qquad m\ge1,                         \hypertarget{tag-004W}{}\label{cdp:eq:A.BP.1}
\end{equation}
with $u$ a holomorphic unit. In particular this excludes
$\prod_{i=1}^d(x-a_i yz)$ for distinct constants $a_i$,
and every positive full-fibre multiplicity of that germ.

The same exclusion holds for the variable-coefficient
equation $f=u h^m$ with
\begin{equation}
 h=\sum_{i=0}^d a_i(x,y,z)x^{d-i}(yz)^i,
 \qquad P_0(X,V)=\sum_{i=0}^d a_i(0)X^{d-i}V^i,
                                                        \hypertarget{tag-0052}{}\label{cdp:eq:A.BP.4}
\end{equation}
provided $P_0$ has $d$ distinct projective roots and
$P_0(1,0)\ne0$. 
\end{proposition}


\begin{proof}

\begin{enumerate}[label=\arabic*.,leftmargin=*]
\item Put $h=P(x,yz)$. It is weighted homogeneous of degree
$2d$ for the positive weights $(2,1,1)$.
On a small simply connected coordinate neighbourhood
take a holomorphic logarithm of the unit and set
\begin{equation}
 \Phi(x,y,z)=\bigl(xu^{1/(md)},
                  yu^{1/(2md)},zu^{1/(2md)}\bigr).
                                                        \hypertarget{tag-0050}{}\label{cdp:eq:A.BP.2}
\end{equation}
Its derivative at the origin is invertible, and direct
homogeneity gives $h\Phi=u^{1/m}h$. Thus
$f=(h\Phi)^m$ exactly. Write $v=u^{1/(2md)}$ and solve
\[
 \tau^{2md}=u(\tau^{-2}X,\tau^{-1}Y,\tau^{-1}Z),
 \qquad \tau(0)=v(0).
\]
The derivative in $\tau$ at the origin is
$2md\,v(0)^{2md-1}\ne0$. The inverse of $\Phi$ is
$\Psi(X,Y,Z)=(\tau^{-2}X,\tau^{-1}Y,\tau^{-1}Z)$:
uniqueness of the chosen root gives $v\circ\Psi=\tau$
and $\tau\circ\Phi=v$, and hence both inverse identities.
It suffices to construct a
compact spherical cycle of nonzero square in $h=1$,
because weighted homotheties then put it inside an
arbitrarily small fibre of \eqref{cdp:eq:A.BP.1}.

\item The surface $h=1$ is smooth. Its derivatives are
$P_x,zP_v,yP_v$. A common zero would, by Euler's identity,
have $P_v\ne0$, hence $y=z=0$; at such a point
$P_x=d/x\ne0$ because $P(1,0)x^d=1$.
The curve $B=\{P(x,v)=1\}\subset\C^2$ is smooth:
Euler's identity $xP_x+vP_v=dP$ precludes a common
zero of its two partial derivatives on that level.
It is also connected. The map
\[
 B\longrightarrow\PP^1\setminus\{P=0\},
                       \quad(x,v)\longmapsto[x:v]
\]
is an \'etale cover of degree $d$. Locally a projective
representative has nonzero $P$-value, and its $d$
rescalings with $P=1$ give the sheets. On the chart $v\ne0$
the coordinate maps are $(x,v)\mapsto(T=x/v,v)$ and
$(T,v)\mapsto(Tv,v)$, with $P(T,1)v^d=1$.
On the chart $x\ne0$ they are $(x,v)\mapsto(x,\eta=v/x)$
and $(x,\eta)\mapsto(x,x\eta)$, with $x^dP(1,\eta)=1$.
These formulas give every local sheet and its inverse.
Around any
one of the simple projective zeros of $P$, that
value has a simple zero, so its inverse $d$th root
has monodromy a full $d$-cycle. This monodromy is
transitive on the sheets, proving connectedness.

\item The function $v$ on $B$ has exactly $d$ distinct simple
zeros: $P(x,0)=P(1,0)x^d$, and $P_x(x,0)=d/x$ at those points.
Apply \cref{lem:II:conic-bundle-matching-sphere} to this smooth
connected curve $B=\{P(x,v)=1\}$. Its surface $X_B$ is exactly
$\{P(x,yz)=1\}$ through the identity map on $(x,y,z)$. The
sphere of \eqref{cdp:eq:A.BP.3} therefore has square $-2$ in $h=1$.

\item The weighted homothety
$(x,y,z)\mapsto(\lambda^2x,\lambda y,\lambda z)$
maps $h=1$ biholomorphically to $h=\lambda^{2d}$
and shrinks the compact cycles and their isotopy
to the origin. Denote this homothety by $\delta_\lambda$;
its inverse is $\delta_{1/\lambda}$.
The map $\Psi\circ\delta_\lambda$ embeds a neighborhood of
the cycles in the fibre $f=\lambda^{2dm}$, with inverse
$\delta_{1/\lambda}\circ\Phi$ on its image.
Choose a sufficiently small nonzero value for which this global
fibre is smooth.
That fibre is a complex two-torus and has zero
$\pi_2$, so the spherical class has zero homology
image. The oriented codimension-zero embedding
preserves the intersection of the compact cycles,
giving both zero and $-2$ for its square. This
contradiction proves the statement.

\item We now prove the variable-coefficient assertion,
including the boundary comparison. Choose a local holomorphic
$m$th root $v$ of $u$. Replace $h$ by $vh$, each $a_i$
by $va_i$, and $P_0$ by $v(0)P_0$; then $f=h^m$ as an
equality of functions. The initial form is multiplied by
the nonzero scalar $v(0)$,
so its projective roots and its nonzero value at $(1,0)$
are preserved.
For a complex parameter $s$ set
\begin{equation}
 h_s=\sum_{i=0}^d a_i(s^2x,sy,sz)x^{d-i}(yz)^i.
                                                        \hypertarget{tag-0053}{}\label{cdp:eq:A.BP.5}
\end{equation}
This is holomorphic also at $s=0$, with
$h_0=P_0(x,yz)$. For $s\ne0$ there is the exact
identity
\begin{equation}
       h(s^2x,sy,sz)=s^{2d}h_s(x,y,z).           \hypertarget{tag-0054}{}\label{cdp:eq:A.BP.6}
\end{equation}
The leading coefficient $a_0$ is a unit. The implicit
function theorem applied to
$\sum_i a_i(s^2x,sy,sz)T^{d-i}$, treating $(s,x,y,z)$
as parameters, gives distinct holomorphic functions
$b_i(s,x,y,z)$ and the exact factorization
\begin{equation}
       h_s=a_0(s^2x,sy,sz)
                     \prod_{i=1}^d(x-b_i(s,x,y,z)yz)
                                                        \hypertarget{tag-0055}{}\label{cdp:eq:A.BP.7}
\end{equation}
on one small parameter and coordinate neighbourhood.
Each factor defines a smooth surface, since its
$x$ derivative is a unit. For $i\ne j$ the ideal
generated by the two factors equals $(x,yz)$:
write $g_i=x-b_i yz$ and $g_j=x-b_j yz$. Then
\[
 yz=\frac{g_i-g_j}{b_j-b_i},\qquad
 x=\frac{b_jg_i-b_ig_j}{b_j-b_i},\qquad
 g_i=x-b_i yz,\quad g_j=x-b_j yz.
\]
These inverse changes of generators prove the ideal equality.
Thus the zero level has precisely the same two
critical curves $D=(x,yz)=0$ for every $s$.

\item Choose a small closed ball $K$ in the coordinate
chart such that the two axes of $D$ meet its boundary
transversely and the zero-level surfaces away from
those axes are transverse to the boundary. At
$s=0$ this follows directly for the smooth graph
surfaces $x=b_i(0)yz$ and sufficiently small radii.
The real weighted radial vector $(2x,y,z)$ is tangent to each
of these graph surfaces, and its derivative on the squared
radius is $4|x|^2+2|y|^2+2|z|^2>0$. This proves the
stated transversality at $s=0$.
It persists for small $s$ by compactness away from
neighbourhoods of the two boundary links.
Near either link, use normal coordinates $(U,V)$
for the smooth critical curve and a parameter $w$
along it. The normal initial term of $h_s$ is a
degree-$d$ binary form with distinct roots.
On the compact parameterized link its normal
gradient therefore satisfies a uniform bound
\begin{equation}
 \|\partial_{U,V}h_s\|\ge c\|(U,V)\|^{d-1},
 \qquad |\partial_w h_s|\le C\|(U,V)\|^d
                                                        \hypertarget{tag-0056}{}\label{cdp:eq:A.BP.8}
\end{equation}
after the neighbourhood is shrunk. The first bound
follows because the gradient of a homogeneous
reduced binary form has no zero on the unit sphere;
higher normal terms have one additional order.
The leading forms vary over a compact set of reduced binary
forms, so their gradient norms have a uniform positive minimum
on the unit sphere. The second estimate follows from
$h_s\in(U,V)^d$.

\item Transversality of the critical curve to $\partial K$
allows a boundary coordinate with one real component
of $w$ solved as a smooth function of $(U,V)$ and
the remaining component. The correction to normal
differentiation tangent to this boundary is a
bounded multiple of $\partial_w h_s$.
By \eqref{cdp:eq:A.BP.8} it is smaller by one normal order than
the surjective complex normal derivative. As a real map
to $\mathbb R^2$, a complex linear functional with coefficient
vector $G$ has Gram matrix $\|G\|^2I_2$.
The boundary correction has norm at most $C'\|(U,V)\|^d$,
whereas both positive singular values of the normal derivative
are at least $c\|(U,V)\|^{d-1}$. Taking
$\|(U,V)\|<c/(2C')$ preserves surjectivity. Thus
$h_s|_{\partial K}$ is a real submersion at every
nearby point with $h_s\ne0$, uniformly in $s$.
Away from the two link neighbourhoods, compactness
and the preceding zero-level transversality give
the same conclusion for all sufficiently small
nonzero values.

\item Remove the proper critical-value discriminant from
the product of the parameter disc and a sufficiently
small punctured value disc. The critical locus over
that domain cannot meet the boundary, so its
projection is proper and its image is complex
analytic by the proper mapping theorem
\cite[Satz~22$'$, pp.~355--356, and Satz~23, p.~356]{Remmert1957Mappings}.
It is a proper
analytic subset because $h_0$ has no critical point on a
nonzero level. Its complement is connected: locally, generic
complex lines meet a proper analytic subset discretely, and
arcs in those lines can be detoured around the intersections.
Applying these local detours along a path in the connected
parameter--value domain gives a path in the complement. The level
sets in $K$ form a proper submersion with boundary
there. Here is the boundary-preserving transport map.
Along a smooth path $(s(t),\epsilon(t))$ in the complement,
let $A$ be the real derivative in the fibre coordinates,
restricted to $T\partial K$ at the boundary. It is surjective.
Choose a metric and use its right inverse
$R=A^*(AA^*)^{-1}$. Extend the boundary map to a smooth
map $\widetilde R$ on a collar; there $A\widetilde R$ remains
invertible after shrinking the collar, and
$\widetilde R(A\widetilde R)^{-1}$ is a right inverse
extending the boundary one. Combine it with interior right
inverses by a partition of unity. The result is tangent to $\partial K$
at the boundary and remains a right inverse. Integrate
\[
 \dot q=R\bigl(\dot\epsilon-(\partial_s h_s)\dot s\bigr).
\]
The chain rule gives $\frac{d}{dt}(h_{s(t)}(q(t))-
\epsilon(t))=0$. Properness and boundary tangency make
the flow exist over the whole compact path. Reversing the
path gives the inverse by uniqueness of solutions.
The vertical complex orientations vary continuously, so
these diffeomorphisms preserve orientation. This is the
horizontal-lift construction of Ehresmann transport
\cite[Proposition and Lemma, p.~154]{Ehresmann1952Connections},
with the boundary condition imposed as above.
Choose a small sphere cycle of the already proved
model $h_0=P_0(x,yz)$ in a fixed sufficiently small level
$\epsilon\ne0$. Compactness at $(0,\epsilon)$ permits a
sufficiently small nonzero $s_0$ for which the segment
$(ts_0,\epsilon)$, $0\le t\le1$, stays in the regular locus.
Transport along this segment to the level $h_{s_0}=\epsilon$. Writing $T$ for this transport, the embedding is
$\delta_{s_0}\circ T$, where
$\delta_{s_0}(x,y,z)=(s_0^2x,s_0y,s_0z)$.
Its inverse on its image is $T^{-1}\circ\delta_{1/s_0}$.
Equation \eqref{cdp:eq:A.BP.6} gives its fibre value
$s_0^{2dm}\epsilon^m$ for $f=h^m$.
It carries the spherical class of square $-2$
into a global torus. The same intersection and
asphericity contradiction proves \eqref{cdp:eq:A.BP.4}.

\end{enumerate}
\end{proof}

\begin{example}[A nodal triple curve with an aspherical affine fibre]\hypertarget{tag-P229}{}\label{ex:II:nodal-triple-affine-fibre}
\StatementRecord{ex:II:nodal-triple-affine-fibre}{proved}{P229}{}
The polynomial
\begin{equation}
                     f=x(yz)(x-yz)
                                                        \hypertarget{tag-008E}{}\label{cdp:eq:A.U.6}
\end{equation}
has ordinary-triple transversals along the smooth part of the nodal curve $D=(x,yz)=0$, and ordinary-double transversals along the punctured line $y=z=0$. Every fibre $f^{-1}(t)$ with $t\ne0$ is aspherical. Writing $f=P(x,yz)$ gives $P(X,V)=XV(X-V)$ and $P(1,0)=0$.
\end{example}
\begin{proof}
Put $v=yz$. The critical equations are
\[
 f_x=v(2x-v),\qquad f_y=zx(x-2v),\qquad f_z=yx(x-2v).
\]
Their common zero set is $D\cup\{y=z=0\}$. Near a point of $D\setminus\{0\}$, either $y$ or $z$ is a coordinate unit, so $(x,v)$ are transverse coordinates and $f=xv(x-v)$ has three distinct tangent lines. At $(x,0,0)$ with $x\ne0$, its quadratic normal term is $x^2yz$, giving the ordinary double transversal.

For $t\ne0$, let $B_t=\{(x,v):xv(x-v)=t\}$. The maps
\begin{equation}
 B_t\times\C^*\longrightarrow f^{-1}(t),
 \quad (x,v,y)\longmapsto(x,y,v/y),\qquad
 (x,y,z)\longmapsto(x,yz,y)
                                                        \hypertarget{tag-008F}{}\label{cdp:eq:A.U.7}
\end{equation}
are inverse holomorphic maps because $t\ne0$ forces $v$, $y$ and $z$ to be nonzero. The equations $v(2x-v)=x(x-2v)=0$ have no solution on $B_t$, so $B_t$ is a smooth affine curve. It and $\C^*$ are aspherical, and their product proves the assertion.

The deformation obtained by replacing $yz$ by $yz-s$ is
$f_s=x(yz-s)(x-yz+s)$. Along $y=z=0$ one has
$\partial f_s/\partial x=-s(2x+s)$. At a fixed nonzero $x$ this is nonzero for sufficiently small $s\ne0$, whereas that point belongs to the critical line for $s=0$. Thus this deformation changes that critical line on the boundary; the fixed-boundary comparison in \cref{cdp:st:A.5.17} requires the critical curves to remain there.
\end{proof}

\begin{example}[A vanishing cubic jet with a spherical fibre class]\hypertarget{tag-P230}{}\label{ex:II:vanishing-cubic-sphere}
\StatementRecord{ex:II:vanishing-cubic-sphere}{proved}{P230}{\cref{cdp:st:A.5.21}}
The homogeneous polynomial
\begin{equation}
                  h=xy\{z(x-y)+xy\}
\hypertarget{tag-0057}{}\label{cdp:eq:A.BP.9}
\end{equation}
has a cubic normal tensor along $(x,y)=0$ which vanishes at $z=0$. Its regular affine fibre $h=1$ contains a spherical class of square $-2$. Consequently every germ $u h^m$, where $u$ is a holomorphic unit and $m\ge1$, is excluded from a map with nearby smooth fibres complex two-tori.
\end{example}
\begin{proof}
On $xy\ne0$ put
\[
 v=x-y,\qquad Z=y(y+v)z+y^2(2y+v).
\]
Together with $y$ these are holomorphic coordinates, with inverse
$x=y+v$, $z=(Z-y^2(2y+v))/(y(y+v))$. Substitution gives
\begin{equation}
                              h=vZ+y^4.
\hypertarget{tag-004X}{}\label{cdp:eq:A.BP.10}
\end{equation}
Choose an arc between two roots of $y^4=1$ avoiding $y=0$ and the other roots. For continuous choices with $AB=1-y^4$ and $|A|=|B|=|1-y^4|^{1/2}$, the circles
\[
 v=\rho A e^{\mathrm i\theta},\qquad Z=\rho^{-1}B e^{-\mathrm i\theta}
\]
collapse at the endpoints to a matching sphere. Choose the positive constant $\rho$ small enough that $|v|<|y|/2$ along the arc. The sphere then lies in the stated coordinate chart. Each endpoint has linking number $-1$ by the coordinate calculation in \eqref{cdp:eq:A.BP.3}, so its square is $-2$.

Since $h$ has degree four, homothety shrinks the compact sphere into arbitrarily small neighbourhoods of the origin and sends its fibre value to the corresponding fourth power. For a unit $u$, choose a holomorphic root $a=u^{1/(4m)}$. The map $\Phi(x,y,z)=a(x,y,z)(x,y,z)$ satisfies $h(\Phi)^m=u h^m$ and has invertible derivative $a(0)I$ at the origin, so it is a local biholomorphism. Thus every full multiplicity has the same small spherical cycle. Its embedding in a nearby torus preserves its square, whereas its image in $H_2$ is zero because the torus has contractible universal cover. This contradicts the square $-2$.
\end{proof}

\begin{example}[A vanishing cubic jet with aspherical affine fibres]\hypertarget{tag-P231}{}\label{ex:II:vanishing-cubic-aspherical}
\StatementRecord{ex:II:vanishing-cubic-aspherical}{proved}{P231}{}
For
\begin{equation}
                     h=xz(y^2-x^2z).
\hypertarget{tag-004Y}{}\label{cdp:eq:A.BP.11}
\end{equation}
the cubic normal tensor along $x=y=0$ vanishes at $z=0$, and every fibre $h^{-1}(t)$ with $t\ne0$ is aspherical. The quadratic surface branch has normalization $(v,u)\mapsto(v,vu,u^2)$ and conductor divisor $2(v=0)+2(u=0)$ in the full reduced surface. The conductor divisors on the other two normalization branches are $2(y=0)+(z=0)$ on $x=0$ and $(x=0)+2(y=0)$ on $z=0$.
\end{example}
\begin{proof}
For $h=t\ne0$, both $x$ and $z$ are nonzero. The coordinate substitution $u=y/x$, with inverse $y=xu$, gives $h=x^3z(u^2-z)$. Forgetting $x$ is a finite etale triple cover of $E=\{(z,u):z(u^2-z)\ne0\}$, since $x^3=t/[z(u^2-z)]$ has three distinct nonzero roots. The map
\begin{equation}
 \C^*\times(\C\setminus\{1,-1\})\longrightarrow E,
                  \qquad(s,v)\longmapsto(s^2,sv)
\hypertarget{tag-004Z}{}\label{cdp:eq:A.BP.12}
\end{equation}
is a finite etale double cover, with inverse branches given by $s^2=z$, $v=u/s$, and deck involution $(s,v)\mapsto(-s,-v)$. Its source is aspherical. Covering maps identify the universal covers, proving the same assertion for $E$ and for $h^{-1}(t)$.

The quadratic branch $y^2=x^2z$ has image ring $R=\C\{v,u^2,vu\}$ in $B=\C\{v,u\}$. It is finite because $u^2=z$, and birational because $u=y/x$ off $x=0$. Write $B=\C\{v,u^2\}\oplus u\C\{v,u^2\}$. Membership of both $b$ and $ub$ in $R$ is equivalent to divisibility of both coefficients by $v$, so the intrinsic conductor is $vB$.

Apply \cref{lem:II:hypersurface-branch-conductor} in $R=\mathbb C\{x,y,z\}$ to the three factors $x,z,y^2-x^2z$. On the quadratic branch its normalization map $(v,u)\mapsto(v,vu,u^2)$ sends the other-factor product $xz$ to $vu^2$. Multiplication by this element sends the intrinsic conductor $vB$ to $v^2u^2B$, exactly as in the lemma. On $x=0$ the other-factor product is $zy^2$, and on $z=0$ it is $xy^2$; these two branches are smooth, so their intrinsic conductors are their full local rings. These ideals give the stated divisors; their reduced supports are normal-crossing pairs.
\end{proof}

\subsection{Tacnode conductor images}\label{subsec:II:normalization-higher-jets}

\begin{lemma}[Curve rings of delta at most two with at most two branches]\hypertarget{tag-P228}{}\label{lem:II:delta-two-curve-rings}
\StatementRecord{lem:II:delta-two-curve-rings}{proved}{P228}{}
A reduced analytic curve germ in a smooth threefold, of delta invariant at most two and with at most two normalization branches, is planar or has one of the two rings and embedded ideals in \eqref{cdp:eq:A.T.14}. Every planar case is a local complete intersection in the threefold. Each of the two nonplanar cases has delta invariant two.
\end{lemma}
\begin{proof}
\begin{enumerate}[label=\arabic*.,leftmargin=*]
\item For a unibranch
reduced curve ring $A\subset\C\{t\}$ of colength at
most two, an element of order one makes $A=\C\{t\}$.
Otherwise its least positive valuation is two or three:
a least valuation at least four leaves three linearly
independent missing coefficients. If it is three and
the colength is two, the inclusion
$A\subset\C+t^3\C\{t\}$ is equality, since the latter
already has colength two. If it is two, choose a
parameter with $t^2\in A$ by taking a square root of
the leading unit. The even part of every convergent
series belongs to $\C\{t^2\}\subset A$. Subtracting
it shows that the least odd element can be written
$t^{2k+1}v(t^2)$; division by its unit puts
$t^{2k+1}$ in $A$. Consequently
$A=\C\{t^2\}\oplus t^{2k+1}\C\{t^2\}$, of colength
$k\le2$. These are the planar cusp and the planar
$A_4$ curve, besides the smooth case.

\item For two branches, write their image rings as $A_1,A_2$.
The two projections of $A\subset A_1\oplus A_2$ are
surjective. Their kernels define a common finite
quotient $Q_0$ and an exact fibre-product description
$A=A_1\times_{Q_0}A_2$. Thus
$\delta=\delta_1+\delta_2+\operatorname{length}Q_0$.
If both branches are smooth, the quotient is
$\C\{t\}/(t^k)$ with $k\le2$; after changing the
second parameter, the two maps identify the same
$k$-jet. The map
\[
 x\longmapsto(t,s),\qquad
 y\longmapsto(t^k,-s^k)
\]
identifies this fibre product with
$\C\{x,y\}/(y^2-x^{2k})$: write a matching pair as
the sum of its average and $y$ times its difference
divided by $2x^k$. These are the node and tacnode.
If one branch is singular, the delta budget forces
one cusp, one smooth branch, and $Q_0=\C$.


\item An abstract plane curve ring in the given smooth
threefold is lci by the transverse-coordinate
construction in \ref{cdp:st:A.5.14}. Thus the only remaining
non-lci analytic rings, necessarily consuming the
whole delta budget at a single point, are
\begin{equation}
\begin{array}{ll}
 A=\C\{t^3,t^4,t^5\},
 &I=(y^2-xz,\ yz-x^3,\ z^2-x^2y),\\[2pt]
 A=\C\{t^2,t^3\}\times_{\C}\C\{s\},
 &I=(xz,\ yz,\ y^2-x^3).
\end{array}                                             \hypertarget{tag-007R}{}\hypertarget{tag-00XK}{}\hypertarget{tag-00XL}{}\SsixEquationAliases{\label{cdp:eq:C.RA.7}\label{cdp:eq:C.RA.8}}\label{cdp:eq:A.T.14}
\end{equation}
These are normal forms of the embeddings.
Their three displayed coordinate functions form a
basis of the cotangent space of each ring. Lift them
to $\mathcal O_{Y,x}$; their differentials are linearly
independent because the embedding dimension is three.
They are therefore ambient coordinates by the inverse
function theorem. The displayed ideals are their
kernels, obtained respectively by the parametrizations
$(t^3,t^4,t^5)$ and the union of $(t^2,t^3,0)$ with
$(0,0,s)$.
\end{enumerate}
\end{proof}
\begin{proposition}[Every horizontal tacnode image is elliptic]\hypertarget{tag-01B8}{}\label{cdp:st:A.5.18}
\StatementRecord{cdp:st:A.5.18}{proved}{01B8}{}
Let $f:Y\to C_0$ be a proper holomorphic map from a smooth
compact complex threefold with $b_2(Y)=0$, with connected
fibres and smooth general fibre a complex two-torus.
Let $mS$ be a full fibre with irreducible reduced support,
$\eta:T\to S$ its normalization, and $\pi:Z\to T$ a
resolution with an irrational ruling $p:Z\to B$.
Suppose every exceptional curve of $\pi$ is vertical and
the conductor divisor on $Z$ is $2C+V$, where $C$ is a
section of $p$ and $V$ is vertical. Let $D=\eta\pi(C)$.
Suppose the generic plane-curve transversal to $D$ in $S$ is an
ordinary tacnode.
Then $D$ is a smooth elliptic curve, $B=C$ is elliptic,
and $C\to D$ is an unramified double cover. The transversal
type along $D$ is everywhere the same ordinary tacnode.
The normalization $T$ is a smooth minimal ruled surface
over this elliptic curve. Its conductor is twice a
section, and $e(S)=0$. In particular its minimal
resolution is the identity.
\end{proposition}
\begin{proof}
\begin{enumerate}[label=\arabic*.,leftmargin=*]
\item If the normalization of $D$ has positive genus, \ref{cdp:st:A.5.14}
and \ref{cdp:st:A.5.5} make $D$ smooth elliptic and its transversal
type constant. There are two generic normalization
branches of a tacnode. A vertical component of $Z$
is rational, so it cannot dominate this elliptic curve.
Thus both branches lie on $C$, and $C\to D$ is an
unramified double cover: the constant local equation
$x^2-y^4=0$ has two separately smooth normalization
branches on every sufficiently small disc in $D$.
Riemann--Hurwitz gives $g(C)=g(B)=1$.
It remains to exclude the case $D^\nu=\Sigma\simeq\PP^1$.
\item The finite map $C\to\Sigma$ has degree at least two,
because $C\simeq B$ has positive genus. Since a generic
tacnode has exactly two normalization branches, this
degree is exactly two. Consequently no vertical
conductor component can also dominate $D$: its generic
points would supply an additional normalization branch.
Write $g=g(C)\ge1$ and let $\mathcal B\subset\Sigma$
be the reduced branch divisor of this double cover.
Riemann--Hurwitz gives
\begin{equation}
                   \deg\mathcal B=2g+2.             \hypertarget{tag-007M}{}\label{cdp:eq:A.T.1}
\end{equation}
\item By \ref{cdp:st:A.5.8}, the ruling descends to
$\widetilde p:T\to B$. Identify $C$ with $B$ using the
section, let $\phi:B\to D$ be its map to the
image curve, and put
\[
                    \psi=\phi\widetilde p:T\to D.
\]
Every vertical conductor component has constant $\psi$
value. Its image under $\eta$ differs from $D$, as just
proved, and meets $D$ in finitely many points. Thus
there is a finite set $A\subset D$ such that, for each
$s\in S$, either $\psi(\eta^{-1}(s))$ is a singleton
or all its values belong to $A$. Include in $A$ the
constant values on the vertical conductor components
and the values of $\phi$ over their finitely many
intersections with $D$. Fibres over the remaining
points of $D$ have their preimages on $C$, where all
the $\psi$ values equal that same point of $D$.

\item Choose an embedded finite tree $J\subset D$ containing
$A$, in a triangulation of the compact analytic curve.
It can be constructed by joining these finitely many
points by arcs, subdividing their intersections and
retaining a spanning tree. The quotient $q:D\to D/J$ contracts the contractible CW subcomplex $J$, so it is a homotopy equivalence by \cite[Propositions~0.16--0.17, pp.~15--16]{Hatcher2002}. Since
$q\psi$ is constant on every fibre of the finite
proper quotient $\eta$, there is a continuous map
\begin{equation}
 a:S\to D/J,\qquad a\eta=q\psi,\qquad a|_D=q.     \hypertarget{tag-007X}{}\label{cdp:eq:A.T.2}
\end{equation}
The last equality holds first on the dense regular
locus of $D$ and then everywhere by continuity.

\item Use the explicit triangulated neighbourhood retraction
$r:N\to S$ of \eqref{cdp:eq:A.Pi.3}. For a sufficiently small disc
$\Delta$ around the fibre value, properness gives
$U=f^{-1}(\Delta)\subset N$. The map $(ar)|_U$ is
surjective on fundamental groups because its restriction
to $D$ is the homotopy equivalence $q$. Perturbing loops
off the real-codimension-two analytic support gives
$\pi_1(U\setminus S)\twoheadrightarrow\pi_1(U)$.
The torus-bundle exact sequence over $\Delta^*$ therefore
gives the quotient maps
\begin{equation}
 \Z^4\rtimes\Z\twoheadrightarrow\pi_1(U)
       \twoheadrightarrow\pi_1(D/J)\simeq\pi_1(D).
                                                        \hypertarget{tag-0082}{}\label{cdp:eq:A.T.3}
\end{equation}
The first group is metabelian, so the last is metabelian.

\item For each singular point $x$ of $D$, let $r_x$ be the
number of its normalization branches and put
$\beta=\sum_x(r_x-1)$. The finite normalization map
identifies exactly those finite sets of points on
$\PP^1$; since it is a compact Hausdorff quotient,
this is also the underlying topological quotient.
Adjoining arcs between the points to be identified
gives a CW presentation homotopy equivalent to
$S^2\vee\bigvee^\beta S^1$. Hence
$\pi_1(D)$ is the free group $F_\beta$.
For $\beta\ge2$ this group surjects onto $S_4$ by
sending two free generators to $(1234)$ and $(12)$;
the commutator $A_4$ of that quotient is nonabelian.
Equation \eqref{cdp:eq:A.T.3} therefore forces $\beta\le1$.
By \ref{cdp:st:A.5.14} the total delta invariant $\delta=p_a(D)$
is at most two. If $s$ denotes the number of points
of $\Sigma$ over $\operatorname{Sing}D$, then
\begin{equation}
 s=\#\operatorname{Sing}D+\beta\le\delta+\beta\le3.
                                                        \hypertarget{tag-0083}{}\label{cdp:eq:A.T.4}
\end{equation}

\item First suppose that $D$ is a local complete intersection.
Use the bundles and jets of \ref{cdp:st:A.5.5}.
Let $\nu:\Sigma\to D$ be normalization, $I=\mathcal I_D$,
\[
 N=(\nu^*(I/I^2))^*,\qquad
 M=\nu^*(\mathcal O_Y(S)|_D).
\]
Thus $M$ is a degree-zero line bundle on $\Sigma$.
Let $\Xi$ be the curve-normalization
conductor divisor, $c=\deg\Xi=2\delta$, and let
$R\subset N$ and $Q=N/R$ be the saturated radical and
quotient of the generic rank-one Hessian. Write
$q_0=\deg Q$. The adjunction and conductor comparisons
\eqref{cdp:eq:A.J.1} give
\begin{equation}
 \deg N=-2+c,\quad \deg R=-2+c-q_0,\quad
 0\ne a\in H^0(\Sigma,M\otimes Q^{-2}),\quad q_0\le0.
                                                        \hypertarget{tag-0084}{}\label{cdp:eq:A.T.5}
\end{equation}
Here $a$ denotes the intrinsic quadratic coefficient,
whose zero divisor has degree $-2q_0$.
The quartic invariant $\Delta=ae-b^2/4$ of
\ref{cdp:st:A.5.5} is meromorphic in
$L^2$, where $L=M\otimes Q^{-1}\otimes R^{-2}$.
Its pole bound \eqref{cdp:eq:A.J.9} gives a nonzero holomorphic
section $\sigma$ of the square line bundle
$L^2(4\Xi)=(L(2\Xi))^2$. Its zero divisor has degree
\begin{equation}
               d_\Delta=8+2q_0.                    \hypertarget{tag-0085}{}\label{cdp:eq:A.T.6}
\end{equation}

\item The double cover $C\to\Sigma$ is obtained by adjoining
$\sqrt{-\Delta}$ to the function field, with the line
trivializations specified by $L$. Indeed over the
generic point a tacnode has normal expansion
$a\ell^2+b\ell r^2+er^4$, and the two normalization
branches are distinguished by the two roots of
$a(\ell/r^2)^2+b(\ell/r^2)+e$. Their discriminant
is $-4\Delta$. Changing a local frame of $L$ multiplies
the representing function of $\Delta$ by a square,
so this describes the same degree-two extension on
every overlap. At a point with local parameter $z$,
write that representing function as $z^n u(z)$,
$u(0)\ne0$. A root of the unit and removal of the
even power of $z$ show that its normalization is
ramified exactly when $n$ is odd. Adding $4\Xi$ changes
each order by an even number. Thus $\mathcal B$ is
exactly the odd-multiplicity support of
$\operatorname{div}(\sigma)$, and
\begin{equation}
                  2g+2\le 8+2q_0.                 \hypertarget{tag-0086}{}\label{cdp:eq:A.T.7}
\end{equation}

\item At a point of $D_{\rm reg}$ where $a\ne0$, \ref{cdp:st:A.5.17}
applies to the local $m$th root of the full-fibre
map. Its nonzero Hessian makes the local equation a
constant transverse product, necessarily of type
$A_3$ because this is the generic type along $D$.
The conductor cover is therefore unramified there.
Every branch point outside $\nu^{-1}(\operatorname{Sing}D)$
is consequently a zero of $a$. Equations \eqref{cdp:eq:A.T.4}--\eqref{cdp:eq:A.T.5}
give the second bound
\begin{equation}
                    2g+2\le -2q_0+s.              \hypertarget{tag-0087}{}\label{cdp:eq:A.T.8}
\end{equation}
Adding \eqref{cdp:eq:A.T.7} and \eqref{cdp:eq:A.T.8} yields
$4g+4\le8+s\le11$. Since $g\ge1$, we have
\begin{equation}
           g=1,\qquad \deg\mathcal B=4,
                    q_0\in\{-2,-1,0\}.            \hypertarget{tag-0088}{}\label{cdp:eq:A.T.9}
\end{equation}

\item Let a smooth point of $D$ be a
branch point at which $\Delta$ has a simple zero.
The preceding nonzero-Hessian argument gives $a(0)=0$.
In coordinates $D=(x,y)=0$ with curve parameter $z$,
choose the holomorphic saturated Hessian factor as the
first normal coordinate. The quadratic jet is then
$a(z)x^2$. The radical cubic vanishes identically,
so a local reduced equation has the exact expansion
\begin{equation}
 f_0=a(z)x^2+x^2(xA+yB)+xy^2b(z)+xy^3C+y^4E,
                                                        \hypertarget{tag-007N}{}\label{cdp:eq:A.T.10}
\end{equation}
with holomorphic coefficients. The identity
$\Delta=ae-b^2/4$, $a(0)=0$ and the simple zero of
$\Delta$ imply $b(0)=0$, $a'(0)\ne0$ and $e(0)\ne0$.
Indeed an order at least two for $a$, together with
$b(0)=0$, would give an order at least two for $\Delta$;
and then $\Delta'(0)=a'(0)e(0)$ is nonzero.
Thus $b/a$ is holomorphic and $E(0)\ne0$.

\item Set $X=x+b(z)y^2/(2a(z))$. In this analytic
coordinate change the coefficient of $Xy^2$ cancels.
All terms created from $x^2(xA+yB)$ lie in
$(X^2)+(Xy^3)+(y^4)$. Therefore, collecting the
terms of $X$-degree zero and one, gives exactly
\[
 f_0=X^2A_1(X,y,z)+Xy^3B_1(y,z)+y^4C_1(y,z),
 \qquad C_1(0)\ne0.
\]
Choose a fourth root of $C_1$ and define
\begin{equation}
 Y=yC_1^{1/4}+
                  \frac{XB_1}{4C_1^{3/4}}.
                                                        \hypertarget{tag-007O}{}\label{cdp:eq:A.T.11}
\end{equation}
The map $(X,y,z)\mapsto(X,Y,z)$ is locally
biholomorphic, and direct expansion gives
$f_0-Y^4\in(X^2)$. Hence $f_0=Y^4+X^2H(X,Y,z)$.
Its quadratic jet along $D$ is still $a(z)X^2$,
so $H(0,0,z)=a(z)$ and $\partial H/\partial z(0)\ne0$.
The further coordinate $Z=H$ gives exactly
\begin{equation}
                         f_0=Y^4+X^2Z.             \hypertarget{tag-007P}{}\label{cdp:eq:A.T.12}
\end{equation}
For the full-fibre map, first take a local
$m$th root of its unit. This multiplies the reduced
equation by a unit, which changes neither zero order
of $a$ nor of $\Delta$. The same coordinate construction
then makes the full map $(Y^4+X^2Z)^m$.
\Cref{cdp:st:A.5.15} excludes it. Consequently every
smooth branch point has order at least three in
$\operatorname{div}(\sigma)$.

\item At a singular point $x\in D$, suppose the intrinsic section $a$ is nonzero at every
normalization point above $x$. Choose a
constant linear combination of the two lci conormal
generators which is nonzero under the normalized
quotient at all those finitely many points. The
corresponding diagonal entry $a_0$ of the symmetric
Hessian matrix belongs to $\mathcal O_{D,x}$ and
has nonzero value on every normalization branch.
It is therefore a unit in $\mathcal O_{D,x}$.
The other normalized factor coefficient is the
quotient $h=q_{XY}/a_0\in\mathcal O_{D,x}$.
Thus square completion and monic polynomial division
in the proof of \eqref{cdp:eq:A.J.7}--\eqref{cdp:eq:A.J.8} both take place
over $\mathcal O_{D,x}$ itself. The regular-sequence
graded pieces inject into their normalization
pullbacks, so their cancellations already descend
on $D$. We may repeat that proof with $d=1$.
It gives a holomorphic section $\Delta$ of $L^2$
above $x$, without a conductor pole. At a branch
point above $x$, this proves the sharper inequality
\begin{equation}
 \operatorname{ord}\sigma\ge
                        4\operatorname{ord}\Xi+1.
                                                        \hypertarget{tag-007Q}{}\label{cdp:eq:A.T.13}
\end{equation}

\item If $q_0=-2$, the divisor of $\sigma$ has degree four
and has four odd-multiplicity points, so each is
simple. At most three lie above singularities of
$D$, by \eqref{cdp:eq:A.T.4}. The remaining smooth branch point
contradicts \eqref{cdp:eq:A.T.12}.
If $q_0=0$, the section $a$ has no zeros and all
four branch points would have to lie above the
singular locus; this contradicts $s\le3$.

\item Finally let $q_0=-1$, so that $\deg\operatorname{div}(a)=2$
and $d_\Delta=6$. If $k$ of the four branch points
lie above $\operatorname{Sing}D$, then the local
obstruction and parity give
\[
               6\ge k+3(4-k)=12-2k.
\]
Hence $k=s=3$, and equality holds: the three singular
preimages have order one in $\operatorname{div}(\sigma)$
and the one smooth branch point has order three.
The relations $s\le\delta+\beta\le3$ now force
$\delta=2$, $\beta=1$ and two singular points of
delta invariant one. The exact delta-one
classification in \ref{cdp:st:A.5.14} identifies them as one
ordinary node and one cusp. Their normalization
conductor orders are respectively $(1,1)$ and $2$;
these follow directly from the displayed node and
cusp rings in that proof.

\item At the cusp preimage, if $a$ were nonzero,
\eqref{cdp:eq:A.T.13} would make the order of $\sigma$ at least
$4\cdot2+1=9$, exceeding its total degree six.
So $a$ vanishes there. It also vanishes at the
smooth branch point. These two distinct points
exhaust the degree-two divisor of $a$. The section
is consequently nonzero at both node preimages.
Apply \eqref{cdp:eq:A.T.13} at the node: each of its two branch
points has order at least $4\cdot1+1=5$ in
$\operatorname{div}(\sigma)$, again exceeding the
total degree six. This excludes every rational lci image.

\item By \cref{lem:II:delta-two-curve-rings}, the remaining non-lci cases are the two analytic rings in \eqref{cdp:eq:A.T.14}. Each consumes the whole delta budget at one point. Use the ambient coordinates and ideals specified there.

The local reduced equation $f_0$ lies in $I^2$ in each case. In the first case put
\[
 g_1=y^2-xz,\quad g_2=yz-x^3,\quad
 g_3=z^2-x^2y,\quad
 F=x^5+xy^3+z^3-3x^2yz.
\]
The exact identities
\begin{equation}
 xF=g_2^2-g_1g_3,\qquad
 yF=xg_1^2+g_2g_3,\qquad
 zF=g_3^2+xg_1g_2                              \hypertarget{tag-007S}{}\hypertarget{tag-00XM}{}\SsixEquationAliases{\label{cdp:eq:C.RA.9}}\label{cdp:eq:A.T.15}
\end{equation}
show that the maximal ideal annihilates the class of
$F$ modulo $I^2$. Moreover
$I^{(2)}=I^2+(F)$ in the convergent local ring.
Here is a finite-module verification of that equality.
Let $J=I^2+(F)$. The quotient $\C\{x,y,z\}/J$ is
finite over the discrete valuation ring $\C\{x\}$,
since its reduction is the finite parametrized curve.
Its generic rank is nine: the curve has degree three
over the $x$ coordinate, and the square of its smooth
codimension-two ideal has length three along each
generic normal slice. Modulo $x$, the quotient has
the nine-element basis
\[
 1,y,y^2,y^3,\ z,yz,y^2z,\ z^2,yz^2,
\]
because its ideal is $(y^4,y^3z,y^2z^2,z^3)$.
A finite module over a discrete valuation ring is
the sum of its free and torsion parts; equality of
generic rank and residue-space dimension here makes
the torsion part zero. Thus $J$ is $x$-saturated.
By \eqref{cdp:eq:A.T.15}, $J$ and $I^2$ agree after inverting $x$.
The curve is lci away from the origin, so its symbolic
square differs from its ordinary square only by a
module supported there. Hence
$I^{(2)}=I^2:x^\infty=J$.

On its normalization let
$a_1=dg_1$ and $b_1=dg_2+t\,dg_1$.
They form a basis of the normalized conormal image;
$dg_3=-2t^2a_1+t b_1$.
The quadratic term of $F$ in this basis is
\begin{equation}
 q_F=\frac3t a_1^2-\frac3{t^2}a_1b_1
                           +\frac1{t^3}b_1^2,
 \qquad \det(q_F)=\frac3{4t^4}.                  \hypertarget{tag-007T}{}\hypertarget{tag-00WS}{}\SsixEquationAliases{\label{cdp:eq:C.RA.11}}\label{cdp:eq:A.T.16}
\end{equation}
Every equation in $I^2$ has a holomorphic quadratic
form in this same basis. If $f_0=cF+I^2$ had
$c(0)\ne0$, the determinant of its quadratic form
would have the noncancellable leading pole
$3c(0)^2/(4t^4)$: the cross terms with a holomorphic
form have pole at most three. It could not have
generic rank one. Thus $c(0)=0$, and \eqref{cdp:eq:A.T.15} gives
$f_0\in I^2$.

For the second ring put $G=y^2-x^3$ and
$P=(z,G)$, $Q_1=(x,y)$. Then $I=P\cap Q_1$ and
\begin{equation}
 I^{(2)}=P^2\cap Q_1^2=I^2+(zG).                \hypertarget{tag-007U}{}\label{cdp:eq:A.T.17}
\end{equation}
Indeed $P^2=(z^2,zG,G^2)$, while $zG,G^2$ already
belong to $Q_1^2$. Multiplication by $z^2$ is
injective modulo $Q_1^2$, so membership of
$z^2A+zGB+G^2C$ in $Q_1^2$ forces $A\in Q_1^2$.
This gives the asserted generating ideal directly.
Furthermore the following products lie in $I^2$:
\begin{equation}
 x(zG)=(xz)G,\qquad y(zG)=(yz)G,\qquad
 z(zG)=(yz)^2-x(xz)^2.
 \hypertarget{tag-00WT}{}\label{cdp:eq:C.RA.12}
\end{equation}
On the cusp branch use $a_1=d(xz)=t^2dz$ and
$b_1=dG=-3t^4dx+2t^3dy$.
The quadratic form of $zG$ is $t^{-2}a_1b_1$,
of determinant $-1/(4t^4)$. The same rank-one
argument forces its coefficient to vanish at the
origin, and the three identities again give
$f_0\in I^2$.

\item At these non-lci points, put
\[
 K=\operatorname{im}\bigl(\nu^*(I/I^2)
                         \longrightarrow\nu^*\Omega_Y^1\bigr),
 \qquad N=K^*.
\]
The coherent image $K$ is locally free of rank two
on the smooth curve $\Sigma$. Let $R_{\rm ram}$ be
the zero divisor of $d\nu$, and let $K_{\rm sat}$
be the kernel of the surjection
$\nu^*\Omega_Y^1\to\Omega_\Sigma^1(-R_{\rm ram})$.
The torsion quotient $K_{\rm sat}/K$ defines a divisor
$J_{\rm jac}$. Taking determinants in these
maps gives
\begin{equation}
 \deg N=-2+\deg(J_{\rm jac}-R_{\rm ram}).         \hypertarget{tag-007V}{}\label{cdp:eq:A.T.18}
\end{equation}
For $(t^3,t^4,t^5)$, the wedge $a_1\wedge b_1$
has coefficients $t^8(3,4t,5t^2)$, while $d\nu$
has common factor $t^2$. Thus the correction is
$8-2=6$. For the cusp branch in the second case,
the wedge is $-t^5(2,3t,0)$ and the tangent common
factor is $t$, giving correction four. On its
smooth branch the wedge is $s^2(dx\wedge dy)$,
giving correction two. In both cases \eqref{cdp:eq:A.T.18}
therefore gives $\deg N=4$.

Because $f_0\in I^2$, its quadratic form is a
holomorphic section of $\operatorname{Sym}^2K\otimes M$,
independent of its expression as a sum of products
of ideal generators. This follows either from its
Hessian or from injectivity of
$\operatorname{Sym}^2K\to\nu^*\operatorname{Sym}^2\Omega_Y^1$.
Its saturated radical and quotient again give line
bundles $R,Q$, with $q_0=\deg Q\le0$ and
$\deg R=4-q_0$. The generic quartic invariant
is meromorphic in $L^2$, where
$L=M\otimes Q^{-1}\otimes R^{-2}$; now
\begin{equation}
                        \deg L^2=-16+2q_0.        \hypertarget{tag-007W}{}\label{cdp:eq:A.T.19}
\end{equation}

\item We bound its poles using transverse affine
coordinate maps, with the curve parameter held
fixed during every normal differentiation. For the
space cusp use, for $t\ne0$,
\begin{equation}
 (u,v,t)\longmapsto
 \left(t^3,\ t^4+\frac{v}{3t^5},\
                   t^5-\frac{u}{t^3}+\frac{2v}{3t^4}\right).
                                                        \hypertarget{tag-007Y}{}\label{cdp:eq:A.T.20}
\end{equation}
Near each nonzero parameter, the inverse is the local cubic
root $t$ of $x$, together with
$u=2t^4y-t^3z-t^8$ and $v=3t^5(y-t^4)$.
Its two transverse derivatives are the dual frame to
$a_1,b_1$.
Each coefficient of a derivative has pole at most
five. The first normal derivatives of all three
$g_i$ are holomorphic, since
$dg_1=a_1$, $dg_2=b_1-ta_1$ and
$dg_3=-2t^2a_1+t b_1$.
Second normal derivatives have pole at most ten.
For $f_0=\sum c_{ij}g_i g_j$, the cubic Taylor
coefficients consequently have pole at most ten:
the only surviving terms use a first and a second
derivative of the $g_i$, or a first derivative of
$c_{ij}$ and two first derivatives of the $g_i$.
Every fourth derivative of a holomorphic ambient
function in \eqref{cdp:eq:A.T.20} has pole at most twenty.

Choose a holomorphic change of the $N$ frame adapted
to its saturated radical and quotient. It does not
increase these bounds. The radical cubic is zero
because the generic transversal is $A_3$.
Writing the relevant Taylor terms as
$a\ell^2+b\ell r^2+er^4$ therefore gives a
holomorphic $a$, a mixed cubic $b$ of pole at most
ten, and $e$ of pole at most twenty. The invariant
$\Delta=ae-b^2/4$ has pole at most twenty.
These are the same intrinsic jets on overlaps:
to weighted order four, where $\ell$ has weight
two and $r$ has weight one, an adapted change is
$\ell'=\alpha\ell+u r^2$, $r'=\beta r$.
Direct substitution multiplies $\Delta$ by
$\alpha^{-2}\beta^{-4}$; changes of the defining
section supply the factor $M^2$. Higher weighted
terms and changes of the curve parameter do not
alter these coefficients. This verifies the
transition law for the line $L^2$ in \eqref{cdp:eq:A.T.19}.

There is only one singular preimage $p$ in this case.
Thus $\Delta$ is a holomorphic section of the square
line $L^2(20p)$, whose degree is $4+2q_0$.
As in \eqref{cdp:eq:A.T.7}, the even pole correction preserves
branch parity, so
\[
 \deg\mathcal B\le4+2q_0,
 \qquad \deg\mathcal B\le-2q_0+1.
\]
The second bound again uses \ref{cdp:st:A.5.17} away from that
single singular image preimage. Adding them gives
$2\deg\mathcal B\le5$, contradicting
$\deg\mathcal B=2g+2\ge4$.

\item For the cusp-plus-line germ, the corresponding
affine coordinate maps are
\begin{equation}
 (u,v,t)\longmapsto
       \left(t^2,t^3+\frac{v}{2t^3},\frac{u}{t^2}\right),
 \qquad
 (u,v,s)\longmapsto\left(\frac us,\frac vs,s\right).
                                                        \hypertarget{tag-007Z}{}\label{cdp:eq:A.T.21}
\end{equation}
The first inverse uses the local square root $t$ of $x$,
with $u=t^2z$ and $v=2t^3(y-t^3)$; the second is
$s=z$, $u=zx$, $v=zy$. Their first generator derivatives
are holomorphic in the respective normalized conormal frames.
The same product differentiation gives cubic pole
bounds six and two, and the fourth derivative bounds
are twelve and four. Thus $\Delta$ has poles bounded
by $12p+4q$ at the two normalization preimages.
The line $L^2(12p+4q)$ is a square and has degree
$2q_0\le0$ by \eqref{cdp:eq:A.T.19}. A nonzero section has no
positive-degree odd divisor, whereas this double
cover has at least four branch points. This is the
last contradiction. Hence the first case applies: $D$ and
$C$ are elliptic and the transversal type is constant.

\item Factor $\pi$ through the minimal
resolution $\pi_{\min}:Z_{\min}\to T$, with the resulting
sequence of exceptional-curve blowdowns
$q:Z\to Z_{\min}$. Every curve contracted by $q$ is
vertical for $p$, by the stated assumption on $\pi$.
The equality of structure-sheaf direct images for each
blowdown, in local coordinates on $B$, gives a descended
ruling $p_{\min}:Z_{\min}\to B$ with
$p=p_{\min}q$. The image $C_{\min}$ is still a section,
and the pushed-forward effective conductor is
$D_{\min}=2C_{\min}+V_{\min}$ with $V_{\min}$ vertical.
By \ref{cdp:st:A.5.8}, the singularities of $T$ are rational.
The degree map on exceptional cycles and nef generation over an affine rational surface singularity are \cite[Proposition~4.14(ii), pp.~91--92, and the lemma in \S4.17, pp.~94--95, arXiv v1]{Reid1997Chapters}. On the present analytic minimal resolution the relative evaluation map is obtained as follows. Take a small Stein
representative $V$ of one singular point and put
$U=\pi_{\min}^{-1}(V)$. For a line bundle $L$ on $U$ with
$d_i=\deg(L|_{E_i})\ge0$ on the exceptional curves, choose
$d_i$ transverse curvettes through distinct smooth points
of each $E_i$, away from a prescribed point $z_0$ of the
exceptional set. Their images are analytic curve germs
in $V$; after shrinking $V$, their strict transforms give
an effective divisor $A$ on $U$ avoiding $z_0$, of degree
$d_i$ on every $E_i$.
Rationality, Leray and the Stein property give
$H^1(U,\mathcal O_U)=H^2(U,\mathcal O_U)=0$, so the
exponential sequence identifies $\operatorname{Pic}(U)$
with $H^2(U,\Z)$. A triangulated regular neighbourhood
$N\subset U$ of the exceptional rational tree retracts
onto that tree. The class of $L\otimes\mathcal O_U(-A)$
therefore vanishes on $N$: its degree on every component
is zero. Properness gives a smaller Stein representative
with inverse image $U'\subset N$. Restriction followed
by the exponential-sequence isomorphism on $U'$ gives
$L|_{U'}\simeq\mathcal O_{U'}(A)$. Its canonical section
is nonzero at $z_0$. Thus the relative evaluation
$\pi_{\min}^*\pi_{\min*}L\to L$ is surjective near the
exceptional set.
On the minimal resolution,
$\deg(\omega_U|_{E_i})=-2-E_i^2\ge0$.
Applying this construction to $L=\omega_U$ and using
$\pi_{\min*}\omega_{Z_{\min}}=\omega_T$ gives the
surjection $\pi_{\min}^*\omega_T\to\omega_{Z_{\min}}$.
Finite normalization duality identifies
$\omega_T=\eta^*\omega_S\otimes\mathcal I$, where
$\mathcal I$ is the conductor ideal. Hence its pullback
image is the invertible ideal $\mathcal O_{Z_{\min}}(-D_{\min})$.
The normalization-canonical comparison gives
$K_{Z_{\min}}+D_{\min}=\pi_{\min}^*\eta^*\omega_S$
as line bundles, and its rational Chern class is zero.

\Cref{cdp:st:C2d} gives a smooth normalization near the
inverse image of the elliptic curve $D$, with normal
line of degree zero on its unramified double cover.
The minimal resolution is an isomorphism there, so
$C_{\min}^2=0$. No other conductor or exceptional
curve meets it. Hence $V_{\min}\cdot C_{\min}=0$.
For a ruling blowdown $\sigma:W\to W'$, apply \cref{lem:II:canonical-boundary-blowup} with $Z_1=W$, $Z_2=W'$, $D_1=D_W$, $D_2=\sigma_*D_W$ and $m=d$. Its formula \eqref{cdp:eq:A.T.22} gives $a=d-1$. The minimal-resolution data above satisfy \cref{lem:II:elliptic-ruled-boundary}: the section has square zero, the effective vertical divisor is disjoint from it, the canonical-boundary class is zero and every exceptional curve is vertical. Hence $V_{\min}=0$, $Z_{\min}$ is a minimal ruled bundle and $\pi_{\min}$ is an isomorphism.
Finally $\eta$ is an isomorphism away from
its conductor section $C_T$ and the elliptic image
$D$. Additivity of the topological Euler number on
these decompositions gives
\begin{equation}
 e(S)=e(T)-e(C_T)+e(D)
      =e(\PP^1)e(B)-0+0=0.                       \hypertarget{tag-0081}{}\label{cdp:eq:A.T.23}
\end{equation}
The full multiple fibre has the same underlying
topological space and therefore the same Euler number.


\end{enumerate}
\end{proof}

\begin{corollary}[Connected double covers of the nonplanar delta-two images]\hypertarget{tag-P227}{}\label{cor:II:nonplanar-tacnode-cover}
\StatementRecord{cor:II:nonplanar-tacnode-cover}{proved}{P227}{\cref{cdp:st:A.5.18,cdp:st:A.5.17}}
Let $f:Y\to C_0$ be a proper map from a smooth compact complex threefold with $b_2(Y)=0$, connected fibres and smooth general fibre a complex two-torus. Let $mS$ be a full fibre with irreducible reduced support. Let $D\subset\operatorname{Sing}S$ have generic ordinary-tacnode transversal, normalization $\Sigma\simeq\PP^1$, and one singular point, of one of the two types in \eqref{cdp:eq:A.T.14}. The normalized conductor over $D$ has no connected component mapping to $\Sigma$ with degree two.
\end{corollary}
\begin{proof}
Use the same conormal-image bundle
$K=\operatorname{im}(\nu^*(I_D/I_D^2)\to\nu^*\Omega_Y^1)$ and $N=K^*$ as in \eqref{cdp:eq:A.T.18}. The quartic invariant gives the degree-two function-field extension by $\sqrt{-\Delta}$, as in the comparison preceding \eqref{cdp:eq:A.T.7}. For a connected cover $C\to\Sigma$, Riemann--Hurwitz gives branch degree $2g(C)+2$.

In the cusp-plus-line case, \eqref{cdp:eq:A.T.19} and the pole correction \eqref{cdp:eq:A.T.21} give a square line of degree $2q_0\le0$. Its nonzero section has no odd divisor, contradicting the positive branch degree. In the space-cusp case, \eqref{cdp:eq:A.T.20} gives
\[
 2g(C)+2\le4+2q_0,\qquad
 2g(C)+2\le-2q_0+1.
\]
Their sum forces $g(C)=0$, and the two inequalities then force the integer $q_0=-1$. The cleared quartic has degree two and exactly two simple branch points. Only one lies over the singular image point. The other is a smooth simple branch point, whose exact coordinate reduction \eqref{cdp:eq:A.T.10}--\eqref{cdp:eq:A.T.12} contradicts \cref{cdp:st:A.5.15}. The nonzero-Hessian input in this branch count is \cref{cdp:st:A.5.17}; all these local maps are those of the stated full torus fibre.
\end{proof}

\subsection{Ordinary-triple conductor images}\label{subsec:II:normalization-triple}

\begin{proposition}[Ordinary-triple images have arithmetic genus at most one]\hypertarget{tag-01B9}{}\label{cdp:st:A.5.19}
\StatementRecord{cdp:st:A.5.19}{proved}{01B9}{}

Let $f:Y\to C_0$ be a proper holomorphic map from a smooth
compact complex threefold with $b_2(Y)=0$, with connected
fibres and smooth general fibre a complex two-torus.
Let $mS$ be a full fibre with irreducible reduced support,
and let $D\subset\operatorname{Sing}S$ be an irreducible
curve with generic ordinary-triple plane transversal.
Then $D$ is a local complete intersection in $Y$ and
$p_a(D)\le1$. The possibilities for its underlying curve
are a smooth rational curve, a rational curve with one
ordinary node or cusp, and a smooth elliptic curve.
In the elliptic case the transversal type is constant.

When $D$ is singular rational, every component of the
normalized conductor curve dominating $D$ is rational,
and its map to $D^\nu\simeq\PP^1$ has degree one.
Consequently, in the irrational ruled setup of \ref{cdp:st:A.5.18}
with the generic transversal replaced by an ordinary
triple point, the image of the noncontracted horizontal
section is either smooth rational or smooth elliptic.
In the elliptic case the horizontal cover is unramified
of degree three, the normalization is smooth minimal
elliptic ruled, and the fibre has Euler number zero.
In the rational case the horizontal cover has degree
two or three.
\end{proposition}


\begin{proof}

\begin{enumerate}[label=\arabic*.,leftmargin=*]
\item By \ref{cdp:st:A.5.14}, $p_a(D)\le2$. If its normalization has
positive genus, that statement and \ref{cdp:st:A.5.5} already give
a smooth elliptic curve of constant transverse type.
Suppose its normalization $\nu:\Sigma\to D$ is
$\PP^1$, with total delta invariant $\delta\le2$.

\item First let $D$ be lci. Put
$N=(\nu^*(I_D/I_D^2))^*$ and
$M=\nu^*(\mathcal O_Y(S)|_D)$.
The reduced equation $f_0$ vanishes on $D$. Its first
normal jet vanishes generically and hence everywhere
in the locally free sheaf $I_D/I_D^2$, so $f_0\in I_D^2$.
Its quadratic jet likewise vanishes in $I_D^2/I_D^3$.
Thus $f_0\in I_D^3$, and its initial cubic $F_3$ is a
section of $\operatorname{Sym}^3N^*\otimes M$.
Adjunction and the curve conductor give
$\deg N=-2+2\delta$, while $\deg M=0$.
Its nonzero binary-cubic discriminant is a section
of $M^4\otimes(\det N)^{-6}$, of degree
\begin{equation}
                         12-12\delta.             \hypertarget{tag-0089}{}\label{cdp:eq:A.U.1}
\end{equation}
Hence $\delta\le1$.

\item We exclude every non-lci germ of delta at most two
by its symbolic cube. The one- and two-branch analytic classification is \cref{lem:II:delta-two-curve-rings}, with its embedded ideals in \eqref{cdp:eq:A.T.14}. There is one
additional possibility when three branches occur:
delta at least two is forced by identifying their
constants, so equality makes the ring exactly
\[
 \{(a,b,c)\in\C\{t_1\}\oplus\C\{t_2\}\oplus\C\{t_3\}:
                      a(0)=b(0)=c(0)\}.
\]
Its embedding in the smooth threefold is the
union of the three coordinate axes, with ideal
$I=(xy,xz,yz)$. Indeed the three branch parameters
give a basis of the cotangent space, so their lifts are ambient coordinates by the inverse-function argument of \cref{lem:II:delta-two-curve-rings}.
Thus there are exactly three non-lci cases to consider;
each consumes the whole delta budget at one point.

\item For the space cusp use $I=(g_1,g_2,g_3)$ and
$F$ from \eqref{cdp:eq:A.T.15}. Its symbolic-cube identity is
\begin{equation}
                         I^{(3)}=I^3+FI.           \hypertarget{tag-008A}{}\label{cdp:eq:A.U.2}
\end{equation}
Here is a verification in the convergent local ring.
The right-hand side agrees with $I^3$ after inverting
$x$, by \eqref{cdp:eq:A.T.15}. Its quotient is finite over
$\C\{x\}$, of generic rank $3\cdot6=18$: the quotient
of a normal plane by the cube of its maximal ideal has
the six monomials of total degree less than three. Modulo $x$, its ideal in $\C\{y,z\}$ is
\[
 (y,z)^6+(y^2z^3,yz^4,z^5).
\]
The fifteen monomials of total degree at most four,
together with $y^5,y^4z,y^3z^2$, give residue dimension
eighteen. The same discrete-valuation-module argument
used for \eqref{cdp:eq:A.T.15} makes the quotient $x$-torsion-free.
It is therefore the saturation of $I^3$, which equals
the symbolic cube since the curve is lci off the origin.
This proves \eqref{cdp:eq:A.U.2}.

\item For the cusp-plus-line ideal $I=(xz,yz,G)$ with
$G=y^2-x^3$, one has
\begin{equation}
                   I^{(3)}=I^3+(zG)I.             \hypertarget{tag-008B}{}\label{cdp:eq:A.U.3}
\end{equation}
To verify this directly, write $P=(z,G)$ and
$Q=(x,y)$. Their powers are primary because these
prime ideals are complete intersections. The
symbolic cube is $P^3\cap Q^3$.
Write an element of $P^3$ as
$z^3A+z^2GB+zG^2C+G^3E$.
The last two terms are already in $Q^3$.
Reducing the other terms modulo $Q^3$ gives
$zA+GB=0$ there. The constant-in-$(x,y)$ coefficient
of $B$ must therefore be divisible by $z$; write
$B=zB'+B_Q$, $B_Q\in Q$.
It follows that $A+GB'\in Q^3$.
Consequently
\[
 P^3\cap Q^3=(z^3Q^3,z^2GQ,zG^2,G^3)
                        =I^3+(zG)I,
\]
as claimed.

\item For the three axes, membership of
$x^a y^b z^c$ in the symbolic cube means
$a+b,a+c,b+c\ge3$.
If one exponent is zero, the other two are at
least three and the monomial belongs to $I^3$.
If all are positive, either at least two are at
least two, in which case it belongs to $(xyz)I$,
or two are one, which violates one of the three
inequalities. Hence
\begin{equation}
                         I^{(3)}=I^3+(xyz)I.       \hypertarget{tag-008C}{}\label{cdp:eq:A.U.4}
\end{equation}

\item In all three cases the generic ordinary-triple
equation $f_0$ belongs to $I^{(3)}$: on the smooth
locus of each curve branch its order in the two
normal variables is at least three. Define the
bundle by its conormal image,
\[
 K=\operatorname{im}\bigl(\nu^*(I/I^2)
                           \to\nu^*\Omega_Y^1\bigr),
 \qquad N=K^*.
\]
For the first two cases \eqref{cdp:eq:A.T.18} gives $\deg N=4$.
For the three axes, on the branch $(t,0,0)$ the
frame is $t\,dy,t\,dz$, with conormal determinant
order two and no tangent ramification. Each of the
three branches has the same correction, for total
six, so again $\deg N=-2+6=4$.

\item Use the affine maps \eqref{cdp:eq:A.T.20}--\eqref{cdp:eq:A.T.21}.
For the three axes add
\begin{equation}
              (u,v,t)\longmapsto(t,u/t,v/t)
                                                        \hypertarget{tag-008D}{}\label{cdp:eq:A.U.5}
\end{equation}
on the first branch; its inverse is $t=x$, $u=xy$, $v=xz$.
On the other branches use $(u/t,t,v/t)$ and $(u/t,v/t,t)$,
with inverses $(t,u,v)=(y,xy,yz)$ and $(z,xz,yz)$,
respectively. In these charts the
first normal derivatives of the ideal generators
are holomorphic in the specified conormal frame.
An $I^3$ term consequently has holomorphic initial
cubic. The extra term in each of \eqref{cdp:eq:A.U.2}--\eqref{cdp:eq:A.U.4}
has initial cubic equal to the quadratic term of
its displayed symbolic-square generator, multiplied
by the first derivative of an ideal generator.

\item For the space cusp that quadratic term is \eqref{cdp:eq:A.T.16},
so the cubic pole is at most three. For the
cusp-plus-line germ it is $t^{-2}uv$ on the cusp
and $s^{-1}v^2$ on the smooth branch, so the pole
bounds are two and one. For the three axes it is
$t^{-1}uv$ in \eqref{cdp:eq:A.U.5}, giving pole bound one on
each branch. Coefficients multiplying these terms
are holomorphic; their derivatives do not enter
the initial cubic because the lower normal jets
of the products vanish.

\item The initial cubic is intrinsic on the dense smooth
part of $D$, as all lower jets are zero. Its
transition under a normal-coordinate change is
therefore the cubic tensor in
$\operatorname{Sym}^3N^*\otimes M$; higher-order
coordinate terms make no contribution. Factoring a binary cubic into three linear forms
expresses its discriminant as the product of the squares
of their pairwise determinants. Multiplication of the
equation by a unit multiplies this discriminant by its
fourth power. Substituting old normal variables $w=P w'$
multiplies it by $(\det P)^6$. These are the coefficient
transitions of a meromorphic section
of $M^4\otimes(\det N)^{-6}$, of degree $-24$,
with pole divisor of total degree at most twelve
in each of the three cases. After this exact pole
correction its line has negative degree, at most
$-12$. It cannot have the required nonzero
holomorphic section. This excludes all non-lci
cases and proves that $D$ is lci with $p_a(D)\le1$.
The list of underlying curves follows from the
delta-one classification in \ref{cdp:st:A.5.14}.

\item We give the finite-cover comparison at a singular
rational image. Here $\delta=1$, so \eqref{cdp:eq:A.U.1} has
degree zero and its nonzero section is nowhere zero
on $\Sigma=\PP^1$. The cubic defines a subscheme
$R_3\subset\PP(N)$ of tangent-root lines. Every
fibre consists of three distinct points. In a
projective coordinate chart, its derivative with
respect to the root coordinate is nonzero, so the
implicit function theorem makes
$R_3\to\Sigma$ a finite etale map of degree three.
There is no component in a whole projective fibre,
since its nonzero discriminant forbids an identically
zero cubic there.

\item Over the generic smooth locus of $D$, blow up its
normal coordinates by $v=\lambda u$. The strict-transform
equation is
$f_0(u,\lambda u,t)/u^3=F_3(1,\lambda,t)+uH(u,\lambda,t)$.
The nonzero derivative in $\lambda$ at each root makes
the strict transform smooth there. Its finite map to $S$
near these fibres is the normalization, and its exceptional
intersection maps a conductor branch to its tangent line
$[1:\lambda]$. The reciprocal chart handles the other
projective coordinate. These maps give an isomorphism
between the function-field algebra of $R_3$ and
that of the normalized conductor curves dominating
$D$, after their maps are factored through $\Sigma$.
Both are normal finite curves over $\Sigma$.
Taking the integral closure of $\mathcal O_\Sigma$
in this same algebra therefore identifies the two
finite covers everywhere. Every connected finite
etale cover of the simply connected Riemann sphere
has degree one. Thus each of these conductor
components is rational and maps isomorphically to
$\Sigma$, proving the second assertion.

\item In the irrational ruled setting the horizontal
section has positive genus. Its finite map to a
rational image normalization must have degree at
least two. The preceding comparison excludes a
singular rational image, and the generic triple
transversal bounds its degree above by three.
The remaining rational image is therefore smooth,
with degree two or three. If the image is elliptic,
no vertical rational curve can dominate it, so all
three branches belong to the horizontal section.
\Cref{cdp:st:C2d} gives an unramified degree-three
cover, normal degree zero and constant type.
The ruling and its vertical exceptional curves satisfy
the hypotheses of the factorization $\pi=\pi_{\min}q$
in the proof of \ref{cdp:st:A.5.18}, where
$q:Z\to Z_{\min}$ and $\pi_{\min}:Z_{\min}\to T$
is the minimal resolution. Put $C_{\min}=q(C)$ and
$V_{\min}=q_*V$; the pushed-forward effective conductor
is $2C_{\min}+V_{\min}$. By \ref{cdp:st:C2d}, the
normalization is smooth near its elliptic conductor,
so $\pi_{\min}$ is an isomorphism there. Thus
$C_{\min}^2=0$, and the vertical conductor and exceptional
curves are disjoint from $C_{\min}$. The canonical
comparison and blowdown formula \eqref{cdp:eq:A.T.22}
therefore give a smooth minimal elliptic ruled
normalization with conductor twice the section.
The normalization Euler formula \eqref{cdp:eq:A.T.23}
gives zero.

\end{enumerate}
\end{proof}

\subsection{Rational normalizations}\label{subsec:II:normalization-rational}

\begin{proposition}[Rational normalization bounds with the exact constant-sheaf map]\hypertarget{tag-01B6}{}\label{cdp:st:A.5.16}
\StatementRecord{cdp:st:A.5.16}{proved}{01B6}{\SeriesRef{cdp:st:A.5.17}, \SeriesRef{cdp:st:A.5.12}, \SeriesRef{cdp:eq:A.RT.5}, \SeriesRef{cdp:st:A.5.9}, \SeriesRef{cdp:eq:A.RT.9}, \SeriesRef{cdp:st:A.5.3}, \SeriesRef{cdp:eq:A.RT.8}}

Let $f:Y\to B$ be a proper holomorphic map from a smooth
complex threefold, with connected fibres and smooth
general fibre a two-torus. Let $mS$ be a full fibre with irreducible
nonnormal reduced support. Write $\eta:T\to S$ for normalization,
$\pi:Z\to T$ for its minimal resolution, and $\Gamma\subset S$
for the conductor-image scheme. Suppose that $Z$ is rational,
$\omega_S$ is torsion and $\chi(S,\OO_S)=0$.
These two identities hold in particular in the compact
case $b_2(Y)=0$.
Put $Q=\eta_*\OO_T/\OO_S$ and $e=h^0(S,\omega_S)$.
Then $e\in\{0,1\}$, and
\begin{equation}
\begin{gathered}
 R^1\pi_*\OO_Z=0,\qquad
 H^1(T;\mathbb C)=H^1(T,\OO_T)=H^2(T,\OO_T)=0,\\
 H^0(S,Q)\xrightarrow{\ \simeq\ }H^1(S,\OO_S),\qquad
 h^0(S,Q)=1+e,\quad h^1(S,Q)=e,\quad
 \chi(\Gamma,\OO_\Gamma)=-1.
\end{gathered}
\hypertarget{tag-0071}{}\label{cdp:eq:A.RT.1}
\end{equation}
In particular $T$ has rational singularities. The coefficient map
\begin{equation}
 H^1(S;\mathbb C)\longrightarrow H^1(S,\OO_S)
\hypertarget{tag-007C}{}\label{cdp:eq:A.RT.2}
\end{equation}
is injective, so $b_1(S)\le1+e\le2$.

There is an effective integral divisor $D$ on $Z$ with
$\OO_Z(D)\simeq K_Z^{-1}$, whose strict-transform part is the
normalization conductor. It satisfies
\begin{equation}
 h^0(D,\OO_D)=h^1(D,\OO_D)=1,
 \quad D_{\rm red}\text{ connected},
 \quad p_a(D_{\rm red})\le1.
\hypertarget{tag-007F}{}\label{cdp:eq:A.RT.3}
\end{equation}
The conductor image is connected,
$p_a(\Gamma_{\rm red})\le2$, and
$e(\Gamma_{\rm red})\ge0$. Set-theoretically
$D_{\rm red}=\pi^{-1}(\eta^{-1}(\Gamma_{\rm red}))$.
Exceptional curves over points outside the conductor have
coefficient zero in $D$. With $E_p=(\pi^{-1}(p))_{\rm red}$ the exact Euler
formula is
\begin{equation}
 e(S)=e(Z)-e(D_{\rm red})+e(\Gamma_{\rm red})
       -\sum_{p\in\operatorname{Sing}T\setminus\eta^{-1}(\Gamma)}
                  \bigl(e(E_p)-1\bigr).
\hypertarget{tag-007G}{}\label{cdp:eq:A.RT.4}
\end{equation}
In the torus smoothing of the statement, \ref{cdp:st:A.5.17} excludes
every point in this sum. Hence $T$ is smooth outside the
normalization conductor, and the sum is zero.
\end{proposition}


\begin{proof}

\begin{enumerate}[label=\arabic*.,leftmargin=*]
\item Let $r_0>0$ trivialize $\omega_S^{\otimes r_0}$.
A nonzero section of $\omega_S$ has nonzero constant $r_0$th
power on the compact connected surface $S$, so it trivializes
$\omega_S$. Consequently $e=0$ or $1$, with $e=1$ precisely
when $\omega_S\simeq\OO_S$. The hypothesis $\chi(\OO_S)=0$
and Serre duality give $h^1(S,\OO_S)=1+e$.

For the asserted compact specialization, apply the canonical
evaluation maps in \eqref{cdp:eq:A.CD.1} to the same restriction
$f^{-1}(\Delta)\to\Delta$ and its same full divisor $mS$.
They use the canonical generator and the principal equation
$(t)=mS$, giving $\omega_S^{\otimes m}\simeq\OO_S$.
The Cartier sequence
$0\to\OO_Y(-S)\to\OO_Y\to\OO_S\to0$ has equal first two
Euler characteristics: $b_2(Y)=0$ makes
$c_1(\OO_Y(S))=0$ in rational cohomology, so their Chern
characters agree in Hirzebruch--Riemann--Roch.
Thus $\chi(\OO_S)=0$. These are the canonical-evaluation
and Cartier-sequence calculations of \cref{cdp:st:A.5.12}.

\item Apply the finite-duality inclusion \eqref{cdp:eq:A.RT.5} of \cref{lem:II:conductor-duality} to $\eta:T\to S$.
A nonzero global section of $\omega_T$ is therefore a nonzero
section of the torsion line $\eta^*\omega_S$ vanishing on its
nonempty conductor. Its $r_0$th tensor power would be a nonzero
holomorphic function on compact connected $T$ vanishing there,
which is impossible. Consequently
$H^2(T,\OO_T)=H^0(T,\omega_T)^\vee=0$.

\item For the proper resolution $\pi$, normality gives
$\pi_*\OO_Z=\OO_T$, and its fibres have dimension at most one.
The low-degree Leray sequence is
\begin{equation}
\begin{split}
0\longrightarrow H^1(T,\OO_T)\longrightarrow H^1(Z,\OO_Z)
 \longrightarrow H^0(T,R^1\pi_*\OO_Z)\\
 \longrightarrow H^2(T,\OO_T)\longrightarrow H^2(Z,\OO_Z)
 \longrightarrow0.
\end{split}
\hypertarget{tag-007I}{}\label{cdp:eq:A.RT.6}
\end{equation}
Rationality of $Z$ makes both its positive-degree groups zero.
Together with \eqref{cdp:eq:A.RT.5}, this gives $H^1(T,\OO_T)=0$ and
$R^1\pi_*\OO_Z=0$, since this last sheaf has finite support.
The fibres of $\pi$ are connected. Proper base change for the
constant sheaf gives $R^0\pi_*\mathbb C_Z=\mathbb C_T$.
The degree-one topological Leray edge map is consequently an
injection $H^1(T;\mathbb C)\hookrightarrow H^1(Z;\mathbb C)=0$.

\item In the normalization exact sequence
\begin{equation}
 0\longrightarrow\OO_S\longrightarrow\eta_*\OO_T
       \longrightarrow Q\longrightarrow0,
\hypertarget{tag-007J}{}\label{cdp:eq:A.RT.7}
\end{equation}
the map on $H^0$ is the identity on constants. The preceding
vanishings therefore make its connecting map
$H^0(Q)\to H^1(\OO_S)$ an isomorphism and its next connecting
map $H^1(Q)\to H^2(\OO_S)$ an isomorphism. These are the
maps and dimensions in \eqref{cdp:eq:A.RT.1}. In particular $\chi(Q)=1$.
The conductor-duality map of \cref{lem:II:conductor-duality}, applied to this normalization, is
\begin{equation}
 Q\simeq\omega_\Gamma\otimes\omega_S^{-1}|_\Gamma.
\hypertarget{tag-007K}{}\label{cdp:eq:A.RT.8}
\end{equation}
The conductor image is a proper Cohen--Macaulay curve, including
any nilpotent structure. Serre duality on it gives
$\chi(Q)=-\chi(\omega_S|_\Gamma)$. A torsion line bundle has
degree zero on every reduced component, and therefore also has
total degree zero on this possibly nonreduced curve. Curve
Riemann--Roch gives
$\chi(\omega_S|_\Gamma)=\chi(\OO_\Gamma)$.
This proves $\chi(\OO_\Gamma)=-1$.

\item Here is the asserted injectivity of the complex coefficient map.
In the category of sheaves of complex vector spaces set
\[
 \mathcal V=\operatorname{coker}
     (\mathbb C_S\longrightarrow\eta_*\mathbb C_T).
\]
The diagram from constants to holomorphic functions induces the
specific sheaf map
\begin{equation}
 \mathcal V\longrightarrow Q,
 \quad[(a_1,\ldots,a_r)]\longmapsto[(a_1,\ldots,a_r)] .
\hypertarget{tag-007L}{}\label{cdp:eq:A.RT.9}
\end{equation}
At a point $s\in S$, the $r$ entries refer to the connected
normalization germs above $s$. A tuple of branchwise constants
that comes from an element of $\OO_{S,s}$ has equal values on
all these germs, because evaluation at their common image point
is one scalar. It is therefore a diagonal tuple. This proves
that \eqref{cdp:eq:A.RT.9} is injective on every stalk.

\item A finite map has zero positive-degree direct images of the
constant sheaf, by proper base change on its finite fibres.
Since $H^1(T;\mathbb C)=0$ and the map on constants is an
isomorphism, the connecting map from
$0\to\mathbb C_S\to\eta_*\mathbb C_T\to\mathcal V\to0$
identifies $H^0(S,\mathcal V)$ with $H^1(S;\mathbb C)$.
Naturality of the connecting maps gives the commutative square
\begin{equation}
\begin{array}{ccc}
 H^0(S,\mathcal V)&\xrightarrow{\ \simeq\ }&H^1(S;\mathbb C)\\
 \big\downarrow&&\big\downarrow\\
 H^0(S,Q)&\xrightarrow{\ \simeq\ }&H^1(S,\OO_S).
\end{array}
\hypertarget{tag-0072}{}\label{cdp:eq:A.RT.10}
\end{equation}
Its left arrow is injective by \eqref{cdp:eq:A.RT.9}, proving \eqref{cdp:eq:A.RT.2}.

\ref{cdp:st:A.5.3} gives the effective integral divisor $D$ and the exact
line-bundle relation
$K_Z\otimes\OO_Z(D)\simeq(\eta\pi)^*\omega_S$.
The torsion line on the right is trivial on the smooth rational
surface: its Picard group is torsion-free and has zero identity
component. Hence $\OO_Z(-D)\simeq K_Z$. The divisor sequence
$0\to K_Z\to\OO_Z\to\OO_D\to0$, rationality of $Z$ and
Serre duality prove \eqref{cdp:eq:A.RT.3}. In detail the first positive
cohomology of $K_Z$ is zero, and $H^2(K_Z)=\mathbb C$;
the positive cohomology of $\OO_Z$ is zero. Thus $D$ is
connected. Restriction $\OO_D\to\OO_{D_{\rm red}}$ is
surjective on $H^1$, because the kernel is supported in
dimension one. Consequently $p_a(D_{\rm red})\le1$.

\item The support is determined by the exceptional intersection equations. On a minimal resolution every exceptional
irreducible curve $E$ has $K_Z\cdot E\ge0$. Indeed
$E^2<0$, and adjunction gives
$K_Z\cdot E=2p_a(E)-2-E^2$; a smooth rational
$(-1)$-curve is excluded by minimality. If $E$ has coefficient
zero in $D$, then
$0=(K_Z+D)\cdot E$ forces $K_Z\cdot E=D\cdot E=0$.
Thus $E$ is disjoint from $D_{\rm red}$.
Every exceptional fibre over the normalization conductor meets
its strict transform. Connectedness of that fibre then forces
all its exceptional curves to have positive coefficient in $D$.
Conversely, a positive-coefficient exceptional curve over a point
outside the conductor would lie in a component of $D_{\rm red}$
disjoint from every strict conductor curve. The connectedness
just proved, and the nonempty strict conductor, exclude this.
This proves the stated support equality and the zero-coefficient
assertion. It also proves connectedness of $\Gamma_{\rm red}$.

\item The bound on its arithmetic genus follows
directly from \eqref{cdp:eq:A.RT.8}. One has
$h^1(\Gamma,\omega_S|_\Gamma)=h^0(Q)=1+e$,
and restriction to $\Gamma_{\rm red}$ is surjective on $H^1$.
For $e=1$ the line is trivial, giving
$p_a(\Gamma_{\rm red})\le2$. For $e=0$ its $H^1$ has
dimension at most one, and degree-zero Riemann--Roch gives
$p_a(\Gamma_{\rm red})=1-h^0(\omega_S|_{\Gamma_{\rm red}})
+h^1(\omega_S|_{\Gamma_{\rm red}})\le2$.
Finally normalization and resolution are isomorphisms off the
conductor and the finitely many exceptional fibres.
Stratifying by these disjoint closed sets gives \eqref{cdp:eq:A.RT.4} by
additivity of the compactly supported Euler characteristic.
At a point outside the conductor, $T\to S$ is an
isomorphism and $S$ is normal. A normal complex surface
has isolated singularities. The reduced local equation
therefore has an isolated critical point whenever that
point is singular. Since the full fibre has equation
$u h^m$, the isolated clause of \ref{cdp:st:A.5.17} excludes it.
This proves the final assertion and eliminates the displayed
sum by a local comparison with the nearby torus.

\item We prove the asserted Euler bound for the conductor image,
including its possible positive-genus component. Put
$C=\Gamma_{\rm red}$. The normalization exact sequence for
the connected reduced curve $D_{\rm red}$ shows that the sum
of its normalized component genera is at most
$p_a(D_{\rm red})\le1$. Any component with genus one is
smooth, is unique, and meets the remaining rational curves
only in trees; a second attachment path joining distinct
points of that elliptic component would increase the
arithmetic genus to at least two. Each component of $C$ is
the finite image of a strict conductor component of
$D_{\rm red}$. Riemann--Hurwitz therefore shows that its
normalization is rational or elliptic, with at most one
elliptic-normalized component.

\item Write $n_0$ for the number of rational-normalized components
of $C$, $n_1\in\{0,1\}$ for the number of its
elliptic-normalized components, and $r_p,\delta_p$ for the
number of branches and the delta invariant at a singular
point. The exact coherent and topological normalization maps
give
\[
 \chi(\OO_C)=n_0-\sum_p\delta_p,
 \qquad e(C)=2n_0-\sum_p(r_p-1).
\]
For each reduced curve germ, independent constants on its
normalized branches give $\delta_p\ge r_p-1$.
The already proved arithmetic-genus bound and connectedness
give $\chi(\OO_C)\ge-1$, hence $e(C)\ge n_0-1$.
If $n_0\ge1$ this is nonnegative. If $n_0=0$, there is
one elliptic-normalized component and
$\sum\delta_p\le1$. A negative value could therefore only
be $-1$, from one ordinary node. Indeed its curve ring has
colength one in the direct sum of its two normalization
rings, and is contained in the subring of pairs with equal
constant value; equality of colengths makes these rings
equal. The two points over that node are distinct on the
elliptic normalization.

\item Here the hypersurface embedding excludes this last case.
At least one point $t\in T$ over the node receives both
local conductor branches. Otherwise partition the local
irreducible hypersurface branches according to which of
the two conductor branches they receive. Both parts are
nonempty. Two surface branches from opposite parts meet
in a curve, since their distinct equations in the regular
three-dimensional ambient local ring generate an ideal
of height two. That intersection curve belongs to the
conductor image, so is one of its two branches. Both
surface branches receive it, contradicting the partition.
For a locally irreducible hypersurface there is only one
normalization point and the assertion is immediate.

\item The connected resolution fibre $\pi^{-1}(t)$ lies in
$D_{\rm red}$ by the support equality proved above.
The strict conductor curves over the two branches are
elliptic, because a rational curve cannot map
nonconstantly to an elliptic curve. There is at most one
such curve on $D_{\rm red}$. If the two strict curves
were different, they would already give total genus two.
If they coincide, they meet $\pi^{-1}(t)$ at distinct
points: the map of that normal elliptic curve factors
through the normalization of its image, so one point
cannot map to both points over the node. The connected
resolution fibre then supplies an additional path joining
those points, outside the elliptic curve. Its components
are rational, since $T$ has rational singularities.
This produces a graph cycle in addition to the elliptic
genus, again forcing $p_a(D_{\rm red})\ge2$.
Both alternatives are impossible. Thus $e(C)\ge0$.

\end{enumerate}
\end{proof}

\subsection{Minimal rational boundary complements}\label{subsec:II:normalization-boundary}

\begin{definition}[Minimal-boundary contraction data]
\hypertarget{tag-P209}{}\label{def:II:minimal-boundary-data}
\StatementRecord{def:II:minimal-boundary-data}{proved}{P209}{}
Let $Z$ be a smooth compact rational surface and let $D$ be an effective
integral divisor with $\OO_Z(D)\simeq K_Z^{-1}$. Fix a sequence
\(Z=Z_N\xrightarrow{\sigma_N}\cdots\xrightarrow{\sigma_1}Z_0\)
of contractions of smooth rational $(-1)$-curves to a minimal
rational surface. Let $D_i$ be the pushforward of $D$ to $Z_i$.
For the exceptional curve $E_i$ of $\sigma_i$, write $a_i$ for
its coefficient in $D_i$ and $r_i$ for the multiplicity of
$D_{i-1}$ at the contraction centre. Put
\(N_0=\#\{i:r_i=1\}\).
The sequence and all its centres are part of these data.
\end{definition}

\begin{lemma}[A proper modification and its conductor pushout]\hypertarget{tag-P294}{}\label{lem:II:proper-conductor-pushout}
\StatementRecord{lem:II:proper-conductor-pushout}{proved}{P294}{\cite[Proposition~0.16, p.~15; \S4.G, p.~459; Theorem~1.20, p.~43]{Hatcher2002}}
Let $h:X\to Y$ be a proper continuous surjection of compact Hausdorff spaces, let $C\subset Y$ be closed, and put $A=h^{-1}(C)$. Suppose that $h$ restricts to a homeomorphism $X\setminus A\to Y\setminus C$. Give $X$ and $C$ finite CW structures for which $A\subset X$ is a subcomplex. With $j:A\hookrightarrow X$ and $k=h|_A:A\to C$, the canonical map
\[
 X\cup_A C\longrightarrow Y
\]
is a homeomorphism, and the double mapping cylinder of $j,k$ has this homotopy type. In particular,
\[
 e(Y)=e(X)-e(A)+e(C).
\]
For nonempty connected $X,A,C$, the maps induced by $h$ and by $C\hookrightarrow Y$ give
\[
 \pi_1(Y)\simeq\pi_1(X)*_{\pi_1(A)}\pi_1(C).
\]
When $j_*$ is an isomorphism, the inclusion $C\hookrightarrow Y$ induces an isomorphism on fundamental groups. Its inverse on the image of $\pi_1(X)$ is $k_*j_*^{-1}$, and on $\pi_1(C)$ it is the identity.

For disconnected $A$ or $C$, form the incidence graph with vertices the connected components of $X$ and $C$ and one edge for each connected component of $A$, using $j$ and $k$ for its endpoints. A spanning forest gives a van Kampen presentation: take the free product of the fundamental groups of the vertex spaces, add one stable letter $t_a$ for each edge outside the forest, and impose
\[
 t_a^{-1}j_{a,*}(u)t_a=k_{a,*}(u)
 \qquad(u\in\pi_1(A_a)),
\]
with $t_a=1$ on forest edges and with the vertex paths used to base the two displayed homomorphisms. This presentation is taken separately on each connected component. In particular, for connected $X$ and nonempty $C$, writing $r=b_0(A)$ and $s=b_0(C)$ gives exactly $r-s$ stable letters. The quotient by all vertex groups is the free group $F_{r-s}$.
\end{lemma}
\begin{proof}
The map from $X\sqcup C$ to $Y$, given by $h$ on $X$ and inclusion on $C$, identifies exactly the equivalence classes defining $X\cup_A C$. Outside $C$ this follows from the stated homeomorphism; over $c\in C$, the whole fibre $h^{-1}(c)\subset A$ is identified with the point $c$. The induced map is therefore a continuous bijection. Its domain is compact and its target is Hausdorff, so it is a homeomorphism.

The CW-pair hypothesis gives the homotopy extension property by \cite[Proposition~0.16, p.~15]{Hatcher2002}. Apply the first paragraph of \cite[\S4.G, p.~459]{Hatcher2002} to the diagram $C\xleftarrow{k}A\overset{j}{\hookrightarrow}X$: the canonical collapse from
\[
 M=\bigl(X\sqcup(A\times[0,1])\sqcup C\bigr)
       /\bigl((a,0)\sim j(a),\ (a,1)\sim k(a)\bigr)
\]
to $X\cup_A C$ is a homotopy equivalence. Composing it with the homeomorphism already proved restricts to $h$ on $X$ and to the prescribed inclusion on $C$; thus this equivalence uses the required maps. The two open cylinder halves retract onto $X,C$ and their intersection onto $A$. Mayer--Vietoris \cite[\S2.2, pp.~149--150]{Hatcher2002} therefore gives $e(M)=e(X)+e(C)-e(A)$.

For connected $X,A,C$, Theorem~1.20 of \cite[p.~43, proof pp.~44--46]{Hatcher2002} applied to the same cover gives the displayed pushout of groups with edge maps $j_*,k_*$. When $j_*$ is invertible, the maps $k_*j_*^{-1}$ on $\pi_1(X)$ and the identity on $\pi_1(C)$ agree on $\pi_1(A)$. They define the stated inverse to the homomorphism induced by $C\hookrightarrow Y$.

For several components, base the edge maps by paths in their vertex spaces and collapse a spanning forest of the cylinder paths. Successive applications of the same van Kampen theorem give the displayed presentation: transport along the remaining cylinder path $t_a$ identifies $t_a^{-1}j_{a,*}(u)t_a$ with $k_{a,*}(u)$. For connected $X$, surjectivity of $k$ makes the incidence graph connected, with $1+s$ vertices and $r$ edges, so it has $r-s$ edges outside a spanning tree. Killing the vertex groups leaves precisely the free group on these letters.
\end{proof}

\begin{proposition}[Minimal boundary complements and conductor images]
\hypertarget{tag-P210}{}\label{prop:II:minimal-boundary-complements}
\StatementRecord{prop:II:minimal-boundary-complements}{proved}{P210}{}
In the setting of \cref{cdp:st:A.5.16}, choose the contraction
data of \cref{def:II:minimal-boundary-data} for its divisor $D$.
Every reverse blowup has coefficient $a_i=r_i-1$. It increases
the complement Euler number by one for $r_i=1$ and leaves it
unchanged for $r_i\ge2$, giving \eqref{cdp:eq:A.RT.12}.
The minimal boundary classes and complement Euler numbers are given in
\eqref{cdp:eq:A.RT.13}--\eqref{cdp:eq:A.RT.16}; their negative
values occur exactly for three concurrent lines on $\mathbf P^2$
and the ruled boundaries in \eqref{cdp:eq:A.RT.17}.

For either of these two minimal-boundary types, $D_{\rm red}$
is simply connected and all its normalized components are
rational. Inclusion of the conductor image in $S$ induces
the group isomorphism in \eqref{cdp:eq:A.RT.19}; this group is
trivial or infinite cyclic, and
\(e(\Gamma_{\rm red})\ge2\).
Under the full CDP Betti hypotheses $b_1(Y)=b_2(Y)=0$,
the group is infinite cyclic.
\end{proposition}

\begin{proof}
\begin{enumerate}[label=\arabic*.,leftmargin=*]
\item At the contraction $\sigma:Z_i\to Z_{i-1}$ apply \cref{lem:II:canonical-boundary-blowup} with $Z_1=Z_i$, $Z_2=Z_{i-1}$, $D_1=D_i$, $D_2=D_{i-1}$ and $m=r$. Both canonical-boundary lines are trivial, so \eqref{cdp:eq:A.RT.11} gives $a=r-1$. Its complement formula increases the Euler number by one exactly when $r=1$ and leaves it unchanged when $r\ge2$. Iterating along the chosen contraction sequence proves the equality below, where $N_0$ counts the centres of multiplicity one. \begin{equation}
 e(Z\setminus D_{\rm red})
 =e(Z_0\setminus (D_0)_{\rm red})+N_0,
 \qquad N_0\ge0,
 \hypertarget{tag-0074}{}\label{cdp:eq:A.RT.12}
\end{equation}
Formula \eqref{cdp:eq:A.RT.12} concerns
the resolution complement. For the torus smoothing,
the isolated clause of \ref{cdp:st:A.5.17} has already removed the
off-conductor exceptional correction in \eqref{cdp:eq:A.RT.4}. For the minimal surface $\mathbb F_2=\PP(\OO\oplus\OO(2))$, its two disjoint bundle sections satisfy $C_0^2=-2$ and $C_\infty\sim C_0+2F$. The divisor $2C_\infty\sim2C_0+4F=-K_{\mathbb F_2}$ is disjoint from $C_0$, while the contraction sequence to this minimal model is empty and $N_0=0$.

\item On $\mathbf P^2$, the divisor $D_0$ is an effective cubic.
Its possible reduced supports and their complement Euler
numbers are
\begin{equation}
\begin{array}{c|c}
 \text{reduced support}&e(\mathbf P^2\setminus(D_0)_{\rm red})\\ \hline
 \text{smooth cubic, nodal cubic, cuspidal cubic}&3,\ 2,\ 1\\
 \text{line and conic, transverse or tangent}&1,\ 0\\
 \text{three lines forming a triangle, or concurrent}&0,\ -1\\
 \text{two distinct lines, or one line}&0,\ 1.
\end{array}
\hypertarget{tag-0075}{}\label{cdp:eq:A.RT.13}
\end{equation}
For example the union of three concurrent lines has Euler
number $6-3+1=4$, whereas $e(\mathbf P^2)=3$.
The last row retains respectively a double line plus a line,
and a triple line. These are all cubic possibilities: an
irreducible singular cubic has arithmetic genus one and
therefore exactly a node or a cusp; the other degrees partition
three as $2+1$ or $1+1+1$.

\item On $\mathbb F_n$, with $n=0$ or $n\ge2$, write $C_0^2=-n$
for a minimal section and $F$ for a ruling fibre. Then
\begin{equation}
 D_0\sim 2C_0+(n+2)F.
 \hypertarget{tag-0076}{}\label{cdp:eq:A.RT.14}
\end{equation}
Its horizontal part has degree two over the ruling base.
If that part consists of two distinct sections $C_1+C_2$,
write $V$ for the vertical part, $d$ for its total degree and
$k=C_1\cdot C_2$. Adjunction on either section gives
$k+d=2$. Let $q$ be the number of distinct intersection
points of $C_1$ and $C_2$, and let $r$ count those points
whose ruling fibre occurs in $V_{\rm red}$. A vertical fibre
meets the two sections in two points unless they coincide at
that fibre. Additivity therefore gives the exact formula
\begin{equation}
 e((D_0)_{\rm red})=4-q+r,
 \qquad e(\mathbb F_n\setminus(D_0)_{\rm red})=q-r\ge0.
 \hypertarget{tag-0077}{}\label{cdp:eq:A.RT.15}
\end{equation}
This includes two sections and a fibre through their common
point: that union has Euler number four and complement Euler
number zero.

If the horizontal part is an irreducible bisection $C$, the
same adjunction calculation gives $p_a(C)=1-d$, so $d=0$
or $1$. For $d=0$, its Euler number is at most two, and the
complement Euler number is at least two. For $d=1$, the
bisection is a smooth rational curve, the vertical part is one
reduced fibre, and their set-theoretic intersection has one or
two points. The complement Euler number is respectively one
or two.

The remaining horizontal part is $2C$, where $C$ is a section.
Write $C\sim C_0+kF$. The vertical divisor has total degree
$d=n+2-2k$, and, with $v$ its number of distinct fibres,
\begin{equation}
 e((D_0)_{\rm red})=2+v,
 \qquad e(\mathbb F_n\setminus(D_0)_{\rm red})=2-v.
 \hypertarget{tag-0078}{}\label{cdp:eq:A.RT.16}
\end{equation}
If $C\ne C_0$, positivity of intersection with $C_0$ gives
$k\ge n$; hence $d\le2-n$ and $v\le2$.
For $n=0$ one also has $d\le2$. Consequently the only
negative values on a minimal ruled surface occur in the
explicit configurations
\begin{equation}
 Z_0=\mathbb F_n,\quad n\ge2,\qquad
 D_0=2C_0+\sum_{j=1}^{v}a_jF_j,\qquad
 a_j\ge1,\quad \sum_j a_j=n+2,\quad v\ge3.
 \hypertarget{tag-0079}{}\label{cdp:eq:A.RT.17}
\end{equation}
The fibres $F_j$ here are distinct.

\item There is a further fundamental-group restriction on precisely
these negative minimal-model configurations. Their reduced
supports are unions of rational curves with simply connected
underlying spaces: three concurrent lines on $\mathbf P^2$,
or the section with its attached ruling fibres in \eqref{cdp:eq:A.RT.17}.
Every reverse blowup in \eqref{cdp:eq:A.RT.11} preserves this property.
A multiplicity-one blowup leaves the support isomorphic to
the preceding one. A positive-coefficient blowup at a smooth
point adds a rational leaf, at an ordinary double point
subdivides the incidence edge, and at the initial concurrent
point replaces that vertex by an exceptional $\mathbf P^1$
with three attached branches. Subsequent supports have only
the resulting ordinary double points. None of these explicit
operations introduces a topological loop. Consequently
$D_{\rm red}$ is simply connected throughout these cases.

Apply \cref{lem:II:proper-conductor-pushout} to $(X,Y,A,C,h)=(Z,S,D_{\rm red},\Gamma_{\rm red},\eta\pi)$. Proposition~\ref{cdp:st:A.5.16} gives the full preimage equality and the isomorphism of complements. A compatible triangulation of the compact analytic pair $(Z,D_{\rm red})$ makes the inclusion a subcomplex \cite[Theorem~2, p.~464]{Lojasiewicz1964}; the compact conductor curve is triangulable as well. The resulting pushout uses the maps
\begin{equation}
 D_{\rm red}\hookrightarrow Z,
 \qquad D_{\rm red}\xrightarrow{\ h|_{D_{\rm red}}\ }
                           \Gamma_{\rm red}.
 \hypertarget{tag-007A}{}\label{cdp:eq:A.RT.18}
\end{equation}
The connected clause of that lemma, with $Z$ rational and $D_{\rm red}$ simply connected as just proved, gives
\begin{equation}
 \pi_1(S)\simeq
 \pi_1(Z)*_{\pi_1(D_{\rm red})}\pi_1(\Gamma_{\rm red})
 \simeq\pi_1(\Gamma_{\rm red}).
 \hypertarget{tag-007B}{}\label{cdp:eq:A.RT.19}
\end{equation}
All normalizations of the image components are $\mathbf P^1$,
by Riemann--Hurwitz. Lemma~\ref{lem:II:curve-normalization-graph} therefore identifies the fundamental group of this connected curve with the free group of rank $b_1(\Gamma_{\rm red})$.

\item Restrict the proper map in \cref{cdp:st:A.5.16} to a disc
about this full fibre $mS$. Its total space is smooth, its other
fibres are two-tori, and $S$ is irreducible and reduced.
Thus \cref{night:P4R.1} applies to the same degeneration and makes
$\pi_1(S)$ virtually abelian. Through \eqref{cdp:eq:A.RT.19},
\cref{lem:II:curve-normalization-graph} applies to the same rationally
normalized conductor image and bounds its free-group rank by one.
We obtain
\begin{equation}
 \pi_1(S)\simeq\pi_1(\Gamma_{\rm red})\in\{1,\mathbb Z\},
 \qquad b_1(\Gamma_{\rm red})\le1.
 \hypertarget{tag-007D}{}\label{cdp:eq:A.RT.20}
\end{equation}
Under the full CDP Betti hypotheses, the nonzero class on
every support from \ref{cdp:st:D.1.9} excludes the trivial group in
\eqref{cdp:eq:A.RT.20}. 

\item For these two boundary classes one can further exclude a
one-component image with first Betti number one. Its normalization
is $\mathbf P^1$, and exactly one point $q$ identifies two distinct
normalization points $\alpha,\beta$. Other singular points are
unibranch and do not change this topological quotient. The universal
cover of the image is an infinite chain of normalized spheres $C_j$,
with $\beta_j$ identified with $\alpha_{j+1}$.

The map $h|_{D_{\rm red}}$ lifts to this cover because
$D_{\rm red}$ is simply connected. Each noncontracted rational
component factors through the normalization of the image and maps
surjectively onto it; its lift therefore covers one whole $C_j$.
Each contracted component lifts to a point. There are finitely many
components, so their connected lifted image is a finite chain
$C_m,\ldots,C_M$.

Choose a preimage in $D_{\rm red}$ of the left endpoint node and
let $t_-$ be its image under $\pi$ in the normalization $T$ of $S$.
The resolution fibre $\pi^{-1}(t_-)$ is connected and lies in
$D_{\rm red}=h^{-1}(\Gamma_{\rm red})$. Its lift is a single point
of the discrete fibre over $q$. Every local conductor branch of $T$
through $t_-$ has a strict transform meeting this resolution fibre.
It can receive only the branch directed into $C_m$: receiving the
other branch would force the lifted image to contain $C_{m-1}$.
At the right endpoint the same argument produces a point $t_+$
receiving only the opposite branch. The two points are distinct,
because their connected resolution fibres lift to different endpoint
nodes, including when $m=M$.

Distinct normalization points over $q$ correspond to distinct local
irreducible hypersurface branches of $S$ in the smooth threefold.
Their two distinct equations have an intersection curve: the ideal
they generate has height two. Every such curve is a
conductor-image branch received by both normalization points.
There are only the two image branches at $q$, whereas $t_-$ and
$t_+$ receive disjoint ones. This contradiction excludes the
one-component, first-Betti-number-one image.

Since every image component is rational-normalized, write
$n_\Gamma$ for their number and $\beta_\Gamma\le1$ for the
first Betti number. Their exact topological normalization
formula now yields, on these negative boundary classes,
\begin{equation}
 e(\Gamma_{\rm red})=n_\Gamma+1-\beta_\Gamma\ge2.
 \hypertarget{tag-007E}{}\label{cdp:eq:A.RT.21}
\end{equation}
For one component the excluded case is the only value below
two; for at least two components the inequality follows
directly from $\beta_\Gamma\le1$.
\end{enumerate}
\end{proof}

\section{The global normalization criterion}\label{sec:II:global-normalization-criterion}

\subsection{Global cohomology tests}

\begin{proposition}[A finite cohomology test and a fixed fibre]\hypertarget{tag-01FM}{}\label{cdp:st:C4c}
\StatementRecord{cdp:st:C4c}{proved}{01FM}{}

Let $f:Y\to\PP^1$ be a proper connected-fibre map from a smooth compact complex threefold with $b_1(Y)=b_2(Y)=0$ and smooth general fibre a complex two-torus. Let $\Sigma$ be its finite set of critical values, and fix a connected component $P$ of $\Pic(Y)$. For $c\in\Sigma$, write $F_c$ for the full Cartier fibre with its multiplicities, and put
\begin{equation}
 Z_c=\{L\in P:H^0(F_c,\Omega_Y^1|_{F_c}\otimes
                  (K_Y\otimes L^{-1})|_{F_c})\ne0\}.
                                                               \hypertarget{tag-008W}{}\label{cdp:eq:C.18}
\end{equation}
Then the following are equivalent: every $Z_c$ is a proper subset of $P$; and $H^2(Y,T_Y\otimes L)=0$ on a dense open subset of $P$. In that case the generic twist in $P$ satisfies
\begin{equation}
 H^q(Y,T_Y\otimes L)=0\quad(q\ne1),\qquad
 \chi(Y,T_Y\otimes L)\le0,\qquad c_3(Y)\le0.          \hypertarget{tag-008X}{}\label{cdp:eq:C.19}
\end{equation}
For a reduced fibre, membership in the complement of $Z_c$ is exactly the two explicitly typed kernel-and-connecting-map conditions of \ref{cdp:st:A.1.1}, with kernel $N^*\otimes\mathcal J^{-1}$ computed in \ref{cdp:st:A.2.1}. For $F_c=mS$, the equivalent complete test is the finite lifting filtration of \ref{cdp:st:A.3.1}. These tests use only sheaves and maps on the finitely many singular full fibres.

When $b_3(Y)=0$, at least one fixed critical fibre $F_c$ satisfies $Z_c=P$. Its support is irreducible, by \ref{cdp:st:A.3.5}. Under the integral assumption $H^2(Y,\Z)=0$, the Picard group is connected, so this is a single fixed fibre for which the ambient group in \eqref{cdp:eq:C.18} is nonzero for every line bundle on $Y$.
\end{proposition}


\begin{proof}
\begin{enumerate}[label=\arabic*.,leftmargin=*]
\item The universal Picard family from \ref{cdp:st:C.1.1} restricts to a holomorphic family on the proper space $F_c\times P$. Its underlying sheaf is flat over $P$, because it is a line-bundle twist of the fixed locally free restriction $\Omega_Y^1|_{F_c}$. Proper semicontinuity therefore makes $Z_c$ a closed complex analytic subset of $P$.

\item Assume that each $Z_c$ is proper. The complement of their finite union is open dense and has full Lebesgue measure in $P\simeq\C^q$. Proposition \ref{cdp:st:D.1.11} supplies a dense set of full measure of line bundles non-torsion on every full fibre. Choose $L$ in the intersection. On every smooth torus fibre, $K_Y|_F$ and the conormal are trivial. The intrinsic cotangent bundle is trivial as well, so the conormal exact sequence expresses the ambient cotangent restriction as an extension of trivial bundles. The line $L|_F$ has zero integral Chern class, since $c_1(L)$ is torsion and $H^2(F,\Z)$ is torsion-free; it is nontrivial by the chosen non-torsion condition. Its cohomology in every degree vanishes, as follows by the unitary Fourier resolution of a topologically trivial line on a torus. The ambient section group is consequently zero at every smooth fibre. At each critical fibre it is zero by the definition of the chosen complement.

\item Serre duality on each Gorenstein Cartier fibre, with $\omega_{F_c}=K_Y|_{F_c}$, and top-degree base change give $R^2f_*(T_Y\otimes L)=0$. Twisting by $f^*\OO(n)$ preserves every fibre restriction. For sufficiently large $n$, it kills $H^1(\PP^1,R^1f_*(T_Y\otimes L)\otimes\OO(n))$. We choose $n$ divisible by the finite order of $c_1(f^*\OO(1))$, so this twist stays in $P$. The two-column Leray sequence gives a line bundle in $P$ with zero $H^2$. Proper semicontinuity makes this zero locus dense open. Intersecting it with the dense open endpoint-vanishing loci in the proof of \ref{cdp:st:C.1.3} proves \eqref{cdp:eq:C.19}, using $\chi=c_3/2$ from \ref{cdp:st:C.1.1}.

\item Conversely, suppose $H^2$ vanishes on a dense open subset of $P$, and choose $L$ in that subset and in the full-measure set of \ref{cdp:st:D.1.11}. The preceding smooth-fibre argument makes $R^2f_*(T_Y\otimes L)$ a torsion sheaf. Leray gives a surjection
\[
 H^2(Y,T_Y\otimes L)\longrightarrow
 H^0(\PP^1,R^2f_*(T_Y\otimes L)).
\]
Thus that torsion sheaf is zero. Top-degree base change and fibre duality imply $L\notin Z_c$ for every $c$. In particular none of the $Z_c$ equals $P$, proving the equivalence.

\item When $b_3=0$, \ref{cdp:st:C.1.3} gives $H^2(Y,T_Y\otimes L)\ne0$ for every $L$. The equivalence therefore forces some $Z_c$ to be all of $P$. There are only finitely many $c$, and $P$ is an irreducible vector-space torsor; a finite union of proper closed analytic subsets cannot exhaust it. The integral hypothesis makes $P=\Pic(Y)$ by the exponential sequence. This proves the assertion about a fixed fibre. Its support is irreducible: if it were reducible, \ref{cdp:st:A.3.5} would give a dense open set of ambient-vanishing twists in the component containing $K_Y\otimes L^{-1}$. Pullback along the biholomorphism $L\mapsto K_Y\otimes L^{-1}$ would make $Z_c$ proper, a contradiction. Finally \cref{cdp:st:A.1.1,cdp:st:A.2.1,cdp:st:A.3.1} give the stated fibre tests including every Cartier layer and connecting map.

\end{enumerate}
\end{proof}


\begin{proposition}[Generic tangent cohomology vanishing]\hypertarget{tag-01FK}{}\label{cdp:st:C4a}
\StatementRecord{cdp:st:C4a}{proved}{01FK}{}

Let $f:Y\to\PP^1$ be a connected-fibre holomorphic map from a smooth compact connected complex threefold with $b_1(Y)=b_2(Y)=0$, and smooth general fibre a complex two-torus. Suppose every full fibre has one of the following forms: a reduced irreducible normal-crossing surface; or a smooth torus family modulo a finite free cyclic action over a cyclic base cover, so that its reduction is smooth and its multiplicity is the covering order. Let
\[
 0\ne\gamma\in H^0(\PP^1,R^1f_*\Z)
\]
be a generator, and suppose that, for every singular reduced full fibre $S$, its pullback to $H^1(\widetilde S,\Q)$ under normalization is nonzero. Then a dense open subset of every component of $\Pic(Y)$ satisfies
\begin{equation}
 H^q(Y,T_Y\otimes L)=0\quad(0\le q\le3),\qquad
 \chi(Y,T_Y\otimes L)=0,\qquad c_3(Y)=0.                \hypertarget{tag-008T}{}\label{cdp:eq:C.15}
\end{equation}
The generator exists under the stated rational Betti-number assumptions; its normalization condition concerns the indicated integral cohomology class.
\end{proposition}


\begin{proof}
\begin{enumerate}[label=\arabic*.,leftmargin=*]
\item The integral Leray sequence gives an injection
\[
H^0(\PP^1,R^1f_*\Z)\longrightarrow H^2(\PP^1,\Z)=\Z
\]
whose cokernel injects into the finite group $H^2(Y,\Z)$. Thus its source is infinite cyclic. By \ref{cdp:st:D.1.8}, the restriction of $\gamma$ to every full fibre is a nonzero integral class: a class in the preimage of a small disc that vanishes on a nearby fibre comes from the meridian of the punctured disc; a normal disc to any irreducible component makes a nonzero multiple of that meridian nullhomotopic, forcing the class to vanish. The real-to-holomorphic cohomology map on a compact connected complex space is injective on real classes, by \ref{cdp:st:D.1.7}, using the pluriharmonic maximum principle. Coherent Leray therefore supplies $\beta\in H^1(Y,\OO_Y)$ whose restriction to every full fibre is the image of $\gamma$.

\item Fix a real irrational number $r$ and put $M_r=\exp(r\beta)$, with the exponential normalized by $2\pi i$. On each fibre this line is represented by the unitary character $\exp(2\pi ir\gamma)$. It is nontrivial, and on a singular fibre its normalization pullback is nontrivial by the hypothesis on $\gamma$. To see this directly, triviality of an $n$th tensor power would place $nr\gamma$ in the integral lattice after applying $H^1(\R)\to H^1(\OO)$; real injectivity makes this an equality in $H^1(\R)$, which is impossible for an irrational multiple of a nonzero integral vector.

\item The same construction works in every Picard component. Each component contains a finite-order line bundle $P$: for a representative $P_0$, a positive power $P_0^m$ lies in the complex vector group $\Pic^0(Y)$, where division by $m$ is possible. Dividing $P_0$ by the resulting line makes its $m$th power trivial. Replacing $M_r$ by $M_r\otimes P$ retains nontrivial unitary restriction on each fibre and on each specified normalization. Indeed an $m$th power eliminates $P$ and reduces a purported trivialization to the preceding irrationality contradiction.

\item Set $L=K_Y\otimes(M_r\otimes P)^{-1}$. Tensoring with $K_Y$ and inversion give a bijection of Picard components, so this choice produces a test line in each desired component. On every full fibre $F$ the Serre-dual bundle for degree two is
\[
 (T_Y\otimes L)|_F^*\otimes\omega_F
       =\Omega_Y^1|_F\otimes(M_r\otimes P)|_F,
\]
because $F$ is a Cartier fibre and $\omega_F=K_Y|_F$. At a smooth torus fibre both the conormal line and the intrinsic cotangent bundle are trivial; nontrivial unitary degree-zero lines on the torus have zero cohomology. The ambient space of sections is therefore zero. At a reduced irreducible normal-crossing fibre, Proposition \ref{cdp:st:A.4.5} proves the torsion-free cotangent vanishing for every nontrivial unitary line. The normal-crossing kernel in \ref{cdp:st:A.1.2} is the normalization pushforward, and its sections vanish because the pulled-back unitary line is nontrivial. A holomorphic section of a unitary line on a compact connected complex manifold has constant norm by the maximum principle and is either zero or a trivialization. Thus the ambient group vanishes there as well.

\item On a cyclic-cover neighbourhood of a smooth multiple fibre, the covering map to $Y$ is finite \'etale. The pulled-back line restricts nontrivially to the covering central torus: finite \'etale pullback is injective on rational degree-one cohomology by its trace, and $r$ is irrational. The scheme-theoretic inverse image of the full fibre is the corresponding thickening of this torus. Each layer in its ideal filtration is the same ambient cotangent restriction tensored with this nontrivial line, since the covering base coordinate trivializes the conormal. Every layer has zero cohomology. The exact filtration sequences therefore give zero sections on the full thickening. Taking finite-group invariants, an exact functor, proves the required ambient vanishing on the original full fibre.

\item Write $A=P$ for the finite-order line of step~3, and let
$\mathscr P$ be the Picard component containing the line $L$
constructed in step~4. For each critical full fibre $F_c$, steps~4--5
give $L\notin Z_c$, where $Z_c$ is the locus
\eqref{cdp:eq:C.18} on $\mathscr P$: its twisting line is
$(K_Y\otimes L^{-1})|_{F_c}=(M_r\otimes A)|_{F_c}$.
Thus every $Z_c$ is proper.
Proposition~\ref{cdp:st:C4c}, applied to the same $f$ and this
Picard component, gives the dense open vanishing in degrees
$q\ne1$ and the index inequality \eqref{cdp:eq:C.19}, including
the endpoint vanishings of \ref{cdp:st:C.1.3}.
In particular $\chi(T_Y\otimes L)=-h^1\le0$ there.

\item Proposition \ref{cdp:st:A.4.6} gives Euler number zero or two for each reduced irreducible normal-crossing fibre; the positive case has rational normalization. Our normalization hypothesis excludes that case, since rational smooth surfaces have $H^1(-,\Q)=0$. Smooth torus fibres and the smooth finite \'etale quotients in the second fibre class have Euler number zero. Therefore \SeriesRef{cdp:st:C2a} gives $c_3(Y)=0$. The index formula in \ref{cdp:st:C.1.1} gives $\chi=0$, forcing $h^1=0$ on the same dense open set. This proves all of \eqref{cdp:eq:C.15}.

\end{enumerate}
\end{proof}






\begin{proposition}[Twisted tangent cohomology on rational homology six-spheres]\hypertarget{tag-01FL}{}\label{cdp:st:C4b}
\StatementRecord{cdp:st:C4b}{proved}{01FL}{}

Let $f:Y\to\PP^1$ be a connected-fibre holomorphic map from a smooth compact complex threefold with $b_1=b_2=b_3=0$. For every $L\in\Pic(Y)$ there is a scheme-theoretic fibre $F$ with
\begin{equation}
 H^2(F,(T_Y\otimes L)|_F)\ne0,\qquad
 H^0(F,\Omega_Y^1|_F\otimes(K_Y\otimes L^{-1})|_F)\ne0. \hypertarget{tag-008U}{}\label{cdp:eq:C.16}
\end{equation}
For torus general fibre the reduced fibre Euler numbers sum to two. For the two fibre forms listed in \ref{cdp:st:C4a}, exactly one fibre has rational normalization; it is an irreducible reduced normal-crossing fibre of Euler number two. For that normalization,
\begin{equation}
H^2(\widetilde S,K_{\widetilde S}\otimes
             \eta^*(K_Y^{-1}\otimes L))=\C
\quad\hbox{for every }L\in\Pic(Y).                    \hypertarget{tag-008V}{}\label{cdp:eq:C.17}
\end{equation}
\end{proposition}


\begin{proof}
\begin{enumerate}[label=\arabic*.,leftmargin=*]
\item The divisor of any fibre is nonzero and effective. Proposition~\ref{cdp:st:C.1.3} therefore gives $H^2(Y,T_Y\otimes M)\ne0$ for every line bundle $M$. Suppose $R^2f_*(T_Y\otimes L)=0$. Twisting by a sufficiently large positive pullback from the base kills $H^1(R^1f_*)$, so Leray would give a line bundle $M=L\otimes f^*\OO(n)$ with $H^2(Y,T_Y\otimes M)=0$, a contradiction. Thus $R^2f_*(T_Y\otimes L)$ is nonzero. A nonzero stalk has nonzero residue by Nakayama's lemma, and top-degree base change gives the first nonvanishing in \eqref{cdp:eq:C.16}. Cartier-fibre adjunction gives $\omega_F=(K_Y\otimes\OO_Y(F))|_F\simeq K_Y|_F$, since the normal line comes from the base point. Serre duality identifies the dual of the first group with the second group displayed there.

\item For torus general fibre, \ref{cdp:st:C2a} applied to
the same $f$ gives $\sum e(F_{\rm red})=e(Y)=2$.
For the two fibre forms in \ref{cdp:st:C4a}, every reduced support
has normal crossings, including the smooth finite torus quotients.
The map $g$ of \ref{cdp:st:C.2.NC.5} is therefore this $f$, with
$C=\PP^1$ and the same full fibres. Since $b_1=b_2=b_3=0$,
that proposition gives exactly one support with rational
normalization and Euler number two. The allowed singular full
fibre is reduced; its hexagonal normalization description is the
one in \ref{cdp:st:A.4.6}. Its normalization map
$\eta:T\to S$ is the map $\eta:\widetilde S\to S$ here.
Under $T=\widetilde S$ and the identity on the ambient line $L$,
the first group of \eqref{cdp:eq:C.NC.11} is precisely
\eqref{cdp:eq:C.17}. The rational Picard pullback and Serre pairing
proved there give its value $\C$ for every $L$.

\end{enumerate}
\end{proof}















\begin{proposition}[A finite transgression criterion for tangent cohomology]\hypertarget{tag-01FN}{}\label{cdp:st:C4d}
\StatementRecord{cdp:st:C4d}{proved}{01FN}{\SeriesRef{cdp:st:D.1.9}, \SeriesRef{cdp:st:A.4.9}, \SeriesRef{cdp:st:A.4.7}, \SeriesRef{cdp:st:A.3.4}, \SeriesRef{cdp:st:A.5.2}}

Let $f:Y\to\PP^1$ satisfy the hypotheses of
\cite[Theorem~2.2, p.~680]{CDP2020}: $Y$ is a connected smooth compact
complex threefold with $H^1(Y,\Z)=H^2(Y,\Z)=0$ and $a(Y)=1$,
and $f$ is its holomorphic algebraic reduction with connected fibres. Let $\beta\in H^1(Y,\OO_Y)$ lift the integral transgression generator of \ref{cdp:st:D.1.9}, and put $M_z=\exp(2\pi i z\beta)$, with the same exponential convention as in that statement. For each critical value $c$, define
\begin{equation}
 B_c=\{z\in\C:
 H^0(F_c,\Omega_Y^1|_{F_c}\otimes M_z|_{F_c})\ne0\},
 \qquad F_c=f^{-1}(c).
 \hypertarget{tag-008Z}{}\label{cdp:eq:C.20}
\end{equation}
When each $B_c$ is a proper subset of $\C$, a dense open subset of $\Pic(Y)$ satisfies the vanishing and the two inequalities in \eqref{cdp:eq:C.19}.
The premise is a finite collection of tests on the full fibres,
given by the conormal criterion \ref{cdp:st:A.1.1} and the
Cartier-thickening filtration \ref{cdp:st:A.3.1}; it is satisfied in particular when every reduced support is smooth.

When $c_3(Y)>0$, there is a fixed critical value $c$ with $B_c=\C$. Its reduced support is nonnormal. The other fibres may have arbitrary multiplicities and singularities.
\end{proposition}


\begin{proof}

\begin{enumerate}[label=\arabic*.,leftmargin=*]
\item The transgression lift exists because the coherent Leray sequence over $\PP^1$ has no degree-one or degree-two base term for $\OO_{\PP^1}$. The integral degree-one Leray sequence identifies the integral invariant group with $\Z$ under the two stipulated integral vanishings. \ref{cdp:st:D.1.9} proves that its restriction to every full fibre, and to every nearby smooth fibre, is nonzero. On a neighbourhood $U$ of any fibre, $M_r$ for real $r$ has the unitary local system
\[
 a\longmapsto\exp(2\pi i r\langle\gamma_U,a\rangle)
\]
as its flat realization. The exponential coefficient map sends
$\gamma_U$ to $\beta|_U$, giving this equality on $U$.

\item The holomorphic family of exponentiated \v Cech cocycles representing $z\beta$ restricts to a flat family of coherent bundles on $F_c\times\C$. Proper semicontinuity makes $B_c$ a closed analytic subset of the complex line. A proper such subset is discrete and countable. There are finitely many critical values. We can therefore choose a real irrational $r$ outside the union of all $B_c$ when all are proper.

\item On a smooth fibre $F$, the character just displayed is the irrational exponential of a nonzero integral class. The explicit calculations for tori, bielliptic surfaces, Kodaira surfaces, Hopf surfaces and Inoue surfaces in \ref{cdp:st:A.4.9} give
\[
 H^0(F,\mathbb A_r)=H^1(F,\mathbb A_r)=0.
\]
The unitary maximum principle gives $H^0(F,M_r)=0$, and the injection \ref{cdp:st:A.4.7} gives $H^0(F,\Omega_F^1\otimes M_r)=0$. The conormal sequence of a smooth full fibre, whose conormal line is trivial, then gives
$H^0(F,\Omega_Y^1|_F\otimes M_r)=0$.
At every critical full fibre this group vanishes by our choice of $r$. Set $L=K_Y\otimes M_r^{-1}$. Cartier-fibre duality and top-degree base change give $R^2f_*(T_Y\otimes L)=0$. A sufficiently positive base twist kills $H^1(R^1f_*(T_Y\otimes L))$, so
\[
 H^2(Y,T_Y\otimes L\otimes f^*\OO(n))=0
\]
for every sufficiently large integer $n$. The exponential sequence identifies $\Pic(Y)$ with a connected complex vector space. Semicontinuity therefore makes degree-two vanishing dense open in this whole group. Intersecting with the endpoint-vanishing loci in the proof of
\cref{cdp:st:C.1.3} and applying \cref{cdp:st:C.1.1} gives
\eqref{cdp:eq:C.19}.

\item For a full fibre $mS$ with smooth reduction, the case $m=1$
follows from the conormal sequence above, and the case $m\ge2$
follows from \ref{cdp:st:A.3.4}, including all nilpotent layers.
Thus the ambient group vanishes for every real irrational $r$,
and $B_c$ is proper.

\item We have $c_3(Y)=2-b_3(Y)$, and \ref{cdp:st:C.1.1}
gives $c_3(Y)/2\in\Z$. Hence $b_3(Y)$ is even, and
$c_3(Y)>0$ implies $b_3(Y)=0$. Proposition~\ref{cdp:st:C.1.3} gives $H^2(Y,T_Y\otimes L)\ne0$ for every line bundle $L$. The preceding proof shows that the finitely many $B_c$ cannot all be proper; hence one fixed $B_c$ is the entire complex line. Its reduced support cannot be smooth by the preceding
full-fibre argument. A connected normal reduced support would be irreducible. Since it is a Cartier divisor in $Y$, its dualizing line has rational first Chern class zero and its structure-sheaf Euler characteristic is zero by divisor Riemann--Roch, exactly as calculated in \ref{cdp:st:A.5.2}. That statement proves it smooth. This contradicts the full-fibre argument above, so the support is nonnormal.

\end{enumerate}
\end{proof}








\subsection{Normalization restrictions and the global criterion}

\begin{lemma}[Persistent sections and the normalization pullback]\hypertarget{tag-P263}{}\label{lem:II:persistent-normalization-pullback}
\StatementRecord{lem:II:persistent-normalization-pullback}{proved}{P263}{\SeriesRef{cdp:st:C.3.11}, \SeriesRef{cdp:st:A.3.5}, \SeriesRef{cdp:st:A.3.4}, \SeriesRef{cdp:st:A.5.2}, \SeriesRef{cdp:st:A.3.1}, \SeriesRef{cdp:st:A.4.9}, \SeriesRef{cdp:st:A.5.4}, \SeriesRef{cdp:st:A.5.3}, \SeriesRef{cdp:st:A.5.12}, \SeriesRef{cdp:st:A.5.13}, \SeriesRef{cdp:st:C2d}}
Let $f:Y\to\PP^1$ satisfy the integral CDP hypotheses stated in
\cref{cdp:st:C4d}, and let $F$ be a full fibre such that
\[
 H^0\bigl(F,(\Omega_Y^1\otimes M)|_F\bigr)\ne0
 \qquad(M\in\Pic(Y)).
\]
Then $F=mS$ for a positive integer $m$ and an irreducible nonnormal
reduced support $S$. For its normalization $\eta:T\to S$, the
pullback $\eta^*(M|_S)$ is trivial for every $M\in\Pic(Y)$.
\end{lemma}

\begin{proof}
\begin{enumerate}[label=\arabic*.,leftmargin=*]
\item By \cref{cdp:st:C.3.11}, the general fibre is a two-torus
or a Hopf surface. In the Hopf case it has $h^1(\OO)=1$ and
$h^2(\OO)=0$, as used in the conductor bound below.
\ref{cdp:st:A.3.5} makes the support of $F$ irreducible, so $F=mS$ for a positive integer $m$. \ref{cdp:st:A.3.4} excludes a smooth support, and \ref{cdp:st:A.5.2} then excludes a normal support. Its filtration in \ref{cdp:st:A.3.1} has finitely many graded factors
\[
 \Omega_Y^1|_S\otimes M|_S\otimes\OO_Y(-jS)|_S,
 \qquad 0\le j<m.
\]
Nonvanishing on the filtered sheaf requires nonvanishing on at least one factor. Each factor's nonvanishing locus is closed analytic in the connected vector group $\Pic(Y)$. A finite union of proper such loci cannot cover that group. Hence one factor has nonzero sections for every $M$. Translating by $\OO_Y(-jS)$ shows
$H^0(S,\Omega_Y^1|_S\otimes M|_S)\ne0$ for every $M$.
\ref{cdp:st:A.4.9} kills the torsion-free quotient on a dense open set of twists. The exact sequence \ref{cdp:st:A.1.1} and proper semicontinuity then show that the inverse-Jacobian kernel itself has sections for every twist.

\item We apply the precise persistent-kernel reduction \ref{cdp:st:A.5.4}. It gives a fixed component of the normalization on whose resolution every twist of one fixed rank-one sheaf has a section. The Gauduchon degree argument there proves that the pullback of every ambient line has degree zero for every Gauduchon metric. Its conductor comparison excludes nonnegative Kodaira dimension, and its class VII argument makes every ambient line trivial on that normalization component. Rational components have the same triviality because their Picard group has no connected part. In either case the projection formula descends this triviality from the resolution to the normalization. It remains to treat an irrational ruled resolution $p:Z\to B$.

\item Let $h:Z\to S\subset Y$. \ref{cdp:st:A.5.3} supplies an effective conductor divisor $D_Z$ with
\begin{equation}
 K_Z+D_Z=h^*(K_Y+S),\qquad D_Z\cdot F=2              \hypertarget{tag-0091}{}\label{cdp:eq:C.22}
\end{equation}
for a general ruling fibre $F$. Its horizontal part is therefore either reduced of total ruling degree two, or $2C$ for a section $C$. For each ruling blowdown $\sigma:Z\to Z'$, apply \cref{lem:II:canonical-boundary-blowup} with $Z_1=Z$, $Z_2=Z'$, $D_1=D_Z$, $D_2=D'=\sigma_*D_Z$, and its multiplicity $m=d=\operatorname{mult}_{\sigma(E)}D'$. Its coefficient $a$ is the coefficient of the same exceptional curve $E$ in $D_Z$. Formula~\eqref{cdp:eq:C.22b} gives $a=d-1$ and identifies $K_Z+D_Z$ with the pullback of $K_{Z'}+D'$, which again has zero rational Chern class.

\item We first treat the reduced horizontal part $H$. If a component is exceptional for the normalization resolution, every ambient line restricts trivially to it. Its normalization maps finitely onto $B$, and the norm composed with pullback on $\Pic^0(B)$ is multiplication by the covering degree. Hence this restriction has finite kernel; the connected ambient Picard image is trivial, giving the asserted triviality. We may therefore assume that the horizontal components map to conductor curves. Adjunction on each of them gives
\[
 2p_a(H_i)-2+(D_Z-H_i)\cdot H_i=0.
\]
The second term is nonnegative. Since $H_i$ dominates a positive-genus base, its arithmetic genus is at least one. Thus the base has genus one, each $H_i$ is smooth elliptic, and the horizontal curves are disjoint from each other and from the vertical part. Contract ruling $(-1)$-curves to a minimal bundle. A contracted curve outside $D_Z$ has $D_Z\cdot E=1$, so cannot identify two horizontal points or create a horizontal node. A contracted curve in the vertical part is disjoint from $H$. Downstairs the horizontal curves remain smooth and disjoint, and any nonzero vertical divisor on a minimal ruled bundle meets them. Consequently the downstairs vertical part is zero. Reversing the contractions, a blowup away from the horizontal conductor would give coefficient $-1$ on its exceptional curve in \eqref{cdp:eq:C.22}, which persists under further strict transform and contradicts effectiveness. A blowup at a smooth point of a coefficient-one horizontal curve gives exceptional coefficient zero. Hence $D_Z=H$ has no vertical component. Every codimension-one singularity of $S$ is then an ordinary
node. Apply \cref{cdp:st:A.2.3} to this reduced hypersurface:
its normalization criterion gives $\mathcal J^{-1}=
\eta_*\mathcal O_T$ on all of $S$. Under the map
$[F]\otimes q\mapsto q\,dF|_S$ of \cref{cdp:st:A.2.1},
the inverse-Jacobian kernel is therefore
$\eta_*\eta^*\mathcal O_Y(-S)$. A section of a topologically trivial line on the K\"ahler resolution has an effective zero divisor of zero K\"ahler volume, so is nowhere zero. Persistence for all twists therefore makes all ambient pullbacks trivial and proves the asserted triviality.

\item We now have $D_Z=2C+V$, where $C$ is a section and $V$ is vertical. If $C$ is exceptional for $Z\to\widetilde S$, every ambient line is trivial on $C$. Restriction $\Pic^0(Z)\to\Pic^0(C)$ is an isomorphism, because $C$ is a ruling section. This again makes all ambient pullbacks trivial. Otherwise let $D\subset S$ be its image and let $D^\nu$ normalize that curve. When $D^\nu$ is rational, every ambient line pulls back trivially to $D^\nu$, hence to $C$, giving the same conclusion.

\item Suppose instead that $g(D^\nu)\ge1$. A vertical rational curve cannot map finitely onto $D^\nu$, so all the normalization branches over the general point of $D$ belong to $C$. Their conductor coefficient is two. We determine the transverse reduced plane curves from these conductor exponents. Write its reduced equation as $f=\prod_i f_i$, with branch ring $R_i$ and normalization $B_i$. Apply \cref{lem:II:plane-branch-conductor} to this same equation, with its unit equal to one. It gives the conductor exponent $2\delta(f_i)+\sum_{j\ne i}I(f_i,f_j)$ on branch $i$. A singular branch has $\delta(f_i)\ge1$, so its intrinsic exponent is at least two.
Exponent two forces $R_i=\C+t^2\C\{t\}=\C\{t^2,t^3\}$,
since an element of order one would again make the ring smooth.
Any additional branch has intersection multiplicity at least
two with a singular branch and is therefore excluded. Thus a
singular branch gives exactly the cusp. When all branches are
smooth, their pairwise positive intersection multiplicities sum
to two on each branch. This gives two branches of contact two
(the tacnode), or three pairwise transverse branches (the
ordinary triple point), and exhausts the possibilities.
The only remaining alternative is a nonconstant pullback
$\Pic(Y)\to\Pic(Z)$; a constant pullback is the zero
homomorphism and already gives the asserted triviality. In that alternative,
the coherent conductor argument \ref{cdp:st:A.5.12} gives $p_a(D)\le1$
for torus general fibre, and \ref{cdp:st:A.5.13} gives the same bound
in the second case.
Since $g(D^\nu)\ge1$, it follows that $D$ itself is smooth
elliptic. We can therefore apply \ref{cdp:st:C2d} directly. It makes the
transverse type constant, gives the \'etale conductor cover
$C\to D$, and proves $C^2=0$. In particular $C$ and the
ruling base are elliptic.

\item Adjunction in \eqref{cdp:eq:C.22} gives $V\cdot C=0$. No ruling $(-1)$-curve meets $C$: for one outside $V$, its intersection with $K_Z+2C+V$ would be at least $-1+2=1$; for one in $V$, the equality $V\cdot C=0$ forbids such an intersection. Thus every contraction to the minimal ruled bundle avoids $C$ and preserves its self-intersection zero. On that bundle $K+2C$ is numerically trivial, so the effective vertical part is numerically trivial and hence zero. A first reverse blowup, necessarily away from $C$, would again produce negative conductor coefficient. There are therefore no reverse blowups: $Z$ is minimal ruled and $D_Z=2C$. There is also no exceptional curve for the normalization resolution: a negative curve $E\ne C$ would satisfy $2p_a(E)-2=E^2-2C\cdot E<0$, hence would be rational and contained in a ruling fibre; a minimal ruled bundle has only self-intersection-zero fibres. Thus the normalization itself is smooth here.

\item By the explicit inverse-Jacobian computations in \ref{cdp:st:C2d}, its pullback allows at most a simple pole along $C$. There are no other divisorial poles. Hartogs extension for the line bundle on the smooth surface gives an injection of the kernel into $h^*\OO_Y(-S)\otimes\OO_Z(C)$. The exact sequence of the section $C$ and $R^1p_*\OO_Z=0$ give
\begin{equation}
 0\longrightarrow\OO_B\longrightarrow p_*\OO_Z(C)
   \longrightarrow N_{C/Z}\longrightarrow0,
 \qquad \deg N_{C/Z}=0.                              \hypertarget{tag-0092}{}\label{cdp:eq:C.22a}
\end{equation}
For a degree-zero line $P$ on the elliptic curve, both end terms after tensoring by $P$ have zero sections unless $P=\OO_B$ or $P=N_{C/Z}^{-1}$. Thus $H^0(Z,\OO_Z(C)\otimes p^*P)$ is supported at at most two points of $\Pic^0(B)$. The ambient pullback image is a connected complex subgroup of this Picard group. Persistence for every ambient twist forces that subgroup to be a point, and it contains the identity. Therefore every ambient line pulls back trivially to $Z$ in this last case as well.

\item Write $\pi:Z\to T$ for the normalization resolution.
The equality $\pi_*\OO_Z=\OO_T$ and the projection formula
identify the sections of a line on $T$ with those of its pullback.
Applying this to the line and its inverse descends the mutually
inverse trivializations obtained on $Z$. Since $S$ is irreducible,
$T$ is the normalization component selected above. Thus
$\eta^*(M|_S)\simeq\OO_T$ for every ambient line $M$.
\end{enumerate}
\end{proof}

\begin{proposition}[The unrestricted normalization target under the integral CDP hypotheses]\hypertarget{tag-01FO}{}\label{cdp:st:C4e}
\StatementRecord{cdp:st:C4e}{proved}{01FO}{\SeriesRef{cdp:st:C.1.1}, \SeriesRef{cdp:st:C.3.11}, \SeriesRef{cdp:st:C4c}, \SeriesRef{cdp:st:C4d}, \SeriesRef{cdp:st:A.5.13}, \SeriesRef{lem:II:persistent-normalization-pullback}}

Let $f:Y\to\PP^1$ satisfy the integral CDP hypotheses and $c_3(Y)>0$. Then a fixed irreducible nonnormal reduced support $S$ has normalization $\eta:\widetilde S\to S$ satisfying
\begin{equation}
 H^2\bigl(\widetilde S,\omega_{\widetilde S}\otimes
           \eta^*(K_Y^{-1}\otimes L)\bigr)\ne0
 \qquad\hbox{for every }L\in\Pic(Y).                  \hypertarget{tag-0090}{}\label{cdp:eq:C.21}
\end{equation}
The full fibre over this support is $mS$ for a positive integer $m$.
\end{proposition}


\begin{proof}
The index formula \cref{cdp:st:C.1.1} gives
$c_3(Y)=2-b_3(Y)$ and $c_3(Y)/2\in\Z$. Consequently
$c_3(Y)>0$ gives $b_3(Y)=0$.
By \cref{cdp:st:C.3.11}, the general fibre is a two-torus or a
Hopf surface. In the torus case \cref{cdp:st:C4c} supplies a fixed
full fibre with a nonzero ambient cotangent section for every
twist. In the Hopf case, $h^1(\OO)=1$ and $h^2(\OO)=0$;
\cref{cdp:st:C4d} gives persistence along the transgression line,
and \cref{cdp:st:A.5.13} identifies that line with $\Pic(Y)$.
In both cases, inversion and tensoring with $K_Y$ identify the
twisting parameter in the cited statements with every
$M\in\Pic(Y)$.

Lemma~\ref{lem:II:persistent-normalization-pullback} therefore
gives $F=mS$ with $S$ irreducible and nonnormal, and makes
$\eta^*(K_Y^{-1}\otimes L)$ trivial for every $L$.
The compact normal surface $T=\widetilde S$ is connected and
Cohen--Macaulay. Serre duality identifies the group in
\eqref{cdp:eq:C.21} with $H^0(T,\OO_T)^\vee=\C$.
\end{proof}








\begin{proposition}[An unrestricted geometric replacement for CDP20 Theorem 2.2]\hypertarget{tag-01FP}{}\label{cdp:st:C4f}
\StatementRecord{cdp:st:C4f}{proved}{01FP}{\SeriesRef{lem:II:persistent-normalization-pullback}, \SeriesRef{cdp:st:C.3.11}, \SeriesRef{cdp:st:C4c}, \SeriesRef{cdp:st:A.3.4}, \SeriesRef{cdp:st:A.3.5}, \SeriesRef{cdp:st:A.5.2}, \SeriesRef{cdp:st:A.2.2}, \SeriesRef{cdp:eq:C.VC.4}, \SeriesRef{cdp:st:A.5.13}, \SeriesRef{cdp:st:C4d}, \SeriesRef{cdp:st:D.1.11}, \SeriesRef{cdp:st:A.4.9}}

Let $f:Y\to\PP^1$ satisfy the integral CDP hypotheses.
For each critical full fibre with
irreducible nonnormal reduced support $S$, let
$\eta_S:T_S\to S$ be its normalization. Then a general
line bundle $L\in\Pic(Y)$ satisfies
\begin{equation}
 H^q(Y,T_Y\otimes L)=0\quad(q\ne1)
 \hypertarget{tag-015G}{}\label{cdp:eq:C.VC.1}
\end{equation}
exactly when, for every such $S$, the holomorphic pullback
homomorphism
\begin{equation}
 \Pic(Y)\longrightarrow\Pic(T_S)
 \hypertarget{tag-015H}{}\label{cdp:eq:C.VC.2}
\end{equation}
is nonzero. Equivalently, its differential
$H^1(Y,\OO_Y)\to H^1(T_S,\OO_{T_S})$ is nonzero.
In this case $\chi(Y,T_Y\otimes L)\le0$ and $c_3(Y)\le0$.
If one map \eqref{cdp:eq:C.VC.2} is zero, then
$H^2(Y,T_Y\otimes L)\ne0$ for every $L\in\Pic(Y)$.
All fibre multiplicities and all reduced-support
singularities are permitted in this criterion.
\end{proposition}


\begin{proof}

\begin{enumerate}[label=\arabic*.,leftmargin=*]
\item The classification \cref{cdp:st:C.3.11} gives the two
cohomological cases: torus general fibre, or a Hopf general fibre
with $h^1(\OO)=1$ and $h^2(\OO)=0$.
The Picard group of $Y$ is the connected exponential image
of $H^1(Y,\OO_Y)$. A holomorphic group homomorphism out of
this group is zero precisely when its differential is zero.
For each fixed critical full fibre, the ambient-section
nonvanishing locus is closed analytic in this group, as
in \ref{cdp:st:C4c}. \ref{cdp:st:A.3.4} makes it proper when the support is smooth,
and \ref{cdp:st:A.3.5} makes it proper when the support is reducible.
A connected normal support is irreducible and is smooth
by \ref{cdp:st:A.5.2}, so these cases leave just the supports in the
statement.

\item Suppose the nonvanishing locus for $F=mS$ is the whole
Picard group. Lemma~\ref{lem:II:persistent-normalization-pullback}
makes $\eta_S^*(M|_S)$ trivial for every $M\in\Pic(Y)$.
This is exactly the vanishing of the homomorphism
\eqref{cdp:eq:C.VC.2}.

\item Conversely, suppose \eqref{cdp:eq:C.VC.2} is zero. The last nonzero
ideal layer of the full fibre is a subsheaf
\begin{equation}
 (\Omega_Y^1\otimes M\otimes\OO_Y(-(m-1)S))|_S
 \hookrightarrow
          (\Omega_Y^1\otimes M)|_{mS}.
 \hypertarget{tag-015I}{}\label{cdp:eq:C.VC.3}
\end{equation}
The canonical normalization submodule in \ref{cdp:st:A.2.2} gives
an injection into the left side from
\begin{equation}
 \eta_{S*}\bigl(\eta_S^*((M\otimes\OO_Y(-mS))|_S)\bigr).
 \hypertarget{tag-015J}{}\label{cdp:eq:C.VC.4}
\end{equation}
Here the first conormal factor $\OO_Y(-S)|_S$ combines
with the displayed $(m-1)$st ideal layer; this comparison
retains every nilpotent layer needed for the injection.
Since $mS$ is a full fibre, $\OO_Y(-mS)|_S$ is a trivial
line, with its one-dimensional scalar factor determined
by a local base parameter. The pullback of $M$ to $T_S$
is trivial by \eqref{cdp:eq:C.VC.2}. Hence \eqref{cdp:eq:C.VC.4} has a nonzero global
section for every $M$, and its injections prove that the
full-fibre nonvanishing locus is the entire Picard group.
This also proves the converse directly for arbitrary $m$.

\item We have proved an exact geometric characterization of
every improper critical-fibre locus. For torus general
fibre, applying \ref{cdp:st:C4c} gives the equivalence with generic
degree-two vanishing. In the second case \ref{cdp:st:A.5.13} identifies
$\Pic(Y)$ with the transgression line, so properness of
all these loci is exactly the premise of \ref{cdp:st:C4d} and gives
generic degree-two vanishing there as well.
The proof of \cref{cdp:st:C.1.3}, applied to this same $Y$ and any nonzero full fibre divisor, gives degree-zero and degree-three vanishing on a dense open subset of $\Pic(Y)$. Intersecting it with the degree-two vanishing locus gives \eqref{cdp:eq:C.VC.1}. The index formula in \cref{cdp:st:C.1.1} then gives $\chi(Y,T_Y\otimes L)=-h^1(Y,T_Y\otimes L)\le0$ and $c_3(Y)\le0$.
For the last assertion in the torus case, generic non-torsion on every
smooth torus fibre makes $R^2f_*(T_Y\otimes L)$ a
torsion sheaf for a dense set of twists, by \ref{cdp:st:D.1.11}.
The fixed full-fibre obstruction makes it nonzero.
Its nonzero global sections are a quotient of
$H^2(Y,T_Y\otimes L)$ by the Leray sequence.
Thus degree-two cohomology is nonzero on a dense set.
Its proper semicontinuity locus is closed analytic in
the connected Picard family, so it is nonzero for every
twist. In the second case use the twists
$L_r=K_Y\otimes M_r^{-1}$ for real irrational $r$.
\ref{cdp:st:A.4.9} and the smooth-fibre conormal sequence make their
$R^2f_*(T_Y\otimes L_r)$ torsion. The fixed full-fibre
obstruction again makes that torsion sheaf nonzero,
so Leray gives $H^2(Y,T_Y\otimes L_r)\ne0$.
These twists have accumulation points in the complex
line $\Pic(Y)=\C\beta$. A closed complex analytic subset
of that line containing them is the whole line.
This proves all-twist nonvanishing in the second case
and completes both implications of the criterion.

\end{enumerate}
\end{proof}






\part{Irreducible supports and Euler equality}\label{part:II:irreducible-euler}

\section{Irreducible supports: topology and canonical orders}
\hypertarget{tag-034F}{}\label{sec:II:night-sign}
\begin{definition}[The standing torus-fibration setting]\hypertarget{tag-P204}{}\label{def:II:standing-torus-fibration}
\StatementRecord{def:II:standing-torus-fibration}{proved}{P204}{}
The global torus-fibration setting consists of a smooth compact complex
threefold $Y$ with $b_1(Y)=b_2(Y)=0$ and a proper surjective holomorphic
map $f:Y\to B$ with connected fibres and smooth general fibre a complex
two-torus. A full fibre is the effective Cartier divisor $f^*(b)$;
its reduced support is denoted by $S$. Euler characteristic means
topological Euler characteristic, while $\chi(\mathcal O)$ denotes the
coherent Euler characteristic. For an irreducible support $S$, the
full fibre is $mS$; we use the normalization $\nu:N\to S$ and a resolution
$\sigma:T\to N$. In the irreducible-or-normal-crossings setting,
every reduced fibre support is irreducible or has normal crossings.
Local degenerations and comparison examples have the hypotheses
specified in their statements.
\end{definition}
\subsection{The normal-crossings sign theorem}

\begin{proposition}[Euler characteristic of a normal crossing divisor]\hypertarget{tag-02XY}{}\label{night:P4.6}\label{stmt:P4.6}
\StatementRecord{night:P4.6}{proved}{02XY}{the definitions and constructions stated in the proof}

Let $X$ be a smooth compact complex threefold with $b_2(X)=0$, and let $F$ be a
reduced compact normal-crossings divisor. We allow different local
branches of one global irreducible component to meet. If $T(F)$ is the
number of ordinary triple points, then
\[
 e(F)=T(F),\qquad T(F)\text{ is even}.
\]
For a fibre $\sum_j m_jF_j$ with this reduced support, the divisor
in the formula is $F=\sum_jF_j$.
\end{proposition}
\begin{proof}

\begin{enumerate}[label=\arabic*.,leftmargin=*]
\item Write $\nu:S=\coprod S_i\to X$ for the normalization immersion, and $D_i$
for the conductor divisor on $S_i$. The surfaces $S_i$ are smooth and
the curves $D_i$ are reduced nodal divisors. The quotient
$N_i=\nu_i^*T_X/T_{S_i}$ is a line bundle. A local equation of $F$ gives
\[
 N_i\simeq \nu_i^*\OO_X(F)\otimes\OO_{S_i}(-D_i).
\]
Indeed, for the coordinate equation $x_1\cdots x_k=0$ the normal
derivative on $x_1=0$ is $x_2\cdots x_k$, whose zero divisor is $D_i$.
These are morphisms of holomorphic vector bundles, including when
$S_i$ is not embedded globally.

\item Since $b_2(X)=0$, both $c_1(T_X)$ and $c_1(\OO_X(F))$ vanish in real
cohomology. Duality gives $b_4(X)=0$, so $c_2(T_X)$ also vanishes. The
tangent exact sequence yields
\[
 K_{S_i}\equiv-D_i,\qquad c_2(S_i)=K_{S_i}^2.
\]
If $a=\sum_i\chi(\OO_{S_i})$, the surface Riemann--Roch formula
$12\chi(\OO_{S_i})=K_{S_i}^2+e(S_i)$ consequently gives
$\sum_i e(S_i)=6a$.

\item Let $\Gamma$ be the normalization of the double-curve locus, with one
sheet for each unordered pair of local branches. The normalization of
$\coprod D_i$ is an unramified degree-two cover of $\Gamma$, recording
the choice of branch. Each triple point contributes three nodes to
$\coprod D_i$. Adjunction gives
$\chi(\OO_{D_i})=-D_i\cdot(D_i+K_{S_i})/2=0$. The normalization
sequence for these nodal curves therefore gives
\[
 0=2\chi(\OO_\Gamma)-3T(F)=e(\Gamma)-3T(F).
\]
The normalization complex of $F$ is the exact alternating restriction
complex of the local coordinate hyperplanes. Permuting the branches
globally replaces its double-curve term by $\OO_\Gamma$ tensored with
the sign line of the two-sheeted branch cover. That line has order at
most two, and hence degree zero on every component; it has the same
coherent Euler characteristic as $\OO_\Gamma$. The triple term has
length one at each triple point. It follows that
\[
 \chi(\OO_F)=a-\chi(\OO_\Gamma)+T(F)=a-T(F)/2.
\]
The exact sequence for the Cartier divisor $F$ and analytic
Hirzebruch--Riemann--Roch give $\chi(\OO_F)=0$, since
$\operatorname{ch}(\OO_X(-F))=\operatorname{ch}(\OO_X)$ in rational
cohomology. Thus $T(F)=2a$. Finally ordinary Euler additivity gives
\[
 e(F)=\sum_i e(S_i)-e(\Gamma)+T(F)=6a-2T(F)=T(F).
\]
At a point with $r=1,2,3$ local branches, the alternating
contribution is $r-\binom r2+\binom r3=1$. This gives the
displayed Euler additivity formula and proves the identity and parity.

\end{enumerate}
\end{proof}

\begin{proposition}[The Chern number of a normal crossing fibration]\hypertarget{tag-02XZ}{}\label{night:P4.7}\label{stmt:P4.7}
\StatementRecord{night:P4.7}{proved}{02XZ}{\SeriesRef{night:P4.6}}

Let $f:X\to B$ be a proper surjective holomorphic map with connected
fibres from a smooth compact complex threefold satisfying $b_1(X)=b_2(X)=0$
to a smooth compact curve. Suppose all reduced fibre supports have
normal crossings. Then
\[
 B=\Pp^1,\qquad
 c_3[X]=\sum_{b\in B}T((f^{-1}(b))_{\mathrm{red}})\in\{0,2\},\qquad
 b_3(X)=2-c_3[X].
\]
\end{proposition}
\begin{proof}
\begin{enumerate}[label=\arabic*.,leftmargin=*]
\item Connected fibres give $f_*\R=\R$. The beginning of the constant-sheaf
Leray sequence injects $H^1(B;\R)$ into $H^1(X;\R)$, so $B$ has genus
zero. A smooth fibre $G$ has trivial normal line bundle. Its tangent
sequence gives $c_2(G)=c_2(X)|_G=0$, hence $e(G)=0$ for any such fibre.
The critical values are finite. Ehresmann's theorem gives a smooth
bundle over their complement. Euler additivity and Proposition~\ref{night:P4.6} give the
displayed sum. Duality gives $e(X)=2-b_3(X)$, and the nondegenerate
alternating cup product on $H^3(X;\R)$ makes $b_3$ even. The
nonnegative sum is therefore zero or two.

\end{enumerate}
\end{proof}

\subsection{Irreducible supports: canonical torsion and the normal case}

\begin{proposition}[Irreducible fibres of torus fibrations]\hypertarget{tag-02WJ}{}\label{night:P4.12}\label{stmt:P4.12}
\StatementRecord{night:P4.12}{proved}{02WJ}{\SeriesRef{night:P4.6}}

Let $Y$ be a connected smooth compact complex threefold with
$b_1(Y)=b_2(Y)=0$, and let $f:Y\to B$ be a proper surjective
holomorphic map with connected fibres to a smooth compact curve.
Assume that the smooth general fibre is a complex two-torus, and let
$f^{-1}(b)=mS$ have irreducible reduced support. Then
\[
 \omega_S^{\otimes m}\simeq\OO_S,\qquad
 \chi(\OO_S)=0,\qquad h^1(\OO_S)\in\{1,2\}.
\]
Every normal such support is smooth. For a singular, necessarily
nonnormal support, a minimal resolved model of its normalization
has Kodaira dimension $-\infty$.
\end{proposition}
\begin{proof}
\noindent\textbf{Step 1. Canonical torsion.}
Choose a disc $\Delta$ with no other critical value, with coordinate
$t$ vanishing at $b$, and write $U=f^{-1}(\Delta)$. The coherent
direct image $f_*\omega_U$ is torsion-free of rank one: nonzero
base functions act injectively on sections, and a nearby torus has
one-dimensional canonical space. It is therefore a line bundle
on the disc. A generator gives a section $v$ of $\omega_U$,
nonzero and nowhere vanishing on every noncentral fibre. Its
divisor is $aS$ for an integer $a\ge0$. If $a\ge m$, then $v/t$
would be holomorphic, contradicting primitivity of the generator
at zero. We have $0\le a<m$ and
\[
 \omega_U\simeq\OO_U(aS),\qquad\OO_U(mS)\simeq\OO_U.
\]
Adjunction gives $\omega_S\simeq\OO_U((a+1)S)|_S$, proving
$\omega_S^{\otimes m}\simeq\OO_S$. The ambient
Riemann--Roch computation in \ref{night:P4.6}, applied to $S$, gives
$\chi(\OO_S)=0$. Reduced irreducibility gives $h^0(\OO_S)=1$.
A nonzero section of a torsion line bundle has a power equal to
a nonzero constant, and hence is nowhere zero; its space of
sections consequently has dimension at most one. Serre duality
on the Gorenstein surface $S$ gives
$h^2(\OO_S)=h^0(\omega_S)\in\{0,1\}$, proving the assertion
about $h^1$.

\noindent\textbf{Step 2. Exceptional curves and discrepancies.}
For a minimal resolution $\pi:T\to S$ of a normal Gorenstein
surface write $K_T=\pi^*K_S+\sum a_jE_j$. By \cref{lem:II:exceptional-intersection},
applied to this resolution $\pi:T\to S$, the exceptional
intersection matrix $M$ is negative definite and $-M^{-1}$
has positive entries.

Adjunction gives
$K_T\cdot E_j=2p_a(E_j)-2-E_j^2\ge0$: the sole negative
possibility is a smooth rational $(-1)$-curve, excluded by
minimality of the resolution. The discrepancy equations $Ma=(K_T\cdot E_j)_j$ imply
$a_j\le0$. Hence $K_T=\pi^*K_S-E$ with $E$ an effective
exceptional integral divisor.

Put $p=\dim H^0(S,R^1\pi_*\OO_T)$. Leray, with this sheaf
supported at finitely many points, gives
\begin{equation}\hypertarget{tag-01OD}{}\label{eq:P4chi-resolution}
 \chi(\OO_T)=\chi(\OO_S)-p=-p.
\end{equation}
For an effective exceptional cycle $Z$, put $\mathcal F=I_Z$.
Fix $s\in S$ and its maximal ideal $\mathfrak m_s$. Each coherent
quotient $\mathcal F/\mathfrak m_s^k\mathcal F$ is supported on a
compact infinitesimal fibre of dimension at most one. Its cohomology
vanishes in degrees at least two. To see this, filter by the nilradical
and reduce to coherent sheaves on the reduced compact curve. Removing
point-supported torsion and pulling back to its finite normalization
leaves vector bundles on compact Riemann surfaces, up to kernels and
cokernels supported at finitely many points. Their Dolbeault resolutions
have length one, while point-supported sheaves have no higher
cohomology. The finite filtration therefore gives the asserted vanishing
on every infinitesimal fibre.
Analytic formal functions and coherence
\cite[\S6, Haupts\"atze~I and~IIa, p.~59]{Grauert1960D} identify
the completed stalk of $R^2\pi_*\mathcal F$ with the inverse limit of
these zero groups. Nakayama's lemma gives $R^2\pi_*I_Z=0$.
The exact sequence of $Z$ gives a surjection
$R^1\pi_*\OO_T\to R^1\pi_*\OO_Z$, and therefore
\begin{equation}\hypertarget{tag-01OE}{}\label{eq:P4exceptional-genus}
 p\ge h^1(Z,\OO_Z).
\end{equation}

\noindent\textbf{Step 3. Exclusion of normal singular supports.}
Suppose first that $E=0$. Then $K_T$ is torsion and the
classification \cite[Chapter~VI, Theorem~1.1, p.~243]{BHPV2004E} gives Kodaira dimension zero and nonnegative
holomorphic Euler characteristic. Equation
\eqref{eq:P4chi-resolution} forces $p=0$ and $\chi(\OO_T)=0$.
The remaining surfaces are tori, bielliptic surfaces and Kodaira
surfaces. They have no rational curves: a map from $\Pp^1$ to
a torus is constant, bielliptic surfaces have finite unramified
torus covers, and a primary Kodaira surface has an elliptic
fibration with smooth elliptic reduced fibres over an elliptic
curve. A map from $\Pp^1$ to a secondary Kodaira surface lifts
to its finite unramified primary Kodaira cover. On a torsion
canonical surface a negative curve is a smooth rational
$(-2)$-curve by adjunction. Thus there are no exceptional
curves, and $S$ is smooth.

Suppose next that $E\ne0$. Since $K_T+E$ is torsion, all
positive plurigenera vanish, so $T$ is rational, ruled, or a
blowup of a class-VII surface. Rationality contradicts
\eqref{eq:P4chi-resolution}, since $\chi(\OO_T)=1$. In class
VII that equation gives $p=0$, whereas
\[
 \chi(\OO_E)=-E\cdot(E+K_T)/2=0.
\]
The nonempty effective compact divisor $E$ has $h^0(\OO_E)\ge1$,
so $h^1(\OO_E)\ge1$, contradicting
\eqref{eq:P4exceptional-genus}.

Finally suppose $T$ is ruled, after point blowups, over a curve
$C$ of genus $g\ge1$. Equation \eqref{eq:P4chi-resolution}
gives $p=g-1$. A general ruling fibre $F$ has $K_T\cdot F=-2$,
so $E\cdot F=2$. An exceptional component $E_j$ therefore
maps onto $C$. Its normalization has genus at least $g$ by
Riemann--Hurwitz, so $h^1(\OO_{E_j})\ge g$. Taking $Z=E_j$
in \eqref{eq:P4exceptional-genus} gives $p\ge g$, a contradiction.

\noindent\textbf{Step 4. The resolved normalization of a nonnormal support.}
Write $\nu:N\to S$ for the finite normalization and $\pi:T\to N$
for its minimal resolution. The residue comparison of dualizing
sheaves gives a pole divisor along the conductor; its strict
transform $D$ on $T$ is nonzero and effective. Pull back a local generator of $\omega_S$ as a meromorphic
two-form. Its canonical divisor has the form $-D+\sum a_jE_j$.
The equations for the exceptional coefficients have right side
$(K_T+D)\cdot E_j\ge0$, because both $K_T\cdot E_j$ and
$D\cdot E_j$ are nonnegative. The same inverse-matrix argument
gives $a_j\le0$. Taking the $m$th power to trivialize the pulled
back $\omega_S$ consequently gives
\[
 mK_T+D'\sim0
\]
for a nonzero effective integral divisor $D'$. This is also the
conductor formula of \cite[Lemma~4.6, p.~685]{CDP2020}. Any
positive pluricanonical section on $T$, multiplied by the
appropriate effective section and raised to a power, would give
a nonzero holomorphic function on compact $T$ with a zero.
It follows that $\kappa(T)=-\infty$. Contraction of exceptional
curves of the first kind preserves the plurigenera, so the minimal
model has the same Kodaira dimension.

\end{proof}



\subsection{Monodromy and Euler contributions}

\begin{proposition}[Normal crossing torus fibres and their normalizations]\hypertarget{tag-02Y0}{}\label{night:P4.8}\label{stmt:P4.8}
\StatementRecord{night:P4.8}{proved}{02Y0}{\SeriesRef{night:P4.12}, \SeriesRef{night:P4.6}, \SeriesRef{night:P4.7}}

For a fibration in Proposition~\ref{night:P4.7} with general fibre
a complex two-torus, an irreducible singular support with
rational normalization has Euler number two; every other irreducible
support has Euler number zero. The supports of Euler number zero have the following forms: a smooth support has a finite unramified
torus cover; a singular irreducible support has smooth minimal
elliptic-ruled or Hopf normalization; a reducible support has smooth
minimal Hopf or non-Kähler elliptic normalization components, with
smooth elliptic double curves and finite unramified branch maps.
Thus $c_3[X]=0$ precisely when every fibre belongs to these zero cases.
\end{proposition}
\begin{proof}

\begin{enumerate}[label=\arabic*.,leftmargin=*]
\item We apply \cite[Chapter~VI, Theorem~1.1 and Table~10, pp.~243--244]{BHPV2004E} to the smooth compact
normalization components, using the identity $e(S_i)=K_{S_i}^2$ proved
in Proposition~\ref{night:P4.6}. The relevant non-Kähler classification is stated in
\cite[arXiv version~3, Theorem~4.2, pp.~17--18; Theorem~5.1 and proof, pp.~27--29]{OVV2020}: a minimal non-Kähler
surface is elliptic or of class VII, and the reduced fibres of a
minimal non-Kähler elliptic surface are smooth elliptic curves.

\item For a component with nonempty conductor, a Kähler form $\omega$ would
give $K_{S_i}\cdot\omega=-D_i\cdot\omega<0$. No positive power of
the canonical bundle can then have a nonzero section. The classification
makes a Kähler component rational or ruled. If it is obtained by $k$
blowups of a ruled surface over a genus-$g$ curve, then
\[
 e=4(1-g)+k,\qquad K^2=8(1-g)-k.
\]
Equality gives $k=2(1-g)$, so $g\le1$ and the genus-one case is minimal.
A rational component has $\chi(\OO)=1$, hence $e=6$ here, including
the case with minimal model $\Pp^2$.

\item A minimal non-Kähler elliptic surface has $e=K^2=0$, by its smooth
elliptic reduced fibres. A minimal class-VII surface has $b_1=1$ and
$p_g=0$, so $\chi(\OO)=0$ and $K^2=-e$, with $e=b_2\ge0$.
Blowups preserve $\chi(\OO)$ and change $(e,K^2)$ by $(1,-1)$.
Consequently $e=K^2$ forces minimality and $e=0$ in all these
non-Kähler cases. The class-VII case has $b_2=0$. The finite-cover classification
\cite[arXiv version~3, Theorem~1, p.~1]{BBK2019} gives a finite
unramified map $p:\widehat S_i\to S_i$ with $\widehat S_i$ Hopf
or Inoue. The inverse image $p^{-1}(D_i)$ is a nonempty compact
curve: surjectivity makes it nonempty, and an unramified map preserves
its local codimension. Inoue surfaces have no compact curves
\cite[journal edition, Theorem~3.1, p.~303]{KhaledTeleman2026}; the type-III
family is covered by the explicit unramified type-II double cover
of \cite[journal edition, proof of Remark~3.5, p.~305]{KhaledTeleman2026}.
Thus $\widehat S_i$ is Hopf. Its universal covering
$\C^2\setminus\{0\}\to\widehat S_i$, composed with $p$, is a
universal covering of $S_i$, which is therefore a Hopf surface.

\item For an irreducible singular support there is one normalization
component. Proposition~\ref{night:P4.6} gives $T=2\chi(\OO_{S_1})$, which is two in the
rational case and zero in the other cases. Proposition~\ref{night:P4.12} proves that the
dualizing sheaf of an irreducible support is torsion. Adjunction on
the normalization therefore makes $K_{S_1}+D_1$ torsion. Since
$D_1\ne0$, every positive plurigenus of $S_1$ vanishes: multiplying
a pluricanonical section by an appropriate positive power of the
conductor section would give a section of a torsion line bundle with
a zero, whose sufficiently high power is a nonzero constant. Thus a
non-Kähler normalization in this irreducible case is of class VII,
and hence is the Hopf case already obtained.

\item For a reducible connected fibre, each normalization component meets
another global irreducible component. Pulling back the divisor of
that component gives a nonzero effective divisor on $S_i$ with zero
real first Chern class. A Kähler form would pair positively with it,
contradicting this zero class. Hence every component is non-Kähler,
has $\chi(\OO)=0$, and $T=0$. The conductor components are disjoint
and smooth, because all their nodes would come from triple points.
For any conductor component $C$, adjunction and $K+D\equiv0$ give
$2g(C)-2=(K+C)\cdot C=0$, since the other conductor components
are disjoint. Thus the double curves are smooth elliptic, and their
branch maps are the unramified maps already described in Proposition~\ref{night:P4.6}.

\item Finally let the support $G$ of a multiple fibre $mG$ be smooth. The
normalized cyclic base change $t=u^m$ is locally a disjoint union of
smooth charts: the original equation is $t=z^m$ times a holomorphic
unit, whose root can be extracted locally. These charts show that
the central fibre of the base-changed family maps finitely and
unramifiedly to $G$, and that its map to the $u$-disc is a proper submersion. Its nearby
fibres are the original tori. Catanese's theorem
\cite[definition of deformation equivalence, p.~261; Theorem~2.1 and proof, pp.~264--268]{Catanese2004} states that every
deformation of a complex torus is a complex torus; its definition
uses precisely a proper holomorphic submersion over a connected
disc. The normalized cyclic base change constructed above is such
a proper holomorphic submersion, so its central fibre is a torus;
Ehresmann's theorem also shows that it is connected. This gives
the finite unramified torus cover of $G$.

\end{enumerate}
\end{proof}

\begin{proposition}[Monodromy traces and minimum-multiplicity strata]\hypertarget{tag-02WK}{}\label{night:P4.13}\label{stmt:P4.13}
\StatementRecord{night:P4.13}{proved}{02WK}{\cite[\S2, pp.~237--243, and \S3, pp.~243--245]{ACampo1975}}


Let $f:U\to\Delta$ be a proper holomorphic map from a smooth complex
threefold, with connected fibres and smooth complex two-torus fibres
over $\Delta^*$. Choose an embedded resolution
$\pi:U'\to U$, unchanged away from the central fibre, and write
the full central divisor of $f\pi$ as $\sum_jN_jE_j$, with simple
normal crossings. Define
\begin{equation}
E_j^\circ=E_j\setminus\bigcup_{i\ne j}E_i,
 \qquad a_n=\sum_{N_j=n}e(E_j^\circ).
 \hypertarget{tag-00RH}{}\label{cdp:eq:C.MS.1}
\end{equation}
Let $T$ be the monodromy on $H^1(F,\Z)$ of a nearby fibre and $M$
the order of its semisimple part. Then $T\in\operatorname{SL}_4(\Z)$
is quasi-unipotent and
\begin{equation}
\begin{aligned}
 L_k&=\det(I-T^k)
       =\sum_{n\mid k}n a_n,\\
 a_n&=\frac1n\sum_{d\mid n}\mu(n/d)\det(I-T^d).
 \end{aligned}                                  \hypertarget{tag-00RJ}{}\label{cdp:eq:C.MS.2}
\end{equation}
where $\mu$ is the integer M\"obius function. On $H^1(F,\Q)$,
\begin{equation}
(T^M-I)^2=0,
                  \qquad\operatorname{rank}(T^M-I)\le2.
 \hypertarget{tag-00RI}{}\label{cdp:eq:C.MS.10}
\end{equation}
Every $L_k$ is nonnegative. The sequence $(a_n)$ vanishes exactly
when $T$ has eigenvalue one.

Write the original full central fibre as $\sum_i m_iS_i$, put
$m_0=\min_i m_i$, and let $S_i^\circ$ be the smooth locus of $S_i$
with its intersections with the other components removed. Then
\begin{equation}
m_0\sum_{m_i=m_0}e(S_i^\circ)
                    =\det(I-T^{m_0})\ge0.        \hypertarget{tag-00RL}{}\label{cdp:eq:C.MS.4}
\end{equation}
For $m_0>1$ the entire sequence $(a_n)$ vanishes. In particular, for
an irreducible full fibre $mS$, with monodromy $A$ on $H_1(F,\Z)$,
\begin{equation}
\begin{gathered}
 m e(S_{\mathrm{reg}})
   =\sum_{q=0}^4(-1)^q\operatorname{Tr}(\Lambda^qT^m)
   =\det(I-A^m),\\
 0\le e(S_{\mathrm{reg}})\le16/m,\qquad
 m>1\ \Longrightarrow\ e(S_{\mathrm{reg}})=0 .
\end{gathered}
 \hypertarget{tag-0150}{}\label{cdp:eq:C.UI.8}
\end{equation}
\end{proposition}
\begin{proof}
\begin{enumerate}[label=\arabic*.,leftmargin=*]
\item The resolution is an isomorphism over
$\Delta^*$, so it identifies the nearby
torus and its monodromy before and after resolution.
The proper global construction in
\cite[\S2, pp.~237--243, and \S3]{ACampo1975}
gives the nearby-cycle cohomology sheaves on the
resolved central support and their spectral sequence
to $H^*(F,\Q)$. In a chart
$z_1^{N_1}\cdots z_r^{N_r}=t$, a local nearby
fibre has $\gcd(N_1,\ldots,N_r)$ components,
each with the homotopy type of $(S^1)^{r-1}$.
Monodromy permutes these components cyclically;
after their permutation is killed its action
is a translation on the torus and induces the
identity on cohomology. A common multiple of all
the $N_j$ therefore acts identically on every
stalk of every nearby-cycle cohomology sheaf.
It acts identically on the corresponding
cohomology groups of those sheaves and hence
on the associated graded of the finite
spectral-sequence filtration of $H^*(F,\Q)$.
Its difference from the identity is nilpotent
on that cohomology. This proves quasi-unipotence
for this proper analytic normal-crossings
model.
In degree one this filtration has only the
two possible terms $H^1(E,\mathcal H^0\psi)$
and $H^0(E,\mathcal H^1\psi)$, where
$E=\bigcup_jE_j$ and $\mathcal H^q\psi$
denotes the nearby-cycle cohomology sheaf
already described. Thus for a
common multiplicity multiple $N$ one has
$(T^N-I)^2=0$. The semisimple order $M$
divides $N$, and $U=T^M$ is unipotent.
Writing
$U^{N/M}-I=(U-I)(I+U+\cdots+U^{N/M-1})$,
the second factor is invertible over $\Q$
and commutes with $U$. Hence $(U-I)^2=0$.
Its image is contained in its kernel on
the four-dimensional space, proving
the rank bound in \eqref{cdp:eq:C.MS.10}.
\item The same local model computes the alternating
stalk trace of the $k$th monodromy power. At a
point of $E_j^\circ$ it is $N_j$ if $N_j\mid k$
and zero otherwise. At a point lying on two or
more components it is zero, since the torus
factor has positive dimension and Euler number
zero. Taking the alternating trace through the
constructible spectral sequence, or its finite
cell filtration, gives
\begin{equation}
\sum_q(-1)^q\operatorname{Tr}
             (T^k;H^q(F,\Q))
       =\sum_{N_j\mid k}N_j e(E_j^\circ).
 \hypertarget{tag-00RN}{}\hypertarget{tag-014Y}{}\SsixEquationAliases{\label{cdp:eq:C.UI.6}}\label{cdp:eq:C.MS.6}
\end{equation}
This is the global form of the explicit local
calculation in \cite[Theorem~1, p.~235,
and proof on pp.~243--244]{ACampo1975}.
The cup-product map
$\Lambda^qH^1(F,\Z)\to H^q(F,\Z)$ is an
isomorphism for a torus and commutes with the
monodromy. Thus the left side of
\eqref{cdp:eq:C.MS.6} is exactly $\det(I-T^k)$.
The complex orientation of the smooth fibre
is preserved by transport. The map on its
top cohomology is consequently the identity;
under the fourth exterior power this says
$\det T=1$. Integer M\"obius inversion now
gives both equations of \eqref{cdp:eq:C.MS.2}.
\item The characteristic polynomial of $T$ has real coefficients.
Its nonreal eigenvalues occur in conjugate root-of-unity pairs,
whose contribution to $\det(I-T^k)$ is $|1-\lambda^k|^2$.
The real eigenvalues are $1$ and $-1$, contributing zero or two.
Thus $L_k\ge0$ and $L_k\le16$. An eigenvalue one makes every
$L_k$, and hence every $a_n$, zero. Conversely, the absence of
that eigenvalue makes every factor of $L_1$ positive, so
$a_1=L_1>0$.
\item For \eqref{cdp:eq:C.MS.4}, choose the resolution with centres
over the singular locus of the original reduced
support. Its first exceptional multiplicities
are at least $2m_0$: at such a centre either
one support component has order at least two
or at least two components occur. All later
exceptional multiplicities are sums of the
orders of the components containing their
centres. If the generic point of a centre is
contained in an existing exceptional component,
that component already contributes at least
$2m_0$. Otherwise the map is locally an
isomorphism at its generic point and the
original reduced support still has order at
least two there. Thus this case has the same
lower bound, including successive resolutions
over disjoint original singular points.
In particular the exceptional components satisfy
\begin{equation}
 N_j\ge2m_0\quad\text{for every exceptional }E_j.
 \hypertarget{tag-014X}{}\label{cdp:eq:C.UI.5}
\end{equation}
The open stratum of a strict
transform is exactly the
$S_i^\circ$ defined in the statement. In \eqref{cdp:eq:C.MS.6} for $k=m_0$, only
strict transforms of multiplicity $m_0$
contribute; this proves \eqref{cdp:eq:C.MS.4}.
\item For $m_0>1$, the resolved divisor has no multiplicity-one
component: strict-transform multiplicities are at least $m_0$,
and the exceptional multiplicities have the bound just proved.
Hence $a_1=0$ and the preceding criterion makes the entire
sequence vanish. For the irreducible full fibre $mS$,
\eqref{cdp:eq:C.MS.4} gives $m e(S_{\mathrm{reg}})=L_m$.
Under the evaluation pairing $H^1(F;\C)=H_1(F;\C)^\vee$, the
cohomology action is $A^T$, since $(T\alpha)(c)=\alpha(Ac)$.
Consequently $L_m=\det(I-(A^m)^T)=\det(I-A^m)$.
The bound $0\le L_m\le16$ and the vanishing for $m>1$ prove
all the remaining assertions.
\end{enumerate}
\end{proof}


\subsection{Finite conductor covers and the normalization differential}

\begin{proposition}[Genus-one curves with nonreduced conductor]\hypertarget{tag-02WL}{}\label{night:P4.14}\label{stmt:P4.14}
\StatementRecord{night:P4.14}{proved}{02WL}{the definitions and constructions stated in the proof}

The reducible-curve inference in \cite[proof of Proposition~7.3,
p.~694]{CDP2020} omits the two plane cubics
\[
 Z(Y^2-XZ)=0,\qquad XY(X-Y)=0.
\]
They are reduced Gorenstein curves of arithmetic genus one with
trivial dualizing sheaf and rational normalization, but have
respectively two tangent components and three concurrent
components. Their conductor order is two on every normalization
branch. These curves refute the asserted nodal-cycle
classification at that step.
\end{proposition}
\begin{proof}

\begin{enumerate}[leftmargin=*]
\item Plane adjunction gives $\omega_C=\OO_C(3-3)=\OO_C$ and
$p_a(C)=1$. The first cubic has a line tangent to a smooth
conic at $[1:0:0]$. Its local equation is $y(y-x^2)=0$.
The normalization restrictions are $(p(t,0),p(t,t^2))$.
Their image is exactly the pairs $(a,b)$ with $b-a$ divisible
by $t^2$: the inverse expression is
$p(x,y)=a(x)+y(b(x)-a(x))/x^2$. Thus the normalization
quotient has length two and the conductor exponent is two on
each branch.

\item For the second cubic the branches are $(t,0),(0,t),(t,t)$.
All constant coefficients must agree, and the third linear
coefficient is the sum of the first two. These three conditions
are sufficient. In every degree $k\ge2$, the monomials
$x^k-x^{k-1}y$, $y^k-xy^{k-1}$ and $x^{k-1}y$ independently
realize the coefficients on the three branches. Hence the
normalization quotient has length three and each conductor
exponent is again two. The normalization Euler characteristics
are two and three; subtracting these quotient lengths gives
$\chi(\OO_C)=0$ in both cases.

\item Let $C_2$ and $C_3$ be the two cubics above. The normalization
of $C_2$ is the disjoint union of two copies of $\Pp^1$, parametrized
by $[s:t]\mapsto[s:t:0]$ on the line and
$[s:t]\mapsto[s^2:st:t^2]$ on the conic. The parameters $t/s$ agree
at the two preimages of the tacnode. Exchanging these copies with
these parameters preserves their local ring
$\{(a,b):b-a\in(t^2)\}$ and descends to an involution $\sigma_2$
of $C_2$. In the local coordinates above it is
$(x,y)\mapsto(x,x^2-y)$. The projective automorphism
\[
 \sigma_3[X:Y:Z]=[-Y:X-Y:Z]
\]
has order three and cyclically permutes the components of $C_3$.
For $r=2,3$, take $E'=\C/(\Z+i\Z)$ and let $\tau_r$ be
translation by $1/r$. Its order is exactly $r$, and every nonidentity
power acts freely. Thus the diagonal action defines the compact
surface and its map to an elliptic curve
\[
 X_r=(C_r\times E')/\langle(\sigma_r,\tau_r)\rangle
       \longrightarrow E=E'/\langle\tau_r\rangle.
\]
The cyclic group acts transitively on the components of
$C_r\times E'$, so $X_r$ is irreducible. Let $\widetilde\sigma_r$
be the lift to the normalization $\nu_r:C_r^\nu\to C_r$ induced
by the displayed parameter maps, so
$\nu_r\widetilde\sigma_r=\sigma_r\nu_r$. The normalization of $X_r$ is
$(\coprod_{j=1}^r\Pp^1\times E')/\langle(\widetilde\sigma_r,\tau_r)\rangle$.
Choose a component $C_{r,0}$ of $C_r^\nu$. Restriction of the quotient
map to $C_{r,0}\times E'$ is an isomorphism. Its inverse sends a class
represented by $(z,e)$ with $z\in\widetilde\sigma_r^j C_{r,0}$ to
$(\widetilde\sigma_r^{-j}z,\tau_r^{-j}e)$; this expression is unchanged
on each orbit and is holomorphic on every component.
The reduced inverse-image conductor is the single section
$\Gamma'=\{p\}\times E'$, its conductor divisor is $2\Gamma'$, and
its map to the singular curve $E$ is the finite unramified cover
$E'\to E$ of degree $r$. Analytically locally on $E$, the image
conductor has algebra $\OO_E[u]/(u^2)$ for $r=2$ and
$\OO_E[u,v]/(u,v)^2$ for $r=3$. These local algebras follow from
the normalization constraints in the first two steps and remain
nonreduced after the free quotient. The constructions give a single
horizontal conductor component of multiplicity two with generic
degree two or three over the image conductor.

\end{enumerate}
\end{proof}

\begin{remark}[Local threefolds with permuted conductor branches]
\hypertarget{tag-P242}{}\label{rem:II:permuted-conductor-branches}
\StatementRecord{rem:II:permuted-conductor-branches}{proved}{P242}{\SeriesRef{night:P4.14}}
The branch permutations in \cref{night:P4.14} give local smooth
threefolds with irreducible divisors. Take an elliptic curve $C$
with a free translation $\tau$ of order two or three. The map
$(x,y)\mapsto(x,x^2-y)$ exchanges the two tacnode branches and
preserves $y(y-x^2)$. The order-three map
$(x,y)\mapsto(-y,x-y)$ cycles the three branches and preserves
$xy(x-y)$. The diagonal action on $C$ times a small invariant
neighbourhood of the origin is free. Its quotient is a smooth
threefold, and the divisor is irreducible because the action
permutes its branches transitively. Its normalization conductor
is one elliptic curve mapping to $C/\langle\tau\rangle$ with
degree two or three.
\end{remark}

\begin{lemma}[Minimal models of negative Kodaira dimension]
\hypertarget{tag-P244}{}\label{lem:II:negative-kodaira-models}
\StatementRecord{lem:II:negative-kodaira-models}{proved}{P244}{}
Let $T_0$ be a connected smooth compact minimal complex surface with
$\kappa(T_0)=-\infty$.
\begin{enumerate}[label=(\roman*),leftmargin=*]
\item A K\"ahler $T_0$ is projective and is biholomorphic to
$\mathbb P^2$ or to a projective-line bundle over a smooth compact curve.
\item A non-K\"ahler $T_0$ is of class VII and satisfies
$b_1(T_0)=h^1(T_0,\mathcal O_{T_0})=1$ and
$\chi(\mathcal O_{T_0})=0$.
\item A $T_0$ in~\textup{(ii)} with $b_2(T_0)=0$ and a nonempty
curve is a Hopf surface, with universal cover
$\mathbb C^2\setminus\{0\}$.
\end{enumerate}
\end{lemma}
\begin{proof}
The hypothesis $\kappa(T_0)=-\infty$ gives $p_g(T_0)=0$.
For a K\"ahler $T_0$, every real degree-two cohomology class is
therefore of type $(1,1)$. A sufficiently close rational perturbation
of a K\"ahler class, followed by integral scaling, gives an integral
K\"ahler class. Kodaira's embedding theorem
\cite[Main Theorem and proof, pp.~314--315]{Kodaira1954} gives an
embedding $\Phi:T_0\hookrightarrow\mathbb P^N$.
Apply \cite[VI.17--VI.18, pp.~111--112]{Beauville1978} to the
projective surface $\Phi(T_0)$: its pluricanonical systems vanish,
so it is ruled. Minimality and
\cite[V.10, pp.~79--80; III.10, p.~40]{Beauville1978} give an
isomorphism of $\Phi(T_0)$ with $\mathbb P^2$ or with a
projective-line bundle. Compose this isomorphism, and in the latter
case its projection to the base curve, with $\Phi$ to obtain the
asserted biholomorphism and ruling of $T_0$.

For a non-K\"ahler $T_0$, \cite[Theorem~5.1, pp.~27--29]{OVV2020}
gives an isotrivial elliptic surface or a class-VII surface, with
overlap permitted. In the elliptic alternative, the local maps in
\cite[Theorem~4.2 and Remark~4.4, pp.~17--21]{OVV2020} are
\[
 (E\times\Delta_s)/\langle(z,s)\mapsto(z+e,\zeta_m s)\rangle
 \longrightarrow\Delta_t,\qquad t=s^m,
\]
where $e$ has exact order $m$. The elliptic action of
\cite[Proposition~4.5, pp.~21--22]{OVV2020} gives a nowhere-zero
vertical vector field. For the fibration $p:T_0\to B_0$, its
quotient line bundle $Q=T_{T_0}/\mathcal O_{T_0}$ maps under $dp$
to $p^*T_{B_0}$ with vanishing divisor
$R_p=\sum_i(m_i-1)E_i$, since $dt=m_i s^{m_i-1}ds$ in the local
model. Thus
\[
 Q\simeq p^*T_{B_0}(-R_p),\qquad
 K_{T_0}\simeq p^*K_{B_0}\otimes\mathcal O_{T_0}(R_p).
\]
Distinct fibres are disjoint and each reduced multiple fibre has
self-intersection zero, so $K_{T_0}^2=0$. The nowhere-zero vector
field gives $c_2(T_0)=0$. Noether's formula and $p_g=0$ give
$h^1(\mathcal O_{T_0})=1$; the identity
$b_1=2h^1(\mathcal O)-1$ in
\cite[Theorem~2.5, pp.~7--9]{OVV2020} gives $b_1(T_0)=1$.
Together with $\kappa(T_0)=-\infty$, this proves that the elliptic
alternative also belongs to class VII. In the class-VII alternative
the same identity and $p_g=0$ give
$h^1(\mathcal O_{T_0})=1$ and $\chi(\mathcal O_{T_0})=0$.

Finally, for the surface in~\textup{(iii)},
\cite[Theorem~1, p.~1]{BBK2019} gives an \'etale cover
$a:T'\to T_0$ with $T'$ Hopf or Inoue. Its compact domain makes
$a$ finite. For a nonempty curve $R\subset T_0$, the inverse image
$a^{-1}(R)$ is a nonempty curve. Inoue surfaces contain no curves
by \cite[Proposition~2(i), pp.~274--275; Proposition~3(i), p.~277;
paragraph following (21), p.~279]{Inoue1974}; therefore $T'$ is Hopf.
The composite
\[
 \mathbb C^2\setminus\{0\}\longrightarrow T'
   \xrightarrow{a}T_0
\]
of its universal cover with $a$ is a covering with simply connected
domain. It is the universal cover of $T_0$, which proves that $T_0$
is Hopf.
\end{proof}


\begin{cited}[Curves on a Hopf surface]\hypertarget{tag-P248}{}\label{lem:II:hopf-curves}
\StatementRecord{lem:II:hopf-curves}{proved}{P248}{\cite[Theorem~9, pp.~232--233]{Kato1975Hopf}; \cite[arXiv v1, \S2.2.2, p.~4; \S2.2.6, p.~5; \S2.3.1, p.~8]{AbashevaDeev2024AB}}
Every irreducible compact curve on a Hopf surface is smooth elliptic and has self-intersection zero. Distinct such curves are disjoint.
\end{cited}
\begin{proof}[Application of the curve descriptions]
Theorem~9 of \cite[pp.~232--233]{Kato1975Hopf} presents a Hopf surface as a bundle over the circle with fibre a finite quotient of $S^3$. That fibre has zero first and second rational cohomology, by transfer from $S^3$, and the bundle cohomology sequence gives $b_2=0$. Thus all divisor intersection numbers vanish.

For a primary Hopf surface, \cite[arXiv v1, \S2.2.2, p.~4; \S2.2.6, p.~5]{AbashevaDeev2024AB} describes the non-diagonal case's single axis curve, the nonresonant diagonal case's two axis curves, and the resonant case's elliptic fibration with smooth reduced fibres. The axis curves are the quotients $\mathbb C^*/\langle\alpha\rangle$ and are smooth elliptic. In the elliptic fibration, a horizontal irreducible curve would have positive intersection with a general fibre, equal to its degree over the base, contrary to $b_2=0$. Every irreducible curve is therefore a reduced elliptic fibre. These descriptions make distinct curves disjoint in each primary case.

For a secondary Hopf surface $H$, use the finite unramified quotient $\pi:H'\to H=H'/G$ from a primary Hopf surface in \cite[arXiv v1, \S2.3.1, p.~8]{AbashevaDeev2024AB}. The inverse image of an irreducible curve $E$ is a union of the disjoint smooth elliptic curves just described. For a component $E'$, the restriction of $\pi$ identifies $E'/G_{E'}$ with $E$, where $G_{E'}$ is its stabilizer: the local quotient charts are biholomorphic because the action is free. The induced map $E'\to E$ is finite unramified, so $E$ is smooth elliptic. Inverse images of distinct curves are disjoint, and hence so are the curves themselves. Finally $(\pi^*E)^2=\deg(\pi)E^2=0$, giving square zero downstairs as well.
\end{proof}

\begin{lemma}[A single conductor branch under finite normalization]\hypertarget{tag-P249}{}\label{lem:II:single-conductor-branch}
\StatementRecord{lem:II:single-conductor-branch}{proved}{P249}{}
Let $S$ be a reduced hypersurface in a smooth complex threefold and $\nu:N\to S$ its finite normalization. Suppose that $C=\operatorname{Sing}(S)_{\rm red}$ is a curve and that $A=(\nu^{-1}C)_{\rm red}$ has exactly one local analytic curve branch at each of its points. Then $C$ has exactly one local analytic branch at each point; its distinct irreducible components are disjoint.
\end{lemma}
\begin{proof}
Fix $p\in C$ and two irreducible curve germs $C_1,C_2$ of $C$ at $p$. Choose normalization preimages whose germs of $A$ map onto these two germs. If the preimages coincide, the unique irreducible germ of $A$ there has one irreducible image, giving $C_1=C_2$. Otherwise the two preimages correspond to distinct irreducible analytic surface branches $S_1,S_2$ of $S$ at $p$. In the regular three-dimensional local ring of the ambient manifold, their nonassociate irreducible equations generate an ideal of height two. Their intersection consequently contains a curve germ through $p$. That intersection lies in the conductor image of each surface branch. Its inverse image on each normalization branch is contained in the unique germ of $A$ there, so its image is both $C_1$ and $C_2$. Again $C_1=C_2$. This proves local unibranchness and excludes intersections of distinct global irreducible components.
\end{proof}

\begin{lemma}[Numerically anticanonical boundaries on irrational ruled surfaces]\hypertarget{tag-P252}{}\label{lem:II:irrational-ruled-boundary}
\StatementRecord{lem:II:irrational-ruled-boundary}{proved}{P252}{\cref{lem:II:canonical-boundary-blowup,lem:II:boundary-blowup-groups}; \cite[III.10, p.~40; III.18, pp.~46--47]{Beauville1978}}
Let $r_0:T_0\to B$ be a minimal ruled surface over a smooth compact curve of genus $g\ge1$, let $\beta:T\to T_0$ be a sequence of point blowups, and let $D$ be an effective integral divisor with $K_T+D\equiv0$. Put $D_0=\beta_*D$ and $A=D_{\rm red}$.

The horizontal part of $D_0$ has ruling degree two, counted with coefficients. Its reduced case has $g=1$: $D_0$ is one smooth elliptic bisection or two disjoint smooth elliptic sections, unramified over $B$. It has no vertical component, and $A$ consists of their disjoint strict transforms, without any blowup-exceptional component.

The other case is
\[
 D_0=2H+\sum_j a_jF_j,\qquad a_j>0,
\]
where $H$ is a smooth section and the $F_j$ are distinct ruling fibres. Here $A$ is the strict transform of $H$ with trees of smooth rational curves attached, with only ordinary double intersections. Inclusion $j:A\hookrightarrow T$ induces an isomorphism on fundamental groups. This isomorphism commutes with the maps induced by each blowdown and with the section and ruling maps on $T_0$.
\end{lemma}
\begin{proof}
The projective-bundle and canonical formulas of \cite[III.10, p.~40; III.18, pp.~46--47]{Beauville1978} write $T_0=\mathbb P_B(V)$ and
\[
 h^2=\deg V,\quad hF=1,\quad F^2=0,\qquad
 K_{T_0}=-2h+(\deg V+2g-2)F.
\]
Apply \cref{lem:II:canonical-boundary-blowup} along $\beta$ to $D$ and its successive pushforwards. It gives $K_{T_0}+D_0\equiv0$, hence $D_0F=2$. It also gives exceptional coefficient $\mu-1$ at every reverse blowup, where $\mu$ is the multiplicity of the current divisor. Effectiveness puts every centre on its support.

For an irreducible horizontal component $H_i$ of coefficient one, adjunction gives
\[
 2p_a(H_i)-2=-(D_0-H_i)H_i\le0.
\]
Its normalization maps finitely to $B$, with degree $d\ge1$, so Riemann--Hurwitz gives $2g(H_i^{\rm norm})-2\ge d(2g-2)$. Since $p_a(H_i)\ge g(H_i^{\rm norm})$, these inequalities force $g=1$, $p_a(H_i)=g(H_i^{\rm norm})=1$ and $(D_0-H_i)H_i=0$. Thus $H_i$ is smooth elliptic, its map to $B$ is unramified, and every other horizontal component is disjoint from it. Every vertical fibre meets it positively, so the vertical part vanishes. The degree-two possibilities are exactly the stated bisection or two sections. A reverse blowup at a smooth coefficient-one point has exceptional coefficient zero and leaves its strict transform smooth. This proves the entire reduced case by induction.

Otherwise the horizontal degree two is twice a single degree-one component $H$. Its normalization maps isomorphically to $B$; composing the inverse with the normalization map into $T_0$ gives a section whose closed image is $H$. Hence $H$ is smooth. Every vertical component on $T_0$ is a fibre, giving the displayed divisor. Initially its reduced support is $H$ with transverse rational fibres attached. Under reverse blowup, a smooth point of coefficient at least two adds a rational leaf, a smooth coefficient-one point adds no boundary component, and a node subdivides its edge. Thus only rational trees are attached to the strict transform of $H$.

On $T_0$, the ruling and its section induce inverse maps on fundamental groups because the ruling fibres are simply connected. Attaching the rational fibres to $H$ adds no generators or relations. At every blowup, \cref{lem:II:boundary-blowup-groups} gives the two fundamental-group isomorphisms induced by that same blowdown and makes their square with the boundary inclusion commute. Their composition gives the asserted isomorphism $j_*$ and its precise compatibility.
\end{proof}

\begin{proposition}[Finite rank drop of a normalization]\hypertarget{tag-02WM}{}\label{night:P4.15}\label{stmt:P4.15}
\StatementRecord{night:P4.15}{proved}{02WM}{\SeriesRef{cdp:st:A.5.3}, \SeriesRef{night:P4.12}, \SeriesRef{night:P4.13}, \SeriesRef{lem:II:negative-kodaira-models}}

Let $f:Y\to B$ be a proper surjective holomorphic map with
connected fibres from a connected smooth compact complex threefold
with $b_1(Y)=b_2(Y)=0$ to a smooth compact curve, and suppose that
its smooth general fibre is a complex two-torus. Let $mS$ be a full
fibre whose reduced support $S$ is irreducible and singular.
Suppose its normalization $\nu:T\to S$ has smooth domain $T$. Set
$j:T\xrightarrow{\nu}S\hookrightarrow Y$ and assume that
$dj:T_T\to j^*T_Y$ has rank two away from finitely many points.
Let $\ell$ be the length of its maximal-minor scheme. Then
\[
 e(T)=K_T^2-\ell.
\]
The alternatives for $T$ are a rational surface with
$e(T)=6-\ell/2$ and $0\leq\ell\leq6$, a minimal ruled surface
over an elliptic curve with $\ell=0$, or a Hopf surface with
$\ell=0$. In either of the last two alternatives $e(S)=0$.
In those last two alternatives, the reduced conductor $R\subset T$
and $C=\operatorname{Sing}(S)_{\mathrm{red}}$ are disjoint unions of
smooth elliptic curves, and the restriction $\nu|_R:R\to C$ is
finite and unramified.
\end{proposition}
\begin{proof}

\begin{enumerate}[label=\arabic*.,leftmargin=*]
\item The differential is injective as a sheaf map: its kernel vanishes
on a dense open set, and its locally free source is torsion-free.
Its rank-one cokernel $Q$ is locally free away from finitely many
points. The depth lemma applied to the length-one locally free
resolution gives depth at least one at those points. Hence $Q$
has no finite-length subsheaf, and all its possible torsion would
be supported there. Thus $Q$ is torsion-free, and $L=Q^{**}$ is a
line bundle. In local frames the signed maximal-minor row of $dj$
annihilates its two columns. Its determinant transition functions
give a map to $L=\det(j^*T_Y)\otimes\det(T_T)^{-1}$.
The induced map $Q\to L$ is generically injective, hence injective
because $Q$ is torsion-free. Its image is the ideal generated by
those minors, tensored with $L$. We obtain the exact sequence
\[
 0\longrightarrow T_T\xrightarrow{dj}j^*T_Y
   \longrightarrow\mathcal I_Z\otimes L\longrightarrow0,
 \qquad L=K_T\otimes j^*(\det T_Y),\qquad
 \operatorname{length}(Z)=\ell.
\]
Poincar\'e duality gives $b_4(Y)=b_2(Y)=0$, so
$c_1(T_Y)=c_2(T_Y)=0$ in rational cohomology. Writing $k=c_1(K_T)$,
the Chern-character identity for this sequence is
\[
 3=\left(2-k+\frac{k^2-2c_2(T_T)}2\right)
     +\left(1+k+\frac{k^2}2-[Z]\right).
\]
Its degree-four part gives $e(T)=K_T^2-\ell$, and Noether's formula
gives $2e(T)+\ell=12\chi(\mathcal O_T)$.

Let $D$ be the full conductor divisor on $T$, and put $R=D_{\mathrm{red}}$.
Propositions~\ref{cdp:st:A.5.3} and~\ref{night:P4.12}, with the
identity as the resolution of $T$, give
$\omega_T(D)\simeq\nu^*\omega_S$ and $\omega_S^{\otimes m}\simeq\mathcal O_S$.
In particular $K_T+D\equiv0$. Proposition~\ref{night:P4.12}
also gives $\kappa(T)=-\infty$ and nonnormality of the singular $S$.
Since the hypersurface $S$ satisfies $S_2$, its nonnormal locus
contains a curve; thus $D$ is nonempty.
Let $q:T\to T_0$ contract $r$ point blowups to a minimal model.
Pluricanonical sections descend across each blowup by Hartogs'
extension on the complement of its centre, so $\kappa(T_0)=-\infty$
and $p_g(T_0)=0$.

Lemma~\ref{lem:II:negative-kodaira-models} applies to this smooth
compact minimal $T_0$. In its K\"ahler case, $T_0$ is
$\mathbb P^2$ or a projective-line bundle; in its non-K\"ahler case,
$T_0$ is of class VII with $\chi(\mathcal O_{T_0})=0$.

For a projective-line bundle over a genus-$g$ curve, followed by
$r$ point blowups, the bundle and blowup formulas
\cite[II.3, pp.~17--18; III.18, pp.~46--47]{Beauville1978} give
\[
 e(T)=4(1-g)+r,\qquad K_T^2=8(1-g)-r.
\]
The Chern identity becomes $2r+\ell=4(1-g)$. For $g\ge1$ it forces
$g=1$ and $r=\ell=0$.
For a class-VII minimal model, write $b_2(T_0)=k_0\ge0$.
The invariants in
Lemma~\ref{lem:II:negative-kodaira-models}\textup{(ii)} give
$e(T)=k_0+r$ and $K_T^2=-k_0-r$, and the Chern identity gives
$2(k_0+r)+\ell=0$. Hence $T$ is minimal with $b_2=0$ and $\ell=0$.
Its reduced conductor $R$ is a nonempty curve, so
Lemma~\ref{lem:II:negative-kodaira-models}\textup{(iii)}, applied to
this $T$, proves that $T$ is Hopf.

For rational $T$, $\chi(\mathcal O_T)=1$ and $e(T)\ge3$ by the
same rational minimal-model and blowup formulas. Hence
$e(T)=6-\ell/2$ and $0\le\ell\le6$.
When $\ell=0$, the maximal-minor ideal is the unit ideal, so $dj$
has rank two everywhere.

\item Consider the minimal elliptic ruled alternative. Apply \cref{lem:II:irrational-ruled-boundary} with $T_0=T$, $\beta=\operatorname{id}_T$ and this same conductor divisor $D$. Its reduced-horizontal case already makes $R$ a disjoint union of smooth elliptic curves. Its other case has $D=2H+\sum_j a_jF_j$, with $H$ a smooth elliptic section and distinct rational fibres $F_j$ meeting it once. Writing $v$ for their number gives $e(R)=v$ and
\[
 e(S_{\rm reg})=e(T\setminus R)=e(T)-e(R)=-v.
\]
Proposition~\ref{night:P4.13} forces $v=0$. Thus $R$ consists of smooth elliptic curves in this case as well.

\item In the Hopf alternative, \cref{lem:II:hopf-curves}, applied to this same $T$, makes the reduced conductor $R$ a disjoint union of smooth elliptic curves.

\item Apply \cref{lem:II:single-conductor-branch} to the finite normalization $\nu:T\to S\subset Y$, with $A=R$ and $C=\operatorname{Sing}(S)_{\rm red}$. The preceding cases give its unique smooth conductor branch at each point of $T$. Since $j:T\to Y$ is immersive, each image branch is smooth and the curve map is a local isomorphism. The lemma identifies all image branches at each point of $C$, so $C$ is smooth there.
Globally the restriction $\nu|_R:R\to C$ is finite and unramified.
For a component of degree $d$, Riemann--Hurwitz gives
$0=d(2g(C_i)-2)$, so every image component is elliptic.
Euler additivity for the normalization and its conductor gives
\[
 e(S)=e(T)-e(R)+e(C)=0.
\]

\end{enumerate}
\end{proof}

\subsection{Spherical classes at nonimmersive normalization points}

\begin{proposition}[Spherical intersection lattices]
\hypertarget{tag-02WN}{}\label{night:P4.16}\label{stmt:P4.16}
\StatementRecord{night:P4.16}{proved}{02WN}{}
\hypertarget{tag-02WO}{}\label{night:P4.17}\label{stmt:P4.17}
\StatementRecord{night:P4.17}{proved}{02WO}{}
For integers $a\ge2$ and $n\ge1$, the affine surface
$F_{a,n}=\{x^a-zy^n=1\}$ is simply connected and has compact
second-homology intersection lattice $A_{a-1}(-n)$.
Every class in this lattice is spherical.
The reduced hypersurface germ $x^a-zy^n=0$, with any positive
common multiplicity, is excluded from a fibre of a proper
holomorphic degeneration of complex two-tori with smooth total
threefold. For $a=2$ the spherical generator has square $-2n$;
this includes the Whitney umbrella at $n=2$.
The case $a=n=3$ gives the order-three cyclic-ramification germ.
The normalization morphisms below also include singular
normalizations.
\end{proposition}
\begin{proof}
\begin{enumerate}[label=\arabic*.,leftmargin=*]
\item Let $\zeta_1,\ldots,\zeta_a$ be the distinct roots of
$x^a=1$. The open subset $U=\{y\ne0\}$ of $F_{a,n}$ is
$\mathbb C^*\times\mathbb C$ in coordinates $(y,x)$.
For $0<\delta<1$, its complement is covered by the disjoint
graph neighbourhoods
\[
 N_i:\quad x=\zeta_i(1+zy^n)^{1/a},\qquad |zy^n|<\delta.
\]
The chosen root is holomorphic on the displayed disc.
Each $(y,z)$ domain contracts by simultaneous real scaling.
The intersection $N_i\cap U$ has coordinates $(y,zy^n)$
in $\mathbb C^*\times\{|w|<\delta\}$, and its generating
$y$-loop maps to the generating $y$-loop in $U$.
Van Kampen \cite[Theorem~1.20, p.~43]{Hatcher2002} gives
$\pi_1(F_{a,n})=0$. Mayer--Vietoris
\cite[\S2.2, pp.~149--150]{Hatcher2002} gives
\[
 H_2(F_{a,n};\mathbb Z)
 =\ker\bigl(\mathbb Z^a
       \xrightarrow{(c_i)\mapsto\sum_i c_i}\mathbb Z\bigr).
\]
The Hurewicz theorem \cite[Theorem~4.32, pp.~366--367]{Hatcher2002}
makes every class spherical.

\item Blow up $\mathbb P^1_x\times\mathbb P^1_y$ successively
$n$ times above each $(\zeta_i,0)$. At that point
$b=y$ and $c=x^a-1$ are local coordinates. The successive
charts are
\[
 c=b^j c_j,\qquad z=\frac{c_j}{b^{n-j}}.
\]
For $j<n$, the next centre is $c_j=b=0$.
At $j=n$, $z=c_n$ is regular. The last exceptional curve,
minus its point on the preceding boundary, is the line
$y=0,x=\zeta_i$ in $F_{a,n}$; every earlier exceptional
curve is a pole divisor of $z$.
Thus $F_{a,n}=T_{a,n}\setminus B$, where $T_{a,n}$ is the
resulting smooth projective surface and $B$ consists of the
two infinity divisors, the strict transform of $y=0$,
and the first $n-1$ exceptional curves on each chain.

\item In the total-transform basis $H_x,H_y,E_{i,j}$,
the nonzero products are $H_xH_y=1$ and $E_{i,j}^2=-1$.
The boundary classes are
\[
 H_x,\ H_y,\ H_y-\sum_iE_{i,1},\qquad
 E_{i,j}-E_{i,j+1}\quad(1\le i\le a,\ 1\le j<n).
\]
Their curves form a connected tree of rational curves, so
$H^1(B;\mathbb Z)=0$. Compact-support duality
\cite[Theorem~3.35, p.~245]{Hatcher2002} and the exact sequence
of $(T_{a,n},B)$ identify compact $H_2(F_{a,n};\mathbb Z)$
integrally with the kernel of intersection against these
boundary classes in $H_2(T_{a,n};\mathbb Z)$.
Orthogonality to $H_x,H_y$ removes the ruling coefficients;
orthogonality to the exceptional differences makes the
coefficients constant along each chain. The strict transform
of $y=0$ makes the sum of these constants zero.
The primitive integral basis is therefore
\[
 v_i=\sum_{j=1}^nE_{i,j}-\sum_{j=1}^nE_{i+1,j}
       \quad(1\le i<a).
\]
It has $v_i^2=-2n$, $v_iv_{i+1}=n$, and all other
off-diagonal products zero, giving $A_{a-1}(-n)$.
For $a=2$, this is precisely the difference of the two
exceptional chains and has square $-2n$.
The integral-kernel computation from the displayed boundary
classes is listed in \cref{tab:II:programs}.

\item The polynomial $g=x^a-zy^n$ has positive weights
$(n+1,a,a)$ and degree $d=a(n+1)$. The map
\[
 (x,y,z)\longmapsto
 (\varepsilon^{n+1}x,\varepsilon^a y,\varepsilon^a z)
\]
takes $F_{a,n}$ to the fibre of value $\varepsilon^d$.
It carries a compact spherical representative of $v_1$
and a compact transverse representative for its square
into an arbitrarily small neighbourhood of the origin.

For the germ-exclusion assertion, apply \cref{cdp:st:A.5.15}
with exponent triple $(a,b,c)=(a,n,1)$. The coordinate map
is the identity on $(x,y,z)$, and its full defining function
$u(x^a-y^bz^c)^m$ is exactly $u(x^a-zy^n)^m$.
The proper holomorphic map, smooth total threefold, torus
nearby fibres and arbitrary positive multiplicity are the
same hypotheses. This gives the asserted exclusion.

\item We specify the normalization morphism. Let $\mu_a$ act on
$\mathbb C^2_{u,v}$ by $(u,v)\mapsto(\zeta u,\zeta^{-n}v)$.
Then
\[
 x=u^nv,\qquad y=u^a,\qquad z=v^a
\]
gives a finite map to $x^a=zy^n$: both $u$ and $v$ satisfy
monic equations over its coordinate ring. The function-field
degree is $a$, with the displayed faithful $\mu_a$ action.
Consequently the induced quotient map is finite and birational.
The invariant ring is integrally closed, so this is the
normalization. For $(a,n)=(2,3)$ it has coordinates
$U=u^2,V=uv,W=v^2$ with $UW=V^2$, and its map is
$(x,y,z)=(UV,U,W)$. Thus it has an $A_1$ singularity.
For $n=a$ the quotient is smooth with coordinates $(y,v)$,
and the reduced conductor $\{y=0\}$ maps to $\{x=y=0\}$
by $z=v^a$. A positive loop about $z=0$ sends $v$ to
$e^{2\pi i/a}v$, cyclically permuting the $a$ branches.

\end{enumerate}
\end{proof}

\subsection{The integral pole divisor on a normalization resolution}

\begin{lemma}[Conductor support in a proper torus degeneration]\hypertarget{tag-P300}{}\label{lem:II:proper-conductor-divisor}
\StatementRecord{lem:II:proper-conductor-divisor}{proved}{P300}{\cref{cdp:st:A.5.3,lem:II:conductor-duality,cdp:st:A.5.17}}
Let $f:X\to\Delta$ be a proper holomorphic map from a smooth
complex threefold to a disc, with connected fibres, smooth
over $\Delta^*$, whose smooth fibres are complex two-tori.
Let its central fibre be $mS$, with $S$ irreducible and reduced.
Write $\nu:N\to S$ for normalization, $\pi:T\to N$ for the
minimal resolution, and $\rho=\nu\pi$. Let $D$ be the canonical
effective divisor of \cref{cdp:st:A.5.3} for these maps.
Every connected exceptional configuration of $\pi$ meets
the strict conductor and has strictly positive coefficients
in $D$. The reduced singular locus of $S$ is its reduced
conductor image, and
\[
       D_{\rm red}=\rho^{-1}((\operatorname{Sing}S)_{\rm red}).
\]
The map $\rho$ is an isomorphism between the complements
of these two reduced closed sets.
\end{lemma}
\begin{proof}
Apply \cref{cdp:st:A.5.3} with its normalization component
$T_i=N$, resolution surface $Z_i=T$, $\eta=\nu$,
$\pi_i=\pi$ and $h_i=\rho$. Its divisor is defined by the
orders of the pullback of the same local generator of
$\omega_S$. It gives the conductor multiplicities and
all exceptional-coefficient assertions once every connected
exceptional configuration meets the strict conductor.
Write $D_{\rm str}$ for the strict-conductor part of $D$,
and let $p\in N$ be a singular point. Its connected
exceptional fibre is $\pi^{-1}(p)$.

The finite-duality identity of \cref{lem:II:conductor-duality} makes
the conductor $\mathfrak c$ a rank-one reflexive ideal on the normal surface
$N$. Such an ideal is the intersection of its codimension-one
localizations in the fraction field. Hence
$\operatorname{Supp}(\mathcal O_N/\mathfrak c)$ has no isolated
point: a conductor ideal equal to $\mathcal O_N$ at all nearby
codimension-one points is equal to $\mathcal O_N$ itself.
A conductor curve through $p$ has strict transform meeting
$\pi^{-1}(p)$, by properness. Thus a block disjoint from
$D_{\rm str}$ has unit conductor at $p$. Put $q=\nu(p)$. Apply the unit-conductor assertion of
\cref{lem:II:conductor-duality} to $\nu:N\to S$, at this
same $p$ above $q$. Its local rings are
$R=\OO_{S,q}$ and the finite semilocal normalization
$B=(\nu_*\OO_N)_q$. The unit conductor at $p$ gives
$B=R$, so $\nu$ is an isomorphism near $q$.
Since $p$ is singular on $N$,
the image $q$ would be a normal isolated hypersurface
singularity.

At such a point choose a coordinate on the base vanishing
at the fibre. Its full local function is $f=u g^m$,
where $u$ is a unit and $g$ is the reduced local equation
of this normal isolated hypersurface singularity. The
critical-point germ of $g$ is isolated: each critical
component through the point has constant value zero,
so lies in the zero-dimensional singular locus of
$\{g=0\}$. The proper map restricted over a disc has
smooth nearby fibres complex two-tori, exactly the
hypotheses of \cref{cdp:st:A.5.15}. The isolated-critical-point
assertion of \cref{cdp:st:A.5.17} therefore applies
to this same $f,u,g,m$.

In that application its root is
$H=\exp((\log u)/m)g$, so $f=H^m$ exactly.
Its local-fibre inclusion is the identity on the chosen
coordinate neighbourhood,
\[
           K\cap H^{-1}(\delta)\hookrightarrow
                          f^{-1}(\delta^m).
\]
The proposition transports the compact Morse sphere
of square $-2$ into that local level and then into the
same torus fibre by this inclusion, which excludes
the isolated point.

Thus every exceptional block meets $D_{\rm str}$, and
\cref{cdp:st:A.5.3} gives strictly positive coefficients
on the whole block. Every normal point of $S$ is smooth by
the same local argument. The singular locus is therefore the
conductor locus, and its inverse image under $\rho$ consists
exactly of the strict conductor and these exceptional fibres.
This proves the support equality and the complement isomorphism.
\end{proof}


\begin{proposition}[Conductor and exceptional pole divisor]\hypertarget{tag-02YI}{}\label{night:P4R.6}\label{stmt:P4R.6}
\StatementRecord{night:P4R.6}{proved}{02YI}{\cref{lem:II:proper-conductor-divisor,night:P4.12}}

Let $f:Y\to B$ be a proper surjective holomorphic map with connected
fibres from a smooth compact complex threefold with $b_1(Y)=b_2(Y)=0$
to a smooth compact curve, and suppose that its smooth general fibre
is a complex two-torus. Let $mS$ be a fibre with irreducible, reduced,
singular support $S$, let $\nu:N\to S$ be its normalization, and let
$\pi:T\to N$ be its minimal resolution. Put $\rho=\nu\circ\pi$.
There is an effective integral divisor $D$ on $T$ such that
\[
 \omega_T(D)\simeq\rho^*\omega_S,
 \qquad \operatorname{Supp}D
       =\rho^{-1}(\operatorname{Sing}S)_{\rm red}.
\]
In particular $K_T+D$ is torsion, by \cref{night:P4.12}.
The strict-transform conductor has its usual positive integral
coefficients. Every connected exceptional configuration over a
singular point of $N$ has strictly positive coefficients in $D$.
Every such configuration meets the strict-transform conductor.
\end{proposition}
\begin{proof}
Restrict $f$ to a sufficiently small disc about the
value of $mS$ and apply \cref{lem:II:proper-conductor-divisor}
to this restriction, with the same $S,N,T,\nu,\pi,\rho$.
Its divisor $D$ is defined by the orders of the pullback
of the same local generator of $\omega_S$. It gives the
line-bundle identity, effectiveness, support equality
and all assertions about conductor and exceptional
coefficients in the statement.

Under the global hypotheses here, \cref{night:P4.12}
gives $\omega_S^{\otimes m}\simeq\OO_S$. Therefore
\[
 (\omega_T(D))^{\otimes m}
   \simeq\rho^*(\omega_S^{\otimes m})\simeq\OO_T,
\]
which proves the torsion conclusion.
\end{proof}


\subsection{Normalization cohomology and conductor schemes}

\begin{proposition}[Normalization cohomology and the conductor Euler characteristic]\hypertarget{tag-02X6}{}\label{night:P4.33}\label{stmt:P4.33}
\StatementRecord{night:P4.33}{proved}{02X6}{}

Let $f:Y\to B$ satisfy the torus-fibration hypotheses of
\cref{night:P4.12}, and let $mS$ be a fibre with irreducible,
reduced, singular support. Write $\nu:N\to S$ for the normalization,
$\rho_N:T\to N$ for its minimal resolution and $\rho=\nu\circ\rho_N$.
Then $H^2(N,\mathcal O_N)=0$. If
$p=\dim H^0(N,R^1\rho_{N*}\mathcal O_T)$ is the sum of the
geometric genera of the normalization singularities, then
\[
 p=h^1(T,\mathcal O_T)-h^1(N,\mathcal O_N).
\]
In particular a rational resolution forces every normalization
singularity to be rational and gives $\chi(\mathcal O_N)=1$.
Let $C_{\rm sch}\subset S$ and $D_{\rm sch}\subset N$ be the image and inverse-image conductor schemes. The nonempty $C_{\rm sch}$ is Cohen--Macaulay of pure dimension one, and
\[
 \chi(\mathcal O_{D_{\rm sch}})=0,\qquad
 \chi(\mathcal O_{C_{\rm sch}})
   =-\chi(\mathcal O_N)=h^1(N,\mathcal O_N)-1.
\]

For the ruled case, suppose that $T$ is obtained by point blowups
from a $\mathbb P^1$-bundle over a curve $E$ of genus $g\ge1$,
that the induced ruling has a section $H$, and that the pole
divisor of \cref{night:P4R.6} has horizontal part $2H$ and
rational vertical trees as its remaining support. A contracted
$H$ gives $h^1(N,\mathcal O_N)=0$ and $p=g$.
For a noncontracted $H$, write $a:T\to E$ for the ruling and $h:E\simeq H\hookrightarrow T$ for the section. There is a holomorphic map $q:N\to E$ and a section $s=\rho_N\circ h$ satisfying
\[
 q\circ\rho_N=a,\qquad q\circ s=\operatorname{id}_E.
\]
Here $p=0$, $R^1\rho_{N*}\mathcal O_T=0$, and $h^1(N,\mathcal O_N)=g$.
\end{proposition}
\begin{proof}

\begin{enumerate}[label=\arabic*.,leftmargin=*]
\item Apply \cref{lem:II:conductor-duality} to the finite normalization $\eta=\nu:N\to S$. The hypersurface $S\subset Y$ is Gorenstein, and the normal surface $N$ is Cohen--Macaulay. The lemma identifies the conductor $I_S$ with $\mathcal Hom_{\mathcal O_S}(\nu_*\mathcal O_N,\mathcal O_S)$ by evaluation at $1$, whose inverse is multiplication by that value. Its finite-duality map and the projection formula give
\[
 \nu_*\omega_N=I_S\otimes\omega_S.
\]
It also makes the image conductor scheme $C_{\rm sch}$ Cohen--Macaulay of pure dimension one.
The dualizing line bundle $\omega_S$ is torsion by \cref{night:P4.12}.
A nonzero section of a torsion line bundle on a compact
irreducible reduced complex space has no zero: raising it
to a trivializing positive power gives a nonzero holomorphic
function, hence a nonzero constant. Such a section cannot
belong to the proper conductor ideal along the nonempty
conductor. Therefore $H^0(N,\omega_N)=0$. Analytic Serre duality
on the compact Cohen--Macaulay surface \cite[Theorem~1, pp.~81--82,
and the depth calculation on p.~89]{RR70push} gives
$H^2(N,\mathcal O_N)=0$.

\item Proposition~\ref{night:P4.12} gives $\kappa(T)=-\infty$,
so $H^0(T,\omega_T)=0$ and $H^2(T,\mathcal O_T)=0$ by Serre duality. The low-degree Leray sequence,
using $\rho_{N*}\mathcal O_T=\mathcal O_N$ and the
finite support of $R^1\rho_{N*}\mathcal O_T$, and the vanishing
$H^2(N,\mathcal O_N)=0$ proved above, gives
\[
 0\longrightarrow H^1(N,\mathcal O_N)
 \longrightarrow H^1(T,\mathcal O_T)
 \longrightarrow H^0(N,R^1\rho_{N*}\mathcal O_T)
 \longrightarrow0.
\]
This proves the formula for $p$. If $T$ is rational,
$h^1(\mathcal O_T)=0$, so $p=0$. Each nonnegative local
geometric genus is zero, giving the asserted rational
singularities and Euler characteristic.

\item Put $L=\nu^*\omega_S$,
and write $I_N=\omega_N\otimes L^{-1}$ for the conductor
on $N$. Serre duality gives $\chi(I_N)=\chi(L)$.
Riemann--Roch on $T$, applied to the torsion line bundle
$\rho_N^*L$, gives
$\chi(\rho_N^*L)=\chi(\mathcal O_T)$.
The projection formula makes
$R^1\rho_{N*}\rho_N^*L=L\otimes R^1\rho_{N*}\mathcal O_T$;
its finite length is still $p$. The resolution fibres have dimension
at most one, so there are no higher direct images in degree two.
Leray therefore gives
$\chi(L)=\chi(\mathcal O_N)$.
Consequently the conductor scheme $D_{\rm sch}\subset N$
has $\chi(\mathcal O_{D_{\rm sch}})=0$.
The conductor square has the exact sequence
\[
 0\longrightarrow\mathcal O_S\longrightarrow
 \nu_*\mathcal O_N\oplus\mathcal O_{C_{\rm sch}}
 \longrightarrow\nu_*\mathcal O_{D_{\rm sch}}\longrightarrow0,
\]
where the first map is diagonal and the second takes the difference
of restrictions. Together with $\chi(\mathcal O_S)=0$, it gives
$\chi(\mathcal O_{C_{\rm sch}})=-\chi(\mathcal O_N)$.
Since $N$ is compact connected and $H^2(\mathcal O_N)=0$,
the remaining equality follows.

\item Finally suppose the ruled pole support is $H$ with
rational vertical trees. If $H$ is contracted, the
composition $H\to T\to N$ is constant. Pullback from
$H^1(N,\mathcal O_N)$ to $H^1(H,\mathcal O_H)$ is zero.
But restriction
$H^1(T,\mathcal O_T)\to H^1(H,\mathcal O_H)$ is an
isomorphism: the ruling has section $H$, a
$\mathbb P^1$-bundle has zero first higher direct image
of its structure sheaf, and the point blowups
preserve these groups. The Leray injection above
therefore forces $H^1(N,\mathcal O_N)=0$, so $p=g$.
If $H$ is not contracted, \cref{night:P4R.6} puts every
exceptional curve among the rational vertical components. Write
$a:T\to E$ for the ruling and $h:E\simeq H\hookrightarrow T$
for its section. Every exceptional curve maps constantly to $E$,
and every fibre of $\rho_N$ is connected. Thus $a$ is constant
on the fibres of $\rho_N$. Properness and
$\rho_{N*}\mathcal O_T=\mathcal O_N$ descend $a$ to a holomorphic
map $q:N\to E$ with $q\circ\rho_N=a$. The map
$s=\rho_N\circ h$ satisfies $q\circ s=\operatorname{id}_E$.
Consequently $s^*q^*=\operatorname{id}$ on $H^1(E,\mathcal O_E)$,
so $h^1(N,\mathcal O_N)\ge g$. The Leray injection into
$H^1(T,\mathcal O_T)$ gives the opposite inequality. Therefore
$h^1(N,\mathcal O_N)=g$ and $p=0$. The sheaf $R^1\rho_{N*}\mathcal O_T$ has finite support, so its zero space of global sections makes the sheaf itself zero.

\end{enumerate}
\end{proof}



\subsection{Unibranch fibre supports}

\begin{proposition}[Euler number for a bijective normalization]\hypertarget{tag-02YD}{}\label{night:P4R.2}\label{stmt:P4R.2}
\StatementRecord{night:P4R.2}{proved}{02YD}{\SeriesRef{night:P4.12}, \SeriesRef{night:P4.8}, \SeriesRef{night:P4R.1}, \SeriesRef{lem:II:negative-kodaira-models}}

In the standing torus-fibration setting of this section, every irreducible singular support
whose normalization is smooth and whose normalization morphism
is bijective has nonnegative topological Euler number.
Its Euler number is zero precisely when that normalization is
a minimal ruled surface over an elliptic curve or a Hopf surface.
\end{proposition}
\begin{proof}

\begin{enumerate}[label=\arabic*.,leftmargin=*]
\item Let $\nu:T\to S$ denote this finite normalization morphism.
A continuous bijection from the compact $T$ to the Hausdorff
$S$ is a homeomorphism. It therefore induces an isomorphism
of fundamental groups and equality of Euler numbers.
\SeriesRef{night:P4R.1} makes $\pi_1(T)$ virtually abelian.

\item The effective-conductor and canonical-torsion argument of
\SeriesRef{night:P4.12} applies to this same $T$ and proves
$\kappa(T)=-\infty$. Write $q:T\to T_0$ for the contraction
to a minimal model, with $k$ point blowups. Pluricanonical sections descend across these point blowups by
Hartogs' extension away from their centres, so $\kappa(T_0)=-\infty$.
Lemma~\ref{lem:II:negative-kodaira-models} therefore applies to
the smooth compact minimal $T_0$: its K\"ahler alternatives are
$\mathbb P^2$ and projective-line bundles, and its non-K\"ahler
alternative is class VII with $b_1(T_0)=1$.
Blowing up a smooth point preserves the fundamental group:
van Kampen compares a four-ball with the disk bundle
of $\OO_{\mathbb P^1}(-1)$, both attached along the same
simply connected boundary $S^3$. For a ruled minimal model
$p:T_0\to C$ over a genus-$g$ curve, the homotopy sequence
of its $\mathbb P^1$-bundle gives the isomorphism $p_*$
on fundamental groups. The resulting isomorphism from
$\pi_1(S)$ to $\pi_1(C)$ is $p_*q_*(\nu_*)^{-1}$.
For $g\ge2$ this group is not virtually abelian. Indeed a
finite-index subgroup corresponds to an unramified cover
of genus $1+d(g-1)>1$, and a genus-$h$ surface group maps
onto the nonabelian free group on $h$ generators by sending
the $b$ generators to the identity in its usual presentation.
Consequently $g\le1$.

\item A point blowup replaces a point by $\mathbb P^1$ and
increases Euler characteristic by one. The cell counts of
$\mathbb P^2$ and a $\mathbb P^1$-bundle over a genus-$g$
curve give $3$ and $4-4g$, respectively. For class VII,
Poincar\'e duality and $b_1(T_0)=1$ give $e(T_0)=b_2(T_0)$.
Thus
\[
 e(T)=\begin{cases}
 3+k,&T_0=\mathbb P^2,\\
 4-4g+k,&T_0\text{ ruled over a genus-}g\text{ curve},\\
 b_2(T_0)+k,&T_0\text{ of class VII}.
 \end{cases}
\]
Since $g\le1$, these formulas give
$e(S)=e(T)\ge0$. Equality in the ruled case is exactly
$g=1$ and $k=0$. Equality in class VII gives a minimal surface with
$b_2=0$. The conductor is nonempty because $S$ is singular
and nonnormal; \SeriesRef{cdp:st:A.5.3} gives its effective
conductor divisor on $T$. The reduced support of this effective divisor is a nonempty
curve. Lemma~\ref{lem:II:negative-kodaira-models}\textup{(iii)}
applies to the resulting minimal class-VII surface $T$ with
$b_2(T)=0$ and proves that $T$ is Hopf. The converse follows from the same Euler
formulas and the homeomorphism $\nu$.

\end{enumerate}
\end{proof}


\section{Double curves and normalization models}\label{sec:II:double-curve-normalizations}

\subsection{Ordinary double curves and normal surface alternatives}

\begin{proposition}[A ramified conductor with vanishing nearby homology]\hypertarget{tag-02WQ}{}\label{night:P4.19}\label{stmt:P4.19}
\StatementRecord{night:P4.19}{proved}{02WQ}{}

For the reduced hypersurface germ $y(x^2-zy^2)=0$, the normalization
conductor has multiplicity two on each component and its reduced
map to the singular locus is ramified. The affine nearby fibre is
$\mathbb C^*\times\mathbb C$. Every compact spherical class
in an open subset of this affine fibre has self-intersection zero.
\end{proposition}
\begin{proof}

\begin{enumerate}[label=\arabic*.,leftmargin=*]
\item One normalization component is the plane $y=0$; the other is
the Whitney normalization $(u,v)\mapsto(uv,v,u^2)$.
The singular locus is the smooth $z$-axis. The plane conductor
is $x^2=0$, obtained by restricting the other component's
equation. Write $R_B=\mathbb C\{v,uv,u^2\}$ and
$N_B=\mathbb C\{u,v\}$. The module $N_B$ is generated over
$R_B$ by $1,u$. A convergent series lies in $R_B$ precisely when
its odd powers of $u$ have coefficients divisible by $v$.
Thus $h,uh\in R_B$ is equivalent to $h\in vN_B$, proving
that the intrinsic conductor is $vN_B$.
The union ring is the fibre product of the component rings over
$\mathbb C\{x,z\}/(x^2)$. On the Whitney component, a
conductor element $h$ must satisfy $hN_B\subset vR_B$; the
intrinsic conductor calculation gives $h\in v^2N_B$.
On the plane, the corresponding kernel is $(x^2)$.
These are the two claimed conductor ideals.
The reduced conductor cover is the identity sheet from the
plane together with $u\mapsto u^2$ from the Whitney component.
The latter is ramified.

\item Put $F=\{y(x^2-zy^2)=1\}$. Here $y\ne0$, and the inverse maps
\[
 (x,y,z)\longmapsto(y,x),\qquad
 (y,x)\longmapsto\left(x,y,\frac{yx^2-1}{y^3}\right)
\]
identify this fibre with $\mathbb C^*\times\mathbb C$.
The retraction $(y,x)\mapsto(y/|y|,0)$ is a homotopy equivalence
with $S^1$, so $\pi_2$ and $H_2$ vanish. For an open inclusion
$j:U\hookrightarrow F$, compact-support duality identifies the
homology map $j_*$ with extension by zero $j_!$ on dual classes.
For the compact-support dual $\alpha$ of a compact spherical class,
\[
 \int_U\alpha^2=\int_F(j_!\alpha)^2=0.
\]
This proves the asserted zero self-intersection in $U$.

\end{enumerate}
\end{proof}

\begin{proposition}[Ordinary double strata and normal crossings]\hypertarget{tag-02WS}{}\label{night:P4.20}\label{stmt:P4.20}
\StatementRecord{night:P4.20}{proved}{02WS}{\SeriesRef{night:P4.16}, \SeriesRef{night:P4.6}, \SeriesRef{night:P4.8}}

Let $mS$ be a fibre with irreducible reduced support $S$ in a proper
holomorphic degeneration of complex two-tori with smooth total threefold.
Suppose the normalization of $S$ is smooth and, outside
finitely many points, $S$ is smooth or has the ordinary-double
germ $xy=0$. Then $S$ has normal crossings everywhere.
Propositions~\ref{night:P4.6}--\ref{night:P4.8} apply to
fibrations in the standing compact setting whose reduced supports
all satisfy these local hypotheses or already have normal crossings.
\end{proposition}
\begin{proof}

\begin{enumerate}[label=\arabic*.,leftmargin=*]
\item Take a small ball $B$ around an exceptional point, with
boundary transverse to the smooth and double strata away
from that point. Write the local defining function as $a g_0^m$,
where $g_0$ is reduced and $a$ is a unit. Choose a holomorphic
root $h^m=a$ and set $g=h g_0$. The local defining function
is then $g^m$.

\item Consider a small holomorphic deformation $G_s$ of $g$ whose
zero set near $\partial B$ retains its smoothly moving smooth
and ordinary-double strata, transverse to $\partial B$.
There is a uniform $\delta>0$ such that
$G_s|_{\partial B}$ is submersive over $0<|t|<\delta$.
At a boundary point on a double curve, choose coordinates
$a,b,c$ with $G_s=ab$, absorbing a unit in $a$.
Transversality of the double curve to $\partial B$ allows one
real $c$-coordinate to be solved from the boundary equation.
The complex coordinates $a,b$ remain free there, so $ab$ is
submersive when nonzero. At a smooth point use $G_s=a$.
Compactness gives uniform bounds; points outside these
neighbourhoods stay away from zero.

\item For each fixed small $s$ there are finitely many nonzero
critical values in this disc. After shrinking, the analytic
critical set in a neighbourhood of the compact closed ball
has finitely many relevant local components. The function is constant on
each analytic critical component. A component at a nonzero
small value cannot meet $\partial B$, so it is compact;
the maximum principle excludes a positive-dimensional
compact analytic subset of a ball. Thus the relevant critical
points are isolated and finite. Paths avoiding these values
join a small regular value to one fixed regular value $t_0$.
Proper Ehresmann transport on the closed ball, with the
boundary submersion from the preceding step, identifies the fibres along
such paths. For sufficiently small deformation the fibres
$G_s=t_0$ also form a smooth proper family to $s=0$, by
compactness. Hence every compact spherical class of nonzero
square in a local nearby fibre of $G_s$ transports into
$g=t_0$. Its inclusion in the compact torus fibre is impossible.

\item At a normalization preimage write the map germ as
$\phi:(\mathbb C^2,0)\to(\mathbb C^3,0)$.
Away from the exceptional point it is immersive with transverse
multiple images. The isolated-instability criterion in
\cite[p.~5]{GCP2024} makes this monogerm finitely determined.
Write $I_2(d\phi)$ for the ideal of its maximal minors.
A rank drop gives a nonempty finite scheme, and
\cite[Lemma~2.3, p.~5]{GCP2024} identifies its positive length
\[
 C(\phi)=\dim_{\mathbb C}\mathbb C\{u,v\}/I_2(d\phi)
\]
with the number of cross-caps in a stable deformation.
Apply this formula separately at each normalization preimage.
Generic polynomial perturbations vary the branch jets independently.
The stable multigerms listed in \cite[p.~4]{GCP2024} place each
cross-cap on a single sheet. Concretely, the incidence of its
zero-dimensional critical set with another two-dimensional image
in the three-dimensional target has negative expected dimension;
generic relative perturbations avoid it.

\item For the cross-cap coordinates, choose the first target
coordinate at a transverse rank-one point as a source
coordinate $u$, and write $\phi=(u,F(u,v),H(u,v))$.
Transversality of the equations $F_v=H_v=0$ gives a linear
combination with nonzero second $v$-derivative. The
holomorphic Morse lemma with parameter $u$, followed by
subtracting the critical value, makes the second coordinate
$v^2$. Subtract the even part of the third coordinate to
write it as $vA(u,v^2)$. Transversality gives $A_u(0,0)\ne0$.
The target change $u\mapsto A(u,Y)$ and the matching source
change give exactly $(u,v)\mapsto(u,v^2,uv)$.
The target coordinate permutation $(x,y,z)=(uv,u,v^2)$ gives
$x^2-zy^2=0$. Proposition~\ref{night:P4.16}, with $a=2$ and $n=2$,
gives a compact spherical class of square $-4$ in its nearby fibre.

\item The deformed images have holomorphic equations $G_s$ with
$G_0=g$. To see this, a generic projection of the original
finite normalization to two target coordinates is finite;
Weierstrass preparation preserves this after shrinking the
deformation. The total image is an analytic hypersurface
in $\mathbb C^3\times\Delta$. Its central equation has
multiplicity one, as seen at a smooth point of $S$; a unit
adjustment gives $G_0=g$. Outside a smaller ball the original
map is a transverse immersion, and this compact boundary
data persist under small deformation. We therefore have
the boundary property above; compare the stable-image discussion
in \cite[\S5.3, p.~33]{GCP2024}. Transporting the
crosscap class contradicts its torus inclusion. The
normalization is consequently immersive everywhere.

\item Write its image germ as a unit times $g_1\cdots g_r$, where
each branch is smooth and $v_i=dg_i(0)\ne0$. Absorb the
unit in one factor. Independent gradients give $r\leq3$
and normal-crossing coordinates by the inverse function
theorem. Suppose instead $\sum_i c_iv_i=0$ is a nonzero
relation, and let $I$ be its nonzero support. Set
$a_i=\varepsilon/c_i$ for $i\in I$ and $a_j=\eta$ for
$j\notin I$, where $0<|\varepsilon|\ll|\eta|\ll1$.
Perturb constants so that $h_i(0)=a_i$. Keep all gradients
except one $i_0\in I$, and change that gradient by
\[
 \Delta v_{i_0}=-a_{i_0}\sum_{j\notin I}\frac{v_j}{a_j}.
\]
This is arbitrarily small and makes $\sum_i dh_i(0)/a_i=0$.
Thus $H=\prod_i h_i$ is critical at the origin and has
nonzero critical value $\prod_i a_i$. Its Hessian divided
by that value is
\[
 \sum_i\frac{\operatorname{Hess}(h_i)}{a_i}
 -\sum_i\frac{dh_i\otimes dh_i}{a_i^2}.
\]
An arbitrarily small generic quadratic change to $h_{i_0}$
makes this matrix invertible: its determinant is a nonzero
polynomial in the freely variable symmetric Hessian.
The branch functions remain smooth and retain their
ordinary-double/smooth boundary strata. We have constructed
a holomorphic Morse point off the divisor.

\item The linear coordinates
$x=w_1$, $y=w_2+\mathrm i w_3$ and $z=-w_2+\mathrm i w_3$
identify its Morse quadratic with $x^2-zy$.
By \cref{night:P4.16}, with $a=2$ and $n=1$, the nearby
fibre has a compact spherical class of square $-2$.
The boundary transport in the third item puts this class
in the original torus fibre, where its square must be zero.
Hence the original branch gradients are independent
everywhere, proving the normal-crossing assertion.

\end{enumerate}
\end{proof}


\subsection{Cuspidal transverse type}

\begin{lemma}[Boundary transport for a holomorphic deformation]\hypertarget{tag-P295}{}\label{lem:II:boundary-deformation-transport}
\StatementRecord{lem:II:boundary-deformation-transport}{proved}{P295}{\cite[\S3 and \S3.1, arXiv v1, pp.~3--4]{SiersmaTibar2017Deformation}}
Let $g_\tau$ be a holomorphic family of functions on a neighbourhood of a closed ball $B\subset\C^3$, for parameters $\tau$ near zero, and suppose that every critical value of $g_0$ in $B$ is zero. Choose the ball transverse to the smooth strata of $g_0^{-1}(0)$. Assume the following product data near the critical set on its boundary. Finitely many coordinate families, depending smoothly on $\tau$ and holomorphically on the source variables, identify $g_\tau$ with a fixed holomorphic function $q(w)$ independent of at least one complex coordinate $z$. Each $q$ has only the critical value zero in the chosen chart, and some real direction in the $z$ coordinates is transverse to $\partial B$. These charts cover all boundary critical points for small $\tau$; on the rest of the zero-level boundary the restriction of $g_0$ to $\partial B$ is a submersion.

After shrinking the parameter neighbourhood and target disc, the deformation is trivial on the boundary in the sense of \cite[\S3, p.~3]{SiersmaTibar2017Deformation}. Each $g_\tau$ has only finitely many critical values in $B$, and there is a common nonzero regular value $b$ for which parameter transport identifies
\[
 B\cap g_\tau^{-1}(b)\simeq B\cap g_0^{-1}(b)
\]
by an orientation-preserving diffeomorphism. A compact class in a local regular fibre at an interior point of $B$, with value in the chosen target disc, can be transported locally to a value regular for the whole ball, included in that fibre in $B$, transported along a path avoiding its critical values to the $b$-fibre, and then transported to $g_0^{-1}(b)$. These maps preserve the intersection form on compact two-dimensional classes. A class represented by a map from $S^2$ is sent to the class of the composite map from $S^2$. The construction composes over nested balls and successive small deformations.

For a local expression $\xi\circ F=g_0^m$ of a map to a smooth curve, its final inclusion is into $F^{-1}(b^m)$. When that fibre is a complex two-torus, every transported spherical class has square zero.
\end{lemma}
\begin{proof}
In a product chart, a real $z$-direction gives a field $W$ with $dg_\tau(W)=0$ and $d\rho(W)\ne0$ near the relevant boundary point, where $\rho$ defines the sphere. Holding the product coordinates fixed while changing a real parameter gives a field $V$ satisfying
\[
 dg_\tau(V)+\partial_\tau g_\tau=0.
\]
Replace it by $V-d\rho(V)W/d\rho(W)$. This still solves the displayed equation and is tangent to the fixed boundary. Away from the critical set, boundary submersivity supplies a solution of the same linear equations. A partition of unity preserves the equations and boundary tangency. Its flow on the compact boundary gives a trivialization preserving the function and its critical strata. The same product directions show uniform transversality of small nonzero levels near the critical boundary; openness and compactness give it on the remaining boundary. This verifies the admissibility hypothesis in \cite[\S3, p.~3]{SiersmaTibar2017Deformation} for the specified maps.

A positive-dimensional critical component meeting the boundary has value zero by the product charts. Any other positive-dimensional component would be a compact positive-dimensional analytic subset of the interior of a coordinate ball, which the maximum principle excludes. The remaining critical components are isolated, and compactness and local finiteness of analytic components make them finite. All their values tend uniformly to zero as $\tau\to0$: otherwise a convergent sequence of critical points would limit to a critical point of $g_0$ with a nonzero value. Choose $b$ small and then the parameters smaller so that every critical value has modulus less than $|b|/2$. The family of $b$-levels is a proper submersion with transverse boundary, so boundary-tangent horizontal fields give the parameter diffeomorphism.

For a compact class on a local regular fibre, smooth local transport near its compact representative permits an arbitrarily small change of target value. Choose the resulting value $b_\tau$ regular also for the whole ball. The excluded set of values is finite, so a path $\gamma$ in the target disc joins $b_\tau$ to $b$ through regular values. Lifting $\dot\gamma$ by boundary-tangent solutions of $dg_\tau(V)=\dot\gamma$ gives the target transport. Thus the maps are the local transport and inclusion, this target transport, and the inverse parameter transport at $b$, in that order. Each diffeomorphism preserves the complex orientation; the inclusions and transports preserve intersection numbers computed with compact representatives in the interiors. Repeating these maps first in an inner ball and then in an outer ball proves the nested assertion. Finally $g_0=b$ lies in $\xi\circ F=b^m$. A sphere in a two-torus lifts to its contractible universal cover, so its homology class and square are zero.
\end{proof}

\begin{lemma}[Two hypersurfaces meeting along a smooth curve]\hypertarget{tag-P299}{}\label{lem:II:smooth-curve-intersection}
\StatementRecord{lem:II:smooth-curve-intersection}{proved}{P299}{}
Let two reduced hypersurface germs $\phi=0$ and $\psi=0$ in a smooth complex threefold have common zero set equal to the smooth curve $\Gamma=\{x=y=0\}$. Let $z$ be its coordinate and let $r$ be their intersection multiplicity on a generic transverse plane. Then
\[
 A=\C\{x,y,z\}/(\phi,\psi)
\]
is a finite free $\C\{z\}$-module of rank $r$, and its fibre at $z=0$ has complex dimension $r$. For $r=1$, $(\phi,\psi,z)$ is a holomorphic coordinate system. For $r=2$, one factor is smooth; after interchanging them if necessary, there are holomorphic coordinates $(X,Y,z)$ in which
\[
                     \phi=X,\qquad \psi=Y^2+XH(X,z).
\]

For $r=2$ these coordinates give the algebra isomorphism
\begin{equation}
 A\simeq\C\{z,\epsilon\}/(\epsilon^2),\qquad
 \epsilon\longmapsto[Y].
\label{eq:II:manual:4.56a}
\end{equation}
When $\psi$ has generic normal order two along $\Gamma$,
the smooth factor is $\phi$, and the same coordinates give
\begin{equation}
             \phi=X,\qquad\psi=Y^2+X^2k(X,z).
\label{eq:II:manual:4.56c}
\end{equation}
Suppose in addition that $\psi$ is irreducible and
$k(0,z)\not\equiv0$. Its finite birational partial
normalization is
\begin{equation}
 N'=\{u^2+k(X,z)=0\},\qquad Y=Xu.
\label{eq:II:manual:4.56d}
\end{equation}
Write $k=a^2h$, with $h$ reduced and with a unit absorbed
into $a$ by a holomorphic square root. Its normalization is
\[
 \widehat N=\{v^2+h(X,z)=0\},\qquad
 u=av,\quad Y=Xav.
\]
The inverse image of $(X,Y)$ is $(X)$, and its divisor
is reduced. At a point where this reduced curve is smooth
and its degree-two map to $\Gamma$ is ramified,
$h_z\ne0$ and $\widehat N$ is smooth. In particular,
$k_z(0,0)\ne0$ makes $N'$ itself the smooth normalization.
\end{lemma}
\begin{proof}
The zero fibre of $A$ over $z=0$ is isolated, so the analytic map to the $z$-disc is finite after taking a sufficiently small representative. The two equations form a regular sequence in the regular ambient local ring. Their quotient is therefore a one-dimensional complete intersection, and its parameter $z$ is a nonzerodivisor. A finite torsionfree module over the discrete valuation ring $\C\{z\}$ is free. Its rank is the generic transverse intersection multiplicity $r$, proving the fibre-length assertion.

For $r=1$, the quotient $\C\{x,y\}/(\phi(x,y,0),\psi(x,y,0))$ has length one, so its ideal is $(x,y)$. The two linear parts are independent. Thus $(x,y,z)\mapsto(\phi,\psi,z)$ has invertible Jacobian and its inverse is the unique local solution of these coordinate equations.

For $r=2$, the two generators of this plane ideal cannot both lie in $(x,y)^2$, since the quotient would retain the three independent classes $1,x,y$. Thus one factor, say $\phi$, is smooth. Set $X=\phi$ and choose a complementary normal coordinate $y$, keeping $z$. On $X=0$ the zero set of $\psi$ is exactly $y=0$; its generic multiplicity is two, so $\psi|_{X=0}$ is a unit times $y^2$. The parameter Morse construction in $y$, keeping $X,z$, gives $Y=(y-\mu(X,z))\sqrt{V(X,y,z)}$, with $V$ a unit and with inverse the unique small solution for $y$. Its centre satisfies $\mu(0,z)=0$, and its remaining critical value vanishes on $X=0$. It is therefore $XH(X,z)$, which proves the exact displayed form.

The quotient by $(X,\psi)$ in the displayed coordinates
is exactly $\C\{Y,z\}/(Y^2)$, proving
\eqref{eq:II:manual:4.56a}. Generic normal order two
means $\psi\in I_\Gamma^{(2)}=I_\Gamma^2$: its normal
Taylor coefficients of degree at most one vanish
identically along the smooth curve. Thus $d\psi=0$
on $\Gamma$, so the smooth factor is $\phi$. In the
preceding Morse coordinates, $Y^2+XH\in(X,Y)^2$
forces $H\in(X)$, giving \eqref{eq:II:manual:4.56c}.

The element $u=Y/X$ satisfies the monic relation
$u^2+k=0$, and $N'\to\{\psi=0\}$ is the identity
on the remaining coordinates, with inverse $u=Y/X$
off $X=0$. This proves finiteness and birationality in
\eqref{eq:II:manual:4.56d}. Likewise $v=Y/(Xa)$ is
integral, with relation $v^2+h=0$. The algebra of
$\widehat N$ is generated by this $v$ over that of
$\{\psi=0\}$, so its map is finite and has the same
fraction field. Substitution gives the exact identity
$\psi(X,Xav,z)=X^2a^2(v^2+h)$.
Here $h$ is a nonunit: a unit $h$ would have a
holomorphic square root and would factor $\psi$ as
$(Y+\mathrm i Xa\sqrt h)(Y-\mathrm i Xa\sqrt h)$,
contrary to its irreducibility.
Since the plane curve $h=0$ is reduced, its singular
locus is isolated. The singular locus of $v^2+h=0$
is contained in $v=0$, $dh=0$. The hypersurface is
Cohen--Macaulay and regular in codimension one, hence
normal. The finite birational map is the normalization.

Its inverse-image ideal is
$(X,Y)\OO_{\widehat N}=(X)$. The quotient
$\C\{v,z\}/(v^2+h(0,z))$ is a one-dimensional
complete intersection, generically reduced because
$h(0,z)\not\equiv0$. It has no embedded associated
prime and is therefore reduced. A unit $h(0,z)$ gives
two unramified points; a positive even order gives
two curve branches, and an odd order at least three
gives a singular curve. Smoothness and ramification
therefore force a simple zero of $h(0,z)$, so $h_z\ne0$
and $\widehat N$ is smooth. Finally $k_z(0,0)\ne0$
makes $N'$ smooth directly; its already proved finite
birational map is then the normalization.
\end{proof}


\begin{proposition}[Contact of two smooth branches]\hypertarget{tag-P298}{}\label{prop:II:smooth-branch-contact}
\StatementRecord{prop:II:smooth-branch-contact}{proved}{P298}{\cref{prop:II:weighted-two-branch-sphere,cdp:st:A.5.17,lem:II:boundary-deformation-transport}}
Let $F$ be a holomorphic map from a smooth complex threefold to a smooth curve, with nearby smooth fibres complex two-tori. Suppose its full local expression is
\[
 F(X,Y,Z)=U(X,Y,Z)(X-H(Y,Z))^a(X+H(Y,Z))^b,
 \qquad a,b\ge1,
\]
where $U$ is a holomorphic unit and $H\ne0$ vanishes at the origin. Then $\{H=0\}_{\mathrm{red}}$ is smooth. If $k$ is its multiplicity in $H$, an invertible holomorphic coordinate change puts the full function in the form
\[
                   (X-Y^k)^a(X+Y^k)^b.
\]
Every reduced local generator of its support has the exact product form $X^2-Y^{2k}$ in holomorphic coordinates.
\end{proposition}
\begin{proof}
\begin{enumerate}[label=\arabic*.,leftmargin=*]
\item Factor $H=\nu\prod_i h_i^{k_i}$ in $\C\{Y,Z\}$, with $\nu$ a unit and the $h_i$ distinct reduced irreducible factors. We use holomorphic deformations $H_s$ which retain these exponents near the boundary, and keep the full family
\[
          F_s=U(X-H_s)^a(X+H_s)^b
\]
with the original $U,a,b$. Choose a small closed ball transverse to the smooth strata of $F_0^{-1}(0)$, meeting the plane curves $h_i=0$ in their smooth disjoint links. Every critical value of $F_0$ in this ball is zero after shrinking: its critical germ has finitely many components, and the function is constant zero on each component through the origin.

We verify the boundary data of \cref{lem:II:boundary-deformation-transport} for this full family. Off the intersection of the two branches, a normal coordinate and a root of the unit give $w^a$ or $w^b$. On the intersection link belonging to $h_i$, write $H_s=\nu_s h_{i,s}^{k_i}$ with $\nu_s$ a unit, and take the coordinate $y=\nu_s^{1/k_i}h_{i,s}$, with a complementary coordinate $z$. The inverse is the unique local solution of these two coordinate equations; their Jacobian is nonzero on the link. Thus
\[
 F_s=U_s(X-y^{k_i})^a(X+y^{k_i})^b.
\]
The polynomial on the right has weights $(k_i,1)$ in $(X,y)$ and degree $d_i=k_i(a+b)$. Multiplying $X$ by $U_s^{1/(a+b)}$ and $y$ by $U_s^{1/d_i}$, keeping $z$, removes the unit exactly. Its inverse is
\[
 (X',y',z)\longmapsto(\tau^{-k_i}X',\tau^{-1}y',z),\qquad
 \tau^{d_i}=U_s(\tau^{-k_i}X',\tau^{-1}y',z),
\]
with the selected nonzero root at the link. The implicit derivative is $d_i\tau^{d_i-1}\ne0$ there. These charts depend holomorphically on $s$ and on the source. Along the link their $z$ direction is tangent to the same smooth curve, which is transverse to the boundary; this transversality persists nearby. The weighted Euler identity shows that the fixed transverse polynomial has only the critical value zero. The smooth-branch charts give the same boundary property when $a>1$ or $b>1$. A finite collection covers the boundary critical set. Hence \cref{lem:II:boundary-deformation-transport}, with $g_s=F_s$ and $m=1$, carries every compact local spherical class constructed below into a fibre of the original $F$, preserving its square.

\item Suppose at least two factors occur. Apply the finite local-degree construction in the proof of \cref{cdp:st:A.5.17} to $(h_i,h_j)$. A small regular value $(s_i,s_j)$ has a nonempty finite inverse image, and the functions $h_i-s_i,h_j-s_j$ are transverse coordinates there. Choose the remaining constants to avoid their zero sets at one selected point. The deformation
\[
                H_s=\nu\prod_i(h_i-s_i)^{k_i}
\]
has there the form $\mu v^{k_i}w^{k_j}$, with $\mu$ a unit. Replacing $v$ by $\mu^{1/k_i}v$, keeping $w$, is invertible and gives $H_s=v^{k_i}w^{k_j}$ exactly. The full $F_s$ is therefore the function of \cref{prop:II:weighted-two-branch-sphere}, with the same $a,b$. Its compact sphere has nonzero square, and the preceding full-family transport gives a contradiction. There is only one factor. Absorb its plane unit and write $H=h^k$, with $h$ reduced and irreducible.

\item We show that $h$ is smooth, retaining the unit $U$ in the full function. Put $N=a+b$ and introduce an independent variable $q$ for a value of $H$. Write $L_X=\partial_X\log U$. The equation
\[
 NC+(a-b)+q(C^2-1)L_X(qC,Y,Z)=0,
 \qquad C(0,Y,Z)=c_0=\frac{b-a}{N},
\]
has a unique holomorphic solution $C=C(q,Y,Z)$ near the origin, since its derivative in $C$ at $q=0$ is $N$. Both $c_0-1=-2a/N$ and $c_0+1=2b/N$ are nonzero. For $q\ne0$, the point $X=qC$ is precisely the nearby solution of
\[
 \partial_X\log\bigl(U(X,Y,Z)(X-q)^a(X+q)^b\bigr)=0
\]
off the two zero factors. Define the holomorphic unit and coordinate
\[
 W(q,Y,Z)=U(qC,Y,Z)(C-1)^a(C+1)^b,\qquad
 \Lambda(q,Y,Z)=qW(q,Y,Z)^{1/N}.
\]
The derivative $\partial_q\Lambda$ at $q=0$ is the nonzero chosen root of $W(0,Y,Z)$. Thus $v=\Lambda(q,Y,Z)$ has a holomorphic inverse $q=R(v,Y,Z)$, characterized by
\[
 \Lambda(R(v,Y,Z),Y,Z)=v,\qquad
 R(\Lambda(q,Y,Z),Y,Z)=q.
\]
In particular $R(v,Y,Z)=vR_1(v,Y,Z)$ with $R_1$ a unit. The value of the full function at its $X$-critical centre is exactly $\Lambda(q,Y,Z)^N$.

\item Put
\[
       \widetilde h=h\,W(H,Y,Z)^{1/(kN)},\qquad
       \Lambda(H,Y,Z)=\widetilde h^{\,k}.
\]
Choose the roots compatibly in these two equalities. The germ $\widetilde h$ is reduced, and it is singular precisely when $h$ is singular. For a singular $h$, apply the finite-gradient construction in the proof of \cref{cdp:st:A.5.17} to $\widetilde h$: small linear terms and a constant give
\[
 \widetilde h_s=\widetilde h-\lambda_1Y-\lambda_2Z-\lambda_0
\]
with an interior Morse critical point $p$, nonzero value $v=\widetilde h_s(p)$, and invertible Hessian $Q$. Define
\[
                 H_s=R(\widetilde h_s^{\,k},Y,Z).
\]
At $s=0$ this is $H$, by the two inverse identities. Since $R(v,Y,Z)=vR_1(v,Y,Z)$, this is a deformation of one reduced plane factor with its original exponent $k$. Its boundary charts are exactly those of the first step. Both analytic factors $X-H_s$ and $X+H_s$ and their full multiplicities remain in $F_s$.

Let $A_s=H_sC(H_s,Y,Z)$. The defining equation for $C$ gives $(F_s)_X(A_s,Y,Z)=0$, and the critical value function is the exact identity
\[
                   F_s(A_s,Y,Z)=\widetilde h_s^{\,kN}.
\]
At $(A_s(p),p)$ the full gradient therefore vanishes and the value is $v^{kN}\ne0$. Moreover
\[
 (F_s)_{XX}(A_s,Y,Z)
 =H_s^{N-2}W(H_s,Y,Z)
 \left(H_s^2\partial_X^2\log U(A_s,Y,Z)
              -\frac{a}{(C-1)^2}-\frac{b}{(C+1)^2}\right).
\]
The bracket has value $-N^3/(4ab)\ne0$ at the origin; at the selected point $H_s\ne0$, so $(F_s)_{XX}\ne0$. Differentiating $(F_s)_X(A_s,Y,Z)=0$ identifies the Schur complement of this entry in the full Hessian with the Hessian of $F_s(A_s,Y,Z)$. At $p$ it is $kN v^{kN-1}Q$. Consequently
\[
 \det\operatorname{Hess}F_s
    =(F_s)_{XX}(kN)^2v^{2kN-2}\det Q\ne0.
\]
The explicit Morse-coordinate and sphere construction in the proof of \cref{cdp:st:A.5.17} gives a local sphere of square $-2$. All parameters, its critical value and its support can be chosen inside the transport neighbourhood of the first step. Its image in the original torus fibre has the same square, a contradiction. Thus $h$ is smooth.

\item Take $Y=h$ and a complementary plane coordinate $Z$, so that $H=Y^k$ exactly. This coordinate map has nonzero Jacobian, and its inverse is the unique local solution of these two equations. The weighted unit map of the first step, now with $d=k(a+b)$, changes the full function to $(X-Y^k)^a(X+Y^k)^b$, with its stated inverse. Any reduced local generator in these coordinates is $V(X^2-Y^{2k})$ for a unit $V$. The further map
\[
             (X,Y,Z)\longmapsto(V^{1/2}X,V^{1/(2k)}Y,Z)
\]
changes that generator to $X^2-Y^{2k}$. Its inverse is the same displayed implicit map with weights $(k,1)$ and degree $2k$. This proves both asserted normal forms.
\end{enumerate}
\end{proof}

\begin{proposition}[Local rigidity of a transverse cusp]\hypertarget{tag-02X0}{}\label{night:P4.28}\label{stmt:P4.28}
\hypertarget{tag-01BT}{}\label{cdp:st:C.2.CR.4}
\StatementRecord{night:P4.28}{proved}{02X0}{\cref{lem:II:boundary-deformation-transport,night:P4.16,night:P4.18,cdp:st:A.5.15}}
\StatementRecord{cdp:st:C.2.CR.4}{proved}{01BT}{\cref{lem:II:boundary-deformation-transport,night:P4.16,night:P4.18,cdp:st:A.5.15}}

Let $p$ lie on a fibre of a holomorphic map $F$ from a smooth complex threefold to a smooth curve, with nearby smooth fibres complex two-tori. Suppose that the reduced singular locus of the fibre support is a smooth curve near $p$, and that its generic transverse plane-curve germ is an ordinary cusp. Every reduced local defining function $g$ then has holomorphic coordinates in which
\begin{equation}
                         g=X^2+Y^3.
\hypertarget{tag-00D1}{}\label{cdp:eq:C.CR.15}
\end{equation}
In particular the germ of the support is the product of that cusp with a smooth curve and its normalization is smooth. A local unit and the common full-fibre multiplicity are retained through the explicit root and coordinate maps in the proof.
\end{proposition}
\begin{proof}

\begin{enumerate}[label=\arabic*.,leftmargin=*]
\item The reduced support is locally irreducible. Otherwise
two local surface components would intersect in a curve
contained in the smooth reduced singular curve $C$,
giving at least two branches at its generic point.
Choose a base coordinate $\xi$ and write
$\xi\circ F=u g_{\rm red}^{\,m}$ locally.
A holomorphic $m$-th root of the unit $u$ gives the
reduced function $g=u^{1/m}g_{\rm red}$ with
$\xi\circ F=g^m$. Its nonzero local level $g=b$
therefore embeds in the original torus fibre
$\xi=b^m$.

Every component of the critical germ of $g$ has constant value, which is zero because it contains $p$. Thus, after shrinking, its full reduced critical locus is precisely the smooth singular curve $C$.

Choose a Milnor ball transverse to this curve. The deformations below retain transverse type $A_2$ near its boundary: the parameter Morse lemma and a cube root of the remaining unit give product charts with function $a^2+b^3$ and the third coordinate along the curve. These are the product data of \cref{lem:II:boundary-deformation-transport}. That lemma transports their compact local classes into the original level $g=b$, and the equality $\xi\circ F=g^m$ identifies its inclusion into the torus fibre with parameter $b^m$. Every transported spherical class consequently has square zero.
\item Choose coordinates $(x,y,z)$ with $C=\{x=y=0\}$.
The quadratic normal part of $g$ is
$a(z)x^2+2b(z)xy+c(z)y^2$, with generic rank one.
Thus $ac=b^2$ in the discrete valuation ring
$\mathbb C\{z\}$. Choose a nonzero diagonal entry
of smallest order, say $a$, interchanging $x,y$
when necessary. Such an entry exists. When $b=0$ the form is already
diagonal; otherwise
$2\operatorname{ord}b=\operatorname{ord}a+
\operatorname{ord}c$ makes $b/a$ holomorphic.
The invertible map
$(x,y,z)\mapsto(x+(b/a)y,y,z)$ completes the square
and writes the quadratic normal term as $h(z)x^2$.
The function $h$ is not identically zero.
\item First suppose the Hessian has rank one at $p$.
Solve $g_x=0$ as $x=\phi(y,z)$.
The zero $xy$ quadratic coefficient gives
$\phi=O(y^2)$. With $u=x-\phi$, Taylor division
writes $g=u^2A(u,y,z)+k(y,z)$, where $A$ is a unit.
The coordinate $X=u\sqrt A$ is locally invertible.
The remaining function satisfies
$k(0,z)=k_y(0,z)=k_{yy}(0,z)=0$, so
\[
 g=X^2+y^3h_1(y,z).
\]
Its cubic coefficient $h_1(0,z)$ is not identically
zero by generic transverse type $A_2$.
When $h_1(0,0)\ne0$, the invertible coordinate
$Y=yh_1(y,z)^{1/3}$ gives $g=X^2+Y^3$.
\item A zero of $h_1(0,z)$ at the origin has finite
positive order. The reduced plane zero set of $y^3h_1(y,z)$ then has a component through the origin besides $y=0$, since $h_1(0,z)$ is not identically zero, and is singular there. For the deformation
$g_\eta=X^2+y^3(h_1+\eta)$, a sufficiently small
generic $\eta$ gives a simple zero $z_\eta$ near
zero. One-variable preparation and avoidance of
the critical values give simplicity; the number
of zeros persists by Rouch\'e's theorem.
At $(0,0,z_\eta)$ the map
$(X,y,z)\mapsto(X,y,Z=h_1(y,z)+\eta)$
has nonzero Jacobian and gives $g_\eta=X^2+y^3Z$.
The coordinates $(x,y,z)=(X,y,-Z)$ identify this
function with $x^2-zy^3$. Proposition~\ref{night:P4.16},
with $a=2$ and $n=3$, gives a compact spherical class
of square $-6$ in its affine nearby fibre. Positive weighted scaling,
with weights $(2,1,1)$ and degree $4$, places
compact representatives in its local Milnor fibre.
On the fixed boundary, $h_1(0,z)$ stays nonzero
for small $\eta$, so the transverse type remains
$A_2$. The transport of item~1 carries this
nonzero square into the original torus, a contradiction.
\item Suppose that the Hessian has rank zero at $p$, so $h(0)=0$. Write $k(z)y^3$ for the cubic coefficient on the kernel of its quadratic normal form. Both $h$ and $k$ are nonzero generically. Consider
\begin{equation}
                  g_{\epsilon,\delta}
                      =g+\epsilon x^2+\delta y^3.
\hypertarget{tag-00D2}{}\label{cdp:eq:C.CR.16}
\end{equation}
Choose a small regular value $-\epsilon$ of $h$ with a simple preimage $z_\epsilon$ near zero, and choose $\delta$ with $k(z_\epsilon)+\delta\ne0$. Such choices exist by one-variable preparation and the density of regular values; only finitely many values of $\delta$ are excluded after fixing $\epsilon$. The quadratic normal coefficient is $h(z)+\epsilon$, and the kernel cubic is a unit at the selected point. On the fixed boundary both coefficients remain nonzero for small parameters, so its product-cusp charts persist.

\item Here is the common preparation map, also used for the second family below. For a function $G$ with the indicated critical axis, quadratic normal form $h_G(z)x^2$ and nonzero kernel cubic at the point, Weierstrass preparation in $y$ gives
\[
 G=UP,\qquad P=y^3+a(x,z)y^2+b(x,z)y+c(x,z),
\]
with $U$ a unit. At $x=0$, one has $P(0,y,z)=y^3$, so $a,b,c$ vanish there. The critical-curve equation gives $c_x(0,z)=0$. Differentiating $UP$ on that curve gives $G_{xy}=U b_x$; its zero mixed quadratic entry therefore gives $b_x(0,z)=0$. Thus $x\mid a$ and $x^2\mid b,c$. The invertible substitution $y=Y-a/3$ gives
\begin{equation}
 P=Y^3+x^2\bigl(\beta(x,z)Y+\gamma(x,z)\bigr).
\hypertarget{tag-00D3}{}\label{cdp:eq:C.CR.17}
\end{equation}
Indeed its two depressed coefficients are $b-a^2/3$ and $c-ab/3+2a^3/27$, both divisible by $x^2$. The function $\gamma(0,z)$ is the quadratic normal coefficient of $G$ divided by $U(0,0,z)$.

For $G=g_{\epsilon,\delta}$ at the selected point, $\gamma(0,z)$ has a simple zero. Hence
\[
 Z=\beta(x,z)Y+\gamma(x,z)
\]
is an invertible change of the curve coordinate and gives $G=U(Y^3+x^2Z)$. In these new coordinates the map
\[
 (\widehat x,\widehat Y,\widehat Z)=U^{1/3}(x,Y,Z)
\]
has invertible Jacobian and satisfies $\widehat Y^3+\widehat x^{\,2}\widehat Z=U(Y^3+x^2Z)$. Thus the exact function form, after renaming the coordinates, is
\begin{equation}
                              Y^3+x^2Z.
\hypertarget{tag-00D4}{}\label{cdp:eq:C.CR.18}
\end{equation}
The permutation and sign change $(x',y',z')=(Y,x,-Z)$ gives $(x')^3-(y')^2z'$. The compact spherical cycle of \cref{cdp:st:A.5.15,night:P4.18}, with exponents $(3,2,1)$, has square $-2(2/\gcd(2,1))(1/\gcd(2,1))=-4$. It is the class of the explicit continuous map from $S^2$; no embedded representative is required. Ordinary degree-three scaling places its compact representative in an arbitrarily small local nonzero level. The direct parameter path in \eqref{cdp:eq:C.CR.16} and \cref{lem:II:boundary-deformation-transport} carry that class into the original torus, a contradiction.

\item The prepared two-stage family uses the same calculation with its full coordinate map. First choose a small $\delta$ for which $g_{0,\delta}=g+\delta y^3$ has nonzero kernel cubic at the origin. Let $\Phi_\delta$ be the preceding Tschirnhaus coordinate change and write
\[
 g_{0,\delta}\circ\Phi_\delta^{-1}
       =U_\delta\{Y^3+x^2H_\delta(x,Y,z)\},\qquad
 H_\delta=\beta_\delta Y+\gamma_\delta.
\]
Here $H_\delta(0,0,z)=h(z)/U_\delta(0,0,z)$ is nonzero generically and vanishes at zero. Keeping $U_\delta$ fixed gives the holomorphic family
\[
 \widehat g_{\delta,\eta}
       =U_\delta\{Y^3+x^2(H_\delta+\eta)\}.
\]
The relation to the first family in the same coordinates is the exact identity
\[
 g_{\epsilon,\delta}\circ\Phi_\delta^{-1}
       =U_\delta\{Y^3+x^2(H_\delta+\epsilon/U_\delta)\}.
\]
At $\eta=0$ the prepared family is exactly the first-stage function through $\Phi_\delta$. For a small generic $\eta$, the function $H_\delta(0,0,z)+\eta$ has a simple zero. The coordinate $Z=H_\delta+\eta$ and the same simultaneous cube-root scaling give \eqref{cdp:eq:C.CR.18}, with the same compact spherical cycle.

These transports use nested balls. Fix the original ball $B_1$ for $g\to g+\delta y^3$, whose boundary retains transverse type $A_2$. There are only finitely many additional critical points by \cref{lem:II:boundary-deformation-transport}. Choose a smaller ball $B_2$ about the origin in the preparation chart, avoiding those points, with boundary transverse to the critical curve where $h(z)\ne0$. The second perturbation retains type $A_2$ there, and the lemma transports its compact class to a regular fibre of $g+\delta y^3$ inside $B_2$. Choose that value regular also for $B_1$. Inclusion and target transport in $B_1$, followed by the inverse first parameter transport, place it in $g=b$ and then in the torus $\xi=b^m$. These are the same interior inclusions and orientation-preserving maps, so the compact square remains $-4$.
\item Both a zero Hessian and a vanishing kernel cubic
coefficient have been excluded. The coordinates
of item~3 give the product cusp germ.
Its normalization is the finite map
\[
 (u,z)\longmapsto (a=u^3,b=-u^2,z),
\]
with smooth source. Off $b=0$ its rational inverse
is $u=-a/b$, and the map on the reduced conductor
is the identity on the $z$ coordinate. For any other reduced defining function $g'=v g$, the unit $v$ is retained through the invertible coordinate change $X'=v^{1/2}X$, $Y'=v^{1/3}Y$, with the curve coordinate fixed. It gives $g'=(X')^2+(Y')^3$ exactly. This also proves the function assertion in \eqref{cdp:eq:C.CR.15} for a reduced generator chosen before the full-map unit was absorbed.
\end{enumerate}
\end{proof}


\subsection{The conductor pushout for ruled normalizations of higher genus}



\begin{proposition}[Conductor topology over a higher-genus ruling]\hypertarget{tag-02YF}{}\label{night:P4R.3}\label{stmt:P4R.3}
\StatementRecord{night:P4R.3}{proved}{02YF}{\SeriesRef{cdp:st:A.5.3}, \SeriesRef{night:P4.12}, \SeriesRef{night:P4.13}, \SeriesRef{night:P4R.1}}

In the standing torus-fibration setting of this section, let an irreducible singular fibre
support $S$ have smooth normalization $\nu:T\to S$ whose
minimal model is ruled over a curve of genus $g\ge2$.
Write $D$ for the conductor divisor on $T$ and
$C=\operatorname{Sing}(S)_{\rm red}$.
The inclusion $D_{\rm red}\to T$ induces an isomorphism on
fundamental groups, and the normalization pushout induces
$\pi_1(S)\simeq\pi_1(C)$. Moreover,
\[
 e(C)\ge0,\qquad e(S)>0.
\]
The assertion includes conductor components on different
ruling fibres that map onto the same curve in $C$.
\end{proposition}
\begin{proof}

\begin{enumerate}[label=\arabic*.,leftmargin=*]
\item Proposition~\ref{cdp:st:A.5.3}, with the identity as the
minimal resolution of the smooth surface $T$, gives
$\omega_T(D)\simeq\nu^*\omega_S$ with $D$ effective and with its
conductor multiplicities. Proposition~\ref{night:P4.12} makes this
line torsion, so $K_T+D\equiv0$.
Let $\beta:T\to T_0$ contract the ruling blowups, and put $D_0=\beta_*D$. Apply \cref{lem:II:irrational-ruled-boundary} with these same $T,T_0,D$ and the ruling base of genus $g\ge2$. The lemma's adjunction and Riemann--Hurwitz inequalities exclude its reduced-horizontal case. Thus $D_0=2H+\sum a_jF_j$, and $D_{\rm red}$ is the section with rational trees. Its inclusion $j:D_{\rm red}\hookrightarrow T$ induces the precise fundamental-group isomorphism given there.

\item The normalization $\nu:T\to S$ is finite and is an isomorphism away from its conductor. Its nonisomorphism locus is exactly $C$: a normal point is identified with a neighbourhood in the smooth normalization, while every smooth point is normal. Thus $D_{\rm red}=\nu^{-1}(C)$ and the complementary open sets are identified. A triangulation of the compact analytic pair $(T,D_{\rm red})$ makes its inclusion a subcomplex \cite[Theorem~2, p.~464]{Lojasiewicz1964}. Apply \cref{lem:II:proper-conductor-pushout} with $(X,Y,A,C,h)=(T,S,D_{\rm red},C,\nu)$. The three spaces in its connected clause are connected, since $D_{\rm red}$ is the section with rational trees and maps onto $C$. The preceding step gives the required isomorphism on the left edge, so
\[
 \pi_1(S)\simeq\pi_1(T)*_{\pi_1(D_{\rm red})}\pi_1(C)
       \simeq\pi_1(C).
\]
The second isomorphism is induced by $C\hookrightarrow S$, with inverse on $\pi_1(T)$ given by $(\nu|_{D_{\rm red}})_*(j_*)^{-1}$ for $j:D_{\rm red}\hookrightarrow T$.

\item Proposition~\ref{night:P4R.1} makes $\pi_1(C)$ virtually
abelian. Apply \cref{lem:II:curve-normalization-graph} to this same
connected reduced curve. Its component genera $g_i$ and incidence
cycle rank $b$ therefore have the alternatives stated there, and
$e(C)\ge0$. Additivity and \cref{night:P4.13} give
\[
 e(S)=e(S_{\rm reg})+e(C)\ge0.
\]

\item Finally equality is impossible here. It would force
$e(C)=0$, so the same formula gives $c=1$, $g_1=1$
and $b=0$. Every irreducible component of
$D_{\rm red}$ maps finitely and nonconstantly onto
this curve. A rational vertical component would
therefore give a nonconstant map from $\mathbb P^1$
to its genus-one normalization, contrary to
Riemann--Hurwitz. There are no vertical conductor
components. The initial conductor is consequently
$2H$, and there can have been no blowup: the first
centre would lie on $2H$ and would create a vertical
conductor component of coefficient one. Thus
$T=T_0$ and $D_{\rm red}=H$. It follows that
\[
 e(S_{\rm reg})=e(T\setminus H)
       =(4-4g)-(2-2g)=2-2g<0,
\]
contradicting Proposition~\ref{night:P4.13}. This proves strict positivity.

\end{enumerate}
\end{proof}

\subsection{Rational incidence trees}

\begin{proposition}[Euler number for a rational incidence tree]\hypertarget{tag-02YG}{}\label{night:P4R.4}\label{stmt:P4R.4}
\StatementRecord{night:P4R.4}{proved}{02YG}{\SeriesRef{night:P4.13}, \SeriesRef{night:P4R.1}, \SeriesRef{night:P4R.3}}
For the support in \SeriesRef{night:P4R.3}, a singular-locus
incidence graph that is a tree of rational normalization
components gives $e(S)\ge3$.
\end{proposition}
\begin{proof}
Put $C=\operatorname{Sing}(S)_{\rm red}$ and let $c\ge1$ be
its number of irreducible components. The normalization
pushout and incidence-graph calculation in
\SeriesRef{night:P4R.3}, with all normalization genera zero
and graph cycle rank zero, give
\[
 \pi_1(S)\simeq\pi_1(C)=1,\qquad e(C)=c+1\ge2.
\]
The cyclic quotient in \SeriesRef{night:P4R.1} forces $m=1$.
Its formula for $b_1(S)$ gives $\det(I-A)\ne0$ for the positive
monodromy $A$. By \SeriesRef{night:P4.13},
$e(S_{\rm reg})=\det(I-A)\ge0$. This nonzero integer is at least
one, and additivity gives $e(S)=e(S_{\rm reg})+e(C)\ge3$.
\end{proof}

\subsection{Hopf normalizations}\label{subsec:II:smooth-hopf-normalizations}





\begin{lemma}[Hopf foliations under point blowups]\hypertarget{tag-P250}{}\label{lem:II:hopf-foliation-blowups}
\StatementRecord{lem:II:hopf-foliation-blowups}{proved}{P250}{\cref{lem:II:hopf-curves}}
Let $T_0$ be a primary Hopf surface, and let $T\to T_0$ be a sequence of $k$ point blowups. There is a nonsingular holomorphic foliation on $T_0$ tangent to every compact curve. Its saturated pullback has tangent line bundle $L\subset T_T$, tangent to every strict transform of a compact curve and every exceptional curve, with exactly $k$ simple saddle zeros. There is an integral divisor $V$ supported on the exceptional locus such that, numerically,
\[
 c_1(L)=-K_T+V,\qquad K_T^2=-k,\qquad K_TV-V^2=0.
\]
At a blowup $\mu$ with exceptional curve $E$,
\[
 K'=\mu^*K+E,\qquad V'=\mu^*V+\varepsilon E,
 \qquad \varepsilon=0\text{ at a regular point},\quad
 \varepsilon=1\text{ at a saddle}.
\]
\end{lemma}
\begin{proof}
\begin{enumerate}[label=\arabic*.,leftmargin=*]
\item The primary normal form in \cite[arXiv v2, formula~(1), p.~3; Theorem~2.1(1), p.~5]{Khaled2024Hopf} is
\[
 T_0=(\mathbb C^2\setminus\{0\})/\langle F\rangle,
 \qquad F(z,w)=(\alpha z+\lambda w^n,\beta w),
\]
where $0<|\alpha|\le|\beta|<1$ and $\lambda(\alpha-\beta^n)=0$.
Use the curve descriptions in \cite[arXiv v1, \S2.2.2, p.~4; \S2.2.6, p.~5]{AbashevaDeev2024AB} to choose an invariant vector field tangent to every compact curve. For $\lambda=0$ and no relation $\alpha^q=\beta^p$ with positive integers $p,q$, the curves are the two axis images; use $z\partial_z+w\partial_w$. When such a relation holds, the invariant map
\[
 (z,w)\longmapsto[z^q:w^p]
\]
descends to $T_0\to\mathbb P^1$. Every compact curve lies in a fibre: a horizontal curve would have positive intersection with a general fibre, whereas $b_2(T_0)=0$. Use $p z\partial_z+q w\partial_w$, which annihilates $z^q/w^p$ and is tangent along the axes as well. For $\lambda\ne0$, the sole curve is the image of $w=0$; use $n z\partial_z+w\partial_w$. In this last case its pushforward under $F$ is $(n(\alpha z+\lambda w^n),\beta w)$, exactly its value at $F(z,w)$. The diagonal fields are invariant by the same direct differential calculation. Each field is nowhere zero on $\mathbb C^2\setminus\{0\}$, so it descends to the required nonsingular holomorphic foliation.

\item Pull this foliation through the point blowups
and saturate its tangent sheaf. It is a line
bundle $L\subset T_T$. Every strict transform
of a compact curve of $T_0$ and every exceptional
curve is invariant. At a regular point use coordinates $(x,y)$ with
leaves $y=\text{constant}$. In the chart
$x=u,y=uv$ the saturated generator is
$u\partial_u-v\partial_v$; there is one
saddle, and the other blowup chart is regular
away from it. At an existing saddle use
$b x\partial_x-a y\partial_y$, with positive
integers $a,b$. The two charts have generators
\[
 b u\partial_u-(a+b)v\partial_v,\qquad
 (a+b)u\partial_u-a v\partial_v.
\]
One saddle is replaced by two. Each has
zero ideal $(u,v)$ and length one.
Consequently there are exactly $k$ saddles.

\item Let $V$ be the divisor measuring the
divisorial zero of the pulled-back transverse
one-form before saturation. Numerically
\[
 c_1(L)=-K_T+V.
\]
The pullback line bundle is topologically
trivial because it comes from $T_0$.
The preceding charts give, at every blowup $\mu$,
\[
 K'=\mu^*K+E,\qquad
 V'=\mu^*V+\varepsilon E,\qquad
 \varepsilon=
 \begin{cases}0&\text{at a regular point},\\
 1&\text{at a saddle}.
 \end{cases}
\]
Starting with $V=0$ and $K_{T_0}^2=0$, we obtain
\[
 K_T^2=-k,\qquad K_TV-V^2=0.
\]
The second identity is preserved because its
change is $\varepsilon^2-\varepsilon=0$.

\end{enumerate}
\end{proof}

\begin{lemma}[Point blowups and a minimal resolution]\hypertarget{tag-P291}{}\label{lem:II:last-blowup-minimal-resolution}
\StatementRecord{lem:II:last-blowup-minimal-resolution}{proved}{P291}{}
Let $\pi:Z\to T$ be the minimal resolution of a normal complex surface, and let $\sigma:Z\to Z_0$ be a sequence of point blowups of a smooth surface. If every curve exceptional for $\sigma$ is exceptional for $\pi$, then $\sigma$ is an isomorphism.
\end{lemma}
\begin{proof}
A nonempty point-blowup sequence has a last exceptional curve, which is a smooth rational curve of self-intersection $-1$ on $Z$. The hypothesis makes it $\pi$-exceptional, contradicting minimality of the resolution. Thus the sequence is empty.
\end{proof}

\begin{proposition}[Hopf normalizations]\hypertarget{tag-02X9}{}\label{night:P4.36}\label{stmt:P4.36}
\StatementRecord{night:P4.36}{proved}{02X9}{\cref{lem:II:hopf-curves,lem:II:single-conductor-branch,lem:II:hopf-foliation-blowups,lem:II:last-blowup-minimal-resolution,lem:II:exceptional-intersection,night:P4R.6}}
\hypertarget{tag-02WW}{}\label{night:P4.24}\label{stmt:P4.24}
\StatementRecord{night:P4.24}{proved}{02WW}{}
\hypertarget{tag-02YH}{}\label{night:P4R.5}\label{stmt:P4R.5}
\StatementRecord{night:P4R.5}{proved}{02YH}{}
\hypertarget{tag-02YK}{}\label{night:P4R.8}\label{stmt:P4R.8}
\StatementRecord{night:P4R.8}{proved}{02YK}{}
\hypertarget{tag-02WZ}{}\label{night:P4.27}\label{stmt:P4.27}
\StatementRecord{night:P4.27}{proved}{02WZ}{}
Let $f:Y\to B$ be a proper surjective holomorphic map with connected fibres from a connected smooth compact complex threefold with $b_1(Y)=b_2(Y)=0$ to a smooth compact curve, with smooth general fibre a complex two-torus. Let $mS$ be a fibre with irreducible, reduced, singular support. Write $\nu:N\to S$ for its normalization and $\sigma:T\to N$ for the minimal resolution. Suppose that $T$ has a Hopf minimal model $T_0$. Then $\sigma$ and the minimal-model contraction $T\to T_0$ are isomorphisms: $N$ is a smooth minimal Hopf surface. The reduced conductor on $N$ and its reduced image $C=\operatorname{Sing}(S)_{\rm red}$ are disjoint unions of smooth elliptic curves. Each map between their components is a finite unramified covering, and
\[
 e(S_{\rm reg})=e(C)=e(S)=0.
\]
The number $k$ of point blowups from $T_0$ to $T$ and the number $v$ of rational components of the pole divisor in \cref{night:P4R.6} are zero. Thus $e(S_{\rm reg})=k-v$ and $e(S)=k-v+e(C)$ as well.
\end{proposition}
\begin{proof}

\begin{enumerate}[label=\arabic*.,leftmargin=*]
\item Let $T\to N\to S\subset Y$ be the minimal normalization
resolution. Write $\sigma:T\to N$ for its minimal resolution
map and $\nu:N\to S$ for the finite normalization, and put
$\rho=(S\hookrightarrow Y)\circ\nu\circ\sigma$.
Let $\pi:T\to T_0$ contract its $k$ blowups
to a minimal Hopf surface. Take a finite unramified
primary Hopf cover $T'_0\to T_0$ and pull back the
sequence of point blowups. This gives a finite unramified
map $T'\to T$ and a blowup morphism $T'\to T'_0$.
If $q:T'\to T$ denotes this map, the ambient map is
$\rho\circ q:T'\to Y$. Each blowup centre has as many
distinct lifts as the degree of $q$, and all the Chern
calculations take place on $T'$. We write $T,T_0,k,\rho$
for these covered objects until we explicitly return to the original resolution.

\item Apply \cref{lem:II:hopf-foliation-blowups} to this primary Hopf surface and this sequence of point blowups. Its line bundle $L\subset T_T$ is tangent to the strict conductor and every exceptional curve. Pull the effective integral divisor of \cref{night:P4R.6} from the original resolution through $q$. Since $q$ is unramified, its pullback $D$ satisfies $K_T+D=q^*(K_{T_{\rm original}}+D_{\rm original})$ and is torsion. Apply \cref{lem:II:canonical-boundary-blowup} at each step to $D$ and its pushforward: a centre of multiplicity $\mu_x$ has exceptional coefficient $\mu_x-1$, so every centre lies on the invariant support. Write $\varepsilon=0$ at a regular point and $\varepsilon=1$ at a saddle. The lemma gives
\[
 c_1(L)=-K_T+V,\qquad V'=\mu^*V+\varepsilon E,
 \qquad K_T^2=-k,\qquad K_TV-V^2=0.
\]
Topologically trivial line-bundle factors from $T_0$ do not change these intersection calculations.

\item Put $B=\rho^*TY$. Its first and second rational Chern
classes vanish: $b_2(Y)=0$ gives $c_1(TY)=0$ rationally,
and Poincar\'e duality gives $H^4(Y;\mathbb Q)=0$ and
hence $c_2(TY)=0$. The derivative $L\to B$ may have a divisorial zero along curves contracted by $\rho$. Denote that effective integral divisor by $Z$.
Dividing by it gives the saturated line injection in
codimension one
\[
 L'=L(Z)\longrightarrow B,
 \qquad c_1(L')=-K_T+V+Z.
\]
The quotient $Q=B/L'$ is torsionfree of rank two.
Write $z_0\ge0$ for the length of $Q^{**}/Q$.
The double dual is locally free on the smooth surface \cite[Tag~0B3N]{StacksProject},
and the Chern sequence gives
\[
 c_2(Q^{**})=(-K_T+V+Z)^2-z_0.
\]

\item Let $E_B\subset B$ be the saturated rank-two image
of $d\rho:TT\to B$. Its determinant is
$\det E_B=\omega_T^{-1}(R)$, where the effective
integral divisor $R$ is the divisorial zero of the
determinant map. Since $E_B$ is saturated in $B$,
the divided foliation line $L'$ is contained in it.
The image of $E_B/L'$ in $Q^{**}$ is saturated in
codimension one. Its rank-one double dual
$J=(E_B/L')^{**}$ is a line bundle. The map into
$Q^{**}$ extends across the finitely many omitted
points by Hartogs, since both sheaves are locally
free there. Its saturation is therefore a line bundle
of first Chern class $R-V-Z$; the rank-one quotient
is an ideal sheaf of a zero-dimensional scheme
times a line bundle of first Chern class $K_T-R$.
Let its length be $z_1\ge0$. Thus
\[
 (-K_T+V+Z)^2-z_0
       =(R-V-Z)(K_T-R)+z_1.
\]
Using $K_T^2=-k$ and $K_TV-V^2=0$, this is precisely
\begin{equation}
 \mathcal B(R,Z)
 :=R^2-K_TR-VR-ZR-K_TZ+2VZ+Z^2
 =k+z_0+z_1.                                      \label{eq:hopf-normalization-chern}
\end{equation}

\item Push $R$ and $Z$ down along the chosen sequence
of point blowups.
At one reversed blowup write
\[
 R'=\mu^*R+cE,\quad Z'=\mu^*Z+dE,
 \quad K'=\mu^*K+E,\quad V'=\mu^*V+\varepsilon E.
\]
Here $c,d$ are integers.
Since $E^2=-1$ and a pulled-back divisor has zero
intersection with $E$, direct expansion gives
\[
 \mathcal B(R',Z')-\mathcal B(R,Z)
 =-c^2+(1+\varepsilon)c+cd+(1-2\varepsilon)d-d^2.
\]
Set $a=c-1$ and $b=d-(1-\varepsilon)$. For each
$\varepsilon\in\{0,1\}$ this is exactly
\begin{equation}
 1-(a^2+b^2-ab)\le1.                            \label{eq:hopf-normalization-increment}
\end{equation}
The quadratic form is
$\frac12(a^2+b^2+(a-b)^2)$; it is nonnegative
and is zero exactly when $a=b=0$.
All divisor intersections on the initial Hopf
surface are zero because $b_2(T_0)=0$.
Summing \eqref{eq:hopf-normalization-increment} therefore gives
$\mathcal B(R,Z)\le k$.
Equation~\eqref{eq:hopf-normalization-chern} forces equality, $z_0=z_1=0$,
and at every single blowup
\begin{equation}
 c=1,\qquad d=1-\varepsilon.                    \label{eq:hopf-normalization-equality}
\end{equation}

\item Every created exceptional curve has positive
coefficient in $Z$. At a regular centre, \eqref{eq:hopf-normalization-equality}
adds one to its nonnegative pulled-back coefficient.
The intermediate divisors are effective because each is
the pushforward of the final effective divisor $Z$.
At a saddle centre it adds zero, but at least one
previously created exceptional curve passes through
that saddle. The original Hopf foliation has no
zero; the two-chart blowup calculation places each
new saddle on the new exceptional curve. Induction
therefore gives a positive pulled-back coefficient
at every later saddle blowup as well.

\item The foliation is tangent to each exceptional curve.
If its image under $\rho$ were a curve, the derivative
along its tangent line would be nonzero at a generic
point: a nonconstant holomorphic map of smooth
normalized curves in characteristic zero is
generically immersive. Positive divisorial order
of $d\rho|_L$ contradicts that fact. Thus every
created exceptional curve is contracted by $\rho$,
and hence by the composite $\sigma\circ q$ to the original
normalization $N$, since its map to $S$ is finite.

\item Return to the original, possibly secondary Hopf surface. Every exceptional curve of $\pi:T\to T_0$ has an exceptional lift in the finite unramified cover, and the preceding argument contracts that lift to a point of $N$. Thus every $\pi$-exceptional curve is $\sigma$-exceptional. Apply \cref{lem:II:last-blowup-minimal-resolution} with its two maps equal to $\sigma:T\to N$ and $\pi:T\to T_0$, respectively. It gives $k=0$. Every compact curve on $T=T_0$ has square zero by \cref{lem:II:hopf-curves}. A nonempty exceptional configuration for $\sigma$ would have a negative definite intersection matrix by \cref{lem:II:exceptional-intersection}, applied to this resolution. Hence $\sigma$ is also an isomorphism, and $N=T=T_0$ is smooth minimal Hopf.

\item On the primary cover after $k=0$, the map $\rho:T\to Y$ is finite. The foliation line $L$ is a subbundle of $T_T$. Every divisorial zero of $d\rho|_L$ would lie on a compact curve tangent to $L$. The finite restriction to that curve is nonconstant and has nonzero derivative at its generic point. Thus $Z=0$. We already proved $z_0=0$, so $Q=Q^{**}$ is locally free. The exact sequence $0\to L\to B\to Q\to0$ splits locally; the map $d\rho|_L$ is therefore nowhere zero. Every conductor component on the primary Hopf cover is smooth elliptic and tangent to $L$, so its restriction to $Y$ is immersive everywhere. The finite unramified covering transfers this conclusion to the original normalization.

The reduced conductor on $N$ consists of disjoint smooth elliptic curves by \cref{lem:II:hopf-curves}. Apply \cref{lem:II:single-conductor-branch} to the finite normalization $\nu:N\to S\subset Y$: the inverse image of $C$ has one smooth branch at each of its points. The lemma identifies every pair of image branches. Each image branch is smooth by the preceding differential calculation, so $C$ is a disjoint union of smooth curves and every conductor map is finite unramified. Riemann--Hurwitz makes each image component elliptic. Since normalization is an isomorphism away from these curves,
\[
 e(S_{\rm reg})=e(N)-e((\nu^{-1}C)_{\rm red})=0,
 \qquad e(C)=0,\qquad e(S)=0.
\]
There are no rational conductor components by \cref{lem:II:hopf-curves}, giving $v=0$ and the asserted identities.

\end{enumerate}
\end{proof}

\subsection{The remaining smooth ruled normalizations}

\begin{proposition}[Positive Euler number for rational normalization]\hypertarget{tag-02WT}{}\label{night:P4.21}\label{stmt:P4.21}
\StatementRecord{night:P4.21}{proved}{02WT}{\SeriesRef{cdp:st:A.5.3}, \SeriesRef{night:P4.12}, \SeriesRef{night:P4.13}, \SeriesRef{night:P4.20}, \SeriesRef{night:P4.6}, \SeriesRef{night:P4R.1}}

Let $f:Y\to B$ be a proper surjective holomorphic map with connected
fibres from a smooth compact complex threefold with $b_1(Y)=b_2(Y)=0$
to a smooth compact curve, with smooth general fibre a complex two-torus.
Let $mS$ be a fibre with irreducible singular reduced support, and let
$\nu:T\to S$ be its smooth rational normalization. Then $e(S)>0$. A normal-crossing such support has $e(S)=2$.
\end{proposition}
\begin{proof}

\begin{enumerate}[label=\arabic*.,leftmargin=*]
\item Let $D$ be the conductor divisor on $T$, with its scheme
multiplicities, and put $C=\operatorname{Sing}(S)_{\mathrm{red}}$.
Propositions~\ref{cdp:st:A.5.3} and~\ref{night:P4.12} give
$\omega_T(D)\simeq\nu^*\omega_S$, a torsion line bundle.
The minimal resolution in the first proposition is the identity of the
smooth surface $T$. Pullback to a common point-blowup resolution of
$T$ and $\mathbb P^2$ injects $\operatorname{Pic}(T)$ into a free
abelian Picard group; see \cite[II.3, pp.~17--18; II.12, p.~24]{Beauville1978}.
Thus $D\sim-K_T$. Contract $T$ to a minimal rational surface $T_0$,
which is $\mathbb P^2$ or a Hirzebruch surface $\mathbb F_n$ with
$n\ne1$ \cite[V.10, pp.~79--80]{Beauville1978}, and push $D$ forward
to $D_0$. At each reversed point blowup $q$, apply
\cref{lem:II:canonical-boundary-blowup} with the anticanonical
divisors $D_{\rm new},D_{\rm old}$ and centre multiplicity $\mu$.
It gives $D_{\rm new}=q^*D_{\rm old}-E$ and exceptional
coefficient $\mu-1$, so every centre lies on the conductor.

\item Apply the minimal-boundary enumeration in
\eqref{cdp:eq:A.RT.13}--\eqref{cdp:eq:A.RT.16} to this same
$(T_0,D_0)$. On $\mathbb P^2$ its divisor is the listed plane
cubic. On $\mathbb F_n$ the tautological class $h$ and minimal
section $C_0$ are related by $h=C_0+nF$; thus
$2h+(2-n)F=2C_0+(n+2)F$, exactly the anticanonical class
in that enumeration. The ruling and fibre class are unchanged.

A smooth elliptic cubic or bisection has fundamental group
$\mathbb Z^2$. Every other listed support has rational
normalization components. The graph formula in
\cref{lem:II:curve-normalization-graph} gives $\mathbb Z$ for
a nodal cubic, a triangle of lines, a transverse conic and line,
a rational cycle formed by two sections and their vertical
components, or a rational bisection and a fibre meeting at two
distinct points. These are reduced nodal divisors. The section
intersection and vertical-degree equality $k+d=2$ in
\cref{prop:II:minimal-boundary-complements} makes every coefficient
in such a cycle one. Coincident intersections, a vertical fibre
through a section-intersection point, a tangency, a cusp, and
the multiple-component cases have tree incidence graph. Their
supports are simply connected. Thus the only nontrivial groups
are $\mathbb Z$ and $\mathbb Z^2$, and both occur only for
reduced nodal conductor.

\item At every blowup, \cref{lem:II:boundary-blowup-groups}
identifies the boundary fundamental groups through the
blowdown map. Thus $\pi_1(D_{\rm red})$ is unchanged along the
chosen sequence.
If the initial group was nontrivial, the conductor was
reduced nodal. A smooth-point blowup adds no conductor
component, and a node blowup inserts one component of
coefficient one. It remains reduced nodal.

\item In this latter case the generic transverse conductor exponent of every normalization branch is one. Apply \cref{lem:II:plane-branch-conductor} to this same reduced transverse plane equation; its exponent-one conclusion gives two smooth transverse branches. Consequently $S$ has ordinary double curves outside finitely
many points. Proposition~\ref{night:P4.20} makes $S$ normal crossing
everywhere. Write $t$ for its number of ordinary triple points.
On each normal-crossing chart the composite $T\to S\subset Y$
is the inclusion of a smooth coordinate plane, hence is immersive.
Apply \cref{night:P4.15} to the same finite normalization $\nu:T\to S$ and its immersion into $Y$. Its maximal-minor ideal is the unit ideal, so $\ell=0$. The rational alternative of that proposition gives $e(T)=6$.
The reduced nodal anticanonical conductor is either smooth elliptic
or a cycle of rational curves. Adjunction on a rational component
gives exactly two node preimages. Its number of components therefore
equals its number of nodes, and its Euler number equals that number;
the smooth elliptic case has Euler number zero and no nodes. Every triple point
of $S$ has three node preimages on $D$, so $e(D)=3t$.
Over the smooth part of $C$ the conductor has two points per fibre,
whereas a triple point has three. Hence
\[
e(D)=2e(C\setminus\{\text{triple points}\})+3t
     =2e(C)+t,
\]
and $e(C)=t$. Additivity gives $e(S)=6-2t$.
Proposition~\ref{night:P4.6}, applied to this reduced divisor $S$,
gives $e(S)=t$. Therefore $t=2$ and $e(S)=2$.

\item In the remaining case $D_{\mathrm{red}}$ is connected
and simply connected, and all its normalized components
are rational. Apply \cref{lem:II:proper-conductor-pushout} with $(X,Y,A,C,h)=(T,S,D_{\rm red},C,\nu)$. The finite normalization identifies the complementary open sets, the boundary is connected, and its inclusion is a cofibration by analytic triangulation. Both $\pi_1(D_{\rm red})$ and $\pi_1(T)$ are trivial, so the inclusion induces an isomorphism and the lemma gives $\pi_1(S)\simeq\pi_1(C)$ through $C\hookrightarrow S$.
Every component of $C$ has rational normalization,
because it is the finite nonconstant image of a
rational component of $D_{\mathrm{red}}$. Proposition~\ref{night:P4R.1} makes
$\pi_1(C)$ virtually abelian. Apply \cref{lem:II:curve-normalization-graph} to this
connected rationally normalized curve. It identifies the group
with $F_b$ and gives $b\le1$. If $c\ge1$ is its number of components, then
\[
 e(C)=c+1-b\ge c\ge1.
\]
Proposition~\ref{night:P4.13} gives $e(S_{\mathrm{reg}})\ge0$. Additivity proves
$e(S)=e(S_{\mathrm{reg}})+e(C)>0$.

\end{enumerate}
\end{proof}

\begin{proposition}[Euler number of an elliptic-ruled normalization]\hypertarget{tag-02WU}{}\label{night:P4.22}\label{stmt:P4.22}
\StatementRecord{night:P4.22}{proved}{02WU}{\SeriesRef{cdp:st:A.5.3}, \SeriesRef{night:P4.12}, \SeriesRef{night:P4.13}, \SeriesRef{night:P4.15}, \SeriesRef{night:P4R.1}, \SeriesRef{night:P4R.3}}

In the standing torus-fibration setting of this section, an irreducible singular support
$S$ whose smooth normalization $T$ is birationally ruled
over an elliptic curve has $e(S)\ge0$. Equality holds
exactly when $T$ is minimal.
\end{proposition}
\begin{proof}

\begin{enumerate}[label=\arabic*.,leftmargin=*]
\item Write $\nu:T\to S$ for normalization, $D$ for its full
conductor divisor, and $C=\operatorname{Sing}(S)_{\mathrm{red}}$.
Propositions~\ref{cdp:st:A.5.3} and~\ref{night:P4.12}, with the
identity as the resolution of $T$, give the torsion line
$\omega_T(D)\simeq\nu^*\omega_S$.
Push $D$ through its ruling blowdowns $\beta:T\to T_0$ to $D_0$ on the minimal elliptic ruled surface. Apply \cref{lem:II:irrational-ruled-boundary} to these same data. Its reduced-horizontal alternative gives one elliptic bisection or two disjoint elliptic sections, no vertical component, and no exceptional component in $D$ after the blowups. Apply \cref{lem:II:plane-branch-conductor} to each generic transverse plane equation of $S$. Its coefficient-one conclusion gives two smooth transverse branches at the general point of each conductor image.
Thus the normalization differential has no divisorial rank drop;
its analytic rank-drop locus is finite.
Proposition~\ref{night:P4.15} makes $T$ minimal,
makes the differential immersive everywhere, and
gives $e(S)=0$.

\item The other possibility is
\[
 D_0=2H+\sum_j a_jF_j,
\]
where $H$ is a smooth elliptic section. Apply the double-section alternative of \cref{lem:II:irrational-ruled-boundary} to this same divisor. It makes $D_{\rm red}$ the elliptic section with rational trees and its inclusion $j:D_{\rm red}\hookrightarrow T$ an isomorphism on fundamental groups. The finite normalization identifies the complementary open sets. Thus \cref{lem:II:proper-conductor-pushout}, with $(X,Y,A,C,h)=(T,S,D_{\rm red},C,\nu)$, gives $\pi_1(S)\simeq\pi_1(C)$, with inverse on $\pi_1(T)$ equal to $(\nu|_{D_{\rm red}})_*j_*^{-1}$.
Proposition~\ref{night:P4R.1} makes this group virtually abelian.
Every rational conductor component has rationally normalized image;
only $H$ can dominate an elliptically normalized component.
The formula in \cref{lem:II:curve-normalization-graph} gives $e(C)\ge0$.
Proposition~\ref{night:P4.13} and additivity therefore give $e(S)\ge0$.

\item If $e(S)=0$, the curve calculation forces $C$ to have
one component, genus-one smooth normalization, and
no point-identification graph cycles. Each conductor
component maps finitely and nonconstantly onto $C$.
Riemann--Hurwitz excludes rational conductor
components. Thus $D_{\mathrm{red}}=H$ and no blowup
occurred: the first blowup on $2H$ would create a
vertical conductor component of coefficient one.
Therefore $T$ is minimal.

\item Conversely suppose $T$ is minimal. In the double-section
case let $v$ be the number of distinct vertical fibres
in the conductor. Since $e(T)=e(H)=0$ and each vertical
$\mathbb P^1$ meets $H$ once,
\[
 e(S_{\mathrm{reg}})=e(T\setminus D_{\mathrm{red}})=-v.
\]
Proposition~\ref{night:P4.13} forces $v=0$, hence $D=2H$. At a general point
of its image curve, every normalization branch has
conductor exponent two. Apply \cref{lem:II:plane-branch-conductor} to that same transverse plane equation. It identifies the branch exponent with twice the intrinsic delta invariant plus the sum of its intersections with the other branches. A singular
branch together with another branch would have
exponent greater than two. Thus either there is one
ordinary cusp branch, or all branches are smooth:
two with intersection multiplicity two, or three
with pairwise intersection multiplicity one. Four
smooth branches would have exponent at least three,
so the list is exhaustive.

\item In the cusp case $H\to C$ has generic degree one.
The normalization of $C$ is therefore $H$, an elliptic
curve. Virtual abelianness of $\pi_1(C)$ excludes any
additional point-identification graph cycle, so
$e(C)=0$. Additivity gives $e(S)=0$. In the other two
cases the normalization differential is immersive at
the generic point of $H$. There is no other conductor
component, so its rank-drop set is finite. Proposition~\ref{night:P4.15}
applies and gives $e(S)=0$.

\end{enumerate}
\end{proof}


\subsection{Anticanonical supports on class VII surfaces}

\begin{proposition}[Anticanonical supports on class VII surfaces]\hypertarget{tag-01F0}{}\label{cdp:st:C.2.VII.1}
\StatementRecord{cdp:st:C.2.VII.1}{proved}{01F0}{\cref{cdp:st:C.3.8}; \cite[arXiv v3, Theorem~0.2, p.~2; Proposition~2.29, p.~19; Theorem~2.38, pp.~24--25]{Dloussky2006}; \cite[\S1, pp.~284--285, and \S3, p.~296]{DOT2003}; \cite[arXiv v3, Corollary~2.25 and Proposition~2.27, pp.~16--19]{Dloussky2006}; \cite[Introduction, p.~283]{DOT2003}; \cite[arXiv v3, Lemma~2.22 and the following paragraph, p.~15; Proposition~2.27, pp.~17--19]{Dloussky2006}}

Let $Z$ be a smooth compact surface of class VII, not
necessarily minimal, and let $D>0$ be an effective integral
divisor such that $K_Z+D$ is torsion in $\Pic(Z)$.
Write $R=D_{\mathrm{red}}$, and let $r(R)$ be its number
of rational irreducible components. Then
\begin{equation}
 e(R)=r(R),\qquad r(Z)\le b_2(Z),
\hypertarget{tag-015K}{}\label{cdp:eq:C.VII.1}
\end{equation}
where $r(Z)$ counts all rational curves of $Z$.
Every connected component of $R$ has arithmetic genus one.
In particular a connected component containing an elliptic
curve consists of that single elliptic curve with trees of
rational curves attached; it cannot contain a second
elliptic curve or a path joining two distinct points of the
first one.

Let $\sigma:Z\to Z_0$ be a minimal-model blowdown and put $D_0=\sigma_*D$. Then $D_0>0$ and $K_{Z_0}+D_0$ is torsion. For $b_2(Z_0)=0$, the surface $Z_0$ is Hopf and $(D_0)_{\rm red}$ is a disjoint union of smooth elliptic curves. For $b_2(Z_0)>0$, the surface has a global spherical shell and exactly $b_2(Z_0)$ rational curves. Its boundary is one of the following:
\begin{enumerate}
\item an intermediate surface, whose boundary support is a connected rational cycle with rational trees attached;
\item an Inoue--Hirzebruch surface, on which $D_0$ is the reduced cycle or union of two reduced cycles;
\item a parabolic Inoue surface, with a smooth elliptic curve $E$ and a disjoint rational cycle $C$, on which $D_0=E+C$, $C^2=0$, $C\equiv0$ and $E^2=-b_2(Z_0)$.
\end{enumerate}
\end{proposition}

\begin{proof}

\begin{enumerate}[label=\arabic*.,leftmargin=*]
\item Let $\sigma:Z\to Z_0$ be the stated minimal-model blowdown. A torsion line is represented by a finite unitary character, and point blowups induce isomorphisms on fundamental groups by \cref{lem:II:boundary-blowup-groups}. Thus $K_Z+D$ descends to a torsion line on $Z_0$. Pushing forward a meromorphic canonical section with poles along $D$, or its torsion-trivializing power, gives the effective divisor $D_0$ and $K_{Z_0}+D_0$ torsion. The divisor $D_0$ is nonzero, since a torsion canonical bundle would give Kodaira dimension zero.

\item For a minimal class VII surface with positive
second Betti number, put $L_0=\mathcal O_{Z_0}(K_{Z_0}+D_0)$. The numerically anticanonical datum of \cite[arXiv v3, Definition~2.20, p.~14]{Dloussky2006} is $m=1$ and $F=L_0^{-1}$, because $K_{Z_0}+D_0+F=0$ in $\operatorname{Pic}(Z_0)$. It gives a global spherical shell by
\cite[arXiv v3, Theorem~0.2, p.~2; Proposition~2.29, p.~19; Theorem~2.38, pp.~24--25]{Dloussky2006}.
The reduced support is one rational cycle with attached
trees, one or two rational cycles, or, in the parabolic
Inoue case, a disjoint rational cycle and a smooth elliptic
curve. The intermediate graph and its curve
configuration are described in
\cite[\S1, pp.~284--285, and \S3, p.~296]{DOT2003}.
The fact that no connected piece of the anticanonical
support is an isolated tree is
\cite[arXiv v3, Corollary~2.25 and Proposition~2.27, pp.~16--19]{Dloussky2006};
Lemma~2.22 there shows that a curve meeting that support
belongs to it. These statements apply to our divisor with
index one, since its twisting line is torsion and therefore
flat. They give, on the minimal model, arithmetic genus
one for each connected component and Euler number equal
to the number of its rational components. A surface with a global spherical shell has exactly $b_2$ rational curves, and every minimal class VII surface has at most $b_2$, as stated in
\cite[Introduction, p.~283]{DOT2003}.

\item On an Inoue--Hirzebruch surface, the reduced cycle or union of two cycles is numerically anticanonical. Its curve intersection matrix is negative definite. Subtracting this divisor from $D_0$ and intersecting with every component gives a vector in the kernel of that matrix; hence all coefficient differences are zero, and $D_0$ is reduced nodal.

For the parabolic Inoue case, write
$E$ for its elliptic curve and $C$ for its reduced rational
cycle. The exact canonical relation is
$K_{Z_0}+E+C=0$; the curves $E$ and $C$ are disjoint,
$E^2=-b_2(Z_0)$, and the intersection matrix of the
components of $C$ has kernel generated by their sum
\cite[arXiv v3, Lemma~2.22 and the following paragraph, p.~15; Proposition~2.27, pp.~17--19]{Dloussky2006}.
The cycle has $C^2=0$ by \cite[arXiv v3, Remark~0.6, pp.~3--4]{Dloussky2006}. The negative-definite intersection form of the class VII surface, described in \cite[arXiv v3, \S1, p.~4]{Dloussky2006}, then gives $C\equiv0$; the canonical relation also gives $E^2=K_{Z_0}^2=-b_2(Z_0)$. Consequently an effective integral divisor with
$c_1(K_{Z_0}+D_0)_{\mathbb R}=0$ has the form
$D_0=E+kC$, $k\ge0$: intersection with $E$ fixes its
coefficient at one, and the cycle intersection matrix
makes all cycle coefficients equal. Our stronger
torsion relation then gives
$K_{Z_0}+D_0=(k-1)C$ torsion. Choose $n>0$ trivializing $\OO((k-1)C)^{\otimes n}$.
A trivialization compared with the canonical meromorphic
section gives a meromorphic function with divisor
$n(k-1)C$. For $k>1$ it has no poles and a nonzero
effective zero divisor, impossible on the compact
connected surface. For $k=0$ its reciprocal has that
property. Hence $k=1$. In particular both connected
pieces of the stated reduced support really occur.

\item For a minimal model with $b_2=0$, \cref{lem:II:negative-kodaira-models} applies to its nonzero effective divisor and makes the surface Hopf. \Cref{lem:II:hopf-curves}, applied to this minimal surface $Z_0$, makes its compact curves smooth elliptic and pairwise disjoint. Thus every connected component of $(D_0)_{\rm red}$ has arithmetic genus one and Euler number zero, and there are no rational components on that model.

\item At each point blowup, \cref{lem:II:canonical-boundary-blowup} applies to the effective boundary and its torsion canonical line. Its coefficient identity gives
\begin{equation}
 D_{\mathrm{new}}=\sigma^*D_{\mathrm{old}}-E.
\hypertarget{tag-015O}{}\label{cdp:eq:C.VII.2}
\end{equation}
Thus the coefficient of $E$ is $\operatorname{mult}_pD_{\mathrm{old}}-1\ge0$, as in that lemma.
If this coefficient is positive, the reduced support
is its entire reduced inverse image: a smooth point
acquires a rational leaf, or an ordinary intersection
is subdivided by the exceptional curve. These operations
preserve arithmetic genus and increase both the Euler
number and the rational-component count by one.
If the coefficient is zero, the centre is a smooth point
of a multiplicity-one component and the new support is
its strict transform, isomorphic to the old support.
Thus all the stated support properties transfer in the
required direction. Finally each blowup increases $b_2$
by one and adds precisely its exceptional rational
curve; every other rational curve is the strict transform
of one on the preceding surface. The curve-count
inequality in \eqref{cdp:eq:C.VII.1} follows.

\end{enumerate}
\end{proof}






\subsection{The reduced resolved-conductor case}



\begin{proposition}[Semi-log-canonical fibre supports]\hypertarget{tag-02WX}{}\label{night:P4.25}\label{stmt:P4.25}
\StatementRecord{night:P4.25}{proved}{02WX}{\SeriesRef{night:P4.16}, \SeriesRef{night:P4.20}, \SeriesRef{night:P4.6}, \SeriesRef{night:P4.8}}

Let $f:Y\to B$ be a proper surjective holomorphic map with connected
fibres from a connected smooth compact complex threefold to a smooth
compact curve. Assume $b_1(Y)=b_2(Y)=0$ and that the smooth general
fibre is a complex two-torus. Every semi-log-canonical reduced fibre
support of $f$ has normal crossings.

For an irreducible reduced support $S$, let $\nu:N\to S$ be its
normalization, let $\pi:T\to N$ be a resolution, and put
$\rho=\nu\circ\pi$. For local residue generators $\Omega$ of $\omega_S$,
define the integral divisor $D=-\operatorname{div}(\rho^*\Omega)$, so that
$\omega_T(D)\simeq\rho^*\omega_S$ with this identification.
A reduced normal-crossing divisor $D$ gives normal crossings on $S$
and a smooth normalization $N$. Proposition~\ref{night:P4.6} applies
to each such support. Propositions~\ref{night:P4.7}--\ref{night:P4.8}
apply to fibrations whose reduced supports are all semi-log-canonical.
\end{proposition}
\begin{proof}

\begin{enumerate}[label=\arabic*.,leftmargin=*]
\item At a generic point of the conductor, an exponent one on every normalization
branch forces the transverse plane curve to be an
ordinary node by \cref{lem:II:plane-branch-conductor}, applied to the reduced transverse plane equation and its normalization branches. The support $S$ is a
Cartier hypersurface, hence is Cohen--Macaulay and
has invertible dualizing sheaf. On its normalization $N$, the pair consisting of $N$
and its reduced conductor is log canonical.
To check this last assertion, pull its canonical
divisor to $T$. The resulting boundary is exactly
$D$, by the displayed isomorphism and the
divisorial orders of a local generator. A smooth
surface with reduced normal-crossing boundary is
log canonical: blowing up a point gives new
boundary coefficient $\mu-1\le1$, and this
bound persists under successive point blowups.
Discrepancy comparison therefore gives the
assertion on $N$. Definition~2.3 and Lemma~2.2 of
\cite[p.~189]{LiuRollenske2012SLC} give semi-log-canonicity in
the complex analytic category.

\item The analytic classification is
\cite[Theorem~1.1, p.~186; Table~1, p.~187;
Propositions~3.1 and~4.1, pp.~192--200]{LiuRollenske2012SLC}.
An isomorphism of the singular hypersurface germs lifts to an ambient
biholomorphism: lift its three coordinate classes to convergent functions;
the induced cotangent isomorphism makes their linear part invertible.
Thus, in the displayed coordinates, the full local fibre function is
$U\prod_i g_i^{m_i}$, with $U$ a unit and the original component
multiplicities $m_i$. The normal singular germs in the table are isolated.
Each is irreducible, since distinct surface branches would meet along
a singular curve. Write its full function as $Ug^m$, choose a holomorphic
$m$th root of $U$, and set $h=U^{1/m}g$, so that $f=h^m$. Apply the isolated-hypersurface clause of \cref{cdp:st:A.5.17}
to this same local map $f=h^m$. That clause gives a vanishing sphere
of square $-2$ in the nearby torus fibre and excludes the germ.

\item The ordinary double point and $xyz=0$ have normal crossings.
The change $z\mapsto-z$ identifies the pinch point $x^2+y^2z=0$
with the case $n=2$ of Proposition~\ref{night:P4.16}.
For a weighted homogeneous full function $G$ of positive weights
$(w_1,w_2,w_3)$ and degree $W$, the unit is absorbed by the
biholomorphism
\[
 (z_1,z_2,z_3)\longmapsto
 (U(z)^{w_1/W}z_1,U(z)^{w_2/W}z_2,U(z)^{w_3/W}z_3),
\]
using one local logarithm of $U$; it satisfies $G\circ\Psi=UG$.
The weights used below are positive. Scaling the coordinates by
$\lambda^{w_i}$ places every compact cycle in $G^{-1}(1)$ in
$G^{-1}(\lambda^W)$ inside an arbitrarily small neighborhood.
For a deformation $G_s$ of an irreducible reduced equation with full
multiplicity $m$, use the full map $G_s^m$. At nonzero values its
levels are the disjoint union of the levels $G_s=\zeta t^{1/m}$,
where $\zeta^m=1$, and
$d(G_s^m)=mG_s^{m-1}dG_s$. Thus the boundary transport on these
levels carries the original full multiplicity.
The remaining double-point equations are
\[
 x^2+y^2z^2,\qquad
 x^2+y^2(z^2+y^{q-2}),\quad q\ge3.
\]
For the first equation the two normalization
branches are the smooth graphs $x=\pm i yz$.
Their tangent planes coincide at the origin,
whereas they meet transversely off that point
along the two axes. Put $U_1=x-iyz$ and $U_2=x+iyz$, and let their full fibre
multiplicities be $a,b>0$. For $d=\gcd(a,b)$ write
$a'=a/d$, $b'=b/d$. On the component $U_1^{a'}U_2^{b'}=1$
of $U_1^aU_2^b=1$, set
\[
 U_1=t^{b'},\qquad U_2=t^{-a'},\qquad
 yz=\frac{t^{-a'}-t^{b'}}{2i}.
\]
The inverse parameter is $t=U_1^rU_2^s$ for integers
$rb'-sa'=1$. Define
\[
 x(t)=\frac{t^{b'}+t^{-a'}}2,\qquad
 v(t)=\frac{t^{-a'}-t^{b'}}{2i}.
\]
The map $t\mapsto(x(t),v(t))$ embeds $\mathbb C^*$ as a smooth
connected curve $B$. Its holomorphic inverse is
$t=(x-iv)^r(x+iv)^s$, by the same B\'ezout identity.
The function $v$ has $a'+b'\ge2$ distinct simple zeros:
they satisfy $t^{a'+b'}=1$, and at each of them
$v'(t)=-(a'+b')t^{-a'-1}/(2i)\ne0$.
The surface $X_B$ of \cref{lem:II:conic-bundle-matching-sphere}
is this component of $U_1^aU_2^b=1$ through the identity map
on $(x,y,z)$. That lemma gives the embedded sphere and its
self-intersection $-2$.
Weights $(2,1,1)$ place this sphere in the full nearby fibre with
the original multiplicities. Every sphere in a complex two-torus
lifts to $\mathbb C^2$ and contracts there, so its homology class,
and hence its self-intersection, is zero. This excludes the germ.

\item For finite $q$, the polynomial $x^2+y^2(z^2+y^{q-2})$ is irreducible. As a quadratic in $x$, a factorization over $\mathbb C[y,z]$ would make $z^2+y^{q-2}$ a square. A monic square of degree two in $z$ has the form $(z+a(y))^2$; its zero $z$ coefficient forces $a=0$, contradicting its nonzero constant term $y^{q-2}$. Gauss' lemma gives irreducibility over the fraction field as well. With weights $(q,2,q-2)$, lowest weighted terms of an analytic factorization would factor this same polynomial, so the analytic germ is irreducible. Thus its full multiplicity is a common
integer $m$. After the unit map above, put
\[
 G_s=x^2+y^2(z^2+y^{q-2}+s).
\]
For $s\ne0$ small, at each point
$(0,0,z_0)$ with $z_0^2=-s$ the function
$Z=z^2+y^{q-2}+s$ is a local coordinate:
its $z$ derivative is $2z_0\ne0$.
The reduced equation in these coordinates is
$x^2+y^2Z$, a pinch point. At the original
boundary double axis $x=y=0$, $z\ne0$,
the transverse Hessian determinant is
$4(z^2+s)$, which stays nonzero on the
compact boundary when $s$ is small.
All boundary smooth points remain smooth.
Thus the zero divisor retains the same
boundary smooth and double strata.
The spherical class of square $-4$
from the pinch-point nearby fibre transports,
by the boundary transport in Proposition~\ref{night:P4.20}, to
the original nearby torus fibre.
This excludes every finite $q$.

\item The remaining triple-point equations are
\[
 xyz+x^p,\quad p\ge3,\qquad
 xyz+x^p+y^q,\quad q\ge p\ge3.
\]
For the first put
$G_s=xyz+x^p+s x^2
=x(yz+x^{p-1}+s x)$.
When $s\ne0$, its two local normalization
branches are smooth: the derivative in $x$
of the second factor at the origin is $s$.
Their tangent planes coincide. The two branches of $uv=0$ have
distinct tangent planes, and an ambient biholomorphism has an invertible
derivative preserving that distinction. Hence these two tangent branches
cannot be a normal crossing. Along
the boundary $y$ axis the transverse
Hessian determinant is $-y^2$, and along
the boundary $z$ axis it is $-z^2$.
Both axes persist as ordinary-double
strata. For full component multiplicities $a,b$, weights $(2,p-1,p-1)$
give degree $2a+2(p-1)b$ and absorb the original unit. Use the deformation
$x^a(yz+x^{p-1}+sx)^b$. At its boundary double strata, free normal
coordinates $u,v$ give $u^av^b$, which is submersive at nonzero
values; thus the same boundary transport applies with these multiplicities.
For fixed nonzero $s$, the coordinate $X=x(s+x^{p-2})$ is
biholomorphic near the origin. The full deformed function becomes
a unit times $X^a(X+yz)^b$. The coordinates
\[
 x'=X+yz/2,\qquad y'=y,\qquad z'=z/(2i)
\]
satisfy $x'-iy'z'=X$ and $x'+iy'z'=X+yz$.
After absorbing the unit with weights $(2,1,1)$, the preceding
two-branch case supplies a sphere of square $-2$. Boundary transport
therefore gives the contradiction in the original full fibre.

\item For the second equation, positive weights
$(q,p,pq-p-q)$ and the primitive degree-one expression in $z$
show analytic irreducibility by the same initial-form argument.
It therefore has one common full multiplicity. Put
\[
 G_s=xyz+x^p+y^q+s x^2+s y^3,
 \qquad 0<|s|\ll1.
\]
Its transverse Hessian along the boundary
$z$ axis has determinant $-z^2$,
so its sole boundary double stratum
persists. We check the smooth normalization
of its germ at the origin explicitly.
Since $\partial^2G_s/\partial x^2(0)=2s$,
the critical equation in $x$ has a unique
analytic solution $x=\xi(y,z)$ with
$\xi$ divisible by $yz$. The parametric
holomorphic Morse change in $x$ writes
\[
 G_s=X^2+y^2H(y,z).
\]
Divisibility follows directly on substituting
$\xi$: each of $\xi yz,\xi^p,s\xi^2$
is divisible by $y^2$, as are the two
pure $y$ terms. Moreover
\[
 H_y(0,0)=
 \begin{cases}1+s&q=3,\\s&q>3,\end{cases}
 \qquad H(0,z)=-z^2/(4s).
\]
The finite
birational map
\[
 \{v^2+H(y,z)=0\}\longrightarrow\{G_s=0\},
 \qquad (v,y,z)\longmapsto(X=yv,y,z),
\]
followed by the inverse Morse coordinate
change, is the normalization. Its source
is smooth because $H_y(0,0)\ne0$,
and the map is finite because $v$
satisfies the displayed monic equation.
It is birational on $y\ne0$ and is
therefore the normalization by normality
of its source. At the origin its
differential has rank one. The normalization of a normal-crossing
germ consists of coordinate-plane inclusions, all of differential rank
two. The normalization isomorphism induced by an ambient biholomorphism
preserves these ranks, so this rank-one map excludes normal crossings. The local construction in Proposition~\ref{night:P4.20} produces
a spherical class of nonzero square, and its boundary transport gives
the contradiction. The transverse Hessians and the restrictions
$H(0,z)$ and $H_y(0,0)$ follow by collecting the degree-two terms
and substituting the critical equation
$yz+p\xi^{p-1}+2s\xi=0$.

\item The remaining germs in Table~1 are smooth, ordinary double,
and triple normal crossing. The linear coordinates
$(u,v)=(x+iy,x-iy)$ give $uv=0$ for the double germ;
the triple germ is $xyz=0$. Their normalization morphisms are the
inclusions of their disjoint smooth coordinate-plane branches.
Proposition~\ref{night:P4.6} applies to each compact support so obtained.
When every reduced support is semi-log-canonical, the same local
conclusion at every fibre supplies the hypotheses of
Propositions~\ref{night:P4.7}--\ref{night:P4.8}.

\end{enumerate}
\end{proof}



\subsection{Excluding the remaining class-VII minimal models}

\begin{proposition}[Class VII minimal models with positive second Betti number]\hypertarget{tag-02XA}{}\label{night:P4.37}\label{stmt:P4.37}
\StatementRecord{night:P4.37}{proved}{02XA}{\cref{cdp:st:C.2.VII.1,night:P4R.6,night:P4.25,night:P4.15,night:P4.36}}
\hypertarget{tag-02WV}{}\label{night:P4.23}\label{stmt:P4.23}
\StatementRecord{night:P4.23}{proved}{02WV}{}

No irreducible singular support in the standing torus-fibration
setting of \cref{night:P4.12} has normalization
whose minimal resolution has a class-VII minimal model with
positive second Betti number. Consequently every class-VII
normalization occurring here is a smooth minimal Hopf surface
and has Euler number zero.
\end{proposition}
\begin{proof}

\begin{enumerate}[label=\arabic*.,leftmargin=*]
\item Let $\nu:N\to S$ be the finite normalization,
let $\sigma:T\to N$ be its minimal resolution, and let
$i:S\hookrightarrow Y$ be the inclusion. Put
$\rho=i\circ\nu\circ\sigma$.
Let $\beta:T\to T_0$ be the minimal-model contraction, and let $D_T$ be the effective integral divisor of \cref{night:P4R.6} for this same $\nu$ and $\sigma$. Its support is $\rho^{-1}(\operatorname{Sing}S)_{\rm red}$ and $K_T+D_T$ is torsion. Apply \cref{cdp:st:C.2.VII.1} to $(Z,D,\sigma)=(T,D_T,\beta)$. It gives the nonzero divisor $D_0=\beta_*D_T$ and the three positive-Betti minimal-model alternatives.

On the Inoue--Hirzebruch and parabolic Inoue alternatives, that proposition makes $D_0$ reduced nodal. At each reverse point blowup, \cref{lem:II:canonical-boundary-blowup}, applied to the successive pushforwards of this same $D_T$, gives exceptional coefficient $\mu-1$ and preserves a reduced normal-crossing boundary. Thus $D_T$ is reduced normal crossing. Its defining residue identity is precisely the one in \cref{night:P4.25}; that proposition gives normal crossings on $S$ and a smooth $N$. The minimal resolution $\sigma$ is therefore an isomorphism. On every normal-crossing chart, $N\to S\hookrightarrow Y$ includes a smooth coordinate plane, so the maximal-minor ideal is the unit ideal. Apply \cref{night:P4.15} to this map, with $\ell=0$, to get $e(T)=K_T^2$.

For the class-VII surface $T$, \cref{lem:II:negative-kodaira-models}\textup{(ii)}, invariance of $\chi(\mathcal O)$ under point blowup, and Noether's formula give
\[
 e(T)=b_2(T)>0,\qquad K_T^2=-b_2(T).
\]
This contradicts the preceding equality. Hence the remaining positive-Betti model $T_0$ is intermediate Kato. Write $n=b_2(T_0)>0$ and let $A_0$ be its maximal reduced curve divisor. Proposition~\ref{cdp:st:C.2.VII.1} identifies $\operatorname{Supp}D_0=A_0$, and $K_{T_0}+D_0$ is torsion.

\item Favre's Theorem~1.2.24(1)--(2),
\cite[p.~35]{FavreThesis2002}, states that the
$n$ rational curve classes span $H^2(T_0;\mathbb Q)$
and that their intersection form is negative definite.
Theorem~1.2.30 and the foliation construction and conormal
calculation in \cite[pp.~42, 44, 46]{FavreThesis2002}
give a holomorphic foliation with conormal
line
\[
 N_{\mathcal F}^*=\mathcal O_{T_0}(-A_0)
                       \otimes\text{a flat line}.
\]
Theorem~1.2.32 and its proof
\cite[pp.~42, 46]{FavreThesis2002} show that every
component of $A_0$ is invariant and that its singularities
are precisely the intersection points, with local
first integrals $z^\alpha w^\beta$, $\alpha,\beta>0$.
The invariant-curve and order-one assertions are also
given by Dloussky--Oeljeklaus
\cite[Theorem~3.1, pp.~1528--1529]{DlousskyOeljeklaus1999}.

\item It follows from the determinant of tangent and normal
lines that the foliation tangent line has numerical
first Chern class $-K_{T_0}-A_0$. The local generator
at an intersection is
$\beta z\partial_z-\alpha w\partial_w$.
At a regular point use the two charts of \cref{lem:II:hopf-foliation-blowups}. At a saddle substitute $a=\alpha$, $b=\beta$ into its coordinate calculation. The charts $x=u,y=uv$ and $x=uv,y=v$ give, respectively,
\[
 \beta u\partial_u-(\alpha+\beta)v\partial_v,
 \qquad (\alpha+\beta)u\partial_u-\alpha v\partial_v.
\]
These equalities follow by applying the differential of each displayed blowdown map to its vector field; they hold for the positive real weights here. Every new zero is a saddle with nonzero opposite-sign eigenvalues. The transverse-form orders in those same charts therefore give $\varepsilon=0$ at a regular centre and $\varepsilon=1$ at a saddle, with the canonical and saturation formulas stated below.

\item Let $T\to T_0$ contain $k$ blowups. Write
$c_1(L)=-K_T+V$ for the saturated foliation line;
initially $V_0=-A_0$, and
$V'=\mu^*V+\varepsilon E$ at a blowup.
The reduced divisor $A_0$ is a rational cycle with
rational trees attached. Its normalization has
$n$ rational components and $n$ nodes, including
self-nodes, so $\chi(\mathcal O_{A_0})=0$.
Adjunction, or the divisor exact sequence on $T_0$,
therefore gives
\[
 (K_{T_0}+A_0)A_0=0.
\]
Hence $K_{T_0}V_0-V_0^2=0$.
Under each blowup this quantity changes by
$\varepsilon^2-\varepsilon=0$.
Since $K_{T_0}^2=-n$, we obtain
\[
 K_TV-V^2=0,\qquad K_T^2=-n-k.
\]

\item Define $B=\rho^*TY$, the divisorial determinant zero
$R$, the divisorial foliation zero $Z$, and the
nonnegative lengths $z_0,z_1$ as in the proof of \cref{night:P4.36}.
Its saturated sheaf calculation, including the
Hartogs extension of the rank-one double dual,
now gives
\begin{equation}
 \mathcal B(R,Z)
 =R^2-K_TR-VR-ZR-K_TZ+2VZ+Z^2
 =n+k+z_0+z_1.                                  \label{eq:class-vii-chern}
\end{equation}
Every reversed blowup changes $\mathcal B$ by at
most one, by the unrestricted integer identity
\eqref{eq:hopf-normalization-increment}.

\item Push $R,Z$ down to effective divisors $R_0,Z_0$
on $T_0$. Their supports are contained in $A_0$:
the derivative is generically injective away from
the resolved conductor, and all
compact curves of the intermediate surface are
components of its maximal divisor. Put $D=D_0$,
$A=A_0$. Using $K=-D$ numerically and $V=-A$ gives
\[
 \mathcal B_0
 =R_0^2+DR_0+AR_0-Z_0R_0+DZ_0-2AZ_0+Z_0^2.
\]
For divisor classes define
$q(X,Y)=X^2-XY+Y^2$.
Adjunction above gives $DA=A^2$, while $D^2=-n$.
Expansion consequently gives the exact completion
\begin{equation}
 \mathcal B_0
       =n+q(R_0+D,Z_0+D-A).                     \label{eq:class-vii-completion}
\end{equation}
The form on divisor classes is negative definite.
Moreover
\[
 q(X,Y)=\tfrac12\bigl(X^2+Y^2+(X-Y)^2\bigr),
\]
so it is nonpositive, and it vanishes only for
$X=Y=0$. The first argument $R_0+D$ is nonzero:
it is a nonzero effective combination of the
$n$ independent curve classes. Thus
$\mathcal B_0<n$.
After the $k$ blowups we still have
$\mathcal B(R,Z)<n+k$, contradicting \eqref{eq:class-vii-chern}.
This excludes the intermediate minimal model.

\item All class-VII cases with positive second Betti number have been excluded. For $b_2(T_0)=0$, the same $D_0$ is a nonzero effective divisor. Apply \cref{lem:II:negative-kodaira-models}\textup{(iii)} to this minimal $T_0$ and its curve $\operatorname{Supp}D_0$; it makes $T_0$ Hopf. Proposition~\ref{night:P4.36}, with the original maps $\nu:N\to S$ and $\sigma:T\to N$, proves that both $\sigma$ and $\beta$ are isomorphisms and gives Euler number zero.

\end{enumerate}
\end{proof}

\subsection{Resolved ruled normalizations}

\begin{proposition}[Euler numbers with an irrational ruled resolution]\hypertarget{tag-02YJ}{}\label{night:P4R.7}\label{stmt:P4R.7}
\StatementRecord{night:P4R.7}{proved}{02YJ}{\cref{lem:II:irrational-ruled-boundary,lem:II:proper-conductor-pushout,lem:II:curve-normalization-graph,night:P4.13,night:P4R.1,night:P4R.6}}
\hypertarget{tag-02YL}{}\label{night:P4R.9}\label{stmt:P4R.9}
\StatementRecord{night:P4R.9}{proved}{02YL}{}
Let $f:Y\to B$ satisfy the torus-fibration hypotheses of \cref{night:P4.12}, and let $mS$ be a fibre with irreducible, reduced, singular support. Write $\nu:N\to S$ for its normalization and $\pi:T\to N$ for its minimal resolution. Suppose that $T$ is birationally ruled over a curve of positive genus. Then $e(S)\ge0$.

Let $D$ be the pole divisor of \cref{night:P4R.6}, let $\beta:T\to T_0$ contract to a minimal ruled model, and put $D_0=\beta_*D$. A reduced horizontal part of $D_0$ makes the ruling genus one and $N$ smooth and minimal, with $e(S)=0$.
\end{proposition}
\begin{proof}
\begin{enumerate}[label=\arabic*.,leftmargin=*]
\item Put $\rho=\nu\circ\pi$, $A=D_{\rm red}$ and $C=\operatorname{Sing}(S)_{\rm red}$. Proposition~\ref{night:P4R.6} gives $K_T+D$ torsion and $A=\rho^{-1}(C)_{\rm red}$. Apply \cref{lem:II:irrational-ruled-boundary} to the identical minimal ruling $T_0$, blowdown sequence $\beta$ and divisor $D$.

\item In its reduced-horizontal case, the ruling genus is one, and $A$ consists of one or two disjoint smooth elliptic strict transforms with coefficient one. Every exceptional curve of $\pi$ belongs to $A$ by \cref{night:P4R.6}. A connected exceptional configuration would therefore be one of these disjoint elliptic components, disjoint from the strict conductor, contrary to the same proposition. Thus $\pi$ is an isomorphism and $N=T$ is smooth. Apply \cref{lem:II:plane-branch-conductor} to each generic reduced transverse plane equation of $S$: coefficient one makes the two branches smooth and transverse. The differential of $T\to S\hookrightarrow Y$ therefore has rank two away from finitely many points. Proposition~\ref{night:P4.15}, on this same smooth elliptic ruled normalization, makes $T$ minimal and gives $e(S)=0$.

\item In the double-section case, the boundary lemma makes $A$ connected and gives an isomorphism $j_*:\pi_1(A)\to\pi_1(T)$ for its inclusion. The image $C$ is connected, and $\rho$ identifies $T\setminus A$ with $S_{\rm reg}$. A triangulation of the analytic pair makes $A\subset T$ a subcomplex \cite[Theorem~2, p.~464]{Lojasiewicz1964}. Apply \cref{lem:II:proper-conductor-pushout} with $(X,Y,A,C,h)=(T,S,A,C,\rho)$, including every exceptional curve of $\pi$. It gives
\[
 \pi_1(S)\simeq\pi_1(T)*_{\pi_1(A)}\pi_1(C)\simeq\pi_1(C).
\]
The last isomorphism is induced by $C\hookrightarrow S$, with inverse equal to the identity on $\pi_1(C)$ and to $(\rho|_A)_*j_*^{-1}$ on $\pi_1(T)$. Proposition~\ref{night:P4R.1} makes this group virtually abelian. Apply \cref{lem:II:curve-normalization-graph} to the same connected curve $C$; it gives $e(C)\ge0$, and its displayed genus and graph formula identifies the equality case. Proposition~\ref{night:P4.13}, applied to this fibre $mS$, gives $e(S_{\rm reg})\ge0$. Thus $e(S)=e(S_{\rm reg})+e(C)\ge0$.
\end{enumerate}
\end{proof}




\subsection{Numerical sign for arbitrary irreducible supports}

\begin{proposition}[Euler numbers of irreducible and normal crossing fibres]\hypertarget{tag-02WY}{}\label{night:P4.26}\label{stmt:P4.26}
\StatementRecord{night:P4.26}{proved}{02WY}{}

In the standing torus-fibration setting of \cref{sec:II:night-sign},
every irreducible reduced fibre support
has nonnegative topological Euler number. For a
singular irreducible support, rational resolution
of the normalization gives strictly positive
Euler number; so does a resolution whose minimal
model is of class VII with positive second
Betti number. These assertions permit singular
normalization.

Every normal-crossing support has the
nonnegative even Euler number of \cref{night:P4.6}.
Consequently, for a fibration all of whose
reduced supports are irreducible or have
normal crossings,
\[
 B=\mathbb P^1,\qquad c_3(Y)\in\{0,2\},
 \qquad b_3(Y)=2-c_3(Y).
\]
The numerical conclusion follows from
these Euler inequalities.
\end{proposition}
\begin{proof}
\begin{enumerate}[label=\arabic*.,leftmargin=*]
\item For an irreducible singular support use its normalization $\nu:N\to S$, minimal resolution $\sigma:T\to N$, and composite $\rho=\nu\circ\sigma$. Put
\[
 A=\rho^{-1}(\operatorname{Sing}S)_{\rm red},\qquad
 C=\operatorname{Sing}(S)_{\rm red}.
\]
Proposition~\ref{night:P4R.6} gives an effective integral divisor $D$ with support $A$ and $K_T+D$ torsion. Proposition~\ref{night:P4.12} gives $\kappa(T)=-\infty$. Apply \cref{lem:II:negative-kodaira-models} to a minimal model $T_0$ of this same $T$.

\item First let $T$ be rational. Push $D$ along a chosen minimal-model contraction $\beta:T\to T_0$ to $D_0$. A rational surface has torsion-free Picard group, by the common-resolution pullback in the proof of \cref{night:P4.21}; thus $D_0\sim-K_{T_0}$. Apply the plane-cubic and ruled degree-two enumeration in \eqref{cdp:eq:A.RT.13}--\eqref{cdp:eq:A.RT.16} to this identical pair $(T_0,D_0)$. The incidence argument in the proof of \cref{night:P4.21} uses only these boundary configurations. It gives a connected support: either its fundamental group is trivial and all its normalized components are rational, or $D_0$ is reduced normal crossing.

At each reverse blowup, apply \cref{lem:II:canonical-boundary-blowup,lem:II:boundary-blowup-groups} to the successive pushforwards of this same $D$. The first gives exceptional coefficient $\mu-1$ and preserves the reduced normal-crossing alternative. The second identifies the boundary groups through the blowdown and preserves connectedness. Every nonexceptional normalized component is the strict transform of its predecessor, and every added exceptional component is rational. Hence the two alternatives hold for $D$ on $T$.

\item In the reduced normal-crossing alternative, \cref{night:P4.25}, with the residue divisor $D$ just constructed, makes $S$ normal crossing and $N$ smooth. Its minimal resolution is the identity. Proposition~\ref{night:P4.21}, applied to the same smooth rational normalization $\nu:T=N\to S$, gives $e(S)=2$.

In the other alternative, $A$ is connected and simply connected, and $T$ is simply connected. The map $\rho$ identifies $T\setminus A$ with $S\setminus C$. A compatible analytic triangulation makes $A\subset T$ a subcomplex \cite[Theorem~2, p.~464]{Lojasiewicz1964}. Apply \cref{lem:II:proper-conductor-pushout} with $(X,Y,A,C,h)=(T,S,A,C,\rho)$. The inclusion $j:A\hookrightarrow T$ induces an isomorphism of the trivial groups, and the lemma identifies
\[
 \pi_1(S)\simeq\pi_1(C)
\]
through $C\hookrightarrow S$, with inverse on $\pi_1(T)$ given by $(\rho|_A)_*j_*^{-1}$.

Every normalized component of $C$ is rational: it is dominated by a noncontracted normalized component of $A$, and the resulting nonconstant map from $\mathbb P^1$ is finite. The image $C$ is connected. Proposition~\ref{night:P4R.1} makes $\pi_1(C)$ virtually abelian. Apply \cref{lem:II:curve-normalization-graph} to this same curve. If $b$ is its incidence-graph cycle rank and $c\ge1$ its number of components, that lemma gives $b\le1$ and
\[
 e(C)=c+1-b\ge c\ge1.
\]
Proposition~\ref{night:P4.13}, for the same full fibre $mS$, gives $e(S_{\rm reg})\ge0$. Thus $e(S)>0$.

\item For a ruled minimal model of positive base genus, apply \cref{night:P4R.7} to these same $\nu,\sigma$ and $D$ to obtain $e(S)\ge0$. For a class-VII minimal model, \cref{night:P4.37} applies to the same normalization resolution. It excludes the positive-second-Betti case altogether and makes the remaining normalization a smooth minimal Hopf surface with $e(S)=0$. In particular the positive-second-Betti case in the statement has no instance. This covers every singular irreducible support in the minimal-model list.

\item A normal irreducible support is smooth by \cref{night:P4.12}. Use the cyclic base-change map and torus deformation theorem in the final paragraph of the proof of \cref{night:P4.8}, with this full fibre $mS$. They give a finite unramified torus cover of $S$. Multiplicativity under that finite covering gives $e(S)=0$.

\item Normal-crossing supports have the nonnegative even Euler number given by \cref{night:P4.6}. For the final global conclusion, the degree-one Leray injection for this connected-fibre map $f:Y\to B$ in the proof of \cref{night:P4.7} gives $B=\mathbb P^1$. The Euler-additivity calculation in the proof of \cref{cdp:st:C.3.5} applies to this same $f$: the complement of the finitely many critical values is a proper smooth bundle with two-torus fibre and Euler number zero. Consequently
\[
 e(Y)=\sum_{b\in B}e\bigl((f^{-1}(b))_{\rm red}\bigr)\ge0.
\]
Apply \cref{cdp:st:C.1.1} with this $Y$ and $L=\mathcal O_Y$. Its Chern--Gauss--Bonnet and Poincar\'e-duality equalities give $c_3(Y)=e(Y)=2-b_3(Y)$, and its alternating degree-three pairing makes $b_3(Y)$ even. Thus $c_3(Y)\in\{0,2\}$ and $b_3(Y)=2-c_3(Y)$.
\end{enumerate}
\end{proof}






\section{Ruled Euler equality and its conductor}\label{sec:II:ruled-euler-equality}

\subsection{Conductor curves at ruled Euler equality}\label{subsec:II:zero-euler-ruled-conductor}

\begin{lemma}[Conductor curves of ruled supports with Euler number zero]\hypertarget{tag-P253}{}\label{lem:II:zero-euler-ruled-conductor}
\StatementRecord{lem:II:zero-euler-ruled-conductor}{proved}{P253}{\cref{lem:II:irrational-ruled-boundary,lem:II:proper-conductor-pushout,lem:II:curve-normalization-graph,lem:II:plane-branch-conductor,night:P4.13,night:P4R.1,night:P4R.6,night:P4.33}}
Let $f:Y\to B$ satisfy the torus-fibration hypotheses of \cref{night:P4.12}, and let $mS$ be a fibre with irreducible, reduced, singular support and $e(S)=0$. Write $\nu:N\to S$ for its normalization, $\pi:T\to N$ for its minimal resolution, $\rho=\nu\circ\pi$, and $D$ for the pole divisor of \cref{night:P4R.6}. Let $\beta:T\to T_0$ contract to a minimal ruled surface over a smooth curve $B_0$ of genus $g\ge1$. Suppose the horizontal part of $D_0=\beta_*D$ is $2H_0$, where $H_0$ is a section. Put
\[
 A=D_{\rm red},\qquad C=\operatorname{Sing}(S)_{\rm red},
\]
and let $H\subset T$ be the strict transform of $H_0$.

The curve $C$ is irreducible, has elliptic normalization $\eta_C:E\to C$, and is unibranch at every point. Every rational component of $A$ is contracted by $\pi$, and $H$ is the sole strict conductor component. The ruling $r:T\to B_0$ descends to $r_N:N\to B_0$. With $h:B_0\simeq H\hookrightarrow T$ and $s=\pi\circ h$ one has
\[
 r_N\circ\pi=r,\qquad r_N\circ s=\operatorname{id}_{B_0},
 \qquad (\nu^{-1}C)_{\rm red}=s(B_0).
\]
In particular $s$ is a closed holomorphic embedding. The finite map $\nu\circ s:B_0\to C$ factors uniquely through $\eta_C$.

All singularities of $N$ are rational, and
\[
 h^1(N,\mathcal O_N)=g,\qquad
 \chi(\mathcal O_{C_{\rm sch}})=g-1,\qquad
 \chi(\mathcal O_{D_{\rm sch}})=0,
\]
where $C_{\rm sch}$ and $D_{\rm sch}$ are the conductor schemes on $S$ and $N$. The image scheme $C_{\rm sch}$ is Cohen--Macaulay of pure dimension one. The generic transverse germ along $C$ is a cusp, a tacnode or an ordinary plane triple point. The degree of $B_0\to E$ is respectively one, two or three.
\end{lemma}
\begin{proof}
\begin{enumerate}[label=\arabic*.,leftmargin=*]
\item Apply \cref{lem:II:irrational-ruled-boundary} to the identical data $(T,T_0,\beta,D)$. It makes $A$ the section $H$ with rational trees attached and identifies the map on fundamental groups induced by $j:A\hookrightarrow T$. Proposition~\ref{night:P4R.6} gives $A=\rho^{-1}(C)_{\rm red}$, and $\rho$ identifies the complementary smooth loci. A compatible analytic triangulation makes $A\subset T$ a subcomplex. Apply \cref{lem:II:proper-conductor-pushout} with $(X,Y,A,C,h)=(T,S,A,C,\rho)$. The connected image $C$ satisfies
\[
 \pi_1(S)\simeq\pi_1(T)*_{\pi_1(A)}\pi_1(C)\simeq\pi_1(C).
\]
The final isomorphism is induced by $C\hookrightarrow S$; its inverse is the identity on $\pi_1(C)$ and $(\rho|_A)_*j_*^{-1}$ on $\pi_1(T)$. Proposition~\ref{night:P4R.1}, applied to the same full fibre $mS$, makes this group virtually abelian.

\item Apply \cref{lem:II:curve-normalization-graph} to this connected curve $C$. With $c$ components, normalization genera $g_i$ and incidence-graph cycle rank $b$, it gives
\[
 \pi_1(C)=\bigl(*_i\pi_1(C_i^{\rm norm})\bigr)*F_b,
 \qquad e(C)=c+1-2\sum_i g_i-b\ge0.
\]
Proposition~\ref{night:P4.13} gives $e(S_{\rm reg})\ge0$, so $e(S)=0$ makes $e(C)=0$. In the all-rational alternative of the graph lemma, $b\le1$ and $e(C)\ge c\ge1$. Its other alternative has one elliptic normalization, $b=0$, and $e(C)=c-1$. Thus $c=1$ and the normalization $E$ is elliptic. For one component the same graph has $b=\sum_p(r_p-1)$, where $r_p$ counts the normalized preimages of $p$. Hence every point of $C$ is unibranch.

\item A nonconstant map from a normalized rational component of $A$ to $C$ would lift through the normalization $\eta_C:E\to C$. Its resulting finite map $\mathbb P^1\to E$ contradicts Riemann--Hurwitz. Every rational component of $A$ is therefore contracted by $\rho$ and, since $\nu$ is finite, by $\pi$. Surjectivity onto the curve $C$ makes $H$ noncontracted and the sole strict conductor component. Its map to $C$ is finite and factors uniquely through $E$.

The section and vertical trees now satisfy the noncontracted clause of \cref{night:P4.33}, with its $(E,a,h,q,s,\rho_N)$ equal to $(B_0,r,h,r_N,s,\pi)$ here. That clause supplies the displayed ruling identities, rationality of all normalization singularities, and $h^1(N,\mathcal O_N)=g$. The same proposition supplies the two conductor Euler characteristics and the Cohen--Macaulay image scheme. The identity $r_Ns=\operatorname{id}_{B_0}$ makes $s$ a closed holomorphic embedding: $s$ is proper and its local restriction map $s^*$ is surjective, with right inverse $r_N^*$. Every connected exceptional configuration of $\pi$ meets the strict conductor by \cref{night:P4R.6}; since the strict conductor is $H$, each exceptional image lies on $s(B_0)$. This proves $(\nu^{-1}C)_{\rm red}=s(B_0)$.

\item The strict-transform coefficient of $H$ in $D$ is two, the coefficient of $H_0$ in $D_0$. At every general point of $C$, all its normalized transverse branches are the points of this same $H$, so all their conductor exponents are two. Apply \cref{lem:II:plane-branch-conductor} to the reduced transverse equation $F=\prod_iF_i$ at that point. For every branch it gives
\[
 2=2\delta(F_i)+\sum_{j\ne i}I(F_i,F_j).
\]
A singular branch therefore has delta invariant one and is the only branch. Its conductor is $t^2\mathbb C\{t\}$. Its ring has no element with nonzero linear term in $t$, since the inverse function theorem would make such an element a local parameter and make the branch smooth. Thus its ring is exactly
\[
 \mathbb C+t^2\mathbb C\{t\}
   =\mathbb C\{t^2,t^3\}
   \simeq\mathbb C\{X,Y\}/(Y^2-X^3),
\]
where the last map sends $(X,Y)$ to $(t^2,t^3)$. Lifting these two generators from the maximal ideal modulo its square to the transverse plane gives an invertible cotangent map, hence ambient holomorphic coordinates. This is the cusp germ.

For smooth branches the formula leaves two branches of intersection multiplicity two, or three with all pairwise intersections one. In the first case write them as $y=0$ and $y=x^2u(x)$, with $u(0)\ne0$. The coordinates $X=x\sqrt{u(x)}$, $Y=2y-X^2$ carry their reduced union to $Y^2-X^4=0$; the inverse uses the inverse of $x\mapsto x\sqrt{u(x)}$ and $y=(Y+X^2)/2$. In the second case write two as $x=0,y=0$ and the third as $y=xu(x)$, $u(0)\ne0$. The coordinates $X=xu(x)$, $Y=y$, with inverse given by the one-variable inverse of $x\mapsto xu(x)$, carry the union to $XY(Y-X)=0$. Four branches would give an intersection sum at least three; a sole smooth branch has conductor exponent zero. These cases exhaust the formula. The generic number of transverse normalization branches is the degree of $H\to C$, hence the degree of $B_0\to E$, giving one, two or three as asserted.
\end{enumerate}
\end{proof}

\subsection{Cuspidal equality on a ruled normalization}

\begin{proposition}[Zero-Euler cuspidal ruled supports]\hypertarget{tag-02X1}{}\label{night:P4.29}\label{stmt:P4.29}
\StatementRecord{night:P4.29}{proved}{02X1}{\cref{lem:II:zero-euler-ruled-conductor,night:P4.28,night:P4.22}}

Let $f:Y\to B$ satisfy the standing torus-fibration hypotheses
of \cref{sec:II:night-sign}, and let $S$ be an irreducible reduced
singular fibre support whose normalization has ruled
resolution. Suppose $e(S)=0$.
In the double-horizontal-divisor case,
when the generic transverse type is
cuspidal $A_2$, the normalization is
a smooth minimal elliptic ruled surface.
The image conductor is a smooth elliptic
curve and every point along it has the
product cusp germ of \cref{night:P4.28}.
\end{proposition}
\begin{proof}
\begin{enumerate}[label=\arabic*.,leftmargin=*]
\item Write $\nu:N\to S$, $\pi:T\to N$ and $\rho=\nu\circ\pi$. Proposition~\ref{night:P4.26} excludes a rational resolution from $e(S)=0$, so the ruling base has genus $g\ge1$. Apply \cref{lem:II:zero-euler-ruled-conductor} to this same full fibre, normalization resolution and double horizontal pole divisor. Its notation gives the sole strict conductor $H$, the irreducible image $C$ with elliptic normalization $E$, and the descended ruling and section.

The generic cusp has normalization $v\mapsto(v^3,-v^2)$ and one branch. Thus the finite map $H\to E$ has degree one. A finite degree-one map between these smooth compact curves is an isomorphism, so $g=1$. The common lemma gives $\chi(\mathcal O_{C_{\rm sch}})=0$ and makes $C_{\rm sch}$ Cohen--Macaulay of pure dimension one.

\item At a generic cusp, the conductor is $v^2\mathbb C\{v\}$, the image of $(x,y)$ under $x=v^3,y=-v^2$. Hence the image conductor scheme is generically the reduced curve $C$. Its nilradical is supported at finitely many points and vanishes by the depth-one property. Therefore $C_{\rm sch}=C$ at every point, including images of exceptional curves.

The normalization sequence for $\eta_C:E\to C$ gives
\[
 0=\chi(\mathcal O_C)=\chi(\mathcal O_E)-\sum_p\delta(C,p)
   =-\sum_p\delta(C,p).
\]
Every delta invariant is zero, so $C$ is smooth.

\item Apply \cref{night:P4.28} at every point of this smooth reduced singular curve, with the same local reduced equation of $S$ and generic cusp germ. Its product coordinates give normalization
\[
 (t,v)\longmapsto(t,v^3,-v^2)
\]
and identify the reduced conductor $v=0$ with $C$. Thus $N$ is smooth along $C$ and equals $S_{\rm reg}$ away from it. Proposition~\ref{night:P4.22}, applied to this same smooth elliptic ruled normalization with $e(S)=0$, makes it minimal.
\end{enumerate}
\end{proof}


\subsection{The conormal calculation and ordinary-triple equality}

\begin{lemma}[Conormal defects of compact curves]\hypertarget{tag-P254}{}\label{lem:II:rational-conormal-defects}
\StatementRecord{lem:II:rational-conormal-defects}{proved}{P254}{\cref{lem:II:plane-branch-conductor}}
Let $C$ be a compact reduced curve in a smooth complex threefold $Y$, with normalization $\nu_C:E=\coprod_{a=1}^c E_a\to C$, inclusion $i:C\hookrightarrow Y$ and $j=i\nu_C$. Let $g_a$ be the genus of $E_a$, and suppose $\deg(j|_{E_a})^*\Omega_Y=0$ for every $a$. Put $\gamma=\sum_{a=1}^c(1-g_a)$. Write $P$ for the ideal of $C$ and $P^{(2)}$ for the functions vanishing to order at least two along every local generic branch. Put
\[
 K=\ker(j^*\Omega_Y\to\Omega_E),\qquad
 \operatorname{im}(dj)=\Omega_E(-R),\qquad Q=P/P^{(2)}.
\]
Differentiation embeds $Q$ in $j_*K$. At $p\in C$ let $\tau_p$ be the length of $j_*K/Q$, let $r_p$ be the sum of the coefficients of $R$ above $p$, and put $d_p=\delta(C,p)+\tau_p-r_p$. Then
\[
 \chi(\mathcal O_Y/P^{(2)})=5\gamma-\sum_p d_p.
\]
At a smooth point $d_p=0$. At a planar point $d_p=5\delta(C,p)$, so an ordinary node has $d_p=5$. At a nonplanar point with at most two analytic branches, $d_p\ge7$.
\end{lemma}
\begin{proof}
\begin{enumerate}[label=\arabic*.,leftmargin=*]
\item Dividing the derivative row on $E$ by its common zeros gives a surjection onto the line bundle $\Omega_E(-R)$. Its kernel $K$ is locally free of rank two and is a direct summand in a local frame of $j^*\Omega_Y$. The degree hypothesis gives
\[
 \deg(K|_{E_a})=2-2g_a+\deg(R|_{E_a}).
\]
Differentiation of a defining function gives its normal derivative along each generic smooth branch of $C$. Its kernel in $P$ is exactly $P^{(2)}$. A section of $Q$ killed away from finitely many points has zero generic derivative on every branch and is therefore zero. Thus the resulting map $Q\to j_*K$ is injective, and its cokernel has finite length $\tau=\sum_p\tau_p$. Write $j_a=j|_{E_a}$ and $K_a=K|_{E_a}$. The direct image is the finite direct sum $j_*K=\bigoplus_a(j_a)_*K_a$. Riemann--Roch on each $E_a$, followed by addition of these Euler characteristics, gives
\[
 \chi(Q)=4\gamma+\deg R-\tau.
\]
The normalization sequence for this same $\nu_C$ gives $\chi(\mathcal O_C)=\gamma-\sum_p\delta(C,p)$. Applying Euler additivity to
\[
 0\longrightarrow Q\longrightarrow\mathcal O_Y/P^{(2)}
   \longrightarrow\mathcal O_C\longrightarrow0
\]
proves the formula. For $E=\mathbb P^1$ it reads $5-\sum_p d_p$; for a connected elliptic $E$ its constant term is zero; for two rational normalized components its constant term is ten. All local lengths below are computed from the same map $Q\hookrightarrow j_*K$, including branches which belong to different global components.

\item On each normalized branch choose a parameter $t$ and ambient coordinates vanishing at its image. If $m=\min_i\operatorname{ord}_t x_i(t)$, the common order of the derivative entries is $m-1$. This is the coefficient of $R$. A generator of $\bigwedge^2K$ has a unit Pl\"ucker coordinate because $K$ is saturated in the free ambient rank-three module. Let $L$ be the $\mathbb C\{t\}$-span of the differentiated defining equations. The elementary-divisor decomposition of this rank-two lattice gives
\[
 \operatorname{length}(K/L)=
 \min\operatorname{ord}_t\{\text{two-by-two gradient minors}\}.
\]
Over all branches above $p$, the inclusions are
$Q\subset j_*L\subset j_*K$. Their finite quotient lengths are complex dimensions, so $\tau_p$ is at least the sum of these branchwise lattice lengths. At an embedding-dimension-three point every defining equation has no linear term. Each gradient entry then has order at least $m$ on this branch, and its two-by-two minors have order at least $2m$.
At a smooth point $m=1$, differentiation identifies the free conormal module with $K$, giving $R=0$, $\tau_p=0$ and $d_p=0$.

\item First consider a planar curve, with ideal $(z,h(x,y))$.
Its conormal is free over $\mathcal O_C$, generated by
$dz,dh$. On each normalized branch, write
$dh=a_i u_i$ for a primitive plane kernel generator,
and put $v_i=\operatorname{ord}a_i$.
Write $r_i$ for the coefficient of $R$ on that branch. The primitive plane kernel has a unit coefficient. At a singular unibranch point of multiplicity $m$, both entries of $dh$ have order at least $m$, so $v_i\ge m$; its conormal lattice quotient has at least that length. This is the same branchwise normal-derivative map used above.
The inclusion of each of these two free summands into
its normalized module has length $\delta(C,p)$.
The second summand has the additional vanishing orders
$v_i$. Hence
\[
 \tau_p=2\delta(C,p)+\sum_i v_i.
\]
Apply \cref{lem:II:plane-branch-conductor} to $A=\mathbb C\{x,y\}/(h)$ and its finite normalization $B$. For its conductor $I=\operatorname{Hom}_A(B,A)$ and branch exponents $c_i$, the lemma gives $\sum_i c_i=\operatorname{length}(B/I)=2\delta(C,p)$ and the identity $\omega_B=I\omega_A$ inside the meromorphic differentials.
Choose a generic plane coordinate $x$ with minimal
valuation on every normalized branch. The dualizing
generator $dx/h_y$ then pulls back with
order $r_i-v_i$. The identity
$\omega_B=I\omega_A$ therefore gives
$c_i=v_i-r_i$. Combining these equalities proves
$\tau_p-r_p=4\delta$ and the stated planar formula.

\item Now consider a nonplanar unibranch curve of multiplicity
$m$. Its embedding dimension is three. By the gradient-lattice calculation above,
all gradient entries have valuation at least $m$, so
$\tau_p\ge2m$, while $r_p=m-1$ and
$\delta\ge m-1$. A nonplanar unibranch curve has
$m\ge3$. For multiplicity two, choose an ambient coordinate and
a normalization parameter with $x=t^2$, and write
each other coordinate as $a(x)+t b(x)$. After subtracting
the even parts, choose the nonzero $b$ of least order.
The other odd coordinate is an analytic multiple of it
over $\mathbb C\{x\}$. This is a smooth ambient equation
containing the curve, so its embedding dimension is at
most two. The case where both odd parts vanish contradicts
the birational normalization parametrization.

\item The estimates give $d_p\ge2m\ge6$. Equality at six
would require $m=3$, $\delta=2$, and $\tau_p=6$.
The valuation semigroup of a unibranch curve has exactly
$\delta$ gaps, as follows by filtering its local ring
and its normalization by order. Below the conductor,
each graded piece of the normalization is one-dimensional;
the corresponding piece of the curve ring has dimension
one precisely at an attained order. Summing the quotient
dimensions counts the gaps. Thus the two gaps here
are one and two, and we can choose generators with
initial orders three, four, and five. They form an
ambient coordinate system: their classes generate the
three-dimensional cotangent space, since the square
of the maximal ideal has order at least six.
In a defining relation, the quadratic monomials
$x^2$ and $xy$ have respectively the unique least
orders six and seven, so their coefficients vanish.
Every gradient with valuation three consequently has
its initial vector in the single $z$ direction
(from a possible $xz$ term). Two such initial
gradient vectors are dependent. Every two-by-two
minor therefore has valuation at least seven,
contradicting $\tau_p=6$. Hence $d_p\ge7$ in this case.

\item Finally consider a nonplanar point with two branches
of multiplicities $m_1,m_2$, with reduced ideals $P_1,P_2$.
Two smooth curve branches
in a smooth threefold are contained in a smooth
surface. A generic linear function is immersive on both
smooth branches. Choose it as the common coordinate
$x$ and write them as graphs
$(y,z)=(a_i(x),b_i(x))$. Divide their difference pair
by its common power of $x$; the equation
\[
 \frac{b_1-b_2}{x^k}(y-a_1)
 -\frac{a_1-a_2}{x^k}(z-b_1)=0
\]
has a nonzero linear derivative and contains both.
Writing its left side as $F$, the map
$(x,y,z)\mapsto(x,y,F)$ is a coordinate change when
$F_z(0)\ne0$; otherwise use $(x,y,z)\mapsto(x,F,z)$.
This makes the containing surface a coordinate plane.
Thus in the nonplanar case $m_1+m_2\ge3$.
The absence of a linear defining equation gives,
on the two branches separately,
$\tau_p\ge2(m_1+m_2)$.
The reduced-curve normalization sequence gives
$\delta(C,p)=\delta_1+\delta_2+
\operatorname{length}\mathcal O_Y/(P_1+P_2)
\ge m_1+m_2-1$.
Also $r_p=m_1+m_2-2$.
It follows that
$d_p\ge2(m_1+m_2)+1\ge7$.


\end{enumerate}
\end{proof}

\begin{lemma}[Primary ordinary-triple conductor schemes]\hypertarget{tag-P256}{}\label{lem:II:triple-conductor-symbolic-square}
\StatementRecord{lem:II:triple-conductor-symbolic-square}{proved}{P256}{\cref{lem:II:conductor-duality}}
Let $S$ be a reduced hypersurface in a smooth complex threefold $Y$, with finite normalization $\nu:N\to S$. Let $C$ be an irreducible curve component of its reduced image conductor, and suppose that its generic transverse germ is the ordinary triple point $xy(x-y)=0$. Write $P\subset\mathcal O_Y$ for the reduced ideal of $C$. The ambient ideal of the primary image conductor component supported on $C$ is $P^{(2)}$, taking order two on every local analytic branch of $C$. This also identifies the full image conductor scheme when $C$ is its whole reduced support.
\end{lemma}
\begin{proof}
Apply \cref{lem:II:conductor-duality} to this finite normalization. The image conductor is Cohen--Macaulay of pure dimension one. Its primary component along $C$ is obtained by intersecting the isolated primary ideals belonging to all local minimal primes of $C$. These ideals glue by their unique localization-and-contraction characterization. The quotient by their intersection injects into the direct sum of their primary quotients. Its associated primes are therefore among these one-dimensional minimal primes; it has no finite-support submodule and has depth one. Write $J$ for this ambient component ideal.

At a generic ordinary triple point $xy(x-y)=0$,
parametrize the three branches by $(v,0),(0,v),(v,v)$.
The values of $x^2,xy,y^2$ are respectively
$(v^2,0,v^2),(0,0,v^2),(0,v^2,v^2)$.
Their coefficient matrix is invertible. More explicitly,
$(x^2-xy)a(x)+(y^2-xy)b(y)+xyc(x)$ has the three
values $(v^2a(v),v^2b(v),v^2c(v))$, for arbitrary
convergent series $a,b,c$. Thus all terms of order
at least two can be prescribed independently.
The first-order image
is spanned by $(v,0,v)$ and $(0,v,v)$; none of the
three individual coordinate vectors belongs to it.
Multiplication by the three branch idempotents
therefore excludes every nonzero first-order
conductor term. A nonzero constant is excluded in
the same way. The conductor is exactly $(x,y)^2$.

Consequently $J_{\mathfrak p}=P_{\mathfrak p}^{2}$ at every local minimal prime $\mathfrak p$ of $C$. This gives $J\subset P^{(2)}$. A section of $P^{(2)}/J$ vanishes at all these generic points and hence has finite support. Depth one of $\mathcal O_Y/J$ makes that submodule zero. Thus $J=P^{(2)}$ at every point, including multibranch self-identifications of the global curve $C$.
\end{proof}

\begin{proposition}[Symbolic double of a unibranch elliptic curve]\hypertarget{tag-02X3}{}\label{night:P4.30}\label{stmt:P4.30}
\StatementRecord{night:P4.30}{proved}{02X3}{}

Let $Y$ be the smooth compact threefold in the torus-fibration
setting of \cref{night:P4.12}, and let $C\subset Y$ be a compact
irreducible reduced curve. Suppose its normalization is an elliptic
curve $E$, and every point of $C$ is unibranch. Write $P$ for its
ideal in $Y$ and $P^{(2)}$ for the ideal of functions vanishing to
order at least two at the generic point of $C$. Then
\[
 \chi(\mathcal O_Y/P^{(2)})\le0.
\]
Equality forces $C$ to be smooth. The assertion uses
$c_1(TY)\cdot C=0$, which follows from $b_2(Y)=0$.
\end{proposition}
\begin{proof}
Write $\nu:E\to C$ for normalization, $i:C\hookrightarrow Y$ for inclusion and $j=i\nu$. Apply \cref{lem:II:rational-conormal-defects} with these same $Y,C,E,j$. The normalization has genus one, and $b_2(Y)=0$ gives $\deg j^*\Omega_Y=0$. Its formula is
\[
 \chi(\mathcal O_Y/P^{(2)})=-\sum_p d_p.
\]
Every point is unibranch. A singular planar point has $d_p=5\delta(C,p)>0$, and a nonplanar point has $d_p\ge7$; smooth points have $d_p=0$. Thus the displayed Euler characteristic is nonpositive, with equality precisely when every normalization defect vanishes. Then $\nu:E\to C$ is an isomorphism and $C$ is smooth.
\end{proof}

\begin{proposition}[Zero-Euler ruled supports with ordinary triple curves]\hypertarget{tag-02X4}{}\label{night:P4.31}\label{stmt:P4.31}
\StatementRecord{night:P4.31}{proved}{02X4}{\cref{lem:II:zero-euler-ruled-conductor,night:P4.30,night:P4.22}}

Let $f:Y\to B$ satisfy the standing torus-fibration hypotheses
of \cref{sec:II:night-sign}, and let $S$ be an irreducible reduced
singular fibre support whose normalization has ruled
resolution. Suppose $e(S)=0$.
In the double-horizontal-divisor case, when the generic
transverse type is the ordinary triple point $xy(x-y)=0$,
the normalization is a smooth minimal elliptic ruled surface.
Its conductor image is a smooth elliptic curve, and its
reduced conductor is an unramified cover of that curve of
degree three. Every germ along the image is analytically
the product of an ordinary plane triple point with a disc.
\end{proposition}
\begin{proof}
\begin{enumerate}[label=\arabic*.,leftmargin=*]
\item Write $\nu:N\to S$, $\pi:T\to N$ and $\rho=\nu\circ\pi$. Proposition~\ref{night:P4.26} excludes rational resolution at Euler number zero. Apply \cref{lem:II:zero-euler-ruled-conductor} to this same positive-genus ruling and double horizontal pole divisor. It gives the irreducible, everywhere unibranch image curve $C$ with elliptic normalization $E$, the sole strict conductor $H$ of genus $g\ge1$, and the ruling maps $r_N:N\to B_0$ and $s=\pi h:B_0\to N$, where $h:B_0\simeq H\hookrightarrow T$, satisfying $r_Ns=\operatorname{id}_{B_0}$. The generic ordinary triple point has three normalization branches, so the finite map $H\to E$ has degree three. The lemma also gives rational normalization singularities and
\[
 \chi(\mathcal O_{C_{\rm sch}})=g-1.
\]
\item Let $C_{\rm sch}$ be the image conductor scheme and $P\subset\mathcal O_Y$ the reduced ideal of $C$. Apply \cref{lem:II:triple-conductor-symbolic-square} to this same finite normalization $\nu:N\to S$. The common reduction has made $C$ the whole irreducible image support, and its generic transverse germ is the stated ordinary triple point. Thus
\[
 \mathcal O_{C_{\rm sch}}=\mathcal O_Y/P^{(2)}.
\]
The equality includes all special points of $C$.

\item The conductor Euler characteristic from the first step is $\chi(\mathcal O_{C_{\rm sch}})=g-1$ for this same scheme, now identified with $\mathcal O_Y/P^{(2)}$.
Proposition~\ref{night:P4.30} applies to the
irreducible unibranch curve with elliptic
normalization and gives $g-1\le0$. Since $g\ge1$,
we have $g=1$ and equality. That proposition
makes $C$ smooth. Riemann--Hurwitz now gives
zero ramification for the degree-three map $H\to C$,
so it is unramified.
\item The three points of $H$ over a point of $C$
remain distinct in $N$: the section identity $r_N\pi h=\operatorname{id}_{B_0}$ in \cref{lem:II:zero-euler-ruled-conductor} makes $\pi|_H$ a closed embedding. Every connected exceptional configuration
meets the strict conductor by \cref{night:P4R.6}.
Thus every point of $\nu^{-1}(C)$ lies on $\pi(H)$,
and the finite normalization gives precisely three
local irreducible branches of $S$.
Each branch is a hypersurface in the smooth threefold
$Y$ and is smooth at a generic point of $C$.
It is smooth away from $C$ as well; after shrinking
at the chosen point, its singular locus is at most
that point. The hypersurface Cohen--Macaulay
property and regularity in codimension one make
the branch normal. Its finite birational
normalization map is therefore an isomorphism,
identifying it with the corresponding component
of the local normalization.
\item Let $f_1,f_2,f_3$ be equations of the three branches.
The intersection of any two has reduced support $C$
and generic multiplicity one, since the generic
transverse triple point is ordinary.
On one normal branch, the equation of the other
is a nonzero divisor. The quotient is a
one-dimensional Cohen--Macaulay ring, generically
reduced with support $C$; its finite-support
nilradical vanishes. The quotient is consequently
$\mathcal O_C$.
If $R$ is the two-dimensional local ring of that
branch and $a$ the other equation, a lift $u$ of
a parameter of $R/(a)=\mathcal O_C$ generates
its maximal ideal together with $a$.
Thus $\operatorname{edim}R\le2=\dim R$,
so $R$ is regular. All three branches are smooth.
\item The pairwise quotient
$\mathcal O_Y/(f_i,f_j)=\mathcal O_C$ is regular
of dimension one. Its cotangent-space dimension
therefore makes $df_i,df_j$ independent along $C$.
Choose $x=f_1,y=f_2$ and a coordinate $z$ along $C$.
The third equation belongs to $(x,y)$, so it is
$f_3=a(x,y,z)x+b(x,y,z)y$.
Pairwise independence makes both $a$ and $b$
units. The map
$(x,y,z)\mapsto(X=ax,Y=-by,z)$
has invertible Jacobian along $C$ and satisfies
$f_3=X-Y$ and
$f_1f_2f_3=-(ab)^{-1}XY(X-Y)$.
It therefore carries the principal ideal of the
reduced union to $(XY(X-Y))$, giving the asserted
product germ.
Its three smooth normalization branches map
their conductor curves isomorphically to the
$z$-disc. Thus $N$ is smooth along the conductor,
which is a finite unramified cover of $C$ of
degree three. Away from the conductor, $N=S_{\rm reg}$.
Proposition~\ref{night:P4.22} now makes this
smooth elliptic ruled normalization minimal.
\end{enumerate}
\end{proof}


\subsection{The Hessian line and the tacnodal equality case}

\begin{lemma}[Rank-one tensors in a conormal lattice]\hypertarget{tag-P259}{}\label{lem:II:rank-one-conormal-lattice}
\StatementRecord{lem:II:rank-one-conormal-lattice}{proved}{P259}{}
Let $\mathcal O=\mathbb C\{t\}$, let $K=\mathcal Oe_1\oplus\mathcal Oe_2$,
and let $E=t^a\mathcal Oe_1\oplus t^b\mathcal Oe_2$ with $0\le a\le b$.
For a primitive vector $\ell=u e_1+v e_2$, put $L=\mathcal O\ell$,
\[
 \alpha=\operatorname{length}(L/(L\cap E)),\qquad
 \beta=\operatorname{length}((K/L)/\operatorname{im}E).
\]
Then $\alpha+\beta=a+b$ and $\beta\ge a$.
For a rank-one symmetric tensor $T=\varepsilon t^d\ell^2$, with
$\varepsilon$ a unit, membership in $\operatorname{Sym}^2E$ gives
$d/2+\beta\ge a+b$. The inclusion in $E$ of every column of $T$
gives $d/2+\beta\ge a+b/2$. The column condition is independent
of the chosen basis of $K$.
\end{lemma}
\begin{proof}
The inclusion and projection maps give the exact sequence
\[
 0\longrightarrow L/(L\cap E)\longrightarrow K/E
 \longrightarrow (K/L)/\operatorname{im}E\longrightarrow0,
\]
so $\alpha+\beta=a+b$. Since $E\subset t^aK$, its image in
the free rank-one quotient $K/L$ is contained in $t^a(K/L)$,
giving $\beta\ge a$. Directly, with the order of a zero coordinate
read as infinity,
\[
 \alpha=\max\{a-\operatorname{ord}u,b-\operatorname{ord}v\},\qquad
 \beta=\min\{a+\operatorname{ord}v,b+\operatorname{ord}u\}.
\]
Membership of $T$ in $\operatorname{Sym}^2E$ gives $d\ge2\alpha$
by its two diagonal coefficients; the mixed coefficient satisfies
the intermediate inequality. This proves the first bound.
For the second, one coordinate of $\ell$ is a unit, so its
corresponding column puts $t^d\ell$ in $E$. Thus $d\ge\alpha$,
and $d/2+\beta\ge(a+b+\beta)/2\ge a+b/2$.
Under a basis change the columns are invertible linear combinations
of the same homomorphism $K^*\to K$; their membership in $E$ is
the intrinsic inclusion of its image in $E$.
\end{proof}

\begin{proposition}[Euler characteristic of a tacnodal conductor]\hypertarget{tag-02Y3}{}\label{night:P4R.10}\label{stmt:P4R.10}
\StatementRecord{night:P4R.10}{proved}{02Y3}{\SeriesRef{lem:II:conductor-duality}, \SeriesRef{night:P4.30}}

Let $S$ be an irreducible reduced fibre support in the standing
torus-fibration setting of \cref{def:II:standing-torus-fibration}.
Suppose that its reduced singular curve $C$ is irreducible and
unibranch at every point, that its normalization $j:E\to C\hookrightarrow Y$
is elliptic, and that its generic transverse germ is a tacnode
$x^2-y^4=0$. Let $C_{\mathrm{sch}}$ be the image conductor scheme.
Then
\[
 \chi(\mathcal O_{C_{\mathrm{sch}}})\leq0.
\]
Equality makes $C$ smooth and makes the transverse Hessian of
the reduced defining equation have rank one at every point of $C$.
\end{proposition}
\begin{proof}

\begin{enumerate}[label=\arabic*.,leftmargin=*]
\item Let $P\subset\mathcal O_Y$ be the reduced ideal of $C$, and let
$I\subset\mathcal O_Y$ be the inverse image of the conductor ideal
of $S$. Apply \cref{lem:II:conductor-duality} to the finite normalization of this reduced hypersurface $S$. Its image conductor scheme has structure sheaf $\mathcal O_Y/I$, which the lemma makes Cohen--Macaulay of pure dimension one.
At the generic tacnode the normalization has two smooth branches;
its conductor is $(x,y^2)$ in the displayed coordinates.
Consequently $I$ contains $P^{(2)}$. Indeed a section of $P^{(2)}$
has zero image in $\mathcal O_Y/I$ at the generic point, and any
nonzero image would generate a finite-length submodule, contrary
to depth one. We obtain
\[
 0\longrightarrow J:=P/I\longrightarrow
 \mathcal O_{C_{\mathrm{sch}}}\longrightarrow\mathcal O_C
 \longrightarrow0.
\]
Here $J^2=0$, because $P^2\subset P^{(2)}\subset I$.
The module $J$ is therefore an $\mathcal O_C$-module. It is
torsionfree of rank one: a finite-support torsion submodule would
again be a finite-length submodule of $\mathcal O_Y/I$.

\item On $E$, put
\[
 K=\ker\bigl(j^*\Omega_Y\longrightarrow\Omega_E\bigr),
 \qquad\operatorname{im}(dj)=\Omega_E(-R).
\]
These are a rank-two bundle and a line bundle respectively.
The coefficient of $R$ at a point with local parametrization
$x_i=x_i(t)$ is $r=\min_i\operatorname{ord}_t x_i'(t)$.
We have $\deg K=\deg R$, since $\deg j^*\Omega_Y=0$
and $E$ is elliptic. Differentiation identifies
$G=P/P^{(2)}$ with a subsheaf of $j_*K$, as in \cref{night:P4.30}.
Write $E_0\subset K$ for the image of its pullback to $E$.
This is a rank-two locally free subsheaf with finite-length
quotient.

\item The Hessian of a local reduced equation $f$ of $S$ annihilates
the tangent derivative of $C$, since all first derivatives of
$f$ vanish on $C$. It thus defines a quadratic tensor in
$\operatorname{Sym}^2K$. Under a change $f\mapsto u f$ its
restriction to $C$ multiplies by $u$, so globally it is a section
\[
 q\in H^0\bigl(E,\operatorname{Sym}^2K\otimes j^*\mathcal O_Y(S)\bigr).
\]
It has generic rank one. Let $L\subset K$ be the saturation of
its generic image line. On the smooth curve $E$ this is a line
subbundle. The tensor lies in $L^{\otimes2}\otimes
j^*\mathcal O_Y(S)$ everywhere, since that is a subbundle and
the assertion holds on a dense open set. Write its zero orders
as $d_p\geq0$. The normal line of $S$ has degree zero on $E$:
locally over the base disc the divisor $mS$ is principal.
It follows that
\[
 2\deg L=\sum_p d_p.
\]

\item Generically $I/P^{(2)}$ is exactly the line of the linear
conductor equation $x$, which is the Hessian line. Projection
to $K/L$ therefore gives an injection
\[
 J\longrightarrow j_*(K/L).
\]
To check the kernel at the exceptional points, a section whose
projection vanishes generically has torsion image in $J$; the
torsionfreeness just proved makes that image zero. Let $b_p$ be
the colength of the projected lattice $E_0\to K/L$ at $p$.
Let $q_p$ be the additional nonnegative length between $J$ and
the direct image of this normalized image lattice. Riemann--Roch
on $E$ now gives the exact formula
\[
 \chi(J)=\deg R-\frac12\sum_p d_p-\sum_p b_p-\sum_p q_p.
\]
We prove locally that $d_p/2+b_p\geq r_p$.

\item At a smooth point of $C$, $r_p=0$, so the claim follows from
nonnegativity. At a singular point put
$m=\min_i\operatorname{ord}_t x_i(t)$, so $r_p=m-1$.
First suppose the embedding dimension is three. Every defining
equation of $C$ has no linear term. Its gradient along the
parametrization has all entries of order at least $m$.
Every two-by-two minor of the conormal image consequently has
order at least $2m$. Since $K$ is saturated in the ambient free
rank-three module, a generator of $\bigwedge^2K$ has a unit
Pl\"ucker coordinate. Thus
\[
 \tau_p:=\operatorname{length}(K/E_0)_p\geq2m.
\]
Apply the column clause and its exact sequence in
\cref{lem:II:rank-one-conormal-lattice} to $E=E_0\subset K$,
the same saturated Hessian line $L$, $d=d_p$ and $\beta=b_p$.
The Hessian columns are differentials of first derivatives of $f$,
which belong to $P$; their images consequently lie in $E_0$.
Writing $a_p=\operatorname{length}(L/(E_0\cap L))_p$, the lemma
gives $d_p\ge a_p$ and $a_p+b_p=\tau_p$. Hence
\[
 \frac{d_p}{2}+b_p\ge\frac{a_p}{2}+b_p
 \ge\frac{\tau_p}{2}\ge m>m-1=r_p.
\]

\item Now suppose the embedding dimension is two. Choose coordinates
with $P=(z,H(x,y))$, where $H$ is an irreducible reduced plane
equation with no linear term. Here $P^{(2)}=P^2$.
For an explicit verification, $z,H$ is a regular sequence, and
$P/P^2$ is free over the domain $\mathcal O_Y/P$.
If $s\notin P$ and $sv\in P^2$, reduction modulo $P$ gives
$v\in P$; reduction in this free conormal module then gives
$v\in P^2$. Thus $P^2$ is $P$-primary and equals its symbolic
square. The equation $f$ belongs to this square, so
\[
 f=A z^2+B zH+C H^2.
\]
Let $u$ be a primitive generator of the plane kernel of the
derivative row. We have $dH=w u$, with
$a=\operatorname{ord}_t w\geq m$, because both plane entries
of $dH$ have order at least $m$. In the basis $(dz,u)$ of $K$
we consequently have
\[
 E_0=\langle dz,w u\rangle,
 \qquad q=2A\,dz^2+2Bw\,dz\,u+2Cw^2u^2.
\]
The displayed tensor lies in $\operatorname{Sym}^2E_0$.
Apply the symmetric-square clause of
\cref{lem:II:rank-one-conormal-lattice} to this same lattice,
whose elementary-divisor orders are $(0,a)$, with $d=d_p$
and $\beta=b_p$. It gives
\[
 d_p/2+b_p\ge a\ge m>m-1=r_p.
\]
Together with the nonplanar calculation this proves the
claimed local inequality at every singular curve point.

\item We conclude that $\chi(J)\leq0$. Since normalization gives
$\chi(\mathcal O_C)=-\sum\delta(C,p)$, the conductor exact
sequence yields
\[
 \chi(\mathcal O_{C_{\mathrm{sch}}})
 =-\sum_p\delta(C,p)+\chi(J)\leq0.
\]
Equality forces every delta invariant to vanish, so $C=E$ is
smooth. Then $R=0$, $E_0=K$, and all normalization descent
defects vanish. The exact formula reduces to
$\chi(J)=-\tfrac12\sum d_p$, so equality also forces
$d_p=0$ everywhere. The Hessian has rank one at every point.
The local Hessian calculations can also be written explicitly.
Put $H=y^2-x^3$ and
$F=(\alpha z+\gamma H)^2+H^4$, with constants
$\alpha\ne0$ and $\gamma$. Along $(t^2,t^3,0)$,
the first-order part of $\alpha z+\gamma H$ gives
\[
 \operatorname{Hess}(F)=2\ell\ell^{\mathsf T},
 \qquad\ell=(-3\gamma t^4,\,2\gamma t^3,\,\alpha)^{\mathsf T}.
\]
This matrix has rank one. For $zx^2+xy^2$, restriction
of the Hessian to $x=y=0$ gives
$\operatorname{diag}(2z,0,0)$, whose quadratic coefficient
has a simple zero at the origin.

\end{enumerate}
\end{proof}

\begin{proposition}[Tacnodal supports of Euler number zero]\hypertarget{tag-02Y4}{}\label{night:P4R.11}\label{stmt:P4R.11}
\StatementRecord{night:P4R.11}{proved}{02Y4}{\cref{lem:II:zero-euler-ruled-conductor,night:P4R.10,night:P4.22}}

Let $S$ be an irreducible singular support in \cref{def:II:standing-torus-fibration} whose
normalization has ruled resolution, and suppose $e(S)=0$.
In the double-horizontal-divisor case with generic transverse
tacnode, the normalization is a smooth minimal elliptic ruled
surface. The conductor image is a smooth elliptic curve, its
reduced preimage is an unramified covering of degree two,
and every local germ is a product tacnode.
\end{proposition}
\begin{proof}

\begin{enumerate}[label=\arabic*.,leftmargin=*]
\item Write $\nu:N\to S$ and $\pi:T\to N$ for normalization and its minimal resolution. Proposition~\ref{night:P4.26} excludes rational resolution. Apply \cref{lem:II:zero-euler-ruled-conductor} to this same positive-genus ruling and double horizontal pole divisor. It supplies the sole strict conductor $H$ of genus $g\ge1$, the irreducible unibranch image $C$ with elliptic normalization, and the descended ruling $r_N:N\to B_0$ with section $s=\pi h$, where $h:B_0\simeq H\hookrightarrow T$. Its conductor-scheme equality is
\[
 \chi(\mathcal O_{C_{\rm sch}})=g-1.
\]
Apply \cref{night:P4R.10} to this same $S\subset Y$ and $C$, whose generic transverse germ is the stated tacnode. It gives $g-1\le0$. Thus $g=1$ and equality holds; the same proposition makes $C$ smooth and gives rank one for the reduced equation's transverse Hessian everywhere on $C$. The finite map $H\to C$ has degree two, the generic number of normalization branches. Riemann--Hurwitz between the two elliptic curves makes it unramified.

\item We verify the local analytic conclusion. Over any point of
$C$, its two preimages on $H$ cannot be identified in $N$:
each exceptional block is vertical for the descended ruling
and hence meets only one point of the section $H$.
There are therefore two local irreducible branches of $S$.
The holomorphic Morse lemma with parameters writes its reduced
equation as $X^2+g(y,z)$. Its two analytic factors are monic
linear polynomials in $X$, by Weierstrass factorization, and
their coefficients sum to zero. Each branch is a graph and
is smooth. Their intersection has support exactly the smooth
curve $C$. Consequently, in coordinates with $C=\{X=y=0\}$,
the factor coefficient is a unit times a power of $y$.
The generic tacnode fixes that power to two. Extracting a square
root of the unit in the $y$ coordinate gives
$X^2-y^4=0$ times the $z$ disc. This also verifies the
normalization and conductor morphisms locally.
The normalization was already smooth away from the conductor
by \cref{night:P4R.6}. It is therefore smooth everywhere, and \cref{night:P4.22} makes
this elliptic ruled surface minimal.

\end{enumerate}
\end{proof}

\subsection{Zero-Euler fibre supports}

\begin{proposition}[The zero-Euler fibre supports]\hypertarget{tag-02X5}{}\label{night:P4.32}\label{stmt:P4.32}
\StatementRecord{night:P4.32}{proved}{02X5}{\cref{lem:II:zero-euler-ruled-conductor,night:P4.26,night:P4.29,night:P4R.11,night:P4.31,night:P4.37,night:P4.36}}

Let $f:Y\to B$ satisfy the torus-fibration hypotheses of
\cref{night:P4.12}. An irreducible reduced fibre support has
Euler number zero exactly in the following cases: it is smooth with a finite unramified
torus cover, or it is singular with smooth minimal elliptic ruled
normalization, or it is singular with smooth minimal Hopf
normalization. In the ruled double-horizontal case the conductor
image is a smooth elliptic curve; the local germs are product
cusps, product tacnodes, or product ordinary plane triple points,
and the reduced conductor maps have degrees one, two, or three
respectively and are unramified.

For a fibration whose supports are all irreducible or normal
crossing, $c_3(Y)=0$ exactly when every irreducible support has
one of these forms and every reducible support has the
zero-Euler normal-crossing form of \cref{night:P4.8}.
\end{proposition}
\begin{proof}
\begin{enumerate}[label=\arabic*.,leftmargin=*]
\item Proposition~\ref{night:P4.12} makes every normal irreducible support smooth. Apply the cyclic base-change construction in the proof of \cref{night:P4.8} to this full fibre $mS$. Its local map is $t=uz^m=w^m$ with $w=u^{1/m}z$, followed by $t=s^m$ and the normalized charts $w=\zeta s$. The resulting finite unramified central cover $S'\to S$ is a torus by the torus-deformation theorem applied there to this proper submersion. Thus $e(S)=e(S')/\deg(S'\to S)=0$. This equality also proves the converse for a support with any finite unramified torus cover.

For singular $S$, write $\nu:N\to S$ and $\pi:T\to N$ for normalization and its minimal resolution. Proposition~\ref{night:P4.12} gives $\kappa(T)=-\infty$, so \cref{lem:II:negative-kodaira-models} applies to a minimal model. Proposition~\ref{night:P4.26} excludes a rational resolution from $e(S)=0$. On a class-VII resolution, \cref{night:P4.37}, applied to these same $\nu,\pi$ and full fibre, makes $N$ a smooth minimal Hopf surface. Conversely \cref{night:P4.36} gives $e(S)=0$ on every such Hopf normalization.

\item The remaining resolution is ruled over a curve of genus $g\ge1$. Let $D$ be its pole divisor and $D_0$ its pushforward to the minimal ruling. Apply \cref{lem:II:irrational-ruled-boundary} to these same data. Its reduced-horizontal alternative has $g=1$, and \cref{night:P4R.7} makes $N$ smooth and minimal with $e(S)=0$.

In the double-horizontal alternative, apply \cref{lem:II:zero-euler-ruled-conductor} to the identical $mS,\nu,\pi,D$ and ruling. It gives the irreducible image curve $C$ with elliptic normalization, the noncontracted section $H$ as the sole strict conductor, and the complete list of generic transverse cusp, tacnode and ordinary-triple germs, with degrees one, two and three. Every exceptional component remains included in the resolved pole divisor used by that lemma.

\item Propositions~\ref{night:P4.29}, \ref{night:P4R.11} and~\ref{night:P4.31} apply, respectively, to these three transverse germs on this same support. They make $N$ smooth minimal elliptic ruled, $C$ smooth elliptic, and every local germ a product with a disc. Their degree-one, degree-two and degree-three reduced conductor maps are respectively an isomorphism and the two stated unramified coverings. Conversely \cref{night:P4.22}, on any smooth minimal elliptic ruled normalization in the original setting, gives $e(S)=0$. This proves both directions of the irreducible classification.

\item The Euler-additivity calculation in the proof of \cref{night:P4.26}, for this same $f:Y\to B$, writes $c_3(Y)=e(Y)$ as the finite sum of exceptional reduced-support Euler numbers. Propositions~\ref{night:P4.26} and~\ref{night:P4.6} make every summand nonnegative. The sum is zero precisely when each summand is zero. The irreducible summands have just been classified. For a reducible normal-crossing support, apply the single-support component calculation in the proof of \cref{night:P4.8} to that divisor in this same $Y$: zero triple points give minimal Hopf or non-K\"ahler elliptic normalization components with smooth elliptic double curves and unramified branch maps. Conversely this form has Euler number zero by \cref{night:P4.6}. This proves the final assertion.
\end{enumerate}
\end{proof}


\subsection{The components outside the pole divisor in a ruled fibre}

\begin{proposition}[Components outside the pole divisor of a ruled fibre]\hypertarget{tag-02Y5}{}\label{night:P4R.12}\label{stmt:P4R.12}
\StatementRecord{night:P4R.12}{proved}{02Y5}{\cref{lem:II:irrational-ruled-boundary,night:P4R.6,night:P4R.7}}

Let $T\to N\to S$ be the minimal normalization resolution of an
irreducible support in \cref{def:II:standing-torus-fibration}, with ruled minimal model of positive
base genus. Let $r:T\to B_0$ be its ruling, and let $D$ be the
effective integral pole divisor of \cref{night:P4R.6}. Every component $E$ of a
reducible ruling fibre which is absent from $\operatorname{Supp}D$
is a smooth rational curve of self-intersection $-1$. It meets
exactly one pole component, transversely once, and that component
has pole coefficient one. In the double-horizontal-section case,
$E$ is disjoint from the horizontal section. Distinct such
components are disjoint; each is a leaf of the reduced ruling fibre.
The full ruling-fibre multiplicities are arbitrary positive integers.
\end{proposition}
\begin{proof}

\begin{enumerate}[label=\arabic*.,leftmargin=*]
\item Write $\beta:T\to T_0$ for the blowdowns to the minimal ruling and $D_0=\beta_*D$. Proposition~\ref{night:P4R.6} gives $K_T+D\equiv0$. Apply \cref{lem:II:irrational-ruled-boundary} to this same $\beta$ and $D$. In its reduced-horizontal case, \cref{night:P4R.7} makes $N$ smooth and minimal; its minimal resolution $T\to N$ is then an isomorphism. Every ruling fibre is irreducible. The stated reducible fibre therefore belongs to the double-horizontal-section case.

A ruling fibre on $T$ is obtained from a smooth $\mathbb P^1$
by the point blowups inverse to $\beta$. Its reduced support is
a normal-crossing tree of smooth rational curves. Each component
of a reducible fibre has negative self-intersection: for the full fibre $F_b=\sum m_iE_i$, with every $m_i>0$, we have
$0=F_bE_i=m_iE_i^2+\sum_{j\ne i}m_jE_jE_i$;
connectedness and reducibility make the last sum positive.

\item For $E$ outside the pole support, effectiveness and adjunction give
\[
 0\le DE=-K_TE=2+E^2.
\]
Thus $E^2$ is $-1$ or $-2$. In the second case $DE=0$,
so $E$ is disjoint from the full pole support. Proposition~\ref{night:P4R.6} identifies a neighbourhood of $E$ under $T\to S$ with a neighbourhood in $S_{\rm reg}$. Write the original full fibre as $mS$ and a base parameter as $t$. On this regular locus the full function has the form $t=uz^m$, where $u$ is a unit. Choose a local root and set $w=u^{1/m}z$; the inverse is the unique local solution of this equation. Normalize the base change $t=s^m$. Its charts are $w=\zeta s$, with $\zeta^m=1$; the projection to the original neighbourhood is unramified and its map to $s$ is a submersion. These charts glue to a finite unramified cover near $E$.

The restriction of this cover to $E\simeq\mathbb P^1$ is a disjoint union of copies of $E$, since $\pi_1(\mathbb P^1)=0$. Fix one copy $E'$. The unramified projection identifies its normal bundle in the central surface with that of $E$, of degree $-2$. Lift a constant real base vector field near this compact $E'$ through the submersion. Its flow is defined uniformly for sufficiently small time on a neighbourhood of $E'$ and gives an orientation-preserving identification with a neighbourhood in a nonzero fibre. That fibre is the original torus $Y_{s^m}$, because the normalized base change over $s\ne0$ has that same fibre. The image of $E'$ is consequently a sphere of square $-2$ in this torus. Its map lifts to the contractible universal cover $\mathbb R^4$, so it is null-homotopic and has zero homology class and square. This contradiction gives $E^2=-1$ and $DE=1$.

\item All pole coefficients and all intersection multiplicities are
positive integers at an intersection. Hence $DE=1$ says exactly
that there is one intersection, transverse and on a coefficient-one
component. A horizontal pole section of coefficient two cannot
be that component.

\item It remains to exclude an intersection of two components $E,E'$
outside the pole divisor. Both have square $-1$, and both meet
the pole support. In the full ruling fibre, let their positive
multiplicities be $m,m'$. The two fibre-intersection equations
would give
\[
 m=m'+a,\qquad m'=m+b,
\]
where $a,b$ are positive sums of the multiplicities of their
other vertical neighbours; their respective pole neighbours make both
sums positive. Adding these equations gives $a+b=0$, a
contradiction. Thus no two outside components meet. Each has
exactly the one pole neighbour already found and is a leaf.

\end{enumerate}
\end{proof}

\subsection{Primary restriction and elliptic conductor components}

\begin{lemma}[Restriction to the horizontal primary conductor]
\hypertarget{tag-P313}{}\label{lem:II:horizontal-primary-conductor}
\StatementRecord{lem:II:horizontal-primary-conductor}{proved}{P313}{\cref{night:P4R.6,night:P4.33}; \cite[Remark~2.1, p.~415]{Greb2011}}
In the standing torus-fibration setting of
\cref{def:II:standing-torus-fibration}, let $mS$ have irreducible,
reduced, singular support. Write $\nu:N\to S$ for normalization,
$\sigma:T\to N$ for its minimal resolution and $\rho=\nu\sigma$.
Suppose that $T$ is obtained by point blowups of a minimal ruled
surface over a curve of genus $g\ge1$ and that the pole divisor
$D$ of \cref{night:P4R.6} has horizontal part $2H$ for the induced
ruling, where the section $H$ is not contracted by $\sigma$.
Let $C_H=\rho(H)$ and let $C_{H,\rm sch}$ be the primary component
of the image conductor scheme $C_{\rm sch}$ supported on $C_H$.
For the conductor ideal $I_N\subset\mathcal O_N$, the residue and
restriction maps give injections of coherent $\mathcal O_S$-algebras
\begin{equation}
 \mathcal O_{C_{\rm sch}}\longrightarrow
 \nu_*(\mathcal O_N/I_N)\longrightarrow\rho_*\mathcal O_D.
 \label{eq:II:elliptic-conductor-algebra-injection}
\end{equation}
Under each point blowup $\beta:T'\to T''$ in the contraction to a
minimal ruled model, with successive pole divisors $D',D''$, one has
\begin{equation}
 \beta_*\mathcal O_{D'}=\mathcal O_{D''},
 \qquad R^1\beta_*\mathcal O_{D'}=0.
 \label{eq:II:elliptic-conductor-blowup-pushforward}
\end{equation}
The primary restriction is injective on $H^0$ and surjective on
$H^1$. In particular,
\begin{equation}
 g-1=\chi(\mathcal O_{C_{\rm sch}})
 \le\chi(\mathcal O_{C_{H,\rm sch}}).
 \label{eq:II:elliptic-primary-Euler-comparison}
\end{equation}
\end{lemma}
\begin{proof}
\begin{enumerate}[label=\arabic*.,leftmargin=*]
\item Apply \cref{lem:II:irrational-ruled-boundary} to this ruling,
blowup sequence and divisor $D$. Its double-section case gives
the rational vertical trees required by \cref{night:P4.33}.
The noncontracted section therefore gives rational singularities
of $N$ and $h^1(N,\mathcal O_N)=g$. Put
$L_N=\nu^*\omega_S$. Finite duality in that proposition identifies
$I_N=\omega_N\otimes L_N^{-1}$. The proper resolution $\sigma$
has $R^1\sigma_*\mathcal O_T=0$ and fibres of dimension at most
one; $N$ is a normal Cohen--Macaulay surface. Apply
\cite[Remark~2.1, p.~415]{Greb2011} to this $N$ and $\sigma$.
Its canonical-sheaf equality is the trace isomorphism
$\sigma_*\omega_T=\omega_N$. The pole identity and projection
formula give
\[
 \sigma_*\mathcal O_T(-D)
 =\sigma_*(\omega_T\otimes\sigma^*L_N^{-1})
 =I_N.
\]
These are the same ideal subsheaves of $\mathcal O_N$: the pole
identity uses the pullback of the residue section defining finite
duality, and the trace is its identity on the birational
isomorphism locus. This agreement determines the maps between
the torsion-free sheaves. The divisor ideal sequence gives
$\mathcal O_N/I_N\hookrightarrow\sigma_*\mathcal O_D$; composing
with the injection from the finite conductor square gives
\eqref{eq:II:elliptic-conductor-algebra-injection}.

\item At a point blowup the canonical-boundary identity is
$D'=\beta^*D''-E_\beta$. Its structure-sheaf direct images are
$\beta_*\mathcal O=\mathcal O$ and $R^1\beta_*\mathcal O=0$.
Moreover,
\[
 \beta_*\mathcal O(E_\beta)=\mathcal O,
 \qquad R^1\beta_*\mathcal O(E_\beta)=0.
\]
For the first equality, a meromorphic function with a pole only
on $E_\beta$ descends off the centre and extends across it by
Hartogs. The exact sequence
$0\to\mathcal O\to\mathcal O(E_\beta)
\to\mathcal O_{E_\beta}(-1)\to0$, together with
$H^0(\mathbb P^1,\mathcal O(-1))=
H^1(\mathbb P^1,\mathcal O(-1))=0$, gives the second equality.
The structure-sheaf calculation follows on the two blowup charts,
or from the exceptional filtration with successive
$\mathcal O_{\mathbb P^1}(j)$, $j\ge0$. Projection by
$\mathcal O(-D'')$ and the two divisor ideal sequences now give
\eqref{eq:II:elliptic-conductor-blowup-pushforward}, with its
algebra map induced by pullback and quotient. Higher direct
images in degree two vanish because the fibres have dimension
at most one.

\item Thus $H^0(D,\mathcal O_D)=H^0(D_0,\mathcal O_{D_0})$,
where on the minimal ruled surface
$D_0=2H_0+V_0$ and $V_0=\sum_j a_jF_j$.
The kernel of restriction to $2H_0$ is
$\mathcal O_{V_0}(-2H_0)$. Its filtration on each thick fibre
$a_jF_j$ has successive layers $\mathcal O_{\mathbb P^1}(-2)$,
since $F_j^2=0$ and $H_0F_j=1$. It has no global section.
The Cartier curve $2H_0$ is Cohen--Macaulay and has no
zero-dimensional associated point. A global function on $D$
which vanishes at the generic doubled section is therefore zero.

\item Proposition~\ref{night:P4.33} makes the image conductor
Cohen--Macaulay of pure dimension one. For
$R=\mathcal O_{S,p}$ and its conductor ideal $I$, take every
minimal prime $\mathfrak p_i$ over $I$ corresponding to a local branch
of the global component $C_H$. Its component ideal is
\[
 J_H=\bigcap_i\bigl(IR_{\mathfrak p_i}\cap R\bigr).
\]
The isolated primary components are uniquely determined, so
these intersections agree on overlaps and define
$C_{H,\rm sch}$, including its multibranch points.
Restriction gives a
surjection $\mathcal O_{C_{\rm sch}}\to
\mathcal O_{C_{H,\rm sch}}$. A global section of its kernel maps
by \eqref{eq:II:elliptic-conductor-algebra-injection} to a function
on $D$ vanishing at the generic doubled section, so is zero by
the preceding item. The kernel sheaf has support of dimension
at most one and hence has zero second cohomology. The restriction
is therefore injective on $H^0$ and surjective on $H^1$.
Its Euler inequality, together with
$\chi(\mathcal O_{C_{\rm sch}})=g-1$ from
\cref{night:P4.33}, proves
\eqref{eq:II:elliptic-primary-Euler-comparison}.
\end{enumerate}
\end{proof}

\begin{proposition}[Elliptic components of the conductor]\hypertarget{tag-02Y6}{}\label{night:P4R.13}\label{stmt:P4R.13}\label{night:P4R.13:elliptic-components}
\StatementRecord{night:P4R.13}{proved}{02Y6}{}

In the standing torus-fibration setting of
\cref{def:II:standing-torus-fibration}, let $mS$ be a fibre
with irreducible, reduced, singular support. Suppose that the
normalization of $S$ has a ruled resolution and that the
reduced image conductor has a component with elliptic
normalization. Then the normalization of $S$ is a smooth
minimal elliptic ruled surface, the image conductor is a
disjoint union of one or two smooth elliptic curves, and
$e(S)=0$. Along each of these curves every local germ is the
product of its generic transverse germ with a disc. No
rational conductor component is attached to them.

When there are two image components, the normalization
conductor consists of two disjoint sections of the minimal
ruling. Each section maps to its own image curve by an
unramified double cover, and every transverse germ is $A_1$.
\end{proposition}
\begin{proof}

Let $\rho_N:T\to N$ be the minimal normalization resolution,
let $\nu:N\to S$ be the finite normalization, and put
$\rho=\nu\circ\rho_N$. Let $D$ be the divisor of
\cref{night:P4R.6}, so that
$\omega_T(D)\simeq\rho^*\omega_S$ and
$\operatorname{Supp}D=\rho^{-1}(\operatorname{Sing}S)_{\rm red}$.
Write $\beta:T\to T_0$ for a sequence of point blowdowns to a
minimal rational or ruled surface and put $D_0=\beta_*D$.
The canonical divisor identity gives $K_{T_0}+D_0\equiv0$.

\begin{enumerate}[label=\arabic*.,leftmargin=*]
\item Suppose first that $T$ is rational. A component $H$ of $D$
dominates the elliptically normalized image component in the
hypothesis. Its normalization has positive genus, so $\beta$
does not contract it. On a Hirzebruch minimal model the equality
$D_0F=2$ forces its image $H_0$ to have coefficient one and
ruling degree two: a curve of ruling degree one is a section
and is smooth rational. Adjunction gives
\[
 0\le 2g(\widetilde H_0)-2
 \le 2p_a(H_0)-2
 =-(D_0-H_0)H_0\le0.
\]
Thus $H_0$ is smooth elliptic and disjoint from the remaining
divisor. Every remaining component would be a fibre meeting
$H_0$ twice, so $D_0=H_0$. On a minimal model $\mathbb P^2$,
the divisor has degree three; a component with positive
normalized genus must be its whole smooth cubic. Again
$D_0=H_0$.

Under a reverse blowup the new pole coefficient is $\mu-1$.
Effectiveness excludes a centre off the smooth coefficient-one
curve, and a centre on it introduces coefficient zero.
Induction gives $D=H$, with $H$ smooth elliptic. Every
$\rho_N$-exceptional curve belongs to $D$, while $H$ dominates
the conductor image. Hence $\rho_N$ has no exceptional curve;
the resulting finite birational map to normal $N$ is an
isomorphism. In particular $N=T$ is smooth rational and its
whole image conductor is the irreducible curve $C_0=\nu(H)$.
Write $n:E\to C_0$ for its elliptic normalization and
$\eta:H\to E$ for the factor of $\nu|_H$.

Apply \cref{lem:II:plane-branch-conductor} to the reduced
generic transverse plane equation of $S$. Exponent one on
every sheet gives two smooth transverse branches, so
$\deg\eta=2$.
Riemann--Hurwitz between the two elliptic curves makes $\eta$
unramified. Finite conductor duality on smooth $N$ identifies
the entire normalization conductor ideal with
$\omega_N\otimes\nu^*\omega_S^{-1}=\mathcal O_N(-H)$ as an
ideal subsheaf. This equality determines its scheme
structure at every point.

We identify the image conductor at every point. If $p\in C_0$
has $r$ inverse images in $E$, the unramified double cover has
exactly $2r$ points over $p$ in $H$. These are all normalization
points over $p$, by the support equality for $D$, and correspond
to the local irreducible hypersurface branches of $S$.
Indeed a split finite normalization fibre of an irreducible
analytic germ would split its finite inverse image after
shrinking and give a nontrivial idempotent in its normalization
field. Each irreducible analytic branch therefore has one
normalization point.
Each such normalization germ has one smooth conductor germ;
denote its irreducible image curve germ by $B_i$.
Any two distinct surface factors intersect in a curve germ:
their ideal has height two in the regular three-dimensional
ambient local ring. That intersection lies in the singular
locus, and its inverse image in either normalization germ lies
in its single conductor curve. Finiteness then gives
$B_i=B_j$. These images exhaust the branches of $C_0$, whence
$r=1$ and there are exactly two surface branches at $p$.

Each branch is regular in codimension one, because the generic
double curve is a node and the complement of $C_0$ is smooth.
It is a hypersurface, hence Cohen--Macaulay; Serre's criterion
\cite[Tag~031S, Lemma~10.157.4]{StacksProject} makes it normal.
Its normalization is smooth, so it is itself smooth. If its
two local equations are $f_1,f_2$ in $R=\mathcal O_{Y,p}$,
factoriality of this local analytic ring
\cite[Chapter~II, Theorem~2.10, pp.~82--83]{Demailly2012E}
gives the exact difference-of-restrictions sequence
\[
0\longrightarrow R/(f_1f_2)
 \longrightarrow R/(f_1)\oplus R/(f_2)
 \longrightarrow R/(f_1,f_2)\longrightarrow0.
\]
Multiplication by the two normalization idempotents identifies
the conductor on the first branch with $(f_2|_{f_1=0})$,
and on the second with $(f_1|_{f_2=0})$. These are the reduced
Cartier ideals of their smooth conductor curves. Thus
$R/(f_1,f_2)$ is the smooth reduced conductor ring, and
$df_1,df_2$ are independent. Completing them by a curve
coordinate $z$ gives the local biholomorphism
$(X,Y,z)=(f_1,f_2,z)$ and the equation $XY=0$.
This also proves globally that the image conductor scheme
is the smooth elliptic curve $C_0=E$ and that its normalization
conductor is exactly $H$.

The full conductor square is therefore
\[
0\longrightarrow\mathcal O_S
 \longrightarrow\nu_*\mathcal O_N\oplus\mathcal O_{C_0}
 \xrightarrow{\ (a,b)\mapsto a|_H-\eta^*b\ }
 \eta_*\mathcal O_H\longrightarrow0.
\]
The last map on global functions is $(a,b)\mapsto a-b$.
The rational surface has $H^1(\mathcal O_N)=H^2(\mathcal O_N)=0$,
and pullback $H^1(\mathcal O_{C_0})\to H^1(\mathcal O_H)$ is an
isomorphism, since trace composed with pullback is multiplication
by two. Hence $h^1(\mathcal O_S)=0$; equivalently Euler
additivity gives $\chi(\mathcal O_S)=1$. Both contradict
\cref{night:P4.12}. The ruling base consequently has positive
genus.

\item Suppose the horizontal part of $D_0$ is reduced.
Apply \cref{lem:II:irrational-ruled-boundary} to the same
minimal ruled model $T_0$, blowdown $\beta$ and divisor $D$.
It gives genus one, one elliptic bisection or two disjoint
elliptic sections, all with coefficient one, and no vertical
pole component. Proposition~\ref{night:P4R.7}, with this
normalization resolution, gives $N=T$ smooth and minimal
and $e(S)=0$.

The coefficient-one conductor gives generic transverse nodes,
so the differential of $N\to S\hookrightarrow Y$ has rank two
away from finitely many points. Proposition~\ref{night:P4.15}
gives minimality, $e(S)=0$, and rank-drop length zero. It also
gives a disjoint union of smooth elliptic image curves and an
unramified finite map from the reduced conductor. Its degree
over every image component is two, as at a generic node.
Thus every image point has precisely two normalization points.
The immersive normalization maps them to smooth surface germs;
the same conductor-square calculation as in the first item
makes their intersection reduced and smooth and gives the
coordinates $(f_1,f_2,z)$ with equation $XY=0$.

There are exactly three possible component maps. One connected
bisection maps to one image curve by an unramified double cover.
Two sections with the same image each map isomorphically to
that curve, since their positive degrees add to two. Two
sections with different images each map by an unramified
double cover to its own image curve. Thus there are one or
two image components; all the asserted conclusions hold
in this case.
\item In the remaining case the minimal-model divisor is
$D_0=2H_0+V_0$, where $H_0$ is a smooth section and $V_0$
is vertical. The pole divisor $D$ of \cref{night:P4R.6}
has the doubled-horizontal form. Its support is a smooth section $H$
with rational trees attached. The pushout and curve-group
calculation of \cref{night:P4R.7} show that the reduced image
conductor $C$ is connected, that its point-identification
graph has no cycles, and that it has exactly one normalized
component of positive genus, of genus one. Denote the
corresponding reduced component by $C_E$ and its normalization
by $j:E\to C_E\hookrightarrow Y$. In particular every point
of $C_E$ is unibranch. Every other normalized component of
$C$ is rational.

\item A rational component of $D$ cannot map nonconstantly to $E$.
Thus $H$ maps nonconstantly to $E$. Write $g=g(H)\ge1$.
At a general point of $C_E$, every normalization branch has
conductor exponent two, because the coefficient of $H$ in
$D$ is two. The formula in \cref{lem:II:plane-branch-conductor}, applied
to this reduced transverse equation, leaves a single branch
with $\delta=1$, two smooth branches with intersection
multiplicity two, or three smooth branches with pairwise
intersection multiplicity one. The first ring is
$\mathbb C\{t^2,t^3\}$ by the exponent-two enumeration in
the proof of \cref{night:P4.32}. These are $A_2,A_3,D_4$, respectively, and $H\to E$
has degree $1,2,3$ in the three cases.

\item Apply \cref{lem:II:horizontal-primary-conductor} with
$\sigma=\rho_N$, the pole divisor $D$ above, and $C_H=C_E$.
The preceding item makes $H$ noncontracted; the ruling has
positive genus and horizontal pole part $2H$. Thus this is the
same conductor ideal $I_N$ and the same primary quotient
$C_{E,\rm sch}=C_{H,\rm sch}$. The restriction map gives
\eqref{eq:II:elliptic-primary-Euler-comparison}.

\item We evaluate the primary quotient on the right side;
the left side is the Euler characteristic of the
entire conductor scheme. In type $A_2$ the primary conductor
is reduced: a pure one-dimensional quotient which
is generically reduced has no nilpotent submodule
of finite support. Thus its Euler characteristic
is $-\sum_p\delta(C_E,p)\le0$. In type $D_4$ the
primary ideal is the symbolic square of the ideal
of $C_E$, by exactly the generic-conductor and
depth-one argument of \cref{night:P4.31}. \Cref{night:P4.30} gives the same
nonpositive bound, with equality only for smooth
$C_E$.

\item In type $A_3$, the primary ideal contains that
symbolic square and its nilpotent module is
torsionfree of rank one on $C_E$. The proof of
\cref{night:P4R.10} applies to this primary quotient: its
differentiation map, saturated Hessian line and
projection to the quotient conormal line depend
only on this ideal and on the equation of
$S$. At intersections with the other conductor components,
these maps are the restrictions of the same differentiation
map, saturated Hessian line and quotient projection. More explicitly its
Euler characteristic is
\[
 -\sum_p\delta(C_E,p)+\deg R
 -\frac12\sum_p d_p-\sum_p b_p-\sum_p q_p\le0,
\]
with the same ramification divisor $R$, Hessian
zero orders $d_p$, projected lattice colengths
$b_p$ and descent defects $q_p$ as in \cref{night:P4R.10}.
The local inequalities there prove this bound
at every singular point of $C_E$. Equality
makes $C_E$ smooth and its Hessian nowhere zero.

\item On the other hand \cref{night:P4.33} gives
$\chi(\mathcal O_{C_{\rm sch}})=g-1$. Thus
\begin{equation}
 0\le g-1
 =\chi(\mathcal O_{C_{\rm sch}})
 \le\chi(\mathcal O_{C_{E,\rm sch}})\le0.                    \label{eq:II:elliptic-primary-zero-Euler-chain}
\end{equation}
All these quantities vanish. We obtain $g=1$,
$C_E=E$ smooth, and an unramified map $H\to E$.
In type $A_3$ its transverse Hessian has rank
one everywhere on $E$.

\item It remains to exclude attachments to this smooth
elliptic component. Put
\[
 K=\mathcal I_E/\mathcal I_E^2,
 \qquad N_E=\mathcal O_Y(S)|_E.
\]
Both $\deg K$ and $\deg N_E$ are zero. The first
follows from the conormal exact sequence,
$c_1(TY)E=0$, and $\deg\Omega_E=0$; the second
follows because a positive multiple of $S$ is
principal in a tube of this fibre.

\item In type $A_2$ let $L\subset K$ be the saturated
quadratic-image line and put $Q=K/L$. The nonzero
quadratic tensor is a section of
$L^2\otimes N_E$, so $2\deg L\ge0$.
The cubic jet on the null direction is an
intrinsic section of $Q^3\otimes N_E$.
Here is the coordinate check, including points
where the quadratic tensor vanishes. A change
of a second-order lift changes the cubic
coefficient by the polarization of the
quadratic tensor. That polarization vanishes
on the null direction. A change of tangential
lift contributes a tangential derivative of
the rank-one quadratic tensor; its evaluation
twice on the null direction is again zero.
Changing a reduced defining equation multiplies
the tensor by its transition unit. Thus the
stated line bundle and its holomorphic section
are defined on all of $E$. The section is
nonzero generically by type $A_2$, so
$-3\deg L\ge0$. Both degrees are zero,
and both sections are nowhere zero. The parametrized Morse lemma gives
$F=a(x,y,z)(x-\chi(y,z))^2+h(y,z)$, with $a$ a unit
and $\chi$ solving $F_x=0$. Here $h=y^3b(y,z)$ with
$b$ a unit. The coordinates
$X=(x-\chi)\sqrt a$ and $Y=y\sqrt[3]b$, together with $z$,
give the product equation $X^2+Y^3=0$ at every point of $E$.

\item In type $A_3$ \eqref{eq:II:elliptic-primary-zero-Euler-chain} already makes the
quadratic tensor nowhere zero, whence
$\deg L=\deg Q=0$. Its cubic null jet is
identically zero. Removing the quadratic
direction gives an intrinsic residual fourth
jet in $Q^4\otimes N_E$. In local coordinates
its relevant terms and coefficient are
\begin{equation}
 a x^2+bxy^2+cy^4,
 \qquad q_4=c-\frac{b^2}{4a},\qquad a\ne0.                 \label{eq:II:elliptic-tacnode-fourth-jet}
\end{equation}
Substitution of the critical coordinate
$x=-by^2/(2a)+O(y^3)$ proves the formula.
The change $x\mapsto x+vy^2$ sends $b$ to
$b+2av$ and $c$ to $c+bv+av^2$, leaving
$q_4$ unchanged. Changes of normal frames
and of equation give the claimed
$Q^4\otimes N_E$ transition. Tangential
lift changes do not affect a function
whose residual order is four. This section
is nonzero generically and has degree zero;
it never vanishes. In the same Morse coordinates the
residual function is $h=y^4b(y,z)$ with $b$ a unit.
The coordinate $Y=y\sqrt[4]b$, together with the quadratic
coordinate $X$ and $z$, gives the product tacnode
$X^2+Y^4=0$ at every point of $E$.

\item In type $D_4$, the equation belongs to
$\mathcal I_E^3$. This ideal is primary
because $E$ is smooth, so generic order
three proves membership everywhere.
The cubic normal jet belongs to
$\operatorname{Sym}^3K\otimes N_E$.
Its discriminant is a section of
$(\det K)^6\otimes N_E^4$, of degree
zero. It is nonzero generically, hence
nowhere zero. The three distinct normal
tangent directions give three smooth
holomorphic branches locally along $E$.
After straightening the first two to
$x=0,y=0$, the third is $ax+by=0$
with both coefficients units. Setting
$X=ax,Y=-by$ gives the product equation
$XY(X-Y)=0$.

\item In all three cases the entire singular
locus near every point of $E$ is $E$
itself. A rational conductor tree could
not meet it. Connectedness of $C$ then
forces $C=E$. The displayed local
product forms make $N$ smooth along
its conductor, and \cref{night:P4R.6} already makes
it smooth elsewhere. Hence $T=N$.
A finite normalization cannot contract
a curve, and a rational curve cannot
map nonconstantly to $E$. Thus $D=2H$
has no vertical component. The first
blowup of $2H_0$ in a nonempty
blowup sequence would create a vertical
pole component of coefficient one;
later blowups never remove its strict
transform from the pole support.
Consequently no blowup occurs and $N$
is minimal elliptic ruled. Finally
\[
 e(S)=e(N)-e(H)+e(E)=0.
\]

\end{enumerate}
\end{proof}

\subsection{Three doubled nodal supports}\label{subsec:II:three-doubled-nodal-supports}

\begin{example}[A compact fibration with three doubled nodal supports]
\hypertarget{tag-P205}{}\label{ii:three-doubled-nodal-fibres}
\StatementRecord{ii:three-doubled-nodal-fibres}{proved}{P205}{\cite[\S7, Satz~3, p.~62, and S\"atze~5--6, p.~64]{Grauert1960D}; \cite[pp.~35--36]{Grauert1963ErratumD}}
There is a connected smooth compact complex threefold $Y$ and a proper
surjective holomorphic map $f:Y\to\mathbb P^1$ with connected fibres such that
\[
 b_1(Y)=b_2(Y)=0,\qquad e(Y)=0,\qquad b_3(Y)=2.
\]
Every regular fibre is a complex two-torus. The exceptional values are
$\infty,0,1$, and each full exceptional fibre is $2S$, where $S$ is reduced
and irreducible. Its normalization is a smooth minimal elliptic ruled
surface $\mathbb P_C(\mathcal O_C\oplus L)$ with $L$ a nontrivial line
bundle of order two. The normalization conductor consists of two disjoint
sections. Each maps by an unramified degree-two isogeny onto a different
smooth elliptic image curve. Every germ along either image curve is the
product of a transverse node with a disc.
\end{example}
\begin{proof}
\begin{enumerate}[label=\arabic*.,leftmargin=*]
\item Let
$B^\circ=\mathbb P^1\setminus\{\infty,0,1\}$. Choose positive meridians
$\gamma_1,\gamma_2,\gamma_3$, in that order, with product one; positivity
at infinity uses $1/t$. Use paths from a base point in the upper
half-plane to its boundary values $\infty,0,1$. On the lattice
$\Lambda\simeq\mathbb Z^4$ take the homology monodromies
\[
 A_1=\begin{pmatrix}1&0&0&0\\0&-1&2&0\\0&0&-1&0\\0&1&0&1\end{pmatrix},\quad
 A_2=\begin{pmatrix}-1&0&0&0\\0&1&0&0\\2&0&-1&0\\0&0&1&1\end{pmatrix},\quad
 A_3=\begin{pmatrix}-1&0&0&0\\0&-1&-2&0\\-2&0&1&0\\2&1&1&1\end{pmatrix}.
\]
Their determinants are one and $A_1A_2A_3=I$. The following integral
matrices all have determinant one:
\[
 C_1=\begin{pmatrix}0&-1&0&0\\-1&0&-2&-1\\0&0&0&1\\0&0&1&0\end{pmatrix},\quad
 C_2=\begin{pmatrix}0&0&0&1\\0&-1&0&0\\-1&0&-2&-1\\0&0&1&0\end{pmatrix},\quad
 C_3=\begin{pmatrix}0&0&0&1\\1&-1&2&-2\\0&1&0&1\\0&0&-1&0\end{pmatrix}.
\]
Direct multiplication gives
\begin{equation}
 A_iC_i=C_iA_{\rm loc},\qquad
 A_{\rm loc}=\begin{pmatrix}1&0&0&0\\0&1&0&0\\-1&0&-1&-1\\0&0&0&-1\end{pmatrix}.
 \label{eq:P205:integral-markings}
\end{equation}
Thus, in the local basis $(a,b,c,d)$, parallel continuation sends
$a$ to $a-c$, fixes $b$, sends $c$ to $-c$, and sends $d$ to $-d-c$.

\item We construct periods for this representation. On the real space
with basis $(e_0,e_1,e_2,e_3)$, set
\[
 B_1(\epsilon)=\begin{pmatrix}1&0&0\\0&-1&\epsilon\\0&0&-1\end{pmatrix},\quad
 B_2(\epsilon)=\begin{pmatrix}-1&0&0\\0&1&0\\\epsilon&0&-1\end{pmatrix},\quad
 B_3(\epsilon)=\begin{pmatrix}-1&0&0\\0&-1&-\epsilon\\-\epsilon&0&1\end{pmatrix}.
\]
They have product one. Their squared monodromies have logarithms
\begin{equation}
 N_1=-2\epsilon E_{23},\qquad
 N_2=-2\epsilon E_{31},\qquad
 N_3=\epsilon^2 E_{21},
 \label{eq:P205:nilpotents}
\end{equation}
where $E_{ij}e_j=e_i$ on the three-dimensional summand. Each $N_i$ has
rank one and square zero for real $\epsilon>0$. On the
four-dimensional space, extend $N_i$ by zero on $e_0$. Let $P$ have columns
$(e_0+e_1)/2,(e_0+e_2)/2,(e_0+e_3)/2,e_0$, and put
\[
 D_\epsilon=\operatorname{diag}(1,1,\epsilon^2/4,\epsilon/2),
 \qquad \Lambda_\epsilon=D_\epsilon P\mathbb Z^4.
\]
For real $\epsilon>0$ this is a full lattice. The identity
\[
 \operatorname{diag}(1,B_i(\epsilon))
 =D_\epsilon\operatorname{diag}(1,B_i(2))D_\epsilon^{-1}
\]
shows that its matrices in the displayed lattice basis are exactly
$A_i$. Thus $D_\epsilon P:\mathbb Z^4\to\Lambda_\epsilon$
intertwines the fixed integral representation with these real
matrices, whose nilpotent entries tend to zero.

\item Consider the fourfold cover
\[
 p:\mathbb P^1_z\longrightarrow\mathbb P^1_t,
 \qquad t=(z+z^{-1})^2/4.
\]
Its deck involutions are $g_1(z)=-z$, $g_2(z)=-1/z$ and $g_3(z)=1/z$.
Their fixed pairs lie over $\infty,0,1$, respectively. The triangle
$\operatorname{Re}z>0$, $\operatorname{Im}z<0$, $|z|<1$ maps
biholomorphically to the upper half-plane. The inverse is
$z=\sqrt t-\sqrt{t-1}$ with the branches there. Its three boundary
vertices are $0,-i,1$. The marked meridians lift at these vertices
to the three deck involutions in the stated order; a positive full
loop at each vertex has monodromy $B_i^2$.

For complex $\epsilon$ near zero form the flat cohomology bundle on
$B^\circ$ with monodromy $\operatorname{diag}(1,B_i)^{-t}$ and pull it
to the punctured $z$-sphere. It has a jointly holomorphic extension
$\mathcal E$ over $\mathbb P^1_z\times\Delta_\epsilon$.
Indeed, in a local coordinate $s$ at a puncture, replace a multivalued
flat cohomology frame by
\[
 e_{\rm can}(s)=e_{\rm flat}(s)
       \exp\bigl((\log s)N^t/(2\pi i)\bigr).
\]
The frame is single-valued: its two factors acquire $\exp(-N^t)$ and
$\exp(N^t)$. Declaring it holomorphic at zero gives transition
functions holomorphic jointly in $s$ on annuli and in $\epsilon$.
The logarithms in \eqref{eq:P205:nilpotents} are finite polynomials,
so this construction applies to the full parameter family. The deck
maps extend because coordinate changes introduce logarithms of
holomorphic units. At a fixed point their homology action is
\[
 \operatorname{diag}(1,B_i(\epsilon))\exp(-N_i/2).
\]
Its square is the identity, since the squared monodromy is
$\exp N_i$ and the two factors commute. The inverse transpose
gives the cohomology action in the canonical frame. This is the
finite part of the half-turn comparison, as follows by inserting
$\log(-s)=\log s+\pi i$ into the frame formula. The group
identities extend from the punctured curve.

At $\epsilon=0$ the extension is $\mathcal O_{\mathbb P^1}^4$.
Apply \cite[\S7, Satz~3, p.~62, and S\"atze~5--6, p.~64]{Grauert1960D},
with \cite[pp.~35--36]{Grauert1963ErratumD}, to the proper product
projection $q:\mathbb P^1\times\Delta_\epsilon\to\Delta_\epsilon$.
The base is reduced and $\mathcal E$ is flat over it: its stalks have
no torsion over the one-variable local discrete valuation ring.
Semicontinuity gives $h^1(\mathcal E_\epsilon)=0$ near zero, constancy
of Euler characteristic gives $h^0(\mathcal E_\epsilon)=4$, and base
change makes $W=q_*\mathcal E$ locally free of rank four with fibre
$H^0(\mathcal E_\epsilon)$. The evaluation map $q^*W\to\mathcal E$
is an isomorphism at zero. Its closed degeneracy set has proper
projection and misses zero, so the map is an isomorphism after
shrinking the parameter disc.

The deck action on $W$ can be trivialized with fixed characters.
For each character $\chi$ of the Klein four-group use
\[
 P_\chi(\epsilon)=\frac14\sum_{g\in G}\chi(g)^{-1}\rho_\epsilon(g).
\]
Its image is a holomorphic line subbundle, since its rank is one at
zero and remains one. Choose a holomorphic frame of $W$ whose values at zero are
the real character basis, and denote its sections by
$e_\chi(\epsilon)$. The section
$P_\chi(\epsilon)e_\chi(\epsilon)$ is holomorphic and remains nonzero
after shrinking. These four sections and the evaluation isomorphism
give an equivariant holomorphic trivialization of $\mathcal E$
with the four fixed character matrices.

\item In that trivialization use the two-plane spanned by the rows
\[
 \Pi_{\rm tw}(\zeta)=
 \begin{pmatrix}1&i&-\zeta&i\zeta\\\zeta&-i\zeta&1&i\end{pmatrix}.
\]
At infinity the rescaled rows tend to $(0,0,-1,i)$ and $(1,-i,0,0)$.
This gives a holomorphic map into the Grassmannian. In the finite
chart its determinant against its conjugate is
$4(1+|\zeta|^2)^2>0$. The character matrices
\[
 S_1=\operatorname{diag}(1,1,-1,-1),\quad
 S_2=\operatorname{diag}(1,-1,1,-1),\quad
 S_3=\operatorname{diag}(1,-1,-1,1)
\]
act by $\zeta\mapsto-\zeta,-1/\zeta,1/\zeta$. More explicitly,
\[
 \Pi_{\rm tw}(\zeta)S_i=K_i(\zeta)\Pi_{\rm tw}(g_i\zeta),
 \qquad
 K_1=\begin{pmatrix}1&0\\0&-1\end{pmatrix},\quad
 K_2=\begin{pmatrix}0&-\zeta\\\zeta&0\end{pmatrix},\quad
 K_3=\begin{pmatrix}0&\zeta\\\zeta&0\end{pmatrix}.
\]
For the last two identities use $\zeta\ne0$; the displayed limiting
frames give the same plane maps at zero and infinity. These invertible row changes identify the planes on the overlapping
charts. Take
\[
 F^1_\epsilon(z)=F_{\rm tw}(-1/z^3).
\]
The cubic map commutes with these three transformations, so this
subbundle is equivariant. At the marked lifts $0,-i,1$, its limiting frames are,
respectively, the displayed frame at $\zeta=\infty$,
$\Pi_{\rm tw}(i)$ and $\Pi_{\rm tw}(-1)$: these are the
values of $-1/z^3$ at those three lifts.

We check transversality uniformly. In continued homology coordinates
the period row matrix satisfies
\[
 \Pi_{\rm per}=\Pi_{\rm can}\exp(ZN),\qquad
 \Pi_{\rm per}^{t}=\exp(ZN^t)\Pi_{\rm can}^{t},\qquad
 Z=\frac{\log s}{2\pi i}.
\]
This agrees with parallel cohomology monodromy $A^{-t}$.
Let $\omega$ be the wedge of the two canonical-frame covectors,
and let $N^t$ act on exterior powers by its induced derivation.
Rank one gives $(N^t\omega)\wedge(N^t\bar\omega)=0$, and trace zero
gives $(N^t\omega)\wedge\bar\omega+\omega\wedge(N^t\bar\omega)=0$.
Consequently the determinant is exactly
\[
 D=D_0+2i\operatorname{Im}Z\,(N^t\omega\wedge\bar\omega).
\]
There is no quadratic logarithmic term. At $\epsilon=0$ the three canonical frames are the twistor
frames just listed. Substituting those frames and the polynomial
matrices $N_i(\epsilon)$ gives, with $y=\operatorname{Im}Z$,
\[
 4(1+2\epsilon y),\qquad16(1+2\epsilon y),\qquad
 16(1+\epsilon^2y).
\]
Divide the coefficient of $y$ by $\epsilon$ at the first two
punctures and by $\epsilon^2$ at the third. The quotients extend
continuously to $\epsilon=0$, because the nilpotents have those exact
factors and the canonical frames depend holomorphically on the
parameter. Their central limits are strictly positive. Thus fixed
puncture discs of radius less than one and a fixed positive parameter
interval make all these coefficients positive. The constant terms
are uniformly positive there by their twistor limits. Since $y>0$
on these discs, $D>0$ throughout them. On their fixed compact
complement transversality follows by continuity from $\epsilon=0$;
the other three punctures follow by deck equivariance. The real transition matrices $\operatorname{diag}(1,B_i)$ have
determinant one, and a change of Hodge-row basis multiplies $D$
by $|\det K|^2$. Hence these inequalities agree on overlaps.

Fix one real $\epsilon$ in this interval. Real evaluation
\[
 H_{\mathbb R}\longrightarrow(F^1)^*,\qquad
 \lambda\longmapsto(\sigma\mapsto\sigma(\lambda))
\]
is an isomorphism. Quotienting the holomorphic bundle $(F^1)^*$ by
the evaluated lattice $\Lambda_\epsilon$ gives a smooth proper
holomorphic group family $J\to B^\circ$ of two-tori, with zero
section and precisely the monodromies $A_i$. The periods are
holomorphic sections. A compact real period parallelepiped over a
compact base neighbourhood proves local properness.

\item At each puncture use the marking $C_i$.
Put $p=2a-c$. On the square-root base with coordinate $s$, the
periods of $p,b,c$ are single-valued and $d$ acquires $2c$ around
a positive loop. At the canonical limit the semisimple part has
real eigenspaces $\langle p,b\rangle$ and $\langle c,d\rangle$.
The transverse Hodge plane has one complex line on each. Evaluation
on the nonzero real vectors $p,c$ gives independent columns.
Normalize those columns to $(1,0),(0,1)$. The periods are then
\[
 a=(1/2,1/2),\quad b=(\tau(s),\beta(s)),\quad c=(0,1),\quad
 d=(\alpha(s),2\log(s)/(2\pi i)+\gamma(s)),
\]
with the four coefficient functions holomorphic at zero.
The logarithm kills $p,b,c$, and its exact increment on $d$ is $2c$.
The half-turn comparison is
\[
 \Pi(-s)=R\Pi(s)A_{\rm loc},\qquad R=\operatorname{diag}(1,-1).
\]
The fixed $p,c$ columns force this $R$. In the endpoint markings the
direct deck map has lattice matrix $A_{\rm loc}^{-1}$; continued
parallel transport first adds $c$ to $d$, and returning by $R$
gives $-d-c$, as in $A_{\rm loc}$. Comparing the other columns gives
\[
 \tau(-s)=\tau(s),\quad\gamma(-s)=\gamma(s),\quad
 \alpha(-s)=-\alpha(s),\quad\beta(-s)=-\beta(s).
\]
In particular $\alpha(0)=\beta(0)=0$.

The sign in this integral marking is $\operatorname{Im}\tau(0)<0$.
Indeed $\det P=-1/8$, $\det D_\epsilon>0$ and $\det C_i=1$,
so the positive determinant in the real $e$-basis becomes negative
in $(a,b,c,d)$. Row normalization changes it by a positive factor.
At the split limit its coefficient of $\operatorname{Im}\ell$,
$\ell=\log(s)/(2\pi i)$, is $4\operatorname{Im}\tau(0)$.
The limiting transverse line gives $\operatorname{Im}\tau(0)\ne0$,
and $\operatorname{Im}\ell>0$ proves the sign. We keep this basis
and shrink so $\operatorname{Im}\tau$ stays bounded away from zero
in the negative half-plane. The odd coordinate change
$r=s\exp(\pi i\gamma(s))$ is locally invertible and changes the last
period to $d=(\alpha(r),2\log(r)/(2\pi i))$. The parity statements
persist, and the original base coordinate is a holomorphic invertible
function of $r^2$.

\item Form a chain of charts $U_k$, $k\in\mathbb Z$, with coordinates
$(z,x_k,y_k)$ and $r=x_ky_k$, glued by
\[
 x_{k+1}=y_k^{-1},\qquad y_{k+1}=r y_k,\qquad z_{k+1}=z_k.
\]
The chain is Hausdorff. For $j>k$ on a nonzero fibre the
transition is $x_j=r^{k-j}x_k$, $y_j=r^{j-k}y_k$. If a
sequence with $r\to0$ converges in both charts,
boundedness of $x_j$ forces a nonzero limiting $y_k$. Then
$x_k\to0$, and $j\ge k+2$ would make $x_j$ unbounded. For
$j=k+1$ the limits satisfy $x_{k+1}=1/y_k$, $y_{k+1}=0$,
so they are already identified by the adjacent transition.
The same-chart limits agree, and these cases also cover
central sequences. Thus distinct points have disjoint
neighbourhoods. On nonzero fibres $u=r^kx_k$ is the common
multiplicative coordinate.
Set
\[
 b(r)=\exp(2\pi i\beta(r)).
\]
Take the commuting actions
\[
 \begin{aligned}
 P:(z,x,y)&\longmapsto(z+1,x,y),\\
 B:(z,x,y)&\longmapsto(z+\tau(r),b(r)x,b(r)^{-1}y),\\
 T:U_k\longrightarrow U_{k+2},\quad(z,x,y)&\longmapsto(z+\alpha(r),x,y).
 \end{aligned}
\]
They preserve the gluing equations. The quotient by $P,B$ is a
chain of line-bundle charts over $\mathbb C/(\mathbb Z+\tau(r)\mathbb Z)$;
these actions are free and properly discontinuous because
$|\operatorname{Im}\tau|$ is bounded below. On its nonzero fibres use
the invariant norm
\[
 \|u\|=|u|\exp\left(
 \frac{2\pi\operatorname{Im}\beta(r)}{\operatorname{Im}\tau(r)}
          \operatorname{Im}z\right).
\]
The $B$-multiplier and the exponential change cancel exactly. The
map $T$ acts by $(z,u)\mapsto(z+\alpha(r),r^2u)$ and multiplies this
norm by
\[
 \kappa(r)=|r|^2\exp\left(
 \frac{2\pi\operatorname{Im}\beta(r)\operatorname{Im}\alpha(r)}
      {\operatorname{Im}\tau(r)}\right)
       =|r|^2\exp(O(|r|^2)).
\]
It lies in $(0,1)$ on a sufficiently small punctured disc.

The invariant chart norms are
\[
 \begin{aligned}
 \|x_k\|&=|x_k|\exp\!\left(
 \frac{2\pi\operatorname{Im}\beta(r)}{\operatorname{Im}\tau(r)}
 \operatorname{Im}z\right),\\
 \|y_k\|&=|y_k|\exp\!\left(
 -\frac{2\pi\operatorname{Im}\beta(r)}{\operatorname{Im}\tau(r)}
 \operatorname{Im}z\right).
 \end{aligned}
\]
They satisfy $\|x_k\|\|y_k\|=|r|$ and
$\|u\|=|r|^k\|x_k\|$. They extend to the respective
line-bundle charts at the centre.

These estimates also prove proper discontinuity at the centre.
A compact set is covered by finitely many charts, say $|k|\le L$,
with both chart norms bounded by $R_0\ge1$. On nonzero
fibres its points satisfy
\[
 |r|^{L+1}/R_0\le\|u\|\le R_0|r|^{-L}.
\]
An intersection with its $n$-th translate forces
$|n\log\kappa(r)|\le(2L+1)|\log|r||+2\log R_0$.
After shrinking, $|\log\kappa(r)|\ge|\log|r||$, giving a uniform
bound on $|n|$. At the centre $T$ shifts the component and node
indices by two, so only finitely many translates of those charts
meet. It has no central fixed point. Thus the quotient $X$ is a
smooth Hausdorff complex threefold.

For properness choose on every nonzero orbit a representative with
$\kappa(r)\le\|u\|\le1$. Use $U_0$ when $\|u\|\ge|r|$ and $U_1$
otherwise. In the first chart both coordinate norms are at most
one; in the second they are bounded by one and $|r|^2/\kappa(r)$.
The latter is uniformly bounded. Closed norm discs over the compact
elliptic base family are compact, and these two images also cover
the central projective-line halves. This proves properness over a
closed subdisc. Additively, quotient by $c$ is $u=\exp(2\pi i w)$;
$P,B,T$ then give the periods $p,b,d$. Hence the generic fibre of
$X$ has period lattice $(p,b,c,d)$.

\item Define
\[
 g(z,u,r)=(z+1/4,r/u,-r).
\]
On the chain this is the holomorphic map
\[
 U_k\longrightarrow U_{-k},\qquad
 (z,x,y)\longmapsto(z+1/4,(-1)^ky,(-1)^{k+1}x).
\]
Substitution verifies the adjacent-chart transition. The parity of
$\tau,\alpha,\beta$ gives $gP=Pg$, $gB_r=B_{-r}g$ and
$gT_rg^{-1}=T_{-r}^{-1}$, so it descends to $X$. Its square is
\[
 h=g^2:(z,x,y,r)\longmapsto(z+1/2,-x,-y,r),\qquad g^4=P.
\]
On a noncentral fibre $h$ translates by $(p+c)/2$, which is outside
the real-independent lattice $(p,b,c,d)$ and has order two. At the
centre $\alpha=\beta=0$ and translation by $1/2$ on the elliptic
base proves freeness. The odd powers exchange nonzero base values;
at the centre their translations by $1/4$ and $3/4$ are nonzero,
including at the nodes. The order-four action is therefore free.

Set $M=X/\langle g\rangle$. It is smooth and proper over the original
disc. The full central divisor is $2S$: the local full equation is
a unit times $(xy)^2$ and the quotient cover is étale. The central
fibre of $X$ is $E_{\tau(0)}\times I_2$, where
$E_\tau=\mathbb C/(\mathbb Z+\tau\mathbb Z)$. The map $g$ exchanges its two
components and preserves its two node parities. Thus $S$ is
irreducible. Include either component $E_{\tau(0)}\times\mathbb P^1$
in $X_0$ and compose with the quotient $X_0\to S$. The component
stabilizer is $\langle h\rangle$, so this map induces
\[
 N=(E_{\tau(0)}\times\mathbb P^1)\big/
       \langle(z,v)\mapsto(z+1/2,-v)\rangle
       \longrightarrow S.
\]
It is finite and bijective away from the two node curves, hence
birational. Its source is smooth, so this is the normalization map.
The surface $N$ is $\mathbb P_C(\mathcal O_C\oplus L)$, where
$C=E_{\tau(0)}/\langle1/2\rangle$ and $L$ is the sign line of that
double cover. Its square is trivial. A trivialization of $L$ would
give a nonzero anti-invariant holomorphic function on the compact
elliptic cover, which is impossible; hence $L$ is nontrivial.
Every rational curve in $N$ lies in a ruling fibre, so there is no
$(-1)$-curve and $N$ is minimal. Its two sections map to two
distinct curves by the degree-two isogenies
\[
 E_{\tau(0)}/\langle1/2\rangle
       \longrightarrow E_{\tau(0)}/\langle1/4\rangle.
\]
The quotient map is a local biholomorphism, so its inverse chart
identifies the reduced germ with $xy=0$ times the $z$-disc. The conductor pushout of \cref{lem:II:proper-conductor-pushout}
gives $e(S)=e(N)-e(D)+e(C_{\rm im})$. Here $D$ is the union
of the two sections and $C_{\rm im}$ the union of their image
curves. All three terms vanish, since $N$ is a projective-line
bundle over an elliptic curve and all four conductor curves
are elliptic. Thus $e(S)=0$.

\item The punctured local model is biholomorphic to the already
constructed $J$, including its affine descent. Quotienting $X$ by
$h$ adds the period $a=(p+c)/2$. Its linear descent is
$g_{\rm lin}(z,w,r)=(z,-w,-r)$, while the affine descent is
\[
 g_{\rm aff}(z,w,r)=(z+1/4,\ell(r)-w,-r),\qquad
 \ell(r)=\log(r)/(2\pi i).
\]
Put $v(r)=(-\ell(r)/2,-\ell(r)/2+1/8)$. A full loop changes $v$
by $-a$, so it is a holomorphic section of the root torus family.
Since $\ell(-r)=\ell(r)+1/2$,
\[
 v(-r)-Rv(r)=-(1/4,\ell(r)),\qquad
 \Phi_r(x)=x+v(r),\qquad
 \Phi_{-r}g_{\rm aff}=g_{\rm lin}\Phi_r.
\]
Thus $\Phi$ descends to a biholomorphism with $J$ on the punctured
original disc. Its inverse is translation by $-v$.

\item We calculate the local restriction maps. First set
$\alpha=\beta=0$ and make $\tau$ constant, so the root tube is
$E\times\mathcal E$, with $\mathcal E$ the displayed $I_2$ family.
Apply \cref{lem:II:compact-analytic-tube} to the proper
map $X\to\Delta_r$ just constructed: its total space is smooth
and its punctured fibres are smooth two-tori. The inclusion of
$E\times I_2$ in a sufficiently small closed root tube is a
homotopy equivalence. This inclusion is equivariant for $g$,
so its induced cohomology isomorphism commutes with $g^*$.
The nearby elliptic fundamental class maps to $[P]+[Q]$, where
$P,Q$ are the two central components: the inverse image of an
oriented radial interval is a triangulated three-chain with
boundary the nearby fibre minus those components, each with
multiplicity one by $xy=r$. For the dual central classes put
$\eta=(\eta_P+\eta_Q)/2$. Its restriction is the primitive elliptic
area class. The graph one-class restricts to $d^*$. In the plumbing
description, collapse the circular $c$ direction in each neck.
The longitudinal path from $U_0$ to $U_2$, closed by $T$,
traverses the two-component graph once and represents $d$.
Thus this graph map sends $c$ to zero and $d$ to its oriented
generator.

The map $g$ translates $E$, exchanges $P,Q$ and reverses the graph
loop. Its invariant central cohomology has dimensions
$(1,2,2,2,1)$. Representatives are the constant; the two $E$
one-classes; the $E$ area and $\eta$; the $E$ one-classes times
$\eta$; and the product of the two area classes. Their restrictions
are independent and nonzero. In the rational basis $(p,b,c,d)$,
local monodromy is identity on the $E$ span and minus a unipotent
matrix on $(c,d)$, so this is exactly its invariant cohomology.
The free degree-four quotient of tubes and the degree-two
translation quotient of each nonzero root fibre give a commuting
restriction square. Rational pullback has inverse transfer divided
by the covering degree. Consequently
\[
 H^q(M;\mathbb Q)\xrightarrow{\sim}
        H^q(F;\mathbb Q)^{A_{\rm loc}}\quad(0\le q\le4).
\]

These maps persist for the moving periods. Deform $\alpha,\beta$
linearly to zero and $\tau(r)$ linearly to $\tau(0)$. Parity persists;
the norm estimates hold uniformly on a single smaller disc, with
$\operatorname{Im}\tau$ in a fixed negative half-plane. The quotient
charts form a smooth proper family of closed tubes with boundary,
with the function $r^2$ independent of this real parameter. Local
lifts of the parameter vector field preserving $r^2$ glue by a
partition of unity. Their complete flows give diffeomorphisms of
the tube, chosen regular fibre and boundary pairs. The period
generators vary continuously with the fixed integral homology marking.
Thus the stated restriction maps hold for the original local model.

\item Apply \cref{lem:II:proper-conductor-pushout} to the finite
normalization $N\to S$, its two disjoint conductor sections and
their two image curves. For the local groups use a sufficiently
small tube in $M$. Applying \cref{lem:II:compact-analytic-tube}
to $M\to\Delta_t$ identifies its integral homology with that
of $S$ through the inclusion. Choose compatible finite CW structures
with those sections as subcomplexes. The incidence graph is a
tree. Projection identifies $\pi_1(N)$ and each section group
with $\langle p/2,b\rangle$, while the two image groups are
$\langle t_1,b\rangle$ and $\langle t_2,b\rangle$, with
$t_i=p/4$. Each section-to-image map sends $(p/2,b)$ to
$(2t_i,b)$. Thus the two image lattices are amalgamated along
this common sublattice. Their abelianization has relation
$2t_1=2t_2$; put $t=t_2$ and $\delta=t_1-t_2$. It is
\[
 H_1(M;\mathbb Z)=\mathbb Z t\oplus\mathbb Z b\oplus\mathbb Z/2,
 \qquad t=p/4.
\]
For the nearby inclusion $j$, the monodromy relations give
$j_*c=0$ and $2j_*d=0$. Since $p=2a-c$ and $j_*p=4t$,
the free part of $j_*a$ is $2t$; the $b$-loop maps to $b$.
Dualizing gives
\[
 \operatorname{im}\bigl(H^1(M;\mathbb Z)\to
           H^1(F;\mathbb Z)^{A_{\rm loc}}\bigr)
       =2\mathbb Z a^*\oplus\mathbb Z b^*.
\]
Under $\Phi^{-1}$ the zero Jacobian section is
$w=-v=(\ell/2,\ell/2-1/8)$. A positive original base loop changes
$\ell$ by $1/2$, giving $z$-increment $1/4$ and no $b$-increment;
its endpoints are identified by $g_{\rm aff}$. Its free tube class
is $t$. The unique rational tube class restricting to
$u a^*+v b^*$ therefore has value $u/2$ on this loop. Throughout
the preceding deformation the same formula for $w$ and the same
identity for $\Phi$ hold. The resulting continuous family of loops
in the marked tube/boundary pairs preserves this value.

\item Choose local base coordinates $t_i$ at $\infty,0,1$, and
$0<r_0<r_1$ so the closed $r_1$-discs are disjoint and lie in the
local construction domains. Use the open cover
\[
 B_0=\mathbb P^1\setminus\bigcup_i\{|t_i|\le r_0\},\qquad
 D_i=\{|t_i|<r_1\}.
\]
Glue $J|B_0$ to $M_i|D_i$ on the entire annuli
$r_0<|t_i|<r_1$ by the explicit $C_i$ markings and $\Phi_i$ maps.
Their inverse translations were just proved. The result is a smooth
Hausdorff complex threefold: different base values separate points,
and a single base chart supplies a Hausdorff model near each fibre.
The map is proper locally on every base chart, hence proper. Since
the base is compact, so is $Y$. Every fibre is connected; thus $Y$
is connected. The central descriptions give its three full fibres
$2S_i$, and all other fibres are the regular tori of $J$.

\item The primitive
common invariant integral cohomology vector and its pullbacks are
\[
 \alpha=(1,1,1,2)^t,\qquad
 C_1^t\alpha=(-1,-1,0,0)^t,\quad
 C_2^t\alpha=(-1,-1,0,0)^t,\quad
 C_3^t\alpha=(1,0,0,0)^t.
\]
Their first coordinates are $u_i=(-1,-1,1)$. By the preceding exact
loop calculation the local rational classes have base-loop values
$(-1/2,-1/2,1/2)$, whose sum is $-1/2$.

Apply Mayer--Vietoris to the inverse image of a sphere with three
holes and the three tubes, with annular overlaps. The degree-zero
difference map $\mathbb Q^4\to\mathbb Q^3$ is onto. The outside
Jacobian is, as a real torus bundle, the flat lattice bundle via
evaluation. Its zero section gives
\[
 H^1(Y_0;\mathbb Q)=H^1(B_0;\mathbb Q)\oplus\mathbb Q\alpha.
\]
A base class has three positive-loop values $x_i$ with $\sum x_i=0$;
the invariant constant fibre class has zero value on the zero-section
loops. Matching a class $c\alpha$ with the tube classes first
forces their fibre restrictions to be $cC_i^t\alpha$, and then
forces $x_i=cu_i/2$. Summing gives $0=-c/2$, so all classes vanish.
This proves $H^1(Y;\mathbb Q)=0$.

An integral
invariant class extends to each tube precisely when its local
$a^*$-coefficient is even. All three $u_i$ are odd, so
$H^0(\mathbb P^1,R^1f_*\mathbb Z)=2\mathbb Z\alpha$.
The base-loop values for $2\alpha$ are $-1,-1,1$. In the obstruction
convention given by this Mayer--Vietoris difference map and the
ordered positive meridians, the transgression is
\[
 d_2(2\alpha)=-[\mathbb P^1]\in H^2(\mathbb P^1;\mathbb Z).
\]
In particular it is primitive and nonzero; rationally
$d_2(\alpha)=-\tfrac12[\mathbb P^1]$.

\item Put $V=R^1f_*\mathbb Q|B^\circ$. The local
restriction maps above identify
\[
 R^qf_*\mathbb Q=j_*\bigwedge^qV.
\]
Here $j:B^\circ\hookrightarrow\mathbb P^1$, and the identifications
are the identity on the smooth base and the proved isomorphisms
on all other stalks. For every invertible $T$, the identity
$T^{-1}-I=-T^{-1}(T-I)$ identifies the fixed kernels of
$T$ and $T^{-1}$. Apply it to each $A_i^t$ and its exterior
powers for the parallel cohomology monodromies $A_i^{-t}$.
For a homology vector $x$, the equations $(A_1-I)x=0$
give $x_2=x_3=0$, and $(A_2-I)x=0$ then gives $x_1=0$;
$A_3$ also fixes $(0,0,0,1)^t$. The common homology fixed
space is therefore $\mathbb Q(0,0,0,1)^t$. For a cohomology vector $u$, the
first transpose gives $u_4=2u_2$ and $u_3=u_2$, and the
second gives $u_1=u_3$. The displayed $\alpha$ satisfies
the third as well. Thus the common cohomology fixed space
is exactly $\mathbb Q\alpha$.
Each local fixed space of $V$ has dimension two. Euler additivity
therefore gives $\chi(j_*V)=-4+3\cdot2=2$. The exact sequence from
$j_!V$ to $j_*V$ has only skyscraper cokernel, so their second
cohomologies agree. Compact-support duality on the oriented
punctured sphere identifies its dimension with the common fixed
dimension of $V^*$, namely one. Since $h^0(j_*V)=1$, this proves
$H^1(\mathbb P^1,j_*V)=0$.

Let $\lambda_1^*,\ldots,\lambda_4^*$ be the basis dual to
the displayed integral lattice basis. Write an exterior
class as $\sum_{i<j}v_{ij}\lambda_i^*\wedge\lambda_j^*$.
Invariance under
$\bigwedge^2A_1^t$ and $\bigwedge^2A_2^t$ implies
\[
 v_{14}=2v_{12},\quad v_{13}=v_{12},\quad
 v_{34}=2v_{24},\quad v_{24}=0,\quad
 v_{23}=-v_{12},\quad v_{14}=2v_{34}.
\]
They force all six coefficients to vanish. In coefficient
order $12,13,14,23,24,34$, the six corresponding linear
forms $-2v_{12}+v_{14}$, $2v_{12}-2v_{13}$,
$-2v_{24}+v_{34}$, $-2v_{24}$,
$-2v_{12}-2v_{23}$ and $-v_{14}+2v_{34}$ have determinant
$-16$. Thus $H^0(\mathbb P^1,j_*\bigwedge^2V)=0$. The three
total-degree-two Leray terms are now
\[
 \operatorname{coker}d_2=0,\qquad H^1(\mathbb P^1,j_*V)=0,
 \qquad H^0(\mathbb P^1,j_*\bigwedge^2V)=0.
\]
Hence $H^2(Y;\mathbb Q)=0$. Every support has Euler number zero,
so Euler additivity gives $e(Y)=0$. Poincar\'e duality on the smooth
compact oriented six-manifold gives $b_3(Y)=2$.
\end{enumerate}
\end{proof}
\section{Local exclusions and uniform transverse types}\label{sec:II:uniform-transverse-types}

\subsection{Conductor graphs and rational Euler numbers}

\begin{lemma}[Two image branches force a conductor cycle]\hypertarget{tag-P255}{}\label{lem:II:conductor-two-branch-cycle}
\StatementRecord{lem:II:conductor-two-branch-cycle}{proved}{P255}{\cref{lem:II:conductor-duality}; \cite[Tags~0AG0 and 0BBZ, Lemmas~15.123.2 and 10.60.12]{StacksProject}}
Let $\nu:N\to S$ be the finite normalization of a reduced hypersurface in a smooth complex threefold. Let $C$ be a compact irreducible component of the reduced image conductor, and let $D\subset N$ be the union, with reduced structure, of the conductor curves which dominate $C$. Put $c=b_0(D^\nu)$ and $k=b_0(D)$, and let $\beta(D)$ be the cycle rank of its normalization-incidence graph. At a point $p\in C$ with exactly two analytic branches and lying on no other image-conductor component, one has
\[
 \sum_{q\in\nu^{-1}(p)}\bigl(|\kappa^{-1}(q)|-1\bigr)\ge c,
 \qquad \beta(D)\ge k\ge1,
\]
where $\kappa:D^\nu\to D$ is normalization. In particular, a conductor with an acyclic normalization-incidence graph has no such image point.
\end{lemma}
\begin{proof}
Let $\eta:C^\nu\to C$ be normalization and $h=\nu|_D$. Integral closure gives the unique factorization
\[
 h\kappa=\eta\phi,\qquad \phi:D^\nu\to C^\nu.
\]
Each normalized component of $D$ maps nonconstantly, because $\nu$ is finite. It is compact, so its map to the connected compact curve $C^\nu$ is finite and surjective. Write $p_+,p_-$ for the preimages of $p$.

The $n$ points of $\nu^{-1}(p)$ correspond to the irreducible analytic hypersurface branches of $S$ at $p$. Every pair of distinct branches meets in a curve: their nonassociate irreducible equations in the regular local ring of dimension three have height two. The height bounds are \cite[Tags~0AG0 and 0BBZ, Lemmas~15.123.2 and 10.60.12]{StacksProject}; a common height-one prime would make the two factors associates. Their intersection is in the nonnormal locus, and hence in $C$ near $p$, since no other image component passes through $p$.

Every point over $p$ lies on $D$. With at least two surface branches this follows from those pairwise intersections. With one surface branch it follows by taking the finite inverse image of either conductor-curve branch through $p$. The conductor purity of \cref{lem:II:conductor-duality} also excludes an isolated conductor component. For each normalization point record the nonempty subset of $\{+,-\}$ consisting of the image branches of its local conductor curves. The pairwise surface intersections make these subsets pairwise intersect. Such a family has a common element: a singleton forces its element into all subsets, and in the absence of a singleton every subset is $\{+,-\}$. Interchange the labels to make the common element $+$.

Let $r_+$ and $r_-$ be the numbers of points of $D^\nu$ above $p_+$ and $p_-$ under $\phi$, and put $r_q=|\kappa^{-1}(q)|$. Every normalization point over $p$ has a plus branch, so $r_+\ge n$. Surjectivity on each of the $c$ normalized components gives $r_-\ge c$. Therefore
\[
 \sum_{q\in\nu^{-1}(p)}(r_q-1)=r_++r_--n\ge c.
\]
The normalization-incidence graph of a reduced compact curve with $k$ connected components and $c$ normalized components has cycle rank
\[
 \beta(D)=k-c+\sum_{q\in D}(r_q-1).
\]
Every term outside $\nu^{-1}(p)$ is nonnegative. The preceding inequality gives $\beta(D)\ge k$, proving the assertions.
\end{proof}

\begin{proposition}[Acyclic normalization conductors]\hypertarget{tag-02X8}{}\label{night:P4.35}\label{stmt:P4.35}
\StatementRecord{night:P4.35}{proved}{02X8}{\cref{lem:II:conductor-two-branch-cycle,night:P4.26,night:P4.33,night:P4R.6}}
\hypertarget{tag-02X7}{}\label{night:P4.34}\label{stmt:P4.34}
\StatementRecord{night:P4.34}{proved}{02X7}{}

In the setting of \cref{night:P4.12}, no irreducible support with rationally resolved normalization
has Euler number one. Consequently, in a fibration whose
supports are all irreducible or normal crossing, every
irreducible singular support with rationally resolved
normalization has Euler number exactly two.

More generally, an irreducible support cannot have reduced
image conductor equal to a rational curve with exactly one
point-identification cycle when its reduced normalization
conductor has no graph cycles. This excludes that conductor
case for every ruled-resolution double-horizontal pole tree.
\end{proposition}
\begin{proof}
\begin{enumerate}[label=\arabic*.,leftmargin=*]
\item Let $\nu:N\to S$ be normalization and $D\subset N$ its reduced conductor. The finite-duality identity of \cref{lem:II:conductor-duality} makes the conductor ideal on the normal surface a maximal Cohen--Macaulay module. The exact sequence from that ideal to $\mathcal O_N$ makes its nonzero quotient Cohen--Macaulay of dimension one by the depth lemma. Thus its reduced support $D$ has no isolated component. Suppose that its normalization-incidence graph is a forest and its image $C$ is an irreducible rationally normalized curve with exactly one identification cycle. The graph formula gives a unique two-branch point $p\in C$, with every other point unibranch. Apply \cref{lem:II:conductor-two-branch-cycle} to this whole image curve and its same inverse conductor $D$. Its inequality is $\beta(D)\ge b_0(D)\ge1$, contradicting the forest. This proves the general conductor assertion.

\item Now let the normalization have rational resolution
and suppose $e(S)=1$. A normal such support would be
smooth by \cref{night:P4.12}; rationality would give
$\chi(\mathcal O_S)=1$, contrary to the same proposition.
Thus $S$ is singular. The normal-crossing alternative
has Euler number two by the single-divisor calculation in
the proof of \cref{night:P4.6}. In the other alternative
\cref{night:P4.26} gives a connected simply connected resolved
pole support $A\subset T$, and
\[
 \pi_1(S)=\pi_1(C)=F_b,\qquad b\le1,
 \qquad e(C)=c_C+1-b.
\]
Its normalized components are rational.
Since $e(S_{\rm reg})\ge0$, Euler number one forces
$c_C=b=1$ and $e(S_{\rm reg})=0$.
It remains to verify the tree property for the
conductor $D\subset N$, allowing singular $N$.

\item Proposition~\ref{night:P4.33} makes every singularity
of $N$ rational. Write $\rho_N:T\to N$ for the resolution.
Every connected exceptional fibre of $T\to N$
is simply connected. Over a small Stein neighbourhood,
$R^1\rho_{N*}\mathcal O_T=0$ gives
$H^1(\mathcal O_T)=0$ on its inverse image.
This proper map has every fibre of dimension at
most one. For a coherent sheaf $\mathcal F$ on its inverse
image, let $Z_j$ be the infinitesimal fibre at $y\in N$
defined by $\mathfrak m_y^j\mathcal O_T$, and put
$\mathcal F_j=\mathcal F/\mathfrak m_y^j\mathcal F$.
Every coherent sheaf on the compact space $Z_j$ has zero
cohomology in degrees at least two. Indeed, filter it by
powers of the nilradical. Each quotient is a coherent
sheaf on the reduced fibre. For a finite fibre the
vanishing is immediate. For a curve fibre, its finite
normalization has zero higher direct images, since inverse
images of sufficiently small Stein neighbourhoods are
Stein. The kernel and cokernel of the natural
pullback--pushforward map have finite support. On each
smooth normalized curve,
a coherent sheaf is an extension of a vector bundle by a
sheaf of finite support. The Dolbeault resolution of that
bundle has length one. These exact sequences give the
vanishing on the reduced fibre and then on $Z_j$.
Grauert's coherence and formal-functions theorems
\cite[\S6, Haupts\"atze I and IIa, p.~59]{Grauert1960D}
identify the completion of $(R^2\rho_{N*}\mathcal F)_y$
with $\varprojlim_j H^2(Z_j,\mathcal F_j)=0$.
The stalk is finite over the Noetherian local ring
$\mathcal O_{N,y}$, so its zero completion makes the stalk
zero. Thus $R^2\rho_{N*}\mathcal F=0$ for every coherent
sheaf $\mathcal F$. Applied to the ideal of the reduced
exceptional curve $E$, Leray and the Stein property
give $H^2(\mathcal O_T(-E))=0$ there. Its ideal
sequence consequently gives $H^1(E,\mathcal O_E)=0$.
The curve normalization exact sequence gives
$h^1(\mathcal O_E)=\sum_i g_i+\beta+
\sum_x(\delta_x-r_x+1)$, where $c_E$ is the number of
normalized components, $g_i$ are their genera,
$\delta_x$ is the normalization-quotient length,
$r_x$ counts branches, and
$\beta=1-c_E+\sum_x(r_x-1)$ is the graph cycle rank.
Each term is nonnegative: distinct branch values modulo
the common value give $\delta_x\ge r_x-1$.
Consequently all normalized components are rational and
their point-identification graph is a tree.
Thus $E$ is simply connected.

\item By \cref{night:P4R.6}, $A$ is the full inverse image of $D$.
The proper map $A\to D$ is the quotient which
collapses these connected exceptional curves to
points and is an isomorphism elsewhere.
The quotient-induced map to $D$ is a continuous
bijection from a compact space to a Hausdorff space,
and therefore a homeomorphism. Compatible triangulation
\cite[Theorem~2, p.~464]{Lojasiewicz1964} makes these curve
inclusions cofibrations. Collapsing a connected
simply connected subcomplex preserves fundamental
groups by van Kampen. Collapsing the finitely many
exceptional fibres therefore gives $\pi_1(D)=0$.
Each normalized component of $D$ is rational,
being birational to a strict pole component in
$T$. Hence its point-identification graph is a
tree. The already proved graph assertion contradicts
the description of $C$. This excludes $e(S)=1$.

\item For a ruled double-horizontal pole tree, the
resolved pole support is $H$ with rational trees attached.
Every exceptional fibre is connected and is a
union of its components. If it does not contain
$H$, it is a connected rational subtree. If it
contains $H$, all the remaining components form
rational trees attached to that component once.
Under $\rho_N$, the normalization-incidence graph
is obtained by contracting the subtree consisting of the
exceptional component vertices and their incident point
vertices, and deleting smooth-point leaves. A strict
component cannot meet that subtree in two distinct
branches, since those two incidences and the subtree path
would already make a cycle. Contracting such a connected
subcurve therefore cannot create a cycle in the graph:
contraction of a connected subgraph of a tree again has a
tree as quotient. The only positive-genus vertex,
if contracted, is replaced by a point, and this
also introduces no graph edge. Thus the image
conductor $D$ has no graph cycles, and the same
contradiction applies.

\item For a fibration all of whose supports are irreducible
or normal crossing, \cref{night:P4.26} makes the total sum of
nonnegative fibre Euler numbers at most two.
The rational-resolution case is strictly positive;
excluding one leaves exactly two.

\end{enumerate}
\end{proof}









\subsection{Two exact local exclusions for ordinary-double curves}

\begin{proposition}[Local exclusions for ordinary double curves]\hypertarget{tag-02XB}{}\label{night:P4.38}\label{stmt:P4.38}
\StatementRecord{night:P4.38}{proved}{02XB}{\cref{lem:II:boundary-deformation-transport,cdp:st:A.5.15,cdp:st:A.5.17}}

An irreducible fibre support in \cref{def:II:standing-torus-fibration} cannot have a singular
planar reduced critical curve at a point at which every generic
transverse germ is an ordinary double point. Nor can its reduced
critical curve there be the transverse union of a plane cusp
and a smooth branch, again with ordinary-double transverse germs
on both components. The latter assertion also holds with the
cusp replaced by a plane tacnode.
These exclusions allow singular normalization.
\end{proposition}
\begin{proof}

\begin{enumerate}[label=\arabic*.,leftmargin=*]
\item Write $F=Uf_0^m$ for the original local fibre map, where
$m$ is the multiplicity of the irreducible support and $U$ is a unit.
Choose $f=U^{1/m}f_0$, so that $F=f^m$ exactly. The reduced
critical-curve germ is unchanged: every critical component through
the origin has constant value zero, and on the zero divisor
$df=U^{1/m}df_0$. We use this root in both calculations below.

Choose a small closed ball transverse to the smooth zero-divisor
strata and to the smooth critical-curve links. At each such link,
the nondegenerate transverse quadratic form gives parameter Morse
coordinates, with the curve coordinate retained, in which the root
function is $XY$. These coordinates and their local inverses vary
holomorphically under the sufficiently small deformations below.
A real direction in the curve coordinate is tangent to every level
and transverse to the fixed boundary. On its remaining compact
part, submersivity persists. These are the boundary hypotheses of
\cref{lem:II:boundary-deformation-transport}. Apply that lemma to
the first root family with the exact equality $F=f^m$. Subsequent
coordinate and unit changes use its nested-ball comparison with
exponent one at the intermediate root levels. Their composition
maps every compact spherical representative and its intersection
push-off into the original full fibre $F=t_0^m$.

\item First suppose that the reduced critical curve has ideal
$P=(z,h(x,y))$, where $h$ is reduced and singular at the
origin. The equation $f$ lies in $P^{(2)}$: its generic
transverse ordinary double point vanishes to order two on
each component. The complete-intersection square $P^2$ is
unmixed, so $P^{(2)}=P^2$. We therefore have the exact expression
\[
 f=A z^2+Bzh+Ch^2.
\]
Adding a small constant multiple of $z^2$ preserves the critical
curve and the ordinary-double boundary condition. We choose it
so that the $z^2$ coefficient at the origin is nonzero. The
Morse lemma in $z$, with parameters $x,y$, then gives
\[
 f=Z^2+h(x,y)^2u(x,y).
\]
Indeed the critical $z$ coordinate vanishes on $h=0$ and
is divisible by $h$, so its substitution has the displayed
square factor. Add a small constant to $u$ and shrink the
ball so that $u$ is a unit. Absorbing its square root in $h$
gives $Z^2+H(x,y)^2$, with the same singular reduced zero
curve. Apply the finite-gradient construction in the proof of
\cref{cdp:st:A.5.17} to this reduced singular plane function $H$.
It gives an arbitrarily small linear and constant perturbation
$H_s$ with a Morse critical point of nonzero critical value. At any such point $Z^2+H_s^2$ has a
nondegenerate isolated critical point: its plane Hessian is
$2H_s\operatorname{Hess}H_s$, and its $Z$ Hessian is $2$.
Its vanishing sphere has self-intersection $-2$.
The moving curve $H_s=0$ remains transverse to the chosen
boundary, with nondegenerate normal quadratic form there.
The preceding boundary transport therefore contradicts
$\pi_2(T^4)=0$.

\item For the second assertion write
\[
 h=y^2-x^k,\qquad P=(xz,yz,h),\qquad k=3\text{ or }4.
\]
This is the union of the plane curve $z=h=0$ and the
transverse line $x=y=0$. Its symbolic square is
\[
 P^{(2)}=(z,h)^2\cap(x,y)^2
  =(x^2z^2,xyz^2,y^2z^2,zh,h^2).
\]
For $k=4$, the first equality uses the two smooth branches of
$h=0$ separately; their complete-intersection square is unmixed
and hence is still $(z,h)^2$. To verify the intersection
formula directly, write an element of $(z,h)^2$ as
$az^2+bzh+ch^2$. The last two terms belong to $(x,y)^2$,
and $z$ is a nonzero divisor modulo this ideal. The first term
belongs to it exactly when $a\in(x,y)^2$.
Consequently
\[
 f=Azh+z^2(Bx^2+Cxy+Dy^2)+Gh^2.
\]
Small constant changes in $A,B$ retain the boundary condition
and make $A(0)B(0)\ne0$. Replace $h$ in this expression by
$h_s=y^2-x^k-sx^2$. The new critical curve contains
$z=h_s=0$ and $x=y=0$, and its boundary ordinary-double
germs persist. On the latter axis the transverse quadratic
matrix has coefficients
\[
 a=-sAz+Bz^2,\qquad b=Cz^2,\qquad c=Az+Dz^2,
\]
with off-diagonal entry $b/2$. Here the four coefficient
functions are restricted to that axis. Its determinant
divided by $z^2$ is
\[
 F(s,z)=(-sA+Bz)(A+Dz)-C^2z^2/4.
\]
At $(s,z)=(0,0)$ we have $F_z=AB\ne0$ and
$F_s=-A^2$. The implicit-function theorem gives a simple root
$z(s)=sA(0)/B(0)+O(s^2)$, nonzero for small nonzero $s$.
There the $y^2$ coefficient is nonzero, the Hessian has rank
one, and its determinant has a simple zero in the axis
direction. The parametrized Morse lemma gives
$Y^2+x^2H(x,z)$, with $H_z\ne0$. The exact coordinate
change $Z=H(x,z)$ yields $Y^2+x^2Z$.
The coordinate map $(X_0,Y_0,Z_0)=(Y,x,-Z)$ identifies this
function with $X_0^2-Y_0^2Z_0$. Apply \cref{cdp:st:A.5.15}
with $(a,b,c)=(2,2,1)$; its compact spherical class has square
$-2(2/\gcd(2,1))(1/\gcd(2,1))=-4$. The nested comparison in
the first item maps that same class into the original full fibre
$F=t_0^m$. Its sphere map lifts to the contractible universal cover
of the complex two-torus, so its homology class and square are
zero, contradicting the preserved square $-4$.

\end{enumerate}
\end{proof}


\subsection{A coefficient-two exceptional curve over an ordinary-triple germ}

\begin{proposition}[Pole coefficients over an ordinary triple curve]\hypertarget{tag-02XC}{}\label{night:P4.39}\label{stmt:P4.39}
\StatementRecord{night:P4.39}{proved}{02XC}{}

Let $S=(f=0)$ be a reduced hypersurface germ in a smooth
threefold, with one-dimensional reduced critical locus $C$.
Suppose that every generic transverse germ along $C$ is an
ordinary triple point. On the minimal resolution $T$ of its
normalization, let $D$ be the integral total pole divisor
defined by the residue of an ambient volume form divided by
$f$. No exceptional curve of positive genus has coefficient
two in $D$.

In particular, the contracted horizontal section of a
positive-genus ruled resolution of an irreducible fibre support
in \cref{def:II:standing-torus-fibration} cannot lie over
a point where all local critical-curve branches have generic
type $D_4$.
\end{proposition}
\begin{proof}

\begin{enumerate}[label=\arabic*.,leftmargin=*]
\item We blow up the ambient origin, write $S_1$ for the strict
transform and $N_1$ for its normalization, and take the
minimal resolution $T_1$ of $N_1$. The proper map $N_1\to S$
lifts to the normalization $N$ by its universal property.
Thus $T_1\to N$ is a resolution and dominates $T$ by a
sequence of point blowups. The strict transform of a
positive-genus exceptional curve $H$ remains present and has
the same pole coefficient two.

\item Put $d=\operatorname{mult}_0f$, and denote the ambient
exceptional divisor by $E$. Along the generic points of $C$,
the function and its derivatives of orders one and two vanish.
They therefore vanish on the reduced curve $C$, so $d\ge3$.
Adjunction gives
\[
 \omega_{S_1}=\pi^*\omega_S\otimes
        \mathcal O((2-d)E)|_{S_1}.
\]
Let $D_1$ be the total conductor divisor on $T_1$ relative
to $\omega_{S_1}$, and let $e$ be the Cartier pullback of
$E$. With the specified residue comparison, the total pole
of the original germ on $T_1$ is
\begin{equation}
                    D_1+(d-2)e.                       \label{eq:II:triple-pole-transform}
\end{equation}
The divisor $D_1$ is effective. Its strict conductor
coefficients are positive integers. For exceptional curves
of the minimal resolution of $N_1$, the negative-matrix
argument of \cref{night:P4R.6} applies locally: $K_{T_1}F\ge0$, since
there is no smooth rational exceptional $(-1)$-curve;
adding the strict conductor retains nonnegative intersection.
The inverse exceptional intersection matrix has nonpositive
entries, so all exceptional coefficients are nonnegative.

\item If $H$ is exceptional over $N_1$, its coefficient in $D_1$
is strictly positive. A zero coefficient would give
\[
 0=(K_{T_1}+D_1)H
   =2p_a(H)-2-H^2+D_1H>0.
\]
If that coefficient is one, adjunction gives
\[
 0=2p_a(H)-2+(D_1-H)H.
\]
The normalization of $H$ has positive genus, so $p_a(H)\ge g(H^\nu)\ge1$. Both terms in the last equality are nonnegative, hence $p_a(H)=1$ and $(D_1-H)H=0$. The normalization formula $p_a(H)=g(H^\nu)+\sum\delta(H,p)$ now gives $g(H^\nu)=1$ and zero delta invariant at every point. Thus $H$ is smooth elliptic and is disjoint from every other positive component of $D_1$. A zero-coefficient exceptional
curve cannot meet a positive one, by its own intersection
equation and $K_{T_1}F\ge0$. The connected exceptional
fibre is consequently $H$ alone and meets no strict conductor.
The entire germ of $S_1$ at its image is normal with an
isolated simple elliptic singularity. Indeed, any additional
hypersurface branch would intersect this branch in a curve
and would give strict conductor meeting $H$.

\item Set $n=\operatorname{ord}_H(e)>0$ and
$a=\operatorname{coeff}_H(D_1)\ge0$. Formula \eqref{eq:II:triple-pole-transform}
gives the exact identity
\begin{equation}
                       2=(d-2)n+a.                    \label{eq:II:triple-pole-order}
\end{equation}
Thus \eqref{eq:II:triple-pole-order} makes $d\ge5$ impossible. For $d=4$ we have
$n=1,a=0$, and $H$ is not exceptional over $N_1$.
At its generic point the normalization is an isomorphism,
so $H$ is birational to a reduced component of the projective
quartic tangent cone. Every second derivative of $f$ vanishes
on $C$. Its leading term along any analytic arc of $C$
therefore makes that quartic have a point of multiplicity
at least three. An irreducible quartic has arithmetic genus
three and loses at least three there. A smooth elliptic cubic
component together with the remaining line could have
multiplicity at most two at the point. All components of
degree at most two are rational. Thus every reduced tangent
component is rational, excluding $H$ for $d=4$.

\item We are left with $d=3$ and
\[
                   (n,a)=(2,0)\quad\hbox{or}\quad(1,1).
\]
Choose a coordinate $u$ with nonzero pure cubic initial
coefficient, apply Weierstrass preparation, and depress the
monic cubic. Multiplication by the preparation unit does not
affect Hessian vanishing on $C$, because $f=df=0$ there.
We obtain
\[
 f=\text{unit}\cdot\bigl(u^3+b(x,z)u+c(x,z)\bigr).
\]
The second $u$ derivative is $6u$, so $C$ lies in the smooth
surface $u=0$. Write its reduced ideal as $(u,h(x,z))$.
The complete-intersection cube is unmixed and equals its
symbolic cube. Generic order three along $C$ and coefficient
comparison therefore give the exact expression
\begin{equation}
 f=\text{unit}\cdot
       \bigl(u^3+h^2B(x,z)u+h^3C(x,z)\bigr).           \label{eq:II:triple-depressed-cubic}
\end{equation}
If $C$ is singular, $\operatorname{ord}h\ge2$ and the
tangent cone is the triple plane $u^3$. The degree-four
term restricts to zero on this plane. After the ambient
blowup the entire reduced exceptional line is critical,
so none of its points is an isolated normal singularity.
If $C$ is smooth, take $x=h$. Every homogeneous term $f_j$
then has degree at least three in $u,x$. The cubic tangent
cone is a binary cubic in these two variables, hence three
lines, counted with multiplicities, through the tangent
direction of $C$.

\item First suppose that $H$ dominates a tangent-cone line.
For $(n,a)=(2,0)$ the normalization is generically an
isomorphism, and $H$ is birational to that line. For
$(n,a)=(1,1)$, a transverse normalization branch has conductor
exponent one. It is smooth and meets exactly one other smooth
branch transversely: the intrinsic conductor exponent of a
singular plane branch is at least two, and each other branch
contributes a positive intersection multiplicity. Thus the
generic transverse germ of $S_1$ is a node. A multiplicity-one
tangent line is generically smooth and cannot have $a=1$.
For a multiplicity-three line, the two geometric branches of
that node have exceptional orders whose weighted sum is three.
A connected double cover would give equal orders and an even
sum. Thus the order-one branch has degree one over the line
and is rational.

\item The remaining possibility is a double line together with a
distinct simple line. Here $C$ is smooth by \eqref{eq:II:triple-depressed-cubic}. Make a
linear change so that $f_3=xu^2$. On the $x$ chart of the
blowup put $u=xv,z=xw$. The strict equation is
\[
 v^2+x f_4(1,v,w)+x^2f_5(1,v,w)+\cdots.
\]
An ordinary double curve on $x=v=0$ requires
$f_4(1,0,w)=0$ identically. Its transverse quadratic form is
\begin{equation}
 v^2+xvB(w)+x^2C(w),
 \qquad \deg B\le1,\quad\deg C\le2.                  \label{eq:II:triple-tangent-quadratic}
\end{equation}
The degree bounds in \eqref{eq:II:triple-tangent-quadratic} follow from the degree-at-least-three
condition in $u,x$ on every $f_j$: here
$B=\partial_v f_4(1,0,w)$ and $C=f_5(1,0,w)$.
The normalization over the exceptional line adjoins a root
of $q^2+B(w)q+C(w)$. Its discriminant has degree at most
two. Write $B(w)=b_0+b_1w$ and
$C(w)=c_0+c_1w+c_2w^2$. For homogeneous coordinates
$[X:Z]$ on the exceptional line, put
\[
 B_h=X^2(b_0X+b_1Z),\qquad
 C_h=X^3(c_0X^2+c_1XZ+c_2Z^2),\qquad
 Q_h=XU^2+B_hU+C_h.
\]
Here $Q_h$ has weight five when $X,Z$ have weight one and
$U$ has weight two. On $X\ne0$ the substitutions $w=Z/X$
and $U=X^2q$ give
\[
 Q_h(X,Xw,X^2q)=X^5\bigl(q^2+B(w)q+C(w)\bigr).
\]
The homogeneous discriminant is therefore
\[
 B_h^2-4XC_h
 =X^4\bigl((b_0X+b_1Z)^2
             -4(c_0X^2+c_1XZ+c_2Z^2)\bigr).
\]
It contains the factor $X^4$; the residual homogeneous factor
has degree two. The resulting quadratic extension of $\mathbb C(w)$
is branched at at most two points, including infinity.
Riemann--Hurwitz makes its compact normalization rational;
the components of a split cover are rational as well.

\item Finally suppose that $H$ is exceptional over a point of $N_1$.
We must have $n=a=1$, and the first part of the proof would
make $S_1$ normal with an isolated simple elliptic singularity
there. This cannot happen. For singular $C$ the whole
exceptional line is critical, as already proved. For smooth
$C$ its strict transform passes through the common point of
the tangent lines, so that point is not an isolated normal
singularity. Away from it, at a line of multiplicity
$e=2$ or $3$, divide by the nonzero factor from any other
line. The strict equation has the form
\[
                         v^e+xG(v,w,x),
\]
where $G(0,w,0)$ is affine linear in $w$, by the same
normal-degree bound. If this function is nonzero at a point,
the surface is smooth. If it is identically zero, the entire
line is critical. At a simple zero, $G_w$ is a unit and the
exact coordinate change $W=G(v,w,x)$ gives $v^e+xW$, a
rational double point of type $A_{e-1}$. A multiplicity-one
line is smooth away from intersections. There is therefore
no isolated simple elliptic singularity. All possibilities
for $H$ have been excluded.

\end{enumerate}
\end{proof}


\Needspace{10\baselineskip}
\subsection{The smooth critical curve of transverse type \texorpdfstring{$A_3$}{A3}}

\begin{proposition}[Local models along a smooth tacnodal critical curve]\hypertarget{tag-02Y7}{}\label{night:P4R.14}\label{stmt:P4R.14}
\StatementRecord{night:P4R.14}{proved}{02Y7}{\cref{cdp:st:A.5.15,cdp:st:A.5.17,night:P4.28,lem:II:boundary-deformation-transport,lem:II:smooth-curve-intersection}}

Let $S$ be a reduced fibre support in the standing torus-fibration
setting of \cref{def:II:standing-torus-fibration}. At a point
where the entire reduced critical locus of its local defining
function is one smooth curve, with generic transverse type
$A_3$, holomorphic coordinates put that function in one of
the forms
\begin{equation}
 X^2+Y^4,\qquad X(Y^2+XZ).                              
\label{eq:II:manual:P4R.14.1}
\end{equation}
The normalization germs are respectively two smooth branches,
or a smooth plane and an $A_1$ quadric. No positive-genus
exceptional curve occurs over such a point.
\end{proposition}
\begin{proof}
\begin{enumerate}[label=\arabic*.,leftmargin=*]
\item Choose a coordinate $z$ on the critical curve and normal coordinates $x,y$. Its quadratic transverse Hessian has generic rank one. Diagonalization over $\C\{z\}$, with the coordinate maps of \cref{night:P4.28}, gives the quadratic term $h(z)x^2$. The cubic coefficient on the null direction vanishes identically. With weights two for $x$ and one for $y$, the weight-four terms are
\begin{equation}
 h(z)x^2+a(z)xy^2+b(z)y^4,\qquad
 \Delta(z)=4h(z)b(z)-a(z)^2.
\label{eq:II:manual:P4R.14.2}
\end{equation}
The remaining terms have larger weight, and $\Delta$ is generically nonzero.

\item We prove that $\Delta$ is a unit, keeping track of the irreducible factors of the full fibre. Every local component of the reduced support contains the critical axis: its intersection with another component has dimension at least one and lies on that axis. The generic transverse germ has two branches, so there are at most two local components.

First consider a locally irreducible reduced equation. The full map is $Uf^M$, and the exact root $f_0=U^{1/M}f$ satisfies $F=f_0^M$. Its critical germ is the same axis: every critical component through the origin has value zero, and on the zero set $df_0=U^{1/M}df$. Its quartic discriminant is the preceding $\Delta$ multiplied by the square of a unit on the axis. We may use this root for the following deformation.

A zero of $\Delta$ persists, counted with multiplicity, under small independent constant changes of $h,a,b$, by the argument principle on a small parameter disc. Generic such changes make its zeros simple and avoid $h=0$: the equations $h=\Delta=0$ require $a=0$, while for $h\ne0$ the derivative with respect to the constant in $b$ is $4h\ne0$. At a simple zero with $h\ne0$, parameter Morse coordinates give
\[
                X^2+y^4V(y,z),\qquad V_z(0,0)\ne0.
\]
The coordinate $Z=V(y,z)$ gives $X^2+y^4Z$. Its local sphere from \cref{cdp:st:A.5.15} has nonzero square. On the fixed boundary the root functions retain the product charts $X^2+Y^4$, by the parameter Morse map and a fourth root of a unit. Apply \cref{lem:II:boundary-deformation-transport} to these root functions, with the original exact equality $F=f_0^M$. The sphere reaches an original torus fibre, a contradiction. Thus the discriminant is a unit in this case.

\item In the reducible case write $f=\phi\psi$ with the two local reduced factors, absorbing the generator unit into $\psi$. Their common zero set is exactly the axis, and their generic transverse intersection multiplicity is two. Apply \cref{lem:II:smooth-curve-intersection} to these same factors and the same curve coordinate $z$. Its finite free intersection algebra has rank two and special-fibre length two; one factor, say $\phi$, is smooth. Its linear term on the axis is a unit times $x$. Write its weight-two term as $\alpha x+\beta y^2$ and that of $\psi$ as $\gamma x+\delta y^2$, with coefficients in $\C\{z\}$ and $\alpha$ a unit. On the smooth surface $\phi=0$, the only zero set of $\psi$ is $y=0$. Therefore the restricted function is a unit times $y^2$, its generic contact order being two. Its leading coefficient $\delta-\gamma\beta/\alpha$ is a unit. Comparing the product with the preceding quartic gives
\[
                 \Delta=-(\alpha\delta-\beta\gamma)^2,
\]
so $\Delta$ is a unit also here. This argument changes neither factor nor its full multiplicity.

\item If $h(0)\ne0$, parameter Morse coordinates and a fourth root of the residual unit give the first form in \eqref{eq:II:manual:P4R.14.1}. Suppose $h(0)=0$. Then $a(0)\ne0$. The function
\[
                    \frac{f(y^2v,y,z)}{y^4}
\]
is holomorphic and restricts at $y=0$ to $h(z)v^2+a(z)v+b(z)$. Its finite root $-b(0)/a(0)$ is simple. The implicit-function theorem gives a holomorphic root $v(y,z)$; division by its monic graph gives $X=x-y^2v(y,z)$ and $f=Xg$. On $X=0$, $g$ is $y^2$ times a unit. Its parameter Morse centre is zero there, and its critical value vanishes there. Parameter Morse coordinates in $y$, keeping $X,z$, therefore give exactly
\begin{equation}
                  f=X\bigl(Y^2+XH(X,z)\bigr).
\label{eq:II:manual:P4R.14.3}
\end{equation}
Here $H(0,z)$ is not identically zero, because the generic transverse Hessian has rank one, while $H(0,0)=0$. This factorization also shows that the locally irreducible case considered above cannot occur. Both forms obtained from a unit discriminant are reducible.

\item Write the full map in these coordinates as
\[
                   F=U X^a\bigl(Y^2+XH(X,z)\bigr)^b,
                   \qquad a,b\ge1.
\]
The two factors are the two local components. In the second form the second factor cannot split further, since every local component contains the axis and there are at most two. Choose $V=U^{1/b}$ and apply the parameter Morse construction in $Y$ to $V(Y^2+XH)$, keeping $X,z$. If its centre is $Y=\mu(X,z)$ and its quadratic unit is $A$, the new coordinate is $\widetilde Y=(Y-\mu)\sqrt A$, with inverse the unique local solution of this equation. The centre vanishes on $X=0$, and the remaining critical value is $XK(X,z)$. Thus the full function becomes exactly
\[
                         F=X^a(\widetilde Y^2+XK)^b.
\]
Differentiating that critical value in $X$ at $X=0$ gives the exact identity
\[
                    K(0,z)=V(0,0,z)H(0,z).
\]
In particular $K_z(0,0)\ne0$ is equivalent to $H_z(0,0)\ne0$. The reduced critical germ of $X(\widetilde Y^2+XK)$ is still the same axis. Indeed its reduced zero divisor is unchanged under this invertible map and multiplication by a unit, and every critical component through the origin has constant value zero, as above. We now rename $\widetilde Y$ as $Y$.

\item Put $N=a+b$ and $G_{a,b}=NK+bXK_X$. Suppose $K_z(0,0)=0$. The germ $K_z$ is nonzero and not divisible by $X$, since $K(0,z)$ has finite order at least two. We prove that $(K_z,G_{a,b})$ has an isolated zero. On a parametrized irreducible branch of $K_z=0$, choose a parameter with $X=t^r$, $r\ge1$. The chain rule there gives
\[
                       XK_X=\frac{t}{r}\frac{dK}{dt}.
\]
If $G_{a,b}$ vanished identically, every coefficient of order $j$ in $K(t)$ would be multiplied by $N+bj/r>0$. Hence $K$ would vanish identically on that branch, and then so would $K_X$. This would give a critical-curve branch of the reduced function off $X=0$, contradicting the hypothesis. The intersection is therefore isolated and has positive local degree.

Deform $K$ to $K_s=K-\eta z-\varepsilon$, retaining the full function $F_s=X^a(Y^2+XK_s)^b$. Its off-divisor critical equations are
\[
       Y=0,\qquad K_z-\eta=0,\qquad
       G_{a,b}-N\eta z-N\varepsilon=0.
\]
The last two are the fibres of
\[
        \Psi_{a,b}(X,z)=
              \left(K_z,\frac{G_{a,b}-NzK_z}{N}\right).
\]
Its isolated zero has positive local degree. On a small compact representative this degree persists for small parameters. Choose a regular parameter value outside the images of $X=0$ and $K_X=0$. These images have real four-dimensional measure zero, by their countable smooth two-dimensional strata. The germ $K_X$ is not identically zero: otherwise $K$ would depend only on $z$ and $(K_z,G_{a,b})$ would vanish on the line $z=0$. Thus the exceptional images do not fill the open parameter ball, and the chosen fibre is nonempty and finite.

\item Write $G_s=G_{a,b}-N\eta z-N\varepsilon$ and
\[
 J=\det\frac{\partial(G_s,K_z-\eta)}{\partial(X,z)}.
\]
At a chosen solution $K_s=-bXK_X/N\ne0$, so
\[
            F_s=\left(-\frac bN\right)^b
                         X^{a+2b}K_X^b\ne0.
\]
At $Y=0$, the logarithmic derivatives are
\[
 (\log F_s)_X=\frac{G_s}{XK_s},\qquad
 (\log F_s)_z=\frac{b(K_z-\eta)}{K_s},\qquad
 (\log F_s)_{YY}=\frac{2b}{XK_s}.
\]
The mixed $Y$ entries vanish at the critical point. The determinant of the $X,z$ logarithmic Hessian block is $bJ/(XK_s^2)$. Since the full gradient vanishes there, this gives
\[
                 \det\operatorname{Hess}F_s
                   =\frac{2b^2F_s^3J}{X^2K_s^3}\ne0.
\]
These are holomorphic Morse points with nonzero critical values, in the full original multiplicities.

For $a=b=1$ and $K=H$, this calculation gives the original equations, with $G=2H+XH_X$,
\begin{equation}
                H_z-\eta=0,\qquad G-2\eta z-2\varepsilon=0,
\label{eq:II:manual:P4R.14.4}
\end{equation}
the map $\Psi=(H_z,(G-2zH_z)/2)$, and the critical value
\begin{equation}
                       f_{\eta,\varepsilon}=-\frac{X^3H_X}{2}.
\label{eq:II:manual:P4R.14.5}
\end{equation}
Its $Y$-Hessian entry is $2X$, its remaining block is
\[
 \begin{pmatrix}
 X(3H_X+XH_{XX})&X^2H_{Xz}\\
 X^2H_{Xz}&X^2H_{zz}
 \end{pmatrix},
\]
and its full determinant is $2X^4J$.

\item Choose a sufficiently small ball transverse to the smooth strata of the full zero divisor. Its original critical values are all zero, by the same critical-component argument. On its fixed boundary near the critical axis, $K_s(0,z)$ is a unit. The coordinate $T=XK_s(X,z)$, keeping $Y,z$, is invertible: its inverse is $X=\chi_s(T,z)$, the unique local solution of $T=XK_s(X,z)$. Write $\chi_s=T\kappa_s$ with $\kappa_s$ a unit. Then
\[
                   F_s=\kappa_s^a T^a(Y^2+T)^b.
\]
Weights $(2,1)$ in $(T,Y)$, of degree $2N$, remove this unit by multiplying $T$ by $(\kappa_s^a)^{1/N}$ and $Y$ by $(\kappa_s^a)^{1/(2N)}$, keeping $z$. Its inverse solves
\[
 \tau^{2N}=\kappa_s(\tau^{-2}T',z)^a,
 \qquad(T,Y,z)=(\tau^{-2}T',\tau^{-1}Y',z).
\]
The implicit derivative at the link is $2N\tau^{2N-1}\ne0$. This gives the fixed full transverse polynomial $T^a(Y^2+T)^b$. Its weighted Euler identity gives only critical value zero. The tangential coordinate is transverse to the chosen boundary. At smooth component points the corresponding fixed charts are a power of a normal coordinate. Hence \cref{lem:II:boundary-deformation-transport} applies to this full family with $m=1$. The Morse sphere constructed in \cref{cdp:st:A.5.17} has square $-2$ and transports into an original torus fibre, a contradiction.

Therefore $K_z(0,0)\ne0$, and the displayed relation between $K$ and $H$ gives $H_z(0,0)\ne0$. The coordinate $Z=H(X,z)$, keeping $X,Y$, puts the original reduced generator into the second form of \eqref{eq:II:manual:P4R.14.1}; its inverse is the unique solution of $H(X,z)=Z$. Its normalization consists of the plane $X=0$ and the normal quadric $Y^2+XZ=0$, whose minimal resolution has one rational $(-2)$-curve. The first form has the two smooth factors $X\pm\mathrm iY^2$. No positive-genus exceptional curve occurs.

\item The following reduced hypersurface has pole coefficient two
and a positive-genus exceptional curve:
\[
 f=x\bigl[y^2+xz(x-z)(x-2z)\bigr].
\]
It has only the smooth $A_3$ critical axis $x=y=0$. The
second normalization branch is the double plane branched
over four distinct lines. Blowing up their common point
and normalizing gives a double cover of the exceptional
$\mathbb P^1$ branched in four points, hence an elliptic
curve of square $-2$. Its discrepancy is $-1$ and the
order of $x$ on it is one, giving pole coefficient two
for the displayed total reduced hypersurface. Its polar
deformation excludes a torus smoothing, as follows.
For $H=z(x-z)(x-2z)$, the deformation $H\mapsto H-1$ has
the exact Gr\"obner basis
\[
 3x-5z^4-14z,\qquad 25z^6+50z^3-2.
\]
It gives six simple solutions; adjoining $x$, $H_x$ or
$J$ to the critical ideal gives the unit ideal in each
case. For $H\mapsto H-\varepsilon$ all six scale by
$\varepsilon^{1/3}$ and converge to the origin.

\end{enumerate}
\end{proof}


\subsection{A nonzero Hessian and the entire critical curve}

\begin{proposition}[Product germs from a nonzero Hessian]\hypertarget{tag-02Y8}{}\label{night:P4R.15}\label{stmt:P4R.15}
\StatementRecord{night:P4R.15}{proved}{02Y8}{\cref{cdp:st:A.5.17,prop:II:smooth-branch-contact}}

At every critical-curve point of a reduced fibre support in the
standing torus-fibration setting of
\cref{def:II:standing-torus-fibration} with nonzero ambient Hessian,
holomorphic coordinates
give the exact product germ
\[
                       X^2+Y^m,\qquad m\ge2.
\]
In particular its entire reduced critical curve is smooth.
At a singular point of the reduced critical curve, the ambient
Hessian vanishes. This assertion allows different generic
transverse types on the different branches before the exclusion.
\end{proposition}
\begin{proof}
A nonzero symmetric Hessian has a direction of nonzero quadratic value. Parametrized holomorphic Morse coordinates give the reduced generator $f=X^2+g(Y,Z)$, where $g\ne0$ because $f$ is reduced. The full local map is a unit times the product of its irreducible factors with their original positive multiplicities.

If $f$ is locally irreducible, the full map is $Uf^M$. The root $H=U^{1/M}f$ satisfies $F=H^M$ exactly, and at the critical point $\operatorname{Hess}H=U(0)^{1/M}\operatorname{Hess}f$. Apply the nonzero-Hessian conclusion of \cref{cdp:st:A.5.17}. It gives $H=X_1^2+Y_1^m$ with $m\ge2$, after multiplying the transverse coordinate by a constant whose $m$th power is $-1$. In these coordinates $f=V(X_1^2+Y_1^m)$ for a unit $V$. The weighted coordinate map
\[
       (X_1,Y_1,Z_1)\longmapsto
          (V^{1/2}X_1,V^{1/m}Y_1,Z_1)
\]
gives the exact stated generator. Its inverse is obtained by solving
\[
 \tau^{2m}=V(\tau^{-m}X',\tau^{-2}Y',Z'),
 \qquad (X_1,Y_1,Z_1)=(\tau^{-m}X',\tau^{-2}Y',Z'),
\]
with the selected nonzero root; the implicit derivative is $2m\tau^{2m-1}\ne0$ at the origin.

If $f$ is reducible, Weierstrass factorization of its monic quadratic in $X$ gives two monic linear factors. Their $X$ coefficients sum to zero, so $g=-H(Y,Z)^2$ and the factors are exactly $X-H$ and $X+H$. Both are smooth. The full map is therefore $U(X-H)^a(X+H)^b$ for its original $a,b$. Apply \cref{prop:II:smooth-branch-contact}. Its full-function transport retains these multiplicities and proves that $H=h^k$ with $h$ smooth after absorbing its plane unit. Thus the original reduced generator is exactly $X^2-Y^{2k}$; multiplying $Y$ by a chosen constant whose $2k$th power is $-1$ gives $X^2+Y^{2k}$.

Both cases give $m\ge2$ and reduced critical curve $X=Y=0$. The factor and gradient deformations used in the two earlier owners permit the different initial exponents on the different branches, so no equality of their generic transverse types has been imposed. The explicit Hessian and moving-node calculation for their uniform-power suspension is recorded in the proof of \cref{cdp:st:A.5.17}.
\end{proof}


\subsection{Local intersection forms and the quartic tangent obstruction}

\begin{lemma}[The conductor graph at positive Euler number]
\hypertarget{tag-P314}{}\label{lem:II:positive-ruled-conductor-graph}
\StatementRecord{lem:II:positive-ruled-conductor-graph}{proved}{P314}{\cref{night:P4.26,night:P4.35,night:P4R.13}}
\hypertarget{tag-P203}{}\label{scope:II:positive-ruled-configuration}
\StatementRecord{scope:II:positive-ruled-configuration}{proved}{P203}{}
In the irreducible-or-normal-crossings torus-fibration setting of
\cref{def:II:standing-torus-fibration}, let $mS$ have irreducible
reduced support with $e(S)>0$, and suppose that its normalization
has a resolution ruled over a curve of positive genus.
Write $\rho=\nu\sigma:T\to N\to S$ for the minimal resolved
normalization. Its pole divisor $D$ has horizontal part $2H$,
and $A=D_{\rm red}=\rho^{-1}(C)_{\rm red}$ is the section $H$
with rational trees attached. The inclusion $A\hookrightarrow T$
induces an isomorphism on fundamental groups, and the map
$T\cup_A C\to S$ induced by $\rho$ and $C\hookrightarrow S$
is a homeomorphism. The reduced conductor in $N$ has an acyclic
normalization-incidence graph.
The reduced conductor image
$C=\operatorname{Sing}(S)_{\rm red}$ has exactly two rational
normalized components and one point-identification cycle.
The inclusion $C\hookrightarrow S$ induces an isomorphism
$\pi_1(C)\simeq\pi_1(S)\simeq\mathbb Z$. Moreover,
\[
 e(S)=2,\qquad e(S_{\rm reg})=0,\qquad
 \sum_{p\in C}(r_p-1)=2,
\]
where $r_p$ is the number of analytic branches of $C$ at $p$.
\end{lemma}
\begin{proof}
\begin{enumerate}[label=\arabic*.,leftmargin=*]
\item Apply \cref{night:P4.26} to the given full fibration.
Its nonnegative fibre Euler numbers sum to $c_3(Y)\le2$, so
$0<e(S)\le2$. The support is singular: a smooth ruled support
would have torsion canonical bundle by \cref{night:P4.12},
whereas its canonical degree on a general ruling fibre is $-2$.
Write $\rho=\nu\sigma$ for the normalization and its minimal
resolution, and let $D$ be the pole divisor of
\cref{night:P4R.6}. Apply
\cref{lem:II:irrational-ruled-boundary} to the same $T,D$ and
minimal ruling. Its reduced-horizontal case is excluded by
\cref{night:P4R.7}, which gives $e(S)=0$ in that case.
Thus $D$ has horizontal part $2H$, with rational trees in its
remaining support $A=D_{\rm red}$.

\item The boundary lemma gives an isomorphism
$j_*:\pi_1(A)\to\pi_1(T)$ for the inclusion $j:A\hookrightarrow T$.
The support equality in \cref{night:P4R.6} is
$A=\rho^{-1}(C)_{\rm red}$, including every exceptional component.
Apply the conductor pushout \cref{lem:II:proper-conductor-pushout}
with $(X,Y,A,C,h)=(T,S,A,C,\rho)$, using a compatible analytic
triangulation. Its quotient map is exactly the displayed
homeomorphism $T\cup_A C\to S$. It gives the isomorphism induced by $C\hookrightarrow S$;
its inverse is the identity on $\pi_1(C)$ and
$(\rho|_A)_*j_*^{-1}$ on the image of $\pi_1(T)$.
Proposition~\ref{night:P4R.1}, applied to a small full tube of
this same $mS$, makes the resulting group virtually abelian.

\item Let $c$ be the number of normalized components of $C$,
let their genera be $g_i$, and let $b$ be its incidence-graph
cycle rank. The connected curve and maps just obtained satisfy
\cref{lem:II:curve-normalization-graph}. Any elliptic normalized
component would satisfy every hypothesis of \cref{night:P4R.13}
for this $f,Y,B,mS$ and ruled normalization resolution, giving
$e(S)=0$. Thus the curve-group alternatives leave all $g_i=0$
and $b\le1$. In particular
\[
 \pi_1(S)\simeq F_b,\qquad e(C)=c+1-b.
\]

\item If $b=0$, the cyclic quotient of \cref{night:P4R.1}
makes $m=1$. Its formula for $b_1(S)$ gives
$\det(I-A_{\rm mon})\ne0$ for the positive small-loop
monodromy. Proposition~\ref{night:P4.13}, with this same $m=1$,
then gives the positive integer
$e(S_{\rm reg})=\det(I-A_{\rm mon})\ge1$.
But $e(C)=c+1\ge2$, so $e(S)\ge3$, a contradiction.
Therefore $b=1$ and $\pi_1(C)\simeq\mathbb Z$.
Now $e(C)=c$, and the nonnegative smooth-locus Euler number in
\cref{night:P4.13} gives $c\le e(S)\le2$.

\item The resolved pole support is the section with rational trees
from the first item. The final ruled application in the proof of
\cref{night:P4.35} contracts precisely the connected exceptional
subtrees of $\sigma$ and shows that the reduced conductor in $N$
has no graph cycles. The same proposition therefore excludes
$c=1$, since $b=1$. Hence $c=2$. The inequalities already proved
give $2=e(C)\le e(S)\le2$, so $e(S)=2$ and $e(S_{\rm reg})=0$.
The same normalization-incidence formula is
$b=1-c+\sum_p(r_p-1)$; substitution of $b=1,c=2$ gives the
last displayed equality. These are the maps of the positive
ruled configuration with $D_T=A$ and $D_N=\sigma(A)_{\rm red}$.
\end{enumerate}
\end{proof}

\begin{proposition}[Vanishing local intersection forms from a nonzero first Betti number]\hypertarget{tag-02XQ}{}\label{night:P4.51}\label{stmt:P4.51}
\StatementRecord{night:P4.51}{proved}{02XQ}{}

Let $f:Y\to B$ satisfy the torus-fibration hypotheses of
\cref{night:P4.12}, and let $mS$ be a fibre with irreducible
reduced support satisfying $b_1(S)>0$. Every local
Milnor fibre of a reduced local equation of $S$ has
identically zero rational intersection form. In particular, consider a fibration all of whose
reduced supports are irreducible or have normal crossings.
Every irreducible support $S$ with $e(S)>0$ whose
normalization has a resolution ruled over a curve of
positive genus has $\pi_1(S)\simeq\mathbb Z$, and the same
local vanishing conclusion applies to it.
\end{proposition}
\begin{proof}

\begin{enumerate}[label=\arabic*.,leftmargin=*]
\item Let $U$ be a sufficiently small full tube and $F_t$ a
nearby compact smooth torus fibre. The inclusion and neighbourhood
retraction in the proof of \cref{night:P4R.1} identify the
homomorphism
\[
 \pi_1(F_t)=\mathbb Z^4\longrightarrow
 \pi_1(U)=\pi_1(S)
\]
and give a finite cyclic cokernel, of order equal to
the fibre multiplicity. Its map on rational first
homology is therefore onto. Write
\[
 W=\ker\bigl(H_1(F_t;\mathbb Q)\longrightarrow
                  H_1(S;\mathbb Q)\bigr).
\]
Then $\dim W=4-b_1(S)\le3$.

\item Choose a coordinate ball about the point, with its
closed ball contained in $U$. The local full fibre
equation is $u g^m$, with $g$ reduced and $u$ a
holomorphic unit. Choose $v=\exp((\log u)/m)$ and put $h=vg$,
so the exact equation is $f=h^m$. The nearby fibre in the ball
is the disjoint union of $h=\zeta t^{1/m}$.
To compare the original reduced equation, put $g_s=\exp(s\log v)g$.
On a nonzero level, the vector field
\[
 Z_s=-\frac{(\log v)g_s\,\overline{\nabla g_s}}{\|\nabla g_s\|^2}
\]
has flow $\Phi_s$ satisfying $g_s\circ\Phi_s=g$.
We justify the required uniform bound. First,
\[
 |g(z)|\le C\|z\|\,\|\nabla g(z)\|
\]
near the origin. We use real analytic curve selection,
which follows from the rectilinearization maps in
\cite[Theorem~0.2, p.~6]{BierstoneMilman1988} as follows.
A sequence in a semianalytic set $X$ tending to a boundary
point lifts, after taking a subsequence, to a convergent
sequence in one of the compact sets $L_i$ of that theorem.
Its limit $y_0$ lies in the closure of one quadrant
$Q\subset\varphi_i^{-1}(X)$. For $y_1\in Q$, the map
$t\mapsto\varphi_i(y_0+t(y_1-y_0))$ is the required
real analytic arc. If the displayed bound failed, applying
this construction to the semianalytic set
\[
 0<|g(z)|,\qquad
 \|z\|^2\|\nabla g(z)\|^2<\epsilon^2|g(z)|^2,
 \qquad\epsilon>0,
\]
would give a real analytic arc $(z(t),\epsilon(t))$ tending
to $(0,0)$ on which the quotient
$\|z\|\,\|\nabla g\|/|g|$ tends to zero. Set
$r=\operatorname{ord}_t z$, $q=\operatorname{ord}_t\nabla g(z(t))$
and $k=\operatorname{ord}_t g(z(t))$. The chain rule gives
$k-1\ge q+r-1$, whereas the strict inequality on the arc
gives $q+r>k$, a contradiction. An identically zero gradient
on the arc would instead make $g(z(t))$ identically zero,
also a contradiction.

Put $L=\log v$ and let $M_L$ bound $\|\nabla L\|$ on the ball.
The exact derivative
\[
 \nabla g_s=e^{sL}(\nabla g+s g\nabla L)
\]
and $CM_L\|z\|\le\tfrac12$ give, uniformly for $0\le s\le1$,
\[
 \frac{|g_s|}{\|\nabla g_s\|}
 \le2C\|z\|,
 \qquad \|Z_s\|\le2C\sup|L|\,\|z\|.
\]
In particular the gradient is nonzero on every nonzero
level in the smaller ball. Thus a sufficiently small compact Milnor
fibre of $g$ flows into a Milnor fibre of $h$ in a slightly
larger ball still contained in $U$. The flow and its inverse
are smooth on these nonzero levels and preserve their complex
orientations by continuous transport from $s=0$. Hence the
original local fibre $M$ has an orientation-preserving
codimension-zero embedding $j:M\hookrightarrow F_t$.
Its inclusion into $U$ factors through the contractible
ball. Thus every loop of $M$ maps to zero in
$H_1(U;\mathbb Q)$, and the image subgroup
\[
 H=\operatorname{im}\bigl(\pi_1(M)\longrightarrow
                                 \mathbb Z^4\bigr)
\]
satisfies $H\otimes\mathbb Q\subset W$.

\item The covering criterion gives a lift $\widetilde j:M\to F_H$
to the connected cover $p_H:F_H\to F_t$ corresponding to $H$,
with $p_H\circ\widetilde j=j$. Writing $F_t=\mathbb R^4/\Lambda$,
this cover is $\mathbb R^4/H$, with $p_H([x]_H)=[x]_\Lambda$.
A basis of $H$ identifies its real
span modulo $H$ with a torus and the complementary factor with
a real vector space. Thus $F_H$ is a $K(H,1)$, and the map
$(p_H)_*$ on second rational homology is exactly
\[
 \bigwedge^2(H\otimes\mathbb Q)\longrightarrow
              \bigwedge^2 H_1(F_t;\mathbb Q).
\]
Every class from $H_2(M;\mathbb Q)$ consequently maps
into $\bigwedge^2W$. The torus intersection pairing is
exterior multiplication followed by the orientation
isomorphism $\bigwedge^4H_1(F_t;\mathbb Q)=\mathbb Q$.
It vanishes on $\bigwedge^2W$ because $\dim W\le3$.

\item This ambient pairing is the local intersection
pairing on $M$. We represent finitely many cycles
inside its interior and make them transverse there;
the orientation-preserving inclusion into $F_t$
retains all intersection points and signs. Thus the
local pairing is zero.

\item For the ruled application, the hypotheses of
\cref{lem:II:positive-ruled-conductor-graph} hold for this
$f:Y\to B$ and this full fibre $mS$: all reduced supports
are irreducible or normal crossing, $e(S)>0$, and its
normalization resolution is ruled over a positive-genus
curve. The lemma identifies $\pi_1(S)$ with $\mathbb Z$
through the conductor inclusion. Hence $b_1(S)=1$, and
the preceding four items apply to every reduced local
equation of this same $S$.

\end{enumerate}
\end{proof}



\subsection{The excess-one curve calculation and the ordinary-double sector}

\begin{lemma}[Reduced curve germs with one excess condition]\hypertarget{tag-P206}{}\label{ii:excess-one-curve-germs}
\StatementRecord{ii:excess-one-curve-germs}{proved}{P206}{}

Let $A$ be the local analytic ring of a reduced complex curve
germ with $1\le r\le3$ branches and $\delta(A)=r$.
Let $B=\bigoplus_{i=1}^r\mathbb C\{t_i\}$ be its normalization
and let $R\subset B$ consist of the tuples with equal constant
terms. Then $A$ is the inverse image in $R$ of a hyperplane
in the space of branch first derivatives. After rescaling
the branch parameters, its analytic ring has one of the
following forms:
\begin{enumerate}[label=\alph*.,leftmargin=*]
\item a plane cusp with $r-1$ independent smooth branches
joined at their constant terms;
\item a plane tacnode with $r-2$ independent smooth branches
joined at their constant terms, for $r\ge2$;
\item an ordinary plane triple point, for $r=3$.
\end{enumerate}
The embedding dimensions in the first two cases are $r+1$
and $r$, respectively. In particular the first case with
$r=3$ does not embed in a smooth threefold. For every model
which does embed, the analytic ring isomorphism is induced
by a local change of ambient coordinates.
\end{lemma}
\begin{proof}
\begin{enumerate}[label=\arabic*.,leftmargin=*]
\item The quotient by equal constant terms has length $r-1$,
so
\[
 \operatorname{length}(R/A)=\delta(A)-(r-1).
\]
This length is one. The simple $A$-module $R/A$
is annihilated by $\mathfrak m_A$, so
$\mathfrak m_A R\subset A$. Choose
$v\in\mathfrak m_R$ spanning $R/A$ and write
$v^2=a+cv$, where $a\in\mathfrak m_A$ and $c\in\mathbb C$.
If $c\ne0$, the convergent square-root series solves
$b^2-cb-a=0$ with $b\in\mathfrak m_A$. Then
$(v-b)(v+b-c)=0$ and the second factor is a unit, forcing
$v=b$, a contradiction. Thus $v^2\in A$ and
\[
                         \mathfrak m_R^2\subset A.
\]
Consequently $A$ is exactly the inverse image of a hyperplane
\[
                         \sum_i a_i f_i'(0)=0
\]
in $\mathfrak m_R/\mathfrak m_R^2=\mathbb C^r$.
The coordinate change $u_i=\lambda_i t_i$ replaces $a_i$
by $a_i\lambda_i$ in this equation. Choosing nonzero
$\lambda_i$ makes the nonzero coefficients any specified
nonzero constants and gives an analytic isomorphism of the
normalization which carries the indicated subring to the
new one.

\item With one nonzero coefficient, that branch has ring
$\mathbb C\{t^2,t^3\}$ and the others are independent smooth
branches joined at their constant terms. Indeed a convergent
series with zero linear coefficient is a sum of a constant,
$t^2h_0(t^2)$ and $t^3h_1(t^2)$. The embedding
dimension is $r+1$: for $r=1$ this is a plane cusp, for
$r=2$ a cusp and transverse line, and for $r=3$ it does
not embed in a smooth threefold.

With two nonzero coefficients, choose them to be $1,-1$,
so the pair has equal first derivatives. Its
ring is generated by $x=(t,t)$ and $y=(t^2,-t^2)$, giving the
plane tacnode $y^2=x^4$. More explicitly, a pair $f_1,f_2$
with the same constant and linear terms is the restriction of
$G(x)+yH(x)$, where
$G(x)=(f_1(x)+f_2(x))/2$ and
$H(x)=(f_1(x)-f_2(x))/(2x^2)$ are convergent.
The kernel of this restriction is $(y^2-x^4)$:
division by that monic polynomial reduces a function to
$G(x)+yH(x)$, and its restrictions $G\pm x^2H$ both vanish
only when $G=H=0$. The other $r-2$ branches are
independent transverse lines, each contributing one further
generator to the cotangent space. The embedding dimension
is therefore $r$.

For $r=3$ and three nonzero coefficients, choose them to be
$1,1,-1$. The first jets are then the restrictions of linear
forms to the lines parametrized by $(t,0)$, $(0,t)$ and
$(t,t)$. Their products span all three second-jet coordinates:
$x^2-xy$, $y^2-xy$ and $xy$ restrict to $(t^2,0,0)$,
$(0,t^2,0)$ and $(0,0,t^2)$, respectively. Their multiples
by convergent functions of $x$, $y$ and $x$, respectively,
give all independent branch terms of order at least two.
Thus the ring is exactly the ordinary plane triple point,
with kernel $(xy(x-y))$. Equivalently,
the latter ring has precisely the same constant and
first-derivative conditions and colength three in $B$.
These cases exhaust $r\le3$.

The analytic ring isomorphisms lift to ambient coordinate
changes. In embedding dimension three, choose minimal
algebra generators and lift them to ambient functions;
their independent classes modulo the square of the maximal
ideal give an invertible linear part, so the inverse-function
theorem supplies the coordinate change. In embedding
dimension two, a defining equation of the containing ideal
has nonzero linear part; first straighten that smooth
containing surface, and then apply the same argument in
its two coordinates. This proves the asserted ambient forms.
\end{enumerate}
\end{proof}

\begin{proposition}[Exclusion of ordinary double curves at positive Euler number]\hypertarget{tag-02XE}{}\label{night:P4.40}\label{stmt:P4.40}\label{night:P4.40:positive-euler}
\StatementRecord{night:P4.40}{proved}{02XE}{}

In the standing torus-fibration setting of
\cref{def:II:standing-torus-fibration}, suppose every reduced
fibre support is irreducible or has normal crossings.
No irreducible support $S$ with $e(S)>0$ and a
positive-genus ruled normalization resolution has generic
transverse type $A_1$ on every component of its critical curve.
\end{proposition}
\begin{proof}
\begin{enumerate}[label=\arabic*.,leftmargin=*]
\item Suppose such a support exists. Apply
\cref{lem:II:positive-ruled-conductor-graph} with the given
$f:Y\to B$, its full fibre $mS$ and its minimal normalization
resolution $\rho:T\to S$. The global support hypothesis,
$e(S)>0$ and the positive ruling genus are exactly its
hypotheses. The pole divisor has a doubled horizontal
section $H$ with rational trees attached, and the image
conductor $C$ has two rational normalized components and
one point-identification cycle.

The section $H$ is contracted by the normalization resolution:
otherwise its generic conductor exponent would be two by
\cref{night:P4R.6}, whereas the reduced transverse node has
exponent one on each branch by
\cref{lem:II:plane-branch-conductor}. Apply the contracted-section
clause of \cref{night:P4.33} to this same divisor and resolution.
It gives $\chi(\mathcal O_{C_{\rm sch}})=-1$.
The image conductor is Cohen--Macaulay of dimension one and
generically reduced in type $A_1$. Its nilradical would have
finite support, so depth one makes it zero. Thus
$C_{\rm sch}=C$ and $\chi(\mathcal O_C)=-1$.

\item Write $r_p$ for the number of normalized preimages of
$p\in C$ and $\delta_p$ for its normalization-quotient length.
The two rational components give
$\chi(\mathcal O_C)=2-\sum_p\delta_p$.
The one-cycle incidence graph gives
$\sum_p(r_p-1)=2-1+1=2$. Therefore
\[
 \sum_p\delta_p=3,\qquad
 \sum_p(r_p-1)=2,\qquad
 \sum_p\bigl(\delta_p-(r_p-1)\bigr)=1.
\]
Each last summand is nonnegative: taking independent
branch values modulo the common value embeds a vector space
of dimension $r_p-1$ into the normalization quotient.
Hence exactly one point has excess one, and $r_p\le3$
at every point. Apply \cref{ii:excess-one-curve-germs} to
that reduced curve germ in the smooth threefold $Y$.
Its analytic isomorphisms give a singular plane curve,
a plane cusp with a transverse smooth branch, or a plane
tacnode with a transverse smooth branch. Each is excluded
by \cref{night:P4.38}, applied to the same reduced
hypersurface germ and its assumed generic transverse
type $A_1$ on every critical-curve component. This is
the required contradiction.
\end{enumerate}
\end{proof}

\subsection{The uniform cuspidal sector}

\begin{proposition}[Euler bound for uniform cuspidal ruled supports]\hypertarget{tag-02XH}{}\label{night:P4.43}\label{stmt:P4.43}
\StatementRecord{night:P4.43}{proved}{02XH}{\SeriesRef{night:P4R.13}, \SeriesRef{night:P4R.7}}

In the irreducible-or-normal-crossings torus-fibration setting
of \cref{def:II:standing-torus-fibration}, an irreducible
reduced fibre support whose normalization has a resolution
ruled over a positive-genus curve cannot have $e(S)>0$ and
generic transverse type $A_2$ on every component of its
critical curve.
\end{proposition}
\begin{proof}

\begin{enumerate}[label=\arabic*.,leftmargin=*]
\item We use the exact normalization map
$\nu:N\longrightarrow S$ and its minimal resolution
$\rho:T\longrightarrow N$. A germ of $S$ cannot have
two distinct analytic hypersurface branches. Indeed,
their intersection has dimension one in the smooth
ambient threefold. A general point of an irreducible
component of that intersection has at least two local
hypersurface branches, so its transverse curve germ is
reducible. The transverse cusp $A_2$ is irreducible.
This contradicts the stated uniform transverse type.
The points of the finite normalization over a germ are
its analytic branches. Consequently $\nu$ is bijective.
It is proper, so its continuous map from a compact space
to the Hausdorff space $S$ is a homeomorphism.

\item Apply \cref{lem:II:positive-ruled-conductor-graph}
to this $f:Y\to B$ and full fibre $mS$. The statement gives
its global support hypothesis, $e(S)>0$ and the positive
ruling genus. It gives the doubled horizontal section $H$
and identifies $\pi_1(S)$ with $\mathbb Z$ through
$C\hookrightarrow S$. Let $B_0$ be the base curve of the
minimal ruled surface, of genus $g\ge1$. Point
blowups preserve the fundamental group, and the ruling
gives $\pi_1(T)=\pi_1(B_0)$. By the fibre-tree calculation
in \cref{night:P4R.7}, every vertical exceptional configuration is a
tree of rational curves. Collapsing any such connected
tree preserves the fundamental group: apply van Kampen
to its tubular neighbourhood and the complement, or
equivalently to the cofibration quotient by that
simply connected tree. These quotients are the topological contractions of the resolution map.

\item If $H$ is not contracted by $\rho$, all resolution
exceptional configurations are of this vertical kind.
It follows that $\pi_1(N)=\pi_1(B_0)$. If $H$ is
contracted, its connected exceptional configuration
contains $H$. The inclusion $H\hookrightarrow T$
induces an isomorphism on fundamental groups, since
the ruling restricts to the identity section and all
point blowups preserve that isomorphism. The quotient
which collapses this configuration therefore kills
the whole group. The other contracted configurations
are rational trees, giving $\pi_1(N)=1$.

\item The normalization homeomorphism from the first item
therefore identifies $\pi_1(N)$ with $\mathbb Z$.
This contradicts the preceding resolution quotients:
their abelianizations have ranks zero or $2g$, whereas
the abelianization of $\mathbb Z$ has rank one.

\end{enumerate}
\end{proof}


\subsection{The ordinary-node defect in transverse type \texorpdfstring{$A_3$}{A3}}

\begin{lemma}[Quadratic conductor algebras and their local fibres]\hypertarget{tag-P257}{}\label{lem:II:quadratic-conductor-normalization}
\StatementRecord{lem:II:quadratic-conductor-normalization}{proved}{P257}{\cref{lem:II:conductor-duality}}
Let $R$ be a reduced complex analytic local ring with maximal ideal
$\mathfrak m$ and finite normalization $B$. A local fibre algebra
$B/\mathfrak mB$ makes $R$ a domain.

More specifically, let $S$ be the regular local analytic ring of a
smooth threefold, let $G,H$ be a regular sequence in its maximal ideal,
and let
\[
 R=S/(f),\qquad f=G^2+UH^2
\]
be reduced. Suppose its conductor ideal in $R$ is $I=(G,H)$.
The finite-duality identification with the total quotient ring gives
\[
 B=\operatorname{Hom}_R(I,R)=R+Rt,\qquad
 t=G/H,\qquad t^2=-U.
\]
For $U(0)=0$, the reduced equation $f$ is irreducible. For $U(0)\ne0$,
the equations $G-H\sqrt{-U}$ and $G+H\sqrt{-U}$ are its two
irreducible factors.
\end{lemma}
\begin{proof}
Every maximal ideal of the finite integral algebra $B$ lies over
$\mathfrak m$, and hence corresponds to one of $B/\mathfrak mB$.
A local fibre makes $B$ local. A normal local ring is a domain,
and the normalization inclusion $R\hookrightarrow B$ then makes
$R$ a domain as well.

For the displayed hypersurface, a common prime divisor of $H$ and
$f$ would divide $G$, contrary to the regular sequence. Thus $H$
is a nonzerodivisor in $R$. The finite-duality multiplier map and
Gorenstein biduality in \cref{lem:II:conductor-duality} identify
$B$ with $\operatorname{Hom}_R(I,R)$, every homomorphism being
multiplication in the total quotient ring. The equations
$tH=G$, $tG=-UH$ and $t^2=-U$ give $R+Rt\subset B$.

Conversely, let a multiplier $\theta$ satisfy $\theta H=r$ and
$\theta G=s$ in $R$. Lifting $r,s$ to $S$, the relation
$rG-sH=0$ in $R$ is
\[
 rG-sH=v(G^2+UH^2),\qquad G(r-vG)=H(s+vUH).
\]
The regular sequence gives $r-vG\in(H)$, so division by the
nonzerodivisor $H$ yields $\theta=vt+w$ with $v,w\in R$.
This proves the displayed algebra equality. Its nonzero fibre is
a quotient of $\mathbb C[t]/(t^2+U(0))$. For $U(0)=0$ it is local,
and the first assertion applies. For $U(0)\ne0$, the two displayed
nonunit factors have no common irreducible factor, since a common
factor would divide both $G$ and $H$. There are at most two
normalization points by the fibre algebra and at least one for
each irreducible factor. Consequently each displayed factor is
irreducible and these are all the factors.
\end{proof}

\begin{proposition}[Hessian defect at a nodal critical curve]\hypertarget{tag-02XF}{}\label{night:P4.41}\label{stmt:P4.41}
\StatementRecord{night:P4.41}{proved}{02XF}{\cref{lem:II:conductor-duality,cdp:st:A.5.17,night:P4R.10,night:P4R.15,prop:II:smooth-branch-contact,lem:II:boundary-deformation-transport}}

At a point of a support in the standing setting of
\cref{def:II:standing-torus-fibration}, suppose the
entire reduced critical curve is the ordinary node
$C=\{z=xy=0\}$ and both branches have generic transverse
type $A_3$. Let $d_i,b_i,q$ be the Hessian zero orders,
projected conormal colengths, and normalization descent
defect of \cref{night:P4R.10}. Then
\begin{equation}
 \delta(C)+\sum_{i=1}^2\left(\frac{d_i}{2}+b_i\right)+q
                         \ge\frac92.                 
\label{eq:II:manual:4.41a}
\end{equation}
Here $R=0$ and $\delta(C)=1$. Equality makes the normalization
of the hypersurface smooth over the point.
\end{proposition}
\begin{proof}

\begin{enumerate}[label=\arabic*.,leftmargin=*]
\item Put $P=(z,xy)$, and let $I$ be the ambient ideal of the
scheme-theoretic conductor image. We have $P^2=P^{(2)}\subset I$,
and $J=P/I$ is torsionfree of rank one on both branches.
After dividing its normalized image by the two branchwise
vanishing orders, it is an $A=\mathcal O_C$ submodule of
$B=\mathbb C\{x\}\oplus\mathbb C\{y\}$ generating both
summands. The conductor of $A$ in $B$ is $\mathfrak m_A$,
so $\mathfrak m_AB\subset J$. Thus $J$ is determined by
a subspace of $\mathbb C^2$ projecting onto both coordinates.
If that subspace is a line, rescaling gives $J\simeq A$ and
$q=1$. If it is the whole space, $J\simeq B$ and $q=0$.
These are module isomorphisms. The divided orders are
the $b_i$ already in the formula.

\item Put $\mathcal P=\mathbb C\{x,y,z\}$, let $f$ be the
reduced local equation of the support, and set
$R=\mathcal P/(f)$ and $\mathfrak c=I/(f)\subset R$.
The multiplier identification in \cref{lem:II:conductor-duality}
gives $\operatorname{Hom}_R(\mathfrak c,R)=\overline R$.
Since $\mathfrak c$ is also an ideal of $\overline R$, its
ring condition is
\begin{equation}
 \operatorname{Hom}_R(\mathfrak c,\mathfrak c)=
 \operatorname{Hom}_R(\mathfrak c,R).
\label{eq:II:manual:4.41b}
\end{equation}
Apply \cite[Lemma~1.9, Corollary~1.11 and Theorem~1.12,
pp.~531--532; Remark~1.13, p.~533]{DeJongVanStraten1990Normalization}
with base $\operatorname{Spec}\mathbb C$, ambient ring $\mathcal P$,
hypersurface equation $f$, ambient ideal $I$, and hypersurface
ideal $\mathfrak c$. The quotient $\mathcal P/I=R/\mathfrak c$
is Cohen--Macaulay of codimension two by
\cref{lem:II:conductor-duality}. A Hilbert--Burch matrix for $I$,
with a coefficient column for $f$ appended, gives a square
presentation matrix $M$ with $\det M=f$, after the corresponding
choice of signs. The cited criterion says that every entry of
$\operatorname{adj}(M)\bmod f$ belongs to $\mathfrak c$.
The quotient map $\mathcal P\to R$ has preimage
$\mathfrak c$ equal to $I$, so this is exactly membership of
every ambient adjugate entry in $I$. We use the matrices below.

\item Suppose first that $q=1$. The surjection
$P/P^2\to J$ is a surjection $A^2\to A$, with primitive
kernel generator $g$. Taking a complementary generator $k$
of $P$ gives $I=(g,k^2)$. A determinant presentation
$f=ag+bk^2$ has two-by-two adjugate containing $a,b,g,k^2$.
Condition \eqref{eq:II:manual:4.41b} says $a,b\in I$, so the exact equation is
\begin{equation}
                   f=Ag^2+Bgk^2+Ck^4.                
\label{eq:II:manual:4.41c}
\end{equation}
We claim that $A(0)$ cannot be a unit. Suppose it is. If $g_z(0)\ne0$, then $dg(0)\ne0$ and the ambient Hessian of \eqref{eq:II:manual:4.41c} is $2A(0)\,dg(0)\otimes dg(0)\ne0$. Apply \cref{night:P4R.15} to this original reduced generator and its original full fibre. Its exact product coordinates make the entire reduced critical curve smooth, contradicting the node. This application includes the original suspension construction: straightening $g$ and completing the square gives the plane residual obtained from $k^4$; the finite-gradient and nonzero-critical-value construction is the one in \cref{cdp:st:A.5.17}. The full-function comparison is supplied by the two exact cases in \cref{night:P4R.15}.

\item Suppose $g_z(0)=0$. Its $xy$ coefficient is a unit, and take $k=z$. The square-root coordinate in the first factor of \eqref{eq:II:manual:4.41c} gives
\[
       G=\sqrt A\,g+\frac{B}{2\sqrt A}z^2,\qquad
       f=G^2+Dz^4,\qquad D=C-\frac{B^2}{4A},
\]
and the conductor is still $I=(G,z^2)$. Here $G(x,y,0)$ is $xy$ times a unit and $dG(0)=0$. We determine the local factors from this conductor before deforming the full function.

Apply \cref{lem:II:quadratic-conductor-normalization} with this
same $R$, $G$, $H=z^2$, $U=D$ and conductor $I=(G,z^2)$.
The sequence $(G,z^2)$ is regular because $G(x,y,0)$ is $xy$
times a unit. Its exact multiplier is $t=G/z^2$, and the lemma gives
\[
 \overline R=R+Rt,\qquad t^2=-D.
\]
For $D(0)=0$ the reduced germ is locally irreducible. For $D(0)\ne0$,
the two original irreducible factors are exactly
$G-z^2\sqrt{-D}$ and $G+z^2\sqrt{-D}$.

\item In the locally irreducible case, the full map is $Vf^M$. Take its exact root $f_0=V^{1/M}f$, writing
\[
 f_0=G_0^2+z^4D_0,\qquad
 G_0=V^{1/(2M)}G,\quad D_0=V^{1/M}D.
\]
It still has $G_0(x,y,0)=xy$ times a unit and $dG_0(0)=0$. The first deformation is $D_0\mapsto D_0+\delta$, with $\delta$ arbitrarily small and chosen so that this is a unit. The coordinate $Z=z(D_0+\delta)^{1/4}$, keeping $x,y$, is invertible; its inverse is the unique small solution for $z$, whose derivative at the origin is nonzero. Parameter Morse coordinates in $x,y$, keeping $Z$, then give $G_0=XY+a(Z)$, where $a\in(Z^2)$. This is the same square-completion construction as in \cref{cdp:st:A.5.17}: the plane Hessian is invertible, and its successive square-root coordinates and their implicit inverses depend holomorphically on $Z$.

A small change $a\mapsto a+\alpha Z^2$ makes $a''(0)\ne0$. Shift $G_0$ by $-\varepsilon$. The resulting root function is
\[
                    (XY+a(Z)-\varepsilon)^2+Z^4.
\]
At the origin its critical value is $\varepsilon^2$, its $XY$ Hessian entry is $-2\varepsilon$, and its $ZZ$ entry is $-2\varepsilon a''(0)$. It is a Morse point. These are the original suspension and moving-node formulas, now applied with the established full-map equality $F=f_0^M$.

\item In the two-factor case, take $Z=zs^{1/2}$ with $s^2=-D$, keeping $x,y$; its implicit inverse is again holomorphic. The reduced function is $G^2-Z^4$, and the full function has its original multiplicities:
\[
                F=U(G-Z^2)^a(G+Z^2)^b,
                \qquad a,b\ge1.
\]
Parameter Morse coordinates in $x,y$, keeping $Z$, give $G=XY+A(Z)$ with $A\in(Z^2)$ and transport the unit to a holomorphic unit still denoted $U$. Retain this unit and both factors. Deform $A$ to $A_\alpha=A+\alpha Z^2$ and choose an arbitrarily small $\alpha$ such that, writing $A_\alpha=cZ^2+O(Z^3)$,
\[
          N=a+b,\qquad T=a(c-1)+b(c+1)=Nc+b-a\ne0.
\]
The full deformation is exactly
\[
 F_{\alpha,\varepsilon}
 =U(X,Y,Z)(XY+A_\alpha(Z)-Z^2-\varepsilon)^a
                (XY+A_\alpha(Z)+Z^2-\varepsilon)^b.
\]
These are holomorphic deformations of the two original analytic factors, with their multiplicities fixed.

For the selected $\alpha$, put $(X,Y,Z)=\varepsilon(r_1,r_2,r_3)$. Each of the two parenthesized factors, divided by $-\varepsilon$, is a holomorphic function of $(\varepsilon,r)$ equal to $1$ at $\varepsilon=0$. Its linear term in $\varepsilon$ is
\[
                 -r_1r_2-(c\mp1)r_3^2.
\]
Consequently the logarithmic gradient
\[
 \mathcal L(\varepsilon,r)
       =\nabla_{X,Y,Z}\log F_{\alpha,\varepsilon}(\varepsilon r)
\]
extends holomorphically to $\varepsilon=0$. If $\ell_j=(\partial_j\log U)(0)$, its limiting value is
\[
            (\ell_1-Nr_2,\ \ell_2-Nr_1,\ \ell_3-2Tr_3).
\]
It has the unique zero
$r_0=(\ell_2/N,\ell_1/N,\ell_3/(2T))$, and its derivative in $r$ is
\[
          \begin{pmatrix}0&-N&0\\-N&0&0\\0&0&-2T\end{pmatrix},
          \qquad \det=2N^2T\ne0.
\]
The implicit-function theorem supplies $r(\varepsilon)$ with $\mathcal L(\varepsilon,r(\varepsilon))=0$ and $r(0)=r_0$. Thus $p_\varepsilon=\varepsilon r(\varepsilon)$ is a critical point of the full function. Its value is $(-\varepsilon)^N U(0)(1+O(\varepsilon))\ne0$. At that critical point the exact Hessian identity is
\[
 \operatorname{Hess}F_{\alpha,\varepsilon}(p_\varepsilon)
   =\frac{F_{\alpha,\varepsilon}(p_\varepsilon)}{\varepsilon}
             D_r\mathcal L(\varepsilon,r(\varepsilon)).
\]
The displayed limiting matrix is invertible, so this is a holomorphic Morse point for sufficiently small nonzero $\varepsilon$.

\item We specify the transport for both constructions. Choose an outer ball transverse to the original zero-divisor strata, small enough that its original critical values are all zero. In the root case, the first family $G_0^2+z^4(D_0+\delta)$ has fixed transverse $A_3$ charts near the boundary. After the selected $\delta$ and its invertible coordinate changes, choose an inner ball about the origin inside that outer ball. The family $(XY+a(Z)+\alpha Z^2-\varepsilon)^2+Z^4$ has the same boundary product charts, along the moving smooth parts of $Z=0,XY=\varepsilon$. The nested transports of \cref{lem:II:boundary-deformation-transport} take its local Morse sphere first to the selected $\delta$ root fibre, then to an original root fibre, and finally into $F^{-1}(b^M)$.

For the two-factor family, near each boundary intersection link use $v=XY+A_\alpha(Z)-\varepsilon$ and $Z$ as the two normal coordinates. A complementary coordinate runs along the moving curve and is transverse to the boundary. The full transverse function is $U(v-Z^2)^a(v+Z^2)^b$. The weighted unit map and inverse from \cref{prop:II:smooth-branch-contact}, with $k=2$, give the fixed function $(v-Z^2)^a(v+Z^2)^b$. At smooth component points, the corresponding charts are powers of a normal coordinate. These charts verify every boundary hypothesis of \cref{lem:II:boundary-deformation-transport} for $g_{\alpha,\varepsilon}=F_{\alpha,\varepsilon}$ and $m=1$.

In each case the local Morse sphere from \cref{cdp:st:A.5.17} has square $-2$ and is carried into an original torus fibre with the same square. Its sphere map lifts to the contractible universal cover of that torus, so its class has square zero. This contradiction proves $A(0)=0$ at the original point.

\item If $g_z(0)$ is a unit, both saturated Hessian lines approach
$dz$, so $b_1=b_2=1$ and $d_i\ge1$. If $g_z(0)=0$,
the differential of $g$ on each branch has exactly the
common factor $t_i$ and a remaining nonzero coefficient
in the other normal direction. Thus $b_i=0,d_i\ge3$.
In either case the left side of \eqref{eq:II:manual:4.41a} is at least five,
including $q=1$.

\item Now let $q=0$. Choose generators $a,b$ of $P$ mapping
to $(1,0),(0,1)$ in $J\simeq B$. Their coefficient matrix
is invertible modulo the maximal ideal, hence invertible
locally. We obtain
\begin{equation}
 I/P^2=(ya,xb),\qquad
 f=xU(x)b^2+yV(y)a^2\pmod {P^3}.                      
\label{eq:II:manual:4.41d}
\end{equation}
The second identity records the Hessian kernel $b$ on the
$x$ branch and $a$ on the $y$ branch. Write the exact
relation $xy=\alpha a+\beta b$. Near a general point of
the $x$ axis, solve this for $y$ using coordinates $x,a,b$.
The null cubic coefficient contains
$V(0)\alpha(x,0)/x$; the remaining contribution from
$P^3$ is holomorphic at the node. Its identical vanishing
in type $A_3$ forces $V(0)\alpha(0)=0$. The other branch
gives $U(0)\beta(0)=0$.

\item If both $a_z(0),b_z(0)$ are nonzero, the relation for $xy$
has $\alpha(0)\beta(0)\ne0$. Hence $U(0)=V(0)=0$,
both $b_i=1$, and both $d_i\ge2$. The defect is at least
five. Otherwise interchange the branches so that $a_z(0)$
is a unit and $b_z(0)=0$. Then $\alpha(0)=0$ and
$\beta(0)$ is a unit, forcing $U(0)=0$. On the $x$ axis,
$db$ has exactly the common factor $x$ and remaining
$dy$ coefficient a unit. On the $y$ axis, $da$ is primitive
with a unit $dz$ coefficient. Thus
\[
 b_x=0,\quad d_x=3+\operatorname{ord}_xU\ge4,
 \qquad
 b_y=1,\quad d_y=1+\operatorname{ord}_yV\ge1.
\]
This proves \eqref{eq:II:manual:4.41a}. Equality requires
$\operatorname{ord}_xU=1$ and $V(0)\ne0$.

\item We identify the normalization in this exact equality case.
Use $a$ as the ambient coordinate $z$. The generator relation
with $\beta$ a unit gives
\[
                    I=(a^2,ya,x^2y-x\alpha a).
\]
Modulo its first two generators, replace $\alpha$ by
$\alpha(x,0,0)=x\ell(x)$. The coordinate
$Y=y-\ell(x)a$ puts the ideal into
$I=(z^2,Yz,x^2Y)$. Renaming $Y$ as $y$, its Hilbert--Burch
matrix and determinant extension are
\[
 \begin{pmatrix}y&0\\-z&x^2\\0&-z\end{pmatrix},
 \qquad
 M=\begin{pmatrix}y&0&a\\-z&x^2&b\\0&-z&c\end{pmatrix},
 \qquad
 \det M=az^2+byz+cx^2y.
\]
The new two-by-two minors are
$yb+za,yc,-zc,-x^2a,-za,x^2c+zb$, up to signs.
Their membership in $I$ gives exactly
\[
             a=yA,\qquad c=zC,\qquad b=-x^2C+zB
                          \pmod I.
\]
Therefore $f\in(yz^2)+I^2$, so
\begin{equation}
           f=Uyz^2+Vz^4+Wx^2y^2z+Tx^4y^2.            
\label{eq:II:manual:4.41e}
\end{equation}
The equality orders $d_y=1,d_x=4$ say precisely
$U(0)T(0)\ne0$.

\item Substitute $y=z^2v$ and divide \eqref{eq:II:manual:4.41e} by $z^4$.
At the finite root its $v$ derivative is $U(0)$.
The implicit-function theorem supplies the smooth
factor $g_1=y-z^2v(x,z)$. Division by it gives
\[
                  g_2=z^2A+x^2yzB+x^4yC,
                  \qquad A(0)C(0)\ne0.
\]
For clarity, this division retains the factors $x^2,x^4$:
subtracting the expression at $y=z^2v$ supplies one factor
$y-z^2v$; every extra substituted $y$ supplies $z^2$ and
is absorbed into the first term. The parameter Morse
centre in $z$ is divisible by $x^2y$. Substitution leaves
$x^4y$ times a unit. A coordinate change therefore
gives $g_2=Z^2+x^4Y$. Its finite birational normalization is
\[
                 (x,u)\longmapsto(x,-u^2,x^2u),
\]
with target coordinates $(x,Y,Z)$. It is smooth; finiteness
follows from $u^2=-Y$ and birationality from $u=Z/x^2$
away from $x=0$. The other branch $g_1$ is smooth already.
Thus the full normalization is smooth over the node.

\end{enumerate}
\end{proof}


\begin{proposition}[Two identifications in the tacnodal ruled case]\hypertarget{tag-02XG}{}\label{night:P4.42}\label{stmt:P4.42}
\StatementRecord{night:P4.42}{proved}{02XG}{\SeriesRef{night:P4.33}, \cref{lem:II:plane-branch-conductor}, \SeriesRef{ii:excess-one-curve-germs}, \SeriesRef{night:P4.41}, \SeriesRef{night:P4R.10}, \SeriesRef{night:P4R.13}, \SeriesRef{night:P4R.14}, \SeriesRef{night:P4R.6}}

In the irreducible-or-normal-crossings torus-fibration setting
of \cref{def:II:standing-torus-fibration}, let $S$ be an
irreducible reduced fibre support whose normalization has
a resolution ruled over a positive-genus curve. Suppose
that its critical curve has generic transverse type $A_3$
on every component. Its point-identification cycle cannot
consist of two two-branch points, including nonordinary
points and additional unibranch singular points of the
critical curve.
\end{proposition}
\begin{proof}

\begin{enumerate}[label=\arabic*.,leftmargin=*]
\item Proposition~\ref{night:P4.26} gives $e(S)\ge0$.
At $e(S)=0$, the uniform $A_3$ case of
\cref{night:P4.32} has a smooth elliptic conductor image.
Its normalization is the identity and has no distinct
branch points to identify. Thus the asserted two-point
configuration has $e(S)>0$. Apply
\cref{lem:II:positive-ruled-conductor-graph} to this
$f:Y\to B$ and $mS$. The conductor then has two rational
normalized components and one graph cycle. In the stated
configuration that cycle consists of two two-branch
points. We compute their local Hessian defects.

\item At a planar curve point, write its ideal as $(z,h)$.
Its symbolic square equals its ordinary square, so the
rank-one Hessian belongs to $\operatorname{Sym}^2 E_0$,
where $E_0$ is the normalized conormal image. On each normalized branch apply
\cref{lem:II:rank-one-conormal-lattice} with $E=E_0$, the same
saturated Hessian line $L$, exponent $d$ and quotient colength
$\beta=b_{\rm proj}$. Its symmetric-square bound is
$d/2+b_{\rm proj}\ge a+b$ for the elementary-divisor orders $(a,b)$.
On normalized branch $i$, write $dh=t^{v_i}w_i$ with
$w_i$ a primitive plane kernel generator. The conormal
image is generated by $dz,dh$, so $a_i+b_i=v_i$.
Apply \cref{lem:II:plane-branch-conductor} to the same
plane equation $h$ and its normalization. Its identity
$\omega_{\widetilde C}=I\omega_C$, evaluated on the
residue $dx/h_y$ of order $r_i-v_i$, gives
$c_i=v_i-r_i$. Since $\sum_i c_i=2\delta$, this yields
$\sum_i(a_i+b_i)=2\delta+\deg R$. The whole local defect
is consequently at least $3\delta+q$.
A nonordinary planar two-branch point has $\delta\ge2$
and contributes at least six.

\item For a nonplanar point, the ambient Hessian vanishes: a
nonzero rank-one Hessian would give $f=X^2+g(Y,Z)$ by
the parameter Morse lemma and would put the entire
critical curve in the smooth plane $X=0$.
On a branch of multiplicity $m$, all conormal generators
have order at least $m$, and every Hessian entry has
order at least $m$. Thus $b_{\rm proj}\ge m$, $d\ge m$,
and, since $r=m-1$,
\[
                  d/2+b_{\rm proj}-r\ge1+m/2\ge1.
\]
When $\delta\ge3$, a nonplanar two-branch point
therefore contributes at least five.
The only remaining possibility has $\delta=2$.
The exact excess-one classification in \cref{ii:excess-one-curve-germs} identifies
it with a cusp and transverse line,
\[
                  P=(xz,yz,h),\qquad h=y^2-x^3.
\]
Here $P^{(2)}/P^2$ is generated by $zh$ and is
annihilated by the maximal ideal. Write $f=Azh+P^2$.
Along the cusp, the normal quadratic determinant is
$4KG-A^2$, with $K(0)=0$. Its rank is one generically,
so $A(0)=0$ and in fact $f\in P^2$.
The conormal lattices on the two normalized branches
are respectively $\operatorname{diag}(t^2,t^3)$ and
$\operatorname{diag}(t,t)$. Their total colength is seven,
while $\deg R=1$ and $\delta=2$.
Applying the same symmetric-square inequality gives
defect at least $2+7-1=8$.
Every two-branch point thus contributes at least $9/2$,
and equality is possible only at the ordinary node of
\cref{night:P4.41}, whose normalization is smooth.

\item Apply the rank-one Hessian calculation of
\cref{night:P4R.10} to the two normalization maps
$j_i:\mathbb P^1\to C\hookrightarrow Y$ and the nilradical
of the full conductor scheme, pulled back to each normalized
component. Here
$K_i=\ker(j_i^*\Omega_Y\to\Omega_{\mathbb P^1}(-R_i))$
is the normalized derivative kernel. Since
$c_1(\Omega_Y)\cdot C_i=0$, it has degree
$2+\deg R_i$.
The quotient $K_i/L_i$ has Euler characteristic
$3+\deg R_i-\deg L_i$. Combining both lines with the
normalization sequence of $C$ gives the exact formula
\begin{equation}
 \chi(\mathcal O_{C_{\rm sch}})
   =8-\sum_p\left[
       \delta_p+\sum_{i\text{ over }p}
          \left(\frac{d_i}{2}+b_i-r_i\right)+q_p
                 \right].                            
\label{eq:II:refined-4-42a}\end{equation}
Every extra unibranch singular point has a strictly positive
bracket by the local conormal inequalities of \cref{night:P4R.10}.
At a smooth point the bracket is nonnegative. The two
two-branch points contribute at least nine together by the
local calculation just given.

\item If the horizontal section is not contracted, \cref{night:P4.33} gives
$\chi(\mathcal O_{C_{\rm sch}})=g-1\ge0$, contradicting
\eqref{eq:II:refined-4-42a}, whose right side is at most $-1$.
If the section is contracted, that Euler characteristic is
$-1$. Equality must then hold at both two-branch points and every
other bracket must vanish. There are no additional
unibranch singular points. Both cycle points are ordinary
nodes and the normalization is smooth over them by \cref{night:P4.41}.
At every smooth point of the
critical curve, \cref{night:P4R.14} gives either the product tacnode
or the displayed $A_1$ normalization degeneration. The
latter has Hessian zero order one and contributes $1/2$,
so it is excluded by the equality in \eqref{eq:II:refined-4-42a}. Thus the
normalization is smooth along all of the conductor, and
it is smooth away from the conductor by \cref{night:P4R.6}. A contracted
positive-genus horizontal section in its minimal resolution
is impossible. This proves the assertion.

\end{enumerate}
\end{proof}


\Needspace{12\baselineskip}
\section{Tacnodal conductor sectors}\label{sec:II:tacnodal-conductor-sectors}

\subsection{Further uniform transverse sectors}

\begin{proposition}[The contracted section and three-branch critical point]\hypertarget{tag-02XI}{}\label{night:P4.44}\label{stmt:P4.44}
\StatementRecord{night:P4.44}{proved}{02XI}{\SeriesRef{night:P4.33}, \SeriesRef{night:P4.42}, \SeriesRef{night:P4R.14}, \SeriesRef{night:P4R.6}}

In the setting of \cref{lem:II:positive-ruled-conductor-graph},
suppose that the generic transverse type along every image-conductor
component is $A_3$. Then the horizontal section $H$ of the pole
divisor is contracted by the minimal normalization resolution
$\sigma:T\to N$. The point-identification cycle of the reduced
image conductor $C$ occurs at a single three-branch point.
\end{proposition}
\begin{proof}

\begin{enumerate}[label=\arabic*.,leftmargin=*]
\item Apply \cref{lem:II:positive-ruled-conductor-graph} to this
same full fibre $mS$ and the maps $T\xrightarrow{\sigma}N\xrightarrow{\nu}S$.
It gives two rational normalized conductor components and
$\sum_{p\in C}(r_p-1)=2$. The one cycle therefore occurs either
at two two-branch points or at one three-branch point.
Proposition~\ref{night:P4.42}, with this same conductor and
its uniform transverse type $A_3$, excludes the first case.
Thus the cycle occurs at one three-branch point $p$. Suppose that the horizontal section $H$ is
not contracted. It maps finitely to one normalized
conductor component $C_H^\nu=\mathbb P^1$.
The generic tacnode has two conductor preimages.
The degree of $H\to\mathbb P^1$ is greater than one,
since $H$ has genus $g\ge1$, and hence is exactly two.
No other, vertical conductor component can map onto
$C_H$: the horizontal map already supplies both generic
preimages.

\item At a point where the entire critical curve is smooth,
\cref{night:P4R.14} gives either the product tacnode or the germ
$X(Y^2+XZ)$. Each has two distinct normalization points
over the critical-curve point. Each point lies on the
strict conductor, by \cref{night:P4R.6}. Thus the two preimages must
both lie on $H$, and the degree-two map is unramified
there. All ramification is over singular points of $C$.

\item By \eqref{eq:II:refined-4-42a} and \cref{night:P4.33} the sum of the local defects is
$9-g\le8$. A nonplanar three-branch point has
$\delta\ge2$, and the branchwise estimate in \cref{night:P4.42}
adds at least one per branch; its defect is at least
five. A planar three-branch point has $\delta\ge3$
and contributes at least $3\delta\ge9$, which is
already impossible. Every additional unibranch singular
point contributes at least two: in the planar case
the bound is $3\delta$, and in the nonplanar case it is
at least $\delta+1$. There is consequently at most one
additional unibranch singular point.

\item The normalization $C_H^\nu$ has at most two points
above $p$, since the other global conductor component
supplies at least one of its three branches. Including
the possible additional unibranch point leaves at most
three possible branch values for $H\to\mathbb P^1$.
Every ramification point of a degree-two map has index
two, and its branch value has only that ramified
preimage. Riemann--Hurwitz requires $2g+2\ge4$ distinct
branch values. This contradiction proves that $H$ is
contracted.

\end{enumerate}
\end{proof}

\begin{lemma}[Cubic normalization and factor lifts]\hypertarget{tag-P280}{}\label{lem:II:ordinary-triple-factor-lifts}
\StatementRecord{lem:II:ordinary-triple-factor-lifts}{proved}{P280}{}
Let $f\in\C\{u,x,y\}$ be reduced, with $f(0)=0$ and reduced
critical locus $C=\{u=h(x,y)=0\}$, where $h$ is a nonzero
reduced germ with $h(0)=0$. Suppose the generic
transverse germ along every branch of $C$ is an ordinary triple
point. There is a coordinate $w=u+\lambda h$, with inverse
$u=w-\lambda h$, and a factorization
\begin{equation}
 f=\varepsilon\prod_{j=1}^s\phi_j,\qquad
 \phi_j(w,x,y)=\Phi_j(w,x,y,h(x,y)),\qquad
 \sum_{j=1}^s n_j=3,
 \label{eq:II:triple-factor-lifts}
\end{equation}
where $n_j\ge1$, $\varepsilon$ is a unit, the $\phi_j$ are precisely the
distinct irreducible factors, and
\[
 \Phi_j(w,x,y,H)=\sum_{k=0}^{n_j}a_{jk}(w,x,y)w^{n_j-k}H^k
\]
has holomorphic coefficients. Each $\phi_j$ has order $n_j$ at the
generic point of every branch of $C$. Consequently every full
local function with this reduced support has the exact expression
\begin{equation}
 G=U\prod_{j=1}^s\Phi_j(w,x,y,h)^{m_j},\qquad m_j\ge1,
 \label{eq:II:triple-full-factor-lifts}
\end{equation}
for a holomorphic unit $U$ and its component multiplicities $m_j$.
\end{lemma}
\begin{proof}
Put $P=(u,h)$ in the regular local ring $\C\{u,x,y\}$.
The quotients $P^k/P^{k+1}$ are free over $\C\{u,x,y\}/P$,
with basis the degree-$k$ monomials in $u,h$. Their finite
filtration makes $\C\{u,x,y\}/P^3$ Cohen--Macaulay of dimension
one. Thus $P^3$ has exactly the minimal primes of $P$ as
associated primes. Generic transverse multiplicity three gives
$f\in P^3$ after localization at each such prime, and hence
$f\in P^3$ itself.

Choose $w=u+\lambda h$ to be a nonzerodivisor modulo $f$.
Choose it also so that the coefficient $a$ in
\begin{equation}
 f=a h^3+b h^2w+c hw^2+d w^3
 \label{eq:II:triple-cubic-expression}
\end{equation}
is nonzero at the generic point of every branch of $C$.
Both choices exclude finitely many constants $\lambda$: the
first avoids the finitely many factors of $f$, and on each
critical-curve branch the second evaluates a nonzero binary
cubic at $(-\lambda,1)$. If the binary cubic at the origin is
nonzero, choose $\lambda$ so that $a(0)\ne0$ as well. Set
$S=\C\{w,x,y\}$ and $R=S/(f)$.

In the total quotient ring of $R$, put $t=h/w$ and
\begin{equation}
 \omega=at,\qquad \nu=at^2+bt.
 \label{eq:II:triple-normalization-generators}
\end{equation}
The relation $at^3+bt^2+ct+d=0$ gives the multiplication table
\begin{equation}
 \omega^2=-b\omega+a\nu,\qquad
 \omega\nu=-ad-c\omega,\qquad
 \nu^2=-bd-d\omega-c\nu.
 \label{eq:II:triple-normalization-products}
\end{equation}
Thus $J=R+R\omega+R\nu$ is a finite birational $R$-algebra.
The relations among these generators include the columns of
\begin{equation}
 M=\begin{pmatrix}
      -ah&-bh&ch+dw\\
      w&-h&0\\
      0&w&h
    \end{pmatrix},\qquad \det M=f.
 \label{eq:II:triple-normalization-matrix}
\end{equation}
The cokernel of $M:S^3\to S^3$ has depth two: the determinant
makes this map injective over the regular ring $S$, and its
cokernel has a free resolution of length one. Its associated
primes are the height-one primes dividing $f$. Inverting $w$
solves its relations as \eqref{eq:II:triple-normalization-generators}
and identifies it with $R[w^{-1}]$. Hence the surjection from
this cokernel to $J$ has a $w$-torsion kernel, which is zero.
In particular $J$ has depth two.

The multiplication table makes $J$ integral over $R$, so it
embeds in the normalization $\overline R$. Outside $C$ the
surface is smooth and $J=R$. At a smooth point of $C$ where
the transverse cubic has three distinct roots and $a$ is a
unit, use coordinates $(w,h,z)$. There $J=R[t]$ is given by
\[
 wt=h,\qquad t^3+(b/a)t^2+(c/a)t+d/a=0.
\]
Substituting $h=wt$, the three simple roots give three
holomorphic solutions $t=t_i(w,z)$ by the implicit function
theorem. These are the three smooth normalization branches,
with maps $(w,z)\mapsto(w,wt_i(w,z),z)$. The generic transverse
hypothesis and the choice of $a$ make all exceptional points
on $C$ isolated; after shrinking, only the origin remains.
Thus $\overline R/J$ has finite length. The normal surface
module $\overline R$ also has depth two. The local cohomology
sequence of $0\to J\to\overline R\to\overline R/J\to0$,
using $H^0_{\mathfrak m}(\overline R)=H^1_{\mathfrak m}(J)=0$,
gives $\overline R=J$.

If $a,b,c,d$ all vanish at the origin, the matrix $M$ has
zero reduction modulo the maximal ideal, so
$\overline R/\mathfrak m\overline R$ has basis $1,\omega,\nu$.
By \eqref{eq:II:triple-normalization-products} its span of
$\omega,\nu$ is a square-zero ideal and the fibre algebra is
local. Apply the local-fibre assertion of
\cref{lem:II:quadratic-conductor-normalization} to this same
$R$ and $B=\overline R=J$. It makes $R$ a domain, so $f$ is irreducible.
In this case take $s=1$, $n_1=3$
and
\[
 \Phi_1=aH^3+bH^2w+cHw^2+dw^3.
\]
Its order along each branch is three by the transverse
hypothesis.

Otherwise $a$ is a unit and $\overline R=R[t]$. Factor the
monic polynomial
\[
 Q(T)=T^3+(b/a)T^2+(c/a)T+d/a
\]
as $Q=\prod_jQ_j$, with
$Q_j(T)\bmod\mathfrak m=(T-r_j)^{n_j}$ and the $r_j$ distinct.
These factors have holomorphic coefficients. Indeed, the
derivative of the coefficient multiplication map is
$(\dot Q_j)_j\mapsto\sum_j\dot Q_j\prod_{k\ne j}Q_k$;
reduction modulo each coprime factor shows its kernel is
zero, and its source and target both have dimension three.
The analytic implicit function theorem therefore gives the
stated factors and their uniqueness.

The corresponding relatively prime factors decompose
$\overline R$ into algebras
\[
 B_j=S[T]/(wT-h,Q_j(T)).
\]
Here $R[t]=S[T]/(wT-h,Q(T))$: the primitive linear polynomial
$wT-h$ is prime, and substituting $T=h/w$ in $Q$ gives
$f/(aw^3)\ne0$ in the fraction field of $S$. Thus the second
ring is a complete intersection of dimension two. Its reduction
modulo $w$ has dimension one, and thus $w$ is a nonzerodivisor;
inverting $w$ gives the same subalgebra of the total quotient
ring. Each $B_j$ has a local residue algebra
$\C[T]/((T-r_j)^{n_j})$ and is a local factor of the normal
ring $\overline R$. It is consequently a domain. Define
\[
 \phi_j=w^{n_j}Q_j(h/w),\qquad
 \Phi_j(w,x,y,H)=w^{n_j}Q_j(H/w).
\]
The formulas are polynomials in $h$ and $H$, respectively,
with coefficients in $S$, and $f=a\prod_j\phi_j$. Sending
$T$ to $h/w$ gives the isomorphism
\[
 B_j[w^{-1}]\simeq S[w^{-1}]/(\phi_j).
\]
It is a domain. Since $w$ is a nonzerodivisor modulo $f$,
it divides none of the $\phi_j$. Factoriality of $S$ then
makes each $\phi_j$ irreducible in $S$; reducedness of $f$
makes these factors distinct.
Their monic $H^{n_j}$ terms give order $n_j$ at every
minimal prime of $(w,h)$. Finally, replacing any given
irreducible equations by their associated $\phi_j$ multiplies
their full product by the corresponding units raised to
$m_j$. Incorporating those units in $U$ gives
\eqref{eq:II:triple-full-factor-lifts} exactly.
\end{proof}

\begin{proposition}[Planar critical curves of ordinary triple type]\hypertarget{tag-02XJ}{}\label{night:P4.45}\label{stmt:P4.45}
\StatementRecord{night:P4.45}{proved}{02XJ}{\cref{lem:II:ordinary-triple-factor-lifts,lem:II:boundary-deformation-transport,cdp:st:A.5.17}}

In a torus degeneration, a point of a planar
singular critical curve cannot have generic transverse
type $D_4$ along all its local branches.
\end{proposition}
\begin{proof}
\begin{enumerate}[label=\arabic*.,leftmargin=*]
\item Choose a local base coordinate and denote its full
fibre function by $G$. The entire reduced critical curve
of a reduced defining function is $C=\{u=h(x,y)=0\}$,
where $h$ is reduced and singular. Apply
\cref{lem:II:ordinary-triple-factor-lifts} to that reduced
function. In its coordinates the full function is
\eqref{eq:II:triple-full-factor-lifts}. Set
$N=\sum_jm_jn_j$. Introduce an independent variable $H$ and
small constant coefficient parameters $\tau=(\tau_{jk})$.
Write $a_{jk,\tau}=a_{jk}+\tau_{jk}$ and let $\Phi_{j,\tau}$
have these coefficients in its displayed degree-$n_j$
monomials, so $\Phi_{j,0}=\Phi_j$. This gives the full family
\begin{equation}
 \mathcal G_\tau(w,x,y,H)=U(w,x,y)
                    \prod_j\Phi_{j,\tau}(w,x,y,H)^{m_j}.
 \label{eq:II:triple-weighted-lift}
\end{equation}
Then $G=\mathcal G_0(w,x,y,h)$, with the same unit and every
component multiplicity.

Choose a sufficiently small closed ball $K$ whose boundary is
transverse to the smooth support strata and to the disjoint
smooth links of $C$. Shrink it so that all critical values
of $G$ are zero: each component of its critical germ has
constant value and contains the origin. The coefficient
perturbations preserve the displayed degree-$n_j$ monomials
and all the powers $m_j$. Choose a small parameter
$\tau=\epsilon$ so that the total reduced binary cubic at the
origin has three distinct finite roots in $v=w/H$, and so
that
\begin{equation}
 P(v)=U(0)\prod_j p_j(v)^{m_j},\qquad
 p_j(v)=\sum_{k=0}^{n_j}a_{jk,\epsilon}(0)v^{n_j-k}
 \label{eq:II:triple-weighted-polynomial}
\end{equation}
has a simple critical point $v_0$ with $P(v_0)\ne0$.
These choices form a dense open set of the coefficient
space. To see the last condition, write the three roots as
$a_i$ and their positive weights as $\mu_i$. Away from the
roots, $P'/P=\sum_i\mu_i/(v-a_i)$ has quadratic numerator.
For three distinct real $a_i$, this sum has one simple
zero in each intervening interval, since its derivative
is $-\sum_i\mu_i/(v-a_i)^2<0$. Thus its discriminant is a
nonzero polynomial in the root parameters; its complement
meets every small open set of distinct root configurations.
The monic factor blocks in the preceding lemma allow those
roots to move independently within their blocks. In its
irreducible zero-cubic case all four coefficients can be
perturbed. In both cases the choices are arbitrarily small.

\item We check boundary transport for the full product.
Along a critical-curve link, the three reduced branches
are smooth and pairwise transverse in a transverse plane.
Take equations $a,b$ for two of them and a complementary
coordinate $z$ along the link. The third has equation
$a\alpha+b\beta$, with $\alpha,\beta$ units. The map
$(a,b,z)\mapsto(X,Y,z)=(a\alpha,b\beta,z)$ has nonzero
Jacobian on the link; its inverse is the unique nearby
solution of these coordinate equations. It identifies the
three branches with $X=0$, $Y=0$ and $X+Y=0$. Thus the full
function is
\[
 W(X,Y,z)X^{\mu_1}Y^{\mu_2}(X+Y)^{\mu_3},
 \qquad \mu_1+\mu_2+\mu_3=N.
\]
Multiplying both $X,Y$ by the same root $W^{1/N}$ removes
the unit. The inverse solves
\[
 (X,Y)=(\zeta^{-1}X',\zeta^{-1}Y'),\qquad
 \zeta^N=W(\zeta^{-1}X',\zeta^{-1}Y',z);
\]
its implicit derivative on $X'=Y'=0$ is
$N\zeta^{N-1}\ne0$. These maps depend holomorphically on the
deformation parameters and keep the $z$ direction along
the moving curve. Their product directions are transverse
to the fixed boundary. At smooth support points the full
function is a power of a normal coordinate after taking
the corresponding unit root. The homogeneous Euler identity
shows that the transverse polynomials have only the critical
value zero. These charts verify
\cref{lem:II:boundary-deformation-transport} for the full
function and the coefficient deformations just chosen.

After fixing $\epsilon$, choose an inner ball $K'\Subset
\operatorname{int}K$ within the coordinate neighbourhood
constructed below. The same charts apply there when $h$
is replaced by a sufficiently small plane deformation
$H_s$ whose zero set is smooth and transverse on its
boundary links: the common curve is $w=H_s=0$ and the
three transverse factors remain distinct. Each application
of the transport lemma has exponent $m=1$, since its
function is the complete weighted product. Its nested
assertion permits these inner and outer balls.

\item Keep $\tau=\epsilon$ fixed and put
\begin{equation}
 \Psi(v,x,y,H)=U(Hv,x,y)
       \prod_j\bigl(H^{-n_j}\Phi_{j,\epsilon}(Hv,x,y,H)\bigr)^{m_j}.
 \label{eq:II:triple-critical-centre}
\end{equation}
The displayed form of $\Phi_{j,\epsilon}$ makes the quotients
holomorphic, including at $H=0$, and $\Psi(v,0,0,0)=P(v)$.
The implicit function theorem
applied to $\Psi_v=0$ gives $v=V(x,y,H)$ near $v_0$, with
$\Psi_{vv}(V,x,y,H)$ a unit. Put
\begin{equation}
 \beta(x,y,H)=\Psi(V(x,y,H),x,y,H),\qquad
 \Lambda(x,y,H)=H\beta(x,y,H)^{1/N}.
 \label{eq:II:triple-critical-value-coordinate}
\end{equation}
Here $\beta$ is a unit. The chosen root gives
$\Lambda_H(0)=\beta(0)^{1/N}\ne0$, so $q=\Lambda(x,y,H)$
has a holomorphic inverse $H=R(x,y,q)=qR_1(x,y,q)$, with
$R_1$ a unit. Both identities $R(x,y,\Lambda(x,y,H))=H$
and $\Lambda(x,y,R(x,y,q))=q$ hold. Choose the inner ball
$K'$ and its transverse boundary inside this neighbourhood.
The function $\mathcal G_\epsilon(w,x,y,h)$ vanishes at the
origin, so the critical-component argument of the first
step makes all its critical values in $K'$ zero after
shrinking. Choose its transport target disc inside the
outer one.

Set $\widehat h=\Lambda(x,y,h(x,y))$. It is a unit times
the same reduced singular plane function $h$. Its gradient
has an isolated zero. The finite-gradient construction in
the proof of \cref{cdp:st:A.5.17}, applied to this plane
function, gives arbitrarily small linear and constant
perturbations $\widehat h_s$ with a Morse critical point
$p$, nonzero value $q_0=\widehat h_s(p)$ and invertible
Hessian $Q$. Define
\begin{equation}
 H_s=R(x,y,\widehat h_s),\qquad
 G_s=\mathcal G_\epsilon(w,x,y,H_s(x,y)).
 \label{eq:II:triple-full-deformation}
\end{equation}
At $s=0$ this is $H_0=h$. Since $H_s=\widehat h_s
R_1(x,y,\widehat h_s)$, the boundary zero sets and their
transversality satisfy the preceding step.

\item The $w$-critical centre is
$A_s(x,y)=H_sV(x,y,H_s)$, and the full critical value is
the exact identity
\begin{equation}
 (G_s)_w(A_s,x,y)=0,\qquad
 G_s(A_s,x,y)=\widehat h_s^{\,N}.
 \label{eq:II:triple-full-critical-value}
\end{equation}
At $(A_s(p),p)$ the full gradient vanishes and its value
is $q_0^N\ne0$. Moreover
\[
 (G_s)_{ww}(A_s(p),p)
   =H_s(p)^{N-2}\Psi_{vv}(V(p,H_s(p)),p,H_s(p))\ne0.
\]
Differentiating the first identity in
\eqref{eq:II:triple-full-critical-value} identifies the
Schur complement of this Hessian entry with the Hessian
of the second identity. At $p$ it is $Nq_0^{N-1}Q$.
Consequently
\begin{equation}
 \det\operatorname{Hess}G_s
   =(G_s)_{ww}N^2q_0^{2N-2}\det Q\ne0.
 \label{eq:II:triple-full-Hessian}
\end{equation}
The local Morse sphere of \cref{cdp:st:A.5.17} has square
$-2$. Choose its ball and regular value inside $K'$.
Apply
\cref{lem:II:boundary-deformation-transport} first to
\eqref{eq:II:triple-full-deformation} in $K'$. Include the
resulting compact class in a regular $\mathcal G_\epsilon$
level in $K$, choosing the value regular for $K$ as in
that lemma. Apply its target transport in $K$ and the
inverse parameter transport for the coefficient family
\eqref{eq:II:triple-weighted-lift}. These maps give an
orientation-preserving embedding of this compact sphere
into a smooth fibre of the initial full function $G$,
preserving its self-intersection. That fibre is contained
in the original complex two-torus. The sphere lifts to
its contractible universal cover and has square zero
there, contrary to its transported square $-2$.
\end{enumerate}
\end{proof}


\begin{proposition}[Euler bound for uniform ordinary-triple ruled supports]\hypertarget{tag-02XK}{}\label{night:P4.46}\label{stmt:P4.46}
\StatementRecord{night:P4.46}{proved}{02XK}{\SeriesRef{night:P4.31}, \SeriesRef{night:P4.33}, \cref{lem:II:rational-conormal-defects,lem:II:triple-conductor-symbolic-square}, \SeriesRef{night:P4.39}, \SeriesRef{ii:excess-one-curve-germs}, \SeriesRef{night:P4.45}}

In the standing setting of \cref{def:II:standing-torus-fibration},
suppose every reduced fibre support is irreducible or has normal
crossings. A support with positive-genus ruled resolution cannot
have positive Euler number and generic transverse type $D_4$
along every component of its critical curve.
\end{proposition}
\begin{proof}

\begin{enumerate}[label=\arabic*.,leftmargin=*]
\item If the horizontal section is contracted, \cref{night:P4.39} already
gives a contradiction. Suppose it is not contracted.
The remaining critical curve $C=C_1\cup C_2$ has two rational normalized
components and one graph cycle. Let $\nu:N\to S$ be normalization,
let $P_i$ be the reduced ideal of $C_i$, and put $P=P_1\cap P_2$.
Apply \cref{lem:II:triple-conductor-symbolic-square} to this same $\nu$
and to each $C_i$. Its primary conductor ideals are $P_i^{(2)}$.
The conductor quotient is Cohen--Macaulay of dimension one by
\cref{lem:II:conductor-duality}, so it has no isolated primary component.
Its full ideal is therefore $P_1^{(2)}\cap P_2^{(2)}=P^{(2)}$;
the equality follows by imposing order two on exactly the same generic branches.
Apply \cref{lem:II:rational-conormal-defects} to this full $C$, with
$E=E_1\amalg E_2$, $E_i=\mathbb P^1$, $j_i:E_i\to Y$ the normalization
followed by inclusion, and $Q=P/P^{(2)}$. The standing first-Chern-class
vanishing gives $\deg j_i^*\Omega_Y=0$ on each component.
Let $\tau_p$ be the full colength of $Q$ in
$(j_1)_*K_1\oplus(j_2)_*K_2$. It includes the branchwise lattice
and descent colengths. The lemma's formula, with $\gamma=2$, gives
\[
 \chi(\mathcal O_{C_{\rm sch}})
       =10-\sum_p(\delta_p+\tau_p-r_p).
\]
\Cref{night:P4.33} identifies the left side with $g-1$, so the total
defect is $11-g\le10$.

\item We prove that a three-branch point contributes at least
eleven. A planar point contributes $5\delta\ge15$ by
\cref{lem:II:rational-conormal-defects}, applied to the same local plane equation and its normalized derivative lattice. Consider a
nonplanar point. If its three branches are smooth, the
curve has multiplicity three and embedding dimension
three. A generic linear parameter has an Artinian
section of length three and embedding dimension two,
with ideal $(v,w)^2$. The defining curve ideal therefore
has exactly three minimal generators. To justify the
lifting, the parameter is a nonzerodivisor on the
Cohen--Macaulay curve ring, so the Tor term vanishes in
the ideal sequence modulo that parameter; Nakayama
then lifts precisely the three generators.
Each branchwise conormal lattice has rank two and
colength at least two, since every defining gradient
has positive order. Their sum has colength at least
six. Modulo the three branch maximal ideals its
normalization has six generator values. The image of
$Q$, generated by three elements over the local curve
ring, supplies a vector space of dimension at most
three. The additional descent colength is therefore
at least three. Since $\delta\ge2$ and $R=0$, the
full defect is at least $2+6+3=11$.

\item If one of the three branches is singular, let $m\ge4$
be total multiplicity. The branchwise gradient bound
gives $\tau\ge2m$, while $\deg R=m-3$.
The general inequality $\delta\ge m-1$ leaves only
$m=4,\delta=3$ as a possible value below the required
bound. That value has excess $\delta-(3-1)=1$.
The exact classification in \cref{ii:excess-one-curve-germs} shows that its only
singular-branch model is a cusp and two independent
lines, of embedding dimension four. It does not embed
in the ambient smooth threefold. Thus $\delta\ge4$
in the present case, and the full defect is at least
$\delta+2m-(m-3)=\delta+m+3\ge11$.
A cycle at one three-branch point is impossible.

\item The other possible cycle has two two-branch points.
Apply the local estimates of \cref{lem:II:rational-conormal-defects} to the same ideal $P$ at each of these points. Each contributes at least five; equality is
possible only for a planar ordinary node. A nonplanar
two-branch point contributes at least seven, and a
nonordinary planar point at least ten. The total bound
of ten therefore forces two ordinary planar nodes
and $g=1$. Each such node is forbidden by \cref{night:P4.45}.
This proves the assertion.

\end{enumerate}
\end{proof}

\subsection{The full quadratic conductor at three axes}

\begin{proposition}[Rational normalization for a nondegenerate quartic square]\hypertarget{tag-02XL}{}\label{night:P4.47}\label{stmt:P4.47}
\StatementRecord{night:P4.47}{proved}{02XL}{\cref{night:P4.41,lem:II:quadratic-conductor-normalization,lem:II:boundary-deformation-transport}}

Consider a torus degeneration at a point where the
entire reduced critical curve is the three coordinate axes,
with generic transverse type $A_3$ on each axis. Suppose
the local Hessian--conductor defect is nine, the three
Hessian orders are two, and the quartic equation is the
square of $a xy+b xz+c yz$ with $abc\ne0$.
Then the normalization has only rational singularities
over that point. In particular it cannot contain the
contraction of a positive-genus horizontal section.
\end{proposition}
\begin{proof}

\begin{enumerate}[label=\arabic*.,leftmargin=*]
\item Let $f$ be the reduced local defining equation and put $R=\mathbb C\{x,y,z\}/(f)$. Set $P=(xy,xz,yz)$ and $h=xyz$. The symbolic square is
$P^{(2)}=P^2+(h)$, and the conductor image ideal
$I$ contains $P^{(2)}$. Also $f\in P^{(2)}$.
Write $f=Ah+P^2$. On the $x$-axis, a nonzero constant
term of $A$ would give a leading rank-two transverse
Hessian, since the quadratic term is then a nonzero
multiple of $xyz$. Generic type $A_3$ has rank one.
Thus $A(0)=0$. Since $\mathfrak m h\subset P^2$,
we have $f\in P^2$.

\item The normalized conormal lattice is $tK$ on each axis,
so its projected colength is one. The curve has
$\delta=2$ and $R=0$. The stated defect therefore
forces the remaining normalization descent colength
to be one. The normalized quotient values form a
two-dimensional subspace of $\mathbb C^3$ with one
relation involving all three coordinates. Multiplying
a source generator that is a unit on one branch by
that branch's axis coordinate generates its positive
orders alone. Hence all positive-order branch functions
occur in the image, and the leading kernel vector
lifts to an element $g\in I$ with quadratic
part $a xy+b xz+c yz$.

\item We prove the exact ideal equality $I=(g,h)$.
On each coordinate plane, $g$ is the product of its
two coordinate functions times a unit. Thus
$I_0=(g,h)$ has radical $P$. Furthermore $P^2\subset I_0$:
modulo $h$, multiplication of $g$ by $xy$ gives a unit
times $(xy)^2$, and the other two squares follow in
the same way; all mixed products are multiples of $h$.
At the generic point of every axis, $I_0$ is the
tacnode conductor ideal, generated by a normal coordinate
and the square of the null coordinate. It agrees
generically with $I$. The ideal $I_0$ is a
codimension-two complete intersection, so
$\mathcal O/I_0$ is Cohen--Macaulay and $P/I_0$ is
torsionfree on the reduced curve. The map
$P/I_0\to P/I$ has finite-length kernel, which must
vanish. This proves $I=I_0$.

\item Apply the conductor ring condition \eqref{eq:II:manual:4.41b}
to $\mathfrak c=I/(f)\subset R$, with the ambient ideal $I=(g,h)$
just proved. Write $f=\lambda g+\mu h$. The determinant and adjugate are
\[
 M=\begin{pmatrix}h&\lambda\\-g&\mu\end{pmatrix},\qquad
 \det M=f,\qquad
 \operatorname{adj}(M)=\begin{pmatrix}\mu&-\lambda\\g&h\end{pmatrix}.
\]
Theorem~1.12, p.~532 of \cite{DeJongVanStraten1990Normalization},
with the same regular ambient ring and Cohen--Macaulay quotient,
therefore gives $\lambda,\mu\in I$. Thus $f\in I^2$.
Write $f=A g^2+Bgh+Ch^2$.
Its quartic part makes $A(0)\ne0$. Completing the square
and absorbing $\sqrt A$ into the first factor gives
the exact equation
\begin{equation}
                         f=G^2+h^2U.                 \label{eq:II:axis-quartic-square}
\end{equation}
The quadratic part of $G$ is nondegenerate, and
$(G,h)=(g,h)$. If $U(0)\ne0$, the two factors
$G\pm h\sqrt{-U}$ have nondegenerate quadratic
Hessians. Each is an isolated $A_1$ surface singularity,
so the full normalization has rational singularities.

\item Suppose $U(0)=0$. Apply
\cref{lem:II:quadratic-conductor-normalization} with the same
$R$, $G$, $H=h=xyz$, $U$ and conductor $I=(G,h)$.
The sequence is regular: on each coordinate plane, $G$ restricts
to the product of its two coordinates times a unit. The normalization
is therefore $R+R(G/h)$, with $(G/h)^2=-U$, and its fibre algebra
is local. The original reduced germ $f$ is irreducible.

Write the original full fibre function as $F=Vf^M$ and take
\[
 f_0=V^{1/M}f=G_0^2+h^2U_0,\qquad
 G_0=V^{1/(2M)}G,\qquad U_0=V^{1/M}U.
\]
Then $F=f_0^M$ exactly. The conductor ideal remains $(G_0,h)$,
the quadratic part of $G_0$ is nondegenerate, and $U_0(0)=0$.
We use $f,G,U$ for these root data in the following deformation,
which keeps the three axes fixed.

Add an arbitrarily small generic
linear form to $U$, and then replace $U$ by
$U-\epsilon$. On a fixed boundary the transverse
$A_3$ coefficient is a unit, so these changes retain
the boundary transverse type. Write the quadratic part as $a xy+b xz+c yz$.
The diagonal substitution $(x,y,z)=(\alpha X,\beta Y,\gamma Z)$,
where $\alpha^2=c/(ab)$, $\beta=1/(a\alpha)$ and
$\gamma=1/(b\alpha)$, is invertible and gives $G_2=XY+XZ+YZ$.
It multiplies $h$ by the nonzero constant $\alpha\beta\gamma$.
Rename these coordinates $x,y,z$.
Let $L$ be the resulting linear part of $U$.
Under $(x,y,z)=\epsilon(X,Y,Z)$, equation \eqref{eq:II:axis-quartic-square} becomes
\[
 \frac{f_\epsilon}{\epsilon^4}
   =G_\epsilon(X,Y,Z)^2+\epsilon^3(XYZ)^2
                          \bigl(L(X,Y,Z)-1+O(\epsilon)\bigr),
 \qquad G_\epsilon=G_2+O(\epsilon).
\]
Near a nonzero point of the smooth cone $G_2=0$,
use $G_\epsilon$ as the normal coordinate. The
parameter Morse centre is $O(\epsilon^3)$.
Substitution leaves the tangential function
$\epsilon^3(\Phi+O(\epsilon))$, where
\[
                 \Phi=(XYZ)^2(L-1)|_{G_2=0}.
\]

\item We verify a nonempty open set of choices of $L$ for
which $\Phi$ has a nondegenerate critical point of
nonzero value. Parametrize the projective conic by
\[
              [X:Y:Z]=[t(t+1):t+1:-t].
\]
For $L=X+Y+Z$, set
$H(t)=-t^2(t+1)^2$ and $\ell(t)=t^2+t+1$.
In radial coordinate $r$,
$\Phi=r^6H(t)^2(r\ell(t)-1)$.
Radial criticality gives $r=6/(7\ell(t))$, and
\[
 \frac{d}{dt}\left(\frac{H}{\ell^3}\right)
  =\frac{t(t-1)(t+1)(t+2)(2t+1)}{(t^2+t+1)^4}.
\]
The three critical points $t=1,-2,-1/2$ are simple
and avoid zeros and poles. Their radial coordinates
are $2/7,2/7,8/7$. At each point,
\[
       \Phi=-\frac{1024}{823543},\qquad
       \det\operatorname{Hess}_{r,t}\Phi
                     =\frac{4718592}{1977326743}.
\]
These numbers are nonzero. Simplicity and avoidance
of zeros and poles are nonvanishing algebraic conditions
on the coefficients of $L$. The explicit example proves
that they hold on a nonempty Zariski open set, so an
arbitrarily small generic linear perturbation of any
initial linear part is sufficient.

\item After eliminating the normal Morse coordinate, divide
the two tangential critical equations by $\epsilon^3$.
They are analytic in $\epsilon$ and have invertible
Jacobian at the displayed points. The implicit-function
theorem lifts them to nondegenerate three-variable
critical points of $f_\epsilon$. Their critical values
have leading term $\epsilon^7$ times the displayed
nonzero value. They tend to the original point.
On each fixed critical-axis boundary link the nonzero transverse
quadratic coefficient and residual quartic coefficient give parameter
Morse coordinates and a fourth root of a unit, keeping the curve
coordinate. Their implicit inverses identify every sufficiently small
root deformation with $X^2+Y^4$ there. The real curve direction is
transverse to the fixed boundary; the remaining smooth support strata
have submersion charts. Thus \cref{lem:II:boundary-deformation-transport}
applies to this root family with the established equality $F=f_0^M$,
including the smaller balls used for the Morse calculation. It maps
the compact Morse sphere of square $-2$ and its intersection push-off
into the original full fibre $F=t_0^M$. A sphere map into the complex
two-torus lifts to its contractible universal cover and has zero
homology class, contradicting that preserved square. Hence $U(0)=0$
is excluded, and the rational normalization conclusion follows.

\end{enumerate}
\end{proof}

\begin{remark}[Boundary nodes of a shifted square]\hypertarget{tag-P258}{}\label{rem:II:shifted-square-boundary-nodes}
\StatementRecord{rem:II:shifted-square-boundary-nodes}{proved}{P258}{}
Let $h=xyz$ and let $G$ vanish on the coordinate axes. Suppose
$G_y,G_z,U$ are nonzero at a point $(t,0,0)$ with $t\ne0$.
The family $f_\eta=(G-\eta)^2+h^2U$, on the transverse slice
$x=t$, has ordinary nodes at the two points $G=\eta$, $h=0$
for every sufficiently small nonzero $\eta$.

At a generic boundary point $(t,0,0)$ of the $x$-axis,
the numbers $t$, $G_y$, $G_z$, and $U$ are nonzero.
For small $\eta\ne0$, the equations $G=\eta$ and
$h=0$ on this transverse slice have two distinct solutions:
one has $y=0$, and the other has $z=0$. The implicit-function
theorem gives each solution because $G_z$ and $G_y$ are units,
respectively. Their remaining coordinates are nonzero since
$G(t,0,0)=0$. Since $h=tyz$, the two Jacobians are
\[
 \left.\det\frac{\partial(G,h)}{\partial(y,z)}\right|_{y=0}
       =-tzG_z,\qquad
 \left.\det\frac{\partial(G,h)}{\partial(y,z)}\right|_{z=0}
       =tyG_y.
\]
Both are nonzero. A local square root of the unit $U$ makes
$(u,v)=(G-\eta,h\sqrt U)$ analytic coordinates at each
point, and the perturbed transverse equation is
\[
 (G-\eta)^2+h^2U=u^2+v^2=(u+iv)(u-iv).
\]
The invertible map $(u,v)\mapsto(u+iv,u-iv)$ therefore identifies
each of these two germs with an ordinary node. The same argument
applies on the other axes.

\end{remark}


\subsection{The three-branch tacnode defect budget}

\begin{proposition}[Conductor defects at three-branch critical points]\hypertarget{tag-02XM}{}\label{night:P4.48}\label{stmt:P4.48}
\StatementRecord{night:P4.48}{proved}{02XM}{\SeriesRef{night:P4.42}, \SeriesRef{night:P4.47}}

For a fibre support in \cref{def:II:standing-torus-fibration}, at a planar three-branch critical-curve point with uniform
generic transverse type $A_3$, the local conductor defect
is at least ten. At an ordinary nonplanar three-branch
point, the defect is at least nine. Equality in the
second assertion forces Hessian orders $(2,2,2)$,
projected colengths $(1,1,1)$, descent colength one,
and a quartic square $(a xy+b xz+c yz)^2$ with at
least two nonzero coefficients. The case $abc\ne0$
has rational normalization by \cref{night:P4.47}.
\end{proposition}
\begin{proof}

\begin{enumerate}[label=\arabic*.,leftmargin=*]
\item For a planar critical curve the ideal is $P=(u,h)$.
The rank-one nilradical $J=P/I$ of the conductor image
is generated by at most two elements over $\mathcal O_C$.
Its normalized image has at least one generator value
on each of the three branches. Modulo the three branch
maximal ideals, the source supplies a vector space of
dimension at most two, while the normalized image
supplies dimension three. Hence its descent colength
$q$ is at least one. The planar symmetric-square and
different calculation in \cref{night:P4.42} gives local defect
$\ge3\delta+q$. Since a planar three-branch point has
$\delta\ge3$, its defect is at least ten.

\item For an ordinary nonplanar point use coordinates
with $P=(xy,xz,yz)$. Apply the first proof item of
\cref{night:P4.47} to this identical ideal and reduced equation:
$P^{(2)}=P^2+(xyz)$, the rank-one generic transverse Hessian
kills the constant coefficient of $xyz$, and
$\mathfrak m xyz\subset P^2$. Thus $f\in P^2$.
This calculation uses only the three axes and their generic
transverse type, before the nonzero-coefficient hypothesis in
that proposition is used. On each axis, its normalized conormal
lattice is exactly $tK$, so the projected colength
is exactly one. Every Hessian order is at least two.
Write the quartic form in the basis $(xy,xz,yz)$ as
\[
 A(xy)^2+C(xz)^2+E(yz)^2
      +2B(xy)(xz)+2D(xy)(yz)+2F(xz)(yz).
\]
The rank-one Hessians give
$AC=B^2$, $AE=D^2$, $CE=F^2$.
If a diagonal coefficient vanishes, its adjacent
off-diagonal coefficients vanish and the remaining
two-by-two block has rank at most one. If all three
diagonal coefficients are nonzero, the null direction on the
$x$-axis is $y=-(B/A)z$. The coefficient linear in $x$ of the
normal cubic on that direction is
$2(B/A)(DB/A-F)xz^3$. Generic transverse type $A_3$ makes it
zero, giving $F=BD/A$. Higher-degree terms contribute higher
powers of $x$ to that normal cubic and cannot alter this coefficient. Thus the whole quartic is a
square $(a xy+b xz+c yz)^2$.

\item If the quartic is zero, all three Hessian orders are
at least three, and the local defect is at least
$2+3+9/2=19/2$. If the square has just one nonzero
coefficient, say $a xy$, the $z$-axis Hessian order
is at least three. In the source-generator basis
$(xy,xz,yz)$, the three leading quotient-value rows
have the forms
\[
                      (0,1,0),\qquad(0,0,1),
                      \qquad(0,*,*).
\]
Their rank is at most two, so $q\ge1$ and again
the defect is at least $19/2$.

\item If at least two coefficients are nonzero, all three
Hessian orders are two. The leading quotient-value
matrix has rank two. For example, up to rescaling
its rows it is
\[
 \begin{pmatrix}-b&a&0\\-c&0&a\\0&-c&b\end{pmatrix}.
\]
It has zero determinant and rank two whenever at
least two of $a,b,c$ are nonzero. Therefore $q\ge1$,
and the local defect is at least $2+3+3+1=9$.
Equality forces $q=1$ and the displayed Hessian and
projection orders. Proposition~\ref{night:P4.47} applies when all three
quadratic coefficients are nonzero.

\end{enumerate}
\end{proof}


\subsection{The remaining ordinary-axis conductor}

\begin{proposition}[Rational normalization at a three-axis critical point]\hypertarget{tag-02XN}{}\label{night:P4.49}\label{stmt:P4.49}
\StatementRecord{night:P4.49}{proved}{02XN}{\SeriesRef{night:P4.41}, \SeriesRef{night:P4.47}, \SeriesRef{night:P4.48}; \cref{lem:II:quadratic-conductor-normalization,lem:II:boundary-deformation-transport}}

For a fibre support in \cref{def:II:standing-torus-fibration}, at an ordinary three-axis critical-curve point with uniform
generic transverse type $A_3$ and local defect nine, the
normalization has rational singularities. This includes the
quartic square with exactly one missing coefficient left
by \cref{night:P4.47,night:P4.48}.
\end{proposition}
\begin{proof}

\begin{enumerate}[label=\arabic*.,leftmargin=*]
\item Write $f$ for the reduced local equation,
$R=\mathbb C\{x,y,z\}/(f)$ and $I$ for its ambient conductor ideal.
Proposition~\ref{night:P4.48} gives a quartic square with at least
two nonzero coefficients, and \cref{night:P4.47} proves the result
when all three are nonzero. In the remaining case, permute the
axes to write the leading conductor form as $a xy+b xz$ with
$ab\ne0$. The map $(x,y,z)\mapsto(x,ay,bz)$, with inverse
$(x,Y,Z)\mapsto(x,Y/a,Z/b)$, gives the form $x(y+z)$.
As in \cref{night:P4.47}, the descent colength one lifts the leading
kernel vector to $g\in I$ with this quadratic part.
Write $g=x(\alpha y+\beta z)+yz\gamma$, where
$\alpha,\beta$ are units and $\gamma(0)=0$.
Replacing $y,z$ by $\alpha y,\beta z$ is an analytic
coordinate change preserving the axes. Put $k=y+z$,
$h=yz$. Then $g=xk-hc$, with $c(0)=0$.
The ideal $I_0=(g,xh,h^2)$ lies in the conductor
$I$ and has Hilbert--Burch matrix
\[
                 \begin{pmatrix}h&0\\-k&h\\c&-x\end{pmatrix}.
\]
It has radical $(xy,xz,yz)$, is Cohen--Macaulay, and
agrees with $I$ generically on each axis. The
torsionfree finite-kernel comparison used in \cref{night:P4.47}
therefore gives the exact equality $I=I_0$.

\item There are exactly two required ideal forms. Since
$\mathcal O/(x,k,h)=\mathbb C\{y\}/(y^2)$, analytic
division writes $c=\lambda y+xA+kB+hC$ for a scalar
$\lambda$. Replace $g$ by $g+xhA+h^2C$, and then
set $x'=x-hB$. These generator and coordinate changes
give
\[
                    I=(x'k-\lambda yh,x'h,h^2).
\]
If $\lambda=0$ this is $(xk,xh,h^2)$. If
$\lambda\ne0$, rescaling $x'=\lambda X$ gives
$(Xk-yh,Xh,h^2)$. Every change is invertible near
the origin. We treat the two forms separately.

\paragraph{The ideal $(xk,xh,h^2)$.}
Appending $(a,b,c)$ to its Hilbert--Burch matrix gives
determinant $axk+bxh+ch^2$. Put $\mathfrak c=I/(f)\subset R$.
Apply the adjugate form of \eqref{eq:II:manual:4.41b}
and \cite[Theorem~1.12, p.~532]{DeJongVanStraten1990Normalization}
through the quotient $\mathbb C\{x,y,z\}\to R$: its ambient
conductor ideal is this $I$, its hypersurface ideal is $\mathfrak c$,
and its Cohen--Macaulay quotient is $R/\mathfrak c$.
Thus each adjugate entry belongs to $I$. The colon ideals
\[
 I:h=(x,h),\qquad I:x=(k,h),\qquad I:k=(x,h^2)
\]
give exactly
$f=A x^2k^2+B x^2kh+Cxh^2+Dh^4$, with $A(0)\ne0$.
Complete the square in $k$, absorb $\sqrt A$ into $x$,
and replace $y$ by $y(1+\gamma z)$ to make
$k+\gamma h$ the new $y+z$. These are analytic
coordinate maps; $h$ changes only by a unit.
Absorbing that unit in the coefficients gives
\begin{equation}
                       f=x^2k^2+xh^2C+h^4D.          \label{eq:II:first-axis-conductor-equation}
\end{equation}
If $C(0)\ne0$, substitute $x=h^2v$ and divide by
$h^4$. The finite $v$-root is simple, giving a smooth
factor $x-h^2v(y,z)$. Its quotient is $xk^2+h^2U$
with $U$ a unit. The coordinate $X=x/U$ turns it
into $Xk^2+h^2$. Its finite birational normalization is
\[
                 u=h/k,\qquad X=-u^2,
                 \qquad u(y+z)=yz.
\]
The last equation is an $A_1$ quadric, so the
normalization has rational singularities.

\item Suppose $C(0)=0$. Append
$(a,b,d)=(xk,hC,h^2D)$ to the Hilbert--Burch matrix of
$I=(xk,xh,h^2)$. Its determinant is the displayed $f$.
The same finite-duality and adjugate identification gives the
normalization generated by $1,v,w$, with
\[
 v=xk/h,\qquad w=xk^2/h^2+C.
\]
Indeed the adjugate fractions are $n_1=-a/h=-v$ and
$n_2=-(ak+bh)/h^2=-w$. The element $h=yz$ is a nonzerodivisor
modulo $f$: its restrictions to $y=0$ and $z=0$ are respectively
$x^2z^2$ and $x^2y^2$. Division in the total quotient ring gives
\[
 v^2=-xC-Dh^2,\qquad vw=-Dkh,\qquad w^2=Cw-Dk^2.
\]
The fibre is a nonzero quotient of
$\mathbb C[v,w]/(v,w)^2$. Apply the local-fibre assertion of
\cref{lem:II:quadratic-conductor-normalization} to this normalization.
It makes the original reduced equation irreducible. Write its full
map as $\mathcal F=Vf^M$ and choose $f_0=V^{1/M}f$. Replacing $x$ by
$X=V^{1/(2M)}x$, with $y,z$ fixed, has nonzero $x$ derivative
at the origin. Its inverse is the unique local solution $x=\chi(X,y,z)$
of this displayed equation. It writes this root
in the same form $x^2k^2+xh^2C+h^4D$. Its new coefficients are the
transported $V^{1/(2M)}C$ and $V^{1/M}D$, so $C(0)=0$ remains true.
Thus $\mathcal F=f_0^M$ before the following root deformation.

Perturb $D$ by a small constant
and the linear part of $C$ generically, retaining
$C(0)=0$, and then replace $C$ by $C-\epsilon$.
On a fixed boundary the transverse tacnode coefficients
are units, so these changes retain boundary $A_3$.
Under $x=\epsilon^3X$, $y=\epsilon Y$,
$z=\epsilon Z$, the leading scaled equation is
\[
 \frac f{\epsilon^8}=F+O(\epsilon),\qquad
 F=X^2K^2+XH^2(L-1)+D_0H^4,
 \quad K=Y+Z,\ H=YZ.
\]
For $L=K$, $D_0=1$, put $W=Y-Z$.
At $W=0$ the $X$-critical coordinate is
$X=-K^2(K-1)/32$, and the remaining derivative is
$K^5(12K^2+7K-3)/512$.
At either root of $12K^2+7K-3$, the critical
value of the displayed polynomial $F$ and its Hessian determinant are respectively
\[
 \frac{K^6(K+1)(3K-1)}{1024},\qquad
 -\frac{3K^{10}(K+1)(3K-1)(28K^2+14K-5)}{32768}.
\]
They are nonzero: their numerator resultants with the
critical polynomial are $8748$ and $4923269424$.
Nondegeneracy, nonzero value and avoidance of the axes
hold on a nonempty algebraic open set of coefficients.
An arbitrarily small generic perturbation is therefore
available. The implicit-function theorem lifts these
Morse points through the higher weighted terms, with
nonzero critical values of order $\epsilon^8$.
The root boundary comparison is explicit: on every fixed axis
link the transverse quadratic and quartic coefficients remain units,
so parameter Morse coordinates and a fourth root give
$X^2+Y^4$, with their implicit inverses retaining the curve coordinate.
The curve direction is transverse to the boundary. Apply
\cref{lem:II:boundary-deformation-transport} with $\mathcal F=f_0^M$
and its nested representative maps. The square-$-2$ Morse sphere
therefore lies in the original full torus fibre, a contradiction.
Thus $C(0)=0$ is excluded.

\paragraph{The ideal $(g,xh,h^2)$, where $g=xk-yh$.}
Append $(a,b,d)$ to
\[
                         \begin{pmatrix}h&0\\-k&h\\y&-x\end{pmatrix}.
\]
Its determinant is $ag+bxh+dh^2$. The new adjugate
entries give the conditions
\[
       ax,\ ah,\ bx+dh,\ kd+by,\ dh-ay,\ ak+bh\in I.
\]
The exact colon identities are
\[
 I:h=(x,h),\qquad I:x=(xk,h,z^2),\qquad I:x^2=(k,h).
\]
The first two conditions give $a=hA$ modulo $I$;
then $dh-ay$ gives $d=Ay+xD+hE$.
Initially $bx+dh$ permits
$b=-Ak+hB+z^2B_2$ modulo $I$.
The condition $ak+bh$ forces $z^2B_2\in(x,h)$,
so $B_2\in(x,y)$. Its $x$-part lies in $I$, since
$xz^2=h^2-xh$ modulo $I$, and its $y$-part is
absorbed in $hB$. Thus $b=-Ak+hB$ modulo $I$.
The remaining condition is $y(D+B)+kE\in(x,h)$.

\item Write $T=D+B$. Modulo $(x,h,k)$, the equality
$yT=0$ in $\mathbb C\{y\}/(y^2)$ forces $T(0)=0$.
Write $T=xW+yU+kV$. Since $k$ is a nonzerodivisor
modulo $(x,h)$, the general solution is
$E=-yU-yV+xE_1+hE_2$. Substitution in the determinant
gives $f\in(p)+I^2$, where
\[
                            p=yh^2(x-h).
\]
The exact identities
$xh^3=ygh^2-yp$ and $h^4=ygh^2-kp$ then put the
equation in the form
\begin{equation}
 f=A g^2+B gxh+Cgh^2+D(xh)^2+\kappa p,
                       \qquad A(0)\ne0.              \label{eq:II:second-axis-conductor-equation}
\end{equation}

\item For both values of $\kappa(0)$, use the appended coefficients
$a=Ag+Bxh+Ch^2$, $b=Dxh+\kappa yh$,
$d=-\kappa yh$. The multiplier identification
$\operatorname{Hom}_R(\mathfrak c,R)=\overline R$, with
$\mathfrak c=I/(f)$ for this conductor ideal, gives the
normalization generated by $1,n_1,n_2$, where the
adjugate fractions are
\[
          n_1=-a/h,\qquad n_2=-(ak+bh)/h^2.
\]
Since $A$ is a unit, the same normalization
algebra is generated by $1,v,w$, where
$v=(n_1+Bx+Ch)/A=-g/h$ and $w=n_2$.
Direct multiplication gives the exact relations
\begin{align*}
 Av^2-Bxv-Chv+Dx^2+\kappa xy-\kappa yh&=0,\\
 xw&=Ayv-Bxy-(\kappa+C)yh,\\
 vw&=\kappa h-Dxy,\\
 w^2+[(\kappa+B)y+Ck]w-A\kappa k+ADy^2
             -D(\kappa+C)yh-B\kappa h&=0.
\end{align*}
Since $A(0)\ne0$, reduction modulo the original maximal ideal
gives $v^2=vw=w^2=0$. The fibre algebra is local, so the first
assertion of \cref{lem:II:quadratic-conductor-normalization}, applied
to this same finite normalization, makes $f$ irreducible.

First take $\kappa(0)\ne0$. The last equation eliminates $k$
analytically, with $k$ in the maximal-ideal square.
Thus $h=y(k-y)$ has quadratic initial term $-y^2$.
Write $a=A_0$, $b=B_0$, $d=D_0$, $u=\kappa_0$.
The three quadratic equations become
\[
 av^2-bxv+dx^2+uxy=0,\qquad
 xw-ayv+bxy=0,\qquad vw+uy^2+dxy=0.
\]
Set $V=av-(b/2)x$, $\delta=ad-b^2/4$,
$Y=\delta x+auy$, $W_0=w+(b/2)y$, and
$W=auW_0+\delta V$. This invertible linear change
turns them into
\[
                       V^2+xY=0,\qquad
                       VW+Y^2=0,\qquad xW-YV=0.
\]
They define the irreducible rational normal cubic cone.
The full tangent-cone ideal contains this prime ideal
and has the same dimension two, so it equals it.
Blowing up the normalization point is smooth along its
exceptional rational normal cubic: the two independent
initial equations have independent derivatives along
that smooth projective curve in every blowup chart.
The exceptional curve is $\mathbb P^1$ with normal
bundle $\mathcal O(-3)$. The exact sequences of its
thickenings have successive terms $\mathcal O(3n)$,
all with $H^1=0$. Formal functions gives
$R^1\rho_*\mathcal O=0$. The singularity is rational.

\item Finally suppose $\kappa(0)=0$. By the same local normalization
fibre, the original full map is $\mathcal F=Vf^M$. Its root $f_0=V^{1/M}f$
has the same conductor ideal and the same displayed ring-condition
form: multiply each of $A,B,C,D,\kappa$ by $V^{1/M}$.
The first coefficient remains a unit and $\kappa(0)=0$ is preserved.
Use this root for the following deformation. Perturb $D$ and the
linear part of $\kappa$ generically, and replace
$\kappa$ by $\kappa-\epsilon$. The fixed axes retain
boundary $A_3$: equation \eqref{eq:II:second-axis-conductor-equation} stays in the same
ring-condition module, and its boundary transverse
quartic coefficients remain units.
Scale $x=\epsilon^2X$, $y=\epsilon Y$, $z=\epsilon Z$.
Put $K=Y+Z$, $H=YZ$, $G=XK-YH$.
After division by $\epsilon^6$, the first terms are
\[
 A_0G^2+\epsilon(B_0GXH+C_0GH^2)
   +\epsilon^2\bigl[D_0(XH)^2+p(L-1)\bigr].
\]
Near a nonzero smooth point of $G=0$, eliminate the
normal Morse coordinate. On that cone the identities
$XH^3=-Yp$, $H^4=-Kp$ absorb the square-completion
corrections into the coefficients, leaving the leading
tangential function
\[
                    \Phi=D'(XH)^2+p(L'-1)|_{G=0}.
\]
Choose $D'=1$, $L'=-Z$ to exhibit a nonempty open
set of suitable coefficients. The parametrization
\[
 K=r,\quad Y=tr,\quad Z=(1-t)r,\quad
 X=t^2(1-t)r^2
\]
gives $\Phi=t^4(1-t)^4r^7[1+r(t^2-t+1)]$.
At $t=1/2$, $r=-7/6$, its value and Hessian are
\[
 -\frac{823543}{573308928},\qquad
 \operatorname{diag}\left(\frac{117649}{1990656},
                          \frac{15647317}{214990848}\right).
\]
The Hessian determinant is
$1840891197733/427972821516288$, nonzero.
The angular critical numerator factors as
$(2t-1)(3t^2-3t-4)$, so this point is simple and
avoids the axes and $K=0$. The generic coefficient
condition is algebraic and nonempty. After normal
elimination, divide the tangential critical equations
by $\epsilon^2$ and apply the implicit-function theorem.
The resulting three-variable Morse points have nonzero critical
values of order $\epsilon^8$ and tend to the origin.
The fixed-axis boundary charts are again the parameter Morse
and quartic-root maps $X^2+Y^4$, with the same transverse curve
direction and inverses. Apply \cref{lem:II:boundary-deformation-transport}
to this root family and the exact original equality $\mathcal F=f_0^M$.
Its nested maps preserve the square-$-2$ sphere and its push-off
in the original full torus fibre. This contradicts the torus's
asphericity and excludes $\kappa(0)=0$, completing both ideal forms.

\end{enumerate}
\end{proof}

\subsection{The tacnode and transverse-line critical curve}

\begin{proposition}[The defect of a tacnode with a transverse line]\hypertarget{tag-02Y9}{}\label{night:P4R.16}\label{stmt:P4R.16}
\StatementRecord{night:P4R.16}{proved}{02Y9}{\SeriesRef{night:P4.33}, \SeriesRef{ii:excess-one-curve-germs}, \SeriesRef{night:P4.42}, \SeriesRef{night:P4.44}}

In the setting of \cref{lem:II:positive-ruled-conductor-graph},
with generic transverse type $A_3$ along every conductor component,
the three-branch point of the critical curve cannot
be the union of a plane tacnode and an independent transverse
line. Its local conductor defect is at least eleven, whereas
the total defect in the contracted-section case is nine.
\end{proposition}
\begin{proof}

\begin{enumerate}[label=\arabic*.,leftmargin=*]
\item Lemma~\ref{ii:excess-one-curve-germs} gives ambient
coordinates for this curve. Its reduced ideal is
\[
 P=(xz,yz,h),\qquad h=y^2-x^4.
\]
Its three parametrizations are $(t,t^2,0)$, $(t,-t^2,0)$,
and $(0,0,t)$. The two tangent branches have intersection
multiplicity two, and their union meets the independent line
with intersection length one. The normalization sequence
therefore gives $\delta=3$. Every branch is immersive, so
its ramification divisor has coefficient zero.

\item Apply the symbolic-square identity in the proof of
\cref{night:P4.38} to the identical ambient coordinates, with
$h=y^2-x^4$ and $k=4$. That identity concerns only this reduced
curve ideal and its generic primes. It gives
\[
 P^{(2)}=(z,h)^2\cap(x,y)^2=P^2+(zh).
\]
The equality on the right uses $P=(xz,yz,h)$: the three
quadratic $z^2$ generators are in $P^2$, and the remaining
two generators are $zh,h^2$.
Furthermore
\[
 xzh=(xz)h,\qquad yzh=(yz)h,\qquad
 z^2h=(yz)^2-x^2(xz)^2.
\]
The class of $zh$ modulo $P^2$ is therefore annihilated
by the ambient maximal ideal. It is nonzero, since it has
a nonzero cubic term and $P^2$ has order at least four.
Thus $P^{(2)}/P^2$ is exactly one copy of the residue field.

\item Let $f$ be the reduced defining function of the support.
Its generic transverse multiplicity is two, so
$f\in P^{(2)}$ and
\[
                       f=\kappa zh+f_0,\quad f_0\in P^2,
                       \quad\kappa\in\mathbb C.
\]
At $(t,t^2,0)$ with $t\ne0$, use $z,h$ as transverse
coordinates. The quadratic transverse part of $f$ has the
form
\[
 t^2a(t)z^2+(\kappa+t b(t))zh+c(t)h^2,
\]
where $a,b,c$ are holomorphic at $t=0$. The transverse
Hessian has rank one, so its determinant vanishes for
all sufficiently small nonzero $t$. Its constant term at
zero is $-\kappa^2$ (up to the common convention
for the quadratic-form matrix). Hence $\kappa=0$, and
we have proved the exact inclusion $f\in P^2$.

\item Let $K_i$ be the saturated rank-two kernel of the derivative
of the $i$th parametrization, and let $E_i\subset K_i$ be
the span of the differentiated generators of $P$. On the
two tangent branches choose the bases
\[
                dz,\quad dy\mp2t\,dx.
\]
The generator matrices have columns
\[
 \begin{pmatrix}t&\pm t^2&0\\0&0&\pm2t^2\end{pmatrix}.
\]
Their elementary-divisor orders are $(1,2)$. On the
independent line use $dx,dy$; the matrix is
\[
                   \begin{pmatrix}t&0&0\\0&t&0\end{pmatrix},
\]
with orders $(1,1)$.

\item The preceding calculation proves $f\in P^2$.
On each branch apply the symmetric-square clause of
\cref{lem:II:rank-one-conormal-lattice} with $E=E_i\subset K_i$,
its primitive Hessian direction, $d=d_i$ and $\beta=b_i$.
The three exact elementary-divisor pairs just computed give
the respective bounds $3,3,2$ for $d_i/2+b_i$.
The local defect
formula of \cref{night:P4.42} consequently gives
\[
 \delta+\sum_i(d_i/2+b_i-r_i)+q
       \ge3+(1+2)+(1+2)+(1+1)=11.
\]
Here $q\ge0$ is the normalization descent colength. \Cref{night:P4.44} contracts
the horizontal section, and \cref{night:P4.33} and \eqref{eq:II:refined-4-42a} give total
defect nine. The displayed lower bound contradicts that
total. This excludes the stated analytic curve germ.

\end{enumerate}
\end{proof}

\subsection{The nonordinary three-branch tacnode sector}

\begin{proposition}[Defect bounds at nonordinary three-branch points]\hypertarget{tag-02YA}{}\label{night:P4R.17}\label{stmt:P4R.17}
\StatementRecord{night:P4R.17}{proved}{02YA}{\SeriesRef{ii:excess-one-curve-germs}, \SeriesRef{night:P4.42}, \SeriesRef{night:P4.48}, \SeriesRef{night:P4R.10}, \SeriesRef{night:P4R.15}, \SeriesRef{night:P4R.16}}

In the setting of \cref{lem:II:positive-ruled-conductor-graph},
with generic transverse type $A_3$ along every conductor component, every nonordinary three-branch critical-curve point has
local conductor defect strictly greater than nine. In particular,
the delta-four nonplanar cases have the explicit lower bounds
$13$, $21/2$, $11$, and $23/2$ below. The only three-branch germ
compatible with the total defect budget nine is the ordinary
nonplanar union of three coordinate axes.
\end{proposition}
\begin{proof}
\begin{enumerate}[label=\arabic*.,leftmargin=*]
\item 
Use the notation of \cref{night:P4R.10} and \cref{night:P4.42}: \[
 D(C,f)=\delta(C)+
        \sum_i\left(\frac{d_i}{2}+b_i-r_i\right)+q,\qquad q\ge0.
\] Here \(d_i\) is the order of the saturated rank-one Hessian, \(b_i\) is the projected conormal colength, and \(r_i=m_i-1\) for a branch of multiplicity \(m_i\).

\item \emph{Bounds from Hessian columns}\label{bounds-without-an-ordinary-square-assumption}

Let \(E_i\subset K_i\) be the normalized conormal lattice, with
elementary-divisor orders \(e_1\le e_2\). Every first derivative of
\(f\) belongs to the radical ideal \(P\) of the critical curve, so
every Hessian column belongs to \(E_i\). Apply the column clause of
\cref{lem:II:rank-one-conormal-lattice} with
\(a=e_1,b=e_2,d=d_i,\beta=b_i\) and the same saturated Hessian line.
It gives
\begin{equation}\label{eq:II:nonordinary-triple-a}
                  \frac{d_i}{2}+b_i
                    \ge e_1+\frac{e_2}{2}.
\end{equation}
This clause applies to every reduced defining equation with the
stated gradient condition, including its symbolic-square terms.

Since the curve is nonplanar, \(P\) has no linear equation. Every defining gradient therefore has order at least \(m_i\) on branch \(i\), and \(b_i\ge m_i\). \Cref{night:P4R.15} gives zero ambient Hessian, so \(f\) has order at least three and \(d_i\ge m_i\). For total multiplicity \(m=\sum m_i\), \begin{equation}\label{eq:II:nonordinary-triple-b}
                  D(C,f)\ge\delta(C)+3+\frac m2.
\end{equation} On a branch of multiplicity \(m_i\), the classes of
\(t_i,\ldots,t_i^{m_i-1}\) in its normalization quotient are linearly
independent, giving \(\delta_i\ge m_i-1\). The three branch values
modulo their common value add two gluing dimensions. Thus
\(\delta\ge\sum_i(m_i-1)+2=m-1\), and \(m\ge5\) gives
\(D\ge(m-1)+3+m/2\ge19/2>9\).
For the following multiplicity-three and multiplicity-four
classification, assume \(\delta=4\); the other delta invariants
are treated in the final item.

\item \emph{Three smooth branches: the complete excess-two classification}\hypertarget{tag-0005}{}\label{three-smooth-branches-the-complete-excess-two-classification}

Put \[
 B=\bigoplus_{i=1}^3\mathbb C\{t_i\},\quad
 F^j=\bigoplus_{i=1}^3t_i^j\mathbb C\{t_i\},\quad
 A_0=\mathbb C_{\rm diag}+F^1.
\] The curve ring \(A\subset A_0\) has \(\operatorname{length}(A_0/A)=2\), since \(A_0\) has delta invariant two. Each branch projection of \(A\) is onto.

Let \(W_1\subset F^1/F^2=\mathbb C^3\) be its first-jet image. Each coordinate projection is nonzero. If \(W_1=\mathbb C^3\), Nakayama first gives \(\mathfrak m_AA_0=F^1\) and then, on the finite \(A\)-module \(A_0/A\), gives \(A=A_0\). This contradicts colength two. Thus \(\dim W_1=1\) or \(2\).

\item \emph{One-dimensional first jets}\label{a.-one-dimensional-first-jets}

Scale the branch parameters so \(W_1=\mathbb C(1,1,1)\). The two missing first jets use the entire colength, hence \(F^2\subset A\). The exact algebra is \[
 A=\mathbb C_{\rm diag}+\mathbb C(t_1,t_2,t_3)+F^2.
\] It is generated by \[
 x=(t_1,t_2,t_3),\quad y=(t_1^2,0,0),\quad z=(0,t_2^2,0).
\] Indeed \(y,z,x^2-y-z\) give all three quadratic branch ideals, and multiplication by \(x\) gives every higher order. The embedded model is exactly \begin{equation}\label{eq:II:nonordinary-triple-c}
       P=\bigl(y(y-x^2),yz,z(z-x^2)\bigr).
\end{equation} The branches are \((t,t^2,0),(t,0,t^2),(t,0,0)\). Its differentiated generators give \(E_i=t_i^2K_i\) on all three branches. For example, on \((t,t^2,0)\) use the primitive basis \(dy-2t\,dx,dz\). Formula \eqref{eq:II:nonordinary-triple-a} gives \begin{equation}\label{eq:II:nonordinary-triple-d}
                         D\ge4+3\cdot3=13.
\end{equation}

\item \emph{Two-dimensional first jets}\hypertarget{tag-003V}{}\label{b.-two-dimensional-first-jets}

The three branch evaluations are nonzero linear forms on \(W_1\). If no two are proportional, their squares are independent: choose the first two as \(s,t\); the third is \(as+bt\), \(ab\ne0\), whose square supplies the nonzero \(st\) coefficient. Thus products give the whole second-jet space. Multiplication by one first jet nonzero on all branches then gives every later jet. This would leave only one missing positive-jet dimension, contrary to colength two.

Two evaluations are therefore proportional. Scale and label the parameters to obtain \(W_1=\{(a,a,b)\}\). Its products give the same two-dimensional second-jet space. The one remaining missing dimension occurs there: if the second-jet image were full, multiplication by a nonvanishing first jet would make every later one full. Hence \(F^3\subset A\), and the pure second-jet image is exactly \(\{(c,c,d)\}\).

A lift of \(W_1\) has quadratic difference \(\lambda a+\mu b\). The analytic changes \(t_i\mapsto t_i+c_it_i^2\), on the first two branches, remove \(\lambda a\). Thus the exact ring has common constants, with first and second coefficients satisfying \begin{equation}\label{eq:II:nonordinary-triple-e}
                     a_1=a_2,\qquad c_1-c_2=\mu a_3,
\end{equation} and arbitrary independent terms of order at least three.

For \(\mu\ne0\), the two elements \[
 x=(t_1,t_2,0),\qquad y=(\mu t_1^2,0,t_3)
\] generate this whole ring. The exact identities \[
 \begin{aligned}
 xy/\mu&=(t_1^3,0,0),\\
 x^3-xy/\mu&=(0,t_2^3,0),\\
 y^3-\mu(xy)^2&=(0,0,t_3^3),\\
 y^2-\mu x^2y&=(0,0,t_3^2)
 \end{aligned}
\] give the independent cubic ideals and every permitted quadratic jet. Multiplication by \(x\) or \(y\) gives all the higher branch series. This ring is planar, contrary to the present hypothesis.

Thus \(\mu=0\). The ring is generated by \((t_1,t_2,0),(0,0,t_3),(t_1^3,0,0)\). Naming the independent transverse coordinate \(z\) gives the exact embedded model \begin{equation}\label{eq:II:nonordinary-triple-f}
                       P=(xz,yz,y(y-x^3)).
\end{equation} Its branches are \((t,0,0),(t,t^3,0),(0,0,t)\), and its conormal elementary divisors are \[
                         (1,3),\quad(1,3),\quad(1,1).
\] On the second branch use \(dy-3t^2dx,dz\); the gradients give \(t^3(dy-3t^2dx)\) and \(t\,dz\). The column estimate \eqref{eq:II:nonordinary-triple-a} gives \begin{equation}\label{eq:II:nonordinary-triple-g}
                  D\ge4+\frac52+\frac52+\frac32
                       =\frac{21}{2}>9.
\end{equation}

These analytic ring isomorphisms give ambient isomorphisms as follows. In each nonplanar case the three displayed generators form a basis of the curve cotangent space. Lift them to the ambient local ring. Their linear Jacobian is invertible, so the inverse-function theorem gives an ambient coordinate isomorphism to the stated model.

\item \emph{Multiplicity four: one cusp and two smooth branches}\hypertarget{tag-02QH}{}\label{multiplicity-four-one-cusp-and-two-smooth-branches}

The singular branch has multiplicity two and delta at least one. Joining three branches costs at least two further delta units. If the branch delta is two, all remaining gluing is minimal, and the local ring is \[
 \mathbb C_{\rm diag}
 +\mathfrak m_{\rm singular}
 +t_2\mathbb C\{t_2\}+t_3\mathbb C\{t_3\}.
\] Its embedding dimension is at least \(2+1+1=4\), impossible in a smooth threefold. A larger branch delta is impossible. Thus this branch is the ordinary cusp \(\mathbb C\{t^2,t^3\}\).

Set \[
 A_0=\mathbb C_{\rm diag}
 +\mathfrak m_{\mathbb C\{t^2,t^3\}}
 +t_2\mathbb C\{t_2\}+t_3\mathbb C\{t_3\}.
\] Now \(A_0\) has delta three, and \(A\subset A_0\) has colength one. The one-dimensional \(A\)-module \(A_0/A\) is killed by \(\mathfrak m_A\), so \(\mathfrak m_AA_0\subset A\). Its elements have zero constant term, giving \(\mathfrak m_AA_0=\mathfrak m_A\). Hence \(\mathfrak m_A\) is an ideal of \(A_0\). The local algebra \(A_0/\mathfrak m_A\) has length two, so its maximal ideal squares to zero. Therefore \begin{equation}\label{eq:II:nonordinary-triple-h}
                         \mathfrak m_0^2\subset A.
\end{equation}

The four cotangent coordinates of \(A_0\), denoted \(a,b,c,d\), are represented by the cusp functions \(t^2,t^3\) and the two smooth parameters. The three-dimensional subspace \(W=\mathfrak m_A/\mathfrak m_0^2\) is a hyperplane \begin{equation}\label{eq:II:nonordinary-triple-i}
                 \lambda a+\mu b+\nu c+\xi d=0.
\end{equation} Its projection to \(a,b\) is onto, since the branch image is the full cusp ring, and the two smooth projections are nonzero. Ambient embedding dimension three gives \(\mathfrak m_A^2=\mathfrak m_0^2\). Thus products of \(W\) must generate \(\mathfrak m_0^2/\mathfrak m_0^3\), whose coordinates are \begin{equation}\label{eq:II:nonordinary-triple-j}
                              a^2,\quad ab,\quad c^2,\quad d^2.
\end{equation}

\item \emph{A nonzero coefficient of the cusp cubic}\hypertarget{tag-0004}{}\label{a.-a-nonzero-coefficient-of-the-cusp-cubic}

For \(\mu\ne0\), solve \eqref{eq:II:nonordinary-triple-i} for \(b\). Surjectivity of the cusp projection forces \((\nu,\xi)\ne(0,0)\). The three minimal ambient coordinate functions have parametrizations \[
 \begin{array}{c|c}
 \text{cusp}&(t^2-\lambda't^3,-\nu't^3,-\xi't^3)\\
 \text{second branch}&(0,t,0)\\
 \text{third branch}&(0,0,t).
 \end{array}
\] Their tangent directions are independent.

By \cref{night:P4R.15}, \(f\in\mathfrak m^3\). Let \(G_3\) be its homogeneous cubic part. Since \(df\) vanishes on all three branches, \(dG_3\) vanishes on each independent tangent axis. Every cubic monomial other than \(xyz\) has a derivative nonzero on an axis. Hence \(G_3=cxyz\). On the \(x\)-axis its normal Hessian is \[
                         x\begin{pmatrix}0&c\\c&0\end{pmatrix}.
\]
On the cusp, \(x=t^2+O(t^3)\) and \(y,z=O(t^3)\). The
\((y,z)\) block is therefore
\(t^2\left(\begin{smallmatrix}0&c\\c&0\end{smallmatrix}\right)
+O(t^3)\). Terms of degree at least four in \(f\) contribute
Hessian entries of order at least four on this branch. Its leading
minor is \(-c^2t^4\). Generic rank one forces \(c=0\), so
\(f\in\mathfrak m^4\).

The Hessian entries have order at least \(2m_i\) on branch \(i\), and \(b_i\ge m_i\). Its defect contribution is at least \(m_i+1\). Therefore \begin{equation}\label{eq:II:nonordinary-triple-k}
                            D\ge4+(4+3)=11>9.
\end{equation}

\item \emph{A zero coefficient of the cusp cubic}\hypertarget{tag-003U}{}\label{b.-a-zero-coefficient-of-the-cusp-cubic}

For \(\mu=0\), surjectivity in \eqref{eq:II:nonordinary-triple-j} forces \(\lambda,\nu,\xi\) all to be nonzero. Otherwise two of \(a,c,d\) are proportional on \(W\), so their squares are dependent, and \eqref{eq:II:nonordinary-triple-j} has rank at most three. That creates an extra minimal generator, contrary to embedding dimension three. Scale the branch parameters so that \(a+c+d=0\).

The pure hyperplane generators belong to \(A\) exactly by \eqref{eq:II:nonordinary-triple-h}. They give the parametrizations \[
                       (t^2,t^2,t^3),\quad(t,0,0),\quad(0,t,0).
\] Their defining ideal is \begin{equation}\label{eq:II:nonordinary-triple-l}
                 P=(z(x-y),xy(x-y),z^2-x^2y).
\end{equation} A Hilbert--Burch matrix is \[
 \begin{pmatrix}
 xy&z\\
 -z&-x\\
 0&-(x-y)
 \end{pmatrix}.
\] Its maximal minors give \eqref{eq:II:nonordinary-triple-l}, with alternating signs. The ideal has codimension two and is Cohen--Macaulay. Its zero set is exactly the cusp and the two axes, and it is generically reduced on each branch. A nonzero nilpotent quotient supported at the origin would violate depth one, so it is the reduced intersection ideal.

On the cusp use the primitive conormal basis \[
                     k_1=dx-dy,\qquad k_2=3t\,dx-2dz.
\] The differentiated generators are \(t^3k_1,t^4k_1,t^3(tk_1-k_2)\). Thus its elementary divisors are \((3,3)\). On each smooth axis they are \((1,2)\). Here \(r=1\) on the cusp and zero on the axes. By \eqref{eq:II:nonordinary-triple-a}, \begin{equation}\label{eq:II:nonordinary-triple-m}
                  D\ge4+\left(\frac92-1\right)+2+2
                       =\frac{23}{2}>9.
\end{equation}

This exhausts multiplicity four. The multiplicity-three and multiplicity-four computations prove the strict local bound for every nonplanar three-branch curve with delta four.

\item \emph{Every nonordinary triple exceeds the budget}\hypertarget{tag-01OO}{}\label{every-nonordinary-triple-exceeds-the-budget}

For a planar three-branch point, \cref{night:P4.48} gives defect at least ten, including the ordinary plane triple point. At a nonplanar point with delta at least five, \eqref{eq:II:nonordinary-triple-b} and \(m\ge3\) give \[
                  D\ge5+3+\frac32=\frac{19}{2}>9.
\] Delta four was exhausted above.

At delta three the excess \(\delta-(r-1)\) is one. Lemma~\ref{ii:excess-one-curve-germs} gives precisely an ordinary plane triple, a plane tacnode with an independent transverse line, or a cusp with two independent lines. The first is excluded by \cref{night:P4.48}. The second has defect at least eleven by \cref{night:P4R.16}. The third has embedding dimension four and cannot occur in the ambient smooth threefold.

Finally a three-branch point with delta two has zero excess. Its ring equals \(\mathbb C_{\rm diag}+F^1\), generated by the independent branch parameters. It is exactly the three coordinate axes, with ideal \((xy,xz,yz)\). Thus every nonordinary three-branch point is excluded from the total defect budget nine. The only remaining local possibility is the ordinary nonplanar triple treated by \cref{night:P4.48}, whose conductor equations and normalizations are determined in \cref{night:P4.47,night:P4.49}.


\end{enumerate}
\end{proof}

\subsection{Completion of the uniform tacnodal sector}

\begin{proposition}[Euler bound for uniform tacnodal ruled supports]\hypertarget{tag-02XP}{}\label{night:P4.50}\label{stmt:P4.50}
\StatementRecord{night:P4.50}{proved}{02XP}{\SeriesRef{night:P4.33}, \SeriesRef{night:P4.42}, \SeriesRef{night:P4.44}, \SeriesRef{night:P4.47}, \SeriesRef{night:P4.48}, \SeriesRef{night:P4.49}, \SeriesRef{night:P4R.14}, \SeriesRef{night:P4R.16}, \SeriesRef{night:P4R.17}, \SeriesRef{night:P4R.6}}

In the irreducible-or-normal-crossings setting of
\cref{def:II:standing-torus-fibration}, every support whose
normalization has a resolution ruled over a positive-genus curve
and whose generic transverse type is $A_3$ along every
critical-curve component has nonpositive Euler number.
\end{proposition}
\begin{proof}

\begin{enumerate}[label=\arabic*.,leftmargin=*]
\item Let $mS$ be such a full fibre with $e(S)>0$, and write
$\rho=\nu\sigma:T\to N\to S$ for its resolved normalization.
The connected ruled surface $T$ is irreducible, and hence so is
its image $S$. Apply \cref{lem:II:positive-ruled-conductor-graph}
to this same $f,Y,B,mS,\nu,\sigma$.
\Cref{night:P4.42} and \cref{night:P4.44} leave only a contracted horizontal
positive-genus section and a cycle at one three-branch
point of a conductor with two rational normalized
components. The total defect is nine by \cref{night:P4.33} and
\eqref{eq:II:refined-4-42a}. \Cref{night:P4.48} excludes planar triples; \cref{night:P4R.16}
excludes the nonplanar excess-one tacnode-plus-line
case; the exact excess-two classification and conormal
calculation in \cref{night:P4R.17} exclude
every delta-four nonplanar triple. The branchwise bound
$\delta+3+m/2$ excludes all larger delta invariants,
where $m\ge3$ is total multiplicity.
The only remaining triple is the ordinary three-axis
point. \Cref{night:P4.48} gives local defect at least nine, so it
uses the entire budget, and \cref{night:P4.47} and \cref{night:P4.49} give
rational normalization there.

\item The equality
$\sum_{p\in C}(r_p-1)=2$ in
\cref{lem:II:positive-ruled-conductor-graph} is already exhausted
by this three-branch point. Every additional singular point of
the critical curve would therefore be unibranch. The local
inequalities in \cref{night:P4R.10} give strictly positive defect there,
so such points are absent.
At a smooth point of the entire critical curve,
\cref{night:P4R.14} gives the product tacnode or its $A_1$
normalization degeneration; both have rational
normalization. Away from the conductor, \cref{night:P4R.6}
excludes exceptional configurations. The normalization
therefore has rational singularities everywhere.
Let $E$ be the reduced exceptional fibre of this same
$\sigma:T\to N$ containing the horizontal positive-genus
section $H$. Apply the exceptional-cohomology argument in
\cref{night:P4.35} over a Stein neighbourhood of $\sigma(H)$.
Rationality gives $R^1\sigma_*\mathcal O_T=0$; the ideal
sequence of $E$ and the proved vanishing of $R^2\sigma_*$
give $H^1(E,\mathcal O_E)=0$. Its curve normalization
sequence then forces every normalized component of $E$
to have genus zero. This contradicts the positive genus
of $H$ and proves the assertion.

\end{enumerate}
\end{proof}


\section{Contracted sections and cubic spectral covers}\label{sec:II:contracted-cubic-covers}

\subsection{The horizontal image in the unrestricted ruled remainder}

\begin{proposition}[Transverse type along a noncontracted horizontal section]\hypertarget{tag-02YB}{}\label{night:P4R.18}\label{stmt:P4R.18}
\StatementRecord{night:P4R.18}{proved}{02YB}{}

In the irreducible-or-normal-crossings torus-fibration setting
of \cref{def:II:standing-torus-fibration}, let $mS$ have
irreducible reduced support with $e(S)>0$ and normalization
resolution ruled over a positive-genus curve. A noncontracted
horizontal pole section maps with degree two or three to a
rational normalized conductor component whose generic
transverse germ is $D_4$.
The conclusion applies to every transverse type and
conductor-exponent pattern on the other image component.
\end{proposition}
\begin{proof}

\begin{enumerate}[label=\arabic*.,leftmargin=*]
\item Let $\nu:N\to S$ be the normalization,
$\sigma:T\to N$ its minimal resolution, and
$\rho=\nu\circ\sigma$. Let $D$ be the effective pole divisor
of \cref{night:P4R.6} on $T$.
Write $H$ for the horizontal section, $g=g(H)\ge1$, and
$C_H$ for its image in the reduced conductor $C$. Apply \cref{lem:II:positive-ruled-conductor-graph} to this
$f:Y\to B$ and this full fibre $mS$. Its global support,
positive-Euler and ruling-genus hypotheses are precisely
those in the statement. Thus $C$ has two rational
normalized components and one point-identification cycle.
The nonconstant map
\[
                         H\longrightarrow C_H^\nu=\mathbb P^1
\]
has degree at least two: a degree-one finite map between smooth
curves is an isomorphism, whereas their genera differ. Each
generic point of $H$ has conductor exponent two, by the pole divisor of \cref{night:P4R.6}. Thus the generic transverse plane curve
along $C_H$ has at least two normalization branches of exponent two.

\item Apply \cref{lem:II:plane-branch-conductor} to the reduced transverse equation and its branches $B_i$:
\[
             c_i=2\delta(B_i)+\sum_{j\ne i}I(B_i,B_j).
\]
A singular branch with $c_i=2$ permits no other branch,
so the two specified branches are smooth. With two total
branches their intersection multiplicity is two, giving $A_3$.
With three total branches, the first two exponents force all
three pairwise intersection multiplicities to equal one;
the third branch is then smooth as well, giving $D_4$.
Four branches would give exponent at least three on every
smooth branch. These possibilities are exhaustive. In type
$A_3$, $H$ has degree exactly two and supplies both generic
normalization preimages. In type $D_4$ its degree is two or
three.

\item Assume that the transverse germ along $C_H$ is $A_3$.
Let $C_{H,\rm sch}$ be the primary component of the image
conductor supported on $C_H$. Apply
\cref{lem:II:horizontal-primary-conductor} with this
$\nu,\sigma,D,H$ and $C_H$. The section is noncontracted
by hypothesis, its pole coefficient is two, and the
positive-genus ruling and its rational vertical trees
are the ones just obtained. The same primary restriction
map therefore gives
\begin{equation}
 g-1=\chi(\mathcal O_{C_{\rm sch}})
 \le\chi(\mathcal O_{C_{H,\rm sch}}).
 \label{eq:horizontal-primary-euler}
\end{equation}

\item We evaluate the right side on $C_H^\nu=\mathbb P^1$.
For the normalization map $j:\mathbb P^1\to C_H\hookrightarrow Y$,
use the differentiated conormal, saturated Hessian line
and rank-one nilradical of the primary conductor as in
\cref{night:P4R.10}. The map is
$j^*\Omega_Y\to\Omega_{\mathbb P^1}(-R)$, so its kernel
$K$ has degree $2+\deg R$, since $c_1(Y)\cdot C_H=0$. The Hessian line has degree
$\tfrac12\sum d_p$. Riemann--Roch and the normalization
sequence give exactly
\begin{equation}
 \chi(\mathcal O_{C_{H,\rm sch}})
 =4-\sum_p\left[\delta(C_H,p)
      +\sum_{i\mapsto p}(d_i/2+b_i-r_i)+q_p\right].
                                                            \label{eq:horizontal-primary-defect}
\end{equation}
All terms are those of this primary curve and its primary
conductor. They remain defined at its intersections with
the other component, because the pure conductor quotient
extends its generic containment of the symbolic square.

\item Every singular point of $C_H$ contributes at least three
to the sum in \eqref{eq:horizontal-primary-defect}. At a planar point the ideal of $C_H$ is
$(u,h)$, its symbolic square equals its ordinary square,
and the reduced equation of $S$ belongs to that square. The
symmetric-square estimate and residue calculation in the
second item of the proof of \cref{night:P4.42}, applied
to this ideal $(u,h)$ of $C_H$, give a contribution at
least $3\delta(C_H,p)+q_p\ge3$.
At a nonplanar point every defining equation has order at
least two. On a branch of multiplicity $m_i$ its projected
conormal order is at least $m_i$. The point is singular
on the entire critical curve, so \cref{night:P4R.15} makes the ambient
Hessian zero and gives $d_i\ge m_i$. Since $r_i=m_i-1$,
the contribution is at least
\[
                      \delta(C_H,p)+\sum_i(1+m_i/2),
\]
which is at least three at a singular curve point. The
same arguments give nonnegative contributions everywhere;
at a smooth primary point its conormal is the full kernel
and its only possible contribution is the nonnegative
Hessian zero order divided by two.
Combining \eqref{eq:horizontal-primary-euler} and \eqref{eq:horizontal-primary-defect} gives total primary
defect at most $5-g\le4$. Thus $C_H$ has at most one
singular point, and two different singular points are
impossible even in the presence of arbitrary transverse
types on the other global component.

\item We count the possible branch values of the degree-two map
$H\to\mathbb P^1$ with all point identifications.
The normalization graph formula is
\[
                         \sum_{p\in C}(r_p-1)=2.
\]
Connectedness leaves either two two-branch points or one
three-branch point. In the two-point case, two intersections
of the global components give two points on $C_H^\nu$.
There can be at most one additional unibranch singularity
of $C_H$, so these supply at most three possible branch
values. The other connected possibility is a self-identification
on one component and one intersection of the two components.
A self-identification on $C_H$ supplies two normalization
points and already costs at least three in the primary
defect. It permits no additional singular point, so with
the intersection there are again at most three possible
branch values. A self-identification on the other component
supplies no additional point on $C_H^\nu$.

\item At a three-branch point, the two global components must
supply two and one branches. Thus $C_H^\nu$ has at most
two points over it. Two points already make $C_H$ singular
there and exclude any further singular point by the defect
bound. One point permits at most one further unibranch
singularity. This case also has at most three possible
branch values.

\item There is no ramification elsewhere. At a point where the
entire critical curve is smooth, \cref{night:P4R.14} gives precisely
the product tacnode or $X(Y^2+XZ)$. Both have two distinct
normalization points over that point of the conductor.
For the quadric branch $Y^2+XZ=0$, its resolution chart is
$(s,v)\mapsto(-sv^2,sv,s)$; the strict conductor line $v=0$
maps with coordinate $s=Z$. Thus the resolution retains the two
distinct strict conductor preimages. Their two generic preimages
lie on $H$, so properness
places both special preimages on $H$ as well. The
degree-two map is therefore unramified there. We have
accounted for every possible ramification value.
Riemann--Hurwitz for a degree-two map from a genus-$g$
curve requires $2g+2\ge4$ distinct branch values; each
ramified fibre consists of its single point of index two.
This contradicts the count of at most three. Thus the
generic germ along $C_H$ must be $D_4$, with degree two
or three as proved at the start.

\end{enumerate}
\end{proof}

\subsection{The quartic tangent obstruction}\label{subsec:II:quartic-tangent-obstruction}

\begin{lemma}[Weighted initial forms and monodromy bounds]\hypertarget{tag-P322}{}\label{lem:II:weighted-initial-monodromy}
\StatementRecord{lem:II:weighted-initial-monodromy}{proved}{P322}{\cite[\S1.2, p.~445; \S\S3.4--3.6, pp.~451--452; \S5.3, p.~457; \S6.3, p.~461; \S7.1, p.~462]{Siersma1991Variation}; \cite[proof of Theorem~2C.3, pp.~180--181]{Hatcher2002}; \cite[Theorem~1.2, p.~180]{ParkSuh1996}}
\begin{enumerate}[label=\arabic*.,leftmargin=*]
\item Let $g:(\C^3,0)\to(\C,0)$ be a holomorphic germ
with weighted initial polynomial $P$ of positive degree $d$
for positive integer weights $w=(w_1,w_2,w_3)$. For
$c\ne0$, every finite collection of compact cycles and
intersection chains in $\{P=c\}$ transports, preserving
oriented intersections, into a local level
$\{g=c\epsilon^d\}$ for all sufficiently small positive
$\epsilon$.

\item Let $g$ be reduced and let $F$ be its compact local
Milnor fibre. Write $T$ for its monodromy on $H_2(F;\Q)$
and $M$ for the sum of the transverse plane-curve Milnor
numbers along the one-dimensional critical branches;
put $M=0$ for an isolated critical point. When the rational
intersection form of $F$ is zero,
\begin{equation}
                         \operatorname{rank}(T-I)\le M.
\label{eq:II:manual:4.52b}
\end{equation}

\item For a reduced weighted homogeneous polynomial $P$
of degree $d$ and positive integer weights $w$,
the affine fibre $F=\{P=1\}$ and a compact local Milnor fibre
have the same homology and weighted scalar monodromy.
For every $\zeta\in\mu_d$, acting with weights $w$,
\[
 \sum_{j=0}^2(-1)^j\operatorname{tr}
       (\zeta\mid H_j(F;\Q))=\chi(F^\zeta).
\]
Suppose that the cyclic group $\mu_d$, acting with weights
$w$, acts freely on this fibre, and put
$Q=\{P=1\}/\mu_d$. Its monodromy satisfies
\[
                  \operatorname{rank}(T-I)\ge(d-1)\chi(Q).
\]
\end{enumerate}
\end{lemma}
\begin{proof}
\begin{enumerate}[label=\arabic*.,leftmargin=*]
\item Put $D_\epsilon(z)=(\epsilon^{w_1}z_1,
\epsilon^{w_2}z_2,\epsilon^{w_3}z_3)$ and
$G_\epsilon=\epsilon^{-d}g\circ D_\epsilon$.
These functions extend at $\epsilon=0$ with value $P$
and converge there with all derivatives on compact sets.
Weighted Euler identity gives $dP\ne0$ on $P=c$.
Choose nested compact neighbourhoods of the cycle and
intersection representatives on that level. The Hermitian
gradient divided by its squared norm solves
\[
            dG_\epsilon(V_\epsilon)=-\partial_\epsilon G_\epsilon
\]
on a neighbourhood of these chains. A cutoff equal to one
along the relevant short trajectories makes the flow
defined for small $\epsilon$. It sends the chains into
$G_\epsilon=c$. The holomorphic map $D_\epsilon$ then
puts them into $g=c\epsilon^d$ in any fixed original
coordinate ball. The flow starts with the identity, and
the scaling has its complex orientation, so all their
oriented intersection numbers are unchanged.

\item Let $\partial_2F$ be the boundary pieces over the
critical-link circles, and write
\[
 j:H_2(F;\Q)\to H_2(F,\partial F;\Q),\qquad
 j_2:H_2(F;\Q)\to H_2(F,\partial_2F;\Q).
\]
The chain-level identity of
\cite[\S3.6, pp.~451--452; \S7.1, p.~462]{Siersma1991Variation}
is
\begin{equation}
                  j_2(T-I)=\operatorname{VAR}^{II}j.
\label{eq:II:manual:4.52a}
\end{equation}
It holds in complex fibre dimension two with the displayed
unmodified maps. Poincar\'e--Lefschetz duality identifies
$j$ with the intersection pairing. Thus $j=0$ makes
$\operatorname{im}(T-I)$ lie in $\ker j_2$, the image of
$H_2(\partial_2F;\Q)$. Each boundary piece is the mapping
torus of a transverse plane-curve Milnor fibre $F_i$.
The dimension-two Wang sequence of
\cite[\S5.3, p.~457; \S6.3, p.~461]{Siersma1991Variation}
gives
\[
 H_2(\partial_2F;\Q)=\bigoplus_i\ker(A_i-I)
                   \subset\bigoplus_iH_1(F_i;\Q),
\]
where $A_i$ is vertical monodromy. This proves
\eqref{eq:II:manual:4.52b}. For an isolated critical
point the same conclusion follows from
$\ker j=\ker(T-I)$ in
\cite[Proposition~3.4, p.~451]{Siersma1991Variation}.

\item Choose an integer $L$ divisible by every $w_i$.
The weighted norm
\[
                  R_w(z)=\sum_i|z_i|^{2L/w_i}
\]
is proper and semialgebraic on $P=1$. Its critical values
form a finite set: they are semialgebraic and have measure
zero by Sard's theorem. Above the largest value, the
normalized gradient supplies a product end. A large
regular sublevel is therefore a compact core invariant
under weighted scalar monodromy. The identity
$R_w(D_\epsilon z)=|\epsilon|^{2L}R_w(z)$ shows that the
critical values on $P=t$ are $|t|^{2L/d}$ times those on
$P=1$. Weighted scaling and the same product collar
identify the core with a compact local Milnor fibre.
Reducedness gives a connected fibre and homology only
in degrees zero, one and two, by the Milnor-fibre
statements in \cite[\S1.2, p.~445; Remark~3.3, p.~451]{Siersma1991Variation}.

To compute the character on this core $C=\{P=1,R_w\le R_0\}$,
choose $R_0$ above all critical values and regular also
on every fixed locus. Its algebraic double is
\[
 \widehat C=\{(z,t):P(z)=1,\ t\in\R,
                         t^2=R_0-R_w(z)\}.
\]
It is compact. The real algebraic action of
$\mu_d\times\{\pm1\}$ acts on $z$ by weighted scalars
and on $t$ by its sign. The equivariant triangulation
of \cite[Theorem~1.2, p.~180]{ParkSuh1996} respects
the isotropy strata. After barycentric subdivision,
the equator $t=0$, fixed by the second factor, is a
subcomplex, and each closed half is a subcomplex.
Projection $(z,t)\mapsto z$ identifies the half $t\ge0$
with $C$, with inverse $z\mapsto(z,\sqrt{R_0-R_w(z)})$.
Thus $C$ has a finite equivariant triangulation.
After one further barycentric subdivision, a simplex
preserved by an element is fixed pointwise. The alternating
cellular trace is therefore the Euler number of the fixed
subcomplex. The trace identity in
\cite[proof of Theorem~2C.3, p.~181]{Hatcher2002}
identifies it with the alternating trace on homology.
The invariant gradient collar restricts to each fixed
locus, so $C^\zeta\hookrightarrow F^\zeta$ is a
homotopy equivalence. This proves the displayed character
formula on $F$.

For a free action every nonidentity character value
is zero. At the identity the value is
$\chi(P^{-1}(1))=d\chi(Q)$.
The character formula consequently gives
\[
 H_0(F;\C)-H_1(F;\C)+H_2(F;\C)=\chi(Q)\C[\mu_d]
\]
as virtual representations. The group acts trivially
on $H_0$. Each nontrivial character therefore occurs
in $H_2$ at least $\chi(Q)$ times. Finite-order
monodromy is semisimple, proving the lower bound.
\end{enumerate}
\end{proof}

\begin{proposition}[Multiplicity of contracted horizontal pole sections]\hypertarget{tag-02XR}{}\label{night:P4.52}\label{stmt:P4.52}
\StatementRecord{night:P4.52}{proved}{02XR}{\cref{lem:II:positive-ruled-conductor-graph,lem:II:weighted-initial-monodromy,night:P4.51,cdp:st:A.5.3,cdp:st:A.5.17,night:P4R.15}}
In the setting of \cref{lem:II:positive-ruled-conductor-graph},
let the horizontal pole section $H$ be contracted by the
minimal normalization resolution. Its image on $S$ has
ambient multiplicity exactly three. The generic transverse
types along the critical curve are arbitrary.
\end{proposition}
\begin{proof}
\begin{enumerate}[label=\arabic*.,leftmargin=*]
\item Use $\nu:N\to S$, $\sigma:T\to N$ and
$\rho=\nu\sigma$ from
\cref{lem:II:positive-ruled-conductor-graph}.
The section $H$ has positive genus and pole coefficient two.
Connected fibres give $H^1(B;\Q)\hookrightarrow H^1(Y;\Q)=0$,
so $B\simeq\PP^1$. The nonzero integral fibre class in the
proof of \cref{cdp:st:D.1.9}, on the underlying space
$|mS|=S$, gives $b_1(S)>0$. Thus the first assertion of
\cref{night:P4.51} applies to this same full fibre $mS$.

At the contracted image write $g=g_d+g_{d+1}+\cdots$
for a reduced local equation. The full local fibration
function is $\Phi=u g^m$ with $u$ a unit. The map in
\cref{night:P4.51} transports a local $g$-level to the
same level of $h=\exp((\log u)/m)g$, and its inclusion
into the full fibre is
$h^{-1}(\delta)\hookrightarrow\Phi^{-1}(\delta^m)$.
It proves that every local $g$-fibre has zero rational
intersection form. Apply the first assertion of
\cref{lem:II:weighted-initial-monodromy} with weights
$(1,1,1)$ and $P=g_d$. The pairing on $\{g_d=1\}$
is consequently zero.

\item Suppose that the projective quartic
$C_4=\{g_4=0\}$ is reduced and has a normalization
component of positive genus. By degree, either it is
an irreducible quartic with normalization genus $g_0\ge1$,
or it is a smooth elliptic cubic and a line. In the first
case $\sum\delta_i=3-g_0\le2$, so the sum of transverse
Milnor numbers satisfies
$M=\sum_i(2\delta_i-r_i+1)\le4$.
In the second case B\'ezout gives total intersection
length three and at least one intersection point.
Thus $\sum\delta_i=3$ and $\sum(r_i-1)\ge1$, giving
$M\le5$.

The map $g_4^{-1}(1)\to\PP^2\setminus C_4$ is the
quotient by the free scalar action of $\mu_4$.
The curve normalization and genus formula give
\[
 \chi(g_4^{-1}(1))=4\chi(\PP^2\setminus C_4)=4(7-M).
\]
The last assertion of \cref{lem:II:weighted-initial-monodromy}
therefore gives
\begin{equation}
             \operatorname{rank}(T-I)\ge3(7-M)\ge6>M.
\label{eq:II:manual:4.52c}
\end{equation}
This contradicts \eqref{eq:II:manual:4.52b}; the difference
of the lower and upper bounds is $21-4M\ge1$.

\item Blow up $Y$ at the image point, with map
$\beta:Y_1\to Y$ and exceptional divisor $E$.
Let $S_1$ be the compact strict transform, put
$\mu=\beta|_{S_1}$, and write
$\rho_1:T_1\to S_1$ for its normalization followed by
minimal resolution. The proper map $\mu\rho_1:T_1\to S$
factors through $N$. Over each singular point of $N$,
minimal-resolution domination gives its factorization
through $T$; over the smooth locus the factorization is
the given map to $N=T$. They agree on their overlap,
giving a proper holomorphic modification $T_1\to T$.
It is a composition of point blowups by
\cite[Propositions~3.10 and~3.12, p.~20]{PopescuPampu2016}.
The strict transform of $H$ therefore has the same
genus and pole coefficient two.

Apply \cref{cdp:st:A.5.3} to the compact Cartier
hypersurface $S_1\subset Y_1$, with normalization
and minimal resolution $\rho_1$. Denote its canonical
effective divisor by $D_1$. In a point-blowup chart
$(r,u_1,u_2)\mapsto(r,ru_1,ru_2)$, the ambient
volume form acquires the factor $r^2$, while the reduced
equation acquires $r^d$. The residue comparison, with
the permuted charts on the other exceptional patches, gives
\[
 \omega_{S_1}=\mu^*\omega_S\otimes
                        \OO_{S_1}((2-d)E).
\]
Comparing pullbacks of the same original local generator
gives $D_{\rm old}=D_1+(d-2)\rho_1^*E$ on $T_1$.
Set $n=\operatorname{ord}_H(\rho_1^*E)>0$ and
$a=\operatorname{coeff}_H(D_1)\ge0$. Their coefficient
identity is
\begin{equation}
                             2=(d-2)n+a.
\label{eq:II:manual:4.52d}
\end{equation}
When $H$ is exceptional for $\rho_1$, its positive genus
and \cref{cdp:st:A.5.3} give $a>0$.

Thus $d\ge5$ is impossible. For $d=4$, the identity
forces $n=1,a=0$. The curve $H$ is neither exceptional
over the normalization of $S_1$ nor a strict conductor
component. At its generic point the normalization is
an isomorphism. Its image in $S_1\cap E$ has multiplicity
one, since the local exceptional parameter has order one,
and is birational to $H$. This positive-genus component
of the quartic has degree three or four, leaving at most
one linear factor. Hence the entire quartic is reduced,
and the previous item excludes it.

\item An isolated critical point of $g$ is excluded by
\cref{cdp:st:A.5.17}, applied to the same full function
$\Phi=u g^m$. Its root $h=\exp((\log u)/m)g$ and the
identity inclusion of its local level into
$\Phi^{-1}(\delta^m)$ are precisely the maps above;
the proposition carries the square-$-2$ Morse sphere
into that torus fibre.

At a nonisolated point with $d=2$, the Hessian is
nonzero. Apply \cref{night:P4R.15} to this same reduced
equation and full local fibration. It gives the product
equation $X^2+Y^k$. For even $k$, its two analytic
branches are smooth; for odd $k$, its normalization is
the smooth parametrization
$(Y,X,Z)=(v^2,\mathrm i v^k,Z)$.
Neither minimal normalization resolution has an
exceptional curve. A point with $d=1$ is smooth.
These cases cannot receive the contracted section $H$,
so the remaining multiplicity is three.
\end{enumerate}
\end{proof}

\subsection{The cubic double line and its weighted leading fibre}\label{subsec:II:cubic-double-line}

\begin{proposition}[Normality after the multiplicity-three blowup]\hypertarget{tag-02XS}{}\label{night:P4.53}\label{stmt:P4.53}
\StatementRecord{night:P4.53}{proved}{02XS}{\cref{night:P4.52,lem:II:weighted-initial-monodromy,lem:II:conductor-duality,lem:II:plane-branch-conductor,cdp:st:A.5.3}}

In the setting of \cref{night:P4.52}, blow up the
multiplicity-three image of the contracted horizontal
section $H$, normalize the strict transform $S_1$, and
take its minimal resolution $T_1$. The strict transform
of $H$ is exceptional over a point of $S_1$. The entire
germ of $S_1$ there is normal and has an isolated simple
elliptic singularity. In particular $H$ cannot dominate
a curve in the ambient exceptional plane. The generic transverse types along the original critical
curve are arbitrary.
\end{proposition}
\begin{proof}
\begin{enumerate}[label=\arabic*.,leftmargin=*]
\item Use the comparison maps of \cref{night:P4.52}:
$\beta:Y_1\to Y$ is the point blowup, $S_1$ is the
compact strict transform, and $\rho_1:T_1\to S_1$ is
normalization followed by minimal resolution.
Apply \cref{cdp:st:A.5.3} to this same compact hypersurface
$S_1\subset Y_1$, and denote its canonical effective divisor
by $D_1$. With $E$ the ambient exceptional plane, put
$n=\operatorname{ord}_H(\rho_1^*E)>0$ and
$a=\operatorname{coeff}_H(D_1)$. The coefficient equation
\eqref{eq:II:manual:4.52d}, now with multiplicity three,
is $2=n+a$. Hence $(n,a)=(2,0)$ or $(1,1)$.

\item Suppose $H$ dominates a reduced component $L$ of the
projective cubic tangent curve. At its generic point
a transverse Weierstrass presentation has degree
$\mu_L$, the multiplicity of that component in the cubic.
Factoring over the completed valuation field of the
exceptional coordinate $r$ gives the exact degree
identity
\begin{equation}
 \sum_{H'}\operatorname{ord}_{H'}(r)
                 [\mathbb C(H'):\mathbb C(L)]=\mu_L.             
\label{eq:II:manual:4.53a}
\end{equation}
The factors group the roots by valuation branches;
the two factors in each summand are the ramification
index and residue-field degree. Residue extensions
are separable in characteristic zero.

\item A positive-genus component of a cubic is a smooth
elliptic cubic of multiplicity one. The strict
transform is smooth along it, with $n=1,a=0$,
contrary to $2=n+a$. Every other reduced component
is rational. The finite map from $H$ then has degree
$h\ge2$. Equation \eqref{eq:II:manual:4.53a} excludes $n=2$, since
$\mu_L\le3$. Thus $n=a=1$.

\item At a generic point of $L$, the $h$ points on $H$
give distinct branches of the transverse plane germ,
each with conductor exponent one. The branch formula of \cref{lem:II:plane-branch-conductor}
\[
             c_i=2\delta_i+\sum_{j\ne i}I_{ij}
\]
forces $h=2$ and exactly two smooth branches meeting
once, with no additional branch. Each has $r$-order
one, so their total Weierstrass degree is two. Hence
$\mu_L=2$, and the tangent polynomial is $y^2\ell$,
where $y=0$ is the double line and $\ell$ a different
linear form.

\item Choose complementary coordinates $x,z$. The
restriction $\ell(x,z)$ is nonzero; any $y$ term in
$\ell$ has higher weight for the weights used below.
In a generic blowup chart along the double line,
write $v$ for the transverse coordinate. The
strict-transform equation begins
\[
                 \ell v^2+r(f_4|_L)+\cdots.
\]
A nonzero $f_4|_L$ would give ramification order two
and residue degree one. Therefore $f_4$ is divisible
by $y$. The two branches meet once, so after
$v=ru$ their distinct leading slopes solve
\[
                         \ell u^2+G_3u+G_5=0,
\]
where $G_3,G_5$ are forms of degrees three and five
in $x,z$. The discriminant
$\Delta=G_3^2-4\ell G_5$ is nonzero and nonsquare
in $\mathbb C(\mathbb P^1)$. The residue-field
extension of $H\to L$ is
$\mathbb C(\mathbb P^1)(\sqrt\Delta)$.

\item Consequently the original equation has weighted
leading polynomial
\begin{equation}
          P_5=\ell(x,z)y^2+G_3(x,z)y+G_5(x,z),               
\label{eq:II:manual:4.53b}
\end{equation}
of degree five for weights $(1,2,1)$ on $(x,y,z)$,
with no smaller-weight terms. This identifies the
normalized projective component birationally with
$H$ through its function field.

\item The local reduced equation is the same $g$ as in
\cref{night:P4.52}; its full fibration function is
$\Phi=u g^m$. That proof applies the first clause of
\cref{night:P4.51}, using $b_1(S)>0$, to this full fibre.
Choose a holomorphic logarithm of $u$, and put
$g_u=\exp((\log u)/m)g$, the function denoted $h$ there,
and $\alpha=\exp((\log u(0))/m)\ne0$.
The proof of \cref{night:P4.51}, with the identity inclusion
$g_u^{-1}(\delta)\hookrightarrow\Phi^{-1}(\delta^m)$,
gives zero intersection pairing on these local $g_u$-levels.
The weighted initial polynomial of $g_u$ is $\alpha P_5$.
Apply the first clause of
\cref{lem:II:weighted-initial-monodromy} to $g_u$, with
\[
 D_\epsilon(x,y,z)=(\epsilon x,\epsilon^2y,\epsilon z),
 \qquad G_\epsilon=\epsilon^{-5}g_u\circ D_\epsilon,
 \qquad c=\alpha.
\]
The functions $G_\epsilon$ converge with all derivatives
on compact sets to $\alpha P_5$. The initial level is
$\{\alpha P_5=\alpha\}=\{P_5=1\}$, so every compact
intersection calculation there transports into
$g_u=\alpha\epsilon^5$, inside the full fibre
$\Phi=(\alpha\epsilon^5)^m$. Its pairing is zero.

\item Let $C_5=\{P_5=0\}\subset\mathbb P(1,2,1)$.
The quadratic is irreducible over $\mathbb C(x/z)$,
since the degree-two residue cover $H$ is connected.
The only possible further factor is a common
vertical factor of $\ell,G_3,G_5$. Since $\ell$
is linear, either $P_5$ is irreducible or
$P_5=\ell Q_4$, with $Q_4$ monic quadratic in $y$.
Both alternatives are reduced.

\item In the irreducible case, let $k_i$ be the
multiplicities of the distinct zeros of the
homogeneous degree-six discriminant. The normalized
double cover has genus
$(\#\{i:k_i\text{ odd}\}-2)/2$. Positive genus
therefore requires at least four odd multiplicities,
and gives
\begin{equation}
              \sum_i(k_i-1)=6-\#\{i\}\le2.                  
\label{eq:II:manual:4.53c}
\end{equation}
Away from $\ell=0$ and the weighted vertex,
completing the quadratic identifies the projective
curve germ with $w^2-s^{k_i}$, up to a unit and
holomorphic coordinates. Its transverse Milnor
number is $k_i-1$. At $\ell=0$, if $G_3\ne0$ the
finite root is simple and the other point is at
the weighted vertex. If $G_3=0$, absence of a
common factor gives $G_5\ne0$, and $\Delta$ has
a simple zero. There is no additional singularity.
In the chart $y=1$ the vertex is the quotient of
$\mathbb C^2$ by $(x,z)\mapsto(-x,-z)$; the
lifted equation $\ell+G_3+G_5$ has nonzero
differential. The original polynomial also has
nonzero derivative in the $\ell$ direction on
the $y$-axis, so that axis is not critical.
The positive-genus multiplicity partitions of the sextic are
$(1,1,1,1,1,1)$, $(2,1,1,1,1)$ and $(3,1,1,1)$;
their transverse Milnor-number sums are respectively
$0,1,2$. Thus $M\le2$.

\item In the reducible case write
$G_3=\ell A_2$, $G_5=\ell A_4$. Then
$Q_4=y^2+A_2y+A_4$ has discriminant of degree four.
A connected positive-genus double cover requires
four distinct simple roots, the partition
$(1,1,1,1)$. Its normalization is
elliptic and $Q_4$ is smooth on the weighted plane;
it avoids the weighted vertex. The rational
component $\ell=0$ meets it in total length two,
read directly from the restricted monic quadratic.
Either there are two transverse intersections,
with total Milnor number two, or one tangency of
order two, with Milnor number three. There are no
other critical rays. Hence in all cases
\begin{equation}
                              M\le3.                       
\label{eq:II:manual:4.53d}
\end{equation}

\item The weighted plane has Euler number three: its
strata $x\ne0$, $x=0,z\ne0$, and $x=z=0$ are
$\mathbb C^2$, $\mathbb C$, and a point.
In the irreducible case $e(C_5)\le e(H)\le0$.
In the reducible case
$e(C_5)=0+2-\#(Q_4\cap\{\ell=0\})\le1$.
Therefore
\begin{equation}
               \chi\bigl(\mathbb P(1,2,1)\setminus C_5\bigr)
                         \ge2.                            
\label{eq:II:manual:4.53e}
\end{equation}

\item The weighted action of $\mu_5$ on $F=\{P_5=1\}$ is
free, and its quotient is $\mathbb P(1,2,1)\setminus C_5$.
Indeed the only possible nontrivial stabilizer of the
weighted $\C^*$ action lies at the weighted vertex,
which belongs to $C_5$. Thus the quotient map is an
unramified fivefold cover. The third clause of
\cref{lem:II:weighted-initial-monodromy}, with
$w=(1,2,1)$ and $d=5$, identifies its weighted scalar
monodromy with local Milnor monodromy and gives
\[
 \operatorname{rank}(T-I)
 \ge4\chi\bigl(\mathbb P(1,2,1)\setminus C_5\bigr)
 \ge8.
\]
Here one may use the invariant proper semialgebraic norm
$R=|y|^2+(|x|^2+|z|^2)^2$: it satisfies
$R(D_\epsilon z)=|\epsilon|^4R(z)$, and the normalized
gradient construction of that lemma applies to it.
The two resulting compact cores both include into $F$
by homotopy equivalences; these inclusions identify their
homology, monodromy and compact intersection calculations.
The polynomial $P_5$ is reduced by the factorization above.
Its zero pairing therefore satisfies the second clause
of the lemma, with the transverse sum just computed:
$\operatorname{rank}(T-I)\le M\le3$.
This also includes an isolated critical point, for which
$M=0$. The bounds differ by at least $8-3=5$, excluding
every dominating $H$.

\item Thus $H$ is exceptional over the normalization of $S_1$.
The positive-genus refinement in \cref{cdp:st:A.5.3}
gives strictly positive coefficients on its entire
connected exceptional block. In particular $a>0$, so
$n=a=1$. Adjunction on this smooth positive-genus curve
gives
\[
                    0=2g(H)-2+(D_1-H)\cdot H.
\]
Every term on the right is nonnegative. Hence $g(H)=1$
and $H$ meets no other component of $D_1$. The whole
exceptional block has positive coefficients and is connected,
so it consists of $H$ alone and is disjoint from the
strict conductor.

Let $q$ be the normalization point below $H$ and
$p_1$ its image in $S_1$. A conductor curve through $q$
would have strict transform meeting the proper exceptional
fibre $H$. Thus the conductor has no codimension-one
support near $q$. It is reflexive by
\cref{lem:II:conductor-duality}, so its stalk at $q$ is
the unit ideal. Apply that lemma to
$A=\OO_{S_1,p_1}$, its finite semilocal normalization $B$,
and this same $q$: it gives $B=A$ on the entire local
germ. There are consequently no additional local branches.
The strict-transform germ is normal and its minimal
resolution has the single smooth elliptic exceptional
curve $H$. This is an isolated simple elliptic singularity.

\end{enumerate}
\end{proof}

\subsection{Exclusion of the contracted horizontal section}\label{subsec:II:elliptic-contraction}

\begin{proposition}[Noncontraction of the horizontal pole section]\hypertarget{tag-02XT}{}\label{night:P4.54}\label{stmt:P4.54}
\StatementRecord{night:P4.54}{proved}{02XT}{\cref{night:P4.53,night:P4R.18,lem:II:weighted-initial-monodromy,cdp:st:A.5.17}; \cite[Theorem~1.4.1, pp.~39--41]{Dolgachev1982}}

In the setting of \cref{lem:II:positive-ruled-conductor-graph},
the horizontal pole section is not contracted by the
normalization resolution. It consequently maps with
degree two or three onto a rational normalized conductor
component whose generic transverse germ is $D_4$.
The transverse types on the other conductor component are arbitrary.
\end{proposition}
\begin{proof}

\begin{enumerate}[label=\arabic*.,leftmargin=*]
\item By \cref{night:P4.53}, a contracted section would be the sole smooth
elliptic exceptional curve $H$ of an isolated normal
hypersurface singularity on the first ambient strict
transform. Its self-intersection is $-d<0$, and the
exceptional coordinate $r$ has order one along
$H$. Write $\pi:\widetilde V\to(V,0)$ for this local
resolution. 

\paragraph{Analytic formal functions and the associated graded ring.}

\item Put \(L=\mathcal O_H(-H)\), of degree \(d\), and
\[
 I_m=(\pi_*\mathcal O_{\widetilde V}(-mH))_0
     =\{a\in\mathcal O_{V,0}:\operatorname{ord}_H(a)\ge m\}.
\]
For \(m\ge1\), the restrictions of
\(\mathcal O(-mH)\) to successive infinitesimal neighbourhoods
of \(H\) have graded factors \(L^{m+j}\), \(j\ge0\).
Each has zero first cohomology, by Serre duality on the elliptic
curve and its strictly positive degree. Their extensions also
have zero first cohomology.

\item Grauert's analytic formal-functions theorem and coherence theorem are
\cite[\S6, Haupts\"atze I and IIa, p.~59]{Grauert1960D}, with the proof
corrections in \cite[pp.~35--36]{Grauert1963ErratumD}. The maximal-ideal pullback and the
ideal of \(H\) have the same radical. On a finite cover of
the compact exceptional curve, Noetherianity gives one uniform
power of the latter contained in the former. Their powers are
therefore cofinal. Applying the theorem with
\(\mathcal O(-mH)\) gives a zero completion for
\(R^1\pi_*\mathcal O(-mH)\). This coherent module is
supported at the origin, so it is zero. Consequently
\begin{equation}
 I_m/I_{m+1}\simeq H^0(H,L^m)\qquad(m\ge1),             
\label{eq:II:manual:4.54w.2}
\end{equation}
through the restriction maps in
\[
 0\longrightarrow\mathcal O(-(m+1)H)
 \longrightarrow\mathcal O(-mH)
 \longrightarrow L^m\longrightarrow0.
\]
These maps respect multiplication. Since normality gives
\(\pi_*\mathcal O_{\widetilde V}=\mathcal O_V\), and a
function vanishes at the contracted point exactly when its
restriction to \(H\) vanishes, \(I_1=\mathfrak m\).
The degree-zero piece is \(\mathbb C\). Thus
\begin{equation}
 \operatorname{gr}_{I}\mathcal O_{V,0}
   =\bigoplus_{m\ge0}H^0(H,L^m)                         
\label{eq:II:manual:4.54w.3}
\end{equation}
as graded algebras.

\item The same formal-functions theorem in degree zero identifies
\begin{equation}
 \widehat{\mathcal O}_{V,0}
      \simeq\varprojlim_m H^0(mH,\mathcal O_{mH})
      \simeq\varprojlim_m\mathcal O_{V,0}/I_m.          
\label{eq:II:manual:4.54w.4}
\end{equation}
The second equality uses the just-proved first-direct-image
vanishing. This supplies the completeness and separation of the valuation filtration
needed in the leading-equation argument below.

\item Because \(\mathfrak m^2\subset I_2\), \eqref{eq:II:manual:4.54w.2} gives a
surjection
\(\mathfrak m/\mathfrak m^2\to H^0(H,L)\).
The source has dimension three and the target dimension \(d\).
Hence \(d\le3\).

\paragraph{The three elliptic section rings.}

\item For \(d=3\), Riemann--Roch shows that \(L\) separates
points and tangent vectors: subtracting any length-two
subscheme leaves a degree-one line bundle with one section.
The complete linear series therefore embeds \(H\) as a smooth
plane cubic. The restriction sequence on \(\mathbb P^2\),
using \cite[Theorem~1.4.1(ii), p.~39]{Dolgachev1982}
with weights $(1,1,1)$, gives its whole
section ring as
\[
 \mathbb C[x,y,z]/(F_3),\qquad \deg(x,y,z)=(1,1,1),
\]
where \(F_3=0\) is that smooth cubic. For example,
$X^3+Y^3+Z^3=0$ is smooth in projective space because
its three gradient equations force $X=Y=Z=0$.

\item For \(d=2\), the complete pencil is base-point free and gives
a degree-two map \(p:H\to\mathbb P^1\), with
\(L=p^*\mathcal O(1)\). Trace splits
\(p_*\mathcal O_H=\mathcal O\oplus B^{-1}\).
Its Euler characteristic is zero, so \(\deg B=2\).
The algebra structure is therefore a double-cover equation
\(z^2=F_4(x,y)\), where the binary quartic has four distinct
zeros because the double cover is smooth. Projection formula
gives the entire section ring with generator weights
\((1,1,2)\) and that single degree-four relation.

\item For \(d=1\), write \(L=\mathcal O_H(p)\).
Riemann--Roch gives a degree-one section \(x\), one new
degree-two section \(y\), and one new degree-three section
\(z\). Dividing by powers of \(x\) gives functions of pole
orders two and three at \(p\). Their Weierstrass relation
has weighted degree six. Equivalently, after graded changes
only for describing this section ring, it is
\[
 z^2=y^3+A x^4y+B x^6,
 \qquad 4A^3+27B^2\ne0,
 \qquad\deg(x,y,z)=(1,2,3).
\]
The pole-order count shows that the monomials with at most one
factor \(z\), reduced by this relation, give every space
\(H^0(H,L^m)\). Thus this is the whole section ring. Its weighted projective curve is the smooth $H$.
The displayed Weierstrass generators describe this graded ring;
the analytic coordinates used below are constructed while fixing $r$.

\paragraph{Coordinates retaining the given exceptional parameter.}

\item Start with analytic ambient coordinates \((r,u,v)\), retaining
the given \(r\). Their images form a basis of
\(\mathfrak m/\mathfrak m^2\). The map to
\(I_1/I_2=H^0(L)\) has the rank just computed.

\item When \(d=3\), it is an isomorphism, so any linear complement
to \(r\) supplies the other two valuation-one coordinates.

\item When \(d=2\), choose a linear coordinate \(s\) whose
leading section complements that of \(r\), and choose the
remaining independent linear coordinate \(v_0\) in the
kernel of the degree-one restriction. Thus
\(\operatorname{ord}_H(v_0)\ge2\). Its degree-two leading
section must be independent of \(r^2,rs,s^2\). Otherwise
all degree-two initial sections of analytic functions would
lie in that three-dimensional span: a Taylor expansion shows
that every term of ordinary order at least three has valuation
at least three, and the linear terms of valuation two are
accounted for by \(v_0\). This contradicts
\(h^0(L^2)=4\). Therefore \(v_0\) has valuation exactly
two and is the required indecomposable degree-two generator.
The coordinates \((r,s,v_0)\) were obtained by an invertible
linear change.

\item When \(d=1\), first choose two independent linear coordinates
\(u_0,v_0\) in the kernel of the degree-one restriction.
At least one linear combination, denoted \(s\), has a
degree-two initial section independent of \(r^2\), since
\(h^0(L^2)=2\). Choose the remaining independent linear
coordinate \(v_1\), and subtract a scalar multiple of \(s\)
so that its degree-two initial term is \(\alpha r^2\).
Set
\[
                         v_2=v_1-\alpha r^2.
\]
This is a triangular analytic coordinate change, and
\(\operatorname{ord}_H(v_2)\ge3\). Its degree-three
initial section must be independent of \(r^3,rs\).
Indeed, if it were absent or in that span, any Taylor series
of valuation three would have degree-three initial section
in that same span: the possible terms are precisely the
linear \(v_2\) term and the weighted-degree-three monomials
\(r^3,rs\); its valuation-one and valuation-two terms
must first vanish. This contradicts \(h^0(L^3)=3\).
Thus \(v_2\) has valuation exactly three and supplies the
indecomposable degree-three generator. The three ambient
coordinates remain analytically invertible, because only
an invertible linear change and the displayed quadratic
translation were performed.

\paragraph{The weighted leading equation.}

\item Use the coordinates $(r,s,v)$ just constructed, with
their computed valuations. Their initial sections generate
the graded ring. Let $F_D$ be its single graded relation,
expressed in these initial generators. The analytic function
\(F_D(r,s,v)\) has valuation at least \(D+1\).
At each successive valuation, \eqref{eq:II:manual:4.54w.3} writes its initial
section as a weighted homogeneous polynomial in the three
generators. Subtract that polynomial. This produces a formal
series \(G=F_D+F_{D+1}+\cdots\) whose image in every
\(\mathcal O_{V,0}/I_m\) is zero. Since all generator
weights are at most three, the correction polynomials tend
to zero in ordinary power-series degree; their sum is a
well-defined formal power series. By \eqref{eq:II:manual:4.54w.4}, its image in
\(\widehat{\mathcal O}_{V,0}\) is zero.

\item The completed hypersurface ideal is principal, generated by
the analytic equation \(f\). Thus \(G=h f\) for
a formal series \(h\). Taking weighted initial forms gives
\(F_D=\operatorname{in}(h)\operatorname{in}(f)\).
The polynomial \(F_D\) is irreducible because its weighted
projective curve is the smooth connected elliptic curve
\(H\). Also \(f\) is nonconstant and vanishes at the
origin. Hence \(\operatorname{in}(h)\) is a nonzero
constant, and \(\operatorname{in}(f)\) is a scalar multiple
of \(F_D\). Thus the chosen analytic coordinates, with the fixed
exceptional parameter $r$, have weights and leading degrees
\begin{equation}
 (1,1,1;3),\qquad(1,1,2;4),\qquad(1,2,3;6).
\label{eq:II:manual:4.54w.1}
\end{equation}
All remaining terms of the analytic equation have strictly
larger weighted degree.

\paragraph{Descent to the original blowdown coordinates.}

\item In a blowup chart, first centre $(u,v)$ at the point
below $H$. The corresponding original coordinate change is
$(x,y,z)\mapsto(x-u(p)z,y-v(p)z,z)$; its inverse adds
the same two linear terms. Thus $r=z$ is fixed.
In the centred chart $x=ru$, $y=rv$, every change above has the form
\[
 s=a u+b v+c r,\qquad
 v'=a'u+b'v+c'r-\alpha r^2,
 \qquad ab'-a'b\ne0.
\]
It descends to the holomorphic coordinate map
\begin{equation}
 \begin{aligned}
 x'&=r s=a x+b y+c r^2,\\
 y'&=r v'=a'x+b'y+c'r^2-\alpha r^3,\\
 r'&=r.
 \end{aligned}                                        
\label{eq:II:manual:4.54w.5}
\end{equation}
Its Jacobian determinant at the origin is \(ab'-a'b\),
so the inverse-function theorem supplies an analytic inverse.
This is the required morphism preserving the original
parameter and blowdown structure.

\paragraph{Leading polynomials in the original coordinates.}
Apply this calculation to the simple elliptic point supplied
by \cref{night:P4.53}. The original equation is $r^3f_1$, and its original
coordinates are $(x,y,z)=(rs,rv,r)$. The three pairs of
original weights and degrees are
\[
 (2,2,1;6),\qquad (2,3,1;7),\qquad (3,4,1;9).
\]
For a strict-transform monomial $r^as^bv^c$, multiplication
by $r^3$ gives $x^by^cz^{3+a-b-c}$. Its original weighted
degree is three plus its strict-transform weighted degree.
Analyticity of the original equation makes every
negative-power coefficient zero. Thus the weighted
leading term is a polynomial of the indicated degree.
On $z\ne0$, the ratios
\[
     \left(x/z^{1+w_s},\ y/z^{1+w_v}\right)
\]
identify its projective component with the strict-transform
elliptic relation on $r\ne0$. This is a birational map
of the residue-field curves; their smooth
normalizations are the same elliptic curve $H$.

\item Write $g$ for the original reduced support equation
and $\Phi=u g^m$ for the full local fibration function.
As in \cref{night:P4.52}, $b_1(S)>0$ gives the first
clause of \cref{night:P4.51} on this same full fibre.
Choose a holomorphic logarithm of $u$ and put
$g_u=\exp((\log u)/m)g$ and
$\alpha=\exp((\log u(0))/m)\ne0$.
The identity inclusion
$g_u^{-1}(\delta)\hookrightarrow\Phi^{-1}(\delta^m)$
is the full-fibre map in \cref{night:P4.51}; its proof
gives zero intersection pairing on these local $g_u$-levels.
For each pair $(w;D)$ just obtained, the weighted initial
polynomial of $g_u$ is $\alpha P$. Apply
\cref{lem:II:weighted-initial-monodromy} to $g_u$, with
$D_\epsilon(x,y,z)=(\epsilon^{w_x}x,
\epsilon^{w_y}y,\epsilon^{w_z}z)$,
$G_\epsilon=\epsilon^{-D}g_u\circ D_\epsilon$ and
$c=\alpha$. Thus the initial level is
$\{\alpha P=\alpha\}=\{P=1\}$, and every compact
pairing there transports into
$g_u=\alpha\epsilon^D$, inside the full fibre
$\Phi=(\alpha\epsilon^D)^m$. Its pairing is zero.
Once its leading polynomial is shown reduced, the second
clause of that lemma gives
$\operatorname{rank}(T-I)\le M$, with $M=0$ for an
isolated critical point.

\paragraph{Degree six.}
The leading polynomial is exactly
\[
                    P_6(x,y,z)=G_3(x,y,z^2),
\]
where $G_3(X,Y,Z)$ defines a smooth projective cubic.
The map $(x,y,z)\mapsto(x,y,z^2)$ gives a finite double
cover from $F=\{P_6=1\}$ to $\{G_3=1\}$, branched over
\[
                       B=\{G_3(x,y,0)=1\}.
\]
The homogeneous cubic fibre has Euler number nine: it
is a threefold unramified cover of the complement of
an elliptic cubic in $\mathbb P^2$, whose Euler number
is three. Let $k_0\in\{1,2,3\}$ be the number of distinct
zeros of the nonzero binary cubic $G_3(x,y,0)$ on
$\mathbb P^1$. The curve $B$ is a threefold cover of
that punctured projective line. Hence
\[
          \chi(B)=3(2-k_0),\qquad \chi(F)=18-\chi(B)=12+3k_0.
\]

\item Every nonzero critical orbit of $P_6$ lies on $z=0$
at a tangency of the cubic with $Z=0$. If the
intersection orders are $e_i$, then $\sum e_i=3$.
At a tangency, smoothness of $G_3$ makes its
$Z$-derivative a unit, and the exact transverse
germ is $u^{e_i}+z^2$, after holomorphic coordinates
and a unit. Consequently
\[
                       M=\sum(e_i-1)=3-k_0\le2.
\]
This also shows reducedness of $P_6$: a repeated
factor would give a two-dimensional critical set,
whereas the derivative calculation gives only
these critical orbits.

\item Weighted monodromy has order six. Its only
nonidentity element with a fixed set is its cube,
which negates $z$ and fixes $x,y$; its fixed set is
$B$. The other four nonidentity elements have no
fixed point on $F$. The fixed-locus character formula of
\cref{lem:II:weighted-initial-monodromy} applies to
this same polynomial, with weights $(2,2,1)$ and degree
six. Thus the character of $H_0-H_1+H_2$ is
$\chi(F)$ at the identity, $\chi(B)$ at the cube,
and zero elsewhere. Its
trivial-character multiplicity is
\[
                       \frac{\chi(F)+\chi(B)}6=3.
\]
Since $H_0$ has only the trivial character, the
nontrivial-character dimension in $H_2$ is at least
$\chi(F)-3=9+3k_0\ge12$.

The variation inequality
\eqref{eq:II:manual:4.52b} contradicts the lower bound twelve for
the nontrivial monodromy rank; the gap is at least $12-2=10$.

\paragraph{Weighted curve cohomology.}
For
$\mathbb P(a,b,c)=\operatorname{Proj}\mathbb C[u,v,w]$
and a reduced equation of weighted degree $e$,
the exact graded sequence sheafifies to
\[
 0\longrightarrow\mathcal O_{\mathbb P}(-e)
 \longrightarrow\mathcal O_{\mathbb P}
 \longrightarrow\mathcal O_D\longrightarrow0.
\]
Apply \cite[Theorem~1.4.1(i)--(iii), pp.~39--41]{Dolgachev1982}
with $S=\C[u,v,w]$, weights $(a,b,c)$ and twist $-e$.
It gives $H^1(\mathbb P,\mathcal O(k))=0$ and identifies
the dimension of $H^2(\mathbb P,\mathcal O(-e))$ with
that of $S_{e-a-b-c}$. The long exact sequence therefore gives
\begin{equation}
 h^1(\mathcal O_D)=
 \#\{\hbox{nonnegative weighted monomials of degree }
                         e-(a+b+c)\},
 \qquad h^0(\mathcal O_D)=1.                         
\label{eq:II:manual:4.54a}
\end{equation}
In particular the arithmetic genus is zero below
degree $a+b+c$, and is one in degrees $a+b+c$
and $a+b+c+1$ when there is exactly one variable
of weight one.

\item At a cyclic quotient vertex of prime order two
or three, with no invariant linear ambient
coordinate, a smooth coarse curve has a smooth
plane pullback. To prove this, its normalized
branches form one orbit. A single singular
branch of multiplicity $m\ge2$ has a faithful
normalization action of order $p$; every invariant
ambient function has order at least $2m>p$,
whereas a regular quotient requires a function
of normalization order $p$. For $p$ transitively
permuted branches the stabilizer is trivial,
and every invariant function has order at least
two, whereas a regular quotient requires order
one. Both alternatives are impossible. The
remaining single branch is smooth, and its
reduced equation has the character of a nonzero
linear ambient coordinate. These arguments
use the analytic invariant rings:
averaging makes the invariants functor exact,
so every quotient-curve function lifts to an
invariant ambient function.

\paragraph{Degree seven.}
For degree seven the coordinate permutation
$(x,y,z)\mapsto(z,x,y)$ identifies the weighted plane
with $\operatorname{Proj}\C[z,x,y]=\mathbb P(1,2,3)$;
its inverse restores $(x,y,z)$.
Formula \eqref{eq:II:manual:4.54a} gives arithmetic genus zero for
curves of degree at most five and arithmetic
genus one in degrees six and seven. Thus a
positive-genus component has degree six or
seven, genus one, and multiplicity one. There
is at most one further degree-one factor;
the polynomial is reduced.

\item In degree seven the whole curve $D$ is
irreducible with arithmetic and normalization
genus one. Its total delta invariant is zero,
so it is a smooth coarse elliptic curve.
At the quotient vertices of orders two and
three, the preceding argument makes its plane
pullback smooth. Elsewhere the ambient plane
is smooth. Thus the affine cone has no
nonzero critical orbit and $M=0$.

\item In degree six write $D=E\cup\{z=0\}$, with
$z$ the weight-one coordinate. The curve $E$
is smooth coarse elliptic. It cannot contain
either quotient vertex: its degree-six
equation has trivial character there, while
both local-coordinate characters are
nontrivial, so a smooth lift would require
an impossible invariant linear equation.
The genus formula for $D$ gives total delta
one. Its two smooth components therefore
meet once transversely in the smooth part
of the ambient weighted plane. There is
one ordinary-double critical orbit and
$M=1$.

\item The weighted plane has Euler number three,
from its affine-cell stratification. In the
two cases $e(D)=0$ or one, so its complement
has Euler number three or two. Both quotient
vertices belong to $D$, since degree seven
is divisible by neither two nor three.
The weighted $\mathbb C^*$ action is free
off $D$, so $\{P_7=1\}$ is an unramified
sevenfold cover of that complement. The
free order-seven monodromy gives nontrivial
$H_2$ rank at least eighteen or twelve,
respectively. This contradicts \eqref{eq:II:manual:4.52b},
whose upper bound is zero or one. The respective gaps are
$18-0=18$ and $12-1=11$.

\paragraph{Degree nine.}
For degree nine the same permutation identifies the plane
with $\operatorname{Proj}\C[z,x,y]=\mathbb P(1,3,4)$.
Formula \eqref{eq:II:manual:4.54a} gives
arithmetic genus zero below degree eight
and one in degrees eight and nine.
A positive-genus component must have one
of these two degrees, and multiplicity one.

\item An elliptic degree-eight component is
impossible. It necessarily contains the
weight-three vertex, since eight is not
divisible by three. The local plane
characters are $1,1$ modulo three, whereas
its equation has character two. The
prime-order quotient argument would give
a smooth lift, whose equation must have
character one. This is a contradiction.

\item Thus the entire degree-nine curve $D$ is
irreducible and smooth coarse elliptic.
It avoids the weight-three vertex: its
equation there has character zero,
incompatible with a smooth plane lift.
At the weight-four vertex the plane
characters are $1,3$ and its equation
has character one. We check the lift
without assuming that four is prime.
The transitive orbit of normalized
branches has size one, two, or four.

\item For four branches the stabilizer is
trivial and every invariant maximal-ideal
function has order at least two on a
branch, contradicting regularity of the
quotient. For two branches the stabilizer
has order two. A singular branch gives
invariant orders at least four, exceeding
two. Two smooth branches would have a
quadratic initial equation, but the
quadratic characters of the ambient
action are only zero and two, never one.
For one branch the normalization action
is faithful of order four. Both ambient
coordinate orders are odd, since their
characters are one and three. A singular
branch therefore has multiplicity at
least three, and all invariant functions
have order at least six, exceeding four.
These contradictions leave a single
smooth branch.

\item The degree-nine cone is therefore smooth
away from its vertex. It is singular
at the vertex, because no variable has
weight nine. Its hypersurface critical point is isolated. Apply the
local clause of \cref{cdp:st:A.5.17} with $H=P_9$:
the resulting compact sphere has square $-2$ in a small
local $P_9$-level. The weighted scaling isomorphism
$D_\lambda:P_9^{-1}(1)\to P_9^{-1}(\lambda^9)$
places its inverse image in $P_9=1$ for a chosen ninth
root $\lambda$ of that level value. The preceding
$G_\epsilon$ transport then puts the same compact
intersection class in $g_u=\alpha\epsilon^9$, inside
$\Phi=(\alpha\epsilon^9)^m$. Its zero pairing gives
the contradiction.

\item The section is therefore noncontracted. Apply
\cref{night:P4R.18} to the identical $f:Y\to B$, full
fibre $mS$, normalization resolution and section $H$.
Its global irreducible-or-normal-crossings, positive-Euler
and positive ruling-genus hypotheses are those of
\cref{lem:II:positive-ruled-conductor-graph} in the statement.
It gives the degree-two or degree-three map onto a rational
normalized conductor component with generic germ $D_4$.

\end{enumerate}
\end{proof}
\subsection{Cubic normal tensors along a critical branch}\label{subsec:II:cubic-normal-tensors}

\begin{proposition}[Cubic normal tensors along a $D_4$ branch]\hypertarget{tag-02YC}{}\label{night:P4R.19}\label{stmt:P4R.19}
\StatementRecord{night:P4R.19}{proved}{02YC}{\cref{lem:II:positive-ruled-conductor-graph,lem:II:monomial-curve-matching-sphere,lem:II:weighted-initial-monodromy,night:P4.51}; \cite[Chapter~II, \S2.A, Theorem~(2.1), p.~80]{Demailly2012E}; \cite[Theorem~3.35, p.~245]{Hatcher2002}}

In the setting of \cref{lem:II:positive-ruled-conductor-graph},
let $f$ be the reduced local equation of $S$, and let
$C_0$ be a smooth local branch of its critical curve
with generic transverse type $D_4$.
At a point where the cubic normal tensor along $C_0$ is
nonzero, that tensor is not the cube of a linear form.
The generic transverse types on the other critical-curve
branches at that point are arbitrary.
\end{proposition}
\begin{proof}

\begin{enumerate}[label=\arabic*.,leftmargin=*]
\item Choose coordinates $(x,y,z)$ with $C_0=(x=y=0)$.
Generic normal multiplicity three gives
$f\in(x,y)^{(3)}$. Powers of this smooth
complete-intersection ideal are primary, so
$f\in(x,y)^3$ at the specified point. Equivalently, each normal
Taylor coefficient of degree at most two is holomorphic in $z$
and vanishes on the generic axis, hence vanishes identically.

\item Suppose its nonzero cubic normal tensor is a cube.
A linear normal-coordinate change makes it $x^3$, up
to a nonzero scalar. Weierstrass preparation in $x$
\cite[Chapter~II, \S2.A, Theorem~(2.1), p.~80]{Demailly2012E}
gives the exact equality
\[
 f=U(x,y,z)\bigl[x^3+a_2(y,z)x^2+a_1(y,z)x+a_0(y,z)\bigr],
\]
where $U$ is a holomorphic unit. The bracket belongs to
$(x,y)^3$, so $a_2,a_1,a_0$ are divisible by $y,y^2,y^3$,
respectively. Depressing the monic cubic by
$x_{\mathrm{old}}=x_{\mathrm{new}}-a_2(y,z)/3$, with $y,z$
fixed, is a holomorphic coordinate change preserving $(x,y)$.
Its inverse is $x_{\mathrm{new}}=x_{\mathrm{old}}+a_2(y,z)/3$,
with $y,z$ fixed. Its Jacobian determinant is one, and
the two lower coefficients are $a_1-a_2^2/3$ and
$a_0-a_1a_2/3+2a_2^3/27$.
It yields
\begin{equation}
 f=U(x,y,z)\bigl[x^3+y^2B(y,z)x+y^3C(y,z)\bigr],
 \qquad B(0,0)=C(0,0)=0.                         \label{eq:normal-cubic-preparation}
\end{equation}

\item Set
\[
 k=\operatorname{ord}_z B(0,z),\qquad
 l=\operatorname{ord}_z C(0,z),\qquad
 \rho=\min(k/2,l/3),
\]
where the order of an identically zero series is infinity.
At least one order is finite: otherwise all generic normal
cubics would be $x^3$, contrary to generic $D_4$.
Thus $0<\rho<\infty$. Choose $L>3\rho$ and positive
rational weights
\[
       \operatorname{wt}(x)=L+\rho,\qquad
       \operatorname{wt}(y)=L,\qquad
       \operatorname{wt}(z)=1.
\]
Multiplication by a common denominator gives positive
integer weights. The weighted degree of the first face
is $3L+3\rho$. Terms of $B$ containing $y$ have weight
gap at least $L-2\rho>0$, and those of $C$ containing
$y$ have gap at least $L-3\rho>0$. Consequently the
leading polynomial, apart from the nonzero factor $U(0)$,
is exactly one of
\begin{equation}
\begin{array}{ll}
 P=x(x^2+a y^2z^k),\quad a\ne0,
       & k/2<l/3,\\[2pt]
 P=x^3+b y^3z^l,\quad b\ne0,
       & l/3<k/2,\\[2pt]
 P=Q(x,yz^r),\quad
 Q(X,V)=X^3+aXV^2+bV^3,\quad a,b\ne0,
       & k=2r, l=3r, r\ge1.
\end{array}                                                   \label{eq:normal-cubic-weighted-models}
\end{equation}
In the last case $r$ is an integer. The depressed cubic
$Q$ is not a cube: a cube with $X^3$ coefficient one and
$X^2V$ coefficient zero would have $a=b=0$. The last case includes cubics $Q$ with a double factor.

\paragraph{The compact class and its finite pullbacks.}
Each curve below is a smooth connected embedded curve
$B\subset\C^2_{x,h}$, with the second coordinate $h$
having simple zeros. Choose two of these zeros and an
embedded joining arc whose interior avoids the others.
Put
\[
 Z_h=\{(b,u,v)\in B\times\C^2:uv=h(b)\},\qquad
 X_B=\{(x,u,v):(x,uv)\in B\}.
\]
The holomorphic maps
\[
 (b,u,v)\longmapsto(x(b),u,v),\qquad
 (x,u,v)\longmapsto((x,uv),u,v)
\]
are inverse. Thus \cref{lem:II:monomial-curve-matching-sphere},
with $b'=c'=1$, gives the same smoothly embedded sphere
on $Z_h$, with self-intersection $-2$.
For the compactly supported Poincar\'e dual $\alpha$,
\begin{equation}
                         \int_{Z_h}\alpha^2=-2.             \label{eq:conic-matching-sphere-square}
\end{equation}

If $\pi:W\to Z_h$ is a finite proper holomorphic map
of degree $d>0$ between smooth complex surfaces, then
$\pi^*\alpha$ is compactly supported and
\begin{equation}
              \int_W(\pi^*\alpha)^2
                       =d\int_{Z_h}\alpha^2=-2d.            \label{eq:finite-map-compact-square}
\end{equation}
The equality follows from the proper oriented degree;
off the analytic branch locus the map has $d$ positive
sheets, and the branch locus has measure zero for
integration of compactly supported top forms.
Poincar\'e duality with compact supports
\cite[Theorem~3.35, p.~245]{Hatcher2002} gives a compact
homology two-class with that square.

\paragraph{The finite maps.}
In the first case of \eqref{eq:normal-cubic-weighted-models}, $P=1$ forces $x\ne0$
and gives
\[
           y^2z^k=h(x),\qquad
           h(x)=\frac{1-x^3}{ax},\qquad x\in\mathbb C^*.
\]
Here $B$ is the graph $\{(x,h(x)):x\in\C^*\}$.
The map
\[
 (x,y,z)\longmapsto(x,u=y^2,v=z^k)
\]
is finite proper of degree $2k$ onto $uv=h(x)$.
Indeed its source algebra is obtained from the target by
adjoining the monic equations $y^2-u=0$ and $z^k-v=0$;
complex roots give surjectivity.
The target is smooth: at $u=v=0$ the zeros of $h$
are simple, and elsewhere a derivative in $u$ or
$v$ is nonzero. The source is smooth by the weighted
Euler identity on $P=1$. Join two of the three cube
roots of one by an arc in $\mathbb C^*$ avoiding the
third. Equations~\eqref{eq:conic-matching-sphere-square} and~\eqref{eq:finite-map-compact-square} give a compact
two-class of square $-4k$ on the leading fibre.

In the second case the exact finite proper map is
\[
 \begin{split}
 \{x^3+b y^3z^l=1\}&\longrightarrow
                 \{uv=(1-x^3)/b\},\\
 (x,y,z)&\longmapsto(x,u=y^3,v=z^l).
 \end{split}
\]
It has degree $3l$; here $B$ is the graph
$\{(x,(1-x^3)/b):x\in\C\}$.
The two adjoining monic equations are $y^3-u=0$ and
$z^l-v=0$, and complex roots again give surjectivity.
The same three simple zeros supply a pulled-back
compact class of square $-6l$.

For the balanced case let
\[
                 B_Q=\{(X,V)\in\mathbb C^2:Q(X,V)=1\}.
\]
It is smooth by the ordinary Euler identity and is
connected even when $Q$ has a double factor. Its
projection to $\mathbb P^1$ minus the distinct zero
directions of $Q$ is an unramified cyclic cover of
degree three. Around a zero direction of multiplicity
one or two its monodromy generates $\mathbb Z/3$.
Such a direction exists because $Q$ is not a cube,
so the cover is connected.

The function $V$ on $B_Q$ has precisely the three
zeros $(\zeta,0)$ with $\zeta^3=1$, and they are
simple because $Q_X(\zeta,0)=3\zeta^2\ne0$.
The surface
\[
             Z_Q=\{(b,u,w):b\in B_Q,\ uw=V(b)\}
\]
is smooth. The exact map
\begin{equation}
 \begin{split}
       \{Q(x,yz^r)=1\}&\longrightarrow Z_Q,\\
       (x,y,z)&\longmapsto
                    \bigl(b=(x,yz^r),u=y,w=z^r\bigr)
 \end{split}                                                   \label{eq:balanced-cubic-finite-map}
\end{equation}
is finite proper of degree $r$: $z$ satisfies the
monic equation $z^r=w$, while $x=X(b)$ and $y=u$. An arc in $B_Q$ joining two $V$-zeros and
avoiding the third supplies a compact class of square
$-2r$ on the source. For
$Q=(X-\lambda V)^2(X+2\lambda V)$, $\lambda\ne0$,
this same degree-$r$ map gives the square $-2r$.

\paragraph{Transport to the full fibre.}
The global data are those of
\cref{lem:II:positive-ruled-conductor-graph}, which gives
$\pi_1(S)\simeq\mathbb Z$ and hence $b_1(S)>0$.
Thus the first assertion of \cref{night:P4.51} applies
to this same full fibre $mS$. Write its local function as
$\Phi=Vf^m$, choose a holomorphic logarithm of the unit
$V$, and put
\[
 f_V=\exp((\log V)/m)f,\qquad
 \alpha_0=\exp((\log V(0))/m),\qquad c=\alpha_0 U(0).
\]
The identity inclusion
$f_V^{-1}(\delta)\hookrightarrow\Phi^{-1}(\delta^m)$
is the full-fibre map in that proof, so the local
$f_V$-levels have zero intersection pairing.

Let $w=(w_x,w_y,w_z)$ be the positive integer weights
obtained above and let $d$ be their degree. The weighted
initial polynomial of $f_V$ is $cP$, including both
the original Weierstrass unit and the root of the
full-function unit. Apply the first clause of
\cref{lem:II:weighted-initial-monodromy} to $f_V$ with
\[
 D_\epsilon(x,y,z)=
 (\epsilon^{w_x}x,\epsilon^{w_y}y,\epsilon^{w_z}z),
 \qquad G_\epsilon=\epsilon^{-d}f_V\circ D_\epsilon.
\]
Its initial level is $\{cP=c\}=\{P=1\}$. The lemma
preserves each compact pairing under transport to
$f_V=c\epsilon^d$ in the original coordinate ball,
and then the identity inclusion places it in
$\Phi=(c\epsilon^d)^m$. The finite-map classes above
therefore have their same nonzero squares $-4k$, $-6l$
or $-2r$ on a local $f_V$-level, contradicting its zero
pairing. This proves the assertion, including the
balanced model with a repeated factor.

\end{enumerate}
\end{proof}

\begin{remark}[Affine fibres with two colliding spectral directions]
\hypertarget{tag-P324}{}\label{rem:II:two-direction-affine-fibres}
\StatementRecord{rem:II:two-direction-affine-fibres}{proved}{P324}{}
The reduced equation $f=y(x^2+y^2z^2)$ has generic
$D_4$ along $x=y=0$ and generic $A_1$ along $x=z=0$.
Two spectral directions collide at their intersection.
Its non-plane normalization sheets $x=\pm\mathrm i yz$
have pole divisor $2(y=0)+(z=0)$, with simple normal
crossings. Its affine fibre $f=1$ is isomorphic to
$(\C^*)^2$ and has zero intersection form. More generally,
for every integer $n\ge1$, the affine fibre of
$y(x^2+y^2z^n)$ has zero intersection form.
\end{remark}
\begin{proof}
For $z\ne0$, the three linear normal factors
$y$, $x+\mathrm i yz$, $x-\mathrm i yz$ are distinct,
giving the ordinary triple point $D_4$. For $y\ne0$,
the last two factors are coordinates transverse to
$x=z=0$: their Jacobian in $(x,z)$ is $-2\mathrm i y$.
They give the transverse node $A_1$.
On the smooth sheet $x=\pm\mathrm i yz$, the residue
of $dx\wedge dy\wedge dz/f$ is
$dy\wedge dz/(\pm2\mathrm i y^2z)$, proving the
stated pole divisor in the coordinates $(y,z)$.

For $n=2$ the maps are
\[
 \begin{split}
 a&=x+\mathrm i yz,\qquad b=x-\mathrm i yz,\\
 x&=(a+b)/2,\qquad y=(ab)^{-1},\qquad
 z=(a-b)ab/(2\mathrm i).
 \end{split}
\]
They are inverse by substitution, and their values satisfy
$y(x^2+y^2z^2)=1$. The product of two unit circles
generates $H_2((\C^*)^2;\Z)$; a radial translation
gives a disjoint homologous copy, so its square is zero.

For general $n$, the map $t=1/y$, $u=x/y$, with $z$
fixed, identifies the affine fibre with
\[
 W_n=\{u^2+z^n=t^3,\ t\ne0\};
 \qquad y=1/t,\quad x=u/t
\]
is its inverse. This is a smooth surface, since the
$t$-derivative is nonzero. Positive real dilation acts by
\[
 (u,z,t)\longmapsto(s^{3n}u,s^6z,s^{2n}t),\qquad s>0.
\]
Both sides of its equation scale by $s^{6n}$. It gives
the product diffeomorphism
\[
 (0,\infty)\times\{(u,z,t)\in W_n:|t|=1\}
 \longrightarrow W_n
\]
with inverse obtained from $s=|t|^{1/(2n)}$ and the
dilation by $s^{-1}$. Compact cycles have bounded
intervals in this positive real coordinate. A sufficiently
large dilation displaces one from another and is isotopic
to the identity, proving that all their pairings vanish.
\end{proof}

\subsection{Genus of residual cubic covers}\label{subsec:II:cubic-spectral-genus}

\begin{lemma}[Genus of residual cubic covers]
\hypertarget{tag-P323}{}\label{lem:II:cubic-spectral-genus}
\StatementRecord{lem:II:cubic-spectral-genus}{proved}{P323}{\cite[III.17--III.18, pp.~45--47]{Beauville1978}}
Let $E$ be a holomorphic rank-two bundle of degree $e$
on $\PP^1$, and let $q\in H^0(\PP^1,\operatorname{Sym}^3E)$
have three distinct generic roots. Use the quotient
projective bundle $p:\PP(E)\to\PP^1$, with tautological
class $\xi$ and fibre class $F$. Remove the common zero
divisor of the coefficients of $q$, of degree $b\ge0$.
Its residual zero divisor $\Gamma$ is reduced, has class
$3\xi-bF$, and satisfies
\begin{equation}
                    p_a(\Gamma)=3e-2-2b.
\label{eq:II:refined-4-55d}
\end{equation}
For irreducible $\Gamma$ with normalization $H$, this is
$g(H)+\delta(\Gamma)$, where $\delta$ is the normalization
colength.

For $\Gamma=\Gamma_H\cup\Gamma_R$, with $\Gamma_H$
irreducible of ruling degree two and $\Gamma_R$ a section,
let $H$ normalize $\Gamma_H$, put
$\delta_H=\delta(\Gamma_H)$ and
$n=\Gamma_H\cdot\Gamma_R$. There is an integer $t$ with
\[
 \Gamma_R=\xi-tF,\qquad
 \Gamma_H=2\xi+(t-b)F,
 \qquad g(H)+\delta_H+n=3e-1-2b.
\]
In degree $e=2$ the individual formulas are
\begin{equation}
 \begin{split}
       g(H)+\delta_H&=p_a(\Gamma_H)=1+t-b,\\
       n&=4-t-b,\\
       g(H)+\delta_H+n&=5-2b.
 \end{split}\label{eq:one-junction-spectral-bound}
\end{equation}
For $e=2$ and $g(H)\ge1$ they give $b\le2$.
When also $b\ge2$,
\begin{equation}
 b=2,\qquad g(H)=1,\qquad\delta_H=0,
                 \qquad\Gamma_H\cdot\Gamma_R=0.
\label{eq:II:manual:4.57a}
\end{equation}
Thus the double component is smooth and disjoint from
the section.
\end{lemma}
\begin{proof}
The tautological quotient $p^*E\twoheadrightarrow\OO(1)$
and its kernel $L$ are the bundles in
\cite[III.17--III.18, pp.~45--47]{Beauville1978}.
Their holomorphic exact sequence gives
$L=p^*\det E\otimes\OO(-1)$ and
$T_{\PP(E)/\PP^1}=\mathcal Hom(L,\OO(1))$.
Consequently
\[
 \xi^2=e,\quad \xi F=1,\quad F^2=0,
 \qquad K_{\PP(E)}=-2\xi+(e-2)F.
\]
Here the first equality is the degree-two Chern identity
$c_2(p^*E)=c_1(L)\xi=0$; the relative tangent formula
gives the canonical equality by determinants. The
tautological class and the fibre class give the integral
cohomology basis, as in the sphere-bundle argument in
the proof of III.18. These computations use the stated
holomorphic bundles and their maps.

Applying the tautological quotient to $q$ gives a section
of $\OO(3)$. Removing its vertical common zero divisor
leaves $3\xi-bF$. Distinct generic roots make every
horizontal component reduced. Adjunction gives
\[
 (3\xi-bF)(\xi+(e-2-b)F)=6e-6-4b,
\]
and hence \eqref{eq:II:refined-4-55d}. Normalization
of an irreducible curve gives $p_a=g+\delta$.

The section has fibre degree one, so its class is
$\xi-tF$ in this same integral basis; subtraction gives
the class of $\Gamma_H$. Adjunction and intersection give
\[
 p_a(\Gamma_H)=e-1+t-b,\qquad n=2e-t-b.
\]
Equivalently, the exact union sequence
\[
 0\longrightarrow\OO_\Gamma\longrightarrow
 \OO_{\Gamma_H}\oplus\OO_{\Gamma_R}
 \longrightarrow\OO_{\Gamma_H\cap\Gamma_R}
 \longrightarrow0
\]
gives $p_a(\Gamma)=g(H)+\delta_H+n-1$, because the
section is a smooth rational curve. Setting $e=2$ proves
\eqref{eq:one-junction-spectral-bound}. Since
$\delta_H,n\ge0$, positive genus gives $b\le2$.
For $b\ge2$, the same identity reads
\[
 (g(H)-1)+\delta_H+n+2(b-2)=0.
\]
All terms are nonnegative, so all vanish. This proves
\eqref{eq:II:manual:4.57a} and the asserted smoothness
and disjointness.
\end{proof}

\section{Conductor junctions and the irreducible classification}\label{sec:II:irreducible-conductor-junctions}

\subsection{The conductor graph and the horizontal spectral cover}

\begin{proposition}[The horizontal conductor cover and its rational third sheet]\hypertarget{tag-02XU}{}\label{night:P4.55}\label{stmt:P4.55}
\StatementRecord{night:P4.55}{proved}{02XU}{}

In the irreducible-or-normal-crossings torus-fibration setting
of \cref{def:II:standing-torus-fibration}, let $mS$ have
irreducible reduced support with $e(S)>0$ and normalization
resolution ruled over a positive-genus curve. The image $A$
of its horizontal pole section $H$ is a smooth rational curve. The map
$H\to A$ has degree exactly two. The third generic
normalization preimage along $A$ lies on one rational
vertical conductor component, mapping isomorphically to $A$.
The transverse type on the other image component is arbitrary. A separated self-identification point on either image
component is impossible.
\end{proposition}
\begin{proof}

\begin{enumerate}[label=\arabic*.,leftmargin=*]
\item Write $\nu:N\to S$ for normalization and
$\sigma:T\to N$ for its minimal resolution. By \cref{night:P4.54,night:P4R.18}, the pole section $H$ is
noncontracted, has genus $g\ge1$, and maps with degree two
or three to $A^\nu=\mathbb P^1$. The full generic transverse
germ along $A$ is $D_4$. The reduced image conductor is
$C=A\cup B$, with two rational normalized components and
one graph cycle.

\item The ruled application in the proof of
\cref{night:P4.35}, with this $\sigma$ and its doubled-section
pole divisor, makes the reduced normalization conductor
$D\subset N$ a tree in its point-identification graph. By \cref{night:P4.33}, the ruling $a:T\to E$ and its
section $h:E\simeq H\hookrightarrow T$ descend to
$q:N\to E$ and $s=\sigma\circ h:E\to N$, with
$q\circ\sigma=a$ and $q\circ s=\operatorname{id}_E$.
Thus $s$ is an embedding, and
$\sigma|_H=s\circ(a|_H)$ identifies the section $H$
with its smoothly embedded image in $N$. Every other normalized conductor
component is rational. Here $D$ includes the vertical conductor components and
all their point identifications.

\item At any point of $C$, points of $N$ over it correspond to
local irreducible hypersurface branches of $S$. Two such
branches intersect along a curve, by the height bound for
two equations in a regular local ring of dimension three.
Their nonempty supports among the local analytic branches
of $C$ therefore pairwise intersect.

\paragraph{Separated self-identifications.}
Suppose $p\in B\setminus A$ has two analytic branches.
Let $D_B$ be the union of conductor components mapping
to $B$, with $c_B$ normalized components and $s_B\ge1$
connected components. Its graph is a forest. Apply \cref{lem:II:conductor-two-branch-cycle} with
$(D,C,p,c,k)=(D_B,B,p,c_B,s_B)$: every normalized component
maps finitely and surjectively to $B^\nu$, and at $p\notin A$
the local hypersurface branches have exactly the two
possible image branches used in that count. Their supports
pairwise intersect by the preceding item. Let
$\kappa:D_B^\nu\to D_B$ and $\eta:B^\nu\to B$ be normalization,
and let $\phi:D_B^\nu\to B^\nu$ satisfy
$\eta\phi=(\nu|_{D_B})\kappa$.
Label $p_+,p_-\in\eta^{-1}(p)$ so that the common
support branch in the cited count is the plus branch, and put
$r_\pm=|\phi^{-1}(p_\pm)|$ and
$r_q=|\kappa^{-1}(q)|$. Writing $n$ for
the number of normalization points over $p$, the count
therefore gives $r_+\ge n$, $r_-\ge c_B$, and
\[
 \sum_{q\in\nu^{-1}(p)}(r_q-1)=r_++r_--n\ge c_B.
\]
The whole forest has contribution $c_B-s_B<c_B$, a
contradiction. The identical substitution with $A$ and all
components mapping to it excludes a self-identification
of $A$ away from $B$. Thus the image cycle consists of
two distinct intersections of $A$ and $B$, or one
three-branch point belonging to both components.

\paragraph{The primary conductor and smoothness of $A$.}
Let $C_{A,\mathrm{sch}}$ be the primary component of the image conductor supported on $A$. Apply \cref{lem:II:horizontal-primary-conductor} with the
normalization $N\to S$, its minimal resolution $T\to N$,
the noncontracted section $H$ and its image $C_H=A$.
Use the pole divisor $D_T$ of \cref{night:P4R.6} on $T$.
The preceding items give its horizontal part $2H$, rational
vertical trees and positive ruling genus. Thus the lemma's
divisor is $D_T$, while $D\subset N$ above is the reduced
normalization conductor. The primary restriction gives
\begin{equation}
                    g-1\le\chi(\mathcal O_{C_{A,\mathrm{sch}}}).
\label{eq:II:refined-4-55a}\end{equation}
The generic ordinary-triple conductor is the transverse
maximal ideal squared. Purity of the conductor
quotient identifies this primary ideal with $P_A^{(2)}$,
as in \cref{night:P4.31}. Apply \cref{lem:II:rational-conormal-defects} with
$C=A$, $E=A^\nu=\mathbb P^1$ and $P=P_A$. Its normalized
derivative lattice and symbolic-square quotient give
\begin{equation}
 \chi(\mathcal O_{C_{A,\mathrm{sch}}})
       =5-\sum_p\bigl(\delta(A,p)+\tau_p-r_p\bigr).
\label{eq:II:refined-4-55b}\end{equation}
A planar point contributes exactly $5\delta(A,p)$, and a
nonplanar singular point with at most two branches contributes
at least seven. There cannot be three branches of $A$ at one
point: the identity
$\sum_p(r_p(C)-1)=2$ would then leave no identification for
the necessary connection with $B$; placing that connection
at the same point gives at least four branches and also
exceeds the identity.

\item Since $g\ge1$, \eqref{eq:II:refined-4-55a}--\eqref{eq:II:refined-4-55b} bound the total defect by
five. If $A$ is singular, it is planar everywhere, has total
delta invariant one, and $g=1$. It is then a locally complete
intersection curve in $Y$. Its conormal
$Q=P_A/P_A^2$ is locally free of rank two. Let
$j:\mathbb P^1\to A$ be its normalization and put $E_0=j^*Q$.
Adjunction gives
\[
             \omega_A=(\omega_Y|_A)\otimes(\det Q)^{-1}.
\]
The degree of $j^*\omega_Y$ is zero because $b_2(Y)=0$.
Also $p_a(A)=1$, so $\deg j^*\omega_A=0$. The degree of a
line bundle is preserved by finite birational normalization:
the normalization sequences for that line bundle and for
$\mathcal O_A$ have the same finite-length correction.
Consequently
\begin{equation}
                              \deg E_0=0.                 
\label{eq:II:refined-4-55c}\end{equation}

\item The reduced equation has generic order three along
$A$. Powers of a complete-intersection curve ideal are
unmixed, so $P_A^{(3)}=P_A^3$. Its third normal term is a global section
\[
 q_3\in H^0\!\left(A,
          \operatorname{Sym}^3 Q\otimes\mathcal O_Y(S)|_A\right).
\]
The line twist has degree zero and becomes trivial on
$\mathbb P^1$. On the quotient projective bundle
$\mathbb P(E_0)$, with tautological class $\xi$ and fibre
class $F$, the pulled-back cubic has zero divisor of class
$3\xi$. Remove its common zero fibres, of total degree
$b\ge0$. The residual spectral curve $\Gamma$ has class
$3\xi-bF$. Its horizontal components are reduced, because
the generic $D_4$ cubic has three distinct roots.

\item At a generic point of
$H$, the first normal derivative of $T\to S\subset Y$
gives the direction of its smooth transverse branch and
hence a point of $\mathbb P(E_0)$. On the smooth compact
curve $H$, its two local coordinates extend after cancelling
their common power. This gives a holomorphic map over the map $H\to\mathbb P^1$. Distinct generic preimages
on $H$ give distinct cubic roots. Hence its map onto the
spectral component $\Gamma_H$ is finite and birational,
and $H$ is that component's normalization. In degree three,
$\Gamma=\Gamma_H$. In degree two, precisely one remaining
rational vertical conductor component supplies the third
root, with degree one over $A$. Its direction curve
$\Gamma_R$ is a section and is isomorphic to $\mathbb P^1$.
Thus $\Gamma=\Gamma_H\cup\Gamma_R$. These exhaust the
three generic roots.

\item Apply \cref{lem:II:cubic-spectral-genus} to this
$E_0$, the pulled-back normal cubic and its common zero
divisor of degree $b$. The preceding derivative maps
identify the same $H$ with the normalization of
$\Gamma_H$, and in degree two identify the third sheet
with $\Gamma_R$. The lemma gives
\eqref{eq:II:refined-4-55d} for this residual divisor.
With \eqref{eq:II:refined-4-55c}, it equals $-2-2b$.
In degree three the normalization formula gives
$p_a(\Gamma)\ge g\ge1$. In degree two the lemma's
union formula gives $p_a(\Gamma)\ge g-1\ge0$.
Both alternatives are impossible. Therefore $A$ is smooth.

\paragraph{The three-branch image point in degree three.}
Suppose now that $H\to A$ has degree three. Then $H$ is
the only conductor component mapping onto $A$. Since $A$
is smooth, a single three-branch image point consists of
one $A$ branch and two $B$ branches, denoted $+$ and $-$.

\item Let $c_B$ be the number of normalized conductor components
mapping to $B$. The entire conductor has $1+c_B$
normalized components and a tree graph, so its total
point-identification contribution is $c_B$. At the proposed
point write $r_A\ge1$, $r_+,r_-\ge c_B$ for the numbers
of conductor branches, and $n$ for the number of normalization
points. Each normalization point has at most one $A$
conductor branch, because $H$ is embedded.

\item A singleton support $\{A\}$ forces every point to carry
$A$. Then $n=r_A$, and the graph contribution is
$r_++r_-\ge2c_B>c_B$. A singleton $\{+\}$ forces
$r_+\ge n$, giving contribution at least
$r_A+r_-\ge1+c_B$; the same applies to $\{-\}$.
Thus every support has at least two elements. Consequently
$r_A+r_++r_-\ge2n$, and its graph contribution is at least
\[
       \frac{r_A+r_++r_-}{2}
                         \ge c_B+\frac{r_A}{2}>c_B.
\]
This is again impossible. Only two distinct $A$--$B$
junctions remain in degree three.

\paragraph{Both degree-three junctions are common cubic zeros.}
Let these junctions be $p,q$. If a local normalization
branch over $p$ has conductor only over $A$, pairwise
hypersurface intersections force every normalization point
over $p$ to carry $A$. Thus $n=r_A$, and the graph
contribution at $p$ is $r_B\ge c_B$, using the entire
budget. At $q$ some normalization point must carry both
$A$ and $B$: otherwise an $A$-only point and a $B$-only
point have disjoint supports. That point contributes
positively to the graph, contradicting the budget. There
is therefore no $A$-only local branch at either junction.

\item Write the smooth curve as $A=\{x=y=0\}$. If its nonzero
cubic normal form has a simple root, use the normal blowup
chart $x=yu$. The implicit-function theorem on the strict
equation gives that root as $u=\alpha(y,z)$. The reduced function has the analytic factor
\[
                         f_1=x-y\alpha(y,z).
\]
This branch is smooth. Exact division writes $f=f_1f_2$,
and restriction of $f_2$ to this branch is $y^2$ times a
unit: the other cubic directions differ from the simple
root at the point. Its intersections with all other
branches are therefore supported only on $A$. It is the
forbidden $A$-only normalization branch.

\item Every nonzero binary cubic has a simple root unless it is
a cube. \Cref{night:P4R.19} excludes a nonzero cube. Hence the
whole cubic tensor vanishes at both $p$ and $q$.

\paragraph{The degree-three genus contradiction.}
For the now smooth curve $A=\mathbb P^1$, adjunction gives
$\deg N_{A/Y}=-2$, so $E=N_{A/Y}^*$ has degree two.
The cubic has at least two common zeros. Its residual
spectral divisor is $3\xi-bF$ with $b\ge2$, and is
irreducible with normalization $H$. Formula \eqref{eq:II:refined-4-55d} gives
\[
                    p_a(\Gamma)=4-2b\le0,
\]
contradicting $p_a(\Gamma)\ge g\ge1$. Degree three is
excluded. The degree-two alternative and its unique rational
third sheet are precisely the possibilities already identified
by the spectral morphism above.

\end{enumerate}
\end{proof}


\subsection{The second conductor component}\label{subsec:II:second-conductor-branches}

\begin{proposition}[At least three branches along the second conductor component]\hypertarget{tag-02XV}{}\label{night:P4.56}\label{stmt:P4.56}
\StatementRecord{night:P4.56}{proved}{02XV}{\cref{night:P4.55,lem:II:smooth-curve-intersection,lem:II:irrational-ruled-boundary}}

In the setting of \cref{night:P4.55}, the second conductor
image component $B$ has at least three generic normalization
branches. Its transverse multiplicity, conductor exponents
and tangency orders are arbitrary.
\end{proposition}
\begin{proof}

\begin{enumerate}[label=\arabic*.,leftmargin=*]
\item Use the smooth curve $A$, its degree-two horizontal
cover $H$, and its degree-one vertical cover $R$ from
\cref{night:P4.55}. All other
normalized conductor components map to $B$. Let $r_B$ be
the generic normalization branch number along $B$, and
$c_B$ the number of global conductor components over it.
Their finite degrees sum to $r_B$. We exclude $r_B=1,2$.
The normalized conductor has $2+c_B$ components and a tree
graph, with total identification contribution $1+c_B$.
Distinct points on one normalized component cannot be
identified, since that would create a graph loop.

For local hypersurface branches $F,G$ meeting only
along the smooth $A$, apply
\cref{lem:II:smooth-curve-intersection} with
$(\phi,\psi,\Gamma)=(F,G,A)$. Generic intersection
length one makes both branches smooth. At length two,
the quotient is \eqref{eq:II:manual:4.56a}; when $G$
has generic normal order two, $F$ is the smooth factor.

\paragraph{Two distinct image junctions.}
At a two-branch image point, let $r_A,r'_B$ count the
normalized conductor points over its two branches, and
let $n$ count normalization points of $S$. Pairwise
intersection of the hypersurface branches gives a common
element in their supports in $\{A,B\}$. Hence
$n\le\max(r_A,r'_B)$, and the graph contribution is at
least $\min(r_A,r'_B)$. Here $r_A=2$ or three, from the
one or two $H$ points and the unique $R$ point.

\item If $c_B=2$, both degrees are one, and $r'_B=2$ at either
junction. The two contributions total at least four,
exceeding the budget three. If $c_B=1$ and its degree is
two, an unramified junction costs at least two and the
other at least one, exceeding the budget two. Thus one
component $V$ remains, with exactly one normalized point
over each junction.

\item If $B$ belongs to every branch support at a junction,
there is only one normalization point there and its
contribution is $r_A\ge2$. The other junction makes this
impossible. Therefore $A$ belongs to every support. The
budget forces contribution exactly one at each junction:
every normalization point has one $A$ branch, and the
sole $V$ point is paired with one of them. Connectedness
forces the graph path $H$--$V$--$R$, with its edges at
the two different junctions.

\item At the $R$--$V$ junction, let $F$ be the corresponding
hypersurface branch. The other branches come from $H$
and carry only $A$. If its two preimages are distinct,
intersection with either such branch has support $A$ and
generic length one, because the full generic $A$ germ
is an ordinary triple point. The preceding calculation
makes $F$ smooth. If the two $H$ preimages coalesce,
their one hypersurface branch $G$ has generic normal
order two along $A$, and the intersection with $F$ has
length two. Formula \eqref{eq:II:manual:4.56a} again makes $F$ smooth.

\item This is impossible. The branch $F$ carries all the nearby
generic $B$ preimages: either the single branch of a
singular unibranch transverse germ, or both branches
when $V\to B^\nu$ has degree two and ramifies at this
point. In either case $F$ is intrinsically nonnormal
along the nearby $B$ germ. It cannot be smooth in a
neighbourhood of the junction.

\paragraph{A three-branch image point with one $B$ component.}
The image branches are $A,B_+,B_-$, since $A$ is smooth.
If $c_B=1$, the normalized points on $V$ over the two
$B$ signs are distinct and cannot be identified in $N$.
No normalization point can therefore carry both signs.
Pairwise hypersurface intersections force every point
carrying either sign to carry $A$ as well.

\item The graph budget two now forces exactly one point over
each sign, each paired with exactly one $A$ branch, and
no other identification. The pairs are $H$--$V$ and
$R$--$V$. An extra $H$ preimage is $A$-only. At the
$R$--$V$ point, the other $H$ branch or branches intersect
its hypersurface branch only along $A$, because their
$B$ supports have the opposite sign or are empty. The
same length-one or length-two calculation makes this
branch smooth, contradicting its intrinsic nonnormality
along the $B$ germ whose generic preimages it contains.

\paragraph{A three-branch point with two $B$ components.}
Write the two degree-one components as $V,W$. If $h=1$
or two is the number of $H$ points over the $A$ value,
the total normalized branch count is
\[
                         T=(h+1)+2+2=h+5.
\]
The graph budget three gives $n\ge h+2$ normalization
points. No singleton image support is possible: a
singleton $\{A\}$ forces $n\le h+1$, and a singleton
$\{B_+\}$ or $\{B_-\}$ forces $n\le2$. Hence every
point has at least two image labels and $n\le T/2$.
This excludes $h=2$. For $h=1$ there are precisely
three points, each with exactly two conductor branches
of different image labels. Up to signs and interchange
of $V,W$, their supports are
\begin{equation}
 \begin{array}{c|c}
 \text{hypersurface branch}&\text{conductor branches}\\ \hline
 F_R&R,\ B_-\text{ on }V\\
 F_H&H,\ B_+\text{ on }W\\
 F_0&B_+\text{ on }V,\ B_-\text{ on }W.
 \end{array}
\label{eq:II:manual:4.56b}
\end{equation}
The tree is the path $H$--$W$--$V$--$R$, and $H\to A$
is ramified with local degree two at its point. There are
exactly four labelled pairings: choose the component
$V$ or $W$ and the sign at the $H$ endpoint; the
opposite endpoint must use the other component and
the opposite sign, and the remaining pair is forced.
Interchanging $V,W$ and the signs gives the displayed form.

\item The intersection of $F_R$ and $F_H$ is supported
on $A$ and has generic length two. The equation of
$F_H$ has generic normal order two. Apply
\cref{lem:II:smooth-curve-intersection} to
$(\phi,\psi,\Gamma)=(F_R,F_H,A)$. It gives the
coordinates $X=F_R$, $Y$, $z$ and the exact equation
\eqref{eq:II:manual:4.56c}, with $\psi=F_H$.
The coordinate map is the parameter-Morse map written
in that lemma, with inverse given by its unique small
solution; it keeps $X,z$ and the ideal of $A$.
At generic $A$, the two $H$ sheets are distinct, so
$k(0,z)\not\equiv0$. The irreducible branch $F_H$
therefore has the finite birational map
\eqref{eq:II:manual:4.56d}, with inverse $u=Y/X$ off
$A$, and the normalization
\[
 \widehat N=\{v^2+h(X,z)=0\},\qquad
 k=a^2h,\quad u=av,\quad Y=Xav
\]
of the same lemma.

\item On $\widehat N$, the inverse-image ideal of
$A$ is $(X)$. Finiteness prevents a curve from mapping
to the junction point, so every component of $X=0$
dominates $A$. Its generic sheets are exactly those
of $H$, the third sheet belonging to $F_R$. Its
reduced support is therefore the embedded horizontal
curve $H_N\subset N$. The descended ruling and its
section in \cref{night:P4.55} identify $H_N$ with
the smooth $H$. Its map to $A$ has local degree two
and is ramified at this point. These are precisely the
smooth-curve and ramification hypotheses in the last
clause of \cref{lem:II:smooth-curve-intersection};
it gives $h_z(0,0)\ne0$ and smoothness of $\widehat N$.
The minimal normalization resolution is the identity
in this neighbourhood.

\item Let $\mu_B\ge1$ be the generic intersection multiplicity of
the two branches along $B$. It is the same on both
local signs, because they belong to the same irreducible
image component. Let $G$ define $F_0$, and set
\[
                      b=\operatorname{ord}_{p,A}(G|_A)>0.
\]
On the smooth branch $F_R$, its divisor is $\mu_B B_-$.
Restricting to $A$ gives
\begin{equation}
                 b=\mu_B I_{F_R}(A,B_-)\ge \mu_B.                
\label{eq:II:manual:4.56e}
\end{equation}
On the smooth normalization $\widehat N$, the divisor of $G$
is $\mu_B W$. Apply \cref{lem:II:irrational-ruled-boundary} to the
same resolved pole divisor, whose horizontal part is $2H$.
It has only ordinary double intersections. Since this
normalization is already smooth, $H$ and $W$ meet
transversely here. Its restriction to
$H$ consequently has order $\mu_B$. But it is the pullback
of the same ambient restriction $G|_A$ under the local
degree-two map. Hence
\[
                              2b=\mu_B,
\]
contradicting \eqref{eq:II:manual:4.56e}, since $2b-\mu_B\ge \mu_B>0$.
Every case with $r_B\le2$ is excluded.

\end{enumerate}
\end{proof}

\subsection{Two distinct conductor junctions}\label{subsec:II:two-conductor-junctions}

\begin{proposition}[Exclusion of two distinct conductor junctions]\hypertarget{tag-02XW}{}\label{night:P4.57}\label{stmt:P4.57}
\StatementRecord{night:P4.57}{proved}{02XW}{\cref{night:P4.55,night:P4.56,night:P4R.19,lem:II:cubic-spectral-genus,lem:II:smooth-curve-intersection,cdp:st:A.5.17}}

In the degree-two configuration of \cref{night:P4.55,night:P4.56},
the image-conductor graph cycle cannot be realized at
two distinct intersections of its two components.
\end{proposition}
\begin{proof}

\begin{enumerate}[label=\arabic*.,leftmargin=*]
\item The component $B$ has at least three generic normalization
branches, hence transverse multiplicity at least three.
Along the smooth $A$, choose coordinates $A=\{x=y=0\}$,
with $z$ on $A$. The equation belongs to $(x,y)^3$;
write $q_3(x,y;z)$ for its normal cubic.

\item A local branch of $B$ has first nonzero normal displacement
$(a,b)t^q$, with $(a,b)\ne0$. All ambient Hessian entries
of the equation vanish identically along $B$, by
its generic multiplicity. The first normal-order term of
its normal Hessian is
\[
       t^q\operatorname{Hess}_{x,y}
                    (q_3(\,\cdot\,,\,\cdot\,;0))(a,b).
\]
Terms of higher normal order and terms from varying $z$
have larger order in $t$. This matrix is therefore zero.
For a homogeneous binary cubic $Q$, a nonzero vector
annihilating this Hessian makes the directional derivative
of $Q$ the zero quadratic. In coordinates with that vector
as first axis, $Q$ is a cube of the other variable, or zero.
\Cref{night:P4R.19} excludes the nonzero cube. Thus the cubic vanishes
at every $A$--$B$ junction, including tangent junctions.

\paragraph{The spectral equality.}
Let $p,q$ be the two distinct junctions. Their cubic
vanishing gives a common zero divisor of degree $b\ge2$.
Use the bundle $E=N_{A/Y}^*$ of degree two and the
normal-direction maps of \cref{night:P4.55}. They identify
this same $H$ with the normalization of $\Gamma_H$ and
the rational third sheet with the section $\Gamma_R$.
Apply \cref{lem:II:cubic-spectral-genus} to this cubic,
$e=2$, and its residual divisor $3\xi-bF$.
Since $g(H)\ge1$, it gives \eqref{eq:II:manual:4.57a}.
The two common cubic zeros are simple, the double spectral
cover is smooth, and it is disjoint from the third section.

\paragraph{The graph support at a junction.}
Let $c_B$ count the global conductor components over $B$.
The graph budget is $c_B+1$. At either junction the
contribution is at least $\min(r_A,r'_B)$, where
$r_A\ge2$ and $r'_B\ge c_B$ count normalized points
over its two image branches.

\item If $c_B=2$, the two contributions total at least four,
exceeding the budget three. If $c_B=1$, a junction with
at least two $B$ preimages costs at least two, and the
other at least one, exceeding the budget two. Thus both
have one $B$ preimage. The two-junction argument of \cref{night:P4.56} gives the
path $H$--$V$--$R$, with one edge at each junction.
At the $R$--$V$ junction the local branch carries all
nearby generic $B$ sheets of $V$, regardless of its
finite degree. Its intersection with the $H$ branch or
branches is a length-one or length-two thickening of
smooth $A$. Apply \cref{lem:II:smooth-curve-intersection} to that branch and
the indicated $H$ branch, with $\Gamma=A$ and $r=1$ or
$r=2$. It makes the branch smooth, contradicting its
intrinsic nonnormality from the at least three $B$ sheets.

\item Therefore $c_B\ge3$. If $A$ belongs to every branch
support at a junction, that point costs at least
$r'_B\ge c_B$, and the other costs at least two.
This exceeds $c_B+1$. The two-element support argument
consequently forces $B$ to belong to every local
hypersurface branch at each junction. This conclusion includes branches carrying only $B$.

\paragraph{Choosing distinct normalization points.}
Write $B_H$ for the ruling base, and use the maps
$q:N\to B_H$ and $h:B_H\simeq H$ in
\cref{night:P4.55}. Put $r_N=h\circ q:N\to H$.
Via its embedded image in $N$, the restriction to $H$
is the identity. The rational curve $R$ maps to a fixed
point $b_0\in H$. If a point of $H$ and a point of $R$ have
the same image in $N$, that $H$ point is $b_0$.
This includes identifications through a contracted
exceptional chain. All such coincidences therefore lie
over the single value $a_0\in A$ which is the image of
$b_0$. Choose a junction different from $a_0$, and call
it $p$. Its $H$ preimages and its $R$ preimage are
distinct normalization points. Different $H$ points
are already distinct because $H$ is embedded.

\item The local equation factors into the corresponding
irreducible hypersurface factors and any factors carrying
only $B$. For an $A$-containing factor, let $b_i\ge0$
be the common $z$-order of its first nonzero normal term.
For a factor not containing $A$, let $c_j\ge1$ be the
order of its nonzero restriction to $A$ at $p$.
Gauss's lemma over $\mathbb C\{z\}$ gives
\begin{equation}
                         \sum_i b_i+\sum_j c_j=1.         
\label{eq:II:manual:4.57b}
\end{equation}
Indeed primitive homogeneous polynomials have nonzero
product modulo $z$, so their common coefficient orders
add exactly. The right side is the simple cubic zero
of \eqref{eq:II:manual:4.57a}.

\item If $H$ is unramified here, there are three
$A$-containing factors of normal order one. At least
two have nonzero linear normal terms, by \eqref{eq:II:manual:4.57b}.
Their directions differ by \eqref{eq:II:manual:4.57a}. The implicit-function
theorem makes their intersection exactly $A$, whereas
both must contain $B$. This is a contradiction.
Thus $H$ is ramified at $p$.

\item There are now two $A$-containing factors, denoted by
$F_H$ and $F_R$, of normal orders two and one. Their
residual spectral directions differ. Equality \eqref{eq:II:manual:4.57b}
has precisely the following three possibilities.

\paragraph{A factor carrying only $B$ uses the zero.}
Here $b_H=b_R=0$ and one such factor restricts to $A$
with order one. The branch $F_R$ is smooth. In coordinates
$X=F_R$, $A=\{X=Y=0\}$, restriction of $F_H$ to $X=0$
is $Y^2$ times a unit, since its residual direction
differs from the $R$ direction. Thus these two branches
intersect only along $A$. They both contain $B$, which
is impossible.

\paragraph{The quadratic factor uses the zero.}
Here $b_H=1$, $b_R=0$, and there is no other factor.
The branch $F_R$ is again smooth, and in its coordinates
exact division gives
\[
                         F_H|_{X=0}=Y^2h(Y,z).
\]
The value $h(0,z)$ has a simple zero: the residual
quadratic is nonzero on the $R$ direction, so its
common scalar factor survives with order one on this
restriction. Thus $h_z(0,0)\ne0$. Its zero curve is the
extra image branch $B$, with intersection multiplicity
one between $F_H$ and $F_R$.

\item The smooth $F_R$ branch supplies one generic $B$
normalization branch. Each branch supplied by $F_H$
has positive intersection multiplicity with it. There
can therefore be at most one such branch. The full
$B$ germ would have at most two branches, contradicting
\cref{night:P4.56}.

\paragraph{The linear factor uses the zero.}
Finally $b_R=1$, $b_H=0$, with no other factor.
The nonzero quadratic form of $F_H$ has a double root.
Choose a complementary normal coordinate $X$ so that,
for $A=\{X=Y=0\}$, the restriction $F_H|_{X=0}$ is
$Y^2$ times a unit. Apply
\cref{lem:II:smooth-curve-intersection} with
$(\phi,\psi,\Gamma)=(X,F_H,A)$ and $r=2$.
Here $X=0$ is the auxiliary smooth hypersurface used
for this calculation. The lemma's coordinate map keeps
$X,z$ and gives the exact equation
\[
                         F_H=Y^2+X^2k(X,z).
\]
Its spectral curve is $u^2+k(0,z)=0$. By
\eqref{eq:II:manual:4.57a} this curve is smooth, and its
double cover ramifies at $p$. Hence $k_z(0,0)\ne0$.
The last clause of \cref{lem:II:smooth-curve-intersection}
identifies its smooth normalization with
$\{u^2+k(X,z)=0\}\to\{F_H=0\}$ by $Y=Xu$,
with inverse $u=Y/X$ off $A$. Thus $F_H$ supplies
one generic $B$ branch.

\item For $F_R$, the simple common zero of its linear normal
term makes its quadratic Taylor part
\[
                       zL(X,Y)+Q_2(X,Y),\qquad L\ne0.
\]
There is no $z^2$ term because $F_R$ vanishes on $A$.
Writing $L=\ell_X X+\ell_Y Y$, the Hessian minors
on $(X,z)$ and $(Y,z)$ are $-\ell_X^2$ and $-\ell_Y^2$.
At least one is nonzero, so its rank is at least two.
Rank three gives an isolated critical point. In rank two,
apply the square-completion coordinate construction in
the proof of \cref{cdp:st:A.5.17} to a nondegenerate
two-dimensional Hessian block, keeping the third coordinate
$W$ as parameter. Each step solves the corresponding
first derivative for its unique holomorphic centre and
uses $U=(w_1-\alpha)\sqrt A$, with the inverse given
by that equation's unique small solution. The nonzero
Schur complement permits the second step. This gives
the exact function $U^2+V^2+h(W)$. A nonisolated critical curve forces
$h'\equiv0$, hence $h=0$ because the factor vanishes
at the point. The resulting factorization
$(U+iV)(U-iV)$ contradicts local irreducibility.
Therefore $F_R$ has at most isolated singularities and
also supplies one generic $B$ branch. Again the full
branch number is two, contradicting \cref{night:P4.56}.

All allocations in \eqref{eq:II:manual:4.57b} have been excluded, proving
the assertion.

\end{enumerate}
\end{proof}

\subsection{A three-branch conductor junction}\label{subsec:II:three-branch-conductor-junction}

\begin{proposition}[Exclusion of a three-branch conductor junction]\hypertarget{tag-02YE}{}\label{night:P4R.20}\label{stmt:P4R.20}
\StatementRecord{night:P4R.20}{proved}{02YE}{\cref{night:P4.55,lem:II:cubic-spectral-genus,lem:II:smooth-curve-intersection,lem:II:hypersurface-branch-conductor,lem:II:plane-branch-conductor,lem:II:irrational-ruled-boundary}}

In the degree-two configuration of \cref{night:P4.55},
suppose that the generic transverse curve along the other
image component $B$ has at least three normalization branches.
The unique image-graph cycle cannot be concentrated at one
point with three local critical-curve branches, namely the
smooth branch $A$ and two local branches of $B$.
\end{proposition}
\begin{proof}

\begin{enumerate}[label=\arabic*.,leftmargin=*]
\item Write these two local branches as $B_+,B_-$ and the
point as $p$. Let $W_1,\ldots,W_c$ be the normalized
conductor components mapping onto $B$, with degrees
$d_i>0$ over $B^\nu$. Put
\[
                         r_B=\sum_i d_i\ge3.
\]
This is the number of generic normalization branches
of the full transverse plane curve along $B$.
The normalization conductor has components
$H,R,W_1,\ldots,W_c$ and a tree graph. Its total
point-identification contribution is $c+1$.
The horizontal section $H$ maps with degree two to
$A$, while the rational vertical curve $R$ maps
with degree one. The latter finite map is an isomorphism
onto the smooth $A$; its factorization through $N$ also
prevents two distinct points of $R$ from being identified.
The embedding of $H$ in $N$ is the section map
of \cref{night:P4.55}. Thus the maps of both $H$ and
$R$ into $N$ are injective. All conductor branches
below are counted on their normalization.

\item Let $E=N_{A/Y}^*$, so $\deg E=2$ by
adjunction and $c_1(Y)\cdot A=0$. The line
$\OO_Y(S)|_A$ has degree zero and is trivial on
$A\simeq\PP^1$. Thus the normal cubic is a section
of $\operatorname{Sym}^3E$. Its derivative maps in
\cref{night:P4.55} identify this same $H$ with the
normalization of $\Gamma_H$ and the rational third
sheet with the section $\Gamma_R$. Let $b$ be the
degree of its common zero divisor. Apply
\cref{lem:II:cubic-spectral-genus} with $e=2$ and
these residual components. With its integer $t$,
$\Gamma_R=\xi-tF$ and $\Gamma_H=2\xi+(t-b)F$;
its $\delta_H$ is the normalization colength, and
$n=\Gamma_H\cdot\Gamma_R$. Formula
\eqref{eq:one-junction-spectral-bound} gives $b\le2$
because $g(H)\ge1$.

\paragraph{The graph when $c\ge2$.}
Let $r_A$ be the number of normalized conductor
points over the branch $A$ at $p$. It is two if
$H\to A$ ramifies there, and three otherwise.
Let $r_+,r_-$ count the points over $B_+,B_-$.
Every $W_i$ is surjective to $B^\nu$, so
$r_+,r_-\ge c$. Let $n_p$ be the number of
normalization points of $S$ over $p$.
Their nonempty supports in $\{A,+,-\}$ pairwise
intersect: distinct local hypersurface components
intersect in a curve.

\item A singleton support $\{A\}$ forces every point to
carry $A$. Then $n_p\le r_A$ and the graph
contribution is at least $r_++r_-\ge2c>c+1$.
A singleton $\{+\}$ forces every point to carry
$+$ and gives contribution at least
$r_A+r_-\ge c+2$; the same applies to $\{-\}$.
Thus every support has at least two elements, and
\[
 \text{contribution at }p
       \ge\frac{r_A+r_++r_-}{2}\ge c+1.
\]
The whole global budget is $c+1$. Equality is
forced throughout:
\begin{equation}
                 r_A=2,\qquad r_+=r_-=c,\qquad n_p=c+1.
                                                               \label{eq:one-junction-graph-equality}
\end{equation}
Every point over $p$ has exactly two normalized
conductor branches, with different image labels,
and all graph identifications occur at $p$.
The horizontal cover ramifies over $p$.
Each $W_i$ has precisely one point over each of
$B_+,B_-$ and has valence two; $H$ and $R$ have
valence one. The connected tree is therefore
the path
\begin{equation}
                    H-W_1-W_2-\cdots-W_c-R.             \label{eq:one-junction-conductor-path}
\end{equation}
Choose the label $+$ at the $H$ endpoint. The next
point on $W_1$ has label $-$. Every internal point
has support $\{+,-\}$, so it joins that point to
the $+$ point of $W_2$. Induction gives the $+$
point on the left of every $W_i$ and its $-$
point on the right. The $R$ endpoint carries $-$.

\item The endpoint hypersurface factors, denoted
$f_H,f_R$, thus have supports $\{A,+\}$ and
$\{A,-\}$. The $c-1$ internal factors $f_j$
have supports $\{+,-\}$. These are all the local
irreducible hypersurface factors.

\paragraph{The length-two intersection.}
The set-theoretic intersection of $f_H$ and $f_R$ is
$A$ alone. Along generic $A$, the first factor supplies
two ordinary-triple branches and the second supplies
the third; their intersection multiplicity is two.
Apply \cref{lem:II:smooth-curve-intersection} with
$(\phi,\psi,\Gamma)=(f_R,f_H,A)$ and $r=2$.
The generic normal order of $f_H$ is two, so $f_R$
is the smooth factor. Its parameter-Morse coordinates,
keeping $u=f_R$ and the curve coordinate $z$, give
\begin{equation}
                    f_H=v^2+u^2g(u,z).
\label{eq:one-junction-quadratic-factor}
\end{equation}
The coordinate $v=(y-\mu(u,z))\sqrt{V(u,y,z)}$ has
the inverse given by its unique small solution for $y$.
Here $\mu(0,z)=0$ and $V$ is a unit; thus the map
preserves $A$ and is the identity on $u,z$.

Write the reduced local equation as
$f=U f_Rf_H\prod_j f_j$, with $U$ a unit. Its cubic
normal tensor in these coordinates is exactly
\begin{equation}
 q_3(u,v;z)=U(0,0,z)\left(\prod_j f_j(0,0,z)\right)
                       u\bigl(v^2+u^2g(0,z)\bigr).
\label{eq:one-junction-cubic-content}
\end{equation}
The coefficient of $uv^2$ in the last factor is one,
and $U(0,0,z)$ is a unit. Therefore the common zero
order is exactly the sum of the orders of $f_j|_A$.

\paragraph{The pole on the smooth rational branch.}
On the smooth branch $f_R=0$, apply
\cref{lem:II:hypersurface-branch-conductor} to the local
factors $f_R,f_H,f_j$ just specified. The unit $U$ gives
the same principal ideal, so its reduced local ring is
the ring in that lemma. The intrinsic conductor of the
smooth branch is the unit ideal. Thus its conductor is
generated by $(f_H\prod_j f_j)|_{f_R=0}$, through the
lemma's multiplication map. The factor $f_H$ contributes
$2A$, and each $f_j$ cuts out only $B_-$ there.

Apply \cref{lem:II:irrational-ruled-boundary} to the
pole divisor of \cref{night:P4R.6} on the same resolved
normalization $T$. Its horizontal part is $2H$, and its
reduced support has only ordinary double intersections.
This local normalization branch is already smooth, so
its minimal resolution is the identity near the point.
Hence $A$ and its $B_-$ curve are smooth
and transverse. Choose the latter's equation
$z=0$ on this surface, with $z$ still a
coordinate on $A$, and extend it to the ambient
threefold. Keeping $u$ and this $z$ fixed, apply the same
parameter-Morse coordinate map of
\cref{lem:II:smooth-curve-intersection}. On $u=0$,
it replaces the complementary coordinate by itself
times a unit, since its centre is zero there.
It therefore preserves $A$, given by that coordinate
zero, and $B_-$, given by $z=0$. Equations
\eqref{eq:one-junction-quadratic-factor} and
\eqref{eq:one-junction-cubic-content} hold in these
adapted coordinates, with their transformed unit $U$. Then
\[
             f_j|_{f_R=0}=\text{unit}\,z^{e_j},
                  \qquad e_j\ge1.
\]
Their sum $\kappa=\sum e_j$ is the generic conductor
exponent of the global component $W_c$ on
this branch. By \eqref{eq:one-junction-cubic-content} it is also the local
common cubic-zero order. The branch on $f_R=0$ is smooth. Apply
\cref{lem:II:plane-branch-conductor} to this same
generic transverse plane curve. Its intrinsic delta
term is zero, so the formula gives
\[
          \kappa=\sum_{\text{other branches}}I
                         \ge r_B-1\ge2.
\]
Combining $b\ge \kappa$ with \eqref{eq:one-junction-spectral-bound} forces
\begin{equation}
     b=\kappa=2,\quad r_B=3,\quad
     g(H)=1,\quad \delta_H=n=0.                         \label{eq:one-junction-exponent-equality}
\end{equation}
Thus $\Gamma_H$ is smooth and disjoint from
$\Gamma_R$, and there is no further common
zero of the cubic.

\paragraph{The normalization at the ramified endpoint.}
The cover $H\to A$ ramifies at $p$ by \eqref{eq:one-junction-graph-equality}.
Equation~\eqref{eq:one-junction-quadratic-factor} presents its spectral cover as
$w^2+g(0,z)=0$. After translating its double
root, smoothness from \eqref{eq:one-junction-exponent-equality} gives
$g_z(0,0)\ne0$. Hence
\begin{equation}
 \begin{split}
        N_H&=\{w^2+g(u,z)=0\},\\
        (u,w,z)&\longmapsto(u,v=uw,z)
 \end{split}                                                   \label{eq:one-junction-normalization}
\end{equation}
is the normalization furnished by the last clause of
\cref{lem:II:smooth-curve-intersection}, applied with
$(X,Y,k)=(u,v,g)$. The hypothesis $g_z(0,0)\ne0$
makes its source smooth. The map is finite and birational
onto \eqref{eq:one-junction-quadratic-factor}, with inverse
$w=v/u$ off $u=0$. Coordinates $(u,w)$ solve
$z=\psi(u,w)$, with $z|_H$ of order two in
$w$. The minimal normalization resolution
has no exceptional curve over this smooth
point, and $H$ is $u=0$.

\item The other conductor branch $W_1$ meets $H$
transversely by \cref{lem:II:irrational-ruled-boundary}
on the same smooth neighbourhood. Thus
$\operatorname{ord}_{W_1}(u)=1$.
By \eqref{eq:one-junction-graph-equality}, $W_1$ has
only one point over the normalized $B_+$ point. Its
local degree is therefore its global degree $d_1$.
For a normalized parameter $t_+$ of $B_+$, pullback
of functions along this finite map gives
\[
 1=\operatorname{ord}_{W_1}(u)
   =d_1\operatorname{ord}_{t_+}(u|_{B_+}).
\]
Both factors are positive integers, so $d_1=1$ and
$u$ has order one on the normalized image.
On the smooth branch $f_R=0$, its normalized conductor
curve over $B_-$ is the same smooth curve $z=0$,
mapped by its ambient inclusion. Its local degree is
one. The unique point of $W_c$ over that normalized
$B_-$ point therefore gives $d_c=1$ as well.
Since $\sum d_i=r_B=3$, the case $c=2$
would have sum two, while $c>3$ is impossible.
Therefore
\begin{equation}
                 c=3,\qquad d_1=d_2=d_3=1.              \label{eq:one-junction-degree-vector}
\end{equation}
There are exactly two internal factors, and
$e_1=e_2=1$.

\paragraph{The conductor mismatch.}
Each $f_j$ intersects $f_H$ only along $B_+$.
Its pullback to the smooth normalization
\eqref{eq:one-junction-normalization} consequently has divisor supported
only on $W_1$. Its restriction to $H$ vanishes
to order two: the order on $A$ is one and
$H\to A$ has ramification index two.
Transversality of $H$ and $W_1$ makes the
divisor exactly $2W_1$. Thus the generic
intersection multiplicity of $f_H$ with
each internal factor along $B_+$ is two.

\item The global marking in
\eqref{eq:one-junction-conductor-path} is
\[
 \begin{array}{c|c|c}
 &\text{local factors}&\text{global components}\\ \hline
 B_+&(f_H,f_1,f_2)&(W_1,W_2,W_3)\\
 B_-&(f_1,f_2,f_R)&(W_1,W_2,W_3).
 \end{array}
\]
At $B_-$, the smooth $R$ branch has conductor exponent
two and two other generic branches. The formula of
\cref{lem:II:plane-branch-conductor} makes its two
intersection multiplicities equal to one, so both other
generic branches are smooth as well. The degree-one maps
$W_i\to B^\nu$ are isomorphisms and identify their
function fields with $\C(B)$. They identify these same
generic transverse branches on the two signs, so their
intrinsic conductor exponents are zero on both.
If $p_+,p_-\ge1$ are the intersections
of the two internal factors, the conductor vectors are
\[
         (4,2+p_+,2+p_+),\qquad
         (1+p_-,1+p_-,2).
\]
The first vector uses the already proved intersection
two of each internal factor with $f_H$; the second
uses its intersection one with $f_R$. The third entry
is the coefficient on the same global component $W_3$:
it equals $2+p_+\ge3$ over $B_+$ and equals two over
$B_-$. Both would be restrictions of the single
divisorial valuation of the conductor ideal along $W_3$,
which is impossible.

\paragraph{The case $c=1$.}
Now the conductor graph has $H,R,W$ and
budget two. The points of $W$ over $B_+$
and $B_-$ cannot be identified in $N$,
since that creates a loop in the conductor
tree. Points carrying the two labels are
distinct; their supports intersect, so
both carry $A$. Each contributes at least
one, exhausting the graph budget. Each
has exactly one $A$ and one $B$ branch,
and no other identification occurs.

\item If $H$ ramifies over $p$, these two
$A$ branches are $H$ and $R$. There are
only the factors with supports
$\{A,+\}$ and $\{A,-\}$. Apply \cref{lem:II:smooth-curve-intersection} to
$(f_R,f_H,A)$ with $r=2$ and normal order two of $f_H$.
It makes $f_R$ smooth.
No other factor contains $B_-$, and this
smooth branch has no singularity there.
Thus $B_-$ cannot be a critical curve.

\item If $H$ is unramified, the three normalized
$A$ preimages are distinct. Two occur
at the $B$ attachments and the third
is an $A$-only normalization point.
There are three factors, each supplying
one generic $D_4$ branch. Every pair
meets only along $A$, with multiplicity
one. Apply \cref{lem:II:smooth-curve-intersection} to each
pair, with $\Gamma=A$ and $r=1$. Its coordinate map
$(\phi,\psi,z)$ makes all three factors smooth near
$p$, with their pairwise intersection equal to $A$. Neither $B_+$ nor $B_-$ can be a
critical branch. This final contradiction
completes the proof.

\end{enumerate}
\end{proof}

\subsection{Projective torus fibrations and Hopf bundles}\label{subsec:II:torus-and-hopf-examples}

\begin{proposition}[A projective torus fibration with Chern number twenty-four]\hypertarget{tag-02WH}{}\label{night:P4.10}\label{stmt:P4.10}
\StatementRecord{night:P4.10}{proved}{02WH}{the definitions and constructions stated in the proof}

There is a smooth projective threefold $Z$ with $b_1(Z)=0$ and a
connected torus fibration $Z\to\Pp^1$ whose reduced special fibres
are irreducible normal-crossings surfaces, but $c_3[Z]=24$.
This example has $b_2(Z)>0$ and the fibre types of
Proposition~\ref{night:P4.7}; its value $c_3[Z]=24$ shows that the
bound in that proposition depends on the hypothesis $b_2=0$.
\end{proposition}
\begin{proof}

\begin{enumerate}[label=\arabic*.,leftmargin=*]
\item On the two affine charts of $\Pp^1$, with $s=t^{-1}$, we take
\[
 y^2z=x^3-3xz^2+t^6z^3,\qquad
 y_\infty^2z=x_\infty^3-3s^4x_\infty z^2+z^3,
\]
glued by $x_\infty=s^2x$ and $y_\infty=s^3y$. These are the
Weierstrass cubics in $\Pp(\OO\oplus\OO(2)\oplus\OO(3))$ for
the line bundle $\OO_{\Pp^1}(1)$. They define an elliptic surface
$\pi:S\to\Pp^1$ with zero section $[0:1:0]$. The discriminant
is a nonzero constant times $t^{12}-4$, or $1-4s^{12}$ on the
other chart. There are twelve simple zeros and a smooth fibre at
infinity. At a discriminant zero the cubic has a double root
$x=1$ or $x=-1$ and a distinct third root, so the fibre has an
ordinary node. Since $6t^5\ne0$ there, the total surface is smooth.
At the zero section the derivative in $z$ is nonzero. The map near
each node is therefore $xy=t$ after holomorphic coordinate changes
and translation of the base.

\item The relative dualizing differential $dx/y$ has transitions
$\pi^*\OO(1)$ and generates the dualizing sheaf also at the nodes.
Thus $\omega_S=\pi^*\OO(-1)$. Relative duality gives
$R^1\pi_*\OO_S=\OO(-1)$, and connected fibres give
$\pi_*\OO_S=\OO$. Leray yields $h^1(\OO_S)=0$, while all
positive plurigenera vanish. Since $S$ is projective, its degree-one
Hodge decomposition gives $b_1(S)=2h^1(\OO_S)=0$.
Apply \cite[Chapter~VI, Theorem~1.1 and Table~10, pp.~243--244]{BHPV2004E}
to its minimal model $S_0$. The contraction $S\to S_0$ preserves
the plurigenera and $b_1$, so $\kappa(S_0)=-\infty$ and
$b_1(S_0)=0$. Table~10 identifies $S_0$ as rational; the
birational contraction then makes $S$ rational as well.
In particular $\pi_1(S)=1$, since the minimal rational models
$\Pp^2$ and the Hirzebruch surfaces are simply connected and
each point-blowup map $\pi_Z:S'\to S$ induces an isomorphism
on fundamental groups by \cref{lem:II:blowup-topology}, applied
to that smooth surface and its point centre $Z$ of codimension two.

\item The projective fibre product $W=S\times_{\Pp^1}S$ has exactly
twelve singularities, at the pair of nodes in each nodal fibre.
Their equations are $xy=uv$, with base function $xy=uv$. We blow
up the coherent ideal of the diagonal $\Delta_S\subset W$. This
divisor is Cartier away from the twelve points, so the blowup is
the identity there. Put $r=x-u$, $h=y-v$ at a singular point.
The equation becomes $ry+uh=0$ and the diagonal ideal is $(r,h)$.
The blowup has charts
\label{layout:II:1}
\[
\begin{aligned}
h&=ar, & y&=-au, & x&=u+r, & f&=-aux,\\
r&=bh, & u&=-by, & x&=-bv, & f&=-byv.
\end{aligned}
\]
The independent coordinates are $(a,u,x)$ and $(b,y,v)$.
Both charts are smooth, with $b=a^{-1}$ on the overlap, and the
exceptional fibre is $\Pp^1$. Thus the resulting $Z$ is smooth
and projective. The displayed equations make every special fibre
reduced and normal crossings. Its support remains irreducible:
the original product of two irreducible nodal cubics is irreducible,
and the exceptional locus has dimension one, so introduces no
surface component. Its normalization is rational, by the
birational map from $\Pp^1\times\Pp^1$.

\item A nodal cubic has Euler number $2-1=1$, so its self-product has
Euler number one. The resolution replaces one point by $\Pp^1$;
each resolved special support therefore has Euler number two.
All other fibres are $E_t\times E_t$ with Euler number zero.
Hence $e(Z)=12\cdot2=24=c_3[Z]$.

\item We compute $\pi_1(Z)$. Every loop can be moved
off the fibre divisors, so the fundamental group of their complement
surjects onto $\pi_1(Z)$. The complement is a bundle with fibre
$E\times E$ and a section over a twelve-punctured sphere. The
section extends over $\Pp^1$, killing the base generators. The
two maps
\[
 S\longrightarrow Z:\quad p\mapsto(p,O(\pi(p))),\qquad
 p\mapsto(O(\pi(p)),p)
\]
are defined because their images miss every conifold point, and
extend the two factor inclusions of $E$. Since $S$ is simply
connected they kill both fibre fundamental groups. Thus
$\pi_1(Z)=1$ and $b_1(Z)=0$. A projective Kähler class is nonzero
in $H^2(Z;\R)$, because its cube has positive integral. Thus $b_2(Z)>0$, as asserted.

\end{enumerate}
\end{proof}

\begin{proposition}[A primary Hopf bundle on a product of spheres]\hypertarget{tag-02WI}{}\label{night:P4.11}\label{stmt:P4.11}
\StatementRecord{night:P4.11}{proved}{02WI}{\SeriesRef{night:P4.7}}

There is a smooth compact complex threefold $C$ with
$b_1(C)=b_2(C)=0$, $e(C)=0$, and a holomorphic submersion
$C\to\Pp^1$ whose fibres are primary Hopf surfaces. Every fibre has fundamental group $\Z$ and therefore has no
finite unramified torus cover. Thus the torus-fibre hypothesis
in the smooth-support equality description of
\cref{cdp:st:C.2.G.1} is necessary even under $b_1=b_2=0$.
The numerical normal-crossings sign statement of
\cref{night:P4.7} applies to this fibration.
\end{proposition}
\begin{proof}

\begin{enumerate}[label=\arabic*.,leftmargin=*]
\item The additive group $\C$ acts on
$V=(\C^2\setminus\{0\})\times(\C^2\setminus\{0\})$ by
\[
 \tau\cdot(z,w)=(e^\tau z,e^{i\tau}w).
\]
For $\tau=a+ib$, the logarithms of the two Euclidean norms change
by $a$ and $-b$. The action is free. It is proper: when both a
point and its translate lie in prescribed compact subsets of $V$,
the two norm ratios bound $a$ and $b$. Each orbit has a unique
representative with both norms one, obtained with
$a=-\log\|z\|$ and $b=\log\|w\|$. This gives an explicit
diffeomorphism $C=V/\C\cong S^3\times S^3$ and proves compactness.
Holomorphic transverse slices, for example $z_j=1$ after choosing
a local logarithm, give quotient charts with holomorphic
transitions. Hence $C$ is a smooth compact complex threefold.
The product cell complex gives free integral homology of ranks
$(1,0,0,2,0,0,1)$ and therefore the stated Betti and Euler numbers.

\item The invariant map $[z,w]\mapsto[z_1:z_2]$ is a holomorphic
submersion to $\Pp^1$. In an affine base chart write
$z=\lambda(1,\xi)$ and set $\lambda=1$ using the action.
The choices of logarithm differ by $2\pi ik$, changing $w$ by
$e^{-2\pi k}$. Thus we obtain the explicit local fibre bundle
with fibre
$H=(\C^2\setminus\{0\})/\langle w\mapsto e^{-2\pi}w\rangle$.
Its universal cover retracts onto the simply connected $S^3$, so
$\pi_1(H)=\Z$. A connected finite unramified torus cover would
inject $\Z^4$ as a finite-index subgroup of $\Z$. Every subgroup
of $\Z$ is cyclic, whereas $\Z^4$ has rank four, which is
impossible.

\end{enumerate}
\end{proof}

\subsection{The irreducible classification and sign theorem}

\begin{proposition}[Classification and Euler numbers of torus-fibre supports]\hypertarget{tag-02XX}{}\label{night:P4.58}\label{stmt:P4.58}
\hypertarget{tag-02Y1}{}\label{night:P4.9}\label{stmt:P4.9}
\StatementRecord{night:P4.9}{proved}{02Y1}{\cref{night:P4.58}}

\StatementRecord{night:P4.58}{proved}{02XX}{\SeriesRef{night:P4.26}, \SeriesRef{night:P4.32}, \SeriesRef{night:P4.35}, \SeriesRef{night:P4.36}, \SeriesRef{night:P4.37}, \SeriesRef{night:P4.54}, \SeriesRef{night:P4.55}, \SeriesRef{night:P4.56}, \SeriesRef{night:P4.57}, \SeriesRef{night:P4.6}, \SeriesRef{night:P4.7}, \SeriesRef{night:P4.8}, \SeriesRef{night:P4R.13}, \SeriesRef{night:P4R.18}, \SeriesRef{night:P4R.20}}

Let $f:Y\to B$ satisfy the standing torus-fibration
hypotheses of \cref{def:II:standing-torus-fibration}. Suppose every reduced fibre support
is irreducible or has normal crossings. Then
\[
 B=\mathbb P^1,\qquad c_3(Y)\in\{0,2\},
                       \qquad b_3(Y)=2-c_3(Y).
\]
Every reduced support has nonnegative Euler number.
An irreducible singular support whose normalization has
a rational resolution has Euler number two. Every other
irreducible support has Euler number zero. Such a support
is either smooth with a finite unramified torus cover,
or singular with smooth minimal elliptic ruled
normalization, or singular with smooth minimal Hopf
normalization. In particular the class-VII normalization
alternative is realized only by a smooth minimal Hopf
surface.

For a normal-crossing support the Euler number equals
its number of triple points, and is even. Its zero case
is exactly \cref{night:P4.8}: in the reducible case the normalization
is a disjoint union of smooth minimal Hopf or non-K\"ahler
elliptic surfaces, with smooth elliptic conductor curves
and finite unramified branch maps. The equality
$c_3(Y)=0$ holds exactly when every support has one of
the corresponding zero-Euler forms.
\end{proposition}
\begin{proof}

\begin{enumerate}[label=\arabic*.,leftmargin=*]
\item \Cref{night:P4.26} proves the sign for every irreducible support,
including singular normalizations, and \cref{night:P4.6,night:P4.7,night:P4.8} give
the full normal-crossing calculation. \Cref{night:P4.35} proves the
exact value two for the rationally resolved singular
case. \Cref{night:P4.36,night:P4.37} show that every class-VII case has
smooth minimal Hopf normalization and Euler number
zero. \Cref{night:P4.32} gives the complete irreducible zero-Euler
classification and its converses.

\item For a positive-Euler support with positive-genus ruled
normalization resolution, apply
\cref{lem:II:positive-ruled-conductor-graph} to this
$f:Y\to B$ and full fibre $mS$. Every reduced support
is irreducible or normal crossing by hypothesis. The lemma
gives a doubled horizontal pole section, two rational
normalized conductor components and one graph cycle.
\Cref{night:P4.54} excludes contraction of that section, using the weighted normalization and blowdown maps.
\Cref{night:P4R.18} reduces its generic image to $D_4$ with degree
two or three. \Cref{night:P4.55} makes that image smooth and excludes
degree three, leaving the degree-two section and one
rational vertical third sheet. \Cref{night:P4.56} forces at least
three generic branches on the other image component.

\item The graph alternatives are exhaustive by \cref{night:P4.55}. \Cref{night:P4.57}
excludes two distinct junctions. \Cref{night:P4R.20} excludes a sole
three-branch image point, with the exact conductor-path
and endpoint normalization calculation. Thus no
positive-Euler ruled support remains. This proves the
asserted value zero for every nonrational irreducible
support, and \cref{night:P4.32} supplies its normalization.

\item Finally Euler additivity gives $c_3(Y)$ as the sum of
these nonnegative support Euler numbers. The base,
Poincar\'e-duality calculation $e(Y)=2-b_3(Y)$, and
parity are proved in \cref{night:P4.26}. They yield the displayed
global values. A sum of nonnegative integers vanishes
exactly when every summand vanishes, giving the equality
classification through \cref{night:P4.32} and \cref{night:P4.8}.

\end{enumerate}
\end{proof}



\part{Periods, tangent cohomology and deformations}\hypertarget{tag-P271}{}\label{part:II:periods-tangent-deformations}
The marked periods determine the Picard resonance loci and
Kodaira--Spencer maps of the constructed torus fibrations.
The fibrewise cup products and normalization maps relate
their coherent cohomology to fibration stability and local
deformation spaces. The source, period and affine attachment
conventions are those of \cref{def:II:marked-conventions}.

\section{Twisted tangent cohomology of the marked families}\hypertarget{tag-0348}{}\label{sec:II:12}

\subsection{Hexagonal fibres and normalization cohomology}

\begin{proposition}[Hexagonal fibres and tangent cohomology in three families]\hypertarget{tag-01FQ}{}\label{cdp:st:C4g}
\StatementRecord{cdp:st:C4g}{proved}{01FQ}{\SeriesRef{cdp:st:A.3.4}, \SeriesRef{cdp:st:A.4.1}, \SeriesRef{cdp:st:C.1.3}, \SeriesRef{cdp:st:D.1.9}; \cite[Theorem~1.1 and Proposition~3.1]{Viaclovsky2026Two}; \cite[Proposition~2.5, Lemma~3.10 and Theorem~5.1]{LiuRollenske2026}}

For every admissible member of the original $(3,4)$ family,
every member of the two-parameter family of
\cite[Theorem~1.1, p.~2]{Viaclovsky2026Two}, and every member of
the Liu--Rollenske family of
\cite[Theorem~5.1, p.~21]{LiuRollenske2026}, there is a reduced
full hexagonal fibre $W$ with normalization
$\eta:\mathrm dP_6\to W$ such that
\begin{equation}
 \eta^*L=\OO_{\mathrm dP_6}\quad(L\in\Pic(Y)),\qquad
 H^2\bigl(\mathrm dP_6,K_{\mathrm dP_6}\otimes
                    \eta^*(K_Y^{-1}\otimes L)\bigr)=\C.
 \hypertarget{tag-012O}{}\label{cdp:eq:C.TEST.1}
\end{equation}
In these three families the singular-fibre Euler contributions are
\begin{equation}
\begin{array}{c|c|c}
\text{family}&\text{reduced singular supports}&
                 \text{their Euler numbers}\\ \hline
(3,4)&W&2\\
\text{two-parameter family}&W,\ W_1,\ W_2&2,0,0\\
\text{Liu--Rollenske family}&W,\ C_{\rm cusp}&2,0 .
\end{array}
\hypertarget{tag-012P}{}\label{cdp:eq:C.TEST.2}
\end{equation}
The other critical full fibres are smooth multiple fibres,
with multiplicities $(3,4)$, $4$, and $4$, respectively.
For every such full multiple fibre $mS$,
$H^0(mS,(\Omega_Y^1\otimes L)|_{mS})=0$ on a dense open
subset of every component of $\Pic(Y)$, including all nilpotent
layers of $mS$. The hexagonal normalization has the nonzero
cohomology in \eqref{cdp:eq:C.TEST.1}, and
$H^2(Y,T_Y\otimes L)\ne0$ for every ambient line bundle
in each family.
\end{proposition}


\begin{proof}

\begin{enumerate}[label=\arabic*.,leftmargin=*]
\item The original local filling is the hexagonal toric quotient
in \cref{cdp:st:A.4.1}. For the two-parameter family,
\cite[Proposition~3.1 and its proof, pp.~6--7]{Viaclovsky2026Two}
gives the hexagonal filling at the zero of the linearization
and the two rank-one fillings at the other nodal values.
For the Liu--Rollenske family,
\cite[Proposition~2.5, p.~5]{LiuRollenske2026} identifies the
hexagonal fibre by gluing opposite boundary curves of
$\mathrm dP_6$. Thus the normalization map in all three cases
is the indicated opposite-side quotient of that surface.

\item For every member in question, $H^2(Y,\Z)$ is a finite
group. For the original affine family,
the integral homology is given by
\cite[Theorems~7.17 and~7.22, pp.~57 and~60]{AC26};
the universal coefficient sequence gives
\[
 H^2(Y,\Z)\simeq\operatorname{Ext}^1(H_1(Y,\Z),\Z)
 \simeq\Z/d\Z,\qquad d=|p|,
\]
because $H_2(Y,\Z)$ is torsion. In the other two families, the six-sphere topology is
\cite[Theorem~1.1, p.~2]{Viaclovsky2026Two} and
\cite[Theorem~5.1, p.~21]{LiuRollenske2026}. In
particular every ambient first Chern class is torsion.
The group $H^2(\mathrm dP_6,\Z)$ is free, so its pullback
to $\mathrm dP_6$ is zero. The exponential sequence there,
together with $H^1(\mathrm dP_6,\OO)=0$, shows that the
pulled-back line bundle is trivial. This applies to both
$L$ and $K_Y^{-1}\otimes L$, and Serre duality on
$\mathrm dP_6$ proves \eqref{cdp:eq:C.TEST.1}.

\item All these threefolds have $b_1=b_2=b_3=0$: use \cite[Theorem~7.22, p.~60]{AC26} for
every original affine member, and the integral homology
of the six-sphere for the other two families. Any full
fibre is a nonzero effective Cartier divisor. Statement
\ref{cdp:st:C.1.3} therefore gives
\[
 h^2(Y,T_Y\otimes L)\geq 1
 \qquad\text{for every }L\in\Pic(Y).
\]
This argument applies equally when $d>1$ and the integral
cohomology has nonzero torsion.

\item The toric strata of $W$ consist of one $(\C^*)^2$ stratum,
three $\C^*$ strata and two triple points. Their compactly
supported Euler numbers give $e(W)=2$.
Each rank-one normalization is a $\PP^1$-bundle over an
elliptic curve, with its zero and infinity sections glued.
Both the normalization and these conductor curves have
Euler number zero, so normalization additivity gives
$e(W_i)=0$. At the Liu--Rollenske cusp,
\cite[Lemma~3.10 and its proof, pp.~11--12]{LiuRollenske2026}
gives a homeomorphism over the base disc from the contraction
quotient to the cuspidal elliptic-surface neighbourhood times
$(S^1)^2$. It is obtained by trivializing the mapping torus
through the displayed isotopy of the translation and circle
multiplier. Restriction over the cusp identifies
$C_{\rm cusp}$ with the cuspidal rational curve times $(S^1)^2$.
Euler multiplicativity gives $e(C_{\rm cusp})=0$.

\item The fourth-root covers are
\cite[Lemma~3.2 and \S3.3, pp.~7--8]{Viaclovsky2026Two} and
\cite[Lemma~2.7 and Corollary~2.8, pp.~6--7]{LiuRollenske2026};
the original root charts have orders three and four.
Their reduced central fibres are smooth finite unramified
torus quotients, and their Euler numbers are zero.
The full-fibre group in \ref{cdp:st:A.3.4} injects under
pullback by the finite \'etale quotient $\pi:\widetilde U\to U$.
Write $t=u^m$ on this cover, with central torus $T$.
Each Picard component contains a finite-order line bundle $A$:
for a representative $A_0$, choose $N>0$ killing its first Chern class,
write $A_0^N=\exp(\alpha)$, and take $A=A_0\exp(-\alpha/N)$.
For the transgression class $\beta$ of \ref{cdp:st:D.1.9}, put
$A_r=A\exp(2\pi i r\beta)$ with $r$ real and irrational.
The nonzero integral transgression character stays nonzero on the
finite torus cover, since transfer composed with pullback is its
degree. Hence $(\pi^*A_r)|_T$ is a nontrivial unitary line bundle.
The conormal sequence is
\[
 0\longrightarrow\OO_T\,du
 \longrightarrow\Omega^1_{\widetilde U}|_T
 \longrightarrow\OO_T^{\oplus2}\longrightarrow0.
\]
Tensoring it by $(\pi^*A_r)|_T$ gives zero in degree zero.
Every conormal power of $T$ is trivialized by a power of $u$;
induction through the $u$-adic filtration therefore gives
$H^0(mT,(\Omega^1_{\widetilde U}\otimes\pi^*A_r)|_{mT})=0$.
Faithfully flat pullback gives the same vanishing on the full fibre
downstairs. Proper semicontinuity, intersected over the finitely many
multiple values, gives the simultaneous dense open set in every
Picard component. Euler localization
therefore gives $c_3(Y)=2$ in all three rows.

\end{enumerate}
\end{proof}




\begin{proposition}[The normalization quotient and its cohomology map]\hypertarget{tag-01FR}{}\label{cdp:st:C5a}
\StatementRecord{cdp:st:C5a}{proved}{01FR}{\citet[Lemma~6.11, p.~81]{DGKXv1}}

Let $Y$ be a smooth compact complex threefold, $S\subset Y$ a reduced normal-crossing divisor, and $\nu:\widetilde S\to Y$ the normalization followed by inclusion. This is a finite holomorphic immersion on each local branch. Put $\mathcal T=T_Y(-\log S)$. There is a canonical exact sequence of coherent $\OO_Y$-modules
\begin{equation}
0\longrightarrow\mathcal T\longrightarrow T_Y
 \longrightarrow\nu_*N_\nu\longrightarrow0,
\qquad N_\nu=K_{\widetilde S}\otimes\nu^*K_Y^{-1}.      \hypertarget{tag-0094}{}\label{cdp:eq:C.3}
\end{equation}
For every holomorphic line bundle $L$ on $Y$, the image of the natural cohomology map is exactly
\begin{equation}
\begin{split}
\im\bigl(H^2(Y,T_Y\otimes L)&\longrightarrow
 H^2(\widetilde S,N_\nu\otimes\nu^*L)\bigr)\\
 &=\ker\bigl(H^2(\widetilde S,N_\nu\otimes\nu^*L)
 \longrightarrow H^3(Y,\mathcal T\otimes L)\bigr).
\end{split}                                             \hypertarget{tag-0095}{}\label{cdp:eq:C.4}
\end{equation}
The target identifies by Serre duality with
$H^0(\widetilde S,\nu^*(K_Y\otimes L^{-1}))^\vee$.
\end{proposition}




\begin{proof}
On a chart $S=(z_1\cdots z_r=0)$, the morphism is
\[
 \sum_i a_i\partial_{z_i}
   \longmapsto (a_i|_{z_i=0})_{1\le i\le r}.
\]
This is the branch-coefficient morphism in the proof of
\cite[Lemma~6.11, p.~81]{DGKXv1}. Its exactness there uses only
this local coordinate description, which is identical on the
present normal-crossing charts. The normal-line transitions
identify its target with $\nu_*N_\nu$, and immersion adjunction
gives $N_\nu=K_{\widetilde S}\otimes\nu^*K_Y^{-1}$.
Tensor by $L$ and apply coherent cohomology. Exact finite
pushforward identifies the normalization term with its
cohomology on $\widetilde S$, giving \eqref{cdp:eq:C.4}.
Serre duality on that smooth compact surface identifies the
target with the stated dual section space. In particular,
$H^3(Y,\mathcal T\otimes L)=0$ makes the map surjective.
\end{proof}






\subsection{Integral and nonintegral resonance}

\begin{definition}[The marked Picard parameter]\hypertarget{tag-P261}{}\label{def:II:marked-Picard-parameter}
\StatementRecord{def:II:marked-Picard-parameter}{proved}{P261}{}
Let $X=X_c$ be the period construction of \cite[Setup and Main
Theorem, p.~2]{AC26}, with admissible parameter $c$, clockwise
affine vectors $v_1=(1,2,-4,0)^t$, $v_2=(-1,-3,3,0)^t$, and
untwisted cusp collar. Thus $p=12l_0-4l_1-3l_2=-1$.
Write $\mathfrak B=\PP(3,4)$, let $\rho_B:\mathfrak B\to B$ be its
coarse map, and let $\bar f:X\to\mathfrak B$ be the root
factorization. Put $\mathcal T=T_X(-\log W)$ and let
$\nu:\mathrm dP_6\to X$ normalize $W$.
The integral root lines are
\[
 M_n=\bar f^*\OO_{\mathfrak B}(n)
     =\OO_X(n(S_1-S_2)),\qquad n\in\Z.
\]
The exponential isomorphism $H^1(X,\OO_X)\to\Pic(X)$ identifies
$\Pic(X)$ with $\C$. Fix its coordinate by
$M_z=\exp(z\gamma^{0,1})$ and $M_1=\OO_X(S_1-S_2)$.
The period rows are the fixed $\tau,\mu,\beta_c$ of the setup;
their relation to the elliptic-quotient coordinates is described
in \cite[\S1.3, pp.~2--3]{Engel2026}.
\end{definition}

\begin{proof}[Well-definedness of the coordinate]
The clockwise transgression sends $12\gamma$ to $-p=1$.
Indeed, evaluation of the finite lifts on $v_1,v_2$ gives the
integral periods $4l_1,3l_2$, and the cusp contributes
$-12l_0$; their sum is $-p$. The exponential sequence therefore
gives $\exp(12\gamma^{0,1})=f^*\OO_B(1)$. Its injectivity and
the root relation identify $\exp(\gamma^{0,1})$ with
$\OO_X(S_1-S_2)$, fixing the sign used in $M_1$.
\end{proof}

\begin{proposition}[The theta determinant and its exact rank locus for $X_c$]\hypertarget{tag-01GX}{}\label{cdp:st:E8}
\StatementRecord{cdp:st:E8}{proved}{01GX}{the definitions and constructions stated in the proof}

Use the period convention
$\Omega=\left(\begin{smallmatrix}\tau&\mu\\\mu&\beta/6\end{smallmatrix}\right)$,
Jacobi's odd theta function with periods $1,\tau$, and
$\beta/6=(i/\pi)\log(\vartheta_1(\mu,\tau)/\eta(\tau))+\text{constant}$.
Writing $h=\partial_\mu\log\vartheta_1$, one has the identity
\begin{equation}\hypertarget{tag-01A1}{}\label{cdp:eq:E-theta-general}
 \det\frac{d\Omega}{d\tau}
 =-\left(\mu'-\frac{i}{2\pi}h\right)^2
       -\frac{\wp(\mu,\tau)}{4\pi^2}.
\end{equation}
For the AC period branch $\wp(\mu,\tau)=0$, this becomes
\[
 \det\frac{d\Omega}{d\tau}
       =-\frac{g_2(\tau)^2}{36\pi^2g_3(\tau)}.
\]
On the smooth root-family base, the Kodaira--Spencer tensor has rank
one at the order-three point and rank two at every other finite
point, including the order-four point. Its determinant has order
two at the order-three point. The logarithmic tensor at the cusp
has rank two.
\end{proposition}


\begin{proof}

\begin{enumerate}[label=\arabic*.,leftmargin=*]
\item The symmetric matrix has a concrete comparison with the raw AC
periods. Adjoin $(\mu,\beta/6)$ to the lattice generated by
$(1,0),(0,1),(\tau,\mu),(6\mu,\beta)$. This enlarges the
lattice by index six; identity on $\mathbb C^2$ induces a finite
holomorphic isogeny of degree six to the torus with period matrix
$(I,\Omega)$. Its differential and pullback on complex first
cohomology are isomorphisms, so the Kodaira--Spencer ranks agree.
This is the comparison used in the determinant calculation.

\item The convention is
$\vartheta_j(z,\tau)=\theta_j(\pi z\mid\tau)$ and
$q=e^{\pi i\tau}$. The normally convergent theta series and
\cite[eqs.~20.2.1 and 20.13.1]{DLMFTheta2026} give
$\partial_\tau\vartheta_1=(4\pi i)^{-1}\partial_\mu^2\vartheta_1$.
Put $Q=q^2$ and
$E_2=1-24\sum_{n\ge1}nQ^n/(1-Q^n)$.
The normally convergent product
\cite[eq.~23.17.8]{DLMFTheta2026} gives
$\eta'/\eta=\pi i E_2/12$ by termwise logarithmic differentiation.
Likewise \cite[eq.~20.5.1]{DLMFTheta2026} gives
\[
 h(\mu)=\frac1\mu-\frac{\pi^2E_2}{3}\mu+O(\mu^3).
\]
Here the coefficient follows from
$\sum_{n\ge1}Q^n/(1-Q^n)^2
=\sum_{n\ge1}nQ^n/(1-Q^n)$, obtained by comparing the absolutely
convergent double series. The quasi-period equation
\cite[eq.~20.2.6]{DLMFTheta2026} makes $\partial_\mu h$ elliptic.
Consequently $-\partial_\mu h-\pi^2E_2/3$ and $\wp$
have identical poles, principal parts and zero Laurent constant.
Their difference is a holomorphic elliptic function with constant
zero, so they agree. Substituting these three
identities into
$\det\Omega'=(\beta/6)'-\mu'^2$ gives
\eqref{cdp:eq:E-theta-general}. On the branch $\wp(\mu)=0$
its second term is zero.

\item For the zero branch differentiate $\wp(\mu,\tau)=0$. Put
$c=\pi^2E_2/3$ so $h'=-\wp-c$. Differentiating the heat equation
once in $\mu$, using $\wp''=6\wp^2-g_2/2$, and using
$E_2'=(\pi i/6)(E_2^2-E_4)$ from
\cite[eq.~(3.5), p.~8]{AtkinGarvan2002}, yields at $\wp=0$
\[
 \partial_\tau\wp=\frac{h\wp'}{2\pi i}
                       +\frac{i g_2}{6\pi}.
\]
Consequently
\[
 \mu'-\frac{i h}{2\pi}=-\frac{i g_2}{6\pi\wp'},
 \qquad (\wp')^2=-g_3,
\]
which proves the second formula wherever $g_3\ne0$. It is an
identity of meromorphic expressions and therefore extends across
the branch values in a local uniformizing parameter.

\item The equation $j(\tau)=1728t$ in
\cite[Remark~3.12, p.~19]{AC26} shows that $\tau-\rho$ has order one
at the order-three root parameter, while $\tau-i$ has order two
at the order-four root parameter. The modular form $g_2$ has a
simple zero at $\rho$ and $g_3$ a simple zero at $i$, with no
other zeros modulo the modular group. Multiplication by
$(d\tau/ds)^2$ therefore gives an order-two zero at the first
point, and cancels the order-two zero of $g_3(\tau(s))$ at the
second. The tensor is nonzero at the first point because
$d\tau/ds\ne0$ there, hence has rank exactly one. Away from the
two root points the coordinate change is unramified and the
formula has nonzero determinant. At the cusp
$\tau=(\log q)/(2\pi i)+O(q)$, $\mu$ has a finite limit, and
$\beta=-\tau+O(1)$; hence the logarithmic derivative has limiting
matrix $\operatorname{diag}(1/(2\pi i),-1/(12\pi i))$.
Its determinant is nonzero.

\end{enumerate}
\end{proof}

\begin{proposition}[Integral twists and the normalization-map exception]\hypertarget{tag-01FS}{}\label{cdp:st:C5b}
\StatementRecord{cdp:st:C5b}{proved}{01FS}{}

For $n\in\Z$ define
\[
 r_1(n)\equiv n\pmod3,\quad 0\le r_1(n)<3,
 \qquad r_2(n)\equiv3n\pmod4,\quad0\le r_2(n)<4,
\]
\begin{gather*}
 d(n)=\frac{n-4r_1(n)-3r_2(n)}{12},\qquad
 \epsilon_n=\begin{cases}0&3\mid n,\\1&3\nmid n.\end{cases}
\\
 u(n)=\max\{d(n)+1,0\},\qquad
 v(n)=\max\{-d(n)-1,0\}.
\end{gather*}
The coarse direct images of $T_X\otimes M_n$ are
\begin{equation}
\begin{aligned}
 f_*(T_X\otimes M_n)&=\OO(d(n))\oplus\OO(d(n-1)),\\
 R^1f_*(T_X\otimes M_n)&=\OO(d(n))\oplus\OO(d(n-1))\oplus\OO(d(n+2)),\\
 R^2f_*(T_X\otimes M_n)&=\OO(d(n+6))\oplus\C_{p_1}^{\epsilon_n}.
\end{aligned}
\hypertarget{tag-0096}{}\label{cdp:eq:C.5}
\end{equation}
Their global cohomology dimensions are
\begin{equation}
\begin{aligned}
 h^0&=u(n)+u(n-1),\\
 h^1&=v(n)+v(n-1)+u(n)+u(n-1)+u(n+2),\\
 h^2&=v(n)+v(n-1)+v(n+2)+u(n+6)+\epsilon_n,\\
 h^3&=v(n+6).                                             
\end{aligned}
\hypertarget{tag-0097}{}\label{cdp:eq:C.6}
\end{equation}
The natural map $H^2(X,T_X\otimes M_n)\to H^2(\mathrm dP_6,K_{\mathrm dP_6})=\C$ has rank one exactly for
\begin{equation}
 n\in\{-6,-3,-2\}\cup\Z_{\ge0};                         \hypertarget{tag-0098}{}\label{cdp:eq:C.7}
\end{equation}
it is zero for all other integers. In particular $M_{-1}$ has tangent cohomology $(0,0,1,0)$ and a zero map to the normalization target.
\end{proposition}



\begin{proof}

\begin{enumerate}[label=\arabic*.,leftmargin=*]
\item The root-bundle formulas in \SeriesRef{fam:st:A.5.2},
with $k=l=1$, give $E=\OO\oplus\OO(-1)$ and
$G=R^1\bar f_*\OO_X=\OO\oplus\OO(-1)$.
The logarithmic tangent sequence maps a vector field to its
base derivative. On a cusp chart $q=z_1\cdots z_r$, this map sends
$\sum a_i z_i\partial_{z_i}$ to $(\sum a_i)q\partial_q$;
its vertical kernel is generated by the differences
$z_i\partial_{z_i}-z_r\partial_{z_r}$ and the remaining relative
coordinate directions. The integral monomial changes of the toric
charts carry these generators to one another. On the finite torus
covers this is the usual relative tangent sequence. Hence the
vertical inclusion is evaluation of the translation-invariant
fields, and its quotient is the pullback of
$T_{\mathfrak B}(-p_0)=\OO(7-12)=\OO(-5)$.

\item The normalization complex of $W$ has global scalar complex
\[
 \C\xrightarrow{0}\C^3
 \xrightarrow{(a,b,c)\mapsto(a-b+c,a-b+c)}\C^2.
\]
Its dimensions are $(1,2,1)$: the second differential has rank
one. These agree with the ranks in \SeriesRef{fam:st:A.5.2},
which also gives $R^2\bar f_*\OO_X=\OO(-1)$.
The cup-product morphism $\det G\to R^2\bar f_*\OO_X$
is an isomorphism on regular tori. Both lines are $\OO(-1)$,
so this morphism is a nonzero constant and hence an isomorphism
everywhere. Together with the logarithmic tangent sequence this gives
\begin{equation}
E=G=\OO\oplus\OO(-1),\quad R^2\bar f_*\OO_X=\OO(-1),
\quad0\to\bar f^*E\to\mathcal T\to\bar f^*\OO(-5)\to0. \hypertarget{tag-0099}{}\label{cdp:eq:C.8}
\end{equation}


\item Choose root sections $x$ of weight $3$, vanishing at $p_2$, and $y$ of weight $4$, vanishing at $p_1$. The Kodaira--Spencer morphism $\OO(-5)\to E\otimes G$ has, by its weights, the form
\[
 \kappa=\begin{pmatrix}0&a y\\b y&c_0x\end{pmatrix}.
\]
The weight-five section space is zero. The isogeny in
\cref{cdp:st:E8} is induced by the identity on $\mathbb C^2$;
its differential and pullback on first complex cohomology
identify the Kodaira--Spencer tensors. At the cusp that
proposition gives rank two, so
$\det\kappa=-ab\,y^2\ne0$ and $ab\ne0$. At $p_1$ one has
$y=0$ and $x\ne0$; its rank-one conclusion gives
$c_0x\ne0$, hence $c_0\ne0$. In frames $e_0,e_1$ for $E$ and $\eta_0,\eta_1$ for $G$,
the invertible changes
$e'_0=ae_0$, $e'_1=c_0e_1$, $\eta'_0=(b/c_0)\eta_0$,
$\eta'_1=\eta_1$ make the three coefficients one; division by
these nonzero constants gives their inverse. Cup product with
this extension class gives
\[
\delta_0=(0,y,y,x)^t,\qquad
\delta_1=\begin{pmatrix}-y&0\\-x&y\end{pmatrix}:
 \OO(-5)\oplus\OO(-6)\to\OO(-1)\oplus\OO(-2).
\]
For $\delta_0$, the change $(a,b,c)\mapsto(a,b-a,c)$ in
$\OO(-1)^2\oplus\OO(-2)$, with inverse
$(a,d,c)\mapsto(a,a+d,c)$, splits off one $\OO(-1)$.
The remaining quotient is
$(a,c)\mapsto xa-yc$ onto $\OO(2)$, with kernel
$h\mapsto(yh,xh)$. Surjectivity follows from the lifts
$(v/x,0)$ on $x\ne0$ and $(0,-v/y)$ on $y\ne0$.
Thus $\operatorname{coker}\delta_0=\OO\oplus\OO(-1)\oplus\OO(2)$.
The second map is injective with determinant $-y^2$.
On the $x$-unit chart, write $u,v$ for the target generators.
The maps $u\mapsto1$, $v\mapsto-y/x$ identify its cokernel
with $\OO/(y^2)$, whose generator has the character of $\OO(-1)$;
the inverse sends $1$ to $u$.
We denote this root-chart module by $Q=\OO(-1)|_{2p_1}$.
The direct-image sequence of \eqref{cdp:eq:C.8} yields
\[
 R^0\bar f_*\mathcal T=E,\quad
 R^1\bar f_*\mathcal T=\OO\oplus\OO(-1)\oplus\OO(2),\quad
 R^2\bar f_*\mathcal T=\OO(-6)\oplus Q.
\]
The last extension splits because
$\operatorname{Ext}^1(\OO(-6),Q)=H^1(Q(6))=0$.

\item \Cref{cdp:st:C5a} gives the quotient $\nu_*K_{\mathrm dP_6}$. The local cusp volume form trivializes $\nu^*K_X$, and every $\nu^*M_z$ is trivial because $H^2(X,\Z)=0$ and $\Pic^0(\mathrm dP_6)=0$. Hence this quotient has cohomology only in degree two, where it is $\C$. The untwisted direct image in \SeriesRef{fam:st:A.5.2}
is locally free of rank one at $p_0$, where the root and coarse
charts coincide. Thus the extension of $\C_{p_0}$ by
$\OO(-6)$ is nonzero. Its middle term is
$\OO(-6)(p_0)=\OO(6)$: in a local coordinate $q$ its
inclusion is multiplication by $q$, with quotient $\C$.
The summand $Q$ splits off because its support is disjoint from
$p_0$. Therefore
\begin{equation}
 R^0\bar f_*T_X=E,\quad
 R^1\bar f_*T_X=\OO\oplus\OO(-1)\oplus\OO(2),\quad
 R^2\bar f_*T_X=\OO(6)\oplus Q.                       \hypertarget{tag-009A}{}\label{cdp:eq:C.9}
\end{equation}

\item For a root line, multiplication by
$x^{r_2(n)}y^{r_1(n)}$ identifies a polynomial of coarse degree
$d(n)$ in $X=x^4,Y=y^3$ with a weighted section of degree $n$.
Every such monomial has uniquely
$(a,b)=(r_2(n)+4i,r_1(n)+3j)$ with $i+j=d(n)$;
division by the displayed factor gives the inverse.
The same formula on the two root charts identifies the invariant
local frames and gives $\rho_{B*}\OO(n)=\OO_B(d(n))$. Averaging the finite stabilizers is exact, so there are no higher direct images of $\rho_B$. Its sections are equivalently monomials $x^ay^b$ with $3a+4b=n$, and Serre duality gives $h^1\OO(n)=h^0\OO(-n-7)$. At $p_1$, the invariant exponent of $Q(n)$ is $n-1$ modulo three. Its basis consists of $y^ju$ with $j=0,1$; invariance means
$j\equiv n-1\pmod3$. There is one basis vector for
$n\equiv1,2\pmod3$ and none for $3\mid n$.
The coarse parameter $y^3$ annihilates both vectors, so the
coarse length is $\epsilon_n$. Apply projection formula to \eqref{cdp:eq:C.9} and then $\rho_{B*}$. This proves \eqref{cdp:eq:C.5}, including the support and length of every torsion summand. Leray over $B$ has columns only zero and one, proving \eqref{cdp:eq:C.6}.

\item The normalization map on $R^2$ is the quotient in
\[
0\longrightarrow\OO(n-6)\longrightarrow\OO(n+6)
 \longrightarrow\C_{p_0}\longrightarrow0.
\]
Its global-section map is nonzero precisely when a section of $\OO(n+6)$ is nonzero at $p_0$. Both $x,y$ are nonzero at that point, so this occurs precisely when $n+6$ lies in the additive semigroup generated by $3,4$. Below six the semigroup consists of $0,3,4$; the seeds
$6,7,8$ and addition by $3$ give every integer at least six.
This proves \eqref{cdp:eq:C.7}. The $H^1(R^1)$ part maps to zero in the cohomology of a sheaf supported at one point, by the natural Leray filtration. At $n=-1$, \eqref{cdp:eq:C.5} has only one nonzero cohomological contribution, the reduced $\C_{p_1}$ summand in degree two, while $H^3(X,\mathcal T\otimes M_{-1})=H^1(\OO(-7))=\C$. The boundary from the normalization target is an isomorphism,
so its preceding cohomology map is zero.

\end{enumerate}
\end{proof}






\begin{proposition}[Nonintegral twists and their resonance lengths]\hypertarget{tag-01FT}{}\label{cdp:st:C5c}
\StatementRecord{cdp:st:C5c}{proved}{01FT}{}

Use the fixed marked functions $\tau,\mu,\beta_c$ and the
marked lifts $z_1,z_2$ of \cref{def:II:marked-conventions}; put
$\rho=\tau(z_1)=e^{\pi i/3}$. For $z\notin\Z$, let $\mathcal R_z$
be the set of regular base points at which integers $(n,p,q,r)$ satisfy
\begin{equation}
 q\tau+r\mu=p,\qquad z=n-6q\mu-r\beta_c.               \hypertarget{tag-008O}{}\label{cdp:eq:C.10}
\end{equation}
Let
\[
 \Lambda_1=\Z+(3\beta_c(z_1)+4-2\rho)\Z,\quad
 \Lambda_2=\Z+(2\beta_c(z_2)+3-3i)\Z.
\]
Set $\ell_1(z)=1+\epsilon_n$ for $z=n+k(3\beta_c(z_1)+4-2\rho)$ with $k\ne0$, and zero off $\Lambda_1$. Set $\ell_2(z)=1$ for $z=n+k(2\beta_c(z_2)+3-3i)$ with $k\ne0$ and $n\equiv1,2\pmod4$, and zero otherwise. Then $\mathcal R_z$ is finite. With $N_z=|\mathcal R_z|$ and $J_z=2N_z+\ell_1(z)+\ell_2(z)$, we have
\begin{equation}
 (h^0,h^1,h^2,h^3)(X,T_X\otimes M_z)=(0,J_z,1+J_z,0),  \hypertarget{tag-008P}{}\label{cdp:eq:C.11}
\end{equation}
and, writing $\mathfrak m_t$ for the maximal ideal of $\OO_{B,t}$,
\begin{equation}
\begin{aligned}
R^1f_*(T_X\otimes M_z)&=
 \bigoplus_{t\in\mathcal R_z}\C_t^2
 \oplus\C_{p_1}^{\ell_1(z)}\oplus\C_{p_2}^{\ell_2(z)},\\
R^2f_*(T_X\otimes M_z)&=
 \C_{p_0}\oplus\bigoplus_{t\in\mathcal R_z}\OO_{B,t}/\mathfrak m_t^2
 \oplus\C_{p_1}^{\ell_1(z)}\oplus\C_{p_2}^{\ell_2(z)}.
\end{aligned}
\hypertarget{tag-008Q}{}\label{cdp:eq:C.12}
\end{equation}
The natural map to $H^2(\mathrm dP_6,K_{\mathrm dP_6})$ is surjective for all these nonintegral twists.
\end{proposition}



\begin{proof}

\begin{enumerate}[label=\arabic*.,leftmargin=*]
\item On a marked smooth torus, the exponential sequence presents $\Pic^0$ as constant complex lattice cocycles modulo the integral cocycles and the two holomorphic period rows. The equation for $z\gamma$ to belong to this sum is
\[
 z\gamma=(n,p,q,r)+a(6\mu,\tau,1,0)+b(\beta_c,\mu,0,1).
\]
The last two coordinates give $a=-q,b=-r$, and the first two give \eqref{cdp:eq:C.10}. If $r=0$, positivity of $\operatorname{Im}\tau$ forces $p=q=0$ and $z=n$. Otherwise
\[
 \operatorname{Im}z=-r\left(\operatorname{Im}\beta_c-
           6(\operatorname{Im}\mu)^2/\operatorname{Im}\tau\right)\ne0.
\]
In particular real nonintegral twists have no regular resonance. The fixed values $\mu(z_1)=(2-\rho)/3$ and $\mu(z_2)=(1-i)/2$ of \SeriesRef{fam:src:C.1.1} give the two displayed lattices by solving the first equation. Their nonreal generators make each representation $n+k\lambda_j$ unique.

\item We give a coherent local complex to calculate every length. On a smooth simply connected torus-family disc use normalized period coordinates $\zeta=Z(s)x+y$, $x,y\in\R^2/\Z^2$, and translation-invariant forms $G_j=(dx_j)^{0,1}$. The invariant horizontal lift is $h=\partial_s+\sum_i(Z'(s)x)_i\partial_{\zeta_i}$. Since $dx_j$ is closed, $\bar\partial G_j=0$ and
$\bar\partial h=\sum Z'_{ij}e_iG_j$. After subtracting the integer cocycle of a resonance, the line multiplier has invariant Dolbeault form $a_1G_1+a_2G_2$, up to the common nonzero factor $2\pi i$. Thus the invariant relative complex has terms
\[
 C^j=(V\oplus\OO h)\otimes\Lambda^jG,\qquad j=0,1,2,
\]
\begin{equation}
 d(v\eta)=v(a\wedge\eta),\qquad
 d(h\eta)=h(a\wedge\eta)+\kappa\wedge\eta,             \hypertarget{tag-008R}{}\label{cdp:eq:C.13}
\end{equation}
where
\[
 a=(z-n+6q\mu+r\beta_c,-p+q\tau+r\mu),\qquad
 \kappa=\begin{pmatrix}6\mu'&\tau'\\\beta_c'&\mu'\end{pmatrix}.
\]
This complex computes the derived direct image. To verify the comparison, decompose relative smooth forms into Fourier characters under real torus translations. The zero character gives \eqref{cdp:eq:C.13}. For every other character the relative Dolbeault differential is exterior multiplication by a nonzero covector; contraction by its normalized dual is a homotopy. On a sufficiently small closed subdisc those covectors are bounded away from zero at bounded frequency and grow linearly at large frequency. The homotopies preserve smooth Fourier series and all their derivatives. The vertical-horizontal filtration of the tangent extension, followed by the filtration in the base antiholomorphic differential, gives a finite filtration with acyclic nonzero-character quotients. Hence their full complexes are acyclic. The zero-character formulas have holomorphic coefficients, so the remaining base Dolbeault resolution is precisely the resolution of \eqref{cdp:eq:C.13}. A finite cyclic stabilizer normalizes translations, so this construction and its averaging projection are equivariant on each root chart.

\item At a regular resonance, $a'=\kappa^t(q,r)^t$ and $(q,r)\ne(0,0)$. Proposition~\ref{cdp:st:E8} shows that $\kappa$ is invertible
at every regular point. Thus $a=s b(s)$ for a primitive holomorphic column $b$.
Permute the two cocycle coordinates so that $b_2$ is a unit and put
\[
 C=\begin{pmatrix}1&-b_1/b_2\\0&1/b_2\end{pmatrix},
 \qquad J=\begin{pmatrix}0&1\\1&0\end{pmatrix},
 \qquad R=J(\kappa C^t)^{-1}.
\]
The coefficient changes $a\mapsto Ca$ and
$\kappa\mapsto R\kappa C^t$ are induced by the cocycle and
vertical frame changes, extended to their exterior powers.
Their inverses are
$C^{-1}=\left(\begin{smallmatrix}1&b_1\\0&b_2\end{smallmatrix}\right)$
and $R^{-1}=\kappa C^tJ$. These holomorphic invertible maps
carry \eqref{cdp:eq:C.13} to
\[
 a=(0,s),\qquad\kappa=J.
\]
In degree one the cycle generators are
$u=e_1G_2$, $v=e_2G_2$, $w=e_2G_1+shG_2$, with boundary
relations $su,sv,u+w$. The maps $u\mapsto(1,0)$,
$v\mapsto(0,1)$, $w\mapsto(-1,0)$ identify their quotient
with $(\C\{s\}/(s))^2$; the inverse uses the classes of $u,v$. In degree two, write
$A=e_1G_1G_2$, $B=e_2G_1G_2$, $C=hG_1G_2$.
The relations are $sA,sB,A+sC,B$, so the mutually inverse maps
$A\mapsto-s$, $B\mapsto0$, $C\mapsto1$ and $1\mapsto C$
identify the quotient with $\C\{s\}/(s^2)$.
These are the regular modules in \eqref{cdp:eq:C.12}.

\item At the order-three point write $s$ for the fixed cubic
root of \cref{def:II:marked-conventions}, and put $t=s^3$,
$\nu=\mu-(2-\tau)/3$ and
$b=\beta_c-\beta_c(z_1)+4(\mu-\mu(z_1))+(2/3)(\tau-\rho)$.
The root characters give $\nu=\nu_2s^2+O(s^3)$ and
$b=b_3s^3+O(s^4)$. The coordinate change
$(x,y)=(-2x_0,2i\sqrt2\,y_0)$, with inverse
$(x_0,y_0)=(-x/2,y/(2i\sqrt2))$, carries the elliptic curve of
\cite[Remark~3.12, p.~19]{AC26} to
$y^2=4x^3-12t^3(t-1)x+8t^4(t-1)^2$.
After $x=s^4X$, $y=s^6Y$ this is smooth at $s=0$.
Consequently its elliptic scale is
$\varphi(s)=Fs^2+O(s^3)$, $F\ne0$, where
$x=\varphi^2\wp$ and $y=\varphi^3\wp'$.
Writing $\tau=\rho+\tau_1s+O(s^2)$, we obtain
\[
 \nu_2^2=9F^2/512,\quad \tau_1=i\pi F^2/2,\quad
 6\mu'^2-\tau'\beta_c'=(3F^2/4)s^2+O(s^3).
\]
For the section $P$ with $x(P)=0$, the chord-and-tangent
formulas give $x(2P)=9t^2/8$ and
$x(3P)=(512-592t+80t^2)/81$.
Since $3\mu$ differs from $3\nu$ by an elliptic period,
$x(3P)(0)=F^2/(9\nu_2^2)$, proving the first equality.
The identities
$\varphi^4g_2=12t^3(t-1)$ and
$\varphi^6g_3=-8t^4(t-1)^2$ give
$g_3(\rho)=-8/F^6$.
Ramanujan's derivative formula, in the normalization of
\cite[eqs.~(3.2) and (3.5), p.~8]{AtkinGarvan2002}, gives
$g_2'(\rho)=-3i g_3(\rho)/\pi$; comparing the linear term of
$g_2(\tau(s))$ proves the stated $\tau_1$.
Proposition~\ref{cdp:st:E8} then gives the determinant coefficient.
In $6\nu'^2-\tau'b'$, the coefficient equation is
$24\nu_2^2-3\tau_1b_3=3F^2/4$, whence $b_3=7i/(32\pi)$.
The cohomology-frame change
$(a_1,a_2)\mapsto-(a_2,a_1+2a_2)$, with inverse
$(c,d)\mapsto(2c-d,-c)$, sends the Picard section to
$3k(\nu,b)$. Adding one third of the first vertical row to the
second gives the Kodaira--Spencer matrix
$\left(\begin{smallmatrix}\tau'&6\nu'\\\nu'&b'\end{smallmatrix}\right)$
after the common sign is absorbed in the vertical frame.
Eliminating the unit $\tau'$ leaves transverse coefficient
\begin{equation}
 3k(b-6\nu\nu'/\tau')
   =\frac{123ik}{64\pi}s^3+O(s^4),\qquad k\ne0.        \hypertarget{tag-008S}{}\label{cdp:eq:C.14}
\end{equation}
It is nonzero. Since $\nu/s^2$ is a unit, subtracting
$(b/\nu)$ times the first cocycle coordinate from the second
aligns the Picard section. The transverse coefficient of the
first Kodaira--Spencer row is
$6\nu'-\tau'b/\nu=-\tau'(b-6\nu\nu'/\tau')/\nu$;
by \eqref{cdp:eq:C.14} it has order one.
After exchanging the two row and column orders and dividing by
the Picard unit, the matrix has the form
$\left(\begin{smallmatrix}A s^2&B s\\C s&D\end{smallmatrix}\right)$,
where $C,D$ are units and the determinant has order two.
Subtract $(A/C)s$ times the second row from the first, then
rescale the first vertical frame by $B-AD/C$, the second by $D$,
and the first cocycle frame by $C/D$.
These units are nonzero; reversing the operations gives the
inverse isomorphism. Their coefficients have the required root
characters. Thus \eqref{cdp:eq:C.13} becomes
\[
 a=(0,s^2),\qquad\kappa=\begin{pmatrix}0&s\\s&1\end{pmatrix}.
\]
For degree one take cycle generators
$u=e_1G_2$, $v=e_2G_2$, $w=e_2G_1+shG_2$.
Their relations are $s^2u,s^2v,su+v+sw$.
The maps $u\mapsto(1,0)$, $v\mapsto(-s,-s)$,
$w\mapsto(0,1)$ and the inverse maps from the two generators
to $u,w$ identify the quotient with
$\C\{s\}/(s^2)\oplus\C\{s\}/(s^3)$.
In degree two the relations in $A,B,C$ are
$s^2A,s^2B,sA+B+s^2C,sB$.
The maps $A\mapsto(1,0)$, $B\mapsto(-s,-s^2)$,
$C\mapsto(0,1)$, inverse to the maps from the generators to
$A,C$, give the same two cyclic modules.

\item The affine characters determine the coarse invariant modules. The integral cocycle at the order-three resonance is $(n,-2k,-k,-3k)$; evaluating on $v_1$ gives $n$. At the order-four resonance it is $(n,-k,-k,-2k)$; evaluation on $v_2$ gives $-n\equiv3n\pmod4$. The invariant exponents at the order-three point are $n+2,n+2$ for the two degree-one generators and $n+2,n$ in degree two. Here an invariant exponent is the required Taylor exponent modulo the stabilizer order. Thus the length-two summand retains precisely the monomials $s^j$ with $0\le j<2$ and $j\equiv n+2\pmod3$: its invariant lengths for $n\equiv0,1,2$ are $0,1,1$. A length-three summand retains one coefficient in every residue class. Both coarse lengths are therefore $1+\epsilon_n$, and multiplication by $t=s^3$ kills them. At the order-four point the nonzero functional
$(1,2)$ and the invertible matrix of \ref{cdp:st:E8} give a
simple zero of the Picard section, so the regular form above applies. The two length-one degree-one generators have invariant exponents $3n+1,3n+2$, and the length-two degree-two generator exponent $3n+2$. Their invariant length is one exactly for $n\equiv1,2\pmod4$, and the coarse parameter $s^4$ kills them. This proves all noncusp summands.

\item At the cusp, the restriction of $M_z$ has normalization pullback trivial and dual-complex character $(a,1)$, $a=e^{2\pi iz}$. Choose the two oriented sides and diagonal of the square in
the dual torus. Its triangular cochains have differentials
$1\mapsto(a-1,0,a-1)$ and
$(u,v,w)\mapsto(u+av-w,-u-v+w)$.
The chain inclusion from the two-edge square sends
$(u,v)\mapsto(u,v,u+v)$ in degree one and
$c\mapsto(c,0)$ in degree two. Its projection sends
$(u,v,w)\mapsto(u,v)$ and $(c,d)\mapsto c+d$.
The homotopy $(c,d)\mapsto(0,0,d)$ contracts their complement.
Thus the normalization complex is homotopy equivalent to
\[
 \C\xrightarrow{(a-1,0)^t}\C^2
 \xrightarrow{(0,a-1)}\C.
\]
For $z\notin\Z$ it is acyclic. Semicontinuity and base change make $R\bar f_*M_z$ vanish on a cusp neighbourhood. The logarithmic sequence \eqref{cdp:eq:C.8} then gives $R\bar f_*(\mathcal T\otimes M_z)=0$ there. Sequence \eqref{cdp:eq:C.3} contributes exactly the reduced skyscraper $\C_{p_0}$ in degree two for the ordinary tangent bundle. Away from resonances the Fourier contractions in step~2
give zero cohomology for the line and for its tangent extension. We verify that this analytic support is proper. Fix a simply connected disc in the regular base. For each fixed integer tuple $(n,p,q,r)$, a resonance has $r\ne0$ and $a^\prime=\kappa^t(q,r)^t\ne0$, so its zeros are isolated. Each such zero set is countable on this disc, and the union over the countably many tuples is countable. The disc has a point outside this union. At that fibre the preceding line-bundle acyclicity and tangent extension give zero cohomology in every degree. Upper semicontinuity gives a neighbourhood with the same vanishing. Each coherent higher direct image therefore has proper analytic support on the compact base curve, hence finite support. Every regular resonance lies in that support by its computed nonzero local modules, proving that $\mathcal R_z$ is finite. Degree zero is zero because a direct image of a bundle is torsion-free and has generic rank zero. These arguments prove \eqref{cdp:eq:C.12}, and Leray proves \eqref{cdp:eq:C.11}. The logarithmic degree-two direct image is a torsion sheaf for these twists, so its first cohomology vanishes. Equation \eqref{cdp:eq:C.4} therefore makes the normalization
map surjective.

\end{enumerate}
\end{proof}



\begin{corollary}[The generic tangent tuple and its jump loci]\hypertarget{tag-P262}{}\label{cor:II:tangent-jump-loci}
\StatementRecord{cor:II:tangent-jump-loci}{proved}{P262}{\SeriesRef{cdp:st:C5b}, \SeriesRef{cdp:st:C5c}, \SeriesRef{cdp:st:C.1.1}}
For the marked family of \cref{def:II:marked-Picard-parameter},
$\chi(X,T_X\otimes M_z)=1$ for every $z\in\C$. Its generic
tangent tuple is $(0,0,1,0)$, attained at every real
nonintegral $z$. For each $q$ and each integer $k$ strictly larger
than the generic value of $h^q$, the set
\[
 \{z\in\C:h^q(X,T_X\otimes M_z)\ge k\}
\]
is a proper locally discrete complex analytic subset of the
Picard line. The nonzero degree-two group, index one and
$c_3(X)=2$ correct respectively parts (a), (b) and (c) of
\cite[Theorem~2.2, p.~680]{CDP2020} for this family.
\end{corollary}

\begin{proof}
The index is $c_3(X)/2=1$ by \cref{cdp:st:C.1.1}.
For real nonintegral $z$, \eqref{cdp:eq:C.10} has no regular
resonance and the nonreal generators of $\Lambda_1,\Lambda_2$
give $\ell_1(z)=\ell_2(z)=0$. Thus
\cref{cdp:st:C5c} gives $(0,0,1,0)$.
Choose a global $\bar\partial$-closed $(0,1)$-form $\alpha$
representing the chosen Picard coordinate. The operators
$\bar\partial+z\alpha$ define its holomorphic universal line
family on $X\times\C$. Proper coherent semicontinuity makes
each displayed cohomology threshold analytic. It is proper
because every real nonintegral parameter has the stated
generic tuple, so it is locally discrete in $\C$.
\end{proof}

\section{Constant-monodromy period comparison and compact \texorpdfstring{\mbox{examples}}{examples}}\hypertarget{tag-0349}{}\label{sec:II:13}

\subsection{Components and the logarithmic period obstruction}

\begin{proposition}[Components under trivial torus monodromy]\hypertarget{tag-01G8}{}\label{cdp:st:D.2.1}
\StatementRecord{cdp:st:D.2.1}{proved}{01G8}{\SeriesRef{cdp:st:D.1.2}, \SeriesRef{cdp:st:C.3.4}}

A smooth compact complex threefold with $b_1=b_2=0$ and a holomorphic
connected-fibre map to a curve, with general fibre a two-torus and
trivial rational monodromy, has reducible fibres. More precisely,
for every nonempty set of $s$ removed values containing all critical
values, their reduced fibres have $r=s+3$ irreducible components in
total. Every component of each reducible fibre has no compact
K\"ahler resolution. Consequently no such fibration can have all
its fibre components admitting compact K\"ahler resolutions.
\end{proposition}


\begin{proof}
\begin{enumerate}[label=\arabic*.,leftmargin=*]
\item 
Pullback on first rational cohomology from the base is injective because
the fibres are connected. Thus the base is $\PP^1$. Trivial monodromy on
the four-dimensional $H_1(F;\Q)$ gives $t'=4$, and the count in
Proposition~\ref{cdp:st:D.1.2} gives $r=s+3$. If all removed fibres were
irreducible one would have $r=s$, which is impossible. The assertion
about each component is Proposition~\ref{cdp:st:C.3.4}, applied to this
threefold. Such a reducible fibre prevents all fibre components from
admitting compact K\"ahler resolutions.

\end{enumerate}
\end{proof}

\begin{proposition}[Logarithmic period obstruction]\hypertarget{tag-01GA}{}\label{cdp:st:D.2.3}
\StatementRecord{cdp:st:D.2.3}{proved}{01GA}{\SeriesRef{cdp:eq:D.18}, \SeriesRef{cdp:eq:D.19}; \cite[\S2.1, p.~234; Theorem~2.18, pp.~237--238; \S2.19 and Proposition~2.20, pp.~238--239; Theorem~2.21, p.~239]{Steenbrink1976}; \cite[Proposition~1.13, pp.~233--234, and proof of Proposition~2.20, pp.~238--239]{Steenbrink1976}}


Let $h:U\to\Delta$ be a proper holomorphic map from a smooth
complex threefold, with connected fibres, smooth two-torus fibres
over $\Delta^*$, and central fibre $F=\sum_i m_iD_i$ with
normal-crossings reduced support $D=\bigcup_iD_i$.
If every normalization component of $D$ has geometric genus zero,
the monodromy on $H^2(U_t;\C)$ is nontrivial. In particular the
monodromy on $H^1(U_t;\Q)$ is nontrivial. The conclusion holds for every choice of positive integers $m_i$.
\end{proposition}


\begin{proof}

\begin{enumerate}[leftmargin=*]
\item Put
\[
\mathscr K^\bullet=
 \Omega_U^\bullet(\log D)\big/
 \bigl(h^*\Omega_\Delta^1(\log0)\wedge
       \Omega_U^{\bullet-1}(\log D)\bigr),\qquad
\mathscr H^2=\mathbb R^2h_*\mathscr K^\bullet.
\]
In \cite[\S2.1, p.~234; Theorem~2.18, pp.~237--238; \S2.19 and Proposition~2.20, pp.~238--239; Theorem~2.21, p.~239]{Steenbrink1976}, the assumptions for these assertions
are precisely a proper holomorphic map from a connected complex
manifold to a disc, smooth off zero, with normal-crossings support.
They give a locally free bundle $\mathscr H^2$,
its Gauss--Manin connection
\[
\nabla:\mathscr H^2\longrightarrow
\Omega_\Delta^1(\log0)\otimes\mathscr H^2,
\]
and the comparison
$\mathscr H^2|_{\Delta^*}\simeq
R^2h_*\C\otimes\OO_{\Delta^*}$.
The residue eigenvalues belong to $[0,1)$ and the extended
monodromy at zero is
$\exp(-2\pi i\operatorname{Res}_0\nabla)$.
Its canonical base-change map identifies the fibre at zero with
$\mathbb H^2(D,\mathscr K^\bullet/t\mathscr K^\bullet)$.
The local calculation on the coefficient ideal uses the monomial
Koszul complex of \cite[Proposition~1.13, pp.~233--234, and proof of Proposition~2.20, pp.~238--239]{Steenbrink1976}.

\item Suppose that the monodromy on $H^2$ were the identity. Its
residue has all eigenvalues zero, and its exponential is the
identity. A nilpotent matrix whose exponential is the identity
is zero. Hence the residue vanishes. In a holomorphic frame
$\nabla=d+A(t)\,dt/t$ with $A(0)=0$, so this is a holomorphic
connection across zero. It has a holomorphic horizontal frame
on the disc. Consequently the coordinates of every holomorphic
section of $\mathscr H^2$ in a flat frame, that is its periods,
are holomorphic and bounded near zero.

\item The second exterior power of the relative logarithmic cotangent bundle is
\[
\mathscr K^2\simeq
 \omega_U(D)\otimes h^*\omega_\Delta(0)^{-1}
 \simeq\omega_{U/\Delta}(D-F).
\]
The direct image $h_*\mathscr K^2$ is a torsion-free
coherent sheaf of rank one, because a smooth torus fibre has
one-dimensional space of holomorphic two-forms. Shrink the disc
and choose a generator $\Omega$ of this line bundle.
Thus $\Omega$ determines a holomorphic section $[\Omega]$ of
$\mathscr H^2$ by the inclusion of the last term of the
relative de Rham complex. Its smooth-fibre comparison is the
ordinary class $[\Omega_t]\in H^2(U_t;\C)$.

\item First suppose that $\Omega$ restricts nontrivially to a normalized
component $T$ of $D$. There is a canonical isomorphism
\[
 \mathscr K^1|_T\simeq\Omega_T^1(\log D_T),\qquad
 \mathscr K^2|_T\simeq\omega_T(D_T),
\]
where $D_T$ is the reduced conductor. To specify the map, on the
branch $z_i=0$ subtract
$\operatorname{res}_i(\eta)(dt/t)/m_i$ from an absolute logarithmic
one-form $\eta$ and restrict its tangential part. The normal
residue of $dt/t$ is $m_i$, so this kills exactly the subbundle
divided out in $\mathscr K^1$ and is independent of the lift.
The exterior square gives the second isomorphism. If every
residue of $\Omega|_T$ vanished, its local logarithmic expression
would be holomorphic on $T$. This contradicts $p_g(T)=0$.
Thus there is a smooth point of a double curve, away from triple
points, where its residue is nonzero.

\item In a small chart about that point, write $h=x^ay^b=t$ and use $z$
along the double curve. The relative two-form is
\[
 \Omega_t=g(x,y,z)\,\frac{dx}{x}\wedge dz,
 \qquad g(0,0,z_0)\ne0.
\]
Choose a fixed small $z$-disc and $\epsilon>0$ such that
$|g(x,y,z)|\ge c>0$ whenever $|x|,|y|<\epsilon$ there.
The map to the annulus
$|t|^{1/a}\epsilon^{-b/a}<|x|<\epsilon$
has $b$ sheets. Integration on its inverse image and the fixed
$z$-disc gives
\[
 \int_{U_t}\Omega_t\wedge\overline{\Omega_t}
 \ \ge\ c'\log\frac{\epsilon^{a+b}}{|t|}
 \ \longrightarrow\ +\infty.
\]
Here the integrand is a nonnegative real volume form on the
complex surface; its local expression is a positive constant
times $|g|^2|x|^{-2}\,dA(x)dA(z)$.

\item On the other hand, $\Omega_t$ is closed, and this integral is
the fixed cup-product pairing of $[\Omega_t]$ with its complex
conjugate. All its periods are bounded by the holomorphic
connection argument, so the pairing is bounded. This is a
contradiction.

\item Suppose that $\Omega$ vanishes on every
branch of $D$. Put
\[
 \mathscr J=\OO_U(-D),\qquad
 \mathscr B^\bullet=
 \mathscr J\mathscr K^\bullet/t\mathscr K^\bullet.
\]
Both $\mathscr J\Omega_U^\bullet(\log D)$ and
$t\Omega_U^\bullet(\log D)$ are subcomplexes for the absolute
exterior derivative: locally the logarithmic derivatives of
$z_1\cdots z_r$ and $t=\prod z_i^{m_i}$ are holomorphic
linear combinations of the $d\log z_i$. Hence the locally split
absolute/relative logarithmic exact sequence remains exact
when its coefficients are restricted to $\mathscr J/(t)$:
\begin{equation}
0\longrightarrow\mathscr B^\bullet[-1]
 \xrightarrow{\ (dt/t)\wedge\ }
 (\mathscr J/(t))\Omega_U^\bullet(\log D)
 \longrightarrow\mathscr B^\bullet\longrightarrow0.
\hypertarget{tag-0195}{}\label{cdp:eq:D.18}
\end{equation}
Signs in the shifted differential are chosen so that the first
arrow is a morphism of complexes. Its connecting morphism
$\rho_{\mathscr B}:\mathscr B^\bullet\to\mathscr B^\bullet$
in the derived category defines degree-preserving endomorphisms
on both cohomology sheaves and hypercohomology. The inclusion
$\mathscr J/(t)\hookrightarrow\OO_U/(t)$ gives a commutative
diagram from \eqref{cdp:eq:D.18} to the analogous exact sequence for
$\mathscr K^\bullet/t\mathscr K^\bullet$. The induced
endomorphism for the latter is precisely the Gauss--Manin
residue. Thus the map on hypercohomology intertwines these
endomorphisms.

\item At a point of $r$ branches,
write
\[
 t=\prod_{i=1}^r z_i^{m_i},\qquad
 e=\gcd(m_1,\ldots,m_r),\qquad
 y=\prod_{i=1}^r z_i^{m_i/e},\qquad t=y^e.
\]
The transverse holomorphic Poincar\'e lemma reduces the relative
complex to the Koszul complex with commuting operators
\[
 D_i=z_i\partial_{z_i}
           -\frac{m_i}{m_1}z_1\partial_{z_1},
 \qquad 2\le i\le r,
\]
on the convergent-power-series module
$(z_1\cdots z_r)/(z_1^{m_1}\cdots z_r^{m_r})$.
The monomials diagonalize these operators. A monomial with some
nonzero eigenvalue has a contracting Koszul homotopy, obtained
by division by that eigenvalue. These divisions preserve
convergence: every nonzero eigenvalue has absolute value at
least $1/m_1$. The surviving monomials have exponent vector
$a(m_1/e,\ldots,m_r/e)$. Membership in $\mathscr J$ requires
$a\ge1$, and quotienting by $t$ requires $a<e$. Consequently
each cohomology stalk is spanned by
\[
 y^a\,d\log z_{i_1}\wedge\cdots\wedge d\log z_{i_q},
 \qquad 1\le a<e,
\]
subject to the relative relation $\sum m_i d\log z_i=0$.
Since $dy^a=(a/e)y^a\,dt/t$, the connecting endomorphism
acts on these representatives by $a/e$. In particular zero is
not an eigenvalue; when $e=1$ the complex is acyclic.

\item The hypercohomology spectral sequence
\[
 H^p\bigl(D,\mathcal H^q(\mathscr B^\bullet)\bigr)
 \ \Longrightarrow\ \mathbb H^{p+q}(D,\mathscr B^\bullet)
\]
is equivariant for this derived endomorphism. There are only
finitely many branch multiplicities, so all eigenvalues on its
terms and abutment lie in a finite subset of $(0,1)$. For
example the product of the corresponding linear factors
annihilates every cohomology sheaf, and a power annihilates the
abutment. Hence $\rho_{\mathscr B}$ is invertible on
$\mathbb H^2(D,\mathscr B^\bullet)$. The residue on
$\mathscr H^2|_0$ is zero. Equivariance therefore proves that
\begin{equation}
 \mathbb H^2(D,\mathscr B^\bullet)\longrightarrow
 \mathbb H^2(D,\mathscr K^\bullet/t\mathscr K^\bullet)
       =\mathscr H^2|_0
 \quad\hbox{is the zero map}.
\hypertarget{tag-0196}{}\label{cdp:eq:D.19}
\end{equation}

\item The central class of $\Omega$ comes from the left side of
\eqref{cdp:eq:D.19}, because $\Omega\in\mathscr J\mathscr K^2$.
Thus $[\Omega]$ vanishes at zero in the locally free bundle
$\mathscr H^2$, and $[\Omega]/t$ is a holomorphic section of
that bundle. Its periods are bounded. On a smooth fibre this
section is the de Rham class of $\Omega_t/t$.

\item Primitivity of the generator $\Omega$ implies that its order
$b_i$ along at least one $D_i$ is less than $m_i$; otherwise
$\Omega/t$ would be a holomorphic section of $\mathscr K^2$
over $U$, contradicting its choice as a generator of
$h_*\mathscr K^2$. Choose a smooth point of that component
away from all other branches and all further zeros. There
are coordinates with $t=z^{m_i}$ in which
\[
 \Omega=z^{b_i}a(z,\xi,\eta)\,d\xi\wedge d\eta,
 \qquad a(0,0,0)\ne0.
\]
On one sheet above a fixed small $(\xi,\eta)$-polydisc, the
positive integral of $\Omega_t/t$ against its conjugate has
the lower bound
\[
 \int_{U_t}\frac{\Omega_t}{t}\wedge
                    \overline{\frac{\Omega_t}{t}}
 \ \ge c''|t|^{2(b_i/m_i-1)}\longrightarrow+\infty.
\]
Its expression as the fixed cup-product pairing of bounded
periods again gives a contradiction.

\item The natural
cup-product isomorphism
$\bigwedge^2H^1(U_t;\Q)\simeq H^2(U_t;\Q)$ for a two-torus
is monodromy-equivariant. Identity on first cohomology would
therefore imply identity on second cohomology.

\end{enumerate}
\end{proof}

\subsection{The Euler upper bound under trivial monodromy}

\begin{proposition}[Euler sign under trivial torus monodromy]\hypertarget{tag-01GC}{}\label{cdp:st:D.2.5}
\StatementRecord{cdp:st:D.2.5}{proved}{01GC}{\SeriesRef{cdp:st:C4c}, \SeriesRef{cdp:st:A.3.5}, \SeriesRef{cdp:eq:D.S.2}, \SeriesRef{cdp:st:D.1.9}, \SeriesRef{cdp:st:A.4.7}, \SeriesRef{cdp:st:A.5.2}, \SeriesRef{cdp:eq:D.S.4}, \SeriesRef{cdp:st:A.5.6}, \SeriesRef{cdp:st:A.5.3}, \SeriesRef{cdp:eq:D.S.5}, \SeriesRef{cdp:st:D.2.3}; \cite[Theorem~1.6 and Remarks~1.7, arXiv v2, pp.~5--6]{BierstoneMilman1997}}


Let $f:Y\to\PP^1$ be a proper surjective holomorphic map with
connected fibres from a compact connected complex threefold with
$b_1(Y)=b_2(Y)=0$. Suppose its smooth general fibre is a complex
two-torus and its monodromy representation on $H_1(F,\Q)$ is
trivial. Then
\begin{equation}
 c_3(Y)\le0,\qquad b_3(Y)\ge2.\hypertarget{tag-019O}{}\label{cdp:eq:D.S.1}
\end{equation}
The fibres may have arbitrary singularities and arbitrary
multiplicities. The construction in \cref{cdp:st:D.2.7} attains
$c_3(Y)=0$ and $b_3(Y)=2$.
\end{proposition}


\begin{proof}

\begin{enumerate}[label=\arabic*.,leftmargin=*]
\item Assume $c_3(Y)>0$. Poincar\'e duality gives
$c_3(Y)=2-b_3(Y)$, and the alternating nondegenerate intersection
pairing on $H^3(Y,\Q)$ makes $b_3(Y)$ even. Hence $b_3(Y)=0$.
Apply \ref{cdp:st:C4c} to the Picard component
$P=K_Y\otimes\Pic^0(Y)$. It gives a fixed critical full fibre
$F$ such that
\begin{equation}
 H^0\bigl(F,\Omega_Y^1|_F\otimes M|_F\bigr)\ne0
 \qquad(M\in\Pic^0(Y)).\hypertarget{tag-019P}{}\label{cdp:eq:D.S.2}
\end{equation}
Indeed the Serre-dual parameter $K_Y\otimes L^{-1}$ runs through
$\Pic^0(Y)$ when $L$ runs through $P$.

\item \Cref{cdp:st:A.3.5} excludes reducible support in \eqref{cdp:eq:D.S.2}, including
all unequal-multiplicity possibilities. Thus $F=mS$ for one
irreducible reduced Cartier divisor $S$ and one positive integer
$m$. We specify why a smooth $S$ is excluded under the present
\emph{rational} Betti-number assumptions. \ref{cdp:st:D.1.9} supplies
$\beta\in H^1(Y,\OO_Y)$ whose restriction near every full fibre
comes from a nonzero integral transgression class. The line
$M_r=\exp(2\pi i r\beta)$, for real irrational $r$, has that
unitary local realization. The normalized root change $t=s^m$
near a smooth $S$ is the finite \emph{\'etale} map of \cref{cdp:st:A.3.4}:
in a chart $t=uz^m$ it has equations $w^m=u$, $s=wz$.
Its central reduction $S'$ is smooth, and its nearby fibres are
the original tori. Proper smooth local triviality transports
the vanishing of degree-zero and degree-one cohomology of the
pulled-back unitary local system from a nearby torus to $S'$.
The injectivity \ref{cdp:st:A.4.7} and the conormal sequence then give zero
ambient cotangent sections on $S'$. Each layer of $mS'$ has the
same twist, since $s$ trivializes its conormal. The finite
filtration and exact invariants functor for the cyclic cover
give $H^0(mS,\Omega_Y^1|_{mS}\otimes M_r)=0$, contradicting
\eqref{cdp:eq:D.S.2}. 
Therefore $S$ is singular. \ref{cdp:st:A.5.2}, applied using
$\omega_S=(K_Y\otimes\OO_Y(S))|_S$ and
$\chi(\OO_S)=0$, makes $S$ nonnormal.

\item The ideal filtration of $mS$ has the finitely many locally free
graded factors on $S$
\begin{equation}
 E_j=(\Omega_Y^1\otimes\OO_Y(-jS))|_S,
 \qquad 0\le j<m.\hypertarget{tag-019Q}{}\label{cdp:eq:D.S.3}
\end{equation}
If all the groups $H^0(S,E_j\otimes M|_S)$ were zero, the exact
filtration sequences would make \eqref{cdp:eq:D.S.2} zero. Thus their
nonvanishing loci cover $\Pic^0(Y)$. Each locus is closed
complex analytic by the Picard family and proper
semicontinuity. The group $\Pic^0(Y)$ is a connected complex
vector space, since $H^1(Y,\Z)=0$. A finite union of proper
closed analytic subsets cannot cover it. Some fixed $j$
therefore satisfies
\begin{equation}
 H^0(S,E_j\otimes M|_S)\ne0
 \qquad(M\in\Pic^0(Y)).\hypertarget{tag-019R}{}\label{cdp:eq:D.S.4}
\end{equation}

\item Let $\eta:T\to S$ be the finite normalization and
$\pi:Z\to T$ its minimal resolution, and put $h=\eta\pi$,
viewed also as a holomorphic map to $Y$.
Pullback of sections in \eqref{cdp:eq:D.S.4} is injective: $S$ is reduced
and irreducible, and a section of the locally free sheaf $E_j$
vanishing on its dense smooth locus is zero. Thus the fixed
vector bundle $h^*(\Omega_Y^1\otimes\OO_Y(-jS))$ has a
section after every twist by $h^*M$, $M\in\Pic^0(Y)$.
For any Gauduchon metric $\omega$ on $Z$, the uniform
negative-degree bound \ref{cdp:st:A.5.6}, applied to all positive and
negative powers of $M$, gives
\[
 \deg_\omega h^*M=0\qquad(M\in\Pic^0(Y)).
\]
Every line bundle $G$ on $Y$ has a positive power in
$\Pic^0(Y)$, because $H^2(Y,\Z)$ is finite. Additivity of
degree therefore gives $\deg_\omega h^*G=0$ for every such
$G$.

\item The conductor comparison \ref{cdp:st:A.5.3} gives an effective divisor
$D_Z$ with
\begin{equation}
 \omega_Z(D_Z)\simeq h^*(K_Y\otimes\OO_Y(S)).\hypertarget{tag-019S}{}\label{cdp:eq:D.S.5}
\end{equation}
Moreover $D_Z\ne0$. Indeed a reduced hypersurface is
Cohen--Macaulay and hence satisfies Serre's condition $S_2$;
its nonnormality must therefore occur in codimension one.
The divisorial conductor is nonzero, and its strict transform
is part of $D_Z$.
Taking Gauduchon degree in \eqref{cdp:eq:D.S.5} yields
\[
 \deg_\omega K_Z=-\deg_\omega\OO_Z(D_Z)<0.
\]
A nonzero holomorphic section of $K_Z$ would have an effective
zero divisor, whose degree is nonnegative. This contradiction
proves $p_g(Z)=h^0(Z,K_Z)=0$.

\item Choose a small closed fibre neighbourhood $U=f^{-1}(\Delta)$.
Embedded resolution of the divisor $F$ gives a proper
holomorphic modification
\[
 q:\widehat U\longrightarrow U
\]
which is a finite composite of blowups with smooth centres in
the successive central divisors, is an isomorphism over
$\Delta^*$, has smooth total space, and makes the reduced
support of $(f\circ q)^{-1}(0)$ a simple normal-crossings
divisor. We use the analytic embedded-resolution theorem in
\cite[Theorem~1.6 and Remarks~1.7, arXiv v2, pp.~5--6]{BierstoneMilman1997}; properness of the
central fibre makes its locally finite construction finite
after shrinking $\Delta$. The strict transform of $S$ is a
resolution of its normalization and has geometric genus zero
by birational invariance of holomorphic two-forms on smooth
compact surfaces. A new exceptional component is $\PP^2$ over
a point or a $\PP^1$-bundle over a smooth compact curve.
Both have geometric genus zero: restriction of a holomorphic
two-form to a ruling chart has negative fibre degree, and
$H^0(\PP^2,K_{\PP^2})=0$. Further strict transforms preserve
this property. Consequently every component of the final
central reduced divisor has geometric genus zero.

\item The map $f\circ q$ remains proper with connected fibres:
a blowup has connected fibres and its structure-sheaf direct
image is the original structure sheaf. Its smooth fibres over
$\Delta^*$ are the original two-tori under the isomorphism
$q|_{\Delta^*}$, so their $H^1$ monodromy remains the identity.
The multiplicities of the resolved central fibre are arbitrary
positive integers and are all permitted in \ref{cdp:st:D.2.3}. That
proposition now gives nontrivial monodromy on $H^2$ and hence
on $H^1$, a contradiction. This proves \eqref{cdp:eq:D.S.1}.

\end{enumerate}
\end{proof}

\subsection{Extension of periods and coherent direct images}

\begin{proposition}[Extension of the marked torus period]\hypertarget{tag-01GD}{}\label{cdp:st:D.2.6}
\StatementRecord{cdp:st:D.2.6}{proved}{01GD}{\SeriesRef{cdp:st:D.2.3}, \SeriesRef{cdp:eq:D.20}}


Let $h:U\to\Delta$ be a proper holomorphic map from a smooth
complex threefold, smooth with connected two-torus fibres over
$\Delta^*$. Suppose that its integral $H^1$ monodromy is the
identity. The marked period map
\[
 \Delta^*\longrightarrow
 \mathcal T_2
 :=\{W\in\operatorname{Gr}(2,H^1(U_t;\C)):
                   W\oplus\overline W=H^1(U_t;\C)\}
\]
extends holomorphically across zero with value in $\mathcal T_2$.
In particular the relative Albanese group family over $\Delta^*$
extends to a smooth proper group family of complex two-tori over
$\Delta$. 
\end{proposition}


\begin{proof}

\begin{enumerate}[leftmargin=*]
\item The punctured family is connected by local triviality and connectedness
of its fibres, so it lies in one open and closed component $U^0$ of $U$.
The inclusion $U^0\hookrightarrow U$ is an isomorphism over $\Delta^*$.
Apply analytic embedded resolution to the compact reduced central divisor
in $U^0$ \cite[arXiv v2, \S0.2(2), p.~3; Theorem~1.6 and Remarks~1.7,
pp.~5--6]{BierstoneMilman1997}. This gives a proper composite
of blowups $\pi:\widetilde U\to U^0$ with smooth centres over that
divisor, with $\widetilde U$ smooth and the reduced central support
of $\widetilde h:=h|_{U^0}\circ\pi$ normal crossing. For $t\ne0$,
$\pi_t:\widetilde U_t\to U_t$ is a biholomorphism;
its pullback identifies holomorphic forms, the integral marking and
cup products. Transporting through these pullbacks, use the relative
logarithmic complex in the proof of Proposition~\ref{cdp:st:D.2.3}
on this model, with its full central multiplicities.
Identity on $H^1$ gives identity on $H^2=\bigwedge^2H^1$.
Consequently its logarithmic Gauss--Manin bundle has zero
residue and a holomorphic flat frame. Let $\Omega$ be a
generator of $\widetilde h_*\mathscr K^2$, as in Proposition~\ref{cdp:st:D.2.3}.
For $t\ne0$, let $\Omega_t$ be the form on $U_t$ determined by
$\pi_t^*\Omega_t=\Omega|_{\widetilde U_t}$.

\item The coefficient-ideal argument in
\eqref{cdp:eq:D.18}--\eqref{cdp:eq:D.19} shows that $\Omega$
cannot vanish on the whole reduced central support: otherwise
its divided cohomology class would have bounded periods while
its positive $L^2$ volume diverges. Thus $\Omega|_T$ is nonzero
on some normalized component $T$. The annular volume estimate
in the same proof makes every conductor residue of every
component restriction zero: a nonzero residue would give a
divergent positive volume while its cup-product expression
stays bounded. The nonzero restriction $\Omega|_T$ is therefore
a holomorphic two-form on $T$.

\item Choose a point of $T$ outside the conductor and the zeros of
this two-form. In a small chart the map is $t=z^m$ and
\[
 \Omega=a(z,\xi,\eta)\,d\xi\wedge d\eta,
 \qquad a(0,0,0)\ne0.
\]
Integrating on one sheet over a fixed small
$(\xi,\eta)$-polydisc yields a constant $c>0$ such that
\begin{equation}
 \int_{U_t}\Omega_t\wedge\overline{\Omega_t}\ge c
 \qquad(0<|t|\ll1).
\hypertarget{tag-0197}{}\label{cdp:eq:D.20}
\end{equation}
The cohomology class $v(t)=[\Omega_t]$, expressed in the flat
frame, is holomorphic across zero. Its cup pairing with its conjugate equals the
integral in \eqref{cdp:eq:D.20}, hence
\[
 v(0)\wedge\overline{v(0)}>0
 \quad\hbox{in }H^4(U_t;\R)\simeq\R,
\]
using the complex orientation. In particular $v(0)\ne0$.

\item Write $V=H^1(U_t;\C)$ with its fixed integral marking.
For $t\ne0$, $v(t)$ is a decomposable element of
$\bigwedge^2V$, spanning $\bigwedge^2H^{1,0}(U_t)$.
The decomposable cone is the closed algebraic set defined by
the Pl\"ucker equations, so $v(0)$ is decomposable. For a
nonzero decomposable $v$, the exact inverse to the Pl\"ucker
map is
\[
 W(v)=\{\alpha\in V:\alpha\wedge v=0\}.
\]
This kernel has dimension two and varies holomorphically near
$v(0)$. The nonzero pairing
$v(0)\wedge\overline{v(0)}$ says exactly that
$W(v(0))\cap\overline{W(v(0))}=0$. The extended period
therefore lies in $\mathcal T_2$.

\item Let $\mathcal W\subset V\otimes\OO_\Delta$ be the extended
rank-two subbundle $W(v(t))$, and let
$\Lambda=H_1(U_t;\Z)$ with its fixed marking. Evaluation of
cohomology on integral cycles gives a holomorphic embedding
of the constant lattice $\Lambda$ in $\mathcal W^*$.
A real cycle annihilating $\mathcal W_t$ also annihilates
$\overline{\mathcal W_t}$, hence all of $V$; it is therefore zero.
Equality of real dimensions makes the real-linear evaluation
map an isomorphism, so the image is a full lattice in every fibre,
uniformly on compact subsets of the disc. Translation by this
lattice is free and properly discontinuous. For a fixed integral
basis, the closed parallelotopes generated by its images cover
each quotient fibre; over a compact subset of the base their
union is compact. Consequently
\[
 \mathcal A=\mathcal W^*/\Lambda\longrightarrow\Delta
\]
is a smooth proper family of complex Lie groups with
two-torus fibres. On $\Delta^*$, integration of holomorphic
one-forms identifies each group fibre with the Albanese torus
of $U_t$, and these identifications vary holomorphically.


\end{enumerate}
\end{proof}

\begin{remark}[Extended torus groups and Albanese maps]
\hypertarget{tag-P213}{}\label{rem:II:extended-torus-groups}
\StatementRecord{rem:II:extended-torus-groups}{proved}{P213}{}
Let $f:Y\to\PP^1$ be a proper holomorphic map from a smooth complex
threefold, with connected two-torus fibres over the complement of
a finite set and globally trivial marked $H^1(-;\Z)$ monodromy.
The local group extensions in \cref{cdp:st:D.2.6} glue to a smooth
proper torus group family over $\PP^1$. The marked
period extensions agree on overlaps because they agree on the
smooth locus; their quotients by the same constant lattice therefore
glue. The period map can be nonconstant, as in \cref{cdp:st:D.1.15}.

For the local filling $U\to\Delta$ in \cref{night:P4.1}, this
group extension is $A\times\Delta$. An extension
$\alpha:U\to A\times\Delta$ of the marked fibre isomorphisms
would induce the composite
\[
 \Lambda=H_1(A;\Z)\longrightarrow H_1(U;\Z)
       \xrightarrow{\alpha_*}H_1(A\times\Delta;\Z)=\Lambda
\]
equal to the identity. The first arrow is the quotient
$\Lambda\to\Lambda/\Z a$ of \eqref{cdp:eq:D.25}, with $a\ne0$.
Its composite therefore kills $a$, a contradiction. Thus the
marked Albanese map on this filling does not extend to its smooth
group extension.
\end{remark}

\begin{proposition}[Coherent and integral cohomology of the trivial-monodromy example]\hypertarget{tag-01GF}{}\label{cdp:st:D.2.8}
\StatementRecord{cdp:st:D.2.8}{proved}{01GF}{\SeriesRef{cdp:st:D.2.7}, \SeriesRef{cdp:st:D.1.11}, \SeriesRef{cdp:st:D.1.12}, \SeriesRef{cdp:eq:D.36}, \SeriesRef{cdp:st:D.1.7}, \SeriesRef{night:P4.3}}


For the threefold and markings of \ref{cdp:st:D.2.7},
\begin{equation}
 R^0f_*\OO_Y=\OO_{\PP^1},\qquad
 R^1f_*\OO_Y=\OO_{\PP^1}^{\oplus2},\qquad
 R^2f_*\OO_Y=\OO_{\PP^1},
 \qquad h^q(Y,\OO_Y)=(1,2,1,0).
\hypertarget{tag-019K}{}\label{cdp:eq:D.33}
\end{equation}
The coherent restriction map to every full fibre is surjective.
The integral comparison is
\begin{equation}
 H^0(\PP^1,R^1f_*\Z)
   =\Z(e_1^*+e_2^*-e_4^*)
   \ \subsetneq\ H^1(A;\Z)^{\pi_1}=\Z^4.
\hypertarget{tag-019L}{}\label{cdp:eq:D.34}
\end{equation}
The map in \eqref{cdp:eq:D.34} is restriction to the smooth base. The exceptional torsion locus
for restrictions of $\Pic(Y)=H^1(Y,\OO_Y)\simeq\C^2$ to
at least one fibre is exactly
\begin{equation}
 \iota\bigl(H^1(A;\Q)\bigr)\subset H^1(A,\OO_A),
 \qquad
 \iota(a,b,c,d)=\tfrac12(a+\mathrm i b,c+\mathrm i d).
\hypertarget{tag-019M}{}\label{cdp:eq:D.35}
\end{equation}
It is countable and dense. Thus the all-fibre non-torsion locus
of \ref{cdp:st:D.1.11} need not be open, even under integral
$H^1(Y;\Z)=H^2(Y;\Z)=0$.
\end{proposition}


\begin{proof}

\begin{enumerate}[leftmargin=*]
\item Write $\mathscr E=R^1f_*\OO_Y$. Proposition~\ref{cdp:st:D.1.12} gives its
local freeness, full-fibre base change, and rank two for this
threefold. The deformation retraction of a local filling $U_i$
onto its central fibre $S_i$ in Proposition~\ref{night:P4.3},
together with \eqref{cdp:eq:D.25}, identifies
\begin{equation}
 H^1(S_i;\Z)=H^1(U_i;\Z)
          =\{\lambda\in\Lambda^*:\lambda(a_i)=0\}.
\hypertarget{tag-019N}{}\label{cdp:eq:D.36}
\end{equation}
The coefficient map into $H^1(U_i,\OO_{U_i})$
corresponds by coherent Leray on the Stein disc to
holomorphic sections of $\mathscr E$. On the punctured
disc, these are constant vectors in $H^1(A,\OO_A)$,
because \eqref{cdp:eq:D.24} has constant marked periods. The explicit
coefficient map is $\iota$ in \eqref{cdp:eq:D.35}: the $(0,1)$ parts
of $dx,dy$ on $E$ are $d\bar z/2$ and
$\mathrm i\,d\bar z/2$.

\item Each annihilator in \eqref{cdp:eq:D.36} has real dimension three, so
its complex span in the two-dimensional period space is
all of that space. The constant complex relation among
three chosen integral generators holds over the
punctured disc and therefore holds as an identity of
holomorphic sections over the entire disc. At zero,
their three real vectors remain independent: the map
$H^1(S_i;\R)\to H^1(S_i,\OO_{S_i})$ is real-linearly
injective by \ref{cdp:st:D.1.7}, and \ref{cdp:st:D.1.12} identifies this target
with the fibre of $\mathscr E$. Three real independent
vectors cannot lie in one complex line. Their central
complex span thus still has dimension two. Two of the
constant period vectors forming a basis consequently
extend as a frame of $\mathscr E$ across zero.

\item These local frames agree with the same constant period
frame on the smooth base. They therefore glue to
$\mathscr E\simeq\OO^{\oplus2}$ on $\PP^1$.
The cup-product isomorphism in \ref{cdp:st:D.1.12} gives
$R^2f_*\OO_Y\simeq\det\mathscr E\simeq\OO$.
Connected fibres give $f_*\OO_Y=\OO$. Leray on $\PP^1$
proves all four dimensions in \eqref{cdp:eq:D.33}, and evaluation
of the trivial rank-two bundle, followed by the
base-change isomorphism, proves coherent surjectivity.

\item The global integral sections are the intersection of
the three annihilators \eqref{cdp:eq:D.36} inside the constant
smooth-fibre lattice. Their defining equations are
$\lambda_3=0$, $\lambda_2+\lambda_4=0$,
$\lambda_1+\lambda_4=0$. The primitive solution is
$(1,1,0,-1)$, giving \eqref{cdp:eq:D.34}. The three singular fibres have two components each. Thus
$r=6$, $s=3$, and $t'=4$ give $6=3-1+4$.

\item Finally $\Pic(Y)=H^1(Y,\OO_Y)$ by the integral vanishing in \cref{cdp:st:D.2.7}
and the universal coefficient theorem. Under its coherent restriction
identification with $H^1(A,\OO_A)$, a line bundle is
torsion on a smooth fibre exactly when some positive
integral multiple of its class belongs to
$\iota(\Lambda^*)$, by the exponential sequence.
This is exactly \eqref{cdp:eq:D.35}. At the singular fibre $S_i$
the same exponential sequence and \eqref{cdp:eq:D.36} replace
$\Lambda^*$ by its sublattice $\operatorname{ann}(a_i)$.
That torsion locus is contained in \eqref{cdp:eq:D.35}, so the union
over all fibres is exactly the displayed set. The
real-linear map $\iota:\R^4\to\C^2$ is an isomorphism;
the image of $\Q^4$ is therefore countable and dense.
Its complement is dense and has full Lebesgue measure,
but has no nonempty open subset. 

\end{enumerate}
\end{proof}

\section{Fibration stability}\hypertarget{tag-034A}{}\label{sec:II:14}

\begin{proposition}[The infinitesimal map and the rank-two boundary]\hypertarget{tag-01GU}{}\label{cdp:st:E5}
\StatementRecord{cdp:st:E5}{proved}{01GU}{the definitions and constructions stated in the proof}

Let $g:Y\to C$ be a proper surjective holomorphic map with connected fibres
from a smooth connected compact complex threefold to a nonsingular compact
connected complex curve. Its smooth general fibre is a two-dimensional complex
torus. There is a canonical exact sequence of vector spaces
\begin{equation}\hypertarget{tag-019Y}{}\label{cdp:eq:E-Leray-map}
0\to H^1(C,T_C)\xrightarrow{g^*}H^1(Y,g^*T_C)
 \to H^0(C,T_C\otimes R^1g_*\mathcal O_Y)\to0.
\end{equation}
On the smooth locus $U\subset C$, the relative tangent sequence
induces a sheaf morphism
\[
 \delta:T_C\otimes R^1g_*\mathcal O_Y|_U
                  \longrightarrow R^2g_*T_{Y/C}|_U.
\]
At a smooth torus fibre $F$, write its Kodaira--Spencer tensor as
$\kappa=u_1\otimes\eta_1+u_2\otimes\eta_2$, where
$\eta_1,\eta_2$ is a basis of $H^1(F,\mathcal O_F)$ and
$u_i\in\operatorname{Lie}(F)$. Then
\[
 \delta(a\eta_1+b\eta_2)
   =(bu_1-au_2)\otimes(\eta_1\wedge\eta_2).
\]
Its rank is the rank of $\kappa$. When this rank is two at one
regular value, the image of $H^1(Y,T_Y)\to H^1(Y,g^*T_C)$ lies
in $g^*H^1(C,T_C)$. For $C=\mathbb P^1$ this image is zero.
\end{proposition}


\begin{proof}

\begin{enumerate}[label=\arabic*.,leftmargin=*]
\item The projection formula gives
$R^qg_*g^*T_C=T_C\otimes R^qg_*\mathcal O_Y$, and the coherent
Leray spectral sequence has only columns zero and one. Its
degree-one edge sequence is exactly \eqref{cdp:eq:E-Leray-map}.
On $U$ the sequence
$0\to T_{Y/C}\to T_Y\to g^*T_C\to0$ is exact. Its extension
class on $F$ is the Kodaira--Spencer tensor. The connecting map is
cup product with that class; applying alternating multiplication to
the two displayed basis vectors gives the formula and its rank.
Exactness of the higher-direct-image sequence says that the image
of a global tangent class restricts to the kernel of $\delta$.
Rank two holds on an open dense subset after it holds at one point.
By \cite[Theorem~2.2, pp.~4--5]{Viaclovsky2026Combs}, the source $T_C\otimes R^1g_*\mathcal O_Y$ is locally free,
so a section zero on that subset is zero everywhere. This proves
the image assertion and the final statement.

\item At a critical point the sheaf map $dg$ need not be onto. Its
image is $\mathcal I_{\rm crit}\otimes g^*T_C$, giving
\[
0\to T_{Y/C}\to T_Y\to
 \mathcal I_{\rm crit}\otimes g^*T_C\to0.
\]
For $t=z^m$, the image ideal is $(z^{m-1})$; for $t=xy$ it is
$(x,y)$; for $t=xyz$ it is $(xy,xz,yz)$. On a root cover of a
smooth multiple fibre the ordinary tangent sequence for $z$ applies.
At a normal-crossing fibre the logarithmic map sends
$\sum a_i z_i\partial_{z_i}$ to
$(\sum a_i)t\partial_t$ and is onto the logarithmic base tangent
line. These formulas describe the image at the special fibres.

\end{enumerate}
\end{proof}

\begin{proposition}[Stability of a rank-two torus fibration]\hypertarget{tag-01GV}{}\label{cdp:st:E6}
\StatementRecord{cdp:st:E6}{proved}{01GV}{\SeriesRef{cdp:st:E5}, \SeriesRef{cdp:eq:E-surj-embedding}, \SeriesRef{cdp:eq:E-Leray-map}, \SeriesRef{cdp:st:E1}; \cite[\S4, equations (4.3)--(4.5), p.~649; \S6, Theorem~6.1 and proof, pp.~660--661]{Horikawa1974}; \cite[Section~9, Theorems~1--2, pp.~83--85]{Douady1966}; \cite[Section~10.1, Proposition~1, p.~86]{Douady1966}; \cite[Section~10.2, Theorem~1, p.~87]{Douady1966}; \cite[Theorem~11, p.~297]{Buchdahl1999}}

Let $g:Y\to\mathbb P^1$ be as in \ref{cdp:st:E5} and suppose its
Kodaira--Spencer tensor has rank two at one regular value. For every
proper deformation $\mathcal Y\to(S,0)$ of $Y$ over a complex
analytic germ, including a nonreduced germ, there is, after shrinking
$S$, a holomorphic map
\[
 G:\mathcal Y\longrightarrow\mathbb P^1\times S
\]
over $S$ whose central restriction is $g$. Every $G_s$ is surjective
with connected fibres and smooth general fibre a complex two-torus.
With $b_2(Y)=0$, every sufficiently small deformation has algebraic
dimension one.
\end{proposition}


\begin{proof}

\begin{enumerate}[label=\arabic*.,leftmargin=*]
\item We use the \v{C}ech construction of holomorphic-map deformations
in \cite[\S4, equations (4.3)--(4.5), p.~649; \S6, Theorem~6.1 and proof, pp.~660--661]{Horikawa1974};
we spell out the portions needed here. The injectivity hypothesis
on degree-two tangent cohomology in its Theorem~6.1 is replaced
here by the explicitly proved unobstructedness of the torus source.

\item First fix a rank-two regular fibre and its embedding $i:F\hookrightarrow
Y$. The normal line is canonically $N=\mathcal O_F\otimes T_c\mathbb
P^1$. The tangent sequence and \ref{cdp:st:E5} imply that
\begin{equation}\hypertarget{tag-01A0}{}\label{cdp:eq:E-surj-embedding}
 H^1(F,T_F)\longrightarrow H^1(F,i^*T_Y)
 \quad\hbox{is surjective}.
\end{equation}
Let $A'\to A$ be a small extension of local Artin complex algebras,
with kernel $I$ annihilated by the maximal ideal. Suppose an embedding
$i_A:F_A\hookrightarrow Y_A$ deforming $i$ and a deformation
$Y_{A'}$ lifting $Y_A$ have been specified. The deformation functor
of a complex torus is smooth: with a real lattice marking its small
complex structures are graphs of arbitrary small complex linear maps
between the two complementary Hodge subspaces; constant Beltrami
tensors have zero bracket. This gives a smooth four-dimensional
miniversal germ for $F$. In particular $F_A$ has a lift $F_{A'}$.

\item Choose local lifts of $i_A$ on sufficiently small coordinate charts.
They exist because the target deformation is smooth over $A'$.
Their differences on overlaps form a one-cocycle in
$i^*T_Y\otimes I$. Changing the lifted transition functions of
$F_{A'}$ by a cocycle $v_{ij}\in T_F\otimes I$ changes this
obstruction by $di(v_{ij})$. Such changes are precisely the action
of $H^1(F,T_F)\otimes I$ on the set of lifts of $F_A$; their
cocycle law has no quadratic term because $I\mathfrak m_{A'}=0$.
Surjectivity \eqref{cdp:eq:E-surj-embedding} therefore makes the map
obstruction zero. A zero-coboundary correction glues the local maps.
The resulting proper map is a closed immersion: the map on local
rings is surjective modulo the nilpotent ideal, hence surjective
by Nakayama, and its underlying central map is the embedding $i$.
Thus the given embedded torus lifts in every ambient small extension.

\item Now suppose $g_A:Y_A\to\mathbb P^1_A$ has been lifted and that
$Y_{A'}$ is prescribed. Choose local lifts $t_i$ of the target
coordinate of $g_A$; on overlaps their discrepancies define the
holomorphic-map obstruction
\[
 \alpha\in H^1(Y,g^*T_{\mathbb P^1})\otimes I.
\]
Coordinate changes on the target give exactly the transition factors
of $g^*T_{\mathbb P^1}$, and the triple-overlap relation gives its
\v{C}ech cocycle law. A coboundary makes the local maps agree, so
$\alpha=0$ is sufficient as well as necessary.

\item For a rank-two regular value $c$, the divisor
$F_A=g_A^{-1}(c)$ is a smooth deformation of the central torus.
The preceding embedding argument provides a Cartier divisor
$F_{A'}\subset Y_{A'}$ lifting it. This forces the restriction
of $\alpha$ to $F$ to vanish, by the following equality of cocycles: choose local equations $t_i-\eta_i$ for
$F_{A'}$, where $\eta_i$ lies in $I\otimes\mathcal O_Y$.
Two such equations differ by a unit. Modulo the central equation
$t=0$, that unit difference disappears, and their overlap equality
reads
\[
 (t_i-t_j)|_F=\eta_i|_F-\eta_j|_F.
\]
This is exactly the assertion that the restricted obstruction
cocycle is a coboundary in $N\otimes I$.

\item By \eqref{cdp:eq:E-Leray-map} and $H^1(\mathbb P^1,T_{\mathbb P^1})=0$,
$\alpha$ is a section of the locally free sheaf
$T_{\mathbb P^1}\otimes R^1g_*\mathcal O_Y$, tensored with $I$.
Its evaluation at every rank-two regular fibre is zero. That set
is open dense; hence the section and $\alpha$ are zero. The map
therefore lifts across every small extension.

\item We pass from the formal assertion to an analytic map and
specify the relative parameter space, including its nilpotents.
Let $H$ be the germ at the graph of $g$ in the absolute Douady
space of $\mathcal Y\times\mathbb P^1$, and let
$\pi:\mathcal R\to H$ be its universal compact flat family.
Douady's existence and finiteness theorems
\cite[Section~9, Theorems~1--2, pp.~83--85]{Douady1966}
apply to this analytic ambient space; its local rings and those
of the parameter space may have nilpotents.

\item After shrinking $H$, one has
\[
 \pi_*\mathcal O_{\mathcal R}=\mathcal O_H,
\]
compatibly with base change. Over the full analytic germ, represent $R\pi_*\mathcal O_{\mathcal R}$
near the central point by a finite complex of finite free analytic
modules. On the central fibre the unit spans
$H^0(Y,\mathcal O_Y)=\mathbb C$. Split off this unit in degree
zero. The differential is injective modulo the maximal ideal on
the complementary degree-zero summand; an invertible minor
cancels that summand with its image in degree one. The resulting
complex begins with $\mathcal O_H$ in degree zero and zero
differential from it. This proves the displayed equality and its
base-change assertion on the full analytic germ.

\item Embed a representative of $(S,0)$ in a polydisc. The coordinate
functions of the map $\mathcal R\to\mathcal Y\to S$ are now
pullbacks of unique holomorphic functions on $H$. They satisfy
the defining equations of $S$, because their pullbacks do and the
unit map just proved is injective. They therefore define a
holomorphic morphism $H\to S$, and the universal subspace lies
scheme-theoretically in
\[
 (\mathcal Y\times_S H)\times\mathbb P^1.
\]
The projection of this subspace to $\mathcal Y\times_S H$ is a
proper morphism between two proper flat families over $H$ and
is an isomorphism at the central point. By the isomorphism-locus
criterion \cite[Section~10.1, Proposition~1, p.~86]{Douady1966},
it is an isomorphism after restricting to an open neighbourhood
of that point. The proof of that criterion uses the finite
direct image, flatness and Nakayama's lemma, and applies to
nonreduced $H$. On this open locus the universal subspace is
the graph of a holomorphic map to $\mathbb P^1$.

\item The universal property of the absolute Douady space consequently
identifies this open locus with the germ $\mathcal H\to S$ of
relative maps from $\mathcal Y/S$ to
$\mathbb P^1\times S/S$. This is the relative version of
Douady's graph construction
\cite[Section~10.2, Theorem~1, p.~87]{Douady1966}; the argument
above supplies the variation of the compact source explicitly.
The just-proved lifting property for every
small Artin extension says that this morphism is formally smooth
at $[g]$. The local analytic smoothness criterion identifies a
formally smooth morphism of finite-dimensional analytic germs with
a projection $S\times(\mathbb C^r,0)\to S$ after shrinking. One
can see the criterion on presentations of their convergent local
algebras: lifting across all square-zero quotients makes the
relative conormal obstruction zero, and the analytic implicit
function theorem eliminates the corresponding equations. The
zero section of this product supplies the required holomorphic
family of graphs. This argument applies to nonreduced $S$.

\item At one regular value, the family $G$ remains a proper submersion
near the central compact fibre. Its nearby fibres are connected
deformations of a torus. Each map $G_s$ is nonconstant, and hence
surjective. Its Stein factorization has finite degree one because
it has a connected regular fibre in this neighbourhood. Thus
every fibre is connected. Over the connected smooth locus of
$G_s$, Ehresmann's theorem identifies the differentiable fibres
with $T^4$. Every compact complex surface with this differentiable
type is a torus: its even first Betti number makes it K\"ahler
by \cite[Theorem~11, p.~297]{Buchdahl1999};
its Albanese has dimension two; the exterior algebra on its four
integral degree-one classes shows the Albanese map has degree one;
and a nontrivial exceptional curve of this degree-one map would
include a $(-1)$-curve, impossible for the even intersection form
of $T^4$. Hence the general fibre is a complex two-torus.
Finally $b_2$ is constant in the differentiable deformation family,
and \ref{cdp:st:E1} proves $a(Y_s)=1$.

\end{enumerate}
\end{proof}

\begin{proposition}[Deformations of the constructed torus fibrations]\hypertarget{tag-01GW}{}\label{cdp:st:E7}
\StatementRecord{cdp:st:E7}{proved}{01GW}{\SeriesRef{cdp:st:E6}, \SeriesRef{cdp:eq:E-log-period}, \SeriesRef{cdp:st:E1}; \cite[Proposition~6.1, pp.~19--21]{Viaclovsky2026Combs}; \cite[Proposition~3.1, pp.~6--7]{Viaclovsky2026Two}; \cite[Proposition~2.5, p.~5]{LiuRollenske2026}; \cite[Proposition~4.1, case~(i), pp.~11--12]{Viaclovsky2026Combs}}

Every member of the following torus-fibred families,
including the period family $X_c$, its members with $|p|>1$,
Viaclovsky's $X_{3k,4l}$, his two-parameter family and his combs,
and the Liu--Rollenske family, satisfies the rank hypothesis of \ref{cdp:st:E6}.
Consequently every sufficiently small deformation retains a
holomorphic torus fibration over $\mathbb P^1$. Every such member
with $b_2=0$, and every sufficiently small deformation of it, has
algebraic dimension one. In particular the fibration-retention
question in \cite[Section~10(1), p.~33]{Viaclovsky2026Combs} has an
affirmative answer.
\end{proposition}


\begin{proof}

\begin{enumerate}[label=\arabic*.,leftmargin=*]
\item At their hexagonal fibre, a marking on the punctured disc puts
the Hodge periods in the form
\begin{equation}\hypertarget{tag-019Z}{}\label{cdp:eq:E-log-period}
  (I_2,Z(q)),\qquad
 Z(q)=\frac{\log q}{2\pi i}B+H(q),
 \qquad \det B\ne0,
\end{equation}
where $H$ is holomorphic at zero. Therefore
\[
 qZ'(q)=\frac{B}{2\pi i}+qH'(q)
\]
is invertible for all sufficiently small nonzero $q$. The
Kodaira--Spencer tensor of a marked torus with periods $(I,Z)$
is represented by $Z'$ after the invertible identifications of
the Lie algebra and $H^1(\mathcal O)$ with the chosen period
coordinates. Thus it has rank two at those regular fibres.

\item The local period models occur in
\cite[Proposition~6.1, pp.~19--21]{Viaclovsky2026Combs},
\cite[Proposition~3.1, pp.~6--7]{Viaclovsky2026Two}, and
\cite[Proposition~2.5, p.~5]{LiuRollenske2026}.
The AC cusp has logarithmic matrix
$B_0=\left(\begin{smallmatrix}0&1\\-1&0\end{smallmatrix}\right)$.
The affine twists and finite logarithmic transformations change
the finite collars and constant gluing data, while the punctured
cusp matrix still has this same invertible logarithmic coefficient.
Engel's two deck transformations at the $I_1$-fibre have independent
valuation vectors $(1,0)$ and $(0,1)$, again giving an invertible
$B$. In the two-parameter family the designated fibre over $p_3$
is the full-lattice rank-two filling; its displayed two valuation
vectors are independent. Liu--Rollenske's compactification over
$I_1$ uses the same two independent valuation directions. Under the covering
$(z_1,z_2)\mapsto(e^{2\pi iz_1},e^{2\pi iz_2})$, lifts of the
angular loops give the two constant period columns, and logarithms
of the two deck multipliers, divided by $2\pi i$, give the columns of $Z(q)$.

\item For the combs, \cite[Proposition~4.1, case~(i), pp.~11--12]{Viaclovsky2026Combs}
uses the flat cohomology basis $\alpha_i,\beta_i$ and the
single-valued holomorphic frame
$e_i=\alpha_i+\ell(q)\beta_i$, $f_i=\beta_i$, with
$\ell(q)=\log q/(2\pi i)$. The extended Hodge subbundle has
the rows $(I,H(q))$ in the frame $(e_i,f_i)$. The fixed value of the
smoothed period section at $q=0$ keeps it transverse to the
$f_i$-plane. Its rows therefore remain $(I,H_s(q))$ with
$H_s$ holomorphic. In the flat basis $(\alpha_i,\beta_i)$ the rows are
$(I,\ell(q)I+H_s(q))$, which is \eqref{cdp:eq:E-log-period} with
$B=I$.

\item All the sources construct the indicated members as compact
connected-fibre holomorphic maps with smooth general two-torus
fibre. \ref{cdp:st:E6} applies, and \ref{cdp:st:E1} gives the algebraic-dimension assertion.
On the differentiable six-sphere these neighbourhoods of the
listed structures constitute an open set on which $a=1$.

\end{enumerate}
\end{proof}

\begin{proposition}[Deformations over a fixed genus-two curve]\hypertarget{tag-01GH}{}\label{cdp:st:E10}
\StatementRecord{cdp:st:E10}{proved}{01GH}{\SeriesRef{cdp:eq:E-Leray-map}, \SeriesRef{cdp:st:E5}}

There is a smooth projective family of compact threefolds
$Y_s$ with connected-fibre torus fibrations $g_s:Y_s\to C_s$
over smooth genus-two curves, whose Kodaira--Spencer tensors
have rank two at a regular value, such that the central map
$g_0$ has no holomorphic extension from this deformation to
the fixed curve $C_0$.
\end{proposition}


\begin{proof}

\begin{enumerate}[label=\arabic*.,leftmargin=*]
\item Let $S\to\mathbb P^1_t$ be the rational elliptic surface obtained
by blowing up the nine base points of the Hesse pencil
$x^3+y^3+z^3-3txyz=0$; see \citet[\S2, pp.~238--241]{ArtebaniDolgachev2009}. Its critical values are
$\{1,\omega,\omega^2,\infty\}$, where $\omega^3=1$, $\omega\ne1$.
Pull another copy back by $\phi(t)=1/(t-4)$; its critical values
are $\{4,5,4+\omega,4+\omega^2\}$, disjoint from the first set.
Their fibre product $Z\to\mathbb P^1$ is a smooth projective
threefold, since at every point at least one elliptic-surface
map is a submersion. Its general fibre is a product of two
elliptic curves. Both elliptic $j$-functions are nonconstant,
so their two Kodaira--Spencer derivatives are simultaneously
nonzero off a finite set. The tensor of their product is then
diagonal of rank two.

\item Let $C_s$ be the smooth genus-two double cover
\[
 y^2=(t-10)(t-11)(t-12)(t-13)(t-14)(t-15-s)
\]
of $\mathbb P^1$, for $s$ in a small disc. All its branch values
avoid the critical values of $Z$. Thus
$Y_s=Z\times_{\mathbb P^1}C_s$ is a smooth projective deformation
of a compact threefold, and $g_s:Y_s\to C_s$ has connected torus
general fibre and rank two at a regular value. The Kodaira--Spencer
class $\theta_C\in H^1(C_0,T_{C_0})$ is nonzero: any infinitesimal
isomorphism of the hyperelliptic covers descends through their
unique hyperelliptic involutions to an infinitesimal projective
transformation fixing five distinct branch values; it is zero,
whereas the sixth branch value moves with derivative one.

\item Put
$L=\pi_C^*\mathcal O_{\mathbb P^1}(1)$ on $C_0$; the
double-cover canonical formula gives $K_{C_0}=L$, of degree two.
For either rational elliptic surface, relative duality gives
$R^0\pi_*\mathcal O=\mathcal O$ and
$R^1\pi_*\mathcal O=\mathcal O(-1)$: its relative canonical
line is $\pi^*\mathcal O(1)$ and every fibre has
$h^\bullet(\mathcal O)=(1,1)$. Relative K\"unneth and the flat
base change to $C_0$ consequently give
\[
 R^qg_*\mathcal O_Y=
 \begin{cases}
  \mathcal O_{C_0}&q=0,\\
  L^{-1}\oplus L^{-1}&q=1,\\
  L^{-2}&q=2,\\
  0&q>2.
 \end{cases}
\]
The vertical tangent sheaf is $g^*(L^{-1}\oplus L^{-1})$:
the invariant vector on each elliptic factor generates its
relative tangent sheaf, also at a node, where the kernel of
$d(xy)$ is generated by $x\partial_x-y\partial_y$.
Riemann--Roch on the genus-two curve gives
$h^0(L^{-a})=0$ and $h^1(L^{-a})=2a+1$ for $a>0$.
Projection and Leray give
\[
\begin{array}{c|rrrr}
 \mathcal F&h^0&h^1&h^2&h^3\\\hline
 \mathcal O_Y&1&2&6&5\\
 g^*T_{C_0}&0&3&10&7\\
 T_{Y/C_0}&0&6&20&14
\end{array}
\]
All higher groups vanish. In particular
$H^0(C_0,T_{C_0}\otimes R^1g_*\mathcal O_Y)=0$ and the
injection $H^1(C_0,T_{C_0})\to H^1(Y,g^*T_{C_0})$ is an
isomorphism between three-dimensional spaces.

\item Compatibility of the \v{C}ech transition equations for the family
$g_s$ gives
\[
 dg(\theta_Y)=g^*\theta_C
       \quad\text{in }H^1(Y_0,g^*T_{C_0}).
\]
The right side is nonzero because the first map of
\eqref{cdp:eq:E-Leray-map} is injective. Hence the first-order
obstruction to extending $g_0$ with fixed target $C_0$ is nonzero.
The obstruction belongs to
the base-deformation summand of \ref{cdp:st:E5}.

\end{enumerate}
\end{proof}

\section{Deformations of explicit torus fibrations}\hypertarget{tag-P268}{}\label{sec:II:explicit-deformations}

\begin{proposition}[All coprime $(3k,4l)$ multiplicities]\hypertarget{tag-01GI}{}\label{cdp:st:E11}
\StatementRecord{cdp:st:E11}{proved}{01GI}{\SeriesRef{cdp:st:E7}; \cite[Theorem~1.1, p.~1]{Viaclovsky2026Log}}

Fix admissible logarithmic-transformation data with multiplicities
$r=3k$, $s=4l$, where $k,l>0$ and $\gcd(r,s)=1$, in the family
constructed from the original elliptic-surface data. For every member
which is diffeomorphic to $S^6$, as in
\cite[Theorem~1.1, p.~1]{Viaclovsky2026Log},
\[
 h^\bullet(X,T_X)=(1,1,1,0),
\]
its Kuranishi germ is a smooth reduced line, and varying the nonzero
linearization scale with those data fixed is locally miniversal.
For the original $(3,4)$ affine-twist family, the same statements
hold whenever $p\ne0$, including $|p|>1$.
\end{proposition}


\begin{proof}

\begin{enumerate}[label=\arabic*.,leftmargin=*]
\item Put $\mathfrak B_{r,s}=\mathbb P(r,s)$. The root morphism
$\sigma:\mathfrak B_{r,s}\to\mathfrak B_{3,4}$ over the identical
coarse $\mathbb P^1$ sends the two root sections to their $k$th and
$l$th powers. In weighted homogeneous coordinates it is
$[x:y]\mapsto[x^l:y^k]$. In line notation
$\sigma^*\mathcal O_{\mathfrak B_{3,4}}(1)
=\mathcal O_{\mathfrak B_{r,s}}(kl)$.
On each fixed smooth root cover, the derivative of an affine
deck map is its linear part. Its translation acts as the identity
on invariant Lie fields and on fibre $\mathcal O$-cohomology.
The induced equivariant bundle comparisons are therefore unchanged. The same
cusp comparison used in \SeriesRef{fam:st:A.5.2} therefore identifies the bundles
\[
 E=G=\mathcal O\oplus\mathcal O(-kl),\qquad
 R^2\bar f_*\mathcal O_X=\mathcal O(-kl),
\]
and
\[
0\to\bar f^*E\to\mathcal T:=T_X(-\log W)
 \to\bar f^*L\to0,
\quad L=\mathcal O(r+s-rs),
\quad K_X=\bar f^*\mathcal O(kl-r-s).
\]
For example the normal-cover coordinate at the first finite point
is $u=v^k$, and differentiation changes the ordinary root tangent
line from the old one to $T_{\mathfrak B_{r,s}}$; the vertical
bundle and its cohomology are pulled back through $u=v^k$.
At the cusp $W$ remains the same
normal-crossing divisor, so its logarithmic base correction is
$-p_0=-rs$ and its normal quotient is $\eta_*K_{\mathrm dP_6}$.

\item Locally $\sigma_*\mathcal O=\mathcal O$ and its positive direct
images vanish: the additional cyclic group has invariants
$\mathbb C\{v\}^{\mu_k}=\mathbb C\{u\}$, and averaging is
exact. Consequently the cohomology of the pulled-back lines
$\mathcal O(-kl)$ and $\mathcal O(-2kl)$ is the cohomology of
the original $\mathcal O(-1)$ and $\mathcal O(-2)$, hence zero.
Projection and Leray give
$h^\bullet(X,\bar f^*E)=(1,1,0,0)$.
For the horizontal line, its coarse direct image is
$\mathcal O_{\mathbb P^1}(-1)$: a vector field on a root disc
must have its coefficient divisible by the root parameter,
so $\rho_*T_{\mathfrak B_{r,s}}=T_{\mathbb P^1}(-p_1-p_2)$,
and the extra $-p_0$ gives the asserted degree. Moreover
$L$ and $L\otimes\mathcal O(-kl)$ have negative orbifold degree,
so have no global sections. Thus
$H^0(X,\bar f^*L)=H^1(X,\bar f^*L)=0$.
The two exact tangent sequences now give
$h^0(X,T_X)=h^1(X,T_X)=1$.

\item On the smooth torus locus,
the boundary on a relative holomorphic one-form is the transpose
of the Kodaira--Spencer tensor. The latter is generically invertible
by \ref{cdp:st:E7}. Hence every ambient holomorphic one-form on an inverse-image
base disc is horizontal. It is the pullback of a base one-form.
At a finite root point the root map is a submersion, so that
one-form extends holomorphically in the root coordinate. At a
generic smooth point of $W$ its local equation is $q=z_1$, which
excludes a pole in its base coefficient. Therefore the canonical
pullback comparison gives
\[
 \bar f_*\Omega_X^1=\Omega_{\mathfrak B_{r,s}}^1.
\]
By Serre duality and the projection formula,
\[
 h^3(X,T_X)=h^0\bigl(\mathfrak B_{r,s},
                         \mathcal O(kl-2r-2s)\bigr)=0.
\]
For the last equality, a weighted monomial of that degree would
give nonnegative integers $a,b$ with
$3k(a+2)+4l(b+2)=kl$. Reduction modulo $k$ makes $b+2$ a positive
multiple of $k$, since $\gcd(k,4l)=1$. The second term alone
is then at least $4kl$, a contradiction. On a six-sphere
$\chi(T_X)=1$, giving $h^2=1$. For the original affine-twist
family, $r=3$, $s=4$, and the horizontal Leray terms are the
acyclic lines $\mathcal O(-5)$ and $\mathcal O(-6)$ on
$\mathbb P(3,4)$. Therefore $h^\bullet(\mathcal T)=(1,1,0,0)$,
and the quotient $\eta_*K_{\mathrm dP_6}$ gives
$h^\bullet(T_X)=(1,1,1,0)$ directly for every admissible affine
twist.

\item Use the period quotient and reference linearization of
\SeriesRef{fam:st:B.1.1} and \SeriesRef{fam:st:B.1.4}. They give
$u=e^{-2\pi ic}$ and
$[\zeta_1,\zeta_2]\mapsto[\zeta_1,e^{2\pi i\zeta_2}]$.
Thus $e=\partial_{\zeta_2}$ maps to $2\pi i$ times the line-bundle
Euler field $\Xi$. Let $\psi$ count the forward generator
$\Phi_u=u\phi$ covering translation by $-6P$, and let
$\gamma=-\psi$. The inverse generator counts the clockwise
deck direction: inversion induces $n\mapsto-n$ on its deck
group $\mathbb Z$, which is its own inverse. For the $(3,4)$
family this is the marking of
\cref{def:II:marked-Picard-parameter}.
For $\kappa(\Phi)=(\partial\Phi)\Phi^{-1}$, differentiation gives
\[
 \kappa(\partial_u)=u^{-1}\Xi\otimes[\psi]
       =-(2\pi iu)^{-1}e\otimes[\gamma],\qquad
 \kappa(\partial_c)=e\otimes[\gamma].
\]
Here the fibre classes extend by the coherent Leray comparison of
\SeriesRef{fam:st:A.5.2}; they are nonzero on every regular torus
fibre. A first-order trivialization of
the compact member would turn the given fibration into a
first-order deformation of the fixed map $f$ with fixed source.
Such deformations are $H^0(X,f^*T_{\mathbb P^1})
=H^0(\mathbb P^1,T_{\mathbb P^1})$, by connected fibres and
the projection formula. Thus they are precisely infinitesimal
projective transformations of the base. The two multiple critical
values and the hexagonal critical value remain the same three
constant sections in this scale family; the local equations
$t=v^r$, $t-1=v^s$, $q=xyz$ show this also over the dual numbers.
The base vector field is therefore zero. The nonzero torus class
cannot be trivialized, proving the nonzero global scale derivative.
Write the Kuranishi local ring as $\mathbb C\{z\}/I$. The scale
classifying map has nonzero derivative, so before quotienting its
substitution into $\mathbb C\{u-u_0\}$ is a local analytic
isomorphism. Every element of $I$ maps to zero; hence $I=0$. The
classifying map is locally an isomorphism, proving smoothness,
reducedness and miniversality.

\end{enumerate}
\end{proof}

\begin{proposition}[The Kuranishi germ of the two-parameter family]\hypertarget{tag-01GK}{}\label{cdp:st:E13}
\StatementRecord{cdp:st:E13}{proved}{01GK}{\SeriesRef{cdp:eq:E13-connection}, \SeriesRef{cdp:eq:E13-residue}, \SeriesRef{cdp:st:E1}, \SeriesRef{cdp:st:E6}; \cite[Proposition~6.4, pp.~19--20]{Viaclovsky2026Two}}

For every member of the family of
\cite[Proposition~6.4, pp.~19--20]{Viaclovsky2026Two}, on the parameter
bidiscs with $\mu$ sufficiently close to zero and nonzero sufficiently
small linearization scale $\lambda$, one has
\[
 (h^0,h^1,h^2,h^3)(X_{\mu,\lambda},T_X)=(1,2,2,0).
\]
Its Kuranishi germ is a smooth reduced surface, and the displayed
$(\mu,\lambda)$-family is locally miniversal. In particular its
sufficiently small deformations all have algebraic dimension one.
\end{proposition}


\begin{proof}

\begin{enumerate}[label=\arabic*.,leftmargin=*]
\item The construction uses the equation and section
\[
 S_\mu:\quad y^2=x^3+tx+\mu t+1,\qquad
 P_\mu=(-\mu,b),\qquad b^2=1-\mu^3,
\]
where the branch of $b$ is close to one. Write $p_1,p_2,p_3$ for its
three distinct nodal values, with the zero of the linearization at
$p_3$. The order-four good-reduction cover at infinity gives the
root curve $\mathfrak B=\mathbb P(1,4)$, with coarse projection
$\rho$ and $\rho^*\mathcal O_{\mathbb P^1}(1)=\mathcal O(4)$.
Let $\bar f:X\to\mathfrak B$ denote the root factorization and
$D=W_1+W_2+W_3$ the reduced divisor of Mumford fibres.

\item The vertical logarithmic Lie bundle is
\[
 E=\mathcal O\oplus\mathcal O(-1).
\]
The first summand is generated by multiplication in the line-bundle
coordinate. For a concrete splitting of the second summand put
$\omega=dx/(2y)$ and use the relative differential
\begin{equation}\hypertarget{tag-01A5}{}\label{cdp:eq:E13-connection}
 \alpha=\frac{y+b}{x+\mu}\,\frac{dx}{2y}.
\end{equation}
Its residues are $1$ at $P_\mu$ and $-1$ at the zero section; there
is no pole at $-P_\mu$, because its numerator vanishes there.
Consequently, in the rational section of
$\mathcal O(P_\mu-O)$ having divisor $P_\mu-O$, the connection
$\nabla(hs)=(dh+\alpha h)s$ is holomorphic along each smooth
elliptic fibre. Here $s$ is that rational section and a point of
the line bundle is written $v s$; the horizontal lift of a vector
$V$ is therefore $\widetilde V=V-\alpha(V)v\partial_v$.
It is a multiplicative connection: the difference between its
pullback under multiplication and the sum of its two pullbacks is a
holomorphic one-form on the product of two elliptic curves, and its
restrictions to both coordinate axes vanish. Every holomorphic
one-form on that product is the sum of forms pulled back from those
axes, so this difference is zero. Its horizontal lift therefore
commutes with the lifted translation defining $X$.

\item At infinity put $t=r^{-4}$, $x=r^{-2}x'$, $y=r^{-3}y'$. Formula
\eqref{cdp:eq:E13-connection}, in the relative cotangent sheaf, becomes
\[
 \alpha=\frac{y'+br^3}{x'+\mu r^2}\frac{dx'}{2y'},
 \qquad y'^2=x'^3+x'+\mu r^2+r^6.
\]
This is the same holomorphic connection at the good-reduction
fibre; the limiting section $(0,0)$ is a nonzero two-torsion point,
so it remains disjoint from the zero section. Moreover
$\omega=r\,dx'/(2y')$, which gives the summand $\mathcal O(-1)$
for its dual Lie vector. At a nodal fibre the connection is
logarithmic at the two branches and its lift is a section of the
vertical logarithmic tangent bundle. These connection and transition maps give the splitting on the
entire root curve.

\item The inverse pairing associated with the nondegenerate alternating
class $\xi=u\wedge w+2\psi\wedge\delta$ in the construction induces a
holomorphic isomorphism
\[
 G:=R^1\bar f_*\mathcal O_X\simeq E.
\]
On a marked smooth fibre the periods
have cross terms in ratio $1:-2$, so the holomorphic one-form
subspace $F^1\subset H^1(F,\mathbb C)$ is Lagrangian for this pairing.
Pairing with $F^1$ therefore gives the isomorphism
$H^1(F,\mathbb C)/F^1\to(F^1)^*=\operatorname{Lie}(F)$.
At a Mumford point, subtract the logarithmic monodromy term from the
periods. In the resulting holomorphic frame the alternating pairing
is still constant and nondegenerate, and the two holomorphic
one-forms are the logarithmic coordinate differentials. Their dual
frame is the logarithmic Lie frame. The quotient frame gives
$R^1\bar f_*\mathcal O_X$: at a rank-one point its first class is
the class of the loop in the normalization diagram and its second
is the class on the elliptic normalization; at the rank-two point
both classes are loops in the normalization diagram. These are
bases, by the normalization exact sequence. Thus the pairing extends
as an isomorphism at these points. The finite good-reduction cover
gives the same conclusion at infinity, and descent uses invariance
of $\xi$. Cup product and the normalization sequences
give $R^2\bar f_*\mathcal O_X=\det G=\mathcal O(-1)$.

\item The logarithmic tangent and canonical sequences are therefore
\[
\begin{split}
 0&\longrightarrow\bar f^*E\longrightarrow T_X(-\log D)
 \longrightarrow\bar f^*\mathcal O(-7)\longrightarrow0,\\
 K_X&=\bar f^*\mathcal O(-4).
\end{split}
\]
Indeed $K_{\mathfrak B}=\mathcal O(-5)$ and
$T_{\mathfrak B}(-p_1-p_2-p_3)=\mathcal O(5-12)$; taking the
determinant of the displayed logarithmic sequence gives the
canonical formula after removing $D$ of degree $12$.
For root line bundles, cohomology is computed by weighted monomials
and Serre duality with $\mathcal O(-5)$. The projection formula and
Leray give
\[
\begin{split}
 h^\bullet(X,\bar f^*E)&=(1,1,0,0),\\
 h^\bullet(X,\bar f^*\mathcal O(-7))&=(0,1,2,1),\\
 h^\bullet(X,T_X(-\log D))&=(1,2,2,1).
\end{split}
\]
The base has dimension one, and the groups of the first displayed
row in degrees two and three vanish.

\item Let $Q=\bigoplus_{i=1}^3\eta_{i*}K_{\widetilde W_i}$, the normal
quotient of $T_X(-\log D)\to T_X$. Here
$\widetilde W_1,\widetilde W_2$ are ruled over elliptic curves and
$\widetilde W_3=dP_6$, so
\[
 h^\bullet(Q)=(0,2,3).
\]
The boundary $H^1(Q)\to H^2(T_X(-\log D))$ is a map between
two-dimensional spaces.
Under Serre duality its dual is the residue homomorphism
\begin{equation}\hypertarget{tag-01A6}{}\label{cdp:eq:E13-residue}
 H^1\bigl(X,\Omega_X^1(\log D)\otimes K_X\bigr)
 \longrightarrow
 H^1(\widetilde W_1,\mathcal O)\oplus
 H^1(\widetilde W_2,\mathcal O).
\end{equation}
The dual logarithmic sequence, tensored with $K_X$, has horizontal
term $\bar f^*\mathcal O(3)$ and vertical quotient
$\bar f^*(\mathcal O(-4)\oplus\mathcal O(-3))$.
The latter has zero cohomology in degrees zero and one: the coarse
direct images of $\mathcal O(-4)$ and $\mathcal O(-3)$ are both
$\mathcal O_{\mathbb P^1}(-1)$, and every line in their tensor
products with $G$ has negative weighted degree, so has no section.
Hence the
source of \eqref{cdp:eq:E13-residue} is canonically
\[
 H^1(X,\bar f^*\mathcal O(3))
 =H^0\bigl(\mathfrak B,\mathcal O(3)\otimes G\bigr)
 =H^0(\mathcal O(3))\oplus H^0(\mathcal O(2))=\mathbb C^2.
\]
On the affine $t$-line the generators of these last two spaces are
constant in the frame just constructed. The residue in each
summand is a nonzero scalar times evaluation followed by
$G_{p_i}\to H^1(\widetilde W_i,\mathcal O)$.

\item Write $e_i$ for the double root in the nodal cubic at $p_i$.
The line killed by this last projection is the graph-loop class.
Under $G\simeq E$ it is the toric Lie line: in the local periods
$(q\alpha,\beta),(\gamma,\lambda\delta)$ it is the direction
of the degenerating first multiplicative coordinate, whereas the
second coordinate is the elliptic normalization. If $V$ is dual to
$\omega$, its horizontal lift $\widetilde V$ has vertical
line-coordinate component $-\alpha(V)$. At the nodal boundary put
\[
 a_i=\frac{b}{e_i+\mu}.
\]
Thus the toric line is generated by
$\widetilde V+a_i v\partial_v$, and the projection to the elliptic
Lie line sends $(v\partial_v,\widetilde V)$ to $(1,-a_i)$, up to
a common nonzero factor at that point. Consequently the matrix of
\eqref{cdp:eq:E13-residue}, after nonzero row rescalings, is
\[
 \begin{pmatrix}1&-a_1\\1&-a_2\end{pmatrix},\qquad
 a_1-a_2=\frac{b(e_2-e_1)}{(e_1+\mu)(e_2+\mu)}\ne0.
\]
The inequality is exact: $b\ne0$; distinct nodal values have
distinct double roots, since $p_i=-3e_i^2$; and $e_i=-\mu$ would
place $P_\mu$ at the node and force $b=0$. The boundary therefore
has rank two.

\item Near the rank-two Mumford fibre put $q=t-p_3$. In the two
multiplicative coordinates its nonconstant periods are
$(qA_1(q),B_1(q))$ and $(A_2(q),\lambda qB_2(q))$, where all four
functions are holomorphic units; see
\cite[equation~(3.1) and Proposition~3.1, p.~6]{Viaclovsky2026Two}.
Normalize the two circle periods to the identity matrix. The remaining
period matrix is $(2\pi i)^{-1}\log q\,I_2+H(q)$ with $H$ holomorphic,
and its derivative is $(2\pi iq)^{-1}I_2+H'(q)$, of rank two for
nonzero $q$ sufficiently close to zero. In these period frames this
derivative is the cotangent Kodaira--Spencer connecting map.
The vertical part of a holomorphic one-form is therefore zero on
this open set and hence throughout the regular locus.
At a reduced singular fibre, a smooth point of one component away
from the conductor is a submersion point, so the coefficient of
the resulting base one-form extends holomorphically. At infinity
the same argument applies on the smooth root cover, and invariance
under the free affine deck action gives descent. Thus
$\bar f_*\Omega_X^1=\Omega_{\mathfrak B}^1$.
It follows that
$h^3(T_X)=h^0(\Omega_X^1\otimes K_X)
=h^0(\mathcal O(-9))=0$. The long exact sequence with quotient
$Q$ now gives $h^0(T_X)=1$, $h^1(T_X)=2$, and $h^2(T_X)=2$.

\item Finally the constructed bidisc has injective Kodaira--Spencer map.
At $\mu=0$, writing $a_i=p_i(0)$, differentiation of
$4t^3+27(\mu t+1)^2=0$ gives $p_i'(0)=-9/(2a_i)$.
An infinitesimal automorphism fixing the multiple value at infinity
has the form $At+B$; the three equalities
$-9/(2a_i)=Aa_i+B$ would make a nonzero quadratic polynomial vanish
at all three $a_i$. The cross-ratio variation is therefore nonzero,
and remains so on a smaller disc. For fixed $\mu$, a normalized
double curve is $\mathbb C^*/\langle c(\mu)\lambda\rangle$.
Its elliptic modulus has nonzero $\lambda$-derivative when
$|c(\mu)\lambda|$ is sufficiently small, as follows from
$j(q)=q^{-1}+744+O(q)$. A trivial first-order deformation of $X$
would trivialize its intrinsic anticanonical pencil, its labelled
exceptional-value configuration, and this double curve; it can
therefore kill neither of the two independent tangent directions.
Thus the bidisc maps to the two-dimensional tangent space of the
Kuranishi germ with invertible differential. Present its local
ring as $\mathbb C\{z_1,z_2\}/I$; the induced substitution into
$\mathbb C\{\mu,\lambda-\lambda_0\}$ is a local analytic
isomorphism before quotienting, by the inverse function theorem.
Every element of $I$ vanishes under it, so $I=0$. The classifying
map is a local analytic isomorphism. \ref{cdp:st:E1} and \ref{cdp:st:E6} give the algebraic-dimension assertion.

\end{enumerate}
\end{proof}

\begin{proposition}[The Liu--Rollenske family and its transverse deformation]\hypertarget{tag-01GL}{}\label{cdp:st:E14}
\StatementRecord{cdp:st:E14}{proved}{01GL}{\SeriesRef{cdp:eq:E14-elliptic-family}, \SeriesRef{cdp:st:E13}, \SeriesRef{cdp:eq:E14-canonical}, \SeriesRef{cdp:eq:E14-jet-map}, \SeriesRef{cdp:eq:E14-log-coho}, \SeriesRef{cdp:st:E6}, \SeriesRef{cdp:st:E1}; \cite[Proposition~2.1, pp.~3--4]{LiuRollenske2026}; \cite[Theorem~5.1, p.~21]{LiuRollenske2026}}

Every member $Y(u)$ of the Liu--Rollenske family, with its scale
sufficiently small and nonzero as in its construction, satisfies
\[
 h^\bullet(Y(u),T_{Y(u)})=(1,2,2,0).
\]
Its Kuranishi germ is a smooth reduced surface. One parameter is
the given scale; a second parameter replaces its cuspidal fibre by
two rank-one nodal fibres. All sufficiently small deformations
admit a torus fibration and have algebraic dimension one.
\end{proposition}


\begin{proof}

\begin{enumerate}[label=\arabic*.,leftmargin=*]
\item In the Weierstrass equation of \cite[Proposition~2.1, pp.~3--4]{LiuRollenske2026}, make the explicit
coordinate comparison
\[
 t_{\rm new}=\frac{27}{4}(t_{\rm old}^{-1}-1),\quad
 x_{\rm new}=x_{\rm old}/t_{\rm old}^2,\quad
 y_{\rm new}=y_{\rm old}/t_{\rm old}^3.
\]
Put $k=27/4$. Near $t_{\rm old}=0$, the new infinity coordinates are
$x_{\rm old}/[k(1-t_{\rm old})]^2$ and
$y_{\rm old}/[k(1-t_{\rm old})]^3$, and its base coordinate is
$t_{\rm old}/[k(1-t_{\rm old})]$. Near $t_{\rm old}=\infty$, put
$v=t_{\rm old}^{-1}$; the new fibre coordinates are
$v^2x_{\rm old},v^3y_{\rm old}$ and the new base coordinate is $k(v-1)$.
These charts extend the displayed maps to an isomorphism of the
Weierstrass models and their minimal resolutions, preserving $O$ and $P$.
They identify $\mathcal O(P-O)$ and its lifted translation; using the
transported linearization preserves the scale $u$.
The new equation is $y^2=x^3+tx+t$, with section $P=(-1,i)$; the multiple
value is now infinity, the cusp is at zero and the remaining nodal
value is $-27/4$. Consider the explicit deformation
\begin{equation}\hypertarget{tag-01A8}{}\label{cdp:eq:E14-elliptic-family}
 y^2=x^3+tx+t+\epsilon,\qquad
 P_\epsilon=(-1,b),\quad b^2=\epsilon-1,
 \quad b(0)=i.
\end{equation}
Its finite discriminant is $4t^3+27(t+\epsilon)^2$, whose discriminant
as a cubic in $t$ is $314928\epsilon^3(1-\epsilon)$.
The nodal value near $-27/4$ remains simple and holomorphically
labelled. The other two nodes coalesce to the type-II cusp at
$\epsilon=0$. Infinity keeps type $III^*$; on its order-four
good-reduction cover the equation is
\[
 y'^2=x'^3+x'+r^2+\epsilon r^6,
 \qquad P'_{\epsilon}=(-r^2,br^3).
\]

\item On a disc containing the cusp
and its two nearby nodal values, the total elliptic surface in
\eqref{cdp:eq:E14-elliptic-family} is smooth over $\epsilon$:
$\partial(y^2-x^3-tx-t-\epsilon)/\partial t=-(x+1)$ is a unit
near the cusp. Translation by $-2P_\epsilon$ is an automorphism
there. Put $M_\epsilon=\mathcal O(P_\epsilon-O)$ and let $c(\epsilon)$
be the simple nodal value continuing $-27/4$. The identity
$\mathcal Hom(t_{2P}^*M,M)\simeq\mathcal O(F)$ in
\cite[Proposition~2.1, pp.~3--4]{LiuRollenske2026} persists in the smooth family
of elliptic surfaces: $H^1(\mathcal O)=H^2(\mathcal O)=0$ centrally
and nearby, so the exponential sequence identifies their Picard
groups with the fixed integral $H^2$ lattice. The two line bundles
have the same Chern class throughout the family. Their ratio is
trivial on every surface; Grauert's theorem makes this trivialization
holomorphic in $\epsilon$. The section of $\mathcal O(F_{c(\epsilon)})$
with divisor $F_{c(\epsilon)}$ therefore gives a holomorphic lifted
translation, agreeing with the original one at $\epsilon=0$.
For a uniformly small nonzero scale its contraction on the
complement of the zero section is uniform on the compact elliptic
fibres. The quotient is consequently free, properly discontinuous,
proper over the disc, and smooth over the parameter. This quotient
itself fills the cusp and the nearby rank-one nodal fibres. At the
separate node the ordinary rank-two Mumford charts vary
holomorphically. At infinity the fixed order-four torsion translation
and good-reduction comparison vary holomorphically in the displayed
equation. Gluing along their marked punctured collars therefore
constructs a proper holomorphic family over $(\epsilon,u)$ whose
central scale slice is $Y(u)$. Every member has the central
differentiable type by Ehresmann. The topology theorem
\cite[Theorem~5.1, p.~21]{LiuRollenske2026} thus
identifies these differentiable manifolds with $S^6$.

\item Let $B$ be the reduced multiple fibre, so $f^*(\infty)=4B$.
The line $\mathcal O_Y(B)$, its canonical section and the isomorphism
$\mathcal O_Y(4B)\simeq f^*\mathcal O_{\mathbb P^1}(\infty)$ define
the root factorization $\bar f:Y\to\mathfrak B=\mathbb P(1,4)$.
For the coarse projection $\rho:\mathfrak B\to\mathbb P^1$ one has
$f=\rho\bar f$ and
$\rho^*\mathcal O_{\mathbb P^1}(1)=\mathcal O_{\mathfrak B}(4)$.
On the central threefold the
canonical form on the contracting quotient is the elliptic-surface
two-form wedged with $dv/v$. At the cusp it extends without a zero
or pole: from $t=(y^2-x^3)/(1+x)$ one gets
\[
 dt\wedge\frac{dx}{2y}=-\frac{dx\wedge dy}{1+x}.
\]
In a hexagonal chart $xyz=t-c(\epsilon)$, the form
$d\log x\wedge d\log y\wedge dt$ is $dx\wedge dy\wedge dz$.
At infinity the substitution
$t=r^{-4}$, $x=r^{-2}x'$, $y=r^{-3}y'$ gives the factor $-4r^{-4}$
in the surface form. The good three-form is invariant under the
free order-four affine action and descends to a nonvanishing frame.
These are the local canonical maps used in \ref{cdp:st:E13}.
Since the elliptic surface has
$K_S=\pi^*\mathcal O_{\mathbb P^1}(-1)$, these comparisons give
\begin{equation}\hypertarget{tag-01A7}{}\label{cdp:eq:E14-canonical}
 K_Y=f^*\mathcal O_{\mathbb P^1}(-1)
     =\bar f^*\mathcal O_{\mathfrak B}(-4).
\end{equation}
The calculation is simultaneous in $\epsilon$. Adjunction on every
full fibre gives trivial dualizing sheaf. \cite[Theorem~2.2, pp.~4--5]{Viaclovsky2026Combs} gives $h^0=1$;
duality gives $h^2=1$; constancy of its Euler characteristic in the
proper flat family gives $h^1=2$. Thus, also on the two-dimensional
base $(t,\epsilon)$, Grauert's theorem gives locally free $R^qf_*\mathcal O$
with base change, for $q=0,1,2$.

\item The multiplicative connection
\[
 \alpha=\frac{y+b}{x+1}\frac{dx}{2y},\qquad
 \nabla(hs)=(dh+\alpha h)s
\]
has precisely the residues of the rational section $s$ with divisor
$P_\epsilon-O$. For a vector $V$ dual to $dx/(2y)$, the horizontal lift is
$\widetilde V=V-\alpha(V)v\partial_v$. The multiplicative
connection commutes with the lifted translation, by the normalized
group-law comparison in \ref{cdp:st:E13}. On the good-reduction
chart its relative form is
\[
 \alpha=\frac{y'+br^3}{x'+r^2}\frac{dx'}{2y'},\qquad
 \frac{dx}{2y}=r\,\frac{dx'}{2y'}.
\]
Together with the scalar vector $v\partial_v$, this transition gives
the logarithmic Lie bundle $E=\mathcal O\oplus\mathcal O(-1)$. At the cusp the vector dual to $dx/(2y)$ is the derivation
$2y\partial_x+3x^2\partial_y$ on the cuspidal fibre, and its
horizontal lift is well defined because
$(y+b)/(x+1)$ is regular near the singular locus.

\item Put $G=R^1\bar f_*\mathcal O$. The inverse pairing associated with the invariant alternating
class gives $G\simeq E$ over the two-dimensional $(t,\epsilon)$
base away from its single cuspidal point: at a smooth fibre it is the pairing
$H^1(F,\mathbb C)/F^1\to(F^1)^*$; at a rank-one nodal fibre its quotient
frame consists of the graph-loop class and the elliptic normalization
class; at the hexagonal fibre both classes are graph-loop classes;
and the good-reduction cover gives the frame at infinity.
These local pairing and normalization maps are the ones in
\ref{cdp:st:E13}, applied to the displayed connection and root transition. Both $G$ and $E$ are locally free across that point, by the
preceding cohomological-flatness and connection calculations.
The matrix entries of the isomorphism and its inverse extend
across that codimension-two set by Hartogs' theorem. Their products
are the identity on a dense open set, hence everywhere. Thus
\[
 G=\mathcal O\oplus\mathcal O(-1),\qquad
 R^2\bar f_*\mathcal O=\mathcal O(-1)
\]
on the central root curve as well; the second equality also follows
directly from relative duality and \eqref{cdp:eq:E14-canonical}.

\item Let $D=C+W$ be the sum of the cuspidal and hexagonal fibres.
For the local tangent calculation choose coordinates in which
$t=y^2-x^3$. They are obtained from the earlier formula by the
invertible change $\widetilde x=x/(1+x)^{1/3}$,
$\widetilde y=y/(1+x)^{1/2}$, with the branches equal to one at
zero, and then renaming the two new coordinates only for this local
logarithmic-module calculation. The logarithmic
tangent module is freely generated by
\[
 2y\partial_x+3x^2\partial_y,\quad
 2x\partial_x+3y\partial_y,\quad\partial_z,
\]
where $z$ is a coordinate along the elliptic singular curve. The
second generator maps to $6t\partial_t$;
the first and third generate the relative logarithmic kernel.
Consequently the global exact sequence is
\[
 0\longrightarrow\bar f^*E\longrightarrow T_Y(-\log D)
 \longrightarrow\bar f^*\mathcal O(-3)\longrightarrow0,
\]
because $T_{\mathfrak B}(-C-W)=\mathcal O(5-8)$.
The vertical cohomology is $(1,1,0,0)$. Every line in the
horizontal Leray terms is $\mathcal O(-3)$ or $\mathcal O(-4)$,
which has no cohomology on $\mathbb P(1,4)$. Hence
\begin{equation}\hypertarget{tag-01AA}{}\label{cdp:eq:E14-log-coho}
 h^\bullet(T_Y(-\log D))=(1,1,0,0).
\end{equation}

\item Return to the Weierstrass coordinates. The preceding coordinate
change preserves the ideal $(x^2,y)$ and induces the identity on
$\mathbb C[x]/(x^2)$, since $\widetilde x\equiv x\bmod x^2$.
The quotient at $W$ is $\eta_*K_{dP_6}$ with cohomology $(0,0,1)$.
At $C$, the map $v\mapsto v(t)|_C$ identifies the quotient with
the Jacobian ideal
\[
 J=(x^2,y)\subset\mathcal O_C.
\]
Indeed its kernel is precisely the logarithmic tangent module and
$N_{C/Y}=\mathcal O_C$. Put $Z=V(J)\subset C$ and consider
the restriction map
\begin{equation}\hypertarget{tag-01A9}{}\label{cdp:eq:E14-jet-map}
 H^1(C,\mathcal O_C)\longrightarrow H^1(Z,\mathcal O_Z).
\end{equation}

\item In the deformation \eqref{cdp:eq:E14-elliptic-family} the critical
base near the cusp is the smooth curve
\[
 \Gamma:\quad t=-3x^2,\qquad
 \epsilon=x^2(3+2x).
\]
Its map to the $\epsilon$-disc is finite flat of degree two; its
central fibre is $\operatorname{Spec}\mathbb C[x]/(x^2)$.
The lifted translation fixes its $x$-coordinate. This follows on
the punctured curve because translation fixes the unique node of
each nodal cubic, and then follows at zero by holomorphic
continuation. On the line coordinate it is
$(x,v)\mapsto(x,c(x)v)$ with $c$ a holomorphic nowhere-zero
contraction. Thus its quotient critical locus is exactly the
elliptic family
\[
 \mathcal Z=(\Gamma\times\mathbb C^*)/
       \langle(x,v)\mapsto(x,c(x)v)\rangle.
\]
In particular $H^0(Z,\mathcal O_Z)$ and
$H^1(Z,\mathcal O_Z)$ are both free rank-one modules over
$A=\mathbb C[x]/(x^2)$, and $H^2(Z,\mathcal O_Z)=0$.

\item Restriction in this proper family is a holomorphic bundle map
$G|_\Gamma\to R^1(\mathcal Z/\Gamma)_*\mathcal O_{\mathcal Z}$.
At a node with coordinate $x$, the toric Lie line is generated by
$\widetilde V+\frac{b}{1+x}v\partial_v$: the horizontal lift has
line-coordinate component $-\alpha(V)=-b/(1+x)$ at that node.
The quotient to the elliptic Lie line, and hence the normalization
restriction under the pairing above, therefore sends
$(v\partial_v,\widetilde V)$ to
$(1,-b/(1+x))$, up to a common nonzero factor. This is the local
map in \ref{cdp:st:E13}. Thus on $\Gamma\setminus\{0\}$ restriction,
after a nonzero scalar change of target frame, has the row
\[
 \left(1,-\frac{b}{1+x}\right).
\]
The normalization of $C$ is a $\mathbb P^1$-bundle over
the elliptic singular curve $E_0$, and its normalization quotient
is $\mathcal O_{E_0}$: the first jet at the cusp transforms with
derivative one under translation on the additive smooth locus.
The normalization exact sequence, together with
$h^1(\mathcal O_C)=2$, $h^2(\mathcal O_C)=1$ and
$h^\bullet(\mathcal O_{\widetilde C})=(1,1,0)$, implies
that $H^1(\mathcal O_C)\to H^1(\mathcal O_{\widetilde C})$
is onto. Restriction from the ruled normalization to its singular
section identifies the latter group with $H^1(\mathcal O_{E_0})$.
Thus restriction to $E_0$ is onto. The common scalar in the row
above cannot have a zero at zero, since both entries would then
vanish and this restriction would be zero. It is a unit.

\item The equality of these holomorphic maps on the punctured curve
extends over all of $\Gamma$, and in particular over its
nonreduced central fibre. Hence \eqref{cdp:eq:E14-jet-map} is,
in the specified bases, the complex-linear map
\[
 \mathbb C^2\longrightarrow A,\qquad
 (a,d)\longmapsto a+d(-i+ix).
\]
Its matrix in $1,x$ has determinant $i$, so it is an isomorphism.

\item The exact sequence $0\to J\to\mathcal O_C\to\mathcal O_Z\to0$
now gives
\[
 h^\bullet(C,J)=(0,1,1).
\]
Together with \eqref{cdp:eq:E14-log-coho} and the hexagonal quotient,
the tangent exact sequence gives
$h^\bullet(T_Y)=(1,2,2,0)$.

\item In the above local trivialization
of the total space the map to the base is
\[
 f_{\epsilon}=\frac{y^2-x^3-\epsilon}{1+x}.
\]
The derivative $\partial_\epsilon f$ has image $-1+x$ in
$\mathcal O_C/J=\mathbb C[x]/(x^2)$ along the singular curve.
If the ambient first-order deformation were trivial, every
first-order variation of its map would be a global section of
$f^*T_{\mathbb P^1}$, hence the pullback of a projective vector
field, because $f_*\mathcal O_Y=\mathcal O_{\mathbb P^1}$.
Its local image in this Jacobian quotient is constant. Changing
the local ambient trivialization adds $df(v)$ and hence adds
zero in that quotient. The nonzero coefficient of $x$ therefore
proves that the $\epsilon$-direction is nontrivial and cannot be
cancelled by the scale direction, which does not change this local
map. With $\epsilon=0$, a normalized singular elliptic curve is
$\mathbb C^*/\langle c_0u\rangle$; its modulus has nonzero
$u$-derivative for sufficiently small $u$. A trivial deformation
of the total space preserves its anticanonical pencil from
\eqref{cdp:eq:E14-canonical}, its three distinguished special values,
and that elliptic singular curve. Hence the scale direction is
nontrivial too. The two directions span $H^1(T_Y)$.
Present the Kuranishi local ring as
$\mathbb C\{z_1,z_2\}/I$. The classifying map of the constructed
bidisc, followed by the embedding in $H^1(T_Y)$, has invertible
differential and is a local analytic isomorphism by the inverse
function theorem. Every element of $I$ pulls back to zero, so $I=0$.
This is the ideal argument of \ref{cdp:st:E13}, and it proves that
the displayed family is locally miniversal over the full reduced germ.
The rank-two hypothesis in \ref{cdp:st:E6} holds for this
Liu--Rollenske construction by \ref{cdp:st:E7}. Thus all sufficiently
small deformations admit torus fibrations.
Their second Betti number remains zero, and \ref{cdp:st:E1}
identifies their meromorphic function field with the pullback of
$\mathcal M(\mathbb P^1)$, giving algebraic dimension one.

\end{enumerate}
\end{proof}

\begin{proposition}[Tangent cohomology and Kuranishi germs of positive combs]\hypertarget{tag-01GM}{}\label{cdp:st:E15}
\StatementRecord{cdp:st:E15}{proved}{01GM}{\SeriesRef{cdp:eq:E15-products}, \SeriesRef{cdp:eq:E15-positive}, \SeriesRef{cdp:st:E5}, \SeriesRef{cdp:st:E13}, \SeriesRef{cdp:st:E1}, \SeriesRef{cdp:st:E6}; \cite[Theorem~1.1, p.~2]{Viaclovsky2026Combs}}

For each positive-comb manifold $X_m$ constructed in
\cite[Theorem~1.1, p.~2]{Viaclovsky2026Combs}, in particular for every $m\geq12$ in
its stated construction, one has
\[
 h^\bullet(X_m,T_{X_m})=(2m+1,\,8m+1,\,6m+1,\,0).
\]
Its Kuranishi germ is smooth and reduced of dimension $8m+1$.
A locally miniversal family is parametrized by all sufficiently nearby
admissible period sections on the root curve, with the affine
attachment classes fixed. Its members have algebraic dimension one
and torus fibrations with the same three critical values,
multiplicities $3,4$, and hexagonal fibre type.
\end{proposition}


\begin{proof}

\begin{enumerate}[label=\arabic*.,leftmargin=*]
\item Let $\mathfrak B=\mathbb P(3,4)$, with coarse projection
$\rho:\mathfrak B\to B$. Write $p_1,p_2$ for the finite-isotropy
points and $p_0$ for the hexagonal value. On the root curve let
$\mathcal H$ be the extended cohomology bundle, $F$ the Hodge
subbundle, $G$ its quotient, and $V$ its dual:
\[
 0\longrightarrow F\longrightarrow\mathcal H
 \longrightarrow G\longrightarrow0,\qquad V=F^*.
\]
The finite-cover comparison identifies
$G=R^1\bar f_*\mathcal O_{X_m}$ by invariant pushforward of the
ordinary torus Hodge sequence. At the hexagon the two deck
characters of the infinite toric quilt map to the two generators
of $H^1(\mathcal O_W)$ in its normalization complex, establishing
the same comparison there. The bundle $\bar f^*V$ is the vertical
logarithmic tangent bundle: its translations are dual to the
holomorphic one-forms on smooth fibres, and at the hexagon the
dual logarithmic coordinate fields extend in every toric chart.
Cup product identifies $R^2\bar f_*\mathcal O$ with
$D:=\bigwedge^2G$. On finite covers this is the ordinary torus
exterior algebra; at the hexagon it is the exterior algebra of
the two cycles in the normalization complex.

\item We specify finite characters by the orders of their invariant
generators. At $p_1$ the nontrivial Hodge line has invariant
generator $w$ times an equivariant frame, while the nontrivial
line of each of $V$ and $G$ has generator $w^2$ times its frame.
At $p_2$ the corresponding orders are $3$ for $F$ and $1$ for
$V,G$. Each bundle also has one trivial-character line. These
characters persist because comb smoothing fixes the limiting
periods. The tautological quotient on each tooth is
$\mathcal O(1)^{\oplus2}$, so each tooth adds $2$ to the coarse
degree of $G$. Therefore
\[
 \deg\rho_*G=2m-1,\qquad D=\mathcal O_{\mathfrak B}(24m-1).
\]
The determinant of $\mathcal H$ has degree zero on the root
curve: its determinant local system is trivial, and its cusp
residue has trace zero. Thus
\[
 \det F=\mathcal O(1-24m),\quad
 \det V=\mathcal O(24m-1),\quad
 \deg\rho_*V=2m-1.
\]
Indeed the finite fractional correction for $V$ and $G$ is
$2/3+1/4=11/12$; subtracting it from $(24m-1)/12$ gives
$2m-1$. For $F$ the correction $1/3+3/4=13/12$ gives
$\deg\rho_*F=-2m-1$.

\item The coarse bundles $V_c=\rho_*V$ and $G_c=\rho_*G$ have
nonnegative splitting. To prove this, evaluate the unchanged
invariant real homology class $\delta$ on $F$. Its resulting
holomorphic section of $V$ never vanishes on a smooth fibre:
a real class annihilating a Hodge plane also annihilates its
conjugate. The same argument holds on the finite good-reduction
covers. At the hexagon, $\delta$ is a nonzero combination of
angular cycles and the logarithmic forms have identity angular
period matrix. Its evaluation is nonzero there too, and gives
a nowhere-zero section of $V_c$. The invariant real cohomology
class $\gamma$ likewise gives a nowhere-zero section of $G_c$;
at the cusp it belongs to $\ker N$, whose projection to the
Hodge quotient is an isomorphism in the transverse graph chart.
At each finite point these nonzero invariant central values give
nonzero values of the coarse sections. A split rank-two bundle
of positive degree on $\mathbb P^1$ with a negative summand
cannot have a nowhere-zero section: every section lies in its
other summand, which has positive degree and hence has zeros.
Since $2m-1>0$, both splittings are nonnegative. In particular
\begin{equation}\hypertarget{tag-01AB}{}\label{cdp:eq:E15-positive}
 h^0(V_c)=h^0(G_c)=2m+1,\qquad H^1(V_c)=H^1(G_c)=0.
\end{equation}

\item Set $A=\rho_*(V\otimes G)$ and $L=\mathcal O_{\mathfrak B}(-5)$.
Multiplication of invariant sections gives exact comparisons
\begin{equation}\hypertarget{tag-01AC}{}\label{cdp:eq:E15-products}
\begin{split}
 0&\longrightarrow V_c\otimes G_c\longrightarrow A
       \longrightarrow\mathbb C_{p_1}\longrightarrow0,\\
 0&\longrightarrow V_c(2m-1)\longrightarrow\rho_*(V\otimes D)
       \longrightarrow\mathbb C_{p_1}\longrightarrow0,\\
 0&\longrightarrow G_c(-1)\longrightarrow\rho_*(L\otimes G)
       \longrightarrow\mathbb C_{p_1}\longrightarrow0.
\end{split}
\end{equation}
Here is the local verification. At $p_1$, multiplying two
nontrivial invariant generators of $V,G$ gives $w^4$, whereas
their tensor character has generator $w$; multiplication is
therefore one factor $t=w^3$. At $p_2$ the product is $w^2$
and the tensor generator also has order two, giving no defect.
Products involving a trivial character have no defect. The
line $D$ has the nontrivial character of $G$ and coarse degree
$2m-1$, proving the second comparison. The invariant-generator
orders of $L$ are $1,1$, and $\rho_*L=\mathcal O_B(-1)$.
Its product with the nontrivial generator of $G$ has the same
single defect at $p_1$, proving the third comparison.

\item All summands on the left in \eqref{cdp:eq:E15-products} have zero
$H^1$, by \eqref{cdp:eq:E15-positive}; for the last sequence their
degrees are at least $-1$. Consequently
\[
\begin{array}{c|c|c}
\text{bundle on }B&h^0&h^1\\ \hline
 A&8m+1&0\\
 \rho_*(V\otimes D)&6m&0\\
 \rho_*(L\otimes G)&2m&0
\end{array}
\]
by Riemann--Roch: their degrees are $8m-3$, $6m-2$, $2m-2$
and their ranks $4,2,2$. Also
$\rho_*(L\otimes D)=\mathcal O_B(2m-1)$.
Finite-group invariant pushforward is exact, so these calculations
give the cohomology on the root curve.

\item The logarithmic tangent sequence is
\[
 0\longrightarrow\bar f^*V\longrightarrow T_{X_m}(-\log W)
 \longrightarrow\bar f^*L\longrightarrow0.
\]
Projection and Leray yield
\[
 h^\bullet(\bar f^*V)=(2m+1,8m+1,6m,0),\qquad
 h^\bullet(\bar f^*L)=(0,2m,2m,0).
\]
Its only possible connecting map between nonzero groups is
$H^0(L\otimes G)\to H^0(V\otimes D)$. This is the map of
sections of the fibrewise cup-product boundary in \ref{cdp:st:E5}. The cusp
has rank two, so its sheaf map is generically injective. A map
from a locally free sheaf on a smooth curve with zero generic
kernel is injective; hence this map on sections is injective.
It follows that
\[
 h^\bullet(T_{X_m}(-\log W))=(2m+1,8m+1,6m,0).
\]
The branch-normal quotient is $\eta_*K_{dP_6}$, since the toric
three-form trivializes the canonical bundle near the fibre.
Its cohomology is $(0,0,1)$, and its exact sequence proves the
asserted tangent-cohomology tuple.

\item We construct a family with this full tangent space. All nearby
root-period sections have normal bundle $V\otimes G$, tangent
space $H^0(A)$, and obstruction space $H^1(A)=0$. This section
germ can be realized using an ordinary coarse Grassmannian bundle
and finite analytic conditions. Split the equivariant bundle as
$F\oplus G$ on a finite-cover chart and write a nearby plane as
a graph matrix. Dividing the invariant-generator factors shows
that a holomorphic lift requires one entry to vanish at $p_1$
and two at $p_2$. At $p_1$ the entry is from the nontrivial Hodge
line to the trivial quotient line. At $p_2$ it is that entry
and the entry between the two nontrivial lines. These are linear
conditions on the values of the local graph matrix. Their normal
sheaf is $A$. The \v{C}ech gluing obstruction across a small Artin
extension lies in $H^1(A)=0$; each local graph chart and its
stated conditions is smooth. The finite-dimensional Douady
section germ with these conditions is therefore smooth of
dimension $8m+1$. 

\item These sections produce compact threefolds even though their
limiting periods are allowed to vary. At a finite point the
equivariant lift is a small deformation on the compact
good-reduction cover. Admissibility is open there. The affine
action has the same real formula $x\mapsto Tx+v/n$; it is
holomorphic for the new equivariant periods and stays free
because its real action is unchanged. At the hexagon the
transverse graph chart is open, and its periods are
\[
 Z(q)=\frac{\log q}{2\pi i}I+H(q)
\]
with $H$ holomorphic and uniformly bounded on a smaller parameter
polydisc. The translations with unit factors
$\exp(2\pi iH_{ij})$ preserve the same toric fan. On logarithmic
radii their action is
$r\mapsto r+(I+O(1/\log|q|))\lambda$; it is uniformly proper
and free for small $q$. A fixed finite collection of toric unit
polydiscs gives compact representatives. Thus the filling is
proper with local equation $q=x_1x_2x_3$, and its total space
is smooth over the parameter. Varying $H(0)$ varies gluing units
and preserves these estimates and local equations. On the
compact complement of the critical discs admissibility is open.
The annular gluings extend holomorphically with their fixed
integral classes: the lattice--Lie--torus exact sequence and
Cartan vanishing on an annular product extend the original
translations. This constructs a proper holomorphic family with
all $8m+1$ period directions and central member $X_m$.

\item Its Kodaira--Spencer map is injective. A first-order trivialization
would make its map variation a section of
$H^0(X_m,f^*T_B)=H^0(B,T_B)$. The three critical values remain
constant by the equations $w^3,w^4,x_1x_2x_3$, so that base
vector field vanishes at three points and is zero. The
trivialization would then be fibre preserving. On a marked
smooth torus the period derivative is the corresponding
constant Dolbeault tensor in
$\operatorname{Lie}(F)\otimes H^1(F,\mathcal O_F)$.
A nonzero constant tensor cannot be a Dolbeault coboundary,
since derivatives of periodic coefficients integrate to zero.
Thus a period section giving a trivial total-space deformation
vanishes on every regular fibre, and hence identically. Its
Kodaira--Spencer map is an isomorphism onto the already computed
$H^1(T_{X_m})$. The inverse-function argument of \ref{cdp:st:E13} then makes
the Kuranishi ideal zero, proving smoothness, reducedness and
miniversality. The topology theorem of the comb construction
identifies these manifolds with $S^6$; \ref{cdp:st:E1} and \ref{cdp:st:E6} give algebraic
dimension one throughout their small deformation germs.

\end{enumerate}
\end{proof}

\begin{corollary}[The inclusion of the comb germ with fixed period values]\hypertarget{tag-P264}{}\label{cor:II:comb-fixed-value-inclusion}
\StatementRecord{cor:II:comb-fixed-value-inclusion}{proved}{P264}{\SeriesRef{cdp:st:E15}; \cite[Lemma~8.2 and Proposition~8.3, pp.~27--28]{Viaclovsky2026Combs}}
For a positive comb of degree parameter $m$ in \cref{cdp:st:E15},
the germ of period sections whose values at the three critical
points are fixed embeds in the root-period germ of that proposition.
The differential is the inclusion in
\[
0\longrightarrow H^0\bigl(B,N(-p_0-p_1-p_2)\bigr)
\longrightarrow H^0(B,A)
\longrightarrow \mathbb C^4\oplus\mathbb C^3\oplus\mathbb C^2
\longrightarrow0,
\]
where $N=(\rho_*F)^*\otimes G_c$ and $A$ is the admissible
period-variation sheaf in the proof of \cref{cdp:st:E15}.
The respective dimensions are $8m-8$, $8m+1$ and $9$.
\end{corollary}

\begin{proof}
The fixed-value germ is constructed in
\cite[Lemma~8.2 and Proposition~8.3, pp.~27--28]{Viaclovsky2026Combs}.
The inclusion sends a marked period section to the same section
on the root curve. Its inverse on its image forgets the root
description and takes the coarse period plane. The local graph
conditions in \cref{cdp:st:E15} identify its differential with
the sheaf inclusion
\[
0\longrightarrow N(-p_0-p_1-p_2)\longrightarrow A
\longrightarrow
 \mathbb C_{p_0}^{4}\oplus\mathbb C_{p_1}^{3}
                     \oplus\mathbb C_{p_2}^{2}
\longrightarrow0.
\]
The first sheaf sets all four graph entries to zero at each
point. The sheaf $A$ imposes respectively zero, one and two
conditions, so the three quotient lengths are $4,3,2$.
The cited Lemma~8.2 gives
$H^1(B,N(-p_0-p_1-p_2))=0$ and the dimension formula
$4+4c_m-2h-4r$, where $c_m=\deg G_c$,
$h=\deg(\rho_*F)+\deg G_c$ and $r=3$.
By \cref{cdp:st:E15}, $c_m=2m-1$ and
$h=(-2m-1)+(2m-1)=-2$. Substitution gives $8m-8$.
The long cohomology sequence is therefore
the asserted exact sequence; its quotient has dimension nine,
and $8m-8+9=8m+1$.
\end{proof}

\section{Rank divisors and theta comparisons}\hypertarget{tag-P269}{}\label{sec:II:period-rank}

\begin{proposition}[Rank divisors from the root Kodaira--Spencer morphism]\hypertarget{tag-01GN}{}\label{cdp:st:E16}
\StatementRecord{cdp:st:E16}{proved}{01GN}{\SeriesRef{cdp:eq:E.16.3}, \SeriesRef{cdp:st:E8}, \SeriesRef{cdp:st:E13}, \SeriesRef{cdp:st:E14}, \SeriesRef{cdp:st:E15}, \SeriesRef{cdp:st:E7}}

We use the logarithmic Kodaira--Spencer morphism at the degenerate fibres and the ordinary morphism on the finite smooth torus covers. For the original affine-twist family, including every $|p|>1$, its only rank defect is rank one at the order-three point. The logarithmic cusp has rank two. For the $(3k,4l)$ logarithmic transforms, on their full root covers the ranks at the two multiple values are respectively
\begin{equation}
 \begin{cases}1&k=1,\\0&k>1,\end{cases}
 \qquad
 \begin{cases}2&l=1,\\0&l>1.\end{cases}\hypertarget{tag-01A2}{}\label{cdp:eq:E.16.1}
\end{equation}
Every other smooth point and the logarithmic cusp have rank two.

For the two-parameter family the logarithmic rank is one at each of its two rank-one nodal values, zero at the order-four root point, and two everywhere else, including the hexagonal fibre. Its determinant has simple coarse zeros at the two nodal values and order four in the root coordinate at the multiple value. For the Liu--Rollenske central family the logarithmic rank is two at both the cuspidal and hexagonal fibres and throughout the remaining smooth locus; at the order-four root point it is zero, and its determinant has order four there.

For a positive comb section $\sigma$ of degree parameter $m$, let $x,y$ be root coordinates of weights $3,4$, with $y=0$ at the order-three point. There is a nonzero homogeneous polynomial $P_\sigma(X,Y)$ of degree $4m$, uniquely determined up to an overall nonzero scalar by the marked period section, such that
\begin{equation}
 \det\kappa_\sigma=y^2P_\sigma(x^4,y^3).\hypertarget{tag-01A3}{}\label{cdp:eq:E.16.2}
\end{equation}
The rank-defect set is exactly its zero divisor. At any such point the rank is zero exactly when all four entries of the Kodaira--Spencer matrix vanish; otherwise it is one. Formula \eqref{cdp:eq:E.16.3} below computes those entries from $\sigma$. The polynomial is nonzero at the hexagonal cusp, so that cusp always has rank two.
\end{proposition}


\begin{proof}

\begin{enumerate}[label=\arabic*.,leftmargin=*]
\item Adding a linearization constant or changing only the admissible affine lifts does not change the derivative of the marked period matrix. \ref{cdp:st:E8} therefore proves the first assertion. On a root cover of a higher logarithmic transform the original local good-reduction coordinate is $s=u^k$ at the first multiple value and $s=u^l$ at the second. The lifted period map is its holomorphic pullback. The chain rule gives $k u^{k-1}\kappa(s)$ or $l u^{l-1}\kappa(s)$. \ref{cdp:st:E8} gives rank one at the original order-three centre and rank two at the original order-four centre. The covering derivative is nonzero off its centre, proving \eqref{cdp:eq:E.16.1} and the absence of further defects. The cusp coordinate is unchanged.

\item For the two-parameter family \ref{cdp:st:E13} proves, with its extension maps,
\[
 E=G=\OO\oplus\OO(-1),\qquad
 \kappa:\OO(-7)\longrightarrow E\otimes G
 \quad\text{on }\PP(1,4).
\]
Thus the four entries have weights $7,6,6,5$. Write $x$ for the weight-one coordinate vanishing at the root point and $y$ for the weight-four coordinate. Every monomial of these respective weights is divisible by $x^3,x^2,x^2,x$. Consequently the entire matrix vanishes at the root point, and its determinant is divisible by $x^4$. The determinant is a section of $\OO(12)$, whose coarse degree is three. At the three Mumford values, the local period matrix has logarithmic
part $(\log q)/(2\pi i)\operatorname{diag}(1,\epsilon)$, where
$\epsilon=0$ at the two rank-one values and $\epsilon=1$ at the
hexagonal value; see \cite[equation~(3.1), p.~6]{Viaclovsky2026Two}.
Its value under $q\partial_q$ is
$(2\pi i)^{-1}\operatorname{diag}(1,\epsilon)$ at the centre.
Thus at each of the two rank-one values the logarithmic matrix has rank one and its determinant vanishes at least once. These two coarse zeros, together with the root zero of order at least four, account for all of the degree of $\OO(12)$. The determinant is nonzero because the logarithmic rank at the hexagonal cusp is two. All the stated orders are therefore exact, with no remaining zeros.

\item For the Liu--Rollenske central member \ref{cdp:st:E14} proves the analogous map from $\OO(-3)$ on $\PP(1,4)$. Its entries have weights $3,2,2,1$; these spaces consist respectively of multiples of $x^3,x^2,x^2,x$. The determinant is a section of $\OO(4)$ and is a constant multiple of $x^4$. The constant is nonzero by the hexagonal local monodromy. Hence the only rank defect is the root point, where all four entries vanish.

\item For a positive comb, \ref{cdp:st:E15} gives
\[
 \det V=\det G=\OO(24m-1),\qquad
 \kappa_\sigma:\OO(-5)\longrightarrow V\otimes G.
\]
The original and comb period planes agree at each critical value.
A local bundle automorphism $\Psi(t)=I+O(t)$ carries one plane to
the other. On the finite covers $t=w^n$ its root lift is given by
\[
 \begin{gathered}
 D_3(w)=\operatorname{diag}(1,1,w,w^2),\qquad
 D_4(w)=\operatorname{diag}(1,1,w,w^3),\\
 \Psi^\sharp(w)=D_n(w)\Psi(w^n)D_n(w)^{-1}.
 \end{gathered}
\]
The poles in $D_n^{-1}$ have order at most $n-1$, so
$\Psi^\sharp(w)=I+O(w)$ is holomorphic and equivariant, with the
same central characters. These are the maps of
\cite[Proposition~6.2, pp.~22--23]{Viaclovsky2026Combs}.
The tautological sequence on the root curve has trivial determinant
for its middle cohomology bundle: the marked flat volume extends
through the finite covers and through the determinant-one logarithmic
frame at the cusp. Taking determinants therefore identifies
$\det V$ with $\det G$. Each comb tooth contributes coarse degree
two to $G$, by \cite[equation~(5.5), p.~16]{Viaclovsky2026Combs}.
The unchanged finite characters make the determinant change the
pullback of $\mathcal O_{\mathbb P^1}(2m)$, which is
$\mathcal O_{\mathfrak B}(24m)$ on $\mathfrak B=\mathbb P(3,4)$.
The determinant is a section of $\det V\otimes\det G\otimes\OO(10)=\OO(48m+8)$. On $\PP(3,4)$ every weighted monomial $x^ay^b$ of this degree satisfies $b\equiv2\pmod3$ and $a\equiv0\pmod4$. Factoring $y^2$ therefore gives exactly the polynomial in \eqref{cdp:eq:E.16.2}, of coarse degree $4m$.

\item On a marked local root disc let the two columns of a holomorphic $4\times2$ matrix $U$ frame the Hodge subbundle of the marked cohomology extension. Let $\nabla$ be its marked Gauss--Manin connection, and let $D$ be the local generator of $T_{\mathfrak B}(-p_0)$: $\partial_s$ away from the cusp and $q\partial_q$ at the cusp. Then
\begin{equation}
 \kappa_\sigma=(\nabla_DU)\bmod\operatorname{im}U,
 \qquad
 \det\kappa_\sigma\ \text{is represented by}\
 \det(U_1,U_2,\nabla_DU_1,\nabla_DU_2).\hypertarget{tag-01A4}{}\label{cdp:eq:E.16.3}
\end{equation}
The first expression uses the quotient frame of $G$ and the dual frame
of $V=F^*$; it specifies every matrix entry and therefore also the
rank-zero test. On a Grassmann chart with invertible upper block $U_0$
of $U$, the frame change $U\mapsto UU_0^{-1}$ gives
$\binom{I}{Z}$. Write $A_{ij}$ for the four blocks of the connection
matrix contracted with $D$. Projection to the quotient by
$(u,v)\mapsto v-Zu$ gives the matrix
\[
 DZ+A_{21}+A_{22}Z-ZA_{11}-ZA_{12}Z.
\]
Its determinant is the four-column determinant of the frame
$UU_0^{-1}$ in \eqref{cdp:eq:E.16.3}. For a Hodge-frame change $U'=UC$, one has
$\nabla_DU'=(\nabla_DU)C+U(DC)$; wedging the four columns
multiplies the coefficient by $(\det C)^2$. Pair the quotient
determinant frame with $U_1\wedge U_2$ to the marked ambient
volume. It changes by $(\det C)^{-1}$, as does the dual Hodge
determinant frame.
For a tangent-frame change $D'=aD$, the coefficient is multiplied
by $a^2$, while the dual differential frame is multiplied by
$a^{-1}$ and its square by $a^{-2}$. The coefficient and frame
therefore give the same section of
$\det V\otimes\det G\otimes T_{\mathfrak B}(-p_0)^{-2}$.
These transition maps determine $P_\sigma$ from the marked period data. In the cusp comparison of \ref{cdp:st:E7} and \ref{cdp:st:E15},
the marked frame satisfies
$\nabla e_i=f_i\,dq/(2\pi iq)$ and $\nabla f_i=0$,
and the Hodge plane is the graph $(I,H_\sigma(q))$;
see \cite[Proposition~4.1, p.~12]{Viaclovsky2026Combs}.
Projection along this graph sends the logarithmic derivative to
$(2\pi i)^{-1}I+qH_\sigma'(q)$. At the cusp its determinant is
$(2\pi i)^{-2}\ne0$.

\end{enumerate}
\end{proof}

\begin{proposition}[Theta determinants at an elliptic half-period]\hypertarget{tag-01GP}{}\label{cdp:st:E18}
\StatementRecord{cdp:st:E18}{proved}{01GP}{\SeriesRef{cdp:eq:E-theta-general}, \SeriesRef{cdp:st:E5}}

For the theta expression in \eqref{cdp:eq:E-theta-general}, the
constant elliptic section $\mu(\tau)=1/2$ admits a holomorphic
contracting linearization over a disc about $\tau=i$ for which
\[
 \mu'-\frac{i}{2\pi}\partial_\mu\log\vartheta_1(\mu,\tau)=0,
 \qquad
 \det\Omega'(i)=-\frac{\wp(1/2;i)}{4\pi^2}\ne0.
\]
The determinant for this theta expression is given by
\eqref{cdp:eq:E-theta-general}. On the period family satisfying
$\wp(\mu)=0$, its Weierstrass term vanishes.
\end{proposition}


\begin{proof}
\begin{enumerate}[label=\arabic*.,leftmargin=*]
\item 
We use $\vartheta_j(z,\tau)=\theta_j(\pi z\mid\tau)$ in the
notation of \cite[eqs.~20.2.1--2, 20.2.6 and 20.2.12]{DLMFTheta2026}.
Thus the nome is $q=e^{\pi i\tau}$, and the Weierstrass periods
are $2\omega_1=1$, $2\omega_3=\tau$. The products below are
\cite[eqs.~20.4.3--5 and 23.17.8]{DLMFTheta2026}, with this
change of theta argument.

\item Let $\Delta$ be a sufficiently small disc about $i$ in the upper
half-plane and choose the holomorphic logarithm of the nonvanishing
function $\vartheta_1(1/2,\tau)/\eta(\tau)$ which is real at $i$.
Put
\[
 \mu(\tau)=\frac12,\qquad
 v(\tau)=\frac{i}{\pi}
       \log\frac{\vartheta_1(1/2,\tau)}{\eta(\tau)}+i,
 \qquad \beta(\tau)=6v(\tau).
\]
At $\tau=i$, the product expansions give
\[
 \frac{\vartheta_1(1/2,i)}{\eta(i)}
 =2e^{-\pi/6}\prod_{n\geq1}(1+e^{-2\pi n})^2>1.
\]
Thus $\operatorname{Im}v(i)>1$, and after shrinking $\Delta$
one has $\operatorname{Im}v>0$. Define
\[
 \Lambda_\tau=\mathbb Z(1,0)+\mathbb Z(0,1)
 +\mathbb Z(\tau,1/2)+\mathbb Z(3,\beta(\tau)),
 \qquad T_\tau=\mathbb C^2/\Lambda_\tau.
\]
The real determinant of these four periods is nonzero, since
$\operatorname{Im}\tau>0$ and $\operatorname{Im}\beta>0$.
They form a locally uniformly full lattice, so the quotients form
a proper holomorphic family of compact complex two-tori over
$\Delta$. Each fibre has
$h^\bullet(\mathcal O)=(1,2,1)$ and
$h^\bullet(T)=(2,4,2)$, with zero higher groups, by its
translation-invariant Dolbeault complex.

\item This family also has the required elliptic line-bundle
description. Quotient first by $(1,0),(0,1),(\tau,1/2)$.
Writing $z$ for the first coordinate and $a=\exp(2\pi iz_2)$
for the second gives the complement of zero in the flat line
bundle on $\mathbb C/(\mathbb Z+\tau\mathbb Z)$ with multipliers
$1$ and $-1$. For $P=1/2$, the meromorphic section
\[
 r(z)=\frac{\vartheta_1(z-P,\tau)}{\vartheta_1(z,\tau)}
\]
has divisor $P-O$ and satisfies $r(z+1)=r(z)$ and
$r(z+\tau)=-r(z)$. These equations identify the displayed
quotient line bundle with $\mathcal O(P-O)$. The last lattice
generator acts by
$(z,a)\mapsto(z+3,\exp(2\pi i\beta)a)$.
Its elliptic translation is $-6P=0$ in the elliptic curve;
its scalar is a holomorphic contracting linearization. This
constructs the family from an elliptic family, a section and a
linearization explicitly.

\item We compare with the symmetric period matrix
$\Omega=\left(\begin{smallmatrix}\tau&\mu\\\mu&v\end{smallmatrix}\right)$.
Adjoining $(\mu,v)$ enlarges $\Lambda_\tau$ to the lattice
generated by $(1,0),(0,1),(\tau,\mu),(\mu,v)$; its index is six,
because $6(\mu,v)=(3,\beta)$ and no smaller multiple belongs
to the original real lattice. Identity on $\mathbb C^2$ induces
a finite holomorphic isogeny of degree six between these families.
Its differential is the identity on the Lie spaces and its
pullback on complex first cohomology is an isomorphism. The
Kodaira--Spencer tensors therefore have the same rank under
these explicitly given comparisons.

\item The identity $\vartheta_1(1/2+z,\tau)=\vartheta_2(z,\tau)$
and evenness of $\vartheta_2$ give
$\partial_\mu\log\vartheta_1(1/2,\tau)=0$. Since $\mu'=0$,
the square term in \eqref{cdp:eq:E-theta-general} is zero. The
determinant, by that heat-equation formula, is
\cite[eqs.~20.7.5, 20.13.1, 23.3.9 and 23.6.2]{DLMFTheta2026}
\[
 \det\Omega'(i)=-\frac{\wp(1/2;i)}{4\pi^2}
 =-\frac{\vartheta_3(0,i)^4+\vartheta_4(0,i)^4}{12}\ne0.
\]
The last inequality follows from the theta products at the real
nome $e^{-\pi}$, whose factors give strictly positive theta
constants. Thus the square term vanishes at $i$ and the determinant is
nonzero there. Adding a holomorphic function $a(\tau)$ to
the lower-right entry of $\Omega$ adds $a'(\tau)$ to
$\det\Omega'$, because its upper-left derivative is one.

\item The direct images and total-space cohomology
in this local example are also determined. Write
$g:\mathscr T\to\Delta$ for the displayed proper family.
Translation-invariant forms, base change, and triviality of
holomorphic vector bundles on a disc give
\[
 R^qg_*\mathcal O_{\mathscr T}
   =\mathcal O_\Delta^{\binom2q},\qquad
 R^qg_*T_{\mathscr T/\Delta}
   =\mathcal O_\Delta^{2\binom2q}\quad(0\leq q\leq2),
\]
with all higher images zero. Shrink once more so that the
nonzero determinant stays nonzero. In the direct-image sequence
of $0\to T_{\mathscr T/\Delta}\to T_{\mathscr T}
\to g^*T_\Delta\to0$, the first boundary sends $1$ to the
nowhere-zero tensor $\kappa$, and the next boundary is the
rank-two isomorphism of \ref{cdp:st:E5}. Thus
\[
 R^qg_*T_{\mathscr T}=
 \begin{cases}
  \mathcal O_\Delta^2&q=0,\\
  \mathcal O_\Delta^3&q=1,\\
  \mathcal O_\Delta&q=2,\\
  0&q>2.
 \end{cases}
\]
Since $\Delta$ is Stein, every total-space group
$H^q(\mathscr T,\mathcal F)$ for these three sheaves is exactly
the space of sections on $\Delta$ of the corresponding displayed
direct image. These groups are the indicated spaces of
holomorphic sections on the disc.

\end{enumerate}
\end{proof}

\section{Canonical pencils in smooth specializations}\hypertarget{tag-034G}{}\label{sec:II:push-stability}

\begin{proposition}[A fractional canonical pencil at a smooth limit]\hypertarget{tag-033D}{}\label{push:src:E.2.5}
\StatementRecord{push:src:E.2.5}{proved}{033D}{\SeriesRef{cdp:st:E2}}

Let $\pi:\mathcal X\to\Delta$ be a proper holomorphic submersion
whose fibres are diffeomorphic to $S^6$.  Suppose that for every
$t\ne0$ the fibre is one of the members of \SeriesRef{push:src:E.2.1}, with the same
coprime multiplicities $(3k,4l)$.  Set
\[
 A=12kl,\qquad D=kl-3k-4l\ne0.
\]
There is a unique holomorphic line bundle $\mathcal L$ on
$\mathcal X$, up to isomorphism, satisfying
\[
 \mathcal L^{D}=K_{\mathcal X/\Delta}^{A}.
\]
Its restriction to $X_t$ for $t\ne0$ is the coarse pencil
$f_t^*\mathcal O_{\mathbb P^1}(1)$.  The coherent sheaf
$\pi_*\mathcal L$ is locally free of rank two, and its central
fibre injects into $H^0(X_0,\mathcal L_0)$.

The central manifold has algebraic dimension one and a
holomorphic algebraic reduction $g_0:X_0\to\mathbb P^1$ with
connected fibres.  If $s_0,s_1$ are the two central sections
obtained from a local frame of $\pi_*\mathcal L$, there are a
nonconstant rational map $R:\mathbb P^1\to\mathbb P^1$ of
degree $e\ge1$ and an effective Cartier divisor $E_0$ such that
\[
 [s_0:s_1]=R\circ g_0,\qquad
 \mathcal L_0\simeq g_0^*\mathcal O(e)\otimes\mathcal O(E_0).
\]
Here $[s_0:s_1]$ denotes the holomorphic extension of the ratio;
the displayed divisor is its common fixed divisor.
\end{proposition}


\begin{proof}

\begin{enumerate}[label=\arabic*.,leftmargin=*]
\item By Ehresmann's theorem the smooth bundle over the contractible
disc is differentiably trivial.  Consequently
$H^1(\mathcal X,\mathbb Z)=H^2(\mathcal X,\mathbb Z)=0$.
The exponential sequence identifies
\[
 \operatorname{Pic}(\mathcal X)=H^1(\mathcal X,\mathcal O_{\mathcal X})
\]
as groups.  This is a complex vector space, so multiplication by
the nonzero integer $D$ is an automorphism.  This proves the
existence and uniqueness of the \emph{relative} root, including
when $D<0$.  The same argument on each fibre makes its Picard
group uniquely divisible.  \SeriesRef{push:src:E.2.1} gives
$K_{X_t}^{A}=(f_t^*\mathcal O(1))^{D}$, so uniqueness identifies
$\mathcal L_t\simeq f_t^*\mathcal O(1)$.

\item Proper direct-image coherence gives $\mathcal V=\pi_*\mathcal L$.
It is torsion-free: a local section killed by a nonzero function
of $t$ vanishes on the inverse image of an open set, hence
vanishes as a section of the invertible sheaf on the smooth
total space.  A coherent torsion-free sheaf on a smooth curve
is locally free.  Its rank is two because the noncentral
algebraic reductions have connected fibres.  Applying $\pi_*$
to
\[
 0\longrightarrow\mathcal L\xrightarrow{\,t\,}\mathcal L
 \longrightarrow i_*\mathcal L_0\longrightarrow0
\]
shows directly that $\mathcal V/t\mathcal V$ injects into
$H^0(X_0,\mathcal L_0)$.  Thus the two central sections are
linearly independent.

\item Apply \cref{cdp:st:E2} to $X_0$. Its hypotheses hold because
$b_2(X_0)=0$ and $\langle c_3(X_0),[X_0]\rangle=2$.
Thus every meromorphic function on $X_0$ extends to a holomorphic
map to $\mathbb P^1$.

\item In particular the ratio of $s_0,s_1$ gives a nonconstant
holomorphic map $h_0:X_0\to\mathbb P^1$.  Also $a(X_0)\le1$:
otherwise two algebraically independent meromorphic functions
would give a holomorphic map $X_0\to\mathbb P^1\times\mathbb P^1$
with two-dimensional image.  At a point of differential rank
two, after choosing finite affine target coordinates $u,v$
centred at its image, the global meromorphic function $u/v$
would have the local indeterminacy of the quotient of two
independent coordinates.  It cannot extend holomorphically
there.  This contradicts the preceding theorem.  Hence
$a(X_0)=1$.

\item Take the Stein factorization $h_0=R\circ g_0$.  Its intermediate
space is a smooth compact curve.  Connectedness of the fibres
injects its $H^1(\ ,\mathbb Z)$ into $H^1(X_0,\mathbb Z)=0$
by the degree-one Leray sequence, so it is $\mathbb P^1$.
Every meromorphic function on $X_0$ factors rationally through
$g_0$: the image of its holomorphic map together with $g_0$
is an analytic curve in $\mathbb P^1\times\mathbb P^1$;
normalizing this curve gives a finite factor of $g_0$, whose
degree is one because $g_0$ has connected fibres.  Thus $g_0$
is the algebraic reduction, and $R$ has some degree $e\ge1$.

\item Finally, on a target chart where the first homogeneous
coordinate of $h_0$ is nonzero, $s_1=h_0s_0$ in a local
trivialization, and the common ideal is generated by $s_0$.
The analogous assertion holds on the other chart.  These local
generators define a nonzero section of
$\mathcal L_0\otimes h_0^*\mathcal O(-1)$ and an effective
Cartier zero divisor $E_0$.  This gives the stated factorization
of $\mathcal L_0$.

\end{enumerate}
\end{proof}

\begin{proposition}[Positive canonical slope preserves the fibration]\hypertarget{tag-033E}{}\label{push:src:E.2.6}
\StatementRecord{push:src:E.2.6}{proved}{033E}{\SeriesRef{push:src:E.2.5}; \citet[proof of Proposition~3.1.2, arXiv v1, p.~13]{AbashevaDeev2024AB}}

Under the hypotheses of \ref{push:src:E.2.5}, suppose $D>0$.  Then $E_0=0$
and $e=1$.  After shrinking the disc and choosing a frame of
$\pi_*\mathcal L$, its evaluation defines a proper flat
holomorphic map
\[
 g:\mathcal X\longrightarrow\mathbb P^1\times\Delta
\]
which is the intrinsic algebraic reduction on every fibre.
Its smooth general central fibre is a complex two-torus.

In particular, two families of the type in \ref{push:src:E.2.5} with positive
canonical slope cannot have the same smooth central manifold
unless their multiplicity pairs agree.  A positive-slope and
a negative-slope family cannot have a common smooth central
manifold either.
\end{proposition}


\begin{proof}

\begin{enumerate}[label=\arabic*.,leftmargin=*]
\item Choose a very general smooth fibre $F$ of $g_0$, avoiding the
images of all vertical components of $E_0$ and the jumping
sets for the pluricanonical direct images below.  Its normal
bundle in $X_0$ is trivial.  Adjunction gives
$K_F=K_{X_0}|_F$, and $c_1(K_F)=0$ integrally because
$H^2(X_0,\mathbb Z)=0$.  In particular $F$ is minimal: a
$(-1)$-curve would have canonical intersection $-1$.
Restriction of the root identity in \ref{push:src:E.2.5} gives
\[
 K_F^{A}\simeq\mathcal O_F(D E_0|_F).
\]
The right hand side is effective since $D>0$, so
$\kappa(F)\ge0$.

\item If $\kappa(F)>0$, some $r>0$ has $h^0(F,K_F^r)\ge2$.
Our very-general choice ensures that this is the generic rank
of the coherent sheaf $g_{0*}K_{X_0}^r$.  A sufficiently
positive twist of this sheaf on $\mathbb P^1$ is globally
generated.  Two of its global sections can therefore be
chosen to be linearly independent on a general fibre.  By
projection formula they are sections of
$K_{X_0}^r\otimes g_0^*\mathcal O(N)$; their ratio is a
meromorphic function on $X_0$ which is nonconstant on $F$.
This contradicts the algebraic-reduction assertion of \ref{push:src:E.2.5}.
The choice of a very general fibre is legitimate: for each
$r$ the exceptional set for the generic cohomology rank is
a proper analytic subset of the base curve, and there are
only countably many $r$.

\item Thus $\kappa(F)=0$.  The classification of minimal compact
complex surfaces makes $K_F$ torsion; the relevant consequence
is also recorded in
\citet[proof of Proposition~3.1.2, arXiv v1, p.~13]{AbashevaDeev2024AB}.
An effective divisor whose line bundle is torsion is zero:
raising its canonical section to a torsion-killing power
gives a holomorphic function on a compact connected manifold,
and that function is a nonzero constant.  Therefore
$E_0|_F=0$.  Every horizontal component of $E_0$ would meet
such a general fibre in a nonempty effective divisor.
Consequently $E_0$ has only vertical components, supported
over finitely many points of $\mathbb P^1$.

\item Choose a linear combination $s=as_0+bs_1$ for which its
moving central divisor is the union of the $e$ distinct
smooth fibres over $R^{-1}([-b:a])$, all disjoint from the
support of $E_0$.  Also choose it to be a regular fibre on
one specified noncentral member.  Both requirements omit
only finitely many choices in the projective line of
combinations.  Extend it by the corresponding combination
of the frame of $\pi_*\mathcal L$.

\item Its total zero divisor $\mathcal D\subset\mathcal X$ is a
relative effective Cartier divisor: the section is not
identically zero on any fibre, by independence of the two
sections, including at zero.  In local coordinates, $t$ is
a non-zero-divisor modulo its defining equation.  Hence
$\mathcal D\to\Delta$ is proper and flat.  Its general
noncentral fibre is a smooth connected torus.  There are no
vertical irreducible components of $\mathcal D$, since such
a component would have to be the whole irreducible
threefold $X_t$.  It follows that $\mathcal D$ is reduced
and irreducible: its unique generic component has multiplicity
one, and a Cartier divisor in a smooth manifold has no
embedded components.  Normalize $\mathcal D$ and take the
Stein factorization over $\Delta$.  The finite factor has
degree one over the punctured generic locus; normality of
the disc makes it the disc itself.  Its central fibre, and
therefore the image $\mathcal D_0$, is connected.

\item But the support of
\[
 \mathcal D_0=E_0+g_0^*R^*([-b:a])
\]
is a disjoint union of the $e$ chosen smooth connected fibres
and, if nonempty, the support of $E_0$ over other base points.
Connectedness forces $E_0=0$ and $e=1$.

\item The central two sections now generate $\mathcal L_0$ at
every point.  The noncentral two sections already generate
their pencils, so the evaluation
$\pi^*\pi_*\mathcal L\to\mathcal L$ is surjective after
shrinking.  It defines the asserted map $g$.  It is proper
because $\pi$ is proper.  On each $X_t$ the nonconstant
map to the smooth curve $\mathbb P^1$ is flat: its local
structure sheaf is torsion-free over the local discrete
valuation ring of the curve.  The total source is flat
over $t$, so the local criterion for flatness applied to
the parameter $t-t_0$ gives flatness over $(z,t)$.
Equivalently, the local fibre equations form a regular
sequence in the smooth source.  At a regular
point of the central base curve, properness makes $g$ a
proper submersion on a neighbourhood of the corresponding
compact fibre.  Ehresmann's theorem identifies that central
surface differentiably with nearby complex two-tori.
The root identity restricts to $K_F^A\simeq\mathcal O_F$, since
$E_0=0$ and $\mathcal L_0=g_0^*\mathcal O(1)$. Hence $\kappa(F)=0$.
Together with $b_1(F)=4$, the classification in
\cite[Chapter~VI, Theorem~1.1 and Table~10, pp.~243--244]{BHPV2004E}
identifies $F$ as a complex two-torus.

\item Suppose now two positive-slope families have the same
central manifold $X_0$.  Its algebraic reduction is intrinsic
up to $\operatorname{PGL}_2$, so the preceding result makes
both central pencil lines the same non-torsion line
$L=g_0^*\mathcal O(1)$.  With indices $i=1,2$ the identities
$K_{X_0}^{A_i}=L^{D_i}$ give
$L^{D_1A_2-D_2A_1}=\mathcal O$.  The line $L$ is non-torsion
because a positive power has a nonzero section with zeros.
Thus $D_1/A_1=D_2/A_2$, and the exact recovery of the ordered
pair from this slope in \SeriesRef{push:src:E.2.3} gives the assertion.

\item Finally, a common smooth limit from opposite signs would
have a nonzero section of a positive power of $K_{X_0}$
and of a negative power, by upper semicontinuity (or by
the same injective direct-image argument).  Raising them
to a common power, their nonzero product is a constant;
both are nowhere zero.  Hence $K_{X_0}$ is torsion.
The exponential sequence and $H^1(X_0,\mathbb Z)=
H^2(X_0,\mathbb Z)=0$ make $\operatorname{Pic}(X_0)$
torsion-free, so $K_{X_0}$ would be trivial.  Its nonzero
holomorphic three-form $\Omega$ is closed and exact because
$H^3(X_0,\mathbb C)=0$.  Stokes' theorem would give
$\int_{X_0}\Omega\wedge\overline\Omega=0$, whereas its
appropriate constant multiple is a positive volume form.
This contradiction excludes the common limit.

\end{enumerate}
\end{proof}

\begin{proposition}[An extended fibration cannot reach the zero period]\hypertarget{tag-033F}{}\label{push:src:E.2.7}
\StatementRecord{push:src:E.2.7}{proved}{033F}{}

Under the hypotheses of \ref{push:src:E.2.5}, suppose in addition that the
intrinsic algebraic reductions extend to a holomorphic map
\[
 g:\mathcal X\longrightarrow\mathbb P^1\times\Delta
\]
whose restriction to $X_0$ is nonconstant.  This hypothesis
holds in \ref{push:src:E.2.6}.  Normalize the punctured family to fixed
discrete logarithmic data with multiplicities $(3k,4l)$, and
write its isomorphism parameter as
\[
 z=\exp\!\left(-\frac{\pi i c}{kl}\right).
\]
The resulting holomorphic function $z:\Delta^*\to\mathbb C^*$
extends across zero and satisfies $z(0)\ne0$.  Its central
value lies strictly inside the admissible period disc.

\end{proposition}


\begin{proof}

\begin{enumerate}[label=\arabic*.,leftmargin=*]
\item Local miniversality in \SeriesRef{push:src:E.2.1} supplies holomorphic local
classifying functions for the family on $\Delta^*$.
\SeriesRef{push:src:E.1.2}--\SeriesRef{push:src:E.1.3} normalize the discrete data by holomorphic
family comparisons.  For fixed data the comparison shifts
are precisely $c\mapsto c+2kl\mathbb Z$, so the displayed
exponential is a single-valued holomorphic function on
$\Delta^*$.  An automorphism with the same marked parameter
has identity linear action on the homology of an ordinary
torus: in the necessity argument of \SeriesRef{push:src:E.1.2} its parameters
are $\alpha=1$ and $h=0$.  Consequently the homology local
system of the punctured family is also the one obtained
from these local classifying maps; possible translation
automorphisms of the comparisons add no homology monodromy.

\item The admissible parameter domain is a punctured disc
$0<|z|<R_{k,l}$, with finite positive radius.  The removable
singularity theorem therefore extends $z$ holomorphically
to $\Delta$.  It remains to exclude a zero at the centre.

\item The three critical-value sections may coincide at zero. Since
$\pi:\mathcal X\to\Delta$ is a submersion, the critical
locus of $g$ is the analytic locus where its differential
has rank less than two.  Properness makes its image
\(\mathcal C\subset\mathbb P^1\times\Delta\) a closed
analytic set.  Its intersection with the central target
curve is the finite critical-value set of the nonconstant
map $g_0:X_0\to\mathbb P^1$.  In particular
$\mathcal C$ has no component equal to that whole curve.
After shrinking $\Delta$, its reduced one-dimensional
part is finite over $\Delta$ and over $\Delta^*$ consists
of the three critical-value sections.

\item Those sections are individually marked by the linear
peripheral monodromy orders $3$, $4$, and infinity.
The normalization of the closure of each section is
therefore finite of degree one over the disc.  It is
isomorphic to the disc, so the three values
$c_1(s),c_2(s),c_0(s)$ extend as holomorphic maps to
$\mathbb P^1$, with collisions allowed at $s=0$.
The unique projective transformation $A_s$ sending this
ordered triple to $(0,1,\infty)$ has meromorphic
coefficients at zero, by its explicit cross-ratio formula.
This is the critical-value argument of \SeriesRef{push:src:B.4.4} in a setting
where no vertical discriminant component is present.

\item Choose a point $v\in\mathbb P^1$ outside all central
critical values.  Properness and closedness of
$\mathcal C$ allow us to shrink the disc so that
$(v,s)\notin\mathcal C$ for every $s$.  Then
\[
 \mathcal T=g^{-1}(\{v\}\times\Delta)\longrightarrow\Delta
\]
is a proper smooth holomorphic family of surfaces.
Its noncentral fibres are ordinary complex two-tori, and
Ehresmann's theorem makes its integral first-homology
monodromy around $s=0$ the identity.

\item In the original marked base, this same family is obtained
by evaluating the period construction at
\[
 t(s)=A_s(v)\in\mathbb P^1\setminus\{0,1,\infty\}.
\]
The function $t(s)$ is meromorphic at zero.  Its loop on a
sufficiently small positive $s$-circle is therefore either
null-homotopic in the marked base or conjugate to a power
of one of its three peripheral loops.  Indeed, if its
limit is ordinary, its image lies in an ordinary small
disc; if its limit is a marked point, a local coordinate
has a zero or pole of finite order and its image lies
in the corresponding punctured disc.  This observation
also covers a constant $t(s)$.

\item Let $\Lambda$ have ordered basis
$(\widehat\gamma,\widehat u,\widehat w,\widehat\delta)$,
put $E=\widehat\delta\otimes\gamma$, and write
$S_a=I+aE$.  The two finite peripheral monodromies are
\[
 A_1=\begin{pmatrix}
 1&0&0&0\\6&0&1&0\\-6&-1&-1&0\\-2&1&0&1
 \end{pmatrix},\qquad
 A_2=\begin{pmatrix}
 1&0&0&0\\0&0&-1&0\\-6&1&0&0\\3&0&1&1
 \end{pmatrix}.
\]
They have orders three and four, respectively, even for
the larger multiplicities: the extra finite deck
transformations are translations and do not affect
$H_1$.  In the cusp marking the third monodromy is
\[
 A_0=I+N,\qquad
 N\widehat u=\widehat w,\quad
 N\widehat\gamma=-\widehat\delta,\quad
 N\widehat w=N\widehat\delta=0.
\]
All peripheral matrices preserve $\widehat\delta$ and
$\gamma$, so commute with every $S_a$.  The change to
the cusp marking preserves the same rank-one operator
$E$; conjugating a peripheral loop consequently does
not alter the cancellation test below.

\item Suppose that $z(s)$ has a zero of order $d>0$ at zero.
Continuation of a logarithm around the positive loop
changes $c$ by $-h$, where
\[
 h=2kl\,d>0.
\]
Since $\Pi_{c+h}=\Pi_cS_h$, a comparison of the identical
complex universal covers from parameter $c-h$ back to
$c$ acts on the marked lattice by $S_{-h}$.
The affine translation part of a closing comparison
acts trivially on $H_1$.  Thus the return map of the
family $\mathcal T$ has homology action, up to a common
conjugation,
\[
 S_{-h},\qquad S_{-h}A_1^n,\qquad
 S_{-h}A_2^n,\qquad\hbox{or}\qquad S_{-h}A_0^n
\]
for an integer $n$.  Reversing the convention for a
peripheral generator only replaces $n$ by $-n$.
The parameter factor and the base factor commute, so
the order in this display has no effect.

\item None of these matrices is the identity.  In the two
finite cases, taking the third or fourth power would
give $S_{-3h}=I$ or $S_{-4h}=I$, impossible because
$E\ne0$ and $h>0$.  In the cusp case, $N^2=EN=NE=0$
and hence
\[
 S_{-h}A_0^n=I+nN-hE.
\]
Its coefficient of $\widehat w$ on $\widehat u$ is
$n$, so identity would force $n=0$, and then $h=0$.
The ordinary case is $S_{-h}\ne I$.  This contradicts
the identity monodromy of the smooth family over
the full disc.  Therefore $z(0)\ne0$.

\item Finally, the extended holomorphic function satisfies
$|z(0)|\le R_{k,l}$.  Equality would make it a constant
of maximal modulus, contradicting $|z(s)|<R_{k,l}$
for $s\ne0$.  Thus $0<|z(0)|<R_{k,l}$, as claimed.

\end{enumerate}
\end{proof}

\begin{proposition}[Horizontal divisors force non-K\"ahler fibre normalizations]\hypertarget{tag-033G}{}\label{push:src:E.2.8}
\StatementRecord{push:src:E.2.8}{proved}{033G}{\SeriesRef{push:src:E.2.5}, \SeriesRef{push:src:E.2.6}; \citet[Lemma~4.2, p.~684]{CDP2020}; \citet[Lemma~3.3 and Corollary~3.4, author version~0.08, p.~6]{Fujino2021E}}

Let $X$ be a compact complex threefold diffeomorphic to $S^6$, and
let $f:X\to\mathbb P^1$ be a surjective holomorphic map with
connected fibres.  Suppose that a smooth general fibre $F$ has
$b_1(F,\mathbb Q)=1$.  Then every reduced fibre of $f$ is
irreducible, and the monodromy on $H^1(F,\mathbb Q)$ is trivial.

If $H\subset X$ is a horizontal prime divisor, the normalization
of every reduced fibre lies outside Fujiki's class $\mathcal C$.
Equivalently, no such normalization admits a K\"ahler resolution.
In particular no fibre has Moishezon normalization.

For every surjective holomorphic map from a complex $S^6$ to
$\mathbb P^1$ with connected fibres, each irreducible component
of a reducible fibre has normalization outside $\mathcal C$.

Consequently, in \ref{push:src:E.2.5} with $D<0$, a horizontal component of
$E_0$ is possible only if every reduced fibre of the central
algebraic reduction $g_0$ is irreducible and has normalization
outside $\mathcal C$.
\end{proposition}


\begin{proof}

\begin{enumerate}[label=\arabic*.,leftmargin=*]
\item Choose a nonempty finite set $\Sigma\subset\mathbb P^1$ containing
all critical values; a regular value can be added if necessary.
Write $s=|\Sigma|$, put $A=f^{-1}(\Sigma)_{\mathrm{red}}$,
and let $r$ be the number of irreducible components of $A$.
The compact analytic set $A$ is triangulable.  Each of its
irreducible two-dimensional components has its complex
fundamental class, and their intersections have real dimension
at most two.  Hence
\[
 H^4(A,\mathbb Q)\simeq\mathbb Q^r.
\]
Alexander duality in the oriented six-sphere gives
\[
 b_1(X\setminus A,\mathbb Q)=r.
\]
On the other hand, $f:X\setminus A\to U:=\mathbb P^1\setminus\Sigma$
is a proper smooth fibre bundle.  The base $U$ has the homotopy
type of a graph with $s-1$ loops.  The degree-one Leray--Serre
sequence therefore gives an exact sequence
\[
 0\longrightarrow H^1(U,\mathbb Q)
 \longrightarrow H^1(X\setminus A,\mathbb Q)
 \longrightarrow H^1(F,\mathbb Q)^{\pi_1(U)}
 \longrightarrow0.
\]
Thus
\[
 r=s-1+\dim H^1(F,\mathbb Q)^{\pi_1(U)}\le s.
\]
Every nonempty fibre over $\Sigma$ has at least one irreducible
component, so $r\ge s$.  Equality follows throughout.  Thus
every fibre over $\Sigma$ has precisely one reduced irreducible
component and the invariant subspace is all of the
one-dimensional $H^1(F,\mathbb Q)$.  Fibres outside $\Sigma$ are
smooth and connected, hence irreducible.  This gives the
component count of \citet[Lemma~4.2, p.~684]{CDP2020}
in the present rank-one case, together with the required
triviality of monodromy.

\item Let $S=f^{-1}(t)_{\mathrm{red}}$.  A horizontal prime $H$ maps
surjectively to the base by properness, so $H\cap S\ne\varnothing$.
Since $S$ is irreducible, $S\subset H$ would imply $S=H$,
contrary to horizontality.  The canonical section of
$\mathcal O_X(H)$ consequently restricts to a nonzero section
on the reduced irreducible space $S$, with a nonempty zero
set.  Its zero divisor is a nonzero effective Cartier divisor
on $S$.  Indeed the section cannot vanish identically on
any local irreducible branch: its pullback would then vanish
on an open subset of the irreducible normalization, hence
everywhere by analytic continuation.  In the reduced local
ring it consequently avoids every minimal prime and is a
non-zero-divisor, including at self-intersections.

\item Let $\nu:S^\nu\to S$ be the normalization.  If $S^\nu$ were
in $\mathcal C$, it would have a compact K\"ahler resolution
$\mu:\widehat S\to S^\nu$; this precise surface statement is
\citet[Lemma~3.3 and Corollary~3.4, author version~0.08, p.~6]{Fujino2021E}.
The pulled-back line bundle
\[
 L=(\nu\mu)^*(\mathcal O_X(H)|_S)
\]
has $c_1(L)=0$ integrally, since $H^2(X,\mathbb Z)=0$.
Its pulled-back section is nonzero and has zeros: normalization
is finite surjective, and modification is proper surjective.
Thus it defines a nonzero effective divisor $B$ on
$\widehat S$.  For a K\"ahler form $\omega$,
\[
 0=\langle c_1(L)\smile[\omega],[\widehat S]\rangle
   =\int_B\omega>0,
\]
a contradiction.  The argument applies to the reduced support
even when the original fibre is nonreduced.  Moishezon spaces
are in $\mathcal C$, so the stated consequence follows.

\item For the assertion about a reducible fibre, let $S_i$ be one
of its reduced irreducible components.  The fibre is connected,
so another component $S_j$ meets $S_i$; otherwise $S_i$ would
be an open-and-closed component of its support.  Distinctness
ensures $S_i\not\subset S_j$.  Apply exactly the preceding
canonical-section and positivity argument to
$\mathcal O_X(S_j)|_{S_i}$.  Its first Chern class is zero
because $H^2(X,\mathbb Z)=0$, whereas its effective Cartier
zero divisor is nonempty.  This proves the assertion for every component of a reducible fibre.

\item For the application, choose a very general smooth fibre $F$
of $g_0$ meeting a horizontal component of $E_0$ and avoiding
its vertical components.  As in \ref{push:src:E.2.6}, adjunction gives
$c_1(K_F)=0$ and $F$ is minimal.  The identity of \ref{push:src:E.2.5} now
restricts to
\[
 K_F^{-A}=\mathcal O_F((-D)E_0|_F),
 \qquad -D>0,
\]
with a nonzero effective divisor on the right.  If a positive
power $K_F^m$ had a nonzero section, multiplying a suitable
power of that section by the $m$th power of the displayed
anticanonical section would give a nonzero holomorphic
function with zeros on $F$.  Compactness forbids this.
Thus $\kappa(F)=-\infty$.  Moreover $F$ is not K\"ahler,
since its nonzero effective divisor $(-D)E_0|_F$ has zero
real Chern class.  The Enriques--Kodaira classification then
places $F$ in class VII and gives $b_1(F)=1$; see
\citet[Chapter~VI, Theorem~1.1 and Table~10, pp.~243--244]
{BHPV2004E}.  The preceding argument applies with $f=g_0$
and any horizontal prime component $H$ of $E_0$.

\end{enumerate}
\end{proof}

\begin{proposition}[No horizontal divisor when all fibres are irreducible]\hypertarget{tag-033H}{}\label{push:src:E.2.9}
\StatementRecord{push:src:E.2.9}{proved}{033H}{\SeriesRef{push:src:E.2.8}; \citet[\S3, Theorem~1, p.~81, and the complete-intersection
calculation on p.~91]{RR70push}; \cite[Chapter~III, Theorem~4.12(ii), pp.~134--135]{BS76push}; \cite[Chapter~III, Theorem~4.1, p.~124]{BS76push}}

Let $X$ be a compact complex threefold diffeomorphic to $S^6$,
and let $f:X\to\mathbb P^1$ be a surjective holomorphic map
with connected fibres.  If every reduced fibre is irreducible,
then $X$ has no horizontal prime divisor for $f$.
\end{proposition}


\begin{proof}

\begin{enumerate}[label=\arabic*.,leftmargin=*]
\item Suppose that $H$ is such a divisor.  It is Cartier because
$X$ is smooth.  Since $H$ is horizontal and every reduced
fibre $S$ is irreducible, $H$ meets $S$ and does not contain
it.  As in \ref{push:src:E.2.8}, restriction of its canonical section is a
non-zero-divisor on every local branch of $S$.

\item Apply \cref{cdp:st:E17} on the reduced irreducible fibre $S$,
with effective Cartier divisor $H|_S$ and a locally free sheaf $V$.
The divisor meets its sole irreducible component, so
\[
 H^0(S,V(-nH))=0\qquad(n\gg0).
\]
For a full fibre $mS$, the ideal filtration has quotients
\[
 V|_S\otimes\mathcal O_S(-jS),\qquad 0\le j<m.
\]
Apply the same result to these finitely many bundles and use the
filtration exact sequences. This gives the vanishing on $mS$.
Multiplication by the canonical section of $H$ is injective on
$mS$: this Cartier hypersurface is Cohen--Macaulay, its only
associated component is $S$, and $H$ contains no local branch of $S$.

\item Let $B$ be any holomorphic vector bundle on $X$.  The map
$f$ is flat, since its local rings are torsion-free over the
discrete valuation rings of the smooth base.  Every analytic
fibre $F_t$ is a Cartier, Cohen--Macaulay, Gorenstein surface,
possibly nonreduced.  Its dualizing line is
\[
 \omega_{F_t}=K_X|_{F_t},
\]
where the constant one-dimensional factor from the target
has been trivialized.  Serre duality on this compact
Gorenstein surface gives
\[
 H^2(F_t,B(nH)|_{F_t})^*
 \simeq H^0\bigl(F_t,(B^*\otimes K_X)(-nH)|_{F_t}\bigr).
\]
The compact-space duality being used includes singular
Cohen--Macaulay spaces; it is the absolute duality theorem
of \citet[\S3, Theorem~1, p.~81, and the complete-intersection
calculation on p.~91]{RR70push}.
The preceding filtration argument makes the right hand
side zero for all sufficiently large $n$, for each $t$.

\item For fixed $n$, the nonvanishing locus
\[
 Z_n=\{t\in\mathbb P^1:
        h^2(F_t,B(nH)|_{F_t})>0\}
\]
is closed by semicontinuity.  Multiplication on the dual
section spaces gives $Z_{n+1}\subset Z_n$.  Their intersection
is empty by fibrewise vanishing.  Compactness of $\mathbb P^1$
therefore supplies a single $N$ with $Z_N=\varnothing$.
Constant-zero cohomology and analytic cohomology-and-base-change
\cite[Chapter~III, Theorem~4.12(ii), pp.~134--135]{BS76push} give
\[
 R^2f_*(B(NH))=0.
\]
Higher direct images vanish in degrees greater than two by
the fibre dimension.  The coherent sheaf $R^1f_*(B(NH))$
has vanishing first cohomology after a sufficiently positive
twist on $\mathbb P^1$.  The projection formula and Leray
sequence consequently give, for some integer $a\gg0$,
\[
 H^2(X,B\otimes M_+)=H^3(X,B\otimes M_+)=0,
 \qquad
 M_+=\mathcal O_X(NH)\otimes f^*\mathcal O(a).
\]

\item Independently, the reduced-space vanishing argument applied
to the connected smooth manifold $X$ and its nonzero divisor
$H$ gives
\[
 H^0(X,B\otimes M_-)=0,
 \qquad M_-=\mathcal O_X(-bH),\quad b\gg0.
\]
We must arrange these different vanishings for one line
bundle.  The exponential sequence gives
\[
 \operatorname{Pic}(X)\simeq H^1(X,\mathcal O_X)
\]
because $H^1(X,\mathbb Z)=H^2(X,\mathbb Z)=0$.
Let $\ell_-,\ell_+$ be the corresponding additive classes.
The affine line
\[
 \ell(z)=(1-z)\ell_-+z\ell_+,\qquad z\in\mathbb C,
\]
is realized by a holomorphic family of line bundles on
$X\times\mathbb C$: choose additive holomorphic \v{C}ech
cocycles representing the two classes and exponentiate their
linear combination.

\item For the proper projection $X\times\mathbb C\to\mathbb C$,
the cohomology jumping loci
\[
 T_j=\{z:h^j(X,B\otimes M_z)>0\},\qquad j=0,2,3,
\]
are closed analytic subsets of $\mathbb C$ by the proper
analytic finite-complex theorem
\cite[Chapter~III, Theorem~4.1, p.~124]{BS76push}: the condition is
expressed by the rank minors in a local finite free complex.
The set $T_0$ omits $z=0$,
while $T_2$ and $T_3$ omit $z=1$.  Each is therefore proper,
and their finite union cannot fill the irreducible complex
line.  Choose $z$ outside that union.  Then
\[
 \chi(X,B\otimes M_z)=-h^1(X,B\otimes M_z)\le0.
\]
Every line bundle on $X$ is topologically trivial, so
Hirzebruch--Riemann--Roch makes this Euler characteristic
equal to $\chi(X,B)$.

\item Finally take $B=T_X$.  Since $H^2(X,\mathbb Z)=H^4(X,\mathbb Z)=0$,
the only contributing term in its Riemann--Roch formula is
\[
 \chi(X,T_X)=\frac12\langle c_3(X),[X]\rangle
            =\frac12\chi_{\mathrm{top}}(S^6)=1.
\]
This contradicts the preceding inequality.  Hence no
horizontal prime $H$ exists.

\end{enumerate}
\end{proof}

\begin{proposition}[Persistence for either canonical slope]\hypertarget{tag-032V}{}\label{push:src:E.2.10}
\StatementRecord{push:src:E.2.10}{proved}{032V}{\SeriesRef{push:src:E.2.5}, \SeriesRef{push:src:E.2.7}, \SeriesRef{push:src:E.2.6}, \SeriesRef{push:src:E.2.8}, \SeriesRef{push:src:E.2.9}}

Under the hypotheses of \ref{push:src:E.2.5}, for every nonzero
$D=kl-3k-4l$ one has $E_0=0$ and $e=1$.
After shrinking the disc, the evaluation of
$\pi_*\mathcal L$ gives a proper flat holomorphic map
\[
 g:\mathcal X\longrightarrow\mathbb P^1\times\Delta
\]
whose restriction to every fibre is its intrinsic algebraic
reduction.  A smooth general central fibre of $g$ is a
complex two-torus.  Its normalized punctured period has
an interior limit as in \ref{push:src:E.2.7}.

Two families satisfying \ref{push:src:E.2.5} cannot have the same smooth
central manifold unless their ordered multiplicity pairs
agree, regardless of the signs of their canonical slopes.
\end{proposition}


\begin{proof}

\begin{enumerate}[label=\arabic*.,leftmargin=*]
\item The case $D>0$ is \ref{push:src:E.2.6}.  If $D<0$ and $E_0$ had a
horizontal prime component, the final part of \ref{push:src:E.2.8} would
give $b_1(F)=1$ for a smooth general fibre of $g_0$ and
would make every reduced fibre of $g_0$ irreducible.
\ref{push:src:E.2.9} would then forbid precisely that horizontal component.
Thus for either sign $E_0$ is vertical.

\item We apply the connected-divisor-limit argument of
\ref{push:src:E.2.6} to this vertical fixed divisor.
Its generic choice of pencil section has central zero
support equal to the disjoint union of the vertical fixed
part and $e$ distinct regular fibres.  The total divisor
is reduced and irreducible with connected central support
after normalization and Stein factorization.  Consequently
$E_0=0$ and $e=1$.  Surjectivity of evaluation, flatness and
the torus assertion follow from the remainder of that same
proof.  The hypothesis of \ref{push:src:E.2.7} is now satisfied for either
sign, and gives the interior period limit.

\item Suppose two such one-sided families have the same central
manifold $X_0$.  Its algebraic reduction is intrinsic, so
both pencil lines are the same line
$L=g_0^*\mathcal O(1)$ up to isomorphism.  With indices
$i=1,2$, the two root identities give
\[
 K_{X_0}^{A_i}=L^{D_i},\qquad
 L^{D_1A_2-D_2A_1}=\mathcal O.
\]
The line $L$ is not torsion, because a positive power has
a nonzero section with zeros.  Hence $D_1/A_1=D_2/A_2$.
\SeriesRef{push:src:E.2.3} recovers the ordered coprime multiplicity pair from
this rational number.

\end{enumerate}
\end{proof}

\section{Extension of critical fibres}\hypertarget{tag-P273}{}\label{sec:II:critical-fibre-extension}

\begin{proposition}[Purity with a smooth torus Jacobian]\hypertarget{tag-032W}{}\label{push:src:E.2.11}
\StatementRecord{push:src:E.2.11}{proved}{032W}{}

Let $B\subset\C^2$ be a sufficiently small bidisc about $0$, put
$U=B\setminus\{0\}$, and let $g:X\to B$ be a proper flat
holomorphic map from a connected smooth complex $(n+2)$-fold,
with connected fibres of dimension $n$.  Let $G\to B$ be a
smooth proper holomorphic family of complex $n$-tori with a
holomorphic zero section and its group law.  Suppose that
$X_U\to U$ is a holomorphic $G_U$-torsor.  Then $g$ is a
submersion on all of $X$, and the given action extends uniquely
to make $X\to B$ a holomorphic $G$-torsor.

\end{proposition}


\begin{proof}

\begin{enumerate}[label=\arabic*.,leftmargin=*]
\item Shrink $B$ so that the holomorphic Lie bundle of $G$ has a
frame $e_1,\ldots,e_n$.  Its period local system has a basis
of holomorphic sections
\[
 \lambda_1(b),\ldots,\lambda_{2n}(b)\in\C^n,
 \qquad
 G_b=\C^n/\sum_{a=1}^{2n}\Z\lambda_a(b).
\]
Indeed the lattice is locally constant over the simply
connected bidisc, and the local inverses of the relative
exponential map make its sections holomorphic.  The periods
are a real basis at every $b$.

\item The infinitesimal torsor action gives holomorphic vector
fields $V_1,\ldots,V_n$ on $X_U$.  Flatness makes
$Z=g^{-1}(0)$ have complex codimension two in $X$.
Hartogs extension in local coordinates therefore extends
each coefficient of each $V_j$ uniquely across $Z$; uniqueness
glues these local extensions to holomorphic vector fields
on $X$.  This use of Hartogs applies to functions on the
complement of an analytic set of codimension at least two
in a complex manifold, hence also coefficientwise to any
locally free sheaf.  The identities
\[
 dg(V_j)=0,\qquad [V_i,V_j]=0
\]
hold on the dense open subset $X_U$ and consequently on $X$.

\item Every $V_j$, and every linear combination with coefficients
pulled back from $B$, is complete.  Its integral curve stays
in a compact fibre; a finite endpoint of a maximal solution
would have an accumulation point in that fibre, where the
holomorphic ordinary differential equation extends the
solution.  The flows thus define a holomorphic action of
$\C^n$ on $X$ over $B$.

\item For $a=(a_1,\ldots,a_{2n})\in\R^{2n}$ take the time-one
flow of the vertical holomorphic vector field whose Lie
coordinate is
\[
 \sum_{j=1}^{2n}a_j\lambda_j(g(x)).
\]
These maps depend smoothly on $a$ and add under composition:
all their coefficients are constant along every vertical
flow.  For $a\in\Z^{2n}$ the map is the identity on $X_U$,
hence everywhere by holomorphic uniqueness.  We obtain an
action of the fixed compact real torus
$T=\R^{2n}/\Z^{2n}$ by holomorphic automorphisms of $X$,
over $B$.  On the fibre over $b$, it is the holomorphic
action of $G_b$ obtained by factoring the $\C^n$ action by
its period lattice.  The action of $T$ is free on $X_U$.
The variable period identities also make the holomorphic
Lie-bundle action descend to a holomorphic action
$G\times_BX\to X$.  Its restriction therefore acts on
each scheme-theoretic fibre, including its nilpotent
structure.

\item We show that its stabilizers have dimension zero everywhere.
Let $x\in X$, $b=g(x)$, and let $H$ be the identity component
of the stabilizer.  Viewed in $G_b$, it is a compact complex
subgroup: it is the identity component of the inverse image
of $x$ under the holomorphic orbit map $G_b\to X_b$.
Put
\[
 K_x=\ker(dg_x)=T_x(X_b),
\]
where the right-hand side is the Zariski tangent space of
the scheme-theoretic fibre.  Differentiating the holomorphic
$G_b$ action on this fibre gives a holomorphic representation
$H\to\operatorname{GL}(K_x)$.  It is trivial, because every
holomorphic function on the compact connected complex group
$H$ is constant.  The induced representation on
$T_xX/K_x$ is also trivial: this quotient injects into
$T_bB$ by $dg_x$, and every group element acts over $B$.
Averaging a Hermitian inner product over the compact group
$H$ gives an invariant complement to $K_x$.  The full
derivative representation on $T_xX$ is therefore trivial.

\item Here the required local linearization has a direct proof.
Choose holomorphic coordinates $z$ centred at $x$ and a
sufficiently small $H$-invariant neighbourhood on which all
$z(hy)$ are defined.  Such a neighbourhood exists by
compactness of $H$ and continuity of the action.  With
normalized Haar measure set
\[
 \Psi(y)=\int_H z(hy)\,dh.
\]
This map is holomorphic, invariant under $H$, and has
derivative the identity at $x$.  It is a holomorphic
coordinate system after shrinking.  Its invariance then
forces every $h\in H$ to be the identity near $x$, and hence
on connected $X$.  This contradicts freeness on $X_U$ unless
$H$ is trivial.  Thus every stabilizer is finite, and the
$n$ vector fields $V_j$ are linearly independent everywhere.

\item Their commuting flows give local holomorphic coordinates
$(z,s)\in\C^n\times\C^2$ in which they span the $z$
directions: choose a transverse complex surface and apply
the holomorphic inverse function theorem to the flow map.
In these coordinates $g$ has the form
\[
 g(z,s)=h(s).
\]
Equidimensionality implies that the germ $h^{-1}(b)$ is
zero-dimensional.  Consequently the germ
$h:(\C^2,0)\to(B,b)$ is finite; this is the local finiteness
criterion obtained by Weierstrass division for an isolated
fibre.  If $dg_x$ failed to have rank two, the determinant
of $dh$ would vanish at $s=0$.  It is not identically zero,
since the map is a submersion away from $Z$.  Its zero set
would therefore contain a complex curve through $0$.
Finiteness sends that curve to a curve of critical values.
This contradicts smoothness over $B\setminus\{0\}$.
Hence $g$ is a submersion everywhere.

\item The $G_b$ orbits now have the full dimension of every smooth
connected fibre and are compact.  Each fibre is a single
orbit.  Its finite stabilizer acts trivially on its tangent
space, because $G_b$ is abelian; it also acts trivially on
the base tangent quotient.  Averaging as above, now over
that finite group, again makes its full derivative trivial
and then makes the action locally trivial.  Freeness on
$X_U$ forces the stabilizer to be trivial.  Thus the extended
action is free and transitive on every fibre.  Local
holomorphic sections of the submersion identify $X$ with
$G$ over small base opens, which proves the torsor assertion.
The extension is unique because it agrees with the original
one on the dense open subset $X_U$.

\end{enumerate}
\end{proof}

\begin{proposition}[No additional central critical values]\hypertarget{tag-032X}{}\label{push:src:E.2.12}
\StatementRecord{push:src:E.2.12}{proved}{032X}{}

In the smooth-specialization setting of \ref{push:src:E.2.10}, let
\[
 g:\mathcal X\longrightarrow\mathbb P^1\times\Delta
\]
be the extended intrinsic fibration, and let
$c_1(s),c_2(s),c_0(s)$ be the three labelled critical-value
sections of \ref{push:src:E.2.7}.  The critical values of $g_0$ are exactly
the distinct points among $c_1(0),c_2(0),c_0(0)$.
If the normalizing projective transformations $A_s$ of
\ref{push:src:E.2.7} have a degenerate limit, the central regular fibres
are isomorphic complex two-tori and $g_0$ has at most two
critical values.  If the constant image of the limiting
rank-one matrix is $0$ or $1$, this torus is the
corresponding finite-root abelian surface.  If that image
is an ordinary source point, it is the ordinary torus at
that point. The cusp
cannot be the constant image.
\end{proposition}


\begin{proof}

\begin{enumerate}[label=\arabic*.,leftmargin=*]
\item \ref{push:src:E.2.7} gives $0<|z(0)|<R_{k,l}$.  A local logarithm
therefore extends the normalized parameter $c(s)$
holomorphically to the disc.  The critical image is closed
and contains the three punctured sections, so it contains
their three limiting points.  Suppose that a different
central critical value $b_0=(v_0,0)$ exists.  Take a small
bidisc $B$ about $b_0$ which avoids all three section
closures and all other central critical values.  Put
$U=B\setminus\{b_0\}$.  The restricted family over $U$
is a proper smooth family of complex two-tori.  We verify
that its Jacobian extends to a smooth holomorphic torus
group family over $B$, so that \ref{push:src:E.2.11} applies.

\item The normalized source coordinate is $t=A_s(v)$.
Represent $A_s$ by a matrix with holomorphic entries,
after clearing poles and removing a common power of
$s$, so its central matrix is nonzero.  If that matrix
is invertible, $t$ extends holomorphically near $b_0$ and
its value there is ordinary: a marked value would put
$b_0$ on the corresponding section closure.  The
ordinary period model with parameter $c(s)$ then gives
the desired smooth Jacobian over $B$.

\item Otherwise the central matrix has rank one.  Denote its
constant image by $t_*$ and its exceptional target point
by $v_*$.  For each marked source value different from
$t_*$, its preimage tends to $v_*$.  At least two of the
three marked values are different from $t_*$, so $v_*$
belongs to at least two labelled section closures.
Consequently $v_0\ne v_*$.  After shrinking $B$, the
function $t(s,v)$ tends uniformly to $t_*$ there.
For an ordinary $t_*$ the ordinary period model again
extends directly.

\item Suppose instead that $t_*=0$ or $1$, and use a local
coordinate $\xi$ at that marked point.  The divisor of
$\xi(t(s,v))$ on $B$ is supported on $s=0$, since $B$
avoids the corresponding critical-value graph.
Consequently
\[
 \xi(t(s,v))=s^n u(s,v),\qquad n>0,
\]
where $u$ is a nowhere-zero holomorphic function.
Choose a nearby $v_1$ for which $(v_1,0)$ is regular.
The family over the full $s$-disc at $v_1$ is
smooth, so its integral degree-one monodromy is the
identity.  The parameter $c(s)$ has no winding.  By
the marked comparison of \ref{push:src:E.2.7}, this monodromy is
$A_1^n$ or $A_2^n$, respectively.  Their orders are
three and four, so $e\mid n$, where $e=3$ or $4$.
A holomorphic $e$th root of $u$ exists on the bidisc,
and
\[
 r(s,v)=s^{n/e}u(s,v)^{1/e}
\]
is a holomorphic root coordinate on $B$.
Pull back the smooth finite-root torus period model
by $(c(s),r(s,v))$.  This produces the required
smooth torus group family $G\to B$.  The larger
multiplicities $3k,4l$ do not change this Jacobian
argument: their additional deck transformations
are translations, and the linear Jacobian model
is the original order-three or order-four root model.

\item The value $t_*=\infty$ is impossible.  In the
cusp coordinate the same divisor argument gives
$q=s^n u$ with $n>0$.  The smooth family at $v_1$
would have monodromy $A_0^n$.  But
$A_0=I+N$, $N^2=0$, and
$N\widehat u=\widehat w\ne0$, so $A_0^n\ne I$.

\item Identify $G|_U$ with the Jacobian as follows.
The punctured bidisc $U$ is simply connected.  Its
integral first-homology local system therefore has
a global marking.  Over $s\ne0$ that marking agrees
with the reference period marking, after one constant
integral change of basis; the transition comparisons
of \SeriesRef{push:src:E.1.3} have identity linear homology action at
fixed marked parameter.  Both holomorphic period
subbundles extend on $U$, and they agree on its dense
open subset $s\ne0$.  They agree everywhere on $U$.
Equivalently, the two quotients of their holomorphic
Lie bundles by the same marked period lattice are
isomorphic torus group families on $U$.

\item The total space is smooth and $g$ is proper and flat.
The only possible nonsmooth fibre over $B$ lies above
$b_0$, and its codimension in the total space is two.
All hypotheses of \ref{push:src:E.2.11} now hold.  That proposition makes
$g$ smooth also above $b_0$, a contradiction.  The
central critical-value set is therefore precisely
the set of labelled limits.

\item For the final assertions, in the rank-one case every
ordinary central target point different from $v_*$
has the same limiting marked torus: it is obtained
from $(c(0),t_*)$, or from the finite-root value
$(c(0),r=0)$.  Equality of the extended period
subbundles proves this assertion.
All marked values except possibly $t_*$ have their
preimages tending to $v_*$.  Hence there are at most
two central critical values; if $t_*$ is ordinary,
there is exactly one.
At $t_*=0$ or $1$, let $B$ be the order-$e$ linear
automorphism of the root torus $F$, with $e=3$ or $4$.
Its complex eigenvalues are $1$ and a primitive
$e$th root of unity.  Define endomorphisms
\[
 N=\sum_{j=0}^{e-1}B^j,\qquad M=[e]-N.
\]
Their images $E_0=\operatorname{im}N$ and
$E_1=\operatorname{im}M$ are elliptic subtori with
complementary tangent lines.  The addition map
\[
 E_0\times E_1\longrightarrow F,\qquad (x,y)\longmapsto x+y,
\]
is surjective because its composite with
$(N,M):F\to E_0\times E_1$ is $[e]$.
Its kernel consists of $(z,-z)$ with
$z\in E_0\cap E_1$.  The identities $N^2=[e]N$
and $NM=0$ show that $[e]z=0$ there, so this
kernel is finite and killed by $e$.  Thus the
addition map is an isogeny from a product of
elliptic curves, and $F$ is an abelian surface.

\end{enumerate}
\end{proof}

\begin{proposition}[Persistence of an isolated finite multiple fibre]\hypertarget{tag-032Y}{}\label{push:src:E.2.13}
\StatementRecord{push:src:E.2.13}{proved}{032Y}{}

In the setting of \ref{push:src:E.2.10}, suppose the central value of
one finite labelled critical section is distinct from
the limits of the other two labelled sections.  Write
its multiplicity as $m$, with $m=3k$ or $4l$, and its
linear monodromy order as $e=3$ or $4$, respectively.
Near that section, normalize the base change taking
an $m$th root of its defining equation.

The resulting total space is smooth, and its map to
the root bidisc is a proper smooth family of two-tori.
The deck group acts freely on the total space.
Consequently the central reduced finite fibre
is smooth, has exactly multiplicity $m$, and belongs
to the same smooth family of finite free affine
quotients as the punctured reduced finite fibres.

This applies to both finite fibres when the three
labelled critical values have distinct limits.  In
the two-critical-value collapsed-target case, it
applies to the isolated finite labelled fibre.
\end{proposition}


\begin{proof}

\begin{enumerate}[label=\arabic*.,leftmargin=*]
\item Choose target coordinates $(v,s)$ in a bidisc about
the section, and let $\Gamma$ be its graph.  Take a
local defining function $\xi$ for $\Gamma$.
On the smooth fourfold the divisor identity is
\[
 g^*\Gamma=m\mathcal D,
\]
where $\mathcal D$ is the prime closure of the
punctured reduced multiple-fibre family.  Indeed
this equality holds away from $s=0$.  The only
possible extra divisorial component would be
$X_0$, but $g_0$ is nonconstant, so $X_0$ is not
mapped into the point $\Gamma(0)$.  The closure
$\mathcal D$ is Cartier because the total space
is smooth.  Thus locally, even if its central
restriction is initially singular or reducible,
\[
 g^*\xi=a^m u
\]
for a local equation $a$ of $\mathcal D$ and a
holomorphic unit $u$.

\item Let $B'\to B$ be the smooth root bidisc
$r^m=\xi$, and let $Y$ be the normalization
of $\mathcal X_B\times_B B'$.  A local $m$th
root of $u$ factors its defining equation as
\[
 r^m-a^m u=
 \prod_{\zeta^m=1}(r-\zeta a u^{1/m}).
\]
For each $\zeta^m=1$, the map from the original
smooth local fourfold to the corresponding normalized
sheet is
\[
 x\longmapsto\bigl(x,\zeta a(x)u(x)^{1/m}\bigr).
\]
Projection to $x$ is its holomorphic inverse.
Thus the normalization is the disjoint union of
these smooth local sheets.  Therefore
$Y\to\mathcal X_B$ is a finite \'etale cover,
and $Y$ is smooth.  The deck group permutes these
branches freely, including over $a=0$.

\item The induced map $g':Y\to B'$ is proper and
equidimensional of relative dimension two.
Since both source and base are smooth, the
two pulled-back base parameters form a regular
sequence in each local ring, giving flatness.
Over $s\ne0$ this is the smooth root
torus family.  Its connectedness makes $Y$
connected and its generic fibre connected.
The Stein factor of $g'$ is finite birational
over the normal bidisc and hence is the
bidisc itself, so all fibres are connected.

\item By \ref{push:src:E.2.12} the original central map has no
critical value away from the labelled limits.
Our choice of the section therefore makes
$g'$ smooth away from the origin of $B'$.
It remains to supply a smooth Jacobian over
the full root bidisc.

\item If the normalizing projective transformations
$A_s$ have an invertible limit, the source
coordinate at the finite marked value has
the form
\[
 \xi_{\rm src}=\xi\,u_1(v,s)
\]
with $u_1$ a unit.  Thus its order-$e$ root
is $r^{m/e}u_1^{1/e}$, and the exact finite-root
period model with parameter $c(s)$ extends
smoothly to the full root bidisc.

\item If $A_s$ has rank-one limit, the isolated
section in the hypothesis must be the
preimage of its constant image $t_*$:
all other marked preimages tend to the
exceptional point.  Thus $t_*$ is precisely
the finite marked value under consideration.
Near the isolated section the source
coordinate has divisor equal to that section
plus a positive multiple of $s=0$, so
\[
 \xi_{\rm src}=s^n\xi\,u_1(v,s),\qquad n>0,
\]
with a nowhere-zero unit $u_1$.  Choose a
nearby central regular target point with
$\xi\ne0$.  Its smooth torus family
has identity parameter monodromy.  As in
\ref{push:src:E.2.12}, $c(s)$ has no winding and hence
$e\mid n$.  The required holomorphic root
coordinate on $B'$ is now
\[
 s^{n/e}r^{m/e}u_1^{1/e}.
\]
It again gives a smooth reference Jacobian
over $B'$.  On the punctured bidisc its
marking and holomorphic period subbundle
agree with the Jacobian by the same
simply-connected argument as \ref{push:src:E.2.12}.

\ref{push:src:E.2.11} makes $g'$ a smooth torsor over the
entire root bidisc.  In particular $r$ is
a submersion coordinate on the central
threefold upstairs.  The unramified deck
cover then shows that the original fibre
has exactly multiplicity $m$.  Restricting
to $r=0$ gives a proper smooth family of
two-tori in the parameter $s$, with free
affine deck action.  Its quotient is a
proper smooth family of the reduced finite
fibres, proving all assertions.

\end{enumerate}
\end{proof}

\begin{proposition}[Extension of the circle field and exclusion of finite-root collapse]\hypertarget{tag-032Z}{}\label{push:src:E.2.14}
\StatementRecord{push:src:E.2.14}{proved}{032Z}{}

In the setting of \ref{push:src:E.2.10}, the normalized vector field
$\delta$ of the punctured family extends to a holomorphic
relative vector field on $\mathcal X$ after shrinking
the parameter disc.  Its time-one flow is the identity,
and its real circle action on the central fibre is
equivariantly diffeomorphic to the prescribed circle
action on the punctured fibres.  In particular its
central fixed locus is a smooth rational curve with
normal bundle $\mathcal O(-1)\oplus\mathcal O(-1)$
and the same two opposite circle weights.

The rank-one limit of the normalizing projective
transformations $A_s$ cannot have constant image
$0$ or $1$.  More generally, a collapsed constant
image is impossible whenever the limiting torus's
complex $\delta$ line has two real-linearly
independent periods.  Thus the finite-root
isotrivial possibilities in \ref{push:src:E.2.12} are excluded.
\end{proposition}


\begin{proof}

\begin{enumerate}[label=\arabic*.,leftmargin=*]
\item Let $\pi:\mathcal X\to\Delta$ be the original
proper holomorphic submersion.  The coherent sheaf
\[
 \mathcal H=\pi_*T_{\mathcal X/\Delta}
\]
is torsion-free: multiplication by a nonzero base
function is injective on the vector bundle upstairs
and hence on its sections.  Its generic rank is
one by \SeriesRef{push:src:E.2.1}.  On the smooth disc it is therefore
locally free of rank one.  Choose a generator $V$
after shrinking the disc.  On $\Delta^*$ it has
the form
\[
 V|_{X_s}=b(s)\delta_s,
\]
with $b$ holomorphic and nowhere zero there.
The normalized fields $\delta_s$ agree on overlaps
of local classifying maps: at fixed discrete data
the comparison maps of \SeriesRef{push:src:E.1.3} have fibre-linear
part the identity and fix the marked period
$\widehat\delta$.  Their translations also
preserve the invariant vector field.
Moreover $V$ is vertical for $g$, because this
holds on the dense punctured family.

\item Choose a regular central target value and a small
bidisc around it on which $g$ is smooth.  The
ordinary or finite-root comparison used in
\ref{push:src:E.2.12} supplies its smooth marked Jacobian, with
the marked period $\widehat\delta$ extending as
a nonzero holomorphic Lie section.  The resulting
invariant field $\delta_B$ on the torsor
is nonvanishing there.  Comparing one nonzero
coefficient of $\delta_B$ with $V$ shows that
$b(s)$ extends holomorphically to $s=0$.
Its independence of the target coordinate follows
from that equality on $s\ne0$ and holomorphic
uniqueness.

\item If $b(0)=0$, the restriction $V|_{X_0}$ vanishes
on a nonempty open subset of connected $X_0$,
and therefore vanishes everywhere.  In local
coordinates for the submersion $\pi$, every
coefficient of $V$ is then divisible by $s$.
The quotients glue uniquely to a global
holomorphic relative vector field $V/s$.
This contradicts the fact that $V$ generates
the free module $\mathcal H$.  Hence $b(0)\ne0$,
and $V/b(s)$ gives the required global extension
of $\delta$.

\item The extended field is complete on each compact
fibre of $\pi$.  Its time-one flow is the
identity on the punctured family and hence on
the whole family by holomorphic uniqueness.
Its real flow consequently defines an action
of $\mathbb R/\mathbb Z$ on $\mathcal X$
over the parameter disc.  Average a smooth
metric and an Ehresmann horizontal distribution
over this compact group.  Parallel transport
then gives equivariant diffeomorphisms of
the fibres.  The central action is therefore
the same effective semifree action, with a
single fixed two-sphere, as on the prescribed
fibres.

\item The fixed set of a holomorphic compact-group
action is a complex submanifold, as follows
from the averaging coordinates used in
\ref{push:src:E.2.11}.  Here it is a proper smooth family
of curves of genus zero, and hence each fibre
is $\mathbb P^1$.  Its two holomorphic normal
weight lines have the same degrees as in
the punctured family.  Those degrees are
$-1,-1$, as follows from the prescribed Mumford model
\cite[Proposition~2.8, pp.~6--7, and Remark~8.3, pp.~22--23, arXiv v1]{Engel2026}.
Its normalization is $P=dP_6$.  The two opposite
boundary curves $D_+,D_-$ fixed by $\delta$ each map
isomorphically to the fixed curve $C$.  At a triple
point, coordinates with $q=xyz$ have $\delta$-weights
$(1,0,-1)$, and $C$ is $x=z=0$.  The differentials
of the two surface-branch inclusions identify
$N_{D_+/P}$ and $N_{D_-/P}$ with the two normal
weight lines of $C$ in the smooth threefold.
The double-point charts $q=uv$ give the same maps
away from the triple points, so these are global
line-bundle isomorphisms over $D_\pm\simeq C$.
Both $D_\pm$ have self-intersection $-1$ on $P$.
The characters
distinguish the lines, so their central
restrictions are exactly
$\mathcal O(-1)\oplus\mathcal O(-1)$.

\item Suppose now that $A_s$ has rank-one limit
with constant image a finite marked value.
\ref{push:src:E.2.12} identifies every central regular
torus, with its marked $\delta$ direction,
with the corresponding finite-root torus
at parameter $c(0)$.  The invariant lattice
of $A_j$ has rank two and spans its
one-dimensional complex fixed subspace.
It contains $\widehat\delta$; its second
primitive period is independent over
$\mathbb R$.  The complex fixed subspace
is precisely the $\delta$ line, since
the other complex eigenvalue is a
primitive third or fourth root of unity.
Consequently, in the normalization in
which the $\delta$ period is $1$, there
is a constant second period
\[
 \kappa\in\mathbb C\setminus\mathbb R.
\]
This is a period of the invariant field
on every central regular fibre.  The
time-$\kappa$ flow is therefore the
identity on their dense union, and
hence on all of $X_0$.

\item The flow of the nonzero field $\delta_0$
thus factors through a nontrivial
holomorphic action of the compact
elliptic curve
$E=\mathbb C/(\mathbb Z+\kappa\mathbb Z)$.
Such an action has no fixed point on
a connected complex manifold unless
it is trivial.  Indeed its derivative
at a fixed point is a holomorphic
representation of the compact complex
group into a general linear group,
so is constant.  Averaging local
coordinates over the compact group
then makes the action trivial near
that point, and holomorphic uniqueness
makes it trivial everywhere.

\item The elliptic action here is nontrivial
because $\delta_0\ne0$.  It therefore
has no fixed point.  This contradicts
the fixed rational curve of the
extended real circle action: a zero
of the generating field is fixed
by its entire complex flow.
The same argument applies to an ordinary
constant-image torus if its $\delta$
line has a second independent period.

\end{enumerate}
\end{proof}

\begin{proposition}[An isotrivial volume obstruction]\hypertarget{tag-0330}{}\label{push:src:E.2.15}\label{push:src:E.2.15:smooth}
\StatementRecord{push:src:E.2.15}{proved}{0330}{}

Let $f:X\to\PP^1$ be a proper holomorphic map with connected
fibres from a compact connected smooth complex threefold.  Suppose
that for some $p\in\PP^1$ the restriction over
$U=\PP^1\setminus\{p\}$ is a smooth family of mutually
isomorphic complex two-tori.  If
\begin{equation}
 K_X^{\otimes A}\simeq f^*\mathcal O_{\PP^1}(D),
 \qquad A\in\Z_{>0},\quad D\in\Z,
\label{eq:II:next-E-2-15-1}\end{equation}
then necessarily $D<-A$.  In particular the ordinary-source
collapsed-target case of \ref{push:src:E.2.12} cannot occur for the
six-sphere families under consideration.
\end{proposition}


\begin{proof}

\begin{enumerate}[label=\arabic*.,leftmargin=*]
\item Choose an affine coordinate $t$ on $U\simeq\C$ and put
$q=t^{-1}$ near $p$.  Since $U$ is simply connected, its
integral first-homology local system has a global marking.
The holomorphic period subbundle is constant: locally the
marked periods belong to the countable integral-change-of-basis
orbit of one fixed torus, and a continuous map from a connected
base into such a countable subset of a complex Grassmannian
is constant.  For the last assertion, any two distinct image
points can be separated by a real continuous coordinate;
its connected image would contain an interval and would be
uncountable.  Thus the Jacobian group family is the product
with one fixed torus $T$.

\item Choose a nonzero translation-invariant holomorphic two-form
$\eta_T$ on $T$.  Local holomorphic origins identify our
family with its Jacobian, and two choices differ by a
translation.  Hence $\eta_T$ defines a global nonvanishing
relative canonical form $\eta$ on $X_U$.  Its fibre volume
is a fixed positive number:
\begin{equation}
 \int_{X_t}\frac14\eta\wedge\overline\eta=V>0.
\label{eq:II:next-E-2-15-2}\end{equation}
The canonical determinant
isomorphism for a submersion gives the nonvanishing
absolute three-form $\eta\wedge dt$ on $X_U$; this notation
does not require choosing an absolute lift of $\eta$.

\item Use \eqref{eq:II:next-E-2-15-1} to pull back a meromorphic section of
$\mathcal O(D)$ whose divisor is $D[p]$.  We obtain a
meromorphic pluricanonical form $\sigma$ with
\[
 \operatorname{div}(\sigma)=D f^*(p).
\]
On $X_U$ it is holomorphic and nowhere zero.  Its quotient
by $(\eta\wedge dt)^{\otimes A}$ is a holomorphic function
constant on each compact connected fibre.  It therefore
descends to an entire nowhere-zero function $\phi$ on $\C$:
\begin{equation}
 \sigma=\phi(t)(\eta\wedge dt)^{\otimes A}.
\label{eq:II:next-E-2-15-3}\end{equation}

\item On a small disc about $p$, the meromorphic section
$q^{-D}\sigma$ is in fact a holomorphic nowhere-zero section
of $K_X^{\otimes A}$ over the whole inverse image.  This
follows directly from \eqref{eq:II:next-E-2-15-1}, or by cancellation of the
divisors, with their multiplicities.  Denote it by
$\Omega_A$.  On the punctured inverse image, since
$dt=-q^{-2}dq$, it is
\begin{equation}
 \Omega_A=(-1)^A P(q)(\eta\wedge dq)^{\otimes A},
 \qquad P(q)=q^{-D-2A}\phi(1/q).
\label{eq:II:next-E-2-15-4}\end{equation}
The function $P$ is single-valued and holomorphic on the
punctured disc.

\item A nowhere-zero holomorphic $A$-canonical form determines
a smooth positive density: if it is
$h(z)(dz_1\wedge dz_2\wedge dz_3)^{\otimes A}$, take
$|h|^{2/A}$ times the ordinary complex-coordinate volume
density.  This is independent of the coordinates.  The
density associated with $\Omega_A$ has finite integral
over the inverse image of a smaller closed disc, because
that inverse image is compact.  Fibre integration of
\eqref{eq:II:next-E-2-15-4}, using \eqref{eq:II:next-E-2-15-2}, consequently gives
\begin{equation}
 \int_{0<|q|<\epsilon}|P(q)|^{2/A}\,d\mathrm{Area}(q)
 <\infty.
\label{eq:II:next-E-2-15-5}\end{equation}
Only a fixed positive normalization constant is suppressed.

\item We spell out the isolated-singularity consequence of
\eqref{eq:II:next-E-2-15-5}, including when $2/A<1$.  The function
$|P|^{2/A}$ is subharmonic.  If $|q|=r$, the mean-value
inequality on the disc of radius $r/2$ centred at $q$ gives
\[
 |P(q)|^{2/A}
 \leq \frac{4}{\pi r^2}
       \int_{r/2<|z|<3r/2}|P(z)|^{2/A}\,d\mathrm{Area}(z).
\]
The last integral tends to zero as $r\to0$, by
integrability.  Hence $r^A|P(q)|\to0$, uniformly on
$|q|=r$.  The function $q^A P(q)$ therefore has a removable
singularity and vanishes at zero.  For some holomorphic
$H$ on the disc,
\begin{equation}
 P(q)=q^{1-A}H(q).
\label{eq:II:next-E-2-15-6}\end{equation}
Equivalently, the pole order of $P$ is at most $A-1$.

\item Combining \eqref{eq:II:next-E-2-15-4} and \eqref{eq:II:next-E-2-15-6} yields
\begin{equation}
 \phi(1/q)=q^{D+A+1}H(q).
\label{eq:II:next-E-2-15-7}\end{equation}
If $D\geq-A$, the right-hand side extends to zero at
$q=0$.  Thus $\phi$ extends holomorphically to $\PP^1$
and vanishes at $p$.  A holomorphic function on $\PP^1$
is constant, so $\phi=0$, contradicting its construction.
This proves $D<-A$.

\item In \ref{push:src:E.2.10} the canonical-root equality has
\[
 A=12kl,\qquad D=kl-3k-4l,
 \qquad
 \frac DA=\frac1{12}-\frac1{4l}-\frac1{3k}
       \geq-\frac12>-1.
\]
If the collapsed constant source point is ordinary,
\ref{push:src:E.2.12} says that the central map has precisely one
critical value and is isotrivial on its complement.
All hypotheses above apply to the central
threefold, and give the required contradiction.

\end{enumerate}
\end{proof}

\begin{proposition}[The central cusp is reduced and irreducible]\hypertarget{tag-0331}{}\label{push:src:E.2.16}
\StatementRecord{push:src:E.2.16}{proved}{0331}{}

In the one-sided smooth specialization of \ref{push:src:E.2.10},
the three labelled critical values remain distinct.
The two finite fibres have their original exact
multiplicities $m=3k$ and $n=4l$.  Every reduced
central critical fibre is irreducible, and the
cusp fibre has multiplicity one.  Its underlying
surface $S$ satisfies
\[
 \chi_{\rm top}(S)=2,\qquad b_1(S;\mathbb Q)=2.
\]
Consequently the central cusp is a reduced irreducible
Cartier surface.
\end{proposition}


\begin{proof}

\begin{enumerate}[label=\arabic*.,leftmargin=*]
\item \ref{push:src:E.2.12}, \ref{push:src:E.2.14} and \ref{push:src:E.2.15:smooth} exclude every degenerate
limit of the normalizing projective transformations.
They therefore extend as invertible projective
transformations, and the three labelled critical
values remain distinct.  \ref{push:src:E.2.13} proves the stated
finite-fibre multiplicities.

\item Normalize the three sections to $0,1,\infty$.
Over $B^\circ=\mathbb P^1\setminus\{0,1,\infty\}$,
the relative map is a proper submersion over
$B^\circ\times\Delta$.  Hence its integral
first-homology monodromies on the central fibre
are the same matrices $A_1,A_2,A_0$ as in
\ref{push:src:E.2.7}.  The common invariant space of the dual
matrices $A_1^t,A_2^t$ has dimension one:
in the displayed marking it is generated by
the covector $\gamma$.  Indeed, put
\[
 v_1=(A_1-I)\widehat u,\qquad
 v_2=(A_2-I)\widehat u,\qquad
 v_3=(A_1-I)\widehat w.
\]
The displayed matrices give the integral identities
\[
 \widehat u=-2v_2-v_3,\qquad
 \widehat w=-v_2-v_3,\qquad
 \widehat\delta=v_1-3v_2-2v_3.
\]
Both difference matrices have first row zero.
Their image lattices therefore generate exactly
$\ker\gamma$, whose annihilator is $\mathbb Z\gamma$.

\item Let $r$ be the total number of irreducible
components of the three reduced critical
fibres, and let $X^\circ$ be their complement
in $X_0$.  Alexander duality on $S^6$ gives
\[
 b_1(X^\circ;\mathbb Q)=r.
\]
The degree-one Leray sequence over the
punctured projective line, whose second
cohomology vanishes, gives instead
\[
 b_1(X^\circ;\mathbb Q)
 =b_1(B^\circ;\mathbb Q)
   +\dim H^1(T;\mathbb Q)^{\rm mon}
 =2+1=3.
\]
There are three nonempty critical fibres,
so each has exactly one reduced component.

\item Write the cusp fibre as $dS$ with $d\ge1$.
We prove $d=1$ by a direct meridian argument.
The smooth fibration over $B^\circ$ has
connected fibres, hence
$\pi_1(X^\circ)\to\pi_1(B^\circ)$ is
surjective.  The inclusion
$X^\circ\hookrightarrow X_0$ is surjective
on fundamental groups, and its kernel is
normally generated by meridians about the
three reduced divisor components.  To see
the latter assertion, make a null-homotopy
disc transverse to their smooth loci and
disjoint from their singular loci, which
have real codimension at least four.
Removing small discs about the finitely
many intersections expresses its boundary
as a product of conjugate meridians.

\item A meridian of a multiplicity-$a$ component
maps to the $a$th power of the corresponding
base meridian, since locally the base
function is a unit times the $a$th power
of a reduced local equation.  Consequently
there is a surjection
\begin{equation}
 \pi_1(X_0)\longrightarrow
 \left\langle a,b,c\ \middle|\
 a^m=b^n=c^d=abc=1\right\rangle.
\label{eq:II:next-E-2-16-1}\end{equation}
This quotient is the orbifold fundamental
group, and the displayed map is induced
by the base map.  Its source is trivial.

\item For clarity the group on the right is
nontrivial for every $d\ge2$ here.
Except for $(m,n,d)=(3,4,2)$, one has
\[
 \frac1m+\frac1n+\frac1d<1.
\]
Indeed for $d\ge3$ the sum is at most
$11/12$; for $d=2$, increasing either
$m=3k$ or $n=4l$ above $3,4$ makes it
at most $23/24$.  Put
$\alpha=\pi/m$, $\beta=\pi/n$,
$\theta=\pi/d$, and choose $\ell>0$ with
\[
 \cosh\ell=
 \frac{\cos\theta+\cos\alpha\cos\beta}
      {\sin\alpha\sin\beta}>1.
\]
Let
\[
 R_\alpha=\begin{pmatrix}
 \cos\alpha&\sin\alpha\\
 -\sin\alpha&\cos\alpha
 \end{pmatrix},\qquad
 D_\ell=\operatorname{diag}(e^{\ell/2},e^{-\ell/2}).
\]
In $\operatorname{PSL}_2(\mathbb R)$ the
elements represented by $R_\alpha$ and
$D_\ell R_\beta D_\ell^{-1}$ have orders
$m,n$, and their product has trace
$-2\cos\theta$ and projective order $d$.
Taking its inverse for $c$ gives a
nontrivial representation of \eqref{eq:II:next-E-2-16-1}.
For the exceptional triple, the
permutations $a=(123)$ and $b=(1324)$
have orders three and four, and
$ab=(24)$ has order two, giving a
nontrivial representation in $S_4$.
Thus $d\ge2$ contradicts \eqref{eq:II:next-E-2-16-1},
and the cusp is reduced.

\item The ordinary torus fibres have Euler
characteristic zero.  The finite reduced
fibres are smooth finite free quotients
of two-tori by \ref{push:src:E.2.13}, so also have
Euler characteristic zero.  Additivity
over the three singular values therefore
gives $\chi_{\rm top}(S)=\chi(X_0)=2$.
Finally, the irreducible-fibre
degree-one invariant-cycle statement
\cite[Corollary~3.7, p.~11, arXiv version~1]{HondaViaclovsky2025D}
gives
\[
 b_1(S;\mathbb Q)=
 \dim H^1(T;\mathbb Q)^{A_0}=2.
\]
Its irreducibility hypothesis was proved above.

\end{enumerate}
\end{proof}

\section{Normal-crossing cusps}\hypertarget{tag-P274}{}\label{sec:II:normal-crossing-cusps}

\begin{proposition}[Persistence of the normal-crossing cusp]\hypertarget{tag-0332}{}\label{push:src:E.2.17}
\StatementRecord{push:src:E.2.17}{proved}{0332}{}

In the smooth fixed-pair specialization of \ref{push:src:E.2.10},
the cusp divisor is a relative normal-crossing divisor
over the parameter disc.  More precisely, after
shrinking that disc, at every point of the total cusp
divisor there are holomorphic coordinates
$(z_1,z_2,z_3,s)$, with $s$ the original parameter,
in which the cusp coordinate of the intrinsic map is
one of
\begin{equation}
 q=z_1,\qquad q=z_1z_2,\qquad q=z_1z_2z_3.
\label{eq:II:next-E-2-17-1}\end{equation}
There are exactly two relative triple-point sections.
Both lie on the fixed rational curve of \ref{push:src:E.2.14}.
In particular the normalization of the total cusp
divisor is smooth and proper over the parameter
disc, and its central fibre is the normalization
of the central cusp.
\end{proposition}


\begin{proof}

\begin{enumerate}[label=\arabic*.,leftmargin=*]
\item By \ref{push:src:E.2.14}--\ref{push:src:E.2.16} the three labelled critical
values remain distinct, the cusp has multiplicity
one, and its reduced central fibre is irreducible.
We normalize its target section to $q=0$.
On a sufficiently small target disc, zero is the
only critical value on every parameter fibre.
The extended field $\delta$ is zero precisely on
the relative fixed rational curve $\mathcal C$,
and the two normal characters there are $+1,-1$.
For $s\ne0$ the two triple points of the prescribed
cusp belong to $\mathcal C_s$; outside that curve
the cusp is smooth or has an ordinary double crossing.

\item We first record exactly the local singularity result
we use.  If $h_s(a,b)$ is a holomorphic family and
$h_0=0$ is a reduced plane-curve germ with an
isolated critical point, choose a sufficiently
small representative of that germ.  The sum of
the multiplicities of all nearby critical points
of $h_s$, regardless of their critical values,
is $\mu(h_0)$.  One can see this directly from
the analytic critical scheme
\[
 \mathbb C\{a,b,s\}/
       ((h_s)_a,(h_s)_b).
\]
After taking a sufficiently small representative,
it is finite over the parameter disc.  Its two
equations are a regular sequence, and its
central fibre is zero-dimensional.  It is thus
Cohen--Macaulay of dimension one, with $s$ a
nonzerodivisor, and is finite flat over the
disc.  Its degree is the central Jacobian-algebra
dimension, namely $\mu(h_0)$.

\item The non-splitting theorem of
\citet[Theorem 2 and its corollary, pp.~274--275]{Lazzeri1973Root}
says that if a deformation of an isolated
hypersurface singularity has more than one
singular point in a single deformed fibre,
their summed Milnor numbers are strictly
smaller than the original Milnor number.
The strict numerical inequality is the
conclusion of the cited corollary's proof.
Consequently, if all the critical points
just counted have critical value zero,
there can be only one such point.  This
application uses a flat deformation of
plane hypersurfaces: the map
$(a,b,s)\mapsto(h_s(a,b),s)$ is flat near
the central germ because its two base
parameters form a regular sequence in
the smooth source.

\item Let $p$ be a central cusp point outside
$\mathcal C_0$.  A relative holomorphic
flow box for the nonzero field $\delta$
gives coordinates $(a,b,v,s)$ in which
$\delta=\partial_v$.  Invariance of the
intrinsic map gives
\[
 q=h_s(a,b).
\]
The central equation is reduced, by
\ref{push:src:E.2.16}.  A reduced plane-curve germ is
either smooth or has an isolated critical
point.  In the smooth case the implicit
function theorem already gives the first
form in \eqref{eq:II:next-E-2-17-1}.

\item In the singular case every critical point
of $h_s$ in a small representative lifts
to a critical point of the map
on $X_s$.  All its critical values are
therefore zero.  For $s\ne0$ each such
point is a node, since the flow box is
disjoint from the fixed curve and hence
from the two triple points.  Conservation
and non-splitting give one node and
$\mu(h_0)=1$.  Thus the central Hessian
is nondegenerate.  The implicit function
theorem gives its unique critical-point
section, whose critical value is zero.
The holomorphic Morse lemma with
parameters then gives $q=ab$ throughout
this neighbourhood, with $s$ unchanged.

\item It remains to treat the fixed curve.
Holomorphic coordinates obtained by
averaging over the extended compact
circle may be chosen in the form
$(z,x,y,s)$, with $\mathcal C$ defined
by $x=y=0$ and with circle weights
$0,+1,-1,0$.  This averaging fixes the parameter $s$.
The holomorphic invariant function $q$
vanishes on $\mathcal C$, so its
convergent power series has the form
\begin{equation}
 q=F_s(T,z)=T g_s(T,z),\qquad T=xy.
\label{eq:II:next-E-2-17-2}\end{equation}
The coefficient
$Q_s(z)=g_s(0,z)$ is intrinsically a
section of the dual of the product of
the two normal weight lines along
$\mathcal C_s$.  By \ref{push:src:E.2.14} this is
$\mathcal O_{\mathbb P^1}(2)$.

\item The central section $Q_0$ is not
identically zero.  Otherwise the
central equation in every such chart
would be divisible by $x^2y^2$,
contradicting reducedness of the
central cusp.  Its zero divisor
therefore has degree two.  Suppose
that both zeros coincide at a point
$z=0$.  The plane equation
$F_0(T,z)=Tg_0(T,z)$ is reduced:
a repeated factor of $g_0$ would
remain repeated after $T=xy$, and
$T$ does not divide $g_0$ because
$Q_0$ is not identically zero.
Thus $F_0$ has an isolated plane
critical point there.

\item For $s\ne0$ the two zeros of $Q_s$
are the two distinct triple points
of the prescribed cusp.  Under the
supposed collision they both lie
in the chosen small plane
representative.  They are distinct
critical points of $F_s$ with value
zero.  Every other critical point
of $F_s$ also has value zero:
this is automatic at $T=0$, and
at $T\ne0$ it follows by lifting
to $xy=T$, where the quotient
map $(x,y,z)\mapsto(T,z)$ is a
submersion.  All sufficiently small
$T$ have such a lift in the original
coordinate chart.  A nonzero
critical value would contradict
the critical-value set of
the intrinsic map.

\item Conservation of all critical
multiplicities and Lazzeri's
non-splitting theorem now give a
contradiction to the two distinct
critical points.  Hence $Q_0$ has
two simple zeros.  They extend
as two disjoint holomorphic
sections over the parameter disc.
Where $Q_0\ne0$, the factor $g_s(xy,z)$ is a unit.
The local biholomorphism
\[
 (x,y,z,s)\longmapsto
 (x,yg_s(xy,z),z,s)
\]
gives $q=xy$ in its target coordinates.
At a simple zero, $\partial g_0/\partial z\ne0$, and
\[
 (x,y,z,s)\longmapsto
 (x,y,g_s(xy,z),s)
\]
is a local biholomorphism giving $q=xyz$.
Both maps retain the original parameter $s$.  This proves
all local forms in \eqref{eq:II:next-E-2-17-1}.

\item Properness makes the central cusp
compact, so finitely many such
neighbourhoods cover it and a
single shrinking of the parameter
disc suffices.  The normalization
of a relative normal-crossing
divisor separates its local
branches.  In the coordinates
above every normalized branch
is smooth over $s$, and these
descriptions agree on overlaps.
Normalization is finite for
complex analytic spaces; its
composition with the proper
cusp map is therefore proper.
Its fibres are precisely the
normalizations of the cusp
fibres in these local charts.
This proves the final assertions.

\end{enumerate}
\end{proof}

\begin{proposition}[The normalized cusp and its marked hexagon]\hypertarget{tag-0333}{}\label{push:src:E.2.18}
\StatementRecord{push:src:E.2.18}{proved}{0333}{}

Let $\nu:\mathcal S\to\mathcal W$ be the
simultaneous normalization of the relative cusp
in \ref{push:src:E.2.17}, and let $\mathcal D\subset\mathcal S$
be its reduced conductor.  After shrinking the
parameter disc, each fibre pair is the standard
toric pair
\[
 (\mathcal S_s,\mathcal D_s)
 \simeq(\mathrm{dP}_6,\hbox{its toric hexagon}).
\]
The six conductor components may be labelled over
the disc so that their incidence, self-intersections
$-1$, opposite-pair involution, and endpoint
identifications are those of the punctured family.
The relative logarithmic tangent bundle has a
global holomorphic frame:
\[
 T_{\mathcal S/\Delta}(-\log\mathcal D)
 \simeq\mathcal O_{\mathcal S}^{\oplus2}.
\]
\end{proposition}


\begin{proof}

\begin{enumerate}[label=\arabic*.,leftmargin=*]
\item In the relative charts of \ref{push:src:E.2.17}, the normalized
branches are smooth over the original parameter.
Their conductor branches and their intersections
are also smooth over that parameter.  The
normalization of the conductor is consequently
a proper smooth family of curves, and the
normalized intersection strata are finite
\'etale over the disc.  Their connected components
and incidence are constant.  Since the disc
is simply connected, label the six curve
families $\mathcal D_1,\ldots,\mathcal D_6$
and the six vertex sections by continuation
from a punctured fibre.  Each curve fibre
is a smooth curve of genus zero.  Its two
distinct vertex sections remain distinct;
the same relative local charts exclude a
new self-intersection or an additional
incidence at the centre.  Thus the central
conductor is a hexagon of six embedded
rational curves.  The degree of each
relative normal line is constant in its
proper curve family, so all six central
self-intersections equal $-1$.

\item Write $S=\mathcal S_0$ and $D_i=\mathcal D_{i,0}$.
Ehresmann transport gives $b_1(S)=0$ and
$b_2(S)=4$.  Contract the three disjoint
curves $D_1,D_3,D_5$.  The analytic
Castelnuovo contraction gives a smooth
compact complex surface $Y$; the needed
analytic statement, including smoothness
of the resulting point, is
\citet[\S4, no.~8(a), pp.~364--365]{Grauert1962E}.
These three contractions take place on the central surface.

\item The images $C_2,C_4,C_6$ of the remaining
curves are smooth rational curves.  Each
meets two contracted curves transversely
once, so has self-intersection
$-1+1+1=1$.  Each pair meets once, at
the image of their intervening contracted
curve; these three intersection points
are distinct and the intersections are
transverse.  Thus they form a triangle.
Also $b_1(Y)=0$ and $b_2(Y)=1$.

\item Any one of these curves, say $C=C_2$,
has nonnegative intersection with every
effective divisor: distinct irreducible
curves have nonnegative local intersection
multiplicities and $C^2=1$.  Adjunction
gives $K_Y\cdot C=-3$.  A nonzero section
of $K_Y^{\otimes r}$ for $r>0$ would
give an effective divisor of intersection
$-3r$ with $C$, impossible.  Hence
$\kappa(Y)=-\infty$.  The compact-surface
classification
\cite[Chapter~VI, Theorem~1.1 and Table~10, pp.~243--244]{BHPV2004E}
together with $b_1(Y)=0$ makes $Y$ rational:
class~VII has first Betti number one,
and a ruled surface has the first Betti
number of its base curve.  To exhibit its map to $\mathbb P^2$, use $H^1(Y,\mathcal O_Y)=0$
and the exact sequence
\[
 0\longrightarrow\mathcal O_Y\longrightarrow\mathcal O_Y(C)
 \longrightarrow\mathcal O_C(1)\longrightarrow0.
\]
The canonical section vanishing on $C$, together with lifts of a
basis of $H^0(C,\mathcal O_C(1))$, has no common zero and defines
\[
 \phi:Y\longrightarrow\mathbb P^2,\qquad
 y\longmapsto[\sigma_0(y):\sigma_1(y):\sigma_2(y)].
\]
Its image has dimension two because $C^2=1$, and its degree is one.
The intersection form of $Y$ has rank one and is generated over
$\mathbb R$ by $[C]$, with square one.  Since $Y$ is projective,
every effective curve has positive intersection with $C$.
Thus $\phi$ contracts no curve, so it is finite of degree one and
hence an isomorphism.  The images of $C_2,C_4,C_6$ are lines because
$\phi^*\mathcal O(1)=\mathcal O_Y(C)$ and $C\cdot C_j=1$.
Their three pairwise intersection points
are distinct, so they form a projective
triangle.  Blowing up its three vertices
recovers exactly the toric degree-six
del Pezzo surface and its hexagonal
boundary.  This proves the assertion
about the central pair; the punctured
pairs were already the prescribed ones.

\item The conductor involution is holomorphic
over the disc.  Its permutation of the
six labelled curve families, and its
permutation of the vertex sections,
are locally constant.  They are therefore
the same opposite-pair and endpoint
identifications as on a punctured fibre.

\item Finally put
$\mathcal E=T_{\mathcal S/\Delta}(-\log\mathcal D)$.
For every fibre pair just identified,
the toric infinitesimal action, for the one-parameter-subgroup lattice
$N\simeq\mathbb Z^2$, sends $n\in N$ to the logarithmic derivation
$\chi^m\mapsto\langle m,n\rangle\chi^m$ for $m\in N^\vee$.
On each smooth toric chart this identifies a lattice basis with
a frame of the logarithmic tangent bundle.  Thus
\[
 \mathcal E|_{\mathcal S_s}\simeq
 \mathcal O_{\mathcal S_s}^{\oplus2},
 \qquad h^0(\mathcal S_s,\mathcal E_s)=2.
\]
Constant-cohomology base change
\cite[\S7, Satz~5, p.~64; hypotheses and proof, pp.~62--64]{Grauert1960D} makes the direct
image a rank-two holomorphic vector bundle on the disc and
identifies its fibres with these section spaces.  Its hypotheses
hold because the map is proper, the disc is reduced,
$\mathcal E$ is locally free and flat over the disc, and the
displayed section dimension is constant.
Choose a frame after shrinking the
disc.  The evaluation map from its
pullback to $\mathcal E$ is an
isomorphism on every fibre, hence
is an isomorphism everywhere by
Nakayama's lemma.  This gives the claimed global holomorphic relative frame.

\end{enumerate}
\end{proof}

\section{Torus actions and analytic cusps}\hypertarget{tag-P275}{}\label{sec:II:analytic-cusp-identification}

\begin{proposition}[Extension of the marked cusp torus action]\hypertarget{tag-0334}{}\label{push:src:E.2.19}
\StatementRecord{push:src:E.2.19}{proved}{0334}{}

Let
\[
 f:\mathcal N\longrightarrow
 B=\Delta_q\times\Delta_s
\]
be the proper cusp neighbourhood of the smooth
specialization in \ref{push:src:E.2.17}, and put
$\mathcal W=f^{-1}(q=0)$.  After shrinking the
bidisc, the rank-two vertical logarithmic tangent
bundle is trivial and is generated by two commuting
complete holomorphic vector fields.  Their action
factors through a holomorphic action of
$(\C^*)^2$ over $B$, with the integral marking
given by the periods $\mathbf w,\boldsymbol\delta$.
On the normalization of every $\mathcal W_s$
this is its usual toric action.

The marked rank-two lattice is primitive in the
period lattice of each ordinary torus fibre and
equals the full invariant lattice of the cusp
monodromy.  At the two relative triple-point
sections, the three nonzero tangent characters
are constant in $s$ and equal to those of the
prescribed Mumford cusp.  In the toric chart
lattice coordinates, transported from that
punctured model by their fixed integral change
of basis from the period marking, they are
respectively
\begin{equation}
 \{(1,0),(0,1),(-1,-1)\},\qquad
 \{(-1,0),(0,-1),(1,1)\}.
\label{eq:II:next-E-2-19-1}\end{equation}
\end{proposition}


\begin{proof}
\begin{enumerate}[label=\arabic*.,leftmargin=*]
\item \emph{The logarithmic bundle and its restriction.}
The coordinates in \ref{push:src:E.2.17} make $\mathcal W$ a
relative normal-crossing divisor over $s$.
Define
\begin{equation}
 \mathcal L=
 \ker\left[
 d\log q:
 T_{\mathcal N/\Delta_s}(-\log\mathcal W)
       \longrightarrow\mathcal O_{\mathcal N}
 \right].
\label{eq:II:next-E-2-19-2}\end{equation}
The map is surjective: if
$q=z_1\cdots z_k$, it sends each
$z_i\partial_{z_i}$, $i\leq k$, to one.
Consequently $\mathcal L$ is locally free of
rank two.  Its sections are holomorphic
vector fields, tangent to the branches, and
annihilate both $q$ and $s$.

Let $\nu:\mathcal S\to\mathcal W$ and
$\mathcal D\subset\mathcal S$ be the
simultaneous normalization and conductor.
Restriction to a branch gives a canonical
isomorphism
\begin{equation}
 \nu^*(\mathcal L|_{\mathcal W})
 \simeq T_{\mathcal S/\Delta_s}(-\log\mathcal D).
\label{eq:II:next-E-2-19-3}\end{equation}
At each stratum the displayed coordinate vector fields give the map. For
$q=xy$, a basis is
$x\partial_x-y\partial_y,\partial_z$;
on the branch $x=0$ it becomes
$-y\partial_y,\partial_z$, a logarithmic
basis there.  For $q=xyz$, use
$x\partial_x-z\partial_z$ and
$y\partial_y-z\partial_z$.  Their
restrictions to $x=0$ have coefficient
matrix of determinant one in the basis
$y\partial_y,z\partial_z$.
The one-branch case is immediate.
The same computation keeps $s$ fixed.

By \ref{push:src:E.2.18} the right-hand side of
\eqref{eq:II:next-E-2-19-3} has a global holomorphic
frame.  On the two preimages of each
conductor curve, the identifications
in \eqref{eq:II:next-E-2-19-3} describe descent by an
invertible two-by-two matrix.  Its
entries are holomorphic on a proper
$\mathbb P^1$ family over $\Delta_s$,
so are constant along each curve
fibre and are holomorphic functions
of $s$.  Denote these matrices by
$G_j(s)$.

For $s\ne0$ both invariant logarithmic
fields in the prescribed Mumford model
descend.  Indeed, its period maps in
multiplicative coordinates multiply
the two coordinates by units depending
only on the base.  Their derivative
on the two vertical logarithmic
fields is the identity.  The toric
small-resolution charts preserve
these same fields.  Thus the entire
two-dimensional space of sections
on the normalized pair descends for
$s\ne0$.  In any chosen frame this
means $G_j(s)=I$, and the identity
theorem gives $G_j(0)=I$ as well.

At a triple point the descent map is explicit.
Let $f_1(y,z,s),f_2(x,z,s),f_3(x,y,s)$ be functions
on the three normalized branches with equal pairwise restrictions
$g_{12}(z,s),g_{13}(y,s),g_{23}(x,s)$ and common restriction $h(s)$
to the triple locus.  The function
\[
 F=f_1+f_2+f_3-g_{12}-g_{13}-g_{23}+h
\]
in $\C\{x,y,z,s\}/(xyz)$ restricts to each $f_i$.
Conversely, restrictions of any function agree on all pairwise
intersections.  The restriction map is injective because
$(x)\cap(y)\cap(z)=(xyz)$.  This proves the equalizer assertion;
the two-branch formula is $f_1+f_2-g_{12}$.
Tensoring these maps with the locally free bundle
$\mathcal L|_{\mathcal W}$ gives the same descent for its sections.
Consequently the chosen normalized
frame descends, and
\begin{equation}
 \mathcal L|_{\mathcal W}
       \simeq\mathcal O_{\mathcal W}^{\oplus2}.
\label{eq:II:next-E-2-19-4}\end{equation}

\item \emph{A frame on the whole neighbourhood.}
The map $f$ is proper and flat; this
also follows locally from the equations
$(q,s)=(z_1\cdots z_k,s)$.
The bundle $\mathcal L$ is flat over
$B$.  On each ordinary fibre it is
the tangent bundle of a complex
two-torus.  On each $q=0$ fibre it
is trivial by \eqref{eq:II:next-E-2-19-4}.  All these
fibres are compact and connected,
so in every case
\[
 h^0(\mathcal N_{q,s},
       \mathcal L|_{\mathcal N_{q,s}})=2,
\]
and those sections generate the
bundle pointwise.  Since $f$ is proper, $\mathcal L$ is flat over the reduced base $B$,
and the displayed dimension is constant,
\cite[\S7, Satz~5, p.~64; hypotheses and proof, pp.~62--64]{Grauert1960D} makes $f_*\mathcal L$
a rank-two vector bundle on $B$ and identifies its fibres with
these spaces of sections.
Its evaluation map
\begin{equation}
 f^*f_*\mathcal L\longrightarrow\mathcal L
\label{eq:II:next-E-2-19-5}\end{equation}
is an isomorphism on every fibre,
and hence everywhere by Nakayama's
lemma.  Shrink the bidisc and
choose a holomorphic frame $V_1,V_2$.

These fields commute.  On every
ordinary compact torus, every
holomorphic vector field is a
translation field, so their
bracket vanishes there.  The
ordinary locus is dense, and
the bracket is holomorphic.
Each field is complete because
it is vertical and every fibre
of the proper map is compact:
the holomorphic local flow can
be continued for every complex
time on that compact invariant
fibre.  Holomorphic dependence
on initial points and parameters
gives a holomorphic action
\[
 \operatorname{Act}:\C^2\times
       \mathcal N\longrightarrow\mathcal N
\]
over $B$.

\item \emph{Holomorphic extension of the two toric periods.}
Choose a smooth point $p_0$ of the
central cusp away from its conductor.
The map $f$ is a submersion there,
so there is a holomorphic section
$p(q,s)$ through $p_0$.
On the central normalized toric
pair the two fields are a basis
of the toric logarithmic Lie
algebra.  Its open orbit is
$(\C^*)^2$.  Let
$v_1^0,v_2^0\in\C^2$ be a
basis of its full period lattice
in the coordinates $V_1,V_2$.
They are complex-linearly
independent.

Near each $v_i^0$, write
$\operatorname{Act}(v,p(q,s))=p(q,s)$
using two vertical coordinates
in a submersion chart around
$p_0$.  The derivative in $v$
is invertible at $(v_i^0,0,0)$,
since the infinitesimal action
is an isomorphism on that open
orbit.  The holomorphic implicit
function theorem gives a unique
holomorphic vector $v_i(q,s)$
near $v_i^0$ solving this equation.
For $q\ne0$, fixing $p(q,s)$
means that this translation is
the identity on the entire
torus fibre.  The equality of
holomorphic maps then extends
over $q=0$.  Thus $v_i(q,s)$
is a period on every fibre.
Their complex-linear
independence persists after
shrinking $B$.

Use these two columns as a
change of the frame.  Because
their coefficients are pulled
back from $B$, the new fields
still commute and are complete.
They have period one, and the
action consequently factors
through
$\C^2/\Z^2\simeq(\C^*)^2$,
using the exponential
$z\mapsto e^{2\pi i z}$
in each coordinate.

\item \emph{Primitivity and the integral marking.}
On $B^\circ=\Delta_q^*\times\Delta_s$
the period lattices of the
translation action are the
integral first-homology local
system of the torus family.
The continuously varying
periods $v_1,v_2$ represent
flat integral classes: in a
local differentiable
trivialization their flow
loops vary continuously, and
integral homology is discrete.
They generate a rank-two
invariant sublattice $L$.

Suppose $L$ were not primitive
at some ordinary fibre.  There
would be a rational pair
$(a_1,a_2)\notin\Z^2$ such
that $a_1v_1+a_2v_2$ is an
integral period there.  This
is an equality of flat
integral classes after
multiplication by a common
denominator, so it persists
throughout the connected
ordinary base.  In particular
the corresponding fixed
finite-order element of
$(\C^*)^2$ acts identically
there.  By the identity
theorem it acts identically
on the central open orbit.
This contradicts the choice
of $v_1^0,v_2^0$ as the full
central toric period lattice.
Thus $L$ is primitive.

The marked cusp monodromy,
constant in $s$ by \ref{push:src:E.2.16},
has the form
\begin{equation}
 \widehat\gamma\longmapsto
       \widehat\gamma-\widehat\delta,
 \qquad
 \widehat u\longmapsto\widehat u+\widehat w,
 \qquad
 \widehat w\longmapsto\widehat w,
 \qquad
 \widehat\delta\longmapsto\widehat\delta.
\label{eq:II:next-E-2-19-6}\end{equation}
For an integral vector
$v=a\widehat\gamma+b\widehat u+c\widehat w+d\widehat\delta$,
the difference operator satisfies
\[
 (T-I)v=b\widehat w-a\widehat\delta.
\]
Its full integral kernel is therefore
$\Z\widehat w\oplus\Z\widehat\delta$;
projection onto the last two coordinates splits its inclusion
in the integral period lattice.  Since $L$ is primitive in that
period lattice and has the same rank as this kernel, they are equal.
There is consequently a matrix $U\in\operatorname{GL}_2(\Z)$ with
$(\mathbf w,\boldsymbol\delta)=(v_1,v_2)U$.
The corresponding change of acting torus is the monomial isomorphism
\[
 (z_1,z_2)\longmapsto
 (z_1^{U_{11}}z_2^{U_{12}},\,
  z_1^{U_{21}}z_2^{U_{22}}),
\]
whose inverse is the monomial map of $U^{-1}$.
It gives precisely the prescribed period marking.

\item \emph{Characters at the triple points.}
The two triple-point sections
are fixed by the torus action:
on the normalization they
are toric vertices.  Restrict
the derivative representation
to the compact real torus
$(S^1)^2$ and average a
Hermitian metric along either
section.  Its characters
are integral pairs and
vary continuously in $s$;
hence their multiset is
locally constant.  Equivalently,
Fourier averaging gives the
individual character
subbundles.  They coincide
at $s=0$ with the known
characters for $s\ne0$.

The two unit
triangles in the prescribed
cocharacter lattice have
height-one rays
\[
 \{(0,0,1),(1,0,1),(0,1,1)\},
 \qquad
 \{(1,0,1),(0,1,1),(1,1,1)\}.
\]
On the dense torus, with coordinates $(z,w,q)$, the corresponding
affine-chart monomials are respectively
\[
 (X_0,X_1,X_2)=\left(\frac{q}{zw},z,w\right),
 \qquad
 (X_0,X_1,X_2)=\left(\frac{q}{w},\frac{q}{z},\frac{zw}{q}\right).
\]
Their inverse monomial maps are
\[
 (z,w,q)=(X_1,X_2,X_0X_1X_2),
 \qquad
 (z,w,q)=(X_0X_2,X_1X_2,X_0X_1X_2).
\]
The exponent rows pair with the listed rays as the identity matrix,
so they give the chart coordinates.  Their first two entries are
the character triples in \eqref{eq:II:next-E-2-19-1}, in the corresponding chart order.
Their sum is zero after restriction to the vertical torus, while
$X_0X_1X_2=q$.  Thus the characters, action and marking have the
asserted values.

\end{enumerate}
\end{proof}

\begin{proposition}[Identification of the analytic cusp]\hypertarget{tag-0336}{}\label{push:src:E.2.20}
\StatementRecord{push:src:E.2.20}{proved}{0336}{}

In the smooth specialization of \ref{push:src:E.2.10}, normalize
the three distinct critical-value sections to
$0,1,\infty$ and write $q=1/t$ near the cusp.
Let $c(s)$ be the holomorphic interior period
lift of \ref{push:src:E.2.7}.  After shrinking the
$(q,s)$-bidisc, its inverse image in the 
family is biholomorphic, over that bidisc, to
the prescribed marked Mumford cusp family at
parameter $c(s)$.  The isomorphism induces
the given identification of the regular
torus Jacobians and their integral markings.
Its affine origin may be fixed by choosing
a holomorphic section through a smooth
point of the reduced cusp.
\end{proposition}


\begin{proof}
Write $f:\mathcal N\to\Delta_q\times\Delta_s$
for the proper cusp neighbourhood and
$\mathcal W=f^{-1}(q=0)$.  \ref{push:src:E.2.16}--\ref{push:src:E.2.18}
give its reduced relative normal-crossing
divisor, its normalized toric hexagon, and
the labelled opposite-side pairing and endpoint identifications.
\ref{push:src:E.2.19} gives a holomorphic action
of $H=(\C^*)^2$ over the bidisc.  Its
period marking is the primitive vanishing
lattice
\begin{equation}
 \Lambda_{\rm tor}
   =\Z\widehat w\oplus\Z\widehat\delta.
\label{eq:II:next-E-2-20-1}\end{equation}
We reconstruct the neighbourhood from
this action, keeping track of the
extension of its characters.

\begin{enumerate}[label=\arabic*.,leftmargin=*]
\item \emph{The topology and the covering.}
The relative coordinates
$q=x,xy,xyz$ give smooth real lifts
of the two real parameter directions
of $s$ which annihilate $dq$ and
are tangent to every cusp stratum.
Away from the cusp such lifts exist
because $f$ is a submersion.
A smooth partition of unity preserves
both these linear requirements.
On the inverse image of a smaller
closed $q$-disc the lifts integrate
along a sufficiently short parameter
path: $q$ remains constant, and
properness prevents escape.
Thus their flows identify the cusp
neighbourhoods, all strata, and
the regular-fibre inclusion maps
by $q$-preserving diffeomorphisms.

For the prescribed Mumford cusp,
\cite[Theorem~4.5, pp.~26--28; Proposition~4.7 and Corollary~4.8,
pp.~30--31]{AC26} give its simply connected toric cover and the free
period quotient, with deck group
\[
 \overline\Lambda=\Lambda/\Lambda_{\rm tor}\simeq\mathbb Z^2.
\]
On a regular fibre the covering is
$\mathbb C^2/\Lambda_{\rm tor}\to\mathbb C^2/\Lambda$.
Its fundamental-group homomorphism is the quotient
$\Lambda\to\overline\Lambda$; the $q$-meridian contracts along
$(x,y,z)=(q,1,1)$. The $q$-preserving diffeomorphisms above
transport this marking, the fibre-inclusion homomorphism and the
meridian contraction to $\mathcal N$.

Take the universal covering
\[
 \widetilde{\mathcal N}\longrightarrow\mathcal N.
\]
Its deck group is
$\overline\Lambda$, and its
central normalized components
form the planar unfolding of
the hexagon.  The $H$ action
lifts.  Indeed the fundamental
loops of $H$ at a central open
orbit lie in one normalized
del Pezzo component, which is
simply connected, and hence die
in $\pi_1(\mathcal N)$.
The identity lift is unique
and satisfies the action law.
Choose a holomorphic section
$p(q,s)$ through a smooth
central cusp point and lift
it to the cover.  On $q\ne0$
its orbit map is an isomorphism
\begin{equation}
 H\times(\Delta_q^*\times\Delta_s)
       \xrightarrow{\ \sim\ }
 \widetilde{\mathcal N}|_{q\ne0}.
\label{eq:II:next-E-2-20-2}\end{equation}
The stabilizer has been removed
exactly by \eqref{eq:II:next-E-2-20-1}, using
the identified regular-fibre
inclusion map.  In particular
the two characters of $H$
are globally defined on this
open covering space.

Connectedness of $H$ makes
its lifted action commute
with every deck transformation:
conjugation acts continuously
on the discrete deck group
and is the identity at the
identity of $H$.  On
\eqref{eq:II:next-E-2-20-2}, the two deck
generators therefore act by
\begin{equation}
 (h,q,s)\longmapsto(P_i(q,s)h,q,s),
 \qquad i=1,2.
\label{eq:II:next-E-2-20-3}\end{equation}
The functions $P_i$ are the
marked period multipliers.
Changing the affine origin
does not change them.
The interior lift in \ref{push:src:E.2.7} and the regular torus family
in \ref{push:src:E.2.12} identify these multipliers on the whole
punctured bidisc.  To check the central parameter explicitly,
trivialize the integral local system on a simply connected open set
with $q\ne0$.  Periods of a holomorphic frame of one-forms give a
holomorphic kernel in $\operatorname{Gr}(2,\Lambda\otimes\C)$.
The prescribed marked matrix is
\[
 \Pi(q,s)=[Z(q,s)\mid I_2],\qquad
 Z(q,s)=\frac{\log q}{2\pi i}B_0+C(q,c(s)),
 \quad
 B_0=\begin{pmatrix}0&1\\-1&0\end{pmatrix},
\]
with $C$ holomorphic at $q=0$
\cite[Proposition~3.21, p.~23; Proposition~4.7, p.~30]{AC26}.
The period kernel of the family and the prescribed kernel agree for $s\ne0$;
their holomorphic Grassmannian maps therefore agree at $s=0$.
The induced isomorphism of the quotient Lie bundles carries every
integral period to the corresponding column of $\Pi$.
These comparisons agree on marking overlaps.  Exponentiating gives
\[
 P_{\bar\lambda}(q,s)
 =q^{B_0\bar\lambda}
   \exp\!\bigl(2\pi i C(q,c(s))\bar\lambda\bigr),
 \qquad \bar\lambda\in\overline\Lambda.
\]
The exponential is a holomorphic torus-valued unit on the full bidisc.

\item \emph{Character extension across adjacent components.}
For a central component $D_v$,
choose a holomorphic base
section $a_v(q,s)$ through
its dense $H$ orbit.  Its
orbit map is an open
holomorphic embedding
\[
 H\times(\Delta_q\times\Delta_s)
       \longrightarrow\widetilde{\mathcal N}.
\]
It is injective on the
central open orbit and on
each regular fibre, and
its differential is invertible
along those loci.  Its
image contains the whole
regular covering.  Two
such origins differ for
$q\ne0$ by a holomorphic
translation $h_{vw}(q,s)$
in $H$.  The following calculation gives its meromorphic extension at $q=0$.

If both origins lie in
the same central component,
their ratio extends as a
holomorphic $H$-valued
unit.  The central open
orbit is free and transitive,
so this follows from the
holomorphic inverse-function
theorem for the orbit map,
fixing both parameters.

For two adjacent components,
use a crossing chart with
$q=ab$ and an orbit coordinate
$z\ne0$.  The stabilizer
of a general conductor
point is a primitive
connected one-dimensional
subtorus of $H$.  Its
two normal characters are
$+1,-1$: on each normalized
toric component these are
the primitive boundary
characters, and invariance
of $q$ makes them opposite.
Compact-circle averaging
gives these exact local
linear coordinates.  A
complementary torus direction
provides the orbit coordinate;
an invariant unit in $q$
can be absorbed while
preserving the weights.

Fix small nonzero constants
$a_0,b_0,z_0$ in this chart
and take the two special
origins
\begin{equation}
 (a,b,z)=(q/b_0,b_0,z_0),\qquad
 (a,b,z)=(a_0,q/a_0,z_0).
\label{eq:II:next-E-2-20-4}\end{equation}
The stabilizer cocharacter
evaluated at $a_0b_0/q$
sends the first to the
second.  This equality
is justified by the local
action itself: a path
in that subgroup can
increase the modulus of
$a$ and decrease that
of $b$, with $ab=q$
fixed, staying inside
the crossing chart;
the phase can be varied
without changing either
modulus.  This path gives the stated action between the two origins.

Comparing \eqref{eq:II:next-E-2-20-4} with
arbitrary chosen origins
adds only holomorphic
$H$-valued units, by the
same-component argument.
The adjacent comparison
therefore has the form
\begin{equation}
 h_{vw}(q,s)=q^{n_{vw}}u_{vw}(q,s),
 \qquad n_{vw}\in\Z^2,
\label{eq:II:next-E-2-20-5}\end{equation}
where $u_{vw}$ is holomorphic
with values in $H$ on the
whole bidisc.  Every central
component is reached by
a finite adjacent path.
Multiplying these expressions
proves the same assertion
for its comparison with
the initial origin.
Different paths agree
on the regular cover,
and hence define identical
meromorphic characters.
The exponents are locally
constant in $s$, since
\eqref{eq:II:next-E-2-20-5} holds with a
nonvanishing unit on the
full bidisc.

\item \emph{The fan and the analytic coordinate maps.}
At either triple-point
section the whole $H$
fixes the point.
Compact-torus averaging
gives three eigen-coordinates
with the constant integer
characters of \ref{push:src:E.2.19}.
Their only primitive
positive invariant relation
is their product.
The reduced normal-crossing
equation therefore writes
$q$ as that product times
a holomorphic invariant
unit.  The unit is a
function of the product
and $s$; absorb it into
one eigen-coordinate to
obtain exactly $q=xyz$.

The orders of the two
global characters along
$D_v$, together with
$\operatorname{ord}_{D_v}q=1$,
give a ray $(v,1)$.
At a triple point the
three incident rays
form a unimodular cone.
Indeed all adjacent
orders and tangent
characters are constant
in $s$ and coincide
with the known ones
for $s\ne0$.
The entire resulting
fan is consequently
the prescribed $A_2$
triangulation, with
its fixed labels.

In a triple chart a
global character is
a Laurent monomial in
$x,y,z$ times a function
of $q,s$.  Evaluate it
on a section keeping
two nonzero coordinates
fixed and solving the
third from $q$.
The origin comparisons
above show that the
remaining function is
a power of $q$ times
a holomorphic unit.
Write the three eigen-coordinates as $z_0,z_1,z_2$ and the
global torus characters as $h_1,h_2$.
Let $m_j=(\alpha_{j1},\alpha_{j2},b_j)$ be the rows dual to the
three incident ray vectors.  The chart map is
\[
 X_j=h_1^{\alpha_{j1}}h_2^{\alpha_{j2}}q^{b_j}
     =z_jU_j(q,s),\qquad j=0,1,2,
\]
where the functions $U_j$ are holomorphic units by the preceding
origin comparison.  The dual-basis identity
$\sum_jm_j=(0,0,1)$ gives
$X_0X_1X_2=q=z_0z_1z_2$ and $\prod_jU_j=1$.
The inverse is the explicit holomorphic map
\[
 z_j=X_j/U_j(q,s),\qquad q=X_0X_1X_2,
 \qquad s=s.
\]
The two maps compose to the identity because they preserve $q,s$.
They give local biholomorphisms through the triple points and,
by their torus translates, along the other strata.
On overlaps they agree on the dense regular cover and hence
everywhere, so they glue.

We finally check the
common neighbourhood
and the deck action.
Choose finitely many
charts covering one
compact normalized
central del Pezzo
component and finitely
many overlaps for
the two deck generators.
Shrink the bidisc once
for this finite data.
On the regular cover
the character map
intertwines the 
deck action with the
reference multipliers
$P_i(q,s)$ in
\eqref{eq:II:next-E-2-20-3}.  Both maps
extend in the charts
just constructed, so
their equality extends
to the central strata.
The reference multipliers
preserve the prescribed
invariant fan.  Deck
translates of the finite
chart collection cover
the entire central
unfolding and use the
same base bidisc.

The chart maps and their displayed inverses intertwine the deck
actions, so they descend to a holomorphic local biholomorphism
$F:\mathcal N\to\mathcal N_{\rm Mumford}$ over the bidisc.
It is proper: for a compact subset of the target, its inverse image
is closed in the compact inverse image under $f$ of the subset's
projection to the bidisc.  Its image is therefore open and closed
in the connected target, and it contains the ordinary locus.
Thus $F$ is a surjective covering.  Its degree is one on the
ordinary locus, so it is a biholomorphism everywhere.
The map carries the chosen section $p$ to the identity section
and induces the given Jacobian identification and integral marking.

\end{enumerate}
\end{proof}

\section{Connected deformation chains}\hypertarget{tag-P276}{}\label{sec:II:connected-deformation-chains}

\begin{proposition}[Global identification of a smooth completion]\hypertarget{tag-0337}{}\label{push:src:E.2.21}
\StatementRecord{push:src:E.2.21}{proved}{0337}{\SeriesRef{push:src:E.2.5}, \SeriesRef{push:src:E.2.7}, \SeriesRef{push:src:E.2.10}, \SeriesRef{push:src:E.2.12}, \SeriesRef{push:src:E.2.14}, \SeriesRef{push:src:E.2.15}, \SeriesRef{push:src:E.2.13}, \SeriesRef{push:src:E.2.20}}

Let $\pi:\mathcal X\to\Delta$ satisfy the hypotheses of
\ref{push:src:E.2.5}, with fixed coprime multiplicities $(3k,4l)$ on
every noncentral fibre.  Then $X_0$ is a member of that
same constructed family.  Its isomorphism parameter is
the interior value $z(0)$ of \ref{push:src:E.2.7}.

After choosing a holomorphic lift $c(s)$ of that parameter
and shrinking the disc, the family is isomorphic
to the prescribed family at $c(s)$.  The ambiguity in
this comparison is the connected translation automorphism
group of \SeriesRef{push:src:E.1.3}.
\end{proposition}


\begin{proof}
\begin{enumerate}[label=\arabic*.,leftmargin=*]
\item 
\ref{push:src:E.2.10} extends the intrinsic pencil to a proper flat map
\[
 g:\mathcal X\longrightarrow\mathbb P^1\times\Delta.
\]
\ref{push:src:E.2.7} supplies its interior nonzero period $z(s)$ and hence
a holomorphic logarithmic lift $c(s)$.  \ref{push:src:E.2.12} excludes
additional critical values, and \ref{push:src:E.2.14}--\ref{push:src:E.2.15} exclude
target collapse.  Normalize the three distinct labelled
sections to $0,1,\infty$ by a holomorphic projective
transformation over the disc.

On the ordinary locus the marked Jacobian has
period matrix $\Pi_{c(s)}(t)$.  More intrinsically, its
holomorphic period kernel in the complexified marked
lattice equals that of this matrix.  This equality holds
for $s\ne0$ by the given family identification and extends
over $s=0$ by the identity theorem.  The local marking
choices agree: an automorphism of a fixed marked member
has identity action on ordinary $H_1$, by \SeriesRef{push:src:E.1.2}.  The
smooth family at any fixed ordinary target point
also has identity homology return around $s=0$.

\ref{push:src:E.2.13} gives smooth normalized root covers at the
two finite values, of degrees $3k$ and $4l$, with their
free cyclic deck actions.  Their marked Jacobians are
the corresponding reference root Jacobians.  \ref{push:src:E.2.20}
identifies the cusp neighbourhood over the full
$(q,s)$-bidisc with the prescribed marked Mumford
neighbourhood, including its period multipliers.
It remains to identify the finite affine actions and
the global translations on the overlaps.

\item \emph{The finite affine classes.}
Choose local holomorphic origins on a finite root torus
family.  Every deck map is affine on the compact torus
fibres.  Its linear part is the fixed marked finite
monodromy.  Subtracting this part gives a holomorphic
translation section.  Lift it to the Lie bundle on
simply connected refinements of the root charts.
The finite-order relation says that the sum of its
successive linear translates is an integral period.
This lattice vector is locally constant in $s$:
the lifted translation and period vectors vary
continuously, whereas the integral coefficients are
discrete.  Changing the lift changes it by precisely
the integral coboundary in \SeriesRef{push:src:E.1.2}.

Consequently the finite affine cocycle has the same
integral class at $s=0$ as at $s\ne0$.  After matching
these classes, the remaining difference is a cocycle
in the complex vector sheaf.  The explicit finite
averaging in \SeriesRef{push:src:E.1.2} removes it, since division by the
finite group order is allowed in that sheaf.
This constructs the finite-root comparisons with
holomorphic dependence on $s$. The reduced finite fibres are the
resulting free quotients.

\item \emph{The integral collar classes.}
Use the two finite-root neighbourhoods, the cusp
neighbourhood, ordinary simply connected charts,
and fixed annular collars to cover the target.
Choose local holomorphic origins on the ordinary
torus charts.  A transition between two such
descriptions is fibrewise affine.  The common
integral marking makes its linear part the
prescribed one, leaving a translation section.

Lift these sections to the Lie bundles on simply
connected refinements of the collars.  Differences
on deck translates and sums on triple overlaps
are integral periods.  Every resulting integral
vector is constant in $s$.  This comparison takes
place on fixed collars away from the singular
centres. The finite and cusp chart
identifications specify exactly which integral
changes of origins are allowed.

For $s\ne0$ these classes are those of the
prescribed logarithmic-transform data, modulo
the integral changes of origins in \SeriesRef{push:src:E.1.2}.
They are therefore the same at $s=0$.
One can normalize them once on the connected
parameter disc.  This includes the cusp
winding class, since \ref{push:src:E.2.20} identifies the
marked period action there.

\item \emph{Solving the holomorphic remainder.}
After this discrete normalization, the remaining
comparison is the additive cocycle in the
extension sheaf $\mathcal V$ of \SeriesRef{push:src:E.1.2}.  Its
finite extension conditions are exactly
those of the already identified root
charts.  Its cusp extension condition
is the toric collar condition
proved in \ref{push:src:E.2.20}.  The calculation in
\SeriesRef{push:src:E.1.2} therefore applies to this cover and gives
\[
 0\longrightarrow\mathcal O_{\mathbb P^1}
       \boldsymbol\delta
 \longrightarrow\mathcal V
 \longrightarrow\mathcal O_{\mathbb P^1}(-\infty)
 \longrightarrow0,
 \qquad H^1(\mathbb P^1,\mathcal V)=0.
\]
The explicit finite averaging and Laurent-series
splitting in \SeriesRef{push:src:E.1.2} kill the additive cocycle.
The resulting translations glue the local
comparisons to a biholomorphism of $X_0$
with the prescribed member at $c(0)$.

The same solution works over a smaller parameter
disc.  On the fixed collars, Cauchy integral
formulas for the Laurent coefficients preserve
holomorphic dependence on $s$, as does finite
averaging.  Thus the solved translations and
their glued comparison are holomorphic in
$s$.  Their ambiguity is the connected
scalar translation group of \SeriesRef{push:src:E.1.3}.

\end{enumerate}
\end{proof}

\begin{proposition}[Propagation from a known open subset of a disc]\hypertarget{tag-0338}{}\label{push:src:E.2.22}
\StatementRecord{push:src:E.2.22}{proved}{0338}{\SeriesRef{push:src:E.2.21}, \SeriesRef{push:src:E.2.5}, \SeriesRef{push:src:E.2.18}, \SeriesRef{push:src:E.2.19}, \SeriesRef{push:src:E.2.20}; \cite[Chapter~II, Theorem~8.8, pp.~118--121]{Demailly2012E}; \cite[Chapter~VI, Theorem~1.1 and Table~10, pp.~243--244]{BHPV2004E}; \cite[Chapter~II, Definition~7.9 and Theorem~7.12, pp.~114--115]{Demailly2012E}}

Let $\pi:\mathcal X\to\Delta$ be a proper holomorphic
submersion with fibres diffeomorphic to $S^6$.  Suppose
there is a nonempty open subset $U\subset\Delta$ such
that every $X_s$, $s\in U$, is a member of the fixed
coprime multiplicity family $(3k,4l)$ of \SeriesRef{push:src:E.2.1}.
Then every fibre over $\Delta$ is a member of that
same family.
\end{proposition}


\begin{proof}
\begin{enumerate}[label=\arabic*.,leftmargin=*]
\item 
Local miniversality in \SeriesRef{push:src:E.2.1} supplies 
holomorphic comparisons with the prescribed family
locally over $U$, with locally constant integral
markings.  Thus the known open set has the relative
root and normal-crossing charts used below.

At each step the exceptional set is a closed discrete analytic
subset of the current parameter domain. Such a set may accumulate
only at points removed from that domain. We apply
\cref{push:src:E.2.21} to these domains in reverse order.

\item \emph{The intrinsic line and the relative pencil.}
Put $A=12kl$ and $D=kl-3k-4l\ne0$.
The argument of \ref{push:src:E.2.5} constructs the unique line
bundle $\mathcal L$ on the whole total space with
\[
 \mathcal L^D=K_{\mathcal X/\Delta}^{A}.
\]
Indeed Ehresmann transport over the contractible disc
gives $H^1(\mathcal X,\Z)=H^2(\mathcal X,\Z)=0$,
so the exponential sequence identifies its Picard
group with a complex vector space.

Proper coherence makes $\mathcal V=\pi_*\mathcal L$
a torsion-free, hence locally free, sheaf on the
disc.  Its rank is two because it has that rank
on $U$.  At every parameter its fibre injects
into $H^0(X_s,\mathcal L_s)$ by the multiplication
exact sequence used in \ref{push:src:E.2.5}.  Choose a frame
of $\mathcal V$ on the disc.  The common zero
locus of its two evaluated sections is closed
analytic in $\mathcal X$.  Its image is closed
analytic by properness and Remmert's theorem
\cite[Chapter~II, Theorem~8.8, pp.~118--121]{Demailly2012E}.
It misses $U$, so is a discrete analytic subset.
Remove it temporarily.

On the remaining connected parameter curve $S$
the sections define a proper holomorphic map
\begin{equation}
 g:\mathcal X_S\longrightarrow\mathbb P^1\times S,
 \qquad \mathcal L=g^*\mathcal O(1).
\label{eq:II:next-E-2-22-1}\end{equation}
Each $g_s$ is nonconstant because its two sections
are independent.  Its fibres have pure dimension
two; the smooth total space and smooth target
then give flatness.  The fibres are connected:
the finite Stein factor is generically of degree
one, as it is so above $U$, and a finite
birational map to the normal target is an
isomorphism.  Its source is irreducible, so
there is no additional vertical component
in the Stein factor.  We retain
\begin{equation}
 K_{\mathcal X_S/S}^{A}=g^*\mathcal O(D).
\label{eq:II:next-E-2-22-2}\end{equation}

\item \emph{Critical values and ordinary tori.}
The relative critical locus of $g$ is analytic
and proper over the target.  Its image has
finite intersection with each target fibre:
the nonconstant proper map $g_s$ to a curve
has finitely many critical values.  Thus the
reduced critical image is finite over $S$.
Its one-dimensional part has degree three,
as can be read over $U$; every such component
dominates $S$.  Remove the images of its
zero-dimensional components and its branch
and singular points.  These are further
proper analytic discrete subsets, absent
over the known open set.

A finite \'etale parameter cover now labels
the three critical sections.  Pointwise
classification after this cover descends;
images of later discrete exceptional sets
remain discrete because this cover is finite.
One may work on each connected component,
each of which meets the inverse image of $U$.
Send the labelled sections to $0,1,\infty$.
There are exactly these three critical values
on the current parameter domain.

The ordinary family over
$B^\circ\times S$, where
$B^\circ=\mathbb P^1\setminus\{0,1,\infty\}$,
is smooth and proper.  Its compact surface
fibres have the topology of the known two-tori
by Ehresmann's theorem and connectedness.
Adjunction gives integral $c_1(K_F)=0$,
since $F$ is a smooth fibre in a complex
six-sphere and its normal line is trivial.
The root identity restricts to $K_F^{\otimes A}\simeq\mathcal O_F$,
so $\kappa(F)=0$. The vanishing of its canonical Chern class
excludes a $(-1)$-curve, and hence $F$ is minimal.
In \cite[Chapter~VI, Theorem~1.1 and Table~10, pp.~243--244]{BHPV2004E},
the minimal surface of Kodaira dimension zero with $b_1(F)=4$
is a complex two-torus.

\item \emph{The finite root charts.}
For any labelled critical section $\Gamma_j$,
every component of the Cartier divisor $g^*\Gamma_j$
dominates $S$.  A vertical divisor component would
be an entire fibre $X_s$, making $g_s$ constant.
Every dominating component meets $U$, where the
reduced fibre is irreducible.  Thus there is
exactly one such component, and its multiplicity
is the known one.  In particular
\begin{equation}
 g^*\Gamma_1=(3k)\mathcal D_1,
 \qquad g^*\Gamma_2=(4l)\mathcal D_2,
 \qquad g^*\Gamma_0=\mathcal W,
\label{eq:II:next-E-2-22-3}\end{equation}
with reduced total hypersurfaces on the right.
For clarity, two distinct dominating components
could not both restrict to the known irreducible
reduced fibre.  Each is a Cartier prime in
the smooth total space, is torsion-free and
hence flat over the smooth parameter curve,
and has pure two-dimensional fibres.  Thus
its fibre over any parameter in $U$ contains
the whole known irreducible reduced surface.
At a general smooth point of
that fibre the known local equation has only
one reduced total branch.  Properness forces
every dominating component to meet that open
part of the known fibre, a contradiction.

For a finite multiplicity $m$, write locally
$q=a^m u$, with $a$ a reduced local equation
of $\mathcal D_j$ and $u$ a unit.  The normalization
of the cyclic root base change $r^m=q$ is
locally the disjoint union of the graphs
\[
 r=\zeta a u^{1/m},\qquad \zeta^m=1.
\]
It is therefore smooth and \'etale over the
total space, with free cyclic deck
action.  On $r\ne0$ its map to the root
disc times $S$ is already smooth.  Its
critical locus is supported over $r=0$,
is proper over $S$, and is empty over $U$.
Delete its discrete analytic image for
each of the two roots.  The resulting
proper smooth root families have connected
fibres by their Stein factors and are
families of two-tori by the same topology,
adjunction and surface argument above.

\item \emph{A finite analytic test for relative normal crossings.}
The total cusp $\mathcal W$ in \eqref{eq:II:next-E-2-22-3}
is a reduced embedded hypersurface in
the smooth fourfold.  Take its finite
analytic normalization
\[
 \nu:\widetilde{\mathcal W}\longrightarrow
       \mathcal W\longrightarrow\mathcal X_S.
\]
Existence and finiteness are the analytic
normalization theorem
\cite[Chapter~II, Definition~7.9 and Theorem~7.12, pp.~114--115]{Demailly2012E}.
The normalization is proper over $S$.
First remove the proper images of its
singular locus, and then of the critical
locus of its map to $S$.  Each is analytic
and misses $U$.  On the resulting domain
the normalization is smooth over $S$.
Next remove the proper image of the
rank-drop locus of
\[
 T_{\widetilde{\mathcal W}/S}
       \longrightarrow\nu^*T_{\mathcal X_S/S}.
\]
This locus too is analytic and absent over
$U$.  The finite map $\nu$ is now an immersion.

For $r=2,3,4$, let $R_r$ be the locus
of ordered, pairwise distinct preimages
in the $r$-fold fibre product of
$\widetilde{\mathcal W}$ over $\mathcal X_S$.
Its partial diagonals are both open and
closed: a finite immersion is locally
an embedding, so its diagonal is open,
and separatedness makes it closed.
Consequently $R_r$ is closed analytic
in the proper fibre product and is
proper over $S$.

At a point of $R_r$ the $r$ local sheets
have relative conormal lines in the same
rank-three bundle $T^*_{\mathcal X_S/S}$.
On $R_2$ and $R_3$, delete the proper
images of the rank-drop loci of these
two or three lines.  Delete also the
image of $R_4$.  These are closed
analytic images and miss $U$, hence
are discrete.  There are only finitely
many such deletions.

At a remaining point the total cusp
has at most three local sheets, each
embedded smoothly and submersive over
$S$, with independent vertical normal
covectors.  Since the \emph{total}
$\mathcal W$ is a reduced embedded
hypersurface, its local ideal is exactly
the product of the distinct local
equations of those sheets.  Independence of the
vertical normals makes these equations,
the parameter on $S$, and the remaining
coordinates a holomorphic coordinate
system.  Thus
\begin{equation}
 q=x,\qquad q=xy,\qquad\hbox{or}\qquad q=xyz
\label{eq:II:next-E-2-22-4}\end{equation}
after absorbing a unit.  These are
relative normal-crossing charts;
in particular the remaining cusp fibres
are reduced.

\item \emph{The normalized toric pair and its action.}
The smooth proper normalization and
its conductor strata have the same
topology, incidence and normal degrees
as over $U$.  A further finite \'etale
cover can label these finitely many
strata.  The argument of \ref{push:src:E.2.18} applies
to every fibre: the six conductor
curves are a hexagon of $(-1)$-curves,
three alternate contractions give a
triangle in $\mathbb P^2$, and the
normalized pair is $(\mathrm{dP}_6,D)$.
Its relative logarithmic tangent bundle
is the evaluation pullback of a rank-two
bundle on the parameter curve.

In this evaluation description the
conductor descent maps are holomorphic
endomorphisms of that rank-two bundle.
They are constant along the compact
conductor curves.  They equal the identity
over $U$, where both toric logarithmic
fields descend.  The holomorphic identity
theorem on the connected parameter curve
makes them the identity everywhere.
This applies the descent proof of \ref{push:src:E.2.19}
to identities holding over the nonempty open set $U$.
Grauert evaluation, commutation on the
ordinary tori, completeness and the
implicit extension of the two primitive
toric periods then give the 
marked cusp torus action locally on
the current parameter curve.  The
integral monodromy and triple-point
characters are constant by transport
and are the prescribed ones.

\item \emph{Continuation of the full period kernel.}
Fix an ordinary target point $t_*$.
The integral local system over
$B^\circ\times S$ has constant marked
peripheral representation, up to a
choice of marking: its matrices are
transported from those over $U$.
Parameter monodromy commutes with the
base representation.  It therefore
preserves the cusp invariant lattice
\[
 \ker(A_0-I)=\Z\widehat w\oplus\Z\widehat\delta.
\]
The wedge of the two corresponding
holomorphic period vectors at $t_*$
has a well-defined zero locus on $S$:
changing its integral basis multiplies
the wedge by a nonzero determinant.
It is nonzero over $U$, so that zero
locus is discrete analytic.  Remove
it as one final exceptional layer.

On a connected universal cover of
the remaining parameter curve, normalize
the Lie frame at $t_*$ by these two
columns.  The $\widehat\gamma$ column
then defines a holomorphic function
$c(s)$ by subtracting the fixed
reference value of $\beta^{\rm part}(t_*)$.
All other column identities at $t_*$
hold over a lift of $U$ and hence
everywhere by the identity theorem.

Now work on the universal cover of
the whole ordinary base times this
parameter cover.  The holomorphic
period kernel is a rank-two subbundle
of the complexified marked rank-four
lattice.  Compare it with
\begin{equation}
 \ker\Pi_{c(s)}(t).
\label{eq:II:next-E-2-22-5}\end{equation}
The latter matrix always has complex
rank two, since its $\widehat w$ and
$\widehat\delta$ columns are the
standard basis.  These two holomorphic
Grassmannian maps agree over the
product with a lift of $U$.  They
therefore agree on the entire
connected ordinary cover.
In a local Lie frame the
identity says that the full period
matrix is $W(t,s)\Pi_{c(s)}(t)$,
where $W$ is its two marked columns.
If $W$ had rank at most one, every
period column would have that rank,
contradicting the real rank four
of the compact torus lattice.

The ordinary and finite-root
fibres are compact tori.  The exact
period determinant is nonzero there.
Its sign is the negative sign on
$U$ and stays negative by connectedness
and nonvanishing.  Equality of the
period kernels also extends to the
smooth finite root charts; their
nonzero determinant preserves this
sign at the root point where the
maximum in \SeriesRef{fam:src:C.2.3} is attained.  Thus
$\operatorname{Im}c(s)+M<0$ by
the precise threshold of \SeriesRef{fam:src:C.2.3}.
This places $c(s)$ in the
admissible interior domain, including
strictness at its boundary.  The
cusp multipliers are the corresponding
marked reference multipliers.  The
proof of \ref{push:src:E.2.20} now reconstructs
the cusp at each parameter:
relative normal-crossing transport
gives its known cover and incidence,
the action gives the same characters,
and adjacent-origin comparison gives
the monomial times unit extensions.
This construction works along any path from $U$ and locally in the
parameter.

The finite affine lattice classes
and integral collar classes are
locally constant on this connected
domain, hence remain those over
$U$.  The additive comparison in
\SeriesRef{push:src:E.1.2} and \ref{push:src:E.2.21} identifies every
remaining threefold with
the reference member at $c(s)$.
The pointwise conclusion descends
from every finite parameter cover
used to label the critical or
conductor strata.

\item \emph{Filling the exceptional layers.}
We have made only finitely many
successive deletions.  At each stage
the deleted subset was closed
analytic and discrete on the
parameter domain then present.
Finite covers do not change that
property after taking their proper
images.  On the final domain all
fibres are known members of the
fixed pair.  At a point of the
last deleted set, choose a small
disc containing no other point of
that set and avoiding all earlier
deleted sets.  Its punctured disc
consists of known members.  \ref{push:src:E.2.21}
identifies the central fibre as
well.  Fill the entire last layer
in this way, then the preceding
one, and continue in reverse order.
This is valid even when a later
layer accumulates at an earlier
deleted point.  Finally every
fibre over the original disc is
a member of the fixed constructed family.

\end{enumerate}
\end{proof}

\Needspace{10\baselineskip}
\begin{theorem}[Separation under arbitrary deformation chains]\hypertarget{tag-0339}{}\label{push:src:E.2.23}
\StatementRecord{push:src:E.2.23}{proved}{0339}{}

In a proper smooth holomorphic family of compact complex
manifolds over a connected reduced complex analytic base,
if one fibre is a sphere member of \SeriesRef{push:src:E.2.1} with coprime
multiplicities $(3k,4l)$, then every fibre is a constructed
member with that same pair.

Consequently two of the constructed sphere members are
deformation equivalent, allowing chains of proper smooth
families over connected complex analytic bases, if and
only if their ordered multiplicity pairs agree.  In
particular the distinct coprime pairs give distinct
deformation classes.
\end{theorem}


\begin{proof}

\begin{enumerate}[label=\arabic*.,leftmargin=*]
\item Fix the pair and let $U$ be the locus of its constructed
members in the base $B$.  It is open by the local
miniversality of \SeriesRef{push:src:E.2.1} and \SeriesRef{push:src:E.2.2}.  We prove that it is
closed.

\item First consider any proper smooth family over a connected
Riemann surface.  If its known-pair locus is nonempty,
Ehresmann's theorem makes all its fibres diffeomorphic
to $S^6$, and that locus is open.  At any point of its closure choose a
coordinate disc meeting that locus.  The intersection
contains a nonempty open subset of the disc, so \ref{push:src:E.2.22}
makes the entire disc known.  Thus the known-pair locus
is closed and equals the connected Riemann surface.

\item Let $b\in\overline U\subset B$.  If $b\in U$ there
is nothing to prove.  Otherwise choose a small local
representative and an irreducible branch $T$ through
$b$ which meets $U$ arbitrarily near $b$.  There are
only finitely many such local branches.  Its dimension
is positive.  The local parametrization theorem gives
a finite proper projection
\begin{equation}
 \rho:T\longrightarrow\Delta^d,
 \qquad \rho^{-1}(0)=\{b\},
\label{eq:II:next-E-2-23-1}\end{equation}
after shrinking the representative and translating
$b$ to the origin.  This is the finite projection
of \citet[Chapter~II, Theorem~4.19 and proof, pp.~95--97]{Demailly2012E}.
The isolated central fibre is part of choosing
the generic projection for the germ; the
Weierstrass representative is then taken small
enough to contain no other point over zero.

\item Choose $a\in U\cap T$ in this representative,
with $\rho(a)\ne0$.  Let $L$ be the complex
line through $0$ and $\rho(a)$, restricted
to the polydisc.  An irreducible component
$C$ of the reduced analytic set
$\rho^{-1}(L)$ which contains $a$ has
dimension one.  Indeed $d-1$ linear
equations give local dimension at least
one, while finiteness of $\rho$ gives
dimension at most one.  Its finite
proper image in $L$ is one-dimensional,
so is all of $L$.  Therefore $C$ contains
the unique point $b$ over zero.

\item Normalize $C$.  Its normalization
$\widetilde C$ is a connected Riemann
surface, finite over $C$; analytic
normalization exists and a normal
curve is smooth by
\cite[Chapter~II, Theorem~7.6 and Corollary~7.7, p.~113; Theorem~7.12, pp.~114--115]{Demailly2012E}.
It has points above both $a$ and $b$.
Pull back the given proper smooth
family to $\widetilde C$.  Its
known-pair locus contains the
inverse image of the open set
$U\cap C$ near $a$.  The preceding
Riemann-surface argument makes
every fibre on $\widetilde C$
known, in particular the fibre
above $b$.  Thus $b\in U$.

\item We have proved that $U$ is open
and closed.  Since $B$ is
connected and $U$ is nonempty,
$U=B$.  A nonreduced base
does not add fibre isomorphism
classes: restriction to its
reduction preserves both the
fibres and smoothness of the
family.  Thus the same conclusion
also applies to such a base
when it is admitted in the
definition of deformation
equivalence.

\item Applying this result to each
link in a finite deformation
chain preserves the pair.
Conversely, \SeriesRef{push:src:E.1.2} and \SeriesRef{push:src:E.1.3} identify
all admissible sphere data
with a fixed pair with the
connected period family,
using its holomorphic
comparisons and its period
lattice $2kl\Z$.
The admissible half-plane
and its punctured-disc
quotient are connected.
Thus all members with
that pair are deformation
equivalent.  This proves
both directions.

\end{enumerate}
\end{proof}

\begin{corollary}[Separation for continuous complex structures]\hypertarget{tag-033A}{}\label{push:src:E.2.24}
\StatementRecord{push:src:E.2.24}{proved}{033A}{\SeriesRef{push:src:E.2.23}; \citet[Section~3.2, Theorem~3.4 and Remark~3.6,
pp.~898--899]{Meersseman2019E}; \cite[Section~2.1, pp.~883--884]{Meersseman2019E}}

Let $\mathcal I(S^6)$ be the space of smooth integrable
almost-complex structures on the fixed smooth sphere,
with its $C^\infty$ topology.  For a fixed coprime pair
$(3k,4l)$, let $\mathcal I_{k,l}$ consist of the structures
whose complex manifolds are biholomorphic to constructed
members with that pair.  Then $\mathcal I_{k,l}$ is open
and closed in $\mathcal I(S^6)$.

Consequently a continuous family of such structures
over any connected topological parameter space cannot
contain constructed members with two different pairs.
The same separation holds in the quotient spaces by
$\operatorname{Diff}_0(S^6)$ or by all orientation-preserving
diffeomorphisms, with their quotient topologies.
\end{corollary}


\begin{proof}

\begin{enumerate}[label=\arabic*.,leftmargin=*]
\item We use the classical Kuranishi local-completeness
theorem at an arbitrary $J_0\in\mathcal I(S^6)$.
The complex orientation is locally constant, so the
oriented formulation applies on each such neighbourhood.
The precise $C^\infty$ version used here is
\citet[Section~3.2, Theorem~3.4 and Remark~3.6,
pp.~898--899]{Meersseman2019E}.  It supplies a connected
open neighbourhood $\mathcal U$ of $J_0$, a
finite-dimensional Kuranishi analytic space $K$
containing $J_0$, and a continuous retraction
$\Xi:\mathcal U\to K$ such that every $J\in\mathcal U$
is obtained from $\Xi(J)$ by a diffeomorphism isotopic
to the identity.  In that theorem the local product
map is first constructed in Sobolev completions;
elliptic regularity gives precisely this restriction
to smooth operators.  The topology on these operators
is the $C^\infty$ topology of
\cite[Section~2.1, pp.~883--884]{Meersseman2019E}.

\item The Kuranishi space carries its proper smooth
holomorphic deformation family.  We can take $K$
connected: in the stated local theorem it is the
continuous retract of the connected $\mathcal U$.
If necessary replace $K$ by its reduction, which
changes neither its points nor the complex manifolds
in its fibres.  Suppose $\mathcal U$ contains some
$J\in\mathcal I_{k,l}$.  Its Kuranishi representative
$\Xi(J)$ has the same complex manifold up to
biholomorphism.  \ref{push:src:E.2.23} applied to the connected
Kuranishi family therefore makes every point of
$K$ a member of the same pair.  Since all points
of $\mathcal U$ are biholomorphic to their
Kuranishi representatives, it follows that
$\mathcal U\subset\mathcal I_{k,l}$.

\item Thus at every $J_0$ the neighbourhood $\mathcal U$
is either contained in $\mathcal I_{k,l}$ or
disjoint from it.  This proves both openness and
closedness.  The inverse image of this set under
a continuous map from a connected parameter space
is then empty or the whole space, proving the
continuous-family assertion.  Finally
$\mathcal I_{k,l}$ is saturated under the indicated
diffeomorphism groups, because its definition uses
only the biholomorphism class.  Its image in either
quotient is open and closed by the definition of
the quotient topology.

\end{enumerate}
\end{proof}

\begin{corollary}[Arbitrary connected deformation chains]\hypertarget{tag-033C}{}\label{push:src:E.2.4}
\StatementRecord{push:src:E.2.4}{proved}{033C}{\SeriesRef{push:src:E.2.5}, \SeriesRef{push:src:E.2.24}, \SeriesRef{push:src:E.2.21}, \SeriesRef{push:src:E.2.22}, \SeriesRef{push:src:E.2.23}}

The ordered coprime pair $(3k,4l)$ is constant in every
proper smooth holomorphic family over a connected complex
analytic base which contains a constructed sphere member.
Two constructed members are deformation equivalent through
such families, allowing arbitrary intermediate members and
finite chains, if and only if their pairs agree.
Moreover different pairs cannot occur in one continuous
family in the $C^\infty$ space of integrable complex
structures on $S^6$.
\end{corollary}


\begin{proof}
For a connected analytic base $B$, pull back the family along
$B_{\mathrm{red}}\hookrightarrow B$. This preserves its complex
fibres and underlying parameter space, and preserves properness and
smoothness by base change. Apply \cref{push:src:E.2.23} to this
family. Its finite-chain conclusion and connected period-family
converse give the asserted equivalence.

By \cref{push:src:E.2.24}, the fixed-pair locus is open and closed
in the $C^\infty$ space of integrable complex structures. Its inverse
image under a continuous family on a connected parameter space is
empty or the whole parameter space, proving the final clause.
\end{proof}

\section{Two independent periods}\hypertarget{tag-P277}{}\label{sec:II:two-independent-periods}

\begin{proposition}[Monodromy, centralizer, and integral lattices]\hypertarget{tag-033O}{}\label{push:src:E.5.2}
\StatementRecord{push:src:E.5.2}{proved}{033O}{\cite{Viaclovsky2026Two}}

Use the source's left-action marking
$H^1(F,\mathbb Z)=\mathbb Z\langle\psi,u,w,\delta\rangle$
and the ordered meridians $*,1,2,3$ with product one.
The cohomological matrices are
\[
\begin{aligned}
 T_*&=\begin{pmatrix}1&0&-2&-1\\0&0&-1&0\\0&1&0&-1\\0&0&0&1\end{pmatrix},&
 T_1&=\begin{pmatrix}1&0&2&2\\0&1&1&1\\0&0&1&0\\0&0&0&1\end{pmatrix},\\
 T_2&=\begin{pmatrix}1&0&0&0\\0&1&0&0\\0&-1&1&0\\0&0&0&1\end{pmatrix},&
 T_3&=\begin{pmatrix}1&0&0&-1\\0&1&1&0\\0&0&1&0\\0&0&0&1\end{pmatrix}.
\end{aligned}
\]
On ordinary dual homology put $H_i=(T_i^{-1})^t$,
$\ell=(1,-1,-1,0)^t$ and $\widehat\delta=(0,0,0,1)^t$.
Then
\[
 \operatorname{Cent}_{M_4(\mathbb Q)}(H_*,H_1,H_2,H_3)
 =\mathbb Q I\oplus\mathbb Q E_{41},
\]
and the cohomological centralizer is
$\mathbb Q I\oplus\mathbb Q E_{14}$.
The fixed and norm lattices at the multiple fibre are
\[
 \begin{aligned}
 \Lambda^{H_*}&=\mathbb Z\ell\oplus\mathbb Z\widehat\delta,\\
 N_4\Lambda&=4\mathbb Z\ell\oplus2\mathbb Z\widehat\delta,\\
 \Lambda^{H_*}/N_4\Lambda&\simeq\mathbb Z/4\oplus\mathbb Z/2.
 \end{aligned}
\]
The primitive vanishing lattices at the three reduced fibres are
\[
 K_1=\mathbb Z(\widehat w+\widehat\delta),\qquad
 K_2=\mathbb Z\widehat u,\qquad
 K_3=\mathbb Z\widehat w\oplus\mathbb Z\widehat\delta.
\]
Their invariant lattices are
\[
 \begin{split}
 \Lambda^{H_1}&=\mathbb Z(\widehat\psi-2\widehat u)
       \oplus\mathbb Z\widehat w\oplus\mathbb Z\widehat\delta,\\
 \Lambda^{H_2}&=\mathbb Z\widehat\psi\oplus
       \mathbb Z\widehat u\oplus\mathbb Z\widehat\delta,\\
 \Lambda^{H_3}&=K_3.
 \end{split}
\]
\end{proposition}


\begin{proof}
The matrices are those of \citet[Proposition~4.2, pp.~11--12]{Viaclovsky2026Two}, whose
derivation uses the elliptic monodromies, the winding of $P$,
and the primitive rank-two cusp.  In particular it determines
the winding integer by the saturation index $d^2=1$.  Direct multiplication gives
$T_*T_1T_2T_3=I$, $T_*^4=I$ and
$(T_i-I)^2=0$ for $i=1,2,3$.
The dual homology matrices are
\[
\begin{aligned}
 H_*&=\begin{pmatrix}1&0&0&0\\-2&0&-1&0\\0&1&0&0\\1&1&0&1\end{pmatrix},&
 H_1&=\begin{pmatrix}1&0&0&0\\0&1&0&0\\-2&-1&1&0\\-2&-1&0&1\end{pmatrix},\\
 H_2&=\begin{pmatrix}1&0&0&0\\0&1&1&0\\0&0&1&0\\0&0&0&1\end{pmatrix},&
 H_3&=\begin{pmatrix}1&0&0&0\\0&1&0&0\\0&-1&1&0\\1&0&0&1\end{pmatrix}.
\end{aligned}
\]
Solving the sixty-four rational equations $CH_i=H_iC$
gives $C=aI+bE_{41}$ with free $a,b$. Transposition gives the
linear bijection between the two centralizers: the equation
$CH_i=H_iC$ is equivalent to $C^tT_i=T_iC^t$, and
$E_{41}^t=E_{14}$.  Solving $(H_*-I)v=0$ gives
$v=(a,-a,-a,d)^t$ integrally, and
\[
 N_4=I+H_*+H_*^2+H_*^3
 =\begin{pmatrix}4&0&0&0\\-4&0&0&0\\-4&0&0&0\\0&2&-2&4\end{pmatrix}.
\]
In the basis $(\ell,\widehat\delta)$ its columns are
$(4,0),(0,2),(0,-2),(0,4)$, giving exactly the displayed
norm group.  Each $K_i$ is the integral image of
$H_i-I$; its Smith form has only unit nonzero entries.
The kernels of these three matrices give the remaining fixed
lattices.

\end{proof}

\begin{cited}[Singular fibres of the two-parameter threefold]\hypertarget{tag-033N}{}\label{push:src:E.5.1}
\StatementRecord{push:src:E.5.1}{proved}{033N}{\SeriesRef{push:src:E.5.2}; \cite[Lemma~2.1, pp.~3--4; Proposition~3.1, pp.~6--7; Theorem~1.1, p.~2]{Viaclovsky2026Two}}
In \cite[Lemma~2.1, pp.~3--4]{Viaclovsky2026Two}, the elliptic surface
$y^2=x^3+tx+1$ has fibre $III^*$ at infinity. The associated
threefold has two rank-one fibres at $p_1,p_2$, a rank-two
hexagonal fibre at $p_3$, and the multiple bielliptic fibre
$4B_*$ at infinity
\cite[Proposition~3.1, pp.~6--7; Theorem~1.1, p.~2]{Viaclovsky2026Two}.
The labels are those of \cref{push:src:E.5.2}; the ranks of
$T_1-I,T_2-I,T_3-I$ in that marking are $1,1,2$.
\end{cited}

\begin{definition}[Two-parameter period data]\hypertarget{tag-P208}{}\label{def:II:two-parameter-period-data}
\StatementRecord{def:II:two-parameter-period-data}{proved}{P208}{}
Let $\mu$ range over a sufficiently small disc about zero, and set
\[
 S_\mu:\ y^2=x^3+tx+\mu t+1,\qquad
 P_\mu=(-\mu,\sqrt{1-\mu^3}),
\]
with the square root equal to one at zero.  Label the three roots
$p_i(\mu)$ of $4t^3+27(\mu t+1)^2$ holomorphically near zero;
$p_3$ is the zero of the chosen linearization.
Choose a holomorphic generator $\sigma_\mu$ of its one-dimensional
space of sections vanishing at $p_3$ and scale it by $\lambda$.
For the added homology generator take the inverse of the
contracting lift.  On the universal cover of the four-punctured
base its period columns have the form
\[
 \Pi_{\mu,c}=
 \begin{pmatrix}2z_\mu&\tau_\mu&1&0\\
 \beta_\mu+c&z_\mu&0&1\end{pmatrix},
 \qquad \lambda=e^{-2\pi i c},\qquad
 \partial_c\Pi=e\otimes\psi.
\]
Here $z_\mu$ is a lifted elliptic logarithm in the chosen marking,
$e=(0,1)^t$, and $\psi$ is the first covector in the marking of
Proposition~\ref{push:src:E.5.2}.
\end{definition}

\begin{theorem}[Biholomorphism classes in the two-parameter family]\hypertarget{tag-033P}{}\label{push:src:E.5.3}
\StatementRecord{push:src:E.5.3}{proved}{033P}{\SeriesRef{push:src:E.5.1}, \SeriesRef{push:src:E.5.2}; \cite{Viaclovsky2026Two}}

There are $\epsilon>0$ and a real constant $C$ such that the
construction of \cite[Proposition~6.4, pp.~19--20]{Viaclovsky2026Two},
with its fixed logarithmic vector $\ell$,
defines a proper holomorphic family
$X_{\mu,c}$ over
$\{ |\mu|<\epsilon\}\times\{\operatorname{Im}c<C\}$.
It has the singular fibres of \ref{push:src:E.5.1} and the markings of
Definition~\ref{def:II:two-parameter-period-data}.
Two of its members are biholomorphic if and only if
\[
 \mu'=\mu,\qquad c'-c\in2\mathbb Z.
\]
Every such biholomorphism covers the identity of the base $\mathbb P^1$.  The quotient parameter in the second variable is
$v=e^{-\pi i c}$, with $\lambda=v^2$.
On a sufficiently small bidisc centred at any
$(0,\lambda_0)$ in \cite[Theorem~1.1, p.~2]{Viaclovsky2026Two}, choose
one logarithm of $\lambda$.  Within that bidisc,
\[
 X_{\mu,\lambda}\simeq X_{\mu',\lambda'}
 \quad\Longleftrightarrow\quad
 (\mu,\lambda)=(\mu',\lambda').
\]
\end{theorem}


\begin{proof}

\begin{enumerate}[label=\arabic*.,leftmargin=*]
\item The equations of \cite[Proposition~6.4, pp.~19--20]{Viaclovsky2026Two}
give the simultaneous elliptic family, line bundle, linearization,
and good-reduction comparison.  On a smaller closed $\mu$-disc
all unit factors in its four local fillings are uniformly bounded
above and below.  Its contraction estimates therefore have a
common nonzero bound $|\lambda|<\eta$. This bound corresponds to
the indicated half-plane under $\lambda=e^{-2\pi i c}$.
At any point of that half-plane, choose a parameter neighbourhood
$U$ on which $(\mu,c)\mapsto(\mu,e^{-2\pi i c})$ is biholomorphic
onto a source bidisc. Its inverse is
$(\mu,\lambda)\mapsto(\mu,-\log\lambda/(2\pi i))$, with the
branch taking the prescribed value of $c$. Transport the integral
marking to this bidisc and use the fixed vector $\ell$ in its
degree-four filling. The integral markings and the four-torsion
section $\ell/4$ continue uniquely on the simply connected
half-plane parameter domain.
On every relatively compact subset the source's finite root
quotients and the fixed rank-one and hexagonal fans give
holomorphic charts with uniform bounds.  Their comparisons are
the specified section-preserving ones, so they agree on
intersections and define a family.  Properness is local on
the parameter domain and follows from the compact four-piece
construction there; smoothness follows from the same free
finite action and unimodular charts.

The independent-variation calculation in
\cite[Proposition~6.4, pp.~19--20]{Viaclovsky2026Two}
gives dimension at least two.
Replace $\sigma_\mu$ by $a(\mu)\sigma_\mu$, where $a$ is a
holomorphic unit.  On the parameter disc choose a holomorphic
logarithm and put $\theta(\mu)=\log a(\mu)/(2\pi i)$.  The maps
\[
 \widetilde c=c+\theta(\mu),\qquad
 \widetilde\beta_\mu=\beta_\mu-\theta(\mu),\qquad
 \widetilde\lambda=\lambda/a(\mu)
\]
have inverses obtained by subtracting $\theta$, adding $\theta$,
and multiplying by $a$, respectively.  They preserve the period
matrix and the linearization
$\lambda\sigma_\mu=\widetilde\lambda(a(\mu)\sigma_\mu)$.
Thus they identify the two presentations over the displayed
parameter-coordinate map, and differences of $c$ at fixed $\mu$
are unchanged.

For the inverse contracting-lift convention, suppress $\mu$ and
write the continued period triples as
\[
 \begin{aligned}
 (\tau_1,z_1,\beta_1)&=(\tau+1,z+1,\beta+2),\\
 (\tau_2,z_2,\beta_2)&=
 \left(\frac{\tau}{1-\tau},\frac{z}{1-\tau},
       \beta+\frac{2z^2}{1-\tau}\right),\\
 (\tau_3,z_3,\beta_3)&=(\tau+1,z,\beta-1),\\
 (\tau_*,z_*,\beta_*)&=
 \left(-\frac1\tau,\frac{z-1}{\tau},
       \beta-1-\frac{2(z-1)^2}{\tau}\right).
 \end{aligned}
\]
Let $\Pi_i$ be the period matrix with the corresponding continued
triple and the same $c$.  Multiplication by the matrices of
Proposition~\ref{push:src:E.5.2} gives
\[
 \Pi_iH_i=G_i\Pi_{\mu,c},\qquad
 G_1=G_3=I_2,\quad
 G_2=\begin{pmatrix}(1-\tau)^{-1}&0\\z(1-\tau)^{-1}&1\end{pmatrix},
 \quad
 G_*=\begin{pmatrix}\tau^{-1}&0\\(1-z)\tau^{-1}&1\end{pmatrix}.
\]
In particular, the second coordinate of the added period jumps
by $-1$ at $p_3$, and
$H_3\widehat\psi=\widehat\psi+\widehat\delta$ cancels that jump.
They express covariance in the chosen inverse-lift convention.

\item On each parameter neighbourhood $U$ constructed above, the
family of \cite[Proposition~6.4, pp.~19--20]{Viaclovsky2026Two}
uses the same smooth-core charts, free degree-four root quotient,
and rank-one and hexagonal fillings as the construction in
\cite[Proposition~6.1, pp.~17--18]{Viaclovsky2026Two}.
Pull back those charts and their line-bundle transition maps along
the local parameter isomorphism
$(\mu,c)\mapsto(\mu,e^{-2\pi i c})$ on $U$, with the transported
marking and the fixed vector $\ell$. Its canonical isomorphism
therefore gives $K_X\simeq f^*\mathcal O_{\mathbb P^1}(-1)$
for the present member. Since $f_*\mathcal O_X=\mathcal O_{\mathbb P^1}$,
projection identifies
$H^0(X,-K_X)=f^*H^0(\mathbb P^1,\mathcal O(1))$.
Thus its base-point-free anticanonical pencil is $f$.
A biholomorphism therefore induces an affine automorphism of
the base $\mathbb P^1$, fixing its distinguished multiple point
at infinity and its distinguished hexagonal point $p_3$.
It may initially interchange the two rank-one points.

\item Put $Q(\mu)=(p_1-p_3)/(p_2-p_3)$.
The source computes $p_i'(0)=-9/(2p_i(0))$.
If $Q'(0)$ vanished, the three velocities would have the form
$Ap_i(0)+B$; the polynomial $Az^2+Bz+9/2$ would vanish
at three distinct numbers, a contradiction.
Also $Q(0)^2-Q(0)+1=0$, so $Q(0)\ne Q(0)^{-1}$.
Shrink $\epsilon$ so that $Q$ is injective and its image
is disjoint from the image of $Q^{-1}$.  Preservation or
interchange of the rank-one values gives respectively
$Q(\mu')=Q(\mu)$ or $Q(\mu')=Q(\mu)^{-1}$.
The second possibility is excluded, and the first gives
$\mu'=\mu$.  The base automorphism now fixes three distinct
points individually and is the identity.

\item On regular fibres a biholomorphism is affine, and its integral
linear map commutes with all four $H_i$.  \ref{push:src:E.5.2} gives
$aI+bE_{41}$.  Comparing the constant third and fourth
period columns makes its complex linear part $aI$;
unimodularity gives $a=\alpha\in\{1,-1\}$.
The first column gives $h=c'-c\in\mathbb Z$ and the
lattice map $A=\alpha(I-hE_{41})$, since
$\Pi_{\mu,c+h}(I-hE_{41})=\Pi_{\mu,c}$.
After the fourth-root base change, the oriented meridian
$\gamma$ satisfies $\gamma^4=\ell$, as in
\cite[equation~(3.14), p.~8]{Viaclovsky2026Two}. A map over the
identity of the base sends it to $k\gamma$ for $k\in\Lambda$.
Thus $A\ell-\ell=N_4k$. Since $A$ commutes with $N_4$,
multiplication by the lattice automorphism $-\alpha A^{-1}$ gives
\[
 (I+hE_{41})\ell-\alpha\ell
 =(1-\alpha)\ell+h\widehat\delta\in N_4\Lambda.
\]
The norm lattice in \ref{push:src:E.5.2} forces $\alpha=1$ and $h$ even.
The congruence uses the full multiple fibre through its oriented
meridian relation.

\item Let $h$ be even and put
\[
 b_*=-\tfrac h4(\widehat u+\widehat w),\qquad
 L_*=-\tfrac h2\widehat u.
\]
Direct multiplication gives
\[
 (I-H_*)b_*=\tfrac h4\widehat\delta+L_*.
\]
Consequently $\zeta\mapsto\zeta+\Pi_{\mu,c}b_*$ on the
complete root torus family conjugates the two affine actions,
modulo the integral arrow $L_*$.  It and its inverse extend
over $r=0$ and descend.  On the punctured root collar its
middle translation includes
$h(\log r)/(2\pi i)e$; this is the derivative-free collar
term for the change of the invariant translation
$\Pi\ell/4$; the central map is the holomorphic root-family
translation just constructed.

\item At a rank-one point the vanishing coordinate is the first
toric coordinate in its chain cover, while the other
direction is the universal cover of the elliptic factor.
An extending translation has integral winding only in $K_j$.
After removing the corresponding power of $q$, its toric
coefficient is a holomorphic unit and its elliptic translation
extends holomorphically.  These are the full analytic collar
conditions.  Necessity follows by lifting a transverse disc
through a smooth central point to this free cover: its residual
defect in $\Lambda/K_j$ is zero.  The coordinate ratios on
the cover then give the asserted unit and holomorphic
remainder.  Conversely an integral power of $q$ translates
the infinite chain by an integer number of components;
the holomorphic remainders give torus and elliptic translations.
These maps and their inverses extend at both branches of every
node.  At the hexagonal point the same description uses two
toric windings in $K_3$, and integer translation of the fan
extends over its three double curves and both triple points.
The criterion requires both the specified integral winding and
the holomorphic unit remainder.

\item Choose the three integral defects
\[
 L_1=0,\qquad L_2=\tfrac h2\widehat u,\qquad
 L_3=-\tfrac h2\widehat w.
\]
They belong to $K_1,K_2,K_3$, respectively, and are realized
by the preceding toric translations with initially constant
unit remainder.  In primitive cocharacter coordinates these
are $1$, $q_2^{h/2}$, and $(q_3^{-h/2},1)$, where the
monodromy defect of $q^a$ is defined as $a$ times the chosen
vanishing generator.  The required sphere relation is
\[
 L_*+H_*L_1+H_*H_1L_2+H_*H_1H_2L_3=0.
\]
Multiplication of the four displayed monodromy matrices gives
zero on the left.
Thus the four local choices have compatible additive defects
and define a torsor of allowed middle translations.

\item The root curve has the order-four
point at infinity and is $\mathfrak C=\mathbb P(1,4)$,
with $K_{\mathfrak C}=\mathcal O(-5)$ and coarse pullback
$\mathcal O_B(1)=\mathcal O_{\mathfrak C}(4)$.
Vertical logarithmic fields on a rank-one chart are generated
by $a\partial_a-b\partial_b$ and the elliptic vector;
on a hexagonal chart they are the two height-zero toric fields.
The root chart has the ordinary vertical tangent frame.
Their linear transition functions depend only on the base;
they define a rank-two bundle $\mathcal E$ on $\mathfrak C$.
The fixed line-coordinate field supplies a subline
$\mathcal O\subset\mathcal E$ everywhere in these logarithmic
frames.  The logarithmic determinant sequence and cancellation
of the three reduced fibre divisors give
$K_X=\bar f^*(K_{\mathfrak C}\otimes(\det\mathcal E)^{-1})$.
The canonical formula and
$\bar f_*\mathcal O_X=\mathcal O_{\mathfrak C}$ therefore
give $\det\mathcal E=\mathcal O(-1)$.
The extension $0\to\mathcal O\to\mathcal E\to\mathcal O(-1)\to0$
splits since $H^1(\mathfrak C,\mathcal O(1))=0$.
In the two-chart Laurent calculation with weights $1,4$,
a first-cohomology monomial has both exponents negative and
hence weighted degree at most $-5$. This gives
$H^1(\mathcal O(k))=0$ for $k=1,0,-1$.
The coarse difference sheaf is $V_B=\rho_*\mathcal E$:
on the root chart this is precisely finite equivariance,
and on each reduced collar it is precisely the holomorphic
remainder just described.  Finite averaging is exact, so
\[
 H^1(B,V_B)=H^1(\mathfrak C,\mathcal E)
 =H^1(\mathcal O\oplus\mathcal O(-1))=0.
\]
The torsor's \v{C}ech class vanishes.  Correcting the four local
translations by its holomorphic zero-cochain gives a global
middle translation and exact equality with all four local
extensions.  Their inverses have the opposite allowed
translations and glue as well.  This constructs the required
biholomorphism over $B$.

\item Choose a single logarithm of $\lambda$ on the
small bidisc.  Its $c$-values can be placed in a strip of
width less than two.  The criterion just proved then says
$c'=c$ and $\mu'=\mu$, which is equivalent to equality of
the original two parameters.  Conversely equality gives the
identity, with the fixed root, rank-one and hexagonal charts.

\end{enumerate}
\end{proof}

\section{Hodge growth and cup-product kernels of combs}\hypertarget{tag-034D}{}\label{sec:II:combs}

\subsection{Comb cup kernels and Hodge numbers}\label{subsec:II:comb-cup-hodge}

\begin{definition}[Cup-product kernel of a positive comb]\hypertarget{tag-P278}{}\label{def:II:comb-cup-kernel}
\StatementRecord{def:II:comb-cup-kernel}{proved}{P278}{\SeriesRef{cdp:st:E15}; \cite[Proposition~6.1, p.~19]{Viaclovsky2026Combs}}
Let $X_m$ be a smooth $m$-tooth member with $m\ge12$, fixed limiting
period planes and standard logarithmic data, as in \cref{cdp:st:E15}.
Write $\bar f_m:X_m\to\mathcal B=\mathbb P(3,4)$ for its root lift,
$\rho:\mathcal B\to B$ for the coarse map, $F$ for the Hodge
bundle of \cref{cdp:st:E15}, $E_m=F^*$ for its vertical bundle,
and $G_m=R^1\bar f_{m*}\mathcal O_{X_m}$.
Exterior product with the logarithmic Kodaira--Spencer tensor defines
\begin{equation}
 K_m=\ker\left\{\rho_*(E_m^*\otimes G_m)
  \xrightarrow{\ \kappa_1\ }
       \rho_*(\mathcal O(5)\otimes\det G_m)\right\}.
 \hypertarget{tag-01TV}{}\label{invII:eq:E.6}
\end{equation}
This is a rank-three vector bundle on $B$.
\end{definition}
\begin{proof}
At the cusp, \cite[Proposition~6.1, p.~19]{Viaclovsky2026Combs}
gives periods $(\log q)I/(2\pi i)+H_m(q)$ with $H_m$ holomorphic.
Their logarithmic derivative is $I/(2\pi i)+qH'_m(q)$ and is
invertible at zero. Thus $\kappa_1$ has generic rank one.
Its kernel is a torsion-free subsheaf of a rank-four bundle on the
smooth curve $B$, hence is locally free of rank three.
\end{proof}

\Needspace{10\baselineskip}
\begin{proposition}[Canonical bundle, exact tangent cohomology and Hodge growth on the combs]\hypertarget{tag-01SV}{}\label{inv:E2.2}
\StatementRecord{inv:E2.2}{proved}{01SV}{}

Let $X_m$ be a smooth $m$-tooth member of the construction in
\cite[Theorem~1.1, p.~2, arXiv version~1]{Viaclovsky2026Combs},
with $m\ge12$, fixed limiting period planes at the three marked
values and the standard logarithmic data. Its root lift
$\bar f_m:X_m\to\mathcal B$ has
vertical bundle $E_m$ and $G_m=R^1\bar f_{m*}\OO_{X_m}$ satisfying
\begin{equation}
 0\longrightarrow\OO\longrightarrow E_m
       \longrightarrow\OO(24m-1)\longrightarrow0,
 \qquad
 0\longrightarrow\OO\longrightarrow G_m
       \longrightarrow\OO(24m-1)\longrightarrow0.      \hypertarget{tag-01TR}{}\label{invII:eq:E.2}
\end{equation}
The canonical bundle and all positive-degree holomorphic forms obey
\begin{equation}
 K_{X_m}=\bar f_m^*\OO(-24m-6),\qquad
 h^{1,0}=h^{2,0}=h^{3,0}=0.                            \hypertarget{tag-01TS}{}\label{invII:eq:E.3}
\end{equation}
The tangent cohomology is that of Proposition~\ref{cdp:st:E15}.
More precisely, writing $\mathcal L=\OO_{\mathcal B}(-5)$,
\begin{equation}
 \begin{split}
 f_{m*}T_{X_m}&=\rho_*E_m,\\
 R^1f_{m*}T_{X_m}&=
  \operatorname{coker}\{\rho_*\mathcal L\xrightarrow{\delta_0}
                                      \rho_*(E_m\otimes G_m)\},\\
 R^2f_{m*}T_{X_m}&\simeq\OO_B(2m)\oplus Q_m,
 \qquad\operatorname{length}Q_m=4m,
 \end{split}                                        \hypertarget{tag-01TT}{}\label{invII:eq:E.4a}
\end{equation}
where $Q_m=\operatorname{coker}\{\rho_*(\mathcal L\otimes G_m)
\xrightarrow{\delta_1}\rho_*(E_m\otimes\det G_m)\}$ has no support
at the cusp. 
For the rank-three bundle $K_m$ of \cref{def:II:comb-cup-kernel},
\begin{equation}
 \boxed{h^{1,1}(X_m)=4m+1+h^0(K_m),\qquad
 h^{1,2}(X_m)=4m+h^0(K_m).}                            \hypertarget{tag-01TU}{}\label{invII:eq:E.5}
\end{equation}
In particular both remaining middle Hodge numbers grow at least
linearly with $m$.
\end{proposition}


\begin{proof}
\begin{enumerate}[label=\arabic*.,leftmargin=*]
\item Put $a=24m-1$ and $A=\OO_{\mathcal B}(a)$.
In the notation of \cref{cdp:st:E15}, the identification
$V=E_m=F^*$ sends a real homology class $v$ to the
functional $\alpha\mapsto\alpha(v)$ on the Hodge plane;
$G=G_m$ is the quotient $\mathcal H/F$.
That proof establishes the nowhere-zero invariant sections
of $V,G$ and their common determinant $A$, including at
the root points and the cusp. Their quotients are therefore
the line bundle $A$, proving \eqref{invII:eq:E.2}.
Exact invariant pushforward gives
\[
 0\longrightarrow\OO_B\longrightarrow\rho_*V
 \longrightarrow\OO_B(2m-1)\longrightarrow0,
 \qquad V=E_m,G_m.
\]
The nonnegative splittings, their cohomology and the
tensor-product comparison are \cref{cdp:eq:E15-positive,cdp:eq:E15-products}
under these same bundle identifications.

\item The logarithmic tangent sequence has vertical term
$\bar f_m^*E_m$ and horizontal term $\bar f_m^*\mathcal L$,
where $\mathcal L=T_{\mathcal B}(-p_0)=\OO(-5)$.
Its determinant gives
$K_{X_m}(W)=\bar f_m^*\OO(-24m+6)$.
Since $\OO_{X_m}(W)=\bar f_m^*\OO(12)$, this proves
the canonical formula in \eqref{invII:eq:E.3}.

By \cref{def:II:comb-cup-kernel}, the logarithmic
Kodaira--Spencer tensor is invertible at the cusp.
The connecting map
$E_m^*\to\OO(5)\otimes G_m$ of the logarithmic
cotangent sequence is injective, so its degree-zero
direct image is $\OO(5)$.
This line has no global section by \SeriesRef{inv:D1.1}.
Hence $h^{1,0}=0$.
The second exterior sequence has subquotients
$\bar f_m^*(E_m^*\otimes\OO(5))$ and
$\bar f_m^*(\det E_m)^{-1}$.
The latter has negative weight $1-24m$.
A map $E_m\to\OO(5)$ is zero on its trivial subline,
since $H^0(\OO(5))=0$, and factors through $A$;
the resulting section has weight $6-24m<0$.
Both spaces of sections vanish, proving $h^{2,0}=0$.
The negative canonical weight gives $h^{3,0}=0$.

\item Write the logarithmic Kodaira--Spencer tensor locally as
$e_1\otimes(ug_1+vg_2)+e_2\otimes(wg_1+zg_2)$.
The connecting maps are multiplication by this tensor
and its exterior product. The matrix of
$\delta_1:\mathcal L\otimes G_m\to E_m\otimes\det G_m$
is $\left(\begin{smallmatrix}-v&u\\-z&w\end{smallmatrix}\right)$,
whose determinant is $uz-vw$.
It is nonzero at the cusp; both $\delta_0$ and
$\delta_1$ are therefore sheaf injections.
The logarithmic direct-image sequence gives
$\bar f_{m*}T(-\log W)=E_m$ and
$R^1\bar f_{m*}T(-\log W)=\operatorname{coker}\delta_0$.

For the ordinary tangent sheaf, differentiation of
$q=x_1x_2x_3$ gives the exact normal quotient
\[
 0\longrightarrow T(-\log W)\longrightarrow T_{X_m}
 \longrightarrow J\longrightarrow0,
 \qquad J=(x_2x_3,x_1x_3,x_1x_2)\subset\OO_W.
\]
Let $\nu:P\to W$ be the normalization, $P=dP_6$,
with toric boundary $D$.
In local frames of $\OO_P(-D)$, the map is
$(h_1,h_2,h_3)\mapsto
h_1x_2x_3+h_2x_1x_3+h_3x_1x_2$ modulo $x_1x_2x_3$,
where $h_i$ is a function on the branch $x_i=0$.
The inverse restricts to that branch and divides
by $x_jx_k$; the two compositions are identities.
These maps identify $J=\nu_*\OO_P(-D)=\nu_*K_P$;
the toric logarithmic two-form identifies $K_P(D)$
with $\OO_P$.
Serre duality on the rational surface $P$ gives
$h^\bullet(J)=(0,0,1)$.
Thus the ordinary and logarithmic direct images
agree in degrees zero and one. Exact invariant
pushforward proves the first two lines of
\eqref{invII:eq:E.4a}.

\item The degree-two logarithmic sequence is
\[
 0\longrightarrow\operatorname{coker}\delta_1
 \longrightarrow R^2\bar f_{m*}T(-\log W)
 \longrightarrow\mathcal L\otimes\det G_m
                 =\OO(24m-6)\longrightarrow0.
\]
The root-degree formula in \SeriesRef{inv:D1.1}
gives degrees $-1$, $2m-1$, $4m-1$ for
$\rho_*\mathcal L$, $\rho_*A$ and $\rho_*A^2$,
and degree $2m-1$ for $\rho_*(\mathcal L A)$.
Using the two extensions in \eqref{invII:eq:E.2},
the degrees of $\rho_*(\mathcal L\otimes G_m)$
and $\rho_*(E_m\otimes\det G_m)$ are consequently
$2m-2$ and $6m-2$. The injective same-rank map
has cokernel $Q_m$ of length $4m$.
It is zero at the cusp, where the cup matrix is invertible.

The quotient $J$ adds a length-one quotient at $p_0$
to the free part of the top direct image.
This is a modification of a line bundle: indeed,
$\omega_W=\OO_W$ and Serre duality identifies
$h^2(W,T_{X_m}|_W)$ with
$h^0(W,\Omega^1_{X_m}|_W)$.
Pull such a one-form to $P$.
Its image in $\Omega^1_P$ is zero, so it is conormal
on each branch. On a double curve, equality of the
two expressions $c_i\,dx_i$ and $c_j\,dx_j$
forces both coefficients to vanish there.
At a triple point this makes $c_i$ divisible by
$x_jx_k$. Hence the form is $h\,dq$, with
$h$ holomorphic on the compact connected $P$.
Thus $h$ is constant, and $dq$ generates this
one-dimensional space. On nearby smooth fibres
the invertible cup matrix likewise gives $h^2=1$.
The map to the smooth base curve is flat, and
$T_{X_m}$ is locally free. Constant-dimension base change
\cite[\S7, Satz~5, p.~64; hypotheses and proof, pp.~62--64]{Grauert1960D}
makes $R^2f_{m*}T_{X_m}$
locally free at $p_0$.
Its free line is therefore
$\rho_*(\mathcal L\otimes\det G_m)(p_0)=\OO_B(2m)$.
The extension by $Q_m$ splits because
$\operatorname{Ext}^1(\OO_B(2m),Q_m)=H^1(Q_m(-2m))=0$.
This proves the last line of \eqref{invII:eq:E.4a}.

\item The map preceding $\kappa_1$ in the cotangent direct-image sequence
\[
 \rho_*E_m^*\xrightarrow{\ \kappa_0\ }
                 \rho_*(\OO(5)\otimes G_m)
\]
is the transpose tensor and is injective.
The source extension has line degrees $-2m-1,0$;
the target extension has line degrees $-1,2m$.
Their total degrees are $-2m-1$ and $2m-1$,
so its cokernel $A_m$ has length $4m$.

\item Put $\mathcal A_m=\Omega^1_{X_m}(\log W)$.
The direct-image sequence gives
$0\to A_m\to R^1f_{m*}\mathcal A_m\to K_m\to0$.
In the residue sequence
$0\to\Omega^1_{X_m}\to\mathcal A_m\to\nu_*\OO_P\to0$,
the form $dq/q$ maps to residue one on every branch.
The degree-zero residue map is surjective and
\[
 f_{m*}\Omega^1_{X_m}=\OO_B(-2),\qquad
 f_{m*}\mathcal A_m=\rho_*\OO(5)=\OO_B(-1).
\]
Since $P$ is rational, the residue quotient has no
higher direct images. Thus
$R^1f_{m*}\Omega^1_{X_m}=R^1f_{m*}\mathcal A_m$.
Its global sections have dimension $4m+h^0(K_m)$,
and Leray adds $h^1(\OO_B(-2))=1$ in total degree one.
This proves the first formula in \eqref{invII:eq:E.5}.
Serre duality gives
$h^{1,3}=h^0(T_{X_m}\otimes K_{X_m})=0$:
its direct image has the two negative root line
subquotients $\OO(-a-7)$ and $\OO(-7)$.
Finally Hirzebruch--Riemann--Roch gives
$\chi(\Omega^1)=c_1c_2/24-c_3/2=-1$ on $S^6$.
Together with $h^{1,0}=h^{1,3}=0$, this yields
$h^{1,1}-h^{1,2}=1$ and the second formula.
\end{enumerate}
\end{proof}

\subsection{Degree bounds and cup duality}\label{subsec:II:comb-kernel-bounds-duality}

\begin{proposition}[Degree and section bounds for comb kernels]\hypertarget{tag-P320}{}\label{prop:II:comb-kernel-bounds}
\StatementRecord{prop:II:comb-kernel-bounds}{proved}{P320}{\SeriesRef{def:II:comb-cup-kernel}, \SeriesRef{cdp:st:E15}, \SeriesRef{inv:E2.2}}
Use the positive-comb data $X_m,E_m,G_m,K_m,\rho$ of
\cref{def:II:comb-cup-kernel}, with $m\ge12$. Put
$e_m=\operatorname{length}\operatorname{coker}\kappa_1$.
The source of $\kappa_1$ has rank four and degree $-2$, its
target is $\mathcal O_B(2m)$, and
\begin{equation}
 \deg K_m=-2m-2+e_m.
 \hypertarget{tag-01TW}{}\label{invII:eq:E.7}
\end{equation}
The section and Hodge numbers satisfy
\[
 0\le h^0(K_m)\le2m+2,\qquad
 4m+1\le h^{1,1}(X_m)\le6m+3,\qquad
 4m\le h^{1,2}(X_m)\le6m+2.
\]
On $B=\mathbb P^1$ with homogeneous coordinates $[X:Y]$, put
$D_0=\mathcal O_B^{\oplus2}\oplus\mathcal O_B(-1)^{\oplus2}$.
The surjections
\[
 \phi_A=(X^{2m},Y^{2m},0,0),\qquad
 \phi_B=(0,0,X^{2m+1},Y^{2m+1}):D_0\longrightarrow\mathcal O_B(2m)
\]
have kernels
\[
 \ker\phi_A\simeq\mathcal O_B(-2m)\oplus\mathcal O_B(-1)^{\oplus2},
 \qquad
 \ker\phi_B\simeq\mathcal O_B^{\oplus2}\oplus\mathcal O_B(-2m-2).
\]
Their source ranks, source degrees, target degrees and cokernel
lengths agree, while their kernel section dimensions are zero
and two, respectively.
\end{proposition}

\begin{proof}
\begin{enumerate}[label=\arabic*.,leftmargin=*]
\item Write $V_c=\rho_*E_m$, $G_c=\rho_*G_m$ and
$\mathcal A_m=\rho_*(E_m^*\otimes G_m)$.
Applying an equivariant homomorphism to invariant sections gives
\[
 \mathcal A_m\longrightarrow\mathcal Hom(V_c,G_c).
\]
This map is an isomorphism off the two finite root points. Its
source is torsion free, so it is injective. At a root point take
a coordinate $w$ with $t=w^r$ on $B$ and equivariant frames
$(e_1,e_2)$ and $(g_1,g_2)$ of $E_m$ and $G_m$.
By \cref{cdp:st:E15}, the invariant frames are
$(e_1,w^ke_2)$ and $(g_1,w^kg_2)$, where
$(r,k)=(3,2)$ or $(4,1)$. The equivariant matrix and its coarse
matrix are, respectively,
\[
 \begin{pmatrix}a(t)&w^{r-k}b(t)\\ w^kc(t)&d(t)\end{pmatrix},
 \qquad
 \begin{pmatrix}a(t)&t b(t)\\ c(t)&d(t)\end{pmatrix}.
\]
These formulas are inverse descriptions of the same morphism:
the sole coarse condition is the vanishing of the upper-right
entry at $t=0$. Hence
\[
 0\longrightarrow\mathcal A_m
 \longrightarrow\mathcal Hom(V_c,G_c)
 \longrightarrow\mathbb C_{p_1}\oplus\mathbb C_{p_2}
 \longrightarrow0.
\]
Both coarse rank-two bundles have degree $2m-1$, so their Hom
bundle has degree zero and $\deg\mathcal A_m=-2$.
On the root curve the determinant of $E_m^*\otimes G_m$ is
trivial, since $\det E_m=\det G_m=\mathcal O(24m-1)$.
The local comparison exhibits the two opposite nontrivial
characters and their one-unit coarse correction at each root point.

The target is $\rho_*\mathcal O(24m+4)$.
In weighted coordinates $[x:y]$ of weights $3,4$, the coarse
map is $[x:y]\mapsto[x^4:y^3]$, so
$\mathcal O(12)=\rho^*\mathcal O_B(1)$.
Multiplication by $y$ identifies
$\mathcal O_B\simeq\rho_*\mathcal O(4)$: it is a frame where
$y\ne0$, and at $y=0$ its order-one coefficient is exactly
the invariant generator at the order-three point.
Division by that local invariant generator is the inverse map.
The projection formula therefore identifies the target with
$\mathcal O_B(2m)$. The generic rank-one assertion of
\cref{def:II:comb-cup-kernel} makes its cokernel a torsion sheaf
of length $e_m$. Its image is a line subsheaf of degree $2m-e_m$.
Additivity of degree in the kernel sequence gives
\eqref{invII:eq:E.7}.

\item Put $S=2m-1$. The nonnegative splittings proved in
\cref{cdp:st:E15} have the form
\[
 V_c=\mathcal O_B(a)\oplus\mathcal O_B(S-a),\qquad
 G_c=\mathcal O_B(b)\oplus\mathcal O_B(S-b).
\]
Exchanging the two summands when necessary gives
$0\le a,b\le S/2$. The four summands of
$\mathcal Hom(V_c,G_c)$ contribute
\[
 \max\{b-a+1,0\},\quad S-a-b+1,\quad
 \max\{a+b-S+1,0\},\quad\max\{a-b+1,0\}
\]
to its number of sections. For $a+b<S$ their sum is at most
\[
 S-a-b+1+|a-b|+2
   =S+3-2\min\{a,b\}\le S+3.
\]
For $a+b=S$ the two chosen inequalities force $a=b=S/2$;
the sum is four and is also at most $S+3$.
The inclusions
$K_m\hookrightarrow\mathcal A_m\hookrightarrow\mathcal Hom(V_c,G_c)$
induce injections on sections, giving
$h^0(K_m)\le S+3=2m+2$.
Substitution into \eqref{invII:eq:E.5}, for this same $X_m$ and
$K_m$, gives the Hodge bounds.
At $m=12$ they give $0\le h^0(K_{12})\le26$,
$h^{1,1}\ge49$ and $h^{1,2}\ge48$; the tangent tuple from
\cref{cdp:st:E15} is $(25,97,73,0)$.

\item The nonzero entries of each $\phi_A,\phi_B$ have no common
projective zero, so both maps are surjective.
The two kernel isomorphisms in the statement are given by
\[
 (s,u,v)\longmapsto(-Y^{2m}s,X^{2m}s,u,v),\qquad
 (u,v,s)\longmapsto(u,v,-Y^{2m+1}s,X^{2m+1}s).
\]
For the first map the inverse coefficient is the second entry
divided by $X^{2m}$ on $X\ne0$, and minus the first entry divided
by $Y^{2m}$ on $Y\ne0$. The kernel equation identifies these
coefficients on the overlap. The last two entries are unchanged.
For the second map the same divisions use the last two entries
and exponent $2m+1$, with the first two entries unchanged.
The respective coefficient lines are
$\mathcal O_B(-2m)$ and $\mathcal O_B(-2m-2)$, so these local
inverses glue to the displayed bundle isomorphisms.
Both kernels have rank three and degree $-2m-2$.
Their displayed summands give zero and two global sections.
For $m=12$ their splitting triples are $(-24,-1,-1)$ and
$(0,0,-26)$.
\end{enumerate}
\end{proof}

\begin{proposition}[Tangent and cotangent cup duality]\hypertarget{tag-P321}{}\label{prop:II:comb-cup-duality}
\StatementRecord{prop:II:comb-cup-duality}{proved}{P321}{\SeriesRef{def:II:comb-cup-kernel}, \SeriesRef{cdp:st:E15}, \SeriesRef{invII:eq:E.4a}}
For the same positive-comb data, put
$F=E_m^*$, $G=G_m$, $D=\det G=\mathcal O(24m-1)$,
$\mathcal L=\mathcal O(-5)$ and
\[
 N=\operatorname{coker}\{\delta_0:\mathcal L\to E_m\otimes G\}.
\]
Define the wedge isomorphism by
$J:G\to G^*\otimes D$, $J(g)(h)=h\wedge g$.
With the same \v{C}ech boundary convention for the tangent and
cotangent sequences, their root cup maps satisfy
\[
 \kappa_{\mathrm{root}}
 =-(\delta_0^*\otimes I_D)(I_F\otimes J).
\]
The restriction of $I_F\otimes J$ gives the isomorphism
\[
 K_m\simeq\rho_*(N^*\otimes D).
\]
\end{proposition}

\begin{proof}
In an ordered frame of $G$, the matrix of $J$ is
$\left(\begin{smallmatrix}0&1\\-1&0\end{smallmatrix}\right)$,
and its inverse is its negative. In local splittings, inverse
transposition of the tangent-extension transition matrix sends
its off-diagonal \v{C}ech cocycle $\beta$ to $-\beta^t$.
Consequently dualizing the extension and then taking exterior
product gives the stated composition. More explicitly, write
\[
 \delta_0(\ell)=u_1\otimes(ag_1+bg_2)
                     +u_2\otimes(cg_1+dg_2),
\]
where $(u_1,u_2)$ and $(f_1,f_2)$ are dual frames of $E_m$ and
$F$. In the ordered basis $(f_1g_1,\allowbreak f_1g_2,\allowbreak f_2g_1,\allowbreak f_2g_2)$,
the row of the displayed composition is $(b,-a,d,-c)$,
which is the contracted dual-extension cup product.

Dualizing the cokernel sequence identifies
\[
 N^*=\ker\{\delta_0^*:F\otimes G^*\to\mathcal L^*\}.
\]
Thus $I_F\otimes J$ maps $\ker\kappa_{\mathrm{root}}$
onto $N^*\otimes D$, with inverse the restriction of
$I_F\otimes J^{-1}$. Both maps are equivariant and glue,
because $J$ is defined by exterior product.
The dual annihilates every torsion section of $N$, since a
homomorphism from a torsion sheaf to the structure sheaf of
the smooth root curve is zero. The identification therefore
uses the stated cokernel $N$ itself.
Taking finite-stabilizer invariants preserves the kernel:
at an order-$r$ point its averaging projection is
$r^{-1}\sum_g g$. Applying $\rho_*$ to the two inverse maps
gives the asserted isomorphism with $K_m$ of
\eqref{invII:eq:E.6}.
By \eqref{invII:eq:E.4a}, $\rho_*N=R^1f_{m*}T_{X_m}$,
and $Q_m$ is the cokernel of the exterior cup product of the
same period tensor. Its support and elementary divisors are
therefore those of that specified morphism.
\end{proof}

\subsection{Vector fields and twelve-tooth smoothings}\label{subsec:II:twelve-tooth-smoothings}

\begin{proposition}[Comb members with more than two vector fields]\hypertarget{tag-01SX}{}\label{inv:E2.4}
\StatementRecord{inv:E2.4}{proved}{01SX}{\SeriesRef{inv:E2.2}; \citet[Corollary~1.3, arXiv v1, p.~2]{HKP2000}; \citet[Remark~1.2, p.~2, and Remark~9.26, p.~83]{AC26}}

The assertion that every complex structure on $S^6$ has
$h^0(T)\le2$, printed in
\citet[Corollary~1.3, arXiv v1, p.~2]{HKP2000}
and repeated in \citet[Remark~1.2, p.~2, and Remark~9.26, p.~83]{AC26}, is false
for the comb members.
\end{proposition}


\begin{proof}
Proposition~\ref{inv:E2.2} gives $h^0(T_{X_m})=2m+1\ge25$ for $m\ge12$.
In the proof of \citet[Corollary~1.3, arXiv v1, p.~2]{HKP2000},
a third field is expressed as a meromorphic linear combination of two fields;
the coefficients are then declared constant using the
Campana--Demailly--Peternell assertion that every meromorphic function
on a complex six-sphere is constant. For the combs, the surjective algebraic reduction
$f_m$ induces an injection $f_m^*:\C(t)\hookrightarrow\mathcal M(X_m)$,
so its pullback of $t$ is nonconstant.
\end{proof}

\begin{proposition}[The kernel on an open family of twelve-tooth smoothings]\hypertarget{tag-03LZ}{}\label{inv:E2.5}
\StatementRecord{inv:E2.5}{proved}{03LZ}{\SeriesRef{inv:E2.2}, \SeriesRef{cdp:st:E15}, \SeriesRef{invII:eq:E.6}; \cite{Viaclovsky2026Combs}}
Use the original marked section of the comb construction on
$\mathcal B=\mathbb P(3,4)$, with coarse base $B=\mathbb P^1$, the
order-three and order-four points $p_1,p_2$, and the hexagonal cusp $p_0$.
Its bundles are $E_0=G_0=\mathcal O\oplus\mathcal O(-1)$ and
$\mathcal L=\mathcal O(-5)$. There is a nonempty open family of
admissible positive-comb manifolds $X_{12}$ obtained by embedded
smoothings of twelve deformed twistor conics attached at distinct
ordinary points, fixing the three marked period planes, for which the
kernel in \eqref{invII:eq:E.6} satisfies
\[
 K_{12}\simeq\mathcal O_B(-8)\oplus\mathcal O_B(-9)^{\oplus2},
 \qquad \operatorname{coker}\kappa_1=0.
\]
On this family,
\[
 h^0(K_{12})=0,\qquad h^1(K_{12})=23,\qquad
 h^{1,1}(X_{12})=49,\qquad h^{1,2}(X_{12})=48.
\]
\end{proposition}

\begin{proof}
\begin{enumerate}[label=\arabic*.,leftmargin=*]
\item 
\textbf{The seed kernel.}
Put $F_0=E_0^*$ and $D_0=\det G_0$. Write $x,y$ for the homogeneous
coordinates of weights $3,4$, respectively. Thus $y$ vanishes at the
order-three point. In row-$G$ and column-$F$ convention the 
logarithmic period derivative has the form
\[
 T_0=\begin{pmatrix}0&ay\\by&cx\end{pmatrix},\qquad ab\ne0.
\]
Indeed its entries have weights $5,4,4,3$; the corresponding section
spaces are $0,\mathbb Cy,\mathbb Cy,\mathbb Cx$. Its invertible
logarithmic residue at the cusp forces $a,b$ to be nonzero.
For a tensor matrix $\eta$ in row-$F$ and column-$G$ convention define
\[
 \langle T,\eta\rangle
   =\sum_{i=1}^2(T_{2i}\eta_{i1}-T_{1i}\eta_{i2})
   =\operatorname{tr}(T\eta J),\qquad
 J=\begin{pmatrix}0&1\\-1&0\end{pmatrix}.
\]
This is the determinant-valued cup pairing. The seed row is
$(by,0,cx,-ay)$, with source weights $0,-1,1,0$ and target weight $4$.
Exact finite invariants give
\[
 \rho_*\mathcal O(4)=\mathcal O_B,
 \qquad \rho_*\mathcal O(1)=\rho_*\mathcal O(-1)=\mathcal O_B(-1).
\]
For the last two identities the invariant sheaves are lines with zero
$H^0$ and $H^1$, as follows from weighted homogeneous sections and
Serre duality with $K_{\mathcal B}=\mathcal O(-7)$.
The invariant section $y$ is the global coarse frame of the target.
Let $\sigma_x:\rho_*\mathcal O(1)\to\rho_*\mathcal O(4)$ be the
multiplication-by-$x$ map. The coarse row is
$(b,0,c\sigma_x,-a)$, which is onto. Its kernel isomorphism is
\[
 (u,v,w)\longmapsto
 (au/b-c\sigma_xv/b,\ w,\ v,\ u),
 \qquad \eta\longmapsto(\eta_{22},\eta_{21},\eta_{12}).
\]
Consequently $K_0=\mathcal O_B\oplus\mathcal O_B(-1)^2$.
This computation includes the possible $cx$ column. The determinant
$-ab y^2$ also proves invertibility of the handle derivative at
every ordinary attachment point.

\item \textbf{The root family and the flat relative kernel.}
Choose twelve distinct ordinary points. The conics chosen below lie
in the fixed-marking embedded smoothing construction of
\cite[Proposition~5.3, pp.~17--18]{Viaclovsky2026Combs}.
Here is the family lift at the finite points. In equivariant
frames $f_i,g_j$, the invariant generators are $w^{a_i^F}f_i$ and
$w^{a_j^G}g_j$. A coarse graph matrix $C(b,s)$ lifts by
\[
 Z_{ji}(w,s)=w^{a_j^G-a_i^F}C_{ji}(w^n,s).
\]
This follows by dividing its invariant graph vector by
$w^{a_i^F}$; multiplying back gives the inverse graph map.
At $p_1$ the order is $3$ and the exponent matrix is
$\left(\begin{smallmatrix}0&-1\\2&1\end{smallmatrix}\right)$.
At $p_2$ the order is $4$ and it is
$\left(\begin{smallmatrix}0&-3\\1&-2\end{smallmatrix}\right)$.
Since the entire period plane stays fixed along each marked section,
the required entries have the joint holomorphic divisions
$C_{12}=bD_{12}$, and also $C_{22}=bD_{22}$ at $p_2$.
Their lifts are respectively
\[
 w^2D_{12}(w^3,s),\qquad
 wD_{12}(w^4,s),\qquad w^2D_{22}(w^4,s).
\]
These are precisely the finite-character conditions of
\ref{cdp:st:E15}; all other exponents are nonnegative.
They construct the family Hodge and quotient bundles, with
their unchanged finite characters. At the cusp the chart in that
proposition is $(\log q)I/(2\pi i)+H(q,s)$ with $H$ jointly
holomorphic and bounded on a smaller parameter polydisc. Its residue
is unchanged. The same finite-cover actions, toric unit translations
and annular attachments give the compact family as in
\ref{cdp:st:E15}. In particular its coarse cusp cup row stays onto.

Choose a smoothing direction nonzero in every node-smoothing space.
Its regular total surface has coordinates $xy=s$ at each node,
with handle $y=0$ and tooth $x=0$. In a flat cohomology
frame write the period plane as the graph of $Z(x,y)$.
The relative dualizing generator $dx/x=-dy/y$ gives the derivative
\[
 \Theta=xZ_x-yZ_y.
\]
The relative cup takes values in
$D\otimes\omega_{\mathcal C/\Delta}(p_0)$, with the finite roots
along the marked sections understood. Put $H=Z_x(0,0)$ and
$B=Z_y(0,0)$. Our conic choices make the two functionals
$\langle H,-\rangle$ and $\langle B,-\rangle$ independent.
The cup-row entries then generate $(x,y)$ by their independent
linear terms and Nakayama's lemma. Invertible local operations
reduce the row to $(x,y,0,0)$, with free kernel generated by
$(-y,x,0,0)$ and the last two coordinate vectors.
Thus the relative kernel is locally free at every node.
Away from the nodes the central coarse row is onto; on the teeth
this follows from the conic calculation in Step 4. Its unit
coefficients persist after shrinking the family. Exact finite
invariants commute with parameter restriction in the character
frames just constructed. Finally $(x,y)$ is the unit ideal on a
noncentral fibre. We obtain a locally free coarse relative kernel
whose smooth restriction is the $K_{12}$ and whose smooth
coarse cokernel is zero.

\item \textbf{The complete node condition.}
Let $\eta_h(x),\eta_t(y)$ be the two branch sections, with common
ordinary value $\eta_0$. Their value equations are
\[
 \langle H,\eta_0\rangle=\langle B,\eta_0\rangle=0.
\]
Write $Z=Z_h(x)+Z_t(y)-Z(0,0)+xyR(x,y)$. The coefficient of $xy$
in $xZ_x-yZ_y$ is zero, including the arbitrary mixed graph term.
The mixed coefficient in the kernel equation is therefore
\begin{equation}
 \langle H,\eta_t'(0)\rangle
       -\langle B,\eta_h'(0)\rangle=0.                 \hypertarget{tag-03M0}{}\label{invII:eq:E12-node}
\end{equation}
This is also sufficient to all analytic orders. Over
$R_0=\mathbb C\{x,y\}$ let $\varphi:R_0^4\to R_0$ be the 
row, with image $(x,y)$. A lift of the branch pair has
$\varphi(\widetilde\eta)=xyh$, and $h(0)$ is exactly the left side
of \eqref{invII:eq:E12-node}. Adding $xy\zeta$ changes $h$ by
$\varphi(\zeta)$, whose values range over $(x,y)$. The displayed
vanishing therefore gives an analytic correction with
$\varphi(\zeta)=-h$. The corrected lift is a section of the
total kernel with the prescribed restrictions. This proves that
the one mixed-jet equation is the entire remaining condition.
The kernel of the central row exceeds the central fibre of the
flat kernel by just one node copy of $\mathbb C$.
Changing a determinant or dualizing generator adds multiples of
the two zero value equations. Tensor-frame changes transport both
derivatives, including the derivatives of the frames.

\item \textbf{Conics and their two-jet map.}
In the flat decomposition $F_b\oplus G_b$ the through-point conics
under consideration have homogeneous inclusion and quotient
\[
 S_{[u:v]}=\binom{uI+vA}{vB},\qquad
 Q_{[u:v]}=(-vB,\ uI+vBAB^{-1}),\qquad \det B\ne0.
\]
Their product is zero. At $v\ne0$ the indicated bottom and left
blocks have full rank, and at $v=0$ the other blocks have full rank.
Thus they give the exact sequence with bundles
$F|_{C_i}=\mathcal O(-1)^2$ and $G|_{C_i}=\mathcal O(1)^2$.
The bottom-block chart recovers $u/v$ from the plane; moreover
$QS'=B$ near $y=v/u=0$. They are embedded immersed conics.
Set $P=I+yA$ and $C=BPB^{-1}$. Then
\[
 Z_t=yBP^{-1},\qquad S=(I;Z_t)P,
 \qquad Q=C(-Z_t,I),\qquad C Z_t'P=B.
\]
The tensor transition from the polynomial conic frames to graph
frames is $\eta_{\rm graph}=P\eta_{\rm conic}C^{-t}$.
The determinant factor is fixed by the exact identity
\[
 \langle CTP,P^{-1}\eta C^t\rangle
       =\det(C)\langle T,\eta\rangle,
\]
which follows from $C^tJC=(\det C)J$.
The ordinary tooth cup is the constant row $\langle B,-\rangle$
on $F\otimes G=\mathcal O^4$, so its kernel is $\mathcal O^3$.
Its global sections are constant in these frames. The derivative
forced by their value at the node is consequently
\[
 \eta_t'(0)=A\eta_0-\eta_0(BAB^{-1})^t.
\]
Substituting this exact frame derivative in \eqref{invII:eq:E12-node}
gives the two conditions on an ordinary handle-kernel section:
\begin{equation}
 \begin{split}
 \langle B,\eta_h(0)\rangle&=0,\\
 \langle B,\eta_h'(0)\rangle
  &=\langle H,A\eta_h(0)-\eta_h(0)(BAB^{-1})^t\rangle.
 \end{split}                                        \hypertarget{tag-03M1}{}\label{invII:eq:E12-jets}
\end{equation}
The eight parameters $(A,B)$, modulo
$(A,B)\mapsto(rA+tI,rB)$ with $r\ne0$, are the 
six-dimensional through-point conic chart: a polynomial basis
of its tautological $\mathcal O(-1)^2$ has constant blocks $I,0$
and linear blocks $A,B$. The source map is $y=r\widetilde y/(1+t\widetilde y)$.
Multiplying its tautological columns by $1+t\widetilde y$ gives
exactly the displayed matrix action. The two-jet kernel in
\eqref{invII:eq:E12-jets} is unchanged by these reparameterizations.
\cite[Lemma~5.1, pp.~15--16]{Viaclovsky2026Combs} gives the same local deformation space and
a nonempty complex-open set of admissible compact conics.

\item \textbf{Dominance onto the cyclic two-jet quotients.}
The perfect pairing $\langle-,-\rangle$ identifies
$B\mapsto\langle B,-\rangle|_{\ker\langle H,-\rangle}$ with the
quotient of a four-dimensional dual space by the line spanned by
$H$. It gives every first covector, subject to a dense open restriction.
For the second coefficient the following frame map is explicit.
Under $F'=FU$, $G'=GV$ one has
\[
 \begin{gathered}
 T'=V^{-1}TU,\qquad \eta'=U^{-1}\eta V^{-t},\\
 \langle T',\eta'\rangle=(\det V)^{-1}\langle T,\eta\rangle,
 \qquad A'=U^{-1}AU,\quad B'=V^{-1}BU.
 \end{gathered}
\]
Choose an eigenvector matrix $R$ for $H^{-1}B$ and put $U=R,V=HR$.
The derivative becomes $I$, the conic direction becomes
$\operatorname{diag}(\beta_1,\beta_2)$, and the tensor becomes
$R^{-1}\eta H^{-t}R^{-t}$. Work on the nonempty open set
\[
 \beta_1\beta_2(\beta_1-\beta_2)(\beta_1+\beta_2)\ne0.
\]
The two value equations then give
$\eta_0=\operatorname{diag}(u,v)$. For
$A=\left(\begin{smallmatrix}p&q\\r&s\end{smallmatrix}\right)$,
the second coefficient in these frames is
\[
 \langle I,A\eta_0-\eta_0(BAB^{-1})^t\rangle
   =r(1+\beta_2/\beta_1)u-q(1+\beta_1/\beta_2)v.
\]
Its Jacobian in $r,q$ has determinant
$-(\beta_1+\beta_2)^2/(\beta_1\beta_2)$, so the two coefficients
are arbitrary. The scalar part of $A$ contributes zero, as required
by the source reparameterization.
In any fixed local handle frame the resulting quotient is
\[
 \eta\longmapsto\ell(\eta(0))+
 t\bigl(\ell(\eta'(0))+\psi(\eta(0))\bigr)\pmod{t^2},
 \qquad \ell\ne0,
\]
where $\psi$ is arbitrary modulo $\ell$. The derivative of
the chosen handle frame translates $\psi$. This proves dominance
from conics to the four-dimensional cyclic length-two
quotient space, with its complete local frame map.

\item \textbf{An attained balanced open set.}
Let $K'$ be the subsheaf of $K_0$ defined by the twelve quotients
\eqref{invII:eq:E12-jets}. Their parameter space is irreducible.
There is an exact witness for its balanced locus: assign four of
the distinct nodes to each summand of
$K_0=\mathcal O\oplus\mathcal O(-1)^2$, and take the coordinate
quotient modulo the square of the local parameter at each assigned
point. Multiplication by each product of four squared linear
factors gives its kernel inclusion. The resulting bundle is
$\mathcal O(-8)\oplus\mathcal O(-9)^2$, so
\begin{equation}
 H^0(K'(7))=0,\qquad H^1(K'(8))=0.                    \hypertarget{tag-03M2}{}\label{invII:eq:E12-open}
\end{equation}
These conditions define a nonempty Zariski-open subset by
semicontinuity. Step 5 pulls it back to a nonempty open subset of
conic data.

The original handle normal bundle has degrees $(-1,0,0,1)$.
With the three fixed marks, positive modifications distributed as
$(4,3,3,2)$ give a globally generated modified normal bundle with
zero $H^1$. The complete normal-sheaf proof in
\cite[Proposition~5.3, pp.~17--18, and \S9.2, p.~30]{Viaclovsky2026Combs}
shows that this is a second nonempty Zariski-open condition in the
tangent directions. The conic tangent map is dominant, so
both conditions pull back to dense open subsets of the same
irreducible conic chart. Their intersection meets the nonempty
complex-open set of admissible compact conics.
For these choices the cited fixed-marking deformation theorem
gives an embedded smoothing with a direction nonzero at every
node. The curve has genus zero and degree one over $B$, so each
smooth curve is a section. It fixes all three marked
period planes and has the root and compact filling of Step 2.

\item \textbf{Collapse, cohomology and the smooth bundle.}
Let $K_C$ denote the central fibre of the locally free relative
kernel. Its handle restriction maps into $K_0$ by the first
elementary modification at each node. On a tooth its image in the
ordinary $\mathcal O^3$ kernel is cut out by the additional nonzero
value functional $\langle H,-\rangle$, and hence
\[
 K_C|_{C_i}=\mathcal O(-1)\oplus\mathcal O^2.
\]
This follows by choosing that functional as the first coordinate:
the first coordinate vanishes at the node and the other two are
constant. Its $H^1$ vanishes, and its two global sections evaluate
injectively into the three-dimensional relative-kernel node fibre.
Their image retains exactly the derivative constraint of Step 3.

Collapse the teeth by the projection $\pi:C\to B$.
The normalization sequence for $K_C$, the zero tooth $H^1$ and the
surjective local evaluation of the handle onto each whole node
fibre give $R^1\pi_*K_C=0$. A handle section extends uniquely over
each tooth precisely by the two equations in
\eqref{invII:eq:E12-jets}, including the sufficient mixed condition.
Thus restriction gives the exact sheaf identification
$\pi_*K_C=K'$. For the pullback line from the base this gives
\[
 H^j(C,K_C\otimes\pi^*\mathcal O_B(k))
       =H^j(B,K'(k)),\qquad j=0,1.
\]
The projection of the embedded family supplies this single line
throughout the smoothing: it is trivial on the central teeth and
is the required $\mathcal O_B(k)$ on each smooth section.
Flatness and semicontinuity therefore transport
\eqref{invII:eq:E12-open} to all sufficiently small noncentral
members of the chosen smoothing.
The central Euler characteristic is
\[
 (1-12)+12\cdot2-12\cdot3=-23,
\]
so its rank-three degree is $-26$. Equivalently, the smooth
surjective cup row has source degree $-2$ and target degree $24$.

These two vanishings determine the bundle by
explicit maps. Riemann--Roch gives $h^0(K_{12}(8))=1$. Its section
is nowhere zero, since a zero would give a section of $K_{12}(7)$.
The quotient of $K_{12}(8)$ by this copy of $\mathcal O$ has rank
two, degree $-2$ and zero $H^0,H^1$. Evaluation of its twist by
one is onto at every point; Riemann--Roch gives two source
dimensions, so that evaluation is an isomorphism. The quotient
is $\mathcal O(-1)^2$, and the extension by $\mathcal O$ splits
because $H^1(\mathcal O(1))^2=0$. Twisting back proves the stated
kernel isomorphism. The two vanishings and surjectivity are open
in the nearby fixed-marking section family, giving the
nonempty open family in the statement. Formula
\eqref{invII:eq:E.5} supplies its two Hodge numbers.

\end{enumerate}
\end{proof}

\section{Uniform non-torsion on every full fibre}\hypertarget{tag-034H}{}\label{sec:II:uniform-nontorsion}

\begin{definition}[Normalized Picard parameter]\hypertarget{tag-P279}{}\label{def:II:normalized-picard-parameter}
\StatementRecord{def:II:normalized-picard-parameter}{proved}{P279}{\SeriesRef{inv:A1.0}, \SeriesRef{inv:C3.1}}
Let $f:X\to B=\mathbb P^1$ be a member of the standard $(3,4)$
family in \SeriesRef{inv:A1.0}, with full special fibres
$3S_1,4S_2,W$ and its fixed period marking.
Use $M_u=\exp(ug)$ for the Picard isomorphism of
\SeriesRef{inv:C3.1}. Let $u_0$ be the unique parameter for
$M_{u_0}=f^*\mathcal O_B(1)$ and set
\[
 \phi(z)=u_0z/12,\qquad \phi^{-1}(u)=12u/u_0,
 \qquad L_z=M_{\phi(z)}.
\]
Thus
\begin{equation}
 L_{12}=f^*\mathcal O_B(1).
 \hypertarget{tag-01TJ}{}\label{inv:wb:eq:9.19}
\end{equation}
For a full fibre $X_b$, write
$\iota_b(z)=(u_0z/12)g|_{X_b}$; its exponential is $L_z|_{X_b}$.
\end{definition}
\begin{proof}
Connected fibres and projection give
$h^0(X,f^*\mathcal O_B(1))=h^0(B,\mathcal O_B(1))=2$, whereas
$h^0(X,\mathcal O_X)=1$. Thus $f^*\mathcal O_B(1)$ is nontrivial,
so $u_0\ne0$ and the displayed scalar maps are inverse isomorphisms.
Under the coherent Leray and base-change maps of \SeriesRef{inv:A1.0},
the generator $g$ restricts to the nonzero invariant class
$\gamma^{0,1}$ on every full fibre. Multiplication by $u_0/12$
preserves that image line and makes every $\iota_b$ injective.
\end{proof}

\begin{theorem}[Uniform non-torsion]
\hypertarget{tag-03DQ}{}\label{wb:thm:uniform-nontorsion-audit}
\hypertarget{tag-01TQ}{}\label{inv:wb:uniform-nontorsion-provenance}
\StatementRecord{wb:thm:uniform-nontorsion-audit}{proved}{03DQ}{\SeriesRef{inv:wb:eq:9.19}, \SeriesRef{inv:wb:eq:9.23}, \SeriesRef{inv:wb:eq:9.22}, \SeriesRef{inv:wb:eq:9.24}, \SeriesRef{inv:wb:eq:9.25}}

For the fibration of \cref{def:II:normalized-picard-parameter}, there is a line bundle \(L\in\operatorname{Pic}(X)\) such that
\(L|_{X_b}\) is non-torsion for every scheme-theoretic fibre
\(X_b=f^*(b)\), and \(L|_{S_j}\) is non-torsion for \(j=1,2\).
In the normalization of \cref{def:II:normalized-picard-parameter}, every \(z\) outside the union of \(\mathbb R\) with a
countable set has these properties.
\end{theorem}



\begin{proof}

\begin{enumerate}[label=\arabic*.,leftmargin=*]
\item Fix a smooth fibre \(F_q\).  Put
\begin{equation}
 \Lambda_q=\iota_q^{-1}\!\left(
 \operatorname{im}[H^1(F_q;\mathbb Z)\to
 H^1(F_q,\mathcal O_{F_q})]\right)\subset\mathbb C.
\hypertarget{tag-01TK}{}\label{inv:wb:eq:9.20}
\end{equation}
The integral image is a lattice in the underlying real vector space of
\(H^1(F_q,\mathcal O_{F_q})\), so \(\Lambda_q\) is discrete.  It contains
\(12\), because \eqref{inv:wb:eq:9.19} is trivial on every fibre.  The exponential
sequence on \(F_q\) shows
\begin{equation}
 L_z|_{F_q}\text{ is torsion}
 \quad\Longleftrightarrow\quad z\in\mathbb Q\Lambda_q.
\hypertarget{tag-01TL}{}\label{inv:wb:eq:9.21}
\end{equation}
If \(\operatorname{rk}\Lambda_q=1\), its nonzero element \(12\) forces
\(\Lambda_q\subset\mathbb R\); all its bad parameters are real.

\item Suppose \(\operatorname{rk}\Lambda_q=2\).  By
\SeriesRef{inv:A1.0}, the image of \(\iota_q\) is the line
generated by the image of \(\gamma\).  Thus there is an integral class
\begin{equation}
 \lambda=a\gamma+bu+cw+d\delta\notin\mathbb Q\gamma
\hypertarget{tag-01TM}{}\label{inv:wb:eq:9.22}
\end{equation}
whose \((0,1)\)-image lies on that line.  On the covering family, the four
marked lattice columns are
\begin{equation}
 \pi_{\widehat\gamma}=(6\mu,\beta_c),\quad
 \pi_{\widehat u}=(\tau,\mu),\quad
 \pi_{\widehat w}=(1,0),\quad
 \pi_{\widehat\delta}=(0,1).
\hypertarget{tag-01TN}{}\label{inv:wb:eq:9.23}
\end{equation}
Every real cohomology class has a unique expression
\[
 2\operatorname{Re}(\bar\xi\,d\zeta_1+
                       \bar\upsilon\,d\zeta_2).
\]
Evaluating it on \eqref{inv:wb:eq:9.23} gives
\begin{equation}
\begin{aligned}
2\operatorname{Re}\bar\xi&=c,&
2\operatorname{Re}\bar\upsilon&=d,\\
2\operatorname{Re}(\bar\xi\tau+\bar\upsilon\mu)&=b,&
2\operatorname{Re}(6\bar\xi\mu+\bar\upsilon\beta_c)&=a.
\end{aligned}
\hypertarget{tag-01TO}{}\label{inv:wb:eq:9.24}
\end{equation}
For \(\gamma\), the first two right sides vanish.  The third equation then
shows that its \((0,1)\)-line is represented, up to a nonzero scalar, by
\((\operatorname{Im}\mu,-\operatorname{Im}\tau)\); the fourth equation
only fixes the scalar.  Write
\(\bar\xi=c/2+ir\), \(\bar\upsilon=d/2+is\).  The condition that the class
in \eqref{inv:wb:eq:9.22} have \((0,1)\)-image on the \(\gamma\)-line is
\[
 \bar\xi\operatorname{Im}\tau+
 \bar\upsilon\operatorname{Im}\mu=0.
\]
Its real part, together with the third equation of \eqref{inv:wb:eq:9.24}, is equivalent
to the single holomorphic equation
\begin{equation}
 c\tau+d\mu-b=0.
\hypertarget{tag-01TP}{}\label{inv:wb:eq:9.25}
\end{equation}
Neither \(a\) nor \(\beta_c\) occurs, as adding a multiple of \(\gamma\)
must not change the condition.

\item The functions \(1,\tau,\mu\) are linearly independent.  Applying the cusp
transformation \((\tau,\mu)\mapsto(\tau-1,\mu)\) to a constant linear
relation first kills the coefficient of \(\tau\).  If the remaining
relation made \(\mu\) constant, the order-four law
\(\mu\circ g_2=1+\mu/\tau\), with nonconstant \(\tau\), would force an
impossible identity.  Hence all coefficients vanish.

\item For every integral \(\lambda\notin\mathbb Q\gamma\), equation \eqref{inv:wb:eq:9.25} is
therefore a nonzero holomorphic function on the covering Riemann surface.
Its zero set is discrete and countable.  There are countably many
\(\lambda\), so their union, and its image in \(B^\circ\), are countable.
This image contains every \(q\) for which
\(\operatorname{rk}\Lambda_q=2\).  Each such fibre contributes the
countable bad set \(\mathbb Q\Lambda_q\); all rank-one fibres together
contribute only parameters in \(\mathbb R\).  Thus the smooth fibres
exclude only \(\mathbb R\) together with a countable set.

\item There remain the three scheme-theoretic special fibres.  By
\SeriesRef{inv:A1.0} and \cref{def:II:normalized-picard-parameter}, restriction
\[
 \mathbb C\longrightarrow H^1(X_b,\mathcal O_{X_b})
\]
is complex-linear and injective for each of them.  The exponential sequence
and countability of \(H^1(X_b;\mathbb Z)\) show that the parameters for
which \(L_z|_{X_b}\) is torsion form a countable set.

\item Finally consider the reductions \(S_j\).  The restriction
\(\varphi_j:\mathbb C\to H^1(S_j,\mathcal O_{S_j})\) is complex-linear.
It is nonzero.  Write \(\mathcal O_X(S_j)=L_{z_j}\).  Since
\[
 \mathcal O_X(m_jS_j)=f^*\mathcal O_B(1)=L_{12},
\]
the additive parameter gives \(z_j=12/m_j\), while
the root model of \SeriesRef{fam:src:A.1.2} identifies
\(\mathcal N_j:=\mathcal O_X(S_j)|_{S_j}\) with the normal line of
the reduced fibre. Its connected degree-\(m_j\) cover
\(\widehat S_j\to S_j\) has free deck generator \(a_j\). In the
root coordinate \(s_j\), the generator acts by
\(s_j\mapsto\zeta_j s_j\), with \(\zeta_j\) primitive of order
\(m_j\). The derivative therefore identifies the line with
\[
 \mathcal N_j\simeq
 (\widehat S_j\times\mathbb C)/
       ((x,v)\sim(a_jx,\zeta_jv)).
\]
A trivializing section of \(\mathcal N_j^{\otimes r}\) pulls back
to a nowhere-zero holomorphic function \(h\) on the compact connected
torus \(\widehat S_j\), with \(h(a_jx)=\zeta_j^r h(x)\).
The function is constant, so \(m_j\mid r\). Conversely the constant
section descends for \(r=m_j\); hence the order is exactly \(m_j\).
In particular,
\[
 L_{12/m_j}|_{S_j}=\mathcal N_j\ne\mathcal O_{S_j}.
\]  If \(\varphi_j\) were zero, the exponential sequence and
the vanishing of the restricted first Chern class would make every
\(L_z|_{S_j}\) trivial, a contradiction.  Thus \(\varphi_j\) is
injective, and the same exponential-sequence argument makes its torsion
parameters countable.  Removing these two additional countable sets proves
the theorem.

\end{enumerate}
\end{proof}
\section{Analytic duality and base change}\hypertarget{tag-P307}{}\label{sec:II:analytic-duality}

\begin{lemma}[Flatness and Cartier fibres over a curve]
\hypertarget{tag-03I8}{}\label{wbII:lem:curve-flatness}
\StatementRecord{wbII:lem:curve-flatness}{proved}{03I8}{\cite[Tag~0AUW]{StacksProject}}

Let $f:X\to B$ be a nonconstant holomorphic map from a connected
complex manifold to a Riemann surface.  Then $f$ is flat, every holomorphic
vector bundle on $X$ is $f$-flat, and every full scheme-theoretic fibre is an
effective Cartier divisor of pure codimension one.  If $X$ is compact, then
$f$ is proper.
\end{lemma}



\begin{proof}

\begin{enumerate}[label=\arabic*.,leftmargin=*]
\item Fix $x\in X$, $c=f(x)$, and put
$R=\mathcal O_{B,c}$, $S=\mathcal O_{X,x}$.  The ring $R$ is a discrete
valuation ring and $S$ is a regular local domain.  The homomorphism
$R\to S$ is injective.  Indeed, if a nonzero germ $h\in R$ mapped to zero,
$h\circ f$ would vanish on a neighbourhood of $x$.  The zero set of a
nonzero one-variable holomorphic germ is isolated, so $f$ would be constant
on that neighbourhood and therefore, by the identity theorem on connected
$X$, constant everywhere.

\item Thus every nonzero element of $R$ acts on the domain $S$ as a nonzero
element and hence as a nonzerodivisor.  Therefore $S$ is torsion-free as an
$R$-module.  A module over a discrete valuation ring is flat if and only if
it is torsion-free \cite[Tag~0AUW]{StacksProject}; hence $S$ is $R$-flat.
This holds at every point, so $f$ is flat.  A vector bundle is locally a
finite free $\mathcal O_X$-module, and a finite direct sum of $R$-flat
modules is flat; hence it is $f$-flat.  If $t$ is a local parameter at
$c$, then $f^*t$ is a nonzero element of the domain $S$ and thus a
nonzerodivisor.  The full scheme-theoretic fibre is therefore locally cut
out by one nonzerodivisor, so it is an effective Cartier divisor of pure
codimension one.  Finally, when $X$ is compact, the underlying continuous
map to the Hausdorff surface $B$ is proper.

\end{enumerate}
\end{proof}

\begin{cited}[Duality for compact Cartier divisors and full fibres]
\hypertarget{tag-03IY}{}\label{wbII:thm:cartier-duality}
\StatementRecord{wbII:thm:cartier-duality}{cited}{03IY}{\cite[\S4, Theorem~2, pp.~82--84; \S7, Proposition~2 and complete-intersection calculation, p.~91]{RR70push}}
\hypertarget{tag-03IJ}{}\label{wbII:prop:cartier-duality-absolute}
\StatementRecord{wbII:prop:cartier-duality-absolute}{cited}{03IJ}{\SeriesRef{wbII:thm:cartier-duality}}
Let $X$ be a smooth complex threefold, let $i:W\hookrightarrow X$
be a compact effective Cartier divisor, and let $E$ be a holomorphic
vector bundle on $X$. Put $F=E|_W$ and
\[
 \omega_W=(\omega_X\otimes\mathcal O_X(W))|_W.
\]
For $q=0,1,2$ the trace pairing gives natural isomorphisms
\begin{equation}
 H^q(W,F)^\vee\simeq H^{2-q}(W,F^\vee\otimes\omega_W),
 \hypertarget{tag-03EB}{}\label{wbII:eq:W:5.11}
\end{equation}
and $H^q(W,F)=0$ for $q>2$.

For the full scheme-theoretic fibre $W=f^{-1}(b)$ of a holomorphic
map $f:X\to B$ to a smooth curve, choose a nonzero vector in
$\mathcal O_B(b)|_b$. It identifies $\omega_W$ with $\omega_X|_W$
and gives a functorial perfect pairing
\[
 H^2(W,E|_W)\times
 H^0(W,E^\vee|_W\otimes\omega_X|_W)\longrightarrow\mathbb C,
\]
or, equivalently,
\begin{equation}
 H^2(W,E|_W)^\vee\simeq
 H^0(W,E^\vee|_W\otimes\omega_X|_W).
 \hypertarget{tag-03EC}{}\label{wbII:eq:W:5.15a}
\end{equation}
A different nonzero vector rescales this last pairing by a nonzero
scalar. All statements use the full analytic structure of $W$.
\end{cited}

\begin{proof}
Apply \cite[\S7, Proposition~2 and complete-intersection calculation,
p.~91]{RR70push} to the codimension-one inclusion $i$. It identifies
the dualizing complex with $\omega_W[2]$. The trace duality of
\cite[\S4, Theorem~2, pp.~82--84]{RR70push}, with the locally free
first argument $F$, then identifies
\[
 H_c^q(W,F)^\vee
 =\operatorname{Ext}^{-q}_W(F,\omega_W[2])
 =H^{2-q}(W,F^\vee\otimes\omega_W).
\]
Compactness of $W$ gives $H_c^q=H^q$ and finite-dimensional coherent
cohomology. The formula is therefore a perfect pairing; negative
Ext degree gives the stated vanishing for $q>2$. Precomposition in
$F$ gives its variance. For a full fibre,
$\mathcal O_X(W)|_W=f|_W^*(\mathcal O_B(b)|_b)$; applying the chosen
vector to this constant line is exactly the asserted adjunction
identification and yields \eqref{wbII:eq:W:5.15a}.
\end{proof}

\begin{cited}[Analytic base change in degree two]
\hypertarget{tag-03II}{}\label{wbII:prop:analytic-base-change}
\StatementRecord{wbII:prop:analytic-base-change}{cited}{03II}{\cite[Chapter~III, \S3, Theorem~3.4 and Corollary~3.5, pp.~116--121]{BS76push}}

Let $f:X\to B$ be proper between complex spaces, let
$\mathcal F$ be coherent and $f$-flat, and fix $c\in B$.  If
\begin{equation}
H^3(X_c,\mathcal F|_{X_c})=0,
\hypertarget{tag-03ED}{}\label{wbII:eq:W:5.16}
\end{equation}
then the canonical fibre map
\begin{equation}
\beta_c^2:
(R^2f_*\mathcal F)_c\otimes_{\mathcal O_{B,c}}\mathbb C(c)
\longrightarrow H^2(X_c,\mathcal F|_{X_c})
\hypertarget{tag-03EE}{}\label{wbII:eq:W:5.17}
\end{equation}
is an isomorphism.
\end{cited}



\begin{proof}
Apply B\u{a}nic\u{a}--St\u{a}n\u{a}\c{s}il\u{a},
Chapter~III, \S3, Corollary~3.5, with $q=3$
\cite[Chapter~III, \S3, Corollary~3.5, pp.~120--121]{BS76push}.  Its hypotheses are precisely
properness, coherence, and flatness.  Condition \eqref{wbII:eq:W:5.16} first forces the
degree-three direct image to vanish near $c$.  The corollary then applies the
equivalent base-change conditions of its preceding Theorem~3.4 in degree
$q-1=2$ and states that the canonical map \eqref{wbII:eq:W:5.17} is bijective.  The top-degree application also appears in
\cite[proof of Lemma~1.5, part~(c), p.~80]{CDP98}.

\end{proof}

\begin{corollary}[Top-degree base change for surface fibres]
\hypertarget{tag-03DV}{}\label{wbII:cor:surface-top-base-change}
\StatementRecord{wbII:cor:surface-top-base-change}{proved}{03DV}{\SeriesRef{wbII:eq:W:5.17-U}, \SeriesRef{wbII:prop:analytic-base-change}; \cite[Chapter~III, \S4, Theorem~4.1, p.~124]{BS76push}}

Let $f:X\to B$ be a proper flat holomorphic map to a Riemann
surface whose full scheme-theoretic fibres have pure complex dimension two,
and let $E$ be a holomorphic vector bundle on $X$.  Then
\begin{equation}
 R^qf_*E=0\quad(q>2)
 \hypertarget{tag-03FU}{}\label{wbII:eq:W:5.17-V}
\end{equation}
and, for every $b\in B$, including a point with nonreduced fibre, the
canonical map
\begin{equation}
 (R^2f_*E)_b\otimes_{\mathcal O_{B,b}}\mathbb C(b)
 \xrightarrow{\ \sim\ }H^2(X_b,E|_{X_b})
 \hypertarget{tag-03FS}{}\label{wbII:eq:W:5.17-T}
\end{equation}
is an isomorphism.
\end{corollary}



\begin{proof}
On a sufficiently small neighbourhood $U$ of $b$, the finite-complex theorem
used in analytic cohomology and base change represents $Rf_*E|_U$ by a
bounded complex $K^\bullet$ of finite free $\mathcal O_U$-modules
\cite[Chapter~III, \S4, Theorem~4.1, p.~124]{BS76push}.  Flatness of
$E$ over $B$ identifies its derived fibre with the ordinary restriction to
the full scheme-theoretic fibre.  Over the discrete valuation ring
$A=\mathcal O_{B,b}$, the fibre map of
\cref{wbII:prop:analytic-base-change} occurs in the universal-coefficient sequence
\begin{equation}
0\longrightarrow H^q(K^\bullet)_b\otimes_A\mathbb C(b)
 \longrightarrow H^q(K^\bullet\otimes_A\mathbb C(b))
 \longrightarrow
 \operatorname{Tor}_1^A(H^{q+1}(K^\bullet)_b,\mathbb C(b))
 \longrightarrow0.
\hypertarget{tag-03FT}{}\label{wbII:eq:W:5.17-U}
\end{equation}
If some $R^qf_*E$ with $q>2$ were nonzero, choose the largest such $q$ and a
point of its support.  Nakayama's lemma makes the left term of \eqref{wbII:eq:W:5.17-U}
nonzero there, while maximality kills the Tor term.  This would give nonzero
coherent cohomology in degree $q>2$ on a compact two-dimensional complex
space, a contradiction.  Thus \eqref{wbII:eq:W:5.17-V} holds.  For $q=2$, the only
obstruction in \eqref{wbII:eq:W:5.17-U} is
$\operatorname{Tor}_1^A((R^3f_*E)_b,\mathbb C(b))$, which is zero by
\eqref{wbII:eq:W:5.17-V}, proving \eqref{wbII:eq:W:5.17-T}.

\end{proof}

\section{Conormal and relative cotangent maps}\hypertarget{tag-P308}{}\label{sec:II:conormal-relative-maps}

\begin{lemma}[Identification of the torsion subsheaf]
\hypertarget{tag-03IF}{}\label{wbII:lem:torsion-identification}
\StatementRecord{wbII:lem:torsion-identification}{proved}{03IF}{the definitions and constructions stated in the proof}

If $S$ is a reduced effective Cartier divisor in a smooth complex manifold,
then
\[
\ker\bigl(\Omega^1_S\longrightarrow j_*\Omega^1_{S_{\rm reg}}\bigr)
  =\operatorname{Tor}(\Omega^1_S).
\]
\end{lemma}



\begin{proof}
Every irreducible component of a reduced hypersurface is generically smooth,
so $S_{\rm reg}$ is dense and its complement is a proper analytic subset.
If a local section of $\Omega^1_S$ vanishes on $S_{\rm reg}$, its support avoids
all minimal primes of the reduced local ring.  Prime avoidance therefore gives
an element of its annihilator outside every minimal prime; that element is a
nonzerodivisor, so the section is torsion.  Conversely,
$\Omega^1_S|_{S_{\rm reg}}$ is locally free, hence has no torsion, and every
torsion section restricts to zero there.

\end{proof}

\begin{theorem}[Conductor-supported ambient differential]
\hypertarget{tag-03IK}{}\label{wbII:prop:conductor}
\StatementRecord{wbII:prop:conductor}{proved}{03IK}{\SeriesRef{cdp:st:A.1.2}, \SeriesRef{cdp:st:A.2.2}, \SeriesRef{wbII:lem:torsion-identification}; \cite[Lemma~10.3, pp.~87--88]{AC26}}
\label{wbII:prop:conductor-differential}
Let $g:\mathcal X\to\Delta$ be a holomorphic map from a smooth complex
threefold to a disc, and let $S=g^{-1}(0)$ be nonempty, reduced and normal crossing.
Thus at each point of $S$ there are coordinates $(x,y,z)$ in which
\begin{equation}
g=ux,\qquad g=uxy,\qquad\text{or}\qquad g=uxyz
\hypertarget{tag-03HS}{}\label{wbII:eq:W:5.3}
\end{equation}
for a holomorphic unit $u$. Let $\eta:\widetilde S\to S$ be the normalization,
$D=S_{\mathrm{sing,red}}$, and $\widetilde D=\eta^{-1}(D)_{\mathrm{red}}$.
Let $A$ be a line bundle on $S$ with
$\eta^*A\simeq\mathcal O_{\widetilde S}$, and choose a nowhere-vanishing
section $e$ of $\eta^*A$.  There is a unique nonzero
\begin{equation}
s\in H^0(S,\Omega^1_{\mathcal X}|_S\otimes A)
\hypertarget{tag-03HZ}{}\label{wbII:eq:W:5.6}
\end{equation}
such that
\begin{equation}
\eta^*s=\eta^*(dg|_S)\otimes e,\qquad s|_D=0.
\hypertarget{tag-03I0}{}\label{wbII:eq:W:5.7}
\end{equation}
Its image $\rho(s)$ in $\Omega^1_S\otimes A$ is supported on $D$ and
therefore torsion; its image in
$\widetilde\Omega^1_S\otimes A$ is zero.  Moreover:
\begin{enumerate}[label=(\alph*)]
\item $\rho(s)=0$ if and only if $e=\eta^*\sigma$ for a nowhere-vanishing
$\sigma\in H^0(S,A)$;
\item if $S$ is compact and connected and $A\not\simeq\mathcal O_S$, then
$H^0(S,A)=0$ and $s$ does not come from the conormal term in \eqref{wbII:eq:W:5.4}.
\end{enumerate}
\end{theorem}



\begin{proof}

\begin{enumerate}[label=\arabic*.,leftmargin=*]
\item The map $\mathcal O_S\to N^*_{S/\mathcal X}$ given by
$h\mapsto h[g]$ is an isomorphism because the fibre ideal is the principal
ideal $(g)$ and $g$ is a nonzerodivisor. Under this isomorphism, the
conormal map of \cref{cdp:st:A.1.1} sends $h$ to $h\,dg|_S$.
Indeed, $d(hg)|_S=h\,dg|_S$. Hence that proposition gives
\begin{equation}
0\longrightarrow\mathcal O_S
\xrightarrow{\,h\mapsto h\,dg\,}
\Omega^1_{\mathcal X}|_S\longrightarrow\Omega^1_S\longrightarrow0
\hypertarget{tag-03HX}{}\label{wbII:eq:W:5.4}
\end{equation}
By \cref{cdp:st:A.1.2}, the normal-crossing normalization gives
the canonical isomorphism
\[
 \eta_*\eta^*N^*_{S/\mathcal X}
 \xrightarrow{\ \sim\ }
 K=\ker(\Omega^1_{\mathcal X}|_S\to\widetilde\Omega^1_S).
\]
In the coordinates \eqref{wbII:eq:W:5.3}, it sends the tuple
$[g]\otimes(a_i)$ to
$u\sum_i a_i(\prod_{j\ne i}x_j)\,dx_i$.
The conductor ideals are $(x,y)$ in the double model and
$(xy,xz,yz)$ in the triple model, with radical pullbacks to the
smooth branches, by the local conductor calculation in
\cite[Lemma~10.2, proof, p.~87]{AC26}. Thus the conductor sequence of
\cref{cdp:st:A.2.2} specializes to
\begin{equation}
0\longrightarrow\mathcal O_S
\xrightarrow{\alpha}
\eta_*\mathcal O_{\widetilde S}\oplus\mathcal O_D
\xrightarrow{\beta}
\eta_*\mathcal O_{\widetilde D}
\longrightarrow0,
\hypertarget{tag-03HY}{}\label{wbII:eq:W:5.5}
\end{equation}
where
\[
 \alpha(h)=(\eta^*h,h|_D),\qquad
 \beta(\widetilde h,k)=\widetilde h|_{\widetilde D}-\eta^*k.
\]
Tensor the canonical isomorphism by $A$ and use the finite
projection formula. The normalized section $[g]\otimes e$ maps to
a unique section $s$ of $K\otimes A$. Its displayed branch formula
gives $\eta^*s=\eta^*(dg|_S)\otimes e$ and $s|_D=0$; it is nonzero
on the smooth locus because $e$ and the conormal vector are nonzero.
Uniqueness also follows from \eqref{wbII:eq:W:5.5} tensored by
$\Omega^1_{\mathcal X}|_S\otimes A$: its two restrictions are the
compatible pair $(\eta^*(dg|_S)\otimes e,0)$.
For proper $g$ this is the construction in
\cite[Lemma~10.3, pp.~87--88]{AC26}.

\item On $S\setminus D$, the normalization is an isomorphism and
$s=dg|_S\otimes e$.  By \eqref{wbII:eq:W:5.4} its intrinsic image vanishes there.  Thus
$\rho(s)$ is supported on $D$ and lies in the torsion subsheaf; its image
modulo torsion is zero.  Here, for $j:S_{\rm reg}\hookrightarrow S$,
$\widetilde\Omega^1_S=\Omega^1_S/\operatorname{Tor}(\Omega^1_S)$
is identified with the image in $j_*\Omega^1_{S_{\rm reg}}$ by
$[\omega]\mapsto\omega|_{S_{\rm reg}}$; its inverse sends a local
restriction to the class of any lift.  The equality of the restriction
kernel and the torsion subsheaf is
\cref{wbII:lem:torsion-identification}.

\item If $\rho(s)=0$, exactness of \eqref{wbII:eq:W:5.4} after tensoring by $A$ gives
$s=dg|_S\otimes\sigma$ for a section $\sigma$ of $A$.  Pulling back and
comparing with \eqref{wbII:eq:W:5.7} on the dense open set where $dg\ne0$ gives
$\eta^*\sigma=e$ there, hence everywhere by holomorphic continuation.
Since $e$ never vanishes and $\eta$ is surjective, $\sigma$ never vanishes
and trivializes $A$.  The converse is immediate.

\item Finally suppose $S$ is compact and connected and $\sigma\in H^0(S,A)$.  On each compact
connected component of $\widetilde S$, write
$\eta^*\sigma=c_i e$; the holomorphic function $c_i$ is constant.  If one
$c_i$ is zero and another is nonzero, connectedness of the finite union of
irreducible components of $S$ provides two such components meeting at a
point.  The value of $\sigma$ at that point would be both zero and nonzero,
a contradiction.  Hence either every $c_i=0$, so $\sigma=0$, or every
$c_i\ne0$, so $\sigma$ is nowhere zero and $A$ is trivial.  When $A$ is
nontrivial only the first case occurs.  Since
$N^*_{S/\mathcal X}\simeq\mathcal O_S$, the conormal global-section space is
$H^0(S,A)=0$, while $s\ne0$.

\end{enumerate}
\end{proof}

\begin{corollary}[Nonzero cotangent connecting kernel]
\hypertarget{tag-03DU}{}\label{wbII:cor:normal-crossing-kernel}
\StatementRecord{wbII:cor:normal-crossing-kernel}{proved}{03DU}{\SeriesRef{wbII:prop:conductor}, \SeriesRef{cdp:st:A.1.1}}

If, in addition, $H^0(S,N^*_{S/\mathcal X}\otimes A)=0$ and
$H^0(S,\widetilde\Omega^1_S\otimes A)=0$, then the section in
\cref{wbII:prop:conductor} proves that $\delta_A$ in
the section-kernel isomorphism of \cref{cdp:st:A.1.1} is not injective.
\end{corollary}



\begin{proof}
The section is nonzero and maps to zero in the quotient
$\widetilde\Omega^1_S=\Omega^1_S/\mathscr T_S$, so it
lies in the left-hand side of the section-kernel isomorphism of \cref{cdp:st:A.1.1}.
Under that isomorphism it is sent to its intrinsic image $\rho(s)$.
The hypothesis $H^0(S,N^*_{S/\mathcal X}\otimes A)=0$ makes this map injective,
and the ambient lift $s$ gives $\delta_A(\rho(s))=0$.  \Cref{cdp:st:A.1.1}
identifies that group with $\ker\delta_A$, which is therefore nonzero.

\end{proof}

\begin{proposition}[Crossing and multiple-fibre modules]
\hypertarget{tag-03IP}{}\label{wbII:prop:relative-local}
\StatementRecord{wbII:prop:relative-local}{proved}{03IP}{the definitions and constructions stated in the proof}

The symbols $R_2$ and $R_3$ denote copies of the same regular local ring,
with different maps from the one-dimensional base:
let $R_2=\C\{x,y,z\}$ with base coordinate $t=xy$, and let
$R_3=\C\{x,y,z\}$ with base coordinate $t=xyz$.  Then
\begin{equation}
\begin{aligned}
 M_2&:=\Omega^1_{R_2/\C\{t\}}
 =\frac{R_2\,dx\oplus R_2\,dy\oplus R_2\,dz}
        {R_2\,(y\,dx+x\,dy)},\\[2pt]
 M_3&:=\Omega^1_{R_3/\C\{t\}}
 =\frac{R_3\,dx\oplus R_3\,dy\oplus R_3\,dz}
        {R_3\,(yz\,dx+xz\,dy+xy\,dz)}.
\end{aligned}
\hypertarget{tag-03E6}{}\label{wbII:eq:CN:6.1}
\end{equation}
Both modules are torsion-free, but neither is locally free at the origin.
They embed in free modules by
\begin{equation}
\begin{aligned}
 M_2&\hookrightarrow R_2e\oplus R_2dz,
 &dx&\longmapsto xe,&dy&\longmapsto-ye,&dz&\longmapsto dz,\\
 M_3&\hookrightarrow R_3e_1\oplus R_3e_2,
 &dx&\longmapsto xe_1,&dy&\longmapsto ye_2,&dz&\longmapsto-z(e_1+e_2),
\end{aligned}
\hypertarget{tag-03E7}{}\label{wbII:eq:CN:6.2}
\end{equation}
with images
\begin{equation}
 R_2dz\oplus (x,y)e,
 \qquad
 R_3xe_1+R_3ye_2+R_3z(e_1+e_2).
\hypertarget{tag-03E8}{}\label{wbII:eq:CN:6.3}
\end{equation}
If $R_m=\C\{s,y_1,y_2\}$ and the base coordinate is $t=s^m$ with
$m\ge2$, then
\begin{equation}
\Omega^1_{R_m/\C\{t\}}
\simeq R_m\,dy_1\oplus R_m\,dy_2
       \oplus R_m/(s^{m-1})\,ds,
\hypertarget{tag-03E9}{}\label{wbII:eq:CN:6.4}
\end{equation}
and the final summand is exactly its torsion submodule.
\end{proposition}



\begin{proof}

\begin{enumerate}[label=\arabic*.,leftmargin=*]
\item Differentiating $t=xy$ and $t=xyz$ gives the presentations in \eqref{wbII:eq:CN:6.1}.  For
$M_2$, the first map in \eqref{wbII:eq:CN:6.2} sends $(a,b,c)$ to $(ax-by,c)$.  If this is
zero, then $c=0$ and $ax=by$.  Since $x$ and $y$ are relatively prime in the
UFD $R_2$, there is an $h$ with $b=xh$ and $a=yh$; hence its kernel is exactly
the displayed relation.  For $M_3$, the second map sends $(a,b,c)$ to
$(ax-cz,by-cz)$.  If it is zero, relative primeness first gives $x\mid c$,
then $y\mid(c/x)$; writing $c=xyh$ gives
$(a,b,c)=h(yz,xz,xy)$, again exactly the defining relation.  Thus both maps
are injective, proving torsion-freeness and \eqref{wbII:eq:CN:6.3}.

\item At the maximal ideal $(x,y,z)$ each defining relation has all coefficients in
that maximal ideal.  Consequently the quotient of either module by the
maximal ideal has dimension $3$ over $\C$, while its generic rank is $2$;
a locally free module of rank $2$ would need only $2$ generators.  Finally, for $t=s^m$ one has
$dt=m s^{m-1}ds=0$ relatively.  The integer $m$ is a unit in $\C$, so the
decomposition \eqref{wbII:eq:CN:6.4} follows.  Its last summand is annihilated by
$s^{m-1}$, while the other two summands are free; after inverting $s$ it
vanishes.  It is therefore exactly the torsion submodule.

\end{enumerate}
\end{proof}

\begin{remark}[Restriction and finite \texorpdfstring{\'etale}{etale} base change]
\hypertarget{tag-03IR}{}\label{wbII:rem:relative-local-scope}
\StatementRecord{wbII:rem:relative-local-scope}{proved}{03IR}{\SeriesRef{wbII:prop:relative-local}, \SeriesRef{cdp:st:B1}}
For $S=\{t=0\}$ in the two normal-crossing charts of
\cref{wbII:prop:relative-local}, the natural base-change map is
\[
 M_i\otimes_{R_i}R_i/(t)\xrightarrow{\ \sim\ }\Omega^1_S,
 \qquad d x_j\otimes1\longmapsto d(x_j|_S).
\]
Its inverse sends each intrinsic coordinate differential to the
same generator in the displayed quotient: both modules have the
presentation with relations $t=0$ and $dt=0$.
In the node $xy=0$, the image of $x\,dy=-y\,dx$ is nonzero and is
annihilated by both $x$ and $y$, as computed in \cref{cdp:st:B1}.
Thus this fibre restriction of the torsion-free total-space module
has the stated conductor-supported torsion.

For a finite surjective \'etale map $h:U'\to U$ over the base curve,
the map $h^*\Omega^1_{U/B}\to\Omega^1_{U'/B}$ induced by
$da\mapsto d(h^*a)$ is an isomorphism. Flat pullback preserves the
kernel of multiplication by a nonzerodivisor, and faithful
flatness reflects its nonvanishing; it therefore transports the
torsion calculation for the multiple-fibre module in
\eqref{wbII:eq:CN:6.4}. For a ramified quotient the transitivity
sequence has the additional term $\Omega^1_{U'/U}$, so the
quotient module is computed with that term and the group action.
\end{remark}

\section{Marked Picard restrictions and boundary cocycles}\hypertarget{tag-P310}{}\label{sec:II:analytic-picard-boundary}

\begin{definition}[Marked periods and cusp normalization]
\hypertarget{tag-P309}{}\label{def:II:analytic-picard-data}
\StatementRecord{def:II:analytic-picard-data}{proved}{P309}{\SeriesRef{def:II:marked-conventions}, \SeriesRef{fam:src:A.1.2}, \SeriesRef{fam:src:A.3.1}, \SeriesRef{push:src:E.3.1}}
Let $f:X\to B=\mathbb P^1$ be the standard family of Paper I,
$B^\circ=B\setminus\{p_0,p_1,p_2\}$, and $F$ is a smooth fibre.
The cusp is $W=f^{-1}(p_0)$ and $i:W\hookrightarrow X$ is its inclusion.
Its normalization is $\eta:\widetilde W\simeq\mathrm{dP}_6\to W$.
Use the ordered dual lattices
\[
 V=H^1(F;\mathbb Z)=\mathbb Z\langle\gamma,u,w,\delta\rangle,
 \qquad
 \Lambda=H_1(F;\mathbb Z)
 =\mathbb Z\langle\widehat\gamma,\widehat u,\widehat w,\widehat\delta\rangle.
\]
Vectors are columns. The cohomology monodromies and their dual homology
actions are
\[
T_1=\begin{pmatrix}1&0&-6&2\\0&-1&1&1\\0&-1&0&1\\0&0&0&1\end{pmatrix},
\quad
T_2=\begin{pmatrix}1&6&0&-3\\0&0&-1&1\\0&1&0&0\\0&0&0&1\end{pmatrix},
\quad A_j=(T_j^{-1})^{\mathsf t}.
\]
Put $T_0=(T_1T_2)^{-1}$ and let $G$ denote the generated monodromy group.
The finite-filling orders, translations and invariant forms are
\[
\begin{gathered}
(m_1,m_2)=(3,4),\qquad
v_1=(1,2,-4,0)^{\mathsf t},\quad v_2=(-1,-3,3,0)^{\mathsf t},\\
\eta_1=2u+w+3\delta,\quad \eta_2=u+w+2\delta,
\qquad q=u\wedge w+6\gamma\wedge\delta.
\end{gathered}
\]
The period functions of Paper I define
\[
\Pi(z)=\begin{pmatrix}6\mu(z)&\tau(z)&1&0\\
\beta(z)&\mu(z)&0&1\end{pmatrix},
\quad
\operatorname{Im}\beta-
\frac{6(\operatorname{Im}\mu)^2}{\operatorname{Im}\tau}<0,
\]
so $\Pi(z):\Lambda_{\mathbb R}\to\mathbb C^2$ is a real isomorphism.
The marking is $\Xi(z,x+\Lambda)=[z,\Pi(z)x]$.
The finite overlap on the $s_j$-cover is translation by
$(\log s_j)/(2\pi i)\,\Pi(s_j)v_j$.
At the cusp use the zero-section gluing, whose vertical translation is zero. These are the data constructed and proved in
\SeriesRef{fam:src:A.1.2}, \SeriesRef{fam:src:A.3.1} and
\SeriesRef{push:src:E.3.1}. In the notation of
\cref{def:II:marked-conventions}, $\beta=\beta^{\rm part}+c$,
$(v_1,v_2)=(\varepsilon,-\varepsilon')$, and the marking
$\Xi$ is induced by the same period matrix $\Pi_c$.
The normalized finite base change is
$y\mapsto(q_j(y),s_j(y))$ over $t_j=s_j^{m_j}$, with deck
$(s_j,\zeta)\mapsto(e^{-2\pi i/m_j}s_j,
R_j(s_j)\zeta+\Pi_c(g_js_j)v_j/m_j)$.
Here
\[
 R_1=\begin{pmatrix}-1/\tau&0\\(1-\mu)/\tau&1\end{pmatrix},
 \qquad R_2=\begin{pmatrix}1/\tau&0\\-\mu/\tau&1\end{pmatrix},
\]
and $R_j(z)\Pi_c(z)=\Pi_c(g_jz)A_j$.
The finite deck and overlap translations are
\SeriesRef{pub:I:eq:PD7} and \SeriesRef{pub:I:eq:PD8}.
At the cusp, use the coordinates $Z=sB_0+C(t)$ of
\SeriesRef{fam:src:A.3.1} and the period action $\Psi_\lambda$
in \SeriesRef{pub:I:eq:PD11}; $\eta$ is the restriction of that
quotient map to the divisor $E_0$.
\end{definition}

\begin{proposition}[Restriction to the cusp normalization]
\hypertarget{tag-03IO}{}\label{wbII:prop:pic0-pullback}
\StatementRecord{wbII:prop:pic0-pullback}{proved}{03IO}{}
\hypertarget{tag-03IG}{}\label{wbII:prop:all-pic-restriction}
\StatementRecord{wbII:prop:all-pic-restriction}{proved}{03IG}{}

Let $i:W\hookrightarrow X$ be the cusp-fibre inclusion and
$\eta:\widetilde W\simeq\mathrm{dP}_6\to W$ its normalization.
Assume that $\omega_X$ is trivial on a neighbourhood of $W$, and put
\[
 A_L=(L^\vee\otimes\omega_X)|_W,
 \qquad
 r=\eta^*i^*:\operatorname{Pic}(X)\longrightarrow
 \operatorname{Pic}(\mathrm{dP}_6).
\]
Then
\begin{equation}
 \eta^*A_L\simeq r(L)^\vee
 \qquad(L\in\operatorname{Pic}(X)).
\hypertarget{tag-03GL}{}\label{wbII:eq:W:5.17b}
\end{equation}
The homomorphism $r$ is zero precisely when every $\eta^*A_L$ is trivial.
For $L\in\operatorname{Pic}^0(X)$ one has
\begin{equation}
 \eta^*A_L\simeq\mathcal O_{\mathrm{dP}_6}.
\hypertarget{tag-03FV}{}\label{wbII:eq:W:5.17a}
\end{equation}
\end{proposition}

\begin{proof}
Pullback commutes with tensor products and duals, and the neighbourhood
trivialization pulls back to
$\eta^*i^*\omega_X\simeq\mathcal O_{\mathrm{dP}_6}$.
Thus $\eta^*A_L\simeq r(L)^\vee$; dualization sends the trivial bundle to
itself and is involutive, proving the asserted equivalence.
The degree-six del Pezzo surface is the blow-up of $\mathbb P^2$ at
three noncollinear points.  Their homogeneous representatives form a
basis of $\mathbb C^3$; the inverse of the matrix of that basis induces
a projective automorphism carrying the points to the three coordinate
points.  Pullback of their ideal sheaves induces the corresponding
isomorphism of blow-ups.  Thus this is the smooth complete toric surface
obtained by subdividing the three two-dimensional cones of the fan of
$\mathbb P^2$.
Demazure vanishing for the zero nef divisor gives
$H^1(\mathrm{dP}_6,\mathcal O)=0$ in algebraic cohomology
\cite[Theorem~9.2.3, 2010 text, p.~410]{CoxLittleSchenck2011};
Serre's comparison theorem gives the same analytic vanishing
\cite[\S12, Theorem~1, p.~19]{Serre1956GAGA}.
Consequently the exponential sequence makes
\[
 c_1:\operatorname{Pic}(\mathrm{dP}_6)\longrightarrow
 H^2(\mathrm{dP}_6;\mathbb Z)
\]
injective.  For $L\in\operatorname{Pic}^0(X)$, naturality gives
$c_1(r(L))=\eta^*i^*c_1(L)=0$.  Hence $r(L)$ and its dual are trivial,
which proves \eqref{wbII:eq:W:5.17a}.
\end{proof}

\begin{proposition}[Boundary cocycles and Chern-class specialization]
\hypertarget{tag-03IH}{}\label{wbII:prop:analytic-ah-boundary}
\StatementRecord{wbII:prop:analytic-ah-boundary}{proved}{03IH}{}
\hypertarget{tag-03IC}{}\label{wbII:lem:picard-crossed-defect}
\StatementRecord{wbII:lem:picard-crossed-defect}{proved}{03IC}{}
\hypertarget{tag-03IE}{}\label{wbII:lem:picard-finite-norm}
\StatementRecord{wbII:lem:picard-finite-norm}{proved}{03IE}{}
\hypertarget{tag-03ID}{}\label{wbII:lem:picard-cusp-defect}
\StatementRecord{wbII:lem:picard-cusp-defect}{proved}{03ID}{}

Let \(P\subset B^\circ\) be the compact pair-of-pants core,
orient its boundary meridians clockwise, and use the zero-section lifts
of \SeriesRef{inv:B3.1}.  The standard family uses the
tautological cusp gluing, whose affine regluing class is $\mu=0$;
thus $\ell_0=\gamma(\mu)=0$.  The based loops satisfy
\begin{equation}
 \rho_1\rho_2\rho_0=1,
 \qquad
 \widehat\rho_1\widehat\rho_2\widetilde a_0=1,
\hypertarget{tag-03EF}{}\label{wbII:eq:W:5.17-A1}
\hypertarget{tag-03GQ}{}\SsixEquationAliases{\label{wbII:eq:W:5.17g}}
\end{equation}
and the second relation has no terminal vertical translation.  For
\(L\in\operatorname{Pic}(X)\), let
\[
 E=c_1(L|_F)\in\operatorname{Alt}^2(\Lambda,\mathbb Z)
   =\bigwedge^2V,
 \qquad E\in(\bigwedge^2V)^G.
\]
be its fibrewise Appell--Humbert form.  Put
\(V_{\mathbb R}=\operatorname{Hom}(\Lambda,\mathbb R)\), and use
\begin{equation}
 A_g=(T_g^{-1})^{\mathsf t},
 \qquad
 \langle T_g\alpha,\lambda\rangle
 =\langle\alpha,A_g^{-1}\lambda\rangle.
\hypertarget{tag-03EY}{}\label{wbII:eq:W:5.17-A2}
\end{equation}
Here $A_g$ is the linear part of the left deck action and $T_g=A_g^{-\mathsf t}$
is its coefficient action.  A clockwise loop ends at the deck translate $g$;
its endpoint return $g^{-1}$ acts on homology by $A_g^{-1}$ and on
cohomology by $A_g^{\mathsf t}=T_g^{-1}$.  The cocycle below uses $T_g$.
There is a crossed homomorphism, unique up to a principal crossed
homomorphism,
\begin{equation}
 d:\pi_1(P)\longrightarrow V_{\mathbb R},
 \qquad
 d(gh)=d(g)+T_gd(h),
 \qquad
 d(g)\sim d(g)+(T_g-I)a.
\hypertarget{tag-03F9}{}\label{wbII:eq:W:5.17-A3}
\hypertarget{tag-03GR}{}\SsixEquationAliases{\label{wbII:eq:W:5.17h}}
\end{equation}
Write $J_P=f^{-1}(P)$, and let $\Omega_E$ be the constant vertical
two-form in the flat marking.  The class of $d$ is the unique preimage of
$c_1(L|_{J_P})-[\Omega_E]$ under the edge injection
\eqref{wbII:eq:W:5.17j}.  We write $z=d$ in that cohomological description.
For an oriented loop $\rho\subset P$ and an $A_\rho$-invariant
$\nu\in\Lambda$, let $a$ be fibrewise translation and
$\widetilde\rho$ the flat-marked lift.  With the base-first convention
for the edge injection, the fibre-first sweep satisfies
\begin{equation}
 \bigl\langle z(\rho),\nu\bigr\rangle
 =-\bigl\langle c_1(L|_{J_P})-[\Omega_E],
 a_*([S^1_\nu]\times[\widetilde\rho])\bigr\rangle.
\hypertarget{tag-03GU}{}\label{wbII:eq:W:5.17j'}
\end{equation}
If \(d_i=d(\rho_i)\), then
\begin{equation}
 d_1+T_1d_2+T_1T_2d_0=0,
 \qquad
 \bar d_1+\bar d_2+\bar d_0=0
 \quad\text{in }V_G\otimes\mathbb R.
\hypertarget{tag-03FM}{}\label{wbII:eq:W:5.17-A4}
\hypertarget{tag-03GS}{}\SsixEquationAliases{\label{wbII:eq:W:5.17i}}
\end{equation}
For the finite filling of order \(m_j\), \(j=1,2\),
\begin{equation}
 N_{T_j}d_j=\iota_{v_j}E,
 \qquad
 m_j\bar d_j=[\iota_{v_j}E],
 \qquad
 N_{T_j}=I+T_j+\cdots+T_j^{m_j-1},
\hypertarget{tag-03FN}{}\label{wbII:eq:W:5.17-A5}
\hypertarget{tag-03GV}{}\SsixEquationAliases{\label{wbII:eq:W:5.17k}}
\end{equation}
where \((\iota_vE)(\lambda)=E(v,\lambda)\).  At the cusp,
\begin{equation}
 \langle d_0,\widehat\delta\rangle=0,
 \qquad
 \bar d_0=0\quad\text{in }V_G\otimes\mathbb R.
\hypertarget{tag-03FO}{}\label{wbII:eq:W:5.17-A6}
\hypertarget{tag-03H3}{}\SsixEquationAliases{\label{wbII:eq:W:5.17n}}
\end{equation}

For every line bundle $L$ on the cusp neighbourhood $N_0$, the nearby-cycle
specialization satisfies
\begin{equation}
 \operatorname{sp}^2\bigl(c_1(L|_W)\bigr)=c_1(L|_F).
\hypertarget{tag-03H4}{}\label{wbII:eq:W:5.17o}
\end{equation}
\end{proposition}



\begin{proof}

\begin{enumerate}[label=\arabic*.,leftmargin=*]
\item Choose an open Riemann-surface neighbourhood \(P^{\rm an}\Subset B^\circ\)
which deformation retracts onto \(P\), and pull the period family to its
universal cover.  Before quotienting by the fibre lattice, the resulting
space is
\[
 \widetilde{\mathcal J}
 =\widetilde P^{\rm an}\times\mathbb C^2.
\]
It is Stein and contractible, so the pullback of \(L\) is holomorphically
trivial.  In a nowhere-zero frame, write the factor belonging to a deck
transformation \(\alpha\) as
\begin{equation}
 a_\alpha\in\mathcal O^*(\widetilde{\mathcal J}),
 \qquad
 a_{\alpha\beta}(x)=a_\alpha(\beta x)a_\beta(x).
\hypertarget{tag-03FP}{}\label{wbII:eq:W:5.17-A7}
\end{equation}
The space is simply connected, hence \(a_\alpha=e^{2\pi i h_\alpha}\) for
a holomorphic logarithm.  Therefore
\begin{equation}
 c(\alpha,\beta)
 =h_\alpha(\beta x)+h_\beta(x)-h_{\alpha\beta}(x)\in\mathbb Z
\hypertarget{tag-03FQ}{}\label{wbII:eq:W:5.17-A8}
\end{equation}
is independent of \(x\), and associativity makes \(c\) an integral group
two-cocycle.

\item Let \(\tau_\lambda\) denote fibre translation by
\(\Pi(z)\lambda\), \(\lambda\in\Lambda\), and let \(\widehat g\) be the
zero-section lift of a base deck transformation.  Period equivariance gives
\begin{equation}
 \widehat g\,\tau_\lambda\,\widehat g^{-1}
 =\tau_{A_g\lambda}.
\hypertarget{tag-03FR}{}\label{wbII:eq:W:5.17-A9}
\end{equation}
On fibre translations, the commutator of \eqref{wbII:eq:W:5.17-A8} is
\begin{equation}
 E(\lambda,\mu)
 =c(\tau_\lambda,\tau_\mu)-c(\tau_\mu,\tau_\lambda).
\hypertarget{tag-03EG}{}\label{wbII:eq:W:5.17-A10}
\end{equation}
The four oriented sides of the \((\lambda,\mu)\)-parallelogram identify
this integer with the exponential connecting class of \(L|_F\),
with the factor-of-automorphy convention of
\cite[equations~(3.3)--(3.10), pp.~3--7]{SchnellLecture3}.
Conjugating the commutator by \(\widehat g\) yields
\begin{equation}
 E(A_g\lambda,A_g\mu)=E(\lambda,\mu).
\hypertarget{tag-03EH}{}\label{wbII:eq:W:5.17-A11}
\end{equation}

\item Push the logarithmic central extension defined by \(c\) out along
\(\mathbb Z\hookrightarrow\mathbb R\), and denote its central element of
weight \(r\) by \(\mathbf c(r)\).  Choose lifts \(U_\lambda\) of the fibre
translations satisfying the real Appell--Humbert normalization
\begin{equation}
 U_\lambda U_\mu
 =\mathbf c\!\left(\frac12E(\lambda,\mu)\right)U_{\lambda+\mu}.
\hypertarget{tag-03EI}{}\label{wbII:eq:W:5.17-A12}
\end{equation}
To see that this normalization exists, start with arbitrary lifts.  The
difference between their real cocycle and \(E/2\) has zero commutator, so
it defines an abelian extension of the free abelian group \(\Lambda\) by
\(\mathbb R\).  Choose lifts of an integral basis and extend
multiplicatively; this splits that abelian extension and gives \eqref{wbII:eq:W:5.17-A12}.
If \(U'_\lambda\) is a second normalized system, then
\(U'_\lambda U_\lambda^{-1}\) is central, say
\(\mathbf c(\langle a,\lambda\rangle)\).  Comparing the two copies of
\eqref{wbII:eq:W:5.17-A12} proves additivity in \(\lambda\), so
\(a\in V_{\mathbb R}\).  Thus every two normalized systems differ by exactly
the character change used below.

\item For \(g\in\pi_1(P)\), choose a lift \(G_g\) of \(\widehat g\) and define
\(d(g)\in V_{\mathbb R}\) by
\begin{equation}
 G_gU_{A_g^{-1}\lambda}G_g^{-1}
 =\mathbf c\bigl(\langle d(g),\lambda\rangle\bigr)U_\lambda.
\hypertarget{tag-03EJ}{}\label{wbII:eq:W:5.17-A13}
\end{equation}
The exponent is additive in \(\lambda\): conjugate \eqref{wbII:eq:W:5.17-A12} and use
\eqref{wbII:eq:W:5.17-A11}.  The lifts of \(gh\) and \(G_gG_h\) differ centrally, so
\[
\begin{aligned}
 \langle d(gh),\lambda\rangle
 &=\langle d(g),\lambda\rangle
   +\langle d(h),A_g^{-1}\lambda\rangle\\
 &=\langle d(g)+T_gd(h),\lambda\rangle.
\end{aligned}
\]
This proves the crossed-homomorphism law.  Replacing
\(U_\lambda\) by
\(\mathbf c(\langle a,\lambda\rangle)U_\lambda\) changes \eqref{wbII:eq:W:5.17-A13} by
\((T_g-I)a\).  Replacing \(G_g\) by a central multiple does not change
\eqref{wbII:eq:W:5.17-A13}, so this proves the stated independence and uniqueness.  Applying
the crossed law to
\eqref{wbII:eq:W:5.17-A1} proves the first identity in \eqref{wbII:eq:W:5.17-A4}; passage to coinvariants
proves the second.

\item The flat marking identifies $J_P$ with the quotient of
$\widetilde P\times\Lambda_{\mathbb R}/\Lambda$ by the left deck action.
The invariance in \eqref{wbII:eq:W:5.17-A11} makes the constant vertical
form $\Omega_E$ descend to $J_P$.
Since $P$ retracts onto a graph, the Serre filtration gives
\begin{equation}
0\longrightarrow H^1(P;V_{\mathbb R})
 \longrightarrow H^2(J_P;\mathbb R)
 \longrightarrow H^0(P;\textstyle\bigwedge^2V_{\mathbb R})
 \longrightarrow0.
\hypertarget{tag-03GT}{}\label{wbII:eq:W:5.17j}
\end{equation}
We identify the edge preimage explicitly.  The universal cover of $J_P$
is contractible and its deck group is
$\Gamma=\Lambda\rtimes\pi_1(P)$, with multiplication
$(\lambda,g)(\mu,h)=(\lambda+A_g\mu,gh)$.
The exponential connecting map sends the logarithmic cocycle
\eqref{wbII:eq:W:5.17-A8} to $c_1(L|_{J_P})$: on an evenly covered
open cover, its integer on a triple overlap is exactly the sum of the
three chosen transition logarithms.  Pushing these integers to
$\mathbb R$ gives the central extension used above.
Choose the lifts $G_g$ multiplicatively; this is possible by choosing
lifts of the two free generators of $\pi_1(P)$ and their inverses.
With section $(\lambda,g)\mapsto U_\lambda G_g$, multiplication in that
extension has real cocycle
\[
 C\bigl((\lambda,g),(\mu,h)\bigr)
 =\frac12E(\lambda,A_g\mu)
   +\langle d(g),A_g\mu\rangle.
\]
Indeed, first move $U_\mu$ across $G_g$ using
\eqref{wbII:eq:W:5.17-A13}, then multiply the two lattice lifts using
\eqref{wbII:eq:W:5.17-A12}.  The first term is the cocycle of
$[\Omega_E]$: integration on the marked lattice triangle with sides
$\lambda,A_g\mu$ gives $E(\lambda,A_g\mu)/2$.
The second term has filtration degree $(1,1)$ and assigns to the
oriented base edge $g$ the covector $d(g)$, with the fibre argument
transported to its endpoint by $A_g$.  Thus the edge class is exactly
$[d]\in H^1(P;V_{\mathbb R})$, and the crossed-homomorphism model
identifies its representatives with $d$ modulo $(T_g-I)a$
\cite[Example~3.5, p.~4]{SchnellLecture3}.  This proves $z=d$ as a choice
of the cohomological representative, with the same full ambiguity.
Explicitly, for $b(\lambda,g)=\langle a,\lambda\rangle$, the group
coboundary is
\[
 (\delta b)\bigl((\lambda,g),(\mu,h)\bigr)
 =\langle a,\mu-A_g\mu\rangle
 =\langle(T_g-I)a,A_g\mu\rangle,
\]
which is exactly the change of the mixed summand under the principal
change of $d$.
For $A_\rho\nu=\nu$, the base-first bar cycle
$[(0,\rho)|(\nu,1)]-[(\nu,1)|(0,\rho)]$ evaluates the mixed term to
$\langle d(\rho),\nu\rangle$.  Reversing the order of the two circles
introduces $(-1)^{1\cdot1}$, proving \eqref{wbII:eq:W:5.17j'}.
The principal term evaluates to zero because
$\langle(T_\rho-I)a,\nu\rangle
=\langle a,(A_\rho^{-1}-I)\nu\rangle=0$.
Restriction of this same cochain to the $m_j$-fold boundary cover gives
\begin{equation}
 z(\rho_j^{m_j})
 =\sum_{r=0}^{m_j-1}T_j^rz(\rho_j)=N_{T_j}d_j.
\hypertarget{tag-03GW}{}\label{wbII:eq:W:5.17l}
\end{equation}
Its value is independent of the chosen representative, since
$N_{T_j}(T_j-I)=T_j^{m_j}-I=0$.

\item For the order-\(m_j\) filling, put
\[
 \widetilde J_j=\Delta_j\times\mathbb C^2,\qquad
 r_j:\widetilde J_j\longrightarrow J'_j
\]
for the lattice quotient, whose deck translations are \(\tau_\lambda\).
On the vector cover define
\[
 Q_{j,\log}:=\widetilde g_j^{\,v_j/m_j},\qquad
 Q_{j,\mathrm{lin}}:=\widetilde g_j^{\,0}.
\]
They satisfy
\begin{equation}
\begin{aligned}
 Q_{j,\log}\tau_\lambda Q_{j,\log}^{-1}
  &=\tau_{A_j\lambda},&
 (Q_{j,\log})^{m_j}&=\tau_{v_j},\\
 Q_{j,\mathrm{lin}}\tau_\lambda Q_{j,\mathrm{lin}}^{-1}
  &=\tau_{A_j\lambda},&
 (Q_{j,\mathrm{lin}})^{m_j}&=1.
\end{aligned}
\hypertarget{tag-03EK}{}\label{wbII:eq:W:5.17-A14}
\end{equation}
The positive sign is fixed by
\(s_j\circ g_j=e^{-2\pi i/m_j}s_j\) and \SeriesRef{pub:I:eq:PD7}: after exactly
\(m_j\) clockwise iterations, the logarithmic vector-cover lift translates
by \(+v_j\).

\item Let \(q_{j,\log},q_{j,\mathrm{lin}}\) be the descended automorphisms of \(J'_j\),
and put \(J_j^{\prime *}=J'_j|_{\Delta_j^*}\).  Both descended automorphisms have
order dividing \(m_j\), because \(\tau_{v_j}\) is trivial after the lattice
quotient; each covers \(g_j\), whose order is exactly \(m_j\), so each
descended order is exactly \(m_j\).
The branch-independent torus translation
\[
 h_j:J_j^{\prime *}\longrightarrow J_j^{\prime *}
\]
constructed in \SeriesRef{pub:I:eq:PD8} satisfies
\[
 h_jq_{j,\log}h_j^{-1}=q_{j,\mathrm{lin}}.
\]
Define
\[
\begin{aligned}
 p_{j,\log}:J_j^{\prime *}&\longrightarrow
  N_j^*:=J_j^{\prime *}/\langle q_{j,\log}\rangle,\\
 p_{j,\mathrm{lin}}:J_j^{\prime *}&\longrightarrow
 J_j^*:=J_j^{\prime *}/\langle q_{j,\mathrm{lin}}\rangle .
\end{aligned}
\]
The conjugacy induces a biholomorphism
\[
 \bar h_j:N_j^*\longrightarrow J_j^*,\qquad
 p_{j,\mathrm{lin}}\circ h_j=\bar h_j\circ p_{j,\log}.
\]
The gluing of \(X\) identifies
\(L|_{N_j^*}\simeq\bar h_j^*(L|_{J_j^*})\); hence the
line-bundle square gives
\begin{equation}
 p_{j,\log}^*(L|_{N_j^*})
 \simeq h_j^*p_{j,\mathrm{lin}}^*(L|_{J_j^*}).
\hypertarget{tag-03EL}{}\label{wbII:eq:W:5.17-A14a}
\hypertarget{tag-03GZ}{}\SsixEquationAliases{\label{wbII:eq:W:5.17l'''}}
\end{equation}

\item On the filled vector cover, let
\[
 P_{j,\log}:J'_j\longrightarrow N_j,\qquad
 \pi_j^{\mathrm{vec}}:=P_{j,\log}\circ r_j:
 \widetilde J_j\longrightarrow N_j .
\]
Here \(P_{j,\log}\) is the full quotient and
\(P_{j,\log}|_{J_j^{\prime *}}=p_{j,\log}\).
The deck group is
\[
 \Gamma_j^{\log}
 =\langle\tau_\lambda,Q_{j,\log}:\lambda\in\Lambda\rangle.
\]
The pullback \((\pi_j^{\mathrm{vec}})^*(L|_{N_j})\) is trivial on the
contractible Stein manifold \(\widetilde J_j\); its factors define an
integral central extension, which we push out along
\(\mathbb Z\hookrightarrow\mathbb R\).  Over \(\Delta_j^*\), the
line-bundle square \eqref{wbII:eq:W:5.17-A14a} identifies each fibre with the global factor
extension by a translation, so its commutator form is \(E\).  That form is
integer-valued and locally constant on the connected disc \(\Delta_j\);
therefore it is \(E\) on the filled cover as well.  Choose vertical lifts
\(U_\lambda^{(j)}\) normalized by
\[
 U_\lambda^{(j)}U_\mu^{(j)}
 =\mathbf c\!\left(\frac12E(\lambda,\mu)\right)
  U_{\lambda+\mu}^{(j)}.
\]
For \(\alpha\in\Gamma_j^{\log}\), define \(A_\alpha,T_\alpha\) by
\[
 \alpha\tau_\lambda\alpha^{-1}=\tau_{A_\alpha\lambda},\qquad
 T_\alpha=(A_\alpha^{-1})^{\mathsf t}.
\]
For a lift \(\widetilde G_\alpha\), define
\(\widetilde d_j(\alpha)\in V_{\mathbb R}\) through
\begin{equation}
 \widetilde G_\alpha U_{A_\alpha^{-1}\lambda}^{(j)}
 \widetilde G_\alpha^{-1}
 =\mathbf c\!\left(\langle\widetilde d_j(\alpha),\lambda\rangle\right)
  U_\lambda^{(j)}.
\hypertarget{tag-03EM}{}\label{wbII:eq:W:5.17-A14c}
\end{equation}
The two lifts in \eqref{wbII:eq:W:5.17-A14c} cover the same translation,
so their quotient is central.
Conjugating the normalized product and using
\[
 E(A_\alpha\lambda,A_\alpha\mu)=E(\lambda,\mu)
\]
shows that its exponent \(r_\alpha(\lambda)\) satisfies
\[
 r_\alpha(\lambda+\mu)=r_\alpha(\lambda)+r_\alpha(\mu).
\]
Thus it is the pairing with a unique element of \(V_{\mathbb R}\), as
asserted.  Since \(A_{\alpha\beta}=A_\alpha A_\beta\), choose the lift of
\(\alpha\beta\) to be \(\widetilde G_\alpha\widetilde G_\beta\); any other
choice differs centrally.  Then
\[
\begin{aligned}
 &\widetilde G_\alpha\widetilde G_\beta
 U_{A_\beta^{-1}A_\alpha^{-1}\lambda}^{(j)}
 \widetilde G_\beta^{-1}\widetilde G_\alpha^{-1}\\
 &\quad=
 \mathbf c\!\left(
  \langle\widetilde d_j(\beta),A_\alpha^{-1}\lambda\rangle
  +\langle\widetilde d_j(\alpha),\lambda\rangle\right)
 U_\lambda^{(j)}.
\end{aligned}
\]
The defining dual action
\(\langle T_\alpha v,\lambda\rangle
=\langle v,A_\alpha^{-1}\lambda\rangle\) therefore proves the crossed law
\begin{equation}
 \widetilde d_j(\alpha\beta)
 =\widetilde d_j(\alpha)+T_\alpha\widetilde d_j(\beta).
\hypertarget{tag-03EN}{}\label{wbII:eq:W:5.17-A14d}
\end{equation}

\item Set $Q_j=P_{j,\log}:J'_j\to N_j$.  The global flat marking
$\Xi$ over the full disc $\Delta_j$ makes the constant vertical form
$\Omega_E$ well defined on $J'_j$.  Radial retraction onto the central
torus makes
\[
 H^2(J'_j;\mathbb R)\longrightarrow H^2(A_0^{(j)};\mathbb R)
\]
an isomorphism.  On the central fibre the quotient is exactly $\pi_j$, and
naturality gives
\begin{equation}
 (Q_j|_{A_0^{(j)}})^*c_1(L|_{S_j})
 =\pi_j^*c_1(L|_{S_j})=E.
\hypertarget{tag-03GX}{}\label{wbII:eq:W:5.17l'}
\end{equation}
The second equality follows from the normalized root cover in
\SeriesRef{fam:src:A.1.2} and the radial restriction above; on a nearby fibre the gluing map $h_j$
is a translation and therefore acts trivially on $H^2$ of that fibre.
Consequently the two classes
$c_1(Q_j^*(L|_{N_j}))$ and $[\Omega_E]$ have the same central restriction,
and the retraction isomorphism proves
\begin{equation}
 c_1(Q_j^*(L|_{N_j}))-[\Omega_E]=0
 \quad\text{in }H^2(J'_j;\mathbb R).
\hypertarget{tag-03GY}{}\label{wbII:eq:W:5.17l''}
\end{equation}


\item The connector has holomorphic translation
$\alpha(s)\Pi(s)v_j$, where $\alpha(s)=\log s_j/(2\pi i)$.
In the full-disc real marking
\[
 \Xi(s,x+\Lambda)=[s,\Pi(s)x],\qquad
 J_s=\Pi(s)^{-1}\circ i\circ\Pi(s),
\]
its translation vector is
\begin{equation}
 a_j(s)=\operatorname{Re}\alpha(s)\,v_j+
        \operatorname{Im}\alpha(s)\,J_s v_j.
 \hypertarget{tag-03H0}{}\label{wbII:eq:W:5.17m}
\end{equation}
The period inequality makes $\Pi(s)$ a real isomorphism, and applying
$\Pi(s)$ to this vector gives exactly $\alpha(s)\Pi(s)v_j$.
A logarithm-branch change adds an integral multiple of $v_j$, which
induces the identity translation on the lattice quotient.
Write $\Omega_E=\frac12\sum_{a,b}E_{ab}\,dx^a\wedge dx^b$ in these
marked coordinates.  Substitution of $x\mapsto x+a_j(s)$ gives
\begin{equation}
 \begin{split}
 h_j^*\Omega_E-\Omega_E
  ={}&\sum_{a,b}E_{ab}\,d a_j^a\wedge dx^b\\
    &+\frac12\sum_{a,b}E_{ab}\,d a_j^a\wedge d a_j^b.
 \end{split}
 \hypertarget{tag-03H1}{}\label{wbII:eq:W:5.17m'}
\end{equation}
Along $s_j(\theta)=r e^{-2\pi i\theta}$, $0\le\theta\le1$, continuous
choice of the logarithm gives $\alpha(1)=\alpha(0)-1$.
The full-disc marking has $J_{s_j(1)}=J_{s_j(0)}$, so
\begin{equation}
 a_j(1)-a_j(0)=-v_j.
 \hypertarget{tag-03H2}{}\label{wbII:eq:W:5.17m''}
\end{equation}

\item Put
\[
 \widetilde J_j^*:=\Delta_j^*\times\mathbb C^2,
 \qquad
 r_j^*:=r_j|_{\widetilde J_j^*}:
 \widetilde J_j^*\longrightarrow J_j^{\prime *}
\]
and define
\[
\begin{aligned}
 \Phi_{j,\log}
  &:=p_{j,\log}\circ r_j^*,\\
 \Phi_{j,\mathrm{glob}}
  &:=p_{j,\mathrm{lin}}\circ h_j\circ r_j^*
   =\bar h_j\circ\Phi_{j,\log}.
\end{aligned}
\]
Both maps are invariant under \(\Gamma_j^{\log}\).  This is immediate for
every \(\tau_\lambda\), while
\[
\begin{aligned}
 \Phi_{j,\log}Q_{j,\log}
 &=p_{j,\log}q_{j,\log}r_j^*
  =\Phi_{j,\log},\\
 \Phi_{j,\mathrm{glob}}Q_{j,\log}
 &=p_{j,\mathrm{lin}}h_jq_{j,\log}r_j^*
  =p_{j,\mathrm{lin}}q_{j,\mathrm{lin}}h_jr_j^*
  =\Phi_{j,\mathrm{glob}}.
\end{aligned}
\]
Thus \(\widetilde J_j^*\) is the same \(\Gamma_j^{\log}\)-cover of the two
biholomorphic quotients.  Pulling the gluing isomorphism to this cover gives
\[
 \Theta_j:
 \Phi_{j,\log}^*(L|_{N_j^*})
 \xrightarrow{\ \sim\ }
 \Phi_{j,\mathrm{glob}}^*(L|_{J_j^*}).
\]
For the comparison of logarithms, pass one further time to the universal
cover
\[
 \widehat J_j^*=\widetilde{\Delta_j^*}\times\mathbb C^2
 \longrightarrow \widetilde J_j^*.
\]
The deck group over the punctured quotient is
\[
 \widehat\Gamma_j=
 \langle\tau_\lambda,q_j\mid
 q_j\tau_\lambda q_j^{-1}=\tau_{A_j\lambda}\rangle.
\]
The map from this universal cover to the filled vector cover induces the
explicit homomorphism
\[
 \sigma_j:\widehat\Gamma_j\longrightarrow\Gamma_j^{\log},
 \qquad \tau_\lambda\longmapsto\tau_\lambda,
 \quad q_j\longmapsto Q_{j,\log}.
\]
Its kernel is generated by
\(q_j^{m_j}\tau_{-v_j}\): this element is one deck turn of
\(\widetilde{\Delta_j^*}\to\Delta_j^*\), while its image on the
filled vector cover is the identity by \eqref{wbII:eq:W:5.17-A14}.
Indeed every word is \(\tau_\lambda q_j^r\); its filled image can be
identity only when \(r=m_jk\), and then only when
\(\lambda=-kv_j\).  Thus the kernel assertion follows directly.

\item Let \(\mathscr E_j\) be the real central extension of
\(\Gamma_j^{\log}\) defined by the filled factors.  Pull it back along
\(\sigma_j\), explicitly as the fibre-product group
\[
 \sigma_j^*\mathscr E_j=
 \{(e,g):e\in\mathscr E_j,\ g\in\widehat\Gamma_j,
           \ p(e)=\sigma_j(g)\}.
\]
The global factor extension restricted to the based boundary group is
an extension of the same \(\widehat\Gamma_j\), with the generator
\(q_j\) identified with \(\widehat\rho_j\) through the fixed whisker.
The endpoint calculation below proves this precise identification.
Pull \(\Theta_j\) to \(\widehat J_j^*\).  This space is simply
connected, so the holomorphic unit comparing the two frames has a
holomorphic logarithm.  Substitution into \eqref{wbII:eq:W:5.17-A8}
changes the two integral cocycles by a coboundary.  Consequently it gives
an isomorphism from \(\sigma_j^*\mathscr E_j\) to the restricted global
extension, sending every \(\mathbf c(r)\) to the same
\(\mathbf c(r)\).  In the comparisons below,
\(U_\lambda^{(j)}\) means its paired lift
\((U_\lambda^{(j)},\tau_\lambda)\), and
\(\widetilde G_{Q_{j,\log}}\) means
\((\widetilde G_{Q_{j,\log}},q_j)\), when the isomorphism
\(\Theta_{j,*}\) is applied.  Their conjugation exponents remain the
ones in the filled extension.  
\item Let \(\gamma_j^J:[0,1]\to J_j^*\) be the zero-section image of the
unwhiskered clockwise boundary meridian, and let \(\omega_j\) be the fixed
zero-section whisker from the global basepoint to \(\gamma_j^J(0)\).  Thus
the based total-space loop in \eqref{wbII:eq:W:5.17-A1} is
\[
 \widehat\rho_j=\omega_j*\gamma_j^J*\omega_j^{-1}.
\]
Let $s_j$ be the root point selected by the lift of the fixed zero-section
whisker, and choose a logarithm branch $\alpha=\log s_j/(2\pi i)$.  For the vector-cover connector
$H_j(s,\zeta)=(s,\zeta+\alpha(s)\Pi(s)v_j)$, set
\[
 \widetilde x_j=(s_j,-\alpha(s_j)\Pi(s_j)v_j),\qquad
 x_j=r_j^*(\widetilde x_j),\qquad x_{j,\mathrm{lin}}=h_j(x_j).
\]
Thus $H_j(\widetilde x_j)=(s_j,0)$ and
$p_{j,\mathrm{lin}}(x_{j,\mathrm{lin}})=\gamma_j^J(0)$.  Let
\(\widetilde\gamma_{j,\mathrm{lin}}\) be the
\(p_{j,\mathrm{lin}}\)-lift of \(\gamma_j^J\) beginning at
\(x_{j,\mathrm{lin}}\).  For the chosen clockwise generator,
\[
 \widetilde\gamma_{j,\mathrm{lin}}(1)
 =q_{j,\mathrm{lin}}(x_{j,\mathrm{lin}}).
\]
Set
\[
 \gamma_j^N:=\bar h_j^{-1}\circ\gamma_j^J,
 \qquad
 \widetilde\gamma_{j,\log}
 :=h_j^{-1}\circ\widetilde\gamma_{j,\mathrm{lin}}.
\]
Then the quotient square gives the point-valued identities
\[
\begin{aligned}
 p_{j,\log}\circ\widetilde\gamma_{j,\log}
  &=\gamma_j^N,\\
 \widetilde\gamma_{j,\log}(0)&=x_j,\\
 \widetilde\gamma_{j,\log}(1)&=q_{j,\log}(x_j).
\end{aligned}
\]
Let $\widetilde\gamma_{j,\mathrm{vec}}$ be the $r_j^*$-lift of
$\widetilde\gamma_{j,\log}$ beginning at the selected $\widetilde x_j$.
For the clockwise lift, $\alpha_1=\alpha_0-1/m_j$.  Writing
$\zeta_0=-\alpha_0\Pi(s_j)v_j$, period equivariance and $A_jv_j=v_j$ give
\[
 \zeta_1=-(\alpha_0-1/m_j)\Pi(g_js_j)v_j
 =R_{g_j}(s_j)\zeta_0+\Pi(g_js_j)v_j/m_j.
\]
By \SeriesRef{pub:I:eq:PD7}, this gives
\[
 \widetilde\gamma_{j,\mathrm{vec}}(1)
 =Q_{j,\log}(\widetilde x_j),
\]
for this selected vector lift.
Fibre loops remain \(\tau_\lambda\), since \(h_j\) is a fibre translation.

\item Lift $\omega_j$ along the zero section on the linear cover and use
$H_j^{-1}$ at its local endpoint; this gives the selected
$\widetilde x_j$ above.  The factor of $L$ along this fixed lift identifies the
target factor over the local boundary segment with the globally based
factor.  Let \(G_{\rho_j}^{(j)}\) be the target lift of the local boundary
deck element whose whiskered image is the chosen global lift
\(G_{\rho_j}\).  Thus it covers \(\widehat\rho_j\) after whiskering and,
by \eqref{wbII:eq:W:5.17-A13}, satisfies
\[
 G_{\rho_j}^{(j)}U_{A_j^{-1}\lambda}(G_{\rho_j}^{(j)})^{-1}
 =\mathbf c(\langle d_j,\lambda\rangle)U_\lambda,
 \qquad d_j=d(\rho_j).
\]
Consequently, for some
\(a_j\in V_{\mathbb R}\) and \(b_j\in\mathbb R\),
\[
 \Theta_{j,*}(U_\lambda^{(j)})
 =\mathbf c(\langle a_j,\lambda\rangle)U_\lambda,
 \qquad
 \Theta_{j,*}(\widetilde G_{Q_{j,\log}})
 =\mathbf c(b_j)G_{\rho_j}^{(j)}.
\]
Here \(U_\lambda\) is the normalized global system in \eqref{wbII:eq:W:5.17-A12}, and
\(\mathbf c(b_j)\) disappears under conjugation.
Therefore
\[
\begin{aligned}
 \langle\widetilde d_j(Q_{j,\log}),\lambda\rangle
 +\langle a_j,\lambda\rangle
 &=\langle d_j,\lambda\rangle
   +\langle a_j,A_j^{-1}\lambda\rangle\\
 &=\langle d_j+T_ja_j,\lambda\rangle,
\end{aligned}
\]
so
\begin{equation}
 \widetilde d_j(Q_{j,\log})=d_j+(T_j-I)a_j.
\hypertarget{tag-03EO}{}\label{wbII:eq:W:5.17-A14e}
\end{equation}

\item For a translation choose
\(\widetilde G_{\tau_v}=U_v^{(j)}\).  The filled normalization above gives
\[
 U_v^{(j)}U_\lambda^{(j)}(U_v^{(j)})^{-1}
 =\mathbf c(E(v,\lambda))U_\lambda^{(j)},
\]
and hence
\begin{equation}
 \widetilde d_j(\tau_v)=\iota_vE,\qquad
 (\iota_vE)(\lambda)=E(v,\lambda).
\hypertarget{tag-03EP}{}\label{wbII:eq:W:5.17-A14f}
\end{equation}
Here
\[
 A_{Q_{j,\log}}=A_j,\qquad T_{Q_{j,\log}}=T_j.
\]
Since \(N_{T_j}(T_j-I)=T_j^{m_j}-I=0\), equations
\eqref{wbII:eq:W:5.17-A14}, \eqref{wbII:eq:W:5.17-A14d}, and \eqref{wbII:eq:W:5.17-A14e} give
\begin{equation}
\begin{aligned}
 N_{T_j}d_j
 &=N_{T_j}\widetilde d_j(Q_{j,\log})\\
 &=\widetilde d_j((Q_{j,\log})^{m_j})\\
 &=\widetilde d_j(\tau_{v_j})\\
 &=\iota_{v_j}E.
\end{aligned}
\hypertarget{tag-03EQ}{}\label{wbII:eq:W:5.17-A15}
\end{equation}
Changing a line-bundle frame gives an isomorphic factor extension; changing
\(\widetilde G_\alpha\) by a central factor does not change conjugation; and
every second normalized vertical system differs by a character, changing
\(\widetilde d_j(\alpha)\) by \((T_\alpha-I)a\).  Finally,
\(\ell\mapsto\ell+2\pi in\) changes a vector-cover lift of \(h_j\) by
\[
 H_j\longmapsto H'_j=\tau_{nv_j}H_j.
\]
At the factor-extension level this means
\[
 \Theta'_{j,*}
 =\operatorname{Ad}(U_{nv_j})\circ\Theta_{j,*}.
\]
Consequently the comparison character changes by
\(n\iota_{v_j}E\), because
\[
 U_{nv_j}
 \bigl(\mathbf c(\langle a_j,\lambda\rangle)U_\lambda\bigr)
 U_{nv_j}^{-1}
 =
 \mathbf c(\langle a_j+n\iota_{v_j}E,\lambda\rangle)U_\lambda.
\]
Since \(A_jv_j=v_j\) and \(E\) is \(A_j\)-invariant,
\[
 \langle T_j(\iota_{v_j}E),\lambda\rangle
 =E(v_j,A_j^{-1}\lambda)
 =E(A_jv_j,\lambda)
 =E(v_j,\lambda),
\]
so this branch change leaves both \eqref{wbII:eq:W:5.17-A14e} and \eqref{wbII:eq:W:5.17-A15} unchanged.
The translation \(\tau_{nv_j}\) also commutes with
\(Q_{j,\mathrm{lin}}\) on the vector cover because \(A_jv_j=v_j\); on the
lattice quotient it is already the identity.  Hence it changes only the
central lift choice there.
The descended \(h_j,\bar h_j\), and the quotient square are branch
independent.  Thus \eqref{wbII:eq:W:5.17-A15} is choice-independent and proves \eqref{wbII:eq:W:5.17-A5},
including its sign.

\item Since
\(A_0\widehat\delta=\widehat\delta\), equation \eqref{wbII:eq:W:5.17-A13} says that
\begin{equation}
 [G_0,U_{\widehat\delta}]
 =\mathbf c\bigl(\langle d_0,\widehat\delta\rangle\bigr).
\hypertarget{tag-03ER}{}\label{wbII:eq:W:5.17-A16}
\end{equation}
This logarithmic commutator is the clutching winding with the base loop
first.  Thus, if \(b\) and \(f\) denote the oriented base and fibre loops,
respectively,
\begin{equation}
 \operatorname{wind}_{b\wedge f}(L)
  =\langle d_0,\widehat\delta\rangle,
 \qquad
 \operatorname{wind}_{\xi_{\widehat\delta}}(L)
 =-\langle d_0,\widehat\delta\rangle.
\hypertarget{tag-03ES}{}\label{wbII:eq:W:5.17-A17}
\end{equation}
The zero-section meridian has constant fibre coordinate, so
$\Omega_E$ vanishes on its sweep
$\xi_{\widehat\delta}=a_*([S^1_{\widehat\delta}]\times[\widetilde a_0])$.
The preceding cocycle-to-edge identification therefore also gives
\begin{equation}
 \langle d_0,\widehat\delta\rangle
 =-\langle c_1(L|_{J_P})-[\Omega_E],
          \xi_{\widehat\delta}\rangle
 =-\langle c_1(L|_{J_P}),\xi_{\widehat\delta}\rangle.
\hypertarget{tag-03H5}{}\label{wbII:eq:W:5.17p}
\end{equation}
The second equality in \eqref{wbII:eq:W:5.17-A17} uses the fibre-first convention
\(\xi_{\widehat\delta}=[S^1_{\widehat\delta}]\times[\rho_0]\), hence the
single sign change.
Use the unimodular cusp chart of \SeriesRef{mapsI:src:A.1.2}
\begin{equation}
 (u,x_1,x_2)=\left(\frac{t}{x_1x_2},x_1,x_2\right),
 \qquad t=ux_1x_2.
\hypertarget{tag-03ET}{}\label{wbII:eq:W:5.17-A18}
\end{equation}
For \(0<\rho\ll1\), the boundary parametrization
\begin{equation}
 (\theta,\phi)\longmapsto
 \bigl(\rho e^{i(\theta-\phi)},1,e^{i\phi}\bigr),
 \qquad t=\rho e^{i\theta},
\hypertarget{tag-03EU}{}\label{wbII:eq:W:5.17-A19}
\hypertarget{tag-03H7}{}\SsixEquationAliases{\label{wbII:eq:W:5.17p0}}
\end{equation}
has base coordinate \(\theta\) and fibre translation
\(\widehat\delta\) in the \(\phi\)-direction. Its \(u\)-winding is
$-1$. Put \(\alpha=\theta-\phi\), let \(a_u\) be the positively
oriented \(u\)-longitude, and let \(m_{x_2}\) be the positively oriented
\(x_2\)-meridian at fixed \(u\).  In the basis
$ (a_u,m_{x_2}) $, the counterclockwise base loop $b_+$ and the fibre loop are
\begin{equation}
 b_+=a_u,
 \qquad
 f=m_{x_2}-a_u,
 \qquad
 b_+\wedge f=a_u\wedge m_{x_2}.
\hypertarget{tag-03EV}{}\label{wbII:eq:W:5.17-A19a}
\end{equation}
The determinant of this change of basis is $+1$.  The clockwise loop in
\eqref{wbII:eq:W:5.17-A17} is $b=-b_+$, so
$b\wedge f=-a_u\wedge m_{x_2}$.  The winding computed below is zero
in either orientation.  The loop $m_{x_2}$ bounds the
analytic disc
\begin{equation}
 z\longmapsto(\rho e^{i\alpha_0},1,z),
 \qquad |z|\le1.
\hypertarget{tag-03EW}{}\label{wbII:eq:W:5.17-A19b}
\end{equation}

\item Choose a thin annulus \(A_u\) about \(|u|=\rho\), a disc \(D_{x_1}\)
about \(1\), and a disc \(D_{x_2}\) containing \(\{|x_2|\le1\}\), small
enough that
\begin{equation}
 A_u\times D_{x_1}\times D_{x_2}
 \subset\{|ux_1x_2|<\epsilon\}.
\hypertarget{tag-03EX}{}\label{wbII:eq:W:5.17-A19c}
\end{equation}
Write \(U=A_u\times D_{x_1}\times D_{x_2}\) for this neighbourhood.
It is Stein and retracts onto \(S^1\), so \(H^2(U,\mathbb Z)=0\).
The exponential sequence is
\[
0\longrightarrow\mathbb Z\longrightarrow\mathcal O_U
 \xrightarrow{\,h\mapsto e^{2\pi i h}\,}\mathcal O_U^*
 \longrightarrow1.
\]
Cartan's theorem~B gives
\(H^1(U,\mathcal O_U)=H^2(U,\mathcal O_U)=0\)
\cite[\S2, fundamental theorem~B, p.~80]{Cartan1957}.
The long cohomology sequence therefore identifies
\(\operatorname{Pic}(U)=H^1(U,\mathcal O_U^*)\) with
\(H^2(U,\mathbb Z)=0\). Pull \(L|_{N_0}\) back to
the cusp cover and choose one holomorphic frame on this entire
neighbourhood.  Compute the clutching cocycle with this extending frame on
both sides of a cut in the \(u\)-longitude.  Its transition unit on
\(D_{x_2}\) is then identically \(F=1\), so by \eqref{wbII:eq:W:5.17-A19a} its
base--fibre commutator winding is
\begin{equation}
 \frac1{2\pi i}\int_{|x_2|=1}\frac{dF}{F}=0.
\hypertarget{tag-03EZ}{}\label{wbII:eq:W:5.17-A20}
\end{equation}
This computes the same alternating logarithmic cocycle as \eqref{wbII:eq:W:5.17-A16}.
Indeed, changing from the global Appell--Humbert frame to the extending
chart frame changes \(c\) by a group coboundary.  Put
\[
 g=\widehat g_0,
 \qquad
 h=\tau_{\widehat\delta}
\]
in the underlying deck group.  Equation \eqref{wbII:eq:W:5.17-A9} and
\(A_0\widehat\delta=\widehat\delta\) give \(gh=hg\).  If
\(c'=c+\delta\beta\) is the cocycle after a frame change, then
\[
 \bigl[c'(g,h)-c'(h,g)\bigr]
 -\bigl[c(g,h)-c(h,g)\bigr]
 =\beta(hg)-\beta(gh)=0.
\]
Thus the alternating cocycle is frame-invariant because $gh=hg$
in the underlying deck group.  Equations \eqref{wbII:eq:W:5.17-A17}, \eqref{wbII:eq:W:5.17-A19a}, and
\eqref{wbII:eq:W:5.17-A20} therefore give
\(\langle d_0,\widehat\delta\rangle=0\).
Finally, the matrices \(A_1,A_2\) fix
\(\widehat\delta\), so the evaluation is unchanged by every principal
change in \eqref{wbII:eq:W:5.17-A3}.  Since
\(V_G=\mathbb Z[\delta]\) and
\(\delta(\widehat\delta)=1\), its vanishing is equivalent to
\(\bar d_0=0\).

\item For \eqref{wbII:eq:W:5.17o}, take $L\in\operatorname{Pic}(N_0)$ and let $\ell=c_1(L)$, let
$\iota_0:W\hookrightarrow N_0$, and let
$\iota_t:F_t\hookrightarrow N_0$ be a sufficiently near smooth cusp fibre.
The neighbourhood realization of the nearby-cycle unit
$\mathbb Z_W\to\psi_{f_0}\mathbb Z$ in \cite[Theorem~B.1, pp.~101--105, especially (N1)--(N2), p.~101]{AC26} gives the commutative square
\begin{equation}
\begin{CD}
 H^2(N_0;\mathbb Z) @>{\iota_0^*}>> H^2(W;\mathbb Z)\\
 @V{\iota_t^*}VV @VV{\operatorname{sp}^2}V\\
 H^2(F_t;\mathbb Z) @= H^2(F_t;\mathbb Z).
\end{CD}
\hypertarget{tag-03H6}{}\label{wbII:eq:W:5.17p'}
\end{equation}
Applying this square to $\ell$ gives
$\operatorname{sp}^2(c_1(L|_W))=c_1(L|_{F_t})$.  Flat parallel transport
in the marked smooth family identifies the latter with $c_1(L|_F)$.
Equivalently, $c_1$ is the connecting morphism in the exponential sequence,
and restriction to the shrinking neighbourhoods and to $F_t$ commutes with
that morphism before passage to the direct limit.

\end{enumerate}
\end{proof}

\section{Line bundles on the toric cusp}\hypertarget{tag-03IU}{}\label{wbII:subsec:local-cusp-picard-obstruction}

\begin{proposition}[An explicit nonzero local cusp restriction]
\hypertarget{tag-03IN}{}\label{wbII:prop:local-cusp-picard-obstruction}
\StatementRecord{wbII:prop:local-cusp-picard-obstruction}{proved}{03IN}{}

Use the cusp coordinates and quotient maps of
\cref{def:II:analytic-picard-data}:
\begin{equation}
 B_0=\begin{pmatrix}0&1\\-1&0\end{pmatrix},
 \qquad
 Z(s)=sB_0+C(t),
 \qquad
 C(t)=\begin{pmatrix}6\mu(t)&h(t)\\ b(t)-h(t)&\mu(t)\end{pmatrix},
 \qquad t=e^{2\pi i s},
\hypertarget{tag-03HB}{}\label{wbII:eq:W:5.17t}
\end{equation}
\begin{equation}
 c_{\lambda}(t)=\exp(2\pi iC(t)\lambda),
 \qquad
 \exp(2\pi iZ(s)\lambda)=t^{B_0\lambda}c_\lambda(t)
 \quad(\lambda\in\bar\Lambda\simeq\mathbb Z^2),
\hypertarget{tag-03HC}{}\label{wbII:eq:W:5.17u}
\end{equation}
and
\begin{equation}
 \Phi_\lambda(x,t)=(t^{B_0\lambda}x,t),
 \qquad
 \Psi_\lambda(p)=(c_\lambda(t(p)),1)\cdot\Phi_\lambda(p),
 \qquad
 \Psi_\lambda\Psi_\mu=\Psi_{\lambda+\mu}.
\hypertarget{tag-03HD}{}\label{wbII:eq:W:5.17v}
\end{equation}
Write \(E_v\), \(v\in\mathbb Z^2\), for the toric prime divisors of
\(Y_\epsilon\), so that \(E_0\simeq\mathrm{dP}_6\), and let
\[
 \pi:Y_\epsilon\longrightarrow N_0=Y_\epsilon/\bar\Lambda,
 \qquad
 \nu=\pi|_{E_0}:E_0\longrightarrow W
\]
be the quotient and the normalization of the central fibre.
Put
\[
 N'=\mathbb Z^2\oplus\mathbb Z,
 \qquad
 M'=\operatorname{Hom}(N',\mathbb Z).
\]
For \(m=(\ell,k)\in M'=\mathbb Z^2\oplus\mathbb Z\), fix the character
convention
\begin{equation}
 \chi^m(x,t)=x^\ell t^k,
 \qquad
 \operatorname{div}(\chi^m)=
 \sum_{v\in\mathbb Z^2}\langle m,(v,1)\rangle E_v.
\hypertarget{tag-03HE}{}\label{wbII:eq:W:5.17w}
\end{equation}

Let
\begin{equation}
 H=\begin{pmatrix}1&0\\0&-6\end{pmatrix},
 \qquad
 P=B_0^{\mathsf t}H=\begin{pmatrix}0&6\\1&0\end{pmatrix},
 \qquad
 Q(v)=v^{\mathsf t}Hv=v_1^2-6v_2^2.
\hypertarget{tag-03HF}{}\label{wbII:eq:W:5.17x}
\end{equation}
Then
\begin{equation}
 \mathcal D_Q=\sum_{v\in\mathbb Z^2}Q(v)E_v
\hypertarget{tag-03HG}{}\label{wbII:eq:W:5.17y}
\end{equation}
is a locally finite Cartier divisor on \(Y_\epsilon\).
The line bundle \(\mathcal O_{Y_\epsilon}(\mathcal D_Q)\) has an explicit
\((\bar\Lambda,\Psi)\)-equivariant structure and descends to a holomorphic
line bundle \(\mathcal L_Q\) on \(N_0\).  If
\[
 (a_1,a_2,a_3,a_4,a_5,a_6)
 =(e_1,e_2,e_2-e_1,-e_1,-e_2,e_1-e_2),
 \qquad C_i=E_0\cap E_{a_i},
\]
then
\begin{equation}
 \nu^*(\mathcal L_Q|_W)
 \simeq
 \mathcal O_{E_0}\!\left(\sum_{i=1}^6Q(a_i)C_i\right),
\qquad
 \left(\deg_{C_i}\nu^*(\mathcal L_Q|_W)\right)_{i=1}^6
 =(-12,2,0,-12,2,0).
\hypertarget{tag-03HH}{}\label{wbII:eq:W:5.17z}
\end{equation}
In particular \(\nu^*(\mathcal L_Q|_W)\) is nontrivial.

On every nonzero fibre
\[
 F_t=\mathbb C^2/(Z(s)\mathbb Z^2+\mathbb Z^2),
\]
the first Chern class of \(\mathcal L_Q|_{F_t}\), in the ordered lattice
basis
\((\widehat\gamma,\widehat u,\widehat w,\widehat\delta)\) indexing the
columns of \([Z\mid I_2]\), is exactly
\begin{equation}
 c_1(\mathcal L_Q|_{F_t})=-2q,
 \qquad
 q=u\wedge w+6\gamma\wedge\delta.
\hypertarget{tag-03FW}{}\label{wbII:eq:W:5.17aa}
\end{equation}
\end{proposition}



\begin{proof}
\begin{enumerate}[label=\arabic*.,leftmargin=*]
\item \emph{Local finiteness and Cartier status.}
The fan of \(Y\) consists of the cones over the cells of the locally finite
\(A_2\)-triangulation at height one.  If
\(U_\sigma\simeq\mathbb C^3\) is the affine chart belonging to a triangle
\(\sigma=\{v_0,v_1,v_2\}\), then \(U_\sigma\) meets precisely
\(E_{v_0},E_{v_1},E_{v_2}\).  Indeed, by the orbit--cone
correspondence, \(U_\sigma\cap E_v\ne\varnothing\) exactly when
\(\mathbb R_{\geq0}(v,1)\) is a face of the cone over \(\sigma\).
Thus \(\{E_v\}_{v\in\mathbb Z^2}\) is locally finite and
\(\mathcal D_Q\) is a well-defined locally finite Weil divisor.

Every triangle of the fan is unimodular, as in
\SeriesRef{mapsI:src:A.1.2}.  If \(z_0,z_1,z_2\) are the
characters dual to
\((v_0,1),(v_1,1),(v_2,1)\), then
\begin{equation}
 \mathcal D_Q|_{U_\sigma}
 =\operatorname{div}\!
   \left(z_0^{Q(v_0)}z_1^{Q(v_1)}z_2^{Q(v_2)}\right).
\hypertarget{tag-03FX}{}\label{wbII:eq:W:5.17ab}
\end{equation}
Negative exponents merely make the displayed monomial meromorphic, as
permitted for a local equation of a Cartier divisor.  Hence
\(\mathcal D_Q\) is Cartier on every chart and therefore on
\(Y_\epsilon\).

\item \emph{The signed divisor pullback identity.}
For \(\lambda\in\bar\Lambda=\mathbb Z^2\), set
\begin{equation}
 b_\lambda=B_0\lambda,
 \qquad
 \ell_\lambda=2Hb_\lambda,
 \qquad
 k_\lambda=Q(b_\lambda),
 \qquad
 m_\lambda=(\ell_\lambda,k_\lambda)\in M'.
\hypertarget{tag-03FY}{}\label{wbII:eq:W:5.17ac}
\end{equation}
The fan automorphism sends \(E_v\) to \(E_{v+b_\lambda}\).
Consequently \(\Phi_\lambda^*E_w=E_{w-b_\lambda}\), and
\begin{align*}
 \Phi_\lambda^*\mathcal D_Q-\mathcal D_Q
 &=\sum_{v\in\mathbb Z^2}
   \bigl(Q(v+b_\lambda)-Q(v)\bigr)E_v\\
 &=\sum_{v\in\mathbb Z^2}
   \bigl(2v^{\mathsf t}Hb_\lambda+
         b_\lambda^{\mathsf t}Hb_\lambda\bigr)E_v
 =\operatorname{div}(\chi^{m_\lambda}).
\end{align*}
The \(t\)-dependent horizontal torus translation
\(p\mapsto(c_\lambda(t(p)),1)\cdot p\) preserves every
torus-invariant prime divisor \(E_v\).  Hence
\begin{equation}
 \boxed{\ 
 \Psi_\lambda^*\mathcal D_Q
 =\mathcal D_Q+\operatorname{div}(\chi^{m_\lambda})\ }.
\hypertarget{tag-03FZ}{}\label{wbII:eq:W:5.17ad}
\end{equation}

\item \emph{The complete \(\Psi\)-cocycle.}
The matrices in \eqref{wbII:eq:W:5.17t} and \eqref{wbII:eq:W:5.17x} give
\begin{equation}
 PC(t)=
 \begin{pmatrix}
 6(b(t)-h(t))&6\mu(t)\\
 6\mu(t)&h(t)
 \end{pmatrix},
\hypertarget{tag-03G0}{}\label{wbII:eq:W:5.17ae}
\end{equation}
which is symmetric.  Thus
\begin{equation}
 A(t):=2PC(t)
\hypertarget{tag-03G1}{}\label{wbII:eq:W:5.17af}
\end{equation}
is a symmetric holomorphic matrix on the full cusp disc.  Define the
holomorphic units \(h_\lambda\) and meromorphic functions \(f_\lambda\) by
\begin{equation}
 h_\lambda(t)=
 \exp\!\left(\pi i\lambda^{\mathsf t}A(t)\lambda\right),
 \qquad
 f_\lambda=h_\lambda(t)\chi^{m_\lambda}.
\hypertarget{tag-03G2}{}\label{wbII:eq:W:5.17ag}
\end{equation}
Equation \eqref{wbII:eq:W:5.17ad} gives
\begin{equation}
 \operatorname{div}(f_\lambda)
 =\Psi_\lambda^*\mathcal D_Q-\mathcal D_Q.
\hypertarget{tag-03G3}{}\label{wbII:eq:W:5.17ah}
\end{equation}

The cocharacter map underlying \(\Phi_\mu\) is
\[
 \varphi_\mu(y,y_3)=(y+y_3b_\mu,y_3),
\]
so its pullback on characters is
\begin{equation}
 \varphi_\mu^*(\ell,k)
 =(\ell,k+\ell^{\mathsf t}b_\mu).
\hypertarget{tag-03G4}{}\label{wbII:eq:W:5.17ai}
\end{equation}
The quadratic identity for \(Q\) gives
\begin{equation}
 m_{\lambda+\mu}=m_\mu+\varphi_\mu^*m_\lambda.
\hypertarget{tag-03G5}{}\label{wbII:eq:W:5.17aj}
\end{equation}
Indeed, the first coordinates add, while the last coordinate on the
right is
\[
 Q(b_\mu)+Q(b_\lambda)+2b_\lambda^{\mathsf t}Hb_\mu
 =Q(b_\lambda+b_\mu).
\]
Using \eqref{wbII:eq:W:5.17u}, \eqref{wbII:eq:W:5.17v}, and \eqref{wbII:eq:W:5.17ai}, one obtains
\begin{equation}\begin{aligned}
 \Psi_\mu^*\chi^{m_\lambda}
 &=\chi^{m_\lambda}(c_\mu(t),1)
   \chi^{\varphi_\mu^*m_\lambda}\\
 &=\exp\!\left(2\pi i\ell_\lambda^{\mathsf t}C(t)\mu\right)
   \chi^{\varphi_\mu^*m_\lambda}\\
 &=\exp\!\left(2\pi i\lambda^{\mathsf t}A(t)\mu\right)
   \chi^{\varphi_\mu^*m_\lambda}.
\end{aligned}\hypertarget{tag-03G6}{}\label{wbII:eq:W:5.17ak}
\end{equation}
The last equality is exact because
\[
 \ell_\lambda^{\mathsf t}C\mu
 =(2HB_0\lambda)^{\mathsf t}C\mu
 =2\lambda^{\mathsf t}B_0^{\mathsf t}HC\mu
 =2\lambda^{\mathsf t}PC\mu
 =\lambda^{\mathsf t}A\mu.
\]
Since \(A=A^{\mathsf t}\), \eqref{wbII:eq:W:5.17ag}, \eqref{wbII:eq:W:5.17aj}, and \eqref{wbII:eq:W:5.17ak} imply
\begin{equation}\begin{aligned}
 f_\mu\,\Psi_\mu^*f_\lambda
 &=h_\mu h_\lambda
   \exp(2\pi i\lambda^{\mathsf t}A\mu)
   \chi^{m_\mu+\varphi_\mu^*m_\lambda}\\
 &=\exp\!\left(
   \pi i(\lambda+\mu)^{\mathsf t}A(\lambda+\mu)\right)
   \chi^{m_{\lambda+\mu}}
 =f_{\lambda+\mu}.
\end{aligned}\hypertarget{tag-03G7}{}\label{wbII:eq:W:5.17al}
\end{equation}

Use the divisor-sheaf convention
\[
 \mathcal O(D)(U)=
 \{g\in\mathcal M_Y(U):
   \operatorname{div}(g)+D|_U\geq0\}.
\]
If \(D'=D+\operatorname{div}(f)\), multiplication by \(f\) maps
\(\mathcal O(D')\) isomorphically to \(\mathcal O(D)\).  Thus
\eqref{wbII:eq:W:5.17ad} and \eqref{wbII:eq:W:5.17ag} define
\begin{equation}
 \theta_\lambda:
 \Psi_\lambda^*\mathcal O(\mathcal D_Q)
 \xrightarrow{\ \times f_\lambda\ }
 \mathcal O(\mathcal D_Q).
\hypertarget{tag-03G8}{}\label{wbII:eq:W:5.17am}
\end{equation}
Because
\((\Psi_\lambda\Psi_\mu)^*=\Psi_\mu^*\Psi_\lambda^*\),
equation \eqref{wbII:eq:W:5.17al} is exactly
\begin{equation}
 \theta_{\lambda+\mu}
 =\theta_\mu\circ\Psi_\mu^*\theta_\lambda.
\hypertarget{tag-03G9}{}\label{wbII:eq:W:5.17an}
\end{equation}

\item \emph{Analytic descent.}
The cusp action of \SeriesRef{pub:I:eq:PD11}, proved in
\SeriesRef{fam:src:A.3.1}, is
free and properly discontinuous.  Every point therefore has an open
neighbourhood \(U\) whose distinct translates are disjoint and on which
\(\pi|_U\) is a biholomorphism onto \(\pi(U)\).  Transport
\(\mathcal O(\mathcal D_Q)|_U\) to \(\pi(U)\), and on overlaps arising
from translates glue by \eqref{wbII:eq:W:5.17am}.  Equation \eqref{wbII:eq:W:5.17an} is precisely the
triple-overlap condition.  The resulting holomorphic line bundle
\(\mathcal L_Q\) on \(N_0\) satisfies
\begin{equation}
 \pi^*\mathcal L_Q
 \simeq\mathcal O_{Y_\epsilon}(\mathcal D_Q).
\hypertarget{tag-03GA}{}\label{wbII:eq:W:5.17ao}
\end{equation}

\item \emph{The normalized central restriction.}
Let \(j_0:E_0\hookrightarrow Y_\epsilon\).  Since
\(\nu=\pi\circ j_0\), \eqref{wbII:eq:W:5.17ao} gives
\begin{equation}
 \nu^*(\mathcal L_Q|_W)
 =j_0^*\pi^*\mathcal L_Q
 \simeq\mathcal O_{E_0}(\mathcal D_Q|_{E_0}).
\hypertarget{tag-03GB}{}\label{wbII:eq:W:5.17ap}
\end{equation}
The only \(E_v\), \(v\ne0\), meeting \(E_0\) are those indexed by the
six \(a_i\) in the statement, and the intersections are the six toric
\((-1)\)-curves \(C_i\).  Since \(Q(0)=0\), there is no
\(E_0|_{E_0}\) term, so
\begin{equation}
 \mathcal D_Q|_{E_0}
 =\sum_{i=1}^6Q(a_i)C_i,
 \qquad
 (Q(a_1),\ldots,Q(a_6))=(1,-6,-5,1,-6,-5).
\hypertarget{tag-03GC}{}\label{wbII:eq:W:5.17aq}
\end{equation}
The hexagonal intersection table is
\begin{equation}
 C_i^2=-1,
 \qquad
 C_i\cdot C_{i\pm1}=1,
 \qquad
 C_i\cdot C_j=0
 \quad(j\notin\{i,i\pm1\}),
\hypertarget{tag-03GD}{}\label{wbII:eq:W:5.17ar}
\end{equation}
with indices modulo six.  Therefore
\[
 \deg_{C_i}\mathcal O_{E_0}
 \!\left(\sum_jQ(a_j)C_j\right)
 =-Q(a_i)+Q(a_{i-1})+Q(a_{i+1}),
\]
and direct substitution gives
\begin{equation}
 \begin{array}{c|rrrrrr}
 i&1&2&3&4&5&6\\ \hline
 -Q(a_i)+Q(a_{i-1})+Q(a_{i+1})
 &-12&2&0&-12&2&0.
 \end{array}
\hypertarget{tag-03GE}{}\label{wbII:eq:W:5.17as}
\end{equation}
The degree on \(C_1=C_{e_1}\) is \(-12\), proving nontriviality.

\item \emph{The generic-fibre factor and its exact Chern class.}
Fix \(t\ne0\), choose \(s\) with \(t=e^{2\pi i s}\), and write
\(x=e^{2\pi i\zeta}\), \(\zeta\in\mathbb C^2\).
Equation \eqref{wbII:eq:W:5.17u} identifies the deck transformation indexed by
\(\lambda\in\mathbb Z^2\) with
\[
 \zeta\longmapsto\zeta+Z(s)\lambda
\]
modulo the horizontal translations
\(\zeta\mapsto\zeta+n\), \(n\in\mathbb Z^2\).  The matrix identities
\begin{equation}
 HB_0=P^{\mathsf t},
 \qquad
 B_0^{\mathsf t}HB_0=PB_0
\hypertarget{tag-03GF}{}\label{wbII:eq:W:5.17at}
\end{equation}
give
\[
 \chi^{m_\lambda}(x,t)
 =x^{2P^{\mathsf t}\lambda}
  t^{\lambda^{\mathsf t}PB_0\lambda},
\qquad
 h_\lambda(t)
 =\exp(2\pi i\lambda^{\mathsf t}PC(t)\lambda).
\]
Since \(Z=sB_0+C\), it follows that
\begin{equation}
 f_\lambda(\zeta)
 =\exp\!\left(
 2\pi i\lambda^{\mathsf t}PZ\lambda+
 4\pi i\lambda^{\mathsf t}P\zeta\right).
\hypertarget{tag-03GG}{}\label{wbII:eq:W:5.17au}
\end{equation}
The matrix \(PZ=sPB_0+PC\) is symmetric:
\[
 PB_0=\operatorname{diag}(-6,1),
\]
and \(PC\) is symmetric by \eqref{wbII:eq:W:5.17ae}.

Under the descent convention \eqref{wbII:eq:W:5.17am}, a pulled-back local section
written as \(\varphi(\zeta)\) in the canonical frame satisfies
\[
 \varphi(\zeta+Z\lambda)
 =f_\lambda(\zeta)^{-1}\varphi(\zeta).
\]
Thus the factor of automorphy for the lattice translation
\(\zeta\mapsto\zeta+Z\lambda+n\) is
\begin{equation}
 a_{(\lambda,n)}(\zeta)
 =\exp(2\pi i g_{(\lambda,n)}(\zeta)),
\qquad
 g_{(\lambda,n)}(\zeta)
 =-\lambda^{\mathsf t}PZ\lambda
  -2\lambda^{\mathsf t}P\zeta.
\hypertarget{tag-03GH}{}\label{wbII:eq:W:5.17av}
\end{equation}
This is independent of \(n\) as a factor of automorphy: replacing
\(\zeta\) by \(\zeta+n\) changes \(g\) by the integer
\(-2\lambda^{\mathsf t}Pn\).

For \(l=(\lambda,n)\), \(l'=(\mu,m)\), the exponential connecting
cocycle is
\begin{equation}\begin{aligned}
 c(l,l')
 &=g_l(\zeta+Z\mu+m)+g_{l'}(\zeta)-g_{l+l'}(\zeta)\\
 &=-2\lambda^{\mathsf t}Pm\in\mathbb Z,
\end{aligned}\hypertarget{tag-03GI}{}\label{wbII:eq:W:5.17aw}
\end{equation}
where the symmetry of \(PZ\) cancels the two mixed \(Z\)-terms.
The corresponding alternating integral form is
\begin{equation}
 E(l,l')=c(l,l')-c(l',l)
 =-2\lambda^{\mathsf t}Pm+2\mu^{\mathsf t}Pn.
\hypertarget{tag-03GJ}{}\label{wbII:eq:W:5.17ax}
\end{equation}
This is \(c_1(\mathcal L_Q|_{F_t})\).  In the ordered column basis of
\([Z\mid I_2]\), the cross-block of the class
\(q=u\wedge w+6\gamma\wedge\delta\) from
\((\widehat\gamma,\widehat u)\) to
\((\widehat w,\widehat\delta)\) is
\begin{equation}
 \begin{pmatrix}
 q(\widehat\gamma,\widehat w)&
 q(\widehat\gamma,\widehat\delta)\\
 q(\widehat u,\widehat w)&
 q(\widehat u,\widehat\delta)
 \end{pmatrix}
 =\begin{pmatrix}0&6\\1&0\end{pmatrix}=P.
\hypertarget{tag-03GK}{}\label{wbII:eq:W:5.17ay}
\end{equation}
Hence \eqref{wbII:eq:W:5.17ax} is \(-2q\).

\end{enumerate}
\end{proof}

\begin{lemma}[The toric support function detects central restriction]
\hypertarget{tag-03I6}{}\label{wbII:lem:analytic-support-specialization}
\StatementRecord{wbII:lem:analytic-support-specialization}{proved}{03I6}{}

Let \(M\in\operatorname{Pic}(N_0)\), let
\(\pi:Y_\epsilon\to N_0\) be the cusp covering, and let
\(\nu:E_0\to W\) be the normalization.  If the Appell--Humbert form of
\(M|_{F_t}\), \(t\ne0\), is \(E_M\), then the three independent degrees
of \(\nu^*(M|_W)\) determine an integral symmetric form
\(B:\mathbb Z^2\times\mathbb Z^2\to\mathbb Z\) satisfying
\begin{equation}
 E_M\bigl((\lambda,0),(0,m)\bigr)
 =-B(B_0\lambda,m)
 \qquad(\lambda,m\in\mathbb Z^2).
\hypertarget{tag-03F0}{}\label{wbII:eq:W:5.17-A21}
\end{equation}
Consequently
\begin{equation}
 E_M=0
 \quad\Longrightarrow\quad
 \nu^*(M|_W)\simeq\mathcal O_{E_0}.
\hypertarget{tag-03F1}{}\label{wbII:eq:W:5.17-A22}
\end{equation}
\end{lemma}



\begin{proof}

\begin{enumerate}[label=\arabic*.,leftmargin=*]
\item Put \(\mathcal M=\pi^*M\).  The identity
\(\pi\circ\Psi_\lambda=\pi\) gives \(\mathcal M\) its \(\Psi\)-equivariant descent structure.  In the notation of
\cref{wbII:prop:local-cusp-picard-obstruction}, order the six rays in the star
of the origin as
\begin{equation}
 (r_1,r_2,r_3,r_4,r_5,r_6)
 =(e_1,e_2,e_2-e_1,-e_1,-e_2,e_1-e_2),
 \qquad C_i=E_0\cap E_{r_i}.
\hypertarget{tag-03F2}{}\label{wbII:eq:W:5.17-A23}
\end{equation}
The quotient identifies opposite conductor branches.  Therefore, for
\(A=\mathcal M|_{E_0}=\nu^*(M|_W)\),
\begin{equation}
 d_i=\deg_{C_i}A,
 \qquad
 (\deg_{C_1}A,\ldots,\deg_{C_6}A)
 =(d_1,d_2,d_3,d_1,d_2,d_3).
\hypertarget{tag-03F3}{}\label{wbII:eq:W:5.17-A24}
\end{equation}

\item Define the integer-valued quadratic function
\begin{equation}
 a(x,y)
 =(d_2+d_3)\binom{x}{2}+d_3xy+(d_1+d_3)\binom{y}{2}.
\hypertarget{tag-03F4}{}\label{wbII:eq:W:5.17-A25}
\end{equation}
Its polarization is
\begin{equation}
 B(v,w)=a(v+w)-a(v)-a(w),
 \qquad
 [B]=
 \begin{pmatrix}
  d_2+d_3&d_3\\ d_3&d_1+d_3
 \end{pmatrix}.
\hypertarget{tag-03F5}{}\label{wbII:eq:W:5.17-A26}
\end{equation}
The locally finite divisor
\begin{equation}
 D_a=\sum_{v\in\mathbb Z^2}a(v)E_v
\hypertarget{tag-03F6}{}\label{wbII:eq:W:5.17-A27}
\end{equation}
is Cartier: on every maximal chart, the three height-one ray generators
form an integral basis, so the three assigned coefficients are the values
of one integral character.  In the standard toric convention the associated
support function has ray values \(-a(v)\); the divisor coefficients here are
the values \(a(v)\), consistent with
\(\operatorname{div}(\chi^m)=\sum_\rho\langle m,u_\rho\rangle D_\rho\)
\cite[2010 text, Propositions~4.1.1--4.1.2, p.~171; Theorems~4.2.8,
p.~181 and~4.2.12, p.~184]{CoxLittleSchenck2011}.
This is the same chartwise Cartier argument as
in \eqref{wbII:eq:W:5.17ab}, and it applies to an arbitrary integer-valued \(a\).

\item Direct substitution gives
\begin{equation}
 (a(r_1),\ldots,a(r_6))
 =(0,0,d_2,d_2+d_3,d_1+d_3,d_1).
\hypertarget{tag-03F7}{}\label{wbII:eq:W:5.17-A28}
\end{equation}
Using \(C_i^2=-1\) and \(C_iC_{i\pm1}=1\), the degree vector of
\(\mathcal O(D_a)|_{E_0}\) is exactly the vector in \eqref{wbII:eq:W:5.17-A24}.  More
generally, restrict the principal identity
\(\operatorname{div}(t)=\sum_wE_w\) to \(E_v\).  Since \(E_v\) meets
precisely \(E_{v+r_i}\), it gives
\begin{equation}
 \mathcal O(D_a)|_{E_v}
 \simeq
 \mathcal O_{E_v}\!\left(
   \sum_{i=1}^6(a(v+r_i)-a(v))C_{v,i}
 \right).
\hypertarget{tag-03F8}{}\label{wbII:eq:W:5.17-A29}
\end{equation}
The coefficient difference is
\begin{equation}
 a(v+r_i)-a(v)=a(r_i)+B(v,r_i).
\hypertarget{tag-03FA}{}\label{wbII:eq:W:5.17-A30}
\end{equation}
The second summand in \eqref{wbII:eq:W:5.17-A30} is the divisor of the toric character
with exponent \(B(v,-)\).  Thus every restriction in \eqref{wbII:eq:W:5.17-A29} has the
same six degrees as \(\mathcal O(D_a)|_{E_0}\).

\item Because \(B_0\bar\Lambda=\mathbb Z^2\), the deck action is transitive on
the components \(E_v\).  The descent structure on \(\mathcal M\)
therefore transports \eqref{wbII:eq:W:5.17-A24}, with its ordered six boundary directions,
to every \(E_v\).  It follows that
\begin{equation}
 \mathcal R:=\mathcal M\otimes\mathcal O(-D_a)
 \quad\text{satisfies}\quad
 \mathcal R|_{E_v}\simeq\mathcal O_{E_v}
 \quad(v\in\mathbb Z^2).
\hypertarget{tag-03FB}{}\label{wbII:eq:W:5.17-A31}
\end{equation}
Indeed, the six toric curves detect \(\operatorname{Pic}(\mathrm{dP}_6)\):
in the blow-up basis \(H,E_1,E_2,E_3\), vanishing of the degrees on the
three exceptional curves kills the three \(E_i\)-coefficients, and
vanishing on any strict transform \(H-E_i-E_j\) kills the remaining
\(H\)-coefficient.  Since
\(H^1(\mathrm{dP}_6,\mathcal O)=0\), equality of those degrees implies
isomorphism.

\item For \(b_\lambda=B_0\lambda\), the torus-translation part of
\(\Psi_\lambda\) preserves each toric prime divisor, while its fan part
translates the index by \(b_\lambda\).  Hence
\begin{equation}
 \Psi_\lambda^*D_a-D_a
 =\sum_v(a(v+b_\lambda)-a(v))E_v
 =\operatorname{div}(\chi^{m_\lambda}),
\hypertarget{tag-03FC}{}\label{wbII:eq:W:5.17-A32}
\end{equation}
where
\begin{equation}
 m_\lambda=(\ell_\lambda,k_\lambda),
 \qquad
 \ell_\lambda=B(b_\lambda,-),
 \qquad
 k_\lambda=a(b_\lambda).
\hypertarget{tag-03FD}{}\label{wbII:eq:W:5.17-A33}
\end{equation}
Thus multiplication by \(\chi^{m_\lambda}\) gives an individual
isomorphism
\(\Psi_\lambda^*\mathcal O(D_a)\simeq\mathcal O(D_a)\).
These individual isomorphisms separate the horizontal winding of the
descent factor of $\mathcal M$ in the tensor-product map below.

\item For a ray \(\mathbb R_{\ge0}(v,1)\), its toric open set in
\(Y_\epsilon\) is
\begin{equation}
 U_v\simeq\Delta_\epsilon\times(\mathbb C^*)^2,
 \qquad E_v^\circ=\{0\}\times(\mathbb C^*)^2.
\hypertarget{tag-03FE}{}\label{wbII:eq:W:5.17-A34}
\end{equation}
Indeed \((e_1,0),(e_2,0),(v,1)\) is an integral basis and \(t\) is the
character dual to the last vector; compare
\cite[2010 text, Example~1.2.21, p.~31]{CoxLittleSchenck2011}.
The set \(U_v\) is Stein and retracts
onto \(E_v^\circ\).  By \eqref{wbII:eq:W:5.17-A31},
\(c_1(\mathcal R|_{U_v})=0\); the exponential sequence and Cartan's
theorem~B therefore trivialize \(\mathcal R|_{U_v}\).  Choose a global
nowhere-zero frame \(s_v\) of \(\mathcal R|_{E_v}\).  A frame on \(U_v\) can be
multiplied by a unit pulled back from \((\mathbb C^*)^2\) so that its
restriction to \(E_v^\circ\) agrees with \(s_v|_{E_v^\circ}\); denote the
resulting frame by \(e_v\).

\item For the comparison of these frames on the generic torus, let \(v,w\) be
adjacent vertices and choose a triangle \(\sigma=\{v,w,u\}\) containing them,
write \(z_v,z_w,z_u\) for the coordinates dual to its three height-one ray
generators, and put
\begin{equation}
 S_\sigma:=U_\sigma\cap Y_\epsilon
 =\{(z_v,z_w,z_u)\in\mathbb C^3:
       |z_vz_wz_u|<\epsilon\}.
\hypertarget{tag-03FF}{}\label{wbII:eq:W:5.17-A34a}
\end{equation}
Because \(\rho_v<\sigma\), localization of the maximal affine chart gives
\begin{equation}
 U_{\rho_v}=D(z_wz_u)\subset U_\sigma,\qquad
 U_{\rho_w}=D(z_vz_u)\subset U_\sigma,\qquad
 U_{\rho_v}\cap U_{\rho_w}=D(z_vz_wz_u)=U_0.
\hypertarget{tag-03I5}{}\label{wbII:eq:analytic-toric-open-directions}
\end{equation}
The ratio \(e_w/e_v\) is defined on \(U_0\cap Y_\epsilon\).
We compare it with a frame on \(S_\sigma\) as follows.

\item The space $S_\sigma$ is the closed graph
\[
 \{(z_v,z_w,z_u,t)\in\mathbb C^3\times\Delta_\epsilon:
   t=z_vz_wz_u\}
\]
in the Stein space \(\mathbb C^3\times\Delta_\epsilon\), and hence is
Stein.  Radial scaling of the three \(z\)-coordinates contracts it to the
origin.  The exponential sequence and Cartan's theorem~B therefore give
\(\operatorname{Pic}(S_\sigma)=0\).  Choose a frame \(e_\sigma\) of
\(\mathcal R|_{S_\sigma}\).

\item Let \(s_v\) be the global frame of
\(\mathcal R|_{E_v}\simeq\mathcal O_{E_v}\) used to normalize \(e_v\).
On \(U_{\rho_v}\cap Y_\epsilon\) put
\[
 A_{\sigma,v}=e_v/e_\sigma.
\]
At \(t=0\), this is the restriction to \((\mathbb C^*)^2\) of the
holomorphic unit
\[
 G_{\sigma,v}
 =\frac{s_v}{e_\sigma|_{E_v\cap S_\sigma}}
 \in\mathcal O^*(E_v\cap S_\sigma),
 \qquad E_v\cap S_\sigma=\{z_v=0\}\simeq\mathbb C^2.
\]
Consequently its winding on every horizontal cocharacter is zero at
\(t=0\).  Under
\(U_{\rho_v}\cap Y_\epsilon\simeq
\Delta_\epsilon\times(\mathbb C^*)^2\), that winding is an integer locally
constant in \(t\), and is therefore zero on every nearby nonzero fibre.
The same argument applies to \(A_{\sigma,w}\).  On their common
domain \(U_0\cap Y_\epsilon\),
\[
 e_w/e_v=A_{\sigma,w}/A_{\sigma,v},
\]
so this ratio has zero horizontal winding.  The one-skeleton of the
\(A_2\)-triangulation is connected, and multiplication along an edge path
gives
\begin{equation}
 \operatorname{wind}_m(e_w/e_v)=0
 \qquad(v,w,m\in\mathbb Z^2)
\hypertarget{tag-03FG}{}\label{wbII:eq:W:5.17-A34b}
\end{equation}
on every nonzero fibre.

\item Let
\(\theta^{\mathcal M}_\lambda:\Psi_\lambda^*\mathcal M\to\mathcal M\)
be the descent isomorphism, let
\(\mu^+_\lambda:\Psi_\lambda^*\mathcal O(D_a)\to\mathcal O(D_a)\)
be multiplication by \(\chi^{m_\lambda}\) from \eqref{wbII:eq:W:5.17-A32}, and define the
inverse-dual isomorphism
\[
 \mu^-_\lambda:=((\mu^+_\lambda)^{-1})^\vee:
 \Psi_\lambda^*\mathcal O(-D_a)\longrightarrow\mathcal O(-D_a),
 \qquad s\longmapsto\chi^{-m_\lambda}s.
\]
The ordinary dual of \(\mu^+_\lambda\) has the opposite arrow; the
inverse-dual above is the map required by the pullback tensor product.  Then
\[
 \theta^{\mathcal R}_\lambda
 :=\theta^{\mathcal M}_\lambda\otimes\mu^-_\lambda:
 \Psi_\lambda^*\mathcal R\longrightarrow\mathcal R.
\]
Since \(\Psi_\lambda(U_v)=U_{v+b_\lambda}\), define
\(r_{\lambda,v}\in\mathcal O^*(U_v)\) by the identity
\begin{equation}
 \theta^{\mathcal R}_\lambda
   (\Psi_\lambda^*e_{v+b_\lambda})
 =r_{\lambda,v}e_v.
\hypertarget{tag-03I4}{}\label{wbII:eq:analytic-residual-descent-factor}
\end{equation}
At \(t=0\) it extends to an isomorphism between two trivial line bundles
on the compact surface \(E_v\).  Its factor in the chosen global frames is
therefore a nonzero constant.  The winding of
\(r_{\lambda,v}(t,-)\) around every horizontal one-parameter subgroup is
an integer locally constant in \(t\); at \(t=0\) it is zero, hence it is
zero for every nearby \(t\ne0\).

\item To express this isomorphism in one frame on the generic torus, put
\(q_{\lambda,v}=e_{v+b_\lambda}/e_v\).  Equation \eqref{wbII:eq:W:5.17-A34b} gives zero
horizontal winding for \(q_{\lambda,v}\); because
\(\Psi_\lambda|_{Y_t}\) is multiplication by a fixed torus point, it
induces the identity on horizontal \(H_1\), and
\(\Psi_\lambda^*q_{\lambda,v}\) also has zero winding.  Relative to the
single frame \(e_v\), the residual multiplier is therefore
\begin{equation}
 h_{\lambda,v}
 =\frac{r_{\lambda,v}}{\Psi_\lambda^*q_{\lambda,v}},
 \qquad
 \operatorname{wind}_m(h_{\lambda,v})=0.
\hypertarget{tag-03FH}{}\label{wbII:eq:W:5.17-A34c}
\end{equation}

\item On \(U_0\) let \(\mathbf1_a\) be the canonical toric frame of
\(\mathcal O(D_a)\), and put \(F_v=e_v\otimes\mathbf1_a\).  Equations
\eqref{wbII:eq:analytic-residual-descent-factor} and \eqref{wbII:eq:W:5.17-A32} give
\[
 \theta^{\mathcal M}_\lambda(\Psi_\lambda^*F_v)
 =h_{\lambda,v}\chi^{m_\lambda}F_v.
\]
On a nonzero fibre, \(\chi^{m_\lambda}=x^{\ell_\lambda}t^{k_\lambda}\).
With the section-descent convention \eqref{wbII:eq:W:5.17am}, the factor of automorphy is
therefore
\[
 \alpha_\lambda=(h_{\lambda,v}\chi^{m_\lambda})^{-1}.
\]
The residual factor has zero horizontal winding by \eqref{wbII:eq:W:5.17-A34c}, and hence
\begin{equation}
 \operatorname{wind}_m(\alpha_\lambda)
 =-\ell_\lambda(m)=-B(B_0\lambda,m).
\hypertarget{tag-03I3}{}\label{wbII:eq:analytic-factor-horizontal-winding}
\end{equation}

\item This winding is exactly the mixed Chern-class entry.  Pull \(F_v\) back to
the universal cover
\[
 p_t:\mathbb C^2\longrightarrow Y_t,
 \qquad \zeta\longmapsto(e^{2\pi i\zeta},t).
\]
The period indexed by \((\lambda,n)\in\mathbb Z^2\oplus\mathbb Z^2\)
acts by \(\zeta\mapsto\zeta+Z(s)\lambda+n\).  The horizontal factor is one,
because \(F_v\) is already defined on \(Y_t\).  Choose a logarithm
\(g_\lambda\in\mathcal O(\mathbb C^2)\) with
\[
 p_t^*\alpha_\lambda=e^{2\pi i g_\lambda}.
\]
For every horizontal period \(n\),
\[
 g_\lambda(\zeta+n)-g_\lambda(\zeta)
 =\operatorname{wind}_n(\alpha_\lambda).
\]
In the exponential connecting cocycle
\[
 c(l,l')=g_l(\zeta+l')+g_{l'}(\zeta)-g_{l+l'}(\zeta),
\]
take the horizontal logarithm to be zero and the logarithm for the combined
period \((\lambda,n)\) to be \(g_\lambda\).  Then
\[
 c((\lambda,0),(0,n))
 =\operatorname{wind}_n(\alpha_\lambda),
 \qquad
 c((0,n),(\lambda,0))=0.
\]
Therefore the alternating exponential connecting integer between
the vertical deck
translation indexed by \(\lambda\) and the horizontal loop indexed by
\(m\) is
\[
 E_M\bigl((\lambda,0),(0,m)\bigr)
 =-\ell_\lambda(m)
 =-B(B_0\lambda,m),
\]
which proves \eqref{wbII:eq:W:5.17-A21}.  This sign also recovers the explicit check in
\cref{wbII:prop:local-cusp-picard-obstruction}: for
\(a(v)=v^{\mathsf t}Hv\), one has \(B=2H\), so \eqref{wbII:eq:W:5.17-A21} is the
cross-block \(-2B_0^{\mathsf t}H=-2P\), namely \(-2q\).

\item If \(E_M=0\), equation \eqref{wbII:eq:W:5.17-A21} and the unimodularity of \(B_0\) force
\(B=0\).  The matrix in \eqref{wbII:eq:W:5.17-A26} then gives
\(d_3=0\), \(d_2=0\), and \(d_1=0\).  Thus all six degrees of \(A\)
vanish, and the Picard detection argument following \eqref{wbII:eq:W:5.17-A31} gives
\(A\simeq\mathcal O_{E_0}\).

\end{enumerate}
\end{proof}

\section{Global Picard restrictions}\hypertarget{tag-03IT}{}\label{wbII:subsec:analytic-all-picard}

\begin{theorem}[Global restriction of line bundles]
\hypertarget{tag-03IW}{}\label{wbII:thm:analytic-all-picard}
\StatementRecord{wbII:thm:analytic-all-picard}{proved}{03IW}{}
\hypertarget{tag-03IM}{}\label{wbII:prop:leray-edge-restriction}
\StatementRecord{wbII:prop:leray-edge-restriction}{proved}{03IM}{}
\hypertarget{tag-03IL}{}\label{wbII:prop:direct-boundary-picard}
\StatementRecord{wbII:prop:direct-boundary-picard}{proved}{03IL}{}

For the standard threefold $f:X\to B$ and every
$L\in\operatorname{Pic}(X)$, one has
\begin{equation}
 \begin{gathered}
 e(c_1(L))=0\in H^0(B,R^2f_*\mathbb Z),
 \qquad c_1(L|_F)=0,\\
 \eta^*i^*L\simeq\mathcal O_{\mathrm{dP}_6}.
 \end{gathered}
\hypertarget{tag-03GM}{}\label{wbII:eq:W:5.17c}
\hypertarget{tag-03H8}{}\SsixEquationAliases{\label{wbII:eq:W:5.17q}}
\end{equation}
In particular $r=\eta^*i^*$ is the zero homomorphism.  The normalized
restriction conclusion is the assertion of
\citet[Theorem~10.5, proof Step~1, equation~(10.5), p.~88]{AC26}.
\end{theorem}

\begin{proof}
\begin{enumerate}[label=\arabic*.,leftmargin=*]
\item By \cite[Proposition~7.26, pp.~61--62]{AC26}, the invariant fibre
class has the form
\begin{equation}
 E=c_1(L|_F)=kq,
 \qquad k\in\mathbb Z,
 \qquad q=uw+6\gamma\delta.
\hypertarget{tag-03FI}{}\label{wbII:eq:W:5.17-A35}
\end{equation}
The central-cover identity \eqref{wbII:eq:W:5.17l'} places $E$ in the
order-four specialization lattice
\[
 \pi_2^*H^2(S_2;\mathbb Z)
 =\{a\gamma\eta_2+bq:a\equiv b\pmod2\}
\]
of \cite[Lemma~A.4, pp.~99--100]{AC26}.  Independence of
$\gamma\eta_2,q$ gives $a=0$, $b=k$, so $k=2n$ for an integer $n$.
In the basis dual to $(\gamma,u,w,\delta)$, the translation vectors
are $v_1=(1,2,-4,0)^{\mathsf t}$ and
$v_2=(-1,-3,3,0)^{\mathsf t}$.  With
$\eta_1=2u+w+3\delta$ and $\eta_2=u+w+2\delta$, contraction gives
\begin{equation}
 \iota_{v_1}q=2\eta_1,
 \qquad
 \iota_{v_2}q=-3\eta_2,
 \qquad
 [\eta_1]=3[\delta],
 \qquad
 [\eta_2]=2[\delta].
\hypertarget{tag-03FJ}{}\label{wbII:eq:W:5.17-A36}
\hypertarget{tag-03H9}{}\SsixEquationAliases{\label{wbII:eq:W:5.17r}}
\end{equation}
Contraction is linear in the form, so the first two identities also give
$\iota_{v_1}(2q)=4\eta_1$ and $\iota_{v_2}(2q)=-6\eta_2$.
Applying \eqref{wbII:eq:W:5.17-A5} with $(m_1,m_2)=(3,4)$ yields
\begin{equation}
 \bar d_1=2k[\delta]=4n[\delta],
 \qquad
 \bar d_2=-\frac32k[\delta]=-3n[\delta].
\hypertarget{tag-03FK}{}\label{wbII:eq:W:5.17-A37}
\hypertarget{tag-03HA}{}\SsixEquationAliases{\label{wbII:eq:W:5.17s}}
\end{equation}

\item The balance and cusp identities
\eqref{wbII:eq:W:5.17-A4} and \eqref{wbII:eq:W:5.17-A6} give
\begin{equation}
 0=\bar d_1+\bar d_2+\bar d_0
   =\frac{k}{2}[\delta].
\hypertarget{tag-03FL}{}\label{wbII:eq:W:5.17-A38}
\end{equation}
Since $V_G\otimes\mathbb R=\mathbb R[\delta]$, one has $k=0$ and hence
$E=0$.  Apply \cref{wbII:lem:analytic-support-specialization} to
$M=L|_{N_0}$.  The normalization map there is
$\nu=\eta:E_0=\mathrm{dP}_6\to W$, and restriction through $N_0\subset X$
gives $\nu^*(M|_W)=\eta^*i^*L$.  Its triviality proves $r(L)=0$.

\item For every $b\in B$,
proper base change and functoriality of the filtered-complex spectral
sequence give a commutative square
\begin{equation}
\begin{CD}
 H^2(X;\mathbb Z) @>{e}>> H^0(B,R^2f_*\mathbb Z)\\
 @V{i_b^*}VV @VV{\operatorname{ev}_b}V\\
 H^2(X_b;\mathbb Z) @= H^2(X_b;\mathbb Z).
\end{CD}
\hypertarget{tag-03GN}{}\label{wbII:eq:W:5.17d}
\end{equation}
Indeed, pull the Leray filtered complex back along $\{b\}\to B$; the
spectral sequence over a point has only column zero, so its edge is the
identity, while the right vertical map is the proper-base-change
identification.  This uses the functoriality and convergence statements in
\cite[Tags~012O, 012U and~012W]{StacksProject} and topological proper base
change \cite[Tag~09V6]{StacksProject}.  The connecting homomorphism in the
exponential sequence is natural under $i_b^*$, so \eqref{wbII:eq:W:5.17d} sends
$c_1(L)$ to $c_1(L|_{X_b})$.


Put $s_L=e(c_1(L))$.  The edge image consists of permanent cycles, giving
\begin{equation}
 s_L\in E_\infty^{0,2}\subseteq\ker
 \bigl(d_2^{0,2}:E_2^{0,2}\to E_2^{2,1}\bigr);
\hypertarget{tag-03GO}{}\label{wbII:eq:W:5.17e}
\end{equation}
by \cite[Tags~011N, 011O and~012U]{StacksProject}.
The vectors $v_1,v_2$ displayed in the first step are exactly the standard
pair in \cite[(A.5), p.~95]{AC26}, explicitly included under the standing
hypotheses of its \S7.7, p.~61.  Proposition~7.26 there identifies
$R^2f_*\mathbb Z$ with a subsheaf of $j_*\mathcal L^2$, where
$j:B^\circ\hookrightarrow B$ and $\mathcal L^2$ is the degree-two
local system on the smooth locus.  Thus restriction and evaluation at a
marked smooth fibre give an injection
\[
 H^0(B,R^2f_*\mathbb Z)\hookrightarrow
 H^0(B^\circ,\mathcal L^2)
 \hookrightarrow\bigwedge^2V.
\]
Under \eqref{wbII:eq:W:5.17d}, the image of $s_L$ is exactly $E=0$.
Injectivity therefore gives $s_L=0$.  At $b=p_0$ the same square also
gives $c_1(L|_W)=0$.
\end{enumerate}
\end{proof}

\begin{remark}[A local bundle without a global extension]
\hypertarget{tag-03IQ}{}\label{wbII:rem:local-cusp-picard-obstruction}
\StatementRecord{wbII:rem:local-cusp-picard-obstruction}{proved}{03IQ}{}
Restriction to the cusp neighbourhood factors the global restriction as
\[
 \operatorname{Pic}(X)\longrightarrow\operatorname{Pic}(N_0)
 \xrightarrow{\,M\mapsto\nu^*(M|_W)\,}
 \operatorname{Pic}(E_0),
 \qquad E_0=\mathrm{dP}_6,
\]
because both composites pull back along
$E_0\xrightarrow{\nu}W\hookrightarrow N_0\hookrightarrow X$.
The second homomorphism sends the bundle $\mathcal L_Q$ of
\cref{wbII:prop:local-cusp-picard-obstruction} to the class with degrees
$(-12,2,0,-12,2,0)$, so that homomorphism is nonzero.
By \cref{wbII:thm:analytic-all-picard}, its composite from
$\operatorname{Pic}(X)$ is zero.  In fact $\mathcal L_Q$ lies outside the
image of $\operatorname{Pic}(X)\to\operatorname{Pic}(N_0)$: its restriction
to a nearby fibre has Chern class $-2q\ne0$ by
\eqref{wbII:eq:W:5.17aa}, whereas the restriction of every global bundle
has class zero by \eqref{wbII:eq:W:5.17q}.
\end{remark}

\section{Tangent cohomology and vanishing criteria}\hypertarget{tag-P315}{}\label{sec:II:analytic-vanishing}

\begin{theorem}[Conductor nonvanishing for the standard family]
\hypertarget{tag-03IX}{}\label{wbII:thm:candidate-conductor}
\StatementRecord{wbII:thm:candidate-conductor}{proved}{03IX}{}

The conductor argument is that of \citet[Lemma~10.3 and Theorem~10.5(a), pp.~87--89]{AC26}.
Assume:
\begin{enumerate}[label=(H\arabic*)]
\item $X$ is a connected compact smooth complex threefold and
$f:X\to B$ is a nonconstant holomorphic map to a smooth compact curve;
\item $W=f^{-1}(p_0)$ is the full scheme-theoretic fibre, is reduced, and
has the normal-crossings local models \eqref{wbII:eq:W:5.3};
\item the normalization $\eta:\widetilde W\to W$ is the degree-six del
Pezzo surface $\mathrm{dP}_6$;
\item $H^2(X;\mathbb Z)=0$.
\end{enumerate}
Then, for every $L\in\operatorname{Pic}(X)$,
\begin{equation}
H^2(W,(T_X\otimes L)|_W)\ne0
\quad\text{and}\quad
R^2f_*(T_X\otimes L)\ne0.
\hypertarget{tag-03HI}{}\label{wbII:eq:W:5.18}
\end{equation}
\end{theorem}



\begin{proof}

\begin{enumerate}[label=\arabic*.,leftmargin=*]
\item Fix $L$ and put
\[
\mathcal F_L=T_X\otimes L,\qquad
A_L=(L^\vee\otimes\omega_X)|_W.
\]
Hypothesis (H4) gives $c_1(L^\vee\otimes\omega_X)=0$.  Pullback therefore
gives $c_1(\eta^*A_L)=0$. The exponential-sequence calculation in the proof of
\cref{wbII:prop:pic0-pullback} makes
$c_1:\operatorname{Pic}(\mathrm{dP}_6)\to H^2(\mathrm{dP}_6;\mathbb Z)$
injective.  Apply that map to the bundle $\eta^*A_L$.
Consequently
\begin{equation}
\eta^*A_L\simeq\mathcal O_{\widetilde W}.
\hypertarget{tag-03HJ}{}\label{wbII:eq:W:5.19}
\end{equation}
\Cref{wbII:prop:conductor} gives
\begin{equation}
H^0(W,\Omega^1_X|_W\otimes A_L)\ne0.
\hypertarget{tag-03HK}{}\label{wbII:eq:W:5.20}
\end{equation}

\item Because $W=f^*(p_0)$, one has
$\mathcal O_X(W)=f^*\mathcal O_B(p_0)$, whose restriction to the fibre is
trivial.  Adjunction therefore gives $\omega_W\simeq\omega_X|_W$.
Apply \cref{wbII:thm:cartier-duality} to $E=\mathcal F_L$:
\[
\begin{split}
H^2(W,\mathcal F_L|_W)^\vee
&\simeq H^0(W,(\mathcal F_L|_W)^\vee\otimes\omega_W)\\
&=H^0(W,\Omega^1_X|_W\otimes A_L).
\end{split}
\]
By \eqref{wbII:eq:W:5.20}, the dual space is nonzero, giving the first nonvanishing in \eqref{wbII:eq:W:5.18}.
The same duality theorem gives
$H^3(W,\mathcal F_L|_W)=0$.

\item Since $X$ is compact, $f$ is proper.  By
\cref{wbII:lem:curve-flatness}, $f$ is flat and the vector bundle
$\mathcal F_L$ is $f$-flat.  Thus every hypothesis of
\cref{wbII:prop:analytic-base-change} is satisfied at $p_0$, and
\[
(R^2f_*\mathcal F_L)_{p_0}\otimes\mathbb C(p_0)
\simeq H^2(W,\mathcal F_L|_W)\ne0.
\]
The stalk, and hence the coherent direct-image sheaf in
\eqref{wbII:eq:W:5.18}, is nonzero.

\end{enumerate}
\end{proof}

\begin{corollary}[Tangent-cohomology nonvanishing on the standard fibre]
\hypertarget{tag-03DR}{}\label{wbII:cor:candidate-conductor-actual}
\StatementRecord{wbII:cor:candidate-conductor-actual}{proved}{03DR}{}
\hypertarget{tag-03DS}{}\label{wbII:cor:candidate-pic0-route}
\StatementRecord{wbII:cor:candidate-pic0-route}{proved}{03DS}{}

For the standard threefold of \SeriesRef{fam:src:A.3.1} and every
$L\in\operatorname{Pic}(X)$, put $A_L=(L^\vee\otimes\omega_X)|_W$.
Then
\begin{equation}
 \begin{gathered}
 H^0(W,\Omega_X^1|_W\otimes A_L)\ne0,
 \qquad H^2(W,(T_X\otimes L)|_W)\ne0,\\
 R^2f_*(T_X\otimes L)\ne0.
 \end{gathered}
\hypertarget{tag-03GP}{}\label{wbII:eq:W:5.17f}
\end{equation}
This includes every $L\in\operatorname{Pic}^0(X)$ and is the
nonvanishing assertion of \citet[Theorem~10.5(a), pp.~88--89]{AC26}.
\end{corollary}

\begin{proof}
The cusp charts in \SeriesRef{push:src:E.3.1}, at $k=l=1$, are
$q=z_0z_1z_2$, with the two-factor and smooth charts on the other strata.
Thus $W=f^{-1}(p_0)$ is the full reduced Cartier fibre and its normalization
is $\eta:\mathrm{dP}_6\to W$.  The isomorphism
$\omega_X\simeq\mathcal O_X(-2S_2)$ of
\SeriesRef{fam:src:A.1.1} trivializes $\omega_X$ on
$X\setminus S_2$, which contains $W$ since $p_0\ne p_2$.
Apply \cref{wbII:thm:analytic-all-picard} and
\eqref{wbII:eq:W:5.17b} to this same $L$, $i$ and $\eta$; they give
$\eta^*A_L\simeq\mathcal O_{\mathrm{dP}_6}$.
The conductor proposition \cref{wbII:prop:conductor} therefore gives
the stated $H^0$ nonvanishing.

Use the last two proof steps of \cref{wbII:thm:candidate-conductor}
with $\mathcal F_L=T_X\otimes L$.  Their maps require precisely the
full Cartier fibre just identified and the nonzero section space just
obtained: adjunction identifies $\omega_W$ with $\omega_X|_W$, and
Cartier-fibre duality identifies $H^2(W,\mathcal F_L|_W)^\vee$ with
$H^0(W,\Omega_X^1|_W\otimes A_L)$.  The same duality gives
$H^3(W,\mathcal F_L|_W)=0$.  Compactness and the nonconstant map to the
smooth curve give properness and $f$-flatness by
\cref{wbII:lem:curve-flatness}.  Hence
\cref{wbII:prop:analytic-base-change} identifies
$(R^2f_*\mathcal F_L)_{p_0}\otimes\mathbb C(p_0)$ with the nonzero
$H^2(W,\mathcal F_L|_W)$.  This proves both remaining assertions.
\end{proof}

\begin{lemma}[Generic endpoints in an analytic parameter space]
\hypertarget{tag-03IA}{}\label{wbII:lem:generic-endpoints}
\StatementRecord{wbII:lem:generic-endpoints}{proved}{03IA}{\cite[Chapter~III, \S4, Theorem~4.1, p.~124]{BS76push}}

Let $Y$ be a connected compact complex manifold of dimension $n$, let $D$ be
a nonzero effective Cartier divisor with $\cO_Y(D)\in\Pic^0(Y)$, and let
$B$ be a holomorphic vector bundle.  Let $\mathcal P\subset\Pic^0(Y)$ be a
connected irreducible reduced complex analytic parameter space carrying a
Poincar\'e line bundle on $Y\times\mathcal P$.  Assume that dualization
preserves $\mathcal P$, and suppose it
  contains points $P_0=\cO_Y(-mD)$ and $Q_0=\cO_Y(-\ell D)$ for some
  $m,\ell\ge1$ with
  $H^0(Y,B\otimes P_0)=0$ and
  $H^0(Y,B^\vee\otimes\omega_Y\otimes Q_0)=0$.  Then there is a dense open subset
$U\subset\mathcal P$ such that, for every $P'\in U$,
\begin{equation}
 H^0(Y,B\otimes P')=H^n(Y,B\otimes P')=0.
 \hypertarget{tag-03E5}{}\label{wbII:eq:CN:5.2}
\end{equation}
\end{lemma}



\begin{proof}
  The chosen point $P_0$ satisfies
  $H^0(Y,B\otimes P_0)=0$.  The Poincar\'e bundle gives a proper flat family over
$\mathcal P$.  Locally, its cohomology is computed by a finite free complex
\cite[Chapter~III, \S4, Theorem~4.1, p.~124]{BS76push}.
The rank conditions on its differentials make
$\{P':h^0(Y,B\otimes P')\ge1\}$ a determinantal locus; these local loci glue
to a closed analytic subset of $\mathcal P$.  It is proper because it misses $P_0$, and irreducibility of
$\mathcal P$ makes its complement dense open.  The other chosen point satisfies
$H^0(Y,B^\vee\otimes\omega_Y\otimes Q_0)=0$; dualization stability puts
$Q_0^\vee$ in $\mathcal P$.  Serre duality then gives
$H^n(Y,B\otimes Q_0^\vee)=0$.  The corresponding top-degree jump locus is
therefore also proper with dense open complement.  Intersecting the two
dense open complements proves the claim.

\end{proof}

\begin{lemma}[Generic endpoints on \texorpdfstring{$\operatorname{Pic}^0$}{Pic0}]\hypertarget{tag-03I7}{}\label{wbII:lem:cdp-generic-endpoints}
\StatementRecord{wbII:lem:cdp-generic-endpoints}{proved}{03I7}{\SeriesRef{cdp:st:E17}; \cite[Chapter~III, \S4, Theorem~4.1, p.~124]{BS76push}}

Let \(Y\) be a connected compact complex manifold, let \(D\) be
a nonzero effective Cartier divisor, and suppose
\(\mathcal O_Y(D)\in\operatorname{Pic}^0(Y)\).  For every holomorphic vector
bundle \(B\), a general \(P\in\operatorname{Pic}^0(Y)\) satisfies
\begin{equation}
 H^0(Y,B\otimes P)=0,\qquad
 H^{\dim Y}(Y,B\otimes P)=0.
\hypertarget{tag-03HL}{}\label{wbII:eq:W:5.23}
\end{equation}
\end{lemma}



\begin{proof}
Apply \cref{cdp:st:E17} with $Z=Y$, $\mathcal V=B$ and the divisor $D$.
The smooth connected space $Y$ has one irreducible component and no
embedded associated components, and $D$ meets that component. Thus
\[
 V_k=H^0(Y,B\otimes\mathcal O_Y(-kD))
 \subseteq H^0(Y,B),
\]
with \(V_k=0\) for all sufficiently large \(k\).
Since
\(\mathcal O_Y(-kD)\in\operatorname{Pic}^0(Y)\), local finite free complexes for the proper family over a
local universal parameter space of topologically trivial line bundles
\cite[Chapter~III, \S4, Theorem~4.1, p.~124]{BS76push} show that
\[
 \{P:h^0(Y,B\otimes P)>0\}
\]
is a proper analytic subset of the connected complex Lie group
\(\operatorname{Pic}^0(Y)\), so its complement is dense open.
Serre duality identifies the dual
of the second group in \eqref{wbII:eq:W:5.23} with
\[
 H^0(Y,B^\vee\otimes\omega_Y\otimes P^\vee).
\]
Apply the first part to \(B^\vee\otimes\omega_Y\) and use that inversion is
an automorphism of \(\operatorname{Pic}^0(Y)\).

\end{proof}

\begin{theorem}[Global vanishing from one degree-two vanishing]
\hypertarget{tag-03J0}{}\label{wbII:thm:corrected-cdp-global}
\StatementRecord{wbII:thm:corrected-cdp-global}{proved}{03J0}{\SeriesRef{wbII:lem:cdp-generic-endpoints}, \SeriesRef{cdp:st:C.1.1}, \SeriesRef{wbII:eq:W:5.28}}
\label{wbII:thm:corrected-global}
Let \(X\) be a smooth compact connected complex threefold such that
\[
 H^1(X;\mathbb Z)=H^2(X;\mathbb Z)=0,\qquad a(X)=1,
\]
and suppose its algebraic reduction is a holomorphic map
\(f:X\to C\).  If there exists \(L_0\in\operatorname{Pic}(X)\) with
\begin{equation}
 H^2(X,T_X\otimes L_0)=0,
\hypertarget{tag-03HM}{}\label{wbII:eq:W:5.24}
\end{equation}
then, for \(M\) in a dense open subset of \(\operatorname{Pic}(X)\),
\begin{equation}
 H^i(X,T_X\otimes M)=0\quad(i\ne1),
\hypertarget{tag-03HN}{}\label{wbII:eq:W:5.25}
\end{equation}
\begin{equation}
 \chi(X,T_X\otimes M)\le0,\qquad
 \left\langle c_3(T_X),[X]\right\rangle\le0.
\hypertarget{tag-03HO}{}\label{wbII:eq:W:5.26}
\end{equation}
\end{theorem}



\begin{proof}

\begin{enumerate}[label=\arabic*.,leftmargin=*]
\item Tensoring the integral vanishings with $\mathbb Q$ gives
$b_1(X)=b_2(X)=0$. Apply \cref{cdp:st:C.1.1} with $Y=X$;
its Picard component map and $H^2(X;\mathbb Z)=0$ give
\begin{equation}
 \operatorname{Pic}(X)=\operatorname{Pic}^0(X)
 \simeq H^1(X,\mathcal O_X).
\hypertarget{tag-03HP}{}\label{wbII:eq:W:5.27}
\end{equation}
Thus this Picard group is a connected irreducible complex vector group.
Choose additive \v{C}ech cocycles $a^{(r)}_{ij}$ representing a basis of
$H^1(X,\mathcal O_X)$. For coordinates $\lambda_r$ on that vector space,
\[
 g_{ij}(\lambda)=\exp\!\left(\sum_r\lambda_r a^{(r)}_{ij}\right)
\]
satisfy the multiplicative cocycle condition and define a holomorphic
line bundle on $X\times H^1(X,\mathcal O_X)$. Its restriction at every
parameter represents the corresponding line bundle, and the projection to the parameter space is proper.  Any fibre divisor \(D\) of
\(f\) is nonzero and effective, and its line bundle lies in
\(\operatorname{Pic}^0(X)\) by \(H^2(X;\mathbb Z)=0\).  Applying
\cref{wbII:lem:cdp-generic-endpoints} to \(B=T_X\) shows that the degree-zero and
degree-three groups in \eqref{wbII:eq:W:5.25} vanish away from two proper analytic subsets
of \(\operatorname{Pic}(X)\).

\item The cohomology-jump locus
\[
 \{M\in\operatorname{Pic}(X):h^2(X,T_X\otimes M)\ge1\}
\]
is locally determinantal in this family
\cite[Chapter~III, \S4, Theorem~4.1, p.~124]{BS76push}.  It is proper because \eqref{wbII:eq:W:5.24} gives a
point outside it.  Its complement meets the two endpoint
vanishing loci in a dense open set, where \eqref{wbII:eq:W:5.25} holds and
\begin{equation}
 \chi(X,T_X\otimes M)=-h^1(X,T_X\otimes M)\le0.
\hypertarget{tag-03HQ}{}\label{wbII:eq:W:5.28}
\end{equation}

\item The same proposition, with $L=M$ and the complex orientation
of $X$, gives
\[
 \chi(X,T_X\otimes M)
 =\frac12\left\langle c_3(T_X),[X]\right\rangle.
\]
Together with \eqref{wbII:eq:W:5.28}, this proves
\eqref{wbII:eq:W:5.26}.

\end{enumerate}
\end{proof}

\begin{corollary}[Relative vanishing criterion]\hypertarget{tag-03DT}{}\label{wbII:cor:corrected-cdp-relative}
\StatementRecord{wbII:cor:corrected-cdp-relative}{proved}{03DT}{\SeriesRef{wbII:thm:corrected-cdp-global}, \SeriesRef{wbII:eq:W:5.24}; \cite[Proposition~2.4, p.~680]{CDP2020}}
\label{wbII:cor:relative}
In the setting of
\cref{wbII:thm:corrected-cdp-global}, it is enough that one line bundle
\(L\in\operatorname{Pic}(X)\) satisfy
\begin{equation}
 R^2f_*(T_X\otimes L)=0.
\hypertarget{tag-03HR}{}\label{wbII:eq:W:5.29}
\end{equation}
The relative criterion is \cite[Proposition~2.4, p.~680]{CDP2020}.
\end{corollary}



\begin{proof}
Let \(A\) be a positive-degree line bundle on the compact Riemann surface
\(C\), and put \(E=T_X\otimes L\).  Projection formula gives
\[
 R^2f_*(E\otimes f^*A^m)=R^2f_*E\otimes A^m=0.
\]
For \(m\gg0\), Serre vanishing on \(C\) gives
\[
 H^1(C,R^1f_*E\otimes A^m)=0.
\]
The Leray spectral sequence has only columns zero and one, so its
total-degree-two terms yield
\[
 H^2(X,E\otimes f^*A^m)=0.
\]
Thus \(L_0=L\otimes f^*A^m\) satisfies \eqref{wbII:eq:W:5.24}, and
\cref{wbII:thm:corrected-cdp-global} applies.

\end{proof}

\begin{theorem}[Reduced-fibre vanishing criterion]
\hypertarget{tag-03IZ}{}\label{wbII:thm:corrected-cdp-fibrewise}
\StatementRecord{wbII:thm:corrected-cdp-fibrewise}{proved}{03IZ}{\SeriesRef{wbII:thm:corrected-cdp-global}, \SeriesRef{wbII:eq:W:5.29}, \SeriesRef{wbII:eq:W:5.25}, \SeriesRef{wbII:eq:W:5.26}, \SeriesRef{cdp:st:A.1.1}, \SeriesRef{wbII:cor:surface-top-base-change}, \SeriesRef{wbII:cor:corrected-cdp-relative}}

In the setting of
\cref{wbII:thm:corrected-cdp-global}, assume every full scheme-theoretic fibre
\(W_c=f^{-1}(c)\) is reduced, and fix \(L\in\operatorname{Pic}(X)\).  Put
\[
 A_c=(\omega_X\otimes L^\vee)|_{W_c},\qquad
 \mathscr T_c=\ker\!\left(
 \Omega^1_{W_c}\longrightarrow
 j_*\Omega^1_{(W_c)_{\rm reg}}\right).
\]
For each \(c\), let
\begin{equation}
 \delta_c:
 H^0(W_c,\mathscr T_c\otimes A_c)
 \longrightarrow
 H^1(W_c,N^*_{W_c/X}\otimes A_c)
\hypertarget{tag-03HT}{}\label{wbII:eq:W:5.30}
\end{equation}
be the connecting map of the cotangent filtration in \cref{cdp:st:A.1.1}.  If, for every
\(c\),
\begin{equation}
 H^0(W_c,\widetilde\Omega^1_{W_c}\otimes A_c)=0,\qquad
 H^0(W_c,N^*_{W_c/X}\otimes A_c)=0,\qquad
 \delta_c\ \text{is injective},
\hypertarget{tag-03HU}{}\label{wbII:eq:W:5.31}
\end{equation}
then \eqref{wbII:eq:W:5.29}, \eqref{wbII:eq:W:5.25}, and \eqref{wbII:eq:W:5.26} hold.
\end{theorem}



\begin{proof}
Apply \cref{cdp:st:A.1.1} to the whole reduced
Cartier fibre \(W_c\).  Conditions \eqref{wbII:eq:W:5.31}
give
\begin{equation}
 H^0(W_c,\Omega_X^1|_{W_c}\otimes A_c)
 =\ker\delta_c=0.
\hypertarget{tag-03HV}{}\label{wbII:eq:W:5.32}
\end{equation}
Apply \cref{wbII:lem:curve-flatness} to $f$: the total space is
connected and smooth, and the algebraic reduction is nonconstant.
Thus $f$ and $E=T_X\otimes L$ are flat over the curve.
The fibre is compact and Cartier, so
\cref{wbII:prop:cartier-duality-absolute}, with this $E$ and a
nonzero vector in $T_cC$, gives
\[
 H^2(W_c,E|_{W_c})^\vee
 \simeq H^0(W_c,\Omega_X^1|_{W_c}\otimes A_c)=0.
\]
The uniform top-degree base-change isomorphism of
\cref{wbII:cor:surface-top-base-change} therefore gives
\[
 (R^2f_*E)_c\otimes_{\mathcal O_{C,c}}\mathbb C(c)=0
\]
for every \(c\).  Nakayama's lemma kills every stalk of the coherent sheaf
\(R^2f_*E\), proving \eqref{wbII:eq:W:5.29}.  The remaining conclusions follow from
\cref{wbII:cor:corrected-cdp-relative}.

\end{proof}

\begin{remark}[The connecting map on a reduced full fibre]
\hypertarget{tag-P311}{}\label{rem:II:reduced-fibre-connecting-scope}
\StatementRecord{rem:II:reduced-fibre-connecting-scope}{proved}{P311}{\SeriesRef{cdp:st:A.1.1}, \SeriesRef{wbII:thm:corrected-cdp-fibrewise}}
Under the two endpoint vanishings in \eqref{wbII:eq:W:5.31}, the
section-kernel isomorphism of \cref{cdp:st:A.1.1} identifies the
ambient group in \eqref{wbII:eq:W:5.32} with $\ker\delta_c$.
Consequently ambient vanishing is equivalent to injectivity of
$\delta_c$, computed on the whole reduced Cartier fibre.
For a nonreduced full fibre $W$, let
$\mathcal J=\ker(\mathcal O_W\to\mathcal O_{W_{\rm red}})$ and
$E_W=(\Omega_X^1\otimes\omega_X\otimes L^\vee)|_W$.
The filtration $E_W\mathcal J^j$ has graded terms
$E_W\otimes_{\mathcal O_W}(\mathcal J^j/\mathcal J^{j+1})$;
its connecting maps enter the computation of $H^0(W,E_W)$.
For $W=mS$, these are the conormal powers and connecting maps of
\cref{cdp:st:A.3.1}, and \cref{cdp:st:A.3.3} gives the finite
Picard-translation vanishing criterion.
The global and relative criteria
\cref{wbII:thm:corrected-cdp-global,wbII:cor:corrected-cdp-relative}
apply to full fibres with their nilpotents and multiplicities.
\end{remark}

\begin{remark}[The invariant-lattice spanning condition]
\hypertarget{tag-P312}{}\label{rem:II:invariant-lattice-spanning}
\StatementRecord{rem:II:invariant-lattice-spanning}{proved}{P312}{\SeriesRef{cdp:st:D.1.10}, \SeriesRef{cdp:st:D.1.7}, \SeriesRef{cdp:st:C.3.11}}
For the algebraic reduction in \cref{cdp:st:D.1.10}, put
$C^\circ=C\setminus\Delta$ and let $F$ be a general smooth fibre.
The coefficient map $\mathbb Z_F\hookrightarrow\mathcal O_F$
gives the map in that proposition. The complex span of its invariant
integral domain equals the whole target precisely when
\begin{equation}
 \operatorname{Span}_{\mathbb C}\operatorname{im}\!\left(
 H^1(F;\mathbb Z)^{\pi_1(C^\circ)}
 \longrightarrow H^1(F,\mathcal O_F)\right)
 =H^1(F,\mathcal O_F).
\hypertarget{tag-03HW}{}\label{wbII:eq:W:5.33}
\end{equation}
This is the exact final spanning condition in
\cite[Theorem~3.1, proof, final paragraph, p.~86]{CDP98}.
Trivial monodromy on $H_1(F;\mathbb Z)$ suffices: by
\cref{cdp:st:C.3.11}, the smooth fibre is a two-torus or belongs to
the branch $b_1(F)=h^1(F,\mathcal O_F)=1$.
For a two-torus the complexified coefficient map is the surjective
projection $H^1(F;\mathbb C)\to H^{0,1}(F)$, and the integral
lattice spans its source. In the second branch,
\cref{cdp:st:D.1.7} sends a nonzero integral generator to a nonzero
vector of the one-dimensional complex target, which therefore spans it.
For the standard marked family, write
$v=a\gamma+bu+cw+d\delta$ in the basis of
\cref{def:II:analytic-picard-data}.
The equation $T_2v=v$ gives $c=b$ and $d=2b$, and the first row of
$T_1v=v$ gives $2d=6c$. Hence $b=c=d=0$, so the common invariant
lattice is $\mathbb Z\gamma$.
Here $h^1(F,\mathcal O_F)=2$, whereas the invariant image has
complex span of dimension one by \cref{cdp:st:D.1.7}.
This family has the integral cohomology and holomorphic algebraic
reduction assumed in \cite[Theorem~2.2, p.~680]{CDP2020}, and its
displayed monodromy matrices are nonidentity. Thus those hypotheses
do not imply either \eqref{wbII:eq:W:5.33} or trivial monodromy.
\end{remark}

\begin{corollary}[Positive Euler number and cohomology nonvanishing]
\hypertarget{tag-03IV}{}\label{wbII:test:cdp-comparison}
\StatementRecord{wbII:test:cdp-comparison}{proved}{03IV}{\SeriesRef{wbII:thm:corrected-cdp-global}, \SeriesRef{wbII:cor:corrected-cdp-relative}, \SeriesRef{wbII:cor:candidate-conductor-actual}; \citet[Section~10.3, Theorem~10.5 and Corollary~10.6, pp.~88--89]{AC26}}

For the standard compact threefold $X$, every $L\in\operatorname{Pic}(X)$
satisfies
\[
 H^2(X,T_X\otimes L)\ne0,
 \qquad R^2f_*(T_X\otimes L)\ne0.
\]
The second conclusion also follows directly from the cusp conductor.
\end{corollary}


\begin{proof}
The standard compact smooth threefold constructed in Paper I has
$H^1(X;\mathbb Z)=H^2(X;\mathbb Z)=0$, holomorphic algebraic reduction
to $\mathbb P^1$, and $\langle c_3(T_X),[X]\rangle=2$.
If the displayed degree-two cohomology vanished for one line bundle,
\cref{wbII:thm:corrected-cdp-global} would give $c_3(X)\le0$, a contradiction.
If the second direct image vanished for one twist,
\cref{wbII:cor:corrected-cdp-relative} would give such a global twist after
tensoring by the pullback of a sufficiently positive base line bundle,
again a contradiction. Independently,
\cref{wbII:cor:candidate-conductor-actual} constructs a nonzero cusp-fibre
cohomology class and identifies it with the fibre of the second direct
image by top-degree base change. The resulting counterexample to the CDP conclusion is due to \citet[Section~10.3, Theorem~10.5 and Corollary~10.6, pp.~88--89]{AC26}.

\end{proof}

\section{The Liu--Rollenske parameter covering}\hypertarget{tag-P316}{}\label{sec:II:lr-parameter-covering}

\begin{theorem}[The parameter double cover]\hypertarget{tag-01RO}{}\label{ii:LR:covering}
\StatementRecord{ii:LR:covering}{proved}{01RO}{}
Fix the scalar $u$ in the Liu--Rollenske action $B_u=uB_0$,
the good fourth root at the $III^*$ source fibre, the logarithmic
vector $a=e_4/4$, its line character, and the periodic cusp fan of
the construction, with its source and line attachment maps. The
resulting construction is a holomorphic family on the universal logarithmic cover of a
sufficiently small punctured $u$-disc. The threefold reached after
one positive $u$-turn is not biholomorphic to the initial threefold.
The family descends to the degree-two parameter cover
\begin{equation}
                         u=w^2.
\label{eq:II:analytic-LR-1}\end{equation}
This degree is minimal. The conclusion holds for every constant
root-line character.
\end{theorem}
\begin{proof}
Write $z=\log u/(2\pi i)$ on the logarithmic parameter cover, and put
$r=e^{\pi iz/2}$ and $w=e^{\pi iz}$. Thus $u=r^4=w^2$.
The source, line bundle and cusp fan remain fixed throughout.

\begin{enumerate}[label=\arabic*.,leftmargin=*]
\item \textbf{The central torus and the logarithmic parameter.}
The source is
\[
 y^2=x^3+\frac{27}{4}t^3(1-t)x+\frac{27}{4}t^5(1-t),
 \qquad P=(-t^2,it^3),\qquad \mathcal M=\mathcal O(P-O).
\]
The lift $B_0$ covers translation by $-2P$ and has the simple
whole-fibre zero at $t=\infty$; see
\cite[Proposition~2.1 and equation~(2.3), pp.~3--4]{LiuRollenske2026}.
On the good root, $P$ specializes to $(1+i)/2$ on
$E=\mathbb C/(\mathbb Z+i\mathbb Z)$. The limiting line has character
\begin{equation}
                  m+in\longmapsto(-1)^{m+n}.          
\label{eq:II:analytic-LR-4}\end{equation}
These identifications, including their line maps, are given in
\cite[Lemma~2.7 and Corollary~2.8, pp.~6--7]{LiuRollenske2026}.
Since $2P$ specializes to zero, $B_0$ acts on the central line
complement by a constant $c\ne0$. This constant is fixed independently
of $u$, so its scalar quotient period is
\begin{equation}
              \alpha=cu,\qquad
              \tau=\frac{\log(cu)}{2\pi i},
              \qquad\operatorname{Im}\tau>0.         
\label{eq:II:analytic-LR-5}\end{equation}
Writing the line coordinate as $v=e^{2\pi i\eta}$ gives
\begin{equation}
 A_\tau=\mathbb C^2/\Lambda_\tau,
 \quad
 \Lambda_\tau=
 \left\langle
  (2,0),(1+i,0),(1,1/2),(0,\tau)
 \right\rangle_{\mathbb Z}.                           
\label{eq:II:analytic-LR-6}\end{equation}
Indeed the first two generators are the periods in the kernel of
\eqref{eq:II:analytic-LR-4}, the third records its other coset, and the fourth is the
quotient by $B_u$. Conversely, exponentiating the second coordinate
recovers this line quotient. In particular
$\delta=(0,1)=2(1,1/2)-(2,0)$ and $e_4=(0,\tau)$.

The source root deck acts by $c_0\mapsto-ic_0$ on $E$. In these flat
coordinates its lift is $(c_0,\eta)\mapsto(-ic_0,\eta+\kappa)$,
where $4\kappa\in\mathbb Z$. The lift in
\cite[proof of Proposition~3.9, p.~10]{LiuRollenske2026} has
$\kappa=1/4$; a constant root-line character adds an integer multiple
of $1/4$. The difference of two lifts is a holomorphic nonzero
function on the compact elliptic curve, hence a scalar, and the
fourth-power condition gives precisely these four choices.
The chosen logarithmic vector is $(0,\tau/4)$ by
\cite[Definition~3.7, p.~10]{LiuRollenske2026}. The reduced multiple
fibre is therefore
\begin{equation}
 B_{\tau,\kappa}=A_\tau/\langle R\rangle,
 \qquad R(c_0,\eta)=(-ic_0,\eta+\tau/4+\kappa).           
\label{eq:II:analytic-LR-7}\end{equation}
Its fourth power is translation by the lattice vector
$(0,\tau+4\kappa)$. For $j=1,2,3$, the imaginary part of the second
coordinate of $R^j$ differs by $j\operatorname{Im}\tau/4$, whereas
the corresponding lattice periods are integral multiples of
$\operatorname{Im}\tau$. Thus the order-four action is free.

In the ordered basis of \eqref{eq:II:analytic-LR-6}, the linear root deck is
\[
 D=\begin{pmatrix}
 1&1&0&0\\-2&-1&-1&0\\0&0&1&0\\0&0&0&1
 \end{pmatrix}.
\]
This is the inverse of the counterclockwise monodromy matrix in
\cite[Proposition~3.6, pp.~9--10]{LiuRollenske2026}: multiplying the
two displayed matrices gives the identity. The equations
$(D-I)b=0$ are $b_2=0$ and $2b_1+b_3=0$, so its integral invariant
lattice is exactly $\mathbb Z\delta\oplus\mathbb Z e_4$.
A positive $u$-turn gives
\begin{equation}
 \tau\longmapsto\tau+1,
 \quad \Lambda_{\tau+1}=\Lambda_\tau,
 \quad a\longmapsto a+\delta/4.                      
\label{eq:II:analytic-LR-8}\end{equation}
The lattice equality has both inclusions because $e_4$ is replaced
by $e_4+\delta$ and $\delta$ already belongs to the lattice.

The same construction gives a holomorphic family on the logarithmic
parameter cover. The action $uB_0$ depends holomorphically on $u$;
its contraction is uniform on compact parameter subsets after the
punctured disc is sufficiently small. The good-root period $e_4$,
the torsion section $e_4/4$, and its free cap quotient then vary
holomorphically. The cusp periods have holomorphic unit coefficients
and use the fixed fan, so their monomial charts and period
identifications also form a holomorphic family. The punctured cap
attachment is the translation constructed in
\cite[Proposition~2.12, pp.~7--8]{LiuRollenske2026}. The proper local
pieces over the compact source base give a proper family.

\item \textbf{An intrinsic invariant of the reduced fibre.}
Projection to the second coordinate gives
\begin{equation}
 \pi_{\tau,\kappa}:B_{\tau,\kappa}\longrightarrow
 E_{\tau,\kappa}:=
 \frac{\mathbb C}{\frac12\mathbb Z+
                       (\frac\tau4+\kappa)\mathbb Z},
 \qquad [(c_0,\eta)]\longmapsto[\eta].                  
\label{eq:II:analytic-LR-9}\end{equation}
The image of the full affine deck group under second-coordinate
translation is generated by $1/2$, $\tau$ and $\tau/4+\kappa$.
The equality
$\tau=4(\tau/4+\kappa)-2(4\kappa)(1/2)$ proves that its image is
exactly the lattice in \eqref{eq:II:analytic-LR-9}. Its zero-second-coordinate subgroup
has source lattice $2\mathbb Z+(1+i)\mathbb Z$: cancellation of the
$\tau$ coefficient forces the exponent of $R$ to be a multiple of
four, after which its fourth-power relation removes it. The fibre
of \eqref{eq:II:analytic-LR-9} is thus the connected elliptic curve with this source
lattice.

Holomorphic one-forms on $A_\tau$ are constant linear forms, since
their lifted coefficients descend to holomorphic functions on the
compact torus. The $R$-invariant subspace is exactly
$\mathbb C\,d\eta$. Its periods on the affine deck group are exactly
the lattice in \eqref{eq:II:analytic-LR-9}. Integration therefore identifies \eqref{eq:II:analytic-LR-9}
with the Albanese map, up to a translation of its target. In
particular the target elliptic curve is intrinsic to the reduced
fibre.

The map $[\eta]\mapsto[2\eta]$, with inverse $[v]\mapsto[v/2]$,
identifies this target with
$\mathbb C/(\mathbb Z+(\tau/2+2\kappa)\mathbb Z)$. Its exponential
parameter is
\begin{equation}
 Q=\exp\bigl(2\pi i(\tau/2+2\kappa)\bigr)
    =e^{4\pi i\kappa}\sqrt{cu}.                     
\label{eq:II:analytic-LR-10}\end{equation}
Put $\omega=\tau/2+2\kappa$, and shrink the punctured disc so that
$\operatorname{Im}\omega>1$. A positive $u$-turn changes $\omega$
to $\omega+1/2$. An isomorphism of these elliptic curves can first
be translated to preserve zero, and then lifted to a complex-linear
map of their universal covers. Comparing the oriented integral
period bases gives integers $A,B,C,D$ with $AD-BC=1$ and
\[
 \omega+\tfrac12=\frac{A\omega+B}{C\omega+D}.
\]
Equality of imaginary parts would give $|C\omega+D|=1$. If $C\ne0$,
this absolute value is at least $|C|\operatorname{Im}\omega>1$.
Thus $C=0$, $A=D=\pm1$, and the change of period is an integer,
which cannot be $1/2$. It follows that
\begin{equation}
            B_{\tau,\kappa}\not\simeq
            B_{\tau+1,\kappa}                       
\label{eq:II:analytic-LR-11}\end{equation}
for every sufficiently small nonzero $u$. The bound is uniform in
the four characters because $\kappa$ is real.

\item \textbf{The algebraic reduction preserves this fibre.}
We specify the source comparison needed for the topological input
\SeriesRef{night:P3F.47}. Its equation is
$Y^2=X^3+\eta X-\eta$, with marked section $G=(1,1)$. The maps
\[
 t=\frac{27}{4\eta+27},\quad x=-t^2X,\quad y=it^3Y;
 \qquad
 \eta=\frac{27(1-t)}{4t},\quad X=-x/t^2,\quad Y=-iy/t^3
\]
are inverse Weierstrass coordinate changes, as substitution shows.
They send $G$ to $P$ and the zero section to the zero section. At
$\eta=\infty$, set $\nu=1/\eta$; then $t/\nu=27/(4+27\nu)$ is a
unit, and the regular coordinates $(\nu^2X,\nu^3Y)$ are changed by
its second and third powers. At $t=\infty$, the regular coordinates
$(t^{-2}x,t^{-3}y)$ are $(-X,iY)$. These maps identify the global
Weierstrass models and lift through their minimal resolutions, with
the same singularities and resolution centres transported by the
local isomorphisms.

They also identify the divisor lines $\mathcal O(G-O)$ and
$\mathcal O(P-O)$. Because zero sections are preserved, the fibre
group laws and translations by $-2G$ and $-2P$ correspond. For the
simple whole-zero linearization in \SeriesRef{night:P3F.47}, use
the transported $B_0$ itself. Thus $uB_0$, the good root, its
logarithmic period, the constant character, and every chosen fan
and attachment map are carried to the stated input without changing
$u$. Here $(a,b)=(1,2)$, $m_*=4$, and the ordinary II cap has
$m_0=1$, $n_0=0$. For the meridian comparison, one positive turn
of the good-root disc projects to four positive base turns;
the attachment $x\mapsto x-(\log\zeta)\beta/(2\pi i)$ contributes
the negative primitive suspension period $-\beta$. Its meridian
is therefore $4q_{\rm base}-\beta$, so $|n_*|=1$ in either
orientation of the suspension generator. Explicitly, reversing
that generator changes the boundary basis by
$\operatorname{diag}(1,-1)$, with the same matrix as inverse.
The two signed tuples are in the domain of \SeriesRef{night:P3F.47},
and $m_0n_*+m_*n_0=n_*=\pm1$. Its sphere conclusion gives
$H^2(Y;\mathbb R)=0$ for every one of these fans and characters.

Let $f:Y\to\mathbb P^1$ be the resulting proper connected torus
fibration. Its smooth fibres are compact complex tori, with
translation-invariant K\"ahler forms. The preceding sphere
calculation gives $H^2(Y;\mathbb R)=0$, so \cref{cdp:st:E1}
applies with this same $Y$ and $f$. Its pullback isomorphism is
\begin{equation}
                   \mathbb C(Y)=f^*\mathbb C(t).      
\label{eq:II:analytic-LR-12}\end{equation}
A biholomorphism of two members induces an automorphism of this
field, hence a M\"obius map of the base: a rational function with a
rational compositional inverse has degree one. The maps to the
base commute with this M\"obius map. Fibre multiplicities are
preserved because its local derivative is nonzero. The unique
multiple fibre has multiplicity four; the ordinary II fibre and
the index-one cusp filling have multiplicity one. Its reduced
support must therefore map biholomorphically to the reduced
multiple fibre of the target. Equation \eqref{eq:II:analytic-LR-11} excludes a
biholomorphism after one $u$-turn.

\item \textbf{The double-turn map and descent.}
Put $F=x+t^2$ and $D_P=P+(-P)-2O$. On a generic source fibre,
$\operatorname{div}(F)=D_P$. In the infinity chart $q=1/t$,
$X=q^2x$ gives $F=q^{-2}(X+1)$, so on the smooth source surface
\[
 \operatorname{div}(F)=D_P+D_0-2F_\infty,
\]
where $D_0$ is supported over $t=0$. The canonical divisor identity
\[
 \operatorname{Hom}(t_P^*\mathcal M,\mathcal M)
   =\mathcal O(D_P)
\]
follows from $t_P^*(P-O)=O-(-P)$. If $s_{D_P}$ denotes its canonical
meromorphic section, $(t^2/F)s_{D_P}$ has divisor $2F_0-D_0$.
It therefore defines a line lift $C_0$ of translation by $-P$
which is holomorphic and invertible away from $F_0$, including
the entire $I_1$ fibre. Translation by the section extends on the
minimal regular elliptic surface, with translation by the opposite
section as inverse; in the nodal chart its action
on the normalization is multiplication by its nonzero smooth-locus
coordinate. The divisor calculation gives the same extension of
the line map and its inverse, including the node.

The quotient $C_0^2/B_0$ is a rational function with vertical divisor
supported at $0$ and $\infty$. Its restriction to each generic
compact elliptic fibre is holomorphic and nonzero, hence constant,
so it is a rational base function. Its order at infinity is $-1$:
$C_0$ is invertible there and $B_0$ has a simple zero. It follows
that $C_0^2/B_0=c_Ct$ for a constant $c_C\ne0$. Choose a square root
of $c_C^{-1}$ and multiply $C_0$ by this fixed scalar, obtaining
\[
                         C^2=tB_0.
\]
This scalar change of the lift $C_0$ gives $C$; $B_0$ and $u$
stay fixed. Set $U=C^{-1}$, a lift of translation by $+P$. Thus
\begin{equation}
              \rho(t)=B_0^{-1}/U^2=t.
\label{eq:II:analytic-LR-13}\end{equation}
At infinity $U$ has order $e_\infty=0$, so
\begin{equation}
                  \operatorname{ord}_{\infty}\rho=-1-2e_\infty=-1.
\label{eq:II:analytic-LR-14}\end{equation}
The square identity also proves that $C$ and $U$ commute with
$B_0$, and consequently with $B_u$.

Write $\zeta$ for the normal-form coordinate on the fixed good root
and put $\theta=\zeta^4$. The global coordinate satisfies $t=T(\theta)$ with
$T(0)=0$ and $T'(0)\ne0$. Choose a holomorphic unit $h$ satisfying
$h(\theta)^4=T(\theta)/\theta$, and put $s=\zeta h(\theta)$.
Then $s^4=t$. The inverse has the form
$\zeta=s\widetilde h(s^4)$, where
$\widetilde h(t)^4=T^{-1}(t)/t$ and
$\widetilde h(0)=h(0)^{-1}$. Both maps commute with multiplication
of the root coordinate by $i$.

The meromorphic pullback of $C$ to the good model has only a central
divisorial order. Its square is $s^4B_0$, and the extending $B_0$
is invertible by \cite[Lemma~2.7(ii)--(iii), p.~6]{LiuRollenske2026}. The central
fibre is smooth and irreducible, so this order is two. Hence
$H=s^{-2}C$ and its inverse extend holomorphically across it;
after the divisorial factor is removed, Hartogs extension in the
smooth line Hom bundle removes any isolated indeterminacy.
Set $L=H^{-1}=s^2U$. If $D$ is the root deck, $U$ commutes with
$D$ because it is defined below the root, while $s^2$ changes sign.
Thus
\begin{equation}
                 DLD^{-1}=-L,
\label{eq:II:analytic-LR-15}\end{equation}
and
\begin{equation}
                              L^2=B_0^{-1}.
\label{eq:II:analytic-LR-16}\end{equation}
A constant root-line character commutes with these scalar factors,
so the identities hold for each one. Explicitly,
$U=\zeta^{-2}h(\theta)^{-2}L$. Its root order is $e_0=-2$, and
\begin{equation}
                    \operatorname{ord}_{t=0}\rho=-e_0/2
                         =1\equiv1\pmod2.
\label{eq:II:analytic-LR-17}\end{equation}

Use the four-turn closing whose exterior multiplier is
$h_0(t)=t^{-1}$. In the fixed $\zeta$ cap its multiplier is the
unit $\theta/t$. These maps agree on their overlap: the cap
attachment after four turns differs by $\theta^{-1}$. Here the
rational factor in the square-root equation is explicit:
\begin{equation}
                             g_0=1,\qquad g_0^2=h_0\rho.
\label{eq:II:analytic-LR-18}\end{equation}
The exterior map over $r\mapsto-r$ is
\begin{equation}
                    \widetilde D_r=-r^{-2}g_0(t)U.
\label{eq:II:analytic-LR-19}\end{equation}
On the fixed good-root cap use $-r^{-2}\zeta^2U$
$=-r^{-2}h(\theta)^{-2}L$. Its inverse is
$-r^2\zeta^{-2}C=-r^2h(\theta)^2H$, so both maps extend across
the central fibre.

To verify the attachment and conjugacy, let $\beta_z$ be the
period for $B_u$ and put $a_z=\beta_z/4$. The cap attachment is
translation by $-\sigma_z$, where
$\sigma_z=(\log\zeta)\beta_z/(2\pi i)$.
Only the linear part $D_\rho$ of the affine deck acts on this
period vector: $\sigma_z(i\zeta)-D_\rho\sigma_z(\zeta)=a_z(i\zeta)$.
In the affine-map conjugacy, the constant line character occurs
on both sides and cancels. A two-turn change gives
$\beta_{z+2}=\beta_z+2\delta$, so the attachment change is scalar
$\zeta^{-2}$. Consequently
\[
 \Phi_{z+2}(-r^{-2}\zeta^2U)=(-r^{-2}U)\Phi_z.
\]
The line lifts $U,L$ induce translations on the smooth tori. To see
this directly, use a degree-zero flat presentation of the line on
an elliptic fibre. A linear lift of a translation has a multiplier
which is periodic on its universal cover; this descends to a
holomorphic nonzero function on the compact elliptic curve and
is constant. Its line logarithm is therefore a translation, with
the opposite translation as inverse. It commutes with translation
by $a_z$. If $A_r$ is the twisted cap deck, \eqref{eq:II:analytic-LR-15} therefore gives
\begin{equation}
               LA_rL^{-1}=-A_r=A_{-r}.
\label{eq:II:analytic-LR-20}\end{equation}
The scalar unit $h(\theta)^{-2}$ is deck invariant. This proves the
conjugacy for the cap map just displayed, with all line characters.
For comparison of the two root coordinates, translation by
$d_z=(\log h)\beta_z/(2\pi i)$ carries the $\zeta$-attachment to
the $s$-attachment; its inverse is translation by $-d_z$, and
$\Phi_z^s\,T_{d_z}=\Phi_z^\zeta$.

At the cusp write the two multiplicative periods as
$(\kappa_1(q)q,\alpha_1(q))$ and
$B_u=(g(q),u q b(q))$, with holomorphic unit coefficients.
Since $U$ extends invertibly over the source $I_1$ fibre, in flat
coordinates it has the form $(Z,V)\mapsto(a(q)Z,d(q)V)$ with
$a,d$ holomorphic units. The multiplier is source-independent by
the flat-line argument above; evaluating it at a smooth point of
the nodal fibre proves the extension as a unit in $q$. The identity
$U^2=qB_0^{-1}$ gives $a^2=g^{-1}$ and $d^2=b^{-1}$, including
the source coordinate. A toric character $Z^jV^kq^\ell$ is thus
multiplied by $a^j(-r^{-2}d)^k$, a holomorphic unit with a
holomorphic reciprocal. This extends \eqref{eq:II:analytic-LR-19} and its inverse
through every chart, cone and face of the chosen periodic fan.
The period maps and the source and line attachments commute with
these multiplications.

The cusp norm can also be read in additive coordinates. Choose
holomorphic logarithms $A=\log a/(2\pi i)$ and
$C_1=\log d/(2\pi i)$, and put
$v=\log r/(2\pi i)$ and $H=\log q/(2\pi i)$. The inverse $B_u$
period and the half-translation are
\[
 \gamma=(2A,2C_1-H-4v),\qquad
 \tau(v)=(H\delta+\gamma)/2=(A,C_1-2v).
\]
The logarithm of \eqref{eq:II:analytic-LR-19} is $\tau(v)+\delta/2$, with the last
term recording its minus sign. The two steps obey
\begin{equation}
 [\tau(v)+\delta/2]+[\tau(v+1/2)+\delta/2]
                       =H\delta+\gamma.
\label{eq:II:analytic-LR-21}\end{equation}
On the exterior this is the same exact norm identity
\begin{equation}
 \widetilde D_{-r}\widetilde D_r
    =r^{-4}g_0^2U^2
    =h_0 B_u^{-1}.
\label{eq:II:analytic-LR-22}\end{equation}
After quotienting by $B_u$, its multiplier is $h_0=t^{-1}$.
On the good-root cap the square is $\theta/t$. At the cusp,
multiplication by $q=t^{-1}$ in the line direction is
$B_u$ followed by multiplication by the units
$(g^{-1},u^{-1}b^{-1})$. Thus the four-turn closing extends through
the same fan, with its period relabelling and inverse.

Quotient the logarithmic family by this four-turn closing to
obtain the family over $u=r^4$. Equation \eqref{eq:II:analytic-LR-22} and its cap
identity make \eqref{eq:II:analytic-LR-19} an order-two action on that family, covering
$r\mapsto-r$. It is free because $r\ne0$. The quotient is a
proper holomorphic family over $w=r^2$, and hence over the
cover $u=w^2$. Properness is local on this parameter base and
follows from the proper local pullbacks; the cap and fan maps
above give its smooth total space. Steps 2--3 exclude a
biholomorphic endpoint after one $u$-turn, so this degree two
is minimal.
\end{enumerate}
\end{proof}

\begin{remark}[Parameter covering and fan returns]\hypertarget{tag-02UH}{}\label{night:P3K.36}\label{stmt:P3K.36}
\StatementRecord{night:P3K.36}{proved}{02UH}{\SeriesRef{ii:LR:covering}}
For the fixed Liu--Rollenske data of \cref{ii:LR:covering},
the minimal parameter cover is $u=w^2$.
The smooth return classes of the reference fan and the fan $T_2$
are computed in \SeriesRef{iv:correction:P3K.36}.
The marked source and cap calculations are
\SeriesRef{night:P3K.73} and \SeriesRef{night:P3K.74}.
\end{remark}

\clearpage
\part{Component geometry and Euler criteria}\label{part:II:component-geometry}
For a degeneration with arbitrary full central fibre, cyclic component
covers carry saturated translation fields, elliptic pencils and integral
monodromy maps. Their surface models and conductor curves determine the
local Euler contributions. The global sign statements use the setting
of \cref{def:II:standing-torus-fibration}; each local degeneration and
comparison construction has the hypotheses stated with its result.

\section{Euler topology and compact modifications}\hypertarget{tag-0342}{}\label{sec:II:06}

\subsection{Euler parity and projective blowup families}\label{subsec:II:unrestricted-euler-parity}

\begin{proposition}[Every torus-degeneration fibre has even Euler number]\hypertarget{tag-01CM}{}\label{cdp:st:C.2.UN.1}
\StatementRecord{cdp:st:C.2.UN.1}{proved}{01CM}{\SeriesRef{cdp:eq:C.UN.3}, \SeriesRef{cdp:eq:C.UN.4}}

Let $f:Y\to B$ be a proper holomorphic map from a smooth complex
threefold to a smooth complex curve, with connected fibres and smooth
general fibre a complex two-torus. For every $b\in B$,
\begin{equation}
                    e((Y_b)_{\rm red})\in2\Z.       \hypertarget{tag-0152}{}\label{cdp:eq:C.UN.1}
\end{equation}
The fibres may have arbitrary multiplicities and singularities.
\end{proposition}


\begin{proof}

\begin{enumerate}[label=\arabic*.,leftmargin=*]
\item Choose a coordinate disc about $b$ containing no other critical value, and apply \cref{lem:II:compact-analytic-tube} to its proper restriction. For a small closed subdisc put $U=f^{-1}(\overline\Delta)$ and $M=\partial U$. The inclusion of $S=(Y_b)_{\rm red}$ is the homotopy equivalence constructed there, so $e(S)=\chi(U)$. The boundary is the mapping torus of the orientation-preserving monodromy of a real four-torus.

\item Use homology with coefficients in $\Q$. Write
$s_i=\operatorname{rank}(H_i(U)\to H_i(U,M))$.
Exactness of the pair sequence, with Poincar\'e--Lefschetz duality,
gives
\begin{equation}
 b_0(M)+b_1(M)+b_2(M)
       =\sum_{i=0}^{5}b_i(U)-s_0-s_3-2s_1-2s_2.
                                                        \hypertarget{tag-0153}{}\label{cdp:eq:C.UN.2}
\end{equation}
Indeed, the summand $b_i(M)$ is the sum of the ranks of
$H_i(M)\to H_i(U)$ and $H_{i+1}(U,M)\to H_i(M)$;
substitution of $b_{i+1}(U,M)=b_{5-i}(U)$ gives the formula.
Here $s_0=0$ and $b_6(U)=0$. The map defining $s_3$ is represented,
after duality, by the middle-dimensional intersection pairing of $U$.
That pairing is skew-symmetric because the middle dimension is three,
so $s_3$ is even. Consequently
\begin{equation}
 \chi(U)\equiv b_0(M)+b_1(M)+b_2(M)\pmod2.            \hypertarget{tag-0154}{}\label{cdp:eq:C.UN.3}
\end{equation}

\item Let $T$ act on $H^1(F,\Q)$ for the four-torus fibre $F$, and set
$a_i=\dim\ker(\Lambda^iT-I)$. Cutting the mapping torus along one
fibre gives the Wang exact sequence; hence
\begin{equation}
 b_0(M)=1,\quad b_1(M)=1+a_1,\quad b_2(M)=a_1+a_2.
                                                        \hypertarget{tag-0155}{}\label{cdp:eq:C.UN.4}
\end{equation}
Orientation preservation says $\det T=1$. The transformation
$G=\Lambda^2T$ preserves the nondegenerate symmetric wedge pairing
on $\Lambda^2\Q^4$ and has determinant $(\det T)^3=1$.
The rank $\operatorname{rank}(G-I)$ is even. Complexify and split off the generalized
$(-1)$-eigenspace. Its dimension is even: all eigenvalues outside
$\{1,-1\}$ occur in reciprocal pairs with equal algebraic
multiplicity, and $\det G=1$. This generalized eigenspace is
nondegenerate and orthogonal to its invariant complement. On it
$G-I$ is invertible, of even rank. On the complement $G+I$ is
invertible, and
\[
                         (G-I)(G+I)^{-1}
\]
is skew-adjoint for the symmetric wedge pairing. Multiplication by
that nondegenerate pairing represents it as a skew-symmetric matrix,
whose rank is even. Its rank equals that of $G-I$ on the complement.
Since $\dim\Lambda^2\Q^4=6$, $a_2$ is
even. Equations \eqref{cdp:eq:C.UN.3}--\eqref{cdp:eq:C.UN.4} prove \eqref{cdp:eq:C.UN.1}.

\end{enumerate}
\end{proof}

\begin{proposition}[All even Euler numbers and projective blowup families]\hypertarget{tag-01CX}{}\label{cdp:st:C.2.UN.2}
\StatementRecord{cdp:st:C.2.UN.2}{proved}{01CX}{\cref{cdp:st:C.2.UN.1}}
\hypertarget{tag-01FH}{}\label{cdp:st:C2b}
\StatementRecord{cdp:st:C2b}{proved}{01FH}{}
The set of Euler numbers of reduced fibres among compact smooth threefolds fibred in complex two-tori is exactly $2\Z$. For every $k\ge1$ and $m\ge0$, there is a smooth projective threefold $Y_{k,m}\to\PP^1$ with connected fibres and smooth general fibre an abelian surface, whose only singular reduced fibre $S_{k,m}$ is normal crossing and satisfies \eqref{cdp:eq:C.UN.5}. It consists of an abelian surface blown up at $m$ distinct points, a chain of $k$ ruled surfaces over a genus-two curve attached to it, and $m$ projective-plane leaves attached to the abelian component. For $m=0$ the map on the abelian component is the identity. The central fibre is reduced. The total space $Y_{k,m}$ has torsion-free integral cohomology with Betti numbers \eqref{cdp:eq:C.UN.6}, and its structure-sheaf cohomology dimensions are $(1,2,1,0)$.

The choices $(k,m)=(1,0)$ and $(1,1)$ give Euler numbers $-2$ and zero, with respectively two and three components. The choices $m=k$ give the zero-Euler family with arbitrarily long chains of ruled components.
\end{proposition}
\begin{proof}
\begin{enumerate}[label=\arabic*.,leftmargin=*]
\item Let $C_2$ be the smooth projective completion of $w^2=z^6-1$. Its six simple branch points and the unramified points at infinity give genus two by Riemann--Hurwitz. Let $A=\operatorname{Jac}(C_2)$ and embed $C_2$ in $A$ by $p\mapsto\OO(p-p_*)$. This is injective because a principal divisor $p-q$ for distinct points would define a degree-one map $C_2\to\PP^1$. It is an immersion because the canonical system of a genus-two curve has no base point; the explicit differentials $dz/w,z\,dz/w$ verify that property, including at infinity. A proper injective immersion is an embedding. Adjunction on the torus gives $N_{C_2/A}=K_{C_2}$.

\item Blow up $A\times\PP^1$ along $C_2\times\{0\}$, and denote the resulting smooth projective threefold by $Y$. Compose its blowdown with projection to $\PP^1$. The exceptional divisor is $E=\PP(K_{C_2}\oplus\OO_{C_2})$. The strict transform of $A\times\{0\}$ is $A$ itself, since blowing up a Cartier divisor in a smooth surface is an isomorphism. These two components meet along a section isomorphic to $C_2$. In local coordinates $(u,v,t)$ with centre $(u,t)=(0,0)$, a blowup chart has $t=u\tau$, so the fibre is the reduced normal-crossing divisor $u\tau=0$. The other chart and the complement give the claimed global map and connectedness of its fibres. 
The smooth horizontal surface $R=C_2\times\PP^1$ has strict transform isomorphic to itself because the centre is a Cartier divisor in $R$. In the exceptional bundle $\PP(K_{C_2}\oplus\OO_{C_2})$, it meets the section corresponding to the base-normal direction. This section is disjoint from the section meeting the strict abelian component. In the chart $u=t\upsilon$ the map is $f=t$ and the strict transform of $R$ is $\upsilon=0$, so the map is a submersion there.

\item Blow up that section and repeat until there have been $k$ curve blowups, always using the intersection of
the strict transform of $R$ with the newest exceptional component.
The next centre is a smooth genus-two curve in the smooth locus of
the reduced central fibre, disjoint from its double curves. Its
normal sequence splits: the normal line to the curve inside $R$
maps isomorphically under $df$ to the tangent line of the base,
and supplies a splitting of the normal sequence of the central
surface. The two line summands give two disjoint sections on the
new exceptional projective bundle. One meets the previous
component; the other is the next centre. This proves by induction
that all $k$ curve blowups are defined and that the reduced central
support has the stated chain. In local coordinates the comparison
is always $t=u\tau$ as in the first blowup chart, so every new component has
multiplicity one and no triple point is introduced.

\item Finally blow up $m\ge0$ distinct points of the strict abelian
component away from the genus-two double curve. They introduce
$m$ projective-plane leaves, each meeting the blown-up abelian
component in one projective line. All maps are proper holomorphic
blowdowns followed by the original projection. Every nonzero
fibre is biholomorphic to $A$.


\item Apply \cref{lem:II:blowup-topology} to the specified curve and point blowdowns. At a genus-two curve the codimension is two, so \eqref{cdp:eq:C.VR.4} gives the increment $e(C_2)=-2$; at a point the codimension is three and the increment is two. Starting from $e(A)=0$ gives
\begin{equation}
                         e(S)=2(m-k).                 \hypertarget{tag-0156}{}\label{cdp:eq:C.UN.5}
\end{equation}
For any integer $n$, choose $k\ge\max(1,-n)$ and $m=k+n$. This realizes $2n$; \cref{cdp:st:C.2.UN.1} gives the reverse inclusion. For $(k,m)=(1,0)$ the fibre is $A\cup\PP(K_{C_2}\oplus\OO_{C_2})$ with intersection $C_2$. For $(1,1)$ it has components $\operatorname{Bl}_p A$, $\PP(K_{C_2}\oplus\OO_{C_2})$ and $\PP^2$, with pairwise intersections $C_2$ and $\PP^1$ and no triple intersection. Taking $m=k$ gives every stated zero-Euler family.

\item The integral Betti numbers of $A\times\PP^1$ are $(1,4,7,8,7,4,1)$. The map \eqref{cdp:eq:C.VR.9} of \cref{lem:II:blowup-topology}, for each curve blowdown of codimension two, adds $H^{j-2}(C_2,\Z)$ in degree $j$. For each point blowdown of codimension three it adds one generator in degrees two and four. These are free summands, so the $k$ curve and $m$ point blowups preserve torsion-freeness and give
\begin{equation}
 (b_0,\ldots,b_6)
      =(1,4,7+k+m,8+4k,7+k+m,4,1).                  \hypertarget{tag-0157}{}\label{cdp:eq:C.UN.6}
\end{equation}
In particular the two initial specializations have Betti numbers $(1,4,8,12,8,4,1)$ and $(1,4,9,12,9,4,1)$. The alternating sum is the value in \eqref{cdp:eq:C.UN.5}. Each blowup has structure-sheaf direct image $\OO$ and zero positive direct images, so Leray preserves the dimensions $(1,2,1,0)$ of $A\times\PP^1$.
\end{enumerate}
\end{proof}



\subsection{Vertical modifications and component classes}\label{subsec:II:vertical-modification-inputs}

\begin{definition}[Marked torus fillings and component classes]\hypertarget{tag-P235}{}\label{def:II:marked-torus-fillings}
\StatementRecord{def:II:marked-torus-fillings}{proved}{P235}{}
Let a proper holomorphic family $p^\times:\mathcal T^\times\to\Delta^\times$
of connected complex two-tori be fixed. A marked smooth filling is a proper
holomorphic map $f:Y\to\Delta$ from a smooth complex threefold with
connected fibres, together with a biholomorphism
$f^{-1}(\Delta^\times)\simeq\mathcal T^\times$ over $\Delta^\times$.
For each irreducible component $S_i$ of $(Y_0)_{\rm red}$, let
$Z_i\to S_i^\nu\to S_i$ be a smooth compact resolution of its
normalization. Write $\mathcal E(f)$ for the finite multiset of
bimeromorphism classes $[Z_i]$ for which $Z_i$ is not uniruled.
Here a smooth compact surface is uniruled when it admits a dominant
meromorphic map $C\times\PP^1\dashrightarrow Z$ for a smooth compact
curve $C$. Repetitions in the multiset record distinct components.


Over a compact base curve $B$, a marking away from a finite set $Q\subset B$ is a biholomorphism with a fixed torus family over $B\setminus Q$. The multiset at $b\in Q$ is defined using the components of $(Y_b)_{\rm red}$ in the same way.
\end{definition}
\begin{proof}[Well-definedness]
Two smooth resolutions of the normalization of the same component are bimeromorphic: the closure of the graph of their common dense normal open set, followed by normalization and resolution, has proper birational projections to both surfaces. Uniruledness is preserved by a bimeromorphism $Z'\dashrightarrow Z$. Compose a dominant meromorphic map $C\times\PP^1\dashrightarrow Z$ with the inverse on its dense isomorphism locus and resolve the graph of the composition. This gives a dominant meromorphic map to $Z'$. Applying the same construction to the inverse proves the converse. Hence both the class and its inclusion in $\mathcal E(f)$ are independent of the chosen resolution.
\end{proof}

\begin{lemma}[Non-uniruled divisors under bimeromorphisms]\hypertarget{tag-P236}{}\label{lem:II:nonuniruled-divisors}
\StatementRecord{lem:II:nonuniruled-divisors}{proved}{P236}{\cref{def:II:marked-torus-fillings}; \cite[\S1.2.4 and Lemma~1.3.1, arXiv v4 p.~6; \S5.4, p.~25]{AKMW2002}}
Let $\phi:M\dashrightarrow N$ be a bimeromorphic map between smooth complex threefolds. It induces a bijection between their compact irreducible divisors whose smooth normalization resolutions are not uniruled, by taking strict transforms; corresponding resolutions are bimeromorphic. For a blowup of a smooth compact point or curve, this correspondence is the strict-transform map, and every new exceptional divisor is uniruled. For holomorphic maps $g:M\to B$ and $h:N\to B$ to a smooth curve with $h\circ\phi=g$ on the domain of $\phi$, corresponding fibre components have the same full-fibre multiplicities. Every smooth compact uniruled surface $Z$ has $H^0(Z,K_Z^{\otimes m})=0$ for $m>0$.
\end{lemma}
\begin{proof}
Fix a compact divisor $P\subset M$ as in the statement. Eliminate the indeterminacies of the inverse map near its compact graph over $P$, using smooth-centre blowups:
\begin{equation}
 \widehat N\xrightarrow{\ a\ }N,
 \qquad b:\widehat N\longrightarrow M,
 \qquad b=\phi^{-1}a.
 \hypertarget{tag-014C}{}\label{cdp:eq:C.UC.1}
\end{equation}
The elimination-of-indeterminacy construction in \cite[\S1.2.4 and Lemma~1.3.1, arXiv v4 p.~6]{AKMW2002} gives this diagram; its analytic form is stated in \cite[\S5.4, arXiv v4 p.~25]{AKMW2002}. The graph projections are proper, so its inverse image over $P$ and its image in $N$ are compact. On neighbourhoods of these compact sets finitely many resolution steps suffice; the displayed maps denote the resulting proper modifications there, composed with the inclusions into $M$ and $N$. For proper disc models one can instead shrink the disc about the compact central fibre, with compact centres above that fibre.

The strict transform $\widehat P$ maps birationally to $P$ under $b$. Indeed, the locus near $P$ over which the proper modification $b$ has positive-dimensional fibres is analytic. It cannot contain a divisor: its preimage would have dimension at least three and be a proper analytic subset of the corresponding irreducible threefold. Thus $b$ is finite over a dense open subset of $P$, where a finite modification of the normal manifold $M$ is an isomorphism. If $a$ contracted $\widehat P$, that compact divisor would be the strict transform of an exceptional divisor created at one of the blowups. In dimension three such a divisor is $\PP^2$ for a point centre, or $\PP(N_{C/Y})$ for a smooth curve centre. Its centre is compact, being the image of the compact divisor. Blowing up a divisor is an isomorphism and creates none.

Both displayed exceptional surfaces are uniruled. On affine charts $(u,v)\mapsto[1:u:v]$ gives a bimeromorphism $\PP^1\times\PP^1\dashrightarrow\PP^2$. For a smooth compact curve $C$, tensoring the rank-two normal bundle by a sufficiently positive line bundle supplies two global sections independent at the generic point. Their projectivization gives a bimeromorphism $C\times\PP^1\dashrightarrow\PP(N_{C/Y})$. The well-definedness proof of \cref{def:II:marked-torus-fillings} preserves uniruledness under their further strict transforms. Thus $a$ cannot contract $\widehat P$. Its image is a compact divisor, and the same generic-isomorphism argument for $a$ makes the map birational there. Applying the construction to $\phi^{-1}$ proves the stated bijection.

For a single blowup the restrictions $S_i'\to S_i$ are proper and are isomorphisms away from the lower-dimensional centre. Their normalized resolution graphs give the asserted bimeromorphisms, while \eqref{cdp:eq:C.EC.2} is the exceptional projection just used. At each generic point the blowdown is an isomorphism. Pullback of the local base parameter therefore preserves its vanishing order on the strict transform. For a bimeromorphism over the base, use the two generic isomorphisms in \eqref{cdp:eq:C.UC.1}; equality of the base functions on their common open set gives the same equality of orders.

A dominant meromorphic map $C\times\PP^1\dashrightarrow Z$ pulls back a nonzero pluricanonical section to a nonzero one, because its generic differential has full rank. The pullback is holomorphic away from the finitely many indeterminacy points and extends over them by Hartogs' theorem. On every $\{c\}\times\PP^1$ its line bundle has degree $-2m$, so its restriction vanishes. This contradicts its being nonzero and proves the final assertion.
\end{proof}

\begin{proposition}[Every even Euler number occurs above each fixed torus filling]\hypertarget{tag-01DJ}{}\label{cdp:st:C.2.UN.4}
\StatementRecord{cdp:st:C.2.UN.4}{proved}{01DJ}{\cref{cdp:eq:C.VR.4}, \cref{cdp:eq:C.VR.6}, \cref{cdp:eq:C.VR.8}, \cref{cdp:eq:C.VR.9}, \cref{cdp:st:C.2.UN.1}}

Let $f:Y\to B$ be a proper holomorphic map from a smooth complex
threefold to a smooth complex curve, with connected fibres and smooth
general fibre a complex two-torus. Fix $b\in B$ and put
$S=(Y_b)_{\rm red}$. For every integer $k$ there is a finite composition
of blowups with smooth compact centres in successive fibres over $b$,
\begin{equation}
 \pi:\widetilde Y\longrightarrow Y,\qquad
 \widetilde f=f\pi,\qquad
 e((\widetilde Y_b)_{\rm red})=2k.                 \hypertarget{tag-015W}{}\label{cdp:eq:C.VR.1}
\end{equation}
The map $\pi$ is an isomorphism over $B\setminus\{b\}$.
In particular every such filling is dominated by a smooth filling
whose central reduced support has Euler number zero. More precisely,
every further smooth model obtained by vertical blowups admits a
zero-Euler refinement by additional vertical blowups. Thus zero-Euler
models are cofinal for this domination relation, in the explicit sense
that every object is dominated by one of them.
Above every given filling there are such zero-Euler refinements with
arbitrarily many irreducible components of their central support.

An elementary operation lowering the Euler number by two consists of
one point blowup and one smooth genus-three curve blowup. On integral
cohomology its additional summands are $\Z^2$ in degrees two and four
and $\Z^6$ in degree three. In particular, whenever the Betti numbers
are finite,
\begin{equation}
 (\Delta b_2,\Delta b_3,\Delta b_4)=(2,6,2),
 \qquad \Delta b_j=0\quad(j\notin\{2,3,4\}).       \hypertarget{tag-015X}{}\label{cdp:eq:C.VR.2}
\end{equation}
The positive full-fibre multiplicities are unchanged on
their strict transforms.
\end{proposition}


\begin{proof}

\begin{enumerate}[label=\arabic*.,leftmargin=*]
\item Apply \cref{lem:II:blowup-topology} to a smooth compact centre $Z\subset S$ of codimension $r$. For $\pi_Z:Y'\to Y$, the new reduced support is the full inverse image $S'=\pi_Z^{-1}(S)$, so the Euler change is \eqref{cdp:eq:C.VR.4}, including intersections with other components.

\item Choose a smooth point $p$ of one component of $S$, away from every
other component. Such points form a nonempty open subset of each
irreducible component. If the full fibre has multiplicity $m>0$ on
that component, local holomorphic coordinates may be chosen with
\begin{equation}
                    t\circ f=x^m.                \hypertarget{tag-0160}{}\label{cdp:eq:C.VR.5}
\end{equation}
Indeed the local equation initially has the form $x^m a$, where $a$
is a holomorphic unit; a local holomorphic $m$th root of $a$ replaces
$x$ by $x a^{1/m}$. Blow up $p$. The exceptional surface
$E\simeq\PP^2$ meets the strict component in the line
$L=\{X=0\}\subset\PP(T_pY)$, with homogeneous coordinates
$[X:Y:Z]$ induced by $(x,y,z)$. It meets no other central component.
In the chart $y=r$, $x=ru$, $z=rv$, the map is
\begin{equation}
                         t=(ru)^m.                \hypertarget{tag-0161}{}\label{cdp:eq:C.VR.6}
\end{equation}
Consequently both $E=(r=0)$ and the strict component $(u=0)$
have multiplicity $m$. Formula \eqref{cdp:eq:C.VR.4}, with $r=3$ and centre a
point, shows that this step increases the central Euler number by two.

\item Inside $E$ take the Fermat quartic
\begin{equation}
 C=\{X^4+Y^4+Z^4=0\}\subset\PP^2.                \hypertarget{tag-0162}{}\label{cdp:eq:C.VR.7}
\end{equation}
Its three partial derivatives cannot vanish simultaneously at a
projective point, so it is smooth. It is connected: if the quartic
were a disjoint union of distinct plane components, two such
components would meet by their positive intersection number; their
intersection would be a singular point. Adjunction gives
$K_C=\OO_C(1)$ and $2g(C)-2=4$, hence $g(C)=3$ and $e(C)=-4$.
On $L$ its equation is $Y^4+Z^4=0$. It has four distinct roots,
none at $Z=0$, and the two gradients of the quartic and $X$ are
independent there. Thus $C$ meets $L$ transversely in exactly four
points.

\item Now blow up the smooth curve $C$ in the threefold obtained after the
first blowup. Its exceptional divisor is the projective line bundle
$\PP(N_{C/Y_1})\to C$. Equation \eqref{cdp:eq:C.VR.4}, with $r=2$, gives Euler
change $e(C)=-4$. The two successive changes sum to $-2$.
The ambient blowup is smooth at the four points of $C\cap L$. Near any
one of them, preserve the local equations $r=0$ of $E$ and $u=0$ of
the strict old component. Transversality gives a third coordinate $v$
such that $C=(r=v=0)$. Absorbing a remaining unit in $r$ preserves
these equations and gives $t=(ru)^m$ as in \eqref{cdp:eq:C.VR.6}. The two charts of
the curve blowup are
\begin{equation}
 \begin{array}{c|c|c}
 \text{chart}&\text{blowdown substitution}&\text{base map}\\ \hline
 \text{$r$ chart}&v=rw&t=(ru)^m\\
 \text{$v$ chart}&r=vw&t=(uvw)^m .
 \end{array}                                      \hypertarget{tag-0163}{}\label{cdp:eq:C.VR.8}
\end{equation}
In the second chart the new exceptional divisor, the strict transform
of $E$, and the strict old component have equations $v=0$, $w=0$,
and $u=0$, respectively. They meet with normal crossings and all
have multiplicity $m$. Away from $C\cap L$ the comparison is the
same two-coordinate blowup of $t=r^m$. The strict transform of $E$
is isomorphic to $E$, since $C$ is Cartier in $E$; the strict old
component is its point blowup at the four transverse intersection
points. These identifications follow directly from the charts in
\eqref{cdp:eq:C.VR.8}. Outside the specified neighbourhoods the maps are
isomorphisms, so any other singularities of the original support
are unchanged.

\item For the same blowdown maps, \cref{lem:II:blowup-topology} gives the integral isomorphism \eqref{cdp:eq:C.VR.9}, with $j$ the exceptional inclusion, $p$ its projection and $\xi=c_1(\OO(1))$. It applies both to the total space and to a compact fibre tube unchanged on its boundary.

\item The point blowup adds $\Z$ in each of degrees two and four. The
genus-three curve blowup adds $H^{q-2}(C,\Z)$, whose ranks in these
three successive degrees are $1,6,1$. This proves \eqref{cdp:eq:C.VR.2} and shows
that pre-existing integral torsion is unchanged.

\item Both steps are proper holomorphic maps from smooth threefolds. A
blowup has connected fibres and pushes its structure sheaf to the
structure sheaf of the original smooth space. Stein factorization
therefore proves that their composition with $f$ again has connected
fibres. The centres lie over $b$, so the resulting map is the same
proper torus family over $B\setminus\{b\}$.

\item Proposition~\ref{cdp:st:C.2.UN.1} gives $e(S)=2n$ for an integer $n$. To obtain $2k$
with $k\ge n$, use $k-n$ point blowups. To obtain $2k$ with $k<n$,
use $n-k$ of the two-step operations just proved. At each stage a
smooth point of a single component exists, and the construction
can be performed there; only finitely many centres are needed.
Applying this same procedure to any given vertical blowup model
proves cofinality with the displayed domination maps.
After reaching Euler number zero, compose one further point blowup
with one of the two-step decreasing operations, at smooth
points. Their net Euler change is zero, while they add three
irreducible central components. Repetition gives arbitrarily many
components above the original filling. Each repetition adds
$(3,6,3)$ to $(b_2,b_3,b_4)$ when those numbers are finite; the
maps are the successive blowdowns.

\end{enumerate}
\end{proof}

\begin{proposition}[The same non-uniruled components occur with every even Euler number]\hypertarget{tag-01CR}{}\label{cdp:st:C.2.UN.14}
\StatementRecord{cdp:st:C.2.UN.14}{proved}{01CR}{\cref{cdp:st:C.2.UN.4}, \cref{cdp:eq:C.EC.2}}

Use the marked fillings and component multisets of \cref{def:II:marked-torus-fillings} for a fixed proper family $p^\times:\mathcal T^\times\to\Delta^\times$ of connected complex two-tori.

For every integer $k$, the following equality holds between sets of
finite multisets:
\begin{equation}
 \begin{gathered}
 \{\mathcal E(f): f\text{ is a marked smooth filling of }p^\times\}\\
 =\{\mathcal E(f): f\text{ is such a filling},\ e((Y_0)_{\rm red})=2k\}.
 \end{gathered}                                                   \hypertarget{tag-00FO}{}\label{cdp:eq:C.EC.1}
\end{equation}
For each filling in the left-hand set, the representative on the
right is obtained by the explicit proper holomorphic vertical
blowup map of \ref{cdp:st:C.2.UN.4}. That map induces a bijection between the
original non-uniruled components and the resulting non-uniruled
components, and a bimeromorphism of their normalization resolutions.
Every original full-fibre multiplicity is preserved by this bijection.
The same assertion holds for compact smooth total spaces over a
fixed compact base curve, with the marked family fixed away from a
finite set $Q$. One can prescribe an arbitrary even Euler number at
each point of $Q$ simultaneously, preserving the multiset and
multiplicity at each point. In particular, the collection of
non-uniruled surface classes which occur in arbitrary compact
torus-fibred fillings is exactly the collection which occur in
compact fillings all of whose reduced fibres have Euler number zero.
\end{proposition}
\begin{proof}
\begin{enumerate}[label=\arabic*.,leftmargin=*]
\item Each blowup in \cref{cdp:st:C.2.UN.4} has a smooth compact point or curve as centre in the central support. Apply \cref{lem:II:nonuniruled-divisors} to its blowdown. The strict-transform restrictions are bimeromorphisms of the original components, with the same full-fibre vanishing orders; their normalization resolutions correspond. The exceptional projection is \eqref{cdp:eq:C.EC.2}, and its surface is uniruled. The resulting bijection therefore preserves exactly the multiset $\mathcal E(f)$, including repetitions and multiplicities. Iterating these specified maps preserves the same data.
\item Starting from any marked filling, \ref{cdp:st:C.2.UN.4} constructs a finite
sequence of exactly these maps whose final reduced Euler number is
$2k$. Every centre lies over zero, so the blowdown is a
biholomorphism over $\Delta^\times$ and transports the given marking
by composition. We have just proved that its multiset is the
original $\mathcal E(f)$. This gives inclusion of the left side of
\eqref{cdp:eq:C.EC.1} in the right. Inclusion in the other direction follows
because each filling with Euler number $2k$ is among the same
marked fillings defining the left side. This also proves the
assertion when there are no fillings: both sets are empty.
\item For the compact version, apply the same construction successively
at the finitely many points of $Q$. Every centre is a compact
analytic submanifold of the current compact threefold. Its global
blowup is compact and smooth, and the composite with the fixed base
map is proper holomorphic with connected fibres, by the
structure-sheaf and Stein-factor comparison in \ref{cdp:st:C.2.UN.4}. Its
restriction off that point is an isomorphism. Thus a construction
at one point preserves every fibre and every marking already
treated at another point. There are only finitely many critical
values of a compact proper map to a curve which is smooth
generically: the image of its compact analytic critical locus is
a proper analytic subset of the curve. Taking $Q$ to contain them
all and taking $k=0$ at each of them proves the final assertion;
the remaining smooth torus fibres already have Euler number zero.
\end{enumerate}
\end{proof}
\subsection{Non-bimeromorphic torus fillings}\label{subsec:II:distinct-torus-fillings}
\begin{proposition}[The same punctured torus family has non-bimeromorphic zero-Euler fillings]\hypertarget{tag-01D8}{}\label{cdp:st:C.2.UN.3}
\StatementRecord{cdp:st:C.2.UN.3}{proved}{01D8}{\cref{cdp:st:D.2.7}, \cref{cdp:eq:C.UC.2}, \cref{cdp:eq:C.UC.1}, \cref{cdp:eq:C.UC.4}, \cref{cdp:eq:C.UC.3}, \cref{cdp:eq:C.UC.6}, \cref{cdp:eq:C.UC.8}, \cref{cdp:eq:C.UC.9}, \cref{cdp:eq:C.UC.10}, \cref{cdp:eq:C.UC.7}, \cref{lem:II:nonuniruled-divisors}}
Let $A=E\times E$, $E=\C/(\Z+\mathrm i\Z)$.
There are two proper smooth-total-space fillings of
$A\times\Delta^*$ over a disc whose reduced central
supports both have Euler number zero, but for which
no bimeromorphic map over the disc exists. One is
$A\times\Delta$. The other is the local filling
\eqref{cdp:eq:D.23}, whose central support consists of a primary
Kodaira surface and a primary Hopf surface meeting
along one elliptic curve. The nonexistence persists
after every finite base change and resolution of the
normalized base change, with the given punctured
family fixed.
These fillings occur in compact smooth threefolds:
the principal $A$-bundle $Y_0$ and the threefold $Y$
constructed in \cref{cdp:st:D.2.7} are isomorphic away from the
three specified fibres. That isomorphism does not
extend bimeromorphically, even after a finite base
change of their base curve. In particular smooth
torus fillings, together with modifications by
blowups and blowdowns over the central fibre, do
not generate all zero-Euler fillings of this
punctured family.
More strongly, the same punctured product has
infinitely many pairwise non-bimeromorphic
zero-Euler fillings. They are obtained by replacing
$\OO_V(R)$ in \eqref{cdp:eq:C.UC.2} by $\OO_V(nR)$,
$n\ge1$. This distinction also persists after
arbitrary finite base change and resolution.
\end{proposition}
\begin{proof}
\begin{enumerate}[label=\arabic*.,leftmargin=*]
\item We use \cref{lem:II:nonuniruled-divisors} for the compact central divisors of these proper smooth disc models. Its diagram \eqref{cdp:eq:C.UC.1} gives the bimeromorphic correspondence of every component whose smooth normalization resolution is not uniruled.
\item Put
$q=e^{-2\pi}$ and
\begin{equation}
 V=\operatorname{Bl}_{(0,0)}(E\times\Delta),
 \qquad R=\operatorname{Exc}(V\to E\times\Delta),
 \qquad U=\OO_V(R)^\times/q^{\Z}.
\hypertarget{tag-014E}{}\label{cdp:eq:C.UC.2}
\end{equation}
The identity on $E\times\Delta$ lifts to the identity on $V$ and identifies the line bundle $L$ of \cref{night:P4.1} with $\OO_V(R)$, with $C=B=E$ and $A=C\times B$. That proposition gives the smooth proper elliptic-bundle quotient. Its fibre is identified with $E$ by $z\mapsto\exp(2\pi\mathrm i z)$: translation by $1$ fixes the exponential, translation by $\mathrm i$ multiplies it by $q$, and its kernel modulo $q^{\Z}$ is exactly $\Z+\mathrm i\Z$. The canonical-section trivialization of the same proposition gives
\begin{equation}
                         U|_{\Delta^*}=A\times\Delta^*.
\hypertarget{tag-014F}{}\label{cdp:eq:C.UC.3}
\end{equation}
This is exactly the construction in \cref{cdp:st:D.2.7},
including its holomorphic gluing to the
compact threefold.
\item The strict transform $E_0$ of $E\times\{0\}$
and the exceptional $R$ meet transversely at
one point. The restrictions of $\OO_V(R)$
to them have degree $+1$ and $-1$,
respectively. Their inverse images in $U$
are
\begin{equation}
 P=\OO_E(0)^\times/q^{\Z},\qquad
 Q=\OO_{\PP^1}(-1)^\times/q^{\Z}
       =(\C^2\setminus\{0\})/\langle q\rangle.
\hypertarget{tag-014G}{}\label{cdp:eq:C.UC.4}
\end{equation}
By \cref{night:P4.1}, $P$ is a primary Kodaira surface and $Q$ is the scalar Hopf surface in the display. They meet in one elliptic fibre $J$, and the map is $xy=t$ along $J$.
\item For the infinite family put
\begin{equation}
 U_n=\OO_V(nR)^\times/q^{\Z},\qquad
 P_n=\OO_E(n[0])^\times/q^{\Z},\qquad
 Q_n=\OO_{\PP^1}(-n)^\times/q^{\Z}.
\hypertarget{tag-014J}{}\label{cdp:eq:C.UC.7}
\end{equation}
If $g_{\alpha\beta}$ are the transition functions of $\OO_V(R)$, those of its $n$th tensor power are $g_{\alpha\beta}^n$. The quotient charts are therefore the same products with $E$ as in \cref{night:P4.1}, with these powered translation transitions; the maps to $V$ are smooth and proper. For local coefficients $s_\alpha$ of the canonical section, the ratios $\lambda_\alpha/s_\alpha^n$ agree on overlaps over the punctured disc. They give the product marking \eqref{cdp:eq:C.UC.3} for $U_n$. Its central support is $P_n\cup_J Q_n$, and the local base map along $J$ is still $xy=t$ because the projection to $V$ is a product in these charts.
Each
$P_n$ contains no rational curve by its
projection to $E$ and then to its elliptic
fibre. For $Q_n$, the usual quotient
$\OO_{\PP^1}(-n)^\times
=(\C^2\setminus\{0\})/\boldsymbol\mu_n$
is explicit: a nonzero vector $v$ represents
$v^{\otimes n}$ in the $n$th tensor power of
the tautological line. Multiplication by $q$
on that tensor power lifts to multiplication
by the positive real number $q^{1/n}$ on
$\C^2\setminus\{0\}$. Hence
\begin{equation}
 Q_n=(\C^2\setminus\{0\})/
       \langle\boldsymbol\mu_n,q^{1/n}\rangle,
 \qquad
 \pi_1(Q_n)=\Z\oplus\Z/n\Z.
\hypertarget{tag-014K}{}\label{cdp:eq:C.UC.8}
\end{equation}
The two commuting generators are independent,
because a nonzero power of $q^{1/n}$ has
absolute value different from one. The action
is free and properly discontinuous, and its
covering space retracts to the simply connected
$S^3$. This proves the fundamental-group
calculation and also proves that every map
$\PP^1\to Q_n$ is constant by lifting it to
the cover.
\item For $P_n$, the corresponding calculation
also determines the extension class. A Hermitian
metric on $\OO_E(n[0])$ identifies its nonzero
total space smoothly with its unit-circle bundle
times $\R_{>0}$. As $q$ is positive real, its
quotient is the unit-circle bundle times
$\R_{>0}/q^{\Z}\simeq S^1$. The circle bundle
has Chern number $n$. Lifting the two generators
of $\pi_1(E)$, the attaching map of the base
two-cell has winding number $n$ in the circle
fibre. With a consistent orientation this gives
the exact presentation
\begin{equation}
 \pi_1(P_n)=
 \langle a,b,c,d\mid
 [a,b]=c^n,\ [a,c]=[b,c]=1,
 [a,d]=[b,d]=[c,d]=1\rangle.
\hypertarget{tag-014L}{}\label{cdp:eq:C.UC.9}
\end{equation}
There are no further relations: the homotopy
sequence of the circle bundle has injective
fibre group because $\pi_2(E)=0$, and the last
factor is the displayed direct $S^1$ factor.
Abelianizing \eqref{cdp:eq:C.UC.8}--\eqref{cdp:eq:C.UC.9} yields
\begin{equation}
 H_1(P_n,\Z)=\Z^3\oplus\Z/n\Z,
 \qquad H_1(Q_n,\Z)=\Z\oplus\Z/n\Z.
\hypertarget{tag-014D}{}\label{cdp:eq:C.UC.10}
\end{equation}
\item Taking $n=1$ in the preceding curve and group calculations gives $H_1(P,\Z)=\Z^3$, $H_1(Q,\Z)=\Z$, and no rational curves on either surface.
Thus their resolutions are not covered by
rational curves. A bimeromorphic comparison
from $U$ to $A\times\Delta$ over $\Delta$
would, by \eqref{cdp:eq:C.UC.1}, have to carry each of
$P,Q$ birationally onto a central divisor.
The product has only the one central
divisor $A$, and neither $P$ nor $Q$ is
birational to it. Indeed, \cref{lem:II:blowup-topology}, applied to a point centre in a smooth compact surface, identifies its fundamental group with that of its blowup. Resolving a surface bimeromorphism into point blowups therefore preserves $H_1$ and $b_1$, whereas $b_1(A)=4$ differs from the ranks given by \eqref{cdp:eq:C.UC.10} at $n=1$. This proves the asserted nonexistence of a bimeromorphic comparison.
\item The principal elliptic-bundle projections of $P_n$ and $Q_n$ are locally trivial proper bundles, so their Euler numbers are the product of the Euler numbers of base and fibre, hence zero. The intersection $J$ also has Euler number zero. Additivity gives $e(P_n\cup_J Q_n)=0$ for every $n$; at $n=1$ this is
\begin{equation}
                    e(P\cup Q)=e(P)+e(Q)-e(J)=0=e(A).
\hypertarget{tag-014H}{}\label{cdp:eq:C.UC.5}
\end{equation}
\item Consider a finite base change.
Locally such a change has equation
$t=u(s)s^m$, $u(0)\ne0$. Choose a local holomorphic logarithm of $u$ and put $\widetilde s=s\exp((\log u(s))/m)$. Its derivative at zero is nonzero, so it is a coordinate and $t=\widetilde s^{m}$. Denote this new coordinate by $s$ in the following charts.
At a general point of either $P$ or $Q$
the original map is a submersion with
coordinate $x=t$. Its base change is the
smooth chart
\begin{equation}
                          x=s^m,
\hypertarget{tag-014I}{}\label{cdp:eq:C.UC.6}
\end{equation}
with the two surface coordinates unchanged.
The component over $s=0$ thus maps with
degree one to the original component.
At their intersection the normalized base
change has equation $xy=s^m$, with one
additional smooth elliptic coordinate.
This is normal, with transverse rational
double point $A_{m-1}$ when $m>1$.
A resolution adds ruled exceptional
surfaces above $J$; the strict transforms
of $P$ and $Q$ remain birational to the
surfaces in \eqref{cdp:eq:C.UC.4}. No choice of further
resolution changes those birational
classes. The target base change remains
the smooth product $A\times\Delta_s$.
The argument using \eqref{cdp:eq:C.UC.1} and their
first Betti numbers again excludes every
bimeromorphic comparison over $\Delta_s$.
\item Finally use the compact gluing of \cref{cdp:st:D.2.7}.
The original principal $A$-bundle $Y_0$
has a smooth abelian central fibre at
each of the three chosen base points.
The replacement $Y$ has the local model
\eqref{cdp:eq:C.UC.2} at each point and is identified
with $Y_0$ elsewhere by \eqref{cdp:eq:C.UC.3}.
An extension of this isomorphism would
respect the base, since the two base
maps agree on its dense open graph.
Its restriction near any of the three
fibres would be the already excluded
comparison. The same argument applies
above every preimage of these points
under a finite cover of the compact base,
using \eqref{cdp:eq:C.UC.6}. This proves the compact
and base-change assertions.
\item Integral first homology, including its torsion,
is invariant under a bimeromorphism of smooth
compact surfaces, by \cref{lem:II:blowup-topology} for each point blowup in a resolution. By \eqref{cdp:eq:C.UC.1},
a bimeromorphism $U_n\dashrightarrow U_{n'}$
over the disc must carry the two non-uniruled
central divisors birationally to the two such
divisors on the target. The ranks in \eqref{cdp:eq:C.UC.10}
force $P_n$ to correspond to $P_{n'}$ and
$Q_n$ to $Q_{n'}$. The torsion orders then
force $n=n'$. After any finite base change,
\eqref{cdp:eq:C.UC.6} still leaves these components
birationally unchanged, and the resolution of
$xy=s^m$ adds only ruled surfaces over $J$.
Further resolutions add only uniruled
exceptional divisors. They cannot replace a
non-uniruled component under \eqref{cdp:eq:C.UC.1}, so
the same torsion argument proves persistence.
Finally all $U_n$ have compact realizations:
replace $V$ by
$\operatorname{Bl}_{(0,0)}(E\times\PP^1)$
and use the same line bundle $\OO(nR)$.
Its elliptic quotient is a compact smooth
threefold, and the projection to $\PP^1$
restricts to \eqref{cdp:eq:C.UC.7} near zero.
\end{enumerate}
\end{proof}
\begin{remark}[Proper models of logarithmic complex tori]\hypertarget{tag-P237}{}\label{rem:II:logarithmic-proper-models}
\StatementRecord{rem:II:logarithmic-proper-models}{proved}{P237}{\cite[\S5.1.1--2, p.~170; Theorems~5.2.1 and 5.3.2, pp.~171 and 180]{KKN2008}}
For a nondegenerate pairing $\langle\ ,\ \rangle:X\times Y\to\mathbb G_{m,\log,S}$ with the admissible monoid presentation of \cite[\S5.1.1--2, p.~170]{KKN2008}, let $A$ be its associated log complex torus over $S$. Here $X$ and $Y$ are the two lattices of the pairing. A complete $Y$-stable fan $\Sigma$ in its associated cone, with finitely many $Y$-orbits, gives
\[
 V(\Sigma)=\bigcup_{\sigma\in\Sigma}V(\sigma),\qquad
 A^{(\Sigma)}=V(\Sigma)/Y\lhook\joinrel\longrightarrow A.
\]
The chart inclusions and the quotient by the lattice action define these maps. The complete fan exists by \cite[Theorem~5.2.1, p.~171]{KKN2008}; \cite[Theorem~5.3.2, p.~180]{KKN2008} makes its quotient a proper log smooth model over the given base. For the fillings in \cref{cdp:st:C.2.UN.3}, the prescribed map on the puncture is \eqref{cdp:eq:C.UC.3}, and the component homology in \eqref{cdp:eq:C.UC.10} obstructs its bimeromorphic extension.
\end{remark}


\section{Component covers and meromorphic fields}\label{sec:II:component-field-geometry}

\subsection{Cyclic component covers}\label{subsec:II:component-field-inputs}

\begin{lemma}[Cyclic covers of a fibre component]\hypertarget{tag-P240}{}\label{lem:II:cyclic-component-cover}
\StatementRecord{lem:II:cyclic-component-cover}{proved}{P240}{\cite[\S1, Universal and canonical desingularization, arXiv v2 p.~11]{BierstoneMilman1997}}
Let $f:U\to\Delta_t$ be a proper holomorphic map from a smooth complex threefold, with connected fibres and smooth complex two-tori over $\Delta_t^*$. Write $f^{-1}(0)=\sum_jm_jS_j$ and fix $S_i$. Normalize the fibre product for $t=s^{m_i}$. It has a resolution
$f':U'\to\Delta_s$ of its total space and central support, isomorphic over $\Delta_s^*$ and at the generic points above $S_i$. Every smooth strict component $E$ above $S_i$ has multiplicity one, and its proper generically finite map to $S_i$ factors through $S_i^\nu$. Its meromorphic function field is algebraic over the pullback of $\mathcal M(S_i^\nu)$, so $a(E)=a(S_i^\nu)$.

The resolution can be chosen equivariant for $\mu_{m_i}$. For such a choice, let $G\subset\mu_{m_i}$ be the stabilizer of $E$. The generic degree of $E\to S_i^\nu$ is $|G|$, and the induced map \eqref{cdp:eq:C.HQ.2} is proper and bimeromorphic. When $E$ has a nonruled minimal model $E_0$, its contraction $\rho:E\to E_0$ is equivariant and induces the proper bimeromorphic map \eqref{cdp:eq:C.HQ.3}. Resolving the graph of these two quotient maps compares a smooth resolution of $S_i^\nu$ with a resolution of $E_0/G$.
\end{lemma}
\begin{proof}
At a generic smooth point of $S_i$ away from the other components, write $t=uz^{m_i}$ with $u$ a holomorphic unit. The normalized root cover is given by
\begin{equation}
 w^{m_i}=u,\qquad z=s/w.
 \hypertarget{tag-010M}{}\label{cdp:eq:C.SC.2}
 \hypertarget{tag-00MQ}{}\SsixEquationAliases{\label{cdp:eq:C.HQ.1}}
\end{equation}
The map to $U$ keeps the other coordinates and sends $(s,w)$ to $z=s/w$; its deck action sends $(s,w)$ to $(\zeta s,\zeta w)$. Since $w$ is a unit and the root equation is \'{e}tale, the total space is smooth here and $s$ has order one on every component above $S_i$.

Resolve the normalized total space and then the central support, preserving the punctured tube and these generic smooth points. The canonical analytic resolution in \cite[\S1, Universal and canonical desingularization, arXiv v2 p.~11]{BierstoneMilman1997} lifts local isomorphisms through its successive blowups. Apply it to the space and then to the central divisor: the finite deck action preserves each centre and lifts to the resulting pair. Properness over the disc makes the central fibre compact; after shrinking the disc finitely many steps suffice. The nearby fibres are the original connected tori. In the Stein factorization of $f'$, the finite map to the normal disc has generic degree one and is therefore an isomorphism. Thus its full fibres are connected.

The strict component $E$ is compact and smooth. The universal property of normalization gives $E\to S_i^\nu\to S_i$. Its Stein factorization $E\to T'\to S_i^\nu$ has a birational first map and a finite second map. For a meromorphic function on $T'$, take the elementary symmetric functions of its values over the finite \'{e}tale locus. Near a point of the normal base, multiplication by a sufficiently high power of a local equation containing the finite image of its polar divisor makes that function holomorphic upstairs. The symmetric functions of this multiple are locally bounded, by finiteness, and extend holomorphically across the branch set by normality. Dividing the $k$th coefficient by the $k$th power of that base multiplier gives a meromorphic extension of the original coefficient. These global coefficients give a monic equation over $\mathcal M(S_i^\nu)$. Birational pullback identifies $\mathcal M(T')$ with $\mathcal M(E)$, proving the function-field and algebraic-dimension assertions.

For the equivariant choice, the reduced inverse image over the connected dense smooth open subset of $S_i$ in \eqref{cdp:eq:C.SC.2} is a $\mu_{m_i}$-torsor. The group permutes its connected components transitively, and the stabilizer of the component contained in $E$ acts freely and transitively on each of its fibres. Thus its degree is $|G|$. Taking the finite quotient and using the proper factorization through normalization gives
\begin{equation}
 E/G\longrightarrow S_i^\nu,
 \hypertarget{tag-00MS}{}\label{cdp:eq:C.HQ.2}
\end{equation}
a degree-one proper map between irreducible normal surfaces, hence a bimeromorphic map. This uses the stabilizer of the selected component even when the root cover over $S_i$ is disconnected.

For a nonruled minimal model $E_0$, each $g\in G$ induces a bimeromorphic automorphism $g_0$ of $E_0$, which is holomorphic by uniqueness of the minimal model. The equality $\rho g=g_0\rho$ on the common open isomorphism extends everywhere. The contraction therefore descends to
\begin{equation}
 E/G\longrightarrow E_0/G.
 \hypertarget{tag-00MT}{}\label{cdp:eq:C.HQ.3}
\end{equation}
It is proper and bimeromorphic since the contraction is so on the invariant open set. Passing to the effective quotient of the action on $E_0$ preserves the quotient space. The two proper quotient maps give the asserted graph comparison.
\end{proof}

\subsection{Semicomplete fields and their restrictions}\label{subsec:II:semicomplete-field-inputs}

\begin{definition}[Semicomplete meromorphic vector fields]\hypertarget{tag-P238}{}\label{def:II:semicomplete-fields}
\StatementRecord{def:II:semicomplete-fields}{proved}{P238}{\cite[Definition~2.3, author preprint p.~7]{Guillot2024}; \cite[Definition~5, arXiv v2 p.~4]{Guillot2019}}
For a meromorphic vector field $V$ on a complex manifold $X$, let $X_V$ be its holomorphic locus. A solution $\phi:\Omega\to X_V$, with $\Omega\subset\C$ connected and open, is \emph{univalent maximal} when its graph map $t\mapsto(t,\phi(t))$ to $\C\times X_V$ is proper. The field is \emph{semicomplete} when every initial point in $X_V$ has such a solution. Equivalently, analytic continuations of a solution germ along two complex-time paths with the same endpoint agree wherever both are defined. We use this equivalence in the sense of \cite[Propositions~2.1--2.2, author preprint pp.~4--7]{Guillot2024}.
\end{definition}

\begin{lemma}[Restriction and specialization of semicomplete fields]\hypertarget{tag-P239}{}\label{lem:II:semicomplete-transport}
\StatementRecord{lem:II:semicomplete-transport}{proved}{P239}{\cref{def:II:semicomplete-fields}; \cite[Theorem~3.3, author preprint p.~9]{Guillot2024}; \cite[Lemma~6, arXiv v2 pp.~4--5]{Guillot2019}}
Semicompleteness is preserved by restriction to an open set and by transport through a bimeromorphic map between smooth manifolds, on the full holomorphic locus of the transported meromorphic field.

Let $f':U'\to\Delta_s$ be a proper holomorphic map from a smooth complex threefold with smooth compact two-torus fibres over $\Delta_s^*$. Let $v$ be a meromorphic vertical field, holomorphic over $\Delta_s^*$, and let $E$ be a smooth central component at whose generic point $v$ is holomorphic. Then $v$ is semicomplete on its holomorphic locus, and its meromorphic restriction $v_E$ is tangent to $E$ and semicomplete on its holomorphic locus.

A nonzero semicomplete meromorphic field on a smooth compact curve is holomorphic, and that curve has genus at most one.
\end{lemma}
\begin{proof}
Restriction preserves single-valued analytic continuation, or equivalently the fourth condition in \cite[Proposition~2.2 and the paragraph following Definition~2.3, author preprint pp.~6--7]{Guillot2024}. The uniformizability locus of a holomorphic field on any complex manifold is closed by \cite[Theorem~3.3, author preprint p.~9]{Guillot2024}. A field transported by a bimeromorphic map agrees with the old one under the biholomorphism of their dense common open sets. Apply restriction there, and then this closedness theorem on the full holomorphic locus of the transported field.

On every punctured fibre of $f'$, the field $v$ is a holomorphic translation field on a compact torus. In its presentation $\C^2/\Lambda$, the flow is $[z]\mapsto[z+\tau a]$ for the corresponding constant vector $a$, and is defined for every $\tau\in\C$. The union of these fibres is dense in the holomorphic locus $W$ of $v$. Closedness therefore makes $v|_W$ semicomplete.

The identity $v(s)=0$ preserves each component ideal of the central divisor wherever $v$ is holomorphic. Indeed, write $s=u\prod_jh_j^{n_j}$ in a smooth local ring. After dividing $v(s)=0$ by $h_j^{n_j-1}$ and reducing modulo $h_j$, the remaining nonzero factors give $v(h_j)=0\bmod h_j$. Thus $E\cap W$ is invariant, and its solutions inherit single-valued continuation under the inclusion into $W$. Its meromorphic field $v_E$ may extend holomorphically across further curves in $E$. Closedness on the holomorphic locus of $v_E$, applied to the dense open subset $E\cap W$, supplies the stated extension.

For a compact curve $C$ and a nonzero field $V$, apply \cite[Lemma~6, arXiv v2 pp.~4--5]{Guillot2019} to $V$ and a univalent maximal solution on $C_V$. The lemma proves that $V$ has no poles. Its nonzero holomorphic section of $T_C$ has an effective zero divisor of degree $\deg T_C=2-2g(C)$, proving $g(C)\le1$.
\end{proof}

\begin{lemma}[Two semicomplete fields in algebraic dimension zero]\hypertarget{tag-P301}{}\label{lem:II:zero-dimensional-field-surface}
\StatementRecord{lem:II:zero-dimensional-field-surface}{proved}{P301}{\cref{lem:II:semicomplete-transport}; \cite[Theorem~9.2, author preprint p.~28]{Guillot2024}; \cite[Theorem~0.3, pp.~641--642]{DOT2001}; \cite[p.~260]{DOT2000}; \cite[Proposition~2(ii), pp.~274--275; Proposition~3(ii), p.~277; Proposition~5, p.~280]{Inoue1974}}
Let $X$ be a smooth compact complex surface with $a(X)=0$ and a two-dimensional complex space of commuting meromorphic fields every member of which is semicomplete. Its minimal model is a complex two-torus or a Hopf surface.
\end{lemma}
\begin{proof}
Choose a complex basis $V_1,V_2$. Apply \cite[Theorem~9.2, author preprint p.~28]{Guillot2024} to each field. Its fibration alternative would give a nonconstant meromorphic function by pullback from its compact base curve. Since $a(X)=0$, each field instead becomes holomorphic on a smooth bimeromorphic model. Such a surface is nonruled: a ruling would again pull back a nonconstant function from its base, and a rational surface has algebraic dimension two. Its nonruled minimal model $X_0$ is unique, and each smooth bimeromorphic model maps to it by point blowdowns.

For such a blowdown $\pi$, the equality $\pi_*\OO=\OO$ descends the derivation $g\mapsto V(\pi^*g)$ to a holomorphic vector field on the base. On the common open isomorphism this is the original meromorphic field. Thus $V_1,V_2$ give two complex-linearly independent commuting holomorphic fields on the same $X_0$.

Apply the complete minimal-surface list in \cite[Theorem~0.3, pp.~641--642]{DOT2001}. Its ruled and elliptically fibred entries have nonconstant meromorphic functions by their indicated projections. The entries remaining in algebraic dimension zero are tori, Hopf surfaces, Inoue surfaces, and minimal class VII surfaces with positive second Betti number. The last have at most one holomorphic field by \cite[p.~260]{DOT2000}. For Inoue's three families the dimensions are respectively $0,1,0$, by \cite[Proposition~2(ii), pp.~274--275; Proposition~3(ii), p.~277; Proposition~5, p.~280]{Inoue1974}. The two independent descended fields therefore leave precisely the torus and Hopf entries.
\end{proof}

\begin{lemma}[The fibration defined by algebraic solution leaves]\hypertarget{tag-P302}{}\label{lem:II:algebraic-leaf-fibration}
\StatementRecord{lem:II:algebraic-leaf-fibration}{proved}{P302}{\cref{lem:II:semicomplete-transport}}
Let $E$ be a smooth projective surface with a nonzero semicomplete meromorphic field $w_E$ whose regular solution leaves have algebraic curves as their Zariski closures. There are a proper birational morphism $\rho:E'\to E$ and a proper map
\begin{equation}
 q:E'\longrightarrow C
 \hypertarget{tag-010Q}{}\label{cdp:eq:C.SC.6}
\end{equation}
to a smooth compact curve with connected fibres, such that the lifted field is tangent to $q$ and its general fibre is rational or elliptic. In particular $\kappa(E)\le1$.
\end{lemma}
\begin{proof}
The irreducible algebraic closures of the regular leaves are invariant curves covering the regular foliation locus. Use the countably many Hilbert polynomials and flattening strata of their universal families to parametrize reduced irreducible curves. After clearing the polar divisor $P$ of $w_E$, tangency is the vanishing of the induced relative conormal map
\[
 \mathcal I/\mathcal I^2\longrightarrow\OO_{\mathcal C}(P),
 \qquad h\longmapsto w_E(h)\bmod\mathcal I
\]
on the universal curve $\mathcal C\subset E\times T$. Coherent presentations and flattening stratification make its fibrewise vanishing constructible. Thus the integral invariant curves have countably many irreducible locally closed parameter strata with constructible incidence images in $E$.

Some incidence image is dense. Otherwise their Zariski closures would be countably many proper closed analytic subsets covering the regular foliation locus, contrary to the Baire theorem. In a stratum with dominant incidence, exactly one generic curve passes through a regular point: two invariant irreducible curves through it contain the same local solution germ and coincide. The parameter stratum therefore has dimension one and its incidence map to $E$ is generically one-to-one. Compactify and normalize this parameter curve and close the incidence graph. Its birational projection to $E$, inverted and followed by projection to the parameter, gives a meromorphic first integral. Resolving its graph and taking the Stein factor gives \eqref{cdp:eq:C.SC.6}.

The meromorphic lift of $w_E$ is tangent to $q$, is nonzero on a general smooth fibre $G$, and is semicomplete on its holomorphic locus by \cref{lem:II:semicomplete-transport}. Restriction to the invariant fibre and the curve clause of that lemma give $g(G)\le1$.

Since $G$ is a smooth fibre, its normal line is trivial and adjunction gives $K_{E'}|_G=K_G$. A negative-degree line on a compact curve has no sections, and a degree-zero line has at most one independent section. Each pluricanonical map therefore contracts the general fibres of $q$. It cannot be generically finite, so $\kappa(E')\le1$; the proper birational map $\rho$ gives $\kappa(E)=\kappa(E')$.
\end{proof}

\subsection{Translation fields on component covers}\label{subsec:II:component-translation-fields}

\begin{proposition}[Commuting semicomplete fields on a component cover]\hypertarget{tag-01CT}{}\label{cdp:st:C.2.UN.16}
\StatementRecord{cdp:st:C.2.UN.16}{proved}{01CT}{\cref{lem:II:cyclic-component-cover}, \SeriesRef{cdp:eq:C.CF.3}, \SeriesRef{cdp:eq:C.CF.5}, \SeriesRef{cdp:eq:C.CF.7}, \SeriesRef{cdp:eq:C.CF.8}; \cite[Theorem~3.3]{Guillot2024}; \cite[Theorem~9.2]{Guillot2024}; \cite[Theorem~0.3]{DOT2001}; \cite[pp.~256 and~260]{DOT2000}; \cite[Propositions~2(ii), 3(ii) and~5]{Inoue1974}}

Let $f:U\to\Delta$ be a proper holomorphic map from a smooth
complex threefold, with connected fibres and with smooth complex
two-tori over $\Delta^*$. Fix an irreducible central component $S_i$
of multiplicity $m_i$. Make the base change $t=s^{m_i}$,
normalize, and resolve as in \cref{lem:II:cyclic-component-cover}. Every chosen smooth strict
component $E$ above $S_i$ has multiplicity one and carries a
two-dimensional complex vector space
\begin{equation}
 \mathfrak a_E\subset
 H^0(E,T_E\otimes_{\mathcal O_E}\mathcal M_E)
 \hypertarget{tag-00BN}{}\label{cdp:eq:C.CF.1}
\end{equation}
of commuting meromorphic vector fields. Every member of
$\mathfrak a_E$ is semicomplete on its holomorphic locus.
The map $E\to S_i$ factors through normalization and is
proper and generically finite.

There are two exhaustive geometric alternatives for this particular
space of fields. Its generic rank is two, in which case any complex
basis is a commuting meromorphic frame; or its generic rank is one,
in which case a resolution of a meromorphic first integral
and its Stein factor give
\begin{equation}
 \widehat E\xrightarrow{\ q\ }B,
 \qquad \widehat E\xrightarrow{\rho}E,
 \hypertarget{tag-00BO}{}\label{cdp:eq:C.CF.2}
\end{equation}
where $\rho$ is proper birational, $B$ is a smooth compact curve,
and the connected general fibre of $q$ is rational or elliptic.

When $a(S_i^\nu)=0$, the unique minimal model of $E$ is a
two-torus or a Hopf surface. 
\end{proposition}


\begin{proof}

\begin{enumerate}[label=\arabic*.,leftmargin=*]
\item Let $f':U'\to\Delta_s$ and $E$ be the resolution and component of \cref{lem:II:cyclic-component-cover}. Its local map \eqref{cdp:eq:C.SC.2} gives multiplicity one, proper generic finiteness and the factorization through $S_i^\nu$. Define
\begin{equation}
 \mathcal T_{U'/\Delta_s}
 =\ker\bigl[df':T_{U'}\longrightarrow f'^*T_{\Delta_s}\bigr].
 \hypertarget{tag-010O}{}\label{cdp:eq:C.SC.4}
\end{equation}

\item Put $\mathcal O=\mathbb C\{s\}$ and
$K=\operatorname{Frac}(\mathcal O)$. Let
\begin{equation}
 M=(f'_*\mathcal T_{U'/\Delta_s})_0,
 \qquad
 L=\{v\in M\otimes_{\mathcal O}K:
                     \operatorname{ord}_E(v)\ge0\}.
 \hypertarget{tag-00BP}{}\label{cdp:eq:C.CF.3}
 \hypertarget{tag-00K4}{}\SsixEquationAliases{\label{cdp:eq:C.FD.3}}
\end{equation}
The order in this formula is the minimum of the coefficient
orders in a local tangent frame at the generic point of $E$.
It is independent of that frame. The direct image is coherent
by properness, is torsion-free because its sections are vector fields on $U'$, and has rank two by base change on the
smooth torus locus. Consequently $M$ is a free rank-two
$\mathcal O$-module. Its sections are identified with holomorphic vertical fields by the definition of direct image.
Any two of them commute: their restrictions to a smooth torus
are translation fields, so their bracket vanishes on the dense
punctured tube and hence everywhere.

\item Choose an $\mathcal O$-basis $v_1,v_2$ of $M$.
At a generic point of $E$, multiplicity one makes $s$ a
coordinate, so the two vertical tangent coordinates $x,y$
give
\begin{equation}
 (v_1\ v_2)=(\partial_x\ \partial_y)A,
 \qquad d=\operatorname{ord}_E(\det A)<\infty.
 \hypertarget{tag-00BQ}{}\label{cdp:eq:C.CF.4}
\end{equation}
All entries of $A$ are holomorphic there. The determinant is
nonzero because the evaluation of the two translation fields
is an isomorphism on each nearby torus. In particular $d\ge0$.
For $v=a_1(s)v_1+a_2(s)v_2\in L$, the column $Aa$ has
nonnegative order along $E$. The identity
$a=(\det A)^{-1}\operatorname{adj}(A)(Aa)$ shows that
each $a_j$ has order at least $-d$. The order of a function
of $s$ along $E$ equals its $s$-adic order, again by
multiplicity one. We obtain inclusions
\begin{equation}
                        M\subset L\subset s^{-d}M.
 \hypertarget{tag-00BR}{}\label{cdp:eq:C.CF.5}
\end{equation}
The set $L$ is an $\mathcal O$-submodule. Since $\mathcal O$
is noetherian, \eqref{cdp:eq:C.CF.5} makes it finite; since it is torsion-free
of rank two over the discrete valuation ring $\mathcal O$,
it is free of rank two.

\item Every member of $L$ is a global meromorphic vertical field
on a sufficiently small representative of $U'$, regular at
the generic point of $E$. Its restriction is therefore a
global meromorphic section of $T_E$. Tangency follows from
$v(s)=0$ at the generic smooth point and then as an identity
of meromorphic sections. Restriction defines the exact map
\begin{equation}
 \operatorname{res}_E:L\longrightarrow
 H^0(E,T_E\otimes\mathcal M_E),
 \qquad
 \ker(\operatorname{res}_E)=sL.
 \hypertarget{tag-00BS}{}\label{cdp:eq:C.CF.6}
\end{equation}
Indeed, zero restriction means order at least one along $E$;
division by $s$ then belongs to $L$. Conversely, every member
of $sL$ restricts to zero. Thus
\begin{equation}
                  L/sL\simeq\mathbb C^2
                    \hookrightarrow
                  H^0(E,T_E\otimes\mathcal M_E).
 \hypertarget{tag-00BT}{}\label{cdp:eq:C.CF.7}
\end{equation}
Define $\mathfrak a_E$ to be its image. Since vertical fields
annihilate functions of $s$, brackets of members of
$M\otimes K$ vanish. Hence the fields in \eqref{cdp:eq:C.CF.7} commute.

\item For $V\in\mathfrak a_E$, choose a representative $v\in L$. It is a meromorphic vertical field on $U'$, holomorphic over $\Delta_s^*$ and at the generic point of $E$. Apply \cref{lem:II:semicomplete-transport} with this $f'$, $v$ and $E$. It proves semicompleteness of $V=v|_E$ on its full holomorphic locus, for every member of $\mathfrak a_E$.

For a particular nonzero section $v\in M$, put $k=\operatorname{ord}_E(v)$ and
\begin{equation}
 w=s^{-k}v.
 \hypertarget{tag-010P}{}\label{cdp:eq:C.SC.5}
\end{equation}
In the vertical frame $(\partial_x,\partial_y)$, $k$ is the minimum of the two coefficient orders. Thus $w\in L$ has order zero and its class in $L/sL$ gives a nonzero field $w_E\in\mathfrak a_E$.

\item Take a complex basis $V_1,V_2$ of $\mathfrak a_E$.
When $V_1\wedge V_2$ is nonzero, the map
$(a,b)\mapsto aV_1+bV_2$ identifies $\mathcal M_E^2$
with $T_E\otimes\mathcal M_E$, giving the meromorphic frame. Otherwise there is a globally defined
meromorphic function $h$ with
\begin{equation}
                  V_2=hV_1,
                  \qquad V_1(h)=0.
 \hypertarget{tag-00BU}{}\label{cdp:eq:C.CF.8}
\end{equation}
The local coefficient ratios agree on their overlaps,
which proves that $h$ is global. It is nonconstant by
the complex-linear independence in \eqref{cdp:eq:C.CF.7}. The second
identity follows from the vanishing bracket.
Resolve the graph of $h:E\dashrightarrow\mathbb P^1$
and take the Stein factor of its projection to
$\mathbb P^1$. This gives exactly \eqref{cdp:eq:C.CF.2}, with a
finite map $B\to\mathbb P^1$. The lifted field is
tangent to $q$ and nonzero on a general fibre.
Apply the bimeromorphic transport in
\cref{lem:II:semicomplete-transport} to this graph resolution.
It preserves semicompleteness on the lifted field's
full holomorphic locus.

\item On a general smooth fibre of \eqref{cdp:eq:C.CF.2}, the lifted field is nonzero, meromorphic and tangent. Its semicompleteness and the compact-curve clause of \cref{lem:II:semicomplete-transport} give a rational or elliptic fibre.

\item When $a(S_i^\nu)=0$, the finite Stein-factor comparison of \cref{lem:II:cyclic-component-cover} gives
\begin{equation}
 a(E)=a(S_i^\nu)=0.
 \hypertarget{tag-00BV}{}\label{cdp:eq:C.CF.9}
\end{equation}
The fields in \eqref{cdp:eq:C.CF.7} satisfy all hypotheses of \cref{lem:II:zero-dimensional-field-surface}. Its two descended holomorphic fields on the common minimal model give the stated torus or Hopf alternative.

\end{enumerate}
\end{proof}

\begin{proposition}[No component of a torus degeneration has general-type normalization]\hypertarget{tag-01CS}{}\label{cdp:st:C.2.UN.15}
\StatementRecord{cdp:st:C.2.UN.15}{proved}{01CS}{\SeriesRef{cdp:eq:C.SC.3}, \SeriesRef{cdp:eq:C.SC.6}; \cite[Theorem~3.3, with its proof]{Guillot2024}; \cite{Guillot2019}}

Let $f:U\to\Delta$ be a proper holomorphic map from
a smooth complex threefold to a disc, with connected
fibres and smooth fibres complex two-tori over
$\Delta^*$. For every irreducible component $S_i$
of the reduced central support, its smooth
normalization resolution $Z_i$ satisfies
\begin{equation}
                            \kappa(Z_i)\le1.
 \hypertarget{tag-010L}{}\label{cdp:eq:C.SC.1}
\end{equation}

\end{proposition}


\begin{proof}
Fix $S_i$ and take the component cover $E\to S_i^\nu$ of \cref{lem:II:cyclic-component-cover}. If $Z_i$ had general type, resolve the graph of the generically finite meromorphic map $E\dashrightarrow Z_i$ to obtain
$a:\widehat E\to E$ birational and $b:\widehat E\to Z_i$ generically finite. The Jacobian pullback gives
\begin{equation}
 K_{\widehat E}=b^*K_{Z_i}+R,\qquad R\ge0.
 \hypertarget{tag-010N}{}\label{cdp:eq:C.SC.3}
\end{equation}
It injects pluricanonical sections. A generically finite pluricanonical map of $Z_i$, composed with $b$, remains generically finite; multiplication by the canonical section of $mR$ gives its subsystem in $|mK_{\widehat E}|$. Thus $K_{\widehat E}$ is big and $\kappa(E)=2$. A smooth compact surface of general type is projective.

Apply \cref{cdp:st:C.2.UN.16} to this same resolved base change. The field $w_E$ given by \eqref{cdp:eq:C.SC.5} is nonzero and semicomplete. A Zariski-dense maximal solution satisfies the hypotheses of \cite[Main Theorem and the following clarification, arXiv v2 p.~2; Proposition~22, p.~25]{Guillot2019}: the surface $E$ is projective, hence K\"ahler, and the solution is univalent maximal on the holomorphic locus of $w_E$ by \cref{def:II:semicomplete-fields}. The resulting birational model is Abelian or admits a rational or elliptic fibration. Its Kodaira dimension is at most one, contrary to $\kappa(E)=2$.

When every regular leaf has an algebraic curve as its Zariski closure, apply \cref{lem:II:algebraic-leaf-fibration} to the projective surface $E$ and the same $w_E$. Its graph and Stein map \eqref{cdp:eq:C.SC.6} give $\kappa(E)\le1$, the same contradiction. These two cases exhaust the regular leaves, so \eqref{cdp:eq:C.SC.1} follows for the chosen component and hence for every $S_i$.
\end{proof}

\subsection{Canonical divisors on component surfaces}\label{subsec:II:component-canonical-support}

\begin{proposition}[Canonical support forces the numerically anticanonical class VII case]\hypertarget{tag-01CU}{}\label{cdp:st:C.2.UN.17}
\StatementRecord{cdp:st:C.2.UN.17}{proved}{01CU}{\SeriesRef{cdp:st:A.5.3}, \SeriesRef{cdp:eq:C.CA.4}, \SeriesRef{cdp:eq:C.CA.3}, \SeriesRef{cdp:eq:C.CA.5}, \SeriesRef{cdp:eq:C.CA.7}, \SeriesRef{cdp:eq:C.CA.9}, \SeriesRef{cdp:eq:C.CA.6}; \cite[Remark~2.7, arXiv v2 p.~9]{DlousskyTeleman2012}; \cite[\S1, pp.~4--5]{Dloussky2006}; \cite[Theorem~0.2, arXiv v3 p.~2; Definition~2.20, p.~14]{Dloussky2006}}

Let $f:U\to\Delta$ be a proper holomorphic map from a smooth
complex threefold, with connected fibres and smooth two-torus
fibres over $\Delta^*$. Write its full central fibre as
$F=\sum_i m_iS_i$. For each component let $\eta_i:T_i\to S_i$ be its normalization,
$\pi_i:Z_i\to T_i$ its minimal resolution, and $h_i=\eta_i\pi_i$.
Let $D_i\ge0$ be the intrinsic resolution conductor of
\cref{cdp:st:A.5.3}, applied to the reduced hypersurface $S_i\subset U$. There are nonnegative integers $a_j$ with
$K_U\simeq\OO_U(\sum_j a_jS_j)$. Put
$E_{ij}=h_i^*(S_j)$ for $j\ne i$, using the effective
Cartier pullback. The exact line-bundle identity is
\begin{equation}
 K_{Z_i}^{\otimes m_i}\otimes\OO_{Z_i}(m_iD_i)
 \simeq
 \OO_{Z_i}\!\left(
  \sum_{j\ne i}\bigl[m_i a_j-m_j(a_i+1)\bigr]E_{ij}
 \right).                                                \hypertarget{tag-00BC}{}\label{cdp:eq:C.CA.1}
\end{equation}
In particular a positive tensor power of the canonical line on the minimal
model of $Z_i$ has a meromorphic section whose divisor is
supported on the images of these interface and conductor
curves.

When this minimal model $M$ is of class VII, it has the
following restrictions. If $b_2(M)=0$, it is a Hopf surface.
If $b_2(M)>0$, there are an integer $n>0$, a nonzero
effective integral divisor $A$ and a complex-flat line
bundle $L$ such that
\begin{equation}
                 K_M^{\otimes n}\otimes\OO_M(A)
                         \simeq L.                       \hypertarget{tag-00BF}{}\label{cdp:eq:C.CA.2}
\end{equation}
Consequently $M$ contains a global spherical shell.
\end{proposition}


\begin{proof}

\begin{enumerate}[label=\arabic*.,leftmargin=*]
\item After shrinking the disc, $f_*K_U$ is a locally free
rank-one sheaf. Coherence follows from properness, and
torsion-freeness follows because a holomorphic canonical
section killed by a power of the base coordinate vanishes
on a dense open subset. On a smooth fibre $F_t$, adjunction identifies
$K_U|_{F_t}$ with $K_{F_t}\otimes T^*_{\Delta,t}$; the base form
$dt$ trivializes its second factor. Thus the direct image has
rank one because a two-torus has a one-dimensional space of
canonical forms. Choose a frame and evaluate it on $U$.
Grauert base change identifies its restriction over
$\Delta^*$, through this adjunction map, with a nonzero
torus canonical form tensored by $dt$, which is nowhere zero.
The evaluation is therefore a nonzero holomorphic section
of $K_U$ with divisor $\sum_j a_jS_j$, where $a_j\ge0$.
This constructs the claimed line-bundle isomorphism.

\item Let $N_i=h_i^*\OO_U(S_i)$. The canonical sections defining
$S_j$ do not vanish identically on $S_i$ for $j\ne i$.
Their pullbacks define the effective divisors
$E_{ij}$, including any exceptional orders. Pullback of
the principal full fibre gives
\begin{equation}
 N_i^{\otimes m_i}\otimes
 \OO_{Z_i}\!\left(\sum_{j\ne i}m_jE_{ij}\right)
                         \simeq\OO_{Z_i}.               \hypertarget{tag-00BG}{}\label{cdp:eq:C.CA.3}
\end{equation}
Adjunction for the individual reduced hypersurface $S_i$,
followed by the conductor comparison of \cref{cdp:st:A.5.3}, gives
\begin{equation}
 \begin{split}
 K_{Z_i}\otimes\OO_{Z_i}(D_i)
   &\simeq h_i^*\bigl((K_U\otimes\OO_U(S_i))|_{S_i}\bigr)\\
   &\simeq N_i^{\otimes(a_i+1)}\otimes
      \OO_{Z_i}\!\left(\sum_{j\ne i}a_jE_{ij}\right).
 \end{split}                                             \hypertarget{tag-00BH}{}\label{cdp:eq:C.CA.4}
\end{equation}
Raising \eqref{cdp:eq:C.CA.4} to the $m_i$-th power and substituting
\eqref{cdp:eq:C.CA.3} gives \eqref{cdp:eq:C.CA.1}.

\item Let $\tau:Z_i\to M$ be the sequence of point blowdowns
to a minimal model. Identity \eqref{cdp:eq:C.CA.1} provides a
meromorphic section of $K_{Z_i}^{\otimes m_i}$ with the
displayed signed divisor, less $m_iD_i$. On the open set
where $\tau$ is an isomorphism this is a meromorphic
$m_i$-canonical section on $M$. In local coordinates at
each omitted point, the coefficient of that section is a
meromorphic function on a punctured smooth surface and
extends meromorphically. Thus the section extends across
those points. Its divisor is precisely the pushforward:
exceptional curves disappear, and every other curve
keeps its order. Equivalently, at an individual blowup
the exact canonical formula
$K_{\widehat M}=\tau^*K_M+E$ removes the corresponding
$m_iE$ from the pulled-back pluricanonical divisor.
We have therefore constructed a signed integral divisor
$B$ on $M$ and an isomorphism
\begin{equation}
                         K_M^{\otimes m_i}\simeq\OO_M(B).
                                                               \hypertarget{tag-00BI}{}\label{cdp:eq:C.CA.5}
\end{equation}
Every component of $B$ is the image of an interface or
intrinsic conductor curve from \eqref{cdp:eq:C.CA.1}.

\item Suppose now that $M$ is of class VII. If it contained no
curve, \eqref{cdp:eq:C.CA.5} would have $B=0$ and would trivialize a
positive power of $K_M$. This contradicts
$\kappa(M)=-\infty$. When $b_2(M)=0$, choose a nonempty
curve. A positive multiple of it is homologically trivial,
since its integral class is torsion. The classification
of such effective divisors in
\cite[Remark~2.7, arXiv v2 p.~9]{DlousskyTeleman2012} gives a Hopf
surface in this zero-Betti case. 

\item For $b_2(M)>0$, let $C_1,\ldots,C_r$ be the distinct
curves in the support of $B$, and write
\begin{equation}
 c_1(K_M)_{\mathbb R}=\sum_{\ell=1}^r c_\ell[C_\ell],
 \qquad c_\ell\in\Q.                                  \hypertarget{tag-00BJ}{}\label{cdp:eq:C.CA.6}
\end{equation}
The intersection form on $H^2(M,\mathbb R)$ is negative
definite and $K_M^2=-b_2(M)$. Moreover
$K_M\cdot C_\ell\ge0$ for every curve on the minimal
surface; these are the class VII facts in
\cite[\S1, arXiv v3 p.~4, and Theorem~1.8, p.~5]{Dloussky2006}.
The cited surface is our minimal $M$ with $b_1(M)=1$ and
$b_2(M)>0$, and its intersection pairing is the one used here.

\item Let $Q=(C_\ell\cdot C_k)$ and
$v_\ell=K_M\cdot C_\ell$. Then
\begin{equation}
                         Qc=v,\qquad v\ge0.             \hypertarget{tag-00BK}{}\label{cdp:eq:C.CA.7}
\end{equation}
The matrix $Q$ is negative semidefinite, and all its
off-diagonal entries are nonnegative. Separate it into
blocks according to the connected components of the
intersection graph. For a negative-definite block, write
$c=c_+-c_-$ with disjoint nonnegative coordinate supports.
If $c_+\ne0$, then
\begin{equation}
 c_+^tQc=c_+^tQc_+-c_+^tQc_-<0,
 \qquad c_+^tv\ge0,                                    \hypertarget{tag-00BL}{}\label{cdp:eq:C.CA.8}
\end{equation}
contradicting \eqref{cdp:eq:C.CA.7}. Thus every coefficient in this
block is nonpositive.

\item For a singular block choose a nonzero kernel vector $u$. Replacing it by its coordinatewise
absolute value cannot decrease its quadratic form, because
the off-diagonal entries are nonnegative. Both quadratic
forms must be zero by negative semidefiniteness. Hence
$|u|$ is again in the kernel. A zero coordinate of $|u|$
would, in the corresponding kernel equation, force every
neighboring coordinate to be zero. Connectedness would
then force the whole vector to vanish. Thus the block has
a strictly positive kernel vector. Pairing \eqref{cdp:eq:C.CA.7} with
this vector gives $v=0$ on the block, and then $Qc=0$.
The real divisor class $\beta=\sum c_\ell[C_\ell]$ of
that block satisfies $\beta^2=0$. The ambient negative
definite form forces $\beta=0$.

\item Discard the singular blocks and negate the coefficients
on the negative-definite blocks. This gives an effective
rational divisor $A_{\Q}$ on the curves of $M$
such that
\begin{equation}
                   c_1(K_M+A_{\Q})_{\mathbb R}=0.        \hypertarget{tag-00BM}{}\label{cdp:eq:C.CA.9}
\end{equation}
It is nonzero, since otherwise $c_1(K_M)_{\mathbb R}=0$
would contradict $K_M^2=-b_2(M)<0$.

\item Choose a positive multiple clearing the denominators in
$A_{\Q}$ and then another positive multiple killing the
torsion integral Chern class in \eqref{cdp:eq:C.CA.9}. We obtain an
integer $n>0$, a nonzero effective integral divisor $A$,
and a topologically trivial holomorphic line
\begin{equation}
                       L=K_M^{\otimes n}\otimes\OO_M(A).
                                                               \hypertarget{tag-00BD}{}\label{cdp:eq:C.CA.10}
\end{equation}
On a class VII
surface $p_g=0$, $c_2(M)=b_2(M)$ and $K_M^2=-b_2(M)$,
so Noether's formula gives $\chi(\OO_M)=0$ and
$h^1(\OO_M)=1$. Also $\dim_{\C}H^1(M,\C)=1$.
The natural map $H^1(M,\R)\to H^1(M,\OO_M)$ is
injective: if a real closed one-form $\alpha$ has
$\alpha^{0,1}=\bar\partial(p+\mathrm i q)$, with $p,q$
real, then
$\alpha=dp+\mathrm i(\bar\partial q-\partial q)$ and
closedness gives $\partial\bar\partial q=0$. The
pluriharmonic maximum principle on the compact connected
surface makes $q$ constant, so $\alpha$ is exact.
Consequently the complex-linear map
$H^1(M,\C)\to H^1(M,\OO_M)$ is an isomorphism.
The exponential sequences for constant and holomorphic
coefficients therefore give the surjection
\begin{equation}
             H^1(M,\C^*)\longrightarrow\Pic^0(M).       \hypertarget{tag-00BE}{}\label{cdp:eq:C.CA.11}
\end{equation}
This is the flat-character map, taking constant
transition functions to the same holomorphic transition
functions. It supplies a complex-flat realization of $L$.

\item The canonical section of $\OO_M(A)$ now becomes a
nonzero section of $K_M^{-n}\otimes L$. In \cite[Definition~2.20, arXiv v3 p.~14]{Dloussky2006},
take $m=n$, $D_m=A$ and the flat line $F=L^{-1}$.
Then $nK_M+F+[A]=0$ is precisely that source's
numerically anticanonical condition. Its
\cite[Theorem~0.2, arXiv v3 p.~2]{Dloussky2006}
therefore gives a global spherical shell.
In particular \eqref{cdp:eq:C.CA.6} and
$K_M^2=-b_2(M)$ force at least one of the displayed
curve classes to be nonzero in real cohomology.

\end{enumerate}
\end{proof}

\subsection{Class VII quotients and fibration descent}\label{subsec:II:component-cover-descent}

\begin{proposition}[Every class VII normalization component has Hopf minimal model]\hypertarget{tag-01CW}{}\label{cdp:st:C.2.UN.19}
\StatementRecord{cdp:st:C.2.UN.19}{proved}{01CW}{\SeriesRef{cdp:eq:C.HQ.2}, \SeriesRef{cdp:st:C.2.UN.17}, \SeriesRef{cdp:st:C.2.UN.16}, \SeriesRef{cdp:eq:C.HQ.3}, \SeriesRef{cdp:eq:C.HQ.4}, \SeriesRef{cdp:eq:C.HQ.6}, \SeriesRef{cdp:eq:C.HQ.7}, \SeriesRef{cdp:eq:C.HQ.8}, \SeriesRef{cdp:eq:C.HQ.10}; \cite[\S1, Universal and canonical desingularization, arXiv v2 p.~11]{BierstoneMilman1997}; \cite[p.~256]{DOT2000}; \cite[\S1.1, p.~2]{Nakamura2008VII}; \cite[Theorem~2.1, arXiv v2 p.~5]{Khaled2024Hopf}; \cite[Remark~2.7, arXiv v2 p.~9]{DlousskyTeleman2012}}

For every proper holomorphic degeneration of smooth complex
two-tori with smooth three-dimensional total space, each
central component whose normalization resolution is of class
VII has a Hopf surface as its minimal model.  In particular no Inoue surface
and no minimal class VII surface with positive second Betti
number is a normalization minimal model of such a component.

The comparison proving this assertion is the cyclic cover
obtained from $t=s^{m_i}$ for that component. Its stabilizer
acts on the minimal model of the covering component, and
the quotient is bimeromorphic to the original normalization.
\end{proposition}


\begin{proof}

\begin{enumerate}[label=\arabic*.,leftmargin=*]
\item Apply \cref{lem:II:cyclic-component-cover} to the chosen component $S_i$ of multiplicity $m_i$, using an equivariant resolution of $t=s^{m_i}$. Write $E$ for its strict component and $G$ for its stabilizer. The root map \eqref{cdp:eq:C.SC.2} and its proper quotient \eqref{cdp:eq:C.HQ.2} give the comparison with $S_i^\nu$.

\item Write $Z_i$ for a smooth normalization resolution and
$M_i$ for its minimal model. The case $b_2(M_i)=0$
already has Hopf minimal model by \ref{cdp:st:C.2.UN.17}. Every
minimal class VII surface with $b_2>0$ has algebraic
dimension zero \cite[p.~256]{DOT2000}. It therefore
suffices to treat $a(M_i)=0$, by an argument which
also covers the zero-Betti case. The function-field comparison in
\cref{lem:II:cyclic-component-cover}, and birational pullback
through $Z_i$ and $M_i$, give $a(E)=a(M_i)=0$.
Now \cref{cdp:st:C.2.UN.16} says that the unique smooth
minimal model $E_0$ is a torus or a Hopf surface.

\item The torus and Hopf alternatives are nonruled. Apply the minimal-contraction clause of \cref{lem:II:cyclic-component-cover} to $\rho:E\to E_0$. The equivariant contraction gives \eqref{cdp:eq:C.HQ.3}; its graph with \eqref{cdp:eq:C.HQ.2} compares $Z_i$ with a resolution of $E_0/G$ by proper bimeromorphic maps.

\item We exclude the torus alternative by its first
cohomology. For a finite holomorphic quotient
$q:A\to A/G$, proper base change gives
$R^jq_*\mathbb Q=0$ for $j>0$, and averaging gives
$(q_*\mathbb Q)^G=\mathbb Q$. Taking invariants is
exact over $\mathbb Q$, so the pullback map
identifies
\begin{equation}
                     H^1(A/G,\mathbb Q)
                          \simeq H^1(A,\mathbb Q)^G.
 \hypertarget{tag-00MU}{}\label{cdp:eq:C.HQ.4}
\end{equation}
For a complex torus the holomorphic action preserves
$H^{1,0}\oplus H^{0,1}$ and conjugation interchanges
the two invariant summands. Hence the dimension in
\eqref{cdp:eq:C.HQ.4} is even. Each singularity of a cyclic surface
quotient has a resolution with a connected chain
of rational curves: after linearization, the
quotient cone is the first quadrant in
$\mathbb Z^2+\frac1n(a,b)\mathbb Z$, and a regular
fan subdivision gives that resolution. Removing
pseudoreflections first only replaces the local
coordinates by their invariant powers. Thus every
exceptional fibre has $H^1=0$ and is connected.
Proper base change and the degree-one part of
Leray show that resolving $A/G$ leaves $H^1$
unchanged. The point-blowup maps in
\cref{lem:II:blowup-topology}, applied to a common
resolution of two smooth compact surfaces, preserve
$H^1$ under their bimeromorphic comparison. The comparisons
\eqref{cdp:eq:C.HQ.2}--\eqref{cdp:eq:C.HQ.3} would therefore make $b_1(Z_i)$
even, contradicting $b_1(Z_i)=1$.

\item It remains to prove that the cyclic quotient of
the Hopf surface $H=E_0$ is again Hopf. 
Let
\begin{equation}
            H=(\mathbb C^2\setminus\{0\})/\Gamma.
 \hypertarget{tag-00MV}{}\label{cdp:eq:C.HQ.5}
\end{equation}
The deck group contains an infinite cyclic
subgroup of finite index generated by a
holomorphic contraction. Equivalently $H$ has
a finite unramified primary Hopf cover; the
normal forms and the finite-cover description
are recalled in \cite[\S1.1, p.~2]{Nakamura2008VII}
and \cite[Theorem~2.1, arXiv v2 p.~5]{Khaled2024Hopf}.
Every deck transformation extends across the
origin by Hartogs, as does its inverse. These
extensions fix the origin and are biholomorphic.
The same is true for any lift $\widetilde g$
of a holomorphic automorphism $g$ of $H$.

\item Define the real character
\begin{equation}
 \delta:\Gamma\longrightarrow\mathbb R,
 \qquad \delta(\gamma)=\log|\det(d\gamma)_0|.
 \hypertarget{tag-00MW}{}\label{cdp:eq:C.HQ.6}
\end{equation}
It is nonzero, since its value on a contraction
is negative. For any automorphism lift,
conjugacy of derivatives gives
\begin{equation}
 \delta(\widetilde g\gamma\widetilde g^{-1})
                         =\delta(\gamma).
 \hypertarget{tag-00MX}{}\label{cdp:eq:C.HQ.7}
\end{equation}
A primary Hopf cover has rational Betti numbers
$(1,1,0,1,1)$: this follows directly from the
covering spectral sequence for
$\mathbb C^2\setminus\{0\}\simeq S^3$ and the
infinite cyclic deck group, whose action on
$H^3(S^3,\mathbb Q)$ preserves orientation.
Transfer from the finite primary cover gives
$b_1(H)\le1$ and $b_2(H)=0$, while \eqref{cdp:eq:C.HQ.6}
gives $b_1(H)\ge1$. Hence $b_1(H)=1$.
Equation \eqref{cdp:eq:C.HQ.7} says that every holomorphic
automorphism acts identically on $H^1(H,\mathbb R)$.
Poincar\'e duality and its complex orientation
then give the identity action on $H^3$ as well.
Its Lefschetz number is therefore
\begin{equation}
                 L(g)=1-1+0-1+1=0.
 \hypertarget{tag-00MY}{}\label{cdp:eq:C.HQ.8}
\end{equation}

\item For a nonidentity finite-order automorphism,
the fixed locus is a disjoint union of smooth
compact curves and isolated points, by finite
holomorphic linearization. Each fixed curve $C$
is smooth. Since $H^2(H,\mathbb Q)=0$, one has
$C^2=K_H\cdot C=0$. The adjunction
isomorphism $K_C=(K_H\otimes\mathcal O_H(C))|_C$
therefore gives $2g(C)-2=0$. Thus all the
curve components of the fixed locus have Euler
number zero. For a finite-order
map the topological Lefschetz number equals
the Euler number of its fixed locus. Indeed,
use an equivariant triangulation and barycentric
subdivision; each invariant simplex is then
fixed pointwise, so the alternating trace
on simplicial cochains counts exactly the
fixed simplices. Equation \eqref{cdp:eq:C.HQ.8} consequently
shows that there are no isolated fixed points.

\item At any point of $H$, the stabilizer in the
effective cyclic group $G$ is cyclic. Its
generator has a fixed curve through that
point and, in linearizing coordinates, acts
by
\begin{equation}
                (x,y)\longmapsto(x,\zeta y),
 \qquad
 \mathbb C\{x,y\}^{\langle\zeta\rangle}
                    =\mathbb C\{x,y^{|G_x|}\}.
 \hypertarget{tag-00MZ}{}\label{cdp:eq:C.HQ.9}
\end{equation}
Here $G_x$ denotes the stabilizer, and
$\zeta$ has its full order; a finite-order
germ with identity derivative is the identity.
This proves that the quotient $Q=H/G$ is
smooth, with the displayed local quotient
map even at its branch curves.

\item The invariant-cohomology comparison in
\eqref{cdp:eq:C.HQ.4}, now in every degree and combined
with \eqref{cdp:eq:C.HQ.7}, gives
\begin{equation}
                         b_1(Q)=1,
                         \qquad b_2(Q)=0.
 \hypertarget{tag-00MR}{}\label{cdp:eq:C.HQ.10}
\end{equation}
The finite map $H\to Q$ injects every space
of pluricanonical sections by the Jacobian pullback. Since a Hopf surface
has $\kappa=-\infty$, this proves
$\kappa(Q)=-\infty$. Also $Q$ is minimal:
an exceptional curve of self-intersection
$-1$ would have a nonzero real $H^2$ class,
contrary to \eqref{cdp:eq:C.HQ.10}.

\item There is a curve on $Q$. Explicitly, the
primary Hopf normal form
$(z_1,z_2)\mapsto(\alpha z_1+a z_2^r,\beta z_2)$
contains the elliptic curve given by $z_2=0$;
its finite image in $H$, and then in $Q$,
is a nonempty curve. A positive multiple
of that curve is homologically trivial by
\eqref{cdp:eq:C.HQ.10}. The divisor classification
\cite[Remark~2.7, arXiv v2 p.~9]{DlousskyTeleman2012} applied
to the minimal class VII surface $Q$ now
gives a Hopf surface: its Enoki alternative
has positive second Betti number and is
excluded by \eqref{cdp:eq:C.HQ.10}. 

\item Finally \eqref{cdp:eq:C.HQ.2}--\eqref{cdp:eq:C.HQ.3} make $Z_i$
bimeromorphic to this smooth Hopf surface
$Q$. Uniqueness of the class VII minimal
model proves that $M_i\simeq Q$. 

\end{enumerate}
\end{proof}

\begin{proposition}[Descent of the rank-one component fibration]\hypertarget{tag-01CY}{}\label{cdp:st:C.2.UN.20}
\StatementRecord{cdp:st:C.2.UN.20}{proved}{01CY}{\SeriesRef{cdp:st:C.2.UN.16}, \cref{lem:II:cyclic-component-cover}, \SeriesRef{cdp:eq:C.FD.6}}

In the setting and with the component cover of \ref{cdp:st:C.2.UN.16},
the generic-rank-one branch gives a fibration on a resolution of
the original normalization. More precisely, for the stabilizer
$G$ of the chosen covering component $E$, there are proper
birational maps and quotient maps in a commutative diagram
\begin{equation}
\begin{array}{ccc}
 \widehat E&\xrightarrow{\ q\ }&B\\
 \downarrow&&\downarrow\\
 \widehat E/G&\longrightarrow&B/G,
\end{array}                                                    \hypertarget{tag-00K2}{}\label{cdp:eq:C.FD.1}
\end{equation}
where $B$ is a smooth compact curve, $q$ has connected general
fibre, and $\widehat E/G$ is bimeromorphic to $S_i^\nu$.
Resolving this last quotient and taking its Stein factor gives
a proper fibration
\begin{equation}
        Z'_i\longrightarrow B/G
        \quad\hbox{with connected rational or elliptic general fibre}.
                                                               \hypertarget{tag-00K3}{}\label{cdp:eq:C.FD.2}
\end{equation}
The finite map from the Stein base of \eqref{cdp:eq:C.FD.2}
to $B/G$ is an isomorphism. For every original component with nonnegative
Kodaira dimension, its general fibre in \eqref{cdp:eq:C.FD.2} is elliptic.
The same conclusion holds for every original component that
is not uniruled.
\end{proposition}


\begin{proof}

\begin{enumerate}[label=\arabic*.,leftmargin=*]
\item Choose the equivariant resolution of \cref{lem:II:cyclic-component-cover}. Its stabilizer $G$ preserves $E$. Use the modules $M$ and $L$ of \eqref{cdp:eq:C.CF.3}, identifying $\mathcal O$ with $\C\{s\}$ and its fraction field with $\C(\{s\})$.
An element $g\in G$ covers $s\mapsto\zeta_gs$. Pushforward of
vector fields is therefore a semilinear automorphism of $M$
and of its meromorphic localization. Since $g(E)=E$, it
preserves the valuation defining $L$. It also preserves $sL$,
because $s$ changes by a nonzero constant. Reduction gives
a complex-linear representation on $L/sL$, and the restriction
map of \eqref{cdp:eq:C.CF.6} is equivariant:
\begin{equation}
 g_*\bigl(v|_E\bigr)=(g_*v)|_E.
                                                               \hypertarget{tag-00K5}{}\label{cdp:eq:C.FD.4}
\end{equation}
Consequently the two-dimensional space $\mathfrak a_E$ and
its rank-one meromorphic tangent distribution are preserved
by $G$.

\item For a basis $V_1,V_2$ write $V_2=hV_1$. The function $h$ is
nonconstant, and both fields annihilate it. If the matrix of
$g_*$ in this basis is
\[
       g_*V_1=aV_1+bV_2,\qquad
       g_*V_2=cV_1+dV_2,
\]
then the coefficient comparison gives
\begin{equation}
              h\circ g^{-1}=\frac{c+dh}{a+bh}.
                                                               \hypertarget{tag-00K6}{}\label{cdp:eq:C.FD.5}
\end{equation}
This identity holds where the displayed denominator is nonzero
and hence as an identity of meromorphic functions. It supplies
the finite projective action on the target of
$h:E\dashrightarrow\mathbb P^1$. Resolve its graph
equivariantly, and then take the Stein factor. Uniqueness of
the Stein factor lifts the action on $\mathbb P^1$ to an action
on the smooth compact curve $B$, and gives an equivariant map
$q:\widehat E\to B$. This is the same connected-leaf
fibration as in \ref{cdp:st:C.2.UN.16}. Its general fibre is rational or
elliptic, by the semicomplete one-dimensional field calculation
there. Taking the finite quotient produces \eqref{cdp:eq:C.FD.1}.

\item Let
$K\subset G$ be the kernel of the action on $B$. Outside
the finitely many fixed points of the nonidentity elements
of $G/K$, a point $b\in B$ has stabilizer exactly $K$.
The inverse image of its orbit in $B/G$ is the disjoint union
of the fibres $q^{-1}(gb)$, and $G$ acts transitively on these
fibres. Hence the quotient fibre is exactly
\begin{equation}
                       q^{-1}(b)/K.
                                                               \hypertarget{tag-00K7}{}\label{cdp:eq:C.FD.6}
\end{equation}
It is connected. A finite quotient of a smooth compact curve
is smooth, as follows locally from the invariant coordinate
$z\mapsto z^r$ at a cyclic stabilizer. Riemann--Hurwitz
shows that the genus of \eqref{cdp:eq:C.FD.6} is at most the genus of
$q^{-1}(b)$, and is therefore zero or one.

\item The generic fibre remains unchanged when the
surface quotient is resolved. A stabilizer of a point in a
general fibre is contained in $K$. It fixes the base coordinate
and, after finite holomorphic linearization, acts by
$(x,y)\mapsto(x,\xi y)$. Its invariant ring is
$\mathbb C\{x,y^r\}$, so the surface quotient is smooth
along a general fibre. A resolution can be an isomorphism
there. The resulting map to $B/G$ has connected general
fibre by \eqref{cdp:eq:C.FD.6}. Its Stein factor is a finite map of normal
compact curves of degree one, hence an isomorphism. This
proves the assertion about the base in \eqref{cdp:eq:C.FD.2}.

\item The proper quotient map \eqref{cdp:eq:C.HQ.2} of
\cref{lem:II:cyclic-component-cover}, composed with the
equivariant graph resolution $\widehat E\to E$, identifies
$\widehat E/G$ bimeromorphically with $S_i^\nu$.
Resolving their graph gives the proper birational comparison
of \eqref{cdp:eq:C.FD.2} with a smooth normalization resolution of the
original component. 

\item Finally a rational general fibre $F$ of \eqref{cdp:eq:C.FD.2} has trivial
normal line in $Z'_i$, and adjunction gives
$K_{Z'_i}|_F=K_F=\mathcal O_{\mathbb P^1}(-2)$.
Every positive pluricanonical section restricts to zero on
every general fibre and therefore vanishes identically.
This makes $\kappa(Z'_i)=-\infty$, so the nonnegative
Kodaira-dimension cases have elliptic general fibre.
For the uniruledness conclusion, the relative anticanonical
evaluation maps a general rational fibre by its degree-two
linear system to a plane conic. The torsion-free direct
image on the compact base curve is a rank-three vector
bundle, and this evaluation gives a meromorphic map of
$Z'_i$ to its projective bundle, birational onto the
closure of the conic family. That projective bundle is
projective because the base is a compact curve. A general
hyperplane section of the conic surface has a component
finite over the base; normalize it to a compact curve
$C'$. On the resulting generic conic over $\mathcal M(C')$,
the diagonal point gives projection to $\PP^1$, a
birational parameterization. Composing its inverse with
the finite base-change projection and the preceding
birational map produces a dominant meromorphic map
$C'\times\PP^1\dashrightarrow Z'_i$.
\Cref{def:II:marked-torus-fillings} and its
birational-invariance proof therefore make the original
component uniruled. This proves the final assertion.

\end{enumerate}
\end{proof}

\subsection{Rank-two frames and elliptic fibrations}\label{subsec:II:rank-two-component-frames}

\begin{definition}[Root curves and normalized fibre charts]\hypertarget{tag-P344}{}\label{def:II:elliptic-root-curve}
\StatementRecord{def:II:elliptic-root-curve}{proved}{P344}{}
For a smooth compact connected curve $B$, distinct points $p_i$ and integers
$m_i>1$, the root curve
\begin{equation}
                    \mathcal B=B[(p_i,m_i)_i]
 \hypertarget{tag-010R}{}\label{cdp:eq:C.SF.1}
\end{equation}
has coarse curve $B$, is equal to $B$ away from the marked points,
and has chart $[\Delta_s/\boldsymbol\mu_{m_i}]$ over $p_i$, with
coarse coordinate $t=s^{m_i}$ and action $s\mapsto\zeta_{m_i}s$.
Its root line $\mathcal R_i$ and section have $m_i$th powers
equal to the pullback of $\OO_B(p_i)$ and its canonical section.
These charts, lines and sections glue through their root-frame
identifications. Automorphisms of $B$ preserving the marked
orders act on this construction by pullback of the same pairs.

Let $q:S\to B$ be a proper surjective holomorphic map with connected fibres
from a smooth compact complex surface. At a marked point suppose
that the full fibre is $m_iD_i$ with $D_i$ reduced. The pairs
$(\OO_S(D_i),s_{D_i})$ define its normalized root lift to
$\mathcal B$. On the root disc at a full fibre $mD$, this lift is
the normalization $V$ of the base change $t=s^m$. The same formulas
for $m=1$ give the identity base change. It is finite
\'etale over the original surface tube $U$, with quotient
$V/\boldsymbol\mu_m=U$; its fibres over the root disc are connected
and its central fibre is reduced.

When $D=C$ is a smooth elliptic curve, $V\to\Delta_s$ is a smooth
elliptic family after shrinking the disc. Its central fibre $C'$
is connected, and the finite quotient $C'\to C$ is a free cyclic
translation quotient of degree $m$.
\end{definition}
\begin{proof}
A change of local frame of $\OO_B(p_i)$ changes its section and
root coordinate by the same unit root. Two choices of that root
differ by the displayed $\boldsymbol\mu_{m_i}$ action. This gives
the stated chart gluing, root lines and section-preserving
pullback action. The equality of full Cartier divisors
$q^*(p_i)=m_iD_i$ identifies
$(\OO_S(D_i),s_{D_i})^{m_i}$ with the pulled-back base pair,
so defines the lift with these same transition functions.

Locally on $S$ write $t=u h^m$, where $u$ is a unit and $h=0$
defines $D$. The normalized base change is
\begin{equation}
                         w^m=u,\qquad s=wh.
 \hypertarget{tag-00FV}{}\label{cdp:eq:C.EH.6}
\end{equation}
The first equation is finite \'etale over the smooth surface,
since $w$ is a unit and $m w^{m-1}\ne0$. On a small simply
connected surface chart it is the disjoint union of the $m$
holomorphic roots of $u$. The map to $s^m=u h^m$ is finite and
is an isomorphism over $h\ne0$, with inverse $w=s/h$ there.
Its source is normal, so it is precisely the normalization.
The generator sends $(w,s)$ to $(\zeta_m w,\zeta_m s)$ and
leaves the two surface coordinates fixed. Its invariant quotient
is $U$, and every nonidentity power acts freely because $w\ne0$.

The proper root family has connected generic fibre, the same
fibre as $q$ at $t=s^m\ne0$. Its normal Stein factor is finite
and generically degree one over the normal $s$-disc, hence is that
disc. All its fibres are connected. Since $h$ is reduced and the
cover is \'etale, $s=wh$ defines a reduced central divisor.

For smooth $C$, its pullback $C'$ is smooth and $s=wh$ has
nonzero transverse derivative along it. Properness permits
shrinking the root disc so that the family is smooth throughout.
The finite \'etale map from the connected compact curve $C'$ to
$C$ makes $C'$ elliptic by Riemann--Hurwitz. Its quotient is
\begin{equation}
                         C'\longrightarrow C'/\boldsymbol\mu_m=C.
 \hypertarget{tag-00FW}{}\label{cdp:eq:C.EH.7}
\end{equation}
Choose an origin on $C'$. An automorphism has form $z\mapsto Lz+a$.
For $L\ne1$, the nonzero endomorphism $1-L$ is surjective, so
$(1-L)z=a$ gives a fixed point. Freeness makes each deck
automorphism a translation. The generator has exact order $m$
by the explicit free action on $w$.
\end{proof}


\begin{lemma}[Elliptic fibrations with identity monodromy]\hypertarget{tag-P306}{}\label{lem:II:identity-monodromy-elliptic}
\StatementRecord{lem:II:identity-monodromy-elliptic}{proved}{P306}{\cite[Theorem~6.2, pp.~564--566; p.~567; \S8(iv)(11), p.~585]{Kodaira1963II}; \cite[\S2.3, arXiv v2 pp.~5--6]{UlmerUrzua2021}}
Let $q:E_0\to B$ be a proper holomorphic map from a smooth compact complex surface to a smooth compact connected curve, with connected fibres and elliptic general fibre. Suppose that $q$ is relatively minimal and that its integral $H_1$ monodromy on the smooth base locus is the identity. Every full fibre is a positive multiple of a smooth elliptic curve. The smooth fibres are torsors under one fixed elliptic curve $F$, and
\[
 \begin{gathered}
 R^1q_*\OO_{E_0}\simeq\OO_B,\qquad
 p_g(E_0)=g(B),\qquad h^1(E_0,\OO_{E_0})=g(B)+1,\\
 \chi(\OO_{E_0})=e(E_0)=0.
 \end{gathered}
\]
At a fibre $mC$, the normalization of the base change $t=s^m$ over a small base disc is a smooth elliptic family and is finite \'etale over the original surface neighbourhood. Its cyclic action on the central elliptic fibre is by translations, and its action on the relative $H^1(\OO)$ line is trivial. The displayed direct-image isomorphism is obtained by taking invariants under this action at every multiple value.


More precisely, the tube of $mC$ is isomorphic over its base disc to
\begin{equation}
(F\times\Delta_s)/\mu_m,\qquad
 (z,s)\longmapsto(z+\tau,\zeta_m s),\qquad
 q=s^m,\qquad\operatorname{ord}(\tau)=m.
\hypertarget{tag-00NE}{}
\hypertarget{tag-00WI}{}
\hypertarget{tag-00IA}{}
\label{cdp:eq:C.HT.4}
\SsixEquationAliases{\label{cdp:eq:C.QT.2}\label{cdp:eq:C.EP.8}}
\end{equation}
The quotient map is finite \'etale, its central cover
$\pi_C:F\to C$ is connected, and the normal line $N_{C/E_0}$
has exact order $m$. The connected cover of unit frames of its
$m$th-power trivialization is this same $F\to C$; replacing the
normal by its dual reverses the cyclic character.
Restriction followed by $\pi_C^*$ is evaluation of the fixed
$H^1(F,\OO_F)$ frame of $R^1q_*\OO_{E_0}$. Thus
\[
 0\longrightarrow H^1(B,\OO_B)\xrightarrow{q^*}
 H^1(E_0,\OO_{E_0})
 \xrightarrow{\ \pi_C^*\circ\operatorname{res}_C\ }
 H^1(F,\OO_F)\longrightarrow0
\]
is exact. For $B=\PP^1$ this restriction composite is an
isomorphism, and
\begin{equation}
q_*\OO_{E_0}=\OO_{\PP^1},\qquad
 R^1q_*\OO_{E_0}=\OO_{\PP^1},\qquad
 (h^0,h^1,h^2)(E_0,\OO_{E_0})=(1,1,0).
\hypertarget{tag-00WJ}{}
\label{cdp:eq:C.QT.3}
\end{equation}

Let $m_bC_b$ denote the full fibres with $m_b>1$. Then
\begin{equation}
K_{E_0}\simeq q^*K_B\otimes
 \OO_{E_0}\Bigl(\sum_b(m_b-1)C_b\Bigr),
 \qquad q_*K_{E_0}\simeq K_B.
\label{eq:II:quasibundle-canonical}
\end{equation}
Fix a nonzero invariant one-form $\alpha_F$ on $F$.
For every meromorphic differential $\theta$ on $B$, the local
expressions
\begin{equation}
q^*\theta\wedge\alpha_F,
\hypertarget{tag-00NF}{}
\label{cdp:eq:C.HT.5}
\end{equation}
glue to a meromorphic two-form on $E_0$. Every meromorphic
two-form with vertical divisor has this form for a unique
$\theta$, and its order is
\begin{equation}
\operatorname{ord}_{C_b}(q^*\theta\wedge\alpha_F)
       =m_b\operatorname{ord}_b(\theta)+m_b-1.
\hypertarget{tag-00NG}{}
\hypertarget{tag-00WM}{}
\label{cdp:eq:C.HT.6}
\SsixEquationAliases{\label{cdp:eq:C.QT.6}}
\end{equation}

\end{lemma}
\begin{proof}
\begin{enumerate}[label=\arabic*.,leftmargin=*]
\item Let $B^\circ$ be the smooth base locus. Trivial integral monodromy gives a global oriented marking $(\gamma_1,\gamma_2)$ of the first homology local system. On a small disc in $B^\circ$, choose a local section and a nonzero holomorphic relative one-form. Its periods $\lambda_1,\lambda_2$ are holomorphic, and $\lambda_1$ is nowhere zero because a nonzero integral cycle on an elliptic curve has nonzero period. The ratio $\tau=\lambda_2/\lambda_1$ is independent of rescaling that form. The marking therefore makes it a single-valued holomorphic map $B^\circ\to\mathbb H$.

The Cayley transform $(\tau-\mathrm i)/(\tau+\mathrm i)$ is bounded. It extends across the finitely many omitted points and is constant on the compact curve $B$, so $\tau$ is constant. Set $F=\mathbb C/(\mathbb Z+\tau\mathbb Z)$. On each small smooth base disc, integration of the relative form normalized by its $\gamma_1$ period, from the chosen section to $x$, identifies the fibre with $F$. These fibrewise integration maps give a holomorphic product with $F$. An overlap preserves the integral marking, so its origin-preserving part is the identity; the overlap is a translation. This identifies the smooth fibres as $F$-torsors and gives a global constant generator of their $H^1(\OO)$ line.

\item Apply Kodaira's classification of relatively minimal proper analytic elliptic fibrations \cite[Theorem~6.2, pp.~564--566; p.~567]{Kodaira1963II}. Every multiplicity-one fibre in the list has a coefficient-one component. At a smooth point of that component away from the other components, the map is a submersion and has a local holomorphic section $\sigma$. On each nearby smooth fibre $F_b$, the map
$x\mapsto\OO_{F_b}(x-\sigma(b))$
identifies $F_b$ with $\operatorname{Pic}^0(F_b)$ and identifies their integral homology local systems. Thus the positive-loop monodromy is the one in \cite[\S2.3, arXiv v2 pp.~5--6]{UlmerUrzua2021}. Identity monodromy leaves only the smooth type $I_0$.

For a multiple fibre, the classification gives type $mI_b$.
Its full divisor is $mD$, with $D$ reduced Cartier. Apply the
general normalized-root construction of
\cref{def:II:elliptic-root-curve} to this same surface tube and
base coordinate $t$. Locally $t=u h^m$, and its maps are
$w^m=u$, $s=wh$, with the original surface coordinates fixed.
The construction gives the finite \'etale cover $\pi:E'\to E_0$
over the tube, the reduced central fibre and the connected
Stein factor over $\Delta_s$. Its punctured monodromy is the
$m$th power of the original monodromy, hence the identity.

Every central prime curve downstairs has canonical degree zero in Kodaira's list. The \'etale equality $K_{E'}=\pi^*K_{E_0}$ gives canonical degree zero on each central prime upstairs. An exceptional rational curve of square $-1$ would have canonical degree $-1$ by adjunction, so the root family is relatively minimal. Its reduced central fibre has coefficient-one components. The preceding section and Abel--Jacobi comparison therefore applies to it; identity monodromy makes that central fibre smooth of type $I_0$, as also in \cite[\S8(iv)(11), p.~585]{Kodaira1963II}. Shrinking the disc makes the root family smooth throughout. Since $\pi$ is \'etale, its reduced image $C$ is smooth elliptic.

\item Write $r:\Delta_s\to\Delta_t$ for $t=s^m$ and $q':E'\to\Delta_s$ for this smooth root family. Its marked integration maps identify it locally with the product by the same elliptic curve $F$. The smooth-elliptic clause of \cref{def:II:elliptic-root-curve}, now applied to the support $C$ proved smooth above, identifies this same central deck action as a free translation action. Its origin-preserving part on neighbouring marked fibres varies in the discrete group $\operatorname{Aut}(F,0)$, so it remains the identity. It consequently acts trivially on the relative $H^1(\OO)$ generator.

Choose a local section of the smooth root family near one point
of its central fibre. The marked integration maps then give a
product $F\times\Delta_s$ after shrinking the disc. Write the
cyclic map as $(z,s)\mapsto(z+\tau+A(s),\zeta_ms)$, where
$\tau$ is its central translation and $A$ is a small holomorphic
complex lift with $A(0)=0$. The order-$m$ law gives
$\sum_{j=0}^{m-1}A(\zeta_m^js)=0$: the sum is lattice-valued
and vanishes at zero. Set
\[
 b(s)=-\frac1m\sum_{j=0}^{m-1}jA(\zeta_m^js),\qquad
 z'=z+b(s),\qquad z=z'-b(s).
\]
Subtracting the shifted sums gives
$b(\zeta_ms)-b(s)=-A(s)$. These inverse coordinate maps
give \eqref{cdp:eq:C.HT.4}. The central action is free, so
$\tau$ has exact order $m$. The conormal frame upstairs is
$ds$, with character $\zeta_m$. If its $a$th tensor power
had a nowhere-zero section downstairs, its pullback would
be $h(ds)^a$ with $h$ holomorphic on the connected compact
$F$. Invariance gives $h(z+\tau)\zeta_m^a=h(z)$.
The function $h$ is a nonzero constant, so $m\mid a$.
The full fibre trivializes the $m$th power, proving the
exact normal order. The unit-frame cover has precisely this
primitive character and is the same connected cyclic cover;
duality inverts that character.


The finite quotient map $\pi:E'\to q^{-1}(\Delta_t)$ gives $\OO_{q^{-1}(\Delta_t)}=(\pi_*\OO_{E'})^{\mu_m}$. Finite direct image has no higher coherent direct images, and averaging over $\mu_m$ makes invariants exact. The direct-image comparison is therefore
\[
 \begin{aligned}
 \left.R^1q_*\OO_{E_0}\right|_{\Delta_t}
 &\simeq\left(r_*R^1q'_*\OO_{E'}\right)^{\mu_m}\\
 &\simeq\left(r_*\OO_{\Delta_s}\right)^{\mu_m}
 \simeq\OO_{\Delta_t}.
 \end{aligned}
\]
The last map identifies a holomorphic function of $t$ with its pullback under $t=s^m$: a power series in $s$ invariant under $s\mapsto\zeta_m s$ has exactly the powers divisible by $m$. These local generators agree on punctured discs with the translation-invariant generator from the global marking. They thus give $R^1q_*\OO_{E_0}\simeq\OO_B$ over the full base, including every multiple value.

\item Connected fibres give $q_*\OO_{E_0}=\OO_B$. The coherent Leray sequence on the curve gives
\begin{equation}
 H^2(E_0,\OO_{E_0})\simeq H^1(B,\OO_B),\qquad
 h^1(E_0,\OO_{E_0})=g(B)+1.
                                                               \hypertarget{tag-00Z7}{}\label{cdp:eq:C.RG.13}
\end{equation}
Serre duality yields $p_g(E_0)=g(B)$, and hence $\chi(\OO_{E_0})=1-(g(B)+1)+g(B)=0$. Euler additivity over the smooth base locus and the finitely many remaining fibres gives $e(E_0)=0$, since each fibre support is an elliptic curve.


The square formed by the root tube, its central inclusion
$F\hookrightarrow E'$, the quotient tube and $C\hookrightarrow E_0$
commutes. On coherent cohomology its upper restriction evaluates
the constant $H^1(F,\OO_F)$ frame. Finite invariants and this
restriction commute, by the same averaging operator. Therefore
restriction on the quotient followed by $\pi_C^*$ equals
evaluation of the frame used above. This proves the stated
restriction map, including at a multiple fibre. The Leray
edge sequence is exact with kernel $q^*H^1(B,\OO_B)$ and
quotient $H^0(B,R^1q_*\OO_{E_0})=H^1(F,\OO_F)$.
For $B=\PP^1$ its left term is zero; together with
\eqref{cdp:eq:C.RG.13} it gives \eqref{cdp:eq:C.QT.3}.

\item Normalize the relative invariant form by the fixed marking.
Translations between product charts change an absolute lift of
$\alpha_F$ only by a base one-form. Wedge with $q^*\theta$
therefore gives the same two-form on overlaps. In the root
chart $t=s^m$, the pullback of $dt\wedge\alpha_F$ is
$m s^{m-1}ds\wedge\alpha_F$. Its divisor there has order
$m-1$ on $s=0$. The finite \'etale map preserves that order,
giving the first canonical isomorphism in
\eqref{eq:II:quasibundle-canonical} and the full formula
\eqref{cdp:eq:C.HT.6}.

A holomorphic two-form upstairs is
$a(s)ds\wedge\alpha_F$, because its coefficient is constant
on every compact elliptic fibre. Root invariance is
$a(\zeta_ms)\zeta_m=a(s)$, so
$a(s)=s^{m-1}c(s^m)$ with $c$ holomorphic. Its descent is
the form corresponding to $m^{-1}c(t)dt$ on the base.
This proves $q_*K_{E_0}=K_B$ across every multiple value.

Choose a nonzero meromorphic base differential to form one
reference two-form. The ratio of any vertical-divisor
meromorphic two-form to that reference has vertical divisor.
On each regular compact fibre outside that divisor it is a
nonzero holomorphic function and hence constant. It descends
to the regular base. On the root tube its Laurent expansion
is a meromorphic function of $s$ invariant under
$s\mapsto\zeta_ms$, so it is a meromorphic function of
$s^m$. This extends the descended ratio through every omitted
value and proves existence and uniqueness of $\theta$ in
\eqref{cdp:eq:C.HT.5}.
\end{enumerate}
\end{proof}

\begin{lemma}[Elliptic first integrals of commuting fields]\hypertarget{tag-P303}{}\label{lem:II:elliptic-first-integral}
\StatementRecord{lem:II:elliptic-first-integral}{proved}{P303}{\cref{lem:II:semicomplete-transport,lem:II:identity-monodromy-elliptic}; \cite[Chapter~1, Proposition~2.2, p.~8]{Peters1990}}
Let $X$ be a smooth compact complex surface and let $V,Y$ be commuting meromorphic vector fields with $V\wedge Y\ne0$, each semicomplete on its holomorphic locus. Resolve the graph of a meromorphic first integral of $V$ and take its Stein factorization
\begin{equation}
 q:\widehat X\longrightarrow B
 \hypertarget{tag-00YU}{}\label{cdp:eq:C.RF.1}
\end{equation}
with connected elliptic general fibre. There is a relatively minimal elliptic fibration $q_0:X_0\to B$ obtained from $\widehat X$ by vertical point blowdowns, with $g(B)\le1$, identity integral $H_1$ monodromy, and every fibre a positive multiple of a smooth elliptic curve. Its smooth fibres are torsors under one fixed elliptic curve. Moreover,
\[
 e(X_0)=K_{X_0}^2=\chi(\OO_{X_0})=\chi(\OO_X)=0.
\]
\end{lemma}
\begin{proof}

\begin{enumerate}[label=\arabic*.,leftmargin=*]
\item Lift the two fields to $\widehat X$. Their brackets vanish on the common open isomorphism with $X$ and hence as meromorphic fields everywhere. Their semicompleteness on the full holomorphic loci follows from \cref{lem:II:semicomplete-transport}.

\item First, $V$ has no horizontal polar curve.  On a general
smooth fibre $F$ of $q$, its restriction is a nonzero
semicomplete meromorphic field on the compact elliptic
curve.  Apply the compact-curve clause of \cref{lem:II:semicomplete-transport} to $V|_F$; it excludes poles on $F$.  Thus this restriction is holomorphic and nowhere
zero, and its flow acts by all translations of $F$.
The commutation relation implies
$(\Phi_V^a)_*Y=Y$ as an identity of meromorphic fields.
Consequently the intersection of the polar divisor of $Y$
with $F$ is invariant under all translations of $F$.
It is a finite set, so it is empty.  This proves that
$Y$ also has no horizontal polar curve.

\item The field $Y$ projects to a nonzero meromorphic field $Z$
on $B$.  In fact, for each local meromorphic function $g$
on $B$,
\[
 V\bigl(Y(q^*g)\bigr)
   =Y\bigl(V(q^*g)\bigr)=0.
\]
The meromorphic function $Y(q^*g)$ is therefore constant
on the connected general fibres; proper Stein descent
makes it the pullback of a meromorphic function on $B$.
These functions define $Z$.  Were $Z$ zero, both fields
would be tangent to the one-dimensional fibres, contrary
to $V\wedge Y\ne0$.

\item Suppose that $Z$ has a pole of order $r\ge1$ at $b$.
The primitive of its time form has a zero of order $r+1$
there.  Its inverse solution consequently has distinct
continuations after a loop about the corresponding time
value.  Choose that loop and its lifted base paths in a
small punctured disc about $b$ avoiding all other bad
values.  The inverse image of the compact union of these
paths is compact.  Both $q$ and $Y$ are holomorphic there,
and $Y$ has no pole there, since all its polar curves are
vertical and the disc was chosen accordingly.  The usual
continuation theorem for holomorphic ordinary differential
equations therefore lifts the whole time loop to a
solution of $Y$.  Its endpoint has a different image in
$B$.  This contradicts semicompleteness of $Y$.
Hence $Z$ is a nonzero holomorphic field on the compact
curve $B$, and $g(B)\le1$.

\item Remove finitely many values so that $q$ is a smooth
elliptic fibration, $V,Y$ are holomorphic, $V$ is nowhere
zero vertically, and $Z$ is nonzero.  On this open base
$B^\circ$, let $\alpha_b$ be the time form of $V|_{F_b}$.
The flow of $Y$, defined for a uniform small time along
each compact fibre, carries whole fibres biholomorphically
to whole fibres and carries $V$ to itself.  It therefore
preserves their time forms.  Integration gives the exact
horizontal map from the integral homology local system
\begin{equation}
 \mathcal H_1(q)\big|_{B^\circ}
   \longrightarrow\underline{\mathbb C},
 \qquad
 \gamma\longmapsto\int_\gamma\alpha_b.
 \hypertarget{tag-00YV}{}\label{cdp:eq:C.RF.2}
\end{equation}
The map is injective on each stalk: the periods of a
nonzero holomorphic one-form on an elliptic curve are a
rank-two lattice in $\mathbb C$.  The preceding flow
comparison proves that its values are locally constant.
Analytic continuation of any integral cycle around a
loop preserves its complex period; injectivity then
proves that every monodromy operator on $H_1(F_b,\mathbb Z)$
is the identity.  The image lattice is fixed, say
$\Lambda\subset\mathbb C$, and the fibres over
$B^\circ$ are torsors under the fixed elliptic curve
$\mathbb C/\Lambda$. Adding back the extra removed regular values preserves trivial monodromy: the inclusion of $B^\circ$ into the full smooth base locus induces a surjection on fundamental groups. The local holomorphic period maps extend across those regular values and agree with the constant period map on the punctured neighbourhoods. Thus all smooth fibres are torsors under the same elliptic curve.

\item Contract the vertical exceptional curves of the first
kind to obtain the relatively minimal elliptic
map
\begin{equation}
 \widehat X\xrightarrow{\ \mu\ }X_0
       \xrightarrow{\ q_0\ }B.
 \hypertarget{tag-00YW}{}\label{cdp:eq:C.RF.3}
\end{equation}
These contractions preserve the smooth fibres and their integral monodromy representation. Apply \cref{lem:II:identity-monodromy-elliptic} with $(E_0,q)=(X_0,q_0)$ and the same base $B$. It gives the smooth multiple fibre supports and $e(X_0)=\chi(\OO_{X_0})=0$. Write $\Lambda=\lambda_1\mathbb Z+\lambda_2\mathbb Z$ in the period marking above. The fixed curve in that lemma, normalized to first period one, is identified with $\mathbb C/\Lambda$ by $[z]\mapsto[\lambda_1z]$; this is the rescaling of the integration map on the common smooth fibres.

\item In particular, for pairwise distinct base values the
fibre supports $C_b$ are disjoint smooth elliptic
curves, and $C_b^2=0$, since some positive multiple
of $C_b$ is a full fibre.  To compute the canonical square, choose a nonzero meromorphic
one-form $\eta$ on $B$.  The meromorphic field $V$
descends through $\mu$, and at the generic point
there is a unique meromorphic two-form $\Omega$
satisfying
\begin{equation}
 \iota_V\Omega=q_0^*\eta.
 \hypertarget{tag-00YX}{}\label{cdp:eq:C.RF.4}
\end{equation}
Since $V$ is tangent to $q_0$, the
right-hand side annihilates $V$, which proves
existence and uniqueness of $\Omega$ by a local
tangent-frame calculation.  Its divisor is
vertical: on a general compact fibre $V$ is
holomorphic and nowhere zero, and $\eta$ is
holomorphic and nonzero at its base value.
Thus a canonical divisor is a finite integral
combination of the disjoint square-zero curves
$C_b$, and $K_{X_0}^2=0$.  Birational invariance under the displayed point blowups and blowdowns \cite[Chapter~1, Proposition~2.2, p.~8]{Peters1990} gives $\chi(\OO_X)=\chi(\OO_{X_0})=0$.

\end{enumerate}
\end{proof}

\begin{proposition}[A component cover with a rank-two translation frame has coherent Euler characteristic zero]\hypertarget{tag-01D0}{}\label{cdp:st:C.2.UN.22}
\StatementRecord{cdp:st:C.2.UN.22}{proved}{01D0}{\cref{cdp:st:C.2.UN.16,lem:II:elliptic-first-integral,lem:II:semicomplete-transport,lem:II:zero-dimensional-field-surface,lem:II:algebraic-leaf-fibration}; \cite[Main Theorem and the following clarification, arXiv v2 p.~2; Proposition~22, p.~25]{Guillot2019}}

Let $X$ be a smooth compact complex surface with
$\kappa(X)\ge0$. Suppose that it carries a two-dimensional
complex vector space $\mathfrak a$ of commuting meromorphic
vector fields, of generic rank two, every member of which
is semicomplete on its holomorphic locus. Then
\begin{equation}
                         \chi(\mathcal O_X)=0.
                                                               \hypertarget{tag-00YY}{}\label{cdp:eq:C.RF.5}
\end{equation}
More precisely, $X$ has a two-torus birational model,
or it has a relatively minimal elliptic fibration
on a proper birational model, over a smooth compact curve
of genus at most one, with identity integral $H_1$ monodromy
and with every fibre a positive multiple of a smooth elliptic
curve. Its smooth fibres are torsors under one fixed elliptic
curve.

This applies in particular to every covering component
$E$ of \ref{cdp:st:C.2.UN.16} with $\kappa(E)\ge0$ in its
generic-rank-two branch. Such a covering component cannot
have a K3 or Enriques minimal model.
\end{proposition}

\begin{proof}
For meromorphic fields $Z_1,Z_2$ on a connected curve with $Z_1\ne0$, write $Z_2=hZ_1$ on the open set where $Z_1$ is nonzero. Then $[Z_1,Z_2]=Z_1(h)Z_1$. A vanishing bracket makes the meromorphic function $h$ locally constant and hence constant on the connected curve. Thus
\begin{equation}
              Z_1(Z_2/Z_1)=0,
              \qquad Z_2=cZ_1\quad(c\in\mathbb C).
                                                               \hypertarget{tag-00Z1}{}\label{cdp:eq:C.RF.8}
\end{equation}

For any proper curve fibration on a smooth birational model $X'$ of $X$, a smooth rational general fibre $F$ has $F^2=0$ and $K_{X'}|_F\simeq\OO_{\mathbb P^1}(-2)$ by adjunction. Every positive pluricanonical section vanishes on every general fibre and hence on their dense union. This would give $\kappa(X')=-\infty$, contrary to $\kappa(X)\ge0$. Thus every genus-at-most-one fibration obtained below has elliptic general fibre.

\begin{enumerate}[label=\arabic*.,leftmargin=*]
\item Suppose $a(X)=0$. Apply \cref{lem:II:zero-dimensional-field-surface} to $X$ and the given two-dimensional commuting semicomplete space $\mathfrak a$. Its minimal model is a two-torus or a Hopf surface. The latter has Kodaira dimension $-\infty$, so $\kappa(X)\ge0$ leaves the torus. Its coherent Euler characteristic is $1-2+1=0$, giving \eqref{cdp:eq:C.RF.5} by birational invariance.

\item Suppose $a(X)=1$. Choose a nonconstant
$h\in\mathcal M(X)$, resolve its graph, and take the
Stein factor $q:\widehat X\to B$. Every
$g\in\mathcal M(X)$ is algebraic over $\mathbb C(h)$.
On a connected general fibre of $q$ it therefore takes
finitely many values and is constant. Proper meromorphic
descent consequently gives the exact field identity
\begin{equation}
                       \mathcal M(X)=q^*\mathcal M(B).
                                                               \hypertarget{tag-00Z2}{}\label{cdp:eq:C.RF.9}
\end{equation}
For example, this descent can be seen by mapping the
proper graph of $g$ to $B\times\mathbb P^1$; its image
is the graph of the descended meromorphic function.
Each field in $\mathfrak a$ differentiates this field
of functions and therefore induces a meromorphic vector
field on $B$. The induced fields commute. They cannot
both vanish, since that would make both original fields
tangent to the one-dimensional fibres and violate
generic rank two. On a curve any two commuting fields
with one nonzero are constant multiples, by \eqref{cdp:eq:C.RF.8}.
A nonzero constant linear combination $V\in\mathfrak a$ is therefore vertical. Its lifted restriction to the compact general fibre satisfies \cref{lem:II:semicomplete-transport}, so that fibre has genus at most one and hence is elliptic by the adjunction argument above. Choose a second member $Y\in\mathfrak a$ with $V\wedge Y\ne0$. The fields $V,Y$ and this graph-and-Stein map $q$ satisfy \cref{lem:II:elliptic-first-integral}, which gives all the asserted elliptic-fibration conclusions.

\item Suppose $a(X)=2$. Choose a smooth projective birational model. The bracket and generic rank are preserved on the common open isomorphism, coherent Euler characteristic is birationally invariant, and \cref{lem:II:semicomplete-transport} preserves semicompleteness on the full holomorphic loci. Work on this projective model and choose a basis $V_1,V_2$ of $\mathfrak a$.

When all regular leaves of $V_1$ have algebraic curves as their Zariski closures, apply \cref{lem:II:algebraic-leaf-fibration}. Its graph resolution and Stein map \eqref{cdp:eq:C.SC.6} have rational or elliptic general fibre. The adjunction argument above excludes a rational fibre. Thus \cref{lem:II:elliptic-first-integral} applies to this map and the meromorphic lifts of $V_1,V_2$.

\item A regular leaf not contained in an algebraic curve is Zariski dense in the projective surface. For such a leaf of $V_1$, \cref{def:II:semicomplete-fields} supplies a univalent maximal solution. The projective surface is K\"ahler, so \cite[Main Theorem and the following clarification, arXiv v2 p.~2; Proposition~22, p.~25]{Guillot2019} applies. On a smooth birational model, either the surface is Abelian, when
\eqref{cdp:eq:C.RF.5} follows at once, or there is a fibration
\begin{equation}
                         q:X'\longrightarrow B
                                                               \hypertarget{tag-00YZ}{}\label{cdp:eq:C.RF.6}
\end{equation}
over a rational or elliptic curve, with rational or elliptic
general fibre, which $V_1$ preserves and whose fibres
contain every polar curve of $V_1$. Again the rational-fibre
case contradicts $\kappa(X')\ge0$. Thus \eqref{cdp:eq:C.RF.6} is an
elliptic fibration. The projection $Z_1=q_*V_1$ is nonzero:
otherwise every leaf of $V_1$ would lie in a compact fibre,
contrary to the selected dense leaf. The lifted fields remain semicomplete on their holomorphic loci by \cref{lem:II:semicomplete-transport}.

\item We prove that the polar divisor of $V_2$ is also vertical.
At a generic point where $V_1$ is holomorphic and nonzero,
the identity $[V_1,V_2]=0$ says that its local flow preserves
$V_2$ as a meromorphic tensor and therefore preserves its
polar divisor. Equivalently, in a flow coordinate
$V_1=\partial_z$, the coefficients of $V_2$ are independent
of $z$. Thus a horizontal polar component would be a
compact invariant multisection of the foliation of $V_1$.


\item Shrink to a nonempty open base $B^\circ$ on which the
fibration is smooth, $V_1$ is holomorphic on each whole
fibre, and $Z_1$ is nonzero. Its flow is defined uniformly
for small time on any one compact fibre, and maps that
whole fibre biholomorphically to a neighboring fibre.
Consequently the elliptic curves are isotrivial over
$B^\circ$. Suppose that $C\subset X'$ were a compact
horizontal invariant curve. Its normalization $\widetilde C$
maps finitely to $B$. In the normalized fibre product with
$\widetilde C$, the diagonal curve supplies a section; it is horizontal for the lifted field. Resolve
this fibre product, taking a resolution which is an
isomorphism above the smooth open locus.

\item Fix an elliptic curve $A$ isomorphic to a smooth fibre.
Over a smaller open subset of $\widetilde C$, the
origin-preserving isomorphisms from $A$ to the fibres
form a finite \'etale torsor with structure group
$\operatorname{Aut}(A,0)$. This is a concrete finite
cover: in a Weierstrass equation with constant $j$, the
change of coordinates is obtained by adjoining the
scalars $\lambda$ with
$g_2=\lambda^4g_{2,A}$ and
$g_3=\lambda^6g_{3,A}$, omitting the equation with
zero constant coefficient when $j=0$ or $1728$.
Its ambiguity is exactly the finite origin-preserving
automorphism group. Normalize the compact curve
$\widetilde C$ in a connected component of this finite
cover. On the resulting compact curve $B'$ there is
a product isomorphism over a nonempty open set,
and resolving its graph gives proper maps
\begin{equation}
 \begin{array}{ccc}
 \widetilde X'&\longrightarrow&X'\\
 \downarrow&&\downarrow q\\
 B'&\longrightarrow&B,
 \end{array}
 \qquad
 \widetilde X'\dashrightarrow A\times B'
 \quad\hbox{birational}.
                                                               \hypertarget{tag-00Z0}{}\label{cdp:eq:C.RF.7}
\end{equation}
Here all curves and surfaces are projective, so the
isomorphism over the algebraic open set defines the
displayed rational map; its graph closure is proper.
The section is carried to the zero section.

\item On the product open set the lifted field has the form
$a(b)\partial_b+c(b)\partial_u$, where $\partial_u$
is a fixed translation field on $A$. Indeed, its vertical
coefficient is holomorphic on each whole compact elliptic
fibre, so is constant on that fibre. Tangency to the zero
section gives $c=0$. Every local leaf therefore lies in
$\{u\}\times B'$. The compact closure of this section
in the proper graph model maps to an algebraic invariant
curve in $X'$. Through a general point this contains
the entire leaf by uniqueness of analytic continuation.
This contradicts the dense leaf of $V_1$. 

\item It follows that $V_2$ has no horizontal polar curve.
Hence it is holomorphic on every fibre over a suitable
nonempty open subset of $B^\circ$. Its normal component
there is a holomorphic function on a compact elliptic
curve and is constant. It consequently projects to a
meromorphic field $Z_2$ on $B$; the coefficients descend
through the proper connected-fibre map, or can be
recovered by the trace along any finite multisection.
The bracket identity gives $[Z_1,Z_2]=0$, so \eqref{cdp:eq:C.RF.8} gives $Z_2=cZ_1$ for a constant $c\in\C$.
Thus $W=V_2-cV_1$ is vertical. It is nonzero and
$W\wedge V_1\ne0$ by the generic-rank-two hypothesis.
It is semicomplete because it is a constant member of
$\mathfrak a$, and it commutes with $V_1$. Apply \cref{lem:II:elliptic-first-integral} to $W,V_1$ and \eqref{cdp:eq:C.RF.6}. Its hypotheses hold by the displayed wedge, constant-linear-combination and commutation identities, so it gives the stated relatively minimal model, fixed period lattice, identity integral monodromy and \eqref{cdp:eq:C.RF.5}.

\item Finally $\chi(\mathcal O)=2$ on a K3 surface and
$\chi(\mathcal O)=1$ on an Enriques surface. These
values and coherent Euler characteristic are preserved
under the point blowups and blowdowns comparing
smooth birational surface models. They exclude those
two minimal models for the rank-two covering component.

\end{enumerate}
\end{proof}

\section{Cyclic quotients and minimal component models}\label{sec:II:component-model-classification}

\subsection{Cyclic quotients of two-tori}\label{subsec:II:cyclic-torus-quotients}

\begin{proposition}[A cyclic torus quotient with nonnegative Kodaira dimension is a torus, bielliptic, or K3]\hypertarget{tag-01D1}{}\label{cdp:st:C.2.UN.23}
\StatementRecord{cdp:st:C.2.UN.23}{proved}{01D1}{\SeriesRef{cdp:st:C.2.UN.16}, \SeriesRef{cdp:eq:C.CQ.3}, \SeriesRef{cdp:eq:C.CQ.4}, \SeriesRef{cdp:eq:C.CQ.6}, \SeriesRef{cdp:eq:C.CQ.9}, \SeriesRef{cdp:eq:C.CQ.11}, \cref{lem:II:cyclic-component-cover}}

Let $A$ be a compact complex two-torus, let $G$ be a finite
cyclic group of holomorphic automorphisms acting effectively
on $A$, and let $Z\to A/G$ be a smooth resolution.
Every plurigenus $P_m(Z)$, $m>0$, is at most one. There are
the following exhaustive alternatives for a generator
$g(z)=Lz+b$:
\begin{equation}
\begin{array}{c|c}
\text{linear part and action}&\text{minimal resolution type}\\ \hline
L=I&\text{a two-torus}\\
\dim_{\C}\ker(L-I)=1,\ G\text{ acts freely}
   &\text{a bielliptic surface}\\
\ker(L-I)=0,\ \det L=1&\text{a K3 surface}\\
\text{all other cases}&P_m(Z)=0\text{ for every }m>0 .
\end{array}                                                    \hypertarget{tag-00CJ}{}\label{cdp:eq:C.CQ.1}
\end{equation}
In particular an Enriques surface is not bimeromorphic to
a cyclic quotient of a complex two-torus.

For the component cover in \cref{cdp:st:C.2.UN.16}, suppose its minimal model is a two-torus. The quotient by its cyclic stabilizer and the proper birational maps of \cref{lem:II:cyclic-component-cover} give
\eqref{cdp:eq:C.CQ.1}. Thus an original component of nonnegative
Kodaira dimension in this branch has minimal model a
torus, a bielliptic surface, or a K3 surface.
\end{proposition}


\begin{proof}

\begin{enumerate}[label=\arabic*.,leftmargin=*]
\item The derivative of a holomorphic torus automorphism is
constant in the translation trivialization: its entries
are holomorphic functions on a compact torus. Lifting to
the universal cover consequently writes the automorphism
as $z\mapsto Lz+b$, with $L\Lambda=\Lambda$ for the
lattice $\Lambda$. For a finite-order automorphism, $L$
is diagonalizable with roots of unity as eigenvalues.

\item Resolve the meromorphic map from $A$ to $Z$ over $A/G$:
\begin{equation}
\begin{array}{ccc}
A'&\xrightarrow{\ v\ }&Z\\
\downarrow{\scriptstyle u}&&\downarrow\\
A&\longrightarrow&A/G .
\end{array}                                                    \hypertarget{tag-00CN}{}\label{cdp:eq:C.CQ.2}
\end{equation}
Here $u$ is a sequence of point blowups and $v$ is
generically finite. Pullback of a pluricanonical form
through $v$ is holomorphic. On the open set where $u$
is an isomorphism it gives a pluricanonical form on
$A$; its coefficients extend across the finitely many
missing points by Hartogs. Hence there is an injective
map
\begin{equation}
 H^0(Z,mK_Z)\longrightarrow H^0(A,mK_A)
                  =\C\,\sigma^{\otimes m},
 \qquad \sigma=dz_1\wedge dz_2 .
                                                               \hypertarget{tag-00CO}{}\label{cdp:eq:C.CQ.3}
\end{equation}
Injectivity follows on the dense locus of the dominant
map. The image is $G$-invariant, so a nonzero image
also requires $(\det L)^m=1$.

\item If the quotient map $A\to A/G$ is ramified along a curve, the quotient is smooth at a generic point of its branch image and the local map is
$(x,y)\mapsto(x^e,y)$ with $e>1$. The Jacobian
pullback of a holomorphic $m$-canonical form vanishes
to order at least $m(e-1)$ along that curve. This is
incompatible with the nonzero constant form in
\eqref{cdp:eq:C.CQ.3}. We have proved
\begin{equation}
 \text{a branch curve in }A\longrightarrow A/G
       \quad\Longrightarrow\quad P_m(Z)=0\quad(m>0).
                                                               \hypertarget{tag-00CP}{}\label{cdp:eq:C.CQ.4}
\end{equation}


\item When $L=I$, the group consists of translations. The
action is free and its quotient is the torus obtained
by enlarging $\Lambda$ by the finitely many translation
vectors. This is the first row of \eqref{cdp:eq:C.CQ.1}.

\item Suppose that $L$ has exactly one eigenvalue equal to
one. Every power of $L$ has that eigenvalue. A
nonidentity element of $G$ with a fixed point therefore
has a fixed curve, unless its derivative is the
identity; in the latter case it is a translation and
has no fixed point. Thus a nonfree action gives
\eqref{cdp:eq:C.CQ.4}. A free action gives a smooth finite
\'etale quotient. Its canonical line has finite order,
since the finite-order character $\det L$ is the
action on $\sigma$; it has no nonzero canonical
form. Averaging over the finite cover gives
\begin{equation}
 h^1(Z,\OO_Z)=1,\qquad h^0(Z,K_Z)=0,\qquad
                  b_1(Z)=2 .
                                                               \hypertarget{tag-00CQ}{}\label{cdp:eq:C.CQ.5}
\end{equation}
There is also an explicit product-elliptic comparison.
Let
$A_+=\ker(L-I)^0$ and $A_-=\operatorname{im}(L-I)$.
They are elliptic subtori, their tangent spaces are
complementary, and the addition map
\begin{equation}
                  A_+\times A_-\longrightarrow A
                                                               \hypertarget{tag-00CR}{}\label{cdp:eq:C.CQ.6}
\end{equation}
is a finite \'etale isogeny. The finite group of lifts
of $G$ to this product acts freely, by translations
on the first factor and affine automorphisms on the
second; its quotient is exactly $A/G$. To see the
lift explicitly, lift the translation vector $b$
through \eqref{cdp:eq:C.CQ.6} and use the linear action
$(1,L|_{A_-})$. Its power equal to the order of $g$
is a deck translation of \eqref{cdp:eq:C.CQ.6}, so the group
generated with the finite deck group is finite.
This proves the stated finite free product quotient,
the bielliptic presentation. In particular
this quotient is projective. This proves the second
row of \eqref{cdp:eq:C.CQ.1}.

\item It remains to treat $\ker(L-I)=0$. The endomorphism
$I-L$ of $A$ is an isogeny, because its derivative
is invertible. Solve $(I-L)p=b$ and translate the
origin to $p$. The generator becomes the linear
map $L$ and has a fixed point. Its order equals
the order of $L$.

\item Assume first that there is no branch curve. At the
fixed origin, no nontrivial power of $L$ can have
exactly one eigenvalue one; such a power would fix
a divisor. If the order is $n$, each eigenvalue
of $L$ is therefore a primitive $n$th root.
After changing the cyclic generator, its local
representation is
\begin{equation}
             (z_1,z_2)\longmapsto
                     (\zeta z_1,\zeta^r z_2),
       \quad 1\le r<n,\quad \gcd(r,n)=1 .
                                                               \hypertarget{tag-00CS}{}\label{cdp:eq:C.CQ.7}
\end{equation}
We give the precise local pluricanonical test. The
quotient is the affine toric surface for the first
quadrant in
\begin{equation}
 N=\Z^2+\Z(1/n,r/n),\qquad
 M=\{(a,b)\in\Z^2:a+rb\equiv0\pmod n\}.
                                                               \hypertarget{tag-00CT}{}\label{cdp:eq:C.CQ.8}
\end{equation}
These are the cocharacter and character lattices:
the invariant monomials are exactly
$z_1^az_2^b$ with $(a,b)\in M$. If
$(\det L)^m=1$, the form in \eqref{cdp:eq:C.CQ.3} descends on
the smooth locus as
\[
 (z_1z_2)^m
       \left(\frac{dz_1}{z_1}\wedge
             \frac{dz_2}{z_2}\right)^{\otimes m}.
\]
Its order on the divisor of a primitive toric ray
$v=(v_1,v_2)\in N$ is
\begin{equation}
                       m(v_1+v_2-1).
                                                               \hypertarget{tag-00CU}{}\label{cdp:eq:C.CQ.9}
\end{equation}
Indeed a unimodular toric chart expresses the
monomial's order as its character pairing with $v$,
while the logarithmic volume form has one simple
pole along that divisor. The same formula holds
on a refinement, by its explicit monomial map.

\item The vector $(1/n,r/n)$ is primitive in $N$: a proper integral divisor of it would have positive first coordinate smaller than $1/n$, whereas every first coordinate in $N$ lies in $(1/n)\Z$. Include its ray in a smooth fan subdivision of the quadrant. To obtain such a subdivision, first
insert this ray and then insert primitive lattice
rays until every adjacent cone has determinant
one in a lattice basis; the Euclidean algorithm
terminates this two-dimensional procedure.
When $r<n-1$, \eqref{cdp:eq:C.CQ.9} is negative on this ray.
The descended form cannot extend to that smooth
resolution. Smooth resolutions have the same
pluricanonical sections under their common
point-blowup resolution, by the Hartogs comparison
already used in \eqref{cdp:eq:C.CQ.3}. Therefore
\begin{equation}
 \ker(L-I)=0,\quad \det L\ne1
      \quad\Longrightarrow\quad P_m(Z)=0\quad(m>0),
                                                               \hypertarget{tag-00CK}{}\label{cdp:eq:C.CQ.10}
\end{equation}
also including the branch-curve case by \eqref{cdp:eq:C.CQ.4}.

\item Finally let $\det L=1$. Every stabilizer acts with
weights $(1,-1)$ after choosing its generator, so
there are no branch curves. For a stabilizer of
order $n$, the rays
\begin{equation}
 (1,0),\ ((n-1)/n,1/n),\ldots,
                 (1/n,(n-1)/n),\ (0,1)
                                                               \hypertarget{tag-00CL}{}\label{cdp:eq:C.CQ.11}
\end{equation}
give a smooth subdivision. Consecutive determinants
are $1/n$, the covolume of $N$, and all these rays
have height $v_1+v_2=1$. Formula \eqref{cdp:eq:C.CQ.9} shows
that the invariant two-form extends without a
zero or pole across every exceptional divisor.
Thus the simultaneous resolution $Z_0$ has
$K_{Z_0}\simeq\OO_{Z_0}$. Its exceptional curves
are the chains of rational curves given by
\eqref{cdp:eq:C.CQ.11}.

\item We compute the fundamental group through the quotient map. Remove the
discrete set of points in $\C^2$ fixed by
nontrivial elements of
$\Gamma=\Lambda\rtimes\langle L\rangle$.
The complement is simply connected: a compact
loop bounds a disc in $\R^4$, and a perturbation
relative to its boundary avoids the finitely
many removed points met in a compact region.
It is the universal cover of the smooth
quotient locus. Filling in a quotient point
kills the normal closure of its finite
stabilizer in $\Gamma$, by van Kampen.
For every $\lambda\in\Lambda$, the element
$(\lambda,L)$ has a fixed point, since $I-L$
is invertible. Thus all $(\lambda,L)$ are
killed, including $(0,L)$; their products
kill every translation $(\lambda,I)$ as well.
The resulting fundamental group is trivial.
Replacing a quotient point by the rational
chain in \eqref{cdp:eq:C.CQ.11} imposes the same relation:
its plumbing neighbourhood is simply connected
and has the same punctured boundary as the
quotient ball. Hence $\pi_1(Z_0)=1$.
Together with $K_{Z_0}=\OO$, this identifies
$Z_0$ as a K3 surface. It is minimal, since
adjunction excludes a rational $(-1)$-curve
on a surface with trivial canonical line.
This proves the third row.

\item An Enriques surface has a nonzero even
pluricanonical form, but has no nonzero
canonical form and has $b_1=0$. None of
the three nonnegative cases in \eqref{cdp:eq:C.CQ.1}
has those properties: the torus and K3
have a canonical form, and the bielliptic
case has $b_1=2$. All the other cases
have no positive pluricanonical form.
This proves the exclusion.

\item Use the equivariant resolution of \cref{lem:II:cyclic-component-cover} for the covering component $E$ in \cref{cdp:st:C.2.UN.16}. Its stabilizer $G$ is cyclic, and \eqref{cdp:eq:C.HQ.2} is the proper birational morphism $E/G\to S_i^\nu$. When the minimal model of $E$ is a torus $A$, the same lemma makes the contraction $\rho:E\to A$ equivariant and gives its proper birational quotient map \eqref{cdp:eq:C.HQ.3}. An element acting trivially on $A$ fixes the dense open subset where $\rho$ is an isomorphism and hence acts trivially on $E$. Thus the induced action on $A$ is effective. The maps give
\begin{equation}
 \begin{array}{ccc}
 E&\longrightarrow&A\\
 \downarrow&&\downarrow\\
 E/G&\longrightarrow&A/G
 \end{array}
 \qquad
             E/G\dashrightarrow S_i^\nu
                                                               \hypertarget{tag-00CM}{}\label{cdp:eq:C.CQ.12}
\end{equation}
where the last birational comparison is the morphism \eqref{cdp:eq:C.HQ.2}. Resolve the graph of these proper quotient maps. Pullback and the same Hartogs comparison as in \eqref{cdp:eq:C.CQ.3} identify the pluricanonical spaces of the resulting smooth models. Applying \eqref{cdp:eq:C.CQ.1} therefore gives the stated minimal models of components with nonnegative Kodaira dimension.

\end{enumerate}
\end{proof}
\begin{remark}[Projective torus quotients and their linear groups]\hypertarget{tag-P304}{}\label{rem:II:projective-torus-quotients}
\StatementRecord{rem:II:projective-torus-quotients}{proved}{P304}{\cref{cdp:st:C.2.UN.23}; \cite[Theorem~2.1, p.~136; Lemma~3.4, p.~139; Example~2.3, p.~137; Lemma~3.6, pp.~141--142]{Yoshihara1995}}
For an Abelian surface $A$ and a finite group $G$ of automorphisms, let $G_0$ be its translation subgroup. The derivative homomorphism $g\mapsto L_g$ has kernel $G_0$ and image $H\simeq G/G_0$. The classification in \cite[Theorem~2.1, p.~136; Lemma~3.4, p.~139]{Yoshihara1995} describes the birational type of a smooth resolution of the same quotient map $A\to A/G$ in terms of $H$ and the fixed loci. Its Enriques cases have $H$ equal to the Klein four group or the dihedral group of order eight, as in \cite[Example~2.3, p.~137; Lemma~3.6, pp.~141--142]{Yoshihara1995}. For a cyclic group $G=\langle g\rangle$, the image is $H=\langle L_g\rangle$. The Klein four group has three elements of order two, and the dihedral group of order eight is nonabelian, so neither is cyclic. This identifies the projective specialization of \cref{cdp:st:C.2.UN.23} through the identity on $A$ and its quotient morphism, and excludes the Enriques rows of the cited classification.
\end{remark}

\subsection{Rank-one fields on elliptic surfaces}\label{subsec:II:rank-one-elliptic-fields}

\begin{proposition}[Every compact elliptic surface satisfies the intrinsic rank-one field condition]\hypertarget{tag-01D4}{}\label{cdp:st:C.2.UN.26}
\StatementRecord{cdp:st:C.2.UN.26}{proved}{01D4}{\SeriesRef{cdp:st:C.2.UN.16}, \SeriesRef{cdp:eq:C.FL.3}; \cite[Theorem~3.3, author preprint p.~9]{Guillot2024}}

Let $q:X\to B$ be a proper holomorphic map from a
smooth compact connected complex surface to a smooth
compact curve, with connected elliptic general fibre.
There are commuting meromorphic vector fields $V_1,V_2$
which are linearly independent over $\mathbb C$, have
generic rank one, and whose every constant linear
combination is semicomplete on its full holomorphic
locus. Their ratio has a graph resolution and
Stein factor birationally equal to $q$.

For a finite cyclic group $G$ acting holomorphically on $X$ and preserving the fibres of $q$ as a family, the two-dimensional space can be chosen $G$-stable.
\end{proposition}


\begin{proof}

\begin{enumerate}[label=\arabic*.,leftmargin=*]
\item The relative canonical line is the invertible
sheaf
\begin{equation}
 \omega_{X/B}=K_X\otimes q^*K_B^{-1},
 \qquad \mathcal E=q_*\omega_{X/B}.
                                                               \hypertarget{tag-00KR}{}\label{cdp:eq:C.FL.1}
\end{equation}
Proper coherence makes $\mathcal E$ coherent. It is
torsion-free: a section of a line bundle annihilated
by a nonzero base function vanishes over the dense
complement of its zero fibre and hence vanishes
identically. On a smooth elliptic fibre the canonical
space has dimension one, so $\mathcal E$ has rank one.
It is therefore a line bundle on $B$.

\item Riemann--Roch on $B$ gives a nonzero holomorphic section of $\mathcal E(D)$ for an effective divisor $D$ of sufficiently large degree. Divide it by the canonical section of $\OO_B(D)$ to obtain a nonzero meromorphic section $\sigma$ of $\mathcal E$.
Evaluation gives a nonzero meromorphic section
\begin{equation}
 a:\mathcal O_X\dashrightarrow\omega_{X/B}
       \quad\hbox{from}\quad q^*\mathcal E\longrightarrow
                                      \omega_{X/B}.
                                                               \hypertarget{tag-00KS}{}\label{cdp:eq:C.FL.2}
\end{equation}
The meromorphic section $a$ is the evaluation of $q^*\sigma$ by the displayed holomorphic sheaf map. On the smooth part of $q$, the differential identification $\omega_{X/B}\simeq\Omega^1_{X/B}$ makes $a$ a relative one-form.

\item Its dual defines a global meromorphic vertical field $V$. On a coordinate disc
of $B$ with coordinate $z$, set
$\Omega_z=a\otimes q^*(dz)$, a meromorphic section
of $K_X$. Define $V$ by
\begin{equation}
                   \iota_V\Omega_z=-q^*(dz).
                                                               \hypertarget{tag-00KT}{}\label{cdp:eq:C.FL.3}
\end{equation}
Contraction by a nonzero meromorphic two-form is an
isomorphism between the meromorphic tangent and
cotangent bundles, so this has a unique meromorphic
solution. Under $z'=\phi(z)$ both sides are multiplied
by $q^*\phi'(z)$; the solution is unchanged and
therefore glues. The alternating identity in
\eqref{cdp:eq:C.FL.3} gives $dq(V)=0$. On a regular fibre its
relative time form is $a$, so $a(V)=1$ with the
usual convention $\Omega_z=q^*(dz)\wedge a$.

\item Remove finitely many base values so that $q$ is
smooth and the evaluated section is holomorphic
and nonzero on every remaining fibre. Such a
removal is possible by generic base change for
$q_*\omega_{X/B}$ and the finite zero and pole
divisors of $\sigma$. A nonzero canonical form
on a compact elliptic curve is nowhere zero.
Thus $V$ restricts on each of these fibres to
a nonzero holomorphic translation field and is
complete there.

\item Choose any nonconstant meromorphic function
$h$ on $B$, and put
\begin{equation}
              V_1=V,\qquad V_2=(q^*h)V.
                                                               \hypertarget{tag-00KU}{}\label{cdp:eq:C.FL.4}
\end{equation}
Their bracket vanishes because $V(q^*h)=0$.
They are independent over $\mathbb C$ since $h$
is nonconstant, and have generic rank one by
the formula. For $c,d\in\mathbb C$, the field
$W=(c+dq^*h)V$ is complete on the dense union
of the preceding compact smooth fibres away
from the poles of $h$. Its coefficient is
constant on each such fibre; zero coefficients
give constant complete solutions. Apply \cite[Theorem~3.3, author preprint p.~9]{Guillot2024} on the full holomorphic locus $X_W$ of $W$. Its uniformizability locus is closed in $X_W$ and contains this dense union of compact fibres. It is therefore all of $X_W$, proving semicompleteness of every $W$.

\item The ratio is $q^*h$, so the resulting map is
the composite $X\xrightarrow q B
\xrightarrow h\mathbb P^1$. The second arrow
is finite. Its Stein factor is $q$ itself:
$q_*\mathcal O_X=\mathcal O_B$, and $B$ is
the normalization in the field of functions
constant on the connected general fibres.


\item Since $G$ preserves the fibre family, $q\circ g$ is constant on each connected fibre for every $g\in G$. The equality $q_*\OO_X=\OO_B$ descends it to a holomorphic map $g_B:B\to B$ with $q\circ g=g_B\circ q$. The corresponding map for $g^{-1}$ is its inverse, so these maps give a finite group of automorphisms of $B$. The finite orbit of $V$ spans a nonzero finite-dimensional representation of $G$ by meromorphic vertical fields. A cyclic complex representation has
an eigenvector $V'\ne0$. On a general smooth
fibre each orbit field is holomorphic and
vertical, and therefore a translation field.
The same is true of $V'$; after removing
finitely many further values it is nonzero
on every remaining fibre. Choose a nonconstant
meromorphic function $h$ on the quotient
curve $B/G$, and pull it back to $B$.
Then the space
\begin{equation}
                  \mathbb C V'+\mathbb C(q^*h)V'
                                                               \hypertarget{tag-00KV}{}\label{cdp:eq:C.FL.5}
\end{equation}
is $G$-stable, with the same eigencharacter
on its two displayed generators. Its members
are semicomplete by the dense compact-fibre
and closedness argument just proved, and its
ratio has the same connected-fibre Stein
factor $q$. 

\end{enumerate}
\end{proof}

\subsection{Compact component cores and geometric genus}\label{subsec:II:component-core-genus}

\begin{lemma}[Compact cycles in smooth component loci]\hypertarget{tag-P305}{}\label{lem:II:compact-component-transport}
\StatementRecord{lem:II:compact-component-transport}{proved}{P305}{\cref{lem:II:cyclic-component-cover}; \cite[Theorem~3.44, pp.~254--255]{Hatcher2002}}
\begin{enumerate}[label=\arabic*.,leftmargin=*]
\item Let $E$ be a smooth compact complex surface and $D\subset E$ a reduced compact analytic set of dimension at most one. Under Poincar\'e--Alexander duality, the map
$H_2(E,\Z)\to H_2(E,E\setminus D;\Z)\simeq H^2(D,\Z)$
is the vector of intersection numbers with the curve components of $D$. The target is free, with one generator for each such curve, and the kernel is the image of $H_2(E\setminus D,\Z)$. With real coefficients the same maps give
\begin{equation}
 H_2(E\setminus D,\mathbb R)\longrightarrow
 H_2(E,\mathbb R)\xrightarrow{\ I_D\ }
 H_2(E,E\setminus D;\mathbb R)
       \simeq H^2(D,\mathbb R).
                                                               \hypertarget{tag-00V4}{}\label{cdp:eq:C.PG.6}
\end{equation}
The holomorphic two-forms determine
\begin{equation}
 P_E=\{[\operatorname{Re}\alpha]:
                       \alpha\in H^0(E,K_E)\}
       \subset H^2(E,\mathbb R),
 \qquad \dim_{\mathbb R}P_E=2p_g(E).
                                                               \hypertarget{tag-00V5}{}\label{cdp:eq:C.PG.7}
\end{equation}
This subspace is positive definite for the intersection form. Its Poincar\'e-dual homology classes lie in $\ker I_D$.

\item Let $f:U\to\Delta_t$ be proper and holomorphic, with $U$ a smooth complex threefold and with connected smooth two-torus fibres over $\Delta_t^*$. Write $f^{-1}(0)=\sum_i m_iS_i$ and $S=\bigcup_iS_i$. For a positive integer $N$ divisible by every $m_i$, normalize the base change $t=s^N$ and resolve its total space and central support. The resulting proper map
\begin{equation}
                         g:\widetilde U\to\Delta_s
                                                               \hypertarget{tag-00V0}{}\label{cdp:eq:C.PG.2}
\end{equation}
can be chosen unchanged over the punctured disc and over the normalized inverse image of every $S_i\setminus\operatorname{Sing}S$. Its nearby fibre is $f^{-1}(s^N)$. The strict components above $S_i$ have multiplicity one, and their union restricts over $S_i\setminus\operatorname{Sing}S$ to a finite \'etale cover of total degree $m_i$, possibly disconnected.

\item On any smooth resolved finite-base-change model of this family, choose finitely many compact finite smooth singular two-cycles in the smooth single-component loci of multiplicity-one components. Include compact transverse perturbations used to compute their intersection numbers. After shrinking the base disc, there are simultaneous smooth embeddings of neighbourhoods of these compact sets into each fibre, preserving the complex orientation and all these intersection numbers. Neighbourhoods belonging to different components have disjoint images. The transported cycles form continuous families over the whole small disc, and their classes on a nearby fibre are invariant under its monodromy.
\end{enumerate}
\end{lemma}
\begin{proof}
\begin{enumerate}[label=\arabic*.,leftmargin=*]
\item The compact analytic set $D$ is triangulable and locally contractible. Apply \cite[Theorem~3.44, pp.~254--255]{Hatcher2002} with $M=E$, $K=D$ and real dimension four. The source theorem identifies relative homology with cohomology of $D$ by cap product with the orientation class. In its naturality diagram, $H_2(E,\Z)\to H_2(E,E\setminus D;\Z)$ corresponds to restriction of $H^2(E,\Z)$ to $D$. Normalization of the curve part changes only identifications at finitely many points, which leave degree-two cohomology unchanged; isolated points have degree-two cohomology zero. Evaluation on the fundamental class of each normalized curve therefore identifies this restriction with the asserted intersection vector. The relative homology sequence gives its kernel. Tensoring this sequence with $\mathbb R$ gives \eqref{cdp:eq:C.PG.6}.

Every holomorphic two-form $\alpha$ on $E$ is closed: its $(3,0)$ derivative vanishes by dimension and its antiholomorphic derivative vanishes by holomorphicity. With the complex orientation,
\begin{equation}
 \int_E(\operatorname{Re}\alpha)^2
       =\tfrac12\int_E\alpha\wedge\overline\alpha>0
                      \quad (\alpha\ne0).
                                                               \hypertarget{tag-00V6}{}\label{cdp:eq:C.PG.8}
\end{equation}
This proves positivity and injectivity of the real-part map on the underlying real vector space of $H^0(E,K_E)$, hence the dimension in \eqref{cdp:eq:C.PG.7}. The pullback of $\alpha$ to the normalization of any curve in $D$ vanishes by dimension. Its real part therefore pairs trivially with each boundary curve, so its Poincar\'e-dual class lifts to the open complement.

\item At a point of $S_i\setminus\operatorname{Sing}S$, write $t=u x^{m_i}$ with $u$ a unit. Factor $t=s^N$ through $t=r^{m_i}$ by $r=s^{N/m_i}$. Substituting this map into the normalized chart \eqref{cdp:eq:C.SC.2} of \cref{lem:II:cyclic-component-cover}, with its coordinate $z$ renamed $x$, gives
\begin{equation}
       w^{m_i}=u,\qquad x=s^{N/m_i}/w.
                                                               \hypertarget{tag-00V1}{}\label{cdp:eq:C.PG.3}
 \hypertarget{tag-00C7}{}\SsixEquationAliases{\label{cdp:eq:C.CL.8}}
\end{equation}
At $s=0$ this equation has exactly $m_i$ local unit roots, and its derivative with respect to $w$ is nonzero. Thus the normalized space is smooth there, $s$ is a normal parameter, and the central cover over the indicated open surface is \'etale of total degree $m_i$. Resolve the total space and central support while preserving these smooth normal-crossing open sets and the punctured tube, using the resolution construction of \cref{lem:II:cyclic-component-cover}. The coordinate $s$ continues to have order one on every strict component. The map to the original fibre over $s\ne0$ is the fibre-product projection, identifying it with $f^{-1}(s^N)$.

\item Let $K_j$ contain the finitely many chosen cycles and transverse perturbations in one smooth single-component locus. The map $g$ is a submersion near $K_j$: its local equation is a unit times a defining coordinate of the multiplicity-one component. Choose smooth horizontal lifts of the two real coordinate vector fields on the base. A partition of unity preserves their prescribed images under $dg$. Multiply by a cutoff outside a slightly larger compact neighbourhood and flow successively for the real and imaginary parts of $s$. For all sufficiently small $s$, these trajectories stay where the horizontal equations hold. The composition maps a neighbourhood of $K_j$ in the central component diffeomorphically onto an open subset of $g^{-1}(s)$; the two equations for $dg$ identify this exact fibre. Continuity from $s=0$ preserves the complex orientation.

The finitely many compact sets in distinct components are disjoint and avoid all intersections of central components. Choose disjoint ambient neighbourhoods for them and one small disc on which all the flows are defined and remain in those neighbourhoods. The images stay disjoint. The local diffeomorphisms preserve the signed transverse intersection numbers of the cycles in each component, and cycles from different components have intersection zero.

These two flows define each parametrized finite cycle over the whole small $s$-disc. Along a circle about zero the family returns to the identical parametrized cycle. In a local differentiable trivialization of the proper smooth family over $s\ne0$, its homology class is locally constant. It is therefore fixed by Gauss--Manin transport around that circle. The construction applies to integral chains and, by real linearity, to real cycles.
\end{enumerate}
\end{proof}


\begin{proposition}[The sum of the normalization geometric genera in one torus fibre is at most one]\hypertarget{tag-01D6}{}\label{cdp:st:C.2.UN.28}
\StatementRecord{cdp:st:C.2.UN.28}{proved}{01D6}{\cref{lem:II:compact-component-transport}; \eqref{cdp:eq:C.PG.5}}

Let $f:U\to\Delta$ be a proper holomorphic map from a
smooth complex threefold, with connected two-torus
fibres over $\Delta^*$. Write its full central fibre
as $\sum_i m_iS_i$, and let $Z_i\to S_i^\nu$ be
smooth resolutions of the normalizations. Then
\begin{equation}
                         \sum_i p_g(Z_i)\le1.
                                                               \hypertarget{tag-00UZ}{}\label{cdp:eq:C.PG.1}
\end{equation}
At most one normalization component in a given fibre has a nonzero holomorphic two-form, and that component has geometric genus one.

The same bound holds for any collection of multiplicity-one components in a smooth resolved finite base change of this family.
\end{proposition}


\begin{proof}

\begin{enumerate}[label=\arabic*.,leftmargin=*]
\item Apply \cref{lem:II:compact-component-transport} to $f$ with a common multiple $N$ of the $m_i$. It gives the resolved map \eqref{cdp:eq:C.PG.2}, whose strict components above each $S_i$ have multiplicity one and whose nearby fibres are the original tori.

\item Select one smooth strict component $E_i$ above each
$S_i$. Its map to $S_i^\nu$ is generically finite.
Resolve the meromorphic comparison to $Z_i$:
\begin{equation}
              E_i\xleftarrow{\rho_i}W_i
                            \xrightarrow{v_i}Z_i.
                                                               \hypertarget{tag-00V2}{}\label{cdp:eq:C.PG.4}
\end{equation}
The first map is a sequence of point blowups,
and the second is generically finite. Pullback
of holomorphic two-forms through $v_i$ is
injective, since its Jacobian is nonzero on
a dense open set. Point-blowup invariance of
holomorphic two-forms follows by restriction
and Hartogs extension across the blown-down
points. Hence
\begin{equation}
                         p_g(E_i)\ge p_g(Z_i).
                                                               \hypertarget{tag-00V3}{}\label{cdp:eq:C.PG.5}
\end{equation}
It suffices to prove the bound for the selected
components $E_i$ of \eqref{cdp:eq:C.PG.2}.

\item Let $D_i$ be the reduced intersection of $E_i$ with the other components of the resolved central support. It is a compact normal-crossing curve, or is empty. Apply the first part of \cref{lem:II:compact-component-transport} with $(E,D)=(E_i,D_i)$. Each Poincar\'e-dual class of $P_{E_i}$ lifts to $H_2(E_i\setminus D_i,\mathbb R)$. Choose a real basis and finite smooth singular representatives of these lifted classes, together with their transverse perturbations. The third part of that lemma embeds neighbourhoods of all these compact cycles into one nearby torus $F=g^{-1}(s)$, with disjoint images for distinct $i$ and with the same intersection pairings. Their images are mutually orthogonal positive subspaces of $H_2(F,\mathbb R)$, of dimensions $2p_g(E_i)$. Preservation of the positive form proves injectivity on their direct sum.

\item The intersection form of an oriented
real four-torus is $3H$, where $H$ is
the hyperbolic plane. This follows
directly from the exterior-algebra
identification of its integral
cohomology: the three complementary
degree-two pairs give three hyperbolic
blocks. Its positive index is therefore
three. We have proved
\begin{equation}
                    2\sum_i p_g(E_i)\le3.
                                                               \hypertarget{tag-00V7}{}\label{cdp:eq:C.PG.9}
\end{equation}
The sum is an integer, so it is at most one. Together with \eqref{cdp:eq:C.PG.5} this gives \eqref{cdp:eq:C.PG.1}. For any collection of multiplicity-one components in a smooth resolved finite base change, apply the first and third parts of \cref{lem:II:compact-component-transport} directly to their complements of the other central components. Their positive subspaces again embed orthogonally in one nearby torus, giving the same inequality.

\end{enumerate}
\end{proof}

\subsection{Elliptic return maps and component genera}\label{subsec:II:elliptic-return-maps}

\begin{proposition}[Rank-one elliptic component covers have only smooth multiple fibres]\hypertarget{tag-01D7}{}\label{cdp:st:C.2.UN.29}
\StatementRecord{cdp:st:C.2.UN.29}{proved}{01D7}{\cref{cdp:st:C.2.UN.16,cdp:st:C.2.UN.22,lem:II:identity-monodromy-elliptic,lem:II:semicomplete-transport,cdp:st:C.2.UN.28,cdp:st:C.2.UN.20,cdp:st:C.2.UN.15,lem:II:cyclic-component-cover}; \cite[Chapter~1, Proposition~2.2, p.~8; \S1, pp.~5--6; \S3, p.~9]{Peters1990}}

Let $E$ be a multiplicity-one covering component
constructed in \ref{cdp:st:C.2.UN.16}, with its saturated two-dimensional
space of limiting translation fields. In the generic-rank-one
branch, take the graph resolution and Stein factor
of their ratio. When its connected general fibre is elliptic,
the resulting relatively minimal elliptic surface
\begin{equation}
                         q:E_0\longrightarrow B
                                                               \hypertarget{tag-00Z3}{}\label{cdp:eq:C.RG.1}
\end{equation}
has identity integral $H_1$ monodromy. Every fibre is a
positive multiple of a smooth elliptic curve, and all smooth
fibres are torsors under one fixed elliptic curve. Moreover
\begin{equation}
 R^1q_*\OO_{E_0}\simeq\OO_B,\qquad
 p_g(E)=g(B)\le1,\qquad \chi(\OO_E)=0.
                                                               \hypertarget{tag-00Z8}{}\label{cdp:eq:C.RG.2}
\end{equation}
The maps to $E_0$ are the point blowups resolving the first integral followed by vertical exceptional-curve contractions.

Every root-cover component of nonnegative Kodaira dimension has coherent Euler characteristic zero; its minimal model is neither K3 nor Enriques. Every original component of Kodaira dimension one has an isotrivial elliptic Iitaka fibration over a curve of genus at most one, obtained through the cyclic quotient and resolution maps of \cref{lem:II:cyclic-component-cover,cdp:st:C.2.UN.20}.
\end{proposition}


\begin{proof}

\begin{enumerate}[label=\arabic*.,leftmargin=*]
\item We work on the root-change and resolution model
$f:U\to\Delta_s$ in \ref{cdp:st:C.2.UN.16}. Its smooth general fibres are complex two-tori. The chosen component $E$
has multiplicity one. Its two ambient meromorphic fields
$V_1,V_2$ are vertical, commute, are holomorphic and
independent on each general torus, and have no pole along
the generic point of $E$. Their restrictions are independent
over $\C$ and have rank one over the meromorphic function
field of $E$.

\item Resolve the meromorphic ratio on $E$ by point blowups
and take its connected-fibre Stein factor $q:E\to B$.
Blow up the same points in the smooth ambient threefold. In local coordinates with $E$ a coordinate divisor, the blowup charts restrict to the point blowup of $E$. Its strict transform still has multiplicity one, since the map is an isomorphism at its generic point, and both fields lift meromorphically. Their generic restrictions and bracket
are unchanged through the common open isomorphism.
We use the same letter $E$ for this modified component.
There is a nonconstant meromorphic function $H$ on $B$
such that
\begin{equation}
                    V_2|_E=(q^*H)V_1|_E.
                                                               \hypertarget{tag-00Z9}{}\label{cdp:eq:C.RG.3}
\end{equation}
The function $H$ is the meromorphic function on the Stein base induced by the ratio. The bracket identity gives $V_1(q^*H)=0$, so $V_1|_E$ is tangent to the fibres of $q$.

\item Remove finitely many values of $B$ so that $q$ has smooth
connected elliptic fibres and $V_1|_E$ is a nonzero
holomorphic translation field on each of them. The latter
assertion follows by applying the threefold specialization and compact-curve clauses of \cref{lem:II:semicomplete-transport} to the ambient field $V_1$, the multiplicity-one component $E$ and its elliptic fibre: the induced field is semicomplete, and its restriction to that compact curve has no pole. After excluding the finitely
many fibres on which its restriction is identically zero,
it is nowhere zero on every remaining fibre. Remove also
the poles and critical points of $H$.

\item Let $D$ be the union on $E$ of its intersections with
the other central divisors and with the ambient
polar divisors of the two fields. Each remaining general
elliptic fibre meets $D$ in finitely many points, after
discarding any finitely many fibres contained in $D$.
The fields are holomorphic on $U^*$; hence all their
ambient polar divisors lie over zero. In particular
outside this explicitly removed set the ambient fields
and the submersion $f$ are holomorphic. At a suitable
point $p$ of a remaining fibre $F_{b_0}$ there are
holomorphic coordinates $(u,z,s)$ with
\begin{equation}
 f=s,\qquad V_1=\partial_u,\qquad
 V_2=h(z,s)\partial_u+A(z,s)\partial_z,
 \qquad h(z,0)=H(z).
                                                               \hypertarget{tag-00ZA}{}\label{cdp:eq:C.RG.4}
\end{equation}
Here $z$ is a local coordinate on $B$ along $E$.
Indeed the holomorphic flow-box theorem gives $V_1=\partial_u$,
verticality preserves $s$, and commutation makes both
coefficients of $V_2$ independent of $u$. The coordinate $z$ in this construction is the local base coordinate in this flow box.

\item Since the ambient fields are independent on the nearby
tori, $A$ is not identically zero. Since their restrictions
to $E$ have rank one, $A(z,0)=0$. Choosing $p$ away
from the zeros of its leading coefficient gives
\begin{equation}
 A(z,s)=s^d\bigl(a(z)+O(s)\bigr),\qquad d\ge1,
 \qquad a(z)\ne0,\qquad H'(z)\ne0
                                                               \hypertarget{tag-00ZB}{}\label{cdp:eq:C.RG.5}
\end{equation}
on a smaller disc. The order $d$ is finite; it is the
order of the ambient wedge along this component in the
chosen nonvanishing frame.

Let $\lambda_1(z),\lambda_2(z)$ be the periods of the
time form of $V_1|_{F_z}$, for an oriented integral
homology basis continued over a small disc about $b_0$.
These are holomorphic functions; a local holomorphic
section and the ordinary integration construction of
an elliptic curve supply the marked periods. Put
$c_i=\lambda_i(b_0)$.

\item The time covering $\C\to F_{b_0}$ of the translation
field pulls the finite bad set $D\cap F_{b_0}$ back to
a discrete subset of $\C$. Choose a piecewise smooth
time path from $0$ to $c_i$ avoiding this discrete set.
Its image is a compact path on which both ambient fields
are holomorphic and $V_1$ is nonzero. A finite chain of
holomorphic flow charts therefore continues the solution
along this entire path for every initial point in a
sufficiently small ambient neighbourhood of $p$.
Its endpoint is again near $p$, since the central path
closes. This constructs a holomorphic return germ $\Phi_i$ on an ambient neighbourhood of $p$.

\item Every local continuation of the flow of $V_1$ preserves
$f,V_1,V_2$, by verticality and the bracket identity.
Thus the return germ in the single flow box has the form
\begin{equation}
 \Phi_i(u,z,s)=\bigl(u+\delta_i(z,s),\phi_i(z,s),s\bigr).
                                                               \hypertarget{tag-00ZC}{}\label{cdp:eq:C.RG.6}
\end{equation}
On $E$ the first integral is constant along its elliptic
orbits, and comparison with the continued period gives
\begin{equation}
 \phi_i(z,0)=z,\qquad
 \delta_i(z,0)=c_i-\lambda_i(z).
                                                               \hypertarget{tag-00ZD}{}\label{cdp:eq:C.RG.7}
\end{equation}
To see the second equality directly, follow the fixed
time path of total time $c_i$ on $F_z$, then return
to the local section by the small flow segment of
time $-\delta_i(z,0)$. The resulting closed path is
the continued integral cycle, so its time integral is
$c_i-\delta_i(z,0)=\lambda_i(z)$. This also specifies
which integral cycle and period enter the return map.

\item Pushing $V_2$ forward by \eqref{cdp:eq:C.RG.6} and equating it to
$V_2$ at the endpoint gives the two exact identities
\begin{equation}
 \begin{split}
 A(\phi_i(z,s),s)&=A(z,s)\,\partial_z\phi_i(z,s),\\
 h(\phi_i(z,s),s)&=h(z,s)+A(z,s)\,\partial_z\delta_i(z,s).
 \end{split}                                                   \hypertarget{tag-00ZE}{}\label{cdp:eq:C.RG.8}
\end{equation}
\item The second identity of \eqref{cdp:eq:C.RG.8}, together with $H'$ being
a unit and \eqref{cdp:eq:C.RG.5}, first gives
$\phi_i-z=O(s^d)$. For instance, factor
$h(\phi_i,s)-h(z,s)=(\phi_i-z)h_1(z,s)$;
the factor $h_1$ is a unit because $h_1(z,0)=H'(z)$.
Write
\begin{equation}
                \phi_i=z+s^d\psi_i(z)+O(s^{d+1}).
                                                               \hypertarget{tag-00ZF}{}\label{cdp:eq:C.RG.9}
\end{equation}
In the first identity of \eqref{cdp:eq:C.RG.8}, all coefficients
of $A$ below order $2d$ cancel. At order $2d$ the
remaining equality is
\begin{equation}
                         a'\psi_i=a\psi_i'.
                                                               \hypertarget{tag-00Z4}{}\label{cdp:eq:C.RG.10}
\end{equation}
Terms $s^{d+k}a_k(z)$ with $k\ge1$ first contribute
to the difference at order $2d+k$, so they do not
enter this calculation. Since $a$ is a unit,
$(\psi_i/a)'=0$. Hence $\psi_i=k_i a$ for a constant
$k_i\in\C$. At order $d$, the second identity gives
\begin{equation}
 H'\psi_i=a\,\partial_z\delta_i(z,0)
          =-a\lambda_i',\qquad
             \lambda_i'= -k_i H'.
                                                               \hypertarget{tag-00Z5}{}\label{cdp:eq:C.RG.11}
\end{equation}
Consequently there are constants $d_i\in\C$ with
\begin{equation}
                         \lambda_i=d_i-k_iH
                    \quad(i=1,2)
                                                               \hypertarget{tag-00Z6}{}\label{cdp:eq:C.RG.12}
\end{equation}
on the small disc.

\item Let $B^\circ$ be a connected complement of finitely
many values containing the disc used above, on which
the elliptic fibration, time form, and $H$ are regular.
Continue the chosen homology basis and periods on its
universal cover. Both sides of \eqref{cdp:eq:C.RG.12} are holomorphic
there, so the identity theorem gives the same formula
everywhere, with the same constants. Since $H$ is a
single-valued meromorphic function on $B$, each
continued period is single-valued on $B^\circ$.
The integration map from $H_1(F_b,\Z)$ to $\C$ is
injective for the time form of a nonzero elliptic
translation field. Thus returning both periods
unchanged forces every integral monodromy matrix to be the identity on $B^\circ$. The inclusion of $B^\circ$ into the full smooth base locus is surjective on fundamental groups, since it adds back only finitely many regular values. Thus the monodromy is the identity on the full smooth base locus as well.

\item Contract the vertical exceptional curves to obtain \eqref{cdp:eq:C.RG.1}. The smooth fibres and their integral monodromy are unchanged. Apply \cref{lem:II:identity-monodromy-elliptic} to this relatively minimal map $q:E_0\to B$. It gives every smooth-multiple-fibre and fixed-elliptic-curve assertion, as well as the full-base direct-image isomorphism in \eqref{cdp:eq:C.RG.2}, with the normalized root and invariant-function descent at every multiple value.

By \eqref{cdp:eq:C.RG.13}, $p_g(E_0)=g(B)$ and $\chi(\OO_{E_0})=0$. The point blowups and blowdowns comparing $E$ and $E_0$ preserve the plurigenera and coherent Euler characteristic \cite[Chapter~1, Proposition~2.2, p.~8]{Peters1990}. Hence $p_g(E)=g(B)$ and $\chi(\OO_E)=0$. The component $E$ is a multiplicity-one component of a smooth resolved finite base change of the original torus family. The finite-base-change clause of \cref{cdp:st:C.2.UN.28} therefore gives $p_g(E)\le1$, and $g(B)\le1$.

\item A root-cover component of nonnegative Kodaira dimension has either generic rank two, covered by \cref{cdp:st:C.2.UN.22}, or rank one. In rank one, \cref{lem:II:semicomplete-transport} gives genus at most one for a connected general first-integral fibre. A rational fibre would give Kodaira dimension $-\infty$ by the adjunction argument of \cref{cdp:st:C.2.UN.20}; hence the general fibre is elliptic and the preceding argument applies. All such covers therefore have coherent Euler characteristic zero. A minimal K3 surface has $p_g=1,q=0$ and coherent Euler characteristic two; an Enriques surface has $p_g=q=0$ and coherent Euler characteristic one \cite[Chapter~1, \S1, pp.~5--6; \S3, p.~9]{Peters1990}. This proves both exclusions.

\item Let an original component have Kodaira dimension one, and let $Z$ resolve its normalization. For its component cover $E$, resolve the meromorphic comparison as $E\leftarrow W\to Z$, with the first map a sequence of point blowups and the second generically finite. The nonzero generic Jacobian makes pullback of every pluricanonical section injective. Birational invariance of the plurigenera then gives $\kappa(E)\ge\kappa(Z)=1$, while \cref{cdp:st:C.2.UN.15} gives $\kappa(E)\le1$.

Use the elliptic fibration $q:E_0\to B$ from \cref{cdp:st:C.2.UN.22} in rank two and from the preceding argument in rank one; both have $g(B)\le1$. It is the connected-fibre Iitaka fibration. Indeed, the canonical line restricts trivially to a smooth elliptic fibre, so the ratios of pluricanonical sections factor through $q$. Each nonconstant pluricanonical map consequently factors through a finite map from $B$ to its normalized image. Since $q_*\OO_{E_0}=\OO_B$, the connected-fibre Stein factor of the pluricanonical map has base $B$. The cyclic stabilizer therefore acts on $q$ through its action on the pluricanonical field. Apply the proper quotient maps \eqref{cdp:eq:C.HQ.2}--\eqref{cdp:eq:C.HQ.3} of \cref{lem:II:cyclic-component-cover} and the fibration descent in \cref{cdp:st:C.2.UN.20}. Resolving their graphs gives the descended fibration on a smooth birational model of the original component. Its base is the quotient of $B$, so Riemann--Hurwitz gives genus at most one.

On a general fibre, the subgroup fixing a general base point acts by translations. Indeed, an automorphism with a nonidentity linear part has a fixed point by the equation $(1-L)x=b$ on the elliptic curve. A quotient by a group containing such an element is ramified, and Riemann--Hurwitz makes the quotient curve rational. On the resulting smooth general fibre, the normal line is trivial, so adjunction restricts every positive pluricanonical line to $\OO_{\mathbb P^1}(-2m)$. All pluricanonical sections would vanish on this dense family of fibres, contrary to Kodaira dimension one.

The translation subgroup in the fixed elliptic curve is locally constant: each of its holomorphic translation parameters has order dividing the finite group order and hence takes values in a finite torsion subgroup. The general quotient fibres are therefore torsors under the fixed quotient elliptic curve. These proper quotient and graph maps give the isotrivial Iitaka fibration on the original component's smooth model.

\end{enumerate}
\end{proof}

\subsection{Cyclic quotients of smooth multiple elliptic fibres}\label{subsec:II:cyclic-component-models}

\begin{lemma}[Local cyclic quotients of elliptic fibrations]\hypertarget{tag-P317}{}\label{lem:II:local-cyclic-elliptic}
\StatementRecord{lem:II:local-cyclic-elliptic}{proved}{P317}{\cref{lem:II:identity-monodromy-elliptic}}
Let $q:E\to B$ be a relatively minimal proper holomorphic map from a smooth compact complex surface to a smooth compact connected curve, with connected elliptic general fibre and identity integral $H_1$ monodromy on its smooth base locus. Let a finite cyclic group $G$ act on $E$ and $B$ with $q$ equivariant. Suppose that the kernel $H$ of the action on $B$ acts by translations on the general elliptic fibre.

The quotient $E_1=E/H$ is smooth. Its map $q_1:E_1\to B$ has connected fibres, identity integral monodromy, one fixed elliptic curve as relative Jacobian, and only positive multiples of smooth elliptic fibres. The remaining group $Q=G/H$ acts faithfully on $B$.

At a base point with stabilizer of order $s$, write $\beta$ for the stabilizer generator's linear action on the fixed elliptic curve. A nonidentity $\beta$ forces the fibre of $q_1$ to have multiplicity one. After quotienting the subgroup acting trivially on the elliptic factor, the local quotient has the product action
\[
 (v,z)\longmapsto(\zeta_d^r v,\zeta_d z),
 \qquad d\in\{2,3,4,6\},\quad r\in\{1,d-1\}.
\]
For $\beta=1$, the quotient surface and the reduced fibre support are smooth. The full fibre remains a positive multiple of that elliptic support.
\end{lemma}
\begin{proof}
\begin{enumerate}[label=\arabic*.,leftmargin=*]
\item By \cref{lem:II:identity-monodromy-elliptic}, the normalized base change $t=z^m$ at a full fibre $mF$ is a smooth elliptic family, finite \'etale over the original surface neighbourhood. Its marked fibrewise integration maps identify it with $F'\times\Delta_z$ for one elliptic curve $F'$. The deck generator has the form
$(v,z)\mapsto(v+\tau+a(z),\zeta_m z)$, where $a(0)=0$. Its free central action makes $\tau$ a point of exact order $m$.

Lift $a$ to a small holomorphic complex-valued function. The finite-order identity gives $\sum_{j=0}^{m-1}a(\zeta_m^jz)=0$: the sum is lattice-valued and vanishes at zero. Put
$b(z)=-m^{-1}\sum_{j=0}^{m-1}j\,a(\zeta_m^jz)$.
Shifting the index in the sum gives $b(\zeta_mz)-b(z)=-a(z)$, so the holomorphic coordinate change $v'=v+b(z)$ removes $a$. Thus the local quotient is
\begin{equation}
 (F'\times\Delta_z)/\langle\delta\rangle,
 \qquad \delta(v,z)=(v+\tau,\zeta_m z),
 \qquad \operatorname{ord}(\tau)=m,
 \qquad t=z^m.
 \hypertarget{tag-00IG}{}\label{cdp:eq:C.EQ.3}
 \hypertarget{tag-00Q4}{}\hypertarget{tag-010W}{}\hypertarget{tag-00HE}{}
\SsixEquationAliases{\label{cdp:eq:C.KX.4}\label{cdp:eq:C.SF.2}\label{cdp:eq:C.EO.1}}
\end{equation}
We also use base-first coordinates $(u,w)$ on $\Delta_u\times F'$. The product permutation $(u,w)\mapsto(v=w,z=u)$ conjugates
$(u,w)\mapsto(\zeta_m u,w+\tau)$ to the displayed deck map and carries $t=u^m$ to $t=z^m$. In particular the base-first and fibre-first charts give the same finite \'etale root cover and the same quotient.

\item A base-fixing element of $H$ lifts functorially to the root chart. Its multiplier on $z$ is a power of $\zeta_m$; compose with the inverse power of $\delta$ to make the lift fix $z$. A power equal to the element's order is then a deck transformation fixing $z$, hence the identity. The lift acts on every smooth elliptic fibre by a translation of bounded finite order. Its translation parameter is holomorphic with values in a finite torsion subgroup, so is constant. The uniquely chosen $z$-fixing lifts therefore form a finite translation subgroup $T\subset F'$ and commute with $\delta$.

First take the quotient by $T$, and write $\bar F=F'/T$ and $m'=\operatorname{ord}(\bar\tau)$; thus $m'\mid m$. The remaining deck map is
$(\bar v,z)\mapsto(\bar v+\bar\tau,\zeta_mz)$.
Its $m'$th power generates the pure disc rotations of order $m/m'$. Quotienting them gives $w=z^{m/m'}$, with remaining action
\begin{equation}
 (\bar v,w)\longmapsto(\bar v+\bar\tau,\zeta_{m'}w),
 \qquad t=w^{m'}.
 \hypertarget{tag-00IH}{}\label{cdp:eq:C.EQ.4}
\end{equation}
That action is free on the central fibre, so $E_1$ is smooth and the full fibre is
$m'(\bar F/\langle\bar\tau\rangle)$.
These charts also identify its general fibres with torsors under the fixed quotient elliptic curve and leave their integral monodromy equal to the identity. The quotient by $H$ is taken over $B$, so its fibres remain connected. The definition of $H$ makes $Q$ faithful on $B$.

\item Let $g\in Q$ generate the stabilizer of a base point, with order $s$. Linearize its base coordinate so that $g_B(t)=\lambda t$, where $\lambda$ is primitive of order $s$. The linear part $\beta$ on the marked elliptic fibres is constant, since it takes values in the discrete group of origin-preserving automorphisms. Suppose $\beta\ne1$, and let $m$ be the full fibre multiplicity.

In the base-first chart from the first item, name the elliptic coordinate $z$ and the root coordinate $u$. Choosing an $m$th root $\rho$ of $\lambda$ gives the functorial lift
\begin{equation}
 \widetilde g(u,z)
       =(\rho u,\beta z+a(u)),\qquad
       \rho^m=\lambda .
 \hypertarget{tag-00Q9}{}\label{cdp:eq:C.KX.9}
\end{equation}
Conjugation by this lift preserves the root deck group. Its action on $u$ preserves the multiplier $\zeta_m$, so it fixes the deck generator. Restriction to $u=0$ gives
\begin{equation}
                         \beta\tau=\tau .
 \hypertarget{tag-00PX}{}\label{cdp:eq:C.KX.10}
\end{equation}
The order of $\beta$ divides $s$. Since $\beta\ne1$, the elliptic endomorphism
$1+\beta+\cdots+\beta^{s-1}$ is zero. Applied to $\tau$, this says $s\tau=0$, so $m\mid s$.

The induced affine map on $F'/\langle\tau\rangle$ has a fixed point: the nonzero endomorphism $1-\beta$ of that elliptic curve is an isogeny and hence surjective. Choose a lift of a fixed point and compose $\widetilde g$ with a deck transformation so that it fixes the chosen lift. Its $s$th power is a deck transformation fixing a central point, hence is the identity by freeness. Therefore $\rho^s=1$. Since $\rho^m=\lambda$ has exact order $s$, the order of $\rho$ is $s$ and $\gcd(m,s)=1$. Together with $m\mid s$ this proves
\begin{equation}
                 m=1\qquad(\beta\ne1).
 \hypertarget{tag-00II}{}\label{cdp:eq:C.EQ.5}
\end{equation}

\item For $\beta=1$, a point stabilizer acts on the reduced elliptic fibre by a translation fixing that point, hence fixes the whole fibre. Averaging local coordinates linearizes its finite action. In coordinates $(x,v)$ with $x$ normal to the fibre, it acts as $(x,v)\mapsto(\zeta_e x,v)$; an element with identity derivative is the identity in these averaged coordinates. The quotient coordinates $(x^e,v)$ prove that the quotient surface and reduced support are smooth. The full map is a power of the normal coordinate times a unit, so its fibre is a positive multiple of the elliptic support.

\item For $\beta\ne1$, the source fibre has multiplicity one by \eqref{cdp:eq:C.EQ.5}. In a marked product chart, after translating a fixed central point to the origin, the action is
$(v,z)\mapsto(\beta v+a(z),\zeta_s z)$ with $a(0)=0$.
Lift $a$ to a small holomorphic complex-valued function, and define
\[
 Pb(z)=\beta^{-1}b(\zeta_sz),\qquad c(z)=-\beta^{-1}a(z).
\]
Here $P^s=1$, and the finite-order identity gives $\sum_{k=0}^{s-1}P^kc=0$. Thus
$b=s^{-1}\sum_{k=0}^{s-1}kP^kc$ satisfies $(P-1)b=c$.
The holomorphic change $v'=v+b(z)$ makes
$a(z)+b(\zeta_sz)-\beta b(z)=0$ and conjugates the action to the product action.

Put $d=\operatorname{ord}(\beta)$. An origin-preserving elliptic automorphism acts on a rank-two integral lattice, with real trace $2\cos(2\pi r/d)$ an integer between $-2$ and $2$. Its nonidentity finite orders are therefore $2,3,4,6$, and the primitive exponents for these orders are $r=1,d-1$. The kernel on the elliptic factor is generated by the $d$th power of $g$ and acts by pure normal rotations. Its quotient coordinate is $z^{s/d}$. Naming that coordinate again $z$ gives
\begin{equation}
             (v,z)\longmapsto(\zeta_d^r v,\zeta_d z),
             \qquad r\in\{1,d-1\}.
 \hypertarget{tag-00IJ}{}\label{cdp:eq:C.EQ.6}
\end{equation}
The quotients by $H$ keep the base function $t$. At a $Q$-stabilizer of order $s$, the quotient base coordinate is $T=t^s$.
\end{enumerate}
\end{proof}

\begin{lemma}[Elliptic fibrations on minimal surface models]\hypertarget{tag-P318}{}\label{lem:II:elliptic-minimal-descent}
\StatementRecord{lem:II:elliptic-minimal-descent}{proved}{P318}{\cite[Chapter~1, Proposition~2.1, p.~8]{Peters1990}}
Let $Z$ be a smooth compact minimal complex surface with $\kappa(Z)\ge0$, and let $\rho:X\to Z$ be a proper birational morphism from a smooth surface. Let $q:X\to C$ be a proper holomorphic map to a smooth compact connected curve, with connected fibres and elliptic general fibre. There is a unique holomorphic map $p:Z\to C$ with $q=p\circ\rho$. It has connected fibres, and pullback gives an injection
$H^1(C,\OO_C)\hookrightarrow H^1(Z,\OO_Z)$.
A smooth compact complex surface of nonnegative Kodaira dimension carrying a relatively minimal elliptic fibration is minimal.
\end{lemma}
\begin{proof}
The canonical class of a minimal surface with nonnegative Kodaira dimension is nef by \cite[Chapter~1, Proposition~2.1, p.~8]{Peters1990}. Apply \cite[Proposition~3.12, p.~20]{PopescuPampu2016} to the proper holomorphic modification $\rho:X\to Z$ between the smooth complex surfaces in the statement, and factor it into point blowups. The one-step Jacobian formula has coefficient one on the exceptional curve, and iteration gives
$K_X=\rho^*K_Z+\sum_j a_jE_j$ with $a_j>0$ on every $\rho$-exceptional prime. For a smooth general elliptic fibre $F$ of $q$, the normal line is trivial, so adjunction gives
\begin{equation}
 0=K_X\cdot F
   =\rho^*K_Z\cdot F+\sum_j a_jE_j\cdot F,
                   \qquad a_j>0.
 \hypertarget{tag-00DL}{}\label{cdp:eq:C.CS.5}
\end{equation}
Every term is nonnegative, hence every $E_j$ is vertical for $q$. The fibres of $\rho$ are connected, and each is therefore sent to one point. This defines $p$ set-theoretically. On a coordinate neighbourhood in $C$, the holomorphic coordinate of $q$ descends through $\rho_*\OO_X=\OO_Z$; these functions give the holomorphic map $p$ and prove uniqueness. Its fibres are the images under $\rho$ of the connected fibres of $q$, so are connected. Consequently $p_*\OO_Z=\OO_C$, and the coherent Leray sequence gives the asserted injection.

For the last assertion, take the minimal contraction $\rho:X\to Z$ of the relatively minimal elliptic surface $q:X\to C$. The argument just given makes every $\rho$-exceptional curve vertical. A nonidentity contraction has a first point blowdown, whose exceptional $(-1)$-curve would therefore lie in a fibre of $q$. Relative minimality excludes that curve. Hence $\rho$ is an isomorphism.
\end{proof}

\begin{lemma}[Canonical characters over an elliptic base]\hypertarget{tag-P319}{}\label{lem:II:elliptic-canonical-character}
\StatementRecord{lem:II:elliptic-canonical-character}{proved}{P319}{\cref{lem:II:identity-monodromy-elliptic,lem:II:local-cyclic-elliptic,cdp:st:C.2.UN.28}}
Let $q:E\to B$ satisfy \cref{lem:II:identity-monodromy-elliptic}, with $B$ an elliptic curve. Write its full fibres as $m_bC_b$. A nonzero base differential and the fibre differential with fixed normalized periods define a generator $\Omega$ of $H^0(E,K_E)$ whose divisor is
$\sum_b(m_b-1)C_b$.
For a finite group $G$ acting on $E$ and $B$ with $q$ equivariant, let $\lambda(g)$ be its character on the base differential and $\beta(g)$ its linear character on the fixed elliptic fibre. Then
$g^*\Omega=\lambda(g)\beta(g)\Omega$.
For every smooth surface $M$ properly birational to $E/G$, the quotient and graph maps give an injection
$H^0(M,K_M)\hookrightarrow H^0(E,K_E)^G$.
\end{lemma}
\begin{proof}
Choose a nonzero differential $\omega_B$ on $B$. On the smooth base locus the marked integration charts of \cref{lem:II:identity-monodromy-elliptic} have translation transition maps. Their fibre differentials differ by base differentials, so $\omega_B\wedge dz$ glues. At a full fibre $mC$, use the base-first root chart of \cref{lem:II:local-cyclic-elliptic}, with root coordinate $u$ and elliptic coordinate $z$. Writing $\omega_B=c(t)\,dt$ gives
\begin{equation}
 \Omega=\omega_B\wedge dz
        =m\,c(u^m)u^{m-1}\,du\wedge dz ,
                  \qquad c(0)\ne0 .
 \hypertarget{tag-00Q7}{}\label{cdp:eq:C.KX.7}
\end{equation}
This form is deck-invariant, since the deck map multiplies $u^{m-1}du$ by $\zeta_m^m=1$ and acts on $z$ by a translation. The root projection is \'etale, so descent preserves the zero order $m-1$ along the reduced central curve. These charts prove the asserted divisor. By \cref{lem:II:identity-monodromy-elliptic}, $p_g(E)=g(B)=1$, so $\Omega$ spans the canonical forms.

The fibre linear part of $g$ is constant because it is a holomorphic function with values in the discrete group of origin-preserving automorphisms of the marked elliptic curve. Its translation term adds a base differential to $dz$. Thus pullback multiplies the wedge by $\lambda(g)\beta(g)$ on the smooth locus, and the equality extends over $E$.

To obtain the last map, resolve the graph of the meromorphic map $E\dashrightarrow M$ induced by the finite quotient and the given birational comparison. In the resulting diagram $E\xleftarrow{\mu}W\xrightarrow{\pi}M$, the map $\mu$ is a sequence of point blowups and $\pi$ is generically finite. The generic Jacobian makes $\pi^*:H^0(M,K_M)\to H^0(W,K_W)$ injective. The point-blowup identification of canonical forms, equivalently Hartogs extension across the finite images of the exceptional curves, identifies the latter space with $H^0(E,K_E)$. This is the same pullback map as in \eqref{cdp:eq:C.PG.4}--\eqref{cdp:eq:C.PG.5}, with these $E,W,M$. It is $G$-invariant on the dense quotient locus and therefore everywhere.
\end{proof}

\subsection{Principal elliptic torsors with zero Chern class}\label{subsec:II:zero-chern-elliptic-torsors}

\begin{lemma}[Principal elliptic torsors with zero Chern class]\hypertarget{tag-P330}{}\label{lem:II:zero-chern-elliptic-torsor}
\StatementRecord{lem:II:zero-chern-elliptic-torsor}{proved}{P330}{}
Let $B=\C/\Gamma$ and $F_1=\C/\Lambda_{F_1}$ be elliptic curves, and let $q:E_1\to B$ be a holomorphic principal $F_1$-torsor. Let $\mathcal F_1$ be the sheaf of holomorphic maps to $F_1$. Its analytic torsor class lies in $H^1(B,\mathcal F_1)$, and its integral Chern class is its image in $H^2(B,\Lambda_{F_1})$ under the connecting map of
\begin{equation}
 0\longrightarrow\Lambda_{F_1}
 \longrightarrow\OO_B
 \longrightarrow\mathcal F_1
 \longrightarrow0
                                                               \hypertarget{tag-00Q0}{}\label{cdp:eq:C.KX.13}
\hypertarget{tag-00WA}{}\SsixEquationAliases{\label{cdp:eq:C.PT.7}}
\end{equation}
Suppose this Chern class is zero. There is a number $a\in\C$ and an isomorphism of principal $F_1$-torsors over the given $B$
\begin{equation}
 E_1\simeq
 \C^2\big/
 \left\langle
 (z,w)\mapsto
       (z+\gamma,w+a\overline\gamma),\
 (z,w)\mapsto(z,w+\lambda)
 \right\rangle_{
       \gamma\in\Gamma,\,
       \lambda\in\Lambda_{F_1}} .
                                                               \hypertarget{tag-00Q1}{}\label{cdp:eq:C.KX.14}
\hypertarget{tag-00WB}{}
\hypertarget{tag-00TF}{}\SsixEquationAliases{\label{cdp:eq:C.PT.8}\label{cdp:eq:C.NS.14}}
\end{equation}
The map to $B$ is induced by $(z,w)\mapsto z$, and the $F_1$-action is translation of the second coordinate. In particular $E_1$ is a complex two-torus. The isomorphism preserves the full analytic torsor class.
\end{lemma}
\begin{proof}
Exactness of \eqref{cdp:eq:C.KX.13} lifts the torsor class to $H^1(B,\OO_B)$. The Dolbeault description of this one-dimensional space represents its lift by a constant form $a\,d\overline z$. On the universal cover its primitive is $a\overline z$, whose changes under $z\mapsto z+\gamma$ are the constants $a\overline\gamma$. These are the additive translation cocycle of the displayed quotient, and their images under $\OO_B\to\mathcal F_1$ represent the original principal torsor.

To identify the total spaces over $B$, compare their translation cocycles $g_{ij}^{\rm mod}$ and $g_{ij}^{E_1}$ on a common base cover. Equality of their classes means that $g_{ij}^{E_1}-g_{ij}^{\rm mod}=c_j-c_i$ for a holomorphic $\mathcal F_1$-valued zero-cochain $(c_i)$. In the convention $w_j=w_i+g_{ij}$, the local translations $w_i\mapsto w_i+c_i$ therefore glue a holomorphic map from the displayed model to $E_1$. The translations by $-c_i$ give its inverse. Both maps commute with the projection to $B$ and the $F_1$-action. A different lift in $H^1(B,\OO_B)$ has the same image and yields the same type of explicit isomorphism.

Choose $\Z$-bases of $\Gamma$ and $\Lambda_{F_1}$. The corresponding four period vectors are real-linearly independent: projection to the first coordinate first kills the coefficients of the two base vectors, and then the two fibre vectors have zero coefficients. Their integral span is therefore a discrete rank-four lattice in $\C^2$ with a compact fundamental parallelepiped. The quotient in \eqref{cdp:eq:C.KX.14} is a complex two-torus.
\end{proof}

\subsection{Elliptic fibre configurations of component surfaces}\label{subsec:II:elliptic-configuration-list}

\begin{proposition}[A properly elliptic normalization component has one of eleven positive-coherent-Euler fibre configurations]\hypertarget{tag-01DC}{}\label{cdp:st:C.2.UN.33}
\StatementRecord{cdp:st:C.2.UN.33}{proved}{01DC}{\cref{lem:II:cyclic-component-cover,cdp:st:C.2.UN.15,cdp:st:C.2.UN.22,cdp:st:C.2.UN.29,lem:II:local-cyclic-elliptic,lem:II:elliptic-minimal-descent,lem:II:elliptic-canonical-character,cdp:st:C.2.UN.28}; \cite[\S2.1, arXiv v2 pp.~3--4; \S2.3, pp.~5--6]{UlmerUrzua2021}; \cite[Theorem~6.2, pp.~564--566; \S8(iv)(11), p.~585]{Kodaira1963II}; \cite[Chapter~1, \S1, pp.~5--6]{Peters1990}}

Let $M$ be the smooth minimal model of a normalization
component of a fibre of an unrestricted two-torus degeneration,
and suppose $\kappa(M)=1$. Its connected elliptic Iitaka map
$p:M\to B_0$ is isotrivial, and $g(B_0)\le1$.
Its coherent Euler characteristic belongs to $\{0,1,2\}$.
When it is zero, every fibre of $p$ is a positive multiple
of a smooth elliptic curve. When it is positive,
$B_0\simeq\PP^1$, $q(M)=0$, and the complete configuration
of the fibres of $p$ with singular reduced support is one
of the following rows:
\begin{equation}
\begin{array}{c|l|c}
\chi(\OO_M)&\text{singular reduced fibres of }p&e(M)\\ \hline
1&\mathrm{II}+\mathrm{II}^*&12\\
1&\mathrm{III}+\mathrm{III}^*&12\\
1&\mathrm{IV}+\mathrm{IV}^*&12\\
1&2\mathrm I_0^*&12\\
1&3\mathrm{IV}&12\\
1&2\mathrm{III}+\mathrm I_0^*&12\\
1&\mathrm{II}+\mathrm{IV}+\mathrm I_0^*&12\\ \hline
2&4\mathrm I_0^*&24\\
2&3\mathrm{IV}^*&24\\
2&2\mathrm{III}^*+\mathrm I_0^*&24\\
2&\mathrm{II}^*+\mathrm{IV}^*+\mathrm I_0^*&24 .
\end{array}                                                   \hypertarget{tag-00IC}{}\label{cdp:eq:C.EQ.1}
\end{equation}
In each row $p_g(M)=\chi(\OO_M)-1$. The table lists the nonsmooth reduced fibres; $p$ may also have smooth multiple fibres.
\end{proposition}


\begin{proof}

\begin{enumerate}[label=\arabic*.,leftmargin=*]
\item Take the component cover $E'\to S^\nu$ and its cyclic stabilizer $G$ from \cref{lem:II:cyclic-component-cover}. Resolve its generically finite comparison with $M$ as $E'\leftarrow W\to M$. Injective pluricanonical pullback and the finite-base-change application of \cref{cdp:st:C.2.UN.15} give $\kappa(E')=1$. The relatively minimal elliptic model of \cref{cdp:st:C.2.UN.22,cdp:st:C.2.UN.29} is minimal by \cref{lem:II:elliptic-minimal-descent}; call it $E$. The birational map from $E'$ is its minimal contraction, so it is the same model in \eqref{cdp:eq:C.HQ.3}. It has the connected elliptic Iitaka map
\begin{equation}
 q:E\longrightarrow B,\qquad g(B)\le1,                 \hypertarget{tag-00IF}{}\label{cdp:eq:C.EQ.2}
\end{equation}
Its smooth fibres have a fixed elliptic modulus and identity integral monodromy, and every full fibre is a smooth multiple. The action of $G$ descends to $E$ by \eqref{cdp:eq:C.HQ.3} and preserves $q$, because the connected-fibre Iitaka map is determined by the pluricanonical field. The same quotient comparison identifies $E/G$ properly birationally with $M$.

The quotient and Stein maps of \cref{cdp:st:C.2.UN.20} give a fibration over $B/G$ on a smooth birational model of $M$. Its general fibre is the quotient of a general fibre of $q$ by the kernel of the base action. A rational general fibre would force every positive pluricanonical section to vanish by adjunction, contradicting $\kappa(M)=1$. Thus that general fibre is elliptic. Apply \cref{lem:II:elliptic-minimal-descent} to a common smooth resolution over the minimal surface $M$; it descends the map to $p:M\to B/G$. Its smooth general fibre is unchanged. This is the connected-fibre Iitaka map, since the pluricanonical ratios factor through these elliptic fibres and their nonconstant image has dimension one. Set $B_0=B/G$; Riemann--Hurwitz gives $g(B_0)\le g(B)\le1$.

\item Let $H=\ker(G\to\operatorname{Aut}(B))$. Every element of $H$ acts on a general elliptic fibre by a translation. Indeed, an affine elliptic automorphism with nonidentity linear part has a fixed point, since its fixed-point equation is an isogeny equation. Riemann--Hurwitz then makes the quotient by any group containing it rational. The resulting general quotient fibre would contradict the preceding pluricanonical restriction argument. The translations have bounded finite order, so form a locally constant subgroup of the fixed elliptic curve.

Apply \cref{lem:II:local-cyclic-elliptic} to this map $q:E\to B$ and its cyclic group $G$. Its root chart \eqref{cdp:eq:C.EQ.3} and the quotient map \eqref{cdp:eq:C.EQ.4} construct the smooth surface $E_1=E/H$ over the same $B$, including every full multiple fibre. The remaining group $Q=G/H$ acts faithfully on $B$. At a point with stabilizer of order $s$, a nonidentity fibre character $\beta$ gives multiplicity one by \eqref{cdp:eq:C.EQ.5}; the primitive product action is \eqref{cdp:eq:C.EQ.6}. A trivial fibre character gives a smooth quotient surface and smooth elliptic reduced support. Those values have Euler contribution zero.

\item Write the fixed elliptic curve as
$y^2=x^3+Ax+B$, so multiplication by $\zeta_d^r$
acts on $x,y$ by $\zeta_d^{-2r},\zeta_d^{-3r}$.
Over the meromorphic function field of the base, the invariant functions
\begin{equation}
          T=z^d,\qquad X=z^{2r}x,\qquad Y=z^{3r}y
                                                               \hypertarget{tag-00IK}{}\label{cdp:eq:C.EQ.7}
\end{equation}
generate the quotient function field. Indeed, adjoining $z$ to their field recovers $x=z^{-2r}X$ and $y=z^{-3r}Y$, and $z^d=T$. Thus the extension has degree at most $d$. The faithful cyclic action of order $d$ gives fixed-field degree $d$, so the invariant field is exactly the displayed one. They give the following minimal Weierstrass equations:
\begin{equation}
\begin{array}{c|c|c|c}
d&r&\text{equation}&\text{Kodaira type}\\ \hline
2&1&Y^2=X^3+AT^2X+BT^3&\mathrm I_0^*\\
3&1&Y^2=X^3+BT^2&\mathrm{IV}\\
3&2&Y^2=X^3+BT^4&\mathrm{IV}^*\\
4&1&Y^2=X^3+ATX&\mathrm{III}\\
4&3&Y^2=X^3+AT^3X&\mathrm{III}^*\\
6&1&Y^2=X^3+BT&\mathrm{II}\\
6&5&Y^2=X^3+BT^5&\mathrm{II}^* .
\end{array}                                                   \hypertarget{tag-00IL}{}\label{cdp:eq:C.EQ.8}
\end{equation}
For $d=3,6$ the curve has $A=0$ and $B\ne0$;
for $d=4$ it has $B=0$ and $A\ne0$.
The equations are minimal: for every row, the coefficient orders are below a simultaneous division by $T^4$ and $T^6$ in the Weierstrass rescaling. The positive $T$-loop lifts to $z\mapsto\zeta_d z$. In the quotient, its endpoint $(v,\zeta_d z)$ is identified with $(\zeta_d^{-r}v,z)$, so its action on the fixed fibre is multiplication by $\zeta_d^{-r}$. This is the inverse of the deck character in \eqref{cdp:eq:C.EQ.6}. In the convention of \cite[\S2.1, arXiv v2 pp.~3--4]{UlmerUrzua2021}, a fixed period $\tau$ and a matrix $\left(\begin{smallmatrix}a_\gamma&b_\gamma\\c_\gamma&d_\gamma\end{smallmatrix}\right)$ send the row cycle $(m,n)$, representing $m\tau+n$, to a cycle representing $(c_\gamma\tau+d_\gamma)(m\tau+n)$; the uniformizing deck multiplier is $(c_\gamma\tau+d_\gamma)^{-1}$. Thus the displayed endpoint identification gives exactly the row-cycle monodromy in \cite[\S2.3, arXiv v2 pp.~5--6]{UlmerUrzua2021}. Its finite-order rows, together with the minimal Weierstrass equations, distinguish the starred and unstarred types in \eqref{cdp:eq:C.EQ.8}; Kodaira's fibre classification is \cite[Theorem~6.2, pp.~564--566]{Kodaira1963II}.
For the displayed minimal equations the discriminant orders are computed
directly as $12r/d$, and the corresponding Kodaira configurations have
the same Euler numbers.
The map \eqref{cdp:eq:C.EQ.7} is birational over the base
after taking the finite quotient. Resolving its
graph and contracting vertical exceptional curves
gives the same relatively minimal elliptic model;
therefore these are the fibre types of the quotient.

\item On $B=\PP^1$, a nontrivial cyclic action is conjugate
to $z\mapsto\zeta_nz$. Its only branch values are
zero and infinity, where positive local generators
are inverse. A nontrivial fibre character of order
$d$ consequently gives the two opposite exponents
in \eqref{cdp:eq:C.EQ.8}. The resulting pairs are
\begin{equation}
 (\mathrm{II},\mathrm{II}^*),\quad
 (\mathrm{III},\mathrm{III}^*),\quad
 (\mathrm{IV},\mathrm{IV}^*),\quad
 (\mathrm I_0^*,\mathrm I_0^*).                         \hypertarget{tag-00IM}{}\label{cdp:eq:C.EQ.9}
\end{equation}
Every pair has Euler sum twelve. A trivial fibre
character gives only smooth multiples.

\item On an elliptic $B$, a finite cyclic group either
acts by translations, in which case it is free,
or has a fixed point. In the latter case translate
that point to the origin. The action is linear of
order $n=2,3,4,6$, as follows from its action on
the rank-two integral period lattice. The branch
indices of $B\to B/Q\simeq\PP^1$ are
\begin{equation}
\begin{array}{c|c}
n&\text{branch indices}\\ \hline
2&(2,2,2,2)\\
3&(3,3,3)\\
4&(4,4,2)\\
6&(6,3,2).
\end{array}                                                   \hypertarget{tag-00ID}{}\label{cdp:eq:C.EQ.10}
\end{equation}
The number fixed by a nontrivial power is $|1-\zeta_n^k|^2$, the determinant of
that power minus one on the lattice. Removing
the points with larger stabilizer and dividing
by the corresponding orbit sizes gives exactly
\eqref{cdp:eq:C.EQ.10}. Each row also satisfies the degree-$n$
Riemann--Hurwitz equation.

\item Write the fibre character as $\zeta_n^a$,
$0\le a<n$. At a branch index $e$, its exponent
is $a$ modulo $e$. Exponent zero gives a smooth
multiple. Otherwise put $c=\gcd(a,e)$; quotienting
the normal reflection kernel gives the entry of
\eqref{cdp:eq:C.EQ.8} with $d=e/c$ and $r=(a\bmod e)/c$.
Enumerating these characters gives
\begin{equation}
\begin{array}{c|c|l}
n&a&\text{nonsmooth fibre configuration}\\ \hline
2&1&4\mathrm I_0^*\\
3&1&3\mathrm{IV}\\
3&2&3\mathrm{IV}^*\\
4&1&2\mathrm{III}+\mathrm I_0^*\\
4&2&2\mathrm I_0^*\\
4&3&2\mathrm{III}^*+\mathrm I_0^*\\
6&1&\mathrm{II}+\mathrm{IV}+\mathrm I_0^*\\
6&2,4&\mathrm{IV}+\mathrm{IV}^*\\
6&3&2\mathrm I_0^*\\
6&5&\mathrm{II}^*+\mathrm{IV}^*+\mathrm I_0^* .
\end{array}                                                   \hypertarget{tag-00IE}{}\label{cdp:eq:C.EQ.11}
\end{equation}
The zero character gives no nonsmooth fibre.


\item A free base orbit has a neighbourhood in the quotient isomorphic to the neighbourhood of one fibre in the orbit. Since the source fibres are smooth multiples, the preceding stabilizer list exhausts the nonsmooth reduced fibres. A free cyclic translation action on an elliptic base also leaves only smooth multiples. Positive Euler number therefore occurs over $B_0=\PP^1$ with nontrivial fibre linear character.

Apply \cref{lem:II:identity-monodromy-elliptic} to $q_1:E_1\to B$; it gives $p_g(E_1)=g(B)$. The quotient comparison $E_1/Q=E/G$ with $M$ gives the injective pullback of canonical forms through a resolved graph as in \eqref{cdp:eq:C.PG.4}--\eqref{cdp:eq:C.PG.5}. For $B=\PP^1$, it follows that $p_g(M)=0$.

For elliptic $B$, apply \cref{lem:II:elliptic-canonical-character} with $(E,q,G)=(E_1,q_1,Q)$. Choose the base generator with derivative $\zeta_n$ and write the fibre character as $\zeta_n^a$, as in \eqref{cdp:eq:C.EQ.11}. The canonical character is then $\zeta_n^{a+1}$. The lemma's exact graph pullback embeds $H^0(M,K_M)$ in its invariant subspace. Hence $p_g(M)=0$ when $a+1\not\equiv0\pmod n$. Every Euler-twelve row of \eqref{cdp:eq:C.EQ.11} has this nontrivial character.

\item On a relatively minimal elliptic surface the canonical
divisor is vertical. A meromorphic differential
on the base and a meromorphic relative differential
give such a divisor; restriction to a general compact
elliptic fibre has no zero or pole. The components
of the displayed reducible singular fibres are $(-2)$
rational curves. A type $\mathrm{II}$ fibre has arithmetic
genus one and square zero, and a smooth reduced multiple
fibre is an elliptic curve of square zero. Adjunction gives
$K_M\cdot C=0$ for each of them, and hence $K_M^2=0$.
Analytic Riemann--Roch \cite[Chapter~1, \S1, pp.~5--6]{Peters1990} gives $12\chi(\OO_M)=e(M)$. Euler localization for the
elliptic map $p$ and \eqref{cdp:eq:C.EQ.8}--\eqref{cdp:eq:C.EQ.11} give
exactly the rows of \eqref{cdp:eq:C.EQ.1}, or zero. For Euler
twelve we proved $p_g=0$, so $\chi=1$ gives $q=0$.
For Euler twenty-four, \ref{cdp:st:C.2.UN.28} gives $p_g\le1$, while
$\chi=2=1-q+p_g$ and $q\ge0$ force $p_g=1,q=0$.
In the zero case
the local list has no nonsmooth reduced fibre,
since every such entry has positive Euler number.


\end{enumerate}
\end{proof}

\subsection{K3 components and cyclic quotients}\label{subsec:II:k3-cyclic-components}

\begin{proposition}[K3 components, cyclic quotients and elliptic pencils]
\hypertarget{tag-01DB}{}\label{cdp:st:C.2.UN.32}
\StatementRecord{cdp:st:C.2.UN.32}{proved}{01DB}{\cref{lem:II:cyclic-component-cover,cdp:st:C.2.UN.22,cdp:st:C.2.UN.29,cdp:st:C.2.UN.23,cdp:st:C.2.UN.28,lem:II:local-cyclic-elliptic,lem:II:elliptic-minimal-descent,lem:II:elliptic-canonical-character,cdp:st:C.2.UN.4}; \cite[Chapter~1, Theorem~3.1, p.~10]{Peters1990}; \cite[\S3.1, p.~7659; \S4.1, Theorem~27, p.~7663]{Roulleau2019GeneralizedKummer}; \cite[Theorem~II, Corollary~II, p.~201]{Lagrange1770}; \cite[\S3, Corollary~3, p.~560; \S6, Theorem~1(3) and its corollary, pp.~573--574]{PiatetskiShapiroShafarevich1971}; \cref{lem:II:zero-chern-elliptic-torsor}}
\hypertarget{tag-01D2}{}\label{cdp:st:C.2.UN.24}
\StatementRecord{cdp:st:C.2.UN.24}{proved}{01D2}{}

The K3 minimal birational classes occurring as
normalization components of proper two-torus
degenerations are exactly the minimal resolutions
of cyclic quotients of complex two-tori by
nontrivial symplectic automorphisms whose generator has no eigenvalue one on the tangent space.
After a translation of the origin, the generator
is linear of order $2,3,4$, or $6$.
Write $n$ for its order.
This criterion holds in every algebraic dimension.

More explicitly, let $f:U\to\Delta$ be a proper holomorphic map
from a smooth complex threefold, with connected two-torus fibres
over $\Delta^*$. Let $S_i$ be an irreducible component of its
central support, and let $M$ be the K3 minimal model of a
resolution of $S_i^\nu$.
For $a(M)>0$ there is an elliptic fibration
$M\to\mathbb P^1$. For $a(M)=0$ the cyclic quotient in the
criterion has algebraic dimension zero as well.
In every algebraic dimension its quotient singularities
and exceptional configurations are
\begin{equation}
\begin{array}{c|c|c|c}
n&\text{quotient singularities}&
 \text{number of exceptional curves}&e(A/\langle g\rangle)\\ \hline
2&16A_1&16&8\\
3&9A_2&18&6\\
4&4A_3+6A_1&18&6\\
6&A_5+4A_2+5A_1&18&6.
\end{array}                                                    \hypertarget{tag-00OY}{}\label{cdp:eq:C.K3.1}
\end{equation}

Every such cyclic quotient K3 minimal class is realized by a
compact smooth threefold over $\mathbb P^1$ with connected torus
general fibre and with both special reduced supports of Euler
characteristic zero. A central component has normalization
resolution with that K3 minimal model.
The total space and construction are projective when the given
torus is projective.
\end{proposition}


\begin{proof}
\begin{enumerate}[label=\arabic*.,leftmargin=*]
\item Use the normalized root base change and equivariant resolution of \cref{lem:II:cyclic-component-cover}. For the chosen component, let $E$
be a smooth strict covering component and $G$ its cyclic
stabilizer. The comparison is
\begin{equation}
 E\longrightarrow S_i^\nu,
 \qquad E/G\dashrightarrow S_i^\nu\dashrightarrow M,
                                                               \hypertarget{tag-00P0}{}\label{cdp:eq:C.K3.2}
\end{equation}
where the first map is generically finite, the quotient
comparison is proper birational, and the remaining arrows
are resolved by their proper graphs. In particular, a
nonzero holomorphic two-form on $M$ pulls back, by its
Jacobian map, to a nonzero holomorphic two-form on
a resolution of the graph in \eqref{cdp:eq:C.K3.2}. Point-blowup
invariance then gives
\begin{equation}
                  p_g(E)\ge1,\qquad \kappa(E)\ge0.
                                                               \hypertarget{tag-00P1}{}\label{cdp:eq:C.K3.3}
\end{equation}
The function-field clause of \cref{lem:II:cyclic-component-cover}, applied to the same generically finite map $E\to S_i^\nu$, gives $a(E)=a(S_i^\nu)=a(M)$. Its monic-polynomial map uses elementary symmetric functions along precisely these finite generic sheets; the last equality is the birational identification through \eqref{cdp:eq:C.K3.2}.
\item Composing the proper birational quotient comparison in \eqref{cdp:eq:C.K3.2} with the birational map to the minimal surface gives
\begin{equation}
                         E/G\dashrightarrow M
                                                               \hypertarget{tag-00PW}{}\label{cdp:eq:C.KX.1}
\end{equation}
as a proper birational comparison after
resolving the graph. Pullback of the K3
two-form is injective, whereas the
compact-core bound \ref{cdp:st:C.2.UN.28} gives $p_g(E)\le1$.
Thus
\begin{equation}
                       p_g(E)=1,\qquad\kappa(E)\ge0.
                                                               \hypertarget{tag-00Q2}{}\label{cdp:eq:C.KX.2}
\end{equation}
The compact-component result in \cref{cdp:st:C.2.UN.15} applies to the resolved finite-base-change family and gives $\kappa(E)\le1$. Pass to its unique smooth minimal model $E_0$; \eqref{cdp:eq:C.HQ.3} of \cref{lem:II:cyclic-component-cover} descends the finite action through the point blowdowns and compares $E/G$ properly birationally with $E_0/G$.
If $E_0$ is a torus, \ref{cdp:st:C.2.UN.23} already gives
the required cyclic K3 quotient.
\item The two-dimensional space of fields on $E$ in \ref{cdp:st:C.2.UN.16}
has generic rank one or two. In rank one, \ref{cdp:st:C.2.UN.20} gives
a rational or elliptic fibration on a birational
model of $M$. The rational alternative would make all
positive pluricanonical sections vanish, by restricting
them to its smooth general fibre. The pulled-back K3
form excludes it. Hence the quotient fibration is elliptic.
\item In rank two, \ref{cdp:st:C.2.UN.22} gives $\chi(\mathcal O_E)=0$ and
its precise torus or elliptic-fibration alternatives.
First suppose $a(M)=a(E)=1$. In the elliptic alternative,
the connected-fibre map $q:E'\to B$ has
\begin{equation}
                     \mathcal M(E')=q^*\mathcal M(B).
                                                               \hypertarget{tag-00P2}{}\label{cdp:eq:C.K3.4}
\end{equation}
This is \eqref{cdp:eq:C.RF.9} for $X=E'$ and the connected-fibre map $q$: the finite-valued restriction to each connected general fibre and the proper graph in $B\times\PP^1$ identify the descended field with $\mathcal M(B)$. The equality makes this field
intrinsic and therefore $G$-invariant. Resolve the graph
of $q$ equivariantly and take its Stein factor. The group
acts on that factor, and the finite quotient is a
map to $B/G$ whose general fibre is $F/K$, with
$K=\ker(G\to\operatorname{Aut}B)$. This is the same
quotient and Stein comparison proved in \ref{cdp:st:C.2.UN.20}.
Its general fibre is connected and rational or elliptic;
the K3 form again excludes the rational case. In the
torus alternative the minimal model of $E$ is a torus
and the cyclic quotient argument below applies.
\item Now suppose $a(M)=a(E)=2$. We work on the smooth
projective minimal model $E_0$, to which the finite
action descends. In rank two, \eqref{cdp:eq:C.K3.3}, $\chi(\mathcal O_{E_0})=0$ and the preceding finite-base-change bound $\kappa(E_0)\le1$ give the following cases of the compact surface classification \cite[Chapter~1, Theorem~3.1, p.~10]{Peters1990}. If $\kappa(E_0)=0$, it is an Abelian
surface: K3 and Enriques surfaces have coherent Euler
characteristic $2$ and $1$, and bielliptic surfaces have
$p_g=0$. If $\kappa(E_0)=1$, its elliptic Iitaka
fibration is intrinsic. More explicitly, $G$ acts on
each space of pluricanonical sections, so a sufficiently
large pluricanonical map, its equivariant graph resolution,
and its Stein factor are all $G$-equivariant. The
connected general fibres are elliptic. Taking the
quotient gives the fibration with general fibre
$F/K$ just described. It is elliptic because its total
quotient has the K3 minimal model. This gives the equivariant elliptic quotient map in the projective non-torus case.
\item On a common smooth resolution, the elliptic quotient map just obtained has the form
\begin{equation}
                  M\xleftarrow{\rho}X\xrightarrow{q}B
                                                               \hypertarget{tag-00P3}{}\label{cdp:eq:C.K3.5}
\end{equation}
with connected elliptic general fibre. Apply \cref{lem:II:elliptic-minimal-descent} with $Z=M$ and the displayed $\rho,q$. Since $M$ is minimal with $K_M=0$, the lemma descends it to a holomorphic map on $M$. Its injection
$H^1(B,\OO_B)\hookrightarrow H^1(M,\OO_M)=0$
gives $B\simeq\PP^1$.

\item In the non-torus cases, \cref{cdp:st:C.2.UN.22,cdp:st:C.2.UN.29} give a relatively minimal elliptic model. Its Kodaira dimension is nonnegative by \eqref{cdp:eq:C.K3.3}, so \cref{lem:II:elliptic-minimal-descent} makes its minimal contraction an isomorphism. The birational comparison with the unique minimal model $E_0$ is therefore an isomorphism, giving the elliptic map
\begin{equation}
                   q:E_0\longrightarrow B
                                                               \hypertarget{tag-00Q3}{}\label{cdp:eq:C.KX.3}
\end{equation}
with identity integral monodromy and
only smooth multiple fibres. Its
relative Jacobian is one fixed elliptic
curve $F$. Apply \cref{lem:II:identity-monodromy-elliptic} to the displayed $q:E_0\to B$. Its full-base direct-image and Leray calculation \eqref{cdp:eq:C.RG.13} gives $p_g(E_0)=g(B)$. Equation \eqref{cdp:eq:C.KX.2} makes $B$ elliptic.
The map can be taken $G$-equivariant.
In the rank-one branch this is \ref{cdp:st:C.2.UN.20}.
In algebraic dimension one, the field identity \eqref{cdp:eq:C.K3.4} makes its connected-fibre base intrinsic. In algebraic dimension
two a non-torus minimal model
with $p_g=1$ and $\chi(\OO)=0$
has Kodaira dimension one, so
its elliptic map is its intrinsic
Iitaka fibration, by the equivariant pluricanonical
comparison above.
In algebraic dimension zero, the limiting field space on the covering component $E$ has generic rank two: two complex-linearly independent fields of generic rank one would satisfy $V_2=hV_1$ with a nonconstant meromorphic function $h$, contrary to $a(E)=0$. The algebraic-dimension-zero branch of \cref{cdp:st:C.2.UN.22} therefore makes its minimal model $E_0$ a torus.

\item Let $H$ be the kernel of the
action of $G$ on $B$.
Every element of $H$ acts
on a general elliptic fibre
by a translation. Indeed
a nonidentity linear part
would make the general
finite quotient fibre
rational. The quotient and Stein map
would then give a rational
fibration on a birational
model of $M$, contradicting
its nonzero canonical form
by restriction to a general
rational fibre.
The translation subgroup
on the fixed $F$ is
constant: its points have
bounded finite order and
belong to a discrete
finite torsion set.
\item Apply \cref{lem:II:local-cyclic-elliptic} to $q:E_0\to B$, $G$ and the translation kernel $H$ just established. It constructs
\begin{equation}
              E_1=E_0/H,\qquad Q=G/H
                                                               \hypertarget{tag-00Q5}{}\label{cdp:eq:C.KX.5}
\end{equation}
as a smooth surface over $B$, with fixed quotient elliptic curve and identity integral monodromy. At a full fibre $mC$, its exact root-quotient map is \eqref{cdp:eq:C.EQ.4}. With $e=m/m'$, the pure normal rotation quotient has local coordinate $x\mapsto x^e$, and the full fibre equation becomes $t=(x^e)^{m/e}$ with $e\mid m$. Thus the full reduced support is smooth elliptic and the multiplicity is the one supplied by that quotient map. The remaining group $Q$ acts faithfully on $B$.

\item Thus $q_1:E_1\to B$ is
again a smooth-multiple
elliptic fibration with
constant Jacobian.
The full-base calculation of \cref{lem:II:identity-monodromy-elliptic}, applied to $q_1:E_1\to B$ and $g(B)=1$, gives
\begin{equation}
 p_g(E_1)=1,\quad q(E_1)=2,\quad
                 \chi(\OO_{E_1})=0 .
                                                               \hypertarget{tag-00Q6}{}\label{cdp:eq:C.KX.6}
\end{equation}
The quotient
$E_1/Q=E_0/G$ remains
birational to $M$.
The group $Q$ acts
faithfully on the base.
Its quotient fibration
has connected elliptic
general fibre. Applying \cref{lem:II:elliptic-minimal-descent} through \eqref{cdp:eq:C.K3.5} makes its base $B/Q\simeq\PP^1$.
Consequently $Q$ is
nontrivial and its
action on $B$ is not
a group of translations.
Write its generator as $b\mapsto Lb+c$ on $B$. Since $L\ne1$, the isogeny $1-L$ is surjective; choose $b_0$ with $(1-L)b_0=c$. Translation by $b_0$ conjugates the generator to $L$.
The derivative character
$\lambda:Q\to\C^*$ is
faithful.
\item Apply \cref{lem:II:elliptic-canonical-character} with $(E,q,G)=(E_1,q_1,Q)$. The base is elliptic, and \eqref{cdp:eq:C.KX.6} gives the unique canonical form. The lemma constructs it by \eqref{cdp:eq:C.KX.7}, with divisor $\sum_b(m_b-1)C_b$, and gives its character $\lambda(g)\beta(g)$.

The proper birational comparison $E_1/Q=E_0/G\dashrightarrow M$ gives the lemma's injective pullback of the nonzero K3 form into this one-dimensional canonical space. The image is invariant, so
\begin{equation}
                   \lambda(g)\beta(g)=1
                       \quad(g\in Q).                \hypertarget{tag-00Q8}{}\label{cdp:eq:C.KX.8}
\end{equation}
The faithful character $\lambda$ consequently makes $\beta$ faithful as well.

\item Let a base point have stabilizer of order $s>1$. By \eqref{cdp:eq:C.KX.8} and the faithfulness of $\lambda$, its fibre character $\beta=\lambda^{-1}$ is nonidentity and has order $s$. Apply the stabilizer clause of \cref{lem:II:local-cyclic-elliptic} to $q_1:E_1\to B$. Its lift \eqref{cdp:eq:C.KX.9}, deck commutation \eqref{cdp:eq:C.KX.10} and multiplicity conclusion \eqref{cdp:eq:C.EQ.5} apply to this same full fibre and give
\begin{equation}
          \text{a branched base value has }m=1 .
                                                               \hypertarget{tag-00PY}{}\label{cdp:eq:C.KX.11}
\end{equation}

\item At a base point
with trivial
stabilizer, the
quotient is locally
an isomorphism
on a neighbourhood
of one fibre in
the orbit. Its
smooth elliptic
support therefore
retains multiplicity
$m$ and the zero
order $m-1$ of
\eqref{cdp:eq:C.KX.7}.
On a common smooth
resolution of
$E_1/Q$ and $M$,
the strict transform
of this curve is
elliptic and cannot
be contracted by
a birational map
of smooth surfaces:
all curves
exceptional for
successive point
blowdowns are
rational.
At its generic
point the pullback
of the nowhere-zero
K3 form must
therefore have
zero order zero.
Its equality with
the descended form
\eqref{cdp:eq:C.KX.7} forces
$m-1=0$.
This treats every
free base orbit.
We have proved
that $q_1$ has
no multiple fibre
at all.
\item The smooth map
$q_1:E_1\to B$
is now a principal
bundle under its
fixed elliptic
curve $F_1$.
Its Chern vector is
\begin{equation}
        c\in H^2(B,\Lambda_{F_1})
                  \simeq\Lambda_{F_1}.
                                                               \hypertarget{tag-00PZ}{}\label{cdp:eq:C.KX.12}
\end{equation}
A bundle
automorphism
covering $g_B$
and having
fibre linear
part $\beta$
gives
$g_B^*c=\beta_*c$.
For the precise
map, lift the
translation
transition
functions on
a base cover
to $\C$.
Their triple
differences are
the lattice
cocycle for $c$.
Conjugating the
trivializations
multiplies that
cocycle by $\beta$,
up to a lattice
coboundary; this
is the displayed
equality.
The holomorphic
automorphism of
the elliptic
base has degree
one, so it acts
as the identity
on $H^2(B,\Z)$.
Thus $c=\beta c$.
For a generator,
$\beta\ne1$, and
$\beta-I$ is
invertible on
$\Lambda_{F_1}
\otimes\R$.
The lattice is
torsion-free, so
$c=0$.
\item Apply \cref{lem:II:zero-chern-elliptic-torsor} to the principal map $q_1:E_1\to B$, its fixed curve $F_1$ and the Chern class $c=0$ just proved in \eqref{cdp:eq:C.KX.12}. Its isomorphism \eqref{cdp:eq:C.KX.14} uses this same base and fibre lattice, and preserves the full analytic torsor class through its local translation cochain. It therefore identifies this very $E_1$ with a complex two-torus over $B$.
\item The group $Q$ is
cyclic and
$E_1/Q=E_0/G$
has the birational
comparison
\eqref{cdp:eq:C.KX.1} to $M$.
\ref{cdp:st:C.2.UN.23} applies
to this torus,
and gives
exactly its
symplectic K3
quotient case.
Thus the torus used below is $A=E_0$ in the initial torus
case and $A=E_1=E_0/H$ in the elliptic-fibration case. In the
latter case its cyclic quotient is the same quotient
$E_1/Q=E_0/G$ through the finite maps already constructed.
\item By
\ref{cdp:st:C.2.UN.23}, either torus $A$ just obtained has a
holomorphic cyclic action whose quotient is
birational to $M$. The K3 alternative there says that
its generator $g(z)=Lz+b$ has no eigenvalue one and
$\det L=1$. Since $I-L:A\to A$ is an isogeny, the
equation $(I-L)p=b$ has a solution. Translation by
$p$ conjugates $g$ to the linear automorphism $L$.
Write its eigenvalues as $\lambda,\lambda^{-1}$,
with $\lambda$ a primitive $n$th root. The integral
action on $H_1(A,\mathbb Z)$ has characteristic polynomial
\begin{equation}
 (x^2-sx+1)^2
   =x^4-2sx^3+(s^2+2)x^2-2sx+1,
 \qquad s=\lambda+\lambda^{-1}\in[-2,2].
                                                               \hypertarget{tag-00P4}{}\label{cdp:eq:C.K3.6}
\end{equation}
Its integral coefficients imply $2s\in\mathbb Z$
and $s^2\in\mathbb Z$. Writing $s=k/2$ shows that
$k$ is even, hence $s\in\mathbb Z$. The value $s=2$
is excluded by the absence of eigenvalue one.
Thus $s=-2,-1,0,1$, giving precisely
$n=2,3,4,6$.
\item For every nonidentity power, its fixed points are the
kernel of an isogeny and their exact number is
\begin{equation}
 \#\operatorname{Fix}(g^k)
 =\det_{\mathbb R}(I-L^k)
 =|\det_{\mathbb C}(I-L^k)|^2
 =\bigl(2-2\cos(2\pi k/n)\bigr)^2.
                                                               \hypertarget{tag-00P5}{}\label{cdp:eq:C.K3.7}
\end{equation}
At a point with stabilizer of order $d$, the local
action is $(u,v)\mapsto(\zeta_du,\zeta_d^{-1}v)$.
Its invariant coordinates $x=u^d,y=v^d,z=uv$ give
exactly $xy=z^d$. The height-one toric resolution
in \eqref{cdp:eq:C.CQ.11} gives its chain of $d-1$ rational
$(-2)$-curves. Formula \eqref{cdp:eq:C.K3.7} now gives the table
by subtracting points with larger stabilizers and
dividing by the remaining orbit sizes. For $n=4$
there are $4$ full-stabilizer points and
$(16-4)/2=6$ order-two orbits. For $n=6$ the
counts are $1$, $(9-1)/2=4$, and
$(16-1)/3=5$ for stabilizer orders $6,3,2$.
The order-two and order-three cases have $16$ and
$9$ fixed points respectively. The finite quotient
Euler formula gives the last column; adding the
exceptional ranks gives $24$ in every row.
The configurations in \eqref{cdp:eq:C.K3.1} are the cyclic rows of the generalized Kummer classification \cite[\S3.1, p.~7659; \S4.1, Theorem~27, p.~7663]{Roulleau2019GeneralizedKummer}. The explicit fixed-point and quotient maps above give the singularity, exceptional-rank and quotient-Euler data for the component's cyclic action.
\item If this quotient K3 $M$ is projective, the table gives
an explicit elliptic pencil criterion. Choose one
exceptional $(-2)$-curve from each of four different
singular points, and call them $C_1,\ldots,C_4$.
They are pairwise disjoint. For an ample integral
class $H$ on $M$, put
\begin{equation}
 H'=H+\sum_{i=1}^4\frac{H\cdot C_i}{2}C_i,
 \quad H'\cdot C_i=0,
 \quad H'^2=H^2+\frac12\sum_{i=1}^4(H\cdot C_i)^2>0.
                                                               \hypertarget{tag-00P6}{}\label{cdp:eq:C.K3.8}
\end{equation}
Lagrange's four-square theorem
\cite[Theorem~II, Corollary~II, p.~201]{Lagrange1770}
expresses the positive rational number $H'^2/2$ as
$\sum_i x_i^2$ with $x_i\in\mathbb Q$: for a
positive fraction $a/b$, apply the integral theorem
to $ab$ and divide the four summands by $b$.
Consequently
\begin{equation}
 D=H'-\sum_{i=1}^4x_iC_i,
 \qquad D^2=0,
 \qquad D\cdot H'=H'^2>0.
                                                               \hypertarget{tag-00P7}{}\label{cdp:eq:C.K3.9}
\end{equation}
Clear denominators to obtain a nonzero isotropic
integral divisor class. The original elliptic-pencil
criterion \cite[\S3, Corollary~3, p.~560; \S6, Theorem~1(3) and its corollary, pp.~573--574]{PiatetskiShapiroShafarevich1971}
then produces a base-point-free elliptic
pencil on $M$.

\item For a quotient of algebraic dimension one, $a(A)=1$ as well. Choose a nonconstant
meromorphic function on $A$, and let $D$ be a
general member of its pencil after removing fixed
components. Its self-intersection is nonnegative,
by intersecting two distinct members. In fact
$D^2=0$. To verify this, average the positive
integration current $[D]$ under translations of
the compact torus. The result is the constant
semipositive $(1,1)$-form representing $c_1(D)$.
If $D^2>0$, this form is positive definite, so
the integral class is a polarization and the
torus is projective, contrary to $a(A)=1$.
The equality $D^2=0$ shows that two distinct
members have no common point. The pencil thus
defines a holomorphic map, and its Stein factor
$q:A\to B$ has connected smooth general fibre
with trivial normal line. Adjunction with
$K_A=0$ gives genus one. The field identity
\eqref{cdp:eq:C.K3.4}, applied to $A$, makes this fibration
$\langle g\rangle$-equivariant;
its quotient and \eqref{cdp:eq:C.K3.5} give the elliptic pencil on $M$.
\item Start with any torus $A$ and a cyclic symplectic action of order $n$ in \eqref{cdp:eq:C.K3.1}. Form the finite analytic quotient
\begin{equation}
 W=(A\times\mathbb P^1_s)/
       \langle(a,s)\mapsto(g(a),\zeta_ns)\rangle,
 \qquad W\longrightarrow\mathbb P^1_t,
 \qquad t=s^n.
                                                               \hypertarget{tag-00OZ}{}\label{cdp:eq:C.K3.10}
\end{equation}
Over $t\ne0,\infty$, the $n$ sheets of the base
change are permuted transitively, and their quotient
fibre is exactly $A$. The reduced fibres over zero
and infinity are exactly $A/\langle g\rangle$.
Choose an analytic resolution of the total quotient
and its two central supports, isomorphic over the
punctured base and at a generic point of each of
these two original components. The strict transform
of the latter maps properly birationally to the
torus quotient. Its normalization resolution
therefore has the K3 minimal model already proved
in \ref{cdp:st:C.2.UN.23}. The resulting smooth compact total
space has connected fibres: its Stein factor over
$\mathbb P^1$ has generic degree one and hence is
the identity finite map to the normal base.
\item Each of its two reduced special supports has even
Euler characteristic by \ref{cdp:st:C.2.UN.1}. Apply the vertical point and quartic-curve blowup operations
of \ref{cdp:st:C.2.UN.4} separately over zero and infinity.
They change the reduced Euler number by $+2$ or
$-2$ until it is zero, preserve the punctured torus
family, and leave the original strict component
birational to the same torus quotient. Resolution
and these blowups are projective morphisms when
$A$ is projective. This supplies all asserted
compact realizations, including every class of
algebraic dimension zero in the criterion.
\end{enumerate}
\end{proof}

\subsection{Minimal models of cyclic component covers}\label{subsec:II:nonnegative-component-covers}

\begin{lemma}[The intrinsic orbit pencil of a bielliptic surface]\hypertarget{tag-P348}{}\label{lem:II:bielliptic-orbit-pencil}
\StatementRecord{lem:II:bielliptic-orbit-pencil}{proved}{P348}{\cite[author notes, Chapter~1, \S3, p.~9]{Peters1990}}
Let $S$ be a bielliptic surface. In a finite free product
presentation
$S=(F_1\times F_2)/\Gamma$, take $\Gamma$ to act by translations
on $F_1$ and with a nontrivial linear character on $F_2$.
Put $K=\ker(\Gamma\to\operatorname{Aut}(F_2))$ and
$\overline F_1=F_1/K$, where $K$ acts by its first translations.
Then $\operatorname{Aut}^0(S)$ is this elliptic translation group,
and its orbit map is
\begin{equation}
 S=(F_1\times F_2)/\Gamma\longrightarrow
                         F_2/\Gamma\simeq\PP^1.
 \hypertarget{tag-00HL}{}\label{cdp:eq:C.EO.16}
\end{equation}
It has connected fibres, fixed elliptic Jacobian
$\overline F_1$, identity integral monodromy and only positive
multiples of smooth elliptic fibres. Every automorphism of $S$
preserves this map and acts on its connected-fibre base.
\end{lemma}
\begin{proof}
The free product presentation is the bielliptic description in
\cite[author notes, Chapter~1, \S3, p.~9]{Peters1990}: its factors
$(E,C)$ correspond to $(F_2,F_1)$ by the product permutation.
The subgroup $K$ embeds in the translations of $F_1$, since an
element trivial on both factors is the identity deck map.
The finite quotient maps give
\[
 (F_1\times F_2)/K=\overline F_1\times F_2,\qquad
 S=(\overline F_1\times F_2)/(\Gamma/K).
\]
These identifications and their inverses send the orbit of
$(x,y)$ to $([x],y)$ and then to its remaining finite orbit.
The action of $\Gamma/K$ is free: a fixed point would lift to
an element of $\Gamma$ whose first translation equals one from
$K$ and whose second action fixes $y$; composing by that element
of $K$ would fix a point of the original product.

Finite \'etale descent identifies $H^0(S,T_S)$ with the invariant
holomorphic vector fields of the product. Such fields are
constant, and the nontrivial second linear character leaves
exactly the first translation line. The action
$F_1\to\operatorname{Aut}^0(S)$ has this nonzero derivative.
Its kernel is exactly $K$: an automorphism of the product that
is the identity on $S$ is a deck transformation, and a first-factor
translation is one precisely when its second factor is the
identity. Its image is compact and has the whole one-dimensional
Lie algebra. It is therefore the connected identity component,
identified with $\overline F_1$.

Two points of $S$ have the same image in $F_2/\Gamma$ precisely
when a deck element first makes their second coordinates equal
and a first translation then carries one to the other. Thus
the displayed map is the orbit map and has connected fibres.
Holomorphic differentials of the quotient curve pull back to
the invariant differentials on $F_2$. The nontrivial linear
character makes that invariant space zero, so the quotient is
the smooth genus-zero curve $\PP^1$.

Use the faithful action of $\Gamma/K$ on $F_2$. At a point of
$F_2$ its stabilizer is cyclic, with a local coordinate on which
a generator acts by $u\mapsto\zeta_m u$. On
$\overline F_1$ it acts by a translation $\tau$ of exact order
$m$: a smaller order would give a nonidentity element fixing a
point of the product. The local quotient is consequently
$(\overline F_1\times\Delta_u)/
\langle(a,u)\mapsto(a+\tau,\zeta_m u)\rangle$, with base
coordinate $t=u^m$. Its full fibre is a smooth elliptic multiple.
The local and ordinary transition maps on the first factor are
translations, giving the stated fixed Jacobian and identity
integral monodromy.

Every automorphism normalizes the intrinsic connected group
$\operatorname{Aut}^0(S)$, hence sends its orbits to orbits.
It therefore acts on their quotient curve. The quotient map is
proper with connected fibres; its descended base map is
holomorphic by the direct-image identity for the structure
sheaf, and the inverse automorphism gives the inverse base map.
\end{proof}


\begin{proposition}[Cyclic covers of components with nonnegative Kodaira dimension]\hypertarget{tag-01DA}{}\label{cdp:st:C.2.UN.31}
\StatementRecord{cdp:st:C.2.UN.31}{proved}{01DA}{\cref{lem:II:cyclic-component-cover,cdp:st:C.2.UN.15,cdp:st:C.2.UN.22,cdp:st:C.2.UN.29,lem:II:identity-monodromy-elliptic,lem:II:elliptic-minimal-descent,cdp:st:C.2.UN.28,cdp:st:C.2.UN.23,cdp:st:C.2.UN.33,cdp:st:C.2.UN.32}; \cite[Chapter~1, \S3, pp.~9--10]{Peters1990}}

Let $f:U\to\Delta$ be an unrestricted proper degeneration
of connected complex two-tori from a smooth threefold.
Let $Z$ be the smooth minimal model of the normalization
of one central component, with $\kappa(Z)\ge0$.
There is a generically finite component cover
$E\to S^\nu$, obtained by the normalized root base change
and resolution, with cyclic stabilizer $G$, such that
$E/G$ is properly birational to $S^\nu$. Its smooth
minimal model $E_0$ satisfies one of these descriptions:

\begin{itemize}
\item $E_0$ is a complex two-torus;
\item $E_0$ has an isotrivial elliptic fibration over a
curve $B$ with $g(B)\le1$, identity integral $H_1$
monodromy, and only positive multiples of smooth
elliptic fibres. In this case
$R^1q_*\mathcal O_{E_0}=\mathcal O_B$,
$p_g(E_0)=g(B)$, and $\chi(\mathcal O_{E_0})=0$.
\end{itemize}
In both cases $\chi(\mathcal O_{E_0})=0$, and the
action of $G$ descends holomorphically to $E_0$.

For the original minimal component the consequences are
\begin{equation}
 p_g(Z)\le1,
 \qquad \chi(\mathcal O_Z)\in\{0,1,2\}.
                                                               \hypertarget{tag-00DH}{}\label{cdp:eq:C.CS.1}
\end{equation}
Every component with $\kappa(Z)=1$ has an isotrivial elliptic Iitaka fibration on $Z$, over a
curve of genus at most one. An Enriques component
has an isotrivial elliptic fibration over $\mathbb P^1$.
Every K3 component is a cyclic torus quotient as in \cref{cdp:st:C.2.UN.32}, with the orders and configurations stated there. In the non-torus-cover branch its comparison also gives an isotrivial elliptic fibration over $\PP^1$.
\end{proposition}


\begin{proof}

\begin{enumerate}[label=\arabic*.,leftmargin=*]
\item Take the component cover $E\to S^\nu$, its cyclic stabilizer $G$ and the proper birational quotient map \eqref{cdp:eq:C.HQ.2} of \cref{lem:II:cyclic-component-cover}. Pluricanonical pullback through a
resolution of the generically finite graph gives
\begin{equation}
             \kappa(E)\ge\kappa(Z)\ge0.
                                                               \hypertarget{tag-00DI}{}\label{cdp:eq:C.CS.2}
\end{equation}
The compact-component result in \cref{cdp:st:C.2.UN.15} applies
also to this finite base-change family, so
$\kappa(E)\le1$. For the saturated limiting field
space, the generic-rank-two branch is covered
by \cref{cdp:st:C.2.UN.22}. In rank one the first integral
has rational or elliptic general fibre. Rational
fibres would force every positive pluricanonical
section to vanish, by the exact adjunction
argument of \cref{cdp:st:C.2.UN.20}, contradicting \eqref{cdp:eq:C.CS.2}.
Thus \cref{cdp:st:C.2.UN.29} applies. These two statements
give the torus or relatively minimal smooth-multiple elliptic alternatives and coherent Euler characteristic zero. In the latter branch, \eqref{cdp:eq:C.CS.2} and \cref{lem:II:elliptic-minimal-descent} make the minimal contraction an isomorphism. Its birational comparison with the unique smooth minimal model $E_0$ is therefore an isomorphism, so the elliptic map is defined on this same $E_0$.

\item Apply \cref{lem:II:identity-monodromy-elliptic} to the relatively minimal elliptic map $q:E_0\to B$ in either field-rank branch. It has the identity integral monodromy just obtained, so the lemma's full-base invariant calculation gives $R^1q_*\OO_{E_0}\simeq\OO_B$ and $p_g(E_0)=g(B)$. The component $E$ belongs to the smooth resolved finite-base-change torus family. The finite-base-change clause of \cref{cdp:st:C.2.UN.28} gives $p_g(E)=p_g(E_0)\le1$, hence $g(B)\le1$.

\item The minimal-model clause of \cref{lem:II:cyclic-component-cover} applies because $\kappa(E)\ge0$. It gives the descended holomorphic action and the proper quotient comparison
\begin{equation}
 \begin{array}{ccc}
 E&\longrightarrow&E_0\\
 \downarrow&&\downarrow\\
 E/G&\longrightarrow&E_0/G
 \end{array}
 \qquad E_0/G\dashrightarrow Z
       \quad\hbox{properly birational}.
                                                               \hypertarget{tag-00DJ}{}\label{cdp:eq:C.CS.3}
\end{equation}
Here the square is \eqref{cdp:eq:C.HQ.3} for this same cover, and its final birational map is the composite with the normalization and minimal model of the chosen component.

\item Suppose $\kappa(Z)=1$. Equation \eqref{cdp:eq:C.CS.2} and \cref{cdp:st:C.2.UN.15} give $\kappa(E_0)=1$. Apply the original-component conclusion of \cref{cdp:st:C.2.UN.29} to this same component cover. Its construction uses $q:E_0\to B$, the action of $G$ on the pluricanonical field, an equivariant graph resolution $E'\to E_0$ and its connected-fibre Stein map. The quotient is
\begin{equation}
                   E'/G\longrightarrow B/G,
                                                               \hypertarget{tag-00DK}{}\label{cdp:eq:C.CS.4}
\end{equation}
The general fibre is $F/K$, where $K=\ker(G\to\operatorname{Aut}(B))$. In the proof of that result, the rational-quotient adjunction argument makes $K$ a finite translation subgroup of the fixed elliptic curve $A$, and its bounded-order translation parameters are locally constant. Thus the general quotient fibre is the fixed curve $A/K$. Its base has genus at most one by Riemann--Hurwitz, and this connected-fibre quotient is the Iitaka map on the smooth birational model of $Z$.

\item Resolve that quotient fibration and its birational comparison with $Z$ as
$Z\xleftarrow{\rho}X\xrightarrow{q}B/G$.
Apply \cref{lem:II:elliptic-minimal-descent} to this common resolution and the minimal surface $Z$ of Kodaira dimension one. Its canonical-intersection identity \eqref{cdp:eq:C.CS.5} makes every exceptional curve vertical and descends the map holomorphically to $Z$. The general fibres and their fixed elliptic modulus are unchanged, giving the asserted isotrivial Iitaka fibration.

\item Now let $Z$ be Enriques. The function-field clause of \cref{lem:II:cyclic-component-cover} applied to $E\to S^\nu$ gives $a(E_0)=a(Z)=2$. Thus the smooth compact surface $E_0$ is projective. If it has Kodaira dimension
one, the preceding intrinsic Iitaka
quotient argument again gives an
isotrivial elliptic fibration on a
birational model of $Z$. In Kodaira dimension zero, the projective cases of the surface classification \cite[Chapter~1, \S3, pp.~9--10]{Peters1990} and $\chi(\mathcal O_{E_0})=0$ make it Abelian or bielliptic.
The Abelian case is excluded by the
exact cyclic quotient result in \cref{cdp:st:C.2.UN.23}.
For the bielliptic case, apply
\cref{lem:II:bielliptic-orbit-pencil} to the same $E_0$ and
its finite free presentation $(A_1\times A_2)/H$.
The map in \eqref{cdp:eq:C.EO.16} is its orbit map with
$(F_1,F_2,\Gamma)=(A_1,A_2,H)$; the finite kernel on $A_2$
gives its fixed quotient elliptic fibre. The lemma gives smooth
multiple fibres, identity monodromy and invariance under every
automorphism of this same $E_0$, hence under the descended $G$.
The quotient
\eqref{cdp:eq:C.CS.4} is available again. Its
general fibre must be elliptic, since
a rational general fibre would force
the nonzero even pluricanonical form
of an Enriques surface to vanish.

\item Apply \cref{lem:II:elliptic-minimal-descent} to the resulting common resolution over the minimal Enriques surface $Z$. Its canonical class is torsion, so the $K_Z$ term of \eqref{cdp:eq:C.CS.5} is zero. The lemma descends the elliptic map to $Z$, and its injection into $H^1(Z,\OO_Z)=0$ makes the base $\PP^1$. This gives the asserted isotrivial Enriques fibration.

\item For a K3 target $Z$, apply \cref{cdp:st:C.2.UN.32} to the same degeneration and component normalization. It identifies $Z$ with the minimal resolution of the stated cyclic symplectic torus quotient in every algebraic dimension. Its proof gives the exact torus $A=E_0$ in the initial torus branch and $A=E_0/H$ in the non-torus branch, with remaining cyclic group $G/H$ and quotient equality
$(E_0/H)/(G/H)=E_0/G$.
In the latter branch, the invariant elliptic map descends through \eqref{cdp:eq:C.K3.5} to an isotrivial elliptic fibration on $Z$ over $\PP^1$. These are the same proper quotient and common-resolution maps as \eqref{cdp:eq:C.CS.3}.

\item The compact-core bound \cref{cdp:st:C.2.UN.28} gives $p_g(Z)\le1$. The preceding cover bound also gives $0\le\kappa(Z)\le1$. In Kodaira dimension zero, the compact surface classification \cite[Chapter~1, \S3, pp.~9--10]{Peters1990} gives coherent Euler characteristic zero for tori, bielliptic and Kodaira surfaces, two for K3 surfaces and one for Enriques surfaces. In Kodaira dimension one, apply \cref{cdp:st:C.2.UN.33} with $M=Z$; its complete local quotient and Noether calculation gives $\chi(\OO_Z)\in\{0,1,2\}$. In particular $\chi(\OO_Z)\ge0$ in every case. On the other hand
\begin{equation}
        \chi(\mathcal O_Z)=1-h^1(Z,\mathcal O_Z)+p_g(Z)
                              \le2.
                                                               \hypertarget{tag-00DM}{}\label{cdp:eq:C.CS.6}
\end{equation}
The integer $\chi(\OO_Z)$ therefore belongs to the set in \eqref{cdp:eq:C.CS.1}.

\end{enumerate}
\end{proof}

\section{Integral component lattices and support topology}\label{sec:II:component-lattices-topology}

\subsection{Integral component lattices and monodromy}\label{subsec:II:component-invariant-lattice}

\begin{proposition}[Compact component cores give a multiplicity-weighted invariant lattice in one nearby torus]\hypertarget{tag-01DF}{}\label{cdp:st:C.2.UN.36}
\StatementRecord{cdp:st:C.2.UN.36}{proved}{01DF}{\cref{lem:II:compact-component-transport,cdp:st:C.2.UN.33,night:P4.13}; \cite[\S3.G, p.~321]{Hatcher2002}; \cite[author draft, Chapter~1, Proposition~3.5, p.~17]{HuybrechtsK3Draft}; \cite[Chapter~1, \S1, pp.~5--6; \S2, p.~8; \S3, p.~9]{Peters1990}}

Let $f:U\to\Delta$ be a proper holomorphic map from
a smooth complex threefold, with connected two-torus
fibres over $\Delta^*$. Write its full central fibre
as $\sum_i m_iS_i$, and put $S=\bigcup_iS_i$ with
its reduced structure. For each component choose a
resolution of its normalization
$h_i:Z_i\to S_i\subset S$ which is an isomorphism
over $S_i\setminus\operatorname{Sing}S$. Set
\begin{equation}
 \begin{gathered}
 D_i=h_i^{-1}(\operatorname{Sing}S)_{\rm red},\qquad
 R_i=\operatorname{span}_{\mathbb R}
       \{[C]:C\text{ is a curve component of }D_i\},
 \\
 K_i=\{v\in H_2(Z_i,\Z)/\operatorname{tors}:v\cdot R_i=0\}.
 \end{gathered}
                                                               \hypertarget{tag-00BW}{}\label{cdp:eq:C.CL.1}
\end{equation}
The set $D_i$ may also have isolated points; they
contribute no curve class to $R_i$. The pairing on
$K_i$ is the restriction of the ordinary oriented
surface intersection form and may have a radical.

Choose a common base-change order $N$ divisible by
all $m_i$ and by the semisimple monodromy order.
Let $F$ be a nearby torus after this base change and let $U_1$ be its monodromy on $H_1(F,\Z)$. There is a homomorphism, obtained by
choosing compact representatives and their transfers,
\begin{equation}
 \Psi:\bigoplus_iK_i\longrightarrow
       H_2(F,\Z)^{\Lambda^2U_1},\qquad
 \Psi(v)\cdot\Psi(w)=\sum_i m_i(v_i\cdot w_i).
                                                               \hypertarget{tag-00C1}{}\label{cdp:eq:C.CL.2}
\end{equation}
Thus every multiplicity appears in the preserved
bilinear form. The kernel is contained in the
radical of that form. If its source is nondegenerate
of rank six, its image has finite integral index
in $H_2(F,\Z)$ and
\begin{equation}
 \left|\det\left(\bigoplus_i m_iQ_{K_i}\right)\right|
                       \text{ is an integer square}.
                                                               \hypertarget{tag-00C2}{}\label{cdp:eq:C.CL.3}
\end{equation}

Let $(p_i,q_i,z_i)$ be the positive, negative and
radical dimensions of the real form on $K_i$.
Let $(p_U,q_U,z_U)$ be those of the intersection
form restricted to
$\ker(\Lambda^2U_1-I)\subset H_2(F,\mathbb R)$.
Then
\begin{equation}
             \sum_i p_i\le p_U\le3,\qquad
             \sum_i q_i\le q_U\le3.                 \hypertarget{tag-00C3}{}\label{cdp:eq:C.CL.4}
\end{equation}
In particular, if $r_i=\dim R_i$, then
\begin{equation}
 \sum_i\max\{b_2^-(Z_i)-r_i,0\}\le q_U\le3,
 \qquad
                 2\sum_i p_g(Z_i)\le p_U.           \hypertarget{tag-00C4}{}\label{cdp:eq:C.CL.5}
\end{equation}
A K3 component whose chosen resolution is obtained
from its minimal model by $n$ point blowups must
therefore have at least $16+n$ independent boundary
curve classes. For an Enriques component or a
properly elliptic component of coherent Euler
characteristic one, the corresponding lower bound
is $6+n$.

The monodromy bound \eqref{cdp:eq:C.MS.10} of
\cref{night:P4.13}, transported through the chosen base change
and the natural cohomology--homology pairing, gives
$(U_1-I)^2=0$. In an oriented real basis of
$H_1(F,\mathbb R)$, all its possible Jordan forms
give the following exact invariant forms:
\begin{equation}
\begin{array}{c|c|c|c}
\operatorname{rank}(U_1-I)&\text{nonzero values of }U_1-I
 &\dim\ker(\Lambda^2U_1-I)&(p_U,q_U,z_U)\\ \hline
0&\text{none}&6&(3,3,0)\\
1&e_2\mapsto e_1&4&(1,1,2)\\
2&e_2\mapsto e_1,\ e_4\mapsto e_3&4&(1,2,1)\\
2&e_2\mapsto e_1,\ e_4\mapsto-e_3&4&(2,1,1).
\end{array}                                                   \hypertarget{tag-00C5}{}\label{cdp:eq:C.CL.6}
\end{equation}
The two last rows record the two oriented real
forms of that Jordan type. In particular rank-one
unipotent monodromy forces $p_g(Z_i)=0$ for every
normalization component.
\end{proposition}


\begin{proof}

\begin{enumerate}[label=\arabic*.,leftmargin=*]
\item The open set $Z_i^\circ=Z_i\setminus D_i$ maps isomorphically to $S_i\setminus\operatorname{Sing}S$, where the reduced support is one smooth embedded surface. Apply the integral part of \cref{lem:II:compact-component-transport} with $E=Z_i$ and $D=D_i$. Its relative homology map is
\begin{equation}
 H_2(Z_i^\circ,\Z)\longrightarrow H_2(Z_i,\Z)
 \longrightarrow H_2(Z_i,Z_i^\circ;\Z)
       \simeq H^2(D_i,\Z).                          \hypertarget{tag-00C6}{}\label{cdp:eq:C.CL.7}
\end{equation}
The target is free on the curve components of $D_i$, and the second arrow gives the vector of intersections with them. Torsion maps to zero. Every element of $K_i$ therefore has an integral lift represented by a compact finite cycle in $Z_i^\circ$. Choose such lifts for an integral basis of each $K_i$, with compact transverse smooth chain representatives and the perturbations needed to compute all their intersections.

\item Apply the root construction of \cref{lem:II:compact-component-transport} to the stated common multiple $N$. Through the isomorphism $Z_i^\circ\simeq S_i\setminus\operatorname{Sing}S$, its central cover becomes
\begin{equation}
        \pi_i:\widetilde Z_i^\circ
                       \longrightarrow Z_i^\circ
                  \quad\text{of total degree }m_i.  \hypertarget{tag-00C8}{}\label{cdp:eq:C.CL.9}
\end{equation}
The cover may be disconnected. In the chart \eqref{cdp:eq:C.PG.3}, the parameter $s$ has order one on every strict component, and the resolution is an isomorphism over this open cover. Thus these cycles lie in smooth single-component loci of multiplicity-one central components.

\item Transfer an integral compact cycle by taking
the sum of all its lifts under \eqref{cdp:eq:C.CL.9},
subdividing its finite singular chain into
evenly covered pieces. This is the ordinary
integral covering transfer
\cite[\S3.G, p.~321]{Hatcher2002}
and is a chain map. For two transverse compact cycles,
every signed intersection point downstairs
has exactly $m_i$ lifts with the same local
orientation. Hence their transferred
intersection is exactly $m_i$ times the
original intersection. This proves the
multiplicity factor directly, including
disconnected root covers.

\item Apply the third part of \cref{lem:II:compact-component-transport} to the transferred cycles and their compact transverse perturbations in these multiplicity-one loci. It gives simultaneous orientation-preserving embeddings into one nearby torus $F=g^{-1}(s)$, preserving intersections within a component and giving disjoint images for distinct strict components.

\item Transport the selected transferred basis
cycles and extend by integer linearity.
This defines the homomorphism $\Psi$.
The transfer intersection calculation and
the disjoint neighbourhoods give the
bilinear equality of \eqref{cdp:eq:C.CL.2}. The choices of integral lifts and their transferred basis cycles determine this homomorphism.

\item The whole-disc cycle families in \cref{lem:II:compact-component-transport} show that every transported basis class is fixed by monodromy on $H_2(F,\Z)$. The natural exterior-algebra identification $H_2(F,\Z)\simeq\Lambda^2H_1(F,\Z)$ intertwines monodromy. Therefore $\Psi$ takes values in the invariant group of \eqref{cdp:eq:C.CL.2}.

\item If $\Psi(v)=0$, the preserved pairing
with every $\Psi(w)$ shows that $v$
is in the radical of the source form.
Choose a real complement of that radical
on which the restricted form is
nondegenerate. Its image is injective
and has the same positive and negative
indices. Maximal positive and negative
subspaces of the target restricted form
therefore give \eqref{cdp:eq:C.CL.4}. If the integral
source is nondegenerate of rank six,
$\Psi$ is an injection of full-rank
free lattices into $H_2(F,\Z)$.
The latter has the unimodular form
$3H$, of determinant $-1$. If the
image has index $c$, the change-of-basis
determinant gives source determinant
$-c^2$. This proves \eqref{cdp:eq:C.CL.3}.

\item Intersect a maximal negative definite
subspace of $H_2(Z_i,\mathbb R)$ with
the codimension-$r_i$ subspace $K_i\otimes\mathbb R$.
Its dimension is at least
$b_2^-(Z_i)-r_i$, so its negative index
is at least the corresponding nonnegative
maximum in \eqref{cdp:eq:C.CL.5}. For the second
inequality there, the real and imaginary
parts of holomorphic two-forms give the
positive $2p_g(Z_i)$-dimensional subspace
of \cref{lem:II:compact-component-transport}, applied with $E=Z_i$ and $D=D_i$. Its Poincar\'e-dual classes pair trivially with every curve component of $D_i$, and hence belong to $K_i\otimes\mathbb R$.
This proves \eqref{cdp:eq:C.CL.5}.

\item A K3 minimal surface has intersection
signature $(3,19)$
\cite[author draft, Chapter~1, Proposition~3.5, p.~17]{HuybrechtsK3Draft};
$n$ point blowups add $n$ orthogonal classes of square $-1$.
The point-blowup splitting, and the Enriques invariants
and signature formula giving signature $(1,9)$, are recorded in
\cite[Chapter~1, \S1, pp.~5--6; \S2, p.~8; \S3, p.~9]{Peters1990}. A properly
elliptic component with $\chi(\OO)=1$
has $q=p_g=0$ and $e=12$ by \ref{cdp:st:C.2.UN.33},
and signature $(1,9)$ as well:
$b_1=0$, $b_2=10$, and the signature
formula with $K^2=0$ gives signature
$-8$. Applying the negative-index
bound to $19+n$ or $9+n$ gives
the asserted boundary ranks.

\item Apply \cref{night:P4.13} to this same degeneration.
Let $T$ be its monodromy on $H^1(F,\Q)$ and $M$ its
semisimple order. The chosen root has $M\mid N$.
Writing $D=T^M-I$, equation \eqref{cdp:eq:C.MS.10} gives
$D^2=0$, so the cohomology monodromy after the root is
$T^N=I+(N/M)D$. The flat perfect pairing with
$H_1(F,\Q)$ identifies $(T^N)^{-1}$ with the dual of
$U_1$. Since $(T^N)^{-1}=I-(N/M)D$, it follows that
$N_0=U_1-I$ has square zero. Its image lies in its kernel
on a four-dimensional space, so its rank is at most two.
In rank one,
choose a vector and its nonzero
image, then complete the basis in
the kernel. Changing the sign of
one unused basis vector makes the
basis positively oriented, giving
the second row of \eqref{cdp:eq:C.CL.6}.
In rank two, choose two image
vectors and their preimages. If
their resulting ordered basis is
negative, reverse one preimage.
This gives the two oriented forms
\begin{equation}
          N_0e_2=e_1,\qquad N_0e_4=\epsilon e_3,
                     \qquad\epsilon\in\{1,-1\}.     \hypertarget{tag-00BX}{}\label{cdp:eq:C.CL.10}
\end{equation}
Write $e_{ij}=e_i\wedge e_j$.
In the oriented wedge basis
$(e_{12},e_{13},e_{14},e_{23},e_{24},e_{34})$
the only nonzero complementary
pairings are
\begin{equation}
 e_{12}\cdot e_{34}=1,\qquad
 e_{13}\cdot e_{24}=-1,\qquad
 e_{14}\cdot e_{23}=1.                              \hypertarget{tag-00BY}{}\label{cdp:eq:C.CL.11}
\end{equation}
In rank one, the only nonzero images under $\Lambda^2U_1-I$ are $e_{23}\mapsto e_{13}$ and $e_{24}\mapsto e_{14}$. Its kernel therefore has basis $(e_{12},e_{13},e_{14},e_{34})$, with one hyperbolic plane and a two-dimensional radical. For \eqref{cdp:eq:C.CL.10}, the nonzero images are
$e_{14}\mapsto\epsilon e_{13}$,
$e_{23}\mapsto e_{13}$ and
$e_{24}\mapsto e_{14}+\epsilon e_{23}+\epsilon e_{13}$.
Thus a vector is invariant precisely when its $e_{24}$ coefficient is zero and its $e_{23}$ coefficient is $-\epsilon$ times its $e_{14}$ coefficient. An invariant basis is
\begin{equation}
       e_{12},\ e_{34},\ e_{13},\ e_{14}-\epsilon e_{23}.
                                                               \hypertarget{tag-00BZ}{}\label{cdp:eq:C.CL.12}
\end{equation}
Its Gram matrix is
\begin{equation}
 \begin{pmatrix}
 0&1&0&0\\
 1&0&0&0\\
 0&0&0&0\\
 0&0&0&-2\epsilon
 \end{pmatrix}.                                               \hypertarget{tag-00C0}{}\label{cdp:eq:C.CL.13}
\end{equation}
These are exactly the rows of
\eqref{cdp:eq:C.CL.6}; the identity row is
the whole form $3H$.
The rank-one row has $p_U=1$, so \eqref{cdp:eq:C.CL.5} forces every $p_g(Z_i)$ to vanish.

\end{enumerate}
\end{proof}

\subsection{Rank-four monodromy stratum signatures}\label{subsec:II:rank-four-monodromy-signatures}

\begin{proposition}[All rank-four monodromy stratum signatures]\hypertarget{tag-01CO}{}\label{cdp:st:C.2.UN.11}
\StatementRecord{cdp:st:C.2.UN.11}{proved}{01CO}{\cref{night:P4.13,cdp:st:C.2.UN.1,cdp:st:C.2.UN.4}; \cite[published version, (0.1)(1), p.~208; Theorem~1.6 and Remarks~1.7(2), pp.~215--216]{BierstoneMilman1997}}

For the proper torus degeneration $f:U\to\Delta$, resolution
$\pi:U'\to U$, monodromy $T$ and stratum sequence $(a_n)$ of
\cref{night:P4.13}, the complete range of
\eqref{cdp:eq:C.MS.2} is the following table. Here $\Phi_d$ is
the $d$th cyclotomic polynomial, $P_T$ is the characteristic
polynomial of $T$, and every unlisted $a_n$ is zero.

\begin{center}
\small
\begin{tabular}{ccl}
$P_T$ & $M$ & Nonzero $a_n$\\\hline
$\Phi_{1}^{4}$ & 1 & $\text{all zero}$ \\
$\Phi_{1}^{2}\Phi_{2}^{2}$ & 2 & $\text{all zero}$ \\
$\Phi_{1}^{2}\Phi_{3}$ & 3 & $\text{all zero}$ \\
$\Phi_{1}^{2}\Phi_{4}$ & 4 & $\text{all zero}$ \\
$\Phi_{1}^{2}\Phi_{6}$ & 6 & $\text{all zero}$ \\
$\Phi_{2}^{4}$ & 2 & $a_{1}=16,\quad a_{2}=-8$ \\
$\Phi_{3}^{2}$ & 3 & $a_{1}=9,\quad a_{3}=-3$ \\
$\Phi_{2}^{2}\Phi_{4}$ & 4 & $a_{1}=8,\quad a_{2}=-4$ \\
$\Phi_{4}^{2}$ & 4 & $a_{1}=4,\quad a_{2}=6,\quad a_{4}=-4$ \\
$\Phi_{5}$ & 5 & $a_{1}=5,\quad a_{5}=-1$ \\
$\Phi_{2}^{2}\Phi_{3}$ & 6 & $a_{1}=12,\quad a_{2}=-6,\quad a_{3}=-4,\quad a_{6}=2$ \\
$\Phi_{2}^{2}\Phi_{6}$ & 6 & $a_{1}=4,\quad a_{2}=-2,\quad a_{3}=4,\quad a_{6}=-2$ \\
$\Phi_{3}\Phi_{6}$ & 6 & $a_{1}=3,\quad a_{2}=3,\quad a_{3}=-1,\quad a_{6}=-1$ \\
$\Phi_{6}^{2}$ & 6 & $a_{1}=1,\quad a_{2}=4,\quad a_{3}=5,\quad a_{6}=-4$ \\
$\Phi_{8}$ & 8 & $a_{1}=2,\quad a_{2}=1,\quad a_{4}=3,\quad a_{8}=-2$ \\
$\Phi_{10}$ & 10 & $a_{1}=1,\quad a_{2}=2,\quad a_{5}=3,\quad a_{10}=-2$ \\
$\Phi_{3}\Phi_{4}$ & 12 & $a_{1}=6,\quad a_{2}=3,\quad a_{3}=-2,\quad a_{4}=-3,\quad a_{6}=-1,\quad a_{12}=1$ \\
$\Phi_{4}\Phi_{6}$ & 12 & $a_{1}=2,\quad a_{2}=5,\quad a_{3}=2,\quad a_{4}=-3,\quad a_{6}=-3,\quad a_{12}=1$ \\
$\Phi_{12}$ & 12 & $a_{1}=1,\quad a_{3}=1,\quad a_{4}=2,\quad a_{6}=2,\quad a_{12}=-2$ \\
\end{tabular}
\end{center}

The table is the range on all orientation-preserving quasi-unipotent
integral rank-four matrices, so every filling in
\cref{night:P4.13} has one of its signatures. One has $a_n=0$
unless $n\mid M$, where
\begin{equation}
M\in\{1,2,3,4,5,6,8,10,12\}.
 \hypertarget{tag-00RK}{}\label{cdp:eq:C.MS.3}
\end{equation}
Every row is realized by a map from a smooth projective threefold
to $\PP^1$ whose general fibre is an abelian surface and every
fibre has Euler number zero. It can be chosen to have exactly
two singular fibres, both with reducible simple-normal-crossings
support.

For an irreducible full fibre $mS$, the table and
\eqref{cdp:eq:C.MS.4} give
\begin{equation}
\begin{gathered}
 m>1\quad\Longrightarrow\quad e(S_{\rm reg})=0,\\
 e(S_{\rm reg})\in
       \{0,1,2,3,4,5,6,8,9,12,16\},\\
 e(S_{\rm reg})>0\quad\Longrightarrow\quad m=1.
 \end{gathered}                                  \hypertarget{tag-00RM}{}\label{cdp:eq:C.MS.5}
\end{equation}
\end{proposition}
\begin{proof}
\begin{enumerate}[label=\arabic*.,leftmargin=*]
\item Apply \cref{night:P4.13} to the same map $f$, embedded
resolution $\pi$, nearby fibre $F$ and its cohomology monodromy
$T$. It gives the determinant and M\"obius identities
\eqref{cdp:eq:C.MS.2}. The characteristic
polynomial of a quasi-unipotent integral rank-four
matrix is a product of cyclotomic polynomials
whose degrees sum to four. The possible individual
orders are
$1,2,3,4,5,6,8,10,12$. Since $\det T=1$,
the multiplicity of the eigenvalue $-1$ is even.
There are exactly nineteen resulting products:
the five containing $\Phi_1$ and the fourteen
remaining products displayed in the table.
This enumeration follows directly from
$\deg\Phi_1=\deg\Phi_2=1$,
$\deg\Phi_3=\deg\Phi_4=\deg\Phi_6=2$,
and degree four for $\Phi_5,\Phi_8,\Phi_{10},
\Phi_{12}$. The sum of the degree-one
multiplicities is even, so the $\Phi_1$
multiplicity is even as well.
\item The determinants $\det(I-T^k)$ depend only on
these eigenvalues, including their multiplicities;
Jordan blocks make no change to them. Each listed
product is represented by the block diagonal
matrix of its integral cyclotomic companion
matrices, which has determinant one. Its traces
therefore compute the exact full range on
matrices. Applying \eqref{cdp:eq:C.MS.2} gives the displayed
integers.
\item The exterior-power identity gives
\begin{equation}
\begin{aligned}
 \exp\!\left(\sum_{k\ge1}\frac{L_k}{k}t^k\right)
 &=\prod_{q=0}^4
       \det\bigl(I-t\Lambda^qT\bigr)^{(-1)^{q+1}}\\
 &=\prod_{n\ge1}(1-t^n)^{-a_n}.
 \end{aligned}                                  \hypertarget{tag-00RO}{}\label{cdp:eq:C.MS.7}
\end{equation}
Every eigenvalue of every exterior power has
order dividing $M$. Hence the first product
factors into cyclotomic polynomials $\Phi_d(t)$
with $d\mid M$. The identity
$\Phi_d(t)=\prod_{n\mid d}(1-t^n)^{\mu(d/n)}$
for $d>1$, with the harmless sign at $d=1$,
expresses it using only indices dividing $M$.
Finite products of the factors $1-t^n$ have
unique integer exponent vectors: the coefficient
at the smallest differing degree in their
formal logarithms detects any difference.
Thus $a_n=0$ for every $n\nmid M$, as claimed.
\item For $m>1$, the vanishing assertion of \cref{night:P4.13}
gives $e(S_{\mathrm{reg}})=0$. For $m=1$,
\eqref{cdp:eq:C.MS.4} gives $e(S_{\mathrm{reg}})=a_1$,
whose table entries are exactly the set in \eqref{cdp:eq:C.MS.5}.
\item We realize every row by a compact family. The fifteen products of two degree-two
blocks are obtained from products of elliptic
curves. The five blocks
\[
                 \Phi_1^2,\quad\Phi_2^2,\quad
                 \Phi_3,\quad\Phi_4,\quad\Phi_6
\]
come respectively from the identity, negation,
and multiplication by a primitive third, fourth
or sixth root of unity on an elliptic curve
with the corresponding triangular or square
lattice. Taking the product automorphism gives
exactly the indicated product on integral
first cohomology.
\item For $\Phi_5$, let $A$ be the Jacobian of
the smooth genus-two curve
\[
                           C:\ y^2=x^5-1.
\]
Its double cover of $\PP^1$ has six simple
branch points, including infinity, so the
genus assertion follows from Riemann--Hurwitz.
The automorphism $(x,y)\mapsto(\zeta_5x,y)$
fixes infinity and acts on the basis
$dx/y,x\,dx/y$ of holomorphic differentials
with eigenvalues $\zeta_5,\zeta_5^2$.
The Abel--Jacobi map based at infinity
intertwines this automorphism with its induced
automorphism of $A$ and induces the
isomorphism on integral $H^1$. Including the
conjugate eigenvalues gives $\Phi_5$.
Composing the Jacobian automorphism with
negation gives order ten and characteristic
polynomial $\Phi_{10}$. Here negation is also
the map induced by the curve's hyperelliptic
involution: the divisor of $x-x(p)$ gives
$[p-\infty]+[\iota(p)-\infty]=0$ on the
Jacobian, verifying that comparison explicitly.
\item For the remaining two rows take
$A=E_i\times E_i$, where
$E_i=\C/(\Z+i\Z)$. The Gaussian-integer
matrices
\begin{equation}
G_8=\begin{pmatrix}0&i\\1&0\end{pmatrix},
 \qquad
 G_{12}=\begin{pmatrix}0&1\\1&i\end{pmatrix}
 \hypertarget{tag-00RP}{}\label{cdp:eq:C.MS.8}
\end{equation}
define automorphisms of this torus.
One has $G_8^2=iI$, and its holomorphic
eigenvalues are $\zeta_8,\zeta_8^5$.
The characteristic polynomial of $G_{12}$
over $\C$ is $x^2-ix-1$, with roots
$(i+\sqrt3)/2=\zeta_{12}$ and
$(i-\sqrt3)/2=\zeta_{12}^5$.
Together with the conjugate eigenvalues,
these give the integral characteristic
polynomials $\Phi_8$ and $\Phi_{12}$ and
the exact orders eight and twelve.
Both matrices preserve the displayed lattice;
their determinants are Gaussian units.
This accounts for all nineteen rows with
finite-order automorphisms $g$ of projective
two-tori $A$.
\item Let $M$ be the order of the selected $g$.
Take the projective finite quotient
\begin{equation}
X=(A\times\PP^1)/
       \langle(a,z)\mapsto(g(a),\zeta_Mz)\rangle,
 \qquad f([a,z])=z^M.                           \hypertarget{tag-00RQ}{}\label{cdp:eq:C.MS.9}
\end{equation}
The action is free over $z\ne0,\infty$.
Over a nonzero finite base value the $M$
slices are permuted transitively, and their
quotient is the original connected torus $A$.
The positive meridian lifts from $z$
to $\zeta_Mz$, so under this identification
the monodromy is $g^{-1}$. Inverting all
the eigenvalues preserves each cyclotomic
polynomial and hence the displayed signature.
Choose a projective embedding $X\hookrightarrow\PP^N$.
Apply \cite[published version, (0.1)(1), p.~208; Theorem~1.6
and Remarks~1.7(2), pp.~215--216]{BierstoneMilman1997}
first to this reduced algebraic $X$, obtaining
$\rho:X_1\to X$ with $X_1$ smooth, and then to the reduced
divisor $D=((f\circ\rho)^{-1}\{0,\infty\})_{\rm red}$
in $X_1$. Denote the second composite by $\sigma:X_2\to X_1$.
The first centers lie over $\operatorname{Sing}X$, which is
contained in the two central fibres; the second centers lie
over $D$. Every blowup is algebraic and projective.
Thus $X_2$ is smooth and projective, the reduced inverse image
of $D$ is SNC, and $f\circ\rho\circ\sigma$ restricts to the
same quotient family over $\PP^1\setminus\{0,\infty\}$. Its Stein factor has generic
degree one over the normal curve $\PP^1$,
so every fibre is connected. When $M=1$
the construction is simply the product,
which already has zero-Euler fibres.
\item For the remaining cases, each of the two
possibly singular supports has even Euler
number by \cref{cdp:st:C.2.UN.1}. Apply the exact vertical
blowups of \cref{cdp:st:C.2.UN.4} to make each number
zero. A point blowup changes it by $+2$;
a point blowup followed by the blowup of
a smooth plane quartic in its exceptional
$\PP^2$ changes it by $-2$. These
operations are projective and are performed
only over the selected value, so they
preserve the punctured family and
the signature \eqref{cdp:eq:C.MS.2}.
\item The operations in \cref{cdp:st:C.2.UN.4} preserve SNC
supports. In that proposition choose the point on one component
of the present central support away from all the others, and
use its Fermat quartic in the exceptional plane. The coordinates
$(r,u,v)$ and blowdown substitutions in
\eqref{cdp:eq:C.VR.6} and \eqref{cdp:eq:C.VR.8} apply to this
same pair of blowups: the base functions are $(ru)^m$ and
$(uvw)^m$, with the components given by their coordinate
hyperplanes. This proves the SNC assertion for these realizations.
After the Euler number is zero, one additional
$+2$ operation and one $-2$ operation over the
same value leave it zero and introduce
exceptional components. Performing this at
both $0$ and $\infty$ ensures precisely two
reducible singular fibres, including the
identity-monodromy row. All other fibres
remain the original smooth tori.
\item The matrix-to-signature map has nontrivial fibres. The matrices
$I_4$ and $I_4+E_{12}$ both have characteristic polynomial
$(x-1)^4$ and zero signature. Their minimal polynomials are
$x-1$ and $(x-1)^2$, respectively: $E_{12}\ne0$ and
$E_{12}^2=0$. Thus the two matrices give distinct conjugacy
classes in the same fibre of the displayed map.
\end{enumerate}
\end{proof}

\subsection{Hopf curves in the fibre boundary}\label{subsec:II:hopf-component-boundaries}

\begin{lemma}[Central primary coverings of Hopf surfaces]\hypertarget{tag-P327}{}\label{lem:II:central-primary-hopf-cover}
\StatementRecord{lem:II:central-primary-hopf-cover}{proved}{P327}{\cite[arXiv v1, \S2.3.1, p.~8]{AbashevaDeev2024AB}; \cite[arXiv v2, normal forms (1), p.~3; Theorem~2.1, p.~5]{Khaled2024Hopf}}
Let $H=W/\Gamma$ be a compact Hopf surface, with universal cover
$W=\C^2\setminus\{0\}$. There is a contraction $F\in\Gamma$
such that $\langle F\rangle$ is central and has finite index in
$\Gamma$. The quotient map
$\pi:P=W/\langle F\rangle\to H$ is a finite unramified primary
Hopf covering. For every finite connected unramified cover $p:P'\to H$, pullback identifies $\mathcal M(H)$ with a subfield of $\mathcal M(P')$ over which $\mathcal M(P')$ is algebraic. Thus $a(P')=a(H)$; in particular $a(P)=a(H)$ for the central primary cover $\pi$.
\end{lemma}
\begin{proof}
The finite primary-cover description in
\cite[arXiv v1, \S2.3.1, p.~8]{AbashevaDeev2024AB}
gives a finite unramified cover $P_0\to H$; use the identity
when $H$ is primary. Its covering injection
$\pi_1(P_0)\hookrightarrow\pi_1(H)=\Gamma$ identifies an infinite
cyclic subgroup of index equal to the covering degree.
The primary normal forms in
\cite[arXiv v2, normal forms (1), p.~3; Theorem~2.1, p.~5]{Khaled2024Hopf}
identify a generator with a contraction of $W$.
The intersection of the finitely many conjugates of this cyclic
subgroup is a normal subgroup of finite index in $\Gamma$ and
of finite index in the original cyclic subgroup. It is therefore
$\langle F\rangle$ for a positive power $F$ of that contraction.

Every deck transformation and its inverse extend holomorphically
across the origin by Hartogs. The extensions are inverse
automorphisms of $\C^2$, since their compositions are the identity
on $W$ and hence everywhere. Each fixes the origin: its restriction
maps $W$ bijectively to $W$, so the remaining point maps to itself.
The chain rule therefore makes
$\gamma\mapsto\log|\det(d\gamma)_0|$ a homomorphism
$\Gamma\to\R$. It is invariant under conjugation and is
nonzero on the contracting $F$. Conjugation on the normal
infinite cyclic group can send $F$ only to $F$ or $F^{-1}$;
the character excludes $F^{-1}$. Thus $F$ is central.
The action of $\Gamma/\langle F\rangle$ on
$P=W/\langle F\rangle$ is free: a fixed orbit would give a deck
transformation $F^{-n}\gamma$ fixing a point of $W$, so
$\gamma=F^n$. This yields the asserted finite unramified map.

For $g\in\mathcal M(P')$, take the elementary symmetric functions of its values on the finitely many local sheets of the given finite unramified cover $p:P'\to H$. They are unchanged by the sheet permutations, so glue to meromorphic functions on $H$. Their monic polynomial annihilates $g$. Since $p$ is surjective, pullback is injective. This proves the asserted algebraic field inclusion for this $p$. Algebraic extensions have the same transcendence degree over $\C$, giving $a(P')=a(H)$. Taking $p=\pi$ gives the central-cover instance.
\end{proof}

\begin{lemma}[Cohomology restrictions on Hopf surfaces]\hypertarget{tag-P337}{}\label{lem:II:hopf-curve-restriction}
\StatementRecord{lem:II:hopf-curve-restriction}{proved}{P337}{\cref{cdp:st:D.1.7,lem:II:central-primary-hopf-cover,lem:II:hopf-curves}; \cite[Theorem~9, pp.~232--233]{Kato1975Hopf}; \cite[arXiv v2, \S4.2 and Remark~4.5, p.~18]{Khaled2024Hopf}}
For every Hopf surface $H$ one has
$b_1(H)=1$, $b_2(H)=0$, $e(H)=0$ and
$h^\bullet(H,\OO_H)=(1,1,0)$.
For every irreducible compact curve $i:C\hookrightarrow H$, the map
\[
 i^*:H^1(H,\OO_H)\longrightarrow H^1(C,\OO_C)
\]
is an isomorphism of one-dimensional complex vector spaces.

For a primary surface $H=W/\langle F\rangle$, where
$W=\C^2\setminus\{0\}$ and $F$ is a contraction, its integral
cohomology is free of ranks $(1,1,0,1,1)$. The additive deck
cocycle $F\mapsto1$ represents a basis $\eta_H$ of
$H^1(H,\OO_H)$. If a connected component of the inverse image
of $C$ has stabilizer $\langle F^d\rangle$, its parametrization
by $\C^*$ identifies $F^d$ with multiplication by a constant
$\lambda$ of modulus less than one, and the pulled-back class
is the cocycle $\lambda\mapsto d$.
For an axis $C=\C^*/\lambda^{\Z}$, inversion $[u]\mapsto[u^{-1}]$
pulls the cocycle $\lambda\mapsto1$ to its negative.

If $\pi:H'\to H$ is a finite unramified primary cover and
$C'$ is a component of $\pi^{-1}(C)$, these restriction maps
commute with pullback along $\pi$ and $\pi_C:C'\to C$;
both pullbacks on $H^1(\OO)$ are isomorphisms.
\end{lemma}

\begin{proof}
Theorem~9 of \cite[pp.~232--233]{Kato1975Hopf} presents $H$
as a bundle over $S^1$ with fibre $S^3/G$, for a finite free
group $G$. Transfer from $S^3$ and the bundle cohomology
sequence give $b_1(H)=1$, $b_2(H)=0$ and $e(H)=0$.
For a primary surface the finite group is trivial; orientation
preservation on the sphere gives the integral groups stated above.

For a primary surface, the character-to-line-bundle isomorphism
in \cite[arXiv v2, \S4.2 and Remark~4.5, p.~18]{Khaled2024Hopf}
is
\[
 \C^*\longrightarrow\Pic(H),\qquad
 \zeta\longmapsto
 [\,(W\times\C)/((x,v)\sim(Fx,\zeta v))\,].
\]
Its differential at $\zeta=1$, using $\zeta=e^t$, sends
$1\in\C$ to the additive deck cocycle $F\mapsto1$.
The tangent space of $\Pic^0(H)$ is $H^1(H,\OO_H)$,
so this is the asserted basis and $h^1(H,\OO_H)=1$.

For an arbitrary Hopf surface choose the finite unramified
primary cover supplied by \cref{lem:II:central-primary-hopf-cover}.
The identity $\operatorname{Tr}\circ\pi^*=(\deg\pi)\operatorname{id}$
makes pullback on $H^1(\OO)$ injective, so its source has
dimension at most one. The coefficient injection of
\cref{cdp:st:D.1.7}, applied to this $H$ with $b_1(H)=1$,
makes the dimension at least one. Thus $\pi^*$ is an
isomorphism and $h^1(H,\OO_H)=1$ for every Hopf surface.
Since $b_2(H)=0$, the rational class $c_1(H)$ has square zero;
also $\int_Hc_2(H)=e(H)=0$. Surface Riemann--Roch gives
$\chi(\OO_H)=0$, and hence $h^2(H,\OO_H)=0$.

By \cref{lem:II:hopf-curves}, the curve $C$ is smooth elliptic.
First suppose that $H$ is primary. The image of
$\pi_1(C)=\Z^2\to\pi_1(H)=\langle F\rangle$
is nonzero: a zero image would lift $C$ to $W$, where both
coordinate functions on the compact curve would be constant.
Write this image as $\langle F^d\rangle$ with $d>0$.
The associated connected covering of $C$ is obtained from its
universal cover $\C$ by a rank-one sublattice. Exponentiation
therefore identifies it with $\C^*$, and the remaining deck
transformation $F^d$ acts by $u\mapsto\lambda u$.
The compact elliptic quotient requires $|\lambda|\ne1$;
replace $u$ by $u^{-1}$ when necessary to have $|\lambda|<1$.
This covering is exactly a connected component of the inverse
image of $C$ in $W$, by the covering subgroup construction.
In the axis coordinates it is the inclusion of that axis;
for a resonant curve it is its parametrization
$(z,w)=(a u^p,b u^q)$.

Restriction of $\eta_H$ to the subgroup generated by $F^d$
is the cocycle with value $d$. If its class on $C$ were zero,
a holomorphic function on $\C^*$ would satisfy
\begin{equation}
                         b(\lambda u)-b(u)=d.
 \hypertarget{tag-00N3}{}\label{cdp:eq:C.HR.4}
\end{equation}
The constant coefficient of the Laurent series on the left is
zero, whereas $d>0$. This proves nonvanishing and hence the
primary restriction isomorphism. For an axis, inversion
conjugates multiplication by $\lambda$ to multiplication by
$\lambda^{-1}$. Pullback therefore evaluates the additive
character on the inverse generator and changes its sign.

For the cover $\pi:H'\to H$, a component $C'$ of
$\pi^{-1}(C)$ maps onto $C$ by a finite unramified elliptic
cover. Naturality gives the commutative square
\[
\begin{array}{ccc}
 H^1(H,\OO_H)&\xrightarrow{\ i^*\ }&H^1(C,\OO_C)\\
 \pi^*\downarrow&&\downarrow\pi_C^*\\
 H^1(H',\OO_{H'})&\xrightarrow{\ (i')^*\ }&H^1(C',\OO_{C'}).
\end{array}
\]
The surface pullback is the isomorphism already proved.
The curve pullback is injective by its own trace identity
$\operatorname{Tr}_{\pi_C}\pi_C^*=(\deg\pi_C)\operatorname{id}$,
and hence is an isomorphism between one-dimensional spaces.
The lower restriction is the primary isomorphism, so the
upper restriction is an isomorphism as well.
\end{proof}

\subsection{Effective anticanonical boundaries on Hopf surfaces}\label{subsec:II:anticanonical-hopf-boundaries}

\begin{lemma}[Effective anticanonical Hopf boundaries]\hypertarget{tag-P325}{}\label{lem:II:reduced-anticanonical-hopf}
\StatementRecord{lem:II:reduced-anticanonical-hopf}{proved}{P325}{\cref{lem:II:central-primary-hopf-cover}; \cite[arXiv v1, \S2.2.2, p.~4; \S2.3.1, p.~8]{AbashevaDeev2024AB}; \cite[arXiv v2, normal forms (1), p.~3; Theorem~2.1 and Remark~2.2, p.~5]{Khaled2024Hopf}}
Let $H$ be a compact Hopf surface and let $D$ be a nonempty
effective integral divisor whose support consists of disjoint smooth
elliptic curves. Suppose $K_H(D)\simeq\OO_H$. Exactly one of the
following alternatives holds:
\begin{enumerate}
\item $D$ is reduced. The pair $(H,D)$ is a diagonal primary or
secondary Hopf surface with its two axes individually marked.
There are coordinates on $W=\C^2\setminus\{0\}$ in which the
full deck group of $W\to H$ is diagonal and the inverse image
of $D$ is $\{zw=0\}\cap W$. The two axes have distinct images
on $H$, and $dz/z\wedge dw/w$ descends to a nowhere-zero
section of $K_H(D)$.
\item $D=dC$ for one smooth elliptic curve $C$ and an integer $d\ge2$.
There is a finite unramified central primary cover $\pi:H'\to H$ for
which $\pi^{-1}(C)$ is one curve. In coordinates marking this
curve by $w=0$, the primary contraction is either
$(z,w)\mapsto(\beta^r z+w^r,\beta w)$ or
$(z,w)\mapsto(\beta^r z,\beta w)$, with $r\ge1$,
$0<|\beta|<1$ and $d=r+1$.
\end{enumerate}
In particular, a coefficient-one component of $D$ forces the
first alternative.
\end{lemma}
\begin{proof}
Write $H=W/\Gamma$ and apply
\cref{lem:II:central-primary-hopf-cover}. It gives a central
contraction $F$ and a finite unramified primary cover
$\pi:P=W/\langle F\rangle\to H$.
Put $D_P=\pi^*D$. The local unramified differential preserves
each divisor coefficient and gives
$K_P(D_P)\simeq\pi^*(K_H(D))\simeq\OO_P$.

Use the complete primary normal-form list in
\cite[arXiv v2, formula~(1), p.~3; Theorem~2.1, p.~5]{Khaled2024Hopf}.
In the nonlinear case write
$F(z,w)=(\beta^a z+w^a,\beta w)$, $a\ge1$.
Its unique curve is $C_P=\{w=0\}/\langle F\rangle$ by
\cite[arXiv v1, \S2.2.2, p.~4]{AbashevaDeev2024AB},
so $D_P=dC_P$ for a positive integer $d$.
The determinant of $dF$ is $\beta^{a+1}$; thus
$w^{a+1}\partial_z\wedge\partial_w$ is an anticanonical
section with divisor $(a+1)C_P$.
The line $\OO_P(C_P)$ has infinite order: a trivial positive
power would turn its canonical section into a holomorphic
function with a nonempty zero divisor on the compact connected
surface. Therefore $K_P(D_P)\simeq\OO_P$ forces $d=a+1\ge2$.
Unramified descent gives the second alternative on $H$.
Here $H'=P$, $r=a$, and the same multiplier $\beta$ give
the asserted central primary cover and its marked contraction.

In diagonal coordinates $F(z,w)=(\alpha z,\beta w)$,
an anticanonical section lifts to
$g(z,w)\partial_z\wedge\partial_w$. Hartogs extends $g$
across zero, and equivariance gives
\begin{equation}
                        g(\alpha z,\beta w)
                                      =\alpha\beta g(z,w).
\hypertarget{tag-00DP}{}\label{cdp:eq:C.CZ.11}
\end{equation}
If there is no positive multiplicative relation between
$\alpha$ and $\beta$, comparison of monomials gives only $zw$.
Indeed a nonzero coefficient of $z^iw^j$ requires
$\alpha^{i-1}\beta^{j-1}=1$; since both multipliers contract,
any solution other than $(i,j)=(1,1)$ gives such a relation.

Otherwise choose positive integers $u,v$ with
$\alpha^u=\beta^v$, and write $u=eq$, $v=ep$ with $(p,q)=1$.
The central power $F^e$ has multipliers
$\alpha_1=\alpha^e$, $\beta_1=\beta^e$ satisfying
$\alpha_1^q=\beta_1^p$.
Choose integers $s,t$ with $sp+tq=1$ and put
$\gamma=\alpha_1^s\beta_1^t$.
Then $\gamma^p=\alpha_1$, $\gamma^q=\beta_1$, and $0<|\gamma|<1$.
The quotient map $W/\langle F^e\rangle\to W/\langle F\rangle$
is finite unramified, and $F^e$ is still central in $\Gamma$.
Pullback to this cover preserves every coefficient of $D_P$.
Writing $F$ for its central generator $F^e$, its coordinate form is
\begin{equation}
                 F(z,w)=(\gamma^p z,\gamma^q w),
       \qquad p,q\ge1,\quad\gcd(p,q)=1,\quad |\gamma|<1 .
                                                               \hypertarget{tag-00AX}{}\label{cdp:eq:C.AS.8}
\end{equation}
The monomials of $g$ now have weight $p+q$. A weighted homogeneous
polynomial factors, after removing its axis monomial, into binomials
$z^q-cw^p$, each of weight $pq$: the ratios of two monomials of
equal weight are powers of $z^q/w^p$, and the resulting one-variable
polynomial factors over $\C$. Consequently, when the zero divisor
is $a$ times an orbit of $m$ such non-axis curves, its weight gives
\begin{equation}
                              ampq=p+q.                 \hypertarget{tag-00AY}{}\label{cdp:eq:C.AS.9}
\end{equation}
The equation $pi+qj=p+q$ is equivalent to
$i=1+q\ell$, $j=1-p\ell$, with $\ell\in\Z$.
For $p,q>1$, nonnegative exponents force $\ell=0$.
For $p>1,q=1$, the full section is
\begin{equation}
                        g=w(az+bw^p).
\hypertarget{tag-00R0}{}\label{cdp:eq:C.ML.9}
\end{equation}
When $a=0$, its divisor is $(p+1)$ times the single axis.
When $a\ne0$, the two factors are distinct and simple;
$z'=z+(b/a)w^p$, $w'=w$ sends them to the axes and commutes
with the contraction. Its inverse is
$z=z'-(b/a)(w')^p$, $w=w'$.
For $p=1,q>1$ use the exchanged coordinates and the exchanged
triangular map. For $p=q=1$, $g$ is a nonzero quadratic.
A repeated linear factor gives twice one curve; two distinct
linear factors give an invertible linear coordinate map,
with inverse its matrix inverse, commuting with the scalar contraction.
Thus on this finite unramified primary cover the divisor is
either reduced with two curves, or a positive multiple of one
curve with coefficient at least two. Since the cover is
surjective and unramified, the same reducedness alternative
and each coefficient descend to $H$; in the latter case the
image of that single curve is all of the support of $D$.
In particular a coefficient-one component excludes the latter case.
In the nonreduced case its pullback to this primary cover also
has a coefficient greater than one, so the classification just
proved leaves exactly one curve. The nonlinear form gives
$d=r+1$ on its axis. In the diagonal calculation a single curve
comes from $w^{p+1}$ when $q=1$, from the exchanged case when
$p=1$, or from a repeated linear factor when $p=q=1$.
The coordinate changes above therefore mark it by $w=0$ and
put the contraction in the second form of the statement,
with $r=d-1$. These are the asserted primary-cover maps.

For the reduced case, the following coordinate maps preserve its full marking.
Use the coordinates just constructed, so the pulled-back
boundary is the two axes and the central contraction is diagonal.
Every deck lift preserves their union and commutes with that
central contraction. A lift preserving each axis has form
$(zU(z,w),wV(z,w))$ with holomorphic units $U,V$.
Commutation makes both units invariant under the contraction;
iteration to zero makes them constant.
A lift interchanging the axes has form $(wU,zV)$.
Comparison of linear terms forces the two multipliers of the
central contraction to agree, and the same iteration makes
$U,V$ constant. It is a diagonal map followed by interchange.
The logarithmic form
\begin{equation}
                             \frac{dz}{z}\wedge\frac{dw}{w}.
                                                               \hypertarget{tag-00AI}{}\label{cdp:eq:C.AS.11}
\end{equation}
is a nowhere-zero section of the pulled-back canonical-boundary
line, and spans its section space because the primary surface
is compact and connected.
Diagonal lifts fix it; interchange negates it.
The nonzero section pulled back from $K_H(D)$ is invariant,
so every deck lift preserves each axis. The full group
$\Gamma$ is diagonal and the two images are distinct.

A diagonal group containing an infinite cyclic subgroup of finite
index is finitely generated abelian, so
$\Gamma=\langle f\rangle\oplus T$ with $T$ finite.
Choose the sign of $f$ so that a positive power is a positive
power of the central contraction times an element of $T$.
Both multipliers of $f$ then have modulus less than one.
The finite group $T$ acts freely on $W$, so each coordinate
character is injective; consequently $T$ is cyclic.
The equality $H=(W/\langle f\rangle)/T$ gives the stated
diagonal primary or secondary presentation with its marked axes.
The logarithmic form descends through this free quotient.
\end{proof}


\begin{proposition}[Hopf curves in the fibre boundary]\hypertarget{tag-01DH}{}\label{cdp:st:C.2.UN.38}
\StatementRecord{cdp:st:C.2.UN.38}{proved}{01DH}{\SeriesRef{cdp:eq:C.HB.2}, \SeriesRef{cdp:eq:C.HB.3}, \SeriesRef{cdp:eq:C.HB.5}, \SeriesRef{cdp:eq:C.HB.4}; \cite[Theorem~2.1]{Khaled2024Hopf}; \cite[\S1.1]{Nakamura2008VII}}

Let $f:U\to\Delta$ be a proper degeneration of
connected complex two-tori from a smooth threefold,
and write its full central fibre as $\sum_i m_iS_i$.
On the open subset of $S_i$ where the reduced
support is smooth, its canonical line has a trivial $m_i$-th power in a neighborhood
of every compact subset. In particular, a
smooth compact elliptic curve $C$ in that open
subset satisfies
\begin{equation}
                         N_{C/S_i}^{\otimes m_i}
                                      \simeq\OO_C.    \hypertarget{tag-00LF}{}\label{cdp:eq:C.HB.1}
\end{equation}

Suppose the normalization of a central component
is itself a smooth minimal Hopf surface of
algebraic dimension zero. Then every elliptic
curve on that normalization is contained in
the inverse image of the singular locus of
the reduced support. For a diagonal primary
Hopf surface this includes both axes. For a
non-diagonal primary Hopf surface it includes
its unique axis; in this case the conclusion
uses the full canonical neighborhood and
the explicit higher-jet calculation below.
The same conclusion descends through every
finite primary Hopf covering.
\end{proposition}

\begin{proof}

\begin{enumerate}[label=\arabic*.,leftmargin=*]
\item Apply the canonical-evaluation assertion of
\cref{cdp:st:C.2.UN.17} to this map $f:U\to\Delta$.
Its generator $\omega$ of $f_*K_U$, evaluated on $U$, has divisor
\begin{equation}
                   \operatorname{div}(\omega)
                               =\sum_i k_iS_i,
                        \qquad k_i\ge0.              \hypertarget{tag-00LG}{}\label{cdp:eq:C.HB.2}
\end{equation}
Let $W\subset U$ be a neighborhood of a
compact subset in the indicated regular
part of $S=S_i$, small enough to avoid all
other components and the singularities of
$S$. On $W$ the fibre equation has
divisor $mS$, where $m=m_i$. Thus it gives
the trivialization of $\OO_W(mS)$, and
\eqref{cdp:eq:C.HB.2} gives
\begin{equation}
 \begin{gathered}
 K_U|_W\simeq\OO_W(k_iS),\\
 K_S|_{S\cap W}\simeq
                  \OO_W((k_i+1)S)|_{S\cap W},\\
 (K_S|_{S\cap W})^{\otimes m}\simeq\OO_{S\cap W}.
 \end{gathered}
                                                               \hypertarget{tag-00LH}{}\label{cdp:eq:C.HB.3}
\end{equation}
These are line-bundle isomorphisms on the
neighborhood, obtained from the canonical and fibre sections. Adjunction
for an elliptic $C\subset S\cap W$ gives
$N_{C/S}=K_C\otimes K_S^{-1}|_C$.
Since $K_C=\OO_C$, \eqref{cdp:eq:C.HB.1} follows.

\item Consider first the diagonal primary
Hopf surface
\begin{equation}
 H_{\alpha,\beta}=(\C^2\setminus\{0\})/
       \langle(z,w)\mapsto(\alpha z,\beta w)\rangle,
                    \qquad0<|\alpha|,|\beta|<1.
                                                               \hypertarget{tag-00LI}{}\label{cdp:eq:C.HB.4}
\end{equation}
The axis $C_\alpha$ is $\C^*/\alpha^{\Z}$,
and its normal line is the quotient
of $\C^*\times\C$ by
$(z,v)\mapsto(\alpha z,\beta v)$.
If a positive tensor power $n$ of this
line is trivial, a nowhere-zero holomorphic
function $a$ on $\C^*$ satisfies
$a(\alpha z)=\beta^n a(z)$.
In its Laurent expansion each nonzero
coefficient of $z^j$ requires
$\alpha^j=\beta^n$. Since $|\alpha|<1$,
there is at most one such exponent, and
there must be one. Its exponent is positive
because $|\beta|<1$ and $n>0$.
Conversely that equality gives the
explicit trivialization $a(z)=z^j$.
It also gives the nonconstant invariant
meromorphic function $w^n/z^j$ on the
Hopf surface. Algebraic dimension zero
therefore makes the normal of each axis
non-torsion. Formula \eqref{cdp:eq:C.HB.1} excludes
an entire such axis from the regular
part of the reduced support.

\item The non-diagonal primary normal form
requires a further argument. The primary
normal forms and finite unramified cover
used here are the explicit ones recalled
in \cite[Theorem~2.1, arXiv v2 p.~5; normal forms (1), p.~3]{Khaled2024Hopf}
and \cite[\S1.1, Theorem~1.1.2, p.~2]{Nakamura2008VII}.
Write the non-diagonal contraction
as
\begin{equation}
 H=(\C^2\setminus\{0\})/\langle F\rangle,
 \qquad F(z,w)=(\beta^r z+w^r,\beta w),
                  \quad r\ge1,\quad0<|\beta|<1,
                                                               \hypertarget{tag-00LJ}{}\label{cdp:eq:C.HB.5}
\end{equation}
and let $C$ be the image of $w=0$.
The invariant meromorphic two-form
$dz\wedge dw/w^{r+1}$ gives the canonical divisor
\begin{equation}
                         K_H\simeq\OO_H(-(r+1)C).
                                                               \hypertarget{tag-00LK}{}\label{cdp:eq:C.HB.6}
\end{equation}
The normal-line quotient has multiplier $\beta$ over
$z\mapsto\beta^r z$; the function $z$ trivializes its $r$th
tensor power. The obstruction uses the canonical line on a
neighbourhood:
\eqref{cdp:eq:C.HB.3} on a neighborhood of $C$ would
make $\OO(m(r+1)C)$ trivial there.
Its canonical divisor section would
then be a holomorphic function with
positive divisor $m(r+1)C$ on that
neighborhood. We show directly that
no such function exists.

\item Lift it to the invariant neighborhood
of $\C^*\times\{0\}$ in the covering
space and expand in the normal variable:
\begin{equation}
 h(z,w)=\sum_{l\ge j}h_l(z)w^l,
       \qquad j>0,\qquad h_j\ne0,
       \qquad h_l\in\OO(\C^*),
       \qquad h\circ F=h.                           \hypertarget{tag-00LL}{}\label{cdp:eq:C.HB.7}
\end{equation}
The radius of convergence may depend
on $z$; the Taylor coefficients are
nevertheless holomorphic on all of
$\C^*$ because the lifted neighborhood
contains that entire axis. The leading
coefficient satisfies
$\beta^j h_j(\beta^r z)=h_j(z)$.
Its Laurent expansion forces
\begin{equation}
                      j=rk,\qquad
                      h_j(z)=c z^{-k},
                      \qquad k\ge1,\quad c\ne0.
                                                               \hypertarget{tag-00LM}{}\label{cdp:eq:C.HB.8}
\end{equation}
At order $j+r$, the same invariant
function equation is exactly
\begin{equation}
 \beta^{j+r}h_{j+r}(\beta^r z)-h_{j+r}(z)
                    =k c\beta^{-r}z^{-k-1}.          \hypertarget{tag-00LN}{}\label{cdp:eq:C.HB.9}
\end{equation}
Indeed the only extra contribution
at that order is the first derivative
of $h_j(\beta^r z+w^r)\beta^j$;
all other Taylor terms start at
higher order. In the Laurent
expansion of the left side, the
coefficient multiplying $z^{-k-1}$
is zero, since
$\beta^{j+r+r(-k-1)}-1=0$.
The right side has the nonzero
coefficient $k c\beta^{-r}$.
This contradiction proves the
neighborhood obstruction for
the non-diagonal axis as well.

\item Let $\nu:H\to S$ be the smooth minimal Hopf
normalization. Apply the conductor purity in
\cref{lem:II:conductor-duality} to the reduced hypersurface $S$:
its nonnormal locus $C\subset S$ is a curve wherever nonempty.
Its finite inverse image in $H$, together with the pullbacks of
all other central components, defines $B\subset H$. Each of the
latter pullbacks is a Cartier divisor on the smooth $H$, so $B$
is a union of compact curves.
On $H\setminus B$, the normalization map is an isomorphism onto
the regular locus of the reduced fibre: normalization is an
isomorphism over the normal locus, and $H$ is smooth.

\item Apply \cref{lem:II:central-primary-hopf-cover} to this
Hopf normalization $H$ and write its finite unramified primary
cover as $\pi:H'\to H$. Put $B'=\pi^{-1}B$.
Its pullback of meromorphic-function fields is exactly the one
in that lemma, so $a(H')=a(H)=0$.

In the diagonal case, a relation $\alpha^q=\beta^p$ with
positive integers $p,q$ would make $z^q/w^p$ descend to a
nonconstant meromorphic function on $H'$. Since $a(H')=0$,
there is no such relation. Both multipliers have absolute value
less than one, so every relation with nonzero integer exponents
can be inverted to have positive exponents. The curve descriptions in
\cite[arXiv v1, \S2.2.2, p.~4; \S2.2.6, p.~5]{AbashevaDeev2024AB}
therefore give exactly the two axis curves in the diagonal case
and the single axis curve in the non-diagonal case.
Apply \cref{lem:II:hopf-curves} to $H'$ for the disjointness of
distinct curves. An axis not contained in $B'$ is thus disjoint
from $B'$.
Since $\pi$ is unramified, $K_{H'}\simeq\pi^*K_H$;
the normalization isomorphism over $H\setminus B$ and this
pullback give the canonical-neighbourhood trivialization
\eqref{cdp:eq:C.HB.3} along that axis. It contradicts the diagonal
normal-line calculation or the non-diagonal higher-jet calculation.
Hence every axis lies in $B'$. A curve on $H$ outside $B$ would
have a compact curve component in its finite inverse image outside
$B'$. The same axis list excludes it. Thus every curve on $H$
lies in the fibre boundary.

\end{enumerate}
\end{proof}

\subsection{Resolution index and Euler number}\label{subsec:II:resolution-euler-index}

\begin{proposition}[Resolution index and Euler number]\hypertarget{tag-01DI}{}\label{cdp:st:C.2.UN.39}
\StatementRecord{cdp:st:C.2.UN.39}{proved}{01DI}{\SeriesRef{cdp:st:C.2.UN.11}, \SeriesRef{cdp:eq:C.ZC.5}, \SeriesRef{cdp:eq:C.ZC.6}, \SeriesRef{cdp:eq:C.ZC.8}, \SeriesRef{cdp:eq:C.ZC.7}, \SeriesRef{cdp:eq:C.ZC.9}, \SeriesRef{cdp:st:C.2.UN.4}; \cite[\S2, pp.~237--243; \S3, pp.~243--245]{ACampo1975}}

Let $f:U\to\Delta$ be a proper holomorphic
two-torus degeneration with smooth total threefold,
and let $S=f^{-1}(0)_{\rm red}$. Choose an embedded
resolution $\pi:\widetilde U\to U$ as a sequence
of blowups with smooth centres contained in the
successive central supports. Write its full
central divisor as $\sum_jN_jE_j$ with simple
normal crossings. Let $\mathcal D$ be the set
of connected double curves of $\bigcup_jE_j$,
and let $t_3$ be its number of triple points.
Count every point centre at every stage of
$\pi$, obtaining an integer $p_\pi$, and list
every connected smooth curve centre as $C_\alpha$.
Define the integer
\begin{equation}
 I_\pi=
     \sum_{D\in\mathcal D}\bigl(g(D)-1\bigr)+t_3
       +p_\pi+\sum_\alpha\bigl(1-g(C_\alpha)\bigr).
                                                               \hypertarget{tag-018W}{}\label{cdp:eq:C.ZC.1}
\end{equation}
It is independent of the chosen embedded
resolution, including additional vertical
blowups used to obtain another such resolution.

Let $T$ be the monodromy on $H^1(F,\Z)$
of a nearby torus, and let $M$ be the order of
its semisimple part. Put
\begin{equation}
             j(T)=\frac1M\sum_{k=1}^{M}\det(I-T^k).
                                                               \hypertarget{tag-018X}{}\label{cdp:eq:C.ZC.2}
\end{equation}
Then
\begin{equation}
                e(S)=j(T)-2I_\pi,\qquad
                e(S)=0\quad\Longleftrightarrow\quad
                         I_\pi=\frac{j(T)}2 .        \hypertarget{tag-018Y}{}\label{cdp:eq:C.ZC.3}
\end{equation}
The constant $j(T)$ belongs to $\{0,4,6,8\}$.
For all nineteen monodromy polynomials its
value and the required zero index are
\begin{equation}
 \begin{array}{c|r|r|r}
\det(xI-T)&M&j(T)&I_\pi\text{ for }e(S)=0\\ \hline
\Phi_{1}^{4}&1&0&0\\
\Phi_{1}^{2} \Phi_{2}^{2}&2&0&0\\
\Phi_{1}^{2} \Phi_{3}&3&0&0\\
\Phi_{1}^{2} \Phi_{4}&4&0&0\\
\Phi_{1}^{2} \Phi_{6}&6&0&0\\
\Phi_{2}^{4}&2&8&4\\
\Phi_{3}^{2}&3&6&3\\
\Phi_{2}^{2} \Phi_{4}&4&4&2\\
\Phi_{4}^{2}&4&6&3\\
\Phi_{5}&5&4&2\\
\Phi_{2}^{2} \Phi_{3}&6&4&2\\
\Phi_{2}^{2} \Phi_{6}&6&4&2\\
\Phi_{3} \Phi_{6}&6&4&2\\
\Phi_{6}^{2}&6&6&3\\
\Phi_{8}&8&4&2\\
\Phi_{10}&10&4&2\\
\Phi_{3} \Phi_{4}&12&4&2\\
\Phi_{4} \Phi_{6}&12&4&2\\
\Phi_{12}&12&4&2\\
\end{array}
                                                               \hypertarget{tag-018Z}{}\label{cdp:eq:C.ZC.4}
\end{equation}
Thus this is a necessary and sufficient
geometric test for the zero Euler number of
every fibre, using only the indicated
curve genera, triple points, resolution centres
and finite monodromy constant. Arbitrary
component and centre multiplicities are
included through the full divisor and the
blowup maps. \end{proposition}

\begin{proof}

\begin{enumerate}[label=\arabic*.,leftmargin=*]
\item An embedded resolution of the specified kind
exists by the analytic resolution theorem used
in \ref{cdp:st:C.2.UN.11}, with all centres over the central
support and with an isomorphism over
$\Delta^*$. Smooth divisor centres need
not be included, since their blowups are
isomorphisms. The remaining connected centres
in a threefold are compact smooth curves or
points. A disconnected centre can be counted
by its disjoint connected components; their
blowups commute. In particular every genus
and count in \eqref{cdp:eq:C.ZC.1} is attached to a compact centre or curve in this resolution.

\item For the full resolved divisor, put
\[
             E_j^\circ=E_j\setminus\bigcup_{i\ne j}E_i.
\]
Apply \cref{night:P4.13} with its map $f$, resolution $\pi$
and cohomology monodromy $T$ equal to those here. The cup-product isomorphisms
$\bigwedge^qH^1(F,\Z)\simeq H^q(F,\Z)$ intertwine its exterior
powers with monodromy on every degree. Thus
\eqref{cdp:eq:C.MS.6} says for every positive $k$ that
\begin{equation}
         \det(I-T^k)=\sum_{N_j\mid k}N_j e(E_j^\circ).
                                                               \hypertarget{tag-0190}{}\label{cdp:eq:C.ZC.5}
\end{equation}
Here the full integer $N_j$ is used for each
component. Choose $L$ divisible by $M$ and
by every $N_j$. The determinant on the left
depends only on the eigenvalues of $T$,
so is periodic with period $M$ even when
its unipotent part is nontrivial. Average
\eqref{cdp:eq:C.ZC.5} over $1\le k\le L$. Exactly $L/N_j$
of these integers are divisible by $N_j$.
It follows that
\begin{equation}
 \begin{split}
 j(T)
 &=\frac1L\sum_{k=1}^{L}\det(I-T^k)\\
 &=\sum_j\frac{N_j}{L}\frac{L}{N_j}e(E_j^\circ)
   =\sum_j e(E_j^\circ).
 \end{split}                                                   \hypertarget{tag-0191}{}\label{cdp:eq:C.ZC.6}
\end{equation}
This calculation includes components whose
multiplicity does not divide $M$; choosing
$L$ retains them before the averaging.

\item Let $\widetilde S=\bigcup_jE_j$. Its
remaining strata are the double curves
with their triple points removed, and
the triple points themselves. Every
triple point lies on exactly three
double curves. Since each connected
smooth compact double curve $D$ has
Euler number $2-2g(D)$, additivity gives
\begin{equation}
 \begin{split}
 e(\widetilde S)
 &=\sum_j e(E_j^\circ)
       +\sum_{D\in\mathcal D}(2-2g(D))-3t_3+t_3\\
 &=j(T)-2\left(
           \sum_{D\in\mathcal D}(g(D)-1)+t_3\right).
 \end{split}                                                   \hypertarget{tag-0192}{}\label{cdp:eq:C.ZC.7}
\end{equation}
The open surface contribution is the finite monodromy average
\eqref{cdp:eq:C.ZC.6}.

\item For each blowup $\sigma:U'\to U_0$ in $\pi$, let
$Z\subset S_0$ be its connected smooth centre and
$S'=(\sigma^{-1}S_0)_{\rm red}$. Apply
\cref{lem:II:blowup-topology} with
$(Y,Z,S,Y',S',r)=(U_0,Z,S_0,U',S',c)$, where
$c=\operatorname{codim}_{\C}(Z,U_0)$. Its full-inverse-image formula
\eqref{cdp:eq:C.ZC.8} gives $e(S')-e(S_0)=(c-1)e(Z)$.
This is $2$ for a point and $2-2g(Z)$ for a smooth connected
curve. The local exceptional factor of order
$\mu=\operatorname{ord}_Z(f_0)$, where $f_0$ is the pullback of $f$
to $U_0$, changes the full divisor coefficient;
the reduced inverse image still contains the entire projective
fibre $\PP^{c-1}$ used in that formula.

\item Iterating the proper maps in
$\pi$ gives
\begin{equation}
 e(\widetilde S)=e(S)+2p_\pi
                        +2\sum_\alpha(1-g(C_\alpha)).
                                                               \hypertarget{tag-0194}{}\label{cdp:eq:C.ZC.9}
\end{equation}
Combining \eqref{cdp:eq:C.ZC.7} and \eqref{cdp:eq:C.ZC.9}
proves \eqref{cdp:eq:C.ZC.3} with exactly the
integer in \eqref{cdp:eq:C.ZC.1}. Every such
resolution is an isomorphism over
the punctured disc and therefore
has the same monodromy.
The formula consequently identifies
$I_\pi$ with $(j(T)-e(S))/2$ for
each choice. This proves independence for every embedded resolution of the
given family.

\item The nineteen possible characteristic
polynomials and their one-period
determinants were computed in \ref{cdp:st:C.2.UN.11}.
Applying the finite average \eqref{cdp:eq:C.ZC.2}
gives the entire table \eqref{cdp:eq:C.ZC.4}.
Equivalently, $j=0$ when the polynomial
has a factor $\Phi_1$; $j=8$ for
$\Phi_2^4$; $j=6$ for
$\Phi_3^2,\Phi_4^2,\Phi_6^2$;
and $j=4$ for the ten other rows.
For each characteristic polynomial, use the direct sum $A_0$
of its integral cyclotomic companion matrices. If $\lambda_1,
\ldots,\lambda_4$ are the eigenvalues with their multiplicities,
triangularization gives
$\det(I-T^k)=\prod_{i=1}^4(1-\lambda_i^k)=\det(I-A_0^k)$.
The finite average in \eqref{cdp:eq:C.ZC.2} therefore gives the
displayed values for the full monodromy, including its unipotent
part. With $a_n=\sum_{N_j=n}e(E_j^\circ)$ as in
\cref{cdp:st:C.2.UN.11}, equation \eqref{cdp:eq:C.ZC.6}
also gives $j(T)=\sum_n a_n$.

\item Start from a smooth torus support with identity monodromy.
A point blowup
creates one rational double curve and
no triple point. The new support has
index $-1$ and Euler number $2$.
Follow it by the genus-three quartic
blowup of \ref{cdp:st:C.2.UN.4}. There are then five
rational double curves, one double
curve of genus three, and four triple
points. Their index is
$-5+(3-1)+4=1$, giving Euler number
$-2$. Viewed instead as a resolution
of the original smooth support, these
two centres contribute
$1+(1-3)=-1$ to the centre term,
so \eqref{cdp:eq:C.ZC.1} returns index zero for
that original support. These exact
charts verify both the sign of the
centre correction and its compatibility
with the universal balancing operations.

\end{enumerate}
\end{proof}

\subsection{Smooth normal-product collars}\label{subsec:II:normal-product-collars}

\begin{lemma}[Smooth normal-product collars]\hypertarget{tag-P326}{}\label{lem:II:normal-product-collars}
\StatementRecord{lem:II:normal-product-collars}{proved}{P326}{}
Let $f:X\to\Delta$ be a proper holomorphic map from a smooth complex threefold whose full central fibre $S$ is reduced and simple normal crossing, with no triple points. For a double curve $C=V_1\cap V_2$, the principal fibre gives an isomorphism
$N_{C/V_1}\otimes N_{C/V_2}\simeq\OO_C$.
A neighbourhood of $C$ is smoothly diffeomorphic over $\Delta$ to a neighbourhood of the zero section in
$N_{C/V_2}\oplus N_{C/V_1}$ with function $(x,y)\mapsto xy$. This map restricts to the identity on $C$ and identifies the two component branches and their normal lines.

The same conclusions hold for a full reduced normal crossing fibre without triple points whose normalization $\nu:\coprod_i V_i\to S$ has smooth components and whose inverse image of each connected double curve $C$ is the disjoint union of two curves, each mapped isomorphically to $C$. These two curves label the local branches; they may lie on the same normalization component. The two normal lines in the conclusions are then the normal lines of these curves in their normalization components, transported to $C$ by the two specified isomorphisms.

These collars can be chosen simultaneously on all double curves. Compact component cores away from them transport to nearby fibres. The two attaching maps on each nearby neck induce the maps of the indicated normal circle bundles, with the reversal $y=t/x$. Hence fundamental-group calculations made with these cores and normal-product necks apply to the given smoothing.
\end{lemma}
\begin{proof}
In the simple normal crossing case the two branch germs along $C$ are the restrictions of its two smooth components. In the normalization case, their two specified inverse curves label these germs globally along $C$. The labels agree on overlaps because each inverse curve maps isomorphically to $C$. The labelled branch germs therefore glue, after shrinking around the compact curve, to two smooth Cartier divisors $B_1,B_2$ in a neighbourhood $N$ of $C$, with $S\cap N=B_1\cup B_2$. The normalization map restricted near either inverse curve is biholomorphic onto this smooth branch. Its derivative identifies the corresponding normal line with $N_{C/B_i}$. This construction applies when the two inverse curves belong to one normalization component as well.

Near $C$ write $L_i=\OO_N(B_i)$ and let $s_i$ be its canonical divisor section. Since the full fibre is reduced and no other component meets $C$, the function $f$ identifies $L_1\otimes L_2$ with the trivial line and carries $s_1s_2$ to $f$. Along $C$ one has
$L_1|_C=N_{C/B_2}$ and $L_2|_C=N_{C/B_1}$.
Choose a smooth retraction $r$ to $C$ and a smooth identification
$L_1\simeq r^*(L_1|_C)$ equal to the identity on $C$. Use the product trivialization to choose the second identification, so that their tensor product preserves $f$. The map
$\Psi=(r,s_1,s_2)$, with the sections expressed in these identifications, has identity derivative along $TC$. Its two normal derivatives are the transverse divisor differentials, which identify the normal bundle with $L_1|_C\oplus L_2|_C$. Its derivative is therefore invertible along $C$. Compactness and the inverse function theorem give the smooth tubular diffeomorphism after shrinking; by construction $f=xy$ and the branches are $x=0$ and $y=0$.

The finitely many double curves are disjoint, so their neighbourhoods may be chosen disjoint. Give the two normal lines dual Hermitian metrics. On $xy=t$, the projection to $x$ identifies a neck with an annulus bundle in $N_{C/B_2}^{\times}$; projection to $y$ is its dual map $y=t/x$. Radial retraction identifies either end with the corresponding normal circle bundle. This proves the asserted maps on fundamental groups, including normal reversal.

Remove small open normal discs about the inverse double curves from every normalization component. Each resulting compact core maps biholomorphically onto its image in the regular part of $S$: normalization is an isomorphism there. Thus the cores are embedded in $X$, including when two boundary circles come from one self-identified normalization. On each such core, $f$ is a submersion because its branch has multiplicity one. Near its boundary the collar coordinates give horizontal lifts of the two real base-coordinate directions that keep the relevant nonzero normal coordinate fixed. Extend these lifts across the compact core with a partition of unity; the right-inverse condition for $df$ is affine and is preserved by this extension. Integrating along the radial segment from $0$ to a sufficiently small $t$ transports the core and agrees with the collar transport at its boundary. The inverse transport gives the inverse diffeomorphism. The attachments therefore use precisely the normal-product maps already described, and van Kampen uses those same homomorphisms for the given nearby fibre.
\end{proof}

\subsection{Marked cycles and their integral lattices}\label{subsec:II:marked-cycle-lattices}

\begin{definition}[Marked normal cycles]\hypertarget{tag-P329}{}\label{def:II:marked-normal-cycle}
\StatementRecord{def:II:marked-normal-cycle}{proved}{P329}{\cite[arXiv v1, \S2.2.7--\S2.2.9, p.~5]{AbashevaDeev2024AB}}
A marked normal cycle consists of an oriented cyclic list of $N\ge1$ compact smooth surfaces $V_1,\ldots,V_N$. Each $V_i$ is either a diagonal primary or secondary Hopf surface, with its two distinct coordinate-axis curves, or a ruled surface
\begin{equation}
                  R_i=\PP(\OO_{E_i}\oplus \mathcal L_i),
                  \qquad E_i\ \hbox{elliptic},\quad
                         \mathcal L_i\in\operatorname{Pic}^0(E_i),
                                                        \hypertarget{tag-00RR}{}\label{cdp:eq:C.MX.1}
\end{equation}
with its zero and infinity sections. We use the lines convention for $\PP$. The orientation orders these curves as incoming $C_i^-$ and outgoing $C_i^+$. Hopf surfaces of every algebraic dimension are allowed. For each edge the data include
\begin{equation}
 \phi_i:C_i^+\xrightarrow{\sim}C_{i+1}^-,
 \qquad
 \theta_i:N_{C_i^+/V_i}\otimes
       \phi_i^*N_{C_{i+1}^-/V_{i+1}}
           \xrightarrow{\sim}\mathcal O_{C_i^+}.
\hypertarget{tag-00LS}{}\label{cdp:eq:C.HC.3}
\end{equation}
The normal crossing surface $S$ is obtained by these elliptic identifications. For $N=1$, this means identifying the two distinct marked curves of the single normalization. The two inverse images of each double curve map isomorphically to it and label its two local branches. Every double curve is disjoint from the others.

The complement of the two marked curves has the logarithmic presentation
\begin{equation}
 G_i=V_i\setminus(C_i^-\cup C_i^+)
       =\mathbb C^2/\Lambda_i,\qquad \Lambda_i\simeq\mathbb Z^3,
\hypertarget{tag-00LO}{}\label{cdp:eq:C.HC.1}
\end{equation}
For $t\ne0$, the edge map $\Phi_i(t):G_i\to G_{i+1}$ sends the elliptic base by $\phi_i$ and sends a nonzero normal vector to its $\theta_i$-dual vector with product $t$. Let $L_i$ be the complex linear part of a logarithmic lift, and let $M_i:\Lambda_i\to\Lambda_{i+1}$ be its restriction. Set
\begin{equation}
 M=M_N\cdots M_1,\qquad L=L_N\cdots L_1.
\hypertarget{tag-00LV}{}\label{cdp:eq:C.HC.6}
\end{equation}
The integral product $M$ is defined up to conjugacy. Its identity condition is independent of the logarithms and of multiplying any $\theta_i$ by a nonzero constant, and it is equivalent to $L=I$.
\end{definition}
\begin{proof}
For a diagonal Hopf surface with contraction $(z,w)\mapsto(\alpha z,\beta w)$, the complement of the incoming axis is the entire normal line bundle of the outgoing axis:
\begin{equation}
 (\mathbb C^*_{w}\times\mathbb C_z)/
       ((w,z)\sim(\beta w,\alpha z)).
\hypertarget{tag-00LT}{}\label{cdp:eq:C.HC.4}
\end{equation}
The specified free finite diagonal quotient gives the secondary case. Its action on the nonzero axis is free: an element with trivial character there would fix its nonzero points in the Hopf cover. Hence the quotient remains a line bundle on the quotient elliptic curve. Coordinatewise exponentiation identifies $\Lambda_i$ with the inverse image of the full diagonal deck group. It contains the two angular periods $2\pi i\Z^2$ and one contracting logarithm, with a finite angular enlargement. These three real periods are independent, since the contracting logarithm has nonzero real part; they form a lattice of rank three. The angular periods span $\C^2$ over $\C$.

For a ruled surface, the chart $v\mapsto[1:v]$ identifies the complement of the infinity section with $\operatorname{Tot}(\mathcal L_i)$; removing the zero section gives $\mathcal L_i^\times$. The other chart uses the inverse fibre coordinate and the dual line. A degree-zero line on an elliptic curve has constant factors of automorphy \cite[arXiv v1, \S2.2.7--\S2.2.9, p.~5]{AbashevaDeev2024AB}. With an oriented base-period basis $\gamma_1,\gamma_2$ and a chosen factor character $\chi_i$, its logarithmic presentation is
\begin{equation}
 \mathcal L_i^\times=\C^2/\Lambda_i,\qquad
 \Lambda_i=\left\langle
       (\gamma_1,\log\chi_i(\gamma_1)),
       (\gamma_2,\log\chi_i(\gamma_2)),(0,2\pi i)
                         \right\rangle_{\Z} .             \hypertarget{tag-00RU}{}\label{cdp:eq:C.MX.3}
\end{equation}
The two projected base periods are real-linearly independent, and the vertical angular period is nonzero. This proves the rank, discreteness and complex spanning assertions. The two positive normal meridians are opposite primitive lattice elements. At a Hopf vertex each angular meridian is primitive as well: a proper fractional angular period would give a nonidentity finite deck element fixing the other nonzero axis.

Each edge is a biholomorphism of the two complements. On universal covers its elliptic base map is affine. The logarithm $h$ of its nowhere-zero normal multiplier has constant increments under both base periods, because the normal lines have constant factors. Thus $h'$ descends to a holomorphic function on the compact elliptic base and is constant. The whole logarithmic lift is affine. In the ordered Hopf coordinates, where $C_i^-$ has coordinate $z_i$ and $C_i^+$ has coordinate $w_i$, it is
\begin{equation}
 \binom{\log z_{i+1}}{\log w_{i+1}}
 =
 L_i\binom{\log z_i}{\log w_i}
 +\binom{d_i}{e_i+\log t},
 \qquad
 L_i=\begin{pmatrix}0&c_i\\-1&k_i\end{pmatrix}.
\hypertarget{tag-00LU}{}\label{cdp:eq:C.HC.5}
\end{equation}
The ruled charts use the same affine construction in their base and line coordinates. In each case $L_i\Lambda_i=\Lambda_{i+1}$, since the lifted map descends to the specified biholomorphism. A change of logarithm adds a lattice translation; a scalar change in $\theta_i$ adds a constant in its normal direction. Neither changes $L_i$. Changing the initial integral basis conjugates $M$. Finally $\Lambda_1$ spans $\C^2$, so $M=I$ holds exactly when $L=I$.
\end{proof}

\subsection{Proper normal products and moving lattices}\label{subsec:II:moving-marked-cycles}

\begin{lemma}[Moving marked cycles and their normal products]\hypertarget{tag-P333}{}\label{lem:II:moving-marked-cycle}
\StatementRecord{lem:II:moving-marked-cycle}{proved}{P333}{\cref{def:II:marked-normal-cycle,lem:II:normal-product-collars}}
For the data of \cref{def:II:marked-normal-cycle}, the prescribed normal-product maps give a smooth proper threefold $\mathcal Y\to\Delta$ whose full central fibre is the marked reduced normal crossing surface $S$. A nearby fibre $Y_t$ has
\begin{equation}
                       \pi_1(Y_t)=\Lambda_1\rtimes_M\mathbb Z.
\hypertarget{tag-00LW}{}\label{cdp:eq:C.HC.7}
\hypertarget{tag-00RV}{}\SsixEquationAliases{\label{cdp:eq:C.MX.4}\label{cdp:eq:C.AD.1}}
\hypertarget{tag-009B}{}
\end{equation}
The same presentation holds for every smooth-total smoothing with this full central fibre and its given branch and normal-product maps, including a single normalization with its two marked curves identified.

When $M=I$, identify the complements with $\C^2/\Lambda_1$ by their complex linear edge maps. If their full affine composite is
\begin{equation}
                    \Phi_N(t)\cdots\Phi_1(t)(v)
                          =v+b_0+v_0\log t.
\hypertarget{tag-00LY}{}\label{cdp:eq:C.HC.9}
\end{equation}
then the integration of the common logarithmic coframe gives
\begin{equation}
                  Y_t\simeq\mathbb C^2/\widehat\Lambda(t)
\hypertarget{tag-00LX}{}\label{cdp:eq:C.HC.8}
\end{equation}
with the exact period lattice
\begin{equation}
                 \widehat\Lambda(t)=
                    \Lambda_1+\mathbb Z(b_0+v_0\log t).
\hypertarget{tag-00LP}{}\label{cdp:eq:C.HC.10}
\hypertarget{tag-00RT}{}\SsixEquationAliases{\label{cdp:eq:C.MX.2}}
\end{equation}

This construction also permits holomorphic motion of the marked components. Choose a basis $\lambda_1,\lambda_2,\lambda_3$ of the common lattice $\Lambda$, and holomorphic vectors $\lambda_j(t)$ with $\lambda_j(0)=\lambda_j$. Transport each lattice vector by its fixed integer coordinates, and write $n_i(t)$ for the transported outgoing normal meridians. Choose holomorphic common affine edge constants $\mathbf b_i(t)\in\C^2$ with the prescribed central values. After shrinking $\Delta$, these data give smooth proper families of the marked Hopf or ruled components and their normal-product filling. Every central component, finite diagonal character, elliptic gluing map and normal-product constant is the prescribed one. Its complete punctured family is the quotient by
\begin{equation}\label{eq:II:moving-cycle-periods}
 \Lambda(t)+\Z\left(b(t)+\frac{\log t}{2\pi i}u(t)\right),
 \qquad b(t)=\sum_i\mathbf b_i(t),\qquad u(t)=-\sum_i n_i(t).
\end{equation}
Here $\Lambda(t)=\sum_j\Z\lambda_j(t)$, and all integral relations are transported with the same integer coefficients. The assertion about moving data assumes $M=I$ at the central value.

For fixed positive integers $h_i$, the same constructed marked
component neighbourhoods admit the compressed equations
$xy=t^{h_i}$ at their edges. They give a proper normal threefold
over $\Delta$, with the same reduced central surface and marked
maps. A neck has the transverse surface singularity $A_{h_i-1}$;
its resolution inserts $h_i-1$ degree-zero ruled elliptic
components. For fixed data its nearby fundamental group is
\eqref{cdp:eq:C.HC.7}. When $M=I$, the complete fixed or moving
punctured family is the coframe quotient with periods
\begin{equation}\label{eq:II:compressed-moving-cycle-periods}
 \Lambda(t)+\Z\left(b(t)+\frac{\log t}{2\pi i}u_h(t)\right),
 \qquad b(t)=\sum_i\mathbf b_i(t),\qquad
 u_h(t)=-\sum_i h_i n_i(t).
\end{equation}
The moving hypotheses and integer transport are the ones above.
The case $h_i=1$ for every edge is the smooth construction already
stated.
\end{lemma}
\begin{proof}
Each marked curve has an entire normal-line chart by \cref{def:II:marked-normal-cycle}. Give its degree-zero normal line a flat Hermitian metric. For an elliptic presentation with periods $\gamma_1,\gamma_2$ and factors $\chi(\gamma_j)$, the real-linear function $h$ determined by $h(\gamma_j)=\log|\chi(\gamma_j)|$ makes $|v|e^{-h(z)}$ invariant. This gives the flat metric explicitly. Choose the metric on its dual at the next component through $\theta_i$, so the product of the two normal norms is $|t|$ on $xy=t$. The two marked curves in every component are disjoint, so their small normal discs can be chosen disjoint.

Fix $0<r_0<R_0$, with $R_0$ small enough for these disjoint charts, and take $|t|<r_0^2$. At each edge use the open subset $\|x\|,\|y\|<R_0$ of the two dual lines with the holomorphic equation $xy=t$. Attach its $x$-end where $\|x\|>r_0$ and its $y$-end where $\|y\|>r_0$ to the respective component charts; the inverse inserts $t/x$ or $t/y$. Both ends cannot overlap because $\|x\|\|y\|<r_0^2$. The two ends attached to the same component remain disjoint even for a single self-identified normalization.

These maps glue the component cores times the disc and the finitely many necks by biholomorphisms. Each equivalence class meets at most one core piece and one neck piece. At either end of an attaching annulus, one of those open pieces omits its limiting point, so the gluing graph is closed relative to the product of the pieces. The finite quotient is therefore Hausdorff.

Choose $r_0<r_1<R_1<R_0$ and a closed base subdisc $|t|\le\delta<r_0^2$. Take in each component the closed core on which both marked normal norms are at least $r_1$, and take the closed neck sets on which both norms are at most $R_1$. These are compact over the compact components or elliptic bases. A point of an open core outside its closed core has one normal norm between $r_0$ and $r_1$; its dual neck coordinate has norm at most $\delta/r_0<r_0$, so the point lies in a closed neck set. A neck point with one norm greater than $R_1$ lies in its corresponding component core; the other marked normal norm there is at least $R_0$ because the two radius-$R_0$ charts are disjoint. It therefore lies in a closed core. These compact sets cover the full inverse image of the closed subdisc and prove properness. The core products are smooth and the neck equation $xy=t$ is smooth, with central branches of multiplicity one. Thus the central fibre is exactly the specified marked surface.

On a ruled complement, the logarithm of the flat norm is a linear height on its logarithmic cover, with compact three-torus levels. For a primary Hopf complement one may take
$h=\log|z|-(\log|\alpha|/\log|\beta|)\log|w|$.
The other normal height is $-(\log|\beta|/\log|\alpha|)h$, so small discs around the two axes bound a compact height interval. These heights are invariant under the finite diagonal group as well. Consequently each core and each annular neck retracts onto a three-torus. They are connected, every attaching annulus is connected, and their incidence cycle is connected, so the nearby fibre is connected. Their collar maps induce exactly the lattice isomorphisms $M_i$. Van Kampen along a spanning tree of the cycle, followed by its last edge, gives \eqref{cdp:eq:C.HC.7}, with the lattice injected. Choose the cycle generator so that conjugation transports the core loops through the oriented edge maps and acts by $M$.

For a given smoothing, apply \cref{lem:II:normal-product-collars} to its full reduced central fibre. Its normalization has the given smooth components, and the two marked curves over each double curve map isomorphically to that curve, so the branch-labelled clause applies also in the self-identification case. The lemma identifies the smooth collars with these same normal circle bundles and product maps, and transports the compact normalization cores compatibly with their boundaries. It therefore supplies exactly the same van Kampen homomorphisms and the same presentation.

Suppose $M=I$. Since the lattices span their complex vector spaces, the complex cycle product is the identity. With $P_1=I$ and $\Lambda=\Lambda_1$, choose constant matrices identifying the complements by
\begin{equation}
                  G=\C^2/\Lambda,\qquad \Lambda\simeq\Z^3,
                  \qquad P_{i+1}L_i=P_i.
\hypertarget{tag-00M8}{}\label{cdp:eq:C.HG.2}
\end{equation}
The common logarithmic differentials $P_i\,d\log(z_i,w_i)$ on Hopf charts, and the corresponding base and logarithmic line differentials on ruled charts, glue across every neck to a closed holomorphic coframe on each nearby compact fibre. The dual fields commute because the forms are closed. They have complete complex-time flows: compactness gives a uniform local existence interval, which can be iterated. Their action of $\C^2$ has open orbits, hence one orbit on the connected fibre, and a discrete cocompact stabilizer. The orbit map proves \eqref{cdp:eq:C.HC.8}.

Integration on the core loops gives exactly $\Lambda_1$, while the oriented cycle loop gives the translation in \eqref{cdp:eq:C.HC.9}. These are all the fundamental-group generators in \eqref{cdp:eq:C.HC.7}; thus the stabilizer is exactly \eqref{cdp:eq:C.HC.10}. A turn of $\log t$ changes its last period by $2\pi i v_0\in\Lambda_1$, because each edge map is single-valued. This proves the full fixed-data period comparison.

For the moving construction, real independence of $\lambda_j(t)$ and complex spanning persist near zero. The incoming normal meridian at vertex $i$ is $-n_{i-1}(t)$. At a Hopf vertex, $n_i(0)$ and $n_{i-1}(0)$ are complex-linearly independent. At a ruled vertex they are equal; extend the primitive $n_i$ to an integral basis $n_i,\lambda_i,c_i$ of the common lattice. The images of $\lambda_i,c_i$ in the elliptic quotient are real-linearly independent, so $\lambda_i$ and $n_i$ are a complex basis. Transport all three by their fixed integer coefficients. The inverse of the complex frame $(\lambda_i(t)/(2\pi i),n_i(t)/(2\pi i))$ then gives the exact ruled presentation
\begin{equation}
        \left\langle
          (0,2\pi i),(2\pi i,0),
                      (\gamma_i(t),b_i(t))
                                  \right\rangle_{\Z} .   \hypertarget{tag-00RS}{}\label{cdp:eq:C.MX.10}
\end{equation}
The projected periods remain real-linearly independent near zero. They define the moving elliptic base, and the factor $\exp b_i(t)$ in this displayed lattice defines its degree-zero line. Its projective completion is a smooth proper family of the original marked ruled surface.

With these ruled frames and the Hopf frames, the common-coordinate identities are
\begin{equation}
 \begin{gathered}
 P_i(t)=\left(\frac{n_i(t)}{2\pi i},
                    -\frac{n_{i-1}(t)}{2\pi i}\right),\\
 \Lambda_i(t)=P_i(t)^{-1}\Lambda(t),
 \quad L_i(t)=P_{i+1}(t)^{-1}P_i(t).
 \end{gathered}
\hypertarget{tag-00M2}{}\label{cdp:eq:C.HG.12}
\hypertarget{tag-009K}{}\SsixEquationAliases{\label{cdp:eq:C.AD.8}}
\end{equation}
where the formula for the two displayed columns of $P_i(t)$ is used at Hopf vertices. A Hopf lattice contains its two fixed angular periods. Its quotient by that angular sublattice has the same abstract integral relations as at zero. A chosen contracting generator has holomorphic logarithms $P_i(t)^{-1}c_i(t)$ with negative real parts on a smaller disc. Every finite generator still has the relation $m v=A(2\pi i,0)+B(0,2\pi i)$ with the same integers. Thus its two finite characters and their freeness conditions are unchanged. Exponentiation gives the marked primary or secondary Hopf family. Uniform contraction gives a compact annulus meeting every orbit over a closed smaller base disc, proving properness of this family.

For every edge, $L_i(t)$ maps its entire lattice to the next one and its outgoing normal to minus the next incoming normal. Hence it induces the elliptic base isomorphism and the reciprocal normal-line map. Adding the chosen common vector $\mathbf b_i(t)$ gives their affine constants, holomorphic in $t$, with the prescribed central maps. The ruled and Hopf cases both extend over their marked curves through their entire line-bundle charts. The component metrics just constructed vary smoothly with $t$: the real-linear equations defining $h_t$ have an invertible base-period matrix. Disjoint cores, flat norms and radii can therefore be chosen uniformly on a smaller disc. The same finite gluing and compact-core argument gives the moving proper filling, still with $xy=t$ and multiplicity one on every central branch.

In common coordinates the logarithmic term on edge $i$ is $-n_i(t)\log t/(2\pi i)$. Summing all edge constants and terms gives \eqref{eq:II:moving-cycle-periods}. On successive cores use their logarithmic maps to the common vector space and subtract the accumulated affine edge constants. Use those same maps on the necks. They agree on overlaps modulo $\Lambda(t)$; the last overlap adds exactly the displayed cycle period. A change of logarithm adds the integral vector $u(t)$. These maps define a holomorphic map of the entire punctured family to the displayed period quotient. On each fibre it is the coframe isomorphism already proved. Its vertical differential is invertible and it is the identity on the base; hence it is locally biholomorphic, and its fibrewise inverses glue holomorphically. This proves the complete moving-family comparison.


For the compressed clause use exactly the same component families,
normal-line maps and metrics. Replace an edge chart by $xy=t^{h_i}$;
on its $x$-end the inverse inserts $y=t^{h_i}/x$, and on its
$y$-end it inserts $x=t^{h_i}/y$. Choose the same four radii
$r_0<r_1<R_1<R_0$ and a closed base radius $\delta$ with
$\delta^{h_i}<r_0^2$ for every edge. Then both attaching ends
cannot meet. An open-core point with normal norm between $r_0$
and $r_1$ has its other neck norm at most
$\delta^{h_i}/r_0<r_0$; a neck point beyond $R_1$ lies in
the same closed component core as before. Thus the identical
closed-graph and finite compact-cover maps prove Hausdorffness
and properness, including for the moving components.

The transverse hypersurface $xy=t^{h_i}$ is normal: it is a
two-dimensional hypersurface smooth away from the origin, hence
regular in codimension one and Cohen--Macaulay. Its resolution
has the following $h_i$ charts $(a_j,b_j)$:
\begin{equation}\label{eq:II:compressed-neck-resolution}
 t=a_jb_j,\qquad x=t^j a_j,\qquad
 y=t^{h_i-j-1}b_j,\qquad 0\le j<h_i.
\end{equation}
The adjacent transition is $a_{j+1}=b_j^{-1}$,
$b_{j+1}=t b_j=a_jb_j^2$. To prove properness, take the closure
of the graph of
$([x:t],[x:t^2],\ldots,[x:t^{h_i-1}])$ in the product of the
hypersurface with $(\PP^1)^{h_i-1}$. In chart $j$ its $k$th
coordinate is $[t^{j-k}a_j:1]$ for $k\le j$ and
$[1:t^{k-j-1}b_j]$ for $k>j$. These formulas extend the
graph map and cover its closure; conversely its consecutive
finite ratio coordinates recover $a_j=x/t^j$ and
$b_j=t^{j+1}/x$ with their displayed regular extensions.
Thus the charts describe that proper graph closure. For
$h_i=1$ the product of projective lines is empty and the
single chart is the identity $x=a_0$, $y=b_0$, $t=a_0b_0$.
For $h_i>1$, over the
origin their zero axes glue by the stated reciprocal transition
to a chain of $h_i-1$ projective lines, and there are no other
exceptional curves. The constant factors of the two normal lines
glue these monomial charts over the elliptic curve. Their inserted
components are projective completions of degree-zero lines,
with two marked sections.
Every resolved central branch has multiplicity one because
$t=a_jb_j$. The resolution is an isomorphism for $t\ne0$.

Each inserted ruled complement has opposite incoming and outgoing
normal meridians. Composing its edge maps gives the same compressed
linear map, while the $h_i$ logarithmic terms sum to
$-h_i n_i(t)\log t/(2\pi i)$. Apply the smooth construction above
to this resolved marked cycle, with those very component and edge
maps. It gives the same injected lattice presentation and coframe
isomorphism on the unchanged nearby fibre. Composing the resolution
isomorphism for $t\ne0$ with that coframe map yields exactly
\eqref{eq:II:compressed-moving-cycle-periods}. The accumulated-edge
coordinate maps and their holomorphic inverses are the previous
ones with each term multiplied by $h_i$; a logarithm change adds
the integral vector $u_h(t)$. This proves the entire compressed
family comparison on its constructed marked neighbourhoods.
\end{proof}



\subsection{Fundamental groups and dual complexes}\label{subsec:II:fibre-fundamental-groups}

\begin{proposition}[Fundamental groups and dual complexes]\hypertarget{tag-01DO}{}\label{cdp:st:C.2.UN.44}
\StatementRecord{cdp:st:C.2.UN.44}{proved}{01DO}{\SeriesRef{cdp:st:C.2.UN.39}, \SeriesRef{cdp:eq:C.FG.9}, \SeriesRef{cdp:eq:C.FG.10}, \SeriesRef{cdp:eq:C.FG.14}, \SeriesRef{cdp:eq:C.FG.15}, \SeriesRef{cdp:eq:C.FG.12}, \SeriesRef{cdp:eq:C.FG.17}, \SeriesRef{cdp:eq:C.FG.18}; \cite[Theorem (LI), p.~775]{Kurdyka1998}; \cite[\S3, Theorem~2]{Lojasiewicz1964}}

Let $f:U\to\Delta$ be a proper holomorphic degeneration of connected
complex two-tori, with smooth total threefold. Write
\begin{equation}
 f^{-1}(0)=\sum_{i=1}^r m_iS_i,\qquad
 S=\bigcup_iS_i,\qquad d=\gcd(m_1,\ldots,m_r).
                                                        \hypertarget{tag-00K8}{}\label{cdp:eq:C.FG.1}
\end{equation}
All fundamental groups below are based; changing the paths to the
chosen base point changes the indicated presentations by their change-of-base-point isomorphisms. Choose a lift $t$ of the positive
generator of $\pi_1(\Delta^*)$, put $\Lambda=\pi_1(F)=\Z^4$, and
write $T$ for conjugation by $t$ on $\Lambda$. There are vectors
$v_i\in\Lambda$, obtained from positively oriented meridians of the
components, for which
\begin{equation}
 \Gamma:=\pi_1(S)
  \simeq
  \bigl(\Lambda\rtimes_T\langle t\rangle\bigr)
       \big/\langle\!\langle x^{v_i}t^{m_i}:1\le i\le r
                    \rangle\!\rangle .                \hypertarget{tag-00KJ}{}\label{cdp:eq:C.FG.2}
\end{equation}
Here $x^v$ denotes the element $v$ of the additive lattice written
multiplicatively. Its image $A$ in $\Gamma$ is abelian, and there is
an exact sequence
\begin{equation}
       1\longrightarrow A\longrightarrow\Gamma
          \longrightarrow\Z/d\Z\longrightarrow1 .    \hypertarget{tag-00KK}{}\label{cdp:eq:C.FG.3}
\end{equation}
In particular the index is exactly $d$ and $A$ is generated by at most
four elements. There is an exact integral description as follows.
Set $n_i=m_i/d$, choose integers $c_i$ with $\sum_i c_in_i=1$, and
put $v_0=\sum_i c_iv_i$. Then
\begin{equation}
 A\simeq\Lambda/R_d,\qquad
 R_d=(T^d-I)\Lambda+
       \left\langle T^jv_i-n_iv_0:
                         1\le i\le r,\ 0\le j<d\right\rangle_{\Z}.
                                                        \hypertarget{tag-00KL}{}\label{cdp:eq:C.FG.4}
\end{equation}
For $d=1$, the whole group is abelian and this specializes to
\begin{equation}
 \pi_1(S)\simeq
   \Lambda\big/\left((T-I)\Lambda+
                   \langle v_i-m_iv_0\rangle_{\Z}\right).
                                                        \hypertarget{tag-00KM}{}\label{cdp:eq:C.FG.5}
\end{equation}
The normalized base change $s^d=f$ is a connected degree-$d$
\emph{\'etale} cover of the total space near $S$. Its reduced central
support $S_d$ is connected, its restriction $S_d\to S$ is the
connected cyclic covering belonging to the subgroup $A$, and
$\pi_1(S_d)=A$ under that covering injection.

For every embedded resolution of the central divisor by smooth
vertical blowups, let $\mathcal D$ be the dual complex of its simple
normal-crossing support. There is a surjection
\begin{equation}
                  \pi_1(S)\longrightarrow\pi_1(\mathcal D).
                                                        \hypertarget{tag-00KN}{}\label{cdp:eq:C.FG.6}
\end{equation}
Its target has an abelian normal subgroup of index dividing $d$,
generated by at most four elements and of rank at most two.
Moreover the cup-product map
$\bigwedge^*H^1(\mathcal D,\Q)\to H^*(\mathcal D,\Q)$ is injective.
When the resolved support has
no triple points, its connected incidence graph is a tree or has
exactly one independent cycle. In particular two components of a simple normal-crossing fibre with no triple points can meet
in at most two connected curves.

For a simple normal-crossing support with no triple points,
write $e_1$ for the number of its connected double curves. The zero
Euler criterion of \ref{cdp:st:C.2.UN.39} takes the following joint form:
\begin{equation}
 \begin{gathered}
 \#\{
      \text{independent incidence cycles}\}\le1,\\
 e(S)=0\ \Longleftrightarrow\
       \sum_{D\ {\rm double}}g(D)=e_1+\frac{j(T)}2 .
 \end{gathered}
                                                        \hypertarget{tag-00KO}{}\label{cdp:eq:C.FG.7}
\end{equation}
Here $j(T)$ is unchanged by replacing the action on $H^1(F)$ by
the corresponding action on $H_1(F)$ or its inverse. If all $m_i$
are equal, $j(T)=0$. If $1$ is not an eigenvalue of $T$, a
no-triple incidence graph must be a tree.
\end{proposition}

\begin{proof}

\begin{enumerate}[label=\arabic*.,leftmargin=*]
\item Apply \cref{lem:II:compact-analytic-tube} to $f$ and
choose a sufficiently small compact closed tube
$U_\epsilon=f^{-1}(|z|\le\epsilon)$. Its inclusion $i:S\hookrightarrow U_\epsilon$ is a
homotopy equivalence. Let $H$ be the strong deformation retraction
provided there. The quantitative gradient estimates for this tube
are as follows. Put $\varphi=|f|^2$ and choose a
smooth Riemannian metric near this tube. The local analytic gradient
inequality, in the form of
\cite[Theorem (LI), p.~775]{Kurdyka1998}, gives near each point of
$S$ an exponent smaller than one. Compactness, a finite coordinate
cover, and uniform equivalence of the coordinate Euclidean metrics
with the chosen metric give constants $c>0$ and $0<\theta<1$ with
\begin{equation}
                 \|\nabla\varphi\|\ge c\varphi^\theta
                   \quad\hbox{on }U_\epsilon\setminus S.
                                                        \hypertarget{tag-00KP}{}\label{cdp:eq:C.FG.8}
\end{equation}
Indeed take the maximum of the finitely many exponents, then shrink
to $0\le\varphi<1$; properness permits a whole smaller tube to be
contained in the resulting neighbourhood of $S$. There is no
critical point of $\varphi$ off $S$, since $df$ is surjective there.

\item  On that complement use the vector field
$V=-\nabla\varphi/\|\nabla\varphi\|^2$. Along its flow the value
of $\varphi$ decreases at speed one. A trajectory starting at level
$a$ therefore exists up to every time smaller than $a$: each
positive closed level band is compact and the vector field is
smooth there. Its remaining length below level $b$ is at most
\begin{equation}
          \int_0^b\frac{du}{c u^\theta}
                 =\frac{b^{1-\theta}}{c(1-\theta)}.    \hypertarget{tag-00KQ}{}\label{cdp:eq:C.FG.9}
\end{equation}
The bound tends uniformly to zero with $b$. The group maps
below use the retraction $H$ and the inverse of the inclusion
$i_*:\pi_1(S)\to\pi_1(U_\epsilon)$ supplied by
\cref{lem:II:compact-analytic-tube}.

\item Write $U^*=U_\epsilon\setminus S$. Proper horizontal transport
over the punctured disc gives a smooth torus bundle. Its homotopy
sequence, and $\pi_2(\Delta^*)=0$, identify
\begin{equation}
       \pi_1(U^*)=\Lambda\rtimes_T\langle t\rangle,
       \qquad \pi_1(U^*)\longrightarrow\pi_1(\Delta^*)=\Z
       \text{ sends }x^v t^a\text{ to }a .             \hypertarget{tag-00K9}{}\label{cdp:eq:C.FG.10}
\end{equation}
The splitting exists because an extension onto the infinite cyclic
group splits after choosing a lift of its generator.

\item  General position makes every loop in the smooth six-manifold
$U_\epsilon$ avoid the divisor $S$. A null-homotopy of a loop from
$U^*$ can be made transverse to the smooth part of $S$ and disjoint
from its singular set: the latter has real dimension at most two,
so a transverse real two-disc misses it. Removing little discs
around the finitely many transverse intersections expresses the
original loop as a product of conjugates of positive or negative
meridians. Conversely each such meridian bounds its normal disc
in $U_\epsilon$.

\item  One meridian for each global irreducible component suffices. The
smooth locus of an irreducible component, with its intersections
with other components removed, is path connected. This follows
also by taking its connected normalization resolution and moving
paths off the proper complex analytic curves and points. A path
there moves positive normal circles through positive normal circles,
giving conjugate based meridians in $U^*$. Near a chosen smooth
point of $S_i$ the equation is $f=u z^{m_i}$ with $u$ a
unit. The meridian has winding $m_i$, hence is exactly
$x^{v_i}t^{m_i}$ in \eqref{cdp:eq:C.FG.10}, for a uniquely determined $v_i$
after the indicated choices. The retraction and this kernel
calculation prove \eqref{cdp:eq:C.FG.2}.

\item  The image $A$ of $\Lambda$ is normal and abelian. Quotienting
\eqref{cdp:eq:C.FG.2} by this image leaves precisely the presentation
$\langle t\mid t^{m_i}=1\text{ for all }i\rangle$.
This is $\Z/d\Z$, which proves the exact index in \eqref{cdp:eq:C.FG.3}.
For later comparison, abelianizing the same presentation
also gives
\begin{equation}
 H_1(S,\Z)\simeq
 \frac{\bigl(\Lambda/(T-I)\Lambda\bigr)\oplus\Z}
      {\langle ([v_i],m_i):1\le i\le r\rangle_{\Z}}.
                                                        \hypertarget{tag-00KA}{}\label{cdp:eq:C.FG.11}
\end{equation}

\item Work momentarily in the open complex neighbourhood of the closed
tube. Since all divisor multiplicities of $f$ are divisible by
$d$, unique factorization in the regular local analytic ring
gives $f=u h^d$ near every point, including singular points of the
support. The normalization of $s^d=u h^d$ is the explicit map
\begin{equation}
       \{v^d=u\}\longrightarrow\{s^d=u h^d\},
                         \qquad s=vh .                \hypertarget{tag-00KB}{}\label{cdp:eq:C.FG.12}
\end{equation}
Its source is finite and \'etale over the smooth original chart:
$v$ is a unit and $d v^{d-1}$ never vanishes. It is normal and is
isomorphic to the base change over $h\ne0$, hence is exactly that
base change's normalization. These normalization maps glue.

\item  Over $U^*$ this is the cyclic covering given by the winding number
modulo $d$ in \eqref{cdp:eq:C.FG.10}. That homomorphism is onto, so the cover
is connected. Its complement in the normalization is a divisor;
the full normalized cover is therefore connected as well.
The tube retraction $H$ lifts uniquely to this covering.
Its lift starts at the identity, fixes the inverse image of $S$
pointwise, and at time one has image in that inverse image. It is
therefore a strong deformation retraction onto $S_d$, proving
connectedness of $S_d$ and the asserted covering injection.
The covering subgroup in $\Gamma$ is the kernel of the winding
modulo $d$ map, which is exactly $A$ by \eqref{cdp:eq:C.FG.3}. The inverse image of an individual $S_i$ is the finite cover
induced by restricting this connected covering of the whole tube.

\item  There is also a useful integral calculation using the punctured
part of this same cover. Its group is
$\Lambda\rtimes_{T^d}\langle w\rangle$, where $w=t^d$.
The kernel of its map to $\pi_1(S_d)$ is normally generated by
all the lifted meridians
\begin{equation}
                x^{T^jv_i}w^{n_i},
                  \qquad 1\le i\le r,\quad0\le j<d .  \hypertarget{tag-00KC}{}\label{cdp:eq:C.FG.13}
\end{equation}
To see completeness of this list, use the coset representatives
$1,t,\ldots,t^{d-1}$ for the winding cover in \eqref{cdp:eq:C.FG.10}.
Conjugating the original normal generators by these representatives,
and then by elements of the covering subgroup, gives exactly its
kernel. These cosets enumerate all lifts over every component $S_i$.

\item  The relations with $j=0$ imply $w^{n_i}=x^{-v_i}$.
Multiplying their $c_i$th powers gives $w=x^{-v_0}$. Consequently
the covering group is abelian, $(T^d-I)\Lambda$ is killed, and
the remaining relations are precisely
$T^jv_i-n_iv_0=0$. Conversely the map $x^v\mapsto[v]$,
$w\mapsto-[v_0]$ from this presentation to $\Lambda/R_d$
respects every semidirect and meridian relation. It is inverse to
the map induced by $\Lambda$. This proves \eqref{cdp:eq:C.FG.4}. For $d=1$
the same argument proves \eqref{cdp:eq:C.FG.5}; in particular it first proves
$t\in A$ and only then imposes $T=I$ on that quotient.

\item Apply \cref{lem:II:blowup-topology} to each smooth point or
curve centre of $\pi:U^{\rm res}_\epsilon\to U_\epsilon$.
The blowdown induces an isomorphism of total-space fundamental
groups. The inclusions $S^{\rm res}\hookrightarrow U^{\rm res}_\epsilon$
and $S\hookrightarrow U_\epsilon$ are homotopy equivalences by
\cref{lem:II:compact-analytic-tube}. They form a commutative square
with the blowdown and its restriction to the full inverse image.
Consequently the latter restriction gives
\begin{equation}
              \pi_1(S^{\rm res})\xrightarrow{\ \simeq\ }\pi_1(S).
                                                        \hypertarget{tag-00KD}{}\label{cdp:eq:C.FG.14}
\end{equation}
The old multiplicities are retained on strict transforms. Every
new multiplicity is the order of the old fibre equation along its
centre and is divisible by $d$, so their gcd remains $d$.

\item  For an SNC support $E=\bigcup E_i$, form the space $P$ from the
disjoint components $E_i$, one $D\times[0,1]$ for each connected
double curve, and one two-simplex for each triple point. The faces
are glued by the intersection inclusions. There are continuous maps
\begin{equation}
                   E\ \xleftarrow{\ p\ }\ P
                       \xrightarrow{\ q\ }\mathcal D,             
                                                        \hypertarget{tag-00KE}{}\label{cdp:eq:C.FG.15}
\end{equation}
where $p$ forgets the simplex coordinate and $q$ collapses every
connected stratum to its associated cell. Compatible triangulation
of the finite analytic family of components and intersections makes
all the inclusions closed subcomplex inclusions; see
\cite[\S3, Theorem~2, p.~464]{Lojasiewicz1964}. Apply
\cref{lem:II:proper-conductor-pushout} successively to
$(E_1\cup\cdots\cup E_{k-1})\sqcup E_k\to E_1\cup\cdots\cup E_k$,
with closed image locus $C=(E_1\cup\cdots\cup E_{k-1})\cap E_k$
and inverse image the two copies of $C$. Its double-mapping-cylinder
comparison applies to these closed subcomplex inclusions. Applying
the same comparison to the intersections gives exactly the
double-curve cylinders and triple-point triangles in $P$.
Thus $p:P\to E$ is a homotopy equivalence; choose a homotopy
inverse $p^{-1}:E\to P$.

\item  Van Kampen on these mapping cylinders gives a particularly precise
description of $q_*$. Start with the vertex fundamental groups and
edge paths, impose the relations induced by the two inclusions of
each $\pi_1(D)$, and then the triangle boundary relations.
Killing the vertex groups leaves exactly the presentation of
$\pi_1(\mathcal D)$. Thus
\begin{equation}
 \pi_1(\mathcal D)\simeq
 \pi_1(E)\big/
       \langle\!\langle\operatorname{im}\pi_1(E_i):i
                              \rangle\!\rangle .       \hypertarget{tag-00KF}{}\label{cdp:eq:C.FG.16}
\end{equation}
The map $q\circ p^{-1}$, where $p^{-1}$ is this homotopy inverse,
and \eqref{cdp:eq:C.FG.14} give \eqref{cdp:eq:C.FG.6}. In particular it is a topological quotient map on fundamental groups, including when
two components have several connected double curves.

\item  The image of $A$ is an abelian normal subgroup whose quotient is
a quotient of $\Z/d\Z$. If there are no triple points,
$\mathcal D$ is a finite connected graph. Let its free-group rank
be $b$ and let the index of that abelian subgroup be $a$.
The corresponding connected graph covering has $a$ times as many
vertices and edges, hence free fundamental group of rank
$1+a(b-1)$. An abelian free group has rank at most one. Thus
$b\le1$, proving the tree-or-one-cycle assertion and its two-vertex
consequence.

\item For any of these tubes, inclusion of a nearby fibre followed by
the constructed retraction is a continuous map
$\sigma:F\to S$. Restriction to its fundamental group has image
$A$ of finite index. Therefore
\begin{equation}
 \sigma^*:H^1(S,\Q)\longrightarrow H^1(F,\Q)
           \quad\hbox{is injective}.                  \hypertarget{tag-00KG}{}\label{cdp:eq:C.FG.17}
\end{equation}
Indeed a homomorphism $\Gamma\to\Q$ vanishing on $A$ factors through
the finite group $\Gamma/A$ and is zero. This proves the stated
direction of the cohomology map. The map $\sigma^*$ in all degrees
is a ring homomorphism because it comes from that continuous
map. Its degree-one image is $T$-invariant, as is also immediate
from the semidirect presentation.

\item  Apply this to the resolved tube $W$ with central support
$E=S^{\rm res}$. The map $q\circ p^{-1}:E\to\mathcal D$ of
\eqref{cdp:eq:C.FG.15} induces an injection on $H^1$ because its fundamental
group map is onto. The composite
\begin{equation}
 H^*(\mathcal D,\Q)\longrightarrow H^*(E,\Q)
               \xrightarrow{\ \sigma^*\ }H^*(F,\Q)
                  =\bigwedge^*H^1(F,\Q)               \hypertarget{tag-00KH}{}\label{cdp:eq:C.FG.18}
\end{equation}
is again a ring homomorphism and is injective in degree one.
The exterior algebra on any linear subspace of $H^1(F,\Q)$
injects into this exterior algebra. Consequently cup product
on $\mathcal D$ is injective on
$\bigwedge^*H^1(\mathcal D,\Q)$. Since $\mathcal D$ is a
two-dimensional cell complex, its third cohomology vanishes,
and three independent degree-one classes would give a
contradiction. Thus $b_1(\mathcal D)\le2$.

\item  The following finite-cover comparison bounds the rank of its
finite-index abelian subgroup. Put
$\Pi=\pi_1(\mathcal D)$ and let $B$ be the image of $A$ in $\Pi$.
The index $a=[\Pi:B]$ divides $d$. Take the connected
degree-$a$ covering $\mathcal D'\to\mathcal D$ belonging to $B$,
and pull it back along the continuous map
$W\to E\to\mathcal D$ just constructed. The pullback is connected
because $\pi_1(W)\to\Pi$ is onto. Its covering subgroup is the
inverse image of $B$ in $\Gamma$. The quotient $\Pi/B$ is cyclic
of order $a$, generated by the image of $t$, so this inverse
image is precisely the kernel of the winding number modulo $a$.
The uniqueness of a based covering with a given subgroup
identifies this pullback, over $W$, with the normalized
root cover
\begin{equation}
                  s^a=f\circ\pi:W_a\longrightarrow W. \hypertarget{tag-00KI}{}\label{cdp:eq:C.FG.19}
\end{equation}
It is \'{e}tale by \eqref{cdp:eq:C.FG.12}, since $a$ divides every resolved
component multiplicity. Its lifted tube retraction has connected
central support $E_a$. The lifted map
$E_a\to\mathcal D'$ is onto on fundamental groups: its image is
exactly $B$, by the inverse-image description of the covering
subgroup. Regarded over the $s$-disc, $W_a$ has connected nearby
fibre equal to the torus $F_{s^a}$. Hence \eqref{cdp:eq:C.FG.17}--\eqref{cdp:eq:C.FG.18}
apply to this map and give $b_1(\mathcal D')\le2$. Since
$\pi_1(\mathcal D')=B$ is finitely generated abelian, this is
exactly $\operatorname{rank}B\le2$, as claimed.

\item  If a no-triple incidence graph had one cycle, its nonzero
$H^1$ would inject by \eqref{cdp:eq:C.FG.18} into the $T$-invariant part of
$H^1(F,\Q)$. Thus $1$ must be an eigenvalue of $T$ in that case.

\item Under the evaluation pairing $H^1(F,\Z)=H_1(F,\Z)^\vee$,
pullback by the homology action $T$ has matrix $T^{\mathsf t}$.
Thus $\det(I-(T^{\mathsf t})^k)=\det(I-T^k)$, as in
\cref{night:P4.13}. The dual parallel-transport convention gives
$T^\vee=(T^{-1})^{\mathsf t}$, and
$\det(I-(T^\vee)^k)=\det(I-T^{-k})=
\det(T)^{-k}\det(T^k-I)=\det(I-T^k)$.
Here $\det T=1$ by preservation of the complex orientation and
$\operatorname{rank}\Lambda=4$ supplies the positive last sign.
For the original support already SNC with no triple
points, \ref{cdp:st:C.2.UN.39} has no resolution-centre terms and gives
$e(S)=j(T)-2\sum_D(g(D)-1)$, proving \eqref{cdp:eq:C.FG.7}.
If all multiplicities are $m$, the full trace formula is
$\det(I-T^k)=m\sum_i e(S_i^\circ)$ when $m\mid k$ and zero
otherwise. Taking $k$ also divisible by the semisimple order makes
$T^k$ unipotent, so this determinant is zero. Hence the common
sum is zero and all traces vanish; in particular $j(T)=0$.

\end{enumerate}
\end{proof}

\subsection{Component groups and geometric genus}\label{subsec:II:component-geometric-genus}

\begin{proposition}[Component groups and geometric genus]\hypertarget{tag-01DR}{}\label{cdp:st:C.2.UN.47}
\StatementRecord{cdp:st:C.2.UN.47}{proved}{01DR}{\SeriesRef{cdp:st:C.2.UN.33}, \SeriesRef{cdp:st:C.2.UN.44}, \SeriesRef{cdp:eq:C.PF.5}, \SeriesRef{cdp:st:C.2.UN.28}, \SeriesRef{cdp:eq:C.PF.4}, \SeriesRef{cdp:eq:C.PF.7}, \SeriesRef{cdp:eq:C.PF.8}, \SeriesRef{cdp:eq:C.PF.6}, \SeriesRef{cdp:eq:C.PF.12}, \SeriesRef{cdp:eq:C.PF.10}}

Let $f:U\to\Delta$ be a proper holomorphic two-torus degeneration
with smooth connected total threefold, let $S=f^{-1}(0)_{\rm red}$,
and put $\Gamma=\pi_1(S)$. For an irreducible component $S_i$ let
$h_i:Z_i\to S_i\hookrightarrow U$ be a resolution of its normalization.
The following are simultaneous constraints on this fibre.

Every $Z_i$ with $p_g(Z_i)>0$ has $p_g(Z_i)=1$, and its map induces a surjection
\begin{equation}
                         \pi_1(Z_i)\longrightarrow\Gamma .
                                                        \hypertarget{tag-00UN}{}\label{cdp:eq:C.PF.1}
\end{equation}
Consequently every SNC dual complex of an embedded resolution of
the fibre is simply connected whenever such a component occurs.
A component with K3 minimal model forces $\Gamma=1$.
More generally,
\begin{equation}
                    \chi(\OO_{Z_i})=2
                         \quad\Longrightarrow\quad\Gamma=1 .
                                                        \hypertarget{tag-00UR}{}\label{cdp:eq:C.PF.2}
\end{equation}
This includes the Kodaira-dimension-one components of holomorphic
Euler characteristic two in \ref{cdp:st:C.2.UN.33}.

For a component with $\chi(\OO_{Z_i})=1$, its fundamental-group
image in $\Gamma$ is either trivial or satisfies

\begin{equation}
 \Gamma\simeq\Z/2\Z,\qquad
 \operatorname{im}\bigl(\pi_1(Z_i)\to\Gamma\bigr)=\Gamma .
                                                        \hypertarget{tag-00US}{}\label{cdp:eq:C.PF.3}
\end{equation}
In that configuration the connected double cover of $Z_i$
induced from the tube has
$\chi(\OO)=2$, $p_g=1$ and $q=0$. At most one component with
$\chi(\OO)=1$ can have nontrivial fundamental-group image.
\end{proposition}

\begin{proof}

\begin{enumerate}[label=\arabic*.,leftmargin=*]
\item Use \cref{lem:II:compact-analytic-tube} on a sufficiently
small tube. Its central inclusion and radial annulus comparison
identify the fundamental groups of the open tube, a smaller
closed tube and $S$ with $\Gamma$. Given a finite quotient
$\rho:\Gamma\twoheadrightarrow Q$ of order $N$, its connected
regular topological covering carries the unique complex structure
obtained by lifting the holomorphic coordinate charts. This gives
a finite \'{e}tale map
\begin{equation}
                     p:U_Q\longrightarrow U,\qquad \deg p=N.
                                                        \hypertarget{tag-00UT}{}\label{cdp:eq:C.PF.4}
\end{equation}
It is proper: a finite covering of a locally compact Hausdorff
space is proper, as follows by covering a compact set with
finitely many relatively compact evenly covered neighbourhoods.
Consequently $f\circ p$ is proper. Its total space is smooth and
connected.

\item The punctured covering $U_Q^*$ is connected, since
$\pi_1(U^*)\to\Gamma\to Q$ is onto by \ref{cdp:st:C.2.UN.44}. Let $A$ be the image
of the nearby fibre group in $\Gamma$ and let $B=\rho(A)$.
The connected components over a fixed nearby fibre are indexed
by $Q/B$. The base generator permutes them transitively, and
$Q/B$ is cyclic, because $\Gamma/A$ is cyclic. If
$a=[Q:B]$, the connected base covering in the Stein factor of
$f\circ p$ is therefore the punctured degree-$a$ cover of the disc.
Its normal completion has one point above zero and is the disc
with coordinate map $t=s^a$. One can obtain this
identification by the classification of coverings of the
punctured disc and then extend the two inverse maps across zero:
each is bounded over a smaller disc, and the completed normal
one-dimensional spaces have the removable-singularity property.
Equivalently it is the normalization of this finite disc cover.
The proper Stein factor is this completion, because it is finite,
normal and agrees over the dense punctured disc. We have thus
constructed a proper map with connected fibres
\begin{equation}
                        f_Q:U_Q\longrightarrow\Delta_s .
                                                        \hypertarget{tag-00UU}{}\label{cdp:eq:C.PF.5}
\end{equation}
Its reduced central support is exactly $(p^{-1}S)_{\rm red}$,
because the completed Stein disc has just one point over zero.
Every smooth general fibre of $f_Q$ is a connected finite \'{e}tale
cover of $F$. Writing $F=\C^2/\Lambda$, its lifted universal
cover identifies it with $\C^2/\Lambda'$ for a finite-index
sublattice $\Lambda'\subset\Lambda$. Hence \eqref{cdp:eq:C.PF.5} is another
two-torus degeneration, to which \ref{cdp:st:C.2.UN.28} applies. The degree-$a$ Stein base cover and the fibre lattice cover
$\C^2/\Lambda^\prime\to\C^2/\Lambda$ are the two maps in this
comparison.

Fix $Z=Z_i$ and take its fibre product with \eqref{cdp:eq:C.PF.4}. It is a
smooth compact, possibly disconnected surface
\begin{equation}
               Z_Q=Z\times_U U_Q\longrightarrow Z
                       \quad\hbox{finite \'{e}tale of degree }N .
                                                        \hypertarget{tag-00UV}{}\label{cdp:eq:C.PF.6}
\end{equation}
Let $H_Q=\rho(h_{i*}\pi_1(Z))$, let $e=|H_Q|$, and put
$k=[Q:H_Q]=N/e$. Path lifting gives exactly $k$ connected
components $Z_1',\ldots,Z_k'$, each a connected cover of $Z$
of degree $e$.

\item These are resolutions of distinct normalization components
of the central support of $U_Q$. To verify this comparison, remove
the singular set of the original support and the exceptional
curves of $Z\to S_i$. Over the remaining dense open set the map
$Z\to U$ is its embedding as the smooth single-component locus.
Its base change is therefore an isomorphism onto the corresponding
open subset of the central support of $U_Q$. Removing those proper
analytic curves and points from a connected smooth $Z_j'$ leaves
it connected. Its proper image is consequently an irreducible
central component, and the displayed isomorphism over the dense
open set proves that $Z_j'$ maps birationally onto its normalization.
Different $Z_j'$ cannot give the same component, since their dense
open isomorphisms have disjoint images. Thus the compact-core
bound of \ref{cdp:st:C.2.UN.28} gives the exact useful inequality
\begin{equation}
                         \sum_{j=1}^k p_g(Z_j')\le1 .
                                                        \hypertarget{tag-00UW}{}\label{cdp:eq:C.PF.7}
\end{equation}
Other central components can only add nonnegative terms to that
bound.

\item Pullback of a holomorphic two-form under a finite \'{e}tale
cover is injective, by restriction to any open set on which
that cover is biholomorphic. If $p_g(Z)>0$, \ref{cdp:st:C.2.UN.28} on the
original degeneration already gives $p_g(Z)=1$. Every
$Z_j'$ then contributes at least one in \eqref{cdp:eq:C.PF.7}, so $k=1$.
We have proved that
\begin{equation}
       \rho(h_{i*}\pi_1(Z))=Q
                 \quad\hbox{for every finite quotient }\rho
                         \hbox{ of }\Gamma .           \hypertarget{tag-00UX}{}\label{cdp:eq:C.PF.8}
\end{equation}

\ref{cdp:st:C.2.UN.44} supplies a normal finitely generated abelian subgroup
$A\subset\Gamma$ of finite index. Every subgroup $H$ of $\Gamma$
is separable by finite quotients in the sense needed here.
If $HA\ne\Gamma$, its image already fails to be all of
the finite quotient $\Gamma/A$. Otherwise, for a proper
$H$, choose $b\in A\setminus H$. The finitely generated
abelian group $A/(H\cap A)$ has a finite quotient detecting
the class of $b$; equivalently some positive integer $n$
satisfies
\begin{equation}
                         b\notin(H\cap A)+nA .
                                                        \hypertarget{tag-00UY}{}\label{cdp:eq:C.PF.9}
\end{equation}
The subgroup $nA$ is invariant under every automorphism of $A$,
hence is normal in $\Gamma$, and $\Gamma/nA$ is finite.
Moreover $HN\cap A=(H\cap A)+nA$ for $N=nA$.
Thus the image of $H$ in this finite quotient misses $b$.
This proves directly that \eqref{cdp:eq:C.PF.8} implies
$h_{i*}\pi_1(Z)=\Gamma$, establishing \eqref{cdp:eq:C.PF.1}.

\item The same elementary argument gives residual finiteness of
$\Gamma$. An element outside $A$ is detected in $\Gamma/A$.
A nonzero element of $A$ is detected in $A/nA$ for some $n$,
and therefore in the finite quotient $\Gamma/nA$.

Choose any embedded resolution of the central divisor and let
$V$ be the strict transform of $S_i$. Both $V$ and $Z$ resolve
the same normal surface $S_i^\nu$. At each of its singular points,
\cite[Proposition~3.10, p.~20]{PopescuPampu2016} factors these maps
through its unique minimal resolution. These local factorizations
glue with the identity over the smooth locus.
Resolve the closure of the graph of $V\dashrightarrow Z$ over
$S_i^\nu$, giving a smooth common resolution $W$.
Its projections $W\to V$ and $W\to Z$ are proper holomorphic
modifications of smooth complex surfaces. By
\cite[Proposition~3.12, p.~20]{PopescuPampu2016}, both are
compositions of point blowups. Pullback and Hartogs extension of holomorphic
two-form coefficients across a point show that each point blowup
preserves $p_g$. The same point-blowup maps induce isomorphisms on $\pi_1$
by \cref{lem:II:blowup-topology}, applied to each smooth surface
and its point centre.
The two maps from the common resolution to $U$ agree on a
dense open set and hence everywhere. Thus the image of the
strict-transform vertex group in $\Gamma$ is exactly the
image in \eqref{cdp:eq:C.PF.1}.

\item By \eqref{cdp:eq:C.FG.14}, the resolved support has group $\Gamma$.
By \eqref{cdp:eq:C.FG.16}, passage to its dual complex kills the image of
every component group. In particular it kills this surjective
vertex group, so the dual-complex fundamental group is trivial.
If $Z$ has K3 minimal model, its fundamental group is zero:
the K3 is simply connected by
\cite[author draft, Chapter~7, Corollary~1.4, pp.~127--128]{HuybrechtsK3Draft},
and \cref{lem:II:blowup-topology} gives the fundamental-group
isomorphisms for its point blowups. Equation \eqref{cdp:eq:C.PF.1} then gives $\Gamma=1$.

The degree in \eqref{cdp:eq:C.PF.6} also gives
\begin{equation}
       \sum_{j=1}^k\chi(\OO_{Z_j'})=N\chi(\OO_Z).
                                                        \hypertarget{tag-00UO}{}\label{cdp:eq:C.PF.10}
\end{equation}
This follows, for example, by applying surface Riemann--Roch to
the \'{e}tale tangent-bundle pullback:
$12\chi(\OO)=\int(c_1^2+c_2)$, and integration of a pulled-back
top-degree class multiplies by the covering degree. Combining
this equality with
$p_g(Z_j')=\chi(\OO_{Z_j'})-1+q(Z_j')$ gives
\begin{equation}
       1\ \ge\ \sum_j p_g(Z_j')
           \ \ge\ N\chi(\OO_Z)-k,\qquad k\le N .
                                                        \hypertarget{tag-00UP}{}\label{cdp:eq:C.PF.11}
\end{equation}
For $\chi(\OO_Z)=2$, the right side is at least $N$.
Every finite quotient of $\Gamma$ therefore has order one.
Residual finiteness proves \eqref{cdp:eq:C.PF.2}.

\item For $\chi(\OO_Z)=1$, a finite quotient in which $H_Q$ is
nontrivial has $e\ge2$ and satisfies
\begin{equation}
                    1\ge N-\frac Ne\ge\frac N2 .
                                                        \hypertarget{tag-00UQ}{}\label{cdp:eq:C.PF.12}
\end{equation}
It follows that $N=e=2$, so the only such finite quotient
is a quotient of order two. If the image
$H=h_{i*}\pi_1(Z)$ in $\Gamma$ is nontrivial, residual
finiteness provides one such quotient
$\rho:\Gamma\to\Z/2\Z$ detecting a nonidentity element of $H$.
If $\ker\rho$ contained a nonidentity element, another finite
quotient would detect that element. The image of the product
of these two quotient maps would then be a finite quotient
of order greater than two still detecting the original
element of $H$, contradicting \eqref{cdp:eq:C.PF.12}. Thus $\ker\rho=1$,
which proves precisely \eqref{cdp:eq:C.PF.3}.

\item In this case \eqref{cdp:eq:C.PF.6} is a connected double cover of $Z$,
and \eqref{cdp:eq:C.PF.10} gives $\chi=2$. Equation \eqref{cdp:eq:C.PF.7} and
$p_g=\chi-1+q$ give $p_g=1$ and $q=0$. If two different
components with $\chi=1$ both had nontrivial image, the
same double cover of the tube would contain their
two distinct normalization resolutions, each of geometric
genus one. This contradicts \ref{cdp:st:C.2.UN.28}. It proves the final
joint assertion.

\end{enumerate}
\end{proof}

\section{Normal supports, class VII and contracted boundaries}\hypertarget{tag-0343}{}\label{sec:II:07}


\subsection{Normal irreducible supports}\label{subsec:II:normal-irreducible-supports}

\begin{proposition}[Every normal irreducible support is a smooth torus quotient]\hypertarget{tag-01DU}{}\label{cdp:st:C.2.UN.5}
\StatementRecord{cdp:st:C.2.UN.5}{proved}{01DU}{\cref{cdp:st:A.5.17}; \cite[Theorem~4.6, pp.~13--14; Definition~3.1 and Theorems~3.2--3.3, p.~9]{Catanese2002}}

Let $f:Y\to B$ be a proper holomorphic map from a smooth
complex threefold to a smooth curve, with connected fibres
and smooth general fibre a complex two-torus. Let $mS$ be
a full fibre whose reduced support is irreducible and normal.
Then $S$ is smooth and $e(S)=0$. More precisely, the
normalized base change $t=s^m$ gives a connected two-torus
$A$ and a finite unramified cyclic map
\[
                         A\longrightarrow S
\]
of degree $m$. Thus $S$ is a two-torus or a bielliptic
surface. In the bielliptic case the order $n$ of $K_S$ is
$2,3,4$, or $6$, and $n\mid m$; in particular a reduced
normal fibre is a torus.

Every two-torus occurs with every positive multiplicity.
Every bielliptic surface of canonical order $n$ occurs
with every multiplicity divisible by $n$. These
realizations can be made in smooth compact threefolds
over $\PP^1$.
\end{proposition}

\begin{proof}

\begin{enumerate}[label=\arabic*.,leftmargin=*]
\item A normal surface hypersurface has only isolated
singularities. At a singular point of $S$, write the
map as $t=u h^m$, with $u$ a holomorphic unit
and $h$ a reduced local equation of $S$. A local
$m$th root of $u$ gives $t=(u^{1/m}h)^m$. The
isolated-hypersurface assertion of \cref{cdp:st:A.5.17} supplies an
embedded vanishing two-sphere of self-intersection
$-2$ in a small nonzero Milnor fibre of $u^{1/m}h$.
This fibre is an open part of the nearby
torus fibre of $f$. A sphere in a torus is nullhomotopic,
and hence has zero homology class and zero
self-intersection. The contradiction proves that $S$
is smooth. 

\item Shrink the base to a disc with coordinate $t$. Near
any point of $S$ write $t=u z^m$ with $z$ a smooth
equation of $S$. The normalization $U'$ of the base
change $t=s^m$ is given locally by the explicit maps
\begin{equation}
        w^m=u,\qquad s=wz,
        \qquad p:U'\longrightarrow U.
 \hypertarget{tag-015B}{}\label{cdp:eq:C.UNN.1}
\end{equation}
Since $u$ never vanishes, $p$ is finite and unramified
of degree $m$; these local normalizations glue to the
normalization of the fibre product. The map
$h:U'\to\Delta_s$ has $ds=w\,dz$ along its central
support, so that support is smooth and reduced.
For each $s\ne0$, the map $p$ identifies $h^{-1}(s)$
with $f^{-1}(s^m)$, a connected torus. The Stein
factor of $h$ has generic degree one over the normal
disc and is therefore the disc. In particular $U'$
and its central fibre $A$ are connected. After
shrinking the disc, $h$ is a proper submersion.
The deck action is $w\mapsto\zeta_m w$, so it is
free also on $A$, and its quotient there is exactly
$S$.

\item Apply \cite[Theorem~4.6 and its proof, pp.~13--14]{Catanese2002}
to the proper submersion $h:U'\to\Delta_s$. Its central fibre is
the compact complex manifold $A$, and every fibre with $s\ne0$ is
a two-torus. A local classifying map to the Kuranishi base of $A$
therefore sends a sequence $s_\nu\to0$ to torus points converging to
the central parameter, exactly the hypothesis of that theorem.
It follows that $A$ is a two-torus. The period-integral map in
the source's proof is the Albanese map of this family; on $A$,
after choosing an origin $a_0$, it is the isomorphism
$a\mapsto[\omega\mapsto\int_{a_0}^a\omega]$.

\item Write a generator of the deck group on
$A=\mathbb C^2/\Lambda$ as $g(z)=Lz+b$.
Every holomorphic torus automorphism has this affine
form: the derivative of a lift is periodic and hence
constant on the compact quotient. The matrix $L$
preserves $\Lambda$ and has finite order. If $1$
were not an eigenvalue of $L$, the equation
$(I-L)z=b$ in $\mathbb C^2$ would give a fixed
point of $g$, contrary to the free deck action.
If $L=I$, the deck group acts by translations and
$S$ is the torus quotient by the resulting finite
subgroup.

\item Otherwise the eigenvalues on $\mathbb C^2$ are
$1$ and a nontrivial root of unity $\zeta$.
The corresponding integral real representation
has characteristic polynomial
$(x-1)^2(x-\zeta)(x-\bar\zeta)$, so the order
$n$ of $\zeta$ is $2,3,4$, or $6$. The images of
$I-L$ and of $I+L+\cdots+L^{n-1}$ are elliptic
subtori with complementary tangent lines. Their
addition map is an isogeny from a product of
elliptic curves to $A$. The kernel of the deck
group's linear representation consists of
translations. Quotienting by it first gives a
torus $A_1$ with a free cyclic action of order
$n$ and quotient $S$. On the canonical form
this action has character $\zeta$, so $K_S$
has exact order $n$, while the invariant
holomorphic one-form space has dimension one.
Equivalently, lifting to the product isogeny
and removing the pure translation kernels gives
the bielliptic presentation
$(E\times F)/G$, with $G$ acting faithfully
by translations on $E$ and with $F/G\simeq\PP^1$.
These are precisely the bielliptic surfaces
of \cite[Definition~3.1 and Theorems~3.2--3.3, p.~9]{Catanese2002}.
The linear order divides the deck order $m$.
In both cases the unramified cover gives
$m\,e(S)=e(A)=0$.

\item We finish with realizations and their multiplicities.
Given any free cyclic torus quotient
$A\to S=A/\langle g\rangle$ of degree $m$, set
\begin{equation}
 Y=(A\times\PP^1)/\langle(a,z)\mapsto
                         (g(a),\zeta_m z)\rangle,
 \qquad f([a,z])=z^m.
 \hypertarget{tag-015C}{}\label{cdp:eq:C.UNN.2}
\end{equation}
The action is free, hence $Y$ is smooth and compact.
For $t\ne0,\infty$ its $m$ slices are permuted
transitively, and the quotient fibre is exactly
$A$. The full fibres at $0$ and $\infty$ are
$mS$. Thus the quotient realizes the displayed support and multiplicity.
For a prescribed torus $S$, an index-$m$ sublattice
with cyclic quotient gives the required cyclic
isogeny $A\to S$ for every $m\ge1$.

\item For a prescribed bielliptic surface, use its
product presentation $(E\times F)/G$. The linear
character $G\to\mu_n$ is surjective, and its
kernel acts by translations on both factors.
Consequently its canonical cyclic cover is a
torus. Moreover the character factors through
$\pi_1(E/G)\simeq\mathbb Z^2$, because $G$
acts on $E$ by translations. A surjective
character $\mathbb Z^2\to\mathbb Z/n$ lifts
to a primitive integral homomorphism
$\mathbb Z^2\to\mathbb Z$: choose integer
lifts of its two coefficients, and change one
by a multiple of $n$ to make their gcd one.
The possible prime divisors of the other
coefficient are avoided by the Chinese
remainder theorem; primes dividing $n$ already
cannot divide both coefficients. Pull this
homomorphism back along the connected elliptic
fibration $S\to E/G$ and reduce it modulo
$m=kn$. Its kernel defines a connected cyclic
unramified cover of degree $m$ that factors
through the canonical torus cover. A connected
finite unramified cover of a torus is the
quotient of $\mathbb C^2$ by a finite-index
sublattice, so this new cover is again a
torus. Applying \eqref{cdp:eq:C.UNN.2} proves every claimed
bielliptic multiplicity.

\end{enumerate}
\end{proof}

\subsection{Local numerical and conductor criteria}\label{subsec:II:local-numerical-conductor-inputs}

\begin{lemma}[Canonical and Chern classes on an irreducible torus-fibre tube]\hypertarget{tag-P281}{}\label{lem:II:torus-tube-numerical}
\StatementRecord{lem:II:torus-tube-numerical}{proved}{P281}{\cref{cdp:st:C.2.UN.17,lem:II:compact-analytic-tube,lem:II:conductor-duality,cdp:st:A.5.3,cdp:st:C.2.UN.5}}
Let $f:Y\to B$ be a proper holomorphic map from a smooth complex threefold to a smooth curve, with connected fibres and smooth general fibre a complex two-torus. Let $mS$ be a full fibre with irreducible reduced support, and let $h=\eta\pi:Z\to T\to S$ be its normalization and minimal normalization resolution. On a sufficiently small tube $U=f^{-1}(\Delta)$ there is an integer $k\ge0$ such that
\begin{equation}
\SsixEquationAliases{\label{cdp:eq:C.UNCR.1}}
\begin{split}
 K_U&\simeq\OO_U(kS),\qquad \OO_U(mS)\simeq\OO_U,\\
 \omega_S&\simeq\OO_U((k+1)S)|_S,\qquad
 \omega_S^{\otimes m}\simeq\OO_S.
\end{split}
\hypertarget{tag-014U}{}\hypertarget{tag-0158}{}
\label{cdp:eq:C.UI.2}
\end{equation}
Moreover,
\begin{equation}
\SsixEquationAliases{\label{cdp:eq:C.UNCR.2}}
\begin{split}
 0&=\chi(\OO_{mS})=m\chi(\OO_S),\\
 c_1(h^*TU)&=0\quad\hbox{in }H^2(Z,\mathbb Q),\qquad
 c_2(h^*TU)=0\quad\hbox{in }H^4(Z,\mathbb Q).
\end{split}
\hypertarget{tag-014V}{}\hypertarget{tag-0159}{}
\label{cdp:eq:C.UI.3}
\end{equation}
The effective integral resolution conductor $D$ satisfies
\begin{equation}
 K_Z(D)\simeq h^*\omega_S.
\hypertarget{tag-014W}{}\label{cdp:eq:C.UI.4}
\end{equation}
Thus $K_Z(D)$ is torsion, and $K_Y$ and $\OO_Y(S)$ have degree zero on every compact curve in this fibre.
\end{lemma}
\begin{proof}
Choose a coordinate $t$ vanishing at the chosen value on a disc containing no other critical value. Apply \cref{cdp:st:C.2.UN.17} to $f|_U:U\to\Delta$, whose full central fibre is the single component $mS$. Its canonical generator has divisor $kS$ with $k\ge0$, giving $K_U\simeq\OO_U(kS)$. The principal fibre equation $(t)=mS$ gives $\OO_U(mS)\simeq\OO_U$. Adjunction for the reduced hypersurface $S$ gives $\omega_S\simeq\OO_U((k+1)S)|_S$; its $m$th power is trivial by the same principal fibre equation. These are the identities in \eqref{cdp:eq:C.UI.2}.

The map is flat, since its local rings are torsion-free over the regular one-dimensional base ring. Proper flatness gives $\chi(\OO_{mS})=\chi(\OO_F)=0$ for a smooth nearby fibre $F$. The Cartier filtration of $\OO_{mS}$ has quotients $\OO_S(-jS)$ for $0\le j<m$, which are torsion line-bundle twists by \eqref{cdp:eq:C.UI.2}. Each has Euler characteristic $\chi(\OO_S)$. To see this on the compact fibre, apply its normalization sequence and the resolution Leray sequence. They express the Euler characteristic in terms of the compact smooth surface $Z$, the conductor quotient on a compact curve, and the finite-support sheaf $R^1\pi_*\OO_Z$. Smooth-surface Riemann--Roch, curve Riemann--Roch applied to the generic-rank filtration, and invariance of the finite lengths show respectively that a torsion twist leaves these three terms unchanged. Summing the filtration proves the first equality of \eqref{cdp:eq:C.UI.3}. Its first Chern-class assertion follows from \eqref{cdp:eq:C.UI.2}.

For the second Chern class, take a small closed tube as in \cref{lem:II:compact-analytic-tube}. The inverse image of a radial segment $[0,t_0]$ is a compact triangulable real five-dimensional space, oriented over the open segment by the complex fibres followed by the radial orientation. At a smooth point of $S$, the local equation $t=u z^m$ gives $m$ radial sheets with the same boundary orientation. The central singular locus has real dimension at most two and contributes no four-dimensional boundary. Its oriented fundamental chain therefore has boundary $[F]-m[S]$, proving $[F]=m[S]$ in $H_4(U,\mathbb Q)$. Since $h$ has degree one onto $S$,
\[
 \int_Zh^*c_2(TU)=\langle c_2(TU),[S]\rangle
 =\frac1m\int_Fc_2(TU|_F)=0.
\]
Here $0\to T_F\to TU|_F\to\OO_F\to0$ and $T_F$ is trivial. Top rational cohomology of the connected closed oriented four-manifold $Z$ is one-dimensional, so $c_2(h^*TU)=0$.

For nonnormal $S$, finite duality in \cref{lem:II:conductor-duality} and the minimal-resolution calculation in \cref{cdp:st:A.5.3} give the effective integral divisor $D$ and the isomorphism \eqref{cdp:eq:C.UI.4}. For normal $S$, \cref{cdp:st:C.2.UN.5} gives $Z=S$ smooth and $D=0$, so the same isomorphism is adjunction. Its torsion and the two curve-degree assertions follow from \eqref{cdp:eq:C.UI.2}.
\end{proof}



\begin{proposition}[Purity of support cohomology]\hypertarget{tag-01C4}{}\label{cdp:st:C.2.MH.1}
\StatementRecord{cdp:st:C.2.MH.1}{proved}{01C4}{\SeriesRef{cdp:eq:C.MH.2}, \SeriesRef{cdp:eq:C.MH.3}, \SeriesRef{cdp:eq:C.MH.5}; \cite[Section~3, Theorem~1 and
Corollary~2]{Steenbrink1999}; \cite[\S4, equation~(3), p.~9]{Steenbrink1999}}

Let $U$ be a smooth complex threefold, let $S\subset U$ be a
compact divisor, and suppose that an embedded resolution
$r:U'\to U$ is a sequence of blowups with smooth compact
centres in $S$, with $D=(r^{-1}S)_{\rm red}$ a simple
normal crossings divisor whose components are projective.
Let $V$ be a finite-monodromy rational local system on $U$.
Suppose that the cohomology of the punctured sufficiently
small fibre tubes around $S$, with coefficients in $V$,
vanishes in every degree. Then the mixed Hodge structure
on $H^k(S,V)$ obtained from a projective hyperresolution
is pure of weight $k$.

In this assertion the support-to-restriction morphism is
the morphism
\begin{equation}
 R\Gamma_D(U',r^*V)\longrightarrow
 \underset{D\subset W}{\mathop{\rm colim}}\,
 R\Gamma(W,r^*V)=R\Gamma(D,r^*V).
 \hypertarget{tag-00QA}{}\label{cdp:eq:C.MH.1}
\end{equation}

\end{proposition}

\begin{proof}

\begin{enumerate}[label=\arabic*.,leftmargin=*]
\item First take constant rational coefficients. Write $D_I$
for the intersection of the components indexed by a
nonempty set $I$, and $D^{[p]}=\coprod_{|I|=p+1}D_I$.
All these spaces are smooth and projective. The ordinary
\v Cech de Rham complex on $D^{[\bullet]}$, with its
\v Cech weight filtration and its differential-form
filtration, is therefore a cohomological mixed Hodge
complex $K(D)$ for $\mathbb Q_D$. Its cohomology in
degree $k$ has weights at most $k$.

\item Here is the precise analytic construction of the other
complex and of the arrow. For each stratum $C=D_I$, put
\begin{equation}
 K(C\log D)_{\mathbb C}
   =\Omega_{U'}^\bullet(\log D)\otimes\OO_C.
 \hypertarget{tag-00QJ}{}\label{cdp:eq:C.MH.2}
\end{equation}
In normal-crossing coordinates the ideal of $C$ is preserved by the logarithmic differential, so \eqref{cdp:eq:C.MH.2} is a complex. Give it the quotient residue filtration and the filtration by form degree. The local filtered comparison in \cite[\S3, Theorem~1 and Corollary~2, pp.~5--7]{Steenbrink1999} identifies the quotient map from the restricted logarithmic complex with a filtered quasi-isomorphism to \eqref{cdp:eq:C.MH.2}; its rational realization is $i_C^*Rj_*\mathbb Q_{U'\setminus D}$, where $j$ is the open inclusion. The source's proof uses the holomorphic normal-crossing coordinates and iterated residue maps. Those maps are defined here on the same divisor line bundles and agree on overlaps. Every residue stratum $D_J\cap C$ is projective by hypothesis, so its shifted de Rham complex, with Tate twist, has the purity used in that proof. This supplies the mixed Hodge complex on the compact stratum in the present analytic ambient space.

\item Totalizing \eqref{cdp:eq:C.MH.2} over $D^{[\bullet]}$ gives the
cohomological mixed Hodge complex $K(D\log D)$, with
rational realization $i_D^*Rj_*\mathbb Q_{U'\setminus D}$.
The natural inclusion $K(D)\to K(D\log D)$ is a
morphism of mixed Hodge complexes. Its mixed cone,
shifted by $[-1]$, computes cohomology with support in
$D$; the cone projection is exactly \eqref{cdp:eq:C.MH.1}. This is
the excision comparison of
\cite[\S4, equation~(3), p.~9]{Steenbrink1999}, applied
on the compact divisor. All complexes just used are
defined on that divisor, and all purity checks are on
its projective strata.

\item The weight bound on the support complex follows from its residue comparison.
The closed-support Mayer--Vietoris complex, together
with the holomorphic Thom isomorphisms, gives
\begin{equation}
 E_1^{-p,q}=
 \bigoplus_{|I|=p+1}
 H^{q-2p-2}(D_I,\mathbb Q)(-p-1)
 \ \Longrightarrow\ H_D^{q-p}(U',\mathbb Q),
 \qquad p\ge0.
 \hypertarget{tag-00QK}{}\label{cdp:eq:C.MH.3}
\end{equation}
The first differential is the alternating sum of Gysin
maps for $D_I\hookrightarrow D_{I\setminus\{i\}}$.
The comparison with the preceding mixed cone can be
verified by induction on the number of components.
For one smooth divisor the residue exact sequence
identifies
\[
 \operatorname{Cone}\bigl[\Omega_{U'}^\bullet
       \longrightarrow\Omega_{U'}^\bullet(\log D_i)\bigr][-1]
       \simeq\Omega_{D_i}^\bullet[-2](-1),
\]
where $(-1)$ records both the shift of the Hodge
filtration by one and the rational $2\pi i$ residue.
For a union $A\cup B$, use the support triangle
$R\Gamma_{A\cap B}\to R\Gamma_A\oplus R\Gamma_B
\to R\Gamma_{A\cup B}$. Iterated residues commute
with its inclusion maps, with the alternating signs
fixed by the ordering of the omitted coordinates.
Induction identifies its totalization with the same
mixed cone. These comparisons agree on overlaps,
since they are the residue maps for the holomorphic
divisor line bundles. Thus \eqref{cdp:eq:C.MH.3} compares the support complex and its arrow to $K(D)$ as filtered complexes. Its displayed term
has pure weight $q$. In total degree $k=q-p$ this is
weight $k+p\ge k$. It follows that $H_D^k(U',\mathbb Q)$
has weights at least $k$.

\item Finite coefficients are handled by taking a finite
\'{e}tale cover on which $V$ is constant and then a
rational direct summand of a finite direct sum of its
regular representation.
The inverse images of the $D_I$ are still projective:
a finite holomorphic \'{e}tale cover of a
projective manifold is projective, since the pullback
of an ample bundle is positive. The deck action acts
on the displayed mixed Hodge complexes. Rational
idempotents consequently give exactly the same
comparison and weight bounds with coefficients in $V$.

\item By the assumed punctured-tube acyclicity, the support
long exact sequence makes \eqref{cdp:eq:C.MH.1} an isomorphism on
cohomology. Properness of $r$ identifies the punctured
tubes upstairs and downstairs, since its centres lie
in $S$. The map is a morphism of mixed Hodge structures,
whose source has weights at least $k$ and whose target
has weights at most $k$. Strictness therefore proves
that $H^k(D,r^*V)$ is pure of weight $k$.

\item We finally make the comparison with the original
support explicit. The proper degree-one map between
the smooth ambient manifolds has the trace map, defined
by oriented Poincar\'e duality and proper pushforward,
and the projection formula gives
\begin{equation}
 r_*r^*=\operatorname{id}:H^k(U,V)\longrightarrow H^k(U,V).
 \hypertarget{tag-00QL}{}\label{cdp:eq:C.MH.4}
\end{equation}
The same identity holds on every smaller ambient
neighbourhood. Passing to the direct limit of these
identities gives an injection
\begin{equation}
                  r^*:H^k(S,V)\hookrightarrow H^k(D,r^*V).
 \hypertarget{tag-00QM}{}\label{cdp:eq:C.MH.5}
\end{equation}
The identifications in this limit are the usual
continuity identifications for a compact triangulable
subset. They can equally be made using cofinal
triangulated regular neighbourhoods retracting onto
$S$ and $D$. Properness ensures that inverse images of
neighbourhoods are cofinal; thus the map in \eqref{cdp:eq:C.MH.5}
is the pullback by $D\to S$.

\item This last map is a mixed Hodge morphism as well. One
may construct the projective hyperresolution of $S$
from its projective normalization resolution and
resolutions of the conductor and exceptional curves.
To compare it with $D^{[\bullet]}$, resolve their
fibre products over $S$. Each such fibre product is a
closed analytic subspace of a product of projective
spaces already occurring in the two resolutions;
Chow's theorem makes it algebraic and projective,
and its resolution is projective. The resulting
common hyperresolution realizes \eqref{cdp:eq:C.MH.5} by pullbacks
between projective strata. Cohomological descent then
identifies this with the stated mixed Hodge structures.
Injectivity into the pure target proves the assertion.

\end{enumerate}
\end{proof}

\begin{lemma}[Finite characters on irreducible torus-fibre tubes]\hypertarget{tag-P289}{}\label{lem:II:torus-tube-character}
\StatementRecord{lem:II:torus-tube-character}{proved}{P289}{\cref{lem:II:compact-analytic-tube,night:P4R.1,lem:II:exceptional-intersection,cdp:st:C.2.MH.1}}
Let $f:Y\to B$ be a proper holomorphic map from a smooth complex threefold to a smooth curve, with connected fibres and smooth general fibre a complex two-torus. Let $mS$ be a full fibre with irreducible reduced support, let $F$ be a nearby smooth fibre, and let $h=\eta\pi:Z\to T\to S$ be a resolution of its normalization. Choose a nonzero class $\gamma_S\in H^1(S,\mathbb Z)$. This class extends to a small tube $U$ and restricts nontrivially to $H^1(F,\mathbb Q)$.

For a prime $\ell$ at which both restrictions remain nonzero, the reduction of the extended class defines a cyclic cover of $U$. Its rational augmentation local system $V$ has rank $d=\ell-1$, is self-dual, satisfies $H^0(S,V)=0$, and has acyclic restriction to $F$ and $U\setminus S$. Support duality identifies $H^k(S,V)$ with $H_{6-k}(S,V)$. The normalization and resolution maps induce a surjection $H^3(S,V)\to H^3(Z,V)$ and an isomorphism $H^4(S,V)\simeq H^4(Z,V)$.

For rational or ruled $Z$, the third-cohomology surjection for the local system associated to $\gamma_S$ is an isomorphism. In these cases
\[
 e(S)=2\quad\text{for rational }Z,\qquad
 e(S)=2-2g\quad\text{for }Z\text{ ruled over a curve of genus }g.
\]
\end{lemma}
\begin{proof}
\begin{enumerate}[label=\arabic*.,leftmargin=*]
\item The homotopy equivalence in \cref{lem:II:compact-analytic-tube} extends $\gamma_S$ to $\gamma_U\in H^1(U,\mathbb Z)$. The regular-fibre homomorphism in \cref{night:P4R.1} has image with quotient $\mathbb Z/m$. A rational degree-one class vanishing on that image factors through this finite group and is zero. Thus restriction to $F$ is injective on $H^1(U,\mathbb Q)$, and $\gamma_U|_F\ne0$.

Choose one loop in $S$ and one in $F$ on which the two integral restrictions have nonzero values. Only finitely many primes divide those values. For any other prime $\ell$, reduction of $\gamma_U$ gives a connected cyclic \'etale cover of $U$, connected also over $S$ and $F$. Let $V$ be the kernel of the trace from its rational regular representation to the constant representation. It has rank $d=\ell-1$, is self-dual, and $H^0(S,V)=0$. Every nontrivial complex character of this prime-order group restricts nontrivially to $F$. One operator $\rho(e_i)-1$ in the Koszul complex for $\pi_1(F)=\mathbb Z^4$ is invertible, and its inverse contracts that complex. The punctured-disc bundle Leray sequence therefore gives
\begin{equation}
 H^j(F,V)=0,\qquad H^j(U\setminus S,V)=0
 \quad\text{for every }j.
\hypertarget{tag-00QP}{}\hypertarget{tag-0151}{}
\SsixEquationAliases{\label{cdp:eq:C.UI.9}}\label{cdp:eq:C.MH.8}
\end{equation}

\item The support sequence, the tube equivalence and oriented local Poincar\'e duality identify the direct-limit maps on shrinking tubes as
\begin{equation}
 H^k(S,V)\simeq H_S^k(U,V)\simeq H_{6-k}(S,V).
\hypertarget{tag-00QE}{}\label{cdp:eq:C.MH.13}
\end{equation}
The support has real dimension four. Consequently $H^1(S,V)=0$, and self-duality gives $h^2(S,V)=h^4(S,V)$.

\item Let $A\subset T$ and $\Gamma\subset S$ be the reduced conductor inverse image and image. Their finite normalization square is a topological pushout. Its cohomological descent sequence is
\begin{equation}
\begin{gathered}
 0\longrightarrow C_V\longrightarrow H^3(S,V)
   \longrightarrow H^3(T,V)\longrightarrow0,\\
 C_V=\operatorname{coker}\bigl[H^2(T,V)\oplus H^2(\Gamma,V)
                   \longrightarrow H^2(A,V)\bigr].
\end{gathered}
\hypertarget{tag-00QB}{}\label{cdp:eq:C.MH.10}
\end{equation}
For normal $S$, its conductor terms are zero. On each exceptional fibre of $\pi$, the pulled-back $V$ is constant. The Gysin classes of the exceptional curves restrict by their intersection matrix, which is invertible by \cref{lem:II:exceptional-intersection}. Applying its inverse separately to each coefficient vector proves that $H^2(Z,V)\to H^2(E,V)$ is surjective. Resolution descent gives $H^3(T,V)\simeq H^3(Z,V)$. The same descent sequences identify $H^4(S,V)$, $H^4(T,V)$ and $H^4(Z,V)$, since their omitted loci have real dimension at most two. These are the stated surjection and fourth-cohomology isomorphism for every $Z$.

\item Suppose now that $Z$ is rational or ruled. The strict transform in an embedded resolution of $(U,S)$ is a blowup of the projective surface $Z$. Every other surface is obtained by blowups from a projective bundle over a compact smooth curve or from $\mathbb P^2$. All centres are compact and lie over $S$, so \cref{cdp:st:C.2.MH.1} applies to the acyclic local system of \eqref{cdp:eq:C.MH.8} and gives
\begin{equation}
 H^3(S,V)\text{ pure of weight }3.
\hypertarget{tag-00QQ}{}\label{cdp:eq:C.MH.9}
\end{equation}
The normalization descent sequence \eqref{cdp:eq:C.MH.10} is then a sequence of mixed Hodge structures obtained from the projective resolutions. The top cohomology $H^2(A,V)$ of a compact curve with a finite local system is pure of weight two; its quotient $C_V$ has the same weight. Strictness and \eqref{cdp:eq:C.MH.9} force $C_V=0$. The maps already constructed consequently give
\begin{equation}
\SsixEquationAliases{\label{cdp:eq:C.MH.11}\label{cdp:eq:C.UI.10}}
\begin{split}
 H^3(S,V)&\simeq H^3(T,V)\simeq H^3(Z,V),\\
 H^4(S,V)&\simeq H^4(T,V)\simeq H^4(Z,V).
\end{split}
\hypertarget{tag-00QC}{}\hypertarget{tag-00QD}{}\hypertarget{tag-014N}{}
\label{cdp:eq:C.MH.12}
\end{equation}

\item Combining \eqref{cdp:eq:C.MH.13} with these isomorphisms gives $\chi(S,V)=2h^4(Z,V)-h^3(Z,V)$. A finite local system of rank $d$ has Euler characteristic $d\,e(S)$. If $Z$ is rational, its pulled-back local system is constant, so $h^4(Z,V)=d$ and $h^3(Z,V)=0$.

For $Z$ ruled over a genus-$g$ curve $B$, point blowups preserve its fundamental group and third cohomology. The local system pulls back from $B$, and the projective-bundle formula gives $H^3(Z,V)=H^1(B,V)(-1)$. When the cyclic character on $B$ is trivial, the two dimensions are $d$ and $2gd$. Otherwise each nontrivial character remains nontrivial on $B$, and the dimensions are $0$ and $(2g-2)d$. Thus
\begin{equation}
 d\,e(S)=\chi(S,V)=2h^4(Z,V)-h^3(Z,V)
 =\begin{cases}
 2d,&Z\text{ rational},\\
 (2-2g)d,&Z\text{ ruled over genus }g.
 \end{cases}
\hypertarget{tag-00QF}{}\hypertarget{tag-014O}{}
\SsixEquationAliases{\label{cdp:eq:C.MH.14}}\label{cdp:eq:C.UI.11}
\end{equation}
This proves both Euler identities.
\end{enumerate}
\end{proof}

\begin{proposition}[Euler number from a ruled normalization]\hypertarget{tag-01C5}{}\label{cdp:st:C.2.MH.2}
\StatementRecord{cdp:st:C.2.MH.2}{proved}{01C5}{\cref{cdp:st:D.1.9,cdp:st:D.1.8,lem:II:torus-tube-character}}

Let $f:Y\to\PP^1$ satisfy the full CDP hypotheses
$b_1(Y)=b_2(Y)=0$ with smooth torus general fibre.
Let $mS$ be a full fibre with irreducible reduced support,
and let $Z\to T\to S$ be a resolution of its
normalization. If $Z$ is rational, then
\begin{equation}
                              e(S)=2.
 \hypertarget{tag-00QN}{}\label{cdp:eq:C.MH.6}
\end{equation}
If $Z$ is ruled over a curve of genus $g\ge1$, then
\begin{equation}
                              e(S)=2-2g.
 \hypertarget{tag-00QO}{}\label{cdp:eq:C.MH.7}
\end{equation}

\end{proposition}

\begin{proof}
The integral class of \cref{cdp:st:D.1.9} restricts nontrivially to $S$ and, by \cref{cdp:st:D.1.8}, to the nearby torus. On a small disc $\Delta$ the integral Leray edge map is the isomorphism
\[
 H^1(f^{-1}\Delta,\mathbb Z)\simeq
 H^0(\Delta,R^1f_*\mathbb Z),
\]
because the two adjacent groups $H^1(\Delta,\mathbb Z)$ and $H^2(\Delta,\mathbb Z)$ vanish. It therefore gives the same tube class as the homotopy equivalence in \cref{lem:II:compact-analytic-tube}: both restrict to the stated class on $S$. Apply \cref{lem:II:torus-tube-character} to this class and to the given normalization resolution. Its rational and ruled alternatives are exactly \eqref{cdp:eq:C.MH.6} and \eqref{cdp:eq:C.MH.7}.
\end{proof}

\subsection{Class VII canonical boundaries}\label{subsec:II:class-vii-canonical-boundaries}

\begin{lemma}[Euler number of a curve with rational and elliptic normalizations]\hypertarget{tag-P290}{}\label{lem:II:curve-euler-delta}
\StatementRecord{lem:II:curve-euler-delta}{proved}{P290}{\cref{lem:II:curve-normalization-graph}}
Let $C$ be a compact reduced complex curve whose normalized irreducible components have genera $g_i\in\{0,1\}$. Let $n_0$ count the components of genus zero, and let $s_x$ and $\delta_x$ be the number of normalization branches and the delta invariant at a singular point. Then
\begin{equation}
\begin{aligned}
 \chi(\OO_C)&=\sum_i(1-g_i)-\sum_x\delta_x,\\
 e(C)&=\sum_i(2-2g_i)-\sum_x(s_x-1),
 \qquad \delta_x\ge s_x-1.
\end{aligned}
\hypertarget{tag-015R}{}\label{cdp:eq:C.VII.5}
\end{equation}
For $\chi(\OO_C)\ge0$ one has
\begin{equation}
 e(C)=\chi(\OO_C)+n_0+
       \sum_x\bigl(\delta_x-(s_x-1)\bigr)\ge0,
\hypertarget{tag-00O8}{}\label{cdp:eq:C.INT.26}
\end{equation}
and $e(C)=0$ precisely when $C$ is a disjoint union of smooth elliptic curves.
\end{lemma}
\begin{proof}
The coherent normalization sequence has finite quotient of length $\delta_x$ at $x$, which gives the first formula. The topological normalization comparison in \cref{lem:II:curve-normalization-graph}, applied to each connected component, gives the second. Independent constants on the normalized branches, modulo their common constant, have dimension $s_x-1$ and form a quotient of the normalization quotient. Hence $\delta_x\ge s_x-1$. Subtracting the two Euler formulas gives \eqref{cdp:eq:C.INT.26}.

Under $\chi(\OO_C)\ge0$, all three terms on its right are nonnegative. Equality forces $n_0=0$, after which $\chi(\OO_C)=-\sum_x\delta_x\ge0$ forces every $\delta_x=0$. Thus each component is smooth elliptic and different components are disjoint. Such a union has Euler number zero, proving the converse. The empty curve satisfies the same formulas.
\end{proof}



\begin{proposition}[Conductor Euler bounds]\hypertarget{tag-01F1}{}\label{cdp:st:C.2.VII.2}
\StatementRecord{cdp:st:C.2.VII.2}{proved}{01F1}{\cref{cdp:st:A.5.12}, \cref{cdp:st:A.5.3}, \cref{cdp:st:C.2.VII.1}, \eqref{cdp:eq:C.VII.5}, \eqref{cdp:eq:C.VII.6}}

Let $S$ be a compact irreducible reduced hypersurface in
a smooth complex threefold, with $\omega_S$ torsion and
$\chi(S,\OO_S)=0$. Let $\eta:T\to S$ be normalization,
let $\Gamma$ be its conductor image scheme, and put
$C=\Gamma_{\mathrm{red}}$. Then
\begin{equation}
 \chi(C,\OO_C)\ge-1.
\hypertarget{tag-015P}{}\label{cdp:eq:C.VII.3}
\end{equation}
Suppose in addition that the minimal resolution
$\pi:Z\to T$ is a class VII surface. Then
\begin{equation}
 e(C)\ge0.
\hypertarget{tag-015Q}{}\label{cdp:eq:C.VII.4}
\end{equation}

\end{proposition}

\begin{proof}

\begin{enumerate}[label=\arabic*.,leftmargin=*]
\item If the conductor is empty, both conclusions are immediate.
We therefore assume that it is nonempty.
Put $L=\omega_S$ and
$\varepsilon=h^0(S,L)\in\{0,1\}$.
A nonzero section of a torsion line on a compact
connected reduced space trivializes it, by taking a
torsion-trivializing power. Serre duality and
$\chi(\OO_S)=0$ therefore give
$h^1(S,\OO_S)=1+\varepsilon$, with
$\varepsilon=1$ exactly when $L=\OO_S$.
The normalization quotient
$\mathcal Q=\eta_*\OO_T/\OO_S$ has
$H^0(\mathcal Q)\hookrightarrow H^1(S,\OO_S)$,
because $S$ and $T$ have the same constant holomorphic
functions. The finite conductor duality of \cref{lem:II:conductor-duality} is
\[
 \mathcal Q=\omega_\Gamma\otimes L^{-1}|_\Gamma .
\]
Serre duality on the Cohen--Macaulay conductor curve
gives $h^1(\Gamma,L|_\Gamma)\le1+\varepsilon$.
The restriction onto $C$ is a quotient of coherent
sheaves on a one-dimensional space. Its kernel has
no second cohomology, so
$h^1(C,L|_C)\le1+\varepsilon$.
Since $L|_C$ is torsion, it has degree zero on every
irreducible component and
$\chi(C,L|_C)=\chi(C,\OO_C)$.
If $\varepsilon=0$, Riemann--Roch gives
$\chi(\OO_C)=h^0(C,L|_C)-h^1(C,L|_C)\ge-1$.
If $\varepsilon=1$, then $L=\OO_S$ and
$h^0(C,\OO_C)\ge1$, while $h^1(C,\OO_C)\le2$.
This proves \eqref{cdp:eq:C.VII.3}, including all nilpotents of
the original conductor scheme in the comparison.

\item Now assume that $Z$ is class VII. \cref{cdp:st:A.5.3} gives a
nonzero effective integral divisor $D$ with
$K_Z+D=(\eta\pi)^*\omega_S$.
Each curve in the reduced strict conductor on $Z$
is rational or elliptic by \ref{cdp:st:C.2.VII.1}. Its finite image
in $C$ has normalization of genus at most one, by
Riemann--Hurwitz. Apply \cref{lem:II:curve-euler-delta}, with $n_0$ the number of rational-normalized components and $r_x=s_x$ the number of branches at $x$. Its formulas \eqref{cdp:eq:C.VII.5}, together with \eqref{cdp:eq:C.VII.3}, give $e(C)\ge n_0-1$.

\item The only possible negative value is consequently
$e(C)=-1$, with every normalized component elliptic.
Then \eqref{cdp:eq:C.VII.3} says $\sum\delta_x\le1$.
Negativity requires two branches, so there is exactly
one singular point, it has two branches and
$\delta=1$, and the local curve ring is
\begin{equation}
 \{(a(u),b(v))\in\C\{u\}\oplus\C\{v\}:a(0)=b(0)\}.
\hypertarget{tag-015S}{}\label{cdp:eq:C.VII.6}
\end{equation}
Indeed that subring already has colength one in the
normalization, and the curve ring is contained in it
with the same colength. Thus the point is an ordinary
node. Its two branches come either from distinct
elliptic normalizations or from two distinct points
of one elliptic normalization.

\item We show that this node is incompatible with the
hypersurface and resolution maps. At least
one normalization point $t\in T$ over the node receives
preimages of both local branches of $C$.
Otherwise partition the irreducible local branches
of the hypersurface germ $S$ according to the one
conductor branch their normalization receives. Both
parts are nonempty. Two hypersurface branches from
opposite parts intersect in a curve: in the regular
three-dimensional local ambient ring, the ideal of
their two distinct equations has height two.
That intersection curve is contained in the conductor
image and therefore is one of its two branches.
Both surface branches must receive it, contrary to
the partition. With only one local surface branch
the assertion is immediate.

\item Every $\pi$-exceptional curve not contained in
$D_{\mathrm{red}}$ is disjoint from it. To see this,
write $E$ for such a curve. Minimality of the
resolution and adjunction give $K_Z\cdot E\ge0$,
while $D\cdot E\ge0$, and
$(K_Z+D)\cdot E=0$ forces equality. In particular
$D\cdot E=0$. A resolution fibre over a point of
the conductor must meet its strict transform, by
properness. Since its fibre is connected, this
disjointness proves
$\pi^{-1}(t)\subset D_{\mathrm{red}}$ for the
normalization point just found.

\item The strict curves over the two elliptic-normalized
branches are elliptic: a rational curve cannot map
nonconstantly onto an elliptic curve. They meet the
connected fibre $\pi^{-1}(t)$. If the curves are
distinct, one connected component of $D_{\mathrm{red}}$
therefore contains two elliptic curves. If they are
the same elliptic curve, their meeting points with
the fibre are distinct. This follows from the exact
factorization of a map from a normal curve through
the normalization of its image: one point cannot
map to both points over the node in \eqref{cdp:eq:C.VII.6}.
The connected resolution fibre consequently gives
a path in $D_{\mathrm{red}}$ joining these distinct
points. Both possibilities contradict \cref{cdp:st:C.2.VII.1}:
the first has two positive-genus components, and
the second adds a graph cycle to an elliptic
component, giving arithmetic genus at least two.
This excludes the only negative possibility and
proves \eqref{cdp:eq:C.VII.4}.

\end{enumerate}
\end{proof}

\begin{proposition}[Euler number for class VII normalization]\hypertarget{tag-01F2}{}\label{cdp:st:C.2.VII.3}
\StatementRecord{cdp:st:C.2.VII.3}{proved}{01F2}{\cref{cdp:st:A.5.12}, \cref{cdp:st:A.5.3}, \cref{cdp:st:C.2.VII.1}, \cref{cdp:st:C.2.VII.2}}

Let $f:Y\to B$ be a proper connected-fibre holomorphic
map from a smooth compact complex threefold with
$b_2(Y)=0$ to a smooth curve, whose general fibre is a
complex two-torus. Let $mS$ be a full fibre with
irreducible reduced nonnormal support, and suppose
that the minimal resolution $Z$ of its normalization
is class VII. Then
\begin{equation}
 e(S)\ge0.
\hypertarget{tag-015T}{}\label{cdp:eq:C.VII.7}
\end{equation}

\end{proposition}

\begin{proof}

\begin{enumerate}[label=\arabic*.,leftmargin=*]
\item Let $\eta:T\to S$ be normalization, let $\pi:Z\to T$ be its minimal resolution, and let $D>0$ be the resolution conductor. By \cref{lem:II:torus-tube-numerical}, $\omega_S$ is torsion, $\chi(\OO_S)=0$ and $K_Z+D=(\eta\pi)^*\omega_S$. Put $C=\Gamma_{\mathrm{red}}$ and $A=(\eta^{-1}C)_{\mathrm{red}}$. Thus \cref{cdp:st:C.2.VII.1,cdp:st:C.2.VII.2} apply to these maps and divisors.

\item We compute $e(T\setminus A)$ through the resolution. As in \ref{cdp:st:C.2.VII.2}, an exceptional curve
$E$ with coefficient zero in $D$ has
$K_Z\cdot E=D\cdot E=0$.
Adjunction and negative self-intersection then give
$p_a(E)=0$ and $E^2=-2$.
These curves are smooth rational and disjoint
from $D_{\mathrm{red}}$.
Their intersection matrix is negative definite by \cref{lem:II:exceptional-intersection}; this makes
each connected configuration a tree: a repeated
edge or a cycle would give a nonnegative square
for a nonzero vector in the matrix. The intersections
are transverse, since their local intersection
multiplicities are one; three curves cannot meet
at a point, since that would give a triangle.
Connectedness of resolution fibres shows that a
fibre either consists entirely of these curves or
lies entirely in $D_{\mathrm{red}}$.
The zero-coefficient fibres lie outside $A$, since
a fibre over $A$ must meet the strict conductor.

\item Let $a$ be the total number of zero-coefficient
exceptional curves, and let $b$ be the number of
exceptional fibres lying in $D_{\mathrm{red}}$
whose images in $T$ lie outside $A$. On
$Z\setminus D_{\mathrm{red}}$ the former fibres
are still present and the latter are absent.
A tree with $a_j$ rational curves has Euler
number $a_j+1$, so replacing it by its image
point subtracts $a_j$. Adding back the image
point of an absent positive-coefficient fibre
adds one. Away from these finitely many fibres,
$\pi$ is the required isomorphism of complements.
Additivity of compactly supported Euler
characteristic therefore gives the exact equality
\begin{equation}
 \begin{aligned}
 e(T\setminus A)
   &=e(Z)-e(D_{\mathrm{red}})-a+b\\
   &=b_2(Z)-r(D_{\mathrm{red}})-a+b
   \ \ge0 .
 \end{aligned}
\hypertarget{tag-015U}{}\label{cdp:eq:C.VII.8}
\end{equation}
The inequality follows because the rational curves
counted by $r(D_{\mathrm{red}})$ and by $a$ are
distinct, and their total is at most $b_2(Z)$ by
\ref{cdp:st:C.2.VII.1}.
Finally $\eta$ is a biholomorphism
$T\setminus A\to S\setminus C$. Hence
\[
 e(S)=e(T\setminus A)+e(C)\ge0
\]
by \ref{cdp:st:C.2.VII.2}. This proves the assertion for the
original reduced support and all full multiplicities.

\end{enumerate}
\end{proof}

\begin{proposition}[Zero Euler number for class VII normalization]\hypertarget{tag-01F3}{}\label{cdp:st:C.2.VII.4}
\StatementRecord{cdp:st:C.2.VII.4}{proved}{01F3}{\cref{cdp:st:C.2.VII.3}, \cref{cdp:st:C.2.VII.2}, \cref{cdp:st:A.5.12}, \cref{cdp:st:C.2.VII.1}, \cref{lem:II:curve-euler-delta,lem:II:last-blowup-minimal-resolution}}

Use the irreducible class VII situation of
\ref{cdp:st:C.2.VII.3}. Let $C=\Gamma_{\rm red}$, let $n_0$
be the number of its rational-normalized components,
and put $v(C)=\sum_x(s_x-1)$, where $s_x$ is the
number of branches at a singular point of $C$.
On the minimal resolution $Z$ let $D$ be the
resolution conductor, let $r$ count the rational
curves in $D_{\rm red}$, and let $a$ count the
zero-coefficient exceptional curves lying over
points outside the conductor. Let $b$ count the
positive-coefficient exceptional fibres over
such points, as in \ref{cdp:st:C.2.VII.3}. Then $e(S)=0$
is equivalent to the following explicit
normalization and curve data:
\begin{equation}
 b=0,\qquad b_2(Z)=r+a,
 \qquad
 \begin{cases}
 n_0=0, v(C)=0,\quad\text{or}\\
 n_0=1, v(C)=2.
 \end{cases}
 \hypertarget{tag-015V}{}\label{cdp:eq:C.VII.9}
\end{equation}
In the first curve alternative all components
have elliptic normalization and the sum of
their delta invariants is at most one. In the
second the total delta invariant is exactly
two, and every singular point satisfies
$\delta_x=s_x-1$.

If the homomorphism
$\Pic^0(Y)\to\Pic^0(Z)$ is nonconstant,
there is the sharper geometric equivalence:
$e(S)=0$ exactly when its normalization is a
smooth minimal Hopf surface and the reduced
conductor image and inverse image are disjoint
unions of smooth elliptic curves, with finite
unramified maps between them.
\end{proposition}

\begin{proof}

\begin{enumerate}[label=\arabic*.,leftmargin=*]
\item The exact comparison already proved in
\ref{cdp:st:C.2.VII.3} is
\begin{equation}
 e(S)=b_2(Z)-r-a+b+e(C),
 \qquad b_2(Z)-r-a\ge0,\quad b\ge0,\quad e(C)\ge0.
 \hypertarget{tag-015L}{}\label{cdp:eq:C.VII.10}
\end{equation}
Every normalized component of $C$ has genus
zero or one. Its topological normalization
formula is $e(C)=2n_0-v(C)$. The stronger
bound in \ref{cdp:st:C.2.VII.2} is $e(C)\ge n_0-1$.
Thus $e(C)=0$ permits precisely $n_0=0,1$,
with $v(C)=0,2$, respectively. This proves
both directions of \eqref{cdp:eq:C.VII.9} for these normalization branch numbers.
The same proof gives
$\chi(\OO_C)=n_0-\sum\delta_x\ge-1$.
For $n_0=0$ this bounds the total delta
invariant by one. For $n_0=1$, the value
$v(C)=2$ and $\delta_x\ge s_x-1$ force
$\sum\delta_x=2$ and equality at every point.

\item Now suppose that the Picard pullback is
nonconstant. The argument of \cref{cdp:st:A.5.12} applies
to $\Pic^0(Y)$ even with only rational
$H^2(Y,\Q)=0$. Namely its nonzero differential
is the composite
\begin{equation}
 H^1(Y,\OO_Y)\longrightarrow H^1(S,\OO_S)
 \longrightarrow H^1(T,\OO_T)
 \longrightarrow H^1(Z,\OO_Z).
 \hypertarget{tag-015M}{}\label{cdp:eq:C.VII.11}
\end{equation}
The middle arrow consequently has positive
rank. Put $\varepsilon=h^0(S,\omega_S)$;
the torsion canonical relation makes
$\varepsilon\in\{0,1\}$ and
$h^1(S,\OO_S)=1+\varepsilon$.
The normalization sequence then gives
$h^0(\mathcal Q)\le\varepsilon$.
Conductor duality and restriction onto $C$
therefore give
\begin{equation}
 h^1(C,\omega_S|_C)\le\varepsilon.
 \hypertarget{tag-015N}{}\label{cdp:eq:C.VII.12}
\end{equation}
If $\varepsilon=0$, degree-zero
Riemann--Roch gives $\chi(\OO_C)\ge0$.
If $\varepsilon=1$, the torsion line
$\omega_S$ is trivial, and
$h^0(C,\OO_C)\ge1$ gives the same inequality.
Lemma~\ref{lem:II:curve-euler-delta} and \eqref{cdp:eq:C.INT.26} imply
that $e(C)=0$ holds precisely when $C$ is
a disjoint union of smooth elliptic curves.

\item Assume $e(S)=0$. Every rational curve in
$D_{\rm red}$ must be resolution-exceptional,
because its image lies in $C$ and cannot
dominate an elliptic curve. A minimal class
VII model with positive second Betti number
and a torsion anticanonical divisor is one
of the proved configurations used in
\ref{cdp:st:C.2.VII.1}. An intermediate or
Inoue--Hirzebruch model has no elliptic
curves, and therefore could not dominate
the nonempty elliptic conductor. A
parabolic Inoue model has a rational cycle
of square zero and exactly $b_2$ rational
curves. Its point blowups therefore still
have exactly $b_2(Z)$ rational curves.
The equality $b_2(Z)=r+a$ makes every one
of them resolution-exceptional, including
any zero-coefficient blowup curves. Thus
the full total transform of the minimal
rational cycle, whose square remains zero,
would be exceptional. This contradicts \cref{lem:II:exceptional-intersection}, since its nonzero exceptional divisor has square zero.
Consequently the minimal model is Hopf.
All rational curves on $Z$ are now the
exceptional curves of its point blowups,
and their number is $b_2(Z)$.
The equality $b_2(Z)=r+a$ makes all of
them resolution-exceptional. For the point-blowup map $\sigma:Z\to Z_0$ to this Hopf minimal model, every $\sigma$-exceptional curve is therefore $\pi$-exceptional. Lemma~\ref{lem:II:last-blowup-minimal-resolution} makes $\sigma$ an isomorphism. \Cref{lem:II:hopf-curves} gives square zero for every curve on this minimal Hopf surface $Z$. Lemma~\ref{lem:II:exceptional-intersection} therefore excludes every $\pi$-exceptional curve. Thus $\pi$ is an isomorphism, and the normalization is already a smooth minimal Hopf surface.
Its reduced conductor curves are elliptic;
their finite maps to the smooth elliptic
curves of $C$ are unramified by
Riemann--Hurwitz.
Conversely those normalization and
conductor data have Euler number zero
by the exact normalization formula.
This proves the sharper equivalence.

\end{enumerate}
\end{proof}

\subsection{Contracted ruled boundaries}\label{subsec:II:contracted-ruled-boundaries}

\begin{lemma}[Duality and ruling genus for a contracted section]\hypertarget{tag-P293}{}\label{lem:II:contracted-section-duality}
\StatementRecord{lem:II:contracted-section-duality}{proved}{P293}{\cref{lem:II:torus-tube-numerical,lem:II:conductor-duality,lem:II:blowup-topology,lem:II:delta-two-curve-rings,cdp:st:A.5.3}; \cite[\S\S2.2.2--2.2.5, pp.~60--61]{vanStraten1987}}
\hypertarget{tag-01BR}{}\label{cdp:st:C.2.CR.2}
\StatementRecord{cdp:st:C.2.CR.2}{proved}{01BR}{}

Let $f:Y\to B_0$ be a proper holomorphic map from a smooth complex threefold to a smooth curve, with connected fibres and smooth general fibre a complex two-torus. Let $mS$ be an irreducible nonnormal full fibre, let $\eta:T\to S$ be its normalization, and let $\pi:Z\to T$ be the minimal resolution. Put $h=\eta\pi$. Suppose that $p:Z\to B$ is an irrational ruling and that its effective adjunction divisor is $D=2C+V$, with $C$ a section and $V$ vertical. Suppose that $\pi$ contracts $C$, and put $g=g(B)$, $p_0=\pi(C)$ and $q=\eta(p_0)$.

Then
\[
 H^1(T;\mathbb R)=H^1(T,\OO_T)=H^2(T,\OO_T)=0.
\]
The sheaf $R^1\pi_*\OO_Z$ is supported only at $p_0$ and has length $g$. After fixing a local generator of $\omega_S$ at $q$, the conductor ideal $I_{p_0}$ of the normalization branch $(T,p_0)$ embeds as an ideal of $R=\OO_{S,q}$; with $J=(\pi_*\omega_Z)_{p_0}$, its residue comparison is \eqref{cdp:eq:C.CR.12} and it satisfies \eqref{cdp:eq:C.CR.13}. One has $g\in\{1,2\}$. For $g=2$, the point $p_0$ is the only point of $T$ over $q$, and the full conductor-image scheme $\Gamma$ is smooth and reduced at $q$. Its component through $q$ has generic transversal an ordinary node or an ordinary cusp.
\end{lemma}
\begin{proof}
\begin{enumerate}[label=\arabic*.,leftmargin=*]
\item Restrict $f$ over a small disc about the given fibre. The same $S,T,Z,C,V$ occur in \cref{lem:II:torus-tube-numerical}; it gives $\omega_S$ torsion, $\chi(\OO_S)=0$, and $K_Z+D=h^*\omega_S$. By the finite-duality inclusion \eqref{cdp:eq:A.RT.5} of \cref{lem:II:conductor-duality}, a section of $\omega_T$ is a section of $\eta^*\omega_S$ vanishing on the nonempty conductor. A nonzero section of a torsion line bundle on the connected compact normal surface $T$ is nowhere zero: a trivializing tensor power is a nonzero constant. Thus $H^0(T,\omega_T)=0$, and Serre duality gives $H^2(T,\OO_T)=0$.

Since $h^1(Z,\OO_Z)=g$ and $h^2(Z,\OO_Z)=0$, the coherent Leray sequence is
\begin{equation}
 0\longrightarrow H^1(T,\OO_T)\longrightarrow H^1(Z,\OO_Z)
 \longrightarrow H^0(T,R^1\pi_*\OO_Z)\longrightarrow0.
\hypertarget{tag-00DB}{}\label{cdp:eq:C.CR.4}
\end{equation}
The higher direct image has finite support. At $p_0$, formal functions identifies its completed stalk with the inverse limit of the $H^1$ of the infinitesimal exceptional fibres. The transition maps on $H^1$ are surjective because their kernels have support of dimension at most one. Restriction from the scheme-theoretic exceptional fibre to $C$ is surjective on $H^1$ for the same reason. The finite-length stalk equals its completion, so its length is at least $h^1(C,\OO_C)=g$. Equation~\eqref{cdp:eq:C.CR.4} bounds the total length by $g$. Equality gives $H^1(T,\OO_T)=0$ and shows that $p_0$ supports the entire higher direct image.

\item Connectedness of resolution fibres makes the constant-sheaf Leray map $H^1(T;\mathbb R)\to H^1(Z;\mathbb R)$ injective. Factor the ruling through its minimal ruled bundle. Its simply connected fibres and the point-blowup isomorphisms in \cref{lem:II:blowup-topology} identify $H^1(Z;\mathbb R)$ with $H^1(B;\mathbb R)$ by pullback. Since $p|_C:C\to B$ is an isomorphism, restriction from $Z$ to $C$ is the inverse identification. Every class pulled back from $T$ restricts to zero on $C$, because $\pi(C)=p_0$. Hence $H^1(T;\mathbb R)=0$.

\item Fix a generator of $\omega_S$ on a small neighbourhood of $q$. The multiplication map in \cref{lem:II:conductor-duality} identifies the branch conductor with
\[
 I_{p_0}=\omega_{T,p_0}\subset\OO_{T,p_0}.
\]
The semilocal normalization over $q$ is the product of its branch rings. Its conductor is the direct sum of their conductor ideals, embedded in $R$ by multiplication; extending by zero on the other branches embeds $I_{p_0}$ itself as an ideal of $R$.

The normal surface $(T,p_0)$ is Cohen--Macaulay. Apply proper duality and Grauert--Riemenschneider vanishing to the resolution over this neighbourhood, in the analytic form of \cite[\S\S2.2.2--2.2.5, pp.~60--61]{vanStraten1987}. The defect of $\OO_T\to R\pi_*\OO_Z$ is $R^1\pi_*\OO_Z$ in degree one. Duality identifies its finite-length dual with $\omega_T/\pi_*\omega_Z$, since the positive direct images of $\omega_Z$ vanish. With $J=(\pi_*\omega_Z)_{p_0}$, this gives
\begin{equation}
 \operatorname{length}(I_{p_0}/J)
  =\operatorname{length}(R^1\pi_*\OO_Z)_{p_0}=g.
 \hypertarget{tag-00CW}{}\label{cdp:eq:C.CR.10}
\end{equation}

\item The residue sequence of the smooth exceptional section is
\begin{equation}
 0\longrightarrow\omega_Z\longrightarrow\omega_Z(C)
   \longrightarrow\omega_C\longrightarrow0.
 \hypertarget{tag-00CX}{}\label{cdp:eq:C.CR.11}
\end{equation}
The same vanishing makes its pushforward quotient at $p_0$ equal to $H^0(C,\omega_C)$, of dimension $g$. The sheaf $\pi_*\omega_Z(C)$ embeds in $\omega_T$: away from the exceptional set its sections are holomorphic two-forms, and the reflexive canonical module of a normal surface extends them across its isolated singularities. Comparing lengths with \eqref{cdp:eq:C.CR.10} proves
\begin{equation}
\hypertarget{tag-015A}{}\SsixEquationAliases{\label{cdp:eq:C.UNCR.3}}
 I_{p_0}=(\pi_*\omega_Z(C))_{p_0},\qquad
 I_{p_0}/J\simeq H^0(C,\omega_C).
 \hypertarget{tag-00CY}{}\label{cdp:eq:C.CR.12}
\end{equation}
Use the adjunction comparison of \cref{cdp:st:A.5.3} with the chosen generator of $\omega_S$. Its pulled-back meromorphic two-form has divisor $-D$, so the comparison identifies $\omega_Z$ with $\OO_Z(-2C-V)$ and $\omega_Z(C)$ with $\OO_Z(-C-V)$ as subsheaves in that trivialization. Thus the quotient in \eqref{cdp:eq:C.CR.12} records the coefficient of order one along $C$, retaining every vertical vanishing condition in $V$.

\item Write $\mathfrak m$ for the maximal ideal of $R$. An element of $\mathfrak m$ pulls back to a function vanishing on $C$, since $h(C)=q$. Every element of $\mathfrak m^2$ therefore vanishes there to order at least two. An element of $I_{p_0}$ already has the prescribed vanishing along $V$ by \eqref{cdp:eq:C.CR.12}; the additional order along $C$ therefore gives
\begin{equation}
                 I_{p_0}\cap\mathfrak m^2\subset J.
 \hypertarget{tag-00CZ}{}\label{cdp:eq:C.CR.13}
\end{equation}
The same comparison puts $I_{p_0}\subset\mathfrak m$. It also shows that $q\in\Gamma$: otherwise $S$ would be normal at $q$, its only branch conductor would be $R$, and its element $1$ would have to vanish along $C$.

Consequently
\begin{equation}
 g=\dim(I_{p_0}/J)
 \le \dim\bigl(I_{p_0}/(I_{p_0}\cap\mathfrak m^2)\bigr)
 \le \dim(\mathfrak m/\mathfrak m^2)\le3.
 \hypertarget{tag-00D0}{}\label{cdp:eq:C.CR.14}
\end{equation}
A middle dimension of three would make $I_{p_0}$ generate the whole cotangent space. Nakayama's lemma for $\mathfrak m/I_{p_0}$ would give $I_{p_0}=\mathfrak m$. The full conductor ideal contains this ideal and is proper, so it would also equal $\mathfrak m$. This contradicts the pure one-dimensional conductor image in \cref{lem:II:conductor-duality}. Thus the middle dimension is at most two, and $1\le g\le2$.

\item For $g=2$, choose $x,y\in I_{p_0}$ with independent linear parts and lift them to ambient coordinates $(x,y,z)$. The full conductor ideal contains $(x,y)$ in $R$. If $F$ is the defining equation of $S$, the one-dimensional conductor quotient forces $F\in(x,y)$; hence $R/(x,y)=\C\{z\}$. Any strictly larger proper ideal would have a zero-dimensional quotient, so purity makes the full conductor ideal exactly $(x,y)$. This proves that $\Gamma$ is smooth and reduced at $q$.

There is no other point of $T$ over $q$. The elements $x,y$ lie in the single branch summand $I_{p_0}$ and vanish identically on every other branch. Such a two-dimensional branch cannot lie in the one-dimensional ambient germ $\{x=y=0\}$.

Finally the conductor-duality isomorphism $\eta_*\OO_T/\OO_S\simeq\omega_\Gamma\otimes\omega_S^{-1}|_\Gamma$ makes this quotient a line bundle near $q$. The generic transverse normalization quotient therefore has length one. For a reduced plane curve with $r$ branches, the branch constants already give $\delta\ge r-1$, so here $r\le2$. The delta-one cases of \cref{lem:II:delta-two-curve-rings} are exactly the ordinary node and ordinary cusp.
\end{enumerate}
\end{proof}

\begin{proposition}[A smooth full critical curve with generic node is a product node]\hypertarget{tag-01EZ}{}\label{cdp:st:C.2.UN.9}
\StatementRecord{cdp:st:C.2.UN.9}{proved}{01EZ}{\cref{lem:II:smooth-curve-intersection,cdp:st:A.5.15,lem:II:boundary-deformation-transport}}

Let a reduced hypersurface germ $\{f=0\}$ in a smooth
complex threefold occur in a fibre of a proper torus-surface
degeneration. Suppose that the full reduced critical locus
of $f$ is a smooth curve and its generic transverse
singularity is an ordinary node. There are local
holomorphic coordinates in which
\begin{equation}
                            f=XY.                  \hypertarget{tag-015D}{}\label{cdp:eq:C.UNNODE.1}
\end{equation}
The conclusion about the reduced hypersurface also holds
when the map has a holomorphic unit and a common
positive multiplicity.
\end{proposition}

\begin{proof}
\begin{enumerate}[label=\arabic*.,leftmargin=*]
\item The entire reduced critical curve is a smooth axis $\Gamma$. Each local irreducible component of a reducible support meets every other component in a curve contained in $\Gamma$, and therefore contains $\Gamma$. The generic ordinary node has two branches, so there are at most two local components. If the generator is reducible, write $f=V\phi\psi$ with the two reduced factors and a unit $V$. Their common zero set is $\Gamma$ and their generic transverse intersection multiplicity is one. Apply \cref{lem:II:smooth-curve-intersection} with $r=1$: $(\phi,\psi,z)$ is a coordinate system. Replacing its first coordinate by $V\phi$, keeping the other two, is invertible, with inverse the unique solution of $X=V(\phi,\psi,z)\phi$. It gives $f=XY$ exactly. This argument permits every original positive full multiplicity on the two components.

\item In the locally irreducible case the full map has the form $F=Uf^M$. Choose the root $f_0=U^{1/M}f$, so $F=f_0^M$ exactly. Its critical-set germ is still $\Gamma$: every critical component through the point has constant value zero, and on that zero set $df_0=U^{1/M}df$. Choose coordinates $(x,y,z)$ with $\Gamma=\{x=y=0\}$. The transverse quadratic term of $f_0$ is
\[
                       (x,y)Q(z)(x,y)^{\mathsf t},
\]
where $Q$ is symmetric and has determinant not identically zero. Shrink so that $\det Q(z)\ne0$ for $z\ne0$. If $Q(0)$ is invertible, the parameter Morse map gives $f_0=X^2+Y^2$. The linear coordinates $X+\mathrm iY,X-\mathrm iY$ give its product form. Absorbing the remaining unit $U^{-1/M}$ into the first factor by the same invertible map as above gives the product form of the original generator. It remains to exclude singular $Q(0)$.

\item Suppose $Q(0)$ is singular. We construct
arbitrarily small constant symmetric matrices $E$
for which $Q(z)+E$ has a rank-one value at a simple
zero of its determinant near the origin. Choose
$z_0\to0$ with $Q'(z_0)\ne0$; such points exist
because the nonconstant determinant makes $Q$
nonconstant. Choose a rank-one symmetric matrix
$R$ close to $Q(0)$ with
\[
                  \operatorname{tr}
                      (\operatorname{adj}(R)Q'(z_0))\ne0.
\]
This choice is possible also when $Q(0)=0$: for a
rank-one symmetric matrix with kernel spanned by
$v$, its adjugate is a nonzero scalar multiple of
$vv^{\mathsf t}$, and the polynomial
$v^{\mathsf t}Q'(z_0)v$ is not identically zero.
Thus the forbidden kernels form a proper finite
set in the projective line, and the rank-one
matrix can be chosen as close to $Q(0)$ as desired.
Set $E=R-Q(z_0)$. Then $E\to0$, and
\begin{equation}
 \det(Q(z_0)+E)=0,\qquad
 \left.\frac{d}{dz}\det(Q(z)+E)\right|_{z=z_0}\ne0.
 \hypertarget{tag-015E}{}\label{cdp:eq:C.UNNODE.2}
\end{equation}

\item Use the full-root perturbation
\[
                  f_E=f_0+(x,y)E(x,y)^{\mathsf t}.
\]
At $(0,0,z_0)$ the rank-one parameter Morse map gives $f_E=X^2+y^2h(y,z)$. The factor $y^2$ follows from vanishing of the function and its normal first derivative along the unchanged axis. Equation~\eqref{cdp:eq:C.UNNODE.2} says that $h(0,z_0)=0$ and $h_z(0,z_0)\ne0$, up to the nonzero coefficient of the split quadratic direction. Thus the invertible coordinate $Z=h(y,z)$ gives
\begin{equation}
                           f_E=X^2+y^2Z.
\hypertarget{tag-015F}{}\label{cdp:eq:C.UNNODE.3}
\end{equation}
The substitution $Z\mapsto-Z$ identifies this with the $(a,b,c)=(2,2,1)$ binomial of \cref{cdp:st:A.5.15}. Its compact spherical class has square $-2(2/\gcd(2,1))(1/\gcd(2,1))=-4$. This class is represented by a continuous map from $S^2$; embeddedness is not required.

\item Choose a small closed ball transverse to $\Gamma$ and to the smooth zero-fibre stratum. Its critical values at $E=0$ are all zero. Along the compact boundary link, $Q$ is invertible, and the parameter Morse maps give fixed product-node charts for all small $E$, retaining the curve coordinate. A real direction in that coordinate is tangent to every function level and transverse to the fixed boundary. On the remaining boundary, submersivity of the original zero level persists for small parameters and small target values. Thus these maps verify all boundary hypotheses of \cref{lem:II:boundary-deformation-transport} for $g_E=f_E$ and the established full-map equality $F=f_0^M$.

That lemma transports the compact class first through regular local values, then through a target path avoiding the finitely many critical values, and then along the parameter family at its common outer regular value $b$. Its final inclusion is into the original fibre $F=b^M$. It preserves the intersection with the compact parallel representative, so the transported sphere has square $-4$. In a complex two-torus its sphere map lifts to the contractible universal cover and has zero homology class, a contradiction. Hence $Q(0)$ is invertible and the product coordinates above apply. The root and final unit-coordinate maps also establish the stated conclusion for a unit times a common positive power.
\end{enumerate}
\end{proof}

\begin{proposition}[Topology and coefficients in the contracted-section case]\hypertarget{tag-01BQ}{}\label{cdp:st:C.2.CR.1}
\StatementRecord{cdp:st:C.2.CR.1}{proved}{01BQ}{\cref{lem:II:contracted-section-duality,lem:II:irrational-ruled-boundary,lem:II:canonical-boundary-blowup,lem:II:proper-conductor-pushout,lem:II:curve-normalization-graph,night:P4R.1,lem:II:proper-conductor-divisor,cdp:st:D.1.9,cdp:st:C.2.VII.2}}

Let $Y$ be a smooth compact complex threefold with
$b_1(Y)=b_2(Y)=0$, and let $f:Y\to\PP^1$ be a proper
holomorphic surjection with connected fibres and smooth
general fibre a complex two-torus. Let $mS$
be a full fibre with irreducible nonnormal reduced support, and
write $h=\eta\pi:Z\to T\to S$ for its normalization and minimal
normalization resolution. Suppose that $p:Z\to B$ is an
irrational ruling, that the effective divisor in \cref{cdp:st:A.5.3} is
$D=2C+V$, where $C$ is a section and $V$ is vertical, and
that $\pi$ contracts $C$. Put $g=g(B)$, let $\Gamma$ be the
conductor-image scheme, and let $n$ be the number of irreducible
components of $\Gamma_{\rm red}$.

One has
\begin{equation}
\begin{gathered}
 H^1(T;\mathbb R)=H^1(T,\OO_T)=H^2(T,\OO_T)=0,
 \qquad \operatorname{length}R^1\pi_*\OO_Z=g,\\
 \pi_1(S)\simeq\pi_1(\Gamma_{\rm red})\simeq\mathbb Z,
 \qquad e(\Gamma_{\rm red})=n,\\
 \chi(\OO_\Gamma)=-1,\qquad
 \chi(\OO_{\Gamma_{\rm red}})\in\{-1,0\},\qquad
 \sum_{x\in\Gamma_{\rm red}}\bigl(\delta_x-r_x+1\bigr)\le1.
\end{gathered}
\hypertarget{tag-00CV}{}\label{cdp:eq:C.CR.1}
\end{equation}
Here $r_x$ is the number of normalization branches at $x$.
The full inverse image $D_{\rm red}$ is a smooth genus-$g$
section with trees of smooth rational curves attached; all
its intersections are ordinary double points. Every component of
coefficient one is a leaf of this tree. For adjacent
components of coefficients $a,b$ one has
\begin{equation}
                         \gcd(a-1,b-1)=1.          \hypertarget{tag-00D6}{}\label{cdp:eq:C.CR.2}
\end{equation}
Choose a sequence of ruling blowdowns to a bundle
$Z_0\to B$. Write its pushed-forward divisor as
$2C_0+\sum_{j=1}^{v}a_jF_j$, with distinct fibres $F_j$
and $a_j>0$, and let $N_0$ count the reverse blowups whose
centre has multiplicity one in the current divisor. Then
\begin{equation}
                     e(S)=2-2g-v+N_0+n.           \hypertarget{tag-00DA}{}\label{cdp:eq:C.CR.3}
\end{equation}
\end{proposition}
\begin{proof}
\begin{enumerate}[label=\arabic*.,leftmargin=*]
\item Restrict $f$ over a small disc about $mS$ and apply
\cref{lem:II:contracted-section-duality} with the same
$S,T,Z,C,V$. It gives the three cohomology vanishings in
\eqref{cdp:eq:C.CR.1}, and $R^1\pi_*\OO_Z$ has length $g$
entirely at $p_0=\pi(C)$. Its Leray sequence is
\eqref{cdp:eq:C.CR.4}. The same numerical input
\cref{lem:II:torus-tube-numerical} gives $\omega_S$
torsion and $\chi(\OO_S)=0$. Its adjunction identity
$K_Z+D=h^*\omega_S$ makes this canonical-boundary
class numerically trivial.

\item Let $\beta:Z\to Z_0$ be the chosen ruling contraction,
and put $D_0=\beta_*D$. Apply the doubled-section case of
\cref{lem:II:irrational-ruled-boundary} with its $T=Z$,
$T_0=Z_0$, $H=C_0$ and this same divisor $D$.
It gives $D_0=2C_0+\sum_j a_jF_j$, the stated tree of
smooth curves, and the isomorphism induced by
$j:D_{\rm red}\hookrightarrow Z$ on fundamental groups.
Under the same blowdowns, this isomorphism is identified
with $\pi_1(C_0)\to\pi_1(Z_0)$ and with the inverse of
the ruling map to $\pi_1(B)$.

At each step $\sigma:W\to W'$, use
\cref{lem:II:canonical-boundary-blowup} with
$Z_1=W$, $Z_2=W'$, $D_1=D_W$, $D_2=\sigma_*D_W$ and
$m=d$, the multiplicity at its centre.
The coefficient $a$ is that of the same exceptional
curve in $D_W$. Formula~\eqref{cdp:eq:C.CR.5} gives
$a=d-1$.

\item Initially each edge of the boundary tree has
coefficients $(2,a_j)$, whose shifted coefficients are
coprime. A node blowup inserts coefficient $a+b-1$,
whose shift is $(a-1)+(b-1)$, so the two new edge gcds
equal the old one. A positive-coefficient blowup at a
smooth point of a coefficient-$a$ component creates
coefficient $a-1$; the new edge has shifted coefficients
$a-1,a-2$. This proves \eqref{cdp:eq:C.CR.2}.

Every coefficient-one component is initially a fibre
leaf. A smooth boundary blowup creates coefficient one
only from coefficient two, and creates a leaf. A node
blowup could create coefficient one only from an edge
$(1,1)$, excluded by \eqref{cdp:eq:C.CR.2}. A blowup
at the unique node of an existing coefficient-one
component subdivides its edge, while a blowup at a
smooth point of that component has exceptional coefficient
zero. Hence every coefficient-one component remains a leaf.

\item Apply \cref{lem:II:proper-conductor-divisor} to
$f$ restricted over the same disc, with
$N=T$, its resolution surface equal to $Z$,
$\nu=\eta$ and $\rho=h$. Its divisor is the same $D$:
both are defined by the orders of the pullback of a local
generator of $\omega_S$, and changing that generator by
a unit changes no order. The lemma identifies the
reduced singular locus of $S$ with its reduced conductor image,
gives
\begin{equation}
 D_{\rm red}=h^{-1}(\Gamma_{\rm red}).
\hypertarget{tag-00DD}{}\label{cdp:eq:C.CR.6}
\end{equation}
It also makes $h$ an isomorphism off these closed sets.
In particular every exceptional configuration meets
the strict conductor and lies in $D_{\rm red}$.

\item Apply \cref{lem:II:proper-conductor-pushout} to
$h:Z\to S$, using $D_{\rm red}\subset Z$ and
$\Gamma_{\rm red}\subset S$ as its closed sets $A,C$,
respectively, with $j$ the inclusion and $k=h|_{D_{\rm red}}$.
The map is proper and surjective, the pairs have compatible
triangulations, and \eqref{cdp:eq:C.CR.6} gives its required
complement isomorphism. Thus
\begin{equation}
 \pi_1(S)=\pi_1(Z)*_{\pi_1(D_{\rm red})}
                         \pi_1(\Gamma_{\rm red})
          \simeq\pi_1(\Gamma_{\rm red}).
\hypertarget{tag-00DE}{}\label{cdp:eq:C.CR.7}
\end{equation}
Here $j_*$ is the isomorphism in the second step. The map
$k_*$ is zero, since the section $C\subset Z$ generates
$\pi_1(D_{\rm red})$ and $h(C)$ is a point.
The resulting isomorphism
$\pi_1(\Gamma_{\rm red})\to\pi_1(S)$ is induced by inclusion.
Every image-conductor component is dominated by a vertical
rational curve; its normalization is therefore $\PP^1$.
\cref{lem:II:curve-normalization-graph} makes
$\pi_1(\Gamma_{\rm red})$ a free group of rank
$b_1(\Gamma_{\rm red})$.

\item Apply \cref{night:P4R.1} to the same full fibre
$mS$ over the disc. It makes $\pi_1(S)$ virtually abelian.
The rational-component case of
\cref{lem:II:curve-normalization-graph} then bounds the
free rank in \eqref{cdp:eq:C.CR.7} by one.
The nonzero integral class on $S$ in \cref{cdp:st:D.1.9},
under the global hypotheses of the statement, excludes
rank zero. The same curve lemma gives
\begin{equation}
 b_1(\Gamma_{\rm red})
   =\sum_x(r_x-1)-n+1=1,\qquad
 e(\Gamma_{\rm red})=2n-\sum_x(r_x-1)=n.
\hypertarget{tag-00DF}{}\label{cdp:eq:C.CR.8}
\end{equation}

\item Put $\epsilon=h^0(S,\omega_S)\in\{0,1\}$ and
$Q=\eta_*\OO_T/\OO_S$. Serre duality and
$\chi(\OO_S)=0$ give $h^1(\OO_S)=1+\epsilon$ and
$h^2(\OO_S)=\epsilon$. Since $S,T$ have the same
constant functions and $H^1(T,\OO_T)=H^2(T,\OO_T)=0$,
the normalization sequence identifies
$H^0(Q)=H^1(\OO_S)$ and $H^1(Q)=H^2(\OO_S)$.
The conductor-duality map of \cref{lem:II:conductor-duality}
is $Q=\omega_\Gamma\otimes\omega_S^{-1}|_\Gamma$.
The torsion twist has degree zero on each component,
so curve duality gives $\chi(\OO_\Gamma)=-\chi(Q)=-1$.
Consequently
\begin{equation}
 h^0(\eta_*\OO_T/\OO_S)=1+\epsilon,\qquad
 h^1(\eta_*\OO_T/\OO_S)=\epsilon,
 \qquad\chi(\OO_\Gamma)=-1.
\hypertarget{tag-00DG}{}\label{cdp:eq:C.CR.9}
\end{equation}
Apply the first assertion of \cref{cdp:st:C.2.VII.2}
to this compact irreducible hypersurface $S$, with its
torsion $\omega_S$ and $\chi(\OO_S)=0$. It gives
$\chi(\OO_{\Gamma_{\rm red}})\ge-1$.
All normalized image components are rational, so
$\chi(\OO_{\Gamma_{\rm red}})=n-\sum_x\delta_x$.
Equation~\eqref{cdp:eq:C.CR.8} and
$\delta_x\ge r_x-1$ give
$\chi(\OO_{\Gamma_{\rm red}})\le0$ and
$\sum_x(\delta_x-r_x+1)=-\chi(\OO_{\Gamma_{\rm red}})\le1$.

\item On the minimal bundle,
$e(Z_0)=4-4g$ and $e((D_0)_{\rm red})=2-2g+v$.
The complement formula of
\cref{lem:II:canonical-boundary-blowup}, applied to the
same successive divisors as in the second step, adds one
exactly at the $N_0$ centres of multiplicity one. Hence
\[
             e(Z\setminus D_{\rm red})=2-2g-v+N_0.
\]
The Euler formula in \cref{lem:II:proper-conductor-pushout},
for the same $h:Z\to S$, $D_{\rm red}$ and $\Gamma_{\rm red}$, and
$e(\Gamma_{\rm red})=n$ give \eqref{cdp:eq:C.CR.3}.
\end{enumerate}
\end{proof}

\begin{proposition}[The genus-two node alternative is excluded]\hypertarget{tag-01BS}{}\label{cdp:st:C.2.CR.3}
\StatementRecord{cdp:st:C.2.CR.3}{proved}{01BS}{\cref{cdp:st:C.2.CR.1,lem:II:contracted-section-duality,lem:II:conductor-two-branch-cycle}}

Under all the hypotheses of \cref{cdp:st:C.2.CR.1}, with $g=2$, the component
of $\Gamma$ through $q=h(C)$ cannot have generic
transversal an ordinary node. Consequently the
remaining genus-two alternative has an ordinary cusp
at its generic point, and exactly one strict conductor
curve of coefficient two and degree one lies above
that component.
\end{proposition}
\begin{proof}
\begin{enumerate}[label=\arabic*.,leftmargin=*]
\item Let $p=\pi(C)$, $q=\eta(p)$, and let $E_0$ be
the connected exceptional fibre containing $C$.
In the boundary tree of \cref{cdp:st:C.2.CR.1}, the
first nonexceptional vertex on every path from $C$
meets $E_0$. Its image in $T$ passes through $p$,
and its further image is a component of $\Gamma$ through
$q$. Apply \cref{lem:II:contracted-section-duality}
to these same data. Since $g=2$, $\Gamma$ is smooth
and reduced at $q$, so there is only one such component,
denoted $G$.

\item Suppose its generic transversal is a node.
Every strict conductor curve above $G$ has coefficient
one, and \cref{cdp:st:C.2.CR.1} makes it a leaf.
Thus each first nonexceptional vertex just found is
a leaf, and every strict conductor curve is such a leaf,
attached at one point to $E_0$. No additional exceptional
tail meets it, since this would raise its boundary
valence. The reduced conductor $D_T\subset T$ is
therefore a union of rational curves through $p$, each
smooth away from $p$ and with one branch at $p$.
Indeed $\pi$ is an isomorphism on each leaf away from
its single attachment point, and contracts the connected
$E_0$ to $p$. Its normalization-incidence graph is a star,
with cycle rank zero.

\item All these image curves pass through $q$, where
$\Gamma$ is smooth. Hence $\Gamma_{\rm red}=G$ is
irreducible. By \eqref{cdp:eq:C.CR.8},
$b_1(G)=1$ and $\sum_x(r_x-1)=1$.
Thus every multibranch point of $G$ has exactly two
branches, and lies away from the smooth point $q$.
For any such point apply
\cref{lem:II:conductor-two-branch-cycle} with
$N=T$, $\nu=\eta$, its curve $C=G$ and its inverse
conductor $D=D_T=(\pi D_{\rm red})_{\rm red}$.
Every component of $D_T$ dominates $G$ under the finite
map $\eta$, and no other image-conductor component
passes through the chosen point. The lemma gives
$\beta(D_T)\ge b_0(D_T)=1$, contradicting the star just
described. Therefore all germs of $G$ would be unibranch.

\item The normalization of $G$ is rational. If all its
germs are unibranch, its normalization is a continuous
bijection from the compact $\PP^1$ to the Hausdorff $G$,
hence a homeomorphism. This gives $\pi_1(G)=1$,
contrary to \eqref{cdp:eq:C.CR.1}.
The remaining generic transversal is a cusp.
Its transverse normalization is unibranch, with
conductor exponent two, so there is one strict
coefficient-two curve, of degree one over $G$.
\end{enumerate}
\end{proof}

\begin{proposition}[Contracted irrational sections and their Euler number]\hypertarget{tag-01CN}{}\label{cdp:st:C.2.UN.10}
\StatementRecord{cdp:st:C.2.UN.10}{proved}{01CN}{\cref{lem:II:contracted-section-duality,cdp:st:C.2.UN.9,cdp:st:C.2.CR.4,cdp:st:C.2.CR.1,cdp:st:C.2.MH.2}}
\hypertarget{tag-01BU}{}\label{cdp:st:C.2.CR.5}
\StatementRecord{cdp:st:C.2.CR.5}{proved}{01BU}{}
\hypertarget{tag-01C6}{}\label{cdp:st:C.2.MH.3}
\StatementRecord{cdp:st:C.2.MH.3}{proved}{01C6}{}
Let $f:Y\to B_0$ be a proper holomorphic map from a
smooth complex threefold to a smooth curve, with connected
fibres and smooth general fibre a complex two-torus.
Let $mS$ be an irreducible nonnormal full fibre, and write
$h=\eta\pi:Z\to T\to S$ for the normalization and its
minimal resolution. Suppose that $p:Z\to B$ is a ruling
over a curve of positive genus, that its effective
adjunction divisor is $D=2C+V$, with $C$ a section and
$V$ vertical, and that $\pi$ contracts $C$. Then
\begin{equation}
                              g(B)=1.              \hypertarget{tag-00D5}{}\label{cdp:eq:C.CR.19}
\end{equation}
For a compact $Y$ with $b_1(Y)=b_2(Y)=0$ and
$B_0=\PP^1$, let $n,v,N_0$ have the meanings in
\cref{cdp:st:C.2.CR.1}. Then
\begin{equation}
\hypertarget{tag-00QG}{}\SsixEquationAliases{\label{cdp:eq:C.MH.15}}
             e(S)=n-v+N_0=0,\qquad n=v-N_0.
\hypertarget{tag-00D7}{}\label{cdp:eq:C.CR.20}
\end{equation}
\end{proposition}
\begin{proof}
\begin{enumerate}[label=\arabic*.,leftmargin=*]
\item Put $g=g(B)$, $p_0=\pi(C)$ and $q=\eta(p_0)$.
Apply \cref{lem:II:contracted-section-duality} to the
map and full fibre in the statement, with the same
$S,T,Z,C,V$. It gives $g\in\{1,2\}$ and the residue
comparison \eqref{cdp:eq:C.UNCR.3}. For $g=2$, the
conductor ideal at $q$ is $(x,y)$, the point $p_0$
is the only normalization point over $q$, and the
smooth reduced conductor has generic transversal
a node or a cusp.

\item Consider a reduced local equation of $S$ at $q$.
Each positive-dimensional component of its critical
locus through $q$ has constant value zero, so lies in
its singular locus. A normal surface is regular in
codimension one, so every singular curve lies in the
nonnormal locus and hence in the conductor. After shrinking the germ,
the full reduced critical locus is the smooth conductor
curve above. In the node case
\cref{cdp:st:C.2.UN.9} gives the exact product
coordinates. Its normalization has two smooth points
over $q$, contradicting the unique $p_0$.
In the cusp case \cref{cdp:st:C.2.CR.4} gives the
product cusp $X^2+Y^3=0$ with normalization
\[
          (u,z)\longmapsto(X,Y,z)=(u^3,-u^2,z).
\]
This normalization is smooth. Its minimal resolution
is the identity near $p_0$, so cannot contract the
positive-genus section $C$ there. Both alternatives
exclude $g=2$, giving \eqref{cdp:eq:C.CR.19}.

\item Under the additional global hypotheses in the
statement, \cref{cdp:st:C.2.MH.2} applies to the
same irreducible full fibre and ruled normalization
resolution. It gives $e(S)=2-2g=0$.
Equation~\eqref{cdp:eq:C.CR.3}, for the same chosen
ruling blowdowns, gives $e(S)=n-v+N_0$ and hence
$n=v-N_0$. These are \eqref{cdp:eq:C.CR.20}.
\end{enumerate}
\end{proof}

\begin{proposition}[The elliptic contraction has a planar conductor and only double transverse branches at its image]\hypertarget{tag-01BV}{}\label{cdp:st:C.2.CR.6}
\StatementRecord{cdp:st:C.2.CR.6}{proved}{01BV}{\cref{cdp:st:C.2.CR.1,lem:II:contracted-section-duality,lem:II:conductor-duality,lem:II:plane-branch-conductor}; \cite[\S2.2.3, p.~60]{vanStraten1987}}

In the case $g=1$ of \cref{cdp:st:C.2.CR.1}, let
$p=\pi(C)$ and $q=\eta(p)$. The conductor-image
scheme is a local complete intersection at $q$.
More precisely, ambient coordinates can be chosen
so that its ideal is
\begin{equation}
                         (x,g(y,z)).               \hypertarget{tag-00D8}{}\label{cdp:eq:C.CR.21}
\end{equation}
Every global conductor component passing through
$q$ has generic transverse type $A_k$ for some
$k\ge1$. For type $A_{2s-1}$, its two transverse
normalization branches have conductor coefficient
$s$ each; their global covering degrees sum to two.
For type $A_{2s}$ there is one coefficient-$2s$
normalization branch, of covering degree one.
\end{proposition}
\begin{proof}
\begin{enumerate}[label=\arabic*.,leftmargin=*]
\item Apply \cref{lem:II:contracted-section-duality}
with $g=1$ and these same $S,T,Z,C,V,p,q$.
Equations~\eqref{cdp:eq:C.CR.12} and
\eqref{cdp:eq:C.CR.13} give an element of the
conductor summand $I_p$ with nonzero ambient linear
part. Call it $x$ and extend it to coordinates
$(x,y,z)$ on the smooth threefold.

\item The full conductor ideal contains $x$, and
\cref{lem:II:conductor-duality} makes its quotient
pure one-dimensional Cohen--Macaulay. Its remaining
ideal in the regular ring $\C\{y,z\}$ is therefore
unmixed of height one. Its primary decomposition
is the intersection of its height-one primary
localizations, which are powers of the corresponding
principal primes in this factorial ring. Their
intersection is generated by their product.
This proves \eqref{cdp:eq:C.CR.21}, with all
multiplicities in $g$.

\item Let $D\subset W$ be a reduced plane curve germ,
let $\nu:\widetilde D\to D$ be its normalization,
and let $\rho:\widetilde W\to W$ be an embedded
resolution with smooth disjoint strict transform
$\widehat D$. Apply the finite-duality comparison
of \cref{lem:II:plane-branch-conductor} to its
equation $f_D$. It identifies
\[
           \mathfrak c_D\,\omega_D=\nu_*\omega_{\widetilde D}
\]
inside the meromorphic differentials on the normalized
branches. The residue sequence on $\widetilde W$ and
\cite[\S2.2.3, p.~60]{vanStraten1987}, applied to this
proper bimeromorphic $\rho$ with smooth source, give
a surjection
$\rho_*\omega_{\widetilde W}(\widehat D)\to
\nu_*\omega_{\widetilde D}$.

Thus for an ambient representative $a$ of a conductor
element, the residue of $a\,du\wedge dv/f_D$ lifts
to a meromorphic two-form on $\widetilde W$ with poles
only along $\widehat D$. Downstairs this lift and
$a\,du\wedge dv/f_D$ have the same residue, so differ
by a holomorphic two-form. Pullback of that difference
is holomorphic. The pullback of the displayed form
therefore has no exceptional pole. On the first
plane blowup, its order along the exceptional curve is
$\operatorname{ord}_0(a)+1-m$, where
$m=\operatorname{mult}_0D$. Hence
\begin{equation}
 \mathfrak c_D\subset\mathfrak m_D^{\,m-1}
 \quad\hbox{in the sense that every ambient
 representative has order at least }m-1.
 \hypertarget{tag-00D9}{}\label{cdp:eq:C.CR.22}
\end{equation}
Adding a multiple of $f_D$ preserves this bound
because $\operatorname{ord}_0(f_D)=m$.

\item At a general point of any conductor component
through $q$, the coordinate $x$ remains in the
conductor and has nonzero transverse linear part.
Indeed its smooth ambient zero plane contains the
curve. A transverse slice and \eqref{cdp:eq:C.CR.22}
give $m\le2$; the slice is singular, so $m=2$.
Weierstrass preparation gives a monic quadratic.
Complete its square and write its nonzero remaining
one-variable term as a unit times $v^{k+1}$.
Taking a $(k+1)$st root of that unit in the $v$
coordinate gives $u^2-v^{k+1}=0$. Each coordinate
change has nonzero derivative and inverse given by
the implicit function theorem.

\item For $k=2s-1$ the branches are $u=\pm v^s$.
They are smooth and have intersection multiplicity
$s$, so \cref{lem:II:plane-branch-conductor} gives
conductor $(v^s)$ on each. For $k=2s$, the
normalization is $v=t^2$, $u=t^{2s+1}$. Its value
semigroup $\langle2,2s+1\rangle$ contains every
integer at least $2s$ and omits $2s-1$, so its
conductor is $(t^{2s})$.
The generic transverse branch count of the finite
normalization equals the sum of its covering
degrees over the image component. These counts are
two and one in the respective cases.
\end{enumerate}
\end{proof}

\section{Horizontal conductors and elliptic interfaces}\label{sec:II:horizontal-conductor-interfaces}

\subsection{Reduced horizontal boundaries}\label{subsec:II:reduced-horizontal-boundaries}

\begin{proposition}[The reduced horizontal case is a minimal elliptic ruled double-curve fibre]\hypertarget{tag-01C3}{}\label{cdp:st:C.2.IRR.1}
\StatementRecord{cdp:st:C.2.IRR.1}{proved}{01C3}{\cref{cdp:st:A.5.3}, \cref{cdp:st:A.5.17}, \cref{lem:II:plane-branch-conductor}, \cref{cdp:st:A.5.15}}

Let $f:Y\to C_0$ be a proper holomorphic map from a smooth compact
complex threefold with $H^2(Y,\mathbb Q)=0$, with connected fibres
and smooth general fibre a complex two-torus. Let $mS$ be a full
fibre with irreducible nonnormal reduced support. Let
$h=\eta\pi:Z\to T\to S$ be the minimal resolution of its
normalization, and suppose that $p:Z\to B$ is an irrational
ruling. Write the effective divisor of \cref{cdp:st:A.5.3} as
$D_Z=H+V$, where $V$ is vertical and the horizontal part $H$
is reduced, of total degree two over $B$.

Then $B$ is elliptic, $T=Z$ is a smooth minimal ruled surface
over $B$, and $D_Z=H$ is a disjoint union of smooth elliptic
curves. The reduced conductor image is a disjoint union of
smooth elliptic curves. Its inverse image $H$ maps onto it
by a finite etale covering of total degree two over each
component. The support $S$ has only ordinary double-curve
normal crossings and
\begin{equation}
e(S)=0.\hypertarget{tag-00OH}{}\label{cdp:eq:C.IR.1}
\end{equation}
\end{proposition}

\begin{proof}

\begin{enumerate}[label=\arabic*.,leftmargin=*]
\item The canonical comparison of
\cref{cdp:st:A.5.3} gives
\begin{equation}
K_Z+D_Z=h^*\omega_S,\qquad c_1(K_Z+D_Z)=0
 \quad\hbox{in }H^2(Z,\mathbb Q).\hypertarget{tag-00OI}{}\label{cdp:eq:C.IR.2}
\end{equation}
If $g(B)\ge2$, that proposition makes the horizontal part
twice a section.  Since the given horizontal part is reduced,
$g(B)=1$.  Contract the vertical ruling blowups to a minimal
ruled surface $\rho:Z\to Z_0$, and write
$D_0=\rho_*D_Z=H_0+V_0$.  Choose a section $C_0$ with
$C_0^2=-n$ and a ruling fibre $F$.  Numerically,
\[
 K_{Z_0}=-2C_0-nF,\qquad H_0=2C_0+bF,\qquad V_0=vF,
 \qquad b+v=n.
\]
Here $v\ge0$.  Adjunction gives $p_a(H_0)=1-v$.
An irreducible $H_0$ has normalization of genus at least one
by Riemann--Hurwitz.  If $H_0$ has two components, they are
sections, and its arithmetic genus is one plus their
intersection number.  Hence $v=0$.  In the irreducible case
$H_0$ is smooth elliptic; in the other case its two smooth
elliptic sections are disjoint.

\item At each ruling blowdown $\sigma:W\to W'$, apply
\cref{lem:II:canonical-boundary-blowup} with
$Z_1=W$, $Z_2=W'$, $D_1=D_W$, $D_2=\sigma_*D_W$,
and $m=d=\operatorname{mult}_{\sigma(E)}D_2$.
Its coefficient $a$ is the coefficient of this same
exceptional curve in $D_W$.
Identity \eqref{cdp:eq:C.IR.4} gives $a=d-1$
and zero rational canonical-boundary class downstairs.
This establishes the numerical class used in the preceding
step through the ruling map $\rho$.

In the reverse sequence, effectivity forces every centre
to lie on the divisor at that step: a centre off it would create
coefficient $-1$, which persists on its strict transform.
Starting with the disjoint smooth coefficient-one curves
$H_0$, each permitted centre therefore has $d=1$.  Its
exceptional coefficient is zero, and its new boundary is
the strict transform of $H_0$.  Induction gives
\begin{equation}
D_Z=H,\qquad V=0.\hypertarget{tag-00OK}{}\label{cdp:eq:C.IR.5}
\end{equation}
In particular $H$ is a disjoint union of smooth elliptic
curves.  Coefficient-zero ruling exceptional curves have
not yet been discarded.

\item Let $E$ be a curve exceptional for
the minimal normalization resolution $\pi$ and not contained
in $H$.  Its negative self-intersection and adjunction give
\[
 2p_a(E)-2=E^2-H\cdot E<0.
\]
Thus $p_a(E)=0$, and $E$ is smooth rational.  Minimality
excludes $E^2=-1$.  Since $H\cdot E=2+E^2\ge0$, it follows
that $E^2=-2$ and $H\cdot E=0$.  A positive-coefficient
exceptional curve is one of the mutually disjoint elliptic
components of $H$.  Consequently every connected exceptional
fibre is disjoint from the strict conductor curves.

The normalization conductor ideal is the rank-one reflexive
sheaf $\omega_T\otimes\eta^*\omega_S^{-1}$ embedded in
$\mathcal O_T$ by finite duality.  Any point on its divisorial
support lies on a conductor curve, whose strict transform
must meet the resolution fibre over that point.  The fibres
just described therefore lie outside that support.  There
the conductor ideal equals $\mathcal O_T$, by reflexivity,
and $\eta$ is an isomorphism.  A nontrivial such resolution
fibre would give an isolated normal hypersurface singularity
of $S$.  \cref{cdp:st:A.5.17} excludes it in the
torus smoothing with its full equation, unit and multiplicity.  Hence $\pi$ is the identity and $T=Z$ is smooth.

\item The conductor ideal of $T$ now
defines the reduced smooth divisor $H$.  Its contraction to
$\mathcal O_S$ is radical, since the quotient embeds into
$\mathcal O_H$.  Let $G$ be this reduced conductor image.
At a generic transverse plane curve, conductor exponent one
on a normalization branch forces that branch to be smooth
and to meet exactly one other smooth branch transversely:
the branchwise conductor formula is twice the intrinsic
delta invariant plus the sum of its intersection
multiplicities with the other branches.  This is the
formula of \cref{lem:II:plane-branch-conductor}.  Thus $H\to G$ has total
generic degree two over each component.

The sheet involution extends to the smooth compact curves
$H$ over the normalization of $G$; a degree-two map of
smooth curves has its holomorphic deck involution, including
the exchange of two degree-one components.  The involution
preserves the finite fibre over every point of $G$, since
its composites to $G$ agree on a dense open subset and
therefore everywhere.

\item We prove the local ring comparison
at every point of $G$.  Let $A$ be the reduced hypersurface
ring, let $B$ be its finite normalization ring, and let $I$
be the conductor.  The ring $B$ is a product of regular
two-dimensional local rings.  Define
\[
 A'=\{b\in B:b\bmod I\in(B/I)^\iota\},
\]
where $\iota$ is the involution of the preceding
step.  There is an inclusion $A\subset A'\subset B$.
The two first rings agree in codimension one by the
ordinary-node calculation, so $A'/A$ has finite length.
The module $A'$ is finite and torsion-free over $A$.
Local cohomology of $0\to A\to A'\to A'/A\to0$ injects
$A'/A$ into $H^1_{\mathfrak m}(A)$.  The latter vanishes
because a hypersurface surface ring has depth two.
It follows that $A'=A$, and therefore
$A/I=(B/I)^\iota$.

The image curve $G$ is consequently the quotient of a
smooth curve by its involution and is smooth.  The exact
unramified local pullback ring is
\[
 \{(a(x,z),b(y,z)):a(0,z)=b(0,z)\}
       \simeq\mathbb C\{x,y,z\}/(xy).
\]
At a ramification point it is
\begin{equation}
\{a(u,t):a(0,t)\in\mathbb C\{t^2\}\}
 =\mathbb C\{u,ut,t^2\}
 \simeq\mathbb C\{x,y,z\}/(x^2-y^2z),\hypertarget{tag-00OL}{}\label{cdp:eq:C.IR.6}
\end{equation}
with $(x,y,z)=(ut,u,t^2)$.  \cref{cdp:st:A.5.15}
excludes the latter local map, including a unit times the
$m$th power.  Therefore all conductor maps are unramified,
their image components are disjoint smooth elliptic curves,
and $S$ has ordinary double-curve normal crossings.

\item Apply the normal-crossings
comparison \cref{cdp:st:C.2.NC.1} to the full fibre
$mS$.  It gives
\begin{equation}
0=e(Y_t)=m\,e(T\setminus H)=m\,e(T),\hypertarget{tag-00OM}{}\label{cdp:eq:C.IR.9}
\end{equation}
because $e(H)=0$.  A minimal ruled surface over the elliptic
curve $B$ has Euler number zero, while every point blowup
adds one.  Thus $\rho$ has no blowups and $T=Z$ is minimal.
Finally the normalization is an isomorphism off
$H$ and $G$, so $e(S)=e(T)-e(H)+e(G)=0$.

\end{enumerate}
\end{proof}


\subsection{Rational boundaries and attachment maps}\label{subsec:II:rational-boundary-attachments}

\begin{proposition}[Positive Euler number for a rational normalization]
\hypertarget{tag-P360}{}\label{prop:II:rational-torus-fibre-euler}
\hypertarget{tag-01CG}{}\label{cdp:st:C.2.RAT.4}
\StatementRecord{cdp:st:C.2.RAT.4}{proved}{01CG}{}
\hypertarget{tag-01CH}{}\label{cdp:st:C.2.RAT.5}
\StatementRecord{cdp:st:C.2.RAT.5}{proved}{01CH}{}
\hypertarget{tag-01CJ}{}\label{cdp:st:C.2.RAT.7}
\StatementRecord{cdp:st:C.2.RAT.7}{proved}{01CJ}{}
\hypertarget{tag-01CK}{}\label{cdp:st:C.2.RAT.8}
\StatementRecord{cdp:st:C.2.RAT.8}{proved}{01CK}{}
\hypertarget{tag-01CL}{}\label{cdp:st:C.2.RAT.9}
\StatementRecord{cdp:st:C.2.RAT.9}{proved}{01CL}{}
\hypertarget{tag-01BY}{}\label{cdp:st:C.2.D4.3}
\StatementRecord{cdp:st:C.2.D4.3}{proved}{01BY}{}
\StatementRecord{prop:II:rational-torus-fibre-euler}{proved}{P360}{\cref{lem:II:torus-tube-numerical,cdp:st:A.5.16,prop:II:minimal-boundary-complements,night:P4.13,lem:II:proper-conductor-pushout,lem:II:torus-tube-character,cdp:st:C.2.UN.1}}
Let $f:Y\to B$ be a proper holomorphic map from a smooth complex
threefold to a smooth curve, with connected fibres and smooth
general fibre a complex two-torus. Let $mS$ be a full fibre with
irreducible reduced support. Suppose its minimal normalization
resolution $h=\eta\pi:Z\to T\to S$ is rational. Then $e(S)\ge2$.
For $H^1(S,\mathbb Q)\ne0$ one has $e(S)=2$.

For the effective resolution conductor $D$ and the reduced
conductor image $G=\Gamma_{\rm red}$, one has
\begin{equation}
\begin{gathered}
 K_Z+D\simeq\OO_Z,\qquad D_{\rm red}\text{ connected},
 \qquad p_a(D_{\rm red})\le1,\\
 G\text{ connected},\qquad p_a(G)\le2,\qquad e(G)\ge0,\\
 D_{\rm red}=h^{-1}(G),\qquad
 S_{\rm reg}\simeq Z\setminus D_{\rm red}.
\end{gathered}
\hypertarget{tag-014P}{}\label{cdp:eq:C.UI.12}
\end{equation}
\end{proposition}
\begin{proof}
\begin{enumerate}[label=\arabic*.,leftmargin=*]
\item A normal support is smooth torus or bielliptic by
\cref{cdp:st:C.2.UN.5}, and cannot have a rational resolution.
Thus $S$ is nonnormal. Restrict $f$ to a disc about its value.
\cref{lem:II:torus-tube-numerical}, applied to the same full
fibre $mS$ and the same maps $\eta,\pi$, gives
$\omega_S^{\otimes m}\simeq\OO_S$ and $\chi(\OO_S)=0$.
These are the two explicit numerical hypotheses of
\cref{cdp:st:A.5.16}. Its conductor is the one from finite
normalization duality and the minimal resolution, hence is
the present $D$, and its image scheme is $\Gamma$.
That proposition gives every assertion of
\eqref{cdp:eq:C.UI.12}. In particular, its isolated-singularity
exclusion removes all exceptional fibres off the conductor;
on that complement both $\eta$ and $\pi$ are isomorphisms.

\item Apply \cref{night:P4.13} to this restriction $f:U\to\Delta$.
Its irreducible minimum-multiplicity formula gives
$e(S_{\rm reg})\ge0$. Euler additivity and
\eqref{cdp:eq:C.UI.12} therefore give
$e(S)=e(S_{\rm reg})+e(G)\ge0$.

Suppose this sum is zero. Both summands then vanish.
First suppose $D_{\rm red}$ has a component of positive
normalization genus. Its arithmetic-genus bound makes that
component smooth elliptic and unique. Choose the contraction
data of \cref{def:II:minimal-boundary-data}.
\cref{prop:II:minimal-boundary-complements} applies to the same
$Z,D$ through \cref{cdp:st:A.5.16}; its blowup formula
\eqref{cdp:eq:A.RT.12} expresses the complement Euler number as
the minimal complement plus $N_0\ge0$.
The elliptic component survives every rational blowdown.
On $\mathbb P^2$ it is the whole smooth anticanonical cubic,
whose complement has Euler number three. On a Hirzebruch
surface it is a coefficient-one bisection $H_0$:
a section is rational, and total horizontal anticanonical
degree is two. If $V_0$ is its vertical part, adjunction gives
$2p_a(H_0)-2+H_0\cdot V_0=0$. Thus $V_0=0$, and the
complement has Euler number four. These are precisely the
elliptic cases of \eqref{cdp:eq:A.RT.13}--\eqref{cdp:eq:A.RT.16}.
Both contradict $e(S_{\rm reg})=0$.

\item All normalized components of $D_{\rm red}$ are therefore
rational. The coherent and topological normalization maps
give $b_1(D_{\rm red})\le p_a(D_{\rm red})\le1$.
Every component of $G$ is rational-normalized, since it is
the finite image of a strict conductor component. Write
$n\ge1$ for its number of irreducible components.
The topological normalization formula for the connected
curve $G$ gives $e(G)=1+n-b_1(G)$; the supposed equality
therefore gives $b_1(G)=n+1\ge2$.

Apply \cref{lem:II:proper-conductor-pushout} to
$(X,Y,A,C,h)=(Z,S,D_{\rm red},G,\eta\pi)$.
The full inverse-image equality and the complement
isomorphism are \eqref{cdp:eq:C.UI.12}; compatible
triangulations make the compact analytic inclusion
$D_{\rm red}\hookrightarrow Z$ a cofibration.
The lemma identifies $S$ with the double mapping cylinder
of the maps $D_{\rm red}\hookrightarrow Z$ and
$h|_{D_{\rm red}}:D_{\rm red}\to G$.
Its cohomology sequence, connectedness, and
$H^1(Z,\mathbb Q)=0$ give
\begin{equation}
 H^1(S,\mathbb Q)\simeq
 \ker\!\left[H^1(G,\mathbb Q)
        \xrightarrow{(h|_{D_{\rm red}})^*}
             H^1(D_{\rm red},\mathbb Q)\right].
\hypertarget{tag-014Q}{}\label{cdp:eq:C.UI.13}
\end{equation}
This space has dimension at least $n>0$.
A nonzero rational class has a nonzero integral multiple
in $H^1(S,\mathbb Z)$. \cref{lem:II:torus-tube-character}
applied to that class and the same normalization resolution
gives $e(S)=2$, contradicting the supposed zero.
Hence $e(S)>0$. The parity theorem
\cref{cdp:st:C.2.UN.1}, for the same proper map and full fibre,
makes $e(S)$ an even integer, proving $e(S)\ge2$.

\item For any $H^1(S,\mathbb Q)\ne0$, the same integral-class
choice and the rational-resolution clause of
\cref{lem:II:torus-tube-character} give $e(S)=2$ directly.
\item For a compact family, restriction to a disc $\Delta$ about
$b=f(S)$ identifies the full fibre and its normalization-resolution
diagram with those of $f|_{f^{-1}(\Delta)}$. On the three reduced
spaces the maps are
\[
 S\times_B\Delta\longrightarrow S,\qquad
 T\times_B\Delta\longrightarrow T,\qquad
 Z\times_B\Delta\longrightarrow Z,
 \qquad (x,b)\longmapsto x,
\]
with inverses $x\mapsto(x,b)$. They commute with $\eta$ and
$\pi$, and the Cartier divisor remains $f^*(b)=mS$.
Consequently the same bound applies to every rational ruled-boundary
configuration, including a degree-two tacnode conductor or a
degree-two or degree-three ordinary-triple conductor, with all
common cubic zeros and their multiplicities.
\end{enumerate}
\end{proof}


\begin{proposition}[Euler signs for two of the rational boundary classes]\hypertarget{tag-01CD}{}\label{cdp:st:C.2.RAT.1}
\StatementRecord{cdp:st:C.2.RAT.1}{proved}{01CD}{\cref{cdp:st:A.5.16}, \eqref{cdp:eq:C.RA.7}, \eqref{cdp:eq:C.RA.8}, \eqref{cdp:eq:C.RA.9}, \eqref{cdp:eq:C.RA.10}, \cref{cdp:st:A.5.5}, \cref{cdp:st:A.5.15}, \cref{cdp:st:A.5.17}, \eqref{cdp:eq:C.RA.14}, \eqref{cdp:eq:C.RA.15}, \eqref{cdp:eq:C.RA.6}, \eqref{cdp:eq:C.RA.13}, \eqref{cdp:eq:C.RA.2}}

Let $f:Y\to B$ be a proper holomorphic map from a smooth compact
complex threefold with $b_2(Y)=0$, with connected fibres and smooth
general fibre a complex two-torus. Let $mS$ be a full fibre with
irreducible nonnormal reduced support. Write
$h=\eta\pi:Z\to T\to S$ for the minimal resolution of its
normalization and suppose that $Z$ is rational. Use the effective
divisor $K_Z+D=0$ of \cref{cdp:st:A.5.16}.

If a contraction to $\mathbf P^2$ sends $D$ to three concurrent
reduced lines, then $e(S)\ge1$. More precisely,
$e(S)=-1+N_0+e(\Gamma_{\rm red})\ge1+N_0$, with $N_0$ as in
\eqref{cdp:eq:A.RT.12}.

Consider also a fixed morphism $\rho:Z\to\mathbb F_n$,
$n\ge1$, which is a sequence of point blowups in reverse and has
\begin{equation}
\rho_*D=2C_0+\sum_{j=1}^r a_jF_j,\qquad
 C_0^2=-n,\qquad
 a_j\ge1,\quad \sum_j a_j=n+2,\hypertarget{tag-00WQ}{}\label{cdp:eq:C.RA.1}
\end{equation}
where the $F_j$ are distinct ruling fibres. Let $C$ be the
strict transform of $C_0$. Suppose $h(C)$ is a curve and the
generic transverse singularity there is a cusp $A_2$.
Then $r\le3$ and $e(S)\ge1$. If $h(C)$ is smooth, then
$r\le2$. These assertions allow singular normalization,
arbitrary vertical conductor multiplicities and a nonreduced
conductor-image scheme.
\end{proposition}

\begin{proof}

\begin{enumerate}[label=\arabic*.,leftmargin=*]
\item Put $G=\Gamma_{\rm red}$. \cref{cdp:st:A.5.16} has already proved
\begin{equation}
e(S)=e(Z\setminus D_{\rm red})+e(G).\hypertarget{tag-00X1}{}\label{cdp:eq:C.RA.2}
\end{equation}
In either boundary class under consideration, $D_{\rm red}$
is simply connected, all its components are rational, and
\eqref{cdp:eq:A.RT.18}--\eqref{cdp:eq:A.RT.20} give
\begin{equation}
\pi_1(S)=\pi_1(G)\in\{1,\mathbb Z\}.\hypertarget{tag-00XC}{}\label{cdp:eq:C.RA.3}
\end{equation}
Write $n_G$ for the number of irreducible components of $G$
and $\beta=b_1(G)\le1$. The normalization maps from their
rational normalizations are homeomorphisms except for the
identification of finitely many points. Consequently
\begin{equation}
e(G)=n_G+1-\beta\ge n_G\ge1.\hypertarget{tag-00XH}{}\label{cdp:eq:C.RA.4}
\end{equation}
The exact two-branch incidence argument \eqref{cdp:eq:A.RT.21} strengthens
this to $e(G)\ge2$ in the present boundary classes.
For three concurrent lines, \eqref{cdp:eq:A.RT.12}--\eqref{cdp:eq:A.RT.13} therefore
prove the first assertion.

\item We prove the cusp assertion. Put $H=h(C)_{\rm red}$.
The transverse $A_2$ curve has a single normalized branch.
Thus $C\to H$ is the normalization map of degree one, and
no other strict conductor component maps onto $H$.
In particular $H^\nu=C=\mathbf P^1$. The restriction
$H^1(G,\OO_G)\to H^1(H,\OO_H)$ is surjective, since its
kernel has one-dimensional support and no second cohomology.
\cref{cdp:st:A.5.16} therefore gives
\begin{equation}
\delta:=p_a(H)\le2.\hypertarget{tag-00XI}{}\label{cdp:eq:C.RA.5}
\end{equation}
Since $H^\nu$ is rational, $\delta$ is the sum of its local
delta invariants. Write $R$ for the union of the other irreducible components of
$G$. The exact constant-sheaf sequence
\[
 0\longrightarrow\mathbb Q_G\longrightarrow
 \mathbb Q_H\oplus\mathbb Q_R\longrightarrow
 \mathbb Q_{H\cap R}\longrightarrow0
\]
gives a surjection
$H^1(G;\mathbb Q)\to H^1(H;\mathbb Q)\oplus H^1(R;\mathbb Q)$,
since $H\cap R$ is finite. Hence $b_1(H)\le b_1(G)\le1$.
If $s$ denotes the number of points of $C$ above singular
points of $H$, normalization of the reduced curve gives
\begin{equation}
s=\sum_{p\in\operatorname{Sing}H}r_p
   =b_1(H)+\#\operatorname{Sing}H\le1+\delta
 \quad(\delta>0),\qquad s=0\quad(\delta=0).\hypertarget{tag-00XJ}{}\label{cdp:eq:C.RA.6}
\end{equation}

\item We need the precise low-delta curve analysis, including its
two possible non-lci cases. A reduced curve germ of delta
invariant at most two has at most three branches.
Three branches require delta two, and the ring is exactly
the subring of its normalization consisting of triples
with equal constant value. This follows by equality of
colengths, both equal to two. Its topological branch
identification alone contributes two loops to a globally
irreducible curve, so $b_1(H)\le1$ excludes it here.

\item The remaining germs have at most two branches.
Apply \cref{lem:II:delta-two-curve-rings} to the same
reduced germ $H\subset Y$ and its normalization.
The two embedded ideals in \eqref{cdp:eq:A.T.14}
are obtained by lifting the fixed normalization
coordinate functions to the ambient regular local ring;
the inverse function theorem makes that lift a
biholomorphic ambient coordinate change. Thus every
remaining non-lci germ is exactly one of those two ideals.

For the space cusp, use the identical generators
$g_1=y^2-xz$, $g_2=yz-x^3$, $g_3=z^2-x^2y$ and
$F=x^5+xy^3+z^3-3x^2yz$ in
\eqref{cdp:eq:A.T.15}. The finite-module calculation there
identifies $I^{(2)}=I^2+(F)$ in this same convergent
local ring: it uses the nine-element residue basis over
$\mathbb C\{x\}$ and generic rank nine, hence
$x$-saturation. The products \eqref{cdp:eq:C.RA.9}
are the identities under the coordinate identity map.
Our reduced equation belongs to $I^{(2)}$, since its
generic normal order is two.

\item Let $N$ be the dual of the torsion-free pullback to $C$ of
the ambient conormal module of $H$. At \eqref{cdp:eq:C.RA.8} the
normalized conormal basis and its dual lifts are
\begin{equation}
\begin{gathered}
 a=dg_1,\qquad b=dg_2+t\,dg_1,\qquad
 dg_3=-2t^2a+tb,\\
 A=-t^{-3}\partial_z,\qquad
 B=(t^{-5}\partial_y+2t^{-4}\partial_z)/3 .
\end{gathered}\hypertarget{tag-00WR}{}\label{cdp:eq:C.RA.10}
\end{equation}
Their wedge is $a\wedge b=t^6\,\iota_{\gamma'}(dx\wedge
dy\wedge dz)$, where $\gamma=(t^3,t^4,t^5)$.
Thus the determinant formula is
$\det N=K_C\otimes h_C^*K_Y^{-1}(6p)$.
Since every ambient line has degree zero on every compact
curve when $b_2(Y)=0$, its degree here is four.

\item Write the local equation as $cF+f_2$, $f_2\in I^2$.
The conormal basis in \eqref{cdp:eq:C.RA.10} is
$(a,b)=(a_1,b_1)$ in \eqref{cdp:eq:A.T.16}.
The quadratic form of $F$ and its determinant are
therefore exactly \eqref{cdp:eq:C.RA.11}.
The local calculation following that identity uses
only the generic rank-one Hessian: its pole-four
term cannot cancel against a regular quadratic
form, so $c(0)=0$ and \eqref{cdp:eq:C.RA.9} gives
$cF\in I^2$. The resulting quadratic form is regular.
This transfers the rank-one calculation to the present
$A_2$ equation; its radical cubic is considered next.

\item The intrinsic cubic on its radical has pole order at most
ten at $p$. For arbitrary holomorphic coefficients, the bound is as follows. In the lifts \eqref{cdp:eq:C.RA.10},
the differentials $dg_i$ are regular, the second derivatives
$D^2g_i(Au+Bv)^2$ have poles at most ten, and a derivative
of a holomorphic coefficient has pole at most five.
The formula for $D^3(b_{ij}g_i g_j)$ therefore has pole
at most ten. Terms containing an undifferentiated $g_i$
vanish on the curve. In these lifts $dx=0$, so explicitly
$D^2g_1=2(dy)^2$, $D^2g_2=2\,dy\,dz$, and
$D^2g_3=2(dz)^2$; their pole orders are at most
ten, nine, and eight, respectively.
The Leibniz rule gives exactly the asserted bound
for every product with an arbitrary holomorphic coefficient.
The radical cubic is independent of a tangent component
of the lift: differentiating the identically zero quadratic
form on its radical along the smooth part of the curve
gives this invariance, and equality extends meromorphically.

\item For the cusp-plus-line ideal \eqref{cdp:eq:C.RA.7},
use $(g_1,g_2,g_3)=(xz,yz,y^2-x^3)$.
The identity map of these generators identifies
$G=g_3$, $P=(z,G)$ and $Q_1=(x,y)$ with
\eqref{cdp:eq:A.T.17}. Its direct ideal-intersection
calculation gives $I^{(2)}=I^2+(zg_3)$,
and \eqref{cdp:eq:C.RA.12} gives the three annihilating
products in the same ring.
On the cusp branch $(t^2,t^3,0)$ use
$a=dg_1$, $b=dg_3$. The quadratic form of $zg_3$
is $t^{-2}ab$. Its determinant has pole four,
while a cross term with a regular form has pole at
most two. The generic rank-one condition therefore
makes its coefficient vanish at the origin;
\eqref{cdp:eq:C.RA.12} then puts the surface equation
in $I^2$.

The determinant corrections are four on the cusp
branch and two on the smooth branch $(0,0,t)$,
so the total determinant degree is again four.
Dual lifts on the cusp branch are
$t^{-2}\partial_z$ and $\frac12t^{-3}\partial_y$;
on the smooth branch they are
$t^{-1}\partial_x,t^{-1}\partial_y$.
The same product derivative formula bounds the
radical-cubic poles by six and two, respectively,
of total degree eight.

\item In either non-lci case let $R\subset N$ be the saturated
radical of the regular rank-one quadratic form and let
$Q=N/R$. Its nonzero square coefficient is a section of
$Q^{-2}\otimes M$, where $M=\OO_Y(S)|_C$ has degree
zero. Hence $\deg Q\le0$ and $\deg R=4-\deg Q\ge4$.
The nonzero intrinsic cubic of a generic $A_2$ is a
section of $R^{-3}\otimes M$ with poles of total degree
at most ten, or eight. It would give respectively
$3\deg R\le10$ or $3\deg R\le8$, both impossible.
This excludes the two non-lci reduced image germs.

\item We may consequently use the lci conormal bundle along $H$.
Let $\Xi$ be its normalization conductor divisor, of degree
$2\delta$. The determinant formula and the
conductor-cleared intrinsic cubic from \cref{cdp:st:A.5.5} give
\begin{equation}
\begin{gathered}
 \det N=K_C(\Xi)\otimes h_C^*K_Y^{-1},\\
 q\in H^0(C,Q^{-2}\otimes M),\qquad
 k\in H^0(C,R^{-3}\otimes M(2\Xi)),\\
 0\ne q^3k^2\in
 H^0\bigl(C,(\det N)^{-6}\otimes M^5(4\Xi)\bigr),\\
 \deg\bigl((\det N)^{-6}\otimes M^5(4\Xi)\bigr)
                         =12-4\delta .
\end{gathered}\hypertarget{tag-00WU}{}\label{cdp:eq:C.RA.13}
\end{equation}
The section is nonzero because the generic transverse
singularity is a cusp.

\item We give the local comparison that strengthens the zero
count at a smooth point of $H$. A simple zero of $q$ at
which $k$ is nonzero is impossible in a torus smoothing.
Choose coordinates with $H=\{x=y=0\}$, radical direction
$\partial_x$, and coordinate $z$ along $H$.
The nonzero cubic permits Weierstrass preparation in $x$:
\begin{equation}
f_S=U\bigl(x^3+a_2(y,z)x^2+a_1(y,z)x+a_0(y,z)\bigr).\hypertarget{tag-00WV}{}\label{cdp:eq:C.RA.14}
\end{equation}
The critical-curve equations and the rank-one Hessian
along it give $a_2\in(y)$ and $a_1,a_0\in(y^2)$.
After depressing the cube by $x'=x+a_2/3$, this is
\begin{equation}
f_S=U'\bigl(x'^3+y^2(B_1(y,z)x'+B_0(y,z))\bigr).\hypertarget{tag-00WW}{}\label{cdp:eq:C.RA.15}
\end{equation}
The simple zero of the Hessian coefficient says
$\partial_zB_0(0,0)\ne0$. Thus
$z'=B_1(y,z)x'+B_0(y,z)$ is an holomorphic
coordinate with $x',y$, and the reduced equation is a
unit times $x'^3+y^2z'$. The full fibre equation is a
unit times its $m$th power. \cref{cdp:st:A.5.15} excludes precisely
this Hirzebruch germ, including that unit and multiplicity.

\item Every point of $C$ above the smooth part of $H$ where a
vertical conductor component meets it has $q=0$.
Indeed a nonzero Hessian would, by \cref{cdp:st:A.5.17}, give a constant
transverse cusp product with smooth normalization and pure
conductor $2C$; its minimal normalization resolution could
have no such vertical component. At a zero of $q$ the
cubic is holomorphic. If it also vanishes, the section
$q^3k^2$ has zero order at least five. If it does not
vanish, \eqref{cdp:eq:C.RA.14}--\eqref{cdp:eq:C.RA.15} force the zero of $q$ to
have order at least two, and the order is at least six.

\item Each of the $r$ original ruling fibres in \eqref{cdp:eq:C.RA.1} has
a positive-boundary total-transform component meeting $C$,
at a different point of $C$. This remains true at each
blowup: a blowup at its intersection with the coefficient-two
section inserts an exceptional curve of positive coefficient.
Using \eqref{cdp:eq:C.RA.6} and the degree in \eqref{cdp:eq:C.RA.13} gives
\begin{equation}
\begin{array}{c|c|c|c}
 \delta& s\text{ at most}&
 \text{bad points above }H_{\rm reg}\text{ at most}&r\text{ at most}\\ \hline
 0&0&2&2\\
 1&2&1&3\\
 2&3&0&3 .
 \end{array}\hypertarget{tag-00WX}{}\label{cdp:eq:C.RA.16}
\end{equation}
This proves the asserted bound on $r$.

\item Finally \eqref{cdp:eq:A.RT.12} and \eqref{cdp:eq:A.RT.16} give
$e(Z\setminus D_{\rm red})=2-r+N_0$.
The forest argument of \eqref{cdp:eq:A.RT.21} applies to this entire
star class, whether its minimal complement is negative
or not, and gives $e(G)\ge2$.
Together with $r\le3$, \eqref{cdp:eq:C.RA.2} gives directly
\begin{equation}
e(S)\ge4-r+N_0\ge1+N_0.\hypertarget{tag-00WY}{}\label{cdp:eq:C.RA.17}
\end{equation}

\end{enumerate}
\end{proof}


\begin{proposition}[Attachments of a rational horizontal boundary]\hypertarget{tag-01CF}{}\label{cdp:st:C.2.RAT.3}
\StatementRecord{cdp:st:C.2.RAT.3}{proved}{01CF}{\cref{cdp:st:C.2.RAT.1}, \eqref{cdp:eq:C.RA.22}, \eqref{cdp:eq:C.RA.23}, \eqref{cdp:eq:C.RA.24}}

Retain the torus-fibre, rational-normalization and ruled-boundary
morphism of \cref{cdp:st:C.2.RAT.1}. Suppose that $C$ is not
contracted, and put $H=h(C)_{\rm red}$.
Let $r_{\rm dom}$ count the original ruling fibres $F_j$
whose positive vertical boundary contains a strict
conductor component dominating $H$.
Among the other fibres select those with a component
not exceptional for $\pi$, and write $K$ for their number.
They meet $C$ in $K$ distinct points. Let $\ell$ be the
number of identifications among their images on $H^\nu$,
so those images have $k=K-\ell$ distinct points.
Then
\[
 K\ge r-N_0-r_{\rm dom},\qquad
 e(S)\ge2-r_{\rm dom}-\ell .
\]
For a degree-one map $C\to H^\nu$, one has $\ell=0$ and
\begin{equation}
0\le r_{\rm dom}\le2,\qquad
       e(S)\ge2-r_{\rm dom}\ge0.\hypertarget{tag-00X3}{}\label{cdp:eq:C.RA.21}
\end{equation}
The generic transverse type and the conductor scheme
are those of the given fibre.
\end{proposition}

\begin{proof}

\begin{enumerate}[label=\arabic*.,leftmargin=*]
\item In the degree-one case the conductor coefficient on $C$ is two, because $C$ is
the strict transform of the coefficient-two section in
\eqref{cdp:eq:C.RA.1} and is not $\pi$-exceptional. At a general point
of $H$, its degree-one map supplies one branch of the reduced
transverse plane curve. If that branch has intrinsic
conductor exponent $c_i^{\rm int}$, its exponent in the
whole plane curve is
\begin{equation}
2=c_i^{\rm int}+\sum_{j\ne i}I(C_i,C_j).\hypertarget{tag-00X4}{}\label{cdp:eq:C.RA.22}
\end{equation}
Here $I(C_i,C_j)$ is the positive intersection multiplicity
of two different plane branches. The residue comparison of \cref{lem:II:plane-branch-conductor} gives this exponent with $c_i=2$. Thus there are at most two other generic
branches. Each vertical strict component dominating $H$
supplies at least one such branch, counted with its covering
degree. This proves $r_{\rm dom}\le2$.

\item For an original fibre $F_j$, let $B_j$ be the positive
vertical boundary in its total transform. Its reduced
support is connected and meets $C$ at a single point.
These assertions follow at each point blowup: a centre
at the intersection of positive components inserts a
positive exceptional component; a centre at a smooth
coefficient-one point creates only a coefficient-zero
leaf; and a centre at a smooth point of coefficient
at least two creates a positive leaf. These are all the
possibilities in the exact divisor formula
$D_{\rm new}=\sigma^*D_{\rm old}-E$.
The $r$ points of $C$ so obtained are distinct.

\item If every component of $B_j$ is $\pi$-exceptional, there
has been at least one multiplicity-one blowup centre
above $F_j$. Indeed, in the absence of such a centre,
every component of the entire total transform of $F_j$
has positive conductor coefficient. That complete
ruling fibre, with its positive multiplicities, has
self-intersection zero. It cannot be supported on the
negative-definite exceptional matrix of a normal surface
resolution. Distinct original fibres require distinct
blowup events. Hence at most $N_0$ of the $B_j$ can be
entirely $\pi$-exceptional.

\item Discard those fibres and the $r_{\rm dom}$ fibres counted
above. The number of selected fibres satisfies
\begin{equation}
K\ge r-N_0-r_{\rm dom},\qquad k=K-\ell\hypertarget{tag-00X5}{}\label{cdp:eq:C.RA.23}
\end{equation}
The selected fibres have a noncontracted strict conductor
component and none dominating $H$. Follow the connected positive
boundary from its point on $C$ to the first component
which is not $\pi$-exceptional. Its image is a curve
different from $H$ meeting $H$ at that point's image.
Every intervening connected chain is contracted by
$\pi$, so its common image is the stated intersection
of curves in $S$.

\item Use the normalization-incidence graph of
$G=\Gamma_{\rm red}$: its vertices are the rational
normalization components and the singular points with
more than one normalization branch, and its edges are
the individual branch preimages. After identifying the selected points with the same
image on $H^\nu$, the $k$ distinct images give $k$
distinct branch-preimage edges at the vertex of $H$.
They remain distinct when their images on $H$ agree
but the points on $H^\nu$ differ. This is precisely
the definition of edges in the normalization-incidence graph. Remove that vertex and all its adjacent
edges. Each of the $c$ remaining connected components
which receives one of these edges contains a
normalization vertex of a curve other than $H$.
In particular $c\le n_G-1$, where $n_G$ is the number
of irreducible components of $G$.

\item Contract a spanning tree in each of those $c$ graph
components. The $k$ selected edges give a connected
multistar with $c+1$ vertices and $k$ edges. Its first
Betti number is $k-c$. For $k=0$ the inequality
$n_G\ge k$ is immediate; otherwise this multistar
gives the displayed bound. Hence the first Betti
number $\beta$ of the full incidence graph satisfies
\begin{equation}
k-c\le\beta\le1,
          \qquad n_G\ge k+1-\beta\ge k.\hypertarget{tag-00X6}{}\label{cdp:eq:C.RA.24}
\end{equation}
The upper bound on $\beta$ is the local torus-group
quotient \eqref{cdp:eq:A.RT.18}--\eqref{cdp:eq:A.RT.20}. Adding genus-zero
normalization components changes no graph first Betti
number, and the normalization quotient gives
$e(G)=n_G+1-\beta\ge n_G$. Combining \eqref{cdp:eq:C.RA.23}--\eqref{cdp:eq:C.RA.24}
with the exact surface Euler formula yields
\begin{equation}
\begin{split}
 e(S)&=2-r+N_0+e(G)\\
     &\ge2-r+N_0+n_G
      \ge2-r+N_0+(r-N_0-r_{\rm dom}-\ell)
       =2-r_{\rm dom}-\ell.
\end{split}\hypertarget{tag-00X7}{}\label{cdp:eq:C.RA.25}
\end{equation}
For degree one $C\to H^\nu$ is an isomorphism, so
$\ell=0$, and the first step gives $r_{\rm dom}\le2$.
This proves both assertions.

\end{enumerate}
\end{proof}


\begin{proposition}[Attachment bounds for rational tacnode covers]
\hypertarget{tag-01CE}{}\label{cdp:st:C.2.RAT.2}
\StatementRecord{cdp:st:C.2.RAT.2}{proved}{01CE}{}
Use the torus-fibre, rational-normalization and ruled-boundary
morphism of \cref{cdp:st:C.2.RAT.1}. Suppose $C$ is not
contracted, $H=h(C)_{\rm red}$ has generic transverse type
$A_3$, and $C\to H^\nu$ has degree two. For the selected
vertical trees of \cref{cdp:st:C.2.RAT.3}, no tree dominates
$H$ and the attachment loss satisfies $\ell\le1$.
This includes singular and non-lci image curves.

For a smooth image $H$, the number $r$ of distinct vertical
fibres in the minimal boundary satisfies $r\le2$, and
\begin{equation}
e(S)\ge2+N_0.\hypertarget{tag-00WZ}{}\label{cdp:eq:C.RA.18}
\end{equation}
\end{proposition}

\begin{proof}

\begin{enumerate}[label=\arabic*.,leftmargin=*]
\item Write $\Sigma=H^\nu\simeq\PP^1$, $G=\Gamma_{\rm red}$,
$\delta=p_a(H)\le2$, and
$s=\#\nu^{-1}(\operatorname{Sing}H)$. The reduced boundary
is a rational tree: every reverse blowup adds a leaf,
subdivides an edge, or leaves its strict support isomorphic
to the old support. The pushout maps
\eqref{cdp:eq:A.RT.18}--\eqref{cdp:eq:A.RT.19} therefore give
the inclusion-induced isomorphism $\pi_1(G)\simeq\pi_1(S)$
for this fixed ruled boundary. Each normalized component
of $G$ is dominated by a noncontracted rational boundary
component and is therefore rational. Its normalization-incidence
graph makes $\pi_1(G)$ free. By \cref{night:P4R.1}, this
free group is virtually abelian, so its rank is at most one.
The normalization graph of $H$ is a subgraph, hence
$b_1(H)\le b_1(G)\le1$. The reduced-curve union sequence
gives $p_a(H)\le p_a(G)\le2$, since every connected component
of the complementary subcurve meets $H$. Consequently
$s=\#\operatorname{Sing}H+b_1(H)\le\delta+1\le3$,
as in \eqref{cdp:eq:C.RA.5}--\eqref{cdp:eq:C.RA.6}. There are exactly two generic
normalization branches above $H$, both supplied by the
degree-two map from $C$. Thus no vertical strict conductor
component dominates $H$.

\item Apply \cref{lem:II:delta-two-curve-rings} to $H$
after excluding the three-branch germ by $b_1(H)\le1$.
Each nonplanar alternative has its single delta-two
singularity and the fixed embedded ideal.
Apply \cref{cor:II:nonplanar-tacnode-cover} with its
rational normalization $\Sigma=H^\nu$, its same
degree-two conductor component $C\simeq\mathbb P^1$,
and its generic $A_3$ equation. All other boundary
components are vertical and supply no additional
branch over $H$. That corollary excludes exactly
these two alternatives, so $H$ is lci.

\item Let $\Xi$ be the normalization conductor divisor of $H$,
and put $I=\mathcal I_{H/Y}$, $N=(\nu^*(I/I^2))^*$ and
$M=\nu^*(\mathcal O_Y(S)|_H)$. Let $R\subset N$ be
the saturated radical, $Q=N/R$, and
$L=M\otimes Q^{-1}\otimes R^{-2}$.
The conductor-cleared construction
\eqref{cdp:eq:A.J.7}--\eqref{cdp:eq:A.J.9}, with
$\deg N=-2+2\delta$ and $\deg M=0$, gives
\begin{equation}
A=\deg\operatorname{div}(a)=-2\deg Q,\qquad
 \sigma\in H^0(\Sigma,L^2(4\Xi)),\qquad
 \deg\operatorname{div}(\sigma)=8-A.\hypertarget{tag-00X8}{}\label{cdp:eq:C.RA.26}
\end{equation}
Here $A$ is a nonnegative even integer and $\sigma$ is
the conductor-cleared version of $\Delta=ae-b^2/4$.
In the generic quadratic algebra
$aw^2+bw+e=0$, the maps $d=aw+b/2$ and
$w=(d-b/2)/a$ are inverse and give $d^2=-\Delta$.
Integral closure extends them to the normalized cover.
For a local generator $c(t)$ of $\mathcal O_\Sigma(-\Xi)$,
the inverse field changes $d_0=c(t)^2d$ and
$d=d_0/c(t)^2$ give $d_0^2=-\sigma$ in the corresponding
local line frames. Write $-\sigma=t^j u(t)$, with $j\ge0$
and $u$ a unit. For odd $j$ its normalized parameter is
$t=s^2$, $d_0=s^j\sqrt{u(s^2)}$, with inverse
$s=d_0/(t^{(j-1)/2}\sqrt{u(t)})$; for even $j$ its
two roots are holomorphic and the cover is unramified.
Thus the degree-two cover $C\to\Sigma$
has two branch points, exactly the odd-order zeros of $\sigma$. Let $B_s$ be
the number of these branch points lying over singular
points of $H$; there are $2-B_s$ over its smooth locus.
Every smooth branch point has zero order at least three
in $\sigma$, by the coordinate exclusion
\eqref{cdp:eq:A.T.10}--\eqref{cdp:eq:A.T.12}.

\item The section $a$ must vanish at every point of $\Sigma$
over a singular point of $H$. To see this, the Hessian
matrix has entries in $\OO_H$. If its normalized scalar
$a$ is nonzero on one branch above $p$, its primitive
rank-one factor has a unit component there. A diagonal
entry consequently has a nonzero value at $p$. That
same entry has the same nonzero value on every branch,
so $a$ is nonzero at all normalization points above $p$.
In that event \eqref{cdp:eq:A.T.13} removes the conductor denominator:
$\Delta$ is holomorphic in $L^2$ there. The orders of
$\sigma$ above $p$ are at least those of $4\Xi$,
which are even. Its total zero divisor must additionally
contain the two parity contributions of the two branch
points, whether or not they lie over $p$. Thus
\begin{equation}
\deg\operatorname{div}(\sigma)
       \ge4\deg(\Xi|_{\nu^{-1}(p)})+2
       =8\delta_p+2\ge10,\hypertarget{tag-00X9}{}\label{cdp:eq:C.RA.27}
\end{equation}
contrary to \eqref{cdp:eq:C.RA.26}. This proves the claimed vanishing.

\item It follows in particular that
\begin{equation}
A\ge2s-2.\hypertarget{tag-00XA}{}\label{cdp:eq:C.RA.28}
\end{equation}
For $s\le1$ this is automatic; for $s=2$ the two
zeros give $A\ge2$. If $s=3$, the bounds
$\delta\le2$ and $b_1(H)\le1$ force exactly two
singular points, each of delta invariant one: one node
and one cusp. The two node preimages contribute at
least one each. At the cusp the ring is
$\mathbb C\{t^2,t^3\}$, so any vanishing Hessian entry
has order at least two. Dividing such a diagonal entry
by the square of a unit component of the primitive
normalized factor shows that $a$ has order at least
two there. Thus $A\ge1+1+2=4=2s-2$.

\item At a smooth point with $a\ne0$,
\cref{cdp:st:A.5.17} gives the constant $A_3$ product and no vertical
attachment. Suppose instead that $a=0$ and $\Delta$
is a unit. In the coordinates of \eqref{cdp:eq:A.T.10}, the
coefficient $b$ is then a unit. Substitute $x=y^2\lambda$
in that equation and divide by $y^4$. The result is
holomorphic in $(\lambda,y,z)$ and has, at the point
in question, the simple root
$\lambda=-e(0)/b(0)$. The implicit-function theorem
gives a holomorphic $\lambda(y,z)$ and a factorization
\begin{equation}
f_0=F_1F_2,\qquad
 F_1=x-y^2\lambda(y,z),\qquad
 F_2|_{F_1=0}=y^2\,u(y,z),\quad u(0)\ne0.\hypertarget{tag-00XB}{}\label{cdp:eq:C.RA.29}
\end{equation}
The last equality is the restriction of
$\partial f_0/\partial x$ to $F_1=0$; its quotient
by $y^2$ has constant value $b(0)$.
The branch $F_1=0$ is smooth. In the normalization
product, an element supported on this branch belongs to
the original hypersurface ring exactly when its value is
divisible by $F_2\bmod F_1$: a lift zero modulo $F_2$
must be a multiple of $F_2$. Its conductor ideal is
therefore $(y^2u)=(y^2)$, giving exactly the divisor
$2\{y=0\}$, as in \eqref{cdp:eq:C.RA.22}. There is therefore no vertical conductor
near this normalization point. Its minimal resolution
is an isomorphism there as well. Since $\Delta$ is a
unit, this is one of the two distinct conductor
preimages over the given point of $\Sigma$.
Consequently a smooth nonbranch value with two vertical
attachments has a positive even zero order of $\sigma$,
at least two.

\item Let $d_{\rm reg}$ count distinct smooth values with
two such attachments. The two branch points and these
distinct nonbranch zeros give
\begin{equation}
3(2-B_s)+B_s+2d_{\rm reg}\le8-A.\hypertarget{tag-00XD}{}\label{cdp:eq:C.RA.30}
\end{equation}
Among the $s$ points above the singular locus, the
$B_s$ branch points each have only one preimage on
$C$. Thus at most $s-B_s$ of these values can have
two attachments. If $d$ counts all values on
$\Sigma$ with two attachments, \eqref{cdp:eq:C.RA.28}--\eqref{cdp:eq:C.RA.30}
give
\begin{equation}
d\le s-B_s+d_{\rm reg}
       \le s+1-\frac A2\le2.\hypertarget{tag-00XE}{}\label{cdp:eq:C.RA.31}
\end{equation}

\item In fact $d\le1$. If $d=2$, \eqref{cdp:eq:C.RA.31} requires
$A\le2s-2$. For $s=0$ this is impossible, and for
$s=1$ it contradicts the vanishing of $a$ over the
singular point. For $s=2$, it forces $A=2$ and
the two singular preimages exhaust the zero divisor
of $a$. For $s=3$, it forces $A=4$, exhausted by
the node and cusp orders $1+1+2$ established above.
In both remaining cases $a$ is nowhere zero over
$H_{\rm reg}$. The constant $A_3$ product there
has neither a cover branch point nor a vertical
attachment. Thus both cover branch points lie
over $\operatorname{Sing}H$, and only the other
$s-2\le1$ singular preimages can have two
attachments. This again contradicts $d=2$.
Consequently
\begin{equation}
d\le1.\hypertarget{tag-00XF}{}\label{cdp:eq:C.RA.31a}
\end{equation}

\item Apply \cref{cdp:st:C.2.RAT.3} to this same
ruled-boundary map and its selected attachments.
Here $r_{\rm dom}=0$, and its loss is $\ell=d'$:
a degree-two fibre loses one distinct image exactly
when both selected points lie over that value.
The preceding bound gives $d'\le d\le1$, so
\begin{equation}
\begin{split}
 e(S)&=2-r+N_0+e(G)\ge2-K+e(G)\\
     &\ge2-K+(K-d')=2-d'\ge1 .
\end{split}\hypertarget{tag-00XG}{}\label{cdp:eq:C.RA.32}
\end{equation}
These are the same maps $Z\to T\to S$ and
$C\to\Sigma\to H$ as in the statement.

\item For the smooth-image refinement, the two branches of a generic tacnode are both supplied by
$C\to H$; no other strict conductor component can dominate $H$.
Put $N=N_{H/Y}$ and $M=\mathcal O_Y(S)|_H$. The condition
$b_2(Y)=0$ and adjunction give $\deg M=0$ and $\deg N=-2$.
Let $R\subset N$ be the saturated radical of the generic
rank-one Hessian, and put $Q=N/R$ and $q_0=\deg Q$.
The normal-jet construction
\eqref{cdp:eq:A.T.5}--\eqref{cdp:eq:A.T.7} applies with
curve conductor divisor $\Xi=0$, since $H\simeq\PP^1$ is smooth.
Thus $\deg R=-2-q_0$ and $q_0\le0$.
The quadratic coefficient $a$ and the quartic invariant
$\Delta=ae-b^2/4$ are nonzero sections with
\begin{equation}
A:=\deg\operatorname{div}(a)=-2q_0,\qquad
 \deg\operatorname{div}(\Delta)=8+2q_0=8-A.\hypertarget{tag-00X0}{}\label{cdp:eq:C.RA.19}
\end{equation}
In particular $A$ is a nonnegative even integer. These are
sections of the line bundles $M\otimes Q^{-2}$ and
$(M\otimes Q^{-1}\otimes R^{-2})^2$, with
$\deg M=0$ and $\deg R=-2-q_0$.

\item Use the same normalized cover and quadratic field
from \eqref{cdp:eq:C.RA.26}: $\Xi=0$ and
$\sigma=\Delta$. The inverse field maps are
$d=aw+b/2$, $w=(d-b/2)/a$, and the odd-order
normalization parameters are exactly those proved above.
Since
$C\simeq\PP^1$ and the map has degree two, Riemann--Hurwitz
gives two distinct branch points on $H$. At each of them
$a$ vanishes: a nonzero Hessian would give the constant
transverse $A_3$ product by \cref{cdp:st:A.5.17} and hence an unramified
conductor cover. The coordinate comparison
\eqref{cdp:eq:A.T.10}--\eqref{cdp:eq:A.T.12} excludes a simple zero of $\Delta$ at
such a point. Its odd zero order is consequently at least
three. Equations \eqref{cdp:eq:C.RA.19} give
\begin{equation}
A\ge2,\qquad 8-A\ge6.\hypertarget{tag-00X2}{}\label{cdp:eq:C.RA.20}
\end{equation}
Therefore $A=2$. The two branch points are exactly the
zeros of $a$, each simple, and the zeros of $\Delta$ have
order three there, with no further zeros.

\item Every original fibre $F_j$ of \eqref{cdp:eq:C.RA.1} gives a positive
vertical boundary component meeting $C$. At its intersection with the
coefficient-two horizontal section, a blowup replaces an
attachment of vertical coefficient $v\ge1$ by one of
exceptional coefficient $2+v-1=v+1$. The strict transform of
the section maps isomorphically to the old section, so the
original $r$ distinct attachment points remain distinct. At any such point
the Hessian coefficient at its image must vanish. Otherwise
\cref{cdp:st:A.5.17} supplies the constant $A_3$ product, whose two
normalization branches are smooth and have no vertical
conductor; the minimal resolution $\pi$ is an isomorphism
near them. Each of the two zeros of $a$ is a branch point
of $C\to H$, and consequently has just one point above it
on $C$. It follows that $r\le2$.

\item The maps \eqref{cdp:eq:A.RT.18}--\eqref{cdp:eq:A.RT.20}
apply to this same rational tree boundary and identify
$\pi_1(G)$ with the virtually abelian group $\pi_1(S)$.
Thus its free normalization-incidence group has rank
$\beta\le1$ and $e(G)=n_G+1-\beta\ge n_G$.
For $n_G=1$ the entire reduced image is the smooth
curve $H\simeq\mathbb P^1$, with Euler number two;
for $n_G\ge2$ the inequality also gives $e(G)\ge2$.
The complement identity \eqref{cdp:eq:A.RT.12} gives
\[
 e(S)=2-r+N_0+e(G)\ge4-r+N_0\ge2+N_0 .
\]

\end{enumerate}
\end{proof}


\subsection{Normal cubics, cyclic covers and rational images}\label{subsec:II:normal-cubic-covers}

\begin{lemma}[Normal cubic covers and compact cycles]
\hypertarget{tag-P361}{}\label{lem:II:normal-cubic-compact-core}
\StatementRecord{lem:II:normal-cubic-compact-core}{proved}{P361}{\cref{lem:II:torus-tube-numerical}}
Let $f:Y\to\Delta$ be a proper holomorphic map from a smooth
complex threefold, with connected fibres, smooth two-torus
fibres over $\Delta^*$, and full central fibre $mS$ with
irreducible reduced support. Let $H\subset S$ be a smooth
rational curve with generic ordinary-triple transversal.
Then $M=\OO_Y(S)|_H$ is trivial and the rank-two normal
bundle $N=N_{H/Y}$ has degree $-2$. A trivialization of $M$
identifies the cubic normal jet with
\begin{equation}
 P\in H^0(H,\operatorname{Sym}^3N^*),\qquad
 U_P=\{v\in N:P(v)=1\}.
\hypertarget{tag-00E1}{}\label{cdp:eq:C.D4.1}
\end{equation}
Let $R_h\subset\PP(N)$ be the horizontal root divisor of $P$.
For a connected normalized conductor cover $C\to H$ of degree
$d$, let $R_d$ be its generic tangent-root component. There is
a finite-cover isomorphism over $H$:
\begin{equation}
C\simeq R_d^\nu
                            \quad\text{over }H.\hypertarget{tag-00Y9}{}\label{cdp:eq:C.RD.2}
\end{equation}
The surface $U_P$ is smooth and has a compact smooth core.
Every compact codimension-zero smooth submanifold of $U_P$
embeds, with its complex orientation and compact intersection
pairing, in a sufficiently close smooth fibre of $f$.
Consequently $b_1(U_P)\le3$ forces its compact intersection
pairing to vanish. Every compact spherical class
in $U_P$ has square zero.
\end{lemma}
\begin{proof}
\begin{enumerate}[label=\arabic*.,leftmargin=*]
\item On a small tube, \cref{lem:II:torus-tube-numerical}
gives $K_Y\simeq\OO_Y(kS)$ and $\OO_Y(mS)\simeq\OO_Y$.
Their degrees on $H$ are zero. Thus $\deg M=0$ and
$M\simeq\OO_{\mathbb P^1}$; the tangent-normal exact sequence
on $H$ gives $\deg N=-2$.
The ideal of the smooth codimension-two curve is generated
by a regular sequence. Its powers have no embedded
associated prime and are primary along $H$. The generic
order-three vanishing of the reduced equation therefore
holds at every point of $H$, and defines the section $P$.
Its generic cubic has three distinct roots.
The radial derivative satisfies $dP_v(v)=3P(v)$, proving
smoothness of $U_P$. A normal fibre on which the whole
cubic vanishes has empty intersection with $U_P$.

\item The cubic has three distinct roots over the generic point.
On a base disc divide its coefficients by their common factor.
The resulting primitive cubic is nonzero on every projective fibre
and defines the horizontal root divisor $R_h$. On overlaps the
two primitive cubics differ by a unit, so their zero schemes agree.
Thus $R_h\to H$ is proper with finite fibres, and is finite of
degree three. Its generic roots are distinct, so it is reduced.
Let $R_d$ be the root component corresponding to $C$.
The function-field algebra of its roots agrees with that of
the normalized conductor curves dominating $H$: a branch
of an ordinary triple is distinguished by its tangent line.
The field map sends the root coordinate $x/y$ to the
tangent ratio on the corresponding normalization branch;
its inverse identifies that branch field with the same
degree-$d$ root field. Taking the integral closure of $\OO_H$
in this field gives the finite-cover identification \eqref{cdp:eq:C.RD.2}.


\item Write $f=u s_0^m$ locally, with $s_0$ a reduced
equation of $S$ and $u$ a unit. The initial normal jet of
the global function $f$ is a nonzero constant multiple
of $P^m$: their ratio is a nowhere-zero holomorphic
function on $H\simeq\mathbb P^1$. Rescale the chosen
trivialization of $M$ by a constant to make that ratio one.
Choose a smooth tubular map $\tau$ from a neighbourhood
of the zero section of $N$ to $Y$, with normal derivative
the given complex-linear normal-bundle identification.
Taylor expansion, including all derivatives on each
compact subset, gives for positive real $s\to0$
\begin{equation}
 s^{-3m}f\bigl(\tau(sv)\bigr)\longrightarrow P(v)^m
       \quad\text{in }C^\infty\text{ on compact subsets of }N.
\hypertarget{tag-00EC}{}\label{cdp:eq:C.D4.2}
\end{equation}
The order $3m$ holds everywhere by the same primary-power
argument for the full equation $f$.

\item On a neighbourhood of a given compact submanifold
$K\subset U_P$, the differential of $P^m$ has a smooth
right inverse as a real map to $\mathbb C$.
Interpolate between $P^m$ and the left side of
\eqref{cdp:eq:C.D4.2}; for sufficiently small $s$ their
differentials remain surjective near $K$. Composing the
parameter derivative with these right inverses and
integrating its negative gives a smooth transport from
the level $P^m=1$ to
$s^{-3m}f(\tau(sv))=1$. The transport and its reverse
are inverse on smaller neighbourhoods of $K$.
Composition with $v\mapsto\tau(sv)$ is the required
embedding into $f^{-1}(s^{3m})$. It preserves orientation,
because its differential deforms the fixed
complex-linear normal identification. The same compact
cycles and their small transverse perturbations compute
their intersections before and after this embedding.

\item The vector bundle $N$ on $\mathbb P^1$ and its
holomorphic section $P$ are algebraic. Resolve the boundary
in a projective compactification of $U_P$. Compatible
triangulations and a sufficiently small regular neighbourhood
of that boundary give, after rounding corners, a compact
smooth core $K\subset U_P$ whose inclusion is a homotopy
equivalence. The radial collars give its homotopy inverse.
Thus one of the preceding embeddings carries all of
$H_1(U_P)$ and $H_2(U_P)$ and their compact intersection form.

\item For its map $K\to T^4$, let
$\Lambda'\subset H_1(T^4,\mathbb Z)=\mathbb Z^4$ be
the fundamental-group image and put
$W=\Lambda'\otimes\mathbb Q$.
The map lifts to $\mathbb R^4/\Lambda'$, homotopy equivalent
to $T^{\operatorname{rank}\Lambda'}$.
The induced second-homology map factors through
$\Lambda^2W\subset\Lambda^2\mathbb Q^4$.
For $b_1(U_P)\le3$, the rank of $\Lambda'$ is at most three,
so $\Lambda^4W=0$. The wedge intersection form vanishes
on this image, and the preserved compact form on $U_P$
therefore vanishes as well.
Finally a spherical class maps to zero in $H_2(T^4,\mathbb Z)$
because $\pi_2(T^4)=0$. Preservation of its square proves
the last assertion.
\end{enumerate}
\end{proof}


\begin{lemma}[Cyclic cubic covers and boundary cohomology]
\hypertarget{tag-P362}{}\label{lem:II:cyclic-cubic-boundary}
\StatementRecord{lem:II:cyclic-cubic-boundary}{proved}{P362}{\cite[\S6, Haupts\"atze I and IIa, p.~59]{Grauert1960D}; \cite[pp.~35--36]{Grauert1963ErratumD}; \cite[Satz~2.3, pp.~274--276]{GrauertRiemenschneider1970}; \cite[Tags~0AU3 and~0A81]{StacksProject}}
Let $N$ be a rank-two bundle of degree $-2$ on
$H=\mathbb P^1$, and let
$P\in H^0(H,\operatorname{Sym}^3N^*)$ have three distinct
generic roots. Use the bundle of lines $p:W=\mathbb P(N)\to H$,
put $L=\OO_W(1)$, $\xi=c_1(L)$, and write $F$ for a fibre class.
These tautological data satisfy \eqref{cdp:eq:C.DD.2}.
Let $B_{\rm tot}$ be the divisor of the induced section
$P\in H^0(W,L^3)$, and let $B=(B_{\rm tot})_{\rm red}$.
Assume that $B_{\rm tot}$ is reduced, or that it is the
sum of reduced weight-one components and twice a single
ruling fibre $F_0$. Put $U_P=\{P=1\}\subset N$.

The normalized cyclic cover $g_0:X_0\to W$ is connected
and has complement $U_P$ over $W\setminus B$.
Let $\rho:X\to X_0$ be a projective resolution on which
$D=X\setminus U_P$ has simple normal crossings, and put
$g=g_0\rho$.
Write $c$ and $b_0$ for the numbers of irreducible and
connected components of $B$, $g_B$ for the sum of their
normalization genera, and
$\nu=\sum_x(r_x(B)-1)$ for its total branch excess.
Let $r_B$ be the rank of the classes $[B_i]$ in
$H_2(W,\mathbb Q)$, put $\varepsilon=c-r_B$, and let
$\lambda$ be the sum of the first Betti numbers of the
reduced exceptional divisors over the singularities of $X_0$.
Then the identities and bounds
\eqref{cdp:eq:C.DD.2}, \eqref{cdp:eq:C.DD.7}--\eqref{cdp:eq:C.DD.10}
hold. The reduced case has algebra \eqref{cdp:eq:C.DD.6},
trivial dualizing sheaf and $\chi(\OO_{X_0})=2$.
Writing $P_g=\sum_xp_g(X_0,x)$, in that case
\eqref{cdp:eq:C.RD.30} holds.
The double-fibre case has the algebra and direct image
in \eqref{cdp:eq:C.DD.16}--\eqref{cdp:eq:C.DD.17}.
At a smooth weight-two branch meeting a smooth weight-one
branch with contact $k\ge1$, the normalized cover has a
rational singularity; the exact resolution and character
maps are \eqref{cdp:eq:C.DD.18}--\eqref{cdp:eq:C.DD.20}.
\end{lemma}
\begin{proof}
\begin{enumerate}[label=\arabic*.,leftmargin=*]
\item The tautological sequence and relative canonical line,
with the stated convention of lines in $N$, give
\begin{equation}
 \xi^2=-\deg N=2,\qquad \xi F=1,\qquad F^2=0,\qquad
 K_W=-2\xi .
\hypertarget{tag-00EW}{}\label{cdp:eq:C.DD.2}
\end{equation}
For reduced $B_{\rm tot}$ the cyclic algebra is
\begin{equation}
 X_0=\operatorname{Spec}_W
       (\OO_W\oplus L^{-1}\oplus L^{-2}),\qquad w^3=P.
\hypertarget{tag-00F2}{}\hypertarget{tag-00EH}{}
\SsixEquationAliases{\label{cdp:eq:C.D4.6}}\label{cdp:eq:C.DD.6}
\end{equation}
Its hypersurface rings are Cohen--Macaulay and regular
in codimension one, hence normal.
A meridian about any weight-one horizontal branch acts
on the three sheets as a three-cycle; thus the cover is
connected. In the double-fibre case the same argument
applies to its normalization.

For $\ell\in W\setminus B$, a nonzero dual vector
$w\in\ell^*$ with $w^3=P|_\ell$ corresponds to the unique
$v\in\ell$ satisfying $w(v)=1$. Conversely, $P(v)=1$
defines the unique linear functional $w$ with $w(v)=1$.
These maps are inverse and commute with every change
of line-bundle frame. They identify the complement
with $U_P$, including above any singularities of $B$.
In the case $N=\OO(-1)^{\oplus2}$ they identify
$W$ with $\mathbb P^1\times\mathbb P^1$ and
$L$ with $\OO(1,1)$, giving the product-bundle version
of \eqref{cdp:eq:C.D4.6}.

\item Fibre cohomology gives $R^jp_*L^{-1}=0$ for all $j$.
Relative duality gives
\begin{equation}
 p_*L^{-2}=0,\qquad R^1p_*L^{-2}=\det N=\OO_H(-2).
\hypertarget{tag-00F3}{}\label{cdp:eq:C.DD.7}
\end{equation}
Consequently
\begin{equation}
 H^1(W,\OO_W)=H^1(W,L^{-1})=H^1(W,L^{-2})=0.
 \hypertarget{tag-00Y2}{}\label{cdp:eq:C.RD.13}
\end{equation}
For the doubled fibre, the extra integral generator
at its generic point in $w^3=s^2u$ is $w^2/s$.
The three character summands of the normal finite
algebra are the rank-one reflexive extensions of
these codimension-one modules. On the smooth surface
$W$ those extensions are line bundles, giving
\begin{equation}
 (g_0)_*\OO_{X_0}=\OO_W\oplus L^{-1}\oplus L^{-2}(F_0).
\hypertarget{tag-00ES}{}\hypertarget{tag-00YO}{}\hypertarget{tag-00EA}{}\SsixEquationAliases{\label{cdp:eq:C.RD.33}\label{cdp:eq:C.D4.18}}\label{cdp:eq:C.DD.16}
\end{equation}
This also identifies the multiplication: it is that
of $1,w,w^2/s$ in the common function field.
Projection formula in \eqref{cdp:eq:C.DD.7} gives
\begin{equation}
 R^1p_*\bigl(L^{-2}(F_0)\bigr)=\det N\otimes\OO_H(1)
                              =\OO_H(-1).
\hypertarget{tag-00ET}{}\label{cdp:eq:C.DD.17}
\end{equation}
Thus $H^1(X_0,\OO_{X_0})=0$ in both cases.
Resolution Leray and Hodge decomposition on $X$ give
\begin{equation}
 b_1(X)=2q(X)\le2\sum_xp_g(X_0,x).
\hypertarget{tag-00F4}{}\label{cdp:eq:C.DD.8}
\end{equation}

\item The exceptional classes on $X$ are independent by
their negative-definite intersection matrix.
Every reduced branch component has ramification
index three, also for weight two. Thus, writing
$\widetilde B_i$ for its strict inverse image and
$E_j$ for the exceptional components,
\[
 g^*B_i=3\widetilde B_i+\sum_j a_{ji}E_j.
\]
A relation $\sum_i\alpha_i[B_i]=0$ lifts to the
relation with strict coefficients $\alpha_i$ and
exceptional coefficients $\sum_i a_{ji}\alpha_i/3$.
Conversely, pushing a relation down to $W$ and
recording its strict coefficients recovers those
$\alpha_i$; independence of the exceptional classes
then recovers its exceptional coefficients.
This proves that there are exactly $\varepsilon$
boundary relations. Pullback from $W$ is injective
and multiplies its intersection form by three.
The boundary span contains $g^*B_{\rm tot}$,
whose square is $3(3\xi)^2=54>0$.
The Hodge index theorem therefore makes that span
nondegenerate: its orthogonal part to this positive
class is negative definite.

\item Put $n_D=\#\{\text{irreducible components of }D\}$.
The relative homology map
$H_2(X,\mathbb Q)\to H_2(X,U_P;\mathbb Q)=\mathbb Q^{n_D}$
is intersection with the boundary components.
Its image has rank $n_D-\varepsilon$; its kernel is
the orthogonal complement of their span, and is
the image of $H_2(U_P)$.
The cokernel of this relative map has dimension
$\varepsilon$, and the pair sequence maps it
injectively to $H_1(U_P)$, with quotient $H_1(X)$.
Thus the pairing on that kernel has dimension
$b_2(X)-n_D+\varepsilon$.
The kernel of $H_2(U_P)\to H_2(X)$ is in the radical
of the compact intersection pairing, so this is
its rank as well. Since
$e(D)=b_0+n_D-b_1(D)$ and
$e(X)=2-2b_1(X)+b_2(X)$, Euler additivity gives
\begin{equation}
\SsixEquationAliases{\label{cdp:eq:C.D4.9}\label{cdp:eq:C.D4.10}}
\begin{split}
 b_1(U_P)&=b_1(X)+\varepsilon,\\
 \operatorname{rank}I_{U_P}
 &=e(U_P)+b_0-2+2b_1(X)-b_1(D)+\varepsilon.
\end{split}
\hypertarget{tag-00F5}{}\hypertarget{tag-00E2}{}\hypertarget{tag-00EK}{}
\label{cdp:eq:C.DD.9}
\end{equation}
For connected $B$, take $b_0=1$; for independent
branch classes also take $\varepsilon=0$.
The two rows of \eqref{cdp:eq:C.DD.9}
give these connected formulas.

\item Start with the normalizations of the components of $B$
and attach each connected exceptional divisor over its
singular point. The attaching-branch graph has
$\nu-c+b_0$ cycles. The total component genus contributes
$2g_B$, and the exceptional divisors contribute at most
$\lambda$. Subsequent boundary blowups subdivide edges
or add trees. Finally $U_P\to W\setminus B$ is an
unramified triple cover, $e(W)=4$ and
$e(B)=2c-2g_B-\nu$. Hence
\begin{equation}
 b_1(D)\le2g_B+\nu-c+b_0+\lambda,\qquad
 e(U_P)=3(4-2c+2g_B+\nu).
\hypertarget{tag-00EM}{}\hypertarget{tag-00E4}{}
\SsixEquationAliases{\label{cdp:eq:C.D4.12}}\label{cdp:eq:C.DD.10}
\end{equation}
The exceptional first Betti numbers equal those of the
corresponding links: the boundary sequence of a plumbed
neighbourhood identifies the intervening second-homology
map with its invertible exceptional intersection matrix.

For later use, each reduced exceptional divisor $E$
satisfies $b_1(E)\le2p_g(X_0,x)$.
To verify this, restrict a good resolution to a Stein
representative. The finite-length stalk
$M=(R^1\rho_*\OO_X)_x$ equals its completion.
The ideals $I=\OO_X(-E)$ and
$J=\mathfrak m_x\OO_X$ have the same radical and
$I^a\subset J\subset I$ for some $a>0$, since $E$
is compact. The containments
$I^{an}\subset J^n\subset I^n$ identify their
cohomology inverse limits. Analytic formal functions
identifies this limit with
$\varprojlim H^1(nE,\OO_{nE})$
\cite[\S6, Haupts\"atze I and IIa, p.~59]{Grauert1960D},
with \cite[pp.~35--36]{Grauert1963ErratumD}.
The kernels of successive thickening restrictions
have one-dimensional support, hence no $H^2$;
all restrictions on $H^1$ are surjective.
Successively lifting a class gives a surjection
$M\to H^1(E,\OO_E)$.
For the connected SNC curve $E$, its normalization
sequence gives $h^1(\OO_E)=\sum g_i+\beta$ and
$b_1(E)=2\sum g_i+\beta$, where $\beta$ is the
rank of its incidence graph. This proves the bound.


\item In the reduced case, finite-cover duality and
\eqref{cdp:eq:C.DD.2} give
$\omega_{X_0}=g_0^*(\omega_W\otimes L^2)\simeq\OO_{X_0}$.
The character algebra and \eqref{cdp:eq:C.DD.7} give
$\chi(\OO_{X_0})=2$. Put $M=R^1\rho_*\OO_X$,
a finite-length sheaf of total length $P_g$.
Grauert--Riemenschneider vanishing applied to an
ample line on the projective surface $X_0$, followed
by projection formula, gives $R^i\rho_*\omega_X=0$
for $i>0$ \cite[Satz~2.3, pp.~274--276]{GrauertRiemenschneider1970}.
Proper duality for $\rho$ and local duality on this
Gorenstein surface, applied to
$\OO_{X_0}\to R\rho_*\OO_X\to M[-1]$, give
\[
 0\longrightarrow\rho_*\omega_X\longrightarrow\omega_{X_0}
 \longrightarrow\mathcal{E}xt^2(M,\omega_{X_0})
 \longrightarrow0
\]
(\cite[Tags~0AU3 and~0A81]{StacksProject}).
The first map restricts forms on the common smooth
locus; the last is the connecting map of the dualized
triangle. The last sheaf is the finite-length Matlis
dual of $M$, with the same local lengths.
Trivializing $\omega_{X_0}$ therefore identifies
$\rho_*\omega_X$ with an ideal of colength $P_g$.
For $P_g>0$ no nonzero constant belongs to this ideal
at a point of positive colength, so $h^0(X,\omega_X)=0$.
For $P_g=0$ it equals $\OO_{X_0}$.
Leray and Serre duality give
\begin{equation}
\begin{gathered}
 \chi(\OO_X)=2-P_g,\qquad q(X)=P_g-1\ (P_g>0),\\
 q(X)=0\ (P_g=0),\qquad q(X)=\max(P_g-1,0).
\end{gathered}
\hypertarget{tag-00YL}{}\hypertarget{tag-00E8}{}\SsixEquationAliases{\label{cdp:eq:C.D4.16}}\label{cdp:eq:C.RD.30}
\end{equation}

\item At the fixed two-branch contact, choose local
coordinates with exact equations
\begin{equation}
 F_0=(s=0),\qquad R=(s-u^k=0),\qquad
 w^3=s^2(s-u^k).
\hypertarget{tag-00EU}{}\hypertarget{tag-00EB}{}\SsixEquationAliases{\label{cdp:eq:C.D4.19}}\label{cdp:eq:C.DD.18}
\end{equation}
Absorb the remaining unit into a holomorphic cube root
of $w$. Blow up the common point $k$ times until the
two strict branches separate. At each step their
weighted multiplicity is $2+1=3$, and normalization
makes the exceptional curve unbranched. The new root
line and weight-two divisor are
\begin{equation}
 L'=\sigma^*L(-E_0),\qquad F'=\sigma^*F-E_0.
\hypertarget{tag-00EV}{}\label{cdp:eq:C.DD.19}
\end{equation}
The chart $s=u s_1$, $w=u w_1$ gives
$w_1^3=s_1^2(s_1-u^{k-1})$ and
$w^2/s=u(w_1^2/s_1)$. The other blowup chart
contains no meeting of the two strict branch curves.
Thus both nontrivial characters have the exact maps
\begin{equation}
 (L')^{-1}=\sigma^*(L^{-1})(E_0),\qquad
 (L')^{-2}(F')=\sigma^*(L^{-2}(F))(E_0).
\hypertarget{tag-00EX}{}\label{cdp:eq:C.DD.20}
\end{equation}
For this point blowup,
$\sigma_*\OO(E_0)=\OO$ and $R^1\sigma_*\OO(E_0)=0$:
extension across the centre gives the first, and
$0\to\OO\to\OO(E_0)\to\OO_{\mathbb P^1}(-1)\to0$
together with the structure-sheaf blowup images
gives the second. Projection formula makes the
positive direct image of each character module zero.

The normalized covers over the blown-up bases map
properly and birationally to $X_0$ through the same
function field. Finite direct image and Leray identify
the pushforward of their $R^1\OO$ with the just computed
zero character modules. Faithfulness of finite pushforward
therefore makes those $R^1\OO$ vanish.
The final normalization is smooth: a single weight-one
or weight-two branch is parametrized by $s=t^3$ and
$w=t$ or $w=t^2$, up to a unit root.
This gives a resolution with $R^1\OO=0$, proving
rationality. Its connected exceptional divisors have
$H^1(\OO)=0$ by formal functions and restriction;
their normalized components are rational and their
incidence graphs are trees. Hence their first
rational Betti numbers also vanish.
\end{enumerate}
\end{proof}


\begin{cited}[Cubic suspensions of plane curves of delta at most five]\hypertarget{tag-01CI}{}\label{cdp:st:C.2.RAT.6}
\StatementRecord{cdp:st:C.2.RAT.6}{proved}{01CI}{}
Let $r(x,y)=0$ be a reduced plane-curve germ with
$\delta(r)\le5$, and let $X_r$ be the isolated surface singularity
$z^3=r(x,y)$. Then $p_g(X_r)\le1$ and
$b_1(\operatorname{link}X_r;\Q)\le2$; the latter equals two
precisely for plane type $A_5$ or $D_4$.
For a collection of these germs with total delta at most five,
\[
 \sum_r p_g(X_r)\le1,\qquad
 \sum_r b_1(\operatorname{link}X_r;\Q)\le2.
\]
\end{cited}

\begin{proof}[Application of the normal-form and spectrum formulas]
The point-blowup formula gives $\delta\ge m(m-1)/2$ for
multiplicity $m$, so $m\le3$; also $\mu=2\delta-b+1\le10$,
where $b$ is the number of branches. Arnold's normal forms
\cite[\S1, I.0 and II$_1$, pp.~92,95; \S2, items~1--12,
pp.~100--101]{Arnold1976NormalForms}
then give
$A_1,\ldots,A_{10}$, $D_4,\ldots,D_9$, and $E_6,E_7,E_8$.
The only further form with nonzero cubic jet and $\mu\le10$
is $J_{2,0}=x^3+a x^2y^2+y^6$, $4a^3+27\ne0$.
It has three smooth branches $x=\lambda_i y^2$, each pair
of intersection multiplicity two, so its delta is six.
An analytic coordinate change $\phi$ with
$r\circ\phi=u r_0$ lifts to the isomorphism
$(X,Y,Z)\mapsto(\phi(X,Y),u(X,Y)^{1/3}Z)$ of the cubic
suspensions; its inverse uses $\phi^{-1}$ and the reciprocal
unit root. A smooth plane germ contributes zero.

The equations, weights of degree one and delta values are
\[
\begin{array}{c|c|c|c}
 \text{type}&r_0&(w_x,w_y)&\delta\\\hline
 A_n&x^2+y^{n+1}&(1/2,1/(n+1))&\lceil n/2\rceil\\
 D_n&x^2y+y^{n-1}&((n-2)/(2(n-1)),1/(n-1))&\lfloor n/2\rfloor+1\\
 E_6&x^3+y^4&(1/3,1/4)&3\\
 E_7&x^3+xy^3&(1/3,2/9)&4\\
 E_8&x^3+y^5&(1/3,1/5)&4
\end{array}
\]
These Newton-nondegenerate forms have third suspension
weight $1/3$. The adjunction formula
\cite[Theorem~2.1.1 and Corollary~2.1.3, pp.~11--13]{MerleTeissier1977}
gives
\[
 p_g(X_r)=\#\{(i,j,k)\in\Z_{\ge0}^3:
       (i+1)w_x+(j+1)w_y+(k+1)/3\le1\}.
\]
The monomial spectrum and suspension formulas
\cite[pp.~165--167]{Steenbrink1989Spectrum} give the
eigenvalue-one multiplicity of the surface monodromy by
counting plane Jacobian monomials whose shifted weights
$(i+1)w_x+(j+1)w_y$ are congruent to $1/3$ or $2/3$
modulo one
\cite[\S1, Theorem~L1 and Example~2.4(a), pp.~1--3]{Nemethi1991Join}.
The bases are
$1,y,\ldots,y^{n-1}$ for $A_n$;
$1,y,\ldots,y^{n-2},x$ for $D_n$;
$\{x^iy^j:0\le i<2,\ 0\le j<3\}$ for $E_6$;
$1,y,\ldots,y^4,x,xy$ for $E_7$; and
$\{x^iy^j:0\le i<2,\ 0\le j<4\}$ for $E_8$.
Weighted monodromy has finite order and is semisimple.
For the simply connected surface Milnor fibre, the
variation and pair sequences identify its eigenvalue-one
dimension with the first Betti number of the link.
Substitution yields
\[
\begin{array}{c|c|c}
 \text{plane type}&p_g(X_r)&b_1(\operatorname{link}X_r;\Q)\\\hline
 A_1,\ldots,A_4&0&0\\
 A_5&1&2\\
 A_6,\ldots,A_{10}&1&0\\
 D_4&1&2\\
 D_5,\ldots,D_9&1&0\\
 E_6,E_7,E_8&1&0
\end{array}
\]
Every nonzero entry costs delta at least three, so at
most one occurs under a total delta bound of five.
\end{proof}


\begin{proposition}[Connected triple covers over a rational image]\hypertarget{tag-01BW}{}\label{cdp:st:C.2.D4.1}
\StatementRecord{cdp:st:C.2.D4.1}{proved}{01BW}{}

Let $Y$ be a smooth compact complex threefold with
$b_2(Y;\Q)=0$, and let $f:Y\to B$ be a holomorphic map to a
smooth curve whose smooth nearby fibres are complex two-tori.
Fix a fibre $mS$, with $S$ reduced and irreducible, and restrict
the base to a small coordinate disc about it. Let
$H\subset S$ be a smooth rational curve along which the generic
transversal of $S$ is an ordinary triple point. Suppose that the
three generic normalization branches form one connected degree-three
cover $C\to H$, where $C$ denotes its smooth complete model.
Then $C$ is rational, and the cubic normal jet of $S$ vanishes
identically on at least one normal fibre over $H$.
In particular, a positive-genus connected degree-three cover is
impossible, and a primitive cubic normal jet is impossible for
any connected degree-three cover.
\end{proposition}

\begin{proof}

\begin{enumerate}[label=\arabic*.,leftmargin=*]
\item Restrict the compact map to the stated disc.
It is proper, with connected fibres because its
nearby two-tori and its irreducible central support
are connected. Apply \cref{lem:II:normal-cubic-compact-core}
to this same $f$, full divisor $mS$, and smooth
rational curve $H$.
Its $N,M,P,U_P$ are the present normal bundle,
normal line, cubic jet and level surface; put $U=U_P$.
The transport \eqref{cdp:eq:C.D4.2} uses the same
function $f$, with its unit and multiplicity.
Thus $b_1(U)\le3$ forces zero compact pairing,
and a spherical class of nonzero square is excluded
by the spherical clause of that lemma.

Write
\begin{equation}
N=\OO_H(a)\oplus\OO_H(-2-a),\qquad a\ge-1.\hypertarget{tag-00EE}{}\label{cdp:eq:C.D4.3}
\end{equation}
If $a>0$, the coefficient of the cube of the first fibre
coordinate vanishes, since it has degree $-3a$. The other
coordinate is then a global linear factor of the cubic.
This contradicts connectedness of the degree-three root
cover. Thus $a=-1$ or $a=0$.

\item For $a=0$, irreducibility of the generic cubic forces the
coefficient of the first coordinate cubed to be a nonzero
constant. Along the trivial summand of $N$, the equation
$P=1$ therefore has three constant sections. Each is a
compact $\PP^1$ in $U$ with normal bundle $\OO_{\PP^1}(-2)$:
the derivative in the trivial summand is nonzero, so the
implicit-function theorem identifies its normal direction
with the other summand. Its homology square is $-2$.
Transporting its tubular neighbourhood by \eqref{cdp:eq:C.D4.2} gives a
sphere of square $-2$ in a torus, whereas the torus has
$\pi_2=0$. This is impossible.

\item We are left with $N=\OO(-1)^{\oplus2}$ and
$E=\PP(N)\simeq\PP^1\times\PP^1$. Let $A$ be the common
zero divisor of the cubic coefficients on the base, of degree
$v$. Removing that scalar factor gives an integral horizontal
root divisor $R$ of bidegree $(3,3-v)$, where the first entry
is its degree over $H$. Its normalization is the given
cover $C$: at a general point the maps identify a normalization
branch with its distinct tangent line, and both complete
normal covers are the integral closure in this same function
field. Thus, writing $g=g(C)$,
\begin{equation}
0\le g\le p_a(R)=4-2v.\hypertarget{tag-00EF}{}\label{cdp:eq:C.D4.4}
\end{equation}
For a primitive cubic $v=0$. For a positive-genus cover with
$v>0$, the bound forces $v=1$. We treat these two cases.
Their total root divisors are respectively
\begin{equation}
B=R_{(3,3)},\qquad
 B=R_{(3,2)}+F_{(0,1)}.\hypertarget{tag-00EG}{}\label{cdp:eq:C.D4.5}
\end{equation}
Both are reduced divisors of class $(3,3)$; the second has
one simple vertical component. They are connected. Let
$c=1$ or $2$ be their number of components, and set
$\nu=\sum_p(r_p(B)-1)$, where $r_p(B)$ is the number of
analytic branches at a singular point. Their total delta
invariants are $4-g$ and $5-g$, respectively, so in the
cases under consideration they are at most four.

\item Apply \cref{lem:II:cyclic-cubic-boundary} with
$W=E$, $L=\OO_E(1,1)$, and this same section
$b=P$ of $L^3$. Its equation \eqref{cdp:eq:C.D4.6}
defines the normal connected cover $X_0$.
Its inverse maps $w(v)=1$ identify the complement
with this $U$ in every tautological frame.

\item Let $\rho:X\to X_0$ resolve the surface singularities and
then the reduced boundary $D=X\setminus U$. Let $Q_p$ denote
the reduced exceptional divisor above a singular point of
$X_0$. We need the following finite computation:
\begin{equation}
\sum_p p_g(X_0,p)\le1,\qquad
  \sum_p b_1(Q_p;\Q)\le2.\hypertarget{tag-00EI}{}\label{cdp:eq:C.D4.7}
\end{equation}
\cref{cdp:st:C.2.RAT.6} gives the geometric-genus
bound and the sum of the link first Betti numbers, since
the total delta here is at most four. The negative-definite
exceptional intersection matrix identifies the rational
first homology of each link with $H_1(Q_p;\Q)$: in the
homology sequence of its plumbed neighbourhood the map on
second homology is that invertible matrix. Thus it gives
exactly the two bounds in \eqref{cdp:eq:C.D4.7}.

\item The cohomology bound
\eqref{cdp:eq:C.DD.8}, for this same resolution, and
\eqref{cdp:eq:C.D4.7} give
\begin{equation}
b_1(X;\Q)=2q(X)\le2.\hypertarget{tag-00EJ}{}\label{cdp:eq:C.D4.8}
\end{equation}
For $c=1$ the branch class is the nonzero class $(3,3)$.
For $c=2$ the classes $(3,2)$ and $(0,1)$ are independent.
Thus the parameters of \cref{lem:II:cyclic-cubic-boundary} are
$b_0=1$, $r_B=c$, $\varepsilon=0$ and $g_B=g$.
Its formula \eqref{cdp:eq:C.D4.9} identifies the same
compact pairing on $U$, and its boundary estimate with
$\lambda\le2$ becomes
\begin{equation}
b_1(D)\le2g+\nu-c+3.\hypertarget{tag-00E3}{}\label{cdp:eq:C.D4.11}
\end{equation}
Its unramified-cover Euler formula is exactly
\eqref{cdp:eq:C.D4.12} for these $c,g,\nu$.
For $c=1$, \eqref{cdp:eq:C.D4.9}--\eqref{cdp:eq:C.D4.12} give intersection rank at
least $3+4g+2\nu>0$. For $c=2$ they give rank at least
$4g+2\nu-2$. In this second case $g\ge1$, and $R$ meets
$F$ with intersection number three, so $\nu\ge1$ and the
rank is again positive. Both cases have $b_1(U)\le2$.
They contradict the compact-core comparison with $T^4$
in \cref{lem:II:normal-cubic-compact-core}, applied at the first step.

\item This excludes $v=0$ for every connected triple cover and
$v=1$ for every positive-genus one. By \eqref{cdp:eq:C.D4.4} these are
all possibilities for a positive-genus cover, proving the
statement. The remaining genus-zero cases have $v=1$ or $2$.

\end{enumerate}
\end{proof}


\begin{proposition}[Double covers over a rational ordinary-triple image]\hypertarget{tag-01BX}{}\label{cdp:st:C.2.D4.2}
\StatementRecord{cdp:st:C.2.D4.2}{proved}{01BX}{}

Let $f:Y\to C_0$ have the torus-fibre and ambient hypotheses
of \cref{cdp:st:C.2.D4.1}. Let $H\subset S$ be a smooth rational curve
with generic ordinary-triple transversal. Suppose two of
the three generic normalization branches form a connected
degree-two cover $C\to H$, and the remaining branch forms
a degree-one cover. Then the smooth complete curve $C$
is rational. The assertion includes all zeros, with their
multiplicities, of the entire cubic normal tensor.

Consequently, in the irrational ruled setup of \cref{cdp:st:A.5.19},
a noncontracted coefficient-two horizontal conductor
with generic ordinary-triple transversal has smooth
elliptic image. Its normalization is smooth minimal
elliptic ruled, the horizontal cover is unramified of
degree three, and $e(S)=0$.
\end{proposition}

\begin{proof}

\begin{enumerate}[label=\arabic*.,leftmargin=*]
\item Suppose $g=g(C)\ge1$. Use the normal tensor and
normal-cone maps of \cref{lem:II:normal-cubic-compact-core},
applied to this same proper disc restriction, $mS$ and $H$.
Write
\label{layout:II:2}
\begin{equation}
\begin{gathered}
N=N_{H/Y}=\OO(a)\oplus\OO(-2-a),\quad a\ge-1,
 \\
\qquad P\in H^0(H,\operatorname{Sym}^3N^*),
 \qquad U=\{P=1\}\subset N.
\end{gathered}\hypertarget{tag-00EL}{}\label{cdp:eq:C.DD.1}
\end{equation}
The comparison embeds a compact core of $U$, including
its compact intersection pairing, into a nearby
torus. As proved there, $b_1(U)\le3$ and a nonzero
intersection pairing contradict this embedding.

\item Put $E=\PP(N)$, using lines in $N$, and let
$L=\OO_E(1)$, $\xi=c_1(L)$ and $F$ be the fibre class.
The tautological-bundle clause of
\cref{lem:II:cyclic-cubic-boundary} uses only $\deg N=-2$.
With $W=E$, $N=N_{H/Y}$ and this tautological $L$,
its identities \eqref{cdp:eq:C.DD.2} give
$\xi^2=2$, $\xi F=1$, $F^2=0$ and $K_E=-2\xi$.
Let $A$ be the common zero divisor of the coefficients
of $P$, and put $v=\deg A$. The divisor of the induced
section of $L^3$ is
\begin{equation}
B_{\rm tot}=R_1+R_2+p^*A.\hypertarget{tag-00EZ}{}\label{cdp:eq:C.DD.3}
\end{equation}
Here $R_1$ is a section and $R_2$ is an integral
degree-two curve over $H$. The section is smooth:
its rational tangent-root map extends over the smooth
complete base into the proper bundle $E$.
The normalization of $R_2$ is precisely $C$. Over the
generic point a normalization branch is identified
with its distinct tangent root; taking the integral
closure in this same quadratic function field identifies
the two complete normal covers, as in \eqref{cdp:eq:C.RD.2}.

\paragraph{All possible common divisors.}
For $a\ge0$ and $v>0$, the pure cube of the coordinate
in $\OO(a)$ has coefficient of degree $-3a-v$, and
therefore vanishes. The other coordinate is a factor.
The quadratic factor is irreducible over the function
field. Its coefficient of the first coordinate squared
must be nonzero; that coefficient has degree
\begin{equation}
2-a-v.\hypertarget{tag-00F0}{}\label{cdp:eq:C.DD.4}
\end{equation}
Thus $a+v\le2$. In particular, $v\le2$, and $v=2$
forces $a=0$.

\item For $a=-1$, identify $E$ with the product of the base
and the projective direction line. Write $d\ge0$ for
the degree of the map defining the section $R_1$.
The bidegrees, with degree over $H$ first, are
\begin{equation}
R_1:(1,d),\qquad R_2:(2,3-v-d),\qquad
 p_a(R_2)=2-v-d.\hypertarget{tag-00F1}{}\label{cdp:eq:C.DD.5}
\end{equation}
The inequality $g\le p_a(R_2)$ gives $v+d\le1$.
Consequently the only common-divisor possibilities
for our positive-genus cover are $v=0$, $v=1$, and
$v=2$ with $N=\OO\oplus\OO(-2)$.

\item For each reduced total divisor below, use
\cref{lem:II:cyclic-cubic-boundary} with
$B=B_{\rm tot}$ and the same cubic $P$.
Its cover \eqref{cdp:eq:C.DD.6} has complement $U$
under the mutually inverse evaluation maps $w(v)=1$.
Write $c$ for the number of branch components,
$b_0=b_0(B)$, $\nu=\sum_x(r_x(B)-1)$,
$r=r_B$ for the rank of their classes, and
$\varepsilon=c-r$. Its total normalization genus
is $g_B=g$, because the section and vertical fibres
are rational. Let $\lambda$ be its exceptional
first-Betti sum. Thus \eqref{cdp:eq:C.DD.7}--\eqref{cdp:eq:C.DD.10}
apply with exactly these symbols, including disconnected
boundary and dependent classes.

\paragraph{No common zero.}
For $v=0$, the divisor $B=R_1+R_2$ is reduced,
has $c=2$, and has arithmetic genus four by
\eqref{cdp:eq:C.DD.2}. Its total delta invariant is therefore
\begin{equation}
\delta(B)=5-g\le4.\hypertarget{tag-00EN}{}\label{cdp:eq:C.DD.11}
\end{equation}
\cref{cdp:st:C.2.RAT.6} and the link comparison
in \eqref{cdp:eq:C.D4.7}
give $\sum p_g\le1$ and $\lambda\le2$.
The rank $r$ is one or two, so
$\varepsilon=2-r\le1$. Thus $b_1(U)\le3$.
Substitution in \eqref{cdp:eq:C.DD.9}--\eqref{cdp:eq:C.DD.10} gives
\begin{equation}
\operatorname{rank}\langle\ ,\ \rangle_U
                        \ge4g+2\nu-r>0.\hypertarget{tag-00EO}{}\label{cdp:eq:C.DD.12}
\end{equation}
This calculation includes a disconnected $B$ and
dependent section classes. It contradicts the
normal-cone comparison.

\paragraph{One simple common zero.}
For $v=1$, one has $B=R_1+R_2+F_0$, still reduced.
It is connected because the fibre meets both
horizontal components. Here $c=3$, $b_0=1$,
$\nu\ge2$, and
$\delta(B)=6-g\le5$.
\cref{cdp:st:C.2.RAT.6} supplies
$\sum p_g\le1$ and $\lambda\le2$.
The section and the fibre have independent classes,
so $r=2$ and $\varepsilon=1$.
Equations \eqref{cdp:eq:C.DD.8}--\eqref{cdp:eq:C.DD.10} yield
\begin{equation}
b_1(U)\le3,\qquad
 \operatorname{rank}\langle\ ,\ \rangle_U
                         \ge4g+2\nu-6\ge2.\hypertarget{tag-00EP}{}\label{cdp:eq:C.DD.13}
\end{equation}
This is the same contradiction.

\paragraph{The exact form when the common divisor has degree two.}
Now $N=\OO\oplus\OO(-2)$ by \eqref{cdp:eq:C.DD.4}.
Let $x$ and $y$ be its projections, valued in
$\OO$ and $\OO(-2)$ respectively. The coefficient
of $x^2$ in the irreducible quadratic is a nonzero
constant. Completing its square is a global bundle
automorphism: the mixed coefficient is a section
of $\OO(2)$ and therefore defines a bundle map
$\OO(-2)\to\OO$. After a nonzero constant rescaling,
the tensor has the form
\begin{equation}
P=a\,y\,(x^2-b y^2),\qquad
 a\in H^0(H,\OO(2)),\quad b\in H^0(H,\OO(4)).\hypertarget{tag-00EQ}{}\label{cdp:eq:C.DD.14}
\end{equation}
Its horizontal root divisors are
$R_1=\{y=0\}$ and $R_2=\{x^2-b y^2=0\}$.
They are disjoint. The quadratic cover has branch
degree at most four. Riemann--Hurwitz and $g\ge1$
force $g=1$, branch degree four, and four distinct
simple zeros of $b$. Thus $R_2$ is already a smooth
elliptic curve. A ruling fibre meets it either at
two transverse points or at a single point of
contact two. This verifies the local models used
in both remaining cases.

\item If $p^*A=F_1+F_2$ with distinct fibres, the reduced
branch divisor has four components. Its only plane
singularities are ordinary nodes and tangencies of
order two, of types $A_1$ and $A_3$.
Their cubic suspensions have geometric genus zero
and exceptional first Betti number zero, by
\cref{cdp:st:C.2.RAT.6} and the link comparison
in \eqref{cdp:eq:C.D4.7}. Consequently $b_1(X)=0$.
The branch divisor is connected, its class rank
is two, and $4\le\nu\le6$.
Thus $\varepsilon=2$, $b_1(D)=\nu-1$, and
\begin{equation}
b_1(U)=2,\qquad
 \operatorname{rank}\langle\ ,\ \rangle_U
                       =2\nu-4\ge4.\hypertarget{tag-00ER}{}\label{cdp:eq:C.DD.15}
\end{equation}
Again the torus comparison is impossible.

\item For $p^*A=2F_0$, apply the weighted clause of
\cref{lem:II:cyclic-cubic-boundary} to
$B_{\rm tot}=R_1+R_2+2F_0$, with the same $L$ and $P$.
Its normalized character algebra and relative direct image
are \eqref{cdp:eq:C.DD.16}--\eqref{cdp:eq:C.DD.17}.
The only intersections are those of $F_0$ with one
smooth horizontal branch, with contact $k=1$ or $2$
by \eqref{cdp:eq:C.DD.14}. Take $s$ as a defining
coordinate of this $F_0$ and $u$ its contact parameter.
Then the local cover is
\eqref{cdp:eq:C.DD.18}; its unit is absorbed by the
cube-root change of $w$ with inverse multiplication
by that unit root. The successive root-line and
character maps are exactly
\eqref{cdp:eq:C.DD.19}--\eqref{cdp:eq:C.DD.20}.
The resolution in \cref{lem:II:cyclic-cubic-boundary} therefore
makes these singularities rational and their
exceptional first Betti numbers zero.

\item In our global double-fibre case it follows that
$b_1(X)=0$ and $\lambda=0$.
The comparison $U\to E\setminus B$ is still an
unramified degree-three cover, since normalization
does not change it off the branch divisor.
Each reduced branch component still has one
strict transform with ramification index three,
so the relation and intersection calculations
\eqref{cdp:eq:C.DD.9}--\eqref{cdp:eq:C.DD.10} apply to this reduced support.
It is connected, $c=3$, its class rank is two,
and $\nu=2$ or $3$. Thus
\begin{equation}
b_1(U)=1,\qquad b_1(D)=\nu,\qquad e(U)=3\nu,
 \qquad\operatorname{rank}\langle\ ,\ \rangle_U
                                      =2\nu>0.\hypertarget{tag-00EY}{}\label{cdp:eq:C.DD.21}
\end{equation}
This contradicts the compact-core embedding. Together with
the other common divisors allowed by
\eqref{cdp:eq:C.DD.4}--\eqref{cdp:eq:C.DD.5}, it proves $g(C)=0$.

\item For the stated irrational ruled consequence,
\cref{cdp:st:A.5.19} leaves a smooth rational or smooth elliptic
image. A rational image would have horizontal
degree two or three. The preceding proof excludes
the first, and \cref{cdp:st:C.2.D4.1} excludes the second.
The elliptic conclusion, its exact minimality
comparison, and its Euler number are those already
proved in \cref{cdp:st:A.5.19}.

\end{enumerate}
\end{proof}


\begin{proposition}[Cubic root divisors and vertical attachments]
\hypertarget{tag-P380}{}\label{prop:II:cubic-root-attachments}
\StatementRecord{prop:II:cubic-root-attachments}{proved}{P380}{\cref{cdp:st:C.2.RAT.3,lem:II:normal-cubic-compact-core,lem:II:cyclic-cubic-boundary,cdp:st:C.2.D4.1,cdp:st:C.2.D4.2}}
Use the torus-fibre, rational-normalization and ruled-boundary
morphism of \cref{cdp:st:C.2.RAT.1}. Suppose the horizontal
section $C$ is not contracted, its image $H$ has generic
ordinary-triple transversal, and $C\to H^\nu$ has degree
$d\in\{2,3\}$. Then $H\simeq\PP^1$.
Put $N=N_{H/Y}$, $W=\PP(N)$ with the convention of lines,
and let $q_3$ be the cubic normal tensor. Its coefficient
line $\OO_Y(S)|_H$ is trivial and $\deg N=-2$.
Let $A$ be the common zero divisor of its coefficients,
let $B_{\rm red}$ be the reduced support of its full root
divisor on $W$, and write
\[
 \nu=\sum_x\bigl(r_x(B_{\rm red})-1\bigr),
\]
where $r_x$ counts local normalization branches. Use the
selected-tree counts and loss $\ell=K-k$ of
\cref{cdp:st:C.2.RAT.3} for this same ruling and $C\to H$.
The dominating-tree count satisfies $r_{\rm dom}\le3-d$.
The following bounds hold:
\[
\begin{array}{c|c|c|c}
d&A&\nu&\ell\\\hline
2&0&\le1&\le1\\
2&\deg A=1&\le3&\le1\\
2&p+q,\ p\ne q&\le5&\le1\\
2&2p&2&0\\
3&\deg A=1&1&0\\
3&p+q,\ p\ne q&3\text{ or }2&\le\nu-2.
\end{array}
\]
These exhaust the possible common divisors. In degree three,
$N\simeq\OO(-1)^{\oplus2}$; the last row has respectively
one or two contacts of order three between the horizontal root
curve and the two vertical fibres. For $A=0$ and a locally
unibranch root component corresponding to $C$, the sharper
bound is $\ell=0$. This applies to a degree-two root component
whose only singularities are one $A_4$ cusp or two $A_2$ cusps.

At every simple root away from $A$, the corresponding
normalization branch is smooth and has conductor ideal $(y^2)$
in the coordinates of \eqref{cdp:eq:C.RD.3}--\eqref{cdp:eq:C.RD.4}.
\end{proposition}

\begin{proof}
Riemann--Hurwitz for $C\simeq\PP^1\to H^\nu$ gives
$-2=d(2g(H^\nu)-2)+\deg\operatorname{Ram}(C/H^\nu)$,
so $g(H^\nu)=0$. By \cref{cdp:st:A.5.19}, a singular
rational ordinary-triple image has only degree-one connected
normalization covers. The degree $d\ge2$ therefore makes $H$
smooth, and $H\simeq\PP^1$. The ambient rational Chern-class
vanishing gives $\deg N=-2$ and
$\OO_Y(S)|_H\simeq\OO_H$. Put $\xi=c_1(\OO_W(1))$
and write $F$ for the class of a fibre of $W\to H$.
The tautological clause of \cref{lem:II:cyclic-cubic-boundary}
gives $\xi^2=2$, $\xi F=1$ and $K_W=-2\xi$.
Set $v=\deg A$. For $d=2$, write the horizontal root divisor
as $R_h=R_1+R_2$, where $R_2$ corresponds to $C$ and the
remaining degree-one component $R_1$ is a section. For $d=3$
the horizontal divisor is the integral curve $R_3$.

\paragraph{Root covers and simple-root attachments.}
\begin{enumerate}[label=\arabic*.,leftmargin=*]
\item Apply the root-cover clause of
\cref{lem:II:normal-cubic-compact-core} to this same full fibre,
curve $H$ and tensor $P=q_3$. Its normalized conductor cover
is the present smooth section $C$, and its horizontal root
component is the present $R_d$. Thus \eqref{cdp:eq:C.RD.2}
identifies the two covers over $H$.
At a locally unibranch point of $R_d$, normalization has one
point in its fibre. Thus two distinct points of $C$ above the
same $p\in H$ give distinct root points whenever $R_d$ is
locally unibranch along that fibre.

\item Over $H\setminus\operatorname{Supp}A$, at a simple tangent
root of $q_3$, choose coordinates
$H=(x,y)=0$, with parameter $z$ on $H$. The local equation
$f_0$ of $S$ belongs to $(x,y)^3$. The lower normal jets
vanish on the reduced smooth curve, and the graded pieces
of its regular-sequence ideal are locally free, which
justifies this membership. Choose the normal coordinates
so the simple tangent root is finite in $x/y$.
Then
\begin{equation}
f_0(y\lambda,y,z)/y^3
       =q_3(\lambda,1,z)+y\,\Phi(\lambda,y,z).\hypertarget{tag-00YK}{}\label{cdp:eq:C.RD.3}
\end{equation}
The simple root gives a holomorphic solution
$\lambda=\lambda(y,z)$ by the implicit-function theorem.
Division by the monic linear factor yields a
factorization
\begin{equation}
f_0=(x-y\lambda(y,z))\,G(x,y,z),\qquad
 G|_{x=y\lambda(y,z)}=y^2 u(y,z),\quad u(0)\ne0.\hypertarget{tag-00YT}{}\label{cdp:eq:C.RD.4}
\end{equation}
Put $\Psi(w,y,z)=q_3(w,1,z)+y\Phi(w,y,z)$. Differentiating
the factorization gives
$u(y,z)=\partial_w\Psi(\lambda(y,z),y,z)$; on $y=0$
this is the derivative of the cubic at its simple root.
The coordinate map $X=x-y\lambda(y,z)$ has inverse
$x=X+y\lambda(y,z)$, so the first factor defines a smooth
normalization branch. In the normalization product, a
tuple supported on this branch belongs to the hypersurface
ring exactly when its value is divisible by $G$ modulo
$X$: a lift zero on the other reduced factor is a multiple
of $G$. Its conductor ideal is consequently
$(G\bmod X)=(y^2u)=(y^2)$. Thus the conductor on this
branch is twice the curve lying over $H$, and the minimal
resolution $\pi$ is an isomorphism nearby. Every vertical
boundary attachment therefore lies over a multiple root.

\item Suppose $A=0$ and $R_d$ is locally unibranch everywhere.
Two distinct $C$-points over $p$ cannot both receive a
vertical boundary attachment. Their root points are
distinct by \eqref{cdp:eq:C.RD.2}. If both could receive attachments,
\eqref{cdp:eq:C.RD.4} would make both roots multiple, requiring total
cubic degree at least four. Thus each point of $H$
receives at most one attachment from $C$.

\item At most $3-d$ generic normalization branches belong to
vertical strict conductor components, so $r_{\rm dom}\le3-d$.
Under the hypotheses $A=0$ and $R_d$ locally unibranch, apply \cref{cdp:st:C.2.RAT.3} to the same
$Z\to T\to S$, its ruling and its $C\to H$.
The preceding simple-root argument proves that
the selected attachment points have distinct images
on $H=H^\nu$; hence its loss $\ell$ is zero.
Its exact attachment inequality gives
\begin{equation}
 e(S)\ge2-r_{\rm dom}\ge d-1>0.
 \hypertarget{tag-00XZ}{}\label{cdp:eq:C.RD.1}
\end{equation} An $A_2$ cusp has equation
$v^2=u^3$, and an $A_4$ cusp has equation $v^2=u^5$.
Their normalization maps are
$t\mapsto(t^2,t^{2k+1})$, for $k=1,2$, with inverse
in the function field $t=v/u^k$. The integral equation
$t^2=u$ gives the finite normalization, with a single
point above the cusp. Applying this at each cusp proves
the final examples.

\end{enumerate}

\paragraph{The degree of the common divisor.}
For $d=2$ we have $v\le2$. Write
$N=\OO(a)\oplus\OO(-2-a)$ with $a\ge-1$.
For $a\ge0$ and $v>0$, the coefficient of the cube
of the first normal coordinate has degree $-3a-v$
and vanishes. Thus the other coordinate is a factor.
The remaining quadratic is irreducible over the
function field, because it describes the connected
degree-two cover. Its first-coordinate-square
coefficient is nonzero: otherwise the quadratic would
also have the other coordinate as a factor.
That coefficient has degree $2-a-v$.
This gives $a+v\le2$.
For $a=-1$, tensoring each line in $N$ with $\OO_H(1)$
identifies $W$ with $\PP^1\times\PP^1$; the inverse
tensors with $\OO_H(-1)$. Under these maps $\xi$
has bidegree $(1,1)$. The section
$R_1$ and quadratic component have bidegrees
$(1,b)$ and $(2,3-v-b)$, with $b\ge0$ and the
first entry denoting degree over $H$. An integral
quadratic component requires $3-v-b\ge1$:
a divisor of bidegree $(2,0)$ is a union of
constant sections over $\C$. Thus $v\le2$ here
as well.
For $d=3$, \cref{cdp:st:C.2.D4.1} gives
$N=\OO(-1)^{\oplus2}$ and $v\in\{1,2\}$ for the same
normal tensor and connected normalization cover.

\paragraph{A primitive degree-two root component.}
\begin{enumerate}[label=\arabic*.,leftmargin=*]
\item Suppose $d=2$ and $A=0$.
The cubic is a section of $\OO_W(3)$.
The assumed absence of a common coefficient zero makes
its root divisor finite of degree three over $H$:
\begin{equation}
R=R_1+R_2\sim3\xi,\qquad
 \xi^2=2,\qquad K_W=-2\xi,\qquad p_a(R)=4.\hypertarget{tag-00Y0}{}\label{cdp:eq:C.RD.11}
\end{equation}
Here $R_1$ is a section and hence is smooth rational;
$R_2$ has normalization $C\simeq\PP^1$, by the exact
integral-closure comparison \eqref{cdp:eq:C.RD.2}. Both are reduced
and they are different. The normalization sequence gives
\begin{equation}
\sum_p\delta_p(R)=5.\hypertarget{tag-00Y1}{}\label{cdp:eq:C.RD.12}
\end{equation}
This equality also holds if the two components are
disjoint: the arithmetic-genus convention is
$p_a(R)=1-\chi(\OO_R)$ and the normalization has Euler
characteristic two.

\item Apply \cref{lem:II:cyclic-cubic-boundary}
with the present $N,W$, $L=\OO_W(1)$,
$P=q_3$, and reduced divisor $B=R_1+R_2$.
Its finite cover is $p_0:X_0\to W$,
its resolution is $X$, and its complement is
$U=\{q_3=1\}$ under the inverse maps $w(v)=1$.
Its three character vanishings are
\eqref{cdp:eq:C.RD.13}; hence
$H^1(X_0,\OO_{X_0})=0$ for this splitting of $N$.
By \eqref{cdp:eq:C.DD.8} and
\cref{cdp:st:C.2.RAT.6}, using the delta-five
identity \eqref{cdp:eq:C.RD.12},
$q(X)\le\sum p_g(X_0,x)\le1$, so $b_1(X)\le2$.

\item Write
\[
 \nu=\sum_p(r_p(R)-1),\qquad
 c=b_0(R),\qquad
 \lambda=\sum_x b_1(\operatorname{link}(X_0,x))\le2.
\]
Normalization of $R$ gives $e(R)=4-\nu$.
Since $U\to W\setminus R$ is \'etale of degree three,
$e(U)=3\nu$. The resolved boundary has the same
number $c$ of connected components. Replacing each
singular point of $R$ by the connected exceptional
divisor of its cubic suspension gives a graph with two
strict-component vertices and one exceptional vertex per
singular point. If there are $s$ such points, it has
$\sum_p r_p(R)=\nu+s$ attachment edges, $2+s$ vertices
and $c$ connected components, hence cycle rank $\nu-2+c$.
Adding the first homology of the connected exceptional
divisors gives
\begin{equation}
b_1(D)=\nu-2+c+\lambda.\hypertarget{tag-00Y3}{}\label{cdp:eq:C.RD.14}
\end{equation}
The exceptional first Betti number is the link's
first Betti number: the exceptional intersection
matrix in its plumbed neighbourhood is invertible
over $\Q$. This is precisely the local comparison
proved after \eqref{cdp:eq:C.D4.7}. Additional boundary-resolution
blowups add rational trees and leave \eqref{cdp:eq:C.RD.14}
unchanged.

\item In the notation of
\cref{lem:II:cyclic-cubic-boundary}, take
$c_{\rm lemma}=2$, $b_0=c$, $g_B=0$, and
$r_B=\dim\operatorname{span}\{[R_1],[R_2]\}$.
Thus
\[
 \varepsilon=2-\dim\operatorname{span}\{[R_1],[R_2]\}
       \in\{0,1\}.
\]
Its relation maps use precisely
$g^*R_i=3B_i+\sum_j a_{ji}E_j$, where
$g=p_0\rho$ and $B_i$ is the strict ramification curve.
Its pair sequence therefore gives
$b_1(U)=b_1(X)+\varepsilon\le3$ and the same compact
intersection pairing. Let $d_D$ be the number of
components of $D$.
Now use $e(X)=2-2b_1(X)+b_2(X)$ and
$e(D)=c-b_1(D)+d_D$. Equations \eqref{cdp:eq:C.RD.14} and
$e(U)=3\nu$ yield the exact identity
\begin{equation}
\begin{split}
 \operatorname{rank}I_U
   &=e(U)-2+2b_1(X)+c-b_1(D)+\varepsilon\\
   &=2\nu+2b_1(X)-\lambda+\varepsilon
     \ge2\nu-2.
\end{split}\hypertarget{tag-00Y4}{}\label{cdp:eq:C.RD.15}
\end{equation}
This formula applies to both connected and disjoint
root components.

\item If $\nu\ge2$, \eqref{cdp:eq:C.RD.15} gives a nonzero compact
intersection pairing on a surface with $b_1(U)\le3$.
Apply \cref{lem:II:normal-cubic-compact-core}
to the same $f,mS,H$ and cubic $q_3$.
Its level is this $U$ and its transport preserves
this compact pairing. Since $b_1(U)\le3$, its
vanishing conclusion contradicts the positive rank.
This contradiction proves $\nu\le1$.

\item Finally suppose a value of $H$ receives vertical
boundary attachments at both its $C$-preimages.
If their tangent-root points in $R_2$ were distinct,
at least one would be a simple root of the cubic,
and \eqref{cdp:eq:C.RD.3}--\eqref{cdp:eq:C.RD.4}
would give conductor ideal $(y^2)$ on that smooth
normalization branch, excluding a vertical attachment. They therefore map to a
single multibranch point of $R_2$. Distinct values
give distinct such points, each contributing at
least one to $\nu$. Hence there is at most one
value with two attachments. The remaining generic triple branch gives
$r_{\rm dom}\le1$. In \cref{cdp:st:C.2.RAT.3}
use the same ruling, map $C\to H$, and selected
vertical boundaries. Their loss $\ell$ is at most
one by the just proved bound. The exact attachment
inequality, with $K\ge r-N_0-r_{\rm dom}$, gives
\begin{equation}
e(S)\ge2-r+N_0+(K-1)
                         \ge1-r_{\rm dom}\ge0.\hypertarget{tag-00Y5}{}\label{cdp:eq:C.RD.16}
\end{equation}

\end{enumerate}

\paragraph{A degree-two component and one simple common zero.}
\begin{enumerate}[label=\arabic*.,leftmargin=*]
\item Suppose $d=2$ and $\deg A=1$.
Write the complete root divisor as
\begin{equation}
B=R_1+R_2+F\sim3\xi,
 \qquad B_h=R_1+R_2\sim3\xi-F.\hypertarget{tag-00Y6}{}\label{cdp:eq:C.RD.17}
\end{equation}
Here $F$ is the ruling fibre over the unique simple
common zero. The horizontal divisor $B_h$ is finite
over $H$, $R_1$ is a section, and the normalization of
$R_2$ is the rational degree-two cover $C$.
Since $K_W=-2\xi$, $\xi^2=2$, and $\xi F=1$,
\begin{equation}
p_a(B_h)=2,
 \qquad \sum_p\delta_p(B_h)=3,
 \qquad I(F,B_h)=3.\hypertarget{tag-00Y7}{}\label{cdp:eq:C.RD.18}
\end{equation}
The full divisor $B$ has total delta six. The local
bounds follow from its decomposition \eqref{cdp:eq:C.RD.17}.

\paragraph{The local cubic-suspension genus along $F$.}
At a point of $F$, put $i=I_p(F,B_h)$. If $i=1$,
the germ of $B$ is an ordinary node. If $i=2$,
the horizontal germ is either smooth with contact
two with $F$, two smooth branches transverse to
$F$, or a multiplicity-two cusp transverse to $F$.
The bound $\delta_p(B_h)\le3$ makes the resulting
plane types respectively
\begin{equation}
A_3,\qquad D_{2k+2}\ (1\le k\le3),\qquad
                 D_{2k+3}\ (1\le k\le3).\hypertarget{tag-00Y8}{}\label{cdp:eq:C.RD.19}
\end{equation}
For two smooth branches $k$ is their contact;
for the cusp $k$ is its delta invariant.
To identify these germs, write the quadratic horizontal
factor as $x^2+b(t)x+c(t)$, where $F=(t=0)$.
Completing the square gives $X=x+b(t)/2$ and
$X^2-t^\mu u(t)$ with $u(0)\ne0$. The inverse
change is $x=X-b(t)/2$. Now put
$X=x\sqrt{u(t)}$, with inverse $x=X/\sqrt{u(t)}$;
this gives $t(x^2-t^\mu)$ up to a unit.
For $\mu=1$, the inverse coordinates
$x=\sqrt2 X$, $t=T+X^2$ give $X^4-T^2$.
For $\mu\ge2$ this is $D_{\mu+2}$, with even
$\mu=2k$ in the two-branch case and odd
$\mu=2k+1$ in the cusp case. A change
$\tau=\zeta t$, $\zeta^\mu=-1$, matches the
plus-sign convention up to a nonzero constant.
For an equation unit $u$, the inverse maps
$z'=z/u^{1/3}$ and $z=u^{1/3}z'$ identify the
cubic suspensions of $ur$ and $r$.
\cref{cdp:st:C.2.RAT.6} gives geometric
genus zero or one for the types in \eqref{cdp:eq:C.RD.19}.

\item If $i=3$, the point contains $R_1$ and all the
intersection of $R_2$ with $F$. The section $R_1$
is transverse to $F$. A smooth branch of $R_2$
with intersection two is tangent to $F$ and
transverse to $R_1$. In coordinates $F=(t=0)$,
$R_1=(x=0)$ write it as $t=\phi(x)$ with
$\operatorname{ord}\phi=2$. Set
$u=x\sqrt{\phi(x)/(2x^2)}$ and $v=t-\phi(x)/2$.
Let $x=\psi(u)$ be the inverse of the first change.
The full inverse is $x=\psi(u)$, $t=v+u^2$.
The branches become $u=0$ and $v=\pm u^2$,
giving $u(v^2-u^4)$, of type $D_6$. If $R_2$ has
two local branches, both are smooth and transverse
to $F$. The horizontal delta bound three forces
all three horizontal branches to be pairwise
transverse. The full germ is therefore four
distinct transverse smooth branches.
The last possibility is one singular branch
of $R_2$. Its multiplicity is two, since its
intersection with $F$ is two. Its intersection
with $R_1$ is at least two; hence \eqref{cdp:eq:C.RD.18}
forces it to be an ordinary cusp, transverse
to both smooth branches $F$ and $R_1$.
These alternatives exhaust $i=3$.

\item For four transverse branches, choose linear axes
transverse to all four tangents. The degree-four
initial plane polynomial then has four distinct
linear factors and is convenient. Its cubic suspension is Newton
nondegenerate with weights $(1/4,1/4,1/3)$;
the only positive lattice point on or below
its compact Newton face is $(1,1,1)$.
The local duality map in \cref{lem:II:cyclic-cubic-boundary}
identifies the colength of the adjunction ideal
$\rho_*\omega_X\subset\omega_{X_0}$ with $p_g$ at
this same isolated hypersurface singularity.
Thus its geometric genus is one by
\cite[Theorem 2.1.1 and Corollary 2.1.3, pp.~11--13]{MerleTeissier1977}.

\item For the cusp and two transverse branches, choose
coordinates $F=(t=0)$ and $R_1=(s=0)$. Write the
cusp tangent as $\alpha s+\beta t=0$, where
$\alpha\beta\ne0$. The inverse linear changes
$u=\alpha s+\beta t$, $v=\beta t$ and
$s=(u-v)/\alpha$, $t=v/\beta$ give the two
smooth branches $v=0$, $u-v=0$ and cusp tangent
$u=0$. Up to units, the equation is
\begin{equation}
r(u,v)=v(u-v)\bigl(u^2+c v^3+
           u\,O((u,v)^2)+O((u,v)^4)\bigr),
                  \qquad c\ne0.\hypertarget{tag-00YA}{}\label{cdp:eq:C.RD.20}
\end{equation}
The restriction of the cubic term to the tangent
$u=0$ is $cv^3$; its nonzero coefficient is the
ordinary-cusp condition. The compact Newton facets of
$z^3-r(u,v)$ have normalized inward vectors
\begin{equation}
(1/4,1/4,1/3),\qquad
             (3/10,1/5,1/3).\hypertarget{tag-00YB}{}\label{cdp:eq:C.RD.21}
\end{equation}
Their plane edges have vertex sets
$\{(3,1),(2,2)\}$ and
$\{(2,2),(0,5)\}$. For the two weights, terms from
$v(u-v)u\,O((u,v)^2)$ have weight at least
$5/4$ and $11/10$, respectively; those from
$v(u-v)O((u,v)^4)$ have weight at least
$3/2$ and $6/5$. Thus all higher terms in
\eqref{cdp:eq:C.RD.20} lie strictly above both faces.
The first edge polynomial is a nonzero multiple
of $u^2v(u-v)$; the second is a nonzero multiple
of $v^2(u^2+c v^3)$. Both have smooth zero sets
in $(\C^*)^2$. Each two-dimensional compact
face additionally contains the term $z^3$,
whose $z$-derivative is nonzero on $(\C^*)^3$.
This proves the required Newton nondegeneracy.
In exponent coordinates $(i,j,k)$, the noncompact
face inequality beyond coordinate positivity is
$j+k/3\ge1$. Every positive integral point has
$j+k/3\ge4/3$, so it satisfies this inequality
strictly.
For either vector in \eqref{cdp:eq:C.RD.21}, the point
$(1,1,1)$ has weight $5/6$, and increasing any
one coordinate makes its weight greater than
one. Thus it is the only positive lattice
point outside the interior of the Newton
polyhedron. The same adjunction theorem gives
$p_g(z^3-r)=1$.

\item Every suspension singularity over $F$ therefore
has geometric genus at most one.
A nonzero value requires $i\ge2$ and
$\delta_p(B_h)\ge1$: the cases with horizontal
delta zero were the ordinary node and $A_3$.
Since the total intersection with $F$ is three,
there is at most one such nonzero value over
$F$. Away from $F$, \cref{cdp:st:C.2.RAT.6} shows that a
nonzero geometric genus costs horizontal delta
at least three. It cannot coexist with a
nonzero value over $F$, by \eqref{cdp:eq:C.RD.18}.
Consequently the cyclic cover $X_0$ branched
over $B$ satisfies
\begin{equation}
\sum_x p_g(X_0,x)\le1.\hypertarget{tag-00YC}{}\label{cdp:eq:C.RD.22}
\end{equation}

\item Apply the exceptional-divisor bound in
\cref{lem:II:cyclic-cubic-boundary} to each cubic-suspension singularity of $X_0$ and its good
resolution. Its $E$ is the corresponding reduced
exceptional divisor, and the same plumbed-neighbourhood
map identifies $H_1(E;\Q)$ with link homology.
Thus $b_1(\operatorname{link}(X_0,x))\le2p_g(X_0,x)$.
Equation \eqref{cdp:eq:C.RD.22} yields
\begin{equation}
\lambda:=\sum_x b_1(\operatorname{link}(X_0,x))
                            \le2.\hypertarget{tag-00YD}{}\label{cdp:eq:C.RD.23}
\end{equation}

\item Apply \cref{lem:II:cyclic-cubic-boundary}
to the same $W,N,L=\OO_W(1),P=q_3$ and reduced
divisor $B=R_1+R_2+F$. Its evaluation maps identify
the complement with $U=\{q_3=1\}$.
The parameters are $c=3$, $b_0=1$, $g_B=0$,
$r_B=2$ and $\varepsilon=1$: the section $R_1$
and fibre $F$ have independent classes, and the
third is in their two-dimensional span.
Together with \eqref{cdp:eq:C.RD.22}, its cohomology
formula gives $b_1(X)\le2$ and $b_1(U)\le3$.
Put $\nu=\sum_p(r_p(B)-1)$.
Normalization of its three rational components
and the resolved-boundary comparison give
\begin{equation}
e(U)=3(\nu-2),\qquad
 b_1(D)=\nu-2+\lambda,
 \qquad
 \operatorname{rank}I_U
       =2\nu-4+2b_1(X)-\lambda\ge2\nu-6.\hypertarget{tag-00YE}{}\label{cdp:eq:C.RD.24}
\end{equation}
This is \eqref{cdp:eq:C.DD.9} with the
displayed parameter substitution and
\eqref{cdp:eq:C.RD.23}. For $\nu\ge4$ it has positive
rank, contrary to \cref{lem:II:normal-cubic-compact-core}
for this same $q_3=1$ and full fibre $mS$.
This proves $\nu\le3$.

\item Let $p\in H$ be the common zero value and
$a=\#(C\to H)^{-1}(p)\in\{1,2\}$.
At every point of $F\cap B_h$, the full divisor
has the branch of $F$ together with the
individual horizontal normalization branches.
It follows exactly that
\begin{equation}
\nu_F:=\sum_{x\in F}(r_x(B)-1)=1+a.\hypertarget{tag-00YF}{}\label{cdp:eq:C.RD.25}
\end{equation}
The contribution one is the section $R_1$;
the other $a$ branches are the normalization
points of $R_2$ above $p$. This identity counts normalization branches when
tangent roots coincide and when the cover ramifies.

\item The loss of distinct attachment values on $H$
at $p$ is at most $a-1=\nu_F-2$. Away from
$p$, the cubic is primitive. Equations
\eqref{cdp:eq:C.RD.3}--\eqref{cdp:eq:C.RD.4} show that each duplicate
attachment there gives a multibranch point of $R_2$,
costing at least one in
$\nu_{\rm out}=\nu-\nu_F$. Hence the total
loss $\ell$ of distinct normalization-incidence
edges satisfies
\begin{equation}
\ell\le\nu_F-2+\nu_{\rm out}
                              =\nu-2\le1.\hypertarget{tag-00YG}{}\label{cdp:eq:C.RD.26}
\end{equation}
The remaining generic branch gives $r_{\rm dom}\le1$.
Apply \cref{cdp:st:C.2.RAT.3} to this same ruled
boundary and selected attachments, with the loss
$\ell$ of \eqref{cdp:eq:C.RD.26}. Its inequality is
\begin{equation}
e(S)\ge2-r_{\rm dom}-\ell\ge0.\hypertarget{tag-00YH}{}\label{cdp:eq:C.RD.27}
\end{equation}

\end{enumerate}

\paragraph{A degree-two component and a common divisor of degree two.}
\begin{enumerate}[label=\arabic*.,leftmargin=*]
\item Suppose $d=2$ and $v=2$. The horizontal divisor is
\begin{equation}
B_h=R_1+R_2\sim3\xi-2F,
       \qquad p_a(B_h)=0,
       \qquad \sum_p\delta_p(B_h)=1.\hypertarget{tag-00YI}{}\label{cdp:eq:C.RD.28}
\end{equation}
Both normalization components are rational. The
normalization sequence gives the total delta in
\eqref{cdp:eq:C.RD.28}. A local quadratic equation
$x^2+b(t)x+c(t)=0$ becomes $X^2=t^\mu$ under
$X=(x+b(t)/2)/\sqrt{d(t)}$, when its discriminant
is $4t^\mu d(t)$, $d(0)\ne0$; the inverse is
$x=\sqrt{d(t)}X-b(t)/2$.
Delta one gives $\mu=2$ or $3$, respectively a node
or cusp. Thus the only possible singularity of
$B_h$ is one ordinary node or one ordinary cusp. An
intersection between $R_1$ and $R_2$, if present,
is the node. Otherwise a node or cusp belongs
to $R_2$. Every other point of the horizontal
divisor is smooth.

\paragraph{Two distinct common zero values.}
First let $A=p+q$ with $p\ne q$. The full branch
divisor is reduced:
$B=R_1+R_2+F_p+F_q\sim3\xi$.
Apply the reduced-cover construction of
\cref{lem:II:cyclic-cubic-boundary} with $P=q_3$
and $L=\OO_W(1)$. Write $p_0:X_0\to W$ for its
finite cover, whose algebra is \eqref{cdp:eq:C.DD.6}.
The local analysis of \eqref{cdp:eq:C.RD.19} and its
intersection-three cases applies at each of the two fibres. A nonzero
local geometric genus there requires positive
horizontal delta. The total horizontal delta
in \eqref{cdp:eq:C.RD.28} is one, and the two fibres are
disjoint. Away from these fibres the local
plane delta is at most one and its cubic
suspension is rational by \cref{cdp:st:C.2.RAT.6}.
Over the two fibres, the positive-genus cases in the
intersection calculations above are $D_4,D_5,D_6$;
each consumes the unique positive horizontal delta.
Their cubic suspensions have $p_g=1$, with link
first Betti number two for $D_4$ and zero for $D_5,D_6$.
All other local cases here are $A_1,A_2,A_3$ and
have both invariants zero. Therefore
\begin{equation}
P:=\sum_x p_g(X_0,x)\le1,\qquad
 \lambda:=\sum_x b_1(\operatorname{link}(X_0,x))\le2.\hypertarget{tag-00YJ}{}\label{cdp:eq:C.RD.29}
\end{equation}

\item Apply the reduced clause of
\cref{lem:II:cyclic-cubic-boundary} to this same
$W$, $L=\OO_W(1)$ and reduced total branch divisor.
Its $P_g$ is the present $P=\sum_xp_g(X_0,x)$.
Formula \eqref{cdp:eq:C.RD.30} therefore gives
$q(X)=0$, since \eqref{cdp:eq:C.RD.29} has $P\le1$.
Thus $b_1(X)=0$.

\item Put $U=X_0\setminus p_0^{-1}(B)$, choose the
resolution so that $D=X\setminus U$ has simple normal
crossings, and put $g=p_0\rho$.
The four branch classes span the two-dimensional
space generated by a section and a fibre.
Apply \cref{lem:II:cyclic-cubic-boundary}
with $c=4$, $b_0=1$, $g_B=0$ and $r_B=2$.
Its pullback and strict-coefficient relation maps
are those of this same $g=p_0\rho$ and these
four branch curves. Hence $\varepsilon=2$,
$b_1(U)=2$, and its compact-pairing identity applies.
The branch divisor and resolved boundary are connected. For
$\nu=\sum_x(r_x(B)-1)$, normalization and the
exact pairing calculation give
\begin{equation}
e(U)=3(\nu-4),\qquad b_1(D)=\nu-3+\lambda,
 \qquad \operatorname{rank}I_U=2\nu-8-\lambda.\hypertarget{tag-00YM}{}\label{cdp:eq:C.RD.31}
\end{equation}
The evaluation maps of
\cref{lem:II:cyclic-cubic-boundary} identify this
$U$ with the same normal level $\{q_3=1\}$.
Apply \cref{lem:II:normal-cubic-compact-core}
to this full fibre and $H$; since $b_1(U)=2$,
the preserved compact pairing has rank zero.
Hence
$\nu=4+\lambda/2\le5$.

\item Let $a_p,a_q\in\{1,2\}$ be the numbers of
$C$-preimages of these two common values.
The exact branch count \eqref{cdp:eq:C.RD.25} gives
$\nu_{F_p}+\nu_{F_q}=2+a_p+a_q$.
The loss of distinct attachment values on
these fibres is at most $a_p+a_q-2$.
Away from them each loss costs a multibranch
point of $R_2$. Consequently
\begin{equation}
\ell\le\nu-4\le1.\hypertarget{tag-00YN}{}\label{cdp:eq:C.RD.32}
\end{equation}
Apply \cref{cdp:st:C.2.RAT.3} to the same curve
$C$, its degree-two map to $H$, and the fixed
vertical trees. Its loss is the present $\ell$,
and the remaining generic branch gives
$r_{\rm dom}\le1$. The resulting inequality is
$e(S)\ge2-r_{\rm dom}-\ell\ge0$.

\item Now let $A=2p$. Apply the weighted clause of
\cref{lem:II:cyclic-cubic-boundary} with
$B_{\rm tot}=R_1+R_2+2F_p$ and $L=\OO_W(1)$.
Its normalized character algebra is
\eqref{cdp:eq:C.RD.33}, with the identical generator
$w^2/s$ along this $F_p$.
Its relative image \eqref{cdp:eq:C.DD.17} is
$\OO_H(-1)$, so $H^1(X_0,\OO_{X_0})=0$.

\item We check that every singularity of this
normalized cover is rational, including the
horizontal node and cusp allowed by \eqref{cdp:eq:C.RD.28}.
For this purpose write a cyclic branch divisor
as $D_1+2D_2$, with $D_1,D_2$ reduced and without
common components, and with root line $L$.
At a point blowup $\sigma$ let
$m_i=\operatorname{mult}D_i$,
$n=m_1+2m_2=3k+j$, $0\le j<3$, and let $E$
be the new exceptional curve. Keeping the
same function field gives root line
$L'=\sigma^*L(-kE)$ and new exceptional
branch weight $j$. Its two character sheaves
are exactly
\begin{equation}
\begin{split}
 (L')^{-1}&=\sigma^*(L^{-1})(kE),\\
 (L')^{-2}(D'_2)
   &=\sigma^*(L^{-2}(D_2))
       \bigl((2k-m_2+\mathbf1_{j=2})E\bigr).
\end{split}\hypertarget{tag-00YP}{}\label{cdp:eq:C.RD.34}
\end{equation}
Here $D'_2=\sigma^*D_2-m_2E+\mathbf1_{j=2}E$.
For a local exceptional equation $e$, write
$d_2=e^{m_2-\mathbf1_{j=2}}d'_2$.
The maps of normalized covers on the two
character generators are
$w_{\rm old}=e^kw_{\rm new}$ and
$w_{\rm old}^2/d_2
=e^{2k-m_2+\mathbf1_{j=2}}w_{\rm new}^2/d'_2$.
They induce the common function-field identification
away from $E$ and give \eqref{cdp:eq:C.RD.34}.

\item At a smooth horizontal branch the contact
with $F_p$ is one or two. The weight-$(2,1)$
contact calculation \eqref{cdp:eq:C.DD.18}--\eqref{cdp:eq:C.DD.20}
separates the branches with character
increments $(1,1)$ at each blowup.
If the horizontal node lies on $F_p$, there
are three smooth local branches of weights
$(2,1,1)$. They are transverse except possibly
for contact two between $F_p$ and the
quadratic branch in the case of an intersection
$R_1\cap R_2$. A blowup has $n=4,m_2=1$,
so both character increments in \eqref{cdp:eq:C.RD.34}
are one. In the transverse case it separates
the three branches. In the contact-two case
one remaining ordinary triple has the same
weights $(2,1,1)$; the same blowup separates it.
For example, write the contact-two equation as
$w^3=s^2x(s-x^2)$. The changes
$s=xu$, $w=xw_1$, then $u=xy$, $w_1=xw_2$
give $w_1^3=xu^2(u-x)$ and $w_2^3=xy^2(y-1)$.
The complementary charts give
$w_1^3=su(1-su^2)$ and $w_2^3=uy(1-y)$,
so every remaining branch crossing is transverse.

\item If the horizontal cusp lies on $F_p$, its
local form is a weight-one cusp transverse
to the weight-two fibre:
\begin{equation}
w^3=s^2(x^2-s^3).\hypertarget{tag-00YQ}{}\label{cdp:eq:C.RD.35}
\end{equation}
The first blowup has $n=4,m_2=1$ and
increments $(1,1)$. The strict cusp becomes
smooth with contact two with the new
weight-one exceptional curve. Blow up that
weight-$(1,1)$ tangency; here $n=2,m_2=0$
and the increments are $(0,1)$.
The remaining ordinary triple has weights
$(1,1,2)$, so a final blowup has increments
$(1,1)$. Explicitly, the three charts
$x=su$, $w=sw_1$; $s=uv$; and
$v=uy$, $w_1=uw_2$ give, successively,
$w_1^3=s(u^2-s)$,
$w_1^3=u^2v(u-v)$, and
$w_2^3=uy(1-y)$.
The complementary charts have equations
$w_1^3=xv^2(1-xv^3)$,
$w_1^3=s^2(sv^2-1)$, and
$w_2^3=vy^2(y-1)$.
These have smooth support or transverse crossings
near the exceptional divisor.
At a crossing of weights $a,b$
in $\{1,2\}$, normalization of
$w^3=x^a y^b$ is the quotient of $\C^2$
under the subgroup
\[
 \{(\zeta^i,\zeta^j):ai+bj=0\pmod3\},
 \qquad x=u^3,\quad y=v^3,\quad w=u^av^b.
\]
The invariant field has degree three over
$\C(x,y)$ and is generated by $w$; the invariant
ring is normal and finite over the displayed raw
ring, so this map identifies it with the integral closure.
For $a=b$ the group has weights $(1,2)$ and its
ring is $\C[u^3,v^3,uv]$, with $(uv)^3=xy$.
Here $uv=w$ for $(a,b)=(1,1)$ and
$uv=w^2/(xy)$ for $(a,b)=(2,2)$.
For $a\ne b$ it has weights $(1,1)$ and the
degree-three Veronese ring
$\C[u^3,u^2v,uv^2,v^3]$; the additional generator
is $w^2/x$ for $(2,1)$ and $w^2/y$ for $(1,2)$.
The resolutions have respectively a two-curve
$(-2)$ chain and one $(-3)$-curve.
The first is the resolution of $z^3=xy$; the second
is $\OO_{\PP^1}(-3)$, whose thickening cohomology
has terms $H^1(\PP^1,\OO(3j))=0$ for $j\ge0$.
Thus both have $R^1\OO=0$.
The horizontal node or cusp away from $F_p$
has a rational cubic suspension by
\cref{cdp:st:C.2.RAT.6}.

\item The preceding base blowups preserve the vanishing
of $R^1\OO$.
Every exceptional character increment in
\eqref{cdp:eq:C.RD.34} was zero or one. A point blowup
has $\sigma_*\OO(E)=\OO$ and
$R^1\sigma_*\OO(E)=0$, from
$\OO(E)|_E=\OO_{\PP^1}(-1)$ and the
structure-sheaf direct images; the same
holds for increment zero. Projection formula
therefore makes the first direct image of
each normalized character algebra vanish.
The character-generator maps after
\eqref{cdp:eq:C.RD.34} are regular, since their
exponents are nonnegative. They give proper
birational maps $f:X'\to X_0$ between the
normalized covers over a base blowup $\sigma$.
For their finite projections $p':X'\to W'$ and
$p_0:X_0\to W$, the equality $p_0f=\sigma p'$
and Leray give
$(p_0)_*R^1f_*\OO_{X'}=R^1\sigma_*(p'_*\OO_{X'})=0$.
Finite pushforward is faithful, so $R^1f_*\OO_{X'}=0$.
Proper birationality and normality also give
$f_*\OO_{X'}=\OO_{X_0}$; iteration uses these
identifications in Leray. Resolving the final rational
cyclic quotients preserves this property.
This proves the claimed local rationality.

\item We consequently have $b_1(X)=0$ and
$\lambda=0$. The reduced branch support
$B=R_1+R_2+F_p$ has three rational
normalization components, is connected,
and has class rank two. Thus
\begin{equation}
b_1(U)=1,\qquad e(U)=3(\nu-2),\qquad
 b_1(D)=\nu-2,\qquad
          \operatorname{rank}I_U=2\nu-4.\hypertarget{tag-00YR}{}\label{cdp:eq:C.RD.36}
\end{equation}
The same compact transport forces rank zero,
hence $\nu=2$.
Equation \eqref{cdp:eq:C.RD.25} then gives exactly one
$C$-preimage over the common value and no
multibranch point of $R_2$ away from it.
There is no loss of distinct attachment
values. Apply \cref{cdp:st:C.2.RAT.3} with this loss $\ell=0$ and $r_{\rm dom}\le1$. It gives
\begin{equation}
e(S)\ge2-r_{\rm dom}\ge1.\hypertarget{tag-00YS}{}\label{cdp:eq:C.RD.37}
\end{equation}

\end{enumerate}

\paragraph{A connected degree-three root component.}
\begin{enumerate}[label=\arabic*.,leftmargin=*]
\item Use the normal-cone construction of \cref{cdp:st:C.2.D4.1}. That result forces
$N=\OO(-1)^{\oplus2}$ and a common cubic-coefficient divisor
of degree $v=1$ or $2$. The horizontal root curve is integral
of bidegree $(3,3-v)$, with rational normalization $C$.
No vertical strict conductor component dominates $H$, since
all three generic branches are supplied by $C$.

\item Apply \cref{cdp:st:C.2.RAT.3} to this same
Hirzebruch-boundary map. Since all three generic branches
come from $C$, its $r_{\rm dom}=0$.
Let $K$ and $k$ be its selected attachment counts
before and after mapping to $H=H^\nu$, and set
$\ell=K-k$. Its exact inequality gives
\begin{equation}
e(S)=2-r+N_0+e(\Gamma_{\rm red})\ge2-\ell.\hypertarget{tag-00E5}{}\label{cdp:eq:C.D4.13}
\end{equation}

\item The cover isomorphism \eqref{cdp:eq:C.RD.2} and the simple-root
formulas \eqref{cdp:eq:C.RD.3} and~\eqref{cdp:eq:C.RD.4}
identify the normalization cover and the conductor ideal at every
simple root outside $A$. The locally-unibranch attachment argument
therefore bounds the loss outside $A$ by zero whenever the
horizontal root curve is locally unibranch.

\paragraph{One simple common zero.}
Here the total reduced branch divisor is
$B=R_{(3,2)}+F_{(0,1)}$, with two rational normalizations
and total delta five. The cubic-suspension bound of \cref{cdp:st:C.2.RAT.6}
gives $P:=\sum p_g\le1$ and
$J:=\sum b_1(Q_p)\le2$. The classes of $R,F$ are
independent. As in the cohomology bound \eqref{cdp:eq:C.D4.8},
the first-Betti and pairing identities \eqref{cdp:eq:C.DD.9},
the boundary bound \eqref{cdp:eq:C.D4.11}, and
the boundary and Euler formulas \eqref{cdp:eq:C.DD.10}, writing
$\nu=\sum(r_p(B)-1)$, we obtain
\begin{equation}
b_1(U)\le2,\qquad e(U)=3\nu,\qquad
 \operatorname{rank}(\hbox{intersection pairing})
                         =2\nu+2b_1(X)-J.\hypertarget{tag-00E6}{}\label{cdp:eq:C.D4.14}
\end{equation}
The torus comparison excludes $\nu\ge2$. Since $R\cdot F=3$,
we have $\nu\ge1$, so $\nu=1$. All branches of $R$ are
then locally unibranch: a multibranch point would contribute
to $\nu$ in addition to the intersection with $F$, or would
raise its contribution at that intersection to at least two.
There is exactly one point of $R\cap F$, and it has just
one preimage on $C$. Thus the common value also receives
at most one attachment. The preceding simple-root argument
handles every other value. Hence $\ell=0$ and \eqref{cdp:eq:C.D4.13}
gives $e(S)\ge2$.

\paragraph{Two distinct common zeros.}
Now $B=R_{(3,1)}+F_1+F_2$. The curve $R$ is smooth rational,
and each $R\cap F_i$ has total intersection number three.
A contact of order $k=1,2,3$ is the plane type $A_{2k-1}$.
Let $P$ be the number of order-three contacts. The table in \cref{cdp:st:C.2.RAT.6} gives precisely
\begin{equation}
\sum p_g=P,\qquad \sum b_1(Q_p)=2P,
                         \qquad 0\le P\le2.\hypertarget{tag-00E7}{}\label{cdp:eq:C.D4.15}
\end{equation}
Apply \cref{lem:II:cyclic-cubic-boundary} to this
reduced divisor, with $W=\mathbb P^1\times\mathbb P^1$,
$L=\OO(1,1)$ and its same cubic section.
Its $P_g$ equals this number $P$ by
\eqref{cdp:eq:C.D4.15}. Formula \eqref{cdp:eq:C.D4.16}
gives $q(X)=\max(P-1,0)$.

Modulo the exceptional classes, there is exactly one relation among
the three strict boundary classes: the pullbacks of $F_1,F_2$ agree,
and each equals three times its strict inverse image in this quotient.
The classes of $R$ and a fibre are independent. The exceptional matrix
is negative definite, so the full boundary span has dimension equal
to the number of exceptional components plus two.
Thus \eqref{cdp:eq:C.D4.10} gives $b_1(U)=2q(X)+1\le3$. With
$\nu=\#(R\cap(F_1\cup F_2))$, the boundary topology
and Euler number give
\begin{equation}
\operatorname{rank}(\hbox{intersection pairing})
          =2\nu-4-2P+4\max(P-1,0).\hypertarget{tag-00E9}{}\label{cdp:eq:C.D4.17}
\end{equation}
We check every possibility. For $P=0$ each fibre has at
least two distinct intersections, so $\nu\ge4$ and the
rank is positive, which is excluded. For $P=1$, one fibre
has a single order-three contact. The other has either
three simple intersections or one contact of order two
and one simple intersection. The first has $\nu=4$ and
positive rank; the only remaining value is $\nu=3$.
For $P=2$, each fibre has one point of intersection and
$\nu=2$. These are all the cases.

\item Since $R$ is smooth, the simple-root argument gives at
most one attachment at any noncommon value. A common
value has as many points on $C$ as the corresponding
intersection with $F_i$. In the remaining $P=1$ case
these numbers are one and two, so $\ell\le1$. In the
$P=2$ case they are both one, so $\ell=0$.
Equation \eqref{cdp:eq:C.D4.13} gives $e(S)\ge1$ and $e(S)\ge2$,
respectively.

\item For the double common zero, the divisor is
$R_{(3,1)}+2F$ on the same $W$.
Apply \cref{lem:II:cyclic-cubic-boundary} with
$L=\OO_W(1,1)$. Its character algebra
\eqref{cdp:eq:C.D4.18} is the one with local basis
$1,w,v=w^2/s$ and multiplication
$w^2=sv$, $wv=sr$, $v^2=wr$, where $s=0$
defines $F$ and $r=0$ defines $R$.
Its three character lines have bidegrees
$(0,0)$, $(-1,-1)$, $(-2,-1)$ and zero $H^1$.

The smooth curve $R$ meets $F$ with contact
$k=1,2$ or $3$. The exact local equation is
\eqref{cdp:eq:C.D4.19}, with $s$ the equation of
this $F$ and $u$ its contact parameter; a local
unit is removed by the cube-root change of $w$
and its inverse. The same successive maps
\eqref{cdp:eq:C.DD.19}--\eqref{cdp:eq:C.DD.20}
are therefore its resolution. They give rational
singularities, so $q(X)=0$ and the exceptional
first-Betti sum is zero.

\item The reduced boundary has two components with independent
classes, and contains the pullback of the ample class
$(3,2)$. If $\nu=\#(R\cap F)$, the parameter substitution
$c=2$, $b_0=1$, $g_B=0$, $r_B=2$, $\varepsilon=0$
in \eqref{cdp:eq:C.DD.9}--\eqref{cdp:eq:C.DD.10} gives
\begin{equation}
b_1(U)=0,\qquad e(U)=3\nu,\qquad
        b_1(D)=\nu-1,\qquad
        \operatorname{rank}(\hbox{intersection pairing})=2\nu>0.\hypertarget{tag-00ED}{}\label{cdp:eq:C.D4.20}
\end{equation}
This contradicts \eqref{cdp:eq:C.D4.2}, and excludes the double common
zero.

\end{enumerate}
\end{proof}


\subsection{Irreducible fibres and zero Euler number}\label{subsec:II:irreducible-sign-criteria}

\begin{proposition}[The irreducible local comparison and its rational and ruled consequences]\hypertarget{tag-01EG}{}\label{cdp:st:C.2.UN.7}
\StatementRecord{cdp:st:C.2.UN.7}{proved}{01EG}{\cref{cdp:st:D.1.9}, \eqref{cdp:eq:C.UI.2}, \eqref{cdp:eq:C.UI.3}, \cref{cdp:st:A.5.3}, \eqref{cdp:eq:C.UI.4}, \eqref{cdp:eq:C.UI.5}, \eqref{cdp:eq:C.UI.6}, \eqref{cdp:eq:C.UI.7}, \cref{cdp:st:C.2.MH.2}, \cref{cdp:st:A.5.11}, \cref{cdp:st:C.2.MH.1}, \cref{cdp:st:C.2.UN.5}, \cref{cdp:st:A.5.16}, \cref{cdp:st:A.5.17}, \eqref{cdp:eq:C.UI.12}, \eqref{cdp:eq:C.UI.11}, \cref{cdp:st:C.2.UN.1}, \cref{cdp:st:C.2.IRR.1}, \cref{cdp:st:A.5.14}, \cref{cdp:st:A.5.18}, \cref{cdp:st:A.5.20}, \cref{cdp:st:A.5.5}, \cref{cdp:st:A.5.19}, \cref{cdp:st:C.2.D4.1}, \cref{cdp:st:C.2.D4.2}, \eqref{cdp:eq:C.CR.5}, \cref{lem:II:proper-conductor-divisor,lem:II:torus-tube-numerical}, \cref{cdp:st:A.5.10}, \eqref{cdp:eq:C.UI.15}; \cite[\S2 and \S3]{ACampo1975}}

Let $f:Y\to B$ be a proper holomorphic map from a smooth
complex threefold to a smooth curve, with connected fibres
and smooth general fibre a complex two-torus. Let $mS$ be a full
fibre with irreducible reduced support, and let $F$ be a nearby
smooth fibre. Write $T_1$ for the local monodromy on
$H^1(F,\mathbb Q)$. Then
\begin{equation}
m\,e(S_{\rm reg})=\det(1-T_1^m)\ge0.\hypertarget{tag-014M}{}\label{cdp:eq:C.UI.1}
\end{equation}
For its normalization and minimal normalization resolution
$h=\eta\pi:Z\to T\to S$, the following statements hold.

\begin{enumerate}
\item If $Z$ is rational, then $e(S)\ge2$. More precisely,
      whenever $H^1(S,\mathbb Q)\ne0$, one has $e(S)=2$.
\item Suppose $Z$ is irrational ruled. A reduced horizontal
      conductor gives a smooth minimal elliptic ruled
      normalization, with disjoint smooth elliptic conductor
      curves and their finite etale double-cover maps to the
      conductor image. The support has ordinary double-curve
      normal crossings and Euler number zero.
\item Suppose the horizontal conductor is twice a noncontracted
      section. The normalization is again smooth minimal
      elliptic ruled. Its conductor is twice that section;
      its image is smooth elliptic, and the map from the
      section is etale of degree $1$, $2$, or $3$.
      The respective constant transverse types are the
      ordinary cusp, tacnode, and ordinary triple point.
      Every such support has Euler number zero.
\item Suppose the horizontal conductor is twice a section
      contracted by $\pi$. Then $e(S)\ge2$ and
      $\pi_1(S)=1$. In particular this case cannot occur
      under the full CDP hypotheses, since \cref{cdp:st:D.1.9} supplies
      a nonzero integral degree-one class on every fibre
      support there.
\end{enumerate}
The alternatives in the last three items exhaust the
irrational ruled case.
\end{proposition}

\begin{proof}

\begin{enumerate}[label=\arabic*.,leftmargin=*]
\item \textbf{Canonical and Chern classes.}
Choose a small disc $\Delta$ around the chosen value and put $U=f^{-1}(\Delta)$. Apply \cref{lem:II:torus-tube-numerical} to the full fibre $mS$ and its minimal normalization resolution. It supplies the torsion canonical and adjunction lines, $\chi(\OO_S)=0$, the two vanishing Chern classes and the degree-zero restrictions in \eqref{cdp:eq:C.UI.2}--\eqref{cdp:eq:C.UI.4}.

\item \textbf{Monodromy.}
Apply \cref{night:P4.13} to $f|_U:U\to\Delta$.
Its cohomology monodromy $T$ is this $T_1$, and
its minimum full multiplicity $m_0$ is $m$.
The resolution is chosen over $\operatorname{Sing}S$,
so the strict open stratum is $S_{\rm reg}$ and
\eqref{cdp:eq:C.UI.5} excludes exceptional contributions
at the power $m$. The local angular fibre is
\eqref{cdp:eq:C.UI.7} of \cref{cdp:st:C.2.NC.1}
with the same multiplicities $a_i$.
Thus \eqref{cdp:eq:C.UI.6} and the torus
cup-product identification in that proposition give
\eqref{cdp:eq:C.UI.8} for this same $T_1$.
Quasi-unipotence and the real/conjugate eigenvalue
products there give \eqref{cdp:eq:C.UI.1}.

\item \textbf{Finite-character comparison.}
For each nonzero class $\gamma_S\in H^1(S,\mathbb Z)$, \cref{lem:II:torus-tube-character} gives its tube extension, its nonzero restriction to the smooth fibre, and the finite augmentation local system. When $Z$ is rational or ruled, its normalization and resolution maps are exactly \eqref{cdp:eq:C.MH.10} and \eqref{cdp:eq:C.MH.12}, and \eqref{cdp:eq:C.UI.11} gives respectively $e(S)=2$ or $e(S)=2-2g$. Since $H^1(S,\mathbb Z)$ is a lattice in $H^1(S,\mathbb Q)$, the nonzero rational condition in item~1 supplies such an integral class.

\item \textbf{Rational normalization.}
Apply \cref{prop:II:rational-torus-fibre-euler} to
this same proper map, its full irreducible fibre
$mS$, and its minimal normalization resolution
$h=\eta\pi$. It proves item~1, including the
nonzero-$H^1$ refinement. The conductor and
complement maps in \eqref{cdp:eq:C.UI.12}--\eqref{cdp:eq:C.UI.13}
are exactly those of these same $\eta,\pi$.

\item \textbf{All noncontracted irrational ruled alternatives.}
On a general ruling fibre $R$, \eqref{cdp:eq:C.UI.4} gives
$D\cdot R=2$. Thus its horizontal part is reduced
of total degree two, or is $2C$ for a smooth
section $C$. Use the proof of
\cref{cdp:st:C.2.IRR.1} with this same $h,Z,D$
and ruling. Its numerical input is precisely
$c_1(K_Z+D)=0$ in \eqref{cdp:eq:C.IR.2};
\cref{lem:II:torus-tube-numerical} gives this
torsion line on the present tube.
The ruling contractions and their divisor
pullbacks are unchanged, so the blowdown identity
\eqref{cdp:eq:C.IR.4} and the resulting boundary equality
\eqref{cdp:eq:C.IR.5}
give the same disjoint elliptic horizontal curves.
The finite conductor rings, their sheet involution,
and the maps $(u,t)\mapsto(ut,u,t^2)$ are local to
these same normalization maps. Their full local
function is still a unit times the $m$th power,
so \cref{cdp:st:A.5.15,cdp:st:A.5.17} give the
same isolated-point and ramified-pinch exclusions.
Finally \cref{cdp:st:C.2.NC.1} applies to this
proper disc restriction and proves
\eqref{cdp:eq:C.IR.9}. Thus the complete blowdown,
ring and Euler comparison proves item~2, including that no horizontal component is contracted.

\begin{enumerate}[label=(\arabic*),leftmargin=1.3em,labelsep=.5em]
\item  Consider $D=2C+V$ with $C$ not contracted.
Every resolution-exceptional curve is vertical.
Indeed a positive-coefficient horizontal one
would have to be $C$. A coefficient-zero
exceptional curve $E$ satisfies
$K_Z\cdot E=D\cdot E=0$ by minimality and
adjunction. A horizontal exceptional curve over
a positive-genus base has $p_a(E)\ge1$ and
$E^2<0$, so $K_Z\cdot E>0$, a contradiction.
These are exactly the exceptional-curve
hypotheses of Propositions~\ref{cdp:st:A.5.18},
\ref{cdp:st:A.5.19} and~\ref{cdp:st:A.5.20}.

\item  Let $H$ be the reduced image of $C$ and
$d=\deg(C\to H^\nu)$. The branchwise conductor
formula is its intrinsic plane-branch conductor
plus the sum of its intersections with the
other branches. Its value on each horizontal
branch is two. If $d=1$, $H^\nu$ has the
positive genus of $C$, and no vertical rational
curve can dominate it. The sole branch with
conductor two is the ordinary cusp. If $d=2$,
the two horizontal branches must be smooth:
a singular branch would contribute at least two
intrinsically and at least one from the other
branch. Their allowed intersections give either
contact two between those branches, or simple
contact between them and with one additional
smooth branch. These are exactly a tacnode and
an ordinary triple point. For $d=3$ all three
branches must be smooth and pairwise transverse,
and there is no further branch. For $d\ge4$
the three other intersection contributions
already exceed two. This proves exhaustion
of the transverse cases used in the following
comparisons.

\item For the conductor genus bound, put
$e=h^0(S,\omega_S)\in\{0,1\}$.
The identities already supplied by
\cref{lem:II:torus-tube-numerical} give
$h^1(\OO_S)=1+e$. For this same finite normalization
the injection $H^0(\eta_*\OO_T/\OO_S)\hookrightarrow
H^1(\OO_S)$ and conductor section duality
\eqref{eq:II:conductor-section-duality}
give \eqref{cdp:eq:A.J.31}.
Restriction to each reduced image component and
degree-zero curve Riemann--Roch give
$p_a(H)\le2$, with the same formula
\eqref{cdp:eq:A.J.32} on its normalization.
These are the maps in the first two steps of
\cref{cdp:st:A.5.14}, with the torsion and
Euler identities established on this tube.
Its delta-one local ring identifications are
unchanged analytic local rings in the same smooth
threefold chart.

For the jet comparisons, take exactly
$I=\mathcal I_{H/U}$, its torsion-free normalized
conormal, its dual $N$, and
$M=\OO_U(S)|_{H^\nu}$.
The degree-zero lines $M$ and $K_U|_{H^\nu}$
come from \cref{lem:II:torus-tube-numerical}.
Thus the determinant, radical, intrinsic cubic,
quartic and conductor-clearing maps in
\cref{cdp:st:A.5.5} are defined on these identical
bundles, with the same orders along the identical
curve conductor. The cusp proof
\cref{cdp:st:A.5.20} and the tacnode proof
\cref{cdp:st:A.5.18} use these genus and degree
identities. Their local normalizations,
root-field maps and unit-preserving coordinate
changes apply to the same equations on this tube.
The rational-image group map in the latter is
the tube quotient of $\mathbb Z^4\rtimes\mathbb Z$,
which this proper family supplies.

For the ordinary-triple proof
\cref{cdp:st:A.5.19}, the same low-delta ideals,
their symbolic cubes and their conormal
determinants apply to $I$; the above genus bound
and zero ambient degrees are its global inputs.
Its singular rational-image conclusion
therefore excludes a positive-genus covering
component here as well. For a smooth rational
image, \cref{lem:II:normal-cubic-compact-core}
gives $M=\OO_H$, $\deg N=-2$ and the transport into this nearby torus.
The complete root-divisor and cover arguments
of \cref{cdp:st:C.2.D4.1,cdp:st:C.2.D4.2}
then apply to this same cubic and exclude
degrees three and two respectively.
The ruled differential comparison
\cref{cdp:st:A.5.10} uses the same map
$h^*T_U\leftarrow T_Z$; its Chern input
$c_2(h^*T_U)=0$ is also supplied by
\cref{lem:II:torus-tube-numerical}.

\item  Consequently the image in each surviving case
is smooth elliptic and the transverse analytic
type is constant. No vertical rational curve
dominates that image, so all its generic
branches belong to $C$; the local
normalizations give etale degrees $1$, $2$,
or $3$, respectively. The normal line of $C$
has degree zero. The divisor comparison
\eqref{cdp:eq:A.T.22}, together with this square-zero
section and the absence of further conductor
along it, contracts no point of $C$ and makes
the minimal ruled boundary exactly $2C$.
A first reverse blowup away from $C$ would
create coefficient $-1$; a blowup on $C$
would create a conductor component meeting
it, contrary to the constant local type.
Thus the normalization is smooth minimal
elliptic ruled, as proved for the tacnode, ordinary-triple and cusp cases in
Propositions~\ref{cdp:st:A.5.18}, \ref{cdp:st:A.5.19}
and~\ref{cdp:st:A.5.20}. The normalization Euler
comparison $e(S)=e(T)-e(C)+e(H)$ gives zero.
This proves item~3.

\end{enumerate}

\item \textbf{A contracted section and its fundamental group.}
Finally let $D=2C+V$ with $C$ contracted and
let $g\ge1$ be the ruling genus. The exact
boundary blowdown relation \eqref{cdp:eq:C.CR.5}, now
supplied by \eqref{cdp:eq:C.UI.4}, shows that $D_{\rm red}$
is the smooth section $C$ with trees of smooth
rational curves attached. All its crossings
are transverse nodes. Its inclusion in $Z$
induces an isomorphism on fundamental groups,
both identified with $\pi_1(C)$.
Apply \cref{lem:II:proper-conductor-divisor} to
$f|_U:U\to\Delta$, with its $N=T$, resolution surface
equal to $Z$, $\nu=\eta$ and $\rho=h$.
The divisor is the same $D$, defined by the pullback
of a local generator of $\omega_S$.
The torsion and numerical line in
\cref{lem:II:torus-tube-numerical} were proved separately
for this same tube and are the input to the boundary
blowdown comparison above. The support and complement
conclusions of \cref{lem:II:proper-conductor-divisor} give
\begin{equation}
D_{\rm red}=h^{-1}(\Gamma_{\rm red}),\qquad
 S_{\rm reg}\simeq Z\setminus D_{\rm red}.\hypertarget{tag-014R}{}\label{cdp:eq:C.UI.14}
\end{equation}
Every component of the connected curve
$\Gamma_{\rm red}$ is dominated by a vertical
rational curve and hence has normalization
$\mathbb P^1$. Apply \cref{lem:II:proper-conductor-pushout}
with $(X,Y,A,C,h)=(Z,S,D_{\rm red},\Gamma_{\rm red},\eta\pi)$.
Its support and complement hypotheses are
\eqref{cdp:eq:C.UI.14}; analytic triangulation
makes $D_{\rm red}\hookrightarrow Z$ a cofibration.
Its double mapping cylinder and van Kampen give
\begin{equation}
\pi_1(S)\simeq
 \pi_1(Z)*_{\pi_1(D_{\rm red})}\pi_1(\Gamma_{\rm red})
 \simeq\pi_1(\Gamma_{\rm red}).\hypertarget{tag-014S}{}\label{cdp:eq:C.UI.15}
\end{equation}
The second amalgamating map is trivial because
$C$ generates $\pi_1(D_{\rm red})$ and is
contracted. This group is free of rank
$\beta=b_1(\Gamma_{\rm red})$.

\begin{enumerate}[label=(\arabic*),leftmargin=1.3em,labelsep=.5em]
\item The quotient map in \cref{night:P4R.1}
makes this $\pi_1(S)$ virtually abelian.
Since \eqref{cdp:eq:C.UI.15} identifies it with
the free incidence-graph group, its rank
$\beta$ is at most one: a finite-index subgroup
of a free group of rank at least two contains
nonzero powers of two free generators, whose
commutator is nontrivial.
If $n\ge1$ is the number
of conductor-image components, its exact
normalization topology gives
\begin{equation}
e(\Gamma_{\rm red})=n+1-\beta\ge n\ge1,
 \qquad e(S)=e(S_{\rm reg})+e(\Gamma_{\rm red})>0.\hypertarget{tag-014T}{}\label{cdp:eq:C.UI.16}
\end{equation}
Here \eqref{cdp:eq:C.UI.1} proves the last inequality.
Parity makes $e(S)\ge2$. If $\beta=1$,
\eqref{cdp:eq:C.UI.15} supplies a nonzero integral
degree-one class on $S$; \eqref{cdp:eq:C.UI.11} then
gives $e(S)=2-2g\le0$, a contradiction.
Thus $\beta=0$, proving $\pi_1(S)=1$.
The transgression of \cref{cdp:st:D.1.9} excludes this
case under the full CDP hypotheses. This
proves item~4.

\end{enumerate}

\end{enumerate}
\end{proof}


\begin{proposition}[The unrestricted irreducible zero-Euler criterion]\hypertarget{tag-01CP}{}\label{cdp:st:C.2.UN.12}
\StatementRecord{cdp:st:C.2.UN.12}{proved}{01CP}{\cref{cdp:st:C.2.UN.5}, \cref{cdp:st:C.2.UN.7}, \cref{cdp:st:C.2.UN.19}, \cref{cdp:st:A.5.20}, \cref{cdp:st:A.5.18}, \cref{cdp:st:C.2.D4.1}, \cref{cdp:st:C.2.D4.2}, \cref{cdp:st:C.2.IRR.1}, \cref{cdp:st:C.2.VII.1}, \cref{cdp:st:C.2.VII.2}, \cref{cdp:st:A.5.17}, \eqref{cdp:eq:C.IZ.6}, \eqref{cdp:eq:C.IZ.5}, \eqref{cdp:eq:C.IZ.8}, \eqref{cdp:eq:C.IZ.7}}

Let $f:Y\to B$ be a proper holomorphic map from a smooth
complex threefold to a smooth complex curve, with connected
fibres and smooth general fibre a complex two-torus. Fix an
irreducible full fibre $mS$, and write
$h=\eta\pi:Z\to T\to S$ for normalization followed by its
minimal resolution. Every such support satisfies $e(S)\ge0$.
Its Euler number is zero precisely in the following geometric
cases, for these normalization and resolution maps.

A normal support is a smooth torus or a smooth bielliptic
surface, with multiplicities as in \cref{cdp:st:C.2.UN.5}.

A nonnormal support with irrational ruled normalization has
smooth minimal normalization ruled over an elliptic curve.
Its conductor image and reduced inverse image are disjoint
unions of smooth elliptic curves. The transversal types and
their finite conductor maps are exactly those in the table:
\begin{equation}
\begin{array}{c|c|c}
 \text{normalization conductor}&\text{transversal type}
       &\text{degree over the image}\\ \hline
 H\text{ reduced, ruling degree }2&xy=0&2\\
 2C,\ C\text{ a section}&x^2-y^3=0&1\\
 2C,\ C\text{ a section}&x^2-y^4=0&2\\
 2C,\ C\text{ a section}&xy(x-y)=0&3
\end{array}\hypertarget{tag-00ON}{}\label{cdp:eq:C.IZ.1}
\end{equation}
Every displayed map of elliptic curves is unramified; the
degree-one map is an isomorphism. The transversal type is
constant along each conductor component. In the first row,
degree two means the total degree over each image component;
its inverse image may have two connected components. There are no
contracted horizontal sections in the zero-Euler case.

For a nonnormal support whose normalization resolution
has a Hopf surface as its minimal model, put $R=D_{\rm red}$,
where $D$ is the effective
resolution conductor with $K_Z+D=h^*\omega_S$. Let $r$ be
the number of rational components of $R$. Put
$C=\Gamma_{\rm red}$ for the conductor image, let $n_0$ be
the number of its rational-normalized components, and set
$v(C)=\sum_x(s_x-1)$, using the number $s_x$ of
normalization branches at each singular point. The exact
zero criterion is
\begin{equation}
b_2(Z)=r,\qquad
 \begin{cases}
 n_0=0,\quad v(C)=0,\quad\sum_x\delta_x\le1,\quad\text{or}\\
 n_0=1,\quad v(C)=2,\quad\sum_x\delta_x=2.
 \end{cases}\hypertarget{tag-00OO}{}\label{cdp:eq:C.IZ.2}
\end{equation}
Every other normalized conductor component is elliptic.
In the second line every singular point satisfies
$\delta_x=s_x-1$. No normalization singularity lies outside
the conductor in either line. An irreducible support with
rational normalization resolution always has $e(S)\ge2$.

\end{proposition}

\begin{proof}

\begin{enumerate}[label=\arabic*.,leftmargin=*]
\item Choose a sufficiently small fibre tube $U$. By \cref{lem:II:torus-tube-numerical}, $K_U$ and $\omega_S$ are torsion, $\chi(\OO_S)=0$, and the resolution conductor satisfies $K_Z+D=h^*\omega_S$. For nonnormal $S$, this effective divisor is nonempty.

\item Every positive pluricanonical section on $Z$ would, after
taking a power trivializing that torsion line, give a
nonzero section of $\OO_Z(-aD)$ for some $a>0$.
Its inclusion in $\OO_Z$ would be a nonzero holomorphic
function on the compact connected surface vanishing along
$D$, a contradiction. Thus $\kappa(Z)=-\infty$.
The surface classification, together with the minimal-model blowdown, makes $Z$ rational, irrational
ruled, or class VII. In the last case \cref{cdp:st:C.2.UN.19} proves,
through the fixed cyclic cover and quotient maps,
that its minimal model is Hopf. Thus the formulation
above includes every class VII possibility for the
degeneration. Normal supports have already been
classified with their finite root covers in \cref{cdp:st:C.2.UN.5}.

\cref{cdp:st:C.2.UN.7} proves strict positivity for rational $Z$.
It also proves strict positivity when an irrational
horizontal section is contracted, and extends every
noncontracted irrational ruled comparison to this tube.
On this tube the inputs are the torsion canonical and normal
lines, the
vanishing of $\chi(\OO_S)$, and
$\int_Zh^*c_2(T_U)=0$, all proved there through restriction
to the nearby torus.

\item Apply \cref{cdp:st:C.2.UN.7} to the same proper
map $f$, full fibre $mS$, minimal resolution $Z$,
and its ruling. Its reduced-horizontal clause is
the first row of \eqref{cdp:eq:C.IZ.1}. Its twice-section
clause has the same branchwise conductor exponent
two and gives exactly the remaining three rows,
with the respective normalized conductor maps
of degrees one, two and three. Its contracted-section
clause has positive Euler number, so contributes
no zero case.
All these comparisons remove vertical conductor and
normalization exceptional curves. Their Euler formula is
\begin{equation}
e(S)=e(T)-e((\eta^{-1}C)_{\rm red})+e(C)=0,\hypertarget{tag-00OP}{}\label{cdp:eq:C.IZ.3}
\end{equation}
since each term vanishes. This proves both directions of
the ruled part of the criterion.

\item Apply \cref{cdp:st:C.2.VII.1} directly to the
compact class VII surface $Z$ and its effective
divisor $D$, whose canonical-boundary line is
the torsion line $h^*\omega_S$. With its
$r(R)=r$ and $R=D_{\rm red}$, it gives
\begin{equation}
e(R)=r,\qquad r\le b_2(Z).\hypertarget{tag-00OQ}{}\label{cdp:eq:C.IZ.4}
\end{equation}
Apply \cref{cdp:st:C.2.VII.2} to this same compact
hypersurface $S\subset U$, finite normalization
$\eta:T\to S$, conductor image $\Gamma$ and
class VII resolution $\pi:Z\to T$.
Its hypotheses $\omega_S$ torsion and
$\chi(\OO_S)=0$ were established in the first step.
For $C=\Gamma_{\rm red}$ it gives
\begin{equation}
\chi(\OO_C)\ge-1,\qquad e(C)\ge0 .\hypertarget{tag-00OR}{}\label{cdp:eq:C.IZ.5}
\end{equation}
Its normalized components are rational or elliptic,
as finite images of the strict conductor curves
classified by \cref{cdp:st:C.2.VII.1}.

\item Apply \cref{lem:II:proper-conductor-divisor}
to $f|_U$, with its $(X,N,T,\nu,\rho)$ equal to
$(U,T,Z,\eta,\eta\pi)$ here.
Its canonical divisor is this same $D$.
It gives $R=h^{-1}(C)$, excludes all normalization
singularities outside the conductor, and identifies
the two restrictions in
\begin{equation}
Z\setminus R\xrightarrow{\ \simeq\ }
 T\setminus(\eta^{-1}C)_{\rm red}
 \xrightarrow{\ \simeq\ }S\setminus C .\hypertarget{tag-00OS}{}\label{cdp:eq:C.IZ.6}
\end{equation}
The first arrow is $\pi$ on this open set and the
second is $\eta$ on its unit-conductor locus.

\item A class VII surface and its point blowups have
$b_1=b_3=1$, so $e(Z)=b_2(Z)$. Euler additivity in
\eqref{cdp:eq:C.IZ.6} yields the exact formula
\begin{equation}
e(S)=b_2(Z)-r+e(C),\qquad
 b_2(Z)-r\ge0,\quad e(C)\ge0.\hypertarget{tag-00OT}{}\label{cdp:eq:C.IZ.7}
\end{equation}
Apply \cref{lem:II:curve-euler-delta} to this
same reduced conductor curve $C$, with its
$r_x=s_x$ and $n_0$ as in the statement. Its
normalization sequence gives
\begin{equation}
e(C)=2n_0-v(C),\qquad
 \chi(\OO_C)=n_0-\sum_x\delta_x,\qquad
 \delta_x\ge s_x-1.\hypertarget{tag-00OU}{}\label{cdp:eq:C.IZ.8}
\end{equation}
Equations \eqref{cdp:eq:C.IZ.5} and \eqref{cdp:eq:C.IZ.8} imply $e(C)\ge n_0-1$.
Consequently $e(C)=0$ allows exactly $n_0=0,v=0$ or
$n_0=1,v=2$. In the first case the coherent bound gives
$\sum\delta_x\le1$. In the second, that bound and
$\delta_x\ge s_x-1$ give $\sum\delta_x=2$, with equality
at every singular point. Conversely either of these
two curve patterns makes $e(C)=0$ in the exact
topological formula. Equation \eqref{cdp:eq:C.IZ.7} proves precisely
\eqref{cdp:eq:C.IZ.2}, and proves the sign for every class VII
support. All three nonnormal surface classes have now
been included, which proves the assertion.

\end{enumerate}
\end{proof}


\begin{proposition}[Every irreducible class VII normalization contributes zero]\hypertarget{tag-01C7}{}\label{cdp:st:C.2.MH.4}
\StatementRecord{cdp:st:C.2.MH.4}{proved}{01C7}{\cref{cdp:st:C.2.VII.4}, \cref{cdp:st:C.2.MH.2}, \eqref{cdp:eq:C.MH.17}, \cref{cdp:st:C.2.UN.1}, \cref{cdp:st:C.2.VII.3}}

Under the full CDP torus-fibration hypotheses, an
irreducible fibre support whose normalization resolution
is a surface of class VII has
\begin{equation}
e(S)=0.\hypertarget{tag-00QH}{}\label{cdp:eq:C.MH.16}
\end{equation}
Consequently every such fibre satisfies the
equality data of \cref{cdp:st:C.2.VII.4}.
\end{proposition}

\begin{proof}

\begin{enumerate}[label=\arabic*.,leftmargin=*]
\item Under the full CDP hypotheses,
\cref{cdp:st:D.1.9} supplies a nonzero class
$\gamma_S\in H^1(S,\mathbb Z)$.
Apply \cref{lem:II:torus-tube-character} to
this same fibre, nearby torus and resolution
$h=\eta\pi$. Choose its prime $\ell$ and
augmentation local system $V$, of rank $d=\ell-1$.
Its support-duality map gives
$h^0(S,V)=h^1(S,V)=0$ and
$h^2(S,V)=h^4(S,V)$.
Its third-cohomology surjection and fourth-cohomology
isomorphism are induced by these same $\eta,\pi$.
Consequently
\begin{equation}
\begin{gathered}
 h^3(S,V)\ge h^3(Z,V),\qquad h^4(S,V)=h^4(Z,V),\\
 d\,e(S)\le 2h^4(Z,V)-h^3(Z,V).
 \end{gathered}\hypertarget{tag-00QI}{}\label{cdp:eq:C.MH.17}
\end{equation}

\item If the cyclic character restricts nontrivially to $Z$,
all its nontrivial powers remain nontrivial because
$\ell$ is prime. Poincar\'e duality on the compact
oriented surface gives $h^4(Z,V)=0$, and \eqref{cdp:eq:C.MH.17}
gives $e(S)\le0$. If its restriction is trivial,
$h^4(Z,V)=d$ and $h^3(Z,V)=d$: a class VII surface
has $b_1=b_3=1$, and point blowups preserve $b_3$.
Thus $e(S)\le1$ in this second case. \cref{cdp:st:C.2.UN.1} proves
that the Euler number of every torus-degeneration
fibre is even, so the second case also gives
$e(S)\le0$. The nonnegative lower bound of \cref{cdp:st:C.2.VII.3}
gives \eqref{cdp:eq:C.MH.16} in both cases. \cref{cdp:st:C.2.VII.4} supplies
its exact normalization-and-conductor equality data.

\end{enumerate}
\end{proof}


\subsection{Elliptic interfaces and reducible fibres}\label{subsec:II:elliptic-interface-signs}

\begin{proposition}[Elliptic images and Picard restriction]\hypertarget{tag-01C0}{}\label{cdp:st:C.2.INT.1}
\StatementRecord{cdp:st:C.2.INT.1}{proved}{01C0}{\eqref{cdp:eq:C.INT.3}, \cref{cdp:st:D.1.7}, \eqref{cdp:eq:C.INT.1}, \eqref{cdp:eq:C.INT.2}; \cite[Lemma~2.2, pp.~1664--1665; \S2, pp.~1666--1667]{BrinzanescuMoraru2005}; \cite[arXiv v2, Remark~2.7, p.~9]{DlousskyTeleman2012}}

Let $Y$ be a smooth compact complex threefold with
$H^2(Y,\Q)=0$, let $S_i$ be an irreducible component of
a connected reducible divisor $S$, and let
$h:Z\to S_i\hookrightarrow Y$ be the minimal resolution
of its normalization. Suppose that the minimal model
$\sigma:Z\to Z_0$ is a Hopf surface or a non-K\"ahler
elliptic surface. Every strict transform $E\subset Z$ of
an elliptic curve on $Z_0$ has a one-dimensional image
under $h$, and the normalization of that image is elliptic.
The restriction homomorphism
$\Pic^0(Y)\to\Pic^0(E)$ is nonconstant for each such $E$.
In particular such a component cannot meet a component
whose minimal model is an Enoki surface.
\end{proposition}

\begin{proof}

\begin{enumerate}[label=\arabic*.,leftmargin=*]
\item Choose an adjacent component $S_j$. The pullback
of its canonical divisor section gives a nonzero effective
Cartier divisor
\begin{equation}
A=h^*S_j>0,\qquad \OO_Z(A)=h^*\OO_Y(S_j).\hypertarget{tag-00NQ}{}\label{cdp:eq:C.INT.1}
\end{equation}
Its first Chern class is torsion. For some $q>0$,
$\OO_Y(qS_j)\in\Pic^0(Y)$, so the restriction homomorphism
$\rho:\Pic^0(Y)\to\Pic^0(Z)$ has in its image the line
$\OO_Z(qA)$. That line has strictly positive degree for
every Gauduchon metric on $Z$.

\item Put $A_0=\sigma_*A$.
The difference $A-\sigma^*A_0$ is exceptional and has
zero real Chern class: both $A$ and its pushforward have
zero real class. The negative definite exceptional
intersection matrix makes every coefficient of this
difference zero. Thus
\begin{equation}
A=\sigma^*A_0,\qquad A_0>0.\hypertarget{tag-00O1}{}\label{cdp:eq:C.INT.2}
\end{equation}
Nonzeroness follows by the same negative definiteness if
$A_0$ were zero. Pullback is an isomorphism
$\Pic^0(Z_0)\simeq\Pic^0(Z)$; the integral blowup
formula and $R\sigma_*\OO_Z=\OO_{Z_0}$ identify both
the lattice and the coefficient space defining it.
We therefore regard $\rho$ as a map into $\Pic^0(Z_0)$.
Its image contains $\OO_{Z_0}(qA_0)$.

\item We first prove the stronger assertion that the composite
restriction to $E$ cannot be zero. Suppose that it is zero.
This situation would in particular arise if $h(E)$ were
a point or had rational normalization.
In the latter case the map from the normal curve $E$
factors as $E\to\PP^1\to Y$. In both cases every line
from $\Pic^0(Y)$ restricts trivially to $E$. Here the
pullback to $\PP^1$ is in $\Pic^0(\PP^1)=0$; this is
the fixed map giving the vanishing. The restriction
of $\sigma$ to $E$ is an isomorphism onto its smooth
elliptic image $E_0$.

\item For a Hopf $Z_0$, apply \cref{lem:II:hopf-curve-restriction}
to this surface and its same elliptic curve $E_0$.
It gives an isomorphism on $H^1(\OO)$ induced by the inclusion.
The coefficient maps commute with this inclusion, and
\cref{cdp:st:D.1.7} makes the real coefficient map on $Z_0$
injective. Consequently the same square gives a nonzero map
\begin{equation}
H^1(Z_0,\R)\longrightarrow H^1(E_0,\R).\hypertarget{tag-00OA}{}\label{cdp:eq:C.INT.3}
\end{equation}
The differential of $\Pic^0(Z_0)\to\Pic^0(E_0)$ is that
$H^1(\OO)$ isomorphism, so its kernel has zero identity
component. Under the assumed zero restriction, the connected
image of $\rho$ lies in this kernel and is zero. This
contradicts the positive-degree line in
\eqref{cdp:eq:C.INT.1}--\eqref{cdp:eq:C.INT.2}.

\item For a non-K\"ahler elliptic $p:Z_0\to B$, every curve
on its minimal model is a reduced elliptic fibre. Indeed
\cite[Lemma~2.2, pp.~1664--1665]{BrinzanescuMoraru2005}
gives degree zero for the restriction of every divisor line to
a smooth fibre; a horizontal effective curve would have
positive intersection with that fibre. The quasibundle
description in \cite[\S2, p.~1663]{BrinzanescuMoraru2005}
makes each reduced fibre smooth elliptic.
The covering construction of
\cite[Lemma~2.2, pp.~1664--1665; \S2, pp.~1666--1667]{BrinzanescuMoraru2005}
provides a finite cyclic map $\psi:Z'_0\to Z_0$ and a
principal elliptic bundle $p':Z'_0\to B'$ such that a
reduced fibre $E'$ above $E_0$ maps to $E_0$ by a finite
isogeny. At a multiple fibre $mE_0$, the normalized base change
$t=s^m$ in \cref{def:II:elliptic-root-curve}, applied to
this same elliptic surface tube, is given by
$w^m=u$, $s=wh$ when $t=u h^m$.
Its inverse over $h\ne0$ is $w=s/h$, and its central
free translation quotient is an isogeny onto $E_0$.
A branching order divisible by $m$ is the base change of
this chart, so gives the same reduced-curve map after
normalization. At the additional ordinary branch values
in the source's construction the reduced-curve map is
an isomorphism. These are precisely the fibre maps
used by $\psi$ above.

\item For this principal bundle,
\cite[\S2, pp.~1666--1667]{BrinzanescuMoraru2005}
gives the presentation and automorphy coordinates
\begin{equation}
Z'_0=\Theta^*/\langle\tau\rangle,
 \quad |\tau|>1,\quad \deg\Theta>0,
 \qquad
 \Pic^0(Z'_0)=p'^*\Pic^0(B')\times\C^*.\hypertarget{tag-00OB}{}\label{cdp:eq:C.INT.4}
\end{equation}
In the last factor a scalar $\alpha$ is the automorphy
factor of a line bundle. Its restriction to
$E'=\C^*/\langle\tau\rangle$ is that same automorphy
line; its kernel is the discrete subgroup
$\langle\tau\rangle$. Consequently the identity
component of the restriction kernel is precisely the
compact group $p'^*\Pic^0(B')$.
Our assumed trivial restriction on $E$ makes
$\psi^*\rho(\Pic^0(Y))$ lie in that identity component.
Gauduchon degree is a continuous additive homomorphism
to $\R$, and is zero on any compact subgroup. On the
other hand this image contains
$\OO_{Z'_0}(q\psi^*A_0)$, the line of a nonzero effective
divisor, which has strictly positive Gauduchon degree.
This contradiction proves the first assertion in both
cases, and proves the stated nonconstant restriction.
Riemann--Hurwitz gives genus at most one for the
normalization of the nonconstant image of $E$, so it is
elliptic.

\item For the last assertion, an effective divisor of zero
real Chern class on a minimal class VII surface is,
by \cite[arXiv v2, Remark~2.7, p.~9]{DlousskyTeleman2012}, an elliptic
divisor on a Hopf surface or a positive multiple of the
rational cycle on an Enoki surface. To apply its integral
homological-triviality hypothesis, multiply the divisor
by an integer killing its torsion integral Chern class
and use Poincar\'e duality. Multiplication changes each
coefficient by that integer and preserves the support,
so the stated support and proportionality conclusion
holds for the original divisor. Formula \eqref{cdp:eq:C.INT.2} also applies to an Enoki
minimal model. Thus the adjacent divisor on its
resolution is the full pullback of a rational cycle;
all its curve images have rational normalization.
On a Hopf or non-K\"ahler elliptic neighbouring
resolution, the adjacent divisor is the full pullback
of a nonzero divisor consisting of elliptic curves.
At least one of their strict transforms must map
into the same intersection. It cannot map to a point
or a rational-normalized curve by the result just
proved. This is the required contradiction.

\end{enumerate}
\end{proof}


\begin{proposition}[Exclusion of Enoki components]\hypertarget{tag-01C1}{}\label{cdp:st:C.2.INT.2}
\StatementRecord{cdp:st:C.2.INT.2}{proved}{01C1}{\cref{cdp:st:C.3.4}, \cref{cdp:st:C.2.INT.1}, \cref{cdp:st:A.5.3}, \cref{cdp:st:C.2.UN.17}, \eqref{cdp:eq:C.INT.5}, \eqref{cdp:eq:C.INT.7}, \eqref{cdp:eq:C.INT.6}; \cite[arXiv v3, Definition~2.20, p.~14; Lemma~2.22 and the following paragraph, p.~15]{Dloussky2006}}

Let $f:Y\to B$ be a proper connected-fibre map from a
smooth compact complex threefold with $H^2(Y,\Q)=0$
and general fibre a complex two-torus. No component
of a reducible fibre has normalization resolution
with Enoki minimal model.
\end{proposition}

\begin{proof}

\begin{enumerate}[label=\arabic*.,leftmargin=*]
\item Write the full fibre as $F=\sum_i m_iS_i$. Its
normalization resolutions $Z_i$ are non-K\"ahler by
\cref{cdp:st:C.3.4}. A neighbouring component gives on each $Z_i$
a nonzero effective divisor $E_{ij}=h_i^*S_j$ of
zero real Chern class. On a class VII minimal model
the precise classification in \cref{cdp:st:C.2.INT.1} makes that
model Hopf or Enoki. The other non-K\"ahler minimal
models are elliptic surfaces, including both types
of Kodaira surfaces. \cref{cdp:st:C.2.INT.1} rules out an interface
between either such type and an Enoki component.
Thus connectedness of the fibre shows that, if an
Enoki component occurs, every component has Enoki
minimal model.

\item Let $D_i^{\rm self}\ge0$ be the effective resolution
conductor for $S_i$ from \cref{cdp:st:A.5.3}, and set
$D_i=D_i^{\rm self}+\sum_{j\ne i}E_{ij}$. This is
also the resolution conductor for the total reduced
support. Adjunction specifies the exact comparisons
\begin{equation}
K_{Z_i}+D_i=h_i^*(K_Y+\sum_j S_j),\qquad
 K_{Z_i}+D_i^{\rm self}=h_i^*(K_Y+S_i).\hypertarget{tag-00OC}{}\label{cdp:eq:C.INT.5}
\end{equation}
Push the effective $D_i$ to the Enoki minimal model.
A positive multiple of its canonical-boundary line is
topologically trivial. On this class VII surface the
coefficient isomorphism and flat-character map
\eqref{cdp:eq:C.CA.11} identify it with a complex-flat
line. Thus the pushed divisor is numerically
anticanonical in Definition~2.20 of
\cite[arXiv v3, p.~14]{Dloussky2006}.
Its minimal model is therefore parabolic Inoue by
\cite[arXiv v3, Definition~2.20, p.~14; Lemma~2.22 and the following paragraph, p.~15]{Dloussky2006}.
Write $E_i^0$ for its elliptic curve and $C_i$ for
its reduced rational cycle. They are disjoint and
\begin{equation}
K_{Z_i^0}+E_i^0+C_i=0.\hypertarget{tag-00OD}{}\label{cdp:eq:C.INT.6}
\end{equation}
The pushforward of $D_i$ has the form
$E_i^0+k_i' C_i$, with $k_i'\ge0$, by intersection
with $E_i^0$ and the kernel of the cycle intersection
matrix. Each nonzero adjacent divisor has the exact
form $E_{ij}=n_{ij}\sigma_i^*C_i$, $n_{ij}\ge1$.
In particular $k_i'\ge\sum_{j\ne i}n_{ij}\ge1$.

\item Apply \cref{cdp:st:C.2.UN.17} to the restriction
$f:f^{-1}(\Delta)\to\Delta$ and the same full divisor
$F=\sum_i m_iS_i$. Its canonical evaluation gives
$K_Y=\OO_Y(\sum_i k_iS_i)$ for integers $k_i\ge0$;
the line from the base is trivial on the chosen disc.
Put $r_i=(k_i+1)/m_i$. The principal full-fibre
relation yields the following equality in
$\Pic(Z_i)\otimes\Q$:
\begin{equation}
K_{Z_i}+D_i
   =\sum_{j\ne i}m_j(r_j-r_i)E_{ij}.\hypertarget{tag-00OE}{}\label{cdp:eq:C.INT.7}
\end{equation}
Indeed subtract $r_i h_i^*F=0$ from the first
formula \eqref{cdp:eq:C.INT.5}. Choose $i$ for which $r_i$ is
maximal. The negative of the right side of
\eqref{cdp:eq:C.INT.7} is an effective rational divisor $A_i$.
Thus $K_{Z_i}+D_i+A_i$ is torsion after clearing
denominators. Push to its minimal model and use
\eqref{cdp:eq:C.INT.6}. Taking Gauduchon degree gives
\begin{equation}
k_i'+\sum_{j\ne i}m_j(r_i-r_j)n_{ij}=1,\hypertarget{tag-00OF}{}\label{cdp:eq:C.INT.8}
\end{equation}
since $\deg\OO(C_i)>0$. All terms after $k_i'$
are nonnegative, and $k_i'\ge\sum n_{ij}\ge1$.
Equality forces $k_i'=1$, exactly one neighbour,
$n_{ij}=1$, and $r_j=r_i$. That neighbour is again
a maximum. Propagating along the connected
intersection graph proves that there are exactly
two components $S_1,S_2$, all ratios are equal,
and on both minimal models
\begin{equation}
(D_i^{\rm self})_0=E_i^0,
 \qquad (E_{ij})_0=C_i.\hypertarget{tag-00OG}{}\label{cdp:eq:C.INT.9}
\end{equation}

\item Every topologically
torsion line descends through a point blowup, by the
integral Picard blowup formula. Apply \cref{lem:II:canonical-boundary-blowup} at each
point blowup, with its upstairs boundary equal to
$(D_i^{\rm self})_{\rm new}$ and its downstairs boundary
the pushforward $(D_i^{\rm self})_{\rm old}$. The second
identity in \eqref{cdp:eq:C.INT.5} has zero real
canonical-boundary class, so the lemma's coefficient
formula gives
\begin{equation}
(D_i^{\rm self})_{\rm new}
   =\sigma^*(D_i^{\rm self})_{\rm old}-E_{\rm exc}.\hypertarget{tag-00NR}{}\label{cdp:eq:C.INT.10}
\end{equation}
Starting with the smooth multiplicity-one elliptic
curve $E_i^0$, effectivity permits centres only on
its strict transform, and the
exceptional coefficient is zero. Hence
$D_i^{\rm self}$ is that elliptic strict transform,
and all these centres are disjoint from $C_i$.
The external cycle $\sigma_i^*C_i=C_i$ is therefore
unchanged, and disjoint from the self conductor.

\item It follows that $S_i$ is normal near the external
intersection. Any
resolution-exceptional curve outside $D_i$ is
disjoint from it: adjunction and minimality give
$K_{Z_i}\cdot E\ge0$, while $D_i\cdot E\ge0$,
so $(K_{Z_i}+D_i)\cdot E=0$ forces $D_i\cdot E=0$.
Connected resolution fibres cannot join the
elliptic self conductor to the external rational
cycle. The conductor on the normal surface is a
reflexive ideal, by the duality calculation in
\cref{cdp:st:A.5.3}, so it has no isolated additional support
on that cycle. Thus normalization is an isomorphism
near its image. Resolution over this image is
crepant by the second formula \eqref{cdp:eq:C.INT.5}, since
$D_i^{\rm self}$ is disjoint from it. Its
exceptional curves are consequently $(-2)$
curves in $C_i$. A whole cycle cannot be contracted,
because its intersection matrix is not negative
definite. The connected fibres which are
contracted are proper chains in that cycle.

\item The intersection scheme $D=S_1\cap S_2$ is a
complete-intersection curve in $Y$. Its pullback
has reduced divisor $C_i$, so it is generically
reduced; its Cohen--Macaulay property then makes
it reduced everywhere. It is connected, since it
is the image of $C_i$. The equality of all $r_i$
means $K_Y+S_1+S_2=rF$ in $\Pic(Y)\otimes\Q$
on the fibre neighbourhood. Adjunction gives
$\omega_D$ torsion, hence $p_a(D)=1$.
Topologically the map $C_i\to D$ contracts just
the proper chains described above and is an
isomorphism elsewhere. It preserves the graph
circle and all its remaining components are
rational. Thus its normalized components are
$\PP^1$, its graph has exactly one cycle, and
$\sum\delta_x=\sum(r_x-1)$. Each contracted
chain has only two ends; every singular point
therefore has two branches and $\delta=1$.
The colength-one ring comparison \eqref{cdp:eq:C.VII.6}
proves that these are ordinary nodes. At least
one node occurs.

\item At such a node the complete-intersection curve
has embedding dimension two. At least one of
the two surface equations has a nonzero linear
part; take it to be $z$. On $z=0$ the other
equation has an ordinary node. The holomorphic
Morse lemma with parameter $z$ therefore writes
it as $xy+h(z)$. Normality of that surface near
$D$ excludes $h\equiv0$, and a change of the
$z$ coordinate writes $h(z)=z^n$, $n\ge1$.
Consequently the original fibration has local
equation
\begin{equation}
f=u(x,y,z)\,z^a(xy+z^n)^b,
 \qquad a,b,n\ge1,\hypertarget{tag-00NS}{}\label{cdp:eq:C.INT.11}
\end{equation}
where $u$ is a unit and $a,b$ are the two full
fibre multiplicities. The unit can be removed
by a biholomorphic change of variables.
Put $d=2(a+nb)$ and choose a local logarithm of
$u$. For $f_0=z^a(xy+z^n)^b$, weighted homogeneity
gives
\begin{equation}
\Phi(x,y,z)=(x u^{n/d},y u^{n/d},z u^{2/d}),
 \qquad f_0\circ\Phi=u f_0.\hypertarget{tag-00NT}{}\label{cdp:eq:C.INT.12}
\end{equation}
The derivative of $\Phi$ at the origin is invertible.
More explicitly, near the chosen root of $u(0)$ solve
\[
 \tau^d=u(\tau^{-n}X,\tau^{-n}Y,\tau^{-2}Z).
\]
Its derivative in $\tau$ at $(X,Y,Z)=0$ is nonzero.
The implicit function theorem gives a holomorphic
solution, and
$(X,Y,Z)\mapsto(\tau^{-n}X,\tau^{-n}Y,\tau^{-2}Z)$
is the inverse of $\Phi$. Thus it identifies the same
full fibres with their original multiplicities $a,b$.

\item For $t\ne0$ small, take a simply connected
sector in the punctured $z$-plane and one
branch of $(t z^{-a})^{1/b}$. An open part
of the smooth fibre is
\begin{equation}
xy=g(z):=(t z^{-a})^{1/b}-z^n.\hypertarget{tag-00NU}{}\label{cdp:eq:C.INT.13}
\end{equation}
Consecutive zeros on this logarithmic branch
have angular separation $2\pi b/(a+nb)<2\pi$;
two such zeros fit in one sector. They are
simple, since at either one
$g'(z)=-(a/b+n)z^{n-1}\ne0$.
Join them by a smooth embedded path with no
other zero and take the circle $|x|=|y|$
over its interior. Explicitly, for a continuous
choice of $\arg g$ along that interior, it is
\begin{equation}
x=|g|^{1/2}e^{i(\arg g/2+\theta)},\qquad
 y=|g|^{1/2}e^{i(\arg g/2-\theta)}.\hypertarget{tag-00NV}{}\label{cdp:eq:C.INT.14}
\end{equation}
Apply \cref{lem:II:monomial-curve-matching-sphere}
with $b'=c'=1$ to the smooth connected graph
$B_t=\{(z,v):v=g(z)\}$ over this sector. The function
$v$ on $B_t$ has the two fixed simple zeros.
Its surface $y_{\rm lem}z_{\rm lem}=v$ is identified
with \eqref{cdp:eq:C.INT.13} by
\[
 (x_{\rm lem},y_{\rm lem},z_{\rm lem})
       \longmapsto (x,y,z)=(y_{\rm lem},z_{\rm lem},x_{\rm lem}),
 \qquad (x,y,z)\longmapsto(z,x,y).
\]
The displayed circle is exactly that lemma's circle
over the chosen path. Its endpoint charts give a
smooth embedded sphere of self-intersection $-2$.
The inverse of $\Phi$ transports it to the original
unit-bearing fibre. These holomorphic coordinate
isomorphisms preserve the complex orientation
and the compact intersection pairing.
The whole sphere lies in the local chart:
$|z|=O(|t|^{1/(a+nb)})$ and
$|x|,|y|=O(|t|^{n/(2(a+nb))})$.
Its image in the compact smooth torus fibre
would be null-homotopic, because a torus has
contractible universal cover. Its integral
homology class would therefore be zero and
its self-intersection zero, contradicting
$-2$. This contradiction excludes the Enoki
alternative.

\end{enumerate}
\end{proof}


\begin{proposition}[Euler sign and equality for reducible fibres]\hypertarget{tag-01C2}{}\label{cdp:st:C.2.INT.3}
\StatementRecord{cdp:st:C.2.INT.3}{proved}{01C2}{\cref{cdp:st:C.2.INT.2}, \cref{cdp:st:C.2.INT.1}, \cref{cdp:st:A.5.3}, \eqref{cdp:eq:C.INT.18}, \eqref{cdp:eq:C.INT.21}, \eqref{cdp:eq:C.INT.20}, \eqref{cdp:eq:C.INT.22}, \eqref{cdp:eq:C.INT.26}, \eqref{cdp:eq:C.INT.27}}

Let $f:Y\to B$ be a proper connected-fibre map from a smooth
compact complex threefold with $H^2(Y;\mathbb Q)=0$ and general
fibre a complex two-torus, as in \cref{cdp:st:C.2.INT.2}.
Let $S$ be a connected reducible fibre support.
Then
\begin{equation}
e(S)\ge0.\hypertarget{tag-00NW}{}\label{cdp:eq:C.INT.15}
\end{equation}
Equality holds precisely when the normalization of $S$
is a disjoint union of smooth minimal Hopf or
non-K\"ahler elliptic surfaces, and the reduced
conductor image and its reduced inverse image are
disjoint unions of smooth elliptic curves. The maps
between these elliptic curves are finite unramified
covers. This criterion allows arbitrary conductor
scheme structure and arbitrary unequal multiplicities
of the full fibre.
\end{proposition}

\begin{proof}

\begin{enumerate}[label=\arabic*.,leftmargin=*]
\item Let $\nu:T=\coprod T_i\to S$ be normalization,
$\pi_i:Z_i\to T_i$ its minimal resolutions, and
$\sigma_i:Z_i\to Z_i^0$ the minimal-model maps.
By \cref{cdp:st:C.2.INT.1}--\cref{cdp:st:C.2.INT.2} the $Z_i^0$ are Hopf or
non-K\"ahler elliptic surfaces. They have no rational
curves and have both Euler characteristics
$e(Z_i^0)=\chi(\OO_{Z_i^0})=0$.
Every rational curve of $Z_i$ is exceptional for
$\sigma_i$; their total number is exactly the number
$N_i$ of point blowups, and $e(Z_i)=N_i$.
The elliptic strict transforms cannot be contracted
by $\pi_i$, by \cref{cdp:st:C.2.INT.1}. Thus every
$\pi_i$-exceptional curve is $\sigma_i$-exceptional.

\item The connected fibres of $\pi_i$ consequently map to
points under $\sigma_i$. Normality and
$(\pi_i)_*\OO_{Z_i}=\OO_{T_i}$ descend the latter
map to a proper birational morphism
$\tau_i:T_i\to Z_i^0$, with
\begin{equation}
\sigma_i=\tau_i\pi_i.\hypertarget{tag-00NX}{}\label{cdp:eq:C.INT.16}
\end{equation}
Since point blowups have $R\sigma_{i*}\OO_{Z_i}
=\OO_{Z_i^0}$, the low-degree Leray sequence gives
\begin{equation}
0\longrightarrow R^1\tau_{i*}\OO_{T_i}
 \longrightarrow 0
 \longrightarrow \tau_{i*}R^1\pi_{i*}\OO_{Z_i}
 \longrightarrow R^2\tau_{i*}\OO_{T_i}=0.\hypertarget{tag-00NY}{}\label{cdp:eq:C.INT.17}
\end{equation}
The last vanishing follows from fibre dimension at
most one. The middle sheaf $R^1\pi_{i*}\OO_{Z_i}$
is supported at finitely many points, and its
pushforward under $\tau_i$ detects it. Hence it is
zero. This proves that $T_i$ has rational
singularities, and in particular
$\chi(\OO_{T_i})=\chi(\OO_{Z_i})=0$.

\item Let $\Gamma$ be the full conductor image scheme,
$A_i\subset T_i$ its full inverse-image conductor,
and let $D_i\ge0$ be the resolution conductor from
\cref{cdp:st:A.5.3}. Its line relation
$K_{Z_i}+D_i=h_i^*(K_Y+S)$ has torsion topological
first Chern class. That line descends through
$\sigma_i$. Consequently the exact canonical
divisor formula gives
\begin{equation}
D_i=\sigma_i^*D_i^0-K_{Z_i/Z_i^0},
 \qquad D_i^0=(\sigma_i)_*D_i.\hypertarget{tag-00NZ}{}\label{cdp:eq:C.INT.18}
\end{equation}
The divisor $D_i^0$ is nonzero, since $D_i$ contains
a nonzero adjacent component divisor whose
minimal-model pushforward is nonzero by \eqref{cdp:eq:C.INT.2}.
It is a sum of positive multiples of disjoint
smooth elliptic curves. At each point blowup,
\eqref{cdp:eq:C.INT.18} has coefficient
$\operatorname{mult}_pD_{\rm old}-1$ on the new
exceptional curve. Effectivity makes the centre
lie on the old support. Starting from the disjoint
elliptic curves, the reduced support consequently
acquires only rational leaves and subdivisions of
edges. Each of its connected components is one
elliptic curve with rational trees attached.
Writing $r_i$ for its number of rational curves,
we obtain
\begin{equation}
e((D_i)_{\rm red})=r_i.\hypertarget{tag-00O0}{}\label{cdp:eq:C.INT.19}
\end{equation}

\item Apply analytic proper duality
\cite[\S2.2.2, p.~60]{vanStraten1987} to
$\pi_i:Z_i\to T_i$ and $\OO_{Z_i}$.
Both spaces are Cohen--Macaulay of pure dimension two,
so their dualizing complexes are $\omega_{Z_i}[2]$
and $\omega_{T_i}[2]$. The already proved identity
$R\pi_{i*}\OO_{Z_i}=\OO_{T_i}$ identifies its two
derived Hom complexes and gives
$R\pi_{i*}\omega_{Z_i}=\omega_{T_i}$.
Apply \cref{lem:II:conductor-duality} to the finite
normalization of this same reduced hypersurface $S$.
On its component $T_i$ it identifies
$\omega_{T_i}\otimes\nu^*\omega_S^{-1}$ with
the conductor ideal $\mathcal I_{A_i/T_i}$.
The divisor sequence on $Z_i$, using
$\OO_{Z_i}(-D_i)=\omega_{Z_i}\otimes h_i^*\omega_S^{-1}$,
therefore gives
\begin{equation}
R\pi_{i*}\OO_{D_i}=\OO_{A_i}.\hypertarget{tag-00O2}{}\label{cdp:eq:C.INT.20}
\end{equation}
The identical calculation for $\sigma_i$, using
\eqref{cdp:eq:C.INT.18}, gives
\begin{equation}
R\sigma_{i*}\OO_{D_i}=\OO_{D_i^0},
 \qquad R\tau_{i*}\OO_{A_i}=\OO_{D_i^0}.\hypertarget{tag-00O3}{}\label{cdp:eq:C.INT.21}
\end{equation}
For the canonical pushforward at each point blowup, its relative canonical
line is $\OO(E)$, and
$0\to\OO\to\OO(E)\to\OO_E(-1)\to0$
has zero cohomology in its last term. It gives
$R\sigma_*\omega=\omega$; composition proves the
version used in \eqref{cdp:eq:C.INT.21}.

\item For $M\in\Pic^0(Y)$ its pullback to $Z_i$ descends
to a line $M_i^0\in\Pic^0(Z_i^0)$, by the Picard
blowup isomorphism. The two lines $\nu^*M|_{T_i}$ and $\tau_i^*M_i^0$
pull back under $\pi_i$ to the same line
$h_i^*M=\sigma_i^*M_i^0$. Applying $\pi_{i*}$ to
this isomorphism and using $\pi_{i*}\OO_{Z_i}=\OO_{T_i}$
and projection formula identifies the two lines on $T_i$.
Now projection formula and
\eqref{cdp:eq:C.INT.20}--\eqref{cdp:eq:C.INT.21} identify the spaces
\begin{equation}
H^0(A_i,\nu^*M|_{A_i})
   =H^0(D_i^0,M_i^0|_{D_i^0}).\hypertarget{tag-00O4}{}\label{cdp:eq:C.INT.22}
\end{equation}
For a component $nE$ of $D_i^0$, the filtration by
its powers of the defining ideal has graded lines
$M_i^0|_E\otimes\OO_E(-jE)$, $0\le j<n$.
They have degree zero, since $E^2=0$ on the
minimal model. \cref{cdp:st:C.2.INT.1} proves that the map from
$\Pic^0(Y)$ to $\Pic^0(E)$ is nonconstant. We can
therefore choose one $M$ avoiding, for each of
the finitely many such curves and integers $j$,
the prescribed value $\OO_E(jE)$.
Each excluded locus is a proper closed analytic
subset of the connected complex Lie group
$\Pic^0(Y)$; their finite union does not fill it.
All the graded lines are nontrivial of degree
zero on an elliptic curve and have no sections.
The filtration and \eqref{cdp:eq:C.INT.22} prove
\begin{equation}
H^0(A_i,\nu^*M)=0\quad\hbox{for every }i.\hypertarget{tag-00O5}{}\label{cdp:eq:C.INT.23}
\end{equation}

\item The conductor inclusion is
$\OO_\Gamma\hookrightarrow\nu_*\OO_A$, where
$A=\coprod A_i$. It gives $H^0(\Gamma,M)=0$.
The normalization sequence and the preceding
rational-singularity computation give
$\chi(\mathcal Q)=0$ for
$\mathcal Q=\nu_*\OO_T/\OO_S$: here
$\chi(\OO_S)=0$ is the Cartier-divisor
Riemann--Roch calculation under $H^2(Y,\Q)=0$.
Conductor duality is
$\mathcal Q=\omega_\Gamma\otimes\omega_S^{-1}|_\Gamma$.
Twisting on this compact Cohen--Macaulay curve by
an ambient topologically torsion line preserves
Euler characteristic; its degree on every
component is zero. Serre duality thus gives
\begin{equation}
\chi(\OO_\Gamma)=-\chi(\mathcal Q)=0.\hypertarget{tag-00O6}{}\label{cdp:eq:C.INT.24}
\end{equation}
The same holds after twisting by $M$, so
$H^0(\Gamma,M)=0$ implies $H^1(\Gamma,M)=0$.
Restriction to $C=\Gamma_{\rm red}$ surjects
on $H^1$, since its kernel is supported in
dimension at most one and has no $H^2$.
Consequently
\begin{equation}
\chi(\OO_C)=\chi(M|_C)=h^0(C,M|_C)\ge0.\hypertarget{tag-00O7}{}\label{cdp:eq:C.INT.25}
\end{equation}

\item Every normalized component of $C$ has genus zero or one, since it is dominated by a curve of $(D_i)_{\rm red}$. Apply \cref{lem:II:curve-euler-delta} with \eqref{cdp:eq:C.INT.25}. Its formula \eqref{cdp:eq:C.INT.26} gives $e(C)\ge0$, with equality precisely for a disjoint union of smooth elliptic curves.

\item We finish the topological comparison on each
normalization component. A $\pi_i$-exceptional
curve outside $(D_i)_{\rm red}$ is a $(-2)$ curve
disjoint from that support, by the adjunction
argument following \eqref{cdp:eq:C.INT.10}. Each connected
configuration of such curves is a tree: a cycle
or a repeated intersection would contradict
negative definiteness of its matrix with
all diagonal entries $-2$. Such fibres lie
off $A_i$, because a fibre over its conductor
meets the strict conductor and is connected.
Let $a_i$ count their irreducible curves, and
let $b_i$ count exceptional fibres contained
in $(D_i)_{\rm red}$ whose image lies outside
$A_i$. Removing the conductor and contracting
these fixed fibres gives the exact equality
\begin{equation}
e(T_i\setminus(A_i)_{\rm red})
       =N_i-r_i-a_i+b_i\ge0.\hypertarget{tag-00O9}{}\label{cdp:eq:C.INT.27}
\end{equation}
A tree of $a$ rational curves has Euler number
$a+1$, so its contraction removes $a$; a
positive-conductor fibre removed upstairs but
represented by a point downstairs adds one.
These are all changes, since normalization is
an isomorphism off the conductor. Nonnegativity
follows because the $r_i$ support curves and
the $a_i$ disjoint exceptional curves are
distinct rational curves among the $N_i$
rational curves of $Z_i$.
The normalization is an isomorphism
$T\setminus A\simeq S\setminus C$. Additivity
with \eqref{cdp:eq:C.INT.26} and \eqref{cdp:eq:C.INT.27} proves \eqref{cdp:eq:C.INT.15}.

\item If $e(S)=0$, every nonnegative summand in these
formulas is zero. In particular $C$ is a
union of smooth elliptic curves, $b_i=0$,
and $r_i+a_i=N_i$. Every rational curve in
$(D_i)_{\rm red}$ must then be $\pi_i$-exceptional:
a rational curve cannot map nonconstantly to
one of these elliptic curves. The $a_i$ curves
are exceptional by definition. Thus all $N_i$
rational curves of $Z_i$ are $\pi_i$-exceptional.
Every curve exceptional for $\sigma_i$ is consequently $\pi_i$-exceptional. \cref{lem:II:last-blowup-minimal-resolution} applies to these two maps and gives $N_i=0$.
There are no rational curves to contract,
and \cref{cdp:st:C.2.INT.1} excludes contracting elliptic
curves. Therefore $Z_i=T_i=Z_i^0$ is smooth
and minimal. Its reduced conductor is a
disjoint union of elliptic curves. A finite
map between smooth elliptic curves is
unramified by Riemann--Hurwitz, proving the
stated map comparison.
Conversely these normalization and conductor
conditions give $e(T_i)=e((A_i)_{\rm red})=e(C)=0$;
the normalization formula gives $e(S)=0$.
The underlying spaces, and all comparisons
just used, are independent of the positive
full-fibre multiplicities. This completes the
proof including the equality criterion.

\end{enumerate}
\end{proof}


\part{Compact occurrence and marked supports}\label{part:II:compact-occurrence}
The component classes above occur in compact torus fibrations through
polyhedral fillings, signed root fibrations and finite quotient maps.
Marked elliptic lattices and normal-line characters determine the
gluings of Hopf and ruled components. The constructions specify the
full fibre multiplicities, punctured families and component inclusions.

\section{Semistable torus fillings}\hypertarget{tag-034B}{}\label{sec:II:all:08}

\subsection{Semistable toric Euler models}\label{subsec:II:euler-monodromy-signatures}
















\begin{proposition}[Euler numbers in semistable toric rank-two models]\hypertarget{tag-01FI}{}\label{cdp:st:C2c}
\StatementRecord{cdp:st:C2c}{proved}{01FI}{\cite[\S6, pp.~264--268, (6.3)--(6.7)]{Mumford1972}}

Consider a rank-two Mumford model with smooth three-dimensional total space obtained from a unimodular periodic triangulation of $\R^2$, divided by a rank-two translation lattice acting freely on the total model. Let $V,E,F$ be the numbers of vertices, edges and triangles in the quotient triangulation. Its reduced central fibre has Euler number
\begin{equation}
e(S)=F=2V\in2\Z_{>0}.                                \hypertarget{tag-008Y}{}\label{cdp:eq:C.2}
\end{equation}
Every positive even number occurs in this class. A rank-one model assembled from $\PP^1$-bundles over elliptic curves with their boundary elliptic sections identified has Euler number zero.
\end{proposition}


\begin{proof}

\begin{enumerate}[label=\arabic*.,leftmargin=*]
\item The toric orbit strata of a rank-two central fibre have dimensions two, one and zero. Its two-dimensional strata are copies of $(\C^*)^2$, its one-dimensional strata copies of $\C^*$, and its zero-dimensional strata correspond to triangles of the quotient triangulation. The first two have compactly supported Euler number zero. Hence the Euler number is $F$. The quotient complex is a real two-torus, so $V-E+F=0$, and every edge has two incident triangles, giving $3F=2E$, with incidences counted even when a face has identified edges. Thus $F=2V$.

\item For existence, use the lattice $\widetilde N=\mathbb Z^2\oplus\mathbb Z$
and the standard triangulation of $\mathbb R^2$ whose triangles are
\[
 [u,u+e_1,u+e_1+e_2],\qquad
 [u,u+e_2,u+e_1+e_2],\qquad u\in\mathbb Z^2.
\]
These are the Delaunay triangles for the positive integral quadratic
form $q(u_1,u_2)=u_1^2-u_1u_2+u_2^2$.
Take the cones over these triangles at height one and all their
faces. They form a fan. Every three-dimensional cone is generated
by a basis of $\widetilde N$, because its triangle is unimodular. Gluing its
affine toric charts therefore gives a smooth complex threefold
$\mathcal X$ with a holomorphic map to the affine line. On each
maximal chart there are coordinates $x_1,x_2,x_3$ in which the map
is exactly $t=x_1x_2x_3$. The central support is consequently
reduced and has normal crossings. Over $t\ne0$ the fibre is the
common algebraic torus $(\mathbb C^*)^2$.

\item Fix an integer $N\ge1$ and let $\Lambda_N=\mathbb Z(Ne_1)\oplus\mathbb Z e_2$.
For $\ell\in\Lambda_N$, the lattice shear
\[
 (u,k)\longmapsto(u+k\ell,k)
\]
preserves the fan and induces an automorphism over the base.
On the punctured torus chart it is
\[
 (z_1,z_2,t)\longmapsto
 (t^{\ell_1}z_1,t^{\ell_2}z_2,t).
\]
Restrict to $|t|<\epsilon<1$. This action is free: away from the
central fibre, a fixed point would imply $t^{\ell_i}=1$ for both
coordinates and hence $\ell=0$; on the central fibre, a fixed
point would require its bounded indexing face of the triangulation
to be invariant under translation by $\ell$, which again forces
$\ell=0$.

\item To prove proper discontinuity, for $t\ne0$,
put $v_i=\log|z_i|/\log|t|$. In the chart of a triangle with
vertices $u_1,u_2,u_3$, write
\[
 v=\sum_{j=1}^3\lambda_j u_j,\qquad
 \lambda_j=\frac{\log|x_j|}{\log|t|},\qquad
 \sum_j\lambda_j=1.
\]
On a compact subset of finitely many charts the $|x_j|$ are
bounded. As $t$ tends to zero, this gives a uniform lower bound
for every $\lambda_j$, and their sum then gives a uniform upper
bound. Thus $v$ remains bounded there. Away from zero the same
conclusion follows from compactness in the torus charts.
Translation by $\ell$ sends $v$ to $v+\ell$. Consequently only
finitely many elements can carry a fixed compact set into itself,
or make two fixed compact sets intersect. On the central fibre
this is also immediate from the finitely many indexing faces.
The action is properly discontinuous, so its quotient is a smooth
complex threefold.

\item The quotient map to the disc is proper. Indeed, translate any
punctured-fibre point so that $v$ belongs to a fixed compact
fundamental parallelogram for $\Lambda_N$. Only finitely many
triangles meet that parallelogram. Choose one containing $v$;
then all its barycentric coordinates are nonnegative, so
$|x_j|=|t|^{\lambda_j}\le1$. These finitely many closed unit
polydiscs provide compact representatives for every sequence
approaching the central fibre. The central fibre itself has
finitely many toric components modulo $\Lambda_N$, and each is
compact: the star of a vertex is a complete two-dimensional fan.
These facts prove compactness over every closed smaller disc.
For $t\ne0$ the quotient fibre is explicitly
\[
 \mathbb C^*/t^{N\mathbb Z}\ \times\
 \mathbb C^*/t^{\mathbb Z},
\]
a product of elliptic curves. These charts and compactness prove that the quotient is a proper
holomorphic degeneration over the complex disc. The comparison
with \cite[\S6, pp.~264--268, (6.3)--(6.7)]{Mumford1972}
is given by the monomials of each chart. For a triangle with
vertices $u_j$, its dual-cone semigroup consists of the pairs
$(\alpha,r)$ satisfying
$r+\langle\alpha,u_j\rangle\ge0$ for every $j$.
Its monomial is $t^rz^\alpha$. Base change along
$\C[t]\to\C[[\pi]]$, $t\mapsto\pi$, sends this monomial to
$\pi^r\mathcal X^\alpha$ in the affine ring of (6.1), with
the same inequalities. Face inclusions give the same monomial
localizations as (6.5), and translation by $\ell$ is the shear
already displayed. The integral affine difference
$q(u+\ell)-q(u)$ gives the polarization data for (6.7).
Thus the chart, face and period maps identify the polyhedral
model used here with the cited construction after that base change.

\item The quotient triangulation has $V=N$ and $F=2N$.
For rank one every normalization component is a $\PP^1$-bundle over an elliptic curve, with Euler number zero, and the identified boundary sections also have Euler number zero. Additivity gives the result.

\end{enumerate}
\end{proof}











\section{Signed roots and compact component constructions}\hypertarget{tag-0345}{}\label{sec:II:09}


\subsection{Signed root fibrations}\label{subsec:II:signed-root-fibrations}

\begin{lemma}[Signed root fibrations over a smooth surface]\hypertarget{tag-P339}{}\label{lem:II:signed-root-fibration}
\StatementRecord{lem:II:signed-root-fibration}{proved}{P339}{}
Let $V$ be a smooth complex surface, and let $D_+,D_-$ be disjoint
smooth reduced divisors, either of which may be empty. Put
$L=\OO_V(D_+-D_-)$, and let $\sigma$ be its canonical meromorphic
section. Fix an integer $m\ge1$ and $q\in\C$ with $0<|q|<1$, and
write $E_q=\C^*/q^{\Z}$.

There is a smooth threefold $U$ with a proper holomorphic map
$\rho:U\to V$ whose fibres have connected elliptic reductions.
Over $V\setminus(D_+\cup D_-)$ this map is the product with $E_q$.
The full inverse image of either divisor is $m$ times its smooth
reduction. The construction is the quotient by the generator
\eqref{cdp:eq:C.PE.7} of the space with the following two charts:
\begin{equation}
 \sigma=a^m\ell,\qquad \ell\in L^\times,
 \qquad v\in V\setminus D_-,
 \hypertarget{tag-00UI}{}\label{cdp:eq:C.PE.5}
 \hypertarget{tag-00IZ}{}\SsixEquationAliases{\label{cdp:eq:C.ER.4}}
\end{equation}
\begin{equation}
 \sigma^{-1}=(a')^m\ell',\qquad \ell'\in(L^{-1})^\times,
 \qquad v\in V\setminus D_+,
 \hypertarget{tag-00UJ}{}\label{cdp:eq:C.PE.6}
\end{equation}
glued by $a'=a^{-1}$ and $\ell'=\ell^{-1}$ where both are defined.
The generator is
\begin{equation}
 F:(\ell,a)\longmapsto(q^m\ell,q^{-1}a),\qquad
 F:(\ell',a')\longmapsto(q^{-m}\ell',qa').
 \hypertarget{tag-00UK}{}\label{cdp:eq:C.PE.7}
 \hypertarget{tag-00J0}{}\SsixEquationAliases{\label{cdp:eq:C.ER.5}}
\end{equation}
Near a point of either divisor, with its equation $x=0$ and a
second coordinate $y$, the map is the free $\boldsymbol\mu_m$
quotient of
\begin{equation}
 \begin{gathered}
 \Delta_z\times E_q\times\Delta_y
   \longrightarrow\Delta_x\times\Delta_y,
 \qquad(z,[w],y)\longmapsto(z^m,y),\\
 \zeta:(z,[w],y)\longmapsto(\zeta z,[\zeta w],y).
 \end{gathered}
 \hypertarget{tag-00UL}{}\label{cdp:eq:C.PE.8}
 \hypertarget{tag-00J2}{}\SsixEquationAliases{\label{cdp:eq:C.ER.7}}
\end{equation}
The product isomorphism away from the two divisors and its inverse are
\begin{equation}
 \begin{gathered}
 (v,[c])\longmapsto[(v,\ell=\sigma(v)c^m,a=c^{-1})],
 \\ [(v,\ell,a)]\longmapsto(v,[a^{-1}]).
 \end{gathered}
 \hypertarget{tag-00UM}{}\label{cdp:eq:C.PE.9}
 \hypertarget{tag-00J3}{}\SsixEquationAliases{\label{cdp:eq:C.ER.8}}
\end{equation}
The higher direct images are
\begin{equation}
 \rho_*\OO_U=\OO_V,\qquad R^1\rho_*\OO_U\simeq\OO_V,
 \qquad R^j\rho_*\OO_U=0\quad(j\ge2).
 \hypertarget{tag-00UF}{}\label{cdp:eq:C.PE.21}
\end{equation}
The additive deck cocycle $F\mapsto1$ gives the specified global
frame of $R^1\rho_*\OO_U$. For compact $V$, the threefold $U$ is compact.
\end{lemma}

\begin{proof}
On the overlap of the two charts, $\sigma$ is nonzero, so both
$a$ and $a'$ are nonzero. Their reciprocal formulas are inverse
holomorphic maps and intertwine the two displayed actions of $F$.
The projection to $L^\times$ is global, using $\ell=(\ell')^{-1}$
on the negative chart. Let $s_\pm$ be the canonical sections of
$\OO_V(D_\pm)$, so $\sigma=s_+/s_-$. For $a=A/B$,
the homogeneous equation
\[
                         s_+B^m=s_-A^m\ell
\]
is an equality in the same line $\OO_V(D_+)$ over $L^\times$.
Since $s_+$ and $s_-$ have no common zero, its closed analytic
subspace in $L^\times\times\PP^1$ is exactly the glued space.
In a local frame near $D_-$ its affine relation is
$x=(a')^m b'$, with $b'\ne0$. Near $D_+$ it is $x=a^m b$,
with $b\ne0$. Each eliminates $x$, and off the two divisors the
coordinates $(v,a)$ with $a\ne0$ do so. The glued space $W$ is
therefore a smooth Hausdorff complex threefold.

Choose a Hermitian norm on $L$. The nonzero global vector $\ell$
has norm multiplied by $|q|^m$ under $F$. That norm is bounded
above and away from zero on every compact subset of $W$.
Thus only finitely many translates of one compact set meet
another, and no nontrivial power of $F$ fixes a point.
The quotient $U=W/\langle F\rangle$ is smooth and Hausdorff.

In the positive chart the coordinate map is
\begin{equation}
 (z,w,y)\longmapsto(x=z^m,y;\ b=w^m,\ a=z/w),
 \qquad w\in\C^*.
 \hypertarget{tag-00J1}{}\label{cdp:eq:C.ER.6}
\end{equation}
An $m$th root $w$ of $b$ determines $z=aw$, and changing
that root multiplies both $z,w$ by the same $\zeta\in\boldsymbol\mu_m$.
These are exactly the fibres of the chart map, and its finite
action is free because $w\ne0$. The map $F$ lifts to
$w\mapsto qw$. In the negative chart the formulas are
$x=z^m$, $b'=w^m$, $a'=z/w$, and $F$ lifts to $w\mapsto q^{-1}w$.
Both cyclic groups on $w$ are $q^{\Z}$, giving
\eqref{cdp:eq:C.PE.8}. Its finite action is free: an equality
$[\zeta w]=[w]$ implies $\zeta=q^k$; absolute values force
$k=0$, and then $\zeta=1$.

The map before that finite quotient is proper, being the product
of the finite disc map with the compact elliptic curve $E_q$.
Its quotient is proper. Away from the divisors,
\eqref{cdp:eq:C.PE.9} is well defined because replacing $c$ by
$qc$ applies $F$, and the displayed inverse proves the product
description. Properness is local on the target, so these charts
prove that $\rho$ is proper. Over $x\ne0$, the finite action
permutes the $m$ root sheets transitively; the fibre is $E_q$.
Over $x=0$ the reduced fibre is $E_q/\boldsymbol\mu_m$, and
the equation $x=a^m b$ or $x=(a')^m b'$ gives its full
multiplicity $m$. The latter reduced elliptic curve is identified
with $\C^*/(q^m)^{\Z}$ by $[w]\mapsto[w^m]$.
The same coordinates show that the reduced inverse-image divisor is smooth.

On a chart \eqref{cdp:eq:C.PE.8}, compute higher direct images
before taking the finite quotient. For the product with $E_q$
they are $\OO$ in degrees zero and one and zero above one.
The finite disc map has exact direct image, and finite-group
invariants are exact over $\C$. Translation by $[\zeta]$ acts
trivially on $H^1(E_q,\OO_{E_q})$, so both surviving invariant
modules are the invariant functions in $z$, namely the
holomorphic functions in $x=z^m$. This proves the local
identifications in \eqref{cdp:eq:C.PE.21}.

The global principal $\Z$ covering $W\to U$ defines the
additive cocycle with value one on $F$. On the product cover of
a positive chart its restriction is the cocycle $w\mapsto qw$
with value one; on a negative chart it is the cocycle for
$w\mapsto q^{-1}w$ with value one. Neither is a coboundary:
the constant Laurent coefficient of $b(q^{\pm1}w)-b(w)$ is zero
for every holomorphic $b$ on $\C^*$, whereas this cocycle has
constant value one. These classes are invariant under the finite
translations, and their pullbacks to the elliptic product covers
are nonzero also at $z=0$. They therefore form nonvanishing
local frames of the computed $R^1$, all induced by the same
global cocycle. This gives its stated trivialization.
Finally, properness of $\rho$ gives compactness over compact $V$.
\end{proof}

\subsection{Elliptic root fillings}\label{subsec:II:elliptic-root-fillings}

\begin{lemma}[An elliptic root surface and its distinguished fibre]\hypertarget{tag-P336}{}\label{lem:II:elliptic-root-surface}
\StatementRecord{lem:II:elliptic-root-surface}{proved}{P336}{}
Let $B$ be an elliptic curve with origin $0$, let
$L=\OO_B([0])$ with canonical section $\sigma$, and let
$m\ge1$ and $0<|\beta|<1$. Put
\begin{equation}
 \begin{gathered}
 \mathcal W_m=\{(v,\ell,a):v\in B,\ \ell\in L_v^\times,
                         \ \sigma(v)=a^m\ell\},\\
 P_m(B,\beta)=\mathcal W_m/\langle g\rangle,\qquad
 g(v,\ell,a)=(v,\beta^m\ell,\beta^{-1}a).
 \end{gathered}
 \label{eq:II:elliptic-root-surface}
\end{equation}
Then $P=P_m(B,\beta)$ is a smooth compact complex surface.
Projection $p:P\to B$ has connected elliptic fibres, is the
product with $E_\beta=\C^*/\beta^{\Z}$ over $B\setminus\{0\}$,
and has full fibre $mJ$ over $0$, where
$J=\C^*/\beta^{m\Z}$. One has
\[
 K_P\simeq\OO_P((m-1)J),\qquad
 p_*\OO_P(kJ)=\OO_B(\lfloor k/m\rfloor[0])\quad(k\ge0),
\]
and $H^\bullet(P,\Z)$ is free of ranks $(1,3,4,3,1)$, while
$h^\bullet(P,\OO_P)=(1,2,1)$ and $e(P)=0$.
For $n\ge1$, writing $d_n=\lfloor n(m-1)/m\rfloor$, the
plurigenus is $1$ when $d_n=0$ and is $d_n$ otherwise.
Thus $\kappa(P)=0$ for $m=1$ and $\kappa(P)=1$ for $m\ge2$.

There is a smooth identification of pairs
$(P,J)\simeq(M_m\times S^1,K_m\times S^1)$, where $M_m$
is the closed oriented three-manifold obtained by the filling
\eqref{cdp:eq:C.ER.16} and $K_m$ is its solid-torus core.
The angular loop of $J$ is zero in $H_1(P,\Z)$ and its
deck-period loop is a primitive generator of the last $S^1$;
in particular $[J]=0$ in $H_2(P,\Z)$.
For $m=1$, projection $(v,\ell,a)\mapsto(v,\ell)$ identifies
$\mathcal W_1$ with $L^\times$, and $P_1(B,\beta)$ with
$L^\times/\beta^{\Z}$ for this degree-one line bundle.
\end{lemma}

\begin{proof}
Apply \cref{lem:II:signed-root-fibration} to
$V=B\times\Delta$, $D_+=\{0\}\times\Delta$, $D_-=\varnothing$,
the line $\operatorname{pr}_B^*L$ and section
$\operatorname{pr}_B^*\sigma$, with its $(m,q)=(m,\beta)$.
The positive chart is defined on all of $V$. The maps
\[
 [(v,y,\ell,a)]\longmapsto([(v,\ell,a)],y),\qquad
 ([(v,\ell,a)],y)\longmapsto[(v,y,\ell,a)]
\]
are inverse holomorphic maps between that threefold and
$P_m(B,\beta)\times\Delta$: the equations and the deck
transformations are identical on every line-bundle chart.
Restricting its projection to $y=0$ therefore gives the map
$p:P\to B$ in the statement. The cyclic chart
\eqref{cdp:eq:C.ER.6} restricts to
$(z,w)\mapsto(x=z^m,b=w^m,a=z/w)$, with the same free
$\boldsymbol\mu_m$ action and infinite map $w\mapsto\beta w$.
It proves smoothness of $P$ and identifies its full fibre as
$mJ$, where $J=\C^*/\beta^{m\Z}$ through $[w]\mapsto[w^m]$.
Properness is preserved by this base change, and $B$ is compact,
so $P$ is compact. The product map and inverse are the restrictions
of \eqref{cdp:eq:C.PE.9}:
$(v,[c])\mapsto[(v,\sigma(v)c^m,c^{-1})]$ and
$[(v,\ell,a)]\mapsto(v,[a^{-1}])$.

Choose a nowhere-zero holomorphic one-form $\omega$ on $B$.
The form on the cyclic chart
\begin{equation}
                         p^*\omega\wedge\frac{dw}{w}
 \hypertarget{tag-00IQ}{}\label{cdp:eq:C.ER.12}
\end{equation}
is invariant under both group actions. On the product part
it is $p^*\omega\wedge dc/c$, since $c=w/z$ and the extra
term from $dz/z$ wedges to zero. A frame change adds a
one-form from the base and hence leaves this two-form unchanged.
If $\omega=h(x)\,dx$, the cyclic chart expression is
\begin{equation}
 m h(z^m)z^{m-1}\,dz\wedge\frac{dw}{w}.
 \hypertarget{tag-00IR}{}\label{cdp:eq:C.ER.13}
\end{equation}
Its zero divisor is $(m-1)J$, giving the canonical isomorphism.

For $k\ge0$ the full direct-image identity is
\begin{equation}
        p_*\OO_P(kJ)=\OO_B\bigl(\lfloor k/m\rfloor[0]\bigr).
 \hypertarget{tag-00IS}{}\label{cdp:eq:C.ER.14}
\end{equation}
A section on the left is constant on each compact fibre
off $0$, so descends there. On the cyclic chart its Laurent
coefficients in $z$ are holomorphic on the compact elliptic
curve and hence constant. Invariance forces the exponents
to be multiples of $m$, and the pole bound $k$ becomes
$\lfloor k/m\rfloor$ in $x=z^m$. This proves both inclusions.
Apply it with $k=n(m-1)$ and use Riemann--Roch on $B$ to
obtain the stated plurigenera and Kodaira dimensions.
All reduced fibres of $p$ are elliptic, so Euler additivity
gives $e(P)=0$. The divisor $mJ$ is a fibre, hence $J^2=0$;
the canonical formula and surface Riemann--Roch give
$\chi(\OO_P)=0$. The plurigenus at $n=1$ is one, and
Serre duality now gives $h^\bullet(\OO_P)=(1,2,1)$.

Choose a Hermitian metric on $L$, constant in the frame at $0$,
and put
\begin{equation}
 M_m=\{(v,\ell,a)\in\mathcal W_m:\|\ell\|=1\}.
 \hypertarget{tag-00IT}{}\label{cdp:eq:C.ER.15}
\end{equation}
The radial line action is transverse to this level, so it is
a smooth oriented three-manifold. On the level,
$|a|^m=\|\sigma\|$ is bounded over the compact base and
the unit line circle bundle is compact; thus $M_m$ is compact.
The maps
$\ell=e^{mu}\ell_0$, $a=e^{-u}a_0$ identify
$\mathcal W_m=M_m\times\R$.
Write $\beta=\rho e^{\mathrm i\theta}$, $0<\rho<1$, and let
$\Phi_t(\ell,a)=(e^{\mathrm i mt}\ell,e^{-\mathrm i t}a)$
be the circle action on $M_m$. The deck map in these
coordinates is $(v_0,u)\mapsto(\Phi_\theta(v_0),u+\log\rho)$.
The explicit invariant map
\[
 (v_0,u)\longmapsto
 \bigl(\Phi_{-\theta u/\log\rho}(v_0),[u]\bigr)
 \in M_m\times(\R/\log\rho\,\Z)
\]
induces a diffeomorphism from its quotient. Its inverse sends
$(v_1,[u])$ to the deck orbit of
$(\Phi_{\theta u/\log\rho}(v_1),u)$; changing the lift of
$[u]$ changes that pair by the deck map. The circle action
preserves the locus $a=0$, so this is an identification of
the stated pairs, including for nonreal $\beta$.

Over the once-punctured base, the canonical section trivializes
$M_m$ with circle coordinate the phase of $a$. Over the
disc the solid-torus coordinates are
$a=\rho_0e^{\mathrm i\vartheta}$, $b=e^{\mathrm i\varphi}$,
and
\begin{equation}
              x=a^m b,\qquad \arg x=m\vartheta+\varphi.
 \hypertarget{tag-00IU}{}\label{cdp:eq:C.ER.16}
\end{equation}
Let $\delta$ be the boundary of the punctured base and let
$h$ be its product circle with the phase-of-$a$ orientation.
The meridian $\vartheta\mapsto\vartheta+2\pi$ maps to
$m\delta+h$; the longitude $\varphi\mapsto\varphi+2\pi$
maps to $\delta$. In the base's first homology $\delta=0$,
so Mayer--Vietoris kills $h$ and the longitude, leaving
$H_1(M_m,\Z)=\Z^2$. The degree-zero part contributes no
class, since both pieces and their boundary torus are connected.
Integral Poincar\'e duality and the universal coefficient
theorem give free homology of ranks $(1,2,2,1)$ for $M_m$;
the displayed product gives the asserted groups of $P$.
The solid-torus core has the longitude class zero in
$H_1(M_m,\Z)$, while the second circle of $J$ is the
primitive deck-period circle. Their product class in
$H_2(P,\Z)$ is therefore zero.

Finally, for $m=1$ the inverse to projection onto $L^\times$
is $(v,\ell)\mapsto(v,\ell,\sigma(v)/\ell)$. It intertwines
the deck maps with multiplication by $\beta$ on $L^\times$,
proving the stated identification for the degree-one bundle.
\end{proof}


\begin{proposition}[Zero-Euler product fillings with essential Kodaira-dimension-one components]\hypertarget{tag-01E5}{}\label{cdp:st:C.2.UN.6}
\StatementRecord{cdp:st:C.2.UN.6}{proved}{01E5}{\SeriesRef{cdp:eq:C.ER.4}, \SeriesRef{cdp:eq:C.ER.5}, \SeriesRef{cdp:eq:C.ER.7}, \SeriesRef{cdp:eq:C.ER.8}, \SeriesRef{cdp:eq:C.ER.14}, \SeriesRef{cdp:eq:C.ER.9}, \SeriesRef{cdp:eq:C.ER.3}}

Put $q=e^{-2\pi}$ and $E=\C^*/q^{\Z}\simeq\C/(\Z+\mathrm i\Z)$.
For every integer $m\ge2$, the fixed punctured product
$(E\times E)\times\Delta^*$ has a proper filling
$f_m:U_m\to\Delta$ with smooth total space and full central fibre
\begin{equation}
                    (U_m)_0=P_m+mQ_m.             \hypertarget{tag-00IN}{}\label{cdp:eq:C.ER.1}
\end{equation}
Here $Q_m$ is a primary Hopf surface, and $P_m$ is a smooth compact
non-K\"ahler elliptic surface over $E$ with precisely one multiple
fibre $mJ$. The surfaces meet transversely along the smooth elliptic
curve $J$. One has
\begin{equation}
 \begin{gathered}
 e((U_m)_0{}_{\rm red})=0,\qquad H_1(P_m,\Z)=\Z^3,\\
 K_{P_m}\simeq\OO_{P_m}((m-1)J),\qquad \kappa(P_m)=1,\\
 h^0(P_m,K_{P_m}^{\otimes r})=
 \begin{cases}
 1,&d_r=0,\\
 d_r,&d_r>0,
 \end{cases}
 \qquad d_r=\left\lfloor\frac{r(m-1)}m\right\rfloor,
 \quad r\ge1.
 \end{gathered}                                   \hypertarget{tag-00IX}{}\label{cdp:eq:C.ER.2}
\end{equation}
Neither $P_m$ nor $Q_m$ can be contracted by a bimeromorphic map
between smooth threefold fillings. The $U_m$ are pairwise
non-bimeromorphic over the disc. After every finite base change and
resolution of its normalization, a central component remains
birational to $P_m$. All these fillings have compact smooth
threefold realizations over $\PP^1$.
\end{proposition}


\begin{proof}

\begin{enumerate}[label=\arabic*.,leftmargin=*]
\item The occurrence of a Kodaira-dimension-one component in an abelian-surface degeneration appears in
\cite[Example~3.4(1), p.~99]{Fujimoto1988Log}. Let
\begin{equation}
 V=\operatorname{Bl}_{(0,0)}(E\times\Delta),\quad
 R=\operatorname{Exc}(V\to E\times\Delta),\quad
 L=\OO_V(R),\quad \sigma\in H^0(V,L)
 \hypertarget{tag-00IY}{}\label{cdp:eq:C.ER.3}
\end{equation}
with $\sigma$ the canonical section. Apply
\cref{lem:II:signed-root-fibration} to this smooth surface $V$,
its smooth divisor $D_+=R$, the empty divisor $D_-$, and the
same $m,q$. Its line and section are exactly the displayed
$L,\sigma$. Because $D_-$ is empty, its positive chart
\eqref{cdp:eq:C.ER.4} is defined on all of $V$; denote that
space by $W_m$ and its quotient by $U_m$. The identity
$(v,\ell,a)\mapsto(v,\ell,a)$ identifies its equation
$\sigma(v)=a^m\ell$ and generator \eqref{cdp:eq:C.ER.5}
with this construction. Thus the lemma gives a smooth
threefold $U_m$ and a proper map $\rho_m:U_m\to V$ with
connected elliptic reductions.

In a frame $\sigma=xe$, $\ell=be$, the cyclic map
\eqref{cdp:eq:C.ER.6} is the same map
$(z,w,y)\mapsto(x=z^m,y;b=w^m,a=z/w)$.
Its simultaneous $\boldsymbol\mu_m$ action and the lift
$w\mapsto qw$ of the infinite generator give exactly the
local projection \eqref{cdp:eq:C.ER.7}. In particular the
proper local models and full multiplicity $m$ are the
ones supplied by that lemma. Compose $\rho_m$ with the
blowdown $V\to E\times\Delta$ and the projection
$E\times\Delta\to\Delta$. All three maps are proper
with connected fibres, so their composite is the required
proper map $f_m$ with connected fibres.

Over $V\setminus R$ the product isomorphism
\eqref{cdp:eq:C.ER.8} is
$(v,[c])\mapsto[(v,\sigma(v)c^m,c^{-1})]$, with inverse
$[(v,\ell,a)]\mapsto(v,[a^{-1}])$.
The blowdown identifies $V|_{\Delta^*}$ with
$E\times\Delta^*$ by its unchanged coordinates, so this
same map identifies the punctured filling with the fixed
$(E\times E)\times\Delta^*$ for every $m$.

\item Let $E_0\simeq E$ be the strict transform of $E\times\{0\}$.
It meets $R\simeq\PP^1$ at one point, and
$L|_{E_0}=\OO_E([0])$, $L|_R=\OO_{\PP^1}(-1)$.
The inverse image of $R$, with its reduction, is $a=0$ in
\eqref{cdp:eq:C.ER.4}, so it is
\begin{equation}
 Q_m=L|_R^\times/(q^m)^{\Z}
     \simeq(\C^2\setminus\{0\})/\langle q^m\operatorname{id}\rangle.
                                                        \hypertarget{tag-00J4}{}\label{cdp:eq:C.ER.9}
\end{equation}
The first isomorphism sends the nonzero tautological vector to
the same vector in $\C^2\setminus\{0\}$. It identifies $Q_m$
as a primary Hopf surface. The reduced inverse image $P_m$ of
$E_0$ is smooth by \eqref{cdp:eq:C.ER.7}, since $E_0$ meets $R$ transversely.
Its map $p:P_m\to E_0$ has smooth elliptic fibres off $[0]$ and
one full fibre $mJ$ over $[0]$. The reduced fibre is
\begin{equation}
                    J=\C^*/(q^m)^{\Z}.
                                                        \hypertarget{tag-00IO}{}\label{cdp:eq:C.ER.10}
\end{equation}
For example the isomorphism $E/\boldsymbol\mu_m\to J$ in
\eqref{cdp:eq:C.ER.7} is $[w]\mapsto[w^m]$, which also specifies the finite
covering between the two elliptic curves.

\item Near $E_0\cap R$, choose coordinates with $R=(x=0)$,
$E_0=(y=0)$ and base map $t=xy$. In the smooth chart before
the free quotient the equation is $t=z^m y$; in \eqref{cdp:eq:C.ER.4} it is
\begin{equation}
                         t=a^m b y,\qquad b\ne0.
                                                        \hypertarget{tag-00IP}{}\label{cdp:eq:C.ER.11}
\end{equation}
It follows that the full central fibre is exactly $P_m+mQ_m$,
and its reduction has the two transverse components stated in
\eqref{cdp:eq:C.ER.1}. The base equation has order $m$ along $Q_m$, including at $J$.

\item Apply \cref{lem:II:elliptic-root-surface} with
$B=E_0$, $\beta=q$ and the same integer $m$.
Restriction of \eqref{cdp:eq:C.ER.4} to $E_0$ is exactly its
space $\mathcal W_m$, since $L|_{E_0}=\OO_{E_0}([0])$
and the restricted section has its simple zero there.
The restriction map $(v,\ell,a)\mapsto(v,\ell,a)$ and its
inverse intertwine \eqref{cdp:eq:C.ER.5} with the lemma's
deck map. They therefore identify this pair $(P_m,J)$
with the pair in that lemma, including the product map
\eqref{cdp:eq:C.ER.8} with parameter $c=a^{-1}$.
The cyclic coordinates are the same $(z,w)$ as in
\eqref{cdp:eq:C.ER.6}; hence the canonical forms
\eqref{cdp:eq:C.ER.12}--\eqref{cdp:eq:C.ER.13}, the
direct-image identity \eqref{cdp:eq:C.ER.14}, and the
Seifert maps \eqref{cdp:eq:C.ER.15}--\eqref{cdp:eq:C.ER.16}
apply with these unchanged charts.
The lemma gives the canonical bundle, plurigenera,
$\kappa(P_m)=1$ and $H_1(P_m,\Z)=\Z^3$ in
\eqref{cdp:eq:C.ER.2}. This odd first Betti number excludes
a K\"ahler metric by Hodge decomposition.

\item The same lemma gives $e(P_m)=0$.
The Hopf surface in \eqref{cdp:eq:C.ER.9} is a principal elliptic bundle over
$\PP^1$, so $e(Q_m)=0$. Finally $e(J)=0$, and consequently
$e(P_m\cup_JQ_m)=0$.

\item Every map $\PP^1\to P_m$ is constant: it first maps constantly
to the elliptic base $E_0$, and then constantly to the reduced
elliptic fibre containing its image. Every map $\PP^1\to Q_m$
lifts to $\C^2\setminus\{0\}$ and is constant there by its
coordinate functions. Thus neither component is covered by
rational curves. Apply \cref{lem:II:nonuniruled-divisors} to these compact central divisors in the smooth proper disc model $U_m$. Its strict-transform correspondence \eqref{cdp:eq:C.UC.1} carries each one birationally onto a divisor in every smooth bimeromorphic filling. The Kodaira dimension of $Q_m$
is $-\infty$: a lifted pluricanonical form extends across the
origin by Hartogs, and its invariance under the contraction
$q^m\operatorname{id}$ forces all its power-series coefficients
to vanish. It cannot correspond birationally to $P_{m'}$.
Therefore any bimeromorphic comparison $U_m\dashrightarrow U_{m'}$
would identify $P_m$ birationally with $P_{m'}$.
Point blowups preserve the plurigenera of a smooth compact surface
\cite[Chapter~1, Proposition~2.2, p.~8]{Peters1990}. Resolving the graph of a surface bimeromorphism by such blowups therefore gives their birational invariance.
The different slopes $(m-1)/m$ in \eqref{cdp:eq:C.ER.2} force $m=m'$.

\item For the persistence assertion, take a finite base change and choose
a local parameter in which it is $t=s^d$; a unit is absorbed by
its local $d$th root. At every point of the nonempty open subset
$P_m\setminus J$, the map $f_m$ is a submersion. Its base change
is the smooth chart $t=s^d$ with the two coordinates along $P_m$
unchanged. Thus a unique central component maps with degree one
onto this open subset. Its closure, followed by normalization and
any resolution, is birational to $P_m$. The preceding
non-uniruled-divisor comparison prevents its contraction under
any further bimeromorphic modification between smooth fillings.
This proves the asserted persistence of the component birational to $P_m$.

\item Finally replace $\Delta$ throughout \eqref{cdp:eq:C.ER.3} by $\PP^1$ and use
the same exceptional divisor, line bundle, and canonical section.
The local proofs of smoothness and properness still apply. Since
the base $V$ is now compact, the quotient is a compact smooth
threefold, with the same central fibre and with every other fibre
the product $E\times E$. This gives the claimed compact realizations.

\end{enumerate}
\end{proof}



\subsection{Multiple elliptic fibrations on Hopf threefolds}\label{subsec:II:multiple-elliptic-hopf-threefolds}

\begin{proposition}[Multiple elliptic fibrations on Hopf threefolds]\hypertarget{tag-P328}{}\label{prop:II:multiple-elliptic-hopf-threefold}
\StatementRecord{prop:II:multiple-elliptic-hopf-threefold}{proved}{P328}{\cref{cdp:st:C.2.UN.6}; \cite[Theorem~2.8 and Corollary~2.10, pp.~96--97]{Fujimoto1988Log}}
For every $m\ge2$, put $q=e^{-2\pi}$ and define
\begin{equation}
 \begin{split}
 H_m&=(\C^3\setminus\{0\})/
       \langle(a,b,c)\longmapsto(qa,q^m b,q^m c)\rangle,\\
 \Phi_m:H_m&\longrightarrow\PP^2,
       \qquad[a,b,c]\longmapsto[a^m:b:c].
 \end{split}                                             \hypertarget{tag-00IV}{}\label{cdp:eq:C.ER.17}
\end{equation}
The threefold $H_m$ is smooth and compact, and $\Phi_m$ is a proper holomorphic map with connected fibres and smooth general fibre an elliptic curve. Every reduced fibre support is a smooth elliptic curve. For the line $B=(X=0)$ and its reduced inverse image $D=\{a=0\}$,
\begin{equation}
             \Phi_m^*B=mD,\qquad B^2=\deg B=1,
             \qquad\pi_1(\PP^2\setminus B)=0.             \hypertarget{tag-00IW}{}\label{cdp:eq:C.ER.18}
\end{equation}
These maps are counterexamples to the divisibility conclusion and the asserted surjection onto $\Z/m$ in
\cite[Theorem~2.8 and Corollary~2.10, pp.~96--97]{Fujimoto1988Log}.
They are the quotient construction \eqref{cdp:eq:C.ER.4}--\eqref{cdp:eq:C.ER.5} with $V=\PP^2$, $L=\OO(1)$ and $\sigma=X$.
\end{proposition}
\begin{proof}
The action is free and properly discontinuous. A compact annulus
between a small sphere and its image under the contraction meets every
orbit, so $H_m$ is compact and smooth. The three projective coordinates
in \eqref{cdp:eq:C.ER.17} are multiplied by the same factor $q^m$, proving descent.
The map is holomorphic, surjective, and proper. Its fibre over
$[1:Y:Z]$ is parametrized by
$a\mapsto(a,Ya^m,Za^m)$ and is $\C^*/q^{\Z}$; its fibre over
$[0:Y:Z]$ is parametrized by
$s\mapsto(0,sY,sZ)$ and is $\C^*/(q^m)^{\Z}$.
Every fibre is connected; its reduced support is the displayed elliptic curve, and the general fibre is smooth. The smooth-complement and constant-modulus assertions are given by the inverse maps
\[
 [a,b,c]\longmapsto([a],b/a^m,c/a^m),\qquad
 ([a],u,v)\longmapsto[a,ua^m,va^m],
\]
between $\Phi_m^{-1}(\PP^2\setminus B)$ and $(\C^*/q^{\Z})\times\C^2$. They are well defined under the displayed contraction. Thus the map meets the elliptic-threefold definition of \cite[pp.~91--92]{Fujimoto1988Log} and the smooth-complement and constant-modulus premises of its Corollary~2.10.
Near a point of the line $B=(X=0)$, use a chart with $b\ne0$ or
$c\ne0$. In the first chart $X/Y=a^m/b$, where $b$ is a unit.
It follows that $\Phi_m^*B=mD$. The line has $B^2=\deg B=1$. Its complement is identified with $\C^2$ by $[X:Y:Z]\mapsto(Y/X,Z/X)$, with inverse $(u,v)\mapsto[1:u:v]$, and is simply connected.
Since $m\ge2$, the self-intersection $1$ is not divisible by $m$, and the trivial fundamental group has no surjection onto $\Z/m$. These give the two stated counterexamples.

The covering construction in \cite[proof of Theorem~2.8, p.~96]{Fujimoto1988Log} has the following normalization. For an
effective divisor $D$ with canonical section $s_D$, the construction
in that proof takes the equation $\xi^m=s_D^m$ in the total space of
$\OO(D)$. Its normalization is the disjoint union of the $m$ global
graphs $\xi=\zeta s_D$, $\zeta^m=1$, each isomorphic to the original
space. The factorization $\xi^m-s_D^m=\prod_{\zeta^m=1}(\xi-\zeta s_D)$ shows that these are all its irreducible components. The normalization map is $(x,\zeta)\mapsto(x,\zeta s_D(x))$; on each component its inverse projection identifies it with the smooth original total space. The total normalization therefore has $m$ connected components. Finally, \eqref{cdp:eq:C.ER.17} is itself a special case of \eqref{cdp:eq:C.ER.4}: take
$V=\PP^2$, $L=\OO(1)$ and $\sigma=X$. A point of $V$ is a line $\lambda\subset\C^3$, and $L_\lambda=\lambda^*$. A nonzero $\ell\in\lambda^*$ determines the unique vector $v\in\lambda$ with $\ell(v)=1$. The equality $\sigma=a^m\ell$ says $X(v)=a^m$, so write $v=(a^m,b,c)$. Conversely $(a,b,c)\ne0$ determines $\lambda=\C(a^m,b,c)$ and its unique functional taking that vector to one. These constructions are holomorphic inverse maps between the hypersurface and $\C^3\setminus\{0\}$. Under \eqref{cdp:eq:C.ER.5}, the vector $v$ changes to $q^{-m}v$ and $a$ to $q^{-1}a$, so $(a,b,c)$ changes to $(q^{-1}a,q^{-m}b,q^{-m}c)$. This is the inverse of the contraction in \eqref{cdp:eq:C.ER.17}, giving the asserted isomorphism of quotient constructions.
\end{proof}

\subsection{Cuspidal self-products}\label{subsec:II:cuspidal-self-products}

\begin{proposition}[A simply connected irreducible rational-normalized fibre]\hypertarget{tag-01ER}{}\label{cdp:st:C.2.UN.8}
\StatementRecord{cdp:st:C.2.UN.8}{proved}{01ER}{\SeriesRef{cdp:st:C.2.MH.2}, \SeriesRef{cdp:eq:C.CP.5}, \SeriesRef{cdp:eq:C.CP.7}, \SeriesRef{cdp:eq:C.CP.6}, \SeriesRef{cdp:eq:C.CP.4}, \SeriesRef{cdp:eq:C.CP.10}; \cite[arXiv version~3, Proposition~4, pp.~10--12, and Remark~5, p.~20]{Rossi2017}}

There is a smooth projective threefold $Y$ with $K_Y\simeq\OO_Y$
and a connected-fibre map $f:Y\to\PP^1$ whose general fibre is a
product of elliptic curves. It has six singular fibres. Each is
reduced and irreducible, has rational normalization resolution, and
satisfies
\begin{equation}
 \begin{gathered}
 \pi_1(S)=1,\qquad e(S)=6,\\
 (H^i(S,\Z))_{i=0}^4=(\Z,0,\Z^4,0,\Z),\qquad
 (h^i(S,\OO_S))_{i=0}^2=(1,2,1).
 \end{gathered}
 \hypertarget{tag-00C9}{}\label{cdp:eq:C.CP.1}
\end{equation}
The ambient cohomology is
\begin{equation}
 \begin{gathered}
 (b_i(Y))_{i=0}^6=(1,0,21,8,21,0,1),\qquad e(Y)=36,\\
 (h^i(Y,\OO_Y))_{i=0}^3=(1,0,0,1),\qquad
 h^{1,1}(Y)=21,\quad h^{2,1}(Y)=3,
 \end{gathered}
 \hypertarget{tag-00CB}{}\label{cdp:eq:C.CP.2}
\end{equation}
where the Betti numbers use rational coefficients. In particular
the nonzero local character used in \ref{cdp:st:C.2.MH.2} need not exist after
the global CDP hypotheses are removed, and its Euler conclusion
$e(S)=2$ does not extend to every irreducible rational-normalized
torus degeneration.
\end{proposition}


\begin{proof}

\begin{enumerate}[label=\arabic*.,leftmargin=*]
\item The following self-product and projective small resolutions are the Namikawa construction described in
\cite[arXiv version~3, Proposition~4, pp.~10--12]{Rossi2017}. Let $B=\PP^1$, let $L=\OO_B(1)$, and take the Weierstrass
surface $g:R\to B$ with zero section and equation
\begin{equation}
                  y^2z=x^3+p(t)z^3,
             \qquad p(t)=t^6-1.                 \hypertarget{tag-00CC}{}\label{cdp:eq:C.CP.3}
\end{equation}
Here $x,y$ have the usual transition weights $L^2,L^3$;
thus $p$ is a section of $L^6$. The projective Weierstrass
bundle supplies the compactification. At infinity the transformed
coefficient is $1-t'^6$, so that fibre is smooth. The six
zeros of $p$ are simple. The only singular fibres are the six
cuspidal cubics, and the surface is smooth: in the affine
Weierstrass chart a possible singular point has $x=y=0$, while
the derivative $p'(t)$ there is nonzero. The zero-section chart
is smooth directly from its equation. Adjunction in the
Weierstrass bundle gives
\begin{equation}
                         K_R=g^*\OO_B(-1).
 \hypertarget{tag-00CD}{}\label{cdp:eq:C.CP.4}
\end{equation}
This formula can also be read from the local canonical form
$dt\wedge dx/y$ and the transition weights just specified.

\item Form the projective fibre product $X=R\times_B R$. Its only
singularities are the six points where both factors are cusp
points. At one of them use the base coordinate $\tau=p(t)$,
and write
\[
 y_1^2-x_1^3=\tau=y_2^2-x_2^3.
\]
Fix a primitive cube root $\rho$. With
\[
 \begin{gathered}
 u=y_1-y_2,\quad v=y_1+y_2,\quad
 a=x_1-x_2,\quad b=x_1-\rho x_2,\\
 c=x_1-\rho^2x_2=\alpha a+\beta b,
 \qquad \alpha=-\rho,\quad\beta=-\rho^2,
 \end{gathered}
\]
the exact threefold equation is
\begin{equation}
                         uv=abc.                \hypertarget{tag-00CE}{}\label{cdp:eq:C.CP.5}
\end{equation}
Both $\alpha$ and $\beta$ are nonzero. We resolve \eqref{cdp:eq:C.CP.5}
by two specified blowups. First blow up the ideal $(u,a)$.
In its $u$-chart, $a=ur$ and
\begin{equation}
                  v=rb(\alpha ur+\beta b),
 \hypertarget{tag-00CF}{}\label{cdp:eq:C.CP.6}
\end{equation}
so $(u,r,b)$ are smooth coordinates. In its $a$-chart put
$u=as$. The equation is $sv=b(\alpha a+\beta b)$; its
only singular point is the origin, an ordinary threefold
node. The strict transform of $(u,b)$ has ideal $(s,b)$
there. Blowing up that ideal gives the two smooth charts
\begin{equation}
 \begin{array}{c|c}
 \text{coordinates}&\text{blowdown formulas}\\ \hline
 (a,s,h)&b=sh,\quad u=as,\quad
                    v=h(\alpha a+\beta sh),\\
 (b,j,v)&s=bj,\quad a=(jv-\beta b)/\alpha,\quad u=abj.
 \end{array}                                    \hypertarget{tag-00CG}{}\label{cdp:eq:C.CP.7}
\end{equation}
The inverse image of the original singular point is a chain
of two projective lines. In the first chart of \eqref{cdp:eq:C.CP.7}
the old and new curves have equations $(a,h)=(0,0)$
and $(a,s)=(0,0)$; they meet once transversely. Charts
\eqref{cdp:eq:C.CP.6}--\eqref{cdp:eq:C.CP.7} complete each to a projective line.
The exceptional locus is this one-dimensional chain. Outside the origin
the respective ideals are Cartier, so these maps are
isomorphisms there.

\item The two ideals are the ideals of global Weil divisors.
The ideal $(u,a)$ is the ideal of the Weil diagonal
$R\hookrightarrow X$. The ideal $(u,b)$ is the Weil
graph of the global automorphism
$(x,y,t)\mapsto(\rho x,y,t)$ of $R$; after the first
blowup take its strict transform. Each of these divisors
is Cartier off the six specified singular points.
Their projective blowups therefore give a smooth
projective threefold $Y$ and a proper small map
$\pi:Y\to X$, with precisely the local charts above.
Write $p_1,p_2:X\to R$ for the two projections and $g_X=g\circ p_1=g\circ p_2:X\to B$. Set $f=g_X\circ\pi:Y\to B$. It agrees with the original product of elliptic
curves away from the six values. Its fibres are connected,
because the product fibres and every fibre of the two
blowdowns are connected.

The coordinate identification with the Weierstrass model of
\cite[arXiv version~3, \S3, (14), p.~9]{Rossi2017} is
$(x_{\rm Rossi},y_{\rm Rossi})=(y,x)$ and $B(t)=t^6-1$.
On the affine Weierstrass open containing the six singular points,
our two blowups are the closure of the graph
$\mathcal R=[u:a]$, $\mathcal S=[u:ab]$: in the $u$-chart these coordinates are
$([1:r],[1:rb])$, and in the $a$-chart the second is $[s:b]$.
The three defining equations are
\[
 \mathcal S_0v=\mathcal S_1c,\qquad \mathcal R_1u=\mathcal R_0a,\qquad
 \mathcal S_1\mathcal R_0=\mathcal S_0\mathcal R_1b.
\]
In the displayed local equations on
\cite[arXiv version~3, p.~11]{Rossi2017}, take the displayed index
$i=5$ and primitive root $\epsilon=\rho$.
Their local variables $(x,y,u,v)$ then become $(v,u,c,a)$,
and their remaining factor is
$(1+\rho)a-\rho c=b$.
Consequently their three equations are exactly the equations above
under $(\mu,\nu)=(\operatorname{sw}\mathcal S,\mathcal R)$ in local frames,
where $\operatorname{sw}[s_0:s_1]=[s_1:s_0]$.
This change has a global bundle interpretation.  With $L=\OO_B(1)$,
$\mathcal R$ takes values in $\PP(L^3\oplus L^2)$ and $\mathcal S$ in
$\PP(L^3\oplus L^4)$.  The canonical map
\[
 \tau_L:\PP(L^3\oplus L^2)
       \longrightarrow\PP(L^4\oplus L^3),\qquad
 \ell\longmapsto\ell\otimes L_b
 \quad(\ell\subset L_b^3\oplus L_b^2)
\]
has inverse tensoring by $L_b^{-1}$.  Thus the maps are
\[
 (\mathcal R,\mathcal S)\longmapsto
   \bigl(\tau_L^{-1}(\operatorname{sw}\mathcal S),\mathcal R\bigr),\qquad
 (\mu,\nu)\longmapsto
   \bigl(\nu,\operatorname{sw}(\tau_L\mu)\bigr).
\]
They are inverse over $X$, respect the three equations and agree
with the identity of $X$ away from the six singular points; hence
they glue to an isomorphism of the resolved threefolds over $X$.
For the global cyclic action (21) on p.~10, choose its primitive
root as $\rho^2$.  Its fifth and zeroth graph divisors have ideals
$(v,c)$ and $(u,a)$, respectively; their successive blowups have
coordinates $\mu=[v:c]$, $\nu=[u:a]$ and the same displayed
equations.  This also identifies the global graph centres directly.

\item The original central support at each of the six values
is the product $S_0=C\times C$ of a cuspidal cubic with
itself. It is reduced and irreducible. Since $\pi$ has
no exceptional divisors, the full fibre upstairs is its
single strict transform, still reduced. Its normalization
resolution is rational: the function field is unchanged
from that of the finite birational parametrization
$\PP^1\times\PP^1\to S_0$. 

\item Topologically a cuspidal cubic is $S^2$. Its finite
normalization is bijective, so compactness makes it a
homeomorphism. Hence $S_0$ is homeomorphic to
$S^2\times S^2$. The restriction $S\to S_0$ is an
isomorphism off the cusp-pair point and replaces that
point by the chain $E$ of two projective lines. The
quotient $S/E\to S_0$ is a continuous bijection from a
compact space to a Hausdorff space, hence a homeomorphism.
Choose a finite triangulation of $(S,E)$. Its cofibration
sequence identifies the reduced homology of $S/E$ with
the relative homology of $(S,E)$. Since
$H_2(E,\Z)=\Z^2$ and $E$ is simply connected, the
long exact sequence gives
\begin{equation}
 \begin{gathered}
 H_4(S,\Z)=\Z,\qquad H_3(S,\Z)=H_1(S,\Z)=0,\\
 0\longrightarrow\Z^2\longrightarrow H_2(S,\Z)
                 \longrightarrow\Z^2\longrightarrow0.
 \end{gathered}                                  \hypertarget{tag-00CH}{}\label{cdp:eq:C.CP.8}
\end{equation}
The last sequence splits as a sequence of abelian groups.
Van Kampen for the same cofibration shows
$\pi_1(S/E)=\pi_1(S)/\langle\!\langle\pi_1(E)\rangle\!\rangle$.
Thus $\pi_1(S)=1$. Universal coefficients now give the
integral cohomology in \eqref{cdp:eq:C.CP.1}, and its Euler number
is $1+4+1=6$.

\item The fibre product is a local complete intersection in
$R\times R$, cut out by the inverse image of the base
diagonal. Adjunction and \eqref{cdp:eq:C.CP.4} give
\begin{equation}
 \omega_X=p_1^*K_R\otimes p_2^*K_R\otimes g_X^*T_B
                         \simeq\OO_X.
 \hypertarget{tag-00CI}{}\label{cdp:eq:C.CP.9}
\end{equation}
The isomorphism over $X$ above identifies $Y$ with the Calabi--Yau resolution in \cite[arXiv version~3, proof of Corollary~2, pp.~19--20; Definition~1, p.~3]{Rossi2017}. Hence $K_Y\simeq\OO_Y$ and $h^1(Y,\OO_Y)=h^2(Y,\OO_Y)=0$, by the stated Calabi--Yau conditions and projective Hodge symmetry.
Adjunction makes $\omega_S=\OO_S$. Proper flatness over
the smooth base preserves the holomorphic Euler
characteristic, so $\chi(\OO_S)=\chi(\OO_{E_t^2})=0$.
Connectedness and Serre duality then give
$(h^0,h^1,h^2)(\OO_S)=(1,2,1)$.

\item A type-II local elliptic monodromy has characteristic
polynomial $x^2-x+1$, so $T-I$ is an integral isomorphism.
This also follows directly from the cyclic sixth-root
trivialization of $y^2=x^3+\tau$. The smooth part of
the fibration has fundamental group generated by its
fibre lattice and base meridians. Its inclusion into
$Y$ is surjective on fundamental groups, by perturbing
loops off the central divisor. The pair of zero sections
extends over every special value and kills each base
meridian in $Y$. Conjugation by one such meridian
therefore kills the image of $T-I$ on each of the two
elliptic lattices. Since those two maps are unimodular,
the whole fibre lattice is killed. Thus $\pi_1(Y)=1$.

\item For this same resolved threefold, \cite[arXiv version~3, Remark~5, p.~20]{Rossi2017} gives $h^{1,1}(Y)=21$ and $h^{2,1}(Y)=3$. Together with the Calabi--Yau conditions above, projective Hodge decomposition and duality give the ambient groups in \eqref{cdp:eq:C.CP.2}. The six computed singular-fibre supports also give $e(Y)=6\cdot6=36$ by Euler localization.

\item Restrict $f$ to a disc about any one of the six special values and apply \cref{lem:II:compact-analytic-tube}. Its inclusion and homotopy inverse identify a sufficiently small closed tube $U$ with $S$ up to homotopy. Hence $\pi_1(U)=1$ and $H^1(U,\Z)=0$. Every finite local system on $U$ is constant, and restriction to a nearby torus sends every integral degree-one class to zero. This gives the asserted obstruction to the nonzero local character of \cref{cdp:st:C.2.MH.2} in this example.

\end{enumerate}
\end{proof}



\begin{remark}[Generic divisor classes of the cuspidal self-product]\hypertarget{tag-P334}{}\label{rem:II:cuspidal-generic-divisors}
\StatementRecord{rem:II:cuspidal-generic-divisors}{proved}{P334}{\cref{cdp:st:C.2.UN.8}}
Let $K=\C(t)$, let $E_K$ be the generic elliptic curve of \eqref{cdp:eq:C.CP.3}, and fix a primitive cube root $\rho$. The two factors, the diagonal and the graph of $\rho$ on $A_K=E_K\times E_K$ span $\operatorname{NS}(A_K)\otimes\Q$. Their intersection matrix, in that order, is
\begin{equation}
 \begin{pmatrix}
 0&1&1&1\\1&0&1&1\\1&1&0&3\\1&1&3&0
 \end{pmatrix}.                                \hypertarget{tag-00CA}{}\label{cdp:eq:C.CP.10}
\end{equation}
\end{remark}
\begin{proof}
Each curve has self-intersection zero. Each factor meets the other three in one point with multiplicity one. The diagonal meets the graph of $\rho$ in the kernel of $1-\rho$; on the triangular elliptic lattice this kernel has order three. These give \eqref{cdp:eq:C.CP.10}, whose determinant is $-3$, so the four classes are independent.

For the upper bound, use $K'=K(u)$ with $u^6=p(t)$. The inverse substitutions
\[
 (X,Y)\longmapsto(u^2X,u^3Y),\qquad
 (x,y)\longmapsto(x/u^2,y/u^3)
\]
identify $E_{K'}$ with the constant elliptic curve $E_0:Y^2=X^3+1$. Base change to an algebraic closure identifies its geometric N\'eron--Severi group with that of $(E_0\times E_0)_{\overline K}$. For the complex abelian surface $E_0^2$, the first Chern class embeds this group into its integral $(1,1)$ classes, and $h^{1,1}(E_0^2)=4$. The Picard scheme and its identity component commute with algebraically closed field extension, and its connected components over $\C$ remain geometrically connected; hence this extension preserves the N\'eron--Severi group. Pullback embeds the group of $K$-defined divisor classes into that geometric group. Its rank is therefore at most four, and the four independent classes already defined over $K$ prove the assertion.
\end{proof}

\subsection{Kummer components and vertical modifications}\label{subsec:II:kummer-occurrence}

\begin{proposition}[A projective zero-Euler filling with an essential Kummer component]\hypertarget{tag-01CQ}{}\label{cdp:st:C.2.UN.13}
\StatementRecord{cdp:st:C.2.UN.13}{proved}{01CQ}{\SeriesRef{cdp:eq:C.KM.4}, \SeriesRef{cdp:eq:C.KM.5}, \SeriesRef{cdp:st:C.2.UN.4}, \SeriesRef{cdp:eq:C.KM.S2}, \SeriesRef{cdp:st:C.2.UN.3}; \cite[Chapter~1, Example~1.3(iii), pp.~8--9]{HuybrechtsK3Draft}}

Fix an abelian surface $A$. There is a smooth projective threefold
$W$ with a proper map $f_W:W\to\PP^1$ whose fibres over
$\C^*$ are isomorphic to $A$. It has precisely two special fibres,
over $0$ and $\infty$. Each has full divisor
\begin{equation}
               2P+\sum_{i=1}^{16}E_i,\qquad
 P=\operatorname{Km}(A),\quad E_i\simeq\PP^2.       \hypertarget{tag-00PI}{}\label{cdp:eq:C.KM.1}
\end{equation}
Here $P$ is the smooth Kummer K3 surface, the $E_i$ are disjoint,
and $P\cap E_i$ is a projective line on $E_i$ and a curve of
self-intersection $-2$ on $P$. The reduced divisor is simple
normal crossing and has Euler number $40$.

There is a composition of smooth vertical blowups $Y\to W$,
disjoint from both K3 components, for which both special reduced
fibres have Euler number zero. Each has $57$ irreducible
components, including its original unchanged Kummer surface.
Every blowup is vertical, so the composition has the same map on the complement of the two special fibres. Their
rational Betti numbers and coherent dimensions are
\begin{equation}
 \begin{array}{c|c|c|c}
 & (b_i)_{i=0}^6 & e & (h^i(\OO))_{i=0}^3\\ \hline
 W&(1,0,39,0,39,0,1)&80&(1,0,1,0)\\
 Y&(1,0,119,240,119,0,1)&0&(1,0,1,0).
 \end{array}                                             \hypertarget{tag-00PL}{}\label{cdp:eq:C.KM.2}
\end{equation}
For either special value, let $S_W$ and $S_Y$ denote its
reduced supports in these two models. Their cohomology is
\begin{equation}
 \begin{gathered}
 (H^i(S_W,\Q))_{i=0}^4=(\Q,0,\Q^{22},0,\Q^{17}),\\
 (H^i(S_Y,\Q))_{i=0}^4=(\Q,0,\Q^{62},\Q^{120},\Q^{57}),\\
 (h^i(S_W,\OO_{S_W}))_{i=0}^2
   =(h^i(S_Y,\OO_{S_Y}))_{i=0}^2=(1,0,1).
 \end{gathered}                                          \hypertarget{tag-00PT}{}\label{cdp:eq:C.KM.S1}
\end{equation}
Neither K3 component can be contracted by a bimeromorphic map
between smooth fillings. For comparison, after the quadratic
base change the normalization of the original quotient is
exactly $A\times\PP^1$, and the corresponding base changes
of $W$ and $Y$ are birational modifications of that product.

\end{proposition}


\begin{proof}

\begin{enumerate}[label=\arabic*.,leftmargin=*]
\item On $A\times\PP^1_s$ take the involution
\begin{equation}
       \iota(a,s)=(-a,-s),\qquad
 X=(A\times\PP^1_s)/\langle\iota\rangle,
       \qquad t=s^2.                                  \hypertarget{tag-00PM}{}\label{cdp:eq:C.KM.3}
\end{equation}
The quotient is projective and normal. Its fixed points are
exactly $A[2]\times\{0,\infty\}$, so there are $32$ of them.
Away from them the action is free. In local coordinates centred
at any fixed point, including the coordinate $s^{-1}$ at
infinity, the action is scalar multiplication by $-1$ on
$\C^3$. Thus the only singularities of $X$ are those $32$
quotient points.

\item Blow up these points in $X$, and call the resulting map
$\pi:W\to X$. Write the coordinates of a
covering chart as $(a,b,s)$. The invariant local ring is
generated by their six quadratic monomials. In the blowup
chart corresponding to $a^2$, put
\begin{equation}
             a=r,\qquad b=ru,\qquad s=rv,
             \qquad R=r^2.                            \hypertarget{tag-00PN}{}\label{cdp:eq:C.KM.4}
\end{equation}
All invariant generators pull back as
\begin{equation}
 (a^2,ab,as,b^2,bs,s^2)
           =(R,Ru,Rv,Ru^2,Ruv,Rv^2).                  \hypertarget{tag-00PO}{}\label{cdp:eq:C.KM.5}
\end{equation}
Consequently $(R,u,v)$ are smooth coordinates on the blowup.
The analogous charts for $b^2$ and $s^2$ cover it: the
exceptional projectivized invariant cone is the quadratic
Veronese image of $\PP^2$, on which at least one of the
three squared coordinates is nonzero. Equivalently, blow up
the origin of $\C^3$ first. The lifted involution fixes its
exceptional $\PP^2$ and sends its normal coordinate $r$ to
$-r$; its quotient has precisely \eqref{cdp:eq:C.KM.4}--\eqref{cdp:eq:C.KM.5}.
The transitions for $R=r^2$ identify its normal line with
$\OO_{\PP^2}(-2)$. This proves that $W$ is smooth and
projective, with exceptional divisors $E_i\simeq\PP^2$.

\item The base map in \eqref{cdp:eq:C.KM.4} is exactly
\begin{equation}
                             t=R v^2.                 \hypertarget{tag-00PP}{}\label{cdp:eq:C.KM.6}
\end{equation}
Here $E_i=(R=0)$ and the strict central component is
$P=(v=0)$. In the chart where $s$ is the exceptional
coordinate the map is $t=R$ and only $E_i$ occurs.
These charts prove the multiplicities and the transverse
intersection in \eqref{cdp:eq:C.KM.1}. They also show that
$N_{E_i/W}|_{P\cap E_i}=\OO_{\PP^1}(-2)$, which is
the normal line of that intersection in $P$.
The same calculation applies over infinity. Off those two
values, choose either square root $s$ of $t$; the map
$a\mapsto[a,s]$ identifies the fibre with $A$, and switching
the square root composes this identification with $a\mapsto-a$.
In particular the map has connected fibres and its monodromy
on $H^1(A,\Z)$ is $-I_4$.

\item The restriction $P\to A/\{\pm1\}$ resolves exactly the
sixteen quotient surface nodes. It is the classical Kummer
construction; see \cite[Chapter~1, Example~1.3(iii), pp.~8--9]{HuybrechtsK3Draft}.
The two-form and coherent cohomology maps in these charts are as follows. A translation-invariant two-form on $A$ is
invariant under $-1$. In a surface chart $a=r$, $b=ru$,
its pullback is
\begin{equation}
                da\wedge db=r\,dr\wedge du
                            =\tfrac12\,dR\wedge du.   \hypertarget{tag-00PQ}{}\label{cdp:eq:C.KM.7}
\end{equation}
It therefore descends to a nowhere-zero form on $P$.
The finite covering $\operatorname{Bl}_{A[2]}A\to P$
injects $H^1(P,\OO_P)$ into the invariant part of
$H^1(A,\OO_A)$, by the averaging trace. This invariant
part is zero because $-1$ acts as $-1$ on $H^1(A,\OO_A)$.
Thus $K_P=\OO_P$ and $h^1(P,\OO_P)=0$, verifying that
$P$ is a K3 surface.

\item The finite quotient map identifies rational cohomology with
invariants. Explicitly, its higher direct images of the
constant rational sheaf vanish by proper base change and
the finiteness of its fibres; averaging identifies the
constant sheaf downstairs with the invariant direct image.
Applied to $A/\{\pm1\}$, this gives Betti numbers
$(1,0,6,0,1)$ and Euler number eight. Resolving each of its
sixteen points replaces that point by $\PP^1$, so
$e(P)=8+16=24$. Hence inclusion--exclusion in \eqref{cdp:eq:C.KM.1} gives
\begin{equation}
          e\left(P\cup\bigcup_{i=1}^{16}E_i\right)
                    =24+16\cdot3-16\cdot2=40.          \hypertarget{tag-00PR}{}\label{cdp:eq:C.KM.8}
\end{equation}
Mayer--Vietoris gives the support groups as well. The restriction
map from the direct sum of $H^2(P,\Q)$ and the sixteen
$H^2(E_i,\Q)$ onto the sixteen $H^2(P\cap E_i,\Q)$ is
surjective: the hyperplane class on each $E_i$ restricts to
the degree-one class of its intersection line. Its kernel
has dimension $22+16-16=22$, and the same exact sequence
gives $H^1(S_W,\Q)=H^3(S_W,\Q)=0$ and
$H^4(S_W,\Q)=\Q^{17}$. The coherent sequence
\[
 0\longrightarrow\OO_{S_W}\longrightarrow
 \OO_P\oplus\bigoplus_i\OO_{E_i}
 \longrightarrow\bigoplus_i\OO_{P\cap E_i}
 \longrightarrow0
\]
has a surjective difference map on $H^0$. The lines and planes
have no positive coherent cohomology, so it gives
$(h^0,h^1,h^2)(\OO_{S_W})=(1,0,1)$.

\item For the finite
quotient $A\times\PP^1\to X$, the same invariant-sheaf
argument gives
\begin{equation}
 H^k(X,\Q)\simeq H^k(A\times\PP^1,\Q)^{\iota},
 \qquad (b_k(X))_{k=0}^6=(1,0,7,0,7,0,1).              \hypertarget{tag-00PS}{}\label{cdp:eq:C.KM.9}
\end{equation}
Indeed $-1$ acts by $(-1)^j$ on $H^j(A)$, while the
base involution acts trivially on both nonzero cohomology
groups of $\PP^1$.

\item Let $j_i:E_i\hookrightarrow W$ and let
$h_i=c_1(\OO_{E_i}(1))$. Collapsing the $32$ exceptional
divisors separately gives exactly $X$. The relative
cohomology sequence of this quotient identifies the
kernel of restriction to the $E_i$ with the pullback
of $H^k(X,\Q)$. The self-intersection formula gives
\begin{equation}
 j_i^*j_{i*}(1)=-2h_i,
       \qquad j_i^*j_{i*}(h_i)=-2h_i^2.                 \hypertarget{tag-00PJ}{}\label{cdp:eq:C.KM.10}
\end{equation}
Thus restriction is surjective in degrees two and four,
with explicit rational sections $-j_{i*}(1)/2$ and
$-j_{i*}(h_i)/2$. In these two degrees pullback is
injective because the preceding odd cohomology of the
exceptional divisors vanishes. The same relative
sequence then leaves the other positive-degree groups
unchanged. Each exceptional point consequently adds one
class in degrees two and four and none in any other
degree. This proves the row for $W$ in \eqref{cdp:eq:C.KM.2}.

\item For coherent cohomology the finite quotient and averaging
give
$H^k(X,\OO_X)=H^k(A\times\PP^1,\OO)^{\iota}$,
of dimensions $(1,0,1,0)$. The local resolution just
used has $R^q\pi_*\OO_W=0$ for $q>0$. Filter the
infinitesimal neighbourhoods of $E_i$ by its ideal.
Their successive quotients are $\OO_{\PP^2}(2r)$,
$r\ge0$, whose positive-degree coherent cohomology
vanishes. Formal functions therefore makes the completed
stalks of those higher direct images zero, and coherence
makes the stalks themselves zero. Normality gives
$\pi_*\OO_W=\OO_X$. Leray proves the coherent row for
$W$.

\item We now perform exactly twenty of the negative-two
operations of \ref{cdp:st:C.2.UN.4} over each of $0$ and $\infty$.
Every operation starts at a smooth point on the
strict transform of one fixed $E_i$, away from its
intersection with $P$ and from all other components.
The first blowup creates a new $\PP^2$, and the second
blows up a smooth plane quartic of genus three inside
it, transverse to its old intersection line. These
centres and their entire modifications are disjoint
from $P$. The charts of \cref{cdp:st:C.2.UN.4} preserve simple
normal crossings and show that both new components
have multiplicity one. Each operation changes the
reduced-fibre Euler number by $-2$ and adds two
components. Thus both final supports have Euler number
$40-20\cdot2=0$ and $17+20\cdot2=57$ components,
with their unchanged multiplicity-two K3 surfaces.

\item For one of these blowups write $\sigma:\widehat W\to W'$, let $C$ be its smooth centre of ambient codimension $c$, and let $S$ be the current reduced support. Apply \cref{lem:II:blowup-topology} with
\[
 (Y,Z,r,\pi_Z,S,S')
  =(W',C,c,\sigma,S,\sigma^{-1}(S)_{\rm red}).
\]
Its restricted ambient Gysin map in \eqref{cdp:eq:C.KM.S2} identifies the cohomology of this full inverse image. This includes the four intersections of the quartic with the other support component.
A point step therefore adds one class in degrees two and four.
The quartic step adds ranks one, six and one in degrees two,
three and four. After twenty blocks, \eqref{cdp:eq:C.KM.S2} gives exactly
the second support row of \eqref{cdp:eq:C.KM.S1}.

\item Each chosen centre is generically contained in just one
reduced component, of multiplicity one in the reduced support.
Consequently the equality of Cartier divisors
$\sigma^*S=S'$ holds, also at the four crossing points as
the explicit equations show. Apply derived pushforward to
$0\to\OO_{\widehat W}(-S')\to\OO_{\widehat W}\to
\OO_{S'}\to0$. The projection formula and
$R\sigma_*\OO_{\widehat W}=\OO_{W'}$ give
\begin{equation}
                         R\sigma_*\OO_{S'}=\OO_S.       \hypertarget{tag-00PV}{}\label{cdp:eq:C.KM.S3}
\end{equation}
Thus all support coherent cohomology groups are unchanged,
proving the remaining assertion of \eqref{cdp:eq:C.KM.S1}.

\item The pullback and exceptional Gysin maps in \ref{cdp:st:C.2.UN.4}
give cohomology increments $(\Delta b_2,\Delta b_3,
\Delta b_4)=(2,6,2)$ per operation. Forty operations
give the second Betti row in \eqref{cdp:eq:C.KM.2}. Smooth-centre
blowups preserve coherent cohomology, by their direct
image $\OO$ and vanishing higher direct images, so
the second coherent row is unchanged. This also checks
$e(Y)=0$ directly from its displayed Betti numbers.

\item Finally, $P$ is not uniruled. A dominating ruled surface
would pull its nonzero holomorphic two-form back to a
nonzero two-form, whereas a ruled surface has no such
form. Apply \cref{lem:II:nonuniruled-divisors} to this compact central divisor in the given smooth threefold filling. Its strict-transform correspondence prevents contraction under a bimeromorphic comparison of smooth fillings. For the quadratic base change, the map
\begin{equation}
 A\times\PP^1_s\longrightarrow
 X\times_{\PP^1_t}\PP^1_s,
         \qquad(a,s)\longmapsto([a,s],s)               \hypertarget{tag-00PK}{}\label{cdp:eq:C.KM.11}
\end{equation}
is finite and generically one-to-one. Its source is
normal, so it is the normalization of the reduced
base change. The normalized base changes of $W$ and
$Y$ factor properly and birationally through this
normalization. This specifies exactly how the Kummer
component is related to a torus after that base change.

\end{enumerate}
\end{proof}






\subsection{Projective covers of signed elliptic roots}\label{subsec:II:projective-root-covers}

\begin{lemma}[Finite product covers of torsion signed roots]\hypertarget{tag-P347}{}\label{lem:II:torsion-root-product-cover}
\StatementRecord{lem:II:torsion-root-product-cover}{proved}{P347}{\cref{lem:II:signed-root-fibration}; \cite[Theorem~2, p.~241]{ColtoiuGasitoiJoita2012}}
Let $B$ be a smooth compact connected curve and let $D_+,D_-$
be disjoint reduced divisors with $D_++D_-$ nonempty.
Suppose $L=\OO_B(D_+-D_-)$ is torsion, and let $\sigma$
be its canonical meromorphic section. Fix $m\ge1$ and
$0<|q|<1$, and write $E_q=\C^*/q^{\Z}$.
Apply the signed-root equations of \cref{lem:II:signed-root-fibration}
on $B$, or equivalently restrict its construction on $B\times\Delta$
with divisors $D_\pm\times\Delta$ to $B\times\{0\}$.
The resulting smooth compact elliptic surface $P\to B$
has a finite unramified cover by a product $C'\times E_q$,
and is projective.

Choose a torsion line $N$ and an isomorphism $N^{\otimes m}\simeq L$.
The normalized cyclic curve $C\to B$ defined by
\begin{equation}
                         y^m=\sigma
\hypertarget{tag-00U4}{}\label{cdp:eq:C.PE.11}
\end{equation}
is connected and smooth, with
$2g(C)-2=m(2g(B)-2)+(m-1)\deg(D_++D_-)$.
There is a connected finite unramified cover $C'\to C$
on which the pullback of $N$ has a nowhere-zero frame $n$.
Writing $y=y_0n$, the product covering is
\begin{equation}
 \begin{split}
 \Phi:C'\times E_q&\longrightarrow P,\\
 (z,[u])&\longmapsto
       [\ell=u^m n(z)^m,\ a=y_0(z)/u].
 \end{split}
\hypertarget{tag-00U5}{}\label{cdp:eq:C.PE.12}
\end{equation}
At a pole it is read in the reciprocal chart
$\ell'=u^{-m}n^{-m}$, $a'=u/y_0$.
\end{lemma}
\begin{proof}
A torsion line has degree zero. Multiplication by $m$ on the
torsion subgroup of the Jacobian of $B$ is surjective, giving
$N$ and the displayed isomorphism. Normalize the cyclic equation
in $\PP(\OO_B\oplus N)$. In local frames it is $y^m=x$
at a positive point and $(y^{-1})^m=x$ at a negative point.
Thus its normalized curve is smooth, fully ramified at these
points and unramified elsewhere. A loop about any one of these
points cyclically permutes all $m$ sheets, proving connectedness.
The degree-$m$ Riemann--Hurwitz formula gives the stated genus.

If $N$ has order $d$, a chosen trivialization $N^d\simeq\OO_B$
defines the finite unramified $\mu_d$-cover of local frames with
$d$th power one. Pull it back to $C$ and take a connected
component. Every such component maps surjectively to $C$;
its tautological frame gives the required $n$ on $C'$.

The equations in \eqref{cdp:eq:C.PE.12} satisfy
$\sigma=a^m\ell$ because $y^m=\sigma$.
Replacing $u$ by $qu$ applies precisely the deck map
$(\ell,a)\mapsto(q^m\ell,q^{-1}a)$.
The reciprocal expressions satisfy the other signed equation
and agree on the common open set. Thus they define the stated
holomorphic map on the entire product. Away from the marked
points, each sheet of $C'\to B$ maps its elliptic fibre
isomorphically to the fibre of $P$. At a marked point the map
of reduced fibres is $[u]\mapsto[u^m]$ from $E_q$ to
$\C^*/q^{m\Z}$, so it is surjective with finite fibres there too.
Compactness makes this surjective map proper and finite.

Absorb the local nowhere-zero frame factor into $u$ and use
the root parameter $\zeta$ on $C'$. In the positive chart,
the map to the smooth surface coordinates $(a,b)$ is
\begin{equation}
   (\zeta,u)\longmapsto(\zeta/u,u^m),\qquad
                    \det d\Phi=m u^{m-2}\ne0.
\hypertarget{tag-00U6}{}\label{cdp:eq:C.PE.13}
\end{equation}
In the negative chart it is
$(\zeta,u)\mapsto(a'=\zeta u,b'=u^{-m})$, with
Jacobian $-m u^{-m}\ne0$. At all other points the base map
is unramified and the fibre map is an isomorphism. This proves
that $\Phi$ is unramified everywhere. Apply
\cite[Theorem~2, p.~241]{ColtoiuGasitoiJoita2012} with
$(X,Y,p)=(C'\times E_q,P,\Phi)$. Its source is projective,
its target is smooth and compact, and every fibre of this
surjective map is zero-dimensional. The theorem gives
projectivity of $P$.
\end{proof}

\subsection{Paired roots and projective elliptic components}\label{subsec:II:paired-root-occurrence}

\begin{proposition}[Opposite logarithmic data produce a projective essential component]\hypertarget{tag-01CV}{}\label{cdp:st:C.2.UN.18}
\StatementRecord{cdp:st:C.2.UN.18}{proved}{01CV}{\SeriesRef{cdp:eq:C.PE.5}, \SeriesRef{cdp:eq:C.PE.6}, \SeriesRef{cdp:eq:C.PE.7}, \SeriesRef{cdp:eq:C.PE.12}, \SeriesRef{cdp:eq:C.PE.8}, \SeriesRef{cdp:eq:C.PE.9}, \SeriesRef{cdp:eq:C.PE.16}, \SeriesRef{cdp:eq:C.PE.18}, \cref{lem:II:signed-root-fibration,lem:II:hopf-curve-restriction,lem:II:nonuniruled-divisors}, \SeriesRef{cdp:st:C.2.UN.6}; \cite[Example~3.4(1), p.~99]{Fujimoto1988Log}}

Put $q=e^{-2\pi}$ and $E=\C^*/q^{\Z}$. Choose distinct
points $p_+,p_-\in E$ whose difference is torsion.
For every $m\ge2$ the fixed punctured family
$(E\times E)\times\Delta^*$ has a smooth proper filling
with full central fibre
\begin{equation}
                         P_m+mQ_++mQ_-.                 \hypertarget{tag-00U2}{}\label{cdp:eq:C.PE.1}
\end{equation}
Here $Q_+$ and $Q_-$ are disjoint primary Hopf surfaces,
$P_m$ is a smooth projective surface of Kodaira dimension
one, and $J_\pm=P_m\cap Q_\pm$ are disjoint smooth
elliptic curves. The support is a simple normal crossing
surface of Euler number zero. The map $P_m\to E$ is an
elliptic fibration with exactly the two multiple fibres
$mJ_+$ and $mJ_-$, and
\begin{equation}
 \begin{gathered}
 K_{P_m}\simeq\OO_{P_m}((m-1)(J_++J_-)),\qquad
 H_1(P_m,\Z)\simeq\Z^4,\\
 h^0(P_m,K_{P_m}^{\otimes r})=
 \begin{cases}1,&d_r=0,\\d_r,&d_r>0,\end{cases}
 \qquad d_r=2\left\lfloor\frac{r(m-1)}m\right\rfloor.
 \end{gathered}                                         \hypertarget{tag-00UD}{}\label{cdp:eq:C.PE.2}
\end{equation}
The surface $P_m$ admits a finite \'etale cover by a
product of two projective curves, one elliptic and the
other of genus at least two. All three components contain
no rational curve and survive birationally in every
bimeromorphic comparison between smooth fillings. The
multiplicity-one component $P_m$ also survives birationally
after every finite base change and resolution. The
fillings for different $m$ are not bimeromorphic over the
disc.

The construction extends to a smooth compact threefold
$Y_m\to\PP^1$ with just this special fibre. All integral
cohomology groups below are free. Writing $S_m$ for its
reduced special support, their ranks and coherent
cohomology dimensions are
\begin{equation}
 \begin{array}{c|c|c}
 &\operatorname{rank}H^k(-,\Z)&h^k(\OO)\\ \hline
 P_m&(1,4,6,4,1)&(1,2,1)\\
 Q_\pm&(1,1,0,1,1)&(1,1,0)\\
 S_m&(1,3,6,7,3)&(1,2,1)\\
 Y_m&(1,3,7,10,7,3,1)&(1,2,1,0).
 \end{array}                                           \hypertarget{tag-00UG}{}\label{cdp:eq:C.PE.3}
\end{equation}
\end{proposition}


\begin{proof}

\begin{enumerate}[label=\arabic*.,leftmargin=*]
\item The logarithmic-transform occurrence in \cite[Example~3.4(1), p.~99]{Fujimoto1988Log} has a Kodaira-dimension-one component. Here the two opposite signs give the following projective component and fixed punctured family. Let $V$ be the blowup of $E\times\Delta$ at $(p_+,0)$ and $(p_-,0)$, with exceptional curves $R_+,R_-$. Put
\begin{equation}
 D=R_+-R_-,\qquad L=\OO_V(D),\qquad
 \operatorname{div}(\sigma)=D,
 \hypertarget{tag-00UH}{}\label{cdp:eq:C.PE.4}
\end{equation}
and apply \cref{lem:II:signed-root-fibration} with
$(D_+,D_-,m,q)=(R_+,R_-,m,e^{-2\pi})$.
Its two equations \eqref{cdp:eq:C.PE.5}--\eqref{cdp:eq:C.PE.6}
are the equations on these same line bundles and sections.
The identity on $(v,\ell,a)$, using $(\ell^{-1},a^{-1})$
in the negative chart, intertwines the gluing and deck maps.
Thus its quotient is the smooth threefold $U_m$ and its proper
projection is $\rho:U_m\to V$. Formula \eqref{cdp:eq:C.PE.9}
identifies its restriction over $\Delta^*$ with
$(E\times E)\times\Delta^*$, with inverse $[c]=[a^{-1}]$.

\item Write $B\simeq E$ for the strict central curve of $V$ and $p:P_m\to B$ for the restriction of $\rho$.
The base function has divisor $B+R_++R_-$. At the two
intersections it has equation $t=xy$, so its pullbacks
are respectively
\begin{equation}
                   t=a^mby,\qquad t=(a')^m b'y.         \hypertarget{tag-00U3}{}\label{cdp:eq:C.PE.10}
\end{equation}
These are transverse products with multiplicities
$m,1$. They prove \eqref{cdp:eq:C.PE.1}, smoothness of $P_m$ and
the stated intersections. Over $R_+$ the component is
$L|_{R_+}^{\times}/q^{m\Z}$; over $R_-$ it is
$(L^{-1}|_{R_-})^{\times}/q^{-m\Z}$. Both line bundles
are $\OO_{\PP^1}(-1)$. Sending a nonzero tautological
vector to the same vector in $\C^2\setminus0$ identifies
both components with
$(\C^2\setminus0)/q^{m\Z}$, a primary Hopf surface.
The two reduced multiple fibres $J_\pm$ are each
$\C^*/q^{m\Z}$.

\item Apply \cref{lem:II:torsion-root-product-cover} with
$B=E$, $D_+=[p_+]$, $D_-=[p_-]$, the same $m,q$, and
$L_B=\OO_E(p_+-p_-)$. This line is torsion because the
difference $p_+-p_-$ is a torsion point of $E$.
Choose the lemma's torsion $m$th root $N$ with its isomorphism
$N^{\otimes m}\simeq L_B$. Its cyclic curve is
\eqref{cdp:eq:C.PE.11}; its genus formula gives $g(C)=m$.
The same cover $C'\to C$, frame $n$, and meromorphic
coordinate $y=y_0n$ give exactly \eqref{cdp:eq:C.PE.12}.
The identity on $(\ell,a)$ intertwines its positive chart
with the restriction of $\rho$, and the reciprocal
$(\ell^{-1},a^{-1})$ does so at $p_-$. Thus the target
is this $P_m$, and the lemma gives its finite unramified
product cover and projectivity. Both branch Jacobians
are those of \eqref{cdp:eq:C.PE.13} with these same frames.

\item Let $\omega$ be a nowhere-zero one-form on $E$.
Off the multiple fibres, the form
$p^*\omega\wedge d\log\ell$ is well defined: a change
of frame adds a base one-form to $d\log\ell$, whose
wedge with $p^*\omega$ is zero. Near $J_+$, write
$\omega=h(x)dx$ and $x=a^m b$; then this two-form is
\begin{equation}
                m h(a^m b)a^{m-1}\,da\wedge db.         \hypertarget{tag-00U7}{}\label{cdp:eq:C.PE.14}
\end{equation}
Near $J_-$ the dual formula has a minus sign and the
same vanishing order $m-1$. It has no further zero.
This proves the canonical line in \eqref{cdp:eq:C.PE.2}. Constancy
on compact elliptic fibres and the invariant local
root charts give the exact direct image
\begin{equation}
 p_*\OO_{P_m}(k_+J_++k_-J_-)
 =\OO_E\!\left(
  \left\lfloor\frac{k_+}m\right\rfloor p_+
 +\left\lfloor\frac{k_-}m\right\rfloor p_-\right).         \hypertarget{tag-00U8}{}\label{cdp:eq:C.PE.15}
\end{equation}
A permitted Laurent term in a root parameter descends
precisely when its order is divisible by $m$.
Substitute $k_+=k_-=r(m-1)$ and use Riemann--Roch
on the elliptic curve. This proves the plurigenera
and $\kappa(P_m)=1$.

\item We give the integral comparisons in \eqref{cdp:eq:C.PE.3}. The
real action
$(\ell,a)\mapsto(\lambda^m\ell,\lambda^{-1}a)$,
$\lambda>0$, preserves the equations and translates
$\log\|\ell\|$. Since $q$ is positive real, its
quotient gives differentiable products
\begin{equation}
                   P_m\simeq M_3\times S^1,\qquad
                   Y_m\simeq M_5\times S^1,             \hypertarget{tag-00U9}{}\label{cdp:eq:C.PE.16}
\end{equation}
where the $M$'s are the level sets $\|\ell\|=1$ in
the corresponding covering spaces. For $Y_m$ use
the construction over
$V=\operatorname{Bl}_{(p_+,0),(p_-,0)}(E\times\PP^1)$.
The local proper maps \eqref{cdp:eq:C.PE.8} make this quotient
compact and smooth. Off $0$ the product comparison
\eqref{cdp:eq:C.PE.9} applies, including at infinity.

\item For $M_3$, remove small discs around $p_+,p_-$.
Over their complement the angular coordinate of
$a$ trivializes its circle fibration. Let $h$ be
the positive angular $a$-loop and let $\delta_\pm$
be positive base boundary loops with $a$ fixed.
In the positive solid torus the phases satisfy
$\delta_+=m\arg(a)+\arg(b)$; in the negative one
they satisfy $\delta_-=m\arg(a')+\arg(b')$ and
$\arg(a)=-\arg(a')$. The two meridians are
\begin{equation}
                         m\delta_++h,\qquad
                         m\delta_- -h.                 \hypertarget{tag-00UA}{}\label{cdp:eq:C.PE.17}
\end{equation}
In the abelianization of the punctured elliptic base,
$\delta_++\delta_-=0$. These relations give
$h=-m\delta_+$ and $\delta_-=-\delta_+$, leaving
$H_1(M_3,\Z)=\Z^3$, generated by the two
elliptic-base loops and $\delta_+$. There is no
torsion. Poincar\'e duality on this closed oriented
three-manifold and \eqref{cdp:eq:C.PE.16} give all the asserted
integral groups of $P_m$.

\item For the compact five-manifold put
$Q=(E\times\PP^1)\setminus(B^4_+\sqcup B^4_-)$,
with two disjoint open balls. The exterior is
$Q\times S^1$, while each filling is
$S^3\times D^2$. Near an exceptional curve regard
a nonzero tautological vector as
$w\in\C^2\setminus0$. On the unit norm slice
$w\in S^3$, and the blowdown is
$u=a^m w$ or $u=(a')^m w$. The boundary maps are
\begin{equation}
 \begin{split}
 (w,e^{\mathrm i\theta})&\longmapsto
       (e^{\mathrm i m\theta}w,e^{\mathrm i\theta}),\\
 (w,e^{\mathrm i\theta})&\longmapsto
       (e^{\mathrm i m\theta}w,e^{-\mathrm i\theta}).
 \end{split}                                             \hypertarget{tag-00UB}{}\label{cdp:eq:C.PE.18}
\end{equation}
The pair sequence for the two removed balls gives
free homology of $Q$ of ranks $(1,2,2,3,0)$; its
two boundary sphere classes are opposite
primitive classes $\delta,-\delta$ in $H_3(Q,\Z)$.
In Mayer--Vietoris for \eqref{cdp:eq:C.PE.18}, the map on
$H_1(S^3\times S^1)^{\oplus2}$ has primitive
rank-one image, the circle generator of the
exterior. The map on $H_3$ is injective with
primitive image, since each of the two filling
$S^3$ summands receives its own identity map.
The map on $H_4$ has primitive rank-one image
generated by $\delta\times[S^1]$. Every kernel
and cokernel used in the sequence is therefore
free, and it gives
\begin{equation}
       \operatorname{rank}H_k(M_5,\Z)=(1,2,5,5,2,1).
                                                               \hypertarget{tag-00UC}{}\label{cdp:eq:C.PE.19}
\end{equation}
K\"unneth and \eqref{cdp:eq:C.PE.16} give the $Y_m$ row of
\eqref{cdp:eq:C.PE.3}, including integral torsion-freeness.

\item Each Hopf surface is differentiably $S^3\times S^1$.
For the support use the union
$S_m=P_m\cup Q_+\cup Q_-$ with intersections
$J_+\sqcup J_-$. In the preceding $M_3$ presentation,
the core angular loop of $J_+$ is $\delta_+$ and
that of $J_-$ is $-\delta_+$. Their radial loops
are the $S^1$ factor of \eqref{cdp:eq:C.PE.16}. Inclusion in
$Q_\pm$ kills the angular loop and takes the
radial loop to its Hopf generator. Consequently
the integral map
\begin{equation}
 \begin{split}
 H_1(J_+\sqcup J_-,\Z)\longrightarrow
 H_1(P_m,\Z)\oplus H_1(Q_+,\Z)\oplus H_1(Q_-,\Z)
 \end{split}                                             \hypertarget{tag-00UE}{}\label{cdp:eq:C.PE.20}
\end{equation}
has primitive rank three: its columns are, up to
orientation signs, $\delta_+$, $z-q_+$,
$-\delta_+$, $z-q_-$. On $H_2$, the two elliptic
fundamental classes map to the same primitive
fibre class of $P_m$ and to zero on the Hopf
surfaces. Primitivity follows directly from
$[J_\pm]=\pm\delta_+\times[S^1]$ with the appropriate
product orientations. More explicitly, the dual chart at
$J_-$ reverses both its angular loop and its radial loop:
its complex orientation is
$(-\delta_+)\times(-z)=\delta_+\times z$.
Using the same global radial generator $z$ at both
ends instead reverses the displayed product orientation
at $J_-$. The restriction map on
cohomology thus has rank one in degree two and
rank three in degree one. Mayer--Vietoris gives
the free support ranks $(1,3,6,7,3)$ and Euler zero.

\item Apply \eqref{cdp:eq:C.PE.21} to this same compact $V$ and $\rho:Y_m\to V$. Its frame is the additive cocycle of the identical generator \eqref{cdp:eq:C.PE.7}. Since $h^k(V,\OO_V)=(1,1,0)$, the two nonzero Leray rows give $h^k(Y_m,\OO_{Y_m})=(1,2,1,0)$: every possible differential has target in $H^j(V,\OO_V)=0$ for $j\ge2$. Restrict the same finite product charts to the strict curve $B=E$. Their base changes are the charts of $P_m\to E$, so the invariant direct-image calculation and this deck frame give $\OO_E$ in degrees zero and one. Leray on the curve gives $h^k(P_m,\OO_{P_m})=(1,2,1)$.

\item Apply \cref{lem:II:hopf-curve-restriction} to
$Q_\pm=(\C^2\setminus0)/q^{m\Z}$ and to the same axis
inclusions $J_\pm\hookrightarrow Q_\pm$. It gives
$h^\bullet(Q_\pm,\OO)=(1,1,0)$ and restriction isomorphisms
$H^1(Q_\pm,\OO)\to H^1(J_\pm,\OO)$, represented by the
cocycle with value one on the generator $q^m$.
The coherent normalization sequence for
$S_m$ therefore has a surjective difference
map in degrees zero and one. Its remaining
terms give $h^k(S_m,\OO)=(1,2,1)$.

\item Finally, a rational curve in $P_m$ would lift
through the finite \'etale map \eqref{cdp:eq:C.PE.12}, since
$\PP^1$ is simply connected. Its two projections
to the positive-genus curves would be constant.
The Hopf components have no rational curve
either: a map from $\PP^1$ lifts to
$\C^2\setminus0$, and its coordinate functions
are constant. Apply \cref{lem:II:nonuniruled-divisors} to these compact non-uniruled central divisors of the smooth proper disc filling $U_m$. Its strict-transform correspondence carries each surface birationally onto a divisor in every smooth bimeromorphic filling. A
finite base change $t=s^d$ is generically,
along $P_m$, the graph $x=s^d$ for the local
transverse coordinate $x=t$. Its normalized
strict component maps with degree one to
$P_m$; resolution preserves that birational
component. The distinct plurigenus slopes
$2(1-1/m)$ distinguish the unique properly
elliptic components, and hence the fillings
for different $m$.

\end{enumerate}
\end{proof}






\section{Kodaira, Enriques and Kummer components}\label{sec:II:kodaira-enriques-occurrence}

\subsection{Secondary Kodaira components}\label{subsec:II:secondary-kodaira-occurrence}

\begin{proposition}[A free half-period quotient realizes secondary Kodaira normalization]\hypertarget{tag-01CZ}{}\label{cdp:st:C.2.UN.21}
\StatementRecord{cdp:st:C.2.UN.21}{proved}{01CZ}{\cref{lem:II:signed-root-fibration,lem:II:hopf-curve-restriction,lem:II:nonuniruled-divisors}, \SeriesRef{cdp:eq:C.SKD.4}, \SeriesRef{cdp:eq:C.SKD.5}, \SeriesRef{cdp:eq:C.SKD.7}, \SeriesRef{cdp:eq:C.SKD.9}, \SeriesRef{cdp:eq:C.SKD.6}, \SeriesRef{cdp:st:C.2.UN.18}, \SeriesRef{cdp:eq:C.SKD.12}}

Let $q=e^{-2\pi}$ and $E=\C^*/q^{\Z}$, and use the
origin of its group structure. There is a smooth compact
threefold $Y\to\PP^1$ whose smooth general fibre is
$E\times E$ and whose two special full fibres are
\begin{equation}
                        2P+2Q\quad\hbox{and}\quad 2B.
                                                            \hypertarget{tag-011N}{}\label{cdp:eq:C.SKD.1}
\end{equation}
The surface $P$ is a secondary Kodaira surface, $Q$ is
a primary Hopf surface, and they meet transversely along
a smooth elliptic curve $J$. The surface $B$ is smooth
bielliptic. Both reduced special supports have Euler
number zero. In particular $P$ occurs as an essential
component in a zero-Euler torus filling.

The canonical bundle of $P$ has exact order two. For
$S=P\cup Q$, the rational cohomology ranks and coherent
dimensions are
\begin{equation}
 \begin{array}{c|c|c}
 &\dim H^k(-,\Q)&h^k(\OO)\\ \hline
 P&(1,1,0,1,1)&(1,1,0)\\
 Q&(1,1,0,1,1)&(1,1,0)\\
 J&(1,2,1)&(1,1)\\
 S&(1,1,1,3,2)&(1,1,0)\\
 B&(1,2,2,2,1)&(1,1,0)\\
 Y&(1,1,2,4,2,1,1)&(1,1,0,0).
 \end{array}                                             \hypertarget{tag-011U}{}\label{cdp:eq:C.SKD.2}
\end{equation}
Each full scheme fibre in \eqref{cdp:eq:C.SKD.1} has coherent dimensions $(1,2,1)$.
The primary Kodaira double cover of $P$ is recovered
by the quadratic base change at the first
special value.
\end{proposition}


\begin{proof}

\begin{enumerate}[label=\arabic*.,leftmargin=*]
\item Apply \cref{lem:II:signed-root-fibration} to the compact surface and line bundle
\begin{equation}
 V=\operatorname{Bl}_{(0,0)}(E\times\PP^1_t),\qquad
 L=\OO_V(R),\qquad
 \widetilde Y=L^\times/q^{\Z},
                                                            \hypertarget{tag-011V}{}\label{cdp:eq:C.SKD.3}
\end{equation}
where $R$ is the exceptional curve, with $(D_+,D_-,m)=(R,\varnothing,1)$ and the same $q$. Let $\sigma$ be the canonical section of $L$. Projection from the hypersurface $\sigma=a\ell$ to $L^\times$ has inverse $\ell\mapsto(\ell,a=\sigma/\ell)$. These maps commute with the deck action and identify the constructed quotient with the displayed $\widetilde Y$. The projection
$\widetilde Y\to V\to\PP^1_t$ is proper and smooth
off its reduced central divisor
$\widetilde P+\widetilde Q$. The first component is
the primary Kodaira surface
$\OO_E([0])^\times/q^{\Z}$, the second is the Hopf
surface $\OO_{\PP^1}(-1)^\times/q^{\Z}$, and their
intersection is an elliptic curve. The restrictions of this same quotient to the strict elliptic curve and the exceptional curve give these two components.

\item The involution $(p,t)\mapsto(-p,-t)$ fixes the blowup
centre and lifts to an involution $\iota$ of $V$.
It has the canonical order-two lift $\widehat\iota$
to $L$ preserving $\sigma$. Define an automorphism
of $\widetilde Y$ by
\begin{equation}
             g[\ell]=[q^{1/2}\widehat\iota(\ell)].
                                                            \hypertarget{tag-011W}{}\label{cdp:eq:C.SKD.4}
\end{equation}
Its square is multiplication by $q$, so $g^2=1$
on $\widetilde Y$. It acts freely. Indeed at a
fixed base point the order-two linearization on
the line $L_v$ has eigenvalue $+1$ or $-1$.
Neither $q^{1/2}$ nor $-q^{1/2}$ belongs to
$q^{\Z}$. Thus \eqref{cdp:eq:C.SKD.4} has no fixed point on
any elliptic fibre, including those over $R$.
The quotient
\begin{equation}
              \pi:\widetilde Y\longrightarrow Y
                        =\widetilde Y/\langle g\rangle
                                                            \hypertarget{tag-011X}{}\label{cdp:eq:C.SKD.5}
\end{equation}
is finite \'etale of degree two, and $Y$ is smooth
and compact. The function $T=t^2$ defines its map
to $\PP^1_T$. For $T\ne0,\infty$, the two values
of $t$ are interchanged by $g$, so each quotient
fibre is exactly one of the original product
tori.

\item The strict central base curve and $R$ are
invariant. The local equation upstairs at their
intersection is $t=xy$, hence $T=x^2y^2$.
Since \eqref{cdp:eq:C.SKD.5} is \'etale, this proves the first
full fibre and the normal-crossing assertion
in \eqref{cdp:eq:C.SKD.1}. At infinity the original smooth
fibre is $E\times E$, and the action is base
inversion on the first elliptic factor and
translation by the class of $q^{1/2}$ on the
second. It is free. Its quotient is a
bielliptic surface $B$, and the transverse
coordinate is squared, giving the full fibre
$2B$.

\item We identify the central quotients more precisely.
The canonical linearization on
$L|_R=\OO_{\PP^1}(-1)$ acts by $-1$ on its
tautological vector: the derivative at the
blown-up point is the scalar $-I$ on $\C^2$.
Thus the quotient of $\widetilde Q$ is
\begin{equation}
                 Q=(\C^2\setminus0)/
                         (-q^{1/2})^{\Z},              \hypertarget{tag-011Y}{}\label{cdp:eq:C.SKD.6}
\end{equation}
again a primary Hopf surface. On $\widetilde P$
the involution covers inversion on $E$.
The nowhere-zero canonical form of the primary
Kodaira surface is the wedge of the base
one-form and the fibre logarithmic differential.
The first changes sign; the second does not.
Consequently $g$ acts by $-1$ on that form.
Its square descends to a nowhere-zero section
of $K_P^{\otimes2}$, while no canonical form
descends. Thus $K_P$ has exact order two.
The finite \'etale cover $\widetilde P\to P$
and the computed $b_1(P)=1$ below identify $P$
as a secondary Kodaira surface.

\item Choose an $\iota$-invariant Hermitian
metric on $L$, by averaging. Let $M_5$ be its
unit circle bundle. The positive-real norm
coordinate gives
$\widetilde Y\simeq M_5\times S^1$. In that
description \eqref{cdp:eq:C.SKD.4} is
$(x,u)\mapsto(\iota_Mx,u+\tfrac12)$, where
the period of $u$ is one. Hence the quotient
is the mapping torus
\begin{equation}
          Y\simeq T(\iota_M)
            =(M_5\times\R)/
               ((x,u+1)\sim(\iota_Mx,u)).               \hypertarget{tag-011Z}{}\label{cdp:eq:C.SKD.7}
\end{equation}
For example the map $[x,u]\mapsto[x,2u]$
identifies the two quotient descriptions.

\item The integral circle-bundle Gysin sequence
over $V$ has Euler class $[R]$. The base has
free cohomology of ranks $(1,2,3,2,1)$.
Cup product with $[R]$ has ranks $1,0,1$
in degrees zero, one and two; the nonzero
images are primitive because $R^2=-1$.
Thus $M_5$ has free homology of ranks
\begin{equation}
                         (1,2,4,4,2,1).                \hypertarget{tag-0120}{}\label{cdp:eq:C.SKD.8}
\end{equation}
Its involution acts as $-I$ on $H_1$ and,
by oriented Poincar\'e duality, on $H_4$.
On $H_2$ it has an integral diagonal
decomposition with two $+1$ and two $-1$
summands. To specify that decomposition,
lift $E\times\{\infty\}$ and
$\{p_*\}\times\PP^1$, where $p_*\ne0$ is
a two-torsion point, using the normalized
canonical section of $L$. These are
invariant cycles, avoid $R$, and generate
the two base summands. The preimages of
the two elliptic-base loops at infinity
are circle tori; they generate the two
remaining summands and are sent to their
negatives. The homological Gysin sequence
identifies these four cycles with an
integral basis. Duality gives the same
two-positive, two-negative decomposition
on $H_3$. The actions on $H_0,H_5$ are
the identity. The Wang sequence of
\eqref{cdp:eq:C.SKD.7} therefore gives the rational
$Y$ row of \eqref{cdp:eq:C.SKD.2}.

\item Writing $T_2=(\Z/2)^2$,
the integral cohomology of $Y$ is
\begin{equation}
 \begin{array}{c|ccccccc}
 k&0&1&2&3&4&5&6\\ \hline
 H^k(Y,\Z)&
 \Z&\Z&\Z^2\oplus T_2&\Z^4\oplus T_2&
 \Z^2\oplus T_2&\Z\oplus T_2&\Z .
 \end{array}                                             \hypertarget{tag-0121}{}\label{cdp:eq:C.SKD.9}
\end{equation}
Indeed the Wang maps are zero or multiplication
by $-2$ on the stated integral summands.
Their kernel terms are free, so the short
exact sequences split as abelian groups;
the universal coefficient sequence gives
\eqref{cdp:eq:C.SKD.9}.

\item On $\widetilde P$, the unit circle bundle
$M_3\to E$ has Euler number one and free
homology ranks $(1,2,2,1)$. Its involution
acts by $-I$ in degrees one and two and
by $+1$ in degrees zero and three.
The quotient $P$ is its mapping torus.
The same Wang maps give
\begin{equation}
 \begin{array}{c|ccccc}
 k&0&1&2&3&4\\ \hline
 H^k(P,\Z)&
 \Z&\Z&T_2&\Z\oplus T_2&\Z .
 \end{array}                                             \hypertarget{tag-011O}{}\label{cdp:eq:C.SKD.10}
\end{equation}
In particular its rational ranks are those
of \eqref{cdp:eq:C.SKD.2}. The Hopf ranks follow from
\eqref{cdp:eq:C.SKD.6}, which is differentiably
$S^3\times S^1$.

\item For the support use the exact union
$S=P\cup Q$ with intersection $J$. The
radial deck loop of $J$ maps to the
mapping-torus generator on $P$ and the
Hopf generator on $Q$. Its angular loop
maps to zero in rational first homology:
upstairs it is the circle of an
Euler-number-one bundle, hence is already
zero in the abelianization. Therefore
the map
\begin{equation}
           H^1(P,\Q)\oplus H^1(Q,\Q)
                         \longrightarrow H^1(J,\Q)
                                                            \hypertarget{tag-011P}{}\label{cdp:eq:C.SKD.11}
\end{equation}
has rank one. Both surfaces have zero
rational $H^2$, so the next restriction
map has rank zero. Mayer--Vietoris gives
the support ranks $(1,1,1,3,2)$.
There is also an exact integral comparison at the only
potential torsion restriction. In the unit circle bundle
$M_3\to E$, the core fibre over zero bounds an
$\iota_M$-invariant oriented surface. To construct it,
take the normalized canonical section over the complement
of a small invariant disc at zero. Choose a local
coordinate and an equivariant frame with
$\sigma=z e$; the lift is
$(z,u)\mapsto(-z,-u)$ in the unit-circle coordinates.
The section on the disc boundary is
$(z,u)=(\varepsilon e^{i\theta},e^{i\theta})$.
Attach the annulus
\[
 (r,\theta)\longmapsto(re^{i\theta},e^{i\theta}),
                 \qquad 0\le r\le\varepsilon .
\]
Its inner boundary is the core circle. The involution
is $\theta\mapsto\theta+\pi$ on this annulus and
preserves its orientation; it also preserves the
orientation of the section over the punctured elliptic
curve. The joined surface, or its invariant rounded
chain, therefore has an oriented mapping torus whose
boundary is $J$. Thus
$H_2(J,\Z)\to H_2(P,\Z)$ is zero integrally,
including its possible torsion target. The angular
first-homology class is already zero in $M_3$;
the other generator maps to a primitive mapping-torus
class on $P$ and to a primitive Hopf class on $Q$.
The integral Mayer--Vietoris sequence now gives
\begin{equation}
 \begin{array}{c|ccccc}
 k&0&1&2&3&4\\ \hline
 H_k(S,\Z)&
 \Z&\Z\oplus T_2&\Z\oplus T_2&\Z^3&\Z^2\\
 H^k(S,\Z)&
 \Z&\Z&\Z\oplus T_2&\Z^3\oplus T_2&\Z^2 .
 \end{array}                                             \hypertarget{tag-011S}{}\label{cdp:eq:C.SKD.14}
\end{equation}
The free kernel terms split as abelian groups.
Any torsion summand in the radial image is absorbed
by eliminating the primitive Hopf generator, so
no choice of radial lift changes these groups.
The invariant cohomology of the free
quotient of $E\times E$ at infinity
gives the row for $B$. In both cases
the displayed Euler number is zero.
The quotient $B$ is the product of an
angular circle with the mapping torus of inversion
on the first elliptic curve. The Wang maps on
that curve are $0,-2I_2,0$. Its integral
cohomology is therefore
\begin{equation}
 H^k(B,\Z)=
 \bigl(\Z,\ \Z^2,\ \Z^2\oplus T_2,\
                  \Z^2\oplus T_2,\ \Z\bigr)_k .
                                                            \hypertarget{tag-011T}{}\label{cdp:eq:C.SKD.15}
\end{equation}
The integral groups of $J$ and $Q$ are free with
the ranks already displayed.

\item For coherent cohomology, the elliptic
line-bundle quotient $\widetilde Y\to V$
has direct images $\OO_V$ in degrees
zero and one and zero above one. The
degree-one frame is the constant
logarithmic deck cocycle, as in
\eqref{cdp:eq:C.PE.21}. Its involution action is
$+1$, while the action on
$H^1(V,\OO_V)=H^1(E,\OO_E)$ is $-1$.
The finite \'etale invariant-sheaf
comparison for \eqref{cdp:eq:C.SKD.5}, applied to
Leray's filtration, gives
$h^k(Y,\OO_Y)=(1,1,0,0)$. The identical
calculation for $\widetilde P\to E$
gives $h^k(P,\OO_P)=(1,1,0)$, consistent
with the canonical character.
Apply \cref{lem:II:hopf-curve-restriction} to the scalar Hopf quotient in \eqref{cdp:eq:C.SKD.6} and its elliptic axis $J$. It gives the Hopf row and the isomorphism $H^1(Q,\OO_Q)\to H^1(J,\OO_J)$ for this inclusion; the elliptic row is $h^\bullet(J,\OO_J)=(1,1)$. The coherent
normalization sequence gives
$h^k(S,\OO_S)=(1,1,0)$. The quotient
at infinity has the same base-negative,
fibre-positive action, giving
$h^k(B,\OO_B)=(1,1,0)$.

\item Upstairs let $\widetilde S$ be the
reduced central fibre of the original
map to $t$. Its coherent normalization
sequence gives $(1,2,1)$, the same
dimensions as a nearby product torus.
Grauert's theorem therefore makes all
$R^qf_*\OO_{\widetilde Y}$ locally free
near $t=0$, of ranks $1,2,1$, with
base change. The inverse image
of the full fibre $2S$ under the
\'etale map \eqref{cdp:eq:C.SKD.5} is $2\widetilde S$,
defined by $t^2=0$. Thus
\begin{equation}
 H^q(2S,\OO_{2S})
  \simeq
 \left((R^qf_*\OO_{\widetilde Y})/
                   t^2(R^qf_*\OO_{\widetilde Y})\right)^g .
                                                            \hypertarget{tag-011Q}{}\label{cdp:eq:C.SKD.12}
\end{equation}
Choose a basis of eigenvectors on the
central fibre of each direct-image
bundle. Lift each vector and average
with $g$ with its specified sign; the
resulting holomorphic sections remain
a frame near zero and satisfy
$g(e)=\pm e$. Since $g(t)=-t$, the
invariants in each summand
$\C[t]/(t^2)e$ have dimension one:
they are generated by $e$ for sign
$+1$ and by $te$ for sign $-1$.
Equation \eqref{cdp:eq:C.SKD.12} gives precisely
$(1,2,1)$. The same argument at
infinity proves the assertion for
$2B$. 

\item No rational curve lies in $P$:
its map from $\PP^1$ would lift to
the primary Kodaira cover, whose
elliptic base and elliptic fibres
force that map to be constant.
The same lifting argument to
$\C^2\setminus0$ treats $Q$.
Apply \cref{lem:II:nonuniruled-divisors} to these compact central divisors in the smooth quotient filling. The proved absence of rational curves gives its non-uniruled hypotheses, and the strict-transform correspondence proves their essentiality. At $T=0$, the finite
map
\begin{equation}
 \widetilde Y\longrightarrow
             Y\times_{\PP^1_T}\PP^1_t,\qquad
 y\longmapsto(\pi(y),f(y)),\quad T=t^2,
                                                            \hypertarget{tag-011R}{}\label{cdp:eq:C.SKD.13}
\end{equation}
is generically one-to-one and has
normal source. It identifies
$\widetilde Y$ with the normalization
of the reduced fibre product.
Consequently the primary Kodaira
component is recovered by this quadratic base change.

\end{enumerate}
\end{proof}






\subsection{Enriques and elliptic components with coherent Euler characteristic one}\label{subsec:II:enriques-family-occurrence}

\begin{proposition}[Enriques and elliptic components in a fixed punctured torus family]\hypertarget{tag-01D5}{}\label{cdp:st:C.2.UN.27}
\hypertarget{tag-01D3}{}\label{cdp:st:C.2.UN.25}
\StatementRecord{cdp:st:C.2.UN.27}{proved}{01D5}{\cref{lem:II:signed-root-fibration,lem:II:torsion-root-product-cover,lem:II:hopf-curve-restriction,lem:II:nonuniruled-divisors,lem:II:blowup-topology,cdp:st:C.2.UN.4,cdp:st:C.2.UN.23}; \cite[Theorem~2.1, p.~136, and Example~2.3, p.~137]{Yoshihara1995}}
\StatementRecord{cdp:st:C.2.UN.25}{proved}{01D3}{}
For every integer $r\ge2$ there is a proper holomorphic map
from a smooth compact complex threefold to $\PP^1$, with
connected fibres and two-torus general fibre, whose central
SNC support contains a smooth minimal projective elliptic
surface $M_r$. Its elliptic map to $\PP^1$ has exactly two
fibres of type $I_0^*$ and $r$ double smooth fibres
$2J_1,\ldots,2J_r$. For a full fibre $F$,
\begin{equation}
 \begin{gathered}
 K_{M_r}\simeq\OO_{M_r}\left(\sum_{j=1}^rJ_j-F\right),
 \qquad 2K_{M_r}\simeq\OO_{M_r}((r-2)F),\\
 P_m(M_r)=\max\{0,1-m+r\lfloor m/2\rfloor\}.
 \end{gathered}
\hypertarget{tag-00GV}{}\label{cdp:eq:C.EN.R1}
\end{equation}
Its cohomology is
\begin{equation}
 (b_0,b_1,b_2,b_3,b_4)=(1,0,10,0,1),\qquad
 (h^0(\OO),h^1(\OO),h^2(\OO))=(1,0,0).
\hypertarget{tag-00GZ}{}\label{cdp:eq:C.EN.R2}
\end{equation}
The surface $M_2$ is Enriques. For every $r\ge3$, $M_r$
has Kodaira dimension one and coherent Euler characteristic one.

Before six vertical modifications, the full central fibre is
\begin{equation}
 2M_r+2A+4\sum_{j=1}^rQ_j+\sum_{i=1}^4E_i,
\hypertarget{tag-00G9}{}\SsixEquationAliases{\label{cdp:eq:C.EN.1}}
\hypertarget{tag-00H0}{}\label{cdp:eq:C.EN.R3}
\end{equation}
where $A\simeq\PP^1\times\PP^1$, the $Q_j$ are primary
Hopf surfaces, and $E_i\simeq\mathbb F_2$. Its reduced
support has Euler number $12$. Six point-then-quartic
modifications, all disjoint from $M_r$, give Euler number zero
and preserve $M_r$ as the same smooth component.
The reduced supports have cohomology
\begin{equation}
 \begin{array}{c|c|c}
 &(b_0,b_1,b_2,b_3,b_4)&(h^0(\OO),h^1(\OO),h^2(\OO))\\ \hline
 S_r&(1,1,r+6,2r,r+6)&(1,1,0)\\
 S_{r,\mathrm{bal}}&(1,1,r+18,2r+36,r+18)&(1,1,0).
 \end{array}
\hypertarget{tag-00GI}{}\SsixEquationAliases{\label{cdp:eq:C.EN.18}}
\hypertarget{tag-00H1}{}\label{cdp:eq:C.EN.R4}
\end{equation}
At $r=2$, the identifications $M=M_2$, $S=S_2$,
$S_{\rm bal}=S_{2,\rm bal}$ and $(Q_a,Q_b)=(Q_1,Q_2)$
give the Enriques fibre and the rows $(1,1,8,4,8)$ and
$(1,1,20,40,20)$, respectively.
All these fillings have the same punctured torus family.
The components $M_r$ are essential under bimeromorphic
comparison of smooth fillings and are pairwise
non-bimeromorphic, with canonical growth slope $r/2-1$.
The corresponding fillings are pairwise non-bimeromorphic
over that punctured family.
\end{proposition}
\begin{proof}
\begin{enumerate}[label=\arabic*.,leftmargin=*]
\item Take the Weierstrass elliptic surface over $\PP^1_T$
\begin{equation}
 y^2=x^3+(T^4-3)x+(T^6+T+2).
\hypertarget{tag-00GJ}{}\label{cdp:eq:C.EN.2}
\end{equation}
The coefficients are sections of $\OO(4)$ and $\OO(6)$.
Its discriminant, up to a nonzero constant factor, is
\begin{equation}
 \begin{split}
 \Delta(T)={}&31T^{12}-36T^8+54T^7+108T^6\\
             &+108T^4+27T^2+108T.
 \end{split}
\hypertarget{tag-00GK}{}\label{cdp:eq:C.EN.3}
\end{equation}
The Euclidean algorithm applied to this polynomial and its
derivative gives $\gcd(\Delta,\Delta')=1$.
Moreover $\Delta'(0)=108$ and the leading coefficient is $31$.
Thus every singular fibre is nodal of type $I_1$, the fibre
at zero is nodal, infinity is smooth, and the total
Weierstrass surface is smooth.

\item Put $\xi=x-1$ and $U=1+T^3(1+\xi)+T^5$. Near the node,
\begin{equation}
 X=\frac{y-\xi\sqrt{3+\xi}}{U},\qquad
 Y=y+\xi\sqrt{3+\xi},\qquad XY=T.
\hypertarget{tag-00GL}{}\label{cdp:eq:C.EN.4}
\end{equation}
The coordinate Jacobian at $\xi=y=T=0$ is nonzero.
Make the quadratic base change $T=t^2$ and resolve $XY=t^2$.
One resolution chart and its deck involution are
\begin{equation}
 (X,Y,t)=(u,uv^2,uv),\qquad h(u,v)=(u,-v).
\hypertarget{tag-00GM}{}\label{cdp:eq:C.EN.5}
\end{equation}
The other interchanges $X,Y$ and replaces $v$ by $v^{-1}$.
The resulting compact smooth surface $V_0$ has central fibre
$R_0+R_1$ of type $I_2$. The strict nodal normalization
$R_0\simeq\PP^1$ is fixed pointwise. The exceptional
$R_1\simeq\PP^1$ has coordinate $z=v$, with $h(z)=-z$,
and meets $R_0$ at $0,\infty$.

\item Blow up the $2r$ points $z=\pm1,\ldots,\pm r$ on $R_1$,
and denote the resulting surface by $V$. The involution lifts
and exchanges their exceptional curves $D_j,D_{-j}$. Set
\begin{equation}
 D_r=\sum_{j=1}^r(D_j-D_{-j}),\qquad L_r=\OO_V(D_r).
\hypertarget{tag-00GN}{}\SsixEquationAliases{\label{cdp:eq:C.EN.6}}
\hypertarget{tag-00H2}{}\label{cdp:eq:C.EN.R5}
\end{equation}
The canonical section $\sigma$ gives the order-two,
section-preserving dualization $h^*L_r\simeq L_r^{-1}$.
Apply \cref{lem:II:signed-root-fibration} with this same $V$,
$D_+=\sum D_j$, $D_-=\sum D_{-j}$, $m=2$, and $0<q<1$.
Its positive and negative equations are $\sigma=a^2\ell$ and
$\sigma^{-1}=(a')^2\ell'$, with
$(\ell',a')=(\ell^{-1},a^{-1})$. Its deck generator is
\begin{equation}
             (\ell,a)\longmapsto(q^2\ell,q^{-1}a).
\hypertarget{tag-00GO}{}\label{cdp:eq:C.EN.7}
\end{equation}
The identity on these line-bundle coordinates identifies the
lemma's smooth proper quotient with $\widetilde Y\to V$.
Its product coordinate is
$\ell=\sigma c^2$, $a=c^{-1}$, with $[c]\in E=\C^*/q^\Z$.
Dualization sends the meromorphic pair to $(\ell^{-1},a^{-1})$
over $h$; in a positive chart mapping to a negative chart its
numerical coordinates are unchanged. It conjugates the deck
generator to its inverse and squares to one. On the product,
\begin{equation}
                  (v,[c])\longmapsto(hv,[c^{-1}]).
\hypertarget{tag-00GP}{}\label{cdp:eq:C.EN.8}
\end{equation}
All blowups and signed roots lie over $t=0$. The blowdowns
$V\to V_0$ restrict to isomorphisms over the punctured base.
Through these maps, the product and involution identify the
punctured family for every $r$ by the identity on $(v,[c],t)$.

\item The restriction of $L_r$ to the strict $R_1$ is trivial.
Choose its frame so that
\begin{equation}
 \phi_r(z)=\prod_{j=1}^r\frac{z-j}{z+j},\qquad
 C_r:\ Y^2=\prod_{j=1}^r(z^2-j^2),\qquad
 y=\frac{Y}{\prod_{j=1}^r(z+j)},\qquad y^2=\phi_r(z).
\hypertarget{tag-00GQ}{}\hypertarget{tag-00GA}{}
\SsixEquationAliases{\label{cdp:eq:C.EN.9}\label{cdp:eq:C.EN.10}}
\hypertarget{tag-00H3}{}\label{cdp:eq:C.EN.R6}
\end{equation}
Apply \cref{lem:II:torsion-root-product-cover} to
$B=R_1$, $D_+=\sum[j]$, $D_-=\sum[-j]$, $m=2$,
$L=N=\OO_{R_1}$, and the section $\phi_r$.
Its normalized cyclic curve is exactly $C_r$, of genus $r-1$;
the identity frame allows $C'=C_r$. If $B_r$ denotes the
component over $R_1$, the covering map is
\begin{equation}
 C_r\times E\longrightarrow B_r,\qquad
 (z,y,[u])\longmapsto[\ell=u^2,a=y/u].
\hypertarget{tag-00GB}{}\label{cdp:eq:C.EN.11}
\end{equation}
The two maps and Jacobians of \eqref{cdp:eq:C.PE.13}, with
$m=2$ and this frame, prove that it is finite unramified
at every root value. In particular $B_r$ is projective.
The deck involution and the lifted dualization are
\begin{equation}
\SsixEquationAliases{\label{cdp:eq:C.EN.12}\label{cdp:eq:C.EN.13}}
 \begin{split}
 \eta(z,Y,[u])&=(z,-Y,[-u]),\\
 k(z,Y,[u])&=(-z,(-1)^rY,[u^{-1}]).
 \end{split}
\hypertarget{tag-00GC}{}\hypertarget{tag-00GD}{}
\hypertarget{tag-00H4}{}\label{cdp:eq:C.EN.R7}
\end{equation}
The translation $[u]\mapsto[-u]$ is nonzero two-torsion,
so $\eta$ is free. The two maps commute and generate a
Klein four group. Since $\phi_r(-z)=\phi_r(z)^{-1}$,
the second map is also $(z,y,[u])\mapsto(-z,y^{-1},[u^{-1}])$.
At $r=2$, the identity on $(z,Y,[u])$ identifies these
curves, covers and involutions with the Enriques construction;
in particular $C_2$ is elliptic and $B_2$ is bielliptic.

\item At infinity, $Y/z^r=\pm1$, and $k_C$ fixes both points
while $jk_C$ exchanges them, where $j(z,Y)=(z,-Y)$.
At zero $Y\ne0$; the roles of $k_C$ and $jk_C$ are
respectively fixed/exchanged for even $r$ and exchanged/fixed
for odd $r$. Thus
\begin{equation}
 \begin{array}{c|cc}
 &\#\operatorname{Fix}(k_C)&\#\operatorname{Fix}(jk_C)\\ \hline
 r\text{ even}&4&0\\ r\text{ odd}&2&2.
 \end{array}
\hypertarget{tag-00H5}{}\label{cdp:eq:C.EN.R8}
\end{equation}
Every fixed-point action on the smooth curve has derivative $-1$.
Both $[u]\mapsto[u^{-1}]$ and $[u]\mapsto[-u^{-1}]$ on $E$
have four fixed points, with derivative $-1$.
The elements $k,\eta k$ therefore have sixteen fixed points
in total, all of type $(-1,-1)$. Their fixed sets are disjoint
because $\eta$ is free, and the remaining free involution
pairs them. Consequently
\begin{equation}
                 X_r=(C_r\times E)/\langle\eta,k\rangle
\hypertarget{tag-00H6}{}\label{cdp:eq:C.EN.R9}
\end{equation}
has exactly eight $A_1$ singularities and no branch curve.
Let $M_r$ be its minimal $A_1$ resolution.

\item The map $w=z^2$ induces $p:M_r\to\PP^1_w$.
Over a general $w$ the four points of $C_r$ form one free
orbit under the group, so the fibre is $E$.
At $w=j^2$, the double-root chart and the translation in
$\eta$ give $2J_j$. At $0,\infty$, the involution on the
elliptic fibre has four fixed points. Its resolution has
full fibre $2D+\sum_{i=1}^4U_i$, of type $I_0^*$.
These are all special fibres; in particular $p$ has connected fibres.

\item On $C_r\times E$ the meromorphic form
\begin{equation}
                         \Omega=dz\wedge\frac{du}{u}
\hypertarget{tag-00GW}{}\label{cdp:eq:C.EN.R10}
\end{equation}
is invariant under both generators. The differential $dz$
has a simple zero at every one of the $2r$ branch points,
a double pole at each point over infinity, and no other
zero or pole. The quotient is free in codimension one,
so the root charts give one zero along each $J_j$.
At infinity the base coordinate $w^{-1}$ pulls back with
order two at both points of $C_r$.

At a quotient point choose $(s,v)$ with action $(-s,-v)$.
In the resolution chart $s^2=r_0$, $sv=r_0\delta$,
the invariant volume form is
$ds\wedge dv=\tfrac12\,dr_0\wedge d\delta$.
Thus the double pole $s^{-2}$ becomes a simple pole on
the exceptional curve; in the other chart it has orders
one and two on the exceptional and strict central curves.
These are the full $I_0^*$ multiplicities. At zero the
form and its crepant extension are nowhere zero. Hence
\begin{equation}
                  \operatorname{div}(\Omega)=\sum_{j=1}^rJ_j-F_\infty.
\hypertarget{tag-00GX}{}\label{cdp:eq:C.EN.R11}
\end{equation}
Since $2J_j\sim F$, this gives both canonical-line identities
in \eqref{cdp:eq:C.EN.R1}.

\item The ratio of a holomorphic $m$-canonical section to
$\Omega^{\otimes m}$ is constant on a general elliptic fibre
and descends meromorphically to its connected-fibre base.
The allowed orders at $2J_j$ and at the full Cartier fibre
$F_\infty$ give
\begin{equation}
 p_*\OO_{M_r}(mK_{M_r})\simeq
 \OO_{\PP^1}\left(\sum_{j=1}^r\lfloor m/2\rfloor[j^2]-m[\infty]\right).
\hypertarget{tag-00GY}{}\label{cdp:eq:C.EN.R12}
\end{equation}
Taking sections proves the plurigenus formula. Its slope is
$r/2-1$; for $r=2$ the even plurigenera are one and the
odd ones zero. The line $2K_{M_r}=(r-2)F$ is nef for all
$r\ge2$. A $(-1)$-curve would have canonical degree $-1$
by adjunction, so $M_r$ is minimal.

\item Apply \cite[Theorem~2, p.~241]{ColtoiuGasitoiJoita2012}
to the finite quotient $C_r\times E\to X_r$. Its source
is projective and its target is normal, so $X_r$ is projective.
The blowup resolving its eight $A_1$ points is projective;
hence $M_r$ is projective. The involution $j$ is
$-I$ on $H^1(C_r,\OO)$: its dual canonical-form basis is
$z^a dz/Y$, $0\le a\le r-2$. Thus $\eta$ removes this
summand and $k$ removes the elliptic summand, giving
$q(M_r)=0$. Resolution of the $A_1$ points preserves
coherent cohomology. Indeed the successive conormal powers
of each exceptional $(-2)$-curve have vanishing $H^1$;
formal functions gives $R^1\pi_*\OO=0$.
Equation \eqref{cdp:eq:C.EN.R12} gives $p_g(M_r)=0$.
The finite-quotient Euler formula gives $(0+0+16)/4=4$;
the eight resolutions add eight. Thus $e(M_r)=12$,
and Hodge symmetry and Poincar\'e duality give
\eqref{cdp:eq:C.EN.R2}.

For $r=2$, both $j$ and $k_C$ negate $dz/Y$, whereas
$jk_C$ is free: over zero it would require $Y=0$, and
at infinity it exchanges the two points. It is a nonzero
two-torsion translation of $C_2$. Both $\eta$ and $\eta k$
are therefore free on $C_2\times E$. Choose a fixed point
of $k_C$ as the origin on $C_2$ and $[1]$ on $E$.
The resulting product group identifies $k$ with inversion
on that same abelian surface. Thus \cref{cdp:st:C.2.UN.23},
applied to this product with these origins, gives that the resolution of
$(C_2\times E)/\langle k\rangle$ is K3.
The induced $\eta$ permutes its sixteen exceptional curves
in pairs and is free off them and on them. Its quotient is
$M_2$, proving that this is an Enriques surface.
This product quotient is one of the Enriques cases in
\cite[Theorem~2.1, p.~136, and Example~2.3, p.~137]{Yoshihara1995}.

\item We realize this resolution inside the compact threefold.
Near $R_0$, the root divisors are absent and
\eqref{cdp:eq:C.EN.8} has four fixed copies of $R_0$,
indexed by $E[2]$. Near either node of $R_0+R_1$, choose
an elliptic coordinate $w$ on which inversion is $w\mapsto-w$.
The action and the base function are
\begin{equation}
              (u,v,w)\longmapsto(u,-v,-w),\qquad T=u^2v^2.
\hypertarget{tag-00GE}{}\label{cdp:eq:C.EN.14}
\end{equation}
Its invariant coordinates satisfy
$A_0=v^2$, $B_0=vw$, $C_0=w^2$, $A_0C_0=B_0^2$,
and $T=u^2A_0$. The resolution charts are
\begin{equation}
 \begin{array}{lll}
 A_0=r_0,&B_0=r_0s,&C_0=r_0s^2,\quad T=u^2r_0,\\
 A_0=r'_0s'^2,&B_0=r'_0s',&C_0=r'_0,\quad T=u^2r'_0s'^2.
 \end{array}
\hypertarget{tag-00GF}{}\label{cdp:eq:C.EN.15}
\end{equation}
They glue by $s'=s^{-1}$, $r'_0=r_0s^2$ and resolve all
four fixed curves. Their restriction to $u=0$ is exactly
the eight $A_1$ resolutions defining $M_r$.
For odd $r$, the fixed lifts over zero are those of
$\eta k$; after the free $\eta$ quotient the local
involution and these charts are unchanged.

The component over $R_0$ is
$A=R_0\times(E/\{\pm1\})\simeq\PP^1\times\PP^1$.
The product trivialization holds on an entire neighbourhood
of $R_0$, and $R_0^2=-2$. Thus the normal bundle of each
fixed curve is $\OO_{R_0}(-2)\oplus\OO_{R_0}$, and its
exceptional projective bundle is $E_i\simeq\mathbb F_2$.
The $2r$ root Hopf components are exchanged in pairs,
giving $Q_1,\ldots,Q_r$. The function $T=t^2$ and
the monomials in \eqref{cdp:eq:C.EN.15} give all
multiplicities in \eqref{cdp:eq:C.EN.R3}.

\item The complete incidence table is
\begin{equation}
 \begin{array}{c|c|c}
 \text{components}&\text{intersection}&\text{number}\\ \hline
 M_r,A&D_0,D_\infty\simeq\PP^1&2\\
 M_r,Q_j&J_j\text{ elliptic}&r\\
 A,E_i&Z_i\simeq\PP^1&4\\
 M_r,E_i&U_{i0},U_{i\infty}\simeq\PP^1&8\\
 M_r,A,E_i&D_\epsilon\cap Z_i\cap U_{i\epsilon}&8.
 \end{array}
\hypertarget{tag-00GG}{}\SsixEquationAliases{\label{cdp:eq:C.EN.16}}
\hypertarget{tag-00H7}{}\label{cdp:eq:C.EN.S1}
\end{equation}
There are no other intersections. On $A$, the curves
$D_0,D_\infty$ are fibres of one projection and the $Z_i$
are sections with a common class. On $E_i$, $Z_i$ is the
negative section and $U_{i0},U_{i\infty}$ are ruling fibres.
The second chart of \eqref{cdp:eq:C.EN.15} has reduced
support $u r'_0s'=0$, proving the SNC assertion, including
every triple point. Inclusion--exclusion gives
\begin{equation}
 e(S_r)=12+4+4\cdot4-(4+4\cdot2+4\cdot4)+8=12.
\hypertarget{tag-00GH}{}\label{cdp:eq:C.EN.17}
\end{equation}

\item At infinity, the base involution fixes a smooth elliptic
fibre. Its four ambient fixed curves are elliptic, and the
same transverse $A_1$ resolutions have elliptic ruled
exceptional surfaces and elliptic intersections, all with
Euler number zero. At every other nodal base value the
product chart is a nodal curve times $E$, also of Euler
number zero. The quotients and resolutions are proper and
compact. For a general $T$, the involution exchanges the
two $t$-sheets, whose quotient is one product of elliptic
curves. The Stein factor has generic degree one over the
normal base $\PP^1$ and is that base, proving connectedness.

\item On $M_r$ the full fibre divisors are
\begin{equation}
 \begin{gathered}
 F_0=2D_0+\sum_{i=1}^4U_{i0},\qquad
 F_\infty=2D_\infty+\sum_{i=1}^4U_{i\infty},\\
 [F_0]=[F_\infty]=2[J_j]\quad(1\le j\le r)
       \quad\text{in }H_2(M_r,\Q).
 \end{gathered}
\hypertarget{tag-00H8}{}\SsixEquationAliases{\label{cdp:eq:C.EN.S2}}
\hypertarget{tag-00GR}{}\label{cdp:eq:C.EN.G1}
\end{equation}
The root chart gives $2J_j$, and the resolutions give
the displayed two $I_0^*$ fibres. Their equality follows
by pullback of point divisors along this same $p$.
The ten rational curves have intersection matrix equal to
the sum of two affine $D_4$ matrices. Each has rank four
and kernel $(2,1,1,1,1)$. Every homology relation lies in
their combined kernel. The common fibre class is nonzero,
since it meets an ample divisor positively. Thus the sole
homology relation is $F_0-F_\infty=0$, and these curves
span dimension nine.

\item Let $S_r^{[p]}$ be the disjoint union of the
$(p+1)$-fold intersections in \eqref{cdp:eq:C.EN.S1}.
The alternating restriction maps give
\begin{equation}
 0\longrightarrow\Q_{S_r}\longrightarrow\Q_{S_r^{[0]}}
 \longrightarrow\Q_{S_r^{[1]}}\longrightarrow\Q_{S_r^{[2]}}
 \longrightarrow0.
\hypertarget{tag-00H9}{}\label{cdp:eq:C.EN.S3}
\end{equation}
At each point this is the augmented cochain complex of
the simplex of its branches, so it is exact. Its spectral
sequence has $E_1^{p,q}=H^q(S_r^{[p]},\Q)$.
The degree-zero row has dimensions $r+6,r+14,8$.
Connectedness gives first rank $r+5$; the second rank is
eight, since each triangle has its own edge $U_{i\epsilon}$.
Its row cohomology is $(1,1,0)$.

Only the $r$ Hopf components contribute to component $H^1$,
and only the $J_j$ to double-curve $H^1$. Their map is
$\Q^r\to\Q^{2r}$ of rank $r$: for the scalar presentation
$(\C^2\setminus0)/q^{2\Z}$ the axis inclusion preserves
the radial deck generator and kills the angular generator.
Each image lies in its own elliptic summand.

\item In degree two, the restriction map is
$\Q^{20}\to\Q^{r+14}$. For its dual curve-class map,
orient each edge from the first component to the second
in \eqref{cdp:eq:C.EN.S1}. Give coefficients
$x_0,x_\infty$ to $D_0,D_\infty$, $y_j$ to $J_j$,
$z_i$ to $Z_i$, and $u_{i\epsilon}$ to $U_{i\epsilon}$.
The $E_i$ equations give $z_i=0$ and
$u_{i\infty}=-u_{i0}$; the $A$ equation gives
$x_\infty=-x_0$. The $M_r$ equation and the sole
rational-curve relation then give exactly
\begin{equation}
 \sum_{j=1}^r y_j=0,\qquad u_{i0}=x_0/2,\qquad
 u_{i\infty}=-x_0/2,\qquad x_\infty=-x_0,\qquad z_i=0.
\hypertarget{tag-00HA}{}\SsixEquationAliases{\label{cdp:eq:C.EN.S4}}
\hypertarget{tag-00GS}{}\label{cdp:eq:C.EN.G2}
\end{equation}
Indeed, modulo $F_0-F_\infty$, a nonzero multiple of
$F_0$ cannot be a linear combination of the differences
$D_0-D_\infty$ and $U_{i0}-U_{i\infty}$: its sums of
the two corresponding coefficient sets would be nonzero.
This proves independence of the half-fibre class and the
four-dimensional difference span. Thus the kernel has
$(r-1)+1=r$ parameters, and the rank is fourteen.

The components contribute dimensions $r$ in degree three
and $r+6$ in degree four; intersections contribute neither.
All higher differentials have zero target: column two
vanishes in positive degree, and $E_2^{2,0}=0$.
Consequently
\begin{equation}
 (b_i(S_r))_{i=0}^4=(1,1,r+6,2r,r+6),\qquad e(S_r)=12.
\hypertarget{tag-00HB}{}\SsixEquationAliases{\label{cdp:eq:C.EN.S5}}
\hypertarget{tag-00GT}{}\label{cdp:eq:C.EN.G3}
\end{equation}

\item The alternating structure-sheaf complex is exact in
the same SNC coordinates. Every component has zero
$H^2(\OO)$; only the Hopf components have nonzero
$H^1(\OO)$. Apply \cref{lem:II:hopf-curve-restriction}
to these same scalar Hopf surfaces and their axis inclusions
$J_j\hookrightarrow Q_j$. The restriction in degree one
is an isomorphism for each $j$, with the additive deck
class as its frame. The degree-zero incidence maps are
unchanged. It follows that
\begin{equation}
                   (h^i(S_r,\OO_{S_r}))_{i=0}^2=(1,1,0).
\hypertarget{tag-00HC}{}\label{cdp:eq:C.EN.S6}
\end{equation}

\item Perform six point-then-quartic operations of
\cref{cdp:st:C.2.UN.4} over zero. Each begins at a smooth
point of an $E_i$ outside its intersections, and its
exceptional quartic and later centre remain disjoint from
$M_r$. Each preserves the smooth total space and SNC
support, changes its Euler number by $-2$, and is the
identity on the punctured family.
Apply \cref{lem:II:blowup-topology} at each point or curve
centre to the full inverse-image support. Its relative-pair
and restricted exceptional Gysin maps are precisely the
maps in \eqref{cdp:eq:C.KM.S2}; a block adds $(2,6,2)$
in degrees two, three and four. The centres have generic
reduced-support multiplicity one, and the ordinary blowup
charts give $\sigma^*S_r=S'_r$ as Cartier divisors,
including the four crossing points at the quartic step.
Thus the same Cartier-sequence and projection-formula
comparison in \eqref{cdp:eq:C.KM.S3}, with these $\sigma,S_r,S'_r$,
gives $R\sigma_*\OO_{S'_r}=\OO_{S_r}$. Hence
\begin{equation}
 \begin{gathered}
 (b_i(S_{r,\rm bal}))_{i=0}^4=(1,1,r+18,2r+36,r+18),\\
 (h^i(S_{r,\rm bal},\OO_{S_{r,\rm bal}}))_{i=0}^2=(1,1,0),
 \qquad e(S_{r,\rm bal})=0.
 \end{gathered}
\hypertarget{tag-00HD}{}\SsixEquationAliases{\label{cdp:eq:C.EN.S7}}
\hypertarget{tag-00GU}{}\label{cdp:eq:C.EN.G4}
\end{equation}
For $r=2$, the identity of the labelled incidence components
and all edge inclusions specializes \eqref{cdp:eq:C.EN.G1}
through \eqref{cdp:eq:C.EN.G4} to the Enriques support.

\item The positive pluricanonical sections in
\eqref{cdp:eq:C.EN.R1} exclude a covering family of rational
curves: a dominant ruled parametrization would pull each
such section back to zero. Apply
\cref{lem:II:nonuniruled-divisors} to the compact divisors
$M_r$ in these smooth proper fillings. Its strict-transform
correspondence carries $M_r$ birationally onto a divisor
in every smooth bimeromorphic filling, proving essentiality.
Distinct canonical growth slopes make these divisors
pairwise non-bimeromorphic. For $r\ge3$, $M_r$ is the
unique component of Kodaira dimension one; all others
are rational, Hopf or ruled. At $r=2$, its Kodaira
dimension is zero and the other components have dimension
$-\infty$. The same divisor correspondence therefore
distinguishes the fillings over the fixed punctured family.
\end{enumerate}
\end{proof}

\subsection{Elliptic components in the Kummer family}\label{subsec:II:kummer-elliptic-occurrence}

\begin{proposition}[The fixed Kummer punctured family has zero-Euler fillings with arbitrarily large canonical growth on a projective elliptic component]\hypertarget{tag-01D9}{}\label{cdp:st:C.2.UN.30}
\StatementRecord{cdp:st:C.2.UN.30}{proved}{01D9}{\cref{lem:II:signed-root-fibration,lem:II:torsion-root-product-cover,lem:II:hopf-curve-restriction,lem:II:nonuniruled-divisors}, \SeriesRef{cdp:st:C.2.UN.23}, \SeriesRef{cdp:st:C.2.UN.13}, \SeriesRef{cdp:st:C.2.UN.25}, \SeriesRef{cdp:eq:C.ET.11}, \SeriesRef{cdp:st:C.2.UN.14}}

Fix the square elliptic curve
$E_0=\C/(\Z+i\Z)$ and $E=\C^*/q^{\Z}$,
where $q=e^{-2\pi}$. For every $r\ge1$ there is
a smooth compact threefold over $\PP^1$,
with the same punctured product-torus family
as the diagonal involution quotient
\begin{equation}
 (E_0\times E\times\PP^1_t)/
      \langle(p,[c],t)\mapsto(-p,[c^{-1}],-t)\rangle,
                                                               \hypertarget{tag-00J5}{}\label{cdp:eq:C.ET.1}
\end{equation}
whose two special reduced supports have Euler
number zero. The support over zero contains an
unchanged smooth minimal projective properly
elliptic surface $M_r$ with four $I_0^*$ fibres
and $r$ double smooth elliptic fibres $2J_j$.
For its connected elliptic map $p_r:M_r\to\PP^1$
and a full fibre $F$,
\begin{equation}
 \begin{aligned}
 K_{M_r}&\simeq\OO_{M_r}\left(\sum_{j=1}^rJ_j\right),
 \qquad 2K_{M_r}\simeq\OO_{M_r}(rF),\\
 (p_r)_*\OO(mK_{M_r})
   &\simeq\OO_{\PP^1}\left(
                \sum_{j=1}^r\lfloor m/2\rfloor[b_j]\right),
 \\
 P_m(M_r)&=1+r\lfloor m/2\rfloor .
 \end{aligned}                                             \hypertarget{tag-00J9}{}\label{cdp:eq:C.ET.2}
\end{equation}
Here $b_j$ is the value of $2J_j$. In particular
$\chi(\OO_{M_r})=2$, $p_g(M_r)=1$, $q(M_r)=0$,
and
\begin{equation}
            (b_0,b_1,b_2,b_3,b_4)(M_r)
                         =(1,0,22,0,1).              \hypertarget{tag-00JA}{}\label{cdp:eq:C.ET.3}
\end{equation}
These essential components, and the associated
fillings over the fixed punctured family, are
pairwise non-bimeromorphic.

Before balancing, the full fibre over zero is
\begin{equation}
                  2M_r+4\sum_{j=1}^rQ_j
                           +\sum_{\nu=1}^{16}P_\nu,
                                                               \hypertarget{tag-00JB}{}\label{cdp:eq:C.ET.4}
\end{equation}
where the $Q_j$ are primary Hopf surfaces and
$P_\nu\simeq\PP^2$. It has Euler number $40$.
Its reduced support and its balanced version
have respectively
\begin{equation}
 \begin{array}{c|c|c}
 &(b_0,b_1,b_2,b_3,b_4)&
                 (h^0(\OO),h^1(\OO),h^2(\OO))\\ \hline
 S_r&(1,0,21+r,2r-1,17+r)&(1,0,1)\\
 S_{r,\mathrm{bal}}&(1,0,61+r,2r+119,57+r)&(1,0,1).
 \end{array}                                             \hypertarget{tag-00JC}{}\label{cdp:eq:C.ET.5}
\end{equation}
\end{proposition}


\begin{proof}

\begin{enumerate}[label=\arabic*.,leftmargin=*]
\item Choose an odd prime $N>2r$ and the torsion
points $p_j=j/N$ of $E_0$, for $1\le j\le r$.
Their pairs $\{p_j,-p_j\}$ are disjoint and
avoid the two-torsion points. Blow up
$E_0\times\PP^1_t$ at $(p_j,0),(-p_j,0)$,
and call the exceptional curves $R_j,R_{-j}$.
The involution $(p,t)\mapsto(-p,-t)$ lifts
to the resulting surface $V$, exchanging
the paired curves. Put
\begin{equation}
       D=\sum_{j=1}^r(R_j-R_{-j}),\qquad L=\OO_V(D).
                                                               \hypertarget{tag-00JD}{}\label{cdp:eq:C.ET.6}
\end{equation}
The canonical meromorphic section $\sigma$
gives $h^*L=L^{-1}$ with its order-two
section-preserving dualization.

\item Apply \cref{lem:II:signed-root-fibration} to the same
$V,L,\sigma$ with $D_+=\sum R_j$, $D_-=\sum R_{-j}$,
$m=2$ and $q=e^{-2\pi}$. The identity in the positive
coordinates $(\ell,a)$ and the reciprocal change
$(\ell',a')=(\ell^{-1},a^{-1})$ identify its two charts
with the restrictions along these signed curves. Its deck
map is $(\ell,a)\mapsto(q^2\ell,q^{-1}a)$, and its
proper quotient is the smooth compact $\widetilde Y\to V$.
Dualization over $h$ sends the meromorphic pair to
$(\ell^{-1},a^{-1})$ and conjugates that deck map to its
inverse. Thus it defines the involution on this same quotient.

\item Away from the root curves the formula
$\ell=\sigma c^2,\ a=c^{-1}$ identifies the
family with $V\times E$. Under that identification the involution is
$(v,[c])\mapsto(hv,[c^{-1}])$.
All modifications lie over $t=0$.
Consequently the quotient family over
$T=t^2$, on the complement of zero and
infinity, is exactly the punctured family
in \eqref{cdp:eq:C.ET.1}, independently of $r$.
The quotient and every resolution are proper over
$\PP^1$; the Stein factor has generic degree one
and is therefore the identity finite map to the
normal base. All fibres are connected.

\item Let $B_r$ be the component above the strict central
elliptic curve $E_0$. It has $2r$ smooth double fibres at
$p_j,-p_j$, and its restricted line is
\begin{equation}
 L|_{E_0}\simeq\OO_{E_0}\left(\sum_{j=1}^r[p_j]-[-p_j]\right).
\hypertarget{tag-00JE}{}\label{cdp:eq:C.ET.7}
\end{equation}
The group sum of this divisor is $2\sum p_j$, a torsion
point, so this line is torsion. Apply
\cref{lem:II:torsion-root-product-cover} with $B=E_0$,
$D_+=\sum[p_j]$, $D_-=\sum[-p_j]$, $m=2$, the same
$q$, and this restricted section. Denote its torsion
square root by $\mathcal N$ (the line $N$ in that lemma);
its connected normalized curve $C$ has
$g(C)=r+1$. Its finite unramified cover $C'\to C$ has
the frame $n$ of $\mathcal N$ and coordinate $y=y_0n$, giving
\begin{equation}
 C'\times E\longrightarrow B_r,\qquad
 (z,[u])\longmapsto[\ell=u^2n^2,a=y_0/u].
\hypertarget{tag-00JF}{}\label{cdp:eq:C.ET.8}
\end{equation}
This is the same map \eqref{cdp:eq:C.PE.12} on the same
root charts. Its two Jacobians are $2$ and $-2u^{-2}$,
the $m=2$ instances of \eqref{cdp:eq:C.PE.13}.
The lemma proves that this map is finite unramified
everywhere and that $B_r$ is projective.

\item The involution on $B_r$ covers inversion
on $E_0$. At each of the four two-torsion
base points, the root curves are absent,
and the product formula gives
inversion on the elliptic fibre.
There are exactly sixteen fixed points
on $B_r$, each with derivative $(-1,-1)$.
There is no fixed curve. Its quotient has sixteen $A_1$ singularities. Apply \cite[Theorem~2, p.~241]{ColtoiuGasitoiJoita2012} to the finite map from the projective $B_r$ to this normal quotient; every fibre has dimension zero. The quotient is projective, and its blowup resolving the sixteen $A_1$ points is the projective surface $M_r$.
The induced base is
$E_0/\{\pm1\}\simeq\PP^1$.
The four fixed base values give $I_0^*$
fibres by the same explicit resolution
as \eqref{cdp:eq:C.EN.S2}. The paired root values
give the $r$ smooth double fibres
$2J_j$. Away from those values the
quotient is a smooth elliptic fibration,
so this is its complete fibre list.

\item Let $\omega_0$ be a nonzero holomorphic
one-form on $E_0$. On $B_r$ the
meromorphic expression
\begin{equation}
                        \Omega=\omega_0\wedge d\log\ell
                                                               \hypertarget{tag-00JG}{}\label{cdp:eq:C.ET.9}
\end{equation}
is global. A change of frame of $L$
adds a base logarithmic differential
to $d\log\ell$; its wedge with
$\omega_0$ is zero. Multiplication
by the constant $q^2$ also leaves
the expression unchanged. In a
positive root chart $z=a^2b$ it is
a unit times $2a\,da\wedge db$;
the negative dual chart has the
opposite sign and the same simple
zero. Thus its divisor on $B_r$
is precisely the sum of the
$2r$ reduced multiple fibres.
The involution changes both
$\omega_0$ and $d\log\ell$ by a
minus sign, up to the harmless
base differential already described.
Hence $\Omega$ is invariant.

\item At its sixteen fixed points the form
is regular and nowhere zero. The
height-one $A_1$ charts of \ref{cdp:st:C.2.UN.23}
extend it without a zero or pole.
Paired root fibres are exchanged
freely and become $J_j$. The
descended form consequently has
the exact divisor
\begin{equation}
                      \operatorname{div}(\Omega)
                                      =\sum_{j=1}^rJ_j .
                                                               \hypertarget{tag-00J6}{}\label{cdp:eq:C.ET.10}
\end{equation}
This proves the first two line-bundle
identities in \eqref{cdp:eq:C.ET.2}.
The ratio of any holomorphic
$m$-canonical section to
$\Omega^{\otimes m}$ is constant
on a connected general elliptic
fibre. It descends meromorphically
to the base. Its allowed pole
order at $b_j$ is exactly
$\lfloor m/2\rfloor$, because
pullback of a base parameter
has divisor $2J_j$; there are
no other allowed poles.
This proves the direct-image equality in \eqref{cdp:eq:C.ET.2}, and hence
all its plurigenera.
Their growth proves $\kappa(M_r)=1$.
The line $2K_{M_r}=rF$ is nef,
so adjunction excludes a rational
$(-1)$-curve; $M_r$ is minimal.

\item The surface $B_r$ has Euler number
zero: every fibre over its elliptic
base has elliptic reduced support,
and the map is locally a smooth
elliptic bundle off finitely many
values. The finite involution
quotient therefore has Euler
number $(0+16)/2=8$. Resolving
sixteen $A_1$ points gives
$e(M_r)=24$. Its canonical
divisor is a sum of disjoint
square-zero fibres, so
$K_{M_r}^2=0$. Surface
Riemann--Roch gives
$\chi(\OO_{M_r})=24/12=2$.
The first plurigenus in \eqref{cdp:eq:C.ET.2}
is one; hence $q(M_r)=0$.
Projectivity, Hodge symmetry and
Poincar\'e duality give \eqref{cdp:eq:C.ET.3}.

\item The ambient quotient has the following resolution.
The only fixed base points of
$h$ lie at $E_0[2]\times\{0,\infty\}$,
and none was blown up. Each has
four fixed elliptic-fibre points.
The thirty-two ambient fixed
points therefore have local
linear action $(-1,-1,-1)$.
The invariant quadratic
blowup charts of \ref{cdp:st:C.2.UN.13} resolve
each quotient point smoothly.
For example, write upstairs
$(x,t,w)=(s,su,sv)$ and put
$R=s^2$. The resolved quotient
chart has coordinates $(R,u,v)$
and
\begin{equation}
                              T=Ru^2 .                \hypertarget{tag-00J7}{}\label{cdp:eq:C.ET.11}
\end{equation}
The exceptional surface is $\PP^2$.
Its intersection with the strict
central component is a line;
restriction is exactly the
minimal $A_1$ surface resolution.
Thus the strict component is
the same smooth $M_r$ proved
above, and the monomial gives
its multiplicity two and the
exceptional multiplicity one.
The paired root Hopf surfaces
give $r$ components of
multiplicity four, exactly as
in \ref{cdp:st:C.2.UN.25}. This proves \eqref{cdp:eq:C.ET.4}
with all full-fibre coefficients.

\item The reduced support is a star:
each $P_\nu$ meets $M_r$ along
one $\PP^1$, each $Q_j$ meets
$M_r$ along the elliptic curve
$J_j$, and there are no other
intersections or triple points.
It is simple normal crossing
by \eqref{cdp:eq:C.ET.11} and the root charts.
Therefore
\begin{equation}
                 e(S_r)=24+16\cdot3-16\cdot2=40.
                                                               \hypertarget{tag-00J8}{}\label{cdp:eq:C.ET.12}
\end{equation}
At infinity the same ambient
resolution gives the original
product-Kummer support
$2K+16\PP^2$ of \ref{cdp:st:C.2.UN.13}, also
of Euler number forty.
Twenty negative-two blocks
on each special support,
all beginning at distinct
smooth points of its plane
components away from the
K3 or $M_r$, make both
Euler numbers zero and
leave those distinguished
surfaces unchanged.

\item The support cohomology in \eqref{cdp:eq:C.ET.5} follows
through exact restriction maps.
The constant-sheaf union
sequence has component term
$M_r\sqcup\bigsqcup Q_j
\sqcup\bigsqcup P_\nu$
and curve term
$\bigsqcup J_j\sqcup
\bigsqcup\PP^1_\nu$.
Its degree-zero incidence
graph is a tree. In degree
one the $r$ Hopf generators
map injectively, each onto
the radial one-dimensional
subspace of its elliptic
attaching curve. The cokernel
has dimension $r$.
The degree-two restriction
has source dimension $38$
and target dimension $16+r$.
Its dual curve-class map
has kernel of dimension
$r-1$: the plane components
force each of their sixteen
edge coefficients to vanish;
the remaining elliptic
classes all equal the
nonzero class $\tfrac12F$
on $M_r$ (a full fibre meets
an ample divisor positively), leaving exactly
the relations whose
coefficients sum to zero.
The map thus has rank
seventeen. Its kernel
and cokernel have dimensions
$21$ and $r-1$.
The components contribute
$r$ in degree three and
$17+r$ in degree four.
The resulting long exact
sequence gives the first
Betti row of \eqref{cdp:eq:C.ET.5}.

\item The analogous structure-sheaf
sequence has isomorphisms
$H^1(Q_j,\OO)\to H^1(J_j,\OO)$,
given by \cref{lem:II:hopf-curve-restriction} for these same primary scalar Hopf surfaces and elliptic axis inclusions. The restriction identifies the additive deck class with the nonzero class on the axis.
The sole component
$H^2(\OO)$ is the one
on $M_r$. Hence the
coherent dimensions are
$(1,0,1)$.
Each of the twenty
balancing blocks adds
$(2,6,2)$ to support
cohomology in degrees
two, three and four,
by the restricted
exceptional Gysin maps
of \eqref{cdp:eq:C.KM.S2}, and
preserves coherent
cohomology by the Cartier pullback
in \eqref{cdp:eq:C.KM.S3}. This
proves the second row.

\item The canonical growth slopes
$r/2$ are distinct and
positive. Therefore these
components are non-uniruled,
essential by \cref{lem:II:nonuniruled-divisors} applied to these compact divisors in the smooth proper fillings, and
pairwise non-bimeromorphic.
They are the unique components
of Kodaira dimension one
in their respective supports;
the other components are
rational, Hopf, K3, or
the ruled blowup components.
The inverse-map resolution
comparison consequently
also distinguishes the
fillings over their same punctured family.

\end{enumerate}
\end{proof}






\section{Equivariant elliptic fillings and Hopf occurrence}\label{sec:II:equivariant-hopf-occurrence}

\subsection{Elliptic torsors and Hopf components}\label{subsec:II:enriques-elliptic-occurrence}

\subsubsection{Elliptic torsor lifts and product fillings}\label{subsubsec:II:elliptic-torsor-lifts}

\begin{lemma}[Equivariant lifts of elliptic torsors]\hypertarget{tag-P345}{}\label{lem:II:equivariant-elliptic-torsor-lift}
\StatementRecord{lem:II:equivariant-elliptic-torsor-lift}{proved}{P345}{\cref{def:II:elliptic-root-curve,lem:II:identity-monodromy-elliptic,lem:II:local-cyclic-elliptic}; \cite[Tag~0F0T, Lemma~59.95.4]{StacksProject}; \cite[Theorem~1.9, p.~51; proof pp.~78--79]{Steinberg1965}}
Let $q:E\to B$ be a proper holomorphic map with connected fibres
from a smooth compact complex surface to a smooth compact connected
curve. Suppose every fibre is a positive multiple of a smooth
elliptic curve, the smooth fibres have one fixed elliptic Jacobian
$F=\C/\Lambda$, and the integral $H_1$ monodromy is the identity.
Let a finite cyclic group $Q$ act on $E$ and $B$, with $q$
equivariant and the action on $B$ faithful; $Q=1$ is allowed.
Write $m_iC_i$ for the multiple fibres and
$\mathcal B=B[(p_i,m_i)_i]$ for \cref{def:II:elliptic-root-curve}.

The given surface is a principal $F$-torsor over $\mathcal B$.
Its local normalized charts are \eqref{cdp:eq:C.SF.2}, with
$F'=F$ through the marked group isomorphism in the proof.
The $Q$-action has a constant linear part
\begin{equation}
                 \beta:Q\longrightarrow\operatorname{Aut}(F,0).
 \hypertarget{tag-00HM}{}\label{cdp:eq:C.EO.2}
\end{equation}
Let $\ell:\Lambda\otimes_\Z\C\to\C$ extend the lattice inclusion.
There is a $Q$-equivariant exact sequence
\begin{equation}
 0\longrightarrow K\longrightarrow
 T=(\Lambda\otimes_\Z\C)/\Lambda
 \xrightarrow{\Phi}F\longrightarrow0,
 \qquad K=\ker\ell\simeq(\C,+),\qquad T\simeq(\C^*)^2.
 \hypertarget{tag-00HN}{}\label{cdp:eq:C.EO.3}
\end{equation}
The given equivariant torsor has a lift to a $T$-torsor $P$ on
$\mathcal B$, with its full $Q$-descent data, and an equivariant
isomorphism
\begin{equation}
                         P\mathbin{\times^T}F\simeq E.
 \hypertarget{tag-00HQ}{}\label{cdp:eq:C.EO.6}
\end{equation}
It admits a meromorphic section satisfying
\begin{equation}
                         g(s(b))=s(gb)\qquad(g\in Q).
 \hypertarget{tag-00HR}{}\label{cdp:eq:C.EO.7}
\end{equation}
For $M=\operatorname{Hom}(T,\C^*)$, let $L_\chi,s_\chi$ be the
associated character line and induced meromorphic section.
Then
\begin{equation}
 D:M\longrightarrow\operatorname{Div}(\mathcal B),\qquad
                 D(\chi)=\operatorname{div}(s_\chi)
 \hypertarget{tag-00HS}{}\label{cdp:eq:C.EO.8}
\end{equation}
is additive and $Q$-equivariant, and the identifications
\begin{equation}
       (L_\chi,s_\chi)\simeq
                   (\OO_{\mathcal B}(D(\chi)),1)
 \hypertarget{tag-00HT}{}\label{cdp:eq:C.EO.9}
\end{equation}
commute with addition of characters and the complete $Q$-action.
The supports of a basis pair of these divisors and all marked
points lie in a finite $Q$-invariant set. Their coarse
coefficients are integral off the marked points and belong to
$m_i^{-1}\Z$ at $p_i$.
\end{lemma}
\begin{proof}
\begin{enumerate}[label=\arabic*.,leftmargin=*]
\item Choose an oriented ordered basis $(\lambda_1,\lambda_2)$ of
the period lattice $\Lambda$. Choose the integral marking on one
smooth fibre to give this basis under its Jacobian identification
with $F$; identity monodromy extends that marking over the smooth
base. Every fibre support is elliptic, so no fibre contains a
rational $(-1)$-curve and $q$ is relatively minimal. Apply
\cref{lem:II:identity-monodromy-elliptic} to this same $q:E\to B$.
For the chosen marking $(\lambda_1,\lambda_2)$, its normalized
fixed curve $A=\C/(\Z+\tau\Z)$, $\tau=\lambda_2/\lambda_1$,
is identified with the specified $F$ by
$[z]\mapsto[\lambda_1z]$, with inverse $[a]\mapsto[a/\lambda_1]$.
These maps carry the normalized periods to the given marking.
The marked integration products therefore give the same
translation torsor on the smooth base.

Apply \cref{lem:II:local-cyclic-elliptic} with $G=1$ and the
same $E,q,B$. Its kernel condition is then the identity
translation. Its root chart \eqref{cdp:eq:C.EQ.3}, including
the averaging coordinate change and its inverse, becomes
\eqref{cdp:eq:C.SF.2} through
$(a',u)\mapsto([\lambda_1a'],u)$.
It carries the exact-order root translation to the same
$\tau_i\in F[m_i]$ and commutes with the map $t=u^{m_i}$.
The normalized root projection in
\cref{def:II:elliptic-root-curve} identifies this quotient
with the original surface tube, with its punctured gluing.
Thus the root-chart torsor has total space the given $E$.
The functorial root lift
of each $g\in Q$ supplies the action on $\mathcal B$.
Its origin-preserving fibre part is constant on the connected
regular base, since $\operatorname{Aut}(F,0)$ is discrete.
This gives \eqref{cdp:eq:C.EO.2}; its translation part remains
part of the torsor descent maps.

\item For this ordered lattice basis, the map
$(v_1,v_2)\mapsto\lambda_1\otimes v_1+\lambda_2\otimes v_2$ identifies
$\C^2/\Z^2$ with $T$. Exponentiating $2\pi i v_j$ gives
the specified isomorphism $T\simeq(\C^*)^2$, in which
\begin{equation}
 \Phi(z_1,z_2)=
 \left[\frac{\lambda_1\log z_1+\lambda_2\log z_2}{2\pi i}\right]
                         \in\C/\Lambda.
 \hypertarget{tag-00HO}{}\label{cdp:eq:C.EO.4}
 \hypertarget{tag-010X}{}\SsixEquationAliases{\label{cdp:eq:C.SF.3}}
\end{equation}
Integral logarithm changes give precisely lattice periods.
The derivative is surjective, so its image is an open subgroup
of the connected group $F$, hence all of $F$. A vector mapping
to a period may be changed by that integral lattice vector to
lie in $\ker\ell$. The intersection $\ker\ell\cap\Lambda$ is
zero. In the displayed coordinates its kernel embedding is
\begin{equation}
 w\longmapsto
     \bigl(\exp(2\pi i\lambda_2w),
                         \exp(-2\pi i\lambda_1w)\bigr).
 \hypertarget{tag-010Y}{}\label{cdp:eq:C.SF.4}
\end{equation}
This is injective because a nonzero integral relation between
the two lattice generators is impossible. It has closed image,
the kernel of $\Phi$, and local logarithms give its holomorphic
inverse. This proves \eqref{cdp:eq:C.EO.3}, including its
additive kernel and local splitting.

The action of $\beta(g)$ on $\Lambda$ gives the corresponding
monomial automorphism of $T$. Its determinant on the oriented
integral lattice is one; its quotient character on $\C$ is
$\beta(g)$, so the character on $K$ is $\beta(g)^{-1}$.
Thus \eqref{cdp:eq:C.EO.3} is equivariant. On
$\mathcal C=[\mathcal B/Q]$ the section sheaves give
\begin{equation}
 0\longrightarrow\mathcal K\longrightarrow\mathcal T
                 \longrightarrow\mathcal F\longrightarrow0,
 \hypertarget{tag-00HP}{}\label{cdp:eq:C.EO.5}
\end{equation}
where $\mathcal K$ is the additive sheaf of its holomorphic
line bundle and the other sheaves carry their indicated finite
linear monodromy. The given equivariant $E$ defines its class
in $H^1(\mathcal C,\mathcal F)$.

\item On a finite root chart, averaging in characteristic zero makes
invariants of coherent sheaves exact. The invariant direct image
to the coarse curve is coherent and has no higher direct images;
the same holds for the finite $Q$ quotient. Coherent cohomology
in degree two on that curve vanishes. Consequently
$H^2(\mathcal C,\mathcal K)=0$, and the long exact sequence
of \eqref{cdp:eq:C.EO.5} lifts the given class to a
$\mathcal T$-torsor on $\mathcal C$. Pull it back to
$\mathcal B$ and call it $P$. Equality of the pushed-out class
with the class of $E$ supplies \eqref{cdp:eq:C.EO.6} in the
equivariant category: its transition and descent maps satisfy
the full cocycle identities.

For $Q=1$, the chosen basis identifies this exact sequence with
\begin{equation}
 0\longrightarrow\OO_{\mathcal B}
 \longrightarrow(\OO_{\mathcal B}^*)^2
 \xrightarrow{\Phi}\mathcal F_{\mathcal B}\longrightarrow0.
 \hypertarget{tag-010Z}{}\label{cdp:eq:C.SF.5}
\end{equation}
Here $\mathcal F_{\mathcal B}$ denotes the sheaf of holomorphic
maps to $F$. The same vanishing gives
\begin{equation}
 \operatorname{Pic}(\mathcal B)^2
                  \longrightarrow H^1(\mathcal B,\mathcal F_{\mathcal B})
                  \quad\hbox{surjective}.
 \hypertarget{tag-0110}{}\label{cdp:eq:C.SF.6}
\end{equation}
The two associated basis-character lines of $P$ are exactly the
pair mapping to the prescribed analytic torsor class. Each has
an expression
\begin{equation}
 L_a\simeq\pi^*M_a\otimes\bigotimes_i\mathcal R_i^{e_{ai}},
                         \qquad a=1,2.
 \hypertarget{tag-0111}{}\label{cdp:eq:C.SF.7}
\end{equation}
Choose $e_{ai}$ to be its inertia character modulo $m_i$.
After removing those root characters, the invariant local
frames descend to an ordinary line $M_a$ on $B$. Their inverse
frame maps prove the displayed line-bundle isomorphisms.

\item Put $k=\C(B/Q)$ and $L=\C(B)$. Faithfulness makes $L/k$
the cyclic extension with group $Q$. At the generic point there
is no root inertia. Monomial descent of $(\mathbb G_m)^2_L$
defines the algebraic $k$-torus $T_k$. The generic fibre of $P$
is an algebraic torsor under it: pull back to $\mathcal B$,
separate the finite root weights as above, and choose meromorphic
sections of the ordinary character lines on the compact curve.
The descent maps are monomials multiplied by meromorphic
functions in $L$, with the already established cocycle identities.

Apply \cite[Tag~0F0T, Lemma~59.95.4]{StacksProject} to the field
extension $k/\C$. The constant field is algebraically closed and
the transcendence degree is one, giving $\operatorname{cd}(k)\le1$;
the same is true for every finite extension of $k$. Characteristic
zero makes $k$ perfect. In Steinberg's cohomological-dimension
convention, the comparison follows also from the Kummer maps:
degree-two cohomology with roots of unity vanishes over those
finite extensions, so their torsion Brauer groups vanish.
Apply \cite[Theorem~1.9, p.~51; proof pp.~78--79]{Steinberg1965}
to this connected linear algebraic group $T_k$. It gives
\begin{equation}
                             H^1(k,T_k)=0.
 \hypertarget{tag-00HU}{}\label{cdp:eq:C.EO.H1}
\end{equation}
The resulting $k$-point of the generic torsor pulls back to the
meromorphic section of $P$ on $\mathcal B$ satisfying
\eqref{cdp:eq:C.EO.7} with its given descent action.

For every character, sending the canonical divisor section to
$s_\chi$ gives \eqref{cdp:eq:C.EO.9}; locally the matching
divisors make its multiplier a nowhere-zero holomorphic unit.
The invariant section gives
$D(g\chi)=gD(\chi)$, for $(g\chi)(t)=\chi(g^{-1}t)$, and
multiplication of character sections proves additivity.
In a root chart a meromorphic section has an integral order in
the root coordinate, so its coarse order lies in $m_i^{-1}\Z$.
There are finitely many zeroes and poles on the compact root
curve. Take the $Q$-orbit of their supports for a basis pair,
together with the marked points, to obtain the asserted set.
\end{enumerate}
\end{proof}


\begin{proposition}[Elliptic surfaces in product torus fillings]\hypertarget{tag-01DD}{}\label{cdp:st:C.2.UN.34}
\StatementRecord{cdp:st:C.2.UN.34}{proved}{01DD}{\cref{lem:II:equivariant-elliptic-torsor-lift,def:II:elliptic-root-curve,cdp:st:C.2.UN.22,cdp:st:C.2.UN.29}; \eqref{cdp:eq:C.SF.8}--\eqref{cdp:eq:C.SF.13}}

Let $q:S\to B$ be a proper holomorphic map from a
smooth compact connected complex surface to a smooth
compact curve with $g(B)\le1$. Assume that every
fibre is a positive multiple of a smooth elliptic
curve, that its smooth fibres have one fixed elliptic
modulus, and that its integral $H_1$ monodromy is
the identity. Then there are a smooth compact complex
threefold $Y$, a proper holomorphic map
$f:Y\to\mathbb P^1$, and a specified isomorphism
from $S$ to a component of its reduced zero fibre,
with the following properties:

\begin{itemize}
\item the component $S$ is unchanged and has
multiplicity one in the full fibre;
\item over $\mathbb P^1\setminus\{0\}$ the map
is a product with a fixed product of two
elliptic curves, and hence has connected torus
fibres and identity monodromy;
\item the reduced zero fibre has simple normal
crossings, all its components are smooth elliptic
quasi-bundles, all its double curves are smooth
elliptic, and it has no triple points;
\item every reduced fibre, including the zero
fibre, has Euler characteristic zero.
\end{itemize}
This realizes every elliptic covering model in
\cref{cdp:st:C.2.UN.22} and \cref{cdp:st:C.2.UN.29}. The component isomorphism preserves the full holomorphic torsor
class of $S$, including its degree-zero parameters and
its individual multiple-fibre translations.
\end{proposition}



\begin{proof}

\begin{enumerate}[label=\arabic*.,leftmargin=*]
\item Write the fixed elliptic curve as $A=\C/\Lambda$ and choose
the ordered basis $(\lambda_1,\lambda_2)$ of its period lattice.
Apply \cref{lem:II:equivariant-elliptic-torsor-lift} with
$E=S$, $F=A$, the given $q:S\to B$, and $Q=1$.
Each fibre is connected because its support is one smooth
elliptic curve, and its positive multiplicity gives surjectivity.
The multiple values and orders are the same $p_i,m_i$, so its
root curve is exactly \eqref{cdp:eq:C.SF.1}. The normalized root
map and marked rescaling in that lemma identify the original
surface tubes with \eqref{cdp:eq:C.SF.2}, including the given
translations $\tau_i$ and the punctured torsor transition maps.

Let $\mathcal T$ be that principal $A$-torsor, with total space
$S$. In the chosen basis the lemma's map $\Phi$ is
\eqref{cdp:eq:C.SF.3}, with kernel \eqref{cdp:eq:C.SF.4}.
Its $Q=1$ sequence \eqref{cdp:eq:C.SF.5} and surjection
\eqref{cdp:eq:C.SF.6} give the two basis-character lines
$(\mathcal L_1,\mathcal L_2)=(L_1,L_2)$ of its lift $P$.
The isomorphism $P\times^T A\simeq S$ identifies
their full analytic torsor class with that of the specified
surface. Write their root-line expressions as
\eqref{cdp:eq:C.SF.7}, with the same $M_a,e_{ai}$ and canonical
root lines.

\item Represent each $M_a$ by an ordinary
signed divisor
\begin{equation}
                        M_a\simeq\mathcal O_B\left(
                                    \sum_p d_{ap}p\right)
                                                               \hypertarget{tag-0112}{}\label{cdp:eq:C.SF.8}
\end{equation}
whose support avoids the finitely
many marked points $p_i$. It may
also avoid any one further chosen
point. Take
an effective divisor away from
that finite set, twist $M_a$ by
a sufficiently large multiple,
and use a section which is nonzero
at all points of the finite set;
curve Riemann--Roch and global
generation provide it. Subtracting
the chosen twist gives \eqref{cdp:eq:C.SF.8}.
This represents degree-zero classes
as well as classes of nonzero
degree.

\item For elliptic $B$, take
$V_0=B\times\mathbb P^1$ with its
projection and selected component
$B=B\times\{0\}$. If $B=\mathbb P^1$,
start with any product
$A_0\times\mathbb P^1$ for an
elliptic curve $A_0$, blow up one
point over zero, and identify its
exceptional curve with the prescribed
$B$. Its single intersection with
the strict central elliptic curve
may be chosen away from all $p_i$;
choose \eqref{cdp:eq:C.SF.8} away from it as well.
In both cases the selected component
has multiplicity one and the
auxiliary central divisor has
simple normal crossings.

\item Blow up the distinct points of
the selected $B$ consisting of
all $p_i$ and all supports in
\eqref{cdp:eq:C.SF.8}. Call the resulting
surface $V$, its map to
$\mathbb P^1$ $h$, and the
exceptional curves $T_i,T_p$.
They are pairwise disjoint and
vertical, each meets the strict
transform of $B$ transversely
at its specified point, and
\begin{equation}
 \mathcal O_V(T_i)|_B=\mathcal O_B(p_i),
 \qquad\mathcal O_V(T_p)|_B=\mathcal O_B(p).
                                                               \hypertarget{tag-0113}{}\label{cdp:eq:C.SF.9}
\end{equation}
These are the restriction
maps of the exceptional divisors.
Every component of $h^{-1}(0)$
still has multiplicity one; its
reduced support is a tree of
smooth curves with no triple
points. The strict $B$ remains
isomorphic to the given curve.

\item Root $V$ along the disjoint
$T_i$ with their specified
orders, obtaining
$\mathcal V=V[(T_i,m_i)_i]$.
The inverse image of the strict
$B$ is exactly \eqref{cdp:eq:C.SF.1}, since
the intersections are transverse.
Write $\widetilde{\mathcal R}_i$
for the corresponding root lines
on $\mathcal V$. The vertical
extensions
\begin{equation}
 \widetilde{\mathcal L}_a
   =\mathcal O_V\left(\sum_p d_{ap}T_p\right)
       \otimes\bigotimes_i
                       \widetilde{\mathcal R}_i^{e_{ai}}
                          \qquad(a=1,2)
                                                               \hypertarget{tag-010S}{}\label{cdp:eq:C.SF.10}
\end{equation}
restrict, through \eqref{cdp:eq:C.SF.9}, to
exactly $\mathcal L_a$. Here
ordinary line bundles are pulled
back to $\mathcal V$ in the
displayed formula.

\item Push the pair in \eqref{cdp:eq:C.SF.10}
out by the group map
\eqref{cdp:eq:C.SF.3}. Explicitly, if their
nonvanishing local transition
functions are $g^1_{\alpha\beta}$
and $g^2_{\alpha\beta}$, glue
the products with $A$ by
\begin{equation}
 (x,a)_\beta=(x,a+
     \Phi(g^1_{\alpha\beta}(x),
          g^2_{\alpha\beta}(x)))_\alpha.
                                                               \hypertarget{tag-010T}{}\label{cdp:eq:C.SF.11}
\end{equation}
The triple-overlap equality follows
from the line-bundle cocycles and
the homomorphism property of $\Phi$.
Local logarithm choices change the
translation only by an element
of $\Lambda$. This produces a
principal $A$-torsor over
$\mathcal V$ whose restriction
is isomorphic to $\mathcal T$
by \eqref{cdp:eq:C.SF.6}--\eqref{cdp:eq:C.SF.10}.

\item Its coarse space $Y$ is smooth.
In a root chart transverse to
$T_i$ it is precisely
\begin{equation}
 (A\times\Delta_w\times\Delta_x)/
 \langle(a,w,x)\mapsto
             (a+\tau_i,\zeta_{m_i}w,x)\rangle,
 \qquad z=w^{m_i}.
                                                               \hypertarget{tag-010U}{}\label{cdp:eq:C.SF.12}
\end{equation}
The translation is the same
$\tau_i$ as in the original
torsor, because the inertia
characters in \eqref{cdp:eq:C.SF.7} are
preserved. Its exact order
is $m_i$, so the finite action
in \eqref{cdp:eq:C.SF.12} is free, including
on $w=0$. Thus the quotient
is a complex manifold. There
are no simultaneous root
intersections. Elsewhere it
is an ordinary elliptic bundle.
These charts also prove
properness over $V$: locally
the map is a finite quotient
of a product with the compact
curve $A$. Since $V$ is
compact, $Y$ is compact.

\item Compose with $h$ to obtain
$f:Y\to\mathbb P^1$.
The component over the strict
$B$ is exactly the coarse
torsor $S$ in \eqref{cdp:eq:C.SF.2}, with
the specified isomorphism
coming from its identical
$H^1$ class. At a generic
point no root divisor is
present and the auxiliary
component has multiplicity
one, so $S$ also has
multiplicity one in the
full fibre of $f$.

\item Every ordinary vertical line
in \eqref{cdp:eq:C.SF.10} has its canonical
meromorphic divisor section,
nonvanishing off the central
fibre. Every root line has
its tautological section,
also nonvanishing there.
Thus both line bundles are
canonically trivial on the
entire complement of zero,
and \eqref{cdp:eq:C.SF.11} gives the
product identification
\begin{equation}
 f^{-1}(\mathbb P^1\setminus\{0\})
   \simeq A\times A_0\times
                         (\mathbb P^1\setminus\{0\}),
                                                               \hypertarget{tag-010V}{}\label{cdp:eq:C.SF.13}
\end{equation}
where $A_0=B$ in the elliptic
base case. The fibres of
$Y\to V$ are connected
elliptic curves, including
the finite translation
quotients over root divisors.
Since $h$ has connected
fibres, so does $f$.

\item Finally the auxiliary central
support is an SNC curve
with no triple points.
Its inverse image in the
free quotient charts
\eqref{cdp:eq:C.SF.12} has coordinate
normal crossings. For
example a crossing whose
root coordinate is $z$
has local full equation
$w^{m_i}x=0$. Every
component of the reduced
support is a smooth
elliptic quasi-bundle over
one auxiliary curve; its
special fibres still have
smooth elliptic reduction.
Its Euler characteristic
is therefore zero by
additivity over the smooth
base stratum and the
finitely many multiple
values. Every double
curve is one smooth
elliptic translation
quotient, and there are
no triple points.
Inclusion--exclusion gives
$e(f^{-1}(0)_{\mathrm{red}})=0$.
All other fibres in
\eqref{cdp:eq:C.SF.13} are smooth tori,
so have the same zero
Euler characteristic.


\end{enumerate}
\end{proof}






\subsubsection{Cyclic quotient fillings}\label{subsubsec:II:cyclic-elliptic-quotients}

\begin{proposition}[Cyclic elliptic quotients in torus fillings]\hypertarget{tag-01DE}{}\label{cdp:st:C.2.UN.35}
\StatementRecord{cdp:st:C.2.UN.35}{proved}{01DE}{\cref{lem:II:equivariant-elliptic-torsor-lift,lem:II:bielliptic-orbit-pencil,lem:II:local-cyclic-elliptic,cdp:st:C.2.UN.31,cdp:st:C.2.UN.23,cdp:st:C.2.UN.32,cdp:st:C.2.UN.1,cdp:st:C.2.UN.4,cdp:st:C.2.UN.20}; \cite[Theorem~13.2, pp.~297--298; Theorem~1.6(1), p.~215]{BierstoneMilman1997}}

Let $q:E\to B$ be a connected compact smooth elliptic
surface with a fixed elliptic Jacobian $F$, identity
integral $H_1$ monodromy, and only smooth multiple
fibres. Suppose $g(B)\le1$. Let $Q$ be a finite
cyclic group acting holomorphically on $E$, preserving
$q$, and acting faithfully on $B$. The trivial group
is allowed. The smooth minimal model of $E/Q$ occurs
as the minimal normalization model of a component of
a fibre in a smooth compact threefold over $\PP^1$
whose general fibre is a complex two-torus. It also
occurs in such a fibre with reduced-support Euler
number zero. Before taking the quotient, the distinguished component of the
constructed covering threefold is $Q$-equivariantly
isomorphic to $E$ over $B$.

Consequently the minimal normalization models of
nonnegative Kodaira dimension occurring in arbitrary
two-torus degenerations are exactly the following
two classes:
\begin{itemize}
\item the minimal resolutions of finite cyclic
quotients of complex two-tori that have nonnegative
Kodaira dimension;
\item the minimal models of $E/Q$ in the preceding
construction that have nonnegative Kodaira dimension.
\end{itemize}
The same two classes occur among zero-Euler fillings.
The torus-quotient class is described completely in
\cref{cdp:st:C.2.UN.23} and \cref{cdp:st:C.2.UN.32}. In the second class, all multiplicities,
translation gluing, and the cyclic action are part
of the input and of the resulting equivariant component isomorphism.
\end{proposition}



\begin{proof}\leavevmode\par

\begin{enumerate}[label=\arabic*.,leftmargin=0pt,labelsep=.5em]
\item \par\medskip\noindent\textbf{The prescribed equivariant torsor.}
Apply \cref{lem:II:equivariant-elliptic-torsor-lift} to the
given $q:E\to B$, fixed curve $F=\C/\Lambda$ and the given
cyclic action $Q$. Its hypotheses are exactly the connected
smooth-multiple, fixed-Jacobian and identity-monodromy data in
the statement, including faithfulness on $B$.
Let $m_iC_i$ lie over $p_i$ and use the same
$\mathcal B=B[(p_i,m_i)]$ of
\cref{def:II:elliptic-root-curve}. The normalized maps in that
definition and the marked product maps of the lemma identify
\eqref{cdp:eq:C.EO.1} with these original surface tubes.

Take the linear character $\beta$, group $T$, additive kernel
$K$ and map $\Phi$ of \eqref{cdp:eq:C.EO.2}--\eqref{cdp:eq:C.EO.4}.
On $\mathcal C=[\mathcal B/Q]$ the exact sheaf sequence is
\eqref{cdp:eq:C.EO.5}. The lemma gives its lift $P$, the
$Q$-equivariant isomorphism \eqref{cdp:eq:C.EO.6} with this
same $E$, and the invariant meromorphic section
\eqref{cdp:eq:C.EO.7}. For every character
$\chi\in M=\operatorname{Hom}(T,\C^*)$, use its associated
line and section $L_\chi,s_\chi$ and the divisor
$D(\chi)$ in \eqref{cdp:eq:C.EO.8}. The isomorphisms
\eqref{cdp:eq:C.EO.9} preserve both multiplication and the
full descent action. Fix two basis characters and let
$\Sigma\subset B$ be the finite $Q$-invariant union of
their divisor supports and the marked points. The coefficient
of each $D(\chi)$ at $p_i$ is in $m_i^{-1}\Z$, and the
other coefficients are integral.

\item \par\medskip\noindent\textbf{An equivariant vertical extension.}
We construct a smooth projective surface
$V_0\to\PP^1_t$ containing $B$ in its
central fibre and extending its given
$Q$-action. Let $n=|Q|$.
For elliptic $B$, choose any faithful
character $\xi:Q\to\C^*$ when $n>1$
(and $\xi=1$ for the trivial group),
and take $V_0=B\times\PP^1_t$
with action $(b,t)\mapsto(gb,\xi(g)t)$,
and use $B\times\{0\}$.
For $B\simeq\PP^1$, first diagonalize
its given cyclic action as
$z\mapsto\lambda(g)z$ and define
$\xi=\lambda^{-1}$. Faithfulness makes
this a faithful disc character too.
For the trivial group use both
characters equal to one. Choose any
elliptic curve $A$ and point $o\in A$.
Blow up $(o,0)$ in $A\times\PP^1_t$,
where the action is
$(a,t)\mapsto(a,\xi(g)t)$.
The exceptional curve is identified with
the given $B$: in a local coordinate $x$
at $o$, its affine coordinate $z=x/t$
has exactly the required action. The
other chart supplies its point at infinity.


\begin{enumerate}[label=(\arabic*),leftmargin=1.3em,labelsep=.5em]
\item  Blow up every point of $\Sigma$ on the
distinguished strict curve $B$, in complete
$Q$-orbits. Call the resulting disjoint
exceptional curves $T_p$, $p\in\Sigma$,
and call the new surface $V$. The strict
curve is still isomorphic to $B$, and
\begin{equation}
 \OO_V(T_p)|_B=\OO_B(p),\qquad
                     g(T_p)=T_{gp}.                \hypertarget{tag-00HF}{}\label{cdp:eq:C.EO.10}
\end{equation}
These identities include the canonical
sections and their equivariant structures.
This restriction map also holds at a fixed point and at a node of the original central fibre. At a node with coordinates $xy=0$
and $B=(x=0)$, its blowup has chart
$x=uv,y=v$. Here the strict $B$ is
$u=0$, $T_p$ is $v=0$, and the
restriction of its canonical section
is the simple parameter $v$ on $B$.
This proves \eqref{cdp:eq:C.EO.10} also in that case.
The other central branch meets $T_p$
at the other point of this exceptional
curve; it does not meet the strict $B$.


\item  Root $V$ along each $T_{p_i}$ to order
$m_i$, obtaining the orbifold surface
$\mathcal V$. Its local chart at a
point of $T_{p_i}$ is
\begin{equation}
 [\Delta_u\times\Delta_v/\mu_{m_i}],
           \qquad x=u^{m_i},\qquad
           \zeta\cdot(u,v)=(\zeta u,v).
                                                               \hypertarget{tag-00HG}{}\label{cdp:eq:C.EO.11}
\end{equation}
The $T_{p_i}$ are disjoint, so there
is only this one inertia group in
each such chart. Its restriction to
$B$ is precisely the previous root
orbicurve $\mathcal B$, including the
functorial $Q$-action. The transverse
intersection in \eqref{cdp:eq:C.EO.10} is what
identifies its root line and root
section with those on $\mathcal B$.


\item  Replace each coarse point $p$ in
$D(\chi)$ by $T_p$, with the same
coefficient. Denote the resulting
orbifold Cartier divisor by
$\widetilde D(\chi)$. All its
coefficients are legitimate integral
root-divisor coefficients on
$\mathcal V$. The map
$\chi\mapsto\widetilde D(\chi)$ is
additive and $Q$-equivariant because
the replacement $p\mapsto T_p$ is.
The character lines
$\OO_{\mathcal V}(\widetilde D(\chi))$
therefore define a principal $T$-torsor
$\widetilde P$ with its exact monomial
$Q$-action. Define the principal
elliptic bundle
\begin{equation}
              \widetilde Y
                =\widetilde P\mathbin{\times^T}F
                     \longrightarrow\mathcal V .   \hypertarget{tag-00HH}{}\label{cdp:eq:C.EO.12}
\end{equation}
Concretely, on every line-bundle
trivialization its transition map is
translation on $F$ by the value under
\eqref{cdp:eq:C.EO.4} of the two multiplicative
transition functions. Equations
\eqref{cdp:eq:C.EO.9}--\eqref{cdp:eq:C.EO.10} give a $Q$-equivariant isomorphism
\begin{equation}
                   \widetilde Y|_{\mathcal B}
                                   \simeq E,         \hypertarget{tag-00HI}{}\label{cdp:eq:C.EO.13}
\end{equation}
over the given base. Locally the meromorphic
frame identification extends
holomorphically because the two
divisors agree: their quotient is
a nowhere-zero holomorphic unit.
The same assertion holds in the
root coordinates. Since both
frames in \eqref{cdp:eq:C.EO.9} are exactly
equivariant, these local isomorphisms
also intertwine the given $Q$-action.
More explicitly, a lift of a base
map to $x=u^{m_i}$ involves a
choice of an $m_i$-th root of
its nonzero local multiplier.
Two choices differ by the
root inertia in \eqref{cdp:eq:C.EO.11}.
The character lines simultaneously
apply the corresponding elliptic
translation in \eqref{cdp:eq:C.EO.14}.
Hence the two lifts induce the
same map on the total coarse
quotient. Compositions of lifts
agree there for the same reason.
This is the finite-chart form
of the canonical section-preserving
action; all root choices induce this same group law.


\end{enumerate}

\item \par\medskip\noindent\textbf{Smoothness, properness and the generic torus.}
Let $(a_1,a_2)\pmod {m_i}$ be the
two weights of $P$ at $p_i$.
In \eqref{cdp:eq:C.EO.11} the total elliptic
bundle has the finite chart
\begin{equation}
 (\Delta_u\times\Delta_v\times F)/\mu_{m_i},
 \quad
 \zeta_{m_i}(u,v,z)
     =(\zeta_{m_i}u,v,z+\tau_i),\quad
 \tau_i=\frac{a_1\lambda_1+a_2\lambda_2}{m_i}
                                           \pmod\Lambda .
                                                               \hypertarget{tag-00HJ}{}\label{cdp:eq:C.EO.14}
\end{equation}
The order of $\tau_i$ is exactly
$m_i$ by \eqref{cdp:eq:C.EO.1} and \eqref{cdp:eq:C.EO.6}.
Every nonidentity inertia element
therefore acts freely on the total
chart, including $u=0$ and every
central crossing $v=0$. Thus the
coarse space in \eqref{cdp:eq:C.EO.12} is a
smooth compact complex threefold;
we continue to call it
$\widetilde Y$. Properness follows
on these finite charts from
properness of the compact factor
$F$ and glues over the compact
surface $V$. The principal-bundle
transition functions in \eqref{cdp:eq:C.EO.12}
give the gluing maps explicitly.


\begin{enumerate}[label=(\arabic*),leftmargin=1.3em,labelsep=.5em]
\item  All divisors used to construct
$\widetilde P$ are vertical over
$t=0$. Their canonical meromorphic
sections are therefore nowhere
zero on the complement of that
fibre. There \eqref{cdp:eq:C.EO.12} is the
product
\begin{equation}
 \widetilde Y|_{t\ne0}
       \simeq A\times F\times(\PP^1\setminus\{0\}),
                                                               \hypertarget{tag-00HK}{}\label{cdp:eq:C.EO.15}
\end{equation}
where $A=B$ in the elliptic-base
case and $A$ is the chosen auxiliary
elliptic curve in the rational-base
case. The action is the prescribed
action on $A$ (the identity in the
rational-base construction),
the linear action $\beta$ on $F$,
and multiplication by $\xi$ on
$t$. The invariant divisor frames identify its translation cocycle
with the identity section on this complement.


\item  Put $X=\widetilde Y/Q$. Apply
\cite[Theorem~13.2, pp.~297--298, with Theorem~1.6(1),
p.~215]{BierstoneMilman1997} to this compact reduced complex
space. Its resolution $\rho:Y\to X$ is a finite composite of
blowups and is an isomorphism over the regular locus of $X$. Its map to
$\PP^1_T$, $T=t^n$, has general
fibre exactly $A\times F$:
the $n$ covering fibres are
permuted simply transitively
for $T\ne0,\infty$. The quotient
is smooth there, so the
resolution leaves those fibres
unchanged. The Stein factor is
a finite map of degree one
to the normal curve $\PP^1$,
and is the identity. Thus
the resulting proper map has
connected fibres.


\item  The distinguished component
before the quotient is exactly
the given $E$ by \eqref{cdp:eq:C.EO.13}.
Its image in the finite quotient
is $E/Q$: taking invariants
is exact in characteristic zero,
so the restriction of the
quotient structure sheaf to
this image is
$(\OO_E)^Q$. It is normal
because $E$ is normal. The
strict component on the ambient
resolution maps properly and
birationally to $E/Q$.
Resolving that component and
contracting its surface
$(-1)$-curves gives exactly
the prescribed minimal model.
Finally \cref{cdp:st:C.2.UN.1} gives even Euler
number for the reduced support.
The vertical point and
quartic blowups of \cref{cdp:st:C.2.UN.4} change
that number by $+2$ and $-2$
respectively. Finitely many
such operations make it zero
and preserve the minimal
normalization model through
their explicit proper
birational maps. 


\end{enumerate}

\item \par\medskip\noindent\textbf{Exhaustion of the nonnegative classes.}
Let $Z$ be a minimal normalization
model of nonnegative Kodaira
dimension in an arbitrary
two-torus degeneration. \cref{cdp:st:C.2.UN.31}
provides its cyclic
component cover and equivariant
minimal contraction. Let $E_0$ be the resulting smooth minimal
model of that cover, and let $G$ be its descended cyclic group. A torus
cover gives exactly the first
class in the statement. In
the remaining case \cref{cdp:st:C.2.UN.22} and
\cref{cdp:st:C.2.UN.29} give a fixed-Jacobian,
identity-monodromy fibration
with only smooth multiples
over a rational or elliptic
curve. 


\begin{enumerate}[label=(\arabic*),leftmargin=1.3em,labelsep=.5em]
\item  For a cover of Kodaira
dimension one it is the
intrinsic connected Iitaka
fibration. In algebraic
dimension one its base is
the intrinsic meromorphic
function field, and the
equivariant Stein factor
of \cref{cdp:st:C.2.UN.20} applies. For the non-torus projective case of Kodaira dimension zero,
the surface classification used in \cref{cdp:st:C.2.UN.31}
and its equality $\chi(\OO_{E_0})=0$ make $E_0$ bielliptic.
Apply \cref{lem:II:bielliptic-orbit-pencil} to its given free
product presentation $(F_1\times F_2)/\Gamma$.
The orbit map \eqref{cdp:eq:C.EO.16} is the resulting map on
this same $E_0$, with fixed fibre
$F_1/\ker(\Gamma\to\operatorname{Aut}(F_2))$ and base
$F_2/\Gamma\simeq\PP^1$. The lemma proves that every
automorphism preserves it, so the given cyclic action does.
It has the required identity monodromy and smooth multiples.


\item  Let $H$ be the kernel of
the cyclic action on the
chosen base. Its action
on a general elliptic
fibre is by translations:
a nontrivial linear
elliptic automorphism
has rational quotient,
which would make the
quotient surface ruled
and annihilate all
positive pluricanonical
sections, contrary to
$\kappa(Z)\ge0$.
Apply \cref{lem:II:local-cyclic-elliptic} to this chosen
$q:E_0\to B$, cyclic group $G$ and kernel $H$.
The elliptic map is relatively minimal, has the stated identity
monodromy, and the preceding argument makes its kernel act by
translations. The lemma's \eqref{cdp:eq:C.EQ.3}--\eqref{cdp:eq:C.EQ.4}
are therefore the root and quotient maps for these same
$E_0,G,H$. They give the smooth surface $E_1=E_0/H$, its fixed
Jacobian and smooth multiples, and the faithful residual action
$Q=G/H$ on $B$. The quotient equality
$(E_0/H)/(G/H)=E_0/G$ is induced by sending an $H$-orbit to
its $G$-orbit, with the inverse on quotient orbits.
Thus its minimal model is the same $Z$, and this is one of
the inputs constructed above.


\item  Conversely, the
construction proves
occurrence for every
member of the second
class. For the torus-quotient class, let $A$ be the prescribed
two-torus and let $G$ be its given finite cyclic automorphism
group. Choose a faithful character $\xi:G\to\C^*$ and take
the compact quotient of $A\times\PP^1_t$ by
$g(a,t)=(g(a),\xi(g)t)$. Over $T=t^{|G|}\ne0,\infty$,
the covering fibres are permuted simply transitively, so the
quotient fibre is this same $A$. Apply the resolution and
strict-component comparison above with
$(\widetilde Y,Q,E)=(A\times\PP^1,G,A\times\{0\})$.
The restriction of the quotient is $A/G$, and the strict
component on the ambient resolution is properly birational
to that specified quotient. Its smooth minimal model is
therefore the prescribed one. The same \cref{cdp:st:C.2.UN.4}
balancing preserves it through the displayed blowup maps.




\end{enumerate}

\end{enumerate}
\end{proof}

\begin{remark}[Characters and transfer for the elliptic torus lift]\hypertarget{tag-P346}{}\label{rem:II:elliptic-torus-characters}
\StatementRecord{rem:II:elliptic-torus-characters}{proved}{P346}{\cref{lem:II:equivariant-elliptic-torsor-lift}; \cite[Theorem~1.9, p.~51]{Steinberg1965}}
For the torus in \cref{lem:II:equivariant-elliptic-torsor-lift},
put $k=\C(B/Q)$ and let $L_\beta=\C(B)^{\ker\beta}$.
Its cyclic splitting action has order $d\in\{1,2,3,4,6\}$.
The following character maps specify its norm and restriction
comparisons. The torsor vanishing is \eqref{cdp:eq:C.EO.H1}.

For every finite extension $L'/k'$ of complex curve function
fields, the norm-one algebraic torus is a connected linear group
over the perfect field $k'$. The same Stacks dimension bound
and Steinberg theorem used in the lemma apply to that group.
The connecting map in the norm sequence therefore gives
\begin{equation}
                   N_{L'/k'}:(L')^*\longrightarrow(k')^*
                         \quad\hbox{surjective}.
 \hypertarget{tag-00HV}{}\label{cdp:eq:C.EO.H2}
\end{equation}

After an integral choice of character basis, a generator acts by
\begin{equation}
\begin{array}{c|c}
d&M_d\\\hline
1&I_2\\
2&-I_2\\
3&\begin{pmatrix}0&-1\\1&-1\end{pmatrix}\\
4&\begin{pmatrix}0&-1\\1&0\end{pmatrix}\\
6&\begin{pmatrix}0&-1\\1&1\end{pmatrix}.
\end{array}
 \hypertarget{tag-00HW}{}\label{cdp:eq:C.EO.H3}
\end{equation}
For orders $3,4,6$, the preserved rank-two lattice is a rank-one
module over $\Z[\zeta_d]$. Choosing a nonzero rational lattice
vector identifies it with a fractional ideal. The Gaussian and
Eisenstein rings are Euclidean for their norms, so this ideal is
principal. The same argument applies to the dual character
lattice and gives the displayed bases; inverse generators are
integrally conjugate. The order-two action is $-I_2$.

For $d=2$, the torus is the product of two quadratic norm-one
tori. For $d=3$, it is the cubic norm-one torus. Indeed,
\begin{equation}
 1\longrightarrow R^1_{L_\beta/k}\mathbb G_m
 \longrightarrow R_{L_\beta/k}\mathbb G_m
 \xrightarrow{N}\mathbb G_m\longrightarrow1
 \hypertarget{tag-00HX}{}\label{cdp:eq:C.EO.H4}
\end{equation}
has character quotient $\Z[C_d]/\Z(1+g+\cdots+g^{d-1})$.
For $d=3$, multiplication by $g$ in the basis $[1],[g]$ is
exactly $M_3$. For $d=2$ the quotient is the sign lattice.
These integral character isomorphisms give the stated torus
isomorphisms by Galois descent.

For $d=4$, let $k_2$ be the quadratic intermediate field.
The comparison is
\begin{equation}
           T_k\simeq R_{k_2/k}
                  \bigl(R^1_{L_\beta/k_2}\mathbb G_m\bigr).
 \hypertarget{tag-00HY}{}\label{cdp:eq:C.EO.H5}
\end{equation}
Induce the sign module of $\langle g^2\rangle$ to
$\langle g\rangle$. The basis $e,ge$ satisfies $g(e)=ge$,
$g(ge)=-e$, giving $M_4$. The associated Shapiro map identifies
degree-one cohomology with that of the norm-one torus over $k_2$;
its norm comparison is \eqref{cdp:eq:C.EO.H2}.

For $d=6$, let $k_2,k_3$ be the intermediate fields of degrees
two and three. The restricted generators act by $M_6^2$ and
$M_6^3=-I_2$. The exact integral change of basis
\begin{equation}
 P=\begin{pmatrix}1&-1\\0&1\end{pmatrix},\qquad
                         P^{-1}M_6^2P=M_3
 \hypertarget{tag-00HZ}{}\label{cdp:eq:C.EO.H6}
\end{equation}
identifies the first restriction with the cubic norm-one torus;
the second is a product of quadratic norm-one tori.

For the transfer comparison, let $k'/k$ be finite and put
$\mathcal G=\operatorname{Gal}(k^{\rm sep}/k)$,
$\mathcal H=\operatorname{Gal}(k^{\rm sep}/k')$ and
$A=T_k(k^{\rm sep})$, written additively here.
Choose left-coset representatives $R$ and write
$\sigma r=r_\sigma h_{\sigma,r}$. For a continuous one-cocycle
$c$ on $\mathcal H$, corestriction is represented by
\begin{equation}
       (\operatorname{cor}c)(\sigma)
           =\sum_{r\in R}r_\sigma c(h_{\sigma,r}).
 \hypertarget{tag-00I0}{}\label{cdp:eq:C.EO.H7}
\end{equation}
The cocycle identity verifies that this is a cocycle. For
$c=\operatorname{res}b$, the equality
$b(\sigma r)=b(\sigma)+\sigma b(r)
 =b(r_\sigma)+r_\sigma b(h_{\sigma,r})$
makes the displayed sum equal to $[k':k]b(\sigma)$ plus the
coboundary of $\sum_{r\in R}b(r)$. Thus
$\operatorname{cor}\operatorname{res}$ is multiplication by
$[k':k]$ on $H^1(k,T_k)$. For the two intermediate fields in
the order-six case these indices are two and three.
\end{remark}



\subsubsection{Hopf surfaces with elliptic reduction}\label{subsubsec:II:elliptic-hopf-occurrence}

\begin{proposition}[Occurrence of every Hopf surface of algebraic dimension one]\hypertarget{tag-01DG}{}\label{cdp:st:C.2.UN.37}
\StatementRecord{cdp:st:C.2.UN.37}{proved}{01DG}{\cref{lem:II:hopf-curve-restriction,lem:II:hopf-curves,def:II:elliptic-root-curve,lem:II:identity-monodromy-elliptic,cdp:st:C.2.UN.34}; \eqref{cdp:eq:C.EH.1}--\eqref{cdp:eq:C.EH.8}}

Let $H$ be any primary or secondary Hopf surface with $a(H)=1$.
There is a proper holomorphic map from a smooth compact complex
threefold to $\mathbb P^1$, with connected fibres and smooth
two-tori away from one value, whose fibre over that value contains
$H$ itself as a smooth component of multiplicity one. Its reduced
special fibre has simple normal crossings and Euler number zero.
Over the complement of the special value the family is the
product of two fixed elliptic curves with the base.

The map used to apply \cref{cdp:st:C.2.UN.34} is the connected elliptic
reduction
\begin{equation}
                         q:H\longrightarrow\mathbb P^1.
\hypertarget{tag-00FQ}{}\label{cdp:eq:C.EH.1}
\end{equation}
Every fibre of $q$ is a multiple of a smooth elliptic curve,
the modulus is constant, and the integral $H^1$ monodromy is
the identity. 
\end{proposition}



\begin{proof}

\begin{enumerate}[label=\arabic*.,leftmargin=*]
\item Apply \cref{lem:II:hopf-curve-restriction} to the given
primary or secondary Hopf surface $H$. Its topological conclusions give
\begin{equation}
                    b_1(H)=1,\qquad H^2(H,\mathbb Q)=0.
\hypertarget{tag-00FR}{}\label{cdp:eq:C.EH.2}
\end{equation}

\item Choose a nonconstant meromorphic function $h$ on $H$.
Its effective zero and pole divisors have no common
irreducible component and have intersection number zero
by \eqref{cdp:eq:C.EH.2}. Local intersection multiplicities of distinct
effective curves on a smooth surface are positive.
Consequently those two divisors are disjoint. In a local
coprime expression $h=f/g$, an indeterminacy point would
belong to both $(f=0)$ and $(g=0)$, so no such point exists.
The function therefore defines a holomorphic map
$h:H\to\mathbb P^1$. Its Stein factorization is
\begin{equation}
               H\xrightarrow{q}B\xrightarrow{\nu}\mathbb P^1,
\hypertarget{tag-00FS}{}\label{cdp:eq:C.EH.3}
\end{equation}
where $B$ is a smooth compact connected curve, $q$ has
connected fibres, and $\nu$ is finite. The degree-one
Leray edge map injects $H^1(B,\mathbb Q)$ into
$H^1(H,\mathbb Q)$, since $q_*\mathbb Q=\mathbb Q$.
Thus $2g(B)\leq1$, and $B\simeq\mathbb P^1$.

\item By \cref{lem:II:hopf-curves}, every irreducible compact
curve $C\subset H$ is smooth elliptic and distinct such curves
are disjoint. For this same curve, adjunction records
\begin{equation}
                 2p_a(C)-2=(K_H+C)\cdot C=0.
\hypertarget{tag-00FT}{}\label{cdp:eq:C.EH.4}
\end{equation}
Since $q$ has connected fibres, each full fibre is
\begin{equation}
                              q^{-1}(b)=m_b C_b,
\hypertarget{tag-00FU}{}\label{cdp:eq:C.EH.5}
\end{equation}
with $C_b$ smooth connected elliptic. The full fibre is this
Cartier divisor: its local ring is a two-dimensional regular
local ring modulo a nonzerodivisor, hence is Cohen--Macaulay
of dimension one and has no embedded component.

\item At a full fibre $m_bC_b$, take the same disc coordinate
$t$, let $m=m_b$, and use \cref{def:II:elliptic-root-curve}
with $D=C_b$. In its local equation $t=u h^m$, the local
defining function $h$ is the $x$ of this curve chart.
Thus its maps \eqref{cdp:eq:C.EH.6} are exactly
$w^m=u$, $s=wx$, with the original surface coordinates fixed.
Its normalization $V\to U$ is finite \'etale, has quotient
$U=V/\mu_m$, and the root family has connected fibres over
$\Delta_s$. Its central quotient
\eqref{cdp:eq:C.EH.7} has $C=C_b$ and is the same map
$C'\to C_b$. Its deck group acts by free translations.

\item Choose a $\mu_m$-invariant smooth horizontal connection
for the submersion $V\to\Delta_s$, by finite averaging.
Radial parallel transport then gives an equivariant
smooth trivialization. A positive meridian in the
$t$-disc acts on $H^1$ by the corresponding deck
translation or its inverse, according to the transport
convention; both induce the identity. We have proved
\begin{equation}
                  T_b=I\quad\hbox{on }H^1(q^{-1}(t),\mathbb Z).
\hypertarget{tag-00FX}{}\label{cdp:eq:C.EH.8}
\end{equation}
The finitely many meridians generate the fundamental
group of the punctured base $\mathbb P^1$, so the entire
integral monodromy representation is trivial.

\item Every fibre support of this $q$ is elliptic, so it contains
no rational $(-1)$-curve and $q$ is relatively minimal.
The connected-fibre map in \eqref{cdp:eq:C.EH.3}, the full
fibres in \eqref{cdp:eq:C.EH.5}, and the preceding identity
monodromy verify all hypotheses of
\cref{lem:II:identity-monodromy-elliptic} with
$(E_0,B,q)=(H,\PP^1,q)$. Its fixed elliptic curve is therefore
the Jacobian for this same map \eqref{cdp:eq:C.EH.1}.
Apply \cref{cdp:st:C.2.UN.34} to $S=H$ and this $q$.
Its root-curve torsor is the one identified by
\eqref{cdp:eq:C.EH.6}--\eqref{cdp:eq:C.EH.7}; the two-line
lift and its analytic torsor isomorphism preserve this class.
Hence the specified component is isomorphic to $H$, with
multiplicity one. The vertical line restrictions,
pushout and free root charts in
\eqref{cdp:eq:C.SF.8}--\eqref{cdp:eq:C.SF.12} give the same
proper compact map and its zero-Euler SNC support, and
\eqref{cdp:eq:C.SF.13} is its full punctured product.

\end{enumerate}
\end{proof}






\Needspace{10\baselineskip}
\subsection{Hopf axes and canonical-order ratios}\label{subsec:II:hopf-axis-adjacency}


\begin{proposition}[Hopf-axis adjacency with arbitrary multiplicities and point blowups]\hypertarget{tag-01DK}{}\label{cdp:st:C.2.UN.40}
\StatementRecord{cdp:st:C.2.UN.40}{proved}{01DK}{\cref{cdp:st:C.2.UN.17,lem:II:central-primary-hopf-cover}; \cite[arXiv v2, normal forms (1), p.~3; Theorem~2.1, p.~5]{Khaled2024Hopf}; \cite[arXiv v1, \S2.2.2, p.~4; \S2.2.6, p.~5]{AbashevaDeev2024AB}}

Let $f:U\to\Delta$ be a proper degeneration of smooth complex
two-tori with smooth total space, and suppose that its reduced
central support has simple normal crossings. Write its full
fibre as $\sum_i m_iE_i$. Choose a generator of the rank-one
sheaf $f_*K_U$ and write its divisor as $\sum_i k_iE_i$.
Set $r_i=(k_i+1)/m_i$.

Suppose that the minimal model $H_i$ of a component $E_i$
is a Hopf surface of algebraic dimension zero. The contraction
$\rho_i:E_i\to H_i$ may consist of arbitrarily many point
blowdowns. Choose a finite unramified primary Hopf
cover $p:P\to H_i$.

If $P$ has diagonal contraction $(z,w)\mapsto(\alpha z,\beta w)$,
the images of both coordinate axes are boundary curves on
$H_i$: their strict transforms in $E_i$ meet another central
component. For each such strict transform, let $E_j$ be its
unique adjacent component at the generic point. Then
\begin{equation}
                 m_i(k_j+1)-m_j(k_i+1)=0,
                 \qquad r_j=r_i.
\hypertarget{tag-00L9}{}\label{cdp:eq:C.HA.1}
\end{equation}
If $P$ has non-diagonal contraction
$(z,w)\mapsto(\beta^r z+w^r,\beta w)$, $r\geq1$, the image
of its unique elliptic curve is a boundary curve. Its unique
generic neighbour satisfies
\begin{equation}
       m_i(k_j+1)-m_j(k_i+1)=-m_i r,
       \qquad m_j(r_j-r_i)=-r<0.
\hypertarget{tag-00LA}{}\label{cdp:eq:C.HA.2}
\end{equation}
These are integral divisor-class consequences on the components. In particular a non-diagonal algebraic-dimension-zero
Hopf component cannot meet another algebraic-dimension-zero
Hopf component along this elliptic curve. Diagonal such
components can meet along their axes only at equal ratios.
\end{proposition}


\begin{proof}

\begin{enumerate}[label=\arabic*.,leftmargin=*]
\item Apply \cref{cdp:st:C.2.UN.17} to the same $f:U\to\Delta$ and full fibre $\sum_i m_iE_i$. Its canonical evaluation is the generator chosen in the statement, with $K_U\simeq\OO_U(\sum_i k_iE_i)$. Each smooth $E_i$ has identity normalization and resolution and intrinsic conductor divisor zero. The principal-fibre and adjunction maps \eqref{cdp:eq:C.CA.3}--\eqref{cdp:eq:C.CA.4} therefore apply directly to $E_i$.
For $j\ne i$ put $D_{ij}=E_j|_{E_i}$ and
$D_i=\sum_{j\ne i}D_{ij}$. Each is the reduced intersection
divisor specified by the simple-normal-crossing coordinates.
The full fibre is principal. Adjunction and its restriction
to $E_i$ therefore give the following integral identity:
\begin{equation}
 m_i(K_{E_i}+D_i)
 \sim\sum_{j\ne i}
       \bigl[m_i(k_j+1)-m_j(k_i+1)\bigr]D_{ij}.
\hypertarget{tag-00LB}{}\label{cdp:eq:C.HA.3}
\end{equation}
Indeed, first use
$K_{E_i}=((k_i+1)E_i+\sum_{j\ne i}k_jE_j)|_{E_i}$,
then replace $m_i E_i|_{E_i}$ by
$-\sum_{j\ne i}m_jD_{ij}$, and finally add $m_iD_i$.


\item Push the divisor class forward through the point blowdowns.
For a proper birational map of smooth surfaces the common
meromorphic function field sends the divisor of a function
to the divisor of the same function. Sending every prime
curve dominating a curve to that curve and every exceptional
prime curve to zero consequently induces a homomorphism
of divisor-class groups. It sends $K_{E_i}$ to $K_{H_i}$:
pull back a meromorphic two-form on $H_i$ and observe that
its additional Jacobian divisor is exceptional. Applying it to \eqref{cdp:eq:C.HA.3} gives
\begin{equation}
 m_i(K_{H_i}+B_i)
 \sim\sum_{j\ne i}
       \bigl[m_i(k_j+1)-m_j(k_i+1)\bigr]\rho_{i*}D_{ij},
 \qquad B_i=\rho_{i*}D_i.
\hypertarget{tag-00LC}{}\label{cdp:eq:C.HA.4}
\end{equation}
The divisor $B_i$ is reduced. A point-blowup morphism is an
isomorphism at the generic point of every curve on its target,
so there is exactly one strict transform dominating that curve.
At its generic point the SNC equation has exactly two branches,
$x_i x_j=0$; an additional branch can meet it only at isolated
triple points. Hence this strict transform occurs in exactly
one $D_{ij}$, with coefficient one.

\item The pullback of \eqref{cdp:eq:C.HA.4} through the specified unramified
cover $p:P\to H_i$ is an identity of integral divisor classes,
and $p^*K_{H_i}=K_P$. The primary normal forms are \cite[arXiv v2, normal forms (1), p.~3; Theorem~2.1, p.~5]{Khaled2024Hopf}. Apply \cref{lem:II:central-primary-hopf-cover} to this given finite connected unramified cover $p:P\to H_i$. Its field inclusion is the pullback by $p$, and gives $a(P)=a(H_i)=0$.

For a diagonal contraction, a relation $\alpha^a\beta^b=1$ for a nonzero integer pair would make $z^aw^b$ a nonconstant meromorphic function on $P$. Thus no such relation holds. The primary curve classification \cite[arXiv v1, \S2.2.6, p.~5]{AbashevaDeev2024AB} then gives precisely the two axis curves $C_1,C_2$. For the nonlinear contraction, \cite[arXiv v1, \S2.2.2, p.~4]{AbashevaDeev2024AB} gives precisely $C=(w=0)$, including resonance $r=1$.

\item Denote by $L_\lambda$ the line bundle
whose lifted fibre coordinate is multiplied by $\lambda$
under the contraction. The map
$\lambda\mapsto L_\lambda$ is injective. A trivialization
would be a nowhere-zero holomorphic eigenfunction on
$\mathbb C^2\setminus\{0\}$. Hartogs extends it and its
inverse across zero, and evaluation there forces
$\lambda=1$. Coordinate functions and the Jacobian
therefore give the exact classes
\begin{equation}
\begin{array}{ll}
 \mathcal O(C_1)=L_\alpha,
 \quad\mathcal O(C_2)=L_\beta,
 &K_P=\mathcal O(-C_1-C_2),\\[2pt]
 \mathcal O(C)=L_\beta,
 &K_P=\mathcal O(-(r+1)C)
 \quad\hbox{in the non-diagonal case}.
\end{array}
\hypertarget{tag-00LD}{}\label{cdp:eq:C.HA.5}
\end{equation}
The first two classes are integrally independent: a relation,
or even a torsion relation, would give the forbidden
positive-power resonance after raising to a power.
The class of $C$ in the second line has infinite order
because $0<|\beta|<1$.

Every curve on $H_i$ has nonempty reduced inverse image
on $P$, since $p$ is unramified and surjective. Thus the
pullback of $B_i$ is a reduced subset of the indicated
axis curves. In the diagonal case, compare coefficients
of their independent classes in the pullback of \eqref{cdp:eq:C.HA.4}.
If an axis were absent, its coefficient on the left
would be $-m_i$ and on the right zero. Each axis is
therefore present. Its coefficient on the left is then
zero, and the coefficient on the right is exactly that
of its unique generic neighbour. This proves \eqref{cdp:eq:C.HA.1},
including when one boundary curve has more than one
component in its finite inverse image.

In the non-diagonal case, absence of $C$ would instead
give the impossible nonzero class $-m_i(r+1)C\sim0$.
It is present, with coefficient one. Comparing its
infinite-order class gives
\begin{equation}
       -m_i r=m_i(k_j+1)-m_j(k_i+1),
\hypertarget{tag-00LE}{}\label{cdp:eq:C.HA.6}
\end{equation}
which is \eqref{cdp:eq:C.HA.2}. The divisor pushforward and the pullback by $p$ prove both assertions on the original component after all its point blowups.

Finally, the shared strict axis is a smooth elliptic
curve. It cannot be contracted by a point-blowdown
morphism on the neighbouring smooth component, since
all exceptional curves of such a morphism are rational.
If that neighbour also has an algebraic-dimension-zero
Hopf minimal model, this same curve is one of its
minimal axes. Its diagonal rule would force equal
ratios, contradicting \eqref{cdp:eq:C.HA.2}; its non-diagonal rule
would force the opposite strict inequality. Neither
is possible. The diagonal--diagonal case merely repeats
the equality \eqref{cdp:eq:C.HA.1}. 

\end{enumerate}
\end{proof}

\subsection{Diagonal Hopf components in compact torus fillings}\label{subsec:II:diagonal-hopf-occurrence}

\begin{proposition}[Diagonal Hopf components in compact torus fillings]\hypertarget{tag-01DL}{}\label{cdp:st:C.2.UN.41}
\StatementRecord{cdp:st:C.2.UN.41}{proved}{01DL}{\SeriesRef{cdp:eq:C.DH.5}, \SeriesRef{cdp:eq:C.DH.6}, \SeriesRef{cdp:eq:C.DH.7}, \SeriesRef{cdp:eq:C.DH.13}, \SeriesRef{cdp:eq:C.DH.14b}, \SeriesRef{cdp:eq:C.DH.12}; \cite[arXiv v3, Definition~3.1, p.~7; \S7.6, pp.~33--34]{SchuttShioda2010}; \cite[Theorem~2.1, p.~2]{Papikian2004Tate}; \cref{lem:II:hopf-curve-restriction,def:II:marked-normal-cycle,lem:II:moving-marked-cycle}; \cite[arXiv v1, \S2.2.6, p.~5]{AbashevaDeev2024AB}}

For every pair $\alpha,\beta\in\C$ with
$0<|\alpha|,|\beta|<1$, the diagonal Hopf surface
\begin{equation}
 H_{\alpha,\beta}=(\C^2\setminus\{0\})/
       \langle(z,w)\mapsto(\alpha z,\beta w)\rangle
                                                               \hypertarget{tag-00F6}{}\label{cdp:eq:C.DH.1}
\end{equation}
occurs as an unchanged smooth component of a
reduced normal-crossing fibre of a proper map
from a smooth compact complex threefold to
$\PP^1$, with general fibre a two-torus.
The distinguished fibre consists of two copies
of \eqref{cdp:eq:C.DH.1}, joined along both axis elliptic
curves by inversion. Every reduced fibre
support in this compact construction has
Euler number zero. This includes all
multiplicatively independent pairs, whose
Hopf surfaces have algebraic dimension zero.

The distinguished support $S$ has torsion-free
integral cohomology of ranks
\begin{equation}
                  (1,2,3,4,2),\qquad
                  h^\bullet(S,\OO_S)=(1,2,1).        \hypertarget{tag-00FG}{}\label{cdp:eq:C.DH.2}
\end{equation}
Each normalization component has integral
cohomology of ranks $(1,1,0,1,1)$ and
$h^\bullet(\OO)=(1,1,0)$.
Locally the nearby torus has the exact
multiplicative period presentation
\begin{equation}
 A_t=(\C^*)^2/
        \langle(\alpha,\beta),(t,t^{-1})\rangle
                    \qquad(0<|t|\ll1).             \hypertarget{tag-00FH}{}\label{cdp:eq:C.DH.3}
\end{equation}
In the compact construction the constants
$\alpha,\beta$ in this formula vary
holomorphically with unchanged limiting
values, and $t$ is replaced by a holomorphic
parameter with a simple zero. The punctured families are identified by \eqref{cdp:eq:C.DH.15}.
\end{proposition}


\begin{proof}\leavevmode\par

\begin{enumerate}[label=\arabic*.,leftmargin=0pt,labelsep=.5em]
\item \par\medskip\noindent\textbf{Normal bundles and plumbing maps.}
The two axes of $H=H_{\alpha,\beta}$ are
$C_\alpha=\C^*/\alpha^{\Z}$ and
$C_\beta=\C^*/\beta^{\Z}$. The open
neighborhood $z\ne0$ of the first is
the entire normal line bundle
\begin{equation}
 L_\alpha=(\C^*\times\C)/
         ((z,x)\sim(\alpha z,\beta x));
                                                               \hypertarget{tag-00FI}{}\label{cdp:eq:C.DH.4}
\end{equation}
the second has the analogous presentation
with $\alpha,\beta$ interchanged.
Inversion on $C_\alpha$ pulls this
normal line to its inverse. Explicitly
the rank-two plumbing bundle over that
elliptic curve is
\begin{equation}
 \mathcal N_\alpha=(\C^*\times\C^2)/
       ((z,x,y)\sim(\alpha z,\beta x,\beta^{-1}y)),
                         \qquad t=xy.               \hypertarget{tag-00FJ}{}\label{cdp:eq:C.DH.5}
\end{equation}
Its branch $y=0$ maps to the first
Hopf copy by $(z,x)$, and its branch
$x=0$ maps to the second by
$(Z,W)=(z^{-1},y)$.
For the other axis use
\begin{equation}
 \mathcal N_\beta=(\C^*\times\C^2)/
       ((w,x,y)\sim(\beta w,\alpha x,\alpha^{-1}y)),
                         \qquad t=xy,                \hypertarget{tag-00FK}{}\label{cdp:eq:C.DH.6}
\end{equation}
whose branches map to the first
copy by $(z,w)=(x,w)$ and to the
second by $(Z,W)=(y,w^{-1})$.
These formulas prove the two dual-normal
identifications, including their
base and fibre transition functions.


\begin{enumerate}[label=(\arabic*),leftmargin=1.3em,labelsep=.5em]
\item  Put $A=-\log|\alpha|>0$ and
$B=-\log|\beta|>0$. On the common
open torus of a Hopf copy define
\begin{equation}
 h=B\log|z|-A\log|w|,\qquad
 \nu_\alpha=|w||z|^{-B/A}=e^{-h/A},\qquad
 \nu_\beta=|z||w|^{-A/B}=e^{h/B}.        \hypertarget{tag-00FL}{}\label{cdp:eq:C.DH.7}
\end{equation}
All three expressions are invariant
under the Hopf deck generator.
The two $\nu$ are the norms of
flat Hermitian metrics on the
normal lines. Neighborhoods
$\nu_\alpha<R$ and $\nu_\beta<R$
are disjoint for $R<1$.
Fix $0<r<R<1$ and $0<\delta<r^2$.
For each copy take the product
of the relatively compact open core
$U_i=\{\nu_\alpha>r,\nu_\beta>r\}$
with $\Delta_\delta$.
In \eqref{cdp:eq:C.DH.5} take
\begin{equation}
 \|x\|=|x||z|^{-B/A}<R,
 \qquad \|y\|=|y||z|^{B/A}<R,
 \qquad |xy|<\delta;                                \hypertarget{tag-00FM}{}\label{cdp:eq:C.DH.8}
\end{equation}
take the analogous open set in
\eqref{cdp:eq:C.DH.6}. Glue it to the first
core where $\|x\|>r$, using the
first branch coordinate and
$t=xy$, and to the second core
where $\|y\|>r$, using the second
branch coordinate. The product
of the two norms is exactly
$|t|$. Thus the two core
overlaps in any one neck are
disjoint, since $|t|<r^2$.
The coordinate maps are
biholomorphisms on the stated
overlaps; their inverses insert
the coordinate $t/x$ or $t/y$.


\item These are the data of \cref{def:II:marked-normal-cycle}
with two copies of $H$, the two axis identifications by inversion,
and the dual-normal maps in \eqref{cdp:eq:C.DH.5} and
\eqref{cdp:eq:C.DH.6}. The flat metrics are exactly the norms
in \eqref{cdp:eq:C.DH.7}--\eqref{cdp:eq:C.DH.8}; on every
neck their product is $|t|$. Apply the fixed-data construction
of \cref{lem:II:moving-marked-cycle} with these same $r,R,\delta$,
cores and inverse maps $t/x,t/y$. Its finite gluing is therefore
this gluing, with the identity on each core and each neck chart.
The lemma gives a Hausdorff smooth threefold $\mathcal Y$, proper
over $\Delta_\delta$, and full central fibre the two copies of
$H$ joined by the indicated two inversions. At the double curves
its holomorphic function is $xy=t$ and both multiplicities are one.


\end{enumerate}

\item \par\medskip\noindent\textbf{The nearby torus.}
For $t\ne0$, the punctured
identifications are
\begin{equation}
 \phi_\alpha(z,w)=(z^{-1},t/w),\qquad
 \phi_\beta(z,w)=(t/z,w^{-1}).                       \hypertarget{tag-00FN}{}\label{cdp:eq:C.DH.9}
\end{equation}
Each conjugates the Hopf deck
generator to its inverse.
Their difference is the
additional period
$(z,w)\mapsto(tz,t^{-1}w)$.
The four additive period
vectors of \eqref{cdp:eq:C.DH.3} are
\begin{equation}
 (2\pi i,0),\ (0,2\pi i),\
 (\log\alpha,\log\beta),\
                         (\log t,-\log t).          \hypertarget{tag-00F7}{}\label{cdp:eq:C.DH.10}
\end{equation}
They are real-linearly
independent: after quotienting
the two imaginary directions,
the real determinant is
\begin{equation}
       -\log|t|\bigl(\log|\alpha|+\log|\beta|\bigr)
                                      \ne0.          \hypertarget{tag-00F8}{}\label{cdp:eq:C.DH.11}
\end{equation}
Their integral span is
therefore a discrete
cocompact lattice in $\C^2$.
This proves directly that
$A_t$ is a compact complex
two-torus. A turn of $t$
adds $(2\pi i,-2\pi i)$
to the last vector, which
is already in the lattice,
so these quotients form
the stated family over
the punctured disc.


\begin{enumerate}[label=(\arabic*),leftmargin=1.3em,labelsep=.5em]
\item  There is a holomorphic comparison
$\Psi_t:\mathcal Y_t\to A_t$.
On the first core use
$[z,w]$, and on the
second use $[Z^{-1},t/W]$.
On the first neck use
$[z,x]$; on the second
use $[x,w]$. On its
second-core overlap
the latter gives instead
$[x/t,tw]$, differing by
the inverse of the period
$(t,t^{-1})$. Thus the
maps agree exactly on
all the overlaps. They
are local biholomorphisms.
The proper source makes
$\Psi_t$ a finite
unramified covering
onto the connected
torus $A_t$.


\item Order the logarithmic coordinates on the first copy as
$(\log z,\log w)$ and on the second as $(\log W,\log Z)$.
Following first the $\beta$-edge and then the $\alpha$-edge,
both linear edge maps are
$L=\left(\begin{smallmatrix}0&-1\\-1&0\end{smallmatrix}\right)$.
Their product is the identity, and their affine composite is
$(\log z,\log w)\mapsto(\log z-\log t,\log w+\log t)$.
Thus \cref{lem:II:moving-marked-cycle} identifies the nearby
fibre with the quotient by the lattice in
\eqref{cdp:eq:C.DH.10}: its last generator is the negative of
the displayed one and generates the same integral subgroup.
Its common-coordinate map on the second core is
$[t/Z,1/W]$, whereas $\Psi_t$ uses $[1/Z,t/W]$.
Their ratio is the lattice period $(t,t^{-1})$, so they
give the same quotient map. On the first core they are
identical, and the neck maps are their restrictions through
the same inverses $t/x,t/y$. Consequently $\Psi_t$ is the
coframe isomorphism of that lemma, including its inverse.


\item  For holomorphic $a(t),b(t)$ with limits $\alpha,\beta$,
choose logarithms on a small disc and move the lattice basis to
$(2\pi i,0)$, $(0,2\pi i)$, $(\log a(t),\log b(t))$.
The two angular meridians keep their integer coordinates;
the ordered second-component frame is the same $L$ as above.
All central dual-normal constants are one. The moving part of
\cref{lem:II:moving-marked-cycle} therefore gives these marked
Hopf components and their two inverse-axis gluing maps.
For a smoothing parameter $Q(t)$ with $Q(0)=0$ and
$Q'(0)\ne0$, let $t=t(\tau)$ be the local holomorphic inverse
of $\tau=Q(t)$. Apply that construction over $\tau$ to
$a(t(\tau)),b(t(\tau))$ and pull it back by $\tau=Q(t)$.
The two inverse base maps identify its neck equations with
$xy=Q(t)$; the resulting total space is smooth and proper.
Its cycle period is $(-\log Q(t),\log Q(t))$, so changing
the sign of that integral generator gives the family isomorphism
\begin{equation}
 \mathcal Y_t\simeq
 (\C^*)^2/
       \langle(a(t),b(t)),(Q(t),Q(t)^{-1})\rangle.
                                                               \hypertarget{tag-00F9}{}\label{cdp:eq:C.DH.12}
\end{equation}


\end{enumerate}

\item \par\medskip\noindent\textbf{A compact suspension.}
We construct a compact elliptic
surface $J\to\PP^1$ with a
section $P$ whose value on a
nodal fibre is the specified
$\alpha$. Start with the
nodal plane cubic
$y^2z=x^2(x+z)$, with origin
$O=[0:1:0]$. On its
normalization put
$v=y/x$; its smooth group
coordinate is
$\zeta=(v-1)/(v+1)$,
with $O$ corresponding to
$\zeta=1$. Choose the
smooth point $P_0$ with
$\zeta=\alpha$.
Take a second cubic
through $O,P_0$, otherwise
general. Its intersection
with the nodal cubic
consists of nine distinct
smooth points, with
transverse intersections,
and misses the node.
These requirements are
nonempty open conditions
on its restriction to
the cubic and its first
jets; the restriction
map of cubic polynomials
is surjective, so a
general such choice
satisfies them.
Blow up the nine
base points of the
pencil. The resulting
smooth rational surface
$J$ has a connected
elliptic pencil, zero
section from $O$, and
section $P$ from $P_0$.
Its central member is
the original nodal
cubic, a fibre of type
$I_1$. At the node the
second cubic is nonzero,
so the pencil parameter
smooths the node to
first order. The full
fibre class is
$F=3H-\sum E_i=-K_J$.
Consequently there is
no vertical rational
$(-1)$-curve, by
adjunction, and this
is the smooth relatively
minimal elliptic model.


\begin{enumerate}[label=(\arabic*),leftmargin=1.3em,labelsep=.5em]
\item The surface $J$ is smooth, projective and relatively minimal over
$\PP^1$, with the section $P$ constructed above. Apply
\cite[arXiv v3, Definition~3.1, p.~7; \S7.6, pp.~33--34]{SchuttShioda2010}
to this $J\to\PP^1$ and this section. Translation on each smooth
fibre extends to an automorphism $T_P:J\to J$ satisfying
$f_J\circ T_P=f_J$; its inverse is translation by $-P$.


\item  Put $q_0=\alpha\beta$.
The quotient
\begin{equation}
 Y^{\rm std}=(\C^*_u\times J)/
        \langle(u,y)\mapsto(q_0u,T_P(y))\rangle
                                                               \hypertarget{tag-00FA}{}\label{cdp:eq:C.DH.13}
\end{equation}
is a smooth compact
threefold. Projection to
$E_{q_0}=\C^*/q_0^{\Z}$
is a locally trivial
holomorphic bundle with
compact fibre $J$:
local choices of a
logarithm give its
trivializations, and
their transition is
the indicated integral
power of $T_P$.
The deck action is
free and properly
discontinuous because
its first coordinate
is multiplication by
$q_0$, with $|q_0|<1$.
It preserves the
elliptic-pencil parameter,
so \eqref{cdp:eq:C.DH.13} maps
properly to $\PP^1$.


\item  On a small punctured
disc at its $I_1$ value
there is the analytic Tate presentation
\begin{equation}
 J_t\simeq\C^*/Q(t)^{\Z},\qquad
 Q(t)=t\,u(t),\quad u(0)\ne0,
 \qquad P(t)=[a(t)],\quad a(0)=\alpha.
                                                               \hypertarget{tag-00FB}{}\label{cdp:eq:C.DH.14}
\end{equation}
We identify the central group coordinate and the family parameter using the explicit uniformization.
The complex uniformization in \cite[Theorem~2.1 and equation~(2.1), p.~2]{Papikian2004Tate} gives
\begin{equation}
 \begin{gathered}
 y^2+xy=x^3+a_4(q)x+a_6(q),\\
 a_4(q)=-5q+O(q^2),\quad
 a_6(q)=-q+O(q^2),\quad
 \Delta(q)=q+O(q^2).
 \end{gathered}                                 \hypertarget{tag-00FC}{}\label{cdp:eq:C.DH.14a}
\end{equation}
Here $a_4=-5\sum_{n\ge1}n^3q^n/(1-q^n)$ and
$a_6=-\sum_{n\ge1}(5n^3+7n^5)q^n/(12(1-q^n))$.
The multiplicative
uniformization has the
central values
\begin{equation}
 x=\frac{u}{(1-u)^2},\qquad
 y=\frac{u^2}{(1-u)^3},\qquad
 X=4x,\quad Y=8y+4x,\qquad
 \frac{Y/X-1}{Y/X+1}=u.                          \hypertarget{tag-00FD}{}\label{cdp:eq:C.DH.14b}
\end{equation}
Its equation becomes
$Y^2=X^2(X+1)$, precisely
the chosen nodal cubic.
The series converge
normally on compact
subsets of the smooth
central chart. Invert
the holomorphic series
$1/j(q)=q+O(q^2)$ at
the given simple zero
of $1/j(t)$ to obtain
$Q(t)$. After completing
the square and cube,
the two short Weierstrass
coefficients are units
at zero. Their equal
$j$-invariants identify
them by a holomorphic
unit fourth root; choose
that root at zero to
match the sixth-weight
coefficient. The identity
of the squared ratios
then matches that
coefficient throughout
the disc. This gives
the Weierstrass
coordinate isomorphism.
Choose its central value
as in \eqref{cdp:eq:C.DH.14b}.
The inverse function
theorem at the smooth
point $u=\alpha$ gives
the holomorphic lift
$a(t)$ with exactly
$a(0)=\alpha$.


\item  Let $b(t)=q_0/a(t)$.
It tends to $\beta$,
and both multipliers
have modulus less than
one on a sufficiently
small disc. The coordinate change
\begin{equation}
                        z=y,\qquad w=u/y
                                                               \hypertarget{tag-00FE}{}\label{cdp:eq:C.DH.15}
\end{equation}
on the multiplicative
cover of \eqref{cdp:eq:C.DH.13}
turns its two deck
generators into
$(a(t),b(t))$ and
$(Q(t),Q(t)^{-1})$.
Thus its punctured
family is exactly
\eqref{cdp:eq:C.DH.12}. Replace
the local filling
of \eqref{cdp:eq:C.DH.13} by
the proper plumbing
just constructed,
using that explicit
isomorphism on the
punctured disc.
This is gluing over
the open cover
$\PP^1\setminus\{0\}$
and a disc at zero.
The total spaces
are identified by
biholomorphisms over
their common base,
so the result is
Hausdorff and smooth;
properness is local
on the base and
holds on each member
of this cover.
The compact base
therefore gives a
compact total space.
The distinguished
central components
are exactly the
prescribed Hopf
surfaces. General
fibres are the
tori of
\eqref{cdp:eq:C.DH.12}, or,
elsewhere, the same
translation suspensions
of smooth elliptic
fibres of $J$.
The latter are tori
by their additive
four-period lattice,
or by the direct
constant-translation
bundle presentation
in \eqref{cdp:eq:C.DH.13}.


\item  The support over zero
has Euler number
$2e(H)-2e(C)=0$.
At every other value
the support is a
locally trivial bundle
over $E_{q_0}$ with
fibre the corresponding
reduced elliptic-pencil
support; its monodromy
is $T_P$. Its Euler
number is the product
with $e(E_{q_0})=0$.
Thus every reduced
support has Euler
number zero.


\end{enumerate}

\item \par\medskip\noindent\textbf{The support cohomology and its maps.}
For the two-component
support there is the
normalization
Mayer--Vietoris sequence
with middle map the
difference of the two
restrictions to
$C_\alpha\sqcup C_\beta$.
In degree zero its
matrix is
$\left(\begin{smallmatrix}1&-1\\1&-1\end{smallmatrix}\right)$,
of rank one and
primitive image.
Choose on each axis
the circle of its
Hopf deck period and
its angular circle.
By \cref{lem:II:hopf-curve-restriction}, each primary component has the integral groups stated above.
Its integral $H^1$ generator restricts
as $(1,0)$ on each
axis. The inversion
gluing reverses these
two generators on
the second side.
Consequently the
degree-one restriction
matrix is
\begin{equation}
                     \begin{pmatrix}
                       1&1\\0&0\\1&1\\0&0
                     \end{pmatrix},                \hypertarget{tag-00FF}{}\label{cdp:eq:C.DH.16}
\end{equation}
again of rank one
with primitive image.
The Hopf components
have $H^2=0$, and
the two elliptic
curves have $H^2=\Z^2$
and no higher
cohomology. Exactness
therefore gives
\eqref{cdp:eq:C.DH.2}, including
torsion-freeness.
For coherent cohomology apply \cref{lem:II:hopf-curve-restriction}
to each of these two primary components and their axis inclusions.
Its basis $\eta_H$, the cocycle with value one on the contracting
deck generator, restricts to the same nonzero class on either axis.
The gluing is inversion, so its pullback on the second restriction
changes that class to its negative. Consequently both rows of
the coherent degree-one difference map are $(1,1)$.
The same lemma gives $h^\bullet(\OO_H)=(1,1,0)$, and the
degree-zero map is the matrix already displayed. The coherent
normalization sequence therefore gives $(1,2,1)$.
All maps in these
calculations are
the stated axis
restrictions and
inversion pullbacks.
For multiplicatively independent multipliers, apply the curve description
\cite[arXiv v1, \S2.2.6, p.~5]{AbashevaDeev2024AB} with
$(\lambda,\mu)=(\alpha,\beta)$. Its two curves are the images
of the same coordinate axes $C_\alpha,C_\beta$ used in the plumbing. A nonconstant meromorphic
function would have infinitely many distinct nonempty generic level divisors,
which is impossible when the only compact curves are those two axes.
Consequently the Hopf component in this case has algebraic dimension zero.


\end{enumerate}
\end{proof}


\Needspace{10\baselineskip}
\subsection{Nonlinear Hopf components over a rational base}\label{subsec:II:nonlinear-hopf-occurrence}

\begin{proposition}[Nonlinear Hopf components over a rational base]\hypertarget{tag-01DM}{}\label{cdp:st:C.2.UN.42}
\StatementRecord{cdp:st:C.2.UN.42}{proved}{01DM}{\SeriesRef{cdp:eq:C.NH.6}, \SeriesRef{cdp:eq:C.NH.5}, \SeriesRef{cdp:eq:C.NH.4}, \SeriesRef{cdp:eq:C.NH.10}, \SeriesRef{cdp:eq:C.NH.12}, \SeriesRef{cdp:eq:C.NH.11}, \SeriesRef{cdp:eq:C.NH.8}, \SeriesRef{cdp:st:C.2.UN.40}, \SeriesRef{cdp:eq:C.NH.14}, \SeriesRef{cdp:eq:C.NH.13}; \cref{lem:II:elliptic-root-surface,lem:II:hopf-curve-restriction}}

For every integer $r\ge1$ and every $\beta$ with $0<|\beta|<1$,
the primary Hopf surface
\begin{equation}
 H_r(\beta)=(\C^2\setminus\{0\})/
   \langle(z,w)\mapsto(\beta^r z+w^r,\beta w)\rangle
                                                        \hypertarget{tag-00S9}{}\label{cdp:eq:C.NH.1}
\end{equation}
occurs unchanged as a smooth multiplicity-one component of a
reduced SNC fibre of a proper map $f:Y\to\PP^1$, where $Y$ is
a smooth compact complex threefold and the general fibre is a
complex two-torus. Every reduced fibre support in this
construction has Euler number zero.

The distinguished full fibre is $S=H_r(\beta)+P_r$, with one
smooth elliptic intersection $J$. The second component is a
smooth elliptic surface over an elliptic curve, with exactly
one multiple fibre $rJ$ when $r>1$, and
\begin{equation}
 \begin{gathered}
 K_{P_r}=\OO_{P_r}((r-1)J),\qquad
 b(P_r)=(1,3,4,3,1),\qquad h^\bullet(\OO_{P_r})=(1,2,1),\\
 H^\bullet(S,\Z)\text{ is torsion-free of ranks }(1,3,5,5,2),
 \qquad h^\bullet(S,\OO_S)=(1,2,1).
 \end{gathered}                                         \hypertarget{tag-00SJ}{}\label{cdp:eq:C.NH.2}
\end{equation}
The Hopf component has integral ranks $(1,1,0,1,1)$ and
coherent dimensions $(1,1,0)$. 
\end{proposition}


\begin{proof}

\begin{enumerate}[label=\arabic*.,leftmargin=*]
\item \par\medskip\noindent\textbf{A compact base with the required translation jet.}
Let $B$ be the elliptic cubic $v^2=u^3-u+1$, with origin the
point at infinity. Its cubic has distinct nonzero roots. Let
$C$ be the smooth projective completion of
\begin{equation}
 v^2=t^{6r}-t^{2r}+1,
 \qquad \psi:C\longrightarrow B,\quad
          (t,v)\longmapsto(t^{2r},v).
                                                        \hypertarget{tag-00SK}{}\label{cdp:eq:C.NH.3}
\end{equation}
The polynomial has distinct roots: a repeated root would be
nonzero and would give a common root of $u^3-u+1$ and its
derivative. At infinity the two charts
$s=t^{-1}$, $v s^{3r}=\pm\sqrt{1-s^{4r}+s^{6r}}$ give two
distinct smooth points $c_+,c_-$. Thus $C$ is a smooth compact
curve of genus $3r-1$. The involution $h_C(t,v)=(t,-v)$ has
quotient $\PP^1_t$, interchanges $c_+,c_-$, and satisfies
$\psi h_C=-\psi$. Its fixed points are the $6r$ finite roots
of the displayed polynomial.

\item The elliptic uniformizer $u/v$ has order $r$ at both
$c_\pm$. Choose a local additive coordinate $p$ on $B$ at
zero. At each of the two points a holomorphic coordinate
$s$ on $C$ can therefore be chosen so that
\begin{equation}
 2\psi(s)=s^r\quad\hbox{in the additive coordinate on }B.
                                                        \hypertarget{tag-00SL}{}\label{cdp:eq:C.NH.4}
\end{equation}
Taking the $r$-th root of its nonzero holomorphic unit proves
this coordinate equality. On the compact surface $B\times C$
consider the commuting automorphisms
\begin{equation}
 A(p,c)=(p+2\psi(c),c),\qquad
 A_{1/2}(p,c)=(p+\psi(c),c),\qquad
 h(p,c)=(-p,h_Cc).
                                                        \hypertarget{tag-00SM}{}\label{cdp:eq:C.NH.5}
\end{equation}
Here $A_{1/2}^2=A$, $h^2=1$, and the equality
$\psi h_C=-\psi$ proves that $h$ commutes with both
translations. All three preserve the pair of points
$(0,c_+),(0,c_-)$.

\par\medskip\noindent\textbf{The weighted blowup and its nonzero line space.}
Perform the weighted blowup of weights $(r,1)$ at both
points, using its smooth orbifold charts. Denote the
resulting compact orbifold surface by $\mathcal V$, its
exceptional divisors by $R_+,R_-$, and put
$L=\OO_{\mathcal V}(R_++R_-)$. Locally its weighted blowup
has the explicit presentation
\begin{equation}
 \left[((\C^2_{z,w}\setminus\{0\})\times\C_x)/\C^*\right],
 \quad \lambda(z,w,x)=(\lambda^r z,\lambda w,\lambda^{-1}x),
 \quad p=x^r z,\quad s=xw.
                                                        \hypertarget{tag-00SN}{}\label{cdp:eq:C.NH.6}
\end{equation}
The exceptional divisor is $x=0$. The line $L$ has fibre
coordinate $\ell$ of weight $-1$ and canonical section
$\ell=x$. The quotient chart $w\ne0$ is the smooth chart
$p=s^r(z/w^r)$; the chart $z\ne0$ has the residual finite
action
\begin{equation}
 (w,x,\ell)\longmapsto
      (\zeta w,\zeta^{-1}x,\zeta^{-1}\ell),\qquad \zeta^r=1.
                                                        \hypertarget{tag-00SO}{}\label{cdp:eq:C.NH.7}
\end{equation}
It is free when $\ell\ne0$. More explicitly, on $L^\times$
use $\lambda=\ell$ in \eqref{cdp:eq:C.NH.6} to set $\ell=1$. This gives
the smooth coordinates
$(z,w,x)\in(\C^2\setminus\{0\})\times\C$, after replacing
the coordinates by $(\ell^rz,\ell w,\ell^{-1}x)$.
Thus the coarse nonzero line space $L^\times$ is a smooth
complex threefold, including the orbifold points.

\item All three maps in \eqref{cdp:eq:C.NH.5} lift canonically. In the
coordinates \eqref{cdp:eq:C.NH.4}, $A$ is $p\mapsto p+s^r$ and its lift
in \eqref{cdp:eq:C.NH.6} is $(z,w,x)\mapsto(z+w^r,w,x)$.
The weighted ideals are therefore preserved. The half
translation replaces the coefficient by $1/2$; $h$
interchanges the two centres and preserves their weighted
orders. The line lifts preserve its canonical divisor
section, so their group laws hold exactly, including
$\widetilde A_{1/2}^{\,2}=\widetilde A$ and
$\widetilde h^{\,2}=1$.

\item Multiply the line fibre by $\beta$ and compose with
$\widetilde A$. On the smooth coordinates just constructed
the resulting automorphism is
\begin{equation}
 g(z,w,x)=
       (\beta^r(z+w^r),\beta w,\beta^{-1}x),
 \qquad s=xw.
                                                        \hypertarget{tag-00SP}{}\label{cdp:eq:C.NH.8}
\end{equation}
It preserves the map to $C$. On $x=0$, the change
$Z=\beta^{-r}z$ turns its first coordinate into
$\beta^r Z+w^r$, proving the exact identification with
\eqref{cdp:eq:C.NH.1}.

\par\medskip\noindent\textbf{Uniform properness for every contraction.}
Choose a Hermitian
orbifold metric on $L$. Write $\sigma$ for its canonical
section. Near either centre, its norm off the exceptional
divisor is comparable, with constants independent of the
point, to
\begin{equation}
       \rho(p,s)=\bigl(d_B(p,0)^2+|s|^{2r}\bigr)^{1/(2r)}.
                                                        \hypertarget{tag-00SQ}{}\label{cdp:eq:C.NH.9}
\end{equation}
Here $d_B$ is a fixed flat elliptic distance. To see the
comparison, normalize the weighted coordinates so that
$|z|^2+|w|^{2r}=1$. Then
$(|p|^2+|s|^{2r})^{1/(2r)}=|x|$, which is the norm of the
canonical section up to a smooth positive bounded factor
on the compact weighted exceptional curve. Away from the
two centres all these norms are bounded above and below.
More precisely this comparison holds on $B$ times a small
$c_\pm$-disc: outside the chosen $p$-coordinate neighborhood,
both $\rho$ and $\|\sigma\|$ are bounded above and away from
zero. It therefore also covers an iterate $A^n$ that moves
$p$ out of that local coordinate chart.

\item For every integer $n$, the triangle inequality on the flat
elliptic curve gives
\[
 d_B(p+2n\psi(s),0)
       \le d_B(p,0)+C_0|n||s|^r.
\]
The same inequality applied to the inverse translation
gives the reverse comparison. Since the canonical lift
sends $\sigma(p,s)$ to $\sigma(p+2n\psi(s),s)$, there is
a global constant $C_1\ge1$ such that
\begin{equation}
 C_1^{-1}(1+|n|)^{-1/r}\|\ell\|
 \le \|\widetilde A^{\,n}\ell\|
 \le C_1(1+|n|)^{1/r}\|\ell\|.
                                                        \hypertarget{tag-00SA}{}\label{cdp:eq:C.NH.10}
\end{equation}
The comparisons first hold off the exceptional divisors.
The ratio of line norms extends continuously and positively
on the entire orbifold line, so they hold there as well.
On the compact part of $C$ away from both centres the
canonical section is nowhere zero and supplies the same
bounds with a fixed constant. This proves the uniform
global assertion.

\item Consequently $\|g^n\ell\|$ tends uniformly to zero on the
unit circle bundle as $n\to+\infty$, and to infinity as
$n\to-\infty$: multiply \eqref{cdp:eq:C.NH.10} by $|\beta|^n$.
The action is properly discontinuous because any two
compact subsets have line norms bounded above and away
from zero. It is free, since a fixed point of a nonzero
power would also be fixed by arbitrarily large positive
powers, contradicting the norm limit.

\item Choose $N$ such that $g^N$ decreases every norm by a factor
strictly less than one. Its minimum norm ratio is a positive
number $c_N$. Every $g^{N\Z}$-orbit meets the compact annulus
$c_N\le\|\ell\|\le1$: take the first iterate whose norm is
at most one. This annulus is compact over the compact
orbifold base. It follows that
\begin{equation}
                 \widetilde Y=L^\times/\langle g\rangle
                                                        \hypertarget{tag-00SB}{}\label{cdp:eq:C.NH.11}
\end{equation}
is a smooth compact threefold, with a proper map to $C$.


\item Off the two exceptional fibres, the canonical section
trivializes $L$ and gives the quotient
\begin{equation}
 (B\times\C^*)/
   \langle(p,\ell)\mapsto(p+2\psi(c),\beta\ell)\rangle.
                                                        \hypertarget{tag-00SC}{}\label{cdp:eq:C.NH.12}
\end{equation}
If $B=\C/\Lambda_B$, its additive covering lattice is
generated by $(\lambda,0)$ for $\lambda\in\Lambda_B$,
$(0,2\pi i)$, and $(a_c,\log\beta)$, where $a_c$ is any
lift of $2\psi(c)$. These four real generators are
independent because $\log|\beta|\ne0$. Changing the lift
of $a_c$ changes a generator by $\Lambda_B$. This proves
that \eqref{cdp:eq:C.NH.12} is precisely a compact two-torus.

\item At either distinguished value the exact function is
$s=xw$. Its full fibre is reduced SNC, with components
$H:x=0$ and $P:w=0$, meeting along
\begin{equation}
          J=(\C^*_z/\beta^{r\Z}),\qquad x=w=0.
                                                        \hypertarget{tag-00SD}{}\label{cdp:eq:C.NH.13}
\end{equation}
Both components and their intersection are compact since
they are closed in \eqref{cdp:eq:C.NH.11}. On $P$ the map to the
unchanged elliptic base $B$ has the exact local form
\begin{equation}
  p=x^r z,\qquad (z,x)\sim(\beta^r z,\beta^{-1}x).
                                                        \hypertarget{tag-00SE}{}\label{cdp:eq:C.NH.14}
\end{equation}
Away from $p=0$ its canonical frame is the product
$B\setminus\{0\}\times E_\beta$; on the overlap the
product coordinate is $c=x^{-1}$ and $z=pc^r$.
Let $L_0=\OO_B([0])$ and let $\sigma_0$ be its canonical section.
In the local frame $\sigma_0=p e$, the map to
\cref{lem:II:elliptic-root-surface} is
$a=x$, $\ell=z e$; its equation is $p=a^r z=x^r z$.
The deck map becomes $(\ell,a)\mapsto(\beta^r\ell,\beta^{-1}a)$.
On $B\setminus\{0\}$ the product parameter in that lemma is
$c=a^{-1}=x^{-1}$, with $\ell=\sigma_0c^r$ and $z=pc^r$.
These maps and their inverses agree on the overlap, so identify
the entire pair $(P,J)$ with
$(P_r(B,\beta),\C^*/\beta^{r\Z})$ from that lemma.
This includes $r=1$, where the surface is
$L_0^\times/\beta^{\Z}$ for the degree-one line bundle.
The lemma gives the full fibre $rJ$ and $e(P)=0$.
The Hopf surface and $J$ also have Euler number zero,
so the entire distinguished support has Euler number zero.
The local holomorphic canonical evaluation form is
$x^r\,dz\wedge dw\wedge dx$. The Jacobian of \eqref{cdp:eq:C.NH.8}
is $\beta^r$, and the multiplier of $x^r$ is $\beta^{-r}$,
so this form is invariant. Its holomorphic vanishing orders
along $(H,P)$ are exactly $(r,0)$. These orders give the strict canonical-order difference in \cref{cdp:st:C.2.UN.40} for this holomorphic form.

\par\medskip\noindent\textbf{The free quotient giving the rational base.}
Fix either square root $\beta^{1/2}$. The commuting
canonical lifts define
\begin{equation}
           G=\beta^{1/2}\widetilde A_{1/2}\widetilde h,
                    \qquad G^2=g.
                                                        \hypertarget{tag-00SF}{}\label{cdp:eq:C.NH.15}
\end{equation}
Thus $G$ induces an involution on $\widetilde Y$. A fixed
point there would have $h_Cc=c$. Such a point is not one
of the two centres. Moreover $\psi(c)$ is a two-torsion
point of $B$, so $2\psi(c)=0$. In the canonical line frame,
\begin{equation}
 G(p,c,\ell)=(-p-\psi(c),c,\beta^{1/2}\ell),
 \qquad g^n(p,c,\ell)=(p,c,\beta^n\ell).
                                                        \hypertarget{tag-00SG}{}\label{cdp:eq:C.NH.16}
\end{equation}
The equality $G(p,c,\ell)=g^n(p,c,\ell)$ is impossible by
the absolute values of the line coordinates: it would
give $|\beta|^{1/2}=|\beta|^n$. The involution is therefore
free, and
\begin{equation}
                 Y=\widetilde Y/\langle G\rangle
                  \longrightarrow C/\langle h_C\rangle=\PP^1
                                                        \hypertarget{tag-00SH}{}\label{cdp:eq:C.NH.17}
\end{equation}
is a smooth compact threefold with the required map.
The two Hopf components over $c_+,c_-$ are exchanged.
Restriction of the quotient to either one is an
isomorphism onto its image, so the prescribed Hopf
surface is unchanged. The entire distinguished fibre is
isomorphic to either $H+P$ above, and remains reduced.
At a fixed value of $h_C$ the full fibre is twice a free
quotient of a two-torus; its reduced support has Euler
number zero. At every other ordinary value it is a
two-torus. All fibres are connected, as is seen directly
from these quotient descriptions and from $H\cap P=J$.


\par\medskip\noindent\textbf{The component and support groups.}
The pair isomorphism above identifies the canonical form, divisor
and integral maps with those of \cref{lem:II:elliptic-root-surface}.
In the coordinates of \eqref{cdp:eq:C.NH.14}, its form is
$dp\wedge d\log(x^{-1})=-x^{r-1}dz\wedge dx$.
Thus the lemma gives $K_P=(r-1)J$,
$h^\bullet(\OO_P)=(1,2,1)$ and the integral ranks $(1,3,4,3,1)$.
Its circle action
$(z,x)\mapsto(e^{\mathrm i r\theta}z,e^{-\mathrm i\theta}x)$
is exactly the phase action in this chart. The resulting
mapping-torus trivialization therefore identifies this pair
with $(M_r\times S^1,\mathrm{core}\times S^1)$ for every
complex $\beta$ in the stated range.

\item By that pair map, the angular loop of $J$ is zero in $H_1(P,\Z)$,
its deck-period loop is primitive in the last $S^1$, and
$[J]=0$ in $H_2(P,\Z)$. In the Hopf component its angular
loop is zero and its radial loop is the primitive deck
generator. The integral normalization sequence of
$S=P\cup_JH$ consequently has the restriction maps
\begin{equation}
 H^0(P\sqcup H,\Z)\to H^0(J,\Z):\ (a,b)\mapsto a-b,
 \qquad
 H^1(P\sqcup H,\Z)\to H^1(J,\Z):
 \begin{pmatrix}0&0&1&-1\\0&0&0&0\end{pmatrix},
                                                        \hypertarget{tag-00SI}{}\label{cdp:eq:C.NH.18}
\end{equation}
up to simultaneous orientation choices for the radial
loops, and its degree-two map is zero. Both displayed
images are primitive. The exact sequence gives precisely
the free groups of ranks $(1,3,5,5,2)$.

\item Apply \cref{lem:II:hopf-curve-restriction} to
$H=H_r(\beta)$ and the axis inclusion
$J=\C^*/\beta^{r\Z}\hookrightarrow H$ in
\eqref{cdp:eq:C.NH.13}. Its generator is the same deck
transformation as in the weighted chart, so the constant
deck basis restricts to the nonzero axis class.
Thus the
coherent degree-one restriction onto $H^1(J,\OO_J)$
is surjective. Together with the degree-zero surjection
and the component groups already computed, the coherent
normalization sequence gives $h^\bullet(S,\OO_S)=(1,2,1)$.
All support maps are the inclusions in \eqref{cdp:eq:C.NH.13}--\eqref{cdp:eq:C.NH.14}.

\end{enumerate}
\end{proof}


\subsection{Secondary Hopf components of algebraic dimension zero}\label{subsec:II:secondary-hopf-occurrence}

\begin{proposition}[Secondary Hopf components of algebraic dimension zero]\hypertarget{tag-01DN}{}\label{cdp:st:C.2.UN.43}
\StatementRecord{cdp:st:C.2.UN.43}{proved}{01DN}{\SeriesRef{cdp:eq:C.SH.3}, \SeriesRef{cdp:st:C.2.UN.41}, \SeriesRef{cdp:eq:C.SH.9}, \SeriesRef{cdp:eq:C.SH.10}, \SeriesRef{cdp:st:C.2.UN.1}, \SeriesRef{cdp:st:C.2.UN.4}, \SeriesRef{cdp:st:C.2.UN.42}; \cite[arXiv v2, Theorem~2.1(1), (3), p.~5]{Khaled2024Hopf}; \cref{lem:II:central-primary-hopf-cover}; \cite[arXiv v3, \S7.6, pp.~33--34]{SchuttShioda2010}}

Every secondary Hopf surface $H$ of algebraic dimension zero
has a finite quotient presentation of one of the forms
\begin{equation}
\begin{array}{ll}
 H_{\alpha,\beta}/\langle(z,w)\mapsto
          (\zeta_n^a z,\zeta_n^b w)\rangle,
 &\gcd(a,n)=\gcd(b,n)=1,\\[2pt]
 H_r(\beta)/\langle(z,w)\mapsto
          (\zeta_n^r z,\zeta_n w)\rangle,
 &\gcd(r,n)=1.
\end{array}
\hypertarget{tag-0114}{}\label{cdp:eq:C.SH.1}
\end{equation}
Here $n$ is a positive integer and $\zeta_n=e^{2\pi i/n}$.
The first primary contraction is $(\alpha z,\beta w)$ with no
multiplicative resonance, and the second is
$(\beta^r z+w^r,\beta w)$, where $0<|\alpha|,|\beta|<1$.
These families give exhaustive quotient presentations, with parameters
satisfying the displayed contraction and coprimality conditions.

Every prescribed such surface occurs itself as a smooth
multiplicity-one component of a reduced SNC fibre of a proper
map from a smooth compact complex threefold to a smooth
compact curve, with connected fibres and general fibre a
complex two-torus. All reduced fibre supports can have Euler
number zero. In the diagonal construction the base is the finite torsion-cover
curve $B'$ in \eqref{cdp:eq:C.SH.6}.
\end{proposition}


\begin{proof}

\begin{enumerate}[label=\arabic*.,leftmargin=*]
\item Write $W=\C^2\setminus\{0\}$ and $H=W/\Gamma$.
Apply \cref{lem:II:central-primary-hopf-cover} to this given
surface. Its central contraction $F\in\Gamma$ gives
\begin{equation}
                 P=W/\langle F\rangle\longrightarrow H.
\hypertarget{tag-0119}{}\label{cdp:eq:C.SH.2}
\end{equation}
The pullback of meromorphic-function fields in that lemma
gives $a(P)=a(H)=0$ for this same covering.

\item Apply \cite[arXiv v2, normal forms (1), p.~3;
Theorem~2.1(1), (3), p.~5]{Khaled2024Hopf} to the primary
contraction $F$. In the diagonal algebraic-dimension-zero case
the relation $\alpha^p=\beta^q$, $p,q>0$, would make
$z^p/w^q$ a nonconstant meromorphic function on $P$.
Thus there is no such relation, and $F$ belongs to
class $\mathrm{II}_c$ and its commuting automorphisms are
$(z,w)\mapsto(Az,Bw)$, with $A,B\ne0$.
For the nonlinear contraction the classes $\mathrm{II}_a$
($r\ge2$) and $\mathrm{II}_b$ ($r=1$) give the full centralizer
\begin{equation}
       (z,w)\longmapsto(A^r z+B w^r,Aw),\qquad A\ne0.
\hypertarget{tag-011A}{}\label{cdp:eq:C.SH.3}
\end{equation}
Here the source parameters $(a,b)$ are $(A,B)$, and its
exponent is the same $r$ as in $F$.
The parameters $(A,B/A^r)$ in this formula have product
$(A_1,c_1)(A_2,c_2)=(A_1A_2,c_1+c_2)$. Both centralizers
are abelian. Since every element of $\Gamma$ commutes
with $F$, the deck group is abelian.

\item A finitely generated abelian group containing $\mathbb Z$
with finite index has the form $\mathbb Z\oplus T$, with $T$
its finite torsion subgroup. Choose a generator $f$ of the
infinite factor with the contracting sign. Some positive power
is $F^k$ times an element of $T$, $k>0$. Thus both diagonal
eigenvalue moduli of $f$ are less than one. In \eqref{cdp:eq:C.SH.3} it has
$|A|<1$ and nonzero additive parameter, since the parameter
of $F$ is nonzero. Rescaling $z$ makes its resonant coefficient
one. We obtain the holomorphic quotient maps
\begin{equation}
       H=(W/\langle f\rangle)/T,\qquad
       \Gamma=\langle f\rangle\times T.
\hypertarget{tag-011B}{}\label{cdp:eq:C.SH.4}
\end{equation}
For diagonal $T\subset(\mathbb C^*)^2$, each coordinate
character is injective: its kernel would fix a nonzero point
of the corresponding axis. Hence $T$ is cyclic and both
weights are units modulo its order. In \eqref{cdp:eq:C.SH.3}, a finite-order
element has additive parameter zero and acts as
$(A^r z,Aw)$. Its finite subgroup is cyclic, and freeness
on the first axis forces $\gcd(r,n)=1$. This proves \eqref{cdp:eq:C.SH.1}
and the finite-primary-cover comparison, including the
exclusion of a nonabelian group in algebraic dimension zero.


\item Conversely the indicated finite maps are unramified.
An equation $\tau^k(z,w)=f^\ell(z,w)$ at a nonzero coordinate
forces $\ell=0$ by the modulus of the contraction multiplier;
the coprimality conditions then force $k=0$ modulo $n$.
In the non-diagonal case use $w$ when it is nonzero and
$z$ on its complement. Thus these are the exact free
finite actions to be extended.

\par\medskip\noindent\textbf{A prescribed torsion section at the nodal fibre.}
For the diagonal case start with the compact elliptic surface
$J\to\mathbb P^1$, zero section, nodal fibre at zero and
section $P$ of \cref{cdp:st:C.2.UN.41}. Its multiplicative value at that
fibre is the prescribed $\alpha$. Explicitly, on
$y^2=x^2(x+1)$ the point with group coordinate $\alpha$ is
\begin{equation}
 \left(\frac{4\alpha}{(1-\alpha)^2},
       \frac{4\alpha(1+\alpha)}{(1-\alpha)^3}\right).
\hypertarget{tag-011C}{}\label{cdp:eq:C.SH.5}
\end{equation}
Use the cubic pencil and its nine base-point blowups in
\cref{cdp:st:C.2.UN.41}, with this point as $P_0$.
The resulting section is the same $P$ and the nodal fibre is unchanged.

\item Let $B^\circ$ be the smooth elliptic-fibre locus.
The exact-order-$n$ part of $J[n]|_{B^\circ}$ is a finite
\'etale algebraic cover. Select a component by the following
local section. The smooth group locus of the nodal
family has central group $\mathbb C^*$; multiplication by
$n$ has nonzero derivative at $\zeta_n$. The holomorphic
implicit-function theorem gives a unique local $n$-torsion
section specializing to $\zeta_n$, necessarily of exact order.
Take the connected cover containing its punctured restriction,
and normalize $\mathbb P^1$ in its function field. This gives
a smooth connected projective curve and finite map
\begin{equation}
                             B'\longrightarrow\mathbb P^1.
\hypertarget{tag-011D}{}\label{cdp:eq:C.SH.6}
\end{equation}
At the selected point $b_0$ this map is unramified: the local
sheet is the graph of the just constructed section and has
degree one over the disc. The $I_1$ fibre, its simple
smoothing parameter and $P(b_0)=\alpha$ remain unchanged.

\item Resolve the compact algebraic base-changed elliptic surface
and contract vertical $(-1)$-curves, away from this unchanged
tube, to obtain a smooth relatively minimal surface $J'\to B'$.
The pulled-back section and the tautological torsion point
extend as sections $P',R$: rational maps from the smooth
projective base curve to the proper regular surface extend by
the valuative criterion. They satisfy
\begin{equation}
                 P'(b_0)=\alpha,\qquad R(b_0)=\zeta_n,
                 \qquad \operatorname{ord}(R)=n.
\hypertarget{tag-011E}{}\label{cdp:eq:C.SH.7}
\end{equation}
Apply \cite[arXiv v3, Definition~3.1, p.~7;
\S7.6, pp.~33--34]{SchuttShioda2010} to the smooth projective
relatively minimal elliptic surface $J'\to B'$ and its sections
$P',R$. Their translations $T_{P'},T_R$ are automorphisms over
$B'$, with inverses $T_{-P'},T_{-R}$. They commute because
their restrictions to the dense smooth-fibre locus are the
commuting translations by these two sections.

\par\medskip\noindent\textbf{The diagonal action and the compact quotient.}
Put $q_0=\alpha\beta$ and form the compact suspension
\begin{equation}
 Y^{\rm std}=(\mathbb C^*_u\times J')/
       \langle(u,y)\mapsto(q_0u,y+P')\rangle.
\hypertarget{tag-011F}{}\label{cdp:eq:C.SH.8}
\end{equation}
Its proper compact bundle comparison is the one in
\cref{cdp:st:C.2.UN.41}. The commuting order-$n$ automorphism
\begin{equation}
               (u,y)\longmapsto
                  (\zeta_n^{a+b}u,y+aR)
\hypertarget{tag-011G}{}\label{cdp:eq:C.SH.9}
\end{equation}
descends. On the selected Tate tube $R=[\zeta_n]$ exactly,
by uniqueness of the torsion section in the smooth group
locus. In the multiplicative coordinate change $z=y$, $w=u/y$,
\eqref{cdp:eq:C.SH.9} is precisely
\begin{equation}
                    (z,w)\longmapsto
                    (\zeta_n^a z,\zeta_n^b w).
\hypertarget{tag-0115}{}\label{cdp:eq:C.SH.10}
\end{equation}
Replace the tube by \cref{cdp:st:C.2.UN.41}'s two-Hopf plumbing using
its exact punctured-family isomorphism. The action extends
on the first Hopf copy by \eqref{cdp:eq:C.SH.10}, on the second by the
inverse weights, and on the first neck by
\begin{equation}
 (z,x,y)\longmapsto
     (\zeta_n^a z,\zeta_n^b x,\zeta_n^{-b}y),
                         \qquad xy=t.
\hypertarget{tag-0116}{}\label{cdp:eq:C.SH.11}
\end{equation}
Interchange $a,b$ and the axes for the other neck.
Substitution in $(z,w)\mapsto(z^{-1},t/w)$ and
$(z,w)\mapsto(t/z,w^{-1})$ proves equivariance of both
gluing maps. The action is free on the whole central
support: on a double curve the equation
$\zeta_n^{ak}z=\alpha^\ell z$, or its counterpart, forces
$\ell=0$ by modulus and then $k=0$ by coprimality.
Off the double curves use the already verified Hopf
freeness. On a nearby torus it is a finite translation,
also free. Hence the quotient is smooth on this entire
tube and its two central components are the prescribed
secondary Hopf surface itself. The inverse weights
generate the same subgroup.

\item Take the finite quotient of the entire compact threefold
and resolve it over the other finitely many bad values,
leaving this smooth tube unchanged. General fibres are
quotients of tori by translations, hence tori. Connectedness
passes through the finite quotient and the connected fibres
of the birational resolution. The distinguished reduced
SNC fibre has two Hopf components and two elliptic double
curves, and therefore Euler number zero. Every other
reduced fibre has even Euler number by \cref{cdp:st:C.2.UN.1}.
Apply \cref{cdp:st:C.2.UN.4} only over those other values to make their
Euler numbers zero; none of these operations changes
the specified Hopf surface or its distinguished tube.
This proves compact diagonal secondary occurrence over
the curve $B'$.

\par\medskip\noindent\textbf{The nonlinear scalar action.}
Use the intermediate compact family $\widetilde Y\to C$ in
\eqref{cdp:eq:C.NH.11} for the prescribed $H_r(\beta)$, before the
additional involution in \eqref{cdp:eq:C.NH.15}. Its smooth nonzero
tautological line space has coordinates
\begin{equation}
 (z,w,x)\in(\mathbb C^2\setminus\{0\})\times\mathbb C,
                   \qquad p-p_0=x^r z,\quad s=xw.
\hypertarget{tag-0117}{}\label{cdp:eq:C.SH.12}
\end{equation}
Multiplication of that line by $\zeta_n$ is the map
\begin{equation}
 (z,w,x)\longmapsto
       (\zeta_n^r z,\zeta_n w,\zeta_n^{-1}x).
\hypertarget{tag-0118}{}\label{cdp:eq:C.SH.13}
\end{equation}
It commutes with the lifted base translations and the
contracting deck map. It is free when $\gcd(r,n)=1$:
a nonzero $w$ or $x$ forces the corresponding root to be
one; if both vanish then $z\ne0$ and its $r$-th power
must be one. It stays free after the infinite cyclic
quotient. Away from the exceptional component the
canonical line trivialization turns equality with a
deck power into $\zeta_n^k=\beta^\ell$, so $\ell=0$.
On the Hopf component use its nonzero $w$ coordinate,
or its $z$ coordinate if $w=0$, for the same modulus
argument. Coprimality then finishes the proof.

\item The resulting quotient is a smooth compact threefold.
Its desired component is exactly the second quotient
in \eqref{cdp:eq:C.SH.1}, with no change to that surface. On every
general torus fibre the action is a line-coordinate
translation, so the quotient is a torus. The distinguished
full fibre remains reduced SNC, and its Euler number
zero is preserved by the free finite quotient. Any
remaining values can be balanced by \cref{cdp:st:C.2.UN.4} away
from that fibre.

The base of the nonlinear quotient is the compact curve $C$ in
\eqref{cdp:eq:C.NH.3}; the construction uses $\widetilde Y\to C$
before the involution in \eqref{cdp:eq:C.NH.15}.

\end{enumerate}
\end{proof}


\subsection{Ruled components over every compact curve}\label{subsec:II:ruled-component-occurrence}

\begin{proposition}[Ruled components over every compact curve]\hypertarget{tag-01DP}{}\label{cdp:st:C.2.UN.45}
\StatementRecord{cdp:st:C.2.UN.45}{proved}{01DP}{\SeriesRef{cdp:st:C.2.UN.4}, \SeriesRef{cdp:eq:C.RU.7}, \SeriesRef{cdp:eq:C.RU.8}, \SeriesRef{cdp:eq:C.RU.5}; \cref{lem:II:blowup-topology}; \cite[Theorem~4D.1, pp.~432--434]{Hatcher2002}}

For every smooth compact connected complex curve $C$, there is a smooth
compact complex threefold $Y$ and a proper map $f:Y\to\PP^1$ with connected
fibres, whose restriction over $\PP^1\setminus\{0\}$ is the product with a
fixed two-torus, and whose reduced zero fibre has Euler number zero and
contains a smooth multiplicity-one component
\begin{equation}
                         R=\PP(N_{C/U})\longrightarrow C.
\hypertarget{tag-010C}{}\label{cdp:eq:C.RU.1}
\end{equation}
Here $C\hookrightarrow U$ is the explicit vertical embedding constructed
below. The component $R$ is retained unchanged by the Euler-balancing
operations. Its birational class is that of $C\times\PP^1$.
Consequently every ruled surface birational class occurs in a compact zero-Euler torus fibre. A rational component is also obtained
unchanged as the exceptional plane of a point blowup followed by balancing
away from that plane.

Writing $g=g(C)$, the integral cohomology of $R$ is torsion-free and has ranks
\begin{equation}
              (b_0,b_1,b_2,b_3,b_4)(R)=(1,2g,2,2g,1),
              \qquad h^\bullet(R,\OO_R)=(1,g,0).
\hypertarget{tag-010D}{}\label{cdp:eq:C.RU.2}
\end{equation}
The rank-two bundle in \eqref{cdp:eq:C.RU.1} is the normal bundle of
the constructed embedding; the conclusion ranges over all ruling bases
and their ruled birational classes.
\end{proposition}


\begin{proof}

\begin{enumerate}[label=\arabic*.,leftmargin=*]
\item Choose a nonconstant meromorphic function $x$ on $C$, by
Riemann--Roch. It gives a finite map $C\to\PP^1$, so the meromorphic
function field of $C$ is a finite separable extension of $\C(x)$.
Choose a primitive element $y$. The meromorphic map
\begin{equation}
                    C\longrightarrow\PP^2,\qquad [1:x:y]
\hypertarget{tag-010E}{}\label{cdp:eq:C.RU.3}
\end{equation}
extends at every point: multiply its three local meromorphic coordinates
by a power of a uniformizer so that their minimum vanishing order is zero.
At least one resulting holomorphic coordinate is nonzero. The image is an
integral projective plane curve $D$, and the map is finite and birational,
because $\C(x,y)=\mathcal M(C)$. Thus $C$ is the normalization of
$D$. For $C=\PP^1$ one may choose $x=s,y=s^2$, giving a smooth conic.
In every case choose projective coordinates so that $D$ is not the fixed
line used below.

\item Start with $A\times\PP^1$, where $A$ is any fixed smooth complex two-torus,
and blow up a point of $A\times\{0\}$. Its exceptional divisor is a $\PP^2$. Place $D$ in this plane. Successively blow up singular
points of its strict transform, performing each point blowup in the
ambient smooth threefold. The strict transform of the plane is precisely
its surface point blowup: in local coordinates this is the strict
transform of the coordinate plane in the three ordinary blowup charts.
The full map to $\PP^1$ is unchanged over its punctured base.

\item For an integral curve $D_0$ on a smooth surface,
normalization gives
\begin{equation}
  p_a(D_0)-g(C)=\operatorname{length}
                (\nu_*\OO_C/\OO_{D_0})\ge0.
\hypertarget{tag-010F}{}\label{cdp:eq:C.RU.4}
\end{equation}
At a singular point the multiplicity $m$ is at least two. Under a surface
point blowup, $D_1=\pi^*D_0-mE$ and
$K_{V_1}=\pi^*K_{V_0}+E$. Adjunction gives
\begin{equation}
                        p_a(D_1)=p_a(D_0)-\frac{m(m-1)}2.
\hypertarget{tag-010G}{}\label{cdp:eq:C.RU.5}
\end{equation}
The normalization is unchanged through this proper birational map.
Therefore its total nonnegative normalization length decreases by at
least one at every singular-point blowup. After finitely many such
blowups the strict transform is smooth and is the specified $C$ itself.
This gives a closed embedding $C\hookrightarrow U$ into the
resulting smooth compact threefold.

\item The strict transform of the original exceptional plane has full-fibre
multiplicity one, since that was its coefficient after the first point
blowup and later strict transforms retain the coefficient. The curve
$C$ is not any boundary curve of this plane: it dominates the original
$D$, which is different from the line of intersection with the original
torus component, whereas every subsequently exceptional boundary curve
maps to a point of that plane. Consequently a generic point of $C$ lies
only in this multiplicity-one component. Blow up the smooth curve $C$ in
$U$. The exceptional divisor is exactly the projectivized normal bundle
in \eqref{cdp:eq:C.RU.1}, with coefficient one in the full fibre. This follows also
from the local charts of the blowup of $(x,y)$: its exceptional fibre at
each point of $C$ is the projective space of normal directions.

\item Choose two meromorphic sections of $N_{C/U}$ which form a basis at the
generic point. Such sections exist by twisting by a sufficiently positive
line bundle on the curve and applying Riemann--Roch. They give a bundle trivialization away from a finite set, and hence a birational map
\begin{equation}
                         \PP(N_{C/U})\dashrightarrow C\times\PP^1.
\hypertarget{tag-010H}{}\label{cdp:eq:C.RU.6}
\end{equation}
The same meromorphic-frame map applies to every rank-two bundle on $C$.
It therefore identifies this birational class with that of every ruled
surface over this specified curve. In genus zero it also identifies the
rational class, for example through the affine-coordinate birational map
$\PP^2\dashrightarrow\PP^1\times\PP^1$.

\item Apply \cref{lem:II:blowup-topology} to the reduced full
inverse image of the zero-fibre support at every blowup.
For a point centre in the threefold, its formula
\eqref{cdp:eq:C.VR.4} has codimension three and adds $2$.
For the smooth curve $C$, it has codimension two and adds
$e(C)=2-2g$. These are the same full inverse images even
when a centre meets several support components.
If $p$ is
the total number of point blowups, including the initial one, the support
after blowing up $C$ has
\begin{equation}
                              e=2p+2-2g.
\hypertarget{tag-010I}{}\label{cdp:eq:C.RU.7}
\end{equation}
It is even. There is a smooth open part of the original torus component
disjoint from $R$, since only finitely many proper curves and points have
been affected there. Perform the positive point blocks or negative
point--quartic blocks of \cref{cdp:st:C.2.UN.4} in such an open part and in the new
exceptional components they create. Each changes Euler number by exactly
$2$ or $-2$. Choosing the indicated finite number makes \eqref{cdp:eq:C.RU.7} zero.
All these centres are disjoint from $R$, so the displayed smooth component,
its normal projection to $C$, and its coefficient one are unchanged.
All other fibres remain the original smooth tori. Every blowup has
connected fibres, so the composite map retains connected fibres as well.
Compactness and smoothness follow at each step from the point or
smooth-curve blowup of the compact smooth threefold.

\item For the cohomology comparison use the lines convention in \eqref{cdp:eq:C.RU.1} and
put $\xi=c_1(\OO_{\PP(N)}(1))$. Pullback and cup product give the integral
projective-bundle map
\begin{equation}
 H^k(C,\Z)\oplus H^{k-2}(C,\Z)
       \longrightarrow H^k(R,\Z),\qquad
 (a,b)\longmapsto\pi^*a+\xi\smile\pi^*b.
\hypertarget{tag-010J}{}\label{cdp:eq:C.RU.8}
\end{equation}
Apply \cite[Theorem~4D.1, pp.~432--434]{Hatcher2002} to
this fibre bundle $\pi:R\to C$, with coefficient ring $\Z$ and
the global classes $1,\xi$. Their restrictions form the integral
basis of each $\PP^1$ fibre, and the source map is precisely
\eqref{cdp:eq:C.RU.8}. For the lines convention, the tautological
sequence $0\to\OO_R(-1)\to\pi^*N\to Q\to0$ has
$c_1(Q)=\pi^*c_1(N)+\xi$ and $c_2(\pi^*N)=0$.
It gives $\xi^2+\pi^*c_1(N)\xi=0$. Since the curve cohomology is
torsion-free, \eqref{cdp:eq:C.RU.8} proves the integral assertion. The fibrewise
calculation $H^0(\PP^1,\OO)=\C$, $H^1(\PP^1,\OO)=0$ gives
$\pi_*\OO_R=\OO_C$, $R^1\pi_*\OO_R=0$, with their restriction and
evaluation maps. Leray therefore gives the coherent tuple in \eqref{cdp:eq:C.RU.2}.

\item There is an explicit all-genus computation supporting the construction.
For $g\ge1$ take the plane curve
\begin{equation}
                 y^2z^{2g-1}=x^{2g+1}-z^{2g+1}.
\hypertarget{tag-010K}{}\label{cdp:eq:C.RU.9}
\end{equation}
It is smooth on $z\ne0$. At its single point at infinity, use
$u=x/y,v=z/y$; the equation is $v^{2g-1}(1+v^2)=u^{2g+1}$.
The holomorphic change $v'=v(1+v^2)^{1/(2g-1)}$ identifies its germ with
the plane branch $v'^{2g-1}=u^{2g+1}$. The two exponents are coprime.
Successive point blowups subtract the smaller exponent from the larger,
and the strict curve is smooth when the smaller exponent becomes one.
Summing \eqref{cdp:eq:C.RU.5} through the subtraction sequence gives total defect
$(2g-2)(2g)/2=2g(g-1)$; subtracting it from the arithmetic
genus $g(2g-1)$ gives $g$. Formula \eqref{cdp:eq:C.RU.7}
then gives the Euler change and the required number of
balancing blocks, and \eqref{cdp:eq:C.RU.8} gives the integral
projective-bundle ranks for every $g$.

\end{enumerate}
\end{proof}


\section{Marked Hopf and ruled cycles}\hypertarget{tag-0346}{}\label{sec:II:10}

\subsection{Logarithmic period families}\label{subsec:II:logarithmic-period-families}\label{subsec:II:marked-cycle-tree-constructions}

\begin{lemma}[A logarithmic family of two-tori]\hypertarget{tag-P331}{}\label{lem:II:logarithmic-period-family}
\StatementRecord{lem:II:logarithmic-period-family}{proved}{P331}{}
Let $\Lambda\subset\C^2$ have the integral basis
\begin{equation}
 \begin{gathered}
 w_1=(a_1,c_1),\qquad w_2=(a_2,c_2),\qquad p=(0,1),\\
 c_1,c_2\in\R,\qquad
 C=\operatorname{Im}(\overline{a_1}a_2)>0.
 \end{gathered}
 \hypertarget{tag-013Q}{}\label{cdp:eq:C.TL.1}
\hypertarget{tag-009H}{}\SsixEquationAliases{\label{cdp:eq:C.AD.5}}
\end{equation}
Let $N$ be a positive integer, and put
\begin{equation}
                         b=(\sigma,i).
 \hypertarget{tag-013R}{}\label{cdp:eq:C.TL.2}
\end{equation}
Here $\sigma\in\C$. For a fixed vector $w=(w^{(1)},w^{(2)})$, define
\begin{equation}
 V_z(w)=
 \bigl(w^{(1)}+z\overline{w^{(2)}},
       w^{(2)}-z\overline{w^{(1)}}\bigr).
 \hypertarget{tag-013S}{}\label{cdp:eq:C.TL.3}
\hypertarget{tag-00TG}{}\SsixEquationAliases{\label{cdp:eq:C.NS.15}}
\end{equation}
On any branch of $\log z$ in $\C^*$, let
\begin{equation}
 \begin{gathered}
 P_1(z)=V_z(w_1),\quad P_2(z)=V_z(w_2),\quad
 P_3(z)=V_z(p),\\
 P_4(z)=V_z(b)+\frac{N}{2\pi i}\log z\,V_z(p).
 \end{gathered}
 \hypertarget{tag-013T}{}\label{cdp:eq:C.TL.4}
\hypertarget{tag-009I}{}\SsixEquationAliases{\label{cdp:eq:C.AD.6}}
\end{equation}
Then $T_z=\C^2/\langle P_1(z),P_2(z),P_3(z),P_4(z)\rangle_{\Z}$ is a smooth proper family of connected complex two-tori over $\C^*$. Changing $\log z$ by $2\pi i$ sends $P_4$ to $P_4+NP_3$. Put $r=|z|$. Its real determinant, in the coordinate order $(\operatorname{Re}w^{(1)},\operatorname{Im}w^{(1)},\operatorname{Re}w^{(2)},\operatorname{Im}w^{(2)})$, is
\begin{equation}
 \begin{aligned}
 D(z)&=\det_{\R}(P_1(z),P_2(z),P_3(z),P_4(z))\\
 &=C(1+r^2)
       \left(1+r^2+\frac{N}{2\pi}(r^2-1)\log r\right)>0
                         \qquad (0<r<\infty).
 \end{aligned}
 \hypertarget{tag-013V}{}\label{cdp:eq:C.TL.6}
\hypertarget{tag-009J}{}\SsixEquationAliases{\label{cdp:eq:C.AD.7}}
\end{equation}
Writing $s=z^{-1}$ and changing the vector frame by $\eta_\infty=\eta/z$ gives
\begin{equation}
 V^\infty_s(w)
       =\bigl(\overline{w^{(2)}},-\overline{w^{(1)}}\bigr)
                          +s w.
 \hypertarget{tag-013W}{}\label{cdp:eq:C.TL.7}
\hypertarget{tag-009L}{}\SsixEquationAliases{\label{cdp:eq:C.AD.9}}
\end{equation}
and the four periods
\begin{equation}
 \begin{gathered}
 V^\infty_s(w_1),\quad V^\infty_s(w_2),\quad V^\infty_s(p),\\
 V^\infty_s(b)
          -\frac{N}{2\pi i}\log s\,V^\infty_s(p).
 \end{gathered}
 \hypertarget{tag-013X}{}\label{cdp:eq:C.TL.8}
\hypertarget{tag-009C}{}\SsixEquationAliases{\label{cdp:eq:C.AD.10}}
\end{equation}
In particular $V_0^\infty(w)=J\overline w$, where $J=\left(\begin{smallmatrix}0&1\\-1&0\end{smallmatrix}\right)$.

For every integer $R\ge1$, replace each $V_z$ in
\eqref{cdp:eq:C.TL.4} by $V_{z^R}$ and keep its coefficient
$N\log z/(2\pi i)$. This also gives a smooth proper family of
connected two-tori on $\C^*$, with determinant
\begin{equation}
 C(1+r^{2R})\left(1+r^{2R}
       +\frac{N}{2\pi}(r^{2R}-1)\log r\right)>0,
                          \qquad r=|z|>0.
\label{eq:II:logarithmic-power-determinant}
\hypertarget{tag-00VJ}{}\hypertarget{tag-0185}{}\hypertarget{tag-00JJ}{}
\SsixEquationAliases{\label{cdp:eq:C.PR.11}\label{cdp:eq:C.WT.24}\label{cdp:eq:C.EW.11}}
\end{equation}
The frame change is division by $z^R$, giving sections
$J\overline w+s^Rw$ and fourth period
$J\overline b+s^Rb-N\log s\,(J\overline p+s^Rp)/(2\pi i)$.
The logarithmic monodromy still adds $N$ times the third period.

For every full rank-four lattice $\Lambda_0\subset\C^2$, the
sections $V_t(\lambda)$, $\lambda\in\Lambda_0$, define a free
translation action on $\OO_{\PP^1}(1)^{\oplus2}$, using
$V_s^\infty(\lambda)=J\overline\lambda+s\lambda$ in the frame
$s=t^{-1}$. Its quotient
\begin{equation}\label{eq:II:compact-twistor-lattice}
 \mathcal T(\Lambda_0)=\OO_{\PP^1}(1)^{\oplus2}/\Lambda_0
                  \longrightarrow\PP^1
\end{equation}
is a smooth proper family of connected complex two-tori with
central fibre $\C^2/\Lambda_0$. Its total space is compact.
In the coordinates $w=(p,z)$ and base parameter $t$, the same
sections are $\ell_t(p,z)=V_t((p,z))$ in
\eqref{cdp:eq:C.NS.15}. Pullback along $t=s^R$, for every
integer $R\ge1$, remains a smooth compact family. 

\end{lemma}
\begin{proof}
For fixed $w$, $V_z(w)$ is holomorphic in $z$. The identity $\langle w,J\overline w\rangle=0$ gives $\|V_z(w)\|^2=(1+|z|^2)\|w\|^2$. Thus $V_z$ is an invertible real-linear map with determinant $(1+|z|^2)^2$; its orientation is positive by continuity from $V_0=I$. Its inverse is $v\mapsto(v-zJ\overline v)/(1+|z|^2)$.
For a full lattice $\Lambda_0$, these invertible real maps
carry any chosen real lattice basis to four real-independent
periods at every finite $t$. In the other frame the sections
are $J\overline\lambda+s\lambda$, whose squared norm is
$(1+|s|^2)\|\lambda\|^2$, by the same orthogonality identity.
They therefore remain a full lattice at $s=0$ as well.
On a closed subdisc in either chart the real period matrices
and their inverses are uniformly bounded. Integral translations
act freely and properly discontinuously there; their closed
real fundamental parallelepipeds cover the quotient over the
closed subdisc and are compact. The frame change
$\eta_\infty=\eta/t$ sends $V_t(\lambda)$ exactly to
$V_s^\infty(\lambda)$, so these local quotients glue as
\eqref{eq:II:compact-twistor-lattice}. The covering quotient is
a holomorphic submersion and is proper by these compact
parallelepipeds. A finite closed cover of $\PP^1$ makes its
total space compact. Base change of a proper holomorphic
submersion preserves both properties; the base-changed total
space is compact because its base is again $\PP^1$.
 For $p=(0,1)$, this gives
\begin{equation}
 V_z^{-1}\bigl(iV_z(p)\bigr)
       =\frac{i(2z,1-r^2)}{1+r^2}.
 \hypertarget{tag-013U}{}\label{cdp:eq:C.TL.5}
\end{equation}
The first three original periods span $\C\times\R$, and their real determinant with a fourth vector $v$ is $C\operatorname{Im}v^{(2)}$. Put $r=|z|$ and $\tau=N\log z/(2\pi i)$. The contribution $\operatorname{Re}\tau\,V_z(p)$ does not change the determinant. The imaginary contribution uses \eqref{cdp:eq:C.TL.5}, together with $\operatorname{Im}\tau=-N\log r/(2\pi)$ and $\operatorname{Im}b^{(2)}=1$. Multiplication by the real determinant of $V_z$ gives exactly \eqref{cdp:eq:C.TL.6}. It is positive because $(r^2-1)\log r\ge0$.

For the power-$R$ family make the same determinant calculation
with the real-linear map $V_{z^R}$ and
$\tau=N\log z/(2\pi i)$. Its inverse and determinant are the
preceding ones at $z^R$; in particular the second coordinate of
$V_{z^R}^{-1}(iV_{z^R}p)$ is $i(1-r^{2R})/(1+r^{2R})$.
The same fourth-column calculation gives
\eqref{eq:II:logarithmic-power-determinant}. It is positive because
$(r^{2R}-1)\log r\ge0$. Division by $z^R$ gives its displayed
infinity frame and the same negative logarithmic coefficient.

For either family, on a simply connected base chart the four holomorphic periods act by translations on $\C^2$. Their real matrices and inverses are uniformly bounded over every compact base subset. Only finitely many integral translation vectors can move one prescribed compact subset into another, so the action is properly discontinuous. The continuously varying closed real parallelepipeds give properness of the quotient over that base subset. The integral change $P_4\mapsto P_4+NP_3$ identifies these quotients on changes of logarithm, constructing the stated family on all of $\C^*$.

Division of each period by $z$ gives \eqref{cdp:eq:C.TL.7}. Since $\log z=-\log s$ up to $2\pi i\Z$, the fourth period becomes the fourth vector in \eqref{cdp:eq:C.TL.8}, with precisely its negative logarithmic coefficient. Its branch changes again add an integral period. This proves the full frame and period comparison at infinity.
\end{proof}

\subsection{The marked-cycle criterion}\label{subsec:II:marked-cycle-criterion}

\begin{proposition}[Marked cycles with torus smoothings]
\hypertarget{tag-01E2}{}\label{cdp:st:C.2.UN.57}
\StatementRecord{cdp:st:C.2.UN.57}{proved}{01E2}{\cref{def:II:marked-normal-cycle,lem:II:moving-marked-cycle,lem:II:logarithmic-period-family,cdp:st:C.2.UN.40,lem:II:central-primary-hopf-cover}; \cite[arXiv v1, \S2.2.6, p.~5]{AbashevaDeev2024AB}}
\hypertarget{tag-01DQ}{}\label{cdp:st:C.2.UN.46}
\StatementRecord{cdp:st:C.2.UN.46}{proved}{01DQ}{}
\hypertarget{tag-01DW}{}\label{cdp:st:C.2.UN.51}
\StatementRecord{cdp:st:C.2.UN.51}{proved}{01DW}{}
\hypertarget{tag-01E0}{}\label{cdp:st:C.2.UN.55}
\StatementRecord{cdp:st:C.2.UN.55}{proved}{01E0}{}
For a marked normal cycle in \cref{def:II:marked-normal-cycle}, its proper normal-product plumbing satisfies
\begin{equation}
 \text{$Y_t$ is a two-torus}\quad(0<|t|\ll1)
       \quad\Longleftrightarrow\quad M=I.
\hypertarget{tag-00LR}{}\label{cdp:eq:C.HC.2}
\end{equation}
Every smooth-total two-torus smoothing with this full reduced central fibre has $M=I$. This identity is also equivalent to occurrence of the specified support as a full reduced fibre of a smooth compact complex threefold over $\PP^1$ with connected two-torus general fibre. The compact construction preserves every marked component and every elliptic gluing map and chooses the normal-product constants. Its only other singular fibre is the complex conjugate marked support. Every reduced fibre support has Euler characteristic zero.

A compatible cycle containing a Hopf component contains at least two Hopf components. For an all-ruled cycle, including the two-section identification of one ruled normalization, the projection of $b_0$ in \eqref{cdp:eq:C.HC.10} to the common elliptic quotient is the full elliptic translation around the cycle. The compact construction preserves that translation.

Every full reduced SNC fibre in a smooth-total two-torus degeneration whose components are unblown diagonal Hopf surfaces of algebraic dimension zero is such a cycle, with no triple points and with both axes of every component in its boundary. For two components, the Hopf surfaces are isomorphic, and a diagonal automorphism of one copy makes both gluing maps inversion. Thus the support is two copies of the same primary or secondary diagonal Hopf surface joined by the two inversion maps. Every such support has the compact realization above. The two normal-product constants remain the parameters of the local quotient \eqref{cdp:eq:C.HC.9}--\eqref{cdp:eq:C.HC.10}.

The identification of the two distinct axes of a single diagonal Hopf surface has $M\ne I$ and admits no smooth-total two-torus smoothing.
\end{proposition}
\begin{proof}
Apply \cref{lem:II:moving-marked-cycle} to the specified components, elliptic identifications and dual-normal maps. It supplies the proper plumbing. For a given smoothing the transverse bilinear term of its defining function gives a dual-normal trivialization on each double curve. Its ratio with the specified trivialization is a nowhere-zero holomorphic function on that compact elliptic curve, hence a constant. By \cref{def:II:marked-normal-cycle} this leaves $M$ unchanged. The lemma therefore gives the same presentation \eqref{cdp:eq:C.HC.7} for the given smoothing, including the branch-labelled self-identification. The injected subgroup $\Lambda_1$ and the cycle generator commute precisely when $M=I$. A two-torus has abelian fundamental group, giving necessity. For $M=I$ the lemma's closed coframe and orbit map give the biholomorphism \eqref{cdp:eq:C.HC.8} with the full lattice \eqref{cdp:eq:C.HC.10}. This proves \eqref{cdp:eq:C.HC.2} with all affine constants retained.

Assume $M=I$ and use the common logarithmic coordinates \eqref{cdp:eq:C.HG.2}. Let $n_i\in\Lambda$ be the positive outgoing normal meridian. The incoming one at vertex $i$ is $-n_{i-1}$. At every ruled vertex the two meridians are opposite, giving
\begin{equation}
                      n_i=n_{i-1}
                  \quad\hbox{at every ruled vertex}.   \hypertarget{tag-00RW}{}\label{cdp:eq:C.MX.5}
\end{equation}
At a Hopf vertex the two complex normal lines are independent. The frames at those vertices and the total normal period of the cycle are
\begin{equation}
 P_i=\left(\frac{n_i}{2\pi i},
                   -\frac{n_{i-1}}{2\pi i}\right),
 \qquad
 2\pi i v_0=-\sum_i n_i=:u\in\Lambda.
\hypertarget{tag-00M9}{}\label{cdp:eq:C.HG.3}
\hypertarget{tag-009E}{}\SsixEquationAliases{\label{cdp:eq:C.AD.2}}
\end{equation}
The sum follows from the edge coefficient $-n_i/(2\pi i)$ in \eqref{eq:II:moving-cycle-periods}. A cycle with one Hopf vertex would force $n_i=n_{i-1}$ there by following \eqref{cdp:eq:C.MX.5} through all the other vertices, contradicting that independence. In particular a single Hopf self-identification is impossible when $M=I$. Equivalently, its sole matrix in \eqref{cdp:eq:C.HC.5} has upper-left entry zero and cannot be the identity.

Put $W=\operatorname{span}_{\R}\Lambda$ and choose a nonzero real linear functional $\ell$ with kernel $W$. A Hopf contraction with logarithms $(\log\alpha_i,\log\beta_i)$ gives the common lattice vector $P_i(\log\alpha_i,\log\beta_i)$. Applying $\ell$ to this vector and using \eqref{cdp:eq:C.HG.3} gives
\begin{equation}
 \log|\alpha_i|\,\ell(i n_i)
       =\log|\beta_i|\,\ell(i n_{i-1}).
\hypertarget{tag-00MA}{}\label{cdp:eq:C.HG.4}
\hypertarget{tag-00RX}{}\SsixEquationAliases{\label{cdp:eq:C.MX.6}\label{cdp:eq:C.AD.3}}
\hypertarget{tag-009F}{}
\end{equation}
Both logarithms are negative. If one transverse value vanished, both would vanish, placing the two independent complex normal lines in the real three-plane $W$. Thus both values are nonzero and have the same sign. Equation \eqref{cdp:eq:C.MX.5} propagates that sign through every ruled vertex. In the all-ruled case, their common normal $n$ has $\C n\cap W=\R n$: projection along $\C n$ gives the rank-two lattice of its elliptic base. Hence the common transverse value is again nonzero. In every case
\begin{equation}
                       \ell(iu)=-\sum_i\ell(i n_i)\ne0.
\hypertarget{tag-00MB}{}\label{cdp:eq:C.HG.5}
\hypertarget{tag-00RY}{}\SsixEquationAliases{\label{cdp:eq:C.MX.7}}
\end{equation}
It follows that $u\ne0$ and $\C u\cap W=\R u$. Write $u=d p$ with $p$ primitive in $\Lambda$ and $d>0$. The common affine edge constants $\mathbf b_i$ have sum $b$, and the fourth period is
\begin{equation}
 b+\frac{\log t}{2\pi i}u,\qquad
             u=-\sum_i n_i\in\Lambda,
                                                        \hypertarget{tag-009G}{}\label{cdp:eq:C.AD.4}
\end{equation}

The real quotient $W/\R p\to\C^2/\C p$ is an isomorphism. Put $K=W\cap iW$, a complex line. Since $ip\notin W$, one has $W=K\oplus\R p$, and a nonzero vector of $K$ together with $p$ is a complex basis of $\C^2$. Send $K$ to the first coordinate line and $p$ to $(0,1)$; this complex linear isomorphism sends $W$ to $\C\times\R$. Extend $p$ to an integral basis and orient its two remaining projected periods so that the lattice has exactly the form \eqref{cdp:eq:C.TL.1}. Every $n_i$ has nonzero second coordinate, by the preceding transverse calculation. Changing one normal-product constant adds an arbitrary complex multiple of this $n_i$ to its edge translation, so choose it to make $b=(\sigma,i)$ as in \eqref{cdp:eq:C.TL.2}. The elliptic gluing map is unchanged because this direction is its normal line. In the all-ruled case $p=-n$ and $u=Np$. Every adjustable direction lies in $\C p$, so the projection of $b$ to $\C^2/(\C p+\Lambda)$, including its full horizontal translation $\sigma$, remains unchanged.

Apply \cref{lem:II:logarithmic-period-family} to this exact lattice, this $b$ and the positive integer $d$ in place of its $N$. Its periods \eqref{cdp:eq:C.TL.4}, with $\Lambda(z)=V_z\Lambda$ and fourth period $B(z)$, are the quotient in \eqref{cdp:eq:C.AD.6}; its determinant \eqref{cdp:eq:C.TL.6} gives \eqref{cdp:eq:C.AD.7}. This constructs the smooth proper family of two-tori on all of $\C^*$.

At zero apply the moving clause of \cref{lem:II:moving-marked-cycle} with $\lambda_j(z)=V_z(\lambda_j)$ and $\mathbf b_i(z)=V_z(\mathbf b_i)$. Every integer coordinate of every normal meridian and finite character is unchanged. Its Hopf frames \eqref{cdp:eq:C.HG.12} are precisely \eqref{cdp:eq:C.AD.8}, while ruled vertices use \eqref{cdp:eq:C.MX.10}. The central components and elliptic maps are the specified ones. The sum of the affine constants is $V_z(b)$, and the normal sum is $dV_z(p)$, so \eqref{eq:II:moving-cycle-periods} is exactly the fourth period of \eqref{cdp:eq:C.AD.6}. The lemma's piecewise logarithmic maps and their holomorphic inverses identify the entire punctured filling with that period family.

At infinity put $s=z^{-1}$ and use the frame $\eta_\infty=\eta/z$ of \cref{lem:II:logarithmic-period-family}. Write $W_s=V_s^\infty$ in \eqref{cdp:eq:C.TL.7}; this is the map \eqref{cdp:eq:C.AD.9}. Choose outgoing normals $n_i^\infty(s)=-W_s(n_i)$ and affine constants $\mathbf b_i^\infty(s)=W_s(\mathbf b_i)$. Their total normal period is $-dW_s(p)$ and their constant sum is $W_s(b)$, giving exactly \eqref{cdp:eq:C.AD.10}.

For a Hopf vertex the frame at $s=0$ satisfies
\begin{equation}
 \begin{gathered}
 P_i^\infty(0)=J\overline{P_i},\\
 (P_i^\infty(0))^{-1}W_0\Lambda
                         =\overline{\Lambda_i}.
 \end{gathered}                                        \hypertarget{tag-009D}{}\label{cdp:eq:C.AD.11}
\end{equation}
Indeed $\overline{2\pi i}=-2\pi i$, so the two columns built from $n_i^\infty(0),-n_{i-1}^\infty(0)$ are $J\overline{P_i}$. Applying its inverse to $W_0\Lambda=J\overline\Lambda$ gives the second identity. Thus the contracting logarithms and finite angular characters are exactly their complex conjugates, still contracting and free.

At a ruled vertex take the separate frame $(\lambda_i/(2\pi i),n_i/(2\pi i))$ used in \eqref{cdp:eq:C.MX.10}. At infinity choose the lattice lift $\lambda_i^\infty(0)=-W_0(\lambda_i)$ as well as $n_i^\infty(0)=-W_0(n_i)$. Its frame is again $J$ times the conjugate original frame, and its component lattice is the conjugate lattice. Exponentiation therefore identifies its elliptic base, degree-zero line and marked projective completion with their complex conjugates. In either component type the coordinate expression of an edge constant is also conjugated:
\[
 (J\overline{P_i})^{-1}J\overline{\mathbf b}
       =\overline{P_i^{-1}\mathbf b}.
\]
Consequently the elliptic gluing maps are precisely the conjugate specified maps. Apply \cref{lem:II:moving-marked-cycle} to these central conjugate data, with moving lattice $W_s\Lambda$, the indicated outgoing normals and the constants $W_s(\mathbf b_i)$. Its full period quotient is \eqref{cdp:eq:C.AD.10}, and its equation $xy=s$ gives a smooth proper filling whose full central fibre is the conjugate marked support.

Glue the two fillings to the period family on $\C^*$ by these punctured-family biholomorphisms and the frame change $\eta_\infty=\eta/z$. Choose disjoint small discs at zero and infinity and a middle annulus. Over every base point the gluing identifies all pieces with one Hausdorff proper local model. A finite closed refinement of the three base charts has compact inverse images, proving that the total space is Hausdorff and proper over $\PP^1$, hence compact. Its charts are smooth. All middle fibres are the connected two-tori of the period family. Each Hopf component has Euler number zero by \cref{lem:II:hopf-curves} and its sphere-bundle calculation; each ruled elliptic component has $e(\PP^1)e(E)=0$. Each double curve is elliptic, and there are no triple points. The normalization and double-curve decomposition therefore gives Euler number zero for both singular supports, including a single ruled self-identification.

For the asserted SNC Hopf branch, apply \cref{cdp:st:C.2.UN.40} to the given smooth unblown components. Both axes lie in the boundary. On a primary diagonal Hopf surface of algebraic dimension zero they are the only curves by \cite[arXiv v1, \S2.2.6, p.~5]{AbashevaDeev2024AB}: a nonzero integer relation between its multipliers would give a nonconstant invariant meromorphic monomial. For a secondary component, \cref{lem:II:central-primary-hopf-cover} preserves algebraic dimension under its finite primary cover. The inverse image of every curve is therefore a union of those axes; the diagonal deck group preserves each axis and their images remain distinct. Thus each component has exactly its two disjoint boundary curves. The SNC equation supplies exactly one other branch at each. The connected incidence graph has degree two at every vertex and is a cycle, with two parallel edges allowed. Its normal-product maps come from the transverse bilinear term of its full reduced fibre equation.

For two Hopf components multiply their edge matrices:
\begin{equation}
 L_2L_1=
 \begin{pmatrix}
   -c_2&c_2k_1\\
   -k_2&-c_1+k_2k_1
 \end{pmatrix}.
\hypertarget{tag-00LQ}{}\label{cdp:eq:C.HC.11}
\end{equation}
The product is the identity exactly when $c_1=c_2=-1$ and $k_1=k_2=0$. Hence $L_1=L_2=-S_0$, where $S_0$ exchanges the two coordinates, and $\Lambda_2=-S_0\Lambda_1$. Coordinatewise exponentiation gives $\Gamma_2=(S_0\Gamma_1)^{-1}=S_0\Gamma_1$. Thus $(z,w)\mapsto(w,z)$ descends to a biholomorphism of the two Hopf surfaces, including their finite diagonal subgroups. Under this map their elliptic identifications have the form $z\mapsto a/z$ and $w\mapsto b/w$. The diagonal automorphism $(Z,W)\mapsto(a^{-1}Z,b^{-1}W)$ commutes with the complete deck group and turns both maps into inversion. The compact construction already proved applies to this same reduced support; \eqref{cdp:eq:C.HC.9}--\eqref{cdp:eq:C.HC.10} preserve each specified local normal-product constant.
\end{proof}

\subsection{Elliptic suspensions realizing marked cycles}\label{subsec:II:cycle-elliptic-suspensions}

\begin{proposition}[Elliptic suspensions realizing marked cycles]\hypertarget{tag-01DS}{}\label{cdp:st:C.2.UN.48}
\StatementRecord{cdp:st:C.2.UN.48}{proved}{01DS}{\cref{cdp:st:C.2.UN.57,lem:II:moving-marked-cycle,cdp:st:C.2.UN.41,lem:II:central-primary-hopf-cover}; \cite[arXiv v3, Definition~3.1, p.~7; \S7.6, pp.~33--34]{SchuttShioda2010}}
Let $S$ be a marked normal cycle with $M=I$, containing a Hopf component, and suppose that every Hopf component has algebraic dimension zero. The other components may be the marked ruled surfaces of \cref{def:II:marked-normal-cycle}. There is a holomorphic map $f:Y\to\PP^1$ from a smooth compact complex threefold with connected two-torus general fibre and full reduced SNC central fibre isomorphic to the specified $S$. Every reduced fibre support has Euler characteristic zero.

This realization is obtained by replacing a neighbourhood of one fibre in the compact elliptic suspension \eqref{cdp:eq:C.HG.10}. Its punctured periods are \eqref{cdp:eq:C.HG.11}, with the affine matching \eqref{cdp:eq:C.HG.13}--\eqref{cdp:eq:C.HG.14}. It preserves every component, every finite secondary character and every elliptic gluing map, and chooses the normal-product constants. In particular it realizes every compatible all-Hopf cycle of arbitrary length in the stated algebraic-dimension-zero class.

The pencil, Tate and suspension construction used here is defined
for every $q_0,a\in\C^*$ with $0<|q_0|<1$, $a\ne1$, and
every integer $d\ge1$. It gives a smooth compact suspension
\eqref{cdp:eq:C.HG.10} over $\PP^1$, whose elliptic-surface
factor has the selected $I_d$ fibre and a section of central
group coordinate $a$. Its selected punctured periods are
\eqref{cdp:eq:C.HG.11}; its general fibre is a connected
two-torus and every reduced support has Euler number zero.
\end{proposition}
\begin{proof}
Use the common lattice $\Lambda$ and vector $u=d p$ from the proof of \cref{cdp:st:C.2.UN.57}, where $p$ is primitive and $d>0$. The nonvanishing \eqref{cdp:eq:C.HG.5} makes $W/\R p\to\C^2/\C p$ an isomorphism of real planes. Extend $p$ to an integral basis $p,\lambda,c$ of $\Lambda$. Choose a complex linear quotient coordinate $X$ with $X(\lambda)=2\pi i$; the two projected periods are real-linearly independent, so $\operatorname{Re}X(c)\ne0$. Reverse $c$ to make this real part negative. Choose the second complex linear coordinate $Z$ by $Z(p)=2\pi i$ and $Z(\lambda)=0$. The vectors $p,\lambda$ are complex-linearly independent, so this determines a complex linear isomorphism and gives
\begin{equation}
 \Lambda=\left\langle
     (0,2\pi i),\ (2\pi i,0),\
     (\log q_0,\log a)\right\rangle_{\Z},
 \qquad 0<|q_0|<1,\qquad a\in\C^*,
 \qquad u=d(0,2\pi i).
\hypertarget{tag-00MC}{}\label{cdp:eq:C.HG.6}
\hypertarget{tag-00RZ}{}\SsixEquationAliases{\label{cdp:eq:C.MX.8}}
\end{equation}
The logarithms are these chosen coordinates of $c$.

One has $a\ne1$. Otherwise subtract an integral multiple of $(0,2\pi i)$ from $c$ so that its second coordinate is zero. Then
\begin{equation}
 \C/\langle2\pi i,\log q_0\rangle_{\Z}
          \longrightarrow G,\qquad [X]\longmapsto[(X,0)]
\hypertarget{tag-00MD}{}\label{cdp:eq:C.HG.7}
\end{equation}
is well defined and injective: equality of its images forces the coefficient of $(0,2\pi i)$ to vanish, leaving exactly the two periods in its source. Its derivative is injective; compactness makes it an embedded elliptic curve in $G$. This is impossible in the complement of the two axes of one of the Hopf components. For a primary component the nonresonant diagonal curve list in \cite[arXiv v1, \S2.2.6, p.~5]{AbashevaDeev2024AB} gives only those axes. For a secondary component, apply \cref{lem:II:central-primary-hopf-cover} to preserve algebraic dimension on its finite primary cover; the inverse image of any curve is a union of its axes, so again there is no curve in the complement.

The nodal cubic and its prescribed smooth point are
\begin{equation}
 \begin{gathered}
 y^2z=x^2(x+z),\qquad O=[0:1:0],\\
 P_0=\left(\frac{4a}{(1-a)^2},
                  \frac{4a(1+a)}{(1-a)^3}\right)
 \end{gathered}
\hypertarget{tag-00ME}{}\label{cdp:eq:C.HG.8}
\end{equation}
in the affine chart $z=1$. Substitution gives the cubic equation, and its smooth group coordinate $\zeta=(y/x-1)/(y/x+1)$ has value $a$. Since $a\in\C^*\setminus\{1\}$, this point is distinct from $O$ and from the node. Apply the nodal-cubic pencil construction in the proof of \cref{cdp:st:C.2.UN.41} to this pair $O,P_0$. Its restriction and first-jet argument uses exactly these two distinct smooth points: it gives a second cubic transverse at nine distinct smooth base points and missing the node. Blowing up those points yields the smooth compact rational elliptic surface with sections $O,P$ and this selected $I_1$ fibre.

Make the global base change $\PP^1_s\to\PP^1$ given by $s\mapsto s^d$ at the selected value. Normalize, resolve and take its smooth relatively minimal elliptic model $J\to\PP^1_s$. The equation at the selected node before resolution is $xy=s^d$. Apply the resolution in
\cref{lem:II:moving-marked-cycle} with $h_i=d$, $t=s$ and
the same $x,y$. Its chart map
\eqref{eq:II:compressed-neck-resolution}, with
$(a_j,b_j)=(u_j,v_j)$, is
$x=u_j^{j+1}v_j^j$, $y=u_j^{d-j-1}v_j^{d-j}$ and $s=u_jv_j$.
It inserts the same chain of $d-1$ projective lines between the
two node branches. Together with the original component its
reduced central equations give the $I_d$ fibre. The two sections avoid the node and retain their values in its identity component. The other operations are proper vertical modifications of the projective base change. The resulting $J$ is smooth, projective and relatively minimal over $\PP^1_s$. Apply \cite[arXiv v3, Definition~3.1, p.~7; \S7.6, pp.~33--34]{SchuttShioda2010} to this $J$ and the section $P$: translation by $P$ extends to an automorphism over the base, with inverse translation by $-P$.

Pull back the exact nodal Tate comparison in the proof of \cref{cdp:st:C.2.UN.41} along the same $s\mapsto s^d$. Its parameter $Q_1$ has a simple zero, so $Q(s)=Q_1(s^d)$ has order $d$. Its smooth central group-coordinate comparison is the coordinate $a$ above; the inverse function theorem at this point gives the section lift in
\begin{equation}
       J_s=\C^*/Q(s)^{\Z},\qquad
       \operatorname{ord}_0Q=d,\qquad P_s=[a(s)],
       \qquad a(0)=a.
\hypertarget{tag-00MF}{}\label{cdp:eq:C.HG.9}
\end{equation}
Choose a holomorphic $d$th root of the unit $Q(s)/s^d$ and put $t=s(Q(s)/s^d)^{1/d}$. This coordinate has nonzero derivative at zero and gives $Q=t^d$ exactly. Write the section lift as $a(t)$ and choose its holomorphic logarithm with value $\log a$.

The construction from the nodal cubic onward uses only
$0<|q_0|<1$, $a\ne0,1$ and $d\ge1$. It therefore also
proves the stated arbitrary-parameter clause. Form the suspension
\begin{equation}
 Y^{\rm std}=
  (\C^*_b\times J)/
      \langle(b,y)\longmapsto(q_0 b,y+P)\rangle
          \longrightarrow\PP^1_s .
\hypertarget{tag-00M0}{}\label{cdp:eq:C.HG.10}
\end{equation}
The action is free and properly discontinuous because its first coordinate is multiplied by $q_0$, with $0<|q_0|<1$. A closed annulus in $\C^*$ times the compact $J$ meets every orbit, so the quotient is compact and smooth. Its general fibre is a connected two-torus: in additive coordinates the elliptic-fibre lattice and a logarithm of its section translation give a rank-four lattice, whose remaining real direction is supplied by $\log|q_0|\ne0$. At the distinguished value its punctured periods are exactly
\begin{equation}
 \Lambda(t)+\Z(0,d\log t),\qquad
 \Lambda(t)=
 \left\langle(0,2\pi i),(2\pi i,0),
                  (\log q_0,\log a(t))\right\rangle_{\Z}.
\hypertarget{tag-00M1}{}\label{cdp:eq:C.HG.11}
\hypertarget{tag-00S0}{}\SsixEquationAliases{\label{cdp:eq:C.MX.9}}
\end{equation}
The map from this additive quotient to the suspension sends $[(X,Z)]$ to $[(e^X,[e^Z])]$. The two angular periods exponentiate to the identity, $(\log q_0,\log a(t))$ is its suspension generator, and $(0,d\log t)$ is the elliptic-fibre generator. Local logarithms give its inverse modulo these same generators.

Transport every $n_i$ and every chosen contracting period by its fixed integer coefficients in the three generators of \eqref{cdp:eq:C.HG.6}, replacing that basis by the basis of $\Lambda(t)$ in \eqref{cdp:eq:C.HG.11}. In particular $\sum_i n_i(t)=-d(0,2\pi i)$. Apply the moving clause of \cref{lem:II:moving-marked-cycle} to this lattice. Its Hopf frames are \eqref{cdp:eq:C.HG.12}, with all finite angular relations unchanged; its ruled frames are \eqref{cdp:eq:C.MX.10}. Thus both sorts of marked component have their prescribed central surfaces.

Choose the affine edge lifts holomorphically, preserving the constants of each central elliptic gluing map. In the all-Hopf case write those fixed constants as $d_i$ and the adjustable normal-product logarithms as $e_i(t)$. Their sum is
\begin{equation}
 \sum_iP_{i+1}(t)
        \binom{d_i}{e_i(t)+\log t}
       =b_0(t)+(0,d\log t).
\hypertarget{tag-00M3}{}\label{cdp:eq:C.HG.13}
\end{equation}
For a ruled edge use the base and line frame \eqref{cdp:eq:C.MX.10}; its adjustable normal constant has the same common coefficient $-n_i(t)/(2\pi i)$. Thus in the mixed case the full sum is likewise $b_0(t)+(0,d\log t)$. Choose a Hopf vertex $i+1$. The coefficients $-n_i(t)/(2\pi i)$ and $-n_{i+1}(t)/(2\pi i)$ of its two adjacent normal-product logarithms are complex-linearly independent near zero. Fix the other constants and solve this invertible two-by-two holomorphic linear system to obtain
\begin{equation}
                              b_0(t)=0.
\hypertarget{tag-00M4}{}\label{cdp:eq:C.HG.14}
\end{equation}
These changes lie in the normal directions and preserve all elliptic gluing maps.

Now apply \cref{lem:II:moving-marked-cycle} with these adjusted affine constants. Its full period formula has constant sum zero and normal sum $d(0,2\pi i)$, so it is exactly \eqref{cdp:eq:C.HG.11}. Its maps send each core logarithm to the common vector space and subtract the accumulated edge constants; on the last overlap the difference is $(0,d\log t)$, and logarithm changes give $\Lambda(t)$. Composing the lemma's punctured-family biholomorphism with the exponential map above identifies the entire local filling with \eqref{cdp:eq:C.HG.10} over the punctured disc, including every affine parameter. Its full central divisor is the prescribed reduced cycle.

Replace the distinguished tube in $Y^{\rm std}$ by this filling. Remove a closed smaller base disc from the suspension and use the filling over a larger disc; glue by the displayed isomorphism on the intervening annulus. At either annular boundary only one piece contains the limiting base point, giving a Hausdorff gluing. Closed smaller pieces in the two proper families cover the result and are compact. This proves compactness, while the same charts preserve smoothness and the holomorphic map to $\PP^1_s$. All fibres are connected.

Every component of the new central support, and every double elliptic curve, has Euler number zero, so its support does also. At any other value the support is unchanged from the suspension. Its projection to $\C^*/q_0^{\Z}$ is a locally trivial bundle with fibre $J_s^{\rm red}$ and transition the automorphism induced by $P$. Its finite cellular complex with these local coefficients has Euler number $e(\C^*/q_0^{\Z})e(J_s^{\rm red})=0$. This gives the asserted Euler number at every value.
\end{proof}

\subsection{Explicit Hopf and ruled cycles}\label{subsec:II:explicit-marked-cycles}

\begin{example}[Three Hopf components and ruled insertions]\hypertarget{tag-P332}{}\label{ex:II:three-hopf-cycle}
\StatementRecord{ex:II:three-hopf-cycle}{proved}{P332}{\cref{cdp:st:C.2.UN.48,cdp:st:C.2.UN.57,cdp:st:C.2.UN.41}}
There is a compact realization with the three primary Hopf contractions and edge matrices
\begin{equation}
 \begin{gathered}
 (e^{-1},e^{-\sqrt2}),\qquad
 (e^{-\sqrt2},e^{-1-\sqrt2}),\qquad
 (e^{-1-\sqrt2},e^{-1}),\\
 L_1=\begin{pmatrix}0&-1\\-1&-1\end{pmatrix},\qquad
 L_2=\begin{pmatrix}0&1\\-1&1\end{pmatrix},\qquad
 L_3=\begin{pmatrix}0&-1\\-1&1\end{pmatrix}.
 \end{gathered}
 \hypertarget{tag-00LZ}{}\label{cdp:eq:C.HG.1}
\end{equation}
The matrix product is the identity, and the full multiplicative cycle period is $(t^{-2},t^2)$. The construction also gives a secondary three-cycle with the finite weights $(1,2),(3,2),(2,4)$ modulo $5$ on these respective components. In both cases the positive divisor of the total normal period in the common lattice is $2$.

For a repeated two-component inversion cycle of length $2k$ with free finite weights $(a,b)$ modulo $n$, so that $\gcd(a,n)=\gcd(b,n)=1$, the positive normal-period divisor is
\begin{equation}
                    d=k\gcd(n,a+b).
\hypertarget{tag-00M7}{}\label{cdp:eq:C.HG.17}
\end{equation}
In particular the weights $(1,4)$ modulo $5$ give divisor $5$ for two components.

For every pair of integers $a,b\ge1$ there is also a marked cycle with two Hopf components and $a+b-2$ ruled components. Its positive normal-period divisor is $\gcd(a,b)$. These cycles have the compact realizations of \cref{cdp:st:C.2.UN.48}.
\end{example}
\begin{proof}
For the three-component example use the common lattice generated by $2\pi i\Z^2$ and the following $c$, with outgoing meridians
\begin{equation}
 n_1=2\pi i(1,0),\quad
 n_2=2\pi i(1,-1),\quad
 n_3=2\pi i(0,-1),\quad
 c=(-1,-\sqrt2).
\hypertarget{tag-00M5}{}\label{cdp:eq:C.HG.15}
\end{equation}
The frames \eqref{cdp:eq:C.HG.3} give the contracting logarithms $P_1^{-1}c$, $-P_2^{-1}c$ and $-P_3^{-1}c$, which are exactly the three pairs in \eqref{cdp:eq:C.HG.1}. Their ratios are irrational, hence their multipliers are multiplicatively independent; the curve calculation in \cref{cdp:st:C.2.UN.41} gives algebraic dimension zero. Matrix multiplication gives $L_3L_2L_1=I$. The corresponding multiplicative maps, with normal-product constants one, are
\begin{equation}
 \begin{gathered}
 (z,w)\mapsto(w^{-1},t z^{-1}w^{-1}),\\
 (z,w)\mapsto(w,t z^{-1}w),\\
 (z,w)\mapsto(w^{-1},t z^{-1}w).
 \end{gathered}
\hypertarget{tag-00M6}{}\label{cdp:eq:C.HG.16}
\end{equation}
Their composite is $(z,w)\mapsto(t^{-2}z,t^2w)$. Thus $u=2\pi i(-2,2)$, whose divisor in the angular lattice is $2$. A nonzero multiple of $c$ has nonzero real part, so it does not enlarge the lattice on this angular line. The marked-cycle criterion and \cref{cdp:st:C.2.UN.48} give its compact realization.

For the finite repeated cycle, in units of $2\pi i$ the angular lattice is $\Z^2+\Z(a,b)/n$ and its total normal period is $k(-1,1)$. A representative $j(a,b)/n$ lies on this line modulo $\Z^2$ precisely when $n\mid j(a+b)$. Put $g=\gcd(n,a+b)$. Then $j$ is a multiple of $n/g$, and the fractions in the first coordinate are the multiples of $a/g$ modulo integers. Since $\gcd(a,g)=1$, they give every multiple of $1/g$. Hence the intersection with the line is $\Z(-1,1)/g$, proving \eqref{cdp:eq:C.HG.17}. The sign follows from the outgoing vectors $(1,0)$ and $(0,-1)$.

In the three-component example enlarge the common angular lattice by $(1,2)/5$. Applying the three inverse frames gives weights $(1,2),(3,2),(2,4)$ modulo $5$. Each coordinate is coprime to $5$, so each finite action is free on both nonzero axes and on their complement. The preceding line calculation has $g=\gcd(5,1+2)=1$ and leaves the divisor of $(-2,2)$ equal to $2$. The same period and normal-product construction therefore gives the asserted secondary compact cycle.

For the mixed examples choose real numbers $A,B<0$ with irrational ratio and use the common lattice
\[
 \Lambda=\langle 2\pi i(1,0),2\pi i(0,1),(A,B)\rangle_{\Z}.
\]
Take the cyclic list of outgoing normal vectors to be $a$ copies of $2\pi i(1,0)$ followed by $b$ copies of $2\pi i(0,-1)$. At each of its two changes, the frame with columns $n_i/(2\pi i),-n_{i-1}/(2\pi i)$ gives a Hopf component; choose the contracting generator $(A,B)$ at the first change and its negative at the second. Their contracting logarithms are $(A,B)$ and $(B,A)$, respectively, so both have algebraic dimension zero. Each repeated normal gives the ruled component constructed from the quotient along its normal line and the lattice of \eqref{cdp:eq:C.MX.10}. The remaining projected real period is nonzero because $A,B<0$, so that quotient is elliptic. This defines all $a+b$ marked components and their normal lines.

Choose the angular frames unimodularly and take each edge linear map to be $P_{i+1}^{-1}P_i$, with affine constants zero. The product telescopes to the identity and identifies each outgoing normal with minus the next incoming normal. Exponentiation gives the elliptic identifications and dual-normal maps of the defined marked cycle. Its normal sum is $u=2\pi i(-a,b)$ and its primitive divisor is $\gcd(a,b)$; the real period $(A,B)$ cannot contribute to a purely angular vector.

For $0<|t|<1$, put $\ell_t=\log|t|$. In the real coordinate order
$(\operatorname{Im}z/(2\pi),\operatorname{Im}w/(2\pi),\operatorname{Re}z,\operatorname{Re}w)$,
the first two angular periods are the first two unit vectors. The determinant of the four periods is therefore exactly
\[
 \det\begin{pmatrix}A&-a\ell_t\\ B&b\ell_t\end{pmatrix}
       =\ell_t(bA+aB)>0.
\]
The normalization of these real coordinates accounts for the two factors $2\pi$ in the angular periods. The constant cycle translation is zero in this example. In an all-ruled cycle, changing normal-product constants adds only vectors in the common normal line, so the projection of the cycle translation to its elliptic quotient remains the prescribed one, as proved in \cref{cdp:st:C.2.UN.57}. The compact realization of every mixed example follows from \cref{cdp:st:C.2.UN.48} with these same components and edge maps.
\end{proof}

\section{Elliptic stars and compact insertion}\label{sec:II:elliptic-stars}

\subsection{Elliptic star data and inverse transforms}\label{subsec:II:elliptic-star-data}\label{subsec:II:star-weighted-constructions}

\begin{definition}[Marked elliptic stars and inverse transforms]\hypertarget{tag-P350}{}\label{def:II:elliptic-star-data}
\StatementRecord{def:II:elliptic-star-data}{proved}{P350}{\cref{lem:II:identity-monodromy-elliptic}}
The leaf pairs are
\begin{equation}
 \begin{gathered}
 (H^{\rm diag}_{r_i}(\beta_i)/\mu_{n_i},J_i),\qquad
 (H^{\rm nl}_{r_i}(\beta_i)/\mu_{n_i},J_i),\\
 r_i\ge1,\qquad 0<|\beta_i|<1,\qquad \gcd(r_i,n_i)=1 ,
 \end{gathered}                                             \hypertarget{tag-00AG}{}\label{cdp:eq:C.AS.1}
\hypertarget{tag-0122}{}
\SsixEquationAliases{\label{cdp:eq:C.SN.1}}
\end{equation}
with marked axis $J_i=\{w=0\}$ and contractions and finite action
\begin{equation}
 (z,w)\longmapsto(\beta_i^{r_i}z,\beta_iw),\qquad
 (z,w)\longmapsto(\beta_i^{r_i}z+w^{r_i},\beta_iw),
 \quad
 \zeta(z,w)=(\zeta^{r_i}z,\zeta w),                         \hypertarget{tag-00AR}{}\label{cdp:eq:C.AS.2}
\end{equation}
A marked elliptic star consists of these pairs, a smooth compact
complex surface with a connected elliptic quasibundle
\begin{equation}
                         q:P\longrightarrow B             \hypertarget{tag-00AS}{}\label{cdp:eq:C.AS.3}
\end{equation}
over a smooth compact connected curve $B$, and pairwise distinct
marked fibres with full divisors $r_iJ_i$. Here a quasibundle has
only positive multiples of smooth elliptic curves as fibres.
Its integral $H_1$ monodromy is the identity, and its smooth
fibres are torsors under one fixed elliptic curve $F$.
The data include the isomorphism from each marked leaf axis to
the corresponding $J_i\subset P$ and a duality between their
normal lines. These identifications form the reduced surface
$S=P\cup H_1\cup\cdots\cup H_n$, with $n\ge1$, disjoint
double curves and no other intersections.

For one marked fibre suppress the index and choose integers $k,l$ with
\begin{equation}
                         rk-nl=1,\qquad \gamma=\beta^n.
                                                        \hypertarget{tag-0125}{}\label{cdp:eq:C.SN.3}
\end{equation}
The paired neighbourhood in $P$, in the axis coordinate $z$
of the leaf cover, has the invariant coordinates
\begin{equation}
 u=z^n,\qquad X=xz^k,\qquad
 p=X^ru^{-l},\qquad
 (u,X)\sim(\gamma^ru,\gamma^lX).
                                                        \hypertarget{tag-0126}{}\label{cdp:eq:C.SN.4}
\hypertarget{tag-00TA}{}
\SsixEquationAliases{\label{cdp:eq:C.NS.1}}
\end{equation}
Here $u\ne0$. Its punctured product map and inverse are
\begin{equation}
 (u,X)\longmapsto(p,W)
       =(X^ru^{-l},\,X^{-n}u^k),\qquad
 u=p^nW^r,\quad X=p^kW^l,\qquad
                       F=\C^*/\gamma^{\Z}.               \hypertarget{tag-0127}{}\label{cdp:eq:C.SN.5}
\hypertarget{tag-00TL}{}\hypertarget{tag-00VV}{}\hypertarget{tag-017U}{}
\SsixEquationAliases{\label{cdp:eq:C.NS.2}\label{cdp:eq:C.PR.6}\label{cdp:eq:C.WT.15}}
\end{equation}
Replace each marked neighbourhood by the product $F\times\Delta_p$
through this punctured isomorphism, keeping every transition map
of $P$ away from the marked values. The resulting map is the
\emph{inverse logarithmic transform}
\begin{equation}
                             q_0:T_0\longrightarrow B.
                                                        \hypertarget{tag-00TN}{}\label{cdp:eq:C.NS.3}
\end{equation}
If there are no unmarked multiple fibres, it is a holomorphic
principal $F$-torsor. Its integral Chern vector is the connecting
class $c(T_0)\in H^2(B,\Lambda_F)$ of the sheaf sequence in
\eqref{cdp:eq:C.KX.13}, with $F_1=F$.
The transform preserves the full analytic torsor transition maps
and is independent, up to an isomorphism over $B$, of the
compatible local coordinates used in its construction.
\end{definition}
\begin{proof}
Every vertical prime of a quasibundle is elliptic, so it has
no vertical exceptional curve of the first kind and is relatively
minimal. Apply \cref{lem:II:identity-monodromy-elliptic} to this
same quasibundle $q:P\to B$ and full fibre $mJ$.
Its chart \eqref{cdp:eq:C.HT.4} gives the fixed-translation
quotient, exact normal order $m$, and the connected
unit-frame cover with its primitive character. These are
the surface germ, root map and normal line used below.

For the marked Hopf pair, the finite quotient acts on the axis
and a transverse coordinate by $(z,x)\mapsto
(\zeta^r z,\zeta^{-1}x)$ on the paired side.
The invariant monomials are $u=z^n$, $X=xz^k$, because
$rk-1=nl$. They have the deck factors in
\eqref{cdp:eq:C.SN.4}, and $p=x^r z=X^r u^{-l}$.
Conversely, every choice of an $n$th root $z$ of $u$ gives
$x=Xz^{-k}$; the $n$ choices are exactly the finite orbits.
This proves the coordinate isomorphism, including its inverse.
The marked curve is $\C^*/\gamma^{r\Z}$, and its paired
normal has factor $\gamma^l$. On the normal bundle of the Hopf axis, the invariant
fibre coordinate is $Y=wz^{-k}$, with factor $\gamma^{-l}$.
The relation $rk-nl=1$ gives $\gcd(l,r)=1$, so these dual
normal lines have exact order $r$.

In the normalized root write $u=y^r$, $X=s y^l$;
then $p=s^r$. The deck action from the root is
$(y,s)\mapsto(\xi y,\xi^{-l}s)$, $\xi^r=1$,
and the infinite period is $y\mapsto\gamma y$.
Writing its positive normal character as $s\mapsto\zeta s$
gives $y\mapsto\zeta^n y$, since $-nl\equiv1\pmod r$.
The unit-frame cover of the dual normal line is therefore
precisely the cyclic elliptic cover of the preceding paragraph.
The given axis and normal-line maps identify the two quotients,
so this is a chart of the given quasibundle germ.

The exponent matrix of \eqref{cdp:eq:C.SN.5} is
$\left(\begin{smallmatrix}-l&r\\k&-n\end{smallmatrix}\right)$,
with determinant $nl-rk=-1$. Direct substitution gives
$u=p^nW^r$, $X=p^kW^l$, so the maps are mutually inverse.
The logarithmic derivative of the same monomials gives
\begin{equation}
\begin{gathered}
 d\log p\wedge d\log W=d\log X\wedge d\log u,\\
 dp\wedge d\log W=X^{r-1}u^{-l-1}\,dX\wedge du.
\end{gathered}
\hypertarget{tag-017V}{}\label{cdp:eq:C.WT.16}
\end{equation}
Indeed the coefficient is $rk-nl=1$. Thus a form of boundary
order $r-1$ in this chart becomes the nonzero product frame
$dp\wedge d\log W$ through the same inverse map.
The infinite deck map sends $W$ to $\gamma W$ and fixes $p$.
They are consequently product isomorphisms over the punctured
disc. For a primary leaf, $n=1,k=0,l=-1$, these are exactly
the coordinates and product maps in
\cref{cdp:eq:C.NS.1,cdp:eq:C.NS.2}.

The product pieces glue to the unchanged outside of $P$ by
these isomorphisms. A finite cover of the compact base by
smaller discs and their compact complement gives compact
inverse images, because every local product and every
unchanged quasibundle piece is proper. The resulting map
is therefore proper. It is smooth at the replaced fibres,
and all its remaining fibre models are those of $q$.
If they are smooth, its overlap maps preserve the integral
marking and hence have identity linear part; it is the
stated principal torsor.

Changing the integers to $k'=k+na$, $l'=l+ra$ changes
$X$ to $X'=Xu^a$ and leaves $p$ and $W$ unchanged in
\eqref{cdp:eq:C.SN.5}. More generally, two compatible root
trivializations differ by a holomorphic translation in the
elliptic product and a constant origin-preserving linear map.
The marked axis and normal cover identify that linear map.
The remaining translation is invariant under the root group,
so descends to a holomorphic translation on the ordinary
$p$-disc. The same map extends across its centre in the
replacement product. Ratios of two normal dualities are
nowhere-zero holomorphic functions on the compact axis,
thus nonzero constants, and give the same extending
translation comparison. These maps glue with the unchanged
outside transitions and have inverse translations. They
prove the asserted choice independence while preserving the
analytic torsor class.
\end{proof}


\subsection{Boundary maps and periods of elliptic stars}\label{subsec:II:star-boundary-periods}

\begin{lemma}[Boundary maps and canonical periods of elliptic stars]\hypertarget{tag-P351}{}\label{lem:II:elliptic-star-boundary}
\StatementRecord{lem:II:elliptic-star-boundary}{proved}{P351}{\cref{def:II:elliptic-star-data,lem:II:normal-product-collars,cdp:st:C.2.UN.17}}
Every leaf pair in \cref{cdp:eq:C.AS.1,cdp:eq:C.AS.2} has
$K_H\simeq\OO_H(-(r+1)J)$, and $\OO_H(J)$ has infinite order.
For a full reduced SNC torus degeneration, let one
component be such a leaf, with its only double curve $J$
on a neighbouring smooth component $P$. The difference
of the canonical evaluation orders and the individual
core expansions are
\cref{cdp:eq:C.NS.8,cdp:eq:C.NS.11,cdp:eq:C.AS.15,cdp:eq:C.AS.16},
with $k_P$ the order of this neighbour. These conclusions
use only that leaf and its double curve; the rest of the
support may contain other components and intersections.
In these expansions $\gamma_{\rm core}$ is the oriented
core of this leaf, and the relative evaluation is divided
by $t^{k_P}$.

For a full star $P\cup_iH_i$ with arbitrary smooth main
component $P$, the main canonical divisor is
\eqref{cdp:eq:C.NS.9}. The global comparison with the
inverse transform and the common core class below use
the following quasibundle and star hypotheses.
Let $q:P\to B$ be a fixed-Jacobian elliptic quasibundle with
identity integral monodromy, and identify the axes of leaf
pairs \eqref{cdp:eq:C.AS.1} with distinct reduced fibres $J_i$.
Suppose their normal lines are dual. The marked full fibres
are $r_iJ_i$, so these data define the inverse transform of
\cref{def:II:elliptic-star-data}.

For a marked fibre use $r,n,k,l,\gamma$ and $(u,X,p,W)$ from
that definition. Let $a$ be its positive angular $u$-loop,
$b$ its deck-period loop, and $\mu_P$ its positive normal
$X$-loop. Let $\delta$ be the positive $p$-loop at fixed $W$,
and let $e,h$ be the angular and deck-period loops of
$F=\C^*/\gamma^{\Z}$. Their integral comparison is
\begin{equation}
 \begin{pmatrix}a&b&\mu_P\end{pmatrix}_{(\delta,e,h)}
   =
 \begin{pmatrix}
   -l&0&r\\
    k&0&-n\\
    0&1&0
 \end{pmatrix},
 \qquad
 \delta=n a+k\mu_P,\quad e=r a+l\mu_P,\quad h=b.
                                                        \hypertarget{tag-0128}{}\label{cdp:eq:C.SN.6}
\hypertarget{tag-00TO}{}
\SsixEquationAliases{\label{cdp:eq:C.NS.4}}
\end{equation}
In particular $\mu_P=r\delta-ne$ is primitive.
The Hopf complement fills the meridian
$\delta_H=na-k\mu_H$ and carries the fibre loops
$e_H=ra-l\mu_H$, $h_H=b$. The normal-product identification
$\mu_H=-\mu_P$ sends them to $\delta,e,h$ respectively.
For the primary choice $n=1,k=0,l=-1$, this is
\begin{equation}
                   e_H=r_i a+\mu_H=e,\qquad h_H=b=h.
                                                        \hypertarget{tag-00TQ}{}\label{cdp:eq:C.NS.6}
\end{equation}
When there are no unmarked multiple fibres, choose a
trivialization over the complement of small marked
base discs, write $\delta_i=\delta_{\rm ext}+\nu_i$,
where $\nu_i\in\Lambda_F$ is its translation winding.
The integral Chern vector and meridians are
\begin{equation}
       \mu_{P,i}=r_i\delta_{\rm ext}
                         +(r_i\nu_i-n_i e_i),\qquad
                         c(T_0)=\sum_i\nu_i.             \hypertarget{tag-0124}{}\label{cdp:eq:C.SN.2}
\hypertarget{tag-00TR}{}
\SsixEquationAliases{\label{cdp:eq:C.NS.7}}
\end{equation}
Thus the rational Seifert vector before replacement is
$\sum_i(\nu_i-n_ie_i/r_i)$, and the replacements add
$\sum_i n_ie_i/r_i$ while preserving all holomorphic
transition maps.

The canonical line and direct image of this $q$ are
\eqref{eq:II:quasibundle-canonical} in
\cref{lem:II:identity-monodromy-elliptic}, with
$(E_0,C_b)=(P,J_b)$.
If $S=P\cup_iH_i$ is the full reduced fibre of a proper
two-torus degeneration from a smooth threefold, a nearby
fibre is smoothly identified with $T_0$. Under this
identification every oriented Hopf core two-torus is the
ordinary fibre class $\gamma_{\rm core}$ of $q_0$. For a canonical
evaluation generator with orders $k_P,k_i$ one has
\begin{equation}
 (k_P-k_i)J_i=K_{H_i}+J_i=-r_iJ_i,
       \qquad k_i-k_P=r_i.
                                                        \hypertarget{tag-00TS}{}\label{cdp:eq:C.NS.8}
\hypertarget{tag-00AJ}{}
\SsixEquationAliases{\label{cdp:eq:C.AS.12}}
\end{equation}
On the main component,
\begin{equation}
                         K_P=\OO_P\Bigl(\sum_i(r_i-1)J_i\Bigr).
                                                        \hypertarget{tag-00TT}{}\label{cdp:eq:C.NS.9}
\hypertarget{tag-00AK}{}
\SsixEquationAliases{\label{cdp:eq:C.AS.13}}
\end{equation}
After division by $t^{k_P}$, the leading relative two-form
on the $i$th Hopf core is $t^{r_i}\eta_i$, where
\begin{equation}
 \eta_i=c_i\,\frac{dz\wedge dw}{w^{r_i+1}},
                         \qquad c_i\ne0.                \hypertarget{tag-00TC}{}\label{cdp:eq:C.NS.11}
\hypertarget{tag-00AL}{}
\SsixEquationAliases{\label{cdp:eq:C.AS.14}}
\end{equation}
For a nonlinear leaf its period is
\begin{equation}
       \int_{\gamma_{\rm core}}\omega_t
          =t^{r_i}\left(\frac{\pm2\pi i c_i\beta_i^{-r_i}}{n_i}
                         +O(t)\right).                 \hypertarget{tag-00AM}{}\label{cdp:eq:C.AS.15}
\hypertarget{tag-00TD}{}
\SsixEquationAliases{\label{cdp:eq:C.NS.12}}
\end{equation}
For a diagonal leaf it satisfies
\begin{equation}
                         \int_{\gamma_{\rm core}}\omega_t=O(t^{r_i+1})
                  \quad\hbox{at a diagonal leaf}.        \hypertarget{tag-00AN}{}\label{cdp:eq:C.AS.16}
\end{equation}
In the full star case with this quasibundle $P$, these are
expansions of its one common holomorphic core period.

For any principal elliptic torsor $q_0:T_0\to B$ over an
elliptic curve, its fundamental group has the exact sequence
and presentation
\begin{equation}
\begin{gathered}
 1\longrightarrow\Lambda_F\longrightarrow\pi_1(T_0)
                    \longrightarrow\Lambda_B\longrightarrow1,\\
 \pi_1(T_0)=\langle a_B,b_B,e,h:\ e,h\text{ central},\
             [a_B,b_B]=e^{c_1}h^{c_2}\rangle,\qquad
                         (c_1,c_2)=c(T_0).
\end{gathered}
\hypertarget{tag-00TE}{}\label{cdp:eq:C.NS.13}
\hypertarget{tag-00W2}{}\SsixEquationAliases{\label{cdp:eq:C.PT.12}}
\end{equation}
The fibre subgroup $\langle e,h\rangle\simeq\Lambda_F$
is injected. Consequently, in a two-torus smoothing its
Chern vector is zero.
\end{lemma}
\begin{proof}
Apply \cref{lem:II:identity-monodromy-elliptic} to the
given quasibundle $q$ and marked full fibre $mJ$.
Its connected root and primitive normal character give
normal order exactly $m$.
The dual Hopf normal has exact order $r$, including its
finite character $\gamma^{-l}$. Hence $m=r$ and the local
maps of \cref{def:II:elliptic-star-data} apply to this quasibundle.

Apply the angular and deck loops to the monomials
$p=X^ru^{-l}$ and $W=X^{-n}u^k$.
The resulting columns are $(-l,k,0)$, $(0,0,1)$,
and $(r,-n,0)$, giving \eqref{cdp:eq:C.SN.6}.
Its determinant is $rk-nl=1$.
The identities for $\delta,e,h$ follow by multiplying
these columns by the displayed integral coefficients.
They also give $\mu_P=r\delta-ne$ and its primitivity.

On the diagonal reference Hopf pair put $Y=wz^{-k}$ and $u=z^n$.
Then $Y^r u^l=w^r/z$ and $Y^n u^k=w^n$.
Use $U=z/w^r$ as the positive complementary base coordinate;
it is the reciprocal of $Y^r u^l$.
The loops $na-k\mu_H$ and $ra-l\mu_H$ therefore change
$(\arg U,\arg(w^n))$ by $(2\pi,0)$ and $(0,2\pi)$,
respectively. The deck loop is the period $w^n\mapsto
\gamma w^n$. The normal product is $XY=xw=t$, so its
positive normal meridians are opposite. Substitution gives
$\delta_H=\delta$, $e_H=e$ and $h_H=h$, including
\eqref{cdp:eq:C.NS.6} in the primary case.

For a diagonal leaf its complement is
$\C_U\times(\C^*/\langle\beta,\mu_n\rangle)$, and
$[w]\mapsto[w^n]$ identifies the second factor with $F$.
For a nonlinear leaf the smooth coordinates
\begin{equation}
 U=\frac{z}{w^{r_i}},\qquad
 V=U-\frac{\beta_i^{-r_i}}{\log|\beta_i|}\log|w|.
                                                        \hypertarget{tag-00TP}{}\label{cdp:eq:C.NS.5}
\end{equation}
make $V$ invariant under the infinite deck map and under
the finite action. They give the same smooth product
$\C_V\times F$. Its compact core fills the complementary
disc and keeps the two fibre circles just calculated.
The normal framing agrees with the diagonal reference
through a smooth comparison of pairs: vary the resonant
coefficient in $(\beta^r z+c w^r,\beta w)$ from $c=1$ to
$c=0$. Its iterates have a uniform polynomial factor times
a contracting exponential, so one compact annulus meets
every orbit throughout this interval. The quotient is a
smooth proper family of pairs with the axis and its normal
characters fixed. A horizontal vector field tangent to the
axis, extended by a partition of unity, has a complete
flow over the compact parameter interval. Its flow gives
the pair diffeomorphism and fixes the displayed peripheral
maps. The product formula above identifies the fibre
circles in that comparison.

Apply \cref{lem:II:normal-product-collars} to the given
full reduced SNC star. Its double curves have two embedded
branches, its normal-product line maps are the given ones,
and there are no triple points. Its compatible core and
collar diffeomorphisms identify the nearby fibre with the
union of the $P$ core and these Hopf fillings. Each Hopf
filling kills $\delta$ and keeps $e,h$, exactly as the
product inserted by \eqref{cdp:eq:C.SN.5}. Gluing these
maps and their inverses gives the claimed diffeomorphism
to $T_0$. Along any path in $B$ joining two marked values
and avoiding the remaining multiple values, the ordinary
elliptic fibres are homologous. Thus all oriented Hopf
cores give the same integral class $\gamma_{\rm core}$.

When $q_0$ is smooth, its translation trivialization on the punctured base gives
$\delta_i=\delta_{\rm ext}+\nu_i$. Substitution into
$\mu_{P,i}=r_i\delta_i-n_ie_i$ gives
\eqref{cdp:eq:C.SN.2}. The obstruction to extending that
trivialization over all base discs is the sum of their
winding vectors, which is the connecting Chern class of
\eqref{cdp:eq:C.KX.13}. This proves both the integral
and rational Seifert comparisons. The full punctured product
maps also preserve the analytic transitions in the preceding
diffeomorphism.

The canonical isomorphism and direct image are the
application of \cref{lem:II:identity-monodromy-elliptic}
to this same $q:P\to B$, its marked and unmarked full
fibres, and its fixed integral elliptic marking.

For any one of the stated leaf occurrences, apply the evaluation construction of
\cref{cdp:st:C.2.UN.17} with all full-fibre multiplicities
equal to one. Its principal fibre identity and adjunction
give $(k_P-k_i)J_i=K_{H_i}+J_i$. The form
$dz\wedge dw/w^{r_i+1}$ is invariant under both deck
groups, for both contractions in \eqref{cdp:eq:C.AS.2}.
It gives $K_{H_i}=-(r_i+1)J_i$.
The line $\OO_{H_i}(J_i)$ has infinite order: a trivial
positive power would make its canonical section a
nonconstant holomorphic function on a compact surface.
Thus the coefficient equality is precisely
\eqref{cdp:eq:C.NS.8}. When the entire support is the star, adjunction on its main
component $P$ gives \eqref{cdp:eq:C.NS.9}.

In a product chart $t=xw$, with $H_i:x=0$ and $P:w=0$,
the total canonical generator is
$x^{k_i}w^{k_P}g\,dz\wedge dw\wedge dx$, with $g$ a
holomorphic unit. Dividing its relative residue by
$t^{k_P}$ gives
$t^{r_i}g\,dz\wedge dw/w^{r_i+1}$.
On the compact leaf, the leading coefficient is a
nowhere-zero section of $K_{H_i}((r_i+1)J_i)$, whose
invariant trivialization just displayed makes it a
nonzero constant. This proves \eqref{cdp:eq:C.NS.11}.
In the logarithmic variables $(U,\log w)$ the nonlinear
core has primitive periods
\begin{equation}
             (0,2\pi i/n),\qquad
             (\beta^{-r},\log\beta).
                                                        \hypertarget{tag-0129}{}\label{cdp:eq:C.SN.7}
\end{equation}
and its form is $c_i\,dU\wedge d\log w$.
Their determinant is nonzero and gives
\eqref{cdp:eq:C.AS.15}, with the sign fixed by orientation.
For a diagonal leaf both periods have zero $U$-component,
so this leading integral is zero, giving
\eqref{cdp:eq:C.AS.16}. Differentiation under the integral
and Stokes' theorem make the integral over the transported
closed cycle holomorphic; the displayed local expressions
give its holomorphic extension and stated orders at zero.

Finally, for a principal elliptic bundle over an elliptic
base, the homotopy sequence injects $\pi_1(F)=\Lambda_F$
because $\pi_2(B)=0$. Its structure group acts by
translations, so this subgroup is central. Lift the two
generators $a_B,b_B$ of a fundamental polygon for $B$ through a
trivialization over its one-skeleton. The attaching map
of its two-cell is the Chern winding $(c_1,c_2)$ in the
fibre lattice, giving \eqref{cdp:eq:C.NS.13}.
An abelian total group forces $e^{c_1}h^{c_2}=1$;
the injected lattice then gives $c_1=c_2=0$.
The nearby two-torus has abelian fundamental group, so
the conclusion follows through the constructed diffeomorphism.
\end{proof}


\subsection{Canonical divisors on Hopf leaf pairs}\label{subsec:II:hopf-leaf-pairs}

\begin{lemma}[A canonical divisor supported on one Hopf curve]\hypertarget{tag-P352}{}\label{lem:II:hopf-single-canonical-curve}
\StatementRecord{lem:II:hopf-single-canonical-curve}{proved}{P352}{\cref{lem:II:hopf-curves,lem:II:hopf-curve-restriction,lem:II:central-primary-hopf-cover,lem:II:reduced-anticanonical-hopf,lem:II:elliptic-star-boundary}; \cref{cdp:eq:C.SH.3,cdp:eq:C.SH.4}}
Let $H$ be a compact primary or secondary Hopf surface and
$J\subset H$ a connected smooth curve. An identity
$K_H\simeq\OO_H(-aJ)$ with $a\in\Z$ holds exactly for
the marked pairs in \eqref{cdp:eq:C.AS.1}, with the full
deck actions \eqref{cdp:eq:C.AS.2}; in these coordinates
$a=r+1$. The isomorphism marks $J$ as the image of $w=0$.
\end{lemma}
\begin{proof}
The curve $J$ is elliptic by \cref{lem:II:hopf-curves}.
By \cref{lem:II:hopf-curve-restriction} and Serre duality,
$H^0(H,K_H)=0$. For $a\le0$ the canonical section of
$\OO_H(-aJ)$ would be nonzero, so $a>0$.
Apply \cref{lem:II:reduced-anticanonical-hopf} with the
effective divisor $aJ$. Its reduced alternative has two
distinct curve images and cannot equal this single-curve
support. Its other alternative gives $a\ge2$ and a finite
unramified central primary cover $\pi:H'\to H$ on which
$D=\pi^{-1}(J)$ is one marked curve. The unramified
differential and divisor pullback give
\begin{equation}
                           K_{H'}=-aD.                  \hypertarget{tag-00AW}{}\label{cdp:eq:C.AS.7}
\end{equation}
The primary-cover clause of that lemma gives either the
nonlinear contraction $(\beta_0^rz+w^r,\beta_0w)$ or the
diagonal contraction $(\gamma^rz,\gamma w)$, with
$D=\{w=0\}$ and $a=r+1$. Its weight comparison
\cref{cdp:eq:C.AS.8,cdp:eq:C.AS.9} and its logarithmic
form \eqref{cdp:eq:C.AS.11} are on this same central-cover
construction; thus all inverse images of the marked curve
and all divisor coefficients are fixed in this application.

In the nonlinear case use the full centralizer map
\eqref{cdp:eq:C.SH.3} for this same central contraction.
Every member of the full deck group commutes with it.
The quotient and splitting maps \eqref{cdp:eq:C.SH.4}
therefore identify $H$ with
$H_r(\beta)/\mu_n$, with finite action
$(z,w)\mapsto(\zeta^rz,\zeta w)$ and $\gcd(r,n)=1$.
Here $\beta$ is the second-coordinate multiplier of the
primitive contracting generator in \eqref{cdp:eq:C.SH.4};
the central generator with multiplier $\beta_0$ is a positive
power of it times a finite deck element. The rescaling of
$z$ used there keeps $w=0$, and the
nonlinear primary curve is unique, so these maps carry
its image precisely to the given $J$.

For the diagonal central contraction, every deck lift
preserves its marked axis. For $r\ge2$, this contraction is
class $\mathrm{III}$ of \cite[arXiv v2, Theorem~2.1(3), p.~5]{Khaled2024Hopf}.
Its source parameters $(a,b,d)=(A,C,B)$ give
\begin{equation}
              (z,w)\longmapsto(Az+Cw^r,Bw).             \hypertarget{tag-00AH}{}\label{cdp:eq:C.AS.10}
\end{equation}
with $A,B\ne0$. For $r=1$, class $\mathrm{IV}$ in the same
theorem gives $\operatorname{GL}(2,\C)$; preservation of
$w=0$ sets its lower-left matrix entry to zero and gives
the same formula. Its action on
$z/w^r$ is affine. The central contraction acts trivially
on that coordinate, so the induced affine group is finite.
The average of its finite orbit of zero is a common fixed
point $c$. The maps $z'=z-cw^r$, $w'=w$ and
$z=z'+c(w')^r$, $w=w'$ conjugate the whole group to
diagonal maps while preserving the marked curve and
commuting with the central contraction.

The invariant form $dz\wedge dw/w^{r+1}$ trivializes
$K_{H'}((r+1)D)$. The section pulled back from
$K_H(aJ)$ differs from it by a constant, since $H'$ is
compact and connected. Invariance under every diagonal
deck element gives $AB=B^{r+1}$, hence $A=B^r$.
The full group is consequently the weighted scalar group
given by its second-coordinate multipliers. It has a
cyclic subgroup of finite index, so it is generated by
one contracting multiplier $\beta$, together with its
finite subgroup $\mu_n$. Freeness on $w=0$ gives
$\gcd(r,n)=1$: a root with $\zeta^r=1$ would fix a
nonzero point of that axis. Thus its quotient is exactly
the diagonal pair of \cref{cdp:eq:C.AS.1,cdp:eq:C.AS.2},
with the same marked curve.

Conversely, for each displayed pair the invariant form
and infinite order of $\OO_H(J)$ in the proof of
\cref{lem:II:elliptic-star-boundary} give
$K_H=-(r+1)J$ and uniqueness of the integer $a$.
This proves both directions with the stated pair maps.
\end{proof}


\subsection{Elliptic pruning and compact restoration}\label{subsec:II:compact-elliptic-star-insertion}

\begin{lemma}[One-mark elliptic quasibundles]\hypertarget{tag-P370}{}\label{lem:II:one-mark-elliptic-torsor}
\StatementRecord{lem:II:one-mark-elliptic-torsor}{proved}{P370}{\cref{def:II:elliptic-root-curve,lem:II:identity-monodromy-elliptic}}
Let $q:T\to\PP^1$ be an identity-monodromy elliptic quasibundle
with fixed Jacobian $F=\C/\Lambda_F$. Mark a fibre with full divisor
$mJ$ over infinity, and suppose every other fibre is smooth and reduced.
The marked surface is either $(F\times\PP^1,F\times\{\infty\})$,
with $m=1$, or a diagonal Hopf surface
$H_m^{\rm diag}(\beta)/\boldsymbol\mu_n$ with marked axis $W=0$,
where $n>0$ and $\gcd(m,n)=1$. The isomorphism is determined by
the torsor class, a primitive oriented fibre basis and the marked
root-chart identification. In this description $n$ is the full
root-line degree.
\end{lemma}
\begin{proof}
The root charts of \cref{lem:II:identity-monodromy-elliptic} identify
$q$ with a principal $F$-torsor on the root curve
\begin{equation}
                 \mathcal B=\PP^1[(\infty,m)].
\hypertarget{tag-00VH}{}\label{cdp:eq:C.PR.1}
\end{equation}
For an oriented basis $\lambda_1,\lambda_2$ of $\Lambda_F$, the
sheaves of holomorphic functions, units and $F$-valued maps fit into
\begin{equation}
\begin{gathered}
0\longrightarrow\OO_{\mathcal B}
\xrightarrow{\ a\mapsto
 (e^{2\pi i\lambda_2a},e^{-2\pi i\lambda_1a})\ }
(\OO_{\mathcal B}^{*})^2
\xrightarrow{\ \Phi\ }F_{\mathcal B}\longrightarrow0,\\
\Phi(u,v)=\frac{\lambda_1\log u+\lambda_2\log v}{2\pi i}
                              \pmod{\Lambda_F}.
\end{gathered}
\hypertarget{tag-00VR}{}\label{cdp:eq:C.PR.2}
\end{equation}
A change of logarithm adds a lattice period, so $\Phi$ is well-defined.
Local logarithms make it surjective. For a pair in its kernel, adjust
these logarithms by integral periods; their weighted sum is then zero,
which gives the displayed additive kernel. This kernel is injected:
$\lambda_1a,\lambda_2a\in\Z$ implies $a=0$, since
$\lambda_2/\lambda_1$ is nonreal.

Finite stabilizer averaging gives
$H^j(\mathcal B,\OO)=H^j(\PP^1,\OO)=0$ for $j=1,2$.
The canonical root line satisfies
\begin{equation}
 \operatorname{Pic}(\mathcal B)=\Z[\mathcal R],\qquad
               \mathcal R^m=\OO_{\PP^1}(\infty).
\hypertarget{tag-00VS}{}\label{cdp:eq:C.PR.3}
\end{equation}
Indeed subtracting the stabilizer weight of a line, as a power of
$\mathcal R$, makes that line descend to $\PP^1$. The displayed
relation then gives its unique integral power of $\mathcal R$.
The cohomology sequence of \eqref{cdp:eq:C.PR.2} identifies the
torsor class with a vector in $\operatorname{Pic}(\mathcal B)^2=\Z^2$.

The zero vector gives $F\times\mathcal B$. The connected root cover
of a smooth multiple fibre is unramified over the surface tube by
\cref{def:II:elliptic-root-curve}. Its inertia therefore acts freely
on the total torsor. In the product this inertia acts trivially on
$F$, so it is trivial: $m=1$ and $T=F\times\PP^1$.

For a nonzero vector, an oriented integral basis change sends it to
$(-n,0)$, with $n>0$. The inverse basis change identifies its primitive
direction with the original fibre basis. Push out
$(\mathcal R^{-n},\OO)$ through \eqref{cdp:eq:C.PR.2} and gauge the
second line coordinate to one. Its remaining ambiguity multiplies
the first coordinate by powers of
$\gamma=\exp(2\pi i\lambda_2/\lambda_1)$, with $|\gamma|<1$.
Thus the surface is
\begin{equation}
                  (\mathcal R^{-n})^\times/\gamma^{\Z}.
\hypertarget{tag-00VT}{}\label{cdp:eq:C.PR.4}
\end{equation}
Write $\mathcal B=\PP(m,1)$, with coarse coordinate $Z/W^m$.
Its nonzero line has the presentation
\begin{equation}
 [((\C^2\setminus0)\times\C^*_\ell)/
  \lambda\cdot(Z,W,\ell)
     =(\lambda^mZ,\lambda W,\lambda^{-n}\ell)].
\hypertarget{tag-00VU}{}\label{cdp:eq:C.PR.5}
\end{equation}
The gauge $\ell=1$ leaves $\boldsymbol\mu_n$ with weights $(m,1)$.
Choosing $\beta^n=\gamma$ lifts the infinite action to
$(Z,W)\mapsto(\beta^mZ,\beta W)$. Freeness on the axis $W=0$
is equivalent to $\gcd(m,n)=1$; on $W\ne0$ it follows from the
weight one. Changing the root $\beta$ multiplies this lift by an
element of the same finite group. These maps give the asserted
marked Hopf isomorphism and its inverse through \eqref{cdp:eq:C.PR.5}.
The integer $n$, together with the inverse integral basis change,
recovers the complete torsor class in \eqref{cdp:eq:C.PR.3}.
\end{proof}

\begin{lemma}[Two-mark elliptic quasibundles]\hypertarget{tag-P371}{}\label{lem:II:two-mark-elliptic-torsor}
\StatementRecord{lem:II:two-mark-elliptic-torsor}{proved}{P371}{\cref{lem:II:one-mark-elliptic-torsor}}
Let $q:T\to\PP^1$ be an identity-monodromy elliptic quasibundle
with fixed Jacobian $F$, whose multiple fibres are among two marked
fibres of orders $a,b$. As a surface with both marked curves and
fibre maps, $T$ is a diagonal primary or secondary Hopf surface, or
an elliptic ruled surface $\PP(\OO_E\oplus L)$ with its two sections,
where $L$ has finite order. The following root-line and quotient
maps determine the full torsor and both marked identifications.
\end{lemma}
\begin{proof}
Put $g=\gcd(a,b)$, $A=a/g$ and $B=b/g$. The map
\begin{equation}
 \PP^1[(0,A),(\infty,B)]
   \longrightarrow\PP^1[(0,a),(\infty,b)],\qquad t=z^g
\hypertarget{tag-00VW}{}\label{cdp:eq:C.PR.7}
\end{equation}
is representable and finite \'etale. At zero its root uniformizers
have map $v=u$ and inertia inclusion $\boldsymbol\mu_A\subset
\boldsymbol\mu_a$; the analogous description holds at infinity.
Since $\gcd(A,B)=1$, its source is $\PP(B,A)$, with coarse coordinate
$z=X^A/Y^B$. The two root divisors have classes $B$ and $A$.
B\'ezout's identity therefore identifies its Picard group with
$\Z[\OO(1)]$. Finite averaging gives the vanishing of its higher
coherent cohomology.

Apply \eqref{cdp:eq:C.PR.2} to this pullback root curve. The pulled
back torsor is represented by an integral vector. For a nonzero
vector use the oriented basis change to $(-n,0)$ and the gauge
\eqref{cdp:eq:C.PR.5}, with weights $(B,A)$ in place of $(m,1)$.
The resulting cover has contraction
$(X,Y)\mapsto(\beta^BX,\beta^AY)$ and finite action
$\boldsymbol\mu_n(B,A)$. Its inertia acts freely exactly when
$\gcd(n,A)=\gcd(n,B)=1$. Cyclic descent preserves the two axes.
On a central primary cover a lift commuting with the contraction
and preserving each axis is $(Xf(X,Y),Yh(X,Y))$. Equivariance
makes $f,h$ invariant under the contraction, hence constant by
iteration to the origin. The finite quotient preserving the axes
is cyclic: its character on either axis is injective by freeness.
These coordinate maps identify the descended marked surface with
a diagonal primary or secondary Hopf surface.

For the zero vector, free inertia gives $A=B=1$ and the pulled back
surface is $F\times\PP^1$. Descent acts by
$(y,z)\mapsto(y+\tau,\zeta_gz)$, since its translation part is a
holomorphic map $\PP^1\to F$ and hence constant. The principal
$F$-action fixes its linear part. Freeness at zero and infinity
makes $\tau$ have exact order $g$. Taking the quotient gives
\begin{equation}
 \PP(\OO_E\oplus L),\qquad
 E=F/\langle\tau\rangle,\qquad
 L=(F\times\C)/((y,v)\sim(y+\tau,\zeta_gv)).
\hypertarget{tag-00VX}{}\label{cdp:eq:C.PR.8}
\end{equation}
The map from $F\times\PP^1$ is the projectivization of the displayed
line quotient. It identifies the axes with the two sections, gives
$L$ exact order $g$, and descends the fibre map $q=z^g$.
All these identifications use the original torsor transition maps.
\end{proof}

\begin{lemma}[Leaf pruning and boundary groups]\hypertarget{tag-P372}{}\label{lem:II:elliptic-tree-pruning}
\StatementRecord{lem:II:elliptic-tree-pruning}{proved}{P372}{\cref{def:II:elliptic-star-data,lem:II:elliptic-star-boundary,lem:II:one-mark-elliptic-torsor,lem:II:two-mark-elliptic-torsor,lem:II:normal-product-collars}}
Let $S$ have a rooted-tree or unicyclic incidence graph. Give its
elliptic components identity-monodromy quasibundle maps, with
all their multiple fibres among the double curves, and give every
nonlinear Hopf component one marked axis as a leaf. All nonroot
vertices off the cycle have base $\PP^1$. The double-curve
identifications and normal dualities determine successive inverse
logarithmic transforms, from the leaves toward the root or cycle.
Every terminal step is a Hopf pair or the product pair
$(F\times\PP^1,F\times\{\infty\})$. The resulting root or marked
cycle, with its torsor atlas, is independent up to marked
isomorphism of the order of disjoint steps. The full integer
secondary framings and primitive fibre bases determine inverse
maps recovering each removed pair.

Replacing the filled leaf complement by its inverse product tube
induces an isomorphism of fundamental groups, equal to the identity
on the unchanged exterior and on its fibre generators. The same
comparison holds for compressed normal-product collars $uv=s^h$.
\end{lemma}
\begin{proof}
For a Hopf leaf of normal order $m$ and secondary integer $n$, set
$r=m$ in \cref{def:II:elliptic-star-data}. Its connected root cover
is identified with the parent's root cover by the given axis map
and dual normal lines. Choose $k,l$ with $mk-nl=1$.
The coordinate map \eqref{cdp:eq:C.PR.6}, together with the quotient
in \eqref{cdp:eq:C.SN.4}, gives
$p=X^mu^{-l}$, $W=X^{-n}u^k$ and the inverse
$u=p^nW^m$, $X=p^kW^l$. Replace this punctured chart by
$\Delta_p\times(\C^*/\gamma^{\Z})$ using that map.
The proof of \cref{def:II:elliptic-star-data} proves properness of
this replacement and extends every compatible coordinate change
as a translation over the product disc. It also identifies changes
$k\mapsto k+na$, $l\mapsto l+ma$ by $X\mapsto Xu^a$ while fixing
$p,W$. Thus the given gluing translation, primitive basis and
integer $n$ remain in the transition atlas.

A product cap is removed by the identity product-disc map on its
parent. At any other terminal vertex, removing its children has
removed their multiple fibres and changed no other chart.
Only the parent mark can remain multiple. Apply
\cref{lem:II:one-mark-elliptic-torsor} to this modified quasibundle;
it is exactly one of the two terminal pairs. This proves that the
procedure continues to the root or cycle. Operations in disjoint
base discs glue to the same complementary atlas and commute.
The same holds in disjoint subtrees. Every order processes a parent
after its children, and any two such orders differ by these
commutations. The resulting isomorphisms fix every remaining mark.
For a cycle, \cref{lem:II:two-mark-elliptic-torsor} identifies each
final component with its marked Hopf or ruled model.

For the group comparison apply the smooth-pair and boundary
clauses of \cref{lem:II:elliptic-star-boundary} with
$(r_i,n_i,\beta_i)=(m,n,\beta)$ to this same marked Hopf pair.
Its equivariant family $(\beta^m z+c w^m,\beta w)$,
$0\le c\le1$, and horizontal flow identify the nonlinear pair
with the diagonal pair, fixing its normal characters and
primitive boundary maps. The diagonal Hopf complement and the
product-cap complement are both $\C\times F$. In the resulting
boundary basis the filling homomorphism is
\begin{equation}
 \Z^3_{\delta,e,h}\longrightarrow\Z^2_{e,h},\qquad
              \delta\longmapsto0,\quad e\longmapsto e,
              \quad h\longmapsto h.
\hypertarget{tag-00VY}{}\label{cdp:eq:C.PR.9}
\end{equation}
Indeed \eqref{cdp:eq:C.SN.6}, with $r=m$, has columns
$(-l,k,0),(0,0,1),(m,-n,0)$ and determinant one.
It sends the reversed Hopf meridian to $\delta$ and its remaining
generators to $e,h$. The product tube $F\times D^2$ has the same
homomorphism on these same three loops. Van Kampen gives identical
pushout maps from the unchanged exterior and the filling; the
identity on that exterior and on $e,h$ induces the asserted group
isomorphism. Repeat the argument with all remaining collars fixed.
For a smoothing, \cref{lem:II:normal-product-collars} identifies
these collars with its normal-product tubes. In $uv=s^h$, fixing
$s\ne0$ still reverses the normal circle through
$v=s^h/u$; the same boundary maps and group comparison apply.
\end{proof}

\begin{lemma}[Compact core families and elliptic period charts]\hypertarget{tag-P373}{}\label{lem:II:elliptic-core-period-charts}
\StatementRecord{lem:II:elliptic-core-period-charts}{proved}{P373}{\cref{lem:II:logarithmic-period-family,lem:II:moving-marked-cycle,cdp:st:C.2.UN.57}}
Let the central core be a complex two-torus with a marked elliptic
subgroup, or a marked Hopf/ruled cycle with $M=I$ whose components
have identity-monodromy elliptic quasibundle maps. For each integer
$R>0$ there is a compact family over $\PP^1$ with that core at zero,
smooth two-torus fibres away from its ends, and a common elliptic
period of exact order $R$. In the torus case every fibre is smooth;
in the cycle case the fibre at infinity is the conjugate marked
cycle. The reduced cycle necks have equation $xy=t$.

Near a regular central elliptic fibre, its rank-two period cover
has the form
\begin{equation}
\begin{gathered}
 \mathcal G=(\C_p\times\C_w^*\times\Delta_t)/
 \langle(p,w,t)\mapsto(p+\delta(t),\gamma(t)w,t)\rangle,\\
 \delta(t)=t^\ell d(t),\qquad d(0)\ne0,\qquad |\gamma(0)|<1.
\end{gathered}
\hypertarget{tag-00VL}{}\label{cdp:eq:C.PR.13}
\hypertarget{tag-0188}{}\SsixEquationAliases{\label{cdp:eq:C.WT.27}\label{layout:II:3}}
\end{equation}
Initially $\ell=R$. An oriented integral change of the two fibre
periods changes $\delta$ by the holomorphic unit in the proof and
preserves its order, the marked central fibre and its lattice.
\end{lemma}
\begin{proof}
For the torus core $\C^2/\Lambda_0$, use the full rank-four clause
of \cref{lem:II:logarithmic-period-family}, pulled back along
$z=t^R$. Its period sections are
\begin{equation}
 V_{t^R}(w)=w+t^RJ\overline w,\qquad
                 J=\begin{pmatrix}0&1\\-1&0\end{pmatrix}.
\hypertarget{tag-00W3}{}\label{cdp:eq:C.PT.13}
\end{equation}
The quotient is the compact smooth family in that lemma, with the
same central lattice and frame $J\overline w+s^Rw$ at infinity.

For a cycle, use the common rank-three complement lattice of
\cref{cdp:st:C.2.UN.57}. Its regular elliptic fibres are translates of
\begin{equation}
 F=(W\cap iW)/(\Lambda\cap W\cap iW),\qquad
                     W=\operatorname{span}_{\R}\Lambda.
\hypertarget{tag-00VI}{}\label{cdp:eq:C.PR.10}
\end{equation}
Indeed the real span of a compact fibre lattice is a complex line
in $W$, hence is $W\cap iW$. Its quotient in each two-boundary
component is $\C^*$, so this is a primitive rank-two sublattice
of the common complement lattice. The connected normal-root maps
identify the fibre groups of adjacent components with this $F$.

Use the power-$R$ clause of \cref{lem:II:logarithmic-period-family}
with the same $\Lambda,b$ and $N=d$. The periods are
$V_{t^R}\Lambda$ and
$V_{t^R}b+d\log t\,V_{t^R}p/(2\pi i)$; their determinant is
\eqref{cdp:eq:C.PR.11}. For the filling at zero, apply the moving
clause of \cref{lem:II:moving-marked-cycle} to
$V_{t^R}(\lambda_j)$ and $V_{t^R}(\mathbf b_j)$, keeping each
integer normal meridian and the neck equation $xy=t$.
Its piecewise logarithmic maps have exactly the preceding fourth
period, so identify the punctured filling with this family.
At infinity divide the vector frame by $t^R$. The sections become
$J\overline w+s^Rw$ and the logarithmic term becomes
$-d\log s/(2\pi i)$. Taking outgoing normals
$-(J\overline{n_j}+s^Rn_j)$ gives the conjugate component frames
and the same reduced necks. The endpoint maps and their inverses
are those of \cref{cdp:st:C.2.UN.57} with this change of sections.
A finite closed cover by the two end discs and a middle annulus
has compact inverse images, proving compactness and properness.

Let $\lambda,\mu=c\lambda$ be an oriented primitive basis of a
regular fibre lattice, with $\operatorname{Im}c>0$. Directly,
\begin{equation}
 \det_{\C}(V_{t^R}\lambda,V_{t^R}\mu)
       =2i\operatorname{Im}(c)\|\lambda\|^2t^R.
\hypertarget{tag-00VK}{}\label{cdp:eq:C.PR.12}
\hypertarget{tag-00W4}{}\hypertarget{tag-0187}{}
\SsixEquationAliases{\label{cdp:eq:C.PT.14}\label{cdp:eq:C.WT.26}}
\end{equation}
Here $\det(\lambda,J\overline\lambda)=-\|\lambda\|^2$.
Choose a holomorphic linear frame in which the first period is
$(0,2\pi i)$. The second is $(\delta(t),\log\gamma(t))$;
its first coordinate is a unit times this determinant.
Exponentiation $(p,V)\mapsto(p,e^V)$ gives
\eqref{cdp:eq:C.PR.13}. Local logarithms give its inverse modulo
the same two periods. The central component charts select a
bounded transverse disc in this cover around the marked fibre.

For a new basis $\lambda'=a\lambda+b\mu$,
$\mu'=c_1\lambda+d_1\mu$, with $ad_1-bc_1=1$, put
\[
 p'=p-\frac{b\delta V}{2\pi ia+b\log\gamma},\qquad
 V'=\frac{2\pi i V}{2\pi ia+b\log\gamma}.
\]
It sends the first period to $(0,2\pi i)$ and the second transverse
period to $2\pi i\delta/(2\pi ia+b\log\gamma)$. The denominator
is a holomorphic unit: at zero it is a nonzero period of the
elliptic curve. These invertible linear maps give the stated
period-basis comparison and preserve $p$ on the central fibre.
\end{proof}

\begin{lemma}[Weighted insertion in an elliptic period chart]\hypertarget{tag-P374}{}\label{lem:II:weighted-elliptic-insertion}
\StatementRecord{lem:II:weighted-elliptic-insertion}{proved}{P374}{\cref{lem:II:elliptic-core-period-charts,def:II:elliptic-star-data}}
Use the period chart \eqref{cdp:eq:C.PR.13}, with parameter $s$.
Let $m,n,h$ be positive integers with $\gcd(m,n)=1$ and
$\ell\ge hm$. The normalized pullback by $t=s^h$ of the
weight-$(m,1)$ base modification, together with the pullback of
$\OO(nE_{\rm old})$, inserts a Hopf pair of normal order $m$
and secondary integer $n$. Its child is nonlinear for $\ell=hm$
and diagonal for $\ell>hm$, and its remaining displacement has
order $\ell-hm$. The complement map is the inverse of
\eqref{cdp:eq:C.PR.6}, with the original fibre basis.

For $m=1$ and $\ell>h$, the normalized blowup of $(p,s^h)$
inserts the product pair $(F\times\PP^1,F\times\{\infty\})$
and gives remaining order $\ell-h$. These replacements glue in
disjoint neighbourhoods of compact marked fibres in a proper
family and preserve properness. Their total spaces are normal,
with the compressed neck $uv=s^h$; for $h=1$ they are smooth
and the central divisor is reduced.
\end{lemma}
\begin{proof}
The normalized weighted chart and the canonical line section are
\begin{equation}
 p=x^mZ,\qquad xW=s^h,\qquad (Z,W)\ne(0,0).
\hypertarget{tag-018F}{}\label{cdp:eq:C.WT.33}
\end{equation}
The section of the pulled back exceptional line is $x^n$.
At a generic point of the new component, $x=s^h/W$, so its
vanishing order is $nh$. This identifies the pullback with its
canonical divisor section, including the multiplicity $h$.
Gauging the nonzero line coordinate to one leaves the finite action
\begin{equation}
 (Z,W,x)\longmapsto(\eta^mZ,\eta W,\eta^{-1}x),
                                  \qquad\eta\in\boldsymbol\mu_n.
\hypertarget{tag-018G}{}\label{cdp:eq:C.WT.34}
\hypertarget{tag-012B}{}\SsixEquationAliases{\label{cdp:eq:C.SN.9}}
\end{equation}
It is free: a nonzero $W$ or $x$ forces $\eta=1$, and on the
remaining axis $Z\ne0$ coprimality gives the same conclusion.
Choose the holomorphic root $\beta(s)^n=\gamma(s)$ with its
prescribed central value. The infinite map is
\begin{equation}
\SsixEquationAliases{\label{cdp:eq:C.PR.14}\label{cdp:eq:C.AS.18}\label{cdp:eq:C.SN.10}\label{cdp:eq:C.NS.19}}
\begin{split}
 (Z,W,x,s)\longmapsto\bigl(&
 \beta(s)^m[Z+d(s)s^{\ell-hm}W^m],\\
 &\beta(s)W,\ \beta(s)^{-1}x,\ s\bigr).
\end{split}
\hypertarget{tag-018H}{}\label{cdp:eq:C.WT.35}
\hypertarget{tag-00VM}{}\hypertarget{tag-00AP}{}
\hypertarget{tag-0123}{}\hypertarget{tag-00TK}{}
\end{equation}
It preserves $xW=s^h$ and changes $x^mZ$ by exactly
$d(s)s^\ell=\delta(s)$. Its terms have the same finite character,
so it commutes with \eqref{cdp:eq:C.WT.34}. For $h=1$,
$s^{\ell-m}W^m=x^{\ell-m}W^\ell$; this is the reduced weighted
chart. For $m=r_i$, $\ell=R$ and $d=d_i$ it is the marked-star
insertion. On the central child it gives the diagonal contraction
for $\ell>hm$ and a nonzero resonant term for $\ell=hm$.
Scaling $Z$ identifies that coefficient with the chosen Hopf
normal form. On $W=0$ its parent chart is
$(Z,x)\mapsto(\beta^mZ,\beta^{-1}x)$.

On the complement of the parent curve put $U=Z/W^m$ and
$w_{\rm child}=W^n$. The induced map is
\begin{equation}
 (U,W^n)\longmapsto
       (U+\delta(s)/s^{hm},\gamma(s)W^n).
\hypertarget{tag-018I}{}\label{cdp:eq:C.WT.36}
\hypertarget{tag-00VN}{}\SsixEquationAliases{\label{cdp:eq:C.PR.15}}
\end{equation}
This proves the order drop. A different marked point uses
$U-U_0$ and the same fibre coordinate. Off the inserted component,
the canonical section identifies the old fibre coordinate with
$x^{-n}$. On the parent, the invariant coordinates
$u=Z^n$, $X=xZ^k$, $mk-nl=1$, give precisely
$u=p^n(x^{-n})^m$, $X=p^k(x^{-n})^l$.
Thus the parent transition and its inverse are
\eqref{cdp:eq:C.PR.6}, including the complete integer $n$.

For a product cap, the ideal $(p,s^h)$ is preserved by the deck
map because $\delta$ is divisible by $s^h$. Its normalized blowup
has charts $p=s^hU$ and $xW=s^h$. Since $\ell>h$, the infinite
action on the exceptional normal directions is the identity.
The exceptional pair is $F\times\PP^1$ with its parent section,
and the other chart has map
\begin{equation}
\begin{gathered}
 p=s^hU,\qquad xW=s^h,\\
 (U,w)\longmapsto(U+\delta(s)/s^h,\gamma(s)w).
\end{gathered}
\hypertarget{tag-018K}{}\label{cdp:eq:C.WT.38}
\hypertarget{tag-00VO}{}\SsixEquationAliases{\label{cdp:eq:C.PR.16}}
\end{equation}
Its restriction to the parent is the identity, since the centre
is a Cartier curve there. These normalized blowups are proper.

We establish the compactness comparison for the Hopf replacement.
The exceptional-section norm is
\begin{equation}
 \rho(p,s)=(|p|^2+|s|^{2hm})^{1/(2m)}
          =|x|(|Z|^2+|W|^{2m})^{1/(2m)},\qquad
                              \rho\ge |s|^h.
\hypertarget{tag-018L}{}\label{cdp:eq:C.WT.39}
\hypertarget{tag-00VP}{}\SsixEquationAliases{\label{cdp:eq:C.PR.17}}
\end{equation}
Since $|\delta(s)|\le C|s|^{hm}$, translation by $j\delta(s)$
changes this norm by factors between
$C^{-1}(1+|j|)^{-1/m}$ and $C(1+|j|)^{1/m}$.
The pulled exceptional line has the corresponding exponents
$-n/m,n/m$. On a smaller closed parameter disc,
$0<\gamma_-\le|\gamma(s)|\le\gamma_+<1$.
These polynomial factors and the exponential factors
$|\gamma(s)|^j$ show that the infinite action is free and properly
discontinuous. A fixed high iterate contracts every line norm,
so each orbit has a representative in a fixed norm annulus.

This estimate also applies over a compact elliptic transverse
base with periods $P_1(s),P_2(s)$ and line factors $e^{Z_1(s)},
e^{Z_2(s)}$. Solve
$\operatorname{Re}(Z_j(s)-c(s)P_j(s))=0$ for $j=1,2$.
The real base-period matrix is invertible, giving a bounded smooth
$c(s)$. The invariant line norm is
$|w|e^{-\operatorname{Re}(c(s)p)}$; the infinite action multiplies
it by $|\gamma(s)|e^{-\operatorname{Re}(c(s)\delta(s))}<1$
after shrinking the disc. Uniform flat metrics give the same
two-sided polynomial estimate for the exceptional line.
A norm annulus over this compact weighted base is compact.

On an affine transverse chart, compare a compact old annular set
with its inverse image. Converting between its old line coordinate
and the new one uses the factor $x^{-n}$ and requires
$O(1+|\log\rho|)$ iterates. The transverse displacement is bounded by
\begin{equation}
                       C|\delta(s)|(1+|\log\rho|).
\hypertarget{tag-018N}{}\label{cdp:eq:C.WT.40}
\end{equation}
For $\rho<1$, use $\rho\ge|s|^h$; this bound tends to zero
as $s\to0$. For $\rho\ge1$, a new representative satisfies
\begin{equation}
 |p_{\rm new}|\le B+C|\delta(s)|\log(2+|p_{\rm new}|).
\hypertarget{tag-00VQ}{}\label{cdp:eq:C.PR.18}
\end{equation}
It is uniformly bounded, since a linear function eventually
dominates this logarithm. The corresponding annulus, including
the compact exceptional Hopf component, is therefore compact.

In a compact Hopf fundamental domain, $Z,W$ remain bounded and
not both zero while $x\to0$. The old base coordinate after line
normalization obeys
\begin{equation}
 p_{\rm old}=O\bigl(|x|^m+
               |s|^\ell(1+|\log|x||)\bigr)\longrightarrow0,
                         \qquad |s|^h=O(|x|).
\hypertarget{tag-018O}{}\label{cdp:eq:C.WT.41}
\end{equation}
The old normal-annulus coordinate has its compact limits on the
marked elliptic fibre. Thus every old cluster point lies on that
fibre. Choose disjoint small neighbourhoods of the marked curves.
The inserted components together with the inverse images of these
punctured neighbourhoods are open, and the cluster comparison
makes the gluing graphs relatively closed. The glued space is
Hausdorff. The bounded-annulus comparison over smaller closed
neighbourhoods, together with the unchanged compact exterior,
proves properness. The chart \eqref{cdp:eq:C.WT.33} gives the
stated compressed necks and normality, and for $h=1$ gives the
smooth reduced normal-product equation.
\end{proof}

\begin{lemma}[Finite stabilizers in weighted elliptic insertion]\hypertarget{tag-P375}{}\label{lem:II:finite-elliptic-insertion-gauge}
\StatementRecord{lem:II:finite-elliptic-insertion-gauge}{proved}{P375}{\cref{lem:II:elliptic-core-period-charts,lem:II:weighted-elliptic-insertion}}
Let a cyclic stabilizer of a marked fibre in
\eqref{cdp:eq:C.PR.13} have order $d_0$, base character $\xi$,
and identity action on its fibre lattice, with affine lift on
$V=\log w$ given by
\begin{equation}
 g(p,V,s)=(a(s)p+b(s),\ V+c(s)p+e(s),\ \xi s).
\hypertarget{tag-0189}{}\label{cdp:eq:C.WT.28}
\end{equation}
Its $d_0$th power is an infinite period to the integer power $k$,
possibly followed by an integral angular period. For every
$hm\le\ell$, holomorphic coordinate gauges make $a(s)=a_0$ and
$b(s)=O(s^{hm})$. They preserve the central marked maps and the
resonant coefficients. The resulting action extends to the
weighted insertion of \cref{lem:II:weighted-elliptic-insertion},
with its full group law, secondary framing and nonlinear translation.
For a product cap, the same assertion holds with $m=1$ and $\ell>h$.
\end{lemma}
\begin{proof}
The $p$-coordinate of the lifted order relation gives
\begin{equation}
 \prod_{j=0}^{d_0-1}a(\xi^js)=1,\qquad
                         g^{d_0}(p)=p+k\delta(s).
\hypertarget{tag-018A}{}\label{cdp:eq:C.WT.29}
\end{equation}
In particular $a_0^{d_0}=1$. Take the logarithm
$A(s)=\log(a(s)/a_0)$ vanishing at zero. Its finite orbit sum is
zero. The function
$u(s)=\exp(-d_0^{-1}\sum_{j=0}^{d_0-1}jA(\xi^js))$
satisfies $u(\xi s)a(s)/u(s)=a_0$ and $u(0)=1$.
The gauge $p\mapsto u(s)p$ therefore makes $a$ constant.
Commutation with the infinite period and the order relation become
\begin{equation}
 \delta(\xi s)=a_0\delta(s),\qquad
 \sum_{j=0}^{d_0-1}a_0^{d_0-1-j}b(\xi^js)=k\delta(s).
\hypertarget{tag-018C}{}\label{cdp:eq:C.WT.30}
\end{equation}
For $j<\ell$, the coefficient of $b_js^j$ in this sum vanishes
unless $\xi^j=a_0$. In that resonant case it equals
$d_0a_0^{d_0-1}b_j$, so $b_j=0$. A nonresonant coefficient is
removed by $p\mapsto p+c_js^j$, whose change in $b$ is
$(\xi^j-a_0)c_js^j$. Removing the finitely many coefficients
below $hm$ gives
\begin{equation}
                             b(s)=O(s^{hm}).
\hypertarget{tag-018D}{}\label{cdp:eq:C.WT.31}
\end{equation}
The central marked point is fixed, so its constant translation
is zero in this centred chart. At an intermediate diagonal Hopf
component, centring a nontrivial finite affine map at its fixed
point is its triangular Hopf automorphism; applying the same
map to all remaining marks gives their coordinate identifications.

For a nonlinear Hopf pair, put $U=z/w^m$, $V=\log w$ on its
universal cover. Its infinite generator is
$(U,V)\mapsto(U+c_0,V+\log\beta)$, with $c_0\ne0$.
Commutation of a parent-axis-preserving lift with a positive
central contraction forces its affine coefficient to satisfy
$a c_0=c_0$, hence its $U$-map is $U\mapsto U+b_0$.
Its lifted order relation is
\begin{equation}
                              d_0b_0=kc_0.
\hypertarget{tag-018E}{}\label{cdp:eq:C.WT.32}
\end{equation}
The angular element fixes $U$. At a terminal insertion
$\ell=hm$, equation \eqref{cdp:eq:C.WT.30} gives
$a_0=\xi^\ell$ and $b_\ell=ka_0d(0)/d_0$.
Dividing by $(\xi s)^{hm}$ leaves $kd(0)/d_0$, exactly the
translation in \eqref{cdp:eq:C.WT.32} after the same scaling of
$U$. For a diagonal component the relation is
$b_0(1+a+\cdots+a^{d_0-1})=0$. It gives $b_0=0$ for $a=1$;
for $a\ne1$ the fixed-point translation gives the centred map.

The gauge \eqref{cdp:eq:C.WT.31} preserves the weighted ideal
$(p,s^{hm})$, hence its normalized weighted modification.
The canonical section lifts this map to the pullback of
$\OO(nE_{\rm old})$. Local choices of an $n$th root of its
nonzero line multiplier differ by
\eqref{cdp:eq:C.WT.34}, so define the same coarse map.
The group law agrees on the unchanged open set, and therefore
on the normal chart. This extends the complete lifted relation,
including $g^{d_0}=G^k$ and its angular integer.

In the child chart $U=p/s^{hm}$, the transverse action is
\begin{equation}
 U\longmapsto\frac{a_0}{\xi^{hm}}U+
                         \frac{b(s)}{\xi^{hm}s^{hm}}.
\hypertarget{tag-018J}{}\label{cdp:eq:C.WT.37}
\end{equation}
Its fibre coordinate is $w_{\rm child}=s^{nh}w=W^n$.
Substituting \eqref{cdp:eq:C.WT.28} gives the holomorphic map
\begin{equation}\label{eq:II:weighted-child-finite-fibre}
 w_{\rm child}\longmapsto
 \xi^{nh}\exp\bigl(c(s)s^{hm}U+e(s)\bigr)w_{\rm child}.
\end{equation}
Its multiplier is a holomorphic unit. On logarithm charts the
source and target contributions $nh\log s$ cancel, leaving the
constant $nh\log\xi$ and the displayed holomorphic terms.
Thus the child has the same affine form
\eqref{cdp:eq:C.WT.28} and the new infinite displacement
\eqref{cdp:eq:C.WT.36}. For a product cap the identical argument
uses the ideal $(p,s^h)$ and the order-$h$ gauge. Finally transport
these local constructions around the entire finite orbit. Their
stabilizer maps agree by the proved group law, so they glue to one
finite action on the replacement family.
\end{proof}

\begin{lemma}[Compact restoration of elliptic trees]\hypertarget{tag-P353}{}\label{lem:II:compact-elliptic-star-insertion}
\StatementRecord{lem:II:compact-elliptic-star-insertion}{proved}{P353}{\cref{lem:II:elliptic-tree-pruning,lem:II:zero-chern-elliptic-torsor,lem:II:elliptic-core-period-charts,lem:II:weighted-elliptic-insertion}}
Consider a reduced marked support and its pruning sequence in
\cref{lem:II:elliptic-tree-pruning}. Its final core is either a
principal elliptic torsor $T_0\to B$ over an elliptic curve with
\begin{equation}
                          c(T_0)=0\in H^2(B,\Lambda_F),
\hypertarget{tag-00AT}{}\label{cdp:eq:C.AS.4}
\end{equation}
or a marked Hopf/ruled cycle with integral lattice product $M=I$.
Let $h_v$ be the sum of edge normal orders from that root or cycle,
with height zero on the core. Choose $R>0$ so that every nonlinear
leaf has height $R$ and every elliptic vertex off the core has
height less than $R$.

There is a proper map $Y\to\PP^1$ from a smooth compact complex
threefold whose full reduced fibre at zero is the given support,
with every component, torsor atlas and elliptic gluing map.
In the rooted case all other fibres are smooth connected two-tori.
In the cyclic case its only other singular fibre is the conjugate
pruned core. Every reduced fibre support has Euler characteristic
zero. For a marked elliptic star the height condition is precisely
\begin{equation}
 \begin{cases}
 r_i=R,&H_i\text{ has the nonlinear form in \eqref{cdp:eq:C.AS.2}},\\
 r_i<R,&H_i\text{ has the diagonal form in \eqref{cdp:eq:C.AS.2}}.
 \end{cases}
\hypertarget{tag-00AU}{}\label{cdp:eq:C.AS.5}
\end{equation}
\end{lemma}
\begin{proof}
For a rooted support apply \cref{lem:II:zero-chern-elliptic-torsor}
to $(E_1,B,F_1)=(T_0,B,F)$. Its local translation cochain and inverse
translations identify the analytic torsor over $B$ with the full
lattice quotient $T_0=\C^2/\Lambda_0$. With
$(z,w)_{\mathrm{KX}}=(p,z)$ this is
\eqref{cdp:eq:C.NS.14}. Its fibre lattice is primitive, since the
quotient base lattice is free. For a cyclic support the pruned
components, their maps and the lattice $M=I$ are those of
\cref{lem:II:elliptic-tree-pruning}. Apply
\cref{lem:II:elliptic-core-period-charts} to this root or cycle and
the chosen $R$. It gives the compact initial family and an
elliptic period of exact order $R$ at each first insertion point.

At a torus root the four periods can also be written in the
coordinates of the marked fibre map. The product map
\eqref{cdp:eq:C.SN.5} identifies its $i$th fibre with
$\C^*/\gamma_i^{\Z}$, where $\gamma_i=\beta_i^{n_i}$.
Choose its basis $\lambda_i=(0,a_i)$, $\mu_i=(0,b_i)$, with
$\gamma_i=\exp(2\pi i b_i/a_i)$, and extend it to an integral
basis of $\Lambda_0$. In the vector coordinates $(P,Z)$ after
$V_{s^R}$, the change
\begin{equation}
 p^{(i)}=P-s^R\frac{\bar a_i}{a_i}Z,\qquad
                         z_i=\frac{2\pi i}{a_i}Z
\hypertarget{tag-00TH}{}\label{cdp:eq:C.NS.16}
\end{equation}
sends the fibre periods to $(0,2\pi i)$ and
\begin{equation}
 (\delta_i(s),\log\gamma_i),\qquad
 \delta_i(s)=s^R\left(\bar b_i-\frac{\bar a_i b_i}{a_i}\right)
                                =d_i s^R,
\hypertarget{tag-00TI}{}\label{cdp:eq:C.NS.17}
\end{equation}
where $d_i\ne0$ by real independence. Write the complementary
periods as $(P_{ij}(s),Z_{ij}(s))$, $j=1,2$.
Their first coordinates remain a real base-period basis near zero.
They define $B_{i,s}=\C/\langle P_{i1}(s),P_{i2}(s)\rangle$
and the degree-zero line $M_{i,s}$ with factors $\exp Z_{ij}(s)$.
Successively quotienting by these periods, by $(0,2\pi i)$ and
by the last fibre period gives
\begin{equation}
 \mathcal T_s=M_{i,s}^{\times}/
 \langle(p,\ell)\mapsto(p+d_is^R,\gamma_i\ell)\rangle,
                              \qquad d_i\ne0.
\hypertarget{tag-00AO}{}\label{cdp:eq:C.AS.17}
\hypertarget{tag-012A}{}\hypertarget{tag-00TJ}{}
\SsixEquationAliases{\label{cdp:eq:C.SN.8}\label{cdp:eq:C.NS.18}}
\end{equation}
The first map is $(p,z)\mapsto(p,e^z)$ with the two line factors;
local logarithms modulo the same four periods give its inverse.
Thus this is the given analytic torus family with its full torsor
parameters. The line $M_{i,s}$ is arbitrary of degree zero; for
complementary periods with zero fibre coordinate it is trivial.

Reverse the pruning sequence. At a vertex of height $h_v$, the
period chart \eqref{cdp:eq:C.PR.13} has remaining order $R-h_v$.
Its primitive period basis is the basis of that inverse pruning
map, transported by the invertible coordinate change in
\cref{lem:II:elliptic-core-period-charts}. At a Hopf step with edge
order $m$ and full secondary integer $n$, apply
\cref{lem:II:weighted-elliptic-insertion} with $h=1$ and
$\ell=R-h_v$. Its formula \eqref{cdp:eq:C.PR.14} is holomorphic
because the child height is at most $R$. Its restriction to the
child is nonlinear precisely when the child height equals $R$;
otherwise it is the required diagonal pair. The chosen root
$\beta(0)^n=\gamma(0)$ and the scaling of $Z$ identify it with
that marked primary or secondary Hopf surface.
Equation \eqref{cdp:eq:C.PR.15} gives the next order $R-h_v-m$.
At a product cap the child is elliptic, so $R-h_v>1$;
\eqref{cdp:eq:C.PR.16} gives the product pair and the same order
drop one. Every further child is treated in its marked coordinate
$U-U_0$ and with its primitive fibre basis.

The compact-line part of \cref{lem:II:weighted-elliptic-insertion}
applies to \eqref{cdp:eq:C.AS.17}, with the same $P_{ij},Z_{ij},
d_i,\gamma_i$; its affine-chart comparison applies at every later
step and at a cyclic core. For a first replacement, the canonical
section gives the isomorphism
\begin{equation}
 \mathcal Y_i\setminus H_i\simeq
                         \mathcal T|_{\Delta_s}\setminus F_i.
\hypertarget{tag-00TM}{}\label{cdp:eq:C.NS.20}
\end{equation}
In its normalized line gauge the old coordinate is $x^{-n_i}$.
The same comparison \eqref{cdp:eq:C.WT.40}, with
$(m,h,\ell)=(r_i,1,R)$ and $s=xw$, gives the transverse displacement
\begin{equation}
                O\bigl(|s|^R(1+|\log|x||)\bigr)\longrightarrow0.
\hypertarget{tag-00AQ}{}\label{cdp:eq:C.AS.19}
\end{equation}
Together with $p-p_i^0=x^{r_i}z$, it sends every old cluster point
to $F_i$. At subsequent steps the same complement isomorphism
uses the preceding family and the next marked fibre. The lemma's
closed-gluing and compact-annulus comparisons therefore give
Hausdorff proper replacements at each of the finitely many stages.

On the strict parent, \eqref{cdp:eq:C.PR.6} is the original inverse
transform in the opposite direction. All other transition maps
are unchanged. The primitive fibre-basis maps give the linear
parts of the elliptic gluing. Their translations are realized by
$(Z,W)\mapsto(a^mZ,aW)$, which commutes with both Hopf contractions
and the finite action and gives every translation on the marked
axis. Product caps use the corresponding elliptic translations
and projective maps fixing the parent section. Consequently each
reverse step recovers the entire marked surface removed at that
step, including the torsor of a component that originally had
several multiple fibres. Induction recovers the given support.

Every insertion has equation $s=xW$ or $s=xy$ in its free coarse
chart, so the total space is smooth and every central component
has multiplicity one. The complement map is the identity on the
punctured family and outside the chosen neighbourhoods. The family
on the rest of $\PP^1$ is the compact initial family, with its
smooth endpoint or conjugate pruned cycle at infinity.
Quasibundle components have Euler number zero by
\cref{lem:II:identity-monodromy-elliptic}; Hopf components have
Euler number zero by \cref{lem:II:hopf-curve-restriction}; product
caps and all double curves also have Euler number zero.
Additivity proves the Euler assertion for every reduced support.
\end{proof}


\subsection{The reduced Hopf-star criterion}\label{subsec:II:reduced-hopf-star-criterion}

\begin{proposition}[Reduced Hopf stars with torus smoothings]
\hypertarget{tag-01DY}{}\label{cdp:st:C.2.UN.53}
\StatementRecord{cdp:st:C.2.UN.53}{proved}{01DY}{\cref{def:II:elliptic-star-data,lem:II:elliptic-star-boundary,lem:II:hopf-single-canonical-curve,lem:II:compact-elliptic-star-insertion,cdp:st:C.2.UN.17,cdp:st:C.2.UN.28,cdp:st:C.2.UN.31,lem:II:cyclic-component-cover,lem:II:zero-chern-elliptic-torsor}}
\hypertarget{tag-01DT}{}\label{cdp:st:C.2.UN.49}
\StatementRecord{cdp:st:C.2.UN.49}{proved}{01DT}{}
\hypertarget{tag-01DV}{}\label{cdp:st:C.2.UN.50}
\StatementRecord{cdp:st:C.2.UN.50}{proved}{01DV}{}
Let $S$ be a compact connected SNC surface with
smooth components $P,H_1,\ldots,H_n$, $n\ge1$, where each
$H_i$ is a primary or secondary Hopf surface. Its intersection
with $P$ is one connected smooth curve $J_i$, the $J_i$ are
disjoint on $P$, and there are no other intersections.
The following assertions are equivalent:
\begin{enumerate}
\item $S$ is the full reduced central fibre of a proper
degeneration of connected two-tori with smooth total threefold.
\item $S$ is the full fibre over zero of a proper map from
a smooth compact complex threefold to $\PP^1$, every other
fibre being a smooth connected two-torus.
\item The leaf pairs are the marked pairs
\cref{cdp:eq:C.AS.1,cdp:eq:C.AS.2}, their normal lines
are dual across every $J_i$, and $P$ admits the
quasibundle \eqref{cdp:eq:C.AS.3} over an elliptic curve,
with fixed Jacobian and identity integral monodromy.
Its marked full fibres are $r_iJ_i$ and it has no
unmarked multiple fibres. The inverse transform of
\cref{def:II:elliptic-star-data} has the zero Chern vector
\eqref{cdp:eq:C.AS.4}, and its leaf exponents satisfy
\eqref{cdp:eq:C.AS.5} for a positive integer $R$.
\end{enumerate}
All components, finite secondary actions, analytic torsor
parameters, primitive angular directions and given
elliptic identifications are preserved in the compact
realization. Every reduced support there has Euler number zero.

When a fixed-Jacobian identity-monodromy quasibundle $q:P\to B$
with the $J_i$ as distinct reduced fibre supports is part
of the given data, the same tests use that given $q$.
For exclusively nonlinear leaves they say that $B$ is
elliptic, every marked full fibre is $r_iJ_i$, there are
no other multiple fibres, the normals are dual, all $r_i$
equal one integer $r$, and the inverse-transform Chern vector
is zero. Arbitrary secondary orders $n_i$ satisfying
$\gcd(r_i,n_i)=1$ remain allowed. The primary specialization
is $n_i=1$, with the local identifications
$k=0,l=-1,u=z,X=x,W=x^{-1}$ in
\cref{cdp:eq:C.SN.4,cdp:eq:C.SN.5}.
If every leaf is diagonal, one can take
$R=1+\max_i r_i$. A nonlinear leaf determines $R=r_i$,
and every diagonal leaf then has strictly smaller exponent.
\end{proposition}
\begin{proof}
Assume the local occurrence. Along each double curve the
defining function has local product equation $xy=t$.
Its bilinear normal term identifies
$N_{J_i/P}\otimes N_{J_i/H_i}$ with $\OO_{J_i}$,
so the given curve identifications have dual normal lines.
Apply the evaluation construction
of \cref{cdp:st:C.2.UN.17} to this proper disc map, with
all full-fibre multiplicities one. Write its canonical
divisor as $k_PP+\sum_i k_iH_i$.
The principal central fibre and adjunction on a leaf give
\begin{equation}
                  K_H=\OO_H(-aJ),\qquad
                  a=k_H-k_P+1\in\Z.                      \hypertarget{tag-00AV}{}\label{cdp:eq:C.AS.6}
\end{equation}
Apply \cref{lem:II:hopf-single-canonical-curve} to this
same $H_i$, its connected curve $J_i$, and
$a=k_i-k_P+1$. It gives the full marked pair
\cref{cdp:eq:C.AS.1,cdp:eq:C.AS.2}, with
$k_i-k_P=r_i\ge1$.
Adjunction on $P$ gives exactly
\eqref{cdp:eq:C.NS.9}. Its effective canonical divisor
makes $\kappa(P)\ge0$ and $p_g(P)\ge1$.
The compact-component bound \cref{cdp:st:C.2.UN.28},
applied to this same fibre and its smooth component $P$,
gives $p_g(P)=1$.
The surface $P$ is minimal: an exceptional rational curve
$E$ is distinct from all the elliptic $J_i$ and intersects
them nonnegatively, whereas adjunction gives $K_P\cdot E=-1$.
This contradicts \eqref{cdp:eq:C.NS.9}.

To identify its fibration, apply the component-cover construction
\cref{lem:II:cyclic-component-cover}. The multiplicity of $P$
in this full fibre is one. In a generic chart of a normalized
root base change $t=s^N$, its equation is $x=s^N$ after
absorbing a unit. There is one strict covering component,
whose map to $P$ has degree one. Its subsequent resolution
remains birational to $P$. The smooth minimal model is
therefore this same $P$, by uniqueness of the minimal model
in nonnegative Kodaira dimension. Thus the cover conclusion
of \cref{cdp:st:C.2.UN.31} applies through this degree-one
map: $P$ is a two-torus or it has a fixed-Jacobian,
identity-monodromy elliptic quasibundle over a curve of
genus at most one. In the latter case
$p_g(P)=g(B)$ by \cref{lem:II:identity-monodromy-elliptic},
so $B$ is elliptic.

We choose this fibration so that all marked curves are fibres.
On a two-torus an embedded elliptic curve is a translate of
an elliptic subgroup: a lift to universal covers has periodic
holomorphic derivative, which is constant. Two such subgroups
with distinct tangent lines have a surjective product map
to the torus with finite kernel. Every two translates would
then meet. The disjoint $J_i$ have one common tangent line
and are therefore fibres of the quotient by that elliptic
subgroup; its connected quotient is an elliptic curve.

For the quasibundle over an elliptic base, use the canonical
map \eqref{eq:II:quasibundle-canonical}.
If there is a multiple fibre, its effective contribution
makes the intersection with every horizontal $J_i$ positive.
But normal duality gives $J_i^2=0$, because the leaf normal
has finite order, and adjunction gives $K_P\cdot J_i=0$.
Thus every $J_i$ is vertical.
If the quasibundle is smooth, it is a principal elliptic
torsor. For a nonzero Chern vector, a horizontal elliptic
curve would map to the elliptic base by a finite unramified
map of nonzero degree. Pullback to that curve has the
diagonal section $j\mapsto(j,j)$ and hence zero Chern class;
on $H^2(B,\Lambda_F)$ the same pullback multiplies the
Chern vector by the degree, a contradiction.
For zero Chern vector, \cref{lem:II:zero-chern-elliptic-torsor}
identifies this same total space with a two-torus through
its full analytic cocycle. The preceding subgroup argument
then supplies its quotient with all $J_i$ as fibres.

If a quasibundle with these marked fibre supports was already
given, keep that map throughout. Its relative minimality
follows from its smooth elliptic fibre supports, and
\cref{lem:II:identity-monodromy-elliptic} gives
$g(B)=p_g(P)=1$ directly. Thus this case has the same
elliptic-base conclusion for the given $q$.

For either resulting map, the normal-order comparison in
\cref{lem:II:elliptic-star-boundary} gives full marked
fibres $r_iJ_i$. Subtracting \eqref{cdp:eq:C.NS.9}
from its canonical formula leaves
\begin{equation}
 q^*K_B+\sum_{\text{unmarked multiple fibres}}(m-1)J=0
                    \quad\hbox{as a line bundle}.       \hypertarget{tag-00TB}{}\label{cdp:eq:C.NS.10}
\end{equation}
Since $B$ is elliptic, this is an effective sum of unmarked
multiple fibres linearly equivalent to zero. A nonzero
such divisor would be the divisor of a holomorphic function
on the compact connected $P$, which must be constant.
There are therefore no unmarked multiple fibres.

The inverse transform is now a principal elliptic torsor
over $B$. The core and collar map of
\cref{lem:II:elliptic-star-boundary} identifies its underlying
four-manifold with a nearby torus, using the unimodular
matrices \eqref{cdp:eq:C.SN.6} and the full analytic maps
\eqref{cdp:eq:C.SN.5}. Its injected fibre lattice in
\eqref{cdp:eq:C.NS.13} forces \eqref{cdp:eq:C.AS.4}.
All oriented Hopf core classes are the same integral
class $\gamma_{\rm core}$. The holomorphic period of the divided
canonical generator therefore has the expansions
\cref{cdp:eq:C.AS.15,cdp:eq:C.AS.16} at every leaf.
A nonlinear leaf gives a nonzero leading coefficient and
finite vanishing order $R=r_i$. Every other nonlinear
leaf has that order, and every diagonal leaf has
$r_i<R$. When all leaves are diagonal, the integer
$1+\max_i r_i$ satisfies \eqref{cdp:eq:C.AS.5}.
This proves every test in the third assertion.

Conversely, use the given leaf presentations, quasibundle,
axis identifications and normal dualities as the marked
star data of \cref{def:II:elliptic-star-data}.
The third assertion supplies exactly the elliptic-base,
zero-Chern and exponent hypotheses of
\cref{lem:II:compact-elliptic-star-insertion}.
Its explicit inverse-transform maps recover the same $P$,
its weighted charts recover each given leaf, and its
axis translations preserve the given gluing maps. It gives
the second assertion with all other fibres smooth tori
and all support Euler numbers zero. Restricting that
proper map to a disc about zero gives the first assertion.

For exclusively nonlinear data, \eqref{cdp:eq:C.AS.5}
is the common-resonance identity $r_i=r=R$.
The matrix \eqref{cdp:eq:C.SN.6} and meridian
\eqref{cdp:eq:C.SN.2} preserve each full secondary integer
$n_i$. For $n_i=1$, the substitutions in the statement
identify those maps with
\cref{cdp:eq:C.NS.1,cdp:eq:C.NS.2,cdp:eq:C.NS.4,cdp:eq:C.NS.7};
the period \eqref{cdp:eq:C.AS.15} becomes
\eqref{cdp:eq:C.NS.12}, and the compact deck map
\eqref{cdp:eq:C.AS.18} becomes \eqref{cdp:eq:C.NS.19}.
These are the complete primary and secondary nonlinear
specializations for the given quasibundle and surface.
\end{proof}

\section{Weighted cycles and finite root actions}\label{sec:II:weighted-root-cycles}

\subsection{Weighted normal data on marked cycles}\label{subsec:II:weighted-cycle-data}\label{subsec:II:weighted-hopf-cycles}

\begin{definition}[Weighted normal data on an oriented marked cycle]\hypertarget{tag-P355}{}\label{def:II:weighted-normal-cycle}
\StatementRecord{def:II:weighted-normal-cycle}{proved}{P355}{\cref{def:II:marked-normal-cycle}}
Let $S=\bigcup_i V_i$ be a compact simple normal-crossing oriented
cycle with at least two smooth embedded components, and with
every edge joining distinct components. Each $V_i$ is a diagonal
primary or secondary Hopf surface with its two marked axes,
or a surface
\begin{equation}
        R_i=\PP(\OO_{E_i}\oplus \mathcal L_i),\qquad
        E_i\ \hbox{elliptic},\qquad \mathcal L_i\in\operatorname{Pic}^0(E_i),
\hypertarget{tag-017F}{}\label{cdp:eq:C.WM.1}
\end{equation}
with its zero and infinity sections. The maps
$\phi_i:C_i^+\to C_{i+1}^-$ identify the outgoing and incoming
marked elliptic curves. Fix positive component multiplicities
$m_i$. Weighted normal data consist of holomorphic lines
$\mathcal N_i$ on $V_i$, sections
\begin{equation}
 \sigma_i\in H^0(V_i,\mathcal N_i^{-m_i}),
 \qquad
 \operatorname{div}(\sigma_i)=
       m_{i-1}C_i^-+m_{i+1}C_i^+,
\hypertarget{tag-016Y}{}\label{cdp:eq:C.WH.1}
\end{equation}
and the line isomorphisms
\begin{equation}
 \begin{aligned}
 \mathcal N_i|_{C_i^+}
   &\simeq \phi_i^*N_{C_{i+1}^-/V_{i+1}},\\
 \phi_i^*(\mathcal N_{i+1}|_{C_{i+1}^-})
   &\simeq N_{C_i^+/V_i}.
 \end{aligned}
\hypertarget{tag-0177}{}\label{cdp:eq:C.WH.2}
\end{equation}
whose leading section coefficients agree on each edge.
Equivalently these coefficients give the same trivialization of
$N_{C_i^+/V_i}^{m_{i+1}}\otimes
 \phi_i^*N_{C_{i+1}^-/V_{i+1}}^{m_i}$.
The data include every nonzero constant in these line maps.
Here $\mathcal L_i$ is the degree-zero ruling line, and
$\mathcal N_i$ is the weighted normal line.

A holomorphic function on a smooth threefold with full fibre
$\sum_i m_iV_i$ and this marked support induces these data.
In local coordinates $x=0$ for $V_i$ and $y=0$ for $V_{i+1}$,
their edge function is $x^{m_i}y^{m_{i+1}}$ after absorbing a unit.
\end{definition}
\begin{proof}
Take $\mathcal N_i=\OO_Y(V_i)|_{V_i}$. Divide the defining
function by the $m_i$th power of the canonical divisor section
and restrict to $V_i$. This gives \eqref{cdp:eq:C.WH.1}.
At a double curve the two divisor sections identify each
restricted ambient normal with the normal of that curve in
the other component, giving \eqref{cdp:eq:C.WH.2}. Their
product is the same defining function, so their leading
coefficients give the stated single weighted trivialization.
A local holomorphic root of its unit gives the displayed
edge equation. These constructions commute with the given
elliptic identifications.
\end{proof}


\subsection{Root covers and inertia of weighted cycles}\label{subsec:II:weighted-cycle-roots}

\begin{lemma}[Weighted cycle plumbing, root maps and inertia]\hypertarget{tag-P356}{}\label{lem:II:weighted-cycle-roots}
\StatementRecord{lem:II:weighted-cycle-roots}{proved}{P356}{\cref{def:II:weighted-normal-cycle,def:II:marked-normal-cycle,lem:II:moving-marked-cycle,lem:II:normal-product-collars}}
Every datum in \cref{def:II:weighted-normal-cycle} gives a
smooth proper local plumbing with full central fibre
$\sum_i m_iV_i$. Set $N=\operatorname{lcm}_i m_i$.
Its normalized base change $t=s^N$ has strict component covers
\begin{equation}
 \mathcal E_i=
 \operatorname{Norm}
 \{\xi\in\mathcal N_i^{-1}:
                         \xi^{m_i}=\sigma_i\}.
\hypertarget{tag-0178}{}\label{cdp:eq:C.WH.3}
\end{equation}
of total degree $m_i$. At an edge put
$a=m_i$, $b=m_{i+1}$, $g=\gcd(a,b)$,
$A=a/g$, $B=b/g$ and $h=N/\operatorname{lcm}(a,b)$.
Its $g$ local branches have the finite normalization maps
\begin{equation}
             uv=s^h,\qquad x=u^B,\quad y=v^A,
             \qquad t=s^N=x^a y^b.
\hypertarget{tag-00JH}{}\label{cdp:eq:C.EW.1}
\hypertarget{tag-0179}{}
\SsixEquationAliases{\label{cdp:eq:C.WH.4}}
\end{equation}
after choosing their constant roots of unity. A connected
strict cover of a Hopf component is a smooth marked diagonal
Hopf surface, and has the same algebraic dimension as that
component. Its two inverse axes are connected.

At a ruled vertex of multiplicity $m$, with neighbouring
multiplicities $a,b$, one has
\begin{equation}
 \begin{gathered}
      m\mid a+b,\qquad
      \deg(\mathcal N|_{\PP^1})=-\frac{a+b}{m},\\
      \gcd(m,a)=\gcd(m,b),\qquad
      q=\frac{m}{\gcd(m,a)}
         =\frac{m}{\gcd(m,b)}.
 \end{gathered}
\hypertarget{tag-017G}{}\label{cdp:eq:C.WM.2}
\end{equation}
Every connected strict cover is a marked
$\PP(\OO_{E'}\oplus\mathcal L')$, with a finite \'{e}tale
elliptic base map $\pi:E'\to E$ and a line isomorphism
\begin{equation}
                   (\mathcal L')^q\simeq\pi^*\mathcal L.
\hypertarget{tag-017H}{}\label{cdp:eq:C.WM.3}
\end{equation}
Its fibre map has degree $q$ and ramification $q$ along each
marked section.

The root support graph is a union of cycles; resolving a
compressed edge inserts $h-1$ degree-zero ruled elliptic
components. Its graph is connected exactly when the nearby
fibre is connected. On a connected graph let $M$ be the
product of its integral complement-lattice maps. Then
\begin{equation}
                  \pi_1(F_s)=\Lambda\rtimes_M\Z.
\hypertarget{tag-017A}{}\label{cdp:eq:C.WH.5}
\end{equation}
with the displayed lattice injected. For $M=I$ the complete
coframe map identifies the nearby fibre with a two-torus.
For a given smooth-total degeneration with this full weighted
support, the same root maps and the same fundamental-group
comparison apply.

For a connected root graph there is one integer $r$ such that,
for each connected strict cover $E_i$ of $V_i$,
\begin{equation}
               \deg(E_i\to V_i)=m_i/r,
                    \qquad r\mid\gcd_i m_i.
\hypertarget{tag-017B}{}\label{cdp:eq:C.WH.6}
\end{equation}
There are $r$ strict covers over each component. Write
$g_i=\gcd(m_i,m_{i+1})$ and
$h_i=N/\operatorname{lcm}(m_i,m_{i+1})$.
Each connected axis cover is \'{e}tale of degree $g_i/r$.
The common strict-component stabilizer $K$ and generic
component inertias $I_i$ satisfy
\begin{equation}
 \begin{gathered}
 |K|=N/r,\qquad |I_i|=N/m_i,\\
 |I_iI_{i+1}|=N/g_i,\qquad
 |I_i\cap I_{i+1}|=h_i.
 \end{gathered}
\hypertarget{tag-016X}{}\label{cdp:eq:C.WH.0}
\end{equation}
At a ruled vertex the base degree is $\gcd(m,a)/r$.

These same characters determine the local quotient of every
constructed moving compressed root plumbing with its finite
root action, base rotation $s\mapsto\zeta_Ns$, and the same
central component, branch and inertia maps. Its quotient is
smooth near the central fibre and has precisely the full
multiplicities $m_i$.
\end{lemma}
\begin{proof}
On a diagonal Hopf component the $m_i$th power of
$\mathcal N_i$ is the inverse of the boundary divisor line.
The two marked axes have constant factors of automorphy in
their diagonal quotient, so this power has a flat holomorphic
connection. Transfer it through \eqref{cdp:eq:C.WH.1}.
In a local frame divide its connection form by $m_i$; the
frame-change formula divides as well. This gives a flat
holomorphic connection on $\mathcal N_i$. Each entire axis
chart is the total space of its degree-zero normal line and
retracts to the axis. The flat line there has the same
monodromy as the pullback of its axis restriction, giving
the required holomorphic line isomorphism.

For a ruled vertex $R=\PP(\OO_E\oplus\mathcal L)$, write
\begin{equation}
  \mathcal N^{m}
       \simeq \OO_R(-aC^- -bC^+),\qquad
       \operatorname{div}(\sigma)=aC^-+bC^+.
\hypertarget{tag-017I}{}\label{cdp:eq:C.WM.4}
\end{equation}
Restriction to a ruling fibre gives
$m\deg(\mathcal N|_{\PP^1})=-(a+b)$. Thus $m\mid a+b$,
$b\equiv-a\pmod m$, and both gcds and the number $q$ in
\eqref{cdp:eq:C.WM.2} follow. On either entire section chart,
the removed divisor is trivial and the remaining zero-section
divisor is pulled back from its degree-zero normal line.
Consequently $\mathcal N^m$ is flat on that chart. Dividing
the transferred connection by $m$, and comparing monodromy
under the chart retraction, identifies $\mathcal N$ there
with its axis pullback. Globally, if
$d_0=\deg(\mathcal N|_{\PP^1})$, then
$\mathcal N\otimes\OO_R(-d_0C^+)$ is trivial on every
ruling fibre. Its direct image is a line $A_0$ on $E$;
fibrewise evaluation is an isomorphism by proper base change
and Nakayama. Hence
$\mathcal N=\OO_R(d_0C^+)\otimes\pi_R^*A_0$.
Equation \eqref{cdp:eq:C.WM.4} and the degree-zero class of
$C^--C^+$ give $\deg A_0=0$. This also describes the same
local pullback maps. The $m_i$th tensor of each parallel
identification agrees with the divisor-line identification
up to a constant. Fix that constant by the leading coefficient
of $\sigma_i$ and choose its $m_i$th root in the line map.
Then on each entire axis chart $\sigma_i$ is exactly that
leading coefficient times the indicated power of the normal
coordinate. Thus the chart comparison also preserves the
full functions used below.

On a component core use the total line $\mathcal N_i$ and
the function $\sigma_i(v^{m_i})$. On each edge use the two
lines of \eqref{cdp:eq:C.WH.2} and their function
$x^{m_i}y^{m_{i+1}}$, where $x$ is normal to $V_i$.
The entire chart maps just proved identify these functions
exactly. Choose the flat metrics through the weighted
trivialization, with bounded component normal norms. The
two ends of an edge cannot meet when both attaching norms
exceed $r_0$ and $|t|<r_0^{m_i+m_{i+1}}$.
Use the same four radii and finite closed-core/closed-neck
cover as \cref{lem:II:moving-marked-cycle}; choose the base
radius smaller than all these weighted bounds. On an
attaching annulus one norm is bounded below, so the equation
bounds the other by the corresponding positive root of
$|t|$. The same maps and closed-graph argument give a
Hausdorff proper space. Core line bundles and edge normal
bundles are smooth; the function has exactly the stated
central multiplicities and is a submersion off zero.

On the normalized root of a component core the formula
$\xi=s^{N/m_i}/v$ gives \eqref{cdp:eq:C.WH.3}, with root
generator $\xi\mapsto\zeta_N^{N/m_i}\xi$. For an edge,
$(x^Ay^B)^g=s^{gABh}$ first gives its $g$ branches. On one
branch the map in \eqref{cdp:eq:C.EW.1} is finite and is
one-to-one off $xy=0$: coprimality of $A,B$ makes the choices
$u^B=x$, $v^A=y$, $uv=s^h$ unique. Its source is the normal
surface $A_{h-1}$ times the axis. Thus this is its finite
normalization; the maps glue by uniqueness of normalization.

The open part of a connected Hopf root component is the
finite-index lattice cover
$\C^2/\Lambda_i'\to\C^2/\Lambda_i$. Let $q_+,q_-$ be
the least positive angular multiples in $\Lambda_i'$.
The local equations $\xi^{m_i}=z^{m_{i+1}}$ and its incoming
version give $q_+=m_i/\gcd(m_i,m_{i+1})$ and the analogous
$q_-$. Divide the two logarithmic coordinates by these
integers and exponentiate. The resulting deck group contains
a contraction and the standard two angular periods. A finite
element fixing a nonzero axis point would give a smaller
positive angular multiple, contradicting the definition of
$q_\pm$. Its quotient is therefore a smooth marked primary
or secondary diagonal Hopf surface. The coordinate-power map
$(z,w)\mapsto(z^{q_+},w^{q_-})$ is equivariant, finite on
compact quotients and equal to the open lattice cover.
Uniqueness of finite normalization identifies it with this
root component, including its two connected axes.
Algebraic dimension is preserved: pullback injects meromorphic
function fields, and multiplication by a meromorphic function
on the finite direct-image module, locally after clearing poles,
gives its monic characteristic polynomial. Its coefficients
glue meromorphically on the base. Thus the field extension
is algebraic and the transcendence degrees are equal.

For a ruled root component remove its sections. Its open group
is $\mathcal L^\times=\C^2/\Lambda$, with primitive normal
$n\in\Lambda$. For its finite-index sublattice $\Lambda'$,
let $q$ be least with $qn\in\Lambda'$. Projection gives
\begin{equation}
 E'=\C^2/(\C n+\Lambda')
      \longrightarrow
 E=\C^2/(\C n+\Lambda).
\hypertarget{tag-017J}{}\label{cdp:eq:C.WM.5}
\end{equation}
The projected rank-two lattices are discrete of finite index,
so this is an \'{e}tale elliptic cover. Its kernel in the open
group is $\C n/\Z qn\simeq\C^*$. Logarithmic base lifts
give constant factors for a degree-zero line $\mathcal L'$.
The open finite map is $z\mapsto z^q$ on this kernel;
exponentiating the two base lifts gives precisely
\eqref{cdp:eq:C.WM.3}. It extends by the homogeneous power map to
\begin{equation}
 \PP(\OO_{E'}\oplus \mathcal L')
          \longrightarrow \PP(\OO_E\oplus\mathcal L),
\hypertarget{tag-017K}{}\label{cdp:eq:C.WM.6}
\end{equation}
This is finite, has ramification $q$ on both sections and agrees
with the root cover on its dense open group. Its source is smooth
and normal, so finite-normalization uniqueness identifies it
with the strict root component. The two local equations give
the stated common $q$. The lattice index is
\begin{equation}
 [\Lambda:\Lambda']
     =q\,[\Lambda:\Lambda'+\Z n].
\hypertarget{tag-017L}{}\label{cdp:eq:C.WM.7}
\end{equation}
where the second factor is the base degree.

Every strict component now has two connected inverse marked
curves, so the root graph is a union of cycles. Apply the
compressed clause of \cref{lem:II:moving-marked-cycle} with
these component maps and the exponents $h_i$ in
\eqref{cdp:eq:C.EW.1}. Its resolution inserts exactly $h_i-1$
degree-zero ruled elliptic components; the full resolved fibre
is reduced. In its proper Stein factor, a connected central
support has one point above zero, and reducedness forces the
local base multiplicity to be one. Its finite map to the disc
therefore has degree one. Applied on each connected graph,
this proves the connectedness assertion. The same lemma gives
\eqref{cdp:eq:C.WH.5} and, for $M=I$, the closed-coframe quotient.
The map over $s\ne0$ to the original fibre over $t=s^N$ is an
isomorphism.

For a given weighted degeneration, its canonical divisor
sections give the data of \cref{def:II:weighted-normal-cycle}.
On a compact component core, a smooth tubular map expresses
its function as $v^{m_i}\sigma_i$ times a unit equal to one
on the zero section; the unique small smooth $m_i$th root
absorbs this unit. On double curves use both canonical
sections and the same weighted product trivialization.
These maps give the stated root covers and collar homomorphisms.
Resolve the explicit root necks and apply
\cref{lem:II:normal-product-collars} to that full reduced SNC
support. It has these same smooth normalization components
and labelled double curves, so its core and boundary maps
are exactly the ones used in \eqref{cdp:eq:C.WH.5}.

The root group preserves each map to the original component
and its two distinct labelled branches. It preserves the
lifted cyclic orientation and acts on the connected graph
by rotations. The stabilizer of a vertex is the kernel of
this rotation action, hence is one group $K$ at every vertex.
It acts transitively on covers of each original component.
Writing $r=N/|K|$, count the $N$ sheets at a generic component
point, whose normal inertia is $N/m_i$. This gives
\eqref{cdp:eq:C.WH.6}. There are $g_i$ points over a generic
axis and $r$ lifted edges, so each connected axis cover has
degree $g_i/r$. The local normalization makes it \'{e}tale at
every axis point. Its pointwise stabilizer has order $N/g_i$.
The product and intersection of the cyclic subgroups of
orders $N/m_i,N/m_{i+1}$ have orders $N/g_i,h_i$; this proves
\eqref{cdp:eq:C.WH.0}. Equation \eqref{cdp:eq:C.WM.7} then
gives the ruled base degree $(m_i/r)/q=\gcd(m_i,a)/r$.

Finally consider a moving compressed root action with these
central maps. At an axis the pointwise stabilizer $J_{\rm axis}$
has order $hAB$. Its order-$h$ subgroup fixes $u,v$ and rotates
$s$; the residual independent characters on $u,v$ have orders
$B,A$. Explicitly choose integers $c_0,d_0$ with
$c_0A+d_0B=1$; a generator has characters
$(\zeta_B^{c_0},\zeta_A^{d_0},\zeta_{hAB})$ on $(u,v,s)$.
Their product is the character of $s^h$.

To obtain these exact coordinates throughout the holomorphic
action, average $u,v$ to their central eigencharacters. Since
the action preserves both branches the averages are $u'=uU$
and $v'=vV$ with units $U,V$. Average an axis coordinate to
an invariant one; its derivative on the axis is unchanged.
The identity $u'v'=s^hUV$ and the character product show that
$UV$ is invariant. Replace $u'$ by $u'/(UV)$. The product is
then exactly $s^h$, the characters and base coordinate $s$
are unchanged, and the differential is invertible on the axis.
The analytic local quotient ring is
\begin{equation}
 \bigl(\C\{u,v,s\}/(uv-s^h)\bigr)^{J_{\rm axis}}
                         =\C\{u^B,v^A\}
                         =\C\{x,y\},
 \qquad t=x^a y^b.
\hypertarget{tag-00JP}{}\label{cdp:eq:C.EW.17}
\hypertarget{tag-0171}{}
\SsixEquationAliases{\label{cdp:eq:C.WH.12}}
\end{equation}
The first invariants, under the order-$h$ subgroup, are
$\C\{u,v\}$; the two residual power invariants give the
displayed ring. The residual axis cover is \'{e}tale. At a general
component point the generic inertia is a normal reflection
of order $N/m_i$. Off the central fibre the base rotation is
free. Thus the quotient is smooth near the whole fibre and
has exactly the stated multiplicities and central component
maps. Finite quotients preserve the properness just proved.
\end{proof}


\subsection{Finite affine data on root cycles}\label{subsec:II:weighted-cycle-affine-data}

\begin{lemma}[Finite root actions and the total normal period]\hypertarget{tag-P357}{}\label{lem:II:weighted-cycle-affine-data}
\StatementRecord{lem:II:weighted-cycle-affine-data}{proved}{P357}{\cref{lem:II:weighted-cycle-roots,def:II:marked-normal-cycle,lem:II:moving-marked-cycle,cdp:st:C.2.UN.57}; \cite[arXiv v2, Theorem~2.1, p.~5]{Khaled2024Hopf}}
Take weighted data whose root graph is connected and has $M=I$.
Identify its open complements by their complex linear edge maps
with $G=\C^2/\Lambda$, where $\Lambda$ has rank three. At root
vertex $j$ let $n_j$ be its outgoing normal meridian and $h_j$
its outgoing compressed-edge exponent. The root generator and
edges have affine expressions
\begin{equation}
 \begin{aligned}
 g(j,v,s)&=(\sigma j,Rv+\tau_j,\zeta_Ns),\\
 e_j(v,s)&=v+\mathbf b_j+l_j\log s,\\
 l_j&=-\frac{h_j n_j}{2\pi i}.
 \end{aligned}
\hypertarget{tag-00JS}{}\label{cdp:eq:C.EW.2}
\end{equation}
where $\sigma$ is a rotation,
$R\Lambda=\Lambda$, $R^N=I$, $Rn_j=n_{\sigma j}$ and
$h_{\sigma j}=h_j$. For lattice elements $\lambda_j,\mu_j$
their complete compatibility and order equations are
\begin{equation}
 \begin{aligned}
 \tau_{j+1}+R\mathbf b_j-\tau_j-\mathbf b_{\sigma j}
             +\frac{h_j n_{\sigma j}}N&=\lambda_j,\\
 \sum_{a=0}^{N-1}R^{N-1-a}\tau_{\sigma^a j}&=\mu_j.
 \end{aligned}
\hypertarget{tag-00JV}{}\label{cdp:eq:C.EW.3}
\hypertarget{tag-017E}{}
\SsixEquationAliases{\label{cdp:eq:C.WH.9}}
\end{equation}
The full vector $u$ and affine constant $b$ are
\begin{equation}
             b=\sum_j\mathbf b_j,\qquad
             u=-\sum_jh_jn_j=d p\in\Lambda,
\hypertarget{tag-00JW}{}\label{cdp:eq:C.EW.4}
\hypertarget{tag-017C}{}
\hypertarget{tag-017N}{}
\SsixEquationAliases{\label{cdp:eq:C.WH.7}\label{cdp:eq:C.WM.9}}
\end{equation}
with $p$ primitive and $d>0$. For
$W=\operatorname{span}_{\R}\Lambda=\ker\ell$, one has
$\ell(iu)\ne0$, $Ru=u$ and $Rp=p$. There are complex
coordinates in which
\begin{equation}
 \begin{gathered}
 W=\C\times\R,\qquad R=\operatorname{diag}(\rho,1),\\
 \Lambda=\langle w_1=(a_1,c_1),\ w_2=(a_2,c_2),\ p\rangle_{\Z},\\
 c_1,c_2\in\R,\qquad
 C=\operatorname{Im}(\overline{a_1}a_2)>0.
 \end{gathered}
\hypertarget{tag-00JX}{}\label{cdp:eq:C.EW.5}
\end{equation}
The eigenvalue $\rho$ has order $1,2,3,4$ or $6$;
$\rho^N=1$ and $R^r=I$, where $r$ is the strict-cover count.
The full affine descent relation is
\begin{equation}
                         (R-I)b-u/N\in\Lambda.
\hypertarget{tag-00JY}{}\label{cdp:eq:C.EW.6}
\end{equation}
Changing normal-product constants on whole rotation orbits
preserves every central component, elliptic gluing and finite
action map. Such a choice can arrange
\begin{equation}
                              b=(\sigma_0,i).
\hypertarget{tag-00JZ}{}\label{cdp:eq:C.EW.7}
\end{equation}
and when $\rho\ne1$ this leaves the first coordinate of $b$
unchanged.
\end{lemma}
\begin{proof}
The orientation and rotation assertion are those of
\cref{lem:II:weighted-cycle-roots}, with each branch mapped
to its original marked branch. An isomorphism of two diagonal
Hopf covers preserving both axes lifts to their universal
$\C^2\setminus0$ covers. Hartogs extends the map and its inverse
across zero. The derivative is diagonal and invertible because
both axes are preserved. The image of a contracting diagonal
deck generator is a diagonal target deck element; conjugating
their derivatives identifies its two multipliers with the
source multipliers. The lift therefore commutes with that
same diagonal contraction. Apply
\cite[arXiv v2, Theorem~2.1, p.~5]{Khaled2024Hopf}:
class $\mathrm{IIc}$ gives a diagonal centralizer; class
$\mathrm{III}$ gives $(Az+Cw^r,Bw)$, whose two-axis
condition forces $C=0$; and class $\mathrm{IV}$ gives
$\operatorname{GL}(2,\C)$, whose two-axis subgroup is diagonal.
These are exactly the marked lifts in all diagonal cases.

For a marked ruled isomorphism, each ruling fibre maps to a
point of the target elliptic base because every holomorphic
map $\PP^1\to E$ is constant. Its base map has an affine lift.
The line multiplier on the universal base cover has constant
automorphy ratios. Its logarithmic derivative is holomorphic
and doubly periodic, hence constant; its logarithm is affine.
Thus both sorts of marked maps are affine in the logarithmic
coordinates of \cref{def:II:marked-normal-cycle}.
Compatibility on the same necks makes their linear part one
map $R$. Finite order and preservation of the entire lattice
give $R^N=I$ and $R\Lambda=\Lambda$. The labelled normal maps
give $Rn_j=n_{\sigma j}$. Composition on a logarithm chart,
with $\log\zeta_N=2\pi i/N$, gives exactly
\eqref{cdp:eq:C.EW.3}; the $N$fold composition gives its
second equation with every affine translation included.

Resolve each compressed edge by
\eqref{eq:II:compressed-neck-resolution}. The inserted ruled
vertices repeat the same outgoing normal. Apply the normal
height calculation in \cref{cdp:st:C.2.UN.57} to this resolved
marked cycle: on Hopf vertices its contracting periods give
\cref{cdp:eq:C.HG.4,cdp:eq:C.AD.3}, and on ruled vertices the
two outgoing normals agree. In the all-ruled case its elliptic
quotient gives the same nonzero sign. Therefore all
$\ell(in_j)$ have one nonzero sign, and the $h_j$ repetitions
give $\ell(iu)\ne0$ for the vector in
\eqref{cdp:eq:C.EW.4}. Rotation preserves this sum, so
$Ru=u$ and $Rp=p$.

The complex line $W\cap iW$ is $R$-invariant and complementary
to $\C p$, since $p\in W$ but $ip\notin W$. Send these lines
to the first and second coordinate axes, with $p=(0,1)$.
Then $W=\C\times\R$. Extend $p$ to an integral basis of
$\Lambda$; its two projected vectors form a real lattice in
the first coordinate, and orient it to make $C>0$. This is
\eqref{cdp:eq:C.EW.5}. The induced integral two-dimensional
matrix of $R$ has determinant one and trace
$2\operatorname{Re}\rho\in\Z$. Its finite order gives exactly
the listed five possible orders. The stabilizer of a root
vertex acts on its finite open-group cover by deck translations.
Since its graph rotation has order $r$, this also gives $R^r=I$.

Sum the first equation of \eqref{cdp:eq:C.EW.3} around the
cycle. The $\tau_j$ cancel, yielding
\eqref{cdp:eq:C.EW.6}. For a whole $\sigma$-orbit, add
$c n_j$ to each $\mathbf b_j$ with the same $c\in\C$.
The first equation is unchanged because $Rn_j=n_{\sigma j}$;
the order equation has no $\mathbf b_j$. These additions are
precisely changes of the normal-product constants in the
marked charts, so every central elliptic map and action stays
fixed. In the chosen coordinates, the second components of
all $n_j$ are real, nonzero and of one sign. Their sum on an
orbit has nonzero second component. Choose $c$ to make the
second component of $b$ equal to $i$. The orbit sum is $R$-fixed;
when $\rho\ne1$ its first component is zero, proving the last
assertion and \eqref{cdp:eq:C.EW.7}.
\end{proof}


\subsection{Finite-equivariant logarithmic periods}\label{subsec:II:finite-equivariant-periods}

\begin{lemma}[Logarithmic period families with a finite root action]\hypertarget{tag-P358}{}\label{lem:II:finite-equivariant-periods}
\StatementRecord{lem:II:finite-equivariant-periods}{proved}{P358}{\cref{lem:II:weighted-cycle-affine-data,lem:II:logarithmic-period-family}}
Use the full lattice and affine data of
\cref{lem:II:weighted-cycle-affine-data}, with
$b=(\sigma_0,i)$ and $R=\operatorname{diag}(\rho,1)$.
Choose $1\le k\le N$ with $\rho=\zeta_N^k$, taking $k=N$
when $\rho=1$. Put
\begin{equation}
 J=\begin{pmatrix}0&1\\-1&0\end{pmatrix},\qquad
 V_s(w)=w+s^kJ\overline w.
\hypertarget{tag-00K0}{}\label{cdp:eq:C.EW.8}
\end{equation}
Then
\begin{equation}
                           R V_s=V_{\zeta_Ns}R.
\hypertarget{tag-00K1}{}\label{cdp:eq:C.EW.9}
\end{equation}
and the periods
\begin{equation}
 \begin{gathered}
 \Lambda_s=V_s\Lambda,\qquad
 B_s=V_s b+\frac{d\log s}{2\pi i}V_s p,\\
                         T_s=\C^2/(\Lambda_s+\Z B_s).
 \end{gathered}
\hypertarget{tag-00JI}{}\label{cdp:eq:C.EW.10}
\end{equation}
define a smooth proper family of connected two-tori on
$\C^*_s$. Its positive real determinant is
\eqref{cdp:eq:C.EW.11}, with $(R,N,z)=(k,d,s)$ in
\cref{lem:II:logarithmic-period-family}.

Transport the full edge and action data by
\begin{equation}
 \begin{aligned}
 \mathbf b_j(s)&=V_s \mathbf b_j,& n_j(s)&=V_s n_j,\\
 l_j(s)&=-h_jn_j(s)/(2\pi i),&
 \tau_j(s)&=V_{\zeta_Ns}\tau_j.
 \end{aligned}
\hypertarget{tag-00JK}{}\label{cdp:eq:C.EW.12}
\end{equation}
Their overlap identity is
\begin{equation}
 \begin{aligned}
 &\tau_{j+1}(s)+R\mathbf b_j(s)-\tau_j(s)-\mathbf b_{\sigma j}(\zeta_Ns)\\
 &\qquad-l_{\sigma j}(\zeta_Ns)\,2\pi i/N
       =V_{\zeta_Ns}\lambda_j.
 \end{aligned}
\hypertarget{tag-00JL}{}\label{cdp:eq:C.EW.13}
\end{equation}
and their $N$fold action has translation $V_s\mu_j$.
The last period satisfies
\begin{equation}
                    R B_s-B_{\zeta_Ns}
                   =V_{\zeta_Ns}\bigl((R-I)b-u/N\bigr)
                   \in\Lambda_{\zeta_Ns}.
\hypertarget{tag-00JM}{}\label{cdp:eq:C.EW.14}
\end{equation}
Thus they define a holomorphic finite action on the entire
punctured family. More precisely, set
$\theta_j(s)=\sum_{a<j}(\mathbf b_a(s)+l_a(s)\log s)$.
In the common coordinate $v_j-\theta_j(s)$ its translation is
\begin{equation}
 R\theta_j(s)+\tau_j(s)-\theta_{\sigma j}(\zeta_Ns).
\hypertarget{tag-00JN}{}\label{cdp:eq:C.EW.15}
\end{equation}

At infinity put $w=s^{-1}$ and divide the vector frame by
$s^k$. Then
\begin{equation}
 W_w(c)=J\overline c+w^k c.
\hypertarget{tag-00JQ}{}\label{cdp:eq:C.EW.18}
\end{equation}
and the periods are $W_w\Lambda$ and
\begin{equation}
                    W_wb-\frac{d\log w}{2\pi i}W_wp.
\hypertarget{tag-00JR}{}\label{cdp:eq:C.EW.19}
\end{equation}
The transformed finite action is exactly
\begin{equation}
 \begin{aligned}
 (j,v,w)&\longmapsto
     (\sigma j,R_\infty v+\tau_{j,\infty}(w),\zeta_N^{-1}w),\\
 R_\infty&=\rho^{-1}R=J\overline R J^{-1},\\
 \tau_{j,\infty}(w)
     &=J\overline{\tau_j}+\rho^{-1}w^k\tau_j
       =W_{\zeta_N^{-1}w}\tau_j.
 \end{aligned}
\hypertarget{tag-00JT}{}\label{cdp:eq:C.EW.20}
\end{equation}
At $w=0$ its lattice and affine data are the conjugates of
the data at zero through $v\mapsto J\overline v$ and the
conjugate base rotation.
\end{lemma}
\begin{proof}
For fixed $c$, $V_s c$ is holomorphic in $s$. The identity
$RJ\overline R^{-1}=\rho J$ gives
\eqref{cdp:eq:C.EW.9}. Apply the power clause of
\cref{lem:II:logarithmic-period-family} with $(R,N,z)=(k,d,s)$
and the same $\Lambda,b,p$. It gives exactly
\eqref{cdp:eq:C.EW.10}, including the unchanged coefficient
$d\log s/(2\pi i)$, the determinant \eqref{cdp:eq:C.EW.11},
logarithm monodromy and the proper quotient. Thus the period
family and its full integer lattice are those in that lemma.

Apply $V_{\zeta_Ns}$ to the first equation of
\eqref{cdp:eq:C.EW.3} and use \eqref{cdp:eq:C.EW.9}.
The correction $h_jn_{\sigma j}/N$ is a real rational
multiple of an integral vector, so it transports to the
correction in \eqref{cdp:eq:C.EW.13}. The translation is in
the target frame, as expressed by
$\tau_j(s)=V_{\zeta_Ns}\tau_j$. For the order equation,
$R^{N-1-a}V_{\zeta_N^{a+1}s}=V_sR^{N-1-a}$; summing gives
the period $V_s\mu_j$. These are the full overlap and
finite-order equations. Likewise \eqref{cdp:eq:C.EW.6}
gives \eqref{cdp:eq:C.EW.14}.

On each successive component chart subtract the accumulated
edge term $\theta_j(s)$. Its transformed coordinate differs
by \eqref{cdp:eq:C.EW.15}. Changing $j$ changes this by an
integral period, using \eqref{cdp:eq:C.EW.13}; the last edge
adds $B_s$. Changing the logarithm adds an integral normal
vector. Thus these maps descend to a single holomorphic
action on the entire punctured period quotient, with the
already proved order and full affine constants.

Multiplying every vector by $s^{-k}$ gives
\eqref{cdp:eq:C.EW.18} and \eqref{cdp:eq:C.EW.19}.
The image vector is multiplied by $(\zeta_Ns)^{-k}$, so the
linear part is $\rho^{-1}R=J\overline R J^{-1}$ and its
translation is exactly \eqref{cdp:eq:C.EW.20}. These are
the frame changes of the same action, hence preserve every
overlap and order identity. At $w=0$ the frame is
$J\overline{(\cdot)}$; it conjugates the lattice, affine
constants and finite action with the displayed inverse base
rotation. This proves the final comparison.
\end{proof}


\subsection{The weighted cycle criterion and compact realization}\label{subsec:II:weighted-cycle-criterion}

\begin{proposition}[Weighted diagonal Hopf and ruled cycles]\hypertarget{tag-01E4}{}\label{cdp:st:C.2.UN.59}
\StatementRecord{cdp:st:C.2.UN.59}{proved}{01E4}{\cref{def:II:weighted-normal-cycle,lem:II:weighted-cycle-roots,lem:II:weighted-cycle-affine-data,lem:II:finite-equivariant-periods,lem:II:moving-marked-cycle}}
For the full marked data of \cref{def:II:weighted-normal-cycle},
the smooth weighted plumbing has connected two-torus fibres
exactly when its normalized $N$th-root graph is connected and
its rank-three integral complement-lattice product satisfies
\begin{equation}
                         M=I
\hypertarget{tag-017M}{}\label{cdp:eq:C.WM.8}
\end{equation}
Every smooth-total two-torus degeneration with this marked
weighted support induces these data and satisfies the test.

Every weighted support satisfying it occurs over $\PP^1$ in
a smooth compact complex threefold, with all marked components,
elliptic gluings and full multiplicities preserved. Its only
other singular fibre is the complex conjugate weighted support.
All reduced fibre supports have Euler number zero. The
realization chooses the normal-product constants. The Hopf
components may have any algebraic dimension, and the cycle
may consist entirely of ruled components.

All component, axis, inertia and ruled-line maps and degrees
are those of \cref{lem:II:weighted-cycle-roots}. The finite
root action, including its linear eigenvalue of order
$1,2,3,4$ or $6$ and all its affine constants, is preserved
through the compact construction.
\end{proposition}
\begin{proof}
Apply \cref{lem:II:weighted-cycle-roots} to these precise
normal lines, sections and maps. Its smooth plumbing has
the full divisor $\sum_i m_iV_i$. On a connected root graph,
the lattice in \eqref{cdp:eq:C.WH.5} is injected. A two-torus
has abelian fundamental group, so the displayed conjugation
action is the identity. Conversely $M=I$ gives its complete
coframe isomorphism to a two-torus. The root fibre over $s$
is isomorphic to the original fibre over $s^N$; connectedness
is exactly the root-graph condition. Its final given-smoothing
comparison proves necessity for the same canonical normal data.

For compact realization use \cref{lem:II:weighted-cycle-affine-data}
on this connected root and choose $b=(\sigma_0,i)$ by its
orbitwise normal changes. The central components, elliptic
gluings and finite action are unchanged. Apply
\cref{lem:II:finite-equivariant-periods} to that entire lattice,
the same $R,\sigma,\tau_j$, and $u=dp$. This gives the
proper punctured two-torus family and its full finite action.

At zero transport the lattice by $V_s$. At a Hopf vertex the
component frame is
\begin{equation}
 P_j(s)=\left(\frac{V_s n_j}{2\pi i},
                        -\frac{V_s n_{j-1}}{2\pi i}\right).
\hypertarget{tag-00JO}{}\label{cdp:eq:C.EW.16}
\end{equation}
This is exactly the Hopf frame in
\cref{lem:II:moving-marked-cycle}, with its fixed integer
angular relations and contracting generator. At a ruled
vertex use that lemma's primitive normal, transverse base
lift and third lattice vector, so the family is the same
$\PP(\OO_{E'(s)}\oplus\mathcal L'(s))$ at $s=0$.
Apply its compressed clause to these components, the constants
$\mathbf b_j(s)$ of \eqref{cdp:eq:C.EW.12}, and the same
integers $h_j$. The period formula
\eqref{eq:II:compressed-moving-cycle-periods} is exactly
\eqref{cdp:eq:C.EW.10}. The lemma's accumulated-edge logarithmic
maps give a biholomorphism of the entire punctured filling
with that period family, with holomorphic inverse.

The action extends over the marked curves and necks.
On a Hopf chart $RP_j(s)=P_{\sigma j}(\zeta_Ns)$, so its
component logarithm changes by a translation; exponentiation
gives the diagonal map of the full primary or secondary
quotient with its fixed angular characters. On a ruled chart
the normal-period map and the affine elliptic base map give
the line-bundle map and its homogeneous projective extension
of \cref{lem:II:weighted-cycle-roots}. In particular a
translation of $\mathcal L^\times$ extends: its base is a
translation of $E$, and its group law gives the fibre-line
isomorphism, extending to zero and infinity. The exact
overlap identity \eqref{cdp:eq:C.EW.13} extends these maps
across $uv=s^{h_j}$; its order identity makes the $N$th
iterate a deck period. Thus this is an action on the whole
normal filling. In its period coordinates the action is
\eqref{cdp:eq:C.EW.15}, so the punctured biholomorphism is
equivariant, including every affine constant.

Apply the quotient clause of \cref{lem:II:weighted-cycle-roots}
to these same central component and inertia maps. Its explicit
eigencharacter coordinates and invariant ring
\eqref{cdp:eq:C.EW.17} prove a smooth quotient and exactly
the original multiplicities. The central root covers and
their maps were fixed throughout, so its reduced quotient
is the prescribed marked surface.

At infinity use $w=s^{-1}$ and the frame $W_w$ in
\eqref{cdp:eq:C.EW.18}. Take outgoing normals $-W_wn_j$,
edge exponents $h_j$, and constants $W_w\mathbf b_j$.
Their logarithmic coefficients are $h_jW_wn_j/(2\pi i)$,
so their complete period is \eqref{cdp:eq:C.EW.19}.
The compressed moving construction applies again. At zero
in this infinity chart its Hopf frame is
$J\overline{P_j(0)}$: complex conjugation changes $2\pi i$
to $-2\pi i$, exactly accounting for the chosen minus sign
on every normal. The contracting logarithms and finite
characters are therefore the conjugate ones. The separate
ruled frame likewise gives the conjugate line and marked
sections. In each component frame the edge constant is
the conjugate of its original value, since
$(J\overline P)^{-1}J\overline b=\overline{P^{-1}b}$.

Use exactly the infinity action \eqref{cdp:eq:C.EW.20}.
It is the frame change of the action already constructed,
so all overlap and finite-order identities hold. Its central
map is conjugated by $v\mapsto J\overline v$ and the base
rotation is conjugate. The quotient clause of
\cref{lem:II:weighted-cycle-roots}, applied to these conjugate
weighted data, gives the conjugate smooth quotient charts
and the same positive multiplicities.

Glue the two proper equivariant fillings to the full family
\eqref{cdp:eq:C.EW.10} by their logarithmic biholomorphisms
and the frame change \eqref{cdp:eq:C.EW.18}. As in the
closed-base cover of \cref{lem:II:moving-marked-cycle}, take
closed smaller discs in each filling and a closed annulus
in the proper punctured family. They cover the inverse image
of $\PP^1_s$ by compact sets. At either overlap boundary only
one open piece includes its limiting base value; the gluing
graphs are closed. The resulting normal root total space is
Hausdorff, proper and compact. Take its finite quotient with
base $t=s^N$. It is smooth at zero and infinity by the proved
invariant rings and elsewhere because the base action is
free. A general quotient fibre is isomorphic to any one of
its $N$ root fibres. Hence it is the specified connected
two-torus, and only the two stated weighted supports are
singular. Every component and every double elliptic curve
of these supports has Euler number zero, so additivity proves
the final Euler assertion.
\end{proof}


\subsection{Weighted cycles of algebraic dimension zero}\label{subsec:II:weighted-cycle-suspensions}

\begin{proposition}[Elliptic suspensions realizing weighted Hopf and ruled cycles]
\hypertarget{tag-01E1}{}\label{cdp:st:C.2.UN.56}
\StatementRecord{cdp:st:C.2.UN.56}{proved}{01E1}{\cref{cdp:st:C.2.UN.59,lem:II:weighted-cycle-affine-data,lem:II:weighted-cycle-roots,cdp:st:C.2.UN.48,lem:II:moving-marked-cycle}}
\hypertarget{tag-01DX}{}\label{cdp:st:C.2.UN.52}
\StatementRecord{cdp:st:C.2.UN.52}{proved}{01DX}{}
Take the data of \cref{def:II:weighted-normal-cycle}, with
at least one Hopf component, and with every Hopf component
of algebraic dimension zero. The connected-root and $M=I$
test in \cref{cdp:st:C.2.UN.59} is necessary and sufficient
for connected two-torus fibres in the smooth weighted plumbing,
and every smooth-total torus degeneration in this marked
class induces data satisfying that test.

Every support satisfying it has a compact realization over
$\PP^1$ obtained by replacing one tube in an elliptic
suspension. It preserves every marked component, elliptic
gluing, finite secondary action and full multiplicity, and
chooses the normal-product constants. Every reduced support
in this compact family has Euler number zero.
All component, axis, inertia and ruled-cover formulas of
\cref{lem:II:weighted-cycle-roots} apply. There are at least
two original Hopf vertices, and every ruling line
$\mathcal L_i$ has infinite order in $\operatorname{Pic}^0(E_i)$.
The all-Hopf case is included. In particular the examples in
\cref{ex:II:weighted-cycle-examples} give multiplicities $(1,n)$
for every $n\ge1$ and a primary three-component support of
multiplicities $(1,2,2)$.
\end{proposition}
\begin{proof}
The local criterion and all root-cover maps are the exact
specialization of \cref{cdp:st:C.2.UN.59,lem:II:weighted-cycle-roots}
to the stated algebraic dimensions. Assume the test passes
and use the same common $G=\C^2/\Lambda$, $u=dp$ and finite
affine data of \cref{lem:II:weighted-cycle-affine-data}.
Projection along $\C p$ gives an elliptic curve $E$ and
identifies $G$ with $\mathcal L^\times$ for a degree-zero
line $\mathcal L$; its logarithmic period construction is
\cref{def:II:marked-normal-cycle}. This line has infinite
order. Otherwise a multiplication isogeny trivializing it
would give a finite map $E'\times\C^*\to G$, whose image of
$E'\times\{1\}$ is a compact curve. One of the identical
open complements is a Hopf axis complement of algebraic
dimension zero. On a primary cover its diagonal contraction
is nonresonant, and \cite[arXiv v1, \S2.2.6, p.~5]{AbashevaDeev2024AB}
gives only the two axes as curves. For a secondary component,
\cref{lem:II:central-primary-hopf-cover} preserves algebraic
dimension on the central primary cover; the inverse image
of a compact curve would be a union of its axes, whose images
are the two marked axes. This excludes a curve in the common
complement.

The finite linear part fixes $p$, so it induces the identity
on the kernel $\C^*$ and an origin-preserving elliptic
automorphism $\beta$ with $\beta^*\mathcal L\simeq\mathcal L$.
For $\beta\ne1$, the nonzero linear map $\beta^*-1$ on the
dual elliptic curve is an isogeny, with finite kernel; it
would make $\mathcal L$ torsion. For $\beta=1$, the two
eigenvalues of $R$ are one and a finite-order shear is the
identity. Hence
\begin{equation}
                              R=I.
\hypertarget{tag-017D}{}\label{cdp:eq:C.WH.8}
\end{equation}
The full affine relation \eqref{cdp:eq:C.EW.6} consequently gives
\begin{equation}
                   u/N\in\Lambda,\qquad
                 d=Nq\quad\hbox{for an integer }q>0.
\hypertarget{tag-016Z}{}\label{cdp:eq:C.WH.10}
\end{equation}
Here $d$ is the primitive divisor of $u=-\sum_jh_jn_j$.
The resolved root cycle has $\sum_jh_j$ vertices: each edge
of exponent $h_j$ contributes its $h_j-1$ inserted ruled
components. The monomial resolution is an isomorphism on
each fibre over $s\ne0$, and the normalized-root projection
identifies that fibre with the weighted fibre over $t=s^N$.
Their common coframe takes the cycle period
$(\log s)u/(2\pi i)$ to $(\log t)(u/N)/(2\pi i)$.
Consequently the resolved graph computes the vector $u$
through these $h_j$ repeated meridians, while the finite
root quotient gives the suspension integer $q=d/N$.

Write $q_i$ for the fibre degree in \eqref{cdp:eq:C.WM.2}
at the $i$th ruled vertex. A torsion ruling line $\mathcal L_i$
would pull back to a
torsion line under the finite \'{e}tale base map in
\eqref{cdp:eq:C.WM.3}; its $q_i$th root $\mathcal L'$ would
then be torsion as well. A trivializing isogeny would give
a compact elliptic curve in this same common open group,
contradicting the previous paragraph. Thus all ruling lines
have infinite order. If there were only one original Hopf
vertex, follow a lifted Hopf vertex to the next one above
that original vertex. Every intermediate strict or resolved
vertex is ruled and repeats the outgoing normal. A deck
rotation takes the first Hopf to the second with $R=I$,
so their outgoing normals agree. The second Hopf would have
equal incoming and outgoing normal directions, contrary to
the independent two-axis frame. This proves the Hopf count.

For the compact suspension adjust edge constants on whole
rotation orbits. With $R=I$, $n_j$ is constant on an orbit.
At a Hopf vertex the two adjacent independent normal vectors
belong to different orbits. Their two orbit-normal sums span
$\C^2$, so solve their two-by-two linear equation to arrange
$\sum_j\mathbf b_j=0$. The full central elliptic and finite
action maps stay fixed by \cref{lem:II:weighted-cycle-affine-data}.
Its compatibility equations reduce exactly to the
$R=I$ specialization \eqref{cdp:eq:C.WH.9}, with
$-l_j2\pi i/N=h_jn_j/N$.

Use the lattice-coordinate isomorphism in
\eqref{cdp:eq:C.HG.6}, with this same primitive $p$. It sends
$u$ to $d(0,2\pi i)$ and gives $0<|q_0|<1$ and $a\ne1$.
The exclusion of $a=1$ is the same embedded elliptic map
\eqref{cdp:eq:C.HG.7} into this common Hopf complement.
Apply the arbitrary-parameter suspension clause of
\cref{cdp:st:C.2.UN.48} with its integer $d$ equal to the
present $q=d/N$, and with these exact $q_0,a$.
It gives a compact suspension whose selected periods are
\begin{equation}
 \begin{gathered}
 \Lambda(t)+\Z(0,q\log t),\\
 \Lambda(t)=
 \langle(0,2\pi i),(2\pi i,0),
                       (\log q_0,\log a(t))\rangle_{\Z} .
 \end{gathered}
\hypertarget{tag-0170}{}\label{cdp:eq:C.WH.11}
\end{equation}
Put $t=s^N$. Every normal, contracting period and finite
angular relation is transported by its fixed integer
coefficients in this rank-three lattice. At Hopf vertices
use the two independent normal frames of
\cref{lem:II:moving-marked-cycle}; at ruled vertices use its
primitive normal, transverse base lift and third period.
Thus the strict central component maps remain the given
finite Hopf power maps or ruled maps \eqref{cdp:eq:C.WM.6}.

The unknown vectors $\mathbf b_j(t),\tau_j(t)$ can be chosen
holomorphically through their central values to satisfy
\eqref{cdp:eq:C.WH.9} and $\sum\mathbf b_j(t)=0$.
Indeed their coefficient matrix has constant integers.
Its right side is a rational linear combination of the
three moving lattice generators. Every rational left-kernel
relation vanishes at zero; the three original generators
are rationally independent, so each rational coefficient
vanishes separately. The same relation therefore vanishes
after transport. A fixed linear right inverse on the matrix
image gives a holomorphic solution. Add its constant kernel
difference from the prescribed central solution to recover
those exact central values.

Apply the compressed construction of
\cref{lem:II:moving-marked-cycle} to these moving components,
edge constants and the original $h_j$. On component cores the
finite maps are translations of the common open group; they
extend to diagonal maps on Hopf quotients and to the line
and projective maps on ruled components. On necks they are
the induced line maps of $uv=s^{h_j}$.
Equation \eqref{cdp:eq:C.WH.9} is their exact overlap identity
and its order equation makes the $N$th iterate a deck period.
Thus these maps give the full finite action. The quotient
clause of \cref{lem:II:weighted-cycle-roots}, with the same
central characters, proves smoothness and precisely the
original weighted support.

The punctured quotient has periods \eqref{cdp:eq:C.WH.11}:
$\log s=(\log t)/N$ and a branch change adds
$u/N\in\Lambda$. The coframe descends because $R=I$.
Its local origins define a holomorphic torsor under this
marked family of tori. The descended coframe identifies the
Lie-algebra bundle with $\OO^{\oplus2}$. On the punctured disc
its sheaf sequence is
\begin{equation}
 0\longrightarrow\mathscr L
   \longrightarrow\OO^{\oplus2}
   \longrightarrow\mathscr A
   \longrightarrow0
\hypertarget{tag-0172}{}\label{cdp:eq:C.WH.13}
\end{equation}
The coherent sheaf $\OO^{\oplus2}$ has vanishing $H^1$ by
\cite[\S2, fundamental theorem~B, p.~80]{Cartan1957}, applied
to the Stein punctured disc. The lattice local system has
vanishing $H^2$: the circle retract has a cellular complex
with cells only in degrees zero and one, also with these
local coefficients. The cohomology sequence therefore gives
$H^1(\mathscr A)=0$. Lifting the translation cocycle and
taking its coboundary produces a holomorphic section and
an isomorphism of this entire punctured quotient with the
period family in \eqref{cdp:eq:C.WH.11}, including its
affine translation data.

Use this isomorphism for the nested-annulus replacement in
\cref{cdp:st:C.2.UN.48}, on the same compact suspension base.
Its closed-base compact cover and closed gluing graphs prove
properness and compactness of the smooth replacement. The
central quotient maps were unchanged, so every original
component, finite action and elliptic gluing is preserved.
At the selected value Hopf and ruled components and double
elliptic curves have Euler number zero. At every other value
the support is the unchanged elliptic-suspension support,
whose Euler number is zero by that same suspension clause.
This proves the stated compact realization.
\end{proof}


\subsection{Explicit weighted cycles and finite linear actions}\label{subsec:II:weighted-cycle-examples}

\begin{example}[Unequal Hopf multiplicities and finite linear actions]\hypertarget{tag-P359}{}\label{ex:II:weighted-cycle-examples}
\StatementRecord{ex:II:weighted-cycle-examples}{proved}{P359}{\cref{cdp:st:C.2.UN.56,cdp:st:C.2.UN.59,ex:II:three-hopf-cycle}}
There are compact smooth-total weighted Hopf-cycle realizations
with multiplicities $(1,n)$ for every $n\ge1$, and one with
three primary Hopf components of multiplicities $(1,2,2)$.
There are also compact realizations with two ruled components
of multiplicity $N$ and finite root linear part of exact order
$N$ for $N=2,3,4,6$.
\end{example}
\begin{proof}
Choose nonresonant contractions $\alpha,\beta$ and put
\begin{equation}
 H=(\C^2\setminus0)/
  \langle(\alpha z,\beta w),
              (\zeta_n z,\zeta_n^{-1}w)\rangle .
\hypertarget{tag-0173}{}\label{cdp:eq:C.WH.14}
\end{equation}
Use the two-copy inversion construction of
\cref{cdp:st:C.2.UN.57} with this same secondary Hopf surface.
Over $s\mapsto\zeta_ns$, take the identity on the first
copy and $(Z,W)\mapsto(Z,\zeta_nW)$ on the second.
The first neck $(z,w)\mapsto(z^{-1},s/w)$ intertwines these
maps. The other neck gives $(Z,W)\mapsto(\zeta_nZ,W)$;
it differs by the anti-diagonal deck element already in
\eqref{cdp:eq:C.WH.14}. These maps therefore give an order-$n$
action on the whole plumbing.

At a double curve the action is $(x,y)\mapsto(x,\zeta_ny)$.
The invariant coordinates are $x,Y=y^n$, with
\begin{equation}
                              t=s^n=x^nY .
\hypertarget{tag-0174}{}\label{cdp:eq:C.WH.15}
\end{equation}
The quotient is smooth, with the first component of
multiplicity one and the second of multiplicity $n$.
The first is the given $H$. On the second the anti-diagonal
finite group and the extra action generate the two independent
scalar groups. The map $(Z,W)\mapsto(Z^n,W^n)$ therefore
identifies its quotient with the primary Hopf surface
$H(\alpha^n,\beta^n)$. The axis maps $z\mapsto z^{-n}$ are
isomorphisms of their elliptic quotients, since their kernels
are exactly the respective finite angular groups.

The explicit compact comparison uses common coordinates
$(b,y)=(zw,w^n)$, so $q_0=\alpha\beta$ and $a=\beta^n$.
The root normal divisor is $d=n$ and the root degree is $N=n$;
thus the suspension index in \eqref{cdp:eq:C.WH.10} is one.
Take the parameter clause of \cref{cdp:st:C.2.UN.48} with
these $q_0,a$ and index one. On its pullback $t=s^n$, choose
$\beta(s)=a(s^n)^{1/n}$ with the given central root and set
$\alpha(s)=q_0/\beta(s)$. The same neck and finite-action
equations hold in these families. On the punctured common
quotient the action is exactly the base deck action.
The smooth finite quotient therefore replaces a tube in
that compact suspension and realizes \eqref{cdp:eq:C.WH.15}.

For the primary triangle of
\cref{ex:II:three-hopf-cycle}, use $s\mapsto-s$ and the
three component actions
\begin{equation}
            (z,w),\qquad(z,-w),\qquad(-z,w).
\hypertarget{tag-0175}{}\label{cdp:eq:C.WH.16}
\end{equation}
Substitution in its three maps \eqref{cdp:eq:C.HG.16}
intertwines the actions. On the middle double curve the
action is a free elliptic translation; on the other two it
is a normal reflection. The quotient has multiplicities
$(1,2,2)$ and primary contraction pairs
\begin{equation}
 (e^{-1},e^{-\sqrt2}),\quad
 (e^{-\sqrt2},e^{-2-2\sqrt2}),\quad
 (e^{-2-2\sqrt2},e^{-1}).
\hypertarget{tag-0176}{}\label{cdp:eq:C.WH.17}
\end{equation}
with axis maps inversion, identity and inversion. The root
normal divisor is two and its root degree is two, so the
same quadratic pullback of the index-one suspension and
the preceding finite quotient give this compact realization.

Finally let $N=2,3,4,6$ and take an elliptic curve $E$ with
an origin-preserving automorphism $\rho$ of exact order $N$.
For $N=2$ take any elliptic curve and $\rho=-1$; for $N=4$
use $\C/(\Z+i\Z)$ and $\rho=i$; for $N=3,6$ use
$\C/(\Z+e^{\pi i/3}\Z)$ and multiplication by
$e^{2\pi i/3},e^{\pi i/3}$, respectively. Start with a root
cycle of $2N$ copies of $E\times\PP^1$, joined by the
identity on their elliptic sections. Its complete data are
\begin{equation}
 \begin{gathered}
 \Lambda=\Lambda_E\oplus\Z p,\quad
 n_j=-p,\quad h_j=1,\quad
 \mathbf b_j=ip/(2N),\quad \tau_j=j p/N,\\
 \sigma(j)=j+2\pmod{2N},\qquad
 R=\operatorname{diag}(\rho,1).
 \end{gathered}
\hypertarget{tag-00JU}{}\label{cdp:eq:C.EW.21}
\end{equation}
In the first equation of \eqref{cdp:eq:C.EW.3} all terms
sum to zero except at the cyclic endpoint, where they give
$-2p$. The $N$fold translation sum is $(N-1)p$ on even
vertices and $Np$ on odd vertices. These are integral periods,
so the full action has the required order. It acts freely
on the central graph. The quotient has two components
$E\times\PP^1$, each of full multiplicity $N$. Identify
vertex two with vertex zero by the inverse generator: one
elliptic gluing is the identity and the other is $\rho^{-1}$.
The root cycle has identity lattice product, and its finite
linear part is this exact $R$. Apply
\cref{cdp:st:C.2.UN.59} to obtain the compact realization and
its conjugate weighted support with these same maps.
\end{proof}







\section{Hopf support graphs and elliptic trees}\label{sec:II:hopf-support-elliptic-trees}

\subsection{Cohomology and residues of support graphs}\label{subsec:II:hopf-support-tree-remainder}

\begin{lemma}[Cohomology and residues on a support graph]\hypertarget{tag-P363}{}\label{lem:II:support-incidence-residues}
\StatementRecord{lem:II:support-incidence-residues}{proved}{P363}{\cref{cdp:st:C.2.UN.17,cdp:st:C.2.UN.44}}
Let $f:U\to\Delta$ be a proper degeneration of connected
complex two-tori from a smooth threefold. Suppose its full
central fibre $S=\bigcup_{v\in V}T_v$ is reduced SNC,
with connected smooth components, disjoint smooth elliptic
double curves $C_e$, and no triple points. Let $G=(V,E)$
be its incidence graph. For an oriented edge $e:v\to w$,
identify its two copies with $C_e$ by the given gluing map,
and write
\begin{equation}
r_{v,e}:H^1(T_v,\OO_{T_v})
                 \longrightarrow H^1(C_e,\OO_{C_e}),
\hypertarget{tag-00N0}{}
\label{cdp:eq:C.HR.1}
\end{equation}
for the endpoint restrictions. Define
\begin{equation}
d_1:\bigoplus_vH^1(T_v,\OO_{T_v})
             \longrightarrow\bigoplus_eH^1(C_e,\OO_{C_e}),
 \qquad(d_1a)_e=r_{v,e}a_v-r_{w,e}a_w
 \quad(e:v\to w).
\hypertarget{tag-00N1}{}
\label{cdp:eq:C.HR.2}
\end{equation}
The normalization sequence is
\begin{equation}
0\longrightarrow\OO_S
 \longrightarrow\bigoplus_{v\in V}\OO_{T_v}
 \xrightarrow{\ d\ }\bigoplus_{e\in E}\OO_{C_e}
 \longrightarrow0.
\hypertarget{tag-00N4}{}
\hypertarget{tag-00W6}{}
\hypertarget{tag-00WK}{}
\label{cdp:eq:C.HR.5}
\SsixEquationAliases{\label{cdp:eq:C.PT.3}\label{cdp:eq:C.QT.4}}
\end{equation}
where $d$ is the difference of the two restrictions at each
edge. Its map on constants is the signed incidence map $d_0$.

\begin{enumerate}[label=(\arabic*),leftmargin=*]
\item Suppose $h^\bullet(T_v,\OO_{T_v})=(1,1,0)$ for every
vertex and every $r_{v,e}$ is nonzero. Then $G$ is unicyclic,
$d_1$ has rank $|V|-1$, and
\begin{equation}
\begin{gathered}
 0\longrightarrow\operatorname{coker}d_0
 \longrightarrow H^1(S,\OO_S)
 \longrightarrow\ker d_1\longrightarrow0,\\
 H^2(S,\OO_S)\simeq\operatorname{coker}d_1.
 \end{gathered}
\hypertarget{tag-00N5}{}
\label{cdp:eq:C.HR.6}
\end{equation}
In any nonzero endpoint and edge bases, the product of
$r_{v,e}/r_{w,e}$ along the oriented cycle is one. It is
independent of those bases, and
\begin{equation}
(h^0,h^1,h^2)(S,\OO_S)=(1,2,1).
\hypertarget{tag-00N2}{}
\hypertarget{tag-00WL}{}
\label{cdp:eq:C.HR.3}
\SsixEquationAliases{\label{cdp:eq:C.QT.5}}
\end{equation}
Every nonzero section of $\omega_S$ has nonzero residues
exactly on the cycle edges. Its component two-forms are
nonzero exactly on the cycle components. A canonical
evaluation generator has its minimum vertical order
exactly on those components.

\item Suppose $G$ is a nontrivial tree with distinguished
vertex $P$, and the preceding dimensions and nonzero
restriction conditions hold at every vertex other than $P$.
The restriction maps at $P$ are arbitrary. Restriction gives
\begin{equation}
H^2(S,\OO_S)\xrightarrow{\sim}H^2(P,\OO_P),
 \qquad H^1(S,\OO_S)\xrightarrow{\sim}H^1(P,\OO_P).
\hypertarget{tag-00W7}{}
\label{cdp:eq:C.PT.4}
\end{equation}
and
\begin{equation}
H^0(S,\omega_S)\xrightarrow{\sim}H^0(P,K_P).
\hypertarget{tag-00W8}{}
\label{cdp:eq:C.PT.5}
\end{equation}
The inverse of the last map puts the given holomorphic
two-form on $P$ and zero on all other components. Every
edge residue is zero. The minimum canonical evaluation
order is attained uniquely at $P$.
\end{enumerate}
\end{lemma}
\begin{proof}
At a double curve the local functions on $xy=0$ are the
pairs of functions on $x=0$ and $y=0$ agreeing on $x=y=0$.
Their difference map is surjective onto the functions on
that curve. These descriptions, with the given gluing maps,
prove \eqref{cdp:eq:C.HR.5}. Connectedness makes
$\operatorname{rank}d_0=|V|-1$. When all component second
cohomology groups vanish, its long exact sequence is
\eqref{cdp:eq:C.HR.6}.

Choose a spanning tree and a nonzero value at one vertex.
The equations $r_{v,e}a_v=r_{w,e}a_w$ determine nonzero
values successively along it. Hence its $|V|-1$ rows are
independent, with a one-dimensional kernel. For a tree
$d_1$ would be surjective and $H^2(S,\OO_S)$ would vanish.
The map $f$ is flat: its nonzero base parameter is a
non-zero-divisor in the regular local rings of $U$, and a
torsion-free module over the base discrete valuation ring
is flat. Proper flat semicontinuity instead gives
$h^2(S,\OO_S)\ge h^2(F_t,\OO_{F_t})=1$.

Apply \cref{cdp:st:C.2.UN.44} to this $f$, with every
component multiplicity equal to one. Its meridian quotient
makes $\pi_1(S)$ abelian. Its map to the SNC dual complex
is a surjection onto $\pi_1(G)$, because the absence of
triple points makes that complex this same graph.
The free graph group is therefore abelian and has rank
at most one. The preceding exclusion of trees gives
$|E|=|V|$. After the spanning-tree equations there is
one remaining row. It vanishes on their kernel precisely
when the endpoint-ratio product around the cycle is one.
A nonzero last row would make $d_1$ invertible and again
give $H^2(S,\OO_S)=0$. Thus its rank is $|V|-1$.
Equation \eqref{cdp:eq:C.HR.6} now proves
\eqref{cdp:eq:C.HR.3}. A vertex-basis scalar occurs once
in a numerator and once in a denominator; an edge-basis
scalar cancels inside its ratio. This proves the stated
basis independence with the given restriction maps.

For the rooted tree, order the nonroot columns from leaves
towards $P$. The submatrix of $d_1$ on these columns is
square. Its leaf pivot is the nonzero restriction at that
leaf, and successive elimination proves it invertible.
Consequently $d_1$ is onto, and projection from its kernel
to $H^1(P,\OO_P)$ is an isomorphism: every root value
has one uniquely determined vector of nonroot values.
The tree map $d_0$ is onto. The normalization long exact
sequence therefore gives both restriction isomorphisms
in \eqref{cdp:eq:C.PT.4}, using the vanishing of
$H^2(T_v,\OO_{T_v})$ away from $P$.

Put $D_v=\bigcup_{e\ni v}C_e$ on $T_v$. Its disjoint
smooth branches give the residue sequence
\begin{equation}
0\longrightarrow K_{T_v}
 \longrightarrow K_{T_v}(D_v)
 \xrightarrow{\operatorname{Res}}\bigoplus_{e\ni v}K_{C_e}
 \longrightarrow0.
\hypertarget{tag-00N6}{}
\label{cdp:eq:C.HR.7}
\end{equation}
Serre duality identifies its residue connecting map with
the duals of the same endpoint restrictions. In a double
chart $xy=0$, a dualizing two-form has opposite residues
on its two branches. Assign the edge residue on its tail
and its negative on its head. When every component has
$H^0(K_{T_v})=0$, the global residue vector is exactly
$\ker d_1^\vee$: the component residue obstruction
vanishes precisely on that kernel, and the component
lifts are unique. Conversely those unique lifts glue by
the opposite-residue rule. These constructions are inverse.

At a leaf of an attached tree, its nonzero dual restriction
coefficient forces the edge residue to vanish. Delete that
leaf and repeat. Thus all tree residues vanish. On the
remaining cycle the two adjacent residues at each vertex
are related by nonzero coefficients; its one-dimensional
kernel has every cycle residue nonzero. A component with
all residues zero has an ordinary holomorphic two-form,
which is zero in the unicyclic case. A cycle component
has its two nonzero residues and hence a nonzero form.
For a rooted tree the same elimination stops at $P$.
Every other component form vanishes, and the form on $P$
has zero residues and is holomorphic. Extension by zero
and restriction are inverse, proving
\eqref{cdp:eq:C.PT.5}.

Take the canonical evaluation section from
\cref{cdp:st:C.2.UN.17}, using $dt$ to identify its
canonical and relative canonical lines. Write its
vertical divisor as $\sum k_vT_v$ and set $k_0=\min k_v$.
Since the full fibre is the reduced divisor of $t$,
division by $t^{k_0}$ is holomorphic and has nonzero
restriction to $S$. Near an edge, with $t=xy$, its
adjunction map is, up to a holomorphic unit,
\begin{equation}
\frac{x^{k_v-k_0}y^{k_w-k_0}\,dx\wedge dy\wedge dz}{xy}
 \quad\longmapsto\quad
 y^{k_w-k_0-1}\,dy\wedge dz\quad\text{on }x=0,
 \qquad k_v=k_0.
\hypertarget{tag-00N7}{}
\label{cdp:eq:C.HR.8}
\end{equation}
The restriction on $x=0$ is zero when $k_v>k_0$.
At a smooth point of a minimum-order component it is
nonzero. Thus its nonzero component forms identify
exactly the minimum-order components. The residue
descriptions just proved give the two asserted minimum
loci. In the tree case they also show that
$H^0(P,K_P)$ contains the resulting nonzero form.
\end{proof}

\begin{proposition}[Hopf support graphs and diagonal cycles]\hypertarget{tag-01DZ}{}\label{cdp:st:C.2.UN.54}
\StatementRecord{cdp:st:C.2.UN.54}{proved}{01DZ}{\cref{lem:II:hopf-curves,lem:II:hopf-curve-restriction,lem:II:support-incidence-residues,cdp:st:C.2.UN.38,cdp:st:C.2.UN.43,cdp:st:C.2.UN.46,cdp:st:C.2.UN.48}}
Let $f:U\to\Delta$ be a proper degeneration of connected
complex two-tori from a smooth threefold, whose full
central fibre $S=\bigcup_vH_v$ is reduced SNC and whose
components are primary or secondary Hopf surfaces.
Its connected double curves are elliptic, it has no
triple points, and its incidence graph is unicyclic.

For these same components and their given curve-gluing
maps, the restrictions \eqref{cdp:eq:C.HR.1} are
isomorphisms, the incidence map \eqref{cdp:eq:C.HR.2}
has rank $|V|-1$, and the endpoint-ratio product around
the cycle is one. In particular
$h^\bullet(S,\OO_S)=(1,2,1)$. A nonzero section of
$\omega_S$ has nonzero residues exactly on the cycle
edges, and its component forms are nonzero exactly on
the cycle components. These are precisely the components
of minimum canonical evaluation order.

When every component has algebraic dimension zero, each
is diagonal and the graph is a cycle with at least two
vertices. More generally, a compact reduced SNC surface
all of whose components are algebraic-dimension-zero
Hopf surfaces is a full fibre of a smooth-total
two-torus degeneration exactly when its components
are diagonal, its graph is such a cycle, its axis
normal lines are dual, and the integral edge product
$M$ of \cref{cdp:eq:C.HC.3,cdp:eq:C.HC.4,cdp:eq:C.HC.5,cdp:eq:C.HC.6}
is the identity. Every such marked support has the
compact realization of \cref{cdp:st:C.2.UN.48}, with
Euler number zero for every reduced fibre support.
\end{proposition}
\begin{proof}
Apply \cref{lem:II:hopf-curves} to every $H_v$.
Its irreducible compact curves are smooth elliptic
and pairwise disjoint. Thus the double curves have
these properties and there are no triple points.
\Cref{lem:II:hopf-curve-restriction} gives
$h^\bullet(H_v,\OO_{H_v})=(1,1,0)$ and a nonzero
restriction for each inclusion $C_e\hookrightarrow H_v$.
At the other endpoint compose its restriction with
pullback by the given elliptic gluing isomorphism.
These are exactly the two endpoint maps in
\eqref{cdp:eq:C.HR.1}. For a secondary component the
same lemma compares them through the finite unramified
primary cover and the component of its inverse curve;
the two finite pullbacks and the restriction square
are its stated maps. The primary constant deck class
restricts with the covering-subgroup multiplicity $d$
in \eqref{cdp:eq:C.HR.4}.

Apply the first clause of
\cref{lem:II:support-incidence-residues} with
$T_v=H_v$ and these $C_e,r_{v,e}$. It gives the
unicyclic graph, rank, exact cycle product, coherent
dimensions, residue support and minimum-order locus.

In algebraic dimension zero,
\cref{cdp:st:C.2.UN.38,cdp:st:C.2.UN.43} give the
nonlinear single elliptic curve and the diagonal
two-axis alternatives, including their finite primary
covers. The inverse image of a quotient curve is a
nonempty union of primary curves, and different
quotient curves have disjoint inverse images.
Since the primary surface has at most two curves,
each quotient component has at most two boundary
curves. A connected unicyclic graph satisfies
$|E|=|V|$, so its average degree is two. Every degree
is therefore two. Its two boundary curves lift to
the two distinct primary axes, so the finite deck
group preserves each axis and the marked component
is diagonal. Thus the graph is a diagonal cycle.
Its components are embedded
and smooth, so an edge joins two distinct components;
there are at least two vertices.

For this cycle the equation $t=xy$ identifies the normal line
bundles of each double curve in its two component surfaces as dual. Removing the axes
gives the open groups and affine logarithmic maps
of \cref{cdp:eq:C.HC.3,cdp:eq:C.HC.4,cdp:eq:C.HC.5,cdp:eq:C.HC.6}.
Their integral product is this same $M$.
\Cref{cdp:st:C.2.UN.46} gives necessity of $M=I$
through the included lattice in the nearby fundamental
group. Conversely its plumbing uses the given component,
axis and normal-product maps, and
\cref{cdp:st:C.2.UN.48} compactifies that family with
the same central support and all gluing constants.
Its Euler conclusion applies to every reduced support
of this compact realization.
\end{proof}













\subsection{Canonical heights and elliptic periods}\label{subsec:II:canonical-tree-heights}\label{subsec:II:hopf-tree-reductions}

\begin{lemma}[Canonical heights and the common elliptic period]\hypertarget{tag-P364}{}\label{lem:II:support-heights-period}
\StatementRecord{lem:II:support-heights-period}{proved}{P364}{\cref{lem:II:support-incidence-residues,lem:II:identity-monodromy-elliptic,lem:II:elliptic-star-boundary,lem:II:hopf-curve-restriction,lem:II:normal-product-collars,cdp:st:C.2.UN.17}}
Let $f:U\to\Delta$ be a proper degeneration of connected
complex two-tori from a smooth threefold, with full
reduced SNC central fibre, disjoint smooth elliptic
double curves, and no triple points. Every component
is either a nonlinear Hopf leaf from
\cref{cdp:eq:C.AS.1,cdp:eq:C.AS.2}, or a smooth elliptic
quasibundle $q_v:T_v\to B_v$ with fixed Jacobian and
identity integral monodromy. Every double curve on a
quasibundle is a fibre support. Consider the following
two domains:
\begin{enumerate}[label=(\arabic*),leftmargin=*]
\item Every $B_v$ is $\PP^1$. Then the incidence graph is
unicyclic. Let $V_0$ be its cycle vertices.
\item The graph is a nontrivial tree, one distinguished
component $P$ has elliptic base $B_P$, and every other
quasibundle base is $\PP^1$. Put $V_0=\{P\}$.
\end{enumerate}
Let $k_v$ be the canonical evaluation orders of
\cref{cdp:st:C.2.UN.17}. Their minimum $k_0$ is attained
exactly on $V_0$. Every multiple fibre of every $q_v$
is a double curve. At an edge $e$, the common exact
normal order $m_e$ is the full fibre multiplicity at
each quasibundle endpoint and the resonance at a
nonlinear endpoint. The dual-normal map identifies
the connected normal-root covers and their Jacobians.
In the unicyclic case the product of these Jacobian
identifications around the cycle is the identity;
the corresponding elliptic product may be a translation.

Root the trees at $V_0$. The heights are
\begin{equation}
h_v:=k_v-k_0=
 \begin{cases}
 0,&v\in V_0,\\
 \displaystyle\sum_{e\text{ on the path from }V_0\text{ to }v}m_e,
       &v\notin V_0.
 \end{cases}
\hypertarget{tag-00NC}{}
\hypertarget{tag-00VZ}{}
\label{cdp:eq:C.HT.2}
\SsixEquationAliases{\label{cdp:eq:C.PT.1}}
\end{equation}
On a quasibundle write the canonical component form as
$\eta_v=q_v^*\theta_v\wedge\alpha_{F_v}$, using the
fixed invariant fibre form of
\cref{lem:II:identity-monodromy-elliptic}. On a cycle
vertex, a coordinate with its two cycle points at
$0,\infty$ gives $\theta_v=c_v\,dz/z$.
On a genus-zero tree vertex, a coordinate with its
parent point at infinity gives $\theta_v=c_v\,dz$.
Each constant is nonzero. In the rooted case
$\theta_P$ is a nowhere-zero holomorphic differential
on $B_P$ and
\begin{equation}
\operatorname{div}(\eta_P)=\sum_{e\ni P}(m_e-1)J_e,
 \qquad K_P\simeq\OO_P\Bigl(\sum_{e\ni P}(m_e-1)J_e\Bigr).
\hypertarget{tag-00W5}{}
\label{cdp:eq:C.PT.2}
\end{equation}

Every regular elliptic fibre of every $q_v$ and every
oriented nonlinear Hopf core represents one transported
integral class $\gamma$ in the nearby two-torus.
For the relative canonical evaluation divided by
$t^{k_0}$, put
\begin{equation}
P(t)=\int_{\gamma_t}\omega_t.
\hypertarget{tag-00ND}{}
\label{cdp:eq:C.HT.3}
\end{equation}
This period extends holomorphically to zero and satisfies
\begin{equation}
P(t)=O(t^{h_v+1})
       \qquad\text{at every quasibundle vertex}.
\hypertarget{tag-00NA}{}
\hypertarget{tag-00WP}{}
\label{cdp:eq:C.HT.11}
\SsixEquationAliases{\label{cdp:eq:C.QT.9}}
\end{equation}
and
\begin{equation}
P(t)=t^{h_v}
 \left(\frac{\pm2\pi i c_v\beta_v^{-m_e}}{n_v}+O(t)\right),
 \qquad c_v\ne0,
 \quad\text{at a nonlinear leaf}.
\hypertarget{tag-00NB}{}
\hypertarget{tag-00W0}{}
\hypertarget{tag-00WE}{}
\label{cdp:eq:C.HT.12}
\SsixEquationAliases{\label{cdp:eq:C.PT.10}\label{cdp:eq:C.QT.10}}
\end{equation}
where $H_{m_e}^{\rm nl}(\beta_v)/\mu_{n_v}$ is the marked
leaf pair and the sign is its orientation sign.
When a nonlinear leaf occurs, $P$ has finite positive
order $R$, every nonlinear leaf has height $R$, and
every quasibundle vertex has height less than $R$.
In the absence of nonlinear leaves, $P$ may vanish
identically.
\end{lemma}
\begin{proof}
Every quasibundle is relatively minimal, since its
vertical curves are elliptic. Apply
\cref{lem:II:identity-monodromy-elliptic} to its
given $q_v$ and full fibre $mJ$. The root quotient
\eqref{cdp:eq:C.HT.4} gives exact normal order $m$
and connected normal-root cover $F_v\to J$.
The nonlinear leaf's normal has order equal to its
resonance, by \cref{def:II:elliptic-star-data}, including
the finite character in \eqref{cdp:eq:C.SN.4}.
The equation $xy=t$ identifies the two normal lines
as dual. Taking their unit-frame covers, and inverting
one normal character, gives the same connected cyclic
cover over this double curve. It identifies their
exact orders and their elliptic Jacobians. Choices
of the root trivialization change these cover maps
only by deck translations.

For a genus-zero quasibundle, \eqref{cdp:eq:C.QT.3}
gives $h^\bullet(\OO)=(1,1,0)$ and nonzero boundary
restrictions through its stated root square.
For a nonlinear leaf these conclusions follow from
\cref{lem:II:hopf-curve-restriction}. In the first
domain apply the unicyclic clause of
\cref{lem:II:support-incidence-residues}; in the
second apply its rooted-tree clause with this $P$.
They give the asserted graph and precisely the
minimum-order set $V_0$.

The adjunction construction of
\cref{cdp:st:C.2.UN.17}, with every full-fibre
multiplicity equal to one, gives a nonzero
meromorphic two-form $\eta_v$ on each component,
with exact divisor
\begin{equation}
\operatorname{div}(\eta_v)
       =\sum_{e=(v,w)}(k_w-k_v-1)J_e.
\hypertarget{tag-00NH}{}
\hypertarget{tag-00WN}{}
\label{cdp:eq:C.HT.7}
\SsixEquationAliases{\label{cdp:eq:C.QT.7}}
\end{equation}
In a double chart this is the residue after dividing
the canonical evaluation by $t^{k_v+1}$.
On the component core the evaluation has no other
zero or pole, so the divisor has exactly this support.
For a quasibundle apply the vertical-form clause of
\cref{lem:II:identity-monodromy-elliptic} to this
particular $\eta_v$. At an unmarked fibre its order
is zero. Formula \eqref{cdp:eq:C.HT.6} would give
$m\operatorname{ord}\theta_v+m-1=0$, which forces
$m=1$. Thus every multiple fibre is marked.
At a marked fibre the same formula gives
\begin{equation}
\begin{gathered}
 d_{v,e}:=\frac{k_w-k_v}{m_e}
       =\operatorname{ord}_{q_v(J_e)}\theta_v+1\in\Z,\\
 \sum_{e\ni v}d_{v,e}=\operatorname{val}(v)-2
       \qquad (B_v=\PP^1).
 \end{gathered}
\hypertarget{tag-00NI}{}
\hypertarget{tag-00NJ}{}
\hypertarget{tag-00WC}{}
\hypertarget{tag-00WO}{}
\label{cdp:eq:C.HT.8}
\SsixEquationAliases{\label{cdp:eq:C.HT.9}\label{cdp:eq:C.PT.9}\label{cdp:eq:C.QT.8}}
\end{equation}
The sum in the second line follows because the
divisor of a rational differential has degree $-2$.
At a smooth unmarked value its order
is zero. Opposite edge orientations negate $d_{v,e}$,
since their common $m_e$ is the exact normal order.

On an edge inside $V_0$ both endpoint orders are
$k_0$, so its difference is zero. A terminal
quasibundle vertex has just its parent edge;
\eqref{cdp:eq:C.HT.9} makes its outward height
increment $m_e$. At a nonlinear leaf apply the
leaf-local clause of \cref{lem:II:elliptic-star-boundary}
with its own neighbouring component as $P$.
Equation \eqref{cdp:eq:C.NS.8} gives the same
increment, with $r=m_e$ and the same evaluation
orders. Now proceed from leaves towards $V_0$.
At a vertex with $c$ child edges of outgoing value
one, \eqref{cdp:eq:C.HT.9} reads
$-d_{\rm parent}+c=c-1$. Therefore the parent
increment also equals $m_e$. This proves
\eqref{cdp:eq:C.HT.2} along every path.

At a cycle vertex the two cycle points have
$d=0$, hence simple poles of $\theta_v$;
each child point has $d=1$ and order zero.
There are no other zeros or poles, giving
$c_v\,dz/z$. At a genus-zero tree vertex the
parent has $d=-1$, hence a double pole at
infinity, and every child has order zero.
This gives $c_v\,dz$. At the elliptic root
all outgoing values are one, so $\theta_P$
has no zero or pole; substitution into
\eqref{cdp:eq:C.HT.7} gives
\eqref{cdp:eq:C.PT.2} as the divisor of that
same canonical section.

In the unicyclic case use the restriction square
of \cref{lem:II:identity-monodromy-elliptic} on
each cycle component. Its composite restriction
and root pullback evaluates the fixed
$H^1(F_v,\OO_{F_v})$ frame. Thus the product of
the given lifted elliptic identifications acts
on this line by the endpoint-ratio product
in \eqref{cdp:eq:C.HR.2}. That product is one
by \cref{lem:II:support-incidence-residues}.
An elliptic isomorphism lifted to universal
covers is $z\mapsto az+b$. It multiplies
the invariant holomorphic one-form by $a$
and its conjugate, representing $H^1(\OO)$,
by $\overline a$. Hence the product has
linear part one and is a translation. This
retains its translation and the given root
cover maps.

To compare integral cycles, pass to the common
normal-root cover at an edge between quasibundles.
Their normal-power trivializations differ by a
nonzero holomorphic function on the compact
double curve, hence by a constant. Multiplying
one base coordinate by that constant gives
the exact collar equations
\begin{equation}
xy=t,\qquad q_v=x^{m_e},\qquad q_w=y^{m_e},
 \qquad x^{m_e}=c\ \longleftrightarrow\ y^{m_e}=t^{m_e}/c.
\hypertarget{tag-00N9}{}
\label{cdp:eq:C.HT.10}
\end{equation}
The two levels are the same connected elliptic
cover with its complex orientation. Varying
the nonzero normal level gives an oriented
isotopy between them. Apply
\cref{lem:II:normal-product-collars} to this
same full reduced SNC support, its normal
line product maps, and these component cores.
Its product-preserving smooth comparison
carries the collar isotopies to the given
smoothing, and its core transport agrees
on both attaching ends. Within each
quasibundle the regular fibres are isotopic
along paths avoiding the finite marked set.
At a nonlinear leaf the boundary map
\eqref{cdp:eq:C.SN.6} in
\cref{lem:II:elliptic-star-boundary} sends
the Hopf fibre generators to the same $e,h$
and its meridian to the product-filling
meridian. Its smooth pair map identifies
the oriented Hopf core with that regular
elliptic fibre. These maps preserve the
full secondary integer $n_v$. Composing
along edges proves equality of the stated
integral classes in the nearby fibre.

On a compact regular fibre in $T_v$, the
leading coefficient of the relative form
is $t^{h_v}\eta_v$. Its restriction to
that holomorphic curve is zero. Smooth
compact-core transport therefore gives
\eqref{cdp:eq:C.HT.11}. At a nonlinear
leaf apply the individual period clause
of \cref{lem:II:elliptic-star-boundary}
with its own parent order $k_{\mathrm{par}}$.
Multiplying its expression by
$t^{k_{\mathrm{par}}-k_0}$ changes the leading
power from $r$ to $k_v-k_0=h_v$ and gives
\eqref{cdp:eq:C.HT.12}, with exactly its
primitive core lattice and coefficient
$\pm2\pi i c_v\beta_v^{-r}/n_v$.
Differentiation under the integral and
Stokes' theorem show that all these
integrals are the same holomorphic period
on the punctured disc. Their compact-core
expressions give a holomorphic extension
at zero and the displayed orders. A
nonlinear coefficient is nonzero, so one
leaf fixes its order $R=h_v>0$; the other
leaf expansions give the same height,
and \eqref{cdp:eq:C.HT.11} makes every
quasibundle height smaller.
\end{proof}

\begin{proposition}[Canonical heights of Hopf support trees]\hypertarget{tag-01E3}{}\label{cdp:st:C.2.UN.58}
\StatementRecord{cdp:st:C.2.UN.58}{proved}{01E3}{\cref{cdp:st:C.2.UN.54,cdp:st:C.2.UN.38,cdp:st:C.2.UN.37,lem:II:support-heights-period,def:II:elliptic-star-data}}
Let $S$ be the full reduced SNC central fibre of a
proper degeneration of connected complex two-tori
from a smooth threefold, with every component a
primary or secondary Hopf surface. A diagonal
algebraic-dimension-zero component forces the whole
support to be a diagonal cycle classified by
\cref{cdp:st:C.2.UN.46,cdp:st:C.2.UN.48}.

In every other case the graph is unicyclic. Every
cycle vertex and every nonterminal tree vertex
has algebraic dimension one and its elliptic
reduction is
\begin{equation}
                    q_v:H_v\longrightarrow\PP^1,        \hypertarget{tag-00N8}{}\label{cdp:eq:C.HT.1}
\end{equation}
with fixed Jacobian, identity integral monodromy
and only smooth multiple fibre supports. Every
other vertex is a nonlinear algebraic-dimension-zero
Hopf leaf or has this same elliptic reduction.
Every multiple fibre is a double curve.

The exact normal order $m_e$ on an edge equals
the full multiplicity at each elliptic endpoint
and the resonance at a nonlinear endpoint.
The given dual-normal map identifies the
connected normal-root covers and their Jacobians.
The Jacobian product around the cycle is the
identity; the elliptic product may be a translation.

For the canonical evaluation orders, put
$k_0=\min k_v$. This minimum is attained exactly
on the cycle, and the heights are
\eqref{cdp:eq:C.HT.2}, with $V_0$ its cycle vertices.
After factoring out the invariant relative
elliptic form, the component base differential
is $c_v\,dz/z$ on a cycle vertex with its cycle
points at $0,\infty$, and $c_v\,dz$ on a tree
vertex with its parent point at infinity,
where $c_v\ne0$.

All regular elliptic fibres and nonlinear
cores represent the integral class $\gamma$
of \cref{lem:II:support-heights-period}.
For its period \eqref{cdp:eq:C.HT.3}, a nonlinear
leaf forces a finite positive order $R$.
Every nonlinear leaf has height $R$, and
every algebraic-dimension-one vertex has
height less than $R$. With no nonlinear
leaf this period may vanish identically.
\end{proposition}
\begin{proof}
Apply the normal-character calculation in
\cref{cdp:st:C.2.UN.38} to the same diagonal
primary surface and its marked axis $w=0$.
Its normal is exactly the quotient line
$(z,v)\sim(\alpha z,\beta v)$ in that
calculation, so it has infinite order in
algebraic dimension zero. Interchanging
the coordinates gives the other axis.
A torsion line on
a secondary surface pulls back to a torsion
line on its finite primary cover, so the
axis normals there are also nontorsion.

For an algebraic-dimension-one Hopf surface,
\cref{cdp:st:C.2.UN.37} gives the stated
elliptic reduction. A horizontal compact
curve would meet a regular fibre, whereas
$H^2(H,\mathbb Q)=0$ makes their intersection
zero. Positivity of local intersections
therefore makes every curve a fibre support.
Apply \cref{lem:II:identity-monodromy-elliptic}
to this reduction: every boundary normal
has finite exact order equal to its full
fibre multiplicity. For a nonlinear Hopf
axis, \cref{def:II:elliptic-star-data} gives
the finite exact resonance order, including
its secondary normal character.

At an SNC edge the map $xy=t$ identifies
the normals as dual and preserves whether
their tensor powers are trivial. Hence
a diagonal algebraic-dimension-zero vertex
can meet only another such vertex.
Connectedness makes all components of that
type, and \cref{cdp:st:C.2.UN.54} gives the
diagonal-cycle conclusion.

Otherwise each algebraic-dimension-zero
component is nonlinear and has its one
elliptic curve, so it is a leaf. The
remaining components have the reductions
\eqref{cdp:eq:C.HT.1}. Their full support
therefore satisfies the first domain of
\cref{lem:II:support-heights-period}, with
$T_v=H_v$ and these given reductions,
double-curve maps and normal dualities.
That lemma proves absence of unmarked
multiple fibres, the normal-root and
Jacobian comparisons, the stated heights
and base differentials, and the integral
core and period conclusions. Its nonlinear
coefficient uses the full primitive lattice
of each given $H_r(\beta)/\mu_n$; it is
\eqref{cdp:eq:C.HT.12} for these same leaves.
\end{proof}













\subsection{Pruning and compact realization of elliptic trees}\label{subsec:II:tree-pruning-remainder}

\begin{proposition}[Unicyclic elliptic supports]\hypertarget{tag-01E6}{}\label{cdp:st:C.2.UN.60}
\StatementRecord{cdp:st:C.2.UN.60}{proved}{01E6}{\cref{lem:II:identity-monodromy-elliptic,lem:II:support-incidence-residues,lem:II:support-heights-period,lem:II:elliptic-tree-pruning,lem:II:compact-elliptic-star-insertion,cdp:st:C.2.UN.57}}
\hypertarget{tag-01E8}{}\label{cdp:st:C.2.UN.62}
\StatementRecord{cdp:st:C.2.UN.62}{proved}{01E8}{}
Let $S$ be a compact connected reduced SNC surface with smooth
embedded components and given elliptic gluing maps. Each component
is a primary or secondary Hopf surface, or is equipped with a
fixed-Jacobian elliptic quasibundle
\begin{equation}
                         q_v:T_v\longrightarrow\PP^1
\hypertarget{tag-00WD}{}\label{cdp:eq:C.QT.1}
\end{equation}
with connected fibres, identity integral monodromy and only smooth
multiple fibre supports. Every boundary curve of such a $T_v$ is
a fibre support of its given map. The analytic torsors and gluing
maps are part of the data.

Occurrence of $S$ as a full reduced fibre of a proper degeneration
of connected complex two-tori from a smooth total space is
equivalent to the following conditions.
\begin{enumerate}
\item The two normal line bundles at each double curve are dual,
and the incidence graph has exactly one independent cycle.
\item A diagonal Hopf component of algebraic dimension zero belongs
to a support consisting entirely of its diagonal cycle, with
integral complement-lattice product $M=I$.
\item In the remaining case every algebraic-dimension-zero Hopf
component is a nonlinear leaf. Every other component has its
elliptic map to $\PP^1$, and all its multiple fibres are marked
double curves. Let $m_e$ be the exact order of the edge normal
line. It is the full multiplicity at an elliptic endpoint and
the resonance at a nonlinear endpoint.
\item In this elliptic case set $h_v=0$ on the cycle and sum the
$m_e$ along the path from the cycle to each other vertex.
There is an integer $R>0$ such that every nonlinear leaf has
$h_v=R$ and every elliptic vertex has $h_v<R$. With no nonlinear
leaf, any integer $R>\max_v h_v$ meets this condition.
\item The pruning maps of \cref{lem:II:elliptic-tree-pruning}, with
their primitive fibre bases, full secondary integers and torsor
transitions, give a marked Hopf/ruled core cycle with $M=I$.
\end{enumerate}
The last three conditions concern the elliptic case. Their
preceding geometric conditions ensure that every pruning step
exists. Disjoint steps give the same marked core up to the
isomorphisms of \cref{lem:II:elliptic-tree-pruning}.

Every such support has a realization over $\PP^1$ with smooth
compact total threefold, preserving all its components, analytic
torsors and elliptic gluing maps. Its only other singular fibre
support is the conjugate pruned core. Every reduced fibre support
has Euler characteristic zero. The quasibundle components may
have Kodaira dimension zero or one as well as the Hopf and ruled
types described in the proof.
\end{proposition}
\begin{proof}
The curves on each Hopf component are disjoint smooth elliptic
curves by \cref{lem:II:hopf-curves}; on a quasibundle they are
distinct smooth fibre supports. Thus there are no triple points.
A Hopf component of positive algebraic dimension has the elliptic
reduction of \cref{cdp:st:C.2.UN.37}. Its boundary curves are fibres:
$H^2(H,\Q)=0$, whereas a horizontal curve would have positive
intersection with a fibre. Include these reductions among the
maps used below. For every map \eqref{cdp:eq:C.QT.1}, apply
\cref{lem:II:identity-monodromy-elliptic} to the same surface,
base and full marked fibres. Its root quotient
\eqref{cdp:eq:C.QT.2} gives their exact normal order, and
\eqref{cdp:eq:C.QT.3} and its commuting root-restriction square
give coherent dimensions $(1,1,0)$ and nonzero endpoint maps.
For a Hopf component use \cref{lem:II:hopf-curve-restriction}
on each of its curve inclusions. Compose these maps with the
given elliptic gluing at each edge.

For a full reduced torus degeneration, apply the unicyclic clause
of \cref{lem:II:support-incidence-residues} to those same component
and restriction maps. Its normalization sequence
\eqref{cdp:eq:C.QT.4} gives one cycle, coherent dimensions
\eqref{cdp:eq:C.QT.5}, product-one coherent holonomy, and the
minimum canonical order precisely on the cycle. The normal-product
function identifies the two normal lines as inverses.
A diagonal algebraic-dimension-zero Hopf axis has infinite normal
order by \cref{cdp:st:C.2.UN.58}. The normal lines at a quasibundle
fibre and at a nonlinear Hopf axis have finite orders, by the
root chart above and \eqref{cdp:eq:C.SN.4}. Since taking the inverse
of a line bundle preserves its order, normal duality propagates
the diagonal algebraic-dimension-zero type along every adjacent
edge. Connectedness gives an all-diagonal support. The cycle
and its complete compact criterion are then
\cref{cdp:st:C.2.UN.57}, with precisely these marked axes and
normal-product maps.

In the remaining case each nonlinear Hopf has its unique elliptic
axis and is a leaf. Apply the first domain of
\cref{lem:II:support-heights-period} to these reductions, the given
maps \eqref{cdp:eq:C.QT.1}, and the same normal-product smoothing.
Its component divisor \eqref{cdp:eq:C.QT.7}, edge equations
\eqref{cdp:eq:C.QT.8} and connected normal-root maps give all the
multiplicities, marked fibres, common Jacobian and path heights.
Its cycle and tree base forms are $c\,dz/z$ and $c\,dz$ on these
same maps. The root-square comparison identifies the lifted cycle's
linear part with the identity while carrying its translation.

The lemma's collar maps, including the full secondary basis in
\eqref{cdp:eq:C.SN.6}, identify the regular fibre classes with the
one integral core class. Its expansions
\cref{cdp:eq:C.QT.9,cdp:eq:C.QT.10}, after division by the cycle
minimum, give one finite order $R$ at every nonlinear leaf and a
strictly smaller height at each elliptic vertex. With no nonlinear
leaf choose an integer above the finitely many heights.
This proves the first four conditions.

Apply \cref{lem:II:elliptic-tree-pruning} to the same component
maps and dual-normal identifications. Its one-mark classification
uses each modified component's full root-line class, including
components that originally had several multiple fibres.
The map \eqref{cdp:eq:C.PR.6} and its inverse carry all other
transitions unchanged. The two-mark classification gives the
marked diagonal Hopf or torsion ruled core. The boundary homomorphism
\eqref{cdp:eq:C.PR.9} identifies the nearby group with the core
plumbing group. The injected semidirect product
$\Lambda\rtimes_M\Z$ of \cref{cdp:st:C.2.UN.57} is therefore
isomorphic to $\Z^4$, forcing $M=I$.

Conversely, in the elliptic case use the same pruning sequence
and its inverse maps in \cref{lem:II:compact-elliptic-star-insertion}.
The core is its cyclic core, the normal order is $m_e$ on each
edge, and its height is exactly $h_v$. Thus its local displacement
at a vertex is of order $R-h_v$ and its strict-parent map is the
inverse of \eqref{cdp:eq:C.PR.6}. Its compact construction restores
every analytic torsor, multiple fibre, marked point and gluing
translation, with multiplicity one. The initial conjugate core
at infinity is unchanged. The diagonal algebraic-dimension-zero
case already has this compact realization by
\cref{cdp:st:C.2.UN.57}. The Euler conclusions follow from these
two constructions, completing both directions.
\end{proof}






\begin{proposition}[Rooted elliptic supports]\hypertarget{tag-01E7}{}\label{cdp:st:C.2.UN.61}
\StatementRecord{cdp:st:C.2.UN.61}{proved}{01E7}{\cref{lem:II:identity-monodromy-elliptic,lem:II:support-incidence-residues,lem:II:support-heights-period,lem:II:elliptic-tree-pruning,lem:II:elliptic-star-boundary,lem:II:zero-chern-elliptic-torsor,lem:II:compact-elliptic-star-insertion,cdp:st:C.2.UN.31}}
Let $S$ be a compact connected reduced SNC surface whose incidence
graph is a nontrivial tree, with a distinguished smooth component
$P$. Every other component is a primary or secondary Hopf surface,
or has a fixed-Jacobian elliptic quasibundle map to $\PP^1$ with
connected fibres, identity integral monodromy and only smooth
multiple supports. The boundary curves of a quasibundle are fibre
supports of its given map. The elliptic gluing maps and analytic
torsors are included in the data.

Occurrence of $S$ as a full fibre of a degeneration of connected
complex two-tori from a smooth total space is equivalent to:
\begin{enumerate}
\item The two normal lines at every double curve are dual.
\item The root has a fixed-Jacobian elliptic quasibundle
$q_P:P\to B$ over a smooth elliptic curve, with identity integral
monodromy and only smooth multiple supports. Its boundary curves
are fibres and every multiple fibre is marked. Every nonroot
component is a nonlinear algebraic-dimension-zero Hopf leaf, or
an elliptic quasibundle over $\PP^1$ of the indicated type whose
multiple fibres are all marked.
\item Each edge normal has exact finite order $m_e$, equal to
its full elliptic multiplicity or nonlinear resonance.
For the heights \eqref{cdp:eq:C.PT.1}, with $V_0=\{P\}$, there is
an integer $R>0$ with $h_v=R$ at every nonlinear leaf and $h_v<R$
at every elliptic nonroot vertex. With no nonlinear leaves any
$R>\max_v h_v$ suffices.
\item The pruning maps of \cref{lem:II:elliptic-tree-pruning} give
a principal elliptic torsor $T_0\to B$ whose integral Chern vector
in $H^2(B,\Lambda_F)$ is zero. This is the torsor defined by the
complete inverse charts, primitive bases and secondary integers.
\end{enumerate}
The connected normal-root maps identify the fixed Jacobian $F$
along the tree. Disjoint pruning steps give the same marked
torsor up to the isomorphisms of \cref{lem:II:elliptic-tree-pruning}.

Every support satisfying these conditions has a realization over
$\PP^1$ with smooth compact total threefold, with the given $S$
as its only singular fibre and all other fibres smooth two-tori.
Every component, analytic torsor and elliptic gluing map is
preserved. All reduced fibre supports have Euler number zero.
In every occurrence the root is minimal and
$h^q(P,\OO_P)=h^q(S,\OO_S)=(1,2,1)$. It is the unique minimum-order
component for canonical evaluation, its canonical divisor is
\eqref{cdp:eq:C.PT.2}, and all order differences are the heights
\eqref{cdp:eq:C.PT.1}. The zero-edge case is the smooth torus
fibre of \cref{cdp:st:C.2.UN.5}.
\end{proposition}
\begin{proof}
Each double curve has a nonroot endpoint. On a Hopf component,
\cref{lem:II:hopf-curves} makes its curves smooth elliptic and
pairwise disjoint. On a quasibundle they are distinct smooth
fibre supports. Hence there are no triple points and the curves
on $P$ are also disjoint.

For every Hopf nonroot component apply
\cref{lem:II:hopf-curve-restriction} to its curve inclusions;
for every quasibundle apply
\cref{lem:II:identity-monodromy-elliptic} to its given map
$q_v:T_v\to\PP^1$. Its coherent formula and root-restriction
square give dimensions $(1,1,0)$ and nonzero endpoint maps through
the connected normal-root covers. Compose with the given gluing
map at each edge. The rooted clause of
\cref{lem:II:support-incidence-residues} applies with the unrestricted
root $P$. It identifies the normalization sequence
\eqref{cdp:eq:C.PT.3} and restriction isomorphisms
\eqref{cdp:eq:C.PT.4} with those of this $S$.

Semicontinuity gives $h^2(S,\OO_S)\ge1$, and
\cref{cdp:st:C.2.UN.28} gives $p_g(P)\le1$; hence $p_g(P)=1$.
Flatness gives $\chi(\OO_S)=0$. Every nonroot component and every
edge has coherent Euler characteristic zero, so the normalization
sequence gives $\chi(\OO_P)=0$ and $h^1(P,\OO_P)=2$.
The same rooted clause gives \eqref{cdp:eq:C.PT.5} by restriction
to $P$ and extension by zero on the other components.
Its canonical evaluation map \eqref{cdp:eq:C.HR.8} makes $P$
the unique minimum-order component. With minimum $k_0$, its
holomorphic two-form has divisor
\begin{equation}
 K_P=\sum_{e=(P,w)}(k_w-k_0-1)J_e,\qquad k_w>k_0.
\hypertarget{tag-00W9}{}\label{cdp:eq:C.PT.6}
\end{equation}
All coefficients are nonnegative. A rational exceptional curve
$E$ with $E^2=-1$ would have $K_P\cdot E=-1$ by adjunction,
whereas its intersections with these distinct elliptic curves
are nonnegative. Thus $P$ is minimal.

The nonzero canonical form gives $\kappa(P)\ge0$. The component
multiplicities are one and the total space is smooth, so the
common-root map of \cref{cdp:st:C.2.UN.31} has degree one over
$P$. Its component is $P$ itself. That theorem makes $P$ a
two-torus or gives it a fixed-modulus, identity-monodromy elliptic
quasibundle with base genus $p_g(P)=1$.

On a two-torus an elliptic curve is a translate of an elliptic
subtorus: the derivative of its lift is periodic and holomorphic,
hence constant. Two distinct elliptic directions give an isogeny
from their product onto the torus, so any translates intersect.
The disjoint boundary curves therefore have one direction and
are fibres of the quotient by that subtorus.
For the other elliptic map $q:P\to B$, a horizontal marked elliptic
curve maps finitely and unramifiedly to $B$ by Riemann--Hurwitz.
At a full multiple fibre the local base function is $s^m$, forcing
ramification on such a curve. Thus the presence of a multiple
fibre makes every marked curve vertical. With no multiple fibre,
$q$ is a principal $F$-torsor. A horizontal elliptic curve gives a
section after its finite base pullback. This pullback multiplies
the integral Chern vector by its positive degree in the torsion-free
group $H^2(B,\Lambda_F)$, so that vector is zero.
Apply \cref{lem:II:zero-chern-elliptic-torsor} to
$(E_1,B,F_1)=(P,B,F)$. Its translation cochain and map
\eqref{cdp:eq:C.PT.8} identify this analytic torsor over $B$ with
a two-torus, to which the preceding quotient argument applies.
In all cases there is a map $q_P$ with every marked curve as a
fibre support.

The connected root chart of \cref{lem:II:identity-monodromy-elliptic},
now on this same $q_P$, gives finite normal order at each root mark.
The dual-normal isomorphism carries that order to its neighbour.
The infinite-order diagonal axes in \cref{cdp:st:C.2.UN.58} therefore
cannot occur at the first such vertex along a path from $P$.
Every nonlinear Hopf has one elliptic axis and is a leaf.
Every other nonroot component has its given genus-zero map or the
Hopf reduction of \cref{cdp:st:C.2.UN.37}; the latter's boundary
curves are fibres by $H^2(H,\Q)=0$ and positivity of a horizontal
intersection with a fibre.

Apply the rooted domain of \cref{lem:II:support-heights-period}
to these maps, $q_P$, and the same normal-product smoothing.
It identifies all connected normal roots and their fixed Jacobian,
marks every multiple fibre, and gives the edge equations
\eqref{cdp:eq:C.PT.9}, hence heights \eqref{cdp:eq:C.PT.1} and the
root divisor \eqref{cdp:eq:C.PT.2}. Its collar map
\eqref{cdp:eq:C.HT.10} and full secondary basis
\eqref{cdp:eq:C.SN.6} identify every regular fibre class, including
the root's class, with the integral core class. Its expansions
\cref{cdp:eq:C.HT.11,cdp:eq:C.PT.10} are therefore expansions of
one period of the evaluation divided by $t^{k_0}$. They give the
common $R$ and every strict height inequality.

Prune using \cref{lem:II:elliptic-tree-pruning} on these same maps.
Each step changes one multiple-fibre tube by
\eqref{cdp:eq:C.PR.6}, or removes a product cap by the identity
product map. All other transitions stay fixed. The final root
is a principal torsor $T_0\to B$ with the resulting analytic atlas.
The boundary map \eqref{cdp:eq:C.PR.9}, carried through the smooth
collars of the given degeneration, gives
\begin{equation}
                    \pi_1(T_0)\simeq\pi_1(F_t)\simeq\Z^4.
\hypertarget{tag-00W1}{}\label{cdp:eq:C.PT.11}
\end{equation}
Apply the principal-bundle group clause of
\cref{lem:II:elliptic-star-boundary} to this $T_0\to B$.
In its exact sequence \eqref{cdp:eq:C.PT.12}, the fibre lattice
is injected and the base commutator is the integral Chern vector.
Equation \eqref{cdp:eq:C.PT.11} makes that vector zero.
The map \eqref{cdp:eq:C.PT.8} of
\cref{lem:II:zero-chern-elliptic-torsor}, with
$(E_1,B,F_1)=(T_0,B,F)$, then identifies the same analytic torsor
with its two-torus quotient. This proves the last necessary
condition, and the pruning lemma gives the stated independence.

Conversely use this zero-Chern inverse torsor, the complete inverse
pruning maps and the heights in
\cref{lem:II:compact-elliptic-star-insertion}. Its root is exactly
$T_0$, its initial period has order $R$, and each next displacement
has order $R-h_v$ by \eqref{cdp:eq:C.PR.15} or
\eqref{cdp:eq:C.PR.16}. Its reverse maps recover $P$ and all other
components with their full torsors and marked gluing maps.
Its proper compact realization has this reduced $S$ as its only
singular fibre and smooth two-tori elsewhere. Its Euler conclusion
applies to the same components and double curves, proving all
remaining assertions.
\end{proof}






\begin{proposition}[Plurigenera and a four-fibre elliptic component]\hypertarget{tag-P376}{}\label{prop:II:elliptic-quasibundle-plurigenera}
\StatementRecord{prop:II:elliptic-quasibundle-plurigenera}{proved}{P376}{\cref{lem:II:identity-monodromy-elliptic,lem:II:compact-elliptic-star-insertion,cdp:st:C.2.UN.62}}
For a genus-zero identity-monodromy quasibundle
$q:T\to\PP^1$ with full multiple fibres $m_iJ_i$, its plurigenera are
\begin{equation}
 h^0(T,\omega_T^{\otimes a})
 =\max\left\{0,1-2a+\sum_i
              \left\lfloor a(1-1/m_i)\right\rfloor\right\},
                                      \qquad a\ge1.
\hypertarget{tag-00WG}{}\label{cdp:eq:C.QT.12}
\end{equation}
There is a compact smooth-total two-torus degeneration whose
reduced singular fibre has a quasibundle component with four
triple fibres and Kodaira dimension one, one elliptic ruled
component, and four nonlinear primary Hopf leaves of resonance
three. All its elliptic gluing maps are the maps below.
\end{proposition}
\begin{proof}
The invariant vertical differential and root charts of
\cref{lem:II:identity-monodromy-elliptic}, applied to this $q$,
identify a pluricanonical section with a meromorphic base
$a$-differential times the $a$th invariant fibre factor.
At $t=s^{m_i}$ its order is
$m_i\operatorname{ord}_{t=0}\theta+a(m_i-1)$.
Holomorphy therefore permits a base pole of order at most
$\lfloor a(1-1/m_i)\rfloor$. Descent in this invariant chart
and pullback by the same root map are inverse to one another.
The resulting line on $\PP^1$ has degree
$-2a+\sum_i\lfloor a(1-1/m_i)\rfloor$, proving
\eqref{cdp:eq:C.QT.12}.

Fix $F=\C^*/\beta^{\Z}$ with $0<|\beta|<1$. Take two copies of
$F\times\PP^1$ and glue their fibres at zero and infinity by the
identity on $F$. This marked ruled cycle has $M=I$.
On one component choose four distinct ordinary base points away
from the two cycle marks. Use four weight-three insertions with
secondary integer one and period order $R=3$ in
\cref{lem:II:compact-elliptic-star-insertion}. They add the
nonlinear primary leaves $H_3(\beta)$.

The modified quasibundle component is explicitly the smooth coarse
surface of $L^\times/\beta^{\Z}$, where
\begin{equation}
 \mathcal B=\PP^1[(p_1,3),\ldots,(p_4,3)],\qquad
 L=\OO_{\mathcal B}(D_1+\cdots+D_4),\qquad 3D_i=p_i.
\hypertarget{tag-00WF}{}\label{cdp:eq:C.QT.11}
\end{equation}
Every inertia acts by a primitive fibre translation, giving a
smooth surface and exactly four triple fibres. In a root chart
$p=s^3$, its line coordinate is $y=sW$, with $W$ the product
coordinate before insertion. The inverse map
\eqref{cdp:eq:C.PR.6}, with $m=3,n=1$, is $W=y/s$.
For the positive divisor line the frames are
$e_{\rm in}=s^{-1}e_{\rm out}$, and the equality
$W e_{\rm out}=y e_{\rm in}$ gives exactly $y=sW$.
Its canonical section has coefficient $s$ in the inner frame.
These frame and inverse-root maps identify the complete line
torsor with the component in the compact construction.

The formula already proved gives
\begin{equation}
 h^0(T,\omega_T^{\otimes a})
       =\max\{0,1-2a+4\lfloor2a/3\rfloor\},\qquad
              \omega_T^{\otimes3}\simeq q^*\OO_{\PP^1}(2).
\hypertarget{tag-00WH}{}\label{cdp:eq:C.QT.13}
\end{equation}
Thus $p_g(T)=0$, $q(T)=1$, $\chi(\OO_T)=0$ and
$h^0(T,\omega_T^{\otimes3b})=2b+1$. The last equality gives
Kodaira dimension one. The other ruled component is unchanged;
the four Hopf leaves and every gluing are those of the displayed
inverse root charts. This is the asserted full compact fibre.
\end{proof}






\begin{definition}[Weighted elliptic support data]\hypertarget{tag-P377}{}\label{def:II:weighted-elliptic-support}
\StatementRecord{def:II:weighted-elliptic-support}{proved}{P377}{\cref{def:II:weighted-normal-cycle}}
Let $S=\bigcup_iS_i$ be a connected SNC surface with smooth embedded
components, with every edge joining distinct components. Assign
positive multiplicities $m_i$. A component is a primary or secondary
Hopf surface, or has a fixed-Jacobian elliptic quasibundle over a
smooth compact curve, with connected fibres, identity integral
monodromy and only smooth multiple supports. Its marked boundary
curves are fibres of that map. The data include the component
surfaces, analytic torsors and elliptic gluing maps $\phi_{ij}$.

Weighted normal data are holomorphic lines $\mathcal L_i$, sections
and edge isomorphisms
\begin{equation}
\begin{gathered}
 \sigma_i\in H^0(S_i,\mathcal L_i^{-m_i}),\qquad
 \operatorname{div}\sigma_i=\sum_{j\sim i}m_jJ_{ij},\\
 \mathcal L_i|_{J_{ij}}\simeq\phi_{ij}^*N_{J_{ji}/S_j},\qquad
 \phi_{ij}^*(\mathcal L_j|_{J_{ji}})\simeq N_{J_{ij}/S_i},
\end{gathered}
\hypertarget{tag-017O}{}\label{cdp:eq:C.WT.1}
\end{equation}
whose leading section coefficients agree at each edge.
On an oriented cycle of the component types in
\cref{def:II:weighted-normal-cycle}, the substitutions
$\mathcal N_i=\mathcal L_i$ and $C_i^\pm=J_{i,i\pm1}$ identify
these sections and edge maps with that definition.
The strict root components are the connected components of
\begin{equation}
 \mathcal E_i=\operatorname{Norm}
 \{\xi\in\mathcal L_i^{-1}:\xi^{m_i}=\sigma_i\},\qquad
                         N=\operatorname{lcm}_i m_i.
\hypertarget{tag-017Z}{}\label{cdp:eq:C.WT.2}
\end{equation}
Their boundary maps and cyclic $\boldsymbol\mu_N$-action use the
same normal-root identifications. For an edge with
$a=gA,b=gB$, $\gcd(A,B)=1$, its compressed neck is
\begin{equation}
 uv=s^h,\qquad x=u^B,\quad y=v^A,\qquad
 h=\frac{N}{gAB},\qquad t=s^N=x^ay^b.
\hypertarget{tag-018B}{}\label{cdp:eq:C.WT.3}
\end{equation}
The root graph has these strict components as vertices and these
lifted curves as edges, with their indicated exponents $h$.
\end{definition}

\begin{proposition}[Weighted elliptic supports]\hypertarget{tag-01E9}{}\label{cdp:st:C.2.UN.63}
\StatementRecord{cdp:st:C.2.UN.63}{proved}{01E9}{\cref{def:II:weighted-elliptic-support,lem:II:weighted-cycle-roots,lem:II:support-incidence-residues,lem:II:support-heights-period,lem:II:elliptic-tree-pruning,lem:II:elliptic-core-period-charts,lem:II:weighted-elliptic-insertion,lem:II:finite-elliptic-insertion-gauge,cdp:st:C.2.UN.59}}
For the data of \cref{def:II:weighted-elliptic-support}, occurrence
as a fibre of a degeneration of connected complex two-tori from
a smooth total space is equivalent to the following conditions.
\begin{enumerate}
\item The root graph is connected. A branch with a diagonal
algebraic-dimension-zero Hopf component is its complete two-axis
cycle and satisfies \cref{cdp:st:C.2.UN.59}. In the remaining case
every elliptic strict root cover has base genus at most one, and
at most one has genus one.
\item With all root bases of genus zero, the compressed root graph
is unicyclic. With a genus-one component, it is a tree rooted there.
Every other root component is a genus-zero elliptic quasibundle
or a nonlinear Hopf leaf. All multiple elliptic fibres are marked,
and the two normal lines at every root edge are dual.
\item Let $m_e$ be their exact torsion order and $h_e$ the exponent
in \eqref{cdp:eq:C.WT.3}. Give the cycle or genus-one root height
zero and put
\begin{equation}
 H_v=\sum_{e\text{ on the rooted path to }v}h_em_e.
\hypertarget{tag-018M}{}\label{cdp:eq:C.WT.4}
\end{equation}
Every nonlinear leaf has one common height $R>0$, and every
elliptic vertex has height less than $R$. With no nonlinear leaf,
$R$ may be chosen larger than all these heights.
\item Pruning by \cref{lem:II:elliptic-tree-pruning} in complete
cyclic-group orbits, with the full torsor maps, secondary integers
and primitive fibre bases, gives a cycle with rank-three complement
product $M=I$, or a principal elliptic torsor over its elliptic
base with integral Chern vector zero.
\item Let $D$ be the finite linear part on the pruned core's
common logarithmic vector space or torus. For the root generator
$s\mapsto\zeta_Ns$ the height satisfies
\begin{equation}
                              \zeta_N^R=\det_{\C}D.
\hypertarget{tag-018R}{}\label{cdp:eq:C.WT.5}
\end{equation}
With no nonlinear leaf, take any sufficiently large positive
$R$ in this character class.
\end{enumerate}
The cycle action may reverse orientation through vertices.
Every support satisfying these conditions has a realization over
$\PP^1$ with smooth compact total threefold, the given weighted
support at zero and its complex conjugate at infinity, and smooth
connected two-tori at every other fibre. It preserves all component
surfaces, analytic torsors, elliptic gluing maps, weighted normal
data and multiplicities. The normal-product constants are chosen
in the construction. Every reduced fibre support has Euler number
zero.
\end{proposition}
\begin{proof}
A fibre function gives \eqref{cdp:eq:C.WT.1} by division by the
$m_i$th power of its component's canonical divisor section.
On the strict root, with normal divisor coordinate $v_i$, one has
$\xi=s^{N/m_i}/v_i$. The root generator is therefore
\begin{equation}
                         \xi\longmapsto\zeta_N^{N/m_i}\xi.
\hypertarget{tag-018S}{}\label{cdp:eq:C.WT.6}
\end{equation}
At an edge, absorb the unit in the local function to obtain
$x^ay^b=s^N$. The monomial normalization of
\cref{lem:II:weighted-cycle-roots}, applied to these same two normal
coordinates, gives its $g$ branches and exactly
\eqref{cdp:eq:C.WT.3}. The leading-coefficient isomorphisms in
\eqref{cdp:eq:C.WT.1} identify their curve covers and normal characters.
The product $uv=s^h$ trivializes the product of the two strict
normal lines on the lifted curve, giving their duality.

Each strict component is smooth. Near a marked curve, write its
root equation as $\xi^m=z^b$ after absorbing a unit. Set
$d=\gcd(m,b)$, $q=m/d$, $B=b/d$. Each branch has normalization
\begin{equation}
                     z=w^q,\qquad\xi=cw^B,
\hypertarget{tag-018T}{}\label{cdp:eq:C.WT.7}
\end{equation}
with $c$ the corresponding root of unity and the curve coordinate
unchanged. Away from these curves the component map is finite
\'{e}tale. The reduced boundary maps are thus finite \'{e}tale
maps of elliptic curves; their stabilizers act by deck translations,
as do their connected normal-root elliptic covers.

For an elliptic component, delete the finitely many marked base
values and take a connected component of its finite surface cover.
Its fibre subgroup is a finite-index lattice
$\Lambda'\subset\Lambda_F$. On its connected-fibre Stein base
cover the fibre is $F'=\C/\Lambda'$, with isogeny $\pi:F'\to F$.
A lifted transition satisfies
\begin{equation}
                                  \pi a=T_b\pi.
\hypertarget{tag-018U}{}\label{cdp:eq:C.WT.8}
\end{equation}
Differentiating through this isogeny gives identity derivative
for $a$, so every lifted transition is a translation, including
those over loops that permuted components of a fibre cover.
At a deleted value lift the root chart of
\cref{lem:II:identity-monodromy-elliptic} for the original component
and apply \eqref{cdp:eq:C.WT.7} to its normal power. The subgroup
trivial on $F'$ acts by normal pseudoreflections; after its quotient,
the remaining inertia acts freely by translations on $F'$.
This gives a smooth multiple chart
$(F'\times\Delta)/\langle(\tau',\zeta_{m'})\rangle$.
Apply the translation gauge and invariant cohomology frame of
\cref{lem:II:identity-monodromy-elliptic} to this compact strict
cover over its smooth Stein base $B'$. It gives
\begin{equation}
                         R^1q_*\OO=\OO_{B'},\qquad p_g=g(B').
\hypertarget{tag-018V}{}\label{cdp:eq:C.WT.9}
\end{equation}
For an algebraic-dimension-one Hopf component the input map is
its reduction in \cref{cdp:st:C.2.UN.37}; the same construction
then applies.

For a nonlinear Hopf component, use its primary contraction
\begin{equation}
 (z,w)\longmapsto(\beta^rz+cw^r,\beta w),\qquad c\ne0.
\hypertarget{tag-017P}{}\label{cdp:eq:C.WT.10}
\end{equation}
Its sole curve is $w=0$. A power of the weighted normal line is
the constant-automorphy axis line. Dividing its connection gives
a flat connection on the weighted line. On $\C^2\setminus0$ the
root section is a constant multiple of $w^b$: its quotient by
$w^b$ extends as a unit at zero, its automorphy factor is one by
evaluation at zero, and iteration of the contraction makes it
constant. Thus \eqref{cdp:eq:C.WT.7} is the map
$(z,v)\mapsto(z,v^q)$. A positive power of the contraction
preserving this branch lifts to
\begin{equation}
 (z,v)\longmapsto(\alpha^{qr}z+c'v^{qr},\alpha v),
                                                    \qquad c'\ne0.
\hypertarget{tag-017Q}{}\label{cdp:eq:C.WT.11}
\end{equation}
Its branch stabilizer is free on $\C^2\setminus0$. A fixed point
would map to a fixed point of the original free deck group; an
element over its identity fixing both $v^q$ and $v^B$ fixes $v$
by $\gcd(q,B)=1$. The finite normalization is therefore the
indicated primary or secondary nonlinear Hopf surface of
resonance $qr$, with the displayed component map.

At a diagonal algebraic-dimension-zero axis, the weighted normal
identity reads $N_1^b\otimes N_2^a\simeq\OO$. Its infinite-order
normal cannot be paired with a finite-order normal: taking a
power that kills $N_2$ would give a finite power killing $N_1$.
Finite coordinate roots preserve nonresonance, since raising a
relation of the lifted multipliers to a common power gives one
for the original multipliers. By \cref{cdp:st:C.2.UN.40} both axes
are in the boundary. This connected branch is the two-axis cycle
of \cref{cdp:st:C.2.UN.59}, which already gives the criterion and
compact realization there. We treat the remaining root components.

Resolve each $uv=s^h$ by the chain in
\eqref{eq:II:compressed-neck-resolution}. The full fibre remains
reduced and the nearby fibres are the same tori. Apply
\cref{cdp:st:C.2.UN.28} to this smooth resolution. It gives
$\sum p_g\le1$ on the strict components, so
\eqref{cdp:eq:C.WT.9} gives precisely the asserted genus bound.
The proper Stein factor of the root family has connected nearby
fibres and one point over zero, so the root graph is connected.
With a genus-one component, \cref{cdp:st:C.2.UN.47} makes the
resolved dual graph simply connected; it is therefore a tree.
Collapsing each inserted chain gives the compressed rooted tree.
With no genus-one component, apply
\cref{lem:II:support-incidence-residues} to this resolution,
using the strict-component restriction squares just constructed
and the identical nonzero restrictions on the inserted ruled
components. It gives the unicyclic graph. In both cases its residue
map gives the stated minimum locus; contraction of a subdivided
edge preserves that locus on the strict components.

The canonical-height calculation on a compressed edge uses the
same curve coordinate $z$ and the relative frame
$du/u\wedge dz=-dv/v\wedge dz$. A component form of curve order
$d_1$ satisfies
\begin{equation}
 s^ku^{d_1+1}\frac{du}{u}\wedge dz
   =-s^{k+h(d_1+1)}v^{-d_1-1}\frac{dv}{v}\wedge dz.
\hypertarget{tag-017R}{}\label{cdp:eq:C.WT.12}
\end{equation}
The vertical-form clause of \cref{lem:II:identity-monodromy-elliptic}
on each strict quasibundle gives
$d_1=m_e\operatorname{ord}\theta+m_e-1$. Hence
\begin{equation}
 \frac{k_w-k_v}{h_em_e}
                           =\operatorname{ord}_{p_e}\theta+1.
\hypertarget{tag-017S}{}\label{cdp:eq:C.WT.13}
\end{equation}
At an unmarked multiple fibre the required order is zero; the
same integer equation forces its multiplicity to be one.
The base-differential degree equation and leaf elimination of
\cref{lem:II:support-heights-period} now apply with edge difference
$(k_w-k_v)/h_e$. At every step the normalized integer in
\eqref{cdp:eq:C.WT.13} is the same integer in that lemma: a child
edge contributes one, a parent edge minus one, and a cycle edge
zero. Thus its outward increment is $h_em_e$, proving
\eqref{cdp:eq:C.WT.4}. The nonlinear leaf form from
\eqref{cdp:eq:C.WT.10} gives this same increment; a product cap
has $m_e=1$.

The connected normal-root collar identifies
$u^{m_e}=c_1$ with $v^{m_e}=s^{h_em_e}/c_1$ as the same oriented
integral two-cycle. Compose it with the full secondary boundary
matrix \eqref{cdp:eq:C.SN.6} at a nonlinear leaf.
The period comparison of \cref{lem:II:support-heights-period}
then uses this same integral class $\gamma_s$. It has order
greater than $H_v$ at an elliptic vertex and exactly $H_v$, with
the nonzero coefficient in \eqref{cdp:eq:C.PT.10}, at a nonlinear
leaf. This proves the common-height and strict inequalities.

Choose a relative canonical generator with constant finite
character by averaging the logarithm of its character unit.
Divide by its minimum power of $s$. The inverse monomials
\eqref{cdp:eq:C.WT.15} and wedge identity
\eqref{cdp:eq:C.WT.16} carry this form to the nonzero canonical
frame on the pruned core. Its character is $\chi=\det_{\C}D$.
The root action preserves the oriented regular-fibre class, so
\begin{equation}
 P(\zeta_Ns)=\chi P(s),\qquad
                            P(s)=\int_{\gamma_s}\Omega_s.
\hypertarget{tag-017T}{}\label{cdp:eq:C.WT.14}
\end{equation}
A nonlinear leaf gives a nonzero leading coefficient of order
$R$, proving \eqref{cdp:eq:C.WT.5}. With no nonlinear leaf choose
a sufficiently large positive exponent in that character class.

Prune by \cref{lem:II:elliptic-tree-pruning} on whole finite orbits.
At each terminal component the one-mark root-line classification
uses its full torsor, primitive basis and secondary integer $n$.
A stabilizer acts on its boundary and connected normal root by
translations, by \cref{cdp:eq:C.WT.7,cdp:eq:C.WT.8}; hence it fixes
the fibre lattice. Transport a chosen basis and $n$ around the
orbit. The maps \eqref{cdp:eq:C.WT.15}, their canonical divisor
sections and their inverses extend its finite action to the
product tubes. The group law agrees on the punctured tubes and
therefore on their extensions. Disjoint steps commute through
unchanged atlases, giving a group-invariant order with children
preceding their parents. The action preserves each original
component, so it cannot exchange the endpoints of an edge between
distinct original components. It can reflect a root cycle through
vertices by interchanging two lifted curves over one original
boundary; parent orientations in a tree are preserved.

For a nonzero $s$, \cref{lem:II:elliptic-tree-pruning} applies its
boundary homomorphism \eqref{cdp:eq:C.PR.9} to $v=s^h/u$.
Its identical pushouts preserve the nearby group throughout.
The final cycle group is $\Lambda\rtimes_M\Z$ with $\Lambda$
injected, so its isomorphism with $\Z^4$ gives $M=I$.
For a genus-one root, \eqref{cdp:eq:C.PT.12} injects the fibre
lattice and identifies the base commutator with the Chern vector;
that vector is zero. The map \eqref{cdp:eq:C.PT.8} identifies this
same principal torsor with its complex two-torus. This proves
the fourth necessary condition.

The finite core action is affine. For a diagonal Hopf pair its
axis-preserving lifts have constant diagonal multipliers, as in
\cref{lem:II:two-mark-elliptic-torsor}; an axis-interchanging lift
is their composition with coordinate exchange. The ruled maps
have the affine logarithmic lifts of
\cref{lem:II:weighted-cycle-affine-data}. Their common edge maps
give one linear part $D$ on the common vector space. For a torus
core the derivative of its lifted map is periodic and holomorphic,
hence constant, giving the same affine conclusion.

Orient a pruned cycle. Let $n_j$ be its outgoing normal meridian
and $h_j$ its neck exponent. Write $\epsilon\in\{1,-1\}$ for the
orientation character and $\sigma$ for the vertex permutation.
The image edge has index $k=\sigma(j)$ for $\epsilon=1$ and
$k=\sigma(j+1)$ for $\epsilon=-1$. The maps are
\begin{equation}
\begin{gathered}
 g(j,v,s)=(\sigma j,Dv+\tau_j,\zeta_Ns),\\
 e_j(v,s)=v+b_j+l_j\log s,\qquad l_j=-h_jn_j/(2\pi i),\\
 Dn_j=\epsilon n_k,\qquad Dl_j=\epsilon l_k.
\end{gathered}
\hypertarget{tag-017W}{}\label{cdp:eq:C.WT.17}
\end{equation}
For a reverse edge, $g_{j+1}e_j=e_k^{-1}g_j$.
Taking $\log\zeta_N=2\pi i/N$, composition and the $N$th-power
relation give
\begin{equation}
\begin{aligned}
 \tau_{j+1}+Db_j-\tau_j-\epsilon b_k
                       -\epsilon l_k\,2\pi i/N&\in\Lambda,\\
 \sum_{a=0}^{N-1}D^{N-1-a}\tau_{\sigma^aj}&\in\Lambda.
\end{aligned}
\hypertarget{tag-017X}{}\label{cdp:eq:C.WT.18}
\end{equation}
Set $b=\sum b_j$ and $u=-\sum h_jn_j=dp$, with $p$ primitive
and $d>0$. The contraction-height calculation
\eqref{cdp:eq:C.AD.3}, applied after subdividing each compressed
edge, gives one nonzero sign for every $\ell(in_j)$, where
$\ker\ell=\operatorname{span}_{\R}\Lambda$.
At ruled vertices this is the elliptic quotient calculation of
\cref{cdp:st:C.2.UN.57}. Therefore $\ell(iu)\ne0$ and
$Du=\epsilon u$. The common elliptic line and $\C p$ give coordinates
\begin{equation}
\begin{gathered}
 W=\operatorname{span}_{\R}\Lambda=\C\times\R,\quad
 p=(0,1),\quad D=\operatorname{diag}(\rho,\epsilon),\\
 \Lambda=\langle(a_1,c_1),(a_2,c_2),p\rangle_{\Z},\quad
 c_1,c_2\in\R,\quad\operatorname{Im}(\overline{a_1}a_2)>0,\\
                     (D-\epsilon I)b-\epsilon u/N\in\Lambda.
\end{gathered}
\hypertarget{tag-017Y}{}\label{cdp:eq:C.WT.19}
\end{equation}
The last relation is the sum of \eqref{cdp:eq:C.WT.18}.
The finite elliptic linear part $\rho$ has order $1,2,3,4$ or $6$
by its integral trace, as in
\cref{lem:II:weighted-cycle-affine-data}. Thus
$\det D=\rho\epsilon$, including a reflection's sign.
Adding $cn_j$ to $b_j$ with the same $c\in\C$ on a whole edge
orbit preserves \eqref{cdp:eq:C.WT.18}. Its orbit sum has nonzero
second coordinate of one sign, so choose $c$ to make the second
coordinate of $b$ equal to $i$. This is the normal-product constant
change along the normal meridian; its restriction to the double
curve is zero, so the elliptic gluing and finite action maps remain
fixed.

For a cycle use these coordinates; for a torus core choose a
finite-invariant Hermitian metric and an orthonormal complex basis.
In both cases $D$ is unitary and
\begin{equation}
 DJ\overline D^{-1}=(\det D)J,\qquad
                       J=\begin{pmatrix}0&1\\-1&0\end{pmatrix}.
\hypertarget{tag-0180}{}\label{cdp:eq:C.WT.20}
\end{equation}
This is the two-dimensional alternating-form identity.
Using \eqref{cdp:eq:C.WT.5}, put
\begin{equation}
 V_s(w)=w+s^RJ\overline w,\qquad DV_s=V_{\zeta_Ns}D.
\hypertarget{tag-0181}{}\label{cdp:eq:C.WT.21}
\end{equation}
Here $V_s$ is the real-linear map of
\cref{lem:II:logarithmic-period-family} evaluated at $s^R$.
Its full rank-four clause gives the compact torus-core family. Transport an affine constant
as $\tau(s)=V_{\zeta_Ns}\tau$ into the target frame. Applying the
real-linear maps to its finite lattice and order equations proves
the same full affine action on this family.

For a cycle use
\begin{equation}
 b_j(s)=V_sb_j,\qquad\tau_j(s)=V_{\zeta_Ns}\tau_j,\qquad
 B_s=V_sb+\frac{d\log s}{2\pi i}V_sp.
\hypertarget{tag-0182}{}\label{cdp:eq:C.WT.22}
\end{equation}
The correction $l_j2\pi i/N=-h_jn_j/N$ is a real rational period
vector, so applying $V_{\zeta_Ns}$ transports
\eqref{cdp:eq:C.WT.18} with its sign into the target frame.
Its summed relation is
\begin{equation}
                         DB_s-\epsilon B_{\zeta_Ns}
                                         \in V_{\zeta_Ns}\Lambda.
\hypertarget{tag-0183}{}\label{cdp:eq:C.WT.23}
\end{equation}
To give the map on the entire period quotient, choose a cut and put
\begin{equation}
\begin{gathered}
 \theta_j(s)=\sum_{a<j}(b_a(s)+l_a(s)\log s),\qquad
 \eta=v_j-\theta_j(s),\qquad
 l_a(s)=-h_aV_s(n_a)/(2\pi i),\\
 \eta\longmapsto D\eta+D\theta_j(s)+\tau_j(s)
                                  -\theta_{\sigma j}(\zeta_Ns).
\end{gathered}
\hypertarget{tag-0184}{}\label{cdp:eq:C.WT.23a}
\end{equation}
Adjacent charts differ by the lattice element in
\eqref{cdp:eq:C.WT.18}. Changing the cut adds $B_s$ on the source
and $\epsilon B_{\zeta_Ns}$ on the target, which is absorbed by
\eqref{cdp:eq:C.WT.23}; changing a logarithm adds integral normal
periods. The second equation of \eqref{cdp:eq:C.WT.18} gives the
$N$th-order relation. These are holomorphic maps of the full family,
including reflections. The normal charts are $uv=s^{h_j}$,
with the coordinates exchanged on a reversed edge. Their moving
Hopf and ruled frames are those of
\cref{lem:II:moving-marked-cycle,lem:II:weighted-cycle-roots}.

Apply the power-$R$ clause of \cref{lem:II:logarithmic-period-family}
to the same $\Lambda,b,d$. Its determinant
\eqref{cdp:eq:C.WT.24} is positive, giving compact connected tori
on every punctured fibre. The preceding compressed charts provide
a proper end filling. At infinity put $z=s^{-1}$ and divide the
vector frame by $s^R$. The transformed data are
\begin{equation}
 W_z(w)=J\overline w+z^Rw,\qquad
              D_\infty=(\det D)^{-1}D=J\overline D J^{-1}.
\hypertarget{tag-0186}{}\label{cdp:eq:C.WT.25}
\end{equation}
The base character is $\zeta_N^{-1}$, and an affine constant is
$J\overline\tau+(\det D)^{-1}z^R\tau
=W_{\zeta_N^{-1}z}\tau$.
The outgoing meridians are $-J\overline{n_j}$ at the endpoint.
The change $\log s=-\log z$ gives the conjugate core with the
same positive neck exponents. It conjugates each affine constant
and reflection, giving the second proper end. A finite closed
cover of the parameter sphere by these end discs and a middle
annulus proves compactness of the core family.

For each insertion orbit, \cref{lem:II:elliptic-core-period-charts}
uses the same primitive regular-fibre lattice. Its determinant
\eqref{cdp:eq:C.WT.26} and exponential quotient
\eqref{cdp:eq:C.WT.27} give initial displacement order $R$.
Its explicit unimodular period-basis map gives each chosen
primitive basis and changes displacement by a unit. At a vertex
of height $H_v$, the remaining order is $\ell=R-H_v$.
For an edge with data $(h_e,m_e,n)$, apply
\cref{lem:II:finite-elliptic-insertion-gauge} to that curve's
stabilizer. The affine core maps above give the lift
\eqref{cdp:eq:C.WT.28}; its fibre-lattice action is the identity by
\eqref{cdp:eq:C.WT.8}, and its lifted $d_0$th-power relation is the
one transported in equivariant pruning. The lemma uses the
stabilizer's base character and all its angular and infinite
integers. It preserves the resonant nonlinear translation in
\eqref{cdp:eq:C.WT.32} and gives the holomorphic maps
\cref{cdp:eq:C.WT.37,eq:II:weighted-child-finite-fibre} on the child.

Now apply \cref{lem:II:weighted-elliptic-insertion} with
$(m,h)=(m_e,h_e)$. The height inequalities give $\ell\ge hm$ for
a Hopf step and $\ell>h$ for a product cap. Its chart
\eqref{cdp:eq:C.WT.33} and pulled canonical section $x^n$ have
coefficient $nh$ on the generic new divisor. Its free finite chart
and infinite map are \cref{cdp:eq:C.WT.34,cdp:eq:C.WT.35}; the
new displacement \eqref{cdp:eq:C.WT.36} has order $\ell-hm$.
Equality gives exactly the nonlinear terminal pair with the
matched finite translation; strict inequality gives the diagonal
pair. A product cap uses \eqref{cdp:eq:C.WT.38} and order drop $h$.
Thus induction subtracts exactly $h_em_e$ at each edge and reaches
order zero exactly at a nonlinear leaf.

The inverse chart \eqref{cdp:eq:C.WT.15} recovers each parent
transition, while the finite gauges carry every remaining mark.
Reversing the complete equivariant pruning sequence therefore
restores all root components, their torsors, and both linear and
translation parts of every gluing. The norm
\eqref{cdp:eq:C.WT.39}, its bound \eqref{cdp:eq:C.WT.40} and cluster
map \eqref{cdp:eq:C.WT.41} in the insertion lemma prove Hausdorff
proper gluing at each stage, with its same exponent $h_e$.
Taking the finite orbit of a norm makes these bounds uniform on
the orbit. There are finitely many disjoint equivariant replacement
stages, and their compressed necks remain $uv=s^{h_e}$.

At a generic point of a strict component over $S_i$, the inertia
has order $N/m_i$, fixes the component and acts by a primitive
character on its normal parameter. Its quotient is smooth and
its full fibre multiplicity is $m_i$. At an original edge
$a=gA,b=gB$, the point inertia has order $hAB=N/g$.
Its subgroup of order $h$ fixes $u,v$ and acts primitively on $s$,
so the quotient of $uv=s^h$ is the smooth $(u,v)$-chart.
The residual group has independent characters of orders $B,A$;
coprimality gives
\begin{equation}
 \bigl(\C\{z,u,v,s\}/(uv-s^h)\bigr)^I
     =\C\{z,u^B,v^A\},\qquad t=(u^B)^a(v^A)^b.
\hypertarget{tag-018P}{}\label{cdp:eq:C.WT.42}
\end{equation}
The curve coordinate $z$ is invariant because its finite map is
\'{e}tale. Finite character averaging gives coordinates
$uU,vV,z+\cdots$ with the same characters; their product is $s^h$
times an invariant unit. Dividing one coordinate by that unit
preserves the characters and gives the same equation. Thus
\eqref{cdp:eq:C.WT.42} is the invariant ring of the full action.
The endpoints lie over distinct original components, so there
is no edge inversion. Away from the double curves the strict
component maps are \'{e}tale after generic normal inertia.
These charts cover the quotient and prove its smoothness.

The weighted lines and sections are recovered as well. On a root
complement $v_i=s^{N/m_i}/\xi$. The character eigensheaves extend
through \eqref{cdp:eq:C.WT.7} as
\begin{equation}
 \mathcal L_i^{\otimes k}\otimes\OO_{S_i}\!\left(
        \sum_j\left\lfloor\frac{k m_j}{m_i}\right\rfloor J_{ij}
                                      \right),\qquad 0\le k<m_i.
\hypertarget{tag-018Q}{}\label{cdp:eq:C.WT.43}
\end{equation}
To identify the tensor power, choose frames
$e_\beta=g_{\alpha\beta}e_\alpha$ of $\mathcal L_i$ and write
$\xi=\xi_\alpha e_\alpha^{-1}$. Then
$\xi_\beta=g_{\alpha\beta}\xi_\alpha$, so the character-$k$
monomial has transition $g_{\alpha\beta}^k$.
Along $J_{ij}$ put $d_{ij}=\lfloor k m_j/m_i\rfloor$.
In \eqref{cdp:eq:C.WT.7}, the function
$\xi_\alpha^k/z_\alpha^{d_{ij}}$ has nonnegative normal order
$kB-q d_{ij}<q$. For $z_\beta=h_{\alpha\beta}z_\alpha$ its
transition is $g_{\alpha\beta}^k h_{\alpha\beta}^{-d_{ij}}$,
which is exactly the frame transition of
$\mathcal L_i^{\otimes k}\otimes\OO_{S_i}(d_{ij}J_{ij})$.
The boundary curves are disjoint, so these factors give the sum
in \eqref{cdp:eq:C.WT.43}. Multiplication, the full section and
the edge identifications in \eqref{cdp:eq:C.WT.1} recover the
original weighted line data and all their integer orders.

At infinity use the conjugate inverse-pruning sequence and
\eqref{cdp:eq:C.WT.25}. Its base character is $\zeta_N^{-1}$,
its canonical character is $(\det D)^{-1}$ and its primitive
elliptic period has order $R$. Hence the same finite-jet and height
maps are the conjugates of those at zero. All compressed charts
and invariant rings \eqref{cdp:eq:C.WT.42} are conjugate too,
giving the conjugate weighted support and a smooth quotient at
that end. The two replacement regions are disjoint in the base.
The compact normal root threefold has a compact finite quotient,
which is smooth by these charts. Off the two ends the base action
is free, and one torus fibre over an orbit maps biholomorphically
to the quotient fibre. All those fibres are connected two-tori.
Every original quasibundle or Hopf component and every double
elliptic curve has Euler number zero. Additivity proves the Euler
assertion for both reduced singular supports and completes the
criterion and compact construction.
\end{proof}






\part{Canonical reduction and the Euler sign}\label{part:II:canonical-reduction}
The divisor of a canonical evaluation form assigns an order to each
component. Boundary intersection equations and normal-line morphisms
control contractions, while cyclic roots compare these orders with full
fibre multiplicities. The resulting reductions determine the Euler
equality cases and the finite generation statements.

\section{Terminal contractions and canonical-flat fibres}\label{sec:II:canonical-flat-classification}

\subsection{Canonical forms under bimeromorphic maps}\label{subsec:II:canonical-bimeromorphic-maps}

\begin{lemma}[Canonical direct images and surviving divisor orders]\hypertarget{tag-P343}{}\label{lem:II:canonical-bimeromorphic-direct-image}
\StatementRecord{lem:II:canonical-bimeromorphic-direct-image}{proved}{P343}{\cref{cdp:st:C.2.UN.17}}
Let $f:X\to\Delta$ and $f':X'\to\Delta$ be proper degenerations
of connected complex two-tori from smooth complex threefolds.
Let $\phi:X\dashrightarrow X'$ be a bimeromorphic map over the disc.
Meromorphic pullback induces an isomorphism of canonical direct images
\begin{equation}
                f'_*K_{X'}\xrightarrow{\ \phi^*\ } f_*K_X.
\hypertarget{tag-00PE}{}\label{cdp:eq:C.KG.6}
\end{equation}
The inverse is the pullback by $\phi^{-1}$. It identifies local
generators up to a holomorphic base unit. For a prime divisor which
survives as a prime divisor under $\phi$, the orders of corresponding
generators along the two divisors are equal.
\end{lemma}
\begin{proof}
A holomorphic top form on $f'^{-1}(A)$, for an open subdisc $A$,
pulls back holomorphically on the domain of $\phi$ in $f^{-1}(A)$.
The indeterminacy set of a meromorphic map from a smooth space
has codimension at least two. In smooth coordinate volume frames,
Hartogs extension of its coefficient therefore gives a holomorphic
top form on all of $f^{-1}(A)$. These extensions agree on chart
overlaps by uniqueness, and commute with restriction in $A$.
Apply the same construction to $\phi^{-1}$. Both composites are
the identity on the common biholomorphic dense open set, hence
on the whole smooth spaces. This proves the displayed sheaf
isomorphism. By \cref{cdp:st:C.2.UN.17}, the canonical direct images
are rank-one torsion-free coherent sheaves on the smooth disc
and thus line bundles, so two
local generators differ by a base unit. At the generic points of
corresponding surviving divisors, the bimeromorphic map identifies
their divisorial local rings and is a biholomorphism on dense open
subsets of those divisors. Its Jacobian is a unit there. Pullback
therefore preserves the vanishing order of the top form.
\end{proof}

\subsection{Section and ruling-fibre contractions}\label{subsec:II:section-fibre-contractions}

\begin{proposition}[Contraction of a section and ruling-fibre boundary]\hypertarget{tag-P335}{}\label{prop:II:section-fibre-contraction}
\hypertarget{tag-01EL}{}\label{cdp:st:C.2.UN.74}
\StatementRecord{prop:II:section-fibre-contraction}{proved}{P335}{\cref{cdp:st:C.2.UN.17,lem:II:blowup-topology,lem:II:canonical-bimeromorphic-direct-image}; \cite[Main Theorem, pp.~483--484; \S3, pp.~499--501]{Nakano1970}; \cite[\S2, Theorem~2\(^{\prime}\) and proof, pp.~640--642; Proposition~9\(^{\prime}\) and proof, p.~642]{FujikiNakano1971}}
\StatementRecord{cdp:st:C.2.UN.74}{proved}{01EL}{}

Let $f:U\to\Delta$ be a proper degeneration of connected
complex two-tori from a smooth threefold, with SNC reduced
central support and arbitrary positive full multiplicities.
Let $T$ be a component with a specified holomorphic
projective-line bundle $p:T\to B$ over a smooth compact
curve. Suppose its boundary in the reduced support is
exactly
\begin{equation}
           C+F_1+\cdots+F_r,\qquad
           C\ \text{a section},\quad F_j=p^{-1}(b_j),
\hypertarget{tag-00XN}{}\label{cdp:eq:C.RC.1}
\end{equation}
where the $b_j$ are distinct. The section is the
intersection with one component $P$, and $F_j$ is the
intersection with a component $Q_j$. The components
$Q_j$ may repeat at different $b_j$, but none is $P$.
All these are reduced transverse intersection
divisors; their intersections $C\cap F_j$ are the
specified triple points. Let the full multiplicities
of $T,P,Q_j$ be respectively $m_T,m_P,n_j$.

There is an ordinary smooth curve contraction along
the ruling, inverse to an ordinary blowup and equal
to an isomorphism off $T$, exactly when
$m_T=m_P$. In particular this always holds for a
full reduced fibre. Set their common value equal
to $m$. The resulting map is
\begin{equation}
              \pi:U\longrightarrow U',\qquad
              f=f'\pi,\qquad \pi(T)=B\subset P'.
\hypertarget{tag-00XR}{}\label{cdp:eq:C.RC.2}
\end{equation}
The threefold $U'$ is smooth and $f'$ is proper with
the same punctured torus family. Its reduced central
support is SNC. The map $P\to P'$ is an isomorphism.
Each $Q_j\to Q'_j$ is the ordinary point blowup at
$b_j\in B$, with all indicated points blown up when
the same component repeats. The curve $B$ lies in
the smooth surface $P'$, avoids every triple point
of the descended support, and meets its double
curves transversely, exactly at the $b_j$.

All full multiplicities are retained: the parent
has multiplicity $m$ and $Q'_j$ has multiplicity
$n_j$. At $b_j$ there are coordinates
$(p_0,q_0,r_0)$ with
\begin{equation}
 \begin{gathered}
 P'=(p_0=0),\quad Q'_j=(q_0=0),\quad
 B=(p_0=r_0=0),\qquad f'=p_0^m q_0^{\,n_j},\\
 p_0=xy,\quad r_0=x,\quad q_0=z,\qquad
                         f=x^m y^m z^{\,n_j}
 \end{gathered}
\hypertarget{tag-00XS}{}\label{cdp:eq:C.RC.3}
\end{equation}
in the blowup chart containing $C\cap F_j$.
Every other component and every neighbourhood away
from the indicated exceptional configurations are
unchanged.

Write
$L=N_{B/P'}$ and $M=N_{P'/U'}|_B$.
The normal extension and full-fibre
trivialization are
\begin{equation}
 \begin{gathered}
 0\longrightarrow L\longrightarrow
       V=N_{B/U'}\longrightarrow M\longrightarrow0,
 \qquad (T,C)=(\PP_B(V),\PP_B(L)),\\
 N_{T/U}\simeq\OO_{\PP_B(V)}(-1),\qquad
 M^{\otimes m}\simeq
             \OO_B\left(-\sum_{j=1}^r n_jb_j\right).
 \end{gathered}
\hypertarget{tag-00XT}{}\label{cdp:eq:C.RC.4}
\end{equation}
The full fibre weights on the ruling-fibre
boundary determine the displayed power of the normal line $M$. In particular $m$ divides
$\sum_jn_j$. For canonical orders $k_P,k_{Q_j}$
on the descended family, the section subline is
\begin{equation}
 L\simeq K_B\otimes M^{-(k_P+1)}
                    \otimes\OO_B\left(-\sum_jk_{Q_j}b_j\right),
             \qquad k_T-k_P=1.
\hypertarget{tag-00XU}{}\label{cdp:eq:C.RC.5}
\end{equation}
The entire extension in \eqref{cdp:eq:C.RC.4}, with its endpoint
identifications, is retained.

For a full reduced fibre, put $Z=\sum_jb_j$. Then
these maps specialize to
\begin{equation}
 \begin{gathered}
 M\simeq\OO_B(-Z),\qquad
 L\simeq K_B\otimes
              \OO_B\left(\sum_j(k_P-k_{Q_j}+1)b_j\right),\\
 df'|_V:\ V\twoheadrightarrow
            \OO_B(-Z)\hookrightarrow\OO_B.
 \end{gathered}
\hypertarget{tag-00XV}{}\label{cdp:eq:C.RC.6}
\end{equation}
If in addition $T$ is a local maximum of the
canonical order and $r>0$, the extension splits:
$V\simeq L\oplus M$, with its specified subline.

Put $S=(f^{-1}(0))_{\mathrm{red}}$ and
$S'=(f'^{-1}(0))_{\mathrm{red}}$. Their invariants satisfy
\begin{equation}
              e(S)=e(S')+2-2g(B),\qquad
              t(S)=t(S')+r,
\hypertarget{tag-00XW}{}\label{cdp:eq:C.RC.7}
\end{equation}
where $t$ counts triple points. The same contraction
exists on a compact total space. Conversely an
ordinary blowup of a smooth curve in a single
central component, transverse to its double curves
and avoiding the triple locus, has exactly the
maps, markings and multiplicities displayed above.
\end{proposition}

\begin{proof}

\begin{enumerate}[label=\arabic*.,leftmargin=*]
\item \par\medskip\noindent\textbf{The normal degree and the corrected contraction theorem.}
Restrict the principal full central fibre to $T$.
For a ruling fibre $\ell$ away from the marked
values, the resulting line isomorphism gives
\begin{equation}
       m_T\deg(N_{T/U}|_\ell)=-m_P,\qquad
       N_{T/U}^{\,m_T}
           \simeq\OO_T\left(-m_PC-\sum_jn_jF_j\right).
\hypertarget{tag-00XX}{}\label{cdp:eq:C.RC.8}
\end{equation}
If $m_T=m_P$, the restriction is $\OO_{\PP^1}(-1)$.
Apply the analytic contraction theorem of
\cite[Main Theorem, pp.~483--484]{Nakano1970}, using the
normal-line criterion established in
\cite[\S2, Theorem~2\(^{\prime}\) and proof, pp.~640--642;
Proposition~9\(^{\prime}\) and proof, p.~642]{FujikiNakano1971}.
In the source's notation, $(\widetilde X,S,M,r)=(U,T,B,2)$:
$T$ is the specified holomorphic projective-line bundle and
$\OO_U(T)|_{p^{-1}(b)}=\OO_{\PP^1}(-1)$ for every $b$.
Its local morphism restricts to this same $p$, and identifies
the neighbourhood with the ordinary blowup of a smooth
codimension-two centre. The local morphisms glue to the identity
off $T$ by \cite[\S3.2, pp.~500--501]{Nakano1970}.
They give the proper blowdown $\pi:U\to U'$ onto a smooth
threefold, with $U=\operatorname{Bl}_B U'$.
Conversely an ordinary such blowup has normal degree minus
one, so \eqref{cdp:eq:C.RC.8} forces $m_T=m_P$.

\item The equality $\pi_*\OO_U=\OO_{U'}$ gives the
unique holomorphic $f'$ in \eqref{cdp:eq:C.RC.2}. Its
properness follows because the inverse image
of a compact set in the disc is the image
under $\pi$ of the corresponding compact
set for $f$. The smooth fibres are unchanged.
The images of the boundary components are determined by
their first normal terms along the centre, as follows.

\par\medskip\noindent\textbf{The parent image is smooth and its whole surface is retained.}
The proper image $P'=\pi(P)$ is a hypersurface
in the smooth target. Its Cartier
pullback has the form
\begin{equation}
                            \pi^*P'=P+aT.
\hypertarget{tag-00XY}{}\label{cdp:eq:C.RC.9}
\end{equation}
Restriction to an ordinary ruling fibre
gives $0=1-a$, since $P\cap T=C$ is a
section and $T|_\ell=\OO_\ell(-1)$.
Thus $a=1$.

\item A local defining section of $P'$ vanishes
to first order along the generic centre.
Its first normal term on $\PP_B(N_{B/U'})$
has zero divisor exactly $P\cap T=C$,
by \eqref{cdp:eq:C.RC.9}. This term cannot vanish at
any whole fibre: $C$ is a section and
contains none. Its differential in the
normal directions is therefore nonzero
at every point of $B$. Hence $P'$ is
smooth along $B$. Since the centre is
a Cartier curve in $P'$, the restriction of its blowup is an
isomorphism $P\to P'$. In particular
it identifies the section $C$ with
the original curve $B$.

\par\medskip\noindent\textbf{Each ruling-fibre boundary becomes a transverse double curve.}
Consider one of the other image components
$Q'=\pi(Q)$. It does not contain $B$:
otherwise its generic first normal
divisor would meet every ruling fibre,
whereas $Q\cap T$ is a finite union
of fibres. More explicitly, writing
$\pi^*Q'=Q+aT$ and restricting to a
generic ruling fibre gives $a=0$.
Thus $\pi^*Q'=Q$.

\item Restrict a local defining section of
$Q'$ to $B$. Its pullback to $T$
vanishes exactly on the specified
reduced fibres belonging to $Q$.
At $b_j$ that section therefore
has a simple zero on $B$. Its
differential is nonzero on the
tangent line of $B$. Consequently
$Q'$ is smooth at $b_j$ and is
transverse to $B$. The independent
normal differential of $P'$ makes
$P'$ and $Q'$ transverse as well.

\item There is only one such other
component at each $b_j$. A second
would have the same fibre in
$T\cap Q$, contradicting the
given SNC intersections and
\eqref{cdp:eq:C.RC.1}. Any further image
component passing through $b_j$
would likewise have a boundary
component on the exceptional
fibre; hence none is missing
from this count. Thus there is
no triple point of the descended
support on $B$.

\item Choose coordinates with $P'$ and
$Q'$ the first two coordinate
planes and $B$ the stated
coordinate curve. Transversality
just proved permits this choice.
The full fibre is a unit times
$p_0^m q_0^{n_j}$; a local root
of that unit absorbs it into
$p_0$ without changing these
marked loci. Blow up $(p_0,r_0)$.
The exact chart $p_0=xy,r_0=x$
gives \eqref{cdp:eq:C.RC.3}. On $Q'$ this
ideal is the maximal ideal
of $b_j$, so $Q$ is precisely
its ordinary point blowup,
with exceptional fibre $F_j$.
On $P'$ the ideal is Cartier,
giving the already established
isomorphism. These are the
claimed original neighbourhood
maps.

\par\medskip\noindent\textbf{The full normal extension and the weighted line.}
The normal sequence of
$B\subset P'\subset U'$
is the first row of \eqref{cdp:eq:C.RC.4}.
The exceptional divisor of
the ordinary blowup is $\PP_B(V)$, where $\PP_B(V)$
parametrizes lines in $V$, with
ambient normal $\OO(-1)$.
The strict parent intersects
it in the specified subline
$\PP_B(L)$. In particular,
with these line conventions,
\begin{equation}
                 \OO_T(C)\simeq
                      \OO_{\PP_B(V)}(1)\otimes p^*M.
\hypertarget{tag-00XO}{}\label{cdp:eq:C.RC.10}
\end{equation}
The principal full fibre on
$U'$ is $mP'+\sum_Qm_QQ'$.
Restriction to $B$ gives
exactly the last isomorphism
of \eqref{cdp:eq:C.RC.4}, since each
$Q'$ meets $B$ simply at
its listed $b_j$. This
also follows by substituting
\eqref{cdp:eq:C.RC.10} into \eqref{cdp:eq:C.RC.8}.
Taking degrees proves the
stated divisibility. The
line $M$ and its power trivialization are
retained, including any
torsion information.

\item Apply \cref{cdp:st:C.2.UN.17} to the descended proper
torus family and \cref{lem:II:canonical-bimeromorphic-direct-image}
to $(X,X',f,f',\phi)=(U,U',f,f',\pi)$.
The pullback in \eqref{cdp:eq:C.KG.6} is ordinary pullback by
this blowdown; its inverse uses exactly the isomorphism
$U\setminus T\simeq U'\setminus B$ and extension across $B$.
Thus a generator downstairs pulls back to a generator upstairs.
Restrict its vertical canonical line to $B$; the resulting isomorphism is
\begin{equation}
 K_{U'}|_B\simeq
       M^{k_P}\otimes
             \OO_B\left(\sum_jk_{Q_j}b_j\right).
\hypertarget{tag-00XP}{}\label{cdp:eq:C.RC.11}
\end{equation}
Other components have a
nonvanishing defining section
on $B$ and contribute their
specified trivializations.
Adjunction for $B\subset P'$
then proves the formula for
$L$ in \eqref{cdp:eq:C.RC.5}. The
ordinary curve-blowup
Jacobian contributes $T$
once to $K_U$; the generic
centre lies only in $P'$.
Thus $k_T=k_P+1$.
Changing the canonical
generator by a power of $f'$
adds the full multiplicities
to the orders, and \eqref{cdp:eq:C.RC.4}
shows that \eqref{cdp:eq:C.RC.5} remains
the same line isomorphism.

\item In the full reduced case
$m=n_j=1$, so $M=\OO_B(-Z)$.
The derivative of
$f'=p_0q_0$ along the
normal bundle of $B$
is the quotient normal
map to $M$, followed by
its inclusion in $\OO_B$
with simple zero at $b_j$.
Away from the marked
points it is surjective.
This proves all maps in
\eqref{cdp:eq:C.RC.6}. For $m>1$ the
ordinary derivative of
$f'$ is zero along the
generic centre; the
quotient in \eqref{cdp:eq:C.RC.4}
instead is the normal map of $P'$,
and has already been
specified independently
of that derivative.

\item For the splitting assertion,
the extension group is the
space $H^1(B,L\otimes
M^{-1})$. Equations \eqref{cdp:eq:C.RC.6}
and Serre duality identify
its dual with
\begin{equation}
 H^0\left(B,\OO_B\left(
                   \sum_j(k_{Q_j}-k_P-2)b_j\right)\right).
\hypertarget{tag-00XQ}{}\label{cdp:eq:C.RC.12}
\end{equation}
At a local maximum
$k_{Q_j}\le k_T=k_P+1$.
Every coefficient in
\eqref{cdp:eq:C.RC.12} is at most
minus one. When $r>0$
the line has negative
degree and no section.
The extension group is
zero, which proves the
splitting with
the marked subline.
When $r=0$ in the full reduced case, the same formulas give
$M=\OO_B$, $L=K_B$ and the endpoint-framed sequence
$0\to K_B\to V\xrightarrow{df'}\OO_B\to0$;
its extension class is the element of $H^1(B,K_B)$
defined by these normal maps.

\par\medskip\noindent\textbf{Euler, triples, the converse and compact descent.}
Apply \cref{lem:II:blowup-topology} to the same
ordinary blowup $\pi$, the smooth compact centre $B$ and
the compact reduced support $S'\supset B$. Its full reduced
inverse image is $S$, with exceptional divisor $T=\PP_B(V)$.
The complement isomorphism $S\setminus T\simeq S'\setminus B$
therefore gives $e(S)-e(S')=e(B)=2-2g(B)$ through that
owner's blowup comparison, for all the stated fibre weights.
At each $b_j$ the chart
\eqref{cdp:eq:C.RC.3} has exactly one
triple point upstairs
and only a smooth double
curve downstairs. There
are no other changes,
which proves \eqref{cdp:eq:C.RC.7}.

\item Conversely, blowing up the
indicated transverse curve
uses exactly these charts.
The exceptional multiplicity
is the generic parent
multiplicity $m$, and
every other weight is
unchanged. The section
and ruling fibres are
exactly the intersections
in \eqref{cdp:eq:C.RC.1}, so all
the displayed maps and
markings are recovered.
The local blowdowns glue to the identity off $T$ by the
same source overlap maps already used above. Thus, on a
compact total space, they define the stated contraction
over the original compact base.

\end{enumerate}
\end{proof}

\subsection{Anticanonical elliptic bisections}\label{subsec:II:anticanonical-bisections}

\begin{lemma}[The residue obstruction to an elliptic anticanonical bisection]\hypertarget{tag-P338}{}\label{lem:II:anticanonical-bisection}
\StatementRecord{lem:II:anticanonical-bisection}{proved}{P338}{}
Let $q:V\to B$ be a holomorphic projective-line bundle over a smooth
compact elliptic curve. Let $C\subset V$ be a smooth connected
elliptic curve such that $h=q|_C:C\to B$ has degree two.
Then $K_V(C)$ is nontrivial.
\end{lemma}
\begin{proof}
The fibre cohomology of a local projective-line bundle chart gives
$q_*\OO_V=\OO_B$ and $R^1q_*\OO_V=0$. Thus Leray identifies
$q^*:H^1(B,\OO_B)\to H^1(V,\OO_V)$ with an isomorphism.
For the inclusion $i:C\hookrightarrow V$, the restriction map is
\begin{equation}
 H^1(V,\OO_V)\xrightarrow{\ i^*\ }H^1(C,\OO_C),
 \qquad i^*q^*=h^*.
\hypertarget{tag-00G6}{}\label{cdp:eq:C.EL.7}
\end{equation}
Riemann--Hurwitz makes $h$ unramified. The finite trace satisfies
$\operatorname{Tr}_h\circ h^*=2\operatorname{id}$, so this
restriction map is nonzero.

Suppose $K_V(C)\simeq\OO_V$. Its divisor residue sequence is
\begin{equation}
\hypertarget{tag-00DO}{}\hypertarget{tag-00QY}{}\SsixEquationAliases{\label{cdp:eq:C.CZ.10}\label{cdp:eq:C.ML.7}}
 0\longrightarrow K_V\longrightarrow K_V(C)
       \xrightarrow{\operatorname{res}}K_C\longrightarrow0.
\hypertarget{tag-00G7}{}\label{cdp:eq:C.EL.8}
\end{equation}
The restriction of $K_V$ to a ruling fibre has degree $-2$, so
$H^0(V,K_V)=0$. The assumed trivialization gives
$H^0(V,K_V(C))=\C$, and $H^0(C,K_C)=\C$ because $C$ is elliptic.
The residue is nonzero: locally a frame of $K_V(C)$ is
$dv/v\wedge du$ with $C=(v=0)$, and its residue is the
nonzero form $du$ along $C$. Hence the residue map on these
one-dimensional spaces is an isomorphism and its connecting
map $H^0(C,K_C)\to H^1(V,K_V)$ is zero. Serre duality for
this residue sequence identifies that connecting map with
the dual of the restriction in \eqref{cdp:eq:C.EL.7}.
This contradicts the nonzero finite pullback proved above.
\end{proof}

\subsection{Terminal ruled contractions}\label{subsec:II:terminal-ruled-contractions}

\begin{proposition}[Contractions and normal extensions of terminal ruled pairs]\hypertarget{tag-01EA}{}\label{cdp:st:C.2.UN.64}
\StatementRecord{cdp:st:C.2.UN.64}{proved}{01EA}{\cref{cdp:st:C.2.UN.17,prop:II:section-fibre-contraction,lem:II:anticanonical-bisection,lem:II:blowup-topology}; \cite[Main Theorem, pp.~483--484]{Nakano1970}; \cite[\S2, Theorem~2\(^{\prime}\), pp.~640--642, and Proposition~9\(^{\prime}\), p.~642]{FujikiNakano1971}; \eqref{cdp:eq:C.PT.13}, \eqref{cdp:eq:C.PR.13}}

Let $f:U\to\Delta$ be a proper degeneration of connected complex
two-tori from a smooth threefold. Assume that the full central fibre
$S$ is reduced and SNC. Let $T$ be one component, meeting the rest
of $S$ in exactly one smooth connected curve $C$, contained in a
single other component $P$. Let $k_T,k_P$ be the orders of the
canonical evaluation along these components. Suppose that either
$q:T\to B$ is a holomorphic $\PP^1$-bundle over a smooth
compact curve, or $T\simeq\PP^2$.

For a ruled $T$, precisely one of the following two numerical and
marked-pair descriptions holds. Either $C$ is a section,
\begin{equation}
 K_T=-2C,\qquad k_T-k_P=1,\qquad
 N_{C/T}\simeq K_B^{-1},\qquad N_{C/P}\simeq K_B,
                                                        \hypertarget{tag-00FY}{}\label{cdp:eq:C.EL.1}
\end{equation}
where the section identifies $C$ with $B$; or $B\simeq\PP^1$,
$C$ is a smooth anticanonical elliptic bisection and
\begin{equation}
 K_T=-C,\qquad k_T=k_P,\qquad
 T\simeq\mathbb F_e\quad\hbox{with }e\in\{0,1,2\}.
                                                        \hypertarget{tag-00G1}{}\label{cdp:eq:C.EL.2}
\end{equation}
For $T\simeq\PP^2$, either $C$ is a line and $k_T-k_P=2$,
or $C$ is a smooth cubic and $k_T=k_P$.

In the section case \eqref{cdp:eq:C.EL.1} there is a proper holomorphic
map $\pi:U\to U'$ onto a smooth threefold, equal to an isomorphism
off $T$, which contracts the ruling to a smooth curve
$B\subset U'$. It identifies $U$ with the ordinary blowup of $U'$
along $B$. The map $f$ descends to a proper holomorphic torus
degeneration $f':U'\to\Delta$, smooth along $B$; the restriction
$P\to P'$ is an isomorphism and identifies $C$ with $B$.
The full new fibre $S'$ is again reduced and SNC. Its normal
bundle and Euler comparisons are
\begin{equation}
\begin{gathered}
 0\longrightarrow K_B\longrightarrow V:=N_{B/U'}
        \xrightarrow{\,df'\,}\OO_B\longrightarrow0,\qquad
 (T,C)\simeq\bigl(\PP_B(V),\PP_B(K_B)\bigr),\\
 N_{T/U}\simeq\OO_{\PP_B(V)}(-1)\simeq\OO_T(-C),\qquad
 e(S)=e(S')+2-2g(B).
\end{gathered}                                                   \hypertarget{tag-00G2}{}\label{cdp:eq:C.EL.3}
\end{equation}
Here $\PP_B(V)$ parametrizes lines. The extension class belongs to the one-dimensional space $H^1(B,K_B)$; the maps in the displayed exact sequence specify its endpoint identifications.

In the plane-line case the corresponding map contracts $T$ to a
smooth point of a smooth reduced central component, identifies
$U$ with its point blowup, and gives
\begin{equation}
       P=\operatorname{Bl}_{p}P',\qquad e(S)=e(S')+2.
                                                        \hypertarget{tag-00G3}{}\label{cdp:eq:C.EL.4}
\end{equation}
These contractions also exist on a compact total threefold over a
compact curve. Iteration on eligible leaves terminates and leaves
the punctured torus family unchanged, with exact Euler comparison
\begin{equation}
 e(S)=e(S_{\rm rem})+2a+\sum_{j=1}^{b}\bigl(2-2g(B_j)\bigr),
                                                        \hypertarget{tag-00G4}{}\label{cdp:eq:C.EL.5}
\end{equation}
where $a$ is the number of point contractions and the
$B_j$ are the curve centres. Reversing those ordinary blowups restores the original marked support and all
its neighbourhood maps.

For every elliptic curve $B$, both the split and the
nonzero-extension pairs in \eqref{cdp:eq:C.EL.3} occur in a compact torus
degeneration over $\PP^1$ with exactly one singular fibre. That
fibre has two components: a two-torus and the prescribed ruled
pair. Its Euler number is zero, and every other fibre is a smooth
two-torus.
\end{proposition}


\begin{proof}

\begin{enumerate}[label=\arabic*.,leftmargin=*]
\item Apply \cref{cdp:st:C.2.UN.17} to this full fibre, with every multiplicity $m_v=1$. For the smooth component $T$, its normalization and resolution maps are the identity and its intrinsic conductor divisor is zero. The intersection divisors with the other components sum to the single curve $C$ in $P$. The canonical evaluation generator of that proposition has divisor $\sum_v k_vT_v$. Its principal-fibre and adjunction maps, \eqref{cdp:eq:C.CA.3} and \eqref{cdp:eq:C.CA.4}, specialize to
\begin{equation}
 \OO_U(T)|_T\simeq\OO_T(-C),\qquad
 K_T\simeq\OO_T\bigl((k_P-k_T-1)C\bigr).
                                                        \hypertarget{tag-00G5}{}\label{cdp:eq:C.EL.6}
\end{equation}
The adjunction map divides the evaluation section by a local equation of $T$ to order $k_T$, so it gives a meromorphic canonical section with the displayed divisor. Multiplication of the evaluation generator by a base power shifts every $k_v$ by the same integer and leaves this formula unchanged. At $C$, the reduced equation $xy=t$ trivializes $N_{C/T}\otimes N_{C/P}$.

\item For a ruling fibre $L\simeq\PP^1$, put $d=C\cdot L$ and
$a=k_P-k_T-1$. Intersecting \eqref{cdp:eq:C.EL.6} with $L$ gives
$-2=ad$. Thus $d>0$, $a<0$, and the integer pairs are exactly
$(d,a)=(1,-2)$ and $(2,-1)$. The map $C\to B$ is finite:
it is a nonconstant proper map of smooth compact curves. In the
degree-one case it is an isomorphism and hence a section.
Adjunction along this section gives
$K_C=\OO_C(-C)$, proving both normal identifications in
\eqref{cdp:eq:C.EL.1}, including their dual-normal map.

\item In the degree-two case $K_T+C$ is trivial, so adjunction gives
$K_C\simeq\OO_C$ and $g(C)=1$. Riemann--Hurwitz restricts
$g(B)$ to zero or one. For $g(B)=1$, apply \cref{lem:II:anticanonical-bisection}
with $(V,B,q,C,h)=(T,B,q,C,q|_C)$. The preceding degree and
adjunction calculations give a smooth connected elliptic
degree-two curve with $K_T(C)\simeq\OO_T$, contradicting that
lemma. Under these identity maps, \eqref{cdp:eq:C.EL.7} is the
same Leray-identified finite pullback and \eqref{cdp:eq:C.EL.8}
is the residue sequence of $C\subset T$.
Thus $B=\PP^1$ and
$T=\mathbb F_e$. The negative section $C_0$ has
$(-K_T)\cdot C_0=2-e$. A distinct irreducible effective
anticanonical curve has nonnegative intersection with $C_0$,
so $e\le2$. This proves \eqref{cdp:eq:C.EL.2}.

\item On the plane, write $C\sim dH$. Equation \eqref{cdp:eq:C.EL.6} instead
gives $ad=-3$. Thus $(d,a)$ is $(1,-3)$ or $(3,-1)$,
giving exactly the stated line and cubic alternatives.

\item In the section case, apply
\cref{prop:II:section-fibre-contraction} to the same family and ruling,
with $(p,m_T,m_P,r)=(q,1,1,0)$. The only boundary is $C$,
so there are no $Q_j$ or marked points $b_j$. The normal lines
in that proposition are $M=\OO_B$ and $L=K_B$ by
\eqref{cdp:eq:C.RC.6}. Its blowdown is therefore the given ruling
contraction, $P\to P'$ is the Cartier blowup isomorphism, and
$df'$ is the normal quotient $V\twoheadrightarrow\OO_B$ with
no zero on $B$. Thus $f'$ is a submersion along the centre.
The exceptional divisor and its marked section are exactly
$(\PP_B(V),\PP_B(K_B))$; \eqref{cdp:eq:C.RC.10} with $M=\OO_B$
gives $N_{T/U}=\OO_{\PP_B(V)}(-1)=\OO_T(-C)$.
These are all the maps and the reduced SNC conclusion in
\eqref{cdp:eq:C.EL.3}, including the Euler equality from
\eqref{cdp:eq:C.RC.7}. The extension keeps its subline $K_B$
and its quotient trivialization $df'$.

Serre duality gives
$H^1(B,K_B)\simeq H^0(B,\OO_B)^*\simeq\C$.
In local normal frames the transition matrices have the form
\begin{equation}
        \begin{pmatrix}k_{ij}&a_{ij}\\0&1\end{pmatrix}.
                                                        \hypertarget{tag-00G8}{}\label{cdp:eq:C.EL.9}
\end{equation}
Changing a lift of the quotient changes $a_{ij}$ by its
\v{C}ech coboundary; rescaling the subline multiplies the
extension class by that scalar. Thus there are the zero and
nonzero orbits when endpoint frames are allowed to change.
The matrices \eqref{cdp:eq:C.EL.9} describe the normal bundle of the contracted centre with these endpoint frames.

\item In the plane-line case the normal line is $\OO_{\PP^2}(-1)$.
Apply \cite[Main Theorem, pp.~483--484; \S3, pp.~499--501]{Nakano1970},
with the normal-line criterion of
\cite[\S2, Theorem~2\(^{\prime}\) and proof, pp.~640--642;
Proposition~9\(^{\prime}\) and proof, p.~642]{FujikiNakano1971},
to $(\widetilde X,S,M,r)=(U,T,\{p\},3)$. It gives an
ordinary point blowdown onto a smooth threefold. The exceptional
coefficient one says that the descended $f'$ has a nonzero
linear term at that point. It is consequently a submersion
there. Its central surface pulls back to its point blowup,
with exceptional curve the specified line $C$, proving
\eqref{cdp:eq:C.EL.4}. The normal-dual map gives $C^2_P=-1$, consistently
with this surface blowdown.

\item The curve-contraction Euler equality is the same
$r=0$ specialization of \eqref{cdp:eq:C.RC.7}. For the point
contraction apply \cref{lem:II:blowup-topology} to the point
centre and the compact full inverse-image support: its
exceptional fibre is $\PP^2$, giving $e(T)-1=3-1=2$.
The same maps on a compact total space give the identical
Euler changes there. The descended family is unchanged over
the punctured base, and its fibres there are the original
marked two-tori. Each contraction removes one central
component; every iteration therefore terminates. Summing
these equalities gives \eqref{cdp:eq:C.EL.5}. The inverse construction uses the
smooth centre, its normal extension, the local ordinary
blowup charts and their unchanged exterior identifications.
Reversing the maps restores the initial threefold near its
support, so this comparison retains all neighbourhood data.

\item Fix the elliptic $B$, choose another elliptic curve
$E$, and set $A=B\times E$. Use the compact smooth twistor
torus family \eqref{cdp:eq:C.PT.13} with period order $R=1$ or $R=2$.
Its central torus is exactly $A$. Fix the curve $B\times\{0\}$
in that central member. In the primitive period cover
near this curve, \eqref{cdp:eq:C.PR.13} writes the deck transformation as
\begin{equation}
 (p,w,t)\longmapsto
       (p+\delta(t),\gamma(t)w,t),\qquad
 \delta(t)=t^R d(t),\quad d(0)\ne0,
                                                        \hypertarget{tag-00FZ}{}\label{cdp:eq:C.EL.10}
\end{equation}
where $B=\C^*/\gamma(0)^{\Z}$ and $0<|\gamma(0)|<1$.
The curve is $p=t=0$. Differentiating these deck
coordinates gives its normal transition over this elliptic
cover:
\begin{equation}
 V:\quad
 \begin{pmatrix}p\\t\end{pmatrix}\longmapsto
 \begin{pmatrix}1&\delta'(0)\\0&1\end{pmatrix}
 \begin{pmatrix}p\\t\end{pmatrix}.
                                                        \hypertarget{tag-00G0}{}\label{cdp:eq:C.EL.11}
\end{equation}
The term $\gamma'(0)w$ is tangent to the central elliptic
curve and hence vanishes in this normal quotient. The other
primitive period has identity normal transition.

\item For $R=2$, \eqref{cdp:eq:C.EL.11} is the identity, so the normal extension
is split. For $R=1$, its upper entry is nonzero. A splitting
would supply a holomorphic $b$ on $\C^*$ satisfying
$b(\gamma(0)w)-b(w)=\delta'(0)$, with one possible opposite
sign from changing the frame convention. The Laurent
expansion on $\C^*$ has zero constant coefficient on the
left and a nonzero one on the right. No such splitting
exists. Thus its extension class is nonzero in $H^1(B,\OO_B)$.
To prescribe the nonzero class while keeping the endpoint
frames and the quotient $df'$ fixed, use instead the period family $V_t(w)=w+a tJ\overline w$, with $a\in\C^*$.
Its central marked torus is unchanged, and differentiation
multiplies the cocycle in \eqref{cdp:eq:C.EL.11} by $a$. Since
$H^1(B,K_B)$ is one-dimensional and the cocycle just proved
is nonzero, this realizes the prescribed class exactly.
The real period determinant is $(1+|at|^2)^2$. At infinity,
dividing the period frame by $at$ gives
$J\overline w+(s/a)w$, with $s=1/t$; this is the same
holomorphic compact completion used in \eqref{cdp:eq:C.PT.13}.

\item Blow up this central curve in the compact smooth
twistor family. The exceptional marked pair is exactly
$\bigl(\PP_B(V),\PP_B(\OO_B)\bigr)$, by the ordinary
normal-bundle description just proved. The strict central
torus is unchanged because its centre is Cartier there.
Ordinary coordinates for the ideal $(p,t)$ have $t=xy$,
so the full new fibre has its two components with
multiplicity one. Its Euler number is zero by \eqref{cdp:eq:C.EL.3},
and all other fibres are the unchanged smooth two-tori.

\end{enumerate}
\end{proof}

\subsection{Canonical-flat central fibres}\label{subsec:II:canonical-flat-fibres}



\begin{proposition}[Trivial local canonical class, triple points and the complete zero case]\hypertarget{tag-01EB}{}\label{cdp:st:C.2.UN.65}
\StatementRecord{cdp:st:C.2.UN.65}{proved}{01EB}{\cref{cdp:st:C.2.UN.17,cdp:st:C.2.UN.5,lem:II:negative-kodaira-models,cdp:st:C.2.UN.19,cdp:st:C.2.UN.64,lem:II:reduced-anticanonical-hopf,lem:II:normal-product-collars,cdp:st:C.2.UN.57}; \cite[Lemmas~2.1--2.2, pp.~959--960]{Kulikov1977}; \cite[Theorem~1\(^{\prime}\), p.~340; Appendix, p.~342]{Kulikov1981}; \cite[Theorem~3.1, author preprint pp.~2--3]{KamenovaTori}}

Let $f:X\to\Delta$ be a proper degeneration of connected complex
two-tori, with smooth three-dimensional total space and full
reduced simple-normal-crossing central fibre
$S=\sum_iV_i$. Suppose that the orders of the canonical
evaluation along all $V_i$ are equal. Equivalently, the canonical
line of a neighbourhood of $S$ is trivial. Let $t$ be the number
of triple points and $r$ the number of rational double curves.
Then
\begin{equation}
                         e(S)=t,\qquad 2r=3t.
\hypertarget{tag-00DN}{}\label{cdp:eq:C.CZ.1}
\end{equation}
In particular $e(S)$ is a nonnegative even integer.

The complete zero-Euler list in this class is as follows.
The fibre is either a smooth two-torus, or is a cycle of smooth
unblown components, each a marked diagonal primary or secondary
Hopf surface, or a marked surface
\begin{equation}
                   \PP(\OO_E\oplus L),
          \qquad E\text{ elliptic},\quad L\in\operatorname{Pic}^0(E).
\hypertarget{tag-00DT}{}\label{cdp:eq:C.CZ.2}
\end{equation}
The boundary curves are exactly the two Hopf axes or the two
ruled sections. Their normal lines are dual across
every edge, and the integral complement-lattice product
is $M=I$.

Conversely every specified marked cycle in this list with
the stated dual-normal data and $M=I$ has a compact realization
over $\PP^1$, retaining all components and elliptic gluing
maps, with smooth total space and trivial canonical line
near the specified fibre. Its other singular fibre is the
conjugate marked cycle and all reduced fibre supports have
Euler number zero. The construction chooses normal-product
constants. A smooth two-torus is realized by its product
with $\PP^1$.
\end{proposition}


\begin{proof}

\begin{enumerate}[label=\arabic*.,leftmargin=*]
\item After shrinking the disc, apply \cref{cdp:st:C.2.UN.17} to the same family with every full-fibre multiplicity equal to one. Its evaluation generator $\Omega$ has divisor $\sum_i k_iV_i$.
If all its orders are $k$, the quotient
$\Omega/f^k$ is a nowhere-zero holomorphic three-form.
This is the trivialization of $K_X$ used below.
Conversely a nowhere-zero three-form generates this
direct image: every other section is its product with
a holomorphic function constant on the compact connected
fibres, hence pulled back from the disc. Its evaluation
orders are all zero. This proves the equivalence in the
statement by these multiplication and division maps. The smooth-total reduced SNC hypotheses are Kulikov's good-degeneration hypotheses, \cite[pp.~957--958]{Kulikov1977}. In \cite[Theorem~1\(^{\prime}\), p.~340]{Kulikov1981}, existence of a good modification with trivial canonical line is characterized by the minimum-order condition on class VII components. Here the equal orders give that trivial line on the given neighbourhood itself through $\Omega/f^k$. The correction to the 1977 Enriques assertion is stated in the later paper's Appendix, p.~342; the principal-divisor and coherent-stratum comparisons used below are the 1977 Lemmas~2.1--2.2.

\item Put $D_i=\sum_{j\ne i}(V_i\cap V_j)$, with each connected
double curve included once, and let $\nu:\coprod_iV_i\to S$ be
the normalization. For each connected double curve
$C\subset V_i\cap V_j$, let $T_C$ be its reduced divisor of
triple points. In the residue construction of
\eqref{cdp:eq:C.NC.6}, take $T_i=V_i$, $C_\alpha=C$ and
$R_\alpha=T_C$.
The two conductor sheets over $C$ are its copies in $V_i$ and
$V_j$; ordering these two global components trivializes their
sign line by $(1,-1)$. The form $\Omega/f^k$ and the principal
divisor $S=(f)$ trivialize $\omega_S=(K_X\otimes\OO_X(S))|_S$.
Thus the same local residue map gives the isomorphism
\begin{equation}
                           K_{V_i}\simeq\OO_{V_i}(-D_i).
\hypertarget{tag-00DU}{}\label{cdp:eq:C.CZ.3}
\end{equation}
Adjunction on $V_i$ gives
\begin{equation}
             K_C=(K_{V_i}+C)|_C\simeq\OO_C(-T_C),
             \qquad 2g(C)-2=-\#T_C.
\hypertarget{tag-00DV}{}\label{cdp:eq:C.CZ.4}
\end{equation}
Every double curve is therefore elliptic and has no
triple points, or is rational and has exactly two.
Each triple point lies on exactly three double curves.
Counting these incidences proves $2r=3t$.

\item Restriction of the trivial line $\OO_X(S)$ to $C$
also gives the exact triple-point comparison
\begin{equation}
 N_{C/V_i}\otimes N_{C/V_j}\otimes\OO_C(T_C)
                                  \simeq\OO_C.
\hypertarget{tag-00DW}{}\label{cdp:eq:C.CZ.5}
\end{equation}
Indeed $\OO_X(V_i)|_C=N_{C/V_j}$,
$\OO_X(V_j)|_C=N_{C/V_i}$, and every remaining
component contributes its transverse triple
point. In degrees this is
$C_{V_i}^2+C_{V_j}^2=-\#T_C$.
This is also the direct line-bundle proof of
\cite[Lemma~2.1, p.~959]{Kulikov1977} in the present notation.

By \eqref{cdp:eq:C.CZ.3}, $K_{V_i}^2=D_i^2$.
The self-intersections of the double curves sum,
by \eqref{cdp:eq:C.CZ.5}, to $-3t$. The cross terms in the
boundary squares contribute $2$ on each of the
three surfaces through a triple point. Consequently
\begin{equation}
                         \sum_iK_{V_i}^2=3t.
\hypertarget{tag-00DX}{}\label{cdp:eq:C.CZ.6}
\end{equation}
With the same globally ordered branches, the alternating map
$(f_i)\mapsto(f_j-f_i)_{i<j}$ and its triple-point relation
from \eqref{cdp:eq:C.NC.7} give exactly
\begin{equation}
 0\longrightarrow\OO_S\longrightarrow\bigoplus_i\OO_{V_i}
  \longrightarrow\bigoplus_C\OO_C
  \longrightarrow\bigoplus_p\C_p\longrightarrow0.
\hypertarget{tag-00DY}{}\label{cdp:eq:C.CZ.7}
\end{equation}
Proper flatness gives $\chi(\OO_S)=\chi(\OO_{F_t})=0$.
Only the rational double curves contribute to
$\sum_C\chi(\OO_C)$, so \eqref{cdp:eq:C.CZ.7} gives
\begin{equation}
                         \sum_i\chi(\OO_{V_i})=r-t.
\hypertarget{tag-00DZ}{}\label{cdp:eq:C.CZ.8}
\end{equation}
The sequence and equality also agree with
\cite[Lemma~2.2 and proof, pp.~959--960]{Kulikov1977}.

\item Hirzebruch--Riemann--Roch and complex Gauss--Bonnet
on each compact complex surface give
$e(V_i)=12\chi(\OO_{V_i})-K_{V_i}^2$.
The normalization and conductor strata in \eqref{cdp:eq:C.NC.5},
with $T_i=V_i$, $C_\alpha=C$ and $\tau=t$, give
\begin{equation}
 \begin{aligned}
 e(S)&=\sum_i e(V_i)-\sum_C e(C)+t\\
     &=12(r-t)-3t-2r+t\\
     &=10r-14t=t,
 \end{aligned}
\hypertarget{tag-00E0}{}\label{cdp:eq:C.CZ.9}
\end{equation}
where the last equality uses $2r=3t$.
The count $2r=3t$ also makes $t$ even.

\item Assume $e(S)=0$. Equations \eqref{cdp:eq:C.CZ.1} and \eqref{cdp:eq:C.CZ.4}
give no triple points and only elliptic double
curves. If there is just one component, the
central fibre is smooth and reduced. The
smooth-family comparison in \ref{cdp:st:C.2.UN.5} makes
it a complex two-torus. We therefore treat
the connected reducible case, where every
$D_i$ is nonempty.

\item Equation \eqref{cdp:eq:C.CZ.3} gives
$H^0(V_i,mK_{V_i})=H^0(V_i,\OO(-mD_i))=0$
for every $m>0$: its image in $H^0(V_i,\OO)=\C$
would be a constant vanishing on $D_i$.
Thus every component has Kodaira dimension $-\infty$. Apply \cref{lem:II:negative-kodaira-models} to its smooth compact minimal model. Its K\"ahler case is rational or a projective-line bundle, and its non-K\"ahler case is of class VII.

\item For a ruled minimal model over a curve $B$
of genus at least two, every elliptic curve
on $V_i$ maps constantly to $B$, by
Riemann--Hurwitz. A general ruling fibre
$F\simeq\PP^1$ avoids their vertical images
and every blowup centre. It would satisfy
$D_i\cdot F=0$, whereas \eqref{cdp:eq:C.CZ.3} and
adjunction give $D_i\cdot F=-K_{V_i}\cdot F=2$.
This contradiction leaves ruled genus zero
or one, and class VII.

\item Their coherent Euler characteristics are,
respectively, one, zero and zero. Equation
\eqref{cdp:eq:C.CZ.8} is now $\sum_i\chi(\OO_{V_i})=0$,
so no rational component occurs. A ruled
component over an elliptic curve has Euler
number equal to its number of point blowups.
A class VII component has $b_1=b_3=1$ and
$e(V_i)=b_2(V_i)\ge0$. Since all double
curves are elliptic, \eqref{cdp:eq:C.CZ.9} is also
$0=e(S)=\sum_i e(V_i)$. Every term is
therefore zero. The ruled components have
no blowups. Every class VII component is
minimal with $b_2=0$, and its minimal-model restriction in \ref{cdp:st:C.2.UN.19} makes
it a Hopf surface. This proves the claimed
absence of blowups and identifies both
remaining kinds of surface.

\item Let $\pi:V\to E$ be one of the minimal ruled
surfaces over an elliptic curve. Each component
of $D$ is elliptic and maps finitely to $E$:
it cannot be contained in a rational ruling
fibre. The map is unramified by
Riemann--Hurwitz. Since $D\cdot F=2$, the
only possibilities are two degree-one
sections or one connected elliptic
bisection $C$.

\item A connected bisection $C$ would be the entire boundary of $V$. Since the fibre is SNC without triple points, it lies in one other component. Apply \cref{cdp:st:C.2.UN.64} with $(T,C,q,B)=(V,C,\pi,E)$. Its degree-two alternative \eqref{cdp:eq:C.EL.2} forces $E\simeq\PP^1$, contradicting that $E$ is elliptic. Under this identity map, \eqref{cdp:eq:C.EL.8} is the residue sequence of $C\subset V$, and its Serre-dual restriction \eqref{cdp:eq:C.EL.7} is the same finite pullback $H^1(E,\OO_E)\to H^1(C,\OO_C)$ through the Leray isomorphism $\pi^*$. Thus the trace and residue exclusion applies to this very bisection.

\item Write the ruling as $V=\PP_E(W)$ with the lines convention. The two disjoint sections give line subbundles $L_0,L_\infty\subset W$. Their sum map $L_0\oplus L_\infty\to W$ is an isomorphism on each fibre, because the two lines are distinct, and hence is a bundle isomorphism. Tensoring it by $L_0^{-1}$ gives
$\PP_E(W)\simeq\PP_E(\OO_E\oplus L)$, where $L=L_\infty\otimes L_0^{-1}$, and sends the two marked sections to the coordinate lines. On the chart with nonzero first coordinate, the ratio of the second coordinate to the first lies in $L$. The remaining marked section is its zero section; removing it gives exactly $L^\times$. Its degree is determined below by the nearby fundamental group.

\item Let $H$ be a Hopf component and let $D$ be its double-curve boundary. It is nonempty, reduced and a disjoint union of smooth elliptic curves, as proved above; \eqref{cdp:eq:C.CZ.3} is exactly $K_H(D)\simeq\OO_H$. Apply \cref{lem:II:reduced-anticanonical-hopf} to this same pair $(H,D)$. It identifies $H$ with a diagonal primary or secondary quotient and $D$ with its two individually marked axes.

\item Every component now has exactly two
boundary curves. The connected
incidence graph is a cycle. At a
ruled vertex its complement is the
nonzero line $L^\times$ over
$E$. For $d=\deg L$ its group is
\begin{equation}
 \pi_1(L^\times)=
 \langle a,b,c\mid [a,c]=[b,c]=1,
                               \ [a,b]=c^d\rangle .
\hypertarget{tag-00DQ}{}\label{cdp:eq:C.CZ.12}
\end{equation}
This is the circle-bundle
extension: trivialize over an
oriented elliptic parallelogram;
the boundary commutator is the
degree-$d$ clutching loop of its
fibre. The fibre generator $c$
has infinite order, since the
homotopy sequence injects
$\pi_1(\C^*)$ when $\pi_2(E)=0$.

\item For a given smoothing take the
component cores and the $xy=t$ collars. A Hermitian radial
retraction identifies a ruled
core with an annulus bundle in
$L^\times$; either boundary
inclusion induces the isomorphism
onto \eqref{cdp:eq:C.CZ.12}. For a diagonal
Hopf core, the same statement is
its rank-three logarithmic
lattice comparison. The dual-normal
maps identify the intervening
circle bundles. Apply \cref{lem:II:normal-product-collars} to this reduced SNC fibre, which has no triple points by \eqref{cdp:eq:C.CZ.1}. Its map $\Psi$ uses the two divisor sections and the normal-dual identification along each of these same double curves. It transports the component cores and gives the attaching map $y=t/x$ on their normal circle bundles. For the ruled vertex, the Hermitian metric is on this same degree-$d$ line $L$; the collar construction applies for every integer $d$.

\item The resulting cycle of group
isomorphisms is a graph of groups
whose vertex groups inject into
the total group. Explicitly cut
one edge, compose the remaining
isomorphisms, and obtain a single
vertex group with one stable
letter implementing the closing
automorphism. Its normal form
contains that vertex group as
a subgroup. Since the nearby group is $\Z^4$, every
injected vertex group is abelian.
Equation \eqref{cdp:eq:C.CZ.12} and the
infinite order of $c$ force
$d=0$ at every ruled vertex.
This proves \eqref{cdp:eq:C.CZ.2} by the
collar and clutching maps.

\item The remaining component groups
are all their rank-three
abelian lattices. The same cycle
calculation gives
$\pi_1(F_t)=\Lambda\rtimes_M\Z$.
Its injected lattice and the
abelian group $\Z^4$ force
$M=I$. We have obtained exactly
the marked data of \ref{cdp:st:C.2.UN.57}, with
all components and their two boundary curves as proved above.

\item For every such marked cycle,
\ref{cdp:st:C.2.UN.57} gives the proper
plumbing, the compact two-ended
period family, and its exact
punctured-family gluing. It
retains all marked components
and elliptic maps. We check
the required local canonical
condition on this construction.

\item On its common period charts
put
\begin{equation}
                       \Omega=dt\wedge d\eta_1\wedge d\eta_2.
\hypertarget{tag-00DR}{}\label{cdp:eq:C.CZ.13}
\end{equation}
Here $\eta$ is the common
vector coordinate obtained
from each component chart
by its moving frame
$P_i(t)$. Thus every overlap
has identity linear part in
this common frame.
Every period and affine edge
translation depends holomorphically
on $t$; its differential changes
the vector differentials only by
multiples of $dt$. Thus this
three-form descends through the
lattice and collar maps.
It is nonzero on component cores.
At a double curve the relative
two-form is a holomorphic unit
times $dx/x\wedge dz$, where
$z$ is an elliptic coordinate
and $xy=t$. Hence
\begin{equation}
             dt\wedge(dx/x\wedge dz)
                               =dy\wedge dx\wedge dz.
\hypertarget{tag-00DS}{}\label{cdp:eq:C.CZ.14}
\end{equation}
The right side is holomorphic
and nowhere zero. The moving
Hopf and ruled frames have
nonzero determinant at the
central value, so the displayed
unit never vanishes after
shrinking the disc. The form
therefore extends across all
of $S$ as a nowhere-zero
three-form. Its evaluation
orders are equal to zero.
The compact realization in
the statement now follows
from \ref{cdp:st:C.2.UN.57}, with the same
verification at its conjugate
fibre. The smooth torus case is its product family.

For a K\"ahler total space, restriction of the K\"ahler form makes every component K\"ahler. A Hopf component has $H^2(H,\R)=0$ by \cref{lem:II:hopf-curves} and its bundle-cohomology proof, whereas a K\"ahler form would have positive square integral; exactness and Stokes give a contradiction. Thus the zero-Euler list in this case consists of a torus or a cycle of the elliptic ruled surfaces in \eqref{cdp:eq:C.CZ.2}. Kamenova's necessary list in \cite[Theorem~3.1 and proof, author preprint pp.~2--3]{KamenovaTori}, under the same smooth-total K\"ahler and trivial-canonical hypotheses, also permits elliptic ruled paths. A terminal vertex of such a path would have the connected bisection excluded above by the identical residue and finite-trace maps \eqref{cdp:eq:C.CZ.10} and \eqref{cdp:eq:C.EL.7}. The marked cycles here carry the compact realizations and volume forms just constructed.

\end{enumerate}
\end{proof}

\section{Torsion-normal boundaries and elliptic cores}\label{sec:II:torsion-boundaries-elliptic-cores}

\subsection{Elliptic fibrations with vertical boundary}\label{subsec:II:elliptic-alignment}

\begin{lemma}[Elliptic alignment on a component]\hypertarget{tag-P260}{}\label{lem:II:vertical-elliptic-alignment}
\StatementRecord{lem:II:vertical-elliptic-alignment}{proved}{P260}{\cref{cdp:st:C.2.UN.16,lem:II:elliptic-first-integral,cdp:st:C.2.UN.29,lem:II:identity-monodromy-elliptic,cdp:st:C.2.UN.28}; \cite[Chapter~1, Proposition~2.2 and the preceding point-blowup comparison, p.~8]{Peters1990}}
Let $f:U\to\Delta$ be a proper holomorphic degeneration of connected complex two-tori from a smooth threefold, with full reduced normal-crossing central fibre $S$. Let $\nu:T\to S_i$ be the smooth normalization of a component and put $D=(\nu^{-1}(\operatorname{Sing}S))_{\mathrm{red}}$. Both inverse branches of a self-double curve are included in $D$. Let $q:T\to B$ be a proper holomorphic map to a smooth compact connected curve, with connected fibres and smooth elliptic general fibre. Suppose every component of $D$ is vertical for $q$.

The map $q$ has identity integral $H_1$ monodromy on its smooth base locus. Contracting its vertical exceptional $(-1)$-curves gives a morphism $\tau:T\to T_0$ and a relatively minimal map $q_0:T_0\to B$, with $q=q_0\circ\tau$. Every fibre of $q_0$ is a positive multiple of a smooth elliptic curve; its regular fibres are torsors under one fixed elliptic curve. Moreover
\begin{equation}
             R^1q_*\OO_T=\OO_B,\qquad
                         p_g(T)=g(B)\le1,
\hypertarget{tag-00RF}{}\label{cdp:eq:C.MR.8}
\end{equation}
and $\chi(\OO_T)=0$. The maps $q$ and $q_0$ have the same smooth fibres away from finitely many values. When $q$ is relatively minimal, $\tau$ is the identity.
\end{lemma}
\begin{proof}
Apply \cref{lem:II:cyclic-component-cover,cdp:st:C.2.UN.16} with $m_i=1$. The root change is the identity. The resolution $\rho:U'\to U$ has a smooth strict multiplicity-one component $E$ and a proper birational factorization
\[
 E\xrightarrow{\ r\ }T\xrightarrow{\nu}S_i,
 \qquad \rho|_E=\nu\circ r.
\]
The map $r$ is a composition of point blowups by \cite[Proposition~3.12, p.~20]{PopescuPampu2016}. In particular $r_*\OO_E=\OO_T$, and it is an isomorphism off finitely many points of $T$.

The saturation and restriction maps \eqref{cdp:eq:C.CF.3}--\eqref{cdp:eq:C.CF.7} give a two-dimensional complex space of commuting meromorphic fields on $E$. Transport them through $r$ on the common open isomorphism. This identifies the meromorphic function fields and derivations, so it preserves their complex-linear independence and brackets. By \cref{lem:II:semicomplete-transport}, every constant linear combination of the resulting fields $V_1,V_2$ on $T$ is semicomplete on its full holomorphic locus.

To locate their poles, push the ambient representatives through $\rho$. They are meromorphic vertical fields on $U$, holomorphic over $\Delta^*$; all their polar divisors are therefore central. The map $\rho$ is an isomorphism at the generic point of $S_i$, where the saturation makes their coefficient orders nonnegative. Thus $S_i$ is not a polar divisor. On $S_i\setminus\operatorname{Sing}S$, this is the only central branch, so the pushed fields are holomorphic near it, by extension of their coefficient functions across codimension two. Their restrictions pull back through the local isomorphism $\nu:T\setminus D\to S_i\setminus\operatorname{Sing}S$ to $V_1,V_2$. Hence their polar curves on $T$ are contained in $D$.

Remove finitely many values of $B$ so that the fibres are smooth elliptic and miss $D$. The fields are holomorphic on each entire compact fibre. In a base coordinate, the coefficient of $dq(V_i)$ is therefore constant along that fibre. Proper connected-fibre descent gives meromorphic vector fields $Z_i$ on $B$ with $dq(V_i)=q^*Z_i$; their brackets vanish.

If $Z_1\ne0$, commutation makes the meromorphic ratio $Z_2/Z_1$ constant, say $c$. Put $Y=V_1$ and $V=V_2-cV_1$. Complex-linear independence makes $V\ne0$; it is vertical, while $Y$ is generically transverse, so $V\wedge Y\ne0$. Choose a nonconstant meromorphic function $h:B\to\PP^1$. The function $h\circ q$ is a first integral of $V$. Its graph is already defined on $T$, and its Stein factor is exactly
$T\xrightarrow{q}B\xrightarrow{h}\PP^1$: the first map has connected fibres and the second is finite. Apply \cref{lem:II:elliptic-first-integral} to this $T,V,Y$ and this Stein map. Its inverse-time continuation and period-injection proof gives identity monodromy on the full smooth base of this same $q$. The case $Z_2\ne0$ uses the exchanged basis.

If $Z_1=Z_2=0$, both fields are vertical. On a regular elliptic fibre they are translation fields, so their ratio is constant there. Proper descent gives
$V_2=(q^*H)V_1$ for a meromorphic function $H$ on $B$. It is nonconstant because the fields are independent over $\C$. Put $q_E=q\circ r$. The equality $r_*\OO_E=\OO_T$ makes $q_E$ a connected-fibre map, and the ratio of the original saturated fields on $E$ is $q_E^*H$. Its graph is already defined, and its Stein factor is exactly $E\xrightarrow{q_E}B\xrightarrow H\PP^1$. Apply \cref{cdp:st:C.2.UN.29} to the multiplicity-one component $E$, these ambient saturated fields and this ratio. Its return-period identity \eqref{cdp:eq:C.RG.12} and period-injection step give identity integral monodromy for $q_E$. Remove also the finitely many images under $q$ of the centres of $r$. On the remaining base, $r$ restricts to biholomorphisms $r_b:(q_E)^{-1}(b)\to q^{-1}(b)$, and their maps $(r_b)_*$ identify the integral homology local systems. The inclusion of this base into the full smooth base of $q$ is surjective on fundamental groups. Thus the integral monodromy of $q$ is the identity as well.

In either case, contract the vertical $(-1)$-curves successively. Each contraction lowers $b_2$, so the process is finite. The map descends at each step because its contracted curve maps to a point; the descended base coordinate is holomorphic through the equality of the direct image of the structure sheaf with the structure sheaf of the smooth blowdown. This gives $q=q_0\circ\tau$ and leaves all smooth fibres outside finitely many values unchanged. The map on fundamental groups from the base with these extra values removed to the smooth base of $q_0$ is surjective, since loops avoid finitely many additional points. Hence the monodromy of $q_0$ is still the identity. Apply \cref{lem:II:identity-monodromy-elliptic} to the compact smooth surface $T_0$, this proper connected-fibre relatively minimal $q_0$, and the full smooth base. It gives the fixed elliptic curve, all full smooth-multiple fibres, $R^1(q_0)_*\OO_{T_0}=\OO_B$, $p_g(T_0)=g(B)$ and $\chi(\OO_{T_0})=0$.

The structure-sheaf direct images of an ordinary point blowup $\mu$ satisfy $\mu_*\OO=\OO$ and $R^1\mu_*\OO=0$. The first equality is Hartogs extension at its smooth centre. For the second, use the two blowup charts $(x,v)$ and $(u,y)$, with $u=1/v$ and $y=xv$. The sets $|x|<\varepsilon,|v|<2$ and $|y|<\varepsilon,|u|<2$ are Stein; their Stein overlap has $1/2<|v|<2$, $|x|<\varepsilon$ and $|xv|<\varepsilon$. They form a cofinal family of neighbourhoods of the exceptional curve as $\varepsilon$ decreases. A holomorphic cocycle on the overlap has a normally convergent Laurent expansion in monomials $x^m v^n$, $m\ge0$. Terms with $n\ge0$ extend to the first chart after decreasing $\varepsilon$; terms with $n<0$ become $y^m u^{m-n}$ and extend to the second. Cauchy estimates on smaller closed annuli give normal convergence of both parts. This splitting makes every cohomology class zero on a smaller neighbourhood and proves the vanishing of the direct-image stalk at the centre. Iteration and the Leray spectral sequence give $R^1q_*\OO_T=R^1(q_0)_*\OO_{T_0}$. The point-blowup invariance of $p_g$ and $\chi$, \cite[Chapter~1, Proposition~2.2, p.~8]{Peters1990}, gives their stated values on $T$. Finally \cref{cdp:st:C.2.UN.28}, applied to this same smooth normalization component $T$, gives $p_g(T)\le1$.
\end{proof}

\subsection{Torsion-normal elliptic curves on rational components}\label{subsec:II:torsion-normal-rational-obstruction}

\begin{proposition}[A torsion-normal elliptic boundary cannot lie on a rational component]\hypertarget{tag-01ED}{}\label{cdp:st:C.2.UN.67}
\StatementRecord{cdp:st:C.2.UN.67}{proved}{01ED}{\cref{lem:II:vertical-elliptic-alignment,lem:II:identity-monodromy-elliptic,cdp:st:C.2.UN.37,cdp:st:C.2.UN.50,cdp:st:C.2.UN.53}; \eqref{cdp:eq:C.RH.1}--\eqref{cdp:eq:C.RH.10}}

Let \(f:U\to\Delta\) be a proper degeneration of connected
complex two-tori, with smooth total threefold and full
reduced normal-crossing central fibre \(S\). Let
\(\nu:P\to S_i\) be a smooth rational normalization, and
let \(D=(\nu^{-1}(\operatorname{Sing}S))_{\mathrm{red}}\),
including both inverse branches of each self-double curve.
There is no smooth elliptic component \(J\) of \(D\) whose
normal line \(N_{J/P}\) is torsion and which is disjoint
from every other component of \(D\).

In particular a component with rational smooth normalization,
including a rational elliptic surface and any point blowups
of a rational surface, cannot be adjacent to a component
whose smooth normalization is elliptic Hopf or nonlinear Hopf.
No Hopf leaf can have a parent with rational smooth normalization. The assertion permits rational double
curves and triple points elsewhere in the central fibre.

\end{proposition}

\begin{proof}

\begin{enumerate}[label=\arabic*.,leftmargin=*]
\item Suppose that such a component and curve exist, and put
\(N=N_{J/P}\), with exact order \(m\ge1\).
The degree of \(N\) is zero, hence \(J^2=0\).
For every \(k\ge1\) use the divisor restriction
sequence
\begin{equation}
 0\longrightarrow\OO_P((k-1)J)
 \longrightarrow\OO_P(kJ)
 \longrightarrow N^{\otimes k}\longrightarrow0.
 \hypertarget{tag-00ZG}{}\label{cdp:eq:C.RH.1}
\end{equation}
For \(1\le k<m\), the last line is a nontrivial
degree-zero line on the elliptic curve. It has
both \(H^0\) and \(H^1\) zero: a nonzero section
would be nowhere zero, and Serre duality gives
the \(H^1\) assertion. Since \(P\) is rational,
\(H^1(P,\OO_P)=0\) and \(H^0(P,\OO_P)=\C\).
Induction in the long exact sequences gives
\begin{equation}
 h^0(P,\OO_P(kJ))=1,\qquad
 H^1(P,\OO_P(kJ))=0\quad(0\le k<m).
 \hypertarget{tag-00ZI}{}\label{cdp:eq:C.RH.2}
\end{equation}
At \(k=m\), the restriction in \eqref{cdp:eq:C.RH.1} is
onto \(H^0(J,\OO_J)=\C\). Thus
\begin{equation}
 h^0(P,\OO_P(mJ))=2,\qquad
 H^0(P,\OO_P(mJ))\longrightarrow H^0(J,\OO_J)
                         \text{ is onto}.
 \hypertarget{tag-00ZJ}{}\label{cdp:eq:C.RH.3}
\end{equation}
One section has divisor \(mJ\); a second restricts
to a nonzero constant on \(J\). Their common
zero locus is empty: on \(J\) the second
does not vanish, and off \(J\) the first
does not vanish. They define the holomorphic map
\begin{equation}
                   q_0:P\longrightarrow\PP^1,\qquad
                             q_0^{-1}(\infty)=mJ .
 \hypertarget{tag-00ZK}{}\label{cdp:eq:C.RH.4}
\end{equation}

\item Use the Stein factor to prove connectedness of the pencil. Write \(P\xrightarrow q B\xrightarrow\beta\PP^1\)
for its Stein factor, and let \(d=\deg\beta\).
The injection \(H^1(B,\OO_B)\to H^1(P,\OO_P)\)
gives \(B\simeq\PP^1\).
Since the support of \eqref{cdp:eq:C.RH.4} is connected,
\(\beta^{-1}(\infty)\) has just one point.
It is totally ramified of order \(d\), and
the fibre of \(q\) there is
\((m/d)J\). Its principal local equation
trivializes \(\OO_P((m/d)J)|_J\), so
\begin{equation}
                   N^{\otimes(m/d)}=\OO_J.
 \hypertarget{tag-00ZL}{}\label{cdp:eq:C.RH.5}
\end{equation}
Exact order \(m\) forces \(d=1\).
Thus \(q=q_0\) has connected fibres.
Adjunction gives \(K_P\cdot J=0\);
a general member \(F\sim mJ\) is smooth
by generic smoothness and satisfies
\begin{equation}
                 2g(F)-2=F\cdot(K_P+F)=0.
 \hypertarget{tag-00ZM}{}\label{cdp:eq:C.RH.6}
\end{equation}
It is a smooth connected elliptic curve.
This constructs the connected elliptic
map on $P$.

\item Every other boundary component \(C\) on
\(P\) is disjoint from \(J\) by hypothesis.
Consequently \(C\cdot q^*(\infty)=mC\cdot J=0\).
If \(q|_C\) were nonconstant, that intersection
would be its positive degree. Hence every
boundary is vertical for this very map \(q\).

\item Apply \cref{lem:II:vertical-elliptic-alignment} with $(T,\nu,D)=(P,\nu,D)$, the connected pencil $q:P\to B=\PP^1$ in \eqref{cdp:eq:C.RH.4}, and the given full reduced normal-crossing torus degeneration. The preceding intersection computation makes every inverse boundary curve vertical. The lemma uses the degree-one component resolution $r:E\to P$ and the fibre identification for $q\circ r$ to prove identity integral monodromy for this precise pencil.

Write $\tau:P\to P_{\min}$ for the lemma's contraction of the vertical exceptional curves. It descends this same $q$ to
\begin{equation}
               q_{\min}:P_{\min}\longrightarrow\PP^1
 \hypertarget{tag-00ZN}{}\label{cdp:eq:C.RH.7}
\end{equation}
Each step is an ordinary point blowdown, and it leaves the marked smooth fibres unchanged. The point-blowup comparison in \cref{lem:II:vertical-elliptic-alignment} gives $\chi(\OO_{P_{\min}})=\chi(\OO_P)=1$.

\item For a full fibre of multiplicity $a$, write its local function as $t=u h^a$, where $u$ is a unit and the divisor of $h$ is primitive. The root chart of \cref{def:II:elliptic-root-curve}, applied in \cref{lem:II:identity-monodromy-elliptic} and used by the alignment lemma on $(T_0,q_0,B)=(P_{\min},q_{\min},\PP^1)$, is
\begin{equation}
                         w^a=u,\qquad s=wh .
 \hypertarget{tag-00ZO}{}\label{cdp:eq:C.RH.8}
\end{equation}
with its exponent $m=a$, the same unit $u$ and the same primitive equation $h$. The root projection is finite \'etale. The canonical-intersection argument in \cref{lem:II:identity-monodromy-elliptic} preserves relative minimality, and its primitive-fibre and translation-inertia maps give the invariant full-fibre coefficient functions. Consequently the conclusion of the alignment lemma for this $q_{\min}$ is
\begin{equation}
 (q_{\min})_*\OO_{P_{\min}}=\OO_{\PP^1},
 \qquad R^1(q_{\min})_*\OO_{P_{\min}}=\OO_{\PP^1}.
 \hypertarget{tag-00ZP}{}\label{cdp:eq:C.RH.9}
\end{equation}
Every full fibre is a positive multiple of a smooth elliptic curve. Coherent Leray gives $\chi(\OO_{P_{\min}})=1-1=0$, contradicting its preserved value one. This proves the asserted obstruction on the original nonminimal component.

\item On a normalization that is an elliptic Hopf
surface every curve
is a fibre support and
has normal of finite
exact order, by the
normalized-root
map in \cref{cdp:st:C.2.UN.37}. A nonlinear
Hopf has its one axis,
whose normal has exact
resonance order, including
the secondary comparison
of \cref{cdp:st:C.2.UN.50}. Distinct Hopf
curves are disjoint.
There can therefore be
no triple point on an
interface with a Hopf
component: such a point
would make two different
boundary curves meet
inside that Hopf surface.
Along an interface between the two distinct global
components, each normalized branch maps isomorphically to
the double curve: the two global component labels preserve
the two local branches of the normal-crossing equation.
The \(xy=t\) derivative identification of their normal
lines therefore gives
\begin{equation}
                     N_{J/P}\otimes N_{J/H}=\OO_J.
 \hypertarget{tag-00ZH}{}\label{cdp:eq:C.RH.10}
\end{equation}
It makes \(N_{J/P}\)
torsion, and the same
absence of triple points
makes \(J\) disjoint
from every other
boundary curve on \(P\).
The obstruction just
proved applies.

\item Every terminal Hopf
pair has one of these
torsion-normal axes,
by the exact
one-boundary canonical
and character comparison
in \cref{cdp:st:C.2.UN.53}. The two
nontorsion axes of
an algebraic-dimension-zero
diagonal Hopf cannot
form such a leaf.
This proves all the
asserted parent and
interface exclusions.

\end{enumerate}
\end{proof}

\subsection{Ruling contractions preserving an elliptic curve}\label{subsec:II:section-preserving-rulings}

\begin{lemma}[Ruling contractions preserving an elliptic curve]\hypertarget{tag-P354}{}\label{lem:II:section-preserving-ruling}
\StatementRecord{lem:II:section-preserving-ruling}{proved}{P354}{\cite[Proposition~3.12, p.~20]{PopescuPampu2016}}
A smooth elliptic curve on $\PP^2$ has square nine. On a Hirzebruch surface its square is eight or nine. In particular none of these curves has square zero.

Let $P$ be a smooth compact nonrational ruled surface and let $J\subset P$ be a smooth elliptic curve with $J^2=0$. Its ruling is a morphism
\begin{equation}
                         p:P\longrightarrow B,
\hypertarget{tag-00T6}{}\label{cdp:eq:C.NP.6}
\hypertarget{tag-013I}{}\SsixEquationAliases{\label{cdp:eq:C.TI.9}}
\end{equation}
where $B$ is elliptic and $p|_J:J\to B$ is an unramified isogeny of degree $d$. The finite base change
\begin{equation}
                 P'=P\times_BJ\longrightarrow P
\hypertarget{tag-0136}{}\label{cdp:eq:C.TI.10}
\end{equation}
has a free deck group $G$, and the inverse image of $J$ is a union of $d$ disjoint sections $J_1',\ldots,J_d'$ with square zero. Their maps to $J\subset P$ are isomorphisms, so their normal lines are the pullbacks of $N_{J/P}$ under these maps.

There is a $G$-equivariant composition of point blowdowns $\pi':P'\to Z'$ over $J$, where $Z'\to J$ is a projective-line bundle. It is an isomorphism on a neighbourhood of every $J_i'$. Its quotient is a composition of point blowdowns $\pi_0:P\to R$ over $B$, with
\begin{equation}
 P\xrightarrow{\pi_0}R=\PP_B(V_0)\xrightarrow{p_0}B,
 \qquad J_0=\pi_0(J),\qquad N_{J_0/R}=N_{J/P}.
\hypertarget{tag-0139}{}\label{cdp:eq:C.TI.13}
\end{equation}
The map $\pi_0$ is an isomorphism on a neighbourhood of $J$. In each reducible ruling fibre of $P'$, the single component meeting one of the $J_i'$ meets all of them and survives; the other components are contracted. These properties specify the exceptional set and the equivariant quotient map. When $d=1$, the base change is the identity and $J_0$ is a section.
\end{lemma}
\begin{proof}
On $\PP^2$, the genus formula $(a-1)(a-2)/2=1$ gives degree $a=3$ and square nine. On $\mathbb F_n$, write the elliptic curve as $aC_0+bf$, where $C_0^2=-n$, $C_0\cdot f=1$, and $f^2=0$. It is distinct from $C_0$, so $a>0$ and $b\ge na$. Adjunction gives
\begin{equation}
                         (a-1)(2b-na-2)=2.
\hypertarget{tag-00RG}{}\label{cdp:eq:C.MR.9}
\end{equation}
For $a=2$ this gives $b=n+2$, $n\le2$ and square eight. The remaining possibility is $a=3$, $n=1$, $b=3$, of square nine. These calculations include $\mathbb F_1$.

For nonrational $P$, compose its point-blowdown map to a minimal ruled surface with the ruling. Every component of a ruling fibre is rational, so $p|_J$ is nonconstant. The base has positive genus. Riemann--Hurwitz therefore gives $g(B)=1$ and makes $p|_J$ unramified. Choosing origins identifies it with an isogeny. The fibre product in \eqref{cdp:eq:C.TI.10} is finite \'etale over $P$. Its deck maps are translations on the second factor, and the inverse image $J\times_BJ$ is the disjoint union of their graphs. This proves the assertions about its sections and their normals.

The surface $P'$ is a point blowup of a projective-line bundle. For a reducible ruling fibre $F=\sum_\alpha m_\alpha E_\alpha$, every $E_\alpha$ is smooth rational, has negative square, and $K_{P'}\cdot F=-2$. Adjunction gives
\begin{equation}
                  \sum_\alpha m_\alpha(2+E_\alpha^2)=2.
\hypertarget{tag-00T7}{}\label{cdp:eq:C.NP.7}
\end{equation}
A section meets a unique multiplicity-one component, at a smooth point of that fibre. If every $(-1)$-component met $J_1'$, there would be at most one such component and it would have multiplicity one. The total positive contribution to the left side would be at most one, and every other contribution would be nonpositive. This contradicts \eqref{cdp:eq:C.NP.7}. Hence a vertical $(-1)$-curve disjoint from $J_1'$ exists. Contract it and continue; the rational ruling descends through each point blowdown, and the number of components of its reducible fibres decreases. Its relatively minimal endpoint is a projective-line bundle.

Every contraction also misses the other $J_i'$. The image of $J_1'$ remains an effective square-zero curve, hence is nef. An image of another section remains disjoint from it and has square at least zero, by the point-blowup square formula. Hodge index on this projective ruled surface makes its square at most zero. It is therefore zero, and the blowdown square formula shows that the centre missed that section as well. Induction proves the neighbourhood assertion for every section.

Two disjoint square-zero sections are numerically equivalent. Their isotropic classes lie in the same zero direction of the orthogonal complement of the nef class of $J_1'$, and intersection with a ruling fibre, equal to one for each section, fixes the proportionality. Thus their intersections with every fibre component agree. The unique multiplicity-one component meeting $J_1'$ consequently meets all the sections. It survives the preceding contraction; every other component of that fibre is contracted.

The deck group permutes the sections and preserves this specified exceptional set. Each connected fibre of $\pi'$ is therefore mapped into a fibre of $\pi'$. The equality $\pi'_*\OO_{P'}=\OO_{Z'}$ descends each deck map holomorphically, and its inverse descends in the same way. This gives the action characterized by $g_{Z'}\pi'=\pi'g_{P'}$. It is free because its action on the ruling base $J$ is free. Hence $R=Z'/G$ is smooth, and $R\to B$ is the projective-line bundle obtained by this finite \'etale descent. The induced proper birational map $\pi_0:P=P'/G\to R$ is an isomorphism near $J$. Apply \cite[Proposition~3.12, p.~20]{PopescuPampu2016} to this morphism of smooth compact complex surfaces to factor it into point blowdowns. These maps give \eqref{cdp:eq:C.TI.13} and preserve the marked normal line.
\end{proof}

\subsection{Vertical fields and exceptional divisors}\label{subsec:II:vertical-field-divisors}

\begin{proposition}[Exceptional divisors in the component boundary]\hypertarget{tag-01ES}{}\label{cdp:st:C.2.UN.80}
\StatementRecord{cdp:st:C.2.UN.80}{proved}{01ES}{\cref{lem:II:vertical-elliptic-alignment,cdp:st:C.2.UN.29,lem:II:elliptic-first-integral,def:II:elliptic-root-curve,lem:II:identity-monodromy-elliptic}}
Let $f:U\to\Delta$ be a proper degeneration of connected complex two-tori from a smooth threefold, with full reduced normal-crossing central fibre $S$. Let $\nu:P\to S_i$ be a smooth normalization and put $D_P=(\nu^{-1}(\operatorname{Sing}S))_{\mathrm{red}}$, with both inverse branches of every self-double curve included. Let $\pi:P\to Z$ be a composition of point blowdowns and let $q_Z:Z\to B$ be a relatively minimal proper connected-fibre elliptic map to a smooth compact curve. Suppose every component of $D_P$ is vertical for $q=q_Z\circ\pi$.

There is a nonzero constant linear combination of the two saturated component fields transported to $P$ in \cref{lem:II:vertical-elliptic-alignment} which is the meromorphic lift of a nowhere-zero holomorphic vertical field
\begin{equation}
               W_Z\in H^0(Z,T_Z),\qquad dq_Z(W_Z)=0.
\hypertarget{tag-00I1}{}\label{cdp:eq:C.EP.1}
\end{equation}
Its complex-time period lattice is the same on every smooth elliptic fibre. It extends with neither a zero nor a pole across every smooth multiple fibre of $q_Z$.

Every prime exceptional for $\pi$ is a polar divisor of its meromorphic lift $W_P$. Consequently
\begin{equation}
                         \operatorname{Exc}(\pi)\subset D_P.
\hypertarget{tag-00I4}{}\label{cdp:eq:C.EP.2}
\end{equation}
In particular every last point-exceptional curve of this morphism is a boundary curve.
\end{proposition}
\begin{proof}
Apply \cref{lem:II:vertical-elliptic-alignment} to $P,q,D_P$. Every curve contracted by $\pi$ is vertical for $q$, so the given sequence $\pi$ is a choice of its relative contraction, with relatively minimal endpoint $q_Z$. The lemma gives identity integral monodromy, one fixed elliptic Jacobian and smooth multiple fibre supports on $Z$. Its degree-one component resolution $r:E\to P$ transports a complex basis $X,Y$ of the saturated fields and proves
\begin{equation}
                X,Y\ \hbox{holomorphic on }P\setminus D_P.
\hypertarget{tag-00I5}{}\label{cdp:eq:C.EP.3}
\end{equation}
They are complex-linearly independent and commute. Remove finitely many base values so that $\pi$ and $r$ are isomorphisms over the remaining connected $B^\circ$, the fibres are smooth, and the fields and their projected coefficients are regular. On every entire compact elliptic fibre these projected coefficients are constant. Thus $X,Y$ project to meromorphic vector fields on $B$.

\begin{enumerate}[label=\arabic*.,leftmargin=*]
\item Suppose their generic meromorphic rank is one. They both project to zero. Otherwise exchange $X,Y$ so that $X$ projects nontrivially. The two commuting projected fields would have a constant ratio $c$, while $Y-cX$ would be a nonzero vertical field and $X$ would be transverse; this would give rank two. Hence
\begin{equation}
                              Y=(q^*H)X
\hypertarget{tag-00I6}{}\label{cdp:eq:C.EP.4}
\end{equation}
for a nonconstant meromorphic function $H$ on $B$. The corresponding ratio on $E$ is $(q\circ r)^*H$. Apply the return-map calculation of \cref{cdp:st:C.2.UN.29} to this same embedded multiplicity-one component and its ambient saturated representatives. On $B^\circ$, the biholomorphisms $r_b$ pull the time form of $X$ back to the time form of its original restriction on $E$. The period identities \eqref{cdp:eq:C.RG.4}--\eqref{cdp:eq:C.RG.12} therefore give, for the same continued integral cycles,
\begin{equation}
             \lambda_i=d_i-k_iH\quad(i=1,2),\qquad
             \lambda_2/\lambda_1=\tau\in\mathbb H,
             \qquad d_i,k_i\in\C.
\hypertarget{tag-00I7}{}\label{cdp:eq:C.EP.5}
\end{equation}
The constants continue on $B^\circ$ by the identity theorem applied to the single-valued $H$. The fixed Jacobian and the continued integral marking make $\tau$ constant. Put
\begin{equation}
                         W=d_1X-k_1Y=\lambda_1X.
\hypertarget{tag-00I8}{}\label{cdp:eq:C.EP.6}
\end{equation}
A nonzero integral cycle has nonzero time period on an elliptic curve. Thus $W$ is nonzero on every fibre over $B^\circ$, and its two periods are exactly $1,\tau$.

\item Suppose the rank is two. The projected fields cannot both vanish. Two commuting nonzero meromorphic fields $a(z)\partial_z,b(z)\partial_z$ on a connected curve satisfy $ab'-ba'=0$, so have a constant ratio. After exchanging the basis there is a constant $c$ with
\begin{equation}
                    W=Y-cX\ne0,\qquad dq(W)=0,
                    \qquad dq(X)\ne0.
\hypertarget{tag-00I9}{}\label{cdp:eq:C.EP.7}
\end{equation}
On the compact fibres over $B^\circ$, $W$ is a nonzero translation field. The transverse flow of $X$ exists for uniform small time along each such fibre. Since $[X,W]=0$, it preserves the time form and its integral periods. This is the map \eqref{cdp:eq:C.RF.2} in the proof of \cref{lem:II:elliptic-first-integral}, with $(V,Y)=(W,X)$ and this same $q$. Its lattice is constant on connected $B^\circ$. A constant rescaling of $W$ makes one nonzero period equal to one.

\item Push $W$ through $\pi$ over $B^\circ$. On each small smooth fibre tube, integration with its fixed periods identifies the tube with a torsor under one fixed elliptic curve $F$; the transition maps have translation linear part and preserve its invariant field.

At a full fibre $mC$ apply
\cref{lem:II:identity-monodromy-elliptic} with
$(E_0,q,B)=(Z,q_Z,B)$. Use its fixed elliptic curve $F$
with the period marking of $W$. Its root quotient
\eqref{cdp:eq:C.EP.8} is written here with elliptic
coordinate $u'$, root base coordinate $z$, and
order-$m$ translation $\eta$. Thus its coordinate map
is $(z,s,\tau)=(u',z,\eta)$ and the
original base function is $t=z^m$. The inverse gauge
maps in that lemma identify the original root family
and this constant-translation quotient over the same
base disc; their derivative in the elliptic direction
is the identity.
The invariant vertical field is fixed by this free action and is nowhere zero. On the punctured root disc the coefficient of our field relative to it is constant, because its nonzero periods are fixed. Hence it extends across $z=0$ and descends through the finite \'etale quotient. The same construction with $m=1$ treats every omitted regular value. Agreement on the common open sets makes these local extensions a global field $W_Z$ satisfying \eqref{cdp:eq:C.EP.1}.

Its meromorphic lift through $\pi$ agrees with the constant combination \eqref{cdp:eq:C.EP.6} or \eqref{cdp:eq:C.EP.7} on the common open isomorphism, and hence everywhere as a meromorphic field. Its polar curves are therefore contained in $D_P$ by \eqref{cdp:eq:C.EP.3}.

\item Let $E$ be any prime exceptional for $\pi$, with image $p\in Z$. A holomorphic flow box $(z,w)$ for the nowhere-zero $W_Z$ at $p$ gives
\begin{equation}
                         W_Z(z)=1,\qquad W_Z(w)=0.
\hypertarget{tag-00IB}{}\label{cdp:eq:C.EP.9}
\end{equation}
At a generic smooth point of $E$, with local equation $s$, both pullbacks have positive order:
\begin{equation}
                  \pi^*w=s^bu,\qquad \pi^*z=s^av,
                  \qquad a,b\ge1,
\hypertarget{tag-00I2}{}\label{cdp:eq:C.EP.10}
\end{equation}
where $u,v$ are holomorphic units. Were $D=W_P$ holomorphic there, divide $D(\pi^*w)=0$ by $s^{b-1}$ and reduce modulo $s$ to obtain
\begin{equation}
                              b\,u\,D(s)=0\pmod{s}.
\hypertarget{tag-00I3}{}\label{cdp:eq:C.EP.11}
\end{equation}
Characteristic zero gives $D(s)\in(s)$. Applying $D$ to the second expression in \eqref{cdp:eq:C.EP.10} then gives $D(\pi^*z)\in(s)$, contrary to the pulled-back identity $D(\pi^*z)=1$. Thus every exceptional prime, including every infinitely-near point divisor, is polar. Together with \eqref{cdp:eq:C.EP.3}, this proves \eqref{cdp:eq:C.EP.2}.
\end{enumerate}
\end{proof}

\subsection{Surface models along a torsion-normal elliptic boundary}\label{subsec:II:torsion-interface-models}

\begin{proposition}[Surface models along a torsion-normal elliptic boundary]\hypertarget{tag-01EV}{}\label{cdp:st:C.2.UN.83}
\hypertarget{tag-01EH}{}\label{cdp:st:C.2.UN.70}
\StatementRecord{cdp:st:C.2.UN.83}{proved}{01EV}{\cref{cdp:st:C.2.UN.17,cdp:st:C.2.UN.67,cdp:st:C.2.UN.15,cdp:st:C.2.UN.19,cdp:st:C.2.UN.31,lem:II:section-preserving-ruling,lem:II:vertical-elliptic-alignment,cdp:st:C.2.UN.80,lem:II:elliptic-star-boundary,lem:II:hopf-single-canonical-curve,lem:II:zero-chern-elliptic-torsor}}
Let $f:U\to\Delta$ be a proper degeneration of connected complex two-tori from a smooth threefold, with full reduced normal-crossing central fibre $S$. Let $\nu:T\to S_i$ be a smooth normalization and let $D=(\nu^{-1}(\operatorname{Sing}S))_{\mathrm{red}}$, with both inverse branches of every self-double curve included. Suppose $J$ is a smooth elliptic component of $D$, disjoint from every other boundary branch, and $N_{J/T}$ has finite exact order $m$.

Write the induced meromorphic canonical section as
\begin{equation}
             K_T=\operatorname{div}\omega_T
                    =a_JJ+\sum_{C\ne J}a_CC.
\hypertarget{tag-013E}{}\label{cdp:eq:C.TI.5}
\hypertarget{tag-00T1}{}\SsixEquationAliases{\label{cdp:eq:C.NP.1}}
\end{equation}
The sum includes every other boundary branch, including those with coefficient zero. Then $T$ is nonrational, and it has one of the following marked surface models.

\begin{enumerate}[label=(\arabic*),leftmargin=*]
\item\label{item:II:torsion-interface-elliptic} There are maps
\begin{equation}
 T\xrightarrow{\pi}Z\xrightarrow{q_Z}B,
                    \qquad q=q_Z\circ\pi,
\hypertarget{tag-0135}{}\label{cdp:eq:C.TI.1}
\hypertarget{tag-00T2}{}\SsixEquationAliases{\label{cdp:eq:C.NP.2}}
\end{equation}
where $\pi$ is a finite composition of point blowdowns and is an isomorphism on a neighbourhood of $J$. The map $q_Z$ is a relatively minimal elliptic quasibundle with one fixed Jacobian, identity integral $H_1$ monodromy, only smooth multiple fibre supports, and $g(B)\le1$. The image $J_0=\pi(J)$ is one such support, with full multiplicity exactly $m$. Every original boundary and every exceptional curve of $\pi$ is vertical for $q$, and every exceptional prime is a boundary curve. Every rational boundary curve is contracted. The inverse point blowups of fibres of $q_Z$, with their centres, positions and normal maps, recover $T$ and its marked curves.

\item There is a composition of point blowdowns $\pi:T\to H$, isomorphic near $J$, onto a nonlinear minimal Hopf surface. Its unique curve is $J_0=\pi(J)$. For an integer $a\ge2$, the divisors of the corresponding meromorphic canonical sections satisfy
\begin{equation}
             K_H=-aJ_0,\qquad
             K_T=-aJ+\sum_Eb_EE,\qquad b_E>0.
\hypertarget{tag-013B}{}\label{cdp:eq:C.TI.2}
\end{equation}
The boundary away from $J$ is exactly $\operatorname{Exc}(\pi)$. If $k_T$ is its ambient canonical evaluation order, the neighbouring order along each such exceptional curve is $k_T+b_E+1\ge k_T+2$.

\item There is a composition of point blowdowns $\pi:T\to R$, isomorphic near $J$, onto a projective-line bundle over an elliptic curve. The curve $J_0=\pi(J)$ is a section, and
\begin{equation}
       K_R=-2J_0,\qquad N_{J_0/R}\simeq\OO_{J_0},\qquad
       K_T=-2J+\sum_Eb_EE,\qquad b_E>0.
\hypertarget{tag-013C}{}\label{cdp:eq:C.TI.3}
\end{equation}
Every other boundary is exactly one of these exceptional primes. After a line twist, $R=\PP(V)$ with its marked section given by
\begin{equation}
              0\longrightarrow\OO\longrightarrow V
                     \longrightarrow\OO\longrightarrow0.
\hypertarget{tag-013D}{}\label{cdp:eq:C.TI.4}
\end{equation}
The extension class is part of the model and may be nonzero. For the component across $J$, of order $k_P$, one has $k_T=k_P+1>0$. The point map records $T$ as successive blowups of this one-section ruled pair, with its other boundary curves the vertical exceptional configurations.
\end{enumerate}

In the ruled birational case, the third description requires that $J$ have degree one for the ruling and that no other horizontal boundary survive the contractions disjoint from $J$. A higher-degree elliptic multisection or a second surviving horizontal boundary gives the elliptic model \ref{item:II:torsion-interface-elliptic}.

A nonnegative coefficient $a_J$ in \eqref{cdp:eq:C.TI.5} forces the elliptic model \ref{item:II:torsion-interface-elliptic}, with all the conclusions just stated. In particular this applies to an embedded component of a full reduced SNC fibre with $a_J\ge0$, and to the parent of a Hopf leaf: its coefficient is $a_J=m-1\ge0$, where $m$ is the exact normal order of the leaf axis. Every model preserves the marked normal line and all boundary branches.
\end{proposition}
\begin{proof}
Choose the canonical evaluation generator of \cref{cdp:st:C.2.UN.17}, with divisor $\sum_i k_iS_i$, $k_i\ge0$. On a branch of $S_i$, the residue of $\Omega/f^{k_i+1}$ gives $\omega_T$. In a normal-crossing chart this is the residue of a unit times a monomial three-form. Along the inverse branch meeting $S_j$ its order is $k_j-k_i-1$; a self-double branch has $j=i$ and order $-1$. This specifies \eqref{cdp:eq:C.TI.5} and the same meromorphic section throughout each point-map comparison. The torsion normal gives $J^2=K_T\cdot J=0$.

\Cref{cdp:st:C.2.UN.67} applies to this smooth normalization and its complete inverse boundary, so excludes a rational $T$. The component restrictions \cref{cdp:st:C.2.UN.15,cdp:st:C.2.UN.19,cdp:st:C.2.UN.31} leave a nonrational ruled, Hopf, or nonnegative-Kodaira-dimension minimal model.

\begin{enumerate}[label=\arabic*.,leftmargin=*]
\item\textbf{Nonnegative Kodaira dimension.}
Let $\pi:T\to Z$ be the point-blowdown map to the minimal model. The component cover of \cref{cdp:st:C.2.UN.31} has degree one because the full multiplicity of $S_i$ is one. Its resolution maps birationally to $T$, and uniqueness of the nonruled minimal model identifies its $E_0$ with this same $Z$. Thus $Z$ is a two-torus, or it has the fixed-Jacobian identity-monodromy elliptic map $q_0:Z\to B_0$ of that proposition, with full fibres $m_iC_i$ and $g(B_0)\le1$.

For this latter map apply \cref{lem:II:identity-monodromy-elliptic}, with $(E_0,q,B)=(Z,q_0,B_0)$. Its canonical line \eqref{eq:II:quasibundle-canonical}, raised to any positive $N$ divisible by all $m_i$, and the Cartier equalities $q_0^*(b_i)=m_iC_i$ give
\begin{equation}
 K_Z^{\otimes N}
   =q_0^*\left(K_{B_0}^{\otimes N}\otimes
       \OO_{B_0}\left(\sum_i N(1-1/m_i)b_i\right)\right).
\hypertarget{tag-00R4}{}\label{cdp:eq:C.MR.12}
\end{equation}
Because $(q_0)_*\OO_Z=\OO_{B_0}$, its pluricanonical sections are exactly the sections of this base line. Consequently
\begin{equation}
                 \delta=2g(B_0)-2+\sum_i(1-1/m_i)\ge0,
\hypertarget{tag-00T3}{}\label{cdp:eq:C.NP.3}
\end{equation}
with $\delta>0$ in Kodaira dimension one and $\delta=0$ in dimension zero. Indeed a negative base degree has no sections, a degree-zero line has at most one, and positive degree has linear section growth by Riemann--Roch. Formula \eqref{cdp:eq:C.MR.12} makes $K_Z$ nonnegative on every curve: its degree on a horizontal curve is its base degree times $\delta$, and it is zero on vertical curves. On a torus $K_Z$ is trivial.

The Jacobian of each point blowup has exceptional order one. For the same meromorphic sections this gives
\begin{equation}
                      K_T=\pi^*K_Z+\sum_Eb_EE,
                      \qquad b_E>0.
\hypertarget{tag-013F}{}\label{cdp:eq:C.TI.6}
\end{equation}
Writing $J_0=\pi(J)$, adjunction yields
\begin{equation}
                 0=K_T\cdot J
                    =K_Z\cdot J_0+\sum_Eb_E(E\cdot J).
\hypertarget{tag-013G}{}\label{cdp:eq:C.TI.7}
\hypertarget{tag-00T4}{}\SsixEquationAliases{\label{cdp:eq:C.NP.4}}
\end{equation}
Every term is nonnegative, and $J$ is not an exceptional rational curve. Hence every exceptional curve misses $J$. The map $\pi$ is an isomorphism on its neighbourhood; $J_0$ is smooth elliptic, has square zero, and has the same normal line.

\item On a two-torus, a lift of an elliptic curve to the universal covers has periodic holomorphic derivative, hence constant derivative. Thus $J_0$ is a translate of an elliptic subgroup, and quotient by that subgroup gives a connected elliptic map with $J_0$ as a fibre. Every disjoint curve is vertical for that quotient.

For the elliptic alternative with $\delta>0$, the equality $K_Z\cdot J_0=0$ and \eqref{cdp:eq:C.MR.12} make $J_0$ vertical for $q_0$. Its smooth irreducible fibre support is exactly $J_0$, so every disjoint curve is vertical as well.

For $\delta=0$ and vertical $J_0$, keep $q_0$; the same fibre-support argument makes every disjoint curve vertical. It remains to adapt the map when $\delta=0$ and $J_0$ is horizontal. If $B_0$ is elliptic, there are no multiple fibres. The map $q_0$ is a principal elliptic torsor. The unramified finite map $J_0\to B_0$ gives its pulled-back torsor the section $j\mapsto(j,j)$. Its integral Chern vector therefore pulls back to zero. On the free group $H^2(B_0,\Lambda_F)$, this pullback multiplies by the positive covering degree, so the vector is zero. Apply \cref{lem:II:zero-chern-elliptic-torsor} with $(E_1,B,F_1)=(Z,B_0,F)$. Its explicit cocycle isomorphism identifies this same torsor with a two-torus, and the preceding subgroup quotient applies.

For $B_0=\PP^1$, the equality $\sum_i(1-1/m_i)=2$ gives
\begin{equation}
             (m_i)=(2,2,2,2),\ (3,3,3),\ (2,4,4),\ (2,3,6).
\hypertarget{tag-00R5}{}\label{cdp:eq:C.MR.13}
\end{equation}
Each summand is at least one half, so there are three or four terms. Four terms give the first list; ordering three denominators and solving $1/m_1+1/m_2+1/m_3=1$ gives the other three. Take the elliptic covering $E\to\PP^1$ with these ramification indices and these marked values: the double cover branched at the four given points in the first case, and the elliptic quotients of orders three, four and six in the other cases, as in \cref{cdp:st:C.2.UN.33}.

Normalize the pullback of $Z$ by this base cover. The local root maps of \cref{def:II:elliptic-root-curve} identify it with a finite \'etale Galois surface cover $Z'\to Z$, and $Z'\to E$ is a principal fixed-elliptic torsor. A connected component of the inverse image of $J_0$ is elliptic and maps finitely and unramifiedly to $E$. Its diagonal section after base pullback kills the Chern vector in the free group $H^2(E,\Lambda_F)$, by the same degree map. \Cref{lem:II:zero-chern-elliptic-torsor}, with $(E_1,B,F_1)=(Z',E,F)$, identifies $Z'$ with a two-torus through that torsor.

The inverse image of $J_0$ is a disjoint union of parallel elliptic translates in $Z'$. Distinct elliptic subgroup directions in a two-torus have a surjective product map with finite kernel, so their translates intersect; disjointness therefore gives one common subgroup $F'$. The deck group $G$ preserves it and acts on $Z'/F'$. The quotient map is
\begin{equation}
                       Z=Z'/G\longrightarrow (Z'/F')/G.
\hypertarget{tag-00R6}{}\label{cdp:eq:C.MR.14}
\end{equation}
It has connected elliptic fibres. The kernel on the base acts freely on a fibre, hence by translations: a finite elliptic automorphism with nonidentity linear part has a fixed point. The same argument for a base stabilizer gives smooth multiple fibre supports in the quotient. In particular $J_0$ is one such support. Every other noncontracted boundary image is disjoint from $J_0$, because $\pi$ is an isomorphism near it, and is therefore vertical for this map. Composing with $\pi$ makes all original boundaries and exceptional curves vertical.

\item\textbf{Hopf minimal models.}
Let $\pi:T\to H$ be the minimal point map. Every curve on $H$ is smooth elliptic with square zero by \cref{lem:II:hopf-curves}. The point-blowup square formula gives
\begin{equation}
                         0=J^2=J_0^2-\sum_\nu r_\nu^2.
\hypertarget{tag-013H}{}\label{cdp:eq:C.TI.8}
\hypertarget{tag-00T5}{}\SsixEquationAliases{\label{cdp:eq:C.NP.5}}
\end{equation}
Thus every centre misses the successive image of $J$, and its entire neighbourhood and torsion normal are preserved. The diagonal algebraic-dimension-zero axes have nontorsion normals by the character comparison of \cref{cdp:st:C.2.UN.38,cdp:st:C.2.UN.58}. An elliptic Hopf minimal model has the elliptic map of \cref{cdp:st:C.2.UN.37}; all its curves are fibre supports, so its composite with $\pi$ has every boundary vertical.

For a nonlinear $H$, its unique curve is $J_0$. Push the meromorphic section in \eqref{cdp:eq:C.TI.5} through the common open isomorphism. Its divisor is $a_JJ_0$, so \cref{lem:II:hopf-single-canonical-curve}, applied with $(H,J,a)=(H,J_0,-a_J)$, gives $-a_J=a=r+1\ge2$ with the precise primary or secondary pair. Its unramified primary map preserves the divisor coefficient, and its invariant canonical form fixes the same meromorphic section up to a constant. Formula \eqref{cdp:eq:C.TI.6} now gives \eqref{cdp:eq:C.TI.2}. Every exceptional prime has its positive discrepancy as coefficient, so it is a boundary curve by \eqref{cdp:eq:C.TI.5}; conversely every other boundary is contracted because $J_0$ is the only curve on $H$. The residue coefficient formula at the start of the proof gives neighbouring order $k_T+b_E+1\ge k_T+2$.

\item\textbf{Ruled models and their marked covers.}
Apply \cref{lem:II:section-preserving-ruling} with $P=T$ and the same $J$. It gives the elliptic ruling $p:T\to B_0$, the isogeny $J\to B_0$ of degree $d$, its finite \'etale base change $T'=T\times_{B_0}J$, and the equivariant contraction $\pi':T'\to R'$ preserving all lifted sections. Its deck group is $G=\operatorname{Deck}(J/B_0)$. Its quotient is $\pi_0:T\to R$, isomorphic near $J$. The equality $R'=R\times_{B_0}J$ is the finite descent map $r'\mapsto([r'],p'(r'))$. Every other surviving boundary is disjoint from $J_0=\pi_0(J)$. Since a ruling fibre meets $J_0$, such a curve is horizontal.

On $R'\to J$, choose the diagonal section $J_1'$ and write $L=N_{J_1'/R'}$. Its order is $m$, by the neighbourhood isomorphism. For $W=p'_*\OO_{R'}(J_1')$, fibrewise evaluation identifies $R'$ with $\PP_J(W^*)$. Pushing forward the section sequence gives $0\to\OO_J\to W\to L\to0$; its last surjection follows from $R^1p'_*\OO_{R'}=0$, since $H^1(\PP^1,\OO)=0$. Dualizing and tensoring by $L$ yields
\begin{equation}
       0\longrightarrow\OO_J\longrightarrow V
             \longrightarrow L\longrightarrow0,
       \qquad R'=\PP_J(V),\qquad J_1'=\PP_J(\OO_J).
\hypertarget{tag-00R2}{}\label{cdp:eq:C.MR.10}
\hypertarget{tag-00T8}{}\SsixEquationAliases{\label{cdp:eq:C.NP.8}}
\end{equation}
The marked subline is the dual of evaluation on $J_1'$, tensored by $L$, so this isomorphism includes the section and its normal.

If $d\ge2$, another lifted section is a compact curve in its complement. If $d=1$ and another boundary survives on $R$, its pulled-back normalization gives such a compact curve mapping finitely to $J$. In either situation the extension splits. For $L\ne\OO_J$, its extension group is $H^1(J,L^{-1})=H^0(J,L)^*=0$. For $L=\OO_J$, the complement of the section is the additive torsor of a class $\xi\in H^1(J,\OO_J)$. A compact curve $A$ in that complement gives a finite map $a:A\to J$ and a section of the pulled-back torsor, so $a^*\xi=0$. The finite trace satisfies $\operatorname{tr}\circ a^*=(\deg a)\operatorname{id}$, including ramified maps, and hence $\xi=0$.

Kill $L$ by a further finite \'etale elliptic base cover $E'\to J$. Choose it Galois over $B_0$, which is possible since finite unramified coverings of elliptic curves are isogenies after choices of origins. The pulled-back surface is $E'\times\PP^1$. Every deck map has form $(b,z)\mapsto(b+\tau,A(b)z)$, and $A:E'\to\operatorname{PGL}_2(\C)$ is constant because its source is compact and its target affine. Thus
\begin{equation}
                     R=(E'\times\PP^1)/G'
                         \longrightarrow\PP^1/G'.
\hypertarget{tag-00R3}{}\label{cdp:eq:C.MR.11}
\hypertarget{tag-00T9}{}\SsixEquationAliases{\label{cdp:eq:C.NP.9}}
\end{equation}
If $K$ is the kernel of the projective action, a regular connected fibre is $E'/K$. Every fibre action is an elliptic translation. At a projective stabilizer the action on $E'/K$ is free of the same order, giving the full smooth-multiple chart \eqref{cdp:eq:C.EP.8}. This proves fixed Jacobian and identity integral monodromy for this map.

The chosen $J_1'$ is the constant section at infinity in this marked product. Every other lifted $J_i'$ misses it, so its projective coordinate is a holomorphic function into $\C$ on the compact $E'$ and is constant. A lifted boundary disjoint from these sections likewise has projective coordinate in their affine complement; that coordinate is holomorphic on its compact normalization and hence constant. Thus every surviving boundary is a fibre support of \eqref{cdp:eq:C.MR.11}, and every contracted boundary is vertical for its composite with $\pi_0$. When $d=1$, the deck maps preserve the section at infinity and have constant affine form $z\mapsto a_gz+b_g$. Their common fixed point is the average $z_0=|G'|^{-1}\sum_g g(0)$. The coordinate $z'=z-z_0$, with inverse $z=z'+z_0$, gives the same quotient while fixing the marked section.

\item\textbf{The multisection pencil and the quotient map.}
For $d\ge2$, this construction also identifies the pencil on the original cover $T'$. Two lifted sections are numerically equivalent by the Hodge-index comparison in \cref{lem:II:section-preserving-ruling}. Their difference is zero in the rational N\'eron--Severi group: its class is orthogonal to an ample class and has square zero, so Hodge index makes it zero in real cohomology. The integral cohomology of a ruled surface and its point blowups is torsion-free by the projective-bundle formula and \cref{lem:II:blowup-topology}, giving zero integral first Chern class as well. The projective-bundle direct images and the point-blowup direct images proved in \cref{lem:II:vertical-elliptic-alignment}, together with the exponential sequence, give the pullback isomorphism
\begin{equation}
          (p')^*:\operatorname{Pic}^0(J)\xrightarrow{\sim}
                         \operatorname{Pic}^0(T').
\hypertarget{tag-0137}{}\label{cdp:eq:C.TI.11}
\end{equation}
Hence $\OO_{T'}(J_1'-J_2')=(p')^*L$; restriction to the diagonal section identifies this $L$ with its normal line. The disjoint divisors $mJ_1',mJ_2'$ are linearly equivalent and define
\begin{equation}
                              q':T'\longrightarrow\PP^1.
\hypertarget{tag-0138}{}\label{cdp:eq:C.TI.12}
\end{equation}
To compare the two maps, put the two constant sections at $0,\infty$ in the product above $J$. The deck subgroup $G_L=\operatorname{Deck}(E'/J)$ acts on this coordinate by the normal character of $L$, whose image has exact order $m$. Its quotient coordinate is $z^m$. The defining ratio of the sections of $\OO(mJ_1')$ pulls back, after multiplying one section by a nonzero constant, to exactly $z^m$. Thus $q'$ is the product-quotient map composed with $\pi'$ on this intermediate cover. Its regular fibre is the connected elliptic quotient of $E'$ by the kernel of that character. The remaining group $G'/G_L=G$ acts on this quotient base, and descent to $T$ gives precisely \eqref{cdp:eq:C.MR.11} composed with $\pi_0$. Every $J_i'$ is the whole reduced support of a fibre: its square is zero, so another component of the same full fibre would be disjoint from it, contrary to connectedness. These identifications preserve the original pencil, its connected fibres, the full normal order and the descending group action.

\item\textbf{The one-section ruled model.}
The remaining ruled case has $d=1$ and no other surviving boundary on $R$. The pushed meromorphic section in \eqref{cdp:eq:C.TI.5} has divisor supported on $J_0$. Intersecting a ruling fibre gives its coefficient $-2$, so $K_R=-2J_0$ as its divisor. Adjunction gives
\begin{equation}
              \OO_{J_0}=K_{J_0}
                     =(K_R+J_0)|_{J_0}=N_{J_0/R}^{-1}.
\hypertarget{tag-013A}{}\label{cdp:eq:C.TI.14}
\end{equation}
Thus $L=\OO$, and \eqref{cdp:eq:C.MR.10} is exactly \eqref{cdp:eq:C.TI.4}, with its extension class unchanged. Applying the point-map Jacobian formula to this same section gives \eqref{cdp:eq:C.TI.3}. Every exceptional prime has positive coefficient and is a boundary; every other boundary is exceptional by construction. The coefficient of $J$ is $k_P-k_T-1=-2$, so $k_T=k_P+1>0$. This gives the third marked model, including all vertical exceptional configurations and their inverse point maps.

\item\textbf{Alignment of the elliptic models.}
In every elliptic construction above, $\pi:T\to Z$ is an isomorphism near $J$ and every boundary is vertical for the connected map $q=q_Z\pi$. The endpoint $Z$ is relatively minimal: it is a minimal nonruled surface, a Hopf surface with no rational curves, or a projective-line bundle over an elliptic base, whose rational curves are square-zero ruling fibres. Apply \cref{lem:II:vertical-elliptic-alignment} to this precise $T,q,D$, using the displayed point map $\pi$ for its relative contraction. It gives identity integral monodromy, smooth multiple fibre supports, $R^1(q_Z)_*\OO_Z=\OO_B$ and $g(B)=p_g(Z)=p_g(T)\le1$. Apply \cref{cdp:st:C.2.UN.80} with $(P,D_P,\pi,q_Z)=(T,D,\pi,q_Z)$. Its constant-combination and polar-prime conclusions give $\operatorname{Exc}(\pi)\subset D$ for this same map.

For the full fibre $nJ_0$, apply
\cref{lem:II:identity-monodromy-elliptic} to this same
$(Z,q_Z,B)$ and $C=J_0$. Its connected root cover and
primitive conormal character give exact normal order
$n$ on $J_0$. The neighbourhood isomorphism $\pi$ pulls
that line back to $N_{J/T}$, whose order is $m$;
therefore $n=m$. A rational boundary cannot map nonconstantly to a smooth elliptic fibre support, and therefore is contracted. The inverse centres and charts of $\pi$ give all the claimed vertical point-blowup data.

Finally, the nonlinear and one-section descriptions have coefficients $a_J=-a\le-2$ and $a_J=-2$ respectively. Thus $a_J\ge0$ gives the elliptic description. In the ruled case, if only $J_0$ survived, its canonical divisor would be $a_JJ_0$, and intersection with a ruling fibre would give $-2=a_Jd$, contradicting that sign. Hence the product construction also has a second surviving boundary in this nonnegative-coefficient case.

For a Hopf leaf $H$ with boundary $J$, adjunction gives $K_H=-(k_H-k_P+1)J$. Apply \cref{lem:II:hopf-single-canonical-curve} to this pair. Its parameter is $r=k_H-k_P\ge1$, and the normal-order calculation of \cref{lem:II:elliptic-star-boundary} gives $m=r$, including the finite secondary action. The coefficient on the parent is therefore $k_H-k_P-1=m-1\ge0$. The inverse normal product along $xy=t$ gives the same order on the parent normalization. This proves the stated Hopf-parent specialization and completes all three models.
\end{enumerate}
\end{proof}

\subsection{Contractions to elliptic cores}\label{subsec:II:elliptic-core-reduction}

\begin{proposition}[Leaf contractions, elliptic cores and Euler number]\hypertarget{tag-01EC}{}\label{cdp:st:C.2.UN.66}
\StatementRecord{cdp:st:C.2.UN.66}{proved}{01EC}{\cref{cdp:st:C.2.UN.17,cdp:st:C.2.UN.64,cdp:st:C.2.UN.65,lem:II:vertical-elliptic-alignment,lem:II:zero-chern-elliptic-torsor,cdp:st:C.2.UN.19,cdp:st:C.2.UN.31,cdp:st:C.2.UN.28,cdp:st:C.2.UN.47,cdp:st:C.2.UN.44,cdp:st:C.2.UN.58,lem:II:reduced-anticanonical-hopf,lem:II:hopf-curves,cdp:st:C.2.UN.50,cdp:st:C.2.UN.53,cdp:st:C.2.UN.54,cdp:st:C.2.UN.61,cdp:st:C.2.UN.62}}

Let $f:U\to\Delta$ be a proper degeneration of connected complex
two-tori from a smooth threefold. Suppose that its full central
fibre $S=\sum_vT_v$ is reduced and SNC, has no triple points,
and that every component is either a smooth minimal compact
complex surface or a holomorphic $\PP^1$-bundle over a smooth
compact curve. The latter convention includes $\mathbb F_1$.

There is a finite sequence of the ordinary holomorphic
contractions of \ref{cdp:st:C.2.UN.64}, each along a terminal component, which
retains the punctured torus family and ends in one of the
following marked supports:
\begin{enumerate}
\item a smooth two-torus;
\item a marked diagonal Hopf and degree-zero elliptic ruled
cycle of \ref{cdp:st:C.2.UN.65}, with its dual normal maps and
complement-lattice product $M=I$;
\item a rooted elliptic tree passing the complete test of
\ref{cdp:st:C.2.UN.61}, including its zero-Chern inverse torsor;
\item a unicyclic genus-zero elliptic quasi-bundle support
passing the complete test of \ref{cdp:st:C.2.UN.62}.
\end{enumerate}
Every original surface and neighbourhood is recovered by
reversing the ordinary blowups. In particular the
normal extension at every curve centre is retained.

No rational component occurs in $S$. Let $\mathcal R_{\ge2}$
be the set of its ruled components whose ruling bases $B_v$
have genus at least two. Then
\begin{equation}
 e(S)=-2\sum_{v\in\mathcal R_{\ge2}}\bigl(g(B_v)-1\bigr)
                                                        \le0.
\hypertarget{tag-00R1}{}\label{cdp:eq:C.MR.1}
\end{equation}
Consequently $e(S)=0$ exactly when $\mathcal R_{\ge2}$ is
empty. In that case all the contractions in the preceding
reduction have elliptic curve centres. Their exceptional
marked pairs are precisely
\begin{equation}
 \left(\PP_B(V),\PP_B(\OO_B)\right),
 \qquad
 0\longrightarrow\OO_B\longrightarrow V
             \longrightarrow\OO_B\longrightarrow0,
 \qquad g(B)=1,
\hypertarget{tag-00R9}{}\label{cdp:eq:C.MR.2}
\end{equation}
with their split or nonzero extension and endpoint
identifications.

Conversely, start with a torus degeneration having
one of the displayed cores and perform any finite sequence of
ordinary blowups along smooth compact elliptic curves in the
smooth locus of its full reduced central fibre, avoiding its
double curves at each stage. The result is again a full
reduced SNC torus fibre without triple points, its components
are in the stated minimal-or-bundle class, and its Euler
number is zero. The reverse maps reconstruct the chosen neighbourhood together with each normal extension.

For every full reduced SNC torus fibre whose connected double curves are all elliptic, $e(S)\ge0$.
Its Euler number is zero exactly when it has no triple
points and every component is minimal and nonrational.
Every such zero-Euler fibre is therefore covered by the
preceding reduction and exact reconstruction.
\end{proposition}

\begin{proof}\leavevmode\par

 
\begin{enumerate}[label=\arabic*.,leftmargin=0pt,labelsep=.5em]
\item For the connected reducible case, apply \cref{cdp:st:C.2.UN.17} to this same family. Every multiplicity is one; the smooth components have identity normalization and resolution and intrinsic conductor zero. Write $k_v$ for its canonical evaluation orders. The principal-fibre and adjunction maps \eqref{cdp:eq:C.CA.3}--\eqref{cdp:eq:C.CA.4} give the meromorphic canonical section with divisor
\begin{equation}
 K_{T_v}=
       \sum_{e=(v,w)}(k_w-k_v-1)C_e.
\hypertarget{tag-00RA}{}\label{cdp:eq:C.MR.3}
\end{equation}
The notation denotes this section and its line isomorphism.
The curves in the sum are pairwise disjoint because there
are no triple points. Along each edge the equation
$xy=t$ gives
\begin{equation}
                 N_{C_e/T_v}\otimes N_{C_e/T_w}
                                        \simeq\OO_{C_e}.
\hypertarget{tag-00RB}{}\label{cdp:eq:C.MR.4}
\end{equation}

\begin{enumerate}[label=(\arabic*),leftmargin=1.3em,labelsep=.5em]
\item  Apply \ref{cdp:st:C.2.UN.64} whenever a terminal ruled component meets its
parent in a section, or a terminal plane meets its parent
in a line. Each is the ordinary blowdown of the
specified smooth curve or point, and the new family is
proper, smooth-total, reduced and SNC. Its other components
and marked curves are unchanged except for the parent
comparison proved there.

\item  This procedure preserves the surface convention in the
statement. At a curve contraction the parent surface is
unchanged. At a plane-line contraction the parent contains
the boundary $(-1)$-curve. A minimal surface has none.
Among the allowed projective-line bundles, a rational
curve on a bundle over a positive-genus base is a ruling
fibre and has square zero. On $\mathbb F_n$ the only
$(-1)$-curve occurs on $\mathbb F_1$ and is its negative
section: an irreducible curve different from the negative
section has class $aC_0+bf$ with $b\ge na$ when $a>0$,
and hence nonnegative square; the fibres have square zero.
Thus that parent is $\mathbb F_1$, and its specified
surface blowdown is $\PP^2$. It remains in the allowed
class. Every other boundary was disjoint from the
contracted curve, so no new intersection or triple point
is introduced.

\item  Each operation removes one component, so it terminates.
The canonical order differences on surviving components
are unchanged: the ordinary curve and point blowup
Jacobians have exceptional orders one and two,
respectively, as in \eqref{cdp:eq:C.EL.1} and \eqref{cdp:eq:C.EL.4}. The remaining
family therefore still has \eqref{cdp:eq:C.MR.3}, and no eligible
\ref{cdp:st:C.2.UN.64} leaf. We call its support $S_0$.
For one smooth reduced component, \ref{cdp:st:C.2.UN.5} already identifies
the fibre as a two-torus. Assume below that $S_0$ is
reducible.

\end{enumerate}

\item At a component $T$ of maximum canonical order, put
$a_e=k_T-k_w+1\ge1$. Equation \eqref{cdp:eq:C.MR.3} reads
\begin{equation}
                         -K_T=\sum_e a_eC_e.
\hypertarget{tag-00RC}{}\label{cdp:eq:C.MR.5}
\end{equation}
Its boundary is nonempty. If a positive pluricanonical
power had a nonzero section, its product with the
corresponding power of the nonzero section in \eqref{cdp:eq:C.MR.5}
would be a nonzero holomorphic function with a nonempty
zero divisor on the compact connected surface. This is
impossible. Thus $\kappa(T)=-\infty$. The surface
classification and the class-VII restriction \ref{cdp:st:C.2.UN.19}
leave a projective-line bundle, $\PP^2$, or a Hopf
surface.

\begin{enumerate}[label=(\arabic*),leftmargin=1.3em,labelsep=.5em]
\item  For a ruled $T\to B$, intersect \eqref{cdp:eq:C.MR.5} with a general
ruling fibre. No boundary curve is vertical: a vertical
curve would be a ruling fibre and would meet every
horizontal boundary, whereas some horizontal boundary
is required by the following equality. Hence
\begin{equation}
              2=\sum_e a_ed_e,\qquad
              d_e=C_e\cdot(\text{ruling fibre})>0.
\hypertarget{tag-00RD}{}\label{cdp:eq:C.MR.6}
\end{equation}
The integer partitions are exactly one section of
coefficient two, one bisection of coefficient one,
or two sections of coefficient one each. The first
is an eligible contraction. The terminal-pair
classification in \ref{cdp:st:C.2.UN.64} makes the second a rational
ruled surface with its smooth anticanonical elliptic
bisection; its neighbour has the same canonical order.
In the third case adjunction along either disjoint
section gives
\[
       2g(B)-2=(K_T+C_e)\cdot C_e=0.
\]
Thus $B$ is elliptic, and both neighbours again have
the same order as $T$. On $\PP^2$ disjoint effective
boundary curves can have only one component, by
B\'ezout. The two possibilities are the eligible
plane-line contraction or the anticanonical cubic
with an equal-order neighbour.

\item  For a Hopf surface of algebraic dimension one use
its elliptic map over $\PP^1$, the connected
normal-root chart and the invariant relative
differential of \ref{cdp:st:C.2.UN.58}. All its curves are fibre
supports. The canonical order formula gives
\begin{equation}
 d_{T,e}:=\frac{k_w-k_T}{m_e}\in\Z_{\le0},
 \qquad
              \sum_{e\ni T}d_{T,e}
                                  =\operatorname{val}(T)-2.
\hypertarget{tag-00RE}{}\label{cdp:eq:C.MR.7}
\end{equation}
An unmarked multiple fibre is excluded by the same
root formula $m\operatorname{ord}\theta+m-1=0$.
Thus the valence is one or two. Valence one gives
$k_T-k_w=m_e>0$; valence two gives two equal-order
neighbours. A nonlinear Hopf surface has its single
elliptic axis, so is a leaf; its canonical
axis form gives the positive drop equal to its
resonance, including the secondary quotient in \ref{cdp:st:C.2.UN.50}.

\item For a diagonal algebraic-dimension-zero Hopf surface,
apply \cref{lem:II:reduced-anticanonical-hopf} to the same
pair $(T,D)$, where $D=\sum_ea_eC_e$ from
\eqref{cdp:eq:C.MR.5}. The support is disjoint and elliptic
by \cref{lem:II:hopf-curves}, and $K_T(D)\simeq\OO_T$ by
that displayed canonical section.
Write $T=W/\Gamma$ with its diagonal deck group and take
the lemma's central primary quotient map
$W/\langle F\rangle\to W/\Gamma=T$, induced by the identity
on $W$. The nonreduced alternative would have
contraction $(z,w)\mapsto(\beta^r z,\beta w)$ and the
nonconstant meromorphic function $h=z/w^r$ on this cover.
For its finite deck group $G$, the coefficients of
$\prod_{g\in G}(u-g^*h)$ descend to meromorphic functions
on $T$. At least one is nonconstant: constant coefficients
would force $h$ to take values in a fixed finite set.
This contradicts algebraic dimension zero. Hence the
lemma's reduced alternative applies, with its two
individually marked axes and $a_1=a_2=1$. Both neighbouring
orders equal $k_T$, by \eqref{cdp:eq:C.MR.5}.

\item  It follows that every noneligible maximum component
which is not a Hopf leaf has all of its neighbours
at that same maximum. An equal-order edge cannot
lead to one of the strict-drop leaves just listed.
Following these edges therefore gives an entire
connected component of the whole incidence graph.
Connectedness of $S_0$ forces all orders to be equal.
In that event \ref{cdp:st:C.2.UN.65} applies; since there are no triple
points, its complete zero-case classification is
exactly the displayed cycle or smooth torus.
If the orders are not all equal, a maximum Hopf
leaf exists. Its boundary is elliptic and its normal
line has finite exact order.

\end{enumerate}

\item Whenever $q:T\to B$ is a connected elliptic map with every boundary curve vertical, apply \cref{lem:II:vertical-elliptic-alignment} to this same component of the remaining full reduced SNC torus family. It gives identity integral monodromy and \eqref{cdp:eq:C.MR.8} for this $q,T,B$. Its relative contraction removes only vertical $(-1)$-curves. On a minimal $T$ there are none. On a projective-line bundle over an elliptic curve, every rational curve is a ruling fibre and has square zero, so there are none there either. For these components the relative contraction is the identity, and the lemma's full smooth-multiple conclusion applies to $q$ itself.

\item Consider a component $T$ meeting an elliptic boundary $C$ with torsion normal. Apply \cref{cdp:st:C.2.UN.83} to the same full reduced SNC family with identity normalization, $J=C$, and its unchanged boundary. The curve is isolated because there are no triple points. That proposition excludes rational $T$ and supplies its three marked surface models.

Every point map in this application is the identity. A minimal $T$ has no $(-1)$-curve. For an allowed projective-line bundle over a positive-genus base, every rational curve is a ruling fibre and has square zero; the rational bundles were just excluded. Thus there is no first curve to contract. The three models consequently describe the given component and its marked curves themselves.

The one-section ruled description would make $T$ an eligible terminal pair for \cref{cdp:st:C.2.UN.64}, already removed in forming $S_0$. A nonlinear Hopf description is its torsion-normal leaf. In every other case the elliptic map \eqref{cdp:eq:C.TI.1} is defined on this same $T$, has every boundary vertical, and has the full identity-monodromy, fixed-Jacobian and smooth-multiple conclusions \eqref{cdp:eq:C.MR.8}.

For a ruled component this is the quotient \eqref{cdp:eq:C.MR.11}: its first cover is $T'=T\times_E C$, its section normal is the same $L$ in \eqref{cdp:eq:C.MR.10}, and its product cover kills that $L$. For nonnegative Kodaira dimension the degree-one component cover identifies its minimal model with $T$, and \eqref{cdp:eq:C.MR.12}--\eqref{cdp:eq:C.MR.14} give the marked quotient through the same principal torsor and Chern-class pullback. These are the ruled and nonruled maps constructed in \cref{cdp:st:C.2.UN.83}; their normal-order identifications carry each boundary to its full fibre support.

\item Start with the maximum Hopf
leaf when the canonical orders
are not constant. Its boundary
has finite normal order, and
\eqref{cdp:eq:C.MR.4} gives the same property
on its neighbour. The preceding
comparisons put that neighbour
in the fixed-Jacobian elliptic
class, with all its boundaries
vertical, unless it is itself
a nonlinear leaf. Every vertical
boundary again has finite exact
normal order by \eqref{cdp:eq:C.QT.2}.
Repeat across the connected
incidence graph. No rational
or higher-genus ruled component
can occur in this propagation,
by \eqref{cdp:eq:C.MR.9} and
Riemann--Hurwitz. A one-boundary
elliptic ruled surface would be
an eligible \ref{cdp:st:C.2.UN.64} section leaf
and is absent from $S_0$.
Consequently the entire $S_0$
consists of the displayed
elliptic quasi-bundles and
nonlinear Hopf leaves.

\begin{enumerate}[label=(\arabic*),leftmargin=1.3em,labelsep=.5em]
\item  Equation \eqref{cdp:eq:C.MR.8} and the
simultaneous bound \ref{cdp:st:C.2.UN.28} give
at most one genus-one base.
If there is one, its component
has $p_g=1$. The component-group and dual-complex
comparison \ref{cdp:st:C.2.UN.47} makes the
incidence graph simply connected.
It is therefore a tree. The
entire support is in the exact
scope of \cref{cdp:st:C.2.UN.61}
and satisfies all its necessary
tests, including its integral
inverse-torsor obstruction.
If there is no genus-one base,
all vertices have coherent
dimensions $(1,1,0)$ and the
nonzero boundary restrictions
of \eqref{cdp:eq:C.QT.3} or \ref{cdp:st:C.2.UN.54}.
The graph quotient of
\ref{cdp:st:C.2.UN.44} has at most one cycle.
Leaf elimination would give
$H^2(\OO_{S_0})=0$ on a tree,
contradicting upper
semicontinuity from a torus.
Thus it is unicyclic and the
whole support is in the exact
scope of \ref{cdp:st:C.2.UN.62}, with all of
its necessary tests satisfied.
This proves that the
contraction procedure reaches
one of the stated cores.
The intermediate families are the ordinary blowdowns of \cref{cdp:st:C.2.UN.64}. Apply \cref{cdp:st:C.2.UN.61,cdp:st:C.2.UN.62} to their final family, with the same components, connected elliptic maps, boundary fibres, normal lines and analytic torsors just identified.

\end{enumerate}

\item We now use the proved endpoint
description to exclude rational
components in the original
support. A nonrational
component in the allowed class
has no rational boundary.
For a Hopf component, apply \cref{lem:II:hopf-curves}
to that same minimal surface: every curve is smooth
elliptic. On a ruled surface
over a positive-genus base a
rational curve is a ruling
fibre. If it were a boundary,
disjointness would leave no
horizontal boundary. A general
ruling fibre would then have
intersection zero with the
right side of \eqref{cdp:eq:C.MR.3}, while
its canonical degree is $-2$.
This is impossible.

\begin{enumerate}[label=(\arabic*),leftmargin=1.3em,labelsep=.5em]
\item  For a minimal component with
$\kappa\ge0$, use its degree-one
\ref{cdp:st:C.2.UN.31} description. A torus has
no rational curve by its
universal covering. In the
elliptic alternative \eqref{cdp:eq:C.MR.12}
has $\delta\ge0$. A rational
curve cannot be vertical in
a smooth elliptic support.
If its map to the base had
degree $d>0$, each local
ramification order over $b_i$
would be divisible by $m_i$,
from the equation
$t=u x^{m_i}$. Riemann--Hurwitz
would therefore give
\begin{equation}
 -2\ \ge\
 d\left(2g(B_0)-2+
              \sum_i(1-1/m_i)\right)
                              =d\delta\ge0,
\hypertarget{tag-00R7}{}\label{cdp:eq:C.MR.15}
\end{equation}
a contradiction. This proves
the asserted absence of
rational boundaries at every
stage of the reduction.

\item  Whenever a rational component
is removed by \ref{cdp:st:C.2.UN.64}, its
section or line boundary is
rational. Its parent must
therefore be rational too.
The operation cannot remove
the last rational component.
Every final core, however,
has no rational component:
this is explicit in \ref{cdp:st:C.2.UN.65},
and in \cref{cdp:st:C.2.UN.61,cdp:st:C.2.UN.62} every
component has coherent
Euler zero. Thus no rational
component existed initially.
In particular no plane-line
contraction occurred.

\end{enumerate}

\item All contractions are now
curve contractions, with
base genus at least one.
All original ruled components
of base genus at least two
must be contracted, since
none occurs in the final
cores. A curve contraction
does not modify its parent
surface, and there are no
point contractions. Hence
these are exactly the
original higher-genus ruled
components, each counted once
with its original ruling base.
Every other contracted
component has elliptic base.
The exact Euler comparison
\eqref{cdp:eq:C.EL.5}, and the zero Euler
number of each final core,
give precisely \eqref{cdp:eq:C.MR.1}.
This also proves the stated
zero criterion independently
of any choice of contraction
order.

\begin{enumerate}[label=(\arabic*),leftmargin=1.3em,labelsep=.5em]
\item  In the zero case every
centre is elliptic.
Equation \eqref{cdp:eq:C.EL.3} then
gives \eqref{cdp:eq:C.MR.2}, including
its full normal extension.
The ordinary
blowups and unchanged
exterior identifications
recover the original
neighbourhood exactly.

\item  For the converse construction,
the displayed cores are in
the stated surface class.
For the quasi-bundle vertices
their smooth-multiple
fibration gives $e=\chi(\OO)=0$.
In nonnegative Kodaira dimension
the argument \eqref{cdp:eq:C.MR.15} excludes
exceptional rational curves.
In Kodaira dimension $-\infty$,
the surface classification,
$\chi(\OO)=0$ and $e=0$
give an unblown elliptic ruled
surface or a Hopf surface.
The root and the canonical-flat
cores are already minimal by
\ref{cdp:st:C.2.UN.61} and \ref{cdp:st:C.2.UN.65}.

\item For each chosen elliptic centre $B$, apply the
ordinary inverse curve map of \cref{cdp:st:C.2.UN.64}, with
its $B$ equal to this centre and its parent $P'$ equal
to the unique central component containing $B$.
The normal sequence is \eqref{cdp:eq:C.EL.3}; a nonzero
holomorphic differential on $B$ identifies $K_B$ with
$\OO_B$ and carries that sequence to \eqref{cdp:eq:C.MR.2},
with the corresponding endpoint identifications.
The map on the strict parent is the blowup of the Cartier
curve $B$, hence an isomorphism. The exceptional component
is its marked projective normal bundle. The centre avoids
all other boundary curves, so the same local map has full
reduced SNC central fibre and no triple points.
Its properness and punctured-family isomorphism are those
of the ordinary blowup. Its Euler change in
\eqref{cdp:eq:C.EL.5} is $2-2g(B)=0$. Iterating these same
maps proves the converse, with every normal extension
and neighbourhood identification preserved.

\end{enumerate}

\item Finally drop the component-minimality and no-triple
assumptions, retaining a full reduced SNC torus fibre
whose connected double curves are all elliptic.
The single-component case is the smooth two-torus
case of \ref{cdp:st:C.2.UN.5}. In the reducible case every component
has a nonempty boundary, and \eqref{cdp:eq:C.MR.3} remains the
adjunction formula even at triple points.

\begin{enumerate}[label=(\arabic*),leftmargin=1.3em,labelsep=.5em]
\item  No component can be birationally ruled over a
curve of genus at least two. Its map through
the contraction to that minimal ruled
surface and then to the base sends each
elliptic boundary to a point, by
Riemann--Hurwitz. A general ruling fibre
avoids those finitely many base values and
all point-blowup centres. Its strict
transform is a rational curve of
canonical degree $-2$ and has zero
intersection with every term of
\eqref{cdp:eq:C.MR.3}, a contradiction.

\item  Write $Z_v$ for the smooth minimal
model of $T_v$, and $n_v$ for the number
of point blowups in that contraction.
If its Kodaira dimension is nonnegative,
the degree-one comparison of \ref{cdp:st:C.2.UN.31} gives
a torus or the smooth-multiple elliptic
minimal model, so $e(Z_v)=0$. If it has
Kodaira dimension $-\infty$, the preceding
paragraph leaves rational, elliptic ruled
or class-VII minimal models. Their Euler
numbers are respectively strictly positive,
zero and zero; the last assertion uses
the Hopf restriction \ref{cdp:st:C.2.UN.19}.
Thus every $e(Z_v)$ is nonnegative,
and it is positive exactly for a
rational component. The point-blowup comparisons give
$e(T_v)=e(Z_v)+n_v$.

\item  Let $t$ denote the number of triple
points. Topological inclusion-exclusion
on the SNC strata gives
\begin{equation}
 \begin{aligned}
 e(S)&=\sum_v e(T_v)-\sum_C e(C)+t\\
     &=\sum_v e(Z_v)+\sum_v n_v+t
                                      \ge0.
 \end{aligned}
\hypertarget{tag-00R8}{}\label{cdp:eq:C.MR.16}
\end{equation}
If $e(S)=0$, all
nonnegative terms vanish: $t=0$,
every $n_v=0$, and no component is
rational. These are exactly the
initial hypotheses of the proved
reduction, with no higher-genus ruled
component. Conversely those same
conditions make every term in
\eqref{cdp:eq:C.MR.16} zero. This proves the
extension and its asserted exact
zero description.

\end{enumerate}

\end{enumerate}
\end{proof}

\subsection{Rational toric components and the dual torus}\label{subsec:II:rational-toric-components}

\begin{proposition}[Toric components and the dual torus]\hypertarget{tag-01EE}{}\label{cdp:st:C.2.UN.68}
\StatementRecord{cdp:st:C.2.UN.68}{proved}{01EE}{\SeriesRef{cdp:st:C.2.UN.65}, \SeriesRef{cdp:eq:C.TC.3}, \SeriesRef{cdp:st:C.2.UN.19}, \SeriesRef{cdp:eq:C.TC.2}, \SeriesRef{cdp:eq:C.TC.5}, \SeriesRef{cdp:eq:C.TC.6}, \SeriesRef{cdp:eq:C.TC.9}, \SeriesRef{cdp:eq:C.TC.8}, \SeriesRef{cdp:eq:C.TC.11}, \SeriesRef{cdp:eq:C.TC.12}; \cite[Definition 1.2 and Proposition 1.3]{GHK2015}}

Let $f:X\to\Delta$ have the full reduced simple-normal-crossing
central fibre $S=\sum_{i=1}^{v}V_i$ of \ref{cdp:st:C.2.UN.65}, and let its
canonical evaluation orders be equal. Suppose that $e(S)>0$.
Then every $V_i$ is rational and its boundary
$D_i=\sum_{j\ne i}V_i\cap V_j$ is the full toric boundary
for an algebraic two-torus action on $V_i$.
The dual complex of this given fibre is a closed real
two-torus. If $r$ counts its connected double curves and
$t$ its triple points, then
\begin{equation}
                  r=3v,\qquad t=2v,\qquad e(S)=2v.
\hypertarget{tag-012C}{}\label{cdp:eq:C.TC.1}
\end{equation}
In particular $v\ge3$ and $e(S)\ge6$ in the embedded
simple-normal-crossing setting of the statement.

\end{proposition}

\begin{proof}

\begin{enumerate}[label=\arabic*.,leftmargin=*]
\item Use the three-form $\Omega/f^k$ of \ref{cdp:st:C.2.UN.65}.
It gives
\begin{equation}
               K_{V_i}+D_i=0,\qquad
               K_C=\OO_C(-T_C)
\hypertarget{tag-012G}{}\label{cdp:eq:C.TC.2}
\end{equation}
by its component residue and the principal-fibre line maps.
Thus a boundary component is either elliptic and disjoint
from the remaining boundary, or rational and meets it
at exactly two transverse points. Consequently the
rational boundary components form disjoint nodal cycles.
The positive Euler identity $e(S)=t$ of \ref{cdp:st:C.2.UN.65} supplies
at least one such cycle.

\item We check how a point blowdown on a component
acts on this divisor. Let $E$ be a smooth rational
$(-1)$-curve and $\pi:V\to V'$ its ordinary blowdown.
Adjunction gives
\begin{equation}
                         D\cdot E=1.
\hypertarget{tag-012H}{}\label{cdp:eq:C.TC.3}
\end{equation}
If $E$ is not in $D$, it meets exactly one smooth
boundary point transversely. It cannot meet a node,
which would contribute at least two to \eqref{cdp:eq:C.TC.3}.
The image $D'=\pi_*D$ is again reduced nodal and
$K_{V'}+D'=0$. Locally this is the ordinary blowup
at a smooth point of $D'$ with $D$ its strict transform.

\item If $E$ is in $D$, then
$(D-E)\cdot E=2$. Its two intersection points are
distinct and transverse. The blowdown identifies
their two transverse tangent directions and creates
one ordinary node on $D'$. The two branches may
belong to the same irreducible curve. This is the
inverse of the blowup of a boundary node;
again $K_{V'}+D'=0$ by the ordinary canonical
Jacobian formula. In particular a rational cycle
persists. When it has become an irreducible nodal
rational curve, it cannot be the smooth exceptional
curve of a further ordinary point blowdown.

\item These two maps preserve the separation between
the rational cycles and the elliptic boundary
components. Contracting the point blowups to
a minimal model therefore retains a rational
cycle whenever one was present. Each component
has $\kappa=-\infty$, because $-K=D$ is
nonempty effective. A nonrational ruled minimal
model has rational curves only in its ruling
fibres: a map from $\PP^1$ to the positive-genus
base is constant. Those fibres are mutually
disjoint smooth rational curves and cannot
contain a nodal rational cycle. The other
nonrational possibility is the Hopf minimal
model of the class-VII component,
by \ref{cdp:st:C.2.UN.19}. \Cref{lem:II:hopf-curves} identifies every compact irreducible
curve on that Hopf minimal model as smooth elliptic. Thus a
rational cycle cannot persist there. Hence a
component whose boundary contains a rational
cycle is rational.

\item On a rational component the exact sequence
\begin{equation}
 0\longrightarrow K_V\longrightarrow\OO_V
                 \longrightarrow\OO_D\longrightarrow0
\hypertarget{tag-012I}{}\label{cdp:eq:C.TC.4}
\end{equation}
has $H^0(K_V)=H^1(K_V)=0$, the second vanishing
being Serre dual to $H^1(\OO_V)=0$.
It follows that $H^0(\OO_D)=\C$, so $D$ is
connected. If $D$ contains a rational curve,
it is therefore exactly one rational nodal
cycle. A neighbouring component contains
that same rational double curve and hence
is rational by the preceding argument.
Propagation through the connected fibre
shows that every $V_i$ is rational and
every $D_i$ is one rational cycle.

Let $\mathcal B$ be the dual complex with
one vertex per $V_i$, one edge per connected
double curve and one triangle per triple
point. Multiple edges and triangles are
retained. Every edge belongs to two
triangles by \eqref{cdp:eq:C.TC.2}, and the link of
each vertex is exactly the boundary cycle
on that component. Thus the topological realization of $\mathcal B$
is a connected closed real surface.

\item All positive-degree structure-sheaf
cohomology of the rational components
and rational double curves vanishes.
The SNC restriction sequence therefore
identifies, by its alternating restriction
maps, the full complex computing
$H^*(S,\OO_S)$ with
\begin{equation}
           \C^{v}\longrightarrow\C^{r}
                                  \longrightarrow\C^{t},
\hypertarget{tag-012J}{}\label{cdp:eq:C.TC.5}
\end{equation}
the cellular cochain complex of
$\mathcal B$. The principal-fibre
adjunction and the volume form
give $\omega_S\simeq\OO_S$.
Serre duality on the compact Gorenstein
surface $S$ consequently gives
\begin{equation}
             H^2(S,\OO_S)\simeq
                         H^0(S,\OO_S)^*\simeq\C.
\hypertarget{tag-012K}{}\label{cdp:eq:C.TC.6}
\end{equation}
Here $H^0(S,\OO_S)=\C$ follows from
compactness, connectedness and reducedness.
Equations \eqref{cdp:eq:C.TC.5}--\eqref{cdp:eq:C.TC.6} make
$\mathcal B$ orientable. Proper flatness
gives $\chi(\OO_S)=0$, hence
$\chi(\mathcal B)=0$. A connected closed
oriented surface with this Euler
characteristic is a real two-torus.

\item The same cochain counts give
$v-r+t=0$. Together with the exact
incidence count $2r=3t$ of \ref{cdp:st:C.2.UN.65},
this proves \eqref{cdp:eq:C.TC.1}. A triple point
in a simple-NC divisor lies on three
distinct embedded components, so
$t>0$ implies $v\ge3$.

For a smooth rational surface with an
anticanonical nodal rational cycle
$D=D_1+\cdots+D_n$, allowing an
irreducible nodal rational curve when
$n=1$, put
\begin{equation}
                        Q(V,D)=e(V)-n=e(V\setminus D).
\hypertarget{tag-012L}{}\label{cdp:eq:C.TC.7}
\end{equation}
Indeed a rational nodal cycle has
Euler characteristic $n$, including
the one-component case. The toric-model comparison below gives
\begin{equation}
 Q(V,D)\ge0,\qquad Q(V,D)=0
       \quad\Longleftrightarrow\quad
       (V,D)\text{ is a toric pair}.
\hypertarget{tag-012M}{}\label{cdp:eq:C.TC.8}
\end{equation}

\item The point-contraction sequence above identifies $V$ as
a sequence of point blowups of $\PP^2$ or a Hirzebruch surface;
in particular it is projective. Hence $(V,D)$ is a Looijenga pair
in \cite[Definition~1.1, pp.~75--76]{GHK2015}. Apply
\cite[Definition~1.2 and Proposition~1.3, pp.~76--77]{GHK2015}
to obtain the proper birational maps
\[
 (V,D)\xleftarrow{\ \rho\ }(\widetilde V,\widetilde D)
       \xrightarrow{\ \tau\ }(\overline V,\overline D),
\]
where $\rho$ is a toric blowup, $(\overline V,\overline D)$ is a
smooth complete toric pair, and $\tau|_{\widetilde D}$ is an
isomorphism onto $\overline D$.
Factor the two morphisms into ordinary smooth point blowdowns.
For the first $(-1)$-curve $E$ contracted by $\rho$, the identity
$\widetilde D=(\rho^{-1}D)_{\rm red}$ implies that $E$ outside
$\widetilde D$ would map outside $D$ and be disjoint from
$\widetilde D$. This contradicts $\widetilde D\cdot E=1$ in
\eqref{cdp:eq:C.TC.3}. Thus $E$ lies in $\widetilde D$, and its
contraction is the boundary-node blowdown described there.
The full reduced inverse-image identity and the anticanonical
nodal pair persist after that contraction, so induction gives
this description for every step of $\rho$.
Each curve contracted by $\tau$ is outside
the boundary, since $\tau|_{\widetilde D}$ is an isomorphism; the
other case after \eqref{cdp:eq:C.TC.3} identifies each step with
the blowdown of an exceptional curve meeting one smooth boundary
point. These descriptions persist at every intermediate pair.
The first kind of step decreases both $e(V)$ and the number of
boundary components by one, while the second decreases only $e(V)$.
The toric orbit decomposition gives $e(\overline V)=\#\overline D$.
Consequently $Q(V,D)=Q(\widetilde V,\widetilde D)$ is the
nonnegative number of point blowdowns in $\tau$.

If this number is zero, $\tau$ is an isomorphism and $\widetilde V$
has its toric action. The action descends along every step of $\rho$:
for the boundary $(-1)$-curve with primitive fan ray $v_i$, the
self-intersection equation is $v_i=v_{i-1}+v_{i+1}$. Deleting
$v_i$ gives a smooth cone because
$\det(v_{i-1},v_{i+1})=\det(v_{i-1},v_i)=1$.
The associated toric morphism is the ordinary point blowdown in
the two charts $y=xu$ and $x=yv$. Its target identifies with the
given blowdown target by the unique isomorphism agreeing off the
exceptional curve and in these local point-blowup charts.
Thus its full boundary is the image of the specified boundary.
Iteration gives the algebraic two-torus action on the original
pair $(V,D)$. Conversely the orbit decomposition of any smooth
complete toric pair gives $Q=0$. This proves
\eqref{cdp:eq:C.TC.8} with its maps to the given pair.

For the ordinary contraction sequence to a minimal rational pair,
the same two Euler changes give
\begin{equation}
 Q(V,D)=Q(V_{\min},D_{\min})+a,
\hypertarget{tag-012N}{}\label{cdp:eq:C.TC.9}
\end{equation}
where $a$ counts the contracted curves outside the successive
boundaries. The minimal equality pairs have the following explicit
coordinates.

\item On $\PP^2$ the boundary has total
degree three, so $n\le3=e(\PP^2)$.
Equality means three distinct lines;
their nodal cycle condition makes
them nonconcurrent. The invertible
matrix of their linear equations is
a projective transformation
to the coordinate triangle.

\item On a minimal rational ruled surface
$q:\mathbb F_b\to\PP^1$, the total
horizontal degree of $D$ is two.
There are at most two horizontal
components. Each vertical prime is
a complete ruling fibre and meets
the horizontal part in two points
counted with multiplicity. The cycle
has only two incidences at each
horizontal component, so there are
at most two vertical components.
Hence $n\le4=e(\mathbb F_b)$.
Equality forces two degree-one
sections and two ruling fibres,
with the sections disjoint.
Their two quotient line maps split
the ruling bundle. A base coordinate
puts the two marked fibres at zero
and infinity. The resulting pair is
\begin{equation}
 \left(\PP_{\PP^1}(\OO\oplus\OO(b)),
        \text{two sections plus the fibres over }0,\infty\right).
\hypertarget{tag-012D}{}\label{cdp:eq:C.TC.10}
\end{equation}
In a line coordinate the transitions
are $w_\infty=z^{-b}w_0$.
The maps $(z,w_0)\mapsto(\lambda z,\mu w_0)$
and $(z^{-1},w_\infty)\mapsto
(\lambda^{-1}z^{-1},\lambda^{-b}\mu w_\infty)$
give its two-torus action
and precisely the four specified
boundary components.

\item Reversing a boundary-node contraction lifts the torus
action through the invariant ideal $(x,y)$; the charts $y=xu$
and $x=yv$ identify its full reduced inverse-image boundary with
the toric boundary. Thus the coordinates above and the maps $\rho,\tau$
describe the action on every charge-zero pair in
\eqref{cdp:eq:C.TC.8}.

Write $n_i$ for the number of
curves in $D_i$. Every double
curve belongs to two component
boundaries, so
\begin{equation}
                         \sum_i n_i=2r=3t.
\hypertarget{tag-012E}{}\label{cdp:eq:C.TC.11}
\end{equation}
The exact HRR and boundary-square
calculation in \ref{cdp:st:C.2.UN.65} gives
\begin{equation}
            \sum_i e(V_i)
             =12(r-t)-3t=3t.
\hypertarget{tag-012F}{}\label{cdp:eq:C.TC.12}
\end{equation}
Subtracting \eqref{cdp:eq:C.TC.11} from
\eqref{cdp:eq:C.TC.12} yields
$\sum_i Q(V_i,D_i)=0$.
Every summand is nonnegative by
\eqref{cdp:eq:C.TC.8}, so each is zero.
The equality morphisms
above make every specified
$(V_i,D_i)$ a toric pair.
This proves the asserted
classification on the given
degeneration.

\end{enumerate}
\end{proof}

\section{Maximum components and canonical-flat toric models}\hypertarget{tag-0347}{}\label{sec:II:11}


\subsection{Rational boundaries and maximum components}\label{subsec:II:rational-boundaries-maxima}

\begin{lemma}[Positive-square rational curves on fibre components]\hypertarget{tag-P365}{}\label{lem:II:positive-rational-curve}
\StatementRecord{lem:II:positive-rational-curve}{proved}{P365}{\cref{lem:II:cyclic-component-cover,cdp:st:C.2.UN.15,cdp:st:C.2.UN.19,cdp:st:C.2.UN.31,lem:II:negative-kodaira-models}; \eqref{eq:II:quasibundle-canonical}; \cite[Proposition~3.12, p.~20]{PopescuPampu2016}}
Let $f:U\to\Delta$ be a proper degeneration of connected complex two-tori from a smooth complex threefold, with full reduced normal-crossing central fibre and smooth component normalizations. Let $T$ be one such normalization, and let $C\subset T$ be a smooth rational curve with $C^2>0$. Then $T$ is rational.

For $C^2=s\in\{1,2\}$, blow up any $s$ distinct points of $C$, writing $\mu:T'\to T$ and $C'$ for its strict transform. The complete system $|C'|$ gives a proper map $q:T'\to\PP^1$ with connected rational general fibre and $q^*\OO_{\PP^1}(1)=\OO_{T'}(C')$. Every irreducible curve on $T$ disjoint from $C$ is unchanged by $\mu$ and is a rational vertical curve for $q$.

For $C^2=1$, the complete system $|C|$ gives a birational morphism $\rho:T\to\PP^2$, isomorphic on a neighbourhood of $C$, with $\rho^*\OO_{\PP^2}(1)=\OO_T(C)$ and $\rho|_C:C\xrightarrow{\sim}\ell$ for a line $\ell$. It factors into point blowdowns away from $C$ and contracts every irreducible curve disjoint from $C$.
\end{lemma}
\begin{proof}
Take a minimal-model map $\pi:T\to Z$. The possible minimal models are given by \cref{cdp:st:C.2.UN.15,cdp:st:C.2.UN.19,lem:II:negative-kodaira-models}. On a nonrational ruled $Z$, every rational curve maps constantly to the positive-genus ruling base. A nonexceptional rational curve on $T$ is therefore the strict transform of a ruling fibre and has nonpositive square; every exceptional curve has negative square. On a Hopf $Z$, every compact curve is elliptic by \cref{lem:II:hopf-curves}, so every rational curve on $T$ is exceptional. These alternatives exclude $C^2>0$.

For the remaining case $\kappa(Z)\ge0$, apply \cref{lem:II:cyclic-component-cover} to the same central component with $m_i=1$. Its root map is $t=s$, its group is $G=\{1\}$ and its degree-one component map $E\to T$ is proper bimeromorphic. Thus the minimal surface $E_0$ in \cref{cdp:st:C.2.UN.31} is this same $Z$. It is a torus or the fixed-Jacobian elliptic model there, with full fibres $m_bC_b$ and identity monodromy. In the torus case $K_Z$ is trivial. In the elliptic case take $\ell_0>0$ divisible by every $m_b$. The canonical map \eqref{eq:II:quasibundle-canonical} and $q_0^*(b)=m_bC_b$ give
\[
 K_Z^{\otimes\ell_0}\simeq q_0^*L,\qquad
 L=K_B^{\otimes\ell_0}\otimes\OO_B\Bigl(\sum_b\ell_0(1-1/m_b)b\Bigr),
\]
with
\begin{equation}
 \frac{\deg L}{\ell_0}=2g(B)-2+\sum_b(1-1/m_b)\ge0.
 \hypertarget{tag-00SW}{}\label{cdp:eq:C.NM.5}
\end{equation}
Indeed a nonzero pluricanonical section has a nonzero power of exponent divisible by $\ell_0$, and $q_{0*}\OO_Z=\OO_B$ identifies it with a section of a positive power of $L$. Hence $\deg L\ge0$, and $K_Z$ is nef. The point-blowup Jacobians give $K_T=\pi^*K_Z+\sum_Eb_EE$ with $b_E>0$. A nonexceptional $C$ then has $K_T\cdot C\ge0$, whereas adjunction gives $K_T\cdot C=-2-C^2<0$. An exceptional $C$ has negative square. The remaining minimal model is rational, proving
\begin{equation}
 \text{a smooth rational curve of positive square on }T
 \quad\Longrightarrow\quad T\text{ is rational}.
 \hypertarget{tag-00SV}{}\label{cdp:eq:C.NM.4}
\end{equation}

For $s=1$ or $2$, the specified point blowups give $C'^2=0$ and leave every curve disjoint from $C$ unchanged. The rational surface $T'$ has $H^1(T',\OO_{T'})=0$. The restriction sequence for $\OO_{T'}(C')$, with $\OO_{C'}(C')\simeq\OO_{\PP^1}$, gives a second section nonzero everywhere on $C'$. Together with the section cutting out $C'$, it generates $\OO_{T'}(C')$ and defines
\begin{equation}
 q:T'\longrightarrow\PP^1,\qquad
 q^*\OO_{\PP^1}(1)=\OO_{T'}(C').
 \hypertarget{tag-00SX}{}\label{cdp:eq:C.NM.6}
\end{equation}
The fibre $C'$ is reduced and connected. In the Stein factorization a finite map of degree greater than one would give either several points over this value or multiplicity greater than one on $C'$. Hence $q$ has connected fibres, and adjunction gives genus zero on its smooth general fibre. Each selected exceptional curve meets $C'$ once and is a section of $q$. On the rational projective surface $T'$, a fibre coordinate using one such section gives a birational map over $\PP^1$ to $\PP^1\times\PP^1$: on the generic genus-zero fibre its degree-one section divisor has two independent sections and defines that coordinate. Normalize and resolve its graph. The two resulting maps are compositions of point blowups by \cite[Propositions~3.10 and~3.12, p.~20]{PopescuPampu2016}. Every vertical component on this common resolution is a strict transform of a rational fibre or an exceptional rational curve. Its image on $T'$ is therefore rational. An irreducible curve disjoint from $C$ has intersection zero with $C'$, so is vertical for the same $q$.

Finally suppose $C^2=1$. The exact sequence
\begin{equation}
 0\longrightarrow\OO_T\longrightarrow\OO_T(C)
 \longrightarrow\OO_C(1)\longrightarrow0
 \hypertarget{tag-009Y}{}\label{cdp:eq:C.AP.8}
\end{equation}
and $H^1(T,\OO_T)=0$ give three sections. They generate the line on $C$, and the section defining $C$ is nonzero off $C$. Thus the complete system gives
\begin{equation}
 \rho:T\longrightarrow\PP^2,\qquad
 \rho^*\OO_{\PP^2}(1)=\OO_T(C),\qquad
 \rho|_C:C\xrightarrow{\sim}\ell.
 \hypertarget{tag-009Z}{}\label{cdp:eq:C.AP.9}
\end{equation}
Its image is a surface since $C^2=1$; the degree formula gives degree one, so $\rho$ is birational. The section cutting out $C$ is the pullback of a linear equation of $\ell$, so $\rho^*\ell=C$ as divisors. Apply \cite[Proposition~3.12, p.~20]{PopescuPampu2016} to this proper morphism of smooth surfaces. Every exceptional curve has zero intersection with $\rho^*\OO(1)=\OO_T(C)$, so is disjoint from $C$, and the displayed divisor equality places its image off $\ell$. The factorization is consequently by point blowdowns away from $C$ and is an isomorphism near $C$. Every irreducible curve disjoint from $C$ has degree zero under $\rho^*\OO(1)$ and is contracted.
\end{proof}

\begin{proposition}[Maximum components without a minimality assumption]\hypertarget{tag-01EF}{}\label{cdp:st:C.2.UN.69}
\StatementRecord{cdp:st:C.2.UN.69}{proved}{01EF}{\cref{cdp:st:C.2.UN.64}, \cref{cdp:st:C.2.UN.50}, \cref{cdp:st:C.2.UN.58}, \cref{cdp:st:C.2.UN.19}, \cref{cdp:st:C.2.UN.15}, \cref{cdp:st:C.2.UN.31}, \eqref{cdp:eq:C.NM.4}, \eqref{cdp:eq:C.NM.7}, \eqref{cdp:eq:C.NM.6}, \cref{cdp:st:C.2.UN.66}, \eqref{cdp:eq:C.NM.9}, \eqref{cdp:eq:C.NM.3}, \eqref{cdp:eq:C.NM.8}, \eqref{cdp:eq:C.NM.10}}

Let $f:U\to\Delta$ be a proper degeneration of connected complex
two-tori from a smooth threefold. Let its full central fibre
$S=\sum_vT_v$ be reduced and simple normal crossing, without triple
points. The components $T_v$ may be nonminimal.
Write $k_v$ for the canonical evaluation orders.

At a local maximum $T_v$ of these orders, precisely one of the
following descriptions applies:
\begin{enumerate}
\item Every adjacent component has order $k_v$.
\item The component is an unblown terminal ruled surface meeting
its parent in a section, with order drop one, or an unblown
terminal plane meeting its parent in a line, with order drop two.
These are the contractions of \cref{cdp:st:C.2.UN.64}.
\item The component is an unblown terminal Hopf surface, with
the exact axis normal order and positive canonical drop of
\cref{cdp:st:C.2.UN.50} and \cref{cdp:st:C.2.UN.58}.
\end{enumerate}
In the second alternative the ruled base has positive genus.
Thus, when the orders are not all equal, a global maximum is
one of the two kinds of strict-drop leaves.

Moreover every rational double curve can be removed by a finite
sequence of the plane-line contractions of \cref{cdp:st:C.2.UN.64}.
At each step the other double curves and their neighbourhoods
are unchanged. If $r_0$ is the number of connected rational
double curves of the original support, the resulting torus degeneration has no rational double curve and satisfies
\begin{equation}
             e(S)=e(S_{\rm nr})+2r_0 .
\hypertarget{tag-00SR}{}\label{cdp:eq:C.NM.1}
\end{equation}
There are exactly $r_0$ point contractions. Reversing their
recorded ordinary blowups recovers the original support,
including its original smooth threefold neighbourhood.
In particular the rational double edges are removed as terminal
edges of the incidence graph; none belongs to a cycle.
\end{proposition}

\begin{proof}

\begin{enumerate}[label=\arabic*.,leftmargin=*]
\item \par\medskip\noindent\textbf{Canonical restrictions at an edge.}
Canonical evaluation, adjunction and the equation $xy=t$
give the line maps \eqref{cdp:eq:C.MR.3}--\eqref{cdp:eq:C.MR.4} on these possibly nonminimal components:
\begin{equation}
 \begin{split}
 K_{T_v}&=\operatorname{div}(\omega_v)
       =\sum_{e=(v,w)}(k_w-k_v-1)C_e,\\
 N_{C_e/T_v}\otimes N_{C_e/T_w}&\simeq\OO_{C_e},\\
 K_{C_e}&\simeq N_{C_e/T_v}^{\,k_w-k_v}.
 \end{split}
\hypertarget{tag-00ST}{}\label{cdp:eq:C.NM.2}
\end{equation}
The last equality uses that the other boundaries are disjoint
from $C_e$. In particular, for $g_e=g(C_e)$,
\begin{equation}
                  (k_w-k_v)C_e^2=2g_e-2.
\hypertarget{tag-00SU}{}\label{cdp:eq:C.NM.3}
\end{equation}
These are the meromorphic-section and dual-normal line maps.

\par\medskip\noindent\textbf{A maximum with a rational boundary is already minimal.}
At a local maximum put $a_e=k_v-k_w+1\ge1$. Then
\begin{equation}
             -K_T=\sum_e a_eC_e,\qquad
             (1-a_e)C_e^2=2g(C_e)-2.
\hypertarget{tag-00SY}{}\label{cdp:eq:C.NM.7}
\end{equation}
The boundary is nonempty in the reducible case. A nonzero
pluricanonical section multiplied by the corresponding
power of this anticanonical section would be a nonzero
holomorphic function having zeros on a compact connected
surface. Consequently $\kappa(T)=-\infty$.

\item For a rational boundary, \eqref{cdp:eq:C.NM.7} gives exactly
\begin{equation}
                     (a_e,C_e^2)=(3,1)\ \text{or}\ (2,2).
\hypertarget{tag-00SZ}{}\label{cdp:eq:C.NM.8}
\end{equation}
Apply \cref{lem:II:positive-rational-curve} to this $(T,C_e)$, with the value $s=C_e^2$ in \eqref{cdp:eq:C.NM.8}. It gives rationality and the pencil \eqref{cdp:eq:C.NM.6}; every other disjoint boundary is a rational vertical curve.
Every such boundary would again have positive square by
\eqref{cdp:eq:C.NM.7}, whereas the Hodge index theorem permits at most
one positive-square class among mutually disjoint curves.
Thus there is only one boundary $C$, and $K_T=-aC$.
Any smooth rational $(-1)$-curve $E$ would satisfy
$a(C\cdot E)=1$, which is impossible for $a=2$ or $3$.
The surface is therefore minimal. The terminal-pair
classification and maps in \cref{cdp:st:C.2.UN.64} give its ruled
section or plane-line contraction.

\par\medskip\noindent\textbf{All other blowdowns occur only at coefficient-one elliptic boundaries.}
Suppose now that no boundary is rational. A $(-1)$-curve
$E$ is not in the boundary, and adjunction gives
\begin{equation}
                     1=(-K_T)\cdot E
                              =\sum_e a_e(C_e\cdot E).
\hypertarget{tag-00T0}{}\label{cdp:eq:C.NM.9}
\end{equation}
Every summand is a nonnegative integer. Exactly one is
one: $E$ meets one boundary transversely once, that
boundary has coefficient one, and $E$ misses all others.
The coefficient-one boundary is elliptic by \eqref{cdp:eq:C.NM.7}.
Its ordinary surface blowdown keeps this boundary smooth
and keeps all boundary curves disjoint. Pushforward of
the meromorphic canonical section gives the same
canonical divisor with the same coefficients on the
new surface. These are point blowdowns of the surface $T$.
Repeat to the minimal model.

\item The minimal possibilities with \eqref{cdp:eq:C.NM.7} are exactly those
computed in \eqref{cdp:eq:C.MR.6}--\eqref{cdp:eq:C.MR.7} and \cref{cdp:st:C.2.UN.64}. On a ruled
surface, intersection with a ruling fibre gives the three
partitions of two: one coefficient-two section, one
coefficient-one bisection, or two coefficient-one sections.
The connected elliptic bisection over an elliptic base is
excluded by the residue/trace map \eqref{cdp:eq:C.EL.7}--\eqref{cdp:eq:C.EL.8}.
On the plane the possibilities are a coefficient-three
line or a coefficient-one cubic. On an elliptic Hopf
surface the invariant relative differential gives
$\sum (k_w-k_T)/m_e=\operatorname{val}(T)-2$ with every
summand nonpositive; thus it is either a strict-drop
one-boundary leaf or a two-boundary equal-height component.
The nonlinear Hopf axis formula gives the same strict
leaf alternative, including its exact secondary framing.
For an algebraic-dimension-zero diagonal Hopf surface,
the two-axis character equation and $|\alpha|,|\beta|<1$
make both coefficients one; a single axis would impose
an integral resonance. These are marked pairs
by \cref{cdp:st:C.2.UN.50}, \cref{cdp:st:C.2.UN.58} and the proof of \cref{cdp:st:C.2.UN.66}.

\item In particular whenever a coefficient is greater than
one, there is just one boundary, with that coefficient.
No coefficient-one boundary then exists. Since the
surface blowdowns above preserve every boundary and
its coefficient, \eqref{cdp:eq:C.NM.9} shows that no such blowdown
could have taken place. The original strict-drop maximum
was already minimal. In all remaining cases every
coefficient is one and every neighbour has the same
height. A component at the global maximum with this
property cannot be adjacent to a strict-drop maximum
leaf. Its equal-height neighbours have the same property,
so connectedness propagates it to the whole support.
This proves the maximum assertion, except for the
rational ruled-section exclusion proved next.

\par\medskip\noindent\textbf{A rational section centre cannot lie in the smooth torus-fibre locus.}
In a ruled-section contraction of \cref{cdp:st:C.2.UN.64} with base
$B\simeq\PP^1$, the descended family is a
submersion near the compact centre $B$. Its containing
central component $P'$ is smooth there and
\begin{equation}
                    N_{B/P'}\simeq K_B\simeq\OO_{\PP^1}(-2).
\hypertarget{tag-00SS}{}\label{cdp:eq:C.NM.10}
\end{equation}
Choose a relatively compact tubular neighbourhood of
$B$ inside that smooth central surface and a slightly
larger compact neighbourhood in the submersion locus.
A smooth metric gives horizontal lifts of the two
coordinate vector fields on the disc. For sufficiently
small time their flows exist on the whole smaller
neighbourhood and embed it in a nearby fibre.
The induced diffeomorphism preserves the oriented
normal bundle and its Euler number. Thus the nearby
two-torus would contain an embedded sphere of
self-intersection minus two. Its universal cover is
the contractible space $\C^2$, so this sphere map is
nullhomotopic, its homology class is zero, and its
self-intersection is zero. This contradiction excludes
the rational ruled-section alternative.

\par\medskip\noindent\textbf{Point contractions along rational edges.}
Take a rational double curve and orient it toward its
higher-height endpoint. Equation \eqref{cdp:eq:C.NM.3} makes the
height difference nonzero; at that higher endpoint
its square is positive, equal to one or two. Apply \cref{lem:II:positive-rational-curve} to this higher endpoint and its incoming curve, whose square is $1$ or $2$. The resulting pencil \eqref{cdp:eq:C.NM.6} makes every other disjoint boundary rational.
Hodge index leaves at most one positive-square
boundary, the incoming one. Every other boundary
therefore has negative square by \eqref{cdp:eq:C.NM.3}, and its
neighbour has strictly higher height. Follow one
of those edges when it exists.

\item The heights strictly increase and there are finitely
many components. The path ends at a rational
component having exactly its one incoming boundary.
It has $K=-aC$, with $(a,C^2)$ as in \eqref{cdp:eq:C.NM.8}.
The same integer-intersection argument excludes every
$(-1)$-curve, so it is a terminal leaf classified in \cref{cdp:st:C.2.UN.64}. Its ruled
section possibility has just been excluded by
\eqref{cdp:eq:C.NM.10}. It is therefore a plane-line leaf, and
\cref{cdp:st:C.2.UN.64} contracts it to a smooth point of the smooth
central locus of its parent.

\item Every other boundary misses that line, so the blowdown is an isomorphism on its neighbourhood.
It removes exactly one rational double curve, creates
no new double curve or triple point, retains the
punctured torus family, and decreases $e(S)$ by two.
Apply the argument to the new degeneration
until no rational double curve remains. Exactly
$r_0$ steps have been made; \eqref{cdp:eq:C.NM.1} and the inverse ordinary-blowup reconstruction follow.

\end{enumerate}
\end{proof}













\subsection{Toric central log spaces}\label{subsec:II:toric-central-log-spaces}

\begin{proposition}[All positive canonical-flat central log spaces over a disc]\hypertarget{tag-01EI}{}\label{cdp:st:C.2.UN.71}
\StatementRecord{cdp:st:C.2.UN.71}{proved}{01EI}{\cref{cdp:st:C.2.UN.68}, \eqref{cdp:eq:C.TG.4}, \cref{cdp:st:C.2.UN.65}, \eqref{cdp:eq:C.TG.8}, \eqref{cdp:eq:C.TG.3}, \eqref{cdp:eq:C.TG.6}, \eqref{cdp:eq:C.TG.10}, \eqref{cdp:eq:C.TG.12}, \eqref{cdp:eq:C.TG.11}, \eqref{cdp:eq:C.TG.13}; \cite[version~3, Theorems~4.14, 5.2 and~5.4, pp.~114--116 and~128--134]{GrossSiebert2006I}; \cite[version~3, Theorem~3.22, pp.~95--97, and Example~3.30, p.~106]{GrossSiebert2006I}}

The positive canonical-flat central fibres of \cref{cdp:st:C.2.UN.68},
with their divisorial log structure and marked
dual complex, have the following complete description.
Take a rank-two lattice $N$, a finite-index sublattice
$\Gamma\subset N$, a $\Gamma$-periodic unimodular
triangulation $\mathscr P$ of $N_{\R}$ with vertex set
$N$, and a homomorphism
\begin{equation}
       c:\Gamma\longrightarrow T_N
          :=\operatorname{Hom}(N^*,\C^*).
\hypertarget{tag-012Q}{}\label{cdp:eq:C.TG.1}
\end{equation}
Require that each triangle have three distinct vertices
in $N/\Gamma$; this is the embedded simple-NC condition.
The central component at $p\in N/\Gamma$ is the toric surface of the star fan
$\{\R_{\ge0}(\sigma-p):p\in\sigma\in\mathscr P\}$.
The homomorphism in \eqref{cdp:eq:C.TG.1} supplies all the remaining
logarithmic gluing constants, modulo the explicit torus
changes of the local charts described below.

Every such datum gives a smooth proper analytic
degeneration of connected complex two-tori over a
sufficiently small disc, with trivial canonical line,
and with exactly the specified marked central log space.
Every fibre in \cref{cdp:st:C.2.UN.68} is obtained in this way.
In a basis of $N$, let the columns of the integral
matrix $Q$ be a basis of $\Gamma$, and choose logarithms
of its two values under $c$ as the columns of $C$.
The nearby torus is
\begin{equation}
 F_t=\C^2\big/
       \left(2\pi i\Z^2+(\log t\,Q+C)\Z^2\right),
       \qquad 0<|t|<\epsilon .
\hypertarget{tag-012X}{}\label{cdp:eq:C.TG.2}
\end{equation}
The central Euler number is $2[N:\Gamma]$.

The construction also permits a holomorphic family of homomorphisms $\lambda\mapsto c_\lambda(t)$ from $\Gamma$ to $T_N$, defined on a disc and satisfying $c_\lambda(0)=c_\lambda$. It gives the same marked central log space, smooth proper torus fibres over the puncture and a nowhere-zero canonical form. In \eqref{cdp:eq:C.TG.2}, its period matrix is obtained by replacing $C$ with a holomorphic matrix $C(t)$ of logarithms of the two values $c_{Qe_j}(t)$. The constant family gives the preceding occurrence.

\end{proposition}

\begin{proof}

\begin{enumerate}[label=\arabic*.,leftmargin=*]
\item \par\medskip\noindent\textbf{Integral-affine charts.}
On each toric component $V_i$ of \cref{cdp:st:C.2.UN.68} write
$N_i=\pi_1(V_i\setminus D_i)$, and let $n_{ij}$
be the primitive angular meridian of its boundary
curve $C_{ij}$. A triple point $ijk$ gives the
unimodular triangle
\begin{equation}
                 0,\ n_{ij},\ n_{ik}\quad\hbox{in }(N_i)_{\R}.
\hypertarget{tag-012Y}{}\label{cdp:eq:C.TG.3}
\end{equation}
To pass from the chart at $j$ to that at $i$,
use the unique integral-affine map which sends
$0$ to $n_{ij}$ and whose linear part $P_{ij}$ satisfies
\begin{equation}
 P_{ij}n_{ji}=-n_{ij},\qquad
 P_{ij}n_{jk}=n_{ik}-n_{ij}.
\hypertarget{tag-012Z}{}\label{cdp:eq:C.TG.4}
\end{equation}
These are the angular maps. Indeed in a
triple chart $f=xyz$, the coordinates on $V_i$
are $(y,z)$ and those on $V_j$ are $(x,z)$.
On the nearby fibre $x=t/(yz)$, a loop varying
$x$ with $z$ fixed has angular vector $-n_{ij}$,
and a loop varying $z$ with $x$ fixed has vector
$n_{ik}-n_{ij}$. Multiplication of these coordinates
by holomorphic units leaves these integral maps
unchanged after the chart is shrunk.

\item There are exactly two triple points on $C_{ij}$.
Write their third vertices as $k,l$, retaining
their two distinct incidences if the same global
component occurs twice. The toric fan relation on
$V_i$ is
\begin{equation}
 n_{ik}+n_{il}=-(C_{ij}^2)_{V_i}n_{ij}.
\hypertarget{tag-0130}{}\label{cdp:eq:C.TG.5}
\end{equation}
The corresponding relation at $j$ agrees with
\eqref{cdp:eq:C.TG.4} at the other endpoint precisely because
\begin{equation}
             (C_{ij}^2)_{V_i}+(C_{ij}^2)_{V_j}=-2,
\hypertarget{tag-0131}{}\label{cdp:eq:C.TG.6}
\end{equation}
the normal-product and two-triple-point
formula of \cref{cdp:st:C.2.UN.65}. Thus there is no affine singularity
in an edge. The complete smooth toric fan at each
vertex makes its star an affine neighbourhood of
the origin. The triangle charts and \eqref{cdp:eq:C.TG.4}
therefore give an integral-affine structure on
the entire dual two-torus $B$.

\par\medskip\noindent\textbf{The angular lattices in the nearby torus.}
The nowhere-zero three-form of \cref{cdp:st:C.2.UN.65} is, in a
triple chart,
\begin{equation}
       \Omega=a(x,y,z)\,dx\wedge dy\wedge dz,
          \quad a(0)\ne0,
          \qquad f=xyz.
\hypertarget{tag-0132}{}\label{cdp:eq:C.TG.7}
\end{equation}
The induced relative two-form on $F_t$ is
$a(x,y,t/(xy))\,d\log x\wedge d\log y$.
Its integral over the real angular two-torus
$|x|=|y|=|z|=|t|^{1/3}$ is exactly
\begin{equation}
 (2\pi i)^2\sum_{q\ge0}a_{q,q,q}t^q
                  =(2\pi i)^2(a(0)+O(t))\ne0.
\hypertarget{tag-0133}{}\label{cdp:eq:C.TG.8}
\end{equation}
This follows by the convergent power-series expansion
and the two circle integrals; only the diagonal
monomials survive.

\item For a fixed component, its toric complement retracts
onto such an angular two-torus. Any finite collection
of its loops, their homotopies and the connecting
paths to the corner lies in a compact subset away
from the boundary. The submersion $f$ near this
compact subset gives smooth horizontal
transports into $F_t$. Consequently its angular
lattice has a well-defined homomorphism
$N_i\to\pi_1(F_t)$, up to conjugation, and the
corner homomorphism is this same map. Since
$\pi_1(F_t)=\Z^4$ is abelian, conjugation makes
no ambiguity. If this map had rank at most one,
the fundamental class of the angular torus would
map to zero in $H_2(F_t,\Z)$: the complex
torus is a $K(\Z^4,1)$ and its degree-two homology
map is the exterior square of its degree-one map.
This contradicts \eqref{cdp:eq:C.TG.8}. Thus every $N_i$
injects into the $\Z^4$.

\item The corner maps and the compact double-curve collars
identify the injected lattices by exactly \eqref{cdp:eq:C.TG.4}.
Transport around any loop in $B$ gives a linear
holonomy $A$ with $\iota_iA=\iota_i$ in $\Z^4$.
Injectivity gives $A=I$. Thus the linear holonomy is trivial, by these injections into the given smooth fibre.

\item The affine charts now differ by translations.
Their Euclidean metrics give a complete flat metric
on compact $B$. Extending straight lines on its
universal cover proves that the developing map is
an affine isomorphism with $N_{\R}$.
The deck group is a full lattice of translations
$\Gamma\subset N$, because all the transition
translations in \eqref{cdp:eq:C.TG.3}--\eqref{cdp:eq:C.TG.4} are integral.
Every developed triangle is unimodular. Conversely
every lattice point of $N$ is a vertex, since a
unimodular triangle contains no other lattice point
and the triangles cover $N_{\R}$. Thus the dual complex is $(N_{\R},\mathscr P)/\Gamma$,
with exactly $[N:\Gamma]$ vertices.

\par\medskip\noindent\textbf{Logarithmic gluing.}
Use the exact toric strata just obtained. The divisorial log structure of $(X,S)$ restricted to $S$ has chart $\mathbb N^3$ at a triple point, with base generator $(1,1,1)$, chart $\mathbb N^2$ along an open double curve, and the one-factor chart elsewhere. The equations $xyz=t$ and $xy=t$ make the log morphism smooth everywhere. Its conductor square is the SNC square, its normalization is the disjoint union of the complete toric $V_i$, and it is Gorenstein.

The gluing classification is \cite[version~3, Theorems~4.14, 5.2 and~5.4, pp.~114--116 and~128--134]{GrossSiebert2006I}. We give the comparison for this analytic central space. Identify every normalization stratum with its toric model. A double curve is $\PP^1$ with its two labelled corner points. Its identification between the two components fixes those points, and hence has the form $z\mapsto a z$ with $a\in\C^*$. Thus all closed-stratum identifications are algebraic toric maps.

\item At a triple point, the log chart restricts on each toric affine plane to a monoid map $p\mapsto h_pz^p$, where $h_p$ is an invertible convergent germ and $h_{p+q}=h_ph_q$ whenever the monomials lie on the same component. Take residues at that corner:
\begin{equation}
       p\longmapsto h_pz^p
          \quad\longmapsto\quad
       p\longmapsto h_p(0)z^p.
\hypertarget{tag-0134}{}\label{cdp:eq:C.TG.9}
\end{equation}
Restriction to a shared axis commutes with this residue map. Consequently the constant monomial maps intertwine the same closed gluing diagram. Apply the colimit construction in the proof of \cite[version~3, Theorem~4.14, pp.~114--116]{GrossSiebert2006I} to these maps. Its identity on the labelled strata gives a toric chart $V(\sigma)$ and the chart map to $S$. For the present standard simplex, the local conductor identity is the equality of $\C\{x,y,z\}/(xyz)$ with the triples of functions on its three planes having equal restrictions on their pairwise axes. Thus this chart map is a local analytic isomorphism; the same two-plane identity applies on an edge.

On overlaps the chart maps differ by the constant monomial maps induced by those closed identifications. The resulting finite algebraic atlas and its algebraic \'etale equivalence relation define an algebraic space $S_{\mathrm{alg}}$, together with the displayed analytic isomorphism $S_{\mathrm{alg}}^{\mathrm{an}}\simeq S$. Its normalization is the same finite union of complete toric surfaces. The inverse image of its diagonal under the product normalization map is the union of the component diagonals and the closed graphs of the conductor identifications, so the diagonal is closed. The finite surjective normalization and the completeness of those surfaces then give properness. This establishes the underlying-space hypotheses of the cited theorem through the specified maps.

\item The log section also descends through this comparison. On a double curve $C$, with corner points $p,q$, its smoothing line is
\[
             N_{C/V_i}\otimes N_{C/V_j}\otimes\OO_C(p+q).
\]
Its degree is zero by \eqref{cdp:eq:C.TG.6}. The given smooth-total equation supplies a nowhere-zero section. In a global trivialization of this algebraic degree-zero line on $\PP^1$, that holomorphic section is a nonzero constant. At a triple point the three restrictions have the multiplicative compatibility of \cite[version~3, Theorem~3.22 and Example~3.23(1), pp.~95--97; Example~3.30, p.~106]{GrossSiebert2006I}: in the standard simplex it is exactly equality of their values at the common point. Thus these same constants and identifications define an algebraic section of the log-smooth-structure torsor.

To check that its analytification is the given log structure, use the units quotient in that theorem: the tuples $(h_p)$ are taken modulo multiplicative unit maps, and their edge ratios are the just specified smoothing sections. The kernel calculation and the gluing of units on the coordinate planes work in the convergent local rings by the same conductor identity. Moreover the ghost group is generated by the base class at generic points. A unit homomorphism killing that class is therefore one on a dense open set and hence everywhere on the reduced space. The local isomorphisms consequently glue uniquely. This is the map underlying \cite[version~3, Proposition~3.20 and its proof, pp.~92--94]{GrossSiebert2006I}; its local invertible-Jacobian step here is the holomorphic inverse function theorem on these SNC charts. The algebraic section thus gives $S_{\mathrm{alg}}^\dagger$ with $(S_{\mathrm{alg}}^\dagger)^{\mathrm{an}}\simeq S^\dagger$ over the standard log point.

The hypotheses of \cite[version~3, Definition~4.3, pp.~108--109]{GrossSiebert2006I} now hold with $k=\C$ and empty log-singular set. The affine singular locus is empty, so every local monodromy coefficient is zero; the positivity condition of Definition~1.54, p.~41, and the simplicity condition of Definition~1.60, pp.~45--46, hold with $p=0$. The nowhere-zero smoothing sections are positive log sections as well.

\item Choose the chart changes \eqref{cdp:eq:C.TG.9} so that the values of the log section at each corner are one, as in \cite[version~3, Definitions~4.23--4.24, p.~120]{GrossSiebert2006I}. Its Proposition~4.25, pp.~120--123, gives Condition~(LC) for these open gluing data. Theorems~5.2 and~5.4, pp.~128--134, then identify the marked log space with lifted gluing data for the constant sheaf $N\otimes\C^*=T_N$ on $B$. The chart comparison of Proposition~4.30, pp.~127--128, makes the isomorphisms explicit: a coboundary has
\[
 s'_{\alpha\beta}=a_\alpha^{-1}s_{\alpha\beta}a_\beta,
\]
and the map on the chart indexed by $\alpha$ is its torus multiplication by $a_\alpha$. The pullback log section is normalized; uniqueness in Theorem~5.2 identifies it with the same section. These are therefore isomorphisms of the marked central log spaces, including their given log structures.

On the contractible universal cover $N_{\R}$, path products trivialize this constant-sheaf cocycle. Changing a path by a deck element records its holonomy, so
\begin{equation}
 H^1(B,N\otimes\C^*)
       =\operatorname{Hom}(\Gamma,T_N).
\hypertarget{tag-012R}{}\label{cdp:eq:C.TG.10}
\end{equation}
The two commuting deck holonomies are exactly the marked values of \eqref{cdp:eq:C.TG.1}.

\par\medskip\noindent\textbf{The toric threefold and its deck maps.}
For the construction take the holomorphic homomorphism $c(t)$ of the statement. All maps below keep the base coordinate $t$ fixed; their restrictions at zero use exactly $c(0)=c$.
For every triangle $\sigma$ of $\mathscr P$, take
the cone over $\sigma\times\{1\}$ in
$N_{\R}\oplus\R$. These cones and their faces
form a fan $\Sigma$. Its three primitive
generators $(p_0,1),(p_1,1),(p_2,1)$ are an
integral basis because the triangle is
unimodular. Its toric chart is therefore
$\C^3$ with coordinates $x_0,x_1,x_2$,
and the character of the last dual coordinate is
\begin{equation}
                  t=x_0x_1x_2,
       \qquad z^m=\prod_{a=0}^2x_a^{\langle m,p_a\rangle}
                       \quad(m\in N^*).
\hypertarget{tag-012S}{}\label{cdp:eq:C.TG.11}
\end{equation}
The charts glue by their Laurent
monomial localizations. The fan intersection
condition makes this a Hausdorff smooth toric
analytic threefold. Although the fan has
infinitely many cones, each central orbit has
the ordinary finite star of its cone, and the
open torus is a single common open chart.
These are analytic neighbourhoods.

\item A translation $\lambda\in\Gamma$ acts on the
lattice $N\oplus\Z$ by the integral shear
$(n,h)\mapsto(n+h\lambda,h)$. It preserves
the fan and induces the monomial map
between its charts. Multiply this map by
$c_\lambda(t)$ on the torus factor. On the
punctured open torus the result is exactly
\begin{equation}
              g_\lambda(z,t)=(c_\lambda(t)t^\lambda z,t).
\hypertarget{tag-012T}{}\label{cdp:eq:C.TG.12}
\end{equation}
For each $t$, the values $c_\lambda(t)$ are a homomorphism in $\lambda$. Their torus multiplications commute with the shears, since they depend only on the fixed base coordinate. Thus $g_\lambda g_\mu=g_{\lambda+\mu}$
on every chart. On the central fibre these
maps identify the component and boundary
charts by the cocycle in \eqref{cdp:eq:C.TG.10}. The
quotient central log space is consequently
the exact space specified earlier, including
its given log section.

\par\medskip\noindent\textbf{Freeness, properness and compact fibres.}
For each $t$, extend $\lambda\mapsto\log|c_\lambda(t)|$ real-linearly to $K(t):N_{\R}\to N_{\R}$. The two generator values are holomorphic and nowhere zero on a smaller closed disc, so $K(t)$ is uniformly bounded there.
For $0<|t|<\epsilon$ put $L=\log|t|$ and
\begin{equation}
                 q=\frac{\log|z|}{L}.
\hypertarget{tag-012U}{}\label{cdp:eq:C.TG.13}
\end{equation}
Then \eqref{cdp:eq:C.TG.12} sends $q$ to
$q+(I+K(t)/L)\lambda$.
Choose $\epsilon<1$ so small that $I+K(t)/L$
and its inverse are uniformly bounded and
invertible for $0<|t|<\epsilon$.
The action on every nonzero fibre is
therefore free and properly discontinuous;
its real logarithmic quotient has a compact
fundamental parallelogram. Its angular
directions already form a compact two-torus.
Taking holomorphic logarithms of the two generator values gives \eqref{cdp:eq:C.TG.2} with $C$ replaced by $C(t)$. Its four real periods are independent because $LQ+\operatorname{Re}C(t)$ is invertible by this same bound. For constant $c$, choose $C(t)=C$.

\item At $t=0$ a nonzero translation cannot fix
any orbit: it would preserve a bounded cell
of the triangulation, which is impossible.
The central action is thus free. To verify
proper discontinuity also at the central
fibre, use a finite collection of fixed
toric charts about a fixed compact set.
For points in those charts tending to zero,
the coordinates in \eqref{cdp:eq:C.TG.11} show that $q$
stays bounded: each $\log|x_a|/L$ is bounded
below by a quantity tending to zero and
the three coefficients sum to one.
If both a point and its $g_\lambda$-translate
stay in such fixed compact sets, \eqref{cdp:eq:C.TG.13}
and the uniform inverse bound force
$\lambda$ to stay bounded. Only finitely
many lattice elements occur. The same
argument at central points uses the finite
sets of their translated cells. Hence the
quotient is Hausdorff and smooth.

\item Finally choose the representatives of $q$
in a uniformly bounded fundamental
parallelogram for $(I+K(t)/L)\Gamma$.
Only finitely many triangles of
$\mathscr P$ meet this bounded set.
In a triangle containing $q$, its
barycentric coordinates are precisely
\begin{equation}
             \ell_a=\frac{\log|x_a|}{L},
             \qquad \ell_a\ge0,\quad\sum_a\ell_a=1.
\hypertarget{tag-012V}{}\label{cdp:eq:C.TG.14}
\end{equation}
Consequently $|x_a|\le1$ in that chart.
Any sequence of fibres tending to zero
has, after its deck translations,
a subsequence in one of these finitely
many closed unit polydiscs, and hence a
convergent subsequence. The central fibre
itself is the finite union of the complete
toric star surfaces. This proves properness
over a smaller disc. The punctured fibres
are connected by their quotient
presentation, and the central dual torus
is connected as well.

\item The triangle-distinctness condition in
the statement prevents a component from
meeting a translate of itself. The
quotient therefore retains smooth embedded
components and the full reduced
simple-NC equation in \eqref{cdp:eq:C.TG.11}.
The differential
\begin{equation}
                  dt\wedge d\log z_1\wedge d\log z_2
\hypertarget{tag-012W}{}\label{cdp:eq:C.TG.15}
\end{equation}
is invariant under \eqref{cdp:eq:C.TG.12}, because every additional term, including $d\log c_\lambda(t)$, is a multiple of $dt$.
In each unimodular chart \eqref{cdp:eq:C.TG.11} it is,
up to the basis determinant $\pm1$,
$dx_0\wedge dx_1\wedge dx_2$.
It extends nowhere vanishingly across
the central fibre. Thus the constructed
total space has trivial canonical line.
The exact central-space comparison and \cref{cdp:st:C.2.UN.68} give the Euler number $2[N:\Gamma]$.

\end{enumerate}
\end{proof}






\subsection{Compact occurrence of toric fibres}\label{subsec:II:compact-toric-occurrence}

\begin{proposition}[The prescribed toric fibre and its conjugate over the projective line]\hypertarget{tag-01EK}{}\label{cdp:st:C.2.UN.73}
\StatementRecord{cdp:st:C.2.UN.73}{proved}{01EK}{\cref{cdp:st:C.2.UN.71}, \eqref{cdp:eq:C.TX.4}, \eqref{cdp:eq:C.TX.8}, \eqref{cdp:eq:C.TX.10}, \eqref{cdp:eq:C.TX.12}}

Every marked central log space classified in \cref{cdp:st:C.2.UN.71} occurs
unchanged in a smooth compact complex threefold mapping
properly to $\PP^1$, with connected complex two-torus
fibres away from two values. Its two full singular fibres
are the prescribed reduced simple-NC log space and its
complex conjugate. Both neighbourhoods have a nowhere-zero
canonical form. If its lattice index is $d=[N:\Gamma]$,
each singular support has Euler number $2d$ and the
compact total space has Euler number $4d$.

The local base coordinates at the two values, and the
maps restoring the prescribed log sections, are part
of the construction below. 
\end{proposition}

\begin{proof}

\begin{enumerate}[label=\arabic*.,leftmargin=*]
\item \par\medskip\noindent\textbf{Initial periods and a positive symmetric part.}
Use the data $(N,\Gamma,\mathscr P,c)$ of \cref{cdp:st:C.2.UN.71}.
Choose $N=\Z^2$, write $\Gamma=Q\Z^2$ with
$Q\in\operatorname{Mat}_2(\Z)$ invertible over $\R$,
and choose the logarithms $C$ of its two torus
holonomies. In the vector coordinate
$(2\pi i)^{-1}Q^{-1}\log z$, the angular lattice is
\begin{equation}
        N_0=Q^{-1}\Z^2\supset\Z^2,
             \qquad b=(2\pi i)^{-1}Q^{-1}C.
\hypertarget{tag-013Y}{}\label{cdp:eq:C.TX.1}
\end{equation}
The two other periods of the local model are
$b_j+(\log u)/(2\pi i)e_j$. Write $b=U+iA$
with $U,A$ real matrices. Choose $R>0$ large enough
that the symmetric part $S$ of $A+RI$ is positive
definite, and put
\begin{equation}
 b'=b+iRI,
 \qquad \kappa=e^{-2\pi R}>0,
 \qquad A+RI=S+aJ,
 \qquad J=\begin{pmatrix}0&1\\-1&0\end{pmatrix}.
\hypertarget{tag-0144}{}\label{cdp:eq:C.TX.2}
\end{equation}
In particular
$\exp(2\pi iQb'_j)=\kappa^{Qe_j}c_{Qe_j}$.
The base-coordinate maps below undo this scale at both ends. The periods apply the real-linear map to the fixed shifted vectors $b'_j$; their logarithmic coefficients multiply outside that map.

\par\medskip\noindent\textbf{A holomorphic family on the entire punctured line.}
Fix any integer $k\ge1$. For each constant
vector $w\in\C^2$, let
\begin{equation}
 V_t(w)=w+t^kJ\overline w,
 \qquad B_t=I+t^kJ,
 \qquad 0\ne t\in\C.
\hypertarget{tag-0145}{}\label{cdp:eq:C.TX.3}
\end{equation}
This is the map $V_{t^k}$ of \cref{lem:II:logarithmic-period-family}, and is holomorphic in $t$ for every fixed $w$.
Define the period lattice by its generators
\begin{equation}
 \Lambda_t=V_t(N_0)
   +\sum_{j=1}^2\Z
       \left(V_t(b'_j)+\frac{\log t}{2\pi i}B_te_j\right).
\hypertarget{tag-0146}{}\label{cdp:eq:C.TX.4}
\end{equation}
Changing the logarithm by $2\pi i$ adds $B_te_j$,
which is an angular period because $\Z^2\subset N_0$.
Thus these lattices are globally single-valued
over $\C^*$, with the displayed integral change
of basis on overlaps of logarithm branches.

\item We prove their real rank is four everywhere.
It suffices to use $\Z^2$ in place of its finite
angular enlargement $N_0$. Subtract the real
angular combinations supplied by $U$ and by
$\operatorname{Re}(\log t/(2\pi i))$ from the
last two columns. Set
\begin{equation}
 x=\operatorname{Re}(t^k),\quad
 w=\operatorname{Im}(t^k),\quad
 r=|t|^k,\quad
 y=-\frac{\log|t|}{2\pi}.
\hypertarget{tag-0147}{}\label{cdp:eq:C.TX.5}
\end{equation}
With $A'=S+aJ=A+RI$, the resulting real
period matrix is exactly
\begin{equation}
 \begin{pmatrix}
 I+xJ & wJ(A'-yI)\\
 wJ   & (I-xJ)A'+y(I+xJ)
 \end{pmatrix}.
\hypertarget{tag-0148}{}\label{cdp:eq:C.TX.6}
\end{equation}
Its determinant, by the displayed block
multiplication, is
\begin{equation}
\hypertarget{tag-016E}{}\SsixEquationAliases{\label{cdp:eq:C.WF.18}}
 \begin{aligned}
 \Delta(t)={}&(1+r^2)^2\det S
       +(1-r^2)(1+r^2)y\operatorname{tr}S\\
       &+y^2(1-r^2)^2
       +\bigl((1+r^2)a+2xy\bigr)^2>0.
 \end{aligned}
\hypertarget{tag-0149}{}\label{cdp:eq:C.TX.7}
\end{equation}
Indeed $(1-r^2)y\ge0$, the first term is
strictly positive, and the other three
terms are nonnegative. The identity uses
$r^2=x^2+w^2$ and holds for every argument
of $t$; it also holds where $B_t$ fails
to be complex-invertible. Thus \eqref{cdp:eq:C.TX.4}
is a discrete rank-four real lattice
for every $t\in\C^*$.

\item The quotient of the holomorphic
vector bundle $\C^2\times\C^*$ by these
locally holomorphic period sections is
therefore a smooth family of connected
complex two-tori. On every compact
annulus, the real period matrices and
their inverses are uniformly bounded.
Their real fundamental parallelepipeds
give compact representatives. This
proves properness of this middle family
over the annulus, and hence over $\C^*$.

\par\medskip\noindent\textbf{The chart and log section at zero.}
Near zero $B_t$ is complex-invertible.
In the middle vector coordinate $v$ use
the exponential map
\begin{equation}
                   z=\exp(2\pi iQ B_t^{-1}v).
\hypertarget{tag-014A}{}\label{cdp:eq:C.TX.8}
\end{equation}
It quotients exactly the angular lattice
$B_tN_0$, since $QB_t^{-1}B_tN_0=\Z^2$.
The other two period translations act by
\begin{equation}
 z\longmapsto c^0_{Qe_j}(t)\,t^{Qe_j}z,
 \qquad
 c^0_{Qe_j}(t)=
       \exp\bigl(2\pi iQ B_t^{-1}V_t(b'_j)\bigr).
\hypertarget{tag-014B}{}\label{cdp:eq:C.TX.9}
\end{equation}
These are holomorphic nonzero vectors
near zero; extend them multiplicatively
to a homomorphism $c^0(t):\Gamma\to T_N$.
Their exact central values are
$c^0_\lambda(0)=\kappa^\lambda c_\lambda$.

\item Use the base coordinate $u=\kappa t$.
The same action is then
\begin{equation}
 z\longmapsto d^0_\lambda(u)u^\lambda z,
 \qquad
 d^0_\lambda(u)=\kappa^{-\lambda}
                    c^0_\lambda(u/\kappa),
 \qquad d^0_\lambda(0)=c_\lambda.
\hypertarget{tag-013Z}{}\label{cdp:eq:C.TX.10}
\end{equation}


\item Apply the holomorphic-unit construction of \cref{cdp:st:C.2.UN.71} with base $u$, the same $(N,\Gamma,\mathscr P)$ and $c_\lambda(u)=d^0_\lambda(u)$ from \eqref{cdp:eq:C.TX.10}. The values are holomorphic and nowhere zero, their multiplicative extension is a homomorphism in $\lambda$, and $d^0_\lambda(0)=c_\lambda$. The resulting full reduced SNC filling is smooth and proper. Equation \eqref{cdp:eq:C.TX.8} divides by exactly its angular lattice, and \eqref{cdp:eq:C.TX.9} identifies its remaining deck maps with the middle period family. It therefore gives the punctured-family biholomorphism, while \eqref{cdp:eq:C.TX.10} identifies the marked central log space in the coordinate $u$ with the prescribed one.

\par\medskip\noindent\textbf{The conjugate filling at infinity.}
Put $s=t^{-1}$ and change the vector
frame by $v_\infty=s^kv$. The angular
periods become
\begin{equation}
 B^\infty_sN_0,
       \qquad B^\infty_s=s^kI+J.
\hypertarget{tag-0140}{}\label{cdp:eq:C.TX.11}
\end{equation}
Reverse the two nonangular period
generators. Their resulting exact
expressions are
\begin{equation}
 -(s^kb'_j+J\overline{b'_j})
                +\frac{\log s}{2\pi i}B^\infty_se_j.
\hypertarget{tag-0141}{}\label{cdp:eq:C.TX.12}
\end{equation}
The integer generator reversal is
permitted because it changes neither
the subgroup nor the quotient.
The matrix $B^\infty_s$ is invertible
near $s=0$. The exponential coordinates
$z_\infty=\exp(2\pi iQ(B^\infty_s)^{-1}v_\infty)$
therefore give the exact deck constants
\begin{equation}
 c^\infty_{Qe_j}(s)=
  \exp\bigl(-2\pi iQ(B^\infty_s)^{-1}
                  (s^kb'_j+J\overline{b'_j})\bigr).
\hypertarget{tag-0142}{}\label{cdp:eq:C.TX.13}
\end{equation}
At zero this is
\begin{equation}
 c^\infty_\lambda(0)
       =\overline{c^0_\lambda(0)}
       =\kappa^\lambda\overline{c_\lambda}.
\hypertarget{tag-0143}{}\label{cdp:eq:C.TX.14}
\end{equation}
The sign in \eqref{cdp:eq:C.TX.12} is essential:
$\exp(-2\pi i\overline{Qb'})$
is the conjugate of $\exp(2\pi iQb')$.

\item Choose the base coordinate $v=\kappa s$
at infinity. Define
$d^\infty_\lambda(v)=\kappa^{-\lambda}
c^\infty_\lambda(v/\kappa)$.
Then $d^\infty_\lambda(0)=\overline{c_\lambda}$,
and the exact action is
$z_\infty\mapsto d^\infty_\lambda(v)v^\lambda z_\infty$.
The two base charts satisfy $uv=\kappa^2$;
they are holomorphic coordinates
at the two points of one projective line.
Apply the same holomorphic-unit clause of \cref{cdp:st:C.2.UN.71}, now with base $v$ and homomorphism $d^\infty(v)$. Its value at zero is $\overline c$. The fan and its monomial chart maps are defined over $\Z$, so complex conjugation changes precisely these gluing constants and the log section. This is the conjugate marked central log space in the stated coordinate $v$.

\par\medskip\noindent\textbf{Compact gluing and canonical neighbourhoods.}
The exponential maps \eqref{cdp:eq:C.TX.8} and
their infinity counterparts are
isomorphisms of the
punctured torus families, since
they first divide by precisely
the angular lattice and then
by the displayed remaining
period group. Glue the two
proper cone fillings to the
middle family along these
isomorphisms. Disjoint smaller
base discs and a compact
intermediate annulus show
directly that the glued total
space is Hausdorff and proper
over $\PP^1$. It is smooth,
and its fibres are connected.
Properness over the compact
base makes the total space
compact. There are no other
singular fibres, because the
middle family is smooth.

\item The canonical form in the holomorphic-unit construction of \cref{cdp:st:C.2.UN.71} is $du\wedge d\log z_1\wedge d\log z_2$ at zero, and $dv\wedge d\log(z_\infty)_1\wedge d\log(z_\infty)_2$ at infinity. Under the two exponential identifications these are the forms on the same end fillings just constructed. That proposition proves their invariance for the homomorphisms $d^0(u)$ and $d^\infty(v)$ and their nowhere-zero extension in every unimodular cone chart.
Thus both end neighbourhoods
are canonical-flat with the
full reduced simple-NC
fibres specified in the
statement. Finally Euler
additivity over the base,
with smooth torus Euler
number zero, gives
$e(Y)=2d+2d=4d$.


\end{enumerate}
\end{proof}




















\section{Canonical reduction and weighted descent}\label{sec:II:canonical-reduction-weighted-descent}

\subsection{Local modifications and overlap maps}\label{subsec:II:modification-records}

\begin{definition}[Modifications with specified overlap maps]\hypertarget{tag-P340}{}\label{def:II:overlap-compatible-modification}
\StatementRecord{def:II:overlap-compatible-modification}{proved}{P340}{\cite[\S4.10.8, pp.~980--982; \S4.14.5, p.~985]{Kulikov1977}; \cite[\S\S2,8, pp.~339,341]{Kulikov1981}}
An \emph{overlap-compatible modification record} consists of finite
analytic chart data, centres with their ideals, local blowup or
contraction maps, and overlap isomorphisms, all over the same base.
In each local replacement $X=U\cup V$, the modified chart $U'$
has an overlap $W'\simeq W=U\cap V$ identified by the restriction
of the specified composite of local maps; the resulting space is
$X'=U'\cup_W V$. The record includes these neighbourhoods and
maps. Its inverse reverses each arrow and each overlap isomorphism.
For degenerations over a disc the centres are central and the
maps restrict to isomorphisms over the punctured disc. The local
$\sigma$-constructions used here are those in the cited sections,
including the additional fibre blowups of \cite[\S8, p.~341]{Kulikov1981}.
\end{definition}

\subsection{Canonical-flat reduction}\label{subsec:II:canonical-flat-reduction}

\begin{proposition}[Canonical-flat cores in the Hopf-minimum case]\hypertarget{tag-01EM}{}\label{cdp:st:C.2.UN.75}
\StatementRecord{cdp:st:C.2.UN.75}{proved}{01EM}{\cref{def:II:overlap-compatible-modification,lem:II:canonical-bimeromorphic-direct-image,lem:II:nonuniruled-divisors,cdp:st:C.2.UN.17,cdp:st:C.2.UN.19,cdp:st:C.2.UN.64,cdp:st:C.2.UN.65,cdp:st:C.2.UN.68,cdp:st:C.2.UN.71,cdp:st:C.2.UN.73,prop:II:section-fibre-contraction}; \cite[Theorem~1$'$, p.~340; \S\S5--10, pp.~340--342]{Kulikov1981}; \cite[\S1, pp.~957--958; \S\S4.2--4.4, pp.~968--970; \S4.10.7--4.10.8, pp.~980--982; \S4.14.5, p.~985; \S4.17, p.~988]{Kulikov1977}}

Let $f:X\to\Delta$ be a proper degeneration of connected
complex two-tori, with smooth total threefold and full
reduced simple-normal-crossing central fibre $S=\sum_iV_i$.
Choose a generator $\Omega$ of $f_*K_X$ near zero and write
\begin{equation}
             \operatorname{div}(\Omega)=\sum_i k_iV_i,
                  \qquad k_i\ge0,\quad \min_i k_i=0.
\hypertarget{tag-00P8}{}\label{cdp:eq:C.KG.1}
\end{equation}
Let $\mathcal H$ be the set of indices whose component has
a Hopf surface as its smooth minimal model. The following
two conditions are equivalent:
\begin{enumerate}
\item $k_i=0$ for every $i\in\mathcal H$.
\item There is a finite overlap-compatible modification
record from $f$ to a smooth proper torus degeneration
$f_0:X_0\to\Delta$ whose full central fibre $S_0$ is
reduced simple NC and whose canonical line is trivial.
\end{enumerate}
The constructed comparison is a bimeromorphic map
over the same disc and induces an isomorphism
\begin{equation}
                X\setminus S\ \simeq\ X_0\setminus S_0
                         \quad\hbox{over }\Delta^*.
\hypertarget{tag-00PA}{}\label{cdp:eq:C.KG.2}
\end{equation}
Intermediate charts retain their normalization branches;
in particular a type-II move may temporarily give a
self-intersection of a global component.

The endpoint has precisely one of the following forms,
with the markings and line maps of the indicated
statements:
\begin{enumerate}
\item a smooth two-torus;
\item a cycle of the marked diagonal Hopf and degree-zero
elliptic ruled pairs in \ref{cdp:st:C.2.UN.65}, with dual edge normals and
the complement-lattice product $M=I$;
\item a positive toric central log space in \ref{cdp:st:C.2.UN.71}, given by
$(N,\Gamma,\mathscr P,c)$, with Euler number
$2[N:\Gamma]$ and the specified full toric component pairs.
\end{enumerate}
Every degeneration in the stated branch is recovered by
reversing its finite record on its endpoint
neighbourhood. Conversely every such reverse record with
a smooth full reduced simple-NC final fibre gives a
degeneration in this branch. These are inverse analytic
constructions, retaining every centre, overlap and normal
extension.

The Euler changes are defined as follows.
Orient a recorded blowup arrow positively from its target
to its blowup, and negatively in the reverse direction.
For its centre $L$ and ideal $I$, put
\begin{equation}
 \begin{split}
  E_I&=\operatorname{Proj}_L\left(\bigoplus_{a\ge0}
                                      I^a/I^{a+1}\right),\\
  \delta(I)&=e_c(E_I)-e_c(L)
       =\int_L\bigl(e((E_I)_x)-1\bigr)\,de_c(x).
 \end{split}
\hypertarget{tag-00PB}{}\label{cdp:eq:C.KG.3}
\end{equation}
Local terms are evaluated on a finite constructible
subdivision subordinate to the retained charts. Restored
overlaps cancel. If $\epsilon_\nu\in\{1,-1\}$ records
the directions in the forward reduction to $S_0$, then
\begin{equation}
             e(S)=e(S_0)-\sum_\nu\epsilon_\nu\delta(I_\nu).
\hypertarget{tag-00PC}{}\label{cdp:eq:C.KG.4}
\end{equation}
For an ordinary smooth centre in a smooth threefold,
$\delta=2$ for a point and $\delta=2-2g(L)$ for a compact
curve. Both type-I and type-II moves have total defect
zero. An ordinary ruled-section contraction in \ref{cdp:st:C.2.UN.74} has
Euler change $-(2-2g(B))$ and removes its specified number
of triple points. Thus the exact zero-Euler criterion is
\begin{equation}
 \sum_\nu\epsilon_\nu\delta(I_\nu)=
 \begin{cases}
 0,&S_0\text{ is of type 1 or 2},\\
 2[N:\Gamma],&S_0\text{ is of type 3}.
 \end{cases}
\hypertarget{tag-00PD}{}\label{cdp:eq:C.KG.5}
\end{equation}
The same normal-cone formula applies to every auxiliary
singular-centre chart.

For a given compact total space over a compact curve,
the reduction and its inverse take place on that same
compact base, by the isomorphism \eqref{cdp:eq:C.KG.2} on an
annulus. A reverse record on a compact core
realization of \ref{cdp:st:C.2.UN.65} or \ref{cdp:st:C.2.UN.73} likewise gives a compact
realization. Its total Euler change is the same as
the central-fibre change in \eqref{cdp:eq:C.KG.4}.
\end{proposition}

\begin{proof}
\begin{enumerate}[label=\arabic*.,leftmargin=*]
\item 
\par\medskip\noindent\textbf{The canonical generator and its normalized orders.}
Apply \cref{cdp:st:C.2.UN.17} to this proper torus family.
Its canonical evaluation gives the vertical effective divisor in
\eqref{cdp:eq:C.KG.1}. Since the full fibre is reduced, positivity
of every $k_i$ would make $\Omega/f$ holomorphic on the smooth
total space. That would contradict primitivity of the generator,
so $\min_i k_i=0$. Multiplication by a holomorphic base unit
does not change these orders.

\par\medskip\noindent\textbf{The composite graphs and canonical pullback.}
In \cite[\S4.10.8, pp.~980--982]{Kulikov1977}, the finite
contractions in $U$ and inverse-order restorations produce $U'$
whose overlap $W'$ is identified with $W$ by their composite.
Section~4.14.5, p.~985, gives the same comparison at the local
singular charts. The additional fibre blowups in
\cite[\S8, p.~341]{Kulikov1981} leave this overlap unchanged.
These are the data of \cref{def:II:overlap-compatible-modification}.

For a finite $\sigma$-chain, take the closure of the common
biholomorphic graph in the proper image of the fibre product of
its arrow graphs. Both projections are proper. Over $W$ it is
the graph of the specified overlap isomorphism. Glue it to the
identity graph over $V$. The two graphs agree on their overlap,
and properness is local on each projection target; thus the
glued graph is proper over $U\cup V$ and $U'\cup_W V$.
It defines mutually inverse bimeromorphic maps. Iteration gives
\eqref{cdp:eq:C.KG.2}, since all maps are over the disc and are
isomorphisms away from the central fibre. The open covers,
normalization branches and any singular restoration charts
remain part of this finite sequence.

For its smooth endpoints apply
\cref{lem:II:canonical-bimeromorphic-direct-image} with this
glued map as $\phi$. Its pullback \eqref{cdp:eq:C.KG.6} identifies
the primitive canonical generators up to a base unit, and
preserves the orders of all surviving prime divisors.

\par\medskip\noindent\textbf{Preservation of the Hopf components.}
A map $\PP^1\to H$ to a Hopf surface lifts to
$\C^2\setminus\{0\}$ and is constant. Thus every rational curve
on a point blowup of $H$ is one of the finitely many curves
exceptional over $H$, and these surfaces are not uniruled.
Apply \cref{lem:II:nonuniruled-divisors} to the resulting proper
bimeromorphic comparison between the smooth endpoints $X$ and $X_0$
and these compact component
divisors. It preserves their strict-transform primes and their
Hopf birational classes. Its generic-point comparison is exactly
the one used in \eqref{cdp:eq:C.KG.6}; hence their canonical
orders are unchanged. At a canonical-flat endpoint a nowhere-zero
three-form generates $f_{0*}K_{X_0}$, so every order is zero.
This proves $k_i=0$ for every $i\in\mathcal H$, including all
point-blown-up Hopf components.

\par\medskip\noindent\textbf{The finite reduction and its local maps.}
By \cref{cdp:st:C.2.UN.19}, every class-VII minimal model
of a component of this torus degeneration is Hopf. Hence every
component to which condition~b) of \cite[Theorem~1$'$, p.~340]{Kulikov1981}
applies has canonical coefficient zero. With the source's map
$\pi:X\to\Delta$ equal to $f$ and its coefficients $r_i=k_i$,
that theorem gives a canonical-flat modification. Its good
degenerations are the proper smooth-total analytic families with
reduced normal-crossing fibre and smooth components defined in
\cite[\S1, pp.~957--958]{Kulikov1977}; the general torus has
trivial canonical line. These are precisely the hypotheses here.

We specify the comparison maps in that proof.
For a nonboundary $(-1)$-curve $A\subset V_i$
meeting one double curve transversely once
and avoiding the remaining boundary, the
principal-fibre normal sequence is
\begin{equation}
 0\longrightarrow\OO_{\PP^1}(-1)
 \longrightarrow N_{A/X}
 \longrightarrow\OO_{\PP^1}(-1)
 \longrightarrow0.
\hypertarget{tag-00PF}{}\label{cdp:eq:C.KG.7}
\end{equation}
It splits because $H^1(\PP^1,\OO)=0$.
Blow up $A$; the exceptional divisor is
$\PP^1\times\PP^1$ with normal
$\OO(-1,-1)$. Contract its other ruling
by \cite[\S2, Theorem~2$'$, pp.~640--642]{FujikiNakano1971}.
The source data are $(\widetilde X,S,M,r)=(\operatorname{Bl}_A X,
\PP^1\times\PP^1,\PP^1,2)$, with the second projection as
the ruling; the normal line has degree $-1$ on each of its fibres. This transfers the point blowup to
the neighbour. The exact adjunction equation
gives $k_i=k_j$ for these two components.
This is the type-I map.

For a double curve which is $(-1)$ on both sides,
the normal bundle is again $\OO(-1)\oplus\OO(-1)$.
On its ordinary blowup $\widetilde X$, the exceptional divisor
$E=\PP^1\times\PP^1$ has normal $\OO(-1,-1)$.
Apply the same normal-line contraction theorem with
$(\widetilde X,S,M,r)=(\widetilde X,E,\PP^1,2)$ and the
second projection as the ruling. It gives the ordinary smooth
curve contraction for these exact normal data. Its two
horizontal boundary sections are the two opposite branches
through the triple points.
The principal-fibre line has exactly two
triple points on that curve. If the opposite
components there have orders $k_3,k_4$,
adjunction gives
\begin{equation}
                         k_1+k_2=k_3+k_4.
\hypertarget{tag-00PG}{}\label{cdp:eq:C.KG.8}
\end{equation}
The ordinary blowup has exceptional full
multiplicity two; the other contraction
removes the old double curve and introduces
the new one between the two opposite
branches. If those branches belong to one
global component, both normalization branches
are retained, as in
\cite[\S4.4, p.~970]{Kulikov1977}. This intermediate
fibre has the two indicated normalization branches over its
self-double curve.

The other contraction charts are the ruled section-and-fibre maps of \ref{cdp:st:C.2.UN.74} and
the plane-line map of \ref{cdp:st:C.2.UN.64}. For a reducible
ruling fibre, the local contractions and
inverse-order restorations in
\cite[\S4.10.7--4.10.8, pp.~980--982; \S4.14.5, p.~985]{Kulikov1977}
give the retained finite chart record.
The construction in
\cite[\S8, p.~341]{Kulikov1981} also blows up the
specified missing fibre components before
contracting the chosen ruled surface.
The new ruled components have canonical
order $k+1$ when the removed component has
order $k$. For a maximum-order component
no such additions are needed.

For termination, apply the finite-level argument in
\cite[\S9, pp.~341--342]{Kulikov1981} with minimum zero
and increment one. The following comparison specifies its levels. Let $R=\max k_i$.
Every removed old component has level
$k\in\{1,\ldots,R\}$. It produces only
finitely many new components, all at level
$k+1$, and produces none when $k=R$.
There are finitely many level-one removals;
therefore finitely many level-two objects
are ever produced. Continuing through the
fixed finite set of levels proves that all
removals terminate. The components left
have order zero. The endpoint
argument in \cite[\S10, p.~342]{Kulikov1981}, using
\cite[\S4.17, p.~988]{Kulikov1977}, gives a smooth
total space and smooth reduced components.
In particular \eqref{cdp:eq:C.KG.8} prevents a minimum
component from acquiring a self-incidence
in a type-II move whose removed edge has
an order above the minimum. The endpoint
generator is therefore nowhere zero.

\par\medskip\noindent\textbf{Endpoint classification and the converse record.}
Apply \ref{cdp:st:C.2.UN.65} to the endpoint.
Its zero-Euler alternatives give types~1 and~2, including all the marked line
and integral-lattice comparisons.
Its positive case is \ref{cdp:st:C.2.UN.68}; \ref{cdp:st:C.2.UN.71}
identifies its central log
space with the specified toric data
by the displayed chart and residue
maps. This proves the stated complete
endpoint list.

Reverse each recorded arrow and each
overlap isomorphism. They reconstruct
the original neighbourhood literally.
Starting instead with any canonical-flat endpoint and such a
reverse record, all off-central maps
remain isomorphisms; the general
fibres are the same connected tori.
For the resulting smooth total space, apply
\cref{lem:II:nonuniruled-divisors} to the inverse proper
bimeromorphic map and a compact Hopf-birational component.
It corresponds to a prime at the canonical-flat endpoint.
The order comparison of \cref{lem:II:canonical-bimeromorphic-direct-image}
therefore makes its canonical order zero. This proves both
directions of the generation statement.
The endpoint is the analytic family with precisely the
neighbourhood and overlap maps used in this construction.

\par\medskip\noindent\textbf{The exact Euler defect, including singular centres.}
For a blowup with centre ideal
$I$, its exceptional inverse image is
the projectivized normal cone in
\eqref{cdp:eq:C.KG.3}. Outside $L$ it is an
isomorphism. Since $L$ is central,
the reduced central inverse image
replaces precisely $L$ by $E_I$.
Additivity of compactly supported
Euler characteristic proves \eqref{cdp:eq:C.KG.3}
for this arrow, including its
constructible-fibre integral.
For these computations refine the
retained analytic charts by relatively
compact subanalytic coordinate
neighbourhoods, with closures inside
the original charts. Compactness of
the central fibre gives a finite
subcover. Subdivide its overlaps and
centres compatibly, and use inverse
images under the retained proper
graphs; all pieces have finite
subanalytic triangulations. This
refinement also defines the local
terms when a displayed centre is an
open part of a curve. All Euler terms are evaluated on this finite
subanalytic subdivision.

In one replacement, evaluated on this
finite-type refinement,
$U\cup V\rightsquigarrow
U'\cup_WV$, inclusion--exclusion gives
\begin{equation}
 \begin{split}
 e(S')-e(S)
   &=e_c(S'\cap U')-e_c(S\cap U)\\
   &=\sum_{\nu\text{ in this chart}}
                         \epsilon_\nu\delta(I_\nu),
 \end{split}
\hypertarget{tag-00PH}{}\label{cdp:eq:C.KG.9}
\end{equation}
because the unchanged $V$ and the
isomorphic restored overlaps cancel.
Equivalently one applies this equality
to the finite subdivision just
described and sums over its disjoint strata. The local
arrow contributions telescope, and
then the global replacements
telescope, proving \eqref{cdp:eq:C.KG.4}.

For a smooth centre of dimension $a$
in a smooth threefold the normal
cone is its rank-$(3-a)$ normal
bundle, so its exceptional fibre
is $\PP^{2-a}$. Thus
\begin{equation}
                \delta(I)=(2-a)e_c(L).
\hypertarget{tag-00P9}{}\label{cdp:eq:C.KG.10}
\end{equation}
For type~I and type~II the common
blowup replaces a compact $\PP^1$
by $\PP^1\times\PP^1$ from either
side, so the two contributions
$+2$ and $-2$ cancel. The ruled
and plane changes follow from
\ref{cdp:st:C.2.UN.74} and \ref{cdp:st:C.2.UN.64}. In an auxiliary
singular-centre chart \eqref{cdp:eq:C.KG.3}
remains the formula: for example,
the blowup of the three coordinate
axes in $\C^3$, with ideal
$(xy,xz,yz)$, has fibre $\PP^1$
over a general point of each axis
and fibre $\PP^2$ over their
common origin. The fibre algebra at the origin is the polynomial algebra
in the three degree-one classes of $xy,xz,yz$, so its
projectivization is $\PP^2$ and its local defect is two.
Substitution of the three endpoint
Euler numbers proves \eqref{cdp:eq:C.KG.5}.

\par\medskip\noindent\textbf{Compact scope.}
Restrict the given compact family
to a small disc at the chosen
value. On a smaller annulus use
the isomorphism \eqref{cdp:eq:C.KG.2}
to glue the new proper filling
to the unchanged complement.
The two proper pieces give a
Hausdorff proper analytic family;
their smooth charts coincide on
the overlap. This produces the
claimed compact replacement,
and the inverse record recovers
the original compact family.
The same construction starts
from the explicit compact cores
of \ref{cdp:st:C.2.UN.65} and \ref{cdp:st:C.2.UN.73} and any reverse
record on their tubes.
Euler additivity cancels the
unchanged complement and annulus,
so the total-space Euler change
is precisely the central one.

\end{enumerate}
\end{proof}

\subsection{Cyclic roots, canonical orders and Euler corrections}\label{subsec:II:cyclic-root-normalization}

\begin{proposition}[Normalized cyclic roots of torus degenerations]\hypertarget{tag-P342}{}\label{prop:II:cyclic-root-normalization}
\StatementRecord{prop:II:cyclic-root-normalization}{proved}{P342}{\cref{lem:II:canonical-bimeromorphic-direct-image,lem:II:blowup-topology,cdp:st:C.2.UN.17}; \cite[arXiv v2, \S0.2(2), pp.~3--4; Theorem~1.6 and Remarks~1.7, pp.~5--6; Theorem~1.10, p.~7]{BierstoneMilman1997}}
Let $f:X\to\Delta_t$ be a proper degeneration of connected complex
two-tori from a smooth threefold, with arbitrary full central fibre.
Choose a generator $\Omega$ of $f_*K_X$ and an embedded resolution
$\pi:Y\to X$ by smooth central point and curve blowups, so that
$\operatorname{div}(f\pi)=\sum_i m_iD_i$ has SNC support and
$\operatorname{div}(\pi^*\Omega)=\sum_i k_iD_i$.
Put $r_i=(k_i+1)/m_i$, and choose an integer $N$ divisible by all $m_i$.
There is a $G$-equivariant normalization and semistable resolution
\begin{equation}
\hypertarget{tag-00A1}{}\hypertarget{tag-016T}{}\SsixEquationAliases{\label{cdp:eq:C.AR.10}\label{cdp:eq:C.WF.6}}
 \begin{gathered}
 Y_N=\operatorname{Norm}(Y\times_{\Delta_t}\Delta_s),
                \qquad t=s^N,\\
 Z\xrightarrow{\ q\ }Y_N\xrightarrow{\ \rho\ }Y,
                 \qquad G=\mu_N,\qquad Y_N/G=Y.
 \end{gathered}
\hypertarget{tag-0101}{}\label{cdp:eq:C.RK.2}
\end{equation}
in which $Z$ is smooth with connected full reduced SNC central fibre.
The map $q$ is a composition of normalized monomial-ideal blowups,
with height-one subdivisions and their coarsening morphisms.
Every component added by $q$ is ruled or rational. The old component
and stratum labels, local unit-root choices and all transition maps
are part of these maps. The quotient in the diagram is the finite
analytic quotient.

For a primitive height-one ray $(v,1)$ the pulled canonical form has
order $\sum_i(k_i+1)v_i-1$. Its minimum is
$a_{\min}=N\min_i r_i-1$, and
$\Omega_Z=s^{-a_{\min}}q^*\rho^*\pi^*\Omega$ generates the canonical
direct image of $Z\to\Delta_s$. For
$\gamma:s\mapsto\zeta s$, its character is
$\gamma^*\Omega_Z=\zeta^{-a_{\min}}\Omega_Z$.

For each nonempty stratum define
\begin{equation}
\hypertarget{tag-0109}{}\hypertarget{tag-00A8}{}\hypertarget{tag-00A9}{}\SsixEquationAliases{\label{cdp:eq:C.RK.7}\label{cdp:eq:C.AR.2}\label{cdp:eq:C.AR.3}}
\begin{gathered}
D_I^\circ=\left(\bigcap_{i\in I}D_i\right)
          \setminus\bigcup_{j\notin I}D_j,\qquad
d_I=\gcd_{i\in I}m_i,\qquad
\nu_I=\frac{N^{|I|-1}d_I}{\prod_{i\in I}m_i},\\
R_N=\sum_I(d_I\nu_I-1)e_c(D_I^\circ),\qquad
\Delta_\pi=2p_\pi+\sum_\beta(2-2g(B_\beta)).
\end{gathered}
\hypertarget{tag-0108}{}\label{cdp:eq:C.RK.6}
\end{equation}
where $p_\pi$ counts the point centres and the $B_\beta$ are the
connected compact curve centres of $\pi$. Then $1\le |I|\le3$,
$\nu_I$ is a positive integer, and the central cover of $D_I^\circ$
in $Y_N$ is finite \'etale of degree $d_I$. The fibre of $q$ over
each point of that cover has Euler number $\nu_I$. With
$S=f^{-1}(0)_{\rm red}$, $S_Y=(f\pi)^{-1}(0)_{\rm red}$ and
$S_Z=(Z_0)_{\rm red}$, one has
$e(S_Z)-e(S_Y)=R_N$ and $e(S_Y)-e(S)=\Delta_\pi$.

Write the original full divisor and generator orders as
\begin{equation}
        \operatorname{div}(f)=\sum_\alpha n_\alpha S_\alpha,
        \qquad \operatorname{div}(\Omega)
                       =\sum_\alpha a_\alpha S_\alpha.
 \hypertarget{tag-00AC}{}\label{cdp:eq:C.AR.6}
\end{equation}
On the normalization $T_\alpha$ of $S_\alpha$, the following finite
algebras glue and have a finite normal cover $C_\alpha\to T_\alpha$:
\begin{equation}
 \begin{gathered}
 f=g_j^{\,n_\alpha}u_j,\qquad
 w_j^{\,n_\alpha}=u_j|_{T_\alpha},\\
 w_k=(g_j/g_k)|_{T_\alpha}\,w_j .
 \end{gathered}
\hypertarget{tag-00ZZ}{}\label{cdp:eq:C.RK.18}
\end{equation}
Here $g_j$ is a local Cartier equation of $S_\alpha$ and
$u_j=f/g_j^{n_\alpha}$. The strict root components over $S_\alpha$
are proper birational resolutions of the connected components of
this cover, through $w_j=s^{N/n_\alpha}/g_j$. Every root component
with a Hopf minimal model is obtained from one of these original
component covers. The exceptional components of $\pi$ and their
finite normal covers are projective, and the additional components
of $q$ are ruled or rational.

The intrinsic density threshold
\begin{equation}
 \tau=\sup\left\{c\in\R:
       |f|^{-2c}|\Omega|^2\text{ is locally integrable
                              near }S\right\}.
\hypertarget{tag-0100}{}\label{cdp:eq:C.RK.19}
\end{equation}
satisfies $\tau=\min_i r_i$. The density $|\Omega|^2$ is the positive
top-degree density of the holomorphic three-form. The strict root
orders are $N(r_i-\tau)$, and the covers and threshold are independent
of the embedded resolution and the divisible exponent.

For a compact family over a compact curve, the same construction
can be made on a connected cyclic base cover ramified to order $N$
at the selected value and \'etale over all other critical values.
Its local root chart is equivariantly isomorphic to $t=s^N$.
The total root family is smooth away from the selected fibre,
where the same equivariant resolution applies.
\end{proposition}

\begin{proof}
\begin{enumerate}[label=\arabic*.,leftmargin=*]
\item \par\medskip\noindent\textbf{Embedded resolution and the exact normalized root charts.}
Apply analytic embedded resolution to the central support
in the smooth threefold: \cite[arXiv v2, \S0.2(2), pp.~3--4;
Theorem~1.6 and Remarks~1.7, pp.~5--6; Theorem~1.10, p.~7]{BierstoneMilman1997}.
The support is compact, so the relatively compact formulation of
Theorem~1.6 gives a finite sequence near this fibre. Its centres
lie over the central support and the maps are isomorphisms over
$\Delta_t^*$.
Divisor centres give identity blowups and are
omitted. The remaining centres in a smooth
threefold are points and compact curves.
The pullback of $\Omega$ is a canonical generator on $Y$:
apply \cref{lem:II:canonical-bimeromorphic-direct-image} with
$(X,X',\phi)=(Y,X,\pi)$ and the corresponding two maps to the disc.

\item Near a stratum with $r$ incident components,
absorb the unit in the fibre equation into
a coordinate. The exact local equation and
volume form have the shape
\begin{equation}
\hypertarget{tag-016R}{}\SsixEquationAliases{\label{cdp:eq:C.WF.4}}
 t=x_1^{m_1}\cdots x_r^{m_r},\qquad
 \pi^*\Omega=u\,x_1^{k_1}\cdots x_r^{k_r}
                         dx_1\wedge dx_2\wedge dx_3,
                         \quad u\text{ a unit}.
\hypertarget{tag-010B}{}\label{cdp:eq:C.RK.9}
\end{equation}
The normalization of $t=s^N$ has $d_I$
local branches. Each is the affine toric
chart with cocharacter lattice
\begin{equation}
\hypertarget{tag-016S}{}\hypertarget{tag-00A2}{}\SsixEquationAliases{\label{cdp:eq:C.WF.5}\label{cdp:eq:C.AR.11}}
 \mathcal L_I=\left\{(v,h)\in\Z^r\oplus\Z:
                         \sum_i m_iv_i=Nh\right\},\qquad
 \sigma_I=\{(v,h)\in(\mathcal L_I)_{\R}:v_i\ge0\}
\hypertarget{tag-00ZR}{}\label{cdp:eq:C.RK.10}
\end{equation}
and cone generated by the height-one rays
$((N/m_i)e_i,1)$. The omitted $3-r$
coordinates are unchanged. Local roots of
the unit, if not absorbed in \eqref{cdp:eq:C.RK.9},
give the same normalization maps and
the finite transition characters.
The group $G$ permutes the $d_I$ branches
and acts torically on each stabilised
branch, over $s\mapsto\zeta s$.

\item The height-one slice has dimension at most
two. Subdivide its edges at every lattice
point, and triangulate each two-dimensional
lattice triangle using all its lattice
points. An empty lattice triangle is
unimodular, by Pick's formula. Thus these
subdivisions give smooth toric charts.
They agree on their subdivided edges.
The refinement can be obtained by
successively inserting the height-one
lattice points as star subdivisions.
To describe one step, write the ray in a simplicial cone
as $v=\sum_i c_i v_i$, with $c_i\ge0$, and let $\ell_i$
be the rational dual basis to its primitive generators $v_i$.
Choose $L>0$ such that every $L\ell_i/c_i$ with $c_i>0$
is an integral character. Blow up the monomial ideal
$(\chi^{L\ell_i/c_i}:c_i>0)$ and normalize.
Its $i$-th chart adjoins all ratios
$\chi^{L\ell_j/c_j}/\chi^{L\ell_i/c_i}$ and then saturates
the character semigroup. Its cone is the region
$\langle\ell_i,w\rangle/c_i\le\langle\ell_j,w\rangle/c_j$
for all $j$ with $c_j>0$, namely the cone spanned by $v$
and the old rays other than $v_i$. These are exactly the
star-subdivision charts and their coarsening maps. On an
overlap, replacing a boundary equation
by a holomorphic unit times that
equation multiplies every monomial
generator by a unit and leaves this
ideal unchanged. Thus the normalized
blowup maps glue through the original
tangent-dependent holomorphic transition units.
The stabilizer of an old coarse cone fixes
each old ray: the original SNC components
are individually labelled, and the group
acts over their identity map in $Y$.
Its action on this lattice slice is
therefore the identity. Choose the
triangulations orbitwise. The toroidal transition maps, which preserve
these labelled rays, glue the subdivisions
to a proper $G$-equivariant map $q$.

\item Every new ray still has height one,
so $s$ has order one on every central
component of $Z$. The resulting toric
boundary is simple NC with smooth
components. The atlas includes every original component and stratum
label together with each subdivision ray.
Transition units fix those labels and
the ray vector; distinct rays in one
old cone are never identified. The
local components in a unimodular
chart therefore belong to distinct
global components, and their smooth
star charts glue as embedded global components. Over a double curve the
added surfaces are projective-line
bundles over finite covers of that
curve, with their indicated modifications;
over a triple point they are rational
toric surfaces. Their minimal models are ruled or rational. The normalized fibre product
over $\Delta_s^*$ is connected because
its compact torus fibres and its base
are connected. The proper Stein factor
has degree one there and hence on the
normal disc. Thus $Z\to\Delta_s$ has
connected full central fibre as well.

\par\medskip\noindent\textbf{Canonical ratios and the explicit minimum test.}
At a primitive height-one ray $(v,1)$,
the toric Jacobian of \eqref{cdp:eq:C.RK.9} gives
\begin{equation}
 \operatorname{ord}_{(v,1)}(q^*\rho^*\pi^*\Omega)
                 =\sum_i(k_i+1)v_i-1.
\hypertarget{tag-00ZS}{}\label{cdp:eq:C.RK.11}
\end{equation}
For a strict ray this is
$N(k_i+1)/m_i-1$. Any height-one
point is a convex combination of the
old vertices, so \eqref{cdp:eq:C.RK.11} is the
same affine combination of their
orders. The minimum among all rays is
therefore
\begin{equation}
                  a_{\min}=N\min_i r_i-1.
\hypertarget{tag-00ZT}{}\label{cdp:eq:C.RK.12}
\end{equation}
It is an integer, since a minimum is
attained at an old vertex. The form
\begin{equation}
            \Omega_Z=s^{-a_{\min}}
                             q^*\rho^*\pi^*\Omega
\hypertarget{tag-00ZU}{}\label{cdp:eq:C.RK.13}
\end{equation}
is holomorphic and has minimum component order zero.
By \cref{cdp:st:C.2.UN.17}, the canonical direct image of
this proper torus family is free of rank one. Relative to
its frame the displayed form is a holomorphic base function.
A zero of that function at the origin would add at least one
to every component order, since the full central fibre is
reduced. The coefficient is therefore a unit, proving that
$\Omega_Z$ is a generator. For the deck generator the resulting character is
$\gamma^*\Omega_Z=\zeta^{-a_{\min}}\Omega_Z$.

\item \par\medskip\noindent\textbf{The finite analytic quotient.}
The quotient of the normal root is $Y$.
Indeed on a
normal local ring $A$ of $Y$,
the root ring is the integral
closure $B$ of
$A[s]/(s^N-f\pi)$. Put $K=\operatorname{Frac}(A)$. The total fraction algebra
is $K[s]/(s^N-f\pi)$; it is reduced and has basis
$1,s,\ldots,s^{N-1}$, on which $G$ acts with the distinct
characters $1,\zeta,\ldots,\zeta^{N-1}$. Its invariant
subalgebra is therefore exactly $K$.
An invariant element of $B$ is
integral over $A$ and lies in
that fraction field, hence lies
in $A$. Thus
\begin{equation}
                              B^G=A.
\hypertarget{tag-00ZV}{}\label{cdp:eq:C.RK.14}
\end{equation}
The same fraction algebra includes every local split branch;
the group permutes them, and the displayed invariant equality
is the local finite-quotient map. 
\item \par\medskip\noindent\textbf{The root correction from stratum fibres.}
Over $D_I^\circ$ the reduced
central support of $Y_N$ is a finite \'etale cover
of degree $d_I$. The local root
equation in \eqref{cdp:eq:C.RK.9} proves the
degree; unit-root transition
maps glue the cover and may
permute its sheets. At a point
of this cover the toric fibre
of $q$ has Euler number $\nu_I$.
To see this, use its torus-orbit
stratification. All positive-
dimensional torus orbits have
Euler number zero. The zero-
dimensional orbits are exactly
the maximal unimodular cones
in the subdivision, so their
number is the normalized
volume of the height-one
simplex.

\item Here is its exact lattice
volume. Set $a_i=m_i/d_I$ and
choose an integral vector $b$
with $\sum_i a_i b_i=1$.
The difference lattice of
the slice is
$\ker(a:\Z^r\to\Z)$.
Appending $b$ to a basis of
this kernel gives a unimodular
basis of $\Z^r$. Consequently
the normalized volume is the
absolute determinant obtained
by appending $b$ to the $r-1$
vertex differences. Expansion
gives
\begin{equation}
 \left|\det\left[
 \frac{N}{m_1}e_1-\frac{N}{m_r}e_r,
 \ldots,
 \frac{N}{m_{r-1}}e_{r-1}-\frac{N}{m_r}e_r,
 b\right]\right|
               =\frac{N^{r-1}d_I}{\prod_i m_i}.
\hypertarget{tag-00ZW}{}\label{cdp:eq:C.RK.15}
\end{equation}
For $r=1$ both sides are one.
This determinant computes $\nu_I$ for every unimodular
triangulation and proves its integrality.

\item Proper stratified Euler
pushforward now gives
\begin{equation}
\hypertarget{tag-00A5}{}\SsixEquationAliases{\label{cdp:eq:C.AR.14}}
 \begin{split}
 e(S_Z)&=\sum_I d_I\nu_I e_c(D_I^\circ),\\
 e(S_Y)&=\sum_I e_c(D_I^\circ),\\
 e(S_Z)-e(S_Y)&=R_N.
 \end{split}
\hypertarget{tag-00ZX}{}\label{cdp:eq:C.RK.16}
\end{equation}
This also checks the finite
quotient directly. A deck
element $\gamma^j$ fixes the
$d_I$-sheeted central cover
over this stratum exactly
when $d_I\mid j$. Therefore
\begin{equation}
 \begin{split}
 e(S_{Y_N}^{\gamma^j})
       &=\sum_{d_I\mid j}d_I e_c(D_I^\circ),\\
 \frac1N\sum_{j=0}^{N-1}
          e(S_{Y_N}^{\gamma^j})
       &=\sum_I e_c(D_I^\circ)=e(S_Y).
 \end{split}
\hypertarget{tag-00ZY}{}\label{cdp:eq:C.RK.17}
\end{equation}
All these equalities hold for
the covers even when
they are globally connected
and their sheet labels have
nontrivial monodromy.

\item For each arrow in $\pi$, apply \cref{lem:II:blowup-topology}
to its smooth compact centre and the compact reduced support
containing it. The full reduced inverse image is the next support.
The Euler increment is two for a point and $2-2g(B_\beta)$
for a connected smooth curve. Their sum is $\Delta_\pi$ in
\eqref{cdp:eq:C.RK.6}, proving $e(S_Y)-e(S)=\Delta_\pi$.

\item \par\medskip\noindent\textbf{Intrinsic covers and the intrinsic minimum.}
The component $S_\alpha$ is Cartier
in the original smooth threefold.
Its local equations satisfy
$g_k=v_{kj}g_j$ for holomorphic
units, and therefore
$u_k=v_{kj}^{-n_\alpha}u_j$.
This is precisely the transition
in \eqref{cdp:eq:C.RK.18}, including its cocycle
identity. The corresponding monic
finite algebras glue before taking
their finite normalizations.
On the open smooth component
away from other components this
is the $\mu_{n_\alpha}$-torsor
of roots of the normal-line
trivialization defined by $f$.
The local Cartier frames and their displayed cocycle define
this cover globally.

\item For any root exponent divisible
by $n_\alpha$, the exact map from
the strict central root cover is
\begin{equation}
                         w_j=\frac{s^{N/n_\alpha}}{g_j}.
\hypertarget{tag-0103}{}\label{cdp:eq:C.RK.21}
\end{equation}
At its generic point numerator
and denominator have the same
order, and their quotient satisfies
\eqref{cdp:eq:C.RK.18}. On a smooth strict
component of $Z$ this meromorphic
quotient is integral over its
holomorphic local ring, by that
monic equation. Normality makes
it holomorphic. The map therefore
extends to a proper map to the
glued finite cover, and factors
through its normalization. It
has generic degree one onto a
connected component of $C_\alpha$.
Every such component is obtained,
since on the dense open both
covers have the same complete
set of $n_\alpha$ roots.
Thus the strict surfaces are
birational resolutions
of these fixed covers for every
choice of embedded resolution
and divisible root exponent.

\item Every exceptional surface of
the original embedded resolution
is projective: initially it is
$\PP^2$ or a projective-line
bundle over a compact curve,
and its subsequent transforms
are projective modifications.
The ideals on these projective
surfaces are algebraic by GAGA.
A finite normal analytic cover
of a projective surface is also
projective: its finite coherent
algebra and multiplication
algebraize by GAGA, and the
finite pullback of an ample
line is ample. Its surface
resolutions remain bimeromorphic
to projective surfaces. Their minimal models are projective, whereas a Hopf surface
has odd first Betti number and hence is not K\"ahler;
the point-blowup invariance of $b_1$ in
\cref{lem:II:blowup-topology} excludes a Hopf minimal
model for these projective resolutions. The new
semistable components have
already been shown ruled or
rational. Hence \eqref{cdp:eq:C.RK.21}
accounts for every Hopf root component.

\item Finally apply the literal change
of variables under $\pi$ to the
density in \eqref{cdp:eq:C.RK.19}, outside
the exceptional sets of measure
zero. Near a resolved crossing
it is a bounded positive unit
times the product of the radial
densities
$|x_i|^{2(k_i-cm_i)}\,dA(x_i)$,
with smooth tangent density.
On finitely many relatively
compact charts about the compact
central fibre, the integral is
finite exactly when
$k_i-cm_i>-1$ for every $i$.
At equality it diverges already
at a general point of that
divisor. Thus
\begin{equation}
\hypertarget{tag-00A6}{}\SsixEquationAliases{\label{cdp:eq:C.AR.15}}
                             \tau=\min_i\frac{k_i+1}{m_i}.
\hypertarget{tag-0104}{}\label{cdp:eq:C.RK.22}
\end{equation}
This is the original density's
integrability threshold, so it
is independent of resolution.
The strict transform of an
original component retains
its original pair
$(a_\alpha,n_\alpha)$.
Together with \eqref{cdp:eq:C.RK.21} and the exceptional-component
comparison, this proves the stated choice independence. Subtracting
$a_{\min}=N\tau-1$ from the strict ray order gives
\begin{equation}
                  N\left(\frac{k_i+1}{m_i}-\tau\right).
 \hypertarget{tag-00A7}{}\label{cdp:eq:C.AR.16}
\end{equation}
with the original strict components retaining
$(k_i,m_i)=(a_\alpha,n_\alpha)$.

\item \par\medskip\noindent\textbf{A compact cyclic cover with the given local root.}
For an original compact base \(B\), let
\(c\) be the selected value and let
\(\Sigma\) be the finite critical-value set.
Choose a point \(p\notin\Sigma\cup\{c\}\).
For sufficiently large \(d\),
\[
 H^1\!\left(B,\OO_B(dp-2c-
                    \sum_{v\in\Sigma\setminus\{c\}}v)\right)=0.
\]
Indeed its degree is greater than \(2g(B)-2\),
so Serre duality gives the vanishing. The
restriction sequence consequently
surjects from \(H^0(B,\OO_B(dp))\) onto
the prescribed first jet at \(c\) and
values at the other critical points.
Choose the meromorphic function \(a\)
with jet \(a=t+O(t^2)\) in the fixed
coordinate at \(c\), and value one at
every other point of \(\Sigma\).
Its only possible pole is \(p\), and
all its other zeros and poles are regular
values of the original family.

\item Normalize the compact cyclic base cover
\(s^N=a\). The simple zero at \(c\)
makes it connected and totally ramified
there. It is \'etale over every other
critical value. Near its additional
branch values the original map is a
submersion, so its base change is locally
a product with the smooth normalized
cover curve and is smooth. Over the
other critical values the base change
is \'etale over the original smooth total
space. The normalized total space is
therefore smooth away from the selected
root fibre. Apply the preceding root resolution near that fibre.
Its coarsening is an isomorphism off the fibre and glues
to the identity on the unchanged annulus and complement.
The local function \(a\) is a coordinate
at \(c\). Write \(a=t u(t)\), with
\(u(0)=1\), and choose its analytic
\(N\)-th root near \(c\). The map
\[
 (t,s)\longmapsto
 \left(t,\sigma=\frac{s}{u(t)^{1/N}}\right),
 \qquad \sigma^N=t,
\]
is an equivariant isomorphism
between the global-cover root chart
and the local root model $t=\sigma^N$;
its inverse is \(s=\sigma u(t)^{1/N}\).

The normalization and resolution maps are proper and equivariant.
Their quotient and the original smooth-centre blowdowns recover
the given compact family through the same local invariant-ring map.
\end{enumerate}
\end{proof}

\subsection{Weighted toric descent and compact occurrence}\label{subsec:II:weighted-toric-descent}

\begin{definition}[Marked weighted toric support]\hypertarget{tag-P349}{}\label{def:II:marked-weighted-toric-support}
\StatementRecord{def:II:marked-weighted-toric-support}{proved}{P349}{}
The marked weighted data of a central fibre consist of its reduced
central analytic space, its component multiplicities, the central
log space of its normalized cyclic root, and the old cone cells
before their semistable subdivision. A coarsening in these data
is a proper toric morphism between the specified monoid charts.
The component gluing maps and the normalized-root/coarsening
maps are included in the datum.
\end{definition}

\begin{proposition}[Exact weighted toric descent and compact occurrence]\hypertarget{tag-01EN}{}\label{cdp:st:C.2.UN.76}
\StatementRecord{cdp:st:C.2.UN.76}{proved}{01EN}{\cref{def:II:marked-weighted-toric-support,prop:II:cyclic-root-normalization,cdp:st:C.2.UN.68,cdp:st:C.2.UN.71,cdp:st:C.2.UN.73}}

Let \(f:X\to\Delta_t\) be a proper degeneration of connected complex
two-tori, with smooth total space and full fibre
\(F=\sum_i m_iV_i\) whose reduced support is simple normal crossing.
Suppose there is a triple point and that the orders of a
generator of \(f_*K_X\) satisfy
\begin{equation}
 \operatorname{div}(\Omega)=\sum_i k_iV_i,\qquad
 \frac{k_i+1}{m_i}=\lambda\quad\hbox{for every }i.
\hypertarget{tag-0165}{}\label{cdp:eq:C.WF.1}
\end{equation}
Then the following conclusions hold.

The normalized cyclic root, followed by the explicit height-one
subdivision below, is a positive canonical-flat toric degeneration
of \cref{cdp:st:C.2.UN.71}.  The original marked weighted support is exactly its
coarse-cone quotient by the root group.  Thus this branch
is described by the toric data \((N,\Gamma,\mathscr P,c)\) of \cref{cdp:st:C.2.UN.71},
the periodic coarse triangles of the marking, and the cyclic
monomial maps \eqref{cdp:eq:C.WF.9}--\eqref{cdp:eq:C.WF.11}, with their smooth
coarse invariant charts \eqref{cdp:eq:C.WF.28}.

There are precisely two possible common ratios:
\begin{equation}
 \lambda=1,\quad k_i=m_i-1;
 \qquad\hbox{or}\qquad
 \lambda=\tfrac12,\quad k_i=\tfrac{m_i}{2}-1.
\hypertarget{tag-016G}{}\label{cdp:eq:C.WF.2}
\end{equation}
In the second case every \(m_i\) is even.  Every positive-dimensional
open stratum of the original support has Euler number zero.
Consequently
\begin{equation}
             e(F_{\mathrm{red}})
               =\#\{\hbox{original triple points}\}>0.
\hypertarget{tag-016Q}{}\label{cdp:eq:C.WF.3}
\end{equation}
In particular this entire weighted branch contains no zero-Euler
support.

Every such marked weighted log support has a realization in a
smooth compact complex threefold over \(\PP^1\), with connected
two-torus fibres away from two values.  The two marked weighted
supports are the given one and its complex conjugate.  The old-cone contractions and finite invariant-ring maps are
part of the marking at both ends, so the reduced analytic support, all its gluing maps,
and every full-fibre multiplicity are unchanged.  The total Euler
number of this compact realization is twice \eqref{cdp:eq:C.WF.3}.

For the given family the coarsening and quotient morphisms
recover its nonreduced fibre and analytic neighbourhood.
For the compact realization the isomorphism of marked weighted
supports is the one of \cref{def:II:marked-weighted-toric-support}.
\end{proposition}

\begin{proof}\leavevmode\par

 
\begin{enumerate}[label=\arabic*.,leftmargin=0pt,labelsep=.5em]
\item \par\medskip\noindent\textbf{The normalized root and its coarsening.}
Choose $M$ divisible by all $m_i$ and put $t=s^M$.
Apply \cref{prop:II:cyclic-root-normalization} with
$(\pi,Y,N,Z,Y_N)=(\operatorname{id}_X,X,M,\widetilde X,X_M)$.
Write $\rho_{\rm tot}=\rho\circ q:\widetilde X\to X$
for the composite finite-root and resolution map.
The original support is already SNC. Its local coordinates and
lattice are \eqref{cdp:eq:C.WF.4}--\eqref{cdp:eq:C.WF.5}, with
the general root exponent $N$ replaced by $M$; all unit roots
and tangent-dependent transition maps are the ones of this $f$.
The same height-one monomial-ideal subdivisions give the proper
equivariant map in \eqref{cdp:eq:C.WF.6}. The finite quotient
$X_M/\mu_M=X$ is its invariant-ring comparison. The subdivision
and its inverse coarsening, all original component labels, and
each new ray's old stratum and lattice position stay in these
maps. In particular the support on $\widetilde X$ is globally
embedded SNC, as proved by that labelled-ray comparison.

\item \par\medskip\noindent\textbf{The canonical orders, with their exact normalization.}
Under that same substitution, the canonical ray formula
\eqref{cdp:eq:C.RK.11} gives, at every height-one ray $(v,1)$,
\begin{equation}
 \operatorname{ord}_{v}(\rho_{\rm tot}^*\Omega)
       =\sum_i(k_i+1)v_i-1
       =M\lambda-1.
\hypertarget{tag-016U}{}\label{cdp:eq:C.WF.7}
\end{equation}
The equality uses $\sum_i m_iv_i=M$ and
$k_i+1=\lambda m_i$. The common value applies to old and new
rays and includes the unchanged tangent differential in a
two-factor chart. The nonzero pulled coefficient is a unit;
the primitive form in \eqref{cdp:eq:C.RK.13} is therefore
\begin{equation}
       \widetilde\Omega
          =s^{-(M\lambda-1)}\rho_{\rm tot}^*\Omega
\hypertarget{tag-016V}{}\label{cdp:eq:C.WF.8}
\end{equation}
and is nowhere zero on $\widetilde X$.
Its exponent is an integer by \eqref{cdp:eq:C.WF.7}.
An old triple cone contains at least one unimodular triangle,
so the root fibre has a triple point.  \cref{cdp:st:C.2.UN.68} and \cref{cdp:st:C.2.UN.71} apply to
this reduced semistable root.  In particular its marked
central log space is the toric space of
\((N,\Gamma,\mathscr P,c)\).  The old cones give a periodic
coarsening \(\mathscr P_0\) of this same developed triangulation.

\begin{enumerate}[label=(\arabic*),leftmargin=1.3em,labelsep=.5em]
\item Positivity of the $k_i+1$ in \eqref{cdp:eq:C.WF.1}
gives $\lambda>0$.  If \(\lambda>1\), every
integer \(k_i-m_i=(\lambda-1)m_i-1\) is nonnegative, so
\(\Omega/t\) is holomorphic on the smooth \(X\).  This
contradicts that \(\Omega\) generates \(f_*K_X\).
Thus \(0<\lambda\le1\).

\end{enumerate}

\item \par\medskip\noindent\textbf{The exact central finite action.}
Write \(\zeta=\exp(2\pi i/M)\).  Lift the root generator
to the developed affine plane and to the path-product
trivialization of the toric cocycle in \cref{cdp:st:C.2.UN.71}.  The resulting
maps on the open torus are
\begin{equation}
 g(z,s)=(a\,s^{b_0}z^A,\zeta s),\qquad
 A\in\operatorname{GL}(N),\quad b_0\in N,\quad a\in T_N,
 \qquad p\longmapsto Ap+b_0.
\hypertarget{tag-016W}{}\label{cdp:eq:C.WF.9}
\end{equation}
They preserve both \(\mathscr P\) and the marked coarsening
\(\mathscr P_0\).  Here \(z^A\) denotes the map whose
pullback on characters is \(A^{\mathsf T}\).

\begin{enumerate}[label=(\arabic*),leftmargin=1.3em,labelsep=.5em]
\item  We justify the central comparison, including its constants.
Each component map preserves its full toric boundary, so on
its dense torus it is a torus translation followed by the
integral map of its fan.  This follows directly by pulling
back a character: its divisor is the corresponding integral
combination of boundary primes, and the quotient of two
functions with that divisor is a holomorphic unit on the
compact component, hence a constant.  In the global
characters at a vertex \(p\), its dense component coordinates
are \(z^m s^{-\langle m,p\rangle}\).  Substitution in
\eqref{cdp:eq:C.WF.9} gives the component constant
\(a\,\zeta^{-(Ap+b_0)}\).
Conversely the double-curve map fixes the tangential
character comparison, and its marked normal-product log
section fixes the other character comparison.  In the
three-factor chart these are exactly the two substitutions
just displayed.  Thus after the \cref{cdp:st:C.2.UN.71} path-product
trivialization the quantities
\(a_p\zeta^{Ap+b_0}\) agree across every edge and equal one
constant \(a\).  This proves \eqref{cdp:eq:C.WF.9} on the entire central
log atlas, with the component maps and normal-product sections
of \cref{def:II:marked-weighted-toric-support}.

\item  The lifted $g$ normalizes the period maps
\(g_\gamma(z,s)=(c_\gamma s^\gamma z,s)\) and satisfies
\begin{equation}
 \begin{split}
 A\Gamma&=\Gamma,\qquad A^M=I,\\
 c_\gamma^A&=\zeta^{A\gamma}c_{A\gamma},\\
 g^M&=g_{\gamma_0}\quad\hbox{for }\gamma_0\in\Gamma .
 \end{split}
\hypertarget{tag-0166}{}\label{cdp:eq:C.WF.10}
\end{equation}
All are equalities of monomial maps.  For clarity the last
one means
\begin{equation}
 \gamma_0=\sum_{j=0}^{M-1}A^{M-1-j}b_0,\qquad
 \prod_{j=0}^{M-1}(a^{A^{M-1-j}})
       \ \zeta^{\sum_{j=0}^{M-1}jA^{M-1-j}b_0}
       =c_{\gamma_0}.
\hypertarget{tag-0167}{}\label{cdp:eq:C.WF.11}
\end{equation}
The central log comparison preserves these relations.
Indeed two monomial maps with the same affine map and base
character differ by a constant torus translation; their
actions on a height-one component determine that constant.
Equality of the central maps therefore implies the
displayed equality on every toric chart.

\item  The relative residue of \eqref{cdp:eq:C.WF.8} transforms under the root
generator by \(\zeta^{-M\lambda}\).  On the developed toric
log atlas it is a common nonzero scalar times
\(d\log z_1\wedge d\log z_2\), and \eqref{cdp:eq:C.WF.9} transforms this
residue by \(\det A\).  The scalar is common because the
developed transition maps are translations and the double
residues agree.  Consequently
\begin{equation}
                  \zeta^{-M\lambda}=\det A\in\{1,-1\}.
\hypertarget{tag-0168}{}\label{cdp:eq:C.WF.12}
\end{equation}
Together with \(0<\lambda\le1\), this proves \eqref{cdp:eq:C.WF.2},
including the evenness of every \(m_i\) in the second case.

\end{enumerate}

\item \par\medskip\noindent\textbf{Euler number on the original strata.}
Let \(V_i^\circ\) be an original component with all double
curves removed.  The strict root components above it have
finite unramified maps to \(V_i^\circ\); locally these are
exactly the branches of \(x_i^{m_i}=s^M\) with the same nonzero unit factors.  Their total degree is \(m_i\).
The semistable subdivisions change nothing there.
On every old root component this inverse image is its full
toric complement \((\C^*)^2\).  The sum of their Euler
numbers is therefore zero, and finite-cover multiplicativity
gives \(e_c(V_i^\circ)=0\).

\begin{enumerate}[label=(\arabic*),leftmargin=1.3em,labelsep=.5em]
\item  For a connected original double curve \(C\), remove its
original triple points and write \(C^\circ\).  The normalized
root over this open curve is locally the two-dimensional
monomial root times \(C^\circ\).  Its subdivision is the
corresponding chain times that curve.  It creates no triple
point over the interior of \(C^\circ\).  Choose a boundary
curve of an old strict root component which dominates $C$.  Every such dominating curve is a toric
\(\PP^1\), and the
preimage of the removed original endpoints is precisely
its two toric endpoints.  Thus its open part is \(\C^*\)
and maps finitely and unramifiedly to \(C^\circ\).
This remains true component by component when the root
cover is disconnected over \(C^\circ\).  Its positive total
degree gives \(e_c(C^\circ)=0\).  Additivity on the original
SNC stratification proves \eqref{cdp:eq:C.WF.3} on the original
coarse cells.

\end{enumerate}

\item \par\medskip\noindent\textbf{The period equations and finite action.}
Choose \(N=\Z^2\), \(\Gamma=Q\Z^2\) and logarithms \(C\)
of \(c\), as in \cref{cdp:st:C.2.UN.73}.  Initially put
\begin{equation}
 \begin{gathered}
 \Lambda= \Z^2,\quad \Lambda_0=Q^{-1}\Z^2,\quad
 B=(2\pi i)^{-1}Q^{-1}C,\\
 D=Q^{-1}AQ,\quad
 d=Q^{-1}b_0,\quad
 a_0=(2\pi i)^{-1}Q^{-1}\log a .
 \end{gathered}
\hypertarget{tag-0169}{}\label{cdp:eq:C.WF.13}
\end{equation}
The additive action is \(w\mapsto Dw+d\,l+a_0\),
where \(l=\log s/(2\pi i)\).  Replace these real lattices
by their images under one real invertible matrix
which makes the finite-order \(D\) orthogonal for an
averaged positive metric.  Replace \(B,D\) by their
conjugates and \(d,a_0\) by their images under the same
matrix.  Denote these images by
\(\Lambda\subset\Lambda_0\), \(B,D,d,a_0\). Thus each integral subgroup is carried by that same
invertible real map, with its lattice inclusion preserved.
Let \(T\) be the real matrix taking this vector
coordinate back to \((2\pi i)^{-1}\log z\).

\begin{enumerate}[label=(\arabic*),leftmargin=1.3em,labelsep=.5em]
\item  For \(\ell\in\Lambda\), \eqref{cdp:eq:C.WF.10} becomes
\begin{equation}
       (DB-BD-D/M)\ell\in\Lambda_0.
\hypertarget{tag-016A}{}\label{cdp:eq:C.WF.14}
\end{equation}
Define
\[
 d_M=\sum_{j=0}^{M-1}D^{M-1-j}d\in\Lambda,\quad
 a_M=\sum_{j=0}^{M-1}D^{M-1-j}a_0,\quad
 e_M=\sum_{j=0}^{M-1}jD^{M-1-j}d .
\]
The order relation is precisely
\begin{equation}
                   a_M+e_M/M-Bd_M\in\Lambda_0 .
\hypertarget{tag-016B}{}\label{cdp:eq:C.WF.15}
\end{equation}
These are lattice-valued identities obtained by taking
the chosen logarithms of \eqref{cdp:eq:C.WF.10}--\eqref{cdp:eq:C.WF.11}.

\end{enumerate}

\item \par\medskip\noindent\textbf{An equivariant compact middle family.}
In the orthogonal coordinates put
\[
 J=\begin{pmatrix}0&1\\-1&0\end{pmatrix},\qquad
 DJ=(\det D)JD.
\]
Since $(\det D)^M=1$, choose $k\ge1$ with
$\zeta^k=\det D$.
Choose \(R>0\) so that the symmetric part \(S\) of
\(\operatorname{Im}B+RI\) is positive definite, and put
\begin{equation}
 B'=B+iRI,\quad a'_0=a_0+iRd,\quad
 \kappa=e^{-2\pi R},\qquad
 V_s(q)=q+s^kJ\overline q,\quad H_s=I+s^kJ.
\hypertarget{tag-016C}{}\label{cdp:eq:C.WF.16}
\end{equation}
The simultaneous shift of \(a_0\) is required: it is
exactly the substitution \(u=\kappa s\) in the original
affine map.  Equations \eqref{cdp:eq:C.WF.14}--\eqref{cdp:eq:C.WF.15} remain true
with \(B',a'_0\), since their added terms cancel.

\begin{enumerate}[label=(\arabic*),leftmargin=1.3em,labelsep=.5em]
\item  On the whole punctured line use the period lattice
\begin{equation}
 \mathcal L_s=V_s(\Lambda_0)+
       \{\,V_s(B'\ell)+l(s)V_s(\ell):\ell\in\Lambda\,\}.
\hypertarget{tag-016D}{}\label{cdp:eq:C.WF.17}
\end{equation}
Changing the logarithm adds an angular period. Write
$\Lambda=P\Z^2$ with $P\in\operatorname{GL}_2(\R)$ and first
use this finite-index sublattice of $\Lambda_0$ for the angular
generators. Right-multiplication of the four real period columns
by $\operatorname{diag}(P^{-1},P^{-1})$ gives the columns in
\eqref{cdp:eq:C.TX.4}, with $(t,b')=(s,B')$. Subtracting their
real angular combinations gives precisely
\eqref{cdp:eq:C.TX.6} with $A'=\operatorname{Im}B'$.
Thus the determinant is $(\det P)^2$ times
\eqref{cdp:eq:C.WF.18}, where
$\operatorname{Im}B'=S+aJ$, $r=|s|^k$,
$x=\operatorname{Re}(s^k)$ and $y=-\log|s|/(2\pi)$.
The factor $(\det P)^2$ is positive, and restoring the finite
angular enlargement $\Lambda\subset\Lambda_0$ preserves rank.
Thus \eqref{cdp:eq:C.WF.17} gives a proper smooth torus family over
\(\C^*\), including the points where \(H_s\) is singular.

\item  The finite action on this family is the map
\begin{equation}
 (v,s)\longmapsto
 \bigl(Dv+V_{\zeta s}(a'_0)
                  +l(s)V_{\zeta s}(d),\ \zeta s\bigr).
\hypertarget{tag-016F}{}\label{cdp:eq:C.WF.19}
\end{equation}
The identity \(DV_s=V_{\zeta s}D\) follows from the
choice of \(k\).  Equation \eqref{cdp:eq:C.WF.14} proves that
\eqref{cdp:eq:C.WF.19} preserves the period subgroups.
Iterating it \(M\) times gives the period
\begin{equation}
 V_s(a'_M+e_M/M)+l(s)V_s(d_M)
       =V_s(B'd_M+n)+l(s)V_s(d_M),\qquad n\in\Lambda_0,
\hypertarget{tag-016H}{}\label{cdp:eq:C.WF.20}
\end{equation}
by \eqref{cdp:eq:C.WF.15}.  Hence its induced action has order
dividing \(M\).  Its base action has order \(M\), so
the action is faithful and free over \(\C^*\).
All these computations use only real coefficients
\(j/M\) inside \(V_s\); the complex number \(l(s)\)
is kept outside, as in the displayed period and action maps.

\end{enumerate}

\item \par\medskip\noindent\textbf{The exact zero-end filling and original action.}
Near zero use
\begin{equation}
                    z=\exp(2\pi iT H_s^{-1}v).
\hypertarget{tag-016I}{}\label{cdp:eq:C.WF.21}
\end{equation}
This divides by exactly the angular lattice.
The other periods give the variable toric holonomies
of \cref{cdp:st:C.2.UN.73}, with central values \(\kappa^\gamma c_\gamma\).
In this coordinate \eqref{cdp:eq:C.WF.19} is
\begin{equation}
 (z,s)\longmapsto
   \bigl(a(s)s^{b_0}z^A,\zeta s\bigr),\qquad
 a(s)=\exp\bigl(2\pi iT H_{\zeta s}^{-1}
                              V_{\zeta s}(a'_0)\bigr).
\hypertarget{tag-016J}{}\label{cdp:eq:C.WF.22}
\end{equation}
The coefficient is a holomorphic nonzero vector and
\(a(0)=\kappa^{b_0}a\).
Using the local base \(u=\kappa s\)
therefore restores exactly \(a\), \(c_\gamma\),
and the given central root action.  Both
\(\mathscr P\) and its marked old coarsening are
preserved. Apply the holomorphic-unit construction of \cref{cdp:st:C.2.UN.71} to the subdivided fan with base $u$ and the period homomorphism $d_\gamma(u)=\kappa^{-\gamma}c^0_\gamma(u/\kappa)$ determined by \eqref{cdp:eq:C.WF.21}. Its value at zero is the same $c_\gamma$. Descend this filling through the marked proper toric coarsening, using the identical units on both atlases. The monomial maps \eqref{cdp:eq:C.WF.22} extend to both and commute with that coarsening morphism.

\item \par\medskip\noindent\textbf{The reflected and conjugate infinity end.}
Set \(q=s^{-1}\), replace the vector frame by
\(v_\infty=q^kv\), and put
\[
 H^\infty_q=q^kI+J,\qquad
 V^\infty_q(w)=q^kw+J\overline w .
\]
Reverse the two nonangular period generators as
in \cref{cdp:st:C.2.UN.73}.  The new period group has generators
\begin{equation}
 H^\infty_q\Lambda_0,\qquad
 -V^\infty_q(B'\ell)+l(q)H^\infty_q\ell
                 \quad(\ell\in\Lambda).
\hypertarget{tag-016K}{}\label{cdp:eq:C.WF.23}
\end{equation}
The exact finite map in this frame is
\begin{equation}
 (v_\infty,q)\longmapsto
 \bigl((\det D)Dv_\infty+
       V^\infty_{\zeta^{-1}q}(a'_0)
       -l(q)V^\infty_{\zeta^{-1}q}(d),\
       \zeta^{-1}q\bigr).
\hypertarget{tag-016L}{}\label{cdp:eq:C.WF.24}
\end{equation}
Since \((\det D)D H^\infty_q=
H^\infty_{\zeta^{-1}q}D\), the exponential
coordinate with frame \(H^\infty_q\) has central
affine vertex map \(p\mapsto Ap-b_0\).
Use the reflected old and subdivided fans
\(-\mathscr P_0\) and \(-\mathscr P\).
After the local base change \(v_{\mathrm{base}}=\kappa q\),
the central map is exactly
\begin{equation}
 (z_\infty,v_{\mathrm{base}})
  \longmapsto
  (\overline a^{-1}v_{\mathrm{base}}^{-b_0}
                   z_\infty^A,\ \zeta^{-1}v_{\mathrm{base}}),
\hypertarget{tag-016M}{}\label{cdp:eq:C.WF.25}
\end{equation}
and its period constants are \(\overline c_\gamma\).
Here \(\overline{a'_0}=\overline{a_0}-iRd\)
and \(-d\,l(q)=-d\,l(v_{\mathrm{base}})+iRd\),
so the recorded shifts cancel exactly.

\begin{enumerate}[label=(\arabic*),leftmargin=1.3em,labelsep=.5em]
\item  Finally apply the torus inversion
\(z_\infty\mapsto z_\infty^{-1}\) and relabel the
period group by \(\gamma\mapsto-\gamma\).
This sends the reflected fans to the original
ones, sends the map \eqref{cdp:eq:C.WF.25} to the conjugate
of \eqref{cdp:eq:C.WF.9}, and leaves the conjugate period
constants correct because
\(\overline c_{-\gamma}^{-1}=\overline c_\gamma\).
Thus the full coarse and subdivided central
log spaces, including the finite action, are
the complex conjugates of the prescribed data.
The two base coordinates satisfy
\(u\,v_{\mathrm{base}}=\kappa^2\).

\end{enumerate}

\item \par\medskip\noindent\textbf{Finite stabilizers and smooth coarse quotients.}
At a point
of an original coarse central chart, take its finite
point stabilizer after quotienting by \(\Gamma\).
The Kummer chart \eqref{cdp:eq:C.WF.4}--\eqref{cdp:eq:C.WF.5} shows that
its inertia fixes each normal ray and is the identity
on the tangent stratum: the latter cover is unramified.
The central log isomorphism identifies these same inertia maps.  Its monomial lattice map is therefore
identity on the ray span and on its quotient, and,
being of finite order, is the identity on the whole
lattice.  We use a sufficiently small invariant
neighbourhood of this point; the group in the calculation
consists exactly of the elements fixing that point.
In the newly constructed chart the point stabilizer's
action on a character \(\chi^m\)
is multiplication by a holomorphic unit \(c_m(s)\).
These coefficients depend only on \(s\), by
\eqref{cdp:eq:C.WF.22} or \eqref{cdp:eq:C.WF.24}, and are multiplicative in \(m\).
Their central constants are the original stabilizer
characters.  For a generator of order \(d\), with
base character \(\xi\), let
\begin{equation}
 L_m(s)=\log\!\frac{c_m(s)}{c_m(0)},\qquad
 B_m(s)=\frac1d\sum_{j=0}^{d-1}jL_m(\xi^js).
\hypertarget{tag-016N}{}\label{cdp:eq:C.WF.26}
\end{equation}
The logarithm is the unique one vanishing at zero.
The finite-order identity gives
\(\sum_jL_m(\xi^js)=0\), and hence
\begin{equation}
 B_m(\xi s)-B_m(s)=L_m(s).
\hypertarget{tag-016O}{}\label{cdp:eq:C.WF.27}
\end{equation}
The coordinate change
\(\chi^m\mapsto e^{-B_m(s)}\chi^m\)
therefore makes every stabilizer coefficient constant.
The functions \(B_m\) are additive in \(m\), so this is
an automorphism of the saturated monoid algebra,
with the displayed inverse.  The base character has
constant coefficient, so its \(B_m\) is zero and the
change fixes $s$. This averaging is over the finite point
stabilizer of order $d$ in the chart.

\begin{enumerate}[label=(\arabic*),leftmargin=1.3em,labelsep=.5em]
\item  We can now use precisely the original coarse invariant
monoid chart determined by the lattice and cone in \eqref{cdp:eq:C.WF.5}
and the finite quotient in \eqref{cdp:eq:C.WF.6}:
\begin{equation}
 \left(\OO_{X_M,\mathrm{coarse}}\right)^{\mathrm{Stab}}
        =\C\{x_1,x_2,x_3\},\qquad
       s^M=x_1^{m_1}\cdots x_r^{m_r}.
\hypertarget{tag-016P}{}\label{cdp:eq:C.WF.28}
\end{equation}
The same equation holds in the new chart after the
explicit unit change \eqref{cdp:eq:C.WF.26}.  This proves smoothness
of the quotient total space and its exact component
multiplicities.  It also proves the embedded SNC
condition for its reduced fibre, including the original
component labels.  Reduction and finite invariants
commute in characteristic zero, so the central
coarsening and quotient identify the original
reduced analytic support and all its gluing maps.

\end{enumerate}

\item \par\medskip\noindent\textbf{Compact gluing.}
The determinant \eqref{cdp:eq:C.WF.18} and real fundamental
parallelepipeds give a proper middle torus family
on every compact annulus. The holomorphic-unit clause of \cref{cdp:st:C.2.UN.71}, applied to the period homomorphisms from \eqref{cdp:eq:C.WF.21} and \eqref{cdp:eq:C.WF.23} after the displayed base rescalings, gives the proper fillings on $\mathscr P$ and its reflected fan. The torus inversion already specified identifies the latter with the conjugate marked data. A toric coarsening has
the same fan support and is proper; the coarse
fillings therefore remain proper.  The exact
exponential maps above identify them with the
middle family on punctured discs.  They glue
to a Hausdorff compact normal root space over
\(\PP^1_s\).  Its finite action is \eqref{cdp:eq:C.WF.19}, extended by the displayed monomial
maps at both ends.  Its quotient is proper
over \(\PP^1_s/\mu_M=\PP^1\) and is smooth
by \eqref{cdp:eq:C.WF.28}; it is free over the punctured
base, where each quotient fibre is one of
the connected two-tori.

\begin{enumerate}[label=(\arabic*),leftmargin=1.3em,labelsep=.5em]
\item  The two marked reduced supports and their
multiplicities are the prescribed one and
its conjugate.  Euler additivity and the
zero Euler number of a smooth torus give
the stated total Euler number.  On the
original family, \eqref{cdp:eq:C.WF.6} recovers its full
nonreduced fibre by its invariant analytic algebra.
The compact family uses the variable period units in
\eqref{cdp:eq:C.WF.17}. At its two ends, the exponential maps,
base rescalings, torus inversion and invariant-chart maps above
give the asserted isomorphisms of marked weighted supports.

\end{enumerate}

\end{enumerate}
\end{proof}

\section{Hopf orders and exceptional configurations}\label{sec:II:hopf-orders-exceptional-configurations}

\subsection{Maximal Hopf component reduction}\label{subsec:II:maximal-hopf-reduction}

\begin{lemma}[Coefficient-one boundaries on Hopf normalizations]\hypertarget{tag-P367}{}\label{lem:II:hopf-maximal-boundary}
\StatementRecord{lem:II:hopf-maximal-boundary}{proved}{P367}{\cref{lem:II:hopf-curves,lem:II:canonical-boundary-blowup,cdp:st:C.2.UN.17}; \cite[\S4.11.2(III), pp.~982--983]{Kulikov1977}}
Let $T$ be the smooth normalization of a component of maximal
canonical order $M$ in a smooth-total torus degeneration with full
reduced normal-crossing central fibre, with two distinct
normalization branches over every self-double curve as in
\cite[\S4.11.2(III), pp.~982--983]{Kulikov1977}. Suppose its smooth minimal
model $H$ is Hopf, and write $\pi:T\to H$ for the point-blowdown map.
The boundary divisor induced by the primitive canonical section is
\[
 -K_T=D_T=\sum_C\bigl(M-k_{\mathrm{other}(C)}+1\bigr)C.
\]
Every rational curve on $T$ is exceptional for $\pi$.
The positive support of $D_T$ is SNC. Two coefficient-one
curves do not meet; each rational coefficient-one curve meets
exactly one other positive boundary curve, and each elliptic
coefficient-one curve is disjoint from the other boundary curves.
Consequently the component has no rational self-double curve and
no triple point at which all three local branches belong to it.
\end{lemma}
\begin{proof}
The canonical restriction, with both preimages of a self-double
curve counted, is
\begin{equation}
 -K_T=D_T=\sum_C a_CC,\qquad
 a_C=M-k_{\mathrm{other}(C)}+1\ge1.
\hypertarget{tag-00MI}{}\label{cdp:eq:C.HM.2}
\end{equation}
It identifies the canonical line and the divisor of its
meromorphic section. A self-double branch has coefficient one.
Apply \cref{lem:II:hopf-curves} to this same minimal surface $H$.
Its curves are disjoint smooth elliptic curves of square zero,
so every rational curve on $T$ is exceptional for $\pi$.

At each point blowup $\sigma$ in the factorization of $\pi$,
\cref{lem:II:canonical-boundary-blowup}, with its two surfaces
equal to these consecutive models and its divisor equal to the
pushforward of $D_T$, gives
\begin{equation}
 D_{\mathrm{new}}=\sigma^*D_{\mathrm{old}}-E.
\hypertarget{tag-00MJ}{}\label{cdp:eq:C.HM.3}
\end{equation}
Every intermediate divisor is effective: a negative coefficient
would persist on its strict transform in $D_T$. Thus all centres
lie on the positive support. Starting with the disjoint elliptic
curves on $H$, that support remains SNC. A smooth-point blowup
of a coefficient-$a$ curve creates coefficient $a-1$; a node
blowup of coefficients $a,b$ creates coefficient $a+b-1$.
Coefficient-zero exceptional curves are omitted from the positive
support and remain curves on the surface.

These updates prove
\begin{equation}
\begin{gathered}
 \text{two coefficient-one positive curves do not meet},\\
 \text{a rational coefficient-one curve is a leaf of the positive boundary},\\
 \text{an elliptic coefficient-one curve is disjoint from the other boundary curves}.
\end{gathered}
\hypertarget{tag-00MK}{}\label{cdp:eq:C.HM.4}
\end{equation}
A smooth blowup creates the pair $(a,a-1)$; a node blowup could
create an edge $(1,1)$ only from an existing edge $(1,1)$.
A rational coefficient-one curve is first created at a smooth
point of a coefficient-two curve and has valence one.
Later node blowups preserve its valence, and a smooth blowup
on it has exceptional coefficient zero. The same smooth-point
calculation leaves an elliptic coefficient-one curve isolated.

For a rational self-double curve, the two normalized preimages
$C_1,C_2$ in \cite[\S4.11.2(III), pp.~982--983]{Kulikov1977}
have coefficient one. Each meets the remaining boundary once,
so the image has one triple point. Their normal-product map gives
\begin{equation}
 C_1^2+C_2^2=-1.
\hypertarget{tag-00ML}{}\hypertarget{tag-013P}{}\SsixEquationAliases{\label{cdp:eq:C.TK.7}}\label{cdp:eq:C.HM.5}
\end{equation}
Both curves are exceptional on a point blowup of $H$ and have
negative integral square, a contradiction. A triple point with
three branches of the component would give intersecting
coefficient-one curves, contrary to \eqref{cdp:eq:C.HM.4}.
\end{proof}

\begin{proposition}[Hopf minimal models at maximal canonical order]\hypertarget{tag-01EO}{}\label{cdp:st:C.2.UN.77}
\StatementRecord{cdp:st:C.2.UN.77}{proved}{01EO}{\SeriesRef{cdp:st:C.2.UN.75}, \SeriesRef{cdp:eq:C.HM.2}, \SeriesRef{cdp:eq:C.HM.4}, \SeriesRef{cdp:eq:C.HM.3}, \SeriesRef{cdp:eq:C.HM.7}, \SeriesRef{cdp:eq:C.HM.9}, \SeriesRef{cdp:eq:C.HM.5}, \SeriesRef{cdp:eq:C.HM.8}; \cite[\S4.11.2(III), pp.~982--983]{Kulikov1977}; \cite[\S4.5, p.~970]{Kulikov1977}; \cite[\S4.10.7--4.10.8, pp.~980--982; Lemma~4.14 and proof, pp.~983--985]{Kulikov1977}; \cite[\S8, p.~341]{Kulikov1981}; \cite[Lemma~2.19, p.~965]{Kulikov1977}; \cite[\S4.11.2, pp.~982--983]{Kulikov1977}; \cite[\S4.13, p.~983]{Kulikov1977}; \cite[Lemma~4.14, pp.~983--984]{Kulikov1977}; \cite[\S4.16.2, p.~986]{Kulikov1977}; \cite[\S9, pp.~341--342]{Kulikov1981}}

Let \(f:X\to\Delta\) be a proper degeneration of connected
complex two-tori with smooth total space and full reduced
simple-NC central fibre.  Normalize its canonical generator
as in \ref{cdp:st:C.2.UN.75} and put
\begin{equation}
        \operatorname{div}(\Omega)=\sum k_iV_i,\qquad
        \min k_i=0,\qquad M=\max k_i .
\hypertarget{tag-00MG}{}\label{cdp:eq:C.HM.1}
\end{equation}
There is a finite overlap-compatible modification record,
over the same disc and unchanged over its punctured part,
whose endpoint has smooth total space and full reduced
normal-crossing fibre, such that every component with
canonical order \(M\) and Hopf minimal model is itself a
smooth minimal Hopf surface.  These Hopf components are
embedded, their boundary curves are disjoint smooth
elliptic curves, and no triple point meets them.
Other endpoint components may retain the two normalized
branches of an NC self-intersection.

All original Hopf-birational prime divisors survive, with
their original orders.  Every whole-component removal block
removes a ruled component of positive order \(q\);
its preliminary additions, when needed, are finitely many
ruled components of order \(q+1\le M\).  There are no such
additions when \(q=M\).  Thus the retained reduced degeneration stages have
no global component of order greater than \(M\).
The full multiplicities and canonical orders of the auxiliary
common blowups are given by divisor pullback and the blowup Jacobian.

The record retains its centres, local contractions,
normalization branches and restored overlap maps.  Its
reverse recovers the original family.  For a given compact
family the same operations take place on that compact
base, by the annular comparison.  Their Euler
changes are exactly the normal-cone terms of \ref{cdp:st:C.2.UN.75}.

\end{proposition}

\begin{proof}\leavevmode\par

 
\begin{enumerate}[label=\arabic*.,leftmargin=0pt,labelsep=.5em]
\item \par\medskip\noindent\textbf{The Hopf normalization and its boundary.}
At each smooth-total NC stage apply \cref{lem:II:hopf-maximal-boundary}
with $T$ equal to the normalization of the chosen order-$M$
Hopf-birational component and $\pi:T\to H$ its minimal point map.
Its divisor \eqref{cdp:eq:C.HM.2} is the restriction of this
family's primitive canonical section. Equations
\eqref{cdp:eq:C.HM.3}--\eqref{cdp:eq:C.HM.4} give the point-blowup
coefficients and coefficient-one branches; the lemma excludes
rational self-double curves. The initial components are embedded,
and the modifications below introduce self-incidences only on
rational exceptional curves. Thus these maximal Hopf components
remain embedded. In the following last-point argument write
$H$ for this embedded source and $\pi:H\to H_{\min}$ for its
minimal point map.

\item \par\medskip\noindent\textbf{The exact last-point alternatives.}
Suppose such an \(H\) is not minimal and choose a last
point blowup in the map \(\pi\).  Its exceptional
curve \(E\) is a smooth \((-1)\)-curve.  Formula
\eqref{cdp:eq:C.HM.3} leaves exactly three cases.

\begin{enumerate}[label=(\arabic*),leftmargin=1.3em,labelsep=.5em]
\item  First, \(E\) is not a boundary curve.  It was created
at a smooth point of a coefficient-one boundary \(C\),
and meets that curve once.  The ambient normal
sequence is
\begin{equation}
 0\longrightarrow\OO_{\PP^1}(-1)
 \longrightarrow N_{E/X}
 \longrightarrow\OO_{\PP^1}(-1)\longrightarrow0 .
\hypertarget{tag-00MM}{}\label{cdp:eq:C.HM.6}
\end{equation}
Its splitting and the ordinary two-ruling diagram of
\ref{cdp:st:C.2.UN.75} give the type-I move.  The neighbour across
\(C\) has order \(M\).  It is not Hopf-birational.
If \(C\) is elliptic, its square on \(H\) is at most
\(-1\) and it contains no triple point; its square on
the neighbour is therefore positive.  If \(C\) is
rational, it was already exceptional before this last
blowup, so its current square is at most \(-2\);
\eqref{cdp:eq:C.HM.4} gives one triple point, and its square on
the neighbour is again positive.  Such squares are
impossible for the corresponding elliptic or rational
curves on a blown-up Hopf.  The move removes one
point blowup from \(H\), adding it to a non-Hopf
component.

\item  Second, \(E\) was created at a node of two positive
boundary curves of coefficients \(a,b\).  It has
coefficient \(a+b-1\) and two triple points.  The
two-normal-bundle equation makes its square on the
other component also \(-1\), so the type-II
move of \ref{cdp:st:C.2.UN.75} applies.  It blows down \(E\) on both
sides and blows up the two opposite normalized
branches.  If an opposite component has order \(M\),
the associated curve \(C\) on \(H\) has coefficient
one.  It is a rational leaf by \eqref{cdp:eq:C.HM.4}, and its
square after the last node blowup is at most \(-2\).
Its square on that opposite component is
\begin{equation}
                 -1-C_H^2\ge1 .
\hypertarget{tag-00MN}{}\label{cdp:eq:C.HM.7}
\end{equation}
Hence that opposite component is not Hopf-birational.
This proves the same assertion when both opposite
branches belong to one global component: their normalization branches are used in \eqref{cdp:eq:C.HM.7}.

\item  Let
\begin{equation}
       h=\sum_{\substack{k_i=M\\
                     V_i\ \mathrm{Hopf\text{-}birational}}}
                              b_2(V_i).
\hypertarget{tag-00MO}{}\label{cdp:eq:C.HM.8}
\end{equation}
At the smooth NC stages just described, this is the
total number of surface point blowups on these Hopfs.
Each of the two moves above strictly decreases \(h\).
The whole-component blocks below are ordered first by their
canonical component counts.

\item  Third, \(E\) was created at a smooth point of a
positive curve of coefficient \(a\ge2\).  It has
coefficient \(a-1\) and one triple point.  Write \(Q\)
for its neighbour and \(P\) for the third component.
The exact normal degrees and canonical orders are
\begin{equation}
 \begin{split}
 E_H^2&=-1,\qquad E_Q^2=0,\\
 q:=k_Q&=M-a+2,\qquad k_P=q-1,\\
                         &1\le q\le M .
 \end{split}
\hypertarget{tag-00MP}{}\label{cdp:eq:C.HM.9}
\end{equation}
The complete rational curve of square zero gives
the ruling on the compact normalization of
\(Q\).  For clarity, its local complete deformations
are disjoint, since their intersection number is zero;
the resulting compact curve family gives the ruling
after the point blowdowns of its singular fibres.
Equivalently this is the ruled-surface comparison in
\cite[\S4.5, p.~970]{Kulikov1977}.  Its one boundary intersection
is the section belonging to \(P\).  Any other
horizontal boundary would meet a general complete
fibre in this same local family, which it does not.
All other boundary curves on \(Q\) are vertical.
This is the precise ruled section-and-fibre
configuration, including any nonminimal fibres.

\end{enumerate}

\item \par\medskip\noindent\textbf{Ruled removal and cleanup of its singular charts.}
We use the finite maps in
\cite[\S4.10.7--4.10.8, pp.~980--982; Lemma~4.14 and proof, pp.~983--985]{Kulikov1977}, with
the missing-fibre correction in
\cite[\S8, p.~341]{Kulikov1981}.  Their overlaps and proper
composite graphs are those constructed in \ref{cdp:st:C.2.UN.75}.
They remove the ruled component in \eqref{cdp:eq:C.HM.9}.
If its boundary is disconnected, they first insert
the specified missing fibre components; these
exceptional ruled surfaces have order \(q+1\).
When \(q=M\), the canonical divisor on its
normalization has only negative boundary
coefficients, and \cite[Lemma~2.19, p.~965]{Kulikov1977}
makes that boundary connected.  No preliminary
addition is then made.  Thus all added orders are
at most the original fixed \(M\).

\begin{enumerate}[label=(\arabic*),leftmargin=1.3em,labelsep=.5em]
\item  The result of this removal need not immediately
have smooth total space.  We retain precisely the
finite sigma-atlas of
\cite[\S4.11.2, pp.~982--983]{Kulikov1977}: its singular sets are
contained in finitely many disjoint charts; in each
chart the recorded local exceptional branches
contract successively to a smooth NC chart.  We
now describe a cleanup, before applying
another last-point argument to a Hopf.

\item  If one of these lists is nonempty, choose a
component of largest canonical order among the
globally occurring last locally blown-up branches.
Use an auxiliary resolution \(a:\widehat X\to X'\)
as in \cite[\S4.13, p.~983]{Kulikov1977} to read its
old strict component \(\widehat V\).  We apply Lemma~4.14 on this auxiliary resolution to the strict
transform of the same global component. The component counts
in \eqref{cdp:eq:C.HM.10} are those of the degeneration $X'$.
Each branch has a complete local exceptional
\(\PP^1\)-fibre.  These complete fibres extend to
a ruling of the old smooth surface
\(\widehat V\).  In the
generic local exceptional chart there is exactly
one horizontal boundary section.  Completeness
of those fibres shows that there can be no
additional global horizontal boundary; all the
other boundary, including the auxiliary-resolution
exceptional curves over the finite exceptional
points, is vertical.  This supplies
conditions (a), (b) and (d) of
\cite[Lemma~4.14, pp.~983--984]{Kulikov1977}, by their normalization and section maps.

\item  If another curve on this chosen component had
been locally blown up, its exceptional branch
would have order one larger.  If that child
is not last, follow its finite local descendant
list.  The order increases at each step, so
a last descendant has strictly larger order
than the chosen component.  This contradicts
the choice of largest last-branch order.
Thus the choice gives condition
(e).  This is the exact comparison in
\cite[\S4.16.2, p.~986]{Kulikov1977}.  The reordering in
its \S4.12 makes its specified branches last
without changing their restored overlaps.
The correction in \cite[\S8, p.~341]{Kulikov1981}
inserts any missing vertical links in the
specified smooth open charts, leaving their
overlap fixed.  Lemma~4.14 then removes this
global ruled component, and the result again
has the same kind of sigma-atlas.

\item  Every cleanup removal has positive order:
a locally exceptional component has the
order of its parent plus one.  It
adds only finitely many components at the
next order, and none at \(M\).  Here we
use precisely the maximal-old-order clause
of \cite[\S8, p.~341]{Kulikov1981}, including its
last-locally-blown-up case: its boundary
is connected and no missing links are
inserted.  A Hopf
component cannot be one of these locally
exceptional global components: its complete
local rational fibres would give rational
curves through an open set of its compact
normalization, contrary to the finite
exceptional-curve description above.
The strict Hopf divisors survive by the
proper graph argument in \ref{cdp:st:C.2.UN.75}.

\item  The component-removal induction of \cref{cdp:st:C.2.UN.75}
applies with its levels $1,\ldots,R$ equal to $1,\ldots,M$
and its counted objects equal to the global components of $X'$.
Each removal deletes one object of order $q>0$, adds finitely
many of order $q+1$ for $q<M$, and adds none for $q=M$.
These are the same updates as in \cite[\S9, pp.~341--342]{Kulikov1981},
so the cleanup terminates.

\item  When cleanup stops, no recorded locally
exceptional branch remains.  The defining
contraction lists of the sigma-atlas are
empty; their terminal smooth NC charts
are therefore the current charts.
The total space is smooth and its full
fibre is reduced NC.  Its normalization
components are smooth.  Other global
components may retain self-incidences.
Apply \cref{lem:II:hopf-maximal-boundary} to the normalization
of each maximal Hopf component in this smooth-total stage.
Its canonical section and orders are unchanged, so the lemma
excludes rational self-incidences.  No elliptic self-incidence
was introduced: type-I and type-II changes
use rational exceptional curves; in a
local ruled contraction the horizontal
parent has order one lower, so a
maximal-order Hopf can only meet that
contraction in vertical rational curves.
Thus the maximal Hopfs are again smooth
embedded components, and their ordinary
point maps to their Hopf minimal models
are available for the next step.

\end{enumerate}

\item \par\medskip\noindent\textbf{Termination of the whole procedure.}
Perform a type-I or type-II step whenever
the chosen last point gives the first or
second case.  In the third case perform
the finite ruled removal and cleanup
just specified.  Across the entire
procedure there are only finitely many
divisorial removal blocks: the same
finite-level induction counts all their
positive component labels together.
Between these blocks every step strictly
decreases the finite integer \eqref{cdp:eq:C.HM.8}.
More explicitly, if \(c_j\) counts the
retained global components of order \(j\),
the block-level potential is
\begin{equation}
                 (c_1,c_2,\ldots,c_M,h)
                 \quad\hbox{in lexicographic order.}
\hypertarget{tag-00MH}{}\label{cdp:eq:C.HM.10}
\end{equation}
A whole-component removal decreases its
first changed \(c_j\), whatever finite
next-level additions or Hopf point
changes occur inside the block.
A type-I or type-II step chosen on a
last Hopf point leaves all the \(c_j\)
unchanged and decreases \(h\).
Lexicographic order on this fixed
finite product of nonnegative integers
has no infinite strictly descending
sequence, by the same successive-coordinate
induction.  We evaluate this potential at the smooth-total endpoints
of the finite blocks.
An infinite procedure would therefore
have either infinitely many divisorial
removals, contradicting that induction, or
an infinite tail of strict decreases of
\(h\).  Both are impossible.

\begin{enumerate}[label=(\arabic*),leftmargin=1.3em,labelsep=.5em]
\item  At the endpoint a nonminimal maximal
Hopf would still have a last point
blowup and one of the three alternatives,
so every such Hopf is minimal.  Its
curves are the disjoint smooth elliptic
curves proved above, and hence no
triple point meets it.  The original
Hopf primes survive with unchanged
orders by \ref{cdp:st:C.2.UN.75}, and each additional
prime used in the record is ruled.
The same proper graphs and
restored overlap isomorphisms give
the asserted reversible comparison.
Their annular gluing and normal-cone
Euler terms are exactly those of \ref{cdp:st:C.2.UN.75}.

\end{enumerate}

\end{enumerate}
\end{proof}

\subsection{A compact zero-Euler fibre with a negative rational curve}\label{subsec:II:compact-negative-curve}

\begin{proposition}[The \((-1,-2)\) curve with canonical orders \((0,0,1)\) occurs in a compact zero-Euler torus fibre]\hypertarget{tag-01EQ}{}\label{cdp:st:C.2.UN.79}
\StatementRecord{cdp:st:C.2.UN.79}{proved}{01EQ}{\cref{lem:II:canonical-bimeromorphic-direct-image,lem:II:blowup-topology}}

There is a smooth compact projective threefold \(X\), a proper map
\(f:X\to\PP^1\) with connected general fibre a complex two-torus,
and a full reduced SNC fibre \(S=f^{-1}(0)\) with \(e(S)=0\),
containing a smooth rational curve \(F\) with the following exact
properties. The curve belongs to just one component \(P\) generically,
is not a double curve, has \(F_P^2=-1\), and meets two other
components transversely, once each, at distinct points. Its entire
analytic neighbourhood, including \(f\), is the germ along the zero
section of
\begin{equation}
 \operatorname{Tot}\bigl(\OO_{\PP^1}(-1)_a
                 \oplus\OO_{\PP^1}(-2)_b\bigr),\qquad
 f=bUV,\quad [U:V]\in\PP^1.                         \hypertarget{tag-00AZ}{}\label{cdp:eq:C.BC.1}
\end{equation}
Here \(P=(b=0)\), and the other two components are the fibres
\(U=0\) and \(V=0\) of the displayed vector bundle.
The generator of \(f_*K_X\) has orders \(0,1,0\)
on these three components in that order. Thus, after ordering the
two neighbours by canonical height, the tuple is \((0,0,1)\).
In particular
\begin{equation}
                    N_{F/X}=\OO(-1)\oplus\OO(-2).            \hypertarget{tag-00B4}{}\label{cdp:eq:C.BC.2}
\end{equation}

There is a composite point-blowdown map
$P\to\PP^1\times E$, where $E$ is an elliptic curve.
It maps \(F\) isomorphically onto \(\PP^1\times\{e_0\}\).
For the elliptic map obtained by composing this blowdown with
\(\PP^1\times E\to\PP^1\), the curve \(F\) is therefore
horizontal. The restriction of the surface blowdown to $F$ is the
displayed isomorphism, so $F$ is not one of its exceptional curves.
\end{proposition}

\begin{proof}

\begin{enumerate}[label=\arabic*.,leftmargin=*]
\item \par\medskip\noindent\textbf{A compact elliptic surface with an \(I_3\) chart.}
Consider the projective Weierstrass surface over $\PP^1_t$
with coefficient sections in $\OO(4)$ and $\OO(6)$ and equation
\begin{equation}
       y^2=x^3+(t^4+1)x+t^6                         \hypertarget{tag-00B5}{}\label{cdp:eq:C.BC.3}
\end{equation}
Its discriminant, up to a nonzero constant, is
$31t^{12}+12t^8+12t^4+4=h(t^4)$, where
$h(q)=31q^3+12q^2+12q+4$.
The identities $h-(q/3)h'=4(q+1)^2$ and $h'(-1)=81$
show that $h$ is square-free; $h(0)=4$ makes its pullback
under $t\mapsto t^4$ square-free as well.
At infinity the equation in coordinates
$(T,X,Y)=(t^{-1},x/t^2,y/t^3)$ is
$Y^2=X^3+(1+T^4)X+1$, with the same discriminant factor
equal to $31$ at $T=0$.
The singular fibres are therefore simple nodal cubics,
and the total surface is smooth.
Translate the coordinate so that one nodal value is $t=0$,
keeping infinity a smooth value, and take the global cubic base change
\(t=s^3\). Normalize and minimally resolve. The resulting
compact smooth projective elliptic surface \(J\to\PP^1_s\)
has an \(I_3\) fibre
\begin{equation}
                         C_0+C_1+C_2 .             \hypertarget{tag-00B6}{}\label{cdp:eq:C.BC.4}
\end{equation}
The components are smooth rational curves, meeting in a triangle
at three distinct nodes, with no triple point. The base change
is unramified at the other singular values; at a smooth value
it is the ordinary pullback of a smooth family. Thus these
operations give a smooth total surface after resolving the
single local \(A_2\) equation at the chosen nodal value.

\item The local comparison is given by two charts covering a
neighbourhood of the whole exceptional curve. At the original nodal point the map is
\(xy=t\), by the holomorphic Morse coordinates. The root is
\(xy=s^3\). Two consecutive charts of its minimal
resolution are
\begin{equation}
 \begin{aligned}
 (x,y,s)&=(b,b^2z^3,bz),\qquad (b,z)\in\C^2,\\
 (x,y,s)&=(u^2c,uc^2,uc),\qquad (u,c)\in\C^2,\\
                 u&=z^{-1},\qquad c=z^2b .
 \end{aligned}                                             \hypertarget{tag-00B7}{}\label{cdp:eq:C.BC.5}
\end{equation}
The curve \(C_1\) is \(b=0\) in the first chart and
\(c=0\) in the second. Its adjacent curves are \(z=0\)
and \(u=0\). The two-chart union is the total space
of \(\OO(-2)\), restricted to a neighbourhood if necessary;
its map is exactly \(s=bz=cu\).
The displayed monomials satisfy $xy=s^3$, glue by the stated
transition and are the successive point-resolution maps.
Their union identifies the entire neighbourhood and its function $s$.

\item On a disc about the chosen nodal value the original
Weierstrass surface is a smooth cubic hypersurface in
$\PP^2\times\Delta_t$. Adjunction gives its canonical line
as $\OO(-3+3)\otimes K_{\Delta_t}$, so its residue two-form
is a nowhere-zero generator on this whole neighbourhood.
At the node it is a nonvanishing multiple of $dx\wedge dy$.
The two-form $dx\wedge ds/x$ is $db\wedge dz$
on the first chart and \(du\wedge dc\) on the second.
Pull back that residue form and divide by $s^2$; the resulting
form is a nonvanishing holomorphic multiple of the
displayed two-form and extends through the unchanged
smooth charts of the original fibre. The three
components have the same original root order two.
This gives a nowhere-zero canonical generator
near \eqref{cdp:eq:C.BC.4}. Hence the canonical evaluation orders on all three components of
$J$ near this fibre are zero. Properness and connectedness
make every holomorphic function on this neighbourhood a base
function; its nowhere-zero canonical form therefore generates
the canonical direct image.

\par\medskip\noindent\textbf{The curve blowup and its canonical orders.}
Fix an elliptic curve \(E\) and a point \(e_0\in E\).
Set \(X_0=J\times E\), with map \(f_0\) induced by \(s\).
Its general fibre is the product of two elliptic curves.
The central components are \(P_i=C_i\times E\), all
of canonical order zero. Let $\pi_0:X_1\to X_0$ be the blowup of the smooth
rational curve
\begin{equation}
                      B_0=C_0\times\{e_0\}\subset P_0 .
                                                               \hypertarget{tag-00B8}{}\label{cdp:eq:C.BC.6}
\end{equation}
The centre is contained generically in just \(P_0\) and
meets each of its two double curves transversely at one
point. The following charts show that its full central fibre is
reduced SNC.
Near any such point write the original map as \(xy\),
the centre as \((x,a)\), and the blowup charts as
\(a=xA\) and \(x=aX\). Their maps are respectively
\(xy\) and \(aXy\). All displayed factors have
multiplicity one. The new exceptional component \(R_0\)
has canonical order one, because the Jacobian of this
smooth curve blowup has a simple exceptional zero.

\item Choose the labelling in \eqref{cdp:eq:C.BC.5} so that \(z=0\) is
the branch of \(C_0\). Let \(a\) be a local coordinate
of \(E\) at \(e_0\). Near the whole curve
\(C_1\times\{e_0\}\), the centre is \((z,a)\)
near its first endpoint and is absent elsewhere.
In the \(z\)-chart of its blowup put
\begin{equation}
               a=zA,\qquad s=zb .                          \hypertarget{tag-00B9}{}\label{cdp:eq:C.BC.7}
\end{equation}
The strict transform \(F\) is \(A=b=0\). In the other
chart of \eqref{cdp:eq:C.BC.5} no modification occurs near this curve;
write its elliptic coordinate as \(a_\infty\). The transition is
\begin{equation}
       u=z^{-1},\qquad A=u a_\infty,\qquad b=u^2c,
             \qquad s=zb=uc .                              \hypertarget{tag-00BA}{}\label{cdp:eq:C.BC.8}
\end{equation}
These are exactly the two vector-bundle charts in \eqref{cdp:eq:C.BC.1}.
They prove \eqref{cdp:eq:C.BC.2} and \(F_P^2=-1\), with the function $s$ and both transition maps. The two boundary intersections occur at
\(z=0\) and \(u=0\), whereas \(F\) is contained
generically only in \(b=0\).

\item Up to the same nonvanishing holomorphic factor, the product
canonical form pulls back as
\begin{equation}
 db\wedge dz\wedge da
       =z\,db\wedge dz\wedge dA .                          \hypertarget{tag-00BB}{}\label{cdp:eq:C.BC.9}
\end{equation}
At the other endpoint this is a nonzero unit times
\(du\wedge dc\wedge da_\infty\). Thus it has a simple
zero exactly on \(R_0=(z=0)\), with zero orders on
\(P\) and its other old neighbour. This gives the canonical tuple for the direct-image generator:
apply \cref{lem:II:canonical-bimeromorphic-direct-image} to
$(X,X',\phi)=(X_1,X_0,\pi_0)$ and the maps to the same disc.
Its pullback is a generator, and the displayed Jacobian
has order one exactly on $R_0$.

\item The surface \(P\) is \(\operatorname{Bl}_p(\PP^1\times E)\),
where \(p\) is the endpoint of
\(\PP^1\times\{e_0\}\). The strict transform of that
rational curve is \(F\). The old double curve through
\(p\) has the other tangent direction in \(P_i\), so
its strict transform meets the exceptional line at a
different point from \(F\). In particular the full
strict transform of \(P_0\) is disjoint from \(F\).
The surface blowdown maps \(F\) isomorphically
to \(\PP^1\times\{e_0\}\); its composition with
the first projection has degree one on \(F\).

\par\medskip\noindent\textbf{Further blowups away from the curve.}
We make the remaining centres explicit. Use
\begin{equation}
 E:\ y^2=x(x-1)(x-2),\qquad e_0=(0,0),\quad e_1=(1,0).
                                                               \hypertarget{tag-00B0}{}\label{cdp:eq:C.BC.10}
\end{equation}
The strict \(P_0\) is still the identical surface
\(\PP^1\times E\): restriction of the first blowup
to this component is the blowup of its Cartier divisor
\(B_0\), hence an isomorphism. Blow up the disjoint
curve \(B_1=C_0\times\{e_1\}\). Again \(P_0\)
is unchanged and all full multiplicities are one.
The second centre is disjoint from \(F\), so this
operation preserves a neighbourhood of the whole
compact curve.

\item Choose the coordinate \(z\) on \(C_0\simeq\PP^1\)
so its two old boundary points are \(0,\infty\). Put
\begin{equation}
 P(z)=z^2+1,\qquad Q(z)=3z^2+z+5,\qquad
 D=\{P(z)x=Q(z)\}\subset\PP^1\times E .
                                                               \hypertarget{tag-00B1}{}\label{cdp:eq:C.BC.11}
\end{equation}
This is the divisor of a section of
\(\OO_{\PP^1}(2)\boxtimes\OO_E(2\cdot\infty_E)\).
Its map to \(\PP^1_z\) is the exact smooth double
cover
\begin{equation}
          Y^2=P Q(Q-P)(Q-2P),\qquad Y=yP^2 .                  \hypertarget{tag-00B2}{}\label{cdp:eq:C.BC.12}
\end{equation}
The four quadratic factors have respectively discriminants
$-4,-59,-31,-11$ and nonzero leading coefficients.
A common root of two factors would satisfy $P=Q=0$;
from $P=0$ one gets $z^2=-1$ and $Q=z+2$, which is impossible.
They are therefore pairwise coprime. Their eight simple roots
give a smooth connected genus-three curve. The chart
at \(P=0\) is the ordinary elliptic infinity chart;
choose its local parameter $\tau$ so that
$x=\tau^{-2}A(\tau)$ and $y=\tau^{-3}B(\tau)$,
where $A(0)=B(0)=1$. The equation $Px=Q$ gives
$P=Q\tau^2/A(\tau)$ and
$Y=Q^2B(\tau)\tau/A(\tau)^2$.
Since $P$ has a simple zero and $Q$ is nonzero there,
these equations express $z$ holomorphically in $\tau$
and make $Y$ a local parameter. They give the same smooth
double cover at each zero of $P$. At \(z=0\)
and \(\infty\) the values \(x=5,3\) are outside
\(\{0,1,2,\infty\}\). Thus \(D\) meets each old
elliptic boundary transversely in two points and avoids
all four existing triple points, which have elliptic
coordinate \(e_0\) or \(e_1\). The simple zeros of
\(Q\) and \(Q-P\) prove transverse intersection
with \(B_0\) and \(B_1\), respectively. These facts
also verify all its boundary intersections directly.

\item Blow up this \(D\) in the threefold. The same two
charts \(xy\), \(aXy\) prove reduced SNC at its
boundary crossings; elsewhere it is a smooth-centre
blowup in one reduced component. It has canonical
order one generically, since its ambient centre lies
in the zero-order component \(P_0\). The centre is
disjoint from \(F\), so \eqref{cdp:eq:C.BC.8}--\eqref{cdp:eq:C.BC.9} remain
unchanged on an open neighbourhood of the entire
curve, including every analytic transverse jet.

\item Initially \(e((C_0+C_1+C_2)\times E)=0\).
Apply \cref{lem:II:blowup-topology} successively with
the smooth centres $B_0,B_1,D$ and the compact reduced
support containing each centre. At every step the next
support is the full reduced inverse image. The codimension
is two, so the Euler increment is the Euler number of
that same centre. Consequently the three blowups give
\begin{equation}
             e(S)=0+e(B_0)+e(B_1)+e(D)
                         =0+2+2-4=0 .                       \hypertarget{tag-00B3}{}\label{cdp:eq:C.BC.13}
\end{equation}
The six final components have Euler numbers $0,4,4,6,6,-8$:
$P_0$ is unchanged, $P_1$ and $P_2$ each acquire four point
blowups, the first two ruled exceptional surfaces each acquire
two point blowups, and the last is ruled over the genus-three
curve $D$. The double curves have total Euler number $24$.
The first two blowups contribute $12$ in total; the last adds
$e(D)=-4$ and eight exceptional rational fibres, contributing
$-4+16=12$. The first four triple points persist, and the
eight boundary intersections of $D$ give eight further triple
points. Thus the final support has twelve triple points.

\item All blowups are along closed smooth curves in a
compact projective threefold. Their total spaces
are smooth, compact and projective. They preserve
the map outside the single fibre; properness and
connectedness follow from their connected blowup
fibres and the original connected elliptic product.
The other fibres remain unchanged products with
\(E\). Finally, every additional point blowup
induced on \(P\) is away from \(F\), so the
same composite surface blowdown and its degree-one
map on \(F\) prove the horizontal section asserted on the final zero-Euler model.

\end{enumerate}
\end{proof}

\subsection{Reduction of exceptional configurations}\label{subsec:II:vertical-elliptic-models}

\begin{definition}[Marked plane pair]\hypertarget{tag-P366}{}\label{def:II:marked-plane-pair}
\StatementRecord{def:II:marked-plane-pair}{proved}{P366}{\cref{lem:II:positive-rational-curve}}
In a smooth-total torus degeneration with full reduced
normal-crossing central fibre, a \emph{marked plane pair}
consists of a smooth component normalization $Q$, its full
reduced boundary $D_Q$, a smooth rational boundary curve
$E\subset Q$ with $E^2=1$ disjoint from every other boundary
curve, and the morphism
$\rho:Q\to\PP^2$ defined by the complete system $|E|$.
The morphism is the one in \cref{lem:II:positive-rational-curve}:
it maps $E$ isomorphically to a line, is an isomorphism near $E$,
and contracts the other boundary curves.
\end{definition}

\begin{proposition}[Elliptic parents at marked Hopf interfaces]\hypertarget{tag-01ET}{}\label{cdp:st:C.2.UN.81}
\StatementRecord{cdp:st:C.2.UN.81}{proved}{01ET}{\SeriesRef{cdp:st:C.2.UN.77}, \SeriesRef{cdp:st:C.2.UN.75}, \SeriesRef{cdp:st:C.2.UN.70}, \SeriesRef{cdp:st:C.2.UN.80}, \SeriesRef{cdp:eq:C.AP.5}, \SeriesRef{cdp:eq:C.AP.3}, \SeriesRef{cdp:eq:C.AP.6}, \SeriesRef{cdp:st:C.2.UN.74}, \SeriesRef{cdp:st:C.2.UN.31}, \SeriesRef{cdp:eq:C.AP.9}, \SeriesRef{cdp:eq:C.AP.10}, \SeriesRef{cdp:eq:C.AP.11}, \SeriesRef{cdp:st:C.2.UN.64}, \SeriesRef{cdp:eq:C.AP.2}, \SeriesRef{cdp:eq:C.AP.4}, \SeriesRef{cdp:eq:C.AP.14}; \cite[\S4.10.7--4.10.8, pp.~980--982; Lemma~4.14 and proof, pp.~983--985]{Kulikov1977}; \cite[\S8, p.~341]{Kulikov1981}}

Let \(f:X\to\Delta\) be a proper degeneration of connected
complex two-tori with smooth total space and full reduced
SNC central fibre. We also allow as the starting recorded model
the smooth-total NC endpoint obtained from such a fibre by
\ref{cdp:st:C.2.UN.77}, retaining its two labelled normalization branches at
each self-double curve. Normalize the canonical generator by
\begin{equation}
 \operatorname{div}(\Omega)=\sum k_iV_i,\qquad
 \min k_i=0,\qquad M=\max k_i.
 \hypertarget{tag-009M}{}\label{cdp:eq:C.AP.1}
\end{equation}
Suppose that the components of order \(M\) with Hopf minimal
model have already been made minimal by the record of
\ref{cdp:st:C.2.UN.77}. Mark any finite collection of their interfaces
\(J_\alpha=H_\alpha\cap P_\alpha\) for which the normal line of
\(J_\alpha\) is torsion and \(k_{P_\alpha}<M\).
Retain repeated markings on the same parent.

There is a finite overlap-compatible modification record,
unchanged over \(\Delta^*\), whose endpoint has smooth total
space and full reduced NC central fibre, such that every
designated parent is an embedded minimal fixed-Jacobian
elliptic quasi-bundle. Its full elliptic fibres are multiples
of smooth elliptic curves, its integral linear monodromy is
the identity, and every one of its boundary curves is a
smooth elliptic fibre support. In particular no triple point
meets a designated parent at the endpoint.

All maximal-order minimal Hopf components and an open
neighbourhood of their union are unchanged by the record.
In particular every marked \(J_\alpha\), its normal map and
its gluing to \(H_\alpha\) are retained. The designated parent
prime divisors survive. Other endpoint components can retain
NC self-incidences; their two normalization branches are
part of the record.

Every retained global component has canonical order between
\(0\) and the original \(M\). Whole-component removals delete
a component of positive order \(q\); their preliminary
additions, when needed, have order \(q+1\le M\), and there
are no additions at \(q=M\). The record retains its centres, local maps, normalization branches and restored
overlaps, so reversing it recovers the original family.
For a given compact family the same modifications are
performed on that compact base. Their Euler changes are
the exact normal-cone terms of \ref{cdp:st:C.2.UN.75}.

\end{proposition}

\begin{proof}\leavevmode\par

 
\begin{enumerate}[label=\arabic*.,leftmargin=0pt,labelsep=.5em]
\item \par\medskip\noindent\textbf{The exact marked parent maps.}
Apply \cref{lem:II:hopf-curves} to each marked minimal
Hopf component $H_\alpha$; its curves are disjoint smooth elliptic curves. Thus each marked curve is isolated from
the other boundaries on both sides. Its normal-product
isomorphism makes the normal on the parent torsion as well.
On the normalization of a designated parent \(P\), of order
\(p\), the marked curve has coefficient
\(M-p-1\ge0\) in the canonical divisor. \ref{cdp:st:C.2.UN.70} therefore
gives a specified point map and elliptic map
\begin{equation}
 P\xrightarrow{\pi_P} Z_P\xrightarrow{q_P}B_P,
 \qquad R^1(q_P)_*\OO_{Z_P}=\OO_{B_P},
 \qquad g(B_P)\le1.
 \hypertarget{tag-009S}{}\label{cdp:eq:C.AP.2}
\end{equation}
The point map is an isomorphism on a neighbourhood of every
marked curve. Every boundary is vertical for \(q_P\pi_P\),
and every rational boundary is exceptional for \(\pi_P\).
The exact constant-period field of \ref{cdp:st:C.2.UN.80}, equations
\eqref{cdp:eq:C.EP.1}--\eqref{cdp:eq:C.EP.11}, proves the additional inclusion
\begin{equation}
                 \operatorname{Exc}(\pi_P)\subset D_P.
 \hypertarget{tag-009T}{}\label{cdp:eq:C.AP.3}
\end{equation}
Here \(D_P\) is the full reduced boundary. It includes
both preimages of a self-double curve at an NC stage.

\begin{enumerate}[label=(\arabic*),leftmargin=1.3em,labelsep=.5em]
\item  These assertions also apply at the intermediate smooth-total
NC stages used below. Indeed the adjunction divisor
on the normalization counts those two preimages; the marked
elliptic neighbourhood stays embedded and unchanged. The
saturated ambient fields restrict as meromorphic fields to
that normalization and are holomorphic off its full boundary.
The surface maps in \ref{cdp:st:C.2.UN.70} and the direct exceptional-valuation
argument in \ref{cdp:st:C.2.UN.80} consequently apply to this smooth normalization with its
two separately marked boundary branches.

\item  The target \(Z_P\) is a minimal surface in its birational
class. Thus
\begin{equation}
 h_P=b_2(P)-b_2(Z_P)\ge0
 \hypertarget{tag-009U}{}\label{cdp:eq:C.AP.4}
\end{equation}
is exactly the number of point blowups in \(\pi_P\). The
minimal second Betti number is unchanged under the subsequent
birational maps: in the ruled case it is two, and in the
other cases the minimal model is unique. We count each
designated parent once in \(h=\sum_P h_P\).

\end{enumerate}

\item \par\medskip\noindent\textbf{The full fibre graph and the receiving degrees.}
The point map \(\pi_P\) is a sequence of point blowups of
smooth elliptic fibre supports. For one such fibre, retain
its entire reduced inverse image, whether or not its strict
elliptic component is a boundary. Its graph is an
elliptic root with rational trees. Write \(C^2=-b_C\) and let
\(n_C\) be the number of its intersections with the other
components of this entire reduced fibre. The point maps give
\begin{equation}
 n_C\le b_C+1\quad\hbox{for a rational component},\qquad
 n_C\le b_C\quad\hbox{for the elliptic component}.
 \hypertarget{tag-009V}{}\label{cdp:eq:C.AP.5}
\end{equation}
The elliptic root starts with \((b,n)=(0,0)\). A new rational
curve has \((b,n)=(1,1)\) at a smooth point and \((1,2)\) at
a node. Blowing up a smooth point of an old component raises
both \(b\) and \(n\) by one; a node blowup raises \(b\) by one
and leaves its valence unchanged. These explicit updates
prove \eqref{cdp:eq:C.AP.5} at every infinitely-near point.

\begin{enumerate}[label=(\arabic*),leftmargin=1.3em,labelsep=.5em]
\item  Choose a last point blowup of \(\pi_P\), with exceptional
curve \(E\). By \eqref{cdp:eq:C.AP.3}, \(E\) is a boundary curve.
It meets zero, one or two other boundary curves.
These are exactly its triple points in the threefold.
The complete fibre graph has at most two neighbours at this
last point, so there is no additional case.

\item  Suppose first that there are two. On the other component
\(Q\), the normal-product formula gives \(E_Q^2=-1\), because
\(E_P^2=-1\). Hence the type-II two-ruling diagram of
\ref{cdp:st:C.2.UN.75} applies. It blows down \(E\) on both adjacent surfaces
and blows up the two opposite normalized branches at their
specified points.

\item  Neither receiving branch belongs to a designated parent.
To prove this, let \(C\) be the opposite boundary on \(P\),
and write \(-b_0\) and \(n_0\) for its square and entire-fibre
valence just before the last node blowup. Afterwards its
square is \(-b_0-1\), and its number \(t_C\) of triple
points is at most \(n_0\). On the opposite receiving surface
\(R\), the exact normal map therefore gives
\begin{equation}
 C_R^2=b_0+1-t_C\ge b_0+1-n_0\ge0,
 \quad\hbox{and } C_R^2\ge1\hbox{ if }C\hbox{ is elliptic}.
 \hypertarget{tag-009W}{}\label{cdp:eq:C.AP.6}
\end{equation}
On a designated parent every rational boundary is exceptional
and has negative square; every elliptic boundary is the
strict transform of a fibre support and has nonpositive
square. Both alternatives contradict \eqref{cdp:eq:C.AP.6}. This argument
uses the individual normalized branch when the two receivers
are branches of one global component. Thus a type-II move
leaves the component counts unchanged and strictly decreases
\(h\); it does not create a self-double curve on a designated
parent.

\end{enumerate}

\item \par\medskip\noindent\textbf{One triple point and ruled removal.}
Suppose $E$ has one triple point. The normal-product map gives
$E_Q^2=0$. Apply the local ruled-surface construction in
\cite[\S4.5, p.~970]{Kulikov1977} with
$(V_1,V_2,L,S_2)=(P,Q,E,Q\cap R)$, where $R$ is the third
component through that point. The curve $E$ is the complete
ruling fibre, $Q\cap R$ is the section, and every other
boundary curve on $Q$ is vertical. Restriction of the common
canonical divisor to a ruling fibre gives
\begin{equation}
 k_R=q-1,\qquad 1\le q=k_Q\le M.
\hypertarget{tag-009X}{}\label{cdp:eq:C.AP.7}
\end{equation}
Apply the finite ruled-removal and singular-chart cleanup from
\cref{cdp:st:C.2.UN.77}, with its removed component, section
parent and fixed maximum equal to $Q,R,M$. On a projective-line
bundle the map is \cref{cdp:st:C.2.UN.74}; the nonminimal
fibres use the local maps of
\cite[\S4.10.7--4.10.8, pp.~980--982; Lemma~4.14 and proof, pp.~983--985]{Kulikov1977}
and the fibre insertions of \cite[\S8, p.~341]{Kulikov1981}.
The same chart composites identify the restored overlaps.
Each removal deletes a positive-order component and adds
only finitely many components of the next order, with none
added at $M$. Thus the component-count tuple strictly decreases
across this finite block, at whose endpoint the total space
is smooth and the full fibre is reduced NC.

\item \par\medskip\noindent\textbf{Zero triple points: the exact plane comparison.}
Suppose that \(E\) has no triple point. Its square on \(Q\)
is \(+1\), and every other boundary on \(Q\) is disjoint
from \(E\). Apply \cref{lem:II:positive-rational-curve} to the smooth normalization with $(T,C)=(Q,E)$. Its hypotheses hold at this full reduced NC stage, and $E^2=1$. It gives rationality and the complete-linear-system morphism \eqref{cdp:eq:C.AP.9}, with the same pair $(Q,E)$ in place of $(T,C)$. This morphism is an isomorphism near $E$, maps $E$ isomorphically to the line $\ell$, and contracts every other boundary because it is disjoint from $E$.

\begin{enumerate}[label=(\arabic*),leftmargin=1.3em,labelsep=.5em]
\item  The canonical restriction to \(E\) gives
\(k_Q=k_P+2\). For this marked plane pair of \cref{def:II:marked-plane-pair},
write $q=k_Q$. Its canonical divisor has the precise form
\begin{equation}
 K_Q=-3E+\sum_C a_CC,
 \qquad \sum_C a_CC=K_Q-\rho^*K_{\PP^2},
 \qquad a_C>0.
 \hypertarget{tag-009N}{}\label{cdp:eq:C.AP.10}
\end{equation}
Indeed choose on \(\PP^2\) the meromorphic two-form whose
divisor is \(-3\ell\). The ratio of the form on \(Q\)
to its pullback has divisor supported on the exceptional
curves. On \(\PP^2\) this same meromorphic function has no
divisorial zero or pole, so it is a nonzero constant. Thus
the equality in \eqref{cdp:eq:C.AP.10} is an equality of divisors,
with the positive discrepancies of the specified point map.
It shows that every exceptional prime is a boundary, and that there is no additional boundary of
coefficient zero. The only coefficient on the marked curve
is \(-3\). A self-double curve away from that marking would
have coefficient \(-1\), so such a pair stays embedded at
the smooth-total NC stages under consideration.

\item  Choose a last point blowup of \(\rho\). Its exceptional
boundary curve \(L\) lies away from \(E\). At a smooth point
of a discrepancy component of coefficient \(a>0\) its new
coefficient is \(a+1\); at a node of coefficients \(a,b>0\)
it is \(a+b+1\); at a point off the discrepancy divisor it
is one. These are the point-map identities for
\(K_Q+3E\). They give the following exhaustive alternatives:
\begin{equation}
\begin{array}{c|c|c}
 \text{triple points of }L&\text{order of its neighbour}
                       &\text{next operation}\\ \hline
 0&q+2&\text{another marked plane pair}\\
 1&q+a+2\ge q+3&\text{a ruled-removal block}\\
 2&q+a+b+2&\text{a type-II move}.
\end{array}
 \hypertarget{tag-009O}{}\label{cdp:eq:C.AP.11}
\end{equation}
In the last row the two receiving orders are
\(q+a+1\) and \(q+b+1\), both at least \(q+2\). The type-II
move decreases \(b_2(Q)\) by one; every other increase in
surface second Betti number occurs at a strictly higher
order. Its centre and all its points on \(Q\) are away from
the marked \(E\), so \eqref{cdp:eq:C.AP.9} persists on the modified \(Q\).
In the first row the same complete linear system constructs
the next marked plane pair, at a strictly greater order.

\item  This yields a finite progress block. Until its first whole
component removal, let
\begin{equation}
                    b_j=\sum_{k_V=j}b_2(V^\nu),
                    \qquad 0\le j\le M,
 \hypertarget{tag-009P}{}\label{cdp:eq:C.AP.12}
\end{equation}
where the sum uses the smooth normalizations. Each
type-II step in \eqref{cdp:eq:C.AP.11} strictly decreases
\((b_0,\ldots,b_M)\) lexicographically. Passing to another
marked plane pair makes no modification and strictly raises
its order. There can be at most \(M\) such successive calls
before a step. Since lexicographic order on this
fixed finite product of nonnegative integers is well founded,
these steps cannot continue indefinitely without a component
removal. At a marked plane pair with no exceptional curve,
the surface is \(\PP^2\), its sole boundary is \(E=\ell\),
and the principal full-fibre restriction gives
\begin{equation}
                 N_{Q/X}=\OO_{\PP^2}(-1).
 \hypertarget{tag-009Q}{}\label{cdp:eq:C.AP.13}
\end{equation}
The point contraction of \ref{cdp:st:C.2.UN.64} removes it to a smooth
point. The adjacent surface is its ordinary point blowup,
and no other component passes through that point. In the
one-triple alternative the positive ruled-removal block
already supplies a component removal.

\item  Stop this progress block at its first whole-component
removal, and perform the finite sigma cleanup just described.
Its component-count tuple strictly decreases. At this
completed block's smooth-total endpoint, reapply the
designated-parent maps \eqref{cdp:eq:C.AP.2}--\eqref{cdp:eq:C.AP.4}. 

\end{enumerate}

\item \par\medskip\noindent\textbf{Preserved maximal Hopf neighbourhoods.}
Every type-II centre used on a designated parent is vertical
for \eqref{cdp:eq:C.AP.2}, hence disjoint from its marked elliptic curves.
An opposite curve of a type-II move cannot lie on a minimal
maximal Hopf: a rational curve is impossible there, and the
elliptic alternative has the positive square \eqref{cdp:eq:C.AP.6}.
For the plane progress moves all opposite boundary curves
are rational. They likewise cannot meet a minimal Hopf.

\begin{enumerate}[label=(\arabic*),leftmargin=1.3em,labelsep=.5em]
\item  A ruled component removed in one of the blocks has
one boundary section, whose neighbouring order is \(q-1\),
and all other boundary is rational and vertical. It cannot
meet a minimal Hopf of order \(M\): that Hopf contains no
rational curve, while its being the section parent would
require \(q=M+1\). The same statement applies to every
cleanup removal and missing-fibre insertion. Plane point
contractions also have rational boundary. Thus these closed centres and local sigma charts avoid the entire
union of maximal minimal Hopfs. Their restored overlap maps
are the identity on an open neighbourhood of that union.

\item  The designated parent prime divisors cannot be removed:
each meets one of these unchanged Hopfs. Their marked
elliptic neighbourhoods therefore persist. At every smooth
NC stage the exact maps \eqref{cdp:eq:C.AP.2}--\eqref{cdp:eq:C.AP.4} can be applied
again to their current normalizations. 

\end{enumerate}

\item \par\medskip\noindent\textbf{The fixed finite termination potential.}
Let \(c_j\) count the retained global components of order
\(j\), for \(1\le j\le M\). Use the block-level potential
\begin{equation}
                  (c_1,c_2,\ldots,c_M,h).
 \hypertarget{tag-009R}{}\label{cdp:eq:C.AP.14}
\end{equation}
All coordinates are nonnegative integers. A direct type-II
step on a last designated-parent point leaves the component counts
unchanged and decreases \(h\) by \eqref{cdp:eq:C.AP.6}. Either of the
other last-point alternatives supplies a finite progress
block with a strict decrease of the first changed component count.
Any finite defect increase inside that whole block is
irrelevant to this lexicographic decrease. The missing-fibre
and cleanup comparisons keep all retained orders at most
the original \(M\); their auxiliary common blowups are not
components of these reduced degeneration stages. Hence \eqref{cdp:eq:C.AP.14} has no infinite
strictly descending sequence.

\begin{enumerate}[label=(\arabic*),leftmargin=1.3em,labelsep=.5em]
\item  When this procedure stops, a positive designated-parent
defect would still have a last point curve. \ref{cdp:st:C.2.UN.80} and the
three exhaustive alternatives would supply another strict
decrease. Thus every such defect is zero. Each designated
normalization is the minimal elliptic model in
\eqref{cdp:eq:C.AP.2}, and every boundary is a fibre support. All new
self-incidences created by the type-I/type-II and ruled
sigma charts are rational exceptional curves. None can
remain on these parents, since their rational boundary
would have to be exceptional for a point map of defect
zero. No elliptic self-incidence was introduced: a ruled
removal has only one horizontal boundary section, and all
its other changes are at its rational fibre curves. The
designated parents are consequently embedded, with disjoint
elliptic boundaries and no triple point.

\item  Every arrow used above is one of the two-ruling,
curve, point, or restored sigma-atlas maps already specified.
The inverse record, including all normalization branches
and overlap restrictions, reconstructs the given family.
The proper graph gluing and annular comparison of \ref{cdp:st:C.2.UN.75}
give the same record on a compact base, and its normal-cone
Euler formula applies at every retained step. 

\end{enumerate}

\end{enumerate}
\end{proof}

\subsection{Maximal canonical order and Hopf leaves}\label{subsec:II:maximal-hopf-leaves}

\begin{proposition}[Every positive maximal Hopf after minimalization is a torsion-normal leaf]\hypertarget{tag-01EU}{}\label{cdp:st:C.2.UN.82}
\StatementRecord{cdp:st:C.2.UN.82}{proved}{01EU}{\cref{cdp:st:C.2.UN.75,cdp:st:C.2.UN.77,cdp:st:C.2.UN.70,cdp:st:C.2.UN.19,lem:II:negative-kodaira-models,lem:II:canonical-boundary-blowup,lem:II:hopf-curves,lem:II:anticanonical-bisection,lem:II:reduced-anticanonical-hopf}; \eqref{cdp:eq:C.ML.3}, \eqref{cdp:eq:C.ML.5}, \eqref{cdp:eq:C.ML.6}, \eqref{cdp:eq:C.ML.8}}

Let \(f:X\to\Delta\) be a proper degeneration of connected
complex two-tori from a smooth threefold, with full reduced
NC central fibre and smooth normalization components. Retain
every normalized branch of a self-double curve, as in \cref{cdp:st:C.2.UN.75}.
For the primitive canonical evaluation write
\begin{equation}
       \operatorname{div}(\Omega)=\sum_i k_iV_i,
       \qquad \min_i k_i=0,\qquad M=\max_i k_i .
\hypertarget{tag-00QR}{}\label{cdp:eq:C.ML.1}
\end{equation}
On the normalization \(T\) of a component of order \(M\),
suppose that one boundary branch \(C\) is smooth elliptic,
is disjoint from all its other boundary branches, and has
coefficient one in its induced anticanonical divisor.
Then every central component has order zero and there
are no triple points.

Consequently consider the endpoint of \cref{cdp:st:C.2.UN.77} for
an original full reduced SNC family with \(M>0\).
Every Hopf component \(H\) of this maximal order is a
terminal component. Its boundary is one smooth elliptic
curve \(C=H\cap P\), and its exact canonical and normal
data satisfy
\begin{equation}
 \begin{gathered}
       -K_H=aC,\qquad a=M-k_P+1\ge2,\\
       N_{C/H}^{\otimes(a-1)}\simeq\OO_C,
       \qquad N_{C/P}\simeq N_{C/H}^{-1},\\
       k_P\le M-1,\qquad
       \operatorname{coeff}_C(K_P)=a-2\ge0 .
 \end{gathered}
\hypertarget{tag-00QT}{}\label{cdp:eq:C.ML.2}
\end{equation}
The normal-line equality uses the marked local branches
on \(P\) when another part of that component is nonnormal.
Each such interface therefore supplies precisely the
marked torsion-normal, nonnegative-coefficient data of
\cref{cdp:st:C.2.UN.70} on that normalization. The curve, its normal
maps and its surrounding neighbourhood are those at the
endpoint of the finite record of \cref{cdp:st:C.2.UN.77}.
\end{proposition}

\begin{proof}

\begin{enumerate}[label=\arabic*.,leftmargin=*]
\item \par\medskip\noindent\textbf{The induced divisor at the maximal order.}
The normalization adjunction map for the reduced NC
fibre, including both branches over every self-double
curve, gives the meromorphic canonical section
on \(T\) with
\begin{equation}
      -K_T=D_T=\sum_B (M-k_{\mathrm{other}(B)}+1)B.
\hypertarget{tag-00QU}{}\label{cdp:eq:C.ML.3}
\end{equation}
Every displayed coefficient is a positive integer.
For a self-double branch its coefficient is one.
This is the restriction of the common volume form and
the principal full-fibre trivialization; together these maps
fix the line isomorphism and the displayed divisor.
The support of \(D_T\) is exactly the full marked
boundary. The curve \(C\) in the statement is one
isolated component of this effective divisor, with
coefficient one.

\item We first prove the following surface consequence of
these data: the whole divisor \(D_T\) is reduced and
is a disjoint union of smooth elliptic curves. The
minimal model of \(T\) is ruled or Hopf. Indeed the
nonzero effective anticanonical divisor makes every
positive pluricanonical section vanish: through
\(K_T^{\otimes n}\simeq\OO_T(-nD_T)\), such a
section would be a holomorphic function on the compact
connected surface vanishing on \(D_T\). Thus
\(\kappa(T)=-\infty\). Apply \cref{lem:II:negative-kodaira-models} to the compact smooth
minimal model. Its ruled/class-VII alternatives and the component
restriction \cref{cdp:st:C.2.UN.19} give the stated ruled or Hopf cases.

\par\medskip\noindent\textbf{Rational surfaces: the entire divisor is the marked elliptic curve.}
For a rational \(T\), the exact divisor sequence is
\begin{equation}
 0\longrightarrow K_T\longrightarrow\OO_T
          \longrightarrow\OO_{D_T}\longrightarrow0.
\hypertarget{tag-00QV}{}\label{cdp:eq:C.ML.4}
\end{equation}
Here \(H^0(K_T)=0\) and
\(H^1(K_T)=H^1(\OO_T)^*=0\) by Serre duality.
Consequently \(H^0(\OO_{D_T})=\C\). An effective
divisor with disconnected support has independent
constant functions on its disjoint components, so
its support is connected. Since \(C\) is isolated
in that support and has coefficient one, the exact
divisor is \(D_T=C\).

\par\medskip\noindent\textbf{Keeping the marked curve through the surface point maps.}
In the remaining cases, take a sequence of point blowdowns
\(\pi:T\to T_0\) to a smooth minimal ruled or Hopf
surface. Push the meromorphic canonical section
through their common function-field identification.
The resulting anticanonical divisors remain effective.
At a reverse point blowup $\sigma$, apply
\cref{lem:II:canonical-boundary-blowup} to the effective
divisors before and after the map. Their canonical-boundary
lines are trivial; its Jacobian comparison, for the pushed
meromorphic section, gives the divisor identity
\begin{equation}
                 D_{\mathrm{new}}=\sigma^*D_{\mathrm{old}}-E.
\hypertarget{tag-00QW}{}\label{cdp:eq:C.ML.5}
\end{equation}
The curve \(C\) is never exceptional, since it is
elliptic. Its image remains a smooth isolated
coefficient-one component at every blowdown. To
check this assertion at a last exceptional curve
\(E\), first suppose its coefficient in the
effective divisor is positive. It is then disjoint
from \(C\), so the contraction is an isomorphism
near \(C\). If its coefficient is zero, \eqref{cdp:eq:C.ML.5}
says that the old divisor has total multiplicity
exactly one at the centre. When \(E\) meets \(C\),
this forces the image of \(C\) to be the only
positive branch there, of multiplicity one. It is
smooth and no other boundary image meets it.
Equivalently, a curve of multiplicity \(r\) at
the centre has total intersection \(r\) with
the exceptional divisor; here \(r=1\).
This proves the assertion by induction through
the specified point maps. Write \(C_0=\pi(C)\).

\par\medskip\noindent\textbf{Irrational ruled surfaces: two coefficient-one sections.}
Let \(p:T_0\to B\) be the minimal ruling. The
elliptic curve \(C_0\) cannot be contained in a
rational ruling fibre. Its map to \(B\) is finite,
so Riemann--Hurwitz gives \(g(B)\le1\).
The rational case has already been proved; hence
\(B\) is elliptic and that finite map is \'etale.
For a ruling fibre \(F\),
\begin{equation}
                 D_{T_0}\cdot F=-K_{T_0}\cdot F=2.
\hypertarget{tag-00QX}{}\label{cdp:eq:C.ML.6}
\end{equation}
Every positive vertical curve would meet \(C_0\),
contrary to its isolation. Thus there is no such
vertical part. If \(C_0\) has degree two over
\(B\), \eqref{cdp:eq:C.ML.6} gives \(D_{T_0}=C_0\).
Apply \cref{lem:II:anticanonical-bisection} with
$(V,B,q,C,h)=(T_0,B,p,C_0,p|_{C_0})$.
The same surface divisor $D_{T_0}=C_0$ gives
$K_{T_0}(C_0)\simeq\OO_{T_0}$, while $C_0$ is smooth,
connected, elliptic and of degree two over the elliptic base.
Under these identifications \eqref{cdp:eq:C.ML.7} is its residue
sequence and \eqref{cdp:eq:C.EL.7} is the Leray-identified finite
pullback. The nonzero trace composite and zero Serre-dual
connecting map exclude this bisection.

\item The remaining possibility in \eqref{cdp:eq:C.ML.6} is
\begin{equation}
                     D_{T_0}=C_0+C'_0,
\hypertarget{tag-00QZ}{}\label{cdp:eq:C.ML.8}
\end{equation}
with two disjoint degree-one sections, each of
coefficient one. A degree-one irreducible curve
over the smooth ruling base is the graph of
its section and is smooth. Thus both curves
are smooth elliptic.

\item Pull \eqref{cdp:eq:C.ML.8} through the reverse point
blowups. Every intermediate divisor is effective:
a negative coefficient would remain negative
on its later strict transform. A centre away
from the positive support would create coefficient
\(-1\) in \eqref{cdp:eq:C.ML.5}, so no such centre occurs.
Starting with the two disjoint coefficient-one
curves, a centre on either has total positive
multiplicity one and creates coefficient zero.
The positive strict transforms remain disjoint
and smooth. Induction shows that no positive
exceptional coefficient ever appears. Therefore
\(D_T\) consists exactly of their two smooth
elliptic strict transforms, still with coefficient
one. All these are comparisons through the surface morphism \(\pi\).

\par\medskip\noindent\textbf{The effective divisor on the Hopf minimal model.}
For the same point map $\pi:T\to T_0$, the pushed-forward
divisor $D_{T_0}$ is nonempty and effective, and
$K_{T_0}(D_{T_0})\simeq\OO_{T_0}$ through the pushed meromorphic
canonical section. Its support consists of disjoint smooth
elliptic curves by \cref{lem:II:hopf-curves}.
The preceding point-map comparison preserves the coefficient-one
curve $C_0$. Apply \cref{lem:II:reduced-anticanonical-hopf} to
$(H,D)=(T_0,D_{T_0})$. Its effective-divisor dichotomy and this
coefficient one force the reduced alternative: every component
of $D_{T_0}$ has coefficient one. The finite unramified cover,
positive-power cover and weight equation used here are exactly
those of that owner, with \eqref{cdp:eq:C.ML.9} its resonant factorization.

\item The same point-map induction following \eqref{cdp:eq:C.ML.8}
now applies to \(T\): every positive coefficient
is one, positive branches never meet, and the
additional exceptional coefficients are zero.
It proves the asserted surface consequence for
every blown-up Hopf component too.

\par\medskip\noindent\textbf{Propagation through the central support.}
The preceding arguments apply to the smooth
normalization of every maximal component carrying
one isolated coefficient-one elliptic branch.
They show that all its boundary branches have
these same properties. By \eqref{cdp:eq:C.ML.3}, the component
across each one again has order \(M\). Since the
branch has no triple point, its counterpart on
that normalization is again an isolated smooth
elliptic branch of coefficient one. The argument
can therefore be repeated across every incident
edge. Self-incidences retain both normalization
branches and introduce no new component order.

\item The finite incidence graph of the connected
central support is connected. The resulting
set of components of order \(M\) has no boundary
edge to its complement; it is consequently the
whole support. Every normalized boundary is
elliptic and isolated, so no triple point occurs.
Primitivity of the original direct-image generator
gives \(\min_i k_i=0\), hence \(M=0\).
This proves the first assertion.

\par\medskip\noindent\textbf{Maximal Hopf leaves.}
At the endpoint of \cref{cdp:st:C.2.UN.77} the maximal Hopf components
are embedded minimal Hopf surfaces, and
their boundary curves are disjoint smooth elliptic
curves. The comparison preserves their orders and
the primitive canonical direct-image line. For
\(M>0\), the first assertion excludes coefficient
one on any such boundary.

\item Apply \cref{lem:II:reduced-anticanonical-hopf} to
the same minimal Hopf $H$ and the effective anticanonical
divisor $D_H=-K_H$ from \eqref{cdp:eq:C.ML.3}.
Its support is the nonempty disjoint elliptic boundary.
The first assertion has excluded coefficient one for $M>0$,
so the reduced alternative is impossible. The other alternative
gives exactly $D_H=aC$ with $a\ge2$.
The coefficient formula then gives the unique neighbour's
order $k_P=M-a+1$.

\item Adjunction on this elliptic curve gives
\begin{equation}
     \OO_C\simeq K_C=(K_H+C)|_C
                          \simeq N_{C/H}^{\otimes(1-a)}.
\hypertarget{tag-00QS}{}\label{cdp:eq:C.ML.10}
\end{equation}
There are no triple points on \(C\), so the
full reduced-fibre trivialization gives
\(N_{C/H}\otimes N_{C/P}\simeq\OO_C\).
Finally the exact coefficient on \(P\) is
\(M-k_P-1=a-2\). These are precisely all the
claims in \eqref{cdp:eq:C.ML.2}. All these identifications use the retained normalized
branches and the specified point maps. In particular the
adjunction isomorphism on $C$ and the full reduced-fibre
normal product are the ones on the given endpoint. The finite
modification record of \cref{cdp:st:C.2.UN.77} supplies its
comparison with the original family.

\end{enumerate}
\end{proof}

\Needspace{10\baselineskip}
\subsection{Reduction above the maximum Hopf order}\label{subsec:II:negative-interface-reduction}

\begin{proposition}[The highest Hopf order can be made the global maximum]\hypertarget{tag-01EW}{}\label{cdp:st:C.2.UN.84}
\StatementRecord{cdp:st:C.2.UN.84}{proved}{01EW}{\SeriesRef{cdp:st:C.2.UN.82}, \SeriesRef{cdp:st:C.2.UN.75}, \SeriesRef{cdp:st:C.2.UN.77}, \SeriesRef{cdp:eq:C.TK.4}, \SeriesRef{cdp:eq:C.TK.5}, \SeriesRef{cdp:st:C.2.UN.19}, \SeriesRef{cdp:eq:C.TK.3}, \SeriesRef{cdp:eq:C.TK.6}; \cite[\S7, p.~341]{Kulikov1981}; \cite[\S4.8.5--4.8.12, pp.~973--976; \S4.16.2--4.16.6, pp.~986--988]{Kulikov1977}; \cite[Lemma~2.19, p.~965]{Kulikov1977}}

Let \(f:X\to\Delta\) be a proper degeneration of connected
complex two-tori with smooth total space and full reduced
SNC central fibre. Use its primitive canonical
generator and write
\begin{equation}
 \operatorname{div}(\Omega)=\sum_i k_iV_i,\qquad
 \min_i k_i=0,\qquad R=\max_i k_i.
 \hypertarget{tag-013J}{}\label{cdp:eq:C.TK.1}
\end{equation}
Let \(\mathcal H\) be the set of prime components whose
smooth minimal model is Hopf, and suppose that
\begin{equation}
                  h=\max_{i\in\mathcal H}k_i>0.
 \hypertarget{tag-013K}{}\label{cdp:eq:C.TK.2}
\end{equation}
There is a finite overlap-compatible modification record,
over the same disc and isomorphic on its punctured part,
whose endpoint has smooth total space and full reduced
NC fibre, such that its global maximum canonical order
is exactly \(h\). Every original Hopf-birational prime
survives with its original order, and no new such prime
appears. Every component of order \(h\) with Hopf minimal
model can, in the same finite record, be made an embedded
minimal Hopf surface. These maximal Hopfs are the
torsion-normal elliptic leaves of \ref{cdp:st:C.2.UN.82}.

At all retained global stages the canonical orders lie
between zero and the original \(R\). A whole-component
removal has positive order \(q\); its possible added
components have order \(q+1\), bounded by the current
maximum, and there are no additions when \(q\) is that
maximum. The termination potential counts every positive order,
including orders at most \(h\).

The record retains its centres, normalization
branches, local sigma maps and restored overlaps.
Reversing it recovers the given analytic family. On
an original compact family it gives the same modification over its compact base. Its Euler changes
are the signed normal-cone terms of \ref{cdp:st:C.2.UN.75}.
\end{proposition}

\begin{proof}

\begin{enumerate}[label=\arabic*.,leftmargin=*]
\item \par\medskip\noindent\textbf{Hopf orders and maximal components.}
For the proper bimeromorphic map between each pair of smooth
endpoints, apply \cref{lem:II:nonuniruled-divisors} to the
strict transforms of the Hopf-birational primes.
Their non-uniruledness follows from \cref{lem:II:hopf-curves}
through their same minimal point maps: all rational curves
are among the finitely many exceptional curves. Thus these
primes survive. The canonical pullback
\eqref{cdp:eq:C.KG.6}, for this same bimeromorphic map, preserves
their orders. Each added prime has its ruling by complete
rational exceptional fibres, so it has no Hopf minimal model.
Hence the indexed set of Hopf canonical orders is unchanged.
In particular its highest value remains \(h\).

\item Consider a current smooth-total NC stage with
maximum \(R'>h\). On the smooth normalization
\(T\) of a maximal component, adjunction is
\begin{equation}
       -K_T=\sum_C (R'-k_{\mathrm{other}(C)}+1)C.
 \hypertarget{tag-013N}{}\label{cdp:eq:C.TK.5}
\end{equation}
Both preimages of a self-double curve are counted;
their coefficients are one. The divisor is nonzero:
the connected support contains a component of Hopf
order \(h<R'\), so no maximal component is an
isolated component of the whole support. Every
coefficient in \eqref{cdp:eq:C.TK.5} is positive. Multiplying a nonzero pluricanonical section by the
corresponding power of the anticanonical section in
\eqref{cdp:eq:C.TK.5} would give a nonzero holomorphic function
on the compact connected surface with a nonempty zero divisor.
This is impossible.
Thus \(\kappa(T)=-\infty\).

\item The component theorem \ref{cdp:st:C.2.UN.19} makes a
nonruled component of this Kodaira dimension
Hopf-birational. No such component has order
\(R'>h\). Hence every current maximal
normalization is ruled or rational, including
the possible plane. Controlled
singular sigma charts are first handled by
their specified last-local ruled branches;
their smooth normalization maps are the input to
\cref{cdp:st:C.2.UN.77}.


\item \par\medskip\noindent\textbf{The precise input to the component selection.}
Apply the component selection in \cite[\S7, p.~341]{Kulikov1981}
with $m=1$, $r_i=k_i$ and its current degeneration equal to
the current family over $\Delta$. Its premise is that the
components of maximal canonical order are ruled. The procedure
in \cite[\S8, p.~341]{Kulikov1981} inserts the missing ruling-fibre
components to connect the chosen boundary. The preceding argument establishes its maximal-component
premise at every stage with maximum greater than $h$.

\item The boundary
analysis at the maximal ruled components in
\cite[\S4.8.5--4.8.12, pp.~973--976; \S4.16.2--4.16.6, pp.~986--988]{Kulikov1977}
produces either a plane with one line or a ruled
normalization with one section and vertical boundary.
The two-section cases are converted by the specified
finite type-I/type-II moves along their section
chains or cycles; a rational tree gives the complete
one-triple ruling curve. At a singular sigma chart,
the last locally blown-up branch of largest order
supplies the same one-section configuration. These
are the two-ruling, complete-curve and
sigma-atlas maps recorded in \ref{cdp:st:C.2.UN.75} and \ref{cdp:st:C.2.UN.77}.

\item For the selected ruled component of order \(q\), restriction of the common canonical
section to a generic complete ruling fibre gives
the section coefficient \(-2\), hence
\begin{equation}
       k_{\mathrm{section\ parent}}=q-1.
 \hypertarget{tag-013L}{}\label{cdp:eq:C.TK.3}
\end{equation}
For a selected plane, the line coefficient \(-3\)
gives the parent order \(q-2\). Consequently every
removed component has positive order. When the selected ruled
component has disconnected boundary, the procedure of \cite[\S8, p.~341]{Kulikov1981} inserts the missing fibre components.
Their curve-blowup Jacobians give order
\(q+1\). If \(q\) is the current maximum, the
effective anticanonical one-section comparison
of \cite[Lemma~2.19, p.~965]{Kulikov1977} makes the boundary
connected and no insertion is made. In every case
\begin{equation}
 \begin{split}
 &\text{one component of order }q>0\text{ is removed},\\
 &\text{only finitely many components of order }q+1
       \text{ may be added},\\
 &q+1\le R_{\mathrm{current}}\text{ whenever an addition occurs}.
 \end{split}
 \hypertarget{tag-013M}{}\label{cdp:eq:C.TK.4}
\end{equation}
We count the global component primes of the reduced
degeneration stages in \eqref{cdp:eq:C.TK.4}.

\item The local contractions are followed by the exact
restoration of the overlap, using the same composite
map. The proper-graph comparison in \eqref{cdp:eq:C.KG.2} and
\eqref{cdp:eq:C.KG.6} therefore gives a bimeromorphic map over the same disc.
If singular sigma charts remain, we use the
largest last locally blown-up branch and its finite
descendant cleanup from \ref{cdp:st:C.2.UN.77}. Such branches are
ruled by their complete exceptional fibres.
Its component choices and maps have the positive-level updates
\eqref{cdp:eq:C.TK.4}. It ends at a smooth-total NC stage with
the labelled normalization branches retained.

\par\medskip\noindent\textbf{The stopped finite-level induction.}
For each \(1\le j\le R\), let \(c_j\) be the
number of retained global components of order
\(j\), and order the tuples
\begin{equation}
                     (c_1,c_2,\ldots,c_R)
 \hypertarget{tag-013O}{}\label{cdp:eq:C.TK.6}
\end{equation}
lexicographically, with the smaller index first.
Each whole-component block \eqref{cdp:eq:C.TK.4} strictly
decreases this tuple: it removes a component of
order \(q\), leaves all smaller orders unchanged,
and only increases the next order. Type-I and
type-II preparatory moves preserve the component
counts. The finite-level induction in \cref{cdp:st:C.2.UN.75},
with its counted objects these components and its upper
level equal to the original $R$, proves termination. The finite cleanup
blocks obey the same comparison.

\item We continue only while the maximum exceeds
\(h\), and finish the finite cleanup belonging
to the last selected block. If the maximum
still exceeded \(h\), \eqref{cdp:eq:C.TK.5} would supply
another selection. The decreasing tuple
therefore forces a finite stop with maximum
at most \(h\). The retained Hopf prime of
order \(h\) forces equality. No order-zero
prime has been removed, by \eqref{cdp:eq:C.TK.3} and its
plane counterpart, so the transported canonical
generator remains primitive with minimum zero.

\par\medskip\noindent\textbf{The maximal Hopf components at the stopping stage.}
The original central components were embedded.
The record has introduced no elliptic self-double
curve on a surviving Hopf component. Type-I and
type-II changes create incidences on their
rational exceptional curves. A removed ruled
component has one horizontal boundary section;
its generic contraction cannot identify two
elliptic horizontal branches. Its other changes
are on its rational fibre curves, and its
restored overlaps retain the original elliptic
maps. Point contractions likewise only change
the indicated rational boundary configuration.

\item At the stopping stage apply \cref{lem:II:hopf-maximal-boundary}
with $M=h$, $T$ the normalization of this order-$h$ Hopf
component, and its point map to the same Hopf minimal model.
The canonical divisor is \eqref{cdp:eq:C.TK.5} at $R'=h$;
the two normalized self-double branches are the ones
preserved by the local charts. The lemma excludes rational
self-double curves and triples of three branches of this
component. The preceding neighbourhood comparison excludes
elliptic self-double curves. These maximal Hopfs are therefore
embedded at the stopped NC stage.

\item All the prerequisites of the last-point and
finite-cleanup proof of \ref{cdp:st:C.2.UN.77} now hold on this
recorded stage. Its input consists of these smooth normalizations,
the labelled NC branches and the embedded maximal Hopfs. Apply those exact moves to remove
their point defects. The retained global orders
stay at most \(h\), every Hopf prime and order
is preserved, and the endpoint maximal Hopfs
are smooth minimal surfaces. \ref{cdp:st:C.2.UN.82} then supplies
their single elliptic boundaries and exact
torsion-normal line maps.

\item Finally the local proper graphs glue with the
recorded overlap isomorphisms, exactly as in
\ref{cdp:st:C.2.UN.75}. They give the punctured-family isomorphism
and the inverse record on its endpoint
neighbourhood. Annular gluing to the unchanged
outside performs the same operation on an
original compact family. The signed normal-cone
Euler formula \eqref{cdp:eq:C.KG.3}--\eqref{cdp:eq:C.KG.4} applies to
every retained arrow. 

\end{enumerate}
\end{proof}

\section{Root descent and finite reduction}\label{sec:II:root-descent-finite-reduction}

\subsection{Canonical-flat root descent}\label{subsec:II:canonical-flat-root-descent}

\begin{proposition}[Root base change and birational descent from canonical-flat cores]\hypertarget{tag-01EP}{}\label{cdp:st:C.2.UN.78}
\StatementRecord{cdp:st:C.2.UN.78}{proved}{01EP}{\cref{prop:II:cyclic-root-normalization,cdp:st:C.2.UN.75,cdp:st:C.2.UN.39}}

Let $f:X\to\Delta_t$ be a proper degeneration of connected
complex two-tori from a smooth threefold, with arbitrary
central support $S=f^{-1}(0)_{\rm red}$. Take an embedded resolution $\pi:Y\to X$ by smooth central
point and curve blowups, and write
\begin{equation}
 \operatorname{div}(f\pi)=\sum_i m_iD_i,
 \qquad \operatorname{div}(\pi^*\Omega)=\sum_i k_iD_i,
 \qquad r_i=\frac{k_i+1}{m_i},
 \qquad N=\operatorname{lcm}_i(m_i),
\hypertarget{tag-00ZQ}{}\label{cdp:eq:C.RK.1}
\end{equation}
where $\Omega$ generates $f_*K_X$. The $D_i$ are the
smooth components of the resolved SNC support. Apply \cref{prop:II:cyclic-root-normalization} with this
resolution $\pi$ and this exponent $N$. Let $Y_N$, $Z$, $q$,
$\rho$ and $G=\mu_N$ be its maps in \eqref{cdp:eq:C.RK.2}.
Then $Z$ is smooth with full reduced SNC central fibre,
and every component added by $q$ is ruled or rational.

The following is an explicit finite-cover condition
on this resolution record: every strict root-cover
component whose smooth minimal model is Hopf lies
over an index $i$ with
\begin{equation}
                            r_i=\min_j r_j.
\hypertarget{tag-0105}{}\label{cdp:eq:C.RK.3}
\end{equation}
For exactly the records satisfying \eqref{cdp:eq:C.RK.3}, the
semistable root $Z$ admits the finite compatible
canonical-flat reduction of \cref{cdp:st:C.2.UN.75}. Denote its
comparison by $\phi:Z\dashrightarrow X_0$.
The flat endpoint has one of the three complete
core types in \cref{cdp:st:C.2.UN.75}, and it carries the finite \emph{birational} action
\begin{equation}
                 \gamma_0=\phi\gamma\phi^{-1}
                 \quad (\gamma\in G).
\hypertarget{tag-0106}{}\label{cdp:eq:C.RK.4}
\end{equation}
Its group laws hold as equalities of meromorphic
maps; it acts over $s\mapsto\zeta s$ and is regular
on the punctured family. Use the full record
$\phi$, the toric coarsening $q$, and $\pi$.
Their inverse and quotient operations reconstruct
the given original degeneration exactly:
\begin{equation}
 X_0\ \dashrightarrow\ Z\ \xrightarrow{q}\ Y_N
       \ \longrightarrow\ Y_N/G=Y
       \ \xrightarrow{\pi}\ X.
\hypertarget{tag-0107}{}\label{cdp:eq:C.RK.5}
\end{equation}
The inverse record first restores the regular action on $Z$;
the finite quotient is then taken on $Y_N$.
Conversely the same operations on any flat-core record with these regularization,
root-coarsening and smooth blowdown maps produce
a torus degeneration in the stated branch.
These maps give finite generation for the resolution-and-root
branch defined by \eqref{cdp:eq:C.RK.3}.

Use the strata, integers $d_I,\nu_I$, and Euler corrections
$R_N,\Delta_\pi$ of \eqref{cdp:eq:C.RK.6} for this resolution.
Let
$\Delta_K$ be the signed normal-cone defect of
the forward \cref{cdp:st:C.2.UN.75} record from $Z$ to $X_0$.
Then the exact original-support formula is
\begin{equation}
                e(S)=e(S_0)-\Delta_K-R_N-\Delta_\pi.
\hypertarget{tag-010A}{}\label{cdp:eq:C.RK.8}
\end{equation}
Thus zero Euler number is equivalent to
$\Delta_K+R_N+\Delta_\pi=e(S_0)$, where the
right side is zero for the smooth or marked
cycle core, and $2[N_0:\Gamma_0]$ for the
toric core. The symbols $N_0,\Gamma_0$ here
are its two-dimensional lattices, separate
from the root exponent $N$ in \eqref{cdp:eq:C.RK.1}.
Formula \eqref{cdp:eq:C.RK.8} agrees with the universal
original-support formula of \cref{cdp:st:C.2.UN.39} by the
root and blowup maps.

The test \eqref{cdp:eq:C.RK.3} is in fact intrinsic to
the original degeneration. Write its original
full fibre and generator orders as
$\sum_\alpha n_\alpha S_\alpha$ and
$\sum_\alpha a_\alpha S_\alpha$.
Use the finite normal covers $C_\alpha\to T_\alpha$
and the density threshold $\tau$ of
\eqref{cdp:eq:C.RK.18}--\eqref{cdp:eq:C.RK.19}, constructed by
\cref{prop:II:cyclic-root-normalization} on this same original family.
Then \eqref{cdp:eq:C.RK.3} is equivalent to
\begin{equation}
 \frac{a_\alpha+1}{n_\alpha}=\tau
 \quad\text{for every }\alpha\text{ for which a connected
 component of }C_\alpha\text{ has Hopf minimal model}.
\hypertarget{tag-0102}{}\label{cdp:eq:C.RK.20}
\end{equation}
Consequently its truth is independent
of the embedded resolution, its
semistable triangulation and the
chosen root exponent divisible by
all resolved multiplicities.

For an original compact family the record can
be made on a finite compact cover of its base,
and \eqref{cdp:eq:C.RK.5} recovers the original compact
family. Reverse generation on compact core records with the displayed finite action
and coarsening likewise remains compact.
The input includes the endpoint analytic neighbourhood and
all the transition maps in the inverse record.
\end{proposition}

\begin{proof}
\begin{enumerate}[label=\arabic*.,leftmargin=*]
\item Apply \cref{prop:II:cyclic-root-normalization} with
$(f,X,\pi,Y,N)$ as in \eqref{cdp:eq:C.RK.1}. Its root charts,
height-one subdivisions, compact support and canonical generator
are exactly those in \eqref{cdp:eq:C.RK.2} and
\eqref{cdp:eq:C.RK.9}--\eqref{cdp:eq:C.RK.13}.
The primitive strict orders are $N(r_i-\tau)$, where
$\tau=\min_jr_j$. All additional semistable components are
ruled or rational. Consequently every Hopf-birational prime of
$Z$ has minimum canonical order exactly when \eqref{cdp:eq:C.RK.3}
holds. Apply \cref{cdp:st:C.2.UN.75} to this full reduced SNC
torus family. Its equivalence proves both directions of the
stated canonical-flat criterion.

\item \par\medskip\noindent\textbf{The finite birational action and its quotient order.}
Apply \cref{cdp:st:C.2.UN.75} to $Z$. Its comparison
$\phi$ is a bimeromorphic map over
$\Delta_s$, with an inverse
record and a punctured isomorphism. Formula \eqref{cdp:eq:C.RK.4}
therefore defines meromorphic maps
on $X_0$. Their products satisfy
the group laws on the dense
punctured set and consequently
as meromorphic maps. They have
order $N$ on the total family,
as seen on the base. Choosing
the endpoint volume by
$\phi^*\Omega_0=\Omega_Z$ also
gives the character $\zeta^{-a_{\min}}$ from
\cref{prop:II:cyclic-root-normalization} under this birational
action.

\item To reconstruct, first reverse
the $\phi$-record.
The action is then the original
regular $G$-action on $Z$.
Apply its equivariant toric
coarsening $q$ to get the normal space $Y_N$. Apply the invariant-ring comparison
\eqref{cdp:eq:C.RK.14} to this same normal root:
$A=\OO_Y$, and $B$ is the integral closure of
$A[s]/(s^N-f\pi)$. It gives $B^G=A$, including the local
split branches permuted by $G$. Thus $Y_N/G=Y$.
The prescribed ordinary blowdowns $\pi$ recover $X$.
These exact maps
prove \eqref{cdp:eq:C.RK.5} and its converse.
The action on the core is the displayed birational conjugate,
and reversal restores its original regular action on $Z$.

\item The Euler comparison of \cref{prop:II:cyclic-root-normalization}
gives $e(S_Z)-e(S_Y)=R_N$ and $e(S_Y)-e(S)=\Delta_\pi$
for these exact proper maps. By definition of the forward
modification defect, $e(S_0)-e(S_Z)=\Delta_K$.
Adding the three equalities proves \eqref{cdp:eq:C.RK.8}.
The values of $e(S_0)$ are those of \cref{cdp:st:C.2.UN.75}.
This is the same original support as in
\cref{cdp:st:C.2.UN.39}, so the two proper-map formulas compute
its same Euler number.

\item The original-cover comparison of
\cref{prop:II:cyclic-root-normalization} is the map
$w_j=s^{N/n_\alpha}/g_j$ with exactly the transition
$w_k=(g_j/g_k)|_{T_\alpha}w_j$.
It accounts for every strict root component and proves that
all Hopf root components come from these original covers.
The original strict transform keeps
$(k_i,m_i)=(a_\alpha,n_\alpha)$, while
$\tau=\min_i(k_i+1)/m_i$ by the same density calculation.
Thus \eqref{cdp:eq:C.RK.3} is exactly \eqref{cdp:eq:C.RK.20}.
The covers and density are defined on the original family,
which proves the stated independence of the resolution,
semistable subdivision and divisible root exponent.

\item For an original compact family, use the compact cyclic
cover in \cref{prop:II:cyclic-root-normalization}. It has the
same selected local root model and is \'etale over every other
critical value. Perform the normal base change and equivariant
resolution globally, and the \cref{cdp:st:C.2.UN.75} modification
near the selected root fibre. Its punctured isomorphism glues
on an annulus to the unchanged complement. Conjugate the global
finite action through this proper bimeromorphism. Reversing
the record makes the action regular, and the proper coarsening,
finite quotient and ordinary blowdowns in \eqref{cdp:eq:C.RK.5}
recover the original compact family. The converse uses these
same maps on the given compact core and its analytic neighbourhood, including every transverse term
in the input.
\end{enumerate}
\end{proof}

\subsection{Contractions of no-triple fibres}\label{subsec:II:no-triple-contractions}

\begin{proposition}[All no-triple fibres are ordinary vertical blowups of the classified elliptic cores]\hypertarget{tag-01EJ}{}\label{cdp:st:C.2.UN.72}
\StatementRecord{cdp:st:C.2.UN.72}{proved}{01EJ}{\SeriesRef{cdp:st:C.2.UN.65}, \SeriesRef{cdp:st:C.2.UN.57}, \SeriesRef{cdp:st:C.2.UN.61}, \SeriesRef{cdp:st:C.2.UN.62}, \SeriesRef{cdp:st:C.2.UN.69}, \SeriesRef{cdp:st:C.2.UN.70}, \SeriesRef{cdp:eq:C.NT.5}, \SeriesRef{cdp:eq:C.NT.6}, \SeriesRef{cdp:st:C.2.UN.64}, \SeriesRef{cdp:st:C.2.UN.5}, \SeriesRef{night:P4.13}, \SeriesRef{cdp:st:C.2.UN.66}}

Let $f:U\to\Delta$ be a proper degeneration of connected complex
two-tori from a smooth threefold. Suppose that its full central
fibre $S=\sum_vT_v$ is reduced and simple normal crossing and
has no triple points. The components may be nonminimal.

There is a finite sequence of ordinary holomorphic
blowdowns over the disc, each the inverse of a point blowup
or a blowup of a smooth compact curve of positive genus in
the smooth locus of the current central fibre, which ends
in a degeneration with one of these marked cores:
\begin{enumerate}
\item a smooth complex two-torus;
\item a marked diagonal Hopf and degree-zero elliptic ruled
cycle with the exact normal maps and lattice product $M=I$
of \ref{cdp:st:C.2.UN.65} and \ref{cdp:st:C.2.UN.57};
\item a rooted elliptic tree satisfying the full criterion
of \ref{cdp:st:C.2.UN.61}, including its original analytic torsors,
secondary framings and zero-Chern inverse principal torsor;
\item a unicyclic support of genus-zero fixed-Jacobian
identity-monodromy elliptic quasi-bundles and nonlinear
Hopf leaves satisfying the full criterion of \ref{cdp:st:C.2.UN.62}.
\end{enumerate}
The original neighbourhood is recovered by reversing the
recorded ordinary blowups on this core degeneration.
Each centre avoids every double curve then present.
For every curve centre $B$ the record retains the exact
normal extension and marked exceptional pair
\begin{equation}
 \begin{gathered}
 0\longrightarrow K_B\longrightarrow V=N_{B/U'}
                 \xrightarrow{\,df'\,}\OO_B\longrightarrow0,\\
 (T,C)=\bigl(\PP_B(V),\PP_B(K_B)\bigr),\qquad
 N_{T/U}\simeq\OO_{\PP_B(V)}(-1).
 \end{gathered}
\hypertarget{tag-00TU}{}\label{cdp:eq:C.NT.1}
\end{equation}
Here $\PP_B(V)$ parametrizes lines. The endpoint
identifications in \eqref{cdp:eq:C.NT.1} are part of the record.

Conversely, take any torus degeneration with one of
these cores and perform any finite sequence of ordinary
point blowups and ordinary blowups of smooth compact
positive-genus curves in the smooth locus of its current
central fibre, avoiding all current double curves. The
result is a full reduced SNC torus fibre without triple
points. Reversing these maps reconstructs the chosen analytic
neighbourhood, including the endpoint identifications in
\eqref{cdp:eq:C.NT.1}.
For a compact total space over a compact curve all the
recorded contractions and blowups have the same compact
version.

If $a$ is the number of point blowups and $B_1,\ldots,B_b$
are the recorded curve centres, then
\begin{equation}
                 e(S)=2a-2\sum_{j=1}^b\bigl(g(B_j)-1\bigr).
\hypertarget{tag-00TV}{}\label{cdp:eq:C.NT.2}
\end{equation}
Thus the complete zero-Euler reconstruction is obtained
exactly by the balance
\begin{equation}
                       a=\sum_{j=1}^b\bigl(g(B_j)-1\bigr).
\hypertarget{tag-00TW}{}\label{cdp:eq:C.NT.3}
\end{equation}
Equivalently, let $r_0$ count the connected rational
double curves of the original support and let
$\mathcal E_{\ge2}$ be its connected double curves of
genus at least two. These intrinsic data satisfy
\begin{equation}
 \begin{split}
 a&=r_0,\qquad
 \{B_j:g(B_j)\ge2\}
       \ \simeq\ \{C_e:e\in\mathcal E_{\ge2}\},\\
 e(S)&=2r_0-2\sum_{e\in\mathcal E_{\ge2}}
                                      \bigl(g(C_e)-1\bigr).
 \end{split}
\hypertarget{tag-00TX}{}\label{cdp:eq:C.NT.4}
\end{equation}
The bijection is the one induced by the contractions
on the double curves. Every nonelliptic edge is a terminal
edge at the stage at which it is removed; in particular
every cycle of the original incidence graph consists of
elliptic double curves.
\end{proposition}

\begin{proof}

\begin{enumerate}[label=\arabic*.,leftmargin=*]
\item \par\medskip\noindent\textbf{First remove every rational double curve by point contractions.}
Apply \ref{cdp:st:C.2.UN.69} to the given degeneration. Its strict-height
ascent ends in plane-line leaves, and each ordinary
blowdown removes one rational double curve and no other
double curve. A rational ruled-section leaf was excluded
there by transporting its square-minus-two centre into
the nearby torus. Thus this first stage makes
exactly $r_0$ point contractions. It leaves a proper
smooth-total torus degeneration, still full reduced
SNC without triple points, and with no rational double
curve. The maps are isomorphisms on neighbourhoods of
every surviving double curve.

\par\medskip\noindent\textbf{An increasing-height path cannot change from higher genus to elliptic.}
Use the canonical orders of the current family.
For every remaining double curve the exact adjunction
map is \eqref{cdp:eq:C.NM.2}. In particular a curve of genus at least
two has endpoints of different heights. Orient that
edge toward its higher endpoint and call that component
$P$.

\item Suppose $P$ has a neighbour $Q$ of strictly greater
height, and write $J=P\cap Q$ for the corresponding
connected boundary. If $J$ were elliptic, \eqref{cdp:eq:C.NM.2}
would give the isomorphism
\begin{equation}
      N_{J/P}^{\,k_Q-k_P}\simeq K_J\simeq\OO_J,
        \qquad a_J=k_Q-k_P-1\ge0 .
\hypertarget{tag-00TY}{}\label{cdp:eq:C.NT.5}
\end{equation}
Its normal is therefore torsion, with the canonical coefficient required by \ref{cdp:st:C.2.UN.70}. All boundary
curves on $P$ are disjoint. That theorem supplies
the map
\begin{equation}
        \pi:P\longrightarrow Z,\qquad
        q_Z:Z\longrightarrow B,
\hypertarget{tag-00TZ}{}\label{cdp:eq:C.NT.6}
\end{equation}
which is a sequence of point blowdowns followed by a
relatively minimal elliptic quasi-bundle with only
smooth multiple supports; every original boundary
is vertical for $q_Z\pi$.

\item No smooth boundary of genus at least two can be
vertical for this map. An exceptional prime of the
point-blowdown map is rational. A nonexceptional
boundary has the same function field as its image
under this birational morphism. A vertical image is
one of the smooth elliptic fibre supports of $q_Z$,
so its smooth strict transform has genus one.
Both alternatives contradict the higher-genus
boundary already on $P$. Hence $J$ is not elliptic.
There are no rational double curves at this stage,
so every strictly higher neighbour is reached by
another higher-genus double curve.

\par\medskip\noindent\textbf{Every higher-genus edge leads to a section-leaf contraction.}
Continue along strictly increasing heights whenever
possible. There are finitely many components, so the
path terminates at a local maximum $T$ having a
strictly lower neighbour along a higher-genus curve.
The maximum theorem \ref{cdp:st:C.2.UN.69} applies to
this possibly nonminimal component. Its equal-height alternative is
impossible because of that lower neighbour. Its Hopf
leaf alternative has an elliptic boundary, and its
plane-line alternative has a rational boundary.
Thus $T$ is exactly an unblown projective-line
bundle whose single boundary is a section of genus
at least two. The \ref{cdp:st:C.2.UN.64} map contracts it
to the smooth curve identified with that section.

\item The parent surface is unchanged by this curve
contraction, since its restriction is the blowup
of a Cartier curve. Every other component misses
the exceptional divisor. Consequently the new
central fibre is still full reduced SNC without
triple points; no rational double curve has appeared,
and every surviving boundary neighbourhood is
unchanged. Canonical order differences on surviving
components are retained by the ordinary blowup
Jacobian, as in \ref{cdp:st:C.2.UN.64}.

\item Repeat this finite procedure. Each step removes
one component and precisely one higher-genus double
curve. It therefore removes all such curves and
gives the exact bijection in the first line of
\eqref{cdp:eq:C.NT.4}. Each curve centre and its normal extension
are exactly \eqref{cdp:eq:C.NT.1}. At this point every remaining
double curve is elliptic.

\par\medskip\noindent\textbf{The remaining support has zero Euler number and is already minimal.}
Denote this remaining full support by $S_{\rm ell}$.
If it is smooth, the smooth-support comparison \ref{cdp:st:C.2.UN.5}
identifies it with a two-torus. Suppose it is
reducible. The nearby-cycle Euler formula
for a reduced SNC fibre is
\begin{equation}
        e(F)=\sum_v e(T_v^\circ),\qquad
        T_v^\circ=T_v\setminus
                          \bigcup_{w\ne v}(T_v\cap T_w).
\hypertarget{tag-00U0}{}\label{cdp:eq:C.NT.7}
\end{equation}
Apply \cref{night:P4.13} to this same family with the
identity resolution and $N_j=1$ for every component.
The least common multiple of the component multiplicities is
$1$. Its nearby-cycle filtration therefore makes the monodromy
$T$ on $H^1(F)$ unipotent, so $\det(I-T)=0=e(F)$.
Equation~\eqref{cdp:eq:C.MS.2} at $k=1$ is therefore
\eqref{cdp:eq:C.NT.7}, with $E_j^\circ=T_v^\circ$.

\item Writing $C_e$ for the double curves, all of which
are elliptic, gives
\begin{equation}
 \begin{split}
 0=e(F)&=\sum_v e(T_v)-2\sum_e e(C_e)
                            =\sum_v e(T_v),\\
 e(S_{\rm ell})&=\sum_v e(T_v)-\sum_e e(C_e)=0.
 \end{split}
\hypertarget{tag-00U1}{}\label{cdp:eq:C.NT.8}
\end{equation}
Now use the all-elliptic extension \eqref{cdp:eq:C.MR.16}, whose
proof applies to these possibly nonminimal
components. Its nonnegative minimal-model Euler
terms and positive point-blowup terms force every
component of $S_{\rm ell}$ to be minimal and
nonrational. Thus the same support satisfies the component and
no-triple hypotheses of \cref{cdp:st:C.2.UN.66}.

\ref{cdp:st:C.2.UN.66} now makes its further ordinary contractions.
All their curve centres are elliptic: the boundary
curves present are elliptic and no rational or
higher-genus edge is introduced. The remaining
support is one of the four cores in the
statement. Its full torsor, lattice, secondary and
normal identifications are those of \ref{cdp:st:C.2.UN.61}, \ref{cdp:st:C.2.UN.62}
or \ref{cdp:st:C.2.UN.65}, as appropriate. 

\par\medskip\noindent\textbf{Exact reversal and the converse construction.}
Every preceding ambient map is one of the ordinary
holomorphic blowdowns of \ref{cdp:st:C.2.UN.64}. Reverse the finite
record on the endpoint family. A point
step replaces a smooth point of one central
component by its exceptional plane. A curve
step has the exact normal sequence \eqref{cdp:eq:C.NT.1}.
These comparisons are the local blowup
charts and glue with the original isomorphism
off each exceptional divisor. They recover the given $U\to\Delta$ with all of
its neighbourhood identifications.

\item For the converse, use the inverse ordinary maps of
\cref{cdp:st:C.2.UN.64} with the chosen point or smooth curve
centre in the current central fibre. The curve map is the
blowup of its ideal $(z,t)$ with quotient normal coordinate
$df'=dt$; the point map is the blowup of $(x,y,t)$.
The principal-fibre pullbacks in these maps have the reduced
equation $t=zs$ or $t=xv$ on their two-component charts.
Each centre avoids the other boundary curves, so these are
the only components through the new double curve. The
maps preserve the smooth total space, full reduced SNC
fibre, absence of triple points and punctured torus family.
Properness follows from the proper ordinary blowup map.

\item On a curve centre $B$ in the smooth central
locus the canonical divisor is supported
on boundaries away from $B$. Its restriction
trivializes $K_{P'}|_B$. Adjunction therefore gives
$N_{B/P'}\simeq K_B$, and the normal sequence
with the quotient $df'$ is precisely \eqref{cdp:eq:C.NT.1}.
No curve of genus zero can be such a centre:
the compact submersion transport \eqref{cdp:eq:C.NM.10}
would again put its square-minus-two sphere
into the nearby torus. Points and smooth curves
of positive genus are exactly the centre
types used in the reconstruction.

\item When the original total space is compact, the
local blowdowns glue to the identity
off their exceptional divisors by \ref{cdp:st:C.2.UN.64}.
Their targets remain compact and the map to
the compact base descends. Conversely every
listed compact centre has its ordinary global
blowup. These are the stated maps on the same compact total
space and base.

\par\medskip\noindent\textbf{Euler balance and the intrinsic double-curve record.}
All four cores have Euler number zero. Apply the exact
Euler comparison \eqref{cdp:eq:C.EL.5} to this sequence of
ordinary contractions, with its $a$ equal to the number
of point centres and its $B_j$ equal to these same curve
centres. Reversing the maps gives \eqref{cdp:eq:C.NT.2};
setting it equal to zero gives \eqref{cdp:eq:C.NT.3}.

\item Every such blowup attaches a single new double
edge: a rational line for a point, or the curve $B$ for a curve centre. It leaves all
pre-existing double curves and their
neighbourhoods unchanged. The cores have only
elliptic double curves. The rational-edge
removal stage made exactly $r_0$ point steps,
and the higher-genus removal stage identifies
each of its centres with exactly the removed
original higher-genus double curve. These
maps prove all of \eqref{cdp:eq:C.NT.4} and its
independence of the chosen contraction order.
They also show that every nonelliptic edge
is removed as a leaf, which proves the
assertion about cycles. 

\end{enumerate}
\end{proof}

\subsection{Finite reduction of reduced normal-crossing fibres}\label{subsec:II:reduced-finite-reduction}

\begin{definition}[Marked collections of surface models]\hypertarget{tag-P368}{}\label{def:II:marked-surface-models}
\StatementRecord{def:II:marked-surface-models}{proved}{P368}{\cref{cdp:st:C.2.UN.83}}
In a smooth-total torus degeneration with full reduced
normal-crossing central fibre, a \emph{marked surface model}
consists of a component normalization $T$, its full reduced
boundary $D_T$, an isolated smooth elliptic boundary curve $J_T$
with torsion normal, and one of the three comparisons of
\cref{cdp:st:C.2.UN.83}, with point map
\begin{equation}
                         \pi_T:T\longrightarrow Z_T.
 \hypertarget{tag-00L2}{}\label{cdp:eq:C.GP.4}
\end{equation}
The comparison includes its elliptic, nonlinear Hopf or
one-section ruled target structure, its isomorphism near
$J_T$, and all its marked boundary branches and normal maps.
Both branches of every self-double curve belong to these data.
Its \emph{point defect} $d_T$ is the number of point blowups
inverse to $\pi_T$.

A \emph{marked collection} $\mathcal A$ is a set of global
component primes with one such chosen model on the normalization
$T_V$ of each $V\in\mathcal A$. A component is \emph{modelled}
when it belongs to this chosen set. Thus changing $\mathcal A$
changes the chosen comparisons on the same components.
\end{definition}

\begin{proposition}[Every reduced SNC torus degeneration has a reversible finite reduction to a classified core]\hypertarget{tag-01EX}{}\label{cdp:st:C.2.UN.85}
\StatementRecord{cdp:st:C.2.UN.85}{proved}{01EX}{\SeriesRef{cdp:st:C.2.UN.65}, \SeriesRef{cdp:st:C.2.UN.68}, \SeriesRef{cdp:st:C.2.UN.71}, \SeriesRef{cdp:st:C.2.UN.61}, \SeriesRef{cdp:st:C.2.UN.62}, \SeriesRef{cdp:st:C.2.UN.75}, \SeriesRef{cdp:st:C.2.UN.84}, \SeriesRef{cdp:st:C.2.UN.57}, \SeriesRef{cdp:st:C.2.UN.73}, \SeriesRef{cdp:st:C.2.UN.82}, \SeriesRef{cdp:st:C.2.UN.77}, \SeriesRef{cdp:st:C.2.UN.81}, \SeriesRef{cdp:st:C.2.UN.83}, \SeriesRef{cdp:eq:C.GP.5}, \SeriesRef{cdp:st:C.2.UN.80}, \SeriesRef{cdp:eq:C.GP.4}, \SeriesRef{cdp:eq:C.GP.8}, \SeriesRef{cdp:st:C.2.UN.74}, \SeriesRef{cdp:st:C.2.UN.50}, \SeriesRef{cdp:st:C.2.UN.28}, \SeriesRef{cdp:st:C.2.UN.47}}

Let \(f:X\to\Delta\) be a proper degeneration of connected
complex two-tori, with smooth total space and full reduced
SNC central fibre. There is a finite overlap-compatible
modification record over the same disc, isomorphic over
\(\Delta^*\), ending in one of the following smooth-total reduced degeneration germs.

\begin{enumerate}
\item A canonical-flat core. Its central fibre is a smooth
two-torus, one of the exact no-triple elliptic-ruled/diagonal
Hopf cycles of \ref{cdp:st:C.2.UN.65}, or one of the marked positive toric
log spaces of \ref{cdp:st:C.2.UN.68} and \ref{cdp:st:C.2.UN.71}. The toric alternative has
Euler number equal to its positive triple count.

\item An embedded no-triple core whose components are
minimal fixed-Jacobian elliptic quasi-bundles with identity
integral monodromy, together with any minimal nonlinear
Hopf leaves. All boundary curves are smooth elliptic fibre
supports or the specified Hopf axes. Every elliptic
quasi-bundle base has genus zero or one, and at most one
has genus one. The exact support, including its original
torsors, normal characters, translations and gluing maps,
is one of the rooted-tree or unicyclic supports classified
in \ref{cdp:st:C.2.UN.61} and \ref{cdp:st:C.2.UN.62}. Its Euler number is zero.
\end{enumerate}

If every Hopf-birational prime has minimum canonical order,
the first alternative is supplied by \ref{cdp:st:C.2.UN.75}. If some such
prime has positive order, the second alternative can be
obtained while retaining a fixed maximal-order minimal
Hopf component and its entire neighbourhood after the
preliminary reduction of \ref{cdp:st:C.2.UN.84}. The maps preserve the nonlinear Hopf neighbourhoods
throughout this reduction.

The record retains all centres, two-ruling maps,
point and curve contractions, local sigma charts,
normalization branches and restored overlap maps.
Its inverse reconstructs the original analytic family
from this endpoint germ. On an originally compact
family all operations take place on that compact base.
If \(S_0\) is the original reduced support, \(S_*\) the
endpoint support, and \(\delta_j\) the signed normal-cone
Euler change of its \(j\)-th retained step as in \ref{cdp:st:C.2.UN.75},
then
\begin{equation}
                 e(S_0)=e(S_*)-\sum_j\delta_j.
 \hypertarget{tag-00KY}{}\label{cdp:eq:C.GP.1}
\end{equation}

The classified core central spaces have the compact
occurrence constructions of \ref{cdp:st:C.2.UN.57}, \ref{cdp:st:C.2.UN.61}, \ref{cdp:st:C.2.UN.62} and \ref{cdp:st:C.2.UN.73}.
The inverse record is defined on the same endpoint
analytic neighbourhood.
\end{proposition}

\begin{proof}\leavevmode\par

 
\begin{enumerate}[label=\arabic*.,leftmargin=0pt,labelsep=.5em]
\item \par\medskip\noindent\textbf{The initial dichotomy and the protected component.}
Normalize the primitive canonical evaluation as
\begin{equation}
       \operatorname{div}(\Omega)=\sum_i k_iV_i,
       \qquad \min_i k_i=0.
 \hypertarget{tag-00L0}{}\label{cdp:eq:C.GP.2}
\end{equation}
If every Hopf-birational prime has order zero, \ref{cdp:st:C.2.UN.75}
gives its canonical-flat endpoint. \ref{cdp:st:C.2.UN.65} and
\cref{cdp:st:C.2.UN.68,cdp:st:C.2.UN.69,cdp:st:C.2.UN.70,cdp:st:C.2.UN.71} give exactly the alternatives in (1),
including the marked log maps in the positive
case. This uses the proved component restriction that
every class-VII minimal model is Hopf. 

\begin{enumerate}[label=(\arabic*),leftmargin=1.3em,labelsep=.5em]
\item  Suppose instead that a Hopf-birational prime has
positive order. Apply the finite truncated record
\ref{cdp:st:C.2.UN.84}, including its maximal Hopf minimalization. At
its smooth-total NC endpoint put
\begin{equation}
                  M=\max_i k_i>0.
 \hypertarget{tag-00L1}{}\label{cdp:eq:C.GP.3}
\end{equation}
The minimum is still zero and all orders are at
most this fixed \(M\). Choose one of the surviving
minimal Hopf components \(H\) of order \(M\).
\ref{cdp:st:C.2.UN.82} shows that it has exactly one elliptic
boundary curve \(J\), with torsion normal, and
that its neighbour has smaller order. The marked
neighbourhood is embedded. No elliptic self-double
curve has been introduced in the preliminary
record: its new incidences are on rational
exceptional curves, as proved in \ref{cdp:st:C.2.UN.84}.

\item  All later whole-component removal blocks will be
the positive-order blocks of \ref{cdp:st:C.2.UN.77} and \ref{cdp:st:C.2.UN.81}. Such
a ruled component of order \(q\le M\) has one
horizontal boundary section, whose neighbour has
order \(q-1\), and every other boundary is rational
and vertical. It cannot meet \(H\): a rational
curve cannot lie on a minimal Hopf, and the
section case would require \(q=M+1\). A plane
contraction also has only rational boundary.
The same argument applies to the cleanup blocks
and their missing-fibre insertions. Direct moves
below will likewise avoid \(H\). Thus its entire
neighbourhood is protected and provides the same
starting interface after every completed block.

\end{enumerate}

\item \par\medskip\noindent\textbf{The chosen surface models.}
At each smooth-total NC stage choose a marked collection
$\mathcal A$ as in \cref{def:II:marked-surface-models}.
Its three types and point maps \eqref{cdp:eq:C.GP.4} are those
of \cref{cdp:st:C.2.UN.83}, including the one-section ruled
model's rank-two bundle extension.

\begin{enumerate}[label=(\arabic*),leftmargin=1.3em,labelsep=.5em]
\item The chosen comparison and its point defect have
the following properties:
\begin{equation}
\begin{gathered}
 d_T\ge0,\qquad
 \operatorname{Exc}(\pi_T)\subset D_T,\\
 \text{every rational boundary is exceptional
 and has negative square},\\
 \text{every elliptic boundary has nonpositive square}.
\end{gathered}
 \hypertarget{tag-00L3}{}\label{cdp:eq:C.GP.5}
\end{equation}
For the elliptic type, the boundary is vertical
for the composed elliptic map; \eqref{cdp:eq:C.GP.5} follows
from \ref{cdp:st:C.2.UN.80} and the fibre squares. For the other
two types the entrance curve has square zero,
all point centres are away from it, and every
other boundary is exactly a positive discrepancy
exceptional curve, by \eqref{cdp:eq:C.TI.2}--\eqref{cdp:eq:C.TI.3}.

\item Initially take $\mathcal A=\{H\}$. \ref{cdp:st:C.2.UN.83} applied
to its isolated torsion-normal axis gives
the elliptic or nonlinear Hopf type, and its
point defect is zero. Let
\begin{equation}
\begin{gathered}
 u=\#\{V_i:V_i\notin\mathcal A\},
 \\
 d=\sum_{V\in\mathcal A}d_{T_V}.
\end{gathered}
 \hypertarget{tag-00L4}{}\label{cdp:eq:C.GP.6}
\end{equation}
Adding a component to $\mathcal A$ means choosing the
maps of \cref{cdp:st:C.2.UN.83} on that same component and
marking its entrance curve.

\end{enumerate}

\item \par\medskip\noindent\textbf{A direct move on a chosen elliptic model.}
Suppose that a modelled elliptic-type component
has positive defect. Choose a last exceptional
curve of \eqref{cdp:eq:C.GP.4}. It is a boundary by
\ref{cdp:st:C.2.UN.80}. Apply the complete-fibre and receiving-degree comparisons
\eqref{cdp:eq:C.AP.5}--\eqref{cdp:eq:C.AP.6} to the inverse fibre of $\pi_T$ containing
this last curve. Their recurrences use its point centres and
normal-product map, so the same equations apply to this
point map for the entrance coefficient supplied by
\cref{cdp:st:C.2.UN.83}.

\begin{enumerate}[label=(\arabic*),leftmargin=1.3em,labelsep=.5em]
\item  Here is the exact comparison with an embedded
curve in the total space at an NC stage. On a
complete fibre graph of the point map, write
\(a_C\) for the coefficient of each component
in the canonical divisor on the
normalization, assigning coefficient zero
also to a component absent from the boundary. Adjacent shifted coefficients
\(a_C+1\) are relatively prime. There is no
initial edge on the smooth elliptic fibre;
a blowup at a smooth point creates consecutive
shifted coefficients, and a node blowup inserts
the sum of the two old shifted coefficients.
Thus the assertion follows through the point map, including negative initial coefficients.
In particular no edge has coefficients
\((-1,-1)\).

\item  A self-double branch has coefficient \(-1\).
If the image of the chosen last \((-1)\)-curve
were nodal at a triple point with two
local branches of its source component, its
two neighbours on the normalization would be
the two self-double branches. The last-node
coefficient formula would also give
\(a_E=(-1)+(-1)+1=-1\), contradicting that
edge invariant. A triple with three branches
of the source gives adjacent \(-1\) branches
directly. With only one local source branch,
two image branches would already make a node
on the chosen curve in that smooth source
normalization, contrary to its being the
smooth last exceptional \(\PP^1\).
The image curve is consequently
smooth. For a self-double its normalized
degree-two cover is unramified, by the local
NC coordinates. A connected such cover by
\(\PP^1\) is impossible by Riemann--Hurwitz;
the cover therefore splits into two distinct
normalized branches. Over the resulting
simply connected image \(\PP^1\), those
labels give two smooth surface germs.
When the triple count is two, the normal-product
map gives the two normal degrees \(-1,-1\).
The type-II neighbourhood and its two
ruling maps therefore apply even if these
two surface germs belong to one global prime.
For the positive-discrepancy models below,
all rational coefficients are positive, so
the same self-double obstruction is immediate.

\item  With two triple points, the type-II move
blows down this curve on its two adjacent
normalized surfaces. Each opposite receiving
boundary has square at least zero, or at least
one when it is elliptic. By \eqref{cdp:eq:C.GP.5} neither
receiver was modelled before this move. Thus
the move adds no point defect to a modelled
surface.

\item  If its second source surface was modelled,
the contracted rational boundary was exceptional
for its chosen point map, again by \eqref{cdp:eq:C.GP.5}.
That map factors through the blowdown:
it is constant on the contracted curve, and
the proper surface contraction and its direct
image of the structure sheaf give this descent.
The resulting birational surface morphism is
the remaining sequence of point blowups. Its
defect decreases by one. The same factorization
applies to the first source. All marked
entrance curves are disjoint from these
exceptional curves and are unchanged. Hence
the marked collection and its models persist,
\(u\) is unchanged, and \(d\) strictly decreases.
If both source germs belong to one global
modelled component, their two normalized
curves are disjoint. The same point map
factors through both ordinary contractions;
its defect drops by two while the global
component is counted only once.

\item  If a last curve has one triple point, use the
positive ruled-removal block of \eqref{cdp:eq:C.AP.7}.
If it has none, use the exact plane progress
block: apply \cref{lem:II:positive-rational-curve} with $T$ the receiving normalization and $C$ the image of this last curve, whose square is $1$, and then use the finite progress in \eqref{cdp:eq:C.AP.10}--\eqref{cdp:eq:C.AP.13}. Each gives a strict component-count decrease after its
finite cleanup. Every last exceptional curve belongs to
the full boundary by \cref{cdp:st:C.2.UN.80}, so these
three alternatives are exhaustive.

\end{enumerate}

\item \par\medskip\noindent\textbf{A direct move on either positive-discrepancy model.}
Now let the chosen model be nonlinear Hopf
or one-section ruled, with positive defect.
The canonical divisor is
\begin{equation}
 K_T=-aJ_T+\sum_E b_EE,
 \qquad b_E>0,
 \qquad a\ge2,
 \hypertarget{tag-00L5}{}\label{cdp:eq:C.GP.7}
\end{equation}
with \(a=2\) for the ruled type. Every point
centre is away from the preserved \(J_T\),
and the full other boundary is exactly the
exceptional set. The graph of each exceptional
tree starts with a rational curve of square
\(-1\) and valence zero. If \(C^2=-b\) and
\(n\) is its valence in this entire exceptional
set, the smooth and node blowups give
\begin{equation}
                              n\le b+1.
 \hypertarget{tag-00L6}{}\label{cdp:eq:C.GP.8}
\end{equation}
At a centre on the exceptional support choose smooth
coordinates identifying that support with $x=0$ or $xy=0$.
The blowup of $(x,y)$ is the same local point morphism whose
square and valence updates prove \eqref{cdp:eq:C.AP.5}.
A centre off the exceptional support starts a rational
tree at $(b,n)=(1,0)$, which also satisfies the bound.
Thus the recurrence in \eqref{cdp:eq:C.AP.5}, applied at these
identified point maps, gives \eqref{cdp:eq:C.GP.8}.
At the last node the opposite
square is \(b+1-t\ge0\), since the triple count \(t\) is at most the old valence.
All opposite curves here are rational.
Thus no receiver is already modelled by
\eqref{cdp:eq:C.GP.5}, and the exact point-map factorization
in the preceding paragraph again strictly
decreases \(d\). The positive discrepancies
also give the precise height comparison:
if the two old coefficients are \(b_1,b_2>0\)
and the current component has order \(p\),
the receiving orders are
\begin{equation}
                 p+b_1+1,\qquad p+b_2+1.
 \hypertarget{tag-00L7}{}\label{cdp:eq:C.GP.9}
\end{equation}

\begin{enumerate}[label=(\arabic*),leftmargin=1.3em,labelsep=.5em]
\item  A last curve with no triple point has a
rational \(+1\) receiver of order \(p+2\),
and the plane-map clause of \cref{lem:II:positive-rational-curve}, applied to that receiving pair, followed by \eqref{cdp:eq:C.AP.10}--\eqref{cdp:eq:C.AP.13}, supplies the finite plane progress block. With one triple point
its receiver is ruled with one section and
has positive order, so the same ruled block
applies. Their centres are in these rational exceptional
configurations, disjoint from the preserved nonlinear
Hopf neighbourhood.

\item  When a one-section ruled model reaches defect
zero, it is the bundle \(\PP(V)\)
with sole boundary section \(J_T\), including
its possibly nonzero extension class. \ref{cdp:st:C.2.UN.74}
gives the ordinary curve contraction along
its ruling. Its order is
\(k_T=k_{\mathrm{neighbour}}+1>0\), by \ref{cdp:st:C.2.UN.83}.
This is another component-count decrease.

\end{enumerate}

\item \par\medskip\noindent\textbf{Extending the marked collection across an elliptic interface.}
We always eliminate positive point defects
before adding another component to $\mathcal A$. We also
remove a defect-zero one-section ruled model
by the preceding contraction. If neither
action is available, every modelled component
is a minimal elliptic quasi-bundle or a minimal
nonlinear Hopf leaf.

\begin{enumerate}[label=(\arabic*),leftmargin=1.3em,labelsep=.5em]
\item  Such a component is embedded and has disjoint
smooth elliptic boundary curves. Indeed no
rational boundary can remain when its point
defect is zero. The original fibre was SNC,
and the retained type-II and one-section
sigma records create new self-incidences only
on rational fibre or exceptional curves.
Their generic one-section contraction cannot
identify two elliptic horizontal branches.
Thus there is no new elliptic self-double
curve on one of these minimal components.

\item  If an unmodelled component remains, the
connected incidence graph of the central
support has an edge from the modelled
region to its complement. Starting with
\(H\) and choosing each new entrance edge
makes that region connected; the direct
moves have retained these entrance curves.
The edge is therefore a smooth elliptic
curve, isolated on its modelled side.
There is no triple point on it, so its
other normalization branch is isolated
as well. Its normal is torsion: on a
quasi-bundle it is the exact normal root
of a smooth multiple fibre, and on a
nonlinear Hopf it is its axis character.
The reduced-fibre normal-product
map gives the inverse torsion line on
the other side.

\item  Apply \ref{cdp:st:C.2.UN.83} to that other normalization.
It supplies one of the three point models
\eqref{cdp:eq:C.GP.4}, with the entrance neighbourhood
unchanged and every original boundary
retained. Add that component and comparison to $\mathcal A$. Then
\(u\) decreases by one. The added point defect is finite and nonnegative;
its value is absorbed by the strict decrease of $u$
in the lexicographic potential below.

\end{enumerate}

\item \par\medskip\noindent\textbf{Resetting the marked collection after a removal.}
After each whole-component removal and its finite sigma
cleanup, set $\mathcal A=\{H\}$ with the original marked
point map and torsion-normal interface on $H$.
The full neighbourhood of $H$ is preserved by the preceding
centre and order argument. This reset therefore gives the
same valid initial model, with finite new values of $u,d$,
at the smooth-total endpoint of the block.

\begin{enumerate}[label=(\arabic*),leftmargin=1.3em,labelsep=.5em]
\item  Let \(c_j\), for \(1\le j\le M\), count
the retained global components of order \(j\).
The full fixed potential is
\begin{equation}
                 (c_1,c_2,\ldots,c_M,u,d).
 \hypertarget{tag-00KZ}{}\label{cdp:eq:C.GP.10}
\end{equation}
Every direct type-II move preserves the component-count
tuple and \(u\), and strictly decreases \(d\).
Adding a model preserves the component counts and strictly
decreases \(u\), regardless of its finite
new defect. Every whole-component block
strictly decreases its first changed component count,
regardless of the finite reset of \(u,d\).
The missing-fibre and cleanup rules
add only finitely many components at the next
order and add none at \(M\). Thus all
retained orders stay at most the fixed \(M\).

\item  Lexicographic order on this fixed finite
product of nonnegative integers is well
founded. The internal plane progress and
sigma cleanup blocks are finite by \ref{cdp:st:C.2.UN.81}
and \ref{cdp:st:C.2.UN.77}. Consequently this procedure
cannot continue indefinitely.

\item  At its endpoint no point defect is
positive, no one-section ruled model
remains, and no component is unmodelled.
Otherwise the preceding paragraphs give
another strict decrease. All components
are therefore the minimal elliptic
quasi-bundles and nonlinear Hopf leaves
in (2), and every boundary is smooth
elliptic and isolated. In particular the
endpoint is embedded SNC and has no
triple point.

\end{enumerate}

\item \par\medskip\noindent\textbf{The classified endpoint and inverse maps.}
Every endpoint component and double curve
has topological Euler number zero, so
\(e(S_*)=0\). The exact invariant relative
cohomology frame gives
\(p_g(Z)=g(B)\) for a quasi-bundle. The
common regular Jacobian is retained by
the normal-root maps: along an
edge joining two quasi-bundle vertices,
the dual normal lines have the same
exact order, and inversion identifies
their connected root covers of the
elliptic double curve. These are exactly
the two regular elliptic fibres, by
\eqref{cdp:eq:C.QT.2}--\eqref{cdp:eq:C.QT.3}; their translations
are retained in the gluing record.
At a nonlinear Hopf leaf the same root
cover is instead its exact period cover
of \ref{cdp:st:C.2.UN.50}. Such leaves do not interrupt
propagation through the connected
quasi-bundle region, and their leaf periods and finite characters are
retained by that comparison.
The
compact-core bound \ref{cdp:st:C.2.UN.28} therefore permits
at most one genus-one base. If there is
one, that normalization has \(p_g=1\).
\ref{cdp:st:C.2.UN.47} gives the surjection from its
fundamental group onto that of the full
support. The graph-of-spaces comparison
gives a surjection from the support
fundamental group to that of its incidence
graph, and kills the image of this
component's fundamental group.
The graph fundamental group is therefore
trivial, so the connected graph is a tree.
The exact coherent restriction and
fundamental-group comparisons of \ref{cdp:st:C.2.UN.61}
then give its rooted-tree support, with the
integral Chern and recursive
period conditions. If there is none,
\ref{cdp:st:C.2.UN.62} gives the unicyclic support and
its exact torsor, height and common-period
conditions. Its nonlinear leaves retain
the secondary finite characters of \ref{cdp:st:C.2.UN.50}.
These applications use the endpoint's original
surfaces and gluing maps in the same smoothing.

\begin{enumerate}[label=(\arabic*),leftmargin=1.3em,labelsep=.5em]
\item  All prior operations are the finite diagrams already used in \ref{cdp:st:C.2.UN.75},
\ref{cdp:st:C.2.UN.77}, \ref{cdp:st:C.2.UN.81} and \ref{cdp:st:C.2.UN.84}. Their proper graphs
glue by the retained overlap maps, so
their inverse recovers this given
analytic family. Equation \eqref{cdp:eq:C.GP.1}
is the telescoping normal-cone formula
on that finite record, including its
restored-overlap subdivision. Annular
gluing to an unchanged compact outside
gives the same modifications
when the starting total space is compact.

\item The compact occurrence maps of
\cref{cdp:st:C.2.UN.57,cdp:st:C.2.UN.61,cdp:st:C.2.UN.62,cdp:st:C.2.UN.73}
realize the classified core central spaces. The inverse
maps just constructed act on this reduction's endpoint
neighbourhood and recover its given torus degeneration.

\end{enumerate}

\end{enumerate}
\end{proof}

\subsection{Finite rooted generation}\label{subsec:II:rooted-generation}

\begin{definition}[Finite rooted reconstruction data]\hypertarget{tag-P341}{}\label{def:II:finite-rooted-reconstruction}
\StatementRecord{def:II:finite-rooted-reconstruction}{proved}{P341}{\cref{def:II:overlap-compatible-modification,prop:II:cyclic-root-normalization,cdp:st:C.2.UN.85}}
Finite rooted reconstruction data comprise the following analytic
neighbourhoods, arrows and transition maps, including the full
transverse terms used by each inverse operation.
\begin{enumerate}
\item A proper smooth-total core family
\(f_*:X_*\to\Delta_s\), with its labelled analytic
charts and overlaps and its punctured torus-family
comparison, is given. Its full reduced central space
is one of the cores in \cref{cdp:st:C.2.UN.85}: a smooth torus; a
canonical-flat marked Hopf/degree-zero-ruled cycle;
a positive periodic toric log space; or an embedded
rooted-tree or unicyclic elliptic quasi-bundle/Hopf
support of \cref{cdp:st:C.2.UN.61,cdp:st:C.2.UN.62}. The core chart functions,
including their transverse terms, belong to the datum.

\item A finite inverse overlap-compatible sigma record
on that core is given. Each arrow includes
the local blowup or contraction map, its centre
and ideal, the target charts, and the isomorphism
on the restored overlap. The maps and their inverse
comparisons are checked on these charts as in \cref{cdp:st:C.2.UN.75}.
The resulting space \(Z\) is smooth and proper over
\(\Delta_s\), with full reduced SNC central divisor.
The record is an isomorphism on the punctured family.

\item A finite cyclic birational action on the core,
over \(s\mapsto\zeta s\), is given, where \(\zeta\)
has order \(N\). Conjugating it by the given
inverse record gives holomorphic maps on every
chart of \(Z\), including the central divisor.
Their compositions satisfy the cyclic group law.
These maps define a regular action of $G=\mu_N$ on $Z$.

\item A labelled equivariant toric coarsening
\(q:Z\to Y_N\) is given. Its local semigroup
charts, units and maps are explicit normalized
monomial-ideal blowups in the reverse direction.
The finite quotient charts of \(Y_N\) have
invariant ring \(\OO_Y\), with smooth coordinates
in which the descended function is
\begin{equation}
                  t=u\prod_{i=1}^r x_i^{m_i},
              \qquad r\le3,\quad u\in\OO_Y^*,
              \quad m_i\mid N.
 \hypertarget{tag-00A0}{}\label{cdp:eq:C.AR.1}
\end{equation}
The normal charts upstairs are exactly the
normalization of \(s^N=t\) in these coordinates.
The labelled height-one subdivisions defining
\(q\) are unimodular, and their overlap maps
and finite action belong to the datum. These ring and
coordinate conditions verify the smooth quotient
\(Y=Y_N/G\) and its full SNC divisor with
the displayed multiplicities.

\item A finite sequence \(\pi:Y\to X\) of
ordinary inverse point or smooth-curve blowups
is given, with their smooth target charts
and gluing maps. Their centres are in the central
fibre. The function \(t\) descends through each
map, as is checked by its displayed pullback.
The final total space \(X\) is smooth. Its
central support is allowed arbitrary singularities
and arbitrary full multiplicities.
\end{enumerate}
The quotient, contraction and overlap identities are
conditions on these finite maps and their coordinate algebras.
An ordinary arrow can equivalently be given by its
Rees-algebra charts and the corresponding proper map.
For compact data, the core family and base, the cyclic
base cover and finite action, and all maps outside these
central charts are given globally; the local overlaps
agree with them on an annulus. This is the compact
version of the same definition.
\end{definition}

\begin{proposition}[Complete rooted generation and the zero-Euler criterion for arbitrary fibres]\hypertarget{tag-01EY}{}\label{cdp:st:C.2.UN.86}
\StatementRecord{cdp:st:C.2.UN.86}{proved}{01EY}{\cref{def:II:finite-rooted-reconstruction,prop:II:cyclic-root-normalization,cdp:st:C.2.UN.75,cdp:st:C.2.UN.78,cdp:st:C.2.UN.84,cdp:st:C.2.UN.85,cdp:st:C.2.UN.57,cdp:st:C.2.UN.61,cdp:st:C.2.UN.63,cdp:st:C.2.UN.71,cdp:st:C.2.UN.73}}

Every proper degeneration of connected complex two-tori
from a smooth threefold, with arbitrary full central
fibre, is the output of the finite construction in
\cref{def:II:finite-rooted-reconstruction}. Conversely every such input produces a proper
degeneration of connected complex two-tori from a smooth
threefold. The inverse record recovers the given
family exactly. For compact threefolds over compact
curves, the same assertions hold for the compact data.

For the data of \cref{def:II:finite-rooted-reconstruction},
put $S_*=f_*^{-1}(0)$ and use the quotient components $D_i$
of \eqref{cdp:eq:C.AR.1}. Define $D_I^\circ,d_I,\nu_I,R_N$
and $\Delta_\pi$ by the common formulas
\eqref{cdp:eq:C.AR.2} in
\cref{prop:II:cyclic-root-normalization}, with the same $N,m_i$
and ordinary blowup record $\pi$.
Let \(\Delta_\sigma\) be the signed normal-cone Euler
change in the forward direction \(Z\dashrightarrow X_*\)
of the sigma record. Then its output satisfies
\begin{equation}
 e\bigl(f^{-1}(0)_{\mathrm{red}}\bigr)
       =e(S_*)-\Delta_\sigma-R_N-\Delta_\pi.
 \hypertarget{tag-00AA}{}\label{cdp:eq:C.AR.4}
\end{equation}
Every zero-Euler fibre is therefore obtained, and only
such fibres are obtained, by the admissible data with
\begin{equation}
 \Delta_\sigma+R_N+\Delta_\pi=
 \begin{cases}
  2[N_0:\Gamma_0],&\text{for the positive toric core},\\
  0,&\text{for each other core}.
 \end{cases}
 \hypertarget{tag-00AB}{}\label{cdp:eq:C.AR.5}
\end{equation}
Here \(N_0\) and \(\Gamma_0\) are the two-dimensional
affine lattices of the toric core; the root exponent
\(N\) in \eqref{cdp:eq:C.AR.1} is a separate integer. The construction
in \eqref{cdp:eq:C.AR.5} specifies spaces, bundles, finite maps,
centres and analytic gluings, as well as their Euler
numbers. Reversing the record recovers its normal extensions and
all transverse jets.

The choice between a canonical-flat core and a positive
Hopf elliptic-network core is intrinsic to the original
degeneration. Use the full original divisor and primitive
generator coefficients in \eqref{cdp:eq:C.AR.6}.
Let \(C_\alpha\to T_\alpha\) be the finite normal
root cover of the normalization of \(S_\alpha\) defined
by \eqref{cdp:eq:C.RK.18}, and let \(\tau\) be the density threshold
\eqref{cdp:eq:C.RK.19}. Set
\begin{equation}
 H=\max\left\{\frac{a_\alpha+1}{n_\alpha}:
       C_\alpha\text{ has a component with Hopf minimal model}
       \right\},
 \hypertarget{tag-00AD}{}\label{cdp:eq:C.AR.7}
\end{equation}
with \(H=\tau\) when this set is empty. Then \(H\ge\tau\).
For every divisible semistable root, its highest
Hopf canonical order, when such a component exists, is
\begin{equation}
                         h=N(H-\tau).
 \hypertarget{tag-00AE}{}\label{cdp:eq:C.AR.8}
\end{equation}
Thus \(H=\tau\) gives the canonical-flat reduction,
and \(H>\tau\) gives the positive-Hopf reduction of
\cref{cdp:st:C.2.UN.84,cdp:st:C.2.UN.85}. The covers, threshold and this branch test
do not depend on an embedded resolution or on its
chosen divisible root exponent.
\end{proposition}

\begin{proof}

\begin{enumerate}[label=\arabic*.,leftmargin=*]
\item \par\medskip\noindent\textbf{Forward construction from the original fibre.}
Apply \cref{cdp:st:C.2.UN.17} and
\cref{prop:II:cyclic-root-normalization} to the given proper
smooth-threefold torus degeneration. Choose its embedded
resolution $\pi:Y\to X$ and $N=\operatorname{lcm}_i m_i$.
The resolved full divisor and canonical orders are
\begin{equation}
 \operatorname{div}(f\pi)=\sum_i m_iD_i,
 \qquad \operatorname{div}(\pi^*\Omega)=\sum_i k_iD_i,
 \qquad N=\operatorname{lcm}_i m_i.
 \hypertarget{tag-00AF}{}\label{cdp:eq:C.AR.9}
\end{equation}
The maps in \eqref{cdp:eq:C.AR.10}, with the lattice
\eqref{cdp:eq:C.AR.11}, are precisely that owner's normalization
and equivariant height-one resolution. The labelled monomial-ideal
maps use the original units and glue to a proper $q$.
Their root $Z$ is connected and smooth, with full reduced
SNC central fibre. The maps $q$ and $\pi$ are isomorphisms
off their central fibres; $\rho$ there is the finite \'etale
base-change map induced by $t=s^N$.

\item Apply the complete reduced-fibre classification in \cref{cdp:st:C.2.UN.85} to
\(Z\), with comparison
\(\phi:Z\dashrightarrow X_*\). This gives
one of the listed cores and all the inverse
chart maps in the input. The root action
becomes the birational action
\begin{equation}
            \gamma_* =\phi\gamma\phi^{-1}
                            \quad(\gamma\in\mu_N).
 \hypertarget{tag-00A3}{}\label{cdp:eq:C.AR.12}
\end{equation}
Its group law holds on the dense punctured
family, hence as meromorphic maps. Reversing
\(\phi\) first restores the original regular
action on \(Z\). Applying the
equivariant coarsening then gives \(Y_N\).
The local invariant algebra is exactly
\eqref{cdp:eq:C.RK.14}: $A=\OO_Y$, and $B$ is the integral
closure of $A[s]/(s^N-f\pi)$, including any split local branches.
Thus $B^{\mu_N}=A$ by the common root owner,
and the full comparison is
\begin{equation}
 X_*\dashrightarrow Z\xrightarrow{q}Y_N
              \longrightarrow Y_N/\mu_N=Y
                       \xrightarrow{\pi}X.
 \hypertarget{tag-00A4}{}\label{cdp:eq:C.AR.13}
\end{equation}
Every arrow and overlap in this diagram has
been included in the input. The inverse record restores
the regular action before its equivariant coarsening and
finite quotient.

\par\medskip\noindent\textbf{The converse construction.}
Begin with \cref{def:II:finite-rooted-reconstruction}. The inverse sigma maps give the
smooth proper \(Z\) of the datum, and the
chart identities give a regular finite action. Its equivariant coarsening and
invariant-coordinate charts construct the
smooth \(Y\) and its function \eqref{cdp:eq:C.AR.1}.
The subsequent ordinary contraction maps
construct the smooth \(X\) and its descended
function. Properness follows from the
proper core and its proper maps;
for a local overlap replacement it is the
proper-graph and annular comparison of
\cref{cdp:st:C.2.UN.75}. The ordinary target maps and finite
quotient are proper, so the final family
is proper after the common disc has been
shrunk if required.

\item For a fixed nonzero \(t\), the \(N\)
choices of \(s\) are permuted freely and
transitively by \(\mu_N\). Inclusion of
any one \(Z_s\) into their union induces
a biholomorphism onto the quotient fibre.
The coarsening and ordinary maps are
isomorphisms there. Consequently the
final general fibre is exactly that
connected complex two-torus. Since the
total space is smooth, the nonconstant
base function is a non-zero-divisor;
the map to the smooth curve is flat.
Its Stein factor has generic degree one
and its base is normal, so all its fibres
are connected. These deductions use only
the maps and coordinate identities in the definition.

\item The inverse elementary record is therefore
a constructive description of the resulting
analytic support, with its ambient
neighbourhood. Forward construction shows
that every original support has such a
record. The converse shows that every
admissible record produces such a support.
Isomorphic labelled atlases and commuting
record maps give biholomorphic outputs
by their explicit glued isomorphism.
The construction associates its output to the atlas and
commuting arrows of each such record.

\par\medskip\noindent\textbf{The Euler changes of the construction maps.}
Apply \cref{prop:II:cyclic-root-normalization} with the
resolution $\pi$, exponent $N$ and coarsening $q$ in
\eqref{cdp:eq:C.AR.10}. Its stratum cover has degree $d_I$,
its toric fibre has Euler number $\nu_I$, and its determinant
and fixed-point computations include every global sheet
permutation. This is exactly
\eqref{cdp:eq:C.AR.14}, $e(S_Z)-e(S_Y)=R_N$.

\item The ordinary blowups give
\(e(S_Y)-e(S)=\Delta_\pi\), and the
forward sigma record gives
\(e(S_*)-e(S_Z)=\Delta_\sigma\).
For an arrow with centre ideal \(I\),
its signed term is the difference between
the Euler numbers of the \(\operatorname{Proj}(\bigoplus I^a/I^{a+1})\)
and its centre, on the finite constructible
subdivision of \cref{cdp:st:C.2.UN.75}. Restored overlaps
cancel. Combining these three exact
comparisons proves \eqref{cdp:eq:C.AR.4}. The core
Euler numbers are those of \cref{cdp:st:C.2.UN.85}, proving
\eqref{cdp:eq:C.AR.5}. Since both construction
directions have already been established,
this selects the full class of zero-Euler
outputs, including singular and nonreduced
original fibres.

\par\medskip\noindent\textbf{The intrinsic endpoint test.}
Use the finite covers $C_\alpha\to T_\alpha$ and threshold
$\tau$ of \cref{prop:II:cyclic-root-normalization} on the same
original family and generator. The strict-cover map is
$w_j=s^{N/n_\alpha}/g_j$, with transitions
$w_k=(g_j/g_k)w_j$. It gives a proper birational resolution
of every connected component of $C_\alpha$. The owner's
projective exceptional-cover and ruled semistable-component
comparisons account for all Hopf root components.

The density equality \eqref{cdp:eq:C.AR.15} and the strict-order
formula \eqref{cdp:eq:C.AR.16} apply with
$(k_i,m_i)=(a_\alpha,n_\alpha)$ on each original strict component.
Their maximum over precisely those Hopf covers is
$N(H-\tau)$, proving \eqref{cdp:eq:C.AR.8} and its integrality.
When there is no Hopf root component, \cref{cdp:st:C.2.UN.75}
applies. Otherwise all its Hopf orders vanish exactly when
$H=\tau$; this is that same canonical-flat criterion. For
$H>\tau$, \cref{cdp:st:C.2.UN.84,cdp:st:C.2.UN.85} give the
positive-Hopf core. Both the covers and the density are defined
on the original family, so this test is independent of the
resolution, subdivision and divisible exponent.

\par\medskip\noindent\textbf{Compact families and their prescribed neighbourhoods.}
Apply the compact-cover construction in
\cref{prop:II:cyclic-root-normalization} at the selected value $c$.
It uses the original critical set $\Sigma$, a pole point
$p\notin\Sigma\cup\{c\}$, and a meromorphic function $a$
with jet $t+O(t^2)$ at $c$ and value one at the other critical
points. The connected normalized cover $s^N=a$ is \'etale
over those other critical values. The explicit map
$\sigma=s/u(t)^{1/N}$, where $a=t u(t)$, identifies its
local root with $\sigma^N=t$, equivariantly and with inverse
$s=\sigma u(t)^{1/N}$. All additional branch values lie
over submersion charts, so the total root family is smooth
away from the selected fibre. Perform the root resolution
and modification there, using their identity maps
on the annulus and complement.

\item The finite action is conjugated by that
global modification and becomes
regular again upon reversal. The quotient
and original blowdowns recover the
original compact family. Conversely the
global compact admissible
data give a compact output by the same
proper maps.

\item The central support parameters of the
flat cycles and elliptic networks are
the complete marked line, torsor,
finite-character and gluing data in
\cref{cdp:st:C.2.UN.57,cdp:st:C.2.UN.61,cdp:st:C.2.UN.62,cdp:st:C.2.UN.63}. The positive
toric core is parametrized by its
periodic affine triangulation and
normalized log gluing data in \cref{cdp:st:C.2.UN.71}.
Their compact occurrence is proved
in those statements and \cref{cdp:st:C.2.UN.73}.
Those constructions supply compact starting families for records
defined on their neighbourhoods.
For a given family, \cref{def:II:finite-rooted-reconstruction}
includes its endpoint atlas, all transverse terms and the
overlap isomorphisms used by the inverse arrows. The diagram
\eqref{cdp:eq:C.AR.13} recovers that analytic neighbourhood
and, for compact data, its global family.

\end{enumerate}
\end{proof}







\section{Euler sign and equality for arbitrary torus fibres}\hypertarget{tag-034C}{}\label{sec:II:all:12}


\begin{theorem}[The Euler sign for every degeneration and its equality criterion]
\hypertarget{tag-01BZ}{}\label{cdp:st:C.2.G.1}
\StatementRecord{cdp:st:C.2.G.1}{proved}{01BZ}{\cref{cdp:st:D.1.9}, \cref{cdp:st:A.3.4}, \cref{cdp:st:A.5.2}, \cref{cdp:st:A.5.3}, \cref{cdp:st:C.2.MH.1}, \cref{cdp:st:C.2.MH.2}, \cref{night:P4.26}, \cref{night:P4.32}, \cref{night:P4.36}, \cref{night:P4.37}, \cref{cdp:st:C.2.INT.1}, \cref{cdp:st:C.2.INT.2}, \cref{cdp:st:C.2.INT.3}, \cref{cdp:st:C2a}; \cite[Proposition~2.1 and Theorem~2.2, pp.~842--843]{Peters2021}}

Let $Y$ be a connected smooth compact complex threefold with
$b_1(Y)=b_2(Y)=0$, and let $f:Y\to B$ be a proper surjective
holomorphic map with connected fibres to a smooth compact curve.
Assume that the smooth general fibre is a complex two-torus.
The full fibres may have arbitrary positive component multiplicities
and their reduced supports may have arbitrary singularities. Then
\begin{equation}
 B\simeq\PP^1,\qquad
 c_3(Y)=e(Y)\in\{0,2\},\qquad b_3(Y)=2-c_3(Y).
 \hypertarget{tag-00KW}{}\label{cdp:eq:C.G.1}
\end{equation}
Every reduced fibre support has nonnegative Euler number.
An irreducible singular support with rational normalization
resolution has Euler number two; every other irreducible support
has Euler number zero.

The equality $c_3(Y)=0$ holds exactly when every support has one
of the following descriptions. A smooth irreducible support has
a finite unramified torus cover. A singular irreducible support
has smooth minimal elliptic ruled or smooth minimal Hopf
normalization. A reducible support has normalization a disjoint
union of smooth minimal Hopf or non-K\"ahler elliptic surfaces;
its reduced conductor image and inverse image are disjoint
unions of smooth elliptic curves, and the conductor branch maps
are finite unramified covers. For $Y_b=f^*(b)=\sum_i m_iS_i$ and $S=(Y_b)_{\rm red}$,
these are the normalization $\nu:\coprod_iT_i\to S$ and the
conductor defined by
$\mathfrak c=\operatorname{Ann}_{\OO_S}
 (\nu_*\OO_{\coprod_iT_i}/\OO_S)$.
Their maps to $Y_b$ compose with the reduction inclusion
$j:S\hookrightarrow Y_b$.
\end{theorem}


\begin{proof}
For each full fibre $Y_b=f^*(b)=\sum_i m_iS_i$, its reduction
$S=(Y_b)_{\rm red}$ and normalization $\nu:\coprod_iT_i\to S$
give the structure-sheaf maps
\[
 \OO_{Y_b}\twoheadrightarrow j_*\OO_S
       \hookrightarrow(j\circ\nu)_*\OO_{\coprod_iT_i}.
\]
Their composite defines $j\circ\nu:\coprod_iT_i\to Y_b$.
The conductor ideal $\mathfrak c$ in the statement defines
$\Gamma=V(\mathfrak c)\subset S$ and its inverse image
$A=V(\mathfrak c\OO_{\coprod_iT_i})$.
The reduced curves used below are $C=\Gamma_{\rm red}$ and
$A_i=(A\cap T_i)_{\rm red}$.

\begin{enumerate}[label=\arabic*.,leftmargin=*]
\item \textbf{The base and the finite character.}
The beginning of connected-fibre Leray injects
$H^1(B,\Q)$ into $H^1(Y,\Q)=0$, so $B=\PP^1$.
\cref{cdp:st:D.1.9} supplies an integral section
of $R^1f_*\Z$ which restricts nontrivially to each reduced
support and to each nearby torus. On a small disc its
Leray lift is an integral class on the entire tube.
Reduction modulo a prime which does not divide its fibre
restriction gives the finite cyclic character used in
\cref{cdp:st:C.2.MH.2}, on the full Cartier fibre $f^*(b)$.

\item \textbf{Smooth supports.}
\cref{cdp:st:A.5.2} excludes a singular normal
support. For a smooth support $S$ in a full fibre $mS$,
the normalized change $t=s^m$ is, in a local equation
$t=uz^m$, the map
\[
 w^m=u,\qquad s=wz.
\]
These charts give a finite unramified central cover of $S$
and a proper holomorphic submersion over the new disc,
as in \cref{cdp:st:A.3.4}. Its nearby fibre is the
original torus. Smooth proper transport gives $b_1=4$
and rational first Chern class zero on the central cover.
It is K\"ahler by the even-$b_1$ criterion in
\cite[Proposition~2.1, p.~842]{Peters2021}.
Its rationally trivial canonical class excludes an exceptional
curve, whose canonical intersection would be $-1$.
Every effective pluricanonical divisor has zero K\"ahler area
and hence is empty, so every nonzero pluricanonical space has
dimension one. In the classification of
\cite[Theorem~2.2, p.~843]{Peters2021}, the remaining
Kodaira-dimension-zero case with $b_1=4$ is a torus;
the ruled alternative has canonical intersection $-2$ with
a ruling fibre and is excluded by the same Chern class.
The degree-$m$ unramified central cover therefore gives
$e(S)=e(S')/m=0$, where $S'$ denotes that cover.

\item \textbf{Normalization resolutions.}
For a singular irreducible support $S$, write $Z\to T\to S$
for the minimal resolution of its normalization. The canonical evaluation and
\cref{cdp:st:A.5.3} give a nonempty effective divisor $D$
with $K_Z+D$ torsion. A positive pluricanonical section,
raised to a power killing this torsion, would give a
nonzero holomorphic function vanishing along a positive
multiple of $D$. Compactness excludes that function.
Hence $\kappa(Z)=-\infty$. The classification in
\cite[Theorem~2.2, p.~843]{Peters2021} gives a rational
resolution, a resolution ruled over a curve of positive
genus, or a resolution of class VII. We apply the following
results to these same maps $Z\to T\to S$.

\item \textbf{Irreducible supports.}
For a rational resolution,
\cref{cdp:st:C.2.MH.2} gives $e(S)=2$. For a resolution
ruled over a curve of genus $g\ge1$, the same proposition
gives $e(S)=2-2g$. The local irreducible inequality in
\cref{night:P4.26} gives $e(S)\ge0$, and therefore $g=1$
and $e(S)=0$. The irreducible-support assertion of that
proposition uses the effective conductor, local trace
inequality and fundamental group of this same $S$. The zero-Euler
comparison in \cref{night:P4.32} then identifies $T$ as
a smooth minimal elliptic ruled surface.

For a class-VII resolution, \cref{night:P4.37} excludes
every minimal model with positive second Betti number
by its saturated-foliation Chern identity and the
negative-definite curve form. The curve-free Inoue
case cannot carry the nonempty effective conductor.
The remaining Hopf case is treated by
\cref{night:P4.36}: its derivative and foliation
calculations make $T$ smooth and minimal, with $e(S)=0$.
These are the single-support conclusions for the
normalization map $T\to S$ fixed in step~3.

\item \textbf{Reducible supports.}
For each normalization resolution, an adjacent global
component gives a nonzero effective divisor with
rationally trivial first Chern class. Integration against
a K\"ahler form would give both positive area and zero
cohomological pairing, so the resolution is non-K\"ahler. \cref{cdp:st:C.2.INT.1}
gives the Picard restriction to each elliptic
curve and excludes a Hopf--Enoki or elliptic--Enoki
interface. \cref{cdp:st:C.2.INT.2} excludes
the remaining all-Enoki possibility using the full
multiplicities, the canonical coefficient maximum,
and the local vanishing sphere of square $-2$.

\cref{cdp:st:C.2.INT.3} applies to this
reduced support. Write $\nu:\coprod_iT_i\to S$ for its
normalization, $A_i$ for the reduced conductor on $T_i$,
and $C$ for the reduced conductor image. Restriction of
$\nu$ gives an isomorphism
$\coprod_i(T_i\setminus A_i)\simeq S\setminus C$, and hence
\[
 e(S)=\sum_i e(T_i\setminus A_i)+e(C),
 \qquad e(T_i\setminus A_i)\ge0,\quad e(C)\ge0.
\]
The scheme conductor enters the coherent-duality argument
of that proposition, and the exceptional loci enter its
topological calculation. Its equality characterization says that
the normalizations are the stated smooth minimal
surfaces, the conductor curves are smooth elliptic,
and their maps are finite unramified covers.

\item \textbf{The global Euler sum.}
The finite-critical-value and Euler-additivity
calculation of \cref{cdp:st:C2a}, together with
complex-oriented Chern--Gauss--Bonnet, gives
\begin{equation}
 c_3(Y)=e(Y)=\sum_{b\in B}e((Y_b)_{\rm red})\ge0.
 \hypertarget{tag-00KX}{}\label{cdp:eq:C.G.2}
\end{equation}
The smooth torus contributes zero over the regular
base. Poincar\'e duality gives $e(Y)=2-b_3(Y)$.
The nonsingular middle cup-product pairing is
skew-symmetric over $\Q$, so $b_3(Y)$ is even.
Its nonnegativity now proves \eqref{cdp:eq:C.G.1}.

\item \textbf{The equality criterion.}
Every term in \eqref{cdp:eq:C.G.2} is nonnegative;
the sum vanishes precisely when every term does.
Steps 2 and 4 give the irreducible forms, and
Step 5 gives the reducible forms with their exact
conductor maps. Conversely each of these forms has
zero Euler number by the same local calculations,
so their occurrence on every support makes the
global sum zero.

\end{enumerate}
\end{proof}

\begin{proposition}[Euler equality under trivial torus monodromy]
\hypertarget{tag-01G9}{}\label{cdp:st:D.2.2}
\StatementRecord{cdp:st:D.2.2}{proved}{01G9}{\cref{cdp:st:C.2.G.1,cdp:st:D.2.5,cdp:st:D.2.7}}
\hypertarget{tag-01GB}{}\label{cdp:st:D.2.4}
\StatementRecord{cdp:st:D.2.4}{proved}{01GB}{}
The class of smooth compact connected complex threefolds with $b_1=b_2=0$,
a proper surjective holomorphic map with connected fibres onto a smooth compact curve,
smooth two-torus general fibre, and trivial rational monodromy is nonempty.
Every member satisfies
\[
 c_3(Y)=0,\qquad b_3(Y)=2.
\]
Every individual reduced fibre support has Euler number zero and
the geometric description of \cref{cdp:st:C.2.G.1}.
Arbitrary singularities and full-fibre multiplicities are allowed.
There is an example with normal-crossing reduced supports,
integral $H_1(Y,\Z)=H_2(Y,\Z)=0$, and algebraic dimension one.
\end{proposition}
\begin{proof}
Apply \cref{cdp:st:C.2.G.1,cdp:st:D.2.5} to the same
threefold $Y$ and map $f$. The first theorem identifies the
base with $\PP^1$, as required by the second proposition.
Their lower and upper bounds give $c_3(Y)=0$.
Poincar\'e duality gives $b_3(Y)=2-c_3(Y)=2$.
The Euler sum and equality criterion of \cref{cdp:st:C.2.G.1}
apply to these same full fibres with their given multiplicities;
each nonnegative summand is zero and has the stated geometric form.
This applies in particular to every member with normal-crossing
reduced supports. The construction in \cref{cdp:st:D.2.7}
has such supports, the stated integral homology and algebraic
dimension, and identity integral monodromy. It therefore proves
nonemptiness of both the full class and this subclass.
\end{proof}












\clearpage
\part{The remaining splitting problem}\label{part:II:remaining-splitting}
\section{The remaining splitting problem}\hypertarget{tag-000B}{}\label{app:II:programs}
\begin{problem}[The remaining comb bundle splitting]\hypertarget{tag-01SW}{}\label{inv:E2.3}
\StatementRecord{inv:E2.3}{open}{01SW}{\SeriesRef{def:II:comb-cup-kernel}, \SeriesRef{prop:II:comb-kernel-bounds}, \SeriesRef{prop:II:comb-cup-duality}, \SeriesRef{cdp:st:E15}, \SeriesRef{inv:E2.5}}
For the positive combs of \cref{def:II:comb-cup-kernel}, determine
the splitting of the rank-three bundle $K_m$ for general tooth
counts and for special members outside the nonempty open
twelve-tooth family of \cref{inv:E2.5}, together with its exact
number of sections.
On that family,
$K_{12}\simeq\mathcal O_B(-8)\oplus\mathcal O_B(-9)^{\oplus2}$
and $h^0(K_{12})=0$, attaining the lower Hodge bounds
$h^{1,1}=49$ and $h^{1,2}=48$.
The degree formula and uniform bounds are
\cref{prop:II:comb-kernel-bounds};
\cref{prop:II:comb-cup-duality} gives
$K_m\simeq\rho_*(N^*\otimes\det G_m)$ for the tangent-cup
cokernel $N$ of the admissible period section.
$h^0(K_m)$ determines the two middle Hodge numbers
through \eqref{invII:eq:E.5}.
For all these combs, \cref{cdp:st:E15} gives the four tangent
cohomology dimensions and the smooth reduced Kuranishi germ
of dimension $8m+1$, with zero obstruction map.
\end{problem}

\clearpage

\section*{Programs}
\addcontentsline{toc}{section}{Programs}
The following programs implement the indicated computations.
The records page cited in the methodology gives their source files.

{\footnotesize
\renewcommand{\arraystretch}{.9}
\setlength{\tabcolsep}{4pt}
\begin{longtable}{@{}>{\raggedright\arraybackslash}p{.57\textwidth}>{\raggedright\arraybackslash}p{.39\textwidth}@{}}
\caption{Programs and their mathematical scope.}\label{tab:II:programs}\\
\toprule Computation and statement & Source file \\
\midrule\endfirsthead
\toprule Computation and statement & Source file \\
\midrule\endhead
\bottomrule\endfoot
Tangent index; integral monodromy, including Wang's four matrices and integral coinvariant quotient; Kodaira and Hopf surfaces; 120 unit quaternions, 14,400 products and a surjective commutator map over $\Q(\sqrt5)$ represented by rational pairs (\cref{cdp:st:C.1.1,cdp:st:C.3.2,cdp:st:C.3.3,cdp:st:D.1.6}) & \nolinkurl{pro_returns/S6_03_CDP_checkpoint_13_20261116T0136Z/S6_03_CDP/code/D_exact_checks.py} \\[3pt]
Normalization, inverse Jacobian module and residue poles (\cref{cdp:st:A.2.5}) & \nolinkurl{pro_returns/S6_03_CDP_checkpoint_13_20261116T0136Z/S6_03_CDP/code/A_vertical_pole.py} \\[3pt]
Curve equations, transverse jets and conormal valuations (\cref{cdp:st:A.2.6}) & \nolinkurl{pro_returns/S6_03_CDP_checkpoint_13_20261116T0136Z/S6_03_CDP/code/A_non_lci_cusp.py} \\[3pt]
Exterior-square parity examples and Euler numbers of blowup families (\cref{cdp:st:C.2.UN.1,cdp:st:C.2.UN.2}) & \nolinkurl{pro_returns/S6_03_CDP_checkpoint_13_20261116T0136Z/S6_03_CDP/code/root_unrestricted_euler.py} \\[3pt]
Euler and integral cohomology of vertical point and curve blowups (\cref{cdp:st:C.2.UN.4}) & \nolinkurl{pro_returns/S6_03_CDP_checkpoint_13_20261116T0136Z/S6_03_CDP/code/E_vertical_refinements.py} \\[3pt]
Exact rational jet restrictions at orders $2,3,4,6,8$ and conductor ideals of the two genus-one plane cubics (\cref{night:P4.14}) & \nolinkurl{pro_returns/S6_06_BLIND_checkpoint_15_20261117T1457Z/S6_06_BLIND/code/part4/genus_one_conductor.py} \\[3pt]
Jet covariance, binary-cubic discriminants and inverse-Jacobian pole coefficients (\cref{cdp:st:C2d}) & \nolinkurl{pro_returns/S6_03_CDP_checkpoint_13_20261116T0136Z/S6_03_CDP/code/C_jets.py} \\[3pt]
Integral compact intersection lattices, primitive boundary kernels, blowup charts and the alternative two-blowup compactification (\cref{night:P4.16}) & \nolinkurl{pro_returns/S6_06_BLIND_checkpoint_15_20261117T1457Z/S6_06_BLIND/code/part4/milnor_intersection.py} \\[3pt]
Coprime exponent matrices, transfer pairings, normalization minors and inverse coordinates on the affine nearby fibre (\cref{night:P4.18,night:P4.19}) & \nolinkurl{pro_returns/S6_06_BLIND_checkpoint_15_20261117T1457Z/S6_06_BLIND/code/part4/binomial_milnor.py} \\[3pt]
Seifert fractions, boundary homology and intersection ranks (\cref{cdp:st:A.5.15}) & \nolinkurl{pro_returns/S6_03_CDP_checkpoint_13_20261116T0136Z/S6_03_CDP/code/E_Hirzebruch.py} \\[3pt]
Integral homology-kernel maps, exponent inverses, rational Seifert elimination and normalization conductor (\cref{cdp:st:A.5.15}) & \nolinkurl{code/S6_SERIES_PUBLISH_II_B2_B5_comparison_certificate.py} \\[3pt]
Endpoint intersection signs, weighted identities, ideal generators and boundary right inverses (\cref{cdp:st:A.5.21}) & \nolinkurl{code/S6_SERIES_PUBLISH_II_B2_BC_comparison_certificate.py} \\[3pt]
Source coordinate maps, the Albanese lattice, and the cap and cusp descent identities (\cref{ii:LR:covering}) & \nolinkurl{code/S6_SERIES_PUBLISH_II_B2_UH_comparison_certificate.py} \\[3pt]
Symbolic squares, tacnode invariants and branch budgets (\cref{cdp:st:A.5.18}; \eqref{cdp:eq:A.T.10}--\eqref{cdp:eq:A.T.12}) & \nolinkurl{code/S6_SERIES_PUBLISH_II_B2_B8_prereading_certificate.py}; \nolinkurl{pro_returns/S6_03_CDP_checkpoint_13_20261116T0136Z/S6_03_CDP/code/E_tacnode_bounds.py} \\[3pt]
Conormal coordinate maps and pole bounds (\cref{cdp:st:A.5.18}) & \nolinkurl{code/S6_SERIES_PUBLISH_II_B2_B8_comparison_certificate.py}; \nolinkurl{pro_returns/S6_03_CDP_checkpoint_13_20261116T0136Z/S6_03_CDP/code/E_low_delta_tacnode.py} \\[3pt]
Symbolic cubes and cubic discriminant valuations (\cref{cdp:st:A.5.19}) & \nolinkurl{code/S6_SERIES_PUBLISH_II_B2_B9_prereading_certificate.py} \\[3pt]
Conormal cubic poles and chart transport (\cref{cdp:st:A.5.19}) & \nolinkurl{code/S6_SERIES_PUBLISH_II_B2_B9_comparison_certificate.py} \\[3pt]
Suspension coordinate inverses and orientation (\cref{cdp:st:A.5.17}) & \nolinkurl{code/S6_SERIES_PUBLISH_II_B2_B7_prereading_certificate.py} \\[3pt]
Power Hessians and quadratic coordinate maps (\cref{cdp:st:A.5.17}) & \nolinkurl{code/S6_SERIES_PUBLISH_II_B2_B7_comparison_certificate.py} \\[3pt]
Tacnode discriminant and ramified coordinate maps (\cref{cdp:st:C.2.RAT.2}) & \nolinkurl{code/S6_SERIES_PUBLISH_II_B2_CE_prereading_certificate.py} \\[3pt]
Quartic invariant and ruled-boundary Euler counts (\cref{cdp:st:C.2.RAT.2}) & \nolinkurl{code/S6_SERIES_PUBLISH_II_B2_CE_comparison_certificate.py} \\[3pt]
Tacnode conormal poles and paired-attachment budget (\cref{cdp:st:C.2.RAT.2}) & \nolinkurl{code/S6_SERIES_PUBLISH_II_B2_CG_prereading_certificate.py} \\[3pt]
Conductor factor and degree-two incidence count (\cref{cdp:st:C.2.RAT.2}) & \nolinkurl{code/S6_SERIES_PUBLISH_II_B2_CG_comparison_certificate.py} \\[3pt]
Cubic root attachments and Euler count (\cref{prop:II:cubic-root-attachments}) & \nolinkurl{code/S6_SERIES_PUBLISH_II_B2_CH_prereading_certificate.py} \\[3pt]
Simple-root conductor derivative (\cref{prop:II:cubic-root-attachments}) & \nolinkurl{code/S6_SERIES_PUBLISH_II_B2_CH_comparison_certificate.py} \\[3pt]
Delta-five cubic monodromy and compact-core rank (\cref{prop:II:cubic-root-attachments}) & \nolinkurl{code/S6_SERIES_PUBLISH_II_B2_CJ_prereading_certificate.py} \\[3pt]
Boundary relations and disconnected Euler count (\cref{prop:II:cubic-root-attachments}) & \nolinkurl{code/S6_SERIES_PUBLISH_II_B2_CJ_comparison_certificate.py} \\[3pt]
Single-zero cubic local bounds and compact pairing (\cref{prop:II:cubic-root-attachments}) & \nolinkurl{code/S6_SERIES_PUBLISH_II_B2_CK_prereading_certificate.py} \\[3pt]
Newton polyhedron and attachment-loss identities (\cref{prop:II:cubic-root-attachments}) & \nolinkurl{code/S6_SERIES_PUBLISH_II_B2_CK_comparison_certificate.py} \\[3pt]
Cubic root divisors, cyclic normalizations and compact pairings (\cref{prop:II:cubic-root-attachments}) & \nolinkurl{code/S6_SERIES_PUBLISH_II_B2_CL_prereading_certificate.py} \\[3pt]
Weighted cyclic blowup charts and local suspension invariants (\cref{prop:II:cubic-root-attachments}) & \nolinkurl{code/S6_SERIES_PUBLISH_II_B2_CL_comparison_certificate.py} \\[3pt]
Spectral bounds and marked conductor exponents (\cref{night:P4R.20}) & \nolinkurl{code/S6_SERIES_PUBLISH_II_A2_YE_original_conductor.py}; \nolinkurl{pro_returns/S6_06_BLIND_checkpoint_15_20261117T1457Z/S6_06_BLIND/code/part4/degree_two_triple_point.py} \\[3pt]
Primary conductor defects and horizontal branch-value counts (\cref{night:P4R.18}) & \nolinkurl{code/S6_SERIES_PUBLISH_II_A2_YB_original_graph.py}; \nolinkurl{pro_returns/S6_06_BLIND_checkpoint_15_20261117T1457Z/S6_06_BLIND/code/part4/horizontal_primary.py} \\[3pt]
Normal-cubic weights, finite maps and two-direction inverses (\cref{night:P4R.19}) & \nolinkurl{code/S6_SERIES_PUBLISH_II_A2_YC_weighted_cubic.py}; \nolinkurl{pro_returns/S6_06_BLIND_checkpoint_15_20261117T1457Z/S6_06_BLIND/code/part4/d4_spectral_triple.py} \\[3pt]
Matching-sphere cap coordinates and orientation (\cref{night:P4R.19}) & \nolinkurl{code/S6_SERIES_PUBLISH_II_A2_YC_original_cap.py} \\[3pt]
Class VII Chern and initial quadratic calculations (\cref{night:P4.37}) & \nolinkurl{code/S6_SERIES_PUBLISH_II_A2_XA_original_completion.py} \\[3pt]
Integral monodromies, centralizers, fixed and norm lattices, root and nodal defects, and period covariance (\cref{push:src:E.5.2,def:II:two-parameter-period-data,push:src:E.5.3}) & \nolinkurl{pro_returns/S6_02_PUSH_final_20261115T2217Z/S6_02_PUSH/code/e_push_two_parameter.py} \\[3pt]
Canonical ratio and translated-integer positivity certificate (\cref{push:src:E.2.15:smooth}) & \nolinkurl{code/S6_SERIES_PUBLISH_II_C2_0330_final_v1_integer_certificate.py} \\[3pt]
Parameter shear, peripheral products and independent cusp coefficients (\cref{push:src:E.2.7}) & \nolinkurl{code/S6_SERIES_PUBLISH_II_C2_033F_final_v1_monodromy_prerelease.py} \\[3pt]
Integral basis transport and cyclotomic invariant-space certificates (\cref{push:src:E.2.7}) & \nolinkurl{code/S6_SERIES_PUBLISH_II_C2_033F_final_v1_basis_cyclotomic_postrelease.py} \\[3pt]
Invariant covectors, cusp invariant space and exact triangle-group representations (\cref{push:src:E.2.16}) & \nolinkurl{code/S6_SERIES_PUBLISH_II_C2_0331_final_v2_exact_certificate.py} \\[3pt]
Cusp coordinate maps, period kernels and chart-unit inverses (\cref{push:src:E.2.20}) & \nolinkurl{code/S6_SERIES_PUBLISH_II_C2_0336_final_v1_coordinate_certificate.py} \\[3pt]
Integral cusp lattice and Laurent-monomial character certificates (\cref{push:src:E.2.19}) & \nolinkurl{code/S6_SERIES_PUBLISH_II_C2_0334_final_v1_exact_certificate.py} \\[3pt]
Original cusp invariant-space and character calculation (\cref{push:src:E.2.19}) & \nolinkurl{pro_returns/S6_02_PUSH_final_20261115T2217Z/S6_02_PUSH/code/d_push_cusp_torus_action.py} \\[3pt]
Root-degree, cup-product and Chern-class identities for the combs (\cref{inv:E2.2}) & \nolinkurl{code/S6_SERIES_PUBLISH_II_root_01SV_prereading_certificate.py} \\[3pt]
Conic coordinate maps, component multiplicities, weighted scalings and oriented sphere Euler numbers (\cref{night:P4.25}) & \nolinkurl{code/S6_SERIES_PUBLISH_II_A_02WX_conic_maps.py} \\[3pt]
Transverse two-jets and critical-value remainders of the double and triple deformations (\cref{night:P4.25}) & \nolinkurl{code/S6_SERIES_PUBLISH_II_A_02WX_deformation_jets_postrelease.py} \\[3pt]
Alternating period pairing, integral monodromy, weighted line cohomology and the branch-value Jacobian (\cref{cdp:st:E13}) & \nolinkurl{code/S6_SERIES_PUBLISH_II_01GK_prereading_certificate.py} \\[3pt]
Restriction-map inverses, residue-frame comparison and the Laurent connection identity (\cref{cdp:st:E13}) & \nolinkurl{code/S6_SERIES_PUBLISH_II_01GK_postrelease_inverse_certificate.py} \\[3pt]
Barycentric retraction identities and the symbolic meridian winding map (\cref{night:P4R.1}) & \nolinkurl{code/S6_SERIES_PUBLISH_II_A_02Y2_neighborhood_certificate.py} \\[3pt]
Integral marking, period and root-coordinate maps, and weighted-line cohomology (\cref{cdp:st:E11}) & \nolinkurl{code/S6_SERIES_PUBLISH_II_01GI_prereading_weight_certificate.py} \\[3pt]
Hilbert-series recurrence and the degree-three Riemann--Roch expansion (\cref{cdp:st:E11}) & \nolinkurl{code/S6_SERIES_PUBLISH_II_01GI_postrelease_hilbert_certificate.py} \\[3pt]
Euler identities from curve and sphere cells, point blowups and class-VII Betti ranks (\cref{night:P4R.2}) & \nolinkurl{code/S6_SERIES_PUBLISH_II_A_02YD_euler_cells.py} \\[3pt]
Meridian relation-space isomorphism, determinant sign and incidence-tree Euler count (\cref{night:P4R.1,night:P4R.4}) & \nolinkurl{code/S6_SERIES_PUBLISH_II_A_02YG_relation_certificate.py} \\[3pt]
Cusp normalization jets, conductor cochains, weighted cohomology and restriction-map inverses (\cref{cdp:st:E14}) & \nolinkurl{code/S6_SERIES_PUBLISH_II_01GL_prereading_exact.py} \\[3pt]
Critical-value norm and exact cusp jet basis changes (\cref{cdp:st:E14}) & \nolinkurl{code/S6_SERIES_PUBLISH_II_01GL_postrelease_critical_norm.py} \\[3pt]
Weighted determinant polynomials, integral period pairings and the Grassmann quotient matrix (\cref{cdp:st:E16}) & \nolinkurl{code/S6_SERIES_PUBLISH_II_01GN_prereading_certificate_v2.py} \\[3pt]
Hodge and tangent frame covariance and the expansion into two-by-two minors (\cref{cdp:st:E16}) & \nolinkurl{code/S6_SERIES_PUBLISH_II_01GN_postrelease_frame_certificate.py} \\[3pt]
Weighted intersection identities, sphere orientation and quadratic Hessian determinant (\cref{night:P4R.6}) & \nolinkurl{code/S6_SERIES_PUBLISH_II_A_02YI_interpolation_certificate.py} \\[3pt]
Quadratic-fibre coordinate inverses and local intersection orientations (\cref{night:P4R.6}) & \nolinkurl{code/S6_SERIES_PUBLISH_II_A_02YI_A1_transport_postrelease.py} \\[3pt]
Curve-incidence Euler coefficients and integral rational-leaf chain cancellation (\cref{night:P4R.7}) & \nolinkurl{code/S6_SERIES_PUBLISH_II_A_02YJ_affine_chain_certificate.py} \\[3pt]
Conductor-cleared jets, binary-cubic discriminants and derivation-graph quotients (\cref{cdp:st:A.5.5,cdp:st:A.5.7}) & \nolinkurl{code/S6_SERIES_PUBLISH_II_BF_prereading_invariant_certificate.py}; \nolinkurl{pro_returns/S6_03_CDP_checkpoint_13_20261116T0136Z/S6_03_CDP/code/E_singular_curve.py} \\[3pt]
Rational centralizer, norm lattice and cross-ratio derivative (\cref{push:src:E.5.3}) & \nolinkurl{code/S6_SERIES_PUBLISH_II_033P_prereading_exact.py} \\[3pt]
Period covariance, integral collar relations and cubic discriminant (\cref{push:src:E.5.3}) & \nolinkurl{code/S6_SERIES_PUBLISH_II_033P_postrelease_exact.py} \\[3pt]
Curve-incidence Euler identities, pole-divisor blowups and rational-character extension through elliptic covering matrices (\cref{night:P4R.5,night:P4R.8}) & \nolinkurl{code/S6_SERIES_PUBLISH_II_A2_Hopf_graph_character.py} \\[3pt]
Hilbert-series numerator, coefficient recurrence and integral-twist Euler identities (\cref{cdp:st:C5b}) & \nolinkurl{code/S6_SERIES_PUBLISH_II_B2_FS_hilbert_certificate.py} \\[3pt]
Chern-root expansion of the tangent Riemann--Roch coefficient (\cref{push:src:E.2.9}) & \nolinkurl{code/S6_SERIES_PUBLISH_II_C2_033H_HRR_postrelease.py} \\[3pt]
Nonintegral-twist quotient maps, theta derivatives, finite characters and cusp contraction (\cref{cdp:st:C5c}) & \nolinkurl{code/S6_SERIES_PUBLISH_II_B2_FT_prereading_certificate.py} \\[3pt]
Tripling-coordinate identity, resonance coefficients and cusp chain homotopy (\cref{cdp:st:C5c}) & \nolinkurl{code/S6_SERIES_PUBLISH_II_B2_FT_comparison_certificate_portable.py} \\[3pt]
Exterior-square inclusion and zero restriction of the four-torus intersection form (\cref{night:P4.51}) & \nolinkurl{pro_returns/S6_06_BLIND_checkpoint_15_20261117T1457Z/S6_06_BLIND/code/part4/local_intersection.py} \\[3pt]
Contraction kernel, rank minors and the zero intersection Gram matrix for a nonzero rational one-class (\cref{night:P4.51}) & \nolinkurl{code/S6_SERIES_PUBLISH_II_A2_XQ_exterior_kernel.py} \\[3pt]
Cyclic-cover monodromy, branched Euler characteristic and normal-clutching residue for the cusp obstruction (\cref{night:P4.28}) & \nolinkurl{code/S6_SERIES_PUBLISH_II_A2_X0_cyclic_cover_square.py} \\[3pt]
Hopf contraction iterates, foliation blowup charts and Chern intersection identities (\cref{night:P4.24}) & \nolinkurl{code/S6_SERIES_PUBLISH_II_A2_WW_original_foliation.py} \\[3pt]
Cubic-cover intersection numbers and compact-core ranks (\cref{cdp:st:C.2.D4.1}) & \nolinkurl{code/S6_SERIES_PUBLISH_II_root_BW_prerelease_intersections.py} \\[3pt]
Cubic-suspension Jacobian spectra by the Koszul Hilbert series (\cref{cdp:st:C.2.RAT.6}) & \nolinkurl{code/S6_SERIES_PUBLISH_II_root_BW_koszul_spectrum.py} \\[3pt]
Conductor pole blowups and curve Euler identities (\cref{cdp:st:C.2.VII.1}) & \nolinkurl{code/S6_SERIES_PUBLISH_II_A2_WV_graph_and_poles.py} \\[3pt]
Cycle intersection kernels and normalization-immersion Chern classes (\cref{night:P4.23}) & \nolinkurl{code/S6_SERIES_PUBLISH_II_A2_WV_original_cycle_chern.py} \\[3pt]
Double-cover root intersections, normalized character algebras and local exceptional matrices (\cref{cdp:st:C.2.D4.2}) & \nolinkurl{code/S6_SERIES_PUBLISH_II_root_BX_prerelease.py} \\[3pt]
Double-cover intersection ranks from Euler strata of the closed cyclic cover (\cref{cdp:st:C.2.D4.2}) & \nolinkurl{code/S6_SERIES_PUBLISH_II_root_BX_postrelease.py} \\[3pt]
Rational conductor configurations and Euler identities (\cref{night:P4.26}) & \nolinkurl{code/S6_SERIES_PUBLISH_II_A2_WY_rational_graphs_and_euler.py} \\[3pt]
Normal-crossing Chern and triple-point identities (\cref{night:P4.26}) & \nolinkurl{code/S6_SERIES_PUBLISH_II_A2_WY_original_chern.py} \\[3pt]
Hopf conductor equality counts (\cref{night:P4.27}) & \nolinkurl{code/S6_SERIES_PUBLISH_II_A2_WZ_equality_character.py} \\[3pt]
Cyclic conductor quotient maps (\cref{night:P4.27}) & \nolinkurl{code/S6_SERIES_PUBLISH_II_A2_WZ_cyclic_quotients.py} \\[3pt]
Quartic Euler characteristic, cyclic projectors and residue valuation (\cref{night:P4.52}) & \nolinkurl{code/S6_SERIES_PUBLISH_II_C2_02XR_final_v1_printed_certificate.py} \\[3pt]
Weighted cubic discriminants, cyclic projectors and rank bounds (\cref{night:P4.53}) & \nolinkurl{code/S6_SERIES_PUBLISH_II_C2_02XS_final_v1_printed_certificate.py} \\[3pt]
Elliptic section rings, weighted cohomology and cyclic characters (\cref{night:P4.54}) & \nolinkurl{code/S6_SERIES_PUBLISH_II_C2_02XT_final_v1_printed_certificate.py} \\[3pt]
Cubic spectral genera and conductor graph bounds (\cref{night:P4.55}) & \nolinkurl{code/S6_SERIES_PUBLISH_II_C2_02XU_final_v1_printed_certificate.py} \\[3pt]
Incidence partitions and the exact quadratic normalization (\cref{night:P4.56}) & \nolinkurl{code/S6_SERIES_PUBLISH_II_C2_02XV_final_v1_printed_certificate.py} \\[3pt]
Cubic Hessians, spectral genus and junction allocations (\cref{night:P4.57}) & \nolinkurl{code/S6_SERIES_PUBLISH_II_C2_02XW_final_v1_printed_certificate.py} \\[3pt]
Rational triple-cover contact partitions, character cohomology and attachment bounds (\cref{prop:II:cubic-root-attachments}) & \nolinkurl{code/S6_SERIES_PUBLISH_II_root_BY_prerelease.py} \\[3pt]
Triple conductor and rational conormal Euler identity (\cref{lem:II:triple-conductor-symbolic-square,lem:II:rational-conormal-defects}) & \nolinkurl{code/S6_SERIES_PUBLISH_II_A2_X7_conormal_gap.py} \\[3pt]
Local conductor defects and gradient minors (\cref{lem:II:rational-conormal-defects}) & \nolinkurl{code/S6_SERIES_PUBLISH_II_A2_X7_original_defects.py} \\[3pt]
Conductor graph colors and Euler reduction (\cref{night:P4.35}) & \nolinkurl{code/S6_SERIES_PUBLISH_II_A2_X8_conductor_graph.py} \\[3pt]
Forest rank and exceptional-curve cohomology counts (\cref{night:P4.35}) & \nolinkurl{code/S6_SERIES_PUBLISH_II_A2_X8_original_forest_count.py} \\[3pt]
Resolved Hopf Chern identities and blowup forms (\cref{night:P4.36}) & \nolinkurl{code/S6_SERIES_PUBLISH_II_A2_X9_foliation_saturation.py} \\[3pt]
Exceptional zero orders under blowups (\cref{night:P4.36}) & \nolinkurl{code/S6_SERIES_PUBLISH_II_A2_X9_original_exceptional_order.py} \\[3pt]
Conductor exponents, quartic Schur complements and cubic exterior weights (\cref{night:P4R.13}) & \nolinkurl{code/S6_SERIES_PUBLISH_II_A2_Y6_primary_jet_comparison.py} \\[3pt]

Excess-one curve rings, their cotangent dimensions and the conductor graph counts (\cref{ii:excess-one-curve-germs,night:P4.40}) & \nolinkurl{code/S6_SERIES_PUBLISH_II_A2_P206_curve_algebra.py}; \nolinkurl{pro_returns/S6_06_BLIND_checkpoint_15_20261117T1457Z/S6_06_BLIND/code/part4/curve_excess.py} \\[3pt]

Integral monodromies, local markings, invariant exterior maps, affine boundary sums and twistor row maps (\cref{ii:three-doubled-nodal-fibres}) & \nolinkurl{code/S6_SERIES_PUBLISH_II_P205_printed_data_check.py} \\[3pt]
Tacnodal conormal inequalities and local Hessians (\cref{night:P4R.10}) & \nolinkurl{pro_returns/S6_06_BLIND_checkpoint_15_20261117T1457Z/S6_06_BLIND/code/part4/tacnode_hessian.py} \\[3pt]
Polar Hessians, exact six-point equations and scaling (\cref{night:P4R.14}) & \nolinkurl{pro_returns/S6_06_BLIND_checkpoint_15_20261117T1457Z/S6_06_BLIND/code/part4/smooth_a3_polar.py} \\[3pt]
Product Hessians and the moving-node family (\cref{night:P4R.15}) & \nolinkurl{pro_returns/S6_06_BLIND_checkpoint_15_20261117T1457Z/S6_06_BLIND/code/part4/hessian_product.py} \\[3pt]
Betti numbers and Euler numbers of two fibre blowups (\cref{cdp:st:C2b}) & \nolinkurl{pro_returns/S6_03_CDP_checkpoint_13_20261116T0136Z/S6_03_CDP/code/C_euler.py} \\[3pt]
Cusp-and-line quadratic matrix and determinant derivatives (\cref{night:P4.38}) & \nolinkurl{pro_returns/S6_06_BLIND_checkpoint_15_20261117T1457Z/S6_06_BLIND/code/part4/critical_curve_node.py} \\[3pt]
Pole-order solutions, normal-degree bounds and homogeneous discriminant (\cref{night:P4.39}) & \nolinkurl{pro_returns/S6_06_BLIND_checkpoint_15_20261117T1457Z/S6_06_BLIND/code/part4/d4_pole.py} \\[3pt]
Normal-crossing branch generators and pinch/cusp module identities (\cref{cdp:st:A.1.2,cdp:st:A.2.4}) & \nolinkurl{pro_returns/S6_03_CDP_checkpoint_13_20261116T0136Z/S6_03_CDP/code/A_local_algebra.py} \\[3pt]
Nodal cubic normalization, descent scalar and cohomology table (\cref{cdp:st:B1}) & \nolinkurl{pro_returns/S6_03_CDP_checkpoint_13_20261116T0136Z/S6_03_CDP/code/B_examples.py} \\[3pt]
Finite character counts for $3\le m\le12$ and all characters (\cref{cdp:st:A.3.2}) & \nolinkurl{pro_returns/S6_03_CDP_checkpoint_13_20261116T0136Z/S6_03_CDP/code/A_multiple_fibres.py} \\[3pt]
Nodal Hessian defects, adjugates and normalization (\cref{night:P4.41}) & \nolinkurl{pro_returns/S6_06_BLIND_checkpoint_15_20261117T1457Z/S6_06_BLIND/code/part4/a3_node.py} \\[3pt]
Uniform transverse-sector Hessians and defect bounds (\cref{night:P4.42,night:P4.44,night:P4.45,night:P4.46}) & \nolinkurl{pro_returns/S6_06_BLIND_checkpoint_15_20261117T1457Z/S6_06_BLIND/code/part4/uniform_sectors.py} \\[3pt]
Conic polar and three-axis conductor matrices (\cref{night:P4.47,night:P4.48,night:P4.49}) & \nolinkurl{pro_returns/S6_06_BLIND_checkpoint_15_20261117T1457Z/S6_06_BLIND/code/part4/a3_three_axes.py} \\[3pt]
Colon ideals, normalization relations and polar Hessians (\cref{night:P4.49}) & \nolinkurl{pro_returns/S6_06_BLIND_checkpoint_15_20261117T1457Z/S6_06_BLIND/code/part4/a3_missing_coefficient.py} \\[3pt]
Tacnode-line conormal elementary divisors (\cref{night:P4R.16}) & \nolinkurl{pro_returns/S6_06_BLIND_checkpoint_15_20261117T1457Z/S6_06_BLIND/code/part4/a3_tacnode_line.py} \\[3pt]
Delta-four curve ideals and conormal lattices (\cref{night:P4R.17}) & \nolinkurl{pro_returns/S6_06_BLIND_checkpoint_15_20261117T1457Z/S6_06_BLIND/code/part4/a3_delta4.py} \\[3pt]
Transverse nodes of a shifted quadratic square (\cref{rem:II:shifted-square-boundary-nodes}) & \nolinkurl{code/S6_SERIES_PUBLISH_II_A2_D09_batch04_boundary_nodes.py} \\[3pt]
Weierstrass resolution charts, Euler number $24$ and product cell ranks (\cref{night:P4.10,night:P4.11}) & \nolinkurl{pro_returns/S6_06_BLIND_checkpoint_15_20261117T1457Z/S6_06_BLIND/code/part4/necessity_examples.py} \\[3pt]
Integral lattice quotients, clutch maps, local cellular differentials and Mayer--Vietoris matrices (\cref{night:P4.3,night:P4.4}) & \nolinkurl{pro_returns/S6_06_BLIND_checkpoint_15_20261117T1457Z/S6_06_BLIND/code/part4/example_mayer_vietoris.py} \\[3pt]
Primitive quotient bases, integral clutchings and Smith forms (\cref{cdp:st:D.2.7}) & \nolinkurl{pro_returns/S6_03_CDP_checkpoint_13_20261116T0136Z/S6_03_CDP/code/D_trivial_monodromy.py} \\[3pt]
Symbolic cubes, cubic discriminant poles and the nodal-image product map (\cref{cdp:st:A.5.19,ex:II:nodal-triple-affine-fibre}) & \nolinkurl{pro_returns/S6_03_CDP_checkpoint_13_20261116T0136Z/S6_03_CDP/code/E_triple_images.py} \\[3pt]
Symbolic-square and symbolic-cube saturations, conormal frames, Hessian and cubic pole orders (\cref{cdp:st:A.5.18,cdp:st:A.5.19,cdp:st:C.2.RAT.1}) & \nolinkurl{pro_returns/S6_03_CDP_checkpoint_13_20261116T0136Z/S6_03_CDP/code/A_low_delta_cusp.py} \\[3pt]
Weighted scaling, simple zeros and the endpoint intersection number for binary products (\cref{cdp:st:A.5.21}) & \nolinkurl{pro_returns/S6_03_CDP_checkpoint_13_20261116T0136Z/S6_03_CDP/code/E_binary_products.py} \\[3pt]
Cubic coefficient degrees, intersection ranks, blowup charts and the scalar-zero coordinate and covering maps (\cref{cdp:st:C.2.D4.2,ex:II:vanishing-cubic-sphere,ex:II:vanishing-cubic-aspherical}) & \nolinkurl{pro_returns/S6_03_CDP_checkpoint_13_20261116T0136Z/S6_03_CDP/code/E_cubic_double.py} \\[3pt]
Return-germ coefficient extractions and integration identities (\cref{cdp:st:C.2.UN.29}) & \nolinkurl{pro_returns/S6_03_CDP_checkpoint_13_20261116T0136Z/S6_03_CDP/code/A_return_germs.py} \\[3pt]
Exterior-square fixed spaces, Gram matrices and signatures of the four unipotent normal forms (\cref{cdp:st:C.2.UN.36}) & \nolinkurl{pro_returns/S6_SERIES_GRIND_checkpoint_10_20261120T0211Z/code/current_audits/audit_sign/code/independent_sign_checks.py} \\[3pt]
Finite cyclic characters, reflection quotients and elliptic-fibre Euler sums (\cref{cdp:st:C.2.UN.33}) & \nolinkurl{pro_returns/S6_03_CDP_checkpoint_13_20261116T0136Z/S6_03_CDP/code/root_cyclic_elliptic_types.py} \\[3pt]
Finite-order determinant and M\"obius sums, cyclic fixed-point and orbit counts, and formal four-curve projection identities (\cref{cdp:st:C.2.UN.11,cdp:st:C.2.UN.24,cdp:st:C.2.UN.39}) & \nolinkurl{pro_returns/S6_SERIES_GRIND_checkpoint_10_20261120T0211Z/code/current_audits/audit_sign/code/UN24_UN39_exact_certificate_08.py} \\[3pt]
Integral meridian and root-cover relation lattices and Smith forms (\cref{cdp:st:C.2.UN.44}) & \nolinkurl{pro_returns/S6_03_CDP_checkpoint_13_20261116T0136Z/S6_03_CDP/code/A_meridian_groups.py} \\[3pt]
Terminal-pair degree equations, blowup Jacobians, normal-extension cocycles, and Euler differences (\cref{cdp:st:C.2.UN.64}) & \nolinkurl{pro_returns/S6_03_CDP_checkpoint_13_20261116T0136Z/S6_03_CDP/code/root_exceptional_leaves.py} \\[3pt]
Ruled degree partitions, elliptic intersections and quotient multiplicities, and blowup Euler sums (\cref{cdp:st:C.2.UN.66}) & \nolinkurl{pro_returns/S6_03_CDP_checkpoint_13_20261116T0136Z/S6_03_CDP/code/A_minimal_reduction.py} \\[3pt]
Rank-four cyclotomic determinant periods, M\"obius signatures, exterior-power zeta identities, and explicit torus automorphism matrices (\cref{cdp:st:C.2.UN.11}) & \nolinkurl{pro_returns/S6_03_CDP_checkpoint_13_20261116T0136Z/S6_03_CDP/code/D_monodromy_signatures.py} \\[3pt]
Root-curve cohomology and sample comb bundle degrees (\cref{inv:E2.2,prop:II:comb-kernel-bounds}) & \nolinkurl{pro_returns/S6_01_INVARIANTS/return/code/d_exact.py} \\[3pt]
Character descent, nonnegative coarse splittings and tangent cohomology for $m=12,13,20,100$ (\cref{cdp:st:E15,inv:E2.2,prop:II:comb-kernel-bounds}) & \nolinkurl{pro_returns/S6_01_INVARIANTS/return/code/d_comb.py} \\[3pt]
Elliptic root coordinates, the integral boundary matrix, sample plurigenera and the Hopf threefold map (\cref{cdp:st:C.2.UN.6,prop:II:multiple-elliptic-hopf-threefold}) & \nolinkurl{pro_returns/S6_03_CDP_checkpoint_13_20261116T0136Z/S6_03_CDP/code/E_elliptic_roots.py} \\[3pt]
Cuspidal self-product charts, the divisor intersection matrix and cohomology arithmetic (\cref{cdp:st:C.2.UN.8,rem:II:cuspidal-generic-divisors}) & \nolinkurl{pro_returns/S6_03_CDP_checkpoint_13_20261116T0136Z/S6_03_CDP/code/root_cusp_product.py} \\[3pt]
Kummer quotient chart cocycles, quadratic generators, Gysin matrices and Euler arithmetic (\cref{cdp:st:C.2.UN.13}) & \nolinkurl{pro_returns/S6_SERIES_GRIND_checkpoint_10_20261120T0211Z/code/current_audits/audit_sign/code/UN13_current_exact_certificate_08.py} \\[3pt]
Enriques filling equations, involution, quotient multiplicities and Euler arithmetic (\cref{cdp:st:C.2.UN.25}) & \nolinkurl{pro_returns/S6_03_CDP_checkpoint_13_20261116T0136Z/S6_03_CDP/code/E_enriques_filling.py} \\[3pt]
Enriques support incidence and curve-class matrices, including the paired-root family (\cref{cdp:st:C.2.UN.25,cdp:st:C.2.UN.27}) & \nolinkurl{pro_returns/S6_03_CDP_checkpoint_13_20261116T0136Z/S6_03_CDP/code/A_enriques_support.py} \\[3pt]
Cyclic integral norm identities and compatible elliptic and kernel characters (\cref{lem:II:equivariant-elliptic-torsor-lift,rem:II:elliptic-torus-characters,cdp:st:C.2.UN.35}) & \nolinkurl{pro_returns/S6_03_CDP_checkpoint_13_20261116T0136Z/S6_03_CDP/code/A_equivariant_torus.py} \\[3pt]
Sample cusp blowup sequences, genus formulas and ruled-support Euler arithmetic (\cref{cdp:st:C.2.UN.45}) & \nolinkurl{pro_returns/S6_03_CDP_checkpoint_13_20261116T0136Z/S6_03_CDP/code/root_ruled_occurrence.py} \\[3pt]
Mixed-cycle edge frames, primitive normal periods and the horizontal translation (\cref{ex:II:three-hopf-cycle,cdp:st:C.2.UN.57}) & \nolinkurl{pro_returns/S6_03_CDP_checkpoint_13_20261116T0136Z/S6_03_CDP/code/A_mixed_cycles.py} \\[3pt]
Rooted restriction matrices, the integral Chern extension and analytic torsor periods (\cref{cdp:st:C.2.UN.61}) & \nolinkurl{pro_returns/S6_03_CDP_checkpoint_13_20261116T0136Z/S6_03_CDP/code/A_rooted_hopf_trees.py} \\[3pt]
Primitive core periods and nested insertion orders (\cref{cdp:st:C.2.UN.61}) & \nolinkurl{pro_returns/S6_03_CDP_checkpoint_13_20261116T0136Z/S6_03_CDP/code/A_tree_periods.py} \\[3pt]
The real-linear period map, its inverse, determinant and frame at infinity (\cref{lem:II:logarithmic-period-family}) & \nolinkurl{pro_returns/S6_03_CDP_checkpoint_13_20261116T0136Z/S6_03_CDP/code/A_twistor_log_cycles.py} \\[3pt]
Weighted periods, finite actions and conjugate frames (\cref{cdp:st:C.2.UN.59}) & \nolinkurl{pro_returns/S6_03_CDP_checkpoint_13_20261116T0136Z/S6_03_CDP/code/D_all_weighted_cycles.py} \\[3pt]
Hopf-cycle matrices, affine translations and finite-secondary period divisors (\cref{ex:II:three-hopf-cycle,cdp:st:C.2.UN.48}) & \nolinkurl{pro_returns/S6_03_CDP_checkpoint_13_20261116T0136Z/S6_03_CDP/code/D_hopf_cycle_matrices.py} \\[3pt]
Inverse root charts and quasibundle plurigenera (\cref{prop:II:elliptic-quasibundle-plurigenera}) & \nolinkurl{pro_returns/S6_03_CDP_checkpoint_13_20261116T0136Z/S6_03_CDP/code/D_quasibundle_trees.py} \\[3pt]
Root multiplicities, inertia groups, resolution charts and finite quotient actions (\cref{cdp:st:C.2.UN.52}) & \nolinkurl{pro_returns/S6_03_CDP_checkpoint_13_20261116T0136Z/S6_03_CDP/code/D_weighted_hopf_cycles.py} \\[3pt]
Weighted restriction matrices and their dual residue vectors (\cref{cdp:st:C.2.UN.54}) & \nolinkurl{pro_returns/S6_03_CDP_checkpoint_13_20261116T0136Z/S6_03_CDP/code/E_hopf_graph_reduction.py} \\[3pt]
Canonical heights, root congruences, finite-orbit gauges and collar maps (\cref{cdp:st:C.2.UN.58}) & \nolinkurl{pro_returns/S6_03_CDP_checkpoint_13_20261116T0136Z/S6_03_CDP/code/E_hopf_tree_heights.py} \\[3pt]
Primitive fibre bases, two-mark Picard lattices and recursive insertion identities (\cref{cdp:st:C.2.UN.60}) & \nolinkurl{pro_returns/S6_03_CDP_checkpoint_13_20261116T0136Z/S6_03_CDP/code/E_hopf_tree_pruning.py} \\[3pt]
Period determinants, conjugate Hopf frames and finite angular characters (\cref{lem:II:logarithmic-period-family,cdp:st:C.2.UN.57}) & \nolinkurl{pro_returns/S6_03_CDP_checkpoint_13_20261116T0136Z/S6_03_CDP/code/root_all_diagonal_cycles.py} \\[3pt]
Hopf-star degree equations, weighted maps, canonical forms and core periods (\cref{cdp:st:C.2.UN.53}) & \nolinkurl{pro_returns/S6_03_CDP_checkpoint_13_20261116T0136Z/S6_03_CDP/code/root_mixed_hopf_stars.py} \\[3pt]
\end{longtable}}

\begin{thebibliography}{99}
\providecommand{\natexlab}[1]{#1}
\providecommand{\doi}[1]{\url{https://doi.org/#1}}
\bibitem[Abasheva and D{\'e}ev(2024)]{AbashevaDeev2024AB}
Anna Abasheva and Rodion D{\'e}ev.

\newblock Complex surfaces with many algebraic structures.

\newblock \emph{International Mathematics Research Notices}, 2024\penalty0
  (9):\penalty0 7379--7400, 2024.

\newblock \doi{10.1093/imrn/rnad190}.

\newblock Cited preprint edition: arXiv:2303.10764v1,
19 March 2023. Page locators refer to this edition.
\newblock URL \url{https://arxiv.org/pdf/2303.10764v1}.

\bibitem[Abramovich et~al.(2002)Abramovich, Karu, Matsuki, and
  W{\l}odarczyk]{AKMW2002}
Dan Abramovich, Kalle Karu, Kenji Matsuki, and Jaros{\l}aw W{\l}odarczyk.

\newblock Torification and factorization of birational maps.

\newblock \emph{Journal of the American Mathematical Society}, 15\penalty0
  (3):\penalty0 531--572, 2002.

\newblock \doi{10.1090/S0894-0347-02-00396-X}.

\newblock Cited preprint edition: arXiv:math/9904135v4,
31 May 2000. Page locators refer to this edition.
\newblock URL \url{https://arxiv.org/pdf/math/9904135v4}.

\bibitem[A'Campo(1975)]{ACampo1975}
Norbert A'Campo.

\newblock La fonction z\^eta d'une monodromie.

\newblock \emph{Commentarii Mathematici Helvetici}, 50:\penalty0 233--248,
  1975.

\newblock \doi{10.1007/BF02565748}.

\newblock URL
  \url{https://www.e-periodica.ch/cntmng?bot=1&pid=com-001%3A1975%3A50%3A%3A27}.

\bibitem[Albanese(2021)]{Albanese2020}
Michael Albanese.

\newblock The {Yamabe} invariants of {Inoue} surfaces, {Kodaira} surfaces, and
  their blowups.

\newblock \emph{Annals of Global Analysis and Geometry}, 59\penalty0
  (2):\penalty0 179--195, 2021.
\newblock \doi{10.1007/s10455-020-09744-3}.
\newblock Citations refer to arXiv:2005.09494v2, 6 November 2020.
\newblock URL \url{https://arxiv.org/pdf/2005.09494v2}.


\bibitem[Arnold(1976)]{Arnold1976NormalForms}
Vladimir I. Arnold.

\newblock Local normal forms of functions.

\newblock \emph{Inventiones Mathematicae}, 35:\penalty0 87--109, 1976.

\newblock \doi{10.1007/BF01390134}.

\bibitem[Artebani and Dolgachev(2009)]{ArtebaniDolgachev2009}
Michela Artebani and Igor~V. Dolgachev.
\newblock The Hesse pencil of plane cubic curves.
\newblock \emph{L'Enseignement Math\'ematique}, 55(3/4):235--273, 2009.
\newblock \doi{10.4171/LEM/55-3-3}.
\newblock URL \url{https://ems.press/content/serial-article-files/44190?nt=1}.

\bibitem[Atkin and Garvan(2002)]{AtkinGarvan2002}
A.~O.~L. Atkin and F.~G. Garvan.
\newblock Relations between the ranks and cranks of partitions.
\newblock arXiv:math/0208050v1, 6 August 2002.
\newblock Published in \emph{The Ramanujan Journal}, 7(1--3):343--366, 2003.
\newblock \doi{10.1023/A:1026219901284}.
\newblock The locators used here refer to the arXiv version.
\newblock URL \url{https://arxiv.org/pdf/math/0208050v1}.

\bibitem[B{\u a}nic{\u a} and St{\u a}n{\u a}{\c s}il{\u a}(1976)]{BS76push}
Constantin B{\u a}nic{\u a} and Octavian St{\u a}n{\u a}{\c s}il{\u a}.

\newblock \emph{Algebraic Methods in the Global Theory of Complex Spaces}.

\newblock Editura Academiei and John Wiley and Sons, 1976.


\bibitem[Barth et~al.(2004)Barth, Hulek, Peters, and {Van de Ven}]{BHPV2004E}
Wolf~P. Barth, Klaus Hulek, Chris A.~M. Peters, and Antonius {Van de Ven}.

\newblock \emph{Compact Complex Surfaces}, volume~4 of \emph{Ergebnisse der
  Mathematik und ihrer Grenzgebiete. 3. Folge}.

\newblock Springer, Berlin, Heidelberg, second edition, 2004.

\newblock \doi{10.1007/978-3-642-57739-0}.

\newblock URL \url{https://link.springer.com/book/10.1007/978-3-642-57739-0}.

\bibitem[Beauville(1978)]{Beauville1978}
Arnaud Beauville.
\newblock \emph{Surfaces alg\'ebriques complexes}.
\newblock Ast\'erisque, volume 54. Soci\'et\'e math\'ematique de France, 1978.
\newblock URL \url{https://www.numdam.org/item/AST_1978__54__R1_0/}.

\bibitem[Bierstone and Milman(1988)]{BierstoneMilman1988}
Edward Bierstone and Pierre~D. Milman.

\newblock Semianalytic and subanalytic sets.

\newblock \emph{Publications Math\'ematiques de l'IH\'ES},
  67:\penalty0 5--42, 1988.

\newblock \doi{10.1007/BF02699126}.

\newblock URL \url{https://www.numdam.org/item/PMIHES_1988__67__5_0.pdf}.


\bibitem[Bierstone and Milman(1997)]{BierstoneMilman1997}
Edward Bierstone and Pierre~D. Milman.

\newblock Canonical desingularization in characteristic zero by blowing up the
  maximum strata of a local invariant.

\newblock \emph{Inventiones Mathematicae}, 128:\penalty0 207--302, 1997.

\newblock \doi{10.1007/s002220050141}.

\newblock URL \url{https://arxiv.org/pdf/alg-geom/9508005v2}.
\newblock Cited preprint edition: arXiv:alg-geom/9508005v2,
16 August 1995; preprint page locators refer to this edition.

\bibitem[Biswas et~al.(2012)Biswas, Mj, and Seshadri]{BMS2012}
Indranil Biswas, Mahan Mj, and Harish Seshadri.

\newblock 3-manifold groups, {K\"ahler} groups and complex surfaces.

\newblock \emph{Communications in Contemporary Mathematics}, 14\penalty0
  (6):\penalty0 1250038, 2012.

\newblock \doi{10.1142/S0219199712500381}.

\newblock URL \url{https://mathweb.tifr.res.in/~mahan/kahler3fld.pdf}.
\newblock Cited author preprint: 27 October 2011; preprint page and statement
locators refer to the linked author edition.

\bibitem[Borcea(1984)]{Borcea1984}
Ciprian Borcea.

\newblock Moduli for {Kodaira} surfaces.

\newblock \emph{Compositio Mathematica}, 52\penalty0 (3):\penalty0 373--380,
  1984.

\newblock URL \url{https://www.numdam.org/article/CM_1984__52_3_373_0.pdf}.


\bibitem[Br{\^i}nz{\u a}nescu and Moraru(2005)]{BrinzanescuMoraru2005}
Vasile Br{\^i}nz{\u a}nescu and Ruxandra Moraru.

\newblock Holomorphic rank-2 vector bundles on non-{K\"ahler} elliptic
  surfaces.

\newblock \emph{Annales de l'Institut Fourier}, 55\penalty0 (5):\penalty0
  1659--1683, 2005.

\newblock \doi{10.5802/aif.2135}.

\newblock URL \url{https://aif.centre-mersenne.org/item/AIF_2005__55_5_1659_0.pdf}.

\bibitem[Brown(1982)]{Brown1982}
Kenneth~S. Brown.

\newblock \emph{Cohomology of Groups}, volume~87 of \emph{Graduate Texts in
  Mathematics}.

\newblock Springer, New York, 1982.

\newblock \doi{10.1007/978-1-4684-9327-6}.

\bibitem[Buchdahl(1999)]{Buchdahl1999}
Nicholas Buchdahl.

\newblock On compact {K\"ahler} surfaces.

\newblock \emph{Annales de l'Institut Fourier}, 49\penalty0 (1):\penalty0
  287--302, 1999.

\newblock \doi{10.5802/aif.1674}.

\newblock URL \url{https://www.numdam.org/article/AIF_1999__49_1_287_0.pdf}.

\bibitem[Buonerba et~al.(2019)Buonerba, Bogomolov, and Kurnosov]{BBK2019}
Federico Buonerba, Fedor Bogomolov, and Nikon Kurnosov.
\newblock Classifying {VII}$_0$ surfaces with $b_2=0$ via group theory.
\newblock arXiv:1709.00062v3, 20 May 2019.
\newblock URL \url{https://arxiv.org/pdf/1709.00062v3}.


\bibitem[Campana et~al.(1998)Campana, Demailly, and Peternell]{CDP98}
Fr{\'e}d{\'e}ric Campana, Jean-Pierre Demailly, and Thomas Peternell.

\newblock The algebraic dimension of compact complex threefolds with vanishing
  second Betti number.

\newblock \emph{Compositio Mathematica}, 112\penalty0 (1):\penalty0 77--91,
  1998.

\newblock \doi{10.1023/A:1000313214795}.

\bibitem[Campana et~al.(2020)Campana, Demailly, and Peternell]{CDP2020}
Fr{\'e}d{\'e}ric Campana, Jean-Pierre Demailly, and Thomas Peternell.
\newblock The algebraic dimension of compact complex threefolds with vanishing
second Betti number.
\newblock \emph{Compositio Mathematica}, 156(4):679--696, 2020.
\newblock \doi{10.1112/S0010437X19007802}.

\bibitem[Cartan(1957)]{Cartan1957}
Henri Cartan.

\newblock Vari\'et\'es analytiques r\'eelles et vari\'et\'es analytiques
  complexes.

\newblock \emph{Bulletin de la Soci\'et\'e Math\'ematique de France},
  85:\penalty0 77--99, 1957.

\newblock \doi{10.24033/bsmf.1481}.

\newblock URL \url{https://www.numdam.org/article/BSMF_1957__85__77_0.pdf}.

\bibitem[Catanese(2002)]{Catanese2002}
Fabrizio Catanese.

\newblock Deformation types of real and complex manifolds.

\newblock In \emph{Contemporary Trends in Algebraic Geometry and Algebraic
  Topology}, volume~5 of \emph{Nankai Tracts in Mathematics}, pages 195--238.
  World Scientific, River Edge, NJ, 2002.
\newblock \doi{10.1142/9789812777416_0009}.
\newblock Citations refer to arXiv:math/0111245v2, 30 January 2002.
\newblock URL \url{https://arxiv.org/pdf/math/0111245v2}.


\bibitem[Catanese(2004)]{Catanese2004}
Fabrizio Catanese.

\newblock Deformation in the large of some complex manifolds, {I}.

\newblock \emph{Annali di Matematica Pura ed Applicata}, 183:\penalty0
  261--289, 2004.

\newblock \doi{10.1007/s10231-003-0096-y}.

\newblock URL
  \url{https://www.mathe8.uni-bayreuth.de/de/team/prof-fabrizio-catanese-old/pdf/085.pdf}.

\bibitem[{Claude} and {L. Alp{\"o}ge}(2026)]{AC26}
{Claude} and {L. Alp{\"o}ge}.
\newblock The $(3,4,\infty)$ modular family of $2$-tori, completed at its three
special points, is a complex structure on $S^6$.
\newblock Manuscript, 108 pages, 2026.
\newblock \url{https://alpo.ge/s6.pdf}.

\bibitem[Col\c{t}oiu et~al.(2012)Col\c{t}oiu, Ga\c{s}i\c{t}oi, and Joi\c{t}a]{ColtoiuGasitoiJoita2012}
Mihnea Col\c{t}oiu, Natalia Ga\c{s}i\c{t}oi, and Cezar Joi\c{t}a.
\newblock On the image of an algebraic projective space.
\newblock \emph{Comptes Rendus Math\'ematique}, 350(5--6):\penalty0 239--241, 2012.
\newblock \doi{10.1016/j.crma.2012.02.002}.
\newblock URL \url{https://www.numdam.org/item/10.1016/j.crma.2012.02.002.pdf}.

\bibitem[Cox et~al.(2011)Cox, Little, and Schenck]{CoxLittleSchenck2011}
David~A. Cox, John~B. Little, and Hal Schenck.

\newblock \emph{Toric Varieties}, volume 124 of \emph{Graduate Studies in
  Mathematics}.

\newblock American Mathematical Society, Providence, RI, 2011.

\newblock \doi{10.1090/gsm/124}.
\newblock Page locators refer to the authors' 2010 text, whose preface is
  dated September 2010.
\newblock URL \url{https://dacox.people.amherst.edu/toric.html}.

\bibitem[de~Jong and van Straten(1990)]{DeJongVanStraten1990Normalization}
Theo de~Jong and Duco van Straten.

\newblock Deformations of the normalization of hypersurfaces.

\newblock \emph{Mathematische Annalen}, 288:\penalty0 527--547, 1990.

\newblock \doi{10.1007/BF01444547}.

\newblock URL
  \url{https://download.uni-mainz.de/mathematik/Algebraische%20Geometrie/DvS-Publikationen/Deformations%20of%20the%20normalization%20of%20hypersurfaces%20-%201990.pdf}.

\bibitem[Demailly(2002)]{DemaillyHodgeAuthor2002}
Jean-Pierre Demailly.

\newblock $L^2$ Hodge theory and vanishing theorems.

\newblock In \emph{Introduction to Hodge Theory}, volume~8 of \emph{SMF/AMS
  Texts and Monographs}. American Mathematical Society and Soci\'et\'e
  Math\'ematique de France, 2002.

\newblock URL
  \url{https://www-fourier.univ-grenoble-alpes.fr/~demailly/manuscripts/hodge-ams.pdf}.

\newblock Author's version; cited page numbers refer to the chapter
  pagination of this PDF.

\bibitem[Demailly(2012)]{Demailly2012E}
Jean-Pierre Demailly.

\newblock Complex analytic and differential geometry, 2012.

\newblock URL
  \url{https://www-fourier.univ-grenoble-alpes.fr/~demailly/manuscripts/agbook.pdf}.

\bibitem[Dloussky(2006)]{Dloussky2006}
Georges Dloussky.

\newblock On surfaces of class {$\mathrm{VII}_0^+$} with numerically
  anticanonical divisor.

\newblock \emph{American Journal of Mathematics}, 128\penalty0 (3):\penalty0
  639--670, 2006.

\newblock \doi{10.1353/ajm.2006.0021}.

\newblock Cited version: arXiv:math/0406387v3, 13 April 2005.
\newblock URL \url{https://arxiv.org/abs/math/0406387v3}.

\bibitem[Dloussky and Oeljeklaus(1999)]{DlousskyOeljeklaus1999}
Georges Dloussky and Karl Oeljeklaus.
\newblock Vector fields and foliations on compact surfaces of class {$\mathrm{VII}_0$}.
\newblock \emph{Annales de l'Institut Fourier}, 49(5):1503--1545, 1999.
\newblock \doi{10.5802/aif.1728}.
\newblock URL \url{https://www.numdam.org/item/10.5802/aif.1728.pdf}.

\bibitem[Dloussky et~al.(2000)Dloussky, Oeljeklaus, and Toma]{DOT2000}
Georges Dloussky, Karl Oeljeklaus, and Matei Toma.

\newblock Surfaces de la classe {VII}$_0$ admettant un champ de vecteurs.

\newblock \emph{Commentarii Mathematici Helvetici}, 75\penalty0 (2):\penalty0
  255--270, 2000.

\newblock \doi{10.1007/PL00000374}.

\newblock URL \url{https://ems.press/content/serial-article-files/42865?nt=1}.

\bibitem[Dloussky et~al.(2001)Dloussky, Oeljeklaus, and Toma]{DOT2001}
Georges Dloussky, Karl Oeljeklaus, and Matei Toma.

\newblock Surfaces de la classe {VII}$_0$ admettant un champ de vecteurs, {II}.

\newblock \emph{Commentarii Mathematici Helvetici}, 76\penalty0 (4):\penalty0
  640--664, 2001.

\newblock \doi{10.1007/s00014-001-8323-x}.

\newblock URL \url{https://ems.press/content/serial-article-files/42912?nt=1}.

\bibitem[Dloussky et~al.(2003)Dloussky, Oeljeklaus, and Toma]{DOT2003}
Georges Dloussky, Karl Oeljeklaus, and Matei Toma.

\newblock Class {$\mathrm{VII}_0$} surfaces with {$b_2$} curves.

\newblock \emph{Tohoku Mathematical Journal, Second Series}, 55\penalty0
  (2):\penalty0 283--309, 2003.

\newblock \doi{10.2748/tmj/1113246942}.

\newblock URL
  \url{https://www.jstage.jst.go.jp/article/tmj/55/2/55_2_283/_pdf/-char/en}.

\bibitem[Dloussky and Teleman(2012)]{DlousskyTeleman2012}
Georges Dloussky and Andrei Teleman.

\newblock Infinite bubbling in non-{K\"ahlerian} geometry.

\newblock \emph{Mathematische Annalen}, 353:\penalty0 1283--1314, 2012.

\newblock \doi{10.1007/s00208-011-0713-9}.

\newblock Cited preprint edition: arXiv:1012.5247v2, 5 October 2011;
page and remark locators refer to this edition.
\newblock URL \url{https://arxiv.org/pdf/1012.5247v2}.

\bibitem[Dolgachev(1982)]{Dolgachev1982}
Igor Dolgachev.
\newblock Weighted projective varieties.
\newblock In James~B. Carrell, editor, \emph{Group Actions and Vector Fields},
Lecture Notes in Mathematics, volume~956, pages~34--71. Springer, Berlin, 1982.
\newblock \doi{10.1007/BFb0101508}.
\newblock URL \url{https://sites.lsa.umich.edu/idolga-new/wp-content/uploads/sites/1467/2024/08/weighted82.pdf}.

\bibitem[Douady(1966)]{Douady1966}
Adrien Douady.

\newblock Le probl\`eme des modules pour les sous-espaces analytiques compacts
  d'un espace analytique donn\'e.

\newblock \emph{Annales de l'Institut Fourier}, 16\penalty0 (1):\penalty0
  1--95, 1966.

\newblock \doi{10.5802/aif.226}.

\newblock URL \url{https://www.numdam.org/item/AIF_1966__16_1_1_0/}.

\bibitem[Du et~al.(2026)Du, Guo, Kusakabe, and Xie]{DGKXv1}
Yun-Heng Du, Bin Guo, Yuta Kusakabe, and Song-Yan Xie.

\newblock Oka geometry and deformations of torus-fibred complex six-spheres,
  2026.

\newblock URL \url{https://arxiv.org/abs/2610.00250v1}.

\bibitem[Ehresmann(1952)]{Ehresmann1952Connections}
Charles Ehresmann.
\newblock Les connexions infinit\'esimales dans un espace fibr\'e diff\'erentiable.
\newblock \emph{S\'eminaire Bourbaki}, 1, expos\'e~24, 153--168, 1952.
\newblock URL \url{https://www.numdam.org/article/SB_1948-1951__1__153_0.pdf}.

\bibitem[Engel(2026)]{Engel2026}
Philip Engel.

\newblock Complex structures on {$S^6$}.

\newblock arXiv:2609.38442v1, 2026.

\newblock URL \url{https://arxiv.org/abs/2609.38442v1}.

\bibitem[Favre(2000)]{FavreThesis2002}
Charles Favre.
\newblock \emph{Dynamique des applications rationnelles}.
\newblock Ph.D. thesis, 2000; author-hosted version dated 22 July 2002.
\newblock URL \url{https://www.cmls.polytechnique.fr/perso/favre.charles/these.pdf}.

\bibitem[Fern\'andez~de Bobadilla and Koll\'ar(2012)]{BobadillaKollar2012}
Javier Fern\'andez~de Bobadilla and J\'anos Koll\'ar.

\newblock Homotopically trivial deformations.

\newblock \emph{Journal of Singularities}, 5:\penalty0 85--93, 2012.

\newblock \doi{10.5427/jsing.2012.5f}.

\newblock URL \url{https://www.journalofsing.org/volume5/bobadilla-kollar.pdf}.

\bibitem[Fujiki and Nakano(1971)]{FujikiNakano1971}
Akira Fujiki and Shigeo Nakano.

\newblock Supplement to {On the Inverse of Monoidal Transformation}.

\newblock \emph{Publications of the Research Institute for Mathematical
  Sciences}, 7\penalty0 (3):\penalty0 637--644, 1971.

\newblock \doi{10.2977/PRIMS/1195193401}.

\newblock URL \url{https://ems.press/content/serial-article-files/41704}.

\bibitem[Fujimoto(1988)]{Fujimoto1988Log}
Yoshio Fujimoto.

\newblock Logarithmic transformations on elliptic fiber spaces.

\newblock \emph{Journal of Mathematics of Kyoto University}, 28\penalty0
  (1):\penalty0 91--110, 1988.

\newblock \doi{10.1215/kjm/1250520560}.

\bibitem[Fujino(2021)]{Fujino2021E}
Osamu Fujino.

\newblock Minimal model theory for log surfaces in {Fujiki}'s class {$\mathcal
  C$}.

\newblock \emph{Nagoya Mathematical Journal}, 244:\penalty0 256--282, 2021.

\newblock \doi{10.1017/nmj.2020.14}.

\newblock Cited author edition: version 0.08,
20 January 2020. Page locators refer to this edition.
\newblock URL \url{https://www.math.kyoto-u.ac.jp/~fujino/surfaces-classC2.pdf}.

\bibitem[Furushima(1993)]{Furushima1993}
Mikio Furushima.

\newblock The complete classification of compactifications of {$\mathbb C^3$}
  which are projective manifolds with the second betti number one.

\newblock Technical Report MPI/93-6, Max-Planck-Institut f\"ur Mathematik,
  1993.

\newblock URL
  \url{https://archive.mpim-bonn.mpg.de/id/eprint/2170/1/preprint_1993_6.pdf}.

\bibitem[Gim\'enez Conejero and Pint\'er(2024)]{GCP2024}
Roberto Gim\'enez Conejero and Gerg\H{o} Pint\'er.
\newblock Signature of the Milnor fiber of parametrized surfaces.
\newblock arXiv:2401.14951v2, 30 January 2024.
\newblock Published in \emph{Journal of the London Mathematical Society},
114(3):e70699, 2026. \doi{10.1112/jlms.70699}.
\newblock The locators used here refer to the arXiv version.
\newblock URL \url{https://arxiv.org/abs/2401.14951v2}.

\bibitem[Granger and Schulze(2014)]{GS2014}
Michel Granger and Mathias Schulze.

\newblock Normal crossing properties of complex hypersurfaces via logarithmic
  residues.

\newblock \emph{Compositio Mathematica}, 150\penalty0 (9):\penalty0 1607--1622,
  2014.

\newblock \doi{10.1112/S0010437X13007860}.

\newblock arXiv:1109.2612v7, 22 May 2014.
\newblock URL \url{https://arxiv.org/pdf/1109.2612v7}.

\bibitem[Grauert(1960)]{Grauert1960D}
Hans Grauert.
\newblock Ein Theorem der analytischen Garbentheorie und die Modulr\"aume
komplexer Strukturen.
\newblock \emph{Publications Math\'ematiques de l'IH\'ES}, 5:5--64, 1960.
\newblock \doi{10.1007/BF02684746}.

\bibitem[Grauert(1962)]{Grauert1962E}
Hans Grauert.

\newblock {\"U}ber Modifikationen und exzeptionelle analytische Mengen.

\newblock \emph{Mathematische Annalen}, 146:\penalty0 331--368, 1962.

\newblock \doi{10.1007/BF01441136}.

\newblock URL \url{https://link.springer.com/article/10.1007/BF01441136}.

\bibitem[Grauert(1963)]{Grauert1963ErratumD}
Hans Grauert.
\newblock Berichtigung: Ein Theorem der analytischen Garbentheorie und die
  Modulr\"aume komplexer Strukturen.
\newblock \emph{Publications Math\'ematiques de l'IH\'ES}, 16:35--36, 1963.
\newblock \doi{10.1007/BF02684291}.

\bibitem[Grauert and Riemenschneider(1970)]{GrauertRiemenschneider1970}
Hans Grauert and Oswald Riemenschneider.

\newblock Verschwindungss\"atze f\"ur analytische Kohomologiegruppen
  auf komplexen R\"aumen.

\newblock \emph{Inventiones Mathematicae}, 11\penalty0 (4):\penalty0
  263--292, 1970.
\newblock \doi{10.1007/BF01403182}.
\newblock URL
  \url{https://www.math.uni-hamburg.de/home/riemenschneider/OR1970c.pdf}.

\bibitem[Greb(2011)]{Greb2011}
Daniel Greb.

\newblock Rational singularities and quotients by holomorphic group actions.

\newblock \emph{Annali della Scuola Normale Superiore di Pisa, Classe di
  Scienze}, ser.~5, 10\penalty0 (2):\penalty0 413--426, 2011.
\newblock \doi{10.2422/2036-2145.2011.2.07}.
\newblock URL \url{https://www.numdam.org/item/ASNSP_2011_5_10_2_413_0.pdf}.


\bibitem[Gross et~al.(2015)Gross, Hacking, and Keel]{GHK2015}
Mark Gross, Paul Hacking, and Sean Keel.

\newblock Mirror symmetry for log {Calabi--Yau} surfaces {I}.

\newblock \emph{Publications Math\'ematiques de l'IH\'ES}, 122:\penalty0
  65--168, 2015.

\newblock \doi{10.1007/s10240-015-0073-1}.

\newblock URL \url{https://www.numdam.org/articles/10.1007/s10240-015-0073-1/}.

\bibitem[Gross and Siebert(2006)]{GrossSiebert2006I}
Mark Gross and Bernd Siebert.

\newblock Mirror symmetry via logarithmic degeneration data {I}.

\newblock \emph{Journal of Differential Geometry}, 72\penalty0 (2):\penalty0
  169--338, 2006.

\newblock \doi{10.4310/jdg/1143593211}.

\newblock Cited version: \emph{arXiv:math/0309070v3}, 20 July 2005; page locators refer to the printed pagination of that PDF.

\newblock URL \url{https://arxiv.org/abs/math/0309070v3}.

\bibitem[Guillot(2019)]{Guillot2019}
Adolfo Guillot.

\newblock Meromorphic vector fields with single-valued solutions on complex
  surfaces.

\newblock \emph{Advances in Mathematics}, 354:\penalty0 106742, 2019.

\newblock \doi{10.1016/j.aim.2019.106742}.

\newblock arXiv:1603.02288v2, 5 August 2019.
\newblock URL \url{https://arxiv.org/pdf/1603.02288v2}.

\bibitem[Guillot(2024)]{Guillot2024}
Adolfo Guillot.
\newblock On the singularities of complete holomorphic vector fields in
dimension two.
\newblock In F.~Cano, J.~L.~Cisneros-Molina, L\^e D\~ung Tr\'ang,
and J.~Seade, editors, \emph{Handbook of Geometry and Topology of
Singularities VI: Foliations}, pages 1--37. Springer, Cham, 2024.
\newblock \doi{10.1007/978-3-031-54172-8_1}.
\newblock Author preprint, 32 pages; theorem numbers and page citations
refer to this preprint.
\newblock URL
\url{https://www.matem.unam.mx/~guillot/guillot-singularities.pdf}.

\bibitem[Hassett(2009)]{Hassett2009Rational}
B.~Hassett, \emph{Rational surfaces over nonclosed fields},
in \emph{Arithmetic Geometry}, H.~Darmon, D.~Ellwood, B.~Hassett and Y.~Tschinkel (eds.),
Clay Mathematics Proceedings \textbf{8}, American Mathematical Society and Clay Mathematics Institute, 2009, pp.~155--209.
\url{https://www.claymath.org/wp-content/uploads/2022/03/cmip08c.pdf}.

\bibitem[Hatcher(2002)]{Hatcher2002}
Allen Hatcher.

\newblock \emph{Algebraic Topology}.

\newblock Cambridge University Press, 2002.

\newblock URL \url{https://pi.math.cornell.edu/~hatcher/AT/AT.pdf}.

\bibitem[Honda and Viaclovsky(2025)]{HondaViaclovsky2025D}
Nobuhiro Honda and Jeff Viaclovsky.

\newblock Some remarks on fibrations in complex geometry.

\newblock \emph{Journal of Mathematical Study}, 58\penalty0 (4):\penalty0
  552--574, 2025.

\newblock \doi{10.4208/jms.v58n4.25.11}.

\newblock Cited preprint edition: arXiv:2504.14420v1,
19 April 2025. Page locators refer to this edition.
\newblock URL \url{https://arxiv.org/pdf/2504.14420v1}.

\bibitem[Horikawa(1974)]{Horikawa1974}
Eiji Horikawa.

\newblock On deformations of holomorphic maps. {II}.

\newblock \emph{Journal of the Mathematical Society of Japan}, 26\penalty0
  (4):\penalty0 647--667, 1974.

\newblock \doi{10.2969/jmsj/02640647}.

\newblock URL
  \url{https://www.jstage.jst.go.jp/article/jmath1948/26/4/26_4_647/_pdf/-char/en}.

\bibitem[Huckleberry et~al.(2000)Huckleberry, Kebekus, and Peternell]{HKP2000}
Alan~T. Huckleberry, Stefan Kebekus, and Thomas Peternell.

\newblock Group actions on {$S^6$} and complex structures on {$\mathbb P^3$}.

\newblock \emph{Duke Mathematical Journal}, 102\penalty0 (1):\penalty0
  101--124, 2000.

\newblock \doi{10.1215/S0012-7094-00-10214-1}.

\newblock URL \url{https://arxiv.org/abs/math/9812076v1}.
\newblock Cited preprint edition: arXiv:math/9812076v1, 13 December 1998;
preprint page locators refer to this edition.

\bibitem[Huybrechts(2016)]{HuybrechtsK3Draft}
Daniel Huybrechts.

\newblock \emph{Lectures on {K3} Surfaces}.

\newblock Cambridge University Press, 2016.

\newblock The cited page numbers refer to the author's 449-page
  pre-publication draft (copyright 2015).

\newblock URL \url{https://www.math.uni-bonn.de/people/huybrech/K3Global.pdf}.

\bibitem[Inoue(1974)]{Inoue1974}
Masahisa Inoue.

\newblock On surfaces of class {$VII_0$}.

\newblock \emph{Inventiones Mathematicae}, 24\penalty0 (4):\penalty0 269--310,
  1974.

\newblock \doi{10.1007/BF01425563}.



\bibitem[Kajiwara et~al.(2008)Kajiwara, Kato, and Nakayama]{KKN2008}
Takeshi Kajiwara, Kazuya Kato, and Chikara Nakayama.

\newblock Logarithmic abelian varieties, part {I}: Complex analytic theory.

\newblock \emph{Journal of Mathematical Sciences, The University of Tokyo},
  15:\penalty0 69--193, 2008.

\newblock URL \url{https://www.ms.u-tokyo.ac.jp/journal/pdf/jms150104.pdf}.

\bibitem[Kamenova(n.d.)]{KamenovaTori}
Ljudmila Kamenova.

\newblock Degenerations of 2-dimensional tori, n.d.

\newblock URL
  \url{https://www.math.stonybrook.edu/~kamenova/homepage_files/degen.pdf}.

\bibitem[Kato(1975)]{Kato1975Hopf}
Masahide Kato.

\newblock Topology of {Hopf} surfaces.

\newblock \emph{Journal of the Mathematical Society of Japan}, 27\penalty0
  (2):\penalty0 222--238, 1975.

\newblock \doi{10.2969/jmsj/02720222}.

\newblock URL
  \url{https://www.jstage.jst.go.jp/article/jmath1948/27/2/27_2_222/_pdf/-char/en}.

\bibitem[Kawai(1969)]{Kawai1969}
Soichi Kawai.

\newblock On compact complex analytic manifolds of complex dimension 3, {II}.

\newblock \emph{Journal of the Mathematical Society of Japan}, 21\penalty0
  (4):\penalty0 604--616, 1969.

\newblock \doi{10.2969/jmsj/02140604}.

\newblock URL
  \url{https://www.jstage.jst.go.jp/article/jmath1948/21/4/21_4_604/_pdf/-char/en}.

\bibitem[Khaled(2024)]{Khaled2024Hopf}
Zahraa Khaled.

\newblock Real structures on primary {Hopf} surfaces.

\newblock \emph{Complex Manifolds}, 11\penalty0 (1), 2024.

\newblock \doi{10.1515/coma-2024-0007}.

\newblock Cited version: arXiv:2310.05265v2, 4 December 2023.
\newblock URL \url{https://arxiv.org/abs/2310.05265v2}.

\bibitem[Khaled and Teleman(2026)]{KhaledTeleman2026}
Zahraa Khaled and Andrei Teleman.
\newblock On the classification of {Inoue} surfaces.
\newblock \emph{Journal of Mathematical Sciences, The University of Tokyo},
  33\penalty0 (3):\penalty0 283--357, 2026.
\newblock Journal edition:
  \url{https://www.ms.u-tokyo.ac.jp/journal/33-3-1.pdf}.
\newblock Cited preprint edition: arXiv:2406.15158v4, 13 July 2026,
  \url{https://arxiv.org/pdf/2406.15158v4}.


\bibitem[Kodaira(1954)]{Kodaira1954}
Kunihiko Kodaira.
\newblock On K\"ahler varieties of restricted type.
\newblock \emph{Proceedings of the National Academy of Sciences of the
United States of America}, 40\penalty0 (5):\penalty0 313--316, 1954.
\newblock \doi{10.1073/pnas.40.5.313}.


\bibitem[Kodaira(1963)]{Kodaira1963II}
Kunihiko Kodaira.

\newblock On compact analytic surfaces: {II}.

\newblock \emph{Annals of Mathematics}, 77\penalty0 (3):\penalty0 563--626,
  1963.

\newblock \doi{10.2307/1970131}.

\newblock URL \url{https://www.jstor.org/stable/1970131}.

\bibitem[Kulikov(1977)]{Kulikov1977}
Viktor~S. Kulikov.

\newblock Degenerations of {$K3$} surfaces and Enriques surfaces.

\newblock \emph{Mathematics of the USSR-Izvestiya}, 11\penalty0 (5):\penalty0
  957--989, 1977.

\newblock \doi{10.1070/IM1977v011n05ABEH001753}.

\newblock URL
  \url{https://www.mathnet.ru/php/getFT.phtml?jrnid=im&option_lang=eng&paperid=1877&what=fullteng}.

\bibitem[Kulikov(1981)]{Kulikov1981}
Viktor~S. Kulikov.

\newblock On modifications of degenerations of surfaces with {$\varkappa=0$}.

\newblock \emph{Mathematics of the USSR-Izvestiya}, 17\penalty0 (2):\penalty0
  339--342, 1981.

\newblock \doi{10.1070/IM1981v017n02ABEH001361}.

\newblock URL
  \url{https://www.mathnet.ru/php/getFT.phtml?jrnid=im&option_lang=eng&paperid=1953&what=fullteng}.

\bibitem[Kurdyka(1998)]{Kurdyka1998}
Krzysztof Kurdyka.

\newblock On gradients of functions definable in o-minimal structures.

\newblock \emph{Annales de l'Institut Fourier}, 48\penalty0 (3):\penalty0
  769--783, 1998.

\newblock \doi{10.5802/aif.1638}.

\newblock URL \url{https://www.numdam.org/article/AIF_1998__48_3_769_0.pdf}.

\bibitem[Lagrange(1868)]{Lagrange1770}
Joseph-Louis Lagrange.

\newblock D{\'e}monstration d'un th{\'e}or{\`e}me d'arithm{\'e}tique.

\newblock In \emph{{\OE}uvres de Lagrange}, volume III, pages 189--201.
  Gauthier-Villars, Paris, 1868.


\bibitem[Lazzeri(1973)]{Lazzeri1973Root}
Fulvio Lazzeri.

\newblock A theorem on the monodromy of isolated singularities.

\newblock In \emph{Singularit{\'e}s {\`a} Carg{\`e}se}, volume 7--8 of
  \emph{Ast{\'e}risque}, pages 269--275. Soci{\'e}t{\'e} Math{\'e}matique de
  France, 1973.

\newblock URL \url{https://numdam.org/item/AST_1973__7-8__269_0/}.

\bibitem[Lehn et~al.(2019)Lehn, Rollenske, and Schinko]{LRS2019}
Christian Lehn, S{\"o}nke Rollenske, and Caren Schinko.

\newblock The complex geometry of a hypothetical complex structure on {$S^6$}.

\newblock Revised version, arXiv:1912.09719v1, 20 December 2019.
\newblock Journal version: \emph{Differential Geometry and its Applications},
  57:\penalty0 121--137, 2018.
\newblock \doi{10.1016/j.difgeo.2017.10.015}.
\newblock Citations refer to the revised arXiv version.
\newblock URL \url{https://arxiv.org/abs/1912.09719v1}.


\bibitem[Liu and Rollenske(2012)]{LiuRollenske2012SLC}
Wenfei Liu and S{\"o}nke Rollenske.
\newblock Two-dimensional semi-log-canonical hypersurfaces.
\newblock \emph{Le Matematiche}, 67(2):185--202, 2012.
\newblock \doi{10.4418/2012.67.2.14}.

\bibitem[Liu and Rollenske(2026)]{LiuRollenske2026}
Wenfei Liu and S{\"o}nke Rollenske.

\newblock Yet another family of complex structures on {$S^6$}.

\newblock arXiv:2609.33826v2, 2026.

\newblock URL \url{https://arxiv.org/abs/2609.33826v2}.

\bibitem[{\L}ojasiewicz(1964)]{Lojasiewicz1964}
Stanis{\l}aw {\L}ojasiewicz.

\newblock Triangulation of semi-analytic sets.

\newblock \emph{Annali della Scuola Normale Superiore di Pisa, Classe di
  Scienze}, ser.~3, 18\penalty0 (4):\penalty0 449--474, 1964.

\newblock URL \url{https://www.numdam.org/item/ASNSP_1964_3_18_4_449_0.pdf}.

\bibitem[Meersseman(2019)]{Meersseman2019E}
Laurent Meersseman.

\newblock The {Teichm\"uller} and {Riemann} moduli stacks.

\newblock \emph{Journal de l'\'Ecole polytechnique---Math\'ematiques},
  6:\penalty0 879--945, 2019.

\newblock \doi{10.5802/jep.108}.

\newblock URL \url{https://www.numdam.org/item/10.5802/jep.108.pdf}.

\bibitem[Michel et~al.(2005)Michel, Pichon, and Weber]{MPW2005Hirzebruch}
Fran{\c c}oise Michel, Anne Pichon, and Claude Weber.

\newblock The boundary of the {Milnor} fiber of {Hirzebruch} surface
  singularities.
\newblock arXiv:math/0509451v1, 20 September 2005.
\newblock Published in \emph{Singularity Theory: Proceedings of the
2005 Marseille Singularity School and Conference}, pages 745--760,
World Scientific, 2007. \doi{10.1142/9789812707499_0032}.
\newblock The locators used here refer to the arXiv version.
\newblock URL \url{https://arxiv.org/pdf/math/0509451v1}.

\bibitem[Milnor(1968)]{Milnor1968}
John Milnor.

\newblock \emph{Singular Points of Complex Hypersurfaces}, volume~61 of
  \emph{Annals of Mathematics Studies}.

\newblock Princeton University Press and University of Tokyo Press, 1968.

\bibitem[Mumford(1961)]{Mumford1961}
David Mumford.

\newblock The topology of normal singularities of an algebraic surface and a
  criterion for simplicity.

\newblock \emph{Publications Mathématiques de l'IHÉS}, 9:\penalty0 5--22,
  1961.

\newblock \doi{10.1007/BF02698717}.

\newblock URL \url{https://www.numdam.org/item/PMIHES_1961__9__5_0.pdf}.

\bibitem[Mumford(1972)]{Mumford1972}
David Mumford.

\newblock An analytic construction of degenerating abelian varieties over
  complete rings.

\newblock \emph{Compositio Mathematica}, 24\penalty0 (3):\penalty0 239--272,
  1972.

\newblock URL
  \url{https://www.dam.brown.edu/people/mumford/alg_geom/papers/1972b--AnalConstrDegAV-Numdam.pdf}.

\bibitem[Nakamura(2008)]{Nakamura2008VII}
Iku Nakamura.

\newblock Survey on {VII}$_0$ surfaces, 2008.

\newblock URL
  \url{https://www.math.sci.hokudai.ac.jp/~nakamura/70surfaces080306.pdf}.

\bibitem[Nakano(1970)]{Nakano1970}
Shigeo Nakano.

\newblock On the inverse of monoidal transformation.

\newblock \emph{Publications of the Research Institute for Mathematical
  Sciences}, 6\penalty0 (3):\penalty0 483--502, 1970.

\newblock \doi{10.2977/PRIMS/1195193917}.

\newblock URL \url{https://ems.press/content/serial-article-files/41677}.

\bibitem[{National Institute of Standards and Technology}(2026)]{DLMFTheta2026}
{National Institute of Standards and Technology}.

\newblock {NIST Digital Library of Mathematical Functions}, chapters 20 and 23,
  2026.

\newblock URL \url{https://dlmf.nist.gov/}.

\bibitem[N{\'e}methi(1991)]{Nemethi1991Join}
Andr{\'a}s N{\'e}methi.

\newblock Generalized local and global {Sebastiani--Thom} type theorems.

\newblock \emph{Compositio Mathematica}, 80\penalty0 (1):\penalty0 1--14, 1991.

\newblock URL \url{https://www.numdam.org/article/CM_1991__80_1_1_0.pdf}.

\bibitem[Ornea et~al.(2020)Ornea, Verbitsky, and Vuletescu]{OVV2020}
Liviu Ornea, Misha Verbitsky, and Victor Vuletescu.
\newblock Classification of non-{K}{\"a}hler surfaces and locally conformally
  {K}{\"a}hler geometry.
\newblock arXiv:1810.05768v3, 26 June 2020.
\newblock Published as M.~S. Verbitsky, V. Vuletescu, and L. Ornea,
\emph{Russian Mathematical Surveys}, 76(2):261--289, 2021.
\newblock \doi{10.1070/RM9858}.
\newblock The locators used here refer to the arXiv version.
\newblock URL \url{https://arxiv.org/pdf/1810.05768v3}.


\bibitem[Papikian(2004)]{Papikian2004Tate}
Mihran Papikian.

\newblock Rigid analytic geometry and abelian varieties.

\newblock Author notes, 13 June 2004.

\newblock URL \url{https://math.stanford.edu/~vakil/snowbird/mihranjun21.pdf}.

\bibitem[Park and Suh(1996)]{ParkSuh1996}
Dae Heui Park and Dong Youp Suh.
\newblock Equivariant semi-algebraic triangulations of real algebraic
$G$-varieties.
\newblock \emph{Kyushu Journal of Mathematics}, 50(1):179--205, 1996.
\newblock \doi{10.2206/kyushujm.50.179}.
\newblock URL \url{https://www.jstage.jst.go.jp/article/kyushujm/50/1/50_1_179/_article}.

\bibitem[Peters(1990)]{Peters1990}
Chris A.~M. Peters.

\newblock Compact complex surfaces, 1990.

\newblock Author's 35-page notes for the Special Session of the
  Canadian Mathematical Society Summer Meeting, June 1990,
  \emph{Differential Geometry, Global Analysis, and Topology};
  citations use the pagination of this author PDF.

\newblock URL
  \url{https://cpeters1.win.tue.nl/ConfsAndSchools/Canada/surface.pdf}.

\bibitem[Peters(2021)]{Peters2021}
Chris Peters.

\newblock On complex surfaces with definite intersection form.

\newblock \emph{New York Journal of Mathematics}, 27:\penalty0 840--847, 2021.

\newblock URL \url{https://nyjm.albany.edu/j/2021/27-32p.pdf}.

\bibitem[Piatetski-Shapiro and
  Shafarevich(1971)]{PiatetskiShapiroShafarevich1971}
I.~I. Piatetski-Shapiro and I.~R. Shafarevich.

\newblock A {Torelli} theorem for algebraic surfaces of type {K3}.

\newblock \emph{Mathematics of the USSR-Izvestiya}, 5\penalty0 (3):\penalty0
  547--588, 1971.

\newblock \doi{10.1070/IM1971v005n03ABEH001075}.

\newblock URL \url{https://www.mathnet.ru/eng/im2021}.

\bibitem[Popescu-Pampu(2016)]{PopescuPampu2016}
Patrick Popescu-Pampu.

\newblock Complex singularities and contact topology.

\newblock \emph{Winter Braids Lecture Notes}, 3, Course~III:1--74, 2016.

\newblock \doi{10.5802/wbln.14}.

\newblock \url{https://www.numdam.org/item/10.5802/wbln.14.pdf}.

\bibitem[Ramis and Ruget(1970)]{RR70push}
Jean-Pierre Ramis and Gabriel Ruget.

\newblock Complexe dualisant et th{\'e}or{\`e}mes de dualit{\'e} en
  g{\'e}om{\'e}trie analytique complexe.

\newblock \emph{Publications Math{\'e}matiques de l'IH{\'E}S}, 38:\penalty0
  77--91, 1970.

\newblock \doi{10.1007/BF02684652}.

\newblock URL \url{https://www.numdam.org/item/PMIHES_1970__38__77_0.pdf}.

\bibitem[Reid(1994)]{Reid1994Nonnormal}
Miles Reid.

\newblock Nonnormal del {Pezzo} surfaces.

\newblock \emph{Publications of the Research Institute for Mathematical
  Sciences}, 30\penalty0 (5):\penalty0 695--727, 1994.

\newblock \doi{10.2977/prims/1195165581}.

\newblock URL \url{https://ems.press/content/serial-article-files/40537}.

\bibitem[Reid(1997)]{Reid1997Chapters}
Miles Reid.
\newblock Chapters on algebraic surfaces.
\newblock In \emph{Complex algebraic geometry (Park City, UT, 1993)},
IAS/Park City Mathematics Series, volume~3, pages 3--159.
American Mathematical Society, Providence, RI, 1997.
\newblock \doi{10.1090/pcms/003/02}.
\newblock Cited preprint edition: arXiv:alg-geom/9602006v1,
6 February 1996.
\newblock Page locators refer to this edition.
\newblock URL \url{https://arxiv.org/pdf/alg-geom/9602006v1}.

\bibitem[Remmert(1957)]{Remmert1957Mappings}
Reinhold Remmert.
\newblock Holomorphe und meromorphe Abbildungen komplexer R\"aume.
\newblock \emph{Mathematische Annalen}, 133:328--370, 1957.
\newblock \doi{10.1007/BF01342886}.
\newblock URL \url{https://gdz.sub.uni-goettingen.de/download/pdf/PPN235181684_0133/PPN235181684_0133.pdf}.

\bibitem[Rossi(2017)]{Rossi2017}
Michele Rossi.

\newblock A small and non-simple geometric transition.

\newblock \emph{Mathematical Physics, Analysis and Geometry}, 20\penalty0
  (2):\penalty0 15, 2017.

\newblock \doi{10.1007/s11040-017-9247-z}.

\newblock Cited preprint: arXiv:1211.6329v3, 31 March 2017.
\newblock URL \url{https://arxiv.org/pdf/1211.6329v3}.

\bibitem[Roulleau(2019)]{Roulleau2019GeneralizedKummer}
Xavier Roulleau.

\newblock On generalized Kummer surfaces and the orbifold Bogomolov--Miyaoka--Yau inequality.

\newblock \emph{Transactions of the American Mathematical Society}, 371(11):7651--7668, 2019.

\newblock \doi{10.1090/tran/7507}.

\bibitem[Schnell(n.d.)]{SchnellLecture3}
Christian Schnell.

\newblock Lecture 3: Cohomology of compact complex tori.

\newblock MAT615 lecture notes, Stony Brook University;\allowbreak{} lecture dated February
  4, year not specified in the document, n.d.

\newblock URL
  \url{https://www.math.stonybrook.edu/~cschnell/mat615/lectures/lecture3.pdf}.

\bibitem[Schr{\"o}er and Siebert(2002)]{SS2002}
Stefan Schr{\"o}er and Bernd Siebert.

\newblock Irreducible degenerations of primary {Kodaira} surfaces.

\newblock In \emph{Complex Geometry: Collection of Papers Dedicated to Hans
  Grauert}, pages 193--222. Springer, 2002.

\newblock \doi{10.1007/978-3-642-56202-0_14}.
\newblock Cited preprint edition: arXiv:math/0001077v2, 28 October 2001;
preprint page and statement locators refer to this edition.
\newblock URL \url{https://arxiv.org/pdf/math/0001077v2}.

\bibitem[Sch{\"u}tt and Shioda(2010)]{SchuttShioda2010}
Matthias Sch{\"u}tt and Tetsuji Shioda.

\newblock Elliptic surfaces.

\newblock In \emph{Algebraic Geometry in East Asia---Seoul 2008}, volume~60
  of \emph{Advanced Studies in Pure Mathematics}, pages 51--160.
  Mathematical Society of Japan, 2010.
\newblock \doi{10.2969/aspm/06010051}.
\newblock Section and page citations refer to arXiv:0907.0298v3,
  9 July 2010.
\newblock URL \url{https://arxiv.org/pdf/0907.0298v3}.


\bibitem[Serre(1956)]{Serre1956GAGA}
Jean-Pierre Serre.

\newblock G\'eom\'etrie alg\'ebrique et g\'eom\'etrie analytique.

\newblock \emph{Annales de l'Institut Fourier}, 6:\penalty0 1--42, 1956.

\newblock \doi{10.5802/aif.59}.

\bibitem[Siersma(1991)]{Siersma1991Variation}
Dirk Siersma.

\newblock Variation mappings on singularities with a 1-dimensional critical
  locus.

\newblock \emph{Topology}, 30\penalty0 (3):\penalty0 445--469, 1991.

\newblock \doi{10.1016/0040-9383(91)90025-Y}.


\bibitem[Siersma and Tib\u{a}r(2017)]{SiersmaTibar2017Deformation}
Dirk Siersma and Mihai Tib\u{a}r.
\newblock Milnor fibre homology via deformation.
\newblock In \emph{Singularities and Computer Algebra}, pages 305--322.
Springer, 2017.
\newblock \doi{10.1007/978-3-319-28829-1_14}.
\newblock The version cited is arXiv:1512.02840v1.
\newblock \url{https://webspace.science.uu.nl/~siers101/SiTi2015DefHomology.pdf}.

\bibitem[Simoni\v{c}(2016)]{Simonic2016Hartogs}
Aleksander Simoni\v{c}.

\newblock Elementary approach to the {Hartogs} extension theorem, 2016.

\newblock URL \url{https://arxiv.org/abs/1608.00950v1}.

\newblock arXiv:1608.00950v1, 1 August 2016;\allowbreak{} Theorem 1, p. 1.

\bibitem[Steenbrink(1976)]{Steenbrink1976}
Joseph Steenbrink.

\newblock Limits of Hodge structures.

\newblock \emph{Inventiones Mathematicae}, 31:\penalty0 229--257, 1976.

\newblock \doi{10.1007/BF01403146}.


\bibitem[Steenbrink(1989)]{Steenbrink1989Spectrum}
Joseph H.~M. Steenbrink.

\newblock The spectrum of hypersurface singularities.

\newblock In \emph{Th\'eorie de Hodge---Luminy, Juin 1987}, number 179--180 in
  Ast\'erisque, pages 163--184. Soci\'et\'e math\'ematique de France, 1989.

\newblock URL
  \url{https://www.numdam.org/article/AST_1989__179-180__163_0.pdf}.

\bibitem[Steenbrink(1999)]{Steenbrink1999}
J.~H.~M. Steenbrink.

\newblock Cup products and mixed {Hodge} structures.

\newblock arXiv:math/9902134v1, 23 February 1999.
\newblock URL \url{https://arxiv.org/pdf/math/9902134v1}.


\bibitem[Steinberg(1965)]{Steinberg1965}
Robert Steinberg.

\newblock Regular elements of semi-simple algebraic groups.

\newblock \emph{Publications Math\'ematiques de l'IH\'ES}, 25:\penalty0 49--80,
  1965.

\newblock \doi{10.1007/BF02684397}.

\newblock URL \url{https://www.numdam.org/item/PMIHES_1965__25__49_0.pdf}.

\bibitem[Teissier and Merle(1976--1977)]{MerleTeissier1977}
Bernard Teissier and Michel Merle.

\newblock Conditions d'adjonction, d'apr\`es {Du Val}.

\newblock \emph{S\'eminaire sur les singularit\'es des surfaces}, pages 1--16,
  1976--1977.

\newblock URL \url{https://www.numdam.org/item/SSS_1976-1977____A16_0.pdf}.

\bibitem[Teleman(2006)]{Teleman2006}
Andrei Teleman.

\newblock The pseudo-effective cone of a non-{K}\"ahlerian surface and
  applications.

\newblock \emph{Mathematische Annalen}, 335\penalty0 (4):\penalty0 965--989,
  2006.

\newblock \doi{10.1007/s00208-006-0782-3}.

\newblock Cited preprint edition: arXiv:0704.2948v3,
30 June 2007.
\newblock Page locators refer to this edition.
\newblock URL \url{https://arxiv.org/pdf/0704.2948v3}.

\bibitem[Teleman(2008)]{Teleman2008}
Andrei Teleman.

\newblock Gauge theoretical methods in the classification of
  non-{K}{\"a}hlerian surfaces, 2008.

\newblock URL
  \url{https://www.i2m.univ-amu.fr/perso/andrei.teleman/documents/postnikov.pdf}.

\bibitem[{The Stacks Project Authors}(2026)]{StacksProject}
{The Stacks Project Authors}.

\newblock The stacks project, 2026.

\newblock URL \url{https://stacks.math.columbia.edu}.

\bibitem[Ulmer and Urz{\'u}a(2021)]{UlmerUrzua2021}
Douglas Ulmer and Giancarlo Urz{\'u}a.

\newblock Bounding tangencies of sections on elliptic surfaces.

\newblock \emph{International Mathematics Research Notices}, 2021\penalty0
  (6):\penalty0 4768--4802, 2021.

\newblock \doi{10.1093/imrn/rnaa222}.

\newblock Cited preprint edition: arXiv:2002.01906v2, 26 May 2020;
the author date on its title page is 27 May 2020. Page locators refer
to this preprint.
\newblock URL \url{https://arxiv.org/pdf/2002.01906v2}.

\bibitem[van Straten(1987)]{vanStraten1987}
Duco van Straten.

\newblock \emph{Weakly Normal Surface Singularities and Their Improvements}.

\newblock PhD thesis, Rijksuniversiteit Leiden, 1987.

\newblock URL
  \url{https://download.uni-mainz.de/mathematik/Algebraische%20Geometrie/DvS-Publikationen/DvS_1987_weakly%20normal%20Surface%20Singularities%20and%20their%20improvements.pdf}.

\bibitem[Viaclovsky(2026{\natexlab{a}})]{Viaclovsky2026Two}
Jeff Viaclovsky.
\newblock A two-parameter family of complex structures on $S^6$.
\newblock arXiv:2609.33785v1, 2026{\natexlab{a}}.
\newblock \url{https://arxiv.org/abs/2609.33785v1}.

\bibitem[Viaclovsky(2026{\natexlab{b}})]{Viaclovsky2026Combs}
Jeff Viaclovsky.

\newblock Combs and complex structures on the six-sphere, 2026{\natexlab{b}}.

\newblock arXiv:2610.10451v1; manuscript dated 7 October 2026.

\newblock URL \url{https://arxiv.org/abs/2610.10451v1}.

\bibitem[Viaclovsky(2026{\natexlab{c}})]{Viaclovsky2026Log}
Jeff Viaclovsky.

\newblock Logarithmic transformations on the six-sphere, 2026{\natexlab{c}}.

\newblock arXiv:2610.39102v1, 29 October 2026; manuscript dated 28 October 2026.
\newblock \doi{10.48550/arXiv.2610.39102}.
\newblock URL \url{https://arxiv.org/abs/2610.39102v1}.
\newblock Cited edition: the version-one preprint; its Theorem~1.1 is on p.~1.


\bibitem[Wang(2026)]{Wang2026}
Zichang Wang.

\newblock A complex structure on a rational homology six-sphere with
  two-torsion.

\newblock arXiv:2609.11541v2, 2026.

\newblock URL \url{https://arxiv.org/abs/2609.11541v2}.
\newblock Cited edition: version~2, 8 November 2026; page and statement locators refer to this edition.

\bibitem[Yoshihara(1995)]{Yoshihara1995}
Hisao Yoshihara.

\newblock Quotients of abelian surfaces.

\newblock \emph{Publications of the Research Institute for Mathematical
  Sciences}, 31\penalty0 (1):\penalty0 135--143, 1995.

\newblock \doi{10.2977/PRIMS/1195164795}.

\newblock URL \url{https://ems.press/journals/prims/articles/3766}.

\end{thebibliography}

\end{document}
